% Version 555. Explicit joint harmonic tests on both sides of every triangular square.
% Additional source comparisons: New/Main.lean lines 121--167,
% New/Endgame/OneSidedCriterion.lean lines 41--99.
% Primary source for the target and the reduction: grh-compact/src/Statements.lean
% and grh-compact/src/Lib/General/NumberTheory/ZeroSymmetry.lean, commit
% 72fead102f663d5c7a407964a13dc0cf9d977207.
\documentclass[11pt,a4paper]{amsart}
\usepackage[T1]{fontenc}
\usepackage{lmodern,mathtools,amssymb,microtype}
\usepackage[margin=29mm]{geometry}
\usepackage[hidelinks]{hyperref}
\hypersetup{pdftitle={Zeros of Dirichlet L-functions},pdfsubject={English mathematical manuscript}}
\newtheorem{theorem}{Theorem}[section]
\newtheorem{lemma}[theorem]{Lemma}
\newtheorem{proposition}[theorem]{Proposition}
\newtheorem{corollary}[theorem]{Corollary}
\theoremstyle{definition}
\newtheorem{definition}[theorem]{Definition}
\newtheorem{conjecture}{Conjecture}
\DeclareMathOperator{\Rp}{Re}
\newcommand{\Z}{\mathbb Z}
\newcommand{\C}{\mathbb C}
\newcommand{\Q}{\mathbb Q}
\newcommand{\R}{\mathbb R}
\newcommand{\e}{\mathrm e}
\newcommand{\dd}{\,\mathrm d}
\title[Zeros of Dirichlet \(L\)-functions]{Zeros of Dirichlet \(L\)-functions\\and one-sided estimates}
\date{}
\begin{document}
\begin{abstract}
We construct the real determinant tests at every level of the triangular array. The ordinary arrow and the arrow between consecutive offsets have the same parameter insertion after reversing the order on the latter. This identifies all inserted weights and proves positivity of each source projection. Dirichlet GRH remains the target conjecture.
\end{abstract}
\maketitle

\begin{conjecture}[Dirichlet GRH]\label{conj:grh}
For every integer \(q\ge1\), every Dirichlet character \(\chi\bmod q\),
and every \(\rho\in\C\), if \(L(\rho,\chi)=0\) and \(0<\Rp\rho<1\),
then \(\Rp\rho=1/2\).
\end{conjecture}

\section{Induction and continuation}
\begin{definition}[Foundational conventions]\label{def:foundations}
We use classical logic, the positive integers, the complete ordered real
field, and \(\C=\R[i]\). We take \(\exp z=\sum_{k\ge0}z^k/k!\),
and use the real logarithm for powers of positive numbers. A holomorphic
function is complex differentiable on an open set. Meromorphic functions
are locally quotients of holomorphic functions with denominator not
identically zero. Integrals on real intervals use length measure.
An integral against a counting measure is a sum. For a positive function
\(g\), \(f=O(g)\) means \(|f|\le Cg\) throughout the stated range,
or eventually if the argument tends to infinity. Subscripts on \(O\)
indicate permitted dependence of \(C\).
\end{definition}

\begin{definition}
\label{def:character}A Dirichlet character modulo \(q\) is a homomorphism
\((\Z/q\Z)^\times\to\C^\times\), extended to a periodic function on
\(\Z\) by zero on nonunits. The principal character is one on the units.
A character is primitive if its restriction to the units does not
factor through the units of any proper positive divisor of \(q\).
For \(q=1\) the unique character has value one on every integer.
For \(\Rp s>1\) write \(L(s,\chi)=\sum_{n\ge1}\chi(n)n^{-s}\),
where \(u^s=\exp(s\log u)\) for \(u>0\).
\end{definition}

\begin{lemma}\label{lem:induce}
Every character \(\chi\bmod q\) is induced by a primitive character
\(\psi\bmod f\) for some \(f\mid q\). For \(\Rp s>1\),
\[
 L(s,\chi)=L(s,\psi)\prod_{\substack{p\mid q\\p\nmid f}}
                   (1-\psi(p)p^{-s}).
\]
Every factor in the finite product is nonzero for \(\Rp s>0\).
\end{lemma}
\begin{proof}
For \(f\mid q\), reduction on units is onto. Indeed, given \(a\) prime
to \(f\), let \(P\) be the product of the primes dividing \(q\) but not
\(f\). Since \((f,P)=1\), choose \(k\) with \(fk\equiv1-a\bmod P\).
Then \(a+fk\) is prime to \(q\) and lifts \(a\). Among divisors through
whose unit group \(\chi\) factors, choose the least. The induced
character there is primitive by minimality.

Values on the units are roots of unity, so the defining series converges
absolutely for \(\Rp s>1\). Expanding finite products of geometric
series and using unique factorisation sums exactly the terms with prime
factors in the chosen finite set. Absolute convergence permits its
increase to all primes. Thus \(L(s,\psi)=\prod_p(1-\psi(p)p^{-s})^{-1}\).
Lemma~\ref{lem:basic-arith} justifies the divisibility steps. Deleting the factors at newly excluded primes gives the displayed
identity. If \(\Rp s>0\), then \(|\psi(p)p^{-s}|\le p^{-\Rp s}<1\).
\end{proof}

\begin{proposition}\label{prop:cont}
The series \(L(s,\chi)\) continues meromorphically to \(\Rp s>0\).
Its only possible pole there is at \(s=1\), a simple pole of residue
\(\varphi(q)/q\) for the principal character. Lemma~\ref{lem:induce}'s
identity holds throughout this half-plane as a meromorphic identity.
\end{proposition}
\begin{proof}
Let \(A(x)=\sum_{1\le n\le x}\chi(n)\) and
\(\mu=q^{-1}\sum_{a=1}^q\chi(a)\). For a nonprincipal character,
multiplying the argument of the complete sum by a unit on which
\(\chi\ne1\) shows that \(\mu=0\). For the principal character
\(\mu=\varphi(q)/q\). Here \(\varphi(q)=|(\Z/q\Z)^\times|\). Complete periods cancel in \(B(x)=A(x)-\mu x\),
so \(|B(x)|\le 2q\). Integrating the series of step functions, for
\(\Rp s>1\) one obtains
\[
 L(s,\chi)=s\int_1^\infty A(x)x^{-s-1}\dd x
 =\frac{\mu s}{s-1}+s\int_1^\infty B(x)x^{-s-1}\dd x .
\]
In the last integral, on \(\Rp s\ge\delta>0\), the integrands before and after one differentiation in \(s\)
are dominated by a constant times
\((1+\log x)x^{-1-\delta}\), an integrable function. Lemmas~\ref{lem:integral-tools}
and~\ref{lem:analytic-tools} give
the continuation, and the explicit term gives the residue. The two sides of the earlier
identity now agree by analytic continuation.
\end{proof}

\section{Reflection in the critical strip}
\begin{lemma}\label{lem:gauss}
Let \(\psi\bmod f\) be primitive, \(e(u)=\exp(2\pi i u)\), and
\(\tau=\sum_{a\bmod f}\psi(a)e(a/f)\). For every integer \(k\),
\[
 \sum_{a\bmod f}\psi(a)e(ak/f)=\bar\psi(k)\tau,\qquad
 |\tau|=\sqrt f .
\]
\end{lemma}
\begin{proof}
For a unit \(k\), change variables \(a\mapsto ak\). Otherwise
\(h=f/(k,f)<f\). There is a unit \(u\equiv1\bmod h\) with
\(\psi(u)\ne1\). Indeed, triviality on the kernel of the surjection
of unit groups proved in Lemma~\ref{lem:induce} would give a
factorisation through \(h\). Multiplication by \(u\) fixes the additive
factor \(e(ak/f)\) and multiplies the sum by \(\psi(u)\), so the sum is zero.
Summing the squared modulus over \(k\bmod f\) gives
\(f\sum_{a\bmod f}|\psi(a)|^2=f\varphi(f)\), since
\(\sum_{k\bmod f} e(k(a-b)/f)\) equals \(f\) for \(a=b\bmod f\) and
zero otherwise by the finite geometric sum. The first identity gives
\(\varphi(f)|\tau|^2\) for the same sum. For \(f=1\) both identities
read \(1=1\).
\end{proof}

\begin{lemma}\label{lem:theta}
Let \(\psi\bmod f\) be primitive and let \(a\in\{0,1\}\) satisfy
\(\psi(-1)=(-1)^a\). For \(t>0\) put
\[
 \Theta_\psi(t)=\sum_{n\in\Z}\psi(n)n^a\e^{-\pi n^2t/f}
 \quad(0^0=1),\qquad
 \omega_\psi=(-i)^a\tau/\sqrt f .
\]
Then \(|\omega_\psi|=1\) and
\(\Theta_\psi(t)=\omega_\psi t^{-a-1/2}\Theta_{\bar\psi}(1/t)\).
\end{lemma}
\begin{proof}
Use the Fourier transform
\(\widehat g(y)=\int_\R g(x)e(-xy)\dd x\). For \(\alpha>0\) the transform
of \(\exp(-\pi\alpha x^2)\) is
\(\alpha^{-1/2}\exp(-\pi y^2/\alpha)\). At \(y=0\) this follows by
squaring the Gaussian integral and using polar coordinates. Differentiation
in \(y\) and integration by parts give the differential equation
\(\widehat g'(y)=-(2\pi/\alpha)y\widehat g(y)\).
Multiplying \(g\) by \(x\) differentiates its transform and multiplies
the result by \(-1/(2\pi i)\). Consequently, for
\(g(x)=x^a\e^{-\pi t x^2/f}\),
\[
 \widehat g(k/f)=(-i)^a(k/t)^a\sqrt{f/t}\,\e^{-\pi k^2/(ft)} .
\]

The periodisation \(\sum_{j\in\Z} g(u+jf)\) is continuous and its
\(k\)-th Fourier coefficient is \(f^{-1}\widehat g(k/f)\), by integrating
one period. These coefficients are absolutely summable. Their Fourier
series recovers the periodisation. Indeed, the continuous difference
has all coefficients zero, so all its convolutions with the Fej\'er
kernels vanish. In the variable \(u/f\), these kernels are
\(\frac1N|\sum_{j=0}^{N-1}e(jv)|^2\). Their integrals on a period are
one and their suprema off any neighbourhood of the integers tend to
zero by the geometric sum formula. Uniform continuity gives convergence
of the convolutions to the difference. Evaluating at each residue \(u\),
multiplying by \(\psi(u)\) and summing now proves the formula by
Lemma~\ref{lem:gauss}.
\end{proof}

\begin{lemma}\label{lem:gamma}
For \(\Rp z>0\) the integral
\(\Gamma(z)=\int_0^\infty \e^{-u}u^{z-1}\dd u\) is holomorphic and nonzero.
\end{lemma}
\begin{proof}
On compact subsets of this half-plane, integrable bounds at zero and
infinity persist after insertion of \(\log u\), giving holomorphy.
Integrating by parts \(n\) times gives
\(\int_0^1 t^{z-1}(1-t)^n\dd t=n!/[z(z+1)\cdots(z+n)]\).
After multiplication by \(n^z\) and substitution \(u=nt\), dominated
convergence applies by \((1-u/n)^n\le\e^{-u}\) on \([0,n]\).
The resulting limit is therefore \(\Gamma(z)\). The logarithm of the
absolute value of the finite expression equals
\[
 -\log|z|+(\Rp z)\log n-\sum_{k=1}^n\log|1+z/k|.
\]
It has a finite limit since
\(\log|1+z/k|=(\Rp z)/k+O_z(k^{-2})\) and
\(\sum_{k\le n}1/k-\log n\) converges by integral comparison.
The modulus of the limit is positive.
\end{proof}

\begin{proposition}\label{prop:reflect}
For every primitive character \(\psi\) and every \(s\) with \(0<\Rp s<1\),
\(L(s,\psi)=0\) if and only if \(L(1-s,\bar\psi)=0\).
\end{proposition}
\begin{proof}
Write \(\Lambda_\psi(s)=(f/\pi)^{(s+a)/2}
\Gamma((s+a)/2)L(s,\psi)\). If \(f>1\), then \(\psi(0)=0\).
Termwise integration in \(\Rp s>1\), using parity to pair \(n\) and
\(-n\), gives
\[
 2\Lambda_\psi(s)=\int_0^\infty\Theta_\psi(t)t^{(s+a)/2}\frac{\dd t}{t}.
\]
Absolute convergence for the interchange follows from the corresponding
sum of \(n^{-\Rp s}\) times the real gamma integral. At infinity the
theta function decreases exponentially. At zero Lemma~\ref{lem:theta}
bounds it by a power of \(1/t\) times an exponentially decreasing
function of \(1/t\). The integral thus continues to an entire function
of \(s\). Substituting \(t=1/u\) gives
\(\Lambda_\psi(s)=\omega_\psi\Lambda_{\bar\psi}(1-s)\) in the strip,
using Proposition~\ref{prop:cont} to identify the continued functions.

For \(f=1\), write \(\Theta(t)=\sum_{n\in\Z}\e^{-\pi n^2t}\) and
\(I(s)=\int_1^\infty(\Theta(t)-1)t^{s/2}\dd t/t\), which is entire.
Using \(\Theta(t)=t^{-1/2}\Theta(1/t)\) in the integral over \((0,1)\)
gives, initially in \(\Rp s>1\),
\[
 2\Lambda_1(s)=I(s)+I(1-s)+\frac2{s-1}-\frac2s .
\]
Here the two rational terms are the integrals of
\((t^{-1/2}-1)t^{s/2}\dd t/t\) over \((0,1)\). Continuation gives
the same identity in the strip, where the right side is invariant under
\(s\mapsto1-s\). All completion factors on both sides in both cases are
finite and nonzero there by Lemma~\ref{lem:gamma}, proving the assertion.
\end{proof}

\begin{theorem}\label{thm:reduction}
Dirichlet GRH is equivalent to the assertion that, for every primitive
Dirichlet character \(\psi\), \(L(s,\psi)\ne0\) whenever
\(1/2<\Rp s<1\).
\end{theorem}
\begin{proof}
The forward implication follows by restriction. Conversely, suppose the
right half of the open strip is zero-free for all primitive characters.
For a zero \(\rho\) in the open strip of any \(L(s,\chi)\), take \(\psi\)
as in Lemma~\ref{lem:induce}. The finite Euler factors have nonzero
product there, so \(L(\rho,\psi)=0\). We have \(\Rp\rho\le1/2\).
If the inequality were strict, Proposition~\ref{prop:reflect} would
give a zero at \(1-\rho\) for \(\bar\psi\) in the excluded half-strip.
The conjugate is primitive because any factorisation through a smaller
modulus would, on conjugation, give the same for \(\psi\). Hence \(\Rp\rho=1/2\).
\end{proof}

\section{Covolume and the missing dimensions}

\begin{lemma}\label{lem:hadamard}
Let \(v_1,\ldots,v_n\) lie in a finite-dimensional complex inner
product space. Write \(G=(\langle v_i,v_j\rangle)_{i,j}\). If
\(\|v_i\|^2\le\e^h\) for every \(i\), with \(h\in\R\), then
\(0\le\det G\le\e^{nh}\).
\end{lemma}
\begin{proof}
If the vectors are dependent, their Gram determinant vanishes. Otherwise
subtract from each vector its orthogonal projection onto the span of its
predecessors. The change of basis is triangular with diagonal entries one.
It preserves the Gram determinant and changes the Gram matrix to a
diagonal matrix with diagonal entries the squared lengths of the
orthogonal components. Each entry is at most the squared length of its
original vector. Their product gives both inequalities. The empty family
has determinant one.
\end{proof}

\begin{lemma}[Loss of rank]\label{lem:rank}
Let \(r>0\) and \(n,H,T,E,\delta,u,j\in\R\). Suppose
\[
 \begin{split}
 r(H-3E)&\le n(H-T+E)+2\delta,\\
 (r-n)(-H+T-E)&\le ru,\qquad 2\delta\le rj .
 \end{split}
\]
Then \(T\le4E+u+j\).
\end{lemma}
\begin{proof}
Writing \(h=H-T+E\) gives \(nh=rh+(r-n)(-h)\).
The three inequalities imply \(r(H-3E)\le r(H-T+E+u+j)\).
Cancel \(rH\) and divide by \(r>0\).
\end{proof}

\begin{proposition}[A logarithm in the covolume estimate]\label{prop:38}
Let \(x\ge2,\ b\ge x^{99},\ f\ge1,\ a\ge b,\ r>0\). Put
\(\ell=\log b,\ B=fa\ell^3,\ E=\ell B\). Suppose \(H,T,\delta,n\in\R\)
satisfy
\begin{align}
 (1-4/b^2)r\le n\le r,\qquad & |H|+|T|\le b^2x^{99},\label{eq:defect}\\
 \delta\le 5r\log(1+b^2)B,\qquad&
 r(H-3E)\le n(H-T+E)+2\delta .\label{eq:volume}
\end{align}
Then
\begin{equation}
 T\le38fa(\log b)^4. \label{eq:strong}
\end{equation}
\end{proposition}
\begin{proof}
Both \(B\) and \(E\) are positive. For \(h=H-T+E\),
\(-h\le |H|+|T|\). Thus
\[
 (r-n)(-h)\le (r-n)b^2x^{99}\le4rx^{99}.
\]
The first multiplication uses \(r-n\ge0\), which is essential when
\(h\) changes sign. By Lemma~\ref{lem:rank},
\[
 T\le4E+4x^{99}+10\log(1+b^2)B .
\]
Now \(\ell>1\) and \(E=fa\ell^4\ge b\ge x^{99}\).
Also \(1+b^2\le b^3\) for \(b\ge2\), whence
\(\log(1+b^2)\le3\ell\). The last display is at most
\((4+4+30)E\), proving \eqref{eq:strong}.
\end{proof}

\begin{corollary}\label{cor:scales}
In Proposition~\ref{prop:38}, let \(m\ge\sqrt x\), let \(U,S>0\), and
suppose \(f=\varphi(M)U\) and \(S\ge U/\ell\) for a positive integer
\(M\). If real numbers \(R,R',R_1'\) satisfy
\[
 T=ma\{S(R-R')+U R_1'\},
\]
then
\[
 R\le R'-(U/S)R_1'+38\varphi(M)(\log b)^5/\sqrt x .
\]
\end{corollary}
\begin{proof}
Divide \eqref{eq:strong} by \(maS>0\).
The quotient \(f/S\) is at most \(\varphi(M)\ell\).
Solving the displayed identity for \(R\) gives the assertion.
\end{proof}

\section{A linear window for two Laplace transforms}
\begin{lemma}\label{lem:moments}
Let \(F:[0,\infty)\to\R\) be measurable with \(|F(v)|\le C\e^v\).
Suppose the function
\[
 \Phi_F(s)=\int_0^\infty F(v)\e^{(1-s)v}\dd v,\qquad \Rp s>2 ,
\]
extends holomorphically to a neighbourhood of \(|s-4|\le3\).
Then there exist \(R>3\) and \(K<\infty\) such that for every
integer \(k\ge0\),
\[
 \left|\int_0^\infty F(v)v^k\e^{-3v}\dd v\right|
       \le K k! R^{-k}.
\]
For two such functions the same \(R,K\) can be used.
\end{lemma}
\begin{proof}
By Lemma~\ref{lem:analytic-tools} and compactness the continuation is holomorphic on a neighbourhood
of a larger closed disk with radius \(R>3\). On its boundary let
\(K\) bound the modulus. The \(k\)-th Taylor coefficient at \(4\)
has modulus at most \(K/R^k\), by integrating the function times
\((s-4)^{-k-1}/(2\pi i)\) around that circle. Differentiating the
defining integral at \(4\) gives \((-1)^k\) times the displayed
moment. The interchange is justified in \(|s-4|<1/2\) by the
integrable majorant \(C v^k\e^{-3v/2}\). Take a common smaller
radius and a larger boundary bound for two functions.
\end{proof}

% The source's window [t,t^2] is IsPolar.exists_window_sub_mul_le.
\begin{theorem}[A linear window]\label{thm:window}
Let \(F,G\) each satisfy Lemma~\ref{lem:moments}, including its
continuation hypothesis. There is \(B_0>1\), depending only on \(F,G\),
such that for all sufficiently large real \(t\) and all
\(0\le X\le\e^t\) there exists \(v\in[t,B_0t]\) with
\[
                        F(v)-XG(v)\le\e^{-t}.
\]
\end{theorem}
\begin{proof}
Fix common \(R,K\) as in Lemma~\ref{lem:moments} and choose \(C\ge1\)
so that \(|F(v)|+|G(v)|\le C\e^v\). Choose positive \(A,B_0\) satisfying
\[
 A\log(R/3)>4,\qquad A\log(4/3)>6,\qquad B_0>A\log3+5 .
\]
For \(t\ge1\) set \(k=\lceil At\rceil\).
Give \(Y\) on \([0,\infty)\) density \(3^{k+1}v^k\e^{-3v}/k!\).
Integration by parts gives \(\int_0^\infty v^k \e^{-cv}\dd v=k!/c^{k+1}\)
for \(c>0\), first for \(k=0\) and then by induction. The density
therefore has mass one and
\(\mathbb E[\e^{sY}]=(3/(3-s))^{k+1}\) for real \(s<3\).
Set \(D_X=F-XG\). The moment bound gives, uniformly in \(X\),
\[
 |\mathbb E D_X(Y)|\le3K(1+\e^t)(3/R)^k=O(\e^{-3t}).
\]

Write \(J=[t,B_0t]\). On \(Y<t\) we have \(\e^Y\le \e^{2t-Y}\).
On \(Y>B_0t\) we have \(\e^Y\le \e^{-B_0t+2Y}\).
The two exponential moments just computed imply
\[
 \mathbb E[\e^Y{\bf1}_{Y\notin J}]
 \le \e^{2t}(3/4)^{k+1}+\e^{-B_0t}3^{k+1}=O(\e^{-4t}).
\]
This also bounds \(\mathbb P(Y\notin J)\) since \(\e^Y\ge1\).
Furthermore
\(\mathbb E[|D_X(Y)|{\bf1}_{Y\notin J}]=O(\e^{-3t})\)
by the growth bound. All constants are independent of \(X,t\).
If \(D_X(v)>\e^{-t}\) throughout \(J\), then
\(\mathbb E D_X(Y)\ge\e^{-t}(1-O(\e^{-4t}))-O(\e^{-3t})\),
contrary to the first moment bound for all sufficiently large \(t\).
\end{proof}

\section{Characters and the boundary of the Euler product}

\begin{lemma}\label{lem:characters}
Let \(G\) be a finite abelian group and let \(\widehat G\) denote its
complex characters. Then \(|\widehat G|=|G|\) and, for \(u,v\in G\),
\[
 |G|^{-1}\sum_{\xi\in\widehat G}\xi(u)\overline{\xi(v)}
      ={\bf1}_{u=v}.
\]
For characters modulo \(M\), with \(v\) a fixed unit, this formula also
holds with \(u\) any residue and the equality interpreted modulo \(M\).
\end{lemma}
\begin{proof}
If \(H\le G\) and the coset of \(g\) has order \(r\) in \(G/H\),
extend a character \(\lambda\) on \(H\) to \(H\langle g\rangle\) by
\(\lambda(hg^j)=\lambda(h)z^j\), where \(z^r=\lambda(g^r)\).
Two representations differ by a multiple of \(r\) in the exponent,
so the formula is well defined. The \(r\) roots give precisely the
\(r\) possible extensions. Starting with the trivial subgroup gives
the count by multiplication of indices. If \(w\ne1\), the character
on \(\langle w\rangle\) sending \(w\) to \(\exp(2\pi i/\operatorname{ord}(w))\)
extends to \(G\). Multiplication of every summand in
\(\sum_\xi\xi(w)\) by this character's value at \(w\) permutes
the summands, forcing the sum to vanish. For \(w=1\) the count just
proved gives the assertion. Take \(w=uv^{-1}\). On a nonunit all
Dirichlet characters vanish by Definition~\ref{def:character}.
\end{proof}

\begin{lemma}[The real singularity of a positive Laplace transform]\label{lem:landau}
Let \(\nu\) be a positive measure on \([0,\infty)\) such that
\(\int \e^{-bv}\dd\nu(v)<\infty\) for some \(b\in\R\).
Put \(\alpha=\inf\{u\in\R:\int \e^{-uv}\dd\nu(v)<\infty\}\).
Then \(J(s)=\int \e^{-sv}\dd\nu(v)\) is holomorphic on
\(\Rp s>\alpha\). If \(\alpha\) is finite, \(J\) has no holomorphic
continuation to a neighbourhood of the real point \(\alpha\).
\end{lemma}
\begin{proof}
The set of convergence on the real line is upward closed.
On each compact subdisk of \(\Rp s>\alpha\), choose a real convergence
point strictly to its left. Its exponential dominates the integrand
and, up to a constant depending on \(k\), its \(k\)-th derivative
there. Lemmas~\ref{lem:integral-tools} and~\ref{lem:analytic-tools}
give holomorphy and the moment formula for derivatives.
If continuation existed on \(|s-\alpha|<4h\), \(h>0\), the Taylor
series about \(b_0=\alpha+h\) would converge at \(b_0-2h\).
There it is the series of nonnegative terms
\[
 \sum_{k\ge0}\frac{(2h)^k}{k!}\int v^k\e^{-b_0v}\dd\nu(v).
\]
By increasing convergence in Lemma~\ref{lem:integral-tools}, its sum is
\(\int \e^{-(\alpha-h)v}\dd\nu(v)\), contradicting the infimum.
\end{proof}

\begin{definition}\label{def:regular}
Write \(\zeta(s)=L(s,1\bmod 1)\). For a character \(\chi\bmod q\)
let \(L^\circ(s,\chi)=L(s,\chi)\) for nonprincipal \(\chi\), and
\(L^\circ(s,\chi)=(s-1)L(s,\chi)\) for principal \(\chi\), with
\(L^\circ(1,\chi)=\varphi(q)/q\) in the latter case. These functions
are holomorphic for \(\Rp s>0\) by Proposition~\ref{prop:cont}.
Put \(\ell_\chi=(L^\circ)'/L^\circ\). Write \(o_\chi(\rho)\) for
the vanishing order of \(L^\circ(\cdot,\chi)\) at \(\rho\), zero at a
nonzero value.
\end{definition}

\begin{proposition}\label{prop:edge}
For every Dirichlet character \(\chi\), \(L^\circ(s,\chi)\ne0\)
on \(\Rp s\ge1\).
\end{proposition}
\begin{proof}
For \(\sigma=\Rp s>1\) the Euler product of Lemma~\ref{lem:induce}
is the exponential of the series
\[
             \sum_{p\ {\rm prime}}\ \sum_{k\ge1}\chi(p)^kp^{-ks}/k .
\]
This series converges absolutely, since
\(\sum_{n\ge2}n^{-\sigma}/(1-2^{-\sigma})<\infty\).
The logarithm formula follows by integrating the geometric series near
zero, or by exponentiating its power series. Hence \(L\) has no zeros
there. For real \(t\) and \(\sigma>1\) the Euler product gives
\[
 \zeta(\sigma)^3\,|L(\sigma+it,\chi)|^4
             |L(\sigma+2it,\chi^2)|\ \ge\ 1 .
\]
Indeed each prime-power contribution to the logarithm, after removing
the positive factor \(p^{-k\sigma}/k\), is \(3\) if \(p\mid q\) and
otherwise \(3+4\cos\theta+\cos(2\theta)=2(1+\cos\theta)^2\), where
\(\exp(i\theta)=\chi(p)^k p^{-ikt}\).
A zero at \(1+it\) would contribute a factor \(O((\sigma-1)^4)\);
\(\zeta(\sigma)^3=O((\sigma-1)^{-3})\) by Proposition~\ref{prop:cont}.
The last factor is bounded unless \(t=0\) and \(\chi^2\) is principal.
In the bounded case the displayed inequality gives a contradiction.

The exceptional nonprincipal \(\chi\) takes only \(0,\pm1\) as
values. The product \(\zeta(s)L(s,\chi)\) in \(\Rp s>1\) has
coefficients \(a_n=\sum_{h\mid n}\chi(h)\). By prime factorization,
\(a_n\) is the product over \(p^j\Vert n\) of
\(1+\chi(p)+\cdots+\chi(p)^j\). Each factor is nonnegative and is
at least one for even \(j\). Hence \(a_{m^2}\ge1\). The convergence
abscissa of \(\sum a_n n^{-s}\) is at least \(1/2\), by the divergent
harmonic sum, and at most one. A zero of \(L(1,\chi)\) would make
\(\zeta L(\cdot,\chi)\) holomorphic in \(\Rp s>0\), including at
this abscissa. Uniqueness in Lemma~\ref{lem:analytic-tools} and
Lemma~\ref{lem:landau}, for the measure with masses \(a_n\) at \(\log n\),
give a contradiction. For a principal character its regularized value
at one is positive by Definition~\ref{def:regular}.
\end{proof}

\begin{corollary}\label{cor:ap}
For every \(M\ge1\) and every unit \(c\bmod M\),
\(\sum_{p\equiv c\bmod M}1/p=\infty\).
\end{corollary}
\begin{proof}
For \(\sigma>1\) use Lemma~\ref{lem:characters} to express the
prime sum in the form
\[
 \sum_{p\equiv c\bmod M}p^{-\sigma}
  =\frac1{\varphi(M)}\sum_{\xi\bmod M}\overline{\xi(c)}
                                  \sum_p \xi(p)p^{-\sigma}.
\]
In each inner sum replacing it by the Euler logarithm in
Proposition~\ref{prop:edge} changes it by \(O(1)\) uniformly as
\(\sigma\downarrow1\), since
\(\sum_p\sum_{k\ge2}p^{-k\sigma}\le 2\sum_{n\ge2}n^{-2}\), with \(k\ge2\)
and \(p\) prime. For a nonprincipal character this logarithm stays
bounded. On a disk about one on which \(L/L(1)\) is close to one,
the convergent series for \(\log(1+z)\) gives a bounded branch after
adding a logarithm of \(L(1)\). The Euler logarithm differs on the real
interval by a constant multiple of \(2\pi i\), by continuity. For the
principal character the real logarithm is
\(\log(1/(\sigma-1))+O_M(1)\) by Proposition~\ref{prop:cont}.
The prime sum therefore tends to infinity, while its summands are at
most \(1/p\).
\end{proof}

\section{A transform of the weighted prime sum}

\begin{definition}\label{def:prime-error}
Let \(M\ge1\) and let \(d\) be an odd prime with \(d\nmid M\).
Write \(\eta_d(n)=(n/d)\), equal to zero on multiples of \(d\), one
on nonzero squares modulo \(d\), and minus one on the other residues.
Squaring on the units has kernel \(\{1,-1\}\), since a field has no
zero divisors. Its image has index two, so \(\eta_d\) is a nonprincipal
character. Let \(\Lambda(p^j)=\log p\) for prime \(p\) and \(j\ge1\),
and set \(\Lambda(n)=0\) otherwise. Define
\[
 \lambda_{d,M}(n)=\Lambda(n)(1-\eta_d(n)){\bf1}_{n\equiv1\bmod M},
 \qquad H_{d,M}(s)=
 \frac{\sum_{\xi\bmod M}(\ell_{\xi\eta_d}(s)-\ell_\xi(s))-1}{\varphi(M)}.
\]
Here \(\xi\eta_d\) is the pointwise product modulo \(Md\). It is
nonprincipal by simultaneous congruences from Lemma~\ref{lem:basic-arith}:
choose a residue equal to \(1\bmod M\) and a nonsquare unit modulo \(d\).
Proposition~\ref{prop:edge} makes \(H_{d,M}\) holomorphic near \(1\).
Write \(\kappa_{d,M}=H_{d,M}(1)\). This is real since character
conjugation permutes the summands and commutes with their evaluation
and differentiation at real arguments, first for \(\Rp s>1\) and
then by continuation. For \(x\ge1\) set
\[
 R_{d,M}(x)=\sum_{n\le x}\lambda_{d,M}(n)(1/n-1/x)
                         -\frac{\log x}{\varphi(M)}-\kappa_{d,M}.
\]
All sums over \(n\) in this section start at \(1\).
\end{definition}

\begin{proposition}\label{prop:polar}
For \(v\ge0\), \(F_{d,M}(v)=R_{d,M}(\e^v)\) is real and continuous and
\(|F_{d,M}(v)|=O_{d,M}((1+v)^2)\). For \(\Rp s>1\),
\[
 \Phi_{d,M}(s)=\int_0^\infty F_{d,M}(v)\e^{(1-s)v}\dd v
              =\frac{H_{d,M}(s)/s-\kappa_{d,M}}{s-1}.
\]
The right side extends holomorphically near every point of \(\Rp s\ge1\)
and meromorphically throughout \(\Rp s>0\), with at most simple poles.
For \(0<\Rp\rho<1\) the residue, interpreted as zero at a regular point, is
\begin{equation}\label{eq:res}
 \frac{\sum_{\xi\bmod M}\big(o_\xi(\rho)-o_{\xi\eta_d}(\rho)\big)}
      {\varphi(M)\rho(1-\rho)} .
\end{equation}
\end{proposition}
\begin{proof}
Each summand enters with weight zero at \(x=n\). The nonnegative
finite sum is at most \(2\log x\sum_{n\le x}1/n\le2\log x(1+\log x)\),
giving continuity, growth and absolute convergence of the integral.
For \(\Rp s>1\), termwise differentiation of the Euler logarithm
from Proposition~\ref{prop:edge} gives
\(-L'(s,\chi)/L(s,\chi)=\sum_n\chi(n)\Lambda(n)n^{-s}\).
Uniform absolute convergence of the derivatives on compact subdisks
follows by comparison with \(\sum_n(\log n)n^{-1-\delta}\) for some
\(\delta>0\). Apply Lemma~\ref{lem:characters} at the identity and
subtract the corresponding twisted sums. Only the untwisted principal
character acquires a term \(1/(s-1)\) when passing from \(L'/L\) to
\(\ell\). It follows that
\[
 \sum_{n\ge1}\lambda_{d,M}(n)n^{-s}
     =H_{d,M}(s)+\frac{s}{\varphi(M)(s-1)} .
\]
With \(x=\e^v\) the transform uses \(x^{-s}\dd x\).
Integrating the summand \(1/n-1/x\) over \(x\ge n\) against this
measure gives \(n^{-s}/[s(s-1)]\). Interchange is justified by the
same formula with \(\Rp s\) and the nonnegative weights, using
Lemma~\ref{lem:integral-tools}. The transforms of \(\log x\) and
\(1\) are \((s-1)^{-2}\) and \((s-1)^{-1}\), by elementary integration.
Substitution gives the formula. Its numerator vanishes at \(1\)
by Definition~\ref{def:prime-error}. All logarithmic derivatives
are regular on the closed right half-plane of Proposition~\ref{prop:edge}.
At an interior zero use Lemma~\ref{lem:analytic-tools} to take their
residues. The minus sign from \(1/(\rho-1)\) reverses the order
of subtraction and gives \eqref{eq:res}.
\end{proof}

\begin{corollary}\label{cor:usable-window}
For each \(d,M\) of Definition~\ref{def:prime-error} there is \(B_0>1\)
such that for sufficiently large \(t\) and every \(0\le X\le\e^t\)
there is \(v\in[t,B_0t]\) satisfying
\(R_{d,M}(\e^v)-X R_{d,1}(\e^v)\le\e^{-t}\).
\end{corollary}
\begin{proof}
The disk \(|s-4|\le3\) lies in \(\Rp s\ge1\). Proposition~\ref{prop:polar}
supplies a holomorphic neighbourhood of that compact disk for both
transforms and its polynomial bound implies the exponential bound in
Lemma~\ref{lem:moments}. Apply Theorem~\ref{thm:window}.
\end{proof}

\begin{theorem}[A one-sided criterion at a prescribed height]\label{thm:signed}
Let \(d,M\) be as in Definition~\ref{def:prime-error}. Fix
\(1/2<\beta<1\) and \(Y\ge0\). Suppose
\(R_{d,M}(x)\le x^{\beta-1}\) for all sufficiently large \(x\), and
\(L^\circ(s,\xi\eta_d)\ne0\) for every \(\xi\bmod M\) when
\(\Rp s\ge\beta\) and \(|\operatorname{Im}s|\le Y\).
Then \(L^\circ(s,\xi)\ne0\) for every \(\xi\bmod M\) when
\(\Rp s>\beta\) and \(|\operatorname{Im}s|\le Y\).
\end{theorem}
\begin{proof}
By the assumed upper bound and continuity in Proposition~\ref{prop:polar}
there is \(C\ge1\) with \(\e^v F_{d,M}(v)\le C\e^{\beta v}\)
for all \(v\ge0\). Give \([0,\infty)\) the positive measure of density
\(C\e^{\beta v}-\e^vF_{d,M}(v)\). In \(\Rp s>1\) its transform equals
\(J(s)=C/(s-\beta)-\Phi_{d,M}(s)\). Its abscissa \(\alpha\) is at
most one by the growth estimate. If \(\alpha>\beta\), the real point
\(\alpha\) is a singularity by Lemma~\ref{lem:landau}. The meromorphic
continuation is regular at \(1\) and has at most simple poles, so
\(\beta<\alpha<1\) and this must be a simple pole. Its residue is
strictly positive: along real \(\sigma\downarrow\alpha\),
\((\sigma-\alpha)J(\sigma)\ge0\), and the residue of a simple pole
is nonzero. On the other hand all twisted orders at \(\alpha\)
vanish. Formula~\eqref{eq:res} gives \(J\) a nonpositive residue
there since \(\alpha(1-\alpha)>0\). This contradiction shows
\(\alpha\le\beta\), also allowing \(\alpha=-\infty\).

The positive measure's integral is consequently holomorphic throughout
\(\Rp s>\beta\). By uniqueness in Lemma~\ref{lem:analytic-tools},
\(\Phi_{d,M}\) is holomorphic there as well. At any point \(\rho\)
of the asserted region with \(\Rp\rho<1\), the residue in
\eqref{eq:res} is zero and all twisted orders vanish. Its numerator
is therefore a sum of nonnegative integers equal to zero. Each
untwisted order is zero. Proposition~\ref{prop:edge} handles
\(\Rp\rho\ge1\).
\end{proof}

\section{A simultaneous choice of quadratic twist}

\begin{lemma}\label{lem:reciprocity}
For \(n\ge1\) let \(m\) be its odd part, obtained by removing its
power of two. For any integer \(v\) set
\(J_v(n)=\prod_{p^j\Vert m}(-v/p)^j\), the empty product being \(1\).
Here \((\cdot/p)\) is the symbol of Definition~\ref{def:prime-error}.
Then \(J_v\) is completely multiplicative in \(n\). If \(d\) is prime
and \(d\equiv7\bmod 8\), then \(J_d(n)=(n/d)\) for all \(n\).
\end{lemma}
\begin{proof}
Multiplicativity follows by adding exponents. We give the reciprocity
calculation. For an odd prime \(p\), put \(h=(p-1)/2\). Multiplication
by a unit \(b\) permutes all the units, so multiplying their values
gives \(b^{p-1}=1\) in the residue field. The \(h\) nonzero squares
are roots of \(X^h-1\) and exhaust its roots, since division by a
linear factor lowers degree by one. Consequently \(b^h=(b/p)\)
in that field. The absolute values of the least signed residues of
\(b,2b,\ldots,hb\) permute \(1,\ldots,h\), since \(bi\equiv\pm bj\)
forces \(i=j\) in this range. Multiplying shows \((b/p)=(-1)^N\),
where \(N\) counts negative signed residues. In particular
\((-1/p)=(-1)^h\). For \(b=2\), \(N=h-\lfloor p/4\rfloor\), even
when \(p\equiv7\bmod8\).

If \(b\) is positive and odd and prime to \(p\), let \(r_i\in[1,p-1]\)
be the positive residue of \(bi\). The sum of the \(r_i\) is congruent
modulo two to \(\sum_{i=1}^h i+N\): each negative signed representative
replaces its absolute value \(j\) by \(p-j\). Summing
\(bi=p\lfloor bi/p\rfloor+r_i\) gives
\(N\equiv\sum_{i=1}^h\lfloor bi/p\rfloor\bmod2\).
For distinct odd primes \(b,p\), this last floor sum and its counterpart
with \(b,p\) interchanged count complementary points in
\(1\le i\le(p-1)/2,\ 1\le j\le(b-1)/2\), according as \(pj<bi\)
or \(pj>bi\). Equality is impossible. Hence
\((b/p)(p/b)=(-1)^{(p-1)(b-1)/4}\).
For \(b=d\equiv7\bmod8\) this gives \((p/d)=(-d/p)\) for odd
\(p\ne d\), and both sides vanish at \(p=d\). The factor at \(2\)
is one by the calculation above. Multiplication proves the formula.
\end{proof}

\begin{lemma}\label{lem:symbol-sum}
There is an absolute \(C\) such that, if the odd part of \(u\ge1\)
is not a square and \(X\ge1\), then
\[
 \left|\sum_{1\le v\le X}J_v(u)\right|
       \le C\sqrt u(1+\log u).
\]
\end{lemma}
\begin{proof}
First let \(\lambda\bmod f\) be primitive and nonprincipal.
By Lemma~\ref{lem:gauss} and Fourier inversion on the residue classes,
\(\lambda(n)\) is \(1/f\) times a Fourier sum whose nonconstant
coefficients have modulus at most \(\sqrt f\), and whose constant
coefficient vanishes. Inversion follows by summing the finite geometric
sum used in that lemma. The modulus of a sum of
\(\exp(2\pi ikn/f)\) over an initial interval is at most
\(1/|\sin(\pi k/f)|\le f/[2\min(k,f-k)]\).
Here \(\sin t\ge2t/\pi\) on \([0,\pi/2]\) follows from concavity.
Summing over \(1\le k<f\) gives \(O(\sqrt f(1+\log f))\),
uniformly in the interval length.

Write \(m\) for the odd part of \(u\).
Up to the constant sign \(\prod_{p^j\Vert m}(-1/p)^j\), the symbol
is a character in \(v\bmod m\). Choose a prime with odd exponent
in \(m\); a nonsquare there and ones at the other primes give a
unit of character value \(-1\) by Lemma~\ref{lem:basic-arith}.
Its primitive inducing character \(\lambda\bmod f\) from
Lemma~\ref{lem:induce} is therefore nonprincipal. Let \(P\) be the
product of the primes dividing \(m\) but not \(f\). Expanding
\({\bf1}_{(v,P)=1}=\prod_{p\mid P}(1-{\bf1}_{p\mid v})\)
expresses the interval sum as at most \(2^{\#\{p:p\mid P\}}\) sums
of \(\lambda(k)\) over \(k\le X/h\), with coefficients of modulus at
most one and \(h\mid P\). For each prime at least \(5\), \(2\le\sqrt p\),
and for the primes \(2,3\) their combined excess factor is less than
two. Thus \(2^{\#\{p:p\mid P\}}\le2\sqrt P\). Since \(fP\mid m\),
the preceding bound proves the claim.
\end{proof}

\begin{definition}\label{def:theta-tail}
For \(j\in\{0,1\}\), \(t>0\), \(T\ge1\) and any complex sequence \(c\)
of modulus at most one, set
\[
 g_j(t)=t^j\e^{-\pi t^2},\quad
 \theta_j(c,t)=\sum_{n\ge1}c(n)g_j(nt),\quad
 I_{j,T}(c)=\int_{1/T}^\infty|\theta_j(c,t)|\,\frac{\dd t}{\sqrt t}.
\]
Write \(G_j(s)=\int_0^\infty g_j(t)t^{s-1}\dd t\) for \(\Rp s>0\).
Define \(\mu(1)=1\), \(\mu(n)=0\) if a prime square divides \(n\),
and otherwise \(\mu(n)=(-1)^k\) for \(n\) a product of \(k\)
distinct primes.
\end{definition}

\begin{lemma}\label{lem:theta-average}
For fixed \(j\), every \(\epsilon>0\), \(T\ge1,\ 1\le X\le T^2\),
and \(|c(n)|\le1\),
\[
 \sum_{v\le X}|\theta_j(cJ_v,t)|^2
   \ll_\epsilon T^2 t^{-1-2\epsilon}\quad(t\ge 1/T),
 \qquad
 \sum_{v\le X} I_{j,T}(cJ_v)\ll_\epsilon T^{2+\epsilon}.
\]
Here and below \(\ll\) has the meaning of \(O\) in
Definition~\ref{def:foundations}, and \(v\) ranges over positive integers.
\end{lemma}
\begin{proof}
For every \(b>0\), \(g_j(t)\ll_b t^{-b}\) because \(t^{b+j}\e^{-\pi
t^2}\) is bounded on \(t>0\). Expand the squared sum in indices \(n,m\).
Absolute convergence follows from the Gaussian decay at each fixed \(t\).
The sum in \(v\) has kernel \(J_v(nm)\) by Lemma~\ref{lem:reciprocity}.
If the odd part of \(nm\) is a square, either \(nm\) or \(2nm\)
is a square. For each such choice \(k=1,2\), the pair has the form
\(kn=a b^2,\ m=a h^2\) for some squarefree \(a\) and positive \(b,h\),
by collecting the primes occurring to odd exponent in \(kn\);
these are exactly the primes occurring to odd exponent in \(m\). For \(\sigma=1/2+\epsilon\) the sum of
\((nm)^{-\sigma}\) over these pairs is at most
\[
 \sum_{k=1,2}k^\sigma\sum_{a,b,h\ge1}(abh)^{-2\sigma}<\infty .
\]
Use \(|\sum_{v\le X}J_v(nm)|\le X\) and
\(g_j(nt)g_j(mt)\ll_\epsilon t^{-2\sigma}(nm)^{-\sigma}\).
These pairs contribute \(O_\epsilon(T^2t^{-1-2\epsilon})\).

For the other pairs Lemma~\ref{lem:symbol-sum} bounds the kernel sum
by \(O_\epsilon((nm)^{1/2+\epsilon/2})\), since
\((1+\log u)u^{-\epsilon/2}\) is bounded on \(u\ge1\).
Use the Gaussian bound with exponent \(3/2+\epsilon\). The double sum
is then at most a constant times
\(t^{-3-2\epsilon}(\sum_n n^{-1-\epsilon/2})^2\), which is
\(O_\epsilon(T^2t^{-1-2\epsilon})\) as \(t\ge1/T\). All numerical
series here converge by comparison with their power integrals.
Finally \(\sum_v|z_v|\le\sqrt{\lfloor X\rfloor\sum_v|z_v|^2}\)
by expanding the sum of squares of differences of the \(|z_v|\).
Integrate the resulting \(O_\epsilon(T^2t^{-1/2-\epsilon})\)
against \(\dd t/\sqrt t\) to obtain the second assertion.
\end{proof}

\begin{lemma}\label{lem:interpolate}
If \(a<b\), \(D\) is bounded and continuous on \(a\le\Rp s\le b\)
and holomorphic inside, and \(|D|\le A,B\) on the left and right
boundary respectively, with \(A,B>0\), then at \(\Rp s=\sigma\)
\[
 |D(s)|\le A^{(b-\sigma)/(b-a)}B^{(\sigma-a)/(b-a)} .
\]
\end{lemma}
\begin{proof}
Divide \(D(s)\) by \(A^{(b-s)/(b-a)}B^{(s-a)/(b-a)}\) to obtain a
bounded function \(D_0\) of boundary modulus at most one.
For any \(\delta>0\) multiply by \(\exp(\delta(s-a)^2)\). On
the two sides the modulus is at most \(\exp(\delta(b-a)^2)\).
On horizontal segments of heights \(\pm H\) it is at most
\(\sup|D_0|\exp(\delta((b-a)^2-H^2))\). The maximum on a closed
rectangle is on its boundary: an interior maximum of a nonconstant
holomorphic function contradicts the circle mean formula of
Lemma~\ref{lem:analytic-tools}, since equality in the bound for the
mean forces the function to be constant on a small circle and then
inside it by the disk formula. The zero maximum case is immediate.
Take \(H\) arbitrarily large and then \(\delta\downarrow0\).
\end{proof}

\begin{proposition}\label{prop:twist-test}
Fix \(0<\epsilon<1/18,\ Y\ge0\). For all sufficiently large \(T\),
depending only on \(\epsilon,Y\), the following holds. Let \(\chi\)
be a nonprincipal primitive character modulo \(Q\le T^2\), with
\(\chi(-1)=(-1)^j\). If
\[
                   I_{j,T}(\chi)+I_{j,T}(\mu\chi)\le T^{2\epsilon},
\]
then \(L(s,\chi)\ne0\) whenever
\(\Rp s\ge1/2+9\epsilon,\ |\operatorname{Im}s|\le Y\).
\end{proposition}
\begin{proof}
The change of variable \(u=\pi t^2\) gives
\(G_j(s)=\tfrac12\pi^{-(s+j)/2}\Gamma((s+j)/2)\). By
Lemma~\ref{lem:gamma} this is holomorphic and nonzero for \(\Rp s>0\).
The integral bounds it uniformly on \(1/2\le\Rp s\le2\) by using
\(t^{-1/2}\) for \(t\le1\) and \(t\) for \(t\ge1\) in the integral.
Moreover \(g_j(z)\ll \e^{-2z}\) for \(z>0\), whence geometric summation
gives \(|\theta_j(c,t)|\ll1/(\e^{2t}-1)\le \e^{-t}/(2t)\)
up to an absolute constant. The last inequality follows from
\(\e^t-\e^{-t}\ge2t\).
For \(\Rp s>1\), termwise integration of the theta series gives
\(G_j(s)\sum c(n)n^{-s}\), justified by absolute summability and
Lemma~\ref{lem:integral-tools}. For \(c=\chi\), Lemma~\ref{lem:theta}
at \(Qt^2\) gives exactly
\[
 \theta_j(\chi,t)=\frac{\omega_\chi}{\sqrt Q\,t}
                  \theta_j(\bar\chi,1/(Qt)).
\]
Indeed \(2\theta_j(\chi,t)=t^j\Theta_\chi(Qt^2)\). This yields
exponential decay at zero, up to powers of \(t\), for each fixed
\(Q\). Consequently the full Mellin integral equals
\(B(s)=G_j(s)L(s,\chi)\) throughout \(\Rp s>0\) by
Proposition~\ref{prop:cont} and Lemma~\ref{lem:analytic-tools}.
Both it and
\(C_T(s)=\int_{1/T}^\infty\theta_j(\mu\chi,t)t^{s-1}\dd t\) are
bounded on \(1/2\le\Rp s\le2\) for each fixed \(T,\chi\), by the
same absolute integrals at its endpoints.

Consider \(D(s)=G_j(s)^2-B(s)C_T(s)\). On the line \(\Rp s=1/2\),
\(|C_T(s)|\le I_{j,T}(\mu\chi)\). The functional equation of the
theta series and the substitution \(u=1/(Qt)\) give
\[
 \int_0^{1/T}|\theta_j(\chi,t)|\,\frac{\dd t}{\sqrt t}
    =\int_{T/Q}^\infty|\theta_j(\bar\chi,u)|\,\frac{\dd u}{\sqrt u}
    \le I_{j,T}(\chi).
\]
Here \(T/Q\ge1/T\) and the kernel is real. Thus \(|B|\le2I_{j,T}(\chi)\)
and \(|D|\ll T^{4\epsilon}\).
On \(\Rp s=2\) the full Mellin integral of \(\mu\chi\) equals
\(G_j(s)/L(s,\chi)\): the absolutely convergent product of its
Dirichlet series with \(L\) has coefficient
\(\chi(n)\sum_{h\mid n}\mu(h)={\bf1}_{n=1}\), by expanding
\(\prod_{p\mid n}(1-1)\). The part omitted in \(C_T\) has modulus
\(O(1/T)\) by the preceding bound \(|\theta|\ll t^{-1}\e^{-t}\).
Also \(|B|\ll1\) there by \(\sum_n n^{-2}<\infty\).
Thus \(|D|\ll T^{-1}\) on this line. The constants on both lines
are absolute, with just two choices of \(j\).
By Lemma~\ref{lem:interpolate},
\[
 |D(s)|\ll T^{4\epsilon(1-u)-u},\qquad
 u=(\Rp s-1/2)/(3/2).
\]
When \(1/2+9\epsilon\le\Rp s\le1\) this is \(O(T^{-2\epsilon})\).
At a zero of \(L\) we would have \(D=G_j^2\). On the compact
rectangle \(1/2\le\Rp s\le1,\ |\operatorname{Im}s|\le Y\),
the minimum of \(|G_j|^2\) over both parities is positive by
Lemma~\ref{lem:gamma}. This contradicts the last bound for sufficiently
large \(T\). Proposition~\ref{prop:edge} handles \(\Rp s\ge1\).
\end{proof}

\begin{theorem}\label{thm:twist}
Let \(M\ge1\), and let \(c\) be a unit modulo \(M\) such that every
prime \(d\equiv c\bmod M\) satisfies \(d\equiv7\bmod8\).
For every \(1/2<\beta<1,\ Y\ge0\) there are infinitely many such
primes for which, simultaneously for every \(\xi\bmod M\),
\[
 L^\circ(s,\xi\eta_d)\ne0
       \quad(\Rp s\ge\beta,\ |\operatorname{Im}s|\le Y).
\]
\end{theorem}
\begin{proof}
Choose \(0<\epsilon<1/18\) with \(9\epsilon\le\beta-1/2\).
For a fixed \(\xi\) take its inducing primitive \(\psi\bmod f\)
from Lemma~\ref{lem:induce}. For a prime of the specified class
we have \(d\nmid M\). The character \(\psi\eta_d\) modulo \(fd\)
is primitive. Indeed a divisor of \(fd\) has the form \(h d^e\)
with \(h\mid f\) and \(e=0,1\). If the product factors through it,
varying the component modulo \(d\) with the component modulo \(f\)
equal to \(1\) shows \(e=1\). Similarly, take the component modulo \(d\)
equal to \(1\) and vary the component modulo \(f\) over the kernel
of reduction modulo \(h\). Primitivity and surjectivity in
Lemma~\ref{lem:induce} give \(h=f\). These independent choices
follow from Lemma~\ref{lem:basic-arith}. The product is nonprincipal
and its parity \(j\) is fixed as \(d\) varies because \((-1/d)=-1\).
The extra factors passing from \(fd\) to \(Md\) are nonzero for
\(\Rp s>0\), by the proof of Lemma~\ref{lem:induce} and
Proposition~\ref{prop:cont}.

For \(X\) large put \(T=\sqrt{fX}\). If \(d\le X\) in this class
has a zero of the twist in the asserted region, Proposition~\ref{prop:twist-test}
and Lemma~\ref{lem:reciprocity} imply
\(I_{j,T}(\psi J_d)+I_{j,T}(\mu\psi J_d)>T^{2\epsilon}\).
Here \(fd\le T^2\). The sum of the left side over all positive integers
\(d\le X\) is \(O_\epsilon(T^{2+\epsilon})\) by Lemma~\ref{lem:theta-average}.
So the number of exceptional primes for \(\xi\) is
\(O_{\epsilon,M,Y}(X^{1-\epsilon/2})\). The same estimate holds for
their union over the finitely many \(\xi\). The sum of their reciprocal
primes in \([2^k,2^{k+1})\) is
\(O_{\epsilon,M,Y}(2^{-k\epsilon/2})\), summable in \(k\).
Corollary~\ref{cor:ap} gives an infinite reciprocal sum in the whole
class. Infinitely many primes survive the exclusion.
\end{proof}

\begin{corollary}\label{cor:prime-suffices}
For every positive integer \(q\), put \(M=18480q,\ c=M/2-1\).
Assume that for every prime \(d\equiv c\bmod M\) and every
\(1/2<\beta<1\) there is \(x_0(d,M,\beta)\) such that
\(R_{d,M}(x)\le x^{\beta-1}\) for \(x\ge x_0(d,M,\beta)\).
If this holds for every \(q\), then \(L(s,\chi)\ne0\) for every
character \(\chi\bmod q\) and \(\Rp s>1/2\), except at \(s=1\) for
the principal character, where it has a simple pole of residue
\(\varphi(q)/q\). In particular Dirichlet GRH follows.
\end{corollary}
\begin{proof}
Since \(16\mid M\), \(c^2-1=M(M/4-1)\) and \(c\equiv7\bmod 8\).
For \(1/2<\Rp s<1\), take \(1/2<\beta<\Rp s\),
\(Y=|\operatorname{Im}s|\). Theorem~\ref{thm:twist} supplies a
single prime \(d\) for which Theorem~\ref{thm:signed} applies.
It excludes all zeros at \(s\) for characters modulo \(M\).
Inducing a character modulo \(q\) to one modulo \(M\) multiplies
its \(L\)-function by nonvanishing finite Euler factors, by the
argument of Lemma~\ref{lem:induce}. This gives the claim in the
half-strip. Proposition~\ref{prop:edge} and
Proposition~\ref{prop:cont} give the assertion on \(\Rp s\ge1\).
Apply Theorem~\ref{thm:reduction}.
\end{proof}

\section{Cyclotomic places and their weights}

\begin{definition}\label{def:levels}
Fix \(q\ge1\), \(M=18480q\), and a prime \(d\equiv M/2-1\bmod M\).
Let \(K=\Q(\sqrt{-d})\subset\C\), choosing the square root with
positive imaginary part. For \(s\in\Z_{\ge1}\), put
\[
 c_s=\lfloor\log s/\log 3\rfloor+5,\quad U=3^{c_s}2^s,\quad
 N_1=MU,\quad E=K(\zeta_{N_1}),\qquad \zeta_h=\exp(2\pi i/h).
\]
An infinite place of \(E\) means a pair of complex-conjugate
embeddings \(E\to\C\). Denote the set of these pairs by \(\Phi\).
\end{definition}

\begin{lemma}\label{lem:cyclo}
For \(h\ge1\), the polynomial
\(P_h(X)=\prod_{a\in(\Z/h\Z)^\times}(X-\zeta_h^a)\)
is irreducible in \(\Q[X]\). Every unit \(a\bmod h\) acts on
\(\Q(\zeta_h)\) by the automorphism \(\zeta_h\mapsto\zeta_h^a\).
\end{lemma}
\begin{proof}
The factorization \(X^h-1=\prod_{k\mid h}P_k(X)\) groups the roots
by their orders. Division by monic polynomials and induction on \(h\)
give \(P_h\in\Z[X]\). Monic rational factors of a monic integral
polynomial are integral. Indeed the product of two integral
polynomials with relatively prime coefficients within each has the
same property, since reduction modulo any prime preserves their
nonzero product. Clearing denominators and extracting the gcd of the
coefficients in a rational factorization proves the claim by comparing
contents and leading coefficients.

Let \(A\) be a monic irreducible factor with root \(\zeta_h\), and
let \(\alpha\) be any root of \(A\). For a prime \(p\nmid h\),
\(\alpha^p\) is again primitive. If it is a root of a different
irreducible factor \(B\), then \(A\mid B(X^p)\). Reduction modulo
\(p\) gives \(\bar A\mid\bar B(X)^p\). Thus \(\bar A,\bar B\) share
an irreducible factor, and since \(AB\mid X^h-1\) the latter has
a repeated factor modulo \(p\). This contradicts its being coprime
to its derivative \(hX^{h-1}\) there. Existence of irreducible
factors and the divisibility property used here follow from polynomial
division over a field, by the least-degree gcd argument of
Lemma~\ref{lem:basic-arith}. Repeatedly applying this conclusion to
prime factors of a positive integer coprime to \(h\) shows that \(A\)
contains every primitive \(h\)-th root. Thus \(A=P_h\). Evaluation
at \(\zeta_h^a\) in \(\Q[X]/(P_h)\) gives the asserted automorphism,
whose inverse uses the inverse exponent. For \(h=1\) the polynomial
is \(X-1\).
\end{proof}

\begin{proposition}\label{prop:places}
We have \(d\nmid N_1\), and \(\eta_d(p)=1\) for every prime \(p\mid N_1\).
The infinite places form a group identified with
\(G=(\Z/N_1\Z)^\times\) by assigning to \(a\in G\) the place of
the embedding fixing \(K\) and sending \(\zeta_{N_1}\) to
\(\zeta_{N_1}^a\). The subgroup \(A=\ker(G\to(\Z/M\Z)^\times)\)
has size \(U\), and \(f=|G|=\varphi(M)U\).
\end{proposition}
\begin{proof}
All prime divisors of \(N_1\) divide \(M\), since \(6\mid M\).
The congruence on \(d\) gives \(d\equiv7\bmod8\), and at each odd
prime \(p\mid M\) it gives \(d\equiv-1\bmod p\).
Thus \(d\nmid N_1\) and \(\eta_d(p)=1\), including \(p=2\), by
Lemma~\ref{lem:reciprocity}. Put
\(\tau=\sum_{a\bmod d}\eta_d(a)\zeta_d^a\).
The Legendre character is primitive since \(d\) is prime and it
is nonprincipal. Lemma~\ref{lem:gauss} gives \(|\tau|^2=d\);
\(\bar\tau=\eta_d(-1)\tau=-\tau\), so \(\tau^2=-d\).
Consequently \(K\subset\Q(\zeta_d)\).
Both \(\zeta_d\) and \(\zeta_{N_1}\) lie in
\(\Q(\zeta_{dN_1})\). For each \(a\bmod N_1\) a unit, choose a
unit exponent equal to \(a\bmod N_1\) and \(1\bmod d\) by
Lemma~\ref{lem:basic-arith}. The induced automorphism of
Lemma~\ref{lem:cyclo} fixes \(K\) and acts as specified on
\(\zeta_{N_1}\). These are all embeddings extending the chosen
inclusion of \(K\), since the image of a primitive root must be
primitive. Each conjugate pair contains just one such embedding,
the two images of \(\sqrt{-d}\) being opposite. Composition transports
the group law on \(G\) to the places. Each unit class modulo \(M\) has
exactly \(U\) lifts modulo \(N_1=MU\), all units since the prime
divisors are the same. Counting proves the last two assertions.
\end{proof}

\begin{lemma}[A sphere in a cyclic subgroup]\label{lem:sphere}
Let \(G,H\) be abelian groups, \(\pi:G\to H\) a homomorphism, \(z,g\in G\),
and \(h\) the finite order of \(\pi(g)\). Let \(n,L,B\) be positive
integers, \(LB\ge2\) and \((LB)^n\le h\). There is a set \(D\subset
z\langle g\rangle\) with \(|D|\ge B^n/(nB^2)\) such that if
\(u,v,w\in D,\ 0<i<j\le L\), and
\(\pi(v/u)^j=\pi(w/u)^i\), then \(u=v=w\).
\end{lemma}
\begin{proof}
In \(\{0,\ldots,B-1\}^n\) some squared Euclidean norm occurs
at least \(B^n/(nB^2)\) times, since there are at most
\(1+n(B-1)^2\le nB^2\) possible norms. For a vector \(a\) on that
sphere put \(e(a)=\sum_{\ell=0}^{n-1}a_\ell(LB)^\ell\), and take
\(zg^{e(a)}\) in \(D\). The exponents are distinct integers in
\([0,h)\), so even their projections after removing \(z\) are distinct.
For \(u,v,w\) corresponding to vectors \(a,b,c\), the group relation
gives \(je(b)\equiv(j-i)e(a)+ie(c)\bmod h\). The quantities on both
sides have base-\(LB\) digits at most \(j(B-1)<LB\) and lie in
\([0,h)\). Equality holds in each digit. Thus \(b=(1-i/j)a+(i/j)c\).
For \(0<t<1\),
\[
 (1-t)\|a\|^2+t\|c\|^2-\|(1-t)a+tc\|^2=t(1-t)\|a-c\|^2
\]
by expansion. Equality of the three norms gives \(a=b=c\).
\end{proof}

\begin{proposition}\label{prop:thick}
Fix \(0<\epsilon<1/100\). For every sufficiently large \(x\) put
\[
 m=2^{\lceil(\log x)/(2\log 2)\rceil},\quad
 r=3^{\lceil5\log x/\log3\rceil},\quad
 s=22\log m/\log2,\qquad \ell=x^\epsilon,\quad b=\exp\ell .
\]
Then \(\sqrt x\le m\le2\sqrt x,\ x^5\le r\le3x^5,\ s\le32\log x\).
For the places in Definition~\ref{def:levels} there is a subset
\(D\subset G\) of size \(S\ge U/\ell\), all of whose elements reduce
to \(-1\bmod M3^{c_s}\). If \(u,v,w\in D,\ 0<i<j\le324s\), and
\((v/u)^j=(w/u)^i\bmod 2^s\), then \(u=v=w\).
\end{proposition}
\begin{proof}
The ceiling bounds give the stated ranges and
\(s\le(11/\log2)\log x+22\le32\log x\) eventually, since
\(\log2=\int_1^2\dd u/u>1/2\).
Put \(K_0=M3^{c_s}\), \(g=1+K_0\), and \(a_2\) equal to the exponent
of \(2\) in \(M\). The integer \(g\) is a unit modulo \(N_1\).
For all sufficiently large \(s\) its order modulo \(2^s\) is
\(h=2^{s-a_2}\). Indeed \(g-1\) has exact exponent \(a_2\ge4\).
Squaring a number congruent to \(1\bmod4\) raises the exponent of
its difference from one by precisely one, by factoring its difference
of squares. Raising to an odd positive power preserves that exponent
by the geometric sum. This proves the order assertion.

Take an integer \(n> \max(2,44/\epsilon)\),
\(L=324s,\ B=\lfloor h^{1/n}/L\rfloor\). For large \(x\),
\(h^{1/n}/L\ge2\), since \(s\to\infty\) and the exponential dominates
a fixed power of \(s\). Apply Lemma~\ref{lem:sphere} to reduction
modulo \(2^s\) with \(z=-1\). Every chosen element is \(-1\bmod K_0\).
Since \(B\ge h^{1/n}/(2L)\), the count gives
\[
 U/S\le n\,3^{c_s}2^{a_2}(2L)^{n-2}h^{2/n}
       \ll_{M,n}(1+\log x)^{n-1}x^{22/n}.
\]
Here \(3^{c_s}\le243s\), \(2^s=m^{22}\le 2^{22}x^{11}\).
The right side is at most \(x^\epsilon\) eventually. In these
comparisons any fixed power of \(\log x\) is \(o(x^\delta)\) for
\(\delta>0\): the exponential series at \(\delta\log x\) dominates
the term of one higher integer degree.
\end{proof}

\begin{definition}\label{def:ask}
Use \(G,A,D,S,U,f\) of Propositions~\ref{prop:places} and~\ref{prop:thick}.
Let \(a\ge1\) be integral and \(1\le x,x_1\le a\). Give \(g\in D\)
weight \(w(g)=a/x\) and \(g\in G\setminus D\) weight \(w(g)=a/x_1\).
Write \(g(n)=-n\bmod N_1\) if \(n\) is a unit, and \(g(n)=1\) otherwise.
With \(y_+=\max(y,0)\) define
\[
 \begin{aligned}
 P(\tau)&=\sum_{n\le a}\lambda_{d,1}(n)\big(a/n-w(g(n)\tau)\big)_+,
 &\quad \bar w&=\frac1f\sum_{g\in G}\log w(g),\\
 P_0&=a(\log a+\kappa_{d,1}-\bar w),
 &\quad \mathcal R&=\sum_{\tau\in A}\big(P(\tau)-P_0\big).
 \end{aligned}
\]
\end{definition}

\begin{proposition}\label{prop:rem-id}
For the weights of Definition~\ref{def:ask},
\[
 \mathcal R=a\{S(R_{d,M}(x)-R_{d,M}(x_1))+U R_{d,1}(x_1)\}.
\]
\end{proposition}
\begin{proof}
If \((n,N_1)>1\) and \(\Lambda(n)\ne0\), then \(n\) is a power
of a prime dividing \(N_1\) and \(\lambda_{d,1}(n)=0\) by
Proposition~\ref{prop:places}. For a unit, translation by \(g(n)\)
maps \(A\) bijectively onto the fiber of \(-n\bmod M\). It therefore
hits all \(S\) elements of \(D\) precisely when \(n\equiv1\bmod M\),
and otherwise hits none. Hence \(\sum_{\tau\in A}P(\tau)/a\) equals
\[
 U\sum_{n\le x_1}\lambda_{d,1}(n)(1/n-1/x_1)
 +S\left\{\sum_{n\le x}\lambda_{d,M}(n)(1/n-1/x)
           -\sum_{n\le x_1}\lambda_{d,M}(n)(1/n-1/x_1)\right\}.
\]
For \(z=x,x_1\), the positive part truncates the corresponding sum
exactly at \(n\le z\), because \(z\le a\). Substitute
Definition~\ref{def:prime-error}. The constants \(\kappa_{d,M}\)
cancel. The other terms in addition to the asserted remainders are
\[
 U(\log x_1+\kappa_{d,1})+
                   S(\log x-\log x_1)/\varphi(M).
\]
They equal \(UP_0/a\), since
\(\bar w=\log a-\log x_1-(S/f)(\log x-\log x_1)\) and
\(f=\varphi(M)U\) by Proposition~\ref{prop:places}. Subtraction
gives the identity.
\end{proof}

\begin{theorem}\label{thm:geometric-reduction}
For each \(M,d,\epsilon\) as in Proposition~\ref{prop:thick}, choose
the places and scales there and set \(a=N_1\lceil xb/N_1\rceil\),
\(t=10^6\log x\). Suppose that for all sufficiently large \(x\)
and every \(v\in[t,t^2]\) the quantity \(\mathcal R\) with \(x_1=\e^v\)
satisfies \(m\mathcal R\le38fa\ell^4\). Then
\begin{equation}\label{eq:one-side}
 R_{d,M}(x)\le 38\varphi(M)x^{-1/2+5\epsilon}+x^{-10^6}
\end{equation}
eventually. If the hypothesis holds for every \(M,d,\epsilon\) in
these ranges, then Dirichlet GRH holds.
\end{theorem}
\begin{proof}
As \(x\to\infty\), \(t^2\le\ell\), so \(1\le x,x_1\le a\)
throughout the window. Proposition~\ref{prop:thick} gives
\(0<U/S\le\ell=x^\epsilon\le\e^t\). By Corollary~\ref{cor:usable-window}
there exists \(v\in[t,B_0t]\) with
\(R_{d,M}(\e^v)-(U/S)R_{d,1}(\e^v)\le\e^{-t}\);
the whole interval lies in \([t,t^2]\) eventually since \(B_0\) depends
only on \(d,M\). Divide the hypothesized estimate at this \(v\) by
\(maS\) and use Proposition~\ref{prop:rem-id}. Since
\(f/(mS)\le\varphi(M)\ell/\sqrt x\), this gives \eqref{eq:one-side}.
For \(1/2<\beta<1\) choose \(0<\epsilon<1/100\) so small that
\(5\epsilon<\beta-1/2\). The right side of \eqref{eq:one-side} divided
by \(x^{\beta-1}\) tends to zero. Corollary~\ref{cor:prime-suffices}
now applies.
\end{proof}

\section{Primes for the auxiliary levels}

\begin{lemma}\label{lem:uniform-disc}
Fix an integer \(B\ge4\) divisible by four and a positive integer
\(A_0\). For all sufficiently small \(\delta>0\), for every
\(Q=B^k\le\delta^{-A_0}\) with \(k\ge2\), every character \(\chi\bmod Q\),
and every \(z\) with \(|z-1|\le\delta\),
\[
          \delta\,|(L^\circ)'(z,\chi)|<|L^\circ(z,\chi)|.
\]
\end{lemma}
\begin{proof}
First suppose \(\chi^2\) is principal. Every unit \(u\equiv1\bmod B^2\)
is a square modulo \(B^k\). At an odd \(p\mid B\) start with the root
\(1\bmod p\), then lift \(y^2\equiv u\bmod p^j\) by adding \(h p^j\)
to \(y\): the new coefficient modulo \(p\) can be chosen arbitrarily
since \(2y\) is invertible. At \(p=2\) start at \(1\bmod 8\).
For \(j\ge3\), changing an odd \(y\) to \(y+2^{j-1}\) changes its
square by \(2^j\bmod 2^{j+1}\). These lifts and
Lemma~\ref{lem:basic-arith} give the assertion (smaller powers follow
by reduction). Therefore \(\chi\) is trivial on the kernel of
reduction to \(B^2\). By surjectivity in Lemma~\ref{lem:induce} it
comes from a character modulo \(B^2\). The two moduli have the same
prime divisors, so the two \(L^\circ\)-functions coincide.
The finitely many functions so obtained are nonzero at \(1\) by
Proposition~\ref{prop:edge}. Continuity of their values and derivatives
gives the assertion for all of them on sufficiently small disks.

For a nonprincipal \(\psi\bmod Q\), and any fixed \(0<r<1/8\), we have
\[
 \sup_{|w-1|\le2r}|L(w,\psi)|\ll_r Q^{2r},\qquad
 \sup_{|w-1|\le r}|L'(w,\psi)|\ll_r Q^{2r}.
\]
Indeed truncate the sum at \(Q\). Its initial terms have absolute sum
at most \(1+\int_1^Q t^{-1+2r}\dd t\). All complete character sums
vanish by the proof of Proposition~\ref{prop:cont}. The remaining
partial sums, starting after \(Q\), are at most \(Q\) in modulus, so
integration as in that proposition bounds the tail by
\(|w|Q\int_Q^\infty t^{-1-\Rp w}\dd t\ll_r Q^{2r}\).
The derivative estimate follows from Lemma~\ref{lem:analytic-tools}
on circles of radius \(r\).

Now suppose \(\chi^2\) is nonprincipal. Choose \(r\) so small that
\(10rA_0<1\), and take \(\delta<r\). If the desired inequality
failed at \(z\), the preceding derivative estimate there and along the
segment to \(1+\delta\) would give
\(|L(1+\delta,\chi)|\ll_r\delta Q^{2r}\). The same size estimate
on the disk gives \(|L(1+\delta,\chi^2)|\ll_r Q^{2r}\).
Since \(\zeta(1+\delta)\le1+1/\delta\), the Euler-product inequality
in Proposition~\ref{prop:edge} at \(t=0\) yields
\(1\ll_r\delta Q^{10r}\le\delta^{1-10rA_0}\), impossible for
sufficiently small \(\delta\).
\end{proof}

\begin{theorem}\label{thm:aux-window}
Let \(4\mid B\), \(B>0\), and fix a positive integer \(A_0\). For all
sufficiently large \(W\), all \(k\ge2\) with \(Q=B^k\le W^{A_0}\),
and all real \(V\in[W,W^{A_0}]\), there is a prime \(p\equiv1\bmod Q\)
with \(V\le\log p\le V+W\).
\end{theorem}
\begin{proof}
Put \(r=W^{-1/2}\), \(a_n=\Lambda(n){\bf1}_{n\equiv1\bmod Q}\), and
\[
 H(z)=\frac{1-z\sum_{\chi\bmod Q}\ell_\chi(1+z)}{\varphi(Q)}.
\]
The Euler logarithm in Proposition~\ref{prop:polar} and
Lemma~\ref{lem:characters} give
\(H(z)/z=\sum_{n\ge1}a_n n^{-1-z}\) for \(\Rp z>0\).
The numerator's \(1\) accounts for regularization of the sole principal
character. Lemma~\ref{lem:uniform-disc} at \(\delta=r\) with exponent
\(2A_0\) shows that \(H\) is holomorphic on a neighbourhood of
\(|z|\le r\) and \(|H|\le2\) there on the closed disk, since
\(|\ell_\chi|<1/r\) and there are \(\varphi(Q)\) summands. Thus for
\(D=1/V\le r\), positivity gives \(\sum a_n n^{-1-D}\le2V\).

Suppose there is no such prime. Subtract the initial terms and put
\[
 P(z)=H(z)-z\sum_{\log n<V}a_n n^{-1-z},\qquad
 B_1(z)=2-Vz-\tfrac{\sqrt W}{4}(1+\sqrt W z)^2 .
\]
On \(|z|=r\), with \(c=\Rp z\), direct expansion gives
\(\Rp B_1(z)=E_1=2-(V+\tau)c\), where
\(\tau=\tfrac W2(1+\sqrt W c)\in[0,W]\).
If \(c\le D\), then \(E_1\ge0\) and \(1-Vc\le E_1\), because
\(\tau c\le W/V\le1\) when \(c>0\). For \(\log n<V\),
\((D-c)\log n\le1-Vc\). Hence
\[
 |P(z)|\le 2+|z|\e^{E_1}\sum_n a_n n^{-1-D}
           \le (2+2r V)\e^{E_1}.
\]
If \(c>D\), the Dirichlet series gives
\(P(z)=z\sum_{\log n\ge V}a_n n^{-1-z}\).
When \(\log n>V+W\),
\((D-c)\log n\le (V+W)(D-c)\le E_1\).
Terms with \(V\le\log n\le V+W\) and \(a_n\ne0\) are proper
prime powers by the supposition. For them
\((-1/4-c)\log n\le -V/4-Vc\le E_1\) for large \(W\), since
\(\tau c\le\sqrt W\le2+V/4\).
The sum of \(\Lambda(n)n^{-3/4}\) over proper prime powers is at most
\((1-2^{-3/4})^{-1}\sum_{m\ge2}(\log m)m^{-3/2}<\infty\).
It follows that \(|P(z)|\le r(2V+O(1))\e^{E_1}\).
Thus \(|P\exp(-B_1)|\ll1+V\) on the whole circle. Its center value
is bounded by the same quantity by Lemma~\ref{lem:analytic-tools}.
Since \(P(0)=1/\varphi(Q)\), this implies
\(\exp(\sqrt W/4-2)\ll\varphi(Q)(1+V)\ll W^{2A_0}\).
The exponential series contradicts this for large \(W\), uniformly
in \(k,V\).
\end{proof}

\begin{lemma}\label{lem:quadratic-norm}
Fix \(d\) as in Definition~\ref{def:levels}. There is an integer
\(H_d\ge1\) such that among any \(H_d+1\) distinct odd primes
\(p\equiv1\bmod d\) there are two distinct \(p_0,p_1\) and an algebraic
integer \(\pi\in K\) with \(\pi\bar\pi=p_0p_1\).
\end{lemma}
\begin{proof}
Put \(\omega=(1+\sqrt{-d})/2\) and \(O=\Z+\Z\omega\).
This is the ring of integers of \(K\). Indeed \(\omega\) satisfies
\(X^2-X+(d+1)/4=0\), and every \(z\in O\) has integral trace and
norm and hence satisfies a monic integral quadratic. Conversely
if \(z=x+y\sqrt{-d}\), \(x,y\in\Q,\ y\ne0\), is a root of a monic
integral polynomial, its monic minimal polynomial
\(X^2-2xX+x^2+dy^2\) is integral by the factor argument of
Lemma~\ref{lem:cyclo}. Thus \(a=2x\) and \(d(2y)^2\) are integers.
The squared denominator of \(2y\) in lowest terms divides the prime
\(d\), so \(b=2y\) is integral. The condition \(4\mid a^2+db^2\)
and \(d\equiv3\bmod4\) show that \(a,b\) have the same parity.
Then \(z=(a-b)/2+b\omega\in O\). Rational algebraic integers are
integers by the same monic factor argument.

For a prime \(p\) in the assertion, \((-d/p)=(p/d)=1\) by the proof
of Lemma~\ref{lem:reciprocity}. Thus the polynomial of \(\omega\)
has two distinct roots modulo \(p\). Evaluation at them identifies
\(O/pO\) with \(\mathbf F_p\times\mathbf F_p\), since a linear
polynomial can attain prescribed values at two distinct points.
Let \(I_p,\overline I_p\) be the two kernels in \(O\), interchanged
by conjugation. They have index \(p\), their sum is \(O\), and
their intersection is \(pO\). Their product is their intersection:
if \(a+b=1\) with \(a\in I_p,b\in\overline I_p\), multiply by
an element of the intersection to express it as a sum of products.
In particular \(I_p\overline I_p=pO\).

Two of the \((1+\lceil\sqrt p\rceil)^2\) elements \(i+j\omega\),
\(0\le i,j\le\lceil\sqrt p\rceil\), agree modulo \(I_p\).
Their difference \(z_p\in I_p\) is nonzero and
\(|z_p|\le2(1+|\omega|)\sqrt p\). Then
\(J_p=(z_p/p)\overline I_p\) is an ideal contained in \(O\).
As a lattice it has index \(|z_p|^2/p\) there.
To verify the calculation, the index of a full sublattice of
\(\Z^2\) with basis columns \(v,w\) is \(|\det(v,w)|\). Integral
row operations from the Euclidean algorithm send \(v\) to \((h,0)\);
successive reduction of the second and then the first coordinate
gives exactly \(|h w'_2|\) cosets, where \(w'\) is the transformed
second column. In the complex plane, multiplication by \(u\)
has real determinant \(|u|^2\), so areas in a fixed lattice basis
are multiplied by that factor. Here \(\overline I_p\) has a basis
given by \(p,\omega-b_p\), for the corresponding root \(b_p\).

The indices of the \(J_p\) are therefore bounded by
\(D_d=\lceil4(1+|\omega|)^2\rceil\). There are only finitely many such
sublattices. An additive subgroup of index \(h\) contains \(hO\),
since the order of each element of the quotient divides \(h\) by
counting cosets of its cyclic subgroup. For \(h\le D_d\) there
are only finitely many subsets of \(O/hO\).
Take \(H_d\) to be the number of sublattices of these indices.
Two distinct primes \(p_0,p_1\) give the same \(J\).
Now \(J I_{p_i}=z_{p_i}O\), hence
\(z_{p_1}I_{p_0}=z_{p_0}I_{p_1}\). Multiplying by
\(\overline I_{p_1}\) gives
\(I_{p_0}\overline I_{p_1}=(p_1 z_{p_0}/z_{p_1})O\).
Its generator belongs to \(O\). Equality of the two indices
\(|z_{p_i}|^2/p_i\) shows that this generator has norm \(p_0p_1\).
\end{proof}

\begin{theorem}[Auxiliary arithmetic data]\label{thm:data}
With \(M,d,\epsilon,\ell,b,x,N_1,s,U,m,r,D\) as in
Definition~\ref{def:levels} and Proposition~\ref{prop:thick}, fix any
\(C>0\). For every sufficiently large \(x\), put \(M_\omega=(Md)^{s+6}\)
and \(t=10^6\log x\). There are distinct primes \(g_0,g_1\in[b^{10},b^{11}]\)
congruent to \(1\bmod M_\omega^2\) and \(1\bmod N_1d\), and an integer
\(\pi\) of \(K\) of norm \(g_0g_1\). For every real \(V\) satisfying
\(100\ell\le V\le2^s\), there is a prime \(p\equiv1\bmod M_\omega^2\)
with \(V\le\log p\le V+\sqrt\ell\). The same data work for all
\(v\in[t,t^2]\), and throughout that interval
\[
 10^{12}C\le\log\ell,\quad 100\log\ell\le\log x,\quad
 10^6\log x\le v,\quad 10^6v\le\ell,\quad 6\log v\le\log\ell .
\]
They also work with \(v=x^{\epsilon/6}\), for which the last
inequality is an equality.
\end{theorem}
\begin{proof}
Put \(W=x^{\epsilon/6}\), so \(\ell=W^6\). Since
\(s\le32\log x\), there is an integer \(A_0>7\), depending only on
\(M,d,\epsilon\), for which \(M_\omega^2\le W^{A_0}\) and
\(2^s\le W^{A_0}\) eventually. Indeed the logarithms of the two
left sides are \(2(s+6)\log(Md)\) and \(s\log 2\), and
\(\log W=(\epsilon/6)\log x\). We have \(N_1d\mid M_\omega\):
after removing the initial \(Md\), the factor \(2^s3^{c_s}\) divides
\((Md)^{s+5}\) since \(6\mid M\) and \(c_s\le s+5\).
Apply Theorem~\ref{thm:aux-window} with base \(Md\) and exponent
\(2(s+6)\). For \(V\in[100\ell,2^s]\), its range hypotheses hold,
and the permitted width \(W\) is at most \(\sqrt\ell\).
For each \(0\le i\le H_d\) use the same theorem with
\(V_i=10\ell+2iW\), where \(H_d\) is from Lemma~\ref{lem:quadratic-norm}.
Eventually all intervals \([V_i,V_i+W]\) are disjoint subsets
of \([10\ell,11\ell]\) with \(V_i\le W^{A_0}\). Their primes satisfy
the two required congruences, and the lemma supplies two and their norm
generator. Finally \(\log\ell=\epsilon\log x\), \(100\epsilon<1\)
and \(t^2\) is a fixed quadratic polynomial in \(\log x\).
The power comparison proved in Proposition~\ref{prop:thick} gives
\(10^6t^2\le\ell\), \(6\log(t^2)\le\log\ell\) and
\(t^2\le W\) eventually. For \(v=W\) we have \(10^6W\le W^6\)
eventually. This proves all the scale conditions.
\end{proof}

\section{The level prime and the arithmetic part of the height}

\begin{lemma}\label{lem:rd}
For the field \(E\) in Definition~\ref{def:levels}, the ring \(O_E\)
of algebraic integers is a free abelian group of rank \(2f\). Write
\(\Delta_E=\det(\operatorname{Tr}_{E/\Q}(b_i b_j))\) for the
discriminant of any integral basis, and
\(\operatorname{rd}_E=|\Delta_E|^{1/(2f)}\). Then
\(1\le\operatorname{rd}_E\le 2f\max(1,|\omega|)^2\), where
\(\omega=(1+\sqrt{-d})/2\).
\end{lemma}
\begin{proof}
The elements \(b_{e,j}=\omega^e\zeta_{N_1}^j\), \(e=0,1,\ 0\le j<f\),
are a rational basis. Indeed the minimal polynomial of \(\zeta_{N_1}\)
over \(K\) has degree \(f\), since Proposition~\ref{prop:places}
gives \(f\) distinct images and Lemma~\ref{lem:cyclo} gives a polynomial
of that degree. Let \(\mathcal B\) be the integer span of this basis.
It is closed under multiplication by the two monic relations; its
elements are integral by the characteristic polynomial of their
multiplication matrices. More generally sums and products of algebraic
integers are integral: the powers of two such integers, reduced by
their monic relations, span a finite integral module containing \(1\)
and stable under both. The determinant of \(XI-A\) annihilates any
element acting by the matrix \(A\) on these generators, by multiplying
by the adjugate. This argument needs no independence of the generators.

Let \(D\) be the matrix with rows the \(2f\) embeddings applied to the
basis. It is invertible. Group the rows by the image of \(\omega\);
in each block all \(f\) distinct primitive roots occur, so a polynomial
of degree less than \(f\) vanishing there vanishes identically. The two
distinct images of \(\omega\) then separate the two coefficients of
each power. Multiplication by a field element is diagonal in these
evaluation coordinates. Its trace is therefore the sum of its images,
and for an algebraic integer this trace is a rational algebraic integer,
hence an integer by Lemma~\ref{lem:quadratic-norm}. The trace matrix
on \(\mathcal B\) is \(D^{\mathsf T}D\), invertible. Thus
\(\mathcal B\subset O_E\subset\{y:\operatorname{Tr}(y\mathcal B)\subset\Z\}\),
the last group a full lattice obtained by applying the inverse trace
matrix to \(\Z^{2f}\). Any subgroup of a finite-rank free abelian
group is free: project to the first coordinate, lift a generator of
its image when nonzero, and apply induction to the kernel. This gives
an integral basis of \(O_E\). Changing between integral bases has
integer determinant with integer inverse and so preserves the
discriminant. The integral inclusion matrix of \(\mathcal B\) into
\(O_E\) has nonzero integer determinant, hence
\(1\le|\Delta_E|\le|\det D|^2\). All entries of \(D\) have modulus
at most \(\max(1,|\omega|)\); apply Lemma~\ref{lem:hadamard} to its
columns and take the \(2f\)-th root.
\end{proof}

\begin{proposition}\label{prop:level-prime}
Use the data of Theorem~\ref{thm:data}, \(v\in[t,t^2]\),
\(x_1=\e^v,\ a=N_1\lceil xb/N_1\rceil\), and Definition~\ref{def:ask}.
Write \(g=g_0g_1\),
\[
 G_0=\log\!\left(\frac{g\,\operatorname{rd}_E}{2\pi}\right),
 \qquad V_0=m\bar w+(m-1)G_0-\tfrac{15}{2}\ell .
\]
For all sufficiently large \(x\), uniformly in this interval of \(v\),
there is a prime \(N\equiv1\bmod M_\omega^2\) such that
\[
 V_0\le\log N\le V_0+\sqrt\ell,\qquad m\ell\le\log N\le30m\ell .
\]
There is \(\pi_N\in O_E\) with \(\pi_N\bar\pi_N=N\).
\end{proposition}
\begin{proof}
By Proposition~\ref{prop:thick}, \(N_1\ll_M (1+\log x)x^{11}\)
and \(f\le N_1\). Hence \(xb\le a\le xb+N_1\le2xb\) eventually.
Since \(t^2=o(\ell)\) by Proposition~\ref{prop:thick},
all weights satisfy \(1\le w(\tau)\le2b\).
The prime sizes in Theorem~\ref{thm:data} and the bounds in
Lemma~\ref{lem:rd} give \(19\ell\le G_0\le23\ell\) eventually:
\(\log g\in[20\ell,22\ell]\) and
\(0\le\log\operatorname{rd}_E=O_{M,d}(1+\log x)=o(\ell)\).
It follows from \(m\to\infty\) and \(0\le\bar w\le2\ell\) that
\(m\ell\le V_0\le25m\ell\). This lies in \([100\ell,2^s]\) eventually
as \(2^s=m^{22}\) and \(\ell=x^\epsilon\). Apply Theorem~\ref{thm:data}
and enlarge the upper bound to \(30m\ell\).

Here \(N\equiv1\bmod4\), so choose \(u\in\Z\) with \(u^2\equiv-1\bmod N\)
by Lemma~\ref{lem:reciprocity}. The ring \(\Z[i]\) admits division
with a smaller squared norm as remainder: round both coordinates of
a complex quotient to their nearest integers, with squared error at
most \(1/2\). Take the Euclidean gcd \(c\) of \(N\) and \(u+i\).
It is a nonunit, for otherwise the Bezout identity multiplied by
\(u-i\) would give \(N\mid u-i\). It is not associated to \(N\)
since \(N\nmid u+i\). Its positive integer norm is therefore a
nontrivial proper divisor of \(N^2\), hence \(N\). Take \(\pi_N=c\).
Since \(4\mid N_1\), \(i=\zeta_{N_1}^{N_1/4}\in O_E\).
\end{proof}

\begin{definition}\label{def:heights}
Continue with the notation of Proposition~\ref{prop:level-prime}.
For \(0\le y\le a\) let \(P(y,\tau)\) be the sum \(P(\tau)\) of
Definition~\ref{def:ask} with \(a/n\) replaced by \(y/n\) and all
other quantities fixed. Put \(k=r+1-m,\ h_2=r^8,\ h=2h_2,\ t_1=5r\).
For integer \(j\ge0\), set \(w_j(\tau)=N_1\lceil w(\tau)/N_1\rceil+jh+h_2\).
Write \(\gamma=\lim_{n\to\infty}(\sum_{k=1}^n 1/k-\log n)\), whose
existence follows by integral comparison as in Lemma~\ref{lem:gamma}, and
\[
 \begin{aligned}
 \Omega&=\frac{2\pi}{\operatorname{rd}_E}
              \exp\!\big(\gamma-L'(1,\eta_d)/L(1,\eta_d)\big),&
 J(u)&=\log(u!)-u\log(\Omega/g)\quad(u\in\Z_{\ge0}),\\
 D_0&=\sum_{\tau\in\Phi}\{J(w_0(\tau))-P(w(\tau),\tau)\},&
 H&=mkD_0,\\
 h_0&=H+m\sum_{\tau\in A}\{J(a)-P(a,\tau)\}-Ua\log N ,&
 B_0'&=fa\ell^3 .
 \end{aligned}
\]
The quotient defining \(\Omega\) is real and finite by
Proposition~\ref{prop:edge}, as \(\eta_d\) is real.
\end{definition}

\begin{proposition}\label{prop:height-cost}
For all sufficiently large \(x\), uniformly in \(t\le v\le t^2\),
\[
 \begin{split}
 0\le w_j(\tau)-w(\tau)&\le x^{42},\qquad
 w_j(\tau)\le2b,\qquad 0\le j\le2m,\\
 h_0&\le H-m\mathcal R+B_0',\qquad
 |H|+m|\mathcal R|\le b^2x^{99}.
 \end{split}
\]
\end{proposition}
\begin{proof}
Rounding adds at most \(N_1\). The largest additional shift is
\((4m+1)r^8\le(8\sqrt x+1)(3x^5)^8\), hence the total is at most
\(x^{42}\) eventually. Also \(x_1\ge x\), so
\(w(\tau)\le a/x\le b+N_1/x\). Since \(b=\exp(x^\epsilon)\),
the asserted upper bound follows.

In Definition~\ref{def:prime-error} at modulus one, the untwisted
regularization is \((s-1)\zeta(s)\). At \(s=1\) its value is \(1\)
and its derivative, by the integral in Proposition~\ref{prop:cont},
is \(1+\int_1^\infty(\lfloor u\rfloor-u)u^{-2}\dd u=\gamma\).
Indeed integration of \(\lfloor u\rfloor=\sum_{k\le u}1\) up to
the integer \(n\) gives \(\sum_{k\le n}1/k-1\) for the integral of
\(\lfloor u\rfloor/u^2\). Therefore
\(\kappa_{d,1}=L'(1,\eta_d)/L(1,\eta_d)-\gamma-1\),
and \(J(a)=\log(a!)+a(G_0+1+\kappa_{d,1})\).
Monotonicity of \(\log\) gives
\(\log(a!)\le\log a+\int_1^a\log u\dd u
=a\log a-a+1+\log a\), also for \(a=1\).
Consequently \(J(a)-P_0\le a(G_0+\bar w)+1+\log a\).
Substitute \(\sum_{\tau\in A}P(a,\tau)=\mathcal R+UP_0\).
The lower bound \(\log N\ge V_0\) now gives
\[
 h_0-H+m\mathcal R
 \le U\{m(1+\log a)+a(G_0+15\ell/2)\}.
\]
We have \(1+\log a\le 3\ell\), \(m\le a\), \(G_0\le23\ell\),
and \(U\le f\). This is at most \(fa\ell^3\) eventually.

For the crude bound, \(0\le P(y,\tau)\le2y\log a(1+\log a)\) by
the harmonic estimate in Proposition~\ref{prop:polar}, and
\(|P_0|\ll_d a(1+\ell)\). For \(1\le u\le2b\),
\(|J(u)|\le u(\log u+|\log(\Omega/g)|)\ll_d b(1+\ell)\),
since \(-\log(\Omega/g)=G_0+1+\kappa_{d,1}\).
It follows that \(|H|+m|\mathcal R|\ll_{M,d}
(mkf b+mUa)(1+\ell)^2\).
Use \(k\le r+1,\ f\le x^{12},\ m\le2\sqrt x,\ a\le2xb\),
\(r\le3x^5\). The last bound is \(O_{M,d}(b x^{18}(1+\ell)^2)\)
and is at most \(b^2x^{99}\) eventually.
\end{proof}

\begin{proposition}\label{prop:local-ask}
For each prime \(p\), \(0\le y\le a\), and \(\tau\in G\), define
\[
 A_p(y,\tau)=\tfrac12\sum_{j\ge1}(1-\eta_d(p)^j)
                (y/p^j-w(g(p^j)\tau))_+ .
\]
Then \(P(y,\tau)=2\sum_{p<x_1}A_p(y,\tau)\log p\).
Define the nonnegative gain
\[
 G_p=m\left(\sum_{\tau\in A}A_p(a,\tau)
                  + k\sum_{\tau\in\Phi}A_p(w(\tau),\tau)\right).
\]
It vanishes unless \(p<x_1\) and \(\eta_d(p)\ne1\).
\end{proposition}
\begin{proof}
Since \(\min w=a/x_1\), a positive summand requires \(p^j<x_1\le a\).
The definition of \(\Lambda\) decomposes \(P(y,\tau)\) into exactly
the indicated prime-power terms with coefficient \(\log p\).
At an inert prime \(\eta_d(p)=-1\) just the odd exponents contribute,
each with coefficient one in \(A_p\). At \(p=d\) all exponents
have coefficient \(1/2\), and at \(\eta_d(p)=1\) every coefficient
vanishes.
\end{proof}

\begin{proposition}\label{prop:loss}
For a prime \(p<x_1\) with \(\eta_d(p)\ne1\), set
\[
 e_p=3fa\ell^2\,p^{-1}
       \begin{cases}m & p\equiv 1\ \text{or }-1\bmod m,\\
                     1 & \text{otherwise}.\end{cases}
\]
Set \(e_N=4mfb\), and set \(e_p=0\) at all other primes.
Eventually these rules have disjoint ranges, and
\(\sum_p e_p\log p\le B_0'\). There is a positive rational \(c\) with
\[
 \log c\le B_0'-\sum_p G_p\log p,\qquad
 p^{G_p-e_p}\le |c|_p<p^{G_p-e_p+1}\quad\text{for every prime }p.
\]
Here \(|c|_p=p^{-\operatorname{ord}_p(c)}\), with
\(\operatorname{ord}_p\) the exponent in the numerator minus the
exponent in the denominator.
\end{proposition}
\begin{proof}
We have \(\log N\ge m\ell>v=\log x_1\) eventually by
Proposition~\ref{prop:level-prime} and Theorem~\ref{thm:data}.
The sum of \(1/p\) up to \(x_1\) is at most \(1+v\). In the two
exceptional classes the primes are of the form \(mj\pm1\) with
\(1\le j\le x_1\) for large \(m\). Then \(m/p\le2/j\),
so their total weighted sum is at most \(4(1+v)\).
Bounding \(\log p\) by \(v\) shows that the contribution from
\(p<x_1\) is at most \(15fa\ell^2 v(1+v)\).
By Theorem~\ref{thm:data}, \(v^6\le\ell\), so its ratio to
\(fa\ell^3\) tends uniformly to zero. The term at \(N\) is at most
\(120m^2fb\ell\le480fa\ell\), by \(m^2\le4x,\ a\ge xb\).
This proves the sum bound. Put \(c=\prod_p p^{\lfloor e_p-G_p\rfloor}\).
Both sequences have finite support by Proposition~\ref{prop:local-ask},
so this is a positive rational. The two inequalities follow from
\(e_p-G_p-1<\lfloor e_p-G_p\rfloor\le e_p-G_p\) and the sum bound.
\end{proof}

\section{The definite unitary group and its columns}

\begin{definition}\label{def:adele}
For a rational prime \(p\), \(\Z_p=\varprojlim_j\Z/p^j\Z\) is the
ring of compatible residues and \(\Q_p=\Z_p[1/p]\).
Put \(E_p=E\otimes_\Q\Q_p,\ O_p=O_E\otimes_\Z\Z_p\). For the rank
\(k=r+1-m\) of Definition~\ref{def:heights}, use the standard
hermitian form \(h(u,v)=\sum_{i=1}^k\bar u_i v_i\).
Use the same involution on \(E_p\), acting trivially on \(\Q_p\).
Put
\[
 \begin{aligned}
 \Gamma&=\{g\in\mathrm{GL}_k(E):\bar g^{\mathsf T}g=1\},&
 G_p^u&=\{g\in\mathrm{GL}_k(E_p):\bar g^{\mathsf T}g=1\},\\
 K_p^0&=\{g\in G_p^u:gO_p^k=O_p^k\}.
 \end{aligned}
\]
Let \(G^u\) consist of tuples \((g_p)_p\) with \(g_p\in G_p^u\)
and \(g_p\in K_p^0\) for all but finitely many \(p\). Write
\(K^0=\prod_p K_p^0\). The diagonal copy of \(\Gamma\) lies in
\(G^u\), since the entries of a rational matrix and its inverse in an
integral basis have only finitely many denominator primes.
\end{definition}

\begin{lemma}\label{lem:split-level}
If a prime \(p\equiv1\bmod N_1d\), then \(O_p\) is a product of
\(2f\) copies of \(\Z_p\), paired without fixed points by conjugation.
Choosing one from each pair identifies \(G_p^u\) with
\(\mathrm{GL}_k(\Q_p)^f\), and \(K_p^0\) with
\(\mathrm{GL}_k(\Z_p)^f\). For \(p=g_0,g_1\) choose the component
in each pair where \(\bar\pi_g\) is divisible by \(p\), with
\(\pi_g\) the integer of norm \(g_0g_1\) from Theorem~\ref{thm:data}.
For \(p=N\) choose the component where \(\pi_N\) is divisible by \(N\).
These choices exist uniquely.
\end{lemma}
\begin{proof}
A finite multiplicative subgroup of a field is cyclic. If its exponent
is \(e\), for each prime power exactly dividing \(e\) an element of
that order is obtained by taking a power of a suitable group element.
Multiplication of these elements of relatively prime orders gives an
element of order \(e\). All group elements are roots of \(X^e-1\), so
the group has at most \(e\) elements and is generated by this element.
Consequently \(\mathbf F_p^\times\) contains \(N_1\)-th roots of exact
order \(N_1\). All the \(N_1\)-th roots are distinct since \(p\nmid N_1\).
The factorization in Lemma~\ref{lem:cyclo}, applied successively to
divisors of \(N_1\), shows that \(P_{N_1}\) has exactly the primitive
ones as roots modulo \(p\). The polynomial of \(\omega\) also has
distinct roots modulo \(p\), as in Lemma~\ref{lem:quadratic-norm}.
A simple root modulo \(p\) of any integral polynomial \(F\) lifts
uniquely at each step from \(p^j\) to \(p^{j+1}\), because
\(F(y+p^jb)\equiv F(y)+p^jb F'(y)\bmod p^{j+1}\).
The compatible sequence gives a simple root in \(\Z_p\).
An element of \(\Z_p\) nonzero modulo \(p\) is a unit, by compatible
inversion at each finite level. Every nonzero element is a unit times
\(p^a\): shift its digits after the first nonzero one. This also shows
that \(\Q_p\) is a field.

In the order \(\mathcal B\) from Lemma~\ref{lem:rd}, reduction at
\(p\) is thus evaluation on the grid consisting of two choices for the
root of the polynomial of \(\omega\) and \(f\) for the root of
\(P_{N_1}\). This evaluation, also over \(\Z_p\), is an isomorphism
onto the product: one-variable interpolation uses only differences
of roots which are units, and can be applied successively in the two
variables. The trace matrix on this product is nondegenerate modulo
\(p\) since it is the ordinary coordinate pairing in its idempotent
basis. Thus \(p\) does not divide the trace determinant on
\(\mathcal B\), and hence does not divide its change-of-basis determinant
in \(O_E\), by Lemma~\ref{lem:rd}. Its matrix is invertible over
\(\Z_p\), giving the same product for \(O_p\). Conjugation interchanges
the two images of \(\omega\) and inverts the chosen root of unity.
A unitary matrix has component \((A^{-1})^{\mathsf T}\) at the partner
of any component \(A\). This proves the identifications. At \(p=g_i\)
the valuations of the two conjugate components of \(\bar\pi_g\) are
nonnegative integers summing to \(1\), since \(g_0,g_1\) are distinct.
The identical argument applies to \(\pi_N\bar\pi_N=N\).
\end{proof}

\begin{definition}\label{def:columns}
Define \(K_c\subset K^0\) as follows. At \(g_0,g_1\) require all the
chosen matrices to reduce to the identity. At \(N\) require the last
row of each chosen matrix to reduce to \((0,\ldots,0,1)\). Elsewhere
use \(K_p^0\). Write \(\mathcal C(K)=\Gamma\backslash G^u/K\) for
any subgroup \(K\) of \(K^0\). At the infinite places \(\tau\in A\)
take the \(a\)-th symmetric power of the conjugate standard
representation in rank \(k\), and at the other places take the trivial
representation. Write \(W\) for their tensor product, with the action
\(\rho\) of \(\Gamma\) given by applying the chosen embedding at each
place to the matrix entries and conjugating. In the symmetric power,
multilinear symmetric products of \(a\) vectors are represented by
homogeneous degree-\(a\) polynomials in formal basis vectors, with
the group acting on each basis vector by its matrix.
The columns form the complex vector space
\[
 {\rm Col}=\{F:G^u\to W:\ F(\gamma g h)=\rho(\gamma)F(g)
                 \text{ for }\gamma\in\Gamma,\ h\in K_c,\ g\in G^u\}.
\]
In the Young diagram convention the rank-\(k\) weight uses heights
\(1,\ldots,k-1\), with determinant exponent zero. Its only nonzero
multiplicity is \(a\) at height one at the places in \(A\).
\end{definition}

\begin{lemma}\label{lem:trace-finite}
There are finitely many double cosets in \(\mathcal C(K^0)\).
For any \(g\in G^u\), the stabilizer \(\Gamma\cap gK^0g^{-1}\)
is finite.
\end{lemma}
\begin{proof}
Let \(L_g=\{x\in E^k:x\in g_pO_p^k\ \text{for all }p\}\).
For some positive integer \(b\), the local modules lie between
\(bO_p^k\) and \(b^{-1}O_p^k\) for every \(p\), by clearing the
finitely many local denominators of \(g_p,g_p^{-1}\). Thus \(L_g\)
is a full \(\Z\)-lattice between \(bO_E^k\) and \(b^{-1}O_E^k\)
and is stable under \(O_E\). Its completion at each \(p\) is precisely
\(g_pO_p^k\). Indeed in the quotient of the two bounding lattices,
multiplication by \(b\) identifies the problem with imposing subgroups
of \((\Z/b^2\Z)^{2fk}\) independently at its prime divisors, which is
possible by Lemma~\ref{lem:basic-arith}. Thus two such lattices are
mapped onto one another by \(\Gamma\) exactly when their tuples are
in the same double coset.

On \(L_g\) put \(B(x,y)=\operatorname{Tr}_{E/\Q}h(x,y)\).
This is symmetric positive definite, since
\(B(x,x)=2\sum_{\tau\in\Phi,i}|\tau(x_i)|^2\).
It is integral: over every \(\Z_p\), the local isometry \(g_p^{-1}\)
identifies the pairing with the trace pairing on \(O_p^k\), and a
rational number integral at every prime is an integer by
Lemma~\ref{lem:basic-arith}. The Gram determinant in an integral
basis of \(L_g\) is independent of \(g\). It is a positive integer
whose exponent at every \(p\) equals that for \(L_1\), since the
two local Gram matrices are related by a basis change with determinant
in \(\Z_p^\times\).

For completeness, positive definite integral lattices of fixed rank
\(n\) and fixed Gram determinant have only finitely many Gram matrices
up to integral basis change. In Euclidean coordinates their covolume
\(V\) is the square root of that determinant and the shortest nonzero
length is at least one. More generally a lattice of rank \(n\) and
covolume at most \(V\) with shortest length at least \(c>0\) has a
basis of lengths bounded in terms of \(n,V,c\). A cube of volume
greater than \(V\) contains two points differing by a nonzero lattice
vector. To see this, sum its translated indicators over any fundamental
parallelotope; the integral is the cube's volume by disjoint tiling,
so at some point the sum exceeds one. A cube of side \(2V^{1/n}\)
therefore bounds the shortest length from above by \(2\sqrt n V^{1/n}\).
A shortest vector \(u\) exists by discreteness and boundedness of a
ball in basis coordinates. Its coordinates have gcd one and thus
integer elementary row operations extend it to a lattice basis.
Projection along \(u\) onto \(u^\perp\) gives a lattice with covolume
at most \(V/c\), by the determinant in that basis. Any nonzero
projected vector lifts, by subtracting an integer multiple of \(u\),
to one with parallel component of length at most \(\|u\|/2\).
The original shortest-length property then gives projected length at
least \(\sqrt3\|u\|/2\). Induction on rank gives a bounded basis there,
whose reduced lifts together with \(u\) give a bounded basis here.
The rank-one assertion follows from the covolume directly.
Integral Gram matrices in a bounded basis have only finitely many
possible entries by Cauchy--Schwarz, which follows by orthogonal
projection to a line.

Fix one trace lattice from each isometry class. Under an identification
with it, multiplication by each member \(z\) of a fixed integral
basis of \(O_E\) has norm at most \(\max_\tau|\tau(z)|\), by the
displayed expression for \(B(x,x)\). Each basis vector of the
trace lattice then has finitely many possible images. Thus there
are finitely many possible \(O_E\)-actions on these representatives.
A trace isometry respecting the action extends rationally to an
\(E\)-linear isometry of \(h\): varying either argument by a field
element detects equality from all traces, since for \(y\ne0\) the
trace of \(y^{-1}y\) is \(2f\). It thus belongs to \(\Gamma\), proving
finiteness of the double cosets. A stabilizer is finite because an
isometry of a positive definite lattice has finitely many choices
for the image of each basis vector.
\end{proof}

\begin{proposition}\label{prop:column-dim}
The column level has index
\[
 I_c=[K^0:K_c]=(N^k-1)^f
          \prod_{i=0}^1\left(\prod_{j=0}^{k-1}(g_i^k-g_i^j)\right)^f .
\]
All stabilizers at the column level are trivial. For representatives
\(g\) of \(\mathcal C(K^0)\) put
\(\mathfrak m_k=\sum_{[g]}|\Gamma\cap gK^0g^{-1}|^{-1}\).
Evaluation at representatives of \(\mathcal C(K_c)\) constructs
all the columns, and
\[
 \dim_{\C}{\rm Col}=\mathfrak m_k I_c \binom{a+k-1}{k-1}^{U}>0.
\]
\end{proposition}
\begin{proof}
Reduction of \(\mathrm{GL}_k(\Z_p)\) onto
\(\mathrm{GL}_k(\mathbf F_p)\) is onto, since any lifts of the entries
of an invertible residue matrix give a unit determinant.
Over \(\mathbf F_p\) successive basis columns have
\(p^k-p^j\) choices after \(j\) independent columns. The stabilizer
of a specified nonzero row has index \(p^k-1\), as the group acts
transitively on those rows by basis extension. Lemma~\ref{lem:split-level}
gives the index by multiplication. The primes \(N,g_0,g_1\) are distinct
because \(\log N\ge m\ell>11\ell\) eventually.

The congruence kernel \(I+pM_k(\Z_p)\) is torsion-free for an odd
prime \(p\). For \(A=I+p^aB\ne I,\ a\ge1\), with a unit entry in
\(B\), raising to a positive power prime to \(p\) preserves the least
entry valuation \(a\) of \(A-I\). Raising to \(p\) raises it by exactly
one: modulo \(p^{a+2}\) the result is \(I+p^{a+1}B\), by the
binomial expansion and \(p\ge3\). Factoring a putative order into
its \(p\)-part and its prime-to-\(p\) part contradicts these two
claims. Now \(\Gamma\cap gK_c g^{-1}\) is finite by
Lemma~\ref{lem:trace-finite}. At \(g_0\) its conjugate lies in a
product of such kernels. It must be trivial there, hence also over
\(E\), since base extension to \(E_{g_0}\) is injective. The same
argument works for any level contained in \(K^0\) imposing the
identity just at the chosen components over \(g_0\).

Above \([g]\in\mathcal C(K^0)\) the stabilizer of \(gK^0\) acts on
the \(I_c\) finer cosets. It acts freely by the preceding argument,
so the number of orbits is \(I_c/|\Gamma\cap gK^0g^{-1}|\).
Summation gives \(|\mathcal C(K_c)|=\mathfrak m_k I_c\).
For a representative \(b_c\) of each finer class, choose any
\(w_c\in W\). The rule \(F(\gamma b_c h)=\rho(\gamma)w_c\) is
well defined by triviality of the stabilizer, and recovers every
column. A monomial basis of the \(a\)-th symmetric power is indexed
by nonnegative \(k\)-tuples summing to \(a\). Inserting \(k-1\)
separators among \(a\) identical letters gives \(\binom{a+k-1}{k-1}\)
possibilities. Taking the tensor product over the \(U\) places proves
the formula.
\end{proof}

\begin{theorem}[Ramified columns]\label{thm:ram-dim}
For every one of the \(f\) chosen components over \(N\), lift a generator
of \(\mathbf F_N^\times\) to a scalar unit there, place the inverse
scalar at its partner and take the identity everywhere else. Denote
the resulting central element by \(z_w\), and by \(T_w\) its right
translation on \({\rm Col}\). Put
\[
 L=\prod_w(1-T_w),\qquad {\rm Col}_{\rm ram}=L({\rm Col}),\qquad
 c_0=\dim_{\C}{\rm Col}.
\]
Then
\[
 \dim_{\C}{\rm Col}_{\rm ram}=c_0\left(1-\frac1{N-1}\right)^f.
\]
For any subspace \(V\) of the columns,
\(\dim L^{-1}V\le\dim V+c_0f/(N-1)\).
\end{theorem}
\begin{proof}
Centrality ensures the translations preserve the level. Their generated
subgroup modulo \(K_c\) is \(C=(\mathbf F_N^\times)^f\): the
condition on the last row forces a scalar to be \(1\bmod N\) in
each component. This group acts freely on \(\mathcal C(K_c)\).
Indeed let \(K'\) be generated by \(K_c\) and all \(z_w\). It still
imposes identity over \(g_0\). If \(gz=\gamma g h\) with \(h\in K_c\)
and \(z\) in the scalar subgroup, then
\(\gamma\in\Gamma\cap gK'g^{-1}=\{1\}\) by
Proposition~\ref{prop:column-dim}, hence \(z\in K_c\).

Choose a class representative at the start of each orbit and take its
translates as representatives along that orbit. Changing scalar lifts
by an element of \(K_c\) does not alter the values. Thus on each orbit
the action in the evaluation description of
Proposition~\ref{prop:column-dim} is ordinary translation on
\(W\)-valued functions on \(C\). For a cyclic group of order
\(h=N-1\), the difference between a function and its shift has image
the subspace of coordinate sum zero. Its kernel consists of constants,
and the sum-zero condition has codimension one, proving the equality
of the two spaces by dimension. Functions on the product are tensor
products of the single-coordinate function spaces, as seen from their
delta bases. The product of the \(f\) differences has image their
tensor product of images, hence rank \((h-1)^f\) per scalar coordinate
of \(W\) per orbit. This gives the equality. For the final inequality,
\(\dim L^{-1}V=\dim\ker L+\dim(V\cap\operatorname{im}L)\), by
extending a basis of the kernel by lifts of an image basis.
The elementary bound \(1-(1-u)^f\le fu\) for \(u\in[0,1]\)
follows by induction.
\end{proof}

\section{Integral columns of bounded size}

\begin{lemma}\label{lem:dual-den}
Write \(D_*=dN_1\), let \(L_g\) and \(B\) be as in
Lemma~\ref{lem:trace-finite}, and put
\(L_g^\#=\{y:B(y,L_g)\subset\Z\}\) in the real span of \(L_g\).
Then \(D_*L_g^\#\subset L_g\), and every nonzero vector of
\(L_g^\#\) has \(B\)-length at least \(1/D_*\).
\end{lemma}
\begin{proof}
We first show \(D_*\{y\in E:\operatorname{Tr}(y\mathcal B)\subset
\Z\}\subset\mathcal B\) for the order of Lemma~\ref{lem:rd}.
Write \(\zeta=\zeta_{N_1}\), and let \(y'\) denote the conjugate
under the automorphism fixing \(\zeta\) and exchanging \(\omega,\bar\omega\).
It exists because the \(2f\) embeddings in that lemma realize both choices
for each primitive root and have image \(E\). In the two sums
\[
 u_0=\sum_{j=0}^{N_1-1}\operatorname{Tr}(y\zeta^{-j})\zeta^j,\qquad
 u_1=\sum_{j=0}^{N_1-1}\operatorname{Tr}(y\omega\zeta^{-j})\zeta^j
\]
all coefficients are integers. Upon expansion of the trace the
finite geometric sum vanishes except at embeddings fixing \(\zeta\),
where it equals \(N_1\). Thus \(u_0=N_1(y+y')\),
\(u_1=N_1(\omega y+\bar\omega y')\). In particular
\(N_1(\omega-\bar\omega)y=u_1-\bar\omega u_0\in\mathcal B\).
Multiplication by \(-(\omega-\bar\omega)\) gives \(D_*y\in\mathcal B\).
The trace dual of \(O_E\) is contained in that of \(\mathcal B\), and
\(\mathcal B\subset O_E\), so the same inclusion holds for \(O_E\).

In a lattice basis its trace dual has coordinate matrix the inverse
of the integral Gram matrix. Consequently the dual commutes with
extension of scalars to \(\Z_p\), and its elements in the real span
have rational coordinates. Conjugation preserves \(O_E\), so the
inclusion just proved holds coordinate by coordinate for \(O_p^k\)
with the hermitian trace pairing. The \(p\)-component \(g_p\) is an
isometry; by Lemma~\ref{lem:trace-finite} the inclusion therefore
holds for each completion of \(L_g\). Intersecting over \(p\) gives
the first claim. A nonzero element of \(L_g\) has positive integral
squared length, giving the second.
\end{proof}

\begin{lemma}[Short vectors spanning a lattice over the reals]\label{lem:gaussian-cover}
Let \(L\) be a full lattice in Euclidean \(n\)-space, \(n\ge1\), and
assume every nonzero vector of its dual lattice \(L^\#\) has length at
least \(c>0\). Then each point of the space is within \(9n/c\) of
\(L\), and the vectors of \(L\) of length at most \(27n/c\) span the
space over \(\R\).
\end{lemma}
\begin{proof}
For \(s>0\), covolume \(V\), and \(x\in\R^n\) the formula needed is
\[
 Z_s(x)=\sum_{y\in L}\e^{-\pi\|y-x\|^2/s^2}
   =\frac{s^n}{V}\sum_{w\in L^\#}\e^{-\pi s^2\|w\|^2}
                                    \exp(2\pi i\langle w,x\rangle).
\]
To prove it, periodize the Gaussian on a fundamental parallelotope.
Absolute uniform convergence there follows by writing the lattice in
a basis, in which Euclidean lengths bound a positive multiple of
coordinate lengths. Its coefficient at \(w\in L^\#\) is the Gaussian
integral over the whole space divided by \(V\); its value is the
product of the one-dimensional transforms in Lemma~\ref{lem:theta}.
The dual basis parametrizes all Fourier indices. The Fourier series
is absolutely convergent and recovers the periodization: convolving
with the product of the one-dimensional Fej\'er kernels in lattice
basis coordinates proves uniqueness by the same concentration
argument as in Lemma~\ref{lem:theta}.

Disjoint open balls of radius \(c/2\) around dual lattice points
within radius \(R\) fit inside the ball of radius \(R+c/2\);
scaling volume bounds their number by \((1+2R/c)^n\).
Take \(s=3\sqrt n/c\). For \(j\ge1\) the shell
\(jc\le\|w\|<(j+1)c\) contributes at most
\((2j+3)^n\exp(-9\pi n j^2)\) to the dual sum in absolute value.
Since \(2j+3\le5^j\), the sum of these bounds is dominated by
\(\sum_{j\ge1}(5\e^{-9\pi})^j<1/2\). Thus
\(s^n/(2V)\le Z_s(x)\le 2s^n/V\), and the same bounds with the
corresponding power hold for any larger \(s\).
If every \(y\) were at distance greater than \(3s\sqrt n\) from \(x\),
we would have
\[
 Z_s(x)\le \e^{-9\pi n/2} Z_{\sqrt2s}(x)
       \le \e^{-9\pi n/2}2^{1+n/2}s^n/V<s^n/(2V).
\]
This proves the distance bound. If all vectors of length at most
\(3R\), \(R=3s\sqrt n\), were in a proper subspace, take a point of
length \(2R\) perpendicular to it. A lattice point within \(R\) of
that point has length at most \(3R\) and lies outside the subspace,
a contradiction.
\end{proof}

\begin{definition}\label{def:weyl-coord}
Use coordinates for the symmetric power of Definition~\ref{def:columns}
in which the pure power at \(z\in\C^k\) has evaluation vector
\(b(z)=(z^\alpha)_{|\alpha|=a}\), with \(z^\alpha=\prod_i z_i^{\alpha_i}\).
Concretely these are the coordinates relative to the formal monomials
scaled by the multinomials \(\binom a\alpha=a!/\prod_i\alpha_i!\).
Give coordinate vectors the hermitian squared norm
\(\sum_\alpha\binom a\alpha |u_\alpha|^2\), and \(W\) the tensor
product norm. For a matrix \(A_1\) write \(S(A_1)\) for the matrix
with coefficients determined by \(b(A_1 z)=S(A_1)b(z)\), expanded
as polynomials. This definition works over every commutative ring,
and polynomial substitution gives \(S(A_1 A_2)=S(A_1)S(A_2)\).

For each \(\tau\in A\) let \(\theta_\tau:E\to E\) be
\(z\mapsto\tau(\bar z)\), extended \(\Q_p\)-linearly to \(E_p\).
This uses the unique embedding at the place fixing \(K\) as in
Proposition~\ref{prop:places}. On finite matrices write
\(\rho_p(u)=\bigotimes_{\tau\in A}S(\theta_\tau(u))\).
A column with all values in the \(E\)-coordinate space is called
integral if for every \(g\in G^u\) and every prime \(p\) the coordinates
of \(\rho_p(g_p^{-1})F(g)\) belong to \(O_p\). Each such condition
gives integral coordinates under every embedding at \(p\), because
multiplication by an element of \(O_p\) in its integral basis gives
a monic equation over \(\Z_p\). Write \(\sigma F(g)\) for \(\sigma\)
applied to the individual coordinates of \(F(g)\).
\end{definition}

\begin{proposition}\label{prop:short-col}
Let \(n=2fk,\ C_0=27nD_*\). The columns are spanned over \(\C\) by
integral columns whose values satisfy, for every embedding
\(\sigma:E\to\C\) and every \(g\in G^u\),
\[
                 \|\sigma F(g)\|^2\le (kaC_0)^{2aU}.
\]
\end{proposition}
\begin{proof}
By Lemmas~\ref{lem:dual-den} and~\ref{lem:gaussian-cover}, for each
representative \(g\) of a column class there are \(k\) vectors
\(z_1,\ldots,z_k\in L_g\) linearly independent over \(E\), each of
trace length at most \(C_0\). Indeed all vectors with this length
bound span the real vector space, and a proper \(E\)-subspace has
real dimension at most \(2f(k-1)<n\). For every \(\tau\in A\),
independently choose \(z_\tau=\sum_{j=1}^k t_{\tau j}z_j\) with
\(t_{\tau j}\in\{0,\ldots,a\}\). Prescribe at \(g\) the value
\[
                     w=\bigotimes_{\tau\in A}b(\theta_\tau(z_\tau)).
\]
Here the automorphism acts coordinatewise on each vector. Extend
to a column supported on its double coset by
Proposition~\ref{prop:column-dim}. The values belong to \(E\).
At \(g\), the moved coordinates are products of monomials in
\(\theta_\tau(g_p^{-1}z_\tau)\). These belong to \(O_p\) since
\(z_\tau\in L_g\), and \(\theta_\tau\) preserves \(O_p\).
At \(\gamma g h\) with \(h\in K_c\), the rational transformation
\(\rho(\gamma)\) cancels, and \(\rho_p(h_p^{-1})\) preserves the
integral coordinates. This proves integrality everywhere.

By the multinomial expansion,
\(\langle b(z),b(y)\rangle=(\sum_i\bar z_i y_i)^a\).
Thus \(b(z)\) has squared norm \(\|z\|^{2a}\), and unitary
matrices preserve the form on their span. Each embedding of \(E\)
is the restriction of a power automorphism of \(\Q(\zeta_{dN_1})\):
to choose its action on \(K\), choose a square or nonsquare exponent
modulo \(d\), since such exponents act on the quadratic Gauss sum
by their Legendre symbol (Lemma~\ref{lem:gauss}); independently
choose its action modulo \(N_1\). In particular these embeddings
commute with complex conjugation. Thus every matrix
\(\sigma\theta_\tau(\gamma)\) for \(\gamma\in\Gamma\) is unitary.
Each \(\sigma\theta_\tau(z_\tau)\) has Euclidean length at most
\(ka C_0\), since the squared trace length bounds the squared length
at every embedding. This gives the asserted bound at \(g\) and at
every rational translate.

For a single tensor factor, a linear functional vanishing on all these
vectors gives a homogeneous polynomial of degree \(a\) which, after
an invertible linear substitution with the vectors \(\theta_\tau(z_j)\)
as columns, vanishes on \(\{0,\ldots,a\}^k\). A polynomial of degree
at most \(a\) in each variable vanishing on this grid is zero, by
successive use of the univariate root bound. Its monomial coefficients
therefore vanish. This proves spanning in each factor; independent
choices span their tensor product by expansion of elementary tensors.
Varying the class gives the entire column space by
Proposition~\ref{prop:column-dim}.
\end{proof}

\begin{corollary}\label{cor:short-ram}
Eventually \({\rm Col}_{\rm ram}\) is spanned over \(\C\) by integral
forms \(F\) with
\(\|\sigma F(g)\|^2\le\exp(fa\ell^3)\) for every \(g,\sigma\).
\end{corollary}
\begin{proof}
Apply \(L\) of Theorem~\ref{thm:ram-dim} to the spanning family of
Proposition~\ref{prop:short-col}. Right translation by \(z_w\)
preserves integrality, since
\(\rho_p(g_p^{-1})F(gz_w)=\rho_p((z_w)_p)
 \rho_p((g_p(z_w)_p)^{-1})F(gz_w)\)
and \((z_w)_p\) has integral entries and inverse. The difference of
two integral columns is integral. In squared norm the triangle
inequality costs at most a factor four for each difference. The
logarithm of the resulting bound is at most
\(f\log4+2aU\log(kaC_0)\).
By Proposition~\ref{prop:thick},
\(k=O(x^5)\), \(N_1,f=O_M(x^{12})\), so
\(\log(kC_0)=O_{M,d}(1+\log x)\).
By Proposition~\ref{prop:level-prime}, \(\log a\le2\ell\)
eventually. Because \(\log x=o(\ell)\) and \(U\le f\), the
asserted bound follows. These coordinates are algebraic; an arbitrary
field automorphism of \(\C\) acts on them through its restriction to
\(E\), so the same estimate holds for all such automorphisms.
\end{proof}

\section{Moments and ranks of projected translates}

\begin{definition}\label{def:row-space}
Put \(\nu=m/2-1\). The rational hermitian space \(Y_0\) has basis
\[
                  e_1,\ldots,e_\nu,\ x_1,\ldots,x_k,\ z,\
                                      e_1^*,\ldots,e_\nu^* .
\]
Here \(x_i\) denotes a basis vector. Its nonzero pairings are
\(h(e_i,e_j^*)=h(e_j^*,e_i)=\delta_{ij}\),
\(h(x_i,x_j)=\delta_{ij}\) and \(h(z,z)=-N\).
Use the integral lattice with this basis except that \(x_k\) is
replaced by \(v=x_k/\bar\pi_N+z/N\). At \(g_0,g_1\) the finite
level is the kernel of reduction on the chosen components of the
stabilizer of this lattice. The corresponding finite group is written
\(\mathcal G\).
\end{definition}

\begin{proposition}\label{prop:row-group}
The hermitian space of Definition~\ref{def:row-space} has rank \(r\)
and signature \((k+\nu,\nu+1)\) at each archimedean place. Its specified
integral lattice has a Gram matrix of determinant a unit.
The level quotient is
\(\mathcal G\simeq\prod_{i=0}^1\mathrm{GL}_r(\mathbf F_{g_i})^f\);
in particular
\[
              |\mathcal G|=\prod_{i=0}^1\prod_{j=0}^{r-1}(g_i^r-g_i^j)^f .
\]
\end{proposition}
\begin{proof}
Each hyperbolic plane has orthogonal eigenvectors \(e_i+e_i^*\)
and \(e_i-e_i^*\) of opposite squared lengths. The middle block has
\(k\) positive directions and one negative direction.
Since \(k+1+2\nu=r\), this gives the signature. The Gram matrix on
\(v,z\) is
\(\left(\begin{smallmatrix}0&-1\\-1&-N\end{smallmatrix}\right)\)
because \(\pi_N\bar\pi_N=N\); all other pairings are unchanged.
Its determinant and the determinants of the other blocks are units.
At each split pair the matrix on one component may be chosen
arbitrarily and the hermitian identity determines the other by
inverse transpose relative to the Gram matrix. Unimodularity
makes this identification valid on integral stabilizers.
Lemma~\ref{lem:split-level} and the reduction and counting arguments
of Proposition~\ref{prop:column-dim} give the formula.
\end{proof}

\begin{lemma}\label{lem:proj-moments}
Let a finite group \(G\) act unitarily by \(\rho\) on a finite-dimensional
complex Hilbert space \(\mathcal H\). Let \(Q:\C^c\to\mathcal H\),
\(c>0\), let \(R\) be an orthogonal projection commuting with
\(\rho\), and set \(W=R\mathcal H,\ d_0=\dim W\).
With standard basis \(e_i\) of \(\C^c\) and Hilbert--Schmidt norm
\(\|T\|_{\rm HS}^2=\sum_i\|Te_i\|^2\) in an orthonormal domain basis, put
\[
 \begin{split}
 E_1&=\sum_i\|(Q^*Q-I)e_i\|+\|(1-R)Q\|_{\rm HS}^2,\\
 E_2&=|G|^{-1}\sum_{s\in G}\|Q^*\rho(s)Q\|_{\rm HS}^2 .
 \end{split}
\]
If \(E_1<c\), then \(d_0>0\). If \(d_0>0,\ E_1\le\alpha c\) and
\(E_2\le(1+\lambda)c^2/d_0\), for \(\alpha,\lambda\ge0\), then for
each subspace \(F\subset W\) there is \(s\in G\) with
\[
 \dim(R\rho(s)Q)^{-1}(F)
 \le c\left(\frac{1+\dim F/d_0}{2}+\alpha+\frac{\lambda}{2}\right).
\]
\end{lemma}
\begin{proof}
Write \(Q_s=\rho(s)Q\), \(P_A\) for orthogonal projection to
a subspace \(A\), and \(D_s=(RQ_s)^{-1}(F)\).
The trace of a product is unchanged on cyclic permutation by summing
its matrix entries with indices reversed. In particular
\[
 \big|\operatorname{tr}\big(P_{D_s}(I-Q^*Q)\big)\big|
 \le\sum_i\|(I-Q^*Q)e_i\|,
\]
by expanding the trace in \(e_i\), since each \(\|P_{D_s}e_i\|\le1\).
In an orthonormal basis beginning with a basis of \(D_s\), the trace
of \(P_{D_s}Q^*Q\) is the sum of \(\|Q_s u\|^2\) over those first
vectors. For each of them \(RQ_su\in F\); hence
\(\|Q_su\|^2=\|P_FQ_su\|^2+\|(1-R)Q_su\|^2\).
Enlarging the two sums to the whole basis bounds \(\dim D_s\) by
\(E_1+\|P_F Q_s\|_{\rm HS}^2\), using commutation of \(R\) and \(\rho\).
If \(R=0\), take \(F=0\); this yields \(c\le E_1\), proving the
first claim. Independence of the Hilbert--Schmidt norm from the basis
follows from its expression as \(\operatorname{tr}(T^*T)\).

Let \(A_0=|G|^{-1}\sum_s Q_sQ_s^*\).
Cyclicity and counting the \(|G|\) pairs with \(s^{-1}t=g\)
give \(\operatorname{tr}(A_0^2)=E_2\).
The average of \(\|P_FQ_s\|_{\rm HS}^2\) is
\(\operatorname{tr}(P_F A_0)\le\sqrt{\dim F}\sqrt{E_2}\) by
Cauchy--Schwarz on the matrix entries. Some \(s\) is at most this
average. With \(u=\dim F/d_0\), the inequality
\(\sqrt{(1+\lambda)u}\le(1+\lambda+u)/2\) gives the result.
\end{proof}

\begin{theorem}[Selection of translates]\label{thm:greedy}
Let \(V,W\) be complex vector spaces of positive finite dimensions
\(c,d_0\). Let \(A_s:V\to W\) be a family of linear maps.
Let \(0<\eta\le1/8\) and suppose that, for each subspace \(F\subset W\),
some \(s\) satisfies
\[
 \dim A_s^{-1}F\le c\left(\frac{1+\dim F/d_0}{2}+\eta\right).
\]
There is a finite set \(T\) of indices such that
\[
 \dim\sum_{s\in T} A_s(V)\ge(1-4\eta)d_0,\qquad
                       |T|c\le3d_0\log(1+1/\eta).
\]
\end{theorem}
\begin{proof}
From \(F=0\) we get \(c(1/2-\eta)\le d_0\) by rank and nullity.
Start with \(F=0\). If its codimension \(k_0\) exceeds \(4\eta d_0\),
choose \(s\) as in the hypothesis and add \(A_s(V)\). The dimension
gain is the rank of \(V\to W/F\), so for the new codimension \(k_1\)
\[
 k_0-k_1\ge c(k_0/(2d_0)-\eta),\qquad
 k_1-2\eta d_0\le(1-c/(2d_0))(k_0-2\eta d_0).
\]
Every such step adds a new index since its gain is positive.
If \(c\ge2d_0\), the first step gives \(k_1\le2\eta d_0\) and
costs \(c\le 8d_0/3\), which satisfies the asserted bound.
If \(c<2d_0\), after \(j\) steps made before stopping the excess
codimension is at most
\(d_0(1-2\eta)\exp(-jc/(2d_0))\) by \(1-u\le\e^{-u}\).
The process therefore stops by
\(j=\lceil(2d_0/c)\log((1-2\eta)/(2\eta))\rceil\).
Its cost is at most
\(2d_0\log(1/(2\eta))+2d_0\).
For \(0<\eta\le1/8\) this is at most \(3d_0\log(1+1/\eta)\):
use \(\log(1/\eta)\ge\log8>1\) and \(2-2\log2<1\).
The target codimension is at most \(4\eta d_0\).
\end{proof}

\begin{definition}\label{def:det-index}
For a commutative integral domain \(O\), two \(O\)-submodules \(M,P\)
in the same module, and \(a\in O\), write \(\mathcal I(M,P;a)\)
if every finite list \(y_1,\ldots,y_n\in M\) admits a matrix
\(U=(u_{ij})\in M_n(O)\) with \(\det U\mid a\) in \(O\) and
\(\sum_j u_{ij}y_j\in P\) for every \(i\). An empty list has
determinant one.
\end{definition}

\begin{lemma}\label{lem:det-index}
The relations \(\mathcal I(M,P;a)\) and \(\mathcal I(P,T;b)\)
imply \(\mathcal I(M,T;ab)\).
If \(f_j:M\to Y\) are \(O\)-linear maps, \(P\subset Y\) is a submodule
and \(\mathcal I(M,f_j^{-1}P;a_j)\) holds for \(1\le j\le h\), then
\[
 \mathcal I\left(P+\sum_{j=1}^h f_j(M),P;\prod_{j=1}^h a_j\right).
\]
For an absolute value \(|\cdot|\) with \(|u|\le1\) on \(O\),
each of the resulting matrices has \(|\det U|\ge|\prod_j a_j|\)
in the second assertion.
\end{lemma}
\begin{proof}
Apply the second recombination matrix to the vectors furnished by the
first. The product of the determinants divides \(ab\).
For the sum of images, write elements of the enlarged module as
\(y_i=y_i'+f_h(x_i)\) with \(y_i'\) in the preceding sum
including \(P\). Recombine all \(x_i\) with one matrix for \(f_h\).
The recombined \(f_h(x_i)\) belong to \(P\), and the recombined
\(y_i'\) still lie in the preceding module. Apply induction and the
product argument. The initial case uses the identity matrix. If
\(\prod_j a_j=t\det U\) with \(t\in O\), multiplicativity gives
the last bound.
\end{proof}

\begin{proposition}\label{prop:selected-family}
Use the arithmetic data above and let \(Q:{\rm Col}\to\mathcal H\),
\(\rho:\mathcal G\to U(\mathcal H)\), \(R,W,d_0,E_1,E_2\) satisfy
Lemma~\ref{lem:proj-moments}, with a chosen Hilbert structure on
\({\rm Col}\), \(c=c_0\). Suppose \(E_1\le\alpha c_0<c_0,\ 
E_2\le(1+\lambda)c_0^2/d_0\) and
\[
                 \alpha+\lambda/2+f/(N-1)\le b^{-2}.
\]
On \({\rm Col}_{\rm ram}\) let \(B_s=c'R\rho(s)Q\) for a common
nonzero scalar \(c'\). Then some set \(T\subset\mathcal G\) with
\(|T|c_0\le3d_0\log(1+b^2)\) supplies a basis of
\(\sum_{s\in T}B_s({\rm Col}_{\rm ram})\), each basis vector of the
form \(B_s z\) with \(z\) from the short integral family of
Corollary~\ref{cor:short-ram}. Its length \(n\) satisfies
\((1-4/b^2)d_0\le n\le d_0\).
\end{proposition}
\begin{proof}
For each \(F\subset W\), choose \(s\) as in
Lemma~\ref{lem:proj-moments}. Multiplication by \(c'\ne0\) preserves
\(F\). Pull the inverse image back through the central difference
operator \(L\); by Theorem~\ref{thm:ram-dim} this adds at most
\(c_0 f/(N-1)\) to its dimension. The hypotheses of
Theorem~\ref{thm:greedy} with \(\eta=b^{-2}\) therefore hold for
the maps \(B_s L\); \(b^{-2}\le1/8\) eventually.
Their images equal the images on the ramified columns. The indicated
short vectors span those columns, so the union of their images spans
the selected sum. Repeatedly choose a vector outside the span of those
already chosen until the span is the whole sum; this takes exactly
its dimension in steps and produces the stated basis.
\end{proof}

\begin{corollary}\label{cor:26}
In Proposition~\ref{prop:38}, retain \eqref{eq:defect} and the
second inequality of \eqref{eq:volume}, and assume
\(\delta\le3r\log(1+b^2)B\). Then \(T\le26fa(\log b)^4\).
If each of the selected translates in Proposition~\ref{prop:selected-family}
costs at most \(c_0B_0'\) in logarithmic determinant index, the total
additive cost is at most \(3d_0\log(1+b^2)B_0'\).
\end{corollary}
\begin{proof}
The proof of Proposition~\ref{prop:38} through Lemma~\ref{lem:rank}
now gives \(T\le4E+4x^{99}+6\log(1+b^2)B\).
Using \(x^{99}\le E,\ \log(1+b^2)\le3\log b\) gives the asserted
constant \(4+4+18=26\). The second assertion follows by multiplication
of the translate count by \(B_0'\); for local determinants the
multiplication of divisors that gives additivity after taking
logarithms is Lemma~\ref{lem:det-index}. With \(T=m\mathcal R\),
the new numerical upper bound implies the one required in
Theorem~\ref{thm:geometric-reduction}.
\end{proof}

\section{Heights of families indexed by finite places}

\begin{definition}\label{def:field-conventions}
We use set theory with the axiom of choice in its maximal-element
form: a partially ordered set in which each chain has an upper bound
has a maximal element if it is nonempty. For a prime \(p\) the ring
\(\Z_p\) consists of compatible residue classes modulo \(p^n\) for
\(n\ge1\). Each element has a unique expansion
\(\sum_{j\ge0}a_jp^j\) in these residues, with \(0\le a_j<p\).
Indeed each lift of a residue modulo \(p^j\) specifies precisely its
next digit. A nonzero element is \(p^k\) times an element with nonzero
constant digit, where \(k\) is its first nonzero digit position.
An element with nonzero constant digit has compatible inverses in the
finite residue rings by Lemma~\ref{lem:basic-arith}. It follows by
multiplication that \(\Z_p\) is a domain with unique maximal ideal
\(p\Z_p\). Its fraction field is \(\Q_p\). On that field the
valuation \(v_p\) sends \(p^ku\), \(u\in\Z_p^\times\), to \(k\),
and \(|x|_p=p^{-v_p(x)}\) for \(x\ne0\), with \(|0|_p=0\).
\end{definition}

\begin{lemma}\label{lem:place-fields}
Each \(\Q_p\) has an algebraic closure \(K_p\) with a multiplicative
absolute value extending \(|\cdot|_p\) and satisfying
\(|x+y|_p\le\max(|x|_p,|y|_p)\). There are field isomorphisms
\(\C\to K_p\). Every embedding of a number field \(F\subset\C\)
into \(K_p\), or into \(\C\), extends to such an isomorphism onto the
specified field.
\end{lemma}
\begin{proof}
A field \(K\) admits an algebraic extension splitting all polynomials
in \(K[X]\): well-order the polynomials, adjoin successively roots
of irreducible factors by taking polynomial-ring quotients, and for
each polynomial repeat until it splits. At a limit take the union.
Repeating this construction for countably many rounds and taking
their union gives an algebraic closure, since the coefficients of any
polynomial in the union occur in one round. Algebraicity is transitive
here: a finite list of algebraic elements generates a finite extension
by the power bases, and multiplication of the dimensions in a tower
follows by multiplying the elements of the two bases. A field
isomorphism extends to an isomorphism of algebraic closures.
In fact order its algebraic partial extensions into the target by
inclusion and take a maximal one by Definition~\ref{def:field-conventions}.
An element outside its domain has a minimal polynomial with a root
in the target; evaluation at that root extends the embedding, a
contradiction. The image is now algebraically closed and the target
is algebraic over it, giving surjectivity.

To extend the absolute value, order local subrings of \(K_p\)
containing \(\Z_p\) by domination: the larger contains the smaller
and its maximal ideal contracts to the smaller maximal ideal. Consider
only the subrings dominating \(\Z_p\). A union over a chain is local
with maximal ideal the union of the maximal ideals, since every element
outside this union has its inverse already in a member of the chain.
Write \(S,\mathfrak m\) for a maximal such pair. For \(u\ne0\) outside
\(S\), one has \(\mathfrak m S[u]= S[u]\). Otherwise a maximal ideal
of \(S[u]\) containing \(\mathfrak m S[u]\) exists by the same
maximal-element principle on proper ideals; it contracts to
\(\mathfrak m\) since no unit can belong to it. Inverting all elements
outside that ideal would give a strictly larger dominating local
subring of \(K_p\). Thus \(1=a_0+\cdots+a_nu^n\) with \(a_i\in\mathfrak m\),
\(n\ge1\). Division by \(u^n(1-a_0)\) gives a monic equation for
\(u^{-1}\). Consequently \(S[u^{-1}]\) is a finitely generated
\(S\)-module. Its product with \(\mathfrak m\) is proper: otherwise
a finite generating column \(v\) satisfies \(v=Cv\) with \(C\) a
matrix over \(\mathfrak m\). The determinant of \(I-C\) lies in
\(1+\mathfrak m\), hence is a unit, and multiplying by its adjugate
would give \(v=0\), although the module contains \(1\). Localizing
at a maximal ideal as above and using maximality of \(S\) proves
\(u^{-1}\in S\).

Write the group \(K_p^\times/S^\times\) additively and denote the
class of \(y\) by \(v(y)\). Order classes by \(v(y)\ge v(z)\) if
\(y/z\in S\). The last paragraph gives a total order, and the ring
laws give \(v(y+z)\ge\min(v(y),v(z))\) for nonzero sums.
The restriction to \(\Q_p^\times\) is \(v_p(\cdot)v(p)\), with
\(v(p)>0\). In a polynomial relation for nonzero \(y\) over
\(\Q_p\), the nonzero terms of least value occur at least twice:
division by a unique least term would make \(1\) the negative of a sum
in \(\mathfrak m\). Thus for two different exponents \(i,j\)
one has \((i-j)v(y)=k v(p)\) for some integer \(k\).
Ordered groups have no torsion, so \(v(y)\) is uniquely identified
with a rational multiple of \(v(p)\), respecting sums and order.
Normalizing \(v(p)=1\) and taking \(p^{-v(y)}\) now gives the required
absolute value.

The complex field is algebraically closed: if a nonconstant
polynomial has no zero, its reciprocal is entire and bounded (use
the leading term outside a disk and compactness inside). The
coefficient bound of Lemma~\ref{lem:analytic-tools} on disks of
arbitrarily large radius makes this reciprocal constant.
The cardinality of \(\Z_p\) is the continuum by its digit expansion:
binary sequences inject into its digit sequences, and each digit
can be encoded in a fixed-length binary block. The real numbers have
the same cardinality, by binary expansions and by injecting binary
sequences into ternary expansions with digits \(0,2\).
The field of fractions and the algebraic closure of \(\Z_p\)'s
fraction field have the same cardinality, since a polynomial is
specified by a finite list and has finitely many roots. Here finite
tuples and countable unions of finite tuples over a set of continuum
size can be encoded in binary sequences by interleaving digits and
prefixing the length. Both \(K_p\) and \(\C\) thus have cardinality
continuum. For a number field \(F\), maximal algebraically independent
sets over \(F\) exist by unions of chains. The whole field is
algebraic over the rational functions in such a set by maximality.
If its cardinal is \(\kappa\), the rational functions and their
algebraic elements number at most \(\max(\aleph_0,\kappa)\): they
are specified by finite lists from an infinite alphabet of this
size and finitely many root choices. For completeness, an infinite
well-ordered cardinal \(\lambda\) satisfies \(|\lambda^2|=\lambda\).
Order ordinal pairs by their maximum and then lexicographically
within each block. By induction on infinite cardinals every proper
initial segment ending at a pair of maximum \(\alpha<\lambda\)
has size \(<\lambda\), since it lies in \((\alpha+1)^2\);
for finite \(\alpha\) this is finite. Their order ranks therefore
inject \(\lambda^2\) into \(\lambda\). Finite powers and the countable
union follow. It follows that the independent sets on the two sides
both have cardinality continuum. Match them by a bijection, extend
the given embedding first to the rational function fields and then
to their algebraic closures by the first paragraph.
\end{proof}

\begin{definition}\label{def:place-height}
For each prime fix one valued field \(K_p\) as in
Lemma~\ref{lem:place-fields}. Let \(\mathcal P\) be the set of pairs
\(\iota=(p,j)\) with \(j:\C\to K_p\) a field isomorphism.
For a family \(y=(y_\iota)_{\iota\in\mathcal P}\) of complex numbers
and \(C\in\R\), write \({\rm ht}(y)\le C\) to mean that there exist
\(h\in\C^\times,\ A\in\R\), and reals \(e_p\) such that
\[
 \begin{aligned}
 |\sigma h|&\le\e^A &&(\sigma\in{\rm Aut}(\C)),\\
 |j(hy_\iota)|_p&\le\e^{e_p} &&(\iota=(p,j)\in\mathcal P),\\
 A+\sum_{p<M}e_p&\le C &&\text{for all sufficiently large integers }M .
 \end{aligned}
\]
The last sum runs over primes. This notation denotes the stated
relation even when an infimum of possible bounds is unattained.
\end{definition}

\begin{proposition}\label{prop:height-product}
Height bounds in Definition~\ref{def:place-height} add under products,
and multiply by \(n\) under a nonnegative integral power \(n\).
If \({\rm ht}(y)\le C,\ C\le C'\) and
\(|j z_\iota|_p\le |j y_\iota|_p\) for every \(\iota=(p,j)\), then
\({\rm ht}(z)\le C'\). For the constant family \(\mathbf 1\) the
height bound \(C\) holds exactly when \(C\ge0\).
\end{proposition}
\begin{proof}
Multiply witnesses and add archimedean and finite exponents.
Monotonicity uses the original witness. Taking \(h=1,A=e_p=0\)
works for \(\mathbf1\) at cost zero and gives the zeroth power;
induction proves the statement for all powers.

Conversely let \(h,A,e_p\) witness a bound for \(\mathbf1\).
The first inequality implies that \(h\) is algebraic over \(\Q\).
Indeed if it were transcendental one could include it in a maximal
algebraically independent set over \(\Q\). For each positive integer
\(n\), translation \(h\mapsto h+n\) fixing the other basis elements
is an automorphism of the rational function field, with inverse
\(h\mapsto h-n\). The extension argument of
Lemma~\ref{lem:place-fields} gives an automorphism of \(\C\).
The values \(|h+n|\) are unbounded. For algebraic \(h\ne0\) of
degree \(D\), let \(a\in\Q^\times\) be the signed constant coefficient
\((-1)^D P_h(0)\) of its monic minimal polynomial. It is the product
of its \(D\) roots, distinct because a characteristic-zero irreducible
polynomial is coprime to its nonzero derivative. Each root, in either
\(\C\) or \(K_p\), defines an embedding of \(\Q(h)\) that extends by
Lemma~\ref{lem:place-fields}. Multiplying the respective bounds gives
\[
        D^{-1}\log|a|\le A,\qquad D^{-1}\log|a|_p\le e_p .
\]
Once \(M\) exceeds all prime divisors in numerator and denominator of
\(a\), prime factorization in Lemma~\ref{lem:basic-arith} gives
\(\log|a|+\sum_{p<M}\log|a|_p=0\). Summing the inequalities and
taking such an \(M\) in the eventual range gives \(0\le C\).
\end{proof}

\section{The determinant of a completed family of rows}

\begin{lemma}\label{lem:mixed-gram}
Let \(S\) be a positive definite symmetric \(h\)-by-\(h\) matrix
over \(\Q\). For an \(h\)-by-\(n\) complex matrix \(X\) define
\[
   C_\sigma(X)=\det\big((\sigma X)^* S(\sigma X)\big)
   \quad(\sigma\in{\rm Aut}(\C)),
\]
where \(^*\) is conjugate transpose. Write \(\sigma'=\mathfrak c
\sigma\mathfrak c\) for complex conjugation \(\mathfrak c\) on both
sides of \(\sigma\). Then
\[
 \big|\sigma\det(X^*SX)\big|^2\le C_{\sigma'}(X)C_\sigma(X).
\]
If \(X\) is enlarged by \(u\) columns, each of squared \(S\)-norm
at most \(\e^\beta\) after application of \(\sigma\), then
\(C_\sigma\) of the enlarged matrix is at most
\(C_\sigma(X)\e^{u\beta}\).
\end{lemma}
\begin{proof}
For two matrices \(Y,Z\) with \(n\) columns, the inequality
\[
 |\det(Y^*SZ)|^2\le\det(Y^*SY)\det(Z^*SZ)
\]
is immediate if either family of columns is dependent. Otherwise
carry out Gram--Schmidt for the \(S\)-inner product on each family.
This writes \(Y=Y_0B,\ Z=Z_0D\) with \(Y_0,Z_0\) isometries from
\(\C^n\), and \(B,D\) invertible. After canceling \(|\det B\det D|^2\)
the assertion is \(|\det(Y_0^* S Z_0)|\le1\).
Each column has Euclidean length at most one, since its squared
length is that of the \(S\)-orthogonal projection of a unit vector.
Lemma~\ref{lem:hadamard} bounds the determinant of its Gram matrix
by one. Its determinant is the squared absolute determinant of the
matrix, as required. Apply this with \(Y=\sigma'X,\ Z=\sigma X\),
since \(\overline{\sigma' x}=\sigma\bar x\) and \(S\) is rational.
For the second assertion apply the Gram--Schmidt product in
Lemma~\ref{lem:hadamard} in the given order, starting with the old
columns. Each further squared residual length is at most the original
squared length. A dependent family has determinant zero, and remains
dependent under enlargement.
\end{proof}

\begin{definition}\label{def:lift-vol}
For a linear map \(u:W\to H\) of finite-dimensional complex Hilbert
spaces let \(Z(u)=\det(u^*u)\). In an orthonormal basis of \(W\) it
equals the Gram determinant of the images of that basis. It is real
and nonnegative, and is positive exactly when \(u\) is injective, by
Lemma~\ref{lem:hadamard} and its proof. For any ordered basis
\(x=(x_1,\ldots,x_D)\) of \(W\), put \(G(u,x)_{ij}=\langle ux_i,ux_j\rangle\),
conjugate linear in the first argument. If \(C\) has the coordinates
of \(x\) as columns in an orthonormal basis, then
\(G(u,x)=C^*(u^*u)C\) in that basis. In particular
\(\det G(u,x)=|\det C|^2 Z(u)\).
\end{definition}

\begin{theorem}[A determinant bound at all finite places]\label{thm:one-det}
Let \(t\) range over a set with a distinguished member \(0\);
let \(u_t:W\to H_t\) be linear maps of finite-dimensional complex
Hilbert spaces. Assume \(u_0\) injective and fix an ordered basis
\(x\) of \(W\), an integer \(k'\ge0\), and real numbers \(V,\delta,g\).
For each \(\iota=(p,j)\in\mathcal P\) of
Definition~\ref{def:place-height} specify \(t_\iota\) and
\(\gamma_\iota\in\C^\times\).
Assume \({\rm ht}((\bar\gamma_\iota\gamma_\iota)^{-1})\le g\), and
\[
         |\sigma\det G(u_0,x)|\le\e^V\quad(\sigma\in{\rm Aut}(\C)).
\]
Let \(e_p\in\R\) satisfy \(\sum_{p\in J}e_p\le\delta\) for every
finite set \(J\) of primes. For each fixed \(\iota=(p,j)\) and for
each \(\varphi\in\{j,j\mathfrak c\}\) suppose we have a set of vectors
\(\mathcal L_{\iota,\varphi}\subset H_{t_\iota}\) and a square matrix
\(U_\varphi\) of size \(\dim W\) over \(\C\) such that
\[
 \begin{aligned}
  |j\langle v,w\rangle|_p&\le1
       &&(v\in\mathcal L_{\iota,j\mathfrak c},
          \ w\in\mathcal L_{\iota,j}),\\
 |\varphi(\det U_\varphi)|_p
       &\ge\e^{-e_p}|\varphi(\gamma_\iota)|_p^{k'},&
 \sum_l(U_\varphi)_{il}u_{t_\iota}x_l&\in\mathcal L_{\iota,\varphi}
                         \quad(\text{all }i).
 \end{aligned}
\]
Then the squared-volume ratio satisfies
\[
  {\rm ht}\!\left(\frac{Z(u_{t_\iota})}{Z(u_0)}\right)
                              \le V+2\delta+k'g .
\]
\end{theorem}
\begin{proof}
For fixed \(\iota=(p,j)\) write \(G=G(u_{t_\iota},x)\). The matrix
of pairings of the two recombined families is
\(\overline{U_{j\mathfrak c}}\,G\,U_j^{\mathsf T}\). Under \(j\),
all its entries have absolute value at most one. Expanding its
determinant as a finite signed sum gives the same bound for its
determinant by Lemma~\ref{lem:place-fields}. From the two lower
determinant bounds, noting \(j(\overline z)=(j\mathfrak c)(z)\),
one obtains
\[
       \left|j\big(\det G\,(\bar\gamma_\iota\gamma_\iota)^{k'}\big)
       \right|_p\le\e^{2e_p}.
\]
By Definition~\ref{def:lift-vol},
\(h=\det G(u_0,x)\) is nonzero and
\(h\,Z(u_{t_\iota})/Z(u_0)=\det G\). Thus \(h,V,(2e_p)_p\)
witness the height bound \(V+2\delta\) for the volume ratio times
\((\bar\gamma_\iota\gamma_\iota)^{k'}\). Multiply by the \(k'\)-th
power of its scalar correction using Proposition~\ref{prop:height-product}.
\end{proof}

\begin{corollary}\label{cor:completion-height}
Suppose independent rows \(x_1,\ldots,x_n\) in a space \(W\) of
dimension \(d_0\) are given. Let \(u_t,t_\iota,\gamma_\iota\),
the pairing bound, and the exponents \(e_p,\delta,g\) be as in
Theorem~\ref{thm:one-det}. Suppose that, in a coordinate
realization of \(H_0\) with positive rational Gram matrix \(S\),
the columns \(X=(u_0x_i)_{i\le n}\) satisfy \(C_\sigma(X)\le\e^V\)
for every \(\sigma\). For each of the two embeddings at \(\iota\)
assume these \(n\) vectors on applying \(u_{t_\iota}\) can be
recombined into \(\mathcal L_{\iota,\varphi}\) by an \(n\)-row
matrix of \(\varphi\)-determinant of absolute value at least
\(\e^{-e_p}\). Suppose \(W\) has a spanning set of vectors \(z\)
such that each \(u_0z\) has squared norm at most \(\e^\beta\) after
every field automorphism in these coordinates, and
\(\gamma_\iota u_{t_\iota}z\in\mathcal L_{\iota,\varphi}\) for
both required embeddings at every \(\iota\). Then
\[
 {\rm ht}\!\left(\frac{Z(u_{t_\iota})}{Z(u_0)}\right)
           \le V+2\delta+(d_0-n)(\beta+g).
\]
\end{corollary}
\begin{proof}
As long as the span of the current rows is proper, some member of
the spanning set lies outside it. In precisely \(k'=d_0-n\) steps
this gives a basis. Lemma~\ref{lem:mixed-gram} gives the archimedean
determinant bound for the completed basis with exponent \(V+k'\beta\).
Use the block diagonal matrix with first block the recombination
of the original rows and second block \(\gamma_\iota I_{k'}\).
Its determinant satisfies the bound of Theorem~\ref{thm:one-det};
the two diagonal blocks give exactly the stipulated integral vectors.
Apply that theorem.
\end{proof}

\begin{lemma}\label{lem:compressed-gram}
Let \(u:U\to H\) be injective and \(A:T\to H\) linear between
finite-dimensional complex Hilbert spaces. Put \(L=u(U)\) and let
\(P_L\) be the orthogonal projection onto \(L\). For \(s\in\C\)
the endomorphism \(M=sP_L A A^*|_L\) satisfies
\[
                 Z(u)\det M=s^{\dim U} Z(A^*u).
\]
If \(v:U\to H'\) is another linear map such that for every \(x,y\in U\)
one has \(\langle vx,vy\rangle=a\langle x,Sy\rangle\) for
\(S\in{\rm End}(U)\), \(a\in\C\), then \(Z(v)=a^{\dim U}\det S\).
\end{lemma}
\begin{proof}
Regard \(u\) as an isomorphism \(\tilde u:U\to L\).
In orthonormal coordinates its squared absolute determinant is
\(Z(u)\). Since \(P_L\) fixes the image of \(u\),
\(\tilde u^*M\tilde u=s(A^*u)^*(A^*u)\).
Take determinants. For the last claim the assumed equality in all
matrix coefficients gives \(v^*v=aS\), so take determinants again.
\end{proof}

\begin{proposition}[Cancellation of the two volume ratios]\label{prop:long-criterion}
Use \(k,m,B_0',D_0,H\) from Definition~\ref{def:heights}, and put
\(E_0=\ell B_0'\). Let \(d_0=\dim W>0\) and assume the hypotheses of
Corollary~\ref{cor:completion-height}. Suppose \(n\ge(1-4/b^2)d_0\),
\(\beta+g\le b^3/4\), and every \(u_{t_\iota}\) is injective.
For \(t=0\) and all \(t_\iota\) suppose there are
\(c_t\in\C^\times,\lambda_j(t)>0\ (1\le j\le m)\) for which
\[
 \begin{split}
 Z(u_t)&=|c_t|^{d_0}\prod_{j=1}^m\lambda_j(t),\\
 {\rm ht}(|c_{t_\iota}/c_0|^{-2})&\le2E_0,\qquad
 {\rm ht}\!\left(\lambda_j(0)/\lambda_j(t_\iota)\right)
      \le d_0 k\left(\frac{E_0}{mk}-D_0\right)\quad(1\le j\le m).
 \end{split}
\]
For the parameter range of Proposition~\ref{prop:38},
\[
                          d_0(H-3E_0)\le V+2\delta .
\]
\end{proposition}
\begin{proof}
Write \(Z_\iota=Z(u_{t_\iota})/Z(u_0)>0\).
The completion term \((d_0-n)(\beta+g)\le d_0 b\), since
\(0\le d_0-n\le4d_0/b^2\) and \(b>0\).
By Corollary~\ref{cor:completion-height},
\({\rm ht}(Z_\iota)\le V+2\delta+d_0 b\).
The \(m\) ratios of \(\lambda_j\) have product of height at most
\(d_0(E_0-mkD_0)=d_0(E_0-H)\) by
Proposition~\ref{prop:height-product}. Square that product and multiply
by \(|c_{t_\iota}/c_0|^{-2d_0}\). The factorization shows the result
is \(Z_\iota^{-2}\), so its height is at most \(2d_0(2E_0-H)\).
Multiplication by \(Z_\iota^2\) gives the constant family \(1\);
the same proposition therefore gives
\(0\le2(V+2\delta+d_0b)+2d_0(2E_0-H)\).
Here \(E_0=fa\ell^4\ge b\) by the parameter hypotheses.
Rearrangement yields the assertion.
\end{proof}

\begin{corollary}\label{cor:row-test}
For all sufficiently large \(x\) and every \(t\le v\le t^2\) in the
range of Theorem~\ref{thm:geometric-reduction}, suppose the hypotheses
of Propositions~\ref{prop:selected-family} and~\ref{prop:long-criterion}
hold for the same space of rows and its selected family. Assume each
selected translate costs at most \(c_0B_0'\) in logarithmic index and
the images of the selected rows under \(u_0\) all have squared norm
at most \(\exp(H-m\mathcal R+E_0)\) in the archimedean coordinates
of Corollary~\ref{cor:completion-height}, for every \(\sigma\).
Here the index costs bound the sums of the exponents \(e_p\) there
by the sum of the costs of the translates. Then
\[
                         m\mathcal R\le26fa\ell^4 .
\]
If the hypotheses hold for all the choices of \(q,d,\epsilon\) in
Theorem~\ref{thm:geometric-reduction}, Dirichlet GRH follows.
\end{corollary}
\begin{proof}
Lemma~\ref{lem:hadamard} permits \(V=n(H-m\mathcal R+E_0)\).
Proposition~\ref{prop:selected-family} gives the dimension range.
Corollary~\ref{cor:26} gives \(\delta\le3d_0\log(1+b^2)B_0'\).
Proposition~\ref{prop:long-criterion} proves the lower volume
inequality, and Proposition~\ref{prop:height-cost} gives the crude
bound on \(H,m\mathcal R\). Thus all the numerical hypotheses of
Corollary~\ref{cor:26} hold with \(r=d_0,\ T=m\mathcal R\);
this gives the claimed bound. It implies the hypothesis with
constant \(38\) in Theorem~\ref{thm:geometric-reduction}, which in
turn gives the assertion through Corollary~\ref{cor:prime-suffices}.
\end{proof}

\section{The different and the additive character for the unitary groups}

\begin{lemma}\label{lem:integer-extension}
Let \(A\) be a complete discrete valuation ring with uniformizer \(p\),
meaning that every nonzero element has the form \(p^j u\), \(u\) a
unit, and residues modulo \(p^l\) determine elements by compatible
limits. Let \(F=A/pA\). If \(P\in A[X]\) is monic and its reduction is
a product of distinct monic irreducibles, then \(A[X]/(P)\) is a
product of complete discrete valuation rings with uniformizer \(p\).
If
\[
 P=X^e+p(c_{e-1}X^{e-1}+\cdots+c_0),\qquad c_i\in A,\quad c_0\in A^\times,
\]
then \(A[X]/(P)\) is a complete discrete valuation ring of uniformizer
the image of \(X\), and \(p\) is a unit times its \(e\)-th power.
Such rings are integrally closed in their fraction fields.
\end{lemma}
\begin{proof}
Lift a factorization \(\bar P=UV\) into coprime monic factors over
\(F\). Given monic lifts whose product is correct modulo \(p^l\), add
corrections \(p^l U_1,p^l V_1\) to them, of degrees less than the
respective factors, to correct the product modulo \(p^{l+1}\).
The map \((U_1,V_1)\mapsto U_1V+UV_1\) is injective since \(U,V\)
are coprime, and domain and codomain (polynomials of degree less
than \(\deg P\)) have equal dimension. It gives the corrections.
Completeness gives two exact monic factors. They generate the unit
ideal in \(A[X]\): the same map for the lifted factors on bounded
coefficient lists has determinant a unit by its reduction and can
therefore represent the polynomial \(1\). For comaximal factors the
quotient by their product maps isomorphically to the product of
the quotients, by their Bezout identity. Repeating splits all factors.
For each resulting ring \(B\) the powers below the degree are an
\(A\)-basis by monic division. Hence \(B\) is complete and
\(\bigcap_l p^l B=0\), by coordinates. An element with nonzero
reduction in its residue field has an inverse: lift a residue inverse,
then invert \(1-pz\) by its convergent geometric series. Every
nonzero element is therefore a power of \(p\) times a unit, and \(p\)
is a nonzero divisor by freeness. This gives the first assertion.

In the second case put \(\pi=X\bmod P\) and \(B=A[\pi]\).
The same power basis and argument about geometric series show that an
element is a unit if its image in \(F[X]/(X^e)\) is a unit, that is,
if its constant coefficient is nonzero modulo \(p\). Thus
\(c_{e-1}\pi^{e-1}+\cdots+c_0\) is a unit and \(\pi^e\) is a unit
times \(p\). In a nonzero sum \(\sum_{i<e}a_i\pi^i\) the nonzero
terms are units times the powers \(e\,v_A(a_i)+i\) of \(\pi\).
These exponents are distinct. Factoring out the least leaves a unit,
since the other terms are multiples of \(\pi\). The element \(\pi\)
is a nonzero divisor since \(\pi^e\) is one; the ring is therefore
a domain with the stated valuation. The powers of \(p,\pi\) give
the same completeness. In a monic equation over either valuation
ring an element of negative valuation would give the leading term a
unique least valuation. Dividing by that term contradicts that the
maximal ideal is proper. This proves integral closure.
\end{proof}

\begin{theorem}\label{thm:integer-ring}
Put \(n=N_1,\ \zeta=\zeta_n,\ \omega=(1+\sqrt{-d})/2\). Then
\[
             O_E=\Z[\omega,\zeta]
                  =\bigoplus_{i=0}^1\bigoplus_{j=0}^{f-1}
                    \Z\,\omega^i\zeta^j .
\]
\end{theorem}
\begin{proof}
Lemma~\ref{lem:rd} gives the indicated order \(\mathcal B\) with this
basis. Fix a rational prime \(p\); write \(n=p^a h,\ p\nmid h\).
The cyclotomic order localized in coefficients to \(\Z_p\) is
generated by \(\zeta_h\) and, for \(a>0\), \(\zeta_{p^a}\).
These are powers of \(\zeta\), and conversely a Bezout identity for
\(h,p^a\) expresses \(\zeta\) as a product of powers of them.
Their separate monic cyclotomic relations present this order:
the resulting map is surjective and the two free modules have the
same rank \(\varphi(h)\varphi(p^a)=\varphi(n)\), by
Lemma~\ref{lem:basic-arith} on units. A kernel in a free module
that vanishes after passage to its fraction field is zero. The same
argument permits adjoining \(\omega\) by its separate quadratic
relation in any order, since the total rank in Lemma~\ref{lem:rd}
is \(2f\).

The polynomial \(P_h\) of Lemma~\ref{lem:cyclo} divides \(X^h-1\);
its reduction is separable because \(p\nmid h\). If \(p\ne d\), the
polynomial of \(\omega\) is separable modulo \(p\) because its
discriminant is \(-d\), nonzero also in characteristic two.
Apply Lemma~\ref{lem:integer-extension} first to \(P_h\) and then
in each resulting factor to the quadratic. If \(a>0\), adjoin
\(\zeta_{p^a}-1\) in each factor. Its polynomial is
\[
  P_{p^a}(1+X)=\sum_{u=0}^{p-1}(1+X)^{u p^{a-1}} .
\]
The constant is \(p\); the reduction is \(X^{\varphi(p^a)}\), since
multiplying by \((1+X)^{p^{a-1}}-1\) gives \((1+X)^{p^a}-1\).
Thus the second case of Lemma~\ref{lem:integer-extension} applies.
If \(p=d\), Proposition~\ref{prop:places} gives \(a=0\), so after
adjoining \(\zeta_h\) use \(\sqrt{-d}\) and its polynomial \(X^2+d\).
This generates the same ring as \(\omega\), since \(d\) is odd.
Again the lemma applies. In every case \(\mathcal B\otimes\Z_p\)
is a product of integrally closed rings, with total fraction algebra
\(E\otimes\Q_p\) by the power basis. A global algebraic integer
satisfies the same monic equation in each local fraction field, hence
lies in every \(\mathcal B\otimes\Z_p\). Its rational coefficients in
the displayed basis lie in \(\Z_p\) for all \(p\), so have no prime
in their reduced denominators by Lemma~\ref{lem:basic-arith}.
They are integers. The opposite inclusion is Lemma~\ref{lem:rd}.
\end{proof}

\begin{lemma}\label{lem:dual-monic}
Let \(\alpha\) have monic irreducible polynomial \(P\in\Z[X]\)
of degree \(h\ge1\), and put \(C=\Z[\alpha]\).
In \(\Q(\alpha)\), the trace dual
\(C^\vee=\{a:\operatorname{Tr}(aC)\subset\Z\}\) is
\(P'(\alpha)^{-1}C\).
\end{lemma}
\begin{proof}
For \(\deg T<h\), interpolation on the distinct complex roots \(z\)
of \(P\) writes \(T(X)\) as the sum of \(T(z)P(X)/((X-z)P'(z))\).
This holds since the polynomials agree on all \(h\) roots and have
degree less than \(h\). Comparison of leading coefficients gives
\(\operatorname{Tr}(T(\alpha)/P'(\alpha))=[X^{h-1}]T\).
Here multiplication in the power basis diagonalizes upon evaluation
at the \(h\) roots, exactly as in Lemma~\ref{lem:rd}.
For products of elements of \(C\) first take the remainder modulo
\(P\). On the power basis the matrix of
\((x,y)\mapsto\operatorname{Tr}(xy/P'(\alpha))\) thus has integer
entries, zeros for \(i+j<h-1\) and ones for \(i+j=h-1\).
Reversing the columns makes it triangular with diagonal one.
Its inverse is integral by its adjugate, proving the duality.
\end{proof}

\begin{proposition}\label{prop:different}
Write \(\mathfrak D_E=\{u\in E:u O_E^\vee\subset O_E\}\), the
different of \(E/\Q\). With \(n,\zeta\) as in
Theorem~\ref{thm:integer-ring}, set
\[
       h_0'=1-\varphi(n)/2-n/4,\qquad
       \varpi=\zeta^{h_0'}P_n'(\zeta)\sqrt{-d}.
\]
Then \(\mathfrak D_E=\varpi O_E,\ \bar\varpi=-\varpi\), and
\[
 \prod_{\tau\in\Phi}|\tau(\varpi)|=\operatorname{rd}_E^f,\qquad
       \operatorname{rd}_E=\sqrt d\, n\prod_{p\mid n}p^{-1/(p-1)} .
\]
\end{proposition}
\begin{proof}
In the basis of Theorem~\ref{thm:integer-ring}, field trace of
a product from the quadratic and cyclotomic factors is the product
of the two traces: the embeddings vary independently by
Proposition~\ref{prop:places}. Apply Lemma~\ref{lem:dual-monic} to
both factors. Products of their dual basis elements are dual to the
product basis. It follows that \(O_E^\vee=(P_n'(\zeta)\sqrt{-d})^{-1}O_E\).
Here \(4\mid n\) and \(\varphi(n)\) is even, since the unit classes
are paired by negation for \(n>2\). The reciprocal relation
\(P_n(X^{-1})=X^{-\varphi(n)}P_n(X)\) follows from inversion of
the primitive roots and the constant coefficient one (their inverse
pairs have product one). Differentiate it at \(\zeta\) to obtain
\(P_n'(\bar\zeta)=-\zeta^{2-\varphi(n)}P_n'(\zeta)\).
As \(\zeta^{2-\varphi(n)-2h_0'}=\zeta^{n/2}=-1\), the first two
factors of \(\varpi\) have real product. Multiplication by this
power of the unit \(\zeta\) preserves the generated ideal.

The trace matrix identifies \(O_E^\vee\) with the inverse trace
matrix applied to \(\Z^{2f}\), so its covolume relative to \(O_E\)
is \(1/|\Delta_E|\). Since \(O_E^\vee=\varpi^{-1}O_E\), diagonalizing
multiplication at the embeddings as in Lemma~\ref{lem:rd} gives
\(|\Delta_E|=\prod_{\sigma:E\to\C}|\sigma(\varpi)|\).
Pairing conjugates proves the first product formula.
For the second, the identity \(X^n-1=\prod_{v\mid n}P_v(X)\) implies
\[
 P_n(X)=\prod_{j\mid n}(X^{n/j}-1)^{\mu(j)}.
\]
Indeed the exponent of each cyclotomic factor on the right is
\(\sum_{j\mid n/v}\mu(j)={\bf1}_{v=n}\), by multiplication over prime
divisors as in Proposition~\ref{prop:twist-test}. Divide by \(X-\zeta\)
and evaluate the rational expression at \(\zeta\), isolating the
factor \(j=1\). At a primitive \(j\)-th root \(\xi\), the absolute
product of \(1-\xi\) over its conjugates is \(P_j(1)>0\).
We have \(P_{p^a}(1)=p\) by the geometric sum, and \(P_j(1)=1\)
for \(j>1\) with at least two prime divisors: evaluating
\((X^j-1)/(X-1)\) at \(1\) gives \(\prod_{h\mid j,h>1}P_h(1)=j\),
and induction gives the claim. Each primitive \(j\)-th root occurs
\(f/\varphi(j)\) times as \(\zeta^{n/j}\) varies over the conjugates
of \(\zeta\), by the surjectivity of unit reduction in
Lemma~\ref{lem:induce}. The product over the \(f\) embeddings fixing
\(K\) of \(|P_n'(\zeta)|\) is consequently
\(n^f\prod_{p\mid n}p^{-f/(p-1)}\). Multiply by \(d^{f/2}\) and
take \(f\)-th roots.
\end{proof}

\begin{proposition}\label{prop:additive}
For \(u\in\Q_p\) let \(\{u\}_p\) be the rational number in \([0,1)\)
specified by its negative-position base-\(p\) digits. Thus
\(u-\{u\}_p\in\Z_p\) and its denominator is a power of \(p\).
Define \(\mathbf e_p(u)=\exp(2\pi i\{u\}_p)\).
For \(E_p=E\otimes\Q_p,\ O_p=O_E\otimes\Z_p\), put
\(\psi_p(y)=\mathbf e_p(\operatorname{Tr}_{E_p/\Q_p}y)\).
For a finite adele \(y=(y_p)\), that is \(y_p\in E_p\) and
\(y_p\in O_p\) for all but finitely many \(p\), put
\[
               \psi_f(y)=\prod_p\psi_p(y_p),\qquad
               \psi^0_f(y)=\psi_f(-y/\varpi).
\]
These are additive characters, and \(\psi_f(a)=\exp(2\pi i
\operatorname{Tr}_{E/\Q}a)\) for diagonal \(a\in E\). The
annihilator of \(O_p\) under the pairing
\((y,z)\mapsto\psi_p(-\bar y z/\varpi)\) is \(O_p\) at every \(p\).
The same holds for \(\prod_p O_p\) under the product pairing.
\end{proposition}
\begin{proof}
Two fractions with power-of-\(p\) denominators differing by an
element of \(\Z_p\) differ by an integer. Therefore \(\mathbf e_p\)
is additive with kernel \(\Z_p\). Integral elements of \(O_p\) have
integral trace from the multiplication matrix in the basis of
Theorem~\ref{thm:integer-ring}. The finite products make sense,
also after division by \(\varpi\), whose inverse has just finitely
many primes in its basis denominators. For \(r\in\Q\),
\(r-\sum_p\{r\}_p\) belongs to every \(\Z_p\), since the fractions
at the other primes are \(p\)-integral. Hence it is an integer,
giving the formula at a diagonal element by taking \(r=\operatorname{Tr}a\).
By the same inverse trace matrix as in Proposition~\ref{prop:different},
the trace dual of \(O_p\) is \(\varpi^{-1}O_p\).
Thus \(\psi_p(-\bar y z/\varpi)=1\) for every \(z\in O_p\)
exactly when \(\bar y\in O_p\), equivalently \(y\in O_p\).
Varying \(z\) at one prime at a time gives the global annihilator.
\end{proof}

\section{Singular frame minors and the second moment}

\begin{definition}\label{def:positive-weight}
On a finite group \(J\) call \(\theta:J\to\R_{\ge0}\) a nonnegative
weight of positive type if there are vectors \(u_a\) in a complex
inner product space such that
\(\langle u_a,u_b\rangle=\theta(a^{-1}b)\).
Write \(\theta(S)=\sum_{a\in S}\theta(a)\) for \(S\subset J\).
If \(J\) acts on a finite set and \(v\) is a member, restricting
\(\theta\) to \({\rm Stab}(v)\), by setting it zero elsewhere, preserves
positive type. Indeed use vectors having the entry \(u_a\) in
coordinate \(av\) of the direct sum over the finite set and zero in
the other coordinates.
\end{definition}

\begin{lemma}[A finite Fourier bound for a singular minor]\label{lem:minor-sieve}
Let \(p\) be prime, \(F=\mathbf F_p\), \(\pi:J\to{\rm GL}_r(F)\),
and \(A,B\) be \(r\)-by-\(j\) matrices over \(F\) with \(j\ge1\).
Write \(\pi^\vee(s)=\pi(s)^{-\mathsf T}\).
For a weight \(\theta\) of Definition~\ref{def:positive-weight} set
\(M_\theta=\theta(J)\). Suppose for all \(x\in F^j\setminus\{0\}\)
that the \(\theta\)-masses of the \(\pi\)-stabilizer of \(Ax\) and
of the \(\pi^\vee\)-stabilizer of \(Bx\) are at most \(Y\ge0\).
Then
\[
 \begin{aligned}
 (p-1)\,\theta\{s:\det(B^{\mathsf T}\pi(s)A)=0\}
 &\le (1-p^{-j})M_\theta\\
 &\quad+p^{-j}(p^j-1)^2p^{r/2}Y .
 \end{aligned}
\]
\end{lemma}
\begin{proof}
Put \(\chi(t)=\exp(2\pi i t/p)\) for a residue represented by an
integer. Geometric summation in one coordinate shows
\(\sum_{w\in F^r}\chi(z^{\mathsf T}w)=p^r{\bf1}_{z=0}\).
In the direct sum of \(p^r\) copies of the space of the \(u_a\), let
\(D_z u\) place \(u\) at coordinate \(z\) and zero elsewhere,
and let \(H_z u\) have entries \(p^{-r/2}\chi(-z^{\mathsf T}w)u\)
at \(w\in F^r\). These two families each have orthogonal coordinate
ranges for distinct \(z\), by the displayed character sum.
For nonzero \(x,y\in F^j\) form the vectors
\[
       V=\sum_{a\in J}H_{\pi^\vee(a)By}u_a,\qquad
       W=\sum_{b\in J}D_{\pi(b)Ax}u_b .
\]
On expanding \(\|V\|^2\), a pair \(a,b\) contributes precisely when
\(a^{-1}b\) fixes \(By\) under \(\pi^\vee\). Every such element
occurs \(|J|\) times and contributes its weight each time. Thus
\(\|V\|^2\le|J|Y\), and the same argument gives \(\|W\|^2\le|J|Y\).
In the mixed pairing each coefficient has the positive phase
\(\chi(y^{\mathsf T} B^{\mathsf T}\pi(a^{-1}b)Ax)\), since the
first variable is conjugated. Cauchy--Schwarz therefore gives
\[
  \left|\sum_s \chi(y^{\mathsf T} B^{\mathsf T}\pi(s)Ax)
                        \theta(s)\right|\le p^{r/2}Y .
\]
For fixed \(s,y\), averaging the phase over \(x\in F^j\) gives the
indicator of \(y\) belonging to the left kernel of
\(B^{\mathsf T}\pi(s)A\). Summing over \(y\ne0\) counts its nonzero
left kernel vectors. A singular matrix has at least \(p-1\) such
vectors. The \(x=0\) term in this restricted sum contributes
\((1-p^{-j})M_\theta\); the remaining pairs number \((p^j-1)^2\)
and each contributes at most \(p^{r/2}Y\) before averaging.
This proves the bound.
\end{proof}

\begin{definition}\label{def:sieve-cell}
At \(p=g_0\) let \(\mathcal W\) index the \(f\) split pairs in
Proposition~\ref{prop:row-group}. In its chosen basis let \(J_w\)
be the matrix of the integral hermitian form on the split pair,
reduced on one factor. In that basis it is the symmetric invertible
matrix in the proof of that proposition. Put \(h=\nu+1\).
Let \(F_j\) be the matrix whose columns are the first \(j\) elements
of the list \(e_1,\ldots,e_\nu,z\), \(1\le j\le h\), in the integral
basis of Definition~\ref{def:row-space}. For \(s\in\mathcal G\)
write \(s_w\in{\rm GL}_r(\mathbf F_p)\) for its component at \(w\).
The leading minors at \(w\) are
\[
                 \Delta_{w,j}(s)=\det(F_j^{\mathsf T}J_w s_w F_j).
\]
At \(w\) a prescribed vector condition means either that \(s_w\)
fixes \(F_h a\), or that \(s_w^{-\mathsf T}\) fixes \(J_w^{\mathsf T}F_h a\),
for a specified \(a\ne0\) in \(\mathbf F_p^h\). The latter condition
expresses fixing the frame vector in the other component, since its
matrix is \(J_w^{-1}s_w^{-\mathsf T}J_w\). Given \(S\subset\mathcal W\)
and one condition \(D_w\) per \(w\in S\), write \(V(S,D)\) for their
common solution set. Let \(\Omega(S,D)\) be the subset of \(V(S,D)\)
where either \(S=\mathcal W\), or at some \(w\in\mathcal W\) all
\(h\) leading minors are nonzero.
\end{definition}

\begin{proposition}\label{prop:stabilizer-index}
For the group and conditions of Definition~\ref{def:sieve-cell},
\([\mathcal G:V(S,D)]=(p^r-1)^{|S|}\).
Eventually \(p^{h+r/2}\le b^{6k}\) and \(3h/p\le b^{-4}\).
\end{proposition}
\begin{proof}
Each of the vectors is nonzero because the frame vectors belong to a
basis and \(J_w\) is invertible. In \({\rm GL}_r(\mathbf F_p)\) any
nonzero vector can be sent to any other by extending both to bases.
The same holds for the inverse-transpose action since inverse transpose
is an automorphism. Cosets of the stabilizer biject to the \(p^r-1\)
nonzero vectors. Different \(w\) occur in independent direct factors
of Proposition~\ref{prop:row-group}. For the bounds,
\(r=k+2\nu+1,\ h=\nu+1,\ k\ge x^5-2\sqrt x\), and
\(\nu\le\sqrt x\) by Proposition~\ref{prop:thick}. It follows that
\(11(h+r/2)=11k/2+22\nu+33/2\le6k\) eventually.
Theorem~\ref{thm:data} gives \(b^{10}\le p\le b^{11}\); also \(h\le b\)
eventually by the exponential comparison in Proposition~\ref{prop:thick}.
These inequalities give both claims.
\end{proof}

\begin{theorem}\label{thm:finite-moment}
Let \(a:\mathcal G\to\R_{\ge0}\), \(D_*>0\), and put
\(C_S=(1+9b^{-4})D_* b^{-6k|S|}\). Work at sufficiently large
\(x\) as in Proposition~\ref{prop:stabilizer-index}. Suppose for every
\(S,D\) in Definition~\ref{def:sieve-cell} that
\[
                         \sum_{s\in\Omega(S,D)}a(s)
                                   \le(1-b^{-4})C_S .
\]
For any nonnegative weight \(\theta\) of positive type bounded by
\(a\) at every element and supported on \(V(S,D)\), one has
\(\theta(\mathcal G)\le C_S\). In particular if \(a\) itself has
positive type then \(\sum_s a(s)\le(1+9b^{-4})D_*\).
\end{theorem}
\begin{proof}
Proceed by descending induction on \(|S|\), proving the conclusion
for all prescribed conditions and weights at once. For
\(S=\mathcal W\) the support lies in \(\Omega(S,D)\), giving the
claim. Otherwise take \(w\notin S\). Suppose \(M_\theta>C_S\).
Restricting to any one additional condition at \(w\) preserves
positive type by Definition~\ref{def:positive-weight}. The new weight
is still bounded by \(a\). Induction bounds its mass by
\(C_{S\cup\{w\}}=b^{-6k}C_S<b^{-6k}M_\theta\).

For \(1\le j\le h\) take \(A=F_j,\ B=J_w^{\mathsf T}F_j\)
and \(\pi(s)=s_w\) in Lemma~\ref{lem:minor-sieve}. Each stabilizer
there is one of the allowed extra conditions by padding vectors with
zeros. Thus we can use \(Y=b^{-6k}M_\theta\); by
Proposition~\ref{prop:stabilizer-index}, \(p^{j+r/2}Y\le M_\theta\).
Lemma~\ref{lem:minor-sieve} bounds the mass with \(\Delta_{w,j}=0\)
by \(2M_\theta/(p-1)\le3M_\theta/p\) since \(p\ge3\).
On the support outside \(\Omega(S,D)\), at least one of these minors
vanishes at this particular \(w\). Summing in \(j\) shows the mass
outside is at most \(3h M_\theta/p\le b^{-4}M_\theta\).
The assumed bound on \(a\) inside the set now yields
\((1-b^{-4})M_\theta\le(1-b^{-4})C_S\), a contradiction. The last
assertion takes \(S=\varnothing,\theta=a\).
\end{proof}

\begin{corollary}\label{cor:concrete-moments}
Use \(\mathcal G,\rho,Q,c_0\) as in Proposition~\ref{prop:selected-family}
and let \(R_i\), \(i=0,1\), be commuting orthogonal projections with
\(R=R_0R_1\) commuting with \(\rho\).
Set \(a(s)=\|Q^*\rho(s)Q\|_{\rm HS}^2\). Suppose the mass hypothesis
of Theorem~\ref{thm:finite-moment} holds for this \(a\) and some
\(D_*>0\). Suppose further that \(d_0=\dim(R\mathcal H)\)
satisfies \(d_0D_*\le|\mathcal G|c_0^2\) and that
\[
 \sum_j\|(Q^*Q-I)e_j\|\le c_0 b^{-4},\qquad
 \operatorname{tr}\big(Q^*(1-R_i)Q\big)\le2c_0 b^{-4}
                                      \quad(i=0,1).
\]
Then the conclusion on the selected family in
Proposition~\ref{prop:selected-family} holds for all sufficiently
large \(x\).
\end{corollary}
\begin{proof}
The weight \(a\) is of positive type. Take, as its vectors,
\(U_t=\rho(t)QQ^*\rho(t)^*\) with inner product
\(\operatorname{tr}(U_t^*U_v)\); cyclic permutation as in
Lemma~\ref{lem:proj-moments} gives \(a(t^{-1}v)\).
By Theorem~\ref{thm:finite-moment},
\(|\mathcal G| E_2\le(1+9b^{-4})D_*\).
The operator \(1-R\) is bounded above as a quadratic form by
\((1-R_0)+(1-R_1)\), because their difference after subtraction of
\(1-R\) is \((1-R_0)(1-R_1)\), an orthogonal projection.
The trace inequalities therefore give \(E_1\le5c_0 b^{-4}<c_0\).
This implies \(d_0>0\) by Lemma~\ref{lem:proj-moments}, and the
dimension hypothesis gives \(E_2\le(1+9b^{-4})c_0^2/d_0\).
Finally \(f\le b\) eventually, since \(f\le x^{12}\) in the proof of
Proposition~\ref{prop:level-prime}, and \(N-1\ge b^9\) by that
proposition. Thus \((5+9/2)b^{-4}+f/(N-1)\le b^{-2}\) eventually.
Proposition~\ref{prop:selected-family} now applies.
\end{proof}

\section{The unipotent covolume at the principal level}

\begin{lemma}\label{lem:skew-covol}
Embed \(E\) diagonally in \(\C^\Phi\) using the embeddings fixing
\(K\) in Proposition~\ref{prop:places}. Real and imaginary coordinates
carry ordinary Lebesgue measure. The covolume of \(O_E\) is
\(V_E=2^{-f}\sqrt{|\Delta_E|}\). The image of \(O_E^-=\{s\in O_E:
\bar s=-s\}\) under \(s\mapsto(\operatorname{Im}\tau s)_\tau\)
has covolume \(|\Delta_E|^{1/4}\) in \(\R^\Phi\).
\end{lemma}
\begin{proof}
The complex embedding matrix of an integral basis has absolute
determinant \(\sqrt{|\Delta_E|}\), by the trace calculation in
Lemma~\ref{lem:rd}. Each pair of conjugate rows changes from real
and imaginary rows by a matrix of determinant of modulus two.
The real coordinate embedding matrix thus gives the stated \(V_E\).

Let \(F=E^{\mathfrak c}\) be the fixed field of conjugation and
\(O_F=F\cap O_E\). Since \(12\mid N_1\), \(E\) contains both
\(i\) and \(\xi=\exp(2\pi i/12)\). We have \(\xi-\bar\xi=i\), a unit.
For every \(x\in O_E\) the elements
\(b=(x-\bar x)/i\) and \(a=x-b\xi\) lie in \(O_F\).
The decomposition \(x=a+b\xi\) is unique by conjugating
and subtracting. Thus \(O_E=O_F\oplus \xi O_F\), and \(O_F\)
has rank \(f\) by the free-subgroup argument of Lemma~\ref{lem:rd}.
In a basis of \(O_F\), its real embedding matrix \(T\) is nonsingular:
the decomposition just obtained gives a full lattice in \(\C^\Phi\).
Its determinant in real coordinates for the two summands gives
\(V_E=|\det T|^2\prod_\tau|\operatorname{Im}\tau\xi|
=2^{-f}|\det T|^2\), since \(\tau(i)=\pm i\).
Finally \(O_E^-=iO_F\), whose imaginary embedding matrix differs
from \(T\) by row signs. Its covolume is \(|\Delta_E|^{1/4}\).
\end{proof}

\begin{lemma}\label{lem:class-covol}
Let \(h\ge1\). Let \(B_p\in{\rm GL}_h(E_p)\) lie in
\({\rm GL}_h(O_p)\) for all but finitely many \(p\), and let
\[
             \Lambda=\{x\in E^h:x\in B_p O_p^h\text{ for every }p\}.
\]
In coordinates \(\C^{h\Phi}\) its covolume is
\(V_E^h\prod_p |\det_{\Q_p} B_p|_p^{-1}\), where the determinant
is taken on the underlying \(\Q_p\)-vector space. Its \(O_E\)-linear
dual lattice of row vectors, locally specified by \(x B_p\in O_p^h\),
has covolume \(V_E^h\prod_p|\det_{\Q_p} B_p|_p\).
\end{lemma}
\begin{proof}
There is a positive integer \(a\) with
\(aO_p^h\subset B_p O_p^h\subset a^{-1}O_p^h\) simultaneously for all
\(p\), by clearing denominators of the finitely many matrices and
inverses in an integral basis. The quotient
\(a^{-1}O_E^h/aO_E^h\) is the product of the corresponding local
quotients, by multiplying by \(a\) and using
Lemma~\ref{lem:basic-arith} coordinatewise modulo \(a^2\).
Consequently these local lattices specify all the reductions of
\(\Lambda\), and it has full rank. Relative indices are the products
of the local relative indices as seen by dividing the two counts
above the common subgroup \(aO_E^h\).
For a matrix over \(\Q_p\), integral row and column operations reduce
it to diagonal form: move an entry of least valuation to the first
position, use its integral entry ratios to clear its column and row,
and repeat on the remaining block. Row and column operations with
unit determinant preserve indices. A diagonal entry \(p^e u\)
contributes relative covolume \(p^e\) because
\(|\Z_p/p^j\Z_p|=p^j\) for \(j\ge0\); negative \(e\) gives the
reciprocal by interchanging the lattices. Multiplying these factors
gives \(|\det_{\Q_p}B_p|_p^{-1}\).
An index \(n\) of two ordinary full lattices multiplies the covolume
by \(n\): coset representatives give a union of \(n\) translates of
a fundamental parallelepiped. This proves the first formula by
Lemma~\ref{lem:skew-covol}. The dual row lattice is specified in
column coordinates by the inverse transpose over \(E_p\).
Its \(\Q_p\)-determinant has inverse absolute value to the original:
after extending scalars to \(K_p\), diagonalize multiplication on
the separable field factors by evaluation at their embeddings as in
Lemma~\ref{lem:rd}, and take the product of ordinary determinants.
This gives the last formula.
\end{proof}

\begin{definition}\label{def:unip}
In \(Y_0\) from Definition~\ref{def:row-space} number the hyperbolic
pairs \(e_a,e_a^*\) from \(a=0\) to \(\nu-1\). Write
\(M_c\) for the middle integral lattice locally obtained by applying
\(c_p\oplus1\) to the middle lattice of that definition, where
\(c=(c_p)_p\) is any finite adelic isometry of the positive
\(k\)-dimensional space, integral at almost every \(p\).
Adjoin the standard integral hyperbolic planes to obtain a lattice
\(\Lambda_c\), using Lemma~\ref{lem:class-covol}.
Let \(U\) be the rational unitary subgroup preserving the flag
\(\langle e_0,\ldots,e_a\rangle\), acting identically on each of
its successive lines and on the middle quotient of the orthogonal
flag. On each dual line it then acts identically by the pairing.
Set
\[
        \Gamma_c=\{u\in U(E^+):(u-I)\Lambda_c\subset g\Lambda_c\},
        \qquad E^+=E^{\mathfrak c},\quad g=g_0g_1 .
\]
Here \(U(E^+)\) denotes the matrices on the \(E\)-space with entries
in \(E\) satisfying these unitary and flag equations. Integrality
of \((u-I)/g\) at every prime is precisely integrality everywhere
and principal congruence at \(g_0,g_1\). At a split pair the two
congruences agree by inverse transpose and unimodularity.
The group is closed under inverses since \(u-I\) is strictly block
upper triangular in the ordered flag and therefore nilpotent.
\end{definition}

\begin{lemma}\label{lem:heisenberg-layer}
Remove the outermost remaining pair \(e_a,e_a^*\) of
Definition~\ref{def:unip}. The intervening space, of dimension
\(n_a=2(\nu-1-a)+k+1\), has hermitian matrix \(J\) and integral
lattice \(L\) consisting of \(M_c\) and the remaining hyperbolic
planes. The projection of \(U\) onto the analogous inner flag group
\(U'\) splits by \(u'\mapsto{\rm diag}(1,u',1)\).
Its kernel consists of the matrices
\[
 H(l,t)=
 \begin{pmatrix}1&l&t\\0&I&-J^{-1}l^*\\0&0&1\end{pmatrix},
 \qquad t+\bar t=-lJ^{-1}l^* .
\]
On arithmetic subgroups the same projection splits. In the kernel
the conditions are precisely
\[
    l=g x,\quad xL\subset O_E,\qquad
    t=-\omega\, lJ^{-1}l^*+g s,\qquad s\in O_E^- ,
\]
where \(\omega\) is from Theorem~\ref{thm:integer-ring}.
\end{lemma}
\begin{proof}
In the outer decomposition the form has antidiagonal corner entries
one and middle block \(J\). A flag element fixes the first basis
vector and preserves its perpendicular. Passing to its quotient by
that vector gives \(u'\). Multiplying by
\({\rm diag}(1,u'^{-1},1)\) leaves an upper triangular block matrix
with identity diagonal. Applying the hermitian identity to its
last column and the middle columns forces its off-diagonal blocks
and trace equation to be exactly the displayed ones. Substitution
proves the converse and the splitting.

The inner lattice is hermitian self-dual at each finite prime:
the initial Gram determinant is a unit by Proposition~\ref{prop:row-group},
the other hyperbolic planes are unimodular, and \(c_p\oplus1\)
is an isometry. Thus \(x L\subset O_E\) is equivalent locally
to \(J^{-1}x^*\in L_p\) for every \(p\), and hence to the global
membership by Lemma~\ref{lem:class-covol}. In particular
\(xJ^{-1}x^*\in O_E\) is fixed by conjugation. Applied to the blocks
of \(H-I\), principal congruence now requires precisely \(l=gx\) as
stated and \(t\in gO_E\). Since \(\omega+\bar\omega=1\), the number
\(-\omega lJ^{-1}l^*\) is one solution of the trace equation in
\(gO_E\) (in fact it belongs to \(g^2O_E\)). The other solutions
differ by \(gO_E^-\). The inner subgroup preserves the corresponding
congruence by the direct-sum lattice and its block diagonal splitting.
\end{proof}

\begin{proposition}\label{prop:unipotent-covol}
For \(\tau\in\Phi\) put
\[
 \kappa_\tau=\frac{g\,\tau(\varpi)}{2\pi i}\in\R^\times,\qquad
 s_{a,\tau}=\kappa_\tau^{\nu-a}/\sqrt N\quad(0\le a\le\nu).
\]
At \(\tau\), scale \(e_a,e_a^*\) by \(s_{a,\tau},s_{a,\tau}^{-1}\),
respectively, the positive middle basis by one, and \(z\) by
\(s_{\nu,\tau}=1/\sqrt N\). These are scales on the original
rational basis \(x_1,\ldots,x_k,z\) before replacing \(x_k\) in
the integral lattice. Use the coordinates of the scaled basis for
the real groups, so matrices transform by \(u\mapsto D_\tau^{-1}uD_\tau\)
with this diagonal \(D_\tau\).
Use Lebesgue measure at each successive kernel of
Lemma~\ref{lem:heisenberg-layer} on the complex row \(l\) and the
real coordinate \(\operatorname{Im}t\). Their product over places
and layers is invariant under both left and right translation.
Put \(d_*= \nu(\nu+k)\). The covolume of \(\Gamma_c\) for this measure
is independent of \(c\) and equals
\[
 \mathfrak v=
  g^{(2d_*+\nu)f}V_E^{d_*}|\Delta_E|^{\nu/4}
      \prod_{\tau\in\Phi}\prod_{a=0}^{\nu-1}
             (s_{a,\tau}^{\,2})^{-(k+2(\nu-a))}.
\]
\end{proposition}
\begin{proof}
Proposition~\ref{prop:different} shows that \(\kappa_\tau\) is real
and nonzero. In a real kernel the real part of \(t\) is determined
by \(l\). Matrix multiplication sends the rows to their sum and adds
the two entries \(t\) with the correction \(-lJ^{-1}l'^*\).
Translation thus preserves Lebesgue measure in \((l,\operatorname{Im}t)\)
by first translating \(l\) and then translating the real coordinate
by a quantity depending on \(l\). The inner flag group acts on \(l\)
on the right by the inverse middle matrix, of complex determinant one,
and fixes \(t\). This preserves the measure; the same statements in
scaled coordinates follow by the diagonal change of basis. In the split
extension the product measure is therefore left and right invariant,
as is seen by integrating first over the kernel and then over \(U'\).
Iteration proves the assertion for the whole group.

For the layer \(a\), the rational row lattice of \(x\) in
Lemma~\ref{lem:heisenberg-layer} has covolume \(N^f V_E^{n_a}\).
Indeed the finite local matrix of its primal lattice is the identity
on the inner hyperbolic planes and the middle matrix
\((c_p\oplus1)B\) on \(x_1,\ldots,x_k,z\), where \(B\) implements the
basis change of Definition~\ref{def:row-space}.
We have \(\det_E B=1/\bar\pi_N\).
The determinant of \(c_p\) times its conjugate is one; the determinant
over \(\Q_p\) is the product over the embeddings of \(E_p\), so
its \(p\)-adic absolute value is one. The product of the finite
absolute values of \(\det_{\Q} B\) is \(N^f\), since at infinity
\(|\det_{\Q}B|=N^{-f}\) by \(\pi_N\bar\pi_N=N\), and the rational
product formula is proved in Proposition~\ref{prop:height-product}.
Apply the row part of Lemma~\ref{lem:class-covol}.

The new row is \(g x\) with its entries at place \(\tau\) multiplied
by \(s_{j,\tau}/s_{a,\tau}\), where \(s_{j,\tau}\) here denotes the
scale of the inner basis vector in column \(j\).
All the paired inner scales cancel in their product; the remaining
factor is \(s_{\nu,\tau}=1/\sqrt N\). Thus the squared absolute
Jacobian over the complex row variables contributes \(N^{-f}\),
canceling the \(N^f\) in the rational row covolume. The row lattice
in the scaled coordinates has covolume
\[
          g^{2n_a f} V_E^{n_a}\prod_\tau (s_{a,\tau}^2)^{-n_a}.
\]
After the row is fixed, the lattice of central translations is given
by \(g\,\operatorname{Im}\tau(s)/s_{a,\tau}^2\) for \(s\in O_E^-\).
By Lemma~\ref{lem:skew-covol} its covolume is
\(g^f|\Delta_E|^{1/4}\prod_\tau(s_{a,\tau}^2)^{-1}\).

Choose half-open fundamental parallelepipeds for these two ordinary
lattices. Every element of the kernel has exactly one left arithmetic
translate with row and central coordinate in their respective
parallelepipeds: first translate the row, using an arithmetic lift
from Lemma~\ref{lem:heisenberg-layer}, then use a unique central
correction. Two possible lifts of the row differ by a central
translation. For the split extension first reduce the projection to
a fixed fundamental domain of the inner arithmetic group, then
reduce the kernel variables. Uniqueness follows by projecting two
putative representatives and then using uniqueness in the kernel.
The measures of the domains multiply by the preceding product
description. Finally \(\sum_{a=0}^{\nu-1}n_a=d_*\) and
\(n_a+1=k+2(\nu-a)\); substituting yields the formula.
\end{proof}

\begin{corollary}\label{cor:unfold-unip}
In the setting of Proposition~\ref{prop:unipotent-covol}, let the
real unipotent group act on a measurable space and let \(\beta\ge0\)
be measurable with \(\sum_{\gamma\in\Gamma_c}\beta(\gamma y)=1\)
for every \(y\). Assume \(u\mapsto\beta(uy)\) is measurable for
each \(y\). Then \(\int_U\beta(uy)\dd u=\mathfrak v\).
\end{corollary}
\begin{proof}
Partition the real group into the left translates of the fundamental
domain constructed in Proposition~\ref{prop:unipotent-covol}.
The lattice is countable, since its elements are matrices over a
number field. Left invariance and monotone convergence from
Lemma~\ref{lem:integral-tools} turn the integral into the integral
of \(\sum_{\gamma\in\Gamma_c}\beta(\gamma uy)=1\) over that domain.
\end{proof}

\section{The unipotent character and Poincar\'e series}

\begin{lemma}\label{lem:unip-approx}
The rational group \(U(E^+)\) of Definition~\ref{def:unip} is dense
in its finite adelic group \(U_f\). Here \(E^+=E^{\mathfrak c}\)
and finite adelic groups have the restricted product topology,
with the integral subgroups prescribed at all but finitely many
primes. For each \(c\), the real group \(U_\infty\) in the scaled
coordinates of Proposition~\ref{prop:unipotent-covol} admits a
compact set \(C_c\) with \(U_\infty=\Gamma_c C_c\), writing the
rational group by its image at infinity.
\end{lemma}
\begin{proof}
The diagonal \(E\) is dense in its finite adeles. In an integral
basis the assertion reduces to \(\Q\). Given finitely many targets
in \(\Q_p\), specified modulo \(p^{a_p}\Z_p\), take an integer \(D\)
supported there, with \(D\) times every target \(p\)-integral and
\(a_p+v_p(D)\ge0\). By Lemma~\ref{lem:basic-arith} choose a single
integer \(b\) in the resulting residue classes modulo
\(p^{a_p+v_p(D)}\). Then \(b/D\) approximates all targets and is
integral away from these primes. For skew elements use the continuous
projection \(x\mapsto(x-\bar x)/2\); it is an identity on skew adeles.
The proof of Lemma~\ref{lem:heisenberg-layer} realizes the kernel over
the finite adeles by the independent coordinates \(l\) and
\(s=t+\omega lJ^{-1}l^*\), which is skew. These polynomial maps
and their inverses are continuous on the restricted products, since
their fixed coefficients are integral outside finitely many primes.
Thus rational kernel elements are dense. The section for the inner
group and multiplication identify the whole unipotent group with the
successive kernels as a topological space. Applying the same argument
to each proves density. For the second claim take closures of the
finitely many bounded parallelepipeds in the proof of
Proposition~\ref{prop:unipotent-covol}. Their compact product has
compact image in \(U_\infty\) under the coordinate reconstruction
in Lemma~\ref{lem:heisenberg-layer}, and contains the fundamental domain.
\end{proof}

\begin{proposition}\label{prop:full-phase}
In the rational basis put \(v_0=e_0,\ldots,v_{\nu-1}=e_{\nu-1},
v_\nu=z\) and \(\alpha(u)=\sum_{a<\nu}u(v_a,v_{a+1})\), with
\(u(i,j)\) denoting row \(i\), column \(j\).
On \(U_f\) the map
\(\chi_f(u)=\psi_f(-\alpha(u)/(g\varpi))\) is a continuous character.
At infinity let \(\alpha_\tau\) denote the same sum of entries in the
scaled basis and put
\[
       \chi_\infty(u)=
       \prod_{\tau\in\Phi}\exp\big(\alpha_\tau(u_\tau)
                                     -\overline{\alpha_\tau(u_\tau)}\big).
\]
For rational \(u\), \(\chi_\infty(u_\infty)\chi_f(u_f)=1\).
For \(u\in\Gamma_c\) both factors are \(1\).
\end{proposition}
\begin{proof}
In a product of two group elements, the entry \((e_a,e_{a+1})\)
for \(a+1<\nu\) is the sum of the corresponding entries of the factors: column \(e_{a+1}\) has support in the flag up to \(a+1\)
with diagonal one, and row \(e_a\) vanishes in earlier flag positions.
The same holds for \((e_{\nu-1},z)\): outside the flag the column
\(z\) equals the unit vector at \(z\) because the group acts as the
identity on the middle quotient and preserves the perpendicular
flag. This proves additivity. At every finite prime
\(\mathbf e_p\) is continuous, being constant on cosets of \(\Z_p\),
so Proposition~\ref{prop:additive} gives the finite character.
For a rational element, consecutive scale ratios along the frame
are \(1/\kappa_\tau\). Thus with \(a'=-\alpha(u)/(g\varpi)\),
\(\alpha_\tau(u_\tau)=-2\pi i\,\tau(a')\).
The product at infinity is \(\exp(-2\pi i\operatorname{Tr}a')\);
Proposition~\ref{prop:additive} gives its inverse at the finite places.
For \(\Gamma_c\), the vectors \(e_{a+1},z\) belong to \(\Lambda_c\),
and the \(e_a\)-coordinate of each \((u-I)v_{a+1}\) belongs to \(gO_E\).
Consequently \(\alpha(u)/g\in O_E\), and
Propositions~\ref{prop:different} and~\ref{prop:additive} give
\(\chi_f(u_f)=1\).
\end{proof}

\begin{definition}\label{def:adelic-row}
Let \(G_\infty\) be the product of the real unitary groups of \(Y_0\)
in the scaled coordinates, and \(G_f\) its finite adelic unitary
group. Let \(K^0_Y\) stabilize the initial lattice in
Definition~\ref{def:row-space}, and \(K_Y\) be its principal
level at \(g_0,g_1\), so \(K_Y\) is normal in \(K^0_Y\).
Write \(\Gamma\) for the rational unitary group and put
\(X=G_\infty\times G_f/K_Y\), with diagonal left action of \(\Gamma\).
For the class representative \(c\) in Definition~\ref{def:unip} let
\(h_c\) act by \(c\) on the positive middle block, and identically
on \(z\) and on the hyperbolic planes. Then
\(\Gamma_c=U(E^+)\cap h_cK_Yh_c^{-1}\) in finite coordinates, by
the definition of \(\Lambda_c\). Right translation \(R_k\) on
functions of \(X\) acts by \((R_k t)(x,hK_Y)=t(x,hkK_Y)\) for
\(k\in K^0_Y\), so only the class of \(k\) in \(\mathcal G\) matters.
\end{definition}

\begin{lemma}\label{lem:quotient-measure}
Give each component of \(X\) the same left and right invariant smooth
volume element on \(G_\infty\), whose positive normalization may be
chosen separately on each archimedean group. The action of \(\Gamma\)
is free, and for compact \(D_1,D_2\subset X\), only finitely many
\(\gamma\in\Gamma\) satisfy \(\gamma D_1\cap D_2\ne\varnothing\).
There is a measurable fundamental domain \(D\) for the left action.
For measurable nonnegative \(b\) on \(X\),
\[
           \int_D\sum_{\gamma\in\Gamma} b(\gamma y)\dd y
                            =\int_X b(y)\dd y .
\]
For invariant functions integrable on the quotient, integration over
\(D\) is independent of the choice of fundamental domain.
\end{lemma}
\begin{proof}
On a real unitary factor with Gram matrix \(J\), the real tangent space
at one is \(\{A:A^\dagger=-A\}\), \(A^\dagger=J^{-1}A^*J\);
\(A\mapsto(I+A)(I-A)^{-1}\) gives a chart near one with inverse
\((g-I)(g+I)^{-1}\). Translate an ordinary volume density on this
tangent space to obtain a smooth left invariant measure by integrating
the density in charts. On overlaps the change of variables preserves
the integral by the Jacobian transformation rule for densities. Right
translations also preserve it: conjugation on the tangent space has
determinant of modulus one. Indeed its complexification is conjugation
on all complex matrices, since each matrix is a sum of a skew-adjoint
matrix and \(i\) times a skew-adjoint matrix. On all matrices left
and right multiplication by \(g,g^{-1}\) have determinants
\((\det g)^r,(\det g)^{-r}\). The real determinant equals the complex
determinant of the complexified operator. Use the product of these
measures at the infinite places and counting measure on \(G_f/K_Y\).

A compact subset meets only finitely many finite components. At each
possible pair \(hK_Y,h'K_Y\), the finite part of a translating
\(\gamma\) belongs to \(h'K_Y h^{-1}\), whose entries lie in
\(a^{-1}O_p\) for all \(p\) and some fixed integer \(a\).
Indeed the initial lattice has fixed rational basis denominators,
the integral stabilizer has bounded entries relative to it, and
\(h,h'\) are integral with integral inverses at almost every prime.
The archimedean matrices of \(\gamma\) are bounded by the compact
sets of the two real projections and their inverses. In the original
rational basis this bounds every conjugate of every entry in
\(a^{-1}O_E\). This is a discrete lattice by Lemma~\ref{lem:skew-covol},
so only finitely many matrices occur. An element fixing a point has
identity archimedean matrix by cancellation and so is identity over
\(E\).

The space is second countable. The archimedean group is a subspace
of finitely many matrix spaces. There are countably many finite
components: \(hK^0_Y\) determines and is determined by its tuple
of local lattices, all of which occur among the finitely many
intermediate modules for some integer denominator bound \(a\), by
Lemma~\ref{lem:class-covol}. The additional level index is finite
by Proposition~\ref{prop:row-group}.
Each point has an open neighbourhood disjoint from its translates
by nonidentity elements: start inside a relatively compact
neighbourhood and shrink to avoid the finitely many possible
intersections, using freeness and the Hausdorff property.
Take a countable cover \(O_j\) by such opens and, successively, keep
in \(O_j\) only points outside all translates of the preceding opens.
Their union is a measurable fundamental domain. Its translates
partition \(X\); the displayed identity follows by additivity and
Lemma~\ref{lem:integral-tools}. Intersecting two fundamental domains
with each other's translates gives the asserted independence.
\end{proof}

\begin{theorem}[Poincar\'e series and their inner products]\label{thm:poincare}
For each of finitely many pairs \((c,i)\) choose an archimedean
continuous function \(W_{c,i}\) satisfying
\(W_{c,i}(u x)=\chi_\infty(u)W_{c,i}(x)\), \(u\in U_\infty\), and
supported in \(U_\infty C\) for some compact \(C\subset G_\infty\).
Such functions exist with \(W_{c,i}(1)\ne0\).
One can choose measurable nonnegative \(\beta_c\) such that
\(\phi_{c,i}=\beta_c W_{c,i}\) is continuous with compact support,
and its left \(\Gamma_c\)-sum equals \(W_{c,i}\).
Set
\[
 P_{c,i}(x,hK_Y)=\sum_{\gamma\in\Gamma}
     {\bf1}_{\gamma_f hK_Y=h_cK_Y}\,\phi_{c,i}(\gamma_\infty x).
\]
This is a locally finite sum defining a continuous \(\Gamma\)-invariant
function with compact support modulo \(\Gamma\).
For each \(c\) there is a choice with \(P_{c,i}(1,h_cK_Y)\ne0\).
With the inner product conjugate linear on the left, the formula
for its translates is
\[
 \begin{split}
 \langle R_k P_{c',j},P_{c,i}\rangle
 &=\sum_{\substack{\gamma\in\Gamma\\
             h_c^{-1}\gamma_f h_{c'}\in kK_Y}}
    \int_{G_\infty}\phi_{c,i}(x)
             \overline{\phi_{c',j}(\gamma_\infty^{-1}x)}\dd x
                     \qquad(k\in K_Y^0).
 \end{split}
\]
The same identity holds for measurable seeds on these finite cosets
when both periodizations exist absolutely almost everywhere, belong to
\(L^2(\Gamma\backslash X)\), and the sum on the right with the absolute
value of the integrand has finite integral.
\end{theorem}
\begin{proof}
Given a continuous compactly supported \(F\) on \(G_\infty\), set
\[
 W(x)=\int_{U_\infty}\overline{\chi_\infty(u)}F(ux)\dd u .
\]
The group \(U_\infty\) is closed by its flag equations. For \(x\)
in a compact neighbourhood the integration is over a fixed compact
subset of that subgroup, proving continuity by dominated convergence.
Right invariance in Proposition~\ref{prop:unipotent-covol} gives
the required equivariance. The support is contained in
\(U_\infty\operatorname{supp}F\), a closed set since the second
factor is compact. Choose \(F\ge0\), positive at \(1\), supported
near enough to \(1\) that \(\operatorname{Re}\chi_\infty(u)>1/2\)
on this support in \(U_\infty\); a continuous bump in a matrix
chart gives such \(F\). Then \(\operatorname{Re} W(1)>0\).

For each \(c\) choose a single compact \(C\) covering the finitely
many supports modulo \(U_\infty\). By Lemma~\ref{lem:unip-approx}
they lie in \(\Gamma_c C_c C\). Take a nonnegative compactly
supported continuous \(\xi\) equal to one on \(C_c C\); for instance
use a cutoff in a matrix norm, since the real groups are closed
matrix subsets with compact bounded closed subsets.
Set \(\beta_c(x)=\xi(x)/\sum_{\delta\in\Gamma_c}\xi(\delta_\infty x)\)
where the denominator is positive, and zero elsewhere. The sum is
locally finite by Lemma~\ref{lem:quotient-measure} at the finite
coset \(h_cK_Y\), and is at least one on the supports of the \(W\)'s.
Thus each \(\phi_{c,i}\) is continuous with compact support, also
at points where the denominator is zero because those points lie
outside the closed supports. Its periodization by \(\Gamma_c\)
gives \(W_{c,i}\), by Proposition~\ref{prop:full-phase} and the
definition of \(\beta_c\).

The seeds on \(X\) have compact support. The finiteness assertion of
Lemma~\ref{lem:quotient-measure} yields locally finite sums;
reindexing gives left invariance and the support descends to the
image of the compact seed support. On each orbit the sum can also
be taken over \(\Gamma_c\backslash\Gamma\) with \(W\) in place
of \(\phi\), by the preceding periodization.
For the nonzero value, fix a relatively compact identity neighbourhood
\(O\) at infinity and the compact \(C_c\subset U_\infty\) above.
There are only finitely many rational \(\gamma\) preserving \(h_cK_Y\)
at finite places whose infinite part lies in \(C_c\bar O\), again
by Lemma~\ref{lem:quotient-measure}. Shrink \(O\) so that \(C_c O\)
avoids all those outside \(U_\infty\). This is possible by compactness
of \(C_c\) and closedness of \(U_\infty\). Any such rational
\(\gamma\) with \(\gamma_\infty\in U_\infty O=\Gamma_c C_c O\)
can be translated by \(\Gamma_c\) into this finite list, where it must
now belong to \(U_\infty\). The rational flag equations then imply
\(\gamma\in\Gamma_c\). Taking the preceding \(F\) with support in
\(O\), only the identity coset in \(\Gamma_c\backslash\Gamma\)
contributes at \((1,h_cK_Y)\), giving the nonzero value \(W(1)\).

For the inner product, unfold the series for \(P_{c,i}\) on a
fundamental domain by Lemma~\ref{lem:quotient-measure}, leaving
its seed paired on \(X\) against the invariant function \(R_kP_{c',j}\).
Only \(h=h_c\) contributes. Expand the conjugated series and invert
its rational index, obtaining precisely the finite condition and
the integral displayed above. For compact seeds only finitely many
overlaps occur. The quotient integrals exist: their absolute square
integrals are finite by this same nonnegative unfolding applied to
the absolute values of the seeds and their finitely many overlaps.
For the measurable extension apply the unfolding identity first to
the absolute values of the two expanded series. The specified finite
absolute integral justifies interchanging all sums and integrals
in the identity by Lemma~\ref{lem:integral-tools}.
\end{proof}

\section{Central projections at the split level}

\begin{lemma}\label{lem:central-proj}
Let \(J\) be a finite group and write \([s]\) for the basis elements
of \(\C[J]\), with multiplication induced by the group law. For a
finite-dimensional irreducible complex representation \(\pi\), put
\(d_\pi=\dim\pi,\ \chi_\pi(s)=\operatorname{tr}\pi(s)\) and
\[
     e_\pi=\frac{d_\pi}{|J|}\sum_{s\in J}\chi_\pi(s^{-1})[s].
\]
In every unitary representation \(e_\pi\) acts as the orthogonal
projection onto the sum of its irreducible subrepresentations
equivalent to \(\pi\), in finite dimension. On an arbitrary unitary
representation it acts as a self-adjoint idempotent. Distinct
irreducible types give mutually orthogonal central idempotents.
The coefficient of \([1]\) in \(e_\pi\) is \(d_\pi^2/|J|\).
\end{lemma}
\begin{proof}
Averaging any positive inner product over the group makes it
invariant. In finite dimension, an invariant subspace has invariant
orthogonal complement, so induction decomposes each representation
into irreducibles. A commuting map between irreducibles is zero
or an isomorphism by its kernel and image. An endomorphism commuting
with an irreducible action is scalar by subtracting one of its
complex eigenvalues, which exists by Lemma~\ref{lem:place-fields}.
On the complex space of matrices \(T:\pi\to\pi'\) the average
of \(T\mapsto\pi'(s)T\pi(s)^{-1}\) over \(s\) is a projection onto
the commuting maps, by reindexing. Its trace equals their dimension.
On matrix units the trace of each summand operator is
\(\chi_{\pi'}(s)\chi_{\pi}(s^{-1})\). The average of these products
is thus one for equivalent irreducibles and zero otherwise.
The displayed \(e_\pi\) is central since traces are conjugacy
invariant. It acts as a scalar on each irreducible, and the trace
just calculated gives the scalar one for \(\pi\) and zero otherwise.
These action identities imply the corresponding products and squares
in the group algebra by applying them to the left regular action:
an element acting as zero there sends \([1]\) to zero and is itself
zero. Finally the involution \([s]^*=[s^{-1}]\) conjugating
coefficients fixes \(e_\pi\) since the character is the trace of
a unitary matrix. This proves the general projection assertion and
the claimed identity coefficient follows by evaluation at \(1\).
\end{proof}

\begin{proposition}\label{prop:integral-proj}
For \(e\) the sum of any subset of the distinct projections in
Lemma~\ref{lem:central-proj}, the coefficients of \(|J|e\) are
algebraic integers. For any \(\sigma\in{\rm Aut}(\C)\), applying
\(\sigma\) to the coefficients leaves a self-adjoint idempotent.
Suppose \(\rho\) is a finite-dimensional unitary action for a
positive definite rational Gram matrix \(S\), and each entrywise
transported action \(s\mapsto\sigma(\rho(s))\) is also \(S\)-unitary.
For any \(v\) and \(s\in J\) one has
\[
       \|\sigma(|J|\rho(e)\rho(s)v)\|_S^2
                                \le |J|^2\|\sigma v\|_S^2 .
\]
For a pair \((p,j)\) as in Definition~\ref{def:place-height},
let \(\mathcal O_j=\{z:|jz|_p\le1\}\). Any \(\mathcal O_j\)-module
\(L\) stable under \(\rho(J)\) is sent into itself by
\(|J|\rho(e)\rho(s)\).
\end{proposition}
\begin{proof}
Every \(s\) has order dividing \(|J|\) by counting its cyclic cosets.
Its matrix is diagonalizable, since it is annihilated by the
polynomial \(X^{|J|}-1\) with distinct roots; the projections on
the separate roots are given by polynomial interpolation.
Thus its trace is a sum of \(|J|\)-th roots of unity.
The coefficients in question are integral by Lemma~\ref{lem:rd}
and lie in the cyclotomic field \(\Q(\zeta_{|J|})\).
An automorphism permutes these roots by a fixed unit exponent, by
Lemma~\ref{lem:cyclo}. On this field it commutes with complex
conjugation, which raises the roots to their inverses. Hence the
idempotent and self-adjoint identities from Lemma~\ref{lem:central-proj}
persist under \(\sigma\). In the transported unitary action the
operator is an orthogonal projection, which is a contraction by
orthogonal decomposition into image and kernel. This gives the norm
bound. At \(p\) each root of unity has absolute value one by
multiplicativity. Hence the images of \(|J|e\)'s coefficients have
absolute value at most one by Lemma~\ref{lem:place-fields}.
Applying the resulting integral combination of group elements proves
the last assertion.
\end{proof}

\begin{definition}\label{def:poly-types}
At \(p=g_i\) write \(A=M_\omega\) from Theorem~\ref{thm:data},
and let \(\mathcal M_r\) be all monic degree \(r\) polynomials over
\(\mathbf F_p\). A label is a squarefree member with nonzero constant
coefficient. It is allowed if none of its roots in \(\mathbf F_p\)
has \(A\)-th power one. It is primitive if it is the minimal polynomial
of a generator of the multiplicative group of the field with \(p^r\)
elements. A good tuple of length \(f\) is a tuple of allowed labels
with at least one primitive member.
\end{definition}

\begin{lemma}\label{lem:count-labels}
For \(2\le r<p\), there are exactly \(\varphi(p^r-1)/r\) primitive
labels, all allowed. Among all \(p^r\) members of \(\mathcal M_r\),
the proportion failing to be an allowed label is at most
\((1+A+r)/p\). The proportion which are primitive is at least
\[
                          \frac{1}{2r(2r\log p+1)} .
\]
\end{lemma}
\begin{proof}
In an algebraic closure of \(\mathbf F_p\) (constructed in the
first paragraph of the proof of Lemma~\ref{lem:place-fields}) the
roots of \(X^{p^r}-X\) form a field: \(p\)-power Frobenius is
additive and multiplicative by the binomial theorem, giving closure
under addition, multiplication and inverses of nonzero elements.
The derivative is \(-1\), so the field has exactly \(p^r\) elements.
Its multiplicative group is cyclic. Indeed for each prime power in
the least common multiple \(L\) of the orders of its elements there
is an element of that prime-power order, by taking a suitable power.
Their product has order \(L\), since the orders are coprime and the
group abelian. All elements are roots of \(X^L-1\), so the group's
order is at most \(L\), proving cyclicity. Its generators, of which
there are \(\varphi(p^r-1)\) by coprime exponents, each generate
the entire field over \(\mathbf F_p\). The field has vector-space
dimension \(r\) by its cardinality. For such a generator \(\xi\)
the elements \(\xi,\xi^p,\ldots,\xi^{p^{r-1}}\) are distinct since
\(p^r-1\) cannot divide \(p^a-1\) for \(0<a<r\). All are roots
of the same minimal polynomial, and give precisely its \(r\) roots.
Thus each polynomial is counted \(r\) times. Since \(r>1\)
it has no root in \(\mathbf F_p\), and is squarefree and allowed.

A monic polynomial with a forbidden root or zero as a root is divisible
by \(X-u\) for one of at most \(1+A\) possible \(u\)'s. This excludes
at most \((1+A)p^{r-1}\) polynomials. A monic polynomial divisible
by the square of a nonconstant monic polynomial of degree \(d\) is
among at most \(p^d p^{r-2d}\) choices, for
\(1\le d\le \lfloor r/2\rfloor\). The sum is at most \(r p^{r-1}\);
division by \(p^r\) gives the claim.
For any integer \(n\ge2\) with distinct prime factors
\(p_1<\cdots<p_t\), reduction to prime powers and
Lemma~\ref{lem:basic-arith} give
\(\varphi(n)/n=\prod_i(1-p_i^{-1})\).
Since \(p_i\ge i+1\), this is at least \(1/(t+1)\), and
\(t\le\log n/\log 2\le2\log n\).
Take \(n=p^r-1\ge p^r/2\) to get the displayed lower bound.
\end{proof}

\begin{proposition}\label{prop:bad-types}
Let \(J_0={\rm GL}_r(\mathbf F_p)\), and suppose that for all
labels \(c\) of Definition~\ref{def:poly-types} there are distinct
irreducible representations \(\pi_c\) of \(J_0\) whose dimension
satisfies
\[
                 \frac{(\dim\pi_c)^2}{|J_0|}
                                  \ge (1-r/p)p^{-r}.
\]
In \(\C[J_0^f]\) form the sum \(e_{\rm good}\) of the tensor products
\(\bigotimes_w e_{\pi_{c_w}}\) over the good tuples. It is a
self-adjoint central idempotent and its complementary share at the
identity is bounded by
\[
 0\le [1](1-e_{\rm good})
   \le \frac{f(1+A+2r)}p+
                    \exp\!\left(-\frac{f}{2r(2r\log p+1)}\right).
\]
For both \(p=g_0,g_1\) and sufficiently large \(x\) this bound is
at most \(b^{-8}/2\).
\end{proposition}
\begin{proof}
The group algebra of the product is the tensor product of the group
algebras by their group-element bases. Distinct tuples give
orthogonal projections by Lemma~\ref{lem:central-proj}, and every
identity coefficient is the product of the individual projection
shares. The identity coefficient of any projection lies between
zero and one: the normalized trace of its action in the regular
representation equals its coefficient at one by the permutation
basis, and the trace of an orthogonal projection is its rank.
The share of \(e_{\rm good}\) is therefore at least
\((1-r/p)^f\) times the counting probability of a uniform good
tuple in \(\mathcal M_r^f\). By Lemma~\ref{lem:count-labels} the
chance of any forbidden coordinate is at most \(f(1+A+r)/p\),
and the chance of no primitive coordinate is at most the indicated
exponential by \(1-t\le\e^{-t}\). Finally
\(1-(1-r/p)^f\le fr/p\) by induction, yielding the bound.

For the scales, \(A=(Md)^{s+6},\,f\le x^{12},\,r\le3x^5\) and
\(s\le32\log x\) give \(f(1+A+2r)\ll_{M,d} x^{C_{M,d}}\)
for a constant independent of \(x\). Theorem~\ref{thm:data} gives
\(p\ge b^{10}\). The first summand is thus at most \(b^{-8}/4\)
eventually by the exponential comparison in Proposition~\ref{prop:thick}.
For the other, \(f=\varphi(M)3^{c_s}2^s\ge2^s=m^{22}\ge x^{11}\)
by Proposition~\ref{prop:places}, and \(\log p\le11\ell\).
Thus \(f/(2r(2r\log p+1))\gg x/\ell\), which exceeds
\(9\ell\) eventually because \(\ell=x^\epsilon,\ 0<\epsilon<1/100\).
This gives \(b^{-9}\le b^{-8}/4\) eventually.
\end{proof}

\begin{lemma}\label{lem:primitive-scale}
Let \(L=p^r-1\), and let \(c\) be primitive with root \(\xi\).
The product of its roots is
\(\xi^{L/(p-1)}\), a generator of \(\mathbf F_p^\times\).
If \(u\in\mathbf F_p^\times,\ u^A=1,\ A^2\mid p-1\), then
\(c^{[u]}(X)=u^r c(X/u)\) is primitive. If \(r\ge3\), \(p\ge5\)
and \(9\mid p-1\), the three polynomials obtained as \(u\) runs
over the cube roots of unity are distinct. In particular one of
them avoids any two prescribed labels.
\end{lemma}
\begin{proof}
The root description of Lemma~\ref{lem:count-labels} gives the
product, fixed by Frobenius and of order \(p-1\).
Write \(u=\xi^e\); since \(u^A=1\) and \(A\mid L\), the exponent
\(e\) is divisible by \(L/A\). Every prime \(b\mid L\) divides
\(L/A\): for \(b\mid A\) this uses \(A^2\mid L\), and otherwise
it follows immediately. Thus \(1+e\) is prime to \(L\), proving
\(u\xi\) a generator. The new polynomial is monic with this root
and degree \(r\), hence is its minimal polynomial. Now take \(A=3\).
The element \(\xi^3\) has order \(L/3\). This exceeds \(p^{r-1}-1\)
for \(p\ge5\), so \(\xi^3\) cannot be fixed by any Frobenius power
of exponent \(1,\ldots,r-1\). Equality of two scaled polynomials
would give \(u\xi=(v\xi)^{p^a}=v\xi^{p^a}\) for one such \(a\), or
for \(a=0\), by the list of roots. Cubing forces \(a=0\), so \(u=v\).
There are three such scalars by cyclicity and \(3\mid p-1\).
\end{proof}

\section{Counting definite classes in the ranks of the lifts}

\begin{lemma}\label{lem:finite-quotient-gens}
Let \(O=O_E\) from Theorem~\ref{thm:integer-ring}, \(h\ge1\),
and \(C\in M_h(O)\) have nonzero determinant. Every \(O\)-submodule
of \(A=C^{-1}O^h/O^h\) is generated by at most \(h\) elements.
In particular the number of such submodules is at most
\(|A|^h\), and \(|A|=|\operatorname{N}_{E/\Q}\det C|\).
\end{lemma}
\begin{proof}
Multiplication by \(C\) identifies \(A\) with \(O^h/CO^h\). On
an integral basis the resulting integer matrix has determinant
\(\operatorname{N}_{E/\Q}\det C\), as seen by evaluating at all
embeddings, as in Lemma~\ref{lem:rd}. Its absolute determinant is
the index by the lattice volume calculation in Lemma~\ref{lem:class-covol}.
Denote this integer by \(a\). It annihilates \(A\), by the adjugate
of that matrix. The ring \(O/aO\) is a product of \(O_p/aO_p\)
over \(p\mid a\) by Lemma~\ref{lem:basic-arith} in an integral basis.
By Theorem~\ref{thm:integer-ring} and its local proof each nonzero
factor is itself a product of quotients \(R/\pi^j R\) of discrete
valuation rings. The action of the coordinate idempotents splits
every \(O/aO\)-module and every submodule into these factors.
Over each \(R\) the module \(A\) is a quotient of \(R^h\).
The inverse image \(M\) in \(R^h\) of any of its submodules is
generated by at most \(h\) vectors: project \(M\) onto the first
coordinate. If the image is nonzero, choose an element of least
valuation there, which generates that ideal, and subtract a multiple
of its lift to kill the first coordinate. Induct on \(h\) for the
kernel, omitting the lift if the image is zero. Project the resulting
generators to the quotient. Pad each factor's list by zeros to length
\(h\), and join corresponding elements by the product decomposition.
These \(h\) elements generate since coefficients in the different
factors can be chosen independently in \(O/aO\). Counting their
ordered lists gives \(|A|^h\); when \(a=1\), the conclusion is immediate.
\end{proof}

\begin{theorem}\label{thm:class-number-bound}
For the standard definite unitary group on \(E^h,\ h\ge1\), let
\(\mathfrak C_h(J)\) denote its rational-by-level double cosets, with
level \(J\) in its finite adelic group. Let \(K_h^0\) be the stabilizer
of the standard lattice. Set \(D_*=dN_1,\ C_h=54fhD_*\).
Then
\[
             |\mathfrak C_h(K_h^0)|
                 \le \big((1+2C_h^2)C_h^2\big)^{2fh^2}.
\]
If \(J\) contains the principal level \(K_h(G)\) of matrices in
\(K_h^0\) congruent to the identity modulo a positive integer \(G\),
then \(|\mathfrak C_h(J)|\) is at most the right side times
\(G^{2fh^2}\).
\end{theorem}
\begin{proof}
The proofs of Lemmas~\ref{lem:trace-finite} and~\ref{lem:dual-den}
apply in any rank \(h\): a coset is the isometry class of its lattice
\(L\), which is integral for the hermitian form everywhere by local
isometry with \(O_p^h\), and \(D_*L^\#\subset L\) for its trace
pairing. Thus the trace dual has shortest nonzero length at least
\(1/D_*\). Lemma~\ref{lem:gaussian-cover} gives \(h\) vectors in \(L\)
of trace length at most \(27(2fh)D_*=C_h\), independent over \(E\)
by the real spanning argument of Proposition~\ref{prop:short-col}.
Let \(Y\in{\rm GL}_h(E)\) have these columns and \(C=Y^*Y\).
Its entries belong to \(O\), where \(O=O_E\), and all their
conjugates have modulus at most \(C_h^2\) by Cauchy--Schwarz
on the column vectors at each embedding. All embeddings commute
with the involution as in the proof of Proposition~\ref{prop:short-col}.
At each embedding \(0<\det(\sigma C)\le C_h^{2h}\) by
Lemma~\ref{lem:hadamard} and invertibility of \(Y\).
Consequently \(|\operatorname{N}_{E/\Q}\det C|\le C_h^{4fh}\).

Distinct algebraic integers differ by an element with nonzero integral
norm, and so some archimedean absolute value of their difference is at
least one. In the sup norm on \(\C^\Phi\), open balls of radius \(1/2\)
about the elements with conjugates bounded by \(C_h^2\) are therefore
disjoint. They lie in the ball of radius \(C_h^2+1/2\). Scaling real
\(2f\)-dimensional volume bounds their number by
\((1+2C_h^2)^{2f}\). There are at most the \(h^2\)-th power of this
number of possible matrices \(C\).

In the basis \(Y\), the original lattice lies between \(O^h\) and
\(C^{-1}O^h\): the columns of \(Y\) lie in \(L\), and pairing them
against every member of \(L\) gives integral coordinates.
For this fixed Gram matrix, Lemma~\ref{lem:finite-quotient-gens}
bounds the number of possible intermediate lattices by \(C_h^{4fh^2}\).
Knowing \(C\) and the intermediate lattice determines the rational
isometry class. In fact if \(Y,Y'\) have the same Gram matrix,
\(Y'Y^{-1}\) is unitary and takes one lattice to the other whenever
the intermediate lattices coincide. Multiplying the two counts yields
the asserted bound.

Reduction of \(K_h^0\) modulo \(G\) has kernel \(K_h(G)\), and its
image has at most \(|O/GO|^{h^2}=G^{2fh^2}\) elements by its matrix
entries and the integral basis of Theorem~\ref{thm:integer-ring}.
Representatives of these cosets refine each double coset into at most
this many classes. The natural map from classes at \(K_h(G)\) onto
those at \(J\) gives the final inequality.
\end{proof}

\begin{corollary}\label{cor:class-exponents}
Put \(t_1=5r,\ n_j=k+j(t_1+1)\) for \(0\le j\le m\). For all
sufficiently large \(x\), for all integers \(1\le h\le2n_m,\ G\ge1\)
and all levels \(J\supset K_h(G)\),
\[
           \log|\mathfrak C_h(J)|\le x^{24}(1+\log G).
\]
In particular a level containing \(K_h(gN)\) has at most
\(\exp(x^{25})\) classes. For the column level \(K_c\) in rank \(k\)
this bound holds as well.
\end{corollary}
\begin{proof}
Proposition~\ref{prop:thick} gives \(h\ll x^{11/2}\),
\(f\le x^{12}\) and \(D_*=dN_1\le d x^{12}\) eventually by the estimates
in Proposition~\ref{prop:level-prime}. Hence \(fh^2\ll x^{23}\) and
\(\log C_h\ll_{d,M}1+\log x\). Taking logarithms in
Theorem~\ref{thm:class-number-bound} yields
\(O_{d,M}(x^{23}(1+\log x+\log G))\).
The claimed inequality follows since \((1+\log x)/x\) tends to zero.
For \(G=gN\), Theorem~\ref{thm:data} and
Proposition~\ref{prop:level-prime} give
\(\log(gN)\le22\ell+30m\ell\ll x^{1/2+\epsilon}\), and eventually
\(1+\log(gN)\le x\). At \(g_0,g_1\) principal congruence modulo
\(gN\) implies the required identity reductions. At \(N\) it fixes
the last row modulo \(N\), so \(K_k(gN)\subset K_c\) by
Lemma~\ref{lem:split-level}.
\end{proof}

\section{Collisions of top places at the nonsplit primes}

\begin{definition}\label{def:collisions}
Use \(A,D,G,M,N_1,U\) of Propositions~\ref{prop:places}
and~\ref{prop:thick}, with \(A\) the subgroup of top places and \(D\)
the distinguished set reducing to \(-1\bmod M\).
For a prime \(p<x_1\) with \(\eta_d(p)\ne1\), put
\[
 h=-p\bmod N_1,\quad o=\operatorname{ord}_{M}(-p),\quad
 L=\operatorname{ord}_{N_1}(-p),\quad
 j_0=\min\{j\ge0:p^j\ge m\}.
\]
These classes are units by Proposition~\ref{prop:places}.
Call \(p\) short if \(p\equiv\pm1\bmod m\).
Call \(p\) served if \(\eta_d(p)=-1\), it is not short, and \(o\) is odd
with \(o\le j_0\). For \(V\subset G\) put
\[
 C_p(V)=
 \sum_{\tau\in A\setminus V}
 \ \sum_{\substack{e\ge1,\ e\ {\rm odd}\ {\rm or}\ p=d\\
              (-p^e)\tau\in A\setminus V}}\frac{a}{p^e}.
\]
The power in \(-p^e\) is taken before multiplying by \(-1\).
This nonnegative sum converges by comparison with \(Ua/(p-1)\).
\end{definition}

\begin{proposition}\label{prop:collision-exact}
For the primes of Definition~\ref{def:collisions} one has
\[
 C_p(\varnothing)=
 \begin{cases}
 Ua\,p^{-o}/(1-p^{-2o})&\eta_d(p)=-1\ \text{and }o\text{ odd},\\
 0&\text{otherwise}.
 \end{cases}
\]
For every unserved prime in this range and all sufficiently large \(x\),
\(m C_p(\varnothing)\le e_p/3\), with \(e_p\) as in
Proposition~\ref{prop:loss}.
\end{proposition}
\begin{proof}
For \(\tau\in A\), membership of \((-p^e)\tau\) in \(A\) is
equivalent to \(-p^e\equiv1\bmod M\).
If \(p=d\), then \(d\equiv M/2-1\bmod M\) and \(d^2\equiv1\bmod M\).
The two possible classes of \(-d^e\) are \(-1\) and \(1-M/2\);
neither is \(1\) since \(M\ge18480\). In the inert case only odd \(e\)
occur. For such \(e\) the class of \(-p^e\) equals \(h^e\) and it
lies in \(A\) exactly when \(o\mid e\), by division by the order.
If \(o\) is odd the exponents are \(o,3o,5o,\ldots\) for each
of the \(U\) starting points. If \(o\) is even there are none.
Summing the series proves the identity.

For \(p\) short, \(C_p(\varnothing)\le (4/3)Ua/p\), since
\(p\ge2,\ o\ge1\). For an unserved prime which is not short, the only
nonzero case has \(o\ge j_0+1\). Then \(p^o\ge pm\), giving
\(m C_p(\varnothing)\le (4/3)Ua/p\). In both cases \(U\le f\)
and \(\ell^2\ge4/3\) eventually, giving the claimed bound.
\end{proof}

\begin{lemma}\label{lem:inert-prime-fibres}
Let \(p\) have \(\eta_d(p)=-1\). In an algebraic closure of
\(\mathbf F_p\) fix a primitive \(N_1\)-th root \(\xi\) and a root
\(z\) of \(X^2-X+(1+d)/4\). Let \(\mathfrak p_\tau\) be the kernel
of the reduction map on \(O_E\) specified by
\(\omega\mapsto z,\ \zeta_{N_1}\mapsto\xi^\tau\), for \(\tau\in G\).
Write \(\overline{\mathfrak p}\) for the image of an ideal under
complex conjugation on \(E\). For \(\sigma,\tau\in G\),
\[
 \begin{aligned}
 \mathfrak p_\sigma=\mathfrak p_\tau
      &\ \Longleftrightarrow\ \sigma/\tau\in\langle p^2\rangle,\\
 \mathfrak p_\sigma=\overline{\mathfrak p_\tau}
      &\ \Longleftrightarrow\ \sigma/\tau\in(-p)\langle p^2\rangle .
 \end{aligned}
\]
All products and subgroups on the right are taken in \(G\).
The assertions hold for every choice of the two residue roots.
\end{lemma}
\begin{proof}
The prime \(p\) divides neither \(d\) nor \(N_1\) by
Proposition~\ref{prop:places}, hence is odd. Lemma~\ref{lem:reciprocity}
gives \((-d/p)=(p/d)=-1\), so the quadratic is irreducible over
\(\mathbf F_p\) and \(z^p=1-z\). The roots of \(X^{N_1}-1\)
are distinct by its derivative and form a cyclic group of order \(N_1\)
by the finite-subgroup argument of Lemma~\ref{lem:count-labels}.
In particular \(\xi\) exists. The factorization by cyclotomic
polynomials in Lemma~\ref{lem:cyclo} shows \(P_{N_1}(\xi)=0\):
a root of a factor \(P_b\), \(b\mid N_1\), has \(b\)-th power one,
forcing \(b=N_1\). The presentation and integral basis in
Theorem~\ref{thm:integer-ring} now give the reduction maps in
the statement. Their images are finite integral domains, hence fields,
since multiplication by a nonzero element is injective and then
surjective on a finite set.

Two such maps into \(\overline{\mathbf F}_p\) have the same kernel
exactly when their values on the generators differ by one common
power of Frobenius. Indeed their common residue field, if the kernels
agree, has \(p^a\) elements for some \(a\). Its image is the full
set of roots of \(X^{p^a}-X\), by the finite-field argument in
Lemma~\ref{lem:count-labels}. On a multiplicative generator the
\(a\) Frobenius images exhaust its distinct conjugates, as in the
proof there. This accounts for every embedding of that residue field.
Conversely a power of Frobenius preserves the kernel. To preserve
the specified image \(z\) the power has to be even. It sends
\(\xi^\tau\) to \(\xi^{p^{2i}\tau}\), proving the first assertion.
After conjugating in \(E\) the two images are \(1-z,\xi^{-\tau}\).
Returning to \(z\) requires an odd power of Frobenius, which gives the
second assertion. Nonnegative powers run through the entire cyclic
subgroups since \(p\) is a unit modulo \(N_1\).
\end{proof}

\begin{theorem}\label{thm:prime-marking}
For a served prime the orbits of \(\langle h\rangle\) that meet \(A\)
correspond bijectively to unordered pairs of distinct primes
\(\{\mathfrak p,\bar{\mathfrak p}\}\) above \(p\) met by the top places.
Their number is \(Uo/L\le U/m\).
Choose one prime in each such pair and call the resulting set
\(\Lambda\). Put \(V=\{\sigma\in G:\mathfrak p_\sigma\in\Lambda\}\).
Then \(V\) and \(hV\) partition these orbits, they avoid \(D\), and
\[
                   |A\cap V|=|A\cap hV|=U/2,\qquad C_p(V)=0 .
\]
More precisely the number of tops inducing each chosen prime is
\(L/(2o)\).
\end{theorem}
\begin{proof}
Write \(m=2^u,\ s=22u\), so \(2^s\mid N_1\) by
Proposition~\ref{prop:thick} and Definition~\ref{def:levels}.
Take \(x\) large enough that \(u\ge3\). The odd integer \(b'=-p\)
is neither \(1\) nor \(-1\bmod 2^u\). Of \(b'-1,b'+1\) one has exact
2-adic order one and the other has order at most \(u-1\). Hence
\(v_2((b')^2-1)\le u\). Odd squares are \(1\bmod8\), so factoring
successive differences of squares gives
\[
 v_2((b')^{2^j}-1)=v_2((b')^2-1)+j-1\qquad(j\ge1).
\]
The order of \(b'\bmod 2^s\) is a power \(2^j\) since the group of
odd residues has \(2^{s-1}\) elements, and cosets of a cyclic
subgroup all have its cardinality. The order exceeds one, and the
identity above yields \(j\ge s-u+1\ge u\). Reduction from \(N_1\)
to \(2^s\) proves \(m\mid L\).

By Lemma~\ref{lem:inert-prime-fibres} the orbit of each \(\tau\)
consists of the two fibres over \(\mathfrak p_\tau,\bar{\mathfrak
p_\tau}\). They are distinct because \(h\) has even order \(L\), so
\(h\notin\langle h^2\rangle\). Consecutive powers of \(h\) alternate
between the two fibres. An orbit through a top returns to \(A\) at
exactly the multiples of \(o\); there are \(L/o\) returns counting
the starting point in one period. Since \(o\) is odd, \(m\mid L/o\)
and half the returns are in each fibre. Dividing the total number
\(U\) of tops by \(L/o\) yields the asserted count and bound.
The reductions modulo \(M\) of such an orbit are in the odd-order
subgroup generated by \(h\bmod M\). That subgroup cannot contain
\(-1\), an element of order two, so the orbit avoids \(D\).
Finally each admissible exponent \(e\) contributing to
\(C_p(\varnothing)\) is odd. Multiplication by \(-p^e=h^e\) switches
fibres. Two tops outside \(V\) in the same orbit cannot therefore
give a term of \(C_p(V)\).
\end{proof}

\section{Elementary divisors and the size of transferred columns}

\begin{lemma}\label{lem:valuation-index}
Let \(O\) be a valuation ring, meaning a domain in whose fraction
field either \(z\) or \(z^{-1}\) belongs to \(O\) for each \(z\ne0\).
For \(c_1,\ldots,c_d\in O\setminus\{0\}\) put
\(P=\bigoplus_{i=1}^d c_i O\subset O^d\). Then
\[
          \mathcal I(O^d,P;z)\quad\Longleftrightarrow\quad
                           \prod_i c_i\mid z .
\]
For a square matrix \(Y\) over \(O\) and \(0\ne u\in O\), consider
\(P'=Y O^d+uO^d\). There exist \(a_i\in O\), possibly zero, with
\(a_i\mid a_{i+1}\), such that in some basis \(P'\) is
\(\bigoplus_i c_i O\), where \(c_i O=a_iO+uO\).
\end{lemma}
\begin{proof}
Apply Definition~\ref{def:det-index} to the \(d\) standard basis
vectors. In any recombination landing in \(P\), the \(i\)-th column
of the matrix is divisible by \(c_i\). This proves necessity.
For sufficiency, start with any finite list of vectors. In coordinate
\(i\), if all entries are zero move to the next coordinate. Otherwise
some nonzero entry divides every other entry: divisibility in \(O\)
is a total preorder by the defining property of a valuation ring.
Use the vector with that entry as pivot. Subtract suitable multiples
of it from the other vectors to clear their \(i\)-th coordinates.
Multiply the pivot vector by \(c_i\). Every \(i\)-th coordinate is now
in \(c_iO\); all previously obtained coordinate conditions persist
under these \(O\)-linear recombinations. The determinant at this
stage is \(c_i\), or one when the stage was skipped. Multiplying
the matrices proves sufficiency, including for the empty list.

In a nonzero matrix first bring an entry dividing all entries to the
upper left by row and column swaps. Clear its row and column by
invertible elementary operations. All entries of the remaining block
are still multiples of that entry. Induction gives a diagonal matrix
with the indicated divisibilities, padding with zeros for a zero
block. Column operations preserve the image module. Row operations
are a change of basis preserving \(uO^d\). In diagonal coordinates
the assertion follows since of \(a_i\) and \(u\) one divides the other.
\end{proof}

\begin{lemma}\label{lem:product-type}
For a domain \(O\), an \(O\)-module \(M\), \(u\in O\) and
endomorphisms \(Y_i\) of \(M\), suppose
\(\mathcal I(M,Y_iM+uM;z_i)\) for \(1\le i\le h\). Then
\[
 \mathcal I(M,Y_1\cdots Y_hM+uM;\ \prod_{i=1}^h z_i).
\]
No commutation between the operators is required.
\end{lemma}
\begin{proof}
The empty product gives \(M\). To combine operators \(Y,Z\), first
recombine a list into \(YM+uM\). Write each resulting vector as
\(Yv_i+uw_i\). Recombine the list of \(v_i\)'s into \(ZM+uM\).
Applied to the first list, this same matrix lands in \(YZM+uM\),
since \(Y(uM)\subset uM\). Determinants multiply, as in
Lemma~\ref{lem:det-index}. Induction now gives the claim.
\end{proof}

\begin{proposition}\label{prop:exact-normalization}
Put \(B=B_0'\) of Definition~\ref{def:heights} and
\(L_{\rm loc}=\sum_p e_p\log p\) with \(e_p\) of
Proposition~\ref{prop:loss}. Uniformly on \(t\le v\le t^2\),
\(L_{\rm loc}=o(B)\) as \(x\to\infty\). Put
\(\beta_p=\lceil G_p-e_p\rceil\), where \(G_p\) is the gain of
Proposition~\ref{prop:local-ask}. The positive rational
\(c=\prod_p p^{-\beta_p}\) satisfies
\[
           \log c\le L_{\rm loc}-\sum_pG_p\log p,\qquad
           |cp^{\beta_p}|_p=1\quad\hbox{for all }p .
\]
For \(\iota=(p,j)\) of Definition~\ref{def:place-height}, let
\(O_j=\{z\in\C:|jz|_p\le1\}\). Let \(M_j\) be an \(O_j\)-module of
columns, \(M'_j\subset M_j\) a submodule, and suppose some submodule
\(S_j\subset M_j\) satisfies \(\mathcal I(M_j,S_j;p^{a_j})\) with
an integer \(0\le a_j\le c_0e_p\). Let \(D_t:M'_j\to H_t\) be
\(O_j\)-linear maps into complex vector spaces with specified
\(O_j\)-submodules \(P_{j,t}\). The allowed \(t\)'s may depend on \(p\).
Assume \(p^{-\beta_p}D_tz\in P_{j,t}\) for all \(z\in S_j\cap M'_j\)
and every allowed \(t\).

Let \(W\) be a complex vector space, let \(Q_s:M'_j\to W\)
be \(O_j\)-linear and let \(u_t:W\to H_t\) be complex linear.
Suppose \(u_tQ_s z=T_{s,t}D_t z\) for linear operators
\(T_{s,t}\) on \(H_t\) preserving \(P_{j,t}\). For every family
\(y_i=cQ_{s_i}z_i,\ z_i\in M'_j\), in which at most \(b_T\)
distinct indices \(s_i\) occur, there is one square matrix \(A_j\)
over \(O_j\) with
\[
 |j(\det A_j)|_p\ge \exp(-b_T c_0 e_p\log p),\qquad
        u_t\Bigl(\sum_l (A_j)_{il}y_l\Bigr)\in P_{j,t}
                      \quad\hbox{for all }i,\ t .
\]
\end{proposition}
\begin{proof}
The proof of Proposition~\ref{prop:loss} bounds \(L_{\rm loc}/B\) by
\(15v(1+v)/\ell+480/\ell^2\). By Theorem~\ref{thm:data},
\(v^6\le\ell\) eventually and \(\ell=x^\epsilon\to\infty\).
The bound tends uniformly to zero. Since
\(-\beta_p=\lfloor e_p-G_p\rfloor\), this \(c\) is the rational
constructed there. Summing the floor inequality proves the logarithmic
bound, and all factors at primes different from \(p\) are \(p\)-adic
units. Thus \(cp^{\beta_p}\) is a unit of \(O_j\), giving
\(cD_t z\in P_{j,t}\) for \(z\in S_j\cap M'_j\).

Group the original list by the index \(s_i\). In each group apply
Definition~\ref{def:det-index} to the original \(z_i\)'s. Every
recombined vector belongs both to \(S_j\) and to \(M'_j\), the latter
by linearity. Its normalized image under \(T_{s,t}D_t\) lies in
\(P_{j,t}\) for every allowed \(t\). The block diagonal matrix of these
recombinations, conjugated by the permutation that grouped the list,
has determinant dividing \(p^{a_j b_T}\). Since elements of \(O_j\)
have absolute value at most one after \(j\), its determinant satisfies
the asserted inequality. This also proves that smaller integers
than \(\beta_p\) need no separate divisibility estimate: multiply
\(p^{-\beta_p}D_tz\) by the corresponding nonnegative power of \(p\).
\end{proof}

\begin{proposition}\label{prop:raw-to-short}
Let \(D_0'\) send the ramified columns to vector-valued functions on
a definite quotient of rank at most \(2n_m\), at a level containing
\(K_h(gN)\) as in Corollary~\ref{cor:class-exponents}. Its squared
norm in archimedean coordinates is the sum of squared norms at class
representatives with rational weights in \((0,1]\).
For every short column \(z\) of Corollary~\ref{cor:short-ram} assume
that the values of \(D_0'z\), after any \(\sigma\in{\rm Aut}(\C)\),
have squared norm at most
\[
             \exp\bigl(h_0+2B+2\sum_pG_p\log p\bigr).
\]
Let \(e=\sum_{w\in\mathcal G}a_w[w]\) be a sum of the orthogonal
central idempotents of Lemma~\ref{lem:central-proj}. Assume that in these coordinates the translated values satisfy
\[
 (u_0Q_s z)(y)=|\mathcal G|\sum_w a_w I_{w,s,y}
                                      (D_0'z)(y_{w,s}) ,
\]
where the points \(y_{w,s}\) may depend on \(y\), and
\(I_{w,s,y}\) has operator norm at most one after every \(\sigma\).
Then for \(c\) of Proposition~\ref{prop:exact-normalization},
eventually, uniformly in \(v,s,z,\sigma\), the squared norm of
\(\sigma(u_0cQ_s z)\) is at most
\(\exp(H-m\mathcal R+6B)\).
\end{proposition}
\begin{proof}
By Proposition~\ref{prop:integral-proj}, \(\sigma(e)\) is a
self-adjoint idempotent in the left regular unitary representation.
Applying it to the basis vector at the identity gives
\(\sum_w|\sigma a_w|^2\le1\). Hence
\(\sum_w|\sigma a_w|\le\sqrt{|\mathcal G|}\) by Cauchy--Schwarz.
The sum of translates increases the squared pointwise bound by at most
\(|\mathcal G|^3\). Summation over class representatives costs at most
\(\exp(x^{25})\) by Corollary~\ref{cor:class-exponents}.
From Proposition~\ref{prop:row-group}, each of the \(f\) split factors
at \(g_i\) has at most \(g_i^{r^2}\) matrices. By
Theorem~\ref{thm:data}, \(3\log|\mathcal G|\le66fr^2\ell\).
Together with the bound on \(r\) in Proposition~\ref{prop:thick},
this gives \(x^{25}+3\log|\mathcal G|\le B\) for large \(x\):
\(a\ge x\exp(x^\epsilon)\) dominates every fixed power of \(x\)
by the exponential series. Every \(\sigma\) fixes \(c\), so its squared
scalar adds \(2\log c\le2B-2\sum_pG_p\log p\), using
Proposition~\ref{prop:exact-normalization} and \(L_{\rm loc}\le B\).
The gain sum cancels. Finally \(h_0\le H-m\mathcal R+B\)
by Proposition~\ref{prop:height-cost}. Adding the five contributions
gives the stated bound.
\end{proof}

\begin{corollary}\label{cor:asymptotic-volume}
In the setting of Corollary~\ref{cor:row-test}, suppose the selected
rows are \(cQ_s z\) of Propositions~\ref{prop:exact-normalization}
and~\ref{prop:raw-to-short}, with their hypotheses for every required
embedding and twist and with the same integral modules as in
Corollary~\ref{cor:completion-height}. Retain the other hypotheses
of Propositions~\ref{prop:selected-family} and~\ref{prop:long-criterion}
on the row maps, completions and factorizations. Then
\[
                m\mathcal R\le (2+o(1))fa\ell^4 .
\]
The \(o(1)\) is uniform for \(t\le v\le t^2\) for fixed \(M,d,\epsilon\).
In particular these estimates, if established for the automorphic
maps for every such choice of parameters, imply Dirichlet GRH.
\end{corollary}
\begin{proof}
Write \(T_0=m\mathcal R,\ E_0=B\ell,\ d_0=\dim W\) and
\(n\) for the number of independent rows selected. The selection
bound and Proposition~\ref{prop:exact-normalization} give
\[
 \delta\le3d_0\log(1+b^2)L_{\rm loc}=o(d_0 E_0),
\]
since \(\log(1+b^2)\le3\ell\) eventually. Put
\(V=n(H-T_0+6B)\), allowed by Proposition~\ref{prop:raw-to-short}
and Lemma~\ref{lem:hadamard}. In the proof of
Proposition~\ref{prop:long-criterion} the volume ratio has height
at most \(V+2\delta+d_0 b\); its inverse square has height at most
\(2d_0(2E_0-H)\). Proposition~\ref{prop:height-product} applied to
their product, with the first ratio squared, yields
\(d_0(H-2E_0)\le V+2\delta+d_0 b\).
The lost dimensions satisfy \(0\le d_0-n\le4d_0/b^2\), and
\(-H+T_0-6B\le b^2x^{99}\) by Proposition~\ref{prop:height-cost}.
Consequently
\[
        T_0\le 2E_0+6B+b+4x^{99}+2\delta/d_0 .
\]
Here \(B/E_0=1/\ell\), \(b/E_0\le1/\ell^4\) since \(a\ge b,\ f\ge1\),
and \(x^{99}/E_0\to0\) by exponential growth of \(b\).
This proves the uniform estimate. Eventually its right side is
at most \(38fa\ell^4\); Theorem~\ref{thm:geometric-reduction}
gives the final implication.
\end{proof}

\section{The diagrams in the chain of lifts}

\begin{definition}\label{def:stage-diagrams}
Use \(h_2,h,t_1,w_j\) of Definition~\ref{def:heights} and
\(n_i=k+i(t_1+1)\) of Corollary~\ref{cor:class-exponents}.
For \(0\le j<m\) and \(\tau\in G\) define a diagram by its
columns of height strictly less than \(n=n_{j+1}\).
There are \(X_\tau=a{\bf1}_A(\tau)\) columns of height one,
\(d_\tau=w_1(\tau)-\lfloor n_1/2\rfloor\) of height \(k\),
and \(h-(t_1+1)\) of height \(n_i\) for each \(1\le i\le j\).
All other multiplicities are zero. Denote the number of columns
of height greater than one by \(r_\tau\), their number of boxes
by \(b_\tau\), and put \(N_j=\sum_\tau(X_\tau+b_\tau)\).
For \(V\subset G\) put \(z_\tau=X_\tau{\bf1}_{\tau\notin V}\).
At a prime \(p\) of Definition~\ref{def:collisions} write
\(I_p=\{e\ge1:e\text{ odd or }p=d\}\),
\(\omega_p=(1-\eta_d(p))/2,\ E_j=\lfloor\log N_j/\log p\rfloor\).
Set \(I_p^0=I_p\) for \(p\ne d\) and \(I_d^0=I_d\cup\{0\}\).
For \(e\in I_p^0,\ \sigma=(-p^e)\tau,\ s_e=p^{-e}\) put
\[
 K(\tau,e)=\min(z_\tau s_e,r_\sigma+z_\sigma)
                  +r_\tau s_e+\min(b_\tau s_e,b_\sigma).
\]
Define the number
\[
 J_p(j,V)=\omega_p\sum_{\tau\in G}\left(
  \sum_{e\in I_p}(X_\tau+b_\tau)s_e
       -\sum_{\substack{e\in I_p^0\\ e\le E_j}}K(\tau,e)\right)
                         -fn^3(E_j+1).
\]
Only positive exponents occur in the geometric sum. At \(p=d\) the
second sum also contains \(K(\tau,0)\), for the pairing with
\(-\tau\). Absolute convergence follows from geometric summation. Its leading coefficient per box is
\(\omega_p\sum_{e\in I_p}p^{-e}\), equal to \(p/(p^2-1)\) at an
inert prime and \(1/(2(p-1))\) at \(p=d\).
\end{definition}

\begin{lemma}\label{lem:diagram-size}
Eventually, uniformly in \(j<m,\tau,p<x_1\) in the preceding ranges,
all displayed multiplicities are nonnegative integers,
\(n^3\le x^{49}\), and
\[
 w(\tau)\le r_\tau\le w(\tau)+x^{49},\qquad
 kw(\tau)\le b_\tau\le kw(\tau)+x^{49},\qquad
 1\le E_j\le 3\ell/\log 2 .
\]
Writing \(u_p=a/(mp)\), one has \(x^{49}=o(u_p)\) uniformly and
\(r_\tau/p\le u_p\) eventually.
\end{lemma}
\begin{proof}
Here \(k>1,\ t_1+1<h\) eventually by Proposition~\ref{prop:thick}.
By Definition~\ref{def:heights}, \(w_1(\tau)\) is an integer at least
\(w(\tau)+h+h_2\), and \(h>\lfloor n_1/2\rfloor\) eventually.
Thus \(d_\tau\ge w(\tau)\), while
\(d_\tau\le w(\tau)+x^{42}\) by Proposition~\ref{prop:height-cost}.
Besides these columns there are exactly \(j(h-t_1-1)\) higher ones,
of height at most \(n_m\). For the upper errors in boxes and columns
it suffices to bound \(kx^{42}+m n_m h\).
Here \(k=O(x^5),\,m=O(\sqrt x),\, n_m=O(x^{11/2}),\,h=O(x^{40})\)
by Proposition~\ref{prop:thick}. This sum and \(n_m^3\) are at
most \(x^{49}\) for large \(x\). Also \(Ua\le N_j\le
 f(a+k\max_\tau w(\tau)+x^{49})\), and \(\max w=a/x\) since
\(x_1\ge x^{10^6}\). The bounds \(f\le x^{12}\), \(a\le2xb\)
in Proposition~\ref{prop:height-cost} and its proof show
\(\log N_j\le 3\ell\). Since \(N_j\ge a\ge x_1>p\), this proves
the claims about \(E_j\).

From \(a\ge x\exp(\ell),\ p<x_1=\exp(v)\) and \(v\le
(10^6\log x)^2=o(\ell)\) it follows that \(u_p\) dominates every
fixed power of \(x\), uniformly. Finally \(w(\tau)\le a/x\) and
\(m/x\to0\), so \(r_\tau\le a/x+x^{49}\le a/m\) eventually.
\end{proof}

\begin{theorem}\label{thm:diagram-gain}
For all sufficiently large \(x\), every prime in
Definition~\ref{def:collisions}, every \(j<m\), and every subset
\(V\subset G\),
\[
 J_p(j,V)\ \ge\ \frac{G_p}{m}-C_p(V)
       +\omega_p\,|A\cap V|\,\frac{a}{x_1}
       -30f\ell\,\frac{a}{mp}.
\]
\end{theorem}
\begin{proof}
Set \(\Delta=x^{49}\) and use \(u_p\) of Lemma~\ref{lem:diagram-size}.
Fix \(\tau,e\) with \(e\in I_p,\ e\le E_j\), and abbreviate
\(s=s_e,\ w=w(\tau),\ w'=w(\sigma)\).
The contribution from the tall boxes satisfies
\[
 b_\tau s-\min(b_\tau s,b_\sigma)=(b_\tau s-b_\sigma)_+
                           \ge k(ws-w')_+-\Delta ,
\]
since \(b_\tau\ge kw,\ b_\sigma\le kw'+\Delta\), and
\((t-\Delta)_+\ge t_+-\Delta\). Also \(r_\tau s\le u_p\).

For the height-one columns, if \(\tau\notin A\) both \(X_\tau\)
and \(z_\tau\) vanish. If \(\tau\in A\cap V\), their contribution
after subtraction of the first minimum in \(K\) is \(as\), equal to
\((as-w')_++\min(as,w')\). At \(e=1\), \(p<x_1\) and \(w'\ge a/x_1\),
so the minimum is at least \(a/x_1\). If \(\tau\in A\setminus V\)
and \(\sigma\notin A\setminus V\), the contribution is
\((as-r_\sigma)_+\ge (as-w')_+-\Delta\).
In the remaining case both places are in \(A\setminus V\). The
contribution is nonnegative, hence at least \((as-w')_+-as\).
Summing this last loss over the indicated pairs gives at most
\(C_p(V)\) after weighting by \(0<\omega_p\le1\).

For \(e>E_j\), the unsubtracted \(X_\tau s\) and \(b_\tau s\)
already bound \({\bf1}_A(\tau)(as-w')_+\) and \(k(ws-w')_+\);
there is no extra \(e=1\) term since \(E_j\ge1\).
Over all \(e,\tau\) these two positive parts with coefficient
\(\omega_p\) give exactly \(G_p/m\) by
Proposition~\ref{prop:local-ask}. At most \(f E_j\) pairs have the
error \(2\Delta+u_p\le3u_p\).
At \(p=d\) there is also the subtracted term at \(e=0\).
Here \(d\) is fixed and \(r_\tau\le a/x+\Delta\le u_p\) eventually,
since \(md/x\to0\) and \(\Delta=o(u_p)\). If \(z_\tau>0\) then
\(\tau\in A\), so \(-\tau\notin A\) and \(z_{-\tau}=0\).
Hence the first minimum in \(K(\tau,0)\) is at most \(u_p\).
The places \(\tau,-\tau\) cannot both reduce to \(-1\bmod M\), so
one of them is outside \(D\) and has weight \(a/x_1\). Thus
\(\min(b_\tau,b_{-\tau})\le ka/x_1+\Delta\le u_p\), using
\(kmd/x_1\to0\). The total extra cost is at most \(3fu_p\).
The final subtraction in
Definition~\ref{def:stage-diagrams} is at most \(f u_p(E_j+1)\)
by Lemma~\ref{lem:diagram-size}. The total of these errors is at most
\(4f(E_j+1)u_p\le30f\ell u_p\), since \(E_j\le3\ell/\log2\),
\(\log2>1/2\), and \(\ell\to\infty\). This proves the inequality.
\end{proof}

\begin{proposition}\label{prop:marked-cost}
For a served prime use \(V,\Lambda\) of Theorem~\ref{thm:prime-marking}.
Put \(\delta_p=j_0 f(n_m+j_0)\). Eventually for every \(j<m\),
\[
     |\Lambda|\delta_p+j_0\sum_{\tau\in V}r_\tau
          \ \le\ |A\cap V|\,a/x_1+3f\,a/(mp).
\]
Consequently the sum over \(0\le j<m\) of
\(J_p(j,V)-|\Lambda|\delta_p-j_0\sum_{\tau\in V}r_\tau\)
is at least \(G_p-(30\ell+3)fa/p\).
For an unserved prime the sum of \(J_p(j,\varnothing)\) is at least
\(G_p-e_p/3-30f\ell a/p\).
\end{proposition}
\begin{proof}
Since \(j_0\le1+\log m/\log2\), the quantity
\((|\Lambda|/f)\delta_p\) is bounded by a fixed power of \(x\),
using \(|\Lambda|\le f,\ f\le x^{12}\) and the bound on \(n_m\)
in Lemma~\ref{lem:diagram-size}. The same is true of \(j_0 x^{49}\).
Both are at most \(u_p=a/(mp)\) eventually by the proof of that lemma.
All \(\tau\in V\) avoid \(D\) by Theorem~\ref{thm:prime-marking},
so their common nominal weight is \(a/x_1\).
If \(p<m\), then \(mpj_0\le m^2j_0\le x_1\) eventually by
\(m^2\le4x,\ x_1\ge x^{10^6}\). Thus \(j_0 a/x_1\le u_p\).
If \(p\ge m\), then \(j_0=1\) and servedness forces \(o=1\).
The element \(h=-p\bmod N_1\) then belongs to \(A\),
so every marked orbit is contained in \(A\). In this case
\(j_0\sum_{\tau\in V}w(\tau)=|A\cap V|a/x_1\).
The bound on \(r_\tau\) in Lemma~\ref{lem:diagram-size} now
gives the claimed cost in both cases. Subtract it from
Theorem~\ref{thm:diagram-gain}. At a served prime \(\omega_p=1\)
and \(C_p(V)=0\), so the terms involving \(|A\cap V|\) cancel.
In the unserved case apply that theorem directly with
\(V=\varnothing\) and use Proposition~\ref{prop:collision-exact}.
\end{proof}

\section{Changing frames to annihilate Frobenius pairings}

\begin{lemma}\label{lem:congruences-frame}
Let \(O\) be a valuation ring as in Lemma~\ref{lem:valuation-index}.
Let \(L_1,\ldots,L_w:O^h\to O\) be linear, and let
\(\alpha_1,\ldots,\alpha_w\in O\setminus\{0\}\). There is a matrix
\(M\in M_h(O)\) with \(\det M\mid\prod_i\alpha_i\) such that every
row \(v\) satisfies \(L_i(v)\in\alpha_i O\) for every \(i\).
\end{lemma}
\begin{proof}
Map the \(h\) standard basis vectors of \(O^h\) into \(O^w\) by the
tuple of the \(L_i\). Apply Lemma~\ref{lem:valuation-index} to this
list and \(P=\bigoplus_i\alpha_iO\). Its recombination matrix has
the stated determinant, and its rows give the desired combinations
of standard basis vectors. When there are no constraints the
identity matrix suffices.
\end{proof}

\begin{definition}\label{def:frame-e}
Fix \(j,p,V\) of Definition~\ref{def:stage-diagrams}. Let
\(\kappa_{\tau,h}\) be the multiplicities there, \(1\le h<n=n_{j+1}\).
Fix any positive integer \(d'\), let \(K_p\) be the valued field of
Lemma~\ref{lem:place-fields} with valuation ring \(O\), and take
arbitrary vectors \(y_{\tau,i}\in O^{d'}\) for \(1\le i<n\).
A frame of height \(h\) at \(\tau\) uses the first \(h\) such vectors.
Let \(\theta\in O\setminus\{0\}\) with \(|\theta|_p<1\) and set
\(\bar O=O/(p,\theta)\).
For \(e\ge0\) define the pairing over \(\bar O\) by
\[
              B_e(u,v)=\sum_{a=1}^{d'}u_a^{p^e}v_a .
\]
It is additive in each variable and semilinear in the first,
by raising sums to \(p^e\) in characteristic \(p\).
\end{definition}

\begin{theorem}\label{thm:null-frames}
With the preceding data there are matrices \(M_{\tau,h}\in M_h(O)\)
for each \(\kappa_{\tau,h}>0\). In the corresponding changed frame
write \(\tilde y_{\tau,h,i}=\sum_{l\le h}(M_{\tau,h})_{il}y_{\tau,l}\).
Their determinants satisfy
\[
 \prod_{\tau,h}|\det M_{\tau,h}|_p^{\kappa_{\tau,h}}
       \ \ge\ |\theta|_p^{\sum_{e\in I_p^0,\ e\le E_j}
                                      \sum_{\tau\in G}K(\tau,e)} .
\]
For each exponent in the sum, \(\sigma=(-p^e)\tau\), and all rows
of contributing source and target frames of heights \(h,h'\), their
reductions to \(\bar O\) satisfy
\[
 B_e(\tilde y_{\tau,h,i},\tilde y_{\sigma,h',i'})=0
\]
except possibly when \(h=1\) and either \(\tau\in V\), or
\(h'=1,\ \sigma\in V\).
\end{theorem}
\begin{proof}
For each of the finitely many directed pairs \((\tau,e)\), write
\(Q=p^e\) and choose \(\theta_e^Q=\theta\) in \(K_p\). The root lies
in \(O\), since \(|\theta_e|_p=|\theta|_p^{1/Q}\). For each component
of \(y_{\sigma,1}\) also choose a \(Q\)-th root in \(O\).
We impose the following linear congruences on the rows \(v\in O^h\)
of each change of basis. Their costs below are the sums of the
exponents of \(|\theta|_p\) in the divisors of
Lemma~\ref{lem:congruences-frame}, counted with the column
multiplicities. Multiplicities of height zero do not occur.

First consider a source of height one with \(\tau\notin V\).
A target is relevant unless it has height one and \(\sigma\in V\).
One choice is to require the single coordinate of the change of
source frame to lie in \(\theta_e O\), at cost \(z_\tau/Q\).
Its \(Q\)-th power vanishes modulo \((p,\theta)\). Another choice
requires of each relevant target frame of height \(h'\) the single
linear congruence
\[
 \sum_a y_{\tau,1,a}^Q\sum_{l\le h'}v_l y_{\sigma,l,a}\in\theta O .
\]
This costs \(r_\sigma+z_\sigma\). It kills the pairing also after any
change of frame at the height-one source, by semilinearity.
Use the cheaper choice. If there is no such source, impose nothing;
the minimum then vanishes.

Next consider sources of height \(h>1\) and a target of height one,
including a marked target. On each source frame impose the single
linear congruence
\[
     \sum_a\Bigl(\sum_{l\le h}v_l y_{\tau,l,a}\Bigr)
                   y_{\sigma,1,a}^{1/Q}\in\theta_e O .
\]
Taking \(Q\)-th powers in \(\bar O\) annihilates its pairing with
\(y_{\sigma,1}\), hence with any multiple of this target.
The cost is at most \(r_\tau/Q\); when no target of this height
occurs the congruence can be omitted.

Finally consider both heights greater than one. Requiring every
coordinate of each tall source row to belong to \(\theta_e O\)
has cost \(b_\tau/Q\). Requiring every coordinate of each tall
target row to belong to \(\theta O\) has cost \(b_\sigma\).
Either condition kills all these pairings; choose the less costly.
Thus the total for this directed pair is at most \(K(\tau,e)\)
of Definition~\ref{def:stage-diagrams}.

For each frame collect the congruences imposed by all directed pairs
at either end, including both sets if the endpoints coincide.
Apply Lemma~\ref{lem:congruences-frame} once to their full list.
Every resulting determinant divides the product of its congruence
scalars and is nonzero. Taking absolute values reverses the
divisibility bound because \(|x|_p\le1\) on \(O\). Multiplication
with the stated multiplicities proves the inequality, and each
imposed vanishing holds simultaneously.
\end{proof}

\section{Integral coordinates for the differential minors}

\begin{lemma}[Straightening flag minors]\label{lem:flag-straighten}
For each \(\tau\) take a matrix of independent variables
\(Y_{\tau,i,a}\), \(1\le i<n,\ 1\le a\le d'\). A minor for a column
of height \(h\) is
\([a_1,\ldots,a_h]_\tau=\det(Y_{\tau,i,a_l})_{1\le i,l\le h}\).
Fix the number of columns of every height at each \(\tau\) and order
the columns there by weakly decreasing height. Fill each box with an
integer in \([1,d']\). A filling is standard if its entries strictly
increase down columns and weakly increase from left to right along
each row. Write \(P_T\) for the product of the column minors of
a filling \(T\). The \(P_T\) for standard \(T\) are a basis over
\(\Z\) of the integer span of all such products. After passage to
any field they remain linearly independent. Over \(K_p\), if a
polynomial \(P=\sum_T c_TP_T\) has coefficients of valuation at
least \(v_0\in\R\), then \(v_p(c_T)\ge v_0\) for each \(T\).
Here \(v_p(z)=-\log|z|_p/\log p\), with \(v_p(0)=+\infty\).
\end{lemma}
\begin{proof}
A repeated entry in a column gives zero. Otherwise sort each
column, multiplying by its permutation sign. If an adjacent pair
\(I=(i_1,\ldots,i_h), J=(j_1,\ldots,j_k)\), \(h\ge k\), violates
the row condition, choose \(t\le k\) with \(i_t>j_t\). Use the
ordered list \(D=(i_t,\ldots,i_h,j_1,\ldots,j_t)\) of length \(h+1\).
The signed sum over its splittings into increasing sublists
\(D_S,D_{S^c}\), with \(|S|=h-t+1\), satisfies
\[
 \sum_S {\rm sgn}(S,S^c)
 [i_1,\ldots,i_{t-1},D_S]_\tau
 [D_{S^c},j_{t+1},\ldots,j_k]_\tau=0 .
\]
Here order within each sublist is inherited from \(D\), and the sign
is that of the permutation listing \(S\) before \(S^c\).
As a function of the \(h+1\) column vectors of the variables with
labels in \(D\), the sum is alternating and multilinear: interchanging
two labels either changes one determinant's sign or pairs two shuffle
terms with opposite signs. The second determinant uses the first
\(k\) coordinates of its vectors, which have \(h\) coordinates.
Expanding all \(h+1\) vectors in the \(h\) basis vectors gives zero
by alternation. This polynomial identity remains valid when entries
among the fixed columns and \(D\) agree.

The term taking \(i_t,\ldots,i_h\) on the left is the initial product,
with coefficient one. Every other nonzero term, after sorting the
columns, has a lexicographically smaller left column: it replaces
one or more labels at least \(i_t\) by labels at most \(j_t<i_t\).
Rearrange the equality to express the initial product in these
terms. All earlier columns stay fixed. Order the finite set of
column words lexicographically, fixing also an order on the places.
Induction on this order proves spanning over \(\Z\).

Order monomials lexicographically by taking variables first by place,
then by increasing row slot \(i\), then by increasing label \(a\),
with an earlier variable having priority for higher exponents.
For a strictly increasing column, the leading monomial is the diagonal
product with coefficient one: in each successive row slot choose
the smallest remaining entry. The leading monomial of a standard
product records for every row of every place the multiset of entries.
The weak increase along rows recovers their order uniquely. Therefore
different standard products have distinct leading monomials. In any
linear relation a largest leading monomial with nonzero scalar
coefficient could not be cancelled. This proves independence over
any field and over \(\Z\). For the valuation claim, recover the
coefficient of the largest leading monomial of the expansion first,
subtract its scalar times the corresponding integral minor product,
and continue in decreasing order. The ultrametric inequality in
Lemma~\ref{lem:place-fields} preserves the lower bound at every step.
\end{proof}

\begin{lemma}\label{lem:poly-test}
For a polynomial \(P\) in finitely many variables over the valued
field \(K_p\), there is a point in its valuation ring \(O\) in each
coordinate at which \(|P|_p\) equals the largest absolute value of
its coefficients.
\end{lemma}
\begin{proof}
For a nonzero polynomial divide by a coefficient of largest absolute
value, obtaining an integral polynomial with nonzero reduction modulo
the maximal ideal \(\{x:|x|_p<1\}\). The residue field is algebraically
closed: lift a monic nonconstant polynomial to \(O\), split it in
\(K_p\) by Lemma~\ref{lem:place-fields}, and note that every root has
absolute value at most one, since a larger root would make the
leading term strictly larger than all other terms. Reducing gives the
factorization. Such a field is infinite, since in a field with
\(q'\) elements the polynomial \(1+\prod_a(X-a)\), over its \(q'\)
elements \(a\), would have no root. A nonzero polynomial over an
infinite field has a point where it does not vanish: in one variable
this follows from the degree bound on the number of roots, and
in several variables first make a nonzero coefficient in the last
variable nonzero by induction. Lift this point coordinatewise to
\(O\). Evaluation of the normalized polynomial now has absolute
value one. The zero polynomial causes no restriction.
\end{proof}

\begin{definition}\label{def:linear-jet}
Keep the diagram \(\kappa_{\tau,h}\), \(N_j,E_j,n\) of
Definition~\ref{def:stage-diagrams}. Fix \(d'\ge1\).
Let \(\Theta\) be a formal power series over \(K_p\) in variables
\(Z_{\tau,a,b}\), \(1\le a\le d',\ 1\le b\le n\). A linear substitution
\(L\) assigns to each variable a homogeneous linear form over \(O\)
in a finite set of auxiliary variables. For \(\varpi,\theta\in O\)
of valuations \(\rho_p,\omega_p\) respectively, put
\[
 \rho_p=\frac{1}{p-1}-\omega_p\sum_{e\in I_p}p^{-e}.
\]
Say that \(\Theta\) has integral coefficients of order \(N_j\) on
the stated linear substitutions if every such substitution satisfying
\[
 \sum_{\tau,a,b} L_{\tau,a,b}^{p^e}L_{(-p^e)\tau,a,b}=0
       \pmod{(p,\theta)}\qquad(e\in I_p^0,\ e\le E_j)
\]
makes all coefficients of degree \(N_j\) in \(\Theta(\varpi L)\)
integral. Congruences are coefficientwise. Only terms of degree
\(N_j\) in \(\Theta\) enter this property.

For an output filling \(T'\) with entries in \([1,n]\) and the
given shape, list its entries in a height \(h\) column \(c\) at
\(\tau\) as \(b_{c,1},\ldots,b_{c,h}\). Define the polynomial
\[
 F_{T'}(Y)=
 \left.
 \left(\prod_{\tau,c}
   \det\left(\sum_aY_{\tau,i,a}
              \frac{\partial}{\partial Z_{\tau,a,b_{c,l}}}
       \right)_{1\le i,l\le h}\right)\Theta
 \right|_{Z=0}.
\]
Partial derivatives are formal derivatives of monomials; they commute.
This polynomial has a unique expansion in the standard products
\(P_T(Y)\) of Lemma~\ref{lem:flag-straighten}, of the indicated shape
and input alphabet \([1,d']\). Indeed each determinant in the product
expands as a sum over distinct increasing \(h\)-tuples of labels \(a\),
of their \(Y\)-minor times the corresponding determinant of derivatives.
To see this, expand by multilinearity in the rows of the derivative
matrix; repeated intermediate labels vanish by alternation, and
collecting each remaining permutation gives the minor. Multiplication
followed by Lemma~\ref{lem:flag-straighten} proves the assertion,
also if \(d'<h\) so that the expansion is zero.
\end{definition}

\begin{lemma}\label{lem:val-factorial}
For a prime \(p\) and \(0\le t\le N\) integers with \(N\ge1\),
write \(s_p(t)\) for the sum of the base-\(p\) digits. Then
\[
 v_p(t!)=\frac{t-s_p(t)}{p-1},\qquad
 \frac{t}{p-1}-(1+\lfloor\log N/\log p\rfloor)
                  \le v_p(t!)\le\frac{t}{p-1}.
\]
\end{lemma}
\begin{proof}
Counting the positive multiples of each power \(p^l\) up to \(t\)
gives \(v_p(t!)=\sum_{l\ge1}\lfloor t/p^l\rfloor\). Insert the
base-\(p\) expansion of \(t\) and sum each digit times its finite
geometric sum to obtain the identity. At most
\(1+\lfloor\log N/\log p\rfloor\) digits each contribute between
zero and \(p-1\), which proves both bounds.
\end{proof}

\begin{theorem}[Divisibility of differential minors]\label{thm:schur-jet}
Suppose \(\Theta\) has the integrality property of
Definition~\ref{def:linear-jet}. Fix nonnegative integers
\(\gamma_{\tau,h,b}\). For \(e\in I_p^0,\ e\le E_j,\ \sigma=(-p^e)\tau\),
require
\[
 \gamma_{\tau,1,b}+\gamma_{\sigma,h',b}\ge1
 \quad\text{whenever }\ \tau\in V\ \text{or }(h'=1,\ \sigma\in V)
\]
for all \(b\) and all contributing source and target heights in this
formula. Put
\(m_{T'}=\sum_{\tau,c}\sum_{l=1}^{h}\gamma_{\tau,h,b_{c,l}}\).
Each coefficient of \(F_{T'}\) in its standard minor basis has valuation
at least
\[
                         J_p(j,V)-m_{T'} .
\]
\end{theorem}
\begin{proof}
Evaluate at arbitrary \(Y_{\tau,i,a}=y_{\tau,i,a}\in O\) and use the
changed frames of Theorem~\ref{thm:null-frames}. For every contributing
\(\tau,h\), \(1\le i\le h,\ 1\le b\le n\), introduce a variable
\(X_{\tau,h,i,b}\). There are at most \(fn\sum_{h<n}h\le fn^3\)
of them. In Definition~\ref{def:linear-jet} take the linear forms
\[
 L_{\tau,a,b}=\sum_{h,i}p^{\gamma_{\tau,h,b}}
                   \tilde y_{\tau,h,i,a}X_{\tau,h,i,b}.
\]
In characteristic \(p\) the \(p^e\)-th power distributes over this
sum. For fixed \(\tau,b\), each term after summing over \(a\) in
a coupling contains a pairing of the two changed frames. It vanishes
modulo \((p,\theta)\) by Theorem~\ref{thm:null-frames} or, in an
exempt pair, by the displayed charge condition because \(p=0\)
in that quotient. Hence every coefficient of degree \(N_j\) of
\(\Theta(\varpi L)\) lies in \(O\).

Apply to this substituted series, at zero, the product of the column
determinants of derivatives
\(\det(\partial/\partial X_{\tau,h,i,b_{c,l}})_{i,l\le h}\).
An order \(N_j\) monomial derivative of multi-index \(\alpha\) gives
\(\prod_i\alpha_i!\) times its degree \(N_j\) coefficient, by successive
differentiation of powers. Lemma~\ref{lem:val-factorial} gives
\(v_p(\alpha_i!)\ge\alpha_i/(p-1)-(E_j+1)\) for every \(i\).
Thus the differentiated value has valuation at
least \(N_j/(p-1)-fn^3(E_j+1)\), also after the signed additions
in the determinants, by Lemma~\ref{lem:place-fields}.

The same differentiated value equals
\[
  \varpi^{N_j}p^{m_{T'}}
       \left(\prod_{\tau,h}(\det M_{\tau,h})^{\kappa_{\tau,h}}\right)
                         F_{T'}(y).
\]
In fact the derivative in \(X_{\tau,h,i,b}\) acts by
\(\varpi p^{\gamma_{\tau,h,b}}\sum_a\tilde y_{\tau,h,i,a}
\partial/\partial Z_{\tau,a,b}\) before substitution. On each
monomial this follows by the product rule, so also for repeated
derivatives in the degree \(N_j\) polynomial. Factoring the
scalars from each column and \(M_{\tau,h}\) from the matrix of
each determinant gives the expression. The determinant factor has
valuation at most \(\omega_p\sum_{e,\tau}K(\tau,e)\) by
Theorem~\ref{thm:null-frames}, with precisely the exponents
\(e\in I_p^0,\ e\le E_j\). Subtract its valuation and
\(N_j\rho_p+m_{T'}\) from the lower bound. By the definitions of
\(J_p\) and \(\rho_p\) the result is \(v_p(F_{T'}(y))\ge
J_p(j,V)-m_{T'}\). Lemma~\ref{lem:poly-test} transfers this bound
from integral evaluations to all monomial coefficients.
Lemma~\ref{lem:flag-straighten} then gives it for all standard
minor coordinates.
\end{proof}

\section{Frobenius on the rows of a tight frame}

\begin{definition}\label{def:tight-field}
Use the representatives \(\tau:E\to E\) fixing \(K\) for the places
in Proposition~\ref{prop:places}. A finite family \(w_i\in
N_1^{-1}O_E\) has the tight frame identity of length \(M'\) if
\(M'\) is a positive integer and
\[
     \sum_i \tau(w_i)\overline{\sigma(w_i)}
                   =M'{\bf1}_{\tau=\sigma}\qquad(\tau,\sigma\in G).
\]
The family \(w_i=\zeta_{N_1}^i,\ 0\le i<N_1\), has this identity
with \(M'=N_1\). Indeed the sum is the finite geometric sum
\(\sum_i\zeta_{N_1}^{(\tau-\sigma)i}\), which vanishes unless
\(\tau=\sigma\bmod N_1\). The length in this example is prime to
each nonsplit \(p\) by Proposition~\ref{prop:places}.
\end{definition}

\begin{proposition}\label{prop:frame-frobenius}
Take any prime \(p<x_1\) with \(\eta_d(p)\ne1\), an embedding
\(\iota:E\hookrightarrow K_p\), and an element \(\theta\) of valuation
\(\omega_p\). Let \(w_i\) satisfy Definition~\ref{def:tight-field}.
For \(\beta\) in the valuation ring \(O\) of \(K_p\), put
\(z_{\tau i}=\beta\iota(\overline{\tau(w_i)})\).
All these entries are integral. With \(\bar z\) denoting reduction
modulo \(\theta O\) in the following formula, for every
\(e\in I_p^0,\ \tau,\sigma\in G\) one has
\[
 \sum_i \bar z_{\tau i}^{p^e}\bar z_{\sigma i}
   =\bar\beta^{p^e+1}M'\,{\bf1}_{\sigma=(-p^e)\tau}.
\]
\end{proposition}
\begin{proof}
Since the images under \(\iota\) of algebraic integers have absolute
value at most one by their monic integral equations, and \(p\nmid N_1\),
each image of \(w_i\) and its conjugates is integral at \(p\).
The quotient \(O/\theta O\) has characteristic \(p\), since
\(v_p(p)=1\ge\omega_p>0\). We claim the two maps
\[
 w\longmapsto\overline{\iota(\overline{\tau(w)})}^{\,p^e},
 \qquad
 w\longmapsto\overline{\iota\bigl(((-p^e)\tau)(w)\bigr)}
\]
from \(O_E[1/N_1]\) to this quotient agree, with the outer bars
denoting reduction. They agree on \(\zeta_{N_1}\) by the exponents.
If \(p\) is inert, \(v_p(\theta)=1\) and \(\theta O=pO\).
The image of \(\omega\) is a root of its irreducible quadratic over
\(\mathbf F_p\), as in Lemma~\ref{lem:inert-prime-fibres}.
In the subfield it generates, an odd power of Frobenius exchanges
\(1-\omega\) and \(\omega\). This applies to every \(e\in I_p^0\)
in the inert case. If \(p=d\), then \(v_p(\iota(\sqrt{-d}))=1/2\)
from its square, so \(\omega\) and \(1-\omega\) both reduce to \(1/2\).
Every Frobenius power fixes that element, including the identity
at \(e=0\). Since \(\omega,\zeta_{N_1}\) generate \(O_E\) by
Theorem~\ref{thm:integer-ring} and \(N_1\) is invertible in the
quotient, the claim follows. Now multiply the claim for \(w_i\)
by the reduction of \(\iota(\overline{\sigma(w_i)})\) and sum.
The tight frame identity with places \((-p^e)\tau,\sigma\)
gives the formula, after inserting the common power of \(\beta\).
\end{proof}

\section{From a one-variable expansion to the theta product}

\begin{lemma}\label{lem:frob-cancellation}
Let \(R\) be a commutative ring of characteristic \(p\), \(H\) a
nilpotent ideal, and \(a_0\in\{1,2\}\). For finitely many
\(l_i,t_i\in H\) suppose \(t_i-t_i^{p^{a_0}}=l_i\). If \(a_0=2\),
suppose \(\sum_i l_i^{p^e}l_i=0\) for all odd \(e\ge1\).
If \(a_0=1\), suppose the same holds for all \(e\ge0\).
Then
\[
                  \sum_i t_i^{p+1}=0,\qquad
                  \sum_i t_i^{p+p^{a_0}}=0 .
\]
\end{lemma}
\begin{proof}
Substitution in \(t_i=l_i+t_i^{p^{a_0}}\), using additivity of
Frobenius, gives \(t_i=\sum_{r\ge0}l_i^{p^{a_0r}}\); the iterated
remainder vanishes and all sums are finite by nilpotence.
For \(A,B\ge0\) one has
\[
  \sum_i l_i^{p^A}l_i^{p^B}=
       \left(\sum_i l_i^{p^{|A-B|}}l_i\right)^{p^{\min(A,B)}}.
\]
Expand the two expressions in the assertion using the finite
series for the \(t_i\). The exponent pairs are respectively
\((a_0r+1,a_0s)\) and \((a_0r+1,a_0s+a_0)\).
For \(a_0=2\) their differences are odd; for \(a_0=1\) all possible
differences occur among the hypotheses. Every term vanishes by the
displayed identity.
\end{proof}

\begin{theorem}\label{thm:product-null}
Use \(p,\varpi,\theta,\rho_p,\omega_p\) of
Definition~\ref{def:linear-jet} at a nonsplit prime, and
\(z_{\tau i},M'\) of Proposition~\ref{prop:frame-frobenius}.
Let \(c\) range over a finite index set. Write \(Z_{\tau,c}\) for
independent variables and \(l_{ic}(Z)=\sum_\tau z_{\tau i}Z_{\tau,c}\).
Suppose there are \(S,\psi\in XO[[X]]\), \(u_0,u_1,u_2\in O\)
and \(a_0\in\{1,2\}\) with
\[
 \begin{aligned}
 S-S^{p^{a_0}}&\equiv X\pmod{\theta},\\
 \psi&\equiv u_0 X^2+u_1 S^{p+1}+u_2 S^{p+p^{a_0}}
                                                  \pmod{\theta}.
 \end{aligned}
\]
Require \(p=d\) if \(a_0=1\), and \(u_0=0\) if \(a_0=2\).
For given \(V_{ic}\in O[[X]]\) and \(\kappa\in K_p^\times\) with
\(v_p(\kappa)=1/(p-1)-\omega_p\), suppose the series \(\Theta\)
satisfies
\[
 \Theta(\varpi Z)=\exp\left(\kappa\sum_{i,c}\psi(l_{ic}(Z))\right)
                              \prod_{i,c} V_{ic}(l_{ic}(Z)).
\]
Then for every \(B\ge1\) and every integral homogeneous linear
substitution \(Z_{\tau,c}=L_{\tau,c}\) satisfying
\[
 \sum_{\tau,c}L_{\tau,c}^{p^e}L_{(-p^e)\tau,c}
              \equiv0\pmod{\theta}\qquad
                (e\in I_p^0,\ e\le\lfloor\log B/\log p\rfloor),
\]
the coefficients through degree \(B\) of \(\Theta(\varpi L)\) belong
to \(O\). For \(B=N_j\) and the indices \(c=(a,b)\) in
Definition~\ref{def:linear-jet}, this supplies the integrality
hypothesis of Theorem~\ref{thm:schur-jet}.
\end{theorem}
\begin{proof}
Take the polynomial ring over \(O/\theta O\) in the substitution
variables modulo all monomials of degree \(>B\). Its positive-degree
ideal is nilpotent and its characteristic is \(p\), because
\(\omega_p\le1\). Every series below can be evaluated in that ring
on positive-degree arguments by finite truncation. Set
\(l_i'=l_{ic}(L),\ t_i'=S(l_i')\), with \(c\) included in the
index \(i'\) in the following sums. By Proposition~\ref{prop:frame-frobenius}
and expansion of the \(p^e\)-th power,
\[
     \sum_{i,c} l_{ic}(L)^{p^e}l_{ic}(L)=
     \bar\beta^{p^e+1}M'\sum_{\tau,c} L_{\tau,c}^{p^e}
                                          L_{(-p^e)\tau,c}.
\]
For \(e\in I_p^0\) up to the stated cutoff this vanishes by hypothesis.
Above the cutoff \(p^e>B\), and each term vanishes by degree.
The first congruence gives \(t_i'-(t_i')^{p^{a_0}}=l_i'\).
Lemma~\ref{lem:frob-cancellation} therefore annihilates both sums
of powers of \(t_i'\) in the second congruence. For \(a_0=1\)
the remaining quadratic sum of the \(l_i'\) vanishes by the
\(e=0\) congruence at \(p=d\). For \(a_0=2\) its coefficient is zero.
It follows that \(\sum_{i,c}\psi(l_{ic}(L))\) vanishes in the
truncated ring.

Thus through degree \(B\) every coefficient of the last sum has
valuation at least \(\omega_p\) and its constant term is zero.
After multiplication by \(\kappa\), the coefficients have valuation
at least \(1/(p-1)\). In the \(\ell'\)-th power each coefficient
of degree at most \(B\) then has valuation at least
\(\ell'/(p-1)\), by finite multiplication. The factorial estimate
\(v_p(\ell'!)\le\ell'/(p-1)\) follows from Lemma~\ref{lem:val-factorial}. Hence division by
\(\ell'!\) leaves these coefficients integral. Summation over
\(\ell'\) in the formal exponential is finite in each degree,
and the other factors in the product have integral coefficients.
This proves the first assertion. Finally \((p,\theta)=\theta O\)
because \(v_p(\theta)\le1\). The congruences of
Definition~\ref{def:linear-jet} are exactly the ones just used.
\end{proof}

\begin{corollary}\label{cor:sparse-derivative}
In Theorem~\ref{thm:product-null} suppose \(p\) is inert, \(\beta\)
and \(u_1\) are units, and \(p\nmid M'\). Set
\(\Psi(Z)=\sum_{i,c}\psi(l_{ic}(Z))\), and now use bars for reduction
to the residue field of \(O\). For every \(\sigma,c\),
\[
 \frac{\partial\bar\Psi}{\partial Z_{\sigma,c}}
   =\bar u_1 M'
       \sum_{\substack{e\ge1\ {\rm odd},\ \tau\in G\\
                        (-p^e)\tau=\sigma}}
              \bar\beta^{p^e+1} Z_{\tau,c}^{p^e}.
\]
Every displayed monomial has nonzero coefficient.
\end{corollary}
\begin{proof}
Here \(a_0=2,\ u_0=0\). After reduction, the solution of
\(\bar S-\bar S^{p^2}=X\) with zero constant term is
\(\bar S(X)=\sum_{r\ge0} X^{p^{2r}}\), by iteration as in
Lemma~\ref{lem:frob-cancellation} in every finite truncation.
Its formal derivative is one. The derivative of
\(\bar S^{p+p^2}\) vanishes in characteristic \(p\), while
that of \(\bar S^{p+1}\) is \(\bar S^p\).
Consequently differentiating \(\bar\Psi\) gives
\(\bar u_1\sum_{i,r}\bar z_{\sigma i}
        (\sum_\tau \bar z_{\tau i} Z_{\tau,c})^{p^{2r+1}}\).
Apply Proposition~\ref{prop:frame-frobenius} in the residue field
to its coefficients. Each exponent determines \(e\), and the map
\(\tau\mapsto(-p^e)\tau\) is a permutation of \(G\). Thus the
displayed monomials are distinct, and their coefficients are units
by the three additional hypotheses.
\end{proof}

\section{A normalized complex multiplication lattice and its theta function}

\begin{lemma}\label{lem:elliptic-count}
Let \(\Lambda=\lambda(\Z+\tau\Z)\) with \(\lambda\ne0,\ \Im\tau>0\),
where \(\Im\) denotes imaginary part. For a nonzero meromorphic
function \(f\), on a closed fundamental parallelogram whose boundary
avoids zeros and poles, the integral of \(f'/f\) on its positively
oriented boundary divided by \(2\pi i\) counts its zeros minus its
poles inside, with multiplicities. In particular a meromorphic
\(\Lambda\)-periodic function that is not identically zero has
equally many zeros and poles per fundamental parallelogram.
An entire \(\Lambda\)-periodic function is constant.
\end{lemma}
\begin{proof}
Zeros and poles are isolated by Lemma~\ref{lem:analytic-tools}, so
there are finitely many in the compact parallelogram. Around one of
order \(n\), positive for a zero and negative for a pole, write
\(f=(z-z_0)^n h\) with \(h\) nonzero holomorphic locally. Then
\(f'/f=n/(z-z_0)+h'/h\). Remove small disjoint disks and use the
triangle integral identity and polygonal approximation in
Lemma~\ref{lem:analytic-tools} on the region between their boundaries
and the parallelogram. Integrating gives \(2\pi i n\) from each
disk. For a periodic logarithmic derivative the integrals on
opposite sides cancel. One can translate a parallelogram slightly to
avoid singularities on the boundary: in a bounded neighbourhood there
are finitely many such points and the forbidden translations lie on
finitely many line segments. An entire periodic function is bounded
on a compact parallelogram and hence everywhere. Its derivative at
any point has modulus at most this global bound divided by the
radius of any circle centered there, by Lemma~\ref{lem:analytic-tools}.
Letting the radius tend to infinity proves constancy.
\end{proof}

\begin{proposition}\label{prop:reduced-sigma}
Put \(y=\Im\tau>0,\ A=|\lambda|^2 y\). On the lattice
\(\Lambda=\lambda(\Z+\tau\Z)\) there is a unique odd entire function
\(\sigma_\Lambda\) with derivative one at zero and transformation law
\[
 \sigma_\Lambda(z+l)=s(l)\exp\!\left(\frac{\pi}{A}
                 (\bar l z+|l|^2/2)\right)\sigma_\Lambda(z)
       \qquad(l\in\Lambda),
\]
where \(s(l)=1\) if \(l/2\in\Lambda\) and \(s(l)=-1\) otherwise.
Its zeros are the lattice points, each simple.
\end{proposition}
\begin{proof}
Set
\[
 F(u)=\sum_{j\in\Z}\exp\left(\pi i(j+1/2)^2\tau+
                                     2\pi i(j+1/2)(u+1/2)\right).
\]
On every bounded set the Gaussian decay in \(|j|^2\) dominates
all powers of \(j\) and the exponential of a constant times \(|j|\).
Thus the series and its differentiated series converge uniformly
there, giving an entire function by the disk formula of
Lemma~\ref{lem:analytic-tools}. Integrating
\(F(u)\exp(-\pi i u)\) on \(0\le u\le1\) isolates its nonzero
\(j=0\) coefficient by elementary integration. Replacing \(j\)
by \(-j-1\) gives \(F(-u)=-F(u)\); shifting the index once gives
\[
 F(u+1)=-F(u),\qquad F(u+\tau)=-\exp(-\pi i\tau-2\pi i u)F(u).
\]
In the integral of \(F'/F\) over a generic parallelogram for
\(1,\tau\) with zero in its interior, the two sides parallel to
\(\tau\) cancel. On translation by \(\tau\) the logarithmic
derivative decreases by \(2\pi i\), so the other two sides
contribute \(2\pi i\). By Lemma~\ref{lem:elliptic-count} there is
one zero with multiplicity inside. Oddness puts it at zero, so it
is simple. The translation laws give the same result in all
translates, whose boundaries contain no zeros.

For \(u=z/\lambda\) put
\(\sigma_\Lambda(z)=\lambda \exp(\pi u^2/(2y))F(u)/F'(0)\).
It is odd and has derivative one at zero. At the generators
\(\lambda,\lambda\tau\) the required sign is negative. Their
transformation factors follow from the two formulas for \(F\),
using \(\tau-2iy=\bar\tau\). For \(l=m\lambda+n\lambda\tau\)
the composed sign is \((-1)^{m+n+mn}\): adding \(\lambda\)
to \(l\) introduces the cross phase
\(\exp(\pi i\Im(\bar\lambda l)/A)=(-1)^n\), and adding
\(\lambda\tau\) gives \((-1)^m\). These recurrences also run
backwards. This sign is positive exactly when \(m,n\) are even.
A second function with the stated properties vanishes at zero and
hence at every lattice point. Its quotient by \(\sigma_\Lambda\)
is entire since the latter zeros are simple. The quotient is periodic,
therefore constant by Lemma~\ref{lem:elliptic-count}, and the two
derivatives at zero fix that constant as one.
\end{proof}

\begin{proposition}\label{prop:scaled-cm}
Let \(O_K=\Z+\omega\Z\), \(\omega=(1+\sqrt{-d})/2\), using the
embedding of positive imaginary part. For a lattice \(\Lambda\)
as in Proposition~\ref{prop:reduced-sigma} define
\[
 \wp_\Lambda(z)=z^{-2}
        +\sum_{0\ne l\in\Lambda}\bigl((z-l)^{-2}-l^{-2}\bigr)
                         \qquad(z\notin\Lambda).
\]
There are \(\lambda\ne0,\ t=\lambda/2,\ a_1\in\C\) such that
\(\Lambda=\lambda O_K\) is stable under \(O_K\),
\(2t\in\Lambda,\ t,\omega t,\bar\omega t\notin\Lambda\), and
\[
       \wp_\Lambda''(t)=2,\qquad a_1^2=12\wp_\Lambda(t).
\]
The functions \(\xi(z)=\wp_\Lambda(z+t)-\wp_\Lambda(t)\) and
\(\phi(z)=(\wp_\Lambda'(z+t)-a_1\xi(z))/2\) then satisfy
\[
 \phi^2+a_1\xi\phi=\xi^3+\xi,\qquad
 \xi(0)=\xi'(0)=\phi(0)=0,\quad \phi'(0)=1 .
\]
\end{proposition}
\begin{proof}
The real linear map \((m,n)\mapsto\lambda(m+\tau n)\) is invertible,
so its norm is bounded below by a positive constant times
\(\max(|m|,|n|)\), by compactness of the unit square's boundary.
There are \(8j\) integer pairs with this maximum \(j\ge1\).
Consequently \(\sum_{l\ne0}|l|^{-3}<\infty\). The summands in the
definition are \(O_C(|l|^{-3})\) on a compact set \(C\) off the
lattice for large \(|l|\), by subtraction of the fractions.
Termwise differentiation is justified locally by the disk formula in
Lemma~\ref{lem:analytic-tools}. Thus \(\wp_\Lambda\) is even and
meromorphic, with principal part \(z^{-2}\) at zero and zero constant
term there. Its derivative, \(-2\sum_{l\in\Lambda}(z-l)^{-3}\),
is periodic by absolute convergence. Each period difference of
\(\wp_\Lambda\) is constant; possible poles cancel by their identical
principal parts. For either primitive generator \(l=\lambda,\lambda\tau\)
evaluate the difference at \(-l/2\), which is not a pole.
Evenness gives zero. This proves full periodicity.

Near zero write \(\wp_\Lambda(z)=z^{-2}+c_2 z^2+c_4 z^4+O(z^6)\)
and set \(g_2=20c_2,\ g_3=28c_4\). Expansion at zero cancels all
negative powers and the constant term in
\(\wp_\Lambda'^2-4\wp_\Lambda^3+g_2\wp_\Lambda+g_3\).
The resulting entire elliptic function is zero by
Lemma~\ref{lem:elliptic-count}. Differentiation yields
\(\wp_\Lambda''=6\wp_\Lambda^2-g_2/2\), first away from the isolated
zeros of \(\wp_\Lambda'\), then throughout off the poles by continuity.
Up to lattice translation, \(\wp_\Lambda'\) has exactly one pole
of order three. By oddness and periodicity it vanishes at the
three nonzero half periods. They are distinct, so the counting
formula of Lemma~\ref{lem:elliptic-count} makes each zero simple.

For the initial lattice \(O_K\) take \(t_0=1/2\), so
\(\wp_{O_K}''(t_0)\ne0\). Choose any fourth root
\(\lambda^4=\wp_{O_K}''(t_0)/2\) and the required square root
\(a_1\) in \(\C\), using Lemma~\ref{lem:place-fields}.
The defining series gives
\(\wp_{\lambda O_K}(\lambda z)=\lambda^{-2}\wp_{O_K}(z)\), hence
\(\wp_\Lambda''(t)=2\). Stability follows since \(O_K\) is a
ring by the monic quadratic relation for \(\omega\).
In the basis \(1,\omega\), the coordinates of \(1/2,\omega/2\)
and \((1-\omega)/2\) are not all integral in each pair, proving
the half-period assertions. Finally put \(e_t=\wp_\Lambda(t)\).
The differential equation gives \(12e_t^2-g_2=4\).
Subtract the cubic equation at \(t\), where \(\wp_\Lambda'(t)=0\),
from the one at \(t+z\). The result is
\((2\phi+a_1\xi)^2=4\xi^3+12e_t\xi^2+4\xi\), which yields the
displayed equation. Differentiation at the half period gives the
asserted initial values.
\end{proof}

\begin{proposition}\label{prop:cm-theta}
Take the lattice of Proposition~\ref{prop:scaled-cm}, with area \(A\).
For an entire function \(F\) and \(y\in\C\) put
\[
 (U_yF)(z)=\exp\!\left(-\frac{\pi}{A}
                          (z\bar y+|y|^2/2)\right)F(z+y).
\]
The normalized theta function
\[
 \theta_\Lambda(z)=\frac{(U_t\sigma_\Lambda)(z)}{(U_t\sigma_\Lambda)(0)}
              =\exp(-\pi z\bar t/A)
                      \frac{\sigma_\Lambda(z+t)}{\sigma_\Lambda(t)}
\]
is even and entire with \(\theta_\Lambda(0)=1\) and simple zeros
at \(t+\Lambda\). Its multiplier at \(l\in\Lambda\) is given by
\[
 \theta_\Lambda(z+l)=s(l)\exp\!\left(\frac{\pi}{A}
    (\bar l z+|l|^2/2+\bar l t-l\bar t)\right)\theta_\Lambda(z).
\]
For a prime \(p\), an isomorphism \(\iota:\C\to K_p\) as in
Definition~\ref{def:place-height}, and torsion shifts \(y_{ic}\)
with \(n'y_{ic}\in\Lambda\) for some \(n'\ge1\) prime to \(p\),
the following formal theta product is well defined. Choose a frame
\(w_i\) of Definition~\ref{def:tight-field} of length prime to \(p\)
and \(b'\in\C\) with \(v_p(\iota b')\ge0\), and put
\[
 \begin{aligned}
 F_y(X)&=\sum_{j\ge0}\iota\!\left((U_y\theta_\Lambda)^{(j)}(0)/j!\right)X^j,\\
 \Theta(Z)&=\prod_{i,c}
 F_{y_{ic}}\left(\sum_{\tau\in G}
       \iota\bigl(b'\overline{\tau(w_i)}\bigr) Z_{\tau,c}\right).
 \end{aligned}
\]
For \(b'=1,\ y_{ic}=0\) its constant term equals one.
\end{proposition}
\begin{proof}
The value \(\sigma_\Lambda(t)\) is nonzero by
Propositions~\ref{prop:scaled-cm} and~\ref{prop:reduced-sigma}.
Their zero set and multiplier formulas give the zeros and the stated
translation formula. Since \(2t\in\Lambda\) but \(t\notin\Lambda\),
one has \(s(2t)=-1\) and
\(\sigma_\Lambda(z-t)=-\exp(-2\pi z\bar t/A)\sigma_\Lambda(z+t)\).
Using oddness of \(\sigma_\Lambda\) gives
\(\theta_\Lambda(-z)=\theta_\Lambda(z)\). Each translate is
entire, so its indicated derivatives exist by
Lemma~\ref{lem:analytic-tools}. Applying \(\iota\) to coefficients
gives formal series. The arguments have zero constant terms and
are homogeneous linear, so in each fixed degree only that degree
of each composed series can contribute before multiplication, and
only finitely many terms enter the product coefficient. The last
claim follows from \(U_0\theta_\Lambda(0)=1\).
\end{proof}

\section{The discriminant along a cycle of index-two lattices}

\begin{lemma}\label{lem:half-lattice-cycle}
Put \(O_K=\Z[\omega]\) as in Lemma~\ref{lem:quadratic-norm}.
For a rank-two complex lattice \(\Lambda\) stable under \(O_K\)
there is exactly one nonzero class \(t\) in
\((\tfrac12\Lambda)/\Lambda\) for which \(\omega t\) and
\(\bar\omega t\) are both nonzero. Write \(T=\omega t,\ S=\bar\omega t\)
in this group, so \(t=T+S\).
The lattice \(\Lambda'=\Lambda+(\omega/2)\Lambda
=\Lambda+\Z T\), with any representative of \(T\), contains
\(\Lambda\) with index two. In \(\C/\Lambda'\) the image of \(t\)
is \(S'=\bar\omega t'\) for the analogous distinguished point \(t'\).
Starting at \(O_K\), after finitely many such enlargements one returns
to a complex multiple of \(O_K\).
\end{lemma}
\begin{proof}
Multiplication by \(\omega\) on the real vector space \(\C\)
has characteristic polynomial \(X^2-X+(d+1)/4\). In an integral
basis of \(\Lambda\) its matrix is integral. Since \(d\equiv7\bmod8\),
the reduced polynomial over \(\mathbf F_2\) has the two distinct
roots 0 and 1. The two-dimensional space
\(\Lambda/2\Lambda\simeq(\tfrac12\Lambda)/\Lambda\) is the direct
sum of two one-dimensional eigenspaces: the idempotents are
\(\omega,1-\omega\) and neither is zero by the characteristic
polynomial. The unique mixed nonzero vector has a nonzero component
in both. The image of \(\omega\) is generated by \(T\), proving
the index. The image of \(t\) in the next quotient is nonzero
and killed by \(\omega\), so it is \(S'\), the unique nonzero such
element of the new two-torsion group.

Write \(I_n=(O_K+(\omega/2)O_K)^n\), the lattice after \(n\) steps
starting from \(O_K\), and \(J_n=2^n I_n=(2,\omega)^n\).
Here products of ideals are finite sums of products of elements.
Since \((2,\omega)(2,\bar\omega)=2O_K\), we have
\(J_n\bar J_n=2^n O_K\): the generators of the first product are
divisible by 2 using \(\omega\bar\omega=(1+d)/4\), and
\(2\omega+2\bar\omega=2\). Each enlargement has index two, so
\(I_n\) has area \(2^{-n}\) times the area of \(O_K\).
Multiplication by \(2^n\) shows \([O_K:J_n]=2^n\), by the
lattice determinant calculation of Lemma~\ref{lem:quadratic-norm}.
The same pigeonhole argument used there supplies nonzero \(z_n\in J_n\)
with \(|z_n|\le 2(1+|\omega|)2^{n/2}\), by comparing the elements
\(a+b\omega,\ 0\le a,b\le\lceil 2^{n/2}\rceil\), modulo \(J_n\).
Then \(z_n\bar J_n/2^n\subset O_K\) has index
\(|z_n|^2/2^n\le 4(1+|\omega|)^2\), again by areas.
There are finitely many sublattices of these indices, since
a sublattice of index \(h\) contains \(hO_K\) by the order of its
finite additive quotient. Equality of two such sublattices at
\(n<n'\), followed by conjugation and rescaling, gives
\(J_{n'}=cJ_n\) for some \(c\ne0\). Multiply by \(\bar J_n\)
and use \(J_{n'}=J_{n'-n}J_n\) to get \(J_{n'-n}=cO_K\).
It follows that \(I_{n'-n}\) is homothetic to the initial lattice.
\end{proof}

\begin{proposition}\label{prop:half-recurrence}
For the lattice and the three classes of
Lemma~\ref{lem:half-lattice-cycle} put \(e_P=\wp_\Lambda(P)\)
for \(P=t,T,S\). They are distinct. Define the nonzero values
\[
 u=\frac{e_T-e_t}{e_S-e_t},\qquad
 V=\frac{16u^2}{1-u},\qquad D=\frac{16(u-1)^2}{u}.
\]
These are unchanged when the lattice and the classes are scaled
by the same nonzero complex number. For the next lattice denote the
corresponding values by \(u',V',D'\). Then
\[
                   V'=-DV,\qquad (48+V+D)V'=-4096 .
\]
\end{proposition}
\begin{proof}
For any nonzero half period \(P\) the elliptic function
\(\wp_\Lambda(z)-e_P\) has a double pole and a zero of order at least
two at \(P\), by Proposition~\ref{prop:scaled-cm} and its proof.
Lemma~\ref{lem:elliptic-count} thus gives exactly that double zero
per cell. In particular the three \(e_P\) are distinct and give
the three roots of \(4X^3-g_2X-g_3\), by the same differential
equation. Their sum is zero. Homothety invariance follows from
the scaling formula for \(\wp\) in Proposition~\ref{prop:scaled-cm}.

Put \(X(z)=\wp_\Lambda(z)-e_T\) and
\(B_T=(e_T-e_t)(e_T-e_S)\). The product \(X(z)X(z+T)\) has
its double zeros paired with its double poles. It is elliptic
and entire, hence constant by Lemma~\ref{lem:elliptic-count}.
Evaluation at \(z=t\), for which \(t+T=S\), gives \(B_T\).
Also
\[
             \wp_{\Lambda'}(z)=\wp_\Lambda(z)+\wp_\Lambda(z+T)-e_T.
\]
Indeed the right side has periods \(\Lambda'\) and the same double
principal parts as the left side. Its constant term at zero in the
Laurent expansion is zero. Their difference is constant and
therefore zero. Since \(e_t+e_S=-e_T\), subtraction of the value
at \(t\) now gives
\[
       \wp_{\Lambda'}(z)-\wp_{\Lambda'}(t)
                       =X(z)+3e_T+B_T/X(z).
\]
On the new torus the image of \(t\) is \(S'\) by
Lemma~\ref{lem:half-lattice-cycle}. The two remaining nonzero half
periods \(t',T'\) have representatives \(y\) with
\(2y=T\bmod\Lambda\): those with \(2y=0\bmod\Lambda\) give only
the identity and \(S'\) after quotienting by the order-two
subgroup generated by \(T\). For such \(y\), \(y+T=-y\bmod\Lambda\)
and \(X(y)^2=B_T\). Let \(r_1,r_2\) be the two values of
\(\wp_{\Lambda'}(y)-\wp_{\Lambda'}(t)\) with \(y=t',T'\) in that order.
They are distinct by the first part of the proof at \(\Lambda'\).
Hence they are the two numbers \(3e_T\pm2\sqrt{B_T}\). In particular
\[
 (r_1-r_2)^2=16 B_T,\qquad r_1r_2=(e_S-e_t)^2 ,
\]
where the product uses \(3e_T=(e_T-e_t)+(e_T-e_S)\). Since
\(u'=(r_1-r_2)/r_1\), we obtain
\[
 V'=16\,\frac{(r_1-r_2)^2}{r_1r_2}
            =256\,u(u-1).
\]
Substitution in the definitions gives \(-DV=256u(u-1)\) and
\(48+V+D=-16/[u(u-1)]\). This proves both recurrences.
\end{proof}

\begin{lemma}\label{lem:cycle-unit}
In a nonarchimedean valued field extending \(\Q_p\), let nonzero
\(V_i,D_i\) have indices modulo a positive integer and satisfy
\(V_{i+1}=-D_i V_i,\ (48+V_i+D_i)V_{i+1}=-4096\).
Then \(|D_i|_p=1\) for every \(i\).
\end{lemma}
\begin{proof}
Put \(r_i=|V_i|_p,\ c=|4096|_p\), with \(0<c\le1\).
When \(r_i\le1\) and \(r_{i+1}>r_i\), the term \(D_i\)
of modulus \(r_{i+1}/r_i>1\) strictly dominates the other
two terms in \(48+V_i+D_i\). The ultrametric inequality in both
directions gives \(r_{i+1}^2=cr_i\). Thus the interval
\(0<r_i<c\) is forward invariant: either the modulus does not
increase or the squared formula applies. If a maximum \(r_i>1\)
occurred, then \(|D_i|_p\le1\) and \(V_i\) would strictly dominate
in the same sum, giving \(r_{i+1}=c/r_i<c\). Forward invariance
would contradict return to the maximum. Thus all \(r_i\le1\).
If any \(r_i<c\), all would be \(<c\) by invariance. For a step
with \(r_{i+1}\le r_i\) the three terms in the sum would all have
modulus at most one, forcing \(r_{i+1}\ge c\). Hence there could
be no nonincreasing step, impossible on a cycle. We have shown
\(c\le r_i\le1\). A strict increase would now give
\(r_{i+1}^2=cr_i\le r_i^2\). Thus all moduli agree, so
\(|D_i|_p=r_{i+1}/r_i=1\).
\end{proof}

\begin{theorem}\label{thm:cm-unit}
For the normalized lattice and \(a_1\) of
Proposition~\ref{prop:scaled-cm}, every \(p\) and every isomorphism
\(\iota:\C\to K_p\) satisfy
\[
              |\iota(a_1^4-64)|_p=1,\qquad |\iota(a_1)|_p\le1 .
\]
\end{theorem}
\begin{proof}
Starting at \(\Lambda=\lambda O_K\), the distinguished class of
Lemma~\ref{lem:half-lattice-cycle} is the previously chosen \(t\)
by Proposition~\ref{prop:scaled-cm}. Its cycle up to scaling is
finite by that lemma, and \(V,D\) are scale invariant. Apply
\(\iota\) to Proposition~\ref{prop:half-recurrence} around
the cycle, then apply Lemma~\ref{lem:cycle-unit}.
It remains to identify \(D\). Set
\(\alpha=e_T-e_t,\ b_1=e_S-e_t\). The derivative at \(e_t\) of
\((X-e_t)(X-e_T)(X-e_S)=X^3-(g_2/4)X-g_3/4\) gives
\(\alpha b_1=3e_t^2-g_2/4=1\), since \(\wp_\Lambda''(t)=2\).
Also \(\alpha+b_1=-3e_t=-a_1^2/4\). It follows that
\[
          a_1^4-64=16(\alpha-b_1)^2
                 =16(u-1)^2/u=D.
\]
The unit assertion follows. Adding \(64\), of \(p\)-adic absolute
value at most one, gives \(|\iota(a_1)|_p^4\le1\).
\end{proof}

\section{Doubling, torsion values and the integral coordinate}

\begin{proposition}\label{prop:cm-duplication}
Use \(\Lambda,\sigma=\sigma_\Lambda,\theta=\theta_\Lambda,t,\xi,\phi,a_1\)
of Propositions~\ref{prop:scaled-cm} and~\ref{prop:cm-theta}. Then
\[
 \begin{aligned}
 \theta(z)^2&=(\wp_\Lambda(z)-\wp_\Lambda(t))\sigma(z)^2,&
 \sigma(z)^2&=\xi(z)\theta(z)^2,\\
 \theta(2z)&=\theta(z)^4-\sigma(z)^4,\\
 \sigma(2z)^2&=(\theta(z)\sigma(z))^2
       \bigl\{4(\theta(z)^4+\sigma(z)^4)
                       +a_1^2(\theta(z)\sigma(z))^2\bigr\}.
 \end{aligned}
\]
Expressions with apparent poles are interpreted by continuation.
The function \(G=\phi\theta^3\) extends to an entire function and
satisfies \(G^2+a_1\sigma^2\theta G=\sigma^6+\sigma^2\theta^4\).
\end{proposition}
\begin{proof}
For \(l\in\Lambda\) the extra phase in the multiplier of
\(\theta^2/\sigma^2\) is \(\exp(4\pi i\Im(\bar l t)/A)\), by
Propositions~\ref{prop:reduced-sigma} and~\ref{prop:cm-theta}.
It equals one since \(2t\) is a period and the oriented area of two
periods is an integer multiple of \(A\). The quotient is therefore
periodic and is even with principal part
\(z^{-2}\) at zero. Subtracting \(\wp_\Lambda\) gives an entire
elliptic function, whose value at \(t\) is \(-\wp_\Lambda(t)\).
Lemma~\ref{lem:elliptic-count} proves the first identity.
As in the proof of Proposition~\ref{prop:half-recurrence}, the
product of \(\wp_\Lambda(z)-e_t\) and \(\wp_\Lambda(z+t)-e_t\)
is the constant \((e_T-e_t)(e_S-e_t)=1\) by
Theorem~\ref{thm:cm-unit} and its proof. This gives the second
identity of the first line.

For fixed \(w\notin\Lambda\) the quotient
\(\sigma(z+w)\sigma(z-w)/(\sigma(z)^2\sigma(w)^2)\)
is periodic by the sigma multiplier, is even in \(z\), and has
principal part \(-z^{-2}\) at zero. Add \(\wp_\Lambda(z)\)
and evaluate at \(z=w\) to get the constant \(\wp_\Lambda(w)\).
Thus the quotient equals \(\wp_\Lambda(w)-\wp_\Lambda(z)\).
Letting \(w\to z\) after division of the corresponding product identity
by \(z-w\) gives \(\sigma(2z)=-\wp_\Lambda'(z)\sigma(z)^4\),
since \(\sigma'(0)=1\). Substituting \(z+t\) for \(z\) in the
product identity and using Proposition~\ref{prop:cm-theta} yields
\[
 \theta(z+w)\theta(z-w)
       =\theta(z)^2\{\theta(w)^2-\xi(z)\sigma(w)^2\}.
\]
At \(w=z\) this gives the stated doubling formula for \(\theta\).
The cubic equation in the proof of Proposition~\ref{prop:scaled-cm}
at \(z\), in the coordinate \(X=\wp_\Lambda(z)-e_t\), reads
\(\wp_\Lambda'(z)^2=4X^3+a_1^2X^2+4X\). Substitute
\(X=\theta(z)^2/\sigma(z)^2\) into the squared duplication formula
for \(\sigma\). This gives the other formula. All resulting identities
extend over the excluded discrete sets by continuity.
Finally \(\phi\) has poles of order at most three at \(t+\Lambda\)
by its definition. The simple zeros of \(\theta\) there cancel
them. Multiplying the cubic for \(\phi,\xi\) by \(\theta^6\) and
using \(\xi\theta^2=\sigma^2\) gives the last identity.
\end{proof}

\begin{theorem}\label{thm:local-cm-height}
For a prime \(p\) and isomorphism \(\iota:\C\to K_p\), define
\[
          H_p(z)=\max\{|\iota(\sigma(z))|_p,|\iota(\theta(z))|_p\}.
\]
Then for all \(z\in\C\), \(H_p(z)>0,\ H_p(2z)=H_p(z)^4\), and
\(|\iota G(z)|_p\le H_p(z)^3\). At every torsion point \(y\)
of \(\C/\Lambda\) one has
\[
 \max\bigl(|\iota((U_y\sigma)(0))|_p,\,
           |\iota((U_y\theta)(0))|_p\bigr)=1 .
\]
\end{theorem}
\begin{proof}
The zeros of \(\sigma,\theta\) are disjoint by
Propositions~\ref{prop:reduced-sigma} and~\ref{prop:cm-theta}.
Choose whichever of their two images has larger absolute value and
call it \(c\). Divide the images by \(c\) to get \(a,b\) with
\(\max(|a|_p,|b|_p)=1\). Put \(A_1=\iota(a_1)\). By
Theorem~\ref{thm:cm-unit}, \(A_1\) is integral and
\(A_1^4-64\) is a unit. Applying Proposition~\ref{prop:cm-duplication}
to the ratios \(y_0=\iota(\theta(2z))/c^4\),
\(x_0=\iota(\sigma(2z))/c^4\) gives
\[
       y_0=b^4-a^4,\qquad
       x_0^2=a^2b^2\{4(b^4+a^4)+A_1^2a^2b^2\}.
\]
Both are integral by these equations. If both vanished in the
residue field, the first would give \(\bar a^4=\bar b^4\ne0\).
Putting \(r=(\bar a/\bar b)^2\) yields \(r^2=1\); the second gives
\(8+\bar A_1^2 r=0\). Squaring would give
\(\bar A_1^4=64\), contradicting the unit property, including in
residue characteristic two. Hence the maximum is one after division
by \(c^4\). For \(G/c^3\), meaning \(\iota(G(z))/c^3\), its monic
quadratic equation has integral coefficients by the last identity
of Proposition~\ref{prop:cm-duplication}. Its absolute value is
therefore at most one, by the dominant-term argument of
Lemma~\ref{lem:poly-test}.

Write \(A\) for the area of \(\Lambda\), and
\(h(y)=H_p(y)/|\iota(\exp(\pi|y|^2/(2A)))|_p\).
At \(y\in n^{-1}\Lambda\), replacing \(y\) by \(y+l\), \(l\in\Lambda\),
in this expression preserves \(h\). In fact the real part of the
exponent in both multipliers is cancelled by the change of
\(\pi|y|^2/(2A)\); the remaining phase is a root of unity because
\(\Im(\bar l y)/A\in n^{-1}\Z\). The extra phase for \(\theta\)
is a root of unity since \(2t\in\Lambda\). Every such root has
\(p\)-adic absolute value one by multiplicativity.
Thus \(h\) is defined on the finite group
\(n^{-1}\Lambda/\Lambda\), and \(h(2y)=h(y)^4\) by doubling.
Two iterates \(2^i y,2^{i'}y\), \(i<i'\), agree in that group, so
\(h(y)^{4^i}=h(y)^{4^{i'}}\). Positivity forces \(h(y)=1\).
The definition of \(U_y\) at zero identifies this with the required
maximum.
\end{proof}

\begin{proposition}\label{prop:cm-parameter}
Use the coefficientwise Taylor map at zero through \(\iota\) of
Proposition~\ref{prop:cm-theta}, and write \(\Phi(Z),\Xi(Z)\)
for the resulting Taylor series of \(\phi,\xi\). Put
\(A_1=\iota(a_1)\) and let \(O\) be the valuation ring of \(K_p\).
There is a unique \(x(X)\in XO[[X]]\) satisfying
\[
              x=X^2+A_1Xx-x^3 .
\]
It belongs to \(X^2O[[X]]\) and \(\Xi=x(\Phi)\).
The subring \(O[[\Phi]]=\{F(\Phi(Z)):F\in O[[X]]\}\)
of \(K_p[[Z]]\) is closed under formal differentiation in \(Z\).
For \(F\in O[[\Phi]]\) and every \(j\ge0\), \(j![Z^j]F\in O\).
Here \([Z^j]\) extracts a coefficient at the indicated degree.
\end{proposition}
\begin{proof}
Theorem~\ref{thm:cm-unit} gives \(A_1\in O\). The map
\(y\mapsto X^2+A_1Xy-y^3\) increases by at least one the order
of the difference of two elements of \(XO[[X]]\), since
\(y^3-z^3=(y-z)(y^2+yz+z^2)\). Starting at zero, the iterates
stabilize coefficientwise to a fixed point in \(X^2O[[X]]\).
The same increase shows uniqueness over \(K_p\) as well.
By Proposition~\ref{prop:scaled-cm}, \(\Phi(Z)=Z+O(Z^2)\).
It has a compositional inverse: in a proposed inverse starting
with \(Z\), at degree \(n\ge2\) exactly the linear term of \(\Phi\)
contributes the unknown coefficient with multiplier one.
This gives a right inverse by recursion and the same uniqueness
gives a two-sided inverse. Substitution in the cubic equation proves
\(\Xi=x(\Phi)\) by uniqueness.

Substituting the Taylor series in the identities of
Proposition~\ref{prop:scaled-cm} gives
\(\Xi'=2\Phi+A_1\Xi\). Differentiating the cubic then gives
\((2\Phi+A_1\Xi)(\Phi'+A_1\Phi-3\Xi^2-1)=0\).
The first factor starts with \(2Z\), so the other factor vanishes.
Hence \(\Phi'=1+3x(\Phi)^2-A_1\Phi\) belongs to \(O[[\Phi]]\).
For \(F_0\in O[[X]]\) one has
\((F_0(\Phi))'=F_0'(\Phi)\Phi'\), by differentiating monomials
coefficientwise. This proves closure. All members have integral
constant coefficient since \(\Phi(0)=0\). Taking the constant
coefficient after \(j\) differentiations proves the last assertion.
\end{proof}

\section{Removing torsion translates from the local expansion}

\begin{lemma}\label{lem:section-polynomials}
For a nonnegative integer \(m\), suppose \(F\) is entire and
\(F/\theta^m\) is \(\Lambda\)-periodic for the functions of
Proposition~\ref{prop:cm-duplication}. There are unique complex
polynomials \(P,Q\) with
\[
          F/\theta^m=P(\xi)+\phi Q(\xi).
\]
Only terms of the form \(\xi^a\phi^b\) with \(b=0,1,\ 2a+3b\le m\)
occur. For each \(\iota:\C\to K_p\) there is a finite constant
\(C_F\) such that \(|\iota(F(z))|_p\le C_F H_p(z)^m\) for all
\(z\).
\end{lemma}
\begin{proof}
The quotient can have poles only at \(t+\Lambda\), of order at
most \(m\). By Proposition~\ref{prop:scaled-cm}, \(\xi\) and \(\phi\)
have leading Laurent terms \((z-t)^{-2}\) and \(-(z-t)^{-3}\) there.
The integers \(2a+3b\) with \(a\ge0,\ b\in\{0,1\}\) attain zero and
each integer at least two exactly once. Successively subtract the
unique scalar multiple of the monomial that cancels a pole of any
such order, starting at the largest. A possible simple pole also
has to disappear: periodicity makes the boundary integral of the
remaining function zero, and its sole possible residue is therefore
zero by the disk excision in Lemma~\ref{lem:elliptic-count}.
The remaining entire elliptic function is constant. The distinct
pole orders give uniqueness. Off the zeros of \(\theta\), multiplying
a monomial by \(\theta^m\) gives
\(\theta^{m-2a-3b}\sigma^{2a}G^b\) by
Proposition~\ref{prop:cm-duplication}. This identity extends
everywhere. Theorem~\ref{thm:local-cm-height} bounds it by
\(H_p(z)^m\) at \(\iota\). The maximum of the absolute values of
the finitely many coefficients, or one if all vanish, suffices
for \(C_F\).
\end{proof}

\begin{proposition}\label{prop:tate-torsion}
Let \(H_p\) be as in Theorem~\ref{thm:local-cm-height} and write
\(g_y(z)=\exp(-\pi(z\bar y+|y|^2/2)/A)\).
For every \(y\in n^{-1}\Lambda\) with \(n\ge1\), and every \(z\in\C\),
\[
  |\iota(g_y(z))|_pH_p(z+y)=H_p(z),\qquad
                         H_p(nz)\le H_p(z)^{n^2}.
\]
\end{proposition}
\begin{proof}
Write \(R_y(z)\) for the left side of the first identity divided
by its right side. It is unchanged if \(y\) is replaced by \(y+l\)
with \(l\in\Lambda\). Indeed the two transformation formulas of
Propositions~\ref{prop:reduced-sigma} and~\ref{prop:cm-theta}
give the same absolute factor on \(H_p(z+y)\), since their
multipliers differ by a sign. After multiplication by
\(g_{y+l}(z)/g_y(z)\), only \(\exp(\pi i\Im(\bar l y)/A)\)
remains, a root of unity at such a torsion point. For
\(f=\sigma,\theta\) the quotient
\((U_y f)^{2n}/\theta^{2n}\) is periodic: under \(z\mapsto z+l\)
the quotient before raising to \(2n\) picks up at most a sign
times \(\exp(2\pi i\Im(\bar l y)/A)\). By
Lemma~\ref{lem:section-polynomials}, \(R_y\) is bounded on all
\(z\), using the larger of the two bounds and taking a \(2n\)-th
root. There are only \(n^2\) classes of \(y\), so a common finite
bound \(C\ge1\) works for all of them.

Now \(R_{2y}(2z)=R_y(z)^4\), by
Theorem~\ref{thm:local-cm-height} and \(g_{2y}(2z)=g_y(z)^4\).
Hence \(R_y(z)^{4^k}\le C\) for all \(k\ge0\), giving \(R_y(z)\le1\).
Applying the same inequality at \(-y,z+y\) yields the reverse bound,
since \(g_{-y}(z+y)g_y(z)=1\). For the second assertion consider
\(S_n(z)=H_p(nz)/H_p(z)^{n^2}\). For \(f=\sigma,\theta\) the function
\(f(nz)^2/\theta(z)^{2n^2}\) is periodic: a period \(l\) at \(z\)
gives the period \(nl\) at \(nz\) with the same exponential raised to
\(n^2\), and squaring removes the signs. Thus \(S_n\) is bounded by
the same lemma. Since \(S_n(2z)=S_n(z)^4\), iteration gives \(S_n(z)\le1\).
\end{proof}

\begin{proposition}\label{prop:bounded-sections-integral}
In Lemma~\ref{lem:section-polynomials}, suppose for a given \(\iota\)
that \(|\iota(F(z))|_p\le H_p(z)^m\) for every \(z\).
All coefficients of \(\iota P,\iota Q\) then belong to the valuation
ring \(O\). In particular, writing \(F_0(Z)\) for the transported
Taylor series of \(\theta\) at zero, the transported Taylor series
of \(F\) belongs to \(F_0(Z)^m O[[\Phi]]\).
\end{proposition}
\begin{proof}
Suppose a coefficient of \(\iota P,\iota Q\) of largest
absolute value \(C\) satisfies \(C>1\). Divide both polynomials
by that coefficient to obtain integral \(P_1,Q_1\) whose reductions
are not both zero. In the residue field put
\(D(X)=4X^3+\bar A_1^2X^2+4X\), \(A_1=\iota(a_1)\).
In odd characteristic this polynomial is nonzero from its cubic term.
In characteristic two it is nonzero because \(\bar A_1\ne0\)
by Theorem~\ref{thm:cm-unit}. The residue field is infinite by
Lemma~\ref{lem:poly-test}. Choose \(x_0\) avoiding the zeros of
\(D\) and of one nonzero polynomial from \(\bar P_1,\bar Q_1\).
Lift \(x_0\) to \(x\in O\). Both roots of
\[
                Y^2+A_1xY-x^3-x=0
\]
exist in \(K_p\) and are integral by their monic equation.
Their squared difference has reduction \(D(x_0)\ne0\), so their
reductions are distinct. At least one root \(y\) gives
\(|P_1(x)+yQ_1(x)|_p=1\), since an affine linear polynomial in
the residue field nonzero at some point cannot vanish at two
different points.

This pair is attained by \(\iota\xi(z),\iota\phi(z)\) for some
\(z\notin t+\Lambda\). In fact \(\xi-\iota^{-1}x\) has a double pole
at \(t\) and therefore has a zero \(z\) by
Lemma~\ref{lem:elliptic-count}. The value of \(\xi\) at \(-z\)
is the same by periodicity and evenness of \(\wp\), while its
derivative changes sign. Thus the two values of \(\iota\phi\) at
\(z,-z\) supply both roots by \(\xi'=2\phi+a_1\xi\) and their
nonzero difference in characteristic zero. Since \(x\in O\),
\(\sigma^2=\xi\theta^2\) in Proposition~\ref{prop:cm-duplication}
gives \(H_p(z)=|\iota(\theta(z))|_p>0\). The chosen pair would give
\(|\iota(F(z)/\theta(z)^m)|_p=C>1\), contradicting the assumption.
This proves integrality. The series conclusion follows by substituting
\(\Xi=x(\Phi)\) from Proposition~\ref{prop:cm-parameter} in the
polynomial expression.
\end{proof}

\begin{theorem}\label{thm:remove-shifts}
Let \(n\ge1,\ p\nmid n,\ ny\in\Lambda\), and retain the series
\(F_y,F_0,\Phi\) of Propositions~\ref{prop:cm-theta}
and~\ref{prop:cm-parameter}. There exists \(V_y\in O[[X]]\) with
\[
                         F_y(nZ)=F_0(Z)^{n^2}V_y(\Phi(Z)).
\]
Consequently, suppose at this \(p\) that \(\Phi(\varpi X)\in O[[X]]\)
and that \(F_0(\varpi X)=\exp(\kappa\psi(X))V_0'(X)\), with
\(V_0'\in O[[X]]\) and the common data \(S,\psi,\varpi,\theta,\kappa\)
of Theorem~\ref{thm:product-null}. Every product of the actual torsion
translates in Proposition~\ref{prop:cm-theta} with denominator \(n\)
then satisfies the conclusions of Theorem~\ref{thm:product-null}
for all \(B\).
\end{theorem}
\begin{proof}
The quotient \((U_y\theta)(nz)/\theta(z)^{n^2}\) is periodic.
The prime-to-\(p\) condition is immaterial for this first calculation.
In the two multipliers a period \(l\) produces the same exponential
of \(n^2(\bar l z+|l|^2/2)\pi/A\). The remaining factor from the
translation by \(y\) is \(\exp(2\pi i\Im(\bar l ny)/A)=1\).
For the sign, write \(\chi(l)=s(l)\exp(2\pi i\Im(\bar l t)/A)\).
One has \(s(nl)=s(l)^{n^2}\) by parity, and the extra phase for
\(t\) is a sign, whose \(n\)-th and \(n^2\)-th powers agree.
Thus \(\chi(nl)=\chi(l)^{n^2}\). By
Proposition~\ref{prop:tate-torsion},
\[
 |\iota((U_y\theta)(nz))|_p\le H_p(nz)\le H_p(z)^{n^2}.
\]
Proposition~\ref{prop:bounded-sections-integral} with \(m=n^2\)
now proves the series equality.

Substitute \(Z=\varpi X/n\). The integral factor remains an
integral series in \(X/n\) since \(\Phi(\varpi X)\) is integral
with zero constant term. The common exponential for all the shifts
in the product is \(\exp(n^2\kappa\,\psi(X/n))\), and the remaining
factors are integral series in \(X/n\). In a tight frame replace
\(\beta=\iota b'\) by \(\beta/n\) when forming the linear arguments,
and replace \(\kappa\) by \(n^2\kappa\). These operations preserve
the required integrality and the prescribed valuation since \(p\nmid n\).
Expansion of the product gives the common exponential of the sum:
the formal identity \(\exp(X)\exp(Y)=\exp(X+Y)\) follows by the
binomial coefficients in each degree for commuting series of zero
constant term. This gives exactly the hypotheses of
Theorem~\ref{thm:product-null}.
\end{proof}

\section{Multiplication on the integral local ring}

\begin{proposition}\label{prop:cm-multiplication}
For \(0\ne\alpha\in O_K\) put \(q'=\alpha\bar\alpha\in\Z_{\ge1}\)
and \(a=\iota(\alpha)\). Then \(H_p(\alpha z)\le H_p(z)^{q'}\) for
every \(z\). In the notation of Theorem~\ref{thm:remove-shifts}, put
\[
       E_\alpha(Z)=F_0(aZ)/F_0(Z)^{q'},\qquad
                         f_\alpha(\Phi(Z))=\Phi(aZ).
\]
One has \(E_\alpha\in O[[\Phi]]^\times,\ f_\alpha(X)\in XO[[X]]\).
If \(q'\) is odd, then the reduction of \(f_\alpha\) to the
residue field is nonzero of order at most \(q'\).
\end{proposition}
\begin{proof}
Let \(\chi(l)=s(l)\exp(2\pi i\Im(\bar l t)/A)\) be the multiplier
apart from the standard exponential in Proposition~\ref{prop:cm-theta}.
For \(l=\lambda(m+n\omega)\) one has
\(\chi(l)=(-1)^{m+mn}\) by the computation of \(s(l)\) in
Proposition~\ref{prop:reduced-sigma} and \(t=\lambda/2\).
This is \((-1)^{|l/\lambda|^2}\), since
\(|m+n\omega|^2=m^2+mn+((1+d)/4)n^2\) and \((1+d)/4\) is even.
Consequently \(\chi(\alpha l)=\chi(l)^{q'}\).
Moreover the basic exponential multiplier at \(\alpha z,\alpha l\)
is the one at \(z,l\) raised to \(q'\).
As in Proposition~\ref{prop:tate-torsion}, this makes
\(f(\alpha z)^2/\theta(z)^{2q'}\) periodic for \(f=\sigma,\theta\).
Lemma~\ref{lem:section-polynomials} bounds
\(H_p(\alpha z)/H_p(z)^{q'}\); the value at \(2z\) equals its
fourth power. Iterating gives the required inequality.

The functions \(\theta(\alpha z)\) and \(G(\alpha z)\) have ratios
with \(\theta(z)^{q'}\), respectively \(\theta(z)^{3q'}\), that are
periodic, since \(\phi\) is elliptic and the signs for \(\theta\)
agree before squaring. They satisfy the bounds required by
Proposition~\ref{prop:bounded-sections-integral}, by the first
assertion and Theorem~\ref{thm:local-cm-height}. The first ratio
has constant term one at zero and hence an inverse in \(O[[\Phi]]\),
by recursively solving the product coefficients. Dividing the
second ratio by the cube of the first gives \(f_\alpha\) in
\(O[[X]]\); its constant term is zero.

For \(q'\) odd the element \(\alpha\) acts as the identity on
\(\Lambda/2\Lambda\): its class in \(O_K/2O_K\simeq
\mathbf F_2\times\mathbf F_2\) is a unit since
\(\alpha\bar\alpha=q'\), and the only unit is \((1,1)\).
The ring decomposition follows directly from the two distinct
roots of the quadratic for \(\omega\) modulo 2. In particular
\(\alpha t=t\bmod\Lambda\).
The ratio \(\sigma(\alpha z)^2/\theta(z)^{2q'}\) is even.
By Proposition~\ref{prop:bounded-sections-integral} it is
\(P(\xi(z))\) with \(\deg P\le q'\) and \(\iota P\) integral.
Indeed its possible term \(\phi Q(\xi)\) must vanish: \(\xi\) is
even, and \(\phi(z)-\phi(-z)=\xi'(z)\) is not identically zero.
Let \(P_{q'}\) be the degree \(q'\) coefficient, taken as zero
if necessary. Since \(\xi\theta^2=\sigma^2\), division by
\(\xi^{q'}\) and passage to \(z=t\) give
\[
            P_{q'}=\sigma(\alpha t)^2/\sigma(t)^{2q'} .
\]
At the two torsion points \(t,\alpha t\) the theta value vanishes.
The torsion formula in Theorem~\ref{thm:local-cm-height}, together
with \(|\alpha t|^2=q'|t|^2\), gives
\(H_p(\alpha t)=H_p(t)^{q'}\). Hence \(|\iota P_{q'}|_p=1\).

In the coordinate \(X=\Phi(Z)\), the polynomial identity now reads
\[
            (\iota P)(x(X))=x(f_\alpha(X))\,E_\alpha(Z)^2 ,
\]
where \(x(X)=X^2+\cdots\) is the integral series of
Proposition~\ref{prop:cm-parameter}. The last factor is a unit
in this coordinate. After reduction, the left side has order
\(2\operatorname{ord}_0(\overline{\iota P})\le2q'\), since its
leading polynomial coefficient is a unit. The right side has order
\(2\operatorname{ord}_0(\bar f_\alpha)\) if \(\bar f_\alpha\ne0\)
and otherwise vanishes. Comparing gives the assertion.
\end{proof}

\begin{proposition}\label{prop:theta-log}
Define \(\mathcal L=-\log F_0\) by the formal expansion of
\(-\log(1+u)\) at \(u=F_0-1\). Then
\(\mathcal L(0)=\mathcal L'(0)=0\) and for some \(c_2\in K_p\)
\[
       \mathcal L''=\Xi+c_2,\qquad
                              v_p\bigl(\iota(\sqrt{-d})c_2\bigr)\ge0 .
\]
For every \(\alpha\) of Proposition~\ref{prop:cm-multiplication},
the derivative of \(q'\mathcal L(Z)-\mathcal L(aZ)\) belongs
to \(O[[\Phi]]\). The same integral-coordinate conclusions about
\(\Phi(aZ)\), \(\exp(v\mathcal L(Z)-\mathcal L(aZ))\) and the
derivative of its exponent hold for every \((a,v)\in O^2\) in
the closure of \(\{(\iota\alpha,\alpha\bar\alpha):0\ne\alpha\in O_K\}\).
Here closure is for the topology of the \(p\)-adic absolute value.
\end{proposition}
\begin{proof}
The series is well defined since \(F_0(0)=1\). Evenness of \(\theta\)
gives vanishing of the linear term. The function
\(-(\sigma'/\sigma)'-\wp_\Lambda\) is elliptic by the sigma multiplier
and has no poles, by the local expansion at each simple zero of
\(\sigma\). It is constant by Lemma~\ref{lem:elliptic-count}.
In the formula for \(\theta\) from Proposition~\ref{prop:cm-theta}
the additional exponential has linear logarithm. Taking two
logarithmic derivatives therefore gives
\(\mathcal L''(Z)-\mathcal L''(0)=\Xi(Z)\).
All these expressions use transported Taylor coefficients of complex
germs. This commutes with addition, multiplication and differentiation
by the power-series rules, and the derivative of the formal logarithm
is the logarithmic derivative by its coefficient expansion.

By Proposition~\ref{prop:cm-multiplication}, \(E_\alpha\) is integral
and invertible in the local coordinate. Its logarithmic derivative
is the derivative of \(q'\mathcal L-\mathcal L(aZ)\), hence integral
in the same coordinate by Proposition~\ref{prop:cm-parameter}.
Differentiate once more and evaluate at zero to see that
\((\alpha\bar\alpha-a^2)c_2\in O\), with the integer \(\alpha\bar\alpha\)
regarded in \(K_p\). Use the two choices \(\omega,\bar\omega\).
The resulting multipliers are
\(-\iota(\omega\sqrt{-d})\) and \(\iota(\bar\omega\sqrt{-d})\).
Subtracting the first from the second gives the asserted bound since
\(\omega+\bar\omega=1\).

To pass to the closure substitute the formal inverse of \(\Phi\) for
\(Z\) in \(\Phi(aZ)\), the stated exponential and the derivative of
its exponent. For fixed degree the coefficient is a polynomial in
\(a,v\) over \(K_p\): each composition uses finitely many powers in
that degree and the exponent has constant term zero. It depends
continuously on the pair. The valuation ring is closed by the
ultrametric inequality, so integrality persists in the closure.
The exponential has constant term one and is again invertible there.
\end{proof}

\section{A formal coordinate at the nonsplit primes}

\begin{lemma}\label{lem:teich-pairs}
Let \(p\) be a prime with \(\eta_d(p)\ne1\) in
Proposition~\ref{prop:places}. Set \(e=2\) if \(p\) is inert and
\(e=1\) if \(p=d\), and put \(q'=p^e,\ Q=q'-1\). Then \(p\ge13\).
There is a primitive \(Q\)-th root of unity \(\zeta\in K_p\), and
for every integer \(i\) the pair
\((\zeta^i,\zeta^{(p+1)i})\) belongs to the closure in
Proposition~\ref{prop:theta-log}.
\end{lemma}
\begin{proof}
The primes \(2,3,5,7,11\) divide \(18480\mid M\), hence all split
by Proposition~\ref{prop:places}. The polynomial \(X^Q-1\) has
\(Q\) distinct roots in the algebraically closed field of
Lemma~\ref{lem:place-fields} by its derivative. The cyclicity
argument of Lemma~\ref{lem:count-labels} gives \(\zeta\).
The reduction map on these roots is injective. Indeed if
\(w\equiv1\) in the residue field then \(1+\cdots+w^{Q-1}\)
reduces to \(Q\ne0\), so \(w^Q-1=0\) forces \(w=1\).
Write \(\xi=\zeta^i\), and use bars on residue classes in this proof.

The \(q'\) roots of \(X^{q'}-X\) in the residue field form a
field of that cardinality by the argument of Lemma~\ref{lem:count-labels}.
The reduced image of \(\iota O_K\) is this field. For \(p\) inert,
the quadratic for \(\omega\) over \(\mathbf F_p\) is irreducible
and Frobenius takes its root to \(1-\bar\omega\), as in
Lemma~\ref{lem:inert-prime-fibres}. For \(p=d\), the identity
\((2\omega-1)^2=-d\) reduces \(\omega\) to \(1/2\).
It follows in both cases that a lift \(\alpha\in O_K\) of
\(\bar\xi\) has complex conjugate reducing to \(\bar\xi^p\).
The lift is nonzero. Thus \(\iota\alpha/\xi\) and
\((\alpha\bar\alpha)/\xi^{p+1}\) reduce to one.

If \(v_p(u-1)=v>0\), the binomial theorem gives
\(v_p(u^p-1)\ge\min(v+1,pv)\), with valuation \(+\infty\) at zero.
While finite, the successive valuations increase by at least
\(\min(1,(p-1)v)>0\). Thus \(u^{q'^j}\to1\).
Since \(\xi^{q'}=\xi\), the norm pairs of \(\alpha^{q'^j}\)
tend to \((\xi,\xi^{p+1})\), as claimed.
\end{proof}

\begin{theorem}\label{thm:frob-parameter}
Use the notation of Lemma~\ref{lem:teich-pairs} and set
\(h=e/2\), so \(h=\omega_p\). There are \(\varpi,\pi\in O\)
and a formal coordinate \(T\in O[[\Phi]]\) with
\[
 v_p(\varpi)=\frac{h}{Q}=\rho_p,\qquad
 \pi=\varpi^Q,\qquad O[[T]]=O[[\Phi]] .
\]
Write \(Z_T\) for the inverse of \(T\) under substitution. Then
\[
 T(Z)=Z+O(Z^{q'}),\quad
 Z_T(X)=X+\sum_{s\ge1}c_sX^{1+Qs},\quad
 Z_T'\in O[[X]],\qquad c_1=-\pi^{-1}.
\]
The series \(S(X)=T(\varpi X)/\varpi\) belongs to \(XO[[X]]\) and
\[
              S-S^{q'}\equiv X\pmod{\pi O[[X]]},\qquad
                  \Phi(\varpi X)\in XO[[X]] .
\]
\end{theorem}
\begin{proof}
Proposition~\ref{prop:theta-log} and Lemma~\ref{lem:teich-pairs}
give \(\Phi(\zeta^iZ)\in O[[\Phi]]\). Since \(p\nmid Q\),
\[
       T(Z)=Q^{-1}\sum_{i=0}^{Q-1}\zeta^{-i}\Phi(\zeta^iZ)
\]
belongs to this ring. The sum of the \(Q\) powers of a nontrivial
\(Q\)-th root is zero by multiplication by the root minus one.
Thus \(T\) selects exactly the terms of \(\Phi\) of degrees
\(1\bmod Q\). It starts with \(Z\). Write \(T=B(\Phi)\) with
\(B\in XO[[X]]\). The coefficient of \(X\) in \(B\) is one,
so the recursive inversion of Proposition~\ref{prop:cm-parameter}
gives \(B^{-1}\in XO[[X]]\) and the asserted equality of rings.
The equality \(T(\zeta Z)=\zeta T(Z)\) gives
\(Z_T(\zeta X)=\zeta Z_T(X)\), yielding the stated support.
Closure under differentiation in Proposition~\ref{prop:cm-parameter}
gives \(T'\in O[[T]]\) with constant term one. Differentiating
\(T(Z_T(X))=X\) therefore gives \(Z_T'=1/T'(Z_T)\in O[[X]]\).
In particular \(v_p(c_s)\ge-v_p(1+Qs)\).

Choose \(\alpha=p\) at an inert prime and \(\alpha=\sqrt{-d}\) at
\(p=d\). Then \(\alpha\bar\alpha=q'\) is odd and \(A=\iota\alpha\)
has valuation \(h\). By Proposition~\ref{prop:cm-multiplication},
\(T(AZ_T(X))=B(f_\alpha(B^{-1}(X)))\) is integral and its
reduction has order at most \(q'\). Indeed a substitution with
unit linear coefficient preserves the first nonzero degree, on
both sides of the composition. This series has only degrees
\(1\bmod Q\) by the covariance under \(\zeta\). Its linear term
reduces to zero. Thus its coefficient of degree \(q'\) is a unit.
The first two terms of \(T\) are \(Z-c_1Z^{q'}\) through that
degree by inversion. It follows that \(c_1(A-A^{q'})\) is a unit,
so \(v_p(c_1)=-h\). Take a \(Q\)-th root
\(\varpi\) of \(-1/c_1\) by Lemma~\ref{lem:place-fields}.
The geometric series in Definition~\ref{def:linear-jet} gives
\(\rho_p=h/Q\). Its positive-exponent sum is \(p/(p^2-1)\)
at inert primes and \(1/(p-1)\) at the ramified prime.

For \(s\ge2\) we have \(v_p(1+Qs)\le(s-1)/2\).
Indeed it vanishes unless \(s\equiv1\bmod p\). In the latter
case \(s\ge p+1\) and \(1+Qs\le sp^2\le p^{(s-1)/2}\).
For the second inequality, at \(s=p+1\) use \(p+1\le p^2\)
and \(p/2\ge4\). Increasing \(s\) multiplies its right side by
\(\sqrt p\) and its left side by \(1+1/s\le2\).
Hence for \(s\ge2\)
\[
 v_p(c_s\varpi^{Qs})\ge sh-(s-1)/2\ge h .
\]
Consequently \(Z_T(\varpi X)/\varpi\) is integral and is congruent
to \(X-X^{q'}\) modulo \(\pi\). Its linear coefficient is one.
Its integral compositional inverse is exactly \(S\). Composing
the congruence with \(S\) gives the first assertion about \(S\).
Finally \(\Phi(\varpi X)=B^{-1}(\varpi S(X))\), which is integral.
\end{proof}

\section{The logarithm of the normalized theta function}

\begin{lemma}\label{lem:project-log}
Retain \(\zeta,Q,T,\mathcal L\) from Lemma~\ref{lem:teich-pairs},
Theorem~\ref{thm:frob-parameter} and Proposition~\ref{prop:theta-log}.
Write \(P_j\) for the projection of a formal series onto degrees
congruent to \(j\bmod Q\), and \(\gamma=p+1\).
Then \(\exp(-\mathcal L+P_\gamma\mathcal L)\)
belongs to \(O[[T]]^\times\) and has constant term one.
\end{lemma}
\begin{proof}
Set \(w=\zeta^\gamma\). In both cases \(w^{p-1}=1\), since
\(Q\mid p^2-1\). For \(0\le i<Q\),
\[
       D_i=\exp(w^i\mathcal L(Z)-\mathcal L(\zeta^iZ))
\]
belongs to \(O[[T]]^\times\) by Proposition~\ref{prop:theta-log},
Lemma~\ref{lem:teich-pairs} and the equality of formal rings in
Theorem~\ref{thm:frob-parameter}. It has constant term one.

Each \(b_i=-w^{-i}/Q\) lies in the closure of the ordinary integers
in \(K_p\). Indeed the residue of \(w^{-i}\) is a root of
\(X^{p-1}-1\) in \(\mathbf F_p\), since that polynomial already
has all its roots among the nonzero elements of \(\mathbf F_p\).
Choose a positive integer \(a\) lifting that root. The binomial
estimate in Lemma~\ref{lem:teich-pairs} gives
\(a^{p^j}\to w^{-i}\). Also, for each \(j\) Bezout's identity in
Lemma~\ref{lem:basic-arith} gives an integer inverse to \(Q\bmod p^j\).
Its difference from \(1/Q\) has valuation at least \(j\), proving
the claim. If \(b\) is a limit of integers, every
\(\binom b k=b(b-1)\cdots(b-k+1)/k!\) belongs to \(O\), since this
is a continuous polynomial for fixed \(k\) and its values at
integers are integers. For a negative integer this follows from
\(\binom{-m}k=(-1)^k\binom{m+k-1}k\) for \(m\ge1\).
Thus \(D_i^{b_i}=\sum_{k\ge0}\binom{b_i}k(D_i-1)^k\) belongs to
\(O[[T]]^\times\). This binomial series equals
\(\exp(b_i\log D_i)\). In an indeterminate \(u\) the binomial series
and \(\exp(b_i\log(1+u))\) both satisfy \((1+u)F'=b_i F,\ F(0)=1\);
coefficient recursion gives uniqueness, and substitution of \(D_i-1\)
is valid. Formal logarithm and exponential are inverse on series
with the indicated constant terms, by differentiation and those
constant terms. The geometric sum proof in
Theorem~\ref{thm:frob-parameter} gives
\(P_\gamma\mathcal L=Q^{-1}\sum_i w^{-i}\mathcal L(\zeta^iZ)\).
Adding the exponents in the product \(\prod_i D_i^{b_i}\)
now gives \(-\mathcal L+P_\gamma\mathcal L\).
\end{proof}

\begin{lemma}\label{lem:two-coefficients}
Use \(q',Q,h,\pi,\varpi,Z_T\) of Theorem~\ref{thm:frob-parameter}.
Put \(\kappa_0=\varpi^{p+1}/p\). Let \(m,n\in K_p[[X]]\) satisfy
\(m\in X^3K_p[[X]],\ n\in X^2K_p[[X]],\ n'\in O[[X]]\) and
\(m'=n Z_T'\). Then for some \(b_1,b_2\in O\),
\[
 \kappa_0^{-1}(P_{p+1}m)(\varpi X)
       \equiv b_1X^{p+1}+b_2X^{p+q'}\pmod{\pi O[[X]]}.
\]
In particular the series on the left is integral. Here \(b_1,b_2\)
may be taken to be its actual coefficients at the displayed degrees.
\end{lemma}
\begin{proof}
Write \(n_i=[X^i]n\) and \(w_J=[X^{QJ}]Z_T'\). All \(w_J\) are
integral, \(w_0=1\) and \(w_1=q'c_1\) has valuation \(e-h\)
by Theorem~\ref{thm:frob-parameter}. For \(i\ge1\) the coefficient
\(i n_i\) of \(n'\) is integral. The only contributing degrees
\(k\ge3\) are \(k=p+1+QI\) for integers \(I\ge0\). In the
ramified case degree two also lies in the congruence class and
has zero coefficient by hypothesis. At a contributing degree
\[
       k[X^k]m=\sum_{J=0}^I n_{p+Q(I-J)} w_J .
\]
Indeed the derivative of \(Z_T\) has support in the multiples
of \(Q\). For \(J>I\) the index of \(n\) is negative or else is
one at \(J=I+1\) in the ramified case, where \(n_1=0\).
Thus each term, after division by \(k\) and scaling, has valuation
at least
\[
       1+hI-v_p(k)-v_p\bigl(p+Q(I-J)\bigr)+v_p(w_J).
\]
For \(I=0\) this is zero. For \(I=1\) one has \(v_p(k)=1\)
because \(k=p+q'\) and \(p\ge13\). For \(J=0\) the other integer
is prime to \(p\), giving lower bound \(h\). For \(J=1\) it is
\(p\), and \(v_p(w_1)=e-h\) gives lower bound \(e-1\ge0\).
This proves integrality at the two displayed degrees.

For \(2\le I<p\), \(k\equiv1-I\bmod p\) is a unit. The other
integer is prime to \(p\) except when \(J=I\), in which case it
is \(p\). The lower bound is therefore at least \(hI\ge h\).
For \(I\ge p\) we have \(k^2\le p^{(I+1)/2}\). At \(I=p\)
one has \(k\le p^3+1<2p^3\) and \(4p^6\le p^{(p+1)/2}\).
Increasing \(I\) multiplies \(k^2\) by at most
\((1+1/I)^2<\sqrt p\), proving the bound by induction.
The two positive integers \(k,p+Q(I-J)\) have product at most
\(k^2\), so the sum of their valuations is at most \((I+1)/2\).
The displayed lower bound is at least
\(1+hI-(I+1)/2\ge h\), since \(h\ge1/2\).
This proves divisibility by \(\pi\) at all other degrees.
\end{proof}

\begin{theorem}\label{thm:actual-theta-normal}
Let \(p,\varpi,\pi,S,q'\) be as in Theorem~\ref{thm:frob-parameter}
and \(F_0\) the transported Taylor series of the normalized theta
function of Proposition~\ref{prop:cm-theta}. There are
\(\psi\in XO[[X]],\ V\in O[[X]]^\times,\ u_0,u_1,u_2\in O\)
and \(\kappa\in K_p^\times\) such that
\[
 \begin{split}
 v_p(\kappa)&=1/(p-1)-\omega_p,\qquad
 F_0(\varpi X)=\exp(\kappa\psi(X))V(X),\\
 \psi(X)&\equiv u_0X^2+u_1S(X)^{p+1}
                          +u_2S(X)^{p+q'}\pmod{\pi O[[X]]}.
 \end{split}
\]
One can take \(u_0=0\) when \(p\) is inert. At every nonsplit \(p\),
the theta products of Proposition~\ref{prop:cm-theta} with any
common torsion denominator prime to \(p\) satisfy the integrality
conclusion of Theorem~\ref{thm:product-null} for every \(B\).
\end{theorem}
\begin{proof}
Take \(\kappa=-\kappa_0\) as in Lemma~\ref{lem:two-coefficients}.
Then \(v_p(\kappa)=(p+1)h/Q-1=1/(p-1)-h\) in both cases.
Write \(d_2=[Z^2]\mathcal L\) and
\(\mathcal L_1=\mathcal L-d_2 Z^2\).
By Proposition~\ref{prop:theta-log}, \(\mathcal L_1''=\Xi\);
it starts in degree four because \(\Xi\) starts in degree two
and the constant and linear terms of \(\mathcal L\) vanish.
Set \(m=\mathcal L_1(Z_T)\) and \(n=\mathcal L_1'(Z_T)\).
The chain rule gives \(m'=nZ_T'\) and \(n'=\Xi(Z_T)Z_T'\).
Proposition~\ref{prop:cm-parameter} and
Theorem~\ref{thm:frob-parameter} make the latter integral.
The orders of \(m,n\) are at least four and three. Thus
\(M(X)=\kappa_0^{-1}(P_{p+1}m)(\varpi X)\) has the form in
Lemma~\ref{lem:two-coefficients}.

The relation \(Z_T(\zeta X)=\zeta Z_T(X)\) gives
\(P_{p+1}m=(P_{p+1}\mathcal L_1)(Z_T)\) by the root-of-unity
average in Lemma~\ref{lem:project-log}. Since
\(Z_T(\varpi S(X))=\varpi X\), it follows that for
\(\psi(X)=\kappa_0^{-1}(P_{p+1}\mathcal L)(\varpi X)\)
we have
\[
                \psi(X)=u_0X^2+M(S(X)).
\]
Indeed, at inert \(p\), \(2\not\equiv p+1\bmod(p^2-1)\), giving
\(u_0=0\). At \(p=d\), degree two is retained and gives
\(u_0=d_2\varpi^2/\kappa_0\). Proposition~\ref{prop:theta-log}
bounds \(v_p(d_2)\ge-1/2\), since \(2\) is a unit and
\(v_p(\iota\sqrt{-d})=1/2\). Here
\(v_p(\varpi^2/\kappa_0)=1-(p-1)h/Q=1/2\), so \(u_0\in O\).
Take \(u_1=b_1,\ u_2=b_2\) from Lemma~\ref{lem:two-coefficients}.
Composition with \(S\in XO[[X]]\) preserves the congruence.
Lemma~\ref{lem:project-log}, evaluated at \(\varpi X\), gives an
integral unit \(V\), since \(T(\varpi X)=\varpi S(X)\).
The definition of \(\psi\) now gives the asserted factorization
of \(F_0=\exp(-\mathcal L)\). Use
Theorem~\ref{thm:frob-parameter} for \(\Phi(\varpi X)\)
and apply Theorem~\ref{thm:remove-shifts} with \(\theta=\pi\)
to obtain the conclusion for all the translates.
\end{proof}

\section{The first exceptional coefficient at an inert prime}

\begin{proposition}\label{prop:first-exceptional}
If \(p\) is inert, the coefficient \(u_1\) constructed in
Theorem~\ref{thm:actual-theta-normal} is a unit. More precisely
\[
              u_1=\frac{[Z^{p-1}]\Xi(Z)}{p+1}\in O^\times .
\]
\end{proposition}
\begin{proof}
Here \(q'=p^2\), so Theorem~\ref{thm:frob-parameter} gives
\(T(Z)=Z+O(Z^{p^2})\). Thus for any \(f\in O[[T]]\) all coefficients
through degree \(p\) in \(Z\) are integral. In particular
\([Z^{p-1}]f'=p[Z^p]f\in pO\). For such \(f\) let
\(\lambda(f)\) denote the residue of \([Z^{p-1}]f\) in the field
\(O/\{v_p>0\}\). The ring \(O[[T]]\) is closed under differentiation
by Proposition~\ref{prop:cm-parameter} and
Theorem~\ref{thm:frob-parameter}, and \(\lambda(f')=0\).

Propositions~\ref{prop:scaled-cm} and~\ref{prop:cm-parameter} give
\[
       (\Xi')^2=4\Xi^3+A_1^2\Xi^2+4\Xi,\qquad
       \Xi''=6\Xi^2+A_1^2\Xi+2 .
\]
The second identity follows from the first by differentiation and
cancellation of the nonzero series \(2\Xi'\). Thus for \(j\ge0\)
\[
 (\Xi^j\Xi')'=(4j+6)\Xi^{j+2}+(j+1)A_1^2\Xi^{j+1}
                                     +(4j+2)\Xi^j .
\]
The term with the derivative of \(\Xi^0\) is zero for \(j=0\).
Apply \(\lambda\) and set \(r=(p-1)/2\). The integer \(4j+6\)
is a unit for \(0\le j\le r-2\), since it is even and strictly
between zero and \(2p\), with \(p\) odd. Now \(\lambda(1)=0\).
If \(\lambda(\Xi)=0\), the recurrence gives
\(\lambda(\Xi^j)=0\) for every \(j\le r\). Yet
\(\Xi(Z)=Z^2+O(Z^3)\) by Proposition~\ref{prop:cm-parameter}
and \(\Phi'(0)=1\); hence \(\lambda(\Xi^r)=1\).
This contradiction proves \([Z^{p-1}]\Xi\in O^\times\).

In Lemma~\ref{lem:two-coefficients} the coefficient at \(p+1\)
is \(p[X^{p+1}]m\). The inverse \(Z_T\) agrees with \(X\)
through degree \(p+1\), so this is \(p[Z^{p+1}]\mathcal L_1\)
in Theorem~\ref{thm:actual-theta-normal}. Taking the coefficient
of degree \(p-1\) in \(\mathcal L_1''=\Xi\) gives the stated formula.
\end{proof}

\begin{corollary}\label{cor:actual-differential}
Use the tight frames of Proposition~\ref{prop:frame-frobenius}.
For the actual theta products of Proposition~\ref{prop:cm-theta}
with one common torsion denominator \(n_0\) prime to \(p\), the
coefficients of the differential minors in Theorem~\ref{thm:schur-jet}
have valuation at least \(J_p(j,V)-m_{T'}\) for every set \(V\)
and every charge system obeying its displayed condition. At an
inert prime with \(p\nmid M'\) and \(\beta\in O^\times\), the
derivative formula of Corollary~\ref{cor:sparse-derivative}
holds for the common exponent of the translates in the normalization
of Theorem~\ref{thm:remove-shifts}, with \(\beta/n_0\) in place of
\(\beta\). Each indicated monomial has a unit coefficient.
\end{corollary}
\begin{proof}
Theorems~\ref{thm:actual-theta-normal} and~\ref{thm:product-null}
establish the linear substitution hypothesis of
Theorem~\ref{thm:schur-jet} for these products, after the simultaneous
change of argument in Theorem~\ref{thm:remove-shifts}.
Thus that theorem gives the stated lower bound for each coefficient
of a differential minor. For the last assertion the argument scale
\(\beta/n_0\) is a unit. The scalar \(n_0^2\) is absorbed into
\(\kappa\) there. Proposition~\ref{prop:first-exceptional} supplies
the unit \(u_1\), so Corollary~\ref{cor:sparse-derivative} applies.
\end{proof}

\section{Integral derivatives at every finite place}

\begin{proposition}\label{prop:quadratic-shifts}
Fix any prime \(p\), any coefficient embedding
\(\iota:\C\to K_p\), and \(0\ne\delta\in O_K\).
Put \(A=\iota(\bar\delta)\), \(N=\delta\bar\delta\).
For every \(y\in\C\) with \(\delta y\in\Lambda\), one has
\[
               F_y(AZ)=F_0(Z)^N V_y(\Phi(Z)),
                      \qquad V_y\in O[[X]] .
\]
\end{proposition}
\begin{proof}
Use the multiplier \(\chi(l)\) of Proposition~\ref{prop:cm-multiplication}.
The quotient \((U_y\theta)(\bar\delta z)/\theta(z)^N\)
is periodic. In fact for a lattice period \(l\), the two quadratic
exponential multipliers match because \(|\delta|^2=N\).
The signs match by \(\chi(\bar\delta l)=\chi(l)^N\).
The remaining phase of the translate is
\(\exp(2\pi i\Im(\bar l\delta y)/{\rm area}(\Lambda))=1\).
Indeed for two lattice vectors
\(\lambda(m+n\omega),\lambda(m'+n'\omega)\) the quotient of the
imaginary part of their hermitian product by the area is the
integer \(mn'-nm'\). The factor \(U_y\) and the signs are those
of Proposition~\ref{prop:cm-theta}.
Since \(\delta y\in\Lambda\) implies \(Ny\in\Lambda\),
Propositions~\ref{prop:tate-torsion} and~\ref{prop:cm-multiplication}
give
\[
       |\iota((U_y\theta)(\bar\delta z))|_p
                    \le H_p(\bar\delta z)\le H_p(z)^N .
\]
The numerator is entire. Proposition~\ref{prop:bounded-sections-integral}
therefore gives the desired equality after taking Taylor series.
\end{proof}

\begin{theorem}\label{thm:all-prime-jets}
Use \(c_2=\mathcal L''(0)\) of Proposition~\ref{prop:theta-log}
and choose \(0\ne t\in O\) with \(t^2c_2\in O\).
Fix a finite set of shifts \(y_i\) with \(\delta y_i\in\Lambda\)
as in Proposition~\ref{prop:quadratic-shifts}. For finitely many
variables \(Z_1,\ldots,Z_s\), take linear forms
\(L_i=\sum_j c_{ij} Z_j\) with \(c_{ij}/(At)\in O\).
Then every mixed partial derivative at zero of
\[
                         \mathcal T(Z)=\prod_i F_{y_i}(L_i)
\]
belongs to \(O\), including the derivative of order zero.
The condition \(t^2c_2\in O\) is necessary for this conclusion
uniformly in the indicated products and forms, for each fixed
\(t\in O\setminus\{0\}\), when \(\delta=1\).
It always suffices to take \(v_p(t)=v_p(d)/4\).
In particular every coefficient system satisfying
\(c_{ij}/(A d)\in O\) obeys the assertion.
\end{theorem}
\begin{proof}
Put \(\mathcal A=O[[\Phi]][t\mathcal L',F_0]\) in \(K_p[[Z]]\).
All its elements have integral constant term, by
Propositions~\ref{prop:cm-parameter} and~\ref{prop:theta-log}.
The derivation \(t\,d/dZ\) preserves this algebra. On its base ring
this follows from Proposition~\ref{prop:cm-parameter} and \(t\in O\).
On the two generators the derivatives are
\[
           t^2(\Xi+c_2)\in O[[\Phi]],\qquad
                             -(t\mathcal L')F_0\in\mathcal A .
\]
The ring rules prove preservation on the whole algebra. Thus
\((t\,d/dZ)^k H(0)\in O\) for every \(H\in\mathcal A,\ k\ge0\),
where evaluation is after differentiation.

Write \(L_i/A=t\sum_j u_{ij} Z_j=t W_i\) with \(u_{ij}\in O\).
By Proposition~\ref{prop:quadratic-shifts} each \(F_{y_i}(L_i)\)
is \(H_i(t W_i)\) with \(H_i\in\mathcal A\). Formal differentiation
with respect to \(Z_j\) changes this to
\(u_{ij}((t\,d/dZ)H_i)(t W_i)\).
The product rule therefore keeps every mixed derivative a finite
\(O\)-linear combination of products of the same kind. Evaluating at
zero proves the assertion.

For necessity take one zero shift, one variable and \(L=tZ\).
The second derivative of \(F_0(tZ)\) at zero is \(-t^2c_2\)
by \(F_0=\exp(-\mathcal L)\) and \(\mathcal L'(0)=0\).
Finally Proposition~\ref{prop:theta-log} gives
\(v_p(c_2)\ge-v_p(d)/2\).
A fourth root of \(d\) exists by Lemma~\ref{lem:place-fields} and
has valuation \(v_p(d)/4\). With that choice, \(t^2c_2\in O\)
and \(d/t\in O\), proving the last two assertions.
\end{proof}

\section{A mean criterion for the prime remainder}

\begin{lemma}\label{lem:narrow-mean}
Fix \(d,M\) as in Definition~\ref{def:prime-error} and
\(0<\eta<1\). There is a constant \(A>0\), depending only on
\(d,M,\eta\), with the following property. For \(t\ge1\) put
\[
 k=\lceil At\rceil,\quad \mu=(k+1)/3,\quad
 J_t=[(1-\eta)\mu,(1+\eta)\mu],\quad
 w_t(v)=\frac{v^k\e^{-3v}}{\int_{J_t}u^k\e^{-3u}\dd u}
                       \quad(v\in J_t).
\]
Eventually \(J_t\subset[t,t^2]\). Uniformly for \(0\le X\le\e^t\),
\[
 \left|\int_{J_t}\bigl(R_{d,M}(\e^v)-X R_{d,1}(\e^v)\bigr)
                      w_t(v)\dd v\right|=O_{d,M,\eta}(\e^{-2t}).
\]
\end{lemma}
\begin{proof}
Put \(F(v)=R_{d,M}(\e^v)\), \(G(v)=R_{d,1}(\e^v)\).
By Proposition~\ref{prop:polar} and Lemma~\ref{lem:moments} there
are common \(R>3,K<\infty\) bounding their moments as in that lemma.
Let \(Y\) have the full probability density
\(3^{k+1}v^k\e^{-3v}/k!\) on \(v\ge0\). Its normalization and
moment generating integral
\(\mathbb E\e^{sY}=(1-s/3)^{-(k+1)}\) for \(s<3\) were proved
in Theorem~\ref{thm:window}. It follows that
\[
             |\mathbb E(F(Y)-XG(Y))|
                   \le 3K(1+\e^t)(3/R)^k .
\]

Set \(h=\min\{\eta-\log(1+\eta),-\eta-\log(1-\eta)\}>0\).
Positivity follows by integration of \(1/u\) on the intervals
from \(1\) to \(1+\eta\) and from \(1-\eta\) to \(1\).
On the upper tail, integrate the inequality
\({\bf1}_{Y\ge(1+\eta)\mu}\le\e^{s(Y-(1+\eta)\mu)}\)
with \(s=3\eta/(1+\eta)\). The moment formula gives
\(\mathbb P(Y\ge(1+\eta)\mu)\le((1+\eta)\e^{-\eta})^{k+1}\).
On the lower tail take \(s=-3\eta/(1-\eta)\) in the same
exponential bound with threshold \((1-\eta)\mu\) and reversed
inequality on \(Y\). This gives
\(\mathbb P(Y\le(1-\eta)\mu)\le((1-\eta)\e^\eta)^{k+1}\).
Thus \(\mathbb P(Y\notin J_t)\le2\e^{-h(k+1)}\).
Integrating \(v^j\) against the density gives
\(\mathbb E Y^j=(k+j)!/(3^j k!)\) by the same elementary integral,
so \(\mathbb E(1+Y)^4=O((k+4)^4)\).
For real square-integrable functions the squared integral of their
product is bounded by the product of their squared norms: expand
the nonnegative integral of \((a f-g)^2\) as a quadratic in \(a\).
Using this inequality with the indicator of the complement, and
using the quadratic bound of Proposition~\ref{prop:polar}, yields
\[
 \mathbb E\!\left[|F(Y)-XG(Y)|{\bf1}_{Y\notin J_t}\right]
       \ll_{d,M}(1+\e^t)(k+4)^2\e^{-h(k+1)/2}.
\]
Choose \(A\) so that \(A\log(R/3)>4,\ Ah/2>4\) and
\((1-\eta)A/3>1\). Both bounds for the integrals are
\(O_{d,M,\eta}(\e^{-2t})\), since every fixed power of \(t\)
is \(o(\e^t)\) by its exponential series.
Conditioning on \(J_t\) divides their difference by
\(\mathbb P(Y\in J_t)\ge1/2\) eventually and gives the stated
density. The endpoints of \(J_t\) now give both inclusions
in \([t,t^2]\) for all sufficiently large \(t\).
\end{proof}

\begin{theorem}\label{thm:mean-covolume}
Use the data of Theorem~\ref{thm:geometric-reduction}. Fix
\(0<\eta<1\), set \(t=10^6\log x\), and choose \(J_t,w_t\) by
Lemma~\ref{lem:narrow-mean}. Write \(\mathcal R_x(v)\) for
\(\mathcal R\) at \(x_1=\e^v\), keeping \(m,a,S,U\) fixed.
Then as \(x\to\infty\)
\[
 R_{d,M}(x)=
  \frac{1}{maS}\int_{J_t}m\mathcal R_x(v)w_t(v)\dd v
                       +O_{d,M,\eta}(x^{-2\cdot10^6}).
\]
In particular the averaged inequality
\[
     \int_{J_t}m\mathcal R_x(v)w_t(v)\dd v\le C f a\ell^4
\]
for fixed \(C=C(d,M,\epsilon,\eta)\) and all sufficiently large
\(x\) gives
\(R_{d,M}(x)\le C\varphi(M)x^{-1/2+5\epsilon}
 +O_{d,M,\eta}(x^{-2\cdot10^6})\).
If this inequality for the average holds for each \(M,d\) of
Definition~\ref{def:levels} and for arbitrarily small fixed
\(\epsilon>0\) in the range of Proposition~\ref{prop:thick},
then \(L(s,\chi)\ne0\) for every Dirichlet character and
\(\Rp s>1/2\) at every regular point. The principal character
modulo \(q\) has a simple pole at \(s=1\) of residue \(\varphi(q)/q\).
Under the same hypothesis Conjecture~\ref{conj:grh} follows.
\end{theorem}
\begin{proof}
For every \(v\in J_t\), Proposition~\ref{prop:rem-id} gives
\[
 R_{d,M}(x)=\frac{\mathcal R_x(v)}{aS}
              +R_{d,M}(\e^v)-(U/S)R_{d,1}(\e^v).
\]
All functions in this formula are continuous in \(v\) by
Proposition~\ref{prop:polar} and the identity itself.
The range of \(x_1\) is admissible by Lemma~\ref{lem:narrow-mean}
and Theorem~\ref{thm:geometric-reduction}. The ratio \(U/S\le
\ell=x^\epsilon\le\e^t\) by Proposition~\ref{prop:thick}.
Integrate with \(w_t\) and use Lemma~\ref{lem:narrow-mean};
this proves the first formula. Since \(m\ge\sqrt{x}\) and
\(f=\varphi(M)U\) by Propositions~\ref{prop:thick}
and~\ref{prop:places}, the averaged inequality gives the
displayed upper bound. Given \(1/2<\beta<1\), choose one of the
available \(\epsilon\) with \(5\epsilon<\beta-1/2\).
Both terms in the bound are then \(o(x^{\beta-1})\), with \(d,M\)
fixed. Corollary~\ref{cor:prime-suffices} and
Proposition~\ref{prop:cont} prove the remaining assertions.
\end{proof}

\section{The size of the main height and the individual volume ratios}

\begin{proposition}\label{prop:positive-main-height}
Use \(t=10^6\log x,\ t\le v\le t^2\), and the data of
Definition~\ref{def:heights}. Uniformly in this range,
\[
 H\ge 16mk\,\ell S\,a/x>0,\qquad
                   \frac{m|\mathcal R_x(v)|}{H}
                         \ll_{d,M}\frac{x(1+v)^2}{k}=o(1)
\]
for all sufficiently large \(x\), with \(\epsilon\) fixed. In
particular \(n(H-m\mathcal R_x(v)+6B)\le d_0(H-m\mathcal R_x(v)+6B)\)
whenever \(0\le n\le d_0\), where \(B=fa\ell^3\).
\end{proposition}
\begin{proof}
In \(P(y,\tau)\) of Definition~\ref{def:heights}, with \(0<y\le a\),
a positive summand requires \(n<yx_1/a\le x_1=\e^v\),
since the least possible weight is \(a/x_1\). By
Definition~\ref{def:prime-error}, \(\lambda_{d,1}(n)\le2\log n\)
for \(n\ge1\). Integral comparison of the harmonic sum gives
\(0\le P(y,\tau)\le2y v(1+v)\).
The formula for \(J\) in Proposition~\ref{prop:height-cost} and
the bound \(G_0\ge19\ell\) in Proposition~\ref{prop:level-prime}
give
\[
     J(w_0(\tau))\ge (G_0+1+\kappa_{d,1})w_0(\tau)
                  \ge18\ell w(\tau)
\]
eventually. Here \(w_0(\tau)\ge w(\tau)\), both are positive,
and \(\log(w_0(\tau)!)\ge0\). Since \(v\le t^2\), we have
\(v(1+v)=o(\ell)\) by Proposition~\ref{prop:thick}.
Consequently \(D_0\ge16\ell\sum_\tau w(\tau)\ge16\ell S a/x\),
using the \(S\) places of weight \(a/x\) alone.
Multiplying by \(mk\) gives the first bound.

By Definition~\ref{def:ask},
\(\log a-\bar w\) is a weighted mean of \(\log x\) and \(v\)
and lies between zero and \(v\). Thus
\[
   |\mathcal R_x(v)|\le Ua\{2v(1+v)+v+|\kappa_{d,1}|\}.
\]
Divide by the lower bound for \(H/m\) and use \(U/S\le\ell\)
from Proposition~\ref{prop:thick}. Finally
\(k=r+1-m\ge x^5/2\) eventually by that proposition.
It follows that \(x(1+v)^2/k\to0\) uniformly.
Thus \(H-m\mathcal R_x(v)>0\) eventually, giving the last inequality.
\end{proof}

\begin{proposition}\label{prop:comparison-volume}
Let \(U,T\) be finite-dimensional complex Hilbert spaces with
\(D=\dim U>0\), and \(A:T\to U\) a surjective linear map.
For a parameter \(t'\) in a set with a distinguished member \(0\),
let \(u_{t'}:U\to V_{t'}\), \(G_{t'}:T\to V_{t'}\) be linear maps
of finite-dimensional complex Hilbert spaces. Suppose
\(u_{t'}^*G_{t'}=\gamma_{t'} A\) with \(\gamma_{t'}\ne0\).
Write \(L_{t'}=u_{t'}(U)\) and \(P_{t'}\) for the orthogonal
projection onto \(L_{t'}\). For a fixed positive real \(h\), let
\[
   \lambda(t')=Z(u_{t'}),\qquad
   \Delta(t')=\det\!\left(hP_{t'}G_{t'}G_{t'}^*|_{L_{t'}}\right).
\]
Both numbers are positive, and
\[
       \frac{\lambda(0)}{\lambda(t')}=
       \frac{\Delta(t')}{\Delta(0)}
                      \left|\frac{\gamma_0}{\gamma_{t'}}\right|^{2D}.
\]
Suppose that parameters \(t_\iota\) have been selected for all
\(\iota=(p,j)\) of Definition~\ref{def:place-height}. If real
\(a_\infty,a_p,C_0,C_1\) satisfy
\[
 \begin{gathered}
 |\sigma\Delta(0)|\le\e^{Da_\infty}\quad(\sigma\in{\rm Aut}(\C)),\qquad
 |j\Delta(t_\iota)|_p\le\e^{Da_p},\\
 a_\infty+\sum_{p<M'}a_p\le C_0
                 \quad\hbox{eventually as }M'\to\infty,\\
 {\rm ht}\!\left(\left|\frac{\gamma_0}{\gamma_{t_\iota}}\right|^2\right)
                                      \le C_1 ,
 \end{gathered}
\]
then \({\rm ht}(\lambda(0)/\lambda(t_\iota))\le D(C_0+C_1)\).
\end{proposition}
\begin{proof}
Surjectivity of \(A\) implies injectivity of \(A^*\): if \(A^*z=0\)
then \(z\) is orthogonal to \(A(T)=U\). Taking adjoints gives
\(G_{t'}^*u_{t'}=\bar\gamma_{t'} A^*\), an injection.
Hence \(u_{t'}\) is injective, and \(G_{t'}^*\) is injective
on \(L_{t'}\). For \(z\in L_{t'}\) the quadratic form of the
compressed operator at \(z\) is \(h\|G_{t'}^*z\|^2\).
It is positive definite, so its determinant is positive by
the Gram determinant formula of Lemma~\ref{lem:hadamard}.
Lemma~\ref{lem:compressed-gram} gives
\[
       \lambda(t')\Delta(t')=
                    h^D|\gamma_{t'}|^{2D}\det(AA^*) .
\]
Indeed \(\det(AA^*)=Z(A^*)>0\) by Definition~\ref{def:lift-vol}.
Divide the identity for \(0\) by the identity for \(t'\).
For the height of \(\Delta(t_\iota)/\Delta(0)\), take \(\Delta(0)\)
as the nonzero witness of Definition~\ref{def:place-height}.
The stated archimedean and finite bounds give cost at most \(DC_0\).
Multiply by the \(D\)-th power of the scalar ratio using
Proposition~\ref{prop:height-product}.
\end{proof}

\section{The comparison before enlarging its error terms}

\begin{theorem}\label{thm:sharper-actual-budget}
Fix \(x,v\) in the range of Proposition~\ref{prop:positive-main-height}
and use the parameters of Definition~\ref{def:heights}. Put
\(B=fa\ell^3,\ b_0=B/(mk)\), and \(L_{\rm loc}=\sum_p e_p\log p\)
as in Proposition~\ref{prop:exact-normalization}.
Let \(W\) have dimension \(d_0>0\) and let \(U_{t'}:W\to H_{t'}\)
be the maps of Corollary~\ref{cor:completion-height}, all injective
for \(t'=0,t_\iota\). Suppose the completions and recombinations
in that corollary exist, with parameters satisfying
\[
 \begin{gathered}
  n\ge(1-4/b^2)d_0,\qquad \beta+g\le b^3/4,\qquad
       \delta\le3d_0\log(1+b^2)L_{\rm loc},\\
 C_\sigma\bigl((U_0x_i)_{i\le n}\bigr)
        \le \exp\{n(H-m\mathcal R_x(v)+6B)\}
                       \qquad(\sigma\in{\rm Aut}(\C)).
 \end{gathered}
\]
For \(1\le j\le m\) suppose the comparison maps of
Proposition~\ref{prop:comparison-volume} exist with \(D=d_0\);
denote their functions by \(\lambda_j,\Delta_j,\gamma_j\).
Assume the determinant bounds of that proposition with
\(a_\infty\le2r b_0\) and \(\sum_{p<M'}a_p\le r b_0\)
eventually for each \(j\). Require
\[
 {\rm ht}\!\left(\left|\frac{\gamma_j(0)}
                        {\gamma_j(t_\iota)}\right|^2\right)
        \le 10m x^{999}+k(6b_0-D_0).
\]
Finally suppose \(Z(U_{t'})=|c(t')|^{d_0}\prod_{j=1}^m\lambda_j(t')\)
with \(c(t')\ne0\) and
\({\rm ht}(|c(t_\iota)/c(0)|^{-2})\le2m^2t_1x^{999}\).
Then, for all sufficiently large \(x\),
\[
 \begin{aligned}
 m\mathcal R_x(v)
 &\le (12+3r/k)B+(10+t_1)m^2x^{999}+b
                            +6\log(1+b^2)L_{\rm loc}  \\
 &\ll_{d,M,\epsilon} fa\ell^3(1+v)^2 .
 \end{aligned}
\]
If these hypotheses hold on the intervals \(J_t\) of
Theorem~\ref{thm:mean-covolume} for fixed \(d,M,\epsilon,\eta\)
and every sufficiently large \(x\), then
\[
           R_{d,M}(x)\le C_{d,M,\epsilon,\eta}
                         x^{-1/2+4\epsilon}(1+\log x)^2
\]
for some positive constant \(C_{d,M,\epsilon,\eta}\).
If they hold for all \(M,d\) of Definition~\ref{def:levels} and
arbitrarily small fixed \(\epsilon>0\) as in
Proposition~\ref{prop:thick}, the nonvanishing and pole assertions of
Theorem~\ref{thm:mean-covolume} and Conjecture~\ref{conj:grh} follow.
\end{theorem}
\begin{proof}
For each stage Proposition~\ref{prop:comparison-volume} gives
\[
 {\rm ht}\!\left(\lambda_j(0)/\lambda_j(t_\iota)\right)
   \le d_0\{3r b_0+10m x^{999}+k(6b_0-D_0)\}.
\]
By Proposition~\ref{prop:height-product} their product has height
at most
\[
      d_0\{(6+3r/k)B+10m^2x^{999}-H\},
\]
using \(mkb_0=B,\ mkD_0=H\). Set \(Z_\iota=Z(U_{t_\iota})/Z(U_0)\).
Squaring the product of ratios and using the scalar height and the
factorization gives
\[
 {\rm ht}(Z_\iota^{-2})
        \le 2d_0\{(6+3r/k)B+(10+t_1)m^2x^{999}-H\}.
\]
For the opposite direction apply Corollary~\ref{cor:completion-height}.
The cost for the missing \(d_0-n\) vectors is at most \(d_0b\).
By Proposition~\ref{prop:positive-main-height}, the exponent for
the given \(n\) vectors is at most \(d_0(H-m\mathcal R_x(v)+6B)\).
Thus
\[
   {\rm ht}(Z_\iota)\le d_0(H-m\mathcal R_x(v)+6B+b)+2\delta .
\]
Multiply the square of this family by its inverse square.
Proposition~\ref{prop:height-product} forces the sum of the corresponding
height bounds to be nonnegative. After division by \(2d_0\),
the terms \(H\) cancel and the claimed explicit inequality follows.

By Proposition~\ref{prop:loss} and its proof,
\[
          L_{\rm loc}\le B\{15v(1+v)/\ell+480/\ell^2\}.
\]
The factor \(\log(1+b^2)\le3\ell\) therefore contributes
\(O(B(1+v)^2)\). Also \(r/k\to1,\ t_1=5r\) and
\(m\ll\sqrt{x},\ r\ll x^5\) by Proposition~\ref{prop:thick}.
Since \(B\ge xb\ell^3\), every fixed power of \(x\) is \(o(B)\)
by the exponential series for \(b=\exp(x^\epsilon)\); also \(b\le B\).
This proves the second bound, uniformly in \(v\).
On \(J_t\) we have \(v=O_{d,M,\eta}(\log x)\) by
Lemma~\ref{lem:narrow-mean}. Integrate the second bound there and
use the identity in Theorem~\ref{thm:mean-covolume}.
Division by \(maS\) costs at most
\(f/(mS)\le\varphi(M)\ell/\sqrt{x}\) in the coefficient of
\(\ell^3(1+\log x)^2\). The error in that identity is
\(O_{d,M,\eta}(x^{-2\cdot10^6})\), proving the one-sided bound.
For any \(1/2<\beta'<1\), choose an available \(\epsilon\)
with \(4\epsilon<\beta'-1/2\); the squared logarithm is
\(o(x^{\beta'-1/2-4\epsilon})\) by the exponential series.
Corollary~\ref{cor:prime-suffices} now proves the last assertion.
\end{proof}

\section{Finite and archimedean bounds for the compressed determinant}

\begin{lemma}\label{lem:max-dual-minor}
Let \(L\) be a \(D\)-dimensional complex vector space, \(D\ge1\),
and \(\ell_a\) a finite family of linear forms on \(L\) whose common
kernel is zero. Fix \(\iota:\C\to K_p\) as in
Definition~\ref{def:place-height}. There exist indices
\(a_1,\ldots,a_D\) and a basis \(z_1,\ldots,z_D\) of \(L\) with
\[
    \ell_{a_i}(z_j)={\bf1}_{i=j},\qquad
                     |\iota(\ell_a(z_j))|_p\le1\quad(a,\ 1\le j\le D).
\]
\end{lemma}
\begin{proof}
In any basis of \(L\), put the coordinate rows of the \(\ell_a\)
in a matrix \(B\). Its columns are independent by the common
kernel condition. Elementary elimination, using a nonzero entry
as a pivot and clearing its column, produces \(D\) independent
rows: each step reduces the column rank of the remaining block
by one. Thus some \(D\)-rowed square minor is nonzero.
Among the finitely many ordered row choices choose a minor with
maximum \(p\)-adic absolute determinant through \(\iota\).
Its forms are a basis of the dual, so there is a basis \(z_j\)
with the first asserted equality, obtained by inverting this
square matrix. In the expansion \(\ell_a=\sum_j\ell_a(z_j)\ell_{a_j}\),
replacing row \(j\) of the chosen determinant by \(\ell_a\)
multiplies its determinant by \(\ell_a(z_j)\), by multilinearity
and alternation. Maximality yields the second assertion.
\end{proof}

\begin{proposition}\label{prop:finite-compressed}
Let \(G:T\to V\) be a linear map of finite-dimensional complex
Hilbert spaces, let \(L\subset V\) have dimension \(D>0\), and let
\(P\) be orthogonal projection onto \(L\). Assume \(G^*|_L\) is
injective. Fix \(h>0,\ \iota:\C\to K_p\) and finitely many
\(f_a\) spanning \(T\). Apply Lemma~\ref{lem:max-dual-minor}
to \(\ell_a(z)=\langle Gf_a,z\rangle\) on \(L\), obtaining indices
\(a_i\) and basis \(z_j\). Then
\[
 \det(hPGG^*|_L)
     =\det\!\left(\left\langle Gf_{a_i},
                        P(hGG^*z_j)\right\rangle\right)_{i,j=1}^D .
\]
In particular suppose \(\mathcal I,\mathcal J\subset V\) satisfy
\(Gf_a\in\mathcal I\) for every \(a\), and
\(hGG^*z\in\mathcal J\) whenever \(z\in L\) and
\(|\iota\ell_a(z)|_p\le1\) for every \(a\). If
\[
    \left|\iota\det(\langle y_i,P y'_j\rangle)_{i,j=1}^D\right|_p
           \le\e^{Db_p}\qquad(y_i\in\mathcal I,\ y'_j\in\mathcal J),
\]
then \(|\iota\det(hPGG^*|_L)|_p\le\e^{Db_p}\). It suffices that
the displayed bound hold for the indicated single pair of families
attached to a maximizing minor.
\end{proposition}
\begin{proof}
If all \(\ell_a(z)\) vanish then
\(\langle f_a,G^*z\rangle=0\) for every \(a\), so \(G^*z=0\) by
spanning, and \(z=0\). Thus the lemma applies. For an endomorphism
\(M\) of \(L\), the \(i,j\) entry of its matrix in the basis \(z_j\)
is \(\ell_{a_i}(Mz_j)\), by duality. Apply this to \(M=hPGG^*|_L\).
By the bounds in Lemma~\ref{lem:max-dual-minor}, every
\(hGG^*z_j\) belongs to \(\mathcal J\). The identity of
determinants proves both versions of the inequality.
\end{proof}

\begin{proposition}\label{prop:arch-compressed}
In Proposition~\ref{prop:finite-compressed} realize \(T=\C^I\)
with norm given by a positive definite rational Gram matrix, and
\(V\) inside a complex coordinate space whose positive definite
rational Gram matrix induces its norm. Use the same matrix for
the norm of the coordinatewise image \(\sigma v\) under any
\(\sigma\in{\rm Aut}(\C)\), for \(T\) and the ambient space of \(V\)
respectively. Suppose \(h\in\Q_{>0}\), and for all \(\sigma\),
\(x\in T\), \(z\in L\) one has
\[
       h\|\sigma(Gx)\|^2\le\e^{b_\infty}\|\sigma x\|^2,\qquad
       \|\sigma(PGG^*z)\|\le\|\sigma(GG^*z)\| .
\]
Then \(|\sigma\det(hPGG^*|_L)|\le\e^{Db_\infty}\) for every
\(\sigma\).
\end{proposition}
\begin{proof}
In these coordinates rationality of the two Gram matrices gives
\[
        \sigma\langle y,z\rangle
             =\langle\sigma' y,\sigma z\rangle,\qquad
             \sigma'=\mathfrak c\sigma\mathfrak c ,
\]
where \(\mathfrak c\) denotes complex conjugation, as in
Lemma~\ref{lem:mixed-gram}. Each eigenvalue \(\alpha\) of
\(hPGG^*|_L\) has a nonzero eigenvector \(z\).
Indeed its characteristic polynomial splits by
Lemma~\ref{lem:place-fields}, and vanishing of the determinant
of \(M-\alpha I\) gives a nontrivial kernel by elimination.
Set \(a=G^*z\), which is nonzero by injectivity.
For a fixed \(\sigma\) choose \(x=(\sigma')^{-1}(\sigma a)\)
coordinatewise in \(T\). Then
\[
      \sigma\langle Gx,z\rangle
      =\sigma\langle x,a\rangle=\|\sigma a\|^2 .
\]
Multiply the left side by \(\sigma\alpha\) and use the eigenvector
equation in the second argument of the pairing. Cauchy--Schwarz
gives
\[
  |\sigma\alpha|\,\|\sigma a\|^2
     \le h\,\|\sigma'(Gx)\|\,\|\sigma(PGa)\|
     \le h\,\|\sigma'(Gx)\|\,\|\sigma(Ga)\| .
\]
Here the last inequality is the projection hypothesis at this
particular \(z\). The pairing inequality follows by subtracting the
orthogonal projection onto the line of one vector as in
Lemma~\ref{lem:hadamard}. Apply the norm hypothesis with
\((\sigma',x)\) and \((\sigma,a)\). Since \(\sigma'x=\sigma a\), their
product bounds the right side by \(\e^{b_\infty}\|\sigma a\|^2\).
Cancellation gives \(|\sigma\alpha|\le\e^{b_\infty}\).
The determinant is the product of the \(D\) eigenvalues counted
with multiplicity by the constant term of the characteristic
polynomial. Multiplication yields the assertion.
\end{proof}

\begin{corollary}\label{cor:compressed-tests}
In Theorem~\ref{thm:sharper-actual-budget} the determinant estimates
are implied by the following tests for each stage. In the notation
of Proposition~\ref{prop:comparison-volume}, use the tests of
Proposition~\ref{prop:arch-compressed} at \(t'=0\) with
\(b_\infty=2r b_0\). For every \(\iota=(p,j')\) use the single
minor of Proposition~\ref{prop:finite-compressed} at \(t'=t_\iota\)
and embedding \(j'\), with \(b_p=r\beta_p\), where
\(\sum_{p<M'}\beta_p\le b_0\) eventually. All other hypotheses
of Theorem~\ref{thm:sharper-actual-budget} have their stated force.
\end{corollary}
\begin{proof}
Injectivity of \(G_{t'}^*\) on \(L_{t'}\) was proved in
Proposition~\ref{prop:comparison-volume}. The two determinant
bounds follow directly from Propositions~\ref{prop:arch-compressed}
and~\ref{prop:finite-compressed}. Their exponents have precisely
the sums required in Theorem~\ref{thm:sharper-actual-budget}.
\end{proof}

\section{The arithmetic lattice of the theta frames}

\begin{proposition}\label{prop:local-cyclo-different}
Put \(N_c=N_1\), \(\zeta=\zeta_{N_c}\) and
\(D_c=P_{N_c}'(\zeta)\) as in Proposition~\ref{prop:different}.
For any prime \(l\) and any embedding \(\varphi:E\to K_l\),
\[
 v_l(\varphi D_c)=
 \begin{cases}
 a_l-1/(l-1),&a_l=v_l(N_c)>0,\\
 0,&l\nmid N_c .
 \end{cases}
\]
If \(l\mid N_c\), the root \(\xi_l=\zeta^{N_c/l}\) satisfies
\(v_l(\varphi(1-\xi_l))=1/(l-1)\). For \(i,k\) distinct modulo
\(l\) the same is true of \(\xi_l^i-\xi_l^k\).
\end{proposition}
\begin{proof}
For a positive integer \(h\) prime to \(l\), reduction to the
residue field is injective on the \(h\)-th roots of unity. Indeed a
root \(u\) reducing to one has \(1+u+\cdots+u^{h-1}\) reducing to
\(h\ne0\), forcing \(u=1\). When \(l\nmid N_c\) this makes every
difference between \(\varphi\zeta\) and a distinct conjugate a unit.
The product formula for the derivative at a simple root gives
\(v_l(\varphi D_c)=0\).

For \(a=a_l>0\) let \(h=N_c/l^a\). In characteristic \(l\)
the only \(l\)-th root of unity is one since \(X^l-1=(X-1)^l\).
If \(\xi\) is primitive of order \(l\), each
\((1-\xi^i)/(1-\xi)\), \(1\le i<l\), therefore reduces to \(i\ne0\).
Evaluating \((X^l-1)/(X-1)=\prod_{i=1}^{l-1}(X-\xi^i)\) at
\(X=1\) shows that \((l-1)v_l(1-\xi)=1\). This proves the
claims about \(\xi_l\) and its powers.

The identity
\[
        P_{N_c}(X)P_h(X^{l^{a-1}})=P_h(X^{l^a})
\]
holds because a root of the polynomial on the right has order
\(h l^j\), \(0\le j\le a\), and the second factor on the left
accounts for exactly \(j<a\); the roots are simple in characteristic
zero. At \(X=\varphi\zeta\) the derivative of the right side has
valuation \(a\). In fact \(P_h'\) evaluated at a primitive
\(h\)-th root is a unit by the first argument of this proof.
The factor \(P_h((\varphi\zeta)^{l^{a-1}})\) has valuation
\(1/(l-1)\). To see this, the argument of \(P_h\) is of order
\(hl\), so by coprime exponents it is uniquely \(u\xi\) with
\(u\) a primitive \(h\)-th root and \(\xi\) a primitive \(l\)-th
root. In the product over primitive \(h\)-th roots only the
factor \(u\xi-u\) reduces to zero. It has the asserted valuation.
Differentiating the polynomial identity at \(\varphi\zeta\) and
subtracting this valuation gives the formula.
\end{proof}

\begin{definition}\label{def:frame-lattice}
Write \(n_K=(d+1)/4=\omega\bar\omega\).
For \(l\mid N_c,\ \varphi:E\to K_l\), define
\(d(l,\varphi)=v_l(\varphi D_c)\) if \(\varphi\omega\) has zero
residue, and \(d(l,\varphi)=0\) otherwise. Set it to zero for
\(l\nmid N_c\). Define the additive subgroup
\[
  \mathcal A=\{x\in E: v_l(\varphi x)\ge-d(l,\varphi)
                      \text{ for every }l,\varphi\}.
\]
For a finite-place embedding \(\iota:\C\to K_p\), put
\(\varphi_\tau(x)=\iota(\overline{\tau(x)})\) for \(\tau\in G\).
All these maps agree on \(K\). Thus
\(d_\iota=d(p,\varphi_\tau)\) is independent of \(\tau\) by
Proposition~\ref{prop:local-cyclo-different}. Choose any
\(s_\iota\in K_p^\times\) of valuation \(d_\iota\). Such a
choice is either \(1\) or \(\varphi_1 D_c\).
\end{definition}

\begin{proposition}\label{prop:frame-lattice-trace}
One has \(O_E\subset\mathcal A\subset N_c^{-1}O_E\). For
\(x,y\in\mathcal A\),
\[
               \sum_{\tau\in G}\tau(x)\overline{\tau(y)}
                             ={\rm Tr}_{E/K}(x\bar y)\in O_K .
\]
For each \(\iota,\tau\) as in Definition~\ref{def:frame-lattice}
and \(x\in\mathcal A\), \(s_\iota\varphi_\tau(x)\) is integral.
\end{proposition}
\begin{proof}
An element of \(E\) is in \(O_E\) precisely when it is integral
at every \((l,\varphi)\) of Definition~\ref{def:frame-lattice}.
A monic integer equation forces all valuations of a root to be
nonnegative: a negative valuation would make the leading term
strictly larger in absolute value than all other terms. Conversely
the elementary symmetric functions of the conjugates of an element
are the rational coefficients, up to signs, of its monic minimal
polynomial. All the conjugates in \(K_l\) extend to embeddings
of \(E\) there by adjoining roots, as in Lemma~\ref{lem:place-fields}.
Hence all coefficients are integral at each \(l\). A rational
number with that property has integer denominator one in lowest
terms by Lemma~\ref{lem:basic-arith}. The characterization follows.
Since \(0\le d(l,\varphi)\le v_l(N_c)\), it implies both inclusions.

Every prime \(l\mid N_c\) divides \(n_K\). Indeed
\(d+1=(M/2)(2u+1)\) for an integer \(u\), every odd prime of
\(N_c\) divides \(M\), and \(16\mid M\). Thus the two reductions of
\(\omega,\bar\omega=1-\omega\), in some order, are \(0,1\) at
every such embedding by their sum and product. Exactly one of
\(\varphi\) and \(\varphi\circ\mathfrak c\) selects the positive
allowance in Definition~\ref{def:frame-lattice} whenever that
allowance is positive. Proposition~\ref{prop:local-cyclo-different}
makes the valuation of \(D_c\) the same at both.
It follows that \(D_c x\bar y\in O_E\) by the characterization
just proved. Write \(D_c x\bar y=F(\zeta)\) with \(F\in O_K[X]\)
and \(\deg F<f\) by Theorem~\ref{thm:integer-ring}. The interpolation
identity in the proof of Lemma~\ref{lem:dual-monic}, applied to
the \(f\) roots of \(P_{N_c}\) over \(K\), gives
\[
       {\rm Tr}_{E/K}(x\bar y)=[X^{f-1}]F\in O_K .
\]
Those roots give exactly the embeddings in \(G\) by
Proposition~\ref{prop:places}. Each commutes with conjugation on
\(E\), as is seen on \(\zeta,\omega\); hence the sum is that
trace. The last assertion is the defining inequality for
\(\mathcal A\) at \(\varphi_\tau\).
\end{proof}

\section{Tight frames of every sufficiently large integral length}

\begin{lemma}\label{lem:quadratic-idempotent}
For \(N_c,n_K,\omega\) in Definition~\ref{def:frame-lattice} there
is \(t_K\in O_K\) with
\[
 \begin{gathered}
 0\le {\rm Tr}_{K/\Q}(t_K)<3N_c,\qquad
 0<{\rm N}_{K/\Q}(t_K)\le4n_K N_c^2,\\
 {\rm Tr}_{K/\Q}(t_K)\equiv1\bmod N_c,\qquad
 {\rm N}_{K/\Q}(t_K)\equiv0\bmod N_c .
 \end{gathered}
\]
For each \(l^a\Vert N_c\) and each \(\varphi:E\to K_l\),
\(\varphi t_K\equiv1\bmod l^aO\) on the branch where
\(\varphi\omega\) has zero residue, and
\(\varphi t_K\equiv0\bmod l^aO\) on the other branch.
\end{lemma}
\begin{proof}
By Proposition~\ref{prop:frame-lattice-trace}, \(l\mid n_K\).
The two roots of \(F(X)=X^2-X+n_K\bmod l\) are \(0,1\).
Their derivatives \(2X-1\) are units, including at \(l=2\).
Given a root \(u\bmod l^j\), the lifts \(u+l^j b\) have
\(F(u+l^jb)\equiv F(u)+l^j bF'(u)\bmod l^{j+1}\).
Exactly one residue \(b\bmod l\) prolongs the root. Thus there
are two roots \(u_0,u_1\bmod l^a\) reducing to \(0,1\), interchanged
by \(u\mapsto1-u\).
Evaluation at this pair identifies \(O_K/l^a O_K\) with
\((\Z/l^a\Z)^2\). Indeed \(O_K=\Z[\omega]\) with basis
\(1,\omega\), by the quadratic ring calculation in
Theorem~\ref{thm:integer-ring}. The evaluation matrix on these
two elements has determinant \(u_1-u_0\), a unit.

Choose the element evaluating to \((1,0)\) for each \(l\).
The Chinese remainder calculation of Lemma~\ref{lem:basic-arith}
combines the two coefficients to give \(t_K=A+B\omega\) with
\(0\le A,B<N_c\). At an embedding \(\varphi\), one of
\(\varphi\omega-u_0,\varphi\omega-u_1\) is a unit. Their product
lies in \(l^a O\), since \(F(X)\equiv(X-u_0)(X-u_1)\bmod l^a\).
This proves the claimed congruences. Conjugation interchanges the
two components, so the trace and norm reduce to \(1,0\bmod N_c\).
Finally \({\rm Tr}(t_K)=2A+B\) and
\({\rm N}(t_K)=|t_K|^2\le N_c^2(1+\sqrt{n_K})^2\le4n_K N_c^2\).
The first component shows \(t_K\ne0\).
\end{proof}

\begin{theorem}\label{thm:every-frame-length}
Put \(N_c=N_1,\ \log_2 u=\log(u)/\log(2)\). For any integers
\(R\ge0\) and
\[
                 L\ge 24d N_c^4(1+\log_2 N_c)(R+1)
\]
there is a finite family \(w_i\in\mathcal A\) with integer
multiplicities \(b_i\ge\max(1,R)\) such that
\[
       \sum_i b_i\,\tau(w_i)\overline{\sigma(w_i)}
                      =L\,{\bf1}_{\tau=\sigma}\qquad(\tau,\sigma\in G).
\]
For every finite-place embedding \(\iota\) there is an index \(i\)
with \(s_\iota\varphi_\tau(w_i)\) a unit for every \(\tau\in G\),
with the notation of Definition~\ref{def:frame-lattice}.
The family and the multiplicities can be taken independent of
\(\iota\). The length condition holds in particular for
\(L\ge(d N_c)^6(R+1)\).
\end{theorem}
\begin{proof}
Let \(t_K\) be as in Lemma~\ref{lem:quadratic-idempotent}.
Start with the \(N_c\) vectors \(\zeta^j\), \(0\le j<N_c\).
They lie in \(O_E\subset\mathcal A\) and have length \(N_c\)
by Definition~\ref{def:tight-field}, with each multiplicity one.
For each \(l^a\Vert N_c\) put \(r_l=N_c/l,\ h_l=N_c/l^a\),
\(\xi=\zeta^{r_l}\) and \(s_l=t_K/l^a\). Form a second family by
taking the following vectors for all such \(l\):
\[
 v_{l,j,k,i}=\zeta^j\left({\bf1}_{k=0}\xi^i+
                           s_l(\xi^i-\xi^k)\right),
          \qquad 0\le j<r_l,\quad 0\le k,i<l .
\]
These are integral at places of residue characteristic other than
\(l\). At the \(l\)-adic embeddings where \(\omega\) has nonzero
residue, \(s_l\) is integral. On the other branch, a nonzero
\(s_l(\xi^i-\xi^k)\) has valuation \(-a+1/(l-1)\) by
Proposition~\ref{prop:local-cyclo-different}. This is the precise
allowance of Definition~\ref{def:frame-lattice}. Hence all the
vectors belong to \(\mathcal A\).

Fix \(\tau,\sigma\in G\), regarded as exponents on \(\zeta\).
Set \(u=\tau(\xi)\), \(w=\sigma(\xi)\). The sums of the \(l\)
powers of \(u\) and \(w\) both vanish, and
\(\sum_{i=0}^{l-1}(u\bar w)^i=l{\bf1}_{u=w}\).
Upon multiplying the two factors in parentheses in the vector
formula, conjugating the second, and summing over \(i,k\), the
constant term is \(l{\bf1}_{u=w}\). The terms linear in \(s_l,\bar
s_l\) sum to \(l(s_l+\bar s_l){\bf1}_{u=w}\), since \(k=0\) in
their other factors. The mixed quadratic sum is
\(2l^2|s_l|^2{\bf1}_{u=w}\); the two cross sums in the product
\((u^i-u^k)(\bar w^{\,i}-\bar w^{\,k})\) vanish separately.
When \(u=w\), \(l\mid\tau-\sigma\). The sum of the remaining
\(\zeta^{j(\tau-\sigma)}\) over \(j<r_l\) is then \(r_l\) for
\(\tau=\sigma\) and zero otherwise. Thus this block has length
\[
 C_l=N_c(1+s_l+\bar s_l+2l|s_l|^2)
      =N_c+h_l\left({\rm Tr}(t_K)+
                        2l\frac{{\rm N}(t_K)}{l^a}\right).
\]
By Lemma~\ref{lem:quadratic-idempotent} this is a positive integer,
congruent to \(h_l\ne0\bmod l\). Every other prime dividing \(N_c\)
divides \(h_l\) and hence \(C_l\). The sum \(C=\sum_{l\mid N_c} C_l\)
over prime divisors is therefore coprime to \(N_c\).
The bounds of the lemma give \(C_l\le N_c+3N_c^2+8n_KN_c^3
\le12d N_c^3\). There are at most \(\log_2 N_c\) prime divisors,
since the product of \(a'\) distinct primes is at least \(2^{a'}\).
Consequently
\(C+N_c\le12d N_c^3(1+\log_2 N_c)\).

Let \(R_0=\max(1,R)\). Our length hypothesis implies
\(L\ge(R_0+N_c)(C+N_c)\), since
\(R_0+N_c\le2N_c(R+1)\). Reduce
\(L-R_0(C+N_c)\) modulo \(N_c\). Since \(C\) is invertible there,
it has the same residue as \(BC\) for some \(0\le B<N_c\).
The difference is nonnegative by the length bound and equals
\(A N_c\) for \(A\in\Z_{\ge0}\). Give the first family multiplicity
\(A+R_0\) everywhere and the second multiplicity \(B+R_0\).
The preceding identities give exactly the asserted length.

If \(d_\iota=0\), the vector \(1\) in the first family supplies
the unit simultaneously at all \(\tau\). Otherwise \(p=l\mid N_c\)
and every \(\varphi_\tau\) sees the branch where \(t_K\) is a unit.
Take the vector with \(j=0,k=1,i=0\) in that block. It is
\(t_K(1-\xi)/l^a\), of valuation \(-d_\iota\) for every such
embedding. Its product with \(s_\iota\) has valuation zero.
Finally \(1+\log_2 N_c\le N_c\) for integral \(N_c\ge2\), since
\(2^{N_c-1}\ge N_c\) by induction, and \(24\le d^5 N_c\)
by Proposition~\ref{prop:places} and the size of \(N_c\).
These inequalities prove the last length comparison.
\end{proof}

\section{One algebraic normalization for all the theta frames}

\begin{theorem}\label{thm:global-frame-scale}
In Definition~\ref{def:frame-lattice} there is a single
\(s_0\in\overline{\Q}^{\,\times}\subset\C^\times\) with
\[
       v_p(\iota s_0)=d_\iota\quad\hbox{for all }(p,\iota),\qquad
       |\sigma s_0|^2=\operatorname{rd}_E/\sqrt d
                                \quad(\sigma\in{\rm Aut}(\C)).
\]
Consequently in Theorem~\ref{thm:every-frame-length} one may use
\(s_\iota=\iota s_0\) at all embeddings simultaneously.
\end{theorem}
\begin{proof}
Take a common positive multiple \(D\) of the integers \(l-1\) for
\(l\mid N_c\) prime, and put
\(e_l=D(v_l(N_c)-1/(l-1))\in\Z_{\ge0}\).
Let \(P_l=(l,\bar\omega)\subset O_K\) and
\[
          I=\prod_{l\mid N_c}P_l^{e_l},\qquad
                       m_0=\prod_{l\mid N_c}l^{e_l}
\]
with both products over primes. For each \(l\),
\(P_l\bar P_l=lO_K\): its generators \(l^2,l\omega,l\bar\omega,n_K\)
belong to \(lO_K\), and \(l\omega+l\bar\omega=l\).
Also \(P_l+\bar P_l=O_K\) since \(\omega+\bar\omega=1\).
Their positive powers remain comaximal. Indeed, if \(a+b=1\)
belongs termwise to two such ideals then every term of
\((a+b)^{2n-1}\) belongs to one of their \(n\)-th powers.
At distinct primes the corresponding powers contain relatively
prime rational integers. Products of ideals each comaximal with
a fixed ideal remain comaximal with it, by multiplying elements
congruent to one modulo the fixed ideal. Hence \(I^k,\bar I^k\)
are comaximal, with product \(m_0^k O_K\), for \(k\ge1\).
For comaximal ideals the product equals the intersection, and
the quotient by their product maps onto the product of quotients:
if \(a+b=1\) as above then \(xb+ya\) realizes the given pair of
residues \(x,y\). As \(O_K\) has rank two over \(\Z\), and
conjugation exchanges the quotients, this gives
\([O_K:I^k]=[O_K:\bar I^k]=m_0^k\).

For every \(k\ge1\) choose a nonzero \(z_k\in I^k\) with
\(|z_k|\le2(1+|\omega|)m_0^{k/2}\) by the square grid
pigeonhole argument of Lemma~\ref{lem:quadratic-norm}.
Then \(J_k=z_k\bar I^k/m_0^k\subset O_K\), by \(I^k\bar I^k=m_0^kO_K\).
By the lattice area calculation of that lemma its index is
\(|z_k|^2/m_0^k\le4(1+|\omega|)^2\).
Only finitely many such sublattices occur, again by
Lemma~\ref{lem:quadratic-norm}. Equality of two \(J_k\)'s
for \(k<k'\) yields \(I^{k'}=b I^k\) for some \(b\in K^\times\)
by conjugation and rescaling. Multiplication by \(\bar I^k\)
and cancellation of the rational factor \(m_0^k\) give
\(I^h=bO_K,\ h=k'-k\ge1\). Thus \(b\in O_K\) and
\(|b|^2=m_0^h\) by the index just computed.

Choose \(s_0^{Dh}=b\) by Lemma~\ref{lem:place-fields}.
At any embedding \(\iota\), the least valuation of the generators
of \(P_l\) is zero except when \(p=l\) and \(\iota\bar\omega\)
has zero residue. In that case \(\iota\omega\) is a unit and
\(\bar\omega\omega=n_K\) shows \(v_l(\iota\bar\omega)\ge1\),
so the least valuation is one. The same minimum holds over all
nonzero members of the ideal, since all integral linear combinations
have valuation at least the minimum. Minima add on products, by
taking products realizing the two minima. It follows from \(I^h=bO_K\)
that \(v_p(\iota b)=h e_p\) on the selected branch at \(p\mid N_c\)
and is zero otherwise. The selected branch is exactly where
\(\varphi_1\omega=\iota\bar\omega\) has zero residue.
Division by \(Dh\) and Proposition~\ref{prop:local-cyclo-different}
give the required valuation.

Every complex automorphism sends \(b\in K\) to \(b\) or \(\bar b\).
Thus
\[
             |\sigma s_0|^2=m_0^{1/D}
                =N_c\prod_{l\mid N_c}l^{-1/(l-1)}
                =\operatorname{rd}_E/\sqrt d
\]
by Proposition~\ref{prop:different}. The last assertion follows from
the valuation formulation of Theorem~\ref{thm:every-frame-length}.
\end{proof}

\begin{corollary}\label{cor:phase-unique}
Any two nonzero complex numbers satisfying the finite-place
valuations of Theorem~\ref{thm:global-frame-scale} and its modulus
condition for automorphisms fixing \(K\) differ by a root of unity.
For any formal series in finitely many variables, substitution of
either scalar times every variable gives at each homogeneous degree
coefficients differing by one common root of unity. In particular
their absolute values agree after every archimedean or finite-place
embedding.
\end{corollary}
\begin{proof}
Write \(w\) for the quotient. Every complex automorphism \(\sigma\)
takes \(K\) to itself. If its restriction there is conjugation,
then \(\mathfrak c\sigma\) fixes \(K\) and
\(|\sigma w|=|\mathfrak c\sigma w|\). Thus \(|\sigma w|=1\)
for every \(\sigma\). The transcendence basis argument in
Proposition~\ref{prop:height-product} shows \(w\) is algebraic.
The finite-place hypotheses give \(v_p(\iota w)=0\) for every
embedding, so \(w\) is integral by the minimal polynomial
argument in Proposition~\ref{prop:frame-lattice-trace}, applied
in \(\Q(w)\). If its degree is \(D'\), each positive power has
degree at most \(D'\) and integer minimal polynomial. Every root
has modulus one, since each embedding extends to an automorphism
of \(\C\) by Lemma~\ref{lem:place-fields}. The elementary symmetric
coefficients are at most \(2^{D'}\) in absolute value by the
triangle inequality. Finitely many such monic integer polynomials
exist, and each has finitely many roots. Two distinct positive
powers of \(w\) therefore agree. Since \(w\ne0\), it is a root
of unity. A homogeneous coefficient of degree \(j\) acquires
the factor \(w^j\), which has absolute value one after any
embedding since it has finite multiplicative order.
\end{proof}

\section{The alternating form of the arithmetic frame lattice}

\begin{lemma}\label{lem:frame-ideal-product}
Let \(\mathcal A,D_c,t_K\) be as in
Definition~\ref{def:frame-lattice},
Proposition~\ref{prop:local-cyclo-different} and
Lemma~\ref{lem:quadratic-idempotent}. Inside \(E\) one has
\[
    \mathcal A=O_E+(t_K/D_c)O_E,\qquad
    \mathcal A\overline{\mathcal A}=D_c^{-1}O_E .
\]
Here the product of two \(O_E\)-modules consists of all finite
sums of products of their elements.
\end{lemma}
\begin{proof}
The valuation inequalities in Definition~\ref{def:frame-lattice}
are stable under addition and multiplication by \(O_E\).
Both \(1\) and \(t_K/D_c\) satisfy them. Indeed at a prime
\(l^a\Vert N_c\) the branch with no allowance has
\(v_l(t_K)\ge a\) by Lemma~\ref{lem:quadratic-idempotent},
while \(v_l(D_c)\le a\) on both branches by
Proposition~\ref{prop:local-cyclo-different}. At primes outside
\(N_c\) the denominator is a unit. Conversely, for
\(z\in\mathcal A\), both \(D_c z\) and \((1-t_K)z\) are
integral at every place. For the latter assertion the congruence
\(t_K\equiv1\bmod l^a\) at each branch with allowance absorbs the
possible denominator of \(z\). The integrality characterization
in Proposition~\ref{prop:frame-lattice-trace} thus puts both
elements in \(O_E\). The identity
\(z=(1-t_K)z+(t_K/D_c)(D_cz)\) proves the first formula.

The same proposition gives \(D_c x\bar y\in O_E\) for all
\(x,y\in\mathcal A\). For the reverse inclusion the product
contains \(1,t_K/D_c,\bar t_K/\bar D_c\).
The derivative reciprocity formula in
Proposition~\ref{prop:different} gives
\(\bar D_c/D_c=-\zeta^{2-f}\in O_E^\times\), so the product
also contains \(\bar t_K/D_c\).
Lemma~\ref{lem:quadratic-idempotent} gives
\(1-t_K-\bar t_K\in N_c\Z\), and
\(N_c/D_c\in O_E\) by Proposition~\ref{prop:local-cyclo-different}
and the same integrality characterization. Therefore
\[
 D_c^{-1}=t_K/D_c+\bar t_K/D_c+(1-t_K-\bar t_K)/D_c
           \in\mathcal A\overline{\mathcal A}.
\]
Closure under \(O_E\) proves the second equality.
\end{proof}

\begin{proposition}\label{prop:frame-polarization}
Fix \(n\ge1\), put \(g=fn\), and on \(E^n\) set
\[
  \mathcal E(x,y)=
     {\rm Tr}_{E/\Q}\!\left(\sum_{k=1}^n\bar x_k y_k/\sqrt{-d}\right)
     =\frac{2}{\sqrt d}
           \Im\!\left(\sum_{\tau\in G,k}
                      \overline{\tau(x_k)}\,\tau(y_k)\right).
\]
Set \(\sigma_0=(2\pi/\sqrt d)^{1/2}>0\), and
\(r(x)=(\sigma_0\tau(x_k))_{\tau,k}\in\C^g\).
On \(\C^g\) use \(H(u,z)=\sum_{\tau,k}\bar u_{\tau k}z_{\tau k}\)
and \(E_H(u,z)=\Im H(u,z)/\pi\).
Then \(E_H(rx,ry)=\mathcal E(x,y)\).
The group \(r(\mathcal A^n)\) is a full lattice, equal to its
dual for \(E_H\). For \(0\ne\delta\in O_K\) set
\(N=\delta\bar\delta\) and \(\Gamma=r(\bar\delta\mathcal A^n)\).
Its dual \(\Gamma^\#=\{v:E_H(v,\Gamma)\subset\Z\}\) satisfies
\[
 \Gamma^\#=r(\delta^{-1}\mathcal A^n)=N^{-1}\Gamma,\qquad
                  [\Gamma^\#:\Gamma]=N^{2g}.
\]
\end{proposition}
\begin{proof}
For \(z\in E\), taking conjugate pairs in the field trace gives
\[ {\rm Tr}_{E/\Q}(z/\sqrt{-d})=2\Im {\rm Tr}_{E/K}(z)/\sqrt d ,\]
since the selected embeddings fix \(K\) by
Proposition~\ref{prop:places}. Applying this identity to
\(\sum_k\bar x_k y_k\) proves both formulas for the alternating
form. The real extension of \(r:E^n\to\C^g\) is an isomorphism.
Indeed, in each copy of \(E\) the embeddings on the powers of
\(\zeta\) have the invertible evaluation matrix at the \(f\)
distinct primitive roots. On the basis \(1,\omega\) the
two embeddings of \(K\) have distinct values at \(\omega\).
The product basis of Theorem~\ref{thm:integer-ring} therefore
has invertible evaluation matrix on all \(2f\) embeddings.
A real linear relation on the selected \(f\) rows implies the
conjugate relation on the other \(f\), hence is zero. Equality
of real dimensions proves the claim.
Proposition~\ref{prop:frame-lattice-trace} sandwiches
\(\mathcal A\) between \(O_E\) and \(N_c^{-1}O_E\).
It is a full free lattice: in integer coordinates any subgroup
of \(\Z^h\) has an integer basis, by projecting to its first
coordinate, choosing a lift of a generator of that image if
nonzero, and adjoining an inductively chosen basis of the kernel.
Its rank here is \(2f\) by the two inclusions.

First test a rational point in one coordinate.
For \(x\in E\), the condition \(\mathcal E(x,\mathcal A)\subset\Z\)
is equivalent to
\(\bar x\mathcal A\subset D_c^{-1}O_E\).
In fact one can multiply any test element of \(\mathcal A\)
by every element of \(O_E\) and stay inside \(\mathcal A\).
The trace test therefore says precisely that
\(\bar x\mathcal A/\sqrt{-d}\subset O_E^\vee\), and
Proposition~\ref{prop:different} gives the stated equivalence.
If \(x\in\mathcal A\), Lemma~\ref{lem:frame-ideal-product}
gives the inclusion. Conversely, multiply the inclusion by
\(D_c\overline{\mathcal A}\) and use the same lemma to obtain
\(\bar x O_E\subset\overline{\mathcal A}\), proving \(x\in\mathcal A\).
Coordinate tests and addition give self-duality among points of
\(r(E^n)\). Every point of the real dual is in this rational
space. Relative to a basis of \(r(\mathcal A^n)\) its coordinates
are obtained by inverting the matrix of \(E_H\) on that basis
and applying it, up to transpose, to an integer column. The
matrix is integral by the rational test and is nonsingular,
since \(E_H(u,iu)=H(u,u)/\pi>0\) for \(u\ne0\). Its inverse
has rational entries by the adjugate formula.
Finally \(E_H(z,\bar\delta\,y)=E_H(\delta z,y)\), so the dual of
\(\Gamma\) is \(\delta^{-1}r(\mathcal A^n)\).
The two descriptions follow from \(\delta\bar\delta=N\);
division by \(N\) in each of \(2g\) independent integral basis
directions gives the stated index.
\end{proof}

\section{Theta products on the polarized torus}

\begin{lemma}\label{lem:cm-section-sign}
For the lattice \(\Lambda=\lambda O_K\) and functions
\(\theta,\sigma,G_\Lambda=\phi\theta^3\) of
Propositions~\ref{prop:cm-theta} and~\ref{prop:cm-duplication}, put
\(A=|\lambda|^2\sqrt d/2\).
For \(a,b,c\in\Z_{\ge0}\) let
\(P=\theta^a(\sigma^2)^bG_\Lambda^c\) and \(k=a+2b+3c\).
For \(y\in\C\) write
\[
 (U_y^{[k]}P)(z)=
       \exp\!\left(-k\pi(\bar y z+|y|^2/2)/A\right)P(z+y).
\]
This function is entire. For \(l=\lambda\alpha\in\Lambda\) it
satisfies \(U_l^{[k]}P=(-1)^{k|\alpha|^2}P\).
If \(y\in\lambda K\), the number \((U_y^{[k]}P)(0)\)
is an algebraic integer.
\end{lemma}
\begin{proof}
Write \(\alpha=u+v\omega\), \(u,v\in\Z\). The sign of
\(\sigma\) is \((-1)^{u+v+uv}\) by Proposition~\ref{prop:reduced-sigma}.
For \(\theta\), Proposition~\ref{prop:cm-theta} gives an additional
factor \(\exp(\pi(\bar l t-l\bar t)/A)=(-1)^v\), since
\(t=\lambda/2\). Its sign is therefore \((-1)^{u+uv}\). Now
\(|\alpha|^2=u^2+uv+n_K v^2\) and \(n_K=(d+1)/4\) is even
by Proposition~\ref{prop:places}. This proves the sign formula
in degree one. The sign in degree two of \(\sigma^2\) is one.
The meromorphic function \(\phi\) is periodic by its defining
formula in Proposition~\ref{prop:scaled-cm}, so \(G_\Lambda\)
has the sign of \(\theta^3\), first off its apparent singularities
and then everywhere by Proposition~\ref{prop:cm-duplication}.
Multiplication proves the law and entireness for \(P\).

For \(y\in\lambda K\), clearing the rational coefficients on the
basis \(1,\omega\) gives \(m y\in\Lambda\) for some \(m\ge1\).
Theorem~\ref{thm:local-cm-height} gives
\[
 |\iota P(y)|_p\le H_p(y)^k
   =|\iota\exp(\pi |y|^2/(2A))|_p^k
\]
at every finite isomorphism in Definition~\ref{def:place-height}.
Here the bound on the factor \(G_\Lambda\) is the cubic bound of
that theorem. Thus \(w=(U_y^{[k]}P)(0)\) has
\(|\iota w|_p\le1\) at every such isomorphism. To verify the
algebraic conclusion, suppose that \(w\) were transcendental.
Extend it to a transcendence basis of \(\C/\Q\). For each
\(j\ge1\) the assignment \(w\mapsto 2^{-j} w\), fixing the
remaining basis elements, gives an automorphism of the rational
function field and extends to \(\C\) by the argument of
Lemma~\ref{lem:place-fields}. Compose a fixed isomorphism at \(2\)
with these automorphisms. The resulting absolute values of the
images of \(w\) are \(2^j|\iota w|_2\), unbounded. Hence \(w\)
is algebraic. Every conjugate of \(w\) in each \(K_p\) occurs under
some such isomorphism by Lemma~\ref{lem:place-fields}. The rational
coefficients of its monic minimal polynomial are consequently
integral at every \(p\), being signed elementary symmetric functions
of these conjugates. Prime factorization in
Lemma~\ref{lem:basic-arith} puts them in \(\Z\).
\end{proof}

\begin{theorem}\label{thm:frame-sections}
Use \(g,r,H,E_H,\Gamma,\delta,N\) of
Proposition~\ref{prop:frame-polarization} and let
\((w_i,b_i)\) be a family in \(\mathcal A\) satisfying the
tight-frame identity of length \(L\) in
Theorem~\ref{thm:every-frame-length}.
Choose \(0\ne\alpha\in O_K\) with \(L|\alpha|^2=N\) and put
\(a'=\lambda\alpha\). For \(1\le k\le n\) set
\[
 B_{ik}(z)=\frac{a'}{\bar\delta\,\sigma_0}
              \sum_{\tau\in G}\overline{\tau(w_i)}z_{\tau k},
 \qquad
 \mathcal S(z)=\prod_{i,k}P_{ik}(B_{ik}(z)),
\]
where \(P_{ik}\) is any monomial of degree \(b_i\) in
Lemma~\ref{lem:cm-section-sign}. Define for \(v,z\in\C^g\)
\[
            (W_v S)(z)=\exp(-H(v,z)-H(v,v)/2)S(z+v).
\]
Then \(\mathcal S\) is entire. For \(u=r(h)\in\Gamma\),
\[
 W_u\mathcal S=\chi_\Gamma(u)\mathcal S,\qquad
 \chi_\Gamma(u)=(-1)^{{\rm Tr}_{E/K}\sum_k\bar h_k h_k}
               =\exp\bigl(\pi i E_H(\bar\omega u,u)\bigr).
\]
For \(\gamma\in E^n\) set \(y_{ik}=-B_{ik}(r\gamma)\).
Then
\[
       (W_{-r\gamma}\mathcal S)(z)
                  =\prod_{i,k}(U_{y_{ik}}^{[b_i]}P_{ik})(B_{ik}(z)).
\]
In particular its value at zero is an algebraic integer.
For every \(t'\in\C\) the Taylor series at zero after the substitution
\(z\mapsto t'z\), viewed as
a formal series in the coordinates of \(z\), is the product of
the one-variable Taylor series of the indicated translates with
\(B_{ik}(t'z)\) substituted. Taking \(P_{ik}=\theta^{b_i}\)
gives \(\mathcal S(0)=1\).
\end{theorem}
\begin{proof}
The tight-frame identity implies
\[
   \sum_{i,k} b_i\,\frac{\pi}{A}
                 \overline{B_{ik}(v)}B_{ik}(z)=H(v,z).
\]
Indeed in each block \(k\) the coefficient of
\(\bar v_{\tau k}z_{\nu k}\) in this sum equals
\(\pi |a'|^2 L{\bf1}_{\tau=\nu}/(A N\sigma_0^2)\).
This equals \({\bf1}_{\tau=\nu}\), by the formulas for
\(A,\sigma_0\) and \(L|\alpha|^2=N\).
It follows by multiplication of exponentials that the formula
for the translate holds for every \(v\) in place of \(-r\gamma\)
and shifts \(B_{ik}(v)\). The product of entire functions on
linear forms is entire. At \(v=-r\gamma\) we have
\[
       y_{ik}=-(a'/\bar\delta)\,{\rm Tr}_{E/K}(\gamma_k\bar w_i)
                  \in\lambda K
\]
by Proposition~\ref{prop:places}. Lemma~\ref{lem:cm-section-sign}
now gives integrality of the product at zero. To see explicitly
that products of algebraic integers are integral, the integer
span of products of powers of the factors, with each exponent
less than its monic degree, is finitely generated, contains one and
is preserved by multiplication by their product. On these
generators multiplication is given by an integral square matrix.
Its characteristic polynomial annihilates each generator by
the adjugate identity for \(XI-C\), and hence annihilates the
product acting on one. Taylor expansion and finite multiplication
give the asserted series identity by the disk formula of
Lemma~\ref{lem:analytic-tools}, or by absolute convergence
of the one-variable expansions on every bounded polydisk.

For \(h=\bar\delta\,x,\ x\in\mathcal A^n\),
\(B_{ik}(rh)=a'{\rm Tr}_{E/K}(x_k\bar w_i)\in\Lambda\) by
Proposition~\ref{prop:frame-lattice-trace}. The translation
formula and Lemma~\ref{lem:cm-section-sign} give the sign with
integer exponent \(\sum_{i,k}b_i|B_{ik}(rh)/\lambda|^2\).
Using the hermitian identity above at \(v=z=rh\) this exponent
is
\[
  \frac{A}{\pi|\lambda|^2}H(rh,rh)
    =\sum_{\tau,k}|\tau(h_k)|^2
    ={\rm Tr}_{E/K}\sum_k\bar h_k h_k .
\]
Since \(\Im\omega=\sqrt d/2\), the first expression is also
\(E_H(\bar\omega rh,rh)\). This proves the claimed multiplier.
Finally \(\theta(0)=1\) by Proposition~\ref{prop:cm-theta}.
\end{proof}

\begin{corollary}\label{cor:nonzero-frame-section}
Let \(R\ge0\) be an integer. For every positive integer \(b\)
such that \(b^2\ge24dN_c^4(1+\log_2 N_c)(R+1)\), the lattice
\(r(b\mathcal A^n)\) admits the section of
Theorem~\ref{thm:frame-sections} with value one at the origin
and every multiplicity at least \(\max(1,R)\).
\end{corollary}
\begin{proof}
Take \(\delta=b,\ L=N=b^2,\ \alpha=1\).
Theorem~\ref{thm:every-frame-length} supplies the frame.
Choose every factor to be the indicated power of \(\theta\)
in Theorem~\ref{thm:frame-sections}.
\end{proof}

\section{Fourier coordinates for the arithmetic theta space}

\begin{lemma}\label{lem:period-coordinates}
Use the data \(g,\Gamma,N,H,E_H\) of
Proposition~\ref{prop:frame-polarization}. There is an integral
basis \(e_1,\ldots,e_g,b_1,\ldots,b_g\) for \(\Gamma\) satisfying
\[
 E_H(e_j,e_k)=E_H(b_j,b_k)=0,\qquad E_H(e_j,b_k)=N{\bf1}_{j=k}.
\]
The \(e_j\) form a complex basis of \(\C^g\).
Write \(e z=\sum_jz_je_j\) and \(b_j=e B_j\), where \(B_j\)
is column \(j\) of a complex matrix \(B=X+iY\) with real
\(X,Y\). The matrices \(X,Y\) are symmetric and \(Y\) is positive
definite. The real Gram matrix \(M=(H(e_i,e_j))\) satisfies
\(MY=\pi N I\).
\end{lemma}
\begin{proof}
First work on the lattice \(\Gamma_0=r(\mathcal A^n)\) of
Proposition~\ref{prop:frame-polarization}. The map
\(x\mapsto E_H(x,\cdot)\) identifies it with its integral dual:
invert the nonsingular real pairing to represent any integer
linear functional and apply self-duality. Take a basis element
\(u\). Its image is primitive since the map and its inverse
are integral. Its values on an integral basis have gcd one,
so Lemma~\ref{lem:basic-arith} supplies \(v\) with \(E_H(u,v)=1\).
For each \(x\) its difference from
\(E_H(x,v)u+E_H(u,x)v\) is perpendicular to both. This splits
off \(\Z u+\Z v\) orthogonally. On the remaining direct summand
the pairing is again unimodular, since any integral linear
functional on it can be extended by zero to the first summand,
and the representing vector of such an extension is perpendicular
to \(u,v\). Induction, with the free subgroup argument of
Proposition~\ref{prop:frame-polarization}, gives \(g\) pairs.
Multiplying the basis by \(\bar\delta\) multiplies \(H\), hence
\(E_H\), by \(N\), producing the required basis for \(\Gamma\).

On the real span \(U\) of the \(e_j\), \(E_H\) vanishes. If
\(u,iu\in U\), then \(0=E_H(u,iu)=H(u,u)/\pi\), so \(u=0\).
Thus \(U\cap iU=0\), and real independence proves complex
independence of the \(g\) vectors. The Gram matrix \(M\) is
real symmetric positive definite. In these coordinates
\(H(eu,ez)=\bar u^t M z\).
The equations \(E_H(e_j,b_k)=N{\bf1}_{j=k}\) give
\(MY=\pi N I\), so \(Y=\pi N M^{-1}\) is symmetric positive
definite. Finally the imaginary part of
\(\bar B^{\,t} MB\) is \(X^t M Y-Y^t M X=\pi N(X^t-X)\).
The vanishing of \(E_H(b_j,b_k)\) proves symmetry of \(X\).
\end{proof}

\begin{theorem}\label{thm:period-dimension}
Let \(\mathcal T(\Gamma)\) denote the complex vector space of
entire functions on \(\C^g\) such that \(W_uS=\chi_\Gamma(u)S\)
for all \(u\in\Gamma\), with \(W,\chi_\Gamma\) of
Theorem~\ref{thm:frame-sections}. Then
\(\dim_\C\mathcal T(\Gamma)=N^g\).
In the basis of Lemma~\ref{lem:period-coordinates} choose a
real column \(t\) with \(\exp(2\pi i t_j)=\chi_\Gamma(e_j)\).
For \(S\in\mathcal T(\Gamma)\), the function
\[
      h(z)=\exp(-z^t M z/2-2\pi i t^t z)\,S(ez)
\]
has a normally convergent expansion
\(\sum_{m\in\Z^g} a_m\exp(2\pi i m^t z)\), meaning absolute
uniform convergence on compact sets. The coefficients
\(\{a_m:0\le m_j<N\}\) are arbitrary independent coordinates.
Thus \(\mathcal T(\Gamma_0)\) has dimension one for
\(\Gamma_0=r(\mathcal A^n)\).
\end{theorem}
\begin{proof}
For \(u,v\in\C^g\), multiplying the exponentials defining \(W\)
gives
\[
             W_u W_v=\exp(\pi i E_H(u,v))\,W_{u+v}.
\]
For \(u=r(h),v=r(h')\) in \(\Gamma\),
\(T={\rm Tr}_{E/K}\sum_k\bar h_k h'_k\in O_K\) by
Proposition~\ref{prop:frame-lattice-trace}, since
\(\bar\delta\mathcal A\subset\mathcal A\).
Write \(T=a+c\omega\) with \(a,c\in\Z\).
The cross term in the integer exponent of
\(\chi_\Gamma(u+v)\) is \(2\Re T=2a+c\);
by Proposition~\ref{prop:frame-polarization}, \(E_H(u,v)=c\).
Hence
\[
     \chi_\Gamma(u+v)
         =\exp(-\pi i E_H(u,v))\chi_\Gamma(u)\chi_\Gamma(v).
\]
These two formulas show that it suffices to impose the equations
for the generators of the integral basis. They also give the
equations for negatives because \(W_{-u}=W_u^{-1}\).

The equations at the \(e_j\) say precisely that \(h\) is periodic
under \(\Z^g\). The equations at the \(b_j\) become
\[
 h(z+B_j)=d_j\exp(-2\pi i N z_j)\,h(z),\qquad
 d_j=\chi_\Gamma(b_j)\exp(-2\pi i t^t B_j-\pi i N B_{jj}).
\]
Indeed the linear exponent contributed by the Gaussian equals
\((\bar B_j-B_j)^t M z=-2\pi i N z_j\), and its constant
contribution is \((\bar B_j-B_j)^t M B_j/2=-\pi i N B_{jj}\),
by Lemma~\ref{lem:period-coordinates}.

For a periodic entire \(h\) define \(a_m\) by integrating
\(h(z)\exp(-2\pi i m^t z)\) over the real unit cube. Moving this
cube to the parallel cube of imaginary part \(y\in\R^g\)
leaves the integral unchanged: in each coordinate integrate on
the boundary of the intervening rectangle, whose vertical sides
cancel by periodicity, and use the disk and contour identities of
Lemma~\ref{lem:analytic-tools}. Iteration over coordinates is valid
by continuity on the compact rectangles and
Lemma~\ref{lem:integral-tools}. Taking \(y_j=-A\,{\rm sign}(m_j)\)
shows for every \(A>0\) that
\[
 |a_m|\le C_A\exp\!\left(-2\pi A\sum_j|m_j|\right)
\]
where \(C_A\) is a maximum of \(|h|\) over a compact box.
The Fourier series therefore converges normally by products
of geometric sums. Its restriction to the real cube has the same
coefficients as \(h\), so it agrees there by the product
Fej\'er-kernel argument of Lemma~\ref{lem:gaussian-cover}.
Successive application of the one-variable identity principle of
Lemma~\ref{lem:analytic-tools} gives agreement on \(\C^g\).

Let \(\mathbf e_j\) be the \(j\)-th integer coordinate column.
The equations at \(B_j\) are now equivalent, by coefficient
comparison, to
\[
  a_{m+N\mathbf e_j}=d_j^{-1}\exp(2\pi i m^t B_j)\,a_m .
\]
For \(m_0\) a residue representative \(0\le (m_0)_j<N\), their
unique solution throughout its residue class is
\[
 \begin{split}
 a_{m_0+N\ell}={}&a_{m_0}\left(\prod_j d_j^{-\ell_j}\right)
    \exp\!\left(2\pi i m_0^t B\ell
        +\pi i N\left(\ell^tB\ell-\sum_j\ell_j B_{jj}\right)\right)
                \qquad(\ell\in\Z^g).
 \end{split}
\]
Substitution at \(\ell+\mathbf e_j\) verifies the recursion
because \(B^t=B\). Each multiplier is nonzero, so the same
recursions determine all negative steps. Conversely give the
finitely many \(a_{m_0}\) arbitrary values and use this formula.
Since \(Y\) is positive definite, its quadratic form has a
positive minimum on the real unit sphere. Thus the coefficients
decay as \(O(\exp(-c\|\ell\|^2+C\|\ell\|))\) with \(c>0\).
On compact subsets the additional Fourier exponent and each
derivative introduce at most an exponential linear in \(\|\ell\|\)
and a polynomial. The series and its differentiated series
converge there by products of one-dimensional Gaussian sums,
as in Proposition~\ref{prop:reduced-sigma}. They give an entire
function with the required translation equations. Reversing
the defining transformation of \(h\) constructs \(S\).
The coordinate count and its specialization \(\delta=1\) follow.
\end{proof}

\section{Orthogonality of the theta translates}

\begin{proposition}\label{prop:theta-shift-irrep}
On \(\mathcal T(\Gamma)\) of Theorem~\ref{thm:period-dimension},
the operators \(W_v\) for \(v\in\Gamma^\#\) preserve the space.
Every nonzero vector spans \(\mathcal T(\Gamma)\) under these
translates. Every complex-linear operator commuting with all of
them is scalar. It suffices in each assertion to use one
representative of each class in \(\Gamma^\#/\Gamma\).
\end{proposition}
\begin{proof}
The addition law proved in Theorem~\ref{thm:period-dimension} gives
\(W_uW_v=\exp(2\pi i E_H(u,v))W_vW_u\). For \(u\in\Gamma\) and
\(v\in\Gamma^\#\) these operators commute, proving preservation.
Changing \(v\) by an element of \(\Gamma\) multiplies its operator
on \(\mathcal T(\Gamma)\) by a nonzero scalar by the same law.

Let \(e_j,b_j,N,t\) be as in Lemma~\ref{lem:period-coordinates}
and Theorem~\ref{thm:period-dimension}.
Write \(U_j=\exp(-2\pi i t_j/N)W_{e_j/N}\) and
\(V_j=W_{b_j/N}\). These shifts are in the dual lattice by
Proposition~\ref{prop:frame-polarization}.
In the Fourier coordinates of the dimension theorem,
\(U_j\) replaces \(h(z)\) by \(h(z+\mathbf e_j/N)\).
Indeed the real translation by \(\mathbf e_j/N\) cancels the
linear and constant terms from the quadratic exponential,
leaving exactly \(\exp(2\pi i t_j/N)\) before normalization.
It follows that the line of coefficients supported on the
residue class \(m\bmod N\) is an eigenspace of each \(U_j\)
with eigenvalue \(\exp(2\pi i m_j/N)\). The lines exhaust the
space by that theorem. Its projection on the line of \(m\) is
\[
      \prod_{j=1}^g\left(
       \frac1N\sum_{h=0}^{N-1}\exp(-2\pi i h m_j/N)\,U_j^h\right)
\]
by the finite geometric sum on each eigenspace. Moreover
\[
 U_i V_j=\exp(2\pi i{\bf1}_{i=j}/N)V_j U_i
\]
by the commutation relation. Each \(V_j\) therefore carries the
line of \(m\) onto the line of \(m+\mathbf e_j\), since it has
inverse \(W_{-b_j/N}\). These moves are transitive on all residue
classes, also when \(N=1\). The span of the translates of a
nonzero vector is stable under all shifts by the addition law,
so the displayed projections recover a nonzero line in that span,
and the moves recover all lines. An operator commuting with the
\(U_j\) preserves each line, and commuting with the \(V_j\) makes
its scalars equal on all of them. This proves both conclusions.
\end{proof}

\begin{theorem}\label{thm:vacuum-sum}
Fix a nonzero \(V_0\in\mathcal T(\Gamma_0)\) with
\(\Gamma_0=r(\mathcal A^n)\), and choose any set \(D\) of
representatives for \(\Gamma^\#/\Gamma\). Then
\(|D|=N^{2g}\) and each \(a\in D\) has the form \(r(x)\) with
\(\delta x\in\mathcal A^n\). For every \(F\in\mathcal T(\Gamma)\)
there are complex numbers \(c_a\) such that
\[
 F=\sum_{a\in D} c_a W_a V_0 .
\]
For every complex-linear functional \(\ell\) on \(\mathcal T(\Gamma)\)
one has the identity in that space
\[
         N^g\sum_{a\in D}\ell(W_a V_0) W_{-a}F
                             =N^{2g}\ell(F)V_0 .
\]
In particular this formula holds for all of the sections in
Theorem~\ref{thm:frame-sections}, with the same set \(D\).
\end{theorem}
\begin{proof}
Proposition~\ref{prop:frame-polarization} gives the cardinality
and the rational representatives. Theorem~\ref{thm:period-dimension}
gives the nonzero one-dimensional space at \(\Gamma_0\).
Since \(\Gamma\subset\Gamma_0\) and the sign in
Theorem~\ref{thm:frame-sections} has the same expression on both
lattices, \(V_0\in\mathcal T(\Gamma)\).
The spanning statement follows from
Proposition~\ref{prop:theta-shift-irrep}.

For a linear operator \(T\) on \(\mathcal T(\Gamma)\), put
\(A(T)=\sum_{a\in D}W_{-a}TW_a\).
The conjugation in a summand depends only on the class of \(a\),
because \(W_{-a}=W_a^{-1}\) and a scalar in \(W_a\) cancels.
Conjugating \(A(T)\) by any \(W_b\), \(b\in\Gamma^\#\), permutes
these classes by translation by \(b\) or \(-b\), by the addition
law in Theorem~\ref{thm:period-dimension}. Thus \(A(T)\) is a
scalar by Proposition~\ref{prop:theta-shift-irrep}.
In coordinates \({\rm tr}(BC)=\sum_{i,j}B_{ij}C_{ji}={\rm tr}(CB)\).
Consequently its trace is \(|D|{\rm tr}(T)\), whereas the trace of
a scalar is \(N^g\) times that scalar by
Theorem~\ref{thm:period-dimension}. Choose
\(T(z)=\ell(z)F\). Its trace is \(\ell(F)\), by the coordinates of
the column \(F\) and the row \(\ell\). Applying the formula for
\(A(T)\) to \(V_0\) and multiplying by \(N^g\) gives the claimed
identity. The final application uses the translation equation of
Theorem~\ref{thm:frame-sections}.
\end{proof}

\section{The finite Heisenberg action on theta sections}

\begin{definition}\label{def:finite-frame-space}
Put \(L_0=\mathcal A^n,\ \Gamma_0=r(L_0)\), and let
\[
        \mathcal F=\bigcup_{m\ge1}\mathcal T(m\Gamma_0).
\]
Here \(\mathcal T(m\Gamma_0)\) uses \(\delta=m\) in
Theorem~\ref{thm:period-dimension}. For \(m\mid m'\) the
larger periods \(m'\Gamma_0\) impose a subset of the equations
at \(m\Gamma_0\), with the same sign. The union is directed by
divisibility and thus is a complex vector space of entire functions.
In the rational space \(E^n\) put
\[
 \widehat L_0=\prod_p (L_0\otimes\Z_p),\qquad
           V_f=\bigcup_{D\ge1} D^{-1}\widehat L_0 .
\]
Its elements are finite adeles, and the subgroups
\(h\widehat L_0,\ h\ge1\) form a basis of open neighborhoods
at zero. Extend the rational alternating form \(\mathcal E\)
of Proposition~\ref{prop:frame-polarization} by scalars to
\(\mathcal E_p\) at each \(p\), and set
\[
 c(X,Y)=\prod_p\mathbf e_p(\mathcal E_p(X_p,Y_p)/2)
                      \qquad(X,Y\in V_f).
\]
\end{definition}

\begin{theorem}\label{thm:finite-frame-action}
There is a unique family of complex-linear operators \(T_X\)
on \(\mathcal F\), indexed by \(X\in V_f\), for which \(T_X S\)
is locally constant in \(X\) for every \(S\), and
\(T_\gamma S=W_{-r(\gamma)}S\) for \(\gamma\in E^n\).
They satisfy
\[
 T_X T_Y=c(X,Y)T_{X+Y},\qquad T_0=1,\qquad T_X^{-1}=T_{-X}.
\]
With \(V_0\) as in Theorem~\ref{thm:vacuum-sum}, \(\mathcal F\)
is spanned by the \(T_X V_0\). The linear map
\[
       S\longmapsto F_S,\qquad F_S(X)=(T_X S)(0)
\]
is injective and has locally constant values in the variable \(X\).
It satisfies \(F_{T_Y S}(X)=c(X,Y)F_S(X+Y)\).
\end{theorem}
\begin{proof}
Relative to a \(\Z\)-basis of \(L_0\) provided by
Proposition~\ref{prop:frame-polarization},
\(\widehat L_0=\prod_p\Z_p^{2g}\). For \(D,h\ge1\), every
\(X\in D^{-1}\widehat L_0\) has a representative
\(\gamma\in D^{-1}L_0\) modulo \(h\widehat L_0\). Indeed reduce
each coordinate of \(DX\) modulo \(Dh\) at the primes dividing
\(Dh\) and combine the residues to an integer coordinate by
Lemma~\ref{lem:basic-arith}. At all other primes \(Dh\) is a unit.
At rational points \(E^n\cap h\widehat L_0=h L_0\), since a
rational coordinate integral at every prime is integral. These
facts prove density and identify the intersections used below.
The form \(\mathcal E_p\) is integral on \(L_0\otimes\Z_p\) by
its integral matrix in Proposition~\ref{prop:frame-polarization}.
Thus almost all local factors of \(c\) equal one by
Proposition~\ref{prop:additive}. The product is a bicharacter,
and on diagonal \(\gamma,\eta\in E^n\) the same proposition gives
\[
             c(\gamma,\eta)=\exp(\pi i\mathcal E(\gamma,\eta)).
\]
It is locally constant on \(V_f^2\): for both arguments in
\(D^{-1}\widehat L_0\), perturbations in \(2D\widehat L_0\)
change \(\mathcal E_p/2\) by \(p\)-integral quantities everywhere.
On rational arguments the displayed law for \(T_\gamma T_\eta\)
therefore follows from Theorem~\ref{thm:period-dimension}.

Suppose \(S\in\mathcal T(m\Gamma_0)\) and
\(\gamma\in D^{-1}L_0\). If \(h\in 2mD L_0\), then
\(W_{-rh}S=S\). In fact this is a period with sign one by
the quadratic formula of Theorem~\ref{thm:frame-sections},
since \(h\) is twice an element of \(L_0\).
Also \(\mathcal E(\gamma,h)\in2\Z\).
The rational addition law thus gives
\(T_{\gamma+h}S=T_\gamma S\).
The function \(T_\gamma S\) lies in
\(\mathcal T(2mD\Gamma_0)\): each period in the latter lattice
already acts on \(S\) by its required sign, and commutes with
\(W_{-r\gamma}\) since its pairing with \(r\gamma\) is integral.
This proves that rational operators carry \(\mathcal F\) into itself.

For arbitrary \(X\in D^{-1}\widehat L_0\), take a rational
representative in \(D^{-1}L_0\) modulo \(2mD\widehat L_0\)
and use its action on \(S\). Two such representatives differ in
\(2mD L_0\), so the definition is independent of the choice.
If \(D\) or the level \(m\) is enlarged to a multiple, the
finer choice has the same action: it differs modulo the older
neighborhood and lies in the old \(D^{-1}L_0\) when \(X\) lies
in \(D^{-1}\widehat L_0\). For the latter assertion, it is a
rational point of \(D^{-1}\widehat L_0\), so the coordinate
test above applies. Passing to common multiples proves independence
for any two choices of \(m,D\). The action is locally constant
on each open \(D^{-1}\widehat L_0\); density gives uniqueness.
Linearity follows by choosing a common level for finitely many
sections and a common representative. The rational addition law
extends to all \(X,Y\) by local constancy and density. Indeed
\(T_Y S\) is constant near each fixed \(Y\), and applying \(T_X\)
to that one section is locally constant near any fixed \(X\).
The identity and inverse statements use alternation of
\(\mathcal E_p\).

Theorem~\ref{thm:vacuum-sum} at \(\delta=m\) spans the whole
\(m\)-th level by rational translates of \(V_0\), proving the
assertion for the union. For rational \(\gamma\), the definition
of \(W\) gives
\[
        F_S(\gamma)=\exp(-H(r\gamma,r\gamma)/2)S(-r\gamma).
\]
The vectors \(r(E^n)\) contain the rational span of a full real
lattice by Proposition~\ref{prop:frame-polarization}, so they
are dense in \(\C^g\). Thus \(F_S=0\) implies \(S=0\) by
continuity. Evaluation of the operator law at zero gives the
translation formula for \(F_S\).
\end{proof}

\begin{corollary}\label{cor:integral-adelic-sections}
For each product \(S\) constructed in Theorem~\ref{thm:frame-sections},
\(F_S(X)\) is an algebraic integer for every \(X\in V_f\).
For \(y\in E^n\) the radial Taylor series of \(T_y S\) at the
origin has the product formula of that theorem.
The finite sum of Theorem~\ref{thm:vacuum-sum} also holds after
applying \(S\mapsto F_S\) to its two sides. In this identity every
shift \(r(x)\) can be represented by \(T_{-x}\) with
\(\delta x\in L_0\).
\end{corollary}
\begin{proof}
Since \(N=\delta\bar\delta\), \(N\Gamma_0\subset\Gamma\).
Hence \(S\in\mathcal T(N\Gamma_0)\).
By Theorem~\ref{thm:finite-frame-action}, for a given \(X\) its
translate of \(S\) equals the translate at one rational
representative \(\gamma\in E^n\).
Theorem~\ref{thm:frame-sections} proves integrality of its value
at zero and gives the entire product formula for every rational
translate. The vacuum identity follows by linearity and the
description of its shifts in Theorem~\ref{thm:vacuum-sum}.
\end{proof}

\section{Counting the nonzero frequencies in the orbital sieve}

\begin{lemma}\label{lem:denominator-series}
Let \(O=O_E\). For \(0\ne J\subset O\) an integral ideal put
\(\mathrm NJ=[O:J]\). For each rational prime \(p\), write
\(O\otimes\Z_p=\prod_{w\mid p} O_w\) as in the proof of
Theorem~\ref{thm:integer-ring}. Denote the valuation with value
one at a uniformizer by \({\rm ord}_w\) and the cardinality of
the residue field by \(q_w\). At each such \(w\) the completion
of \(J\) is the ideal with some integral exponent \(b_w\ge0\).
The exponents determine \(J\), are zero except at finitely many
places, and \(\mathrm NJ=\prod_w q_w^{b_w}\).
For real \(s>1\),
\[
     \sum_{0\ne J\subset O}(\mathrm NJ)^{-s}
             \le\left(\sum_{h\ge1}h^{-s}\right)^{2f}.
\]
For \(0\ne y\in E\), put \(J(y)=\{z\in O:yz\in O\}\). Its
exponents are \(b_w=\max(0,-{\rm ord}_w y)\), and
\[
                    \prod_{\tau\in G}|\tau y|^2
                       =\prod_w q_w^{{\rm ord}_w y}.
\]
\end{lemma}
\begin{proof}
Choose \(0\ne v\in J\). Multiplication by \(v\) on the integral
basis in Theorem~\ref{thm:integer-ring} has nonzero integer
determinant. The positive absolute value \(a\) of that determinant
satisfies \(a O\subset vO\subset J\) by the adjugate. Thus \(O/J\)
is finite. The ring \(O/aO\) is the product of \(O_w/aO_w\)
over \(w\mid p,\ p\mid a\), by the Chinese remainder calculation in
Lemma~\ref{lem:finite-quotient-gens}. Coordinate idempotents split
ideals of that product. In each factor every nonzero preimage ideal
is a uniformizer power, by taking a member of smallest valuation.
Each successive quotient of two adjacent powers has \(q_w\) elements
by multiplication by a uniformizer power. This proves the claimed
description modulo \(aO\), and the index formula. In comparing two
ideals take a common multiple of their two integers \(a\); the
description proves uniqueness from all the exponents. At most
\(2f\) factors occur over \(p\), since each contributes at least
one dimension to \(O/pO\) over \(\mathbf F_p\). Also \(q_w\ge p\).
For any finite set of ideals, summing over all possible exponents
at the primes in its support bounds its sum by the finite product
of \((1-p^{-s})^{-2f}\) there. Expanding the geometric series
and using prime factorization in Lemma~\ref{lem:basic-arith}
bounds this by the displayed power of \(\sum h^{-s}\). The latter
sum converges by integration of \(t^{-s}\) from one to infinity.
The inequality for the nonnegative infinite sum follows.

Choose an integer \(a\ge1\) with \(ay\in O\). Membership of
\(z\in O\) in \(J(y)\) says precisely that \((ay)z\in aO\).
Testing this finite congruence in each of the factors just
described gives the asserted exponents of \(J(y)\).
For integral \(v\ne0\) the index of \(vO\) is also
\(\prod_\tau|\tau v|^2\): on the real span, multiplication
by \(v\) is diagonal on the complex embeddings, of real determinant
this product, by Proposition~\ref{prop:frame-polarization};
index equals the ratio of fundamental volumes as in
Lemma~\ref{lem:class-covol}. The formula for \(J=vO\) gives
\(\prod_w q_w^{{\rm ord}_w v}\) for the same index.
Taking the quotient of these identities at \(ay\) and \(a\)
proves the formula for \(y\).
\end{proof}

\begin{lemma}\label{lem:ideal-polydisk}
For \(J\) as in Lemma~\ref{lem:denominator-series} put
\(J^{-1}=\{y\in E:yJ\subset O\}\).
For real numbers \(R_\tau\ge1\), the number of \(y\in J^{-1}\)
with \(|\tau y|\le R_\tau\) for all \(\tau\) is at most
\[
                     9^f\,\mathrm NJ\prod_\tau R_\tau^2 .
\]
\end{lemma}
\begin{proof}
For distinct \(y,y'\) in this set put \(z=y-y'\).
Then \(zJ\subset O\) is a full sublattice and
\[
             1\le[O:zJ]=\mathrm NJ\prod_\tau|\tau z|^2
\]
by the same real determinant and fundamental-volume calculation
as in Lemma~\ref{lem:denominator-series}.
After dividing coordinate \(\tau\) by \(R_\tau\), the largest
coordinate distance between the two points is at least
\(d_*=(\mathrm NJ\prod_\tau R_\tau^2)^{-1/(2f)}\).
Open balls of radius \(d_*/2\) in this maximum norm are disjoint
and lie in the ball of radius \(1+d_*/2\).
Their real volumes are proportional to the radius to power \(2f\),
because the ball is a product of \(f\) complex disks.
Thus every finite subset has cardinality at most
\((1+2/d_*)^{2f}\le (3/d_*)^{2f}\), since \(0<d_*\le1\).
This gives the stated finite count for the whole set.
\end{proof}

\begin{theorem}\label{thm:denominator-frequency}
Let \(T\) be any finite set of finite places of \(E\).
For a nonzero \(T\)-integer \(y\), meaning
\({\rm ord}_w y\ge0\) at \(w\notin T\), put
\[
 G_T(y)=\prod_{w\in T}q_w^{{\rm ord}_w y},\qquad
 D_T(y)=\prod_{w\in T}q_w^{\max(0,-{\rm ord}_w y)} .
\]
For \(s,t>2,\ j\ge0\) real, \(c_\tau\in\C^\times\), and
\(X_\tau\ge |c_\tau|^{-1}\), put
\(B_\tau=\max(1,|c_\tau|^{-1})\) and
\(d_{j,t}(u)=u^j\min(1,u)^t\) for \(u>0\).
The sum over all nonzero \(T\)-integers satisfies
\[
 \sum_y\frac{G_T(y)^{j/2}}{D_T(y)^s}
             \prod_\tau d_{j,t}(|c_\tau\tau y|^{-1})
 \le \left(\frac{9t}{t-2}
            \left(\sum_{h\ge1}h^{-(s-1)}\right)^2\right)^f
                                  \prod_\tau B_\tau^2 X_\tau^j .
\]
For \(s=t=6\) the factor in parentheses is less than \(15\).
In particular, if \(B\ge1\) and \(X_\tau\le B\) for all \(\tau\),
the sum for these exponents is at most
\((15B^2)^f\prod_\tau X_\tau^j\).
\end{theorem}
\begin{proof}
The generalized index of multiplication by \(y\) on the local
integer rings at \(T\) is \(G_T(y)\): multiplication by a
uniformizer scales additive volume by \(q_w^{-1}\), because its
image in \(O_w\) has that index by
Lemma~\ref{lem:denominator-series}. Multiplication by a unit
preserves the ring and its index. The denominator index, the
index in \(\prod_{w\in T}O_w\) of the elements whose product by
\(y\) stays integral, is \(D_T(y)\) by the same lemma.
The \(T\)-integer condition gives
\[
 D_T(y)=\mathrm N J(y),\qquad
       G_T(y)^{1/2}\le\prod_\tau|\tau y|.
\]
The second inequality follows by inserting the factors with
nonnegative exponents outside \(T\) in its norm formula.
All factors containing power \(j\) in a summand can therefore
be bounded together by \(\prod_\tau X_\tau^j\).

It remains to sum \(D_T(y)^{-s}\prod_\tau
\min(1,B_\tau/|\tau y|)^t\). Fix \(J(y)=J\); all those \(y\)
belong to \(J^{-1}\). For each \(u>0,\ B'>0\) we have
\[
    \min(1,B'/u)^t=\int_1^\infty t v^{-t-1}
                               {\bf1}_{u\le B'v}\dd v
\]
by integration of a power. On a finite set of points with
denominator \(J\), apply this identity in each coordinate and
interchange the finite sum and the integrals. The resulting count
in the product of disks of radii \(B_\tau v_\tau\) is at most
\(9^f\mathrm N J\prod_\tau B_\tau^2 v_\tau^2\) by
Lemma~\ref{lem:ideal-polydisk}. The integrals now separate,
each giving \(t/(t-2)\). Multiplication by \((\mathrm N J)^{-s}\)
and summation over \(J\) use the ideal-series estimate of
Lemma~\ref{lem:denominator-series} at \(s-1>1\).
This proves the bound on every finite sum, hence on the full
nonnegative sum. With \(s=t=6\),
\(\sum_{h\ge1}h^{-5}\le1+2^{-5}+\int_2^\infty u^{-5}\dd u=67/64\).
The numerical inequality \((27/2)(67/64)^2<15\) proves
the last two assertions.
\end{proof}

\section{Positive Fourier kernels on the adelic quotient}

\begin{definition}\label{def:poisson-coordinates}
Write \(E_\infty=\C^G,\ \sigma a=(\tau a)_\tau\),
\(\widehat O=\prod_w O_w,\ E_f=\bigcup_{D\ge1}D^{-1}\widehat O\),
and \(\mathbb A_E=E_\infty\times E_f\). Embed \(a\in E\)
diagonally as \([a]=(\sigma a,a)\).
Use Lebesgue measure on \(E_\infty\) and additive measure on
\(E_f\) normalized by \(\mu(\widehat O)=1\). For \(y\in E\) put
\[
   \mathcal P_y(z,x)=
    \exp\!\left(4\pi i\,\Re\sum_\tau(\tau y)z_\tau\right)\psi_f(-yx).
\]
For integrable \(G_1\) on \(\mathbb A_E\) write
\(\widehat G_1(y)=\int_{\mathbb A_E}\mathcal P_y G_1\).
The mass of a fundamental parallelepiped for \(\sigma O_E\) is
\(V_E=2^{-f}\sqrt{|\Delta_E|}\) by Lemma~\ref{lem:skew-covol}.
\end{definition}

\begin{lemma}\label{lem:poisson-finite-torus}
The measure of each coset of \(m\widehat O\) is \(m^{-2f}\)
for integral \(m\ge1\). Every \(x\in E_f\) is congruent to
some \(a\in E\) modulo \(m\widehat O\), uniquely modulo \(m O_E\).
Thus
\[
  p_m(z,x)=z-\sigma a\bmod\sigma(m O_E)
\]
defines an additive map to a real torus which is trivial on
the diagonal \(E\).
The integral Fourier indices of this torus are parametrized by
\(y\in(m\varpi)^{-1}O_E\), with \(\varpi\) from
Proposition~\ref{prop:different}. The corresponding exponentials
pull back along \(p_m\) to the functions \(\mathcal P_y\) of
Definition~\ref{def:poisson-coordinates}.
\end{lemma}
\begin{proof}
In an integral basis of \(O_E\) the quotient by \(m\widehat O\)
inside \(\widehat O\) has \(m^{2f}\) elements, by the coordinate
Chinese remainder argument of Theorem~\ref{thm:finite-frame-action}.
The measure used here can be constructed by successively partitioning
\([0,1)\) into equally long half-open intervals labelled by compatible
residue classes modulo \(j!\widehat O\), \(j=1,2,\ldots\).
At each step divide an interval equally among the lifts of its
class. Compatible classes specify a unique point of \(\widehat O\)
by the \(p\)-adic digit description in Lemma~\ref{lem:place-fields}.
Pushing forward length gives probability \(m^{-2f}\) for every
coset modulo \(m\widehat O\), by choosing a divisible factorial.
This measure is translation invariant. Indeed translation permutes
cosets at each level. Equality for their finite unions extends to
all Borel sets: the class of equality for two probability measures
is closed under complementation and countable disjoint union.
For a family of sets closed under intersections, its smallest
class with these two closure properties is also closed under
intersections. To see this, fix one set of the family and apply
the two closure properties to the sets whose intersection with
it is in the class, using differences of nested sets. Then fix
an arbitrary set of the class and repeat. Such a class is a
sigma algebra, since countable unions can now be made disjoint.
The cosets generate the Borel sets because they form a countable
base. Extend to \(E_f\) by summing translates of this measure on
its countably many \(\widehat O\)-cosets. Their countability, the
congruence assertion for \(x\), and its uniqueness follow by the
proof of Theorem~\ref{thm:finite-frame-action} with any integral
basis of \(O_E\). These facts give the asserted properties of \(p_m\).

The real pairing \((y,z)\mapsto 2\Re\sum_\tau y_\tau z_\tau\)
is nondegenerate, by testing a single coordinate against its
complex conjugate. On diagonal vectors it is the rational trace
pairing. An integral Fourier index means a real linear form with
integer values on \(\sigma(mO_E)\). In the basis of \(mO_E\) its
dual for this pairing is rational: the rational trace matrix is
invertible by Proposition~\ref{prop:different}. The same proposition
gives that lattice of indices as \((m\varpi)^{-1}O_E\).
For such \(y\), \(\psi_f(-y(x-a))=1\) since \(x-a\in m\widehat O\)
and \(my\) belongs to the trace dual at each place, by
Proposition~\ref{prop:additive}. The same proposition gives
\(\psi_f(-ya)=\exp(-2\pi i{\rm Tr}_{E/\Q}(ya))\).
This is precisely the phase in the pullback of the exponential
at \(z-\sigma a\).
\end{proof}

\begin{theorem}\label{thm:one-sided-adelic}
Let \(G_1:\mathbb A_E\to\R_{\ge0}\) be integrable, and \(A_0\subset E\)
finite. Suppose \(G_1\) is lower semicontinuous at each \([a]\),
\(a\in A_0\), meaning that for every \(\eta>0\) its value nearby
is at least \(G_1([a])-\eta\). Then
\[
 \sum_{a\in A_0}G_1([a])
       \le V_E^{-1}\left(\widehat G_1(0)
                   +\sum_{0\ne y\in E}|\widehat G_1(y)|\right).
\]
The last sum is the supremum of its finite subsums. With the
measure normalization of Definition~\ref{def:poisson-coordinates}
the leading constant is exactly \(V_E^{-1}\).
\end{theorem}
\begin{proof}
Fix \(\eta>0\). Choose \(m\) so that each neighborhood in the
semicontinuity hypothesis contains
\([a]+B\times m\widehat O\) for a common open ball \(B\) about zero
in \(E_\infty\). This is possible by the finite-adelic neighborhood
basis of Definition~\ref{def:finite-frame-space} in integral
coordinates. Shrink \(B\) so that these neighborhoods are disjoint
and its image is injective on the torus \(E_\infty/\sigma(mO_E)\).
Distinct \(a\)'s have distinct archimedean images. A full lattice
has a shortest nonzero length bounded away from zero by invertibility
of its real basis matrix.

In real lattice coordinates on that torus form the product kernel
\[
 K_h(u)=\prod_{j=1}^{2f}h^{-1}
              \left|\sum_{l=0}^{h-1}\exp(2\pi i l u_j)\right|^2
                        \qquad(h\ge1).
\]
It is nonnegative. Its constant Fourier coefficient is one and
all coefficients lie in \([0,1]\): the coefficient at \(k\) in
one coordinate is \(\max(0,1-|k|/h)\) by multiplying the two
finite sums. Its integral for torus probability is one, and the
integral outside any neighborhood of zero tends to zero by the
product concentration argument of Lemma~\ref{lem:gaussian-cover}.
In Lebesgue measure on \(B\) we therefore have
\(\int_B K_h\to m^{2f}V_E\), the real fundamental volume.
On \([a]+B\times m\widehat O\) its pullback under \(p_m\) is
\(K_h(z-\sigma a)\) for archimedean coordinate \(z\), by
Lemma~\ref{lem:poisson-finite-torus}. The finite-coordinate mass
is \(m^{-2f}\). Thus nonnegativity and the disjoint neighborhoods
give
\[
 V_E\sum_{a\in A_0}\max(0,G_1([a])-\eta)
          \le\liminf_{h\to\infty}\int_{\mathbb A_E}G_1(K_h\circ p_m).
\]
Each integral on the right is a finite linear combination of
\(\widehat G_1(y)\) with constant coefficient one and the other
coefficients in \([0,1]\), by Lemma~\ref{lem:poisson-finite-torus}.
It is at most the parenthesis in the assertion by the triangle
inequality. Let \(\eta\downarrow0\).

For sharpness take a half-open fundamental parallelepiped \(P\) of
\(\sigma O_E\) centered at zero, whose interior contains zero, and
\(G_1={\bf1}_{P\times\widehat O}\), \(A_0=\{0\}\).
The value at zero is one and the integral is \(V_E\).
Integration of a nontrivial additive character on \(\widehat O\)
vanishes: translating by a point of nontrivial value multiplies
the integral by that value and preserves the measure. Thus only
indices in \(\varpi^{-1}O_E\) can contribute by
Proposition~\ref{prop:additive}. For nonzero such indices, integration
of the corresponding nonconstant integer exponential over \(P\)
gives zero by the elementary one-coordinate integral. Equality
holds with these choices.
\end{proof}

\begin{corollary}\label{cor:orbital-sampling}
Let \(T\) be a finite set of finite places,
\(C=E_\infty\times\prod_{w\in T}E_w\), and give each \(O_w\)
additive mass one. Let \(l_\tau>0\),
\(m_l(z,x)=((l_\tau z_\tau)_\tau,(x_w)_{w\in T})\),
and \(U_T=\{(z,x):x_w\in O_w\text{ for }w\notin T\}\).
For \(F:C\to\R_{\ge0}\) integrable, lower semicontinuous at the
images under \(m_l\) of a finite set of diagonal \(T\)-integers
\(A_0\), define \(G_1={\bf1}_{U_T}(F\circ m_l)\).
Put
\[
 \begin{gathered}
 d_l=(V_E\prod_\tau l_\tau^2)^{-1},\qquad
 \psi_C(u)=\prod_\tau\exp(u_\tau-\bar u_\tau)\,
                       \psi_f(-u_T^0/\varpi),\\
 \xi(y)=\bigl((2\pi i\,\tau y/(l_\tau\,\tau\varpi))_\tau,
                                     (y_w)_{w\in T}\bigr).
 \end{gathered}
\]
Here \(u_T^0\) means extension by zero to the other finite places.
Then
\[
 \sum_{a\in A_0}F(m_l([a]))
 \le d_l\left(\int_C F+
         \sum_{\substack{y\ne0\\ y\ {\rm a}\ T{\rm\text{-}integer}}}
               \left|\int_C\psi_C(\xi(y)u)F(u)\dd u\right|\right).
\]
\end{corollary}
\begin{proof}
The set \(U_T\) is open and closed in the restricted product
topology of Definition~\ref{def:poisson-coordinates}. Its
finite-coordinate measure factors into mass one on
\(\prod_{w\notin T}O_w\) and the specified measure on
\(\prod_{w\in T}E_w\). This follows first on cosets of a product
of open subgroups by taking their finite quotient indices, and
then on Borel sets by the measure argument in
Lemma~\ref{lem:poisson-finite-torus}. Integration in these
coordinates and in the real coordinates follows first for
indicators of measurable rectangles, then for nonnegative
measurable functions by the same intersection argument on
bounded boxes and increasing simple-function approximation
in Lemma~\ref{lem:integral-tools}. Real coordinate scaling in
each complex factor divides the integral by \(l_\tau^2\).
Consequently \(G_1\) is integrable and meets the hypotheses of
Theorem~\ref{thm:one-sided-adelic}.

Index its transforms by \(y/\varpi\), \(y\in E\).
If \(y\) is nonintegral at some \(w\notin T\), translation by
an element in \(O_w\) changes the character nontrivially
while preserving \(G_1\). Such an element exists by the local
trace dual in Proposition~\ref{prop:additive}, applied with the
other components zero; the product trace pairs the field factors
separately. Thus the transform vanishes. For a \(T\)-integer
the factors outside \(T\) are trivial by the same trace dual.
In the remaining variables put \(u_\tau=l_\tau z_\tau\).
The archimedean exponent at \(y/\varpi\) is
\(\sum_\tau((\xi(y)u)_\tau-\overline{(\xi(y)u)_\tau})\),
since \(i v-\overline{i v}=2i\Re v\).
The finite exponent has exactly the factor in \(\psi_C\) by
Definition~\ref{def:poisson-coordinates}. The change of variables
therefore gives the asserted integrals and the factor \(d_l\)
after applying Theorem~\ref{thm:one-sided-adelic}.
\end{proof}

\section{Iteration of the one-sided orbital sieve}

\begin{lemma}\label{lem:difference-pair}
For a finite set \(A\) in a torsion-free additive group, nonnegative
weights \(b(s)\) and a nonzero group element \(h\),
\[
             \sum_{\substack{s,t\in A\\s-t=h}} b(s)b(t)
                    \le\frac14\left(\sum_{s\in A}b(s)\right)^2 .
\]
The factor \(1/4\) is attained for a suitable two-element set.
\end{lemma}
\begin{proof}
In each coset of \(\Z h\) meeting \(A\), choose one representative.
Each element of the coset is that representative plus a unique
integer multiple of \(h\), by torsion freeness. Color the elements
by the parity of that integer. A pair with difference \(h\) joins
opposite colors, and a given unordered pair occurs in at most one
orientation, since \(2h\ne0\). Its total weighted sum is bounded
by the product of the total weights \(B_0,B_1\) of the two colors.
The inequality \(4B_0B_1\le(B_0+B_1)^2\) follows from
\((B_0-B_1)^2\ge0\). Give both points of \(\{0,h\}\) weight one
to obtain equality.
\end{proof}

\begin{definition}\label{def:orbital-tower}
Let \(C,\psi_C\) be as in Corollary~\ref{cor:orbital-sampling},
\(U=C^\times\), and \(u:C\to U\) replace each zero component by
one, retaining all nonzero components. Fix an integer \(g_0\ge1\)
with \(g_0^{-1}\) a \(T\)-integer. Define
\(I(v)=((v_\tau^{-1})_\tau,(g_0^{-1}v_w^{-1})_{w\in T})\)
for \(v\in U\). Fix \(n\ge0\).
A family of kernels \(J_i:U^i\times U^i\to\C\),
\(0\le i\le n+1\), has bounds \(E_0,\beta,\Lambda_i\) if
\(E_0\ge0,\ \beta:U\to\R_{\ge0},\
\Lambda_i:U^i\to\R_{\ge0}\) and the following properties hold.
Write \((v,z)\) for a list with \(z\) appended to \(v\).
\begin{enumerate}
\item \(|J_0|\le E_0\). For \(0\le i\le n,\ v,w\in U^i\), the
function \(K(x)=J_{i+1}((v,u(x)),(w,u(x)))\) is integrable and
\[
 \int_C K=J_i(v,w),\qquad
 \left|\int_C\psi_C(\xi x)K(x)\dd x\right|
              \le\beta(\xi)\Lambda_i(v)\Lambda_i(w)\quad(\xi\in U).
\]
\item \(J_{n+1}\) is continuous. For every finite list
\(v_1,\ldots,v_b\in U^{n+1}\) with common last entry,
the matrix \((J_{n+1}(v_j,v_k))_{j,k}\) is positive semidefinite,
meaning hermitian with nonnegative quadratic form.
\item For \(1\le i\le n\), \(v,w\in U^i\) with last entries
\(z,z'\) and \(\xi=I(z)-I(z')\in U\),
\[
                  |J_i(v,w)|\le\beta(\xi)\Lambda_i(v)\Lambda_i(w).
\]
\end{enumerate}
\end{definition}

\begin{theorem}\label{thm:orbital-tower-sieve}
Use Definition~\ref{def:orbital-tower}. Let \(A\subset E^\times\)
be a finite set of \(T\)-units, meaning \(s,s^{-1}\) are
\(T\)-integers for all \(s\in A\). For \(0\le i\le n\) and
\(\tau\in G\) let \(l_{\tau i}>0\), and write
\[
 P_i(s)=((l_{\tau i}\tau s)_\tau,(s_w)_{w\in T}),\quad
 c_i=(V_E\prod_\tau l_{\tau i}^2)^{-1},\quad C_i=\prod_{a<i}c_a .
\]
For \(p=(s_0,\ldots,s_{i-1})\in A^i\) put
\(p^*=(P_0(s_0),\ldots,P_{i-1}(s_{i-1}))\) and set
\[
 M_i=\sum_{p\in A^i}\Lambda_i(p^*),\quad Q_i=M_i^2/C_i,\qquad
  H_A=\{g_0^{-1}(s^{-1}-t^{-1}):s,t\in A,\ s\ne t\}.
\]
Define the two frequency sums, using the supremum of finite subsums
for the first,
\[
\begin{aligned}
 e_i&=\sum_{\substack{y\ne0\\ y\ {\rm a}\ T{\rm\text{-}integer}}}
  \beta\!\left(\big((2\pi i_{\C}\,\tau y/(l_{\tau i}\,\tau\varpi))_\tau,
                                         (y_w)_{w\in T}\big)\right),\\
 d_i&=\sum_{y\in H_A}\beta\!\left(\big((g_0\tau y/l_{\tau i})_\tau,
                                         (y_w)_{w\in T}\big)\right).
\end{aligned}
\]
Here \(i_{\C}=\sqrt{-1}\) in the complex exponential convention of
Definition~\ref{def:foundations}. Suppose the \(e_i\) are finite.
For each \(A'\subset A\) and \(0\le i\le n\),
\[
 \begin{split}
 \frac1{C_{i+1}}\Re\sum_{s\in A'}\sum_{p,q\in A^i}
     J_{i+1}((p^*,P_i(s)),(q^*,P_i(s)))
 \ \le\ E_0+\sum_{a=0}^i Q_a e_a+\frac14\sum_{a=1}^iQ_a d_{a-1}.
 \end{split}
\]
In particular, if \(Q_a\le H\) and all the displayed frequency
sums are at most \(\epsilon\), the last upper bound is at most
\(E_0+(1+5i/4)H\epsilon\).
\end{theorem}
\begin{proof}
For \(x\in C\) let \(F_i(x)\) be the sum on \(p,q\in A^i\)
of \(J_{i+1}((p^*,u(x)),(q^*,u(x)))\). The first requirement of
Definition~\ref{def:orbital-tower} gives integrability,
\(\int F_i=\sum_{p,q}J_i(p^*,q^*)\), and the Fourier bound
\(\left|\int\psi_C(\xi x)F_i(x)\dd x\right|\le\beta(\xi) M_i^2\).
The function is real and nonnegative, and lower semicontinuous
at \(P_i(s)\) for \(s\in A\). To prove these properties, consider
any finite list of tuples of length \(j\), summing \(J_j\)
over every ordered pair of tuples. At \(j=n+1\), for tuples
sharing the final component the sum is real and nonnegative
by the quadratic form at the vector of ones. It is continuous
in the tuples by the same definition. For \(j\le n\) the sum
is the integral of the sum at \(j+1\) with one common entry
\(u(z)\) appended. Induct downwards. Nonnegativity passes
through the integral. For a convergent sequence of lists the
limiting sum is at most the lower limit of the successive sums:
at each fixed \(z\) apply the preceding stage, then apply
\(\int\liminf H_j\le\liminf\int H_j\) for nonnegative functions.
This inequality follows by applying Lemma~\ref{lem:integral-tools}
to the increasing sequence \(\inf_{j\ge a}H_j\).
This proves lower semicontinuity on the product of the unit
groups, whose topology has a countable neighborhood basis at
each point from the real and finite valuation balls.
Finally \(x\mapsto u(x)\) is continuous at \(P_i(s)\), since
all components there are nonzero. This proves the assertions
about \(F_i\).

Write \(D_i(A')\) for the unnormalized left side, including
its real part. Corollary~\ref{cor:orbital-sampling} gives
\(D_i(A')\le c_i(\int F_i+e_i M_i^2)\), because \(P_i(s)=m_l([s])\).
At \(i=0\) the integral is \(J_0\le E_0\), proving the initial
bound. At \(i\ge1\), group the terms of \(\int F_i\) according
to the two final entries \(s,t\in A\) of \(p,q\).
The contribution with \(s=t\) is \(D_{i-1}(A)\).
For \(s\ne t\), set \(h=g_0^{-1}(s^{-1}-t^{-1})\).
It is a nonzero \(T\)-integer, and
\[
 I(P_{i-1}(s))-I(P_{i-1}(t))
          =((g_0\tau h/l_{\tau,i-1})_\tau,(h_w)_{w\in T})\in U .
\]
Use the last bound in Definition~\ref{def:orbital-tower}.
For each \(s\), let \(a(s)\) be the sum of \(\Lambda_i(p^*)\)
over paths \(p\) ending at \(s\). The map \(s\mapsto g_0^{-1}s^{-1}\)
injects \(A\) into the torsion-free additive group of \(E\).
For a fixed \(h\ne0\), Lemma~\ref{lem:difference-pair} bounds
the sum of \(a(s)a(t)\) at that difference by \(M_i^2/4\).
Consequently the real off-diagonal contribution to \(\int F_i\)
is at most \(d_{i-1}M_i^2/4\). We have shown
\[
 \frac{D_i(A')}{C_{i+1}}\le\frac{D_{i-1}(A)}{C_i}
                        + Q_i(e_i+d_{i-1}/4).
\]
Induction from \(i=0\) proves the displayed estimate,
and summing the two finite constant bounds proves its corollary.
\end{proof}

\section{The orbital sieve at the arithmetic scales}

\begin{lemma}\label{lem:individual-kappas}
Use the arithmetic data of Propositions~\ref{prop:level-prime}
and~\ref{prop:unipotent-covol}. Fix \(M,d,\epsilon\) as in
Proposition~\ref{prop:thick}. Put \(h=\nu+1=m/2\),
\(l_{\tau i}=|\kappa_\tau|^{2(\nu-i)}/N\) for \(0\le i\le\nu\),
and \(\rho_\tau=|\kappa_\tau|^{2\nu+1}/(gN)\).
Uniformly on \(t\le v\le t^2\) for sufficiently large \(x\),
\[
  b^{19}\le|\kappa_\tau|\le b^{23},\quad
  N\le b^{60h},\quad V_E\le b^f,\quad k/h^9\longrightarrow\infty .
\]
The frequency scales satisfy
\[
   f^{-1}\sum_\tau\log\rho_\tau
     =(m-1)G_0-\log(gN)
     \le -m\bar w-\tfrac{25}2\ell .
\]
\end{lemma}
\begin{proof}
By the differentiated product for \(P_{N_1}\) in
Proposition~\ref{prop:different}, for \(z=\tau\zeta_{N_1}\)
one has
\[
 |P_{N_1}'(z)|=N_1
    \prod_{\substack{j\mid N_1\\j>1}}|1-z^{N_1/j}|^{\mu(j)}.
\]
Every \(z^{N_1/j}\) is primitive of order \(j\). Its distance
to one is between \(1/N_1\) and \(2\): reducing its angle to
the shorter arc gives the lower bound \(2\sin(\pi/j)\ge4/j\)
for \(j\ge2\), by the chord inequality
\(\sin u\ge 2u/\pi\) on \([0,\pi/2]\), which follows from
\(\sin''u=-\sin u\le0\). Only squarefree \(j\) contribute, and
their number is bounded in terms of \(M\) since \(N_1\) has the
same prime divisors as \(M\) by Proposition~\ref{prop:places}.
It follows that \(\big|\log|\tau\varpi|\big|
=O_{M,d}(1+\log N_1)=O_{M,d}(1+\log x)\) by the definition of
\(\varpi\) in Proposition~\ref{prop:different} and
Proposition~\ref{prop:thick}. The range \(20\ell\le\log g\le22\ell\)
in the proof of Proposition~\ref{prop:level-prime} and the formula
for \(\kappa_\tau\) in Proposition~\ref{prop:unipotent-covol} prove
the first bound. The estimate on \(N\) is the upper bound on
\(\log N\) of Proposition~\ref{prop:level-prime}. By
Lemma~\ref{lem:skew-covol}, \(V_E=(\operatorname{rd}_E/2)^f\),
and Lemma~\ref{lem:rd} gives the stated estimate. The bounds
\(h\le\sqrt x,\ k=r+1-m\ge x^5/2\) eventually by
Proposition~\ref{prop:thick} give the asserted divergence.
The product formula for \(\prod_\tau|\tau\varpi|\) in
Proposition~\ref{prop:different} gives
\(f^{-1}\sum_\tau\log|\kappa_\tau|=G_0\). Substitute this and
\(\log N\ge m\bar w+(m-1)G_0-15\ell/2\) from
Proposition~\ref{prop:level-prime}; then \(\log g\ge20\ell\)
proves the last inequality.
\end{proof}

\begin{theorem}\label{thm:quantitative-orbital-sieve}
Use the scales of Lemma~\ref{lem:individual-kappas}. In
Theorem~\ref{thm:orbital-tower-sieve} set \(n=\nu,\ g_0=g\)
and take any finite \(T\) containing every finite place dividing
\(g\). Put \(j_*=k-h^7,\ K_*=h^6\). Let \(B_\tau>0\) be constants.
For \(\xi=(\xi_\infty,\xi_T)\in U\), define
\(G_T(\xi_T),D_T(\xi_T)\) by the local products in
Theorem~\ref{thm:denominator-frequency} with \(y_w\) replaced by
\(\xi_w\). Suppose the bounds in
Definition~\ref{def:orbital-tower} satisfy
\[
 \beta(\xi)\le b^{2001h^9 f}
   \frac{G_T(\xi_T)^{j_*/2}}{D_T(\xi_T)^6}
       \prod_\tau B_\tau\,d_{j_*,6}(|\xi_\tau|^{-1}).
\]
For \(p=(s_0,\ldots,s_{i-1})\in A^i\) put
\(z_j=\prod_{a<j}s_a,\ L_{\tau j}=\prod_{a<j}l_{\tau a}\).
Assume
\[
 \Lambda_i(p^*)\le
   \prod_{j=1}^i\left(\frac{G_T(z_j)^{K_*/2}}{D_T(z_j)^6}
                       \prod_\tau d_{K_*,6}(|L_{\tau j}\tau z_j|^{-1})\right).
\]
For \(A'\subset A\), let \(D_i(A')\) be the real sum of the
kernel entries in Theorem~\ref{thm:orbital-tower-sieve} and
\(C_i\) the product there. Uniformly over these finite sets and
\(0\le i\le\nu\), at all sufficiently large \(x\) one has
\[
        \frac{D_i(A')}{C_{i+1}}\le
                  E_0+b^{2063h^9 f}\prod_\tau B_\tau\rho_\tau^{j_*}.
\]
In particular, if a positive number \(\mathfrak d\) satisfies
\[
          \mathfrak d\prod_\tau B_\tau\rho_\tau^{j_*}
                          \le b^{-7kf},
\]
the error is at most \(\mathfrak d^{-1}b^{-6kf-10}\).
\end{theorem}
\begin{proof}
The lower bound on \(k\) in Lemma~\ref{lem:individual-kappas}
ensures \(j_*\ge0\). For the transforms defining \(e_i\) in
Theorem~\ref{thm:orbital-tower-sieve}, the reciprocal of the
absolute value of the coefficient multiplying \(\tau y\) is
\(l_{\tau i}|\kappa_\tau|/g\) by
Proposition~\ref{prop:unipotent-covol}. For \(d_i\) it is
\(l_{\tau i}/g\). Both are at most \(\rho_\tau\), since
\(|\kappa_\tau|\ge1\), and \(\max(1,\rho_\tau)\le b^{46h}\)
by Lemma~\ref{lem:individual-kappas}. Every member of \(H_A\)
is a nonzero \(T\)-integer because \(g^{-1}\) and the inverses
of the entries are \(T\)-integers. Theorem~\ref{thm:denominator-frequency}
bounds both sums by
\[
     15^f b^{(2001h^9+92h)f}\prod_\tau B_\tau\rho_\tau^{j_*}.
\]
Here the stronger version with individual \(B_\tau\) in the
frequency theorem is used only to bound the damping radii
\(\max(1,\rho_\tau)\); the constants \(B_\tau\) in the present
statement simply multiply the result.

For the prefix bound, \(L_{\tau j}^{-1}\le N^j\) by
Lemma~\ref{lem:individual-kappas}. A tuple is determined by
its successive \(z_j\), \(1\le j\le i\), because
\(s_{j-1}=z_j/z_{j-1}\), with \(z_0=1\). Thus the sum of
the products bounding \(\Lambda_i\) is bounded by the product
of independent sums over all possible nonzero \(T\)-integers
at each prefix. The \(j\)-th sum is at most
\(15^f N^{j(K_*+2)f}\) by Theorem~\ref{thm:denominator-frequency}.
In particular
\[
    M_i\le15^{if}N^{(K_*+2)i(i+1)f/2}.
\]
This also holds if the prefix weight vanishes when a \(z_j\)
is nonintegral at \(T\), and carries \(G_T(z_j)^{K_*/2}\)
without a denominator when integral: there \(D_T(z_j)=1\).

By definition \(C_i^{-1}=V_E^i\prod_{a<i,\tau}l_{\tau a}^2\).
Each \(l_{\tau a}^2\le b^{92(h-1)}\).
Using \(N\le b^{60h}\), \(i(i+1)\le h^2\), \(15\le b\), and
\(V_E\le b^f\), we obtain
\[
          Q_i\le b^{(60h^9+120h^3+92h^2+3h)f}
                         \le b^{61h^9 f}
\]
eventually, since \(h\to\infty\). Apply
Theorem~\ref{thm:orbital-tower-sieve}; its prefactor
\(1+5\nu/4\) is at most \(b\) eventually by
Proposition~\ref{prop:thick}. Including \(15\le b\), the error
is at most \(b^{(2062h^9+92h+2)f}\prod_\tau B_\tau\rho_\tau^{j_*}\),
which proves the first estimate. Finally \(k/h^9\to\infty\)
and \(f\ge1\), so eventually \(2063h^9 f+10\le k f\).
This proves the last error bound under its stated numerical
hypothesis.
\end{proof}

\section{The row parameters and the factorial cancellation}

\begin{definition}\label{def:row-vandermonde}
Use \(r,k,m,h_2,w_j\) of Definition~\ref{def:heights}. Put
\(\Delta_w=2r^8,\ W_\tau=N_1\lceil w(\tau)/N_1\rceil\).
For \(0\le i<r\) the row parameters at \(\tau\) are the integers
\[
 t_{i,\tau}=
 \begin{cases}
     k/2-1-i,& 0\le i<k-1,\\
     h_2-w_{i+2-k}(\tau)=-W_\tau-(i+2-k)\Delta_w,& k-1\le i<r .
 \end{cases}
\]
Here \(r\) is a power of three, \(m\) is even for the sufficiently
large \(x\) under consideration, and \(k=r+1-m\) is even.
Define the normalized Vandermonde product and the numerical
constants
\[
 \begin{gathered}
 \mathcal D_\tau=\prod_{0\le i<l<r}
                 \frac{t_{i,\tau}-t_{l,\tau}}{l-i},\qquad
 \mathcal D=\prod_{\tau\in G}\mathcal D_\tau,\\
 U_j=\prod_{a=0}^{j-1}\frac{2\pi^{a+1}}{a!},\qquad
 C_{\rm B}=\frac{U_k}{U_r}\pi^{(m-1)k}.
 \end{gathered}
\]
\end{definition}

\begin{proposition}\label{prop:exact-row-degree}
The integers \(t_{i,\tau}\) in Definition~\ref{def:row-vandermonde}
are strictly decreasing in \(i\), and \(\mathcal D_\tau\) is
a positive integer. Its product with \(C_{\rm B}\) equals
\[
 \frac{2^{1-m}}{(k-1)!}
   \left(\frac{\Delta_w}{\pi}\right)^{m(m-1)/2}
   \left(\prod_{j=1}^{m-1}j!\right)
   \prod_{\substack{0\le i<k-1\\1\le j\le m}}
                   (W_\tau+j\Delta_w+k/2-1-i).
\]
For all sufficiently large \(x\), uniformly for
\(\tau\in G,\ t\le v\le t^2\), we have
\[
 \begin{aligned}
 m(k-1)\log w(\tau)
 &\le \log\mathcal D_\tau+\sum_{l=k-1}^{r-1}\log(l!),\\
 \log(C_{\rm B}\mathcal D_\tau)
 &\le m(k-1)\log w(\tau)
               +O\!\left(mk/x+m^2\log x\right).
 \end{aligned}
\]
\end{proposition}
\begin{proof}
The first \(k-1\) entries have successive differences one and
the final \(m\) entries have successive differences \(\Delta_w\).
Each gap between the two groups is one of the factors in the
displayed product. The smallest is
\(W_\tau+\Delta_w+1-k/2\ge w(\tau)\ge1\) by
Proposition~\ref{prop:level-prime} and \(\Delta_w=2r^8\).
Strict decrease follows.

Put \(n_i=t_{i,\tau}-t_{r-1,\tau}\ge0\). The matrix with entries
\(\binom{n_i}{j}\), \(0\le i,j<r\), is integral by the binomial
recursion, with \(\binom{u}{j}=0\) for \(0\le u<j\).
The polynomial \(\binom{X}{j}\) has leading coefficient \(1/j!\).
Change its columns to powers of \(X\) by triangular elimination.
The absolute determinant becomes
\(\prod_{i<l}(n_i-n_l)/\prod_{j<r}j!\).
Indeed the determinant of monomial evaluations in the last
variable vanishes at all previous nodes and has their determinant
as leading coefficient; induction yields the product of pairwise
differences. Since \(\prod_{i<l}(l-i)=\prod_{l<r}l!\), this
absolute determinant equals \(\mathcal D_\tau\), proving the
integer assertion.

In the numerator Vandermonde of \(t_{i,\tau}\), the upper group
contributes \(\prod_{l=0}^{k-2}l!\), canceling exactly this part
of its denominator. The lower group contributes
\(\Delta_w^{m(m-1)/2}\prod_{j=1}^{m-1}j!\).
Consequently all the remaining numerator factors are at least
one, and the \(m(k-1)\) mixed factors are at least \(w(\tau)\),
proving the first logarithmic inequality.
By \(r=k+m-1\) and summation of the arithmetic progression
\(\sum_{l=k}^{r-1}(l+1)\), the constant in the definition equals
\[
              C_{\rm B}=2^{1-m}\pi^{-m(m-1)/2}
                                \prod_{l=k}^{r-1}l! .
\]
This cancels all the remaining factorial denominators except
\((k-1)!\). It proves the exact formula.

Each mixed gap is at most \(w(\tau)+N_1+m\Delta_w+k\).
By Proposition~\ref{prop:thick} the additive error is \(O_M(x^{42})\).
Here \(w(\tau)\ge a\exp(-t^2)\ge xb\exp(-t^2)\) in the stated
range by Definitions~\ref{def:ask} and~\ref{def:heights} and
Proposition~\ref{prop:level-prime}. Since \(t^2=o(\ell)\) and
\(b=\exp(\ell)\), the relative error is at most \(1/x\) eventually
with \(d,M,\epsilon\) fixed. The inequality \(\log(1+u)\le u\)
for \(u\ge0\), by integration of \((1+u)^{-1}\), bounds the
logarithm of the mixed product by \(m(k-1)(\log w(\tau)+1/x)\).
Every difference within the lower group is at most \(m\Delta_w\)
and \(\log(m\Delta_w)=O(\log x)\) by Proposition~\ref{prop:thick}.
The factors \(2^{1-m},\pi^{-m(m-1)/2},1/(k-1)!\) are at most
one. This proves the second logarithmic bound.
\end{proof}

\begin{theorem}\label{thm:orbital-degree-budget}
Use the data of Theorem~\ref{thm:quantitative-orbital-sieve} and
Definition~\ref{def:row-vandermonde}. For every fixed
\(0<\eta<25/2\), uniformly for \(t\le v\le t^2\),
\[
          b^{2063h^9 f}\mathcal D
              \prod_{\tau\in G} C_{\rm B}\rho_\tau^{\,k-h^7}
                       \le b^{-(25/2-\eta)kf}
\]
for all sufficiently large \(x\), where \(h=m/2\) as in
Lemma~\ref{lem:individual-kappas}. In particular, whenever the
two bounds on \(\beta,\Lambda_i\) in
Theorem~\ref{thm:quantitative-orbital-sieve} hold with
\(B_\tau=C_{\rm B}\), its conclusion takes the form
\[
        \frac{D_i(A')}{C_{i+1}}
                 \le E_0+\mathcal D^{-1}b^{-(25/2-\eta)kf}
                      \quad (A'\subset A,\ 0\le i\le\nu).
\]
\end{theorem}
\begin{proof}
Sum the last inequality of Proposition~\ref{prop:exact-row-degree}
and the estimate on the frequency scales of
Lemma~\ref{lem:individual-kappas}, the latter multiplied by
\(f(k-h^7)\). After adding \(2063h^9 f\ell\) and dividing by
\(f\), the logarithm of the left side is at most
\[
 -\tfrac{25}2k\ell+2063h^9\ell+\tfrac{25}2h^7\ell
       +m(h^7-1)\bar w+O(mk/x+m^2\log x).
\]
Here \(\bar w=f^{-1}\sum_\tau\log w(\tau)\) by
Definition~\ref{def:ask}. We have \(0\le\bar w\le2\ell\)
by the bounds on the weights in Proposition~\ref{prop:level-prime}.
Since \(m=2h\), the terms involving \(h^9,h^7,mh^7\) divided by
\(k\) tend to zero by Lemma~\ref{lem:individual-kappas}.
Also \(m/(x\ell)\to0,\ m^2\log x/(k\ell)\to0\) by
Proposition~\ref{prop:thick}. These comparisons are uniform in
the window. The terms after the leading term are therefore at
most \(\eta k\ell\) eventually. Exponentiating proves the
first assertion. The second follows by inserting this estimate
in Theorem~\ref{thm:quantitative-orbital-sieve}.
\end{proof}

\section{Finite norm congruences for the column lattice}

\begin{lemma}\label{lem:real-place-rings}
Let \(F=E^{\mathfrak c}\) and \(O_F=O_E\cap F\) as in
Lemma~\ref{lem:skew-covol}. For each prime \(p\),
\(O_F\otimes\Z_p\) is a product of at most \(f\) complete discrete
valuation rings \(A\). For each such factor, write \(Q=|A/\pi A|\)
for the residue cardinality, where \(\pi\) is a uniformizer.
Then \(Q\ge p\). The corresponding algebra \(B=O_E\otimes_{O_F}A\)
is either \(A\times A\) with its factors interchanged, or a complete
discrete valuation ring with uniformizer \(\pi\) and residue field
of cardinality \(Q^2\). In the latter case the involution induces
\(z\mapsto z^Q\) on residues.
\end{lemma}
\begin{proof}
By Lemma~\ref{lem:skew-covol}, \(O_E=O_F\oplus\xi O_F\) with
\(\xi\bar\xi=1,\ (\xi-\bar\xi)^2=-1\).
This decomposition holds over \(\Z_p\), where subtracting the
conjugate of \(a+b\xi\) still detects \(b\). Thus the fixed ring
of \(O_E\otimes\Z_p\) is \(O_F\otimes\Z_p\).
Theorem~\ref{thm:integer-ring} and Lemma~\ref{lem:integer-extension}
decompose \(O_E\otimes\Z_p\) into complete discrete valuation rings
with finite residue fields. Conjugation permutes their identity
idempotents, since the only idempotents in a domain are zero and
one. On an exchanged pair the invariant ring projects isomorphically
onto either factor. On a fixed factor it is the valuation ring of
the fixed subfield: the condition of lying in the original
valuation ring is nonnegative valuation, and the restricted value
group is a nonzero subgroup of \(\Z\), since it contains the value
of \(p\). A fixed element of least positive valuation is a
uniformizer there. The fixed ring is closed under compatible
limits, and its residue field embeds in the finite residue field
of the original ring. This proves the decomposition asserted
for \(O_F\otimes\Z_p\). Its rank over \(\Z_p\) is \(f\) by
Lemma~\ref{lem:skew-covol}. Modulo \(p\) each factor contributes
at least one dimension, and each residue field contains \(\mathbf F_p\).

Over any such \(A\), the algebra \(B\) is given by the monic
polynomial \(X^2-(\xi+\bar\xi)X+1\), by the same basis. Its
discriminant is \(-1\), so it has distinct roots over the algebraic
closure of the residue field, including in characteristic two.
Apply Lemma~\ref{lem:integer-extension}. In the split case
conjugation exchanges the roots. In the field case the two residue
roots are interchanged by \(z\mapsto z^Q\): this map fixes every
element of the field of size \(Q\) by cyclicity in the proof of
Lemma~\ref{lem:count-labels}, and its fixed points satisfy the
degree-\(Q\) equation \(z^Q=z\). It cannot fix the quadratic
generator.
\end{proof}

\begin{definition}\label{def:norm-density}
At a factor \(A\) of Lemma~\ref{lem:real-place-rings}, set
\(H(z)=\sum_{i=1}^k \bar z_i z_i\) on \(B^k\). For \(t\in A\)
and \(n\ge1\) put
\[
 A_n(t)=Q^{-(2k-1)n}\#
 \{z\in(B/\pi^nB)^k:H(z)\equiv t\bmod\pi^n A\},\qquad A_0(t)=1 .
\]
In this section \(k\ge1\). Write \(\sigma=1\) in the split case and
\(\sigma=(-1)^k\) in the field case. On \(A,B^k\) take the
translation-invariant probability measures determined by uniform
counting on finite quotients. Thus \(Q^{-n}A_n(t)\) is the
probability of the indicated congruence.
\end{definition}

\begin{proposition}\label{prop:residue-norm-count}
In Definition~\ref{def:norm-density} one has
\[
 A_1(t)=
 \begin{cases}
  1-\sigma Q^{-k},&t\notin\pi A,\\
  1+\sigma(Q-1)Q^{-k},&t\in\pi A .
 \end{cases}
\]
The reduction of \(H\) has a surjective differential at every
residue vector other than the zero vector.
\end{proposition}
\begin{proof}
Write \(C_k(u)\) for the unnormalized count over the residue field
at \(u\). Begin with \(C_0(u)={\bf1}_{u=0}\). In the field case the
one-coordinate norm is \(z^{Q+1}\) by Lemma~\ref{lem:real-place-rings}.
The cyclic group of order \(Q^2-1\) has kernel of size \(Q+1\)
under this power. Its image is the full group of nonzero fixed
elements by its size \(Q-1\). Only zero has zero norm. Summing the
one-coordinate counts at \(u-H(z)\) gives
\(C_{k+1}(u)=(Q+1)Q^{2k}-Q C_k(u)\).
In the split case the norm is \(xy\). There are \(Q-1\) pairs
at \(u\ne0\) and \(2Q-1\) at zero, giving
\(C_{k+1}(u)=(Q-1)Q^{2k}+Q C_k(u)\).
Substitution in these recurrences, starting at \(k=1\) with the
counts just given, proves the formulas after division by \(Q^{2k-1}\).

The differential is \(v\mapsto\sum_i(\bar z_i v_i+\bar v_i z_i)\).
For an inert nonzero coordinate its image is the image of the
residue trace. This trace has nonzero image: \(s+s^Q\) cannot
vanish on all \(Q^2\) elements, being a nonzero polynomial of
degree \(Q\). The trace is linear over the fixed field and hence
onto it. In the split case the derivative of a product \(xy\)
is \(x\,dy+y\,dx\), which is onto whenever either entry is nonzero.
\end{proof}

\begin{theorem}\label{thm:exact-local-density}
Set \(u=1-\sigma Q^{-k}\) and
\(v=1+\sigma(Q-1)Q^{-k}-Q^{1-2k}\). For \(n\ge2\) one has
\[
 A_n(t)=
 \begin{cases}
 u,& t\notin\pi A,\\
 v,&t\in\pi A\setminus\pi^2 A,\\
 v+Q^{2-2k}A_{n-2}(t/\pi^2),&t\in\pi^2 A .
 \end{cases}
\]
For \(k\ge2\), \(|A_n(t)-1|\le 2Q^{1-k}\) for every \(n,t\),
and \(\alpha_A(t)=\lim_{n\to\infty}A_n(t)\) exists for every \(t\);
at \(t=0\) it equals \(v/(1-Q^{2-2k})\).
In all ranks and for \(t\ne0\) the sequence is constant as soon
as \(n>v_A(t)\), where \(v_A(\pi)=1\). In rank one at \(t=0\)
the split values are \(A_n(0)=1+n(1-Q^{-1})\); in the field
case the values are \(1\) at even depths and \(Q^{-1}\) at odd depths.
\end{theorem}
\begin{proof}
Above any nonzero residue solution, to lift a solution modulo
\(\pi^j\) to one modulo \(\pi^{j+1}\) we add \(\pi^j z'\) to a
representative. The linear term modulo \(\pi^{j+1}\) is
\(\pi^j\,dH(z')\), with the differential at that residue solution.
The quadratic term is divisible by \(\pi^{2j}\).
Proposition~\ref{prop:residue-norm-count} gives exactly
\(Q^{2k-1}\) choices at each step. This part of the normalized
count is thus \(u\) at a unit value and \(v\) at a nonunit, since
at zero one residue solution has been removed. Above the zero
vector write \(z=\pi z'\). Its norm is \(\pi^2H(z')\).
If \(\pi^2\mid t\), the vector \(z'\) is specified modulo
\(\pi^{n-1}\) and its congruence only uses it modulo
\(\pi^{n-2}\). This extra digit has \(Q^{2k}\) choices, giving
the normalized contribution \(Q^{2-2k}A_{n-2}(t/\pi^2)\);
for \(n=2\) this interpretation gives \(A_0=1\).

For \(k\ge2\) put \(c=Q^{1-k}\le1/2\). One has \(|u-1|\le c\) and
\(|v-1|\le(1-1/Q)c+Q^{-k}c\le c\).
Proposition~\ref{prop:residue-norm-count} handles \(n=1\).
Inductively \(c+c^2(1+2c)\le2c\), proving the bound in the last
case as well. Iterating the recurrence at a nonzero \(t\) ends at
a unit or at a value of valuation one, at a level at least one
or two, respectively, when \(n>v_A(t)\).
At \(t=0\) both parities converge to the sum of the same geometric
series when \(k\ge2\), since \(Q^{2-2k}<1\).
In rank one the coefficient \(Q^{2-2k}\) is one.
At zero the increment in two steps is \(2(1-Q^{-1})\) in the
split case and zero in the field case. The initial values at
depths zero and one now give the two formulas asserted.
\end{proof}

\section{The finite Fourier integral for one represented element}

\begin{lemma}\label{lem:norm-gauss}
Use Definition~\ref{def:norm-density}, and let \(\chi\) be an additive
character of \(A\) of conductor \(n\ge1\), meaning it is trivial on
\(\pi^n A\) and nontrivial on \(\pi^{n-1}A\). For uniform probability
on \(B^k\),
\[
    G(\chi)=\int_{B^k}\chi(H(z))\dd z,\qquad
                  |G(\chi)|=Q^{-kn}.
\]
For the trivial character \(G(1)=1\). There are
\(Q^n-Q^{n-1}\) characters of conductor \(n\). For any \(t\in A\)
the sum of \(G(\chi)\chi(-t)\) over characters of conductor at most
\(n\), including the trivial character, is \(A_n(t)\).
\end{lemma}
\begin{proof}
For a finite abelian group, characters on a subgroup extend across
a cyclic quotient generated by \(a\) of order \(d\) modulo the
subgroup: choose any \(d\)-th root of the old value at \(da\) as
the value at \(a\). This gives exactly \(d\) extensions. Beginning
with the trivial subgroup and adjoining one element at a time shows
the dual of a finite abelian group has the group's order. A
character nontrivial at a given nonzero \(b\) is obtained by extending
one from \(\langle b\rangle\). Multiplication by this character
shows the sum of all characters at \(b\) is zero. Applied to
\(A/\pi^n A\), this counts the characters and writes the indicator
of a residue class as \(Q^{-n}\sum_\chi\chi(H(z)-t)\).
Definition~\ref{def:norm-density} gives the last assertion.

In the squared absolute value of \(G(\chi)\) write the two
variables as \(z+y,z\). The inner average in \(z\) is the average
of the character with argument
\(\sum_i\operatorname{Tr}_{B/A}(\bar y_i z_i)\), times
\(\chi(H(y))\). The trace here is \(b+\bar b\).
The bilinear pairing \((y,z)\mapsto\operatorname{Tr}(\bar y z)\)
on \(B\) has determinant a unit. Indeed in \(1,\xi\) the matrix has
entries \(2,\xi+\bar\xi,\xi+\bar\xi,2\) and determinant
\(-(\xi-\bar\xi)^2=1\). If any coefficient of \(y\) in this basis
has valuation less than \(n\), the image ideal of the linear form
on \(B^k\) therefore contains \(\pi^{n-1}A\). The average of
its character vanishes by translation by an argument with character
different from one. For \(y\in\pi^nB^k\), both characters are one.
The measure of the latter set is \(Q^{-2kn}\), giving the claimed
absolute value.
\end{proof}

\begin{theorem}\label{thm:finite-norm-product}
Let \(k\ge3\). Index the factors of Lemma~\ref{lem:real-place-rings}
by \(v\), with residues \(Q_v\), and write \(A_v,B_v\) for their
rings. For any set \(T\) of such factors and \(t_v\in A_v\) at
each \(v\notin T\), the product
\(P_T(t)=\prod_{v\notin T}\alpha_{A_v}(t_v)\) converges to a
positive number. Its logarithm, and the logarithm of any product
with the factors replaced by \(A_{n_v}(t_v)\) at arbitrary
levels \(n_v\ge0\), have absolute value at most \(12 f\,2^{1-k}\).
At \(t_v=1\) for all \(v\notin T\) one has the exact product
\[
           P_T(1)=\prod_{v\notin T}(1-\sigma_v Q_v^{-k}),
       \qquad |\log P_T(1)|\le4f\,2^{-k}.
\]
For \(T=\varnothing\) and \(t=(t_v)\in\prod_v A_v\), put
\(F_f=\bigcup_{m\ge1}m^{-1}\prod_v A_v\).
Give \(\prod_v B_v^k\) probability measure. Normalize additive
measure on \(F_f\) to give the trace dual of \(\prod_v A_v\)
mass one. With the character of Proposition~\ref{prop:additive},
the absolutely convergent integral satisfies
\[
 \int_{F_f}\left(\int_{\prod_v B_v^k}
       \psi_f\big(z(H(y)-t)/2\big)\dd y\right)\dd z
                   =P_{\varnothing}(t).
\]
\end{theorem}
\begin{proof}
Theorem~\ref{thm:exact-local-density} bounds the distance of each
factor from one by \(2Q_v^{1-k}\le1/2\). Integrating \(1/u\)
between it and one bounds its logarithm by \(4Q_v^{1-k}\).
There are at most \(f\) factors above each \(p\) and \(Q_v\ge p\)
by Lemma~\ref{lem:real-place-rings}.
For \(s>1\) monotonicity gives
\(\sum_{j\ge2} j^{-s}\le2^{-s}(1+2/(s-1))\).
With \(s=k-1\) this proves the uniform bound and convergence.
At one the exact value is Theorem~\ref{thm:exact-local-density},
and the same argument uses \(2Q_v^{-k}\) and \(s=k\).

For \(z,w\in F_f\) the character \(\psi_f(zw/2)\) equals the
product of the rational finite characters at
\(\operatorname{Tr}_{F\otimes\Q_p/\Q_p}(z_pw_p)\): the trace
from \(E\) of a scalar in \(F\) is twice its trace from \(F\).
Thus no invertibility of two in any \(A_v\) is required.
Write \(D^\vee\) for the annihilator of \(\widehat A=\prod_v A_v\).
It is also the intersection with \(F_f\) of the trace dual of
\(\widehat O_E=\prod_v B_v\) under the full \(E\)-trace, because
\(\operatorname{Tr}_{E/F}(\xi/i)=1\) and
\(\operatorname{Tr}_{E/F}(\widehat O_E)=\widehat A\).
In an integral basis of \(O_F\) the trace matrix is nonsingular:
its real form on nonzero vectors is \(\sum_\tau x_\tau^2>0\)
by the embedding calculation in Lemma~\ref{lem:skew-covol}.
It identifies \(F_f/D^\vee\) with the additive characters of
\(\widehat A\) factoring through some \(\widehat A/m\widehat A\).
Indeed in the resulting dual basis the classes killed by \(m\)
are \(m^{-1}\widehat{\Z}^{\,f}/\widehat{\Z}^{\,f}\), giving
\(m^f\) distinct characters by the rational character calculation
of Proposition~\ref{prop:additive}; this exhausts the characters
by Lemma~\ref{lem:norm-gauss}. Here
\(\widehat{\Z}=\prod_p\Z_p\). The finite quotients decompose over
the factors \(v\), so each character tuple is trivial except at
finitely many \(v\), and every such tuple occurs by choosing \(m\)
divisible by its finitely many conductors in the local rings.

The inner integral is constant on each coset of \(D^\vee\)
and equals \(\prod_v G(\chi_v)\chi_v(-t_v)\) there by the
finite-quotient calculation of Lemma~\ref{lem:norm-gauss}.
Its absolute integral is finite. The sum of the absolute local
terms at positive conductor is at most
\(\sum_{n\ge1}Q_v^{(1-k)n}\le(4/3)Q_v^{1-k}\).
The sum of these bounds over \(v\) is finite by the preceding
series estimate. Expanding finite products and bounding
\(\prod(1+c_v)\le\exp(\sum c_v)\) justifies the absolutely
convergent expansion over all tuples.
Each local character sum has value \(\alpha_{A_v}(t_v)\) by
Lemma~\ref{lem:norm-gauss} and Theorem~\ref{thm:exact-local-density}.
This proves the integral formula.
\end{proof}

\begin{corollary}\label{cor:off-places-norm}
Use the arithmetic data of Proposition~\ref{prop:level-prime} and
a finite set \(T\) of factors as above. Give every \(B_v^k\) and
\(A_v\) mass one. Off \(T\) impose integrality on vectors. At
finitely many \(v\notin T\) impose also
\(H(z_v)\in t_v+\pi_v^{n_v}A_v\), with \(n_v\ge0\) and
\(t_v\in\operatorname{Frac}(A_v)\). Multiply the resulting measure
by \(\prod_v Q_v^{n_v}\), taking \(n_v=0\) at unused places.
For every fixed \(D>0\) and all sufficiently large \(x\) the result
is at most \(1+b^{-D}\), uniformly in \(T\), the values and the
exponents. Integrating an additional nonnegative test depending
only on \(T\) multiplies its integral by that same factor after
normalization off \(T\).
\end{corollary}
\begin{proof}
A nonintegral target at a constrained place makes that congruence
empty. Otherwise independence of finite coordinate sets and
Definition~\ref{def:norm-density} give exactly
\(\prod_{v\notin T} A_{n_v}(t_v)\) after normalization.
The identity with a test follows first for indicators by
multiplication of measures and then by increasing nonnegative
step functions, as in Lemma~\ref{lem:integral-tools}.
By Proposition~\ref{prop:level-prime} and
Lemma~\ref{lem:individual-kappas}, \(f\le x^{12}\) eventually,
\(k/h^9\to\infty,\ h=m/2\ge\sqrt x/2\).
Since \(\ell=x^\epsilon\) with \(0<\epsilon<1/10\) in
Proposition~\ref{prop:thick}, \(12f2^{1-k}=o(\exp(-D\ell))\).
Theorem~\ref{thm:finite-norm-product} and
\(\exp u\le1+2u\) for sufficiently small \(u\ge0\), by the
exponential series, prove the estimate.
\end{proof}

\section{Unit vectors and the masses of the definite lattices}

\begin{lemma}\label{lem:norm-one-isometry}
On \(E^j\) write \(h(u,v)=\sum_i\bar u_i v_i,\ H(u)=h(u,u)\).
If \(x,y\in E^j\) have the same nonzero norm, a rational isometry
sends \(x\) to \(y\). At every factor \(B\) of
Lemma~\ref{lem:real-place-rings}, the integral isometry group of
\(B^j\) is transitive on vectors with \(H(x)=1\).
\end{lemma}
\begin{proof}
For \(x\ne y\) put \(u=x-y,\ a=h(u,x)\). Then
\(a+\bar a=H(u)\ne0\), since the form is positive definite in every
complex embedding. In particular \(a\ne0\). The map
\(z\mapsto z-u h(u,z)/a\) sends \(x\) to \(y\).
Its two cross terms and its quadratic term in \(h\) cancel because
\(1/a+1/\bar a=H(u)/(a\bar a)\). Hence it is an isometry, since a
linear map preserving a nondegenerate form in finite dimension is
injective and onto. When \(x=y\) use the identity.

For split \(B=A\times A\), a unit vector is a column and covector
\((v,\lambda)\) satisfying \(\lambda(v)=1\) over \(A\).
The kernel of \(\lambda\) has a free basis: use an invertible
coefficient of \(\lambda\) to solve for one coordinate. Adjoining
\(v\) to that basis gives an invertible matrix, and the action by
a matrix on columns and its inverse transpose on covectors sends
a standard pair to this pair.

In the field case every unit of \(A\) is a norm of a unit of \(B\).
The residue assertion was proved in
Proposition~\ref{prop:residue-norm-count}. Given a solution modulo
\(\pi^n\), correcting it by a factor \(1+\pi^n w\) corrects its
norm at the next digit by \(\operatorname{Tr}(w)\) times the old
unit residue. This trace is onto by the same proposition.
Compatible corrections converge by
Lemma~\ref{lem:real-place-rings}. A norm-one \(x\) involves a unit
coordinate, so adjoining suitable coordinate vectors gives a basis.
Subtract \(x h(x,z)\) from those other basis elements. They form
a basis of \(x^\perp\), with unimodular hermitian Gram matrix
since the full basis determinant is a unit and the first block
is \(1\). Any positive-rank free lattice with unimodular hermitian
Gram matrix over \(B\) has a vector of unit norm. Otherwise all
residue norms would vanish and expansion at \(u+cw\) would give
\(\operatorname{Tr}(c\,h(u,w))=0\) for all residues \(c,u,w\).
Trace surjectivity would force the entire residue matrix to vanish.
Scale a lift of a vector of unit norm to have norm one by the unit
norm surjectivity just proved. Continuing on its complement produces
an integral orthonormal basis starting with the given vector.
\end{proof}

\begin{definition}\label{def:marked-mass}
For \(j\ge1\) let \(\mathcal L_j\) consist of the rational isometry
classes of full \(O_E\)-lattices in \(E^j\) locally isometric at
every finite place to the standard \(B_v^j\).
Write \(U(L)\) for the rational isometries preserving \(L\).
Set \(\mathfrak m_j=\sum_{[L]\in\mathcal L_j}|U(L)|^{-1}\),
and give the zero lattice mass \(\mathfrak m_0=1\).
For \(t\in F\) totally positive put
\(r_L(t)=\#\{x\in L:H(x)=t\}\).
\end{definition}

\begin{theorem}\label{thm:marked-mass}
The classes and stabilizers in Definition~\ref{def:marked-mass}
are finite, each \(r_L(t)\) is finite, and for all \(j\ge1\)
\[
              \sum_{[L]\in\mathcal L_j}\frac{r_L(1)}{|U(L)|}
                       =\mathfrak m_{j-1}.
\]
Passing from the full integral level to a subgroup of finite index
\(I\) multiplies both the class mass and the weighted sum of
\(r_L(t)\) by \(I\).
\end{theorem}
\begin{proof}
Use Lemma~\ref{lem:trace-finite} in rank \(j\). Its lattice
description applies to every class here: a lattice and the standard
lattice lie between two common scalar multiples of the latter.
Away from the divisors of that scalar the identity can be used as
local isometry. A vector of norm \(t\) satisfies
\(\sum_{\tau,i}|\tau(x_i)|^2=\sum_\tau\tau(t)\), so it lies in a
bounded subset of a real lattice and there are only finitely many
possibilities.

Mark a vector \(x\) of norm one in \(L\). Its pairing with any
\(y\in L\) is integral at every finite place by the local
isometries. Thus \(h(x,y)\in O_E\) by
Theorem~\ref{thm:integer-ring}, and
\(L=O_E x\perp (L\cap x^\perp)\) by subtracting \(xh(x,y)\).
Use Lemma~\ref{lem:norm-one-isometry} to take \(x\) to the last
standard vector rationally. Its orthogonal complement lattice is
locally standard by the integral assertion of that lemma. Its
rational class is independent of the chosen isometry since two
choices differ by a rational isometry fixing the last vector.
Conversely adjoining the standard last vector to any lattice of rank
\(j-1\) constructs a marked lattice, inverse to this assignment up
to isometry. The stabilizer of a marked pair is precisely the
isometry group of its complement, by restriction and extension by
the identity on the last line. In a finite group an orbit of
size \(a\) has stabilizer of order \(|U(L)|/a\), by the bijection
with its cosets. Summing orbit contributions \(a/|U(L)|\) proves
the formula.

For a fixed full-level class the additional level consists of a
finite coset set of size \(I\), on which \(U(L)\) acts. The
stabilizer of a point is exactly the isometry group preserving that
choice of level. Applying the same orbit identity to this set with
the constant weights \(1\) and \(r_L(t)\) proves the last assertion.
\end{proof}

\begin{proposition}\label{prop:norm-arch-factor}
For Lebesgue measure on \(\C^j\) the image under
\(z\mapsto\sum |z_i|^2\) has density
\(\pi^j u^{j-1}/(j-1)!\) on \(u>0\).
For \(j\ge3\) define
\[
 B_j(t)=\frac{|\Delta_E|^{1/4}}{V_E^j}
    \prod_\tau\frac{\pi^j\,\tau(t)^{j-1}}{(j-1)!}\,
                      P_{\varnothing}(t)
\]
at totally positive \(t\in O_F\), using
Theorem~\ref{thm:finite-norm-product}.
If the normalized class average of \(r_L(1)\) equals \(B_j(1)\),
then the consecutive masses satisfy the exact formula
\[
 \frac{\mathfrak m_j}{\mathfrak m_{j-1}}
   =\left(\frac{{\rm rd}_E^{\,j-1/2}(j-1)!}{(2\pi)^j}\right)^f
        \prod_v(1-\sigma_v Q_v^{-j})^{-1}.
\]
The absolute error in its logarithm after omitting the product
is at most \(4f\,2^{-j}\).
\end{proposition}
\begin{proof}
For one complex coordinate the disk of squared radius \(u\) has
area \(\pi u\), by integrating in polar coordinates. Summing
independent nonnegative squared radii convolves their densities:
integration from \(0\) to \(u\) sends
\(\pi(\pi^i v^{i-1}/(i-1)!)\) to \(\pi^{i+1}u^i/i!\).
Repeated integration proves the density. In the stated implication
the normalized class average is
\(\mathfrak m_{j-1}/\mathfrak m_j\) by
Theorem~\ref{thm:marked-mass}. Substitute
\(V_E=2^{-f}|\Delta_E|^{1/2}\) from Lemma~\ref{lem:skew-covol}
and \({\rm rd}_E=|\Delta_E|^{1/(2f)}\) from Lemma~\ref{lem:rd}.
Taking the reciprocal yields the displayed formula, and
Theorem~\ref{thm:finite-norm-product} bounds the logarithm.
\end{proof}

\section{The rank-one Hecke identity at one split place}

\begin{definition}\label{def:rank-one-hecke}
Let \(A\) be a complete discrete valuation ring of the kind in
Lemma~\ref{lem:real-place-rings}, \(L=\operatorname{Frac}(A)\),
\(Q=|A/\pi A|\), and \(|\pi|=Q^{-1}\).
Fix \(j\ge2\). On compactly supported locally constant complex
functions on \(M_{j,2}(L)\) put
\[
 U_b\phi(X)=|\det b|^{-1}\phi(b^{-1}X),\qquad
 V_c\phi(X)=|\det c|^{j/2}\phi(Xc)
\]
for \(b\in{\rm GL}_j(L),\ c\in{\rm GL}_2(L)\), and put
\(\phi_0={\bf1}_{M_{j,2}(A)}\). These commuting actions preserve
the integral of \(|\phi|^2\) for matrix measure giving
\(M_{j,2}(A)\) mass one. Indeed a diagonal multiplication has
modulus given by the product of its entry moduli by the coset
counts, and an invertible matrix is a product of integral basis
changes and a diagonal matrix by the elimination in
Lemma~\ref{lem:class-covol}. There are two columns and \(j\) rows.
\end{definition}

\begin{lemma}\label{lem:split-hecke-eigen}
Let \(\mathcal D\) be the set of index-\(Q\) inverse images in
\(A^2\) of hyperplanes of its residue space, using row coordinates.
Choose \(c_D\) with \(A^2c_D=D,\ \det c_D=\pi\).
For every linear functional \(\lambda\) invariant under all \(U_b\)
and every \(g\in{\rm GL}_2(L)\), define \(F(g)=\lambda(V_g\phi_0)\).
Then
\[
       \sum_{D\in\mathcal D} F(gc_D^{-1})
                =\big(Q^{j/2}+Q^{\,1-j/2}\big)F(g).
\]
\end{lemma}
\begin{proof}
For rank \(s\) let \([s]_Q=(Q^s-1)/(Q-1)\), with \([0]_Q=0\).
The hyperplanes in residue dimension \(s\) are the kernels of the
nonzero linear forms modulo scalar, so there are \([s]_Q\).
Their inverse images have bases of determinant \(\pi\) up to a
unit. Lift a residue basis putting the hyperplane in standard
position (the lifted determinant is a unit), and multiply a
complementary basis vector by \(\pi\). Adjust one basis element
by a unit to prescribe determinant \(\pi\). This constructs the
\(c_D\), and constructs \(b_H\) for index-\(Q\) column lattices
\(b_H A^j\) indexed by hyperplanes \(H\) in residue dimension \(j\).

For an integral matrix \(X\) of residue rank \(r\), the number of
these column lattices containing its two columns is \([j-r]_Q\).
The number of the row lattices \(D\) containing all its rows is
\([2-r]_Q\). Indeed row rank equals column rank by elimination
over the residue field. From
\([j-r]_Q=[j-2]_Q+Q^{j-2}[2-r]_Q\) we get the identity of functions
\[
 Q^{-1}\sum_H U_{b_H}\phi_0
    =[j-2]_Q\phi_0+Q^{j-2-j/2}\sum_D V_{c_D^{-1}}\phi_0 .
\]
Outside the integral matrices every indicator occurring here is
zero. Apply \(\lambda V_g\), which is left invariant by
Definition~\ref{def:rank-one-hecke}. The left side gives
\(Q^{-1}[j]_Q F(g)\). Now
\(Q^{-1}[j]_Q-[j-2]_Q=Q^{j-2}+Q^{-1}\), giving the eigenvalue.
\end{proof}

\begin{theorem}\label{thm:bounded-split-hecke}
If \(j>2\) and \(F\) from Lemma~\ref{lem:split-hecke-eigen}
is bounded on \({\rm SL}_2(L)\), then
\(F(g)=0\) on \({\rm SL}_2(L)\). In particular \(\lambda(\phi_0)=0\).
\end{theorem}
\begin{proof}
For \(a\in L^\times\) one has \(V_{a I_2}=U_{a^{-1}I_j}\) by
Definition~\ref{def:rank-one-hecke}, so \(F(ag)=F(g)\).
Apply Lemma~\ref{lem:split-hecke-eigen} twice and set
\(\Lambda=Q^{j/2}+Q^{1-j/2}\). For \(g\in{\rm SL}_2(L)\),
\[
 \Lambda^2 F(g)=\sum_{D,D'}F(g\,\pi c_D^{-1}c_{D'}^{-1}).
\]
Each argument on the right has determinant one. There are
\((Q+1)^2\) terms. With \(u=Q^{j/2-1}>1\) one has
\(\Lambda-(Q+1)=(u-1)(Q-u^{-1})>0\).
Taking absolute values and the supremum over \(g\) yields
\(\Lambda^2\sup|F|\le(Q+1)^2\sup|F|\). Its finiteness forces
it to be zero.
\end{proof}

\section{Reduction for the real subfield in rank one}

\begin{definition}\label{def:real-adeles}
For \(F=E^{\mathfrak c}\) set \(\widehat A=\prod_v A_v\),
\(F_\infty=\R^\Phi\) and \(\mathbb A_F=F_\infty\times F_f\),
using the rings of Lemma~\ref{lem:real-place-rings} and the finite
adeles of Theorem~\ref{thm:finite-norm-product}. Use the usual real
absolute values and \(|z|_v=Q_v^{-v_A(z)}\) at finite places.
An idele is a tuple of nonzero components which are units at almost
all finite places. Set \(s(a)=\prod_\tau |a_\tau|\prod_v |a_v|_v\)
on ideles. For \(r>0\) write \(d(r)\) for the idele equal to \(r\)
at all infinite places and \(1\) elsewhere.
For \(a\) an idele and \(u\in\mathbb A_F\), put
\(m(a)=\operatorname{diag}(a,a^{-1})\) and
\(n(u)=\left(\begin{smallmatrix}1&u\\0&1\end{smallmatrix}\right)\).
The group \({\rm SL}_2(\mathbb A_F)\) consists of the tuples of
determinant-one matrices integral at almost all finite places.
Neighborhoods have unrestricted integral components away from
finitely many places. The same convention uses unit components for
the group of ideles.
\end{definition}

\begin{lemma}\label{lem:real-volume-compact}
The real image of \(O_F\) has covolume \(V_F=|\Delta_E|^{1/4}\).
For \(n\ge1\) and an invertible adelic \(n\)-by-\(n\) matrix \(g\)
with integral inverse and entries at almost all finite places, the
lattice
\[
 \{xg_\infty:x\in F^n,\ xg_f\in\widehat A^n\}
              \ \subset\  (F_\infty)^n
\]
has covolume \(V_F^n s(\det g)\). In particular \(s(a)=1\)
for every \(a\in F^\times\). The quotient of \(\mathbb A_F\)
by the additive group \(F\) has a compact set of representatives
in the sense that every class meets the set.
There exists a compact set of ideles meeting every \(F^\times\)-orbit
of ideles with \(s(a)=1\).
\end{lemma}
\begin{proof}
Lemma~\ref{lem:skew-covol} gives the determinant \(V_F\) in an
integral basis of \(O_F\). For the finite congruences repeat the
integer-basis count in Lemma~\ref{lem:class-covol}: choose one
integer clearing all finite denominators of \(g,g^{-1}\). The
local quotient indices multiply by the Chinese remainder argument
in that lemma. An integral elementary row or column operation over
\(A_v\) preserves the index, and a diagonal entry \(\pi_v^e\)
has index \(Q_v^e\) when \(e\ge0\), with the inverse index for
\(e<0\). The preimage by \(g_v\) thus contributes
\(|\det g_v|_v\) to the covolume. The real multiplier contributes
\(\prod_\tau|\det g_\tau|\) by change of basis. For rational scalar
\(g=a\), the lattice \(xa\) in the assertion is \(O_F\) itself;
here integrality at all places is seen in its integral basis.
The scalar product formula follows.

Given a finite adele, the same basis and the Chinese remainder
argument in Lemma~\ref{lem:poisson-finite-torus} supply a rational
element of \(F\) congruent to it modulo \(\widehat A\).
Translation by \(O_F\) then puts the infinite part in a fixed closed
fundamental parallelepiped. That parallelepiped times \(\widehat A\)
is compact. Indeed the integral rings are inverse limits of finite
quotients, and a sequence in their countable product has a convergent
subsequence by successive extraction on the countably many residue
coordinates. The digit metrics give the product topology, for example
by summing bounded coordinate distances with coefficients \(2^{-i}\).

If \(s(a)=1\), the first assertion with \(n=1\) gives a lattice of
covolume \(V_F\), independently of \(a\). By the cube argument in
Lemma~\ref{lem:trace-finite}, there is \(0\ne b\in F\) with
\(ba_f\in\widehat A\) and all \(|\tau(b)a_\tau|\le C\)
for a fixed \(C>1\). The product formula now gives
\[
    1\le\prod_v |ba_v|_v^{-1}
       =\prod_\tau|\tau(b)a_\tau|\le C^f .
\]
In particular every real modulus is at least \(C^{-(f-1)}\).
At finite places the nonnegative valuation exponents have product
of their \(Q_v\)-powers at most \(C^f\). There are only finitely
many possibilities by \(Q_v\ge p\) and the bound on the number of
places over \(p\) in Lemma~\ref{lem:real-place-rings}. Each choice
of exponents gives a translate of \(\prod_v A_v^\times\), compact
by the same digit argument. The indicated archimedean bounds complete
the construction of a compact set.
\end{proof}

\begin{theorem}\label{thm:rank-one-reduction}
There are \(r_0>0\) and a compact
\(\mathcal K\subset {\rm SL}_2(\mathbb A_F)\), depending only on
\(F\), such that each \(g\in{\rm SL}_2(\mathbb A_F)\) has the form
\[
        g=\gamma\,m(d(r))\,\kappa,\qquad
        \gamma\in{\rm SL}_2(F),\quad r\ge r_0,\quad
        \kappa\in\mathcal K .
\]
\end{theorem}
\begin{proof}
By Lemma~\ref{lem:real-volume-compact} the row lattice defined by
\(g\) has covolume \(V_F^2\). The cube argument in
Lemma~\ref{lem:trace-finite} gives a nonzero \(w\in F^2\) with
\(wg_f\) integral everywhere and
\(\|(wg)_\tau\|_2\le C_F\) at all infinite places for a constant
\(C_F\ge1\). Complete \(w\) to the bottom row of a matrix
\(\gamma_1\) of determinant one over \(F\), by solving the single
nontrivial linear equation for its first row.
At a real place a rotation sends that bottom row of \(\gamma_1g\)
to \((0,b_\tau)\) with \(b_\tau=\|(wg)_\tau\|_2>0\).
At a finite place take a component \(b_v\) of the row of maximal
absolute value. Dividing the row by \(b_v\) gives an integral
row containing a unit. Complete it to an integral determinant-one
matrix by solving the determinant equation using that unit.
Its inverse takes the original row to \((0,b_v)\).
These changes of basis give a factorization
\(\gamma_1 g=n(u)m(a)k\), where \(k\) lies in the compact product
of the rotation groups and \({\rm SL}_2(A_v)\), and \(a=b^{-1}\).
At almost every finite place \(b_v\) is a unit because \(w\) has
a unit entry there and \(g_v\) is integrally invertible. Therefore
\(s(a)\ge C_F^{-f}\).

Put \(r=s(a)^{1/f}\ge C_F^{-1}\).
The idele \(a/d(r)\) has size one. By
Lemma~\ref{lem:real-volume-compact}, left multiplication of
\(\gamma_1g\) by a rational \(m(c)\) puts \(ca/d(r)\) in a fixed
compact set of ideles. The unipotent parameter becomes \(c^2u\).
Left multiplication by a rational upper unipotent now puts that
parameter in the fixed compact set of additive representatives from
the same lemma. Write the result \(n(u')m(d(r))m(a')k\).
Moving \(m(d(r))\) to the left replaces \(u'\) by \(d(r)^{-2}u'\).
The finite coordinates are unchanged, while \(r^{-2}\) is bounded.
These parameters all lie in a compact set by taking the continuous
image of the compact set of \(u'\) times the interval
\([0,C_F^2]\) of real multipliers on the infinite coordinates.
The remaining factors lie in compact sets as well. Inverting the
rational multipliers gives the asserted representation.
\end{proof}

\section{Quadratic Fourier transformations in rank one}

\begin{definition}\label{def:quadratic-fourier}
Write \(H(y)=\sum_{i=1}^j\bar y_i y_i\) on \(\mathbb A_E^j\), and
\(B(x,y)=H(x+y)-H(x)-H(y)\). On \(\mathbb A_F\) use
\[
       \Psi(u)=\exp(2\pi i\textstyle\sum_\tau u_\tau)\psi_f(-u_f/2).
\]
It equals one on \(F\) by Proposition~\ref{prop:additive}.
A test is a finite linear combination of products of real
Schwartz functions (smooth with all derivatives decreasing faster
than any inverse power) and locally constant compactly supported
finite functions. Each finite function is a finite sum of
indicators of cosets of \(m\widehat O_E^j\) for a common integer
\(m\ge1\): cover its compact support by finitely many cosets of
constancy and refine their levels. For the proof below the real
functions needed are \(\prod_\tau\exp(-\pi z_\tau H(y_\tau))\)
with \(\Re z_\tau>0\).
Put \(w=\left(\begin{smallmatrix}0&1\\-1&0\end{smallmatrix}\right)\).
On the generators of Definition~\ref{def:real-adeles} use
\[
 \begin{split}
  W(n(u))\phi(x)&=\Psi(-uH(x))\phi(x),\\
  W(m(a))\phi(x)&=s(a)^j\phi(ax),\\
  W(w)\phi(x)&=\int_{\mathbb A_E^j}\phi(y)\Psi(-B(x,y))\dd y .
 \end{split}
\]
In this Fourier integral take Lebesgue measure at infinity and
give \(\widehat O_E^j\) finite mass \(V_E^{-j}\). In particular
the diagonal copy of \(E^j\) has covolume one.
\end{definition}

\begin{lemma}\label{lem:quadratic-poisson}
Each of the three operators in Definition~\ref{def:quadratic-fourier}
is a product of local unitary operators. The Fourier operator
squared sends \(\phi(x)\) to \(\phi(-x)\). For rational generators
the functional \(\Theta(\phi)=\sum_{x\in E^j}\phi(x)\) is invariant.
\end{lemma}
\begin{proof}
Use local Fourier measure at a real place equal to \(2^j\) times
Lebesgue measure on \(\C^j\), since the real bilinear form is
\(2\Re\sum_i\bar x_i y_i\). At a finite place give \(B_v\) mass
\([B_v^\vee:B_v]^{-1/2}\) in each coordinate, where \(B_v^\vee\)
is its annihilator under the finite character of \(B(x,y)\).
The trace matrix in Proposition~\ref{prop:additive} gives
\(\prod_v [B_v^\vee:B_v]=|\Delta_E|\).
For any full additive lattice \(D\) at that place its annihilator
has reciprocal mass: for cosets of two nested lattices their
dual quotient has the same size by the invertible trace matrix,
or by using its inverse transpose in a \(\Z_p\)-basis.
Thus the local transform of a coset indicator is its mass times a
character times \({\bf1}_{D^\vee}\), by translation orthogonality.
Applying this twice gives reflection in zero, and multiplying two
such formulas and integrating gives preservation of the inner
product, by the orthogonality in Lemma~\ref{lem:norm-gauss}.
The product of these measures is the measure prescribed in
Definition~\ref{def:quadratic-fourier}, since
\(2^f/\sqrt{|\Delta_E|}=1/V_E\) by Lemma~\ref{lem:skew-covol}.

For ordinary real Fourier transformation with phase
\(\exp(-2\pi i x\cdot y)\), integration by parts and differentiation
under the integral give rapid decay of the transform and all its
derivatives on Schwartz functions. Inserting
\(\exp(-\pi\delta |y|^2)\) into the inverse integral yields convolution
with \(\delta^{-n/2}\exp(-\pi|x|^2/\delta)\) in real dimension \(n\).
This follows by multiplying the one-dimensional calculation of
Lemma~\ref{lem:theta} and interchanging the absolutely integrable
integrands. These Gaussian kernels have mass one and their mass
outside each fixed neighborhood of zero tends to zero by scaling.
Their convolution converges to the original Schwartz function, by
boundedness and uniform continuity. Dominated convergence also
removes the inserted Gaussian on the Fourier side. This proves
inversion, and substitution into the integral of the product of
two transforms proves preservation of the inner product by absolute
integration. The coordinate change for \(B\) gives precisely the
stated real measure. The other generators are unitary by their
unit-modulus multiplier and the Jacobian \(s(a)^{2j}\) respectively.

For rational \(u,a\) the character is trivial on rational
arguments, and \(a\) permutes the rational vectors with \(s(a)=1\)
by Lemma~\ref{lem:real-volume-compact}. To prove invariance under
\(w\) take a finite coset indicator at \(c+m\widehat O_E^j\),
choosing \(c\in E^j\) by Lemma~\ref{lem:poisson-finite-torus}.
The rational vectors lie at \(c+mO_E^j\), with archimedean
covolume \(m^{2fj}V_E^j\). Periodize the real test on this lattice
as in Lemma~\ref{lem:theta}. Rapid decay gives uniform convergence
of the periodization and absolute summability of its Fourier
coefficients, since the number of lattice points in a real cube
of side \(R\ge1\) is \(O((1+R)^{2fj})\) by a basis change.
The product Fej\'er kernels from Lemma~\ref{lem:poisson-finite-torus}
(or the one-dimensional kernels in Lemma~\ref{lem:theta} in each
coordinate) recover this series. Its dual indices are exactly the
rational elements of the finite annihilator, by the trace-dual
calculation of Lemma~\ref{lem:poisson-finite-torus} with conjugation
in the first argument. The real phase at \(c\) is the inverse of
the finite phase there, since \(\Psi\) is trivial on \(F\).
The reciprocal covolume is the finite mass of \(m\widehat O_E^j\)
in Definition~\ref{def:quadratic-fourier}. This gives
\(\Theta(W(w)\phi)=\Theta(\phi)\), with absolute convergence.
\end{proof}

\begin{theorem}\label{thm:rank-one-weil}
The local versions of these operators define \(W(g)\) up to a
scalar of modulus one for \(g\in{\rm SL}_2(\mathbb A_F)\), with
\(W(g)W(g')\) such a scalar times \(W(gg')\).
Products of rational generators representing the same rational
matrix give exactly the same operator. Thus their choice defines
an ordinary action of \({\rm SL}_2(F)\) preserving \(\Theta\).
All \(W(g)\) commute with local unitary isometries acting by
\(\phi(y)\mapsto\phi(h^{-1}y)\).
On any compact set of \(g\), the finite transforms of a fixed
finite test range over finitely many functions up to unit scalars.
At a real place a fixed radial Gaussian transforms to a scalar
times \(\exp(-\pi zH)\). On a compact set the scalars are bounded
in modulus and \(z\) ranges over a compact subset of \(\Re z>0\).
\end{theorem}
\begin{proof}
At a single place use the operators
\[
 D(u,v)\phi(x)=\Psi\big(B(v,x)+B(u,v)/2\big)\phi(x+u)
\]
with the local character. Direct substitution, and change of
variables in the Fourier integral, show that conjugation by
\(W(n(b)),W(m(a)),W(w)\) sends the column \((v,u)^{\mathsf T}\)
indexing \(D\) to \((v+bu,u)^{\mathsf T},(av,a^{-1}u)^{\mathsf T}\),
and \((u,-v)^{\mathsf T}\), respectively.
These matrices generate the local group: a nonzero bottom pivot
permits elementary triangular reduction and \(w\) exchanges pivots.
Consequently two words for the same matrix differ by a unitary
operator commuting with all translations and character multipliers.

At a finite place such a word sends test functions to test functions.
The image of \({\bf1}_D\), for an additive lattice \(D\), is supported
in \(D\) by invariance under all multipliers from its annihilator,
using the double annihilator computation in
Lemma~\ref{lem:quadratic-poisson}. Translation invariance by \(D\)
makes the image constant there. The scalar is the same for every
coset by translation, and for all lattices by subdividing across
a common smaller one. Thus it is scalar on tests.
At a real place the word maps Schwartz functions continuously to
Schwartz functions: the integrals and integration by parts in
Lemma~\ref{lem:quadratic-poisson} bound each derivative times a
polynomial by finitely many such bounds for the original function.
Differentiate translation and modulation commutation to get
commutation with every \(\partial_i\) and multiplication by \(x_i\)
in real coordinates. The Gaussian \(h=\exp(-\pi |x|^2)\) is killed
by each \(\partial_i+2\pi x_i\). Its image is therefore \(ch\), by
differentiating its quotient by \(h\). The span of the modulations
of \(h\) is dense in \(L^2\): if \(u\) is orthogonal to all of them
then the Fourier transform of \(u\bar h\in L^1\) vanishes.
Convolving with the Gaussian kernels of Lemma~\ref{lem:quadratic-poisson}
then gives zero, by Fourier integration of the Gaussian, while
these convolutions converge to \(u\bar h\) in \(L^1\).
For the latter assertion translations are continuous in \(L^1\):
first use finite sums of rectangular indicators, whose translated
symmetric differences have volumes tending to zero, then
approximate an integrable function by those sums. Such approximation
follows by truncating the function and its domain and approximating
its level sets in measure by finite unions of rectangles, from the
definition of Lebesgue measure by countable box covers.
The orthogonal-complement criterion used here follows by minimizing
distance to a closed subspace: the parallelogram identity makes a
minimizing sequence Cauchy, and varying the limit in each subspace
direction gives orthogonality. Thus the word equals the same scalar
on every Schwartz function by unitarity. In both cases the scalar
has modulus one.

At every finite place away from the finitely many divisors of the
different, the local vacuum \({\bf1}_{B_v^j}\) is fixed by the
chosen \(w\), integral unipotents and diagonal units. On integral
arguments \(\Psi_v\) is trivial, by the trace calculation for \(F\)
in Theorem~\ref{thm:finite-norm-product}. An integral
determinant-one matrix has a unit bottom pivot after a possible
exchange, so it is a product of these integral generators.
Use words fixing this vacuum at all such places. This defines the
product on every test: almost all its factors are vacua. Changes
of words can now give nontrivial scalars at only finitely many
places. On rational generators use the products in
Definition~\ref{def:quadratic-fourier}. By
Lemma~\ref{lem:quadratic-poisson}, two rational words giving the
same matrix preserve \(\Theta\); their scalar ratio is one since
\(\Theta\) is nonzero on a positive Gaussian times the integral
finite indicator. Commutation with isometries holds on the three
generators by preservation of \(B,H\) and the additive measure.
For the last point a unitary determinant has norm one, giving
absolute Jacobian one also in the split case.

At each finite place a fixed test and its Fourier transform are
bounded in support and constant on cosets of some open lattice.
The test is therefore fixed by sufficiently small upper unipotents
and by diagonal units sufficiently close to one. The same assertion
for its Fourier transform gives small lower unipotents by conjugation.
Near the identity the upper and lower unipotent and diagonal
factorization uses a unit pivot and continuous parameters, so the
test is fixed projectively by an open subgroup. A compact set meets
finitely many cosets of that subgroup. For compact adelic sets,
all but a fixed finite set of components lie in the integral groups:
cover by finitely many basic restricted-product neighborhoods.
At the real places the three changes of Gaussian parameter are
\(z\mapsto z+2iu,\ z\mapsto a^2z,\ z\mapsto4/z\), with scalar
multipliers \(1,|a|^j,(2/z)^j\). The last formula follows from
Lemma~\ref{lem:theta} for real \(z>0\) and analytic continuation
under the integral for \(\Re z>0\) by Lemma~\ref{lem:analytic-tools}.
In a neighborhood of any given matrix the pivot and word can
be chosen continuously. A finite cover of a compact set gives
the bounds on modulus and parameter.
\end{proof}

\section{A bound in the rank-one cusp}

\begin{lemma}\label{lem:real-frequencies}
For \(t\in F\) set \(D(t)=\prod_v Q_v^{\max(0,-v_A(t))}\)
if \(t\ne0\), and \(D(0)=1\). For real \(j>2,\ R\ge1\),
\[
 \sum_{t\in F}D(t)^{-j}
       \prod_\tau (1+|\tau(t)|/R)^{-j}
 \ \le\
 \left(\frac{3j}{j-1}\sum_{n\ge1} n^{1-j}\right)^f R^f .
\]
\end{lemma}
\begin{proof}
The ideal \(J(t)=\{a\in O_F:at\in O_F\}\) has index \(D(t)\).
Indeed the proof of Lemma~\ref{lem:denominator-series} applies to
\(O_F\), using its integral basis and the factors in
Lemma~\ref{lem:real-place-rings}. In particular the exponents
determine the ideal and the number of places above \(p\) is at
most \(f\), so that the sum of \((\mathrm NJ)^{1-j}\) over
nonzero integral ideals \(J\) is at most \((\sum_{n\ge1}n^{1-j})^f\).
For \(t=0\) take \(J=O_F\). If \(t,t'\) are distinct members of
\(J^{-1}\), the inclusion \((t-t')J\subset O_F\) gives
\(1\le\mathrm NJ\prod_\tau|\tau(t-t')|\) by scaling real
covolumes as in Lemma~\ref{lem:real-volume-compact}. In the box
\(|\tau(t)|\le R_\tau\), \(R_\tau\ge1\), rescale each coordinate
by \(R_\tau^{-1}\). Distinct points have maximum distance at
least \(\delta=(\mathrm NJ\prod R_\tau)^{-1/f}\). Cubes of side
\(\delta\) about them are disjoint up to boundaries and lie in a
cube of side \(2+\delta\). Hence there are at most
\(3^f\mathrm NJ\prod R_\tau\) points.
Bound each weight \((1+|\tau(t)|/R)^{-j}\) by
\(\int_1^\infty j b^{-j-1}{\bf1}_{|\tau(t)|\le Rb}\dd b\).
Apply the box estimate and integrate. Summing over the possible
\(J(t)\) yields the result by the ideal sum above.
\end{proof}

\begin{definition}\label{def:theta-and-eisenstein}
Fix \(j\ge3\). Take a test \(\phi\) of Definition~\ref{def:quadratic-fourier}
with radial Gaussians at the infinite places, allowing finite linear
combinations. Let \(K\) be an open subgroup of the full finite
integral unitary level \(K^0\) in rank \(j\) that fixes its finite
test. Such a subgroup exists by the coset description of tests:
matrices sufficiently close to the identity preserve each coset at
the finitely many exceptional places. Write \(I=[K^0:K]\) and let
\(h\) range over the finite class representatives at this level
in Definition~\ref{def:columns}. Put \(\Gamma_h=\Gamma\cap hKh^{-1}\)
and \((h\phi)(y)=\phi(y_\infty,h^{-1}y_f)\). Set
\[
 \mathcal T(\phi)=\frac{1}{I\mathfrak m_j}
       \sum_{[h]}\frac{\Theta(h\phi)}{|\Gamma_h|},\qquad
 I_t(\phi)=\int_{\mathbb A_E^j}\phi(y)\Psi(-tH(y))\dd y,\qquad
 \mathcal E(\phi)=\phi(0)+\sum_{t\in F} I_t(\phi).
\]
For \(I_t\) one can also take \(t\in\mathbb A_F\).
\end{definition}

\begin{proposition}\label{prop:rank-one-integrals}
The series defining \(\mathcal E\) converges absolutely. Both
\(\mathcal T,\mathcal E\) are independent of the representatives
and level in Definition~\ref{def:theta-and-eisenstein} and invariant
under \(W(\gamma)\) for \(\gamma\in{\rm SL}_2(F)\).
For \(r\ge1\) and any compact set of \(\kappa\in{\rm SL}_2(\mathbb A_F)\),
write \(\phi'=W(\kappa)\phi\), with any of the unit-scalar choices of
Theorem~\ref{thm:rank-one-weil}. Then uniformly on that set,
\[
 \sum_{t\in F}|I_t(W(m(d(r)))\phi')|=O_\phi(r^{(2-j)f}).
\]
The constant may also depend on \(F,j\) and the compact set.
\end{proposition}
\begin{proof}
For a finite test and finite adele \(s\) the absolute value of its
finite quadratic integral is at most \(C D(s)^{-j}\), with \(D\)
defined by the same finite product as in Lemma~\ref{lem:real-frequencies}
(take factor one at zero). To check this at almost all places take
the integral vacuum, where Lemma~\ref{lem:norm-gauss} gives the
exact modulus with conductor \(\max(0,-v_A(s))\).
At a remaining place write the test as a sum of indicators of
\(c+\pi^e B_v^j\), \(e\ge0\). Substituting \(y=c+\pi^e z\)
adds a linear character and a constant to the quadratic phase
with parameter \(s\pi^{2e}\).
The local character \(\Psi_v\) has annihilator \(\pi^{-d_v}A_v\)
against \(A_v\), for some \(d_v\ge0\): this is the \(F\)-trace
dual calculation in Theorem~\ref{thm:finite-norm-product}.
The induced conductor on \(A_v\) is
\(n=\max(0,-v_A(s)-2e-d_v)\). The squared integral calculation
in Lemma~\ref{lem:norm-gauss} still bounds the modulus of the
probability average by \(Q_v^{-jn}\): the extra linear character
evaluated at the difference has modulus one and does not change
the orthogonality in the other variable. The case \(n=0\) uses
the bound one. Restoring the fixed volume of the coset proves the
bound with a constant depending on the test. Multiplication proves
the finite assertion, uniformly for finitely many tests.

At infinity substitute \(z=ry\) into the integral. The measure
and the multiplier in Definition~\ref{def:quadratic-fourier}
give a factor \(r^{-jf}\). The remaining Gaussian integral at
\(\tau\) is \((z_\tau+2i\tau(t)/r^2)^{-j}\) using ordinary
complex Lebesgue measure and the Gaussian normalization
of Lemma~\ref{lem:theta}. If \(z_\tau\) ranges over a compact
subset of the right half-plane this has modulus at most
\(C(1+|\tau(t)|/r^2)^{-j}\); use a positive lower bound for
\(\Re z_\tau\) on bounded arguments and the triangle inequality
on large ones. The finite tests and Gaussian parameters after
\(\kappa\) satisfy exactly these uniform conditions by
Theorem~\ref{thm:rank-one-weil}. Apply
Lemma~\ref{lem:real-frequencies} with \(R=r^2\) to prove the
bound and convergence.

For \(\mathcal T\), a change of left representative permutes rational
vectors and preserves the real radial test by the rational unitary
matrix at infinity. Right multiplication by \(K\) preserves the
finite test. Refining \(K\) multiplies the class sum and its total
weight by the same index by the orbit counting in
Theorem~\ref{thm:marked-mass}. Use an intersection for any two
choices. Theorem~\ref{thm:rank-one-weil} and
Lemma~\ref{lem:quadratic-poisson} give invariance under rational
matrices. For \(\mathcal E\), let \(P\) be the upper triangular
subgroup of \({\rm SL}_2(F)\).
The cosets \(P\backslash{\rm SL}_2(F)\) are determined by the
projective bottom row, since matrices with proportional bottom rows
differ by a left upper triangular matrix. Representatives are
\(1\) and \(w n(t)\), \(t\in F\). Every \(W(p)\) for \(p\in P\)
fixes evaluation at zero by the scalar product formula in
Lemma~\ref{lem:real-volume-compact}. The sum of \((W(\gamma)\phi)(0)\)
over these representatives equals \(\mathcal E(\phi)\).
The ordinary rational action in Theorem~\ref{thm:rank-one-weil}
and absolute convergence permit reindexing by right translation.
\end{proof}

\begin{theorem}\label{thm:bounded-theta-difference}
For every test of Definition~\ref{def:theta-and-eisenstein},
\[
      \sup_{g\in{\rm SL}_2(\mathbb A_F)}
                 |(\mathcal T-\mathcal E)(W(g)\phi)|<\infty .
\]
\end{theorem}
\begin{proof}
Use Theorem~\ref{thm:rank-one-reduction} and rational invariance
in Proposition~\ref{prop:rank-one-integrals}. Scalar changes in
Theorem~\ref{thm:rank-one-weil} have modulus one. It suffices to
bound the difference at \(W(m(d(r)))\phi'\) with \(r\ge r_0\)
and \(\phi'=W(\kappa)\phi\) for compact \(\kappa\).
There are finitely many finite components up to scalar, so one
level suffices and there are finitely many \(h\) to consider by
Theorem~\ref{thm:marked-mass}. For each of these, the finite
support restricts the rational vectors to a subset of
\(M^{-1}O_E^j\) for a fixed integer \(M\), by clearing finite
denominators. The real factors for nonzero vectors have modulus
at most \(C r^{jf}\exp(-c r^2\sum_\tau H(x_\tau))\), \(c>0\).
The lattice is discrete, has positive least length, and has
\(O((1+H_0)^{2fj})\) points of length at most \(H_0\) by a basis
change. Summing over unit intervals in length therefore bounds
its nonzero contribution by \(O(r^{-1})\) for \(r\ge1\): factor
out half the Gaussian at the least positive length and sum the
remaining fixed Gaussian. The zero term of the normalized class
sum is \(r^{jf}\phi'(0)\) by Theorem~\ref{thm:marked-mass}.
This is exactly the first term of \(\mathcal E(W(m(d(r)))\phi')\).
Proposition~\ref{prop:rank-one-integrals} bounds the other terms
by \(O(r^{(2-j)f})\).
The range \(r_0\le r\le1\), when present, only enlarges the compact
set of transforms, and the same bounds at \(r=1\) suffice.
\end{proof}

\section{The rank-one Siegel--Weil identity}

\begin{lemma}\label{lem:split-mixed-model}
Let \(v\) lie above a completely split rational prime outside the
different. Write \(L=\operatorname{Frac}(A_v)\), \(A=A_v\), and
write a vector in \((E\otimes_F L)^j\) as two columns \(x,y\in L^j\).
Its norm is \(x^{\mathsf T}y\). With \(A^j\) of mass one define
\[
  \widetilde\phi(x,z)=\int_{L^j}\phi(x,y)\Psi_v(-z^{\mathsf T}y)\dd y,
           \qquad X=[\,x\ z\,].
\]
This is an invertible map on compact locally constant tests and
sends the integral vacuum to \(\phi_0\) of
Definition~\ref{def:rank-one-hecke}. Under this map the unitary
isometry acting by \(\phi(b^{-1}x,b^{\mathsf T}y)\) induces \(U_b\),
and the local generators \(n(u),m(a),w\) induce the operators
\(V_{n(u)},V_{m(a)},V_w\) of that same definition.
\end{lemma}
\begin{proof}
The finite character \(\Psi_v\) has annihilator \(A\) against \(A\)
since the place is unramified of rational residue degree one,
by Proposition~\ref{prop:additive} and the identification of
the character in Theorem~\ref{thm:finite-norm-product}. Hence
the finite Fourier inversion and the transform of the indicator
of \(A^j\) follow from Lemma~\ref{lem:quadratic-poisson}. Substitution
\(y'=b^{\mathsf T}y\) gives
\(|\det b|^{-1}\widetilde\phi(b^{-1}x,b^{-1}z)=U_b\widetilde\phi(X)\).
The norm phase for \(n(u)\) replaces \(z\) by \(z+ux\).
For \(m(a)\), substitution \(y'=a y\) cancels its Jacobian against
the factor \(|a|^j\) and replaces \((x,z)\) by \((ax,a^{-1}z)\).
For \(w\) the full Fourier kernel is
\(\Psi_v(-x^{\mathsf T}y'-y^{\mathsf T}x')\); the second Fourier
transform in the variable paired with \(x'\) evaluates the
\(x'\)-argument at \(-z\) by reflection. The result is
\(\widetilde\phi(-z,x)\). These three changes of columns are right
multiplication of \(X\) by their stated matrices.
\end{proof}

\begin{theorem}\label{thm:one-vector-siegel-weil}
For every \(j\ge3\) and every test of
Definition~\ref{def:theta-and-eisenstein} one has
\[
                \mathcal T(\phi)
                       =\phi(0)+\sum_{t\in F}
                          \int_{\mathbb A_E^j}\phi(y)\Psi(-tH(y))\dd y .
\]
\end{theorem}
\begin{proof}
A finite test is a finite sum of pure products with integral vacuum
at almost all places, by its coset description in
Definition~\ref{def:quadratic-fourier}. Work with one such product.
Choose a prime \(p\equiv1\bmod d N_1\) outside its exceptional
places and the different by Corollary~\ref{cor:ap}. By
Lemma~\ref{lem:split-level} we can choose \(v\mid p\) with the
split model of Lemma~\ref{lem:split-mixed-model}. Freeze the test
away from \(v\). Let \(\lambda\) on local matrix tests be
\(\mathcal T-\mathcal E\) applied to this product after inverting
the partial Fourier transform at \(v\).

Both functionals are invariant under each local isometry
\((x,y)\mapsto(bx,b^{-\mathsf T}y)\). For \(\mathcal E\), this
preserves norm, zero and the product measure, since the two
absolute determinants multiply to one. For \(\mathcal T\), write
\(h\) for the isometry supported at \(v\). Choose a product level
\(K\) fixing the test, small enough that \(K\) and \(hKh^{-1}\)
are both contained in \(K^0\). Right translation \(g\mapsto gh\)
carries the classes at \(hKh^{-1}\) bijectively to those at \(K\).
The stabilizers have equal order since
\(g(hKh^{-1})g^{-1}=(gh)K(gh)^{-1}\). The two indices in
\(K^0\) are also equal. At the only relevant place the group is
\({\rm GL}_j(L)\); on invertible matrices the measure
\(|\det X|^{-j}\dd X\), with \(\dd X\) additive matrix measure,
is invariant under left and right multiplication. This follows
by the row or column Jacobian as in Definition~\ref{def:rank-one-hecke}.
An open integral subgroup has positive finite mass by a compact
cover with finite-index cosets. Conjugation preserves its measure,
so the index equality follows by counting these cosets. Apply
the independence of the level in
Proposition~\ref{prop:rank-one-integrals}. By
Lemma~\ref{lem:split-mixed-model} the resulting \(\lambda\) is
invariant under every \(U_b\).

For \(g\in{\rm SL}_2(L)\) embed \(g\) at \(v\) and put the
identity elsewhere. Then \(|\lambda(V_g\phi_0)|\) is precisely the
absolute value of \((\mathcal T-\mathcal E)(W(g)\phi)\) by
Lemma~\ref{lem:split-mixed-model} and Theorem~\ref{thm:rank-one-weil}.
It is bounded on this group by Theorem~\ref{thm:bounded-theta-difference}.
Theorem~\ref{thm:bounded-split-hecke} thus gives
\(\lambda(\phi_0)=0\). Since the original test at \(v\) is the
integral vacuum, this is the claimed equality. Linear combination
treats the general test.
\end{proof}

\section{Siegel's formula for one represented element}

\begin{theorem}\label{thm:one-vector-class-average}
Let \(j\ge3,\ \alpha\in F\) be totally positive. Fix a compactly
supported locally constant complex test \(\Phi_f\) on \(E_f^j\),
and take any open \(K\subset K^0\) fixing it. Use the classes
\([h]\), stabilizers \(\Gamma_h\) and index \(I\) of
Definition~\ref{def:theta-and-eisenstein}. Give \(\widehat O_E^j\)
mass one for \(\dd x_f\), and the trace dual of \(\widehat A\)
mass one for \(\dd z_f\). Then
\[
\begin{split}
 \frac{1}{I\mathfrak m_j}
 \sum_{[h]}\frac{1}{|\Gamma_h|}
       \sum_{\substack{y\in E^j\\H(y)=\alpha}}\Phi_f(h^{-1}y)
 ={}& \frac{V_F}{V_E^j}
       \prod_\tau\frac{\pi^j\,\tau(\alpha)^{j-1}}{(j-1)!}\\
 &{}\cdot\int_{F_f}\!
       \left(\int_{E_f^j}\Phi_f(x)
       \psi_f\big(z(H(x)-\alpha)/2\big)\dd x_f\right)\dd z_f .
\end{split}
\]
Every inner rational sum is finite on the support of the test
and the outer finite-adelic Fourier integral converges absolutely.
\end{theorem}
\begin{proof}
Use \(\phi(y)=\exp(-\pi\sum_\tau H(y_\tau))\Phi_f(y_f)\).
On \(\mathbb A_F\) take infinite Lebesgue measure and finite
measure giving \(\widehat A\) mass \(1/V_F\). A half-open real
fundamental parallelepiped of \(O_F\) times \(\widehat A\) is a
fundamental domain \(\Omega\) of volume one. Indeed
Lemma~\ref{lem:real-volume-compact} gives representatives modulo
\(\widehat A\) in the finite coordinates, unique modulo \(O_F\),
and the half-open convention chooses the remaining translate.
The measure on the quotient is invariant by translation: partition
a translate of \(\Omega\) by its intersections with the disjoint
sets \(a+\Omega,\ a\in F\) and translate the pieces back. A character
\(u\mapsto\Psi(\beta u)\), \(\beta\in F\), integrates to zero there
unless \(\beta=0\). In fact a nonzero \(\beta\) admits a translation
where its character is different from one, by varying a single real
coordinate; the character descends by
Definition~\ref{def:quadratic-fourier}.

Apply Theorem~\ref{thm:one-vector-siegel-weil} to \(W(n(u))\phi\),
multiply by \(\Psi(\alpha u)\), and integrate over \(\Omega\).
The level on the theta side can be used for every \(u\) because the
phase depends only on the norm. Its summands are absolutely
dominated by a fixed decaying Gaussian on finitely many fractional
lattices, as in Theorem~\ref{thm:bounded-theta-difference}.
Orthogonality retains the rational vectors of norm \(\alpha\).
Their real Gaussian factor is \(\exp(-\pi\sum_\tau\tau(\alpha))\);
they lie in bounded balls in these lattices, so their counts on
the support are finite. The term \(\phi(0)\) on the other side
integrates to zero as \(\alpha\ne0\).
The terms of the series are \(I_{t+u}(\phi)\Psi(\alpha u)\),
\(t\in F\). Since \(\Psi(\alpha t)=1\), unfolding the translates
of \(\Omega\) gives \(\int_{\mathbb A_F}I_u(\phi)\Psi(\alpha u)\dd u\).

For absolute integration in this step, the finite bound in the proof of
Proposition~\ref{prop:rank-one-integrals} is \(C D(u_f)^{-j}\).
Over \(A_v\) of mass one the shell \(v_A(u)=-n,\ n>0\), has mass
at most \(Q_v^n\). The integral over finite adeles with all
\(A_v\) of mass one is therefore bounded by
\(\prod_v(1+\sum_{n\ge1}Q_v^{(1-j)n})\). This is finite by
\(j>2\) and the prime sum in Theorem~\ref{thm:finite-norm-product}.
At a real coordinate the bound is \(C(1+|u_\tau|)^{-j}\), also
integrable. Thus the unfolding and separation into the real
coordinate factors and the finite factor are justified by absolute
integration after the inner quadratic integral has been taken.

At a real coordinate set \(a=\tau(\alpha)>0\).
By Proposition~\ref{prop:norm-arch-factor} its Fourier integral
at \(a\) is the inverse transform of the density
\({\bf1}_{t>0}\pi^j t^{j-1}\exp(-\pi t)/(j-1)!\). It has value
\(\pi^j a^{j-1}\exp(-\pi a)/(j-1)!\): insert the Gaussian cutoff
in the outer integral as in Lemma~\ref{lem:quadratic-poisson};
the convolution converges at \(a\) by boundedness of the density
and continuity there. Removing the cutoff is justified by the
integrable transform bound just proved. The Gaussian factors now
cancel on both sides. The finite vector measure in
Definition~\ref{def:quadratic-fourier} gives \(\widehat O_E^j\)
mass \(V_E^{-j}\). In an integral basis of \(O_F\) its trace
matrix is \(T^{\mathsf T}T\) with real embedding matrix \(T\) of
determinant of absolute value \(V_F\) by Lemma~\ref{lem:skew-covol}.
Thus \([\widehat A^\vee:\widehat A]=V_F^2\) by the same basis
calculation as in Proposition~\ref{prop:additive}; the finite
measure on the outer variable gives \(\widehat A^\vee\) mass \(V_F\).
The remaining finite phase is precisely the one displayed, by
the sign in Definition~\ref{def:quadratic-fourier}. The two changes
to the stated measures give \(V_F/V_E^j\).
\end{proof}

\begin{corollary}\label{cor:actual-mass-ratio}
For every \(j\ge3\) and every totally positive integral
\(\alpha\in O_F\),
\[
   \frac{1}{\mathfrak m_j}
       \sum_{[L]\in\mathcal L_j}\frac{r_L(\alpha)}{|U(L)|}
                              = B_j(\alpha)
\]
with \(B_j\) of Proposition~\ref{prop:norm-arch-factor}. In particular
\[
 \frac{\mathfrak m_j}{\mathfrak m_{j-1}}
 =\left(\frac{{\rm rd}_E^{\,j-1/2}(j-1)!}{(2\pi)^j}\right)^f
   \prod_v (1-\sigma_v Q_v^{-j})^{-1}.
\]
The logarithm of the product has absolute value at most \(4f\,2^{-j}\).
\end{corollary}
\begin{proof}
Take \(K=K^0,\ \Phi_f={\bf1}_{\widehat O_E^j}\) in
Theorem~\ref{thm:one-vector-class-average}. By
Lemma~\ref{lem:trace-finite} the inner sum is the representation
count of its lattice. The finite integral equals
\(P_{\varnothing}(\alpha)\) by Theorem~\ref{thm:finite-norm-product},
and \(V_F=|\Delta_E|^{1/4}\) by Lemma~\ref{lem:real-volume-compact}.
This is exactly the definition of \(B_j(\alpha)\).
The special case \(\alpha=1\) satisfies the hypothesis proved
sufficient in Proposition~\ref{prop:norm-arch-factor}, giving the
ratio and its logarithmic bound.
\end{proof}

\section{The column level and local norm measures}

\begin{proposition}\label{prop:marked-local-level}
Extend a finite isometry of the definite space by the identity on
\(z,e_a,e_a^*\) in Definition~\ref{def:row-space}. The inverse
image of the full integral stabilizer in this extension is a
subgroup of \(K^0\) of index \((N^k-1)^f\). The inverse image
of the row level imposing identity reduction also at \(g_0,g_1\)
is exactly \(K_c\) of Definition~\ref{def:columns}.
In \(\mathcal G\) let \(\bar H\) be the image of the definite
integral isometries preserving the row lattice, and \(\bar S\)
the image of the finite integral subgroup generated by these
isometries and the unipotent group of Definition~\ref{def:unip}.
Then
\[
 |\bar H|=\prod_{i=0}^1 |{\rm GL}_k(\mathbf F_{g_i})|^f,\qquad
 |\bar S|=|\bar H|\,g^{f\nu(2k+2\nu+1)}.
\]
\end{proposition}
\begin{proof}
Away from \(N\) the change of middle lattice basis is integral
with integral inverse because \(\pi_N\bar\pi_N=N\).
At \(N\) use the components of Lemma~\ref{lem:split-level},
first taking one with \(\pi_N\) divisible by \(N\).
Projection of the middle lattice to \(X_0=\langle x_1,\dots,x_k\rangle\)
along \(z\) is the full integral lattice in this component because
\(\bar\pi_N\) is a unit. An isometry fixing \(z\) and preserving
the middle lattice therefore acts integrally invertibly on
\(X_0\) here. Its partner is the inverse transpose and is also
integrally invertible.
Write \(A\) for the matrix on \(X_0\). For \(i<k\) the image of
\(x_i\) in the new basis has \(z\)-coordinate \(-A_{ki}/\pi_N\)
and \(v\)-coordinate \(\bar\pi_N A_{ki}\). The image of
\(v=x_k/\bar\pi_N+z/N\) has coordinates \(A_{ik}/\bar\pi_N\)
for \(i<k\), \(A_{kk}\) at \(v\), and \((1-A_{kk})/N\) at \(z\).
The chosen component thus requires precisely that the last row
of \(A\) reduce to the last standard row. In the partner it
requires the last column of the inverse transpose to reduce to
the last standard column, the very same condition. Inverses
preserve the condition as well. Proposition~\ref{prop:column-dim}
counts the index, and at \(g_0,g_1\) integral change of basis
carries identity congruence to itself. This proves the level equality.

At a chosen component above either \(g_i\), all the matrices on
\(X_0\) occur by integral lifting and the split description in
Lemma~\ref{lem:split-level}, independently across components.
In the ordered flag use the line blocks of the \(e_a\), the middle
block of size \(k+1\), and the dual lines in reverse order.
In one split component the unipotent matrices are precisely the
upper block triangular matrices with identity diagonal blocks.
Indeed the flag and dual flag give the two sequences of invariant
subspaces, and identity on their quotients gives the diagonal
blocks. There are no further equations in that component: on the
partner the inverse adjoint is forced, and satisfies the flag
conditions by orthogonality. All residues lift by lifting arbitrary
off-block entries. There are
\(\binom r2-\binom{k+1}2=\nu(2k+2\nu+1)\) such entries.
The definite matrices normalize this unipotent group and intersect
it trivially by the middle diagonal block. Multiplication of the
two finite subgroups is therefore a bijection from their Cartesian
product onto their generated subgroup. Multiplying the orders
over the two primes and the \(f\) components gives the formula.
\end{proof}

\begin{definition}\label{def:finite-fiber}
Let \(T\) be a nonempty finite set of places of \(F\) and put
\(F_T=\prod_{v\in T}\operatorname{Frac}(A_v)\),
\(E_T=\prod_{v\in T}(B_v\otimes_{A_v}\operatorname{Frac}(A_v))\).
Give the integral rings and their finite Cartesian powers mass
one. Write \(L_n=\prod_{v\in T}\pi_v^n A_v\).
For \(\gamma\in F_T\), compact locally constant
\(\varphi:E_T^k\to\C\), and \(k\ge3\), put
\[
 d_T(\gamma,\varphi)=\lim_{n\to\infty}
   \frac{1}{|L_n|}\int_{H(x)-\gamma\in L_n}\varphi(x)\dd x .
\]
Here and below \(|L|\) for an additive compact open set denotes its
measure in its specified ambient space. Let \(K_{c,T}\) be the
projection of the column level to the factors over \(T\).
\end{definition}

\begin{lemma}\label{lem:fiber-continuity}
The limit in Definition~\ref{def:finite-fiber} exists uniformly
in \(\gamma\) and is a continuous function of \(\gamma\).
It is locally constant wherever every component of \(\gamma\)
is nonzero. It is a positive linear functional on nonnegative tests.
\end{lemma}
\begin{proof}
Use the additive character \(\psi_f(z\,\cdot/2)\) restricted to
these places and give the dual of \(\prod_{v\in T}A_v\) mass one
in the \(z\)-variable. Integrating a character over \(L_n^\vee\)
gives \(|L_n|^{-1}{\bf1}_{L_n}\) by the trace pairing and the
finite orthogonality in Lemma~\ref{lem:norm-gauss}. The approximants
therefore have Fourier expressions with \(z\) restricted to
\(L_n^\vee\). By the proof of Proposition~\ref{prop:rank-one-integrals},
the absolute value of the quadratic Fourier integral of the
test is bounded by \(C D(z)^{-k}\) independently of \(\gamma\).
It is integrable by the shell calculation in
Theorem~\ref{thm:one-vector-class-average}. The tail outside
\(L_n^\vee\) tends to zero and is independent of \(\gamma\) in
this bound. Every finite-level approximant is locally constant,
so uniform convergence implies continuity. Positivity follows
from the original integrals.

For a scalar \(u\in F_T\) sufficiently close to one at all places
we can write \(u=a\bar a\) with \(a\) arbitrarily close to one in
\(E_T\). In the field case apply the successive trace correction
of Lemma~\ref{lem:norm-one-isometry} beginning with \(1\) at any
prescribed positive digit depth. In a split pair take \((u,1)\).
Multiplication by all such \(a\) sufficiently close to one fixes
the test, by compact support and its finite coset decomposition.
It preserves the vector measure and multiplication by \(u\)
preserves each \(L_n\). The substitution \(x=ay\) therefore gives
\(d_T(u\gamma,\varphi)=d_T(\gamma,\varphi)\).
When every component of \(\gamma\) is nonzero these \(u\gamma\)
form a neighborhood of it.
\end{proof}

\begin{theorem}\label{thm:localized-class-equation}
Let \(k\ge3\) be the column rank, let \(h_c\) represent the
classes at \(K_c\), and let \(\beta\in F\) be totally positive
and integral off \(T\). For a \(K_{c,T}\)-invariant compact locally
constant test \(\varphi\) on \(E_T^k\) one has
\[
\begin{split}
 \frac1{|\mathcal C(K_c)|}\!
 \sum_c\ \sum_{\substack{y\in E^k,\ H(y)=\beta\\
                    (h_c^{-1}y)_v\in B_v^k\ (v\notin T)}}
       \varphi((h_c^{-1}y)_T)
    ={}&\frac{V_F}{V_E^k}
       \prod_\tau\frac{\pi^k\tau(\beta)^{k-1}}{(k-1)!}\\
    &{}\cdot P_T(\beta)d_T(\beta_T,\varphi).
\end{split}
\]
Here \(P_T\) is the product outside \(T\) in
Theorem~\ref{thm:finite-norm-product}.
For every fixed \(D>0\) it lies between
\((1+b^{-D})^{-1}\) and \(1+b^{-D}\) at all sufficiently large
arithmetic scales, uniformly in \(\beta,T\).
\end{theorem}
\begin{proof}
Use \(\varphi\) at \(T\) and the integral vacuum outside \(T\)
in Theorem~\ref{thm:one-vector-class-average}. The product test
is fixed by \(K_c\). All its class stabilizers have order one
and \(|\mathcal C(K_c)|=I_c\mathfrak m_k\) by
Proposition~\ref{prop:column-dim}.
The Fourier integral over \(T\) equals \(d_T(\beta_T,\varphi)\)
by the proof of Lemma~\ref{lem:fiber-continuity}. The complementary
integral factors absolutely and is \(P_T(\beta)\) by the proof
of Theorem~\ref{thm:finite-norm-product}. The logarithmic bound
there is \(o(b^{-D})\) by the scale estimate in
Corollary~\ref{cor:off-places-norm}. Exponentiating gives both
inequalities.
\end{proof}

\section{The one-sided Siegel bound on the middle quadric}

\begin{definition}\label{def:top-norm-fiber}
Take \(T\) as in Definition~\ref{def:finite-fiber} containing every
place over \(N\). On the middle space of Definition~\ref{def:row-space}
use its integral basis \(x_1,\ldots,x_{k-1},v,z\) at \(T\) and give
the integral coordinate lattice mass one for \(\dd e\).
Write \(q(e)=h(e,e)\), \(s(e)=h(z,e)\). For \(u\in E_T,\ y\in E_T^k\)
put \(P(u,y)=\sum_i y_i x_i-u z/N\) in the rational middle basis.
For \(b\in F_T\) and compact locally constant \(G\) on the middle
space define
\[
 \delta_b(G)=\lim_{n\to\infty}|L_n|^{-1}
                    \int_{q(e)-b\in L_n}G(e)\dd e ,
\]
using \(L_n\) in the scalar space \(F_T\) from
Definition~\ref{def:finite-fiber}.
\end{definition}

\begin{proposition}\label{prop:top-norm-jacobian}
The above limit defines a Borel measure supported on \(q(e)=b\),
finite on compact sets. Put \(\beta_b(u)=b+u\bar u/N\). For compact
locally constant \(G\) one has
\[
 \delta_b(G)=N^f\int_{E_T}
          d_T\big(\beta_b(u),G(P(u,\cdot))\big)\dd u .
\]
If every component of \(\beta_b(t)\) is nonzero then for all
sufficiently small compact open \(L\subset E_T\) about zero,
and also for measurable subsets of those \(L\),
\[
 \int_{s(e)\in t+L}G(e)\dd\delta_b(e)
       =N^f |L|\,d_T\big(\beta_b(t),G(P(t,\cdot))\big).
\]
\end{proposition}
\begin{proof}
In the integral basis the coordinates of \(P(u,y)\) are
\((y_1,\ldots,y_{k-1},\bar\pi_N y_k,-y_k/\pi_N-u/N)\).
Consequently \(s(P(u,y))=u,\ q(P(u,y))=H(y)-u\bar u/N\);
the determinant from \((y,u)\) to \(e\) is \(-1/\pi_N\).
In the notation of Lemma~\ref{lem:denominator-series},
the product of the finite absolute values of \(\pi_N\) is
\(\prod_\tau|\tau\pi_N|^{-2}=N^{-f}\). All factors off \(T\)
are one by integrality of \(\pi_N,\bar\pi_N\) and their product
\(N\). The change of variables thus gives
\(\dd e=N^f\dd y\,\dd u\), by the determinant and index computation
in Lemma~\ref{lem:class-covol}. The composed test
\(G(P(u,y))\) is a finite sum of products \(\chi_i(u)\varphi_i(y)\)
of compact open coset indicators with scalar coefficients, since
\(P\) is a linear isomorphism. Lemma~\ref{lem:fiber-continuity}
gives uniform convergence of the tube averages for each \(\varphi_i\).
Integration over the compact supports in \(u\) yields the first
formula. For nonnegative tests the limit is nonnegative, and it
is additive on disjoint compact open sets.

To obtain an actual measure, partition the space into a countable
union of compact open shells using increasing scalar dilates of
the integral lattice. Each shell has finite mass by the formula
just proved. Refine its finitely many coordinate boxes successively
into smaller residue boxes with diameters tending to zero.
Partition a real interval of length equal to the shell mass into
intervals of the respective child masses at each stage. Compatible
boxes specify a unique point by completeness; endpoints of the
intervals form a null set. The resulting image of length realizes
the prescribed masses of all boxes, including when the shell mass
is zero. The uniqueness argument of Lemma~\ref{lem:poisson-finite-torus}
extends equality on the boxes to all Borel sets. Outside \(q=b\),
continuity of \(q\) gives a neighborhood where a sufficiently
small tube is empty, and a countable base covers this complement
by null sets. This proves support on the stated fiber.

In the product decomposition of \(G(P(u,y))\), each \(\chi_i\)
is constant near \(t\), and \(\beta_b(u)\) is continuous there.
Lemma~\ref{lem:fiber-continuity} now makes the integrand constant
near \(t\) when \(\beta_b(t)\) has all its components nonzero.
The first formula also applies with an extra compact open cutoff
in \(s(e)=u\). It gives the indicated equality of measures near
\(t\) on compact opens and hence on Borel subsets by the preceding
uniqueness argument, and on their completions.
\end{proof}

\begin{theorem}\label{thm:top-siegel-bound}
Let \(t\in E\) satisfy \(|\tau t|>N\) at every infinite place,
and put \(\beta=(t\bar t-N^2)/N\). Use the column class
representatives in Theorem~\ref{thm:localized-class-equation}.
Let \(\mathcal S\) be a finite set of distinct pairs \((c,y)\)
with \(y\in E^k,\ H(y)=\beta,\ (h_c^{-1}y)_v\in B_v^k\) for
\(v\notin T\). Give these pairs weights \(a_{c,y}\ge0\).
Assume \(\beta\) integral off \(T\), which holds if \(\mathcal S\)
is nonempty. Let \(G\ge0\) be a Borel function on the middle
space. Suppose for each pair there is a neighborhood of
\(P(t,(h_c^{-1}y)_T)\) throughout which
\(G(k'e)\ge a_{c,y}\) for all \(k'\) in the image of \(K_{c,T}\)
acting on the middle space. Then for all sufficiently small
measurable neighborhoods \(L\subset E_T\) of zero,
\[
\begin{split}
 |L|\sum_{(c,y)\in\mathcal S} a_{c,y}
 \ \le\ & |\mathcal C(K_c)|\,c_*\,P_T(\beta)
   \prod_{\tau}\big(|\tau t|^2-N^2\big)^{k-1}
   \int_{s(e)\in t+L}G(e)\dd\delta_{-N}(e),\\
 & c_*=\frac{V_F}{V_E^k}
             \left(\frac{\pi^k}{(k-1)!\,N^k}\right)^f .
\end{split}
\]
In particular one may replace \(P_T(\beta)\) by \(1+b^{-6}\)
at all sufficiently large arithmetic scales, uniformly in \(T,t\)
subject to these conditions.
\end{theorem}
\begin{proof}
For each pair choose a compact open neighborhood within the
neighborhood in the hypothesis and saturate it by the compact
group \(K_{c,T}\). Its saturation is open as a union of translates,
and compact as the image of the group times the original compact
set. All its points have \(G\)-value at least the designated
weight. Take \(\Phi(e)\) to be the maximum of the weights of
these saturations containing \(e\), taking zero elsewhere.
Then \(\Phi\le G\) is compact, locally constant and invariant
under \(K_{c,T}\), with \(\Phi(P(t,(h_c^{-1}y)_T))\ge a_{c,y}\).
Compactness of the group here follows from integrality and the
closed unitary and congruence equations by
Proposition~\ref{prop:marked-local-level} and the compactness in
Lemma~\ref{lem:real-volume-compact}.

The test \(\varphi(y')=\Phi(P(t,y'))\) is invariant under the
definite group \(K_{c,T}\), since its action extends by the
identity on \(z\). It is compactly supported and locally constant
by Proposition~\ref{prop:top-norm-jacobian}. Apply
Theorem~\ref{thm:localized-class-equation} and retain the pairs
in \(\mathcal S\) by nonnegativity. All components of \(\beta\)
are nonzero at the finite places, since it is a totally positive
field element. Proposition~\ref{prop:top-norm-jacobian} at
\(b=-N\) identifies \(|L|\,d_T(\beta_T,\varphi)\) with
\(N^{-f}\int_{s(e)\in t+L}\Phi(e)\dd\delta_{-N}(e)\).
The remaining scalar \(\prod_\tau\tau(\beta)^{k-1}\) contributes
\(N^{-(k-1)f}\prod_\tau(|\tau t|^2-N^2)^{k-1}\).
Combining both powers of \(N\) gives exactly \(c_*\).
Use \(\Phi\le G\) and the exterior product bound in
Theorem~\ref{thm:localized-class-equation}.
\end{proof}

\section{The normalization of the orbital packets}

\begin{definition}\label{def:packet-amounts}
Retain \(k,r,\nu\), put \(h=\nu+1,\ c=r-k=2\nu+1,\ d_*=\nu(k+\nu)\),
and use \(U_j\) of Definition~\ref{def:row-vandermonde}. Write
\[
 \begin{gathered}
  C=|\mathcal C(K_c)|,\quad J_c=\binom{a+k-1}{k-1}^{U},
  \quad u=g^{(2d_*+\nu)f},\quad
  S_0=|\bar S|,\quad C_{\rm op}=\frac{U_k}{U_r}
                            \frac{2\pi^k}{(k-1)!},\\
  P_{\rm pkt}= S_0\,u\, C\,c_*,\qquad
  Z_{\rm sc}=\prod_{\tau}\prod_{a'=0}^\nu
        (s_{a',\tau}^2)^{k+2(\nu-a')},\qquad
  Z_{\rm pkt}=\frac{C^2}{\mathfrak m_r}
              C_{\rm op}^{\,f}V_E^h(\mathfrak v Z_{\rm sc})^2.
 \end{gathered}
\]
Here \(c_*\) is the constant of Theorem~\ref{thm:top-siegel-bound},
\(\mathfrak v,s_{a',\tau}\) are from
Proposition~\ref{prop:unipotent-covol}, and \(S_0\) is computed in
Proposition~\ref{prop:marked-local-level}.
For each integer \(j\ge3\) put
\(P_j=\prod_v(1-\sigma_{v,j}Q_v^{-j})\) with
\(\sigma_{v,j}=1\) at split places and \((-1)^j\) at inert ones.
\end{definition}

\begin{theorem}\label{thm:packet-normalization}
The two packet constants of Definition~\ref{def:packet-amounts}
satisfy the exact identity
\[
      \frac{P_{\rm pkt}}{Z_{\rm pkt}}=
                    (1-N^{-k})^{-f}\prod_{j=k+1}^{r}P_j^{-1}.
\]
For \(k\ge3\) the absolute value of its logarithm is at most
\(2 f N^{-k}+4 f\,2^{-k}\). In particular, for any fixed \(D>0\)
and all sufficiently large arithmetic scales,
\(P_{\rm pkt}\le(1+b^{-D})Z_{\rm pkt}\), uniformly throughout the
linear window of Theorem~\ref{thm:mean-covolume}.
\end{theorem}
\begin{proof}
Write \(H_0=|\bar H|\). Propositions~\ref{prop:column-dim}
and~\ref{prop:marked-local-level} give
\(C=H_0(N^k-1)^f\mathfrak m_k,\ S_0=H_0u\).
The omitted last scale in the product of
Proposition~\ref{prop:unipotent-covol} contributes
\((s_{\nu,\tau}^2)^k=N^{-k}\) at every place. Hence
\(\mathfrak v Z_{\rm sc}=uN^{-kf}V_E^{d_*}V_F^\nu\), using
Lemma~\ref{lem:real-volume-compact}. Substitution gives
\[
 \frac{P_{\rm pkt}}{Z_{\rm pkt}}
 =(1-N^{-k})^{-f}\frac{\mathfrak m_r}{\mathfrak m_k}
      \left(\frac{U_r}{2U_k}\right)^f
        V_F^{1-2\nu}V_E^{-(k+h+2d_*)}.
\]
The factor \(2\) here comes from the value \(2\pi^k/(k-1)!\) in
\(C_{\rm op}\): differentiating the real ball volume with respect to
radius at radius one gives twice the density with respect to
squared radius in Proposition~\ref{prop:norm-arch-factor}.

Set \(J=\sum_{l=k}^{r-1}l=c(k+\nu)=2d_*+k+\nu\), by summing the
arithmetic progression. By Corollary~\ref{cor:actual-mass-ratio},
\[
 \frac{\mathfrak m_r}{\mathfrak m_k}
       \left(\frac{U_r}{U_k}\right)^f
  =\big({\rm rd}_E^{\,J+c/2}2^{-J}\big)^f
        \prod_{j=k+1}^r P_j^{-1}.
\]
Indeed the factor from \(U_{l+1}/U_l\) is
\(2\pi^{l+1}/l!\); it cancels the factorial in the mass quotient
at each \(l\). By Lemmas~\ref{lem:skew-covol}
and~\ref{lem:real-volume-compact},
\(V_E=({\rm rd}_E/2)^f,\ V_F={\rm rd}_E^{f/2}\).
Together with \(k+h+2d_*=J+1,\ c=2\nu+1\), these equalities
cancel the powers of \(2\) and of \({\rm rd}_E\) exactly.

Corollary~\ref{cor:actual-mass-ratio} gives
\(\sum_{j=k+1}^r|\log P_j|\le4 f\sum_{j=k+1}^\infty2^{-j}
=4f\,2^{-k}\). Also
\(-\log(1-N^{-k})\le2N^{-k}\) by integration of \(1/(1-v)\)
on \(0\le v\le N^{-k}\le1/2\). These estimates prove the bound.
Proposition~\ref{prop:thick} gives \(k=r+1-m\ge x^5-2\sqrt x\);
\(f\le x^{12}\) eventually by Proposition~\ref{prop:level-prime}.
Thus the bound is \(o(b^{-D})\), as in
Corollary~\ref{cor:off-places-norm}. This proves the last assertion.
\end{proof}

\section{Comparison of the two numerical energies}

\begin{definition}\label{def:two-energy-numbers}
For the quantities of Definitions~\ref{def:packet-amounts}
and~\ref{def:row-vandermonde} put
\[
 E_1=S_0 C J_c,\qquad E_2=\frac{(C J_c)^2}{\mathcal D\mathfrak m_r},
 \qquad T_0=b^{kf/4},\qquad Z_0=2(b/g)^{kf}T_0^2 .
\]
The ordinal subscripts in this section distinguish two real numbers.
The column dimension here is \(C J_c\) by
Proposition~\ref{prop:column-dim}. The positive integer \(\mathcal D\)
is defined from the row parameters.
\end{definition}

\begin{lemma}\label{lem:logarithmic-counts}
Uniformly for \(t\le v\le t^2\) in the arithmetic construction,
\[
\begin{split}
 f^{-1}\log\mathcal D
 &=m(k-1)\bar w-\sum_{l=k-1}^{r-1}\log(l!)
             +O(mk/x+m^2\log x),\\
 \frac{\log J_c}{fk\ell}&=\frac1{\varphi(M)}+o(1).
\end{split}
\]
For every fixed \(D>0\) and at all sufficiently large \(x\), the
constructed primes satisfy
\(J_c\le b^{2kf}\), \( (N^k-1)^f\ge(N^k/b)^f\), and
\[
 \prod_\tau C_{\rm B}\,b^{7k}\rho_\tau^{\,k-h^7}
                        \le \mathcal D^{-1}b^{-D}.
\]
\end{lemma}
\begin{proof}
In the exact Vandermonde of Proposition~\ref{prop:exact-row-degree}
the first \(k-1\) rows cancel the factorials with indices
\(0,\ldots,k-2\). Each of the \(m(k-1)\) mixed differences
lies between \(w(\tau)\) and \(w(\tau)(1+1/x)\) eventually by
the mixed-factor estimate in its proof. The \(m(m-1)/2\) remaining
differences have logarithms between zero and \(\log(m\Delta_w)=
O(\log x)\). Taking logarithms and averaging proves the first
formula with both error bounds.

For \(J_c\) write
\[
 \log\binom{a+k-1}{k-1}
    =(k-1)\log a-\log((k-1)!)
                      +\sum_{i=1}^{k-1}\log(1+i/a).
\]
The sum is between zero and \(k^2/a\). Also
\(\log((k-1)!)\le k\log k\); we have
\(xb\le a\le2xb,\ k=O(x^5),\ \log x=o(\ell)\)
by Propositions~\ref{prop:level-prime} and~\ref{prop:thick}.
Now use \(U/f=1/\varphi(M)\) from Proposition~\ref{prop:places}.
This proves the second formula and \(J_c\le b^{2kf}\) eventually.
Since \(N\ge2\) and \(1/b\to0\), also \(1-N^{-k}\ge1/b\).
Finally Theorem~\ref{thm:orbital-degree-budget} with \(\eta=1\)
absorbs \(b^{7kf+D}\), since \((25/2-1-7)kf>D\) eventually.
\end{proof}

\begin{theorem}\label{thm:two-energies}
For any auxiliary primes satisfying the ranges of
Theorem~\ref{thm:data} and Proposition~\ref{prop:level-prime},
uniformly for \(t\le v\le t^2\),
\[
 \frac{\log(E_2/E_1)}{kf\ell}
     =\frac{\log g}{\ell}-\frac{15}{2}+\frac{1}{\varphi(M)}+o(1).
\]
For the primes produced in the proof of Theorem~\ref{thm:data}, this
expression tends to \(25/2+1/\varphi(M)\). For every allowed choice
in those ranges, for all sufficiently large \(x\) one has
\[
          T_0^{44} E_1\le E_2,\qquad
          T_0^{18} Z_0 E_2\le E_1,\qquad Z_0\le1 .
\]
\end{theorem}
\begin{proof}
Use \(c=m-1=2\nu+1,\ J=\sum_{l=k}^{r-1}l=c(k+\nu)\).
The two class and subgroup identities in the proof of
Theorem~\ref{thm:packet-normalization} give
\[
  E_2/E_1=(N^k-1)^f\,\mathfrak m_k J_c/
                                  (u\mathcal D\mathfrak m_r).
\]
Taking the product of the mass ratios of
Corollary~\ref{cor:actual-mass-ratio} over \(k+1,\ldots,r\) gives
\[
 \frac1f\log\frac{\mathfrak m_r}{\mathfrak m_k}
  =(J+c/2)\log{\rm rd}_E+\sum_{l=k}^{r-1}\log(l!)
                   -(J+c)\log(2\pi)+O(2^{-k}).
\]
The sum of the absolute logarithmic errors after division by \(f\)
is at most \(4\,2^{-k}\) by Theorem~\ref{thm:packet-normalization}.
Use \(\log u/f=(J-k)\log g\). From
Lemma~\ref{lem:logarithmic-counts}, writing
\(G_0=\log(g\,{\rm rd}_E/(2\pi))\) as in
Proposition~\ref{prop:level-prime}, the logarithm of \(E_2/E_1\)
divided by \(f\) is
\[
\begin{split}
 &k\log N-m(k-1)\bar w-J G_0+k\log g\\
 &\quad+c\big(\log(2\pi)-\tfrac12\log{\rm rd}_E\big)
                +\log((k-1)!)+f^{-1}\log J_c\\
 &\quad+O(mk/x+m^2\log x+2^{-k}).
\end{split}
\]
Here \(|\log(1-N^{-k})|\le2N^{-k}\). The same level-prime
proposition gives \(\log N=m\bar w+cG_0-15\ell/2+O(\sqrt\ell)\),
\(G_0=O(\ell)\), and \(\bar w=O(\ell)\); also
\(\log{\rm rd}_E=O(\log x)\) by Lemma~\ref{lem:rd} and
Proposition~\ref{prop:thick}. Since \(J=ck+c\nu\), the leftover
terms \(m\bar w,\ c\nu G_0,\ c\log{\rm rd}_E\) are \(o(k\ell)\)
by \(m=O(\sqrt x),\ k\ge x^5-2\sqrt x\).
The factorial contributes at most \(k\log k=o(k\ell)\).
The displayed errors and \(k\sqrt\ell\) are also \(o(k\ell)\).
The limit for \(J_c\) in Lemma~\ref{lem:logarithmic-counts}
now gives the asserted expansion.

In the proof of Theorem~\ref{thm:data} each \(\log g_i\) is
\(10\ell+O_d(\ell^{1/6})\), since only the fixed finite
number \(H_d+1\) of intervals of width \(\ell^{1/6}\) are used.
Hence \(\log g/\ell\to20\) for that choice. For arbitrary primes
in the stated intervals \(\log g/\ell\ge20\).
Consequently the first ratio has normalized logarithm at least
\(25/2+o(1)>11\), proving \(E_2\ge T_0^{44}E_1\).
The normalized logarithm of \(T_0^{18}Z_0 E_2/E_1\) equals
\(-3/2+1/\varphi(M)+o(1)\), since
\(\log Z_0/(kf\ell)=3/2-\log g/\ell+o(1)\).
It is negative eventually as \(\varphi(M)\ge1\).
The expression for \(Z_0\) also gives \(Z_0\le1\).
\end{proof}

\section{Hyperbolic coordinates on the finite norm fiber}

\begin{lemma}\label{lem:middle-hyperbolic-chart}
Use the finite set \(T\), the scalar and vector measures and the middle
space of Definitions~\ref{def:finite-fiber} and~\ref{def:top-norm-fiber}.
Write \(A_T=\prod_{v\in T} A_v,\ B_T=\prod_{v\in T} B_v\) for the
integral rings and \(E_T^-=iF_T\) for the skew elements. Give \(iA_T\)
mass one in \(E_T^-\). Put \(\xi=\exp(2\pi i/12)\) and
\(\eta=\xi/i\). The integral middle space has an orthogonal sum
decomposition into the hyperbolic plane spanned by \(v,v'\), with
\(v'=-z+N\eta v\) and \(h(v,v')=1\), and the span of
\(x_1,\ldots,x_{k-1}\) with its standard form. For
\(e=uv+t v'+\sum_j y_jx_j\) put \(H(y)=\sum_j\bar y_j y_j\).
For \(t\in E_T^\times\), write
\[
 U(\gamma,t,y,\alpha)
      =\bar t^{-1}\eta(\gamma-H(y))+t\alpha ,
 \qquad \gamma\in F_T,\quad \alpha\in E_T^- .
\]
For each fixed \(t,y\) this gives an affine bijection
\((\gamma,\alpha)\mapsto u\) carrying \(\dd\gamma\,\dd\alpha\)
to \(\dd u\); its first inverse coordinate is \(q(e)\).
\end{lemma}
\begin{proof}
Lemma~\ref{lem:skew-covol} gives \(\xi-\bar\xi=i\) and
\(O_E^-=iO_F\). In the basis \(1,\xi\) the same formulas hold after
completion. The elements with zero trace are precisely \(iA_T\)
in \(B_T\) and \(iF_T\) in \(E_T\), since division by the skew
unit \(i\) makes them fixed. The element \(\eta\in B_T\) has
\(\eta+\bar\eta=1\), so every \(b\in B_T\) has the unique form
\(b=\eta(b+\bar b)+b_-\) with \(b_-\in iA_T\).
The matrix on \(v,z\) in Proposition~\ref{prop:row-group} gives
\(h(v',v')=-N+N(\eta+\bar\eta)=0\). The change from \(v,z\)
to \(v,v'\) has unit determinant. Thus in the new coordinates
\[
              q(e)=\bar t u+t\bar u+H(y),\qquad
                        \dd e=\dd u\,\dd t\,\dd y .
\]
For fixed invertible \(t\), let \(n(t)=t\bar t\in F_T^\times\).
Multiplication of \(E_T\) by \(\bar t\) has additive volume factor
\(\prod_{v\in T}|n(t)_v|_v\). Indeed its determinant on the
two-dimensional algebra over each fixed field is \(n(t)_v\), by
the basis \(1,\xi\) of Lemma~\ref{lem:real-place-rings}; the
measure factor follows by the diagonal index calculation in
Lemma~\ref{lem:real-volume-compact}. For the proposed chart
\(\bar t u=\eta(\gamma-H(y))+n(t)\alpha\).
The trace splitting just proved identifies the two integral
lattices with \(B_T\), hence takes the product of their normalized
Haar measures to the normalized measure on \(E_T\), by counting
residue boxes. Scaling \(\alpha\) by \(n(t)\) contributes the
same factor as multiplication of \(u\) by \(\bar t\).
The factors cancel, and tracing this formula proves both
bijectivity and the assertion on the first coordinate.
\end{proof}

\begin{theorem}\label{thm:middle-fiber-chart}
Let \(k\ge3\), and let \(\delta_\gamma\) be the measure of
Proposition~\ref{prop:top-norm-jacobian} in the ranks of that
proposition. For every \(\gamma\in F_T\) and every nonnegative
Borel test \(G\) one has
\[
 \int G(e)\dd\delta_\gamma(e)=
 \int_{E_T^\times}\int_{E_T^{k-1}}\int_{E_T^-}
  G\!\left(U(\gamma,t,y,\alpha)v+t v'+\sum_j y_jx_j\right)
                 \dd\alpha\,\dd y\,\dd t .
\]
In particular the set where \(t\) is a nonunit of the algebra
\(E_T\) has \(\delta_\gamma\)-measure zero. Every \(E_T\)-linear
isometry of the middle space preserves \(\delta_\gamma\).
\end{theorem}
\begin{proof}
First take \(G\) compact, locally constant and nonnegative.
Choose integers \(a_v\ge0\) such that all coordinates \(u,t,y_j\)
on its support lie in \(D=\prod_{v\in T}\pi_v^{-a_v}B_v\).
There is a constant \(C_D\) such that for every \(c\in F_T\)
and integer \(n\ge1\),
\[
                \int_{D^{k-1}}{\bf1}_{H(y)-c\in L_n}\dd y
                              \le C_D |L_n| .
\]
Indeed, after \(y_v=\pi_v^{-a_v}y'_v\), the target at \(v\)
becomes \(\pi_v^{2a_v}c_v\) and the depth is \(n+2a_v\).
A nonintegral target gives zero. For an integral target
Theorem~\ref{thm:exact-local-density} in rank \(k-1\ge2\)
bounds the local quotient by
\(Q_v^{(2k-4)a_v}(1+2Q_v^{2-k})\) after division by \(Q_v^{-n}\).
The product of these finitely many bounds suffices.

For a positive integer \(B\), let \(D_B\subset D\) consist of
those \(t\) whose valuation in every field component of \(E_T\)
is at most \(B\), with valuation of zero taken as infinity.
This is compact open and consists of invertible elements.
The volume of \(D\setminus D_B\) tends to zero, since in each
field component a ball of depth \(B+1\) has volume its residue
cardinality to the power \(-B-1\), and there are finitely many
components. At fixed \(u,t\), apply the preceding bound with
\(c=\gamma-\bar t u-t\bar u\). The normalized mass of the part
of the tube with \(t\in D\setminus D_B\), weighted by \(G\), is
at most \(\sup G\,|D|C_D |D\setminus D_B|\) for every \(n\).

On \(D_B\) apply Lemma~\ref{lem:middle-hyperbolic-chart} at fixed
\(t,y\). The restricted tube integral divided by \(|L_n|\) is
the mean, for \(s\in\gamma+L_n\), of the chart integral in the
statement with \(t\in D_B\) and \(s\) in place of \(\gamma\).
For \(s\in\gamma+L_1\) its integrand is supported in a fixed
compact set of \(t,y,\alpha\): \(y\in D^{k-1}\), \(t\in D_B\),
and \(\alpha=t^{-1}(u-\bar t^{-1}\eta(s-H(y)))\) with \(u\in D\).
The inverses of \(D_B\) form a compact set by continuity.
On this compact set the integrand is locally constant in all
variables. A finite cover by product neighborhoods at
\(s=\gamma\) shows that every sufficiently small such mean is
its value at \(\gamma\). Hence the restricted measure is the
displayed chart integral over \(D_B\). Let \(B\) tend to infinity.
The uniform tube bound removes the omitted region on the left.
Monotone convergence gives the integral over invertible \(t\) on
the right. This proves the formula for compact open indicators
and proves local finiteness of the image measure.

Compact open boxes form a countable generating family.
Uniqueness on increasing compact shells as in
Proposition~\ref{prop:top-norm-jacobian} now proves equality of
the two Borel measures and the formula for all the asserted
tests. For an isometry \(g\) the identity
\(\overline{\det g}\det g=1\) follows by taking determinants of
the isometry relation in the nonsingular Gram matrix. The
additive Jacobian over \(F_T\) is consequently one: extend the
determinant calculation of Lemma~\ref{lem:middle-hyperbolic-chart}
to matrices by elementary elimination, whose elementary shears
have unit Jacobian. In split factors one performs the elimination
on the two components separately. Each tube around \(q=\gamma\)
is \(g\)-stable. Taking limits on compact open tests proves
invariance on boxes and hence on their generated Borel sets.
\end{proof}

\begin{corollary}\label{cor:one-vector-chart-bound}
In Theorem~\ref{thm:top-siegel-bound} the integral over the
slab \(s(e)\in t+L\) can be evaluated in the hyperbolic chart of
Theorem~\ref{thm:middle-fiber-chart} on the fiber with
\(\gamma=-N\), with exactly the factor \(c_*P_T(\beta)\) of that
theorem. The same chart formula holds with any integral trace-one
element in place of \(\eta\).
\end{corollary}
\begin{proof}
Apply Theorem~\ref{thm:middle-fiber-chart} to the product of \(G\)
and the indicator of the slab for Borel \(L\). Sets in its
completion have Borel representatives with null difference, using
the local equality of measures in
Proposition~\ref{prop:top-norm-jacobian} for the compact open
tests and exhaustion by such tests. The same conclusion also
follows by taking the completion of the equal measures in the
chart formula. For a second integral trace-one element the trace
splitting in Lemma~\ref{lem:middle-hyperbolic-chart} still holds,
as does the change to the corresponding integral hyperbolic
basis. The tube measure uses only \(q\) and the integral volume
on the ambient space. The preceding proof therefore gives the
same measure.
\end{proof}

\section{Integral characters from finite-field eigenvalues}

\begin{definition}\label{def:finite-character-ring}
For a finite group \(J\), a class function is a complex function
constant on conjugacy classes. Put
\(\langle F,H\rangle_J=|J|^{-1}\sum_{g\in J}\overline{F(g)}H(g)\).
Let \(R(J)\) be the integer span of the characters of finite
dimensional complex representations, called the virtual characters.
It is a ring for pointwise addition and multiplication, by direct
sums and tensor products. An elementary subgroup is a direct
product of a cyclic group of order prime to a prime \(\ell\) and
a group of order a power of \(\ell\), allowing trivial factors.
For \(H\le J\) and a class function \(\lambda\) on \(H\) define
\[
  ({\rm Ind}_H^J\lambda)(g)=
     |H|^{-1}\sum_{\substack{x\in J\\x^{-1}gx\in H}}\lambda(x^{-1}gx).
\]
\end{definition}

\begin{lemma}\label{lem:brauer-test-evaluation}
A class function \(F\) on a finite group \(J\) belongs to \(R(J)\)
if and only if its restriction to every elementary subgroup \(H\)
belongs to \(R(H)\).
\end{lemma}
\begin{proof}
Induction in Definition~\ref{def:finite-character-ring} sends a
character to a character. Given an \(H\)-space \(V\), take a copy
of \(V\) for each left coset \(xH\) with a chosen representative.
If \(gx=x'h\), let \(g\) carry the \(x\)-copy to the \(x'\)-copy
by \(h\). The group law follows by composing the two coset
decompositions. Its trace is the displayed sum, since precisely the
fixed cosets contribute. The same formula gives
\[
 F\,{\rm Ind}_H^J\lambda
       ={\rm Ind}_H^J\big((F|_H)\lambda\big).
\]
The restriction of a virtual character is virtual by restricting
the representations. It will therefore suffice to prove that the
ideal \(I\subset R(J)\) generated by \({\rm Ind}_H^J\lambda\)
for elementary \(H\) and one-dimensional characters \(\lambda\)
contains \(1\).

Write \(m=|J|,\ O=\Z[\exp(2\pi i/m)]\), and put \(S=O^b\) with
one coordinate for each of the \(b\) conjugacy classes. The
eigenvalues of every representing matrix of \(g\) are \(m\)-th
roots of unity by the order argument in
Proposition~\ref{prop:integral-proj}. Thus \(R(J)\subset S\).
The ring \(S\) is a finite module over \(\Z\), generated by
the first \(m\) powers of the chosen root in each coordinate.
These elements generate \(S\) over \(R(J)\) as well. If \(IS=S\),
express each of these generators as a combination of them with
coefficients in \(I\). The adjugate of the identity minus the
coefficient matrix shows that its determinant kills every generator
and hence \(1\). This determinant is \(1\) modulo \(I\), so
\(1\in I\).

If \(IS\) is proper, its projection to \(O\) is proper in at
least one coordinate \(g\), since multiplication by the coordinate
idempotents separates the components. Choose a maximal ideal of
\(O\) containing that proper ideal. For existence enumerate the
countable ring \(O\), adjoining successive elements to the ideal
whenever the resulting ideal is proper. The union is proper and
maximal. The quotient \(K'\) is a field generated over \(\Z\) as
a module by finitely many elements. Every element of this field
satisfies a monic polynomial over \(\Z\), by applying the same
adjugate argument to its multiplication on a list of module
generators. It has positive characteristic \(\ell\): in
characteristic zero, substituting \(1/2\) in a monic integer
polynomial and clearing denominators would give an odd integer
equal to zero. Since \(K'\) is a field, \(\ell\) is prime. In
this coordinate every generator of \(I\) must vanish modulo
the maximal ideal.

Write \(g=cw=wc\) with \(c\) of order prime to \(\ell\) and \(w\)
of \(\ell\)-power order. Such factors are powers of \(g\) by
the coprime congruences in Lemma~\ref{lem:basic-arith}. For every
character its values at \(g,c\) have the same reduction in \(K'\):
diagonalize \(g\), so \(c,w\) are diagonal in the same basis, and
every \(\ell\)-power root reduces to one by
\(X^{\ell^s}-1=(X-1)^{\ell^s}\) in characteristic \(\ell\).
The centralizer \(Z_J(c)\) has a subgroup \(P\) whose order is the
largest \(\ell\)-power \(\ell^a\) dividing \(|Z_J(c)|\). Indeed
let any group of order \(\ell^a u\), with \(\ell\nmid u\), act by
left translations on subsets of size \(\ell^a\). Their number
is congruent to \(u\) modulo \(\ell\), by the coefficient of
\(X^{\ell^a}\) in \((1+X^{\ell^a})^u\). Some orbit has size
prime to \(\ell\). Its stabilizer acts freely on the subset and
has order a multiple of \(\ell^a\) by counting cosets, so has
order exactly \(\ell^a\).

Take \(C=\langle c\rangle,\ H=CP=C\times P\). For any character
\(\lambda\) of \(C\) extend it trivially on \(P\). Counting the
preimages of each conjugate in the induction formula gives
\[
 ({\rm Ind}_H^J\lambda)(c)
 =\frac{|Z_J(c)|}{|H|}
       \sum_{z\in C\cap\{xcx^{-1}:x\in J\}}\lambda(z).
\]
Here an element in \(H\) conjugate to \(c\) has trivial
\(P\)-component by its order. The prefactor is an integer prime
to \(\ell\). The set summed over contains \(c\). Multiply these
sums by \(\lambda(c)^{-1}\) and sum over the \(|C|\) characters
of \(C\). By the geometric sum for powers of a cyclic generator
the result is \(|C|\), which is nonzero in \(K'\). Consequently
some induced character has nonzero reduction at \(c\) and at
\(g\). This contradicts properness of the projected ideal.
Thus \(1\in I\). Multiplying any expression for \(1\) in these
generators by \(F\), the displayed induction identity proves the
assertion whenever the restrictions are virtual.
\end{proof}

\begin{lemma}\label{lem:monomial-ell-group}
Let \(P\) be a finite group of prime-power order \(\ell^a\).
Over an algebraically closed field of characteristic different
from \(\ell\), every irreducible representation of \(P\) is
induced from a one-dimensional character of a subgroup.
\end{lemma}
\begin{proof}
Induct on \(|P|\). For a nontrivial kernel apply induction to
the quotient, and pull the inducing character back to the
preimage of its subgroup. Cosets and their actions in the two
induction constructions agree. For abelian \(P\) the commuting
matrices are simultaneously diagonalizable since each is
annihilated by \(X^{|P|}-1\) with distinct roots. To diagonalize
a commuting family, split into eigenspaces successively, all
of them invariant for the remaining matrices. An irreducible
space then has dimension one.

Now take a faithful irreducible representation of a nonabelian
\(P\). Every nontrivial \(\ell\)-group has nontrivial center,
since counting the conjugacy orbits modulo \(\ell\) shows that
its center has size divisible by \(\ell\). Apply this to
\(P/Z(P)\). Some noncentral \(a\) has all its conjugates in
\(a Z(P)\), so \(A=\langle a,Z(P)\rangle\) is abelian and normal.
It has at least two weight spaces, since in a faithful action an
element acting by a scalar is central. The group \(P\) permutes
these spaces transitively by irreducibility. A weight space is
irreducible for its stabilizer \(H\). Indeed, translating any of
its proper nonzero \(H\)-invariant subspaces over the weight-space
orbit would give a proper nonzero \(P\)-invariant direct sum.
The whole space, expressed as this orbit of weight spaces, is
the induced representation from \(H\), which is proper.
Induct inside \(H\). Induction from a subgroup of \(H\), followed
by induction from \(H\), is induction from the original subgroup,
since products of the two sets of coset representatives represent
its cosets in \(P\) and give the same action.
\end{proof}

\begin{theorem}\label{thm:brauer-power-sum}
Let \(J\) be a finite group, \(p\) a prime, and \(V\) a
finite-dimensional representation of \(J\) over an algebraic
closure of \(\mathbf F_p\). Let \(m\) be any positive integer prime
to \(p\) divisible by the prime-to-\(p\) part of \(|J|\).
Fix an isomorphism \(\iota\) of the group of \(m\)-th roots of
unity in this closure with the complex \(m\)-th roots of unity.
For each integer \(b\ge0\),
\[
               \beta_b(g)=\sum_{\lambda}\iota(\lambda)^b
                              \quad(g\in J)
\]
is a virtual character. The sum ranges over the eigenvalues of
\(g\) on \(V\) counted with their algebraic multiplicities.
\end{theorem}
\begin{proof}
All eigenvalues have order dividing \(m\), by the order of \(g\)
and \(X^{p^n}-1=(X-1)^{p^n}\) in characteristic \(p\).
Both groups of \(m\)-th roots have \(m\) elements by separability
and are cyclic by the elementary exponent proof in
Lemma~\ref{lem:count-labels}. For \(b=1\) take an elementary
subgroup \(H=C\times P\) as in
Definition~\ref{def:finite-character-ring}, with \(|P|\) an
\(\ell\)-power. In \(H\) the elements of \(p\)-power order form
a direct factor \(H_p\) commuting with a complement \(D\) of
order prime to \(p\). If \(\ell=p\), take \(D=C\); otherwise
split off the \(p\)-primary part of the cyclic factor \(C\).
For \(h=dy\) with \(d\in D,\ y\in H_p\), the eigenvalues of \(h\)
are those of \(d\) with the same multiplicities. In fact \(d\)
diagonalizes, and on each eigenspace \(y\) is unipotent and
commutes with \(d\).

Over the given field the restriction to \(D\) is a direct sum
of irreducibles. Average any linear projection onto an invariant
subspace over \(D\) to get an invariant complement, since
\(|D|\) is invertible. If \(D\) is cyclic all summands have
dimension one. Otherwise its central cyclic prime-to-\(p\) factor
acts by a single weight on each irreducible, and the remaining
\(\ell\)-group acts irreducibly there. By
Lemma~\ref{lem:monomial-ell-group}, each such summand has a coset
basis induced from a one-dimensional character, with the cyclic
factor acting by scalars. Lift the nonzero entries in this basis
by \(\iota\), preserving products and permutations. These entries
have order dividing \(m\), being values of characters of subgroups
of \(D\). This gives complex representations whose eigenvalues at
\(d\) are precisely the lifts of the old eigenvalues with multiplicity.
Indeed a permutation cycle of length \(e\) with product of its
entries equal to \(h_0\) has characteristic polynomial \(X^e-h_0\).
Here \(p\nmid e\) as the cycle length divides the order of \(d\).
The distinct roots lift bijectively to those of
\(X^e-\iota(h_0)\), since all the roots are eigenvalues of \(d\)
and \(\iota\) is injective. Inflating the direct sum along
\(H\to D\) realizes \(\beta_1|_H\) as a character.

For general \(b\) perform the power operation on the eigenvalues
of this complex representation of \(H\). If its elementary
symmetric functions are \(e_j\), they are the characters of the
exterior powers by diagonalizing each group element. For
\(b\ge1\) its \(b\)-th power sum is an integer polynomial in
these \(e_j\): equate coefficients of \(X^b\) in
\[
 \left(\sum_{j\ge0}(-1)^j e_j X^j\right)
           \left(\sum_{d\ge1}p_d X^d\right)
        =-X\frac{\dd}{\dd X}\prod_i(1-\lambda_i X),
 \qquad p_d=\sum_i\lambda_i^d .
\]
The identity follows by expanding each \(\lambda_i X/(1-\lambda_i X)\);
the constant coefficient of the first factor is \(1\), so the
recursion uses only integer operations. For \(b=0\) the function
is the dimension times the trivial character. In each case
\(\beta_b|_H\in R(H)\). Lemma~\ref{lem:brauer-test-evaluation}
proves virtuality on \(J\).
\end{proof}

\section{Parabolic induction of finite class functions}

\begin{definition}\label{def:parabolic-operations}
Fix a finite field \(K\) of size \(Q\) and characteristic \(p\),
and a rank bound \(R\ge1\). Set \(G_j={\rm GL}_j(K)\) for
\(0\le j\le R\), with \(G_0\) the trivial group. Consider
sequences \(F\) of class functions on the \(G_j\).
Write \(F(g)\) for the value on an automorphism of any such
\(K\)-space, using a basis. For \(g\) on \(V\) and \(a,b\) on
\(A,B\) respectively, define
\[
\begin{split}
 (F*H)(g)&=\sum_{\substack{W\subset V\\gW=W}}
                           F(g|W)H(g|V/W),\\
 (\Delta F)(a,b)&=\frac{1}{|\operatorname{Hom}_K(B,A)|}
           \sum_{u:B\to A}F\!\begin{pmatrix}a&u\\0&b\end{pmatrix}.
\end{split}
\]
In the first formula \(g|V/W\) denotes the induced quotient map.
The second formula is used when \(\dim A+\dim B\le R\).
For functions of two matrix arguments, the product \( * \) sums
independently over invariant subspaces of both spaces, restricting
in both arguments in the first factor and taking quotients in
both in the second. Put \(\mathbf e(g)={\bf1}_{\dim V=0}\)
and \(\mathbf u(g)=1\). A class function supported in positive
rank is cuspidal if \((\Delta F)(a,b)=0\) whenever both ranks are
positive. For functions of ranks \(i,j\), the pairing on two
matrix arguments is the normalized product-group pairing in
Definition~\ref{def:finite-character-ring}.
\end{definition}

\begin{lemma}\label{lem:flag-adjunction}
The product \( * \) is associative with unit \(\mathbf e\).
If \(F,H\) are supported in ranks \(i,j\), their product at rank
\(i+j\) is induction from the subgroup preserving an \(i\)-plane,
of the product class function on the plane and its quotient.
Consequently \( * \) preserves virtual characters in each rank and
\[
 \langle F*H,B'\rangle_{G_{i+j}}
    =\langle F\otimes H,\Delta B'\rangle_{G_i\times G_j}.
\]
\end{lemma}
\begin{proof}
Counting flags \(U\subset W\subset V\) stable under \(g\)
gives the same three factors on \(U,W/U,V/W\) in either
parenthesization. The unique zero subspace and the full space
give the unit assertions. All \(i\)-planes form one orbit of
\(G_{i+j}\) by choosing bases. Their stabilizer consists of the
block upper triangular matrices in
Definition~\ref{def:parabolic-operations}. The fixed-plane sum
for \(g\) is the trace formula in
Definition~\ref{def:finite-character-ring}, since its fixed cosets
are precisely the invariant planes. Ordinary induction preserves
virtual characters by Lemma~\ref{lem:brauer-test-evaluation}
and its proof. For any subgroup \(P\le J\) that trace formula
also gives
\(\langle{\rm Ind}_P^J f,h\rangle_J=\langle f,h|_P\rangle_P\):
interchange the sums and substitute \(g=x y x^{-1}\), \(y\in P\);
each \(x\) then gives the same sum. Averaging over the upper
right block in \(P\) first gives exactly \(\Delta B'\), as
required.
\end{proof}

\begin{proposition}\label{prop:block-mackey}
For the operations of Definition~\ref{def:parabolic-operations},
\[
             \Delta(F*H)=(\Delta F)*(\Delta H).
\]
\end{proposition}
\begin{proof}
Fix \(a,b\). Choose the upper right block \(u:B\to A\) uniformly.
For an invariant subspace \(W\subset A\oplus B\) put
\(X=W\cap A,\ Y=\operatorname{pr}_B W\). Then \(aX=X,\ bY=Y\).
At fixed \(X,Y\), possible subspaces before the stability
condition are parametrized by all maps \(l:Y\to A/X\), whose
graphs describe \(W/X\) in \((A/X)\oplus B\).
Extend \(l\) to a map \(B\to A\) by lifting in bases and taking
zero on a complement of \(Y\). Conjugating by the shear with
that off-diagonal block changes \(X\oplus Y\) to \(W\) and
translates \(u\) by a fixed map. Hence the average at \(W\),
including its stability indicator, equals the one for
\(W=X\oplus Y\).

Choose complements \(A=X\oplus A'\), \(B=Y\oplus B'\).
The condition is now \(u_{A',Y}=0\), which has probability
\(Q^{-\dim A'\dim Y}\). Conditional on it, the two blocks
\(u_{X,Y}\) and \(u_{A',B'}\) are independent uniform maps.
The diagonal actions on \(W\) are \(a|X,b|Y\), and those on
the quotient are \(a|A/X,b|B/Y\). The remaining block of \(u\)
does not enter the two induced actions. Thus their expected
product is the product of the corresponding two block averages.
The number \(Q^{\dim A'\dim Y}\) of possible graphs cancels the
stability probability. Summing over \(X,Y\) gives the assertion.
\end{proof}

\begin{definition}\label{def:cuspidal-virtual}
Take \(m\) to be the prime-to-\(p\) part of \(|G_R|\).
It is divisible by the prime-to-\(p\) part of \(|G_j|\) for
\(j\le R\), by embedding \(G_j\) on a direct summand and
counting cosets. Fix an embedding \(\iota\) of the \(m\)-th roots
of unity as in Theorem~\ref{thm:brauer-power-sum}. For any
multiplicative character \(\theta=\iota^b\) of these roots with
\(b\ge0\), put \(\beta_\theta(g)=\sum_\lambda\theta(\lambda)\)
over the eigenvalues with multiplicity. Define
\[
 \mu=\sum_{n=0}^{R}(-1)^n(\mathbf u-\mathbf e)^{*n},
 \qquad D_\theta F=\beta_\theta F,\qquad
                         C_\theta=(D_\theta\mu)*\mathbf u .
\]
Multiplication in \(D_\theta\) is pointwise. A rank part means
the sequence restricted to the specified rank and zero in all
other ranks.
\end{definition}

\begin{theorem}\label{thm:virtual-cuspidal}
Each positive rank part of \(C_\theta\) is a virtual character
and is cuspidal. On the ranks in
Definition~\ref{def:parabolic-operations} one has
\[
 \mu*\mathbf u=\mathbf u*\mu=\mathbf e,\qquad
 \Delta\mu=\mu\otimes\mu,\qquad
 \Delta C_\theta=C_\theta\otimes\mathbf e+
                         \mathbf e\otimes C_\theta .
\]
\end{theorem}
\begin{proof}
A product of \(R+1\) functions vanishing at rank zero vanishes
at every rank at most \(R\), by the stable flag description in
Lemma~\ref{lem:flag-adjunction}. The finite geometric sum for
\(\mu\) is therefore a two-sided inverse. Since
\(\Delta\mathbf u=\mathbf u\otimes\mathbf u\) and
\(\Delta\mathbf e=\mathbf e\otimes\mathbf e\),
Proposition~\ref{prop:block-mackey} shows that \(\Delta\mu\)
is the inverse of \(\mathbf u\otimes\mathbf u\) under the product
in two arguments, truncated at total rank \(R\). That inverse is
\(\mu\otimes\mu\) by the separated sums. On a triangular block
matrix \(\beta_\theta\) is the sum of its values on the diagonal
blocks, since its characteristic polynomial is their product.
Consequently
\[
 \Delta(D_\theta\mu)
    =(D_\theta\mu)\otimes\mu+\mu\otimes(D_\theta\mu).
\]
Multiply by \(\Delta\mathbf u\) and apply
Proposition~\ref{prop:block-mackey} once more. This proves the
formula for \(\Delta C_\theta\), its value in positive bidegrees,
and hence cuspidality. Its value in rank zero is zero since an
empty matrix has no eigenvalues.

Both \(\mathbf e\) and \(\mathbf u\) are characters or zero in
each rank. The formula defining \(\mu\) and
Lemma~\ref{lem:flag-adjunction} give virtuality of \(\mu\).
Theorem~\ref{thm:brauer-power-sum} gives the same for
\(\beta_\theta\), so pointwise multiplication followed by
parabolic induction proves virtuality of \(C_\theta\).
\end{proof}

\section{The torus calculation for cuspidal characters}

\begin{definition}\label{def:finite-torus-label}
Retain \(K,Q,m,\iota\) and \(G_n\) from
Definitions~\ref{def:parabolic-operations} and~\ref{def:cuspidal-virtual}.
Let \(K_n=\{a\in\overline K:a^{Q^n}=a\}\) for \(1\le n\le R\).
Embed \(T_n=K_n^\times\) in \(G_n\) by multiplication in a
\(K\)-basis of \(K_n\). For \(x\in T_n\) define \(b_x\) in
\(0\le b_x<Q^n-1\) by
\(\iota(x)=\exp(2\pi i b_x/(Q^n-1))\). Put
\(\theta_x=\iota^{b_x}\) on the \(m\)-th roots of unity.
For a monic irreducible \(P\in K[X]\) of degree \(n\) with
\(P(0)\ne0\), choose a root \(x\) and put
\(\tau_P=(-1)^n C_{\theta_x}\) in rank \(n\), extended by zero
in all other ranks.
\end{definition}

\begin{lemma}\label{lem:torus-primary}
The set \(K_n\) is a field of size \(Q^n\) containing \(K\).
A monic irreducible of degree \(e\) has the distinct roots
\(z,z^Q,\ldots,z^{Q^{e-1}}\) in \(K_e\). One has \(K_e\subset K_n\)
if \(e\mid n\), and \(Q^n-1\mid m\).
On \(G_n\), \(C_\theta(g)=0\) whenever the characteristic
polynomial of \(g\) has two distinct irreducible factors.
Also \(\mu(g)=0\) unless \(g\) is diagonalizable over \(\overline K\).
\end{lemma}
\begin{proof}
The roots of \(X^{Q^n}-X\) form a field by the Frobenius power
identities, and the derivative is \(-1\), as in
Lemma~\ref{lem:count-labels}. The elements of \(K\) are fixed by
\(Q\)-powering since their nonzero elements form a group of order
\(Q-1\). Thus the fixed field of \(Q\)-powering in the closure is
exactly \(K\) by degree. For a root \(z\) of degree \(e\) the
\(K\)-algebra \(K[z]\) is a field of size \(Q^e\), by the basis
\(1,z,\ldots,z^{e-1}\) from polynomial division. Its elements
satisfy \(a^{Q^e}=a\). The orbit of \(z\) under \(Q\)-powering
thus has size at most \(e\); its orbit polynomial has coefficients
fixed by \(Q\)-powering and is divisible by the minimal polynomial.
It follows that the orbit has size \(e\). Iterating \(Q^e\)-powering
gives \(K_e\subset K_n\) when \(e\mid n\). The embedding of the
group \(T_n\) of order \(Q^n-1\) into \(G_n\) gives the asserted
divisibility by the definition of \(m\). This justifies the choices
in Definition~\ref{def:finite-torus-label}.

If there are at least two primary factors, split \(V=A\oplus B\)
over \(K\) so that the characteristic polynomials of the actions
\(a,b\) on the two nonzero summands are coprime. Explicitly,
factor the characteristic polynomial as \(D_1D_2\) with both
nonconstant, coprime, and with disjoint root sets. The adjugate
identity for \(XI-g\), whose adjugate commutes with \(g\)
coefficientwise, gives \(D_1(g)D_2(g)=0\) by substitution.
B\'ezout polynomials for \(D_1,D_2\), obtained by polynomial
division, split \(V\) as the sum of the two kernels, with
zero intersection. Both kernels are nonzero, as seen after
triangularizing over the algebraic closure by successively
taking an eigenvector in the quotient. On maps \(l:B\to A\)
the operator \(l\mapsto al-lb\) is invertible. Its kernel is
zero, since \(al=lb\) implies \(f(a)l=lf(b)\) and one may
choose \(f\) equal to one modulo the polynomial on \(A\) and
zero modulo that on \(B\). Thus every upper right block is
obtained by conjugating \(a\oplus b\) by an upper triangular
shear. Its \(C_\theta\)-value equals the block average, zero by
Theorem~\ref{thm:virtual-cuspidal}.

For the last assertion, let \(D\) be the product of the distinct
irreducible factors and set \(S=\ker D(g)\). It is nonzero in
positive rank by the eigenvector argument. Since \(D\) splits
with distinct roots over \(\overline K\), \(S=V\) precisely when
\(g\) diagonalizes there: interpolation in the roots gives the
eigenspace projections when \(D(g)=0\). Every invariant subspace
with diagonalizable restriction lies in \(S\). If \(S\ne V\),
induct on dimension in
\(\sum_{W\subseteq V,\,gW=W}\mu(g|W)=0\), which is
\(\mu*\mathbf u=\mathbf e\) in Theorem~\ref{thm:virtual-cuspidal}.
All proper subspaces contributing to this sum lie in \(S\).
Their sum over \(W\subseteq S\) is zero by the same identity
on \(S\ne0\). This forces \(\mu(g)=0\).
\end{proof}

\begin{proposition}\label{prop:torus-induction}
For \(1\le n\le R\) and every \(\theta\) of
Definition~\ref{def:cuspidal-virtual}, the pointwise product
on \(G_n\) satisfies
\[
                    \mu C_\theta={\rm Ind}_{T_n}^{G_n}(\theta|_{T_n}).
\]
At the identity, \((-1)^n C_\theta(1)=\prod_{j=1}^{n-1}(Q^j-1)\).
\end{proposition}
\begin{proof}
Multiplication by \(z\in K_n^\times\) is diagonalizable over
\(\overline K\), with minimal polynomial that of \(z\) over \(K\).
Its characteristic polynomial is a power of that polynomial:
on a basis over \(K[z]\) it is the direct sum of copies of
multiplication by \(z\) on \(K[z]\). Conversely, a diagonalizable
matrix \(g\) with characteristic polynomial \(P^s\), where
\(e=\deg P,\ n=es\), makes \(V\) into a vector space of dimension
\(s\) over \(K[z]\) through \(K[X]/P\) acting by \(X\mapsto g\).
This uses \(P(g)=0\) by diagonalizability. Choosing a basis over
this field identifies \(g\) with multiplication by any root \(z\)
of \(P\) on \(K_n\), by Lemma~\ref{lem:torus-primary}. Its
centralizer consists of the invertible \(K[z]\)-linear maps.
The elements of \(T_n\) conjugate to \(g\) correspond exactly
to the roots of \(P\). Thus on every other class both sides of
the asserted identity vanish, by Lemma~\ref{lem:torus-primary}
and the induction formula.

Write \(q'=Q^e\) and \(S_\theta(z)=\sum_{h=0}^{e-1}\theta(z^{Q^h})\).
On the current class the stable subspaces are exactly the
\(K[z]\)-subspaces. If \(b_{s,j}\) counts the \(j\)-dimensional
ones in a space of dimension \(s\) over a field of size \(q'\),
then
\[
        \sum_{j=0}^s b_{s,j}(-1)^j(q')^{j(j-1)/2} X^j
                          =\prod_{j=0}^{s-1}(1-(q')^j X).
\]
Indeed, relative to a fixed hyperplane the subspaces inside it
contribute \(b_{s-1,j}\), and the others contribute
\((q')^{s-j} b_{s-1,j-1}\). For the latter choose the intersection
of dimension \(j-1\), and lift a vector with last coordinate
one modulo that intersection. With out-of-range counts zero
this recurrence gives multiplication of the polynomial for
\(s-1\) by \(1-(q')^{s-1}X\), starting at \(1\).
Setting \(X=1\) and using \(\mu*\mathbf u=\mathbf e\) recursively
shows that
\(\mu(g)=(-1)^s(q')^{s(s-1)/2}\). On a \(j\)-dimensional
\(K[z]\)-subspace the eigenvalue sum is \(j S_\theta(z)\).
Differentiating the polynomial at \(X=1\) therefore gives
\[
             C_\theta(g)=-S_\theta(z)\prod_{j=1}^{s-1}(1-(q')^j).
\]
The identity case has \(z=1,\ e=1,\ s=n\), giving the degree.

Counting ordered bases gives
\(|{\rm GL}_s(K[z])|=(q')^{s(s-1)/2}\prod_{j=1}^s((q')^j-1)\)
as in Proposition~\ref{prop:row-group}. The product
\(\mu(g) C_\theta(g)\) is therefore
\(|{\rm GL}_s(K[z])|S_\theta(z)/(Q^n-1)\).
Each conjugate in \(T_n\) has \(|{\rm GL}_s(K[z])|\)
conjugators, by the centralizer just computed. This is exactly
the value of the induced character.
\end{proof}

\begin{theorem}\label{thm:individual-green}
The function \(\tau_P\) of Definition~\ref{def:finite-torus-label}
is independent of the root chosen. It is the character of an
irreducible cuspidal representation of \(G_n\) of degree
\(\prod_{j=1}^{n-1}(Q^j-1)\). Different degree-\(n\) monic
irreducibles with nonzero constant terms give inequivalent
representations.
\end{theorem}
\begin{proof}
For two lifts \(\theta,\eta\), apply
Lemma~\ref{lem:flag-adjunction} to
\(C_\eta=(D_\eta\mu)*\mathbf u\) in the second argument of the
pairing with \(C_\theta\), taking complex conjugates of the
adjunction if needed. Every term involving two positive ranks is
zero by Theorem~\ref{thm:virtual-cuspidal}. The zero first rank
term also vanishes since \(\beta_\eta\) is zero there. The surviving
term is \(\langle C_\theta,\mu\beta_\eta\rangle_{G_n}\).
Since \(\mu\) has integer values by its flag sum, moving it to the
first slot gives \(\langle\mu C_\theta,\beta_\eta\rangle_{G_n}\).
By Proposition~\ref{prop:torus-induction} and the subgroup
adjunction in Lemma~\ref{lem:flag-adjunction}, we obtain
\[
 \langle C_\theta,C_\eta\rangle_{G_n}
    =\frac1{|T_n|}\sum_{z\in T_n}\overline{\theta(z)}
                                   \sum_{h=0}^{n-1}\eta(z^{Q^h})
    =\#\{0\le h<n:\theta|_{T_n}=(\eta|_{T_n})^{Q^h}\}.
\]
The inner sum lists all eigenvalues of multiplication by \(z\)
with their multiplicities, even when \(z\) belongs to a proper
subfield, by Lemma~\ref{lem:torus-primary}. The last equality is
the finite geometric sum on \(T_n\).

The map \(x\mapsto\theta_x|_{T_n}\) bijects \(T_n\) with its
complex characters and carries \(Q\)-powering to \(Q\)-powering,
since both sides use the same exponents modulo \(Q^n-1\).
For roots \(x,y\) of degree \(n\), the count in the last display
is consequently \(1\) if their minimal polynomials agree and \(0\)
otherwise. In particular the norm squared for \(x\) is one.
By Theorem~\ref{thm:virtual-cuspidal} the signed function is an
integer linear combination of irreducible characters. By the
orthogonality proved in Lemma~\ref{lem:central-proj}, its norm
squared is the sum of the squares of those integer coefficients.
Exactly one is \(1\) or \(-1\), and the degree computed in
Proposition~\ref{prop:torus-induction} makes the sign positive.
The inner product between two root choices now proves independence
and inequivalence as asserted. Cuspidality follows from
Theorem~\ref{thm:virtual-cuspidal}.
\end{proof}

\section{Squarefree labels and the marked projection}

\begin{definition}\label{def:squarefree-green}
Let \(c\in K[X]\) be squarefree, monic of degree \(r\le R\) with
\(c(0)\ne0\). Write it as a product of distinct monic irreducibles
\(P_1,\ldots,P_h\), and put
\(\Xi_c=\tau_{P_1}*\cdots *\tau_{P_h}\) in rank \(r\).
The empty product in rank zero is \(\mathbf e\).
Factorization exists by successively factoring reducible polynomials
of smaller degree. Euclidean division and B\'ezout's identity give
Euclid's lemma for an irreducible divisor, and hence uniqueness of
the set of monic factors.
\end{definition}

\begin{theorem}\label{thm:squarefree-green}
The character \(\Xi_c\) is irreducible, independent of the order
of the factors, and distinct for distinct labels \(c\) of the same
degree. If \(n_P=\deg P\) its degree and its share in the regular
representation satisfy
\[
\begin{split}
 d_c=\Xi_c(1)&=\frac{\prod_{j=1}^r(Q^j-1)}
                         {\prod_{P\mid c}(Q^{n_P}-1)},\\
 \frac{d_c^2}{|G_r|}
 &= Q^{-r}\frac{\prod_{j=1}^r(1-Q^{-j})}
                     {\prod_{P\mid c}(1-Q^{-n_P})^2}
 \ \ge\ Q^{-r}\left(1-\sum_{j=1}^r Q^{-j}\right).
\end{split}
\]
The products indexed by \(P\) use only the monic irreducible
divisors of \(c\).
\end{theorem}
\begin{proof}
Each \(\tau_P\) is a character by Theorem~\ref{thm:individual-green}.
Repeated induction in Lemma~\ref{lem:flag-adjunction} realizes
\(\Xi_c\) on a complex representation. For another list of
distinct irreducibles forming \(c'\), pair its product with the
product for \(c\). Repeatedly apply the adjunction of
Lemma~\ref{lem:flag-adjunction} to the latter. In the resulting
\(h\) matrix slots each factor belonging to \(c'\) is assigned
entirely to one slot. Indeed \(\Delta\tau_{P'}=
\tau_{P'}\otimes\mathbf e+\mathbf e\otimes\tau_{P'}\)
by Theorem~\ref{thm:virtual-cuspidal}. Applying
Proposition~\ref{prop:block-mackey} and iterating gives a sum over
all such assignments; the factors assigned to a slot are multiplied
there in their original order. The other side in slot \(i\) is
\(\tau_{P_i}\), of positive rank. Empty slots therefore give zero.
If a slot receives two or more factors, the pairing in it vanishes
by splitting off its first factor with the same adjunction and
using cuspidality. For the remaining assignments each slot receives
exactly one factor. The product of pairings in these slots is
one if every polynomial matches and zero otherwise by
Theorem~\ref{thm:individual-green}. Distinctness of the factors
gives exactly one matching assignment for equal sets and none for
unequal sets. The rank zero convention agrees with these conclusions.
By Lemma~\ref{lem:central-proj}, the inner product of a character
with itself is the sum of the squares of its irreducible
multiplicities. It is one here. The other pairings give both
independence of the order and inequivalence.

Let \(n_i=\deg P_i\). The stabilizer of a flag with successive
dimensions \(n_i\) has order \(Q^{\sum_{i<j}n_i n_j}\prod_i |G_{n_i}|\):
the diagonal blocks are arbitrary invertible matrices and the
strictly upper blocks are free. Transitivity on flags follows by
choosing a basis successively. Each flag contributes the product
of the dimensions from Theorem~\ref{thm:individual-green} at the
identity. Multiplying by the flag count and substituting
\(|G_j|=Q^{j(j-1)/2}\prod_{\nu=1}^j(Q^\nu-1)\) from the proof of
Proposition~\ref{prop:torus-induction} gives the asserted degree,
since
\(\binom r2=\sum_i\binom{n_i}2+\sum_{i<j}n_i n_j\).
Squaring and dividing by \(|G_r|\) gives the second identity.
Its denominator product is at most one. Multiplying the bounds
\(0\le1-Q^{-j}\le1\) inductively gives
\(\prod(1-Q^{-j})\ge1-\sum Q^{-j}\).
\end{proof}

\begin{corollary}\label{cor:actual-good-share}
For \(p=g_i\) in Definition~\ref{def:poly-types}, the characters
\(\Xi_c\) supply the representations of
Proposition~\ref{prop:bad-types}. Let \(e_i\) denote its projection
onto the good tuples at \(p\), formed from these representations.
For \(2\le r<p\) one has the sharper estimate
\[
  0\le [1](1-e_i)\le
   f\left(\frac{A+1}{p}+\frac{2}{p-1}\right)
             +\exp\!\left(-\frac{f}{2r(2r\log p+1)}\right).
\]
On the entire group \(\mathcal G\) of
Proposition~\ref{prop:row-group}, let \(e_0^{\rm good}\) be the
sum of the projections for the tuples good over both primes. It is
a central self-adjoint idempotent. For all sufficiently large
arithmetic scales, \(0\le[1](1-e_0^{\rm good})\le b^{-8}\).
\end{corollary}
\begin{proof}
Take \(K=\mathbf F_p,\ R=r\) in Theorem~\ref{thm:squarefree-green}.
Since \(\sum_{j=1}^r p^{-j}\le r/p\), its degree bound implies the
hypothesis of Proposition~\ref{prop:bad-types}. To obtain the
sharper estimate, use \(\sum_{j=1}^r p^{-j}\le1/(p-1)\).
For a uniformly chosen degree-\(r\) monic polynomial, the
probability of a zero or a forbidden root is at most \((A+1)/p\).
The chance of divisibility by the square of some monic positive
degree polynomial is at most
\(\sum_{j=1}^{\lfloor r/2\rfloor}p^{-j}\le1/(p-1)\), by the
count in Lemma~\ref{lem:count-labels}. Bound the probability
that an \(f\)-tuple fails either test in any slot by \(f\) times
their sum. Failure to contain a primitive member has probability
at most the exponential term, as in Proposition~\ref{prop:bad-types}.
On the retained tuples the product of the degree shares is at least
\(p^{-rf}(1-1/(p-1))^f\), and
\(1-(1-1/(p-1))^f\le f/(p-1)\). This gives the displayed bound.
Finally the two-prime projection is the tensor product of the
two one-prime projections by distributivity. Their identity
coefficients \(t_0,t_1\) lie in \([0,1]\) by
Lemma~\ref{lem:central-proj} and the normalized regular trace.
Its complementary share is \(1-t_0t_1\le(1-t_0)+(1-t_1)\).
Each summand is at most \(b^{-8}/2\) eventually by
Proposition~\ref{prop:bad-types}, now with its hypotheses proved.
\end{proof}

\section{Determinant twists and hyperplane averages}

\begin{proposition}\label{prop:green-named-twist}
Let \(c\) be a label of degree \(r\) in
Definition~\ref{def:squarefree-green}. For \(u\in K^\times\)
write \(c^{[u]}(X)=u^r c(X/u)\) and
\(\chi_u=\theta_u|_{K^\times}\), using the exponent at rank one
in Definition~\ref{def:finite-torus-label}. The map
\(u\mapsto\chi_u\) is a group isomorphism from \(K^\times\)
onto its complex characters, and
\[
                   \Xi_{c^{[u]}}(g)=\chi_u(\det g)\,\Xi_c(g).
\]
In rank one \(\Xi_{X-u}=\chi_u\).
\end{proposition}
\begin{proof}
The exponent on \(K^\times=T_1\) identifies the cyclic group and
its dual as in Theorem~\ref{thm:individual-green}. The formula
for \(C_\theta\) in Proposition~\ref{prop:torus-induction} on a
one-dimensional space gives the rank-one assertion. Suppose
first \(c=P\) is irreducible of degree \(d\) and \(x\in K_d^\times\)
is one of its roots. The roots of \(P^{[u]}\) are the conjugates
of \(ux\), distinct and numbering \(d\), so this is its minimal
polynomial by Lemma~\ref{lem:torus-primary}.
Write \(u_d=\iota^{-1}(\exp(2\pi i/(Q^d-1)))\) and
\(D=(Q^d-1)/(Q-1)\). Since \(u_1=u_d^D\), the exponent of
\(ux\) at rank \(d\) equals that of \(x\) plus \(D\) times the
rank-one exponent of \(u\), modulo \(Q^d-1\). Hence
\[
       \theta_{ux}(z)=\theta_x(z)\,\chi_u(z^D)
                                  \qquad(z\in K_d^\times).
\]
Here \(z^D\in K^\times\), since its \(Q\)-th power equals itself.

At matrices whose characteristic polynomial has two different
irreducible factors, both cuspidal characters vanish by
Lemma~\ref{lem:torus-primary}. At a matrix of characteristic
polynomial \(P_0^s\), with \(\deg P_0=e,\ es=d\), every eigenvalue
\(z\) belongs to \(K_e^\times\subset K_d^\times\). The product
of its \(d\) successive Frobenius powers equals the product of
the \(e\) roots of \(P_0\) raised to the \(s\)-th power, hence
is \(\det g=z^D\), with no condition on diagonalizability.
Every invariant subspace uses a submultiset of these eigenvalues,
by its triangular block matrix. In the defining sum
\(C_{\theta_{ux}}(g)=\sum_W\beta_{\theta_{ux}}(g|W)\mu(g|W)\)
the last display thus multiplies every summand for \(x\) by the
same scalar \(\chi_u(\det g)\). This proves the twist identity
for an irreducible factor. For a general label take the flag
product of Definition~\ref{def:squarefree-green}:
the product of the determinants on the successive quotients of
each invariant flag equals \(\det g\). This gives the identity
on the whole class function.
\end{proof}

\begin{theorem}\label{thm:green-reflection}
Let \(\ell:K^r\to K\) be a nonzero linear functional and put
\(H_\ell=\{g\in G_r:g|_{\ker\ell}=1\}\). In \(\C[G_r]\) set
\[
             B_{\ell,u}=\sum_{g\in H_\ell}\chi_u(\det g)^{-1}[g].
\]
For a label \(c\) with \(c(u)\ne0\) one has
\(B_{\ell,u}\, e_{\Xi_c}=0\), where \(e_{\Xi_c}\) is the central
projection of Lemma~\ref{lem:central-proj} for the representation
in Theorem~\ref{thm:squarefree-green}. In particular, for the
group \(\mathcal G\) and \(A=M_\omega\) of
Definition~\ref{def:poly-types}, any pure group-algebra tensor
with one factor \(B_{\ell,u}\) satisfying \(\chi_u^A=1\)
annihilates \(e_0^{\rm good}\).
\end{theorem}
\begin{proof}
The matrices in \(H_\ell\) are exactly the invertible maps
\(I+a\ell,\ a\in K^r\), since their difference from the identity
factors through the one-dimensional quotient by \(\ker\ell\).
Choose a basis of that kernel extended by a vector of
\(\ell\)-value one. Such a matrix has diagonal blocks \(I_{r-1}\)
and \(t\in K^\times\), and arbitrary upper right column.
In any representation, \(|H_\ell|^{-1}B_{\ell,u}\) acts as
the projection onto vectors transforming by \(\chi_u\circ\det\)
under \(H_\ell\). Translation of the averaging sum proves the
transformation law, and the average fixes each vector with that
law. Its trace on the representation with character \(\Xi_c\) is
\[
  \frac{1}{Q-1}\sum_{t\in K^\times}
       \overline{\chi_u(t)}\,(\Delta\Xi_c)(I_{r-1},t).
\]
By Proposition~\ref{prop:block-mackey} and
Theorem~\ref{thm:virtual-cuspidal}, in this bidegree exactly one
degree-one irreducible factor of \(c\) must go into the second
slot, all others going to the first. Thus the expression is a
sum over the linear factors \(X-v\) of \(c\) of a constant
times \(\langle\chi_u,\chi_v\rangle_{K^\times}\), using
Proposition~\ref{prop:green-named-twist}. This is zero if
\(c(u)\ne0\), by the geometric sum in Lemma~\ref{lem:characters}.
The trace of a projection is its dimension of image, so the
operator vanishes. On all other irreducibles \(e_{\Xi_c}\) acts
as zero by Lemma~\ref{lem:central-proj}. Decomposing the faithful
regular representation as in that lemma proves the group-algebra
identity. In the two-prime tensor, \(\chi_u^A=1\) implies \(u^A=1\)
by the group isomorphism. Each good tuple avoids that root in the
chosen factor, by Definition~\ref{def:poly-types}; each tensor
summand of Corollary~\ref{cor:actual-good-share} is therefore
annihilated.
\end{proof}

\begin{corollary}\label{cor:green-cubic-alternatives}
Suppose \(K=\mathbf F_p,\ 9\mid p-1,\ r=3^a,\ a\ge1\). For a
primitive label \(c\) of degree \(r\) and any two prescribed labels
there is a primitive \(P\) avoiding both such that
\[
       [g]e_{\Xi_P}=[g]e_{\Xi_c}\quad\text{for every \(g\in G_r\)
                    with reducible characteristic polynomial}.
\]
Here \([g]e\) denotes the coefficient of \([g]\) in \(e\).
\end{corollary}
\begin{proof}
By Lemma~\ref{lem:primitive-scale} at least one of the three
primitive labels \(c^{[u]}\), \(u^3=1\), avoids both prescribed
labels. Select such a \(P\). Proposition~\ref{prop:green-named-twist}
gives \(\Xi_P=\chi_u(\det)\Xi_c\) and \(\chi_u^3=1\). As the label
is irreducible, \(\Xi_c\) vanishes on a matrix with two distinct
irreducible characteristic factors by Lemma~\ref{lem:torus-primary}.
In the remaining reducible case the characteristic polynomial is
\(D^s\) for one irreducible \(D\), \(s>1\) dividing \(3^a\).
Its determinant is the \(s\)-th power of the product of the
roots of \(D\) in \(K^\times\), hence a cube, and the twist factor
is one. Thus the characters agree at every stated matrix. The
same matrices have inverses with reducible characteristic
polynomial: an invertible matrix with irreducible characteristic
polynomial acts as multiplication by a degree-\(r\) field
generator on \(K_r\), by Lemma~\ref{lem:torus-primary} and its
proof. Its inverse also generates that same field, and hence
has irreducible characteristic polynomial. Applying this argument
to the inverse proves the needed implication. Finally the dimensions
agree at the identity by the twist law. The formula for central
projections in Lemma~\ref{lem:central-proj} now proves equality
of their coefficients.
\end{proof}

\section{Flag minors for the compact row weights}

\begin{definition}\label{def:compact-flag-space}
Let \(n\ge1\) and let \(\lambda_1\ge\cdots\ge\lambda_n\) be
integers. Put \(\kappa_i=\lambda_i-\lambda_n\). Draw \(\kappa_i\)
boxes in row \(i\), starting at the left, and read its columns
in order from left to right, with respective heights
\(d_1\ge\cdots\ge d_m>0\); the empty diagram is allowed.
For \(I=(i_1<\cdots<i_d)\subset\{1,\ldots,n\}\) put
\[
       M_I(X)=\det(X_{i_a,b})_{1\le a,b\le d}
                         \qquad(X\in{\rm GL}_n(\C)).
\]
Let \(V_\lambda\) be the complex span of the functions
\((\det X)^{\lambda_n}\prod_{j=1}^m M_{I_j}(X)\), where
\(|I_j|=d_j\). If the entries of each \(I_j\) are written from
top to bottom in column \(j\), call the filling semistandard
when its entries weakly increase from left to right in each row.
The group acts on functions by \((\rho(g)F)(X)=F(g^{\mathsf T}X)\).
\end{definition}

\begin{proposition}\label{prop:standard-flag-minors}
The action in Definition~\ref{def:compact-flag-space} preserves
\(V_\lambda\). The functions indexed by semistandard fillings
form a basis of \(V_\lambda\).
\end{proposition}
\begin{proof}
The order of transposes gives \(\rho(g)\rho(h)=\rho(gh)\).
A \(d\)-row minor of \(g^{\mathsf T}X\) is the sum over \(d\)-row
subsets of \(X\) of their minors times the corresponding minor
of \(g^{\mathsf T}\). To verify this expansion, expand the
determinant by multilinearity in its rows. Equal chosen row
indices give zero, and ordering the remaining indices gives
the coefficient determinant. The common determinant power
contributes a scalar. This proves stability.

Omit the common nowhere vanishing determinant factor for the
basis calculation. Consider adjacent columns \(I,J\) of respective
lengths \(a\ge b\) with \(I_s>J_s\) for some \(1\le s\le b\).
Regard the last \(a-s+1\) entries of \(I\) and the first \(s\)
entries of \(J\) as \(a+1\) independent row vectors of length
\(a\), using only their first \(b\) coordinates in the second
minor. Alternate the product of the two minors on these
\(a+1\) vectors. The result is zero, since every alternating
multilinear form on \(a+1\) vectors in dimension \(a\) vanishes
by expansion in a basis. Group the permutations by the choice
of vectors put into the first column. Up to the common positive
factor \((a-s+1)!\,s!\), each choice occurs once with sign, since
permutations within either group have their signs canceled by
the determinant signs. Substituting the original rows of \(X\)
therefore writes the original product as a signed sum of the
other choices.

Every new first column has received at least one of
\(J_1,\ldots,J_s\), all strictly smaller than
\(I_s,\ldots,I_a\). If a row is repeated its minor vanishes.
Otherwise sorting each column introduces only a sign and makes
the first column lexicographically smaller than \(I\). The bound
on the two sets of exchanged entries ensures that no other
choice gives the original first column. Reading all indexed
columns from left to right gives a finite lexicographically ordered
set of lists. Straightening thus terminates and proves spanning.

For independence, order the variables first by matrix column,
and in a column give smaller row index higher priority. In the
lexicographic order on exponent vectors for this list, the
leading monomial of \(M_I\) is \(\prod_{a=1}^d X_{i_a,a}\).
The determinant chooses the smallest available row successively,
which proves this assertion and its coefficient one. Leading
monomials multiply, since comparison of two exponent vectors is
invariant under adding a fixed vector. In a product the exponents
in matrix column \(a\) recover the multiset of entries in row
\(a\) of the filling. For a semistandard filling their sorted
list determines that entire row, so all the leading monomials
are distinct. The polynomials are independent also as functions
on the invertible matrices. Indeed \(X+tI\) is invertible for
all but at most \(n\) scalars \(t\), and polynomial continuation
in \(t\) then extends any relation to every \(X\). A polynomial
zero at all matrices has zero coefficients by induction in its
variables using the one-variable root bound. This proves the
basis statement.
\end{proof}

\begin{lemma}\label{lem:tableaux-binomial}
For a partition \(\alpha_1\ge\cdots\ge\alpha_n\ge0\), the number
of semistandard fillings of its diagram with entries
\(1,\ldots,n\), strictly increasing down columns, is
\[
  \det\!\left[\binom{\alpha_i+n-i}{n-j}\right]_{1\le i,j\le n}
    =\prod_{1\le i<j\le n}
                  \frac{\alpha_i-\alpha_j+j-i}{j-i}.
\]
\end{lemma}
\begin{proof}
The assertion is one in dimension one. Removing the entries
equal to \(n\) leaves a diagram \(\gamma\) of length at most
\(n-1\) with \(\alpha_i\ge\gamma_i\ge\alpha_{i+1}\) for
\(1\le i<n\). Indeed those entries form row suffixes, and
a removed box has no box below it. Every box in row \(n\) is
removed by strict column increase. Conversely, under these
inequalities the boxes of \(\alpha\) outside \(\gamma\) lie at
the bottoms of distinct columns, so filling them with \(n\)
extends each filling of \(\gamma\) uniquely.

Write \(x_i=\alpha_i+n-i\) and \(y_i=\gamma_i+n-1-i\).
The ranges \(x_{i+1}\le y_i<x_i\) are independent and imply
the required ordering of \(\gamma\). Sum the determinant formula
in rank \(n-1\) over these ranges, independently in its rows.
For \(h\ge0\) the binomial recursion telescopes to
\[
       \sum_{y=x_{i+1}}^{x_i-1}\binom yh
                    =\binom{x_i}{h+1}-\binom{x_{i+1}}{h+1}.
\]
The resulting determinant is the one in the statement: in that
\(n\)-row determinant subtract row \(i+1\) from row \(i\), working
from top to bottom, and expand along its last column, equal after
these operations to the last coordinate vector. This proves the
first formula by induction. For the second, as in
Proposition~\ref{prop:exact-row-degree}, use the leading
coefficient \(1/h!\) of \(\binom{X}{h}\) and the Vandermonde
identity. With the degrees and nodes both in decreasing order
the result is \(\prod_{i<j}(x_i-x_j)/\prod_{h=0}^{n-1}h!\),
which is the stated product.
\end{proof}

\begin{corollary}\label{cor:realized-compact-degree}
For the odd row rank \(r\) of Definition~\ref{def:row-vandermonde},
put \(\rho_i^0=(r+1)/2-i\) for \(1\le i\le r\), and
\(\lambda_{i,\tau}=t_{i-1,\tau}-\rho_i^0\). These are decreasing
in the weak sense and are integers. The dimension of the
space in Definition~\ref{def:compact-flag-space} at \(\tau\)
is \(\mathcal D_\tau\). The tensor product over the \(f\)
infinite places has dimension \(\mathcal D\).
\end{corollary}
\begin{proof}
Adjacent \(t\)'s decrease by an integer at least one by
Proposition~\ref{prop:exact-row-degree}, and successive \(\rho_i^0\)
differ by one. The dimension by Proposition~\ref{prop:standard-flag-minors}
and Lemma~\ref{lem:tableaux-binomial} for
\(\alpha_i=\lambda_{i,\tau}-\lambda_{r,\tau}\) has numerator
differences \(\lambda_{i,\tau}-\lambda_{j,\tau}+j-i
=t_{i-1,\tau}-t_{j-1,\tau}\). This is exactly the product in
Definition~\ref{def:row-vandermonde}. Tensor products of bases
multiply these dimensions.
\end{proof}

\section{The compact row representation and its matrix coefficient}

\begin{definition}\label{def:compact-coefficient}
For the weight \(\lambda\) and column heights of
Definition~\ref{def:compact-flag-space}, let
\[
 \mathcal T_\lambda=
  (\det)^{\lambda_n}\otimes
    \bigotimes_{j=1}^m \bigwedge^{d_j}\C^n,\qquad
 v_\lambda=1\otimes\bigotimes_{j=1}^m(e_1\wedge\cdots\wedge e_{d_j}).
\]
The first factor is the line on which \(g\) acts by
\((\det g)^{\lambda_n}\). On each exterior power use the
inner product making its increasing wedges of standard basis
vectors orthonormal, and give the tensor product the product
inner product, conjugate linear in the first variable. Denote
this representation by \(T\). Write \(W_\lambda\) for the span
of \(T(k)v_\lambda\) for \(k\in U(n)=\{k:k^*k=I\}\).
\end{definition}

\begin{theorem}\label{thm:compact-highest-vector}
The \(U(n)\)-space \(W_\lambda\) is irreducible, isomorphic to
\(V_\lambda\), and extends to the rational
\({\rm GL}_n(\C)\)-action there. Its upper unitriangular-fixed
line is spanned by \(v_\lambda\), on which a diagonal matrix
\(t\) acts by \(\prod_i t_i^{\lambda_i}\). Its unit vector
\(v_\lambda\) has matrix coefficient
\[
       m_\lambda(k)=\langle v_\lambda,T(k)v_\lambda\rangle
       =(\det k)^{\lambda_n}\prod_{j=1}^m
                           \det k_{\{1,\dots,d_j\},\{1,\dots,d_j\}} .
\]
\end{theorem}
\begin{proof}
The exterior-power matrix entries are minors by expanding the
wedge in basis vectors. This also shows
\(T(g^{\mathsf T})=T(g)^{\mathsf T}\) and \(T(g^*)=T(g)^*\).
Thus \(T(k)\) is unitary and the coefficient formula follows
by selecting the initial coordinate in each wedge. For any vector
\(a\) in the tensor coordinates define the linear function of
\(a\) with values in the space of functions on the group by
\[
                 F_a(X)=a^{\mathsf T}T(X)v_\lambda.
\]
It sends the standard tensors to the spanning products in
Definition~\ref{def:compact-flag-space} and intertwines \(T(g)\)
and \(\rho(g)\) by the transpose identity.

Every complex \(U(n)\)-invariant subspace of \(\mathcal T_\lambda\)
is \({\rm GL}_n(\C)\)-invariant. Indeed it is finite-dimensional
and closed. Differentiation of \(T(\exp(tA))\) for
\(A^*=-A,\ t\) real, shows it is stable under the infinitesimal
matrix \(dT(A)\). The entries of \(T\) are rational functions
regular on the group, so \(dT\) is complex linear. Every complex
matrix is a complex linear combination of skew-Hermitian
matrices by its Hermitian and skew-Hermitian parts. The subspace
is therefore invariant under \(T(\exp A)=\exp(dT(A))\) for
all \(A\). For the equality, differentiation of the group law
gives the constant-coefficient differential equation on
\(t\mapsto T(\exp(tA))\); multiplying by \(\exp(-t\,dT(A))\)
gives derivative zero and value \(I\) at zero. The elementary
shears \(I+aE_{ij}\) and nonzero diagonal scalings are exponentials
by \(E_{ij}^2=0\) for \(i\ne j\) and the complex logarithm of a
nonzero scalar. Row reduction expresses every invertible matrix
in these generators: adding a row with nonzero pivot when needed
and rescaling permits elimination without row exchanges. This
proves the invariance assertion.

In particular the orbit under \({\rm GL}_n(\C)\) spans \(W_\lambda\).
It is stable under coordinate conjugation, since conjugation
fixes \(v_\lambda\) and conjugates each representing matrix to
the matrix at the conjugate element. Consequently the kernel
of \(a\mapsto F_a\), the bilinear annihilator of \(W_\lambda\),
is its Hermitian orthogonal complement. Restricting to \(W_\lambda\)
now identifies it equivariantly with \(V_\lambda\).

Every \(F\in V_\lambda\) satisfies
\(F(Xb)=\prod_i b_{ii}^{\lambda_i}F(X)\) for \(b\) upper triangular
invertible, since the first \(d\) columns are multiplied on the
right by the first \(d\) diagonal block of \(b\). If \(F\) is
also fixed by \(\rho(u)\) for all upper unitriangular \(u\),
its values are determined by \(F(I)\). In fact every matrix with
all initial principal minors nonzero is a product of a lower
unitriangular matrix and an upper triangular invertible matrix:
eliminate below each successive pivot; the next pivot equals
the ratio of consecutive leading minors. On these matrices the
two transformation laws determine \(F\). They form a dense set,
since each initial principal minor of \(X+tI\) is a monic
polynomial in \(t\). Thus at most a line is fixed.
The nonzero vector \(v_\lambda\) is fixed by these upper
unitriangular matrices as they preserve the initial wedges; its
diagonal weight is the one stated.

Finally every nonzero invariant subspace contains an upper
unitriangular-fixed vector. In the full tensor model the diagonal
action is diagonal with finitely many integer weights. Any
invariant subspace is the direct sum of its intersections with
the weight spaces. Indeed choose a single positive diagonal
element separating the finitely many monomials (avoid the finitely
many proper real hyperplanes in the logarithms of its entries),
then use polynomial interpolation on its distinct eigenvalues.
A finite union of such hyperplanes is proper since the product
of their nonzero linear polynomials is a nonzero polynomial and
cannot vanish everywhere, by the argument of
Proposition~\ref{prop:standard-flag-minors}. Choose a vector of
lexicographically maximal weight in the subspace. For \(i<j\),
\(dT(E_{ij})\) raises a weight by +1 in coordinate \(i\) and
\(-1\) in coordinate \(j\), as follows by conjugating \(E_{ij}\)
by diagonal matrices. Its value on the chosen vector is zero.
Exponentiating shows all upper elementary shears fix this vector,
and these generate the upper unitriangular group by successive
elimination. If \(W_\lambda\) had a proper nonzero \(U(n)\)-invariant
subspace, it and its nonzero orthogonal complement would both
give fixed vectors, contradicting the one-dimensional fixed space
just proved. This establishes irreducibility.
\end{proof}

\begin{proposition}\label{prop:compact-formal-degree}
For invariant probability measure \(\dd k\) on \(U(n)\),
\[
     \int_{U(n)}|m_\lambda(k)|^2\dd k
       =\frac{1}{\dim W_\lambda}.
\]
In particular, for the compact row representations of
Corollary~\ref{cor:realized-compact-degree} the product of these
integrals over the infinite places is \(1/\mathcal D\).
\end{proposition}
\begin{proof}
To construct the measure, give \(n\)-by-\(n\) complex matrices
probability density proportional to
\(\exp(-\operatorname{tr}(X^*X))\) times Lebesgue measure.
It is normalizable by the Gaussian integral in
Lemma~\ref{lem:theta} and is invariant under left and right
unitary multiplication by real orthogonal change of variables.
The singular matrices have measure zero: a nonzero polynomial
in independent complex coordinates has length-area product
measure zero on its zero set, by induction in the number of
variables, first excluding the common zero set of the nonzero
leading coefficient and then the finitely many roots in each
remaining fiber. On invertible matrices let \(k=X(X^*X)^{-1/2}\).
For a positive Hermitian matrix the positive inverse square root
is defined on its eigenspaces by the positive real inverse square
roots of its eigenvalues. Orthogonal diagonalization follows by
taking an eigenvector and passing to its invariant complement,
and positivity gives positive eigenvalues. This matrix function
is continuous locally: on a fixed positive interval containing the
nearby eigenvalues, expand \(z^{-1/2}\) into its convergent binomial
series about a positive number larger than the upper endpoint,
uniformly on the interval. Such nearby bounds follow by bounding
the change of the quadratic form on unit vectors by the matrix
operator norm. The scalar expansion follows by differentiating
the power series solution to \(2(1-t)f'(t)=f(t)\), \(f(0)=1\).
The image \(k\) is unitary and the map is equivariant on the
left and right, since unitary conjugation carries each eigenspace
to the conjugate eigenspace. The image probability is therefore
bi-invariant. For uniqueness, integrate \(F(xy)\) on the product
of this measure in \(x\) and any left-invariant probability in
\(y\): integrating in the two orders gives equality of their
integrals for every bounded Borel \(F\).

On \(W_\lambda\) average the projection onto each unit vector
\(T(k)v_\lambda\). By left invariance the resulting operator
commutes with the group. It is scalar, since any eigenspace of
a commuting complex matrix is invariant and the space is
irreducible by Theorem~\ref{thm:compact-highest-vector}. Its trace
is one, so its scalar is \(1/\dim W_\lambda\). Evaluating the
quadratic form of the average at \(v_\lambda\) proves the formula.
Multiply over the places and use
Corollary~\ref{cor:realized-compact-degree}.
\end{proof}

\section{Definite automorphic forms at the row level}

\begin{definition}\label{def:definite-comparison-forms}
Let \(F=E^{\mathfrak c}\), and let \(H\) be the unitary group of
\(\sum_{a=1}^r\bar x_a x_a\) on \(E^r\). Write \(H_f\) for its
finite adelic group, \(H(F)\) for its rational isometries,
\(K^0\) for the stabilizer of the standard lattice at every
finite place, and \(K(g)\) for its subgroup reducing to the
identity on the chosen factors over \(g_0,g_1\). Write
\(K_\infty=\prod_{\tau\in\Phi}U(r)\). For the weight of
Corollary~\ref{cor:realized-compact-degree} put
\(W=\bigotimes_\tau W_{\lambda_\tau}\),
\(v=\bigotimes_\tau v_{\lambda_\tau}\); let \(R_\infty(k)\)
act on \(W\) by \(\bigotimes_\tau T_\tau(\bar k_\tau)\).
This uses conjugation on the compact argument, as for the
column convention in Proposition~\ref{prop:column-dim}.
Let \(\mathcal A^+\) be the functions \(L:H_f\to W\) satisfying
\[
    L(\gamma h u)=R_\infty(\gamma_\infty)L(h)
       \qquad(\gamma\in H(F),\ h\in H_f,\ u\in K(g)).
\]
Here \(\gamma_\infty\) is the tuple of its matrices at the chosen
complex embeddings. On each compact group take the probability
measure of Proposition~\ref{prop:compact-formal-degree}.
\end{definition}

\begin{proposition}\label{prop:definite-scalar-projection}
The space \(\mathcal A^+\) is linearly isomorphic to the continuous
scalar functions on \(K_\infty\times H_f\) that are invariant on
the left by \(H(F)\), invariant on the right by \(K(g)\), and
fixed by the operator
\[
    (\mathcal P f)(k,h)=\mathcal D\int_{K_\infty}
            \langle R_\infty(k)v,R_\infty(t)v\rangle f(t,h)\dd t .
\]
This operator is the orthogonal projection in the compact variable
onto the space of functions \(k\mapsto\langle R_\infty(k)v,w\rangle\)
with \(w\in W\). The isomorphism sends \(L\) to
\(f_L(k,h)=\langle R_\infty(k)v,L(h)\rangle\).
\end{proposition}
\begin{proof}
Complex conjugation preserves the compact group probability,
since its image is again invariant and the probability is unique
by Proposition~\ref{prop:compact-formal-degree}. The average of the
orthogonal projections onto \(R_\infty(k)v\) therefore factors
over \(\tau\) as the tensor product of the averages computed in
that proposition. For every \(w\in W\) this gives
\[
 \int_{K_\infty} R_\infty(t)v\,\langle R_\infty(t)v,w\rangle\dd t
                  =w/\mathcal D .
\]
In particular the coefficient map is injective, and the
displayed kernel fixes each of its functions. Its image on any
continuous \(f\) is the coefficient function of
\(\mathcal D\int R_\infty(t)v f(t,h)\dd t\) at fixed \(h\).
It is self-adjoint since its kernel is conjugate symmetric
on exchanging \(k,t\). The same formulas apply on the compact
\(L^2\) space by Cauchy--Schwarz, proving the orthogonal projection
assertion.

Unitarity and the defining rule for \(L\) give
\(f_L(\gamma_\infty k,\gamma h u)=f_L(k,h)\).
Continuity follows from local constancy on finite cosets and
continuity of the compact coefficient. Conversely, for an invariant
\(f\) fixed by \(\mathcal P\) use the vector integral just given
to define \(L(h)\). Substitution \(t=\gamma_\infty t'\) and invariance
of measure give the defining rule of
Definition~\ref{def:definite-comparison-forms}. Since \(\mathcal P f=f\)
the coefficient recovers \(f\). The vector average recovers
\(L\) in the opposite direction as well.
\end{proof}

\begin{theorem}\label{thm:definite-comparison-rank}
For the two marked primes in the arithmetic construction one has
\[
    K^0/K(g)\simeq\mathcal G,\qquad
            \dim\mathcal A^+=|\mathcal G|\,\mathfrak m_r\,\mathcal D.
\]
Let each \(e_i\) of Corollary~\ref{cor:actual-good-share} act by
right translation through its factor
\(J_i={\rm GL}_r(\mathbf F_{g_i})^f\) of \(\mathcal G\).
Write \(R_i\) for its action and
\(\delta_i=[1](1-e_i)\). Then
\[
 \operatorname{rank}(1-R_i)=\delta_i\dim\mathcal A^+,\qquad
 (1-\delta_0-\delta_1)\dim\mathcal A^+
       \le \dim(e_0^{\rm good}\mathcal A^+)\le\dim\mathcal A^+ .
\]
At all sufficiently large scales \(\delta_0+\delta_1\le b^{-8}\).
\end{theorem}
\begin{proof}
At a marked rational prime the field splits completely by
Lemma~\ref{lem:split-level}. On each pair an integral isometry
is specified by an arbitrary matrix in \({\rm GL}_r(\Z_{g_i})\)
on the chosen component, with inverse transpose on its partner.
Reduction onto \({\rm GL}_r(\mathbf F_{g_i})\) is surjective by
lifting the entries of an invertible matrix as in
Proposition~\ref{prop:column-dim}. The other places impose the same
subgroup in numerator and denominator. This gives the quotient.

The double cosets of \(H(F)\) by \(K^0\) are finite by
Theorem~\ref{thm:class-number-bound} in rank \(r\). For a representative
\(h\), its stabilizer \(\Gamma_h=H(F)\cap hK^0h^{-1}\) is finite,
since it acts faithfully on its positive definite integral lattice
with trace metric as in Lemma~\ref{lem:trace-finite}. Any finite-order
element of \(H(F)\cap hK(g)h^{-1}\) is the identity. At a chosen
component of \(g_0\) its conjugate belongs to the torsion-free
principal congruence kernel proved in Proposition~\ref{prop:column-dim}.
The embedding of \(E\) into that completion is injective.
Thus \(\Gamma_h\) acts freely on the \(|\mathcal G|\) cosets of
\(K(g)\) in \(K^0\) over \(h\). Counting orbits and summing gives
\[
       |H(F)\backslash H_f/K(g)|
          =\sum_{[h]}|\mathcal G|/|\Gamma_h|
          =|\mathcal G|\,\mathfrak m_r .
\]
Evaluation at these finer representatives takes arbitrary values
in \(W\), by triviality of their stabilizers and the transformation
rule for \(L\). Use Corollary~\ref{cor:realized-compact-degree}
for its dimension. The sum of squared norms of these values gives
a Hilbert norm independent of the representative choices by
unitarity. Right translations act unitarily in this norm.

For any \(s\in J_i\), choose a lift \(\widetilde s\in K^0\) whose
reduction is trivial in \(J_{1-i}\). If right translation by \(s\)
fixes a finer class \([h]\), then \(h\widetilde s=\gamma h u\)
for \(\gamma\in H(F),\ u\in K(g)\). Thus \(\gamma\in\Gamma_h\)
and its conjugate has trivial reduction at the other marked prime.
The same torsion argument at \(g_{1-i}\) forces \(\gamma=1\),
then \(s=1\). Each \(J_i\)-orbit of the finer class set is
therefore free. Evaluate \(L\) at the points \(h\widetilde t\),
\(t\in J_i\), over one representative per orbit. Right translation
just permutes these entries regularly, since different lifts
of the same product differ by \(K(g)\). On each orbit the
action is a direct sum of \(\dim W\) copies of the regular
representation. Its normalized trace on \(1-e_i\) equals the
identity coefficient \(\delta_i\) by the proof of
Proposition~\ref{prop:bad-types}. This trace equals relative rank
since it is a projection. Finally \(R_0,R_1\) commute and their
product is the two-prime projection of
Corollary~\ref{cor:actual-good-share}. The codimension of the
intersection of two images is at most the sum of the codimensions,
by mapping to the direct sum of the quotient spaces. The
quantitative bound follows from Proposition~\ref{prop:bad-types}
and Theorem~\ref{thm:squarefree-green}.
\end{proof}

\begin{corollary}\label{cor:definite-energy-capacity}
With \(C,J_c,E_2\) as in
Definitions~\ref{def:packet-amounts} and~\ref{def:two-energy-numbers},
\[
              E_2\,\dim\mathcal A^+=|\mathcal G|(C J_c)^2 .
\]
For any vector space whose dimension is at most that of the
marked definite subspace \(e_0^{\rm good}\mathcal A^+\), the
left side with its dimension in place of \(\dim\mathcal A^+\)
is bounded above by this same right side.
\end{corollary}
\begin{proof}
Insert the dimension from Theorem~\ref{thm:definite-comparison-rank}
in \(E_2=(C J_c)^2/(\mathcal D\mathfrak m_r)\). The remaining
assertion uses positivity and the subspace inclusion.
\end{proof}

\section{A finite reduction of the marked trace comparison}

\begin{definition}\label{def:elliptic-tests}
Use the two-prime finite group \(\mathcal G\) and the labels of
Corollary~\ref{cor:actual-good-share}. Let \(\mathcal E_0\subset\mathcal G\)
be the elements whose matrix at some chosen coordinate over \(g_0\)
has irreducible characteristic polynomial. For a label \(c\) at
one place write \(e(c)=e_{\Xi_c}\), and for a tuple use
\(e(\mathbf c)=\bigotimes_w e(c_w)\). In a representation at one
coordinate, \(e(c_w)\) means the identity in every other factor.

Let \(H_1,H_2\) be finite-dimensional unitary representations of
\(\mathcal G\). Let \(\mathcal B\) be a unital subalgebra of
\({\rm End}(H_1)\times{\rm End}(H_2)\) whose elements commute with
the respective \(\mathcal G\)-actions. Write \(B_j\) for the
component of \(B\in\mathcal B\) on \(H_j\). Let \((A_i)_{i\in I}\)
be a commuting family in \(\mathcal B\) whose members and their
adjoints commute with one another on \(H_1\oplus H_2\). A joint
eigensystem is a tuple \(\omega=(\omega_i)\) with a nonzero
joint eigenspace on this direct sum; \(H_{j,\omega}\) denotes the
corresponding joint eigenspace on \(H_j\). Say a label at \(w\)
occurs on \(H_{j,\omega}\) if its one-coordinate projection
has nonzero image there.
\end{definition}

\begin{lemma}\label{lem:common-hecke-interpolation}
The direct sum in Definition~\ref{def:elliptic-tests} is the
orthogonal sum of the finitely many joint eigenspaces.
For every joint eigensystem \(\omega\) a polynomial in finitely
many \(A_i\) is the projection \(\Pi_\omega\) onto it on both
spaces. Each \(H_{j,\omega}\) is \(\mathcal G\)-invariant.
\end{lemma}
\begin{proof}
A normal complex matrix has an orthonormal eigenbasis.
Indeed it has an eigenvalue by Lemma~\ref{lem:place-fields}.
For an eigenvector with eigenvalue \(z\), use
\(\|(A^*-\bar z)v\|=\|(A-z)v\|\), obtained by expanding both
squares and using normality. Its orthogonal complement is
invariant under both maps, and induction proves the claim.
Diagonalize a nonscalar member of the family. The other
operators and their adjoints preserve every eigenspace by
commutation. Repeat on the proper summands so obtained until
all restrictions are scalar, which terminates by decreasing
dimension along every branch. Collect equal joint eigenvalues.
For every \(\omega'\ne\omega\) select \(i(\omega')\) where they
differ. Then
\[
 \Pi_\omega=\prod_{\omega'\ne\omega}
 \frac{A_{i(\omega')}-\omega'_{i(\omega')}I}
      {\omega_{i(\omega')}-\omega'_{i(\omega')}} .
\]
This acts as one on the indicated block and zero on all others.
The empty product works for a single block. Commutation with
\(\mathcal G\) gives its invariance assertion.
\end{proof}

\begin{theorem}\label{thm:elliptic-test-transfer}
In Definition~\ref{def:elliptic-tests} assume that
\[
       \operatorname{tr}_{H_1}(B_1s)
        =\operatorname{tr}_{H_2}(B_2s)
             \qquad(B\in\mathcal B,\ s\in\mathcal E_0),
\]
where \(s\) acts by the respective representations. Suppose
also that for each \(\omega\) and each good tuple \(\mathbf c\)
there is a chosen coordinate \(w\) over \(g_0\) at which
\(c_w\) is primitive and one of the three labels
\(c_w^{[u]},\ u^3=1\), occurs on neither \(H_{1,\omega}\) nor
\(H_{2,\omega}\). Then for every \(B\in\mathcal B\) the restrictions
of \(B_1,B_2\) to \(e_0^{\rm good}H_1,e_0^{\rm good}H_2\)
have equal trace and equal characteristic polynomial.
For the occurrence hypothesis it suffices that on each of the
two spaces at most one primitive label occurs at each such
coordinate for each eigensystem.
\end{theorem}
\begin{proof}
The ranks and primes in the construction satisfy
\(r=3^a,\ a\ge1,\ 9\mid g_0-1\) at sufficiently large \(x\)
by Proposition~\ref{prop:thick} and Theorem~\ref{thm:data}.
For fixed \(\omega,\mathbf c\), choose the indicated \(w\) and
a missing cubic alternative \(P=c_w^{[u]}\). Replace the projection
in this single coordinate in \(e(\mathbf c)\) by \(e(P)\),
calling the resulting product \(e'\). It annihilates both
\(\omega\)-spaces, since \(e(P)\) does so and the other coordinate
factors commute with it. The coefficients of
\(e(\mathbf c)-e'\) vanish at every group element whose \(w\)-matrix
has reducible characteristic polynomial: the coefficient at that
element is the product of the unchanged coefficients times the
difference of the two at \(w\). This difference vanishes by
the proof of Corollary~\ref{cor:green-cubic-alternatives}, which
applies to each of the three alternatives.

Summing \(e(\mathbf c)-e'\) over the good tuples gives a group
algebra element \(T_\omega\) supported on \(\mathcal E_0\).
On both \(\omega\)-spaces it acts exactly as \(e_0^{\rm good}\).
Now \(\Pi_\omega\in\mathcal B\) by
Lemma~\ref{lem:common-hecke-interpolation}. Apply the trace
hypothesis to \(B\Pi_\omega\) and extend linearly in \(s\) to
\(T_\omega\). On each space \(\Pi_\omega T_\omega=
\Pi_\omega e_0^{\rm good}\), since all finite group actions commute
with \(\Pi_\omega\). Summing over \(\omega\) yields equality
of the full traces of \(B_j e_0^{\rm good}\). These are precisely
the traces on the indicated ranges by commutation and the
idempotent identity.

At \(B=1\) the dimensions agree. Apply the trace equality also
to every positive power of \(B\), all of which lie in \(\mathcal B\).
Over \(\C\), triangularization by successive invariant eigenlines
shows that the trace of \(B^a\) is the power sum of the eigenvalues
with algebraic multiplicity. The polynomial identity in the proof
of Theorem~\ref{thm:brauer-power-sum} determines every elementary
symmetric coefficient recursively from the power sums: the
coefficient of degree \(a\) on the differentiated side gives
a nonzero integer multiple of the \(a\)-th elementary symmetric
function. This proves equality of characteristic polynomials.
Finally every good tuple has a primitive coordinate over \(g_0\).
The three alternatives there are distinct and primitive by
Lemma~\ref{lem:primitive-scale}; two singleton occurrence sets
exclude at most two of them.
\end{proof}

\begin{corollary}\label{cor:conditional-indefinite-dimension}
Suppose a principal-level space \(H_1\) on the indefinite form
of Definition~\ref{def:row-space} and the definite space
\(H_2=\mathcal A^+\) satisfy Theorem~\ref{thm:elliptic-test-transfer}
for a common algebra \(\mathcal B\) and commuting family as specified
there. Put \(d_0=\dim(e_0^{\rm good}H_1)\). Then
\[
  (1-b^{-8})|\mathcal G|\,\mathfrak m_r\mathcal D
           \le d_0\le |\mathcal G|\,\mathfrak m_r\mathcal D,
  \qquad
  d_0 E_2\le|\mathcal G|(CJ_c)^2
\]
at sufficiently large scales.
\end{corollary}
\begin{proof}
Take the identity in Theorem~\ref{thm:elliptic-test-transfer}.
Theorem~\ref{thm:definite-comparison-rank} bounds the resulting
common dimension on both sides.
Corollary~\ref{cor:definite-energy-capacity} gives the last inequality.
\end{proof}

\section{Elliptic tests on the definite group}

\begin{theorem}\label{thm:compact-alternant}
Let \(k\in U(n)\) have eigenvalues \(z_1,\ldots,z_n\). For distinct
eigenvalues the character of \(W_\lambda\) is
\[
        \chi_\lambda(k)=
           \frac{\det[z_j^{\lambda_i+n-i}]_{i,j=1}^n}
                {\det[z_j^{n-i}]_{i,j=1}^n}.
\]
At repeated eigenvalues use its continuous Laurent-polynomial
extension. For every \(k\),
\[
 \chi_\lambda(k)=\dim W_\lambda
                     \int_{U(n)}m_\lambda(x^{-1}kx)\dd x .
\]
The character of \(k\mapsto T(\bar k)\) is
\(\overline{\chi_\lambda(k)}\).
\end{theorem}
\begin{proof}
Every unitary matrix diagonalizes in an orthonormal basis by the
normal-matrix proof of Lemma~\ref{lem:common-hecke-interpolation}.
For a partition \(\alpha\) of length at most \(n\), let \(S_\alpha(z)\)
be the sum of the monomials of all semistandard fillings of its
diagram, with one factor \(z_i\) for each entry \(i\). For the
diagram \(\kappa\) in Definition~\ref{def:compact-flag-space}, the
basis of Proposition~\ref{prop:standard-flag-minors} diagonalizes
the torus action. Each row index in its minors contributes
the corresponding diagonal coordinate, so the character equals
\((\prod_i z_i)^{\lambda_n}S_\kappa(z)\).

For distinct nonzero variables write
\(\Delta_n(z)=\det[z_j^{n-i}]=\prod_{j<l}(z_j-z_l)\) by the
Vandermonde calculation in Proposition~\ref{prop:exact-row-degree}.
We prove by induction that
\(\Delta_n S_\alpha=\det[z_j^{a_i}]\) with \(a_i=\alpha_i+n-i\).
It holds in one variable. In the removal of entries \(n\) in
Lemma~\ref{lem:tableaux-binomial}, write \(s=z_n\) and
\(u_i=\gamma_i+n-1-i\) for the remaining diagram \(\gamma\).
Removed entries contribute \(s^{\sum_i a_i-\sum_i u_i-(n-1)}\).
Summing the determinants from the induction hypothesis in their
independent rows therefore gives
\[
 \Delta_{n-1}(z_1,\ldots,z_{n-1})\,S_\alpha(z)
  =s^{\sum_i a_i-(n-1)}
        \det\!\left[\sum_{b=a_{i+1}}^{a_i-1}(z_j/s)^b\right]_{i,j<n}.
\]
Evaluate each geometric sum and multiply by
\(\prod_{j<n}(z_j-s)\). The result is \(s^{\sum_i a_i}\) times
the determinant with entries
\((z_j/s)^{a_i}-(z_j/s)^{a_{i+1}}\).
This is \(\det[z_j^{a_i}]_{i,j\le n}\): take \(s^{a_i}\)
out of each row, then subtract the next row successively from
top to bottom, leaving the last column the last coordinate vector.
This proves the formula for \(\alpha\), and multiplication by the
determinant power gives \(\lambda_i+n-i\) in place of
\(\kappa_i+n-i\). The tableau expression is a Laurent polynomial
on the nonzero coordinates and therefore also gives the continuous
value at repetitions.

For the integral, write \(m_\lambda(x^{-1}kx)=
\langle T(x)v_\lambda,T(k)T(x)v_\lambda\rangle\).
The average of the projections on \(T(x)v_\lambda\) is
\(I/\dim W_\lambda\) by Proposition~\ref{prop:compact-formal-degree}.
Taking the trace after composition with \(T(k)\) gives the identity.
The eigenvalues of \(\bar k\) are \(\bar z_i\); the tableau polynomial
has integer coefficients, giving the last assertion.
\end{proof}

\begin{proposition}\label{prop:definite-elliptic-trace}
In Definition~\ref{def:definite-comparison-forms} set \(K=K(g)\)
and choose a representative \(h_c\) for each class at level \(K\).
Let \(\varphi:H_f\to\C\) have finite support modulo right \(K\)
and be invariant on both sides by \(K\). Define
\[
    (\mathcal T_\varphi L)(h)
        =\sum_{yK\in H_f/K}\varphi(y)L(hy).
\]
Writing \(\chi_W\) for the character of \(R_\infty\), one has
\[
    \operatorname{tr}_{\mathcal A^+}\mathcal T_\varphi
       =\sum_c\ \sum_{\gamma\in H(F)}
          \varphi(h_c^{-1}\gamma h_c)\,\chi_W(\gamma_\infty).
\]
At each \(c\) at most \(|{\rm supp}(\varphi)/K|\) rational elements
have a nonzero coefficient. Suppose \(\varphi\) is supported on
local matrices which, at some chosen component above \(g_0\),
belong to \({\rm GL}_r(\Z_{g_0})\) and reduce to a matrix with
irreducible characteristic polynomial. Then each contributing
\(\gamma\) has irreducible characteristic polynomial over \(E\)
and distinct eigenvalues at every infinite place. Its compact
factor is the product over \(\tau\) of the conjugates of the
alternant ratios in Theorem~\ref{thm:compact-alternant}.
\end{proposition}
\begin{proof}
Right invariance of \(L\) makes the operator independent of
representatives in its sum. Left \(K\)-invariance of \(\varphi\)
makes its output right \(K\)-invariant by reindexing, and its
rational transformation law follows from that for \(L\).
Evaluate in the direct sum of copies of \(W\) constructed in
Theorem~\ref{thm:definite-comparison-rank}. For the \(c\)-th diagonal
block, \(yK\) contributes precisely when \(h_c y=\gamma h_c k'\)
for \(\gamma\in H(F),\,k'\in K\). Such a \(\gamma\) is unique by
triviality of \(H(F)\cap h_c K h_c^{-1}\) from that theorem.
Conversely its coset \(h_c^{-1}\gamma h_c K\) is uniquely determined
by \(\gamma\). The operator on that copy of \(W\) is
\(\varphi(y) R_\infty(\gamma_\infty)\), giving the trace and the
finite count.

The characteristic polynomial at the chosen place equals the image
of the rational characteristic polynomial, by conjugation and the
determinant definition. For a monic polynomial over a discrete
valuation field with integral coefficients and irreducible reduction
of unchanged degree, a factorization into two positive-degree
monic polynomials is impossible. In fact the minimum coefficient
valuation of a product is the sum of the two minima: divide the
factors by their minimum powers of a uniformizer and use that
the product of their nonzero residue polynomials is nonzero.
Each monic factor has minimum at most zero, so for a monic
integral product both minima must be zero. Reducing the factors
would factor the reduction in their original positive degrees.
Thus factorization over \(E\) is also impossible by the embedding
of \(E\) into this local field. In characteristic zero a monic
irreducible polynomial is relatively prime to its lower degree,
nonzero derivative. Its roots are therefore distinct at each
embedding. The eigenvalues have modulus one because the matrices
are unitary at the infinite places. Apply
Theorem~\ref{thm:compact-alternant} and multiply the traces on the
tensor factors of \(W\).
\end{proof}

\section{The field of an elliptic term}

\begin{theorem}\label{thm:elliptic-trace-fields}
Let \(E/F\) be the CM extension of Definition~\ref{def:levels}
with its conjugation, and let \(\gamma\in{\rm GL}_r(E)\) have
irreducible characteristic polynomial. Assume that at every
embedding of \(E\) in the chosen CM type all complex eigenvalues
of \(\gamma\) have absolute value one, and that \(\gamma\)
preserves some nonsingular hermitian form \(h\) on \(E^r\).
Put \(L=E[\gamma]\). It is a degree-\(r\) field over \(E\).
The involution of \(E\) extends to an involution \(a\mapsto a^*\)
on \(L\) with \(\gamma^*=\gamma^{-1}\). Its fixed field \(K\) is
totally real, and \(L=K(i)\) is totally imaginary quadratic over
\(K\). In particular \([K:F]=r\).

For any identification \(E^r=L\) intertwining \(\gamma\) with
multiplication by that element, the nonsingular hermitian forms
preserved by \(\gamma\) are exactly
\[
              h_a(u,v)=\operatorname{Tr}_{L/E}(a u^* v),
                                  \qquad a\in K^\times .
\]
Here \(\operatorname{Tr}_{L/E}\) is the trace of multiplication
on the \(E\)-space \(L\). The automorphisms of this space commuting
with \(\gamma\) are the multiplications by \(z\in L^\times\);
they send \(h_a\), by pullback, to \(h_{a z^*z}\).
The sign on the eigenline at an embedding
\(\sigma:L\to\C\) over a chosen embedding of \(E\) is
\(\operatorname{sign}\sigma(a)\). There is one such eigenline for
each of the real embeddings of \(K\).
\end{theorem}
\begin{proof}
Write \(P\) for the characteristic polynomial. The adjugate
argument of Lemma~\ref{lem:torus-primary} gives \(P(\gamma)=0\).
The kernel of evaluation \(E[X]\to E[\gamma]\) is generated by
a nonconstant divisor of \(P\), by division by a monic polynomial
of least degree in that kernel. Irreducibility gives
\(L\simeq E[X]/P\). The action on any nonzero vector gives an
injective \(E\)-linear map \(L\to E^r\), since nonzero field
elements act invertibly. Equality of dimensions gives the
identification in the statement. A linear operator commuting
with \(\gamma\) commutes with every polynomial in \(\gamma\),
hence acts on \(L\) as multiplication by its value at \(1\).

The adjoint for \(h\) sends \(\gamma\) to \(\gamma^{-1}\) and
scalars to their conjugates. Existence and uniqueness of the
adjoint follow from the invertible Gram matrix. Since \(P(0)\ne0\),
\(\gamma^{-1}\) is a polynomial in \(\gamma\). Thus the adjoint
restricts on the commutative algebra \(L\) to the field involution
specified on its generators. For any complex embedding \(\sigma\)
of \(L\) over the chosen type, \(\sigma(\gamma)\) is a root of
the image of \(P\), so
\(\sigma(\gamma^{-1})=\overline{\sigma(\gamma)}\).
The same intertwining with complex conjugation holds on \(E\)
and therefore on \(L\). For an embedding over the other type
compose with complex conjugation to get the same identity.
Every embedding of \(K\) into \(\C\) extends to \(L\), by taking
successive roots of minimal polynomials in \(\C\) as in
Lemma~\ref{lem:place-fields}. It thus has real image. For
\(z\in L\) the elements \((z+z^*)/2,\ (z-z^*)/(2i)\) belong to
\(K\), and \(i^*=-i\). This proves \(L=K\oplus iK\) and also
\(E=F\oplus iF\). Comparing degrees gives \([K:F]=r\).
Each embedding of \(K\) has two extensions to \(L\), interchanged
by complex conjugation since the two send \(i\) to opposite roots
of \(-1\). Precisely one extends the chosen type on \(E\).

Given any other preserved hermitian form, its adjoint on \(L\)
is the same by the equations on \(E,\gamma\). Let \(f(v)=h(1,v)\).
Then \(h(u,v)=f(u^*v)\). The pairing
\((a,v)\mapsto\operatorname{Tr}_{L/E}(av)\) identifies \(L\)
with its \(E\)-linear dual: if \(a\ne0\), test at \(v=a^{-1}\)
to get the nonzero trace \(r\). The dimensions agree. Hence
\(f(v)=\operatorname{Tr}_{L/E}(av)\) for a unique \(a\).
Conjugating the entries of the multiplication matrix and changing
to the starred basis gives
\(\overline{\operatorname{Tr}_{L/E}v}=\operatorname{Tr}_{L/E}(v^*)\).
The hermitian condition consequently says \(a^*=a\), using
nondegeneracy of the trace pairing. Every nonzero such \(a\)
gives a nonsingular form with \(\gamma\) unitary, by
\(\gamma^*\gamma=1\), and multiplication by \(z\) gives the
displayed pullback law.

At an embedding \(\tau:E\to\C\), evaluation at all the roots of
\(\tau P\) identifies \(L\otimes_{E,\tau}\C\) with \(\C^r\).
Its matrix in the power basis is a Vandermonde at distinct roots
by separability, as in Proposition~\ref{prop:definite-elliptic-trace}.
In those coordinates trace is summation, and
\(\sigma(u^*)=\overline{\sigma(u)}\). The orthogonal diagonal
entries of the form are therefore precisely \(\sigma(a)\).
\end{proof}

\begin{corollary}\label{cor:elliptic-norm-sign-condition}
With \(L/K\) as in Theorem~\ref{thm:elliptic-trace-fields},
fix signs \(\varepsilon_\sigma\in\{\pm1\}\) at the real
embeddings of \(K\). Suppose \(d\in K^\times\) has these
signs and there is a finite idele \(z\) of \(L\) with
\(z^*z=d\). Here a finite idele has nonzero components at every
finite place and integral unit components at all but finitely
many places. Multiplication \(h_a\mapsto h_{da}\)
changes precisely the specified eigenline signs. All finite
forms and isometries with their commuting matrix \(\gamma\) are
carried to each other by the single finite adelic multiplication
by \(z\). In particular if finite lattice frames are chosen on
the first form they can be chosen on the second with the same
framed Gram matrix and the same matrix of \(\gamma\).
\end{corollary}
\begin{proof}
At each infinite place the sign of \(\sigma(da)\) is the product
of the two signs, by Theorem~\ref{thm:elliptic-trace-fields}.
At the finite places the identity
\(h_a(zu,zv)=h_{da}(u,v)\) follows by the same trace calculation
after scalar extension. Thus \(z^{-1}\) carries frames on
\(h_a\) to frames on \(h_{da}\) with the same Gram matrix.
It commutes with \(\gamma\), so the framed matrix of \(\gamma\)
is unchanged at every place. Since \(z\) is a finite idele these
are legitimate adelic changes of frames. The inverse change
uses \(d^{-1},z^{-1}\).
\end{proof}

\section{Ideals in the fields of elliptic elements}

\begin{definition}\label{def:general-number-field}
For a finite extension \(P/\Q\) of degree \(D\), let \(O\) be its
ring of algebraic integers. Use all real embeddings and one
embedding from each nonreal conjugate pair to form
\(P_\infty=\R^{r_1}\times\C^{r_2}\), with the ordinary product of
length and area. Put \(e_i=1\) at a real and \(e_i=2\) at a complex
coordinate. A nonzero fractional ideal is a nonzero \(O\)-submodule
of \(P\) contained in \(m^{-1}O\) for some positive integer \(m\).
Those contained in \(O\) are called integral. Products of ideals
are the additive spans of pairwise products. For an integral
ideal \(\mathfrak a\ne0\) its norm is \(N\mathfrak a=[O:\mathfrak a]\).
For \(a\in P\) write \(N_{P/\Q}a,\operatorname{Tr}_{P/\Q}a\)
for the determinant and trace of multiplication.
\end{definition}

\begin{lemma}\label{lem:number-field-lattice}
The image of \(O\) is a full lattice of some covolume \(V_P>0\)
in \(P_\infty\), and \(D=r_1+2r_2\). For every \(a\ne0\),
\[
                  |N_{P/\Q}a|=\prod_i|\sigma_i(a)|^{e_i}.
\]
Every nonzero integral ideal \(\mathfrak a\) has finite index
and lattice covolume \(V_P N\mathfrak a\).
\end{lemma}
\begin{proof}
Each minimal polynomial in characteristic zero has distinct roots
since it is relatively prime to its derivative. Adjoining a finite
list of generators of \(P\) gives exactly \(D\) embeddings into
\(\C\), by extending at each stage by sending a generator to every
root of the transported minimal polynomial. Some rational linear
combination of the generators has distinct values under all
embeddings: the equality for a pair cuts out a proper linear
condition on its coefficients, and a finite product of the resulting
nonzero linear polynomials cannot vanish on all rational tuples
by the root argument in Proposition~\ref{prop:standard-flag-minors}.
This element has degree \(D\) and generates \(P\). Multiplying
it by an integer clears the denominators in the rescaled monic
polynomial, so we may take an integral generator \(u\).
The embedding matrix on \(1,u,\ldots,u^{D-1}\) is a nonsingular
Vandermonde. In the complex embedding coordinates every
multiplication matrix diagonalizes with its conjugates as
diagonal entries. This gives the norm formula and trace as the
sum of embeddings. Conjugate embeddings pair up and replacing
each conjugate pair by real and imaginary coordinates gives
a real nonsingular matrix.

Algebraic integers form a ring: for integral elements \(a,b\)
their powers below the degrees of their two monic equations
span \(\Z[a,b]\) as a finite \(\Z\)-module. Multiplication on
this module and the adjugate determinant trick of
Lemma~\ref{lem:brauer-test-evaluation} give a monic integer
equation for each member. A rational algebraic integer is an
integer by clearing denominators in its monic equation.
Consequently the trace of an algebraic integer is an integer.
Writing \(L_0=\Z[u]\), we have
\(L_0\subset O\subset\{x\in P:\operatorname{Tr}_{P/\Q}(xL_0)
\subset\Z\}\). The right side is a lattice by the invertible
trace matrix in the power basis, the product of the transpose
of the embedding matrix and that matrix. It contains \(L_0\)
with finite index. It follows that \(O\) is a full lattice as
well; the free-basis argument for subgroups of integer lattices
appears in Lemma~\ref{lem:rd}. Any nonzero ideal containing \(a\)
contains \(aO\), a full sublattice of index \(|N_{P/\Q}a|\)
by the multiplication determinant. The assertion on covolumes
then follows by counting the cosets in a fundamental region,
as in Lemma~\ref{lem:class-covol}.
\end{proof}

\begin{proposition}\label{prop:general-ideal-factorization}
Every nonzero fractional ideal is uniquely a product of integer
powers of nonzero prime ideals with only finitely many nonzero
exponents. An ideal is integral exactly when all exponents are
nonnegative. Norm extends multiplicatively to fractional ideals,
with \(N((a))=|N_{P/\Q}a|\). At a prime \(\mathfrak p\), the
local ring \(O_{\mathfrak p}\), obtained by inverting
\(O\setminus\mathfrak p\), is a discrete valuation ring. Its
valuation on \(P^\times\) is the exponent of \(\mathfrak p\)
in the ideal of the element.
\end{proposition}
\begin{proof}
A nonzero prime ideal is maximal since its quotient is a finite
domain by Lemma~\ref{lem:number-field-lattice}, hence a field.
Every nonzero integral ideal contains a product of nonzero
prime ideals, allowing the empty product for \(O\). Otherwise
choose a counterexample maximal by inclusion: every proper
enlargement has strictly smaller finite index. It is not prime,
so some \(a,b\) outside it have product in it. The two larger
ideals formed by adjoining \(a\) or \(b\) have product contained
in it. Using the assertion for both gives a contradiction.

For a nonzero prime \(\mathfrak p\) set
\(\mathfrak p^{-1}=\{x\in P:x\mathfrak p\subset O\}\).
It is fractional, since a nonzero integer belongs to
\(\mathfrak p\) by its finite index. Fix \(0\ne a\in\mathfrak p\)
and choose a product of as few primes as possible contained in
\((a)\). One factor must lie in \(\mathfrak p\) by primality,
hence equals it. The product of the remaining factors contains
some \(b\notin(a)\) by minimality, so
\(b/a\in\mathfrak p^{-1}\setminus O\).
Now \(\mathfrak p\subset\mathfrak p\mathfrak p^{-1}\subset O\).
If the first inclusion were equality, multiplication by \(b/a\)
would preserve the full integer lattice \(\mathfrak p\).
Its integral matrix satisfies a monic polynomial by the adjugate
identity, contradicting \(b/a\notin O\). Thus
\(\mathfrak p\mathfrak p^{-1}=O\).

For a proper nonzero ideal \(\mathfrak a\) take a maximal ideal
\(\mathfrak p\) containing it, possible by finiteness of the
quotient. Then \(\mathfrak a\mathfrak p^{-1}\) is integral and
strictly contains \(\mathfrak a\), by the same lattice argument
using \(b/a\notin O\) above. Iterate at decreasing positive
indices and multiply back the invertible primes. This proves
factorization. If two products coincide, any prime in one contains
the product in the other, hence contains and equals a factor
there. Cancel the invertible factors successively to get uniqueness.
Clearing a rational integer denominator proves the assertion for
fractional ideals and its integrality test.

For any invertible integral ideal \(J\), choose \(x\in J,\ y\in J^{-1}\)
with \(xy\notin\mathfrak p\), since \(1\) is a sum of such products.
Multiplying \(y\) by a suitable integral coefficient makes
\(xy\equiv1\bmod\mathfrak p\). Then multiplication by \(x\)
induces \(O/\mathfrak p\simeq J/\mathfrak p J\), with inverse
multiplication by \(y\), by the displayed congruence on both
modules. This gives
\([J:\mathfrak pJ]=|O/\mathfrak p|\); successive indices
show norm multiplicativity and extension to integer exponents.
The statement for principal ideals follows from
Lemma~\ref{lem:number-field-lattice} after clearing denominators.

Write \(v(x)\) for the indicated exponent. The element \(x\)
lies in \(O_{\mathfrak p}\) exactly when \(v(x)\ge0\):
its denominator ideal \(\{b\in O: bx\in O\}\) has prime exponents
\(\max(0,-v_{\mathfrak q}(x))\), by the integrality test.
When \(v(x)\ge0\) this ideal is not contained in \(\mathfrak p\),
and conversely a denominator outside \(\mathfrak p\) has exponent
zero there. Units correspond to \(v(x)=0\). Any element of
\(\mathfrak p\setminus\mathfrak p^2\) has value one, and such
elements exist by the index just proved. Thus every element
of the local ring is a power of such a uniformizer times a
unit, and any nonzero ideal there is generated by an element
of least valuation.
\end{proof}

\begin{definition}\label{def:signed-ray}
The ordinary class group is the group of nonzero fractional
ideals modulo principal fractional ideals. For a totally real
field \(P\), let \(\mathcal C_4(P)\) be the group of fractional
ideals with exponent zero at every prime over \(2\), modulo the
principal ideals \((a)\) with \(\sigma_i(a)>0\) for every real
embedding and \(a\equiv1\bmod4O_{\mathfrak p}\) for every
\(\mathfrak p\mid2\).
\end{definition}

\begin{lemma}\label{lem:bounded-ideal-classes}
There are finitely many integral ideals of bounded norm. The
two class groups in Definition~\ref{def:signed-ray} are finite,
and each of their classes has an integral representative, prime
to \(2\) in the case of \(\mathcal C_4(P)\).
\end{lemma}
\begin{proof}
An ideal of norm \(j\) contains \(jO\), since its index-\(j\)
additive quotient is annihilated by \(j\). Ideals of that norm
are thus among finitely many subgroups of \(O/jO\).
Given an integral ideal \(\mathfrak a\ne0\), the cube argument
of Lemma~\ref{lem:trace-finite} on its real embedding lattice of
covolume \(V_P N\mathfrak a\) gives \(0\ne a\in\mathfrak a\)
with each real and imaginary coordinate bounded by
\(C_P(N\mathfrak a)^{1/D}\). Indeed two points of a cube with
volume larger than the covolume are congruent modulo the lattice,
by integration over its fundamental parallelepiped. Their
difference has the asserted size. By the norm formula,
\(|N_{P/\Q}a|\le C'_P N\mathfrak a\).
The integral ideal \((a)\mathfrak a^{-1}\) then has norm at most
\(C'_P\) by Proposition~\ref{prop:general-ideal-factorization}
and represents the inverse ordinary class. Clearing denominators
shows all ordinary classes occur this way.

An element \(x\) integral at every prime over \(2\) can be written
\(b/m,\ b\in O,\ m\) a positive odd integer. Its denominator ideal
in Proposition~\ref{prop:general-ideal-factorization} has odd
index: the norm of a prime over an odd rational prime is odd,
since its finite residue field has this characteristic. Every
residue characteristic is a rational prime since the quotient
is a finite field. The denominator ideal contains its own index,
which gives such an \(m\). Thus \(x\) has a well-defined residue
in \(O/4O\). A unit at all primes over \(2\) has invertible
residue by the same fact for \(x^{-1}\). A residue of one is
equivalent to the local congruences of
Definition~\ref{def:signed-ray}, by factorization of \(4O\)
and the just constructed odd denominator. On the group of
principal ideals prime to \(2\), the further ray relation has
finite index: the signs and invertible residue of a generator
take only finitely many joint values, and two generators with
the same values have quotient totally positive and congruent to
one. Together with finiteness of ordinary classes this proves
finiteness of \(\mathcal C_4(P)\). Finally a fractional ideal
prime to \(2\) can have its denominators cleared by a positive
odd integer, by applying the denominator argument to a finite
integer generating set obtained from
Lemma~\ref{lem:number-field-lattice}. Multiplying by the square
of that integer clears denominators and is positive and
congruent to one modulo \(4\), so preserves its signed ray class.
\end{proof}

\section{Counting ideals with signs and a residue condition}

\begin{lemma}\label{lem:unit-log-domain}
In the notation of Definition~\ref{def:general-number-field} put
\[
     \mathcal H=\{y\in\R^{r_1+r_2}:\sum_i e_i y_i=0\},\qquad
     l(u)=(\log|\sigma_i u|)_i\quad(u\in O^\times).
\]
The image of \(l\) is discrete and spans \(\mathcal H\). The same
holds on any finite-index subgroup \(U\) of \(O^\times\).
One can choose in \(U\) units whose logarithms are a real basis
of \(\mathcal H\); the group \(U_0\) they generate has finite
index in \(U\) and its logarithms form a lattice with a bounded
half-open fundamental parallelepiped.
\end{lemma}
\begin{proof}
The norm of a unit has absolute value one by
Proposition~\ref{prop:general-ideal-factorization}, so its logarithm
vector lies in \(\mathcal H\). Bounded logarithms imply bounded
coordinates in the lattice \(O\subset P_\infty\) of
Lemma~\ref{lem:number-field-lattice}. Only finitely many units
satisfy a given such bound; this applies also to the kernel of \(l\).

For \(y\in\mathcal H\), multiply each archimedean coordinate of
\(O\) by \(\exp(y_i)\). Its real covolume is unchanged by the
defining equation of \(\mathcal H\), using two real dimensions
per complex place. By the cube argument of
Lemma~\ref{lem:bounded-ideal-classes} there is \(0\ne a\in O\)
with all \(|\sigma_i(a)|\exp(y_i)\le C\), for a fixed \(C>1\).
Their product with the indicated multiplicities is the positive
integer \(|N_{P/\Q}a|\), between \(1\) and \(C^D\). This gives
a positive lower bound for each factor, depending only on \(C,D\).
The possible principal ideals \((a)\) form a finite set by
Lemma~\ref{lem:bounded-ideal-classes}. Fix one generator per
ideal; then \(a\) equals one of these generators times a unit.
It follows that \(y\) lies within a fixed bounded distance of
the negatives of the logarithms of units. Their linear span
must be all of \(\mathcal H\), since multiples of a vector
orthogonal to a proper span would violate this bound. If a
subgroup has index \(h\), the \(h\)-th power of every unit lies
in it by counting cosets in the abelian quotient. Its logarithms
therefore have the same span. Choose a basis from these logarithms
and choose its unit lifts. Integer combinations are uniquely
represented in this basis, with a half-open parallelepiped for
the quotient. Every coset of \(U_0\) in \(U\) has a representative
whose logarithm is in the bounded parallelepiped. Finiteness of
units with bounded logarithm proves the index assertion, including
when \(\mathcal H=0\).
\end{proof}

\begin{lemma}\label{lem:log-sector-count}
Choose a parallelepiped \(P_0\subset\mathcal H\) as in
Lemma~\ref{lem:unit-log-domain}. For a point of \(P_\infty\)
with all coordinates nonzero put
\[
       t(x)=\left(\prod_i |x_i|^{e_i}\right)^{1/D},\qquad
       y_i(x)=\log |x_i|-\log t(x).
\]
Let \(\mathcal R\) be the set with \(0<t(x)\le1\) and
\(y(x)\in P_0\), with either all real signs allowed or all real
signs prescribed positive. Its ordinary real volume \(v\)
is finite and positive. For every full lattice \(\Lambda\)
in \(P_\infty\), its translate \(z+\Lambda\), and \(Y\ge1\),
\[
   \#\big((z+\Lambda)\cap Y^{1/D}\mathcal R\big)
       =\frac{vY}{\operatorname{covol}\Lambda}
                     +O_{\mathcal R,\Lambda}(1+Y^{1-1/D}).
\]
\end{lemma}
\begin{proof}
The region is parametrized by \(t\in(0,1]\), \(y\in P_0\),
one sign for each real place, and an angle for each complex
place. A real coordinate is \(\pm t\exp y_i\), a complex coordinate
\(t\exp(y_i+i\theta_i)\). The parameter maps on the closures,
with angles in \([0,2\pi]\), are Lipschitz by boundedness of
their first derivatives on compact boxes. The interior contains
a nonempty open set, since \(t,y\) are recovered by the formulas
above and the angles by continuous local arguments at nonzero
coordinates. This also shows the boundary is contained in the
images of finitely many faces of these parameter boxes, including
\(t=0\). Indeed away from those faces the parameter map and
its local inverse give an open neighborhood within the region.
The closures cover all limit points by compactness of the parameter
boxes. Thus the region is bounded and has the stated positive,
finite volume. The number of continuous real parameters is
\(1+(r_1+r_2-1)+r_2=D\).

Write \(T=Y^{1/D}\). Divide each boundary parameter face into
\(O_{\mathcal R}(1+T^{D-1})\) boxes of side at most \(1/T\)
in its at most \(D-1\) coordinates. Their images after dilation
by \(T\) have uniformly bounded diameters. Tile by translates
of a fixed half-open fundamental parallelepiped of \(\Lambda\),
one for each point of \(z+\Lambda\), using these points as
corners of the tiles. The discrepancy between the count times
the covolume and the volume is bounded by the total volume of
tiles whose closures meet the boundary. The closures of all
these tiles lie within their fixed diameter of the bounded
diameter images just obtained. Disjoint interiors and a uniform
upper bound on the volumes of these neighborhoods give the asserted
error, uniformly in the translate. This proof also handles \(D=1\),
where the boundary consists of finitely many points.
\end{proof}

\begin{theorem}\label{thm:partial-ideal-counts}
In any ordinary class of the degree-\(D\) field \(P\), the number
of nonzero integral ideals of norm at most \(Y\ge1\) is
\[
                     c_PY+O_P(1+Y^{1-1/D}),\qquad c_P>0,
\]
with the same constant for every class.
For totally real \(P\), in each class of \(\mathcal C_4(P)\)
the number of integral ideals prime to \(2\) of norm at most
\(Y\) is \(c_{P,4}Y+O_P(1+Y^{1-1/D})\), again with a common
positive constant.
\end{theorem}
\begin{proof}
For the ordinary class choose an integral ideal \(\mathfrak b\)
in its inverse class by Lemma~\ref{lem:bounded-ideal-classes}.
The ideals being counted correspond to \(0\ne a\in\mathfrak b\)
with \(|N_{P/\Q}a|\le YN\mathfrak b\), modulo all units, by
\(\mathfrak a=(a)\mathfrak b^{-1}\) and
Proposition~\ref{prop:general-ideal-factorization}. Choose \(U_0\)
and \(P_0\) for \(U=O^\times\) in Lemma~\ref{lem:unit-log-domain}.
For a unit \(u\), \(y(ua)=l(u)+y(a)\). Thus each \(U_0\)-orbit
has precisely one representative with \(y(a)\in P_0\), since
its logarithms are exactly the basis lattice and \(l\) is
injective on \(U_0\). Multiplication of a nonzero element by
units is a free action, so exactly \([U:U_0]\) such representatives
occur per ideal. Apply Lemma~\ref{lem:log-sector-count} to
\(\Lambda=\mathfrak b,\ z=0\) with all signs allowed, using
\(YN\mathfrak b\) for its argument. Its covolume is
\(N\mathfrak b\,V_P\) by Lemma~\ref{lem:number-field-lattice}.
This gives leading constant \(v/(V_P[U:U_0])\), independent of
\(\mathfrak b\).

For a signed ray class choose an integral \(\mathfrak b\) in its
inverse class prime to \(2\). Now \(\mathfrak a\mathfrak b=(a)\)
with \(a\) totally positive, \(a\in\mathfrak b\), and
\(a\equiv1\bmod 4O\). This equivalence follows from
Definition~\ref{def:signed-ray} and the local congruence criterion
in Lemma~\ref{lem:bounded-ideal-classes}. These last two finite
conditions specify one coset of \(4\mathfrak b\). Indeed
\(\mathfrak b+4O=O\), since any maximal ideal containing both
would contradict the prime factorization; choose \(a_0\in\mathfrak b\)
congruent to one. Also \(\mathfrak b\cap4O=4\mathfrak b\) by taking
prime exponents. The unit group preserving the generator conditions
is \(U=\{u\in O^\times:u>0\text{ at every real place},\
u\equiv1\bmod4O\}\), of finite index by the finite residue and
sign maps. Use its \(U_0,P_0\) in the positive-sign region for
the coset \(a_0+4\mathfrak b\), still with argument \(YN\mathfrak b\).
This time the covolume is \(4^D N\mathfrak b\,V_P\) and the leading
constant is \(v/(4^D V_P[U:U_0])\). The error is uniform over
classes in each assertion by their finiteness.
\end{proof}

\begin{corollary}\label{cor:ray-series-at-one}
The series \(\zeta_P(s)=\sum_{\mathfrak a}(N\mathfrak a)^{-s}\)
over nonzero integral ideals converges absolutely for \(\Re s>1\)
and has a simple pole of positive residue at \(1\). If \(P\) is
totally real and \(\chi\) a nontrivial complex character of
\(\mathcal C_4(P)\), the series
\[
           L_4(s,\chi)=\sum_{(\mathfrak a,2)=1}
                        \chi([\mathfrak a])(N\mathfrak a)^{-s}
\]
over integral ideals is holomorphic on \(\Re s>1-1/D\) by continuation.
Both series have their Euler products over the allowed prime
ideals in \(\Re s>1\).
\end{corollary}
\begin{proof}
For the full sum use the first count of
Theorem~\ref{thm:partial-ideal-counts} and add over all classes.
For the character sum the main terms in the second count cancel:
multiplication by a class on which \(\chi\ne1\) preserves the
complete class sum, so that sum is zero. The respective counting
functions are \(cY+O_P(1+Y^{1-1/D})\), \(c>0\), and
\(O_P(1+Y^{1-1/D})\). Integrate each against
\(sY^{-s-1}\dd Y\) on \([1,\infty)\) as in
Proposition~\ref{prop:cont}. The errors give holomorphic
functions for \(\Re s>1-1/D\) by Lemma~\ref{lem:analytic-tools}.
The first main term gives \(cs/(s-1)\). Absolute convergence
to the right of one follows from the unrestricted count, and
expanding finite geometric products there recovers the series
by unique ideal factorization in
Proposition~\ref{prop:general-ideal-factorization}.
\end{proof}

\section{Split prime classes of an elliptic field}

\begin{lemma}\label{lem:elliptic-finite-norms}
Let \(K\) be a totally real number field and \(L=K(i)\).
Suppose \(L\) contains \(\xi=\exp(2\pi i/12)\) in a fixed complex
embedding, with \(^{*}\) denoting complex conjugation over \(K\).
Then \(O_L=O_K\oplus O_K\xi\). For every prime \(\mathfrak p\)
of \(O_K\), with residue field \(\kappa\) of size
\(Q=N\mathfrak p\), the primes of \(L\) above it are either two
of norm \(Q\) (the split case) or one of norm \(Q^2\)
(the inert case). All have ramification index one. If \(Q\)
is odd the split case occurs exactly when \(Q\equiv1\bmod4\).

Let \(K_{\mathfrak p}\) be the fraction field of the completion
of \((O_K)_{\mathfrak p}\) and let \(L_{\mathfrak p}=
L\otimes_K K_{\mathfrak p}\). In the split case every element
of \(K_{\mathfrak p}^\times\) has the form \(z^*z\); in the
inert case precisely the elements of even valuation have this
form. In both cases a unit has a unit preimage.
\end{lemma}
\begin{proof}
The identity \(\xi-\xi^*=i\) of Lemma~\ref{lem:skew-covol}
holds in the chosen embedding and hence in \(L\).
For \(x\in O_L\), \(b=(x-x^*)/i\) and \(a=x-b\xi\) are fixed
by \(^{*}\) and integral, so belong to \(O_K\).
Conversely all \(a+b\xi\) with integral coefficients are
integral by Lemma~\ref{lem:number-field-lattice}. Independence
follows from \(\xi\ne\xi^*\).
Thus \(O_L=O_K[X]/(P)\), with
\(P=X^2-(\xi+\xi^*)X+1\) of discriminant \(-1\).
Modulo \(\mathfrak p\) this gives two distinct linear factors
or one irreducible quadratic. The prime ideals above
\(\mathfrak p\) correspond to the field factors of
\(\kappa[X]/(\bar P)\), with the stated residue sizes.
Their product is \(\mathfrak pO_L\), since their intersection
is the kernel of reduction and distinct maximal ideals are
comaximal. This proves the ramification assertion by
Proposition~\ref{prop:general-ideal-factorization}.
At an odd prime, \(i=2\xi-(\xi+\xi^*)\) and \(2\) is invertible,
so the same algebra is \(\kappa[X]/(X^2+1)\).
By the cyclicity proof in Lemma~\ref{lem:count-labels},
\(-1\) has a square root in this finite field precisely when its
multiplicative group has order divisible by four.

Put \(A=\widehat{(O_K)_{\mathfrak p}}\), with uniformizer
\(\pi\) as in Proposition~\ref{prop:general-ideal-factorization}.
Taking compatible residues gives the complete algebra \(A[X]/(P)\)
because \(1,X\) are a basis. At each power \(\mathfrak p^j\)
this is the product of the quotients by the \(j\)-th powers of
the primes above \(\mathfrak p\), by the factorization just proved
and the ordinary Chinese remainder argument: comaximal ideals
have comaximal powers by the binomial expansion of \(1=(u+v)^n\),
and two comaximal ideals have intersection equal to their product.
Localizing at a prime leaves each of its prime-power quotients
unchanged since the denominators are invertible there.
The complete algebra is therefore the product of the completed
integer rings above \(\mathfrak p\). By Lemma~\ref{lem:integer-extension}
its split factors are \(A,A\), interchanged by \(^{*}\);
its inert factor is a discrete valuation ring with uniformizer
\(\pi\). In the split case the norm of \((t,1)\) is \(t\).

In the inert residue extension the conjugation interchanges the
roots of \(\bar P\). It is the \(Q\)-th power map, which fixes
\(\kappa\) and has at most \(Q\) fixed elements by the root bound.
Cyclicity in the group of order \(Q^2-1\) makes
\(w\mapsto w^{Q+1}\) onto the group of order \(Q-1\).
The additive map \(w\mapsto w+w^Q\) onto \(\kappa\) is nonzero
since its polynomial has degree \(Q<Q^2\); its image is a
nonzero \(\kappa\)-subspace of \(\kappa\).
Given a unit \(t\in A\), first lift a residue with norm \(\bar t\).
If the norm of a lift \(z_j\) is \(t\) modulo \(\pi^j\), multiply
by \(1+\pi^j c\). Its norm changes by the factor
\(1+\pi^j(c+c^*)\) modulo \(\pi^{j+1}\) for \(j\ge1\).
Surjectivity of residue trace corrects the next digit. The
successive lifts converge to a unit of exact norm \(t\).
Finally \(\pi^*=\pi\), so extracting a power of \(\pi\)
gives exactly the even valuations.
\end{proof}

\begin{theorem}\label{thm:split-classes-kernel}
For \(L/K\) as in Lemma~\ref{lem:elliptic-finite-norms}, let
\(C=\mathcal C_4(K)\) be the signed ray group of
Definition~\ref{def:signed-ray}. Write
\(\chi_4(1)=1,\ \chi_4(3)=-1\) on \((\Z/4\Z)^\times\).
On fractional ideals prime to \(2\) put
\(\omega(\mathfrak a)=\chi_4(N\mathfrak a\bmod4)\), reducing
positive rational norms with odd numerator and denominator.
This descends to \(C\), and
\[
       \ker\omega
       =\big\langle[\mathfrak p]:\mathfrak p\nmid2,\
                         \mathfrak p\text{ splits in }L\big\rangle .
\]
\end{theorem}
\begin{proof}
Norm multiplicativity is Proposition~\ref{prop:general-ideal-factorization}.
If \(a\equiv1\bmod4\) at all dyadic primes, write \(a=1+4c/m\)
with \(c\in O_K\) and \(m\) a positive odd integer by the
denominator argument of Lemma~\ref{lem:bounded-ideal-classes}.
In an integer basis of \(O_K\), the multiplication matrix of \(a\)
is the identity modulo \(4\) over the ring of rationals with odd
denominator. Its determinant \(N_{K/\Q}a\) is one modulo \(4\).
For \(a\) totally positive this determinant is the positive ideal
norm of \((a)\) by Lemma~\ref{lem:number-field-lattice} and
Proposition~\ref{prop:general-ideal-factorization}. Thus \(\omega\)
descends. By Lemma~\ref{lem:elliptic-finite-norms} its value is
\(1\) at the odd split primes and \(-1\) at the odd inert primes.
The subgroup \(S\) on the right is contained in its kernel.

For real \(\sigma>1\), Corollary~\ref{cor:ray-series-at-one}
gives the Euler product for \(\zeta_L(\sigma)\).
Taking the real logarithm and using
\(-\log(1-u)=\sum_{j\ge1}u^j/j\) for \(0\le u<1\)
gives
\[
          \log\zeta_L(\sigma)=
          2\sum_{\mathfrak p\nmid2\ {\rm split}}
                      (N\mathfrak p)^{-\sigma}+O_{K}(1)
                         \quad(\sigma\downarrow1).
\]
Indeed the contribution at exponent one from inert primes is
bounded by \(\sum_{\mathfrak p}(N\mathfrak p)^{-2}
\le\zeta_K(2)\). The higher powers at primes of \(L\) are
bounded by \(2\sum_{\mathfrak P}(N\mathfrak P)^{-2}\le2\zeta_L(2)\).
The finitely many dyadic factors are bounded, by factorization
of \(2O_L\). The positive-residue pole of
Corollary~\ref{cor:ray-series-at-one} now forces divergence of the
displayed split-prime sum as \(\sigma\downarrow1\).

Let \(\lambda\) be a character of \(C/S\), viewed on \(C\).
If \(\lambda\ne1,\omega\), both \(\lambda\) and
\(\lambda\omega\) are nontrivial and their series \(L_4\)
are bounded near one by Corollary~\ref{cor:ray-series-at-one}.
For \(\sigma>1\), multiplying their absolutely convergent Euler
products, with \(Q=N\mathfrak p\) in each factor, gives
\[
 L_4(\sigma,\lambda)L_4(\sigma,\lambda\omega)=
 \prod_{\mathfrak p\nmid2\ {\rm split}}(1-Q^{-\sigma})^{-2}
 \prod_{\mathfrak p\nmid2\ {\rm inert}}
                      (1-\lambda([\mathfrak p])^2 Q^{-2\sigma})^{-1}.
\]
The logarithm of the absolute value of the second product has
absolute value at most \(2\sum_{\mathfrak p}Q^{-2}\), by the
convergent logarithmic power series and \(|\lambda|=1\).
The first product tends to infinity by the split-prime sum.
This contradicts boundedness. Thus every character of \(C/S\)
is \(1\) or \(\omega\). Both are trivial on \(\ker\omega/S\).
Characters of a finite abelian group separate its elements by
Lemma~\ref{lem:characters}, and \(C\) is finite by
Lemma~\ref{lem:bounded-ideal-classes}. Hence
\(\ker\omega/S\) is trivial.
\end{proof}

\section{Changing the real signs of elliptic trace forms}

\begin{theorem}\label{thm:even-sign-norms}
Let \(L/K\) satisfy the hypotheses of Lemma~\ref{lem:elliptic-finite-norms}.
A family \(s_\sigma\in\{1,-1\}\), indexed by the real embeddings
of \(K\), occurs as the signs of an element \(d\in K^\times\)
of the form \(d=z^*z\) at all finite places for a finite idele
\(z\) of \(L\) if and only if \(\prod_\sigma s_\sigma=1\).
When this product equals one, \(d\) may be chosen with
\(d\equiv1\bmod4\) at all primes over \(2\), and with its
fractional ideal supported on split primes away from \(2\).
\end{theorem}
\begin{proof}
Suppose first that \(\prod_\sigma s_\sigma=1\).
The translated full lattice \(1+4O_K\) in
\(K_\infty=\R^{[K:\Q]}\) has a point \(\alpha\) in the open
orthant with signs \(s_\sigma\). Indeed its bounded translated
fundamental tiles cover this real space by
Lemma~\ref{lem:number-field-lattice}, and the point with
coordinates \(Ts_\sigma\) for sufficiently large \(T\) lies
farther from each coordinate hyperplane than a tile diameter.
The norm of \(\alpha\) is positive by the product of its real
embeddings, and is one modulo \(4\) by the determinant argument
of Theorem~\ref{thm:split-classes-kernel}. Thus
\(\omega((\alpha))=1\). The same theorem gives integers
\(e_{\mathfrak p}\), zero except at finitely many odd split
primes, such that \([(\alpha)]=
[\prod_{\mathfrak p}\mathfrak p^{e_{\mathfrak p}}]\)
in \(\mathcal C_4(K)\). By Definition~\ref{def:signed-ray}
their quotient is \((b)\) with \(b\) totally positive and
congruent to one at the indicated modulus. Then \(d=\alpha/b\)
has the required signs, is one modulo \(4\) at the dyadic primes,
and has ideal \(\prod\mathfrak p^{e_{\mathfrak p}}\).
At each nonsplit prime \(d\) is a unit and at each split prime
it is a local norm by Lemma~\ref{lem:elliptic-finite-norms}.
That lemma gives unit preimages except where \(d\) has nonzero
valuation. The choices together are a finite idele.

Conversely, let \(d=z^*z\) at the finite places. We first choose
\(u\in L^\times\) so that \(d'=d/(u^*u)\) is one modulo \(4\)
at each dyadic prime of \(K\). For this it suffices to approximate
\(z\) at all primes \(\mathfrak P\) of \(L\) above \(2\):
if each \(u/z_{\mathfrak P}\) belongs to
\(1+4\widehat{(O_L)_{\mathfrak P}}\), the product with its
conjugate is likewise congruent to one. On the base field this
is the congruence in \(K\), since the ramification indices are
one by Lemma~\ref{lem:elliptic-finite-norms}. For the approximation,
choose \(m\ge0\) with all \(2^mz_{\mathfrak P}\) integral.
For each \(\mathfrak P\) their residues modulo arbitrarily high
powers are represented by elements of \(O_L\): completion
preserves the successive residue rings, and elements outside
\(\mathfrak P\) have inverses modulo \(\mathfrak P^j\) by
comaximality. The Chinese remainder argument in
Lemma~\ref{lem:elliptic-finite-norms} specifies all these residues
by one \(c\in O_L\). Choose the powers strictly above the
valuations of \(2^m z_{\mathfrak P}\), with the differences in
valuation at least \(v_{\mathfrak P}(4)\). Then \(c\ne0\) and
\(u=c/2^m\) satisfies the required relative congruences.

The element \(d'\) is still a norm at every finite place and
has the signs of \(d\), since \(\sigma(u^*u)>0\) at every real
embedding of \(K\). At every odd inert prime the valuation of
\(d'\) is even by Lemma~\ref{lem:elliptic-finite-norms}. All
other odd prime factors have \(\omega=1\), so
\(\omega((d'))=1\) by Proposition~\ref{prop:general-ideal-factorization}
and Theorem~\ref{thm:split-classes-kernel}. Hence the positive
rational number \(|N_{K/\Q}d'|\) is one modulo \(4\).
Also \(N_{K/\Q}d'\equiv1\bmod4\) by the multiplication
determinant computation in that theorem, which uses only the
dyadic congruence for the signed determinant. An odd rational
number and its negative have distinct residues modulo \(4\).
Consequently \(N_{K/\Q}d'>0\), which is precisely the required
product of signs.
\end{proof}

\begin{corollary}\label{cor:elliptic-even-sign-change}
Use the fields and the unitary elliptic element \(\gamma\) of
Theorem~\ref{thm:elliptic-trace-fields}, and fix a real sign
on each of its \(r\) eigenlines at each infinite place of the
chosen CM type. Relative to the signs of \(h_a\), suppose
an even number of the prescribed signs differ. There is a
rational \(\gamma\)-invariant trace form \(h_{a'}\) with exactly
the prescribed signs and a finite change of frames commuting
with \(\gamma\) that preserves both the finite Gram matrix and
the finite matrix of \(\gamma\). In particular, starting with
positive signs everywhere, this applies to any prescribed
selection of \(j\) negative eigenlines at each place, for any
\(0\le j\le r\), when the number of infinite places is even.
\end{corollary}
\begin{proof}
The field \(L\) contains \(E\), which contains \(\xi\) by
Lemma~\ref{lem:skew-covol}. The real embeddings of \(K\)
correspond to these eigenlines by Theorem~\ref{thm:elliptic-trace-fields}.
The quotients of the prescribed signs by the initial ones
have product one. Apply Theorem~\ref{thm:even-sign-norms} to
these quotients and take \(a'=da\). The frame change is given
by Corollary~\ref{cor:elliptic-norm-sign-condition}.
In the last assertion the number of sign changes is \(jf\) where
\(f\) is the even number of infinite places. For the arithmetic
data of Proposition~\ref{prop:places}, \(f=\varphi(M)U\)
is even since \(6\mid M\) and negation pairs the distinct units
modulo \(M\).
\end{proof}

\section{Finite framed classes with prescribed eigenline signs}

\begin{definition}\label{def:fixed-poly-frames}
Fix a polynomial \(P\) occurring in Theorem~\ref{thm:elliptic-trace-fields},
write \(L=E[X]/(P)\), \(y=X\bmod P\), and let \(m_z\) denote
multiplication by \(z\) on \(L\) or its scalar extensions.
Use the resulting fields \(L/K\) and order the extensions of
each embedding in the chosen CM type. Write \(\mathbb A_{E,f}\)
for the finite adelic ring, formed with the completed integral
rings at almost all places, and \(L_f=L\otimes_E\mathbb A_{E,f}\).
The basis of \(L\) over \(E\) identifies this module with the
finite adelic product of the completions of \(L\). Indeed, at
almost all primes a fixed integral power basis spans the full
integer ring locally, since the two full integer lattices have
finite index after clearing denominators by
Lemma~\ref{lem:number-field-lattice}. At a remaining prime the
decomposition follows by taking compatible residues in the
prime-power factorization as in Lemma~\ref{lem:elliptic-finite-norms},
using the ramification exponents in
Proposition~\ref{prop:general-ideal-factorization}.

Choose a nonsingular finite adelic hermitian matrix \(J_f\) and
a compact open subgroup \(K_0\) of
\(G_f=\{g:g^*J_fg=J_f\}\), where the star on matrices is conjugate
transpose. Given a sign assignment \(s\) on the embeddings of \(K\),
consider pairs \((a,x)\), where \(a\in K^\times\) has signs \(s\)
and \(x:\mathbb A_{E,f}^r\to L_f\) is an invertible finite
adelic frame with \(x^*h_a x=J_f\). Take their equivalence
classes under
\[
 (a,x)\ \longmapsto\
 \big(a/(w^*w),\,m_w x k\big)
                 \qquad(w\in L^\times,\ k\in K_0).
\]
Write \(Y_x=x^{-1}m_yx\), \(T(K)=\{u\in L^\times:u^*u=1\}\)
and \(T_x=x^{-1}m_{T(K)}x\subset G_f\).
\end{definition}

\begin{proposition}\label{prop:elliptic-class-description}
Definition~\ref{def:fixed-poly-frames} parametrizes precisely the
rational simultaneous isomorphism classes, up to right \(K_0\)
on the frame, of triples consisting of a nonsingular hermitian
form on an \(r\)-dimensional \(E\)-space, a preserving operator of
characteristic polynomial \(P\), and a finite adelic frame with
Gram matrix \(J_f\) and the prescribed real eigenline signs of the
form. The rational centralizer preserving the form is \(T(K)\),
and its subgroup preserving the frame modulo \(K_0\) is identified
with \(T_x\cap K_0\). This latter group is finite.
\end{proposition}
\begin{proof}
The cyclic module identification in Theorem~\ref{thm:elliptic-trace-fields}
carries the operator to \(m_y\). Its adjoint induces the unique
involution on \(L\) extending conjugation and sending \(y\) to
\(y^{-1}\), so this is the same involution for all triples.
The theorem identifies every form obtained this way with exactly
one \(h_a\) on the fixed model. Its choices of intertwining
isomorphism differ by \(m_w,\ w\in L^\times\). Transporting the
form by \(m_w\) changes \(a\) to \(a/(w^*w)\), and transports
the frame to \(m_wx\), by the pullback rule in the theorem.
The remaining choice of frame modulo level is right \(K_0\).
Conversely each pair gives a triple with the required signs,
so these changes give the full equivalence relation. The same
pullback rule gives the centralizer and the frame stabilizer
as stated.

If \(u\in T(K)\) and \(x^{-1}m_ux\in K_0\), the finite vector
\(u=m_u(1)\) belongs to a compact subset of \(L_f\), by
continuity of evaluation and compactness of \(xK_0x^{-1}\).
Every compact subset of the finite adeles of \(L\) is contained
in \(m^{-1}\widehat O_L\) for one positive integer \(m\), where
\(\widehat O_L\) is the product of the completed integer rings:
the open sets \((j!)^{-1}\widehat O_L\) form an increasing cover
by the restricted product definition. Thus \(mu\in O_L\) by
Proposition~\ref{prop:general-ideal-factorization}. All its
complex embeddings have absolute value \(m\) since \(u^*u=1\)
and \(L/K\) is CM by Theorem~\ref{thm:elliptic-trace-fields}.
The integral lattice of Lemma~\ref{lem:number-field-lattice}
has only finitely many points in this bounded region.
\end{proof}

\begin{theorem}\label{thm:elliptic-class-bijection}
For sign assignments \(s,t\) in Definition~\ref{def:fixed-poly-frames}
with \(\prod_\sigma s_\sigma=\prod_\sigma t_\sigma\), their
framed classes are in bijection. Representatives can be paired
with exactly the same \(Y_x\) and exactly the same subgroup
\(T_x\) of \(G_f\). Hence any right-level invariant weight
depending only on the simultaneous \(K_0\)-conjugacy class of
\((Y_x,T_x)\), including a factor \(1/|T_x\cap K_0|\), has the
same sum on the two sets whenever it is nonnegative or absolutely
summable. The bijection also respects every right translation on
frames before quotienting by \(K_0\).
\end{theorem}
\begin{proof}
Use Theorem~\ref{thm:even-sign-norms} to choose \(d\in K^\times\)
with signs \(t_\sigma/s_\sigma\), and a finite idele \(z\)
with \(z^*z=d\). The map
\[
                (a,x)\longmapsto (da,m_z^{-1}x)
\]
preserves \(J_f\) by Corollary~\ref{cor:elliptic-norm-sign-condition}.
It commutes with the left rational changes in
Definition~\ref{def:fixed-poly-frames} because all multiplications
in \(L_f\) commute, and it commutes with right multiplication on
the frame by associativity. It has inverse given by \(d^{-1},z^{-1}\).
Both \(m_y\) and each \(m_u,\ u\in T(K)\), commute with \(m_z\).
Their conjugates by the new frame are therefore their old
conjugates, term by term. Changing the representative by \(k\)
conjugates all these matrices by the same \(k\).
Proposition~\ref{prop:elliptic-class-description} gives finite
stabilizer orders. Reindexing a nonnegative or absolutely summable
family along a bijection proves the last sum identity.
\end{proof}

\section{Norms and involutions on idele classes}
\begin{definition}\label{def:herbrand-two-quotients}
For an abelian group \(A\) written additively with involution
\(\sigma\), set \(d=1-\sigma,\ N=1+\sigma\) and define
\[
 A^0=\ker d/NA,\qquad A^{-1}=\ker N/dA,\qquad
                    h(A)=|A^0|/|A^{-1}|
\]
when the quotients are finite. In multiplicative notation the
maps are \(d(x)=x/\sigma x,\ N(x)=x\sigma x\).
\end{definition}

\begin{lemma}\label{lem:involution-quotient-rules}
For a short exact sequence \(0\to A\to B\to C\to0\) compatible
with involutions, finiteness of the two quotients for any two
groups implies finiteness for the third and
\(h(B)=h(A)h(C)\). For finite \(A\), \(h(A)=1\).
For finitely generated \(A\),
\(h(A)=2^{r_+-r_-}\), where \(r_\pm\) are the respective
dimensions of the two eigenspaces on \(A\otimes_{\Z}\Q\).
\end{lemma}
\begin{proof}
The endomorphisms \(d,N\) in
Definition~\ref{def:herbrand-two-quotients} have both composites
zero. In the following cyclic sequence the last term is followed
by the first:
\[
         A^0\to B^0\to C^0\to A^{-1}\to B^{-1}\to C^{-1}\to A^0 .
\]
The unnamed connecting arrows take a representative \(c\) killed
by \(D=d\) or \(D=N\), lift it to \(b\in B\), and use the class
of \(Db\in A\). Changing the lift changes this by \(DA\);
changing \(c\) by the image of the other endomorphism can be
lifted by that endomorphism and contributes zero. This proves
well-definedness. At \(B\), if a cycle becomes a boundary in
\(C\), subtract a lift of that boundary using the same endomorphism
to obtain a cycle in \(A\). At \(C\) vanishing of the connecting
class means \(Db=Da\) for some \(a\in A\), so \(b-a\) is a
cycle lifting \(c\). At \(A\), if a cycle becomes a boundary
\(Db\) in \(B\), the image of \(b\) in \(C\) is a cycle with
that connecting class. Consecutive arrows have zero composite
by construction. These arguments prove exactness at all six
positions. Finiteness follows by bounding each remaining term
by its kernel and image. Alternating multiplication of their
orders gives the stated quotient rule.

For a finite group the two identities
\(|\ker D|\,|DA|=|A|\) for \(D=d,N\) give \(h(A)=1\).
For finitely generated \(A\) first remove its finite torsion
subgroup using the exact sequence just proved. In a free lattice
the sum of its integral \(+1\) and \(-1\) eigensubgroups has
finite index because it contains \(2A\), via
\(2x=(x+\sigma x)+(x-\sigma x)\), and their intersection is zero.
The individual contributions for free rank one with the two
involutions are \(2\) and \(1/2\) by direct substitution in the
definition. The torsion-free groups have bases by the subgroup
argument in Lemma~\ref{lem:rd}; the rational eigenspaces are
spanned by clearing denominators. Multiplication gives the formula.
\end{proof}

\begin{lemma}\label{lem:unramified-local-unit-norm-cohomology}
At an unramified quadratic extension of nonarchimedean local
fields \(K'/K\), all units of \(O_K\) are norms of units and
every norm-one unit has the form \(z/\sigma z\) for a unit \(z\).
The same conclusions hold for the split algebra \(K\times K\).
For \(\C/\R\), the norm image is \(\R_{>0}\) and every norm-one
element is of the form \(z/\bar z\).
\end{lemma}
\begin{proof}
In the unramified case one may take the same uniformizer \(\pi\)
in both fields. On the quadratic residue field the norm is
surjective on units and the trace surjective on the additive
group by Lemma~\ref{lem:elliptic-finite-norms} and its
proof. Lift a residue unit of the required norm. To improve a
lift \(z\) correct through \(\pi^m\) to one correct through
\(\pi^{m+1}\), multiply by \(1+\pi^m y\), \(m\ge1\).
Its norm modulo \(\pi^{m+1}\) is \(1+\pi^m{\rm Tr}(\bar y)\),
where \(\bar y\) here means reduction. Trace surjectivity realizes
the needed correction. The successive lifts converge to a
unit of the prescribed norm by completeness. Over any quadratic
field extension a norm-one element \(c\ne-1\) equals
\((1+c)/(1+\sigma c)\). For \(c=-1\) take any nonzero
\(z=a-\sigma a\). At the unramified place divide this solution
\(z\) by the corresponding power of the base uniformizer to
make it a unit while preserving \(z/\sigma z\). In the split
case the norm of \((u,1)\) is \(u\), and a norm-one pair
\((c,c^{-1})\) is the quotient of \((c,1)\) by its conjugate.
At a real place the norm is the squared complex absolute value,
and the quotient construction above applies to every point of
the unit circle.
\end{proof}

\begin{theorem}\label{thm:unramified-quadratic-idele-herbrand}
Let \(M/P\) be a quadratic extension of number fields unramified
at every finite place. Let \(\sigma\) be its nonidentity automorphism
and \(r\) the number of real places of \(P\) becoming complex.
For any number field \(A\) write \(J_A=\mathbb A_A^\times\)
for its ideles and \(C_A=J_A/A^\times\). Then both quotients in
Definition~\ref{def:herbrand-two-quotients} for \(C_M\) are finite,
\[
       h(C_M)=2,\qquad (C_M)^\sigma=C_P,\qquad
       (C_M)^0=C_P/N_{M/P}C_M .
\]
Here \((C_M)^\sigma\) denotes fixed classes under the natural
embedding of \(C_P\); the superscript \(0\) is the quotient from
that definition.
\end{theorem}
\begin{proof}
Let \(J_M^0=M_\infty^\times\prod_{w\ {\rm finite}}O_{M,w}^\times\);
this superscript denotes the subgroup of finite units in this
proof. The ideal map from finite valuations has this kernel,
is onto the fractional ideals, and sends the principal ideles
onto principal ideals by Proposition~\ref{prop:general-ideal-factorization}.
It gives an exact sequence
\[
    1\longrightarrow J_M^0/O_M^\times\longrightarrow C_M
                            \longrightarrow{\rm Cl}(M)\longrightarrow1.
\]
The last group is finite by Lemma~\ref{lem:bounded-ideal-classes}.
For \(J_M^0\) compute the two endomorphisms separately at each
place of \(P\). The finite factors contribute trivial quotients
by Lemma~\ref{lem:unramified-local-unit-norm-cohomology}, including
for the infinite product because solutions to the unit equations
can be chosen at every place inside the product. At a split
infinite place the two factors are exchanged and both quotients
are trivial by the split formula of that lemma. A real place
becoming complex gives quotient orders two and one.
Thus \(h(J_M^0)=2^r\).

By Lemma~\ref{lem:unit-log-domain}, the logarithms of units span
the weighted-sum-zero hyperplane of the real space on the
archimedean places of \(M\), have finite kernel and form a full
lattice in that hyperplane. The real representation of \(\sigma\)
on this space permutes coordinates. Each split infinite place
of \(P\) gives one pair and hence one dimension in both
eigenspaces. Each of the \(r\) other places gives one fixed
coordinate. Passing to the weighted-sum-zero hyperplane removes
one fixed dimension, because the sum maps onto the trivial
real representation. The difference of the two dimensions
is \(r-1\). The rational eigenspaces of the unit lattice have
the same dimensions by solving the rational linear equations
for the two eigenvalues and extending scalars. Therefore
\(h(O_M^\times)=2^{r-1}\) by Lemma~\ref{lem:involution-quotient-rules}.
Applying that lemma twice in the displayed sequence and
\(1\to O_M^\times\to J_M^0\to J_M^0/O_M^\times\to1\)
gives the asserted finiteness and \(h(C_M)=2\).

If \([x]\in C_M\) is fixed, \(x/\sigma x\) is a rational
norm-one element of \(M^\times\). The algebraic quotient
construction in Lemma~\ref{lem:unramified-local-unit-norm-cohomology}
gives \(y\in M^\times\) with \(y/\sigma y=x/\sigma x\).
Now \(x/y\) is a fixed idele, hence lies in \(J_P\) by the
local field embeddings and the diagonal descriptions at split
places. Conversely every idele of \(P\) is fixed, and
\(J_P\cap M^\times=P^\times\), by invariance under \(\sigma\).
This proves the fixed-class identification and its norm-quotient
consequence.
\end{proof}

\section{A quadratic norm principle}
\begin{lemma}\label{lem:split-prime-bound-for-norm-index}
In Theorem~\ref{thm:unramified-quadratic-idele-herbrand} one has
\[
                         [C_P:N_{M/P}C_M]\le 2 .
\]
\end{lemma}
\begin{proof}
Let \(H=C_P/N_{M/P}C_M\), which is finite by
Theorem~\ref{thm:unramified-quadratic-idele-herbrand}. At a finite
place of \(P\) all units are local norms by
Lemma~\ref{lem:unramified-local-unit-norm-cohomology}, and so is
their full product. At infinity every tuple positive at the real
places and arbitrary at the complex places is a norm by the same
lemma. From an ideal prime to \(2\) construct a class in \(H\)
by taking the uniformizer at each prime to its ideal exponent,
and one elsewhere, including infinity. The class is independent
of uniformizers and multiplicative. If the ideal is generated
by a totally positive \(z\equiv1\) at the prescribed modulus
\(2\), these finite valuations, including the zero ones above
\(2\), agree with the principal idele of \(z\); their ratio differs
only by finite units and infinite norms. Thus this map factors
through the signed ray class group \(C_2\) of
Definition~\ref{def:general-signed-ray}. It is onto: for any
idele \(t\) choose \(z\in P^\times\) of the same valuations
as \(t\) at the primes above \(2\) and the same signs at the
real places. Such a \(z\) exists by the open approximation
of Lemma~\ref{lem:arch-finite-approx}. Dividing \(t\) by \(z\)
gives a representative of the required kind after removing
the finite units and infinite norms.

Every class in \(H\) has square one, since the norm of a
base idele embedded in \(J_M\) is its square. Hence \(H\) is
a finite vector space over \(\mathbb F_2\). Its characters
\(\xi:H\to\{1,-1\}\) are obtained by assigning either sign
on a basis. Their average at a class is one at the identity,
zero elsewhere by pairing opposite character values.
For real \(\sigma>1\) write \(L_\sigma=\log(1/(\sigma-1))\)
and \(Q_{\mathfrak p}=N\mathfrak p\). At primes of \(P\)
away from \(2\), let \([\mathfrak p]_H\) be the class just
defined. For \(\xi\ne1\), its pullback to \(C_2\) is nontrivial,
so its ideal series is holomorphic at one by
Proposition~\ref{prop:general-ray-norm-count}. The real logarithm of its absolute Euler product differs from the
sum of its degree-one terms by at most \(2\zeta_P(2)\), since
\(\sum_{k\ge2}Q^{-k\sigma}/k\le2Q^{-2}\) for every prime norm
\(Q\ge2\). Boundedness at one therefore gives
\[
     \sum_{\mathfrak p\nmid2}\xi([\mathfrak p]_H)
                           Q_{\mathfrak p}^{-\sigma}\le O_\xi(1)
                           \qquad(\sigma\downarrow1).
\]
For the trivial character the sum is \(L_\sigma+O(1)\) by
Corollary~\ref{cor:ray-series-at-one}. Averaging these inequalities
over all characters bounds the sum for \([\mathfrak p]_H=1\)
by \(L_\sigma/|H|+O(1)\).

Every split prime lies in that class since a uniformizer on one
branch has the prescribed norm. Apart from finitely many primes,
the sum on split primes of \(P\) is
\(\tfrac12 L_\sigma+O(1)\). Indeed the logarithm of
\(\zeta_M(\sigma)\) equals \(L_\sigma+O(1)\) by
Corollary~\ref{cor:ray-series-at-one}. At a split prime there
are two primes of \(M\) of the same absolute norm, and at an
inert prime the one absolute norm is \(Q_{\mathfrak p}^2\).
This follows from the residue algebra of the quadratic
extension. The sum of all inert contributions and all higher
prime powers is \(O(1)\) by absolute ideal convergence at two.
Comparison of the two bounds as \(L_\sigma\to\infty\)
gives \(|H|\le2\).
\end{proof}

\begin{theorem}\label{thm:unramified-quadratic-local-global-norm}
Let \(M/P\) be a quadratic number field extension unramified at
every finite place. A nonzero \(x\in P\) is the norm of an
element of \(M^\times\) if and only if it is a norm at each
place of \(P\). Equivalently it is positive at every real place
becoming complex and has even valuation at every inert finite
place.
\end{theorem}
\begin{proof}
The norm of a rational element gives its local norms by taking
each embedding. At a split place the norm is the product
of the two components. At an inert place a base uniformizer
has norm its square and all units are norms by
Lemma~\ref{lem:unramified-local-unit-norm-cohomology}. Conversely
the valuation of every local norm there is even since conjugation
preserves the valuation and the common uniformizer has valuation
one in both fields. At a real place becoming complex use the
positive norm assertion in the same lemma. These observations
give the explicit local criterion.

If all local conditions hold, the local choices can be combined
into \(j\in J_M\) with \(Nj=x\), choosing unit solutions where
\(x\) is a finite unit. The class \([j]\) then lies in the norm
kernel on \(C_M\). Theorem~\ref{thm:unramified-quadratic-idele-herbrand}
gives \(h(C_M)=2\). Lemma~\ref{lem:split-prime-bound-for-norm-index}
bounds the order of its numerator by two. Its denominator is
therefore one. By Definition~\ref{def:herbrand-two-quotients}
there exists \([k]\in C_M\) with \([j]=[k]/[\sigma k]\).
For idele representatives this says
\(j=(k/\sigma k)y\) for some \(y\in M^\times\). Taking norms
gives \(x=N_{M/P}y\), as required.
\end{proof}

\begin{proposition}\label{prop:mixed-split-kernel}
Let \(P\) be any number field containing \(a_0\) with \(a_0^2=3\)
and suppose that \(i^2=-1\) has no solution in \(P\).
Put \(M=P(i)\) and let \(\sigma\) fix \(P\) and negate \(i\).
Then \(M/P\) is unramified at all finite places. Every local
unit there is a norm of a unit. One has
\(O_M=O_P\oplus O_P\xi\) for \(\xi=(a_0+i)/2\).
\end{proposition}
\begin{proof}
The element \(\xi\) satisfies \(X^2-a_0X+1=0\) and
\(\xi^4-\xi^2+1=0\), hence is integral. Its difference from its
conjugate is \(i\), a unit. For \(x\in O_M\), both
\(b=(x-\sigma x)/i\) and \(x-b\xi\) are fixed and integral,
so belong to \(O_P\). Subtracting conjugates also proves
uniqueness. Over every completed integer ring the discriminant
of this integral basis is \(-1\). Consequently the quadratic
polynomial has distinct residue roots over the algebraic closure
at every finite prime. Lemma~\ref{lem:integer-extension}
gives either two copies of the base ring or the unramified
quadratic integer ring. The product decomposition into the
completions at the primes of \(M\) is the comaximal ideal argument
of Lemma~\ref{lem:elliptic-finite-norms}. The unit norm assertion
follows from Lemma~\ref{lem:unramified-local-unit-norm-cohomology}.
\end{proof}

\begin{theorem}\label{thm:rational-two-squares}
Let \(P\) be any number field containing a square root of \(3\)
and let \(a\in P^\times\). If \(i\notin P\), there exist
\(x,y\in P\) with \(a=x^2+y^2\) if and only if \(a\) is
positive at every real embedding and is a norm from
\(P(i)\otimes_P P_{\mathfrak p}\) at every finite prime.
If \(i\in P\) every \(a\) has a representation.
\end{theorem}
\begin{proof}
If \(i\in P\), take \(x=(a+1)/2\) and \(y=(a-1)/(2i)\).
Otherwise each element of \(P(i)\) has the form \(x+iy\)
with norm \(x^2+y^2\). Proposition~\ref{prop:mixed-split-kernel}
gives unramifiedness at all finite primes. At a real embedding
the norm image is precisely the positive reals, and at a complex
place the norm is the product in a split algebra. Theorem~\ref{thm:unramified-quadratic-local-global-norm}
now gives both implications.
\end{proof}

\section{Rational vectors of prescribed hermitian norm}

\begin{lemma}\label{lem:arch-finite-approx}
For a number field \(M\), its diagonal image is dense in the
product of \(M_\infty\) and \(M_{\mathfrak p}\) for any finite
set of nonzero prime ideals \(\mathfrak p\). The assertion
also holds for any fixed finite number of coordinates over \(M\).
\end{lemma}
\begin{proof}
Let the targets be \(\eta_\infty,\eta_{\mathfrak p}\).
Choose a positive integer \(m\) clearing the denominators of
all \(m\eta_{\mathfrak p}\) in the corresponding completed
integer rings, using powers of their residue characteristics.
Choose positive integers \(k_{\mathfrak p}\) so large that
a congruence modulo \(\mathfrak p^{k_{\mathfrak p}}\) after
multiplication by \(m\) ensures the required finite precision.
The argument on residues in Theorem~\ref{thm:even-sign-norms}
applies at these primes: every residue in a completion lifts
to \(O_M\) by invertibility of the localization denominators,
and the Chinese remainder argument in
Lemma~\ref{lem:elliptic-finite-norms} combines the congruences.
Thus elements \(b\) congruent to all \(m\eta_{\mathfrak p}\)
form one lattice coset \(b_0+\mathfrak a\), with
\(\mathfrak a=\prod\mathfrak p^{k_{\mathfrak p}}\).
It tiles \(M_\infty\) by bounded translates of a fundamental
parallelepiped by Lemma~\ref{lem:number-field-lattice}.
Choose positive integers \(n\to\infty\) congruent to one
modulo \(\mathfrak a\), for instance \(n=1+jN\mathfrak a\)
since \((N\mathfrak a)O_M\subset\mathfrak a\).
Choose \(b\) in the coset at bounded distance from
\(mn\eta_\infty\). Then \(b/(mn)\) tends to \(\eta_\infty\)
and meets every finite precision since \(n\equiv1\) and is
a local unit there. The empty finite set uses \(\mathfrak a=O_M\).
Apply the scalar construction to each coordinate for the last
assertion.
\end{proof}

\begin{lemma}\label{lem:binary-hermitian-norm}
Let \(F\) be a totally real number field containing \(\sqrt3\),
let \(E=F(i)\) with conjugation \(^{*}\), and let \(h\) be a
hermitian form on an \(E\)-space, linear in the second variable.
Write \(H(w)=h(w,w)\). Let \(u,v\) be two vectors, with
\(\alpha=H(u)\) totally positive and \(\beta=H(v)\ne0\).
Assume \(v_{\mathfrak p}(\beta)=0\) whenever
\(v_{\mathfrak p}(\alpha)\ne0\). Then the \(E\)-span of
\(u,v\) contains a vector of norm one.
\end{lemma}
\begin{proof}
Set \(B=(h(u,v)+h(v,u))/2\in F\), \(\mu=\alpha\beta-B^2\).
If \(\mu\) is a square in \(F\), including zero, choose \(m\in F\)
with \(m^2=\mu\) and put \(M=F\). Otherwise take the quadratic
field \(M=F(m)\) with \(m^2=\mu\).
We show \(\alpha=x^2+y^2\) in \(M\). It is positive at every
real embedding of \(M\) by restriction. When \(i\notin M\)
consider a finite prime of \(M\) above \(\mathfrak p\) of \(F\).
The restriction of its valuation on \(F\) is a positive multiple
of \(v_{\mathfrak p}\): the contraction is a nonzero prime since
the residue characteristic belongs to it, and local units of
\(F\) remain units there, whereas a uniformizer has positive
valuation by contraction and clearing a denominator outside
\(\mathfrak p\). If \(\alpha\) is a unit at \(\mathfrak p\) it
is consequently a unit upstairs and a local norm by
Proposition~\ref{prop:mixed-split-kernel}. Otherwise \(\beta\)
is a local unit and thus a local norm, and
\[
                \alpha\beta=B^2+m^2=(B+im)^*(B+im)
\]
in the quadratic algebra over the completion of \(M\). The
product is nonzero, so \(B+im\) is invertible there; its inverse
is its conjugate divided by this norm. Dividing by the norm
\(\beta\) proves the local norm condition for \(\alpha\).
Theorem~\ref{thm:rational-two-squares} supplies \(x,y\in M\)
in both cases \(i\notin M\) and \(i\in M\).

There are \(p,q,r,s\in F\) satisfying
\[
         p^2+q^2+\mu(r^2+s^2)=\alpha,\qquad pr+qs=0.
\]
For \(M=F\) take \(p=x,q=y,r=s=0\). For \([M:F]=2\) write
\(x=p+rm,\ y=q+sm\) and compare the coefficients of \(1,m\).
Let \(C=p+iq,\ D=-s+ir,\ A=(C-BD)/\alpha\).
The imaginary part of \(\bar C D\) is \(pr+qs=0\), so
\(\bar A D\) lies in \(F\). The cross terms in
\(H(Au+Dv)\) therefore sum to \(2B\bar A D\). Writing
\(|z|^2=\bar z z\) for \(z\in E\), expansion gives
\[
 \begin{split}
    \alpha H(Au+Dv)
       &=|\alpha A+BD|^2+(\alpha\beta-B^2)|D|^2\\
       &=p^2+q^2+\mu(r^2+s^2)=\alpha .
 \end{split}
\]
Division by \(\alpha\) proves the assertion even if the two
chosen vectors are linearly dependent.
\end{proof}

\begin{theorem}\label{thm:hermitian-represent-one}
Let \(E/F\) be as in Lemma~\ref{lem:binary-hermitian-norm}
and \(h\) a hermitian form on a finite-dimensional \(E\)-space.
Suppose it represents \(1\) over \(E\otimes_F F_v\) for every
real and finite completion \(F_v\) of \(F\). Then it represents
\(1\) over \(E\).
\end{theorem}
\begin{proof}
Use a rational \(E\)-basis. At each real embedding choose a
complex vector representing \(1\). Approximate these vectors by
one rational vector \(u\) using Lemma~\ref{lem:arch-finite-approx}
with \(M=E\). The set of vectors with positive norm is open at
each real place by continuity of the hermitian polynomial.
Thus \(\alpha=H(u)\) is totally positive. The primes of \(F\)
at which \(v_{\mathfrak p}(\alpha)\ne0\) form a finite set by
Proposition~\ref{prop:general-ideal-factorization}.
At each of them choose a local vector representing \(1\).
Having unit norm is an open condition there, since the units of
a completed discrete valuation ring are the union of its nonzero
residue classes. The finite algebra of \(E\) is a field or a
product of the completions above \(\mathfrak p\) by
Lemma~\ref{lem:elliptic-finite-norms}; here \(\xi=(\sqrt3+i)/2\)
is integral and gives the basis as in Proposition~\ref{prop:mixed-split-kernel}.
Apply Lemma~\ref{lem:arch-finite-approx} with \(M=E\), including
all the primes above this finite set and the real-place complex
coordinates already specified. We obtain a rational vector \(v\)
with \(\beta=H(v)\) totally positive and a unit at each prime
in the set. Lemma~\ref{lem:binary-hermitian-norm} gives the
required rational vector in their span.
\end{proof}

\section{Rational classification of the hermitian spaces}

\begin{lemma}\label{lem:four-phase-cancellation}
Let \(F\) be a field of characteristic zero, and
\(D=F[X]/(X^2+1)\) with \(i=X\bmod(X^2+1)\) and
\(\bar i=-i\). On a finite free \(D\)-module with nonsingular
hermitian form \(h\), any two vectors \(u,v\) of the same norm
\(a\in F^\times\) are carried to one another by an isometry.
Each splits off a free line with nonsingular orthogonal complement.
\end{lemma}
\begin{proof}
The algebra \(D\) is a quadratic field or a product of two fields.
The four scalars \(1,-1,i,-i\) are distinct in each field factor
and have norm one. Choose one of them, \(\eta\), so that
\[
       b=h(u-\eta v,u)=a-\bar\eta h(v,u)
\]
is a unit of \(D\). In each field factor a nonzero \(h(v,u)\)
forbids at most one choice; a zero forbids none. Thus at most
two choices are forbidden. For \(w=u-\eta v\), we have
\(h(w,w)=b+\bar b\) by equality of the norms of \(u,\eta v\).
The formula
\[
                        Tz=z-w b^{-1}h(w,z)
\]
sends \(u\) to \(\eta v\). In the expansion of \(h(Tz,Tz')\) the
coefficient of \(h(z,w)h(w,z')\) is
\(-b^{-1}-\bar b^{-1}+h(w,w)/(b\bar b)=0\), as in
Lemma~\ref{lem:norm-one-isometry}. It follows that \(T\) preserves
\(h\); the Gram identity gives a left inverse for its square
matrix, hence an inverse by taking determinants in each field
factor. Multiplication by \(\eta^{-1}\) completes the isometry.
The projection \(z\mapsto u a^{-1}h(u,z)\) has image the line
\(Du\) and kernel its orthogonal complement. Its two summands
are free, by taking bases over each field factor, since the
line has rank one at each factor. In such bases nonsingularity
of the whole Gram matrix and of the scalar \(a\) gives
nonsingularity on the complement.
\end{proof}

\begin{theorem}\label{thm:hermitian-local-global}
Let \(F\) be a totally real number field containing \(\sqrt3\),
\(E=F(i)\), and let \(h,h'\) be nonsingular hermitian forms on
\(E\)-spaces of the same finite dimension. If they are isometric
over \(E\otimes_F F_v\) at every finite place and have the same
signature at every real embedding of \(F\), there is an isometry
over \(E\). A rational isometry implies these local conditions.
\end{theorem}
\begin{proof}
Over \(E\), or over \(\C\) at a real place of \(F\), a nonzero
hermitian form has some vector of nonzero norm. Otherwise
expanding the norms of \(x+y\) and \(x+iy\) gives respectively
zero real and imaginary parts of \(h(x,y)\) for all \(x,y\).
Splitting off each nonzero diagonal line by the projection in
Lemma~\ref{lem:four-phase-cancellation} gives an orthogonal
basis for a nonsingular form. At a real place rescale each basis
vector by the inverse square root of the absolute value of its
diagonal entry. This gives a diagonal of \(+1\)'s and \(-1\)'s.
The number of positive entries equals the maximal dimension of
a positive definite subspace, since projection onto the positive
coordinates is injective on any such subspace and those coordinates
themselves give equality. The same holds for negative entries.
This proves both invariance of signature and existence of local
isometries between forms of equal signature.

We induct on the common dimension, the zero-dimensional assertion
being immediate. Choose \(v\) for \(h'\) with norm \(a\ne0\).
Then \(a\in F^\times\), and \(a^{-1}h\) represents \(1\) over
all completions by the given and just constructed local isometries.
Theorem~\ref{thm:hermitian-represent-one} gives \(u\) with
\(h(u,u)=a\) rationally. At each completion any chosen isometry
from \(h'\) to \(h\) sends \(v\) to a vector of norm \(a\).
Compose it with the isometry in Lemma~\ref{lem:four-phase-cancellation}
to send \(v\) to \(u\). Its restriction is an isometry of the
two local orthogonal complements. These are the scalar extensions
of the rational complements, as is seen from the explicit line
projections. Applying the induction hypothesis to the nonsingular
complements and taking the orthogonal sum with \(v\mapsto u\)
gives the required rational isometry. Conversely a rational
isometry extends to each completion by its matrix equation.
\end{proof}

\begin{corollary}\label{cor:fixed-group-elliptic-classes}
Use the notation of Definition~\ref{def:fixed-poly-frames}.
For sign assignments \(s,t\) as in Theorem~\ref{thm:elliptic-class-bijection},
suppose given rational hermitian spaces \(V_s,V_t\), each
with a finite frame of Gram matrix \(J_f\), and having at each
real place the numbers of positive and negative entries specified
respectively by \(s,t\). The bijection of that theorem identifies
the framed conjugacy classes on the two specified rational unitary
groups, subject to these eigenline signs and fixed \(P\). Their
finite matrices and embedded rational centralizers can be paired
as in that theorem.
\end{corollary}
\begin{proof}
The base field \(F\) of the construction is totally real, and
\(E=F(i)\) contains \(\xi=\exp(2\pi i/12)\) by
Lemma~\ref{lem:skew-covol}. Hence \(F\) contains
\(\sqrt3=\xi+\xi^*\). Given \((a,x)\) of signs \(s\), the
trace form \(h_a\) and the specified \(V_s\) are locally
isometric at every finite place by their two frames for \(J_f\).
At infinity the line signs of Theorem~\ref{thm:elliptic-trace-fields}
give precisely the required signature. By
Theorem~\ref{thm:hermitian-local-global} a rational isometry
\(b:(L,h_a)\to V_s\) exists. It carries the operator to
\(b m_y b^{-1}\) and the frame to \(b x\). Another such \(b\)
differs by a rational isometry of \(V_s\), precisely the
equivalence in its framed conjugacy classes. Conversely every
framed class on \(V_s\) arises by
Proposition~\ref{prop:elliptic-class-description}. The finite
matrix in the new frame is still \(x^{-1}m_yx\) and the
centralizer in finite coordinates is still \(T_x\). The same
argument for \(t\) composes with Theorem~\ref{thm:elliptic-class-bijection}
to give the assertion.
\end{proof}

\section{Counting rational operators in a moving lattice}

\begin{definition}\label{def:moving-elliptic-lattice}
Let \(V\) be a nonsingular rank-\(r\) hermitian space for the
\(E/F\) of Definition~\ref{def:levels}, \(r\ge1,\ f=[F:\Q]\).
Fix archimedean frames with Gram matrices \(S_\tau\) diagonal
with entries \(1,-1\), and a finite adelic frame with Gram
matrix \(1\). Let \(H=\prod_\tau U(S_\tau)\), and \(G_f=U(1_r)(\mathbb
A_{F,f})\), the finite isometry group of the identity matrix.
Let \(K^0\) preserve the standard integral finite lattice and
\(K_0\subset K^0\) be compact open. Write
\(X=H\times G_f/K_0,\ \Gamma=U(V)(F)\). Compose the fixed frames
on the right by \(x\in X\), writing \(X_\infty,X_f\) for the
resulting frames, using any lift of the finite coordinate.
For \(\gamma\in\Gamma\) put
\[
   c_\infty(x,\gamma)=X_\infty^{-1}\gamma X_\infty,\qquad
   c_f(x,\gamma)=X_f^{-1}\gamma X_f,\qquad
   m(h)=\prod_\tau\|h_\tau\|_{\rm op}.
\]
Use the ordinary Euclidean norm in each \(\C^r\), and its sum
of squares on their direct sum to measure \(\|v\|_x=
\|X_\infty^{-1}v\|\) for \(v\in V\).
Let \(\Lambda_x=\{v\in V:X_f^{-1}v\in\widehat O_E^r\}\).
The local finite matrix is defined up to \(K_0\)-conjugacy;
\(\Lambda_x\) is independent of the lift.

For a compact set \(\Omega\subset G_f\) invariant under
\(K_0\)-conjugacy and real \(R\ge1\), let \(N_R(x)\) count
the \(\gamma\in\Gamma\) with irreducible characteristic polynomial
over \(E\), \(c_f(x,\gamma)\in\Omega\), and
\(m(c_\infty(x,\gamma))\le R\).
\end{definition}

\begin{lemma}\label{lem:moving-short-vector}
The image of \(\Lambda_x\) in the archimedean frame is a full
real lattice of dimension \(D=2fr\) and covolume independent of
\(x\). Its shortest nonzero length \(\lambda_x\) is attained and
is at most a constant independent of \(x\). For all \(x\) and
all \(\gamma\) with \(c_f(x,\gamma)\in\Omega\) there is one
positive integer \(c=c(\Omega)\), independent of \(x,\gamma\),
such that \(c\gamma\Lambda_x\subset\Lambda_x\).
\end{lemma}
\begin{proof}
Apply Lemma~\ref{lem:class-covol} in a fixed rational basis.
For a finite unitary matrix reducing the comparison of two
frames to that same basis, the determinant over \(E\) has norm
one to \(F\). Its determinant over each \(\Q_p\) therefore has
absolute value one: over a splitting field that determinant
is the product over all embeddings of the \(E\)-determinant,
whose pairs over \(F\) have product one. This can also be seen
by the embedding diagonalization in Lemma~\ref{lem:number-field-lattice}
applied to the finite separable algebra. Thus right movement
of the finite frame by \(G_f\) leaves the covolume unchanged.
At infinity every unitary matrix has determinant of complex
absolute value one by taking the determinant of \(h_\tau^*S_\tau
h_\tau=S_\tau\). Hence its real determinant also has absolute
value one. A bounded ball in a full lattice has finitely many
points, and some ball contains a nonzero point, giving attainment
and positivity. The cube argument of Lemma~\ref{lem:bounded-ideal-classes}
bounds a shortest length in terms of dimension and covolume.
The entries in the compact family \(\Omega\) have denominators
bounded by a positive integer \(c\), by the increasing open
cover used in Proposition~\ref{prop:elliptic-class-description}.
Thus \(c X_f^{-1}\gamma X_f\) preserves the standard integral
module by inclusion, giving the last assertion.
\end{proof}

\begin{proposition}\label{prop:cyclic-orbit-count}
With constants depending on the fixed data and \(\Omega\), one
has uniformly for \(R\ge1,\ x\in X\)
\[
                     N_R(x)\le C R^{D r(r+1)/2}.
\]
If \(N_R(x)>0\) there is an \(E\)-basis \(v_0,\ldots,v_{r-1}\)
of \(V\) in \(\Lambda_x\) with
\(\|v_j\|_x\le C(cR)^j\) for \(0\le j<r\).
\end{proposition}
\begin{proof}
Choose \(u\in\Lambda_x\) of shortest nonzero length
\(\lambda_x\) by Lemma~\ref{lem:moving-short-vector}.
For every \(h_\tau\in U(S_\tau)\) the inverse is
\(S_\tau h_\tau^*S_\tau\) and has the same operator norm, since
\(S_\tau\) and its inverse are ordinary unitary matrices.
Consequently \(1\le\|h_\tau\|_{\rm op}\) by submultiplicativity,
and every factor is at most \(m(h)\).
For a counted \(\gamma\) the vectors \(c^j\gamma^j u\) are
in \(\Lambda_x\) and have norm at most \((cR)^j\lambda_x\)
by iteration. The vectors with \(0\le j<r\) are \(E\)-independent.
Indeed any relation would give a polynomial \(Q\) of degree
less than \(r\) with \(Q(\gamma)u=0\). Its characteristic
polynomial \(P\) is irreducible and annihilates \(\gamma\) by
the adjugate identity used in Theorem~\ref{thm:elliptic-trace-fields}.
For nonzero \(Q\), polynomial division gives \(A Q+B P=1\),
so \(u=0\), a contradiction. These vectors therefore give the
stated basis, using the uniform upper bound on \(\lambda_x\).

Moreover the tuple \((c^j\gamma^j u)_{1\le j\le r}\) determines
\(\gamma\), since the images of the basis
\((\gamma^j u)_{0\le j<r}\) are given by this tuple.
For every \(a\ge1\), at most \((1+2a)^D\) vectors of
\(\Lambda_x\) have norm at most \(a\lambda_x\).
Indeed the open balls of radius \(\lambda_x/2\) about them are
disjoint in the enclosing ball of radius \((a+1/2)\lambda_x\);
compare their Euclidean volumes by dilation. Multiply these
bounds for \(a=(cR)^j,\ 1\le j\le r\), and use \(1+2a\le3a\)
to obtain the claimed exponent.
\end{proof}

\begin{lemma}\label{lem:real-unitary-ball}
For \(S\) any diagonal matrix of \(r\) signs, Haar volume on
\(U(S)\) satisfies
\[
        \operatorname{vol}\{g:\|g\|_{\rm op}\le T\}
                                  \le C_{S} T^{4r^2}\qquad(T\ge1)
\]
for a constant depending also on its fixed normalization.
\end{lemma}
\begin{proof}
The invariant volume in Lemma~\ref{lem:quotient-measure} is
finite on compact sets by its continuous density in local charts.
Choose a relatively compact identity neighborhood containing all
members of \(U(S)\) within a fixed positive operator norm distance
from identity. In the \(2r^2\) real matrix coordinates, entries of
the matrices being counted lie in a box of side \(2T\).
Divide it into boxes of side at most \(c_r/T\), with number
at most \(C_r T^{4r^2}\); choose a regular subdivision with
the least integer number of sides meeting this bound. Two group
elements \(g,g'\) in a small box satisfy
\(\|g^{-1}g'-1\|_{\rm op}\le T\|g'-g\|_{\rm op}\le 2r c_r\),
since \(\|g^{-1}\|_{\rm op}=\|g\|_{\rm op}\) as in
Proposition~\ref{prop:cyclic-orbit-count}. For small fixed
\(c_r\) each nonempty box therefore lies, on the group, in
one left translate of the chosen neighborhood. Adding the
volumes proves the bound.
\end{proof}

\section{Absolute convergence of the geometric trace}

\begin{lemma}\label{lem:cyclic-frame-volume}
In Definition~\ref{def:moving-elliptic-lattice} use the product
of invariant volume on \(H\) and counting measure on \(G_f/K_0\).
Let \(J=r(r-1)/2\) and fix constants \(C,c\ge1\).
For an ordered \(E\)-basis \(B=(v_0,\ldots,v_{r-1})\) of \(V\)
write \(W_R(B)\) for the set of \(x\in X\) such that
\(v_i\in\Lambda_x,\ \|v_i\|_x\le C(cR)^i\) for every \(i\).
For \(R\ge1\) the following bounds hold with constants
independent of \(R,B\).
\[
\begin{aligned}
 \#\{(h(v_i,v_j))_{ij}:W_R(B)\ne\varnothing\}
       &\ll R^{4frJ},\\
 \operatorname{vol}_X W_R(B)&\ll R^{4frJ+4fr^2J}.
\end{aligned}
\]
\end{lemma}
\begin{proof}
Write \(T=(h(v_i,v_j))\). If \(W_R(B)\ne\varnothing\), the
finite gram matrix of the frame is the identity, giving
\(T_{ij}\in O_E\) by Proposition~\ref{prop:general-ideal-factorization}.
At each real place \(|\tau(T_{ij})|\le C^2(cR)^{i+j}\) by
Cauchy--Schwarz with the diagonal sign matrix of operator norm one.
In the full real lattice \(O_E\) of Lemma~\ref{lem:number-field-lattice}
there are \(O_E(t^{2f})\) elements with each conjugate of absolute
value at most \(t\ge1\), by the disjoint-ball argument of
Proposition~\ref{prop:cyclic-orbit-count}, with a fixed shortest
length. Multiply these upper bounds over \(i,j\), using
\(\sum_{i,j}(i+j)=2rJ\).

In coordinates given by \(B\), every possible finite lattice lies
between \(O_E^r\) and \(T^{-1}O_E^r\). The second inclusion holds
because pairing its elements with each \(v_i\) is integral.
The quotient between the two bounds has order
\(n_T=|N_{E/\Q}\det T|\), by applying the determinant of the
integer endomorphism \(T\) to the full lattice \(O_E^r\).
Evaluation under all embeddings splits its rational determinant
as a product of the determinants of \(\tau T\), by the
diagonalization in Lemma~\ref{lem:number-field-lattice}.
There are at most \(n_T^r\) intermediate \(O_E\)-modules.
Indeed at each prime the preimage of a submodule of the finite
quotient is free of rank \(r\) over the local discrete valuation
ring: project onto its first coordinate, lift a generator of
that ideal, and continue on the kernel, as in the integer
lattice proof of Lemma~\ref{lem:rd}. Thus each local submodule
has a generating tuple of length \(r\) in that local quotient.
Its number of choices is at most the local quotient order to
the power \(r\). The local quotient orders multiply to \(n_T\):
the finite module is killed by \(n_T O_E\), and the Chinese
remainder idempotents for the powers of primes in this ideal
decompose it and every submodule into the corresponding local
factors. The same lattice determines its tuple of finite local
lattices by the congruence and localization argument of
Lemma~\ref{lem:class-covol}. Hence it determines at most one
coset modulo \(K^0\). There are \([K^0:K_0]<\infty\) lifts at
level \(K_0\), by compactness and openness.

At every real place put \(M_\tau=X_\tau^{-1}B_\tau\), where
\(B_\tau\) is the matrix of these rational basis vectors in a
fixed rational basis. Its columns have lengths at most \(C(cR)^i\)
in the indicated order and \(M_\tau^* S_\tau M_\tau=\tau T\).
The determinant bound of Lemma~\ref{lem:hadamard} gives
\[
 |\det M_\tau|\ll R^J,\qquad
 |\tau(\det T)|=|\det M_\tau|^2,\qquad
 n_T\ll R^{4fJ}.
\]
This bounds the number of finite coordinates by \(O(R^{4frJ})\).

For the real volume choose a reference archimedean frame \(X^0_\infty\)
in the permitted region, with matrices \(M^0_\tau\).
Then \((X^0_\tau)^{-1}X_\tau=M^0_\tau M_\tau^{-1}\) has operator
norm \(\ll R^J/|\det M_\tau|\). Indeed row \(j\) of the adjugate
of \(M_\tau\) consists of signed minors using all columns except
\(j\), of size \(O(R^{J-j})\). Multiply that row by column \(j\)
of \(M^0_\tau\) and sum. The denominator is fixed for \(T,\tau\)
by its Gram identity. Moreover
\(\prod_\tau|\det M_\tau|^2=|N_{F/\Q}\det T|\ge1\),
since \(\det T\) is a nonzero integer of \(F\). By
Lemma~\ref{lem:real-unitary-ball}, left translation by the
reference frame and multiplication of the bounds at the \(f\)
places give archimedean volume \(O(R^{4fr^2J})\). The individual
radii can be taken at least one by enlarging the fixed constant,
using the determinant bound above. Multiplying by the finite
count proves the result.
\end{proof}

\begin{theorem}\label{thm:elliptic-unfold-convergence}
Let \(D_0\) be a measurable fundamental domain for \(\Gamma\)
in \(X\) of Definition~\ref{def:moving-elliptic-lattice}, with
the measures of Lemma~\ref{lem:cyclic-frame-volume}. Then
\[
             \int_{D_0}N_R(x)\dd x \ll R^{8f(r+1)^4}\qquad(R\ge1).
\]
Let \(\phi\) be a bounded measurable finite test, invariant by
\(K_0\)-conjugacy and supported in \(\Omega\). For a measurable
\(k:H\to\C\) with \(|k(h)|\ll m(h)^{-B}\) and \(B>8f(r+1)^4\),
the sum over the rational operators of irreducible characteristic
polynomial can be integrated absolutely. For representatives of
their rational conjugacy classes it satisfies
\[
 \int_{D_0}\sum_\gamma \phi(c_f(x,\gamma))\,k(c_\infty(x,\gamma))
                 \dd x
   =\sum_{[\gamma]}\int_{\Gamma_\gamma\backslash X}
           \phi(c_f(x,\gamma))\,k(c_\infty(x,\gamma))\dd x ,
\]
where \(\Gamma_\gamma\) is the centralizer in \(\Gamma\).
When the finite test forces irreducibility, these are the full
geometric sums. The bound applies to \(B=r^8\) for the arithmetic
parameters at all sufficiently large scales.
\end{theorem}
\begin{proof}
The construction of the fundamental domain in
Lemma~\ref{lem:quotient-measure} applies here. Countably many
finite coordinates follow from the lattice inclusions with
bounded denominators and finite level index. Properness follows
from bounded rational matrix entries at all places on intersections
of compact sets, and the archimedean action is free. The same
proof applies to any subgroup of \(\Gamma\). In particular, the
integral of a periodized nonnegative function on \(D_0\) equals
its integral on \(X\).

Where \(N_R(x)>0\), Proposition~\ref{prop:cyclic-orbit-count}
supplies a basis \(B\) with \(x\in W_R(B)\) for fixed \(C,c\).
Two rational ordered bases of the same Gram matrix differ by a
unique rational isometry, since the linear map sending basis to
basis preserves the entries of the form. For each possible
Gram matrix choose one such \(B_T\). Hence the indicator of
\(N_R(x)>0\) is bounded by
\(\sum_T\sum_{\delta\in\Gamma}{\bf1}_{W_R(B_T)}(\delta^{-1}x)\).
Integrate and unfold these nonnegative indicators. By
Lemma~\ref{lem:cyclic-frame-volume} their integral is
\(O(R^{8frJ+4fr^2J})\) with \(J=r(r-1)/2\).
The uniform point count of Proposition~\ref{prop:cyclic-orbit-count}
costs the further factor \(R^{fr^2(r+1)}\). Their total exponent
is \(f(2r^4+3r^3-3r^2)\le8f(r+1)^4\), proving the first bound.

On \(2^a\le m(c_\infty)<2^{a+1}\), \(a\ge0\), bound the
absolute sum by \(\|\phi\|_\infty C 2^{-aB}N_{2^{a+1}}\).
The integrated bounds sum as a convergent geometric series.
Thus we can group by conjugacy and integrate termwise by
Lemma~\ref{lem:integral-tools}. For each fixed \(\gamma\),
the translates of \(D_0\) by representatives of
\(\Gamma_\gamma\backslash\Gamma\) tile a fundamental domain
for \(\Gamma_\gamma\). The simultaneous conjugation and left
translation law for \(c_f,c_\infty\) now gives the equality.
If every element of \(\Omega\) has irreducible characteristic
polynomial at some completion, a proper factorization over
\(E\) would persist there. Residue irreducibility of an integral
matrix suffices by Proposition~\ref{prop:definite-elliptic-trace}.
Finally \(f\le x^{12}\) and \(r\ge x^5\) for large scales by
Propositions~\ref{prop:thick} and~\ref{prop:level-prime}.
Thus \(8f(r+1)^4/r^8\le128x^{-8}<1\) eventually.
\end{proof}

\section{Vanishing away from the elliptic set}

\begin{lemma}\label{lem:orbital-cycle}
Use \(H,X,\Gamma\) of Definition~\ref{def:moving-elliptic-lattice}.
For fixed \(\gamma\in\Gamma\), let \(Z=\Gamma_\gamma\backslash X\),
\(C(x)=c_\infty(x,\gamma)\). Right multiplication \(x\mapsto xa\)
by \(a\in H\) on the archimedean coordinate preserves the measure
of \(Z\) and satisfies \(C(xa)=a^{-1}C(x)a\). The finite part
\(c_f(x,\gamma)\) modulo \(K_0\)-conjugacy stays fixed.
For product measure on \(Z\times H\) or \(Z\times H^2\), the
following changes preserve integrals:
\[
 (x,a)\longmapsto (xa,a^{-1}C(x)),\qquad
 (x,a,b)\longmapsto (xa,b,(ab)^{-1}C(x)).
\]
In the first case the two factors \(a,a^{-1}C(x)\) are
cyclically interchanged, and in the second the three factors
\(a,b,(ab)^{-1}C(x)\) are cyclically permuted when the last
factor is recovered from the new base point.
\end{lemma}
\begin{proof}
The rational centralizer commutes with \(\gamma\), proving
invariance of \(C(x)\) on its quotient; the formulas follow from
multiplication of the frames. Right \(H\)-translation preserves
measure on \(X\) by Lemma~\ref{lem:quotient-measure}, commutes
with the left action, and sends a fundamental domain to another.
The independence of the domain, as in
Theorem~\ref{thm:elliptic-unfold-convergence}, proves preservation
on \(Z\). In each product first replace \(x\) by \(x'=xa\)
with \(a,b\) fixed. Now \(a^{-1}C(x)=C(x')a^{-1}\).
For the first transformation invert \(a\) and translate it on
the left at fixed \(x'\). For the second put \(b\) first, and
replace \(a\) by \(b^{-1}C(x')a^{-1}\) at fixed \(x',b\).
These are measure preserving since the measure on \(H\) is
invariant on both sides and under inversion: inversion exchanges
left and right translations and has Jacobian of absolute value
one at identity on its tangent space. Product integration as in
Lemma~\ref{lem:integral-tools} proves the transformations for
nonnegative and absolutely integrable tests. Multiplying the
new factors and using \(C(xa)=a^{-1}C(x)a\) identifies the
recovered final factor as \(a\) in both cases.
\end{proof}

\begin{theorem}\label{thm:orbital-idempotent-vanishing}
In Theorem~\ref{thm:elliptic-unfold-convergence} let \(\gamma\)
have irreducible characteristic polynomial over \(E\).
Let \(B>8f(r+1)^4\) and let \(e:H\to\C\) be measurable, with
\(|e(g)|\le C_0m(g)^{-B}\) on \(H\). Suppose
\[
 \int_H |e(g)|\, m(g)^{B+1}\dd g<\infty,\qquad
        \int_H e(a)e(a^{-1}g)\dd a=e(g)\quad(g\in H).
\]
If some eigenvalue of \(\gamma\) at an archimedean factor has
absolute value different from one, then for every finite test
\(\phi\) as in that theorem
\[
        \int_{\Gamma_\gamma\backslash X}
                   \phi(c_f(x,\gamma)) e(c_\infty(x,\gamma))\dd x=0.
\]
\end{theorem}
\begin{proof}
Write \(Z,C(x)\) as in Lemma~\ref{lem:orbital-cycle} and
\(\Phi(x)=\phi(c_f(x,\gamma))\). The preceding absolute
convergence theorem gives
\(\int_Z|\Phi(x)|m(C(x))^{-B}\dd x<\infty\) by taking the
real archimedean test \(m^{-B}\) and the absolute finite test.
At a place with the given eigenvalue \(\lambda\), choose a nonzero
left eigenrow \(v\) for the matrix \(\gamma_\tau\) in the fixed
archimedean frame. All eigenvalues there are simple by irreducibility
and characteristic zero as in Theorem~\ref{thm:elliptic-trace-fields}.
Consequently a commuting invertible matrix acts by a nonzero
scalar on the line spanned by \(v\).
For \(x\) represented by \(h\in H\) in its real coordinates, set
\[
    D(x,a)=\log\|v h_\tau a_\tau\|-\log\|v h_\tau\|,
              \qquad l=\log|\lambda|.
\]
The scalar just described cancels in the difference, making
\(D\) well defined on \(Z\). Direct cancellation gives
\(D(x,ab)=D(x,a)+D(xa,b)\), and \(D(x,C(x))=l\).
By the inverse norm equality in
Proposition~\ref{prop:cyclic-orbit-count},
\(|D(x,a)|\le\log m(a)\).

Put \(I=\int_Z\Phi(x)e(C(x))\dd x\) and define
\[
 T=\int_{Z\times H}
       \Phi(x)e(a)e(a^{-1}C(x))D(x,a)\dd x\dd a .
\]
All the integrals used below are absolutely convergent.
For example \(m\) is submultiplicative and at least one.
Thus for \(d=(ab)^{-1}C(x)\) one has
\[
             |e(d)|\le C_0 m(a)^B m(b)^B m(C(x))^{-B}.
\]
Either \(|D(x,a)|\), \(|D(xa,b)|\), or their sum is at most
\(2m(a)m(b)\). This bounds the product with
\(|\Phi(x)e(a)e(b)|\) and \(|e(d)|\) by a constant times the
product of the three integrable functions already specified
on \(Z,H,H\). The analogous two-variable bound follows by
omitting \(b\). A constant factor \(l\) also preserves
integrability.

In the formula for \(I\), replace \(e(C(x))\) by its double
convolution integrand using idempotence. For \(b=a^{-1}C(x)\)
we have \(D(x,a)+D(xa,b)=l\). The transformation of the
pair in Lemma~\ref{lem:orbital-cycle} interchanges the two
increment integrals. Hence \(lI=2T\).

In the first occurrence of \(e\) defining \(T\), expand by
idempotence and change variables from the intermediate factor
and the original argument to the two successive factors \(a,b\).
By left invariance the measure on them is \(\dd a\dd b\).
The three kernel factors are now \(e(a)e(b)e(d)\) and the
increment is \(D(x,ab)\). Split this increment as
\(D(x,a)+D(xa,b)\). The first term integrates to \(T\) by
idempotence in the last two factors. The second has the same
integral as the first by the triple transformation of
Lemma~\ref{lem:orbital-cycle}, which leaves \(\Phi\) unchanged.
Therefore \(T=2T\), so \(T=0\). Since \(l\ne0\), the formula
\(lI=2T\) gives the desired vanishing.
\end{proof}

\section{Comparison of the elliptic geometric sums}

\begin{proposition}\label{prop:elliptic-product-trace}
In Theorem~\ref{thm:elliptic-unfold-convergence}, consider a class
\([\gamma]\) whose eigenvalues have absolute value one at every
archimedean place. In the finite cosets \(G_f/K_0\), let \(gK_0\)
run over representatives of the \(\Gamma_\gamma\)-orbits through
its action in the fixed finite frame. Write \(\gamma_{f,g}\) for
its matrix in that frame moved by \(g\), and \(A_{\gamma,g}\) for
the stabilizer of \(gK_0\) in \(\Gamma_\gamma\). The contribution
to the integral at this finite orbit is
\[
   \frac{\phi(\gamma_{f,g})}{|A_{\gamma,g}|}
         \int_H k(h^{-1}\gamma_\infty h)\dd h ,
\]
where \(\gamma_\infty\) is in the fixed real frames. This formula
holds whenever the finite weight is nonzero and the corresponding
absolute orbital integral converges. Zero finite weights contribute
zero. For a product test \(k=\prod_\tau k_\tau\) with absolutely
convergent individual orbital integrals, the real factor is the
product of those integrals. At a given \(\tau\) the matrix in each
such integral can be taken diagonal, with each eigenvalue placed
in a coordinate of the sign of its eigenline.
\end{proposition}
\begin{proof}
Use the decomposition over the countable discrete set of finite
cosets in Definition~\ref{def:moving-elliptic-lattice}.
Above a single orbit the quotient by \(\Gamma_\gamma\) identifies
with \(A_{\gamma,g}\backslash H\), since any finite coset in the
orbit can be carried to the chosen representative and the remaining
changes are exactly its stabilizer. The group \(A_{\gamma,g}\)
is finite by Proposition~\ref{prop:elliptic-class-description},
applied to the trace form and this finite frame. Its action on
\(H\) by left translation is free, since the rational embedding
at infinity is injective. The full real Haar integral is its
order times the quotient integral, by the finite partition into
translates of a measurable fundamental domain as in
Lemma~\ref{lem:quotient-measure}. The finite test is constant on
the chosen slice, giving the formula. Product integration follows
from Lemma~\ref{lem:integral-tools}. The simultaneous diagonal
description in Theorem~\ref{thm:elliptic-trace-fields} makes the
eigenlines orthogonal; rescale each eigenvector by the reciprocal
square root of its absolute diagonal entry and order the resulting
vectors according to the sign matrix. The change is unitary for
this hermitian matrix, and left translation in the real orbital
integral gives the asserted diagonal representative.
\end{proof}

\begin{theorem}\label{thm:geometric-stable-reduction}
At all \(f\) real places let the fixed rational unitary group of
Definition~\ref{def:moving-elliptic-lattice} have the same
signature \((p,q')\), \(p+q'=r\), \(f\) even. Identify all its
finite groups with those of the definite comparison space using
the fixed finite frames. Normalize the compact real groups to
have mass one and take \(K_0=K(g)\) of
Definition~\ref{def:definite-comparison-forms}.
Let \(\phi\) be bi-\(K_0\)-invariant of compact support and assume
each finite matrix in its support has irreducible characteristic
polynomial at some completion. Use counting measure on \(G_f/K_0\)
and suppose \(e=\prod_\tau e_\tau\) satisfies the hypotheses of
Theorem~\ref{thm:orbital-idempotent-vanishing} on the indefinite group.

At place \(\tau\) write \(I_\tau(z_1,\ldots,z_r)\) for the integral
of \(e_\tau(u^{-1}\operatorname{diag}(z_1,\ldots,z_r)u)\)
over \(U(p,q')\) in coordinates of the fixed signs. Assume these
integrals converge absolutely on distinct unit circle arguments
and satisfy
\[
 \frac{(-1)^{pq'}}{p!\,q'!}
      \sum_{\pi\in{\mathfrak S}_r}
                I_\tau(z_{\pi(1)},\ldots,z_{\pi(r)})
             =\overline{\chi_{\lambda_\tau}
                         (\operatorname{diag}(z_1,\ldots,z_r))},
\]
where \(\chi_{\lambda_\tau}\) is the compact character of
Theorem~\ref{thm:compact-alternant} for the row parameters.
It suffices to assume this identity at the eigenvalues of
the elliptic polynomials occurring with nonzero finite weight.
Then the global geometric integral on the initial group equals
the trace of \(\mathcal T_\phi\) on the definite space
\(\mathcal A^+\) in Proposition~\ref{prop:definite-elliptic-trace}.
\end{theorem}
\begin{proof}
By Theorem~\ref{thm:elliptic-unfold-convergence} the integral is
an absolutely convergent sum over rational conjugacy classes.
All with a nonunit-circle eigenvalue have zero contribution by
Theorem~\ref{thm:orbital-idempotent-vanishing}. For each remaining
polynomial \(P\), order its roots \(z_{\tau,1},\ldots,z_{\tau,r}\)
at each real place. Proposition~\ref{prop:elliptic-product-trace}
writes the contribution as a sum over the framed classes with
those roots. At each such term \(\phi(\gamma_{f,g})/|A_{\gamma,g}|\)
is determined by the finite matrix and the embedded centralizer.
Choosing \(q'\) of the eigenlines negative at each place changes
\(fq'\) signs from the all-positive pattern. This number is even.
Therefore Corollary~\ref{cor:fixed-group-elliptic-classes} and
Theorem~\ref{thm:elliptic-class-bijection} identify the framed
classes for every such pattern with those of the definite group
for \(P\), preserving precisely this finite weight. There are
no additional rational choices associated to the identification:
the quotient by rational conjugation is the one in
Proposition~\ref{prop:elliptic-class-description}.

Permuting eigenvalues within equal-sign coordinates leaves
\(I_\tau\) unchanged, since the permutation matrix is unitary
for the sign matrix and can be absorbed by left translation.
Summing over all assignments of the negative lines thus gives
\((p!q'!)^{-1}\sum_\pi I_\tau(z_{\tau,\pi(1)},\ldots,z_{\tau,\pi(r)})\)
at this place. The choices at different places are independent.
Their finite sums distribute over the product, yielding
\((-1)^{fpq'}\prod_\tau\overline{\chi_{\lambda_\tau}
(\operatorname{diag}(z_{\tau,1},\ldots,z_{\tau,r}))}\)
by the assumed identity. The sign is one since \(f\) is even.

On the compact group, the integral of the test
\(\dim W_{\lambda_\tau}\,\overline{m_{\lambda_\tau}}\) on the
same conjugates is exactly this conjugate character at \(\tau\)
by Theorem~\ref{thm:compact-alternant}. Hence the definite
geometric sum with this compact test gives the same finite
weights and archimedean factors. To identify it with the trace,
choose the finite double coset representatives \(h_c\) of
Proposition~\ref{prop:definite-elliptic-trace}. Their rational
stabilizers are trivial by Theorem~\ref{thm:definite-comparison-rank}.
Integration over each full compact real slice of mass one
therefore produces exactly its sum of finite coefficients times
the character of \(R_\infty\), the formula of
Proposition~\ref{prop:definite-elliptic-trace}. Absolute convergence
on the initial group justifies the reindexing, and on the definite
group the sum has finite support by that proposition.
\end{proof}

\section{The compact-dual unitary measure}

\begin{definition}\label{def:unitary-density}
Let \(S=\operatorname{diag}(\varepsilon_1,\ldots,\varepsilon_n)\),
\(\varepsilon_j=\pm1\), \(n\ge1\). On the real space of
skew-hermitian \(n\)-by-\(n\) matrices take the coordinate
Lebesgue density \(\dd Z\), with the imaginary part of each
diagonal entry and the two real coordinates of each strictly
lower triangular entry. On \(U(S)\) use the density
\(\dd\nu_S\) obtained at \(g\) by applying this density to the
skew-hermitian tangent coordinate \(g^*S\dd g\). Set
\[
                  u_n=\prod_{j=1}^n\frac{2\pi^j}{(j-1)!},
                      \qquad u_0=1,\qquad
                  \dd h_S=u_n^{-1}\dd\nu_S .
\]
\end{definition}

\begin{lemma}\label{lem:unitary-cayley-density}
The measures \(\nu_S,h_S\) are invariant under left and right
translations. For \(Z^*=-Z\) put \(A=SZ/2,\ D=1-A\). On
\(\det D\ne0\) the map \(g=(1+A)D^{-1}\) is a smooth chart
onto \(\{g\in U(S):\det(g+1)\ne0\}\) and
\[
                 \dd h_S(g)=u_n^{-1}|\det D|^{-2n}\dd Z .
\]
Its complement in \(U(S)\) has measure zero.
\end{lemma}
\begin{proof}
Differentiating \(g^*Sg=S\) makes \(g^*S\dd g\) skew-hermitian.
Conversely the chart at zero below realizes every such tangent
coordinate at identity, and translation does the same everywhere.
The form is unchanged by left translation and is pulled back
by right translation by \(b\) to its congruence by \(b\).
For any invertible complex \(C\), the real map \(Z\mapsto C^*ZC\)
on the skew-hermitian space has determinant of absolute value
\(|\det C|^{2n}\). For unitary \(C\) it preserves the real
Euclidean norm \((\operatorname{tr} Z^*Z)^{1/2}\); hence it
preserves every constant multiple of its Euclidean density,
including the chosen one. For positive diagonal entries \(s_i\)
the diagonal coordinate scales by \(s_i^2\) and each of the two
coordinates for \(i>j\) scales by \(s_i s_j\). Their product
is \(\prod_i s_i^{2n}\). For general \(C\), diagonalize
\(C^*C>0\) with a unitary matrix by the spectral argument of
Lemma~\ref{lem:common-hecke-interpolation}; dividing \(C\) by
its positive square root on the right gives a unitary matrix.
This proves the formula by composition. Since
\(|\det b|=1\) for \(b\in U(S)\), right invariance follows.

The identity \(A^*S=-SA\) gives
\((1+A)^*S(1+A)=D^*S D\). Thus the chart takes values in the
group. Conversely \(A=(g-1)(g+1)^{-1}\) satisfies
\(A^*S=-SA\) by multiplying the equation by \(g+1\) and its
adjoint and using \(g^*Sg=S\). Then \(Z=2SA\) is its inverse
coordinate and \(D=2(g+1)^{-1}\). These rational maps are
smooth on their stated domains. From \(g=2D^{-1}-1\) we get
\(\dd g=D^{-1}S\dd Z D^{-1}\). Also
\(g^*S D^{-1}S=(D^*)^{-1}\), since
\(D^*=S(1+A)S\) and \(A\) commutes with \(D^{-1}\).
The tangent coordinate is consequently
\((D^{-1})^*\dd Z D^{-1}\), giving the claimed density by
the congruence determinant. Finally for each \(g\) only finitely
many circle scalars \(t\) satisfy \(\det(tg+1)=0\), by the
polynomial root bound and its nonzero constant coefficient.
Integrating the exceptional indicator over the circle and \(U(S)\)
in either order by Lemma~\ref{lem:integral-tools}, left invariance
under \(g\mapsto tg\) shows that the exceptional set has measure zero.
\end{proof}

\begin{theorem}\label{thm:compact-dual-unitary-volume}
For \(S=1_n\), \(\nu_S(U(n))=u_n\); thus \(h_S\) is the
compact probability used in Proposition~\ref{prop:compact-formal-degree}.
For \(0\le k\le n\) the volume of the space of \(n-k\) ordered
orthonormal vectors with the quotient density of
\(\nu_{1_n}\) by \(\nu_{1_k}\) is \(u_n/u_k\).
In particular \(u_n=U_n\) in Definition~\ref{def:row-vandermonde}.
\end{theorem}
\begin{proof}
Consider the first-column map to the unit sphere in \(\C^n\).
It is onto by completion to an orthonormal basis: adjoin independent
standard basis vectors and successively subtract projections
onto the preceding normalized vectors, then normalize.
Locally it has smooth sections by the same procedure, using a
fixed list of independent columns at the point and the positivity
of each normalizing length on a neighborhood. The stabilizer of
the first column at the first standard vector is \(U(n-1)\)
in the remaining coordinates. At identity the differential sends
a skew-hermitian matrix to its first column. The chosen coordinate
density is the product of the ordinary tangent Euclidean density
of that column on the sphere with the density on its
\((n-1)\)-dimensional square block. Thus the quotient density
is the sphere's Euclidean surface density there. This identity
holds at every point by left invariance and invariance of the
sphere density under \(U(n)\). In local product coordinates given
by a section and right movement in the fiber, any change of
section adds vertical components to lifts of base tangents and
translates the fiber. These changes have determinant one for
the product of the quotient and fiber densities, by block
triangularity and the invariance in Lemma~\ref{lem:unitary-cayley-density}.
Ordinary integration in charts, partitioning a finite cover of
the compact sphere into disjoint measurable pieces, now gives
the volume as the sphere area times \(\nu_{1_{n-1}}(U(n-1))\).

The integral of \(\exp(-\|z\|^2)\) on \(\C^n\) is \(\pi^n\)
by planar polar coordinates and multiplication of integrals.
In radius \(\rho\) and surface coordinates the Jacobian is
\(\rho^{2n-1}\), because each tangent direction to the unit
sphere is scaled by \(\rho\) and the radial unit normal is
orthogonal to them. The radial integral is \((n-1)!/2\) after
substitution \(t=\rho^2\) and integration by parts as in
Theorem~\ref{thm:window}. This gives sphere area
\(2\pi^n/(n-1)!\). Induction starting with the trivial group
in dimension zero proves the stated value \(u_n\).
Stopping the column construction with \(k\) free columns gives
the ordered frame quotient with fiber \(U(k)\), and successive
integration gives \(u_n/u_k\). The displayed definition of
\(u_n\) agrees factor by factor with \(U_n\). By uniqueness of
compact invariant probability in Proposition~\ref{prop:compact-formal-degree}
the compact normalizations agree.
\end{proof}

\section{Principal minors and the polar coefficient}

\begin{definition}\label{def:polar-row-coefficient}
For an integer weight \(\alpha=(\alpha_1,\ldots,\alpha_n)\) order
the indices \(i_1,\ldots,i_n\) by decreasing weight and put
\(A_j=\{i_1,\ldots,i_j\}\). For an invertible matrix \(M\) set
\[
        m_\alpha(M)=(\det M)^{\alpha_{i_n}}
           \prod_{j=1}^{n-1}\det M[A_j,A_j]^{
                              \alpha_{i_j}-\alpha_{i_{j+1}}}.
\]
The zero exponents make this independent of tie breaking.
Let \(t_1>\cdots>t_n\) belong to \((n-1)/2+\Z\) and let
\(S=\operatorname{diag}(\varepsilon_i)\) have signs \(\pm1\).
Write \(\rho_i=(n+1)/2-i,\ \mu_i=t_i+\rho_i\), and define
\[
 \Lambda_i=\mu_i-\#\{j>i:\varepsilon_j=\varepsilon_i\}
                    +\#\{j<i:\varepsilon_j=\varepsilon_i\}.
\]
Let \(K=U(S)\cap U(n)\), the product of the unitary groups on the
two sign subspaces, with probability measure \(\dd u\) from
Theorem~\ref{thm:compact-dual-unitary-volume}. Set
\(c_\Lambda=\int_K |m_\Lambda(u)|^2\dd u\).
For \(g\in U(S)\) write \(g=kP\) with \(k\) ordinary unitary
and \(P\) positive definite hermitian. Define its polar coefficient
by
\[
       \psi(g)=\frac1{c_\Lambda}\int_K
       \frac{m_\Lambda(ku)\,\overline{m_\Lambda(u)}}
                    {m_{-\mu}(u^*P u)}\dd u .
\]
\end{definition}

\begin{lemma}\label{lem:sign-run-determinant}
Let \(P>0\) be hermitian with \(P^{-1}=SPS\), and write
\(d(A)=\det P[A,A]\), \(d(\varnothing)=1\).
Let \(J_+=\{i:\varepsilon_i=1\}\), \(A_j=\{1,\ldots,j\}\)
and \(J=\{1\le j<n:\varepsilon_j\ne\varepsilon_{j+1}\}\).
Then
\[
 \begin{gathered}
      d(A)=d(A^c)\ge1,\qquad
      d(J_+)\le\prod_{j\in J}d(A_j),\\
      d(J_+)\ge\frac{\|P\|_{\rm op}+\|P\|_{\rm op}^{-1}}2 .
 \end{gathered}
\]
For any \(c\ge0\) such that \(\mu_j-\mu_{j+1}\ge c\) on \(J\),
one has \(m_{-\mu}(P)\ge d(J_+)^c\).
\end{lemma}
\begin{proof}
All principal blocks have positive determinant, and \(\det P=1\)
since \(\det P=\det P^{-1}>0\). For a partition into two principal
blocks \(A,B\) with cross block \(C\), elimination gives
\(\det P=\det A\det(B-C^*A^{-1}C)\), and the lower diagonal
block of \(P^{-1}\) is \((B-C^*A^{-1}C)^{-1}\), by multiplying
the two triangular elimination factors. As \(P^{-1}=SPS\),
this proves \(d(I)=d(I^c)\) for every subset \(I\).
Also \(\det(B-C^*A^{-1}C)\le\det B\), since the conjugated
positive matrix \(B^{-1/2}(B-C^*A^{-1}C)B^{-1/2}\) has
eigenvalues in \((0,1]\). The argument applies to every
principal submatrix. Thus \(d(U\cup V)\le d(U)d(V)\) for disjoint
sets and \(1\le d(I)d(I^c)=d(I)^2\). For nested \(A\subset B\),
taking complements gives \(d(B\setminus A)\le d(A)d(B)\).
Decompose \(J_+\) into its maximal consecutive intervals.
For each such interval with endpoints \(l,h\) its determinant
is at most \(d(A_{l-1})d(A_h)\). All interior boundaries
between opposite signs occur exactly once in the product over
these intervals. The empty and full determinants are one.
This proves the upper bound on \(d(J_+)\).

The block diagonal part of \(P\) on the sign subspaces is
\((P+P^{-1})/2\), of determinant \(d(J_+)d(J_+^c)=d(J_+)^2\).
The positive eigenvalues of \(P\) occur with their reciprocals
since \(P^{-1}=SPS\). If \(T=\|P\|_{\rm op}>1\), \(T\) and
\(T^{-1}\) each supply a factor \((T+T^{-1})/2\) to that
determinant, and all other factors are at least one. If \(T=1\)
all factors are one. This proves the required lower bound.
Finally \(\mu\) is decreasing, so the distinct proper sets
occurring for \(-\mu\) are tails of this order. Complementarity
and \(\det P=1\) give
\(m_{-\mu}(P)=\prod_{j<n} d(A_j)^{\mu_j-\mu_{j+1}}\).
The exponents are positive, every base is at least one, and
the exponents at the indicated boundaries are at least \(c\).
The upper bound on \(d(J_+)\) yields the last assertion.
\end{proof}

\begin{theorem}\label{thm:polar-row-decay}
The function \(\psi\) is continuous. It takes value one at
identity and satisfies \(\psi(g^{-1})=\overline{\psi(g)}\). Under the
gap condition of Lemma~\ref{lem:sign-run-determinant}, if
\(T=\|g\|_{\rm op}\), it satisfies
\[
          |\psi(g)|\le \left(\frac{T+T^{-1}}2\right)^{-c}.
\]
If \(c>B+1+4n^2\) for \(B\ge0\), then
\(\int_{U(S)}|\psi(g)|\,T^{B+1}\dd h_S(g)<\infty\) and
\(0<\int_{U(S)}|\psi|^2\dd h_S<\infty\).
\end{theorem}
\begin{proof}
By the spectral argument of Lemma~\ref{lem:common-hecke-interpolation}
the positive square root \(P=(g^*g)^{1/2}\) exists and is unique.
Indeed any positive square root commutes with its square and
on each eigenspace has only the positive square root eigenvalue.
It depends continuously on \(g\): for any convergent sequence
of positive definite squares near a fixed one, the roots are
bounded and stay uniformly positive by the quadratic-form bounds
on the squares. Every convergent subsequence of roots then has
the unique root of the limit. Finite-dimensional compactness of
bounded closed sets proves continuity. Since
\(S(g^*g)S=(g^*g)^{-1}\), uniqueness gives \(SPS=P^{-1}\).
Thus \(P\in U(S)\) and \(k=gP^{-1}\in K\). The polar decomposition
is unique and \(T=\|P\|_{\rm op}\). The integer weights in
Definition~\ref{def:polar-row-coefficient} give continuous minors
on invertible matrices, with a strictly positive denominator on
each positive definite argument. The compact integration is
continuous by uniform continuity on compact neighborhoods, and
\(c_\Lambda>0\) since \(m_\Lambda(1)=1\) and every open identity
neighborhood has positive smooth measure. At identity the
numerator integrates to \(c_\Lambda\).

Conjugating \(P\) by \(u\in K\) preserves its sign-block determinants.
Apply Lemma~\ref{lem:sign-run-determinant} in the denominator.
Cauchy--Schwarz and left invariance give
\(\int_K|m_\Lambda(ku)m_\Lambda(u)|\dd u\le c_\Lambda\).
The asserted inequality follows. For inversion write
\(g^{-1}=k^{-1}(kP^{-1}k^{-1})\) and substitute \(u=kv\).
The numerator is the complex conjugate of the original one at \(v\).
The denominator agrees, since \(P^{-1}=SPS\), \(v\) commutes
with \(S\), and diagonal sign conjugation preserves every
principal minor. This proves the symmetry.

Finally use \(T\ge1\) and the last bound in the weaker form
\(2^c T^{-c}\). On \(2^j\le T<2^{j+1}\),
Lemma~\ref{lem:real-unitary-ball} bounds the integral of
\(|\psi|T^{B+1}\) by a constant times
\(2^{j(B+1+4n^2-c)}\). Summation gives finiteness. Since
\(|\psi|\le1\), this also bounds its squared integral. Positivity
of the latter follows from continuity and its value at identity.
\end{proof}

\begin{corollary}\label{cor:row-polar-moment}
In Definition~\ref{def:row-vandermonde}, reindex the row entries
from \(1\) to \(r\) in decreasing order. Suppose the first \(k-1\)
entries have a common sign; the remaining signs are arbitrary.
For all sufficiently large arithmetic scales take \(B=r^8\).
At each place put \(d_\tau=(\int |\psi_\tau|^2\dd h_{S_\tau})^{-1}\)
and \(e_\tau=d_\tau\overline{\psi_\tau}\). The product
\(e=\prod_\tau e_\tau\) is continuous and satisfies
\[
          |e(g)|\ll m(g)^{-B},\qquad
                \int_H |e(g)|\,m(g)^{B+1}\dd g<\infty .
\]
It satisfies \(e(g^{-1})=\overline{e(g)}\) and
\(\int_H e(a)e(a^{-1})\dd a=e(1)\).
\end{corollary}
\begin{proof}
In the top \(k-1\) positions there is no change of sign. The
gap between the last top entry and the next is
\(1-k/2+W_\tau+2r^8\); all subsequent gaps are \(2r^8\) by
Definition~\ref{def:row-vandermonde}. Here \(W_\tau\ge0\).
The \(\mu\)-gaps exceed the \(t\)-gaps by one. Since \(k\le r+1\)
and \(r\to\infty\) by Proposition~\ref{prop:thick}, all
sign-change gaps are at least \(c=r^8+4r^2+2\) eventually.
Apply Theorem~\ref{thm:polar-row-decay} and multiply the
pointwise bounds and the finite collection of integrals. At
identity the convolution formula follows from inversion symmetry
and \(\int |e_\tau|^2=d_\tau=e_\tau(1)\).
\end{proof}

\section{Independence of the block splitting}

\begin{definition}\label{def:block-splitting-integral}
Retain \(S,K,\mu,\Lambda,c_\Lambda\) of
Definition~\ref{def:polar-row-coefficient}. Let \(B\) be the
complex group of invertible matrices block diagonal for the signs,
\(\mathcal T\) its group of invertible diagonal matrices and
\(T=K\cap\mathcal T\). Upper and lower triangular matrices in
\(B\) are defined by the given order of indices in each block.
Write \(\chi_\alpha(b)=\prod_i b_{ii}^{\alpha_i}\) for a
triangular \(b\). A complex matrix \(Y\) is strictly accretive
if \((Y+Y^*)/2\) is positive definite. For \(g\in{\rm GL}_n(\C)\)
and \(a\in B\) with \(a^{-1}g\) strictly accretive put
\[
  \Psi_a(g)=c_\Lambda^{-1}\int_K
   \frac{m_\Lambda(av)m_\Lambda(v^{-1})}
        {m_{-\mu}(v^{-1}a^{-1}gv)}\dd v .
\]
A strictly accretive matrix and each of its principal compressions
are invertible, since a nonzero kernel vector would have quadratic
form of real part zero. This also makes the denominator in the
display nonzero for unitary \(v\).
\end{definition}

\begin{lemma}\label{lem:flag-cycle-jacobian}
On the complex manifold \(B/\mathcal T\), use the holomorphic
volume form \(\Theta\) obtained by wedging, in a fixed order,
all strictly off-diagonal entries within blocks of \(v^{-1}\dd v\).
Its restriction to \(K/T\) is a nonzero constant multiple of the
oriented invariant probability density. For \(b\in B\), where
\(v^{-1}bv=l u^{-1}\) admits a lower triangular \(l\) and a
unipotent upper triangular \(u\), put \(w=v l=bv u\).
This gives a holomorphic map \(F_b:[v]\mapsto[w]\) on that
open locus. Set
\(\delta_i=\#\{j>i:\varepsilon_j=\varepsilon_i\}
              -\#\{j<i:\varepsilon_j=\varepsilon_i\}\). Then
\[
                       F_b^*\Theta
                =\chi_{\delta}(u)\chi_{\delta}(l)^{-1}\Theta .
\]
\end{lemma}
\begin{proof}
The quotient parametrizes ordered lines forming a basis inside
each block. On the open set where one specified entry of each
column is nonzero, scale the column to put that entry equal to
one. These are complex charts with holomorphic transition maps.
Under \(v\mapsto v d\) with \(d\in\mathcal T\) variable,
\(v^{-1}\dd v\) changes by conjugation by \(d\) and the addition
of a diagonal form. The factors \(d_{jj}/d_{ii}\) on the
off-diagonal entries multiply to one. These forms also vanish on
vertical diagonal directions, so their wedge defines \(\Theta\)
and is nonzero in every chart. On the compact cycle the columns
are mutually orthogonal; their phases are the only ambiguity
in their unit representatives. Choosing phases to make specified
nonzero entries positive gives local smooth sections of \(K\to K/T\).
Tangent coordinates at identity modulo the diagonal are the
off-diagonal coordinates of a skew-hermitian matrix, with each
opposite pair negatives of complex conjugates. Their two-form
wedge is a nonzero constant times the real two-coordinate
density. Left translation proves that \(\Theta\) restricts to
a constant nonzero multiple of an oriented density everywhere.
This orientation may be chosen at identity and translated since
the diagonal isotropy acts by rotations on the pairs. Integration
against that density after normalization is the pushforward of
Haar probability from \(K\): local sections identify the invariant
density on \(K\) with base density times invariant fiber density,
as in Theorem~\ref{thm:compact-dual-unitary-volume}.

Elimination with the successive upper-left pivots in each block
gives the unique \(l,u\) on the locus where these minors are
nonzero; all their entries are rational with denominators products
of such minors. Diagonal change of the representative conjugates
\(l,u\) by the same diagonal matrix. Thus \(F_b\) is well defined
and holomorphic. For \(X=v^{-1}\dd v,\ Y=w^{-1}\dd w\) we have
\[
            Y=u^{-1}Xu+u^{-1}\dd u
              =l^{-1}Xl+l^{-1}\dd l .
\]
Modulo upper triangular matrices the lower entries of \(Y\) are
the image of the lower entries of \(X\) under conjugation by \(u\).
The determinant of this map is
\(\prod_{i>j,\ \varepsilon_i=\varepsilon_j}u_{jj}/u_{ii}
=\chi_\delta(u)\). Indeed image terms from the elementary matrix
at \((i,j)\) occur only in rows at most \(i\) and columns at
least \(j\), so all off-diagonal corrections remaining below the
diagonal have strictly smaller distance to it. The dual argument
modulo lower triangular matrices, using \(l\), gives determinant
\(\prod_{i<j,\ \varepsilon_i=\varepsilon_j}l_{jj}/l_{ii}
=\chi_\delta(l)^{-1}\) on the upper entries. The two lists of
entries depend separately on the corresponding lists of \(X\).
Multiplication of their determinants gives the formula.
\end{proof}

\begin{theorem}\label{thm:accretive-splitting}
For fixed \(g\in{\rm GL}_n(\C)\), \(\Psi_a(g)\) is independent
of \(a\in B\) whenever \(a^{-1}g\) is strictly accretive.
For \(g\in U(S)\), if \(a\) is its block diagonal part, then
\(\Psi_a(g)=\psi(g)\) of Definition~\ref{def:polar-row-coefficient}.
\end{theorem}
\begin{proof}
Inside each sign block \(\Lambda\) is decreasing in the index
order. For \(i<j\) of the same sign, \(t_i-t_j\ge j-i\) and
\(\rho_i-\rho_j=j-i\). The difference \(\delta_i-\delta_j\) is
twice the difference of their positions within that block, at
most \(2(j-i)\); now use \(\Lambda=\mu-\delta\).
Thus the leading sets defining \(m_\Lambda\) on block diagonal
matrices can be ordered to intersect each block in a prefix.
The determinant on such a prefix multiplies by the product of
the corresponding diagonal entries under right upper triangular
multiplication or left lower triangular multiplication. Telescoping
the exponents in Definition~\ref{def:polar-row-coefficient}
gives multiplier \(\chi_\Lambda\). For \(m_{-\mu}\) on arbitrary
matrices the defining sets are tails of the full order. Left
upper or right lower triangular multiplication therefore gives
multiplier \(\chi_{-\mu}\). These identities include the full
determinant power, and allow vanishing numerator minors.

For \(v\in B\) write \(f_a(v)\) for the quotient of minors in
Definition~\ref{def:block-splitting-integral}, continued from
\(K\) by the same formula with \(v^{-1}\). It is invariant under
right diagonals and holomorphic where its denominator is defined
and nonzero. For \(w=bv u=v l\) as in Lemma~\ref{lem:flag-cycle-jacobian},
use \(w=bv u\) on the first numerator and \(w^{-1}=l^{-1}v^{-1}\)
on the second. In the denominator use \(w^{-1}=u^{-1}v^{-1}b^{-1}\)
on the left and \(w=v l\) on the right. The multiplication laws
just proved give
\[
 f_a(w)=
 \chi_{\Lambda-\mu}(u)\chi_{\Lambda-\mu}(l)^{-1}f_{ab}(v).
\]
Since \(\Lambda-\mu=-\delta\), the Jacobian in
Lemma~\ref{lem:flag-cycle-jacobian} gives the identity of top
forms \(F_b^*(f_a\Theta)=f_{ab}\Theta\). The denominator at \(w\)
is nonzero whenever the denominator at \(v\) for \(ab\) is
nonzero, since all triangular diagonal factors are invertible.

Strict accretivity of \(Y\) is equivalent to that of \(Y^{-1}\),
since the hermitian part of \(Y^{-1}\) is the congruence of the
hermitian part of \(Y\) by \(Y^{-1}\). Thus admissible \(a\)
are exactly the block matrices with \(g^{-1}a\) strictly
accretive. This is a convex open set, all of whose elements
are invertible. Take the line segment between two such matrices.
Subdivide it into finitely many intervals on each of which,
writing the path as \(a b_s\) from its initial matrix \(a\),
all \(b_s\) are sufficiently close to identity. This is possible
since the inverse is bounded along the compact segment.
Uniformly for \(v\in K\), all leading pivots of \(v^{-1}b_s v\)
are then nonzero by compactness and continuity. The maps
\(F_{b_s}\) give a smooth deformation of \(K/T\) starting at
the identity inclusion. Along this deformation \(f_a\Theta\)
is holomorphic in a neighborhood of the image: the denominator
for \(ab_s\) is nonzero on \(K\) by strict accretivity and
the invariance formula preserves every required denominator
minor up to an invertible factor.

The integral of this form on the deformed cycle is constant.
For detail, a holomorphic top form is closed by differentiation
in complex charts. Pull back to the interval times the compact
oriented manifold \(K/T\) and write the result
\(\omega_s+\dd s\wedge\eta_s\), with \(\omega_s\) of top real
degree on the cycle. Closedness gives
\(\partial_s\omega_s=\dd_{K/T}\eta_s\).
The integral of this latter exact form is zero. Cover the
compact manifold by finitely many coordinate balls whose smaller
balls still cover it. Smooth nonnegative functions positive on
the smaller balls and supported inside the larger ones exist by
translating and rescaling the function
\(\exp(-1/(1-|x|^2))\) on the unit ball, extended by zero.
Divide their sum into each function to obtain a partition of
unity. The integral of the exterior derivative of each localized
form is zero by integration of compactly supported coordinate
derivatives and the fundamental theorem of calculus. Summing
gives the assertion, including the zero-dimensional case where
the integral is constant along a deformation in a discrete space.
Thus the two endpoint integrals of \(f_{ab_s}\Theta\) on \(K/T\)
agree. Its normalized restriction is Haar integration for these
diagonal-invariant functions by Lemma~\ref{lem:flag-cycle-jacobian}.
Applying this to each subinterval proves independence.

For the polar decomposition \(g=kP\) in
Theorem~\ref{thm:polar-row-decay}, both \(k\) and
\(a=k(P+P^{-1})/2\) are admissible. Indeed
\((P+P^{-1})/2>0\) commutes with \(P\), so the remainder for
the latter matrix is positive definite. This matrix is the
block diagonal part of \(g\), since \(SPS=P^{-1}\) and \(k\)
commutes with \(S\). At \(a=k\) the formula is exactly the
polar integral, using \(m_\Lambda(v^{-1})=\overline{m_\Lambda(v)}\)
for \(v\) unitary. Independence gives the claimed equality.
\end{proof}

\section{The compact projection of the real coefficient}

\begin{lemma}\label{lem:circle-leading-minors}
Use \(K,\Lambda\) of Definition~\ref{def:polar-row-coefficient}.
For \(s\in\C^\times\) put \(D(s)=\operatorname{diag}(s^1,\ldots,s^n)\)
and \(L=\sum_{i=1}^n i\Lambda_i\). For \(v,j\in K\),
\[
             s^{-L}m_\Lambda(vD(s)j^{-1})
\]
extends to a polynomial in \(s\), with value
\(m_\Lambda(v)m_\Lambda(j^{-1})\) at zero.
\end{lemma}
\begin{proof}
Since \(\Lambda\) is decreasing within each sign block by
Theorem~\ref{thm:accretive-splitting}, its minor expression on a
block matrix is the product of the expressions of the restricted
weights on the two blocks. Indeed in each principal determinant
in Definition~\ref{def:polar-row-coefficient} the matrix is a
direct sum, and combining successive equal subsets in each block
telescopes the exponents to the restricted weight differences.
For one block with indices in increasing order, let \(A\) be
the first \(l\) indices. Multilinearity of the determinant in
the columns gives
\[
  \det(vD(s)j^{-1})[A,A]
     =\sum_{\substack{I\ {\rm in\ the\ block}\\|I|=l}}
          \det v[A,I]\,\det j^{-1}[I,A]\,s^{\sum_{i\in I}i}.
\]
Terms with repeated indices cancel by alternation, and grouping
permutations of a fixed set gives the second determinant. The
least possible exponent belongs uniquely to \(I=A\), whether
or not its coefficient vanishes. Dividing by that power gives a
polynomial with constant term the product of the two principal
minors. The full block determinant, which can occur to a negative
power, is the single monomial given by the full index set and
has nonzero coefficient. Multiplying its factor and the
nonnegative powers of the preceding polynomials accounts for
exactly \(s^L\) by the same telescoping. The product of their
constant coefficients gives the displayed value.
\end{proof}

\begin{theorem}\label{thm:real-compact-product}
For any \(x,y\in U(S)\), the coefficient of
Theorem~\ref{thm:accretive-splitting} satisfies
\[
    \int_K m_\Lambda(k^{-1})\psi(xky)\dd k
                         =c_\Lambda\,\psi(x)\psi(y).
\]
\end{theorem}
\begin{proof}
Fix \(x\), take \(a\) its block diagonal part and put \(Z=a^{-1}x\).
This remainder is strictly accretive by Theorem~\ref{thm:accretive-splitting}.
For \(y\) sufficiently near identity in \(U(S)\), \(ak\) is an
admissible splitting of \(xky\) for every \(k\in K\), since
\(k^{-1}Zk y\) is strictly accretive uniformly by compactness.
The identity matrix is also an admissible splitting of \(y\).
In the left side apply Definition~\ref{def:block-splitting-integral}
with \(ak\); write \(v\in K\) for the integral variable there,
and set \(j=kv\). Bi-invariance of compact probability gives
that the left side equals \(1/c_\Lambda\) times the integral
over \(j,v\in K\) of
\[
  \frac{m_\Lambda(vj^{-1})m_\Lambda(aj)m_\Lambda(v^{-1})}
       {m_{-\mu}(Z_jY_v)},\qquad Z_j=j^{-1}Zj,\quad Y_v=v^{-1}yv .
\]
By elimination starting with its last pivot, factor
\(Y_v=U_v C_v\) with \(U_v\) unipotent upper triangular and
\(C_v\) lower triangular, now for the full order of indices.
This factorization exists uniformly near identity because the
trailing principal pivots are nonzero there; the factors are
continuous rational functions there and both tend uniformly
to identity as \(y\) tends to identity. The multiplication law
for trailing minors in Theorem~\ref{thm:accretive-splitting}
gives \(m_{-\mu}(Y_v)=\chi_{-\mu}(C_v)\).

Replace \(v\) in the integrand by \(vD(s)\) with \(|s|=1\).
This leaves its Haar integral invariant. We compute its circle
mean first, holding \(j,v\) fixed. The factor
\(m_\Lambda((vD(s))^{-1})\) becomes
\(s^{-L}m_\Lambda(v^{-1})\) by diagonal multiplication. In the
denominator write
\[
 D(s)^{-1}Y_vD(s)=U_v(s)\,(D(s)^{-1}C_vD(s)).
\]
The second factor is lower triangular with the same diagonal as
\(C_v\). Here \(U_v(s)\) has unit diagonal and upper entries
\(s^{j'-i}(U_v)_{ij'}\) for \(i<j'\). It is a polynomial matrix
with \(U_v(0)=1\). For \(|s|\le1\), \(U_v(s)\) stays uniformly
near identity with \(y\) sufficiently near identity. All principal
minors of \(Z_j U_v(s)\) are nonzero on this closed disk because
these matrices are strictly accretive, uniformly over \(j,v\).
The circle integrand can therefore be written as the boundary
values of a holomorphic function on a neighborhood of the disk:
use Lemma~\ref{lem:circle-leading-minors} for
\(s^{-L}m_\Lambda(vD(s)j^{-1})\) and use the inverse of
\(m_{-\mu}(Z_j U_v(s))\chi_{-\mu}(C_v)\) for its varying
denominator. The compact circle average is its value at zero,
by the Cauchy formula in Lemma~\ref{lem:analytic-tools}.
The resulting integrand is the product
\[
        \frac{m_\Lambda(aj)m_\Lambda(j^{-1})}{m_{-\mu}(Z_j)}
          \ \frac{m_\Lambda(v)m_\Lambda(v^{-1})}{m_{-\mu}(Y_v)} .
\]
All compact interchanges are legitimate by continuity and the
nonzero denominators. Integrating and dividing by \(c_\Lambda\)
gives \(c_\Lambda\Psi_a(x)\Psi_1(y)\), as asserted near identity.

For arbitrary \(y\), use a real analytic path \(y(t)\) from
\(1\) to \(y\). Explicitly write \(y=kP\) as in
Theorem~\ref{thm:polar-row-decay}, and diagonalize \(k\) in
its two unitary blocks to write \(k=\exp A\) with \(A\)
skew-hermitian and block diagonal, by choosing arguments of
its eigenvalues. By positive spectral calculus \(P=\exp H_0\)
with \(H_0\) hermitian, and \(SH_0S=-H_0\) since \(SPS=P^{-1}\).
Thus \(y(t)=\exp(tA)\exp(tH_0)\) belongs to \(U(S)\) for real
\(t\), and its matrix entries are entire in complex \(t\) by
their convergent exponential series. For real \(t_0\), take as
splittings of \(xk y(t)\) the block diagonal parts of
\(xk y(t_0)\), fixed as \(t\) varies. The remainders at \(t_0\)
are uniformly strictly accretive for \(k\in K\) by compactness,
so the same holds on a complex disk about \(t_0\). There the
integral with these fixed splittings is holomorphic in \(t\):
the integrand is rational in the matrices with nonvanishing
denominators, uniformly bounded on smaller compact disks and
compact integration variables. Integrating its Cauchy power
series term by term on a smaller disk as in
Lemma~\ref{lem:analytic-tools} proves holomorphy.
On the real interval it gives precisely the required function
by Theorem~\ref{thm:accretive-splitting}. The same argument
applies to \(\psi(y(t))\). The two real analytic functions agree
near zero; successive overlapping disks and the identity theorem
of Lemma~\ref{lem:analytic-tools} give equality at \(t=1\).
\end{proof}

\section{Orthogonality of the real kernel}

\begin{theorem}\label{thm:row-real-orthogonality}
Let \(G=U(S)\) and suppose the moment bound of
Theorem~\ref{thm:polar-row-decay} holds for some \(B\ge0\).
Write \(d=(\int_G|\psi|^2\dd h_S)^{-1}\) and
\(e=d\,\overline{\psi}\). Then for every \(x,y\in G\)
\[
 d\int_G \overline{\psi(z)}\,\psi(xzy)\dd h_S(z)=\psi(x)\psi(y).
\]
In particular \(e(1)=d>0,\ e(g^{-1})=\overline{e(g)}\) and
\[
     \int_G e(z)e((zx)^{-1}y)\dd h_S(z)
                        = \frac{e(x^{-1})e(y)}{e(1)} .
\]
There exists a continuous compactly supported \(\beta:G\to\C\)
such that, for every \(y\in G\),
\[
                 \int_G\beta(z)e(z^{-1}y)\dd h_S(z)=e(y).
\]
\end{theorem}
\begin{proof}
Write \(m=m_\Lambda\) in this proof and \(c=c_\Lambda\). If
\(w\in L^1(G)\) satisfies \(c w(z)=\int_K m(k^{-1})w(k^{-1}z)\dd k\),
then
\[
       \int_G w(z)\psi(xzy)\dd h_S(z)
               =\psi(x)\int_G w(z)\psi(zy)\dd h_S(z).
\]
Indeed replace \(cw\) by its compact average on the left, use
\(z=kz'\) and apply Theorem~\ref{thm:real-compact-product}
to the middle \(k\)-integral at \(x,z'y\). Divide by \(c>0\).
All these integrals converge absolutely since \(\psi\) is
bounded by one by Theorem~\ref{thm:polar-row-decay} and \(m\)
is bounded on \(K\). The function \(w(z)=\overline{\psi(z)}
=\psi(z^{-1})\) is integrable by the same theorem and satisfies
the required identity by Theorem~\ref{thm:real-compact-product}
at \(z^{-1},1\). Thus the displayed factorization holds for this \(w\).

For \(I(x,y)=\int_G\overline{\psi(z)}\psi(xzy)\dd h_S(z)\)
we have obtained \(I(x,y)=\psi(x)I(1,y)\).
In the integral \(I(1,y)\) put \(u=(zy)^{-1}\).
Translation and inversion preserve the measure by
Lemmas~\ref{lem:unitary-cayley-density} and~\ref{lem:orbital-cycle}.
Inversion symmetry of \(\psi\) then gives \(I(1,y)=I(y,1)
=\psi(y)I(1,1)=\psi(y)/d\). This proves the first identity.
Since
\(e((zx)^{-1}y)=d\,\psi(y^{-1}zx)\), the second integral is
\(d^2 I(y^{-1},x)\), giving the stated formula.

Choose a nonnegative continuous function \(\eta\) of
\(\operatorname{tr}(z^*z)\), equal to one near its value \(n\)
at identity and zero above \(n+2\). This is compactly supported:
the trace bound bounds all entries, and \(G\) is a closed subset
of the matrices by its Gram equation. It is left \(K\)-invariant.
The number \(A=\int_G\eta|\psi|^2\dd h_S\) is positive by
continuity and \(\psi(1)=1\). Thus
\(\beta=\eta\,\overline{\psi}/A\) is continuous, compactly
supported and satisfies the same compact averaging identity for
\(w\), by left \(K\)-invariance. Apply the first displayed
factorization of the proof to \(w=\beta,x=y^{-1}\), with the
last argument \(1\). Since \(\int_G\beta(z)\psi(z)\dd h_S(z)=1\)
and \(e(z^{-1}y)=d\psi(y^{-1}z)\), the required formula follows.
\end{proof}

\begin{proposition}\label{prop:actual-real-row-projection}
For the sign patterns and parameters of
Corollary~\ref{cor:row-polar-moment} use
\(e=\prod_\tau d_\tau\overline{\psi_\tau}\) on \(H\). At any
of the levels in Definition~\ref{def:moving-elliptic-lattice} let
\(\mathcal L=L^2(\Gamma\backslash X)\), using the quotient
measure there and \(\dd h_{S_\tau}\) at each real place. Then
\[
             (\mathcal P F)(x)=\int_H e(a)F(xa)\dd a
\]
defines an orthogonal projection on \(\mathcal L\).
Its range consists precisely of the square-integrable functions
fixed by convolution with this explicit real kernel. In the
geometric integral of Theorem~\ref{thm:elliptic-unfold-convergence}
for this \(e\), all terms with a nonunit-circle eigenvalue give
zero for all sufficiently large arithmetic scales.
\end{proposition}
\begin{proof}
Products of the identities of Theorem~\ref{thm:row-real-orthogonality}
give \(\int_H e(a)e(a^{-1}g)\dd a=e(g)\), the symmetry under
inversion, and a compactly supported reproducing function.
All products are finite and the needed integrability follows
from Corollary~\ref{cor:row-polar-moment}. The quotient measure
is preserved by right \(H\) as in Lemma~\ref{lem:orbital-cycle},
whose proof applies to the whole rational group.
For \(F\in\mathcal L\), Cauchy--Schwarz with weight \(|e(a)|\)
and integration over the quotient give
\[
  \int_{\Gamma\backslash X}
       \left(\int_H |e(a)F(xa)|\dd a\right)^2\dd x
                     \le \|e\|_1^2\|F\|_2^2 .
\]
This gives an almost everywhere defined bounded operator independent
of representatives. Applying the same estimate twice justifies
the iterated integral; explicitly \(|e|*|e|\) has \(L^1\)-norm
\(\|e\|_1^2\) by left invariance. Hence \(\mathcal P^2=\mathcal P\)
by the convolution identity. For the inner product with
\(F'\in\mathcal L\), the integral of the absolute product is
bounded by \(\|e\|_1\|F\|_2\|F'\|_2\). The changes
\(x\mapsto xa,\ a\mapsto a^{-1}\) and inversion symmetry give
\(\mathcal P^*=\mathcal P\). An idempotent fixes its range and
a self-adjoint idempotent has range orthogonal to its kernel,
proving the projection assertion. At \(B=r^8\) the size
requirement follows from Theorem~\ref{thm:elliptic-unfold-convergence}.
The actual kernel now satisfies all the hypotheses of
Theorem~\ref{thm:orbital-idempotent-vanishing}, which gives
the last claim.
\end{proof}

\section{Finite rank of the compact type projections}

\begin{definition}\label{def:cyclic-row-module}
Use the measure and the parameters of
Theorem~\ref{thm:row-real-orthogonality}, with the moment bound
there. Put \(v_0=\sqrt d\,\psi\) in \(L^2(G,h_S)\).
Write \(\pi(g)f(z)=f(zg)\), and let \(W\) be the closed linear
span of \(\pi(g)v_0\), \(g\in G\). The inner product is conjugate
linear in its first variable. A compact dominant weight is
\(\nu\in\Z^n\) decreasing weakly in index order within each
of the two sign sets. On \(K\) take the tensor product \(T_\nu\)
of the two compact representations of
Theorem~\ref{thm:compact-highest-vector} for the restricted
weights and their inherited order. A block of size zero supplies
the trivial line. Put \(s_\nu=\dim T_\nu\),
\(\chi_\nu=\operatorname{tr}T_\nu\), and write \(m_\nu\) for
the coefficient of the tensor product of their distinguished
unit vectors. It agrees on \(K\) with the minor formula of
Definition~\ref{def:polar-row-coefficient}, by telescoping the
two block factors as in Lemma~\ref{lem:circle-leading-minors}.
\end{definition}

\begin{lemma}\label{lem:cyclic-row-isometry}
For \(a,b\in W\),
\[
 d\int_G|\langle a,\pi(g)b\rangle|^2\dd h_S(g)
                 =\|a\|^2\|b\|^2,\qquad
 \langle v_0,\pi(g)v_0\rangle=\psi(g).
\]
The action on \(W\) is unitary, strongly continuous and irreducible.
\end{lemma}
\begin{proof}
Right invariance of \(h_S\) gives unitarity. The second identity
is Theorem~\ref{thm:row-real-orthogonality} with \(x=1\).
In particular \(\|v_0\|=1\), and
\(\|\pi(g)v_0-v_0\|^2=2-2\Rp\psi(g)\) gives continuity at
identity by Theorem~\ref{thm:polar-row-decay}. Translation, linearity
and approximation by finite sums of translates give strong
continuity on \(W\). For \(u_i,v_i\in G\), \(i=1,2\), change
variables to \(z=u_1^{-1}gv_1\) and use the two-sided identity
of Theorem~\ref{thm:row-real-orthogonality}. It gives
\[
 d\int_G\overline{\psi(u_1^{-1}gv_1)}\,
              \psi(u_2^{-1}gv_2)\dd h_S(g)
       =\psi(u_2^{-1}u_1)\psi(v_1^{-1}v_2).
\]
These are the inner products of the corresponding translates,
in reverse order for the \(u\)'s. Expanding finite sums gives
the norm identity there. For sequences of finite sums tending
to \(a,b\), subtract coefficients by changing one vector at a
time. The finite-sum norm identity makes the resulting coefficients
Cauchy in \(L^2\). They also converge uniformly pointwise to
\(\langle a,\pi(g)b\rangle\) by Cauchy--Schwarz and unitarity.
On compact sets the two limits therefore agree in \(L^2\);
the norm balls exhaust \(G\). This proves the identity.
If a nonzero proper closed invariant subspace existed, take
nonzero \(a\) in it and \(b\) in its invariant orthogonal
complement. All the integrands would vanish, contradicting
the positive right side.
\end{proof}

\begin{theorem}\label{thm:row-compact-finite-rank}
For a compact dominant weight let
\[
 P_\nu=s_\nu\int_K\overline{m_\nu(k)}\,\pi(k)\dd k .
\]
The integral, taken on each vector of \(W\), defines an
orthogonal projection of finite rank. Its rank at \(\nu=\Lambda\)
is one. A positive rank requires
\[
 \sum_{i=1}^n(\nu_i-\Lambda_i)=0,\qquad
 \sum_{i=1}^j(\nu_i-\Lambda_i)\ge0\quad(1\le j<n).
\]
\end{theorem}
\begin{proof}
Compact integration of a continuous vector-valued function is
the norm limit of finite sums: partition a finite cover by sets
on which its oscillation is small and weight samples by the
measures of the pieces. Uniform continuity and completeness
give existence and the usual norm bound. To compute the compact
convolution, average the rank-one orthogonal projector onto
the distinguished line in each factor over conjugation by that
factor. The average commutes with the representation; the
eigenspaces of the self-adjoint average are invariant and
irreducibility in Theorem~\ref{thm:compact-highest-vector} forces
a scalar. Its trace is one, so the scalar is the reciprocal
dimension, as in Proposition~\ref{prop:compact-formal-degree}.
The averages multiply for the two factors. Taking the coefficient
between the distinguished vector and its translate gives
\(\int_K m_\nu(k)m_\nu(k^{-1}j)\dd k=m_\nu(j)/s_\nu\).
Inversion conjugates the coefficient. Integration using these
identities shows \(P_\nu^*=P_\nu=P_\nu^2\).

Put \(F_\nu(x,y)=\int_K\overline{m_\nu(k)}\psi(xky)\dd k\).
For \(x,y\) in sufficiently small fixed neighborhoods of identity
perform the change and elimination of
Theorem~\ref{thm:real-compact-product}. Choose any integers
\(o_1<\cdots<o_n\) and put \(D_o(s)=\operatorname{diag}(s^{o_i})\)
and \(N=\sum o_i(\nu_i-\Lambda_i)\). In the notation of that
proof the circle integrand, apart from \(1/c_\Lambda\), is
\[
 s^N H_\nu(s;v,j)
 \frac{m_\Lambda(aj)m_\Lambda(v^{-1})}
      {m_{-\mu}(Z_j U_v(s))\,m_{-\mu}(Y_v)},\qquad
 H_\nu=s^{-\sum o_i\nu_i}m_\nu(vD_o(s)j^{-1}).
\]
Here \((U_v(s))_{il}=s^{o_l-o_i}(U_v)_{il}\) for \(i<l\).
In the block minor expansion of Lemma~\ref{lem:circle-leading-minors}
the least exponent on an index set of prescribed size still
belongs to its initial segment. Thus \(H_\nu\) is a polynomial
with coefficients finite sums of products of functions of \(v\)
and functions of \(j\). A negative power of a full block
determinant merely contributes separated nonzero determinant
factors. The displayed denominator stays nonzero on the unit
disk after shrinking the neighborhoods, by the strict accretivity
argument in Theorem~\ref{thm:real-compact-product}.

The circle mean vanishes if \(N>0\). If \(N\le0\), put \(h=-N\).
The mean extracts coefficient \(h\) from the holomorphic factor
after \(s^N\). Expand \(1/m_{-\mu}(Z_j(I+V))\) at \(V=0\)
in the strictly upper triangular entries of \(V\).
It is rational and regular there, so its Taylor terms of total
degree at most \(h\) are a finite sum of monomials in \(V\)
with coefficients depending only on \(Z_j\). This also follows
by expanding each nonzero denominator about its constant term
as a geometric series through degree \(h\). Substitution of
\(V=U_v(s)-I\) makes all higher terms vanish through degree \(h\)
since every entry is divisible by \(s\). After multiplication
by \(H_\nu\) the selected coefficient is a finite sum of
products of functions of \((x,j)\) and functions of \((y,v)\).
The number \(L\) of products can be chosen solely from the
dimensions, weights and \(o\). Compact integration in \(j,v\)
gives \(F_\nu(x,y)=\sum_{l=1}^L A_l(x)B_l(y)\) in these
neighborhoods. The zero case is included with \(L=0\).

Every \((L+1)\)-rowed determinant
\(\det[F_\nu(x_i,y_j)]\) consequently vanishes there.
For arbitrary arguments join each \(x_i,y_j\) to identity
by its analytic path constructed in
Theorem~\ref{thm:real-compact-product}. Along these simultaneous
paths each entry is real analytic: use the fixed splitting
at a real parameter value, uniform over the compact variable,
and integrate the local Cauchy series as in that theorem.
The identity theorem in Lemma~\ref{lem:analytic-tools} extends
the vanishing to the prescribed endpoints. By
Lemma~\ref{lem:cyclic-row-isometry} the entries of the
projected Gram matrix for any translates \(\pi(y_i)v_0\) equal
\(s_\nu F_\nu(y_i^{-1},y_j)\). Thus every \(L+1\) such
projected vectors are dependent: their Gram determinant is
zero precisely when the linear combinations have a nonzero
null vector, by the squared norm formula. They span a space
of dimension at most \(L\), which is closed since a finite
orthonormal basis recovers the coordinates continuously.
Density of the translates gives the rank bound. For \(\Lambda\),
Theorem~\ref{thm:real-compact-product} factors the Gram matrix
and \(s_\Lambda c_\Lambda=1\) by the compact average above,
so its rank equals one by evaluation at identity.

For positive rank, vanishing when \(N>0\) shows
\(\sum o_i(\nu_i-\Lambda_i)\le0\) for every such \(o\). Adding
an arbitrary integer to each \(o_i\) gives equality of the
total weights. Summation by parts gives the negative of
\(\sum_{j<n}(o_{j+1}-o_j)\sum_{i\le j}(\nu_i-\Lambda_i)\).
Taking any single gap arbitrarily large proves the prefix
inequalities. Conversely those inequalities and the vanishing
total imply all these scalar inequalities.
\end{proof}

\section{The compact character averages of the orbital integral}

\begin{proposition}\label{prop:row-integral-multiplicity}
For a compact dominant weight as in
Definition~\ref{def:cyclic-row-module}, the function
\[
 B_\nu(g)=\int_K\overline{\chi_\nu(k)}
                    \psi(g^{-1}kg)\dd k
\]
is nonnegative and integrable on \(G\), and
\(d\int_G B_\nu(g)\dd h_S(g)=\operatorname{rank}P_\nu\).
\end{proposition}
\begin{proof}
Let \(b_1,\ldots,b_m\) be an orthonormal basis of the range of
\(P_\nu\), finite by Theorem~\ref{thm:row-compact-finite-rank}.
For \(w_g=\pi(g)v_0\),
\[
 \langle w_g,P_\nu w_g\rangle
   =\sum_{i=1}^m|\langle b_i,\pi(g)v_0\rangle|^2
   =s_\nu\int_K\overline{m_\nu(k)}\,\psi(g^{-1}kg)\dd k .
\]
Lemma~\ref{lem:cyclic-row-isometry} gives integral \(m/d\)
for this nonnegative function. Replace \(g\) by \(ug\) and
average in \(u\in K\). On the right conjugate the compact variable
by \(u\). The conjugacy average of \(m_\nu\) is
\(\chi_\nu/s_\nu\), by Theorem~\ref{thm:compact-alternant}
on each factor. The average on the left is consequently \(B_\nu\).
Integrating it over \(G\) leaves the integral unchanged by
left invariance and nonnegative iterated integration as in
Lemma~\ref{lem:integral-tools}.
\end{proof}

\begin{definition}\label{def:opposite-polynomial}
For diagonal coordinates \(z=(z_1,\ldots,z_n)\) put
\[
 Q(z)=\prod_{\substack{i<j\\\varepsilon_i\ne\varepsilon_j}}(z_j-z_i),
 \qquad V(z)=\prod_{\substack{i<j\\\varepsilon_i=\varepsilon_j}}(z_i-z_j).
\]
An empty product is one. Write \(Q(k)\) on \(K\) for the
value on its two lists of eigenvalues in the respective sign
positions. Up to a fixed sign \(Q(z)\) is
\(\prod_{\varepsilon_i=+1,\varepsilon_j=-1}(z_i-z_j)\);
it is symmetric in each list.
\end{definition}

\begin{lemma}\label{lem:elliptic-resultant-bound}
For a constant depending only on \(n\) and any \(k\in K,\ g\in G\),
\[
       |Q(k)|\,\|g\|_{\rm op}^2
                \le C_n(1+\|g^{-1}kg\|_{\rm op})^{n-1}.
\]
The function \(Q\) is continuous on \(K\).
\end{lemma}
\begin{proof}
With just one sign, \(G\) is compact unitary,
\(\|g\|_{\rm op}=1\) and \(Q=1\). Otherwise let \(P_+,P_-\)
be the monic characteristic polynomials of the two blocks of
\(k\), of degrees \(p,q'>0\). The map on polynomials of degree
less than \(p\) given by multiplication by \(P_-\) modulo \(P_+\)
has determinant \(R=\prod_{\varepsilon_i=+1}P_-(z_i)\). For
distinct roots of \(P_+\), evaluation at them identifies this map
with a diagonal matrix, by division and polynomial interpolation.
For repeated roots the identity follows by continuity in the
roots, since division by a monic polynomial is polynomial in
its coefficients. Hence \(R\) agrees with \(Q(k)\) up to a
fixed sign and is continuous. The coefficients of the map and
its adjugate are bounded in terms of \(n\), since both input
polynomials have roots on the unit circle. Apply the adjugate
to the constant polynomial \(1\) to get \(B\) of degree \(<p\)
with \(BP_-\equiv R\bmod P_+\). Then \(H=BP_-\) has degree
at most \(n-1\) and bounded coefficients. Its value at every
positive-block eigenvalue is \(R\), and at every negative-block
eigenvalue is zero. By block unitary diagonalization
\(H(k)=R E_+\) with \(E_+=(I+S)/2\).

Write \(T=\|g\|_{\rm op}\). Since \(g^{-1}Sg=Sg^*g\),
\(2\|g^{-1}E_+g\|_{\rm op}\ge T^2-1\).
Evaluation of \(H\) at \(g^{-1}kg\) bounds
\(|R|\,\|g^{-1}E_+g\|_{\rm op}\) by a constant times the
right-hand power in the statement. Also \(|R|\le2^{pq'}\).
Adding these two estimates gives the bound for \(T^2\).
\end{proof}

\begin{theorem}\label{thm:regularized-row-character}
Set \(N=4n^2+1\) and assume the moment hypothesis of
Theorem~\ref{thm:row-real-orthogonality} and the sign-change
gap of Lemma~\ref{lem:sign-run-determinant} with
\(c\ge (n-1)N\). The function
\[
 H(k)=d\int_G Q(k)^N\psi(g^{-1}kg)\dd h_S(g)
\]
is continuous and conjugation invariant on \(K\).
Whenever the two blocks have disjoint spectra the orbital
integral of \(\psi\) itself is absolutely convergent and
equals \(H(k)/(dQ(k)^N)\). These conclusions apply to
the arithmetic rows of Corollary~\ref{cor:row-polar-moment}
at all sufficiently large scales.
\end{theorem}
\begin{proof}
Both \(\|g^{-1}kg\|_{\rm op}\) and \(\|g\|_{\rm op}\) are at
least one because elements of \(U(S)\) have determinant of
modulus one. Combine Lemma~\ref{lem:elliptic-resultant-bound}
with Theorem~\ref{thm:polar-row-decay}. The modulus of the
integrand, before its scalar \(d\), is bounded by a constant
depending on \(n,c\) times \(\|g\|_{\rm op}^{-2N}\), uniformly
in \(k\). Its integral is finite by summing dyadic shells with
Lemma~\ref{lem:real-unitary-ball}, since \(2N>4n^2\).
Dominated convergence from Lemma~\ref{lem:integral-tools}
and continuity of the integrand give continuity of \(H\).
Left translation of \(g\) gives conjugacy invariance for all
\(k\). Division by the fixed nonzero \(Q(k)^N\) gives the
unregularized assertion, including at repeated roots within
a single block. Finally the gap \(r^8+4r^2+2\) used in
Corollary~\ref{cor:row-polar-moment} exceeds
\((r-1)(4r^2+1)\) for all sufficiently large \(r=n\).
\end{proof}

\section{The Laurent polynomial on the compact Cartan}

\begin{definition}\label{def:compact-alternation}
Let \(\mathfrak W\) be the permutations preserving the two sign
sets. Write \(z^\alpha=\prod_i z_i^{\alpha_i}\),
\((w\alpha)_i=\alpha_{w^{-1}i}\), and
\(A_\alpha(z)=\sum_{w\in\mathfrak W}\operatorname{sgn}(w)
z^{w\alpha}\). Put \(\delta_i=\#\{j>i:\varepsilon_j=\varepsilon_i\}\).
For an exponent \(\alpha\in\Z^n\) with distinct coordinates
within each sign set there is a unique \(w\in\mathfrak W\)
such that \(w\alpha=\nu+\delta\) with \(\nu\) compact dominant:
sort each block decreasingly and subtract \(\delta\).
Put \(X_\alpha=\operatorname{sgn}(w)\chi_\nu\) and
\(t(\alpha)=\operatorname{sgn}(w)\operatorname{rank}P_\nu\).
For an exponent with a repeated coordinate within a block put
\(X_\alpha=t(\alpha)=0\).
\end{definition}

\begin{lemma}\label{lem:compact-character-density}
The \(\chi_\nu\) form an orthonormal system in \(L^2(K,\dd k)\).
Their finite linear combinations are uniformly dense in the
continuous class functions on \(K\). On its torus
\(V X_\alpha=A_\alpha\).
\end{lemma}
\begin{proof}
For two single-block modules average a map between them by
\[
       M\longmapsto\int T_\nu(k)M T_\lambda(k)^{-1}\dd k .
\]
Invariance shows that the average is a projection onto the
space of intertwining maps. A nonzero intertwiner of irreducibles
is an isomorphism by its kernel and image. An endomorphism is
scalar since it has a complex eigenvalue and its eigenspace
is invariant. Distinct highest weights cannot give isomorphic
modules: by differentiating an intertwining equation their
complexified actions intertwine as in
Theorem~\ref{thm:compact-highest-vector}, and the upper
unitriangular-fixed lines have their specified diagonal weights.
The trace of this averaging operator is therefore \(1\) for
equal weights and \(0\) otherwise. Its trace in the elementary
matrix basis is \(\int\chi_\nu\overline{\chi_\lambda}\).
Multiplying the two block identities gives orthonormality.

Expansion of the two block determinants gives \(A_\alpha\),
and the ordinary Vandermonde product gives \(V=A_\delta\).
Theorem~\ref{thm:compact-alternant} on the two blocks gives
\(V X_\alpha=A_\alpha\), including repeated coordinates by
pairing the terms of opposite sign. Any continuous class
function is determined by its restriction to the torus by
unitary diagonalization. Approximate the restriction uniformly
by trigonometric polynomials with the product of the
one-dimensional Fej\'er kernels of
Lemma~\ref{lem:poisson-finite-torus}. Indeed their nonnegative
integrals are one and their mass outside any prescribed torus
neighborhood of identity tends to zero, so uniform continuity
gives the approximation. Average a polynomial approximant in
\(\mathfrak W\) to make it invariant. Multiplying it by \(V\)
gives a finite alternating sum of monomials. A repeated
coordinate forces zero coefficient by transposition, and collecting
each other orbit expresses it in the basis \(A_{\nu+\delta}\)
with coefficient that of \(z^{\nu+\delta}\). Divide by \(V\)
where the same-block coordinates are distinct, and continue
to the whole torus by continuity. Thus the symmetric approximants
are restrictions of finite character sums, with the same uniform
error on \(K\).
\end{proof}

\begin{theorem}\label{thm:row-orbital-polynomial}
Under the hypotheses of Theorem~\ref{thm:regularized-row-character}
there is an alternating Laurent polynomial \(P\) for
\(\mathfrak W\), with integer coefficients, such that
\[
 P(z)=V(z)H(\operatorname{diag}z),\qquad
 [z^\alpha]P=\sum_\beta [z^\beta]Q^N\,t(\alpha-\beta).
\]
Here \([z^\alpha]\) denotes the Laurent coefficient. In particular
\(t(\Lambda+\delta)=1\). If \(t(\gamma)\ne0\), some permutation
of its two blocks equals \(\nu+\delta\) with
\(\sum_{i\le j}(\nu_i-\Lambda_i)\ge0\) for all \(j\) and
total sum zero. If a Laurent polynomial agrees with \(P\) on
a nonempty open subset of the unit torus, it agrees everywhere.
\end{theorem}
\begin{proof}
Write \(c_\beta=[z^\beta]Q^N\), integers invariant under
\(\mathfrak W\) by Definition~\ref{def:opposite-polynomial}.
On the torus
\(\overline{Q}^N X_\alpha=\sum_\beta c_\beta X_{\alpha-\beta}\):
multiply by \(V\), alternate
\(z^\alpha\sum_\beta c_\beta z^{-\beta}\), and use
Lemma~\ref{lem:compact-character-density}. By block conjugation
and continuity the identity holds on \(K\). For fixed \(g\)
multiply this identity after conjugation by \(\psi(g^{-1}kg)\)
and integrate over \(k\). On integrating over \(g\), the left
side permits exchange by the uniform integrable bound in
Theorem~\ref{thm:regularized-row-character}. Each compact
average on the right is integrable by
Proposition~\ref{prop:row-integral-multiplicity}, including its
possible sign. Consequently
\[
 \langle X_\alpha,H\rangle_{L^2(K)}
                      =\sum_\beta c_\beta\,t(\alpha-\beta).
\]
In particular all \(b_\nu=\langle\chi_\nu,H\rangle\) are
integers. For any finite set of distinct \(\nu\), expanding the
squared norm of \(H-\sum b_\nu\chi_\nu\) and using orthonormality
gives \(\sum|b_\nu|^2\le\|H\|^2\). Thus only finitely many
are nonzero. Subtracting their finite character sum gives
a continuous class function orthogonal to all compact
characters. Its squared integral is zero by uniform approximation
in Lemma~\ref{lem:compact-character-density}. It vanishes everywhere,
since every nonempty open subset of the compact group has
positive measure by the smooth density of
Definition~\ref{def:unitary-density}. Now
\(P=\sum_\nu b_\nu A_{\nu+\delta}\) has all the stated coefficients,
as \([z^\alpha]P=\langle X_\alpha,H\rangle\).
The assertions about \(t\) follow from
Theorem~\ref{thm:row-compact-finite-rank}. Finally an open
torus set contains a product of open arcs. After clearing
negative powers, hold all but one variable fixed in these arcs.
A one-variable polynomial vanishing there has infinitely many
roots and is zero. Repetition in all variables proves uniqueness.
\end{proof}

\section{Compression along a signed unit vector}

\begin{lemma}\label{lem:signed-column-measure}
Put \(\eta=\varepsilon_1\), \(S'=\operatorname{diag}
(\varepsilon_2,\ldots,\varepsilon_n)\). In the first-column map
\(g\mapsto w=ge_1\) from \(U(S)\), the range is
\(w^*Sw=\eta\) and the fiber has group \(U(S')\). Use
\(h_S,h_{S'}\) of Definition~\ref{def:unitary-density} and
mass one for the trivial group in dimension zero. After integrating
the coordinate phases, the quotient measure in \(h_S\) is
\((n-1)!\prod_{i\ne j}\dd s_i\), where \(s_i=|w_i|^2>0\)
and \(s_j\) is determined by \(\sum_i\varepsilon_i s_i=\eta\).
Every choice of eliminated index \(j\) gives this formula.
\end{lemma}
\begin{proof}
The complement of a vector of norm \(\eta\) for the signed
form is nonsingular and has the residual signature: decompose
orthogonally into this line and its complement, diagonalize the
hermitian form on the complement and count positive and negative
dimensions, as in Theorem~\ref{thm:hermitian-local-global}.
Thus \(w\) extends to a signed orthonormal frame. The stabilizer
of \(e_1\) has the stated block form. Local smooth choices of
frame follow by projecting a fixed basis of the complement
at one point and performing successive signed orthogonalization;
the pivots keep their nonzero signs in a neighborhood.

On the quadric define a density by taking twice ambient Lebesgue
volume per unit change of \(w^*Sw\). In local coordinates
\((t,y)\) with \(t=w^*Sw\) this means twice the coefficient of
\(\dd t\,\dd y\) in ambient volume at \(t=\eta\), times \(\dd y\).
The derivative in a real coordinate of a nonzero entry can be
made nonzero by a phase, so these coordinates cover the quadric.
Changes of the \(y\)'s give the ordinary density transformation.
This density is invariant since \(G\) preserves \(t\) and
has complex determinant of modulus one. At \(w=e_1\) it is
ordinary tangent coordinate measure, as the absolute derivative
in the first real direction is \(2\). At identity the map on
the skew-hermitian tangent coordinate in
Definition~\ref{def:unitary-density} takes its first column
multiplied by \(S\), leaving a free lower square block.
Signs do not change this coordinate density. Thus the quotient
density for the two unnormalized measures is precisely the
density just constructed. The product formula in local sections
follows by the same block triangular change of coordinates as
in Theorem~\ref{thm:compact-dual-unitary-volume}. A countable
chart cover suffices here: use the countable rational balls
in matrix coordinates and partition into successive differences.
For nonnegative integrands one can therefore integrate first
over the fiber and then over the column.

For normalized measures multiply the quotient density by
\(u_{n-1}/u_n=(n-1)!/(2\pi^n)\). In planar polar coordinates
ordinary complex volume is \(\dd s_i\,\dd\theta_i/2\).
Eliminating one \(s_j\) in favor of \(t\) has Jacobian of
absolute value one. All phases range from \(0\) to \(2\pi\);
their integration contributes \(2\pi^n\) including the initial
factor two. This gives the stated quotient density. A vanishing
coordinate occupies a set of measure zero: near such a point
eliminate a real coordinate belonging to a different nonzero
entry, leaving two free real coordinates that have to vanish.
The case \(n=1\) gives the usual probability on the circle.
\end{proof}

\begin{theorem}\label{thm:signed-compression}
Let \(A\) be positive diagonal with signed eigenvalues
\(x_i=\varepsilon_i A_{ii}\) pairwise distinct. For a first
column \(w\), put \(B=(g^*Ag)[2{:}n,2{:}n]\) in any signed
orthonormal completion. Write \(y\) for the signed eigenvalues
of \(S'B\), and \(\mathcal V(x)=\prod_{i<j}|x_i-x_j|\).
The image measure on \(y\), ordered by decreasing magnitude
within each sign, has density
\[
                (n-1)!\frac{\mathcal V(y)}{\mathcal V(x)}
                         \prod_l\dd y_l
\]
with positive coordinate volume. For roots of sign \(\eta\)
there is one new magnitude between each consecutive pair of
old magnitudes. For roots of the opposite sign there is one
between each consecutive pair and one above the largest old
magnitude if that sign occurs. These intervals vary independently,
and their endpoints may be omitted. The empty list has
Vandermonde product one.
\end{theorem}
\begin{proof}
Since \(B>0\), the matrix \(S'B\) is similar to the hermitian
matrix \(B^{1/2}S'B^{1/2}\), whose signature is that of \(S'\)
by the change of quadratic form coordinates. Its eigenspaces
of distinct eigenvalues are orthogonal for \(S'\) since
\((S'B)^*S'=S'(S'B)\). On a \(\lambda\)-eigenspace,
\(v^*Bv=\lambda v^*S'v\), so the restriction of \(S'\) is
definite of the sign of \(\lambda\). Choosing normalized bases
there shows that \(B\) is congruent by \(U(S')\) to the positive
diagonal of the magnitudes \( |y_l| \) in sign order. In particular
its signed eigenvalues are independent of the completion.

In the basis \(g\), the bottom square block of \(g^{-1}SA g\)
equals \(S'B\). Its resolvent has first diagonal entry given
both by its cofactor and by \(e_1^*g^{-1}=\eta w^* S\).
Thus, with \(z\) an indeterminate,
\[
   \frac{\prod_l(z-y_l)}{\prod_i(z-x_i)}
                       =\sum_i\frac{u_i}{z-x_i},\qquad
              u_i=\eta\varepsilon_i s_i,\quad \sum_i u_i=1 .
\]
Multiply all signed eigenvalues by \(\eta\) to discuss their
order, so the removed sign is positive and the \(u_i\) have
the signs of these rescaled roots \(x_i\). Between consecutive
poles of the same sign the real rational function tends to
infinities of opposite signs at the two ends. If negative
poles occur, it is positive immediately to the left of the
leftmost one and negative far to its left, since it is
asymptotic to \(1/z\). These intervals account for all
\(n-1\) roots, exactly one simple root apiece. Conversely any
points in those intervals give the residues
\[
             u_i=\frac{\prod_l(x_i-y_l)}
                              {\prod_{k\ne i}(x_i-x_k)}.
\]
In decreasing signed order the number of roots above \(x_i\)
equals the number of poles above it for a positive pole, and
is one less for a negative pole. The residues therefore have
the required signs. Their sum is one by the leading term at
infinity, using the partial fractions obtained by subtracting
the displayed residues from the rational function. Thus
the correspondence with interior \(s\)-coordinates is bijective.

Omit the last residue, which is determined by their sum.
Its remaining derivative matrix has entries
\(-u_i/(x_i-y_l)\), \(1\le i,l\le n-1\).
The absolute Cauchy determinant for the entries \(1/(a_i-b_j)\)
is \(\mathcal V(a)\mathcal V(b)/\prod_{ij}|a_i-b_j|\).
Indeed clearing the denominators leaves a polynomial alternating
in each list, divisible by both Vandermonde products by setting
two arguments equal and dividing off their difference. The
degrees agree, so the quotient is constant. Multiplying the
uncleared determinant by \(a_m-b_m\) and setting \(a_m=b_m\)
reduces it to the determinant in size \(m-1\), giving constant
of modulus one by induction. Substitution of the residue
formula now gives the absolute Jacobian
\(\mathcal V(y)/\mathcal V(x)\); all \(|x_i-y_l|\) cancel.
The change between \(u_i\) and \(s_i\), any reorderings and
changes of signs of the variables preserve the absolute
Jacobian. Apply Lemma~\ref{lem:signed-column-measure}. In
dimension one both coordinate lists of the density are empty
after omission and the conclusion is immediate.
\end{proof}

\section{Evaluation of the positive principal-minor integral}

\begin{definition}\label{def:positive-flag-integral}
Let \(b_1>\cdots>b_n\) be integers, \(\mu_i=b_i-i+1\),
and \(A=\operatorname{diag}(a_i)\) with \(a_i>0\).
For \(S\) as in Definition~\ref{def:unitary-density} put
\[
 J_S(\mu,A)=\int_{U(S)}
             \frac{\dd h_S(g)}{m_{-\mu}(g^*Ag)},\qquad
 D(b)=\prod_{i<j}\frac{b_i-b_j}{j-i}.
\]
The integrand is positive. For each sign \(\sigma\), list its
indices \(I_\sigma\) in their original order, and write \(a^\sigma\)
for the list of magnitudes of that sign in decreasing order.
For these two lists set
\(E_\sigma=\det[(a^\sigma_j)^{b_i}]_{i\in I_\sigma,j}\).
The empty determinant is one.
\end{definition}

\begin{theorem}\label{thm:positive-flag-evaluation}
When the \(a_i\) are distinct within each sign, the positive
integral of Definition~\ref{def:positive-flag-integral} is finite
and equals
\[
               J_S(\mu,A)=
          \frac{E_+E_-}{D(b)\,\mathcal V((\varepsilon_i a_i)_i)} .
\]
\end{theorem}
\begin{proof}
In dimension one the group is compact of mass one and the
integrand is \(a_1^{b_1}\). For higher dimension the trailing
minors in Definition~\ref{def:polar-row-coefficient} give
\[
 m_{-\mu}(g^*Ag)=(\det A)^{-\mu_1}m_{-\widehat\mu}(B),
       \qquad \widehat\mu_{i-1}=\mu_i-\mu_1\quad(2\le i\le n),
\]
where \(B=(g^*Ag)[2{:}n,2{:}n]\). Indeed the full determinant
has modulus factor \(|\det g|^2=1\); its exponent is \(-\mu_1\).
All shorter trailing minors have their successive weight
differences, including the full determinant of \(B\) with
exponent \(\mu_1-\mu_2\). On the first-column fiber, changing
the remaining basis replaces \(B\) by \(h^*Bh\).
Theorem~\ref{thm:signed-compression} conjugates this positive
form by \(U(S')\) to the diagonal with signed eigenvalues \(y\).
Thus by invariance the fiber integral is \(J_{S'}\) with these
magnitudes and weights \(\widehat\mu\). Its integer exponents
in Definition~\ref{def:positive-flag-integral} are
\(\widehat b_{i-1}=b_i-b_1-1\), still strictly decreasing.
Nonnegative integration in Lemma~\ref{lem:signed-column-measure}
permits this step before finiteness is known.

Use the induction formula in the eigenvalue density of
Theorem~\ref{thm:signed-compression}. The factor
\(\mathcal V(y)\) cancels, leaving one determinant in the
magnitudes of each sign. In each determinant its \(j\)-th
column depends only on one independent interval variable.
Integrate by columns: expanding the determinant gives a finite
sum of products of absolutely integrable powers. The finite
intervals have positive endpoints and all powers at infinity
are \(b_i-b_1-1<-1\). This proves that the manipulations are
absolutely integrable.

For the removed sign write its old magnitudes \(a_1'>\cdots>a_m'\).
The entry with row \(i\ne1\) and interval
\((a_{j+1}',a_j')\) integrates to
\[
       \frac{(a_{j+1}')^{b_i-b_1}-(a_j')^{b_i-b_1}}{b_1-b_i}.
\]
Extract the denominator once per row. The remaining determinant
equals the \(m\)-rowed determinant on the \(a_j'\) with row
exponents \(b_i-b_1\) including its initial row of ones. Indeed
in the latter matrix replace each column after the first by
itself minus its predecessor, working from right to left, then
expand in its first row. This includes \(m=1\).

For the opposite sign write its old magnitudes
\(h_1>\cdots>h_e\), \(h_0=\infty\). The retained variables
lie in \((h_j,h_{j-1})\). After taking the same row denominators
out, the column entries are \(h_j^{b_i-b_1}-h_{j-1}^{b_i-b_1}\),
with \(\infty^{b_i-b_1}=0\). These differences have determinant
\(\det[h_j^{b_i-b_1}]\), by the same triangular column
operation. If \(e=0\) its contribution is one. The original
prefactor \((\det A)^{\mu_1}=(\prod a_i)^{b_1}\) puts back
the missing power in every column of both determinants. Their
product is \(E_+E_-\). The remaining factor multiplying
\(1/\mathcal V((\varepsilon_i a_i)_i)\) is
\[
       \frac{(n-1)!}{D(\widehat b)\prod_{i>1}(b_1-b_i)}
                           =\frac1{D(b)}.
\]
This uses cancellation of the common factors in all pairs with
\(i,j>1\) and \(\prod_{j=2}^n(j-1)=(n-1)!\). It completes
the induction.
\end{proof}

\section{From positive diagonal matrices to the elliptic torus}

\begin{lemma}\label{lem:accretive-diagonal-integral}
For the data of Definition~\ref{def:positive-flag-integral} put
\[
 F(a)=\int_G\frac{\dd h_S(g)}{m_{-\mu}(g^*\operatorname{diag}(a)g)},
 \qquad
 W(a)=\prod_{\substack{i<j\\\varepsilon_i=\varepsilon_j}}(a_i-a_j)
       \prod_{\substack{i<j\\\varepsilon_i\ne\varepsilon_j}}(a_i+a_j).
\]
For \(\Rp a_i>0\) the integral is absolutely convergent and is
holomorphic separately in all coordinates. With the block
alternant of Definition~\ref{def:compact-alternation} one has
\(D(b)W(a)F(a)=A_b(a)\) on this product of half-planes.
\end{lemma}
\begin{proof}
Fix a positive diagonal \(A_0\) with distinct entries. For any
full-column-rank submatrix \(C\) of the columns of \(g\), put
\(L=C^*C>0\). If \(\min_i\Rp a_i\ge h>0\), the matrix
\(L^{-1/2}C^*\operatorname{diag}(a)C L^{-1/2}\) has real
part at least \(hI\), in the sense of hermitian forms. Each
eigenvalue has real part at least \(h\) by evaluation on its
eigenvector. It follows that the modulus of its determinant
is at least \(h^{\dim L}\), since the determinant is the product
of its eigenvalues with multiplicity. The same argument for the
hermitian matrix with \(A_0\) bounds its determinant above by
\(\|A_0\|_{\rm op}^{\dim L}\). Therefore the modulus of each
trailing subdeterminant at \(a\) is bounded below by its value
at \(A_0\) times \((h/\|A_0\|_{\rm op})^{\dim L}\).
All exponents of the proper minors of \(m_{-\mu}\) are
nonnegative. Its full determinant factor is
\((\prod a_i)^{-\mu_1}\) and is uniformly bounded above and
away from zero on compact subsets of the domain. Thus the
absolute integrand there is bounded by a fixed multiple of the
positive integrand at \(A_0\), integrable by
Theorem~\ref{thm:positive-flag-evaluation}. Pointwise it is a
rational holomorphic function with nonzero denominators.
The Cauchy integral argument of Lemma~\ref{lem:analytic-tools}
gives separate holomorphy.

On real positive entries decreasing within each sign, \(W(a)\)
equals \(\mathcal V((\varepsilon_i a_i)_i)\) and \(A_b(a)\)
is the product \(E_+E_-\) of
Definition~\ref{def:positive-flag-integral}. Hence the identity
holds on an open real box by Theorem~\ref{thm:positive-flag-evaluation}.
Keeping all other variables in their box intervals, continue it
in one variable to its connected half-plane by
Lemma~\ref{lem:analytic-tools}. Repeat with those already extended
variables held fixed anywhere in their half-planes. This gives
the asserted identity throughout, even when \(W(a)=0\).
\end{proof}

\begin{theorem}\label{thm:elliptic-unnormalized-character}
Use the hypotheses of Theorem~\ref{thm:row-orbital-polynomial},
with \(b_i=t_i+(n-1)/2=\mu_i+i-1\). Then
\[
       P=\frac{d}{D(b)}Q^{N-1} A_b .
\]
In particular for pairwise distinct unit \(z_i\),
\[
 \int_G\psi(g^{-1}\operatorname{diag}(z)g)\dd h_S(g)
                  =\frac{A_b(z)}{D(b)\,V(z)Q(z)} .
\]
\end{theorem}
\begin{proof}
First take \(z_i\) on the unit circle with \(\Rp(\varepsilon_i z_i)>0\)
and with all \(z_i\) distinct. Put \(a_i=\varepsilon_i z_i\)
and \(A=\operatorname{diag}(a)\). For \(h=g^{-1}\operatorname{diag}(z)g\)
we have \(S h=g^* A g\), strictly accretive. Hence \(S\) is an
admissible block splitting for \(h\) in
Theorem~\ref{thm:accretive-splitting}. The diagonal multiplication
of the numerator gives
\[
 \psi(h)=\frac{\sigma}{c_\Lambda}
       \int_K\frac{|m_\Lambda(v)|^2}
                   {m_{-\mu}((gv)^*A(gv))}\dd v,\qquad
 \sigma=\prod_i\varepsilon_i^{\Lambda_i}.
\]
The integral of the absolute value over \(g,v\) is finite by
Lemma~\ref{lem:accretive-diagonal-integral}, right invariance in
\(g\), and the definition of \(c_\Lambda\). Integrating first
in \(g\) gives \(\sigma F(a)\).

In the sign substitution \(z_i=\varepsilon_i a_i\), the two
block determinants in \(A_b\) acquire together the factor
\((-1)^{\sum_{\varepsilon_i=-1}b_i}\). The factors of \(VQ\)
acquire, relative to \(W(a)\), one minus sign for each
\(i<j\) with \(\varepsilon_j=-1\). Indeed a pair of negative
signs uses \(z_i-z_j\) in \(V\), and an opposite pair ending
in a negative sign uses \(z_j-z_i\) in \(Q\). Thus the total
sign in \(A_b(z)/(VQ)(z)\) relative to \(A_b(a)/W(a)\) is
\((-1)^{\sum_{\varepsilon_i=-1}(b_i+i-1)}=\sigma\).
For the last equality use \(b_i=\mu_i+i-1\) and
\(\sum_{\varepsilon_i=-1}\mu_i=\sum_{\varepsilon_i=-1}\Lambda_i\);
the earlier and later same-sign pairs cancel in
Definition~\ref{def:polar-row-coefficient}. By
Lemma~\ref{lem:accretive-diagonal-integral} this proves the
orbital formula on the prescribed open set of the torus.

By Theorems~\ref{thm:regularized-row-character}
and~\ref{thm:row-orbital-polynomial}, multiplication by
\(dVQ^N\) gives \(P\) there. Both sides of the resulting
identity are Laurent polynomials. It therefore holds everywhere
by the uniqueness in Theorem~\ref{thm:row-orbital-polynomial}.
Dividing on the regular torus proves the last formula with the
convergence asserted in Theorem~\ref{thm:regularized-row-character}.
\end{proof}

\section{The formal degree and the elliptic character identity}

\begin{lemma}\label{lem:extreme-orbital-coefficient}
Use \(b,N,Q,\delta,\Lambda\) of
Theorems~\ref{thm:row-orbital-polynomial}
and~\ref{thm:elliptic-unnormalized-character}. Put
\(u_i=\#\{j<i:\varepsilon_j\ne\varepsilon_i\}\) and
\(\kappa=b+(N-1)u\). Suppose \(\kappa\) is strictly decreasing
within each sign in its index order. Then
\[
           [z^\kappa]P=[z^\kappa](Q^{N-1} A_b)=1.
\]
\end{lemma}
\begin{proof}
Put \(o_i=i\), with \(\langle o,\beta\rangle=\sum i\beta_i\).
In \(Q^j\), for \(j\ge0\), the unique monomial with maximum
\(\langle o,\cdot\rangle\) is \(z^{ju}\), with coefficient one.
Indeed in each factor for \(i<h\) the term \(z_h\) has strictly
greater level than \(-z_i\). The support of \(Q^j\) is invariant
under \(\mathfrak W\) by Definition~\ref{def:opposite-polynomial}.
For \(w\in\mathfrak W\),
\(\langle o,w\kappa\rangle\ge\langle o,\kappa\rangle\)
with equality only for \(w=1\). In each sign block, interchanging
a smaller preceding value with a larger following one strictly
decreases its scalar product with the increasing indices.
Repeatedly removing such inversions uniquely restores the
strictly decreasing list.

By Definitions~\ref{def:polar-row-coefficient}
and~\ref{def:compact-alternation},
\(\Lambda_i+\delta_i=\mu_i+\#\{j<i:\varepsilon_j=\varepsilon_i\}
=b_i-u_i\). Suppose \(\beta\) contributes to the convolution
for \([z^\kappa]P\) in Theorem~\ref{thm:row-orbital-polynomial}.
Choose \(w\) so \(w(\kappa-\beta)=\nu+\delta\) is sorted.
The prefix inequalities and equal total in that theorem give
\(\langle o,\nu-\Lambda\rangle\le0\) by summation by parts.
Since \(w\beta\) occurs in \(Q^N\), it follows that
\[
 \langle o,w\kappa\rangle
       \le\langle o,\Lambda+\delta+Nu\rangle
        =\langle o,\kappa\rangle .
\]
Thus \(w=1\); equality also forces \(\beta=Nu\). Its coefficient
is one and \(t(\kappa-Nu)=t(\Lambda+\delta)=1\).
Conversely this term occurs, giving the first assertion.
For a term \(\beta+w b=\kappa\) in \(Q^{N-1}A_b\),
apply \(w^{-1}\). Since \(w^{-1}\beta\) occurs in \(Q^{N-1}\),
\(\langle o,w^{-1}\kappa\rangle\le\langle o,b+(N-1)u\rangle\).
The same strictness gives \(w=1\), leaving only the leading
monomial of \(Q^{N-1}\) with positive coefficient one.
\end{proof}

\begin{theorem}\label{thm:real-row-stable-character}
Under the hypotheses of Lemma~\ref{lem:extreme-orbital-coefficient},
the reciprocal squared norm \(d\) of the unit-normalized-at-identity
coefficient \(\psi\) is \(D(b)\). With \(p\) positive and \(q'\)
negative signs, \(e=d\,\overline\psi\) and
\(\lambda_i=t_i-(n+1)/2+i\), for pairwise distinct unit \(z_i\)
we have
\[
 \frac{(-1)^{pq'}}{p!\,q'!}\sum_{\pi\in\mathfrak S_n}
       \int_G e\bigl(g^{-1}\operatorname{diag}
                  (z_{\pi(1)},\ldots,z_{\pi(n)})g\bigr)\dd h_S(g)
          =\overline{\chi_\lambda(\operatorname{diag}z)} .
\]
Both assertions apply at every real place of
Corollary~\ref{cor:row-polar-moment} for all sufficiently large
arithmetic scales, and their degrees multiply to \(\mathcal D\).
\end{theorem}
\begin{proof}
Take the coefficient of \(z^\kappa\) in
Theorem~\ref{thm:elliptic-unnormalized-character}.
Lemma~\ref{lem:extreme-orbital-coefficient} gives \(1=d/D(b)\).
The product \(VQ\) is \((-1)^{pq'}\) times the full Vandermonde
\(\Delta(z)=\prod_{i<j}(z_i-z_j)\). Permuting the arguments
of \(\Delta\) gives the sign of that permutation. Alternating
\(A_b\) over all permutations gives \(p!q'!\) times the full
alternant of \(b\): in its monomial expansion each inner block
permutation gives the same full alternating sum after changing
the outer index, with the two inner signs cancelling. Hence
after summation the unnormalized orbital formula gives, on
multiplying by \((-1)^{pq'}d/(p!q'!)\), the full alternant
divided by \(\Delta(z)\). This is \(\chi_\lambda\) by
Theorem~\ref{thm:compact-alternant} since \(b_i=\lambda_i+n-i\).
Conjugating gives the identity for \(e\). Absolute convergence
follows from Theorem~\ref{thm:regularized-row-character}.

For the arithmetic parameters let \(n=r\). Within a consecutive
run of one sign the \(u_i\) are constant and the \(b_i\)
strictly decrease. If \(i<j\) of the same sign cross a change
of sign, at least one intervening \(b_h-b_{h+1}\) is at least
\(r^8+4r^2+1\) by the proof of
Corollary~\ref{cor:row-polar-moment}, since the \(\mu\)-gap
is larger by one. The other gaps are positive. As
\(|u_i-u_j|\le r\) and \(N-1=4r^2\), this gives
\(\kappa_i>\kappa_j\) for sufficiently large \(r\).
The remaining convergence and regularization hypotheses hold
by Theorem~\ref{thm:regularized-row-character}. Finally
\(D(b)=\mathcal D_\tau\) by Definition~\ref{def:row-vandermonde}
because \(b_i-b_j=t_i-t_j\). Take the product over the real
places.
\end{proof}

\begin{corollary}\label{cor:actual-global-elliptic-traces}
For the arithmetic row parameters, finite frames, rational groups
and elliptic finite tests in Theorem~\ref{thm:geometric-stable-reduction},
assume the first \(k-1\) row entries have a common sign at
every real place. At all sufficiently large scales, the global
geometric integral formed from the actual kernel
\(e=\prod_\tau\mathcal D_\tau\overline{\psi_\tau}\) equals the
finite test trace on the definite space \(\mathcal A^+\).
\end{corollary}
\begin{proof}
Proposition~\ref{prop:actual-real-row-projection} supplies the
kernel hypotheses, and Theorem~\ref{thm:real-row-stable-character}
supplies the required orbital identity and degrees in the
specified measures. The compact highest weights are those of
Corollary~\ref{cor:realized-compact-degree}. Apply
Theorem~\ref{thm:geometric-stable-reduction}.
\end{proof}

\section{Lattice minima on the indefinite quotient}

\begin{lemma}\label{lem:reduced-real-basis}
Let \(L\) be a full lattice in a real Euclidean space of
dimension \(m\ge1\). Write \(\lambda_j(L)\) for the least radius
of a closed ball containing \(j\) linearly independent lattice
vectors. Let \(L^*\) be its Euclidean dual, defined by integral
pairings with \(L\). There is an ordered integral basis
\(v_1,\ldots,v_m\) with \(\|v_i\|\le C_m\lambda_i(L)\) and
\[
       \lambda_{m+1-i}(L^*)\le C_m/\lambda_i(L)
                                   \qquad(1\le i\le m).
\]
If \(\lambda_1(L)\ge1/T,\ T\ge1\), the lengths of this basis
are bounded by \(C_m\operatorname{covol}(L)T^{m-1}\).
\end{lemma}
\begin{proof}
Successively choose a shortest nonzero vector in the lattice
orthogonally projected off the preceding vectors. Such vectors
are primitive, and after lifting extend the preceding vectors
to a primitive sublattice, by the basis extension and projection
argument in Lemma~\ref{lem:bounded-ideal-classes}. The lifted
vectors thus form an integral basis. In Gram--Schmidt coordinates
the resulting column matrix is \(D C\), where
\(D=\operatorname{diag}(r_i)\) consists of the positive successive
projected lengths and \(C\) is upper unitriangular. Reduce each
entry of \(C\) above the diagonal to absolute value at most
\(1/2\) by subtracting integer multiples of earlier basis
vectors, in decreasing row order in each column. This preserves
all projected choices. In the projection after \(i-1\) vectors,
the reduced lift of the next projected vector has norm squared
at most \(r_{i+1}^2+r_i^2/4\). It is nonzero so has norm
at least \(r_i\). Hence \(r_{i+1}\ge\sqrt3\,r_i/2\).
In the first \(i\) columns every entry is bounded by a constant
times \(r_i\). Conversely among \(i\) independent vectors at
least one projects nontrivially off the first \(i-1\) vectors
and hence has length at least \(r_i\). This proves
\(r_i\le\lambda_i(L)\le C_m r_i\) and the basis bound.

The Euclidean dual basis has matrix \(D^{-1}C^{-\mathsf T}\).
The entries of \(C^{-1}\) are uniformly bounded in terms of
\(m\), by summing the finite power series in the strictly upper
triangular \(I-C\). In its \(j\)-th column the dual basis matrix
has nonzero entries only in rows \(h\ge j\), bounded there
by \(C_m/r_h\). Thus its last \(m+1-i\) columns have lengths
at most \(C'_m/r_i\) by the neighboring bound on the \(r_h\).
This proves the dual assertion. Finally \(r_1=\lambda_1(L)\),
all \(r_h\ge c_m r_1\), and their product is the covolume.
Consequently every \(r_i\le C_m\operatorname{covol}(L)T^{m-1}\),
as required.
\end{proof}

\begin{proposition}\label{prop:unitary-thick-volume}
In Definition~\ref{def:moving-elliptic-lattice} let
\(\lambda_j(x)\) be the \(j\)-th minimum of \(\Lambda_x\) for
\(\|\cdot\|_x\), where \(m=2fr\). There is a constant depending
only on the fixed field and rank such that
\[
                  \lambda_i(x)\lambda_{m+1-i}(x)\le C.
\]
For fixed arithmetic data and level, the image of
\(\{\lambda_1(x)\ge1/T\}\) in \(\Gamma\backslash X\), \(T\ge1\),
is contained in a compact set. Its quotient volume is at most
\(C'(1+T)^K\) for fixed finite \(C',K\).
\end{proposition}
\begin{proof}
In real frame coordinates write \(\mathbf S\) for the direct
sum of the sign matrices. The rational symmetric pairing
\(B(v,w)=\operatorname{Tr}_{E/\Q}h(v,w)\) is
\(2\langle v,\mathbf S w\rangle_{\R}\) in those coordinates.
If \(L\) denotes the framed lattice, its \(B\)-dual is
\(\tfrac12\mathbf S L^*\). Elements of this dual are rational
vectors in the original coordinates by inverting the nonsingular
rational pairing in a lattice basis. For such a \(w\), each
\(h(v,w)\) with \(v\in\Lambda_x\) belongs to the trace dual
of \(O_E\), by testing \(\alpha v\) for all \(\alpha\in O_E\).
There is a fixed positive integer \(a\) sending this trace
dual into \(O_E\): invert the integral trace matrix and clear
its nonzero determinant as in Lemma~\ref{lem:rd}. The hermitian
\(O_E\)-dual of \(\Lambda_x\) equals \(\Lambda_x\). In fact
its local completions are the ones in the finite frame by
Lemma~\ref{lem:class-covol}, and in that frame the Gram matrix
is the identity, so the dual condition is precisely integrality
at every place. Thus \((a/2)\mathbf S L^*\subset L\).
Since \(\mathbf S\) preserves Euclidean lengths, the first
assertion follows from Lemma~\ref{lem:reduced-real-basis}.

By Lemmas~\ref{lem:moving-short-vector} and~\ref{lem:reduced-real-basis},
on the indicated set there is an integral basis of lengths at most
\(R=C_0(1+T)^{m-1}\), increasing \(C_0\) so \(R\ge1\). Choose
\(r\) of these vectors giving an \(E\)-basis \(B_0\) of the
rational space; the full list spans it already over \(\Q\).
Let \(T_0\) be their hermitian Gram matrix. It has integral
entries, with all conjugates bounded by \(R^2\). The lattice
count in Lemma~\ref{lem:cyclic-frame-volume} gives at most
\(C_E(1+R)^{4fr^2}\) choices, with \(C_E\) also depending on
\(r\). All rational bases of the same Gram matrix lie in one
\(\Gamma\)-orbit by their unique basis isometry. Fix one
representative for each Gram matrix. The argument of that
lemma bounds the number of possible finite coordinates for
it by \([K^0:K_0]|N_{E/\Q}\det T_0|^r\).
For \(M_\tau=X_\tau^{-1}(B_0)_\tau\), the column lengths
are at most \(R\), and \(|\det M_\tau|^2=|\tau\det T_0|\).
Thus \(|N_{E/\Q}\det T_0|\le R^{4fr}\).

At infinity fix one frame in the region, when it is nonempty.
After translation by its inverse every other frame has
operator norm at \(\tau\) at most \(C_r R^r/|\det M_\tau|\),
by the column bound and the adjugate formula, exactly as in
Lemma~\ref{lem:cyclic-frame-volume}. These balls are compact
since they are closed and bounded in a space of matrices and
the signed unitary equation is closed. Their product volume
is bounded by \(C(1+R)^{4fr^3}\) by
Lemma~\ref{lem:real-unitary-ball} since
\(\prod_\tau|\det M_\tau|^2=|N_{F/\Q}\det T_0|\ge1\).
We have a finite union of compact sets of volume polynomial
in \(R\) whose images cover the required set. The quotient
volume of their image is at most their original volume by
unfolding the sum of the translated indicators as in
Theorem~\ref{thm:elliptic-unfold-convergence}. Substitution of
\(R\) proves the assertion.
\end{proof}

\section{Small rational lifts at a marked split place}

\begin{lemma}\label{lem:dual-short-flag}
Fix the arithmetic data of Proposition~\ref{prop:unitary-thick-volume}
and let \(m=2fr\). For \(\lambda=\lambda_1(x)>0\) sufficiently
small there exist a nonzero proper totally isotropic \(E\)-subspace
\(U\subset V\), and integral bases \((a_i)\) of \(\Lambda_x\cap U\)
and \((v_j)\) of \(\Lambda_x\cap U^{\perp_h}\) as \(\Z\)-modules,
such that
\[
             \max_{i,j}\|a_i\|_x\|v_j\|_x
                              \le C\,\lambda^{1/(2m)} .
\]
\end{lemma}
\begin{proof}
Put \(\epsilon=\lambda^{1/(2m)}\) and use the successive minima
\(\lambda_i=\lambda_i(x)\). The last minimum is bounded below
by a positive constant because \(m\) independent vectors of
length \(\lambda_m\) span a sublattice of covolume at least
\(\operatorname{covol}(\Lambda_x)>0\); their parallelepiped
has volume at most \(\lambda_m^m\). The covolume is fixed by
Lemma~\ref{lem:moving-short-vector}. Thus for sufficiently small
\(\lambda\) there is a first \(j<m\) with
\(\lambda_j/\lambda_{j+1}<\epsilon\). Otherwise
\(\lambda_m\le\lambda\epsilon^{-(m-1)}\) would tend to zero.
We have \(\lambda_j\le\lambda\epsilon^{-(j-1)}\le\sqrt\lambda\).

The \(\Q\)-span \(U\) of the vectors of length at most
\(\lambda_j\) has dimension \(j\) by the strict gap. Every
vector of length less than \(\lambda_{j+1}\) lies in \(U\).
Multiplication by each member of a fixed integral basis of
\(O_E\) preserves the lattice and increases lengths by a fixed
factor at most the largest absolute value of its conjugates.
It therefore sends the generators of \(U\) into \(U\).
This proves that \(U\) is an \(E\)-subspace. In particular its
rank over \(E\) lies strictly between zero and \(r\).
By Proposition~\ref{prop:unitary-thick-volume},
\(\lambda_{m-j}\le C/\lambda_{j+1}\). Vectors of length at most
\(\lambda_{m-j}\) pair to zero under the integral trace form
\(B=\operatorname{Tr}_{E/\Q}h\) with all the short generators
of \(U\), since the absolute value of the pairing is at most
\(2C\lambda_j/\lambda_{j+1}<1\). The trace orthogonal complement
of \(U\) equals its hermitian orthogonal complement: pair with
every scalar multiple of a given member of \(U\) and use
nondegeneracy of the field trace as in Lemma~\ref{lem:number-field-lattice}.
These vectors span that complement by its rational dimension
\(m-j\). Likewise the generators of \(U\) pair to zero under
\(B\) with one another by \(2\lambda_j^2<1\). This makes \(U\)
totally isotropic for \(h\).
The two saturated sublattices therefore have last minima at most
\(\lambda_j\) and \(C/\lambda_{j+1}\), respectively. Apply
Lemma~\ref{lem:reduced-real-basis} in their real spans to get
the two bases and the required product bound.
\end{proof}

\begin{theorem}\label{thm:small-unipotent-lifts}
At the level \(K(g)\) of Definition~\ref{def:definite-comparison-forms}
choose one factor \(w\) of the marked group
\({\rm GL}_r(\mathbf F_p)^f\) for \(p=g_0\) or \(p=g_1\), using
the same finite frames for the indefinite space. For \(\lambda_1(x)\)
as in Lemma~\ref{lem:dual-short-flag}, let \(U_w\) be the
reduction in this frame of \(\Lambda_x\cap U\) at \(w\).
It has dimension \(\dim_E U\). For every matrix in
\[
       N_w=\{I+L:\operatorname{im}L\subset U_w,\ L|_{U_w}=0\}
\]
there is a rational unitary element \(\gamma\) with
\(c_f(x,\gamma)\in K^0\) reducing to that matrix at \(w\) and
to identity at all other marked factors. It satisfies
\[
       \max_\tau\|c_\infty(x,\gamma)_\tau-I\|_{\rm op}
                              \le C\,\lambda_1(x)^{1/(2m)}
\]
for a constant independent of \(x\) and of the matrix.
\end{theorem}
\begin{proof}
Write \(a,v\) for vectors in the respective integral bases of
Lemma~\ref{lem:dual-short-flag}. The lattices obtained by
intersecting either \(E\)-subspace with \(\Lambda_x\) are
saturated: if a nonzero integral multiple of a lattice vector
is in the subspace then the vector belongs to it. Their rational
spans are the full subspaces. A saturated subgroup of a finite
rank free abelian group extends from a basis to an integral
basis of the group: its finitely generated torsion-free quotient
is free by integer elimination as in Lemma~\ref{lem:rd}, so
lift a quotient basis. Consequently after completion at \(p\)
the intersections are direct summands over \(\Z_p\) and equal
the intersections with the completed rational subspaces. The
idempotents of the split algebra \(O_E\otimes\Z_p\) give the
same assertion in each factor, as in
Lemma~\ref{lem:split-level}. In particular their local reductions
span the full subspaces in the residue lattice, with unchanged
dimensions. The chosen factors have residue field \(\mathbf F_p\)
by Theorem~\ref{thm:data}. The conjugate place \(\bar w\) is
the partner of \(w\) in the unitary group. The hermitian pairing
identifies that partner's reduction perfectly with the dual of
the \(w\)-lattice because the frame has Gram matrix \(I\).
The reductions of the \(v\)'s at \(\bar w\) thus give all
linear functionals on the \(w\)-space vanishing on \(U_w\):
they vanish there, span the orthogonal space, and its dimension
is \(r-\dim_E U\). The reductions of the \(a\)'s span \(U_w\).
Their rank-one tensors span the maps from the quotient to
\(U_w\) by expanding in bases of the range and of the dual
of the quotient.

Choose \(\theta\in 2O_E\) equal to one modulo \(w\) and zero
at every other prime of \(E\) above \(g_0g_1\), including
\(\bar w\). All these primes are odd and distinct, so the
Chinese remainder argument of Proposition~\ref{prop:general-ideal-factorization}
applies. Write the prescribed \(L\) as a sum of the tensors
just described with coefficients in \(\mathbf F_p\); lift the
coefficients to integers \(l_{av}\) between \(0\) and \(p-1\).
In the rational space put
\[
 C y=\theta\sum_{a,v}l_{av}\,a\,h(v,y),\qquad
 T=C-C^\dagger,\qquad
                 \gamma=I+T+\tfrac12 T^2 ,
\]
where the dagger denotes adjoint for \(h\), conjugate linear in
the first entry. Its image for \(C\) is in \(U\), and
\(C|_U=0\). The adjoint has image in \(U^{\perp_h}\) and
kills \(U^{\perp_h}\), since \(C\) has image in \(U\).
In particular \(C^2=(C^\dagger)^2=C^\dagger C=0\);
hence \(T^2=-CC^\dagger\) and \(T^3=0\). Here
\(C^\dagger C=0\) uses \(U\subset U^{\perp_h}\).
As \(T^\dagger=-T\), multiplication shows that
\(\gamma^\dagger=I-T+T^2/2=\gamma^{-1}\).
The integral pairing and the choice of \(\theta\) show that
both \(C,C^\dagger\) send \(\Lambda_x\) into \(2\Lambda_x\).
Thus \(\gamma\) and its inverse preserve \(\Lambda_x\). At
\(w\), \(C\) reduces to \(L\) and \(C^\dagger\) to zero by
the conjugate scalar. At the other chosen factors both vanish.
The partner action at \(\bar w\) is the forced inverse transpose
and represents the same single unitary factor.

In each real frame, the operator norm of \(a\,h(v,\cdot)\)
is at most \(\|a\|_x\|v\|_x\) since its local sign matrix is
unitary. The same holds with the vectors exchanged for its
adjoint. The fixed scalar, the number of terms and their
coefficients are bounded in terms of the fixed field, primes
and rank. Lemma~\ref{lem:dual-short-flag} gives the required
bound on \(T\), and consequently on \(T+T^2/2\) for small
\(\lambda_1(x)\).
\end{proof}

\section{Rapid decay at a cuspidal marked place}

\begin{definition}\label{def:marked-bounded-forms}
At \(K_0=K(g)\) use \(Y=\Gamma\backslash X\) and the real
kernel \(e\) of Corollary~\ref{cor:actual-global-elliptic-traces}.
For an irreducible label \(c\) in one chosen factor \(w\) over
a marked prime, let \(E_{w,c}\) be its central projection in
the finite right action of \(\mathcal G=K^0/K(g)\), given
by Lemma~\ref{lem:central-proj}. Let \(\mathcal H_{w,c}\)
consist of the bounded continuous functions \(F:Y\to\C\)
satisfying \(R(e)F=F,\ E_{w,c}F=F\), where
\(R(e)F(x)=\int_H e(t)F(xt)\dd t\). The local character is
the cuspidal character of Theorem~\ref{thm:individual-green}.
Set \(h(x)=\max(1,\lambda_1(x)^{-1})\).
\end{definition}

\begin{lemma}\label{lem:marked-cusp-oscillation}
For \(F\in\mathcal H_{w,c}\) and \(x\) as in
Theorem~\ref{thm:small-unipotent-lifts},
\[
   |F(x)|\le\max_{u\in N_w}|F(xb_u^{-1})-F(x)| ,
\]
where \(b_u=c_\infty(x,\gamma_u)\) are the small lifts in
that theorem. For every real \(A\ge0\) there is a constant
\(C_A\), depending only on the fixed data, such that
\[
               |F(x)|\,h(x)^A\le C_A\|F\|_\infty
                                      \qquad (x\in Y).
\]
\end{lemma}
\begin{proof}
In the finite irreducible module, averaging the representation
over \(N_w\) is the orthogonal projection onto its invariants:
left and right invariance of the finite average give its
idempotence and inversion its self-adjointness. Its trace
is zero, by the proper-parabolic cuspidality in
Theorem~\ref{thm:individual-green} evaluated at the two
identity diagonal blocks in Definition~\ref{def:parabolic-operations}.
Thus this average acts by zero on the label. In every copy of
it the same holds, so the average kills \(E_{w,c}\) by the
finite decomposition in Lemma~\ref{lem:central-proj}. The
unipotent group for any conjugate subspace has the same property.
In particular \(\sum_{u\in N_w}F(xu)=0\). By
Theorem~\ref{thm:small-unipotent-lifts}, the rational translation
\(\gamma_u\) carries \(xb_u^{-1}\) to \(xu\) in \(X\):
the finite frames differ by the prescribed reduction at \(w\)
and by the identity at all other marked factors. Subtracting
the common value \(F(x)\) before averaging proves the inequality.

Theorem~\ref{thm:row-real-orthogonality} supplies a continuous
compactly supported \(\beta\) on \(H\) with \(\beta*e=e\) by
taking the product over the real places. Absolute integration
of the bounded \(F\) gives \(R(\beta)F=F\). For \(b\) in a
fixed relatively compact identity neighborhood use
\[
 F(xb^{-1})-F(x)
      =\int_H(\beta(bt)-\beta(t))F(xt)\dd t .
\]
The support of every difference lies in a fixed compact set \(C\),
and its \(L^1\)-norm tends to zero as \(b\to1\) by uniform
continuity there. For \(t\in C\) we have \(h(x)\le D h(xt)\)
with \(D\ge1\) independent of \(x,t\). Indeed the norm of
every framed lattice vector changes by at most
\(\max_\tau\|t_\tau^{-1}\|_{\rm op}\) under the new frame.
The same inequality holds for \(h_T=\min(h,T)\), \(T\ge1\).

Put \(M_T=\sup_x h_T(x)^A|F(x)|\), finite by boundedness.
The integral for the difference at each \(b_u\) multiplied by
\(h_T(x)^A\) is at most
\(D^A M_T\int|\beta(b_ut)-\beta(t)|\dd t\).
The quantitative estimate of
Theorem~\ref{thm:small-unipotent-lifts} makes this at most
\(M_T/2\) for every \(u\) once \(h(x)>H_A\), uniformly in
\(F,T\). On \(h(x)\le H_A\) the corresponding value
\(h_T(x)^A|F(x)|\) is at most \(H_A^A\|F\|_\infty\).
Taking the supremum gives \(M_T\le H_A^A\|F\|_\infty\),
since otherwise both bounds are strictly smaller than \(M_T\).
Let \(T\) increase to infinity.
\end{proof}

\begin{theorem}\label{thm:finite-marked-projected-space}
The space \(\mathcal H_{w,c}\) is finite dimensional, and every
member lies in \(L^1(Y)\cap L^2(Y)\). As a subspace of \(L^2(Y)\)
it equals the range of \(\mathcal P E_{w,c}\), where \(\mathcal P\)
is Proposition~\ref{prop:actual-real-row-projection}.
For any finite set of such marked factors and cuspidal labels,
the sum of these spaces is also finite dimensional.
\end{theorem}
\begin{proof}
On \(2^j\le h(x)<2^{j+1}\) the volume is bounded by a constant
times \(2^{jK}\) by Proposition~\ref{prop:unitary-thick-volume}.
Choose \(A>K+1\) in Lemma~\ref{lem:marked-cusp-oscillation}.
Both the integral and the squared integral of the absolute
value converge by summing the resulting geometric bounds.
Continuous functions have unique representatives in \(L^2\)
within this space: every nonempty quotient open set contains
a chart of positive measure by the construction in
Lemma~\ref{lem:quotient-measure}.

By the decay with \(A=1\), outside \(\{h\le2C_1+1\}\) each
\(F\) is bounded by \(\|F\|_\infty/2\). The latter set is
covered by the image of a compact subset of \(X\) by
Proposition~\ref{prop:unitary-thick-volume}. The integral
difference in Lemma~\ref{lem:marked-cusp-oscillation} gives
\(|F(xb)-F(x)|\le\|F\|_\infty/2\) for all \(b\) in a single
open identity neighborhood in \(H\). Since the finite coordinate
is discrete, the compact set can be covered by finitely many
right neighborhoods of this kind at points \(x_1,\ldots,x_s\).
If \(F(x_i)=0\) for all \(i\), both on this compact image
and beyond it the supremum is at most \(\|F\|_\infty/2\).
Thus evaluation at these finitely many points is injective,
which proves finite dimension.

The finite right action preserves the quotient measure and
commutes with \(\mathcal P\), since it acts on the finite
coordinate. The bounded continuous forms satisfy both
projections, hence lie in their joint Hilbert range. Conversely,
bounded square-integrable functions are dense in \(L^2(Y)\)
by truncating their moduli, using Lemma~\ref{lem:integral-tools}.
Project such a bounded function first by \(E_{w,c}\) and then
by \(\mathcal P\). The result is bounded by the finite group
sum and \(\|e\|_1\). It has a continuous representative since
it is fixed by \(R(\beta)\). More explicitly, pull a bounded
measurable representative back to \(X\) and integrate it
against \(\beta\) on the right. Translation continuity in
\(L^1\) gives the continuous representative on each real
component, and the finite coordinate is discrete. Changes on
a quotient null set pull back to null sets in each component
by unfolding with the countable rational group. These changes
give zero after convolution at every point since the archimedean
action is transitive on each component of \(X\). The projection
identities therefore hold for the continuous representative,
first almost everywhere and then everywhere. Thus a dense
subspace of the joint Hilbert range is in \(\mathcal H_{w,c}\).
This finite-dimensional subspace is closed (convergent vectors
have convergent coordinates in an orthonormal basis). Equality
follows, and a finite sum of the indicated spaces has dimension
at most the sum of the separate dimensions.
\end{proof}

\section{The trace on the marked space}

\begin{definition}\label{def:marked-spectral-projector}
Let \(\mathcal E\) be the finite group algebra projection onto
the sum of the types with at least one primitive label at a
chosen factor over \(g_0\), with no restriction on the other
factors. Explicitly it is \(1-\prod_w(1-\sum_c E_{w,c})\),
where \(w\) ranges over the chosen factors over \(g_0\) and
\(c\) over the primitive labels there. Their characters are
cuspidal by Theorem~\ref{thm:individual-green} and
Definition~\ref{def:poly-types}. Let \(\mathcal H\) be the
range of \(\mathcal P\mathcal E\) on \(L^2(Y)\). Normalize
the finite Haar measure by \(\operatorname{vol}(K(g))=1\).
The algebra of compactly supported bi-\(K(g)\)-invariant
functions acts by \(R(\phi)F(x)=\int\phi(t)F(xt)\dd t\).
It contains the finite group algebra by assigning \([s]\)
the indicator of \(sK(g)\) for \(s\in K^0/K(g)\).
\end{definition}

\begin{proposition}\label{prop:marked-poincare-kernel}
Use the arithmetic rows at sufficiently large scale and a finite
test \(\phi\) in Definition~\ref{def:marked-spectral-projector},
with \(\mathcal E*\phi=\phi*\mathcal E=\phi\). Let
\(u_1,\ldots,u_l\) be an orthonormal basis of \(\mathcal H\).
For every \(x,y\in X\) the following identity holds:
\[
 \sum_{\gamma\in\Gamma}
 e\bigl(X_\infty(x)^{-1}\gamma X_\infty(y)\bigr)
 \phi\bigl(X_f(x)^{-1}\gamma X_f(y)\bigr)
       =\sum_{j=1}^l (R(\phi)u_j)(x)\,\overline{u_j(y)} .
\]
The series on the left converges absolutely and locally uniformly.
\end{proposition}
\begin{proof}
The finite central projectors for different factors commute,
and distinct labels at one factor give orthogonal projections
by Lemma~\ref{lem:central-proj}. Thus the range of the stated
finite projection is the sum of their ranges, as can be seen
by decomposing the identity into the products of each projection
and its complement. Theorem~\ref{thm:finite-marked-projected-space}
makes \(\mathcal H\) finite dimensional, with continuous, bounded
representatives. Any finite test is a finite combination of
indicators of right level cosets, as its support is compact
and the cosets are open and disjoint. Right integration defines
a bounded operator on \(L^2(Y)\): at the full finite group it
is a weighted integral of right translations, preserving the
quotient measure by translation invariance as in
Lemma~\ref{lem:orbital-cycle}. In more detail, lift the functions
to the full finite coordinate with right \(K(g)\)-invariance,
using the product disintegration by \(K(g)\) of mass one.
Cauchy--Schwarz with weight \(|\phi|\), followed by right
translation in this full quotient, gives operator norm at most
\(\|\phi\|_1\). Right translation there preserves measure
because the finite group is unimodular: the local conjugation
Jacobian on its unitary Lie algebra has modulus one. Indeed
the trace pairing \(\operatorname{Tr}_{E_v/F_v}\operatorname{tr}(AB)\)
on skew-hermitian tangent matrices is nonsingular. After splitting
the quadratic algebra over a field extension, these matrices
are identified with arbitrary matrices in the first factor
and their negative adjoints in the second, so the pairing
is twice the ordinary nonsingular matrix trace pairing.
Nonsingularity descends by its determinant, and conjugation
preserves the pairing, hence has determinant of absolute value
one. In a local chart this gives its Haar scaling at identity
by the change of variables in Lemma~\ref{lem:unitary-cayley-density},
using the rational Cayley chart over the local field and the
additive Haar Jacobian of Lemma~\ref{lem:class-covol}. Equivalently,
on sufficiently small congruence neighborhoods the Cayley map
is a conjugation-equivariant bijection with tangent lattices
and preserves their relative indices. Almost all finite factors
preserve their integral subgroup. Left quotienting by \(\Gamma\)
preserves this translation identity by the fundamental-domain
argument in Lemma~\ref{lem:orbital-cycle}.
Associativity of \(R\) and the convolution product follows by
left translation in the inner integral. The real and finite
operators commute by separate variables. Thus
\(T=\mathcal P R(\phi)\) has
\(T=(\mathcal P\mathcal E)T=T(\mathcal P\mathcal E)\).
On this finite-dimensional range it acts by \(R(\phi)\), which
preserves continuity since the finite test acts by finitely
many right translates on level-invariant functions.

On compact subsets of \(X\times X\), finite coordinates
belong to finite sets. The finite test confines the entries
of possible rational \(\gamma\) in a fixed rational basis to
\(c^{-1}O_E\) for one positive integer \(c\), by the compact
denominator argument of Lemma~\ref{lem:moving-short-vector}.
If \(m(X_\infty(x)^{-1}\gamma X_\infty(y))\le R\), all their
archimedean conjugates have size at most \(C R\) on these
compact sets: every unitary factor has norm at least one and
the outer frame matrices and their inverses range in a compact
set. The lattice packing bound of
Lemma~\ref{lem:cyclic-frame-volume} gives \(O((1+R)^{2fr^2})\)
possible matrices. The decay at \(B=r^8>2fr^2\) by
Corollary~\ref{cor:row-polar-moment} and the parameter bounds
in Theorem~\ref{thm:elliptic-unfold-convergence} now give
absolute uniform convergence by dyadic shells. The sum is
continuous and descends in both variables by reindexing
rational matrices and using the two finite level invariances.

Call the resulting kernel \(K(x,y)\). For each bounded function
\(F\) supported on a compact part of \(Y\), unfolding on \(Y\)
gives \(TF(x)=\int_Y K(x,y)F(y)\dd y\): in the integral on
\(X\) change variables from the translating elements to their
action on \(x\). The original integral of the absolute values
is at most \(\|e\|_1\|\phi\|_1\|F\|_\infty\), which justifies
all rearrangements. The right side is continuous in \(x\):
a compact set in the quotient is covered by the images of
finitely many relatively compact local charts, so local uniform
convergence and local boundedness dominate on its support.
The same operator has the finite expansion
\(TF=\sum_j R(\phi)u_j\,\langle u_j,F\rangle\) in \(L^2\).
Its two continuous representatives therefore agree everywhere.
For each \(x\) this holds for all bounded measurable tests on
compact supports. Their integrals separate continuous functions
of \(y\): integrate the conjugate of a proposed difference
on a small relatively compact open set where it is nonzero.
The positive quotient measure there gives a contradiction.
This proves the pointwise identity.
\end{proof}

\begin{theorem}\label{thm:marked-actual-spectral-trace}
Under the hypotheses of Proposition~\ref{prop:marked-poincare-kernel},
suppose \(\phi\) is supported on finite matrices having irreducible
characteristic polynomial at some completion. For the groups and
signs of Corollary~\ref{cor:actual-global-elliptic-traces} we have
\[
           \operatorname{tr}_{\mathcal H}R(\phi)
                         =\operatorname{tr}_{\mathcal A^+}R(\phi).
\]
Here the definite action is the finite test action of
Proposition~\ref{prop:definite-elliptic-trace}.
\end{theorem}
\begin{proof}
Set \(x=y\) in Proposition~\ref{prop:marked-poincare-kernel}.
The finite sum is integrable by Cauchy--Schwarz on \(Y\);
its integral is \(\sum_j\langle u_j,R(\phi)u_j\rangle\).
On the geometric side absolute termwise integrability on the
diagonal follows from Theorem~\ref{thm:elliptic-unfold-convergence}
and the support hypothesis. This is exactly the integral of
Corollary~\ref{cor:actual-global-elliptic-traces}, which identifies
it with the definite trace.
\end{proof}

\section{Hecke tests on the actual marked spaces}

\begin{lemma}\label{lem:row-spherical-hecke}
At a split unmarked place let \(L\) be the local field, \(O\)
its discrete valuation ring with uniformizer \(\varpi\), and
\(K_L={\rm GL}_r(O)\). With \(K_L\) of measure one, the algebra
of compactly supported bi-\(K_L\)-invariant functions on
\({\rm GL}_r(L)\) is commutative and stable under
\(\phi^*(g)=\overline{\phi(g^{-1})}\).
Their actions on both \(\mathcal H\) and \(\mathcal E\mathcal A^+\)
are commuting normal operators, also across distinct such places.
\end{lemma}
\begin{proof}
By the split unitary identification of Lemma~\ref{lem:split-level},
the group and level on one member of the local pair have the
stated form. Every invertible matrix is in a double coset of
a diagonal matrix with entries powers of \(\varpi\). Multiply
by a scalar to make the matrix integral, move an entry of
minimum valuation to the first diagonal position by permutation
of rows and columns, and scale its unit factor away. Its
divisibility into every other entry permits integral elementary
row and column elimination of its column and row. The remaining
square block is nonsingular, so induction diagonalizes it.
Remove the initial common scalar. Thus transpose preserves
every double coset. It is a measure-preserving antiautomorphism:
it interchanges left and right Haar measures, which agree
by unimodularity as in Proposition~\ref{prop:marked-poincare-kernel},
and preserves the mass-one subgroup. Its pullback reverses
convolution by inversion of order in the group and a change
of the integration variable. Since it fixes every bi-invariant
function, convolution is commutative. Inversion and conjugation
preserve compact support and level. Right integration on the
full finite quotient and unitarity as in
Proposition~\ref{prop:marked-poincare-kernel} give
\(R(\phi)^*=R(\phi^*)\), including for the definite space
with its invariant norm from
Theorem~\ref{thm:definite-comparison-rank}. Hence these operators
commute with their adjoints. Distinct places act on separate
coordinates and commute under iterated integration. Operations
away from the marked primes commute with \(\mathcal E\) and
with the real projection, so preserve the two indicated spaces.
\end{proof}

\begin{theorem}\label{thm:actual-label-reduction}
Take the finite spaces \(\mathcal H\) and
\(\mathcal H^+=\mathcal E\mathcal A^+\) of
Definition~\ref{def:marked-spectral-projector}. Let
\(\mathcal B\) be any unital algebra of finite bi-\(K(g)\)
tests commuting with \(\C[\mathcal G]\), supported at the
chosen factors over \(g_0\) on the kernels of reduction.
Let \((A_i)\) in \(\mathcal B\) act normally and commute with
one another and their adjoints on the two spaces. In particular
one can use the spherical tests of
Lemma~\ref{lem:row-spherical-hecke} at split places away from
the marked level. For every joint eigensystem and every good
tuple over \(g_0\), suppose one of its primitive coordinates
has a cubic alternative occurring on neither joint eigenspace,
in the sense of Definition~\ref{def:elliptic-tests}. Then
the characteristic polynomials and traces of each \(B\in\mathcal B\)
agree on \(e_0^{\rm good}\mathcal H\) and
\(e_0^{\rm good}\mathcal A^+\). Thus the dimension \(d_0\)
of the good indefinite space satisfies
\[
 (1-b^{-8})|\mathcal G|\mathfrak m_r\mathcal D
       \le d_0\le |\mathcal G|\mathfrak m_r\mathcal D,\qquad
                      d_0 E_2\le|\mathcal G|(CJ_c)^2 .
\]
\end{theorem}
\begin{proof}
Fix an eigensystem \(\omega\) of the direct sum. Finiteness
was proved for these spaces in
Theorems~\ref{thm:finite-marked-projected-space}
and~\ref{thm:definite-comparison-rank}.
Lemma~\ref{lem:common-hecke-interpolation} gives its projection
on both spaces by a polynomial \(\Pi_\omega\) in finitely many
\(A_i\). Evaluate this polynomial by convolution in \(\mathcal B\).
For every good tuple choose a missing primitive alternative
\(P\) at its indicated coordinate \(w\). Subtract from its
tensor of label projectors the same tensor with its \(w\)-factor
replaced by \(E_{w,P}\), then sum these differences over all
the good tuples. Denote the resulting finite group algebra
element by \(T_\omega\). On both \(\omega\)-spaces it acts by
\(e_0^{\rm good}\) since the subtracted tensor annihilates
the space. By Corollary~\ref{cor:green-cubic-alternatives}
the coefficients of each difference vanish unless the matrix
at its indicated \(w\) has irreducible characteristic
polynomial. Both tensors in the difference also carry a
primitive label at \(w\). Hence
\(\mathcal E T_\omega=T_\omega\mathcal E=T_\omega\) already
in the finite group algebra, by the expression for \(\mathcal E\).

For \(B\in\mathcal B\), the test
\(\phi=B*\Pi_\omega*T_\omega\) consequently satisfies the two
projection identities of Proposition~\ref{prop:marked-poincare-kernel}.
At every point of its support the matrix at some chosen
\(g_0\)-coordinate remains integral with irreducible reduction:
the \(\mathcal B\)-factors have identity reduction in all these
coordinates, and convolution support lies in the product of
the compact supports. In the finite union of cosets contributing
to the product, use the coordinate supplied by the coefficient
of \(T_\omega\). Irreducible reduction implies irreducibility
over its local field by the monic factor argument in
Proposition~\ref{prop:definite-elliptic-trace}. Apply
Theorem~\ref{thm:marked-actual-spectral-trace} to \(\phi\).
Its definite trace is already its trace on \(\mathcal H^+\)
by the projection identities. As \(\Pi_\omega\) commutes with
the finite group algebra, on both spaces we have
\(\Pi_\omega T_\omega=\Pi_\omega e_0^{\rm good}\).
Summing over \(\omega\) gives equality of the traces of
\(B e_0^{\rm good}\). The range of \(e_0^{\rm good}\) lies
in the range of \(\mathcal E\) by
Definition~\ref{def:poly-types}. Apply the trace equality also
to the identity and to all positive powers of \(B\). The
power-sum argument of Theorem~\ref{thm:elliptic-test-transfer}
gives equality of characteristic polynomials on the two good
spaces. The dimension and energy estimates now follow from
Theorem~\ref{thm:definite-comparison-rank} and
Corollary~\ref{cor:definite-energy-capacity}.
\end{proof}

\section{Recovery of a primitive label}

\begin{lemma}\label{lem:large-root-label}
Let \(p\) be a prime, \(r\ge1\), and \(A\ge1\) with
\(p-1\nmid A\). Let \(T:\mathbf F_{p^r}^{\times}\to L^\times\)
be an injective homomorphism into a field of characteristic zero.
Write \(\mu_A(L)=\{x\in L:x^A=1\}\) and fix
\(\eta\in\mathbf F_p^\times\) with \(\eta^A=1\). For a primitive
label \(c\) of degree \(r\), with roots taken in \(\mathbf F_{p^r}\),
put
\[
            Z_c=\{T(\eta x): c(x)=0\}.
\]
The sets \(Z_c\) for different primitive \(c\) are disjoint.
Their products satisfy
\[
              \biggl(\prod_{z\in Z_c} z\biggr)^A\ne1 .
\]
Consequently, for a nonempty finite multiset \(B\) in \(L^\times\),
there is at most one primitive \(c\) satisfying both
\[
    B\subset Z_c\cup\mu_A(L),\qquad
           \frac{\prod_{\beta\in B}\beta}
                {\prod_{z\in Z_c}z}\in\mu_A(L).
\]
The first inclusion refers to individual elements; the numerator
of the quotient counts their multiplicities.
\end{lemma}
\begin{proof}
If two sets meet, equality \(T(\eta x)=T(\eta x')\) implies
\(x=x'\), whose unique monic minimal polynomial over
\(\mathbf F_p\) is the label. By Lemma~\ref{lem:primitive-scale},
the roots of \(c\) are the Frobenius conjugates of a generator
of \(\mathbf F_{p^r}^\times\), and their product generates
\(\mathbf F_p^\times\). The \(r\) images forming \(Z_c\) are
distinct by injectivity. Thus their product is \(T(\eta^r)\)
times an element of exact order \(p-1\). Raising to \(A\)
kills the first factor and preserves a nontrivial value from
the second since \(p-1\nmid A\). If \(B\) satisfies the product
condition, its product also has nontrivial \(A\)-th power.
At least one \(\beta\in B\) lies outside \(\mu_A(L)\).
For every admissible label this same \(\beta\) must belong
to \(Z_c\). Disjointness proves uniqueness.
\end{proof}

\begin{proposition}\label{prop:one-inertia-matrix-suffices}
For the actual finite spaces and commuting family in
Theorem~\ref{thm:actual-label-reduction}, take a marked
\(p=g_0\), \(A=M_\omega\) of Definition~\ref{def:poly-types}
(the level exponent, independent of a Hecke eigensystem).
Fix one of the two spaces, one joint eigensystem and one
chosen coordinate over \(p\). Suppose there are an algebraically
closed characteristic-zero field \(L\), an injective \(T\)
and a scalar \(\eta\) as in Lemma~\ref{lem:large-root-label},
and a single invertible matrix \(M\) of positive size over \(L\).
Suppose for every primitive label \(c\) occurring there one can
choose a list \(\alpha_1,\ldots,\alpha_a\in\mu_A(L)\), with
repetitions allowed, such that
\[
 \begin{split}
 &\left(\prod_{z\in Z_c}(X-z)\prod_{j=1}^a(X-\alpha_j)\right)(M)=0,\\
 &\det M=\prod_{z\in Z_c}z\prod_{j=1}^a\alpha_j .
 \end{split}
\]
Then at most one primitive label occurs there. It suffices also
to allow a different matrix for each occurring label, provided
all those matrices have the same characteristic polynomial.
If the hypothesis holds for both spaces at all the indicated
systems and coordinates, the dimension and trace conclusions of
Theorem~\ref{thm:actual-label-reduction} follow.
\end{proposition}
\begin{proof}
Let \(B\) be the multiset of roots of the common characteristic
polynomial. It is nonempty and contains no zero. For each root
\(\beta\), the matrix \(M-\beta I\) is singular; choose a nonzero
vector in its kernel. Applying the displayed annihilating
polynomial of \(M\) to this vector proves that \(\beta\) is one
of its roots. The product of the multiset is \(\det M\) by
the constant coefficient of \(\det(XI-M)\). Hence both conditions
of Lemma~\ref{lem:large-root-label} hold. At the marked primes
\(p-1>A\) by Theorem~\ref{thm:data} and its prime size bound,
so \(p-1\nmid A\). The lemma proves the assertion. Each good
tuple has a primitive member, and by Lemma~\ref{lem:primitive-scale}
it has three distinct cubic alternatives. The two spaces exclude
at most two choices by uniqueness on each, giving a common
absent alternative as required.
\end{proof}

\begin{corollary}\label{cor:primitive-determinant-invariance}
In the arithmetic ranks \(3\mid r\), for every primitive label
\(c\) and every \(u\in\mathbf F_p^\times\) with \(u^3=1\),
the products of the roots of \(c\) and \(c^{[u]}\) agree.
Thus the determinant prescription in
Proposition~\ref{prop:one-inertia-matrix-suffices} can take the
same value for all three distinct cubic alternatives with
identical auxiliary scalars.
\end{corollary}
\begin{proof}
The second list of roots, counted with multiplicity, is obtained
by multiplying the first by \(u\) by Lemma~\ref{lem:primitive-scale}.
Their product is multiplied by \(u^r=1\). The same remains true
after applying the fixed twist \(\eta\), the homomorphism \(T\)
and multiplication by the prescribed auxiliary scalars.
\end{proof}

\section{A tame element at the marked prime}

\begin{lemma}\label{lem:marked-unramified-tower}
Let \(p\) be a prime and fix the valued algebraic closure \(K_p\)
of Lemma~\ref{lem:place-fields}, with valuation \(v(p)=1\).
Every \(\Q_p\)-automorphism of \(K_p\) preserves its valuation.
For \(a\ge1\), \(L_a=\Q_p(\mu_{p^a-1})\subset K_p\) has
degree \(a\), valuation group \(\Z\), and residue field of
order \(p^a\). Reduction identifies its indicated group of
roots of unity with the whole multiplicative residue group.
The field \(U=\bigcup_{j\ge1}L_{j!}\) has value group \(\Z\)
and the same residue field as \(K_p\).
\end{lemma}
\begin{proof}
On any finite extension in \(K_p\), fix a \(\Q_p\)-basis and
use the maximum of the absolute values of the coordinates.
The given field absolute value is at most a constant times
this maximum by the triangle inequality. It is continuous
for the coordinate norm by this bound. It has positive minimum
on the coordinate unit sphere, since that sphere is compact:
a sequence in \(\Z_p^m\) has a subsequence with all coordinates
stabilizing modulo successively higher powers of \(p\) by
diagonal extraction, and completeness gives convergence.
The unit sphere is closed there. Scaling by a power of \(p\)
now bounds the coordinate norm by a constant times the field
absolute value. A \(\Q_p\)-automorphism of \(K_p\) acts by
a bounded linear map on a finite normal extension containing
any prescribed element, for instance the splitting field of
its minimal polynomial. Apply its norm bound to successive
powers of that element and take real \(k\)-th roots as
\(k\to\infty\), and then apply the inverse automorphism.
This proves valuation invariance.

Construct the field with \(p^a\) elements as in
Lemma~\ref{lem:torus-primary} and lift the monic minimal
polynomial of a generator of its multiplicative group to a
monic \(P\in\Z_p[X]\). Its residue degree is \(a\), since
the generator and zero generate the entire field. It is
irreducible over \(\Q_p\) by the monic factor argument of
Proposition~\ref{prop:definite-elliptic-trace}. In the resulting
field take the ring \(O=\Z_p[X]/P\). Its coefficient module in
the power basis is free and complete, and \(O/pO\) is the
chosen field. Every element with nonzero reduction is a unit:
lift an inverse modulo \(p\) and invert the remaining factor
\(1-pz\) by its convergent geometric series. Every nonzero
member of \(O\) is a power of \(p\) times such a unit by its
basis coefficients. Under an embedding of the field in \(K_p\)
all members of \(O\) have absolute value at most one: the
power-basis generator satisfies the monic integral polynomial
\(P\), so a value of absolute value greater than one would
give a unique largest term. Applying this also to inverses of
units gives valuation zero on the units. The field is \(O[1/p]\);
it has the asserted value group and residue field.

Each nonzero residue class lifts to a root of
\(X^{p^a-1}-1\). Its derivative there is a unit. Given a root
modulo \(p^j\), adding \(p^j t\) changes the polynomial value
modulo \(p^{j+1}\) by \(p^j t\) times that derivative. This
selects a unique digit at each step; completeness gives the
root. All \(p^a-1\) roots are obtained since they are distinct.
The lift of a multiplicative generator has powers of degrees
\(0,\ldots,a-1\) linearly independent over \(\Q_p\):
scale a putative dependence to integral coefficients with one
unit, and reduce modulo \(p\). The residue of the generator
has minimal polynomial of degree \(a\), a contradiction. Thus the
constructed field equals \(L_a\). If \(a\mid b\) then
\(p^a-1\mid p^b-1\), so \(L_a\subset L_b\).
Any integral element of \(K_p\) has all conjugates integral
by valuation invariance and the extension of embeddings from
Lemma~\ref{lem:place-fields}. The coefficients of its monic
minimal polynomial are therefore in \(\Z_p\). Its residue is
algebraic over \(\mathbf F_p\) and belongs to a finite field:
adjoining it has dimension at most that polynomial's degree.
The unique subfield of order \(p^a\) in the residue field, if
present, consists of the roots of \(X^{p^a}-X\) by
Lemma~\ref{lem:torus-primary}. The residues of the \(L_{j!}\)
therefore contain all residues, and the value group of their
union is still \(\Z\).
\end{proof}

\begin{theorem}\label{thm:marked-tame-one}
Let \(w\) be a chosen split place above \(p=g_0\), embed \(E\)
into its completion \(\Q_p\), and take its algebraic closure
\(\overline E\) inside \(K_p\). For \(r\ge1\) put
\(m=p^{r!}-1\). Identify the field with \(p^{r!}\) elements
inside the residue field, and fix any injective homomorphism
\(T\) from its multiplicative group to \(\C^\times\).
There are an automorphism \(s\) of \(\overline E/E\) preserving
the chosen valuation and acting trivially on its residue field,
and \(\xi\in\overline E\) with
\[
       \xi^m=p,\qquad
       T\bigl(\overline{s(\xi)/\xi}\bigr)=\exp(2\pi i/m).
\]
The bar here denotes reduction of a valuation unit.
For every \(h\mid m\), the ratio for \(\xi^{m/h}\) has image
\(\exp(2\pi i/h)\) under \(T\). In particular the same \(s\)
tests all the degrees dividing \(r!\).
\end{theorem}
\begin{proof}
By Lemma~\ref{lem:marked-unramified-tower} the roots of
\(X^m-1\) lie in \(U\) and reduce isomorphically to the residue
group of order \(m\). Its image under \(T\) consists of all
complex \(m\)-th roots of unity by injectivity and their number.
Choose \(\zeta\) in this group of local roots reducing to the
specified generator under \(T\). Take an \(m\)-th root \(\xi\)
of \(p\) in \(K_p\). Its valuation is \(1/m\). The powers
\(1,\xi,\ldots,\xi^{m-1}\) are independent over \(U\), since
their nonzero multiples by coefficients in \(U\) have distinct
fractional parts of their valuations. A sum with a unique
term of least valuation cannot vanish by the ultrametric
inequality. Hence \(X^m-p\) is its minimal polynomial over \(U\).
The map fixing \(U\) and sending \(\xi\) to \(\zeta\xi\) is an
automorphism of \(U(\xi)\) by evaluation at this other root,
with inverse given by \(\zeta^{-1}\). Extend to \(K_p\) by the
algebraic-closure argument of Lemma~\ref{lem:place-fields}.
It preserves the valuation by Lemma~\ref{lem:marked-unramified-tower}
and is trivial on residues because it fixes their representatives
in \(U\). The elements algebraic over \(E\) form a stable
algebraically closed subfield, and \(E\subset\Q_p\) is fixed.
Restriction gives \(s\) with the claimed inertia action.
Taking \((m/h)\)-th powers gives the smaller ratios. If \(a\mid r!\)
then \(p^a-1\mid m\), which gives the last assertion.
\end{proof}

\section{A logarithm estimate for compressed matrices}

\begin{definition}\label{def:compressed-logarithm}
Use the ordinary hermitian inner product on \(\C^n\).
Let \(\|X\|_2=(\operatorname{tr}X^*X)^{1/2}\) denote the
Hilbert--Schmidt norm in this section. If \(P\) is an orthogonal
projection, \([P,X]=PX-XP\). For \(\|X\|_{\rm op}<1\) put
\[
 L(X)=\sum_{l\ge1}\frac{(-1)^{l+1}}l X^l,\qquad
 l_P(X)=\operatorname{tr}_{\operatorname{im}P}
               L\bigl((PXP)|_{\operatorname{im}P}\bigr).
\]
A zero-dimensional trace is zero. The compressed determinant
\(\det((I+PXP)|_{\operatorname{im}P})\) is \(\exp l_P(X)\).
Indeed any complex matrix can be put in upper triangular form
by choosing an eigenline and inducting on the quotient. In
that basis both the convergent series and its exponential
have the scalar functions on the diagonal. Each scalar
\(\exp(\sum_{l\ge1}(-1)^{l+1}z^l/l)=1+z\) for \(|z|<1\)
by differentiation and evaluation at zero.
\end{definition}

\begin{lemma}\label{lem:two-crossing-logarithm}
In the orthogonal decomposition of \(P\) and \(I-P\), write
\(X=\left(\begin{smallmatrix}A&B\\ C&D\end{smallmatrix}\right)\).
If \(\|X\|_{\rm op}\le\delta\le1/2\), then
\[
  \left|l_P(X)-\operatorname{tr}(P L(X))
                     -\tfrac12\operatorname{tr}(BC)\right|
                         \le 3\delta\|B\|_2\|C\|_2 .
\]
\end{lemma}
\begin{proof}
Write \((X^l)_{ij}\) for the four blocks. Multiplying on
the left gives
\((X^l)_{11}=A(X^{l-1})_{11}+B(X^{l-1})_{21}\) and
\((X^l)_{21}=C(X^{l-1})_{11}+D(X^{l-1})_{21}\).
Iterating twice gives, for \(l\ge2\),
\[
 (X^l)_{11}-A^l
   =\sum_{i=0}^{l-2}\sum_{j=0}^{l-2-i}
            A^i B D^j C (X^{l-2-i-j})_{11}.
\]
Each term has absolute trace at most
\(\|B\|_2\|C\|_2\delta^{l-2}\).
Indeed \(|\operatorname{tr}(YZ)|\le\|Y\|_2\|Z\|_2\)
by entrywise Cauchy--Schwarz, and multiplication on either
side by a matrix of operator norm \(s\) multiplies the
Hilbert--Schmidt norm by at most \(s\), by summing squared
column lengths and then applying the same fact to adjoints.
All four compressions have operator norm at most that of
the original matrix. There are \(l(l-1)/2\) terms. At \(l=2\)
the difference is exactly \(BC\). The logarithmic series in
Definition~\ref{def:compressed-logarithm} now give the quadratic
term and a remainder bounded by \(\|B\|_2\|C\|_2\) times
\[
   \sum_{l\ge3}\frac{l-1}{2}\delta^{l-2}
       =\frac{\delta(2-\delta)}{2(1-\delta)^2}\le3\delta .
\]
The last bound follows by monotonicity of the nonnegative
series after division by \(\delta\) up to \(\delta=1/2\);
at \(\delta=0\) the claim is immediate.
\end{proof}

\begin{theorem}\label{thm:moving-compressed-log}
Let \(P,Q\) be orthogonal projections and
\(\|X\|_{\rm op}\le\delta\le1/2\). Suppose \(Y\) is a matrix
and \(c\ge0\) satisfies
\[
 \|P-Q\|_2,\ \|Y\|_2\le c,\quad
 \|[P,X]-Y\|_2\le\delta c,\quad
                   \operatorname{tr}([P,Q]Y)=0.
\]
Then
\[
    \left|l_P(X)-\operatorname{tr}(Q L(X))
                     +\tfrac14\operatorname{tr}(Y^2)\right|
                              \le12\delta c^2.
\]
\end{theorem}
\begin{proof}
Put \(J=[P,X]\), \(E=Q-P\), and use the two blocks of
Lemma~\ref{lem:two-crossing-logarithm}. Since \(J\) has off-diagonal
blocks \(B,-C\), we have
\(\operatorname{tr}(J^2)=-2\operatorname{tr}(BC)\) and
\(2\|B\|_2\|C\|_2\le\|J\|_2^2\). The lemma therefore gives
error at most \(3\delta\|J\|_2^2/2\) with \(Q=P,\ Y=J\).
Also \(\|J\|_2\le(1+\delta)c\le3c/2\). Substituting \(Y^2\)
for \(J^2\) in the trace costs at most
\(\delta c(\|J\|_2+\|Y\|_2)/4\), by cyclicity and the same
trace bound.

The diagonal blocks of \(E\) have signs \(E_{11}\le0\),
\(E_{22}\ge0\). Their absolute traces sum to
\(\operatorname{tr}P+\operatorname{tr}Q-2\operatorname{tr}(PQ)
=\|P-Q\|_2^2\), using \(P^2=P,\ Q^2=Q\).
Since \(\|L(X)\|_{\rm op}\le2\delta\) by its series,
their contribution to \(\operatorname{tr}(E L(X))\) has modulus
at most \(2\delta c^2\), by diagonalizing those hermitian
blocks. Write \(E_{\rm off}\) for the remainder.
Then \(\operatorname{tr}(E_{\rm off}X)=-\operatorname{tr}([P,Q]J)\)
by block multiplication. The trace hypothesis and
\(\|[P,Q]\|_2\le\|E\|_2\) bound this by \(\delta c^2\).
For \(l\ge2\) the derivation rule gives
\(\|[P,X^l]\|_2\le l\delta^{l-1}\|J\|_2\).
The Hilbert--Schmidt norm of the off-diagonal blocks of a
matrix equals that of its commutator with \(P\).
Thus the off-diagonal part of \(L(X)-X\) has norm at most
\(2\delta\|J\|_2\), by the geometric series. Pairing with
\(E_{\rm off}\) costs at most \(3\delta c^2\).
Adding the bounds gives at most
\((27/8+5/8+2+1+3)\delta c^2\le12\delta c^2\).
\end{proof}

\section{The moving principal minors near identity}

\begin{lemma}\label{lem:triangular-flag-shift}
Let \(g\in U(S)\), \(X=g-I\), \(\delta=\|X\|_2<10^{-3}\).
Let \(B_l,B_u\) retain respectively the strictly lower and
strictly upper entries of \(X\) whose two indices have the
same sign, setting other entries to zero. Put
\(Y=(I-B_l)g(I-B_u)=I+Z\). Then \(\|Z\|_2\le4\delta\) and
\(\psi(g)=\Psi_1(Y)\), in the notation of
Definition~\ref{def:block-splitting-integral}. If \(R\) projects
onto any initial interval of indices, including the entire
set, then
\[
                     l_R(Z)=l_R(X).
\]
\end{lemma}
\begin{proof}
Expanding the product gives
\[
 Z=X-B_l-B_u-B_lX-XB_u+B_lB_u+B_lXB_u .
\]
The Hilbert--Schmidt product bounds of
Lemma~\ref{lem:two-crossing-logarithm} give
\(\|Z\|_2\le3\delta+3\delta^2+\delta^3\le4\delta\).
All the intermediate matrices \(g,(I-B_l)g,Y\) and their
adjoints are strictly accretive because their operator
distances to identity are less than one. For any such
admissible left multiplication by a unit lower triangular
block matrix \(L\), Theorem~\ref{thm:accretive-splitting}
gives \(\Psi_1(LM)=\Psi_L(LM)=\Psi_1(M)\). In the second
equality \(m_\Lambda(Lv)=m_\Lambda(v)\) by the triangular
minor law in that theorem. Also
\(\Psi_1(M^*)=\overline{\Psi_1(M)}\), since its numerator
is \(|m_\Lambda(v)|^2\) and principal determinants conjugate
under adjoint. This proves invariance for right unit upper
triangular block multiplication whenever the arguments are
accretive, by applying left multiplication to the adjoint.
Now \(\psi(g)=\Psi_1(g)\) by Theorem~\ref{thm:accretive-splitting}.
Apply the two invariances to the stated \(Y\).

An initial principal block of \(Y\) is the product of the
corresponding three blocks, by lower triangularity on the
left and upper triangularity on the right. Its determinant
is the same as that of the initial block of \(g\). The
identity also holds on replacing \(X,B_l,B_u\) simultaneously
by \(sX,sB_l,sB_u\), \(0\le s\le1\). The norm bounds on the
whole path justify the logarithms in
Definition~\ref{def:compressed-logarithm}. Their exponentials
agree, and their difference is continuous, takes values in
\(2\pi i\Z\) and is zero at \(s=0\). This proves the
claimed equality of logarithms.
\end{proof}

\begin{theorem}\label{thm:actual-tail-minors}
In Lemma~\ref{lem:triangular-flag-shift} let \(P\) be the
coordinate projection onto a final interval and \(k\in K\).
There exists \(r_P(k)\in\C\) with
\[
 \begin{split}
 l_{kPk^*}(Z)&=\operatorname{tr}(P L(X))
                          +\tfrac14\|[P,X]\|_2^2+r_P(k),\\
 |r_P(k)|&\le600\delta
                 \bigl(\|[P,X]\|_2^2+\|[P,k]\|_2^2\bigr).
 \end{split}
\]
In particular the exponential of the first line is the actual
principal determinant of \(k^*Yk\) on that interval.
\end{theorem}
\begin{proof}
Put \(R=I-P,\ Q=kPk^*,\ A=[P,X]\), \(a=\|A\|_2\) and
\(v=\|Q-P\|_2=\|[P,k]\|_2\) by unitary invariance. Write \(A_o\)
for its opposite-sign entries and \(A_c=A-A_o\). Since \(P\)
is diagonal, \([P,B_l]\) and \([P,B_u]\) are the corresponding
masks of \(A\). Their sum is \(A_c\), the diagonal of \(A\)
being zero. In the expression for \(Z\) in
Lemma~\ref{lem:triangular-flag-shift}, each quadratic commutator
has norm at most \(2\delta a\) by the product rule, and the
cubic one at most \(3\delta^2a\). Therefore
\[
 \|[P,Z]-A_o\|_2\le7\delta a,\qquad
 \|[Q,Z]-A_o\|_2\le8\delta(a+v).
\]
The second inequality uses \(\|Z\|_{\rm op}\le4\delta\).
The trace of \([Q,P]A_o\) is zero: \(Q,P\) preserve the sign
blocks while \(A_o\) exchanges them. Use
Theorem~\ref{thm:moving-compressed-log} with first projection
\(Q\), second \(P\), comparison matrix \(A_o\), size \(4\delta\)
and \(c=2(a+v)\). It gives
\[
 l_Q(Z)=\operatorname{tr}(P L(Z))-\tfrac14\operatorname{tr}(A_o^2)
                    +e_1,\qquad |e_1|\le384\delta(a^2+v^2).
\]
Use the same theorem on \(Z\) with both projections \(R\),
comparison \(-A_o\) and \(c=2a\); its error is at most
\(192\delta a^2\). On \(X\) with both projections \(R\)
use \(-A\), size \(\delta\) and \(c=a\); its error is at most
\(12\delta a^2\). By Lemma~\ref{lem:triangular-flag-shift}
the two resulting compressed logarithms coincide and the full
logarithms have equal trace. Subtract these equations from the
full traces and insert in the formula for \(l_Q(Z)\). This
gives
\[
 l_Q(Z)=\operatorname{tr}(P L(X))
            +\tfrac14\operatorname{tr}(A^2-2A_o^2)+e_2,\qquad
 |e_2|\le588\delta(a^2+v^2).
\]

We have
\(\operatorname{tr}(A^2-2A_o^2)=\operatorname{tr}(SAS A)\),
since \(SAS=A_c-A_o\) and mixed products have zero trace.
The signed unitary identity \(g^*=Sg^{-1}S\) gives exactly
\(A^*=Sg^{-1}A g^{-1}S\), using \(P^*=P,\ SP=PS\) and
\([P,g^{-1}]=-g^{-1}[P,g]g^{-1}\).
The geometric inverse series gives
\(\|g^{-1}-I\|_{\rm op}\le2\delta\).
Hence \(\|A^*-SAS\|_2\le(4\delta+4\delta^2)a\le5\delta a\).
Taking its trace product with \(A\) costs at most \(5\delta a^2\).
After the factor \(1/4\) this is within the stated total.
Finally restriction of \(Y\) to the subspace after orthogonal
compression by \(Q\) is unitarily identified with the
compression of \(k^*Yk\) by \(P\); their determinant is
\(\exp l_Q(Z)\) by Definition~\ref{def:compressed-logarithm}.
\end{proof}

\section{Phase and quadratic decay near identity}

\begin{lemma}\label{lem:compact-gaussian-weight}
For a real weight \(\alpha\) put
\(E_\alpha(X)=\sum_{i,j}|\alpha_i-\alpha_j|\,|X_{ij}|^2\).
If \(i_1,\ldots,i_n\) are in decreasing weight order and \(P_j\)
is the projection onto the first \(j\) such coordinates, then
\[
 E_\alpha(X)=\sum_{j<n}(\alpha_{i_j}-\alpha_{i_{j+1}})
                                      \|[P_j,X]\|_2^2.
\]
For integral \(\alpha\) and ordinary unitary \(k\),
\(\displaystyle |m_\alpha(k)|^2\le\exp(-E_\alpha(k)/2)\).
\end{lemma}
\begin{proof}
The squared norm of the commutator counts \(|X_{ab}|^2\) once
when exactly one of \(a,b\) belongs to the projection. The
weights for which this happens sum to \(|\alpha_a-\alpha_b|\)
by telescoping. For a fixed \(P=P_j\) put \(B=(I-P)kP\).
On \(\operatorname{im}P\) the squared absolute determinant
of \(PkP\) is \(\det(I-B^*B)\), since \(k^*k=I\).
The eigenvalues of \(B^*B\) are between zero and one, so
\(\det(I-B^*B)\le\exp(-\operatorname{tr}B^*B)\) by
\(1-s\le e^{-s}\). Also \(\|Pk(I-P)\|_2^2=\|B\|_2^2\):
both are \(\operatorname{rank}P-\|PkP\|_2^2\) by unitarity.
Thus \(\operatorname{tr}B^*B=\|[P,k]\|_2^2/2\).
Take the powers in Definition~\ref{def:polar-row-coefficient}
and multiply, recalling \(|\det k|=1\).
\end{proof}

\begin{theorem}\label{thm:row-gaussian-phase}
Use the real coefficient of
Definition~\ref{def:polar-row-coefficient}. Let \(L_0\ge0\)
bound every \(2|t_i|\), and let \(T\) satisfy
\[
       T\ge\max\{1,n^2,n^2\log(n+L_0+1)\}.
\]
Suppose \(g\in U(S)\), \(X=g-I\), \(\delta=\|X\|_2\) and
\(\delta T\le10^{-8}\). With
\(q_X=E_\mu(X)/4\) and \(D_\Lambda=\operatorname{diag}(\Lambda_i)\),
the phase \(\exp(\operatorname{tr}(D_\Lambda L(X)))\) has
modulus one and
\[
 \left|\psi(g)\exp\bigl(-\operatorname{tr}(D_\Lambda L(X))\bigr)
                         - e^{-q_X}\right|
                      \le \tfrac1{10}e^{-q_X}+e^{-100T}.
\]
\end{theorem}
\begin{proof}
Lemma~\ref{lem:triangular-flag-shift} writes \(\psi(g)=\Psi_1(Y)\).
In \(m_{-\mu}(k^*Yk)\) the proper minors are the final intervals
after position \(j\), occurring with power \(\mu_j-\mu_{j+1}\).
The full determinant has power \(-\mu_1\) and its logarithm
is \(\operatorname{tr}L(X)\) by that lemma.
Insert Theorem~\ref{thm:actual-tail-minors} for each proper
minor. Weighted summation of the projected traces gives
\(-\operatorname{tr}(\operatorname{diag}(\mu)L(X))\):
the coefficient in position \(i\) is
\(-\mu_1+\sum_{j<i}(\mu_j-\mu_{j+1})=-\mu_i\).
A final projection and its complementary initial projection have
commutators of the same norm. Lemma~\ref{lem:compact-gaussian-weight}
therefore sums the quadratic terms to \(q_X\) and bounds the
sum of the errors by \(600\delta(4q_X+E_\mu(k))\).

For \(k\in K\), the inequality
\(E_\mu(k)\le E_\Lambda(k)+2n^2\) follows from
\(|\mu_i-\Lambda_i|\le n\) and \(\sum_{ij}|k_{ij}|^2=n\)
by unitarity. The same individual weight bound and
\(\|L(X)\|_{\rm op}\le2\delta\) show that replacing the phase
weight \(\mu\) by \(\Lambda\) costs at most \(2n^2\delta\) in
the exponent. The signed unitary equation implies
\(L(X)^*=-S L(X) S\): near identity the convergent logarithm
of the inverse is minus the logarithm by substitution of the
scalar power series, and \(g^*=Sg^{-1}S\). Each diagonal entry
is therefore purely imaginary, proving the phase assertion.
Writing \(z(k)=E_\Lambda(k)/2\), the integrand in the formula
for the phase-corrected \(\psi\) is
\[
      c_\Lambda^{-1}|m_\Lambda(k)|^2 e^{-q_X+r(k)},\qquad
      |r(k)|\le3000\delta\bigl(q_X+z(k)+n^2\bigr).
\]
This follows directly by exponentiating the sums above, so no
choice of a logarithm of the whole minor product is needed.

The compact dimension computation in
Proposition~\ref{prop:compact-formal-degree} and
Theorem~\ref{thm:compact-highest-vector} gives
\(c_\Lambda^{-1}=s_\Lambda\) for the two blocks. By the
dimension formula of Proposition~\ref{prop:standard-flag-minors}
and Lemma~\ref{lem:tableaux-binomial}, each block contributes
a product of factors
\((\Lambda_i-\Lambda_j+d_{ij})/d_{ij}\), where \(d_{ij}\ge1\)
is the distance of the two indices within the block.
Since \(\Lambda_i=t_i+(n+1)/2-i-(\mu_i-\Lambda_i)\),
every factor is at most \(L_0+4n+1\). There are at most
\(n^2/2\) factors and \(L_0+4n+1\le(n+L_0+1)^3\).
In particular \(\log s_\Lambda\le3T\).

Since \(c_\Lambda^{-1}|m_\Lambda|^2\dd k\) is a probability,
the portion of the integrated error from \(|r(k)|\le1/20\)
is at most \(e^{-q_X}/10\), using
\(|e^r-1|\le e^{|r|}-1\le1/19\) there.
On its complement put \(s=q_X+z(k)\). The bound for \(r\)
and \(\delta T\le10^{-8}\) imply
\(s+n^2>(5000/3)T\), hence \(s>1600T\).
They also give \(|r(k)|\le s/2\), using \(n^2\le T\) and
\(T\ge1\). Lemma~\ref{lem:compact-gaussian-weight} bounds
the absolute error integrand on this set by
\[
             2\exp\bigl(3T-s+|r(k)|\bigr)\le e^{-100T}.
\]
Integrate against compact probability. Adding the two bounds
proves the theorem, including \(X=0\) by the value at identity
in Theorem~\ref{thm:polar-row-decay}.
\end{proof}

\section{The common compact block of a period}

\begin{definition}\label{def:paired-phase-data}
Let \(H\) be a group, let \(S,S'\) be diagonal sign matrices of
sizes \(n,m\ge1\), and let \(j=(g,g'):H\longrightarrow U(S)\times
U(S')\) be a homomorphism. Use parameters, weights and coefficients
\((\mu,\Lambda,\psi),(\mu',\Lambda',\psi')\) as in
Definition~\ref{def:polar-row-coefficient}, with the moment
hypothesis of Theorem~\ref{thm:row-real-orthogonality} for both.
Identify a coordinate space \(V\) isometrically with a coordinate
subspace in each of the two positive sign spaces. Write
\(D=\operatorname{diag}\Lambda,\ D'=\operatorname{diag}\Lambda'\).
Assume that \(D,D'\) vanish on the two copies of \(V\), that
\(-D\) has spectrum in \([B,3B]\) on \(V^\perp\), for \(B\ge1\),
and that \(D'\ge0\). If both signs occur in \(S'\), also assume
\(D'\le 3B I\). Suppose that \(j(H)\) contains the simultaneous
copies of every \(u\in U(V)\), extended by identity on both
complements. Set \(Z_s=g(s)-I,\ Z'_s=g'(s)-I\). Require
\[
 \operatorname{tr}Z_s=\operatorname{tr}Z'_s,\qquad
 \phi(s)=\operatorname{tr}(DZ_s)-\operatorname{tr}(D'Z'_s)
                    \in i\R,\qquad \phi(st)=\phi(s)+\phi(t).
\]
Put \(\theta(s)=\psi(g(s))\,\overline{\psi'(g'(s))}
e^{-\phi(s)}\). Write \(o(g)\) for the matrix obtained from
\(g\) by retaining only entries between opposite signs. All
norms with subscript \(2\) in this section are Hilbert--Schmidt
norms.
\end{definition}

\begin{lemma}\label{lem:paired-log-phase}
In Definition~\ref{def:paired-phase-data}, if
\(\max(\|Z_s\|_{\rm op},\|Z'_s\|_{\rm op})<1\), then
\[
 \operatorname{tr}\bigl(DL(Z_s)\bigr)-
           \operatorname{tr}\bigl(D'L(Z'_s)\bigr)=\phi(s),
 \qquad L(Z)=\sum_{b\ge1}(-1)^{b-1}Z^b/b .
\]
If \(s_0\) maps to one of the simultaneous copies of \(u\in U(V)\),
then \(\theta(s_0s)=\theta(s)\) for all \(s\).
\end{lemma}
\begin{proof}
Apply the linear functional
\(\ell(A,A')=\operatorname{tr}(DA)-\operatorname{tr}(D'A')\)
to \((j(s)-I)^b\) in the direct product matrix algebra.
For \(b\ge1\) the binomial coefficients with alternating signs
sum to zero. As \(\phi(s^a)=a\phi(s)\), the result equals
\(\phi(s)\sum_{a=0}^b(-1)^{b-a}\binom ba a\).
This sum is the derivative at one of \((t-1)^b\), hence is one
for \(b=1\) and zero for \(b>1\). Summing the absolutely
convergent logarithm series proves the identity.

For \(k=u\oplus I\) on either factor the nested minors in
Definition~\ref{def:polar-row-coefficient} give \(m_\Lambda(k)=1\)
with the weight of that factor. Indeed every proper minor with
positive exponent contains either all or none of \(V\).
After including the full determinant, the total exponent
of \(\det u\) is its weight zero, by telescoping. If \(V=0\)
the assertion holds directly. For a compact \(k\) the polar
formula reduces to \(m_\Lambda(k)\) by compact coefficient
orthogonality in Proposition~\ref{prop:compact-formal-degree}
applied to its two sign blocks. Thus the unit vector \(v_0\) in
Lemma~\ref{lem:cyclic-row-isometry} satisfies
\(\langle v_0,\pi(k)v_0\rangle=1\), so
\(\|\pi(k)v_0-v_0\|^2=0\). Left multiplication by \(k\) fixes
the coefficient on the whole group. This holds in both factors.
Finally \(DZ_{s_0}=D'Z'_{s_0}=0\), so \(\phi(s_0s)=\phi(s)\).
\end{proof}

\begin{proposition}\label{prop:small-simultaneous-representative}
Under Definition~\ref{def:paired-phase-data}, suppose
\(0\le r\le1/100\) and
\(\|o(g(s))\|_2^2+\|o(g'(s))\|_2^2\le r^2\).
Some \(s_0\in H\) in the simultaneous compact block satisfies
\[
 \theta(s_0s)=\theta(s),\qquad
 \max\bigl(\|Z_{s_0s}\|_2,\|Z'_{s_0s}\|_2\bigr)\le 4r.
\]
\end{proposition}
\begin{proof}
Suppress \(s\) in the notation. Taking the \(i\)-th diagonal
entry of \((I+Z)S(I+Z)^*=S\) gives
\[
            -2\operatorname{Re} Z_{ii}
                 =\sum_j\varepsilon_i\varepsilon_j |Z_{ij}|^2 .
\]
Multiply by \(-D_{ii}\) and add the corresponding primed
identity multiplied by \(D'_{ii}\). The left sides sum to
zero by the imaginary phase condition. Moving the terms with
opposite signs shows that the sum of all squared entries
weighted by these nonnegative row weights is at most \(6Br^2\):
only opposite-sign entries occur on the right, each counted
twice with row weight at most \(3B\). In particular the bottom
rows of \(Z\) relative to \(V\oplus V^\perp\) have squared
norm at most \(6r^2\).

Write \(g=\left(\begin{smallmatrix}A&F\\ C&H_0\end{smallmatrix}\right)\)
in this splitting and \(S=I\oplus S_0\). Thus
\(\|C\|_2^2+\|H_0-I\|_2^2\le6r^2\).
From \(g^*Sg=S\) we have
\(A^*A=I-C^*S_0 C\) and \(A^*F=-C^*S_0H_0\).
If \(V\ne0\), \(A\) is invertible since \(\|C\|_2^2<1\).
Write \(A=uP\) with \(P=(A^*A)^{1/2}\) and \(u\) unitary,
as in Theorem~\ref{thm:polar-row-decay}. Multiplication by
\((I+P)^{-1}\) gives \(\|P-I\|_2\le\|C\|_2^2\), and
\[
 \|F\|_2\le\frac{1+\sqrt6\,r}{\sqrt{1-6r^2}}\|C\|_2
                         \le\tfrac{11}{10}\|C\|_2 .
\]
The last inequality follows on squaring for \(r\le1/100\).
Take \(s_0\) over the simultaneous element \(u^*\); for
\(V=0\) take identity and omit the empty corners.
For \(\widetilde g=g(s_0s)\), summing its four squared block
differences gives
\(\|\widetilde g-I\|_2^2
 \le(1+(11/10)^2+6r^2)6r^2\le14r^2\).
The trace equality at \(s_0s\) now controls its second factor.
Indeed the unweighted row identity for either matrix \(h\)
gives \(\|h-I\|_2^2=-2\operatorname{Re}\operatorname{tr}(h-I)
+2\|o(h)\|_2^2\). Multiplication by these compact blocks
preserves the norm of \(o(h)\). Thus
\(\|Z'_{s_0s}\|_2^2\le14r^2+2r^2\).
Lemma~\ref{lem:paired-log-phase} gives the coefficient equality.
\end{proof}


\section{Norm balls and decay on the period subgroup}

\begin{lemma}\label{lem:period-ball-cover}
Let \(H\) be a locally compact group with a right invariant Borel
measure \(\nu\), finite on compact sets. Let \(g_i:H\to U(S_i)\)
be continuous homomorphisms of sizes \(n_i\ge1\), \(1\le i\le l\).
Put \(d=2\sum_i n_i^2\) and
\(C_\delta=\{s:\max_i\|g_i(s)-I\|_{\rm op}\le\delta\}\).
For \(y_i\in U(S_i)\), \(R\ge1\) and \(0<\delta\le1\),
\[
 \nu\{s:\max_i\|g_i(s)y_i\|_{\rm op}\le R\}
                \le(1+2R^2/\delta)^d\,\nu(C_\delta).
\]
If \(l=1\), \(g=g_1\) is proper, \(n=n_1\), and \(c>4n^2\),
then for every \(y\in U(S_1)\)
\[
 \int_H\left(\frac2{\|g(s)y\|_{\rm op}}\right)^c\dd\nu(s)
          \le 2^c3^{2n^2}\frac{c}{c-4n^2}\,\nu(C_1)<\infty .
\]
\end{lemma}
\begin{proof}
Use the maximum of the operator norms in the real vector space
of matrix tuples of dimension \(d\). From the tuples in the
sublevel set choose a maximal family having pairwise distances
greater than \(\delta/R\). The balls of radius \(\delta/(2R)\)
around the chosen tuples have disjoint interiors and are contained
in the ball of radius \(R+\delta/(2R)\). Lebesgue volume scales
as the \(d\)-th power of the radius by scalar change of variables.
Thus any such family has size at most \((1+2R^2/\delta)^d\);
successive choices stop after finitely many steps. Choose a
preimage \(t\) of each center. Every \(s\) in the sublevel set
has a center within \(\delta/R\). For every \(i\),
\[
 g_i(st^{-1})-I=(g_i(s)y_i-g_i(t)y_i)(g_i(t)y_i)^{-1}.
\]
A signed unitary matrix \(h\) satisfies \(h^{-1}=S_i h^*S_i\),
so \(\|h^{-1}\|_{\rm op}=\|h\|_{\rm op}\). The right side has
norm at most \(\delta\), giving a cover by the right translates
\(C_\delta t\). Right invariance gives the first bound, also
when the sublevel set is empty.

In positive dimension \(\|h\|_{\rm op}\ge1\) follows by applying
the preceding inverse identity to \(hh^{-1}=I\).
For \(t\ge1\) write \(t^{-c}=\int_t^\infty cR^{-c-1}\dd R\).
The integrals on \(H\) and in \(R\) can be interchanged here.
Indeed put \(u=R^{-c}\) and approximate \(t^{-c}\) from below by
\(2^{-j}\sum_{a=1}^{2^j}{\bf1}_{\{t^{-c}\ge a2^{-j}\}}\).
These functions increase to \(t^{-c}\); the corresponding step
functions of \(u\) increase to the sublevel indicator except
at its endpoint, or equivalently give the same integral there.
Lemma~\ref{lem:integral-tools} gives the identity by passage to
the limit (at atoms the exceptional \(u\)'s are countable on
each finite sublevel set). The first bound
with \(\delta=1\) gives the result on integrating
\((3R^2)^{2n^2}cR^{-c-1}\) from one to infinity. The inverse
image \(C_1\) is compact by properness since the target ball
is closed and bounded in a finite-dimensional matrix space.
\end{proof}

\begin{lemma}\label{lem:period-offblock-decay}
In Definition~\ref{def:polar-row-coefficient} let \(h\ge0\)
and assume \(\mu_j-\mu_{j+1}\ge h\) at every change of sign.
Then for every \(g\in U(S)\)
\[
        |\psi(g)|\le (1+\|o(g)\|_2^2)^{-h/4}.
\]
\end{lemma}
\begin{proof}
Use \(g=kP\) from Theorem~\ref{thm:polar-row-decay}.
Let \(a=g-o(g)\) and \(O=o(g)\). The diagonal sign blocks of
\(g^*Sg=S\), after multiplication by their signs, give
\(a^*a=I+O^*O\). The determinant on the right is the product
of \(1+s_j\) over the nonnegative eigenvalues of \(O^*O\).
It is at least \(1+\sum_j s_j=1+\|O\|_2^2\).
Since \(k\) preserves the two sign subspaces,
\(|\det a|\) is the product of the two sign-block determinants
of \(P\). These determinants coincide by
Lemma~\ref{lem:sign-run-determinant}; denote their common value
by \(d_+\). Hence \(d_+^4\ge 1+\|O\|_2^2\).
The same lemma bounds each denominator in the polar integral
below by \(d_+^h\), unchanged under conjugation of \(P\) by \(K\).
Cauchy--Schwarz bounds the integral of the absolute numerator
by \(c_\Lambda\) as in Theorem~\ref{thm:polar-row-decay}.
This proves the assertion including the one-sign case.
\end{proof}


\section{A lower bound for a real period}

\begin{theorem}\label{thm:paired-real-period}
Give the group \(H\) of Definition~\ref{def:paired-phase-data} a
locally compact topology and a right invariant Borel measure
\(\nu\), finite on compact sets and positive on nonempty open
sets. Assume \(g,g'\) are continuous and \(g\) is proper.
Let \(\beta\ge10^{12}\) satisfy
\[
 n+m+\max\{\max_i2|t_i|,\max_j2|t'_j|\}\le\beta,\qquad
 \|g'(s)\|_{\rm op}\le\beta+1\quad
           \text{when }\|g(s)-I\|_{\rm op}\le1 .
\]
Put \(N=(n+m)^2,\ T=N\log\beta\), and suppose that the
successive \(\mu\)-gaps at sign changes in each factor are at
least \(h=10^{22}T^3\). With
\(C=\nu\{s:\|g(s)-I\|_{\rm op}\le1\}\), one has
\[
         \operatorname{Re}\int_H\theta(s)\dd\nu(s)
                         \ \ge\ \beta^{-15N} C>0.
\]
If \(g'(H)\) is ordinary unitary, the exponent \(15N\) can
be replaced by \(11N\).
\end{theorem}
\begin{proof}
The first-factor decay of Theorem~\ref{thm:polar-row-decay}
and \(|\psi'|\le1\) give, for \(0\le a\le h\),
\[
       |\theta(s)|\le\min\{1,(2/\|g(s)\|_{\rm op})^a\}.
\]
For \(\|g(s)\|_{\rm op}\ge2\) this follows by decreasing the
decay exponent to \(a\), and otherwise the bound by one suffices.
Lemma~\ref{lem:period-ball-cover} with \(a=8N>4n^2\) gives
integrability of the continuous function \(\theta\).
It also gives \(C<\infty\), and the interior of this set
contains identity, so \(C>0\).

Write \(r=10^{-9}/T\) and \(e_0=\exp(-100T)\). If
\(\|o(g(s))\|_2^2+\|o(g'(s))\|_2^2\le r^2\), use the
representative \(s_0s\) of
Proposition~\ref{prop:small-simultaneous-representative}.
Each matrix difference then has Hilbert--Schmidt norm at most
\(4r\), so its product with \(T\) is at most \(10^{-8}\).
The size of \(T\) meets the lower bounds in
Theorem~\ref{thm:row-gaussian-phase} for both factors. By that
theorem and Lemma~\ref{lem:paired-log-phase}, \(\theta(s)\) is
the product of a number within \(A/10+e_0\) of \(A\) and the
conjugate of a number within \(A'/10+e_0\) of \(A'\), where
\(A,A'\in[0,1]\) are the corresponding positive Gaussians.
Here the two logarithmic phases cancel \(e^{-\phi(s_0s)}\) exactly.
Expansion bounds \(|\theta(s)-AA'|\) by
\(0.21AA'+4e_0\). In particular
\(\operatorname{Re}\theta(s)\ge|\theta(s)|/2-6e_0\).
On the complementary region Lemma~\ref{lem:period-offblock-decay}
in the two factors gives
\[
  |\theta(s)|\le(1+r^2)^{-h/4}\le\exp(-hr^2/8)\le e_0 .
\]
We used \(\log(1+r^2)\ge r^2/2\), obtained by integration
of \(1/(1+t)\) for \(0\le t\le r^2\le1\).

The positive part of \(|\theta|/2-\operatorname{Re}\theta\)
is thus bounded by \(6e_0\) everywhere and also by
\(3|\theta|/2\). Taking the geometric mean and using the
decay estimate with \(a=16N\) bounds it by
\[
       3e^{-50T}(2/\|g(s)\|_{\rm op})^{8N}.
\]
The last inequality of Lemma~\ref{lem:period-ball-cover}
bounds its integral by \(3\beta^{-50N}16^{8N}C\).
Indeed \(8N/(8N-4n^2)\le2\) and
\(2^{8N}3^{2n^2}\cdot2\le16^{8N}\).

Let \(U_0\) consist of the elements for which both operator
norms of the differences from identity are at most
\(\delta=\beta^{-4}\). Their Hilbert--Schmidt norms are at
most \(\sqrt\beta\,\delta\), since the squared norm is the
sum of the squared singular values. Multiplication by
\(T\le\beta^2\log\beta\) therefore gives at most \(10^{-8}\);
indeed \(\beta^{-3/2}\log\beta\) is decreasing on this range.
Each weight difference \(|\mu_i-\mu_j|\) and its primed
analogue is at most \(2\beta\) by
Definition~\ref{def:polar-row-coefficient}. Thus the two
Gaussian exponents \(E_\mu(Z_s)/4,E_{\mu'}(Z'_s)/4\) are
each at most \(\beta^{-6}\). By
Theorem~\ref{thm:row-gaussian-phase} each coefficient has
modulus at least \(0.9\exp(-\beta^{-6})-e_0>3/4\).
It follows that \(|\theta|\ge1/2\) on \(U_0\).

Apply Lemma~\ref{lem:period-ball-cover} to the pair with
\(R=\beta+1\), both translating matrices equal to identity.
The set defining \(C\) lies in this sublevel set. Its covering
factor is at most \(\beta^{14N}\), since
\(1+2(\beta+1)^2\beta^4\le\beta^7\) and \(d\le2N\).
If the second image is ordinary unitary, take \(R=2\)
and use \(1+8\beta^4\le\beta^5\), giving \(\beta^{10N}\).
Write \(u=14\) or \(10\), respectively. Then
\[
 \operatorname{Re}\int_H\theta\,\dd\nu
    \ge \tfrac14\beta^{-uN}C-3\beta^{-50N}16^{8N}C
    \ge \beta^{-(u+1)N}C .
\]
For the final inequality \(16^8<10^{10}\le\beta\),
\(3\beta^{-49N}\le \beta^{-uN}/8\), and
\(\beta^N\ge8\). This proves both bounds.
\end{proof}


\section{A signed interpolation basis at the Bessel point}

\begin{lemma}\label{lem:signed-bessel-interpolation}
Fix \(\nu\ge0\), put \(l=2\nu+2\), and choose decreasing integers
\(0=x_0>x_1>\cdots>x_l\). On \(\C^l\), numbered \(1,\ldots,l\),
let \(S=\operatorname{diag}((-1)^{i-1})\) and
\([v,w]=v^*Sw\). Put
\[
 c_i^2=\prod_{\substack{0\le j\le l\\j\ne i}}|x_i-x_j|,
 \quad c_i>0,\qquad
 P(z)=\prod_{i=1}^l(1-x_i z)^{1/2}=\sum_{a\ge0}p_a z^a ,
\]
taking the square roots equal to one at zero. Define real vectors
\(v_0,v_1,\ldots,v_{l-1}\) by
\[
 (v_0)_i=c_0/c_i,\qquad
 (v_j)_i=c_i^{-1}\sum_{a=0}^{j-1}p_a x_i^{\,j-a}
                       \quad(1\le j<l).
\]
Then \([v_0,v_0]=1,\ [v_0,v_j]=0\) for \(j\ge1\), and
\([v_j,v_k]=-{\bf1}_{j+k=l}\) for \(1\le j,k<l\).
For \(T_x=\operatorname{diag}(x_i)\) one has
\[
 T_xv_0=c_0v_1,\qquad T_xv_j=v_{j+1}-p_jv_1\quad(1\le j<l-1).
\]
\end{lemma}
\begin{proof}
Let \(F(t)=\prod_{i=0}^l(t-x_i)\). Multiplication by \(F(t)\)
shows that \(\sum_i(F'(x_i)(t-x_i))^{-1}=1/F(t)\):
the difference of the numerators has degree at most \(l\) and
vanishes at all \(l+1\) nodes. A nonzero polynomial of degree
at most \(l\) has at most \(l\) roots by successive division by
its linear factors. Comparing the coefficients of \(t^{-d-1}\)
at infinity gives for \(d\ge0\), with \(0^0=1\),
\[
 \sum_{i=0}^l\frac{x_i^d}{F'(x_i)}
      =[z^{d-l}]\prod_{i=1}^l(1-x_i z)^{-1}.
\]
A coefficient of negative index is zero. Since
\(F'(x_i)=(-1)^i c_i^2\), the signed moments from indices
\(i\ge1\) are \(0^d/c_0^2\) minus the coefficient on the right.
This proves the assertions involving \(v_0\), since \(p_0=1\)
and \(1\le j<l\). For \(j,k\ge1\) apply the formula to
each monomial of degree \(j+k-a-b\) in their inner product.
All those degrees are positive. If \(a\ge j\) or \(b\ge k\)
and the other index is nonnegative, then
\(j+k-a-b-l<0\), since \(j,k<l\).
Consequently the sum for \(-[v_j,v_k]\) equals the unrestricted
coefficient of degree \(j+k-l\) in
\[
             P(z)^2\prod_{i=1}^l(1-x_i z)^{-1}=1.
\]
The identities for \(T_x\) follow by multiplying each entry by
\(x_i\) and isolating the last summand of \(v_{j+1}\).
\end{proof}

\begin{proposition}\label{prop:bounded-bessel-isometry}
Let \(X\) be a coordinate space of dimension \(k-1\), \(k\ge1\),
with standard positive form. Extend the space of
Lemma~\ref{lem:signed-bessel-interpolation} by \(X\), and
write \(S_1=I_X\oplus S\). In the target use the complex
space of Definition~\ref{def:row-space} in real-place coordinates
with \([z,z]=-1\). Its middle positive space is
\(X\oplus\C u\) with unit \(u\). Denote its Gram matrix by
\(J\), with hyperbolic pairs \(e_i,e_i^*\), \(1\le i\le\nu\).
Let \(f_i=e_i\) for \(1\le i\le\nu\) and \(f_{\nu+1}=z\).
There exists a specified isometry \(Q\) from \(S_1\) to \(J\)
which is identity on \(X\) and satisfies
\[
 Qv_0=u,\qquad Qv_j=f_j\ (1\le j\le\nu+1),\qquad
 Qv_{l-j}=-e_j^*\ (1\le j\le\nu).
\]
For \(b\ge8(\nu+1)\) and \(K_0=\max_i |x_i|\le 2b^2\),
\[
             1+\max(\|Q\|_{\rm op},\|Q^{-1}\|_{\rm op})
                               \ \le\ b^{5(\nu+1)} .
\]
All operator norms here use the standard positive Euclidean
norm on each coordinate space.
\end{proposition}
\begin{proof}
The Gram matrix computed in Lemma~\ref{lem:signed-bessel-interpolation}
is invertible, so the \(l\) prescribed vectors form a basis.
The target vectors have the same Gram matrix, including the
minus sign on each dual and the negative square of \(z\).
Consequently the stated prescriptions uniquely define \(Q\)
and \(Q^*JQ=S_1\).

Each \(c_i\ge1\) by integrality and distinctness, and
\(c_0\le K_0^{l/2}\). On \(|z|=1/(2K_0)\) the product
\(P(z)\) has modulus at most \((3/2)^{l/2}\). Indeed each
factor has modulus at most \((3/2)^{1/2}\) and its Taylor
series there is obtained from the branch taking value one at
zero. This branch exists on \(|x_i z|<1\), for instance by
exponentiating half the convergent logarithm series of
\(1-x_i z\). For \(j\ge1\), \(c_i(v_j)_i\) is the
coefficient of \(z^{j-1}\) in \(x_iP(z)/(1-x_i z)\).
The coefficient estimate in Lemma~\ref{lem:analytic-tools}
bounds \(|(v_j)_i|\) by \((2K_0)^j(3/2)^{\nu+1}\).
The same uniform upper bound at \(j=l-1\) covers \(v_0\).
The norm of the \(l\)-dimensional block of \(Q^{-1}\)
is at most \(l\) times the maximum absolute entry by
Cauchy--Schwarz on the rows. Thus that norm is at most
\[
   l(4b^2)^{2\nu+1}(3/2)^{\nu+1}
                     =12(\nu+1)24^\nu b^{4\nu+2}.
\]
For \(\nu\ge2\) we have \(24\le b\); for \(\nu=0,1\) one
has \(24^\nu\le(3/2)b^\nu\). In all cases the last expression
is at most \((9/4)b^{5\nu+3}\le b^{5(\nu+1)}-1\),
using \(b\ge8\). The identity block on \(X\) obeys the same
bound. Finally \(Q^{-1}=S_1 Q^*J\) and \(S_1,J\) are ordinary
unitary, so the two norms in the statement are equal.
\end{proof}


\section{The character on the real Bessel subgroup}

\begin{definition}\label{def:real-bessel-graph}
Use the matrix \(J\), vectors \(f_j\) and positive middle space
\(P_0=X\oplus\C u\) of Proposition~\ref{prop:bounded-bessel-isometry}.
Use the successive blocks
\[
 e_1,\ldots,e_\nu,\quad P_0\oplus\C z,\quad
                              e_\nu^*,\ldots,e_1^* .
\]
Let \(\mathcal U\subset U(J)\) consist of the upper block
triangular elements with identity diagonal blocks, the real
group of Definition~\ref{def:unip} in the scaled basis.
Embed \(K_0=U(P_0)\) in \(U(J)\) by identity on all the other
coordinates, writing \(\iota(k)\) for the image. Set
\(H=\{(\iota(k)n,k): n\in\mathcal U,\ k\in K_0\}\). For a
matrix \(y\in U(J)\) put
\(\sigma(y)=\sum_{j=1}^{\nu} f_j^*y f_{j+1}\), where a star
on a column in a matrix product denotes its ordinary conjugate
transpose. Use product measure on \(H\) in the coordinates
\((n,k)\), with \(\dd n\) as in
Proposition~\ref{prop:unipotent-covol} at this place and compact
dual measure \(\dd k\) on \(K_0\) from
Theorem~\ref{thm:compact-dual-unitary-volume}.
\end{definition}

\begin{lemma}\label{lem:real-graph-character}
The group \(H\) is closed, its first projection is proper,
and its measure is right invariant, finite on compact sets
and positive on every nonempty open set. On its first
projection \(\sigma\) is additive under multiplication.
For \(A=\sum_{j=1}^\nu f_{j+1}f_j^*\) and
\(C=A+JA^*J\), every \((y,k)\in H\) satisfies
\[
          \operatorname{tr} C(y-I)=\sigma(y)-\overline{\sigma(y)},
       \qquad \operatorname{tr}(y-I)=\operatorname{tr}(k-I).
\]
\end{lemma}
\begin{proof}
The embedding of the compact group normalizes \(\mathcal U\)
by block diagonal conjugation. Hence \((n,k)\mapsto(\iota(k)n,k)\)
is a continuous bijection onto a subgroup with continuous inverse.
Its first projection is the closed matrix group specified by
upper block triangularity, identity on the one-dimensional
diagonal end blocks and middle diagonal block \(k\oplus1_z\),
\(k\in K_0\). The second coordinate is read continuously from
that middle block. Thus the first projection is a homeomorphism
onto a closed subgroup, hence is proper. Local finiteness and
positivity follow from the smooth positive factors of the
product measure. Right translation by \((\iota(k_0)n_0,k_0)\)
sends \((n,k)\) to
\((\iota(k_0)^{-1}n\iota(k_0)n_0,kk_0)\). Conjugation here
preserves the unipotent measure, as each matrix row coordinate
of Lemma~\ref{lem:heisenberg-layer} changes by a unitary on the
middle positive coordinates and its central imaginary coordinate
is fixed. The factor measures are right invariant by
Proposition~\ref{prop:unipotent-covol} and
Theorem~\ref{thm:compact-dual-unitary-volume}.

For \(y_1,y_2\) in the first projection, an entry in row
\(e_j\) and the next flag column \(f_{j+1}\) of \(y_1y_2\)
is the sum of their two corresponding entries. Only the
diagonal at \(e_j\) and the entry in column \(f_{j+1}\)
can contribute. At \(f_{\nu+1}=z\) this uses that each middle
block fixes \(z\). Thus \(\sigma(y^{-1})=-\sigma(y)\), also for
\(\nu=0\) by the empty sum. Now \(\operatorname{tr}(Ay)=\sigma(y)\).
Since \(y^{-1}=Jy^*J\), cyclicity of the trace gives
\(\operatorname{tr}(JA^*Jy)=\overline{\operatorname{tr}(Ay^{-1})}\).
Subtract the identity contribution, which is zero. The second
trace formula follows by adding the traces of the diagonal blocks.
\end{proof}

\begin{theorem}\label{thm:bessel-graph-phase}
Suppose a weight \(D\) of Definition~\ref{def:polar-row-coefficient}
on the \(S_1\)-space vanishes on \(X\). Assume its remaining
diagonal entries \(\lambda_1>\cdots>\lambda_l\), in the
alternating signs of Lemma~\ref{lem:signed-bessel-interpolation},
belong to \([-3B,-B]\) with \(B\ge1\).
Fix an integer \(a_0\ge0\) and on \(P_0\) take
\(D'=a_0 uu^*\). Choose \(Q\) of
Proposition~\ref{prop:bounded-bessel-isometry} with
\(x_i=\lambda_i-a_0\). Define
\[
       j(y,k)=(Q^{-1}yQ,k),\qquad
       \phi(y,k)=\sigma(y)-\overline{\sigma(y)} .
\]
Then the trace, weight and compact block conditions of
Definition~\ref{def:paired-phase-data} hold. If the parameter and
gap bounds of Theorem~\ref{thm:paired-real-period} hold for these
weights and \(\beta\), its lower bound with exponent
\(11(r+k)^2\), \(r=k+2\nu+1\), holds for this actual graph.
\end{theorem}
\begin{proof}
The matrices \(Q^{-1}yQ\) preserve \(S_1\). Let \(D_0=QDQ^{-1}\),
which satisfies \(D_0^*J=JD_0\). Since \(D\) restricts to
\(a_0 I+T_x\) on the alternating block, the transitions in
Lemma~\ref{lem:signed-bessel-interpolation} give, for \(\nu\ge1\),
\[
    D_0e_j=a_0 e_j+f_{j+1}-p_j f_1\quad(1\le j\le\nu).
\]
For every \(\nu\) one also has \(D_0u=a_0u+c_0 f_1\), and
\(D_0X=0\); when \(\nu=0\), \(f_1=z\). In
Lemma~\ref{lem:real-graph-character}, \(C\) is \(J\)-self-adjoint
and \(Ce_j=f_{j+1}\), since \(Je_j=e_j^*\) is perpendicular
in Euclidean coordinates to all the frame vectors. Hence
\(D_0-C\) preserves every initial isotropic flag span and,
by \(J\)-self-adjointness, its \(J\)-orthogonal complement.
These are precisely the successive proper partial sums of the
ordered blocks in Definition~\ref{def:real-bessel-graph}.
Thus \(D_0-C\) is upper block triangular, also for \(\nu=0\)
where there is a single block. The product with
\(y-\iota(k)\) is strictly upper block triangular, with trace zero.

The ordinary compression of \(D_0-C\) to \(P_0\) equals \(D'\).
Indeed \(C\) vanishes on \(P_0\), and \(f_1\) has zero projection
there. As \(\iota(k)-I\) is supported on \(P_0\), the preceding
trace calculation and Lemma~\ref{lem:real-graph-character} give
\[
 \operatorname{tr}\bigl(D(Q^{-1}yQ-I)\bigr)
     =\operatorname{tr}(D'(k-I))+\sigma(y)-\overline{\sigma(y)}.
\]
The unweighted traces agree by the same lemma. This phase is
imaginary and additive. An ordinary unitary on \(X\), extended
by identity, commutes with the change of basis on its complement
and thus occurs in both coordinates of \(j\) at \(n=I\).
The other weight bounds follow from the assumptions and
\(D'\ge0\) on the one-sign second factor.
Continuity and properness follow from
Lemma~\ref{lem:real-graph-character} and conjugation by \(Q\).
The second matrix always has operator norm one, so all group
and measure hypotheses of Theorem~\ref{thm:paired-real-period}
are satisfied.
\end{proof}


\section{A Bessel function at each real place}

\begin{proposition}\label{prop:actual-bessel-weights}
Use the sufficiently large scales of Proposition~\ref{prop:thick},
a place \(\tau\in\Phi\) and the row of
Definition~\ref{def:row-vandermonde}. Write \(a_\tau=a\) for
\(\tau\in A\) and \(a_\tau=0\) otherwise. Use the positive signs
on the first \(k-1\) entries and alternating signs beginning with
positive on the remaining \(m=2\nu+2\). Put
\(\beta=b^3\). For this row and the compact \(U(k)\)-weight
\(D'=a_\tau uu^*\), the hypotheses of
Theorem~\ref{thm:bessel-graph-phase} and all bounds of
Theorem~\ref{thm:paired-real-period} hold, uniformly on the
interval \(t\le v\le t^2\) of Proposition~\ref{prop:level-prime}.
Moreover \(3B+a_\tau\le2b^2\) with \(B=W_\tau\).
\end{proposition}
\begin{proof}
Number the first \(r\) coordinates starting at zero in this
calculation. For \(0\le i<k-1\) there are \(i\) previous
positive signs out of \(k+\nu\) positive signs in total.
In Definition~\ref{def:polar-row-coefficient}
the compact correction subtracted from \(\mu_i\) is
\(k+\nu-1-2i\); this equals
\((k/2-1-i)+(r-1)/2-i=\mu_i\) by \(r=k+2\nu+1\).
Hence \(\Lambda_i=0\). For all other indices
\(|\Lambda_i-t_i|\le 3r/2\) by the same formula.
Their consecutive \(t_i\) decrease by \(2r^8\), so the
corresponding \(\Lambda_i\) decrease strictly. They lie
between \(-W_\tau-m\,2r^8-3r/2\) and
\(-W_\tau-2r^8+3r/2\).
By the bounds in Proposition~\ref{prop:exact-row-degree},
\(W_\tau\ge w(\tau)\ge xb\exp(-t^2)\), which exceeds every fixed
power of \(x\) eventually since \(t^2=o(\log b)\).
As \(m r^8=O(x^{41})\), these entries lie in \([-3B,-B]\).

Proposition~\ref{prop:level-prime} and
Definition~\ref{def:row-vandermonde} give \(B\le2b+N_1\),
\(a_\tau\le a\le 2xb\) and \(N_1=O_M(x^{11}(1+\log x))\).
Thus \(3B+a_\tau\le2b^2\) and \(8(\nu+1)\le b\) eventually.
In the compact factor take \(t'_i=\Lambda'_i+(k+1)/2-i\)
in decreasing order of \(\Lambda'\), numbered from one.
These entries are distinct and yield \(D'\), since with one
sign the compact correction is \(k+1-2i\).
The bounds above imply
\(r+k+\max(2\max |t_i|,2\max |t'_i|)\le b^3\).
There are no sign changes in the compact factor. In the
first factor every sign-change \(\mu\)-gap is at least \(r^8\)
eventually by Corollary~\ref{cor:row-polar-moment}. Furthermore
\[
 \frac{r^7}{((r+k)^2\log(b^3))^3}
       \ \ge\ \frac{r}{12^3 x^{3\epsilon}}\ \longrightarrow\ \infty
\]
by Proposition~\ref{prop:thick}. These facts give the required
gap and moment bounds in both factors; the second group is compact.
\end{proof}

\begin{theorem}\label{thm:actual-real-bessel-vector}
With the data of Proposition~\ref{prop:actual-bessel-weights},
let \(G_1=U(J)\), and choose an ordinary unitary isometry \(p\)
from \(S_1\) to \(J\). Put \(K(y)=\psi(p^{-1}yp)\),
\(e_J(y)=d\,\overline{K(y)}\) with \(d=\mathcal D_\tau\)
from Theorem~\ref{thm:real-row-stable-character}. Give
\(G_1\) the measure transported from \(h_{S_1}\) by \(p\).
Write \(\pi=\operatorname{Sym}^{a_\tau}P_0\) for the unitary
representation on symmetric tensors, also writing \(\pi(k)\)
for its action. There exist \(x\in G_1\) and a bounded
continuous function \(W:G_1\to\pi\) with
\(1+\|x\|_{\rm op}\le b^{5(\nu+1)}\) and \(W(x)\ne0\).
For every \((s,k)\in H\) and \(y\in G_1\) one has
\[
 \begin{split}
 W(sy)&=\exp(\sigma(s)-\overline{\sigma(s)})\,\pi(k)W(y),\\
 W(y)&=\int_{G_1}e_J(a)W(ya)\dd a,\qquad
 \|W(y)\|\le\int_H|K(x^{-1}s y)|\dd\nu_H(s,k).
 \end{split}
\]
Here \(\nu_H\) is the graph measure of
Definition~\ref{def:real-bessel-graph}. With \(c=8r^2+1\)
there is also the uniform bound
\[
 \int_H(2/\|x^{-1}s y\|_{\rm op})^c\dd\nu_H(s,k)
            \le16^c b^{33(r+k)^2}\|W(x)\|.
\]
\end{theorem}
\begin{proof}
Such a \(p\) is obtained by taking the orthonormal sums and
differences of each pair \(e_j,e_j^*\), together with the
middle coordinates and a permutation matching the signs.
Put \(x=Qp^{-1}\) with \(Q\) from
Theorem~\ref{thm:bessel-graph-phase}, giving the stated size
by Proposition~\ref{prop:bounded-bessel-isometry}. Equip the
tensor power of \(P_0\) with its product inner product and
take the subspace fixed by permutations. Its compact action
is unitary. The unit vector \(v_0=u^{\otimes a_\tau}\) has
coefficient \(\langle v_0,\pi(k)v_0\rangle=(u^*ku)^{a_\tau}\),
where tensor power zero is the trivial line. This equals
\(m_{\Lambda'}(k)=\psi'(k)\) by the minor formula of
Definition~\ref{def:polar-row-coefficient} and compact
orthogonality as in Lemma~\ref{lem:paired-log-phase}. Its
conjugate is the coefficient at \(k^{-1}\).
This compact action pairs with the conjugate symmetric power
specified in Definition~\ref{def:columns}: the ordinary tensor
inner product gives a linear pairing with the conjugate space,
invariant under the two actions.

For \(s_H=(s,k)\), abbreviate \(\phi(s_H)=\sigma(s)-\overline{\sigma(s)}\).
Define
\[
 W(y)=\int_H K(x^{-1}s y)e^{-\phi(s_H)}
                               \pi(k^{-1})v_0\,\dd\nu_H(s_H).
\]
The integrand has norm \(|K(x^{-1}s y)|\). Put
\(g(s_H)=Q^{-1}sQ\), \(Y=p^{-1}x^{-1}yp\) and
\(C_1=\nu_H\{\|g(s_H)-I\|_{\rm op}\le1\}\).
The norm of \(g(s_H)Y\) is \(\|x^{-1}s y\|_{\rm op}\).
By Theorem~\ref{thm:polar-row-decay} the integrand is bounded
by \((2/\|g(s_H)Y\|_{\rm op})^c\), since \(c\le r^8\)
eventually. Lemma~\ref{lem:period-ball-cover} bounds the
integral of this majorant by \(16^c C_1\), uniformly in \(y\),
as \(c>8r^2\). On every compact set of \(y\) another integrable
majorant is \((2/\|g(s_H)\|_{\rm op})^c\) times
\(\sup\|Y^{-1}\|_{\rm op}^c\), by submultiplicativity. Thus
\(W\) is bounded and continuous by Lemma~\ref{lem:integral-tools}.

At \(y=x\), pairing with \(v_0\) gives exactly
\(\int_H\theta\,\dd\nu_H\) for the graph in
Theorem~\ref{thm:bessel-graph-phase}. Its real part is at least
\(b^{-33(r+k)^2}C_1>0\) by
Theorem~\ref{thm:paired-real-period}. Cauchy--Schwarz gives the
asserted bound on the majorant in terms of \(\|W(x)\|\).
Right translation in \(H\) in the defining integral, using
Lemma~\ref{lem:real-graph-character}, proves the first
transformation formula, including its phase sign.
For the reproducing formula, substitute the definition of
\(W(ya)\) in the integral over \(G_1\). The iterated absolute
integral is at most \(16^c C_1\|e_J\|_1<\infty\) by the
uniform bound already proved and Corollary~\ref{cor:row-polar-moment}.
Interchange the integrals and apply
Theorem~\ref{thm:row-real-orthogonality} in the form
\(\int e_J(a)K(z a)\dd a=K(z)\) to \(z=x^{-1}s y\).
This yields \(W(y)\).
\end{proof}


\section{The norm and right translates of the Bessel vector}

\begin{lemma}\label{lem:product-integration}
On a product of two countably exhausted finite-measure spaces,
with their product measure on the product sigma algebra, a
nonnegative measurable function has equal product and iterated
integrals in either order. An integrable complex or finite-dimensional
vector-valued function also has equal integrals in either order.
\end{lemma}
\begin{proof}
Restrict first to finite-measure factors. For indicator functions
of measurable rectangles the two iterated integrals give the
product of the two measures by the definition of product measure.
The sets for which both section integrals are measurable and give
the product measure comprise a family containing the whole space,
closed under relative complements of nested members and under
countable disjoint unions by Lemma~\ref{lem:integral-tools}.
To see that it contains every product-measurable set, let
\(\mathcal L\) be the smallest family with these closure properties
containing the rectangles. For a fixed rectangle \(A\) the sets
\(B\) with \(B\cap A\in\mathcal L\) form such a family containing
the rectangles, hence contain \(\mathcal L\). Fix next
\(B\in\mathcal L\) and repeat with the roles interchanged.
This proves closure of \(\mathcal L\) under finite intersections.
Differences and disjointification now give closure under arbitrary
countable unions, proving the assertion about measurable sets.
Increasing simple functions and then increasing finite-measure
exhaustions give the result for nonnegative functions by
Lemma~\ref{lem:integral-tools}. For complex functions separate
the positive and negative parts of their two real coordinates;
absolute integrability makes the resulting integrals finite.
Apply the same reasoning in each coordinate of a vector.
\end{proof}

\begin{lemma}\label{lem:bessel-quotient-density}
Use \(G_1,\mathcal U,H\) from Definition~\ref{def:real-bessel-graph}
and the Haar measures used in Theorem~\ref{thm:actual-real-bessel-vector}.
There exists a positive continuous density \(a(y)\) relative to
\(\dd y\) on \(G_1\) such that \(\int_{\mathcal U}a(ny)\dd n=1\).
Write \(\dd\bar y=a(y)\dd y\). For every nonnegative Borel function
\(h\) and \(x_0\in G_1\) one has
\[
 \int_{G_1}\int_H h(x_0^{-1}s y)\dd\nu_H(s,k)\dd\bar y
                  =\int_{G_1}h(y)\dd y .
\]
Integrals against \(\dd\bar y\) of nonnegative left
\(\mathcal U\)-invariant functions are independent of the choice
of such a density and invariant under right translation. One can
choose \(\dd\bar y\) left invariant under \(\iota(K_0)\).
\end{lemma}
\begin{proof}
Put \(\eta(y)=\exp(-\|y\|_2^2)\) and
\(A(y)=\int_{\mathcal U}\eta(ny)\dd n\). The inclusion
\(\mathcal U\hookrightarrow G_1\) is proper because it is a
closed matrix subgroup. Up to conjugation by the ordinary unitary
matrix \(p\) in Theorem~\ref{thm:actual-real-bessel-vector},
Lemma~\ref{lem:period-ball-cover} applies to it. On compact
sets of \(y\), submultiplicativity gives
\(\|ny\|_2\ge\|n\|_{\rm op}/\sup\|y^{-1}\|_{\rm op}\).
For \(t\ge1\) the exponential of a negative constant times
\(t^2\) is bounded by a constant times \(t^{-8r^2}\),
as follows from the exponential series. The covering lemma
therefore gives an integrable majorant. Positivity of measure
on open sets and Lemma~\ref{lem:integral-tools} show that \(A\)
is finite, continuous and positive. It is left
\(\mathcal U\)-invariant by right invariance of \(\dd n\).
The density \(a=\eta/A\) has the required orbit integral one.
Both numerator and denominator are left \(\iota(K_0)\)-invariant:
compact multiplication preserves the Euclidean norm and compact
conjugation preserves \(\dd n\) by Lemma~\ref{lem:real-graph-character}.

The unipotent measure is invariant under inversion. Indeed the
product density in Proposition~\ref{prop:unipotent-covol} is
bi-invariant, so its pullback by inversion is also bi-invariant
and agrees at identity since the derivative there is minus
identity. The matrix groups are exhausted by compact norm balls,
and all the measures in question are locally finite. By
Lemma~\ref{lem:product-integration}, for any of the stated
densities,
\[
 \int_{G_1}a(y)\int_{\mathcal U}h(ny)\dd n\dd y
  =\int_{G_1}h(w)\int_{\mathcal U}a(n^{-1}w)\dd n\dd w
  =\int_{G_1}h(w)\dd w .
\]
Apply this first to \(h(x_0^{-1}\iota(k)\,\cdot)\) and then
integrate in \(k\), using the compact mass one of
Theorem~\ref{thm:compact-dual-unitary-volume}. This gives the
graph formula. For two densities \(a_1,a_2\) and an invariant
\(h\), insert \(1=\int_{\mathcal U}a_2(ny)\dd n\) in
\(\int h(y)a_1(y)\dd y\) and use the unipotent formula to get
\(\int h(w)a_2(w)\dd w\). A right translate of \(a_1\) also has
orbit integral one. Right invariance of the group volume now
proves invariance of the quotient integral.
\end{proof}

\begin{theorem}\label{thm:bessel-orthonormal-schur}
Put \(D=\dim\pi\) and let \(W\) be the function of
Theorem~\ref{thm:actual-real-bessel-vector}. There is a
positive constant \(\alpha\) such that \(F=\alpha W\) satisfies
\[
       \int\|F(y)\|\dd\bar y<\infty,\qquad
       \int F_i(y)\overline{F_j(y)}\dd\bar y=\delta_{ij}
\]
in every orthonormal basis of \(\pi\). For every
\(\varphi\in L^1(G_1)\cap L^2(G_1)\),
\[
 \int \left\|\int_{G_1}\varphi(g)F(yg)\dd g\right\|^2\dd\bar y
                   \le \frac{D}{d}\int_{G_1}|\varphi(g)|^2\dd g .
\]
\end{theorem}
\begin{proof}
The graph bound on \(W\) integrates to at most \(\|K\|_1\)
by Lemma~\ref{lem:bessel-quotient-density}; this is finite by
Corollary~\ref{cor:row-polar-moment}. Since \(W\) is bounded
and nonzero on a neighborhood of \(x\), its squared norm has
finite positive integral \(q\). Put \(\alpha=\sqrt{D/q}\).
The endomorphism \(M=\int F(y)F(y)^*\dd\bar y\) commutes with
every \(\pi(k)\) by compact equivariance and the invariant
density of Lemma~\ref{lem:bessel-quotient-density}. It is scalar.
Indeed torus invariance makes it diagonal on symmetric tensor
monomials, whose distinct diagonal characters are indexed by
the nonnegative exponent tuples with sum \(a_\tau\).
Commutation with the derivative of the unitary rotation
generated by \(E_{ij}-E_{ji}\) makes the diagonal entries equal
whenever one exponent unit moves from position \(j\) to \(i\).
The corresponding coefficient is nonzero when the exponent at
\(j\) is positive, by differentiating the action on the tensor
factors. All tuples are connected by these moves, including the
single tuple when \(a_\tau=0\) or \(k=1\). As
\(\operatorname{tr}M=D\), this proves \(M=I\).

Write \(R(g)F(y)=F(yg)\), and use the quotient inner product
with conjugate linearity in its first variable. These translates
have the same left unipotent character, so their inner products
are integrals of left unipotent invariant scalars.
Lemma~\ref{lem:bessel-quotient-density} gives a unitary right
action on their closed span. Let \(P=\int e_J(g)R(g)\dd g\).
More explicitly this defines a bounded operator on the span of
the translates, with norm at most \(\|e_J\|_1\), by
Cauchy--Schwarz with the absolute kernel and quotient invariance.
It extends to the closure; approximate the integral by finite sums
there. To justify this approximation, \(g\mapsto R(g)F\) is norm
continuous: as \(g\to1\) its inner product with \(F\) converges
to \(\|F\|^2\) by dominated convergence against the integrable
\(\|F\|\) and boundedness and continuity of \(F\). Its norm is
constant. Translating gives the same conclusion for each vector
of the span, and approximation gives it on the closure.

By symmetry \(P^*=P\). The second convolution identity of
Theorem~\ref{thm:row-real-orthogonality} gives
\(P R(g)P=K(g)P\): in the double integral the density at \(y\)
is \(\int e_J(z)e_J((zg)^{-1}y)\dd z\). Absolute convergence
follows from the squared \(L^1\)-norm of \(e_J\), and that identity
evaluates the density as \(K(g)e_J(y)\).
Theorem~\ref{thm:actual-real-bessel-vector} gives \(PF=F\).
It follows that
\(\langle R(g)F,R(h)F\rangle=D K(g^{-1}h)\).
This is also the Gram matrix of
\(\sqrt D\,\pi_1(g)v_1\) in the cyclic representation
\((\pi_1,v_1)\) of Lemma~\ref{lem:cyclic-row-isometry}, transported
by \(p\). Thus \(R(g)F\mapsto\sqrt D\,\pi_1(g)v_1\) extends
by linearity and completion to an isometry of the two cyclic spans.
For any unit vector \(v\) in the latter span, Cauchy--Schwarz
and the integral identity in that lemma give
\[
 \left|\int\varphi(g)\langle v,\pi_1(g)v_1\rangle\dd g\right|
                                      \le d^{-1/2}\|\varphi\|_2.
\]
Taking the supremum over \(v\) gives the stated convolution
bound after rescaling by \(\sqrt D\). The Hilbert space integral
agrees with the displayed pointwise integral by
Lemma~\ref{lem:product-integration}, first against bounded test
vectors of finite support in measure and then by density.
\end{proof}


\section{Truncation of the real Bessel function}

\begin{lemma}\label{lem:actual-integrated-row-tail}
Use \(G_1,b,r,d\) of Theorem~\ref{thm:actual-real-bessel-vector}.
Put \(c=8r^2+1\), \(V_1=\operatorname{vol}\{\|g-I\|_{\rm op}\le1\}\),
and \(I_j=\int_{G_1}\omega_j(g)\dd g\), where
\(\omega_j(g)=(2/\|g\|_{\rm op})^j\).
For all sufficiently large scales,
\[
 0<V_1<\infty,\qquad I_c\le16^c V_1,\qquad
             I_{2c}\le16^{2c}V_1,\qquad d I_{2c}\le b^{30r^2}.
\]
\end{lemma}
\begin{proof}
The first three bounds follow from
Lemma~\ref{lem:period-ball-cover} and positivity of the smooth
volume near identity. On \(\|g-I\|_{\rm op}\le b^{-4}\),
Theorem~\ref{thm:row-gaussian-phase} gives \(|K(g)|\ge1/2\)
eventually. Indeed take \(T=(r+k)^2\log(b^3)\) as in
Proposition~\ref{prop:actual-bessel-weights}.
The Hilbert--Schmidt norm is at most \(\sqrt r\,b^{-4}\),
the largest weight difference for \(\mu\) is at most \(2b^3\),
and \(r,\log b\) are bounded by fixed powers of \(x\), whereas
\(b\) grows faster than every such power by
Proposition~\ref{prop:thick}. Thus both the Gaussian exponent and
the product of the matrix difference norm with \(T\) tend to zero.
The integral of \(|K|^2\) is \(1/d\) by
Theorem~\ref{thm:row-real-orthogonality}. The small ball therefore
has volume at most \(4/d\).
In Lemma~\ref{lem:period-ball-cover} take \(R=2,\delta=b^{-4}\),
noting that the set defining \(V_1\) lies in this sublevel set.
This gives \(d V_1\le4(1+8b^4)^{2r^2}\le4 b^{10r^2}\).
Now \(2c=16r^2+2\), so \(4\cdot16^{2c}b^{10r^2}\le b^{30r^2}\)
for \(r,b\) sufficiently large.
\end{proof}

\begin{theorem}\label{thm:actual-bessel-truncation}
Retain \(F,\alpha,D,x\) of
Theorems~\ref{thm:actual-real-bessel-vector}
and~\ref{thm:bessel-orthonormal-schur}. Put \(n=\nu+1\),
\(R=b^{6n+2}\), \(L=3r^3 f\), where \(f=|\Phi|\) is the
number of real places, and
\[
       S(y)=1+\inf_{(s,k)\in H}\|s y\|_{\rm op},\qquad
       E=\{y:S(y)\le R\}.
\]
There is a bounded continuous vector \(V(y)\) with the same
left transformation formula as \(F\), vanishing off \(E\), with
\[
       \int\|V\|^2\dd\bar y\le D,\qquad
       \int\|F-V\|^2\dd\bar y\le b^{-10r^2}D .
\]
There are a nonnegative Borel function \(A\) invariant under
left translation by the first projection of \(H\), and a positive
finite Borel measure \(\lambda\) on \(G_1\), satisfying
\[
 \begin{gathered}
  \max(\|F(y)\|,\|V(y)\|)
       \, b^{\,L{\bf1}_{F(y)\ne V(y)}}\le A(y)
       \le\lambda\{g:yg\in E\},\\
                  \lambda(G_1)\int A(y)\dd\bar y\le b^{2r^3}D .
 \end{gathered}
\]
The support set \(E\) has compact image modulo \(\mathcal U\).
\end{theorem}
\begin{proof}
The function \(S\) is invariant under the first projection.
For \(y'\to y\), submultiplicativity gives
\(S(y')-1\le(S(y)-1)\|y^{-1}y'\|_{\rm op}\) and the reverse
bound with \(y,y'\) exchanged. This proves continuity. If
\(S(y)\le R\), some \(s\) satisfies \(\|sy\|_{\rm op}\le R+1\).
Its compact middle component is ordinary unitary, so
\(\|n_0 y\|_{\rm op}\le R+1\) for some \(n_0\in\mathcal U\).
The group norm ball is compact. Since \(E\) is closed and
\(\mathcal U\)-invariant, its image is the image of its closed
intersection with that compact ball, proving compactness.

Use \(c,\omega_j,I_j,V_1\) from
Lemma~\ref{lem:actual-integrated-row-tail}. Write
\[
  t_s=x^{-1}s y,\quad
  H_2(y)=\int_H\omega_{2c}(t_s)\dd\nu_H(s,k),\quad
  P=b^{33(r+k)^2}16^c .
\]
The integral with exponent \(c\) is at most \(P\|W(x)\|\) by
Theorem~\ref{thm:actual-real-bessel-vector}. Since
\(\omega_{2c}\le 2^c\omega_c\), \(H_2\) is finite, Borel and
invariant by right translation in the graph integral when \(y\)
is left translated by the first projection. By
Lemma~\ref{lem:bessel-quotient-density}, \(\int H_2\dd\bar y=I_{2c}\).

If \(\|t_s g-I\|_{\rm op}\le1\), then \(\|t_sg\|_{\rm op}\le2\),
so \(S(yg)\le1+2\|x\|_{\rm op}\le R\) by
Theorem~\ref{thm:actual-real-bessel-vector}. Moreover
\(\|g\|_{\rm op}\le2\|t_s\|_{\rm op}\), since a \(J\)-unitary
has \(\|t_s^{-1}\|_{\rm op}=\|t_s\|_{\rm op}\).
The left translate \(\{\|t_s g-I\|_{\rm op}\le1\}\) has volume
\(V_1\). Integration on it therefore gives
\[
  \int_{\{g:yg\in E\}}\omega_c(g)\dd g
                       \ge V_1\,2^{-c}\omega_c(t_s).
\]
Multiply by \(2^c\omega_c(t_s)/V_1\) and integrate in \(s,k\).
It follows that \(H_2(y)\le\lambda_0\{g:yg\in E\}\) for the measure
\[
 \dd\lambda_0(g)=2^c P\,\|W(x)\|\,\omega_c(g)\dd g/V_1 .
\]
Writing \(p_0=150r^2\), Lemma~\ref{lem:actual-integrated-row-tail}
gives \(\lambda_0(G_1)\le b^{p_0}\|W(x)\|\) eventually:
\(r+k\le2r\) and \(2^c16^{2c}\le b^{2c}\) for sufficiently
large \(b\).

By Corollary~\ref{cor:row-polar-moment} and the pointwise
bound of Theorem~\ref{thm:polar-row-decay},
\(|K(t)|\le\omega_j(t)\) whenever \(0\le j\le r^8\),
also when \(\|t\|_{\rm op}<2\) by \(|K|\le1\).
Use \(j=2c\) in the reproducing identity at \(x\) and apply
Cauchy--Schwarz with measure \(\omega_{2c}(t)\dd t\).
Then unfold with Lemma~\ref{lem:bessel-quotient-density},
using the invariance of \(\|W\|\) under the graph, to obtain
\[
 \begin{aligned}
 \|W(x)\|^2
   &\le d^2 I_{2c}\int_{G_1}\omega_{2c}(t)\|W(xt)\|^2\dd t\\
   &=d^2 I_{2c}\int H_2(y)\|W(y)\|^2\dd\bar y
      \le d^2 I_{2c} b^{p_0}\|W(x)\|\,D/\alpha^2 .
 \end{aligned}
\]
Since \(W(x)\ne0\), it follows for
\(A=\alpha H_2,\ \lambda=\alpha\lambda_0\) that
\[
  \lambda(G_1)\int A\dd\bar y
      \le D b^{2p_0}(d I_{2c})^2
      \le D b^{360r^2}\le D b^{2r^3}
\]
eventually. Also \(\int A^2\dd\bar y\le\lambda(G_1)\int A\dd\bar y\).

Set \(\zeta(y)=\max(0,\min(1,R-S(y)))\) and \(V(y)=\zeta(y)F(y)\).
The cutoff is continuous and invariant under the graph.
If \(F(y)\ne V(y)\), then \(S(y)>R-1\). In that case
\(\|t_s\|_{\rm op}>2b\) for all \(s\), since otherwise
\(S(y)\le1+2b\|x\|_{\rm op}\le R-1\) for \(b\) sufficiently large.
From \(f\le N_1=O_M(x^{11}(1+\log x))\) and \(r\ge x^5\) in
Propositions~\ref{prop:thick} and~\ref{prop:level-prime}
one has \(2c+L\le r^8\) eventually. Apply the pointwise bound
with \(j=2c+L\) on this set, and with \(j=2c\) elsewhere, in
the integral for \(W\). This gives the required majorization,
including that for \(V\) since \(\zeta\in[0,1]\).
The norm of \(V\) is bounded by that of \(F\), and the squared
difference integrates to at most
\(b^{-2L}\int A^2\dd\bar y\le b^{-10r^2}D\) because
\(2L-2r^3\ge4r^3\) and \(r\to\infty\).
\end{proof}


\section{Right derivatives of the real Bessel functions}

\begin{lemma}\label{lem:row-holomorphic-ball}
Retain \(\psi\) of Definition~\ref{def:polar-row-coefficient}
in rank \(r\), and put \(H_0=1+r+\max_i |t_i|\).
There is a holomorphic extension of \(\psi\) near identity
in the complex matrix space, of modulus at most two on
\(\|g-I\|_{\rm op}\le1/(32rH_0)\).
The function \(\psi\) is smooth on the whole real group.
\end{lemma}
\begin{proof}
For \(\delta=\|g-I\|_{\rm op}<1\) the real part of the matrix
is positive definite in the sense of
Definition~\ref{def:block-splitting-integral}. Use the splitting
\(a=I\) there. Its numerator is \(|m_\Lambda(k)|^2\), and the
denominator is \(m_{-\mu}(k^{-1}gk)\). At
\(\delta\le1/(32rH_0)\), the logarithm series on every principal
compression differs from zero by at most \(2\delta\) in operator
norm, since \(\sum_{j\ge1}\delta^j/j\le\delta/(1-\delta)\).
Its trace, of absolute value at most \(2r\delta\), exponentiates
to the determinant by Definition~\ref{def:compressed-logarithm}
and its determinant calculation. In the product of minors the
sum of the absolute proper exponents is the variation of
\(-\mu\), at most \(2H_0\); the absolute exponent of the full
determinant is at most \(H_0\). Hence the reciprocal denominator
has modulus at most \(\exp(6rH_0\delta)<2\). Integrate with
the probability \(c_\Lambda^{-1}|m_\Lambda(k)|^2\dd k\).
The integral agrees with \(\psi\) on the real group by
Theorem~\ref{thm:accretive-splitting}. Holomorphy follows
by integrating the Taylor series of the rational integrands
on compact interior polydisks, with the coefficient bound
of Lemma~\ref{lem:analytic-tools} successively in the variables.
Around a general point \(g_0\) use its fixed admissible
block-diagonal splitting from Theorem~\ref{thm:accretive-splitting}.
It stays admissible on a complex neighborhood; the compact
set of all denominators stays disjoint from zero after shrinking
that neighborhood. The identical argument yields a local holomorphic
extension there, proving smoothness on the real group.
\end{proof}

\begin{theorem}\label{thm:bessel-all-right-derivatives}
Let \(F,\pi,D,d\) be as in Theorem~\ref{thm:bessel-orthonormal-schur}
and \(n=\nu+1\). Index words \(w\) by ordered pairs of matrix
coordinates \((i,j)\) in \(G_1\), and write \(|w|\) for the
length. There exist bounded smooth \(\pi\)-valued functions \(F_w\)
with \(F_\varnothing=F\), all obeying the left transformation rule
of Theorem~\ref{thm:actual-real-bessel-vector}. If \(\gamma\)
is a real differentiable group curve with
\(\dot\gamma(t)=\gamma(t)X\), then
\[
 \frac{\dd}{\dd t}F_w(\gamma(t))
                   =\sum_{i,j}X_{ij}F_{(i,j)w}(\gamma(t)).
\]
For all sufficiently large scales, all \(|w|\le 2n^2\), and
every integrable \(\varphi\) on \(G_1\),
\[
  \int\left\|\int_{G_1}\varphi(g)F_w(yg)\dd g\right\|^2\dd\bar y
      \le (b^5)^{2|w|}\frac{D}{d}\int_{G_1}|\varphi(g)|^2\dd g .
\]
\end{theorem}
\begin{proof}
The compactly supported reproducing function \(\beta_0\) in
Theorem~\ref{thm:row-real-orthogonality} can be taken smooth.
In its proof choose the real scalar cutoff \(\eta\) to be a
smooth function of \(\operatorname{tr}(g^*g)\) supported below
\(r+2\) and positive at \(r\); for instance take
\(\exp(-1/(r+2-t))\) for \(t<r+2\) and zero beyond.
Repeated differentiation gives polynomials in \(1/(r+2-t)\)
times the same exponential, hence all derivatives vanish at
the endpoint by the exponential series. The remaining factor
is smooth by Lemma~\ref{lem:row-holomorphic-ball}.
Transport the function by \(p\), writing \(B_0\) for it on \(G_1\).
Then \(B_0*e_J=e_J\) and
\[
                F(y)=\int_{G_1} B_0(z)F(yz)\dd z
\]
by reproduction and Lemma~\ref{lem:product-integration}.
For any tangent matrix \(X^*J+JX=0\), replacing \(y\) by
\(y\exp(tX)\) transfers the derivative to
\(B_0(\exp(-tX)z)\). Thus every successive right derivative
is continuous and bounded, given by the same integral with a
smooth compactly supported derivative of \(B_0\). This follows
by differentiation under the integral, since on compact
parameter sets the supports lie in one compact set and \(F\)
is bounded. It also holds in quotient Hilbert norm for the
right action \(R\) of Theorem~\ref{thm:bessel-orthonormal-schur}.
Indeed the norm of the integral with any integrable scalar
kernel \(h\) is at most \(\|h\|_1\|F\|_2\) by unitarity,
and the difference quotients of a smooth compact kernel
converge in \(L^1\) to its derivative by uniform convergence
on that compact support. They lie in the closed cyclic span.
These smooth derivatives in a basis of tangent directions give
ordinary smoothness in group coordinates, also seen directly
by replacing \(\exp(tX)\) by a local smooth matrix chart.

For any complex matrix \(Z\), put \(Z^\dagger=JZ^*J\) and
define its derivative by complex linear extension from the
real tangent space, using
\[
       Z=(Z-Z^\dagger)/2+i(Z+Z^\dagger)/(2i).
\]
The two matrices multiplying \(1,i\) are tangent, each of
operator norm at most \(\|Z\|_{\rm op}\). The decomposition is
unique by the involution, so this extension is complex linear.
Define \(F_{(i,j)w}\) as the derivative of \(F_w\) for the
matrix unit \(E_{ij}\). Smoothness, boundedness, the chain
rule in the statement and the transformation under the graph
follow from the preceding construction and commutation of
left and right translations.

Take a word of length \(l\ge1\) in real tangent matrices
\(X_1,\ldots,X_l\) of operator norms at most one. Its mixed
derivative vector is the derivative, once in each variable
at zero, of \(R(g(t))F\) for
\(g(t)=\exp(t_1X_1)\cdots\exp(t_lX_l)\).
By the Gram identity of Theorem~\ref{thm:bessel-orthonormal-schur}
its squared Hilbert norm is \(D\) times the derivative in all
\(2l\) real variables of
\(K(g(t)^{-1}g(s))\) at zero. The expression inside \(K\) is a
product of \(2l\) exponentials, with the first \(l\) in reverse
order and with negative parameters. For complex parameters of
moduli \(\le u=(128lrH_0)^{-1}\) its difference from identity
has operator norm at most \(\exp(2lu)-1\le4lu\). This follows
by expanding the exponentials and bounding the sum of all
nonconstant terms. Conjugation by the ordinary unitary \(p\)
preserves the bound. Lemma~\ref{lem:row-holomorphic-ball}
gives a holomorphic scalar of modulus at most two throughout
this polydisk. Applying Lemma~\ref{lem:analytic-tools} in
each variable bounds the mixed derivative by \(2u^{-2l}\);
its factorial factor is one since each variable occurs once.
The original word therefore has norm at most \(\sqrt{2D}u^{-l}\).
Expanding complex matrix units into their two tangent parts
shows
\[
       \|F_w\|_2\le\sqrt D(512 l r H_0)^l\qquad(l=|w|\ge1).
\]
For \(l\le 2n^2\), \(512lrH_0\le b^5\) eventually by
Propositions~\ref{prop:actual-bessel-weights} and~\ref{prop:thick}:
\(H_0\le b^3+r+1\), and \(r,n\) grow at most polynomially in \(x\).

Under the cyclic isometry in Theorem~\ref{thm:bessel-orthonormal-schur}
any such vector \(F_w\) corresponds to a vector of the same norm
in the irreducible cyclic representation. For \(\varphi\in L^1\cap L^2\)
and a unit test vector there, the two-vector formula of
Lemma~\ref{lem:cyclic-row-isometry} and Cauchy--Schwarz bound the
absolute integrated coefficient by
\(d^{-1/2}\|F_w\|_2\|\varphi\|_2\). Taking the supremum over
unit vectors and squaring proves the desired bound. Length
zero is Theorem~\ref{thm:bessel-orthonormal-schur}; for
\(\varphi\in L^1\) outside \(L^2\) the pointwise integral exists
by boundedness, and the asserted extended-real inequality holds.
\end{proof}


\section{Coordinates for two Bessel frames}

\begin{definition}\label{def:open-bessel-frame}
In the coordinates of Definition~\ref{def:real-bessel-graph}
put \(n=\nu+1\), let \(F_0=(e_1,\ldots,e_\nu,z)\), and let
\(L_0=(e_1^*,\ldots,e_\nu^*,u)\) be the indicated matrices of
columns, with \(u\) the fixed unit coordinate in \(P_0\). Write
\(P=e_n e_n^*\) on the separate coordinate space \(\C^n\),
\(Q=I-P\), and \(E_0=(I_\nu\ \ 0)\).
For \(g\in G_1\) put \(Z(g)=F_0^*JgF_0\). For any complex
\(n\)-square matrix \(Z\) put
\[
  \Delta(Z)=|\det Z|^2-|\det Z_{1:\nu,1:\nu}|^2,\qquad
  \Omega=\{Z:\Delta(Z)>0\},
\]
with empty determinant one. Write \(H_0\) for the stabilizer
of the individual columns of \(F_0\) in \(G_1\).
For a \(\nu\)-by-\(k\) complex matrix \(B\) and a
skew-hermitian \(\nu\)-square matrix \(T\), let \(u(B,T)\)
fix \(z\) and on the ordered blocks \(e,P_0,e^*\) act by
\[
   \begin{pmatrix}I&B&T-BB^*/2\\0&I&-B^*\\0&0&I\end{pmatrix}.
\]
\end{definition}

\begin{lemma}\label{lem:frame-point-stabilizer}
The map \((B,T,k_0)\mapsto u(B,T)\iota(k_0)\) parametrizes
\(H_0\) bijectively, smoothly in both directions. A left
and right invariant measure \(\dd h_0\) there is
\(\dd B\,\dd T\,\dd k_0\), with ordinary complex entry measures
for \(B\), real imaginary diagonal and ordinary complex strictly
lower entries for \(T\), and probability on \(U(k)\).
\end{lemma}
\begin{proof}
An isometry fixing \(z\) preserves its orthogonal complement.
Fixing all \(e_i\) gives identity on their columns and on the
bottom quotient by the hermitian pairing. The middle diagonal
block is then some \(k_0\in U(k)\). After right multiplication
by \(\iota(k_0)^{-1}\), the first middle row block \(B\) is free,
the middle last block is \(-B^*\), and the last upper block
\(C\) satisfies \(C+C^*=-BB^*\), by multiplying the Gram
identity. Thus \(T=C+BB^*/2\) gives the claimed unique
parameters. Conversely the displayed matrix in
Definition~\ref{def:open-bessel-frame} satisfies that identity.
Multiplication of two unipotent matrices adds their \(B\)'s and
adds the \(T\)'s with a skew-hermitian correction depending only
on the two \(B\)'s. Left and right translations preserve
\(\dd B\,\dd T\) by successive translations. Compact conjugation
acts unitarily on the \(B\) coordinates and fixes \(T\).
These facts prove both invariances of the product measure.
\end{proof}

\begin{proposition}\label{prop:exchanged-frame-section}
For \(Z\in\Omega\) put \(\xi=(Z^{-1})_{nn}\), \(v=e_n^*Z^{-1}\),
\(s=(1-|\xi|^2)^{1/2}\). Then
\(\Delta(Z)=|\det Z|^2 s^2\) and \(s>0\). Put
\[
 A_Z=-P-Qv^*v(Q/2+P),\qquad B_Z=Q+sP,\qquad
 Y_Z=(F_0 A_Z+L_0 B_Z)Z .
\]
There is an isometry \(\Gamma_Z\in G_1\) carrying \(F_0\) to
\(Y_Z\), carrying \(Y_{Z^*}\) to \(F_0\), and fixing the
complement of the columns of \(F_0,L_0\) pointwise.
It depends real analytically on \(Z\in\Omega\) and
\(Z(\Gamma_Z)=Z\). Every \(g\) with \(Z(g)=Z\) can be written
\(h_1\Gamma_Z h_2^{-1}\) with \(h_1,h_2\in H_0\).
\end{proposition}
\begin{proof}
Positive \(\Delta(Z)\) implies \(\det Z\ne0\), and the cofactor
formula gives \(\Delta(Z)=|\det Z|^2(1-|\xi|^2)\).
In the Euclidean orthonormal columns \((F_0,L_0)\) the Gram
matrix has blocks \(\left(\begin{smallmatrix}-P&Q\\Q&P\end{smallmatrix}\right)\).
Consequently \(F_0^*JY_Z=Z\). To compute the square Gram of
\(Y_Z Z^{-1}\), the first diagonal term is \(-P\), the last
is \(s^2P\), and the two mixed terms sum to
\(-v^*v+|\xi|^2P\). Their sum is \(-v^*v\), giving
\(Y_Z^*JY_Z=-P\). The bases \((F_0,Y_{Z^*})\) and \((Y_Z,F_0)\)
therefore have equal Gram matrix. They are bases of the same
subspace since \(B_Z Z\) and its counterpart for \(Z^*\)
are invertible. On this subspace the required isometry has
matrix, in the basis \((F_0,L_0)\),
\[
 \Gamma_Z=T_Z
       \begin{pmatrix}0&I\\I&0\end{pmatrix}T_{Z^*}^{-1},
       \qquad
 T_Z=\begin{pmatrix}I&A_Z Z\\0&B_Z Z\end{pmatrix}.
\]
Its extension by identity on the positive complement is an
isometry. This expression is real analytic, as the square root
is evaluated at a strictly positive real number.

For the last assertion put \(Y=gF_0\). Let \(C,L,w,D_0\)
be its row blocks on \(e,P_0,z,e^*\), respectively, with
\(w\) a single row. Then \(D_0=E_0 Z,\ w=-e_n^* Z\).
Write \(C Z^{-1}=(C_1,c)\), \(L Z^{-1}=(L_1,u')\), and
\(v=(v_1,\xi)\), splitting off the last column or entry.
The identity \(Y^*JY=-P\), multiplied by \(Z^{-*}\) and
\(Z^{-1}\) on its two sides respectively, gives
\[
  C_1+C_1^*=-v_1^*v_1-L_1^*L_1,\qquad
  c=-v_1^*\xi-L_1^*u',\qquad (u')^*u'=s^2 .
\]
Here \(Z^{-*}=(Z^{-1})^*\).
Conversely these equations impose precisely that Gram identity.
Choose \(k_0\in U(k)\) with \(k_0u=u'/s\); extend a unit vector
to an orthonormal basis by repeated orthogonal projection and
normalization of coordinate vectors. For \(Y_Z\) the corresponding
free data are \(L_1=0,\ u'=su,\ C_1=-v_1^*v_1/2\).
Applying \(u(B,T)\iota(k_0)\) to \(Y_Z\) gives the desired
\(L_1=-B^*,\ u'=s k_0u\) and
\(C_1=-v_1^*v_1/2+T-BB^*/2\). Choose
\(B=-L_1^*\), \(T=(C_1-C_1^*)/2\). The intervening equation
then also gives exactly the prescribed \(c\), so \(Y=h_1Y_Z\).
It follows that \(g^{-1}h_1\Gamma_Z\) fixes \(F_0\).
This proves the decomposition, including \(\nu=0\) with empty
unipotent parameters.
\end{proof}


\section{The measure on the space of frames}

\begin{lemma}\label{lem:regular-fiber-density}
Let \(\Psi\) be a continuously differentiable map from an open
subset of \(\R^a\) to \(\R^b\). Near a point where its derivative
has rank \(b\), choose \(a-b\) coordinate functions \(u\) so that
\((u,\Psi)\) has invertible derivative. With Lebesgue densities on
both spaces, the density on \(\Psi^{-1}(t)\) is given locally by
\[
              |\det D(u,\Psi)|^{-1}\dd u .
\]
This is independent of the complementary coordinate functions,
including smooth replacements for which the derivative remains
invertible. Integrating these densities in \(\dd t\) gives ambient
Lebesgue integration there. Under invertible linear changes with
absolute determinants \(a_0,b_0\) in the ambient and target spaces,
respectively, the density transforms by \(a_0/b_0\) in the new
coordinates.
\end{lemma}
\begin{proof}
Complete a nonzero \(b\)-row minor to an invertible matrix with
coordinate rows. A continuously differentiable square map whose
derivative is invertible has continuously differentiable local
inverse. To verify this, translate and compose with the inverse
derivative at the point. The map now has form \(x+R(x)\),
\(R(0)=0\), with derivative norm of \(R\) at most \(1/2\) on a
small closed cube in the maximum norm. The same Lipschitz bound
holds by integrating along segments. For targets in the concentric
half cube, iteration of \(x\mapsto y-R(x)\) stays in the cube and
has successively halved differences, giving the unique inverse.
The inverse is Lipschitz with constant two. Substituting this
bound in the derivative expansion yields derivative
\((I+DR(x))^{-1}\) for the inverse, which is continuous.
Differentiating this formula gives the smooth version.

For completeness, the local change of variables for these coordinates
uses the absolute determinant. For an invertible linear map the
volume factor is \(|\det|\) by elimination into permutations,
dilations and shears, using translated one-dimensional integration
for a shear. Partition a compact cube inside a diffeomorphism
chart into small cubes. On each subcube with center \(x\), after
subtracting the value at \(x\) and applying the inverse derivative
at \(x\), the map is identity plus a function vanishing at the
center with Lipschitz constant \(\eta\to0\) uniformly under
refinement. This follows from uniform continuity of the derivative
and boundedness of its inverse on the compact cube. Its image
lies in the concentric cube expanded by \(1+\eta\), and contains
the one contracted by \(1-\eta\) by the iteration just proved.
After the linear change back these give the two volume factors
\(|\det| (1\pm\eta)^a\). Images of faces have zero \(a\)-volume
by covering their Lipschitz images by \(O(j^{a-1})\) cubes of
side \(O(1/j)\). Summing over disjoint interiors and passing to
the limit gives the determinant integral for cubes. Finite unions,
subdivision into cubes and the measurable-set closure argument of
Lemma~\ref{lem:product-integration} extend the equality to Borel
integrals in a countable collection of such charts.

Apply that change of variables to \((u,\Psi)\) and integrate in
\(t\) by Lemma~\ref{lem:product-integration}. Comparing two
admissible coordinate systems, the transition fixes \(t\), so
its Jacobian is precisely the Jacobian of its \(u\) coordinates
along the fiber. This proves independence on each fiber, including
the indicated point. Linear ambient and target changes have the
stated ratio by the same determinant rule.
\end{proof}

\begin{proposition}\label{prop:real-frame-orbit-measure}
Let \(\mathcal V=G_1 F_0\), with matrices of full column rank.
Give the space of all complex \(r\)-by-\(n\) matrices coordinate
Lebesgue measure and the hermitian \(n\)-square matrices the
real diagonal and complex strictly lower coordinate measures.
On \(\mathcal V\) write \(\dd\lambda(Y)\) for the fiber density
of \(Y\mapsto Y^*JY\) at \(-P\) from
Lemma~\ref{lem:regular-fiber-density}. Put \(u_j=U_j\) of
Definition~\ref{def:row-vandermonde}, including \(u_0=1\).
For \(h\ge0\) Borel on \(G_1\),
\[
   \int_{G_1}h(g)\dd g =
    \frac{2^n u_k}{u_r}
       \int_{\mathcal V}\int_{H_0} h(g_Y q)\dd h_0(q)
                                      \dd\lambda(Y),
                \qquad g_Y F_0=Y .
\]
The inner integral is independent of the choice of the lift.
\end{proposition}
\begin{proof}
For every full column rank \(Y\), \(Y^*J\) is surjective. The
derivative \(B\mapsto Y^*JB+B^*JY\) reaches an arbitrary
hermitian \(C\) by choosing \(Y^*JB=C/2\), column by column.
Thus the density is defined near every frame. An isometry
of \(J\) has determinant of modulus one and acts on each
ambient column with real volume determinant one. It preserves
the Gram map, hence preserves its level density by
Lemma~\ref{lem:regular-fiber-density}.

At identity, write \(U=J\dot g\), an arbitrary skew-hermitian
matrix in the density of Definition~\ref{def:unitary-density}.
Indeed ordinary unitary conjugation by \(p\) preserves this
density as in Lemma~\ref{lem:unitary-cayley-density}.
The derivative of \(g\mapsto gF_0\) gives exactly the columns
\(UF_0=J\dot Y\). In the coordinates with the \(F_0\) indices
first, these columns have an arbitrary skew-hermitian upper
\(n\)-square block and arbitrary complex lower block. Every
such tuple extends to a skew-hermitian \(U\). These are exactly
the tangent constraints for \(Y^*JY=-P\) at \(F_0\).
The local inverse construction of Lemma~\ref{lem:regular-fiber-density},
applied on a linear complement of the stabilizer directions in
the Cayley chart of Lemma~\ref{lem:unitary-cayley-density},
therefore gives a local section of the frame map. Left translations
give local sections on the entire orbit, open in the full-rank
Gram level. Each \(g\) over such a section has a unique expression
\(g_Yq\) by Lemma~\ref{lem:frame-point-stabilizer}. Two
sections differ by an element of \(H_0\) depending smoothly on
\(Y\); left invariance of \(\dd h_0\) proves the claimed
independence and measurability.

We compute the ratio in these coordinates at identity. The
ambient change \(\dot Y\mapsto J\dot Y\) preserves real volume.
If \(V\) is the upper \(n\)-square block in these coordinates,
the Gram derivative is \(V+V^*\). Passing to \(V+V^*\) and
\((V-V^*)/2\) has absolute determinant \(2^n\) in their free
hermitian and skew-hermitian coordinates. On the diagonal each
real part doubles; in a pair of complex entries the map
\((a,b)\mapsto(a+\bar b,(a-\bar b)/2)\) has real determinant
of absolute value one. Thus \(\dd\lambda\) is \(2^{-n}\)
times the independent first-column-block coordinates of \(U\).

The remaining independent skew-hermitian coordinates of \(U\)
give its block on \(P_0,e^*\). On the tangent of \(H_0\) it is
\(\left(\begin{smallmatrix}\dot k&-\dot B^*\\
\dot B&\dot T\end{smallmatrix}\right)\) by
Lemma~\ref{lem:frame-point-stabilizer}. Their coordinate
density is \(u_k\,\dd h_0\), since the compact probability at
identity has density \(u_k^{-1}\) by
Theorem~\ref{thm:compact-dual-unitary-volume}. Choosing a different
tangent lift of the frame only adds a linear function of the
first coordinates to these remaining entries, a shear of determinant
one. Finally the full group density is \(1/u_r\) times all
independent coordinates of \(U\). This gives \(2^n u_k/u_r\)
at identity. Both densities are left invariant on \(G_1\)
by the preceding section comparison and invariance of the level
density, so the same ratio holds everywhere. Integrating in
countably many charts, made disjoint by successive differences,
gives the formula by Lemma~\ref{lem:product-integration}.
\end{proof}


\section{The exact density of the real open cell}

\begin{theorem}\label{thm:bessel-cell-exact-density}
With the normalizations of
Lemmas~\ref{lem:frame-point-stabilizer} and~\ref{lem:unitary-cayley-density},
write
\(C_{\rm op}=u_k^2/(u_{k-1}u_r)\).
For every nonnegative Borel function \(h\) on \(G_1\),
\[
\begin{split}
 \int_{Z(g)\in\Omega} h(g)\dd g
   = C_{\rm op}\int_\Omega \Delta(Z)^{k-1}
        \iint_{H_0^2}h(q_1\Gamma_Z q_2^{-1})\dd h_0(q_1)
                           \dd h_0(q_2)\,\dd Z ,
\end{split}
\]
where \(\dd Z\) is ordinary Lebesgue measure on its complex entries.
\end{theorem}
\begin{proof}
Use the frame density of Proposition~\ref{prop:real-frame-orbit-measure}
and the notation of Proposition~\ref{prop:exchanged-frame-section}.
For \(Y\) in the ambient matrix space keep \(Z=F_0^*JY\)
as an independent coordinate. Its rows are exactly \(D_0,-w\),
so this step has absolute determinant one. For invertible \(Z\)
replace the remaining rows \(C,L\) by
\[
               C=(C_1,c)Z,\qquad L=(L_1,u')Z .
\]
At fixed \(Z\) the ambient real volume factor is
\(|\det Z|^{2(\nu+k)}\) since there are \(\nu+k\) independent
complex rows. In the constraint replace \(Y^*JY+P\) by
\(Z^{-*}(Y^*JY+P)Z^{-1}\). Equivalently the original output is
the congruence by \(Z\) of this new output. The determinant of
this congruence on hermitian \(n\)-square matrices is
\(|\det Z|^{2n}\), by multiplying by \(i\) in the skew-hermitian
calculation of Lemma~\ref{lem:unitary-cayley-density}.
As \(Z\) itself is an independent coordinate the derivative is
block triangular in this change. Thus the residual factor for
the level density is \(|\det Z|^{2(k-1)}\) by
Lemma~\ref{lem:regular-fiber-density}.

The three blocks of the new output are
\[
 C_1+C_1^*+L_1^*L_1+v_1^*v_1,\qquad
 c+L_1^*u'+v_1^*\xi,\qquad
                    (u')^*u'-(1-|\xi|^2).
\]
Eliminating the first block leaves \(2^{-\nu}\dd T\), with
\(T=(C_1-C_1^*)/2\), by the paired-entry calculation in
Proposition~\ref{prop:real-frame-orbit-measure}.
Eliminating the second gives a complex translation of \(c\),
with determinant one. For the last block, use
\(u'=s\omega\), where \(\omega\) runs over the unit sphere
in \(\C^k\). Radius \(s\) contributes \(s^{2k-1}\dd s\) by
Theorem~\ref{thm:compact-dual-unitary-volume}, and the constraint
uses \(\dd(s^2)=2s\dd s\). At the specified positive
\(s=(1-|\xi|^2)^{1/2}\) the remaining factor therefore is
\[
                 \frac{u_k}{2u_{k-1}}s^{2(k-1)}\dd\omega ,
\]
where \(\dd\omega\) is probability on the sphere. The changes
in the first two output blocks as \(L_1,u'\) vary do not affect
these determinants, since those variables were held fixed when
the first two blocks were eliminated.
Together with \(\dd L_1\,\dd Z\), these are all the free coordinates.
All frames satisfying these equations at \(Z\in\Omega\) lie
in the orbit: the explicit \(q_1Y_Z\) in the proof of
Proposition~\ref{prop:exchanged-frame-section} realizes every
solution. In that formula \(B=-L_1^*\) is volume preserving,
\(T\) is unchanged and \(\omega=k_0u\). The distribution of
\(k_0u\) for compact probability is \(\dd\omega\) by the
first-column calculation in Theorem~\ref{thm:compact-dual-unitary-volume}.

Finally \(\Delta(Z)=|\det Z|^2s^2\), and \(n=\nu+1\) cancels
\(2^{-\nu-1}\) against \(2^n\) in
Proposition~\ref{prop:real-frame-orbit-measure}. The remaining
constant is \(C_{\rm op}\). In its inner integral choose
\(g_Y=q_1\Gamma_Z\). Different choices of \(k_0\) carrying
\(u\) to \(\omega\) give different lifts of the same \(Y\),
so that inner integral is unchanged. The \(H_0\) measure
is inversion invariant by bi-invariance and the derivative
of inversion at identity, as in Lemma~\ref{lem:bessel-quotient-density}.
Changing the second integration variable to its inverse gives
the formula.
\end{proof}

\begin{corollary}\label{cor:bessel-pivot-weights}
Outside a set of Lebesgue measure zero, write
\(Z=l^*\operatorname{diag}(c_1,\ldots,c_n)u\) with \(l,u\)
upper unitriangular complex matrices and \(c_i\ne0\).
Use ordinary complex measure on each independent upper entry
of \(l,u\) and on each \(c_i\). The first two factors of the
measure on the right of Theorem~\ref{thm:bessel-cell-exact-density}
in these coordinates are
\[
 C_{\rm op} (|c_n|^2-1)^{k-1}
      \prod_{j=1}^{n-1}|c_j|^{2(k-1)+4(n-j)}
               \dd l\,\dd u\,\prod_{j=1}^n\dd c_j
                       \qquad(|c_n|>1).
\]
\end{corollary}
\begin{proof}
The set where any initial principal determinant vanishes has
zero Lebesgue measure. Each such polynomial is nonzero at the
identity. More generally, for a nonzero polynomial in complex
coordinates apply induction on the number of variables to its
nonzero leading coefficient in the last variable. Off the null
set of that coefficient there are finitely many roots in the
last variable. Lemma~\ref{lem:product-integration} on bounded
boxes and countable exhaustion prove the claim.
Eliminate the first nonzero pivot \(c_1\). The remaining part
of the first column is \(c_1\) times the conjugates of the
independent first-row entries of \(l\), and the remaining
part of the first row is \(c_1\) times those of \(u\).
The lower-right block is its Schur complement translated
by their outer product times \(c_1\). The first step thus
has volume determinant \(|c_1|^{4(n-1)}\).
The initial determinants of the Schur complement are the
remaining ones divided by \(c_1\), by subtracting the first
row multiples. Repeating gives unique factors, and the total
Jacobian is \(\prod_j|c_j|^{4(n-j)}\). The full and penultimate
determinants are \(\prod_j c_j\) and \(\prod_{j<n}c_j\).
Substitution in \(\Delta\) gives the asserted power and range,
also when \(n=1\).
\end{proof}


\section{Integration of a convolution over nearby frame matrices}

\begin{lemma}\label{lem:bessel-frame-action}
For \(n_0\in\mathcal U\) define \(A(n_0)\) by
\(n_0F_0=F_0 A(n_0)\). This is a surjective homomorphism to
the complex upper unitriangular group \(N_n\), with kernel
\(\mathcal U_0=\{u(B,T)\}\) of
Definition~\ref{def:open-bessel-frame}. Its Haar integration
with the normalization of Proposition~\ref{prop:unipotent-covol} is
\[
 \int_{\mathcal U}h(n_0)\dd n_0
       =\int_{N_n}\int_{\mathcal U_0}h(v\,s(A))\dd v\,\dd A ,
\]
where \(s(A)\) is the section described below,
\(\dd v=\dd B\,\dd T\), and \(\dd A\) is Lebesgue measure in
the strictly upper complex entries. Moreover
\(Z(n_0 y)=A(n_0)^{-*}Z(y)\).
\end{lemma}
\begin{proof}
The defining flag and the identity middle quotient in
Definition~\ref{def:unip} give \(n_0F_0=F_0 A\) with \(A\)
upper unitriangular. The action on these independent columns
is multiplicative. The kernel is the intersection with \(H_0\),
which by Lemma~\ref{lem:frame-point-stabilizer} is \(\mathcal U_0\)
since the positive middle quotient is fixed. Write
\(A=\left(\begin{smallmatrix}a&v\\0&1\end{smallmatrix}\right)\).
Take \(s(A)\) to fix the middle positive space and set its matrix
on \(e,z,e^*\) to be
\[
  \begin{pmatrix}
    a&v& vv^*a^{-*}/2\\
    0&1&v^*a^{-*}\\
    0&0&a^{-*}
  \end{pmatrix}.
\]
The \(e,z,e^*\) Gram pairings are preserved: the two terms from
the \(e,e^*\) pairing cancel the negative middle square on
the last block, the mixed \(z,e^*\) terms cancel, and the
\(e,e^*\) identity block is preserved by \(a^*a^{-*}=I\).
This section lies in \(\mathcal U\) and maps to \(A\).

On \(N_n\), left or right multiplication adds to an entry a
constant and terms from smaller distances above the diagonal,
so \(\dd A\) is bi-invariant. Every \(n_0\) has unique smooth
coordinates \(v_0s(A)\) with \(v_0\in\mathcal U_0\). The
product density is right invariant: multiplying on the right
changes \(A\) by right translation and at fixed \(A\) changes
\(v_0\) by right translation, with factor
\(s(A)n_0s(A A(n_0))^{-1}\) in the kernel. At identity the
coordinate tangent entries are the strictly upper \(e,e\)
entries and the \(e,z\) entries for \(A\), the \(e,P_0\)
entries for \(B\), and the \(e,e^*\) entries for \(T\) with
their skew constraint. They are exactly the independent row
entries and the imaginary central entries in the successive
layers of Lemma~\ref{lem:heisenberg-layer}. Choosing the
opposite half of the skew entries conjugates and changes sign,
with absolute determinant one. The two densities agree at
identity and hence everywhere by right invariance. Lastly use
\(n_0^{-1}F_0=F_0A(n_0)^{-1}\) and \(J\)-unitarity in the
definition of \(Z\) to obtain its transformation law.
\end{proof}

\begin{proposition}\label{prop:frame-multiplicity-bound}
Let \(W_1,m_0\) be nonnegative Borel functions on \(G_1\)
with \(W_1\) left \(\mathcal U\)-invariant. Let \(C\ge0,\ R\ge1\).
Assume that \(\int_{\mathcal U_0}m_0(vy)\dd v\le C\) for all \(y\),
and that \(m_0(y)m_0(n_0y)>0\) implies
\(\max_{ij}|A(n_0)_{ij}|\le R\) for \(n_0\in\mathcal U\).
Then
\[
    \int_{G_1}m_0 W_1\,\dd g
       \le C(\pi R^2)^{n(n-1)/2}\int W_1(y)\dd\bar y .
\]
\end{proposition}
\begin{proof}
Unfold the left integral by the unipotent formula in
Lemma~\ref{lem:bessel-quotient-density}, taking the invariant
\(W_1\) outside the \(\mathcal U\) integral.
Fix \(y\). If no translate has \(m_0(n_0y)>0\), the inner
integral is zero. Otherwise fix one such translate \(n_1y\).
For every other one, \(A(n_0)A(n_1)^{-1}\) has its entries
bounded by \(R\). This restricts \(A(n_0)\) to a right
translate of a product of complex disks in the free entries,
of volume \((\pi R^2)^{n(n-1)/2}\). Integrate first over
\(\mathcal U_0\) in Lemma~\ref{lem:bessel-frame-action} and
use the assumed bound \(C\). This proves the inequality.
\end{proof}

\begin{corollary}\label{cor:open-slice-quotient-bound}
Let \(P_*\ge8,\ c_*>0\) and let \(\mathcal Q\) be a Borel set of
frame matrices with \(\Delta(Z)\ge c_*\) and
\(\|Z\|_{\rm op},\|Z^{-1}\|_{\rm op}\le P_*\). For a nonnegative
Borel \(W_1\) invariant on the left under \(\mathcal U\) and
\(\iota(K_0)\),
\[
 C_{\rm op} c_*^{\,k-1}\int_{\mathcal Q}\int_{H_0}
       W_1(\Gamma_Z h^{-1})\dd h_0(h)\dd Z
                  \le P_*^{3n^2}\int W_1(y)\dd\bar y .
\]
\end{corollary}
\begin{proof}
Use the construction in Lemma~\ref{lem:bessel-quotient-density}
with \(\mathcal U_0\) in place of \(\mathcal U\) to get a
density \(a_0\) with unit \(\mathcal U_0\) orbit integral,
invariant on the left under \(\iota(K_0)\). That proof applies
since \(\mathcal U_0\) is closed by Lemma~\ref{lem:bessel-frame-action}
and compact conjugation preserves \(\dd v\) by
Lemma~\ref{lem:frame-point-stabilizer}. In particular
\(\int_{H_0}a_0(hy)\dd h_0(h)=1\), by compact normalization.
Apply Theorem~\ref{thm:bessel-cell-exact-density} to
\(m_0W_1\) with \(m_0(y)=a_0(y){\bf1}_{\mathcal Q}(Z(y))\).
Both \(W_1\) and \(Z\) are invariant under the left \(H_0\)-action,
so this integral bounds the left side from above.
If \(m_0(y)m_0(n_0y)>0\), Lemma~\ref{lem:bessel-frame-action}
gives \(A(n_0)^{-*}=Z(n_0y)Z(y)^{-1}\).
Taking inverse adjoints bounds \(\|A(n_0)\|_{\rm op}\le P_*^2\).
Also \(\int_{\mathcal U_0}m_0(vy)\dd v\le1\).
Proposition~\ref{prop:frame-multiplicity-bound} applies with
\(C=1,\ R=P_*^2\). The factor
\((\pi P_*^4)^{n(n-1)/2}\) is at most \(P_*^{3n^2}\),
since \(\pi\le4\le P_*\) from the unit disk in its enclosing square.
\end{proof}


\section{Uniform estimates from two bounded frames}

\begin{lemma}\label{lem:stable-open-frame-section}
Write \(s(Z)=(1-|(Z^{-1})_{nn}|^2)^{1/2}\) on \(\Omega\).
Suppose \(M\ge8n,\ |Z_{ij}|\le M^2,\ \Delta(Z)\ge M^{-1}\)
and \(P_*\ge M^{9n}\).
If \(\|Z'-Z\|_{\rm op}\le P_*^{-6}\), then
\[
 \Delta(Z')\ge\Delta(Z)/2,\qquad
   \max(\|Z'\|_{\rm op},\|(Z')^{-1}\|_{\rm op},s(Z')^{-1})
                                                        \le P_*.
\]
For any \(Z'\in\Omega\) with the last three bounds and \(P_*\ge8\),
\[
       \|\Gamma_{Z'}\|_{\rm op}\le P_*^8,\qquad
       \|\dd\Gamma_{Z'}[H]\|_{\rm op}\le P_*^{16}\|H\|_{\rm op}
\]
for every real tangent perturbation \(H\) of the complex matrix.
\end{lemma}
\begin{proof}
The norms of \(Z,Z'\) are at most \(nM^2+1\le M^4\).
For either the full determinant or the initial minor, multilinearity
and Lemma~\ref{lem:hadamard} bound the change by
\(n M^{4(n-1)}P_*^{-6}\), on replacing one column at a time.
The two absolute determinants are at most \(M^{4n}\); for
the empty minor its change is zero. It follows that
\(|\Delta(Z')-\Delta(Z)|\le4n M^{8n}P_*^{-6}\le1/(2M)\).
In particular \(|\det Z'|\ge(2M)^{-1/2}\).
The product of the singular values equals the absolute determinant
by taking determinants in the polar decomposition as in
Theorem~\ref{thm:polar-row-decay}. Thus the inverse norm is at
most \(M^{4(n-1)}(2M)^{1/2}\), and
\(s(Z')^{-1}=|\det Z'|/\sqrt{\Delta(Z')}\le M^{4n}(2M)^{1/2}\).
This proves the first assertions.

For the derivative normalize \(\|H\|_{\rm op}=1\); the zero
direction is immediate. Use the matrices in
Proposition~\ref{prop:exchanged-frame-section} at \(Z'\) and
at its adjoint. Differentiating an inverse gives
\(\dd v\) of norm at most \(P_*^2\), while
\(|\dd s|\le|\xi|\,|\dd\xi|/s\le P_*^3\).
Thus \(\|A_{Z'}\|\le2P_*^2,\ \|\dd A_{Z'}\|\le2P_*^3\),
\(\|B_{Z'}\|\le1,\ \|B_{Z'}^{-1}\|\le P_*,\
\|\dd B_{Z'}\|\le P_*^3\), in operator norm.
The triangular expression for \(T_{Z'}\) gives
\(\|T_{Z'}\|\le4P_*^3,\ \|\dd T_{Z'}\|\le4P_*^4\).
Its inverse has upper right block \(-A_{Z'} B_{Z'}^{-1}\)
and lower right block \((Z')^{-1}B_{Z'}^{-1}\); its norm
is at most \(4P_*^3\). The derivative of this inverse has
norm at most \(64P_*^{10}\) by the inverse derivative identity.
In the formula for \(\Gamma\) all these bounds also apply at
the adjoint. Multiplication gives norm at most \(16P_*^6\),
and the product rule bounds its derivative by
\(16P_*^7+256P_*^{13}\le P_*^{16}\), since \(P_*\ge8\).
The action on the orthogonal complement is constant identity.
\end{proof}

\begin{proposition}\label{prop:bessel-two-frames-bound}
For \(\rho\ge8n\) put
\(\mathcal B_\rho=\{y:\|y^{-1}f_j\|\le\rho,\ 1\le j\le n\}\),
with Euclidean column norm. Use \(S(y)\) from
Theorem~\ref{thm:actual-bessel-truncation}. Suppose
\(Z\in\Omega,\ |Z_{ij}|\le\rho^2\), and put
\(M=\rho^{9n}/\Delta(Z)\). Then \(M\ge\rho\). If
\(y\in\mathcal B_\rho\) and \(S(\Gamma_Z y)\le\rho\), then
\(\|y\|_{\rm op}\le M^{24n}\).
The entry bound assumed here holds when \(Z=Z(\gamma)\) and
\(\mathcal B_\rho\cap\gamma\mathcal B_\rho\ne\varnothing\).
\end{proposition}
\begin{proof}
In the last assertion write \(w,\gamma^{-1}w\in\mathcal B_\rho\).
The cross Gram of \(w^{-1}F_0\) and \(w^{-1}\gamma F_0\)
is \(Z(\gamma)\). The form matrix \(J\) has norm one, so
Cauchy--Schwarz gives the entry bound. Hence
\(\Delta(Z)\le|\det Z|^2\le(n\rho^2)^{2n}\le\rho^{6n}\)
by Lemma~\ref{lem:hadamard}. This gives \(M\ge\rho\) and
\(\Delta(Z)\ge M^{-1}\).

Choose \(v\in\mathcal U\) with \(\|v\Gamma_Z y\|_{\rm op}\le\rho\).
Indeed an element of the graph with this bound exists by the
definition of the infimum with the additive term one, and
removing its compact factor preserves the norm. Write
\(A=A(v)\) and \(\widetilde Z=Z(v\Gamma_Z)=A^{-*}Z\) by
Lemma~\ref{lem:bessel-frame-action}. The two column frames
\((v\Gamma_Z y)^{-1}F_0,\ y^{-1}F_0\) have cross Gram
\(\widetilde Z\); their individual column norms are at most
\(\rho\). Thus \(\|\widetilde Z\|_{\rm op}\le n\rho^2\le M^3\).
Left multiplication by a lower unitriangular matrix preserves
each leading determinant, so \(\Delta(\widetilde Z)=\Delta(Z)\).
It follows for \(U=Z,\widetilde Z,Z^*\) that
\[
  \|U\|_{\rm op}\le M^3,\qquad
  \|U^{-1}\|_{\rm op}\le M^{3(n-1)}\sqrt M\le M^{3n},\qquad
  s(U)^{-1}\le M^{3n}\sqrt M .
\]
Here the inverse and \(s\) estimates use the same singular value
and determinant formulas as Lemma~\ref{lem:stable-open-frame-section}.
Consequently \(\|A\|_{\rm op}\le M^{3+3n}\le M^{6n}\),
since \(A^*=Z\widetilde Z^{-1}\). The inverse block formula
in that lemma gives \(\|T_{Z^*}^{-1}\|_{\rm op}\le M^{11n}\):
its upper right block has norm at most
\((1+M^{6n})M^{3n}\sqrt M\), its lower right at most
\(M^{6n}\sqrt M\).

Now \(\Gamma_Z^{-1}F_0=Y_{Z^*}\) and \(v F_0=F_0 A\).
Hence \(\|y^{-1}Y_{Z^*}\|_{\rm op}\le\rho\|A\|_{\rm op}\),
whereas \(\|y^{-1}F_0\|_{\rm op}\le\sqrt n\,\rho\).
Multiplying the two-frame matrix by \(T_{Z^*}^{-1}\) bounds
\(y^{-1}\) on the span of the columns of \(F_0,L_0\) by
\[
             (\sqrt n\,\rho+\rho M^{6n})M^{11n}\le M^{20n}.
\]
This span contains the entire negative eigenspace of \(J\):
the complement is the positive space \(X\) in
Proposition~\ref{prop:bounded-bessel-isometry}.
Let \(P_-\) be its negative eigenspace projection, so
\(J=I-2P_-\). For \(z_0=y^{-1}\), the identity
\(z_0Jz_0^*=J\) gives
\(z_0z_0^*=J+2z_0P_-z_0^*\).
Consequently \(\|y\|_{\rm op}^2=\|z_0\|_{\rm op}^2
\le1+2M^{40n}\le M^{48n}\). This proves the claimed bound.
\end{proof}


\section{Evaluation of a convolution at a frame matrix}

\begin{lemma}\label{lem:bessel-cube-chain}
Let \(d_0\ge1,\ \eta>0\), and let a finite tree of vector-valued
functions on \([0,\eta]^{d_0}\) have depth \(d_0\), root \(f\),
and at most \(m_0\) children per vertex. Suppose each function
is continuously differentiable and, at depths \(<d_0\), each of
its first coordinate derivatives is a linear combination of its
children with scalar coefficients of modulus \(\le K_0\).
Coefficients may depend on the point and the coordinate.
Allow an additional measure parameter \(\omega\) with all functions
measurable and these bounds holding at each \(\omega\). Then
\[
 \int\|f(0,\omega)\|^2\dd\omega
   \le (2/\eta+2\eta m_0^2 K_0^2)^{d_0}
       \max_j\iint_{[0,\eta]^{d_0}}\|f_j(t,\omega)\|^2\dd t\,\dd\omega ,
\]
where the maximum is over every vertex, including the root.
\end{lemma}
\begin{proof}
For a vector-valued \(C^1\) function \(v\) of one coordinate,
write \(v(0)=v(t)-\int_0^t v'(u)\dd u\). Cauchy--Schwarz and
averaging over \(0\le t\le\eta\) give
\[
       \|v(0)\|^2\le (2/\eta)\int_0^\eta\|v\|^2
                          +2\eta\int_0^\eta\|v'\|^2 .
\]
The squared derivative is at most
\(m_0 K_0^2\sum_j\|v_j\|^2\) by the assumed expansion and
Cauchy--Schwarz. Apply this inequality successively in the
coordinates, integrating the ones already treated and holding the
others at zero. At every step the sum of the nonnegative scalar
multipliers of the new integrals increases by a factor at most
\(2/\eta+2\eta m_0^2K_0^2\). At depth \(i\) of this iteration
only vertices of tree depth at most \(i\) have been used.
Integrate also in \(\omega\), using
Lemma~\ref{lem:product-integration}, and bound the final integrals
by their maximum. This gives the formula.
\end{proof}

\begin{theorem}\label{thm:bessel-section-convolution}
Use \(F,D,d\) of Theorem~\ref{thm:bessel-orthonormal-schur}.
Suppose \(P_*\ge\max(8r,b^5)\), \(c_*>0\), and throughout
\(\|Z'-Z\|_{\rm op}\le P_*^{-6}\) suppose
\[
 \Delta(Z')\ge c_*,\qquad
 \max(\|Z'\|_{\rm op},\|(Z')^{-1}\|_{\rm op},s(Z')^{-1})\le P_* .
\]
For \(\varphi\in L^1(G_1)\cap L^2(G_1)\) supported where
\(\|y\|_{\rm op}\le P_*^{20}\), one has
\[
  C_{\rm op} c_*^{\,k-1}
    \int_{\substack{h\in H_0\\ \|h\|_{\rm op}\le P_*^{20}}}
       \left\|\int \varphi(y)F(\Gamma_Z h^{-1}y)\dd y\right\|^2
                     \dd h_0(h)
       \le P_*^{430n^2}\,\frac{D}{d}\|\varphi\|_2^2 .
\]
\end{theorem}
\begin{proof}
Vary \(Z\) in its \(d_0=2n^2\) real matrix coordinates by
\(Z_t=Z+\sum_{a=1}^{d_0}t_a E_a\), where \(E_a\) are the
ordinary real and imaginary matrix units. Use the cube of
side \(\eta=P_*^{-8}\). Its operator displacement is at most
\(2n\eta\le P_*^{-6}\), by the entry bound and Cauchy--Schwarz.
Put \(G(t,h)=\Gamma_{Z_t}h^{-1}\). For differentiation in
any coordinate the right tangent matrix
\[
       T=G(t,h)^{-1}\partial_a G(t,h)
          =h\Gamma_{Z_t}^{-1}(\partial_a\Gamma_{Z_t})h^{-1}
\]
has norm at most \(P_*^{64}\), by
Lemma~\ref{lem:stable-open-frame-section} and the equality of
the norms of each isometry and its inverse.

The root of the tree in Lemma~\ref{lem:bessel-cube-chain} is the
integral in the statement with \(Z_t\) in place of \(Z\).
Every vertex of depth \(l\) has form
\(\int\varphi_j(y)F_w(G(t,h)y)\dd y\), \(|w|=l\),
with \(|\varphi_j|\le|\varphi|\).
For a coordinate derivative the tangent at the argument of
\(F_w\) is \(y^{-1}Ty\). Use
Theorem~\ref{thm:bessel-all-right-derivatives} and expand its
\((i,j)\)-entry as
\(\sum_{a,b}(y^{-1})_{ia}T_{ab}y_{bj}\).
A child has the new word \((i,j)w\) and the new kernel obtained
by multiplying the parent kernel by
\((y^{-1})_{ia}y_{bj}/P_*^{40}\). This multiplier has modulus
at most one on the support. The corresponding scalar coefficient
is \(P_*^{40}T_{ab}\), bounded by \(K_0=P_*^{104}\).
There are at most \(m_0=r^4\) children. The derivative passes
under the integral because each derivative is bounded on the
specified region by a constant times \(|\varphi(y)|\), using
boundedness of the finitely many functions of each depth and
the smoothness on compact neighborhoods of the section. This
also gives continuity of the partial derivatives by dominated
convergence, including those at depth \(d_0\) using smoothness
at the next depth. The parameter \(h\) ranges over the
measurable subset indicated in the statement.

For each vertex use on the whole group the nonnegative function
\[
            W_1(g)=\left\|\int\varphi_j(y)F_w(gy)\dd y\right\|^2.
\]
Its left invariances follow from the unitary transformation
formula for \(F_w\). Theorem~\ref{thm:bessel-all-right-derivatives}
bounds its quotient integral by
\(P_*^{2l}(D/d)\|\varphi\|_2^2\) for \(l\le d_0\).
Corollary~\ref{cor:open-slice-quotient-bound} on this cube
then bounds its full integral in \(t,h\), multiplied by
\(C_{\rm op}c_*^{k-1}\), by \(P_*^{3n^2+2l}(D/d)\|\varphi\|_2^2\).
Restriction of \(h\) decreases this nonnegative integral.
Lebesgue measure of the entries of \(Z_t\) is exactly \(\dd t\).

Finally \(m_0^2\le P_*^8\) and
\(2/\eta+2\eta m_0^2K_0^2\le2P_*^8+2P_*^{208}\le P_*^{210}\).
Lemma~\ref{lem:bessel-cube-chain} therefore gives total exponent
at most \(210(2n^2)+3n^2+4n^2\le430n^2\), as required.
\end{proof}


\section{The real open-cell second moment}

\begin{lemma}\label{lem:bessel-unipotent-tube}
For every \(T_0\ge8r\) there exist a compact
\(\Theta\subset\mathcal U_0\) of positive finite volume and a
compact \(\mathcal T\subset G_1\) such that
\[
 \Theta\{g:\|g\|_{\rm op}\le T_0\}\subset\mathcal T
 \subset\{g:\|g\|_{\rm op}\le T_0^{12}\},\qquad
 \int_{\mathcal U_0}{\bf1}_{\mathcal T}(v y)\dd v
                                      \le2\operatorname{vol}\Theta
\]
for every \(y\in G_1\).
\end{lemma}
\begin{proof}
Use the coordinates \(u(B,T)\) and measure in
Lemma~\ref{lem:frame-point-stabilizer}. Put
\(K_1=\sqrt r\,T_0^2\), \(L=8r^2K_1\) and take \(\Theta\)
to be \(\|B\|_2\le L,\ \|T\|_2\le L^2\) in Hilbert--Schmidt
norms. It is compact with positive finite volume, also if a
parameter space is zero-dimensional. Set \(\mathcal T\) equal
to its product with the closed group ball of radius \(T_0\).
That ball is compact, and the explicit block matrix bounds
\(\|u(B,T)\|_{\rm op}\le5L^2\) on \(\Theta\). Since
\(L=8r^{5/2}T_0^2\le T_0^5\), the second inclusion follows.

If \(\|u(B',T')\|_{\rm op}\le T_0^2\), reading the top
middle block gives \(\|B'\|_2\le K_1\). The skew-hermitian
part of the top last block gives \(\|T'\|_2\le K_1\).
Multiplication on the right by such an element takes
\(\Theta\) inside the two balls of radii
\(L+K_1,\ L^2+K_1+LK_1\). Indeed the central correction is
a half difference of the two products \(B'B^*,BB'^*\).
Thus the volume of their entire product set is at most
\((1+2K_1/L)^{2r^2}\operatorname{vol}\Theta
\le e^{1/2}\operatorname{vol}\Theta\le2\operatorname{vol}\Theta\),
by scaling of Lebesgue measure in real dimension
\(2\nu k+\nu^2\le2r^2\).
If \(v_1y=w_1g_1\in\mathcal T\), \(w_1\in\Theta\),
\(\|g_1\|_{\rm op}\le T_0\), then any other
\(vy=wg\in\mathcal T\) satisfies
\(w^{-1}v v_1^{-1} w_1=g g_1^{-1}\in\mathcal U_0\) of norm
at most \(T_0^2\). The set of such \(v\)'s is contained in a
right translate of the product just estimated. Right invariance
gives the last inequality. If there is no \(v_1\), the integral
is zero.
\end{proof}

\begin{theorem}\label{thm:actual-bessel-open-cell}
Let \(F\) be the actual normalized real Bessel function of
Theorem~\ref{thm:bessel-orthonormal-schur}. Let \(\rho\ge b\),
\(\gamma\in G_1\),
\(\mathcal B_\rho\cap\gamma\mathcal B_\rho\ne\varnothing\),
and \(\Delta=\Delta(Z(\gamma))>0\). Put \(M=\rho^{9n}/\Delta\).
At all sufficiently large arithmetic scales there exist a compact
set \(\mathcal T\subset G_1\), a random \(t\in H_0\) with a
specified Borel probability law, and a positive finite measure
\(\lambda_\gamma\) on \(G_1\) satisfying
\[
 \begin{split}
 \mathbb E_t\left\|\int_{\mathcal T}\varphi(y)
             F(\gamma t y)F(t y)^*\dd y\right\|_2^2
               &\le \int |\varphi(y)|^2\dd\lambda_\gamma(y),\\
 \lambda_\gamma(G_1)&\le
     M^{15000 n^3}\frac{D^2}{d\,C_{\rm op}(\Delta/2)^{k-1}} .
 \end{split}
\]
Here \(\|\,\cdot\,\|_2\) on endomorphisms is the Hilbert--Schmidt
norm, and \(\varphi\) is any bounded Borel scalar function.
For every \(t\) in the construction, if \(y\in\mathcal B_\rho\)
and \(S(\gamma t y)\le\rho\), then \(y\in\mathcal T\).
\end{theorem}
\begin{proof}
Write \(Z=Z(\gamma)\), \(\gamma=h_1\Gamma_Z h_2^{-1}\) by
Proposition~\ref{prop:exchanged-frame-section}. By
Proposition~\ref{prop:bessel-two-frames-bound} the entries of
\(Z\) are bounded by \(\rho^2\) in modulus, \(M\ge\rho\),
and \(\Delta\ge M^{-1}\). Set \(T_0=M^{24n}\), \(P_*=M^{30n}\)
and \(c_*=\Delta/2\). Since \(b\) dominates every fixed power
of \(r\) by Proposition~\ref{prop:thick}, \(P_*\ge\max(8r,b^5)\)
and \(T_0\ge8r\) eventually.
Lemma~\ref{lem:stable-open-frame-section} verifies the
neighborhood hypotheses of Theorem~\ref{thm:bessel-section-convolution}.
Use \(\Theta,\mathcal T\) of Lemma~\ref{lem:bessel-unipotent-tube}.
Let \(v\) have uniform probability on \(\Theta\), let \(k_0\)
have independent Haar probability on \(U(k)\), and put
\(\omega=v\iota(k_0)\), \(t=h_2\omega^{-1}\).
These are continuous functions of the two variables and have
values in \(H_0\). If the frame and size conditions on \(y\)
hold, then \(\omega^{-1}y\in\mathcal B_\rho\) since \(\omega F_0=F_0\).
Also \(S(\Gamma_Z\omega^{-1}y)\le\rho\) by the invariance of
\(S\) and \(\gamma t=h_1\Gamma_Z\omega^{-1}\).
Proposition~\ref{prop:bessel-two-frames-bound} gives
\(\|\omega^{-1}y\|_{\rm op}\le T_0\), hence
\(\|v^{-1}y\|_{\rm op}\le T_0\) by compact unitarity.
The first inclusion in Lemma~\ref{lem:bessel-unipotent-tube}
now gives \(y\in\mathcal T\).

For each fixed \(\omega\), use the unitary left transformation
of \(F\) in Theorem~\ref{thm:actual-real-bessel-vector} to remove
\(h_1\) from the first vector in the product and \(t\) from the
second. The matrices and scalar phases removed are independent
of \(y\), and preserve the Hilbert--Schmidt norm. In an orthonormal
basis of \(\pi\), the resulting squared norm is
\[
       \sum_{j=1}^D \left\|\int \varphi_j(y)
                         F(\Gamma_Z\omega^{-1}y)\dd y\right\|^2 ,
 \qquad \varphi_j={\bf1}_{\mathcal T}\varphi\,\overline{F_j}.
\]
Each test is bounded, supported on a compact set where
\(\|y\|_{\rm op}\le T_0^{12}\le P_*^{20}\).
Also \(\|\omega\|_{\rm op}=\|v\|_{\rm op}\le T_0^{12}\),
since \(\Theta\subset\mathcal T\). By
Lemma~\ref{lem:frame-point-stabilizer} integrating this sum
in the probability law is integration in the indicated subset
of \(H_0\) divided by \(\operatorname{vol}\Theta\).
Theorem~\ref{thm:bessel-section-convolution} term by term thus
proves the first inequality with
\[
  \dd\lambda_\gamma(y)=
     \frac{P_*^{430 n^2} D}{d C_{\rm op}c_*^{\,k-1}
                               \operatorname{vol}\Theta}\,
             {\bf1}_{\mathcal T}(y)\|F(y)\|^2\dd y .
\]
To bound its mass, if \(y,n_0 y\in\mathcal T\) with
\(n_0\in\mathcal U\), then
\(\|A(n_0)\|_{\rm op}\le\|n_0\|_{\rm op}\le P_*^{40}\).
The first inequality uses \(n_0F_0=F_0A(n_0)\) and the
Euclidean orthonormal columns of \(F_0\). The second uses the
two norm bounds and \(\|y^{-1}\|_{\rm op}=\|y\|_{\rm op}\).
Proposition~\ref{prop:frame-multiplicity-bound} and
Lemma~\ref{lem:bessel-unipotent-tube} now give
\[
     \int_{\mathcal T}\|F(y)\|^2\dd y
       \le 2\operatorname{vol}\Theta
               (\pi P_*^{80})^{n(n-1)/2} D
       \le \operatorname{vol}\Theta\,P_*^{42n^2}D .
\]
Plugging this into the measure and using
\(P_*^{472n^2}=M^{14160n^3}\le M^{15000n^3}\)
proves its asserted mass bound.
\end{proof}


\section{Support at the last real pivot}

\begin{definition}\label{def:real-pivot-representative}
Use \(F_0,L_0,G_1,\iota,P_0,z,u\) of
Definitions~\ref{def:real-bessel-graph} and~\ref{def:open-bessel-frame}.
For \(p=(p_1,\ldots,p_n)\in(\C^\times)^n\), put
\(a(c)=c/\min(1,|c|)\) and \(s(c)=\sqrt{|a(c)|^2-1}\).
The matrix \(w(p)\) sends \(e_j\) to \(p_j e_j^*\) and \(e_j^*\)
to \(e_j/\overline{p_j}\) for \(1\le j\le\nu\), fixing the middle
space. On the ordered orthonormal pair \(u,z\) put
\[
 B(c)=\begin{pmatrix}\overline{a(c)}&s(c)\\-s(c)&-a(c)\end{pmatrix}.
\]
Extend \(B(c)\) by the identity on the other coordinates, and put
\(g_p=w(p)B(p_n)\).
\end{definition}

\begin{lemma}\label{lem:real-middle-double-coset}
The map \(p\mapsto g_p\) is continuous into \(G_1\), and
\(Z(g_p)=\operatorname{diag}(p_1,\ldots,p_\nu,a(p_n))\).
If \(\sigma\in G_1\) fixes all the \(e_j,e_j^*\) and
\(z^*J\sigma z=c\), then \(|c|\ge1\) and
\(\sigma=\iota(k_1)B(c)\iota(k_2)\) for some
\(k_1,k_2\in U(P_0)\).
\end{lemma}
\begin{proof}
In a hyperbolic pair the mixed pairing after applying \(w(p)\)
is \(\overline{p_j}/\overline{p_j}=1\) and both diagonal pairings
are zero. The two columns of \(B(c)\) have squared signed norms
\(1,-1\) and zero mixed pairing since \(|a(c)|^2-s(c)^2=1\).
The diagonal frame pairings are now as stated and the coordinates
are continuous everywhere on the indicated domain.

Since \(\sigma\) preserves the orthogonal complement of the fixed
hyperbolic pairs, write \(\sigma z=v-cz\), \(v\in P_0\). Its signed
norm gives \(\|v\|^2=|c|^2-1\). If \(v\ne0\), take \(k_1\) sending
\(u\) to \(v/\|v\|\), by completing both unit vectors to orthonormal
bases by successive orthogonal projections. If \(v=0\), take
\(k_1=1\). Then \(B(c)^{-1}\iota(k_1)^{-1}\sigma\) fixes \(z\)
and all hyperbolic coordinates. On their orthogonal complement
\(P_0\) it is an ordinary unitary \(k_2\), giving the factorization.
\end{proof}

\begin{proposition}\label{prop:top-pivot-support}
Put \(R=b^{6n+2}\) and use \(S\) of
Theorem~\ref{thm:actual-bessel-truncation}. Suppose
\(S(y),S(g_p^{-1}y)\le R\). Write \(d_j=\prod_{a=1}^jp_a\)
for \(0\le j\le\nu\), with \(d_0=1\). Then
\[
 |d_j|\le R^{2j}\quad(1\le j\le\nu),\qquad
 |d_\nu|\max(1,|p_n|)\le R^{2n}.
\]
Moreover, with maxima over an empty index set omitted, put
\[
 L=\max\bigl(\{1\}\cup\{R^{2j}/|d_j|:1\le j\le\nu\}\bigr),
 \quad
 H=\max\bigl(\{1\}\cup\{|p_j|^{-1}:1\le j\le\nu\}\bigr).
\]
Then \(\|y\|_{\rm op}\le1+2\sqrt n(1+H)R(1+nL)^n\).
\end{proposition}
\begin{proof}
Since \(S\) is one plus an infimum, choose \(t,h\in\mathcal U\) with
\[
          \max(\|ty\|_{\rm op},\|hg_p^{-1}y\|_{\rm op})\le R .
\]
Indeed the compact part of each element in the defining infimum
is ordinary unitary, as in Theorem~\ref{thm:actual-bessel-truncation}.
For every \(T\in G_1\), \(T^{-1}=JT^*J\), so the norms of \(T\)
and its inverse agree. Therefore \(m=(ty)(hg_p^{-1}y)^{-1}=t g_p h^{-1}\)
has norm at most \(R^2\). The frame columns are Euclidean orthonormal,
so \(X=Z(m)\) has norm at most \(R^2\). By
Lemma~\ref{lem:bessel-frame-action},
\[
 X=A(t)^{-*}\operatorname{diag}(p_1,\ldots,p_\nu,a(p_n))A(h)^{-1}.
\]
The right multiplication law here follows directly by applying
\(h^{-1}\) to \(F_0\). Each leading determinant thus gives the
corresponding product of diagonal entries. Every \(j\)-row minor
of \(X\) has absolute value at most \(R^{2j}\): project the
selected columns to the selected coordinates, then successively
remove their projections onto earlier columns; the absolute
determinant is the product of the residual lengths. This proves
the two product bounds.

Write the displayed factorization as \(X=L_0' D_0' U_0'\).
For \(i>j\), the entry \((L_0')_{ij}\) equals the determinant of
the submatrix on rows \(1,\ldots,j-1,i\) and columns \(1,\ldots,j\)
divided by \(d_j\). Indeed only the first \(j\) inner indices
contribute because \(U_0'\) is upper triangular, and the remaining
lower-factor determinant is its last diagonal entry in this minor.
The transposed argument gives \((U_0')_{ji}\) via the
column-replaced minor with the same denominator. All these entries
have modulus at most \(L\), involving only \(j\le\nu\).
For either triangular factor \(I+N\), strict triangularity gives
\(N^n=0\) and \(\|N\|_{\rm op}\le nL\) by entrywise
Cauchy--Schwarz. The finite geometric inverse then has norm at
most \((1+nL)^n\). This bounds both \(A(t)\) and \(A(h)\).

Now \(y^{-1}F_0=(ty)^{-1}F_0A(t)\) and
\(y^{-1}g_pF_0=(hg_p^{-1}y)^{-1}F_0A(h)\) both have norm at most
\(C=R(1+nL)^n\). Hence the columns \(y^{-1}e_j,y^{-1}z\) have
norm at most \(C\), and the columns \(y^{-1}e_j^*\) have norm at
most \(HC\) because \(g_pe_j=p_je_j^*\). The negative eigenspace
of \(J\) has orthonormal basis \(z,(e_j-e_j^*)/\sqrt2\).
For its ordinary orthogonal projection \(P_-\) this yields
\(\|y^{-1}P_-\|_{\rm op}\le\sqrt n(1+H)C\), by Cauchy--Schwarz
in these \(n\) coordinates. Finally \(J=I-2P_-\) and
\(y^{-1}J(y^{-1})^*=J\) give
\(\|y^{-1}\|_{\rm op}^2\le1+2\|y^{-1}P_-\|_{\rm op}^2\).
Taking square roots and using the inverse norm identity proves
the bound.
\end{proof}


\section{The real kernel at a common last pivot}

\begin{definition}\label{def:real-top-correlation}
Let \(V\) be the cut function of Theorem~\ref{thm:actual-bessel-truncation}
and \(D=\dim\pi\). For \(g\in G_1\) for which the entries are
absolutely integrable, set
\[
       T(g)=\int_{G_1}V(y)V(g^{-1}y)^*\dd y\in\operatorname{End}(\pi).
\]
Here \(v w^*\) carries \(h\) to \(v\langle w,h\rangle\), with
the inner product conjugate-linear in its first entry. Put
\[
 \begin{split}
 \omega(p)&=\frac{\sqrt{C_{\rm op}}}{D}
       \max(0,|p_n|^2-1)^{(k-1)/2}
           \prod_{j=1}^{\nu}|p_j|^{k-1+2(n-j)},\\
 \Xi(p)&=\omega(p)T(g_p),\qquad
 J_n(p,q)=\operatorname{tr}\bigl(\Xi(p)^*\Xi(q)\bigr).
 \end{split}
\]
Here \(p,q\in(\C^\times)^n\), \(g_p\) is from
Definition~\ref{def:real-pivot-representative}, \(k=\dim P_0\ge2\)
at the scales of Proposition~\ref{prop:actual-bessel-weights}, and
\(C_{\rm op}\) is as in Theorem~\ref{thm:bessel-cell-exact-density}.
\end{definition}

\begin{theorem}\label{thm:real-top-gram-kernel}
The integral \(T(g_p)\) exists for every \(p\).
The maps \(\Xi\) and \(J_n\) are continuous. Each finite matrix
\((J_n(p_i,p_j))_{i,j}\) is positive semidefinite.
If \(p_n=q_n=c=z^*J\sigma z\) for an isometry \(\sigma\in G_1\)
fixing each \(e_j,e_j^*\), one has
\[
       J_n(p,q)=\omega(p)\omega(q)
        \operatorname{tr}\bigl(T(w(p)\sigma)^*T(w(q)\sigma)\bigr).
\]
For fixed \(p_1,\ldots,p_\nu,q_1,\ldots,q_\nu\), put
\(p'=(p_1,\ldots,p_\nu,c)\), \(q'=(q_1,\ldots,q_\nu,c)\).
The function \(c\mapsto J_n(p',q')\) extends by zero at \(c=0\)
to a continuous function on \(\C\), supported in
\[
      1\le |c|\le
         R^{2n}/\max\left(\left|\prod_{j=1}^{\nu}p_j\right|,
                            \left|\prod_{j=1}^{\nu}q_j\right|\right).
\]
\end{theorem}
\begin{proof}
On any compact subset of pivot space all \(|d_j|^{-1}\) for \(j\le\nu\)
and \(|p_j|^{-1}\) for \(j\le\nu\) are bounded. By
Proposition~\ref{prop:top-pivot-support} a single closed group norm ball
contains every \(y\) for which the integrand is nonzero, uniformly
over this subset. Such a ball is compact because the ambient
matrix space is finite dimensional, \(G_1\) is closed there by
\(g^*Jg=J\), and every solution of that equation is invertible.
It has finite measure by the smooth density of
Lemma~\ref{lem:unitary-cayley-density} in finitely many translated
charts. The integrand is continuous in \(p,y\) by
Lemma~\ref{lem:real-middle-double-coset} and bounded in norm by
\(\sup_y\|V(y)\|^2\), which is finite by
Theorem~\ref{thm:actual-bessel-truncation}. Dominated convergence
in Lemma~\ref{lem:integral-tools} now proves existence and continuity
of \(T(g_p)\). The prefactor is continuous, including at \(|p_n|=1\)
since \(k\ge2\). By Corollary~\ref{cor:bessel-pivot-weights},
\(D^2\omega^2\) is exactly the density multiplying Lebesgue
measure in the pivot and unitriangular coordinates on \(|p_n|>1\).

The trace is a finite sum of products of entries, so \(J_n\) is
continuous. For scalars \(t_i\) its quadratic form is
\(\|\sum_i t_i\Xi(p_i)\|_2^2\), where \(\|\cdot\|_2\) is the
Hilbert--Schmidt norm; this gives the positivity including all
repetitions of pivots. Write \(\sigma=\iota(k_1)B(c)\iota(k_2)\)
by Lemma~\ref{lem:real-middle-double-coset}. Since \(w(p)\)
commutes with these compact matrices,
\(w(p)\sigma=\iota(k_1)g_p\iota(k_2)\). The left compact action
on \(V\) from Theorems~\ref{thm:actual-real-bessel-vector}
and~\ref{thm:actual-bessel-truncation} is \(\pi\): the simple-root
coordinates of a pure compact element vanish. Substitution
\(y=\iota(k_1)y'\) in the defining integral gives
\[
   T(w(p)\sigma)=\pi(k_1)T(g_p)\pi(k_2).
\]
This proves absolute integrability there. Apply the same formula
to \(q\); both outer factors are unitary, giving the trace identity.
Finally the prefactors vanish for \(|c|\le1\). Outside the stated
outer radius one of the two support intersections is empty by
Proposition~\ref{prop:top-pivot-support}. These statements give the
claimed extension, compact support and hence integrability on \(\C\).
\end{proof}


\section{The measure of a pair of real Bessel orbits}

\begin{theorem}\label{thm:real-double-pivot-integration}
Let \(S_0=\{\iota(k_0)t:k_0\in U(P_0),\ t\in\mathcal U\}\).
Use the product measure \(\dd s_0=\dd k_0\,\dd t\) given by
Definition~\ref{def:real-bessel-graph} on its first projection.
For every nonnegative Borel function \(h\) on \(G_1\) one has
\[
 \int_{G_1}h(g)\dd g
   =D^2\int_{(\C^\times)^n}\omega(p)^2
     \int_{S_0^2} h(s_1g_p s_2^{-1})\dd s_1\dd s_2\,\dd p .
\]
Here \(\dd p\) is ordinary Lebesgue measure in the \(n\) complex
entries. For \(|p_n|\le1\) the integrand on the right is assigned
zero. The same identity holds for integrable complex \(h\).
\end{theorem}
\begin{proof}
First the set where a leading minor of \(Z(g)\) or \(\Delta(Z(g))\)
vanishes has measure zero. On the full-measure Cayley chart of
Lemma~\ref{lem:unitary-cayley-density}, transported to \(J\) by an
ordinary unitary change of basis, the product of the squared
absolute values of the leading minors times \(\Delta\) is a rational
function of real coordinates with a nowhere vanishing denominator.
Its numerator is a nonzero real polynomial. Indeed take \(g_p\) with
all the first \(\nu\) pivots one and \(p_n=2\). By
Lemma~\ref{lem:real-middle-double-coset} the product is nonzero,
and replacing \(g_p\) by a scalar circle multiple preserves its
nonvanishing. Some such multiple lies in the chart, as only finitely
many scalars can give eigenvalue \(-1\). A nonzero polynomial on a
real coordinate space has null zero set by the successive slicing
argument in Corollary~\ref{cor:bessel-pivot-weights}, which applies
unchanged to real variables. The chart density is absolutely
continuous, proving the assertion.

Whenever \(Z=Z(g)\) is invertible, put \(v=gF_0Z^{-1}e_n\).
Then \(F_0^*Jv=e_n\), so the \(e_j^*\)-coordinates of \(v\) are
zero and its \(z\)-coordinate is \(-1\). Computing its signed
norm in two ways gives
\[
      \|\operatorname{proj}_{P_0}v\|^2-1
             =v^*Jv=-|(Z^{-1})_{nn}|^2 .
\]
Since \((Z^{-1})_{nn}\) is the penultimate leading determinant
divided by \(\det Z\), this gives \(\Delta(Z)\ge0\).
Consequently \(\Delta>0\) and every leading minor is nonzero on
a set of full Haar measure.

Use the decomposition \(Z=l^*\operatorname{diag}(p_1,\ldots,p_n)u'\)
of Corollary~\ref{cor:bessel-pivot-weights} on this set. For
\(|p_n|>1\), the representative
\(s(l)^{-1}g_p s(u')\), where \(s(\cdot)\) is the section of
Lemma~\ref{lem:bessel-frame-action}, has precisely this \(Z\).
By Proposition~\ref{prop:exchanged-frame-section}, it and \(\Gamma_Z\)
belong to the same double coset of \(H_0\). The inner integral
in Theorem~\ref{thm:bessel-cell-exact-density} is independent of
the choice of representative there, by left and right invariance
from Lemma~\ref{lem:frame-point-stabilizer}. This replacement has
a measurable integrand since the new representative is continuous
in its coordinates.

In the resulting formula the first factor on the left of \(g_p\)
is \(h_1s(l)^{-1}\); the factor on its right is the inverse of
\(h_2s(u')^{-1}\), where \(h_1,h_2\in H_0\).
As \(h,l\) vary, \(h s(l)^{-1}\) has measure \(\dd s_0\).
Indeed \(h=v_0\iota(k_0)\) with \(v_0\in\mathcal U_0\).
Conjugation by \(k_0\) preserves its measure in the \(B,T\)
coordinates of Lemma~\ref{lem:frame-point-stabilizer}. At
\(A=l^{-1}\) the lift \(s(l)^{-1}\) differs from \(s(A)\)
on the left by an element of \(\mathcal U_0\). Right invariance
of the measure on this kernel absorbs the difference. Inversion
preserves the entry measure on \(N_n\): in the finite geometric
inverse each entry has leading term minus itself plus a polynomial
in entries of smaller distance above the diagonal. The asserted
measure is now the disintegration of
Lemma~\ref{lem:bessel-frame-action} and compact probability on
\(U(P_0)\). The same reasoning applies on the right. The remaining
pivot density is \(D^2\omega(p)^2\) by
Corollary~\ref{cor:bessel-pivot-weights}. This proves the identity
for nonnegative functions. The four nonnegative parts of a complex
integrable function give the final assertion.
\end{proof}

\begin{corollary}\label{cor:relative-pivot-pair}
For a nonnegative Borel function \(h(y,z)\) invariant under
independent left multiplications by \(S_0\) in its two arguments,
\[
 \iint h(y,z)\dd\bar y\,\dd\bar z
     =D^2\int_{|p_n|>1}\omega(p)^2
            \int_{G_1}h(y,g_p^{-1}y)\dd y\,\dd p .
\]
The integrations in the first \(\nu\) pivots are over
\((\C^\times)^\nu\). Absolute integrability on either side
permits the same formula for complex \(h\).
\end{corollary}
\begin{proof}
Choose the compact-invariant density \(a\) of
Lemma~\ref{lem:bessel-quotient-density}. Then
\(\int_{S_0}a(sy)\dd s_0(s)=1\) for every \(y\).
On the left insert \(a(y)a(z)\) against ordinary Haar measures.
At fixed \(y\) change \(z\) to \(g=yz^{-1}\); inversion preserves
the bi-invariant density by the derivative at the identity, and
left translation preserves it by Lemma~\ref{lem:unitary-cayley-density}.
Use Theorem~\ref{thm:real-double-pivot-integration} in \(g\).
Set \(y=s_1 y'\), giving \(z=s_2g_p^{-1}y'\).
The independent invariances of \(h\) remove \(s_1,s_2\).
The two density integrals are one by the orbit identity for \(a\).
Lemma~\ref{lem:product-integration} justifies all changes of order
for these nonnegative integrands. Applying this conclusion first
to the modulus gives the integrable complex version.
\end{proof}


\section{Haar measure in successive real floors}

\begin{definition}\label{def:successive-real-floors}
Number the hyperbolic pairs \(1,\ldots,\nu\) from outside to
inside. Write \(V_j\) for the space remaining after removing
the first \(j-1\) pairs, and \(G_j\) for its isometries acting
as the identity on the removed pairs. Thus \(V_n=P_0\oplus\C z\).
For \(j<n\) let \(\eta_j\) be the measure on the kernel of the
\(j\)-th layer in Lemma~\ref{lem:heisenberg-layer}, given by
Lebesgue measure on its complex row and imaginary central entry.
Let \(w_j(c)\) be \(w(p)\) acting on the \(j\)-th pair only with
parameter \(c\). Put \(d_j=k-1+2(n-j)=\dim_\C V_{j+1}\) and set
\[
 \alpha_j=(c\mapsto w_j(c)^{-1})_*(|c|^{2d_j}\dd c),\qquad
 \theta_j=\alpha_j*\eta_j \quad(1\le j<n).
\]
Here \(*\) is the image of the product of measures under group
multiplication in the order written. Write \(\kappa\) for probability
on \(\iota(U(P_0))\). Define
\[
 \begin{split}
 \alpha_n&=(c\mapsto B(c)^{-1})_*
       \left(\frac{u_k}{u_{k-1}}(|c|^2-1)^{k-1}
                                      {\bf1}_{|c|>1}\dd c\right),
       \qquad\theta_n=\alpha_n*\kappa,\\
 \mu_{\rm up}&=\kappa*\theta_n*\theta_{n-1}*\cdots*\theta_1,\\
 \mu_{\rm all}&=\kappa*\theta_n*
               (\eta_{n-1}*\theta_{n-1})*\cdots*(\eta_1*\theta_1).
 \end{split}
\]
All pushforwards can be read by integrating nonnegative Borel
tests on the indicated products with the displayed densities.
\end{definition}

\begin{theorem}\label{thm:exact-real-floor-haar}
Put \(\lambda=u_r/u_k\), where \(u_j\) are the compact volumes
in Definition~\ref{def:unitary-density}. The measures of
Definition~\ref{def:successive-real-floors} satisfy
\[
       \mu_{\rm all}=\lambda\,\dd g,\qquad
       \dd t*\mu_{\rm up}=\mu_{\rm all}.
\]
Here \(\dd t\) is Haar measure on \(\mathcal U\) in
Proposition~\ref{prop:unipotent-covol}. Thus every nonnegative
Borel function \(A\) invariant under left \(\mathcal U\) satisfies
\(\int A\,\dd\mu_{\rm up}=\lambda\int A\,\dd\bar y\).
\end{theorem}
\begin{proof}
Conjugation by \(h\in G_{j+1}\) preserves \(\eta_j\). Indeed the
row in Lemma~\ref{lem:heisenberg-layer} becomes \(l h^{-1}\) and
the corner entry \(t\) stays fixed. The absolute real determinant
of \(l\mapsto l h^{-1}\) is \(|\det h|^{-2}=1\), since taking
determinants in the hermitian identity gives \(|\det h|=1\).
Also \(w_j(c)\) commutes with \(G_{j+1}\) since each map fixes
the space on which the other acts.

For fixed pivots, integrate the defining product for \(\mu_{\rm all}\)
from inside to outside. If the inner product has current value
\(h\in G_{j+1}\), the next expression is
\(h v_j w_j(c_j)^{-1} v'_j\) with independent measures \(\eta_j\)
on \(v_j,v'_j\). Move \(v_j\) to the left of \(h\) using the
conjugation just proved. By induction the integrand in the
inner group has argument \(t_{\rm in}g_{\rm in}t'_{\rm in}\)
where the middle factor is the product of the inner Weyl matrices
and \(B(c_n)^{-1}\), and the others are the two inner unipotent
products with compact factors. Since \(w_j(c_j)\) commutes with
\(t'_{\rm in}\), this step inserts its inverse into the middle
product and appends the new independent unipotents to the outer
left and right. The initial inner expression is
\(k_1 B(c_n)^{-1} k_2\) with the two compact probabilities.
This proves by induction that after all steps the two outer
factors have measure \(\dd s_0\) on \(S_0\). Indeed the left
order is outside to inside followed by compact probability, and
the right order is its reverse; the successive semidirect
decomposition of Lemma~\ref{lem:heisenberg-layer} with the
conjugation invariance above gives precisely the product measure
of Proposition~\ref{prop:unipotent-covol}. The compact group
preserves this measure by the same conjugation. Inversion also
preserves \(\dd s_0\), by inversion on the compact group and
the unipotent group together with that invariance.

The product in the middle is \(g_{\bar c}\), where conjugation
is componentwise. This follows from \(w_j(c)^{-1}=w_j(\bar c)\)
and \(B(c)^{-1}=B(\bar c)\) by multiplying the matrices of
Definition~\ref{def:real-pivot-representative}. Coordinatewise
conjugation preserves the pivot densities. Their product is
\(D^2\omega(c)^2(u_k/u_{k-1})/C_{\rm op}\), and this last
constant ratio is \(u_r/u_k\) by
Theorem~\ref{thm:bessel-cell-exact-density}. Comparison with
Theorem~\ref{thm:real-double-pivot-integration} proves the first
identity.

For the second identity move only the left \(v_j\)'s to the front
of the product for \(\mu_{\rm all}\). At each step all original
factors preceding \(v_j\) belong to \(G_{j+1}\), so the same
change of variable applies at fixed preceding factors. Starting
with \(j=1\) and continuing inward leaves the front product
\(v_1\cdots v_\nu\) with measure \(\dd t\); the remaining expression
is exactly the defining product for \(\mu_{\rm up}\).
These integrations all use Lemma~\ref{lem:product-integration}.
Finally apply the two measure identities to \(a(y)A(y)\), with
\(a\) of Lemma~\ref{lem:bessel-quotient-density}. On the left,
\(A(ty)=A(y)\) and \(\int a(ty)\dd t=1\); on the right one obtains
\(\lambda\int a(y)A(y)\dd y\).
\end{proof}


\section{The real energy on sets of finite measure}

\begin{proposition}\label{prop:truncated-pivot-energy}
Let \(V,F\) be as in Theorem~\ref{thm:actual-bessel-truncation}.
For \(B=V\) or \(F\) and \(x\in G_1\) set
\[
 f_x^B(y)=\langle V(y),B(yx)\rangle,\qquad
 c_B(x)=\int f_x^B(y)\dd\bar y .
\]
For every Borel \(E\subset G_1\) of finite Haar measure define
\[
 H_E(p)=\omega(p)^2\int_E\int_{G_1}
        f_x^V(y)\overline{f_x^V(g_p^{-1}y)}\dd y\,\dd x .
\]
On \(|p_n|\le1\) take \(H_E=0\). Then \(H_E\) is integrable
and
\[
                \int H_E(p)\dd p=D^{-2}\int_E |c_V(x)|^2\dd x .
\]
For fixed \(p\), let \(C(p)\) be the upper bound for \(\|y\|_{\rm op}\)
in Proposition~\ref{prop:top-pivot-support}. If
\(E\supset\{x:\|x\|_{\rm op}\le C(p)^2\}\), then
\(H_E(p)=\|\Xi(p)\|_2^2\).
\end{proposition}
\begin{proof}
The functions \(f_x^B\) are invariant under \(S_0\) on the left by
the common unitary transformation in
Theorems~\ref{thm:actual-real-bessel-vector}
and~\ref{thm:actual-bessel-truncation}. Their quotient integrals in
absolute value are at most \(D\): use Cauchy--Schwarz, the squared
norm bounds in the latter theorem and right invariance in
Lemma~\ref{lem:bessel-quotient-density}. For fixed \(p\) the
integrand defining \(H_E\) can be nonzero only when both \(y\)
and \(yx\) satisfy the two support conditions of
Proposition~\ref{prop:top-pivot-support}. Thus \(y\) belongs to
a compact group norm ball and so does \(x=y^{-1}(yx)\), the
latter with radius \(C(p)^2\). Boundedness gives absolute
integrability for each \(p\).

Apply Corollary~\ref{cor:relative-pivot-pair} for each \(x\) to
\(h(y,z)=f_x^V(y)\overline{f_x^V(z)}\). Applied to \(|h|\),
the same identity bounds the integral, including its pivot
density \(D^2\omega^2\), by \(D^2\) uniformly in \(x\).
Integration on \(E\) thus permits all changes of order by
Lemma~\ref{lem:product-integration}. This gives the stated
energy identity. Under the last containment, integrate over
all \(x\) for this particular \(p\). Set \(y'=yx\), so
\(\dd x=\dd y'\) at fixed \(y\). Expanding in an orthonormal
basis the resulting two integrals factor coordinate by coordinate
to give \(\omega(p)^2\sum_{i,j}|T(g_p)_{ij}|^2\) by
Definition~\ref{def:real-top-correlation}. Both \(y\) and \(y'\)
are confined to the same compact ball by the support condition.
\end{proof}

\begin{theorem}\label{thm:finite-real-energy-bound}
Use the preceding notation and \(d>0\) of
Theorem~\ref{thm:bessel-orthonormal-schur}. One has
\[
   \int_{G_1}|c_F(x)|^2\dd x\le D^2/d,\qquad
               |c_V(x)-c_F(x)|\le D b^{-5r^2}\quad(x\in G_1).
\]
For every Borel \(E\) of finite Haar measure \(|E|\),
\[
   D^{-1}\left(\int_E|c_V(x)|^2\dd x\right)^{1/2}
                   \le d^{-1/2}+b^{-5r^2}|E|^{1/2}.
\]
In particular \(D^{-2}\int_E|c_V|^2\le(1+b^{-10})/d\) whenever
\(b\ge3\) and \(d|E|\le b^{10r^2-22}\).
\end{theorem}
\begin{proof}
The previous proof gives \(|c_F|\le D\). For a compact group norm
ball \(A\) use \(\varphi={\bf1}_A\overline{c_F}\) in
Theorem~\ref{thm:bessel-orthonormal-schur}. It belongs to \(L^1\)
and \(L^2\) since the ball has finite measure. Pair the resulting
convolution with \(V\) on the quotient. The integral of the
absolute scalar integrand is bounded by \(D\|\varphi\|_1\),
as follows from the same quotient Cauchy--Schwarz estimate.
Lemma~\ref{lem:product-integration} identifies the pairing with
\(\int_A|c_F|^2\). Quotient Cauchy--Schwarz and the Schur bound give
\[
       \int_A|c_F|^2\le\sqrt D\,(D/d)^{1/2}
                                      \left(\int_A|c_F|^2\right)^{1/2}.
\]
This bounds the integral by \(D^2/d\), also if it vanishes.
Increasing the balls to the whole group and applying
Lemma~\ref{lem:integral-tools} gives the first assertion.
For the second use \(\|V\|_2\le\sqrt D\), right invariance and
\(\|V-F\|_2\le b^{-5r^2}\sqrt D\) on the quotient by
Theorem~\ref{thm:actual-bessel-truncation}. The triangle inequality
on \(E\), obtained by expanding the square and applying
Cauchy--Schwarz to the mixed term, now gives the finite-measure
bound. In the indicated range the right side is at most
\(d^{-1/2}(1+b^{-11})\), whose square is at most \((1+b^{-10})/d\)
since \(2+b^{-11}\le b\).
\end{proof}


\section{Real floor quadrics and their measures}

\begin{definition}\label{def:real-flat-quadric}
For \(1\le j\le n\) use the space \(V_j\) of
Definition~\ref{def:successive-real-floors}. Put \(h=n-j\) and
\(m=k-1\). In \(P_0\oplus\C z\) replace the two coordinates
\(u,z\) by the hyperbolic pair
\((u+z)/\sqrt2,(u-z)/\sqrt2\). This change is ordinary unitary.
Write a vector in \(V_j\) as \((d,t,A,w,B)\), with \(d,t\) the
new pair, \(A,B\in\C^h\) on the remaining hyperbolic pairs and
\(w\in\C^m\) on the positive complement. Its signed norm is
\[
             Q(v)=2\operatorname{Re}(\bar t d+A^*B)+\|w\|^2 .
\]
For every \(\gamma\in\R\), define \(\delta_\gamma^{(j)}\) as the
image of the product measure \(\dd t\,\dd A\,\dd w\,\dd B\,\dd a\)
for \(t\in\C^\times,\ a\in\R\) under the map retaining
\(t,A,w,B\) and putting
\[
               d=\frac{\gamma-2\operatorname{Re}(A^*B)-\|w\|^2}
                              {2\bar t}+ita .
\]
Here complex coordinates have ordinary two-dimensional Lebesgue
measure. Each such image is supported on \(Q=\gamma\).
\end{definition}

\begin{lemma}\label{lem:real-flat-disintegration}
For compactly supported continuous \(\phi\) on \(V_j\), the function
\(\gamma\mapsto\int\phi\,\dd\delta_\gamma^{(j)}\) is continuous.
For every nonnegative Borel function \(B_1\) on \(\R\times V_j\),
\[
  \int_{V_j} B_1(Q(v),v)\dd v
       =\frac12\int_\R\int_{V_j} B_1(\gamma,v)
                                       \dd\delta_\gamma^{(j)}(v)\dd\gamma .
\]
The measure at every \(\gamma\) is invariant under \(G_j\);
the same flat formula may be used with any one of the displayed
hyperbolic pairs as its distinguished pair.
\end{lemma}
\begin{proof}
At fixed \(t\ne0,A,w,B\), write \(d=t(x+ia)\) with \(x,a\in\R\).
The ordinary Jacobian is \(|t|^2\) and
\(\dd\gamma=2|t|^2\dd x\). The discarded set \(t=0\) has zero
ambient measure by Lemma~\ref{lem:product-integration}. These
facts prove the integral identity. If \(\phi\) vanishes outside
the coordinate box of radius \(L\), the free complex coordinates
in its flat formula have moduli \(\le L\) where it is nonzero,
and \(|a|\le L/|t|\) since \(a=\operatorname{Im}(d/t)\).
The integral \(\int_{|t|\le L}|t|^{-1}\dd t=2\pi L\) by
polar coordinates. Thus the parameter integrands are dominated
by the integrable product of these indicators times
\(\|\phi\|_\infty\). Continuity follows from
Lemma~\ref{lem:integral-tools}.

Each \(G_j\)-isometry preserves \(Q\) and has complex determinant
of modulus one. Its real Jacobian equals the squared modulus of
this determinant, for instance by triangularizing a complex matrix
by unitary changes of coordinates as in
Lemma~\ref{lem:unitary-cayley-density}. For continuous compactly
supported \(\phi\), insert \({\bf1}_I(\gamma)\phi(v)\) in the
integral identity for any bounded interval \(I\). After an
isometry change in the ambient integral, the two continuous
functions of \(\gamma\) thus have the same integral on every \(I\).
They agree at every real point by continuity. The same comparison
works on exchanging two hyperbolic coordinate pairs, which
preserves the ambient measure. Equality on continuous compact
tests implies equality of the corresponding locally finite Borel
measures here: on a bounded open set approximate its indicator
from below by the continuous minima of one and a positive integer
times distance to the complement. Monotone convergence and open
exhaustion give equality on opens, and differences and countable
disjoint unions on bounded boxes then give equality on their
generated Borel sets, as in Lemma~\ref{lem:product-integration}.
\end{proof}

\begin{proposition}\label{prop:actual-floor-quadrics}
For \(j<n\), the map \(g\mapsto g^{-1}e_j\) sends \(\theta_j\)
to \(\delta^{(j)}_0\). At \(j=n\), the map \(g\mapsto g^{-1}z\)
sends \(\theta_n\) to \(\delta^{(n)}_{-1}\).
The same assertions hold after applying any fixed member of \(G_j\)
to the resulting vector.
\end{proposition}
\begin{proof}
At \(j<n\) use \(e_j,e_j^*\) as the distinguished pair by
Lemma~\ref{lem:real-flat-disintegration}. In the notation of
Lemma~\ref{lem:heisenberg-layer} a point integrated by \(\theta_j\)
is \(g=w_j(c)^{-1}H(l,t_0)\). The outside coordinates of
\(g^{-1}e_j\) are \(c\bar t_0,c\), and the inner vector is
\(v'=cJ' l^*\), where \(J'\) is the hermitian matrix on \(V_{j+1}\).
This follows by triangular inversion and \(t_0+\bar t_0=-lJ'l^*\).
In particular
\[
     c\bar t_0=-\frac{(v')^*J'v'}{2\bar c}
                         -ic\,\operatorname{Im}t_0 .
\]
Changing the row \(l\) to \(v'\) multiplies its real measure by
\(|c|^{2d_j}\), because \(J'\) is ordinary unitary. This is exactly
the pivot weight in Definition~\ref{def:successive-real-floors}.
Taking \(a=-\operatorname{Im}t_0\) leaves ordinary real measure,
which gives the required flat formula.

At the last floor \(g=B(c)^{-1}\iota(k_0)\) and
\(g^{-1}z=(|c|^2-1)^{1/2}k_0^{-1}u-cz\) on \(|c|>1\).
The vector \(k_0^{-1}u\) has the probability of a uniformly
distributed unit vector by Theorem~\ref{thm:compact-dual-unitary-volume}.
In ordinary coordinates \(v'=v_+-cz\), \(Q(v')=\|v_+\|^2-|c|^2\).
The polar formula of that theorem in the \(k\) positive coordinates,
with \(\gamma=\|v_+\|^2-|c|^2\), gives radial measure
\[
        \frac12\,\frac{u_k}{u_{k-1}}
          (\gamma+|c|^2)^{k-1}{\bf1}_{\gamma+|c|^2>0}
              \dd\gamma\,\dd\omega\,\dd c ,
\]
where \(\dd\omega\) is the unit-vector probability. For compact
continuous tests the integral in \(\omega,c\) is continuous in
\(\gamma\), since \(k\ge2\), \(c\) stays in a bounded set and the
positive radial power tends to zero at the boundary. Comparing
the two disintegrations on every interval as in
Lemma~\ref{lem:real-flat-disintegration} identifies this measure
at \(\gamma=-1\) with the required flat measure (after the common
factor \(1/2\) is canceled). It agrees with \(\theta_n\) by its
definition. All fixed isometries preserve the resulting measures
by the lemma.
\end{proof}


\section{Fourier bounds on a bounded part of a real floor}

\begin{lemma}\label{lem:real-complex-plancherel}
Write \(\psi(t)=\exp(t-\bar t)\) for \(t\in\C\).
For \(m\ge0,\ \xi\in\C^\times,\ H\in L^1(\C^m)\cap L^2(\C^m)\),
put \(\widehat H_\xi(v)=\int\psi(\xi v^*w)H(w)\dd w\). Then
\[
           \int|\widehat H_\xi(v)|^2\dd v
                       \le(\pi/|\xi|)^{2m}\int |H(w)|^2\dd w .
\]
\end{lemma}
\begin{proof}
For \(\epsilon>0\) integrate the squared transform against
\(\exp(-\epsilon\|v\|^2)\). Absolute integration of its expansion
follows from the Gaussian integral and \(\|H\|_1^2<\infty\).
In \(2m\) real coordinates the integral of the Gaussian times
\(\psi(\xi v^*z)\) equals
\((\pi/\epsilon)^m\exp(-|\xi|^2\|z\|^2/\epsilon)\).
Indeed the linear form \(2\operatorname{Im}(\xi v^*z)\) in those
real coordinates has coefficient vector of squared length
\(4|\xi|^2\|z\|^2\); apply the one-dimensional Gaussian transform
proved in Lemma~\ref{lem:theta} to each coordinate.
Use \(z=w-w'\) in the expanded squared transform.
For each \(z\), Cauchy--Schwarz gives
\(\int |H(w')H(w'+z)|\dd w'\le\|H\|_2^2\). The remaining
integral of \(\exp(-|\xi|^2\|z\|^2/\epsilon)\) is
\((\pi\epsilon/|\xi|^2)^m\). This proves the bound uniformly
in \(\epsilon\), and monotone convergence in
Lemma~\ref{lem:integral-tools} as \(\epsilon\downarrow0\)
proves the assertion. For \(m=0\) the measures have mass one.
\end{proof}

\begin{theorem}\label{thm:real-flat-fourier}
In Definition~\ref{def:real-flat-quadric} let \(\mu,\nu'\) each
be either ambient Lebesgue measure or a measure
\(\delta_\gamma^{(j)}\), with arbitrary real level in either slot.
Write \(\mathcal Q_L\) for the box of vectors whose flat coordinate
entries \(d,t,A,w,B\) have modulus at most \(L>0\). Put
\[
 M_{\rm amb}(L)=(\pi L^2)^{2h+2},\qquad
 M_{\rm quad}(L)=4\pi L^2(\pi L^2)^{2h},
\]
and denote the appropriate quantity by \(M_\mu\) or \(M_{\nu'}\).
Take Borel \(f\in L^2(\mu)\) vanishing off \(\mathcal Q_L\), a
shift \(y_0\in V_j\), and \(L'>0,\ \xi\in\C^\times\). Then
\[
 \int_{\mathcal Q_{L'}}
  \left|\int\psi\bigl(\xi [v+y_0,z]\bigr) f(z)\dd\mu(z)\right|^2
           \dd\nu'(v)
 \le M_{\nu'}(L')M_\mu(L)(\pi/|\xi|)^{2m}
                               \int |f|^2\dd\mu .
\]
Here \([v,z]=v^*Jz\) and \(m=k-1\).
\end{theorem}
\begin{proof}
For either measure separate the variables \(p=(t,A)\), \(w\), and
the fibre. The fibre variables are \((B,d)\) for ambient measure
and \((B,a)\) for a quadric. All parameter measures are products
by Definition~\ref{def:real-flat-quadric}. At fixed \(p,w\) the
fibre volume in \(\mathcal Q_L\) is at most \(m_{\mu,L}(p)\).
For ambient measure take \((\pi L^2)^{h+1}\), and for a quadric
take \((\pi L^2)^h 2L/|t|\), in both cases times the indicator
that the entries of \(p\) have modulus at most \(L\). The latter
bound follows from \(|a|\le L/|t|\) in
Lemma~\ref{lem:real-flat-disintegration}. Set its value zero at
\(t=0\), a null set. Direct integration gives
\(\int m_{\mu,L}(p)\dd p=M_\mu(L)\). This argument applies also
to the target measure. In particular each box has finite measure
since the remaining \(w\) range in a finite box. Thus \(f\in L^1(\mu)\).

Use subscripts \(v,z\) to distinguish output and input coordinates.
The pairing in the phase is a sum of three terms:
\[
  \bar d_v t_z+B_v^*A_z,\qquad
  w_v^*w_z,\qquad
  \bar t_v d_z+A_v^* B_z+[y_0,z].
\]
At fixed \(p_v,p_z,w_z\) integrate \(f\) times the phase of the
third term over the input fibre; call the resulting function
\(H_{p_v,p_z}(w_z)\). Its absolute value is bounded by the fibre
integral of \(|f|\). Cauchy--Schwarz on that fibre gives
\[
 |H_{p_v,p_z}(w_z)|^2
      \le m_{\mu,L}(p_z)\int_{\rm fibre}|f(z)|^2 .
\]
In the third term \(d_z\) may depend on \(w_z\); throughout this
fibre integral \(p_v\) is fixed, and the third term is independent
of \(w_v\). Thus the second phase produces precisely
\(\widehat H_{p_v,p_z,\xi}(w_v)\) of
Lemma~\ref{lem:real-complex-plancherel}. The first term has phase
of modulus one depending only on the output and \(p_z\), so the
absolute value of the full integral is at most
\(\int|\widehat H_{p_v,p_z,\xi}(w_v)|\dd p_z\), uniformly in the
output fibre. All these decompositions use
Lemma~\ref{lem:product-integration}.

Apply weighted Cauchy--Schwarz in \(p_z\) to square this last
bound, using weight \(m_{\mu,L}(p_z)\) and omitting its zero set,
where the corresponding integrals vanish almost everywhere.
For each fixed \(p_v\), integration in \(w_v\) over the whole
space and Lemma~\ref{lem:real-complex-plancherel} bound the result by
\[
 \begin{aligned}
 M_\mu(L)(\pi/|\xi|)^{2m}
       \iint\frac{|H_{p_v,p_z}(w_z)|^2}{m_{\mu,L}(p_z)}
                    \dd w_z\dd p_z
       &\le M_\mu(L)(\pi/|\xi|)^{2m}\|f\|_{L^2(\mu)}^2 .
 \end{aligned}
\]
The fibre estimate ensures the Fourier lemma applies for almost
every \(p_z\); \(w_z\) ranges in a bounded box which also gives
the \(L^1\) condition. On the output box integrate the uniform
fibre bound first, at cost \(m_{\nu',L'}(p_v)\), then integrate
in \(p_v\). Its integral \(M_{\nu'}(L')\) gives the claimed factor.
\end{proof}


\section{An exact negative moment for the real floor weights}

\begin{lemma}\label{lem:real-projective-moment}
Let \(m=s+1+c\), \(s,c\ge0\) integers, and let \(\Pi\) be ordinary
orthogonal projection of rank \(s+1\) on \(\C^m\). For \(T,H>0\),
\[
 \int_{\|w\|\le T}\left(\frac{H}{\|\Pi w\|^2}\right)^s\dd w
           =\frac{H^s\pi^m}{s!(c+1)!}T^{2(c+1)} .
\]
Values on the null kernel may be assigned arbitrarily.
\end{lemma}
\begin{proof}
Choose orthonormal bases in the range and kernel, so that
\(w=(w',w'')\) and \(\Pi w=(w',0)\). This ordinary unitary change
preserves real volume. The ball in \(\C^c\) above \(w'\), with
\(u=\|w'\|^2\), has volume \(\pi^c(T^2-u)^c/c!\).
By Theorem~\ref{thm:compact-dual-unitary-volume}, polar integration
in \(\C^{s+1}\) gives \(\pi^{s+1}u^s\dd u/s!\), including the
angular probability. Their product times \(H^s u^{-s}\) integrates
over \(0<u<T^2\). Integrating \((T^2-u)^c\) gives
\(T^{2(c+1)}/(c+1)\), proving the formula, also for \(s=0\) or \(c=0\).
\end{proof}

\begin{proposition}\label{prop:real-two-position-gauge}
Let \(m=k-1=s+1+3n\) and \(s\ge0\). Let \(E_+\) denote the
inclusion of the \(m\) orthonormal positive coordinates of
Definition~\ref{def:real-flat-quadric} into the full space. Let
\(C\) be any complex matrix with \(m\) columns and at most \(n\)
rows, and let \(g,h\in G_1\). Call a pair \((z,z')\) close if
\(\|z\|<1/4\) and \(\|z-z'\|<1/4\), with ordinary Euclidean norms.
Define
\[
  W=\{w\in\ker C:(g^{-1}E_+w,h^{-1}E_+w)\text{ is close}\}.
\]
Let \(|\cdot|\) on subsets of \(\C\) denote ordinary area, with
outer area permitted for arbitrary subsets. There is an
orthogonal projection \(\Pi\) of rank \(s+1\) such that the
function \(\phi(v)=36/\|\Pi v\|^2\), valued \(+\infty\) on the
kernel, has the following properties.
\begin{enumerate}
\item \(\bigl|\{v^*w:w\in W\}\bigr|^{-1}\le\phi(v)\), with
\(1/0=+\infty\).
\item The set of \(\alpha\in\C\) such that for each \(w\in W\)
there exist close pairs \((z_1,z_2),(z_1',z_2')\) with
\[
        (\alpha v)^*w=[z_1,z'_1]-[z_2,z'_2]
\]
has outer area at most \(\phi(v)\).
\item For integers \(j\ge0,\ n\ge48\), and \(L\ge\max(m,4)\),
\[
   \int_{\max_i|v_i|\le L^j}\phi(v)^s\dd v
                                  \le L^{(2j+3)(3n+1)} .
\]
\end{enumerate}
\end{proposition}
\begin{proof}
Write \(P_-\) for the Euclidean projection onto the negative
eigenspace of \(J\), of rank \(n\) by
Proposition~\ref{prop:top-pivot-support}. Stack the maps
\(C,P_-g^{-1}E_+,P_-h^{-1}E_+\). Their combined rank is at most
\(3n\), so their common kernel has dimension at least \(s+1\).
Choose \(\Pi\) onto any \((s+1)\)-plane in it. For \(w\) in this
plane the two images \(g^{-1}E_+w,h^{-1}E_+w\) lie in the positive
eigenspace, where Euclidean and signed squared norms agree.
Since \(g,h\) are isometries, both Euclidean norms equal \(\|w\|\).
Consequently \(\{w\in\operatorname{ran}\Pi:\|w\|<1/8\}\subset W\)
by the triangle inequality.

Put \(a=\|\Pi v\|\). If \(a>0\), each scalar \(t\) with \(|t|<a/8\)
equals \(v^*w\) for \(w=t\Pi v/a^2\) in this ball. Thus the
image has area at least \(\pi a^2/64\). It is measurable since
\(W\) is open in the subspace \(\ker C\) and a linear functional
on that space is either zero or open onto \(\C\).
For any close pair the second vector has norm \(<1/2\).
Since \(J\) is ordinary unitary, the difference of the two pairings
in the second assertion has modulus at most
\[
       \|z_1-z_2\|\|z'_1\|+\|z_2\|\|z'_1-z'_2\|<3/16 .
\]
It follows from the disk just constructed that
\(|\alpha|a/8\le3/16\), hence \(|\alpha|\le3/(2a)\).
Its disk has area \(9\pi/(4a^2)\). Both this and the first
reciprocal-area bound are at most \(36/a^2\), using \(2<\pi<4\).
When \(a=0\) the extended inequalities follow directly.

The entry box lies in the ball of radius \(T=\sqrt m\,L^j\).
Use Lemma~\ref{lem:real-projective-moment} with \(c=3n,H=36\).
From the exponential series \((36\pi)^s/s!\le\exp(36\pi)\).
As \(L\ge4\), \(\log L\ge1\), and \(3n+1\ge145>36\pi\).
Thus \(\exp(36\pi)\le L^{3n+1}\), and
\(\pi^{3n+1}/(3n+1)!\le L^{3n+1}\).
Finally \(T^{2(3n+1)}\le L^{(2j+1)(3n+1)}\).
Multiplication gives the third bound.
\end{proof}


\section{A pair of negative unit vectors at the last pivot}

\begin{lemma}\label{lem:middle-haar-cutoff}
On \(M=G_n\) the measure \(\dd m=\kappa*\theta_n\) of
Definition~\ref{def:successive-real-floors} is bi-invariant and
inversion invariant. For \(t\ge1\), with
\(c_k=\pi u_k/(k u_{k-1})\),
\[
 \begin{split}
 \delta_{-1}^{(n)}\{v:\|v\|\le t\}
                 &=c_k\bigl((t^2-1)/2\bigr)^k,\\
 v_2:=m\{h:\|h\|_{\rm op}\le2\}&=c_k(9/16)^k .
 \end{split}
\]
For \(A\ge1\) write \(B_G(A)=\{g\in G_1:\|g\|_{\rm op}\le A\}\).
There is a nonnegative Borel function \(a_A\le v_2^{-1}\) on
\(G_1\), supported in \(B_G(2A)\), with
\(\int_M a_A(hg)\dd m(h)=1\) for all \(g\in M B_G(A)\).
\end{lemma}
\begin{proof}
For a nonnegative Borel test \(F_1\), first average \(F_1(h)\)
on the left over \(\kappa\). This gives a function only of
\(v=h^{-1}z\), since two representatives of the same \(v\) differ
on the left by a matrix fixing \(z\), hence by \(\iota(U(P_0))\).
The function of \(v\) is Borel. Indeed for \(v=v_+-cz\) with
\(Q(v)=-1\) choose a unitary carrying \(u\) to \(v_+/\|v_+\|\)
when \(v_+\ne0\), and the identity otherwise. A Borel choice is
obtained by appending the coordinate vectors in order, subtracting
their projections to the vectors retained and normalizing each
nonzero residual. The inverse of the resulting compact rotation
times the boost, in the order \((\iota(k_0)B(c))^{-1}\), is a Borel
representative with \(h^{-1}z=v\). Proposition~\ref{prop:actual-floor-quadrics}
identifies the distribution of \(v\) under \(\kappa*\theta_n\)
with \(\delta_{-1}^{(n)}\). Right translation acts on \(v\) by the
inverse isometry, preserving that measure by
Lemma~\ref{lem:real-flat-disintegration}. This proves right
invariance. The product \(\kappa*\alpha_n*\kappa\) is also invariant
under inversion by \(B(c)^{-1}=B(\bar c)\), radiality of its density,
and inversion invariance of compact probability. Thus it is left
invariant as well.

For \(v=v_+-cz\) on the negative unit quadric,
\(\|v\|^2=2|c|^2-1\). In the radial formula of
Proposition~\ref{prop:actual-floor-quadrics} integrate
\((u_k/u_{k-1})(|c|^2-1)^{k-1}\dd c\) from \(|c|=1\) to
\(\sqrt{(t^2+1)/2}\). In polar coordinates this gives the first
formula. The norm of \(B(c)\) on \(|c|>1\) is
\(|c|+\sqrt{|c|^2-1}\), since its two nontrivial squared singular
values are \((|c|\pm\sqrt{|c|^2-1})^2\) by direct multiplication.
Compact factors preserve this norm, which is at most two exactly
when \(|c|\le5/4\). The same radial integration gives \(v_2\).
It also shows that \(m\) is finite on all norm balls.

The set \(M B_G(A)\) is closed. For a convergent sequence \(h_i b_i\)
with \(b_i\in B_G(A)\), compactness gives a subsequence with \(b_i\)
convergent and then \(h_i\) has a limit in the closed subgroup \(M\).
On this set put \(a_A(g)={\bf1}_{B_G(2A)}(g)/N_A(g)\), where
\(N_A(g)=\int_M{\bf1}_{B_G(2A)}(hg)\dd m(h)\), and set \(a_A=0\)
off it. The denominator is left \(M\)-invariant, Borel, and finite
since \(\|h\|_{\rm op}\le2A\|g\|_{\rm op}\) on its integration set.
If \(\|g\|_{\rm op}\le A\) it is at least \(v_2\); invariance
extends this to \(M B_G(A)\). The asserted orbit integral and
all the other bounds follow directly.
\end{proof}

\begin{proposition}\label{prop:last-real-oscillatory-pair}
Fix heads \(p,q\in(\C^\times)^\nu\). Write \(w_p=w((p,c))\),
\(P(p)=\prod_{j=1}^\nu|p_j|^{d_j}\) for \(d_j\) of
Definition~\ref{def:successive-real-floors}, and \(A_p=C((p,c))\)
for the bound in Proposition~\ref{prop:truncated-pivot-energy};
it is independent of \(c\). Put \(X=A_p A_q\),
\(\delta=\delta_{-1}^{(n)}\). For \(g,x\in G_1\) define functions
of \(v=h^{-1}z\) and \(v'=\tau^{-1}z\), \(h,\tau\in M\), by
\[
 \begin{split}
  U_{g,x}(v)&=\langle V(hg),V(hgx)\rangle,\\
  W_{g,x}(v')&=\langle V(\tau w_p^{-1}g),V(\tau w_q^{-1}gx)\rangle .
 \end{split}
\]
For every \(\xi\in\C\), the integral
\(I_\xi=\int_\C\psi(\xi c)J_n((p,c),(q,c))\dd c\) equals
\[
 \frac{u_kP(p)P(q)}{u_rD^2}
 \int_{B_G(X)}\int_{B_G(2A_p)} a_{A_p}(g)
       \iint\psi(\xi[v,v'])U_{g,x}(v)\overline{W_{g,x}(v')}
              \dd\delta(v)\dd\delta(v')\,\dd g\,\dd x .
\]
The quadrics in the last integral may be restricted respectively to
\(\|v\|\le2R A_p\) and \(\|v'\|\le2R A_p\|w_p\|_{\rm op}\).
\end{proposition}
\begin{proof}
The two functions on the quadrics are well defined by the compact
equivariance of \(V\) and the commutation of both Weyl elements
with \(\iota(U(P_0))\). They are Borel by the representative in
Lemma~\ref{lem:middle-haar-cutoff}. By
Theorem~\ref{thm:real-top-gram-kernel} and the middle radial density,
\(I_\xi\) equals \(u_kP(p)P(q)/(u_rD^2)\) times
\[
   \int_M \psi(\xi[z,\sigma z])
       \operatorname{tr}\bigl(T(w_p\sigma)^*T(w_q\sigma)\bigr)
                                                    \dd m(\sigma).
\]
Indeed the middle Haar product integrates a pair of independent
compact factors around \(B(c)^{-1}\). These factors preserve the
displayed trace by that theorem and preserve \([z,\sigma z]\).
The change \(c\mapsto\bar c\) replaces the inverse boost by \(B(c)\).
The ratio of the pivot and radial constants is
\(C_{\rm op}u_{k-1}/u_k=u_k/u_r\).

Expand the trace in its two independent integration variables
\(y,y'\). Its integrand is
\[
  \langle V(y),V(y')\rangle\,
       \overline{\langle V(\sigma^{-1}w_p^{-1}y),
                             V(\sigma^{-1}w_q^{-1}y')\rangle}.
\]
Whenever this product is nonzero,
\(\|y\|_{\rm op}\le A_p,\ \|y'\|_{\rm op}\le A_q\). In fact
the middle double-coset decomposition of
Lemma~\ref{lem:real-middle-double-coset} reduces each pair to
Proposition~\ref{prop:top-pivot-support} by compact translations,
which preserve both norm and \(S\). It also gives
\(|[z,\sigma z]|\le R^{2n}/|\prod_j p_j|\).
Thus \(\sigma\) lies in a norm ball in \(M\), by the boost formula
of Lemma~\ref{lem:middle-haar-cutoff}. The integral of the absolute
value of the expanded product is finite, as all three variables
are bounded and \(V\) is bounded.

Set \(y'=yx\), so \(\|x\|_{\rm op}\le X\) on this support. Insert
\(1=\int_M a_{A_p}(h^{-1}y)\dd m(h)\) there and make the changes
\(y=hg,\ \sigma=h\tau^{-1}\) at fixed \(h\).
The measures are preserved by translation and inversion. Now
\([z,\sigma z]=[h^{-1}z,\tau^{-1}z]\) and \(w_p,w_q\) commute with
\(M\). The two scalar factors become exactly the displayed
\(U_{g,x},\overline{W_{g,x}}\). The change of order is justified
also for their absolute values by the original absolute integral
and the nonnegative insertion, using
Lemma~\ref{lem:product-integration}. Projection from \(h,\tau\)
to their negative unit vectors gives the two copies of \(\delta\)
by Proposition~\ref{prop:actual-floor-quadrics}.

If \(V(hg)\ne0\), some \(t\in\mathcal U\) satisfies
\(\|thg\|_{\rm op}\le R\) by
Theorem~\ref{thm:actual-bessel-truncation}. The vector
\(t^{-1}z\) equals \(z\) plus a linear combination of the outer
\(e_j\) by Lemma~\ref{lem:bessel-frame-action}. Hence
\(\|h^{-1}z\|\le\|h^{-1}t^{-1}z\|\le R\|g\|_{\rm op}\),
where the first inequality is ordinary orthogonal projection
onto \(V_n\). Apply this also to \(\tau w_p^{-1}g\) to bound
the second quadric support. Lemma~\ref{lem:middle-haar-cutoff}
confines \(g\) to the required norm ball, proving the last claims.
\end{proof}


\section{Fourier decay at the last real pivot}

\begin{theorem}\label{thm:actual-top-pivot-fourier}
Use the notation of Proposition~\ref{prop:last-real-oscillatory-pair},
and let \(T_1=2R A_p,\ T_2=T_1\|w_p\|_{\rm op}\).
Write \(b_G(t)=\int_{B_G(t)}\dd g\),
\(v(t)=c_k((t^2-1)/2)^k\) for \(t\ge1\), and
\(M_V=\sup_y\|V(y)\|\).
For every \(\xi\ne0\) one has
\[
 \begin{split}
 |I_\xi|&\le C(p,q)(\pi/|\xi|)^{k-1},\\
 C(p,q)&=\frac{u_kP(p)P(q)}{u_rD^2}
      \frac{b_G(X)b_G(2A_p)}{v_2}\,
      4\pi T_1T_2\,M_V^4\sqrt{v(T_1)v(T_2)} .
 \end{split}
\]
At every \(\xi\), including zero, one also has
\[
 |I_\xi|\le\frac{u_kP(p)P(q)}{u_rD^2}\,
        \frac{b_G(X)b_G(2A_p)}{v_2}\,M_V^4 v(T_1)v(T_2).
\]
\end{theorem}
\begin{proof}
All the constants are finite by
Lemmas~\ref{lem:unitary-cayley-density} and~\ref{lem:middle-haar-cutoff}
and Theorem~\ref{thm:actual-bessel-truncation}.
For fixed \(g,x\) in Proposition~\ref{prop:last-real-oscillatory-pair},
the two amplitudes have modulus at most \(M_V^2\). Their squared
integrals on the quadric are at most \(M_V^4v(T_1)\) and
\(M_V^4v(T_2)\) by their respective supports and
Lemma~\ref{lem:middle-haar-cutoff}. The Euclidean norm balls of
these radii are contained in the corresponding coordinate boxes
of Definition~\ref{def:real-flat-quadric}, because the coordinate
basis is ordinary orthonormal. Pair \(U_{g,x}\) by Cauchy--Schwarz
with the Fourier integral of \(\overline{W_{g,x}}\) on the first
box. Theorem~\ref{thm:real-flat-fourier} with both measures
\(\delta_{-1}^{(n)}\), \(h=0\), \(m=k-1\) and shift zero gives
\[
  \left|\iint\psi(\xi[v,v'])U_{g,x}(v)\overline{W_{g,x}(v')}
                                \dd\delta(v)\dd\delta(v')\right|
   \le (\pi/|\xi|)^{k-1}(4\pi T_1T_2)
                      M_V^4\sqrt{v(T_1)v(T_2)} .
\]
Here the factor \(4\pi T_1T_2\) is the square root of the
product of the two fibre bounds. Integrating in \(x,g\) and
using \(0\le a_{A_p}\le v_2^{-1}\) gives the first assertion.
The direct integral bound by the product of the two \(L^1\) norms
gives the second assertion.
\end{proof}

\begin{corollary}\label{cor:actual-penultimate-gram}
Define \(J_\nu(p,q)=I_0\) on \((\C^\times)^\nu\) in each argument.
This is a continuous positive semidefinite kernel. For all \(p,q\),
the Fourier transform at frequency \(\xi\) of the common-last-entry
function integrating to \(J_\nu(p,q)\) satisfies
Theorem~\ref{thm:actual-top-pivot-fourier}. Its constants are
bounded uniformly when the two heads range in compact sets.
\end{corollary}
\begin{proof}
For a compact set of heads the last variable is confined to a common
bounded disk by Theorem~\ref{thm:real-top-gram-kernel} and the
lower bounds on all first pivot moduli. Its integrand vanishes
for \(|c|\le1\), and elsewhere it is continuous jointly in the
heads and \(c\). It is bounded on this compact range by continuity.
Lemma~\ref{lem:integral-tools} gives continuity of the mean.
For heads \(p_i\) and scalars \(s_i\), the quadratic form of
\(J_\nu\) is \(\int_{\C}\|\sum_i s_i\Xi((p_i,c))\|_2^2\dd c\);
the integral is finite and nonnegative by
Theorem~\ref{thm:real-top-gram-kernel}. The Fourier assertion is
that theorem's common-last-entry function in
Proposition~\ref{prop:last-real-oscillatory-pair}. Its constants
are locally bounded because all \(A_p,A_q,P(p),P(q),\|w_p\|_{\rm op}\)
are bounded on compact head sets, by their explicit formulas in
Proposition~\ref{prop:top-pivot-support} and
Definition~\ref{def:real-pivot-representative}.
\end{proof}


\section{Bounded representatives at partial real flags}

\begin{lemma}\label{lem:bounded-signed-complement}
Let \(J=J^*=J^{-1}\) on a finite-dimensional complex Euclidean
space decomposed as \(E\oplus F\), with both spaces \(J\)-invariant
and Euclidean orthogonal. Let \(T^*JT=J\) and
\(B=\max(1,\|T|_E\|_{\rm op})\). Some isometry \(H\) acting
identically on \(E\) satisfies
\[
              \|TH\|_{\rm op}\le B+\sqrt{1+B^2}.
\]
\end{lemma}
\begin{proof}
Write \(T_E:E\to V\) for \(T|_E\), \(J_E=J|_E\) and
\(\mathcal V=T F\). The projection onto \(T E\) along
\(\mathcal V\) is \(T_EJ_E T_E^*J\). Indeed this map is the
identity on \(T E\) and zero on \(T F\) by \(T^*JT=J\).
Its complementary projection \(Q\) has norm at most \(1+B^2\).
Write \(P\) for ordinary orthogonal projection onto \(\mathcal V\).
The self-adjoint operator \(P J|_{\mathcal V}\) on \(\mathcal V\)
has inverse \(QJ|_{\mathcal V}\). To check the left inverse, for
\(y\in\mathcal V\) compute \(QJPJy=Qy-QJ(I-P)Jy=y\).
Here \(J(\mathcal V^\perp)=T E\), since it is the signed
orthogonal complement of \(T F\). In finite dimension the left
inverse is the inverse. Every eigenvalue of \(P J|_{\mathcal V}\)
therefore has absolute value at least \(1/(1+B^2)\).

Diagonalize this operator in an ordinary orthonormal basis, as in
Lemma~\ref{lem:common-hecke-interpolation}. The numbers of either
sign coincide with those of \(J|_F\). Indeed in a diagonal signed
basis the largest dimension of a subspace on which the form is
positive equals the number of positive entries, by projecting to
their span, and similarly for negative entries. The isometry \(T\)
preserves these maximal dimensions. Send a unit signed eigenbasis
of \(F\) to the matched eigenvectors in \(\mathcal V\) divided by
the square roots of the absolute eigenvalues. This signed isometry
\(U:F\to\mathcal V\) has ordinary norm at most \(\sqrt{1+B^2}\).
Define \(H|_F=T^{-1}U\) and \(H|_E=I\). The two spaces are also
signed orthogonal, so \(H\) is an isometry. On the two summands
\(TH\) is \(T_E,U\), which gives its norm bound by the triangle
inequality.
\end{proof}

\begin{theorem}\label{thm:partial-real-head-sizes}
For \(1\le j\le\nu\) write \(E=(e_1,\ldots,e_j)\),
\(E'=(e_1^*,\ldots,e_j^*)\), and \(Q_j=\{s\in G_1:sE=E A_s,\
A_s\in N_j(\C)\}\), with \(N_j\) upper unitriangular. Let \(w\)
flip these \(j\) pairs with parameters \(p_i\in\C^\times\) as in
Definition~\ref{def:real-pivot-representative} and fix the other
coordinates. Write \(D_i=\prod_{l=1}^i p_l,\ D_0=1\). Suppose
\[
 s,t\in Q_j,\qquad \|sg\|_{\rm op},\ \|tw^{-1}g\|_{\rm op}\le R,
 \qquad L\ge\max\bigl(4,r,R,\max_{0\le i\le j}|D_i|^{-1}\bigr).
\]
Then \(|D_i|\le R^{2i}\) for \(i\le j\), and some \(h\in G_{j+1}\)
satisfies
\[
       \|hg\|_{\rm op}\le L^{4j+2},\qquad
                   \|hw^{-1}g\|_{\rm op}\le L^{6j+3}.
\]
Also \(\|sE'\|_{\rm op}\le L^{4j+1}\). If \(s\) acts trivially
on the intervening signed space \(V_{j+1}\) modulo
\(\operatorname{span}E\), then
\(\|g\|_{\rm op}\le L^{6j+4}\) and
\(\|w^{-1}g\|_{\rm op}\le L^{8j+5}\).
\end{theorem}
\begin{proof}
As in Proposition~\ref{prop:top-pivot-support}, take
\(X=E^*J(s w t^{-1})E\). It has norm at most \(R^2\), since
\(s w t^{-1}=(s g)(t w^{-1}g)^{-1}\). Its triangular
factorization is \(X=L_1 D U_1\), where \(L_1=A_s^{-*}\),
\(U_1=A_t^{-1}\), \(D=\operatorname{diag}(p_i)\). This gives
the bound on the leading determinants. For \(b\le i\le j\),
\[
 (L_1^{-1})_{ib}=(-1)^{b+i}
    \frac{\det X[\{1,\ldots,i\}\setminus\{b\},\,1{:}i-1]}{D_{i-1}} .
\]
Indeed the signed cofactors form a row annihilating the first
\(i-1\) columns in the leading \(i\) rows, by a determinant with
a repeated column. Their last entry after division is one.
The corresponding row of \(L_1^{-1}\) has these properties, and
the first \(i-1\) entries are uniquely determined by the invertible
leading minor. Transpose the argument to obtain the analogous
formula for the entries of \(U_1^{-1}\). All the numerator minors
have modulus at most \(R^{2(i-1)}\) by the column projection proof
in Proposition~\ref{prop:top-pivot-support}. Bounding an operator
norm by \(j\) times its largest entry now gives
\(\|A_s\|_{\rm op},\|A_t\|_{\rm op}\le L^{2j}\).

Write \(T=g^{-1}\). Then \(TE=(s g)^{-1}EA_s\) and
\(TwE=(t w^{-1}g)^{-1}EA_t\) have norms at most \(L^{2j+1}\).
Since \(wE=E'D\) and
\(|p_i|^{-1}=|D_{i-1}/D_i|\le L^{2j-1}\), we obtain
\(\|TE'\|_{\rm op}\le L^{4j}\) and
\(\|sE'\|_{\rm op}\le R\|TE'\|_{\rm op}\le L^{4j+1}\).
On the span of \(E,E'\), \(T\) therefore has norm at most
\(2L^{4j}\). Apply Lemma~\ref{lem:bounded-signed-complement} to
this span and \(V_{j+1}\), and take \(h=H^{-1}\). Then
\(\|hg\|_{\rm op}=\|(hg)^{-1}\|_{\rm op}\le6L^{4j}\le L^{4j+2}\).
Since \(|p_i|=|D_i/D_{i-1}|\le L^{2j+1}\), the norm of \(w\) and
of its inverse is at most \(L^{2j+1}\). It commutes with \(h\),
giving the second assertion.

Under the additional hypothesis, in the blocks \(E,V_{j+1},E'\)
the matrix \(s\) has diagonal \(A_s,I,A_s^{-*}\), zero lower
blocks and upper blocks \(B_1,C_1,D_1\) in positions \(12,13,23\).
This follows by preserving \(E\) and its signed perpendicular,
and by evaluating its mixed \(E,E'\) pairing.
The signed inner products of middle and last columns give
\(B_1=-A_s D_1^*J|_{V_{j+1}}\). The whole last block column has
norm at most \(L^{4j+1}\) by the preceding bound on \(sE'\).
Thus \(\|s\|_{\rm op}\le6L^{6j+1}\le L^{6j+3}\) by adding the
norm bounds of its six blocks. Multiplying \(sg\) by \(s^{-1}\),
and then multiplying \(g\) by \(w^{-1}\), gives the two final
inequalities.
\end{proof}


\section{Uniform integrals on the successive real floors}

\begin{definition}\label{def:real-near-floor}
Write \(f_j=e_j\) for \(1\le j\le\nu\), \(f_n=z\).
For \(0\le j\le n\) let \(Q_j\) consist of those \(s\in G_1\) with
\(s f_a-f_a\in\operatorname{span}(f_1,\ldots,f_{a-1})\) for
all \(a\le j\). Thus \(Q_0=G_1\), and for \(1\le j\le\nu\)
this agrees with Theorem~\ref{thm:partial-real-head-sizes}.
Say \(g\) is near at \(j\) when \(Q_j g\) meets \(B_G(R)\);
write \(\mathcal N_j\) for this set. For \(1\le j<n\) and \(x\in G_1\)
let \(v_j(g,x)\) be the ordinary area of the set
\[
 \{[f_j,w]:[f_a,w]=0\ (a<j),\
                  (g^{-1}w,(gx)^{-1}w)\ \text{is close}\},
\]
using Proposition~\ref{prop:real-two-position-gauge}.
\end{definition}

\begin{lemma}\label{lem:real-near-floor-invariance}
Each \(\mathcal N_j\) is closed, \(v_j(g,x)\) is a positive finite
Borel function invariant under \(g\mapsto qg\) for \(q\in Q_j\),
and for every nonnegative Borel function \(H_1\) invariant on
the left by \(Q_j\) one has
\[
     \int H_1(\sigma q g)\dd\theta_j(\sigma)
           =\int H_1(\sigma g)\dd\theta_j(\sigma)
                 \qquad (q\in Q_{j-1},\ 1\le j\le n).
\]
If \(g\in B_G(R),\ \sigma\in G_j,\ \sigma g\in\mathcal N_j\),
then \(\|\sigma^{-1}f_j\|\le R^2\).
\end{lemma}
\begin{proof}
The flag conditions are closed matrix equations and specify a
subgroup since the flag matrices are upper unitriangular. The
product \(Q_j B_G(R)\) is closed by the compact subsequence
argument of Lemma~\ref{lem:middle-haar-cutoff}. The close set in
the defining subspace for \(w\) is open, convex, invariant under
scalars of modulus at most one, and bounded by the condition on
\(g^{-1}w\). Its image by \(w\mapsto[f_j,w]\) is a bounded open
disk centered at zero with positive radius. Indeed \(e_j^*\)
is an allowed direction with image one. Membership of a scalar
\(t\) in this disk is an open condition on \((g,x,t)\), since one
can perturb a witnessing \(w\) by a multiple of \(e_j^*\) to
realize any nearby scalar and the norm inequalities are strict.
Integrating its Borel indicator in \(t\) proves measurability.
Under \(w\mapsto q w\), the first \(j-1\) orthogonality equations
are equivalent by triangularity of \(q^{-1}\), and \([f_j,q w]\)
equals \([f_j,w]\) on that subspace. The close pair is preserved
on simultaneously replacing \(g\) by \(qg\).

Any \(q\in Q_{j-1}\) factors as \(h u\), where \(h\in G_j\) and
\(u\) has upper block triangular form in the blocks consisting
of the first \(j-1\) vectors, \(V_j\), and their duals, with
identity middle block and upper unitriangular first block. In
fact \(q\) preserves the span of the first vectors and its signed
perpendicular and thus acts on their intermediate quotient as
an isometry. Lift this isometry diagonally to \(h\) to obtain the
factorization. The subgroup of such \(u\) is normalized by
\(G_j\) and lies in \(Q_j\), because the middle images are
unchanged modulo the first \(j-1\) vectors. Hence in the integral
the left invariance of \(H_1\) reduces \(q\) to \(h\).
For fixed \(g\), \(H_1(\sigma g)\) on \(G_j\) depends only on
\(\sigma^{-1}f_j\), since the stabilizer of \(f_j\) there is
contained in \(Q_j\). A Borel representative outside a null set
is the triangular inverse parametrization at \(j<n\) used in
Proposition~\ref{prop:actual-floor-quadrics} when the dual
coordinate \(t\ne0\). At \(j=n\) use
Lemma~\ref{lem:middle-haar-cutoff}. The law of \(\sigma^{-1}f_j\)
is invariant under \(h^{-1}\) by
Proposition~\ref{prop:actual-floor-quadrics}, so the null exception
also has zero measure after this change. This proves the integral
identity. Finally choose \(q'\in Q_j\) with \(\|q'\sigma g\|_{\rm op}
\le R\). Projecting \(\sigma^{-1}(q')^{-1}f_j\) orthogonally
onto \(V_j\) gives \(\sigma^{-1}f_j\), since \(\sigma\) fixes
the earlier frame vectors. The unprojected vector has norm at
most \(\|g\|_{\rm op}\|(q'\sigma g)^{-1}\|_{\rm op}\le R^2\).
\end{proof}

\begin{theorem}\label{thm:real-floor-mass}
At sufficiently large arithmetic scales, uniformly in \(g,x\in G_1\),
\[
 \begin{aligned}
 \int_{\sigma g\in\mathcal N_j}v_j(\sigma g,x)^{-s}
                           \dd\theta_j(\sigma)&\le R^{29n}
                                             &&(1\le j<n),\\
 \theta_j\{\sigma:\sigma g\in\mathcal N_j\}&\le R^{5k}
                                             &&(1\le j\le n).
 \end{aligned}
\]
Here \(m=k-1=s+1+3n\), as in
Proposition~\ref{prop:real-two-position-gauge}.
\end{theorem}
\begin{proof}
If \(\sigma g\in\mathcal N_j\) for some \(\sigma\in G_j\), then
\(g\in\mathcal N_{j-1}\), because \(G_j,Q_j\subset Q_{j-1}\).
Otherwise both relevant integrals vanish by the support of
\(\theta_j\). We can now take \(g\in B_G(R)\) using the integral
invariance of Lemma~\ref{lem:real-near-floor-invariance}.
Put \(T=R^2,\ h=n-j\). The orbit vector \(v=\sigma^{-1}f_j\)
on the integration set has norm at most \(T\).

For the first bound apply Proposition~\ref{prop:real-two-position-gauge}
to \(g,gx\) with no row constraints. If \(w\) belongs to its set
\(W\), then \(\sigma E_+ w\) is signed orthogonal to the preceding
\(f_a\), and the close pair required for \(v_j(\sigma g,x)\) is
\((g^{-1}E_+w,(gx)^{-1}E_+w)\). Its scalar at \(f_j\) is
\([v,E_+w]=w_v^*w\), where \(w_v\) is the \(m\)-component of
\(v\) in Definition~\ref{def:real-flat-quadric}. Thus
\(v_j(\sigma g,x)^{-1}\le\phi(w_v)\) for the projection there.
By Proposition~\ref{prop:actual-floor-quadrics} and the fibre
bound proved in Theorem~\ref{thm:real-flat-fourier},
the integral is at most
\[
     M_{\rm quad}(T)
          \int_{\|w\|\le T}\left(\frac{36}{\|\Pi w\|^2}\right)^s\dd w .
\]
Indeed at each fixed \(w\) one integrates all remaining parameters
in the box using \(m_{\mu,T}(p)\) from that proof. By
Lemma~\ref{lem:real-projective-moment} and the exponential estimate
in Proposition~\ref{prop:real-two-position-gauge}, the ball
integral is at most \(R^{6(3n+1)}\). Also
\(M_{\rm quad}(R^2)=4\pi R^4(\pi R^4)^{2h}\le R^{6+10h}\)
as \(R\ge4\). Their product is at most \(R^{29n}\) when \(n\ge48\).

For the unweighted bound the identical fibre calculation applies
to the constant one, including at the last floor by the negative
quadric assertion of Proposition~\ref{prop:actual-floor-quadrics}.
The remaining ball has volume \(\pi^m R^{4m}/m!\).
Since \(\pi^m/m!\le\exp(\pi)\le R^4\), the resulting bound is
\(R^{4m+10h+10}\le R^{5k}\) for the actual ranks
\(k\ge n^7,\ n\ge48\) by Proposition~\ref{prop:thick}.
\end{proof}


\section{Recovery of a large position from its flag coordinates}

\begin{lemma}\label{lem:real-dual-flag-reconstruction}
Let \(u\in\mathcal U\), \(\nu\ge1\), and put \(Z_0\ge1\).
Suppose that for each \(a=1,\ldots,\nu\) there is a vector \(w_a\)
with \([e_b,w_a]=\delta_{ba}\) for \(1\le b\le a\) and
\(\|w_a\|,\|u w_a\|\le Z_0\).
With \(Y_0=\nu Z_0(1+\nu Z_0)^\nu\) one has
\[
                     \|u^{\pm1}\|_{\rm op}\le12(1+Y_0)^3 .
\]
\end{lemma}
\begin{proof}
Let \(W,W'\) have columns \(w_a,u w_a\). Their bottom blocks,
in the rows of \(e_1^*,\ldots,e_\nu^*\), are lower unitriangular
matrices \(Z,Z'\). For \(W\) this is the assumed pairing
condition. For \(W'\), \(u^{-1}e_b-e_b\) lies in the span of the
earlier vectors and all their pairings with \(w_a\) vanish for
\(b\le a\). The norms of \(W,W'\) are at most \(\sqrt\nu Z_0\)
by columnwise Cauchy--Schwarz. Each strictly lower part of \(Z,Z'\)
has norm at most \(\nu Z_0\) by the entry bound, so the two inverse
norms are at most \((1+\nu Z_0)^\nu\) by the finite geometric
series.

Write \(u\) in the three blocks \(e,V_n,e^*\) with diagonal
\(A,I,A^{-*}\) and upper blocks \(B,C,D'\). The identities
\(Z'=A^{-*}Z\) and \(Z=A^*Z'\) give
\(\|A\|_{\rm op},\|A^{-*}\|_{\rm op}\le Y_0\).
Both \(W Z^{-1}\) and \(W' Z^{-1}\) have norm at most \(Y_0\),
and the first has bottom block \(I\). Multiplication by \(u\)
sends the first matrix to the second. In the middle block
this gives \(\|D'\|_{\rm op}\le2Y_0\). The mixed signed pairing
of the middle and last columns of \(u\) gives
\(B=-A(D')^*J|_{V_n}\), so \(\|B\|_{\rm op}\le2Y_0^2\).
The first block of the same matrix equation now gives
\(\|C\|_{\rm op}\le Y_0+Y_0^2+2Y_0^3\). Summing the six block
bounds proves the stated inequality for \(u\), and
\(u^{-1}=Ju^*J\) gives the other one.
\end{proof}

\begin{theorem}\label{thm:real-large-position-cap}
Let \(0\le j\le\nu\), \(p,q\in(\C^\times)^j\), and let \(w_p,w_q\)
flip these first \(j\) pairs. Write \(D_{p,i}=\prod_{a\le i}p_a\)
for \(1\le i\le j\) and similarly for \(q\). Let \(g,x\in G_1\) satisfy
\[
  S(g),S(gx)\le R,\qquad
               w_p^{-1}g,\ w_q^{-1}gx\in\mathcal N_j .
\]
Choose \(L\ge\max(4,r,R)\) also bounding the inverses
\(|D_{p,i}|^{-1},|D_{q,i}|^{-1}\), and \(T\ge1\). Put
\[
    A=\max\bigl(R,r,4,T,H_j\bigr),\qquad
             H_0=1,\quad H_j=L^{4j+1}\quad(j>0).
\]
If \(\|x-I\|_{\rm op}>A^{16n}\), then there is an \(a\) with
\(j<a\le\nu\) and \(v_a(g,x)<\pi/T^2\). When \(j=\nu\)
the assumptions imply \(\|x-I\|_{\rm op}\le A^{16n}\).
\end{theorem}
\begin{proof}
Suppose all the indicated areas are at least \(\pi/T^2\).
As in Proposition~\ref{prop:top-pivot-support}, choose
\(u,v\in\mathcal U\) such that both \(ug\) and \(vgx\) lie
in \(B_G(R)\). Put \(s=v u^{-1}\).
For \(a\le j\) take \(w_a=u e_a^*\). Its first \(a\) dual
coordinates have the required values by the unitriangular action
of \(u^{-1}\) on the first frame. Its image under \(s\) is
\(v e_a^*\). The two near conditions give left multipliers in
\(Q_j\) as required by Theorem~\ref{thm:partial-real-head-sizes}
applied separately to \(g,p\) and \(gx,q\). Consequently both
norms are at most \(L^{4j+1}\). No such vectors are needed
when \(j=0\).

For \(a>j\), the disk of attainable values in
Definition~\ref{def:real-near-floor} has radius at least \(1/T\)
by Lemma~\ref{lem:real-near-floor-invariance}. Choose a witnessing
vector \(w\) for a value \(c\) of modulus \(1/(2T)\), and put
\(w_a=u w/c\). The first \(a-1\) dual coordinates are zero
and the \(a\)-th is one by the flag action of \(u^{-1}\).
The close conditions give
\(\|g^{-1}w\|<1/4,\ \|(gx)^{-1}w\|<1/2\).
Multiplying by \(ug\) and \(vgx\) shows that the norms of
\(w_a\) and \(s w_a\) are at most \(RT\).
We can apply Lemma~\ref{lem:real-dual-flag-reconstruction}
with \(Z_0=\max(H_j,RT)\le A^2\). Since \(\nu\le A\),
\(Y_0\le A^{3+4\nu}\). Hence
\[
       \|x-I\|_{\rm op}
           \le 1+R^2\,12(1+Y_0)^3
           \le A^{16n}.
\]
For the last inequality \(n=\nu+1,\ A\ge4\), so
\(1+96 A^{11+12\nu}\le A^{15+12\nu}\le A^{16n}\).
The first inequality uses \(x=(ug)^{-1}s^{-1}(vgx)\).
This proves the contrapositive, including the empty range at
\(j=\nu\).
\end{proof}


\section{Integration of the scalar-area weights along a real radical}

\begin{definition}\label{def:real-radical-weights}
For \(0\le j\le\nu\) let \(N^j\subset Q_j\) be the subgroup acting
identically on the middle quotient \(V_{j+1}\) with respect to
the first \(j\) pairs. Thus \(N^0=\{1\}\). Give \(N^j\) the product
of the first \(j\) layer measures in Proposition~\ref{prop:unipotent-covol},
using unit mass at \(j=0\). For \(a\le\nu\) put
\[
 a_a(g,y)=v_a(g,g^{-1}y),\qquad
 \mathcal W_j(g,y)=\prod_{a=1}^j a_a(g,y)^s .
\]
The two images in the close-pair condition for \(a_a(g,y)\)
are \(g^{-1}w,y^{-1}w\). For \(t\in V_{j+1},\ \beta\in\R\)
write \(n_j(t,\beta)\) for the \(j\)-th unipotent element sending
\[
 \begin{split}
 e_j^*&\longmapsto e_j^*+t+(i\beta-[t,t]/2)e_j,\\
 z'&\longmapsto z'-[t,z']e_j\qquad(z'\in V_{j+1}),
 \end{split}
\]
and fixing \(e_j\) and the previous pairs.
\end{definition}

\begin{lemma}\label{lem:partial-real-radical-measure}
For \(1\le j\le\nu\), to integrate on \(N^j\) one may integrate
at \(u=u_0 n_j(t,\beta)\) or at \(u=n_j(t,\beta)u_0\) against
\(\dd u_0\,\dd t\,\dd\beta\), \(u_0\in N^{j-1}\). Also
\(Q_j=N^j G_{j+1}=G_{j+1}N^j\), and conjugation by \(Q_j\)
preserves the measure on \(N^j\). For \(q\in Q_j\),
\(\mathcal W_j(qg,qy)=\mathcal W_j(g,y)\).
\end{lemma}
\begin{proof}
The layer matrices are those of Lemma~\ref{lem:heisenberg-layer}
with row \(-t^*J|_{V_{j+1}}\) and imaginary corner coordinate
\(\beta\). This row change has absolute real determinant one since
the Gram matrix is ordinary unitary. In the decomposition by
the first \(j-1\) pairs the middle block of \(u\in N^j\) ranges
over exactly these \(n_j\). The projection to that block is a
homomorphism with kernel \(N^{j-1}\), split by the indicated
embedding. The layer product in outside-to-inside order gives
\(\dd u_0\,\dd t\,\dd\beta\). Conjugation by any \(h\in G_j\)
preserves the measure on \(N^{j-1}\), since in each of the earlier
layers it acts on the vector variable as a signed isometry of
complex determinant of modulus one and fixes the real central
variable. This is the conjugation computation of
Theorem~\ref{thm:exact-real-floor-haar}. In particular, reversing
these two factors has the same measure. The same argument shows
that \(G_{j+1}\) preserves measure on \(N^j\). Haar measure there
is bi-invariant by the iterated semidirect products with these
unit determinants, so conjugation by \(N^j\) also preserves it.
The factorization of \(Q_j\) is obtained by extracting the middle
block, as in Lemma~\ref{lem:real-near-floor-invariance}. That
lemma and \(Q_j\subset Q_a\) for \(a\le j\) give the last identity.
\end{proof}

\begin{theorem}\label{thm:real-positive-radical-weight}
Let \(0\le j\le\nu,\ g\in\mathcal N_j,\ y\in G_1\), and
\(L\ge\max(4,r,R)\), at sufficiently large arithmetic scales.
Then
\[
 \int_{N^j}{\bf1}_{\|g^{-1}u y-I\|_{\rm op}\le L}
                      \mathcal W_j(g,u y)\dd u\le L^{70 n j}.
\]
\end{theorem}
\begin{proof}
Simultaneously multiply \(g,y\) on the left by an element of
\(Q_j\), conjugating the integration variable by that element.
Lemma~\ref{lem:partial-real-radical-measure} shows that the integral
is unchanged, so assume \(g\in B_G(R)\). If the indicated set is
empty the bound holds. Otherwise translate \(y\) on the left by
one of its admissible elements in \(N^j\); right invariance in
\(u\) then permits \(\|g^{-1}y-I\|_{\rm op}\le L\). For any
contributing \(u\),
\[
       \|u\|_{\rm op}\le R^2(1+L)^2\le L^6 .
\]
In fact \(u=g(g^{-1}uy)(g^{-1}y)^{-1}g^{-1}\), and the norm of
a signed isometry equals that of its inverse.

We proceed by induction from \(j=0\), where the assertion is
the bound by one for an indicator. For \(j>0\) write
\(u=u_0n_j(t,\beta)\) by Lemma~\ref{lem:partial-real-radical-measure}.
The middle block obtained by removing \(j-1\) pairs is \(n_j(t,\beta)\).
Its \(e_j^*\) column gives \(\|t\|,|\beta|\le L^6\) wherever
the integrand is nonzero. Write \(t_+\) for its component in
the positive coordinate subspace \(E_+\C^m\), identifying that
subspace with \(\C^m\). Use \(W,\phi\) of
Proposition~\ref{prop:real-two-position-gauge} for the pair of
isometries \(g,y\) with no row constraints.
We claim \(a_j(g,u y)\le\phi(t_+)\).

Indeed let \(\alpha=[e_j,w]\) be attained in its defining disk,
and let \(z=E_+w_0\) with \(w_0\in W\). Direct subtraction gives
\[
 [g^{-1}w,g^{-1}z]-[(uy)^{-1}w,y^{-1}z]
                  =[w,z-u z]=\bar\alpha\,t_+^*w_0 .
\]
For the last equality, \(n_j(t,\beta)z=z-[t,z]e_j\);
\(u_0\) fixes both \(z\) and \(e_j\) modulo the span of the
earlier \(e_a\), and these discarded components pair to zero
with \(w\). The pairs formed by the two images of \(w\) and
by the two images of \(z\) are both close. Thus every
\(\alpha\) lies in the scalar set in part 2 of
Proposition~\ref{prop:real-two-position-gauge}, giving the claim.

At fixed \(t,\beta\) apply the induction hypothesis to \(g\) and
\(n_j(t,\beta)y\), using \(\mathcal N_j\subset\mathcal N_{j-1}\).
The result is bounded by \(L^{70n(j-1)}\) times the integral
of \(\phi(t_+)^s\) over the above \(t,\beta\) ranges.
The Euclidean complement of \(E_+\) in \(V_{j+1}\) has complex
dimension \(2(n-j)\). Its ball of radius \(L^6\) has volume at
most \((\pi L^{12})^{2(n-j)}\le L^{26n}\).
By Lemma~\ref{lem:real-projective-moment}, the integral on the
positive ball of radius \(L^6\) is at most \(L^{14(3n+1)}\);
the exponential and factorial factors are bounded as in
Proposition~\ref{prop:real-two-position-gauge}.
The real central interval has length at most \(2L^6\le L^7\).
Together these three bounds are at most \(L^{70n}\) for
\(n\ge48\). Nonnegative integration by
Lemma~\ref{lem:product-integration} completes the induction.
\end{proof}


\section{Reciprocal weights on the opposite real flag}

\begin{theorem}\label{thm:real-inverse-radical-weight}
For \(0\le j\le\nu\), let \(w_p\) flip the first \(j\) pairs,
with parameters \(p_1,\ldots,p_j\in\C^\times\).
Assume \(g\in\mathcal N_j,\ x\in G_1,\ L\ge\max(4,r,R)\) and
\(|p_1\cdots p_i|^{-1}\le L\) for \(1\le i\le j\).
At sufficiently large arithmetic scales,
\[
 \int_{N^j}\frac{{\bf1}_{\,w_p u g\in\mathcal N_j}}
                 {\mathcal W_j(w_p u g,w_p u g x)}\,\dd u
       \le L^{125 n^2 j}\prod_{a=1}^j |p_a|^{-2d_a-1}.
\]
Here \(d_a=\dim_\C V_{a+1}\) as in
Definition~\ref{def:successive-real-floors}.
\end{theorem}
\begin{proof}
The integral is unchanged under \(g\mapsto q g\) for
\(q\in Q_j\). For \(q\in N^j\) this uses right invariance in
\(u\). For \(q\in G_{j+1}\), conjugate \(u\) by \(q\), commute
\(q\) with \(w_p\) and use the invariance of both \(\mathcal N_j\)
and \(\mathcal W_j\) under simultaneous left multiplication.
These cases suffice by Lemma~\ref{lem:partial-real-radical-measure}.
Thus suppose \(g\in B_G(R)\).

We use induction on \(j\), the empty product giving the assertion
at \(j=0\). Decompose \(u=n_j(t,\beta)u_0\) in the inner-to-outer
order of Lemma~\ref{lem:partial-real-radical-measure}. If
\(w_p u g\in\mathcal N_j\), apply
Theorem~\ref{thm:partial-real-head-sizes} at \(u g\), with left
multiplier \(u^{-1}\) and Weyl element \(w_p^{-1}\), whose
parameters are \(\bar p_a\). Its proof bounds
\(\|u^{\pm1}\|_{\rm op}\le L^{6j+3}\), since \(u^{-1}\) is
trivial on the middle quotient and \(g\in B_G(R)\). More
explicitly its mixed middle and last blocks are bounded by
\(L^{6j+1},L^{4j+1}\) and its diagonal blocks by
\(L^{4j+1}\), as shown in the last paragraph of that proof.
Also \(|p_j|\le L^{2j+1}\) by its leading-minor estimate.
The middle block of \(u\) on \(V_j\) is \(n_j(t,\beta)\).
Consequently
\[
         |p_j|\|t\|,\ |p_j||\beta|\le L^{8j+4}\le L^{12n}
\]
on the integration set.

For fixed \(u_0\), the image \(u_0 E_+ w_0\) equals \(E_+w_0\)
plus earlier \(e_1,\ldots,e_{j-1}\) with coefficients linear in
\(w_0\). Put these coefficients in the matrix \(C\) of
Proposition~\ref{prop:real-two-position-gauge}; it has at most
\(j-1\) rows. Apply that proposition with isometries \(g,gx\).
If \(w_0\in W\), the vector \(w=w_p n_j(t,\beta)u_0 E_+w_0\)
has zero earlier dual coordinates and \(j\)-th dual coordinate
\(-p_j t_+^*w_0\), where \(t_+\) is the projection of \(t\)
onto the indicated positive coordinates. Indeed \(u_0 E_+w_0=E_+w_0\);
then \(n_j\) subtracts \([t,E_+w_0]e_j\), and \(w_p e_j=p_j e_j^*\).
The two inverse images of \(w\) under \(w_p u g,w_p u g x\) are
\(g^{-1}E_+w_0,(gx)^{-1}E_+w_0\), a close pair. Therefore part 1
of that proposition gives
\[
 a_j(w_p u g,w_p u g x)^{-1}\le \phi(-\bar p_j t_+).
\]
Dependence of the projection in \(\phi\) on \(u_0\) causes no
integration issue, since we use its bound on the integral at
each fixed \(u_0\) and the resulting scalar bound is uniform.

Write \(w_{<j}\) for the flip on the first \(j-1\) pairs.
The two arguments of each earlier area \(a_i,\ i<j\), are obtained
from \(w_{<j}u_0g,w_{<j}u_0gx\) by simultaneous left multiplication
by \(w_j(p_j)n_j(t,\beta)\in G_j\subset Q_{j-1}\). These earlier
areas are unchanged by Lemma~\ref{lem:partial-real-radical-measure}.
Also the original near condition implies
\(w_{<j}u_0g\in\mathcal N_{j-1}\). It remains to bound the
integration in \(t,\beta\) of the new factor. Set
\(t'=-\bar p_j t,\ \beta'=|p_j|\beta\). The change gives the
factor \(|p_j|^{-2d_j-1}\). Both the norm of \(t'\) and the
modulus of \(\beta'\) are bounded by \(T=L^{12n}\) on the
original set. Enlarge to these bounds. As in
Theorem~\ref{thm:real-positive-radical-weight}, the integral of
\(\phi(t'_+)^s\) is at most
\[
          L^{2(3n+1)} T^{2(3n+1)}
                    (\pi T^2)^{2(n-j)}\,2T
                      \le L^{125 n^2}.
\]
Indeed the exponent after \(\pi\le L,\ 2\le L\) is at most
\((24n+2)(3n+1)+(24n+1)2n+12n+1\le125n^2\).
Now integrate the remaining earlier factors in \(u_0\) by the
induction hypothesis, valid since \(g\in B_G(R)\).
All original integrands are nonnegative and measurable by
Lemma~\ref{lem:real-near-floor-invariance}. Their iterated
integrals give the asserted bound.
\end{proof}

\begin{corollary}\label{cor:real-radical-test}
With the head and scale of Theorem~\ref{thm:real-inverse-radical-weight},
let \(A\) be nonnegative Borel, invariant under left \(N^j\), and
zero outside \(\mathcal N_j\). For a positive sigma-finite Borel
measure \(\rho'\) on \(G_1\),
\[
 \int\frac{A(y){\bf1}_{w_p y\in\mathcal N_j}}
             {\mathcal W_j(w_p y,w_p y x)}
           \dd(\dd u*\rho')(y)
      \le L^{125n^2j}\prod_{a=1}^j |p_a|^{-2d_a-1}
                                                \int A\,\dd\rho' .
\]
Here \(\dd u\) denotes the integration measure on \(N^j\).
\end{corollary}
\begin{proof}
At \(y=u y_0\), \(A(y)=A(y_0)\). If this value is positive,
\(y_0\in\mathcal N_j\), so the integral in \(u\) with \(g=y_0\)
is bounded by Theorem~\ref{thm:real-inverse-radical-weight}.
Integrate in \(y_0\) by Lemma~\ref{lem:product-integration}.
For a possibly infinite \(A\) the same argument first applies to
its truncations and then Lemma~\ref{lem:integral-tools} gives
the result.
\end{proof}


\section{Successive averages of a real Bessel pairing}

\begin{definition}\label{def:real-successive-pairing}
Fix \(0\le j_0\le\nu\) and two heads of length \(j_0\).
Let \(d_*=w_q^{-1}w_p\). On \(e_a\) for \(a\le j_0\) its
diagonal entry is \(r_a=p_a/q_a\) and on \(e_a^*\) it is
\(1/\bar r_a\), by Definition~\ref{def:real-pivot-representative}.
Set \(r_a=1\) for \(j_0<a\le n\). Let \(B\) be any bounded
Borel function taking values in \(\pi\) and transforming on
\(S_0\) by the same rule as \(V\) in
Theorem~\ref{thm:actual-bessel-truncation}. Write
\(P_i=\kappa*\theta_n*\cdots*\theta_{i+1}\) for \(0\le i\le n\),
with \(P_n=\kappa\). Define
\[
 K_i(g,y)=\int\langle V(\sigma g),B(d_*\sigma y)\rangle
                                      \dd P_i(\sigma).
\]
In this section write \(\ell(u)=\sigma(u)\) for the sum of the
simple-root entries in Definition~\ref{def:real-bessel-graph}.
\end{definition}

\begin{lemma}\label{lem:real-pairing-floor-laws}
All \(K_i\) are bounded Borel and vanish unless \(g\in\mathcal N_i\).
For \(1\le i\le n\),
\[
 K_{i-1}(g,y)=\int K_i(\sigma g,\sigma y)\dd\theta_i(\sigma),
 \qquad K_n(g,y)=\langle V(g),B(d_*y)\rangle .
\]
The iterated integrals are absolutely convergent for all \(g,y\).
For \(i\ge j_0\), \(K_i(g,y)\) is invariant under simultaneous
left translation of both arguments by \(Q_i\cap G_{j_0+1}\).
For \(1\le i\le n,\ u,u'\in N^{i-1}\),
\[
 K_i(u g,u'y)
   =\psi\bigl(\ell(d_*u'd_*^{-1})-\ell(u)\bigr)K_i(g,y).
\]
For \(j_0<i\le n\), \(|K_i|\) is invariant under simultaneous
left multiplication by \(Q_i\).
\end{lemma}
\begin{proof}
The compact factor commutes with \(d_*\), and both functions
have the same compact unitary action. The formula for \(K_n\)
follows, also with the modulus inside the integral.
If this pairing is nonzero, \(S(g)\le R\) implies that a left
unipotent translate of \(g\) belongs to \(B_G(R)\). In particular
\(g\in\mathcal N_n\). Use backward induction in the product
measures for the integral of the absolute pairing. When
\(\sigma g\in\mathcal N_i,\ \sigma\in G_i\), one has
\(g\in\mathcal N_{i-1}\). By Theorem~\ref{thm:real-floor-mass},
each integration increases the uniform bound by at most \(R^{5k}\).
The root bound is \(M_V\sup\|B\|\), and thus all the claimed
integrals are finite. Lemma~\ref{lem:product-integration} gives
the recurrence and Borel measurability.

For the simultaneous invariance begin at \(i=n\).
Here \(Q_n=\mathcal U\,\iota(U(P_0))\), since after removing all
outer pairs the flag condition on \(z\) means precisely fixing
\(z\) on the middle quotient. On the indicated intersection
\(d_*\) commutes with the multiplier, and both components undergo
the same unitary action. In the induction step take
\(j_0\le i-1\), \(q'\in Q_{i-1}\cap G_{j_0+1}\) and factor
\(q'=h u\) with \(h\in G_i,\ u\in N^{i-1}\cap G_{j_0+1}\) by
Lemma~\ref{lem:partial-real-radical-measure}. For \(\sigma\in G_i\),
conjugation of \(u\) by \(\sigma h\) stays in that intersection and
in \(Q_i\), so the induction hypothesis removes it from the
integrand. The remaining integrand on \(G_i\) is invariant on the
left by \(G_i\cap Q_i\), since this stabilizer lies in \(G_{j_0+1}\).
It factors measurably through \(\sigma^{-1}f_i\) outside a null
set by the chart argument in Lemma~\ref{lem:real-near-floor-invariance}.
Right multiplication by \(h\) preserves that law by
Proposition~\ref{prop:actual-floor-quadrics}. This proves the
claimed invariance down to \(i=j_0\).

The measure \(P_i\) is carried by isometries fixing the first
\(i\) frame vectors and their hyperbolic partners when present.
For \(i=n\) these are the compact middle rotations. Conjugating
\(u\in N^{i-1}\) by such a \(\sigma\) preserves that radical
and all its simple-root entries. Indeed the only possible nonzero
ones have indices \(a\le i-1\), and both frame coordinates in each
such entry are among the first \(i\). For \(a<i\) the ordinary
\(e_a\)-row is fixed as well by the split block decomposition.
Conjugation by \(d_*\) multiplies the \(a\)-th entry by
\(r_a/r_{a+1}\). Applying the separate Bessel transformation laws
inside the integral therefore gives the displayed phase.

Finally suppose \(j_0<i<n\). For a simultaneous \(u\in N^i\)
the difference of the two phases in the integral involves only
the first \(j_0\) simple-root entries, because their multipliers
\(r_a/r_{a+1}\) equal one for \(a>j_0\). These entries are
unchanged by \(\sigma\) as \(j_0+1\le i\). The phase is constant
in the integral, so its modulus is unchanged. Each \(q'\in Q_i\)
is a product of such \(u\) and an element of \(G_{i+1}\),
which lies in \(Q_i\cap G_{j_0+1}\). At \(i=n\), use the separate
phase law for \(N^\nu\) and compact equivariance. This proves
the last assertion.
\end{proof}

\begin{proposition}\label{prop:real-radical-fourier-phase}
Let \(j_0<i\le\nu\) and \(g,y\in G_1\). For \(\sigma\in G_{i+1}\)
write \(v=\sigma^{-1}f_{i+1}\). There is a Borel amplitude \(F(v)\)
on the corresponding quadric satisfying
\(F(v)=K_{i+1}(\sigma g,\sigma y)\) almost everywhere for
\(\theta_{i+1}\), and
\[
 K_i(g,n_i(t,\beta)y)=
    \int\psi(-[t,v]) F(v)\dd\delta_\gamma^{(i+1)}(v)
                           \qquad(t\in V_{i+1},\ \beta\in\R).
\]
Here \(\gamma=0\) for \(i+1<n\) and \(\gamma=-1\) for \(i+1=n\).
\end{proposition}
\begin{proof}
The stabilizer of \(f_{i+1}\) in \(G_{i+1}\) belongs to
\(Q_{i+1}\cap G_{j_0+1}\), so
Lemma~\ref{lem:real-pairing-floor-laws} and the measurable lifts
of Proposition~\ref{prop:actual-floor-quadrics} give \(F\).
Its integral in absolute value is finite by the same lemma.
In the recurrence for \(K_i\), conjugation by \(\sigma\) changes
\(n_i(t,\beta)\) into \(n_i(\sigma t,\beta)\) by the signed
isometry property. This is an element of \(N^i\) commuting with
\(d_*\), since \(i>j_0\) and its action is confined to \(V_i\).
Its only possible simple-root entry is the \(e_i\)-coordinate
of the image of \(f_{i+1}\), which equals
\(-[\sigma t,f_{i+1}]=-[t,v]\) by
Definition~\ref{def:real-radical-weights}. Apply the independent
translation rule in Lemma~\ref{lem:real-pairing-floor-laws} to
the second argument of \(K_{i+1}\). The orbit law in
Proposition~\ref{prop:actual-floor-quadrics} gives the formula.
\end{proof}


\section{A Fourier estimate centered at an admissible translate}

\begin{definition}\label{def:real-position-indicator}
For \(L\ge1,\ 0\le i\le\nu\) write \(\chi_i(g,y)\) for the
indicator of the existence of \(u\in N^i\) satisfying
\(\|g^{-1}u y-I\|_{\rm op}\le L\). This is a Borel indicator.
Indeed if \(g_a,y_a\) converge and \(u_a\) are witnesses, then
\(g_a^{-1}u_a y_a\) lies in a compact set. Thus a subsequence
of \(u_a\) converges in the closed subgroup \(N^i\) and its
limit is a witness.
\end{definition}

\begin{lemma}\label{lem:real-centered-radical}
Take \(1\le i\le\nu,\ g\in\mathcal N_i,\ y\in G_1\),
\(L\ge\max(4,r,R)\). Choose \(s\in Q_i\) with \(s g\in B_G(R)\)
and denote by \(h\in G_{i+1}\) its action on the middle quotient.
If \(\sigma\in G_{i+1}\) and \(\sigma g\in\mathcal N_{i+1}\),
then \(\|h\sigma^{-1}f_{i+1}\|\le R^2\). Suppose
\(\chi_{i-1}(g,n_i(t_0,\beta_0)y)=1\).
For every \((t,\beta)\) with \(\chi_{i-1}(g,n_i(t,\beta)y)=1\),
\(\|h(t-t_0)\|\le L^6\). For each fixed \(t\) the admissible
central coordinates \(\beta\) belong to an interval of length
at most \(2L^6\).
\end{lemma}
\begin{proof}
Choose \(s'\in Q_{i+1}\) with \(s'\sigma g\in B_G(R)\).
Then \(s\sigma^{-1}(s')^{-1}=(sg)(s'\sigma g)^{-1}\) has norm
at most \(R^2\). Applied to \(f_{i+1}\), its ordinary projection
to \(V_{i+1}\) is \(h\sigma^{-1}f_{i+1}\).
In fact \((s')^{-1}f_{i+1}-f_{i+1}\) is in the span of the first
\(i\) vectors, fixed by \(\sigma^{-1}\) and preserved by \(s\).
Also \(s|_{V_{i+1}}\) agrees with \(h\) modulo the same span.
This proves the first bound.

For an admissible \((t,\beta)\) choose its witness
\(u_t\in N^{i-1}\) and use \(u_0\) for the chosen base witness.
In the expression
\[
 s\,(u_t n_i(t,\beta))(u_0 n_i(t_0,\beta_0))^{-1}s^{-1}
\]
insert \(g,g^{-1}\) at the two extremes and \(yy^{-1}\) between
the two middle factors. The two bounded positions and \(sg\in B_G(R)\)
give norm at most \(R^2(1+L)^2\le L^6\).
On the middle quotient obtained by removing the first \(i-1\)
pairs, \(u_t,u_0\) are the identity. The image of \(s\) there
is \(h n_i(t_s,\beta_s)\) for some \(t_s,\beta_s\), by the
layer splitting of Lemma~\ref{lem:heisenberg-layer}. Multiplication
in the layer adds the vector variables and sends the central
coordinate of a product to
\(\beta+\beta'-\operatorname{Im}[t,t']\), by multiplying the
corner entries of Definition~\ref{def:real-radical-weights}.
Conjugating by \(h\) sends \(t\) to \(ht\) and fixes the real
coordinate. Thus the quotient of the displayed expression is
\(n_i(h(t-t_0),\beta-\beta_0+b(t))\), for a real function
\(b(t)\) independent of \(\beta\). This quotient is an ordinary
diagonal block of the full matrix, hence has norm at most \(L^6\).
Its last column bounds both \(\|h(t-t_0)\|\) and the modulus of
the imaginary corner coordinate. The latter yields the claimed
interval.
\end{proof}

\begin{theorem}\label{thm:real-centered-plancherel}
For \(j_0<i\le\nu\) in
Definition~\ref{def:real-successive-pairing}, \(g,y\in G_1\)
and \(L\ge\max(4,r,R)\), one has
\[
 \begin{split}
 \int\chi_{i-1}(g,n_i(t,\beta)y)|K_i(g,n_i(t,\beta)y)|^2
                                     \dd t\,\dd\beta
    \le {}&\pi^{2m}L^{38n}\chi_i(g,y)\\
      &{}\cdot\int|K_{i+1}(\sigma g,\sigma y)|^2
                                      \dd\theta_{i+1}(\sigma).
 \end{split}
\]
Here \(m=k-1\), and if the indicator on the right vanishes the
right side is taken to be zero.
\end{theorem}
\begin{proof}
If there is an admissible pair on the left then \(\chi_i(g,y)=1\)
by Lemma~\ref{lem:partial-real-radical-measure}. If \(g\notin
\mathcal N_i\) the left side vanishes by
Lemma~\ref{lem:real-pairing-floor-laws}. In the remaining
case use \(h,t_0,\beta_0\) of Lemma~\ref{lem:real-centered-radical}.
By Proposition~\ref{prop:real-radical-fourier-phase} the pairing
has the form \(\int\psi(-[t,v])F(v)\dd\delta_\gamma^{(i+1)}(v)\),
independent of \(\beta\). It has amplitude zero unless
\(\sigma g\in\mathcal N_{i+1}\), so after the change \(v'=hv\)
its support lies in \(\|v'\|\le R^2\). This change preserves
the quadric measure by Lemma~\ref{lem:real-flat-disintegration}.
The function \(F(h^{-1}v')\) is bounded and has compactly
bounded support, hence belongs to \(L^2\). Its squared integral is
the \(\theta_{i+1}\)-integral on the right by
Proposition~\ref{prop:actual-floor-quadrics}.

Integrate the indicator first in \(\beta\), at a cost of
\(2L^6\), and put \(t'=h(t-t_0)\), a change of real determinant
one. The resulting output range lies in \(\|t'\|\le L^6\).
The factor \(\psi(-[ht_0,v'])\) has modulus one and can be
absorbed in the amplitude without changing its squared norm.
Apply Theorem~\ref{thm:real-flat-fourier} on \(V_{i+1}\) with
ambient output, quadric input and frequency \(-1\).
The Euclidean balls are contained in the flat boxes of the
same radii. The cost multiplying the squared amplitude is at most
\[
  2L^6 M_{\rm amb}(L^6) M_{\rm quad}(R^2)\pi^{2m}
       \le L^{7+26(n-i)+6+10(n-i-1)}\pi^{2m}
       \le L^{38n}\pi^{2m}.
\]
This proves the assertion, including cases in which the integral
on the left is zero.
\end{proof}


\section{Squared norms of successive pairings on a ball}

\begin{definition}\label{def:real-ball-energies}
Use Definition~\ref{def:real-successive-pairing} and its twist \(d_*\).
Write \(\tau_i=\theta_i*\cdots*\theta_1\), with \(\tau_0\) unit mass
at the identity. For \(0\le l\le\nu\) put
\(\rho_l=\eta_{l+1}*\cdots*\eta_\nu*\mu_{\rm up}\), so
\(\rho_\nu=\mu_{\rm up}\). With the radius \(L\) used to define
\(\chi_l\), set
\[
 \begin{aligned}
 A_l&=\iint\chi_l(g,y)|K_{l+1}(g,y)|^2\dd\rho_l(y)\dd\tau_{l+1}(g),\\
 Z_l(B,d_*;x)&=\int\mathcal W_l(g,gx)|K_{l+1}(g,gx)|^2
                                                    \dd\tau_{l+1}(g).
 \end{aligned}
\]
These are nonnegative integrals. The definition of \(Z_l\) is
independent of \(L\).
\end{definition}

\begin{lemma}\label{lem:partial-real-quotient}
For \(0\le l\le\nu\) one has
\(\dd u_{N^l}*\rho_l=\lambda\dd y\).
If \(h\) normalizes \(N^l\) and the change \(u\mapsto huh^{-1}\)
has Haar Jacobian \(\delta(h)\), every nonnegative left
\(N^l\)-invariant Borel function \(A\) satisfies
\[
          \int A(hy)\dd\rho_l(y)=\delta(h)\int A(y)\dd\rho_l(y).
\]
For \(h\in G_{l+1}\) this Jacobian equals one. For \(l=\nu\)
and \(h=d_*\) it is
\(\delta(d_*)=\prod_{a=1}^{j_0}|r_a|^{\,2d_a+2}\).
\end{lemma}
\begin{proof}
Outside-to-inside multiplication of the first \(l\) layer measures
gives the measure of \(N^l\) by
Lemma~\ref{lem:partial-real-radical-measure}. Including the remaining
layers gives \(\dd u_{\mathcal U}*\mu_{\rm up}\), which is
\(\lambda\dd y\) by Theorem~\ref{thm:exact-real-floor-haar}.
To obtain a cutoff on \(N^l\), divide
\(\zeta(y)=\exp(-\|y\|_2^2)\) by \(\int_{N^l}\zeta(uy)\dd u\).
The denominator is positive, finite and left \(N^l\)-invariant.
Indeed the matrix equations define a closed subgroup, and the
polynomial ball estimate of Lemma~\ref{lem:period-ball-cover}
gives convergence locally uniformly in \(y\), exactly as in
Lemma~\ref{lem:bessel-quotient-density}. Denote the quotient
by \(a_l\); then \(\int a_l(uy)\dd u=1\).
Applying the convolution formula to \(A a_l\) gives
\(\int A\dd\rho_l=\lambda\int A(y)a_l(y)\dd y\).
Use it first for \(A(h\,\cdot)\), change variables by left Haar
invariance, and use the convolution formula again:
\[
 \int A(hy)\dd\rho_l(y)
  =\lambda\int A(z)a_l(h^{-1}z)\dd z
  =\int A(z)\int a_l(h^{-1}u z)\dd u\,\dd\rho_l(z).
\]
Putting \(u=h v h^{-1}\) makes the inner integral \(\delta(h)\).
The argument applies also at infinite integral by
Lemma~\ref{lem:product-integration}. For \(G_{l+1}\) the Jacobian
is the one calculated in Lemma~\ref{lem:partial-real-radical-measure}
(and is one trivially if \(l=0\)).
On layer \(a\), conjugation by \(d_*\) sends the vector and central
coordinate to \((\bar r_a\,d_*|_{V_{a+1}} t, |r_a|^2\beta)\),
by evaluating on \(e_a^*\) in Definition~\ref{def:real-radical-weights}.
The middle isometry has complex determinant of modulus one;
the complex vector space has dimension \(d_a\).
Multiplying the layer Jacobians, including the central real
coordinate of each, proves the last assertion.
\end{proof}

\begin{theorem}\label{thm:real-ball-energies}
At sufficiently large arithmetic scales take \(L\ge\max(4,r,R)\).
Suppose that \(B\) in Definition~\ref{def:real-successive-pairing}
has squared quotient norm at most \(e^2D\). For each twist one has
\[
 \int_{\|x-I\|_{\rm op}\le L}
     \sum_{l=j_0}^{\nu} Z_l(B,d_*;x)\,\lambda\dd x
 \le n\,\pi^{2m\nu}L^{70n^2}\lambda^2D^2e^2\,\delta(d_*),
 \qquad m=k-1 .
\]
\end{theorem}
\begin{proof}
At \(l=\nu\), Lemma~\ref{lem:real-pairing-floor-laws} and
Cauchy--Schwarz in the value space give
\[
 A_\nu\le\int\|V(g)\|^2\dd\tau_n(g)
                         \int\|B(d_*y)\|^2\dd\rho_\nu(y)
          \le\lambda^2D^2e^2\delta(d_*).
\]
Indeed \(\kappa*\tau_n=\mu_{\rm up}\) and \(\|V\|\) is compact
invariant. Theorem~\ref{thm:exact-real-floor-haar} converts each
squared norm to its quotient integral with factor \(\lambda\);
Lemma~\ref{lem:partial-real-quotient} removes the twist from
the second. Use Theorem~\ref{thm:actual-bessel-truncation} for
the first quotient norm.

For \(j_0<i\le\nu\) write \(\rho_{i-1}=\eta_i*\rho_i\) in
\(A_{i-1}\). Theorem~\ref{thm:real-centered-plancherel} bounds
the inner integral over \(\eta_i\) by
\(\pi^{2m}L^{38n}\chi_i(g,y)
\int |K_{i+1}(\sigma g,\sigma y)|^2\dd\theta_{i+1}(\sigma)\).
Here \(\chi_i(\sigma g,\sigma y)=\chi_i(g,y)\) as
\(\sigma\in G_{i+1}\) normalizes \(N^i\).
For fixed \(\sigma,g\), the function
\[
       y'\longmapsto\chi_i(\sigma g,y')
                                |K_{i+1}(\sigma g,y')|^2
\]
is left \(N^i\)-invariant: absorb the multiplier into the
witness for the indicator and use the unit-modulus phase of
Lemma~\ref{lem:real-pairing-floor-laws}. Hence
Lemma~\ref{lem:partial-real-quotient} permits the replacement of
\(\sigma y\) by \(y'\) under \(\rho_i\). The measure on \(\sigma g\)
is \(\theta_{i+1}*\tau_i=\tau_{i+1}\).
This proves \(A_{i-1}\le\pi^{2m}L^{38n}A_i\).

Finally substitute \(y'=gx\) at fixed \(g\) in the integral of
\(Z_l\) over the stated ball. Decompose
\(\lambda\dd y'=\dd u_{N^l}*\rho_l\) by
Lemma~\ref{lem:partial-real-quotient}. The modulus of \(K_{l+1}(g,uy)\)
equals that of \(K_{l+1}(g,y)\). When this is positive,
\(g\in\mathcal N_{l+1}\subset\mathcal N_l\) by
Lemma~\ref{lem:real-pairing-floor-laws}.
Theorem~\ref{thm:real-positive-radical-weight} bounds the integral
of the remaining indicator times \(\mathcal W_l(g,uy)\) in \(u\)
by \(L^{70nl}\chi_l(g,y)\); if \(\chi_l=0\) it vanishes.
The contribution is therefore at most \(L^{70nl} A_l\).
All integrals rearranged here are nonnegative, so
Lemma~\ref{lem:product-integration} applies without a preliminary
integrability assertion. Iterate the bound for \(A_l\) and sum,
using \(70nl+38n(\nu-l)\le70n^2\) and \(\nu+1=n\).
\end{proof}


\section{Stopping the real integrals at the first small area}

\begin{definition}\label{def:partial-real-stopping-energy}
Fix \(B,d_*\) as in Definition~\ref{def:real-successive-pairing},
\(0\le j\le\nu\), and \(g,x\in G_1\). Write
\[
 U_j(g,x)=\int |K_{j+1}(h g,h g x)|\dd\theta_{j+1}(h).
\]
For \(j\le l\le\nu\) let
\[
 E_{j,l}(g,x)=
 \int \mathcal W_l(hg,hgx)|K_{l+1}(hg,hgx)|^2
       \dd(\theta_{l+1}*\cdots*\theta_{j+1})(h).
\]
In particular \(E_{0,l}(1,x)=Z_l(B,d_*;x)\) in
Definition~\ref{def:real-ball-energies}.
\end{definition}

\begin{proposition}\label{prop:real-area-stopping}
At sufficiently large arithmetic scales, for every \(g,x\) and \(j\)
in Definition~\ref{def:partial-real-stopping-energy} one has
\[
 \mathcal W_j(g,gx) U_j(g,x)^2
       \le R^{29n(\nu-j)+5k} E_{j,\nu}(g,x).
\]
Let \(u>0\). Suppose that for every \(h\in G_{j+1}\) for which
\(\langle V(hg),B(d_*h g x)\rangle\ne0\), there is some \(j<i\le\nu\)
with \(v_i(hg,x)\le u\). Then
\[
 \mathcal W_j(g,gx) U_j(g,x)^2
       \le R^{29n(\nu-j)}u^s\sum_{l=j}^{\nu-1}E_{j,l}(g,x).
\]
An empty sum is zero.
\end{proposition}
\begin{proof}
Lemma~\ref{lem:real-pairing-floor-laws} gives \(|K_i(g,gx)|\le
U_i(g,x)\) for \(0\le i\le\nu\). In each integral defining
\(U_{i-1}\) the integrand vanishes unless \(h g\in\mathcal N_i\).
For \(i\le\nu\) insert the reciprocal factors
\(v_i(hg,x)^{-s/2}\) and \(v_i(hg,x)^{s/2}\).
Cauchy--Schwarz and Theorem~\ref{thm:real-floor-mass} give
\[
 U_{i-1}(g,x)^2
    \le R^{29n}\int v_i(hg,x)^s |K_i(hg,hgx)|^2\dd\theta_i(h).
\]
For \(i=n\) use Cauchy--Schwarz with unit weight on the support.
The cost is at most \(R^{5k}\) by the same theorem. These
inequalities hold also when a nonnegative right side is infinite,
by applying Cauchy--Schwarz on increasing bounded sets.
At each successive integration \(h\in G_i\subset Q_a\) for
\(a<i\). The preceding factors \(v_a(\,\cdot\,,x)\) are unchanged
under this left multiplication by Lemma~\ref{lem:real-near-floor-invariance}.
Starting at \(i=j+1\), iterating up to \(n\) and multiplying by
\(\mathcal W_j(g,gx)\) now yields the first inequality.

For the second, at each \(i\le\nu\) split the preceding weighted
integral according as its newly encountered area is at most \(u\)
or exceeds \(u\). On the first set bound the new factor by \(u^s\)
and retain \(|K_i|^2\) without expanding it at another stage.
Removing the previous restrictions increases this nonnegative
contribution. Its bound, including the initial weight, is
\(R^{29n(i-j)}u^s E_{j,i-1}(g,x)\). On the second set keep the
new weight and continue. The restrictions on earlier areas are
preserved under subsequent left multiplications as above.
If all the areas encountered up to \(\nu\) exceed \(u\), they
still do so after the final \(\theta_n\)-multiplication. The full
left multiplier is a product of elements of the nested groups
\(G_{j+1},\ldots,G_n\), hence lies in \(G_{j+1}\).
The hypothesis then makes \(K_n\) zero by
Lemma~\ref{lem:real-pairing-floor-laws}. This also proves the
claim directly when \(j=\nu\). Summing the stopped contributions
and bounding \(i-j\) by \(\nu-j\) gives the result.
\end{proof}

\begin{corollary}\label{cor:large-real-translation-energy}
For the empty head let \(c_B(x)=\int\langle V(y),B(yx)\rangle
\dd\bar y\), with \(B=V,F\), or \(V-F\). At sufficiently large
scales, for every \(x\in G_1\),
\[
 \lambda^2 |c_B(x)|^2\le R^{29n\nu+5k} Z_\nu(B,1;x).
\]
If \(T\ge\max(4,r,R)\) and \(\|x-I\|_{\rm op}>T^{16n}\), then
\[
 \lambda^2 |c_V(x)|^2\le R^{29n\nu}(\pi/T^2)^s
                                  \sum_{l=0}^{\nu-1}Z_l(V,1;x).
\]
\end{corollary}
\begin{proof}
In the empty head \(P_0=\mu_{\rm up}\) and \(d_*=1\).
The two factors of the inner product have the same unitary
transformation on \(S_0\). Its quotient integral is absolutely
convergent by the two quotient norm bounds in
Theorem~\ref{thm:actual-bessel-truncation} and right invariance
in Lemma~\ref{lem:bessel-quotient-density}. Applying
Theorem~\ref{thm:exact-real-floor-haar} to the absolute value and
then to the scalar gives \(K_0(1,x)=\lambda c_B(x)\).
Use \(|K_0(1,x)|\le U_0(1,x)\) and the first inequality of
Proposition~\ref{prop:real-area-stopping}. For \(B=V\), a nonzero
raw pairing at \(h,hx\) implies \(S(h),S(hx)\le R\) by
Theorem~\ref{thm:actual-bessel-truncation}. At \(j=0\) both
near conditions are automatic since \(Q_0=G_1\). Thus
Theorem~\ref{thm:real-large-position-cap} gives exactly the
hypothesis of the second stopping inequality with \(u=\pi/T^2\).
\end{proof}


\section{The full real energy of the cut coefficient}

\begin{lemma}\label{lem:shells-real-cut-mixed}
At sufficiently large scales set \(U=(2R)^{16n}\) and
\[
 \mathcal A=n(1+2^{70n^2+1})R^{29n\nu+5k}U^{70n^2},
 \qquad C_{\rm B}=\frac{u_k}{u_r}\pi^{(2\nu+1)k}.
\]
Let \(B\) satisfy Theorem~\ref{thm:real-ball-energies} for the empty
head. Define \(c_B\) by the same quotient integral as in
Corollary~\ref{cor:large-real-translation-energy}. Then
\[
            \int_{G_1}|c_V(x)c_B(x)|\dd x
                           \le D^2 e\, C_{\rm B}\mathcal A .
\]
\end{lemma}
\begin{proof}
The integral for \(c_B\) is absolutely convergent by the quotient
Cauchy--Schwarz argument in Corollary~\ref{cor:large-real-translation-energy}.
Its first pointwise bound extends to this \(B\): that proof uses
only the common equivariance and the finite squared quotient norm.
Write \(m_+=k-1\) and \(C=R^{29n\nu+5k}\).
On the ball of radius \(U\) about the identity, apply that pointwise
bound to both coefficients, followed by Cauchy--Schwarz in \(x\)
and Theorem~\ref{thm:real-ball-energies} for each function.
Since \(e=1\) is valid for \(V\), the integral on this ball is
at most
\[
            D^2e\,\lambda^{-1} n\pi^{2m_+\nu} C U^{70n^2}.
\]
The factor \(\lambda^{-1}\) comes from using \(\lambda\dd x\)
in the squared norm estimate and from the \(\lambda^2\) on each
left side of Corollary~\ref{cor:large-real-translation-energy}.

On the shell \(2^aU<\|x-I\|_{\rm op}\le2^{a+1}U\), \(a=0,1,\ldots\),
use \(T_a=(2^aU)^{1/(16n)}\). Here \(T_a\ge2R\ge\max(4,r,R)\)
eventually by Proposition~\ref{prop:thick} and
Theorem~\ref{thm:actual-bessel-truncation}.
Apply the second pointwise bound of
Corollary~\ref{cor:large-real-translation-energy} to \(c_V\) and
the first to \(c_B\). Enlarge the integrals of the nonnegative
\(Z_l\)'s to the ball of radius \(2^{a+1}U\).
The preceding estimate for the inner ball, multiplied by
\[
    2^{(a+1)70n^2}(\pi/T_a^2)^{s/2}
    \le 2^{70n^2}2^{-a(s/(16n)-70n^2)},
\]
bounds the mixed integral on this shell, since
\(\sqrt\pi/(2R)\le1\). Now \(s=k-3n-2\) by
Theorem~\ref{thm:real-floor-mass}, \(r=k+2n-1\),
\(r\ge x_{\rm ar}^5\) and \(n\le\sqrt{x_{\rm ar}}\), where
\(x_{\rm ar}\) denotes the arithmetic scale of Proposition~\ref{prop:thick}.
Hence \(s/(16n)-70n^2\ge1\) eventually. The geometric sum is at
most two. Finally \(\lambda^{-1}\pi^{2m_+\nu}\le C_{\rm B}\)
by Theorem~\ref{thm:exact-real-floor-haar}. Adding the inner ball
and all shells gives the claim by Lemma~\ref{lem:integral-tools}.
\end{proof}

\begin{theorem}\label{thm:unrestricted-real-cut-energy}
Let \(d=\mathcal D_\tau\) be the formal degree of the actual row
at a real place \(\tau\) in Theorem~\ref{thm:bessel-orthonormal-schur}.
For each fixed \(0<\eta<1\) and at all sufficiently large arithmetic
scales, uniformly in the place and the permitted window parameter,
\[
 \frac{1}{D^2}\int_{G_1}|c_V(x)|^2\dd x
       \le\frac{1+b^{-(5-\eta)r^2}}{d}
       \le\frac{1+b^{-10}}{d}.
\]
\end{theorem}
\begin{proof}
Take \(B=V-F\). Its admissible norm parameter is \(e=b^{-5r^2}\)
by Theorem~\ref{thm:actual-bessel-truncation}. Expanding
\(|c_F|^2=|c_V-c_B|^2\) gives
\(|c_V|^2\le |c_F|^2+2|c_Vc_B|\).
Use Theorem~\ref{thm:finite-real-energy-bound} on the first term
and Lemma~\ref{lem:shells-real-cut-mixed} on the second. This gives
\[
 \frac{d}{D^2}\int|c_V|^2
                  \le 1+2 b^{-5r^2}d C_{\rm B}\mathcal A .
\]
We estimate the constants using \(X\) for the arithmetic scale.
Since \(R=b^{6n+2}\), \(\log_b U=16n\log_b(2R)=O(n^2)\).
Consequently
\(\log_b(2\mathcal A)=O(n^4+nk+n^3)=o(r^2)\) by
Proposition~\ref{prop:thick}, since \(n\le\sqrt X\) and
\(k\le r,\ r\ge X^5\). All constants here are absolute.
The constant \(C_{\rm B}\) equals the one in
Definition~\ref{def:row-vandermonde}, whose parameter \(m\) is
\(2n\). By Proposition~\ref{prop:exact-row-degree},
\[
 \log_b(d C_{\rm B})
 \le 2n(k-1)\log_b w(\tau)
       + O\left(\frac{nk/X+n^2\log X}{\log b}\right).
\]
Proposition~\ref{prop:actual-bessel-weights} and its proof give
\(1\le w(\tau)\le W_\tau\le 2b^2\). Thus the last upper bound
is \(o(r^2)\) uniformly in the indicated data. For the fixed
\(\eta\) this gives \(2d C_{\rm B}\mathcal A\le b^{\eta r^2}\)
eventually. The second asserted inequality follows from
\((5-\eta)r^2\ge10\) at the same scales.
\end{proof}


\section{Integration over the preceding hyperbolic pairs}

\begin{definition}\label{def:real-preceding-measure}
For \(0\le j\le\nu\) put
\[
 \begin{split}
 \mathfrak m_j&=(\eta_j*\theta_j)*\cdots*(\eta_1*\theta_1),\\
 \mathfrak h_j&=\kappa*\theta_n*
          (\eta_\nu*\theta_\nu)*\cdots*(\eta_{j+1}*\theta_{j+1}).
 \end{split}
\]
An empty product of measures is unit mass at \(1\).
For a kernel of initial length \(j_0\le j\), write
\(S_j(g,x)=\sum_{l=j}^{\nu}E_{j,l}(g,x)\) and
\(Z_{\ge j}(x)=\sum_{l=j}^{\nu}Z_l(B,d_*;x)\), using
Definition~\ref{def:partial-real-stopping-energy}. Write
\(W_j(g,x)=\mathcal W_j(g,gx)\).
\end{definition}

\begin{lemma}\label{lem:real-preceding-translation}
For every \(j\) in Definition~\ref{def:real-preceding-measure},
\(\mathfrak m_j=\dd u_{N^j}*\tau_j\) and
\(\mathfrak h_j*\mathfrak m_j=\lambda\dd g\).
The measure \(\mathfrak h_j\) is carried by \(G_{j+1}\) and is right
invariant under that group. For a nonnegative Borel function \(A\)
invariant on the left by \(G_{j+1}\) and any \(w\) commuting with
\(G_{j+1}\), one has
\[
                 \int A(wg)\dd\mathfrak m_j(g)
                            =\int A(g)\dd\mathfrak m_j(g).
\]
The particular choices \(w=w_p^{\pm1}\) are permitted for a head
of length \(j\).
\end{lemma}
\begin{proof}
In the defining convolution for \(\mathfrak m_j\) move the first
\(j-1\) radicals to the left of \(\eta_j*\theta_j\) by induction:
both factors are supported on \(G_j\), which preserves measure on
\(N^{j-1}\) by Lemma~\ref{lem:partial-real-radical-measure}.
The leftmost factors now give the product on \(N^j\) of that lemma,
and the rest give \(\tau_j\). The full convolution is exactly
\(\mu_{\rm all}=\lambda\dd g\) of Theorem~\ref{thm:exact-real-floor-haar}.

At \(j=\nu\), right invariance of \(\mathfrak h_\nu\) is
Lemma~\ref{lem:middle-haar-cutoff}. Descend from \(j=i\) to \(i-1\).
The measure \(\mathfrak h_i*\eta_i\) is right invariant under the
stabilizer of \(e_i\) in \(G_i\). Indeed the stabilizer is
\(G_{i+1}\) times the \(i\)-th layer by
Lemma~\ref{lem:partial-real-radical-measure} on \(V_i\):
an isometry fixing \(e_i\) acts by an isometry on its perpendicular
modulo that line, and the kernel consists of exactly the matrices
\(n_i(t,\beta)\). Right translation by this kernel preserves
\(\eta_i\). Right translation by \(G_{i+1}\) can be moved past
\(\eta_i\) by its determinant-one conjugation and then preserves
\(\mathfrak h_i\) by induction. Averaging a nonnegative Borel test
on the left by \(\mathfrak h_i*\eta_i\) therefore gives a function
of the orbit vector \(g^{-1}e_i\). This function is Borel on the
full-measure chart by the section in
Proposition~\ref{prop:actual-floor-quadrics}. Right translation by
an element of \(G_i\) preserves the law of this vector under
\(\theta_i\), including the null complement of the chart. Thus
\(\mathfrak h_{i-1}=\mathfrak h_i*\eta_i*\theta_i\) has the required
right invariance. Support follows from the nested groups.

For the last identity choose a strictly positive Haar integrable
Borel function \(\phi\) on \(G_1\). For example sum the indicators
of the norm balls of positive integral radii \(a=1,2,\ldots\), with
coefficients \(2^{-a}/(1+|B_G(a)|)\); the balls are compact and have
finite Haar measure by Lemma~\ref{lem:unitary-cayley-density}.
Put \(H(g)=\int\phi(hg)\dd\mathfrak h_j(h)\), which is positive
since the defining product has positive measure. It is left
\(G_{j+1}\)-invariant by the right invariance just proved.
Both \(H(g)\) and \(H(wg)\) are finite for \(\mathfrak m_j\)-almost
every \(g\). Indeed their integrals equal \(\lambda\int\phi\dd g\),
using the full convolution, commutation with \(w\), and left Haar
invariance. Define \(a(g)=\phi(g)/H(g)\) when \(H(g)<\infty\) and
zero elsewhere. The averages of \(a(hg)\) and \(a(whg)\) over
\(\mathfrak h_j\) are consequently one for almost every such \(g\).
Inserting the respective averages and using
Lemma~\ref{lem:product-integration} identifies the two integrals
in the statement with \(\lambda\int A(wy)a(wy)\dd y\) and
\(\lambda\int A(y)a(y)\dd y\). Left invariance identifies them.
All \(w_p^{\pm1}\) act as the identity on \(V_{j+1}\) and preserve
its complementary pairs, proving the commutation.
\end{proof}

\begin{lemma}\label{lem:real-tail-square-invariance}
For fixed \(x\), \(S_j(\,\cdot\,,x)\) is nonnegative Borel, is
left \(Q_j\)-invariant and vanishes outside \(\mathcal N_j\).
Its integral against \(\tau_j\) equals \(Z_{\ge j}(x)\).
\end{lemma}
\begin{proof}
In the definition of \(E_{j,l}\), the function at its last stage is
\(W_l(y,x)|K_{l+1}(y,yx)|^2\). Its invariance under \(Q_{l+1}\)
follows from Lemmas~\ref{lem:partial-real-radical-measure}
and~\ref{lem:real-pairing-floor-laws}, since \(l+1>j_0\).
Both factors are Borel and the pairing requires \(y\in\mathcal N_{l+1}\).
Applying Lemma~\ref{lem:real-near-floor-invariance} to each
successive integration gives invariance under \(Q_j\) at the end.
The support condition descends because
\(Q_i,G_i\subset Q_{i-1}\) and \(\theta_i\) is carried by \(G_i\).
Finally the product with \(\tau_j\) is precisely \(\tau_{l+1}\).
Integrate and sum using Lemma~\ref{lem:product-integration}.
\end{proof}

\begin{proposition}\label{prop:real-crossed-square-bound}
Fix heads \(p,q\) of length \(j\). In
Definition~\ref{def:real-preceding-measure}, let \(S,Z\) denote
\(S_j,Z_{\ge j}\) for \(V,1\), and \(S',Z'\) the corresponding
functions for \(B,d_*\) with \(d_*=w_q^{-1}w_p\).
Let \(L\ge\max(4,r,R)\) bound every
\(|p_1\cdots p_a|^{-1}\), \(1\le a\le j\), at sufficiently large scales.
Fix \(x\in G_1\). Suppose \(f\) is nonnegative Borel, vanishes
unless both \(g,w_p^{-1}g\in\mathcal N_j\), and on its support obeys
\[
 f(g)^2\le c^2\,
     \frac{S(g,x)}{W_j(g,x)}
     \frac{S'(w_p^{-1}g,x)}{W_j(w_p^{-1}g,x)},\qquad c>0.
\]
Whenever \(Z(x),Z'(x)<\infty\),
\[
 \int f\dd\mathfrak m_j
    \le c L^{125n^2j}\prod_{a=1}^j|p_a|^{-2d_a-1}
                                      \sqrt{Z(x)Z'(x)} .
\]
\end{proposition}
\begin{proof}
In Cauchy--Schwarz pair \(S(g,x)\) with the reciprocal weight
and the near indicator at \(w_p^{-1}g\). Pair \(S'(w_p^{-1}g,x)\)
with the reciprocal weight and the near indicator at \(g\).
For the first integral use
Lemma~\ref{lem:real-preceding-translation} and
Corollary~\ref{cor:real-radical-test} with \(\rho'=\tau_j\)
and the head \(\bar p\), since \(w_p^{-1}=w_{\bar p}\).
This gives \(L^{125n^2j}\prod |p_a|^{-2d_a-1} Z(x)\) by
Lemma~\ref{lem:real-tail-square-invariance}. In the second integral
the integrand is \(A(w_p^{-1}g)\) with
\[
      A(y)=S'(y,x)\,{\bf1}_{\mathcal N_j}(w_p y)/W_j(w_p y,x).
\]
This function is left \(G_{j+1}\)-invariant since this group commutes
with \(w_p\), preserves the weights and near sets, and fixes \(S'\)
by the same lemma. Apply Lemma~\ref{lem:real-preceding-translation}
to replace its argument by \(g\), then apply
Corollary~\ref{cor:real-radical-test} with head \(p\).
The second squared integral has the identical constant times \(Z'(x)\).
The measure \(\tau_j\) is sigma-finite here: \(\mathfrak m_j\)
integrates the positive \(H\) in the preceding proof to a finite
number, and \(\mathfrak m_j=\dd u_{N^j}*\tau_j\). Integrating \(H(ug)\)
in \(u\) gives a positive function with finite \(\tau_j\)-integral.
Both applications of the corollary are therefore valid, including
when a squared integral is zero. Taking the two square roots
proves the claim.
\end{proof}


\section{Absolute convergence for arbitrary real heads}

\begin{definition}\label{def:real-pivot-forms}
Fix \(0\le j\le\nu\), \(p,q\in(\C^\times)^j\), and put \(i=j+1\).
In Definition~\ref{def:real-successive-pairing} use \(B=V\)
throughout. Use \(K_a\) for the empty head, and \(K'_a\) for
the twist \(d_*=w_q^{-1}w_p\). Set \(P_j(p)=\prod_{a\le j}|p_a|^{d_a}\),
\(\kappa_j=P_j(p)P_j(q)\), and \(C_j=\kappa_j/(\lambda D)^2\).
For functions of \(\sigma,t\in G_i\) write
\[
 H(\sigma)=K_i(\sigma g,\sigma gx),\qquad
 H'(t)=K'_i(tw_p^{-1}g,tw_p^{-1}gx),
\]
with \(g,x,p,q\) understood. Define the following functions of
\(g,x\), using \(f_i=e_i\) if \(i<n\), \(f_n=z\):
\[
 \begin{split}
 B_j(g,x)&=\iint |H(\sigma)H'(t)|\dd\theta_i(t)\dd\theta_i(\sigma),\\
 B_j^\xi(g,x)&=\iint\psi\bigl(\xi[\sigma^{-1}f_i,t^{-1}f_i]\bigr)
       H(\sigma)\overline{H'(t)}\dd\theta_i(t)\dd\theta_i(\sigma).
 \end{split}
\]
Here \(\xi\in\C\). These integrals are finite at every \(g,x\)
by Lemma~\ref{lem:real-pairing-floor-laws} and
Theorem~\ref{thm:real-floor-mass}. Write
\(\mathcal I_j^\xi(p,q)=C_j\iint B_j^\xi(g,x)\dd\mathfrak m_j(g)\dd x\)
whenever absolutely convergent. Use \(J_j^*(p,q)=\mathcal I_j^0(p,q)\).
\end{definition}

\begin{lemma}\label{lem:real-head-ball-majorant}
At sufficiently large scales suppose
\[
 L\ge\max(4,r,R,|D_{p,a}|^{-1},|D_{q,a}|^{-1}:0\le a\le j),
 \qquad L'\ge\max(4,r,R),
\]
where \(D_{p,a}=p_1\cdots p_a\), \(D_{p,0}=1\), and similarly for
\(q\). For any Borel \(X\subset\{x:\|x-I\|_{\rm op}\le L'\}\),
\[
 C_j\int_X\int B_j(g,x)\dd\mathfrak m_j(g)\dd x
 \le n C_{\rm B}L^{125n^2j+1}(L')^{70n^2}R^{29n\nu+5k}.
\]
If \(T\ge\max(4,r,R,L^{4j+1})\) and every \(x\in X\) satisfies
\(\|x-I\|_{\rm op}>T^{16n}\), the same bound holds multiplied by
\((\pi/T^2)^{s/2}\). If either head has \(|D_{p,a}|>R^{2a}\)
or \(|D_{q,a}|>R^{2a}\) for some \(a\le j\), then \(B_j\) vanishes
identically.
\end{lemma}
\begin{proof}
Write \(S,Z,S',Z'\) as in Proposition~\ref{prop:real-crossed-square-bound}.
The two factors of \(B_j\) are exactly the \(U_j\)'s of
Definition~\ref{def:partial-real-stopping-energy} at \(g\) and
\(w_p^{-1}g\). Positivity of the first yields some \(h\in G_i\)
with both \(V(hg),V(hgx)\ne0\). To see this, dominate the averaged
pairing by the integral of the absolute raw pairing and use that
the upper multipliers lie in \(G_i\). Thus \(g,gx\in\mathcal N_j\).
Positivity of the second yields \(h'\in G_i\) with
\(V(h'w_p^{-1}g),V(h'w_q^{-1}gx)\ne0\): the twist and both Weyl
matrices commute with \(G_i\), and \(d_*w_p^{-1}=w_q^{-1}\).
This gives \(w_p^{-1}g,w_q^{-1}gx\in\mathcal N_j\).
Theorem~\ref{thm:partial-real-head-sizes} at \(g,p\) and at \(gx,q\)
now proves the assertion about large minors (vacuous at \(j=0\)).

Proposition~\ref{prop:real-area-stopping} gives the pointwise
hypothesis of Proposition~\ref{prop:real-crossed-square-bound} for
\(B_j\) with \(c=R^{29n\nu+5k}\).
For the assertion with \(T\), consider the first factor on the
set where \(B_j>0\). If the raw pairing at \(h g,hgx\) is nonzero
for \(h\in G_i\), both sizes are at most \(R\). The two required
near conditions still hold after this translation because
\(h\in Q_j\) commutes with \(w_p,w_q\). Apply
Theorem~\ref{thm:real-large-position-cap} at \(h g\).
The second stopping estimate improves \(c\) by the claimed factor.
When \(j=\nu\) its hypothesis forces this product to be zero.

Apply Proposition~\ref{prop:real-crossed-square-bound} in \(g\)
and Cauchy--Schwarz in \(x\). Theorem~\ref{thm:real-ball-energies}
shows the two sums of squared norms are finite almost everywhere
on the ball and gives
\[
 \int_X\sqrt{Z(x)Z'(x)}\dd x
     \le n\pi^{2(k-1)\nu}(L')^{70n^2}\lambda D^2\,\delta(d_*)^{1/2}.
\]
The empty-head sum starting at \(j\) is bounded by its sum starting
at zero. The twisted sum starts at precisely \(j\); use \(e=1\)
for both functions. On multiplying by \(C_j\), its remaining
pivot factor is
\[
 \kappa_j\prod_{a=1}^j |p_a|^{-2d_a-1}
                \prod_{a=1}^j |p_a/q_a|^{d_a+1}
                    =|D_{q,j}|^{-1}\le L
\]
by Lemma~\ref{lem:partial-real-quotient}. Finally
\(\lambda^{-1}\pi^{2(k-1)\nu}\le C_{\rm B}\) by
Lemma~\ref{lem:shells-real-cut-mixed}. This proves both bounds.
\end{proof}

\begin{theorem}\label{thm:real-head-absolute-integrability}
The forms \(\mathcal I_j^\xi(p,q)\) of Definition~\ref{def:real-pivot-forms}
are absolutely convergent at every fixed pair of heads and every
\(\xi\in\C\), at sufficiently large arithmetic scales. More precisely,
with \(L\) as in Lemma~\ref{lem:real-head-ball-majorant} and
\(U=(2L^{4j+1})^{16n}\), one has
\[
 C_j\iint B_j(g,x)\dd\mathfrak m_j(g)\dd x
 \le n C_{\rm B}(1+2^{70n^2+1})L^{125n^2j+1}
                       U^{70n^2}R^{29n\nu+5k}.
\]
\end{theorem}
\begin{proof}
Use Lemma~\ref{lem:real-head-ball-majorant} on the inner ball
of radius \(U\). On the shell of inner radius \(2^aU\) and outer
radius \(2^{a+1}U\), \(a\ge0\), take \(T=(2^aU)^{1/(16n)}\).
It exceeds \(L^{4j+1}\). Relative to the inner ball bound, the
shell contributes at most
\(2^{70n^2}2^{-a(s/(16n)-70n^2)}\), since
\(\sqrt\pi/(2L^{4j+1})\le1\). Sum this series as in
Lemma~\ref{lem:shells-real-cut-mixed}. The modulus of each
oscillatory integrand, before either floor integration, equals
the integrand of \(B_j\). Lemma~\ref{lem:product-integration}
therefore gives all asserted absolute integrals.
\end{proof}


\section{Descent from the real Gram kernel to its total mass}

\begin{lemma}\label{lem:real-lower-pivot-descent}
At sufficiently large scales take \(0\le j<\nu\) and arbitrary
heads \(p,q\) of length \(j\). For \(i=j+1\), the function
\(c\mapsto J_i^*((p,c),(q,c))\) is absolutely integrable on
\(\C^\times\). For every \(\xi\in\C\),
\[
 \int_{\C^\times}\psi(\xi c)J_i^*((p,c),(q,c))\dd c
                         =\mathcal I_j^\xi(p,q).
\]
\end{lemma}
\begin{proof}
Use \(H,H'\) of Definition~\ref{def:real-pivot-forms}.
For fixed \(\sigma,g,x\), the integrand in \(t\) on the right
depends on \(t\) only through \(t^{-1}f_i\). Indeed
\(K'_i\) is invariant under simultaneous multiplication by the
stabilizer of \(f_i\) in \(G_i\), since it is contained in
\(Q_i\cap G_i\); apply Lemma~\ref{lem:real-pairing-floor-laws}
with initial length \(i-1\). Right multiplying \(t\) by \(\sigma\)
preserves integrals of such functions by
Proposition~\ref{prop:actual-floor-quadrics}. Replacing \(t\) with
\(t\sigma\) changes the phase to \(\psi(\xi[f_i,t^{-1}f_i])\)
and changes the second pairing to \(H'(t\sigma)\).

Disintegrate \(t=w_i(c)^{-1}u\), with
\(|c|^{2d_i}\dd c\) on \(\C^\times\) and \(\dd\eta_i(u)\) on its
layer, by Definition~\ref{def:successive-real-floors}.
As \(u f_i=f_i\), one has
\([f_i,t^{-1}f_i]=[f_i,w_i(c)f_i]=c\).
The untwisted pairing \(H(\sigma)\) equals
\(K_i(u\sigma g,u\sigma gx)\) by
Lemma~\ref{lem:real-pairing-floor-laws} with empty head, since
\(u\in Q_i\). The matrices \(u,\sigma,w_i(c)\) commute with
\(w_p,w_q\), since they act only on \(V_i\).
Set \(z'=u\sigma g\). The resulting product of the two pairings is
\[
 K_i(z',z'x)\,
  \overline{K'_i((w_p w_i(c))^{-1}z',(w_p w_i(c))^{-1}z'x)}.
\]
The new Weyl matrices are \(w_p w_i(c),w_q w_i(c)\); their ratio
is still \(d_*\). Thus \(K'_i\) is exactly the pairing needed for
the extended heads. By Definition~\ref{def:real-preceding-measure}
the law of \(z'\) is \(\eta_i*\theta_i*\mathfrak m_j=\mathfrak m_i\).
Also \(C_j |c|^{2d_i}\) is the normalization constant \(C_i\) for
these two heads. At zero frequency \(B_i^0\) is the product of
\(K_i\) and the conjugate translated \(K'_i\), by
Lemma~\ref{lem:real-pairing-floor-laws}. This identifies the
integral in \((z',x)\) with \(J_i^*((p,c),(q,c))\) after normalization.

Each change applies first to the modulus, whose original total
integral is finite by Theorem~\ref{thm:real-head-absolute-integrability}.
Consequently the transformed product is absolutely integrable
jointly in \(c,z',x\). The kernels are Borel in their head variables:
their defining integrands are compositions of Borel functions
with continuous matrix maps, and the fixed parameter measures
give measurability under successive integrations by
Lemma~\ref{lem:product-integration}. The same lemma gives absolute
integrability in \(c\) after integrating the pairings and permits
the preceding changes of order for the complex integral.
\end{proof}

\begin{lemma}\label{lem:real-final-pivot-descent}
For heads \(p,q\) of length \(\nu\), at the same scales and
every \(\xi\in\C\), one has
\[
  \int_{\C^\times}\psi(\xi c)J_n((p,c),(q,c))\dd c
                                      =\mathcal I_\nu^\xi(p,q),
\]
where \(J_n\) is the Gram kernel of
Definition~\ref{def:real-top-correlation}. Its integrand is
absolutely integrable.
\end{lemma}
\begin{proof}
Use \(A_p,X,M\) of Proposition~\ref{prop:last-real-oscillatory-pair}.
At \(i=n\), Lemma~\ref{lem:real-pairing-floor-laws} gives exactly
the raw pairings \(U_{g,x}, W_{g,x}\) of that proposition.
Proposition~\ref{prop:actual-floor-quadrics} converts the two
\(\theta_n\)-integrals in \(B_\nu^\xi\) to its two quadric integrals.
For fixed \(x\), \(B_\nu^\xi\) is invariant under \(g\mapsto h g\)
for \(h\in M\). Indeed replace the two orbit vectors by their
simultaneous images under \(h^{-1}\), using the invariant orbit
law and signed pairing and commutation of \(h\) with both Weyl
matrices. The same holds for the nonnegative majorant \(B_\nu\).

If \(B_\nu(g,x)>0\), choose \(\sigma,t\in M\) at which both
pairings are nonzero. Put \(y=\sigma g,\ y'=yx,\ h=\sigma t^{-1}\).
Then \(V(y),V(h^{-1}w_p^{-1}y),V(y'),V(h^{-1}w_q^{-1}y')\)
are all nonzero. The compact-rotation reduction in
Proposition~\ref{prop:last-real-oscillatory-pair} gives
\(\|y\|_{\rm op}\le A_p,\ \|y'\|_{\rm op}\le A_q\).
In particular \(g\in M B_G(A_p)\) and \(x\in B_G(X)\).
Insert the cutoff \(a_{A_p}\) of Lemma~\ref{lem:middle-haar-cutoff};
its integral along \(M\) is one on this support. Since
\(\mathfrak h_\nu=\kappa*\theta_n\) is exactly the measure used
in that lemma, Lemma~\ref{lem:real-preceding-translation} gives
\[
 \int B_\nu^\xi(g,x)\dd\mathfrak m_\nu(g)
                       =\lambda\int a_{A_p}(y)B_\nu^\xi(y,x)\dd y
\]
for almost every \(x\). This follows first for the majorant and
then for the complex integrand by absolute integrability in
Theorem~\ref{thm:real-head-absolute-integrability}. The cutoff
on the right vanishes outside \(B_G(2A_p)\). Finally
\(C_\nu\lambda=u_kP_\nu(p)P_\nu(q)/(u_rD^2)\).
The asserted formula now follows from
Proposition~\ref{prop:last-real-oscillatory-pair}, which also gives
absolute integrability at the last pivot.
\end{proof}

\begin{theorem}\label{thm:real-gram-total-mass}
Define \(J_j=J_j^*\) for \(0\le j\le\nu\) and use \(J_n\) of
Definition~\ref{def:real-top-correlation}. At every fixed pair
of heads of length \(j<n\), integration of \(J_{j+1}\) in one
common appended complex coordinate gives \(J_j\) absolutely.
The Fourier integral there at any \(\xi\) equals
\(\mathcal I_j^\xi\). All \(J_j\) are positive semidefinite.
For \(0<\eta<1\) fixed, at sufficiently large scales one has
\[
 \int_{(\C^\times)^n}\|\Xi(p)\|_2^2\dd p
   = J_0
   =D^{-2}\int_{G_1}|c_V(x)|^2\dd x
   \le (1+b^{-(5-\eta)r^2})/\mathcal D_\tau .
\]
\end{theorem}
\begin{proof}
Lemmas~\ref{lem:real-lower-pivot-descent}
and~\ref{lem:real-final-pivot-descent} give the Fourier identities.
Set \(\xi=0\) for ordinary integration. In particular the definition
at \(\nu\) agrees with Corollary~\ref{cor:actual-penultimate-gram}.
The top kernel is positive semidefinite by
Theorem~\ref{thm:real-top-gram-kernel}. For a finite list of heads
and complex coefficients, multiply the descent identity by the
corresponding coefficient products and sum. The quadratic form
is the integral of a nonnegative quadratic form at the next stage.
This proves positive semidefiniteness by downward induction; all
finite sums can be integrated termwise by the absolute convergence
just proved. On the diagonal the integrands are nonnegative, so
Lemma~\ref{lem:product-integration} identifies the iterated integral
with the product integral of \(\|\Xi\|_2^2\).
Finally \(\mathfrak m_0\) is unit mass at \(1\), \(C_0=(\lambda D)^{-2}\),
and \(B_0^0(1,x)=|K_0(1,x)|^2=\lambda^2|c_V(x)|^2\) by
Corollary~\ref{cor:large-real-translation-energy} and the pairing
recurrence. Apply Theorem~\ref{thm:unrestricted-real-cut-energy}.
\end{proof}


\section{Fourier integration at an arbitrary real pivot}

\begin{lemma}\label{lem:two-real-amplitude-supports}
Use \(j,i,p,q,L,H,H'\) of Definition~\ref{def:real-pivot-forms}
and Lemma~\ref{lem:real-head-ball-majorant}. For fixed \(g,x\)
the two amplitudes \(H(\sigma),H'(t)\) are bounded Borel functions
of \(v=\sigma^{-1}f_i,\ v'=t^{-1}f_i\) on their orbit, up to orbit
null sets. If \(B_j(g,x)>0\), there is a single \(h_0\in G_i\) such
that their respective supports, up to null sets, satisfy
\[
             \|h_0 v\|\le L^{6j+4},\qquad
             \|h_0 v'\|\le L^{6j+4}.
\]
\end{lemma}
\begin{proof}
For both amplitudes the stabilizer of \(f_i\) in \(G_i\) belongs
to the subgroup of simultaneous invariance of \(K_i\) or \(K'_i\)
in Lemma~\ref{lem:real-pairing-floor-laws}. The Borel sections
and orbit laws in Proposition~\ref{prop:actual-floor-quadrics}
give the factorization and that lemma on the pairings gives
boundedness. If \(B_j>0\), then \(g,w_p^{-1}g\in\mathcal N_j\)
by Lemma~\ref{lem:real-head-ball-majorant} and its proof.
For \(j>0\), Theorem~\ref{thm:partial-real-head-sizes} gives
\(h_0\in G_i\) for which both \(\|h_0g\|_{\rm op}\) and
\(\|h_0w_p^{-1}g\|_{\rm op}\) are at most \(L^{6j+3}\).
For \(j=0\) take \(h_0=g^{-1}\).

Where \(H(\sigma)\ne0\), choose \(a\in Q_i\) with
\(\|a\sigma g\|_{\rm op}\le R\) using
Lemma~\ref{lem:real-pairing-floor-laws}. The projection to \(V_i\)
of \(h_0\sigma^{-1}a^{-1}f_i\) is \(h_0\sigma^{-1}f_i\):
\(a^{-1}f_i-f_i\) is in the span of the preceding frame entries,
fixed pointwise by \(G_i\). The unprojected norm is at most
\(R\|h_0g\|_{\rm op}\), by writing the matrix as
\((h_0g)(a\sigma g)^{-1}\).
Use the same proof with \(w_p^{-1}g\) at a nonzero value of \(H'\).
Since \(R\le L\), this gives both bounds with the same \(h_0\).
\end{proof}

\begin{proposition}\label{prop:real-frequency-on-ball}
For \(\xi\ne0\), and \(L,L',X\) as in
Lemma~\ref{lem:real-head-ball-majorant}, one has
\[
 C_j\int_X\int |B_j^\xi(g,x)|\dd\mathfrak m_j(g)\dd x
 \le n C_{\rm B}L^{125n^2j+50n^2+1}(L')^{70n^2}|\xi|^{-(k-1)} .
\]
\end{proposition}
\begin{proof}
For \(B_j(g,x)>0\), apply the common isometry of
Lemma~\ref{lem:two-real-amplitude-supports} to the two orbit
vectors. The measure on each is \(\delta_\gamma^{(i)}\) by
Proposition~\ref{prop:actual-floor-quadrics}; both the measures
and the phase pairing are preserved by this change.
The two amplitudes are supported on the flat box of radius
\(T=L^{6j+4}\). These boxes have finite measure by
Theorem~\ref{thm:real-flat-fourier}, so the bounded amplitudes
are square integrable. Cauchy--Schwarz against the first amplitude
and that Fourier theorem applied to the conjugate second amplitude
give
\[
 |B_j^\xi(g,x)|^2\le
    M_{\rm quad}(T)^2(\pi/|\xi|)^{2(k-1)}
        \int|H|^2\dd\theta_i\int|H'|^2\dd\theta_i .
\]
Here \(M_{\rm quad}(T)=4\pi T^2(\pi T^2)^{2(n-i)}\).
Since \(L\ge4,\ j\le n-1\), one has
\(\log_L M_{\rm quad}(T)\le 2+12n+2n(1+12n)\le50n^2\).
The first squared amplitude integral is
\(E_{j,j}(g,x)/W_j(g,x)\), and the second is its twisted
counterpart at \(w_p^{-1}g\), as \(G_i\subset Q_j\) preserves
\(W_j\) by Lemma~\ref{lem:partial-real-radical-measure}.
Each numerator is one summand of the corresponding \(S_j\).
If \(B_j=0\), its oscillatory integral vanishes.

Apply Proposition~\ref{prop:real-crossed-square-bound} with
\(c=L^{50n^2}(\pi/|\xi|)^{k-1}\). Integrate over \(X\) by the
squared-norm and determinant calculation of
Lemma~\ref{lem:real-head-ball-majorant}. Before its final
replacement by \(C_{\rm B}\), the factor involving \(\pi,\lambda\)
in that calculation is \(\lambda^{-1}\pi^{2(k-1)\nu}\).
Including the extra factor from the Fourier integration gives
\(\lambda^{-1}\pi^{(2\nu+1)(k-1)}\le C_{\rm B}\).
The remaining factors give the asserted bound.
\end{proof}


\section{Uniform real Fourier bounds in the initial minors}

\begin{definition}\label{def:explicit-real-band-weights}
For \(t>0,\ h\ge0\) put \(d_h(t)=t^h\min(1,t)^6\).
For a head \(p\in(\C^\times)^j\) let
\[
       \mathcal H_j(p)=\prod_{a=1}^j d_{n^6}(|D_{p,a}|^{-1}),
       \qquad K=k-n^7 .
\]
Here \(K\ge0\) at the scales under consideration by
Proposition~\ref{prop:thick}.
\end{definition}

\begin{lemma}\label{lem:combined-real-band-estimate}
Fix \(\delta>0\). At all sufficiently large arithmetic scales take
heads \(p,q\) of length \(j\le\nu\), \(\xi\ne0\), and set
\[
 L=\max\bigl(\{R,4,r\}\cup
  \{|D_{p,a}|^{-1},|D_{q,a}|^{-1}:1\le a\le j\}\bigr).
\]
Let \(0\le f(g,x)\le B_j(g,x)\) be Borel. Suppose that
for every \(L'\ge\max(4,r,R)\) the integral of \(C_j f\)
on \(\|x-I\|_{\rm op}\le L'\), with measures
\(\dd\mathfrak m_j(g)\dd x\), is at most
\(n C_{\rm B} L^{n^5}(L')^{70n^2}|\xi|^{-(k-1)}\).
Then
\[
 C_j\iint f(g,x)\dd\mathfrak m_j(g)\dd x
   \le C_{\rm B} b^{(12+\delta)n^9}
           d_K(|\xi|^{-1})\,\mathcal H_j(p)\mathcal H_j(q).
\]
\end{lemma}
\begin{proof}
By Lemma~\ref{lem:real-head-ball-majorant} we need only treat
\(|D_{p,a}|,|D_{q,a}|\le R^{2a}\).
Write \(P=70n^2,\ A=16nP,\ t=\max(1,R^{10}|\xi|)\), and
\(U=(2L^{4j+1}t)^{16n}\). On the ball of radius \(U\), take the
smaller of the hypothesized estimate and the unconditional
estimate of Lemma~\ref{lem:real-head-ball-majorant}.
On each shell of inner radius \(2^aU\) and outer radius \(2^{a+1}U\),
use that lemma with \(T=(2^aU)^{1/(16n)}\).
Since \(\sqrt\pi/(2L^{4j+1})\le1\), the additional factor satisfies
\((\pi/T^2)^{s/2}\le t^{-s}2^{-a s/(16n)}\).
As in Lemma~\ref{lem:shells-real-cut-mixed}, the sum over \(a\)
costs at most \(2^{P+1}t^{-s}\) relative to the unconditional
inner ball estimate. Let \(E=n^5+(4j+1)A\). Substituting \(U^P\)
and using \(125n^2j+1\le n^5\) eventually gives the total bound
\[
 n C_{\rm B}2^{P+A+1} L^E R^{29n\nu}
 t^A\left(\min(R^{5k},|\xi|^{-(k-1)})+R^{5k}t^{-s}\right).
\]

Set \(y=|\xi|^{-1}\). Eventually \(5k\le10K,\ 5k\le10s\)
and \(s-A\ge K+6\), as \(k\ge X^5-2\sqrt X,\ n\le\sqrt X\),
\(s=k-3n-2,\ A=1120n^3\), for the scale \(X\) in
Proposition~\ref{prop:thick}. Also \(k-1\ge s\).
If \(t=1\), then \(y\ge R^{10}\) and the parenthesized expression
is at most \(2R^{5k}\le2d_K(y)\). If \(t>1\), then
\(R^{5k}t^{-s}\le y^s\) and \(t^A=R^{10A}y^{-A}\).
When \(y\le1\), the resulting sum is at most
\(2R^{10A} y^{K+6}\). When \(1<y<R^{10}\), use
\(y^{-A}\le1\) and \(y^{k-1}\le y^K R^{10n^7}\).
In every case the factor beginning with \(t^A\) costs at most
\(2 R^{10A+10n^7}d_K(y)\).

It remains to translate the single scale \(L\) into the two head
weights. For each occurrence of \(z=|D_{p,a}|^{-1}\) or
\(|D_{q,a}|^{-1}\) below one, its factor in the product of head
weights is \(z^{n^6+6}\ge R^{-2a(n^6+6)}\). The factors for
\(z\ge1\) are at least one, and any such factor attaining \(L\)
is \(L^{n^6}\). Since \(R\ge4,r\) eventually and
\(E\le n^6\) eventually, whether \(L=R\) or some \(z=L\) this gives
\[
 L^E\le R^{\,n^6+2j(j+1)(n^6+6)}
                               \mathcal H_j(p)\mathcal H_j(q).
\]
The remaining factor apart from \(C_{\rm B}\) and the three
prescribed weights has logarithm to base \(b\) at most
\[
 \log_b(n2^{P+A+2})+
 (6n+2)\bigl(29n\nu+10A+10n^7+n^6+2j(j+1)(n^6+6)\bigr).
\]
Divide by \(n^9\). Since \(j\le n-1,\ n\to\infty,\ \log b\to\infty\),
the upper limit is at most \(12\), uniformly in both heads and
the frequency. This proves the assertion for each fixed \(\delta\).
\end{proof}

\begin{theorem}\label{thm:uniform-common-real-frequency}
For each fixed \(\delta>0\) at all sufficiently large arithmetic
scales, for every \(0\le j\le\nu\), every \(p,q\in(\C^\times)^j\),
and every \(\xi\in\C^\times\), one has
\[
 \left|\int_{\C^\times}\psi(\xi c)J_{j+1}((p,c),(q,c))\dd c\right|
 \le C_{\rm B} b^{(12+\delta)n^9}
      d_{k-n^7}(|\xi|^{-1})\,\mathcal H_j(p)\mathcal H_j(q).
\]
\end{theorem}
\begin{proof}
The integral is \(\mathcal I_j^\xi(p,q)\) by
Theorem~\ref{thm:real-gram-total-mass}. In
Lemma~\ref{lem:combined-real-band-estimate} take \(f=|B_j^\xi|\),
which is bounded by \(B_j\). Proposition~\ref{prop:real-frequency-on-ball}
implies its ball hypothesis since \(125n^2j+50n^2+1\le n^5\)
eventually. The triangle inequality for the defining integral
gives the result. In particular \(\delta=1\) supplies the
frequency bound with the larger exponent \(2000n^9\).
\end{proof}


\section{The real unipotent integral for unequal last entries}

\begin{lemma}\label{lem:unequal-real-pivot-phase}
Take \(1\le j\le\nu\) and heads \(p,q\) of length \(j\) with
\(p_j\ne q_j\). Set \(\xi=p_j^{-1}-q_j^{-1}\).
Use the twisted \(K'_a\) of Definition~\ref{def:real-pivot-forms}.
For \(g,x\in G_1\), write \(F_g(\sigma^{-1}f_{j+1})
=K'_{j+1}(\sigma g,\sigma gx)\) for \(\sigma\in G_{j+1}\)
on the orbit of \(f_{j+1}\). With
\(\delta'=\delta^{(j+1)}_\gamma\), where \(\gamma=0\) for
\(j+1<n\) and \(\gamma=-1\) for \(j+1=n\), one has
\[
  K'_j(n_j(t,\beta)g,n_j(t,\beta)gx)
      =\int\psi\bigl(p_j\xi[t,v]\bigr)F_g(v)\dd\delta'(v).
\]
\end{lemma}
\begin{proof}
The stabilizer of \(f_{j+1}\) in \(G_{j+1}\) fixes the pairing,
and its Borel orbit sections give \(F_g\) as in
Lemma~\ref{lem:two-real-amplitude-supports}. Unfold the recurrence
for \(K'_j\) in Lemma~\ref{lem:real-pairing-floor-laws}.
Conjugation of \(n_j(t,\beta)\) by \(\sigma\in G_{j+1}\) gives
\(n_j(\sigma t,\beta)\in N^j\), as in
Proposition~\ref{prop:real-radical-fourier-phase}.
Its only possibly nonzero simple-root entry is \(-[\sigma t,f_{j+1}]\).
In the second argument conjugation by \(d_*\) multiplies that
entry by \(r_j=p_j/q_j\), since \(r_{j+1}=1\).
The two separate transformation laws therefore give the phase
of \((1-r_j)[t,\sigma^{-1}f_{j+1}]=p_j\xi[t,\sigma^{-1}f_{j+1}]\).
Proposition~\ref{prop:actual-floor-quadrics} yields the integral.
The recurrence is absolutely convergent by
Lemma~\ref{lem:real-pairing-floor-laws}.
\end{proof}

\begin{proposition}\label{prop:unequal-real-layer-bound}
Let \(L\) be as in Lemma~\ref{lem:real-head-ball-majorant}, with
the heads of Lemma~\ref{lem:unequal-real-pivot-phase}. For every
\(g,x\) at sufficiently large scales one has
\[
 \begin{aligned}
 &\int {\bf1}_{\mathcal N_j}(w_p u g)|K'_j(ug,ugx)|^2\dd\eta_j(u)\\
 &\qquad\le L^{100n^2}|p_j|^{-2d_j-1}(\pi/|\xi|)^{2(k-1)}
                  \int |K'_{j+1}(\sigma g,\sigma gx)|^2
                                                \dd\theta_{j+1}(\sigma).
 \end{aligned}
\]
\end{proposition}
\begin{proof}
If the integral on the left is positive there is a layer element
\(u_0\) with \(g'=u_0g\in\mathcal N_j\) and
\(w_p g'\in\mathcal N_j\), by Lemma~\ref{lem:real-pairing-floor-laws}.
Replace \(g\) by \(g'\) there, using right invariance of the layer
measure. The squared amplitude integral on the right is unchanged:
\(\sigma u_0\sigma^{-1}\in N^j\) contributes only a phase to
\(K'_{j+1}\) by the same lemma. Thus take \(g,w_p g\in\mathcal N_j\).
Choose \(a\in Q_j\) with \(\|ag\|_{\rm op}\le R\), and write
\(h_0\in G_{j+1}\) for its action on the middle quotient.
For the amplitude in Lemma~\ref{lem:unequal-real-pivot-phase},
\(\|h_0v\|\le R^2\) on its nonzero set by the first assertion of
Lemma~\ref{lem:real-centered-radical}: its proof only uses
\(\sigma g\in\mathcal N_{j+1}\).

For every \(u=n_j(t,\beta)\) with \(w_p u g\in\mathcal N_j\),
apply Theorem~\ref{thm:partial-real-head-sizes} to \(ug\),
the Weyl matrix \(w_p^{-1}=w_{\bar p}\), and the near multiplier
\(a u^{-1}\). It gives
\(\|a u^{-1}e_j^*\|\le L^{4j+1}\). Use it also at the identity
layer element to obtain \(\|a e_j^*\|\le L^{4j+1}\).
By Definition~\ref{def:real-radical-weights},
\(u^{-1}e_j^*-e_j^*=-t+(-i\beta-[t,t]/2)e_j\).
Projecting its image under \(a\) to \(V_{j+1}\) gives \(-h_0t\);
in particular \(\|h_0t\|\le 2L^{4j+1}\).
For fixed \(t\), the difference of the two images for admissible
\(\beta,\beta'\) is \(-i(\beta-\beta')a e_j\). Its norm is at
most \(2L^{4j+1}\), and \(\|a e_j\|\ge1\) by the unit diagonal of
the flag action. The admissible central coordinates have length
at most \(2L^{4j+1}\), in the sense of their Lebesgue measure.

Apply the phase formula and integrate first in the center.
Change \(t,v\) by the same \(h_0\); this preserves ambient and
orbit measures by Proposition~\ref{prop:actual-floor-quadrics}
and \(|\det h_0|=1\). Use ambient output of radius
\(T=2L^{4j+1}\), quadric input of radius \(R^2\), and frequency
\(p_j\xi\) in Theorem~\ref{thm:real-flat-fourier}. The squared
amplitude integral on the orbit is exactly the one in the
statement. The other factors are bounded by
\[
  T\,M_{\rm amb}(T)M_{\rm quad}(R^2)
                 (\pi/|p_j\xi|)^{2(k-1)} .
\]
Here the two box factors have \(h=n-j-1\). Since \(T\le L^{4j+2}\)
and \(L\ge4\), their product with \(T\) has logarithm to base
\(L\) at most
\((4j+2)+(2h+2)(1+8j+4)+2+4+10h\le70n^2\).
Also \(|D_{p,j}|\le R^{2j}\) by the head-size theorem at \(g\);
hence \(|p_j|\le L^{2j+1}\). As
\(2d_j+1-2(k-1)=4(n-j)+1\), replacing \(|p_j|^{-2(k-1)}\)
by \(|p_j|^{-2d_j-1}\) incurs a factor at most
\(L^{(2j+1)(4(n-j)+1)}\le L^{30n^2}\).
These estimates give the result, including the zero case.
\end{proof}


\section{The real family of kernels in the orbital sieve}

\begin{proposition}\label{prop:unequal-real-frequency-ball}
For the heads and nonzero \(\xi=p_j^{-1}-q_j^{-1}\) of
Lemma~\ref{lem:unequal-real-pivot-phase}, and \(L,L',X\) of
Lemma~\ref{lem:real-head-ball-majorant}, one has
\[
 C_j\int_X\int |B_j^0(g,x)|\dd\mathfrak m_j(g)\dd x
 \le n C_{\rm B} L^{125n^2j+50n^2+1}(L')^{70n^2}|\xi|^{-(k-1)}.
\]
\end{proposition}
\begin{proof}
Fix \(x\), writing \(W_a(g)=W_a(g,x)\),
\(I_a(g)={\bf1}_{\mathcal N_a}(g)/W_a(g)\), and put \(w=w_p\).
Let \(w_0\) flip the first \(j-1\) pairs with their entries from \(p\).
Write \(Z=Z_{j-1}(V,1;x)\) and \(Z'=Z_{\ge j}(x)\) for the
twisted kernel. Both are finite for almost every \(x\) in the ball
by Theorem~\ref{thm:real-ball-energies}.
The integrand on the left is
\(|K_j(g,gx)K'_j(w^{-1}g,w^{-1}gx)|\), by
Lemma~\ref{lem:real-pairing-floor-laws}.
Its nonzero set requires both \(g,w^{-1}g\in\mathcal N_j\).
In Cauchy--Schwarz in \(g\), use the squared factors
\[
 \begin{split}
 A(g)I_j(w^{-1}g),\qquad
 &A(g)=W_{j-1}(g)|K_j(g,gx)|^2,\\
 \Phi(w^{-1}g),\qquad
 &\Phi(y)={\bf1}_{\mathcal N_j}(wy)W_j(y)I_{j-1}(w_0y)
                                               |K'_j(y,yx)|^2 .
 \end{split}
\]
Their product is the squared integrand wherever the latter is
nonzero. Indeed \(w_0w^{-1}=w_j(p_j)^{-1}\in G_j\subset Q_{j-1}\)
preserves the earlier weights and near condition by
Lemma~\ref{lem:partial-real-radical-measure} and
Definition~\ref{def:real-near-floor}. All weights are positive.

Put \(H_p=\prod_{a\le j}|p_a|^{-2d_a-1}\).
For the first squared factor, \(A\) is left \(Q_j\)-invariant
and supported in \(\mathcal N_j\) by
Lemma~\ref{lem:real-pairing-floor-laws} for the empty head.
Its integral against \(\tau_j\) equals \(Z\) by definition.
Write \(\mathfrak m_j=\dd u_{N^j}*\tau_j\) by
Lemma~\ref{lem:real-preceding-translation} and apply
Corollary~\ref{cor:real-radical-test} with head \(\bar p\).
This bounds the first integral by \(L^{125n^2j}H_pZ\).

For the second squared factor, \(\Phi\) is left \(G_{j+1}\)-invariant:
the two Weyl matrices commute with that group, the weights and
near sets are invariant, and \(K'_j\) is simultaneously invariant
there by Lemma~\ref{lem:real-pairing-floor-laws}. Hence
Lemma~\ref{lem:real-preceding-translation} replaces \(\Phi(w^{-1}g)\)
by \(\Phi(g)\) in the integral. In
\(\mathfrak m_j=\eta_j*\dd u_{N^{j-1}}*\tau_j\), permitted by
Lemma~\ref{lem:partial-real-radical-measure}, integrate first
over \(g=u y\), \(u\) in layer \(j\).
The factors \(W_j(uy)\) and \(I_{j-1}(w_0uy)\) are constant in
\(u\), since \(u\in Q_j\subset Q_{j-1}\) commutes with \(w_0\).
Proposition~\ref{prop:unequal-real-layer-bound} bounds the remaining
inner integral by
\[
 L^{100n^2}|p_j|^{-2d_j-1}(\pi/|\xi|)^{2(k-1)}
                         \int |K'_{j+1}(\sigma y,\sigma yx)|^2
                                                  \dd\theta_{j+1}(\sigma).
\]
After multiplying by \(W_j(y)\), the squared integral is
\(E_{j,j}(y,x)\) for the twisted kernel. Enlarge this to \(S_j(y,x)\).
It is left \(Q_j\)-invariant and supported in \(\mathcal N_j\) by
Lemma~\ref{lem:real-tail-square-invariance}. It is thus left
\(N^{j-1}\)-invariant and supported in \(\mathcal N_{j-1}\).
Apply Corollary~\ref{cor:real-radical-test} to this function,
the head of length \(j-1\) giving \(w_0\), and \(\rho'=\tau_j\).
The remaining integral against \(\tau_j\) equals \(Z'\). The
second squared factor consequently has integral at most
\[
 L^{100n^2+125n^2(j-1)}H_p(\pi/|\xi|)^{2(k-1)} Z'.
\]
Cauchy--Schwarz gives the product of the two square roots.
Integration of \(\sqrt{ZZ'}\) over \(X\) is bounded by
Theorem~\ref{thm:real-ball-energies} as in
Lemma~\ref{lem:real-head-ball-majorant}; \(Z\) is one summand of
the empty-head sum. The pivot powers cancel to \(|D_{q,j}|^{-1}\le L\)
as in that lemma. Finally the total power of \(\pi\) is absorbed
by \(C_{\rm B}\) exactly as in Proposition~\ref{prop:real-frequency-on-ball}.
\end{proof}

\begin{theorem}\label{thm:actual-real-sieve-family}
At every sufficiently large scale and every actual real place,
use the functions \(J_0,\ldots,J_n\) of Theorem~\ref{thm:real-gram-total-mass}.
For each fixed \(\delta>0\), these satisfy all the kernel conditions
of Definition~\ref{def:orbital-tower} over the single factor
\(C=\C\), with length parameter \(\nu=n-1\), phase
\(\psi(z)=\exp(z-\bar z)\), \(I(v)=v^{-1}\), and bounds
\[
 \begin{split}
 E_0&=(1+b^{-10})/\mathcal D_\tau,\qquad
 \Lambda_j(p)=\mathcal H_j(p),\\
 \beta(\xi)&=C_{\rm B}b^{(12+\delta)n^9}
                         d_{k-n^7}(|\xi|^{-1})\qquad(\xi\ne0).
 \end{split}
\]
The scale threshold may depend on \(\delta\) and the fixed arithmetic
parameters. At the top the kernel is the Gram kernel of the actual
vector \(\Xi\) in Definition~\ref{def:real-top-correlation}, including
its identity with the matrix integral for any middle isometry with
the specified last entry. One may in particular use \(2000n^9\)
as the exponent of \(b\) in \(\beta\).
\end{theorem}
\begin{proof}
Theorem~\ref{thm:real-gram-total-mass} gives the bottom bound,
absolute common-coordinate integration and its value at zero
frequency. Assigning one to an appended zero coordinate changes
none of these Lebesgue integrals. Theorem~\ref{thm:uniform-common-real-frequency}
gives the bound at nonzero common-coordinate frequencies.
Theorem~\ref{thm:real-top-gram-kernel} proves top continuity,
positive semidefiniteness and the claimed identity at every
admissible middle isometry.

It remains to bound \(J_j(p,q)\) when \(1\le j\le\nu\) and
\(\xi=p_j^{-1}-q_j^{-1}\ne0\).
Use \(f(g,x)=|B_j^0(g,x)|\le B_j(g,x)\). The ball hypothesis
of Lemma~\ref{lem:combined-real-band-estimate} follows from
Proposition~\ref{prop:unequal-real-frequency-ball}, since
\(125n^2j+50n^2+1\le n^5\) eventually. Apply that lemma, then
the triangle inequality in the definition of \(J_j^*\).
This gives precisely \(\beta(\xi)\Lambda_j(p)\Lambda_j(q)\),
as required. Taking \(\delta=1\) proves the last statement.
\end{proof}


\section{Finite-place phase functions and congruence cosets}

\begin{definition}\label{def:finite-bessel-level}
Let \(T\) be a finite set of finite places of the fixed field \(F\);
over \(E\) take all factors above \(T\), a set stable under
conjugation. Write \(A_T=\prod_{v\in T}A_v\) and
\(B_T=\prod_{v\in T}B_v\) for the integer rings of
Lemma~\ref{lem:real-place-rings}. Their fraction algebras are
\(F_T,E_T\), allowing the two factors at a split place. Let
\(G_T\) be the \(E_T\)-linear isometries of the row form in
Definition~\ref{def:row-space}. Denote the specified integral row
lattice over \(B_T\) by \(\mathcal L\), its stabilizer by \(K_T\),
and the kernel of the action on \(\mathcal L/g\mathcal L\) by
\(K_g\), where \(g=g_0g_1\) as in Proposition~\ref{prop:level-prime}.
Let \(U_T\) be the flag unipotent group of Definition~\ref{def:unip}
over \(F_T\), and set
\[
        \chi_g(u)=\psi_T^0(\alpha(u)/g).
\]
Here \(\psi_T^0\) is the restriction of \(\psi_f^0\) of
Proposition~\ref{prop:additive}, and \(\alpha\) is from
Proposition~\ref{prop:full-phase}. Write \(|\pi_v|_v=Q_v^{-1}\)
for a uniformizer of \(F_v\).
\end{definition}

\begin{lemma}\label{lem:finite-bessel-radical-volume}
The subgroup \(K_T\) is compact open, \(K_g\) is open and normal
in it, and \(K_T/K_g\) is finite. The character \(\chi_g\) is
continuous and trivial on \(U_T\cap K_g\).
There is a bi-invariant measure \(\dd u\) on \(U_T\) assigning
\(U_T\cap K_T\) volume one, for which
\[
 v_g:=|U_T\cap K_g|
   =\prod_{v\in T}|g|_v^{\,\nu(2k+2\nu+1)} .
\]
If \(T\) contains every place above \(g_0,g_1\), this equals
\(g^{-f\nu(2k+2\nu+1)}\).
\end{lemma}
\begin{proof}
In the integral basis, the Gram matrix and its inverse are integral
by Proposition~\ref{prop:row-group}. Thus an integral matrix
satisfying the isometry equation has integral inverse. The
stabilizer is the closed solution set of that equation among
integral matrices, hence compact, since the rings are inverse limits
of finite residue rings. It is open by the open condition that
the entries be integral. The congruence subgroup is the kernel
of a homomorphism to a finite group and is open by the entry
conditions. The simple-root sum is additive on the unipotent
group by the matrix computation of Proposition~\ref{prop:full-phase}.
For an element of \(U_T\cap K_g\) its value is in \(gB_T\),
since the frame vectors are integral and their outer coordinates
are integral coordinate functionals. Proposition~\ref{prop:additive}
gives \(\chi_g=1\) there and continuity everywhere.

At layer \(a=1,\ldots,\nu\), write \(l\) for the row on the
intervening space, of rank \(d_a=k-1+2(n-a)\). Use the coordinates
\((l,s)\) with corner entry \(-\omega lJ^{-1}l^*+s\), where
\(s+\bar s=0\) and \(\omega+\bar\omega=1\) by
Theorem~\ref{thm:integer-ring}. By Lemma~\ref{lem:heisenberg-layer},
integrality in this layer is exactly integrality of the row on
the intervening lattice and of \(s\); at principal level these
are multiplied by \(g\). Indeed the quadratic term is then in
\(g^2B_T\). The free row lattice has \(2d_a\) coordinates over
\(A_T\), and the skew lattice is \(i A_T\) by
Lemma~\ref{lem:skew-covol} after completion.
Give each of these additive lattices volume one. Layer multiplication
translates \(l\), then translates \(s\) at each fixed \(l\), on
both the left and the right. The inner group acts on \(l\) by an
isometry whose determinant has norm one, and leaves its scalar
quadratic term invariant. By the diagonal index computation in
Lemma~\ref{lem:class-covol}, this preserves the measure. The
semidirect products of the layers therefore give a bi-invariant
measure with the required normalization. The measure ratio at
layer \(a\) is \(\prod_v|g|_v^{2d_a+1}\); summing the exponents
gives \(\nu(2k+2\nu+1)\). Each \(g_i\) splits into \(f\) places
of \(F\) with \(Q_v=g_i\) and exponent one by
Lemma~\ref{lem:split-level}, proving the last formula.
\end{proof}

\begin{definition}\label{def:finite-bessel-vector}
For \(y\in G_T\) put \(\phi_g(y)=\chi_g(u)\) if \(y=u a\) with
\(u\in U_T,\ a\in K_g\), and put it equal to zero otherwise.
For \(K_g\subseteq Z\subseteq K_T\) a subgroup, let \(D_Z=K_g\backslash Z\)
be the finite set of cosets \(K_g a\). In the Hilbert space
\(\C^{D_Z}\) with the sum inner product define
\[
                    W_Z(y)_{K_g a}=\phi_g(ya^{-1}).
\]
\end{definition}

\begin{proposition}\label{prop:finite-bessel-unfold}
The functions of Definition~\ref{def:finite-bessel-vector} are
well defined, locally constant, and bounded. One has
\(\phi_g(u y h)=\chi_g(u)\phi_g(y)\) for \(u\in U_T,\ h\in K_g\).
Left multiplication by \(u\) acts on \(W_Z\) by \(\chi_g(u)\);
right multiplication by \(z\in Z\) permutes its coordinates via
\(K_g a\mapsto K_g a z^{-1}\). For every \(y,x\in G_T\),
\[
 v_g\langle W_Z(y),W_Z(yx)\rangle
    =\int_{U_T}{\bf1}_Z(uy)
                   \phi_g\bigl((uy)x(uy)^{-1}\bigr)\dd u .
\]
In particular,
\[
 \|W_Z(y)\|^2=[U_T\cap Z:U_T\cap K_g]\,{\bf1}_{U_T Z}(y).
\]
\end{proposition}
\begin{proof}
Two decompositions \(y=u a=u'a'\) have
\((u')^{-1}u\in U_T\cap K_g\), where
Lemma~\ref{lem:finite-bessel-radical-volume} gives trivial
character. The support is unchanged by the indicated left and
right multiplications and the formula on it follows from the
character law. In particular \(\phi_g\) is locally constant by
right \(K_g\)-invariance, and its modulus is at most one.
Changing \(a\) to \(h a\), \(h\in K_g\), multiplies its inverse
on the right by \(h^{-1}\), so the coordinate formula is
independent of the representative. Right translation by \(z\)
uses \(z a^{-1}=(a z^{-1})^{-1}\). Normality in \(Z\) implies
that \(K_g\) fixes every coordinate. This proves local constancy
and the stated equivariances.

Partition the integration set by \(uy\in K_g a\) for \(K_g a\in D_Z\).
This set in \(u\) is nonempty exactly when \(W_Z(y)_{K_g a}\ne0\);
in that case it is a right translate of \(U_T\cap K_g\) and has
measure \(v_g\). On it write \(uy=h a\) with \(h\in K_g\). Then
\(\phi_g(ya^{-1})=\chi_g(u)^{-1}\), and by the left and right
transformation laws
\[
  \phi_g\bigl((uy)x(uy)^{-1}\bigr)
    =\chi_g(u)\phi_g(yx a^{-1})
    =\overline{W_Z(y)_{K_g a}}\,W_Z(yx)_{K_g a}.
\]
Integration on the finitely many cosets gives the identity.
For \(x=1\) the integral is the measure of
\(\{u:uy\in Z\}\), either empty or a translate of \(U_T\cap Z\).
This compact subgroup is a finite union of cosets of \(U_T\cap K_g\)
since it maps into the finite group \(Z/K_g\). Dividing its
measure by \(v_g\) proves the squared-norm formula.
\end{proof}


\section{The finite coefficient and its congruence contribution}

\begin{lemma}\label{lem:finite-unitary-level-haar}
For any \(c>0\) the group \(G_T\) of Definition~\ref{def:finite-bessel-level}
has a left and right invariant Borel measure \(\mu\), finite on compact
sets, with \(\mu(K_T)=c\). Its normalized restriction to \(K_T\) is
the compatible uniform probability on the finite congruence images.
\end{lemma}
\begin{proof}
First take one place \(v\), and put \(A=A_v,\ B=B_v\),
with uniformizer \(\pi\) of \(A\). In the integral basis write
\(J\) for the Gram matrix and \(K(\pi^h)\) for isometries congruent
to the identity modulo \(\pi^h\). The compact group \(K=K_T\) in
this one-place notation is the inverse limit of its finite images
on the lattices modulo \(\pi^h\): a compatible tuple of matrices
lifts entrywise by completeness, and the closed isometry equation
passes to the limit. Uniform probabilities on these images are
consistent because all fibres of a surjective group homomorphism
have the same size. To construct the probability on \(K\), choose
representatives for each successive kernel and lifts for the previous
image. One may sample each successive digit from length measure on
\([0,1)\) by successively subdividing each interval into the same
number of equal half-open intervals as its digit choices. The
resulting compatible residues give a Borel map to \(K\) and its
cylinder probabilities are the indicated ones. They determine the
measure, since the congruence cylinders form a countable basis;
the generating-family argument of Lemma~\ref{lem:product-integration}
gives uniqueness. Both left and right multiplication preserve these
probabilities at each depth.

Extend this probability by left translation to each left coset of
\(K\) and sum. The cosets are countable because the finite-dimensional
ambient local field has a countable basis of residue balls, obtained
from the finite residue digits and powers of \(\pi\).
Changing the representative of a coset preserves the measure there.
Thus the result is left invariant and finite on compact sets,
which meet only finitely many open cosets. Any left invariant Borel
measure finite on compact sets is determined by its value on \(K\):
in each congruence quotient all cosets have equal measure.

For right invariance put \(X^\dagger=J^{-1}\bar X^{\mathsf T}J\)
and \(\mathfrak l=\{X\in M_r(B):X^\dagger=-X\}\).
On \(\pi^h\mathfrak l\), \(h\ge1\), the Cayley map
\(C(X)=(1+X)(1-X)^{-1}\) is a bijection to \(K(2\pi^h)\), where
the argument of \(K\) denotes the congruence ideal. Indeed its
inverse is \((y-1)(y+1)^{-1}\); when \(y=1+2\pi^h b\) the matrix
\((y+1)/2=1+\pi^h b\) has integral inverse by its geometric series.
Taking adjoints proves the two defining equations. The difference
\[
         C(X)-C(Y)=2(1-X)^{-1}(X-Y)(1-Y)^{-1}
\]
shows that residue cylinders on this lattice correspond exactly
to congruence cosets on its image, with an extra factor \(2\)
in the congruence ideal. Hence on this chart the group measure
is a constant multiple of additive measure on the skew space:
all cylinders of the same depth have equal mass on both sides,
and their partitions determine both measures.

For any isometry \(y\), conjugation \(X\mapsto yXy^{-1}\) has
determinant one on the skew space over \(F_v\). After extension
to \(E_v=F_v\otimes_F E\) it is conjugation on the full matrix
algebra: every matrix decomposes uniquely over \(F_v\) as a skew
matrix plus \(i\) times a skew matrix by taking its adjoint
symmetric and antisymmetric parts. The scalar extension map is
thus an isomorphism, also for a split pair. On all matrices left
multiplication has determinant \((\det y)^r\) by columns, and
right multiplication by \(y^{-1}\) has its reciprocal by rows.
Scalar extension is injective on \(F_v\), giving the assertion
there. Conjugate a sufficiently deep chart so that it and its
image lie in the chart at \(h=1\); this is possible by clearing
the finitely many denominators of \(y,y^{-1}\).
The Cayley map commutes with conjugation and the additive
determinant-one map preserves volume, by the diagonal lattice
calculation of Lemma~\ref{lem:class-covol} in coordinates over
\(\Q_p\). Conjugation therefore preserves the mass of a single
open congruence subgroup. Its image measure is left invariant;
the uniqueness above and finite index in \(K\) show that the
ratio is one everywhere. Right invariance follows by composing
conjugation and left translation. For arbitrary \(T\) take the
finite product and multiply by \(c\); when \(T\) is empty use
mass \(c\) at the identity.
\end{proof}

\begin{definition}\label{def:finite-cut-correlation}
With the two measures of Lemmas~\ref{lem:finite-bessel-radical-volume}
and~\ref{lem:finite-unitary-level-haar}, put \(v_Z=|U_T\cap Z|\).
For a left \(U_T\)-invariant function \(h\) supported on \(U_TZ\)
use the quotient integral
\(v_Z^{-1}\int_Z h(y)\dd\mu(y)\). Define the coefficient
\[
 c_Z(x)=v_Z^{-1}\int_Z\langle W_Z(y),W_Z(yx)\rangle\dd\mu(y),
 \qquad m_Z=\mu(Z)\mu(K_g)/v_g^2 .
\]
The integral is the quotient integral of the displayed inner
product by Proposition~\ref{prop:finite-bessel-unfold}.
\end{definition}

\begin{proposition}\label{prop:finite-congruence-energy}
Let \(C_Z(x)\) be the mean of \(\phi_g(yxy^{-1})\) for \(y\in Z\)
under the probability \(\dd\mu(y)/\mu(Z)\).
Then \(v_g c_Z=\mu(Z)C_Z\), with \(C_Z=1\) on \(K_g\).
The function \(C_Z\) is invariant on both sides by \(K_g\) and
under conjugation by \(Z\). One has the exact identities
\[
 \begin{split}
 v_Z^{-1}\int_Z\|W_Z(y)\|^2\dd\mu(y)&=\mu(Z)/v_g,\\
 \frac{1}{m_Z}\int_{G_T}|c_Z(x)|^2\dd\mu(x)
      &=\mu(Z)\left(1+
         \sum_{\substack{xK_g\in G_T/K_g\\xK_g\ne K_g}}|C_Z(x)|^2\right).
 \end{split}
\]
The latter equality holds in the extended nonnegative reals.
\end{proposition}
\begin{proof}
The first identity is immediate from the norm formula in
Proposition~\ref{prop:finite-bessel-unfold}, since \(v_Z/v_g\)
is the index there. For \(y\in Z\) the condition \(u y\in Z\) in
its unfolding is \(u\in U_T\cap Z\). Integrating in \(y\) over
\(Z\), left translation by such \(u\) removes it from the conjugation
argument. Absolute integrability follows from boundedness and
compactness of the two groups. Lemma~\ref{lem:product-integration}
then yields the asserted formula for \(c_Z\).
Right translation of the integration variable in \(Z\) gives
conjugacy invariance there. Since \(Z\) normalizes \(K_g\),
right translation of \(x\) by \(K_g\) only introduces a right
\(K_g\) multiplier into the phase function and preserves it.
Left invariance by \(K_g\) follows since \(h x\) is conjugate
by \(h\in K_g\) to \(x h\). On \(K_g\) every conjugate belongs
to \(K_g\) and has phase one. Finally all cosets \(xK_g\) have
the same positive mass \(\mu(K_g)\); there are countably many
by the basis argument of Lemma~\ref{lem:finite-unitary-level-haar}.
Summing the constant squared modulus on these cosets and dividing
by \(m_Z\) gives the second identity.
\end{proof}


\section{A diagonal average at and away from the level}

\begin{definition}\label{def:finite-simple-torus}
In \(G_T\) consider the diagonal subgroup \(\mathcal T_g\) with
entries \(d_a\) on \(e_a\) and \(1/\bar d_a\) on \(e_a^*\)
for \(1\le a\le\nu\), fixing the middle space, where
\(d_a\in H_g=\{v\in B_T^\times:v-1\in gB_T\}\).
This subgroup lies in \(K_g\) by Definition~\ref{def:finite-bessel-level}.
Give it product probability and set \(d_{\nu+1}=1\).
For \(u\in U_T\) denote by \(b_a(u)\) the \(e_a\)-coordinate of
\(u e_{a+1}\) when \(a<\nu\), and of \(u z\) when \(a=\nu\).
Their sum is \(\alpha(u)\) from Proposition~\ref{prop:full-phase}.
Write \(w\mid T\) for each field component of \(E_T\);
let \(O_w\) be its integer ring, \(\pi_w\) a uniformizer,
\(q_w\) its residue cardinality and \(\psi_w^0\) its factor of
\(\psi_T^0\).
\end{definition}

\begin{proposition}\label{prop:finite-diagonal-phase}
Define \(\Phi_g(y)=\int_{\mathcal T_g}\phi_g(t y t^{-1})\dd t\).
It vanishes outside \(U_TK_g\). For \(y=u h\) with \(u\in U_T\)
and \(h\in K_g\) it is
\[
       \Phi_g(u h)=\prod_{a=1}^{\nu}\prod_{w\mid T}
                                  A_w\bigl(b_a(u)_w\bigr),
\]
where for \(g\in\pi_w O_w\) one has
\[
              A_w(b)={\bf1}_{O_w}(b)\,\psi_w^0(b/g),
\]
and for \(g\in O_w^\times\) one has
\[
 A_w(b)=
 \begin{cases}
  1,& b\in O_w,\\
  -1/(q_w-1),& b\in\pi_w^{-1}O_w\setminus O_w,\\
  0,& b\notin\pi_w^{-1}O_w .
 \end{cases}
\]
For every intermediate \(Z\) of Definition~\ref{def:finite-bessel-vector},
\(C_Z(x)\) of Proposition~\ref{prop:finite-congruence-energy} equals
the mean of \(\Phi_g(yxy^{-1})\) for \(y\in Z\).
\end{proposition}
\begin{proof}
The group \(\mathcal T_g\) normalizes \(U_T\) by the flag equations
and \(K_g\) by Lemma~\ref{lem:finite-bessel-radical-volume}, so it
preserves the support under conjugation. For \(y=u h\), the last
factor remains in \(K_g\), whereas the simple-root entries of
\(t u t^{-1}\) are \((d_a/d_{a+1}) b_a(u)\) by diagonal
multiplication. The map taking \((d_a)\) to \((d_a/d_{a+1})\)
is an automorphism of \(H_g^\nu\): its inverse takes successive
products from the last entry. At every finite congruence quotient
it permutes the points and preserves uniform probability, hence
it preserves the compact probability by the cylinder argument of
Lemma~\ref{lem:finite-unitary-level-haar}. The components over
\(w\mid T\) are independent as \(B_T\) is their direct product.
The desired integral reduces to the average of
\(\psi_w^0(s b/g)\) over \(s\in O_w^\times,\ s-1\in gO_w\).

Proposition~\ref{prop:additive} applied in one component shows
that the additive annihilator of \(O_w\) for multiplication and
\(\psi_w^0\) is \(O_w\). Indeed the involution permutes the
component rings, and the test elements in its statement can be
varied independently in each component. A nontrivial character
on a compact additive subgroup has mean zero, since translation
by a point with character different from one multiplies the mean
by that value. The trivial character has mean one.
If \(g\) is a nonunit, write \(s=1+g z\) with \(z\in O_w\).
The uniform additive probability on \(z\) gives Haar probability
on these units: multiplication by \(1+g z_0\) induces the
measure-preserving affine map \(z\mapsto z_0+(1+g z_0)z\).
The mean is thus \(\psi_w^0(b/g)\) times the additive mean at \(b\).
If \(g\) is a unit, \(s\) ranges over all of \(O_w^\times\).
Its probability is additive measure restricted there and divided
by \(1-q_w^{-1}\), since multiplication by a unit preserves every
residue count. Subtracting the integral on \(\pi_w O_w\) from
that on \(O_w\) gives the numerator
\({\bf1}_{O_w}(b)-q_w^{-1}{\bf1}_{\pi_w^{-1}O_w}(b)\).
This evaluates all factors. Finally left translation by
\(t\in\mathcal T_g\subset Z\) preserves probability on \(Z\).
Integrate first over \(t\) in the conjugation argument using
Lemma~\ref{lem:product-integration}, which applies to these
bounded functions on compact groups.
\end{proof}

\begin{corollary}\label{cor:finite-integral-congruence-sum}
Put \(\bar K=K_T/K_g,\ \bar Z=Z/K_g\), and let \(\bar U\)
be the image of \(U_T\cap K_T\). Its character induced by
\(\chi_g\) is denoted by \(\bar\chi\). Then the part of the normalized
energy of Proposition~\ref{prop:finite-congruence-energy} with
\(x\in K_T\) is the explicit finite sum
\[
    \mu(Z)\sum_{a\in\bar K}
      \left|\,|\bar Z|^{-1}\sum_{\substack{z\in\bar Z\\zaz^{-1}\in\bar U}}
                                 \bar\chi(zaz^{-1})\right|^2 .
\]
The complementary part is supported only on those cosets \(xK_g\)
where some \(y\in Z\) has \(yxy^{-1}=u h\) with \(h\in K_g\) and
with every \(b_a(u)_w\in O_w\) for \(w\mid g\) and
\(b_a(u)_w\in\pi_w^{-1}O_w\) for the other \(w\mid T\).
\end{corollary}
\begin{proof}
If \(x=u h\in K_T,\ h\in K_g\), then \(u\in U_T\cap K_T\).
Conversely membership of the residue of \(x\) in \(\bar U\) gives
such a decomposition, with \(\phi_g(x)\) equal to its residue
character by Lemma~\ref{lem:finite-bessel-radical-volume}.
Uniform probability on \(Z\) pushes to uniform probability on
\(\bar Z\) because all its cosets have equal measure.
Proposition~\ref{prop:finite-congruence-energy} gives the sum.
For \(C_Z(x)\ne0\), the last average in
Proposition~\ref{prop:finite-diagonal-phase} has a nonzero integrand
at some \(y\). Its product formula gives exactly the stated
coordinate requirements.
\end{proof}


\section{Integral hyperbolic coordinates on the finite quadrics}

\begin{lemma}\label{lem:finite-auxiliary-planes}
Let \(a_*=3\nu+6,\ m_*=k-1-2a_*\), and \(1\le i\le n\).
At all sufficiently large scales \(m_*>0\). After removing the
first \(i-1\) outer pairs, the remaining finite lattice over \(B_T\)
is an orthogonal sum of \(\ell_i=n-i+a_*+1\) integral hyperbolic
planes and a standard lattice of rank \(m_*\). In these coordinates
the norm and product additive measure are
\[
  Q(a,b,w)=\operatorname{Tr}_{E_T/F_T}(a^*b)+w^*w,\qquad
             \dd a\,\dd b\,\dd w .
\]
Here \(a,b\in E_T^{\ell_i},\ w\in E_T^{m_*}\), every integer
coordinate ring has volume one, and \(a^*b=\sum_h\bar a_h b_h\).
The plane on \(v,v'\) in Lemma~\ref{lem:middle-hyperbolic-chart}
can be retained among the displayed planes.
\end{lemma}
\begin{proof}
That lemma gives the plane on \(v,v'\) and \(k-1\) mutually
orthogonal standard lines. At each \(v_0\in T\) choose
\(\epsilon\in B_{v_0}^\times\) with \(\epsilon\bar\epsilon=-1\).
For a split pair take \((-1,1)\). For an inert place choose a
residue solution by the cyclic norm calculation in
Proposition~\ref{prop:residue-norm-count}. A unit solution modulo
\(\pi^h\) lifts modulo \(\pi^{h+1}\): varying by \(\pi^h d\)
changes the next digit by the trace of \(\bar\epsilon d\) on
the residue field, which is onto by the differential calculation
in that proposition. Completeness gives the desired unit.
On two standard lines \(x_1,x_2\) put
\(e=x_1+\epsilon x_2,\ e'=x_1-\eta e\), with
\(\eta+\bar\eta=1\) from Lemma~\ref{lem:middle-hyperbolic-chart}.
Then \(e,e'\) have zero norm and pairing one; their determinant
in the given basis is \(-\epsilon\), a unit. Apply this to
\(a_*\) disjoint pairs and retain all outer planes. There are
enough lines by Proposition~\ref{prop:thick}. These integral basis
changes and their inverses preserve additive measure by
Lemma~\ref{lem:class-covol}. Taking the product over the places
gives the assertion.
\end{proof}

\begin{definition}\label{def:finite-linear-slice}
At a place \(v\in T\), write \(A=A_v,\ B=B_v,\ Q_v=|A/\pi A|\).
For \(a\in E_v^\ell\), \(\ell\ge1\), \(a\ne0\), put
\(d_v(a)=\max\{d\in\Z:a\in\pi^d B^\ell\}\).
Write \(L_a(b)=\operatorname{Tr}_{E_v/F_v}(a^*b)\).
Choose an integral \(A\)-basis of \(B\) as in
Lemma~\ref{lem:skew-covol}. Among the \(2\ell\) coefficients
of this linear form over \(F_v\), let \(c(a)\) be the first of
largest absolute value. On the hyperplane \(L_a(b)=h\) define
\(\lambda_{a,h}\) by solving for that coordinate and taking
\(|c(a)|_v^{-1}\) times product additive measure in the other
coordinates. For \(T\) take the product of these measures,
omitting the null set where \(a_v=0\) for some \(v\).
\end{definition}

\begin{lemma}\label{lem:finite-linear-slice-mass}
One has \(|c(a)|_v=Q_v^{-d_v(a)}\). For \(\theta\in F_v^\times\)
and all \(h\in F_v\),
\[
 \lambda_{a,h}(\theta B_v^\ell)=
    |\theta|_v^{2\ell-1}Q_v^{d_v(a)}
              {\bf1}_{\theta\pi^{d_v(a)}A_v}(h).
\]
Write \(\Delta_T(\theta)=\prod_{v\in T}|\theta_v|_v\) for
\(\theta\in F_T^\times\), and define
\[
       M_T(\ell)=\prod_{v\in T}
                    \frac{1-Q_v^{-2\ell}}{1-Q_v^{1-2\ell}}.
\]
For the product measures in Definition~\ref{def:finite-linear-slice},
\[
 \int_{a\in\theta B_T^\ell}
              \sup_{h\in F_T}\lambda_{a,h}(\theta B_T^\ell)\dd a
                       =\Delta_T(\theta)^{4\ell-2}M_T(\ell).
\]
\end{lemma}
\begin{proof}
The image \(L_a(B_v^\ell)\) is \(\pi^{d_v(a)}A_v\).
After division of \(a\) by \(\pi^{d_v(a)}\), at least one coordinate
has nonzero residue. Its trace pairing with a freely chosen
coordinate is onto on the residue field: in the field case use
the surjective trace in Proposition~\ref{prop:residue-norm-count},
and in the split case the two coordinates on the other side
are independent. All coefficients are integral and one is
therefore a unit, proving the asserted image and the value of
\(|c(a)|_v\). For the largest coefficient the solved coordinate
lies in \(\theta A_v\) precisely when \(h\in c(a)\theta A_v\),
provided the other coordinates lie in \(\theta A_v\). This proves
the first formula.

On \(B_v^\ell\) the measure of \(d_v(a)\ge h\) is \(Q_v^{-2\ell h}\)
for integers \(h\ge0\), since \(B_v^\ell\) is free of rank \(2\ell\)
over \(A_v\). The point \(a=0\) has zero measure by letting \(h\)
increase. Thus the integral of \(Q_v^{d_v(a)}\) on this lattice is
\((1-Q_v^{-2\ell})\sum_{h\ge0}Q_v^{(1-2\ell)h}\).
On dilating \(a\) by \(\theta\), the factor \(Q_v^{d_v(a)}\)
is multiplied by \(|\theta|_v^{-1}\) and its integration measure
by \(|\theta|_v^{2\ell}\). The supremum in the statement is attained
simultaneously at \(h=0\). Multiplying the single-place integrals
and using Lemma~\ref{lem:product-integration} proves the formula.
\end{proof}

\begin{proposition}\label{prop:finite-floor-norm-disintegration}
On the finite space in Lemma~\ref{lem:finite-auxiliary-planes}, let
\(\delta_\gamma^{[i]}\), \(\gamma\in F_T\), be the measure on
\(Q=\gamma\) given in the variables \((a,b,w)\) by
\[
                     \dd a\,\dd w\,\dd\lambda_{a,\gamma-w^*w}(b).
\]
It is finite on compact sets, integrates continuously in \(\gamma\)
on compact locally constant tests, and is invariant under every
\(E_T\)-linear isometry of this space. For every nonnegative Borel
test \(H(\gamma,z)\),
\[
 \int H(Q(z),z)\dd z
      =\int_{F_T}\int H(\gamma,z)\dd\delta_\gamma^{[i]}(z)\dd\gamma .
\]
For \(i=n\) the measure is the middle norm measure in
Theorem~\ref{thm:middle-fiber-chart}.
\end{proposition}
\begin{proof}
All chosen coordinate maps are Borel, since the ordering of
the finitely many coefficients partitions the nonzero \(a\)'s
into Borel sets. At fixed \(a,w\) and all but one coordinate
of \(b_v\), changing its last coordinate to
\(L_a(b_v)+(w^*w)_v\) gives exactly the factor \(|c(a)|_v^{-1}\).
The excluded zero sets have zero additive measure. Applying
Lemma~\ref{lem:product-integration} gives disintegration.

For a test bounded by \(C\) and supported in a product box of
radius \(\theta B_T\) on every entry, the solved-coordinate
integral is dominated, independently of \(\gamma\), by \(C\)
times the upper fibre volume of Lemma~\ref{lem:finite-linear-slice-mass}.
The \(w\) and \(a\) range over the displayed finite boxes, and
that lemma proves integrability of this majorant. More explicitly,
before integrating the free \(b\)-coordinates one can take
\(C\prod_v |c(a_v)|_v^{-1}\) times the indicators of their
boxes and the \(a,w\) boxes. For fixed \(a,w\) and free coordinates,
the solved coordinates are continuous in \(\gamma\).
Dominated convergence in Lemma~\ref{lem:integral-tools} thus gives
continuity for compact locally constant tests and local finiteness
for their indicator majorants. An isometry preserves \(Q\) and
ambient additive measure: its determinant has norm one from the
hermitian matrix identity, and one uses the determinant computation
of Lemma~\ref{lem:class-covol}. For a compact open indicator its
two disintegrations before and after the change have equal
integrals over every compact open set of \(\gamma\)'s. Their
continuous difference must vanish at every point, since every
open ball has positive measure. The countable-basis uniqueness
argument of Lemma~\ref{lem:product-integration}, on increasing
compact boxes, extends equality to all Borel sets. The same
argument compares at \(i=n\) with the continuous disintegration
of Theorem~\ref{thm:middle-fiber-chart} and
Proposition~\ref{prop:top-norm-jacobian}.
\end{proof}


\section{Fourier integration on the finite norm fibres}

\begin{lemma}\label{lem:finite-positive-fourier}
For \(\xi\in E_T^\times\) put
\(m_T(\xi)=\prod_{v\in T}|{\rm N}_{E_v/F_v}(\xi_v)|_v\).
Give \(E_T^b\) the additive measure for which \(B_T^b\) has volume
one, for any integer \(b\ge0\). If \(H\in L^1\cap L^2(E_T^b)\), then
\[
 \int_{E_T^b}\left|\int_{E_T^b}
       \psi_T^0(\xi z^*w)H(w)\dd w\right|^2\dd z
                         \le m_T(\xi)^{-b}\int |H|^2 .
\]
\end{lemma}
\begin{proof}
For a box \(D=\alpha B_T^b,\ \alpha\in F_T^\times\), the integral
over \(z\in D\) of \(\psi_T^0(\xi z^*(w-w'))\) is
\(|D|{\bf1}_{\xi^{-1}\alpha^{-1}B_T^b}(w-w')\).
Indeed conjugation preserves the integer rings and
Proposition~\ref{prop:additive} gives the annihilator coordinatewise;
translation gives zero for any nontrivial character as in
Proposition~\ref{prop:finite-diagonal-phase}. The box
\(\xi^{-1}\alpha^{-1}B_T^b\) has volume \(m_T(\xi)^{-b}/|D|\)
by the determinant and diagonal lattice calculation of
Lemma~\ref{lem:class-covol}. Expand the square on \(D\) and integrate
in \(z\) first, using \(|D|\|H\|_1^2<\infty\).
For each difference \(d=w-w'\), Cauchy--Schwarz and additive
invariance give \(\int |H(w)H(w-d)|\dd w\le\|H\|_2^2\).
The product of the two box volumes now gives the assertion on
\(D\). Enlarge \(D\) to exhaust the space and use
Lemma~\ref{lem:integral-tools}. For \(b=0\) the same calculation
has both volumes equal to one.
\end{proof}

\begin{theorem}\label{thm:finite-flat-fourier}
Use the space of Lemma~\ref{lem:finite-auxiliary-planes} at floor
\(i\), put \(\ell=\ell_i\), and let \(\mu,\mu'\) be either ambient
additive measure or any of the respective quadric measures
\(\delta_\gamma^{[i]}\). For \(\theta\in F_T^\times\) write
\(D_\theta\) for the box in which every coordinate is in
\(\theta B_T\), and set
\[
 A_{\rm amb}(\theta)=\Delta_T(\theta)^{4\ell},\qquad
 A_{\rm quad}(\theta)=\Delta_T(\theta)^{4\ell-2}M_T(\ell).
\]
Use \(A_\mu\) and \(A_{\mu'}\) with these meanings. If
\(h\in L^2(\mu)\) vanishes outside \(D_\theta\),
\(\theta'\in F_T^\times,\ \xi\in E_T^\times\), and \(z_0\)
is any fixed vector, then
\[
 \int_{D_{\theta'}}
 \left|\int \psi_T^0\bigl(\xi[z+z_0,y]\bigr)h(y)\dd\mu(y)\right|^2
                  \dd\mu'(z)
 \le m_T(\xi)^{-m_*}\, A_{\mu'}(\theta')A_\mu(\theta)
                                             \int |h|^2\dd\mu .
\]
Here \([\, ,\,]\) is the hermitian pairing with norm \(Q\).
\end{theorem}
\begin{proof}
Keep \(a,w\) as independent variables in both measures. For ambient
measure the volume of the \(b\)-fibre restricted to \(D_\theta\)
is at most \(M(a)={\bf1}_{\theta B_T^\ell}(a)\Delta_T(\theta)^{2\ell}\).
For a quadric use the same indicator times the supremum of the
fibre volume in Lemma~\ref{lem:finite-linear-slice-mass}.
In both cases \(\int M(a)\dd a=A_\mu(\theta)\). The identical
construction in the output gives \(M'(a')\) and \(A_{\mu'}(\theta')\).
The \(w\) boxes have finite measure, so \(h\) is also integrable.

At fixed \(a_z,a_y,w_y\), integrate \(h(y)\) on its \(b_y\)-fibre
with the phase corresponding to \(a_z^*b_y+[z_0,y]\), including
the factor \(\xi\). Denote the result as a function of \(w_y\)
by \(H(w_y)\), suppressing the two fixed \(a\)'s. Its squared
modulus is bounded by \(M(a_y)\int_{\rm fibre}|h|^2\)
by Cauchy--Schwarz. The other terms of \([z,y]\) are
\(w_z^*w_y+b_z^*a_y\). The latter term contributes a phase of
modulus one independent of \(w_y\). Therefore the modulus of the
original transform is bounded by the integral in \(a_y\) of the
modulus of the transform of \(H\) in Lemma~\ref{lem:finite-positive-fourier}.
This bound is independent of \(b_z\).
Weighted Cauchy--Schwarz in \(a_y\), with weight \(M(a_y)\),
bounds its square by \(A_\mu(\theta)\int |\widehat H(w_z)|^2
M(a_y)^{-1}\dd a_y\). Omit the zero set of \(M\), where the
corresponding fibres make no contribution. For almost every
\(a_y\), \(H\) has finite squared integral by its fibre bound
and is integrable since its \(w_y\) box has finite measure.
At each \(a_z\), integrate \(w_z\) over the full space and apply
Lemma~\ref{lem:finite-positive-fourier} with \(b=m_*\).
The result is at most \(A_\mu(\theta)m_T(\xi)^{-m_*}\int|h|^2\)
by cancellation of the fibre weight. Integration of the output
fibre upper bound \(M'(a_z)\) gives the other factor. The changes
of order follow from Proposition~\ref{prop:finite-floor-norm-disintegration}
and Lemma~\ref{lem:product-integration}.
\end{proof}

\begin{corollary}\label{cor:finite-integral-fibre-product}
At all sufficiently large arithmetic scales, uniformly over finite
sets \(T\), one has
\[
             1\le\prod_{i=1}^n M_T(\ell_i)\le1+2^{-5n}.
\]
For every \(\gamma\), the mass under \(\delta_\gamma^{[i]}\)
of the integral lattice box \(D_1\) is at most \(M_T(\ell_i)\).
\end{corollary}
\begin{proof}
The second assertion follows by integrating the upper fibre volume
of Lemma~\ref{lem:finite-linear-slice-mass} at \(\theta=1\)
over the \(w\) lattice, which has mass one.
For the first, put \(h=2\ell_i-1\ge6n+7\) by
Lemma~\ref{lem:finite-auxiliary-planes}. Each local factor is at
least one and its logarithm is at most \(2Q_v^{-h}\), since
\(-\log(1-t)=\int_0^t(1-u)^{-1}\dd u\le2t\) for \(0\le t\le1/2\).
By Lemma~\ref{lem:real-place-rings} at most \(f\) places lie over
each rational prime and each \(Q_v\) is at least that prime.
Also \(\sum_{m\ge2}m^{-h}\le
2^{-h}+\int_2^\infty t^{-h}\dd t\le2\cdot2^{-h}\).
Hence the logarithm of the product is at most
\(4nf\,2^{-(6n+7)}\).
Propositions~\ref{prop:thick} and~\ref{prop:level-prime}
give \(f\le X^{12},\ n\ge\sqrt X/2\) eventually at scale \(X\).
Thus \(4nf\,2^{-(6n+7)}\le 2^{-5n-1}\) eventually.
Using the exponential series, \(\exp(t)\le1+2t\) for
\(0\le t\le1/2\), which proves the bound.
\end{proof}


\section{Sections of the finite unitary orbits}

\begin{definition}\label{def:finite-floor-sections}
Number the outer pairs of Definition~\ref{def:row-space} from
\(1\) through \(\nu=n-1\). Let \(V_i\) be the orthogonal complement
of the first \(i-1\) pairs over \(E_T\), and put
\(\mathcal L_i=\mathcal L\cap V_i\). Let \(G_i\) consist of the
isometries of \(V_i\), extended by the identity on the removed
pairs, and let \(K_i\) be the stabilizer of \(\mathcal L_i\) in
\(G_i\). Set \(f_i=e_i\) for \(i<n\) and \(f_n=z\). For \(i<n\)
let \(N_i\) be the group of matrices \(H(l,t)\) of
Lemma~\ref{lem:heisenberg-layer} on the \(i\)-th pair with
intervening space \(V_{i+1}\) and Gram matrix \(J\). With
\(t=-\omega lJ^{-1}l^*+s\), give \(N_i\) the measure
\(\eta_i=\dd l\,\dd s\) of Lemma~\ref{lem:finite-bessel-radical-volume}.
For \(c\in E_T^\times\), let \(w_i(c)\) exchange the two basis
vectors by
\(e_i\mapsto c e_i^*\), \(e_i^*\mapsto\bar c^{-1}e_i\),
and fix \(V_{i+1}\). These assignments preserve the pairing.
Use \(d_i=\dim_{E_T} V_{i+1}=k-1+2(n-i)\) for the common
rank in the field factors. Define the measure on \(G_i\)
\[
  \theta_i=(c\mapsto w_i(c)^{-1})_*
                 \bigl(m_T(c)^{d_i}\dd c\bigr)*\eta_i .
\]
Here \(m_T\) is from Lemma~\ref{lem:finite-positive-fourier}.
The convolution is the pushforward by ordered multiplication.
Write \(u_T=|B_T^\times|\) for the additive volume of the
integral units. Empty products of measures have mass one.
\end{definition}

\begin{lemma}\label{lem:finite-lower-orbit-section}
For \(1\le i<n\) the image of \(\theta_i\) under
\(\sigma\mapsto\sigma^{-1}f_i\) is \(\delta_0^{[i]}\) of
Proposition~\ref{prop:finite-floor-norm-disintegration}. In the
coordinates \(d e_i+y+c e_i^*\), this map on the parameters of
\(\theta_i\) is a Borel bijection onto the open part of the
quadric with \(c\in E_T^\times\). Its complement is null for
\(\delta_0^{[i]}\), and \(\theta_i(K_i)\ge u_T\).
\end{lemma}
\begin{proof}
Take \(\sigma=w_i(c)^{-1}H(l,t)\).
The inverse of \(H(l,t)\) is \(H(-l,\bar t)\) by the trace
equation of Lemma~\ref{lem:heisenberg-layer}. Acting on \(ce_i^*\)
gives \(d=c\bar t,\ y=cJ^{-1}l^*\). Conjugation preserves additive
measure since it preserves each integral lattice and its residue
balls. The integral invertible matrix \(J\) preserves it as well,
by Lemma~\ref{lem:class-covol}. Hence at fixed \(c\) the change
from \(l\) to \(y\) gives \(\dd y=m_T(c)^{d_i}\dd l\).
Writing \(Q(y)=[y,y]\), the missing coordinate is
\[
             d=-\bar c^{-1}\bar\omega Q(y)-c s.
\]
For general \(\gamma\in F_T\) replace this by
\(d=\bar c^{-1}\bar\omega(\gamma-Q(y))-cs\). Then
\(\operatorname{Tr}(\bar c d)+Q(y)=\gamma\).
The lattice splitting
\(B_T=\bar\omega A_T\oplus iA_T\) follows by subtracting
\(\bar\omega\) times the trace, as in
Lemma~\ref{lem:middle-hyperbolic-chart}. In this splitting
\(\bar c d=\bar\omega(\gamma-Q(y))-c\bar c s\).
On the skew line the multiplier \(-c\bar c\) changes volume
by \(m_T(c)\), the same factor as multiplication by \(\bar c\)
on \(E_T\). This proves that \((\gamma,s)\mapsto d\) is an affine
bijection carrying product measure to \(\dd d\).
All inverse formulas on the open part are continuous in the
field coordinates.

To identify the slice at every \(\gamma\), first take a compact
open test on \(d,y,c\) supported where \(c,c^{-1}\) are bounded
in every field component. For \(\gamma\) in a compact neighborhood,
the inverse formula bounds \(s\) on the support by a fixed compact
box. Dominated convergence, Lemma~\ref{lem:integral-tools}, gives
continuity of its slice integral in \(\gamma\). Integrating over
any compact open set of \(\gamma\)'s gives its ambient integral
there by the preceding Jacobian and Lemma~\ref{lem:product-integration}.
Proposition~\ref{prop:finite-floor-norm-disintegration} gives the
same integral and continuity, hence equality at every \(\gamma\).
In the coordinates of Lemma~\ref{lem:finite-auxiliary-planes}
we may take the \(e_i^*\)-coefficient as the first \(a\)
coordinate, exchanging the two sides of this hyperbolic plane
if necessary. The set where any field component of \(c\) vanishes
has measure zero in \(\dd a\), hence in that disintegration
even on \(Q=0\). Exhaust the complement by compact open boxes
with \(c,c^{-1}\) bounded and use the uniqueness on generating
boxes from Lemma~\ref{lem:product-integration}. This proves the
measure assertion. For integral unit \(c\) the Weyl element is
integral and \(m_T(c)=1\). Restricting also \(H(l,t)\) to its
integral subgroup of volume one gives the lower bound.
\end{proof}

\begin{proposition}\label{prop:finite-top-orbit-section}
On \(Q=-N\) in \(V_n\), a Borel section \(\sigma(e)\) satisfies
\(\sigma(e)^{-1}z=e\) for \(\delta_{-N}^{[n]}\)-almost every \(e\).
Its pushforward measure \(\theta_n\) on \(G_n\) satisfies
\[
       (\sigma\mapsto\sigma^{-1}z)_*\theta_n=\delta_{-N}^{[n]},
                         \qquad \theta_n(K_n)\ge u_T .
\]
\end{proposition}
\begin{proof}
In the notation of Lemma~\ref{lem:middle-hyperbolic-chart} write
\(e=dv+y+cv'\), with \(y\) in the \(k-1\) standard lines. By
Theorem~\ref{thm:middle-fiber-chart} and
Proposition~\ref{prop:finite-floor-norm-disintegration} we may take
\[
  c\in E_T^\times,\qquad
  d=\bar c^{-1}\eta(-N-Q(y))+c\alpha,\qquad
  \dd\delta_{-N}^{[n]}=\dd c\,\dd y\,\dd\alpha ,
\]
where \(\alpha\in iF_T\). Let \(D\) act diagonally by
\(-\bar c^{-1},-c\) on \(v,v'\), fixing the standard lines.
It is an isometry and \(Dz=d_0v+cv'\) with
\(d_0=-N\eta/\bar c\), since \(z=N\eta v-v'\).
Put \(u=y/c,\ h=(d-d_0)/c\). Equality of the two norms
\(Q(e)=Q(Dz)\) gives \(h+\bar h=-Q(u)\).
The map \(A\) fixing \(v\), sending \(v'\) to \(v'+u+hv\), and
sending a standard-space vector \(y'\) to \(y'-[u,y']v\) is
triangular with unit diagonal. The added term in the pairing
of \(Av'\) and \(Ay'\) cancels, and
\(Q(Av')=Q(u)+h+\bar h=0\). All remaining basis pairings are
unchanged. Thus \(A\) is an isometry and \(ADz=e\).
Take \(\sigma(e)=(AD)^{-1}\). These are continuous matrix
formulas on the open chart, whose complement is null. For
\(c\in B_T^\times\), integral \(y\) and \(\alpha\in iA_T\),
each displayed matrix and its inverse is integral, since
\(N,\eta\) are integral. The product volume of this subset
of the chart is \(u_T\). This proves both assertions.
\end{proof}


\section{Group measure obtained from the finite quadrics}

\begin{definition}\label{def:finite-floor-haar}
In the groups of Definition~\ref{def:finite-floor-sections}, let
\(S_i\) be the stabilizer of \(f_i\) in \(G_i\) and put \(H=S_n\).
Take on \(H\) its bi-invariant measure \(\mu_H\) with
\(\mu_H(H\cap K_n)=1\). Indeed \(z\) has nonzero norm in every
field factor, so an isometry fixing it acts freely as an arbitrary
isometry of its perpendicular, the \(k\) original standard lines.
Conversely every such isometry extends by fixing \(z\).
Lemma~\ref{lem:finite-unitary-level-haar}, applied in that standard
basis, supplies bi-invariant measure. The subgroup \(H\cap K_n\)
is compact since \(H\) is closed, and open in \(H\) since \(K_n\)
is open. It has positive finite mass by the congruence cylinder
construction in that lemma, so the specified scaling exists.
Use \(\theta_i\) from Lemma~\ref{lem:finite-lower-orbit-section}
and Proposition~\ref{prop:finite-top-orbit-section} and define
\[
 \mathcal M_n=\mu_H*\theta_n,\qquad
 \mathcal M_i=\mathcal M_{i+1}*\eta_i*\theta_i\quad(i<n),
 \qquad \mathcal M_{\rm up}=\mu_H*\theta_n*\cdots*\theta_1 .
\]
\end{definition}

\begin{theorem}\label{thm:finite-floor-haar}
Each \(\mathcal M_i\) is a bi-invariant measure on \(G_i\),
positive and finite on \(K_i\). In particular the measure
\(\mu=\mathcal M_1\) on \(G_T\) satisfies
\[
       u_T^n\ \le\ \mu(K_T)\ \le\
              \prod_{i=1}^{n}M_T(\ell_i)\ \le\ 1+2^{-5n}
\]
at sufficiently large arithmetic scales, uniformly over \(T\).
\end{theorem}
\begin{proof}
For \(i<n\), the equation of an isometry fixing \(e_i\) forces
the bottom row in this three-block basis to be \((0,0,1)\), by
pairing with \(e_i\). The induced matrix on
\(e_i^\perp/E_Te_i\) is an isometry of \(V_{i+1}\).
Dividing out its block diagonal lift leaves exactly \(N_i\)
by the computation of Lemma~\ref{lem:heisenberg-layer}.
Thus \(S_i=G_{i+1}N_i\) with a unique factorization in these
two subgroups and with \(N_i\) normal. Conjugation by \(h\in
G_{i+1}\) changes the row \(l\) to \(lh^{-1}\), fixing \(s\):
the norm \(lJ^{-1}l^*\) is preserved. The determinant of \(h\)
has norm one, so \(\eta_i\) is invariant under this action by
Lemma~\ref{lem:finite-bessel-radical-volume}. It is also invariant
on both sides in \(N_i\). If \(\mathcal M_{i+1}\) is bi-invariant,
the product \(\nu_i=\mathcal M_{i+1}*\eta_i\) is consequently
bi-invariant on \(S_i\). Indeed multiplication by a member of
either subgroup either translates that factor alone or also
conjugates the other integration variable with its measure
preserved. Its mass on \(S_i\cap K_i\) is
\(\mathcal M_{i+1}(K_{i+1})\): integrality of an integral inverse
pair forces integrality of the inner diagonal block and its
inverse, after which the remaining layer must be integral, and
the converse holds by multiplication. The integral layer has
volume one. At \(i=n\) use \(\nu_n=\mu_H\), with corresponding
mass one.

To prove right invariance in both cases, take a nonnegative
Borel test \(B\) on \(G_i\). The function
\(g\mapsto\int_{S_i} B(sg)\dd\nu_i(s)\) depends only on
\(e=g^{-1}f_i\), since two representatives differ on the left
by \(S_i\) and \(\nu_i\) is right invariant.
On the full-measure open chart of
Lemma~\ref{lem:finite-lower-orbit-section} or
Proposition~\ref{prop:finite-top-orbit-section} it has a Borel
representative \(b(e)\) obtained by inserting the explicit
section. Extend \(b\) by zero off that chart. The pushforward
laws in those two results give \(\int B\dd\mathcal M_i=\int b
\dd\delta^{[i]}_{Q(f_i)}\). For each \(y\in G_i\), replacing
\(g\) by \(gy\) replaces \(e\) by \(y^{-1}e\). Both \(e\) and
\(y^{-1}e\) lie in the open chart almost everywhere because
the quadric measure is invariant by
Proposition~\ref{prop:finite-floor-norm-disintegration}.
Integration of \(b(y^{-1}e)\) therefore gives the same value.
The measurability and iterated integrations here follow from
the Borel parameter formulas and Lemma~\ref{lem:product-integration}.

For \(B={\bf1}_{K_i}\), the set of \(s\in S_i\) with
\(sg\in K_i\) is empty or a right coset of \(S_i\cap K_i\):
if \(s_0g\in K_i\), the condition is \(ss_0^{-1}\in S_i\cap K_i\).
In the latter case \(e=(s_0g)^{-1}f_i\in\mathcal L_i\).
The integral is thus at most
\(\nu_i(S_i\cap K_i)\delta^{[i]}_{Q(f_i)}(\mathcal L_i)\).
It is at least \(\nu_i(S_i\cap K_i)\theta_i(K_i)\), by restricting
\(g\) to \(K_i\) before pushing forward. The first bound is at
most \(\nu_i(S_i\cap K_i)M_T(\ell_i)\) by
Corollary~\ref{cor:finite-integral-fibre-product}, and the second
bound is at least \(\nu_i(S_i\cap K_i)u_T\) by the two section
results. The number \(u_T\) is positive because each of the finitely
many field factors contains an open ball of integral units.
Right invariance now gives local finiteness since every compact
set meets only finitely many open right cosets of \(K_i\).
Inversion carries this measure to a left invariant measure of
the same positive finite mass on \(K_i\). Uniqueness by residue
cylinders as in Lemma~\ref{lem:finite-unitary-level-haar} identifies
the inverse measure with the bi-invariant measure of that mass
in the integral basis of \(V_i\). The latter is preserved by
inversion, by the same uniqueness. This proves bi-invariance
at each induction step. Multiply the successive bounds, starting
at \(i=n\), and apply Corollary~\ref{cor:finite-integral-fibre-product}.
\end{proof}

\begin{corollary}\label{cor:finite-actual-quotient}
For the measure \(\dd u\) of Lemma~\ref{lem:finite-bessel-radical-volume},
one has \(\mu=\dd u*\mathcal M_{\rm up}\).
If \(A\) is a nonnegative Borel left \(U_T\)-invariant function
supported on \(U_T Z\) for \(K_g\subseteq Z\subseteq K_T\), then
\[
             \int A(y)\dd\mathcal M_{\rm up}(y)
                  =v_Z^{-1}\int_Z A(y)\dd\mu(y).
\]
In particular the integrals of \(\|W_Z(y)\|^2\) and of
\(\langle W_Z(y),W_Z(yx)\rangle\) against \(\mathcal M_{\rm up}\)
equal \(\mu(Z)/v_g\) and \(c_Z(x)\), respectively, with
\(\mu=\mathcal M_1\) in Definition~\ref{def:finite-cut-correlation}.
\end{corollary}
\begin{proof}
In the product defining \(\mathcal M_1\), first move the independent
\(N_1\) factor immediately preceding \(\theta_1\) to the front.
Every preceding factor lies in \(G_2\) and preserves \(\eta_1\)
by the conjugation calculation of Theorem~\ref{thm:finite-floor-haar}.
Repeat inside the remaining product with \(N_2,\ldots,N_\nu\).
The resulting leading product has measure \(\dd u\) by the
successive layer decomposition in
Lemma~\ref{lem:finite-bessel-radical-volume}. The other factors
are exactly \(\mathcal M_{\rm up}\). Lemma~\ref{lem:product-integration}
proves the convolution identity on nonnegative tests.
At \(y\in U_T Z\) the set of \(u\) with \(uy\in Z\) is a right
coset of \(U_T\cap Z\), of volume \(v_Z\). Elsewhere it is empty.
Apply the identity to \(A{\bf1}_Z\) to obtain the quotient formula.
Proposition~\ref{prop:finite-bessel-unfold} gives the required
support and invariance for both displayed tests. For the complex
test apply the formula first to its absolute value, which is bounded
on \(Z\) as both factors are bounded. Taking the real and imaginary
parts and using Proposition~\ref{prop:finite-congruence-energy}
gives the two values.
\end{proof}


\section{Products of coefficient integrals at vector stabilizers}

\begin{definition}\label{def:finite-normalized-energy}
Use \(\mu=\mathcal M_1\) from Theorem~\ref{thm:finite-floor-haar}.
For \(K_g\subseteq Z\subseteq K_T\), define
\[
 \begin{aligned}
 \mathcal E_T(Z)&=
 \frac{1}{\mu(Z)\mu(K_g)}
       \int_{G_T}\left|\int_Z\phi_g(yxy^{-1})\dd\mu(y)\right|^2
                                                        \dd\mu(x),\\
 \mathcal R_T(Z)&=\mu(Z)\!
       \sum_{\substack{xK_g\in G_T/K_g\\xK_g\ne K_g}}|C_Z(x)|^2 .
 \end{aligned}
\]
The sum runs over the other left cosets \(xK_g\) of the level
subgroup, and \(C_Z\) is the mean in
Proposition~\ref{prop:finite-congruence-energy}.
\end{definition}

\begin{proposition}\label{prop:finite-local-energy-product}
One has \(\mathcal E_T(Z)=\mu(Z)+\mathcal R_T(Z)
=m_Z^{-1}\int|c_Z|^2\dd\mu\).
For a partition \(T=T_1\sqcup\cdots\sqcup T_a\), write \(\mu_j\)
for the same group measure constructed over \(T_j\). Then
\(\mu=\bigotimes_j\mu_j\). If \(Z=\prod_j Z_j\) is a product of
intermediate groups, then \(\mathcal E_T(Z)=\prod_j
\mathcal E_{T_j}(Z_j)\), also in the extended nonnegative reals.
\end{proposition}
\begin{proof}
Proposition~\ref{prop:finite-congruence-energy} identifies the
inner integral with \(v_g c_Z(x)\) and gives the coset sum.
The normalizing denominator is \(v_g^2m_Z\), proving the first
identities. The coordinate sections of
Lemma~\ref{lem:finite-lower-orbit-section} and
Proposition~\ref{prop:finite-top-orbit-section} act componentwise.
The chart variables, additive measures and Weyl densities are
products, as are the layer measures. The groups \(H,H\cap K_n\)
in Definition~\ref{def:finite-floor-haar} split as products, and
the product of the respective \(\mu_{H,j}\) is bi-invariant with
the prescribed mass one on \(H\cap K_n\). Uniqueness on compact
open cylinders in Lemma~\ref{lem:finite-unitary-level-haar} gives
exactly \(\mu_H\). Products and group multiplications now commute
in all the convolutions defining \(\mathcal M_1\) by
Lemma~\ref{lem:product-integration}. This proves the measure
identity, independent of the integral bases used to describe the
quadric measures, by Proposition~\ref{prop:finite-floor-norm-disintegration}.

The group \(U_TK_g\) is the product of the corresponding sets.
On it the phase \(\chi_g\) is the product of its local phases by
Proposition~\ref{prop:additive} and the componentwise simple-root
sum in Definition~\ref{def:finite-bessel-level}. Off that product
at least one factor is zero. Thus \(\phi_g\) itself factors.
The compact inner integral factors by iterated integration and
the normalizing masses factor by the measure identity. The outer
integrand is nonnegative, so the outer integrals factor as well.
Each factor is at least \(\mu_j(Z_j)>0\). For the empty set of
places all masses and the identity coefficient are one.
\end{proof}

\begin{theorem}\label{thm:marked-energy-reduction}
Let \(T\) contain the places above \(g_0\) and \(g_1\), and let
\(S\) be any subset of the split pairs above \(p=g_0\) with one
prescribed condition per pair as in Definition~\ref{def:sieve-cell}.
Identify each pair with its underlying place of \(F\).
There is a product \(Z\) between \(K_g\) and \(K_T\) whose
members are exactly the integral isometries over \(T\) satisfying
these conditions. Its contribution from \(K_g\) is
\[
       \mu(Z)=\mu(K_T)(p^r-1)^{-|S|}.
\]
Put \(T_0=T\setminus S\), and write \(Z_v\) for its factor at a
selected \(v\in S\). The following estimates suffice at every
sufficiently large scale for
\(\mathcal E_T(Z)\le (1+b^{-6})b^{-6k|S|}\):
\[
 \mathcal R_{T_0}(K_{T_0})\le2^{-n},
       \qquad \mathcal R_{\{v\}}(Z_v)\le p^{-3k/4}
                         \quad (v\in S).
\]
\end{theorem}
\begin{proof}
At \(v\in S\) take the inverse image of the indicated vector or
covector stabilizer in the split reduction
\({\rm GL}_r(\mathbf F_p)\). The reduction is onto by
Proposition~\ref{prop:row-group}. At other places take the entire
integral subgroup. The product contains \(K_g\) by its defining
congruence, and is exactly the prescribed inverse image. Its
index follows from Proposition~\ref{prop:stabilizer-index}, since
the conditions have independent factors and each prescribed
vector is nonzero. Partition into cosets of equal group measure
to obtain the identity.

Use Proposition~\ref{prop:finite-local-energy-product} with one
block \(T_0\) and one for each selected pair. On \(T_0\) the
assumed remainder and Theorem~\ref{thm:finite-floor-haar} give
\[
       \mathcal E_{T_0}(K_{T_0})\le 1+2^{-5n}+2^{-n}\le1+b^{-6}.
\]
The last inequality holds eventually since \(n\ge\sqrt X/2\)
and \(\log b=X^\epsilon\), \(0<\epsilon<1/100\), by
Proposition~\ref{prop:thick}, where \(X\) denotes the arithmetic
scale. At one selected place \(\mu_v(K_v)\le2\) at those scales.
Thus the contribution from the level subgroup is at most
\(2/(p^r-1)\le4p^{-k}\), since \(r\ge k\).
Our other hypothesis gives
\(\mathcal E_{\{v\}}(Z_v)\le4p^{-k}+p^{-3k/4}\).
Proposition~\ref{prop:level-prime} gives \(p\ge b^{10}\),
so the last bound is at most \(b^{-6k}\) for sufficiently large
\(X\), independently of \(v\) and of its prescribed vector.
Multiplying proves the conclusion, including \(S=\varnothing\).
\end{proof}


\section{Denominators in a relative position of two finite flags}

\begin{definition}\label{def:finite-head-bounds}
Use the groups of Definition~\ref{def:finite-floor-sections}. For
\(0\le j\le\nu\) put \(U^{(j)}=N_1\cdots N_j\), and
\(P_j=U^{(j)}G_{j+1}\). The empty unipotent product is the identity.
Each inner group normalizes the outer layers by the proof of
Theorem~\ref{thm:finite-floor-haar}, so these are subgroups.
For \(p=(p_1,\ldots,p_j)\in(E_T^\times)^j\) write
\[
 w_p=\prod_{a=1}^j w_a(p_a),\quad
 D_a(p)=\prod_{h=1}^a p_h,\quad
 B(p)=\prod_{a=1}^jD_a(p),\quad d(p)=B(p)\overline{B(p)}.
\]
All empty scalar products are one. For a self-conjugate
\(d\in B_T\cap E_T^\times\) write
\(\mathcal B(d)=\{x\in G_T:d x\mathcal L\subseteq\mathcal L\}\).
This is a compact open set. Indeed in the integral basis it
is the intersection of the isometry equations with the product
of the compact open entry boxes \(d^{-1}B_T\); the equations
force invertibility since the Gram matrix is invertible in
every field factor. The identity \(x^{-1}=J_0^{-1}x^*J_0\)
with \(J_0,J_0^{-1}\) integral by Proposition~\ref{prop:row-group}
also shows that \(\mathcal B(d)\) is stable under inversion.
\end{definition}

\begin{theorem}\label{thm:finite-head-minors}
Let \(u_1,u_2\in U^{(j)},\ l\in G_{j+1}\). If
\(k_0=u_1 w_p l u_2\) belongs to \(K_T\), then
\[
 D_a(p)\in B_T\quad(1\le a\le j),\qquad
                  u_1,u_2,l\in\mathcal B(d(p)).
\]
More precisely, \(D_j(p)\) clears the entries of the block of
\(l\) on \(V_{j+1}\) when \(j>0\).
\end{theorem}
\begin{proof}
When \(j=0\) the assertion is simply that \(l=k_0\) is integral.
Order the column basis in blocks as
\((e_1,\ldots,e_j,V_{j+1},e_1^*,\ldots,e_j^*)\), and order the
row basis as \((e_1^*,\ldots,e_j^*,V_{j+1},e_1,\ldots,e_j)\).
In both occurrences of \(V_{j+1}\) use its integral basis.
In these coordinates the matrix is
\[
                 L\,\operatorname{diag}(p_1,\ldots,p_j,
                        l|_{V_{j+1}},\bar p_1^{-1},\ldots,
                        \bar p_j^{-1})\,R ,
\]
where \(L\) is \(u_1\) in the row basis and \(R\) is \(u_2\) in
the column basis. Relative to the three blocks, \(L\) is lower
and \(R\) upper triangular with identity middle blocks. The
first diagonal block of \(L\) is lower unitriangular and that
of \(R\) is upper unitriangular. Indeed the flag action in
Definition~\ref{def:unip} is upper unitriangular on the first
\(j\) isotropic vectors and inverse conjugate transpose on their
partners, by the pairing. This also follows by multiplying
the layer matrices of Lemma~\ref{lem:heisenberg-layer}.

The determinant of the first \(a\) rows and columns of the
displayed matrix is \(D_a(p)\): the two triangular corners have
determinant one, and terms using an intermediate index beyond
\(a\) vanish. Since \(k_0\) is integral, each \(D_a(p)\) is
integral in every component. For \(s>a\), take rows
\(1,\ldots,a-1,s\) and columns \(1,\ldots,a\).
The determinant is \(L_{sa}D_a(p)\): in the left minor the first
\(a-1\) rows have zero last column and determinant one on the
other columns. Using instead columns \(1,\ldots,a-1,s\) and
the first \(a\) rows gives \(D_a(p)R_{as}\). Thus each entry
of the first \(j\) columns of \(L\), respectively rows of \(R\),
has denominator dividing \(D_a(p)\) in column or row \(a\).
Bordering the first \(j\) rows and columns by one middle row
and one middle column gives \(D_j(p)\) times the corresponding
entry of \(l|_{V_{j+1}}\). Indeed in that bordered matrix one
may subtract combinations of the first \(j\) rows and columns
to remove the off-diagonal triangular blocks. The surviving
middle entry is the stated one because both middle blocks
of \(L,R\) are the identity. This argument is over \(E_T\);
the bordered determinant itself belongs to \(B_T\).

It remains to control entries outside the indicated rows and
columns. In the column basis write
\[
        R=\begin{pmatrix}A&C&E\\0&I&F\\0&0&A^{-*}\end{pmatrix},
                   \qquad F=-J^{-1}(A^{-1}C)^* ,
\]
where \(J\) is the integral unimodular Gram matrix on \(V_{j+1}\)
and \(A^{-*}=(A^{-1})^*\). These equations follow by pairing
the first and last block columns, then the middle and last block
columns. Write \(A=I+U\) with \(U\) strictly upper triangular.
In \(A^{-1}=\sum_{h=0}^{j-1}(-U)^h\), every entry of \(U^h\)
is a sum indexed by strictly increasing chains. A term uses
one entry per distinct row, with denominator the corresponding
\(D_a(p)\). In \(U^h C\) the additional \(C\)-factor uses the
last index of the chain, still distinct from the preceding rows.
Thus \(B(p)\) clears both \(A^{-1}\) and \(A^{-1}C\), as well
as every entry already known in the first \(j\) rows. Since
\(B(p)\) and its conjugate are integral, \(d(p)\) clears every
block of \(R\). For \(L\) take its conjugate transpose in the
row basis. This upper block matrix preserves the inverse Gram
matrix in that basis, which again has identity cross blocks and
an integral unimodular middle Gram matrix, by the isometry
equation. Its first \(j\) rows have the conjugate denominators.
The same chain argument clears all its entries by the same
self-conjugate \(d(p)\). Finally \(d(p)/D_j(p)\) is integral,
so the middle bound also gives \(l\in\mathcal B(d(p))\).
\end{proof}

\begin{corollary}\label{cor:finite-simultaneous-representative}
If \(y,w_p^{-1}y\in P_jK_T\), then every \(D_a(p)\) is integral.
There is \(s\in G_{j+1}\) with \(sy\in\mathcal B(d(p))\) and
\(s w_p^{-1}y\in\mathcal B(d(p)^2)\).
If \(j=\nu,\ l\in G_n\) and both \(W_Z(y)\) and
\(W_Z(lw_p^{-1}y)\) are nonzero, then \(y,l,lw_p^{-1}y\)
all belong to \(\mathcal B(d(p))\).
\end{corollary}
\begin{proof}
Write \(y=u h k,\ w_p^{-1}y=u'h'k'\), where \(u,u'\in U^{(j)}\),
\(h,h'\in G_{j+1}\) and \(k,k'\in K_T\).
Then \(k(k')^{-1}=u_1 w_p l_0 u_2\) with
\[
      u_1=h^{-1}u^{-1}h,\qquad l_0=h^{-1}h',
                     \qquad u_2=(h')^{-1}u'h'.
\]
Here \(w_p\) commutes with \(G_{j+1}\), since their nonidentity
blocks act on disjoint summands. Apply Theorem~\ref{thm:finite-head-minors}.
With \(s=h^{-1}\) the two representatives are \(u_1^{-1}k\)
and \(l_0 u_2 k'\). The bounds follow from inversion stability,
integrality of the compact factors and multiplication of
denominators. For the last assertion, the support formula
of Proposition~\ref{prop:finite-bessel-unfold} gives
\(y=u k,\ lw_p^{-1}y=v k'\), where \(u,v\in U_T=U^{(\nu)}\)
and \(k,k'\in Z\). Now \(k(k')^{-1}=u^{-1}w_p l^{-1}v\).
Apply the same theorem and inversion stability to \(u,v,l\),
and multiply by \(k,k'\) to obtain the claim.
\end{proof}


\section{The last finite pivot as a Gram kernel}

\begin{definition}\label{def:finite-top-correlation}
Use \(d_a\) and \(m_T\) of Definition~\ref{def:finite-floor-sections}
and \(\mu=\mathcal M_1\) of Theorem~\ref{thm:finite-floor-haar}.
For \(p\in(E_T^\times)^\nu,\ l\in G_n\) put
\[
 H(p)=\prod_{a=1}^\nu m_T(p_a)^{d_a/2},\qquad
 \mathcal J(p,l)=H(p)\int_{G_T}
                      W_Z(y)W_Z(lw_p^{-1}y)^*\dd\mu(y).
\]
Vectors are columns indexed by \(D_Z\) from
Definition~\ref{def:finite-bessel-vector}; the product is a
matrix, integrated entry by entry. Put \(\tau(l)=[z,l^{-1}z]\).
For a compact open additive subgroup \(L\subseteq E_T\) of volume
\(|L|>0\) and \(t\in E_T\), set
\[
       V(p,t;l)=\frac{{\bf1}_{t+L}(\tau(l))}
                      {\sqrt{m_Z|L|}}\,\mathcal J(p,l).
\]
Here \(m_Z\) is from Definition~\ref{def:finite-cut-correlation}.
Use the Hilbert--Schmidt inner product
\(\langle A,B\rangle=\operatorname{tr}(A^*B)\).
\end{definition}

\begin{lemma}\label{lem:finite-correlation-compact}
The integral defining \(\mathcal J(p,l)\) is absolutely convergent
for all \(p,l\). It vanishes unless every \(D_a(p)\) is integral,
and otherwise is supported in \(l\in G_n\cap\mathcal B(d(p))\).
As a function of \(p,l\) it is locally constant. For each \(p\)
there is a neighborhood on which \(\mathcal J(p',l)=\mathcal J(p,l)\)
for all \(l\in G_n\).
\end{lemma}
\begin{proof}
By Corollary~\ref{cor:finite-simultaneous-representative}, a nonzero
integrand requires the asserted integrality and, on that domain,
\(l,y\in\mathcal B(d(p))\). Multiplication by \(K_T\) preserves
the integral lattice, and the ball in
Definition~\ref{def:finite-head-bounds} is compact.
The integrand is bounded by Proposition~\ref{prop:finite-bessel-unfold}
and Haar measure is finite on compact sets by
Theorem~\ref{thm:finite-floor-haar}. Each entry is therefore
absolutely integrable. Give \(p\) a compact open neighborhood
on which the valuations of all \(p_a\) in every field factor
are fixed. There the integrality conditions are constant, \(H(p')\)
is constant, and the common bounding ball is fixed because its
defining principal fractional ideal \(d(p')^{-1}B_T\) is fixed.
If some minor fails integrality the whole function vanishes there.
Otherwise put \(D=\mathcal B(d(p))\).
The matrix-valued integrand is locally constant jointly in \(p,l,y\)
by the explicit formula for \(w_p\), continuity of multiplication,
and local constancy in Proposition~\ref{prop:finite-bessel-unfold}.
On the compact sets \(y\in D,\ l\in G_n\cap D\), cover the slice
at fixed \(p\) by finitely many product neighborhoods of constancy.
Intersect their neighborhoods of \(p\); then the integrand agrees
with its value at \(p\) for every \(y,l\) in these compact sets.
Outside them it vanishes whenever the product defining it could
be nonzero. Integrating proves the claimed uniformity. At fixed
\(p,l\), the same argument with \(y\in D\) gives constancy in a
neighborhood of \(l\). This proves joint local constancy.
\end{proof}

\begin{theorem}\label{thm:finite-top-gram}
For heads \(p,q\) of length \(\nu\) and last entries \(t,t'\in E_T\)
define
\[
 J_n((p,t),(q,t'))=\int_{G_n}
             \langle V(p,t;l),V(q,t';l)\rangle\dd\mathcal M_n(l).
\]
This is a jointly locally constant function, and its matrices on
all finite lists of heads and last entries are positive semidefinite.
For common last entry \(t=t'\) it is
\[
 \frac{1}{m_Z |L|}\int_{\tau(l)\in t+L}
           \langle\mathcal J(p,l),\mathcal J(q,l)\rangle
                                                  \dd\mathcal M_n(l).
\]
For fixed \(p,q\) the common-entry function is compactly supported
and integrable on \(E_T\). Write
\(L^\perp=\{\xi\in E_T:\psi_T^0(\xi a)=1\ \text{for all }a\in L\}\).
For every \(\xi\in E_T\),
\[
 \begin{split}
 &\int_{E_T}\psi_T^0(\xi t)J_n((p,t),(q,t))\dd t\\
 &\qquad=\frac{{\bf1}_{L^\perp}(\xi)}{m_Z}\int_{G_n}
      \psi_T^0(\xi\tau(l))\langle\mathcal J(p,l),\mathcal J(q,l)\rangle
                                                  \dd\mathcal M_n(l).
 \end{split}
\]
These formulas apply to the terminal pivot of
Definition~\ref{def:orbital-tower} after restricting the entries to
\(E_T^\times\).
\end{theorem}
\begin{proof}
The preceding lemma bounds the supports in \(l\) by compact sets
of finite \(\mathcal M_n\)-mass. On them the locally constant
functions \(\mathcal J\) are bounded by compactness. The varying
heads are locally constant as functions with values in this space
of bounded functions, by that lemma. The indicators are constant
when \(t,t'\) vary in their respective cosets of \(L\).
Consequently the defining integral exists, is jointly locally
constant, and its quadratic form on coefficients \((a_j)\) is
the integral of \(\|\sum_j a_j V(p_j,t_j;l)\|_{\rm HS}^2\).
The common-entry formula follows by squaring the indicator.
Its support in \(t\) is contained in the sum of the compact
image under \(\tau\) of a common support in \(l\) and the compact
group \(L\). Thus it is integrable, and
Lemma~\ref{lem:product-integration} interchanges the two integrals
after multiplication by the bounded character.
At fixed \(l\), change \(t\) to \(\tau(l)+a\) with \(a\in L\).
The integral of \(\psi_T^0(\xi a)\) on \(L\) is its volume if
\(\xi\in L^\perp\), and otherwise zero by translating by an
element with character different from one. This proves the
Fourier identity. The complement of \(E_T^\times\) is a finite
union of coordinate zero sets of additive measure zero.
Replacing individual zero components by one under the integration
therefore preserves the common-entry formulas. Both occurrences of
the terminal kernel in Definition~\ref{def:orbital-tower} use a
common last entry, for integration or for a matrix on a common
slice. Multiplying the two indicators in the definition gives
exactly the single indicated restriction on every such slice.
\end{proof}

\begin{corollary}\label{cor:finite-top-orbit-energy}
Fix \(p_1,\ldots,p_a\), coefficients \(b_1,\ldots,b_a\in\C\)
and a common \(t\). With \(\sigma(e)\) from
Proposition~\ref{prop:finite-top-orbit-section}, one has
\[
 \begin{split}
 &m_Z|L|\sum_{i,j=1}^a \bar b_i b_j J_n((p_i,t),(p_j,t))\\
 &\quad=\int_{[z,e]\in t+L}\int_H
       \left\|\sum_{i=1}^a b_i\mathcal J(p_i,h\sigma(e))\right\|_{\rm HS}^2
                              \dd\mu_H(h)\dd\delta^{[n]}_{-N}(e).
 \end{split}
\]
The inner integral is independent of the chosen orbit representative.
\end{corollary}
\begin{proof}
Apply the common-entry formula of Theorem~\ref{thm:finite-top-gram}
to the quadratic combination, giving a squared norm. Substitute
\(\mathcal M_n=\mu_H*\theta_n\) from Definition~\ref{def:finite-floor-haar}.
The vector \((h\sigma)^{-1}z\) is \(\sigma^{-1}z\) because
\(h\in H\). A second representative for the same vector differs
on the left by a member of \(H\), which may be absorbed using
right invariance of \(\mu_H\). Push forward \(\theta_n\) by the
orbit map using Proposition~\ref{prop:finite-top-orbit-section}
and Lemma~\ref{lem:product-integration} to obtain the identity.
\end{proof}


\section{A common integral subspace for two finite positions}

\begin{definition}\label{def:finite-close-volume}
Call two vectors \(z_1,z_2\in V_1\) close if \(z_1\in\mathcal L\)
and \(z_1-z_2\in g\mathcal L\), with the level
\(g\in A_T\cap F_T^\times\) of Definition~\ref{def:finite-bessel-level}.
Let \(I:E_T^{m_*}\to V_1\) include the \(m_*\) standard lines
of Lemma~\ref{lem:finite-auxiliary-planes}. Set
\[
        \mathcal P_T=\{t^{-1}u t:t\in G_T,\ u\in U_T\}K_g.
\]
For \(x\in K_T,\ y\in\mathcal P_T\) and a matrix \(C\) over \(E_T\)
with \(m_*\) columns and at most \(\nu\) rows, define the module
\[
  W=\{w:Cw=0,\ (x^{-1}Iw,(xy)^{-1}Iw)\text{ is close}\}.
\]
For \(v\in E_T^{m_*}\) write \(\mathfrak a_v=\{v^*w:w\in W\}\).
Throughout this section the arithmetic scale is large enough
that \(m_*>5n\), as follows from Proposition~\ref{prop:thick}.
\end{definition}

\begin{lemma}\label{lem:finite-content-summand}
The set \(W\) is a compact \(B_T\)-submodule contained in
\(B_T^{m_*}\), and it contains \(Y B_T^d\), where
\(d=m_*-3\nu\) and \(Y\) consists of the last \(d\) columns
of some \(V\in{\rm GL}_{m_*}(B_T)\).
For every \(v\), \(\mathfrak a_v\) is a product of compact
fractional ideals, allowing a zero ideal in some factors.
\end{lemma}
\begin{proof}
The closeness conditions are closed and \(B_T\)-linear, as is \(Cw=0\).
The first integrality condition implies \(Iw\in\mathcal L\),
because \(x\) stabilizes \(\mathcal L\). The integral direct
decomposition of Lemma~\ref{lem:finite-auxiliary-planes} then
gives \(w\in B_T^{m_*}\). This proves compactness.
Write \(y=t^{-1}u t h\) with \(u\in U_T,\ h\in K_g\).
In each field component \(\operatorname{rank}(u^{-1}-1)\le2\nu\):
in the three blocks consisting of all outer isotropic vectors,
the middle space and the dual vectors, the only potentially
nonzero entries lie in the first \(\nu\) rows or the last
\(\nu\) columns. This block form follows from the successive
matrices in Lemma~\ref{lem:heisenberg-layer}, also for \(u^{-1}\).
Stack \(C\) with the matrix \((u^{-1}-1)t x^{-1}I\).
Their combined rank is at most \(3\nu\) componentwise.

Over a single discrete valuation ring, multiply this stacked
matrix by a common nonzero denominator. Move an entry of least
valuation to the upper left by row and column permutations.
Its quotients into the entries of its column and row are integral,
so elementary row and column subtractions clear that column and
row apart from the pivot. Repeating on the smaller block gives
a diagonal matrix up to integral invertible row and column
operations, stopping when the remaining block is zero. The
number of nonzero pivots is the rank over the field.
Thus an invertible integral change of columns puts at least the
last \(d\) columns in the kernel of the original stacked matrix.
Choose these changes independently in all field components and
combine them to give \(V,Y\). If \(w\in YB_T^d\), then \(Cw=0\),
\(x^{-1}Iw\in\mathcal L\), and
\((xy)^{-1}Iw=h^{-1}x^{-1}Iw\). Since \(h^{-1}\in K_g\), the
two vectors are close. This proves the inclusion.

The linear image \(\mathfrak a_v\) is a compact \(B_T\)-module
in \(E_T\). Its components split independently by multiplication
by the coordinate idempotents of \(B_T\). A nonzero compact
submodule of a local field is bounded, so the integer valuations
of its nonzero elements have a minimum. Any member of minimum
valuation generates the whole submodule over the integer ring,
by division into every other member. This proves the ideal
description.
\end{proof}

\begin{theorem}\label{thm:finite-projective-moment}
Use \(W,Y,d\) from Lemma~\ref{lem:finite-content-summand}, and define
\[
       \phi(v)=|\{c\in E_T:cY^*v\in B_T^d\}|.
\]
Use the value \(+\infty\) when this set has infinite measure.
Then \(1/|\mathfrak a_v|\le\phi(v)\), with \(1/0=+\infty\).
Moreover the set of \(\alpha\in E_T\) such that for every \(w\in W\)
there are close pairs \((z_1,z_2),(z'_1,z'_2)\) satisfying
\[
       (\alpha v)^* w
                    =g^{-1}\bigl([z_1,z'_1]-[z_2,z'_2]\bigr)
\]
has volume exactly \(1/|\mathfrak a_v|\). For
\(\beta\in E_T^\times\) and real \(0<s<d\) one has
\[
 \int_{\beta^{-1}B_T^{m_*}}\phi(v)^s\dd v
      =m_T(\beta)^{-(m_*-s)}
                \prod_{w\mid T}\frac{1-q_w^{-d}}{1-q_w^{-(d-s)}}.
\]
Here \(w\) ranges over the field components of \(E_T\), of residue
cardinality \(q_w\). In particular one may take \(s=d-2n\).
Then the \(n\)-th power of the product on the right is at most
\(1+2^{-n}\) at sufficiently large scales, uniformly in \(T\).
\end{theorem}
\begin{proof}
For close pairs all four vectors lie in \(\mathcal L\).
Subtract the two pairings by changing one argument at a time.
The differences of the arguments belong to \(g\mathcal L\),
and \(\bar g=g\), so the divided difference is integral by
Proposition~\ref{prop:row-group}. Conversely take \(z_1=z_2=e_1\),
\(z'_1=gb e_1^*\), \(z'_2=0\), for any \(b\in B_T\).
Both pairs are close and the divided difference is \(b\).
Thus the asserted scalar set is exactly
\(\{\alpha:\bar\alpha\mathfrak a_v\subseteq B_T\}\).
For a nonzero ideal \(\pi_w^a O_w\) in a field component,
the scalars multiplying it into \(O_w\) form \(\pi_w^{-a}O_w\),
of reciprocal additive volume. A zero component gives the
full field of infinite volume. All nonzero components give
open sets of positive volume. Taking products, and using the
lattice-preserving conjugation, proves the asserted equality.
Since \(YB_T^d\subseteq W\), the subset \(v^*YB_T^d\subseteq
\mathfrak a_v\) has the corresponding scalar set of volume
\(\phi(v)\), giving the inequality.

The change of variables \(v\mapsto V^*v\) preserves every
box \(\beta^{-1}B_T^{m_*}\) and additive measure, as its
matrix and inverse are integral and the map is \(E_T\)-linear.
This follows by partition into residue boxes as in
Lemma~\ref{lem:class-covol}. Its last \(d\) coordinates are
\(Y^*v\). On the unit box, at component \(w\) let \(h\) be the
minimum of their valuations. For \(j\ge0\), the fraction with
\(h\ge j\) is \(q_w^{-dj}\), since all \(d\) coordinates must
lie in \(\pi_w^j O_w\). The set with all zero has volume zero.
The scalar bound there is \(q_w^h\), so the integral of its
\(s\)-th power is
\((1-q_w^{-d})\sum_{j\ge0}q_w^{-(d-s)j}\).
These formulas also give measurability, including on the
exceptional zero sets. Under \(v=\beta^{-1}v_0\) the scalar
volume becomes \(m_T(\beta)\phi(v_0)\), and the vector measure
changes by \(m_T(\beta)^{-m_*}\). Lemma~\ref{lem:product-integration}
and the geometric sum give the formula.

For \(d-s=2n\), bound the logarithm of each factor by \(2q_w^{-2n}\)
as in Corollary~\ref{cor:finite-integral-fibre-product}.
Each rational prime lies under at most \(2f\) field components
by Lemma~\ref{lem:real-place-rings}, and \(q_w\) is at least
that prime. Thus the logarithm of the \(n\)-th power is at most
\(8nf\,2^{-2n}\), using the series comparison in that corollary.
Since \(f\le X^{12}\) and \(n\ge\sqrt X/2\) eventually, this
is at most \(2^{-n-1}\) eventually. Exponentiation as there
proves the stated bound.
\end{proof}


\section{Congruence displacement detected by flag ideals}

\begin{definition}\label{def:finite-flag-ideal}
For the remainder of this section take a single finite place
\(v\) of \(F\), of residue cardinality \(Q\), with
\(A=A_v,\ B=B_v\) and a uniformizer \(\pi\in A\).
Use the measures and groups above with \(T=\{v\}\).
For \(a,b'\in G_T\), let
\[
 \begin{split}
 \mathcal L(a,b')&=\{w:(a^{-1}w,(b')^{-1}w)\text{ is close}\},\\
 I_i(a,b')&=\{[e_i,w]:w\in\mathcal L(a,b'),\
                             [e_j,w]=0\text{ for }j<i\},\\
 A_i(a,b')&=|I_i(a,b')|\qquad(1\le i\le\nu).
 \end{split}
\]
Let \(h(x)\) be the least integer \(h\ge0\) with
\((x-1)\mathcal L\subset g\pi^{-h}\mathcal L\).
Thus \(h(x)=0\) exactly when \(x\in K_g\): if \(x-1\) has entries
in \(gB\), then \(x\) is integral and its inverse is integral
by the isometry equation and Proposition~\ref{prop:row-group}.
\end{definition}

\begin{lemma}\label{lem:finite-flag-invariance}
Each \(I_i(a,b')\) is a principal fractional \(B\)-ideal with
invertible generator in \(E_v\). Its positive finite volume is
a Borel function, invariant under simultaneous left multiplication
by \(P_i\) of Definition~\ref{def:finite-head-bounds}.
If \(a\in P_iK_T\), then \(I_i(a,b')\subseteq B\).
If \(\langle W_Z(a),W_Z(ax)\rangle\ne0\), there exist
\(u,v_1\in U_T\) with \(k_1=ua,\ k_2=v_1 ax\in Z\) and
\(k_1^{-1}k_2\in K_g\). In particular \(x\in\mathcal P_T\)
and every \(I_i(a,ax)\subseteq B\).
\end{lemma}
\begin{proof}
The close-vector module is a compact open full lattice: it is the
intersection of \(a\mathcal L\) with the inverse image of
\(g\mathcal L\) under \(a^{-1}-(b')^{-1}\). It is open, closed
and contained in a compact set. Its intersection with the stated
orthogonal subspace therefore contains sufficiently small multiples
of the entire lattice in that subspace, in particular of \(e_i^*\).
Projection gives a compact module containing a nonzero open ideal.
The last argument of Lemma~\ref{lem:finite-content-summand} makes
its ideals principal, with no zero field components.

Membership of a scalar \(c\) in this image is open in
\((a,b',c)\): perturb a witness by \((c'-c)e_i^*\) at a nearby
scalar \(c'\); the earlier orthogonalities persist, and the two
lattice conditions are open also in \(a,b'\).
Integration of the resulting indicator in \(c\) gives Borel
measurability by Lemma~\ref{lem:product-integration}.
For \(q\in P_i\), the vector \(q^{-1}e_j\) differs from \(e_j\)
by a combination of earlier frame vectors for \(j\le i\).
The change \(w\mapsto q w\) therefore preserves the scalar and
earlier orthogonalities, and preserves closeness after left
multiplication of \(a,b'\) by \(q\). This proves invariance.
At \(a\in K_T\) the first vector condition forces \(w\in\mathcal L\),
giving an integral scalar. The same follows at \(a\in P_iK_T\).

A nonzero inner product has a coordinate with both entries
nonzero, say at \(K_g c,\ c\in Z\). By
Definition~\ref{def:finite-bessel-vector}, the points \(a,ax\)
then belong respectively to \(U_T K_g c\) and \(U_T K_g c\).
Choose the indicated unipotent multipliers taking them to
\(k_1,k_2\in K_g c\). Normality of \(K_g\) in \(Z\) gives
\(k_1^{-1}k_2\in K_g\), and
\(x=k_1^{-1}u v_1^{-1}k_1(k_1^{-1}k_2)\in\mathcal P_T\) by
Definition~\ref{def:finite-close-volume}. Finally \(U_T\subset P_i\):
after removing \(i\) outer layers the remaining factors belong to
\(G_{i+1}\). Thus \(a\in P_iK_T\) for all \(i\).
\end{proof}

\begin{theorem}\label{thm:finite-flag-displacement}
If \(\langle W_Z(a),W_Z(ax)\rangle\ne0\), then
\[
             h(x)\le2\sum_{i=1}^\nu
                                  \bigl(-\log_Q A_i(a,ax)\bigr).
\]
In particular for \(x\notin K_g\) some \(i\) satisfies
\[
       A_i(a,ax)\le Q^{-\max\{1,h(x)/(2\nu)\}} .
\]
\end{theorem}
\begin{proof}
Take \(u,v_1,k_1,k_2\) from Lemma~\ref{lem:finite-flag-invariance},
and put \(S=v_1u^{-1}\in U_T\).
For each \(i\) choose a generator \(t_i\) of \(I_i(a,ax)\) and
a witness \(w_i\) with \([e_i,w_i]=t_i\) and the earlier
coordinates zero. Put \(\Delta=\prod_i t_i\).
Use the matrices of columns \(Y=(uw_i)_i,\ Y'=(v_1w_i)_i\).
They are integral and \(Y'-Y\) has entries in \(gB\). Indeed,
writing \(z=a^{-1}w_i,\ z'=(ax)^{-1}w_i\), the two columns are
\(k_1z,k_2z'\), with \(z,z'\) integral and \(z-z'\in g\mathcal L\).
Also \(k_1^{-1}k_2\in K_g\).
In the decomposition into all outer vectors, \(V_n\), and all
dual vectors (the duals in the same order), the bottom blocks
\(D,D'\) of \(Y,Y'\) are lower triangular with diagonal
\(t_1,\ldots,t_\nu\). This follows from the flag action of \(u,v_1\)
as in Lemma~\ref{lem:finite-flag-invariance}. Both determinants
are \(\Delta\); the adjugate formula makes
\(\Delta D^{-1},\Delta (D')^{-1}\) integral. In particular
\(\Delta\) is an invertible scalar in the field algebra and
is integral.

Write \(YD^{-1}=(X,Y_0,1)^{\mathsf T}\), meaning three stacked
blocks with identity bottom block. Then \(\Delta X,\Delta Y_0\)
are integral, and
\((S-1)YD^{-1}=(Y'-Y)D^{-1}\) has entries in \(g\Delta^{-1}B\).
In the same blocks put
\[
             S=\begin{pmatrix}A_0&B_0&C_0\\0&1&F\\0&0&A_0^{-*}
                                              \end{pmatrix}.
\]
Its bottom equation gives \(A_0^{-*}=D'D^{-1}\); hence
\(A_0^*=D(D')^{-1}\) and
\(A_0^*-1=(D-D')(D')^{-1}\in g\Delta^{-1}M_\nu(B)\).
In particular \(A_0\) is bounded entrywise by the fractional
ideal \(\bar\Delta^{-1}B\), and \(A_0-1\) by \(g\bar\Delta^{-1}B\).
The middle difference is \(F\), so its entries belong to
\(g\Delta^{-1}B\). The isometry equation gives
\(B_0=-A_0 F^*J\), for the integral unimodular middle matrix
\(J\), exactly as in Theorem~\ref{thm:finite-head-minors}.
Its entries belong to \(g\bar\Delta^{-2}B\).
The top difference is \((A_0-1)X+B_0Y_0+C_0\); this puts
the entries of \(C_0\) in \(g(\Delta\bar\Delta^2)^{-1}B\).
The bottom difference \(A_0^{-*}-1=(D'-D)D^{-1}\) also has the
required bound. All entries of \(S-1\) consequently lie in
\(g(\Delta\bar\Delta)^{-2}B\). This last ideal is conjugation
invariant, so \(S^{-1}-1=J_0^{-1}(S-1)^*J_0\) has entries in it
as well, with \(J_0\) the full Gram matrix.
Since \(x=k_1^{-1}S^{-1}k_2\) and \(k_1^{-1}k_2-1\) has entries
in \(gB\), the same bound holds for \(x-1\).

The valuation of \(\Delta\bar\Delta\) in \(F_v\) is
\(-\log_Q m_T(\Delta)\) by Lemma~\ref{lem:finite-positive-fourier}.
Each \(|t_i B|=m_T(t_i)=A_i(a,ax)\), by the additive determinant
calculation of Lemma~\ref{lem:middle-hyperbolic-chart}.
This proves the first inequality. The individual exponents
\(-\log_Q A_i\) are nonnegative integers by integrality of
the ideals. When \(h(x)>0\) their maximum is positive and at
least \(h(x)/(2\nu)\), which proves the second assertion.
\end{proof}


\section{Integrable pairings on a finite flag}

\begin{definition}\label{def:finite-averaged-pairings}
At one place as in Definition~\ref{def:finite-flag-ideal},
extend \(P_0,\ldots,P_\nu\) by \(P_n=U_T H\) and put
\(\mathcal N_i=P_iK_T\). Write \(M_i=M_T(\ell_i)\),
\(s=m_*-3\nu-2n\), and
\[
 C_v=\prod_{w\mid v}
 \frac{1-q_w^{-(m_*-3\nu)}}{1-q_w^{-2n}}.
\]
For two lists \(p,q\) of the same length \(j_0\le\nu\), put
\(d_*=w_q^{-1}w_p\) and
\(\lambda_a=p_a/q_a\) for \(a\le j_0\), \(\lambda_a=1\) for
\(j_0<a\le n\). This diagonal isometry acts by \(\lambda_a\)
on \(e_a\), by \(\bar\lambda_a^{-1}\) on \(e_a^*\), and fixes
\(V_n\). Set
\[
 T_i=\mu_H*\theta_n*\cdots*\theta_{i+1},\quad T_n=\mu_H,\qquad
 K_i(a,b')=\int\langle W_Z(\sigma a),W_Z(d_*\sigma b')\rangle
                                               \dd T_i(\sigma).
\]
In \(T_0\) the last factor is \(\theta_1\), so
\(T_0=\mathcal M_{\rm up}\). All these definitions allow empty
lists. The notation \(K_i(a,b')\) denotes a complex pairing;
\(K_i\) without arguments in Definition~\ref{def:finite-floor-sections}
denotes a compact group.
\end{definition}

\begin{proposition}\label{prop:finite-average-mass}
For all \(a\in G_T\) and \(1\le i\le n\),
\[
 \theta_i\{\sigma:\sigma a\in\mathcal N_i\}\le M_i.
\]
For \(i<n\) and \(x\in\mathcal P_T\) one also has
\[
 \int_{\sigma a\in\mathcal N_i}
                    A_i(\sigma a,\sigma ax)^{-s}\dd\theta_i(\sigma)
                                                        \le M_iC_v .
\]
Finally \(\mu_H\{h:ha\in U_TK_T\}\le1\) for every \(a\).
\end{proposition}
\begin{proof}
For \(i\le n\), the quotient map taking the middle block of
\(P_{i-1}=U^{(i-1)}G_i\) has kernel \(U^{(i-1)}\) and sends
\(P_i\) onto \(S_i=\operatorname{Stab}_{G_i}(f_i)\).
For \(i<n\) this follows from the factorization \(S_i=G_{i+1}N_i\)
in Theorem~\ref{thm:finite-floor-haar}; at \(i=n\) it is the
definition of \(P_n\). These groups contain the outer unipotents
and are nested. Their sets \(\mathcal N_i\) are Borel since they
are unions of open cosets of \(K_T\).
If \(\sigma a\in\mathcal N_i\), then \(a\in P_{i-1}K_T\), as
\(\sigma\in G_i\). For a \(P_i\)-invariant nonnegative Borel
function \(F\), the integral \(\int F(\sigma a)\dd\theta_i\)
is left \(P_{i-1}\)-invariant as a function of \(a\).
Indeed a member of \(U^{(i-1)}\) conjugates through \(\sigma\)
into the same subgroup of \(P_i\). For \(h_0\in G_i\),
\(\sigma\mapsto F(\sigma a)\) depends only on \(\sigma^{-1}f_i\),
since \(S_i\subset P_i\); on the full-measure orbit chart its
representative is Borel by Lemma~\ref{lem:finite-lower-orbit-section}
and Proposition~\ref{prop:finite-top-orbit-section}. Replacing
\(\sigma\) by \(\sigma h_0\) preserves that integral by quadric
invariance in Proposition~\ref{prop:finite-floor-norm-disintegration}.
Thus in both asserted \(\theta_i\)-integrals one may assume
\(a\in K_T\), since the weights have the stated invariance by
Lemma~\ref{lem:finite-flag-invariance}.
If \(\sigma a\in\mathcal N_i\), there is \(q'\in P_i\) with
\(q'\sigma a\in K_T\), hence \(q'\sigma\in K_T\cap P_{i-1}\).
Its middle block lies in \(K_i\), with inverse integral by the
inverse block equation, and is \(s'\sigma\) for \(s'\in S_i\).
Consequently \(\sigma^{-1}f_i=(s'\sigma)^{-1}f_i\in\mathcal L_i\).
The orbit laws and Corollary~\ref{cor:finite-integral-fibre-product}
give the first estimate.

For \(i<n\) use \(W\) of Definition~\ref{def:finite-close-volume}
with integral first position \(a\), translation \(x\), and no
row equations. For \(w\in W\), the vector \(\sigma Iw\) is
orthogonal to all earlier frame vectors and is admissible in
\(\mathcal L(\sigma a,\sigma ax)\). Its scalar is
\([f_i,\sigma Iw]=v_+^*w\), where \(v_+\) is the projection of
\(\sigma^{-1}f_i\) onto the \(m_*\) standard lines. Hence the
reciprocal of \(A_i(\sigma a,\sigma ax)\) is at most \(\phi(v_+)\)
by Theorem~\ref{thm:finite-projective-moment}. Integrate on the
integral quadric box just obtained. In the disintegration of
Proposition~\ref{prop:finite-floor-norm-disintegration}, integrate
the \(b\)-fibre and then the hyperbolic \(a\)-coordinates using
Lemma~\ref{lem:finite-linear-slice-mass}. Their upper bound is
\(M_i\) independently of \(v_+\). The moment integral on the
remaining integral box is \(C_v\) by
Theorem~\ref{thm:finite-projective-moment} with \(\beta=1\).

For the final bound, suppose \(h_0a=u_0k_0,\ ha=uk\) with
\(h_0,h\in H,\ u,u_0\in U_T,\ k,k_0\in K_T\). Then
\(u^{-1}hh_0^{-1}u_0=kk_0^{-1}\) is integral, and its middle
diagonal block is the action of \(hh_0^{-1}\) on \(V_n\),
since the two unipotents have identity middle diagonal block.
The inverse matrix gives integrality of the inverse block.
Thus \(h\in(H\cap K_n)h_0\), of mass one by
Definition~\ref{def:finite-floor-haar}. An empty set has measure
zero, proving the general assertion.
\end{proof}

\begin{theorem}\label{thm:finite-pairing-laws}
The integrals defining \(K_i\) are absolutely convergent and
Borel for all arguments. Each vanishes unless \(a\in\mathcal N_i\)
and has absolute value at most
\((v_Z/v_g)\prod_{b=i+1}^{n}M_b\). One has
\[
 K_{i-1}(a,b')=\int K_i(\sigma a,\sigma b')\dd\theta_i(\sigma),
 \qquad K_0(1,x)=c_Z(x)\quad(j_0=0)
\]
for the recurrence \(1\le i\le n\) and the indicated empty
head, respectively. For \(u,u'\in U^{(i-1)}\), \(1\le i\le n\),
\[
 K_i(ua,u'b')=
       \overline{\chi_g(u)}\chi_g(d_*u'd_*^{-1})K_i(a,b').
\]
For \(i\ge j_0\), \(K_i\) is invariant under simultaneous left
multiplication of both arguments by \(P_i\cap G_{j_0+1}\).
For \(i>j_0\) its modulus is invariant under simultaneous
left multiplication by all of \(P_i\).
\end{theorem}
\begin{proof}
The pointwise norm formula of Proposition~\ref{prop:finite-bessel-unfold}
bounds the raw pairing by \(v_Z/v_g\) and restricts its first
argument to \(U_TZ\). At \(H\) the last bound of
Proposition~\ref{prop:finite-average-mass} applies, and positivity
of the support measure implies \(a\in U_THK_T=\mathcal N_n\)
because \(H\) normalizes \(U_T\).
Apply its first bound successively to the integral of the
absolute value. The support at each step lies in
\(\mathcal N_{i-1}\) by the block observation in its proof.
This gives the asserted estimates, Borel measurability and the
recurrence by Lemma~\ref{lem:product-integration}.
Corollary~\ref{cor:finite-actual-quotient} identifies the empty-head value.

For the separate translation formula, \(T_i\) is carried by
\(G_{i+1}\) for \(i<n\), and by \(H\) at \(i=n\).
These multipliers normalize \(U^{(i-1)}\) and leave the
simple-root coordinates of its elements unchanged. Indeed a
multiplier in \(G_{i+1}\) fixes both frame vectors entering
each of the first \(i-1\) coordinates. All later simple entries
of an element of that subgroup are zero by the block equations
in Lemma~\ref{lem:heisenberg-layer}. At the last step \(H\)
also fixes the successor \(z\). Conjugation by \(d_*\) multiplies
the \(a\)-th simple entry by \(\lambda_a/\lambda_{a+1}\);
here \(\lambda_n=1\). Thus the two characters from translating
the arguments in Proposition~\ref{prop:finite-bessel-unfold}
are constant in the integral, giving the formula.

For simultaneous invariance at \(i=n\), the \(H\)-factor is
absorbed by right invariance of \(\mu_H\); \(d_*\) commutes
with \(H\). A simultaneous unipotent multiplier in
\(U_T\cap G_{j_0+1}\) has equal first and second character.
For the induction from \(i\) to \(i-1\ge j_0\), factor the given
multiplier as a member of \(G_i\) times one of
\(U^{(i-1)}\cap G_{j_0+1}\) using the middle block. Conjugating
the latter by a member of \(G_i\) stays in the same subgroup
and it can be removed by the induction hypothesis. As a function
of \(\sigma\in G_i\) the resulting simultaneous integrand is
left invariant by \(S_i\), which lies in \(P_i\cap G_{j_0+1}\).
The orbit-section argument of Proposition~\ref{prop:finite-average-mass}
therefore absorbs the \(G_i\)-factor. This proves the first
simultaneous invariance.

For \(j_0<i<n\) a member of \(U^{(i)}\) conjugated by the upper
multiplier in \(G_{i+1}\) can change its \(i\)-th simple entry,
but that entry has character multiplier
\(\lambda_i/\lambda_{i+1}=1\). All entries of index at most
\(j_0\) are preserved, as both frame vectors there are fixed.
The ratio of the two characters for a simultaneous translation
is thus constant in the integral and has absolute value one.
The \(G_{i+1}\)-factor belongs to the simultaneous invariance
already proved. At \(i=n\) use the separate formula for
\(U^{(\nu)}\) and the \(H\)-invariance. This proves the modulus
assertion.
\end{proof}


\section{The weighted unipotent integral on a congruence ball}

\begin{definition}\label{def:finite-layer-ball}
Continue at a single finite place and put
\(\mathcal W_j(a,b')=\prod_{i=1}^j A_i(a,b')^s\) with
\(s\) from Definition~\ref{def:finite-averaged-pairings}.
For an integer \(H_0\ge0\) let
\(\Omega_{H_0}=\{x\in\mathcal P_T:h(x)\le H_0\}\).
It is open, since \(\mathcal P_T\) is a union of open cosets
of \(K_g\) and the entry conditions are open. Let
\(\chi_j(a,b')\) indicate that \(a^{-1}u b'\in\Omega_{H_0}\)
for some \(u\in U^{(j)}\). This is Borel since the defining
condition is a union of open sets. For \(i\le\nu\) define
\[
             v_{g,i}=m_T(g)^{d_i}|g|_v
                         =|U_T\cap K_g\cap N_i|_{\eta_i}.
\]
The equality and \(\prod_i v_{g,i}=v_g\) are the layer calculation
of Lemma~\ref{lem:finite-bessel-radical-volume}.
Use \(n_i(t,c)\in N_i\) for the element with row
\(-t^*J\) and corner \(-\omega Q(t)+c\) on \(V_i\);
it sends \(e_i^*\) to \(e_i^*+t+(-\omega Q(t)+c)e_i\).
Its measure is \(\dd t\,\dd c\) on \(V_{i+1}\times iF_v\),
since the row change preserves the integral lattices.
Here the \(i\) in \(iF_v\) denotes the skew unit \(\sqrt{-1}\).
\end{definition}

\begin{lemma}\label{lem:finite-congruence-difference}
If \(h(x),h(y)\le H_0\), then the entries of \(xy^{-1}-1\)
belong to \(g\pi^{-2H_0}B\). Conjugation by \(K_T\) preserves
this matrix condition. In \(n_i(t,c)\) it implies
\(t\in g\pi^{-2H_0}\mathcal L_{i+1}\) and restricts \(c\),
at each fixed \(t\), to a set of volume at most
\(|g|_v Q^{2H_0}\).
\end{lemma}
\begin{proof}
The relation \(y^{-1}=J_0^{-1}y^*J_0\) and integrality of the
Gram matrices give the same entry bound for \(y^{-1}-1\) as for
\(y-1\). Now \(xy^{-1}-1=(x-y)y^{-1}\); the entries of \(x-y\)
lie in \(g\pi^{-H_0}B\), those of \(y^{-1}\) in \(\pi^{-H_0}B\)
because \(g\) is integral. Multiplication on both sides by
integral invertible matrices preserves the ideal condition.
Read the last column of \(n_i(t,c)-1\). Besides the claimed
condition on \(t\), the corner belongs to \(g\pi^{-2H_0}B\).
Two admissible \(c\)'s at fixed \(t\) differ by an element of
the intersection with the skew line, \(g\pi^{-2H_0}iA\) by
Lemma~\ref{lem:skew-covol}. This group has exactly the stated
additive volume.
\end{proof}

\begin{theorem}\label{thm:finite-positive-radical-weight}
For \(0\le j\le\nu,\ a\in\mathcal N_j,\ b'\in G_T\), one has
\[
 \int_{U^{(j)}}{\bf1}_{\Omega_{H_0}}(a^{-1}u b')
                    \mathcal W_j(a,u b')\dd u
 \ \le\ \chi_j(a,b')\, C_v^j Q^{60n j H_0}
                                \prod_{i=1}^j v_{g,i}.
\]
The measure on the partial unipotent group assigns integral
unipotents mass one. The threshold on the arithmetic scale is
independent of the residue cardinality \(Q\).
\end{theorem}
\begin{proof}
If there is no admissible \(u\), the integral is zero. Simultaneous
left multiplication of \(a,b'\) by a member of \(P_j\), and
conjugation of \(u\) by that member, preserve the integral:
\(P_j\) normalizes \(U^{(j)}\), its inner isometry factors preserve
the row and skew measures by Theorem~\ref{thm:finite-floor-haar},
and the unipotent measure is bi-invariant as in
Lemma~\ref{lem:finite-bessel-radical-volume}. The weight is unchanged
by Lemma~\ref{lem:finite-flag-invariance} and the nesting of the
subgroups. We may thus take \(a\in K_T\). Replace \(b'\) by
\(u_*b'\) for one admissible \(u_*\), using right invariance in
the integral. We may now also take \(a^{-1}b'\in\Omega_{H_0}\).
For any contributing \(u\), conjugate the quotient of
\(a^{-1}u b'\) and \(a^{-1}b'\) by \(a\). By
Lemma~\ref{lem:finite-congruence-difference}, \(u-1\) has entries
in \(g\pi^{-2H_0}B\).

Induct on \(j\), the initial assertion being the bound for one
indicator. For \(j>0\) write \(u=u_0n_j(t,c)\), with the product
of the \(U^{(j-1)}\)-measure and \(\dd t\,\dd c\) by the split
projection in Lemma~\ref{lem:heisenberg-layer}. The induced middle
block after removing \(j-1\) pairs is \(n_j(t,c)\), so the bounds
on \(t,c\) in Lemma~\ref{lem:finite-congruence-difference} apply.
Let \(t_+\) be the component of \(t\) on the \(m_*\) standard
lines of Lemma~\ref{lem:finite-auxiliary-planes}.
Take the module \(W\) in Definition~\ref{def:finite-close-volume}
at \(x=a,\ y=a^{-1}b'\), with no row equations. The choice of
\(b'\) guarantees \(y\in\mathcal P_T\).
We claim
\(A_j(a,u b')\le\phi(t_+/g)\), with \(\phi\) given by
Theorem~\ref{thm:finite-projective-moment}. For
\(\lambda=[e_j,w]\in I_j(a,u b')\) choose a witness \(w\), and
for \(v\in W\) put \(z=Iv\). Then
\[
 g^{-1}\bigl([a^{-1}w,a^{-1}z]-[(u b')^{-1}w,(b')^{-1}z]\bigr)
              =g^{-1}[w,z-u z]=(\lambda t_+/g)^*v .
\]
Indeed \(n_j(t,c)z=z-[t,z]e_j\) and \(u_0\) acts trivially
on this vector modulo the earlier frame vectors, to which \(w\)
is orthogonal. Both argument pairs on the left are close by the
definitions of \(I_j,W\). The scalar-volume assertion of that
theorem implies the claim.

At fixed \(t,c\), use this bound for the last weight factor
and the induction hypothesis at \(a,n_j(t,c)b'\) for the other
\(j-1\) factors, dropping its indicator \(\chi_{j-1}\) from
the upper bound. This use is valid since
\(\mathcal N_j\subseteq\mathcal N_{j-1}\). It remains to integrate
\(\phi(t_+/g)^s\) over the indicated range of \(t,c\).
Integrating the central variable costs \(|g|_v Q^{2H_0}\).
Put \(t=g t_0\), giving the factor \(m_T(g)^{d_j}\).
Now \(t_0\in\pi^{-2H_0}\mathcal L_{j+1}\). The
\(d_j-m_*\) other coordinates have volume
\(Q^{4H_0(d_j-m_*)}\), and the moment on the standard coordinates
is \(C_v Q^{4H_0(m_*-s)}\) by
Theorem~\ref{thm:finite-projective-moment}. All changes of order
are for nonnegative integrands by Lemma~\ref{lem:product-integration}.
The cost of this step is thus at most
\[
                 v_{g,j} C_v Q^{(2+4(d_j-s))H_0}.
\]
From \(d_j=k-1+2(n-j)\), \(a_*=3\nu+6\), \(m_*=k-1-2a_*\)
and \(s=m_*-3\nu-2n\), one gets \(d_j-s=13n+3-2j\).
Hence \(2+4(d_j-s)\le60n\) at the indicated scales.
Multiplication over the steps completes the induction.
\end{proof}


\section{A Fourier integral centered on a finite flag layer}

\begin{lemma}\label{lem:finite-layer-fourier-phase}
In Definition~\ref{def:finite-averaged-pairings} let
\(j_0<i\le\nu\) and \(a,b'\in G_T\). Put \(f=f_{i+1}\),
\(\delta=\delta^{[i+1]}_{Q(f)}\), where \(Q(f)\) here denotes the
hermitian norm. There is a bounded Borel amplitude \(F\) on the
quadric such that \(F(\sigma^{-1}f)=K_{i+1}(\sigma a,\sigma b')\)
for \(\theta_{i+1}\)-almost every \(\sigma\). It satisfies
\[
       K_i(a,n_i(t,c)b')=\int\psi_T^0(-[t,e]/g)F(e)\dd\delta(e).
\]
If \(a\in\mathcal N_i\) and \(q a\in K_T,\ q\in P_i\) has
middle block \(h\in G_{i+1}\) on the quotient after removing
\(i\) pairs, then \(F(e)=0\) almost everywhere outside
\(h^{-1}\mathcal L_{i+1}\). In particular \(F\in L^1\cap L^2(\delta)\).
For \(a\notin\mathcal N_i\) one may take \(F=0\) almost everywhere.
\end{lemma}
\begin{proof}
The stabilizer \(S_{i+1}\) is contained in \(P_{i+1}\cap G_{j_0+1}\).
By Theorem~\ref{thm:finite-pairing-laws} the indicated function
on \(G_{i+1}\) therefore factors through \(\sigma^{-1}f\).
Use the Borel sections of Lemma~\ref{lem:finite-lower-orbit-section}
and Proposition~\ref{prop:finite-top-orbit-section} and extend by
zero off their charts. In the recurrence for \(K_i\), conjugation
of \(n_i(t,c)\) by \(\sigma\in G_{i+1}\) gives \(n_i(\sigma t,c)\):
the corner is unchanged and the middle column transforms by
\(\sigma\). This unipotent commutes with \(d_*\) since \(i>j_0\).
Its only possibly nonzero simple-root coordinate is that at \(i\),
namely the \(e_i\)-entry of the image of \(f\), equal to
\(-[\sigma t,f]=-[t,\sigma^{-1}f]\). Use the separate translation
law of Theorem~\ref{thm:finite-pairing-laws} at \(i+1\) and push
forward by the orbit map. This gives the displayed formula.

If the amplitude is nonzero at \(\sigma^{-1}f\), then
\(\sigma a\in\mathcal N_{i+1}\) by that theorem, forcing
\(a\in\mathcal N_i\) as in
Proposition~\ref{prop:finite-average-mass}. Choose \(p'\in P_{i+1}\)
with \(p'\sigma a\in K_T\). Since \(q a\in K_T\), the matrix
\(p'\sigma q^{-1}\) lies in \(P_i\cap K_T\).
Its middle block after removing \(i\) pairs lies in \(K_{i+1}\)
and has the form \(s'\sigma h^{-1}\) with \(s'\in S_{i+1}\).
Taking its inverse on \(f\) gives \(h\sigma^{-1}f\in\mathcal L_{i+1}\).
The support assertion follows. The transformed lattice is compact,
and \(\delta\) is locally finite by
Proposition~\ref{prop:finite-floor-norm-disintegration}. Boundedness
from Theorem~\ref{thm:finite-pairing-laws} proves integrability
of the amplitude and its square.
\end{proof}

\begin{theorem}\label{thm:finite-centered-fourier}
For \(j_0<i\le\nu\) and every \(a,b'\in G_T\), with the indicators
of Definition~\ref{def:finite-layer-ball}, one has
\[
 \begin{split}
 &\int_{N_i}\chi_{i-1}(a,n b')|K_i(a,n b')|^2\dd\eta_i(n)\\
 &\quad\le \chi_i(a,b')\,v_{g,i}M_{i+1}Q^{40nH_0}
     \int_{G_{i+1}}|K_{i+1}(\sigma a,\sigma b')|^2
                                                   \dd\theta_{i+1}(\sigma).
 \end{split}
\]
\end{theorem}
\begin{proof}
The inner square integral is finite by
Lemma~\ref{lem:finite-layer-fourier-phase} and the orbit law.
The left side is zero if \(a\notin\mathcal N_i\), or if no
\(n_i(t_0,c_0)\) has \(\chi_{i-1}(a,n_i(t_0,c_0)b')=1\).
Otherwise fix such a pair, choose \(q a=k\in K_T\) with
\(q\in P_i\), and use its middle block \(h\) from that lemma.
Necessarily \(\chi_i(a,b')=1\) by the layer factorization
\(U^{(i)}=U^{(i-1)}N_i\).
For every admissible \(t,c\) there are \(u,u_0\in U^{(i-1)}\)
so that \(a^{-1}u n_i(t,c)b'\) and \(a^{-1}u_0n_i(t_0,c_0)b'\)
belong to \(\Omega_{H_0}\). Their quotient, conjugated by \(k\),
is
\[
       q u n_i(t,c)n_i(t_0,c_0)^{-1}u_0^{-1}q^{-1}.
\]
Its difference from the identity lies in the matrix ideal
\(g\pi^{-2H_0}B\) by Lemma~\ref{lem:finite-congruence-difference}.
Take its block on the quotient after the first \(i-1\) pairs.
There \(u,u_0\) disappear and \(q\) acts as an element of
\(N_i G_{i+1}\) with \(G_{i+1}\)-block \(h\).
In the product of two layer matrices their vector parameters
add and the corner entries add with one extra term determined
by their vector parameters, by multiplying the three-block
matrices of Lemma~\ref{lem:heisenberg-layer}. It follows that
the quotient block has vector parameter \(h(t-t_0)\), and its
corner depends on \(c\) additively with coefficient one at
fixed \(t\). Thus
\[
       h(t-t_0)\in g\pi^{-2H_0}\mathcal L_{i+1},\qquad
       |\{c:\chi_{i-1}(a,n_i(t,c)b')=1\}|\le |g|_v Q^{2H_0}.
\]
For the second bound use that differences of the admissible skew
parameters lie in \(g\pi^{-2H_0}iA\). This reasoning uses only
a fixed translate in estimating the domain.

The Fourier formula in Lemma~\ref{lem:finite-layer-fourier-phase}
is independent of \(c\). Integrate that parameter first. Change
\(t\) to \(t'=h(t-t_0)/g\), giving a factor \(m_T(g)^{d_i}\);
the isometry has determinant of norm one. In the amplitude
integral send \(e\) to \(e'=he\), preserving \(\delta\) by
Proposition~\ref{prop:finite-floor-norm-disintegration}. The phase
from \(t_0\) has modulus one and may be absorbed into the amplitude.
Now the input amplitude is supported on the integral quadric box
and the output \(t'\) ranges in \(\pi^{-2H_0}\mathcal L_{i+1}\).
Apply Theorem~\ref{thm:finite-flat-fourier} there with frequency
\(-1\), ambient output and quadric input. Its factor is at most
\(Q^{8\ell_{i+1}H_0} M_{i+1}\). Its input squared norm is the
square integral in the statement by the orbit law. Multiplying by
the central and vector Jacobians gives
\(v_{g,i}M_{i+1}Q^{(8\ell_{i+1}+2)H_0}\).
Finally \(\ell_{i+1}=4n+3-i\) by
Lemma~\ref{lem:finite-auxiliary-planes}, so \(8\ell_{i+1}+2\le40n\)
at sufficiently large scales. All integrations of squared
absolute values use Lemma~\ref{lem:product-integration}.
\end{proof}


\section{Squared inner products on the finite unipotent quotient}

\begin{definition}\label{def:finite-partial-quotients}
In the notation of Definition~\ref{def:finite-averaged-pairings},
put \(\tau_i=\theta_i*\cdots*\theta_1\) for \(1\le i\le n\)
and let \(\tau_0\) be unit mass at the identity. Put
\(\rho_l=\eta_{l+1}*\cdots*\eta_\nu*\mathcal M_{\rm up}\) for
\(0\le l\le\nu\), so \(\rho_\nu=\mathcal M_{\rm up}\).
Write \(\dd u_l\) for the measure on \(U^{(l)}\) of
Theorem~\ref{thm:finite-positive-radical-weight}.
For the diagonal \(d_*\) of that definition of the pairings, set
\[
                 D_*=\prod_{a=1}^{j_0}m_T(p_a/q_a)^{d_a+1}.
\]
\end{definition}

\begin{lemma}\label{lem:finite-partial-quotient-law}
One has \(\dd u_l*\rho_l=\mu\). Suppose \(h_1\in G_T\)
normalizes \(U^{(l)}\) and conjugation by \(h_1\) multiplies
volumes there by \(\Delta(h_1)>0\). For every nonnegative Borel
function \(F\) invariant under left \(U^{(l)}\)-multiplication,
\[
                 \int F(h_1 y)\dd\rho_l(y)
                         =\Delta(h_1)\int F(y)\dd\rho_l(y).
\]
For \(h_1\in G_{l+1}\) one has \(\Delta(h_1)=1\). At
\(l=\nu,\ h_1=d_*\) one has \(\Delta(h_1)=D_*\).
\end{lemma}
\begin{proof}
The successive semidirect group products from
Lemma~\ref{lem:finite-bessel-radical-volume} give
\(\dd u_l*\eta_{l+1}*\cdots*\eta_\nu=\dd u\).
Corollary~\ref{cor:finite-actual-quotient} gives the first
identity. A compact set \(C\) has finite \(\rho_l\)-measure:
restrict \(u\) to the compact mass-one subgroup \(U_0=U^{(l)}
\cap K_T\) and bound \(\rho_l(C)\) by \(\mu(U_0 C)\) using
that identity. Measures of the other products can throughout
be integrated in their defining coordinates by
Lemma~\ref{lem:product-integration}.

Choose one \(y_b\) in each double coset \(U^{(l)}y_bK_T\);
there are at most countably many since \(G_T/K_T\) is countable
by Lemma~\ref{lem:finite-unitary-level-haar}. On \(y_bK_T\)
put \(a_l=1/|U^{(l)}\cap y_bK_T y_b^{-1}|\), and put it equal
to zero outside their union. The denominator is a positive
finite volume since the intersection is compact open.
For \(y\in U^{(l)}y_bK_T\), the set of \(u\) taking \(y\) into
\(y_bK_T\) is a right coset of that intersection. Hence
\(\int a_l(uy)\dd u_l=1\) everywhere and \(a_l\) is Borel.
Applied to \(F a_l\), the convolution identity yields
\(\int F\dd\rho_l=\int F a_l\dd\mu\).
Apply this also to \(F(h_1\,\cdot)\), which is left invariant
under \(U^{(l)}\), and use left Haar invariance to get
\(\int F(z)a_l(h_1^{-1}z)\dd\mu(z)\).
Expanding the convolution once more, the inner cutoff integral
is \(\int a_l(h_1^{-1}u z)\dd u_l\); substituting
\(u=h_1v h_1^{-1}\) gives \(\Delta(h_1)\).
This proves the asserted formula.

The determinant-one action for \(h_1\in G_{l+1}\) follows on
each earlier layer from Theorem~\ref{thm:finite-floor-haar}.
For \(h_1=d_*\), conjugation on layer \(a\) sends its vector
and skew coordinate to
\((\bar\lambda_a d_*|_{V_{a+1}}t,\lambda_a\bar\lambda_a c)\),
by applying the matrix to \(e_a^*\) and using the corner
parametrization in Definition~\ref{def:finite-layer-ball}.
The middle isometry preserves vector volume; the scalar factors
give \(m_T(\lambda_a)^{d_a}\) on the row and
\(m_T(\lambda_a)\) on the skew line. Multiplying gives \(D_*\).
\end{proof}

\begin{proposition}\label{prop:finite-squared-pairing}
For the same group and level,
\[
 \iint|\langle W_Z(a),W_Z(b')\rangle|^2
            \dd\mathcal M_{\rm up}(a)\dd\mathcal M_{\rm up}(b')
                                    =m_Z\,v_Z/v_g .
\]
Consequently the last averaged pairing of
Theorem~\ref{thm:finite-pairing-laws} satisfies
\[
 \iint |K_n(a,b')|^2\dd\tau_n(a)\dd\rho_\nu(b')
                                     \le D_*m_Z\,v_Z/v_g .
\]
\end{proposition}
\begin{proof}
The first integrand is left \(U_T\)-invariant independently
in each variable, and vanishes off \((U_TZ)^2\).
By Corollary~\ref{cor:finite-actual-quotient} it is enough to
integrate on \(Z^2\) with measure
\(\dd\mu(a)\dd\mu(b')/v_Z^2\).
Let \(\bar Z=Z/K_g\) and let \(B_0\) be the image of \(U_T\cap Z\)
in \(\bar Z\). It has size \(b_0=v_Z/v_g\). The character on
it induced by \(\chi_g\) is well defined by
Lemma~\ref{lem:finite-bessel-radical-volume}.
At \(a\in Z\) the support of the vector in \(\bar Z\) is
the set \(B_0\bar a\): indeed
\(\phi_g(a c^{-1})\ne0\) exactly when \(\bar a\bar c^{-1}\in B_0\).
On this support the entries have modulus one and are the values
of that character at \(\bar a\bar c^{-1}\).
Two such supports are disjoint or equal. If they are equal, the
two vectors differ by a scalar of modulus one, because
\(\bar b'\bar a^{-1}\in B_0\) multiplies every character
argument on the left. Thus the squared inner product is \(b_0^2\)
on the coincident supports and zero otherwise. For independent
uniform elements the probability of coincidence is
\(b_0/|\bar Z|=b_0\mu(K_g)/\mu(Z)\). Uniformity of the residues
follows by equality of the measures of their cosets.
The resulting quantity is
\((\mu(Z)/v_Z)^2 b_0^3\mu(K_g)/\mu(Z)=m_Z v_Z/v_g\), as claimed.

For \(K_n(a,b')\), the \(H\)-integration has support in
\(\{h:h a\in U_TK_T\}\) of measure at most one by
Proposition~\ref{prop:finite-average-mass}. Cauchy--Schwarz bounds
its square by the \(H\)-integral of
\(|\langle W_Z(ha),W_Z(d_*h b')\rangle|^2\).
At fixed \(h,a\) the corresponding function of \(h b'\) is
left \(U_T\)-invariant, since \(d_*\) normalizes \(U_T\) and
its character has modulus one. Lemma~\ref{lem:finite-partial-quotient-law}
with \(h\in H\subset G_n\) permits the replacement of \(h b'\)
by \(b'\) under \(\rho_\nu\). The remaining measure in \(h a\)
is \(\mu_H*\tau_n=\mathcal M_{\rm up}\). Now apply the same
lemma with \(h_1=d_*\) in the second variable. The first equality,
with the resulting factor \(D_*\), gives the assertion.
All these changes of order use Lemma~\ref{lem:product-integration}
on nonnegative functions.
\end{proof}


\section{Second moments on finite congruence balls}

\begin{definition}\label{def:finite-second-moments}
With Definitions~\ref{def:finite-layer-ball}
and~\ref{def:finite-partial-quotients}, for \(j_0\le l\le\nu\) set
\[
 \begin{split}
 B_l&=\iint \chi_l(a,b')\,|K_{l+1}(a,b')|^2
                                  \dd\rho_l(b')\dd\tau_{l+1}(a),\\
 Z_l(x)&=\int\mathcal W_l(a,ax)|K_{l+1}(a,ax)|^2\dd\tau_{l+1}(a).
 \end{split}
\]
Here \(B_l\) uses the fixed integer \(H_0\) in the indicator,
and \(Z_l\) is independent of it. Both are nonnegative integrals.
\end{definition}

\begin{lemma}\label{lem:finite-second-moment-recurrence}
For \(j_0<i\le\nu\) one has
\[
   B_{i-1}\le v_{g,i}M_{i+1}Q^{40nH_0}B_i,\qquad
   B_\nu\le D_* m_Z v_Z/v_g.
\]
For every \(j_0\le l\le\nu\),
\[
  \int_{\Omega_{H_0}} Z_l(x)\dd\mu(x)
                  \le C_v^l Q^{60nlH_0}\left(\prod_{i=1}^l v_{g,i}\right) B_l .
\]
\end{lemma}
\begin{proof}
The terminal estimate follows from Proposition~\ref{prop:finite-squared-pairing}
after dropping \(\chi_\nu\). Write \(\rho_{i-1}=\eta_i*\rho_i\)
in \(B_{i-1}\) and apply Theorem~\ref{thm:finite-centered-fourier}.
For \(\sigma\in G_{i+1}\), normalizing \(U^{(i)}\) gives
\(\chi_i(\sigma a,\sigma b')=\chi_i(a,b')\).
At fixed \(\sigma,a\), the function
\[
       b''\longmapsto \chi_i(\sigma a,b'')
                                     |K_{i+1}(\sigma a,b'')|^2
\]
is invariant under left \(U^{(i)}\) because its indicator absorbs
the unipotent in the witness and the separate translation law
in Theorem~\ref{thm:finite-pairing-laws} has unit modulus.
Use Lemma~\ref{lem:finite-partial-quotient-law} with this
\(\sigma\) to replace \(\sigma b'\) by \(b''\) under \(\rho_i\).
The measure on \(\sigma a\) is \(\theta_{i+1}*\tau_i=\tau_{i+1}\).
This gives \(B_i\) times the asserted factor.

For the last inequality change \(x\) to \(y=ax\) at fixed \(a\).
Write \(\dd\mu(y)=\dd u_l*\rho_l\) by
Lemma~\ref{lem:finite-partial-quotient-law}. The modulus of
\(K_{l+1}(a,u b')\) is independent of \(u\in U^{(l)}\) by the
separate translation formula, also when \(l+1=n\).
A nonzero value forces \(a\in\mathcal N_{l+1}\subseteq\mathcal N_l\).
The integral over \(u\) of the remaining indicator and weight is
bounded by Theorem~\ref{thm:finite-positive-radical-weight},
including its factor \(\chi_l(a,b')\). This gives the claimed
multiple of \(B_l\). All integrals interchanged are nonnegative
by Lemma~\ref{lem:product-integration}.
\end{proof}

\begin{theorem}\label{thm:finite-translation-second-moment}
At all sufficiently large arithmetic scales, uniformly in the
finite place, intermediate group \(Z\), heads of fixed common
length \(j_0\le\nu\), and integer radius parameter \(H_0\ge0\),
\[
       \int_{\Omega_{H_0}}\sum_{l=j_0}^\nu Z_l(x)\dd\mu(x)
                 \le 2n\,m_Z\,v_Z\,D_*\,Q^{60n^2 H_0}.
\]
\end{theorem}
\begin{proof}
For a fixed \(l\), iterate Lemma~\ref{lem:finite-second-moment-recurrence}
from \(\nu\) down to \(l\) and insert the result into its last
inequality. The product of the level factors is exactly \(v_g\)
by Definition~\ref{def:finite-layer-ball}; it cancels the
denominator \(v_g\) in \(B_\nu\).
The remaining local density factor is
\[
              C_v^l\prod_{i=l+1}^{\nu}M_{i+1}
                    \le(1+2^{-n})(1+2^{-5n})\le2 .
\]
Here Theorem~\ref{thm:finite-projective-moment} and
Corollary~\ref{cor:finite-integral-fibre-product} apply to this
single place and all factors are at least one. The exponent of
\(Q^{H_0}\) is at most \(60nl+40n(\nu-l)\le60n^2\).
There are at most \(n\) indices. Summing the bounds proves the
theorem.
\end{proof}


\section{The finite coefficient away from the principal level}

\begin{lemma}\label{lem:finite-first-small-ideal}
Use the empty heads of Definition~\ref{def:finite-averaged-pairings}
and \(Z_l(x)\) of Definition~\ref{def:finite-second-moments}.
For \(x\in\mathcal P_T\setminus K_g\) let
\(e_x=Q^{-\max\{1,h(x)/(2\nu)\}}\). At sufficiently large scales,
\[
                |c_Z(x)|^2\le2 e_x^s\sum_{l=0}^{\nu-1} Z_l(x).
\]
For \(x\notin\mathcal P_T\), \(c_Z(x)=0\).
\end{lemma}
\begin{proof}
The integral for \(K_0(1,x)\) equals \(c_Z(x)\) by
Theorem~\ref{thm:finite-pairing-laws}.
If any raw inner product \(\langle W_Z(a),W_Z(ax)\rangle\)
is nonzero, Lemma~\ref{lem:finite-flag-invariance} gives
\(x\in\mathcal P_T\), proving the last assertion.
For \(x\in\mathcal P_T\setminus K_g\), apply Cauchy--Schwarz
to the recurrence at the \(i\)-th step, \(1\le i\le\nu\).
The integrand vanishes unless \(\sigma a\in\mathcal N_i\).
Using the reciprocal \(s\)-th power bound of
Proposition~\ref{prop:finite-average-mass} on that set gives
\[
 |K_{i-1}(a,ax)|^2\le M_i C_v
  \int A_i(\sigma a,\sigma ax)^s|K_i(\sigma a,\sigma ax)|^2
                                                      \dd\theta_i(\sigma).
\]
Start this inequality at \(a=1,\ i=1\). At each step split the
right integral where \(A_i\le e_x\) and where \(A_i>e_x\).
On the former subset bound \(A_i^s\) by \(e_x^s\), keeping the
squared pairing there. On the latter keep the weight and expand
the pairing by the inequality at the next step. All earlier
weights and comparisons with \(e_x\) are unchanged by that next
left multiplier, because \(G_{i+1}\subset P_b\) for \(b\le i\)
and Lemma~\ref{lem:finite-flag-invariance} applies. The resulting
contribution from the first small ideal occurring at \(i\) is
at most
\[
           \left(\prod_{b=1}^i M_b C_v\right)e_x^s Z_{i-1}(x):
\]
indeed the preceding weights give \(\mathcal W_{i-1}\), the
product of the first \(i\) measures gives \(\tau_i\), and the
restrictions may be dropped from this nonnegative integral.
This procedure uses the recurrence only finitely many times.

On the remaining set at \(i=\nu\), all \(\nu\) ideal volumes
exceed \(e_x\). This still holds after every left translation
by an element of \(G_n\). However \(T_\nu=\mu_H*\theta_n\) is
supported there, and Theorem~\ref{thm:finite-flag-displacement}
makes every raw pairing in \(K_\nu(a,ax)\) zero on this set.
Thus only the stopped contributions remain. Their products of
density factors are at most two by
Theorem~\ref{thm:finite-projective-moment} and
Corollary~\ref{cor:finite-integral-fibre-product}. Summing proves
the claim, including cases of infinite individual upper integrals.
\end{proof}

\begin{theorem}\label{thm:finite-one-place-remainder}
At every sufficiently large arithmetic scale, for each finite
place \(v\) of \(F\) of residue cardinality \(Q\) and every
intermediate group \(K_g\subseteq Z\subseteq K_{\{v\}}\),
\[
                 \mathcal R_{\{v\}}(Z)\le Q^{-k+120n^3}
                                             \le Q^{-9k/10}.
\]
In particular the full integral in
Definition~\ref{def:finite-normalized-energy} is finite.
\end{theorem}
\begin{proof}
By Proposition~\ref{prop:finite-congruence-energy} the remainder
equals \(m_Z^{-1}\int_{x\notin K_g}|c_Z(x)|^2\dd\mu(x)\).
At height \(h(x)=h\ge1\) use Lemma~\ref{lem:finite-first-small-ideal},
then enlarge the domain for the nonnegative \(Z_l\)'s to
\(\Omega_h\). Theorem~\ref{thm:finite-translation-second-moment}
applies with \(D_*=1\) and \(v_Z\le1\) by
Lemma~\ref{lem:finite-bessel-radical-volume}. The bound becomes
\[
       4n\sum_{h=1}^\infty
                        Q^{60n^2h-s\max\{1,h/(2\nu)\}} .
\]
Here \(s=k-11n-4\), by Lemma~\ref{lem:finite-auxiliary-planes}
and Definition~\ref{def:finite-averaged-pairings}.
For \(h\le2\nu\) every term is at most \(Q^{120n^2\nu-s}\).
Beyond \(2\nu\) the terms form a geometric progression with
ratio \(Q^{60n^2-s/(2\nu)}\le1/2\): indeed \(Q\ge2\),
\(n\le\sqrt X\) and \(k\ge X^5/2\) eventually by
Proposition~\ref{prop:thick}. Their sum is bounded by one
further copy of \(Q^{120n^2\nu-s}\). Countable additivity
therefore bounds the remainder by \(4n(2\nu+1)Q^{120n^2\nu-s}\).
Since \(\nu=n-1\) and
\(11n+4+\log_2(4n(2\nu+1))\le120n^2\) eventually, this is at
most \(Q^{-k+120n^3}\). Also \(120n^3\le k/10\) eventually
by the same scale bounds, proving the two estimates uniformly.
\end{proof}

\begin{theorem}\label{thm:finite-marked-bottom}
For every finite set \(T\) of places of \(F\),
\(\mathcal E_T(K_T)\le\mu(K_T)+2^{-n}\) at all sufficiently
large scales. For \(T,S\) and the actual vector stabilizer
\(Z\) of Theorem~\ref{thm:marked-energy-reduction},
\[
                  \mathcal E_T(Z)\le(1+b^{-6})b^{-6k|S|}.
\]
\end{theorem}
\begin{proof}
At \(v\) write \(K_v\) for its integral isometry group.
The lower bound in Theorem~\ref{thm:finite-floor-haar} gives
\(\mu_v(K_v)\ge (1-Q_v^{-1})^{2n}\ge2^{-2n}\).
Indeed the split integer algebra has unit volume
\((1-Q_v^{-1})^2\), and the inert one has unit volume
\(1-Q_v^{-2}\), by deleting the maximal ideals.
Consequently Theorem~\ref{thm:finite-one-place-remainder} gives
\[
 \sum_{v\in T}\frac{\mathcal R_{\{v\}}(K_v)}{\mu_v(K_v)}
       \le 2^{2n} f\sum_{h=2}^\infty h^{-9k/10}
       \le 2^{2n+1}f\,2^{-9k/10}.
\]
Here at most \(f\) places lie above each rational prime by
Lemma~\ref{lem:real-place-rings}, and we used the decreasing
integral estimate of Corollary~\ref{cor:finite-integral-fibre-product}.
The last bound is at most \(2^{-n-2}\) eventually because
\(f\le X^{12}\), \(n\le\sqrt X\) and \(k\ge X^5/2\).
By Proposition~\ref{prop:finite-local-energy-product} the ratio
\(\mathcal E_T(K_T)/\mu(K_T)\) is the product of one plus each
local ratio in this sum. Using \(1+t\le\exp(t)\) and
\(\exp(t)-1\le2t\) for \(0\le t\le1/2\), its excess over one
is at most \(2^{-n-1}\). Multiplying by
\(\mu(K_T)\le1+2^{-5n}\le2\) from Theorem~\ref{thm:finite-floor-haar}
gives the first assertion. In particular the unselected block
in Theorem~\ref{thm:marked-energy-reduction} has the remainder
bound required there. At each selected place \(Q_v=g_0\),
Theorem~\ref{thm:finite-one-place-remainder} gives
\(g_0^{-9k/10}\le g_0^{-3k/4}\) for its intermediate group.
That reduction proves the second assertion.
\end{proof}


\section{A simultaneous second moment over the finite places}

\begin{definition}\label{def:finite-product-pairings}
Let \(T\) be an arbitrary finite set of places of \(F\) and
\(Z=\prod_{v\in T} Z_v\) a product of intermediate groups, in
particular the group in Theorem~\ref{thm:finite-marked-bottom}.
Use all groups, orbit and layer measures, partial quotients,
averaged pairings \(K_i\), and flag ideals as in
Definitions~\ref{def:finite-flag-ideal},
\ref{def:finite-averaged-pairings} and~\ref{def:finite-partial-quotients},
now over the product algebras \(E_T,F_T\).
Use \(\mathcal P_T=\prod_v\mathcal P_{\{v\}}\) and
\(\mathcal N_i=P_iK_T\), also at \(i=n\).
For an integer tuple \(\mathbf H=(H_v\ge0)_{v\in T}\), put
\[
   R(\mathbf H)=\prod_{v\in T}Q_v^{H_v},\qquad
   \Omega_{\mathbf H}=\prod_{v\in T}\Omega_{H_v},\qquad
   R(x)=\prod_{v\in T}Q_v^{h(x_v)} .
\]
Here and below \(\mathcal W_j(a,b')=\prod_{i=1}^j A_i(a,b')^s\)
uses \(s=m_*-3\nu-2n\), and
\(A_i(a,b')=\prod_v A_{i,v}(a_v,b'_v)\) is the additive volume
of the product ideal. Put \(C_T=\prod_v C_v\) and use
\(D_*=\prod_{a=1}^{j_0}m_T(p_a/q_a)^{d_a+1}\).
Write \(Z_l(x)\) as in Definition~\ref{def:finite-second-moments}
with product measures. For \(j_0\le j\le\nu\) put
\[
 \begin{split}
 U_j(a,x)&=\int|K_{j+1}(h a,h ax)|\dd\theta_{j+1}(h),\\
 S_j(a,x)&=\sum_{l=j}^{\nu}\int
   \mathcal W_l(h a,h ax)|K_{l+1}(h a,h ax)|^2
                             \dd(\theta_{l+1}*\cdots*\theta_{j+1})(h).
 \end{split}
\]
An empty tuple of places carries the singleton probability on
every group and scalar set.
\end{definition}

\begin{proposition}\label{prop:finite-product-second-moment}
The orbit laws, recurrence, separate and simultaneous translation
laws, and support conclusions of
Theorem~\ref{thm:finite-pairing-laws} hold with these product
measures. One has
\[
 \int_{\Omega_{\mathbf H}}\sum_{l=j_0}^\nu Z_l(x)\dd\mu(x)
                       \le 2n\,m_Z v_Z D_*\,R(\mathbf H)^{60n^2}.
\]
In addition \(S_j(\,\cdot\,,x)\) is left \(P_j\)-invariant,
vanishes off \(\mathcal N_j\), and
\(\int S_j(a,x)\dd\tau_j(a)=\sum_{l=j}^\nu Z_l(x)\).
\end{proposition}
\begin{proof}
The orbit and unipotent measures are products by the coordinate
calculations of Proposition~\ref{prop:finite-local-energy-product}.
All the iterated group measures are products by ordered
multiplication in each component. Also \(D_Z=\prod_vD_{Z_v}\),
and the vector entry is the product of the respective entries
by the factorization of \(\phi_g\) in that proposition. Thus the
inner product, then its absolutely convergent successive averages,
factor. The local translation and support laws give their stated
product versions. Each \(Z_l\) factors separately by
Lemma~\ref{lem:product-integration} on nonnegative functions.
Apply Lemma~\ref{lem:finite-second-moment-recurrence} at each
place for one fixed \(l\), iterating from \(\nu\). Its level
factors cancel its terminal \(v_{g,v}^{-1}\).
Before taking any sum over layers the resulting bound is
\[
 m_{Z_v}v_{Z_v}D_{*,v}\ C_v^l\!
       \prod_{i=l+1}^\nu M_{\{v\}}(\ell_{i+1})
                           \ Q_v^{60n^2 H_v}.
\]
The normalizations \(m_Z,v_Z,D_*\) are the products of their
one-place quantities by their definitions and
Proposition~\ref{prop:finite-local-energy-product}.
After multiplication the two density factors have product at
most \(C_T^n\prod_i M_T(\ell_i)\le2\) at all sufficiently large
scales by Theorem~\ref{thm:finite-projective-moment} and
Corollary~\ref{cor:finite-integral-fibre-product}. Summing over
at most \(n\) indices proves the bound.

For the last assertions fix \(l\). The innermost test
\(\mathcal W_l(y,yx)|K_{l+1}(y,yx)|^2\) is \(P_{l+1}\)-invariant
since \(l+1>j_0\), by the modulus law and
Lemma~\ref{lem:finite-flag-invariance}. Its support lies in
\(\mathcal N_{l+1}\). The orbit integration in the proof of
Proposition~\ref{prop:finite-average-mass} descends the invariance
and support through one layer; its proof uses the orbit law and
applies verbatim to a product via its explicit product sections.
In detail a preceding unipotent conjugates into the invariance
subgroup and the remaining isometry acts on the invariant orbit
measure. A nonzero contribution at a layer forces membership in
the preceding near set as both the integrating group and the next
subgroup lie in the preceding subgroup. Iterating proves the
claims concerning \(S_j\), also for infinite values. Finally
\((\theta_{l+1}*\cdots*\theta_{j+1})*\tau_j=\tau_{l+1}\).
Nonnegative integration gives the identity for its integral.
\end{proof}

\begin{lemma}\label{lem:finite-product-stopping}
For \(a\in G_T,\ x\in\mathcal P_T,\ j_0\le j\le\nu\), one has
\[
                 \mathcal W_j(a,ax)U_j(a,x)^2\le 2 S_j(a,x).
\]
If \(e>0\) and for every \(h\in G_{j+1}\) with
\(\langle W_Z(h a),W_Z(d_*h ax)\rangle\ne0\) some
\(j<i\le\nu\) has \(A_i(h a,h ax)\le e\), then the right side
may be replaced by \(2e^s S_j(a,x)\).
\end{lemma}
\begin{proof}
For \(i\le\nu\), Proposition~\ref{prop:finite-average-mass} gives
\[
 \int_{\sigma a\in\mathcal N_i}A_i(\sigma a,\sigma ax)^{-s}
                         \dd\theta_i(\sigma)\le M_T(\ell_i)C_T
\]
by product integration. At \(i=n\) the integral of one there is at
most \(M_T(\ell_n)\). In \(U_{i-1}\), the integrand vanishes
off the indicated set by Proposition~\ref{prop:finite-product-second-moment}.
Weighted Cauchy--Schwarz, using \(A_i^{s/2}\) and its reciprocal
inside the integral of the absolute value, gives
\[
 U_{i-1}(a,x)^2\le M_T(\ell_i)C_T
  \int A_i(\sigma a,\sigma ax)^s
                |K_i(\sigma a,\sigma ax)|^2\dd\theta_i(\sigma)
 \qquad(i<n).
\]
At \(i=n\) use the same inequality with weight one and the
stated terminal mass. The recurrence gives
\(|K_i(a,ax)|\le U_i(a,x)\). Iterate the inequalities starting
at \(i=j+1\), multiplying by \(\mathcal W_j\). Later multipliers
preserve the earlier ideals because they belong to the corresponding
\(P_b\). The terminal contribution is at most the \(l=\nu\)
term of \(S_j\) times
\(C_T^\nu\prod_i M_T(\ell_i)\le2\), proving the first estimate.
For the second, stop as soon as a new \(A_i\) is at most \(e\),
bounding \(A_i^s\) by \(e^s\) and retaining \(|K_i|^2\).
Dropping the earlier restrictions bounds that stopped term by
the \(l=i-1\) summand of \(2e^s S_j\) just as in
Lemma~\ref{lem:finite-first-small-ideal}. If no such index is
reached, the raw upper factors, including \(H\), are all in
\(G_{j+1}\), and all previous comparisons persist under them.
The hypothesis makes each remaining raw integrand zero. This
argument also covers \(j=\nu\), when the hypothesis requires
all such integrands to vanish.
\end{proof}


\section{Reciprocal ideal volumes on the opposite finite flag}

\begin{definition}\label{def:finite-head-denominator}
For a list \(p\) of length \(j\le\nu\) with all
\(D_a(p)\in B_T\), use \(B(p),d(p)\) of
Definition~\ref{def:finite-head-bounds} and put
\[
     L(p)=m_T(B(p))^{-1},\qquad
     L_v(p)=Q_v^{h_v},\quad h_v=\operatorname{ord}_v d(p)_v .
\]
Then \(L(p)=\prod_v L_v(p)\ge1\) by the norm determinant formula
in Lemma~\ref{lem:middle-hyperbolic-chart}.
Put \(\mathfrak m_0\) equal to unit mass at one and
\(\mathfrak m_j=(\eta_j*\theta_j)*\cdots*(\eta_1*\theta_1)\).
\end{definition}

\begin{theorem}\label{thm:finite-inverse-radical}
For \(a\in\mathcal N_j,\ x\in\mathcal P_T\) and any \(p\) as in
Definition~\ref{def:finite-head-denominator}, one has
\[
 \int_{U^{(j)}}\frac{{\bf1}_{\mathcal N_j}(w_p u a)}
                       {\mathcal W_j(w_p u a,w_p u ax)}\dd u_j
 \le C_T^j L(p)^{60nj}
                    \prod_{b=1}^j m_T(p_b)^{-d_b-1/2}.
\]
If some \(D_b(p)\) fails integrality, the integration domain
in the numerator is empty.
\end{theorem}
\begin{proof}
It suffices to prove the result at one place: all domains,
weights and measures factor before integration by
Proposition~\ref{prop:finite-product-second-moment}.
Replacing \(a\) by an element of \(U^{(j)}a\) leaves the integral
unchanged by right invariance in \(u\). Replacing it by \(h a\),
\(h\in G_{j+1}\), conjugates \(u\) and commutes \(h\) past \(w_p\),
which preserves the integral by Lemma~\ref{lem:finite-flag-invariance}
and the determinant-one conjugation in
Theorem~\ref{thm:finite-floor-haar}. Thus take \(a\in K_T\)
locally. On the domain there is \(q'\in P_j\) with
\(q'w_pu a\in K_T\). Factoring \(q'=u_1 l\), \(l\in G_{j+1}\),
gives \(u_1 w_p l u\in K_T\); thus
Theorem~\ref{thm:finite-head-minors} proves all the leading minors
integral and bounds \(u\in\mathcal B(d(p))\). This proves also
the empty-domain assertion.

Induct on \(j\), the empty case having bound one. Write
\(u=n_j(t,c)u_0\), \(u_0\in U^{(j-1)}\), in this order. It
gives the product of their measures since conjugation by \(N_j\)
preserves the earlier layers by Theorem~\ref{thm:finite-floor-haar}.
The middle block after removing \(j-1\) pairs is \(n_j(t,c)\).
Writing \(h_v=\operatorname{ord}_v d(p)\), we therefore have
\(t\in\pi^{-h_v}\mathcal L_{j+1}\), and at fixed \(t\) the
admissible \(c\)'s have volume at most \(Q^{h_v}\) by the corner
calculation of Lemma~\ref{lem:finite-congruence-difference}.
The valuations of \(p_j=D_j(p)/D_{j-1}(p)\) in each field
component are between \(-h_v\) and \(h_v\): all the initial
minors are integral and their product and its conjugate give
\(d(p)\). Consequently \(m_T(p_j)^{1/2}\le Q^{h_v}\).

At fixed \(u_0\), its action on \(Iv_0\), for \(v_0\in E_v^{m_*}\),
adds a combination of \(e_1,\ldots,e_{j-1}\). Take the
coefficients to be the rows of \(C v_0\) in
Definition~\ref{def:finite-close-volume}, with first position
\(a\) and translation \(x\). For \(v_0\) in its module \(W\)
we have \(u_0Iv_0=Iv_0\). The vector \(w=w_p u Iv_0\) has
earlier dual coordinates zero and \(j\)-th dual coordinate
\(-p_j t_+^*v_0\): the layer inserts \(-[t,Iv_0] e_j\)
and the \(j\)-th flip sends \(e_j\) to \(p_j e_j^*\).
The two inverse images of \(w\) under \(w_p u a,w_p u ax\)
are close by the choice of \(W\). Hence
\[
 A_j(w_p u a,w_p u ax)^{-1}\le \phi(-\bar p_j t_+)
\]
by Theorem~\ref{thm:finite-projective-moment}.
Use \(t'=-\bar p_j t\), which gives volume factor
\(m_T(p_j)^{-d_j}\) and puts \(t'\) in
\(\pi^{-2h_v}\mathcal L_{j+1}\).
The moment formula of that theorem on the standard coordinates,
and the box volumes for the other \(d_j-m_*\) coordinates, bound
the integral of \(\phi(t'_+)^s\) by \(C_vQ^{4(d_j-s)h_v}\).
For the center use \(Q^{h_v}\le m_T(p_j)^{-1/2}Q^{2h_v}\).
Since \(2+4(d_j-s)\le60n\) as computed in
Theorem~\ref{thm:finite-positive-radical-weight}, the cost of
this step is at most \(m_T(p_j)^{-d_j-1/2} C_v Q^{60nh_v}\).
The bound is uniform in \(u_0\); all original integrands are
Borel by Lemma~\ref{lem:finite-flag-invariance}.

If \(p'\) is the prefix of length \(j-1\), then
\(w_pu a=w_j(p_j)n_j(t,c)w_{p'}u_0a\).
The first two factors belong to \(G_j\), hence preserve the
earlier weights and may be removed from the necessary condition
\(w_{p'}u_0a\in\mathcal N_{j-1}\).
Apply the induction hypothesis to that prefix and use
\(h_v(p')\le h_v(p)\), which follows since \(D_j(p)\) is integral.
Nonnegative integration now gives the asserted one-place bound,
and multiplication gives the claimed product bound.
\end{proof}

\begin{lemma}\label{lem:finite-preceding-translation}
For \(0\le j\le\nu\),
\(\mathfrak m_j=\dd u_j*\tau_j\) and
\(\mu=\mathcal M_{j+1}*\mathfrak m_j\).
Consequently if a nonnegative Borel function \(A\) is left
\(G_{j+1}\)-invariant and \(w\) commutes with \(G_{j+1}\), then
\[
                 \int A(wy)\dd\mathfrak m_j(y)
                              =\int A(y)\dd\mathfrak m_j(y).
\]
\end{lemma}
\begin{proof}
The second convolution identity follows by expanding
Definition~\ref{def:finite-floor-haar} down to layer \(j\).
For the first identity use induction on \(j\): in
\((\eta_j*\theta_j)*\dd u_{j-1}*\tau_{j-1}\), the first
two factors lie in \(G_j\) and their conjugations preserve
\(\dd u_{j-1}\) by Theorem~\ref{thm:finite-floor-haar}.
Move it to the front; \(\dd u_{j-1}*\eta_j=\dd u_j\), and
\(\theta_j*\tau_{j-1}=\tau_j\).
To obtain the last assertion use a cutoff \(a(y)\) with
\(\int_{G_{j+1}}a(hy)\dd\mathcal M_{j+1}(h)=1\) for every \(y\).
Construct it on one \(y_bK_T\) per double coset
\(G_{j+1}y_bK_T\), with value the reciprocal volume of
\(G_{j+1}\cap y_bK_Ty_b^{-1}\), and set it to zero elsewhere.
The proof of Lemma~\ref{lem:finite-partial-quotient-law} applies
since \(\mathcal M_{j+1}\) is bi-invariant, finite and positive
on compact open subgroups by Theorem~\ref{thm:finite-floor-haar}.
Insert \(a\) using the second convolution identity and the
left invariance of \(A\). Replacing \(y\) by \(wy\) uses left
invariance of \(\mu\); the inner cutoff integral is still one
by commutation with \(w\). This proves the formula also for
infinite integrals.
\end{proof}


\section{Two finite flags with prescribed initial pivots}

\begin{definition}\label{def:finite-pivot-forms}
For \(p,q\) of length \(j\le\nu\) put \(i=j+1\), and extend
\(H(p)=\prod_{a=1}^j m_T(p_a)^{d_a/2}\) to these lengths.
Use \(K_b,S_j,Z_l\) of Definition~\ref{def:finite-product-pairings}
with empty initial lists, and write \(K'_b,S'_j,Z'_l\) for the
pairings and sums with \(d_*=w_q^{-1}w_p\). At fixed \(a,x\) let
\[
 F(\sigma)=K_i(\sigma a,\sigma ax),\qquad
 F'(t)=K'_i(t w_p^{-1}a,t w_p^{-1}ax).
\]
Define
\[
 \begin{split}
 B_j(a,x)&=\iint|F(\sigma)F'(t)|\dd\theta_i(\sigma)\dd\theta_i(t),\\
 B_j^\xi(a,x)&=\iint\psi_T^0\bigl(\xi[\sigma^{-1}f_i,t^{-1}f_i]\bigr)
       F(\sigma)\overline{F'(t)}\dd\theta_i(\sigma)\dd\theta_i(t),\\
 \mathcal I_j^\xi(p,q)&=\frac{H(p)H(q)}{m_Z}
             \iint B_j^\xi(a,x)\dd\mathfrak m_j(a)\dd\mu(x)
 \end{split}
\]
for \(\xi\in E_T\), the last expression whenever absolutely
convergent. The two inner integrals exist at every \(a,x\) by
the finite support-mass bounds in
Proposition~\ref{prop:finite-average-mass} and the product
pairing laws of Proposition~\ref{prop:finite-product-second-moment}.
\end{definition}

\begin{lemma}\label{lem:finite-partial-flag-displacement}
If \(B_j(a,x)>0\) then \(x\in\mathcal P_T\),
\[
    a,\ ax,\ w_p^{-1}a,\ w_q^{-1}ax\in\mathcal N_j,
       \qquad D_b(p),D_b(q)\in B_T\quad(1\le b\le j).
\]
Let \(L(p),L(q)\) be as in Definition~\ref{def:finite-head-denominator}
and put \(A=m_T(g)^{-1}L(p)^2L(q)^2\). If \(R(x)>A^{4n}\), then
for every \(h\in G_{j+1}\) with
\(\langle W_Z(ha),W_Z(hax)\rangle\ne0\) there is some \(j<i'\le\nu\)
such that \(A_{i'}(ha,hax)\le R(x)^{-1/(4n)}\).
\end{lemma}
\begin{proof}
Positivity in Definition~\ref{def:finite-pivot-forms} yields
a product of upper multipliers in \(G_{j+1}\) making the raw
untwisted pairing nonzero. Lemma~\ref{lem:finite-flag-invariance}
at each place gives \(x\in\mathcal P_T\), and the two points
then lie in \(U_TZ\subseteq\mathcal N_j\).
The upper multipliers can be removed since \(G_{j+1}\subseteq P_j\).
The same support argument for the twisted factor gives
\(w_p^{-1}a,w_q^{-1}ax\in\mathcal N_j\), since \(d_*\)
commutes with \(G_{j+1}\) and \(d_*w_p^{-1}=w_q^{-1}\).
Corollary~\ref{cor:finite-simultaneous-representative} gives the
integrality assertion.

Replace \(a\) by \(ha\) with its raw pairing nonzero.
All four near conditions persist by the commutation and subgroup
inclusion just used. Choose \(u,v_1\in U_T\) with
\(u a=k_1,\ v_1 ax=k_2\in Z,\ k_1^{-1}k_2\in K_g\), now at
these replaced arguments, by Lemma~\ref{lem:finite-flag-invariance}
at each place. Write \(u=u_{\rm out}u_{\rm in}\) with
\(u_{\rm out}\in U^{(j)},\ u_{\rm in}\in G_{j+1}\), by splitting
the first \(j\) layers. Since \(w_p^{-1}a\in P_jK_T\), write
\(a=w_p l u' k'\) with \(l\in G_{j+1},u'\in U^{(j)}, k'\in K_T\).
Thus \(u_{\rm out} w_p(u_{\rm in}l)u'=k_1(k')^{-1}\) is integral,
so \(u_{\rm out}\in\mathcal B(d(p))\) by
Theorem~\ref{thm:finite-head-minors}. Since \(u_{\rm in}\) fixes
the first \(j\) dual vectors, \(d(p)u e_b^*\in\mathcal L\) for
\(b\le j\). Using \(ax,q,v_1\) in the identical factorization
gives \(d(q)v_1e_b^*\in\mathcal L\).

Put \(\alpha=g d(p)d(q)\), so \(m_T(\alpha)^{-1}=A\).
In the determinant argument of Theorem~\ref{thm:finite-flag-displacement}
choose \(w_b=\alpha e_b^*\) for \(b\le j\).
Their two images \(u w_b,v_1w_b\) both lie in \(g\mathcal L\),
their dual entries before position \(b\) vanish, and their \(b\)-th
entries equal \(\alpha\). For \(b>j\) choose witnesses for generators
\(t_b\) of \(I_b(a,ax)\) at every component as in that theorem.
The two matrices of columns are thus integral, congruent modulo
\(g\), and their lower triangular blocks have the same determinant
\(\Delta=\alpha^j\prod_{b>j}t_b\). The adjugate and block
calculation of that theorem, which used precisely these three
facts about the columns, gives at each \(v\)
\(h(x_v)\le2\operatorname{ord}_v(\Delta\bar\Delta)\).
Multiply \(Q_v^{h(x_v)}\) over \(v\) to obtain
\[
         R(x)\le A^{2j}\left(\prod_{b>j} A_b(a,ax)\right)^{-2}.
\]
In the stated range \(A^{2j}\le R(x)^{1/2}\), whence the product
of the remaining volumes is at most \(R(x)^{-1/4}\).
For \(j<\nu\), at least one factor is at most
\(R(x)^{-1/(4(\nu-j))}\le R(x)^{-1/(4n)}\).
For \(j=\nu\) the empty product makes such a raw pairing impossible.
This proves the assertion at the originally replaced arguments.
\end{proof}

\begin{proposition}\label{prop:finite-crossed-weights}
Fix \(x\in\mathcal P_T\) and two integral-minor heads \(p,q\) as
in Lemma~\ref{lem:finite-partial-flag-displacement}. Let \(f(a)\)
be a nonnegative Borel function supported where
\(a,w_p^{-1}a\in\mathcal N_j\). Suppose on its support
\[
 f(a)^2\le c^2\,
  \frac{S_j(a,x)}{\mathcal W_j(a,ax)}
  \frac{S'_j(w_p^{-1}a,x)}{\mathcal W_j(w_p^{-1}a,w_p^{-1}ax)}
\]
with \(c>0\). Write \(Z=\sum_{l=j}^\nu Z_l(x)\) and
\(Z'=\sum_{l=j}^\nu Z'_l(x)\) in this statement only.
If these are finite, then
\[
 \int f\dd\mathfrak m_j
 \le c C_T^j L(p)^{60nj}
         \prod_{b=1}^j m_T(p_b)^{-d_b-1/2}\sqrt{ZZ'} .
\]
\end{proposition}
\begin{proof}
Use Cauchy--Schwarz, pairing \(S_j(a,x)\) with the reciprocal
weight and near indicator at \(w_p^{-1}a\), and \(S'_j(w_p^{-1}a,x)\)
with those at \(a\). In the first integral use
\(\mathfrak m_j=\dd u_j*\tau_j\) of
Lemma~\ref{lem:finite-preceding-translation}. The \(S_j\)-factor
is left \(P_j\)-invariant and supported on \(\mathcal N_j\) by
Proposition~\ref{prop:finite-product-second-moment}. Apply
Theorem~\ref{thm:finite-inverse-radical} in \(u_j\) with head
\(\bar p\), since \(w_p^{-1}=w_{\bar p}\) by the pair matrices.
Its factor is the displayed product times \(C_T^j L(p)^{60nj}\);
the integral of \(S_j\) against \(\tau_j\) is \(Z\).
The other integral has the form \(\int A_1(w_p^{-1}a)\dd\mathfrak m_j(a)\)
with
\[
       A_1(y)=S'_j(y,x)
         \frac{{\bf1}_{\mathcal N_j}(w_p y)}
              {\mathcal W_j(w_p y,w_p yx)} .
\]
This function is left \(G_{j+1}\)-invariant because that group
commutes with \(w_p\) and preserves each factor under simultaneous
translation. Lemma~\ref{lem:finite-preceding-translation} removes
\(w_p^{-1}\) from its argument. The same unipotent integration
now applies with head \(p\), giving the same factor times \(Z'\).
Taking square roots of the product proves the assertion.
\end{proof}


\section{Absolute convergence on finite height shells}

\begin{definition}\label{def:finite-translation-shell}
For a nonnegative integer tuple \(\mathbf h=(h_v)_{v\in T}\) put
\(X_{\mathbf h}=\{x\in\mathcal P_T:h(x_v)=h_v\ (v\in T)\}\),
\(R=R(\mathbf h)\) and \(P=60n^2\). For heads having integral
initial minors write \(A=m_T(g)^{-1}L(p)^2L(q)^2\), and put
\[
 \kappa(R)=
 \begin{cases}
  1,&R\le A^{4n},\\
  R^{-s/(8n)},&R>A^{4n}.
 \end{cases}
 \qquad
 C_j=\frac{H(p)H(q)}{m_Z}.
\]
All measures and pairings in this section are those of
Definition~\ref{def:finite-pivot-forms}. For \(T=\varnothing\)
there is just one shell and \(R=1\).
\end{definition}

\begin{proposition}\label{prop:finite-shell-majorant}
For integral-minor heads of length \(0\le j\le\nu\) one has
\[
 C_j\int_{X_{\mathbf h}}\int B_j(a,x)\dd\mathfrak m_j(a)\dd\mu(x)
 \le 8n L(p)^{60nj}L(q)^{1/2}R^P\kappa(R).
\]
At all sufficiently large arithmetic scales the integral defining
\(\mathcal I_j^\xi\) is absolutely convergent for all \(p,q,\xi\).
If either list has a nonintegral initial minor this integral is zero.
\end{proposition}
\begin{proof}
For fixed \(x\), \(B_j\) is the product of \(U_j(a,x)\) for the
empty head and \(U'_j(w_p^{-1}a,x)\) for the prescribed heads.
By Lemma~\ref{lem:finite-partial-flag-displacement} it vanishes
off the indicated near sets, and is identically zero for a
nonintegral initial minor. In the remaining case apply
Lemma~\ref{lem:finite-product-stopping} to both factors.
For \(R>A^{4n}\) the displacement lemma allows the stronger bound
for the first squared factor with \(e=R^{-1/(4n)}\). Taking square
roots gives the hypothesis of Proposition~\ref{prop:finite-crossed-weights}
with \(c=2\kappa(R)\). Its \(Z,Z'\) are finite almost everywhere
on \(X_{\mathbf h}\subseteq\Omega_{\mathbf h}\), by
Proposition~\ref{prop:finite-product-second-moment}. That proposition
and Cauchy--Schwarz in \(x\) bound the integral of \(\sqrt{ZZ'}\)
on this shell by \(2n\,m_Zv_Z D_*^{1/2}R^P\).
The head factors simplify exactly:
\[
 H(p)H(q)D_*^{1/2}\prod_{b\le j}m_T(p_b)^{-d_b-1/2}
       = m_T(D_j(q))^{-1/2}\le L(q)^{1/2}.
\]
The final inequality uses integrality of every initial minor,
with \(D_0=1\). The facts \(C_T^j\le2\) and \(v_Z\le1\)
follow from Theorem~\ref{thm:finite-projective-moment} and the
integral unipotent normalization in
Proposition~\ref{prop:finite-bessel-unfold}. This proves the estimate.

By the place count of Lemma~\ref{lem:real-place-rings},
\[
 \sum_{\mathbf h}R(\mathbf h)^{-2}
 =\prod_{v\in T}(1-Q_v^{-2})^{-1}\le\exp(2f).
\]
Indeed \(-\log(1-y)\le2y\) for \(0\le y\le1/2\) by integration,
and \(\sum_{m\ge2}m^{-2}\le\int_1^\infty t^{-2}\dd t=1\).
Since \(k=r+1-2n\), Proposition~\ref{prop:thick} gives
\(s/(8n)\ge P+2\) eventually. For \(R\le A^{4n}\) one has
\(R^{P+2}\le A^{4n(P+2)}\); for \(R>A^{4n}\),
\(R^P\kappa(R)\le R^{-2}\). Summing the shell estimate
therefore bounds the integral of the full absolute majorant.
Lemma~\ref{lem:finite-partial-flag-displacement} covers its
complement and \(|B_j^\xi|\le B_j\), as required.
\end{proof}

\begin{lemma}\label{lem:finite-two-orbit-boxes}
Fix \(a,x\) with \(B_j(a,x)>0\) and write \(i=j+1,\ d=d(p)\).
There is \(h\in G_i\) such that the two functions \(F(\sigma)\)
and \(F'(t)\) in Definition~\ref{def:finite-pivot-forms} factor
through \(v=h\sigma^{-1}f_i,\ v'=h t^{-1}f_i\), respectively.
As functions on the two quadrics they are bounded and vanish
almost everywhere outside \(d^{-1}\mathcal L_i\) and
\(d^{-2}\mathcal L_i\), respectively. On the product of these
boxes, \([v,v']\in d^{-3}B_T\).
\end{lemma}
\begin{proof}
The near conditions of Lemma~\ref{lem:finite-partial-flag-displacement}
allow Corollary~\ref{cor:finite-simultaneous-representative} to give
one \(h\in G_i\) with \(ha\in\mathcal B(d)\) and
\(hw_p^{-1}a\in\mathcal B(d^2)\). As in
Lemma~\ref{lem:finite-layer-fourier-phase}, the stabilizer of \(f_i\)
in \(G_i\) is contained in \(P_i\cap G_{j+1}\). The translation
laws of Proposition~\ref{prop:finite-product-second-moment} give
bounded Borel amplitudes through the orbit sections for both
lists. The simultaneous factor \(h\) preserves the quadric
measure by Proposition~\ref{prop:finite-floor-norm-disintegration}.

If \(F(\sigma)\ne0\), choose \(u'\in P_i\) with
\(u'\sigma a=k'\in K_T\) by the same support law. Thus
\(ha(k')^{-1}f_i=h\sigma^{-1}(u')^{-1}f_i\) lies in
\(d^{-1}\mathcal L\). The difference \((u')^{-1}f_i-f_i\)
belongs to the span of \(e_1,\ldots,e_j\): on the quotient
by these preceding frame vectors \(P_i\) fixes \(f_i\), by its
block definition, also for \(f_n=z,\ P_n=U_TH\).
Since \(h,\sigma\in G_i\) preserve that span and fix its
vectors, projection onto the middle summand gives
\(v\in d^{-1}\mathcal L_i\). Apply this argument with first
position \(w_p^{-1}a\) to obtain the other bound. The integral
basis in Lemma~\ref{lem:finite-auxiliary-planes} has integral
hermitian pairings, and \(d=\bar d\). This proves the last claim.
\end{proof}

\begin{proposition}\label{prop:finite-common-shell}
For the heads in Proposition~\ref{prop:finite-shell-majorant} set
\(d=d(p)\). If \(\xi\in E_T,\ \eta\in E_T^\times\) and
\(\xi-\eta\in d^3B_T\), then
\[
 \begin{split}
 C_j\int_{X_{\mathbf h}}\int |B_j^\xi|\dd\mathfrak m_j\dd\mu
 &\le 8n L(p)^{60nj+30n}L(q)^{1/2}R^P\\
 &\hspace{13mm}{}\cdot\min\{\kappa(R),m_T(\eta)^{-m_*/2}\}.
 \end{split}
\]
\end{proposition}
\begin{proof}
Where \(B_j>0\), Lemma~\ref{lem:finite-two-orbit-boxes} and
Proposition~\ref{prop:additive} allow replacement of \(\xi\)
by \(\eta\) in the phase since the argument changes by an
integer at each component. Isometry preserves the pairing
in the two new orbit coordinates. Apply Cauchy--Schwarz against
the first amplitude on its support and apply
Theorem~\ref{thm:finite-flat-fourier} with the second amplitude
conjugated, input \(d^{-2}\mathcal L_i\) and output
\(d^{-1}\mathcal L_i\), both with quadric measure. The bound for
\(|B_j^\xi|\) is the product of the two amplitude norms in
\(L^2(\delta_i)\) times
\[
 m_T(\eta)^{-m_*/2} M_T(\ell_i)L(p)^{3(4\ell_i-2)/2}.
\]
Here \(\Delta_T(d^{-1})=L(p)\);
\(M_T(\ell_i)\le2\) by Corollary~\ref{cor:finite-integral-fibre-product}
and \(3(4\ell_i-2)/2\le30n\) eventually by
Lemma~\ref{lem:finite-auxiliary-planes}. The square of each
amplitude norm times its corresponding \(\mathcal W_j\) is the
\(l=j\) summand of \(S_j\) or \(S'_j\), using the orbit measure
identity and Definition~\ref{def:finite-product-pairings}.
The hypotheses of Proposition~\ref{prop:finite-crossed-weights}
therefore hold with \(c=2L(p)^{30n}m_T(\eta)^{-m_*/2}\).
Integrate this bound on the shell using the same second moments
and exact head cancellation as in Proposition~\ref{prop:finite-shell-majorant}.
That proposition also bounds \(|B_j^\xi|\) using \(\kappa(R)\).
Since \(L(p)\ge1\), taking the smaller estimate gives the result.
\end{proof}

\section{Denominators in the common-pivot frequency bound}

\begin{definition}\label{def:finite-explicit-band}
Write \(K=k-n^7\) and let \(w\mid T\) run over the field components
of \(E_T\). For \(\xi\in E_T^\times\) and \(p\in(E_T^\times)^j\) put
\[
 \begin{split}
 D_T(\xi)&=\prod_{w\mid T}q_w^{\max(0,-\operatorname{ord}_w\xi_w)},\\
 \beta_T(\xi)&=b^{n^9 f}\, m_T(\xi)^{-K/2}D_T(\xi)^{-6},\\
 \Lambda_{T,j}(p)&=
    \prod_{a=1}^j {\bf1}_{B_T}(D_a(p))\,m_T(D_a(p))^{-n^6/2}.
 \end{split}
\]
Here the empty head has weight one, and \(D_T\) agrees with the
finite denominator in Theorem~\ref{thm:denominator-frequency}.
The identity \(m_T(\xi)=\prod_{w\mid T}q_w^{-\operatorname{ord}_w\xi_w}\)
follows from the determinant of multiplication on the local
integer rings in Lemma~\ref{lem:middle-hyperbolic-chart}.
\end{definition}

\begin{lemma}\label{lem:finite-shell-interpolation}
At all sufficiently large scales, for any \(u>0\),
\[
 \sum_{\mathbf h} R^P\min\{\kappa(R),u^{-m_*/2}\}
 \le e^{2f}A^{4n(P+2)}
                  \min(1,u^{-1/2})^{K+12}.
\]
Moreover, for any integral-minor heads, \(0\le e\le160n^2\),
\(0\le e'\le1\), one has
\[
 \begin{split}
 8n e^{2f} L(p)^{e+36}L(q)^{e'} A^{4n(P+2)+6}
 \ \le\ b^{11000n^3 f}\bigl(L(p)L(q)\bigr)^{500n^3}.
 \end{split}
\]
\end{lemma}
\begin{proof}
For \(m_*=k-6n-7,\ s=k-11n-4\), put
\(\theta=1-(K+12)/m_*\). Proposition~\ref{prop:thick}
implies \(K>0,\ m_*>0,\ s/m_*\ge1/2\) eventually, since
\(r\ge X^5,\ n\le\sqrt X,\ k=r+1-2n\) at scale \(X\).
Also \(m_*-(K+12)=n^7-6n-19\ge n^7/2\) eventually.
Hence \(0<\theta<1\) and
\(\theta s/(8n)\ge n^6/32\ge P+2\).
For positive \(a,b'\) their minimum is at most
\(a^\theta(b')^{1-\theta}\), by comparison with each factor.
Since \(\kappa\le1\), use \(a=\kappa,\ b'=\min(1,u^{-m_*/2})\).
On \(R>A^{4n}\) one has \(R^P\kappa^\theta\le R^{-2}\);
on the remaining shells it is at most \(A^{4n(P+2)}R^{-2}\).
The sum of \(R^{-2}\) was bounded in the proof of
Proposition~\ref{prop:finite-shell-majorant}. This proves the first claim.

By Theorem~\ref{thm:data} one has \(g=g_0g_1\le b^{22}\).
The local lattice indices give \(m_T(g)^{-1}\le g^{2f}\):
the complete product of the local indices is \(g^{2f}\),
the determinant of multiplication by \(g\) on the rank \(2f\)
integer lattice, by Lemma~\ref{lem:class-covol}. All omitted
factors are at least one. Thus \(A\le b^{44f}L(p)^2 L(q)^2\).
With \(P=60n^2\), both \(160n^2+36+2(4n(P+2)+6)\) and
\(1+2(4n(P+2)+6)\) are at most \(500n^3\) eventually.
Also \(44(4n(P+2)+6)+n\le11000n^3\) eventually.
Finally \(8ne^{2f}\le b^{nf}\) for large \(n,b\), as
\(f\ge1\). These comparisons prove the second bound.
\end{proof}

\begin{theorem}\label{thm:finite-common-frequency}
For \(0\le j\le\nu,\ \xi\in E_T^\times\) and every pair
of heads of length \(j\), one has, eventually in the arithmetic scale,
\[
 C_j\iint |B_j^\xi(a,x)|\dd\mathfrak m_j(a)\dd\mu(x)
               \le\beta_T(\xi)\Lambda_{T,j}(p)\Lambda_{T,j}(q).
\]
For heads with integral initial minors the right side can be
replaced by the following quantity:
\[
 b^{11000n^3 f}\bigl(L(p)L(q)\bigr)^{500n^3}
                         m_T(\xi)^{-K/2}D_T(\xi)^{-6}.
\]
\end{theorem}
\begin{proof}
Proposition~\ref{prop:finite-shell-majorant} gives zero outside
the asserted integral-minor case. With \(d=d(p)\) integral and
nonzero on every component, choose
\[
 \eta_w=\begin{cases}
  \xi_w,&\operatorname{ord}_w\xi_w<\operatorname{ord}_w d^3,\\
  d^3,&\operatorname{ord}_w\xi_w\ge\operatorname{ord}_w d^3 .
 \end{cases}
\]
Then \(\eta\) is a unit of \(E_T\) and \(\eta-\xi\in d^3B_T\).
The minimum of the two valuations gives \(m_T(\eta)\ge m_T(\xi)\).
For each component of negative valuation of \(\xi\),
its factor in \(m_T(\eta)\) is exactly the denominator factor
of \(\xi\) there. At every other component it is at least
\(q_w^{-\operatorname{ord}_w d^3}\), and these latter local
lower bounds are at most one at every component. Hence
\[
     m_T(\eta)\ge D_T(\xi)m_T(d^3)=D_T(\xi)L(p)^{-6},
\]
where \(d=B(p)\overline{B(p)}\). In particular
\[
 \min(1,m_T(\eta)^{-1/2})^{K+12}
  \le m_T(\xi)^{-K/2}D_T(\xi)^{-6}L(p)^{36}.
\]
Proposition~\ref{prop:finite-common-shell} and
Lemma~\ref{lem:finite-shell-interpolation} now yield the improved
bound: the exponent \(60nj+30n\) is at most \(160n^2\),
and \(A\ge1\). Finally \(11000n^3\le n^9\) and
\(500n^3\le n^6/2\) for all sufficiently large \(n\).
Definition~\ref{def:finite-explicit-band} and
\(\Lambda_{T,j}(p)=L(p)^{n^6/2}\) on the surviving heads give
the stated inequality.
\end{proof}

\section{The finite unipotent integral at unequal last entries}

\begin{lemma}\label{lem:finite-unequal-pivot-phase}
Let \(1\le j\le\nu\), let \(p,q\) have length \(j\), and suppose
\(\xi=g^{-1}(p_j^{-1}-q_j^{-1})\in E_T^\times\). For all
\(a,x\in G_T\) there is a bounded Borel amplitude \(F_a\) on the
quadric of \(f_{j+1}\), with measure \(\delta\), given by
\(F_a(\sigma^{-1}f_{j+1})=K'_{j+1}(\sigma a,\sigma ax)\)
for \(\theta_{j+1}\)-almost every \(\sigma\). One has
\[
 K'_j(n_j(t,c)a,n_j(t,c)ax)=
     \int\psi_T^0\bigl(p_j\xi[t,v]\bigr)F_a(v)\dd\delta(v).
\]
If \(a\in\mathcal N_j\) and \(A_0a\in K_T,\ A_0\in P_j\)
has middle quotient \(h\in G_{j+1}\), the amplitude is supported
almost everywhere on \(h^{-1}\mathcal L_{j+1}\). If
\(a\notin\mathcal N_j\) it may be taken identically zero.
\end{lemma}
\begin{proof}
The stabilizer of \(f_{j+1}\) in \(G_{j+1}\) preserves the
twisted pairing by Proposition~\ref{prop:finite-product-second-moment}.
The orbit sections and their measure formulas in
Lemma~\ref{lem:finite-lower-orbit-section} and
Proposition~\ref{prop:finite-top-orbit-section} give the stated
amplitude. If it is nonzero at a representative \(\sigma\), choose
\(A_1\in P_{j+1}\) with \(A_1\sigma a\in K_T\) by the pairing
support law. Then \(a\in\mathcal N_j\). On the middle quotient,
the integral matrix \(A_1\sigma A_0^{-1}\in P_j\cap K_T\) is
\(s_1\sigma h^{-1}\) with \(s_1\) fixing \(f_{j+1}\). Its
inverse preserves \(\mathcal L_{j+1}\), so
\(h\sigma^{-1}f_{j+1}\in\mathcal L_{j+1}\). Local finiteness
of \(\delta\) and boundedness of the pairing give integrability
and square integrability.

Use the pairing recurrence at level \(j+1\). Conjugation of
\(n_j(t,c)\) by \(\sigma\) gives \(u=n_j(\sigma t,c)\), as in
Lemma~\ref{lem:finite-layer-fourier-phase}. Its only possible
simple-root entry is \(-[\sigma t,f_{j+1}]\). Conjugation of \(u\)
by \(d_*=w_q^{-1}w_p\) multiplies that entry by \(p_j/q_j\),
since \(d_*\) is the identity on the middle quotient. The two
translation characters at level \(j+1\) therefore give the phase
with argument
\[
       g^{-1}(1-p_j/q_j)[\sigma t,f_{j+1}]
                             =p_j\xi[t,\sigma^{-1}f_{j+1}].
\]
The assertion follows by the orbit measure formula, also for the
last layer with \(f_{j+1}=z\).
\end{proof}

\begin{proposition}\label{prop:finite-unequal-layer}
Under the hypotheses of Lemma~\ref{lem:finite-unequal-pivot-phase},
assume \(p\) has integral initial minors and put \(L_p=L(p)\).
For all \(a,x\) one has
\[
 \begin{split}
 &\int_{N_j}{\bf1}_{\mathcal N_j}(w_p u a)
                   |K'_j(u a,u ax)|^2\dd\eta_j(u)\\
 &\quad\le 2 L_p^{50n}m_T(p_j)^{-d_j-1/2}m_T(\xi)^{-m_*}
       \int |K'_{j+1}(\sigma a,\sigma ax)|^2\dd\theta_{j+1}(\sigma).
 \end{split}
\]
\end{proposition}
\begin{proof}
If the left integral is positive, there exists \(u_0\in N_j\)
with \(w_pu_0a,u_0a\in\mathcal N_j\). Replace \(a\) by \(u_0a\).
The left integral is unchanged by right invariance of \(\eta_j\)
from Theorem~\ref{thm:finite-floor-haar}. On the right,
\(\sigma u_0\sigma^{-1}\in N_j\subseteq U^{(j)}\) contributes
only factors of modulus one by Proposition~\ref{prop:finite-product-second-moment}.
Thus assume \(w_p a,a\in\mathcal N_j\). Choose \(A_0,h\)
as in Lemma~\ref{lem:finite-unequal-pivot-phase} and put \(d=d(p)\).

For \(u=n_j(t,c)\) with \(w_p u a\in\mathcal N_j\), choose
\(A_1\in P_j\) with \(A_1w_p u a\in K_T\). Multiplication gives
\[
 (A_0u^{-1})w_p^{-1}A_1^{-1}=(A_0a)(A_1w_pua)^{-1}\in K_T .
\]
Write \(A_0u^{-1}=u_1 h,\ A_1^{-1}=h_1u_2\)
with the unipotents in \(U^{(j)}\) and \(h_1\in G_{j+1}\).
Both middle factors commute with \(w_p^{-1}=w_{\bar p}\).
By Theorem~\ref{thm:finite-head-minors}, \(u_1\in\mathcal B(d)\).
Hence \(A_0u^{-1}e_j^*=u_1e_j^*\in d^{-1}\mathcal L\).
This holds also at \(u=1\). Their difference projects onto
\(-ht\) in \(V_{j+1}\), since
\(u^{-1}e_j^*-e_j^*\) equals \(-t\) plus a multiple of \(e_j\)
by Lemma~\ref{lem:heisenberg-layer}. Thus \(ht\in d^{-1}\mathcal L_{j+1}\).
At fixed \(t\), the difference between the images for \(c,c'\)
is \(-(c-c')A_0e_j\), whose \(e_j\) coordinate is \(-(c-c')\).
Consequently every nonempty admissible central fibre lies in one
translate of the integral skew lattice scaled by \(d^{-1}\).
Its measure is at most \(\Delta_T(d^{-1})=L_p\), by the skew
coordinate normalization of Lemma~\ref{lem:finite-congruence-difference}
in each field place.

The phase formula is independent of \(c\). Integrate the center
and send \(t,v\) to \(ht,hv\). This preserves ambient measure
because the determinant of an isometry has norm of absolute value
one. It preserves \(\delta\) by
Proposition~\ref{prop:finite-floor-norm-disintegration}.
The input is supported on the integral quadric box and the
output lies in the ambient box \(d^{-1}\mathcal L_{j+1}\).
Theorem~\ref{thm:finite-flat-fourier} gives a squared-norm factor
at most
\[
 M_T(\ell_{j+1})L_p^{1+4\ell_{j+1}}\,
                m_T(p_j)^{-m_*}m_T(\xi)^{-m_*}
\]
including the central measure. Its squared input norm is the
last integral of the assertion. At each \(v\), every field
valuation of \(p_j\) has absolute value at most
\(\operatorname{ord}_v d_v\) by the minor calculation in
Theorem~\ref{thm:finite-inverse-radical}. There are two degrees
over \(F_v\), giving \(m_T(p_j)\le L_p^2\).
Also \(d_j+1/2-m_*>0\), and
\(1+4\ell_{j+1}+2(d_j+1/2-m_*)\le50n\) eventually by
Lemma~\ref{lem:finite-auxiliary-planes} and
Definition~\ref{def:finite-floor-sections}. Use this inequality
to extract \(m_T(p_j)^{-d_j-1/2}\) and apply
Corollary~\ref{cor:finite-integral-fibre-product} to bound
\(M_T(\ell_{j+1})\) by two. This proves the result.
\end{proof}

\section{Uniform decay for unequal finite pivots}

\begin{proposition}\label{prop:finite-unequal-shell}
Use the heads and frequency of Lemma~\ref{lem:finite-unequal-pivot-phase}
and assume the initial minors of both heads are integral. For
\(1\le j\le\nu\) and each shell of Definition~\ref{def:finite-translation-shell},
\[
 \begin{split}
 C_j\int_{X_{\mathbf h}}\int |B_j^0(a,x)|\dd\mathfrak m_j(a)\dd\mu(x)
 &\le 8n L(p)^{60nj+25n}L(q)^{1/2}R^P\\[-2pt]
 &\hspace{10mm}{}\cdot\min\{\kappa(R),m_T(\xi)^{-m_*/2}\}.
 \end{split}
\]
\end{proposition}
\begin{proof}
Use \(w=w_p\), and let \(w_0\) be the product of the flips for the
first \(j-1\) entries of \(p\). At fixed \(x\in X_{\mathbf h}\)
abbreviate
\[
 W_i(y)=\mathcal W_i(y,yx),\quad
 I_i(y)={\bf1}_{\mathcal N_i}(y)/W_i(y),\qquad
 H_p=\prod_{b=1}^j m_T(p_b)^{-d_b-1/2}.
\]
The weights are finite and positive by Lemma~\ref{lem:finite-flag-invariance}.
By the pairing recurrence, \(B_j^0(a,x)\) equals
\(K_j(a,ax)\overline{K'_j(w^{-1}a,w^{-1}ax)}\). Its nonzero
set requires \(a,w^{-1}a\in\mathcal N_j\). Use Cauchy--Schwarz
in \(a\) with squared factors \(A(a) I_j(w^{-1}a)\)
and \(\Phi(w^{-1}a)\), where
\[
 \begin{split}
 A(y)&=W_{j-1}(y)|K_j(y,yx)|^2,\\
 \Phi(y)&={\bf1}_{\mathcal N_j}(wy)W_j(y) I_{j-1}(w_0y)
                                           |K'_j(y,yx)|^2.
 \end{split}
\]
Indeed their product equals \(|B_j^0|^2\) on that set:
\(w_0w^{-1}=w_j(p_j)^{-1}\in G_j\) preserves \(W_{j-1}\)
and \(\mathcal N_{j-1}\) by Lemma~\ref{lem:finite-flag-invariance}.

For the first squared factor, \(A\) is \(P_j\)-invariant and
supported on \(\mathcal N_j\) by the empty-head pairing laws of
Proposition~\ref{prop:finite-product-second-moment}. Also
\(\int A\dd\tau_j=Z_{j-1}(x)\) for the empty head, by
Definition~\ref{def:finite-second-moments}. Factor
\(\mathfrak m_j=\dd u_j*\tau_j\) using
Lemma~\ref{lem:finite-preceding-translation}. Then
Theorem~\ref{thm:finite-inverse-radical} for \(\bar p\) bounds
the first squared integral by \(C_T^j L(p)^{60nj}H_p Z_{j-1}(x)\).

The function \(\Phi\) is \(G_{j+1}\)-invariant. Indeed this group
commutes with both Weyl matrices, preserves the earlier ideals
and the near sets, and preserves the simultaneous twisted
pairing at \(j\) by Proposition~\ref{prop:finite-product-second-moment}.
Lemma~\ref{lem:finite-preceding-translation} therefore replaces
\(\Phi(w^{-1}a)\) by \(\Phi(a)\) in its integral. In the resulting
integral disintegrate \(\mathfrak m_j=\eta_j*\dd u_{j-1}*\tau_j\):
this follows from the same lemma since conjugation by \(N_j\)
preserves the preceding radical measure, as in
Theorem~\ref{thm:finite-floor-haar}. At \(a=u y,\ u\in N_j\),
both \(W_j(a)\) and \(I_{j-1}(w_0 a)\) are constant in \(u\);
here \(u\) commutes with \(w_0\). Apply
Proposition~\ref{prop:finite-unequal-layer} to the other two factors.
After multiplication by \(W_j(y)\), the squared amplitude
integral is the \(l=j\) summand of \(S'_j(y,x)\). Enlarge it
to \(S'_j(y,x)\), which vanishes off \(\mathcal N_j\) and is
invariant under the preceding radical. Integrate \(I_{j-1}(w_0y)\)
over that radical by Theorem~\ref{thm:finite-inverse-radical} with
the prefix of \(p\). Its \(L\) is at most \(L(p)\).
Proposition~\ref{prop:finite-product-second-moment} evaluates the
integral of \(S'_j\) against \(\tau_j\) as
\(Z'(x)=\sum_{l=j}^\nu Z'_l(x)\). The second squared integral is
thus at most
\[
   2C_T^{j-1} L(p)^{50n+60n(j-1)}
                 H_p\,m_T(\xi)^{-m_*}Z'(x).
\]
Both moments are finite almost everywhere on this shell by the
same proposition. Take square roots and integrate in \(x\).
The integral of \(\sqrt{Z_{j-1}Z'}\) is at most
\(2n\,m_Zv_ZD_*^{1/2}R^P\), by that proposition and
Cauchy--Schwarz. The exact head cancellation in
Proposition~\ref{prop:finite-shell-majorant}, \(C_T^j\le2\)
and \(v_Z\le1\) give the bound involving \(m_T(\xi)\).
The other bound follows from \(|B_j^0|\le B_j\) and that
shell majorant. Taking the minimum proves the assertion.
\end{proof}

\begin{theorem}\label{thm:finite-unequal-frequency}
For every \(1\le j\le\nu\), heads \(p,q\) of length \(j\) with
\(\xi=g^{-1}(p_j^{-1}-q_j^{-1})\in E_T^\times\), and all
sufficiently large arithmetic scales, one has
\[
  C_j\iint |B_j^0(a,x)|\dd\mathfrak m_j(a)\dd\mu(x)
                      \le\beta_T(\xi)\Lambda_{T,j}(p)\Lambda_{T,j}(q).
\]
On heads with integral initial minors one may use the improved
bound of Theorem~\ref{thm:finite-common-frequency}.
\end{theorem}
\begin{proof}
By Proposition~\ref{prop:finite-shell-majorant} only the
integral-minor case contributes. The valuation estimate for
\(p_j\) in Theorem~\ref{thm:finite-inverse-radical} implies that
\(d(p)/p_j\in B_T\), and likewise \(d(q)/q_j\in B_T\).
Hence \(\alpha\xi\in B_T\) for \(\alpha=g d(p)d(q)\in B_T\).
At each field component the denominator exponent of \(\xi\)
is therefore at most \(\operatorname{ord}_w\alpha\). Thus
\(D_T(\xi)\le m_T(\alpha)^{-1}=A\).
Summing Proposition~\ref{prop:finite-unequal-shell} with
Lemma~\ref{lem:finite-shell-interpolation}, use
\[
 \min(1,m_T(\xi)^{-1/2})^{K+12}
      \le m_T(\xi)^{-K/2}A^6D_T(\xi)^{-6}.
\]
Here \(K\ge0,\ A/D_T(\xi)\ge1\); the same inequality holds
whether \(m_T(\xi)\) exceeds one or lies below it.
Now \(60nj+25n\le160n^2\), so the second assertion of
Lemma~\ref{lem:finite-shell-interpolation} proves the improved
bound. Its comparison to \(\beta_T\Lambda_{T,j}(p)\Lambda_{T,j}(q)\)
is that of Theorem~\ref{thm:finite-common-frequency}.
\end{proof}

\section{The finite-place kernels in the orbital sieve}

\begin{lemma}\label{lem:finite-lower-descent}
For \(0\le j\le\nu\) set \(J_j(p,q)=\mathcal I_j^0(p,q)\).
For \(0\le j<\nu\), and heads \(p,q\) of length \(j\), the
function \(c\mapsto J_{j+1}((p,c),(q,c))\) is integrable on
\(E_T^\times\) for additive measure. For every \(\xi\in E_T\),
\[
   \int_{E_T^\times}\psi_T^0(\xi c)J_{j+1}((p,c),(q,c))\dd c
                                      =\mathcal I_j^\xi(p,q).
\]
\end{lemma}
\begin{proof}
Put \(i=j+1\). In the inner integral in \(t\) of
\(B_j^\xi(a,x)\), both the phase and the second pairing depend
on \(t\) only through \(t^{-1}f_i\), by the stabilizer argument
of Lemma~\ref{lem:finite-two-orbit-boxes}, also before applying
the isometry used there. The orbit law of
Lemma~\ref{lem:finite-lower-orbit-section} is invariant under
\(G_i\) by Proposition~\ref{prop:finite-floor-norm-disintegration}.
Thus for fixed \(\sigma\) one may replace \(t\) by \(t\sigma\).
The phase argument becomes \(\xi[f_i,t^{-1}f_i]\), by isometry.

Disintegrate \(t=w_i(c)^{-1}u\) according to
Definition~\ref{def:finite-floor-sections}, with measure
\(m_T(c)^{d_i}\dd c\,\dd\eta_i(u)\).
Then \([f_i,t^{-1}f_i]=[u f_i,w_i(c)f_i]=c\).
The first pairing \(K_i(\sigma a,\sigma ax)\) is unchanged by
simultaneously multiplying its arguments by \(u\in P_i\), by
Proposition~\ref{prop:finite-product-second-moment} for the empty
head. Put \(y=u\sigma a\). The product of pairings is now
\[
 K_i(y,yx)\,
   \overline{K'_i\bigl((w_p w_i(c))^{-1}y,(w_p w_i(c))^{-1}yx\bigr)} .
\]
Indeed \(w_p,w_q\) commute with \(G_i\) and the ratio of the
extended Weyl matrices for \((p,c),(q,c)\) is still \(d_*\).
This product equals the zero-frequency \(B_i^0\) for the
extended lists by the pairing recurrence. The measure in \(y\)
is \(\eta_i*\theta_i*\mathfrak m_j=\mathfrak m_i\), and the
normalization is \(C_jm_T(c)^{d_i}=H(p,c)H(q,c)/m_Z\).

Every change just made also applies to the modulus before either
orbit integration. Its original integral is finite by
Proposition~\ref{prop:finite-shell-majorant}. The transformed
product is therefore absolutely integrable jointly in \(c,y,x\).
The parameterized functions are Borel since their defining
matrix maps are continuous and their integrands and fixed
measures are Borel, by Lemma~\ref{lem:product-integration}.
That lemma permits integration first in \(x,y\), giving the
claimed identity and integrability in \(c\).
\end{proof}

\begin{lemma}\label{lem:finite-last-descent}
For \(j=\nu\) use \(J_n\) of Theorem~\ref{thm:finite-top-gram}
with any compact open additive subgroup \(L\subset E_T\).
For \(\xi\in E_T\) and all heads \(p,q\) of length \(\nu\),
\[
  \int_{E_T}\psi_T^0(\xi c)J_n((p,c),(q,c))\dd c
                     ={\bf1}_{L^\perp}(\xi)\mathcal I_\nu^\xi(p,q).
\]
\end{lemma}
\begin{proof}
Start with the Fourier identity of Theorem~\ref{thm:finite-top-gram}.
Expand \(\langle\mathcal J(p,l),\mathcal J(q,l)\rangle\) by
Definition~\ref{def:finite-top-correlation}, using integration
variables \(y,y'\in G_T\). The trace of the product of the two
rank-one matrices, with the first adjointed, is
\[
  \langle W_Z(y),W_Z(y')\rangle\,
       \overline{\langle W_Z(lw_p^{-1}y),W_Z(lw_q^{-1}y')\rangle}.
\]
In this expansion, a nonzero term forces \(y,l\) into the compact
box for \(p\), and \(y',l\) into the compact box for \(q\), by
Corollary~\ref{cor:finite-simultaneous-representative}. A failing
minor makes every term zero. In the remaining case the integrand
is bounded by Proposition~\ref{prop:finite-bessel-unfold} and
the product of the compact boxes has finite measure. All changes
below therefore apply to the absolute integral as well.

Set \(y'=yx\) using left invariance of \(\mu\), and write \(y=b a\)
with \(\dd\mu=\mathcal M_n*\mathfrak m_\nu\) from
Lemma~\ref{lem:finite-preceding-translation}. Change \(l\) to
\(t=l b\) by right invariance of \(\mathcal M_n\) in
Theorem~\ref{thm:finite-floor-haar}. Both \(b\) and \(t\) lie in
\(G_n\), commuting with \(w_p,w_q\). The phase is
\(\psi_T^0(\xi[b^{-1}z,t^{-1}z])\) since \(l=tb^{-1}\).
The two pairings now have arguments
\[
    (b a,b ax),\qquad (t w_p^{-1}a,t w_q^{-1}ax).
\]
Independently in \(b,t\), use
\(\mathcal M_n=\mu_H*\theta_n\) of
Definition~\ref{def:finite-floor-haar}. The left \(H\)-factors
fix \(z\) in the phase, and the respective averages of the two
pairings are \(K_n(\sigma a,\sigma ax)\) and
\(K'_n(t' w_p^{-1}a,t'w_p^{-1}ax)\), with
\(\sigma,t'\) now distributed by \(\theta_n\). Indeed
\(d_*w_p^{-1}=w_q^{-1}\) and \(d_*\) commutes with \(G_n\).
Thus the normalized integral is exactly
\({\bf1}_{L^\perp}(\xi)\mathcal I_\nu^\xi(p,q)\) by
Definition~\ref{def:finite-pivot-forms}. Integrability in \(c\)
was proved in Theorem~\ref{thm:finite-top-gram}.
\end{proof}

\begin{theorem}\label{thm:actual-finite-sieve-family}
At all sufficiently large arithmetic scales take \(T,S,Z\) of
Theorem~\ref{thm:marked-energy-reduction}. For every compact open
additive subgroup \(L\subset E_T\), use \(J_0,\ldots,J_\nu\)
of Lemma~\ref{lem:finite-lower-descent} and the kernel \(J_n\)
of Theorem~\ref{thm:finite-top-gram}, \(n=\nu+1\).
They satisfy Definition~\ref{def:orbital-tower} over \(C=E_T\)
with length parameter \(\nu\), additive character \(\psi_T^0\),
\(I(z)=g^{-1}z^{-1}\), bounds \(\beta_T,\Lambda_{T,j}\) of
Definition~\ref{def:finite-explicit-band}, and
\[
 E_0=(1+b^{-6})b^{-6k|S|},\qquad
 J_0=\frac1{m_Z}\int_{G_T}|c_Z(x)|^2\dd\mu(x).
\]
The top common-entry quadratic form is the middle-orbit integral
of Corollary~\ref{cor:finite-top-orbit-energy}.
\end{theorem}
\begin{proof}
By the recurrence and the empty lists,
\(B_0^0(1,x)=|K_0(1,x)|^2=|c_Z(x)|^2\), using
Proposition~\ref{prop:finite-product-second-moment}.
The preceding measure is unit mass at one. Hence the formula
for \(J_0\) follows from Definition~\ref{def:finite-pivot-forms}.
It equals \(\mathcal E_T(Z)\) by
Proposition~\ref{prop:finite-congruence-energy}, giving the bottom
bound by Theorem~\ref{thm:finite-marked-bottom}.
Lemmas~\ref{lem:finite-lower-descent} and~\ref{lem:finite-last-descent}
give absolute integrability and the required values at zero
frequency, since \(0\in L^\perp\). Replacing zero components by
one changes the integrand only on an additive null set by
Theorem~\ref{thm:finite-top-gram}. At invertible frequencies
Theorem~\ref{thm:finite-common-frequency} bounds the Fourier
integrals at every level, including the last, since the factor
\({\bf1}_{L^\perp}\) has modulus at most one. Theorem~\ref{thm:finite-unequal-frequency}
bounds \(J_j=\mathcal I_j^0\) at each required inverse
difference. Theorem~\ref{thm:finite-top-gram} supplies the
continuity and positive semidefiniteness of the top kernel.
Its common-entry formula gives the asserted orbit integral.
\end{proof}

\section{The rational pivot sum for the constructed kernels}

\begin{lemma}\label{lem:product-local-sieve}
Let \(C=\prod_{\alpha\in I}C_\alpha\) be a finite product of the
local scalar algebras in Definition~\ref{def:orbital-tower},
with product additive measure and character, and componentwise
maps \(u,I(\,\cdot\,)\) as in that definition. At each factor let
\(J_{j,\alpha}\) satisfy the three kernel conditions with length
\(\nu\) and bounds \(E_{0,\alpha},\beta_\alpha,\Lambda_{\alpha,j}\).
Then the kernels
\[
           J_j(p,q)=\prod_{\alpha\in I}
                              J_{j,\alpha}(p_\alpha,q_\alpha)
\]
satisfy those conditions on \(C\), with the products of the
corresponding bounds.
\end{lemma}
\begin{proof}
At fixed heads the absolute value of each common-entry function
factors, and every factor is integrable. Integration of the
nonnegative product gives the product of its finite integrals
by Lemma~\ref{lem:product-integration}. The same lemma evaluates
the complex Fourier integral as the product of its local
Fourier integrals. At zero this is the preceding product kernel.
At a unit frequency every local component is a unit. The same
observation applies to the difference of the two inverses in
the last condition, giving both product bounds. Multiplication
gives the bottom inequality and preserves continuity at the top.

For positivity, represent each finite matrix on a common slice
as a Gram matrix. Explicitly give the formal span of its column
indices the sesquilinear pairing with that matrix. Quotient by
the vectors of zero square length. They pair to zero with every
vector by positivity of the quadratic polynomial obtained by
adding any complex multiple of a second vector, so the quotient
has a positive inner product. In an orthonormal basis there,
obtained by successive orthogonal subtraction and normalization,
the images of the index vectors have the prescribed pairings.
For each index take the array of products of its vector
coordinates over all factors. Summing products of its conjugate
coordinates against the array for a second index factors into
the product of the prescribed pairings. The resulting product
matrix is therefore itself a Gram matrix, proving positivity.
\end{proof}

\begin{theorem}\label{thm:actual-rational-orbital-sum}
Take the arithmetic scales and \(l_{\tau i}\) of
Lemma~\ref{lem:individual-kappas}, write \(n=\nu+1\), and
let \(T\) contain all places above \(g\), taking both conjugate
components over split places. For \(S\) a subset of the marked
places and a prescribed vector stabilizer as in
Theorem~\ref{thm:marked-energy-reduction}, choose any compact open
additive subgroup \(L\subset E_T\). Use as \(J_i\) the product
of the \(f\) real kernels of Theorem~\ref{thm:actual-real-sieve-family}
and the finite kernel of Theorem~\ref{thm:actual-finite-sieve-family}.
For a finite set \(A\) of rational \(T\)-units in \(E^\times\)
and a subset \(A'\subset A\), form \(D_i(A')\) and \(C_{i+1}\)
from these kernels by Theorem~\ref{thm:quantitative-orbital-sieve}.
Uniformly in these choices and throughout the arithmetic window,
for \(0\le i\le\nu\) and all sufficiently large \(x\),
\[
 0\le\frac{D_i(A')}{C_{i+1}}
       \le\frac{1+b^{-6}+b^{-8}}{\mathcal D\,b^{6k|S|}}.
\]
\end{theorem}
\begin{proof}
The product character is exactly that of
Corollary~\ref{cor:orbital-sampling}, by
Definition~\ref{def:finite-bessel-level}; the inverse map at
the finite component is \(v\mapsto g^{-1}v^{-1}\). Apply
Lemma~\ref{lem:product-local-sieve}. At unit frequencies use
Theorem~\ref{thm:actual-real-sieve-family} with \(\delta=1\)
and Definition~\ref{def:finite-explicit-band}. The product of
all the scalar powers of \(b\) in the frequency bound is
\(b^{14n^9 f}\), since each real factor contributes \(13n^9\)
to the exponent and the finite product contributes \(n^9 f\).
The other factors are exactly
\[
  m_T(\xi_T)^{-(k-n^7)/2} D_T(\xi_T)^{-6}
              \prod_{\tau} C_{\rm B}\,d_{k-n^7,6}(|\xi_\tau|^{-1}).
\]
For a rational prefix product \(z_j\) in
Theorem~\ref{thm:quantitative-orbital-sieve}, the finite head
weight is zero if any \(z_j\) is nonintegral at \(T\).
Otherwise \(D_T(z_j)=1\) for every prefix, and its finite
weight factor is \(G_T(z_j)^{n^6/2}\) by
Definition~\ref{def:finite-explicit-band}. The real factors
are exactly \(d_{n^6,6}(|L_{\tau j}\tau z_j|^{-1})\) from
Definition~\ref{def:explicit-real-band-weights}. Thus both
hypotheses of Theorem~\ref{thm:quantitative-orbital-sieve} hold
with its \(h=n,\ B_\tau=C_{\rm B}\). Its sums are nonnegative
since each common-entry sum is a nonnegative quadratic form
at the top and successive common-entry integrals preserve
nonnegativity, as proved in Theorem~\ref{thm:orbital-tower-sieve}.

Use Theorem~\ref{thm:orbital-degree-budget} with \(\eta=1/2\).
It bounds the error by \(\mathcal D^{-1}b^{-12kf}\). The product
bottom bound is
\(\mathcal D^{-1}b^{-6k|S|}(1+b^{-6})(1+b^{-10})^f\), by
the two local theorems and Definition~\ref{def:row-vandermonde}.
Here \(|S|\le f\) by the split-place description in
Proposition~\ref{prop:marked-local-level}. Also \(f\le x^{12}\le b\)
eventually by Proposition~\ref{prop:level-prime}. Thus
\((1+b^{-10})^f\le\exp(fb^{-10})\le1+2b^{-9}\) eventually,
using the exponential series. Relative to
\(\mathcal D^{-1}b^{-6k|S|}\), the error is at most \(b^{-6kf}\).
For sufficiently large \(x\),
\(2b^{-9}(1+b^{-6})+b^{-6kf}\le b^{-8}\). Adding proves the bound.
\end{proof}

\section{The weighted Siegel inequality for finite correlation entries}

\begin{definition}\label{def:finite-packet-energy}
Fix a finite index set \(\mathcal A\), heads \(p_i\in(E_T^\times)^\nu\)
for \(i\in\mathcal A\), and a positive semidefinite matrix
\(R=(R_{ij})_{i,j\in\mathcal A}\).
For a coordinate \(d\in D_Z\times D_Z\) of the correlation
matrix of Definition~\ref{def:finite-top-correlation}, let
\[
 v_d(l)=(\mathcal J(p_i,l)_d)_{i\in\mathcal A},\qquad
 Q_R(l)=\sum_d v_d(l)^*R v_d(l)\quad(l\in G_n),\qquad H_0=H\cap K_g .
\]
For \(e=l^{-1}z\) define \(E_R(e)=\int_H Q_R(h l)\dd\mu_H(h)\);
set \(E_R(e)=0\) elsewhere on the quadric of norm \(-N\).
\end{definition}

\begin{lemma}\label{lem:finite-orbit-energy-continuity}
The function \(E_R\) is well-defined, nonnegative, compactly
supported and locally constant on the entire quadric. It is
invariant under \(H_0\). For every compact open additive subgroup
\(L\subset E_T\) and \(t\in E_T\) one has
\[
 \int_{[z,e]\in t+L} E_R(e)\dd\delta^{[n]}_{-N}(e)
   = m_Z|L|\sum_{i,j\in\mathcal A} R_{ij}
                        J_n((p_i,t),(p_j,t)).
\]
\end{lemma}
\begin{proof}
By Lemma~\ref{lem:finite-correlation-compact}, \(Q_R\) is locally
constant and supported on a compact subset of \(G_n\); positivity
of \(R\) gives \(Q_R\ge0\). Two representatives of \(e\) differ by
left multiplication by \(H\). Right invariance of \(\mu_H\) and
its finiteness on compact sets show that \(E_R\) is independent
of the representative and finite. There is a single open subgroup
\(V\) of \(G_n\) such that \(Q_R(l v)=Q_R(l)\) for all \(v\in V\).
Indeed its nonzero set is compact and open. Cover it by finitely
many right neighborhoods \(l_i V_i\) with \(V_i\) open subgroups
and \(Q_R\) constant there. Such subgroups form a basis at one
by the congruence neighborhoods in Lemma~\ref{lem:finite-unitary-level-haar}.
Their intersection works on the union and preserves its complement.
It follows that \(l\mapsto\int_H Q_R(h l)\dd\mu_H(h)\) is locally
constant as well.

The chart in Proposition~\ref{prop:finite-top-orbit-section}
contains \(z=N\eta v-v'\), whose last coefficient is \(-1\).
Its formula gives \(\sigma(z)=1\). Around \(e_0=l_0^{-1}z\)
in the quadric, \(l(e)=\sigma(l_0e)l_0\) is therefore a
continuous section, with \(l(e_0)=l_0\). This proves local
constancy at each represented point. All nonzero values are on
the compact image of \(\operatorname{supp}(Q_R)\) under
\(l\mapsto l^{-1}z\). This image is closed and consists of
represented points. Its complement has zero value, proving
the support and continuity assertions everywhere.

For \(h_0\in H_0\), the Bessel vector satisfies
\(W_Z(h_0y)=W_Z(y)\). Indeed \(h_0\) normalizes \(U_T\), fixes
the frame pairs and \(z\), and preserves the simple-root sum
there. In each coordinate \(\phi_g(h_0y a^{-1})\), the factor
\(h_0\in K_g\) can be moved past the unipotent in the support
\(U_TK_g\); the phase is unchanged. Its inverse gives the same
support in the opposite direction, proving the identity by
Definition~\ref{def:finite-bessel-vector}. It follows that
\(\mathcal J(p,h_0 l)=\mathcal J(p,l)\). Substituting \(y=h_0y'\)
in the matrix integral and using commutation with \(w_p\) also gives
\(\mathcal J(p,l h_0^{-1})=\mathcal J(p,l)\). The second identity
shows \(E_R(h_0e)=E_R(e)\).
Finally expansion of \(Q_R\) gives
\(\sum_{i,j}R_{ij}\langle\mathcal J(p_i,l),\mathcal J(p_j,l)\rangle\).
The group disintegration \(\mathcal M_n=\mu_H*\theta_n\) and
the pushforward in Proposition~\ref{prop:finite-top-orbit-section}
give the displayed formula by Theorem~\ref{thm:finite-top-gram}.
\end{proof}

\begin{proposition}\label{prop:finite-entry-packets}
For a represented point \(e\) fix a finite list of pairs
\((l_\theta,d_\theta)\) with \(l_\theta^{-1}z=e\) and
\(d_\theta\in D_Z^2\). Suppose all pairs
\((H_0 l_\theta,d_\theta)\) are distinct. Then
\[
     \mu_H(H_0)\sum_\theta v_{d_\theta}(l_\theta)^*R v_{d_\theta}(l_\theta)
                                                       \le E_R(e).
\]
The same lower bound holds on the \(H_0\)-translates of a
neighborhood of \(e\) in the norm quadric.
\end{proposition}
\begin{proof}
Fix \(l\) over \(e\) and put \(h_\theta=l_\theta l^{-1}\in H\).
For each fixed coordinate \(d\), those right cosets \(H_0h_\theta\)
having \(d_\theta=d\) are pairwise disjoint. On each, the \(d\)-term
of \(Q_R(h l)\) equals the term at \(l_\theta\) by the left
\(H_0\)-invariance of the matrix correlations proved in
Lemma~\ref{lem:finite-orbit-energy-continuity}. Its integral over
that coset is the term times \(\mu_H(H_0)\), by right invariance.
Use nonnegativity to retain these cosets for each \(d\), then sum
over coordinates. This proves the first assertion. Local constancy
and \(H_0\)-invariance in the same lemma prove the second.
\end{proof}

\begin{theorem}\label{thm:finite-entry-siegel}
Suppose \(T\) contains all places above \(gN\). Take \(\mathcal S,t\)
and \(\beta\) as in Theorem~\ref{thm:top-siegel-bound}. At the
point \(e_{c,y}=P(t,(h_c^{-1}y)_T)\) of each \((c,y)\in\mathcal S\)
choose a list as in Proposition~\ref{prop:finite-entry-packets},
allowing an empty list. Use one \(R\) and one list of heads
for all pairs. Write \(G(t)=\prod_\tau(|\tau t|^2-N^2)^{k-1}\).
For all sufficiently small compact open additive subgroups \(L\)
of \(E_T\),
\[
 \begin{split}
 \mu_H(H_0)\sum_{(c,y)\in\mathcal S}\sum_\theta
        v_{d_\theta}(l_\theta)^*R v_{d_\theta}(l_\theta)
 \ \le\ &|\mathcal C(K_c)|c_*P_T(\beta)G(t)\,m_Z\\[-2pt]
 &{}\cdot\sum_{i,j\in\mathcal A}R_{ij}J_n((p_i,t),(p_j,t)).
 \end{split}
\]
The double sum on the right is real and nonnegative.
\end{theorem}
\begin{proof}
Use the individual bounds from Proposition~\ref{prop:finite-entry-packets}
as weights, with weight zero for an empty list. The image of the
local column level in the middle group is \(H_0\) by
Proposition~\ref{prop:marked-local-level}. To apply the ambient
neighborhood condition in Theorem~\ref{thm:top-siegel-bound},
extend \(E_R\) from the closed quadric to the rest of the middle
space by the largest of the finitely many weights, taking zero
for an empty collection. This extension is Borel and invariant
under \(H_0\). Each relative neighborhood on which the
weight is attained or exceeded is the intersection of an open
ambient set with the quadric. Thus that ambient inequality holds
there and after translation by any element of \(H_0\). Theorem~\ref{thm:top-siegel-bound}
and the integral identity in Lemma~\ref{lem:finite-orbit-energy-continuity}
give the assertion after cancellation of \(|L|>0\).
The same identity proves the last statement.
\end{proof}

\section{The unipotent expansion for the cut Poincar\'e columns}

\begin{definition}\label{def:cut-global-columns}
Use \(G_\infty,G_f,\Gamma,K_Y^0,K_Y,h_c\) of
Definition~\ref{def:adelic-row}; the rational unipotent group is
\(U(F)\), \(F=E^{\mathfrak c}\). For \(\gamma\in\Gamma\) write
\(\gamma_\infty,\gamma_f\) for the archimedean matrix in the
scaled basis and the finite adelic matrix. Set
\[
 V_\infty(y)=\bigotimes_\tau V_\tau(y_\tau),\qquad
 A_c(u)=h_c^{-1}u_f h_c,\qquad
 \Delta_{c,c'}(\gamma)=h_c^{-1}\gamma_f h_{c'} .
\]
Here \(V_\tau\) is the bounded continuous cut vector of
Theorem~\ref{thm:actual-bessel-truncation}. The tensor space is
given the product orthonormal basis. A finite matrix \(x\in G_f\)
is integral off \(T\) if it preserves the specified row lattice
at every place outside \(T\). For \(k_0\in K_Y^0\) with image
\(s\in\mathcal G\), let
\(\mathcal T(c,s,c')=\{\gamma\in\Gamma:\Delta_{c,c'}(\gamma)\in k_0K_Y\}\).
\end{definition}

\begin{proposition}\label{prop:actual-cut-poincare-columns}
For every column class \(c\) there is a nonnegative Borel
function \(\beta_c\le1\) on \(G_\infty\) such that
\(\phi_c=\beta_c V_\infty\) is continuous and compactly supported
and \(\sum_{\delta\in\Gamma_c}\phi_c(\delta_\infty y)=V_\infty(y)\).
Periodizing the coordinates of \(\phi_c\) at \(h_cK_Y\) as in
Theorem~\ref{thm:poincare} gives continuous square-integrable
automorphic functions. For two classes put
\[
 M_{c,c'}(\gamma)
    =\int_{G_\infty}\phi_c(y)\phi_{c'}(\gamma_\infty^{-1}y)^*\dd y.
\]
For \(s\in\mathcal G\) their matrix of translated inner products
in that theorem is
\[
                   \sum_{\gamma\in\mathcal T(c,s,c')}M_{c,c'}(\gamma).
\]
Entry \((a,a')\) pairs the translate of the \(a'\)-th function
for \(c'\) against the \(a\)-th function for \(c\).
Only finitely many rational matrices in this sum can have
anywhere nonzero integrand.
\end{proposition}
\begin{proof}
Under \(u\in U_\infty\) the transformation law in
Theorems~\ref{thm:actual-real-bessel-vector} and~\ref{thm:actual-bessel-truncation}
gives \(V_\infty(u y)=\chi_\infty(u)V_\infty(y)\).
Its support lies in \(U_\infty C\) for a compact \(C\) by the
last assertion of the truncation theorem, taking the finite
product of the compact representatives in its proof.
The cutoffs in Theorem~\ref{thm:poincare} therefore apply
simultaneously to all its coordinates. In that construction
\(\beta_c=\xi_c/\sum_{\delta\in\Gamma_c}\xi_c(\delta_\infty\,\cdot\,)\)
where the denominator is positive, and zero elsewhere.
The identity term gives \(\beta_c\le1\). A change of index shows
\(\sum_{\delta\in\Gamma_c}\beta_c(\delta_\infty y)=1\) throughout
the support of \(V_\infty\). The character is one on \(\Gamma_c\)
by Proposition~\ref{prop:full-phase}. These facts prove the
periodization identity and all the asserted regularity properties
by Theorem~\ref{thm:poincare}. Continuous functions of compact
support on the quotient are square-integrable by the local
fundamental-domain measure of Lemma~\ref{lem:quotient-measure}.

For the finiteness assertion put \(B_c=\operatorname{supp}\phi_c\).
If the integrand is nonzero at \(y\), then
\(\gamma_\infty\in B_c B_{c'}^{-1}\), which is compact.
The finite matrix belongs to \(h_c k_0K_Y h_{c'}^{-1}\), also
compact. The bounded denominator and bounded conjugate argument
of Lemma~\ref{lem:quotient-measure} gives finitely many rational
matrices in their product. The inner product formula in
Theorem~\ref{thm:poincare} now gives the stated matrix entry by entry.
\end{proof}

\begin{lemma}\label{lem:global-finite-phase-expansion}
Let \(T\) contain every place above \(g\). For every \(x\in G_f\)
integral off \(T\) and \(y\in G_\infty\),
\[
  \sum_{\substack{u\in U(F)\\ A_c(u)x\in K_Y}}\phi_c(u_\infty y)
                              =\phi_g(x_T)V_\infty(y).
\]
The sum has only finitely many nonzero terms at each \(y\).
\end{lemma}
\begin{proof}
If \(A_c(u)x\in K_Y\), then at \(T\) its component lies in \(K_g\),
so \(\phi_g(x_T)\ne0\). Conversely, if \(\phi_g(x_T)\ne0\), some
\(v\in U_T\) has \(v x_T\in K_g\) by
Definition~\ref{def:finite-bessel-vector}. Extend \(v\) by the
identity off \(T\). Then \(vx\in K_Y\), since \(x\) was integral
there. The open condition \(v'x\in K_Y\) on the finite adelic
unipotent group therefore holds on a nonempty open set.
Conjugation by \(h_c\) is a homeomorphism of this group because
it acts on the middle space fixing the distinguished vector
and all the outer pairs. Lemma~\ref{lem:unip-approx} gives
\(u_0\in U(F)\) with \(A_c(u_0)x\in K_Y\).
All such elements form \(\Gamma_c u_0\): their pairwise quotients
satisfy \(A_c(u)A_c(u_0)^{-1}\in K_Y\), and conversely
multiplication on the left by any element of \(\Gamma_c\)
preserves the level condition by Definition~\ref{def:adelic-row}.

In that case \(A_c(u_0)\) is integral off \(T\), since both
\(A_c(u_0)x\) and \(x^{-1}\) are integral there. The simple-root
sum is unchanged by conjugation by \(h_c\), which fixes each
frame vector and its dual and fixes \(z\). Its character is
trivial on the integral unipotent group outside \(T\) since
\(g\) is a unit there. By Proposition~\ref{prop:full-phase}
and the phase identity of Proposition~\ref{prop:finite-bessel-unfold},
\[
 \chi_\infty((u_0)_\infty)=\chi_f((u_0)_f)^{-1}
        =\chi_g(A_c(u_0)_T)^{-1}=\phi_g(x_T).
\]
Summing the left side of the assertion over \(\Gamma_c u_0\)
now gives \(V_\infty((u_0)_\infty y)\) by
Proposition~\ref{prop:actual-cut-poincare-columns}, the desired
right side. The sum is locally finite in \(y\) because \(\phi_c\)
has compact support and \(\Gamma_c\) preserves the finite coset
\(h_cK_Y\), by Lemma~\ref{lem:quotient-measure}. In the remaining
case both sides vanish by the support criterion already proved.
\end{proof}

\section{Unfolding one rational unipotent double coset}

\begin{definition}\label{def:rational-middle-pivots}
Let \(G_{\rm mid}(F)\) be the group of rational isometries fixing
each \(e_j,e_j^*\), acting on the middle space of
Definition~\ref{def:row-space}. For a rational head
\(p=(p_1,\ldots,p_\nu)\in(E^\times)^\nu\) use the rational flip
\(w_p\) with the same pair matrices as in
Definition~\ref{def:finite-floor-sections}. For
\(\gamma=w_p\sigma,\ \sigma\in G_{\rm mid}(F)\), define
\[
 \begin{split}
 B^{T}_{c,s,c'}(\gamma)
  &=\sum M_{c,c'}(u\gamma v^{-1}),\\
 \mathsf T_\infty(\gamma)&=\int_{G_\infty}
       V_\infty(y)V_\infty(\gamma_\infty^{-1}y)^*\dd y .
 \end{split}
\]
The sum ranges over \(u,v\in U(F)\) such that \(A_c(u)\) is
integral off \(T\) and \(u\gamma v^{-1}\in\mathcal T(c,s,c')\).
Write \(B_{c,s,c'}(\gamma)\) for the same sum with the
exterior integrality condition omitted.
\end{definition}

\begin{lemma}\label{lem:rational-open-terms}
For fixed \(\gamma=w_p\sigma\) the map
\((u,v)\mapsto u\gamma v^{-1}\) on \(U(F)^2\) is injective.
Both sums in Definition~\ref{def:rational-middle-pivots} have
only finitely many terms with anywhere nonzero integrand.
Given finitely many choices of \(c,s,c',\gamma\), one has
\(B^T_{c,s,c'}(\gamma)=B_{c,s,c'}(\gamma)\) for all of them
as soon as \(T\) contains a sufficiently large fixed finite set.
The integral defining \(\mathsf T_\infty(\gamma)\) converges
absolutely, and equals the tensor product of the single-place
matrix correlations.
\end{lemma}
\begin{proof}
The middle group normalizes \(U\), by the flag equations of
Definition~\ref{def:unip}. An element of
\(U\cap w_p U w_p^{-1}\) preserves both the flag of the \(e_j\)
and the flag of the \(e_j^*\), with diagonal entries one on each.
It preserves the two spans and their common orthogonal
complement, the middle space. The matrices on the two spans
are upper triangular in their listed orders, and each is the
inverse conjugate transpose of the other by the pairing. They
are therefore diagonal and equal to the identity. The induced
middle action is already the identity for \(U\). It follows
that the whole intersection is trivial. Equality of two products
\(u\gamma v^{-1}\) puts the quotient of their left unipotents
in this intersection, proving injectivity.
Proposition~\ref{prop:actual-cut-poincare-columns} gives the asserted
finiteness of the rational matrices that can contribute. For
their finitely many corresponding \(u\)'s each \(A_c(u)\)
is integral off finitely many places, by the definition of
the finite adeles. Taking the union proves stabilization.

At \(\tau\), conjugation by the diagonal basis change of
Proposition~\ref{prop:unipotent-covol} sends the \(j\)-th flip to
\(w_j(s_{j-1,\tau}^2\tau p_j)\): apply it to the scaled pair.
It sends \(\sigma\) to a middle isometry. The integral of
the norm product of the two finite-dimensional vectors is
finite at each place by the compact support-intersection
argument in Theorem~\ref{thm:real-top-gram-kernel}. Finite
product integration using Lemma~\ref{lem:product-integration}
gives absolute convergence and the tensor product formula.
\end{proof}

\begin{theorem}\label{thm:actual-finite-coefficient-unfolding}
Let \(T\) contain the places above \(g\), and suppose
\(\Delta=\Delta_{c,c'}(w_p\sigma)\) is integral off \(T\).
Choose \(k_0\in K_Y^0\) of class \(s\) with \((k_0)_T\in Z\).
Let \(a,b'\in D_Z=K_g\backslash Z\) satisfy
\(a(b')^{-1}=(k_0)_T\bmod K_g\), using the quotient group,
and let \(p_T\) be the head evaluated at \(T\). Then with
\(l=(h_c^{-1}\sigma_f h_{c'})_T^{-1}\in G_n\),
\[
 \mathcal J(p_T,l)_{a,b'}\,\mathsf T_\infty(w_p\sigma)
    =H(p_T)\mu(K_g)\,B^{T}_{c,s,c'}(w_p\sigma).
\]
\end{theorem}
\begin{proof}
Write \(\gamma=w_p\sigma\). Since \(h_c,h_{c'}\) act on the
middle space, \(\Delta_T^{-1}=l w_{p_T}^{-1}\).
In the defining integral for the \((a,b')\) entry of
\(\mathcal J/H(p_T)\), let the group variable be \(x_T a\),
using any representative of the coset and right invariance of
\(\mu\). Right \(K_g\)-invariance of \(\phi_g\) in
Proposition~\ref{prop:finite-bessel-unfold} gives the integrand
\(\phi_g(x_T)\overline{\phi_g(\Delta_T^{-1}x_T(k_0)_T)}\).
Let \(x_f\) be \(x_T\) at \(T\) and one elsewhere, and put
\(x'_f=\Delta^{-1}x_f k_0\). Both are integral off \(T\).
Multiply the last integrand by
\(V_\infty(y)V_\infty(\gamma_\infty^{-1}y)^*\).
Lemma~\ref{lem:global-finite-phase-expansion} in the two arguments
identifies this product with the sum over \(u,v\in U(F)\) of
\[
 {\bf1}_{K_Y}(A_c(u)x_f)\,{\bf1}_{K_Y}(A_{c'}(v)x'_f)
      \phi_c(u_\infty y)\phi_{c'}(v_\infty\gamma_\infty^{-1}y)^* .
\]

Put \(\delta=u\gamma v^{-1}\). Its finite position is
\(A_c(u)\Delta A_{c'}(v)^{-1}\). When
\(J=A_c(u)x_f\in K_Y\), the other finite condition is
\(\Delta_{c,c'}(\delta)^{-1}J k_0\in K_Y\).
Since \(k_0\) normalizes \(K_Y\), this is equivalent to
\(\Delta_{c,c'}(\delta)\in k_0K_Y\), independent of \(J\).
The first condition requires \(A_c(u)\) integral off \(T\);
if this holds, its allowed \(x_T\)'s form the single left
translate \(A_c(u)_T^{-1}K_g\) of mass \(\mu(K_g)\).
Change \(y\) to \(u_\infty y\) by left invariance. The archimedean
matrix integral becomes \(M_{c,c'}(\delta)\).
By Lemma~\ref{lem:rational-open-terms}, there are finitely many
pairs for which this integrand can take a nonzero value with
both finite conditions. For each, its finite set of integration
is a compact coset and its archimedean integrand is bounded
with compact support. Thus one may integrate the expansion
term by term. The separated left side is absolutely integrable
by that lemma and by the compact estimate in
Lemma~\ref{lem:finite-correlation-compact}. Integrating proves
the claimed identity.
\end{proof}

\section{Separation and multiplicity of the rational middle terms}

\begin{lemma}\label{lem:rational-middle-separation}
Let \(T\) contain all places above \(gN\), and fix a rational
\(t\) with \(|\tau t|>N\) for all \(\tau\). Take a finite family
of distinct tuples \((c_i,s_i,c_i',\sigma_i)\), with column classes
\(c_i,c_i'\), lifts \(k_i\in K_Y^0\) of \(s_i\), \((k_i)_T\in Z\),
and \(\sigma_i\in G_{\rm mid}(F)\) satisfying
\([z,\sigma_i z]=t\). Assume \(\Delta_i=h_{c_i}^{-1}(\sigma_i)_f h_{c_i'}\)
integral off \(T\). Write
\[
 \sigma_i z=\sum_{a=1}^k y_{i,a}x_a-\frac{t}{N}z,\qquad
 y_i=(y_{i,a})_a,\qquad l_i=(\Delta_i)_T^{-1}.
\]
For each \(a\in D_Z\), let \(d_{i,a}=(a,\bar k_i^{-1}a)\)
where \(\bar k_i\) denotes \((k_i)_T\bmod K_g\).
Then \((c_i,y_i)\) satisfies the conditions of
Theorem~\ref{thm:top-siegel-bound}, and
\(l_i^{-1}z=P(t,(h_{c_i}^{-1}y_i)_T)\).
In each fibre of \(i\mapsto(c_i,y_i)\) the pairs
\((H_0l_i,d_{i,a})\), indexed by \(i,a\), are distinct.
\end{lemma}
\begin{proof}
The coefficient of \(z\) follows from its norm \(-N\).
Taking norms gives \(H(y_i)=(t\bar t-N^2)/N\).
The representatives \(h_c\) fix \(z\) and act on its
perpendicular in the middle space by the definite matrix.
This gives the formula for \(l_i^{-1}z\) using
Definition~\ref{def:top-norm-fiber}.
Off \(T\), the integral middle basis differs from the rational
one by a matrix integral in both directions since \(N\) is a
unit and \(\pi_N\bar\pi_N=N\). The vector \(\Delta_i z\)
is integral, since \(\Delta_i\) preserves the lattice.
Taking its definite coordinates therefore gives precisely
the exterior integrality required of \(h_{c_i}^{-1}y_i\).

Suppose two pairs agree in a fibre. Their first coordinate in
\(d_{i,a}\) gives \(a=a'\) and the second gives
\(\bar k_i=\bar k_{i'}\). Since all places of the global
level lie in \(T\), \(s_i=s_{i'}\) by Proposition~\ref{prop:row-group}.
Also \(c_i=c_{i'}\) and \(\sigma_i z=\sigma_{i'}z\).
Consequently \(\delta=\sigma_i^{-1}\sigma_{i'}\) is a rational
definite isometry extended by the identity on \(z\) and all
outer pairs. At \(T\),
\[
    (h_{c_i'}^{-1}\delta_f h_{c_{i'}'})_T=l_i l_{i'}^{-1}\in H_0 .
\]
Outside \(T\) this position equals \(\Delta_i^{-1}\Delta_{i'}\)
and preserves the row lattice. It therefore belongs globally to
the extended \(K_c\), by Proposition~\ref{prop:marked-local-level}.
The two right classes coincide, \(c_i'=c_{i'}'\), and the
representatives for that class are the same. Its rational
stabilizer at \(K_c\) is trivial by Proposition~\ref{prop:column-dim}.
This gives \(\delta=1\) and \(\sigma_i=\sigma_{i'}\).
Distinctness of the original tuples yields \(i=i'\), as claimed.
\end{proof}

\begin{lemma}\label{lem:exact-packet-level-mass}
In the preceding setting write \(u, S_0\) as in
Definition~\ref{def:packet-amounts}. One has the exact identity
\[
      \frac{m_Z}{\mu_H(H_0)\,|D_Z|\,\mu(K_g)^2}=S_0 u .
\]
\end{lemma}
\begin{proof}
By Definition~\ref{def:finite-cut-correlation},
\(m_Z=\mu(Z)\mu(K_g)/v_g^2\).
Partitioning \(Z\) into its congruence cosets gives
\(\mu(Z)=|D_Z|\mu(K_g)\).
Since \(T\) contains all places above \(g\), the explicit volume
in Lemma~\ref{lem:finite-bessel-radical-volume} is \(v_g=u^{-1}\).
The measure on \(H\) gives \(H\cap K_n\) volume one by
Definition~\ref{def:finite-floor-haar}. Reduction on \(H\cap K_n\)
has kernel \(H_0\) and image of order \(|\bar H|\), as in
Proposition~\ref{prop:marked-local-level}: all reductions at the
marked pairs lift independently, and the other components may
be set to the identity. Thus \(\mu_H(H_0)=|\bar H|^{-1}\).
Substitute and use \(S_0=|\bar H|u\) from the same proposition.
\end{proof}

\begin{theorem}\label{thm:unfolded-rational-packets}
With the tuples of Lemma~\ref{lem:rational-middle-separation},
fix a finite family \(\mathcal A\) of rational heads \(p\) of
length \(\nu\). Assume the positions
\(\Delta_{c_i,c_i'}(w_p\sigma_i)\) are integral off \(T\).
Suppose \(\rho_p,r_0>0\) satisfy
\(\rho_pH(p_T)=r_0\sqrt{G(t)}\) for every \(p\), and that the
Gram matrices of
\[
       \bigl(\rho_p\,\mathsf T_\infty(w_p\sigma_i)\bigr)_{p\in\mathcal A}
\]
in the Hilbert--Schmidt inner product are the same matrix \(R\)
for every \(i\). Then for all sufficiently small compact open
additive subgroups \(L\subset E_T\),
\[
 \begin{split}
 r_0^2 \sum_i\left\|\sum_{p\in\mathcal A}
                      B^T_{c_i,s_i,c_i'}(w_p\sigma_i)\right\|_{\rm HS}^2
 \le P_{\rm pkt}\,P_T(\beta)
    \sum_{p,q\in\mathcal A} R_{pq} J_n((p_T,t),(q_T,t)),
 \end{split}
\]
where \(P_{\rm pkt}\) is the constant of Definition~\ref{def:packet-amounts}.
Here \(t\) in the finite kernel denotes its image at \(T\).
\end{theorem}
\begin{proof}
Take the distinct pairs \((c_i,y_i)\) in Theorem~\ref{thm:finite-entry-siegel}.
Use in each fibre the representatives \(l_i\) with every coordinate
\(d_{i,a}\) supplied by Lemma~\ref{lem:rational-middle-separation}.
The matrix \(R\) is positive semidefinite as a Gram matrix.
The scalar in that theorem at \(i,a\) is the squared norm of
\[
 \sum_{p\in\mathcal A}
       \mathcal J(p_T,l_i)_{d_{i,a}}\,\rho_p\mathsf T_\infty(w_p\sigma_i).
\]
The classes \(a,b'=\bar k_i^{-1}a\) satisfy
\(a(b')^{-1}=\bar k_i\) in the finite quotient. By
Theorem~\ref{thm:actual-finite-coefficient-unfolding} the display is
\[
    \mu(K_g)\,r_0\sqrt{G(t)}
          \sum_{p\in\mathcal A}B^T_{c_i,s_i,c_i'}(w_p\sigma_i).
\]
This quantity is independent of \(a\), giving the multiplicity
\(|D_Z|\). Cancel \(G(t)>0\) from the Siegel inequality and use
Lemma~\ref{lem:exact-packet-level-mass}. The remaining constant is
\(|\mathcal C(K_c)|c_* S_0 u=P_{\rm pkt}\), as asserted.
\end{proof}

\section{The squared mass of the rational open-cell blocks}

\begin{lemma}\label{lem:sharper-marked-pivot-sum}
For every fixed \(D>0\), the bound of
Theorem~\ref{thm:actual-rational-orbital-sum} holds with
\(1+b^{-D}\) in place of \(1+b^{-6}+b^{-8}\), at all
sufficiently large scales depending also on \(D\).
\end{lemma}
\begin{proof}
On the unselected finite places, Theorems~\ref{thm:finite-marked-bottom}
and~\ref{thm:finite-floor-haar} give bottom value at most
\(1+2^{-5n}+2^{-n}\). At a selected place the proof of
Theorem~\ref{thm:marked-energy-reduction} bounds the principal
contribution by \(4g_0^{-k}\); the remainder is at most
\(g_0^{-9k/10}\) by Theorem~\ref{thm:finite-one-place-remainder}.
Their sum is at most \(b^{-6k}\) for large scales since \(g_0\ge b^{10}\).
Multiply using Proposition~\ref{prop:finite-local-energy-product}.
The factor \(2^{-5n}+2^{-n}\) is \(o(b^{-D})\), since
\(n\ge\sqrt x/2\) and \(\log b=x^\epsilon,\ \epsilon<1/100\)
by Proposition~\ref{prop:thick}. At the real places
Theorem~\ref{thm:real-gram-total-mass}, with \(\eta=1/2\), gives
an additional factor bounded by \((1+b^{-9r^2/2})^f=1+o(b^{-D})\),
using \(f\le b\), \(r\to\infty\) and the exponential series.
Thus relative to \(\mathcal D^{-1}b^{-6k|S|}\) the product
bottom is at most \(1+o(b^{-D})\). The frequency error in
Theorem~\ref{thm:orbital-degree-budget} with \(\eta=1/2\) is
at most \(b^{-6kf}\) relative to the same quantity, as in
Theorem~\ref{thm:actual-rational-orbital-sum}. This too is
\(o(b^{-D})\), proving the claim uniformly in all the choices.
\end{proof}

\begin{lemma}\label{lem:rational-packet-archimedean-factor}
For \(l_{\tau a}=s_{a,\tau}^2\) and \(C_n=\prod_{a=0}^{\nu}
(V_E\prod_\tau l_{\tau a}^2)^{-1}\) put
\[
 r_0=\frac{C_{\rm op}^{\,f/2} Z_{\rm sc}}
               {J_c\prod_{\tau,a=0}^{\nu}l_{\tau a}},
 \qquad
 \rho_p(t)=\prod_\tau\omega_\tau\bigl(
     (l_{\tau,j-1}\tau p_j)_{1\le j\le\nu},l_{\tau,\nu}\tau t\bigr).
\]
The product in the denominator of \(r_0\) ranges over all \(\tau\)
and \(0\le a\le\nu\); \(J_c,Z_{\rm sc}\) are from
Definition~\ref{def:packet-amounts}. If every \(p_j\) is a
rational \(T\)-unit and \(|\tau t|>N\) for all \(\tau\), then
\[
 \rho_p(t)H(p_T)=r_0\sqrt{G(t)},\qquad
 \frac{P_{\rm pkt} C_n}{r_0^2\,\mathfrak v^2\,\mathcal D}
                         =\frac{P_{\rm pkt}}{Z_{\rm pkt}}E_2 .
\]
For any \(\sigma\in G_{\rm mid}(F)\) with \([z,\sigma z]=t\),
the Gram matrix of \((\rho_p(t)\mathsf T_\infty(w_p\sigma))_p\)
is exactly the product of the real top kernels at these scaled
heads with common last entry \(l_{\tau,\nu}\tau t\).
\end{lemma}
\begin{proof}
In the scaled basis the last frame pairing is
\(s_{\nu,\tau}^2\tau t=\tau t/N\), and the other flips acquire
the scales in Lemma~\ref{lem:rational-open-terms}.
The top Gram identity in Theorem~\ref{thm:real-top-gram-kernel}
and the tensor product formula of that lemma give the final
assertion since the trace of a tensor product is the product
of the traces, by diagonal entries in a product basis.

In Definition~\ref{def:real-top-correlation} the factor at
\(\tau\) is \(\sqrt{C_{\rm op}}/\dim\pi_\tau\) times
\[
  N^{-(k-1)}(|\tau t|^2-N^2)^{(k-1)/2}
       \prod_{j=1}^\nu (l_{\tau,j-1}|\tau p_j|)^{d_j}.
\]
By Proposition~\ref{prop:actual-bessel-weights}, \(\pi_\tau\) is
the symmetric power of degree \(a\) at the selected embeddings
and degree zero elsewhere, so \(\prod_\tau\dim\pi_\tau=J_c\)
by the monomial count of Proposition~\ref{prop:column-dim}.
By the product formula in Lemma~\ref{lem:denominator-series},
\(m_T((p_j)_T)\prod_\tau|\tau p_j|^2=1\); all finite
components off \(T\) are units. Thus \(H(p_T)\) cancels the
powers of the rational pivots. For \(a<\nu\),
\(d_{a+1}=k+2(\nu-a)-1\), and
\(N^{-(k-1)}=l_{\tau,\nu}^{k-1}\).
The remaining scale product is \(Z_{\rm sc}/\prod_{\tau,a}l_{\tau a}\),
giving the first identity. Finally
\(C_n=V_E^{-n}(\prod_{\tau,a}l_{\tau a})^{-2}\).
Insert the formulas for \(Z_{\rm pkt}\) and \(E_2\) from
Definitions~\ref{def:packet-amounts} and~\ref{def:two-energy-numbers}
to obtain the second identity.
\end{proof}

\begin{theorem}\label{thm:actual-rational-link-square}
Fix \(S\) and its vector conditions as in
Theorem~\ref{thm:marked-energy-reduction}. Let \(\mathcal I\)
be any finite set of distinct \((c_i,s_i,c_i',\sigma_i)\) with
\(s_i\in\mathcal G\) satisfying those conditions and
\(\sigma_i\in G_{\rm mid}(F)\) such that, for
\(t_i=[z,\sigma_i z]\), every \(|\tau t_i|>N\).
For every fixed \(D>0\), uniformly in \(\mathcal I,S\), at all
sufficiently large scales,
\[
 \frac{1}{\mathfrak v^2}\sum_{i\in\mathcal I}
   \left\|\sum_{p\in(E^\times)^\nu}
                  B_{c_i,s_i,c_i'}(w_p\sigma_i)\right\|_{\rm HS}^2
             \le (1+b^{-D}) E_2\, b^{-6k|S|}.
\]
The sum over \(p\) is finite after deleting its zero terms.
The bound also holds for the sum over all such tuples, defined
as the supremum of finite subsums.
\end{theorem}
\begin{proof}
Fix one tuple. If a rational matrix \(\delta\) occurs for a head
\(p\), its frame matrix \(([e_a,\delta e_b])_{a,b\le\nu}\) has
initial minor of order \(j\) equal to \(\prod_{a\le j}p_a\).
Indeed \(\delta=u w_p\sigma v^{-1}\), and in the two ordered
bases of Theorem~\ref{thm:finite-head-minors} its first \(j\)
rows and columns have the two unitriangular factors and diagonal
\(p_1,\ldots,p_j\). This computation is over \(E\) and uses
no integrality. Consecutive minor ratios thus recover the entire
head. By Proposition~\ref{prop:actual-cut-poincare-columns}
only finitely many \(\delta\) can have nonzero integrand,
giving the assertion about the pivot sum.

Choose a finite set \(A\subset E^\times\) containing all entries
of these contributing heads for the finite family, as well as
all \(t_i\). The empty family already satisfies the assertion.
Sum over the full \(A^\nu\) for each tuple. Choose \(T\)
containing all places above \(gN\), the zeros and poles of
\(A\), and the finitely many places required in
Lemma~\ref{lem:rational-open-terms} for stabilization for every
tuple and every head in \(A^\nu\). Enlarge it also to ensure
the exterior integrality of all their finite middle and flipped
positions, possible for a finite list of finite adelic matrices.
Take all conjugate components at each place. Lifts of the
\(s_i\) have \(T\)-components in the subgroup \(Z\) of
Theorem~\ref{thm:marked-energy-reduction}.

Group the finite family by the exact rational value \(t_i=t\).
Lemma~\ref{lem:rational-packet-archimedean-factor} supplies the
weights \(\rho_p(t),r_0\) and the common matrix \(R\) for
Theorem~\ref{thm:unfolded-rational-packets}. Take a single
compact open \(L\) sufficiently small for all these finitely many
inequalities, using a product of sufficiently small ideal balls.
For any fixed \(D'>D\), Theorem~\ref{thm:localized-class-equation}
bounds every exterior factor \(P_T(\beta)\) by \(1+b^{-D'}\).
The right double sum in each inequality is nonnegative by
Theorem~\ref{thm:finite-entry-siegel}. After substituting the
real kernels for \(R\), the sum of these double sums over the
distinct \(t\)'s is \(D_\nu(A')\) of
Theorem~\ref{thm:actual-rational-orbital-sum} with \(A'=\{t_i\}\).
There is no multiplicity in this set of last entries; all
multiplicities on the left were retained within their respective
Siegel sums. Lemma~\ref{lem:sharper-marked-pivot-sum} bounds
\(D_\nu(A')/C_n\) by
\((1+b^{-D'})/(\mathcal D b^{6k|S|})\).
Use the second identity of Lemma~\ref{lem:rational-packet-archimedean-factor}
and \(P_{\rm pkt}/Z_{\rm pkt}\le1+b^{-D'}\) from
Theorem~\ref{thm:packet-normalization}. The relative factor is
at most \((1+b^{-D'})^3\le1+b^{-D}\) eventually by expansion.
Every bound is uniform in the finite family. Taking the
supremum gives the final assertion.
\end{proof}

\section{Rational frame intersections and open double cosets}

\begin{definition}\label{def:rational-frame-invariants}
Set \(f_a=e_a\) for \(1\le a\le\nu,\ f_n=z,\ n=\nu+1\),
and let \(L_f=\sum_{a=1}^n E f_a\). For \(\gamma\in\Gamma\) put
\(Z_{ab}(\gamma)=[f_a,\gamma f_b]\), and let \(D_j(\gamma)\)
be the leading minor of size \(j\) of \(Z(\gamma)\), with
\(D_0=1\). Write
\[
 \mathfrak d(\gamma)=D_n(\gamma)\overline{D_n(\gamma)}
               -N^2D_\nu(\gamma)\overline{D_\nu(\gamma)}\in F.
\]
Call the condition \(D_1\cdots D_\nu\ne0\) the open flag
condition. These symbols apply to rational matrices. The
initial minors \(D_j(p)\) of a pivot list are the scalar products
of Definition~\ref{def:finite-head-bounds}.
\end{definition}

\begin{theorem}\label{thm:rational-frame-discriminant}
For \(1\le b\le n\) let \(c_b=(-1)^{n+b}\det Z_{\widehat n,\widehat b}\),
deleting the indicated row and column. Let \(Y\in E^k\) be the
definite coordinates of \(w=\sum_b c_b\gamma f_b\). Then
\[
                 \mathfrak d(\gamma)=N H(Y),\qquad
 \mathfrak d(\gamma)=0\quad\Longleftrightarrow\quad
                         L_f\cap\gamma L_f\ne\{0\}.
\]
In particular a nonzero discriminant is totally positive.
All \(D_j,\mathfrak d\) are invariant under left and right
multiplication by \(U(F)\). If \(\gamma\in\mathcal T(c,s,c')\),
its frame entries and all \(D_j\) are algebraic integers.
Their reductions at the marked places depend only on \(s\).
\end{theorem}
\begin{proof}
Expansion along the last row gives \([z,w]=D_n\). Using a repeated
row in that expansion gives \([e_a,w]=0\) for \(a\le\nu\),
so \(w\) has no dual coordinates. Its coefficient at \(z\) is
\(-D_n/N\) by Definition~\ref{def:row-space}. Thus
\([w,w]=H(Y)-D_n\bar D_n/N\). On the other hand, the original
frame has Gram matrix \(\operatorname{diag}(0,\ldots,0,-N)\),
so \([w,w]=-N c_n\bar c_n=-N D_\nu\bar D_\nu\). This proves
the identity. Since \(H\) is the sum of coordinate norms, it
is positive at every real embedding for nonzero \(Y\).

Suppose the first \(\nu\) rows of \(Z\) are independent.
Then \(c=(c_b)_b\) spans their common nullspace and is nonzero.
Indeed one of their minors of order \(\nu\) is nonzero and
the nullspace has dimension one by elimination. The equalities
just proved show \(w\in L_f\) exactly when \(Y=0\); and \(w\ne0\)
since the \(\gamma f_b\) are independent. Conversely any vector
of \(L_f\cap\gamma L_f\) is orthogonal to the first \(\nu\) frame
vectors, so its coordinates in \((\gamma f_b)\) are a multiple
of \(c\). This proves the equivalence in this case. If these rows
are dependent, all \(c_b=0\) and the discriminant vanishes.
Choose a nonzero \(v\) in the span of the \(e_a\) orthogonal to
all \(\gamma f_b\) using a nonzero left relation among the rows,
conjugating its coefficients. Then \(\gamma^{-1}v\) has no dual
coordinates and no \(z\) coordinate. Its norm is zero, so
definiteness forces all its \(x_a\) coordinates to be zero.
Hence \(v\) lies in the span of the \(\gamma e_a\) as well,
proving the equivalence in this case.

On \(L_f\) a unipotent in \(U\) is upper unitriangular in the
listed basis, by Definition~\ref{def:unip}; in particular the
image of \(z\) differs from \(z\) by earlier frame vectors.
Left and right unipotents therefore transform \(Z\) on the
left by a lower unitriangular matrix and on the right by an
upper unitriangular matrix. Each initial determinant is invariant.
For the arithmetic assertion, all \(h_c\) fix the frame, so the
finite entries for \(\gamma\) are precisely those of the frame
matrix for \(\Delta_{c,c'}(\gamma)\). That matrix preserves the
row lattice and its prescribed marked reduction is \(s\).
The frame vectors and their pairings are integral in that
lattice, by Proposition~\ref{prop:row-group}. Integrality at
all finite places gives membership in \(O_E\) by
Lemma~\ref{lem:denominator-series}. Reduction and determinants
commute by their polynomial formulas.
\end{proof}

\begin{theorem}\label{thm:rational-open-decomposition}
A rational isometry \(\gamma\) satisfies the open flag condition
if and only if it can be written
\(\gamma=u w_p\sigma v^{-1}\), with \(u,v\in U(F)\),
\(p\in(E^\times)^\nu,\ \sigma\in G_{\rm mid}(F)\).
All four factors are unique, and then
\[
 p_j=D_j(\gamma)/D_{j-1}(\gamma),\qquad
 t=[z,\sigma z]=D_n(\gamma)/D_\nu(\gamma),\qquad
 \mathfrak d(\gamma)=D_\nu(\gamma)\overline{D_\nu(\gamma)}
                                      (t\bar t-N^2).
\]
Thus in this open flag condition a nonzero discriminant is
equivalent to \(|\tau t|>N\) at every infinite place.
\end{theorem}
\begin{proof}
In the asserted decomposition the matrix \(Z(w_p\sigma)\)
is diagonal with first \(\nu\) entries \(p_1,\ldots,p_\nu\) and
last entry \(t\). Theorem~\ref{thm:rational-frame-discriminant}
therefore gives all the displayed identities and necessity of
the open flag condition.

For existence remove one outer pair at a time. In the first
pair write \(\gamma e_1=d e_1+y+c e_1^*\), where \(c=D_1\ne0\)
and \(y\) belongs to the intervening space. In its first
Heisenberg layer send \(e_1^*\) to
\(e_1^*+y/c+(d/c)e_1\). This is an allowed element \(u_1\):
isotropy of \(\gamma e_1\) gives
\(d/c+\overline{d/c}=-[y/c,y/c]\), the necessary and sufficient
trace condition of Lemma~\ref{lem:heisenberg-layer}.
The isometry \(w_1(c)^{-1}u_1^{-1}\gamma\) fixes \(e_1\).
Its induced isometry on \(e_1^\perp/Ee_1\) is a middle matrix
\(g'\). After removal of its diagonal extension, the matrix
has first column \(e_1\), last row the last coordinate functional
by isometry, and identity intervening block. The same Heisenberg
calculation puts the remainder in the first layer. Thus
\(\gamma=u_1 w_1(c) g' v_1^{-1}\) for some \(v_1\) there.
The invariance proved above shows that each remaining leading
minor in \(g'\) is the next one in \(\gamma\) divided by \(c\),
by the diagonal first block. Repeat this construction on \(g'\).
All factors from deeper layers commute with the earlier flips.
Collect their unipotent factors on the two sides to get the
assertion. In the case of zero outer pairs the middle group
is already the whole group.

For uniqueness the formulas recover \(p\). If two decompositions
with this \(p\) agree, conjugating the quotient of their left
unipotents by \(w_p^{-1}\) gives an element simultaneously in
\(w_p^{-1}Uw_p\) and \(G_{\rm mid}U\). The former preserves the
flag of dual vectors with diagonal one; the latter preserves
the original flag with diagonal one. The proof of
Lemma~\ref{lem:rational-open-terms} forces identity on both
spans. Their common orthogonal complement is preserved, and the
first of these groups has identity induced middle action, giving
identity everywhere. In the second group the middle projection
now gives equality of the two middle factors, after which the
right unipotents agree also. Finally the discriminant is totally
positive whenever it is nonzero by
Theorem~\ref{thm:rational-frame-discriminant}. The last displayed
identity gives the claimed equivalence at infinity.
\end{proof}

\begin{corollary}\label{cor:rational-open-block-identification}
For fixed \(c,s,c'\), the cut Poincar\'e coefficient coming from
terms \(\delta\in\mathcal T(c,s,c')\) with the open flag condition
and \(\mathfrak d(\delta)\ne0\) is the sum, over all
\(\sigma\in G_{\rm mid}(F)\) with \(|\tau[z,\sigma z]|>N\) for
every \(\tau\), of the matrices \(\sum_p B_{c,s,c'}(w_p\sigma)\).
Only finitely many terms can contribute. If at each \(j\le\nu\)
the marked reductions of \(D_j\) prescribed by \(s\) include a
nonzero component, the open flag condition holds for every
\(\delta\in\mathcal T(c,s,c')\).
\end{corollary}
\begin{proof}
By Theorem~\ref{thm:rational-open-decomposition} each indicated
rational term determines exactly one quadruple \(u,p,\sigma,v\)
in the definition of \(B_{c,s,c'}\), and conversely every
summand on the right meets the two stated conditions. The
finiteness assertion of Proposition~\ref{prop:actual-cut-poincare-columns}
therefore permits this regrouping of its coefficient formula.
The last assertion follows from the integrality and reduction in
Theorem~\ref{thm:rational-frame-discriminant}, since a zero
algebraic integer reduces to zero in every component.
\end{proof}

\section{Cancellation by the character on a common unipotent subgroup}

\begin{definition}\label{def:rational-character-agreement}
For \(\delta\in\Gamma\), put
\[
 U_\delta=U(F)\cap\delta U(F)\delta^{-1},\qquad
 \epsilon_\delta(v)=
 \chi_\infty(v_\infty)\overline{\chi_\infty((\delta^{-1}v\delta)_\infty)}
 \quad(v\in U_\delta).
\]
Say the rational characters agree at \(\delta\) when
\(\epsilon_\delta=1\) on \(U_\delta\). By Proposition~\ref{prop:full-phase}
this is equivalent to agreement of the two finite adelic
characters on the same rational elements.
For \(\delta\in\mathcal T(c,s,c')\), let
\(\Gamma_c\delta\Gamma_{c'}\) be its arithmetic unipotent
double orbit, and put \(\Lambda=\Gamma_c\cap\delta\Gamma_{c'}\delta^{-1}\).
\end{definition}

\begin{lemma}\label{lem:cut-orbit-unfold}
With the preceding notation put
\(F_\delta(y)=V_\infty(y)V_\infty(\delta_\infty^{-1}y)^*\).
It is left \(\Lambda\)-invariant and integrable on
\(\Lambda\backslash G_\infty\), where integration means integration
on any measurable domain of representatives for left \(\Lambda\).
One has
\[
       \int_{\Lambda\backslash G_\infty}F_\delta(y)\dd y
          =\sum_{\delta'\in\Gamma_c\delta\Gamma_{c'}}
                                             M_{c,c'}(\delta').
\]
\end{lemma}
\begin{proof}
For \(v\in U_\delta\) equivariance in
Proposition~\ref{prop:actual-cut-poincare-columns} gives
\(F_\delta(v_\infty y)=\epsilon_\delta(v)F_\delta(y)\).
The character is trivial on \(\Lambda\) since the two unipotents
in its definition belong to \(\Gamma_c,\Gamma_{c'}\), respectively.
The left action of \(\Lambda\) on \(G_\infty\) is free and
properly discontinuous, since \(\Lambda\subset\Gamma_c\) and
\(\Gamma_c\) preserves \(h_cK_Y\) at finite places. The construction
of Lemma~\ref{lem:quotient-measure} gives a measurable domain
\(D\) tiling \(G_\infty\) under its left translates, and
independence of the domain for nonnegative or integrable tests.

Choose one pair \((u,v)\in\Gamma_c\times\Gamma_{c'}\) for each
distinct matrix \(\delta'=u\delta v^{-1}\) in the orbit.
All pairs giving that matrix are
\((u\lambda,v\delta^{-1}\lambda\delta)\) for \(\lambda\in\Lambda\).
The entire orbit stays in \(\mathcal T(c,s,c')\), since multiplication
by the finite positions of \(\Gamma_c,\Gamma_{c'}\) preserves the
level coset by normality. In \(M_{c,c'}(\delta')\), replace the
variable by \(u_\infty y\). Its integrand becomes
\[
       \beta_c(u_\infty y)
       \beta_{c'}(v_\infty\delta_\infty^{-1}y)F_\delta(y).
\]
Tile the integral over \(\lambda_\infty D\) and use invariance
of \(F_\delta\). Summing over \(\lambda\) and the representative
pairs gives every pair in \(\Gamma_c\times\Gamma_{c'}\) once.
At a point where \(F_\delta\ne0\) both sums over cutoffs are
one by the proof of Proposition~\ref{prop:actual-cut-poincare-columns}.
The same calculation applies first to the nonnegative norm of
the integrand by Lemma~\ref{lem:product-integration}. Its original
sum is finite by the last assertion of that proposition, since
the associated integrands have compact support and only finitely
many matrices can contribute. Thus \(\int_D\|F_\delta\|<\infty\).
Applying the calculation entry by entry gives the identity.
\end{proof}

\begin{theorem}\label{thm:incoherent-rational-orbit-zero}
If the rational characters fail to agree at
\(\delta\in\mathcal T(c,s,c')\), the sum of the cut coefficients
on \(\Gamma_c\delta\Gamma_{c'}\) vanishes. The same conclusion
holds for a union of such orbits at fixed \(c,s,c'\).
\end{theorem}
\begin{proof}
In \(G_f\) put
\[
 K' = h_c K_Y h_c^{-1}\cap
     \delta_f h_{c'}K_Y h_{c'}^{-1}\delta_f^{-1}.
\]
This group is compact open and
\(\Lambda=\{a\in U_\delta:a_f\in K'\}\) by
Definition~\ref{def:adelic-row}. Choose \(v\in U_\delta\)
with \(\epsilon_\delta(v)\ne1\). Conjugation by \(v\) preserves
\(U_\delta\), since it belongs to that group. Hence
\(\Lambda'=\Lambda\cap v^{-1}\Lambda v\) is characterized there
by the finite condition \(K'\cap v_f^{-1}K'v_f\). It has finite
index \(m\) in \(\Lambda\) and finite index \(m'\) in
\(v^{-1}\Lambda v\): the two maps on left coset sets inject
into the respective coset sets of compact groups by an open
subgroup. Those coset sets are finite since the open cosets
cover a compact set.

Write \(I=\int_{\Lambda\backslash G_\infty}F_\delta\), absolutely
convergent by Lemma~\ref{lem:cut-orbit-unfold}. The function is
invariant also under \(v^{-1}\Lambda v\), since it transforms
on \(U_\delta\) by a scalar character trivial on \(\Lambda\).
A domain for \(\Lambda'\) is a disjoint union of \(m\) translates
of a domain for \(\Lambda\), using representatives of the right
cosets \(\Lambda'\backslash\Lambda\). Thus its integral is \(mI\).
If \(D\) is a domain for \(\Lambda\), then \(v_\infty^{-1}D\)
is a domain for \(v^{-1}\Lambda v\) and its integral is
\(\epsilon_\delta(v)^{-1}I\), by left invariance of group measure.
Tiling a domain for \(\Lambda'\) using that subgroup instead
gives \(m'\epsilon_\delta(v)^{-1}I\). Independence of the
integral from the chosen domain gives equality of the two
expressions. If \(I\ne0\), a nonzero matrix coordinate shows
\(\epsilon_\delta(v)=m'/m\), a positive real of modulus one.
This contradicts the choice of \(v\). Thus \(I=0\), proving
the first claim. At fixed \(c,s,c'\) there are only finitely
many contributing matrices by Proposition~\ref{prop:actual-cut-poincare-columns},
so sums over unions of these orbits also vanish.
\end{proof}

\section{The shared flag of an agreeing rational term}

\begin{definition}\label{def:rational-shared-flag}
Write \(L_a=\sum_{b=1}^a E f_b\) for \(0\le a\le n\), where
\(L_0=0,\ L_n=L_f\). A vector shared at index \(a\) with
coefficients \(c,c'\) is a vector
\(\sum_{b\le a}x_b f_b=\sum_{b\le a}y_b\gamma f_b\) with
\(x_a=c,\ y_a=c'\). The height of a nonzero vector in \(L_f\)
is the largest index with nonzero coefficient, and the height
in \(\gamma L_f\) uses the basis \((\gamma f_b)\).
\end{definition}

\begin{lemma}\label{lem:rational-character-flag-step}
Suppose the rational characters agree at \(\gamma\). If there
is a vector shared at \(a<n\) with coefficients \(c,c'\), then
there is a vector shared at \(a+1\) with coefficients \(c,c'\).
At any index \(a\), existence of such a vector implies
\(c\bar c=c'\bar c'\).
\end{lemma}
\begin{proof}
Let \(w\) be a shared vector at \(a<n\) and let \(m\) be any
\(E\)-linear functional vanishing on \(L_a+\gamma L_a\).
Nondegeneracy of the row matrix in Proposition~\ref{prop:row-group}
gives \(v\) with \(m=[v,\,\cdot\,]\) by inverting that matrix
and conjugating coefficients. For \(r\in E\) put \(X=rw\)
and choose \(\eta\in E\) with \(\eta+\bar\eta=1\), for instance
\(1/2\). We have \([X,X]=[v,X]=0\), since \(a<n\). Define
\[
       U_r y=y+ X[v,y]-v[X,y]-\eta[v,v]X[X,y].
\]
This is a rational isometry. In expanding \([U_r y,U_r y']\),
the two terms involving \(\overline{[v,y]}[X,y']\) cancel,
as do those involving \(\overline{[X,y]}[v,y']\).
The remaining excess is
\((1-\eta-\bar\eta)[v,v]\overline{[X,y]}[X,y']=0\).
Preservation of a nondegenerate form implies injectivity and
hence invertibility in finite dimension.

For every frame vector \(f_b\) the last two added terms vanish.
When \(b\le a\), the first one does too; when \(b>a\),
it belongs to \(L_a\). On \(L_\nu^\perp\) the action is the
identity modulo \(L_\nu\), since \([X,y]=0\) there. These facts
give \(U_r\in U(F)\) by Definition~\ref{def:unip}. The same
proof in the frame \(\gamma f_b\), using \(X\in\gamma L_a\)
and \(m(\gamma L_a)=0\), shows that
\(U_r\in\gamma U(F)\gamma^{-1}\).
For \(w=\sum_{b\le a} x_b f_b\), the simple-root sum is
\(r\sum_{b\le a}x_b m(f_{b+1})=r c\,m(f_{a+1})\).
In the other frame it is \(r c'\,m(\gamma f_{a+1})\).
By Definition~\ref{def:rational-character-agreement}, the
finite characters of these two values agree for every \(r\).
Writing \(d=c\,m(f_{a+1})-c'\,m(\gamma f_{a+1})\), this says
\(\psi_f(-r d/(g\varpi))=1\) for every \(r\in E\).
If \(d\ne0\), take \(r=g\varpi/(4fd)\). Proposition~\ref{prop:additive}
evaluates the character by the rational trace; the trace here
is \(-1/2\) since \([E:\Q]=2f\). Its value is \(-1\), impossible.

Thus every functional annihilating \(L_a+\gamma L_a\)
annihilates \(c f_{a+1}-c'\gamma f_{a+1}\). Extending a basis
of that subspace to a full basis and testing the remaining
coordinate functionals proves that the difference belongs to
\(L_a+\gamma L_a\). Moving the earlier terms to the two sides
gives a shared vector at \(a+1\) with the asserted coefficients.
Iterate up to \(n\). There the norms of the two descriptions
are \(-N c\bar c\) and \(-N c'\bar c'\); their equality proves
the last assertion, also when \(a=n\) initially.
\end{proof}

\begin{theorem}\label{thm:rational-agreeing-levels}
Assume the rational characters agree at \(\gamma\), and set
\(M=L_f\cap\gamma L_f\). If \(M\ne0\), then there is a unique
\(0\le j\le\nu\) and a unique \(\lambda\in E^\times\) with
\(\lambda\bar\lambda=1\) such that \(M\) has a basis
\((w_a)_{j<a\le n}\) of shared vectors at \(a\) with coefficients
\(\lambda,1\). The heights of every nonzero vector of \(M\)
in the two flags coincide. At this level \(D_j(\gamma)\ne0\)
and \(\mathfrak d(\gamma)=0\). If \(M=0\), then
\(\mathfrak d(\gamma)\ne0\) and \(D_n(\gamma)\ne0\).
\end{theorem}
\begin{proof}
Given a nonzero vector of \(M\), take the larger of its heights
in the two flags. At that index at least one of its leading
coefficients is nonzero. The equality of their norms in
Lemma~\ref{lem:rational-character-flag-step} forces both to be
nonzero, proving equality of heights. At a fixed occurring
height two nonzero leading coefficient pairs are proportional:
subtracting a suitable multiple eliminates the leading
coefficient on one side, and the same lemma forces the other
to vanish. Hence their ratio is determined at that height
and has norm one.

Let \(j+1\) be the minimum occurring height, normalize the second
leading coefficient to one and call the first \(\lambda\).
Lemma~\ref{lem:rational-character-flag-step} supplies vectors
with these same leading coefficients up to height \(n\).
They are independent by their distinct heights and span \(M\):
subtract suitable multiples to eliminate successive leading
coefficients in the second flag, and a nonzero remainder of
height at most \(j\) is excluded by minimality. This also proves
uniqueness of \(j,\lambda\).

If \(j>0\) and \(D_j(\gamma)=0\), then a nonzero vector
\(y\in\gamma L_j\) is orthogonal to \(L_j\) by the singularity
of the pairing matrix. It is orthogonal also to the \(w_a\),
since \(j\le\nu\), \(y\in\gamma L_\nu\) and \(w_a\in\gamma L_f\).
The first leading coefficients show
\(L_f=L_j+\sum_{a>j}Ew_a\), so \(y\perp L_f\).
In row coordinates this eliminates its dual and \(z\) entries.
It is isotropic, being in \(\gamma L_j\), so definiteness
eliminates its \(k\) middle coordinates as in
Theorem~\ref{thm:rational-frame-discriminant}. This puts
\(y\in M\), although its height in the second flag is at most
\(j\), a contradiction. For \(j=0\) the minor is one.
The assertion about \(\mathfrak d\) in both cases follows from
Theorem~\ref{thm:rational-frame-discriminant}. In the case
\(M=0\) it is positive at every real embedding, so
\(D_n\bar D_n=\mathfrak d+N^2D_\nu\bar D_\nu\) gives
\(D_n\ne0\).
\end{proof}

\section{The middle relation of two rational flags}

\begin{definition}\label{def:global-middle-relation}
In the space of Definition~\ref{def:row-space} set
\(P=L_\nu,\ M_0=X_0\oplus Ez,\ X_0=\sum_{a=1}^k E x_a\).
Write \(h_0\) for the restriction of the hermitian form to \(M_0\)
and \(\pi:P^\perp=P\oplus M_0\longrightarrow M_0\) for the projection.
For a rational isometry \(\gamma\), put
\[
 W_\gamma=(P+\gamma P)^\perp,\qquad
 \mathcal L_\gamma=
 \{\,(\pi v,\pi\gamma^{-1}v):v\in W_\gamma\,\}
       \ \subset M_0\oplus M_0 .
\]
Equip this last direct sum with the hermitian form
\(B=h_0\oplus(-h_0)\). A subspace is called totally isotropic
if the form vanishes on every pair of its elements.
\end{definition}

\begin{proposition}\label{prop:global-middle-relation}
Let \(c=\dim_E(P\cap\gamma P)\) and let \(a\) be the rank of the
\(\nu\)-by-\(\nu\) matrix \(([e_i,\gamma e_j])_{i,j\le\nu}\).
Then \(\mathcal L_\gamma=\mathcal L_\gamma^{\perp_B}\) has
dimension \(k+1\), and each coordinate projection on
\(\mathcal L_\gamma\) has kernel of dimension \(\nu-a-c\in\{0,1\}\).
It is the graph of an isometry of \(M_0\) exactly when
\(a+c=\nu\). Moreover, for \(u,v\in U(F)\),
\[
        \mathcal L_{u\gamma v^{-1}}=\mathcal L_\gamma,\qquad
        \mathcal L_{\gamma^{-1}}
           =\{(y,x):(x,y)\in\mathcal L_\gamma\}.
\]
\end{proposition}
\begin{proof}
The form is nondegenerate by Proposition~\ref{prop:row-group}.
For any subspace its independent basis vectors give independent
orthogonality conditions, by the invertible matrix of the form.
It follows that \(\dim_E W_\gamma=r-(2\nu-c)=k+1+c\).
Both coordinates of its map to the relation vanish precisely
on \(P\cap\gamma P\), all of which is indeed in \(W_\gamma\)
by isotropy of the two outer spaces. Eliminating that kernel
gives dimension \(k+1\). For \(y,y'\in P^\perp\) one has
\([y,y']=h_0(\pi y,\pi y')\) because \(P\) is the radical
there. Applying this identity to elements of \(W_\gamma\)
and to their inverse images under \(\gamma\) gives total
isotropy of the relation for \(B\). Its orthogonal complement
has dimension \(2(k+1)-(k+1)=k+1\), proving equality.

The vectors of \(W_\gamma\) with first coordinate zero are
\(P\cap(\gamma P)^\perp\), of dimension \(\nu-a\) by the
pairing matrix. Quotienting by the simultaneous kernel gives
\(\nu-a-c\). Interchanging the flags conjugate-transposes the
pairing matrix, so the second kernel has the same dimension.
In either kernel the other component is injective and has
totally isotropic image in \(M_0\). Projection of that image
on \(Ez\) is injective, since in \(X_0\) a nonzero vector
has strictly positive norm at each complex embedding. This
bounds its dimension by one. When the kernels vanish both
projections are isomorphisms by the dimension calculation.
Their composite, with the first inverted, preserves \(h_0\)
by total isotropy. Conversely every graph of an isometry has
zero kernels. For the last identities, \(W_{u\gamma v^{-1}}
=uW_\gamma\). The action of \(U\) on \(P^\perp/P\) is the
identity by Definition~\ref{def:unip}, so the two projected
coordinates of \(u y\) agree with the old ones of \(y\).
For inversion use \(y\mapsto\gamma^{-1}y\).
\end{proof}

\begin{theorem}\label{thm:singular-graph-middle}
On an open flag orbit \(u w_p\sigma v^{-1}\) the relation
\(\mathcal L_\gamma\) is the graph of \(\sigma^{-1}|_{M_0}\).
Two such relations are equal exactly when the two rational
middle factors are equal. If the characters agree and
\(\mathfrak d(\gamma)=0\), let \(j,\lambda\) be as in
Theorem~\ref{thm:rational-agreeing-levels}. Then
\(\mathcal L_\gamma\) is also a graph. Its isometry \(A\)
satisfies \(A z=\lambda^{-1}z\) and restricts to an isometry
of \(X_0\).
\end{theorem}
\begin{proof}
By Proposition~\ref{prop:global-middle-relation} remove the
two outer unipotents. For \(w_p\sigma\) the second outer
space is the span of the dual vectors, so \(W_{w_p\sigma}=M_0\).
The Weyl representative acts identically on this middle space
and \(\sigma\) fixes the hyperbolic pairs, as in
Theorem~\ref{thm:rational-open-decomposition}. This proves the
first two assertions.

In the singular case Theorems~\ref{thm:rational-frame-discriminant}
and~\ref{thm:rational-agreeing-levels} give a basis
\(w_{j+1},\ldots,w_n\) for \(L_f\cap\gamma L_f\). The vectors
with both heights at most \(\nu\) form \(P\cap\gamma P\) and
are exactly the span of the first \(\nu-j\) of these vectors.
In fact the coefficient at \(z\) of \(w_n\) is \(\lambda\ne0\)
and all earlier ones have zero coefficient there. Thus
\(c=\nu-j\). The nonzero leading minor \(D_j\) gives \(a\ge j\)
with the empty minor included. The kernel computation in
Proposition~\ref{prop:global-middle-relation} gives
\(a\le\nu-c=j\), so that proposition gives a graph.
The vector \(w_n\) belongs to both perpendicular spaces,
since \(L_f\subset P^\perp\) and
\(\gamma L_f\subset(\gamma P)^\perp\). Its two projected
coordinates are \(\lambda z,z\). This proves the claimed
action on the line and, by preservation of pairings, on its
orthogonal complement in \(M_0\).
\end{proof}

\begin{proposition}\label{prop:canonical-singular-frame}
Under the singular hypotheses of Theorem~\ref{thm:singular-graph-middle}
put \(D=D_j\), \(Z=Z(\gamma)\), and write
\[
 P_j=\begin{pmatrix}Z_{1:j,\,1:j}&Z_{1:j,\,j+1:n}\\0&I_{n-j}\end{pmatrix}.
\]
For each \(b>j\) there is a unique shared vector \(v_b\) whose
last \(n-j\) coordinates in the transported frame are those of
the \(b\)-th coordinate vector \(\varepsilon_b\) of \(E^n\).
Its entire transported coordinate column is \(P_j^{-1}\varepsilon_b\).
It has height \(b\) in both flags and first-frame coefficient
\(\lambda\) at \(b\). If \(\gamma\in\mathcal T(c,s,c')\), the
coordinates of \(D v_b\) in both frames belong to \(O_E\).
\end{proposition}
\begin{proof}
The shared basis from Theorem~\ref{thm:rational-agreeing-levels}
has upper unitriangular matrix on the final \(n-j\) coordinates
in the transported frame. Successively removing its last
nonzero coordinate proves unique solvability for any such
tail. For the indicated tail this procedure uses \(w_b\) with
coefficient one and only earlier vectors besides it. This
proves the heights and leading coefficient assertion. If \(y\)
denotes the full transported coordinate column, the equality
\(v_b=\sum y_i\gamma f_i\in L_f\) gives \((Zy)_i=0\)
for \(1\le i\le j\) by orthogonality to those \(e_i\).
Together with the tail, this is \(P_j y=\varepsilon_b\).
Expansion along the last rows gives \(\det P_j=D\ne0\).

For an arithmetic term, the entries of \(P_j\) are integral
by Theorem~\ref{thm:rational-frame-discriminant}. Thus
\(D y=\operatorname{adj}(P_j)\varepsilon_b\) is integral
by the cofactor formula. At every finite place the vector
\(D v_b\) belongs to \(\gamma\Lambda_{c'}=\Lambda_c\), since
the transported frame vectors are in the transported lattice.
In \(\Lambda_c\) the \(f_i\) are part of a local integral basis
by Definitions~\ref{def:row-space} and~\ref{def:unip}:
apply the local class matrix to the displayed integral basis,
which fixes each \(e_i\) and \(z\). The first-frame coordinates
are therefore integral at every finite place. Apply
Lemma~\ref{lem:denominator-series} to each rational entry.
\end{proof}

\section{The determinant of the pairing between two middle relations}

\begin{definition}\label{def:middle-lift-disc}
Put \(t(x,x')=([z,x],[z,x'])\) on \(M_0\oplus M_0\). On that
space let
\[
 Q=\tfrac12(h_0\oplus h_0)+N^{-1}t^*t ,
\]
where \(t^*t\) denotes the pullback of the usual hermitian
product on \(E^2\). In the coordinates of
Definition~\ref{def:row-space} its diagonal is
\(\tfrac12(I_k,N,I_k,N)\), so it is positive definite at
every infinite place. Write \(\widehat O=\prod_{\mathfrak p}
O_{E,\mathfrak p}\). For a family \(\alpha=(\alpha_{\mathfrak p})\)
of nonzero finite local scalars, units at almost all primes, put
\[
 I(\alpha)=\prod_{\mathfrak p}(\#\,O_E/\mathfrak p)^{
                         \operatorname{ord}_{\mathfrak p}\alpha_{\mathfrak p}} .
\]
It is multiplicative, and \(I(\bar\alpha)=I(\alpha)\), since
conjugation permutes the places with their residue degrees.
For \(a\in E^\times\) one has \(I(a)=|\operatorname{Nm}_{E/\Q}a|\):
for \(a\) integral both sides give the index of multiplication
by \(a\) in \(O_E\), by the integral-basis determinant and
Lemma~\ref{lem:class-covol}; division gives the general case.

A finite lattice of a rational space is a family of full
local integral lattices equal to the span of a fixed rational
basis at almost all places. For such a lattice \(A\) and a
totally positive definite hermitian form \(Q'\) on its rational
span, put \(\Delta_{Q'}(A)=I(\det(A^*Q'A))^{1/2}\). In this
formula \(A\) denotes any tuple of local basis matrices in
the ambient rational coordinates. The conjugate transpose uses
both factors at a split pair. A basis change multiplies the
determinant by units, so the value is independent of the basis.
Empty Gram determinants are one. For a rational
\(\gamma\in\mathcal T(c,s,c')\) let \(A_\gamma\) be the image
of \(W_\gamma\cap\widehat\Lambda_c\) under the map in
Definition~\ref{def:global-middle-relation}, extended at each
finite place; \(\widehat\Lambda_c\) is the finite completion
of the row lattice of Definition~\ref{def:unip}. This is a
finite lattice on \(\mathcal L_\gamma\). Indeed kernels,
images and rational spans can be obtained by Gaussian
elimination, and, off the finite set of denominators of the
bases and nonzero pivots and their inverses, the same operations
give local integral bases. At the other places elimination
over the discrete valuation rings gives full lattices by
Lemma~\ref{lem:valuation-index}, whose pivot argument also
applies to rectangular matrices.
\end{definition}

\begin{lemma}\label{lem:common-definite-relation}
Let \(C=\mathcal L_\gamma\cap\mathcal L_\delta\),
\(i=k+1-\dim_E C,\ l=\operatorname{rank}(t|_C)\), and let \(T\)
be the first projection of \(C\cap\ker t\). Then \(T\subset X_0\)
and
\[
          k-\dim_E T=i+l-1,\qquad l\le2 .
\]
On this subspace of pairs the first projection is an isometry
from \(Q\) to the ordinary positive hermitian form on \(T\);
the second projection has the same property. On the whole
\(\mathcal L_\gamma\), \(t\) has rank two if
\(\mathfrak d(\gamma)\ne0\), and rank one if the characters
agree and \(\mathfrak d(\gamma)=0\). In the latter case
\(k-\dim_E T\le i\).
\end{lemma}
\begin{proof}
In \(\ker t=X_0\oplus X_0\), the two ordinary inner products
on pairs of members of \(C\) are equal because \(B\) vanishes
by Proposition~\ref{prop:global-middle-relation}. Their
half-sum is \(Q\); positivity gives injectivity of each
projection. Thus \(\dim T=\dim C-l\), proving the identity.

If \(t\) on \(\mathcal L_\gamma\) has rank below two, the two
functionals \(y\mapsto[z,y]\), \(y\mapsto[\gamma z,y]\) are
linearly dependent on \(W_\gamma\). The double orthogonal of a
subspace for a nondegenerate form equals that subspace by the
dimension computation in Proposition~\ref{prop:global-middle-relation}.
Thus \(\alpha z+\beta\gamma z\in P+\gamma P\) for some
\((\alpha,\beta)\ne(0,0)\), conjugating the coefficients of
the functional relation. Write this as
\(p+\alpha z=\gamma(p'-\beta z)\) with \(p,p'\in P\).
Norms give \(\alpha\bar\alpha=\beta\bar\beta\), so both
coefficients are nonzero. This is a nonzero common frame
vector and Theorem~\ref{thm:rational-frame-discriminant}
gives \(\mathfrak d(\gamma)=0\).
For the agreeing singular case use the graph of
Theorem~\ref{thm:singular-graph-middle}; the two \(z\)
functionals are then proportional and nonzero. Their rank on
a subspace \(C\) is at most one.
\end{proof}

\begin{proposition}\label{prop:middle-lattice-congruence}
If \(\gamma,\delta\in\mathcal T(c,s,c')\), then
\[
                    B(A_\gamma,A_\delta)\ \subset\ g\widehat O .
\]
The induced pairing on
\((\mathcal L_\gamma/C)\times(\mathcal L_\delta/C)\) is
nondegenerate.
\end{proposition}
\begin{proof}
For lifts \(y\in W_\gamma,\ y'\in W_\delta\), the pairing is
\[
 [y,y']-[\gamma^{-1}y,\delta^{-1}y']
                 =[y,(1-\gamma\delta^{-1})y'] .
\]
The first equality uses \(P^\perp=P\oplus M_0\) on both sides,
as in Proposition~\ref{prop:global-middle-relation}. For finite
integral lifts, \(\gamma\delta^{-1}\) preserves the class-\(c\)
lattice and acts as the identity modulo \(g\) because the two
maps from the class-\(c'\) lattice have the same reduction.
At the conjugate split components equality of the inverse
transposes follows from Proposition~\ref{prop:row-group}. The
pairing on the lattice is integral, giving the asserted divisibility.
A member of \(\mathcal L_\gamma\) perpendicular to all of
\(\mathcal L_\delta\) belongs to \(\mathcal L_\delta\), by
Proposition~\ref{prop:global-middle-relation}. The converse
holds for every member of \(C\), and the same argument applies
to the other variable. These are exactly the two radicals.
\end{proof}

\begin{theorem}\label{thm:unequal-relation-determinant}
In Proposition~\ref{prop:middle-lattice-congruence} let
\(A_\gamma^\perp,A_\delta^\perp\) denote the images of the
two lift lattices under the \(Q\)-orthogonal projections
away from \(C\), within their respective relations. Then
\[
 \Delta_Q(A_\gamma^\perp)\Delta_Q(A_\delta^\perp)
                                  \ \ge\ (g/2)^{2fi}.
\]
For either of the two lattices one has the exact identity
\(\Delta_Q(A)=\Delta_Q(A\cap C)\Delta_Q(A^\perp)\), with
intersections and projections taken locally.
\end{theorem}
\begin{proof}
The projection is rational because solving the orthogonality
equations uses the inverse of the rational Gram matrix on
\(C\), which is invertible by positivity even at one infinite
place. The local intersection is saturated: an element of
the lattice whose nonzero scalar multiple lies in the
subspace lies there itself. Diagonal elimination as in
Lemma~\ref{lem:valuation-index} extends a basis of this
intersection to a basis of the full local lattice. The added
columns project to a basis of the image. Subtract their
\(C\)-components over the local fields. This is a triangular
basis transformation of determinant one, and now the Gram
determinant factors, proving the identity.

Choose rational bases \(Y,Y'\) for the two orthogonal
complements of \(C\), with local basis matrices \(YH,Y'H'\)
for the projected lattices. Their dimension is \(i\), and
\(D=\det(Y^*B Y')\) is nonzero by
Proposition~\ref{prop:middle-lattice-congruence}. Subtracting
a \(C\)-component leaves \(B\) on each cross pair unchanged.
Thus all entries of \(H^*Y^*B Y'H'\) belong to \(g\widehat O\).
Taking determinants gives
\[
                g^{2fi}\le I(\det H)\,|\operatorname{Nm}_{E/\Q}D|\,
                                       I(\det H').
\]
Here an integral finite scalar has \(I\ge1\), and \(I(g)=g^{2f}\)
by Definition~\ref{def:middle-lift-disc}.

The rational operator \(J=Q^{-1}B\) is diagonal with entries
\(2\) or \(-2\); thus \(J^*QJ=4Q\). At an infinite place apply
the mixed Gram determinant bound in the proof of
Lemma~\ref{lem:mixed-gram} to \(Y,JY'\) and the positive
diagonal form \(Q\). This gives
\[
 |\tau D|^2\le 4^i\det(\tau(Y^*QY))\det(\tau({Y'}^*QY')) .
\]
Multiplication over the \(f\) chosen embeddings, using each
conjugate pair once, gives
\[
 |\operatorname{Nm}_{E/\Q}D|\le 2^{2fi}
     |\operatorname{Nm}_{E/\Q}\det(Y^*QY)|^{1/2}
     |\operatorname{Nm}_{E/\Q}\det({Y'}^*QY')|^{1/2}.
\]
After restoring the two finite basis determinants, the right
side is \(2^{2fi}\Delta_Q(A_\gamma^\perp)
\Delta_Q(A_\delta^\perp)\). Combining the inequalities proves
the theorem, including \(i=0\) by the empty determinant convention.
\end{proof}

\section{An exact discriminant formula for integral lifts}

\begin{definition}\label{def:arithmetic-relation-indices}
Fix \(\gamma\in\mathcal T(c,s,c')\). If
\(\mathfrak d(\gamma)\ne0\), set \(j=n\). Otherwise assume
character agreement and take the \(j\le\nu\) of
Theorem~\ref{thm:rational-agreeing-levels}. Set \(d=\min(j,\nu)\),
\(h=\nu-d\), \(V=P+\gamma P\), and \(P_O=\sum_{a\le\nu}\widehat O e_a\),
with the sum starting at \(a=1\). For lattices \(L'\subset L\)
of the same rank, let \([L:L']\) denote the cardinality of
the finite quotient. Write \(R_d(\gamma)\) for the
\(\widehat O\)-span of the first \(d\) rows of \(Z(\gamma)\)
and define
\[
 \begin{split}
 s_\gamma&=[\,V_f\cap\widehat\Lambda_c:P_O+\gamma P_O\,],\\
 i_\gamma&=[\,R_\nu(\gamma):R_d(\gamma)\,],\\
 m_\gamma&=
 \begin{cases}
 |\operatorname{Nm}_{E/\Q}D_j|,& j\le\nu,\\
 |\operatorname{Nm}_{E/\Q}D_n|/N^{2f},&j=n .
 \end{cases}
 \end{split}
\]
Here \(V_f\) means scalar extension to the finite adeles.
Use the same \(j\) with the two frames exchanged to define
\(i_{\gamma^{-1}}\); its matrix is \(Z(\gamma)^*\).
All these indices are finite by the following lemma.
\end{definition}

\begin{lemma}\label{lem:outer-saturation-factors}
Let \(C_0=P\cap\gamma P\) and
\(F=(e_1,\ldots,e_\nu,\gamma e_1,\ldots,\gamma e_d)\).
The columns of \(F\) are a basis of \(V\). The projections
of \(C_0\) to the last \(h\) outer coordinates in the original
and transported frames are both isomorphisms onto \(E^h\).
The images of the integral intersection
\(C_{0,f}\cap\widehat\Lambda_c\) have indices \(i_\gamma\)
and \(i_{\gamma^{-1}}\) in the two respective copies of
\(\widehat O^h\). In particular
\[
 [\,V_f\cap\widehat\Lambda_c:F\widehat O^{\,\nu+d}\,]
                            =s_\gamma\,i_{\gamma^{-1}} .
\]
\end{lemma}
\begin{proof}
If \(j=n\), Theorem~\ref{thm:rational-frame-discriminant}
gives \(C_0=0\). Both row indices are then one since \(d=\nu\).
If \(j\le\nu\), the first \(\nu-j\) shared vectors of
Theorem~\ref{thm:rational-agreeing-levels} form a basis of
\(C_0\) by the proof of Theorem~\ref{thm:singular-graph-middle}.
Their two tail matrices are upper triangular with nonzero
diagonal. In either case \(\gamma L_d\) meets \(P\) in zero
because a member of this intersection has zero transported
tail. The remaining transported columns are combinations of
the displayed columns because any transported tail can be
removed with a vector of \(C_0\). This proves the basis and
projection assertions.

For \(h>0\) the first \(d\) rows of \(Z\) are independent
by \(D_d\ne0\) (for \(d=0\) the list is empty). The
coordinates of every \(v\in C_0\) give a conjugate-linear
relation on its first \(\nu\) rows since \(v\perp\gamma L_f\).
Prescribing a tail determines such a relation uniquely:
existence uses the tail isomorphism and uniqueness uses the
independence just proved. For \(t=(t_b)_{d<b\le\nu}\in
\widehat O^h\), the sum of \(\bar t_b\) times the corresponding
rows gives a surjection of additive groups onto
\(R_\nu/R_d\). The kernel consists exactly of the tails whose
unique relation has integral head coefficients. This is exactly
the image of \(C_{0,f}\cap\widehat\Lambda_c\), because the
outer vectors are part of a basis of every class lattice by
Proposition~\ref{prop:canonical-singular-frame} and its
basis argument. This is a full lattice, giving finiteness and
the asserted index. On exchanging the frames the same proof
computes the other index: \(\gamma\widehat\Lambda_{c'}
=\widehat\Lambda_c\) and \(Z(\gamma^{-1})=Z(\gamma)^*\)
by preservation of the form. When \(h=0\) the tail assertions
hold with trivial quotients.

The quotient of \(P_O+\gamma P_O\) by the span of \(F\) is
generated by the last \(h\) transported outer vectors.
An integral combination with tail \(t\) is zero in this
quotient exactly when one can adjust it by an integral
transported head to lie in \(P_O\). This says precisely that
\(t\) is the transported tail of an integral vector in \(C_0\).
The index of the span of \(F\) in \(P_O+\gamma P_O\) is
therefore \(i_{\gamma^{-1}}\). Both are full lattices in \(V\),
so multiplication of the two finite indices proves the formula.
\end{proof}

\begin{lemma}\label{lem:relation-bordered-minor}
In the notation of Lemma~\ref{lem:outer-saturation-factors} let
\(B_1=(e_{d+1}^*,\ldots,e_\nu^*,z,\gamma z)\), with \(h+2\)
columns, and let \(J\) be the rational row form. Write
\(H=[F\ B_1]^*J[F\ B_1]+\operatorname{diag}(0_{\nu+d},NI_{h+2})\).
Then
\[
                |\operatorname{Nm}_{E/\Q}\det H|
                                =(N^{2f}m_\gamma)^2 .
\]
\end{lemma}
\begin{proof}
Reorder the columns to put \((e_1,\ldots,e_\nu,z)\) first and
\((\gamma e_1,\ldots,\gamma e_d,e_{d+1}^*,\ldots,e_\nu^*,
\gamma z)\) last. Apply the same permutation to the rows. The
first diagonal block vanishes: the one nonzero norm \(-N\) of
\(z\) is canceled by the added \(N\). Write \(K\) for the upper
right block. Interchanging the two column blocks then gives a
block triangular matrix, so \(\det H=(-1)^n\det K\,\overline{\det K}\).
The intermediate dual columns of \(K\) are the corresponding
identity columns. When \(j=n\) there are none and \(K=Z\),
giving the formula. In the singular case expansion along these
columns gives the minor of \(Z\) on rows and columns
\(1,\ldots,j,n\). Subtract from its last column the unique
combination of the first \(j\) columns that clears its first
\(j\) entries. The resulting vector of column coefficients,
extended by zeros at \(j<b<n\), is \(P_j^{-1}\varepsilon_n\)
of Proposition~\ref{prop:canonical-singular-frame}. Hence the
last entry is \([z,v_n]=-N\lambda\), and the remaining
determinant is \(-N\lambda D_j\). Since \(\lambda\bar\lambda=1\),
its absolute field norm is \(N^{2f}|\operatorname{Nm}_{E/\Q}D_j|\).
This also proves the assertion when \(j=0\) or \(j=\nu\).
\end{proof}

\begin{theorem}\label{thm:exact-integral-lift-discriminant}
For the positive metric \(Q\) and the entire lattice \(A_\gamma\)
of integral lifts from Definition~\ref{def:middle-lift-disc},
\[
                \Delta_Q(A_\gamma)\,s_\gamma\,i_\gamma\,
                                i_{\gamma^{-1}}=m_\gamma .
\]
\end{theorem}
\begin{proof}
Work in finite local bases of \(\widehat\Lambda_c\), grouping
conjugate places. Its Gram matrix \(J_0\) is integral unimodular
by Proposition~\ref{prop:row-group}. Use \(F,B_1\) also for the
coordinates of the preceding columns in these bases. Take a
basis \(F_0\) of the saturation of \(F\) in the local lattice,
and write \(F=F_0R\). Such bases and their extensions to full
lattice bases are obtained by diagonal elimination over each
valuation ring, as in Lemma~\ref{lem:valuation-index}; saturation
replaces the diagonal elementary divisors by ones. Thus
\(I(\det R)=s_\gamma i_{\gamma^{-1}}\) by
Lemma~\ref{lem:outer-saturation-factors}. The map \(F_0^*J_0\)
from the lattice onto the coordinate module of rank \(\nu+d\)
is surjective: \(F_0\) extends to an invertible integral
matrix and so does its adjoint, and \(J_0\) is unimodular.
Choose integral lifts \(D\) of the standard basis, and a basis
\(W\) of its kernel. Then \([D\ W]\) is an integral basis and
\(W\) spans \(W_{\gamma,f}\cap\widehat\Lambda_c\). All choices
can be made integral in fixed rational coordinates off finitely
many places by the pivot argument of
Definition~\ref{def:middle-lift-disc}.

Compute the determinant of
\[
 \begin{pmatrix}
 J_0&J_0F&J_0B_1\\
 F^*J_0&0&0\\
 B_1^*J_0&0&-NI_{h+2}
 \end{pmatrix}.
\]
Eliminating its first diagonal block gives \(\det J_0\)
times \(\det(-H)\). Eliminating its last diagonal block first
instead replaces \(J_0\) by
\(J'=J_0+N^{-1}J_0 B_1B_1^*J_0\) and produces a factor
\((-N)^{h+2}\). These eliminations subtract row multiples
to clear the entries below or above the pivot block. In the
remaining matrix use the basis \([D\ W]\) on the original
coordinates. Since \(F^*J_0[D\ W]=[R^*\ 0]\), expansion along
the last rows, and then along the matching columns, leaves
\(\det(R^*)\det R\,\det(W^*J'W)\) up to sign. All integral
basis changes and \(\det J_0\) have \(I=1\). Hence
\[
 I(\det H)=I(N)^{h+2}I(\det R)^2 I(\det(W^*J'W)).
\]

For \(y\in W_{\gamma,f}\), write \(a(y)\) for the last \(h\)
coefficients along \(e_{d+1},\ldots,e_\nu\), and \(b(y)\)
for its image in \(\mathcal L_\gamma\). The map
\(y\mapsto(a(y),b(y))\) is an isomorphism onto
\(E_f^h\oplus\mathcal L_{\gamma,f}\). Indeed the kernel of
\(b\) is \(C_{0,f}\), which maps isomorphically to the first
summand by Lemma~\ref{lem:outer-saturation-factors}, and \(b\)
is surjective by definition. In these coordinates \(J'|_{W_\gamma}\)
is \(N^{-1}I_h\oplus Q|_{\mathcal L_\gamma}\). To see this,
the entries of \(B_1^*J_0 y\) are \(a(y)\) and \(t(b(y))\).
The original pairing of two lifts equals the middle pairing
on either of their two coordinates, hence their half-sum.
Adding the displayed correction gives exactly the claimed form.

The image of the integral kernel of \(b\) in the first
summand has index \(i_\gamma\) in \(\widehat O^h\).
Extend its basis by lifts of a basis of \(A_\gamma\).
Their first components can be removed over the fields by
subtracting kernel columns, a shear of determinant one.
The Gram determinant in the two orthogonal summands therefore
gives
\[
           I(\det(W^*J'W))
              =I(N)^{-h}i_\gamma^2\Delta_Q(A_\gamma)^2 .
\]
Using \(I(N)=N^{2f}\) in the preceding determinant identity
and substituting Lemma~\ref{lem:relation-bordered-minor} leaves
the claimed equality after taking positive square roots. This
calculation uses \(I(\overline{\det R})=I(\det R)\) on the
product of all finite places; it includes the split factors.
\end{proof}

\section{A fixed frame vector at a marked prime}

\begin{definition}\label{def:finite-agreeing-reduction}
Call \(s\in\mathcal G\) agreeing if, for every \(u,u'\in
U_f\cap K_Y^0\) whose reductions satisfy \(\bar u s=s\bar u'\),
one has \(\chi_f(u)=\chi_f(u')\), with \(\chi_f\) as in
Proposition~\ref{prop:full-phase}. A set of rational terms is
stable if it is invariant under left and right multiplication
by \(U(F)\). Write \(T_f(x,x')=([z,x],[z,x'])\) for the
finite extension of \(t\) of Definition~\ref{def:middle-lift-disc}.
\end{definition}

\begin{proposition}\label{prop:finite-agreement-support}
At a nonagreeing \(s\), for each \(c,c'\) and each stable set
\(\mathcal B\) of rational terms, the sum
\(\sum_{\delta\in\mathcal T(c,s,c')\cap\mathcal B}M_{c,c'}(\delta)\)
is zero.
\end{proposition}
\begin{proof}
Choose \(u,u'\) as in Definition~\ref{def:finite-agreeing-reduction}
with unequal characters. Lemma~\ref{lem:unip-approx} gives
\(v,v'\in U(F)\) whose finite positions \(A_c(v),A_{c'}(v')\)
lie in \(uK_Y,u'K_Y\), respectively. Indeed the two inverse
images of these cosets are nonempty open sets in the finite
unipotent group: the respective conjugates by \(h_c,h_{c'}\)
supply points in them. The class matrices normalize the
unipotent by preserving the flag and acting invertibly on
its middle quotient. They fix every frame vector and leave
unchanged the coefficients along its outer lines. Conjugation
therefore preserves the simple-root character, as in
Lemma~\ref{lem:global-finite-phase-expansion}. On a principal
integral unipotent the character is one by the same integrality
argument as in Proposition~\ref{prop:full-phase}. Consequently
\(\chi_\infty(v_\infty)\ne\chi_\infty(v'_\infty)\) by that
proposition.

The map \(\delta\mapsto v\delta {v'}^{-1}\) permutes the terms
in the stated sum, since \(\bar u s=s\bar u'\). It permutes
their arithmetic double orbits as well: \(v,v'\) normalize
the principal arithmetic unipotents by normality of \(K_Y\)
in \(K_Y^0\). Apply Lemma~\ref{lem:cut-orbit-unfold} on
each orbit. Its integral is multiplied by
\(\chi_\infty(v_\infty)\overline{\chi_\infty(v'_\infty)}\)
under this permutation, by the substitution \(y\mapsto
v_\infty y\). This carries the old stabilizer domain to a
domain for the conjugate stabilizer. Only finitely many orbit
integrals can be nonzero by Proposition~\ref{prop:actual-cut-poincare-columns}.
Thus the total equals itself times a scalar different from one.
\end{proof}

\begin{lemma}\label{lem:finite-agreeing-flag-step}
Let \(s\) be agreeing. Choose either component \(\mathfrak p\)
of a marked split pair at \(p\in\{g_0,g_1\}\). Work over its
residue field \(\mathbf F_p\) in the fixed integral basis,
write \(S\) for the action of \(s\) there, and let \(L_a\)
denote the span of the first \(a\) reduced frame vectors.
If \(1\le a<n\) and
\[
       \sum_{b\le a}x_b f_b=\sum_{b\le a}y_b S f_b
                 \qquad (x_a=c_1,\ y_a=c_2),
\]
there is such a shared vector at \(a+1\) with the same two
leading coefficients.
\end{lemma}
\begin{proof}
Let \(w\) be the displayed vector and \(m\) any linear
functional annihilating \(L_a+S L_a\).
For \(r\in\mathbf F_p\) the matrix \(I+rw m\) in the
\(\mathfrak p\)-component has conjugate by \(S^{-1}\)
equal to \(I+r(S^{-1}w)(mS)\). These matrices lift on the
two sides to integral finite unipotents, with all other
pairs trivial. For the first lift choose \(W\) in the integral
span of the first \(a\) outer basis vectors with the prescribed
coefficients, and a lift \(M\) of \(m\) annihilating those
basis vectors exactly. Use \(I+\tilde r WM\) over \(O_{E,\mathfrak p}\)
with \(\tilde r\) lifting \(r\). Its inverse is \(I-\tilde r WM\)
since \(M(W)=0\). On the other component take the forced
inverse adjoint, which is integral by unimodularity in
Proposition~\ref{prop:row-group}. In the chosen component
the addition has image in \(L_a\) and is zero on \(L_a\).
Thus it preserves each outer flag, acts identically on its
successive quotients and on \(P^\perp/P\). In the partner
component the addition annihilates all of \(P^\perp\),
since it pairs with \(W\in P\) in the first component.
These are precisely the requirements of Definition~\ref{def:unip}.
In the second lift use the coefficients of \(S^{-1}w\) in
the outer frame and lift \(mS\), which also annihilates \(L_a\).
The resulting reductions intertwine by \(s\) at both components.

In the first component the simple-root sum of the first
lift reduces to \(rc_1 m(f_{a+1})\): the terms with indices
less than \(a\) vanish because \(m(L_a)=0\), and those above
\(a\) because \(w\in L_a\). In the second lift it reduces to
\(rc_2 m(Sf_{a+1})\). The partner simple-root sums vanish
because both partner actions fix the entire frame. By
Proposition~\ref{prop:additive} the character
\(x\mapsto\psi_f(-x/(g\varpi))\) on the chosen integral
component induces a nontrivial additive character of
\(\mathbf F_p\): the self-dual integral lattice there has
annihilator \(O_{E,\mathfrak p}\), and \(g\) is a uniformizer
times a unit. More explicitly the component is \(\Q_p\);
dividing an integral variable by \(g\varpi\) at this split
unramified prime gives conductor \(p\), since \(\varpi\) is
a unit there by Proposition~\ref{prop:different}. All other
components of each simple-root sum are zero. Equality for
every \(r\) implies
\(c_1m(f_{a+1})=c_2m(Sf_{a+1})\), because multiplication
by a nonzero residue permutes the field. Thus the vector
\(c_1 f_{a+1}-c_2 S f_{a+1}\) lies in \(L_a+S L_a\),
by extending a basis and testing all annihilating functionals.
Rearrangement gives the required shared vector.
\end{proof}

\begin{theorem}\label{thm:marked-middle-evaluations}
Suppose \(s\) is agreeing. Let \(\mathcal P\) be a set of
distinct primes of \(E\) in the marked split pairs, possibly
including both partners, such that \(s\) fixes a nonzero
vector in the reduced frame span at each \(\mathfrak p\in
\mathcal P\). Put \(X=\prod_{\mathfrak p\in\mathcal P}
\#(O_E/\mathfrak p)\), with empty product one. If
\(\gamma\in\mathcal T(c,s,c')\), then
\[
 T_f(A_\gamma)\subset\widehat O^2,\qquad
 a_1\equiv a_2\pmod{\bar{\mathfrak p}}\quad
 ((a_1,a_2)\in T_f(A_\gamma),\ \mathfrak p\in\mathcal P).
\]
In particular for any two columns of \(T_f(A_\gamma)\) whose
determinant \(\alpha\) is a finite idele, \(I(\alpha)\ge X\).
\end{theorem}
\begin{proof}
At a prime in \(\mathcal P\) take a fixed nonzero frame
vector \(\sum_{b\le a}x_b f_b=S\sum_{b\le a}x_b f_b\)
of height \(a\). Its two leading coefficients agree and
are nonzero. Iteration of Lemma~\ref{lem:finite-agreeing-flag-step},
or this vector itself if \(a=n\), gives a shared vector at
\(n\) with the same nonzero leading coefficient in the two
frames. Dividing by that coefficient gives
\(z-Sz\in P+SP\) in the residue space. Applying the finite
class frames, which fix all the \(f_b\), this says
\(z-\gamma z\) modulo \(\mathfrak p\widehat\Lambda_c\)
is in the reduction of \(P_O+\gamma P_O\).

For an integral lift \(y\in W_{\gamma,f}\cap\widehat\Lambda_c\)
the two entries of the evaluation are \([z,y]\) and
\([\gamma z,y]\). Both are integral by
Proposition~\ref{prop:row-group}. The lift is orthogonal
to \(P_O+\gamma P_O\). Pairing the congruence in the first
argument therefore gives
\([z,y]-[\gamma z,y]\equiv0\bmod\bar{\mathfrak p}\);
the conjugate prime occurs because the first variable is
conjugate-linear. This proves the inclusions. The two rows
of any evaluation matrix of two columns then agree modulo
each \(\bar{\mathfrak p}\). Its determinant is integral
everywhere and has valuation at least one at each of these
distinct primes. Multiplication of residue cardinalities,
which are unchanged by conjugation, gives the asserted
bound by Definition~\ref{def:middle-lift-disc}.
\end{proof}

\section{Discriminant bounds for pairs of congruent terms}

\begin{definition}\label{def:congruent-relation-gauge}
For a term at a level of Definition~\ref{def:arithmetic-relation-indices}
put
\[
 b_\gamma=\prod_{\tau\in\Phi}|\tau\mathfrak d(\gamma)|,\qquad
 G_\gamma=\frac{\max(b_\gamma,N^f m_\gamma)}
                         {i_\gamma i_{\gamma^{-1}}}.
\]
For a rational subspace \(T\subset X_0\) let \(L_c(T)=
T_f\cap\widehat\Lambda_c\), \(d_c(T)=\Delta_{h_0}(L_c(T))\).
For \(a=1\) or \(g\) write
\[
 \begin{split}
 F_\gamma(v)&=([\gamma f_b,v])_{1\le b\le n},\\
 o_{\gamma,a}(T)&=
 [\,F_\gamma(L_c(T)):F_\gamma(L_c(T))\cap aF_\gamma(P_O)\,],
 \qquad P_\gamma(T)=d_c(T)o_{\gamma,1}(T)s_\gamma .
 \end{split}
\]
The index is used below only when finite. All local coordinates
in these definitions are those in the rational row space. The
positive form on \(X_0\) is the ordinary one of
Definition~\ref{def:row-space}.
\end{definition}

\begin{lemma}\label{lem:middle-evaluation-gram}
On the relation \(\mathcal L_\gamma\) the two coordinates of
\(t\) have representing vectors \(r_1,r_2\) for \(Q\), meaning
\(t(v)=(Q(r_1,v),Q(r_2,v))\). There is \(u\in E\) such that
their Gram matrix is
\[
       N\begin{pmatrix}1&-\bar u\\-u&1\end{pmatrix},\qquad
       1-u\bar u=\frac{\mathfrak d(\gamma)}{D_j\bar D_j}.
\]
When \(j=n\), \(u=ND_\nu/D_n\).
\end{lemma}
\begin{proof}
First construct \(r_1=(a-z,uz)\) in the relation, with
\(a\in X_0\). For the regular level take the vector \(w\)
from the cofactor construction in
Theorem~\ref{thm:rational-frame-discriminant}. Its inverse
image under \(\gamma\) is in \(L_f\), and \(w\perp P\).
Thus \(w\in W_\gamma\). Multiply it by \(N/D_n\). The first
middle coordinate now has \(z\)-coefficient \(-1\), while
the second is \(ND_\nu z/D_n\). The discriminant identity
gives the indicated \(1-u\bar u\). In the singular case use
the vector \(w_n\) of Theorem~\ref{thm:rational-agreeing-levels}
multiplied by \(-\lambda^{-1}\); then \(a=0,\ u=-\lambda^{-1}\)
and both sides of the discriminant identity are zero.
Interchanging the two frames and using
Proposition~\ref{prop:global-middle-relation} gives a member
\(r_2=(u'z,a'-z)\) for some \(a'\in X_0\).
The hypotheses for the inverse follow because
\(Z(\gamma^{-1})=Z(\gamma)^*\), and rational agreement is
symmetric by Definition~\ref{def:rational-character-agreement}.
From \(B(r_1,r_2)=0\) we get \(N(u'-\bar u)=0\).

For any \(v=(x,x')\) in the relation the two terms
\(h_0(a-z,x),h_0(uz,x')\) coincide because \(B(r_1,v)=0\).
Their half-sum in \(Q(r_1,v)\) therefore equals
\(\bar u[z,x']\). The contribution of \(N^{-1}t^*t\) is
\([z,x]-\bar u[z,x']\), giving the first representing
property. Exchanging the slots gives the second one.
Evaluating on \(r_1,r_2\) gives the displayed matrix.
\end{proof}

\begin{proposition}\label{prop:common-evaluation-lattice}
Take two terms in \(\mathcal T(c,s,c')\) at levels as in
Definition~\ref{def:arithmetic-relation-indices}, and denote them by
\(\gamma,\delta\). With the notation of
Lemma~\ref{lem:common-definite-relation} set \(C_1=C\cap\ker t\),
and let \(D\) be the \(Q\)-orthogonal projection of
\(A_\gamma\cap C\) away from \(C_1\) inside \(C\). Assume
\(s,\mathcal P,X\) satisfy Theorem~\ref{thm:marked-middle-evaluations},
and put \(X_l=X\) for \(l=2\) and \(X_l=1\) otherwise. Then
\[
 \begin{split}
 \Delta_Q(A_\gamma\cap C_1)&\ge d_c(T)o_{\gamma,1}(T)>0,\\
 \Delta_Q(D)&\ge
       \frac{X_l\,m_\gamma}{\max(b_\gamma,N^f m_\gamma)} .
 \end{split}
\]
In particular both \(o_{\gamma,1}(T)\) and \(o_{\gamma,g}(T)\)
are finite.
\end{proposition}
\begin{proof}
For any member of \(A_\gamma\cap C_1\), an integral lift
\(y\) has the form \(p+x\) with \(p\in P_O,\ x\in L_c(T)\),
since the outer planes split off integrally in
Definition~\ref{def:unip}. Its second middle entry is also
orthogonal to \(z\), so \(y\perp\gamma L_f\) and
\(F_\gamma(x)=-F_\gamma(p)\). The first projection is a
\(Q\)-isometry on \(C_1\) by Lemma~\ref{lem:common-definite-relation}.
Its image lattice \(L'\) spans \(T\) and lies inside the
preimage of \(F_\gamma(P_O)\) in \(L_c(T)\). Thus
\([L_c(T):L']\ge o_{\gamma,1}(T)\) by the induced quotient
map. This full lattice has finite index by local elimination
and Lemma~\ref{lem:class-covol}. Its discriminant equals
\(d_c(T)[L_c(T):L']\), by taking the determinant of the basis
change in Definition~\ref{def:middle-lift-disc}. This proves
the first assertion. For \(a=g\) the preimage contains
\(gL'\), again of finite index.

A projection away from \(C_1\) preserves \(t\), so \(t(D)\)
is integral at every place and satisfies the congruences of
Theorem~\ref{thm:marked-middle-evaluations}. For \(l=0\) its
source has discriminant one, which suffices since \(N\ge1\).
For \(l=1\) choose a nonzero evaluation coordinate on the
rational span of \(D\). On local basis vectors it gives a
finite idele \(\eta\) with \(I(\eta)\ge1\).
By Lemma~\ref{lem:middle-evaluation-gram} its representing
vector has \(Q\)-norm squared \(N\). Applying the mixed Gram
bound of Lemma~\ref{lem:mixed-gram} over all embeddings,
exactly as in Theorem~\ref{thm:unequal-relation-determinant},
gives \(I(\eta)\le N^f\Delta_Q(D)\). This yields the stated
inequality.

For \(l=2\) both evaluations are independent on the rational
span of \(D\). The determinant of their values on local basis
vectors is an idele \(\eta\) and \(I(\eta)\ge X\) by
Theorem~\ref{thm:marked-middle-evaluations}. This case requires
\(j=n\) by Lemma~\ref{lem:common-definite-relation}. By
Lemma~\ref{lem:middle-evaluation-gram}, the discriminant of
the lattice generated by the two rational representing vectors
in their own span is
\[
 \left|\operatorname{Nm}_{E/\Q}
       \frac{N^2\mathfrak d(\gamma)}{D_n\bar D_n}\right|^{1/2}
                         =\frac{b_\gamma}{m_\gamma}.
\]
The same mixed Gram bound gives \(I(\eta)\le
(b_\gamma/m_\gamma)\Delta_Q(D)\). Indeed in a rational
basis \(Y\) for the span of \(D\) the evaluation determinant
is \(\det((r_1,r_2)^*QY)\) times the determinant of the finite
basis change; this latter factor occurs on both sides after
taking \(I\). Combining the inequalities proves the claim.
\end{proof}

\begin{theorem}\label{thm:actual-congruent-pair-gauge}
For the two terms in Proposition~\ref{prop:common-evaluation-lattice}
one has
\[
 \begin{split}
  \Delta_Q(A_\gamma^\perp)\,P_\gamma(T)X_l&\le G_\gamma,\\
  P_\gamma(T)P_\delta(T) X_l^2&\le
                        (2/g)^{2fi}G_\gamma G_\delta .
 \end{split}
\]
The two flag levels may be different. The number \(i\) is
zero precisely when the relations are equal. For a regular
term with \(i=0\) one has \(l=2\); for each singular term
one has \(k-\dim T\le i\), and in all cases \(k-\dim T\le i+1\).
\end{theorem}
\begin{proof}
Factor the discriminant of \(A_\gamma\) along \(C\), then
factor the discriminant of \(A_\gamma\cap C\) along \(C_1\).
The lattice factorization in
Theorem~\ref{thm:unequal-relation-determinant} applies to
both orthogonal projections. Use both lower bounds of
Proposition~\ref{prop:common-evaluation-lattice} and the
exact value of \(\Delta_Q(A_\gamma)\) from
Theorem~\ref{thm:exact-integral-lift-discriminant}. Cancelling
\(m_\gamma>0\) gives the first inequality. Exchange \(\gamma\)
and \(\delta\). The image \(T\) of \(C_1\), its codimension,
and \(X_l\) are unchanged by this interchange. Multiply
the two bounds and apply the lower bound for the two
complementary discriminants in
Theorem~\ref{thm:unequal-relation-determinant}. This proves
the second inequality. Both relations have the same dimension
\(k+1\) by Proposition~\ref{prop:global-middle-relation}, and
the remaining rank assertions follow from
Lemma~\ref{lem:common-definite-relation}.
\end{proof}

\section{Counting definite subspaces by their discriminants}

\begin{definition}\label{def:coarse-definite-subspace}
For \(u\in G^u\) of Definition~\ref{def:adele} let
\(L_u=u\widehat O^k\) be the finite lattice in \(X_0\), and
write \(L_u^{\rm rat}=L_u\cap E^k\). For a rational subspace
\(T\subset X_0\) put \(d_u^0(T)=\Delta_{h_0}(L_u\cap T_f)\).
Set \(\|v\|_\infty^2=\sum_{\tau\in\Phi}\sum_{a=1}^k
|\tau v_a|^2\) for \(v\in E^k\); this is the Euclidean norm
on the product of the complex spaces at infinity.
\end{definition}

\begin{lemma}\label{lem:coarse-orthogonal-disc}
For every \(T\) one has \(d_u^0(T)=d_u^0(T^{\perp_{h_0}})>0\).
If \(A\) is a finite lattice on an \(E\)-subspace \(V\subset
X_0\) of dimension \(s\), the Euclidean covolume of
\(V\cap A\) in \(\prod_{\tau\in\Phi}\tau V\) equals
\(V_E^s\Delta_{h_0}(A)\), where \(V_E=2^{-f}{\rm rd}_E^f\).
\end{lemma}
\begin{proof}
In an arbitrary rational basis \(Y\) of \(V\), write \(A\)
as \(YH\widehat O^s\). The coefficient lattice has covolume
\(V_E^s I(\det H)\) by Lemma~\ref{lem:class-covol} and
Definition~\ref{def:middle-lift-disc}. At each infinite place
the Jacobian of \(v\mapsto(\tau Y)v\) between the underlying
real Euclidean spaces is \(\det(\tau(Y^*Y))\): in an
orthonormal basis for the image the complex matrix is square,
and its real determinant is its absolute determinant squared,
by replacing each complex scalar by its real two-dimensional
matrix. Multiplying over \(\tau\) gives exactly the stated
discriminant by its definition. The value of \(V_E\) follows
from the integral embedding determinant in
Lemma~\ref{lem:skew-covol}.

Locally \(L_u\) is self-dual for \(h_0\) since its local
defining maps are isometries of the standard form and the
standard lattice is self-dual. The orthogonal projection of
\(L_u\) onto \(T_f\) is the \(h_0\)-dual of \(L_u\cap T_f\).
Indeed the projected lattice pairs integrally with this
intersection. Conversely an integral functional on the
intersection extends to the full lattice by completing an
integral basis as in Theorem~\ref{thm:unequal-relation-determinant}.
Self-duality represents that functional by a member of the
full lattice; its orthogonal projection is the required vector.
At a split pair the functional on one component is represented
in the partner, by the same nondegenerate pairing. In a local
basis with Gram matrix \(H_1\) the dual basis is obtained by
multiplication by \(H_1^{-1}\) and its Gram matrix is
\(H_1^{-1}\); the two discriminants therefore multiply to one.
Factor \(\Delta_{h_0}(L_u)=1\) along \(T^{\perp_{h_0}}\) using
the basis and projection argument in
Theorem~\ref{thm:unequal-relation-determinant}. The projected
factor is \(1/d_u^0(T)\), giving the equality. All these
determinants are positive after taking \(I\), since each
restricted rational form is positive definite at infinity.
\end{proof}

\begin{lemma}\label{lem:definite-orthogonal-normals}
Let \(0<s\le k\) and let \(T\) have codimension \(s\).
There are \(E\)-independent \(v_1,\ldots,v_s\in L_u^{\rm rat}\)
whose common orthogonal space is \(T\), such that
\[
             \left(\prod_{\alpha=1}^s\|v_\alpha\|_\infty\right)^{2f}
                \le (fs\,{\rm rd}_E)^{fs}\,d_u^0(T).
\]
\end{lemma}
\begin{proof}
Take the real Euclidean lattice
\(A=L_u^{\rm rat}\cap T^{\perp_{h_0}}\), of dimension \(D=2fs\)
in its real span, by Lemma~\ref{lem:coarse-orthogonal-disc}.
Choose real independent lattice vectors successively of least
length outside the previous span, possible by discreteness
and full rank. Write \(\lambda_j\) for their lengths, which
are nondecreasing. In their successive real flag take the
orthogonal lines, and scale the \(j\)-th line by \(1/\lambda_j\).
A nonzero lattice vector first belonging to the \(j\)-th
flag space had length at least \(\lambda_j\); hence after
scaling its length is at least one. In the scaled lattice
the open balls of radius \(1/2\) centered at lattice points
are disjoint. The volume of one ball is at most the covolume,
by integrating its periodization on a fundamental
parallelotope. It contains an open cube of side \(1/\sqrt D\).
Thus \(\prod_{j=1}^D\lambda_j\le D^{D/2}\operatorname{covol}(A)\).

Now choose the \(v_\alpha\) successively of least length
outside the \(E\)-span of their predecessors. At step
\(\alpha\) that span at infinity has real dimension
\(2f(\alpha-1)\), so one of the first \(2f(\alpha-1)+1\)
real independent lattice vectors lies outside it. Consequently
\(\|v_\alpha\|_\infty\le\lambda_{2f(\alpha-1)+1}\).
Grouping the nondecreasing \(\lambda_j\) into consecutive
groups of \(2f\) gives the asserted product bound since
\(D^{D/2}V_E^s=(fs\,{\rm rd}_E)^{fs}\).
By construction the vectors span the \(s\)-dimensional
orthogonal complement over \(E\), so their common
orthogonal space is \(T\).
\end{proof}

\begin{theorem}\label{thm:definite-subspace-moment}
For every \(0\le s\le k\), with \(k\ge1\),
\[
      \sum_{\substack{T\subset E^k\\\operatorname{codim}_E T=s}}
                      d_u^0(T)^{-(k+1)}
              \le (25fk\,{\rm rd}_E)^{f(k+1)s}.
\]
The bound is independent of the finite unitary isometry \(u\).
\end{theorem}
\begin{proof}
The summation set is countable since rational bases are finite
lists in a number field. For nonzero \(v\in L_u^{\rm rat}\),
\(h_0(v,v)\) is a totally positive algebraic integer by the
local isometries and Lemma~\ref{lem:denominator-series}.
The product of its \(f\) real conjugates is a positive integer
by the integral multiplication determinant in the fixed
field. Hence at least one conjugate is at least one, and
\(\|v\|_\infty\ge1\). The same gives unit separation of
distinct lattice points. In real dimension \(D_0=2fk\),
disjoint open balls of radius \(1/2\) around the points
of length at most \(R\) fit inside the ball of radius
\(R+1/2\). Thus their count is at most \((1+2R)^{D_0}\).
Put \(p=2f(k+1)>D_0\). Using
\(y^{-p}=p\int_y^\infty R^{-p-1}\dd R\) for \(y\ge1\)
and summing nonnegative terms gives
\[
 \sum_{0\ne v\in L_u^{\rm rat}}\|v\|_\infty^{-p}
 \le p\,3^{D_0}\int_1^\infty R^{D_0-p-1}\dd R
 =(k+1)3^{2fk}\le 25^{f(k+1)} .
\]
For the last inequality use \(k+1\le2^k\), by induction,
and \(2^k9^{fk}\le 18^{fk}\) since \(f\ge1\).

For \(s>0\) choose one ordered normal basis from
Lemma~\ref{lem:definite-orthogonal-normals} for each \(T\).
Distinct subspaces give different tuples because their
orthogonal equations recover \(T\). Raise that lemma's
inequality to power \(k+1\) and divide by the product of
vector lengths to power \(p\). Summing the resulting upper
bounds for \(d_u^0(T)^{-(k+1)}\) is bounded by summing
over all ordered \(s\)-tuples of nonzero vectors. The product
of the \(s\) one-vector sums and \(s\le k\) give the formula.
For \(s=0\) the only subspace is the whole space, whose
finite lattice has discriminant one by
Lemma~\ref{lem:coarse-orthogonal-disc} and the zero-space
convention.
\end{proof}

\section{The subspace discriminant at the column level}

\begin{lemma}\label{lem:marked-definite-intersection}
Let \(u\) represent a class for \(K^0\) in Definition~\ref{def:adele},
and \(a\in K^0\). The class represented by \(u a\) at the column
level has the following lattice on \(X_0\):
\[
 \widehat\Lambda_{u a}\cap X_{0,f}
   =\{v\in L_u:\ e_k^{\mathsf T}a_{\mathfrak p}^{-1}
                   u_{\mathfrak p}^{-1}v_{\mathfrak p}
                        \equiv0\bmod\mathfrak p
          \text{ for every chosen }\mathfrak p\mid N\}.
\]
Here \(e_k^{\mathsf T}\) is the last standard row in the
\(k\) definite coordinates. On the finite right coset space
\(K^0/K_c\) with uniform measure, the covectors
\(e_k^{\mathsf T}a_{\mathfrak p}^{-1}\bmod\mathfrak p\)
are independent and uniformly distributed among the nonzero
covectors of \(\mathbf F_N^k\).
\end{lemma}
\begin{proof}
At the identity class write the middle lattice in the basis
of Definition~\ref{def:row-space}. A linear combination of
\(v_*=x_k/\bar\pi_N+z/N\) and \(z\) with integral coefficients
\(b,c\) has zero \(z\)-coordinate exactly when \(b=-Nc\).
Its resulting \(x_k\)-coordinate is \(-\pi_N c\) because
\(\pi_N\bar\pi_N=N\). Thus on \(X_0\) the first \(k-1\)
coordinates are arbitrary integral elements and the last is
in \(\pi_N\widehat O\). Lemma~\ref{lem:split-level} identifies
this ideal as the unit ideal off the chosen primes above \(N\),
and as the maximal ideal at each of those primes. Applying the
isometry \(u a\), which fixes \(z\), gives the displayed lattice.
At \(N\) the right cosets of \(K_c\) record the nonzero
covectors in precisely these components by
Definition~\ref{def:columns}. In a component \(aK_c\mapsto
e_k^{\mathsf T}a^{-1}\) gives a bijection onto the nonzero
covectors, since its stabilizer is the last-row stabilizer
and every nonzero row extends to a basis. The finite factors
at the two other marked primes are independent factors in the
product of coset spaces from Proposition~\ref{prop:column-dim}.
Counting in the product proves all distribution assertions.
\end{proof}

\begin{theorem}\label{thm:exact-marked-subspace-average}
For \(0\le s\le k\) write
\[
 \theta_s=\frac{N^s-1}{N^k-1}
       +\frac{N^k-N^s}{N^k-1}\,N^{-(k+1)} .
\]
Let \(A_u\) be the finite stabilizer
\(\Gamma^u\cap uK^0u^{-1}\), where \(\Gamma^u\) means the
rational isometries of \(X_0\). Let \(I_c=[K^0:K_c]\).
The exact sum for the marked lattice of
Definition~\ref{def:congruent-relation-gauge} is
\[
 \sum_{c\in\mathcal C(K_c)}\ \sum_{\operatorname{codim}T=s}
                     d_c(T)^{-(k+1)}
   = I_c\,\theta_s^{\,f}
       \sum_{[u]\in\mathcal C(K^0)}\frac1{|A_u|}
           \sum_{\operatorname{codim}T=s}d_u^0(T)^{-(k+1)} .
\]
Both inner sums run over all rational subspaces of \(X_0\)
of the given codimension.
\end{theorem}
\begin{proof}
Fix \(u,T\) on the right. In the \(u_{\mathfrak p}\)-frame
the lattice \(L_u\cap T_f\) reduces at a chosen \(\mathfrak p\mid N\)
to a space of dimension \(k-s\): the intersection with a rational
subspace is saturated, and extension of a local basis as in
Theorem~\ref{thm:unequal-relation-determinant} preserves its rank
on reduction. Among the \(N^k-1\) nonzero covectors, exactly
\(N^s-1\) vanish on this space, by extending its basis over
the residue field. In this case imposing the row congruence
has index one in the local intersection. Every other covector
maps it onto \(\mathbf F_N\) and gives index \(N\). An index
multiplies its discriminant by that index by the finite basis
determinant of Definition~\ref{def:middle-lift-disc}. The mean
of the factor in its inverse \((k+1)\)-st power is therefore
\(\theta_s\); the \(f\) independent factors multiply by
Lemma~\ref{lem:marked-definite-intersection}.

The finer classes above \([u]\) are the orbits of \(A_u\) on
\(K^0/K_c\), acting via conjugation by \(u\) on the left.
Each orbit has \(|A_u|\) elements since all fine stabilizers
are trivial by Proposition~\ref{prop:column-dim}. The inner
sum over \(T\) is constant on any such orbit: a rational
isometry changes \(T\) bijectively and preserves the finite
and infinite hermitian determinants. Thus the sum over the
finer classes is obtained by multiplying the coset mean by
\(I_c/|A_u|\). All terms are nonnegative, so one may exchange
the mean and the subspace sums by increasing finite sums.
Theorem~\ref{thm:definite-subspace-moment} gives their finiteness.
Summing over the finite set of coarse classes proves the formula.
\end{proof}

\begin{corollary}\label{cor:marked-subspace-moment}
Put \(C=|\mathcal C(K_c)|\). For \(k\ge2,\ 0\le s\le k\),
\[
   \sum_{c,T:\,\operatorname{codim}T=s}d_c(T)^{-(k+1)}
      \le C\,(2N^{s-k})^f(25fk\,{\rm rd}_E)^{f(k+1)s}.
\]
At the actual arithmetic scales, for every fixed \(\eta>0\)
and all sufficiently large \(x\), uniformly on the window
and in the integer \(s\), the sum is at most
\[
               C N^{-kf}\, b^{\eta k f(s+1)} .
\]
In particular it is at most \(C N^{-kf} T_0^{s+1}\) for
\(T_0=b^{kf/4}\). At \(s=0\) its exact value is
\(C N^{-(k+1)f}\).
\end{corollary}
\begin{proof}
Since \(k\ge2,\ N\ge2\), one has
\[
 \theta_s\le\frac{N^s+N^{-1}}{N^k-1}
       \le N^{s-k}\frac{1+N^{-1}}{1-N^{-k}}
       \le 2 N^{s-k}.
\]
Counting the coarse fibres in the proof of
Theorem~\ref{thm:exact-marked-subspace-average} gives
\(\sum_{[u]} I_c/|A_u|=C\). That theorem and
Theorem~\ref{thm:definite-subspace-moment} give the first
inequality. For the second take logarithms after removal
of \(C N^{-kf}\) and divide by \(kf(s+1)\log b\). The result
is bounded above by
\[
 \frac{\log 2}{k\log b}+
 \frac{\log N}{k\log b}+
 (1+1/k)\frac{\log(25fk\,{\rm rd}_E)}{\log b}.
\]
By Lemma~\ref{lem:individual-kappas} one has \(\log N\le
60n\log b\) with \(k/n^9\to\infty\). The final quotient tends
to zero since \(f,k\) grow at most polynomially in \(x\) by
Proposition~\ref{prop:thick}, and \(\log{\rm rd}_E=O(1+\log x)
=o(\log b)\) by Proposition~\ref{prop:level-prime}. These
estimates are uniform throughout its prescribed window.
For \(s=0\) the single coarse discriminant is one and
\(\theta_0=N^{-(k+1)}\), giving the exact result.
\end{proof}

\section{Arithmetic representatives with bounded real frames}

\begin{definition}\label{def:boxed-real-frames}
In the scaled real coordinates of Proposition~\ref{prop:unipotent-covol}
let \(F_0\) denote the matrix of the \(n=\nu+1\) standard
frame columns. Use indices \(1,\ldots,n\), so the rational
scales are \(d_{\tau a}=\kappa_\tau^{\,n-a}/\sqrt N\).
For \(y\in G_\infty\) write \(Y_{\tau b}(y)=y_\tau^{-1}F_{0,b}\)
and \(S_\tau(y)=S(y_\tau)\) of
Theorem~\ref{thm:actual-bessel-truncation}. Put
\[
 R=b^{6n+2},\qquad \rho=b^{2n^2},\qquad
 \mathcal E=\{y:S_\tau(y)\le R\ \text{for every }\tau\}.
\]
Let \(\mathcal B\) consist of \(y\in\mathcal E\) with
\(\|Y_{\tau b}(y)\|\le\rho\) for all \(\tau,b\), in the
ordinary Euclidean norm of the scaled coordinates. Denote
by \(\Gamma_c^0\) the subgroup of \(\Gamma_c\) fixing every
rational frame vector; it is the kernel of the action
of \(\Gamma_c\) on the frame.
\end{definition}

\begin{lemma}\label{lem:boxed-arithmetic-shears}
Uniformly in \(c\), for \(1\le i<j\le n\) and \(a\in gNO_E\)
there is an element of \(\Gamma_c\) acting on the rational
frame by \(1+aE_{ij}\). For every \(y\in G_\infty\) and
all sufficiently large arithmetic scales there is \(v\in\Gamma_c\)
such that
\[
     \|Y_{\tau b}(v_\infty^{-1}y)\|\le
          b^{60n+26}S_\tau(y)\qquad\text{for every }\tau,b .
\]
Consequently every \(y\in\mathcal E\) has a \(\Gamma_c\)-translate
in \(\mathcal B\).
\end{lemma}
\begin{proof}
For \(j<n\) use the indicated elementary matrix on \(P\),
its inverse adjoint on the dual and identity on the middle.
The matrix and its inverse preserve the class lattice and are
identity modulo \(g\). For \(j=n\) use the \(i\)-th layer of
Lemma~\ref{lem:heisenberg-layer} with row \(l\) equal to \(a\)
on \(z\) and zero on the other intervening basis vectors.
Its adjoint column there is \((\bar a/N)z\) after the sign
in that lemma, and its top-right entry can be taken to be
\(\eta a\bar a/N\), where \(\eta+\bar\eta=1,\ \eta\in O_E\)
as in Lemma~\ref{lem:middle-hyperbolic-chart}. The trace
condition holds since the form at \(z\) is \(-N\). The
\(z\)-coordinates of integral vectors of the middle lattice
belong to \(N^{-1}\widehat O\) by its basis in
Definition~\ref{def:row-space}, unchanged under the class
isometry on \(X_0\). Thus this row takes values in
\(g\widehat O\), the adjoint column lies in \(g\widehat\Lambda_c\),
and the top-right entry is a multiple of \(g\). This gives
the required principal arithmetic element, inducing precisely
the indicated addition on the frame.

Every tuple \((z_\tau)_{\tau\in\Phi}\in\C^\Phi\) can be
approximated in each coordinate to within
\(f\max(1,|\omega|)\) by a tuple \((\tau a_0)_\tau\)
with \(a_0\in O_E\). Indeed the \(2f\) elements
\(\omega^e\zeta_{N_1}^j\) in the proof of Lemma~\ref{lem:rd}
give a real basis of the embedding space by the nonzero
embedding determinant. Round its real coefficients to nearest
integers. Each basis element has norm at most
\(\max(1,|\omega|)\) at every place, giving the bound.

At each \(\tau\) take \(u_\tau\) in the real unipotent with
\(\|u_\tau y_\tau\|_{\rm op}\le S_\tau(y)+1\) by the infimum
and the compact-middle argument in
Theorem~\ref{thm:actual-bessel-truncation}. Let \(A_\tau\)
be its upper unitriangular matrix on the scaled frame.
Right multiplication of \(A_\tau\) by a rational elementary
shear changes its \(ij\)-entry by
\((d_{\tau j}/d_{\tau i})gN\,\tau a_0\). The other changes
are confined to rows \(<i\) in column \(j\). For each
column in turn, processing the entries above its diagonal
from bottom to top, approximate the negative entry after
division by the nonzero scaled step. Lemma~\ref{lem:individual-kappas}
gives \(|d_{\tau j}/d_{\tau i}|\le1\); also \(N\le b^{60n}\),
\(g\le b^{22}\) and \(f\max(1,|\omega|)\le b\) eventually
by Propositions~\ref{prop:level-prime} and~\ref{prop:thick}.
Thus every final entry of \(A_\tau A(v_\infty)_\tau\) has
absolute value at most \(b^{60n+23}\), with \(v\) the product
of the chosen shears. The frame at \(v_\infty^{-1}y\) equals
the frame at \(u y\) times this product of triangular matrices,
since \(y^{-1}v_\infty F_0=(u y)^{-1}u v_\infty F_0\).
Every column of the frame at \(u y\) has norm at most
\(S_\tau(y)+1\), by equality of the operator norms of a
\(J\)-unitary and its inverse. The triangle inequality,
\(S_\tau(y)\ge1\) and \(2n\le b^3\) give the result.
Finally \(b^{60n+26}R\le\rho\) eventually, and \(S_\tau\)
is unipotent-invariant.
\end{proof}

\begin{proposition}\label{prop:boxed-arithmetic-cutoff}
There are Borel functions \(\omega_c,\beta_c:G_\infty\to[0,1]\)
such that \(\omega_c\) is invariant on the left by the real
unipotents fixing \(F_0\) pointwise, \(\beta_c\le\omega_c\),
and
\[
        \sum_{\gamma\in\Gamma_c}\beta_c(\gamma_\infty y)=1,\qquad
        y\in\mathcal E,\ \omega_c(y)>0\ \Longrightarrow\
                                                y\in\mathcal B .
\]
For any nonnegative Borel function \(H\) invariant under
the full real unipotent \(U_\infty\),
\[
       \int_{G_\infty}\beta_c(y)H(y)\dd y
            =\mathfrak v\int_{U_\infty\backslash G_\infty}
                                        H(y)\dd\bar y .
\]
With these cutoffs, the cut matrix integral summed over
each arithmetic double orbit still has the value in
Lemma~\ref{lem:cut-orbit-unfold}.
\end{proposition}
\begin{proof}
Let \(J(y)\) indicate the frame bounds of
Lemma~\ref{lem:boxed-arithmetic-shears} at \(y\) itself. Set
\[
       m(y)=\sum_{v\in\Gamma_c^0\backslash\Gamma_c}J(v_\infty y).
\]
This is independent of coset representatives and invariant
under \(\Gamma_c\) by right multiplication on the left cosets.
It is at least one by that lemma. For fixed \(y\) finiteness
follows from
\[
                 (v_\infty y)^{-1}F_0
                         =y^{-1}F_0 A(v_\infty)^{-1}.
\]
The frame \(y_\tau^{-1}F_0\) is injective, so its matrix has
a left inverse at each place. Hence the bound in \(J\),
with the fixed \(S_\tau(y)\), bounds every entry of
\(A(v_\infty)^{-1}\). In rational coordinates these entries
are integral since \(v^{-1}\) preserves the class lattice
and the frame is part of its local basis. The scaled entries
differ by fixed nonzero factors across the embeddings. The
embedding lattice \(O_E\) is discrete by Lemma~\ref{lem:rd},
so only finitely many rational frame matrices occur. Their
fibres are exactly the left cosets of \(\Gamma_c^0\).

Put \(\omega_c=J/m\), Borel by the countable sum and
continuity of the frames and \(S_\tau\). Both \(J\) and
each summand of \(m\) are invariant under a real unipotent
that fixes the frame, since it and its conjugate by a real
unipotent fix the frame by Lemma~\ref{lem:bessel-frame-action}.
The sum of \(\omega_c\) over the cosets is one. Choose a Borel
domain of left translates for the discrete group
\(\Gamma_c^0\) on \(G_\infty\) by the construction of
Lemma~\ref{lem:quotient-measure}, and let \(\beta_c\) be its
indicator times \(\omega_c\). Summing first over \(\Gamma_c^0\)
proves the partition assertion. Its support implication follows
from Lemma~\ref{lem:boxed-arithmetic-shears} and the definition
of \(J\).

For fixed \(y\), tile \(U_\infty\) by the left translates
of a domain for \(\Gamma_c\). Summing \(\beta_c\) at the
translates gives
\(\int_{U_\infty}\beta_c(u y)\dd u=\mathfrak v\) by
Proposition~\ref{prop:unipotent-covol}. Integrating over the
unipotent quotient, using Lemmas~\ref{lem:bessel-quotient-density}
and~\ref{lem:product-integration}, gives the stated formula.
For an arithmetic double orbit apply the same two tilings as
in Lemma~\ref{lem:cut-orbit-unfold}, first to the norm.
Its quotient integral there is finite. The two new partitions
sum to one everywhere, so the integrals for the new weights
converge absolutely on that orbit and unfolding gives its
previous value. The same argument for the norm shows that
an orbit of zero absolute mass before the change still has
zero absolute mass.
\end{proof}

\section{Separation of bounded rational frames}

\begin{lemma}\label{lem:boxed-equal-classes}
Let \(c,d,d'\) be column classes, \(s\in\mathcal G\),
\(\gamma\in\mathcal T(c,s,d)\), \(\delta\in\mathcal T(c,s,d')\).
If \(\gamma f_b=\delta f_b\) for all \(1\le b\le n\), then
\(d=d'\) and \(\gamma^{-1}\delta\in\Gamma_d^0\).
\end{lemma}
\begin{proof}
The element \(v=\gamma^{-1}\delta\) has
\(h_d^{-1}v_f h_{d'}\in K_Y\) by Definition~\ref{def:cut-global-columns}.
It fixes the frame, so it stabilizes \(P\subset P^\perp\) and
acts on \(P^\perp/P=M_0\) by an isometry fixing \(z\). This is
the extension of a rational isometry \(a\) of \(X_0=z^\perp\cap M_0\).
Projection of \(P^\perp\cap\widehat\Lambda_d\) on \(M_0\)
gives exactly its full middle lattice, since the two outer summands
split integrally by Definition~\ref{def:row-space}. The induced
map therefore carries the middle lattice for \(d'\) onto that
for \(d\), with principal congruence in representative coordinates
modulo \(g\). The same follows for the inverse. Extending it by
the identity on the outer pairs places \(h_d^{-1}a_f h_{d'}\)
in the column level by Proposition~\ref{prop:marked-local-level}.
The two double classes consequently coincide. For a single class
the rational isometry \(a\) lies in its stabilizer, which is
trivial by Proposition~\ref{prop:column-dim}. The action of
\(v\) on the flag and its middle quotient is now the identity,
including every outer graded line. Hence \(v\in U(F)\) by
Definition~\ref{def:unip}. Its finite position lies in \(K_Y\)
and it fixes the frame, so \(v\in\Gamma_d^0\) by
Definition~\ref{def:boxed-real-frames}.
\end{proof}

\begin{proposition}\label{prop:boxed-frame-separation}
In the situation of Lemma~\ref{lem:boxed-equal-classes},
suppose \(y\in G_\infty\) and both \(\gamma_\infty^{-1}y\)
and \(\delta_\infty^{-1}y\) belong to \(\mathcal B\).
Unless \(d=d'\) and \(\gamma^{-1}\delta\in\Gamma_d^0\),
\[
 \max_{\tau,b}\big\|Y_{\tau b}(\gamma_\infty^{-1}y)
                 -Y_{\tau b}(\delta_\infty^{-1}y)\big\|
                   \ \ge\ \varepsilon_0,\qquad
                      \varepsilon_0=(N\rho)^{-1}.
\]
This holds uniformly for all sufficiently large arithmetic scales.
\end{proposition}
\begin{proof}
Write \(z_0=\gamma_\infty^{-1}y,\ v=\gamma^{-1}\delta\) and
\(\Delta_{\tau b}=z_{0,\tau}^{-1}(v_{\infty,\tau}-1)F_{0,b}\).
Suppose the displayed maximum is smaller than \(\varepsilon_0\).
Both \(v f_b\) and \(f_b\) belong to the integral row lattice for
\(d\), since \(v_f\) carries the lattice for \(d'\) onto it.
Thus every \([f_a,(v-1)f_b]\) belongs to \(O_E\) by the basis
in Definition~\ref{def:row-space}. The scaled hermitian matrix
has operator norm one: it exchanges the outer paired coordinates
and acts by signs on the middle. Invariance under \(z_{0,\tau}\)
and the scales of Definition~\ref{def:boxed-real-frames} imply
\[
 |\tau[f_a,(v-1)f_b]|
 \le |d_{\tau a}d_{\tau b}|^{-1}
             \|Y_{\tau a}(z_0)\|\,\|\Delta_{\tau b}\|
 < N\rho\varepsilon_0=1 .
\]
Here \(|\kappa_\tau|\ge1\) by Lemma~\ref{lem:individual-kappas}.
An algebraic integer with this bound at every CM embedding is zero:
the absolute value of its rational norm is the product of the
squared moduli and, for a nonzero integer of \(E\), is at least
one by Lemma~\ref{lem:rd}. Hence \(w_b=(v-1)f_b\) is orthogonal
to \(P\) and \(z\), and so belongs to \(P\oplus X_0\).
Also \([w_b,w_b]=0\) by \([v f_b,v f_b]=[f_b,f_b]\) and
\([w_b,f_b]=0\). Positivity on \(X_0\) gives \(w_b\in P\).

Write \(w_b=\sum_i a_{ib}e_i\). Each \(a_{ib}\in O_E\) because
the outer summand has the standard lattice in every class.
For any nonzero column choose its largest index \(i\) with
\(a_{ib}\ne0\). Choose \(u_\tau\in\mathcal U\) with
\(\|u_\tau z_{0,\tau}\|_{\rm op}\le R+1\) as in
Theorem~\ref{thm:actual-bessel-truncation}. By the upper
unitriangular frame action of Lemma~\ref{lem:bessel-frame-action},
the \(i\)-th coordinate of \(u_\tau z_{0,\tau}\Delta_{\tau b}\)
is \((d_{\tau b}/d_{\tau i})\tau a_{ib}\). Therefore
\[
    |\tau a_{ib}|\le b^{23n}(R+1)\varepsilon_0<1
\]
at every embedding. The first inequality uses
\(|\kappa_\tau|\le b^{23}\); the second holds eventually since
\(n\to\infty,\ R=b^{6n+2},\ \rho=b^{2n^2}\), and \(N\ge1\)
by Lemma~\ref{lem:individual-kappas} and Proposition~\ref{prop:thick}.
The norm argument gives \(a_{ib}=0\), a contradiction. Every
column is therefore fixed by \(v\), and Lemma~\ref{lem:boxed-equal-classes}
applies.
\end{proof}

\begin{theorem}\label{thm:boxed-weighted-frame-count}
Fix \(c,s\), two real points \(y,y'\in G_\infty\), and
\(x\ge0\). Take distinct pairs \((d_p,\delta_p)\) of a column
class and a rational group element, with weights \(w_p\ge0\).
For each positive weight suppose there is
\(\gamma_p\in\mathcal T(c,s,d_p)\) with
\(\delta_p\in\gamma_p\Gamma_{d_p}\),
\(\gamma_{p,\infty}^{-1}y\in\mathcal B\), and
\(w_p\le x\beta_{d_p}(\delta_{p,\infty}^{-1}y')\), using
Proposition~\ref{prop:boxed-arithmetic-cutoff}. Suppose all
their frames \((Y_{\tau b}(\gamma_{p,\infty}^{-1}y))_{\tau,b}\)
belong to a single translate of a real subspace of dimension
\(t\). Then, uniformly for all sufficiently large scales,
\[
                         \sum_p w_p\le b^{6n^2t}x .
\]
\end{theorem}
\begin{proof}
Use the maximum of the Euclidean column norms. If any positive
weight occurs, fix its frame as center. Every frame is within
\(2\rho\) of this one. By Proposition~\ref{prop:boxed-frame-separation}
distinct frames have distance at least \(\varepsilon_0\), since
membership of \(\gamma_p^{-1}\gamma_q\) in the indicated
pointwise stabilizer gives identical frames. In the affine
\(t\)-plane take disjoint open norm balls of radius
\(\varepsilon_0/2\) about any finite set of distinct frames.
They fit inside the ball of radius \(2\rho+\varepsilon_0/2\).
Lebesgue measure in linear coordinates scales as the \(t\)-th
power of the radius by change of variables. The unit ball has
finite positive measure since the coordinate and the column
norms are equivalent in finite dimension, as follows by taking
the maximum and positive minimum on a Euclidean unit sphere.
Thus there are at most \((1+4N\rho^2)^t\le b^{6n^2t}\) distinct
frames, using \(N\le b^{60n}\) in Lemma~\ref{lem:individual-kappas}.
For \(t=0\) there is at most one.

In the fibre of any frame all column classes equal \(d\) by
Lemma~\ref{lem:boxed-equal-classes}, and, fixing one representative
\(\gamma_*\), every \(\delta_p\) belongs to
\(\gamma_*\Gamma_d\). Distinct pairs in this fibre give distinct
elements of that coset. Their weights sum to at most
\(x\sum_{u\in\Gamma_d}\beta_d(u_\infty^{-1}\gamma_{*,\infty}^{-1}y')=x\)
by Proposition~\ref{prop:boxed-arithmetic-cutoff}. Adding over the
frames proves the bound, including the empty family.
\end{proof}

\section{Full Poincar\'e columns and the error of archimedean truncation}

\begin{definition}\label{def:full-arithmetic-amplitudes}
Let \(F_\infty(y)=\bigotimes_{\tau\in\Phi}F_\tau(y_\tau)\) for
the normalized vectors of Theorem~\ref{thm:bessel-orthonormal-schur}.
Write \(V_\infty\) as in Definition~\ref{def:cut-global-columns}.
For \(W=F_\infty\) or \(V_\infty\) set \(K^W_\delta(y)=
W(y)W(\delta_\infty^{-1}y)^*\) and \(z_\delta=\delta_\infty^{-1}y\).
For any set \(B\) of rational group elements put
\[
 A^W_B(c,s,d)=
  \sum_{\delta\in B\cap\mathcal T(c,s,d)}
     \int_{G_\infty}\beta_c(y)\beta_d(z_\delta)K^W_\delta(y)\dd y
\]
whenever the integrals and sum converge absolutely in the
Hilbert--Schmidt norm. Here and below \(\beta_c\) is chosen by
Proposition~\ref{prop:boxed-arithmetic-cutoff}. For a family of
matrices \(A(c,s,d)\) write
\[
 \mathcal E_\Omega(A)=\mathfrak v^{-2}
       \sum_{c,d}\sum_{s\in\Omega}\|A(c,s,d)\|_{\rm HS}^2
                \qquad(\Omega\subset\mathcal G).
\]
\end{definition}

\begin{proposition}\label{prop:full-frame-majorant}
Let \(A_\tau,\lambda_\tau\) be as in
Theorem~\ref{thm:actual-bessel-truncation}. Put
\(A(y)=\prod_\tau A_\tau(y_\tau)\), \(\lambda=\bigotimes_\tau\lambda_\tau\).
For all sufficiently large scales and all \(c,s,y\),
\[
 \sum_d\sum_{\delta\in\mathcal T(c,s,d)}
       \beta_d(z_\delta)A(z_\delta)
       \le b^{12rn^3f}\lambda(G_\infty).
\]
The integral of the left side against \(\beta_c(y)A(y)\dd y\)
is at most \(\mathfrak v J_c\, b^{(12rn^3+2r^3)f}\).
\end{proposition}
\begin{proof}
The upper bound in Theorem~\ref{thm:actual-bessel-truncation} and
the rectangle formula for the product measure give
\(A(z_\delta)\le\int{\bf1}_{\mathcal E}(z_\delta h)\dd\lambda(h)\).
Fix \(h\). Each nonzero summand of weight
\(w_{d,\delta}=\beta_d(z_\delta){\bf1}_{\mathcal E}(z_\delta h)\)
has a representative \(u\in\Gamma_d\) with
\(u_\infty z_\delta h\in\mathcal B\) by
Lemma~\ref{lem:boxed-arithmetic-shears}. The rational matrix
\(\gamma=\delta u^{-1}\) belongs to \(\mathcal T(c,s,d)\)
and has \(\gamma_\infty^{-1}yh\in\mathcal B\).
Apply Theorem~\ref{thm:boxed-weighted-frame-count} at the two
points \(yh,y\), with \(x=1\) and the whole real frame space.
There are \(fn\) complex columns of length \(r\), so its dimension
is \(2frn\). This bounds the sum for this \(h\) by \(b^{12rn^3f}\).
The weights themselves are measurable in \(h\). Integrating the
pointwise bound and using Lemma~\ref{lem:integral-tools} gives
the first assertion. For the second,
Proposition~\ref{prop:boxed-arithmetic-cutoff} integrates
\(\beta_c A\) to \(\mathfrak v\prod_\tau\int A_\tau\dd\bar y_\tau\)
since \(A\) is unipotent invariant. Multiply by the mass of
\(\lambda\) and use the integral bound of
Theorem~\ref{thm:actual-bessel-truncation} in each factor.
The product of the dimensions there is \(J_c\) by
Definition~\ref{def:two-energy-numbers}.
\end{proof}

\begin{theorem}\label{thm:full-global-truncation}
The amplitudes of Definition~\ref{def:full-arithmetic-amplitudes}
converge absolutely. Uniformly in \(c,s\), for all sufficiently
large scales,
\[
 \sum_d\sum_{\delta\in\mathcal T(c,s,d)}
 \int \beta_c(y)\beta_d(z_\delta)
       \|K^{F_\infty}_\delta(y)-K^{V_\infty}_\delta(y)\|_{\rm HS}\dd y
              \le \tfrac12\mathfrak v b^{-k^3f}.
\]
For each fixed \(D>0\), uniformly in \(B,\Omega\), eventually
\[
   \mathcal E_\Omega(A^{F_\infty}_B-A^{V_\infty}_B)\le E_1 b^{-D}.
\]
\end{theorem}
\begin{proof}
For column vectors \(a,b'\), summing the squared absolute values
of the entries gives \(\|a{b'}^*\|_{\rm HS}=\|a\|\,\|b'\|\).
The norm of a pure tensor is the product of its factor norms by
the same coordinate calculation. Each kernel is bounded by
\(A(y)A(z_\delta)\), which gives absolute convergence by
Proposition~\ref{prop:full-frame-majorant}. If the kernels differ
then at least one real place has different full and cut vectors
at one of \(y,z_\delta\). At such an argument both local norms
are at most \(b^{-3r^3f}A_\tau\) by
Theorem~\ref{thm:actual-bessel-truncation}. Thus the norm of the
difference is at most \(2b^{-3r^3f}A(y)A(z_\delta)\), everywhere.
The integrated bound becomes
\[
              2\mathfrak v J_c b^{(12rn^3-r^3)f}.
\]
Use \(J_c\le b^{2kf}\) from Lemma~\ref{lem:logarithmic-counts}.
One has \(r=k+2n-1,\ k/n^9\to\infty,\ n\to\infty\) by
Proposition~\ref{prop:thick}. Thus
\(r^3-k^3=(2n-1)(r^2+rk+k^2)\) exceeds \(12rn^3+2k+1\)
eventually. This gives the first bound since \(b^f\ge4\) eventually.

For fixed \(c,s\) the sum over \(d\) of the squared norms
of the differences is bounded by the square of the sum of
the norms, hence by \(\mathfrak v^2 b^{-2k^3 f}/4\).
There are \(C\) choices of \(c\). The finite level groups embed
in the product of the \(f\) matrix spaces of rank \(r\) over
each of \(\mathbf F_{g_0},\mathbf F_{g_1}\) by
Proposition~\ref{prop:row-group}, so
\(|\mathcal G|\le g^{fr^2}\le b^{22fr^2}\) by
Lemma~\ref{lem:individual-kappas} and the level-prime bounds.
Since \(E_1=|\bar S|CJ_c\ge C\) by
Definition~\ref{def:two-energy-numbers}, the claim follows from
\((2k^3-22r^2)f\to\infty\).
\end{proof}

\begin{theorem}\label{thm:actual-full-poincare-columns}
Periodize each coordinate of the seed
\({\bf1}_{h_cK_Y}\beta_cF_\infty\) by the formula of
Theorem~\ref{thm:poincare}. This gives bounded square-integrable
automorphic columns. Their matrix of translated inner products
at reduction \(s\) is \(A^{F_\infty}_{\Gamma}(c,s,d)\), with
counting measure on \(G_f/K_Y\). The full columns are reproduced
by right convolution with the product of the kernels \(e_{J,\tau}\)
of Theorem~\ref{thm:actual-real-bessel-vector}.
\end{theorem}
\begin{proof}
Fix a finite argument \(hK_Y\). If no rational element takes
it to \(h_cK_Y\), the series vanishes. Otherwise fix one such
\(\gamma_0\); all the other elements are \(v\gamma_0\) with
\(v_f h_cK_Y=h_cK_Y\). These \(v\) constitute
\(\mathcal T(c,1,c)\), a group closed under inversion. The sum
of the norms of the seed terms is bounded uniformly in \(h,y\)
by Proposition~\ref{prop:full-frame-majorant} at
\(\gamma_{0,\infty}y\). Absolute convergence gives measurability
and left rational invariance by reindexing. To obtain the
square integral, periodize the norm of the seed, expand its square
on a domain as a nonnegative double sum, and unfold the first
index by Lemma~\ref{lem:quotient-measure}. The seed has support
in a single finite component. The resulting single sum of norm
overlaps is precisely the \(d=c,s=1\) subfamily in the second
bound of Proposition~\ref{prop:full-frame-majorant}, with each
\(A\) replaced by the smaller \(\|F_\infty\|\). Thus the
periodization belongs to \(L^2\). The measurable-seed assertion
of Theorem~\ref{thm:poincare}, with absolute integrability from
Theorem~\ref{thm:full-global-truncation}, gives the stated matrix.

Group the original series by left \(\Gamma_c\)-cosets of the
rational elements taking \(hK_Y\) to \(h_cK_Y\). The character
is trivial on \(\Gamma_c\) by Proposition~\ref{prop:full-phase},
and the cutoff sums to one by
Proposition~\ref{prop:boxed-arithmetic-cutoff}. Thus each
coset gives the corresponding full vector, and its norm equals
the sum of the seed norms on that coset. The sum of these norms
is uniformly bounded by the preceding argument even with any
real right translate of \(y\). The local kernels are integrable
by Corollary~\ref{cor:row-polar-moment}. Multiplying their
product by this uniform bound permits interchange of convolution
and the absolutely summable coset series by
Lemmas~\ref{lem:integral-tools} and~\ref{lem:product-integration}.
The local reproduction identity in
Theorem~\ref{thm:actual-real-bessel-vector} applies also after
normalization, hence to every coset tensor. This proves reproduction.
\end{proof}

\begin{corollary}\label{cor:full-bessel-identity}
Let \(S(F)=U(F)H(F)\), where \(H(F)\) consists of the rational
isometries of \(X_0\) extended by the identity on \(z\) and
the outer pairs. Then
\[
          A^{F_\infty}_{S(F)}(c,1,d)
                       ={\bf1}_{c=d}\,\mathfrak v I .
\]
For either choice of \(W\), the amplitudes on a set stable
under rational unipotents on both sides vanish at every
nonagreeing finite reduction.
\end{corollary}
\begin{proof}
If \(\delta=u a\in S(F)\) belongs to \(\mathcal T(c,1,d)\),
the map it induces on \(P^\perp/P\) is \(a\). The integral
projection and stabilizer argument of Lemma~\ref{lem:boxed-equal-classes}
gives \(c=d\), \(a=1\). The terms for \(c=d\) are exactly
\(\Gamma_c\). Summing the second cutoff and using triviality
of its phase gives \(\int\beta_c(y)F_\infty(y)F_\infty(y)^*\dd y\).
Apply Proposition~\ref{prop:boxed-arithmetic-cutoff} entrywise,
by decomposition into real and imaginary parts, which are integrable
by Theorem~\ref{thm:bessel-orthonormal-schur}. The last theorem
gives the identity endomorphism in each factor after quotient
integration. This proves the formula.

For finite disagreement take the two rational unipotents used
in Proposition~\ref{prop:finite-agreement-support}. On each
arithmetic double orbit unfold the two cutoffs on their left
domains as in Lemma~\ref{lem:cut-orbit-unfold}. For the full
kernel the identical tiling is justified first for its norm,
whose sum of weighted integrals is finite by
Theorem~\ref{thm:full-global-truncation}. The rational
unipotents permute the double orbits and multiply their
integrals by the same character
value different from \(1\). Absolute convergence permits summing
over the permutation. Thus the total vanishes.
\end{proof}

\section{Archimedean sizes of the arithmetic frame coordinates}

\begin{definition}\label{def:near-labelled-terms}
Call a rational element \(\gamma\) near if there exists
\(y\in\mathcal B\) with \(\gamma_\infty^{-1}y\in\mathcal B\).
For \(\gamma\in\mathcal T(c,s,c')\) choose one representative of
each double coset \(\Gamma_c^0\gamma\Gamma_{c'}^0\), fixing the
choice for each pair \(c,c'\). Nearness is constant on these
cosets: left multiplication of a real argument by an element
of \(\mathcal U_0\) preserves its frame and its size \(S\) by
Definition~\ref{def:boxed-real-frames} and Lemma~\ref{lem:bessel-frame-action}.
Fix also a section \(r\) of the finite set of left cosets
\(\mathcal G/\bar S\). The label of \(s\) is
\(t=r(s\bar S)^{-1}s\in\bar S\). At level \(j\) consider the
near chosen representatives of the levels in
Definition~\ref{def:arithmetic-relation-indices}. Thus \(j=n\)
means \(\mathfrak d\ne0\), and at a singular level \(0\le j<n\)
we use the agreeing flags of Theorem~\ref{thm:rational-agreeing-levels}.
\end{definition}

\begin{lemma}\label{lem:near-volume-inverse}
For a place \(\tau\) and \(y\in\mathcal B\) the matrix
\(M=y_\tau^{-1}F_0\) has operator norm at most \(\sqrt n\,\rho\).
On its range its inverse into \(\C^n\) has operator norm at most
\[
                       (R+1)^n(\sqrt n\,\rho)^{n-1}.
\]
\end{lemma}
\begin{proof}
The sum of the squared column norms is at most \(n\rho^2\).
Cauchy--Schwarz gives the first assertion. Choose \(u\in\mathcal U\)
with \(\|uy_\tau\|_{\rm op}\le R+1\) by
Theorem~\ref{thm:actual-bessel-truncation}. The map
\(M'=(uy_\tau)^{-1}F_0\) satisfies
\(\|M'a\|\ge \|a\|/(R+1)\), since \(F_0\) is an isometric
inclusion in the Euclidean norm. By Lemma~\ref{lem:bessel-frame-action},
\(M'=M A(u)^{-1}\) with \(\det A(u)=1\). Consequently
\(\det(M^*M)=\det({M'}^*M')\ge(R+1)^{-2n}\).
Indeed the positive Gram matrix on the right has all its eigenvalues
at least \((R+1)^{-2}\) by the norm inequality. A positive
hermitian matrix is unitarily diagonalizable: maximize its
quadratic form on the unit sphere to obtain an eigenvector
by differentiation in orthogonal directions and apply induction
on the invariant orthogonal complement. The square roots of
the Gram eigenvalues of \(M\) have product at least
\((R+1)^{-n}\) and each is at most \(\sqrt n\,\rho\).
This bounds the smallest below by
\((R+1)^{-n}(\sqrt n\,\rho)^{1-n}\) and proves the claim.
\end{proof}

\begin{proposition}\label{prop:near-minor-coordinate-bounds}
Uniformly for near arithmetic terms, at all sufficiently large
arithmetic scales and for all \(\tau\in\Phi,\ a,b,\ 0\le j\le n\),
\[
       |\tau Z_{ab}(\gamma)|\le b^{5n^2},\qquad
       |\tau D_j(\gamma)|^2\le b^{11n^3}.
\]
At singular level \(j<n\), the coordinates of \(D_jv_b\)
in the original frame, for every \(b>j\) in
Proposition~\ref{prop:canonical-singular-frame}, have modulus
at most \(b^{10n^3}\) at each embedding. The same bound holds
for these vectors after multiplication of \(\gamma\) on the
left by any rational unipotent fixing the original frame pointwise.
\end{proposition}
\begin{proof}
Take \(y,\gamma_\infty^{-1}y\in\mathcal B\). Pair their
scaled frame columns as in Proposition~\ref{prop:boxed-frame-separation}.
This gives \(|\tau Z_{ab}|\le N\rho^2\le b^{5n^2}\) eventually,
since \(\log N\le60n\log b\) by Lemma~\ref{lem:individual-kappas}.
Expansion in permutations gives
\(|\tau D_j|\le j! b^{5jn^2}\); \(j\le n\) and
\(\log(n!)\le n\log n=o(n^3\log b)\) give the asserted square
bound, uniformly and including \(j=0\).

Write \(a'\) and \(a\) for the coordinate columns of \(D_jv_b\)
in the transported and original frames. By
Proposition~\ref{prop:canonical-singular-frame},
\(a'=\operatorname{adj}(P_j)\varepsilon_b\).
The same permutation estimate bounds the modulus of every
conjugate entry of \(a'\) by \(b^{6n^3}\).
For each \(\tau\), put \(M=y_\tau^{-1}F_0\) and
\(M'=(\gamma_\infty^{-1}y)_\tau^{-1}F_0\).
The equality of the two rational expressions for \(D_jv_b\)
gives, in scaled coordinates,
\[
      M(\tau a_i/d_{\tau i})_i=M'(\tau a'_i/d_{\tau i})_i .
\]
Lemma~\ref{lem:near-volume-inverse} applied to \(M\) and \(M'\)
therefore bounds each \(|\tau a_i|\) by
\[
 b^{23n}\,n\rho\,(R+1)^n(\sqrt n\,\rho)^{n-1} b^{6n^3}.
\]
Here \(b^{23n}\) bounds every ratio \(|d_{\tau i}/d_{\tau l}|\)
by Lemma~\ref{lem:individual-kappas}. The logarithm to base \(b\)
of this expression is \(8n^3+O(n^2+n\log n)\) since
\(R=b^{6n+2},\ \rho=b^{2n^2}\). This proves the bound.
A left unipotent fixing every \(f_i\) preserves \(Z\) and
fixes \(v_b\), since the latter belongs to their span. Thus
it preserves all the coordinates in question.
\end{proof}

\begin{corollary}\label{cor:near-integral-colours}
Among near arithmetic terms there are at most \(b^{12fn^4}\)
different matrices \(Z\), eventually. Put \(M_b=\lceil3b^{10n^3}\rceil\).
If two such terms of singular level \(j\) have the same \(Z\)
and, after possibly different left frame-fixing unipotents,
the outer coordinates of \(D_jv_b\) for the two frames agree
modulo \(M_b O_E\) for \(b>j\), then all these outer coordinates
are equal over \(E\).
\end{corollary}
\begin{proof}
Two distinct elements of \(O_E\) have distance at least one in
the maximum of their conjugate moduli, by the norm argument of
Proposition~\ref{prop:boxed-frame-separation}. The real dimension
of the space of conjugates is \(2f\). The disjoint-ball argument
of Theorem~\ref{thm:boxed-weighted-frame-count}, now centered
at zero with radius \(b^{5n^2}\), bounds the number of choices
for one entry by \((1+2b^{5n^2})^{2f}\).
There are \(n^2\) entries. Integrality from
Theorem~\ref{thm:rational-frame-discriminant} and the first
bound of Proposition~\ref{prop:near-minor-coordinate-bounds}
prove the first assertion. For the second, each difference
is \(M_b a\) with \(a\in O_E\), and all conjugate moduli of
\(a\) are at most \(2b^{10n^3}/M_b<1\) by the same proposition.
Its norm forces \(a=0\).
\end{proof}

\section{Counting integral maps by their finite images}

\begin{lemma}\label{lem:finite-image-moment}
Let \(R_0=\widehat O\), let \(K\subset R_0^h\) be an \(R_0\)-submodule
containing \(aR_0^h\) for a nonzero integral finite idele \(a\), and
write \(Q_0=R_0^h/K\). Let \(L=L_1\oplus L_2\) be a free
\(R_0\)-module with the two summands free of ranks \(e,v\), respectively.
Consider an indexed family of linear maps \(Y_p:L\to R_0^h\)
and integers \(m_p\ge1\) divisible by
\(|(Y_p(L_1)+K)/K|\). Suppose that at most \(A_0\) indices induce
each particular map \(L\to Q_0\). For every integer \(\ell\ge e\),
\[
            \sum_p m_p^{-\ell}
              \le A_0\,|Q_0|^{h+v}
              \le A_0\,I(a)^{h(h+v)} .
\]
\end{lemma}
\begin{proof}
Every submodule \(H\subset Q_0\) is generated over \(R_0\)
by at most \(h\) elements. To see this, work at the finitely many
primes of \(a\). The inverse image of \(H\) in each local free
module contains \(aO_{E,\mathfrak p}^h\); diagonal elimination
of Lemma~\ref{lem:valuation-index} provides \(h\) generators
there. More explicitly the inverse image is finitely generated
by lifts of its finite quotient together with the \(h\) generators
of \(aO_{E,\mathfrak p}^h\). Projection onto each prime factor
is an \(R_0\)-idempotent, so the submodule is the product of its
local components. Combine their generator lists coordinatewise,
padding by zeros. It follows that there are at most \(|Q_0|^h\)
possible \(H\), by counting ordered \(h\)-tuples. For a fixed image
\(H\) of \(L_1\), there are at most \(|H|^e|Q_0|^v\) quotient maps
from \(L\): record the values on bases of its two summands.
Every such map occurs at most \(A_0\) times, and every associated
weight is at most \(|H|^{-\ell}\). Since \(\ell\ge e\), the contribution
per image is at most \(A_0|Q_0|^v\). Finally \(Q_0\) is a quotient of
\(R_0^h/aR_0^h\), which has order \(I(a)^h\) by
Definition~\ref{def:middle-lift-disc}.
\end{proof}

\begin{definition}\label{def:middle-frame-congruence}
Fix \(c\), a level \(j\), and an occurring Gram matrix \(Z\)
as in Definition~\ref{def:near-labelled-terms}.
Decompose the rational matrix of the transported frame as
\[
       (\gamma f_1,\ldots,\gamma f_n)=
                \begin{pmatrix} A_\gamma^{\rm top}\\ X_\gamma\\
                                       t_Z\\ P_Z\end{pmatrix}
            \quad\text{in }P,\ X_0,\ Ez,\ P^* .
\]
Here \(P_Z=Z_{1:\nu,*}\), \(t_Z=-Z_{n,*}/N\).
Write \(Y=P_j\) as in Proposition~\ref{prop:canonical-singular-frame},
also using \(Y=Z\) at \(j=n\), and set \(d=\min(j,\nu)\),
\(D=\det Y=D_j\ne0\). Put \(L_M=M_{0,f}\cap\widehat\Lambda_c\),
\(L_X=X_{0,f}\cap\widehat\Lambda_c\), and define
\[
 K_Z=gP_Z^*\widehat O^\nu+gY^*\widehat O^n
                          \ \subset\widehat O^n .
\]
Since \(Z\) is integral by Theorem~\ref{thm:rational-frame-discriminant},
the cofactor identity gives \(g\bar D\widehat O^n\subset K_Z\).
For a term \(\gamma\in\mathcal T(c,s,c')\) use the integral map
\(L_M\longrightarrow\widehat O^n,\ v\mapsto F_\gamma(v)\)
from Definition~\ref{def:congruent-relation-gauge}.
Its integrality follows by pairing two vectors in the integral row
lattice. On \(P_O\) its image is \(P_Z^*\widehat O^\nu\).
\end{definition}

\begin{lemma}\label{lem:negative-last-middle-column}
The columns of \(P_ZY^{-1}\) of index \(b>d\) are zero.
Its first \(d\) rows are \((I_d\ \ 0)\).
For every such \(b>d\), \(X_\gamma Y^{-1}\varepsilon_b=0\)
at a singular level. At \(j=n\), for \(b=n\), this vector has
Euclidean norm at most \(N^{-1/2}\) at each embedding.
\end{lemma}
\begin{proof}
The first \(d\) rows agree with those of \(Y\), giving the
second assertion. At a singular level \(d=j\) and
\((\gamma f_i)_i Y^{-1}\varepsilon_b\) belongs to the original
frame span by Proposition~\ref{prop:canonical-singular-frame}.
Both its \(P^*\)-coordinates and its \(X_0\)-coordinates therefore
vanish. At \(j=n\), \(P_Z\) is the first \(\nu\) rows of \(Y=Z\),
giving the claim about its last column. For \(a=Y^{-1}\varepsilon_n\)
the vector \((\gamma f_i)_i a\) has all its \(P^*\)-coordinates
zero and its \(z\)-coordinate equal to \(-1/N\), since \(Za=\varepsilon_n\).
Its norm for the row form is \(-N a_n\bar a_n\) because
the transported frame has the same Gram matrix as the original one.
Thus at every embedding
\(\|\tau(X_\gamma a)\|^2=1/N-N|\tau a_n|^2\le1/N\).
\end{proof}

\begin{proposition}\label{prop:middle-congruence-to-shear}
Let \(\gamma,\delta\) have the same row class \(c\), level \(j\) and
matrix \(Z\), with possibly different column classes and reductions.
If \(F_\delta-F_\gamma\) sends \(L_M\) into \(K_Z\), then
\[
                 X_\delta-X_\gamma=gL P_Z
\]
for a rational matrix \(L\) whose columns belong to \(L_X\)
at the finite places. There is \(u\in\Gamma_c^0\) for which
the frames \(u\gamma f\) and \(\delta f\) agree on
\(X_0\oplus Ez\oplus P^*\). Here \(f\) denotes the matrix of
the original frame columns.
\end{proposition}
\begin{proof}
Set \(X=X_\delta-X_\gamma\), considered also as a middle matrix
with zero \(z\)-coordinate. For the form \(h_0\) the middle
lattice \(L_M\) is self-dual by Proposition~\ref{prop:row-group}
and its decomposition into outer and middle parts. In a local
basis the difference of evaluation maps on that lattice is
\(X^*h_0\), evaluated on the chosen basis vectors. Express each
of these columns in \(gP_Z^*\widehat O^\nu+gY^*\widehat O^n\).
Taking adjoints, the rows of coefficients give vectors in the
dual middle basis. That basis consists of integral vectors since
the Gram matrix is unimodular, also at conjugate local pairs.
In particular, for some finite adelic matrices with columns in \(L_M\)
one has \(X/g=L_0 P_Z+L'_0 Y\).
There is an integral projection of \(L_M\) onto \(L_X\):
locally the latter is the intersection with a vector subspace and
thus saturated. The basis extension by diagonal elimination in
Lemma~\ref{lem:valuation-index} extends a basis of this intersection
to one of the whole lattice. Project along the remaining vectors,
at each place. Its linear extension fixes \(X_0\), so applying
it gives the same equality with all columns of \(L_0,L'_0\) in \(L_X\).

Choose a rational \(X_0\)-valued matrix \(\mu\), with columns
in \(L_X\), so close to \(L_0\) at finite places that
\((L_0-\mu)P_ZY^{-1}\) also has columns in \(L_X\).
These are open lattice conditions containing \(L_0\).
For clarity, clear denominators by an integer, approximate the
integral coordinates at the finitely many exceptional places by
Chinese remaindering in \(O_E\), and divide back by that integer,
including every prime of the denominator among the exceptional places.
At all other places integrality suffices. This is the additive
approximation in the proof of Lemma~\ref{lem:unip-approx}.
Now \(B=(X/g-\mu P_Z)Y^{-1}\) is rational with columns in \(L_X\).
For \(b>d\), Lemma~\ref{lem:negative-last-middle-column} gives
\(B\varepsilon_b=0\) at singular levels. At \(j=n\) the only
such column has norm at most \(2/(g\sqrt N)<1\) at every embedding
by the same lemma. Its squared hermitian length lies in \(O_E\),
since the column belongs to the integral middle lattice.
The norm argument of Proposition~\ref{prop:boxed-frame-separation}
forces the length to vanish, hence the column is zero by definiteness.
Since the first \(d\) rows of \(Y\) are the first \(d\) rows of \(P_Z\),
\(B Y\) is a combination of those rows with coefficients the
first \(d\) columns of \(B\). Adding them to the corresponding
columns of \(\mu\) gives the required rational \(L\).

In the three blocks \(P,M_0,P^*\), write \(C_0=gL\) as a middle
matrix with zero \(z\)-entry, put \(C_0^\dagger=C_0^*h_0\), and take
\[
 u=\begin{pmatrix}
 I&-C_0^\dagger&-\eta C_0^\dagger C_0\\0&I&C_0\\0&0&I
 \end{pmatrix},\qquad \eta+\bar\eta=1,\quad\eta\in O_E ,
\]
using Lemma~\ref{lem:middle-hyperbolic-chart}. Multiplying by
the Gram matrix of the hyperbolic and middle summands verifies
\(u^*Ju=J\): the middle cross terms cancel and the remaining
corner is \((1-\eta-\bar\eta)C_0^\dagger C_0\).
Each off-diagonal block takes the class lattice into \(g\)
times that lattice by integrality of \(h_0\) there. Its inverse
does likewise by the finite series for the inverse of a unitriangular
matrix. It fixes \(P,z\) since \(C_0^\dagger z=0\), so
\(u\in\Gamma_c^0\) by Definition~\ref{def:boxed-real-frames}.
Its effect on the middle frame rows is addition of \(C_0P_Z\).
The \(z\)- and dual rows already agree by their formulas in
Definition~\ref{def:middle-frame-congruence}.
\end{proof}

\section{Congruence classes of rational frame orbits}

\begin{lemma}\label{lem:near-label-frame-injective}
Fix \(c,j,t\) as in Definition~\ref{def:near-labelled-terms}.
Let \((c',\gamma),(c'',\delta)\) be chosen representatives
with this label, and \(u,v\in\Gamma_c^0\). If \(u\gamma f=v\delta f\),
then \(c'=c''\) and \(\gamma=\delta\).
\end{lemma}
\begin{proof}
Write \(k,k'\) for the two finite positions of \(u\gamma,v\delta\)
from Definition~\ref{def:cut-global-columns}. All the class representatives
\(h_c,h_{c'},h_{c''}\) fix the frame, so \(k^{-1}k'\) fixes the
standard finite frame and preserves the full integral row lattice.
Its action on the quotient \(P^\perp/P\) fixes \(z\) and
preserves the middle lattice. The diagonal extension of that
quotient action, identity on the outer pairs, is integral with
integral inverse by the direct lattice decomposition of
Definition~\ref{def:row-space}. Removing it leaves an integral
unipotent by the flag characterization of Definition~\ref{def:unip}.
Hence the reduction of \(k^{-1}k'\) is in \(\bar S\) by
Proposition~\ref{prop:marked-local-level}. If \(s,s'\) are the
reductions of \(k,k'\), they have the same left coset and the same
label \(t\), so \(s=r(s\bar S)t=s'\). The left corrections have
principal congruence, hence the original reductions are also
\(s,s'\). Apply Lemma~\ref{lem:boxed-equal-classes} to \(u\gamma,v\delta\).
It gives \(c'=c''\) and \((u\gamma)^{-1}v\delta\in\Gamma_{c'}^0\).
Thus the two original terms have the same arithmetic double coset.
The choice in Definition~\ref{def:near-labelled-terms} identifies them.
\end{proof}

\begin{proposition}\label{prop:near-outer-congruence}
Fix \(Z,j\) and put \(a=M_b gD\bar D\), with
\(M_b\) from Corollary~\ref{cor:near-integral-colours}
and \(D\) from Definition~\ref{def:middle-frame-congruence}.
Take two near rational terms of these data and apply arbitrary
left corrections in \(\Gamma_c^0\) to them. Suppose their
corrected middle rows agree and the difference \(\Xi\) of their
corrected top rows has entries in \(aO_E\).
There exists a further left correction in \(\Gamma_c^0\) making
their entire frames equal.
\end{proposition}
\begin{proof}
The lower rows of both corrected frames are \(P_Z\): a correction
fixes the original frame, hence preserves \(Z\). Equality of
their Gram matrices therefore gives
\[
                    \Xi^*P_Z+P_Z^*\Xi=0.
\]
Write \(P=P_Z\) in this proof. Its rank is \(d=\min(j,\nu)\)
by Lemma~\ref{lem:negative-last-middle-column}, since \(PY^{-1}\)
has its first \(d\) rows \((I_d\ 0)\) and its final \(n-d\) columns
zero. In particular \(\ker P\) is spanned by the last \(n-d\)
columns of \(Y^{-1}\).

We claim that \(\Xi\ker P=0\). For \(j=n\), \(d=\nu\) and
\(P^*\) is injective, so the Gram identity applied to \(\ker P\)
gives this claim. For \(j<n\) the columns
\(\Xi\,\operatorname{adj}(Y)\varepsilon_b\), \(b>j\), are
differences of the outer coordinates of \(D_jv_b\) for the
two corrected frames. They lie in \(M_b O_E^\nu\) since
\(\Xi\in aO_E^{\nu\times n}\), \(gD\bar D\in O_E\), and \(Y\)
is integral. Corollary~\ref{cor:near-integral-colours} makes
them zero. Division by \(D\) proves the claim in this case too.

Define \(C_0:E^\nu\to E^n\) by first retaining only the first
\(d\) coordinates, padding them to length \(n\) by zero, and
then multiplying by \(Y^{-1}\). For each \(x\in E^n\) the first
\(d\) coordinates of \(Y(x-C_0Px)\) vanish. Thus \(PC_0P=P\)
by the preceding description of \(\ker P\), and
\(\Xi C_0P=\Xi\). Put \(J_0=PC_0,\ L=\Xi C_0\), and set
\[
                         H=L-L^*-J_0^*L .
\]
Since \(J_0^*L=C_0^*P^*\Xi C_0\), this last summand is
skew-hermitian by the Gram identity. Thus \(H^*=-H\).
Also \(L^*P=-J_0^*\Xi\), and consequently \(HP=\Xi\).
Every entry of \(DC_0\) is integral. The entries of \(L,L^*\)
are multiples of \(g\) in \(O_E\), because \(\Xi\) is divisible
by \(gD\bar D\). The same is true of \(J_0^*L\), whose possible
denominator is \(D\bar D\). In the three blocks \(P_{\rm out},
M_0,P_{\rm out}^*\), where \(P_{\rm out}\) denotes the outer
isotropic space, the matrix with identity diagonal and sole
off-diagonal block \(H\) in the top right is therefore integral
and principal in every class. Skewness is precisely its isometry
condition, by multiplying the hyperbolic Gram matrix. It fixes
the original frame and adds \(HP=\Xi\) to the top rows.
Choosing its sign to match the order of subtraction aligns
the two frames.
\end{proof}

\begin{theorem}\label{thm:actual-near-index-moment}
Fix a row class \(c\), a finite label \(t\), and a level \(j\).
Let \(\mathcal A\) be any subset of the pairs \((c',\gamma)\)
of near chosen representatives with this label and level as in
Definition~\ref{def:near-labelled-terms}. Fix an \(E\)-subspace
\(T\subset X_0\) of codimension at most \(s\), \(0\le s\le k\).
Assume \(o_{\gamma,g}(T)\) is finite for every member and
choose a positive integer \(m_\gamma'\) divisible by it.
Then at all sufficiently large arithmetic scales, uniformly in
these choices,
\[
    \sum_{(c',\gamma)\in\mathcal A}(m_\gamma')^{-k}
                 \le b^{100 f n^4(n+s+1)} .
\]
\end{theorem}
\begin{proof}
By Corollary~\ref{cor:near-integral-colours} there are at most
\(b^{12fn^4}\) choices of \(Z\). Fix one of them and use
\(a=M_b gD\bar D,\ K=K_Z\).
On the full middle lattice \(L_M\) use the maps \(F_\gamma\)
of Definition~\ref{def:middle-frame-congruence}. Its submodule
\(L_1=L_c(T)\) is saturated locally, so its local bases extend
to bases of \(L_M\) by Lemma~\ref{lem:valuation-index}.
Combining these yields a free summand over \(\widehat O\) of
rank \(e=\dim T\) with a complementary free summand of rank
\(k+1-e\le s+1\). Since \(K\) contains \(g F_\gamma(P_O)\),
the order of the image of \(L_1\) modulo \(K\) divides
\(o_{\gamma,g}(T)\): the former is the order of a quotient of
the finite group in Definition~\ref{def:congruent-relation-gauge}.

For each nonempty fibre of the induced maps on \(L_M\) modulo
\(K\), fix one member. Use Proposition~\ref{prop:middle-congruence-to-shear}
to translate every member on the left by \(\Gamma_c^0\) to have
the same middle frame rows as this member. Its top rows are
integral over \(O_E\) since the full frame is integral in the
class lattice at all finite places and its outer lattice is standard.
There are at most \(I(a)^{\nu n}\) choices of their residue
matrix modulo \(a\), by Definition~\ref{def:middle-lift-disc}.
Proposition~\ref{prop:near-outer-congruence} and
Lemma~\ref{lem:near-label-frame-injective} give at most one
original pair for each such residue matrix. Thus we may use
\(A_0=I(a)^{\nu n}\) in Lemma~\ref{lem:finite-image-moment}.
Indeed \(a\widehat O^n\subset K\) by
Definition~\ref{def:middle-frame-congruence}. With \(\ell=k\ge e\)
the weighted count for this \(Z\) is at most
\[
                         I(a)^{\nu n+n(n+s+1)} .
\]
There is a uniform bound \(I(a)\le b^{48fn^3}\) eventually:
at every \(\tau\) one has \(M_b\le b^{11n^3}\),
\(g\le b^{22}\) by the level-prime bounds of
Lemma~\ref{lem:individual-kappas}, and
\(|\tau D|^2\le b^{11n^3}\) by
Proposition~\ref{prop:near-minor-coordinate-bounds}. Multiply
the squares of the resulting \(f\) bounds on \(|\tau a|\),
using Definition~\ref{def:middle-lift-disc}. Since \(\nu=n-1\),
the total exponent after summing over \(Z\) is at most
\(12fn^4+48fn^4(2n+s)\le100fn^4(n+s+1)\).
\end{proof}

\section{Summing the weights of congruent pairs}

\begin{lemma}\label{lem:pair-evaluation-divisors}
Suppose \(\gamma,\delta\in\mathcal T(c,s_0,c')\) have levels
as in Definition~\ref{def:arithmetic-relation-indices}. Let
\(T\) be the first projection of
\(\mathcal L_\gamma\cap\mathcal L_\delta\cap\ker t\) as in
Lemma~\ref{lem:common-definite-relation}. Write \(d=\min(j,\nu)\)
for the level \(j\) of \(\gamma\), and put
\(T'=T\cap\ker(F_\gamma|_{X_0})\). Then
\[
 \begin{gathered}
    o_{\gamma,g}(T)\ \mid\ g^{2fd}o_{\gamma,1}(T),\qquad
    o_{\delta,g}(T')\ \mid\ s_\delta,\\
    \operatorname{codim}_{X_0}T'\le
                      \operatorname{codim}_{X_0}T+n .
 \end{gathered}
\]
The saturation index \(s_\delta\) uses the full sum of the
two isotropic spaces as in Definition~\ref{def:arithmetic-relation-indices}.
\end{lemma}
\begin{proof}
Write \(B=F_\gamma(P_O)\), \(A=F_\gamma(L_c(T))\).
The image \(B\) has local rank \(d\) by
Lemma~\ref{lem:negative-last-middle-column} and its expression
in Definition~\ref{def:middle-frame-congruence}. It is locally
free, by Lemma~\ref{lem:valuation-index}. Hence \([B:gB]=
I(g)^d=g^{2fd}\) by Definition~\ref{def:middle-lift-disc}.
The index \([A:A\cap gB]\) is \([A:A\cap B]\) times the order
of the image of \(A\cap B\) in \(B/gB\). They are finite:
for every \(x\in T\), the defining common pair has a lift
\(p+x\) with \(p\in P\) and \(\gamma^{-1}(p+x)\in P\oplus X_0\).
Hence \(F_\gamma(x)=-F_\gamma(p)\). Choose such \(p\) on
a rational basis of \(T\) and extend linearly. After multiplication
by a fixed positive integer this map sends \(L_c(T)\) into
\(P_O\), by the bounded denominators of these finite lattices
in Definition~\ref{def:middle-lift-disc}. Thus \(A\cap B\)
has finite index in \(A\). The
order of a subgroup of a finite group divides the group order,
as its cosets partition the group into sets of that size.
This proves the first divisibility. Intersecting \(T\) with
the kernels of the \(n\) coordinates of \(F_\gamma\) proves
the asserted codimension bound.

For \(x\in T'\) one has \(\gamma^{-1}x=p_\gamma+x'\)
with \(p_\gamma\in P,\ x'\in X_0\), since \(F_\gamma x=0\)
means orthogonality to \(\gamma P\) and \(\gamma z\).
Thus \(x\in W_\gamma\) and \((x,x')\in\mathcal L_\gamma\cap\ker t\).
The first projection of this last space is injective by
Lemma~\ref{lem:common-definite-relation} applied with two
copies of \(\gamma\). By the definition of \(T\), the pair
over \(x\) in the common relation is therefore \((x,x')\).
Its membership in \(\mathcal L_\delta\) gives \(p,p'\in P\)
with \(p+x=\delta(p'+x')\) by Definition~\ref{def:global-middle-relation}.
Consequently
\[
                    \delta\gamma^{-1}x-x
                            =p+\delta(p_\gamma-p')\in P+\delta P .
\]
This subspace inclusion is rational and persists under extension
to the finite adeles by applying it to a rational basis of \(T'\).

For \(x\in L_c(T')\) the finite vector
\(z_x=(x-\delta\gamma^{-1}x)/g\) thus belongs to \(P_f+\delta P_f\).
It belongs to \(\widehat\Lambda_c\) because \(\delta\gamma^{-1}\)
acts as identity modulo \(g\) there by
Proposition~\ref{prop:middle-lattice-congruence}.
Moreover \(F_\delta(\delta\gamma^{-1}x)=F_\gamma(x)=0\).
It follows that \(F_\delta(x)=g F_\delta(z_x)\). On
\(P_O+\delta P_O\) the image of \(F_\delta\) is exactly
\(F_\delta(P_O)\), since the original frame is orthogonal to \(P\).
Thus the image of \(F_\delta(L_c(T'))\) modulo \(gF_\delta(P_O)\)
is a subgroup of the image of the finite quotient
\[
       ((P_f+\delta P_f)\cap\widehat\Lambda_c)/(P_O+\delta P_O)
\]
under \(gF_\delta\) modulo \(gF_\delta(P_O)\). Its order divides
that quotient's order \(s_\delta\), as claimed.
\end{proof}

\begin{definition}\label{def:summed-pair-weights}
Use the level \(j\) and near representatives of
Definition~\ref{def:near-labelled-terms}. In each sum below
take the ordered pairs \(\gamma,\delta\) having that level
and lying in the same \(\mathcal T(c,s_0,c')\). Both
representatives use the same chosen double-coset section.
Write \(T=T(\gamma,\delta)\) for the subspace of
Lemma~\ref{lem:pair-evaluation-divisors}, \(h=k-\dim_E T\),
\(d=\min(j,\nu)\), and set
\[
 \begin{split}
   M_h&=|\bar S|\,C\,N^{-kf}(2T_0^2)^h,\\
   w(\gamma,\delta;c)&=
   \frac{1}{M_h d_c(T)^{k+1}
                    \big(g^{2fd}o_{\gamma,1}(T)s_\delta\big)^k}.
 \end{split}
\]
Here \(T_0=b^{kf/4}\) as in Definition~\ref{def:two-energy-numbers}.
\end{definition}

\begin{theorem}\label{thm:summed-actual-pair-weights}
For each fixed \(\epsilon>0\), at all sufficiently large scales,
uniformly in the level \(j\),
\[
  \sum_{c,c'}\sum_{s_0\in\mathcal G}
     \sum_{\gamma,\delta} w(\gamma,\delta;c)\le 2b^{\epsilon kf}.
\]
In particular the sum is at most \(2T_0^2\).
\end{theorem}
\begin{proof}
Group by the finite label of \(s_0\), by the row class \(c\),
and by the rational subspace \(T\) of codimension \(h\).
Fix these data. For a fixed first pair \((c',\gamma)\), the
second term has the same label and the same \(c'\), and its
reduction equals that of the first. The space
\(T'=T\cap\ker(F_\gamma|_{X_0})\) is fixed throughout this
inner sum. By Lemma~\ref{lem:pair-evaluation-divisors},
\(o_{\delta,g}(T')\mid s_\delta\). The inner sum of \(s_\delta^{-k}\)
is therefore at most \(B_{\min(k,h+n)}\), where
\(B_v=b^{100fn^4(n+v+1)}\), by
Theorem~\ref{thm:actual-near-index-moment}. For first terms
having at least one partner the lemma gives
\(o_{\gamma,g}(T)\mid g^{2fd}o_{\gamma,1}(T)\).
The same theorem bounds the sum of the corresponding inverse
\(k\)-th powers by \(B_h\). Each pair \((c',\gamma)\) at
fixed \(c\) determines its reduction uniquely by
Definition~\ref{def:cut-global-columns}, so no multiplicity
in \(s_0\) is incurred. All weights are nonnegative, allowing
these successive sums by Lemma~\ref{lem:integral-tools}.

For any fixed \(\eta>0\), Corollary~\ref{cor:marked-subspace-moment}
gives, eventually,
\[
 \sum_{c}\sum_{\operatorname{codim}T=h}d_c(T)^{-(k+1)}
                \le C N^{-kf} b^{\eta kf(h+1)} .
\]
Summing over the \(|\bar S|\) labels and dividing by \(M_h\)
bounds the contribution at \(h\) by
\[
           (2T_0^2)^{-h}\,B_h B_{\min(k,h+n)}
                                   b^{\eta kf(h+1)} .
\]
Uniformly for \(0\le h\le k\), the logarithm to base \(b\) of the
two \(B\)-factors is at most \(100fn^4(3n+2h+2)=o(kf(h+1))\)
by \(k/n^9\to\infty\) in Proposition~\ref{prop:thick}.
Take \(\xi=\min(\epsilon,1/2)\) and \(\eta=\xi/2\), and take
the scale large enough to absorb the two \(B\)-factors into
\(b^{\xi kf(h+1)/2}\). The resulting bound is at most
\(2^{-h}b^{\xi kf}\), since \(T_0^2=b^{kf/2}\).
The geometric series has sum two. Taking \(\epsilon=1/2\)
proves the last assertion.
\end{proof}

\section{A determinant integral for the real frame volumes}

\begin{definition}\label{def:negative-frame-volume}
At one infinite place use \(F=F_0,\ J,\ G_1=U(J)\) from
Definition~\ref{def:open-bessel-frame}. Let \(X\) be the isometric
inclusion of the \(k\) definite middle coordinates, and let
\(F_-=(I-J)F\). For \(y\in G_1\) put
\[
          V_y=F_-^*y^{-1}F,\qquad v(y)=|\det V_y|.
\]
For \(j\le n\) write \(E_j:\C^j\to\C^n\) for the first-coordinate
inclusion, and \(I_j'=F E_j\). For \(g_*\in G_1\) let
\(Z=F^*Jg_*F,\ z_j=E_j^*ZE_j\). Say the frames have the required
shared columns at \(j\) if for each \(b>j\) there exist \(a_b,x_b\in\C^n\)
vanishing in coordinates larger than \(b\), with
\((a_b)_b=1,\ |(x_b)_b|=1,\ F x_b=g_*F a_b\).
\end{definition}

\begin{lemma}\label{lem:negative-frame-ratio}
The function \(v\) is positive and continuous, and invariant under
left multiplication by \(\mathcal U_0\). Suppose \(z_j\) is invertible
and the required shared columns exist at \(j\). Set
\[
   B=V_y^{-1}F_-^*y^{-1},\qquad t=g_* I_j' z_j^{-1},
   \qquad p=F^*Jt .
\]
Then
\[
             \frac{v(g_*^{-1}y)}{v(y)}
                               =|\det(p^* B t)|\,|\det z_j|.
\]
At \(j=0\) the determinants on the right are one.
\end{lemma}
\begin{proof}
The columns \(e_i-e_i^*,\,2z\) of \(F_-\) span the negative
eigenspace of the hermitian involution \(J\). If \(V_y a=0\),
then \(y^{-1}F a\) belongs to its positive eigenspace. Its
\(J\)-norm is \(-|a_n|^2\), so the vector vanishes and \(a=0\).
Continuity follows by inversion and determinants. An element of
\(\mathcal U_0\) fixes \(F\) pointwise, giving left invariance.
We have \(BF=I_n,\ E_j^*p=I_j\). Put \(M=B g_* I_j'-E_j\).
For each \(b\le j\), \(I_n+Mp^*\) on \(\varepsilon_b\) agrees
with \(B g_*F\) there by \(p^* E_j=I_j\).
For a shared column at \(b>j\), one has
\(p^* x_b=t^*JF x_b=z_j^{-*}{I_j'}^*JF a_b=0\):
when \(j<n\) these \(j\) isotropic columns are orthogonal to
all of \(F\). Hence \((I_n+Mp^*)x_b=x_b=B g_*F a_b\).
The matrices with columns \(\varepsilon_b\) up to \(j\), then
\(x_b\) or \(a_b\), are upper triangular with determinants of
modulus one. Therefore
\[
 |\det(B g_*F)|=|\det(I_n+Mp^*)|
      =|\det(I_j+p^*M)|=|\det(p^*Bt z_j)|.
\]
The middle identity follows by eliminating either off-diagonal
block in the matrix with diagonal identities and off-diagonal
blocks \(M,-p^*\). The left side is precisely the ratio of
the volumes by the definition of \(B\).
\end{proof}

\begin{lemma}\label{lem:complex-determinant-row-integral}
For \(0\le d\le n\), \(k\ge1\), an invertible complex \(d\)-square
matrix \(S\), and a complex \(k\)-by-\(d\) matrix \(A_0\), one has
\[
 \int_{\C^{k\times d}}
       \det\!\big(I_d+(AS+A_0)^*(AS+A_0)\big)^{-(k+n)}\,\dd A
        \ \le\ |\det S|^{-2k}\big(\pi^d/d!\big)^k .
\]
Here measure is Lebesgue measure in the complex entries, and a
zero-dimensional space has mass one.
\end{lemma}
\begin{proof}
Translation and right multiplication on each row remove \(A_0,S\);
the real Jacobian for a complex invertible matrix \(S\) is
\(|\det S|^2\), by triangular elimination into shears, swaps
and complex diagonal scalings. For an unshifted matrix let \(H\)
be \(I_d\) plus the sum of row outer products except the last one.
Writing \(u\) for that last row gives
\(\det(H+u^*u)=\det H(1+uH^{-1}u^*)\), by the block determinant
identity in Lemma~\ref{lem:negative-frame-ratio}. With the change
\(u=wH^{1/2}\) the Jacobian is \(\det H\), as follows by the
positive matrix diagonalization in Lemma~\ref{lem:near-volume-inverse}.
Integration thus lowers the exponent by one and extracts
\(\int_{\C^d}(1+\|w\|^2)^{-a}\dd w\) at the current exponent \(a\).
In one complex coordinate, polar coordinates give
\(\int_{\C}(h+|z|^2)^{-a}\dd z=\pi h^{1-a}/(a-1)\)
for \(h>0,\ a>1\). Iterating gives
\(\pi^d/\prod_{i=1}^d(a-i)\) at \(a>d\).
The smallest exponent used in removing the \(k\) rows is \(n+1\),
so each such factor is at most \(\pi^d/d!\).
All successive integrals are valid for nonnegative functions
by Lemma~\ref{lem:product-integration}.
\end{proof}

\begin{theorem}\label{thm:real-frame-volume-majorant}
Fix \(y\in G_1,\ \theta>0\) with
\(\|u y\|_{\rm op}\le\theta\) for some \(u\in\mathcal U\).
Fix \(Z\in\C^{n\times n}\) with \(z_j=E_j^*ZE_j\) invertible,
\(0\le d\le j\le n\), an invertible \(d\)-square matrix \(R_1\),
and a \(k\)-by-\(j\) matrix \(V\). Let \(E:\C^d\to\C^j\) be
the first-coordinate inclusion. There is a nonnegative continuous
function \(K\) on \(\C^{k\times d}\) with
\[
 \int K(A)\dd A \le
  \left(\frac{(2\theta^2)^j}{|\det z_j|}\right)^{k+n}
  \left(\frac{\pi^d}{d!\,\theta^{2d}|\det R_1|^2}\right)^k.
\]
For every \(g_*\in G_1\) with \(F^*Jg_*F=Z\), the required
shared columns at \(j\), and
\(
 X^*g_* I_j'=V+A R_1 E^* z_j,
\)
one has \(\big(v(y)/v(g_*^{-1}y)\big)^{k+n}\le K(A)\).
\end{theorem}
\begin{proof}
Let \(B,p\) be as in Lemma~\ref{lem:negative-frame-ratio};
\(p=ZE_j z_j^{-1}\) is already fixed by \(Z\).
Use the upper unitriangular frame action \(A(u)\) of
Lemma~\ref{lem:bessel-frame-action}, and put
\[
 T=\theta E_j^* A(u) E_j,\quad C_0=T p^* B,\quad
 Q_1=\theta^{-1}(uy)^{-1} I_j'.
\]
Since \(J\) is an ordinary unitary involution, a \(J\)-isometry
\(h\) has \(h^{-1}=Jh^*J\) and \(\|h^{-1}\|_{\rm op}=\|h\|_{\rm op}\).
Thus \(\|Q_1\|_{\rm op}\le1\). Upper triangularity gives
\(Q_1 T=y^{-1} I_j'\), and therefore \(C_0 y Q_1=I_j\) by
\(BF=I_n,\ p^*E_j=I_j\). It follows that
\((C_0y)(C_0y)^*\ge I_j\), since taking adjoints gives
\(\|z\|\le\|(C_0y)^*z\|\) for every \(z\). Here and below
matrix inequalities refer to the nonnegative hermitian order.
As \(F_-\) is negative for \(J\), one has
\(BJB^*=-(By)(By)^*\) directly from
\(y^{-1}Jy^{-*}=J,\ JF_-=-F_-\). Consequently
\(C_0 J C_0^*=-(C_0 y)(C_0 y)^*\le-I_j\).

Write \(t=g_* I_j' z_j^{-1}\) and set
\(Y_0=JtT^*-C_0^*,\ W=X^*Y_0\). Then \(F^*Y_0=0\)
since \(C_0F=T p^*\). Its columns have zero outer \(e_i\)
coordinates and zero \(z\) coordinate. On such columns the
\(J\)-Gram matrix is just that of their definite components,
so \(Y_0^*JY_0=W^*W\). Also
\(t^*Jt=z_j^{-*}{I_j'}^*JI_j'z_j^{-1}\le0\).
Expanding \(Y_0^*JY_0\) and applying these inequalities gives
\[
       -C_0 tT^*-(C_0tT^*)^*\ \ge\ I_j+W^*W .
\]
For a complex matrix with positive hermitian part \(H\),
write it as \(H^{1/2}(I+iS_1)H^{1/2}\) with \(S_1\)
hermitian. Its determinant has modulus at least \(\det H\)
by diagonalizing \(S_1\). If \(H\ge H_0>0\), conjugation by
\(H_0^{-1/2}\) gives \(\det H\ge\det H_0\) by the same
diagonalization. Applying both facts to \(-C_0tT^*\), and using
\(|\det T|=\theta^j\), now gives by Lemma~\ref{lem:negative-frame-ratio}
\[
 \frac{v(y)}{v(g_*^{-1}y)}
      \le\frac{(2\theta^2)^j}
                {|\det z_j|\,\det(I_j+W^*W)} .
\]
This derivation also holds for \(j=0\) by empty determinants.

Since \(X^*J=X^*\), the hypothesis on the definite components
gives \(WE=A S+A_0\), where
\[
       S=R_1 E^*T^*E,\qquad A_0=(V z_j^{-1}T^*-X^*C_0^*)E .
\]
These are fixed, and \(|\det S|=|\det R_1|\theta^d\) by
triangularity. Since \(WW^*\ge(WE)(WE)^*\), the determinant
identity of Lemma~\ref{lem:negative-frame-ratio} and the
preceding determinant monotonicity show
\(\det(I_j+W^*W)\ge\det(I_d+(AS+A_0)^*(AS+A_0))\).
In the ratio bound replace the denominator determinant by
this smaller one and raise to \(k+n\). The resulting continuous
function \(K(A)\) has the claimed integral by
Lemma~\ref{lem:complex-determinant-row-integral}.
\end{proof}

\section{Averaging the cut orbital integrals}

\begin{definition}\label{def:positive-schur-orbits}
For \(\gamma\in\mathcal T(c,s,c')\) let
\(\mathcal O=\Gamma_c^0\gamma\Gamma_{c'}^0\). Write
\(\beta_c^0={\bf1}_{\mathcal B}\beta_c\), and use the positive
function \(v\) of Definition~\ref{def:negative-frame-volume} at
each real place. For \(\delta\in\mathcal O\) put
\[
 \begin{split}
  q_\delta(y)&=\prod_\tau
      \left(\frac{v(y_\tau)}{v((\delta_\infty^{-1}y)_\tau)}\right)^{k+n},\\
  H_{\mathcal O}(y)&=\sum_{\delta\in\mathcal O}
          \beta_{c'}^0(\delta_\infty^{-1}y)q_\delta(y),\\
  S(\mathcal O;c)&=
        \int_{G_\infty}\beta_c^0(y)\|V_\infty(y)\|^2H_{\mathcal O}(y)\dd y .
 \end{split}
\]
Use \(S(\mathcal O^{-1};c')\) for the same expression on the
inverse orbit. The quantities are nonnegative extended integrals.
\end{definition}

\begin{lemma}\label{lem:two-sided-schur-orbit}
Let \(a_\delta(c,c')=\int\beta_c(y)\beta_{c'}(\delta_\infty^{-1}y)
\|K^{V_\infty}_\delta(y)\|_{\rm HS}\dd y\).
If the two \(S\)-integrals are finite, then
\[
      \left(\sum_{\delta\in\mathcal O}a_\delta(c,c')\right)^2
                       \le S(\mathcal O;c)S(\mathcal O^{-1};c').
\]
If no member is near, all \(a_\delta\) vanish.
\end{lemma}
\begin{proof}
Where the norm integrand is positive, both arguments lie in
\(\mathcal E\) by the cut support in Theorem~\ref{thm:actual-bessel-truncation}.
Positive cutoffs there force them to belong to \(\mathcal B\)
by Proposition~\ref{prop:boxed-arithmetic-cutoff}. The cutoffs may
therefore be replaced by \(\beta^0\), and such a point witnesses
nearness. By the rank-one norm computation of
Theorem~\ref{thm:full-global-truncation} the norm is
\(\|V_\infty(y)\|\,\|V_\infty(\delta_\infty^{-1}y)\|\).
On the space of \((\delta,y)\) with counting times Haar measure
multiply the first norm by \(q_\delta^{1/2}\), the second by
\(q_\delta^{-1/2}\), and each by the square root of the product
of the two cutoffs. Cauchy--Schwarz gives the claimed product:
for the second squared integral substitute
\((\delta^{-1},\delta_\infty^{-1}y)\) for \((\delta,y)\).
This change preserves the measure and exchanges the two ratios.
\end{proof}

\begin{lemma}\label{lem:frame-stabilizer-average}
Put \(U_{0,\infty}=\prod_\tau\mathcal U_0\). Let \(H\) be a
nonnegative Borel function on \(G_\infty\) invariant under
left \(\Gamma_c^0\). Suppose \(A\) is a measurable parameter set
with a finite measure \(m\) of mass at least one, and
\(h:A\to U_{0,\infty}\) is measurable. If
\[
           \int_A H(u h(x)y)\dd m(x)\le D_0
                \quad (u\in U_{0,\infty},\ y\in\mathcal B),
\]
then
\[
       \int\beta_c^0(y)\|V_\infty(y)\|^2 H(y)\dd y
                                      \le D_0\,\mathfrak v J_c .
\]
\end{lemma}
\begin{proof}
Write \(\beta_c=\omega_c{\bf1}_{D_c}\) as in
Proposition~\ref{prop:boxed-arithmetic-cutoff}, with \(D_c\) a
Borel domain for \(\Gamma_c^0\) on the real group.
The factor \(p={\bf1}_{\mathcal B}\omega_c\|V_\infty\|^2\)
is left \(U_{0,\infty}\)-invariant: the frame and size are
unchanged, \(\omega_c\) is invariant there by the cutoff construction,
and the vector transforms by unitary scalars.
For a fixed \(y\), \(D(y)=\{u:u y\in D_c\}\) is a domain for
\(\Gamma_c^0\) acting by left translation on \(U_{0,\infty}\).
Haar measure on \(U_{0,\infty}\) is invariant on the right too:
the triangular group multiplication in the coordinates of
Lemma~\ref{lem:frame-point-stabilizer} has unit Jacobian on both
sides. The translate \(D(y)h(x)\) is again a domain. Partition
one domain by translates of the other as in
Lemma~\ref{lem:quotient-measure}. Left \(\Gamma_c^0\)-invariance
of \(H\) gives
\[
        \int_{D(y)}H(u y)\dd u=\int_{D(y)}H(u h(x)y)\dd u .
\]
Integrate in \(x\). If \(y\in\mathcal B\), the hypotheses and
nonnegative integration give an upper bound \(D_0\int_{D(y)}1\dd u\)
for the left integral, since \(m(A)\ge1\).
If \(D_0=0\), the left integral is zero by integrating the
zero function in \(u\). Multiplication by \(p(y)\), first on
the set \(p(y)>0\), makes the same comparison valid everywhere.

To integrate the orbitwise inequality, choose a positive Borel
density \(\alpha\) on \(G_\infty\) with
\(\int_{U_{0,\infty}}\alpha(u y)\dd u=1\).
It is obtained by dividing a Gaussian in the matrix norm by its
orbit integral as in Lemma~\ref{lem:bessel-quotient-density}:
the subgroup is closed and the polynomial volume bound from
Lemma~\ref{lem:period-ball-cover} yields finiteness. By inversion,
left Haar invariance on \(G_\infty\) and
Lemma~\ref{lem:product-integration}, for every nonnegative \(L\),
\(\int\alpha(y)\int L(uy)\dd u\dd y=\int L(y)\dd y\).
Use \(L={\bf1}_{D_c}pH\) and \(L={\bf1}_{D_c}p\) on the two sides.
It remains to bound \(\int\beta_c^0\|V_\infty\|^2\).
Drop the box indicator and use Proposition~\ref{prop:boxed-arithmetic-cutoff}
and \(\int\|V_\tau\|^2\dd\bar y_\tau\le\dim\pi_\tau\) from
Theorem~\ref{thm:actual-bessel-truncation}, giving \(\mathfrak vJ_c\).
\end{proof}

\begin{proposition}\label{prop:thin-orbit-definite-cosets}
Suppose \(\gamma\) has level \(j\), and put \(d=\min(j,\nu)\).
All \(\delta\in\mathcal O\) have the same \(Z\) and level. They
can be partitioned into at most \(i_\gamma^k\) subsets such
that within each, for a fixed member \(\delta_0\), one has
\[
 X_\delta-X_{\delta_0}=g L_\delta (P_Z)_{1:d,*}.
\]
Here \(X_\delta,P_Z\) are the rational frame blocks of
Definition~\ref{def:middle-frame-congruence}. The \(k\)-by-\(d\)
matrix \(L_\delta\) is unique, rational, with finite columns
in \(L_X\). For every real place \(\tau\), the difference of the
scaled definite blocks of the frames is
\[
 g\,(\tau L_\delta)\operatorname{diag}_{i\le d}
      (1/\bar d_{\tau i})\,(Z_\tau^{\rm sc})_{1:d,*},
 \quad Z_\tau^{\rm sc}=
      \operatorname{diag}(\bar d_{\tau i})\,(\tau Z)\,
                                  \operatorname{diag}(d_{\tau i}).
\]
\end{proposition}
\begin{proof}
Write \(\delta=u\gamma v\) with \(u,v\) in the two pointwise
stabilizers. The second fixes \(f\); the first fixes the original
frame and consequently preserves all its pairings with \(\gamma f\).
It also fixes all the shared vectors themselves; the level is
unchanged by Definition~\ref{def:arithmetic-relation-indices}
and Theorem~\ref{thm:rational-agreeing-levels}. The action of \(u\)
in coordinates \(P,M_0,P^*\) has identity diagonal, and its middle
frame correction is \(C_0 P_Z\), where \(C_0:P^*\to M_0\) is its
middle-right block. The isometry condition gives the upper-middle
block \(-C_0^\dagger\) as in Proposition~\ref{prop:middle-congruence-to-shear}.
Since \(u\) fixes \(z\), the columns of \(C_0\) are orthogonal to \(z\);
they belong to \(gL_X\) by the principal congruence of \(u\)
on each integral dual basis vector. The matrix is rational.

Consider the modules of matrices with rows in the full span
\(R_\nu(\gamma)\) or the initial span \(R_d(\gamma)\) of
Definition~\ref{def:arithmetic-relation-indices}, with coefficient
columns in \(L_X\). In local bases of \(L_X\) these are respectively
\(k\) copies of the two row modules; hence their quotient has
order \(i_\gamma^k\). This description commutes with change
of the local basis by an invertible integral matrix. Record
the class of \((X_\delta-X_\gamma)/g\) in this quotient.
For equal classes the difference is a combination of the first
\(d\) rows with coefficient columns locally integral in \(L_X\).
Those rows are independent by Lemma~\ref{lem:negative-last-middle-column}.
Solving a nonsingular minor makes the coefficients rational
and unique and so identifies them with the local coefficients.
The scaled definite block at \(\tau\) is
\((\tau X_\delta)\operatorname{diag}(d_{\tau i})\): the definite
coordinates have scale one by Proposition~\ref{prop:unipotent-covol}.
Pairing the scaled frames gives the displayed \(Z_\tau^{\rm sc}\).
Substitution gives the final formula, including the empty set
of rows when \(d=0\).
\end{proof}

\section{The cut amplitude at small discriminant}

\begin{proposition}\label{prop:actual-thin-schur-integral}
For any fixed \(\eta>0\), all sufficiently large scales, and
every orbit \(\mathcal O\) of level \(j\),
\[
 S(\mathcal O;c)\le
 \mathfrak v J_c\, b^{\eta kf}(b^4g)^{-dkf}
                       i_\gamma^k m_\gamma^{-(k+1)/2},
       \qquad d=\min(j,\nu).
\]
\end{proposition}
\begin{proof}
The function \(H_{\mathcal O}\) is left \(\Gamma_c^0\)-invariant
since its terms are permuted and \(v\) is left \(\mathcal U_0\)-invariant
by Lemma~\ref{lem:negative-frame-ratio}. In
Lemma~\ref{lem:frame-stabilizer-average} use entrywise complex unit
disks for the parameters \(A=(A_\tau)\), \(A_\tau\in\C^{k\times d}\),
with product Lebesgue measure of mass \(\pi^{kdf}\ge1\).
Set \(\theta=R+1\) and
\[
       R_\tau=\operatorname{diag}_{i\le d}
                          \frac{\pi g}{b\bar d_{\tau i}} .
\]
Choose \(h(A)_\tau\in\mathcal U_0\) so that the middle-right block
of its inverse has first \(d\) columns \(A_\tau R_\tau\) in the
definite subspace and all other entries zero. In the formula for
\(u(B,T)\) of Definition~\ref{def:open-bessel-frame}, take
\(T=0\) and choose \(B\) by its adjoint, then invert if required.
This gives a continuous choice.

Fix \(y\in\mathcal B,\ u\in U_{0,\infty}\) and one of the
subsets of Proposition~\ref{prop:thin-orbit-definite-cosets},
with base \(\delta_0\). Put \(L_\delta'=( (b/\pi)\tau L_\delta)_\tau\).
The frame of
\(g_{*,\tau}=h(A)_\tau^{-1}u_\tau^{-1}\delta_{\infty,\tau}\)
has Gram pairings \(Z_\tau^{\rm sc}\). Its first \(j\) definite
columns are
\[
         V_\tau+(A_\tau+L'_{\delta,\tau})R_\tau E^* z_{\tau,j},
 \qquad z_{\tau,j}=(Z_\tau^{\rm sc})_{1:j,1:j},
\]
where \(V_\tau\) consists of the first \(j\) definite columns
of \(u_\tau^{-1}\delta_{0,\infty,\tau}F_0\) and
\(E:\C^d\to\C^j\) is the initial inclusion. Indeed \(u_\tau^{-1}\)
adds a fixed matrix times the common dual rows. The difference
of definite columns for \(\delta,\delta_0\) is therefore the
scaled difference in Proposition~\ref{prop:thin-orbit-definite-cosets}.
The inverse of \(h(A)_\tau\) adds exactly the remaining term
with \(A_\tau\). At a singular level the shared columns of
Theorem~\ref{thm:rational-agreeing-levels} remain shared under
both left operations since both fix the original frame. The
change to the scaled frame multiplies both leading coefficients
at a shared height by the same factor, so their ratio has modulus
one by that theorem. At \(j=n\) there are no required shared
columns. Also \(y_\tau\) has a representative of norm at most
\(\theta\) as in Theorem~\ref{thm:actual-bessel-truncation}.
Thus Theorem~\ref{thm:real-frame-volume-majorant} supplies
a kernel \(K_\tau\), independent of \(\delta\) in this subset,
for which
\[
       q_\delta(uh(A)y)\le
                     \prod_\tau K_\tau(A_\tau+L'_{\delta,\tau}).
\]
Here the numerator volumes are unchanged under \(u,h(A)\).

In each summand translate \(A\) by \(L'_\delta\), with new
variable \(A'\). At a fixed \(A'\) at most one value of \(L_\delta\)
occurs among the terms whose translated unit disks contain \(A'\).
Indeed if two occur, the difference of any corresponding column
of \(L_\delta\) belongs to \(L_X\) and has squared norm at most
\(4\pi^2 k/b^2<1\) at every embedding. Its hermitian length is
a rational element of \(O_E\) by the integral class lattice and
vanishes by the norm argument of Proposition~\ref{prop:boxed-frame-separation}.
Definiteness then forces columnwise equality.
For this single \(L_\delta\), the point \(Y=uh(A'-L'_\delta)y\)
is fixed. The remaining weights are \(\beta_{c'}^0(\delta_\infty^{-1}Y)\).
In the support all the frames \(\delta_\infty^{-1}Y\), evaluated
by Definition~\ref{def:boxed-real-frames}, belong to the required
box. Their pairwise column differences belong to \(Y_\tau^{-1}P_\tau\)
at each place since the definite, \(z\)- and dual coordinates
of \(\delta f\) agree throughout this fibre. Thus the real
dimension of the difference space is at most \(2f\nu n\).
Theorem~\ref{thm:boxed-weighted-frame-count} bounds the sum of
the weights by \(b^{12f\nu n^3}\). Summing and integrating are
justified for nonnegative functions by
Lemmas~\ref{lem:integral-tools} and~\ref{lem:product-integration}.
Including the number of subsets gives the uniform bound
\[
 \int H_{\mathcal O}(uh(A)y)\dd A\le
       i_\gamma^k b^{12f\nu n^3}\prod_\tau\int K_\tau(A_\tau)\dd A_\tau ,
\]
where the left integral is over the unit disks and those on
the right are unrestricted. More precisely, when different
subsets use different \(K_\tau\), use their common numerical
bound from Theorem~\ref{thm:real-frame-volume-majorant} on the right.

Write \(\ell_{\tau i}=|d_{\tau i}|^2\), \(e=j-d\).
At each \(\tau\) this numerical bound is at most
\[
 \left(\frac{(2\theta^2)^j}
       {|\tau D_j|\prod_{i\le j}\ell_{\tau i}}\right)^{k+n}
 \left(\frac{(b/g)^{2d}\prod_{i\le d}\ell_{\tau i}}
                 {\theta^{2d}}\right)^k .
\]
We used \(\pi^{-d}/d!\le1\). When \(e=1\) the only additional
factor in the first product is \(\ell_{\tau n}=N^{-1}\).
Thus after taking the product over \(\tau\), the factors with
exponent \(k\) give
\[
            m_\gamma^{-k/2}
               \big(2^j\theta^{2e}(b/g)^{2d}\big)^{kf}
\]
by Definition~\ref{def:arithmetic-relation-indices}. The base
in parentheses is at most \((b^4g)^{-d}\) eventually.
For \(j=d\) this follows from \(2^d\le(g/b^6)^d\).
For \(j=n,\ d=n-1\), it follows from
\(2^n\theta^2\le b^{14(n-1)}\le(g/b^6)^{n-1}\),
since \(n\to\infty,\ \theta=b^{6n+2}+1,\ b\to\infty\)
and \(g\ge b^{20}\) by Theorem~\ref{thm:data}.

Multiplying the remaining product with exponent \(n\) by
\(m_\gamma^{1/2}\) contributes at most \(b^{O(fn^3)}\).
In fact \(\prod_\tau|\tau D_j|\ge1\) by integrality, and
its exponent after this multiplication is \(1-n\). The
additional norm factor \(N^{-ef}\) is at most one. Each
\(\ell_{\tau i}^{-1}\le N\le b^{60n}\) by
Lemma~\ref{lem:individual-kappas}. The bound follows from
\(j\le n\) and \(\theta\le b^{6n+3}\).
Finally the packing loss \(b^{12f\nu n^3}\) and this last
\(b^{O(fn^3)}\) together are at most \(b^{\eta kf}\) eventually
by \(k/n^9\to\infty\). Apply Lemma~\ref{lem:frame-stabilizer-average}
to obtain the result.
\end{proof}

\begin{definition}\label{def:actual-orbit-size}
With \(M_0\) here denoting the number \(M_h\) at \(h=0\)
in Definition~\ref{def:summed-pair-weights}, put \(d=\min(j,\nu)\).
For \(x>0\) define
\[
  D_j^\circ=
  \begin{cases}
     T_0^{13-32j} E_1,&0\le j\le\nu,\\
     E_2,&j=n,
  \end{cases}
  \qquad
  a_j(x)=T_0^2
      \left(\frac{D_j^\circ}{M_0 g^{2dkf}x^{k+1}}\right)^{1/2}.
\]
\end{definition}

\begin{theorem}\label{thm:actual-small-discriminant-orbit}
Suppose \(\gamma\) has level \(j\) and
\(b_\gamma\le N^f m_\gamma\), with the quantities of
Definition~\ref{def:congruent-relation-gauge}. At all sufficiently
large arithmetic scales its entire pointwise frame-stabilizer
orbit has
\[
                 \sum_{\delta\in\mathcal O} a_\delta(c,c')
                                  \le\mathfrak v a_j(G_\gamma).
\]
In particular this holds for every singular orbit.
\end{theorem}
\begin{proof}
Inversion exchanges the two frame matrices by conjugate transpose;
hence \(D_j(\gamma^{-1})=\bar D_j(\gamma)\) and \(m_{\gamma^{-1}}=m_\gamma\).
The level is unchanged by the symmetric shared-flag condition of
Theorem~\ref{thm:rational-agreeing-levels} and the invariant
of Definition~\ref{def:rational-frame-invariants}. Apply
Proposition~\ref{prop:actual-thin-schur-integral} with \(\eta=1/2\)
in both orientations. The product bound of
Lemma~\ref{lem:two-sided-schur-orbit}, after taking its square
root and using \(i_\gamma,i_{\gamma^{-1}}\ge1\), gives at most
\[
 \mathfrak v J_c T_0^2 (b^4g)^{-dkf}
     \left(\frac{i_\gamma i_{\gamma^{-1}}}{m_\gamma}\right)^{(k+1)/2}.
\]
The ratio in parentheses equals \(N^f/G_\gamma\) by hypothesis.
To bound its square by \(\mathfrak v^2 a_j(G_\gamma)^2\) it
suffices after cancellation to prove
\(E_1 J_c N^f T_0^{-32d}\le D_j^\circ\), since
\(M_0=|\bar S|C N^{-kf}\) and \(E_1=|\bar S|C J_c\).
Lemma~\ref{lem:logarithmic-counts} gives \(J_c\le T_0^8\),
and \(N^f\le T_0\) eventually by Lemma~\ref{lem:individual-kappas}.
At a singular level \(d=j\) this gives the prescribed power.
At \(j=n\) use \(E_2\ge T_0^{44}E_1\) from
Theorem~\ref{thm:two-energies}. Finally \(\mathfrak d=0\)
at each singular level, hence \(b_\gamma=0\) and the hypothesis
holds there.
\end{proof}

\section{Normalization of the test on the regular real orbit}

\begin{lemma}\label{lem:open-cell-product-test}
Suppose \(\gamma\) is near and \(\mathfrak d(\gamma)\ne0\).
For each \(\tau\) use \(\mathcal T_\tau,t_\tau,\lambda_\tau\)
of Theorem~\ref{thm:actual-bessel-open-cell} for
\(\gamma_{\infty,\tau}\) with \(\rho=b^{2n^2}\).
Under the product of the probability laws of the \(t_\tau\),
for every bounded Borel function \(\phi\) on \(G_\infty\),
\[
 \mathbb E_t\left\|
   \int_{\prod_\tau\mathcal T_\tau}\phi(y)
       F_\infty(\gamma_\infty t y) F_\infty(t y)^*\dd y
           \right\|_{\rm HS}^2
           \le \int|\phi(y)|^2\dd\!\bigotimes_\tau\lambda_\tau(y).
\]
\end{lemma}
\begin{proof}
With \(\ell_{\tau i}=|d_{\tau i}|^2\) put \(l_\tau=\prod_{i=1}^n\ell_{\tau i}\).
Taking determinants in the scaled pairings of
Proposition~\ref{prop:thin-orbit-definite-cosets} gives
\(\Delta(Z_\tau^{\rm sc})=l_\tau^2\,\tau\mathfrak d(\gamma)>0\):
the missing last factor in the penultimate minor is \(N^{-1}\).
Positivity follows from Theorem~\ref{thm:rational-frame-discriminant}.
Nearness supplies the two bounded frames required by
Theorem~\ref{thm:actual-bessel-open-cell} at every place.

Regard each integral in that theorem as a linear operator into
the Hilbert space of square-integrable Hilbert--Schmidt matrices
over the law of \(t_\tau\). Its norm on bounded tests in
\(L^2(\lambda_\tau)\) is at most one. On a finite sum of
products of bounded one-place tests, the operator of this lemma
is their tensor product by Lemma~\ref{lem:product-integration}.
A contraction acting in one tensor factor remains contractive
on finite sums: in an orthonormal basis of the span of the
other factors, expand and add the squared norms of its coefficients.
Applying the one-factor operators successively proves the
inequality on these finite sums.

For passage to a general bounded test, on each group take
the finite measure \(\lambda_\tau+{\bf1}_{\mathcal T_\tau}\dd y_\tau\).
Finite linear combinations of indicators of rectangles are
dense in \(L^2\) of their product. Indeed the sets approximable
in measure by finite unions of rectangles form a sigma-algebra:
complements and finite unions preserve approximation, and for
countable unions first truncate the increasing finite unions
using finiteness of the measure and countable additivity.
They contain all rectangles and hence all Borel sets, since
the finite product has a countable base of rectangles.
Approximate simple functions first and then bounded functions
by subdividing their real and imaginary ranges. This gives
simultaneous \(L^2\) approximation for the two relevant
product measures. The actual integrals against Haar on the
compact product converge uniformly in \(t\), by boundedness
of \(F_\infty\), Cauchy--Schwarz and the finite Haar mass of
that compact set. This identifies the tensor-product limit
and proves the stated inequality.
\end{proof}

\begin{proposition}\label{prop:actual-open-cell-mass-budget}
In Lemma~\ref{lem:open-cell-product-test}, write \(B=b_\gamma\).
Uniformly for the near regular terms at all sufficiently large
scales one has
\[
    b^{-k^3f}\le a_n(B),\qquad
    \prod_\tau\lambda_\tau(G_1)
                        \le \tfrac14\,\mathfrak v^2 a_n(B)^2 .
\]
\end{proposition}
\begin{proof}
Integrality and total positivity of the nonzero \(\mathfrak d\)
give \(B\ge1\) by the integral norm argument of
Theorem~\ref{thm:rational-frame-discriminant}. Also
\(\tau\mathfrak d\le|\tau D_n|^2\), so
\(B\le b^{11fn^3}\) by Proposition~\ref{prop:near-minor-coordinate-bounds}.
In Definition~\ref{def:actual-orbit-size}, \(E_2/M_0\ge
E_1/M_0=J_c N^{kf}\ge1\) by Theorem~\ref{thm:two-energies}.
Since \(g\le b^{22}\), the negative base-\(b\) logarithm of
\(a_n(B)\) is at most
\(f(22\nu k+11n^3(k+1)/2)\le k^3 f\) eventually.
Here \(k/n^9\to\infty\) by Lemma~\ref{lem:individual-kappas}.
This proves the first assertion.

Use \(l_\tau,\ell_{\tau i}\) of the preceding proof, put
\(L=\prod_\tau l_\tau\), and \(h_*=15000n^3\).
Theorem~\ref{thm:actual-bessel-open-cell} and
\(\prod_\tau\Delta(Z_\tau^{\rm sc})=L^2 B\) bound the product of
the measures by
\[
       \frac{J_c^2\,2^{(k-1)f}}{\mathcal D C_{\rm op}^{\,f}}\,
          \rho^{9nf h_*}(L^2B)^{-(k-1+h_*)}.
\]
The degree and dimension products here are those of
Definitions~\ref{def:packet-amounts} and~\ref{def:two-energy-numbers}.
Divide by \(\mathfrak v^2 a_n(B)^2\), using
\(E_2=(C J_c)^2/(\mathcal D\mathfrak m_r)\).
The resulting bound is
\[
 \frac{\mathfrak m_r M_0 g^{2\nu kf}}
      {C^2 C_{\rm op}^{\,f}\mathfrak v^2}\,
 T_0^{-4}2^{(k-1)f}\rho^{9nf h_*}
                  L^{-2(k-1+h_*)}B^{2-h_*}.
\]
Since \(B\ge1,\ h_*\ge2\), the last \(B\)-factor is at most one.
By the definitions in Theorem~\ref{thm:packet-normalization},
the initial fraction is exactly
\[
 \frac{P_{\rm pkt}}{Z_{\rm pkt}}\,
        \frac{V_E^n Z_{\rm sc}^2 g^{2\nu kf}}
                         {u\,c_*N^{kf}} .
\]
Here \(M_0=|\bar S|C N^{-kf}\), so this identity follows by
canceling \(C,|\bar S|\) in the displayed packet formulas.
That theorem bounds \(P_{\rm pkt}/Z_{\rm pkt}\le2\) eventually.
One has \(g^{2\nu kf}/u=g^{-\nu(2\nu+1)f}\le1\).
The formula for \(c_*\) in Theorem~\ref{thm:top-siegel-bound} gives
\[
       \frac{V_E^n}{c_*N^{kf}}
          =\frac{V_E^{n+k}}{V_F}
                    \left(\frac{(k-1)!}{\pi^k}\right)^f
                      \le b^{o(kf)} .
\]
Indeed \(V_E=({\rm rd}_E/2)^f,\ V_F={\rm rd}_E^{f/2}\ge1\)
by Theorem~\ref{thm:packet-normalization}, and
\(\log{\rm rd}_E,\log k=O(1+\log x)=o(\log b)\) by
Propositions~\ref{prop:level-prime} and~\ref{prop:thick}.
Also \((k-1)!\le k^k\) and \(n\le k\).

The remaining scale factor is
\[
    Z_{\rm sc}^2 L^{-2(k-1+h_*)}
         =\prod_{\tau,i}\ell_{\tau i}^{\,2(1+2(n-i)-h_*)}.
\]
Every exponent is negative. Since \(\ell_{\tau i}^{-1}\le
N\le b^{60n}\), this product is at most \(b^{O(fn^5)}\).
Also \(\rho^{9nf h_*}=b^{O(fn^6)}\). These are \(b^{o(kf)}\)
by Lemma~\ref{lem:individual-kappas}, and
\(2^{(k-1)f}=b^{o(kf)}\). Finally \(T_0^{-4}=b^{-kf}\).
The ratio under consideration is at most \(b^{-(1-o(1))kf}\)
and thus at most \(1/4\) eventually.
\end{proof}

\section{Bounds for the cut amplitudes on a fixed level}

\begin{lemma}\label{lem:regular-left-action-and-weight}
Fix \(c,s,c'\) and a set \(B\) of regular rational elements
stable under left and right rational unipotents. For
\(\delta\in B\cap\mathcal T(c,s,c')\) one has
\(\Gamma_c\cap\delta\Gamma_{c'}^0\delta^{-1}=\{1\}\).
In every \(\Gamma_c\delta\Gamma_{c'}^0\) choose one element
\(\delta_R\), taking a near representative of
Definition~\ref{def:near-labelled-terms} if the orbit contains
a near element. On \(B\cap\mathcal T(c,s,c')\) set
\[
            p(\delta,y)={\bf1}_{\delta=\delta_R}
                                     \omega_{c'}(\delta_\infty^{-1}y)
\]
with \(\delta_R\) chosen on the orbit of \(\delta\) and
\(\omega_{c'}\) as in Proposition~\ref{prop:boxed-arithmetic-cutoff}.
Then for every \((\delta,y)\),
\[
         \sum_{u\in\Gamma_c,\,v\in\Gamma_{c'}}
                  p(u\delta v^{-1},u_\infty y)=1 .
\]
Moreover \(\omega_{c'}(h y)=\omega_{c'}(y)\) whenever
\(h\in\prod_\tau H_0\), the pointwise real frame stabilizer.
\end{lemma}
\begin{proof}
If \(u\) is in the stated intersection, it fixes \(\delta f\)
pointwise. Let \(A\) be its upper unitriangular action on the
original rational frame. Pairing the two frames gives
\(A^{-*} Z(\delta)=Z(\delta)\). Since
\(\mathfrak d(\delta)>0\) at infinity,
\(Z(\delta)\) is invertible by Theorem~\ref{thm:rational-frame-discriminant}.
Thus \(A=I\). The middle-right block \(C_0\) of \(u\) as in
Proposition~\ref{prop:thin-orbit-definite-cosets} now satisfies
\(C_0 P_Z=0\) since \(u\delta f=\delta f\).
The matrix \(P_Z\) has full row rank as it consists of the
first \(\nu\) rows of an invertible matrix. Hence \(C_0=0\)
and the only possible off-diagonal block of \(u\) is the
top-right block \(H\). The same fixed-frame equation gives
\(H P_Z=0\), so \(u=1\).

Fix a left coset \(\Gamma_{c'}^0v_*\subset\Gamma_{c'}\).
As \(u,v\) vary with \(v\) in this coset, \(u\delta v^{-1}\)
runs over precisely one orbit of the type used for \(\delta_R\),
each element exactly once by the trivial intersection just proved,
applied at \(\delta v_*^{-1}\). Thus there is one contributing
pair and its weight is \(\omega_{c'}(v_{*,\infty}\delta_\infty^{-1}y)\)
by the pointwise invariance in Proposition~\ref{prop:boxed-arithmetic-cutoff}.
Summing over the cosets gives one by its construction.
Nearness is constant on the smaller double cosets of the two
arithmetic pointwise stabilizers, so the specified choice is possible.

For the last assertion recall that \(\omega_{c'}\) is the
indicator \(J\) divided by the sum of its translates in that
proposition. At each place \(h\in H_0\) fixes the frame
pointwise and belongs to the graph projection defining \(S\)
by Lemma~\ref{lem:frame-point-stabilizer}. Consequently \(S(hy)=S(y)\).
For every arithmetic unipotent \(v\), \(v^{-1}F_0\) is a
linear combination of \(F_0\) by Lemma~\ref{lem:bessel-frame-action}.
Thus \(h^{-1}v^{-1}F_0=v^{-1}F_0\), so \(J(vhy)=J(vy)\);
also \(S(vhy)=S(vy)=S(y)\). These identities prove the
invariance of the quotient defining \(\omega_{c'}\).
\end{proof}

\begin{proposition}\label{prop:regular-global-orbit-bound}
With \(B\) as in Lemma~\ref{lem:regular-left-action-and-weight},
one has
\[
 \|A^{V_\infty}_B(c,s,c')\|_{\rm HS}
      \le \mathfrak v\!
        \sum_{\substack{\gamma\in B\cap\mathcal T(c,s,c')\\
                         \gamma\ \text{near chosen representative}}}
                                           a_n(b_\gamma).
\]
\end{proposition}
\begin{proof}
Write \(p_0(\delta,y)=\beta_c(y)\beta_{c'}(\delta_\infty^{-1}y)\).
It sums to one under the same action as \(p\) by the two
partitions of Proposition~\ref{prop:boxed-arithmetic-cutoff}.
Counting measure times Haar measure is invariant under the
action. For any nonnegative invariant \(K(\delta,y)\), insert
the sum of translates of \(p\) into \(\int p_0 K\), change
variables termwise, and sum the translates of \(p_0\). This
gives \(\int p K=\int p_0 K\) by
Lemma~\ref{lem:integral-tools}. It also holds for integrable
vector-valued \(K\) by applying it to its norm first. Full and
cut kernels are invariant under these arithmetic unipotents
by their transformation laws and Proposition~\ref{prop:full-phase}.
Apply the identity to the two kernels, their norms and the
norm of their difference. The relevant integrals are finite by
Theorem~\ref{thm:full-global-truncation}. For the new weight
the total difference integral is still at most
\(\mathfrak v b^{-k^3 f}/2\), so this bounds the nonnegative
contribution of any one chosen \(\delta_R\) as well.

If the cut integrand with weight \(p\) is nonzero at \(y\),
then \(y\in\mathcal E\) and \(\delta_{R,\infty}^{-1}y\in\mathcal B\).
Indeed the latter point is in \(\mathcal E\) by cut support and
its positive \(\omega_{c'}\) enforces the frame bounds of
Proposition~\ref{prop:boxed-arithmetic-cutoff}. Choose \(u\in\Gamma_c\)
so that \(u_\infty y\in\mathcal B\) by
Lemma~\ref{lem:boxed-arithmetic-shears}. Thus \(u\delta_R\)
is near and in the same left-right orbit, which by the rule
of choice means that \(\delta_R\) itself is a near chosen
representative. Only these terms need to be estimated.

Fix one such \(\gamma=\delta_R\). Apply
Lemma~\ref{lem:open-cell-product-test} with \(\phi=\omega_{c'}\).
Its values lie in \([0,1]\), so some \(t\) in the probability
construction makes the norm of the displayed integral at most
\(\big(\prod_\tau\lambda_\tau(G_1)\big)^{1/2}\).
Indeed otherwise its squared norm would exceed this bound
squared on a set of full probability, and integration would
contradict the lemma. The compact-support implication in
Theorem~\ref{thm:actual-bessel-open-cell} holds for every such \(t\).

In the cut integral at \(\gamma\) change variables to
\(y=\gamma_\infty t z\). The weight becomes
\(\omega_{c'}(t z)=\omega_{c'}(z)\) by
Lemma~\ref{lem:regular-left-action-and-weight}. Wherever the
resulting integrand is nonzero, \(t z\in\mathcal E\) and
\(z\) has all real frame columns bounded by \(\rho\), since
\(t\) fixes the frame and the weight is positive.
Also \(S_\tau((\gamma_\infty t z)_\tau)\le R\le\rho\).
Hence \(z\in\prod_\tau\mathcal T_\tau\) by the cited compact-support
implication. On this product replace the cut kernel by the
full one. The norm of the integrated error is at most
\(\mathfrak v b^{-k^3f}/2\) by the difference integral at
this representative, since restricting an integral of a norm
only decreases it. The full integral there has the bound
just obtained from the choice of \(t\). Proposition~\ref{prop:actual-open-cell-mass-budget}
bounds the sum of these two contributions by \(\mathfrak v a_n(b_\gamma)\).
Finally, the weighted series giving \(A_B^{V_\infty}\) is
absolutely convergent. Sum its terms in the new weight and
enlarge the chosen representative list to all the near representatives
of the smaller double cosets, which includes it.
\end{proof}

\begin{theorem}\label{thm:actual-level-orbit-bound}
Suppose a set \(B\) of rational terms is stable under rational
unipotents on both sides, and every member has a fixed level
\(j\) in Definition~\ref{def:arithmetic-relation-indices}.
For agreeing \(s\) and all sufficiently large arithmetic scales,
\[
 \|A_B^{V_\infty}(c,s,c')\|_{\rm HS}
     \le\mathfrak v
        \sum_{\substack{\gamma\in B\cap\mathcal T(c,s,c')\\
                         \gamma\ \text{near chosen representative}}}
                                          a_j(G_\gamma).
\]
At nonagreeing \(s\) the left side is zero.
\end{theorem}
\begin{proof}
Finite disagreement is covered by
Corollary~\ref{cor:full-bessel-identity}. Split \(B\) at
\(b_\gamma=N^fm_\gamma\), assigning equality to the smaller
side. The discriminant and the leading minor, and hence the
two sides, are invariant under rational unipotents by
Theorem~\ref{thm:rational-frame-discriminant}.
On the smaller side partition by
\(\Gamma_c^0,\Gamma_{c'}^0\) double cosets and apply
Theorem~\ref{thm:actual-small-discriminant-orbit} to their
absolute sums. Cosets without near representatives give zero
by Lemma~\ref{lem:two-sided-schur-orbit}. The larger side
consists entirely of regular terms; there \(d=\nu,\ j=n\)
and both row indices are one by
Definition~\ref{def:arithmetic-relation-indices}. Hence \(G_\gamma=b_\gamma\).
Proposition~\ref{prop:regular-global-orbit-bound} applied
to this stable subset completes the proof by the triangle
inequality.
\end{proof}

\section{The contribution from pairs outside the regular diagonal}

\begin{lemma}\label{lem:orbit-sizes-pair-comparison}
Suppose two near chosen representatives \(\gamma,\delta\) of
level \(j\) belong to the same \(\mathcal T(c,s,c')\), with
\(s\) agreeing. Use \(i,l,h,T\) from
Lemma~\ref{lem:common-definite-relation} and
Definition~\ref{def:summed-pair-weights}, so \(h=k-\dim T\).
Let \(X\ge1\) be a lower bound on the product of the norms
of distinct marked primes at each of which \(s\) fixes a
nonzero reduced frame vector. Set \(X_l=X\) for \(l=2\),
and \(X_l=1\) otherwise. Put \(e=1\) at \(j=n\) and \(e=0\)
otherwise. Then
\[
 \frac{a_j(G_\gamma)a_j(G_\delta)}
      {\sqrt{w(\gamma,\delta;c)w(\delta,\gamma;c)}}
 \le T_0^4 D_j^\circ (2T_0^2)^e Z_0^i X_l^{-(k+1)} ,
\]
where \(Z_0=2(b/g)^{kf}T_0^2\) as in
Definition~\ref{def:two-energy-numbers}.
\end{lemma}
\begin{proof}
The common subspace \(T\) is symmetric in the two terms.
Abbreviate \(d_T=d_c(T),\ P_1=P_\gamma(T),\ P_2=P_\delta(T)\)
from Definition~\ref{def:congruent-relation-gauge}.
Multiplying the two ordered weights gives
\[
 \sqrt{w(\gamma,\delta;c)w(\delta,\gamma;c)}
      =\frac{1}{M_h\,g^{2fdk}\,d_T (P_1P_2)^{k/2}} .
\]
All indices in \(P_1,P_2\) are positive integers, so
\(d_T\le\sqrt{P_1P_2}\). Substitution of
Definition~\ref{def:actual-orbit-size} bounds the ratio in the
statement by
\[
       T_0^4 D_j^\circ (2T_0^2)^h
                    \left(\frac{P_1P_2}{G_\gamma G_\delta}
                         \right)^{(k+1)/2}.
\]
By Theorem~\ref{thm:actual-congruent-pair-gauge} the last
power is at most \((2/g)^{fi(k+1)}X_l^{-(k+1)}\), and \(h\le i+e\).
Here Theorem~\ref{thm:marked-middle-evaluations} applies with
the chosen prime set, whose product is at least \(X\).
Since \(g\ge2,\ b\ge2\), we may increase the factor
\((2/g)^{fi(k+1)}\) to \((b/g)^{kfi}\). Combining its \(i\)
powers with \(2T_0^2\) gives the claimed \(Z_0^i\).
\end{proof}

\begin{theorem}\label{thm:full-energy-from-regular-diagonal}
Let \(B\) be a set of rational group elements stable under
rational unipotents on both sides. Fix \(\Omega\subset\mathcal G\)
and \(X\ge1\). At every agreeing \(s\in\Omega\) assume the
marked fixed-vector condition of Lemma~\ref{lem:orbit-sizes-pair-comparison}
with lower bound \(X\). Let \(\mathcal R\) be a collection of
middle relations, and let \(\mathcal L\) denote the collection
of all rational middle relations. Choose \(H\ge E_1\) satisfying
the following two conditions.
\begin{enumerate}
\item If an agreeing rational term of \(B\) at level zero lies
in \(\mathcal T(c,s,c')\) for an agreeing \(s\in\Omega\), then
\(H\ge T_0^{20}E_1\).
\item If \(\mathcal R\ne\mathcal L\), then
\(H\ge T_0^{12}E_2/X^k\).
\end{enumerate}
For \(\Lambda\in\mathcal L\) write
\(B_\Lambda=\{\gamma\in B:\mathfrak d(\gamma)\ne0,\
\mathcal L_\gamma=\Lambda\}\).
Suppose \(w_0\ge0\) satisfies
\[
        \sum_{\Lambda\in\mathcal R}
                  \mathcal E_\Omega(A_{B_\Lambda}^{V_\infty})\le w_0 .
\]
Then, for every fixed \(D>0\), at all sufficiently large scales
uniformly in the specified families,
\[
        \mathcal E_\Omega(A_B^{F_\infty})
                        \le(1+b^{-D})w_0+b^{-D} H .
\]
\end{theorem}
\begin{proof}
At nonagreeing reductions both amplitudes vanish by
Corollary~\ref{cor:full-bessel-identity}. For the cut amplitudes
discard the rationally disagreeing orbits: their individual
integrals are zero by Theorem~\ref{thm:incoherent-rational-orbit-zero}
and the orbit identity in Proposition~\ref{prop:boxed-arithmetic-cutoff}.
It is enough to discard such terms at singular levels. For the
remaining singular terms the levels \(j=0,\ldots,\nu\) are
exhaustive by Theorems~\ref{thm:rational-agreeing-levels}
and~\ref{thm:rational-frame-discriminant}.
These level sets are rationally unipotent stable: left and
right flag action is upper unitriangular, so preserves the
heights and leading coefficients of shared vectors. Rational
agreement is preserved since the two characters are homomorphisms
and conjugating their common unipotents just conjugates within
their respective groups. The regular level and every middle
relation are stable by Theorem~\ref{thm:rational-frame-discriminant}
and Proposition~\ref{prop:global-middle-relation}.

At any one level group the cut amplitude at each \(c,s,c'\)
by its middle relation. All these coefficient series are
absolutely convergent by Theorem~\ref{thm:full-global-truncation}.
Expand the squared norm of their sum into inner products and
bound each off-diagonal term by the product of the two norms.
At the regular level keep the squares with equal relation in
\(\mathcal R\) exactly. At singular levels keep none. For every
remaining ordered pair of relations, Theorem~\ref{thm:actual-level-orbit-bound}
bounds the product by the sum over corresponding near chosen
term pairs of \(\mathfrak v^2 a_j(G_\gamma)a_j(G_\delta)\).
All upper bounds are nonnegative and sum by
Lemma~\ref{lem:integral-tools}. Each pair of terms determines
its relations and its finite reduction once the two column
classes are fixed. Thus it occurs at most once in this upper
bound, even on an omitted diagonal.

The factor multiplying the geometric mean in
Lemma~\ref{lem:orbit-sizes-pair-comparison} is at most \(H/T_0^3\)
on these pairs, eventually. At \(j=0\) it is at most
\(T_0^{17}E_1\) by \(Z_0\le1\) in Theorem~\ref{thm:two-energies};
if such a pair occurs the first assumption applies. For
\(1\le j\le\nu\) it is at most \(T_0^{17-32j}E_1\).
At \(j=n\) with unequal relations \(i\ge1\) by
Theorem~\ref{thm:actual-congruent-pair-gauge}, so it is at most
\(2 T_0^6 E_2 Z_0\le2T_0^{-12}E_1\) by
Theorem~\ref{thm:two-energies}. On an omitted equal regular
relation \(i=0,\ l=2\) by Theorem~\ref{thm:actual-congruent-pair-gauge}.
Here the bound is \(2T_0^6 E_2/X^{k+1}\le H/T_0^3\) eventually
by the second assumption and \(X\ge1\).

Use \(\sqrt{uv}\le(u+v)/2\) from \((\sqrt u-\sqrt v)^2\ge0\).
The sum of the geometric means over the relevant ordered pairs
is at most the sum over all ordered pairs of one level of
their weights. By Theorem~\ref{thm:summed-actual-pair-weights}
with \(\epsilon=1/4\) this is at most \(2T_0\).
After dividing by \(\mathfrak v^2\), the extra energy at each
level is therefore at most \(2H/T_0^2\).

The square root of \(\mathcal E_\Omega\) is a Euclidean norm
on the entries restricted to \(\Omega\), divided by \(\mathfrak v\).
Adding the \(n\) singular pieces to the regular piece and using
\(\sqrt{w_0+a}\le\sqrt{w_0}+\sqrt a\) gives an upper bound
\(\sqrt{w_0}+(n+1)\sqrt{2H}/T_0\) for the cut norm. Add the
full-to-cut error using Theorem~\ref{thm:full-global-truncation}
with exponent \(4D+8\). Since \(E_1\le H\) and \(T_0=b^{kf/4}\),
the full norm is at most \(\sqrt{w_0}+b^{-2D-2}\sqrt H\)
eventually, using Proposition~\ref{prop:thick}. Finally
\((a+b')^2\le(1+u)a^2+(1+u^{-1}){b'}^2\) for \(u>0\)
by the nonnegativity of \((\sqrt u\,a-b'/\sqrt u)^2\).
Use \(u=b^{-D}\) to obtain the result.
\end{proof}

\section{The second moment of the full Poincar\'e columns}

\begin{definition}\label{def:normalized-actual-poincare}
Give the quotient of Definition~\ref{def:adelic-row} the measure
of Lemma~\ref{lem:quotient-measure}, with counting measure on
the finite components. Let \(\mathcal H\) be its complex
\(L^2\)-space, with \(\rho\) the unitary right action of
\(\mathcal G\) through the finite component. This is well-defined
by normality of \(K_Y\) in \(K_Y^0\) and preserves the measure
since it permutes counting measure and commutes with rational
left translation. Let \(Q:\C^{C J_c}\to\mathcal H\) take the
coordinate vector indexed by a class and a tensor coordinate
to its column in Theorem~\ref{thm:actual-full-poincare-columns}
divided by \(\sqrt{\mathfrak v}\). Put
\(a(s)=\|Q^*\rho(s)Q\|_{\rm HS}^2\).
\end{definition}

\begin{theorem}\label{thm:actual-fixing-set-moment}
For the vector conditions and sets \(\Omega(S,D)\) of
Definition~\ref{def:sieve-cell}, uniformly on the arithmetic
window and in the specified conditions, for every fixed \(A>0\)
and all sufficiently large scales,
\[
           \sum_{s\in\Omega(S,D)}a(s)
                    \le (1+b^{-A}) E_2\, b^{-6k|S|}.
\]
The same statement holds with \(a(s)\) defined using the
inverse translate.
\end{theorem}
\begin{proof}
Set \(Y=b^{-6k|S|},\ H=Y E_2\). Since \(|S|\le f\),
\(Y\ge T_0^{-24}\) and Theorem~\ref{thm:two-energies} gives
\(H\ge T_0^{20}E_1\). The prime product in
Theorem~\ref{thm:full-energy-from-regular-diagonal} can be
taken to be \(X=g_0^{|S|}\), choosing at each prescribed pair
the prime on the side of the fixed vector. The inverse-transpose
condition specifies the partner, as in Definition~\ref{def:sieve-cell}.

If \(S\) is the whole set of \(f\) pairs, take no retained
relations and \(w_0=0\) in that theorem. Indeed
\(T_0^{12}E_2/X^k\le b^{-7kf}E_2\le H\)
since \(g_0\ge b^{10}\) by Theorem~\ref{thm:data}.
Otherwise take all relations. For each \(s\in\Omega(S,D)\)
there is a marked component at which all leading minors are
nonzero. Hence every rational term above \(s\) satisfies the
open flag condition by Corollary~\ref{cor:rational-open-block-identification}.
By Theorem~\ref{thm:singular-graph-middle}, on this open set
the relation is the graph of the inverse middle isometry
\(\sigma^{-1}\). In the decomposition of
Theorem~\ref{thm:rational-open-decomposition}, the common
unipotents may be removed from the discriminant. At the
representative \(w_p\sigma\) the Gram matrix is diagonal in
the outer pivots and has last entry \([z,\sigma z]\).
Thus \(\mathfrak d=|D_\nu|^2([z,\sigma z]\overline{[z,\sigma z]}-N^2)\),
where \(|D_\nu|^2\) denotes \(D_\nu\bar D_\nu\).
Regularity implies \(|\tau[z,\sigma z]|>N\) for every \(\tau\)
by Theorem~\ref{thm:rational-frame-discriminant}.
The cut amplitudes for each regular relation are precisely the
sums over the pivots in Theorem~\ref{thm:actual-rational-link-square}
by Corollary~\ref{cor:rational-open-block-identification}.
Changing cutoffs has preserved these stable orbit integrals by
Proposition~\ref{prop:boxed-arithmetic-cutoff}. Consequently the
sum of their diagonal energies is at most
\(w_0=(1+b^{-A-3})H\) by Theorem~\ref{thm:actual-rational-link-square}.

Apply Theorem~\ref{thm:full-energy-from-regular-diagonal}
in both cases, with \(B=\Gamma\) and exponent \(A+3\).
In the second case the resulting multiplier of \(H\) is
\((1+b^{-A-3})^2+b^{-A-3}\le1+b^{-A}\) eventually.
The first case satisfies the same bound. By
Theorem~\ref{thm:actual-full-poincare-columns} the unfolded
matrix entries are the inner products with conjugate linearity
in the first variable. Written as
\(\langle\rho(s)Q e_{c',j},Q e_{c,i}\rangle\), each is the
complex conjugate of the corresponding entry of \(Q^*\rho(s)Q\).
Thus the squared Hilbert--Schmidt norm is the sum of the
squared absolute entries in either convention. Finally
\(Q^*\rho(s^{-1})Q=(Q^*\rho(s)Q)^*\), so the value of \(a\)
is the same at \(s^{-1}\) and \(s\).
\end{proof}

\begin{corollary}\label{cor:actual-second-moment}
For all sufficiently large scales on the window,
\[
             \sum_{s\in\mathcal G}\|Q^*\rho(s)Q\|_{\rm HS}^2
                                            \le(1+9b^{-4})E_2 .
\]
More generally, for every nonnegative weight of positive type
dominated pointwise by \(a\) and supported on \(V(S,D)\),
its total mass is at most \((1+9b^{-4})E_2 b^{-6k|S|}\).
\end{corollary}
\begin{proof}
The weight \(a\) is of positive type by the finite-rank
Hilbert--Schmidt Gram calculation in
Corollary~\ref{cor:concrete-moments}, applied to this \(Q\).
Theorem~\ref{thm:actual-fixing-set-moment} with \(A=4\) supplies
the hypothesis of Theorem~\ref{thm:finite-moment} with \(D_*=E_2\)
since \(1+b^{-4}\le (1-b^{-4})(1+9b^{-4})\) eventually.
That theorem proves both assertions.
\end{proof}

\begin{proposition}\label{prop:actual-bessel-complement-energy}
Let \(S(F)=U(F)H(F)\) as in Corollary~\ref{cor:full-bessel-identity}.
For \(a=0,1\) let \(\Omega_a\) consist of the \(s\in\mathcal G\)
whose components above \(g_a\) are those of a single member
of \(\bar S\). Then for every fixed \(A>0\), eventually
\[
       \mathcal E_{\Omega_a}(A_{\Gamma\setminus S(F)}^{F_\infty})
                                      \le b^{-A}E_1 .
\]
\end{proposition}
\begin{proof}
At every prime above \(g_a\) all members of \(\bar S\) fix the
first frame vector \(e_1\): the definite group acts on \(X_0\)
and the unipotent frame action is upper unitriangular, also at the
conjugate split component by the rational flag definition of
the group in Proposition~\ref{prop:marked-local-level}. Hence
\(X=g_a^{2f}\) is allowed for every \(s\in\Omega_a\).
No rationally agreeing term of level zero above this set can
belong to the complement. Indeed Theorem~\ref{thm:rational-agreeing-levels}
at \(j=0\) says \(\gamma\) preserves the entire frame flag
with each graded scalar equal to \(\lambda\), including
\(\gamma e_1=\lambda e_1\), and \(\lambda\bar\lambda=1\).
The same scalar occurs in the finite position because the
two class matrices fix \(e_1\). It is integral by the integral
row lattices, and congruent to one at all primes above \(g_a\).
As the prime splits completely by Lemma~\ref{lem:split-level},
\((\lambda-1)/g_a\in O_E\). This algebraic integer has every
conjugate of modulus at most \(2/g_a<1\), so vanishes by
Proposition~\ref{prop:boxed-frame-separation}. Thus \(\lambda=1\).
The action on \(P^\perp/P\) fixes \(z\) and restricts to an
isometry of \(X_0\). Removing its diagonal extension leaves
a unipotent by Definition~\ref{def:unip}, so \(\gamma\in S(F)\).

Take \(H=E_1\), retain no regular relations, and take \(w_0=0\)
in Theorem~\ref{thm:full-energy-from-regular-diagonal}. Its first
condition is vacuous by the preceding argument. By
Theorem~\ref{thm:two-energies} and the formula for \(Z_0\),
\[
    E_2/E_1\le\tfrac12(g/b)^{kf}T_0^{-20}
                        \le\tfrac12 b^{16kf}.
\]
The last inequality uses \(g\le b^{22}\) by Theorem~\ref{thm:data}.
Since \(X^k\ge b^{20kf}\) and \(T_0^{12}=b^{3kf}\),
the second condition also holds. The complement is stable on
both sides since \(U(F)\subset S(F)\), so the asserted bound follows.
\end{proof}

\section{Translation by the finite Bessel subgroup}

\begin{definition}\label{def:normalized-bessel-blocks}
Write \(c_0=CJ_c\) for the dimension of the domain of \(Q\) in
Definition~\ref{def:normalized-actual-poincare}. For a union \(B\) of
double cosets of \(S(F)\) in the rational row group \(\Gamma\), let
\(D_B(s)\) be the \(c_0\)-by-\(c_0\) matrix whose \((c,d)\)-block is
\(\mathfrak v^{-1}A_B^{F_\infty}(c,s,d)\) in
Definition~\ref{def:full-arithmetic-amplitudes}. The sums converge
absolutely by Theorem~\ref{thm:full-global-truncation}. Write
\(D_{\rm out}\) for \(D_{\Gamma\setminus S(F)}\).
\end{definition}

\begin{lemma}\label{lem:bessel-class-permutation}
For each \(t\in\bar S\) there is a permutation \(p_t\) of the
column classes with the following property. For every \(d\) there
is an element \(\sigma_d\in S(F)\cap\mathcal T(d,t,p_t(d))\).
All choices for this pair of classes lie in the same right coset
\(\sigma_d\Gamma_{p_t(d)}\), and
\(\sigma_d\Gamma_{p_t(d)}\sigma_d^{-1}=\Gamma_d\).
\end{lemma}
\begin{proof}
Use \(H\) for the definite middle group and choose an integral
finite lift \(k_t\in S_f\) of \(t\) by
Proposition~\ref{prop:marked-local-level}. Factor
\(h_d k_t=u_f a_f\), with \(u_f\in U_f\) and \(a_f\in H_f\).
Indeed two such factors multiply by the semidirect product law,
since \(H_f\) normalizes the unipotent flag, and \(h_d\in H_f\).
The double coset definition of the column classes in
Proposition~\ref{prop:column-dim} gives
\(a_f=\eta_f h_e k\), with \(\eta\in H(F)\) and \(k\in K_c\),
where \(K_c\) denotes the column level. Its extension lies in
\(K_Y\) by Proposition~\ref{prop:marked-local-level}. Hence
\(h_d^{-1}u_f\eta_f h_e=k_t k^{-1}\in k_t K_Y\).
This membership is open as a condition on \(u_f\).
Lemma~\ref{lem:unip-approx} supplies a rational unipotent \(u\)
in the same open set. Then \(\sigma_d=u\eta\) belongs to
\(\mathcal T(d,t,e)\).

The finite condition for \(\mathcal T\) in
Definition~\ref{def:cut-global-columns} respects multiplication
with matching middle class and inversion with the classes
interchanged. This follows by multiplying the level cosets,
since \(K_Y\) is normal in \(K_Y^0\). If two choices lead to
\(e,e'\), their quotient \(\sigma_d^{-1}\sigma'_d\) belongs to
\(S(F)\cap\mathcal T(e,1,e')\). The proof of
Corollary~\ref{cor:full-bessel-identity} identifies this set as
empty for unequal classes and as \(\Gamma_e\) for equal classes.
Thus \(e=p_t(d)\) is unique, with the stated right coset.
If the images of \(d,d'\) agree, use
\(\sigma_d\sigma_{d'}^{-1}\in S(F)\cap\mathcal T(d,1,d')\)
to deduce \(d=d'\). Injectivity implies surjectivity since
the class set is finite. Finally conjugation by \(\sigma_d\)
preserves \(U(F)\) and carries the level condition at
\((e,1,e)\) to that at \((d,1,d)\). Conjugation by its inverse
gives the reverse inclusion.
\end{proof}

\begin{proposition}\label{prop:right-bessel-amplitude}
For each \(t\in\bar S\) there is a unitary matrix \(U_t\) on the
column coordinate space such that for every \(B\) in
Definition~\ref{def:normalized-bessel-blocks} and every
\(s\in\mathcal G\),
\[
                           D_B(st)=D_B(s)U_t .
\]
\end{proposition}
\begin{proof}
For \(\sigma_d=u\eta\) in Lemma~\ref{lem:bessel-class-permutation}
put \(L_d=\chi_\infty(u_\infty)\bigotimes_\tau\pi_\tau(\eta_\tau)\).
Here the definite isometry acts on the positive definite middle
space in the real coordinates of
Theorem~\ref{thm:actual-real-bessel-vector}. Its transformation
law implies \(F_\infty(\sigma_d y)=L_d F_\infty(y)\), also after
the scalar normalizations of
Theorem~\ref{thm:bessel-orthonormal-schur}. The character has
absolute value one and the tensor factors are unitary. Put the
single nonzero block in block row \(d\) of \(U_t\) in column
\(e=p_t(d)\), with value \(L_d\). This matrix is unitary by the
permutation property.

Right multiplication by \(\sigma_d\) bijects
\(B\cap\mathcal T(c,s,d)\) with \(B\cap\mathcal T(c,st,e)\),
because \(B S(F)=B\). In the latter sum write a term as
\(\delta\sigma_d\) and put \(z=\delta_\infty^{-1}y\). The
kernel in Definition~\ref{def:full-arithmetic-amplitudes} then
becomes \(F_\infty(y)F_\infty(z)^* L_d\). The second cutoff
is \(\beta_e(\sigma_{d,\infty}^{-1}z)\). Its sum over the left
action of \(\Gamma_d\) on \(z\) is one by the conjugation
identity of Lemma~\ref{lem:bessel-class-permutation} and the
partition of Proposition~\ref{prop:boxed-arithmetic-cutoff}.
The same holds for \(\beta_d(z)\). Within each right
\(\Gamma_d\)-coset of \(\delta\) the kernel is constant because
the character is trivial on \(\Gamma_d\) by
Proposition~\ref{prop:full-phase}. Summing the nonnegative
cutoffs on that coset therefore gives the same integral with
either cutoff. Absolute convergence permits the regrouping.
This proves the asserted equality block by block.
\end{proof}

\section{The first moment at the marked primes}

\begin{lemma}\label{lem:actual-complement-convolution}
Let \(i\in\{0,1\}\), put \(J_i=\prod_{w\mid g_i}{\rm GL}_r(\mathbf F_{g_i})\),
with \(w\) running through the \(f\) split pairs, and let
\(\iota_i:J_i\to\mathcal G\) insert this factor. For a self-adjoint
idempotent \(x=\sum_{t\in J_i} x_t[t]\in\C[J_i]\) put
\(Z_x=\sum_t x_t D_{\rm out}(\iota_i(t)^{-1})\). For every fixed
\(A>0\), at all sufficiently large scales, uniformly in \(i,x\),
\[
                           \|Z_x\|_{\rm HS}^2\le c_0 b^{-A}.
\]
\end{lemma}
\begin{proof}
Give matrix-valued functions on \(\mathcal G\) the squared norm
\(\sum_s\|f(s)\|_{\rm HS}^2\). Define the convolution by
\[
             L_x f(s)=\sum_t x_t f(\iota_i(t)^{-1}s).
\]
Each translation preserves counting measure, its adjoint is its
inverse, and translations respect multiplication. Thus \(L_x\)
is a self-adjoint idempotent. Its range and kernel are orthogonal,
so \(\|L_x f\|\le\|f\|\). The set \(\Omega_{1-i}\) in
Proposition~\ref{prop:actual-bessel-complement-energy} contains
\(\bar S\) and is invariant under left translation by \(\iota_i(J_i)\).
Apply the contraction to \(f={\bf1}_{\Omega_{1-i}}D_{\rm out}\).
For every \(s\in\bar S\), Proposition~\ref{prop:right-bessel-amplitude}
gives \(L_x f(s)=Z_x U_s\). Its Hilbert--Schmidt norm is
\(\|Z_x\|_{\rm HS}\), since the columns of a unitary matrix form an
orthonormal basis. Summing just over \(s\in\bar S\) on the left
and using Proposition~\ref{prop:actual-bessel-complement-energy}
on the right gives
\[
 |\bar S|\|Z_x\|_{\rm HS}^2\le
       \sum_{s\in\Omega_{1-i}}\|D_{\rm out}(s)\|_{\rm HS}^2
       \le b^{-A}E_1=b^{-A}|\bar S|c_0
\]
by Definitions~\ref{def:full-arithmetic-amplitudes}
and~\ref{def:two-energy-numbers}.
\end{proof}

\begin{lemma}\label{lem:actual-marked-bessel-trace}
With \(J_i,\iota_i\) as in Lemma~\ref{lem:actual-complement-convolution},
let \(U_i\subset J_i\) be the image of the integral unipotents at
the split pairs above \(g_i\), and let \(\xi_i\) denote the
character induced there by \(\chi_g\) of
Definition~\ref{def:finite-bessel-vector}. Define
\(h_i(t)=\xi_i(t)\) on \(U_i\), and \(h_i(t)=0\) elsewhere. Then
\[
 c_0^{-1}\operatorname{tr}D_{S(F)}(\iota_i(t)^{-1})=h_i(t),
 \qquad
 \sum_{t\in J_i}\left|\,|J_i|^{-1}
                 \sum_{j\in J_i}h_i(jtj^{-1})\right|^2\le 2
\]
at all sufficiently large scales.
\end{lemma}
\begin{proof}
For a diagonal block \(d\) suppose \(\delta=u\eta\in S(F)\)
lies above \((d,\iota_i(t)^{-1},d)\). Projection to \(P^\perp/P\)
as in Lemma~\ref{lem:boxed-equal-classes} shows that the definite
\(\eta\) stabilizes the middle lattice in class \(d\). Its conjugate
by \(h_d\) lies in \(K^0\) by Proposition~\ref{prop:marked-local-level}.
At every marked pair above \(g_{1-i}\) it is the identity modulo
\(g_{1-i}\), by the same projection. This element of the finite
definite stabilizer is trivial by the odd-prime torsion argument
of Proposition~\ref{prop:column-dim}, applied at \(g_{1-i}\).
It follows that \(\delta\in U(F)\) and \(t\in U_i\).

Conversely, for \(t\in U_i\), lift \(t^{-1}\) to an integral
finite unipotent with identity at all places outside \(g_i\).
The open level condition on the variable \(u_f\),
\(h_d^{-1}u_f h_d\in k_{\iota_i(t)^{-1}}K_Y\), holds at the
conjugate of this lift by \(h_d\). This conjugation preserves
\(U_f\), so Lemma~\ref{lem:unip-approx} gives a rational
unipotent \(\delta\) in the same open set. As the simple-root
entries are invariant under conjugation by \(h_d\), the character
of \(\delta_f\) is \(\xi_i(t^{-1})\). Indeed it is trivial
on integral unipotents off \(g\) and on the principal unipotent
kernel above \(g\), by the calculation in
Lemma~\ref{lem:global-finite-phase-expansion}. Consequently
\(\chi_\infty(\delta_\infty)=\xi_i(t)\) by
Proposition~\ref{prop:full-phase}. For this translate the
permutation of Lemma~\ref{lem:bessel-class-permutation} fixes
every \(d\), and the corresponding block of \(U_{\iota_i(t)^{-1}}\)
in Proposition~\ref{prop:right-bessel-amplitude} is
\(\chi_\infty(\delta_\infty) I\). Apply that proposition with
\(B=S(F)\), \(s=1\), using \(D_{S(F)}(1)=I\) from
Corollary~\ref{cor:full-bessel-identity}. Summing the diagonal
entries proves the trace formula, including vanishing off \(U_i\).

For the squared sum take \(T\) to consist of all places over
\(g_0g_1\), and \(Z=K_T\) in
Corollary~\ref{cor:finite-integral-congruence-sum}. Its quotient
is \(J_0\times J_1\) by Proposition~\ref{prop:row-group}.
Membership and the phase in its integral unipotent image are
products of the independent local conditions at the split pairs.
Hence at \(\iota_i(t)\) the conjugation average in that corollary
is precisely the displayed average of \(h_i\). The other factor
has argument \(1\) and phase one. Restrict the nonnegative
integral sum to this slice. It is bounded by
\(\mathcal E_T(K_T)/\mu(K_T)\), since the remainder is
nonnegative. The product estimate in the proof of
Theorem~\ref{thm:finite-marked-bottom} bounds this ratio by
\(1+2^{-n-1}\le2\).
\end{proof}

\begin{theorem}\label{thm:actual-marked-first-moments}
With \(Q,\rho\) of Definition~\ref{def:normalized-actual-poincare},
let \(R_i=\rho(\iota_i(e_i))\) for the one-prime good idempotents
of Corollary~\ref{cor:actual-good-share}, extending \(\iota_i\)
linearly. At all sufficiently large scales, uniformly on the window,
\[
 \sum_{j=1}^{c_0}\|(Q^*Q-I)e_j\|\le c_0 b^{-4},\qquad
             \operatorname{tr}(Q^*(1-R_i)Q)\le 2c_0 b^{-4}
                         \quad(i=0,1).
\]
Here \(e_j\) in the sum denotes the standard coordinate basis.
\end{theorem}
\begin{proof}
By Theorem~\ref{thm:actual-fixing-set-moment}'s inner product
calculation \(D_\Gamma(s)=\overline{Q^*\rho(s)Q}\), with the bar
taken entrywise. At \(s=1\) subtract \(D_{S(F)}(1)=I\) from
Corollary~\ref{cor:full-bessel-identity}. The result is
\(D_{\rm out}(1)\), whose squared norm is at most \(c_0 b^{-8}\)
by Lemma~\ref{lem:actual-complement-convolution} at \(x=[1]\).
The sum of the column norms is at most \(\sqrt{c_0}\) times the
Hilbert--Schmidt norm by Cauchy--Schwarz. This proves the first bound.

For the second let \(x=1-e_i\). Its coefficients are invariant
under conjugacy, and \(\sum_t |x_t|^2=x_1\) by taking the
identity coefficient of \(x^*x=x\). The trace identity
\[
 \operatorname{tr}(Q^*(1-R_i)Q)
      =\sum_{t\in J_i}x_t\operatorname{tr}D_\Gamma(\iota_i(t)^{-1})
\]
follows because \(Q^*\rho(s^{-1})Q=(Q^*\rho(s)Q)^*\).
The complementary part has absolute value at most
\(\sqrt{c_0}\|Z_x\|_{\rm HS}\le c_0 b^{-4}\) by
Lemma~\ref{lem:actual-complement-convolution} at \(A=8\) and
Cauchy--Schwarz on the diagonal entries. For the \(S(F)\)-part
use Lemma~\ref{lem:actual-marked-bessel-trace}. In the sum
\(\sum_t x_t h_i(t)\) one can replace \(h_i(t)\) by its
conjugation average since conjugation permutes \(J_i\) and fixes
\(x_t\). Cauchy--Schwarz bounds the result in absolute value by
\(\sqrt{2x_1}\). The one-prime estimate in the proof of
Corollary~\ref{cor:actual-good-share} gives \(x_1\le b^{-8}/2\).
The \(S(F)\)-part therefore contributes at most \(c_0b^{-4}\).
Adding the two bounds proves the theorem.
\end{proof}

\section{Independent translates of the ramified columns}

\begin{definition}\label{def:actual-marked-span}
Let \(D=|\mathcal G|\mathfrak m_r\mathcal D\), a positive real
number, and put \(\epsilon=b^{-4}\), \(\theta=f/(N-1)\).
Let \(R=R_0R_1\) be the product of the two projections of
Theorem~\ref{thm:actual-marked-first-moments}. Identify the columns
with \(\C^{c_0}\) by evaluation at the representatives and by the
orthonormal tensor coordinates of
Definitions~\ref{def:normalized-actual-poincare}
and~\ref{def:full-arithmetic-amplitudes}. Write \(V_{\rm ram}\)
for the subspace \({\rm Col}_{\rm ram}\) of Theorem~\ref{thm:ram-dim}
in these coordinates and \(P\) for its orthogonal projection.
Write \(Q_s=\rho(s)Q\) and
\[
 W_*=\sum_{s\in\mathcal G}RQ_s(V_{\rm ram}),\qquad
 W=R\mathcal P L^2(Y),
\]
where \(\mathcal P\) is the real convolution projection of
Proposition~\ref{prop:actual-real-row-projection}. The quotient
has counting measure at the finite level as in
Definition~\ref{def:marked-spectral-projector}.
\end{definition}

\begin{theorem}\label{thm:actual-ramified-rank-lower}
At every sufficiently large scale on the window one has
\(W_*\subseteq W\), \(\dim W<\infty\), and
\[
 \dim W_*\ \ge\ \frac{(1-5\epsilon-\theta)^2}{1+9\epsilon}\,D
                \ \ge\ (1-19\epsilon-2\theta)D .
\]
\end{theorem}
\begin{proof}
The full columns are reproduced by \(\mathcal P\) by
Theorem~\ref{thm:actual-full-poincare-columns}. The finite and
real translations commute. Each two-prime good tuple is primitive
in at least one coordinate over \(g_0\), so its projection is
contained in \(\mathcal E\) of
Definition~\ref{def:marked-spectral-projector}. Thus
\(W_*\subset W\subset\mathcal P\mathcal E L^2(Y)\), which is
finite dimensional by Theorem~\ref{thm:finite-marked-projected-space}.

All the following operators on \(L^2\) have range in the
finite-dimensional invariant space spanned by \(W\) and the
\(Q_s(\C^{c_0})\); the traces can be taken there. Put
\[
 A=|\mathcal G|^{-1}\sum_s Q_s Q_s^*,\qquad
 A'=|\mathcal G|^{-1}\sum_s RQ_s P Q_s^* R .
\]
Cyclic permutation in the trace and reindexing pairs
\((s,t)\) by \(s^{-1}t\) give, by
Corollary~\ref{cor:actual-second-moment} and
Definition~\ref{def:two-energy-numbers},
\[
  \|A\|_{\rm HS}^2=|\mathcal G|^{-1}\sum_t a(t)
                       \le (1+9\epsilon)c_0^2/D .
\]
One has
\(0\le A'\le RAR\) as quadratic forms since \(P\) is a
projection. For positive semidefinite \(H,K\) their product
has nonnegative trace, as seen in an orthonormal eigenbasis
of \(H\). Expanding \(\operatorname{tr}((A'+(RAR-A'))^2)\)
therefore gives \(\|A'\|_{\rm HS}\le\|RAR\|_{\rm HS}\le
\|A\|_{\rm HS}\). The last inequality follows in a basis
adapted to \(R\) by deleting matrix entries.

Theorem~\ref{thm:ram-dim} gives \(\dim V_{\rm ram}\ge
(1-\theta)c_0\). Expanding in the coordinate basis,
\[
 |\operatorname{tr}(P(Q^*Q-I))|
 \le\sum_j\|(Q^*Q-I)e_j\|\,\|Pe_j\|\le\epsilon c_0
\]
by Theorem~\ref{thm:actual-marked-first-moments}. Since \(R_0,R_1\)
are commuting projections, \(1-R\le(1-R_0)+(1-R_1)\). The same
theorem gives \(\|(1-R)Q\|_{\rm HS}^2\le4\epsilon c_0\).
Commutation with translations and orthogonality now give
\[
 \operatorname{tr} A'=\|RQP\|_{\rm HS}^2
 =\operatorname{tr}(P Q^*Q)-\|(1-R)QP\|_{\rm HS}^2
 \ge (1-5\epsilon-\theta)c_0 .
\]
The last factor is positive eventually since \(\theta\le b^{-8}\)
by the bounds in Corollary~\ref{cor:concrete-moments}. The range
of \(A'\) lies in \(W_*\). Cauchy--Schwarz on its diagonal
in a basis adapted to \(W_*\) gives
\((\operatorname{tr}A')^2\le\dim(W_*)\|A'\|_{\rm HS}^2\).
This proves the first inequality. The second follows by expanding
both sides after multiplication by \(1+9\epsilon\).
\end{proof}

\begin{proposition}\label{prop:actual-numerical-selection}
For all sufficiently large scales and any \(a\in\C^\times\)
there is a set \(T\subseteq\mathcal G\) with
\[
 |T|c_0\le3D\log(1+b^2)
\]
such that at least \(\lceil(1-4/b^2)D\rceil\) of the vectors
\(aRQ_s z\), \(s\in T\), with \(z\) from the short integral
family of Corollary~\ref{cor:short-ram}, are linearly independent.
\end{proposition}
\begin{proof}
Fix a subspace \(F\subseteq W\) and write \(l=\dim F\).
From the orthogonal-projection argument in
Lemma~\ref{lem:proj-moments}, applied in the finite invariant
space used in Theorem~\ref{thm:actual-ramified-rank-lower},
\[
 \dim (RQ_s)^{-1}F\le5\epsilon c_0+\|P_F Q_s\|_{\rm HS}^2.
\]
The average of the last term is \(\operatorname{tr}(P_F A)\),
at most \(\sqrt l\,\|A\|_{\rm HS}\le
c_0\sqrt{(1+9\epsilon)l/D}\) by the proof of that theorem.
Choose \(s\) at most this average, and apply
\(2\sqrt{(1+9\epsilon)l/D}\le1+9\epsilon+l/D\).
The map \(L\) onto \(V_{\rm ram}\) of Theorem~\ref{thm:ram-dim}
has kernel of dimension at most \(\theta c_0\). Pullback by
\(L\) adds at most this amount to an inverse-image dimension.
Multiplication by \(a\) preserves \(F\). Consequently the map
\(a RQ_s L\) has inverse image of \(F\) of dimension at most
\[
       c_0\left((1+l/D)/2+\eta\right),\qquad
                     \eta=b^{-2},
\]
since \(5\epsilon+9\epsilon/2+\theta\le\eta\) eventually.

For the rounding cost we verify \(c_0\le D\).
The proof of Proposition~\ref{prop:actual-bessel-complement-energy}
gave \(E_2/E_1\le b^{16kf}/2\).
Every element of \(\bar S\) fixes the first frame vector at each
of the \(2f\) chosen factors. The full group at each factor acts
transitively on nonzero vectors by basis extension. Thus
\[
 [\mathcal G:\bar S]\ge\prod_{i=0}^1(g_i^r-1)^f\ge b^{18rf}
\]
eventually by Theorem~\ref{thm:data}. The identities of
Definition~\ref{def:two-energy-numbers} give
\(D/c_0=[\mathcal G:\bar S] E_1/E_2\ge1\).

Start with \(F=0\). While \(u=D-\dim F>4\eta D\), adjoin to \(F\)
the whole image of the supplied \(a RQ_s L\). The dimension gain
is at least \(c_0(u/(2D)-\eta)>0\), by the inverse-image bound.
The new deficit \(u'\) therefore satisfies
\[
          u'-2\eta D\le(1-c_0/(2D))(u-2\eta D).
\]
Each step selects a new \(s\). By \(1-v\le\exp(-v)\), at most
\(\lceil(2D/c_0)\log((1-2\eta)/(2\eta))\rceil\) steps suffice
when \(0<\eta\le1/8\). Using \(c_0\le D\), their cost is at most
\(2D\log(1/(2\eta))+D\le3D\log(1+1/\eta)\).
The images of the indicated short vectors span the resulting
space by Corollary~\ref{cor:short-ram}. Extract a basis of it
by successively adjoining vectors outside the span. Its
dimension is an integer at least \((1-4\eta)D\).
\end{proof}

\begin{corollary}\label{cor:actual-selection-upper-only}
If the actual row space \(W\) of Definition~\ref{def:actual-marked-span}
satisfies \(\dim W\le D\), Proposition~\ref{prop:selected-family}
holds for the normalized columns \(Q\), these good projections
and \(d_0=\dim W\), with its cost measured in units of \(d_0\).
\end{corollary}
\begin{proof}
Work in \(W+\sum_s Q_s(\C^{c_0})\), a finite invariant space
by Theorem~\ref{thm:actual-ramified-rank-lower}. Its \(R\)-image
is \(W\), since \(R\) fixes \(W\) and every \(Q_s\) is in the
real projection. The mass hypothesis in Corollary~\ref{cor:concrete-moments}
holds with \(D_*=E_2\) by Theorem~\ref{thm:actual-fixing-set-moment}
and the comparison with the required constant in
Corollary~\ref{cor:actual-second-moment}. The two first-moment
hypotheses hold by Theorem~\ref{thm:actual-marked-first-moments}.
Finally \(d_0 E_2\le D E_2=|\mathcal G|c_0^2\) by
Definition~\ref{def:two-energy-numbers}. The corollary gives the
asserted basis and cost.
\end{proof}

\section{Ray characters and Frobenius at primes of good reduction}

\begin{definition}\label{def:general-signed-ray}
Let \(K\) be a number field of degree \(d\) over \(\Q\), with
\(O=O_K\), and let \(m\ge2\) be an integer. Let \(I_m\) be the group
of fractional ideals with valuation zero at every prime over \(m\).
Let \(C_m\) be its quotient by principal ideals \((z)\) with
\(\tau(z)>0\) at every real embedding and \(z-1\in mO_{\mathfrak p}\)
for each \(\mathfrak p\mid m\). The numerator of an ideal norm in
\(I_m\) and its denominator, in lowest terms, are prime to \(m\),
by Proposition~\ref{prop:general-ideal-factorization}. We write
\(N\mathfrak a\bmod m\) for the resulting unit residue.
\end{definition}

\begin{proposition}\label{prop:general-ray-norm-count}
The group \(C_m\) is finite. The map
\(\mathfrak a\mapsto N\mathfrak a\bmod m\) on \(I_m\) factors
through \(C_m\). For every nontrivial complex character \(\lambda\)
of \(C_m\), the series
\[
  L_m(s,\lambda)=\sum_{\substack{\mathfrak a\subseteq O,\ \mathfrak a\ne0\\
                              (\mathfrak a,m)=1}}
                     \lambda([\mathfrak a])(N\mathfrak a)^{-s}
\]
has an absolutely convergent Euler product for \(\Re s>1\)
and extends holomorphically to \(\Re s>1-1/d\).
The trivial character has a simple pole at one of positive residue.
\end{proposition}
\begin{proof}
For \(z\) integral at the primes of \(m\), the denominator ideal
in Proposition~\ref{prop:general-ideal-factorization} has index
prime to \(m\). It contains its index, so \(z=a/k\) with \(a\in O\)
and a positive integer \(k\) prime to \(m\). The local congruences
in Definition~\ref{def:general-signed-ray} say \(a-k\in mO\),
by comparing valuations above \(m\) and observing integrality at
the other primes. The matrix of multiplication by \(z\) on an
integer basis of \(O\) therefore reduces to the identity modulo
\(m\) when these congruences hold. Its determinant is \(N_{K/\Q}z\)
by Lemma~\ref{lem:number-field-lattice}. This determinant is
positive under the sign condition, because the nonreal embeddings
occur in conjugate pairs. Hence \(N((z))\equiv1\bmod m\).
This proves factorization of the norm. Finiteness follows from
finiteness of the ordinary class group in
Lemma~\ref{lem:bounded-ideal-classes}: among principal ideals
in \(I_m\), two generators of equal signs and equal invertible
residue have quotient satisfying the ray conditions. There are
finitely many possible pairs of these data.

Choose an integral ideal \(\mathfrak b\in I_m\) representing
the inverse of any prescribed ray class. To obtain one, clear
the denominators of an integer basis of a fractional representative
by an integer prime to \(m\), using the denominator argument
above, then raise this positive integer to the order of its
residue in \((\Z/m\Z)^\times\). Its principal ideal represents
the identity in \(C_m\). Now \(\mathfrak b+mO=O\) since no maximal
ideal contains both. Thus the members \(a\in\mathfrak b\) with
\(a\equiv1\bmod mO\) form one translate of
\(\mathfrak b\cap mO=m\mathfrak b\), by prime factorization.
Take \(U\subset O^\times\) to consist of the units with this
same residue and positive real signs. It has finite index by
the finite sign and residue maps. The ideals in the class of
\(\mathfrak b^{-1}\) being counted are \((a)\mathfrak b^{-1}\)
for the nonzero elements of this translate with positive real
signs, modulo \(U\). This description works with empty sign
conditions for a totally imaginary field.

Apply Lemmas~\ref{lem:unit-log-domain} and~\ref{lem:log-sector-count}
to that translate in the positive-sign sector at
\(Y N\mathfrak b\). If \(v\) is the sector volume, \(U_0\) its
unit basis subgroup, and \(V_K\) the covolume of \(O\), the
leading constant in the ideal count is
\(v/(m^d V_K[U:U_0])\). Indeed the lattice covolume is
\(m^d V_K N\mathfrak b\), and each admissible generator orbit
gives exactly \([U:U_0]\) representatives in the sector as
in Theorem~\ref{thm:partial-ideal-counts}. The error is
\(O_{K,m}(1+Y^{1-1/d})\), uniformly over the finitely many
classes. Summing with the character weights and integrating
the counting function as in Corollary~\ref{cor:ray-series-at-one}
gives the continuations and the stated pole. Unique factorization
and the unweighted count justify the Euler products there.
\end{proof}

\begin{lemma}\label{lem:finite-extension-residue-actions}
Let \(L/K\) be a finite Galois extension of number fields and
\(G={\rm Gal}(L/K)\). Outside a finite set of rational primes
the following hold for every prime \(\mathfrak p\) of \(K\).
Put \(k=O_K/\mathfrak p\), \(q=|k|\), and choose a homomorphism
\(h:O_L\to\overline k\) extending reduction. There is a unique
\(\sigma_h\in G\) with
\[
                         h(\sigma_h x)=h(x)^q\quad (x\in O_L).
\]
All such \(h\) form one free transitive \(G\)-set by composition.
The elements \(\sigma_h\), called the arithmetic Frobenius elements
above \(\mathfrak p\), form a conjugacy class. Writing \(e\) for
their common order, \(\mathfrak p O_L\) is a product of
\(|G|/e\) distinct primes of norm \(q^e\). For a Galois intermediate
extension \(K'/K\), restriction gives the Frobenius at the
restriction of \(h\), outside its finite exceptional set as well.
\end{lemma}
\begin{proof}
Choose an integral generator \(u\) of \(L/\Q\) by
Lemma~\ref{lem:number-field-lattice}. The minimal polynomial
\(P\) of \(u\) over \(K\) has coefficients in \(O_K\), being
symmetric functions of integral conjugates lying in \(K\).
Since the extension is normal, its distinct roots are the
\(\gamma u\) for \(\gamma\in G\). Indeed an embedding over \(K\)
is obtained by choosing any root of \(P\), and normality maps
\(L\) onto itself. Thus \(\deg P=|G|\). Let \(M>0\) be the
index of \(\Z[u]\) in \(O_L\), finite by the same lemma. After
inverting \(M\), the rings \(O_L\) and \(O_K[u]\) coincide.
Exclude the divisors of \(M\) and of the nonzero integer norm
of the discriminant of \(P\). For every other \(\mathfrak p\)
we have
\[
 O_L/\mathfrak p O_L \simeq k[X]/(\bar P),
\]
by monic division in \(O_K[X]\), localization and reduction.
The reduction \(\bar P\) has distinct roots since the product
of the squared root differences remains a unit. Homomorphisms
\(h\) correspond to the roots of \(\bar P\) in \(\overline k\).
For any one \(h\) the values \(h(\gamma u)\) are distinct and
exhaust the roots. This proves freeness and transitivity.
The \(q\)-power of \(h\) is another such map, which defines
\(\sigma_h\), and composition with \(\gamma\) replaces
\(\sigma_h\) by \(\gamma^{-1}\sigma_h\gamma\).

Factor \(\bar P\) into irreducibles over \(k\). Coprime factors
give a product of fields by the polynomial Euclidean algorithm.
Each factor of degree \(a\) has as its roots one orbit of
length \(a\) under the \(q\)-power, as in
Lemma~\ref{lem:torus-primary}: successive Frobenius powers
fix its coefficients, and the product over a full root orbit
is fixed by Frobenius and hence has coefficients in \(k\).
That fixed field is \(k\) since \(X^q-X\) has at most \(q\)
roots. The orbit through \(h\) has length \(e\) by freeness.
The kernels of the projections to the field factors are all
the primes above \(\mathfrak p\), and their intersection is
their product by pairwise coprimality. This proves the claimed
factorization and norms. Finally restrict the defining equation
for \(\sigma_h\) to \(O_{K'}\) and use uniqueness in \(K'/K\).
\end{proof}

\section{Degree-one primes with prescribed Frobenius}

\begin{lemma}\label{lem:cyclotomic-prime-norms}
Let \(K\) be any number field, \(p\) a rational prime,
\(\zeta\) a primitive \(p\)-th root of unity, and
\(H={\rm Gal}(K(\zeta)/K)\subset(\Z/p\Z)^\times\) by its action
on \(\zeta\). For every \(a\in H\), after deleting any finite
set of rational primes, one has as \(\sigma\downarrow1\)
\[
 \sum_{\substack{\mathfrak q\ {\rm over\ the\ remaining\ primes}\\
           N\mathfrak q\ {\rm prime},\ N\mathfrak q\equiv a\bmod p}}
            (N\mathfrak q)^{-\sigma}
       = |H|^{-1}\log\!\left(\frac1{\sigma-1}\right)+O_{K,p}(1).
\]
The bounded term may also depend on the deleted set.
\end{lemma}
\begin{proof}
The extension \(L=K(\zeta)\) is Galois since it splits
\(X^p-1\). Its group acts faithfully on \(\zeta\), and \(H\)
is cyclic by the finite multiplicative subgroup argument of
Lemma~\ref{lem:split-level}. If \(B\le H\), its fixed field
\(L^B\) has \([L:L^B]=|B|\). Indeed for an integral primitive
element \(v\) of \(L/\Q\), its \(|B|\) distinct images under
\(B\) are roots of its minimal polynomial over \(L^B\), and
the product of the corresponding monic linear factors belongs
to \(L^B[X]\). It has exactly that degree, and \(v\) generates
\(L\) over \(L^B\). Since \(H\) is abelian, restriction to
\(L^B\) has kernel \(B\). Its image has size \([L^B:K]\),
so \(L^B/K\) is Galois with group \(H/B\). Here equality of
the automorphism count with the degree gives all embeddings
by the separability argument of Lemma~\ref{lem:number-field-lattice}.

Take \(m\) divisible by \(p\), the excluded primes and all
exceptional rational primes in
Lemma~\ref{lem:finite-extension-residue-actions} for \(L/K\)
and its subextensions \(L^B/K\). At every \(\mathfrak q\nmid m\),
reduction of \(\zeta\) has order \(p\), since
\(\prod_{j=1}^{p-1}(1-\zeta^j)=p\) by evaluating
\((X^p-1)/(X-1)\) at \(1\). Hence the arithmetic Frobenius
acts on \(\zeta\) by \(h_{\mathfrak q}=N\mathfrak q\bmod p\).
In particular this norm belongs to \(H\). The elements
\(h_{\mathfrak q}\), \(\mathfrak q\nmid m\), generate \(H\).
Otherwise, if their subgroup were \(B\subsetneq H\), all these
primes would split completely in \(L^B\) by the restriction and
norm formulas of Lemma~\ref{lem:finite-extension-residue-actions}.
The Euler products after omitting the divisors of \(m\) would give
\[
       \zeta_{L^B}^{(m)}(s)
              =\bigl(\zeta_K^{(m)}(s)\bigr)^{[L^B:K]} .
\]
Both zeta functions have simple poles of positive residue at
one by Corollary~\ref{cor:ray-series-at-one}. Removing a finite
set of Euler factors multiplies by a function holomorphic
and nonzero at one. The equality on real \(s>1\) contradicts
these pole orders.

For every character \(\chi\) of \(H\) let
\(\lambda_\chi(\mathfrak a)=\chi(N\mathfrak a\bmod p)\)
on \(I_m\). This is a character of the signed ray group by
Proposition~\ref{prop:general-ray-norm-count} and \(p\mid m\).
The generation just proved makes it nontrivial if \(\chi\ne1\).
For \(\Re s>1\) one has
\[
                     \prod_{\chi}L_m(s,\lambda_\chi)
                                      =\zeta_L^{(m)}(s).
\]
In fact at \(\mathfrak q\nmid m\) with Frobenius of order \(e\),
its values under all characters of the cyclic \(H\) run through
the \(e\)-th roots of unity, each \(|H|/e\) times. Multiplying
the factors gives \((1-(N\mathfrak q)^{-es})^{-|H|/e}\),
the local product for \(L\) by
Lemma~\ref{lem:finite-extension-residue-actions}.
Now divide by \(L_m(s,1)=\zeta_K^{(m)}(s)\). The ratio of zeta
functions has a finite positive limit at one. Every remaining
factor on the left is holomorphic there by
Proposition~\ref{prop:general-ray-norm-count}. Thus each such
factor has a nonzero value at one.

For real \(\sigma>1\) expand the Euler logarithms and use
orthogonality of the characters of \(H\), obtained by summing
the geometric series on a generator as in Lemma~\ref{lem:characters}.
The terms with exponent at least two are bounded uniformly by
\(2\sum_{\mathfrak q}(N\mathfrak q)^{-2}\le2\zeta_K(2)\).
The principal logarithm is
\(\log(1/(\sigma-1))+O(1)\). All nonprincipal Euler logarithms
are bounded as \(\sigma\downarrow1\) by their nonzero values and
the branch argument in Corollary~\ref{cor:ap}. This proves the
formula with all prime ideals of \(K\) prime to \(m\).
Above a rational prime \(\ell\) there are at most \([K:\Q]\)
primes by \(N(\ell O_K)=\ell^{[K:\Q]}\) and factorization.
Those with norm other than \(\ell\) have norm at least \(\ell^2\),
so their whole sum is bounded by
\([K:\Q]\sum_{\ell\ge2}\ell^{-2}\). A finite change of the
omitted primes changes the result by a bounded amount.
\end{proof}

\begin{theorem}\label{thm:actual-finite-chebotarev}
Let \(L/K\) be a finite Galois extension of number fields, and
fix \(g\in{\rm Gal}(L/K)\). Infinitely many primes \(\mathfrak q\)
of \(K\) of rational prime norm admit a residue map on \(O_L\)
whose arithmetic Frobenius is \(g\) in the sense of
Lemma~\ref{lem:finite-extension-residue-actions}. All these
primes can be required to avoid any specified finite set of
rational primes. The same assertion holds for geometric
Frobenius, defined as the inverse of arithmetic Frobenius.
\end{theorem}
\begin{proof}
Write \(G={\rm Gal}(L/K)\) and \(n\) for the order of \(g\).
Choose a finite normal extension \(L_1/\Q\) containing \(L\)
by adjoining the conjugates of a primitive generator from
Lemma~\ref{lem:number-field-lattice}. Choose an odd prime
\(p\equiv1\bmod n\) outside the exceptional set for \(L_1/\Q\)
in Lemma~\ref{lem:finite-extension-residue-actions}, using
Corollary~\ref{cor:ap}. In particular \(p\) has valuation one
at each prime of \(L_1\) over it. The polynomial
\[
               \frac{(1+X)^p-1}{X}
\]
is irreducible over \(L_1\). At one such valuation it is monic,
has all lower coefficients divisible by \(p\), and constant
term \(p\) of valuation one. In a monic factorization over
the field both factors would have integral coefficients, since
the least coefficient valuation is additive on multiplication
by reduction of primitive polynomials over the residue field.
Both reductions would then be powers of \(X\) of positive
degree, making the constant term divisible by the square of
the uniformizer. This proves irreducibility and hence
\([L(\zeta_p):L]=p-1\).

Put \(E'=L(\zeta_p)\) and \(J=(\Z/p\Z)^\times\). Restriction
embeds \({\rm Gal}(E'/K)\) in \(G\times J\).
The compositum is normal and separable over \(K\) and its
degree is \(|G|(p-1)\), so this inclusion is an isomorphism.
Choose \(a\in J\) of order \(n\) by cyclicity as in
Lemma~\ref{lem:split-level}. Put \(\Delta=\langle(g,a)\rangle\),
\(K_0=(E')^\Delta\). The primitive-element fixed-field argument
in Lemma~\ref{lem:cyclotomic-prime-norms} gives
\([E':K_0]=|\Delta|=n\) and identifies its Galois group with
\(\Delta\). The images of \(\zeta_p\) under \(\Delta\) are
distinct, so it has degree at least \(n\) over \(K_0\).
Thus \(E'=K_0(\zeta_p)\), and its cyclotomic Galois group
over \(K_0\) identifies with \(\langle a\rangle\) by the
second coordinate.

By Lemma~\ref{lem:cyclotomic-prime-norms} for \(K_0,p,a\),
infinitely many primes \(\mathfrak Q\) of \(K_0\) have prime
norm \(\ell\equiv a\bmod p\). Remove the specified primes and
the finitely many exceptions for all three Galois extensions
\(E'/K_0,E'/K,L/K\) in
Lemma~\ref{lem:finite-extension-residue-actions}.
The contraction \(\mathfrak q\) of \(\mathfrak Q\) to \(K\)
also has residue field \(\mathbf F_\ell\), since its residue
field injects into that of \(\mathfrak Q\) and contains
\(\mathbf F_\ell\). Take a residue map for \(E'/K_0\) at
\(\mathfrak Q\). Its Frobenius in \(\Delta\) has second
coordinate \(a\) and so equals \((g,a)\). As both residue
fields have the same size, the very same equation
\(h(\gamma x)=h(x)^\ell\) identifies this element as Frobenius
over \(K\) too. Restriction to \(L\) gives \(g\).
Each \(\ell\) supports finitely many \(\mathfrak Q\), hence the
result gives infinitely many distinct \(\mathfrak q\). Using
\(g^{-1}\) proves the statement for geometric Frobenius.
\end{proof}

\section{Determination of a characteristic polynomial at inertia}

\begin{definition}\label{def:frobenius-data}
Fix an algebraic closure \(\overline K\) of a number field \(K\),
write \(O_{\overline K}\) for its algebraic integers, and put
\(G_K={\rm Aut}(\overline K/K)\). Give \(G_K\) the topology
whose basic neighborhoods prescribe the restriction to a finite
normal subextension of \(\overline K/K\). A Frobenius datum is a
triple \((v,h,\gamma)\) where \(v\) is a prime of \(K\) of
rational prime norm \(q\), \(h:O_{\overline K}\to
\overline{\mathbf F}_q\) extends reduction at \(v\), and
\(\gamma\in G_K\) satisfies
\[
              h(\gamma^{-1}x)=h(x)^q\qquad(x\in O_{\overline K}).
\]
The exponent \(-1\) specifies geometric Frobenius.
\end{definition}

\begin{lemma}\label{lem:lift-residue-frobenius}
Every nonempty open \(V\subset G_K\) contains geometric
Frobenius elements at degree-one primes outside any prescribed
finite set of primes of \(K\). More precisely, infinitely many
such primes \(v\) admit a residue map \(h\) for which there is
at least one datum \((v,h,\gamma)\), and every such datum has
\(\gamma\in V\).
\end{lemma}
\begin{proof}
We first give the lifting argument at arbitrary finite
extensions of number fields. If \(K_2/K_1\) is finite, each
residue map on \(O_{K_1}\) at a prime \(\mathfrak p\) extends to
\(O_{K_2}\), with values in the same algebraic closure of the
residue field. Indeed localize \(O_{K_2}\) at
\(O_{K_1}\setminus\mathfrak p\). It is a nonzero finite
module over the local ring: an integer basis of \(O_{K_2}\)
generates it there by Lemma~\ref{lem:number-field-lattice}.
The maximal ideal of the base cannot generate the unit ideal
in this algebra. Otherwise for module generators \(z\) one
would have \(z=Bz\) with \(B\) a matrix over that maximal
ideal. The adjugate of \(I-B\) and its unit determinant would
give \(z=0\). A maximal ideal containing the extended ideal
therefore lies over \(\mathfrak p\). Its quotient is a finite
field extension of the base residue field, which embeds in the
chosen algebraic closure by adjoining roots successively.

If \(K_2/K_1\) is Galois with group \(H\), then \(H\) is
transitive on residue maps extending a fixed one, including
at ramified primes. To see this, take an integer basis
\(z_1,\ldots,z_m\) of \(O_{K_2}\) and independent
indeterminates \(T_1,\ldots,T_m\). The coefficients of
\[
          \prod_{\tau\in H}
             \left(X-\sum_{j=1}^m T_j\tau z_j\right)
\]
belong to \(O_{K_1}[T_1,\ldots,T_m]\). They are integral and
fixed by \(H\), and its fixed field is \(K_1\) by the
primitive-element argument in Lemma~\ref{lem:cyclotomic-prime-norms}.
For two such maps \(b,b'\) reduction therefore gives the same
polynomial. Substitute \(X=\sum_jT_j b'(z_j)\). The product
using \(b'\) vanishes. The product using \(b\) must then have
a zero factor in the polynomial ring over the residue field
closure, an integral domain. Comparing all coefficients of \(T_j\)
gives \(b'=b\circ\tau\) on the basis for some single \(\tau\in H\).

Now choose a finite normal \(L/K\) and \(g\in{\rm Gal}(L/K)\)
such that all extensions of \(g\) lie in \(V\), by the topology
in Definition~\ref{def:frobenius-data}. Theorem~\ref{thm:actual-finite-chebotarev}
gives infinitely many \(v\) of prime norm \(q\) with a residue
map \(b\) on \(O_L\) whose arithmetic Frobenius is \(g^{-1}\).
We may exclude the underlying rational primes of the given
finite set. Take an increasing sequence of finite normal fields
\(L=L_0\subseteq L_1\subseteq\cdots\) exhausting \(\overline K\).
Such a sequence exists since the polynomials over the countable
field \(K\) form a countable set and adjoining their splitting
fields successively exhausts the closure. Extend \(b\) along
the sequence by the first argument, giving \(h\).
At each level the two maps \(h\) and \(h^q\) agree on \(O_K\).
By transitivity choose an automorphism realizing \(h^q\) by
composition with \(h\) on that level. Extend it to \(G_K\)
by the algebraic embedding extension in Lemma~\ref{lem:place-fields}.
From these choices at successively higher levels extract an
infinite subsequence with constant restriction to \(L_0\), then
one constant on \(L_1\), and continue. This works since the
restriction sets are finite. The eventual restrictions on the
diagonal subsequence are compatible and define an automorphism
of the union, with inverse given by their inverses. Its inverse
\(\gamma\) gives the desired datum. Every such \(\gamma\) must
restrict to \(g\) by uniqueness of finite Frobenius with
residue map \(b\) in Lemma~\ref{lem:finite-extension-residue-actions}.
\end{proof}

\begin{definition}\label{def:frobenius-approximation}
A property of the data in Definition~\ref{def:frobenius-data}
holds for most Frobenius data if for each nonempty open \(V\)
in \(G_K\) there are a nonempty open \(U\subseteq V\) and a finite
set \(S\) of primes of \(K\) such that the property holds whenever
\(\gamma\in U\) and \(v\notin S\). Let \(\ell\) be a rational prime.
Give the algebraic closure \(K_\ell\) of \(\Q_\ell\) the absolute
value of Lemma~\ref{lem:place-fields}. Fix polynomials
\(P_v\in K_\ell[X]\) at the degree-one primes outside a finite
set, assigning arbitrary polynomials at those exceptions.
A continuous homomorphism
\(r:G_K\to{\rm GL}_d(K_\ell)\), with the entrywise topology,
has Frobenius polynomials \(P_v\) in the approximation sense
if, for every \(j\ge0\) and \(\eta>0\), for most Frobenius data
\[
   \left|[X^j]\det(X I-r(\gamma))-[X^j]P_v\right|_\ell\le\eta .
\]
Here \([X^j]\) denotes extraction of one coefficient.
\end{definition}

\begin{theorem}\label{thm:common-inertia-characteristic}
Suppose \(r,r'\) have the same Frobenius polynomials in the
approximation sense of Definition~\ref{def:frobenius-approximation}.
Then \(\det(XI-r(s))=\det(XI-r'(s))\) for every \(s\in G_K\).
In particular this holds for the marked inertia element of
Theorem~\ref{thm:marked-tame-one} when \(K=E\).
\end{theorem}
\begin{proof}
If the coefficients at degree \(j\) differ at \(s\), let their
difference have absolute value \(\delta>0\). Each coefficient
is a finite polynomial in the matrix entries, hence continuous.
There is a nonempty open neighborhood \(V\) of \(s\) where
the absolute difference exceeds \(\delta/2\).
Take \(\eta=\delta/4\) in
Definition~\ref{def:frobenius-approximation}. First apply the
condition for \(r\) within \(V\), then that for \(r'\) within
the resulting open subset. On the resulting nonempty open \(U\)
both conditions hold away from the union of their finite
exceptions. Lemma~\ref{lem:lift-residue-frobenius} supplies a
datum there. Subtraction of the two estimates at its common
\(P_v\) bounds the absolute difference by \(\eta\) by the
ultrametric inequality. This contradicts the choice of \(V\).
The argument applies coefficient by coefficient, including
the leading coefficients, so it also identifies the degrees.
\end{proof}

\begin{corollary}\label{cor:marked-label-from-local-lifts}
For a fixed actual Hecke eigenspace in
Proposition~\ref{prop:one-inertia-matrix-suffices}, fix its
chosen marked prime, and use \(s\) of
Theorem~\ref{thm:marked-tame-one} in \(G_E\).
Suppose every occurring primitive label \(c\) supplies a
continuous \(r_c:G_E\to {\rm GL}_{d_c}(K_2)\), \(d_c>0\),
with the same family \(P_v\) in the approximation sense of
Definition~\ref{def:frobenius-approximation}.
Suppose \(r_c(s)\) obeys the annihilating polynomial and the
determinant identity of Proposition~\ref{prop:one-inertia-matrix-suffices},
for that label and the same \(T,\eta,A\) there. Then the
primitive label is unique if one occurs.
\end{corollary}
\begin{proof}
Theorem~\ref{thm:common-inertia-characteristic} gives the same
characteristic polynomial on all the matrices \(r_c(s)\).
Apply Proposition~\ref{prop:one-inertia-matrix-suffices}
with \(L=K_2\). Its image of the finite cyclic residue group
can be taken in \(K_2\) by composing the chosen \(T\) into
\(\C^\times\) with the fixed field isomorphism of
Lemma~\ref{lem:place-fields}.
\end{proof}

\section{Reconstruction from traces}

\begin{definition}\label{def:trace-identities}
Fix an algebraically closed field \(k\) of characteristic zero, a
group \(G\), and a positive integer \(d\). A function \(t:G\to k\)
satisfies the trace identities of degree \(d\) if, for every finite
list \(g_1,\ldots,g_m\in G\) and \(B\in k[X_1,\ldots,X_m]\),
\[
 B(\operatorname{tr}r(g_1),\ldots,\operatorname{tr}r(g_m))=0
                   \quad\hbox{for every }r:G\to{\rm GL}_d(k)
\]
implies \(B(t(g_1),\ldots,t(g_m))=0\). Here and below such \(r\)
are group homomorphisms. Write \(k[G]\) for the algebra of finite
sums of symbols \([g]\) with \([g][h]=[gh]\), and \(\tau\) for
the linear extension of \(t\) to this algebra.
\end{definition}

\begin{lemma}\label{lem:trace-pairing-quotient}
For \(t\) of Definition~\ref{def:trace-identities} satisfying the
identities, let \(I=\{x\in k[G]:\tau(xy)=0\text{ for every }y\in k[G]\}\).
Then \(I\) is a two-sided ideal and \(A=k[G]/I\) has dimension
at most \(d^2\). The functional \(\tau\) descends to \(A\),
takes \(1\) to \(d\), and has nondegenerate symmetric pairing
\((x,y)\mapsto\tau(xy)\). Its value on each nilpotent is zero.
On each idempotent it takes a value in \(\{0,1,\ldots,d\}\).
\end{lemma}
\begin{proof}
Cyclicity \(\tau(xy)=\tau(yx)\) and \(\tau(1)=d\) follow from
the corresponding matrix trace identities by linear expansion.
Right multiplication preserves \(I\) by associativity and left
multiplication by cyclicity. Taking \(y=1\) gives \(I\subset
\ker\tau\). These facts and the definition give the claimed
pairing. For \(m=d^2+1\) and any \(x_i,y_j\in k[G]\),
\(\det[\tau(x_i y_j)]_{i,j=1}^m=0\). This is a polynomial in
finitely many values of \(t\). For any representation it vanishes
because the \(m\) matrices \(r(x_i)\), with \(r\) extended linearly,
are dependent in the \(d^2\)-dimensional matrix space. If the
images of the \(x_i\) were independent in \(A\), the \(m\) linear
functionals \(y\mapsto\tau(x_i y)\) would be independent by
nondegeneracy. Their joint map onto \(k^m\) is surjective,
since a proper subspace of \(k^m\) has a nonzero annihilating
linear form by basis extension. Choosing preimages of the
coordinate vectors would yield a determinant of one. This
proves the dimension bound.

For \(x\in A\) form the formal power series
\[
       D_x(Z)=\exp\!\left(-\sum_{m\ge1}\tau(x^m)Z^m/m\right).
\]
Each coefficient is a polynomial in finitely many traces of
powers, since the exponent has zero constant term. Its
coefficients past degree \(d\) vanish by the trace identities
applied to a lift of \(x\). Indeed for a \(d\)-by-\(d\) matrix
\(M\), choose an eigenvector over \(k\) and triangularize by
induction on the quotient dimension. The resulting diagonal
\(\lambda_j\) gives trace powers \(\sum_j\lambda_j^m\) and
\(\det(I-ZM)=\prod_j(1-\lambda_j Z)\). Formal differentiation
and the constant term one identify this determinant with the
displayed exponential for matrix trace, degree at most \(d\).
Uniqueness here follows by equating the successive coefficients
in the logarithmic derivative, since the field has characteristic
zero. Thus \(D_x\) is a polynomial with constant term one.

If \(x\) is nilpotent, \(S(Z)=\sum_{m\ge1}\tau(x^m)Z^{m-1}\)
is a polynomial and \(D_x'=-S D_x\). A nonzero \(S\) makes the
degree of the right side at least \(\deg D_x\), incompatible
with the derivative, including when the derivative is zero.
Therefore \(S=0\), in particular \(\tau(x)=0\).
If \(x^2=x\), put \(u=\tau(x)\). Then
\((1-Z)D_x'=-uD_x\). Starting with constant term one, the
coefficient of \(Z^{d+1}\) is
\((-1)^{d+1}u(u-1)\cdots(u-d)/(d+1)!\), by comparison of
successive coefficients. Its vanishing gives the assertion.
\end{proof}

\begin{lemma}\label{lem:simultaneous-matrix-image}
Let \(B\) be a unital \(k\)-algebra and \(S_1,\ldots,S_m\)
pairwise nonisomorphic, nonzero finite-dimensional simple
left \(B\)-modules. Simplicity means there is no proper nonzero
invariant subspace. The joint action maps \(B\) onto
\(\prod_i{\rm End}_k(S_i)\).
\end{lemma}
\begin{proof}
A nonzero homomorphism between two simple modules is an
isomorphism, by its kernel and image. Any endomorphism of
one of them has an eigenvalue over \(k\). Subtracting that
scalar leaves an endomorphism with kernel, hence zero.
Choose a basis \(v_{i1},\ldots,v_{im_i}\) of each \(S_i\).
The lists \((b v_{ij})_{i,j}\) form a \(B\)-submodule of
\(V=\bigoplus_i S_i^{m_i}\). If it were proper, choose a
proper submodule containing it of largest dimension. The
quotient is simple, and at least one coordinate summand maps
nontrivially to it. Identify this quotient with that \(S_i\).
The quotient map on \(V\) is then zero on the summands of
other types, and on the copies of \(S_i\) it is a list of
scalars, at least one nonzero, by the preceding argument.
Taking \(b=1\) gives a nontrivial dependence among the basis
vectors of \(S_i\), since the list for \(b=1\) maps to zero.
This is impossible. Thus all lists of images are realized,
which gives all tuples of endomorphisms.
\end{proof}

\begin{theorem}\label{thm:trace-reconstruction}
Every \(t\) satisfying Definition~\ref{def:trace-identities}
is the trace of a semisimple \(d\)-dimensional representation
of \(G\) over \(k\), unique up to change of basis among
semisimple representations. Here semisimple means a direct
sum of simple representations. If \(k\) is a topological field,
\(G\) a topological group, and \(t\) continuous, this
representation is continuous in the matrix entries.
\end{theorem}
\begin{proof}
Take \(A\) of Lemma~\ref{lem:trace-pairing-quotient}. Choose a
chain \(0=A_0\subset\cdots\subset A_h=A\) of left ideals
with simple successive quotients, by taking a proper submodule
of largest dimension repeatedly. Let \(S_i\) represent the
distinct isomorphism classes of its quotients. Their joint
action has kernel \(J\), a two-sided ideal. For all \(x\in J\)
one has \(x A_j\subset A_{j-1}\), so \(J^h A=0\) and hence
\(J^h=0\). For \(x\in J\), \(y\in A\), the product \(xy\)
is in \(J\) and nilpotent. Lemma~\ref{lem:trace-pairing-quotient}
gives \(\tau(xy)=0\) for every \(y\), hence \(x=0\). By
Lemma~\ref{lem:simultaneous-matrix-image} the action identifies
\(A\) with \(\prod_i{\rm End}_k(S_i)\).

On a matrix factor, cyclicity makes \(\tau\) zero on every
off-diagonal unit \(E_{ab}\), since \(E_{aa}E_{ab}=E_{ab}\)
and \(E_{ab}E_{aa}=0\) for \(a\ne b\). Its values on the
diagonal units agree by comparing \(E_{ab}E_{ba}\) and
\(E_{ba}E_{ab}\). Call this common value \(m_i\). The lemma
on the trace pairing makes \(m_i\) an integer between zero
and \(d\), and nondegeneracy excludes zero, since that entire
factor would then pair to zero. Consequently
\[
       \tau((a_i)_i)=\sum_i m_i\operatorname{tr}(a_i),
             \qquad \sum_i m_i\dim S_i=\tau(1)=d.
\]
The action of \([g]\) on \(\bigoplus_i S_i^{m_i}\) is invertible
with inverse the action of \([g^{-1}]\). Its trace is \(t(g)\).
Each summand is simple under \(G\) because the images of \([g]\)
span \(A\).

Conversely, in any semisimple representation with trace \(t\),
apply Lemma~\ref{lem:simultaneous-matrix-image} to the distinct
constituents. The span of the group has this full matrix
product as its action algebra. Its trace pairing is nondegenerate:
in each factor a nonzero entry is detected by multiplying by
the opposite matrix unit, and the positive multiplicity is
nonzero in \(k\). The kernel of the map from \(k[G]\) is
therefore exactly \(I\). After identifying this algebra with
\(A\), the possible simple modules are the defining modules
of its factors. Indeed the central factor idempotents split
any module, and in a simple module precisely one acts nontrivially.
For a matrix factor choose a nonzero vector fixed by some
\(E_{aa}\), possible since their sum is the identity. Its
images under \(E_{ba}\) span an invariant copy of the defining
module, independent by applying \(E_{ab}\) one at a time.
Simplicity makes it the whole module. The values of \(\tau\)
on its diagonal units now prescribe the same multiplicities
\(m_i\), proving uniqueness.

Finally choose a basis of \(A\) from the spanning images of
group elements \([g_j]\). The matrix
\((\tau([g_i][g_j]))_{i,j}\) is invertible by nondegeneracy.
The coordinates of \([g]\) in this basis are fixed linear
combinations of \(\tau([g][g_j])=t(gg_j)\) by inversion of
that matrix. These functions are continuous when \(t\) is,
so the action of \([g]\) on the stated direct sum is continuous.
\end{proof}

\section{Uniform approximation of eigenvalues and traces}

\begin{definition}\label{def:trace-approximation-relation}
Put \(k=K_\ell\) with the absolute value of
Lemma~\ref{lem:place-fields}. Let \(I_j\) be a finite set for
each \(j\ge1\), \(I=\coprod_j I_j\). An element of \(k\times k^I\)
has coordinates \(z_*,z_i\) for \(i\in I\). Give this set the
product topology and coordinatewise multiplication. We use a
closed submonoid \(N\) containing the all-ones tuple, weights
\(w_i\in k\), and real numbers \(\epsilon_j>0\) tending to zero
such that
\[
               \left|z_*-\sum_{i\in I_j} w_i z_i\right|_\ell
                    \le\epsilon_j\qquad(z\in N,\ j\ge1).
\]
\end{definition}

\begin{lemma}\label{lem:trace-contraction-eigenvalues}
Let \(Y\) be a set, let \(f:Y\to k\) be nonzero, and suppose
\(f\) is uniformly approximable by finite linear combinations
of functions from a family \(\mathcal F\). There exist data of
Definition~\ref{def:trace-approximation-relation} and choices
\(f_i\in\mathcal F\) for \(i\in I\) with the following property.
Whenever a \(k\)-linear operator \(B\) on \(k^Y\) preserves
every uniform upper bound for the absolute values of a function
and \(Bf=z_* f,\ B f_i=z_i f_i\) for all \(i\), one has \(z\in N\).
Zero functions can be omitted from the chosen lists.
\end{lemma}
\begin{proof}
Fix \(y_0\) with \(f(y_0)\ne0\) and put \(F=f/f(y_0)\).
For \(\epsilon_j=2^{-j}\), choose \(I_j\) and scalars \(a_i\)
so that \(|F(y)-\sum_{i\in I_j}a_i f_i(y)|_\ell\le\epsilon_j\)
at every \(y\). Multiplying a uniformly approximating sum for
\(f\) by \(f(y_0)^{-1}\) gives the required sums.
Let \(N_0\) consist of the tuples obtained from operators as
in the statement and set \(N=\overline{N_0}\).
The identity operator gives the all-ones tuple and composition
gives products since all the indicated functions are simultaneous
eigenvectors. The closure is still a monoid, since continuity of
multiplication gives product neighborhoods mapping into any given
neighborhood of a product, and each meets \(N_0\).
Apply \(B\) to the uniform error and evaluate at \(y_0\). This
gives the inequality in
Definition~\ref{def:trace-approximation-relation} on \(N_0\),
with \(w_i=a_i f_i(y_0)\). The inequality defines a closed
set, since it uses finitely many coordinates, so it holds on \(N\).
\end{proof}

\begin{lemma}\label{lem:trace-polynomial-transfer}
For the data of Definition~\ref{def:trace-approximation-relation},
let \(B\in k[X_1,\ldots,X_m]\), and \(z^{(1)},\ldots,z^{(m)}\in N\).
Write \(z_u=(z^{(1)}_u,\ldots,z^{(m)}_u)\) for \(u=*\) or \(u\in I\).
Then
\[
        \left|B(z_*)-\sum_{i\in I_j}w_i B(z_i)\right|_\ell
                        \le \epsilon_j\max_\alpha |B_\alpha|_\ell ,
\]
where \(B_\alpha\) are its coefficients and the maximum for the
zero polynomial is zero.
\end{lemma}
\begin{proof}
For each monomial \(X^\alpha\), the coordinate tuple of its
values is the product of the corresponding powers of the
\(z^{(m')}\), hence belongs to \(N\). This includes the constant
monomial by the identity requirement. Apply the defining bound
to each such tuple, multiply by the coefficient and sum.
The ultrametric triangle inequality gives the displayed estimate.
\end{proof}

\begin{theorem}\label{thm:trace-transport-from-primes}
Write \(y=(v,h,\gamma)\) for the Frobenius data of \(K\) in
Definition~\ref{def:frobenius-data} and use \emph{most} in the sense of
Definition~\ref{def:frobenius-approximation}. Suppose \(b(y)\) and
\(b_i(y)\) are \(k\)-valued functions and \(t_i:G_K\to k\) are
continuous. Assume the following conditions.
\begin{enumerate}
\item For most data one has \((b(y),(b_i(y))_{i\in I})\in N\)
and \(b(y)\in M\), where \(M\subseteq k\) is complete in the
metric of \(k\).
\item For each \(i\in I\) and each \(\delta>0\), for most data
\(\left|b_i(y)-t_i(\gamma)\right|_\ell\le\delta\).
\end{enumerate}
Then there is a continuous \(t:G_K\to k\) with values in \(M\)
such that \((t(g),(t_i(g))_{i\in I})\in N\) for every \(g\).
For each \(\delta>0\), for most data
\(|b(y)-t(\gamma)|_\ell\le\delta\). These last approximation
conditions determine \(t\) uniquely among continuous functions.
Every finite-dimensional \(\Q_\ell\)-subspace of \(k\) is an
eligible complete \(M\).
\end{theorem}
\begin{proof}
Finite intersections of properties holding for most data still
hold for most data, by passing to the supplied open subset
successively and taking the union of finite exclusions.
A bound \(|u(\gamma)|_\ell\le a\) for most data, for continuous
\(u:G_K\to k\) and \(a\ge0\), holds everywhere. Otherwise start
in the nonempty open set where the bound fails, refine once
using the property, and use Lemma~\ref{lem:lift-residue-frobenius}
to find a datum there outside the exceptions.

Set \(t_j'(g)=\sum_{i\in I_j}w_i t_i(g)\), a continuous function.
For fixed \(j\), finite intersection of the second hypothesis at
precision \(\epsilon_j/(1+\max_{i\in I_j}|w_i|_\ell)\) gives,
by the first hypothesis and the bound on \(N\),
\[
                    |b(y)-t_j'(\gamma)|_\ell\le\epsilon_j
\]
for most data. An empty list uses maximum zero. By subtraction
and the preceding continuity argument, for any \(j,l\) one has
\[
       |t_j'(g)-t_l'(g)|_\ell\le\max(\epsilon_j,\epsilon_l)
                              \qquad (g\in G_K).
\]
For fixed \(g,j\) take a neighborhood in which \(t_j'\) varies
by at most \(\epsilon_j\) from its value at \(g\). Refinement
using the estimates and
Lemma~\ref{lem:lift-residue-frobenius} gives some \(m_j=b(y)\in M\)
with \(|m_j-t_j'(g)|_\ell\le\epsilon_j\). These \(m_j\) form
a Cauchy sequence by the uniform inequality. Let \(t(g)\) be
their limit in \(M\). The same inequality gives
\(|t(g)-t_j'(g)|_\ell\le\epsilon_j\) at every \(g\), by letting
\(l\) tend to infinity. This also proves independence of choices.
The uniform limit is continuous. For a prescribed tolerance
use a single sufficiently large \(j\) and continuity of \(t_j'\).
The bound on \(b(y)-t(\gamma)\) for most data follows by the same
choice of \(j\). Two continuous candidates agree by applying
the first paragraph to their difference with arbitrarily small
tolerance.

For a neighborhood of \((t(g),(t_i(g))_i)\) in the product
topology it suffices to control the first and finitely many
other coordinates to a given positive precision. By continuity
start with a neighborhood of \(g\) controlling these finitely
many continuous functions. Inside it impose their prime
approximations and the first hypothesis, then take a datum
by Lemma~\ref{lem:lift-residue-frobenius}. Its tuple in \(N\)
approximates the chosen coordinates by the ultrametric inequality.
Closedness now proves membership in \(N\).

Finally in a finite \(\Q_\ell\)-basis of \(M\) the field norm
is at most a fixed multiple of the maximum coordinate norm.
On the coordinate unit sphere it is continuous and takes a
positive minimum by the compactness argument of
Lemma~\ref{lem:marked-unramified-tower} and linear independence.
Scaling gives the reverse inequality up to a fixed factor.
Completeness of the coordinates proves completeness of \(M\),
including the zero subspace.
\end{proof}

\section{Characteristic polynomials from approximating representations}

\begin{theorem}\label{thm:reconstruct-attached-from-forms}
Use \(k,N,I_j,w_i\) of Definition~\ref{def:trace-approximation-relation},
a number field \(K\), and \(d\ge1\). At its degree-one primes
let \(P_v\) and \(P_{i,v}\) for \(i\in I\) be monic of degree
\(d\) over \(k\). Suppose the following properties hold.
\begin{enumerate}
\item Each \(i\) supplies a continuous
\(r_i:G_K\to{\rm GL}_d(k)\) with Frobenius polynomials \(P_{i,v}\)
in the approximation sense of
Definition~\ref{def:frobenius-approximation}.
\item For each \(a=0,\ldots,d-1\), for most Frobenius data one has
\[
       \bigl(-[X^a]P_v,\,(-[X^a]P_{i,v})_{i\in I}\bigr)\in N .
\]
For most data \(-[X^{d-1}]P_v\) belongs to a fixed complete
subset of \(k\).
\end{enumerate}
Then a continuous semisimple \(r:G_K\to{\rm GL}_d(k)\) has
Frobenius polynomials \(P_v\) in the same sense. It is unique
up to isomorphism among such semisimple representations.
For every \(s\in G_K\) one has
\[
                  (\operatorname{tr}r(s),(\operatorname{tr}r_i(s))_i)
                              \in N .
\]
\end{theorem}
\begin{proof}
Apply Theorem~\ref{thm:trace-transport-from-primes} to
\(b=-[X^{d-1}]P_v\), \(b_i=-[X^{d-1}]P_{i,v}\), and
\(t_i=\operatorname{tr}r_i\). The coefficient of \(X^{d-1}\)
in a characteristic polynomial is minus the trace by expansion
of the determinant, so the theorem applies. It gives a
continuous \(t\) with the stated tuple condition. If a
polynomial \(B\) in the traces at \(s_1,\ldots,s_h\in G_K\)
vanishes for every degree-\(d\) representation, it vanishes
for all \(r_i\). Lemma~\ref{lem:trace-polynomial-transfer},
applied to the trace tuples at these \(s_j\), bounds
\(|B(t(s_1),\ldots,t(s_h))|_\ell\) by
\(\epsilon_j\max_\alpha |B_\alpha|_\ell\) for every
approximation index \(j\). Sending \(j\) to infinity gives zero.
Thus \(t\) is the trace of a continuous semisimple \(r\) by
Theorem~\ref{thm:trace-reconstruction}.

For \(1\le h\le d\) let \(B_h\) be the coefficient of \(Z^h\) in
\(\exp(-\sum_{n=1}^h Y_n Z^n/n)\), a fixed polynomial in
\(Y_1,\ldots,Y_h\). On traces of powers of a matrix this gives
its characteristic coefficient of degree \(d-h\), by the
triangularization in Lemma~\ref{lem:trace-pairing-quotient}.
Put \(C_h=\max_\alpha |(B_h)_\alpha|_\ell\) and write
\(c_h(r,s)=[X^{d-h}]\det(XI-r(s))\). At every \(s\in G_K\),
Lemma~\ref{lem:trace-polynomial-transfer} at \(s,s^2,\ldots,s^h\)
gives
\[
      \left|c_h(r,s)-\sum_{i\in I_j}w_i c_h(r_i,s)\right|_\ell
                                                   \le C_h\epsilon_j.
\]
Fix a positive tolerance \(\delta\). Take \(j\) so large that
\(\max(1,C_h)\epsilon_j\le\delta\). For most data, the sum over
the finite \(I_j\) on the left at \(s=\gamma\) differs from
\(\sum_i w_i[X^{d-h}]P_{i,v}\) by at most \(\delta\), since the
first assumption can be imposed simultaneously to tolerance
\(\delta/(1+\max_{i\in I_j}|w_i|_\ell)\).
The second assumption and the defining inequality for \(N\)
compare this last sum with \([X^{d-h}]P_v\) to within
\(\epsilon_j\). Their minus signs occur on both sides of the
inequality. The ultrametric inequality gives the required
estimate. Leading coefficients already coincide, and
coefficients beyond degree \(d\) vanish on both sides.
Finally any two such representations have the same
characteristic polynomials by
Theorem~\ref{thm:common-inertia-characteristic}. Their traces
agree and uniqueness follows by
Theorem~\ref{thm:trace-reconstruction}.
\end{proof}

\begin{proposition}\label{prop:local-identities-through-approximation}
In Theorem~\ref{thm:reconstruct-attached-from-forms} fix \(s\in G_K\).
Let \(Q_*,Q_i\in k[X]\) of degree at most \(m\) be given.
For each \(0\le n\le m\) suppose there is \(u_n\in k^\times\)
such that
\[
      (u_n[X^n]Q_*,(u_n[X^n]Q_i)_i)\in N .
\]
If \(Q_i(r_i(s))=0\) for every \(i\), then \(Q_*(r(s))=0\).
If, separately, \(\det r_i(s)=b_i\) for every \(i\) and
\((u b_*,(u b_i)_i)\in N\) for some \(u\ne0\), then
\(\det r(s)=b_*\).
\end{proposition}
\begin{proof}
For each \(g\in G_K\) and \(0\le n\le m\) take the trace tuple
at \(s^n g\) and the indicated coefficient tuple. The polynomial
in their coordinates given by \(\sum_n u_n^{-1}Y_n Z_n\) has
value \(\operatorname{tr}(Q_i(r_i(s))r_i(g))=0\) on each
\(i\)-coordinate. By Lemma~\ref{lem:trace-polynomial-transfer}
it has value zero also on the distinguished coordinate on
letting \(j\) tend to infinity. Thus
\(\operatorname{tr}(Q_*(r(s))r(g))=0\) for all \(g\).
The span of \(r(G_K)\) contains \(Q_*(r(s))\) and its trace
pairing is nondegenerate by the construction in
Theorem~\ref{thm:trace-reconstruction}. This yields the first
conclusion. For the determinant, use the polynomial
\((-1)^d B_d\) in the traces at \(s,\ldots,s^d\) from the
preceding theorem, subtracting \(u^{-1}\) times the supplied
last coordinate. The same transfer gives the second conclusion.
\end{proof}

\begin{corollary}\label{cor:inertia-root-containment-limits}
Keep the representation and trace tuple of
Theorem~\ref{thm:reconstruct-attached-from-forms}, and fix \(s\in G_K\).
Let \(S\subset k^\times\) be finite and \(a\ge1\). If for every
\(i\) all eigenvalues of \(r_i(s)\) belong to \(S\) and
\((\det r_i(s))^a=c\) for the same \(c\in k\), then all
eigenvalues of \(r(s)\) belong to \(S\) and
\((\det r(s))^a=c\).
\end{corollary}
\begin{proof}
Put \(H(X)=\prod_{\alpha\in S}(X-\alpha)^d\). Each monic
degree-\(d\) characteristic polynomial of \(r_i(s)\) divides \(H\).
For division of a fixed \(H\) by any monic degree-\(d\)
polynomial, long division makes the coefficients of the remainder
polynomials in the divisor coefficients: at each step one
subtracts the current leading coefficient times a monomial
times the divisor. This also covers \(\deg H<d\).
The divisor coefficients are themselves the polynomials \(B_h\)
of the preceding theorem in the traces of the powers of \(s\).
Each remainder coefficient is therefore a polynomial vanishing
on every approximating trace tuple. By
Lemma~\ref{lem:trace-polynomial-transfer} and \(\epsilon_j\to0\),
it vanishes on the reconstructed traces. Thus the characteristic
polynomial of \(r(s)\) divides \(H\), proving root containment.
The polynomial \(((-1)^d B_d)^a-c\) in the same traces vanishes
on every \(i\)-coordinate, hence on the distinguished one by
the identical argument. This proves the determinant identity.
\end{proof}

\section{Polynomial approximation and a flag over a local extension}

\begin{lemma}\label{lem:local-residue-indicators}
Let \(L\subset K_\ell\) be a finite extension of \(\Q_\ell\).
Write \(O=\{x\in L:|x|_\ell\le1\}\). Its maximal ideal has a
generator \(\pi\), and the integer valuation \(v\) normalized by
\(v(\pi)=1\) has \(v(\ell)\ge1\). Its residue field is finite.
Fix \(h\ge1\), a complete set \(A\subset O\) of representatives
modulo \(\pi^h\), and \(a\in A\). Put
\[
        H_a(X)=\prod_{\substack{c\in A\\ c\ne a}}\frac{X-c}{a-c}.
\]
For every \(x\in O\), \(H_a(x)\) is integral and congruent
modulo \(\pi\) to the indicator of \(a+\pi^h O\). The polynomials
\(H_a(X)^{\ell^j}\) converge uniformly on \(O\) to that indicator
as \(j\to\infty\).
\end{lemma}
\begin{proof}
For \(x\ne0\) in \(L\), all its \(\Q_\ell\)-conjugates in
\(K_\ell\) have the same absolute value by
Lemma~\ref{lem:marked-unramified-tower} and the extension of embeddings
in Lemma~\ref{lem:place-fields}. Their product with
multiplicity is the determinant of multiplication by \(x\) on \(L\):
in the power basis of \(\Q_\ell(x)\) its characteristic polynomial
is its minimal polynomial, repeated \([L:\Q_\ell(x)]\) times on
\(L\) by a basis over this smaller field. Hence
\(-\log_\ell|x|_\ell\in [L:\Q_\ell]^{-1}\Z\). These values form
a subgroup containing one, so its positive generator is attained
by some \(\pi\). Division by powers of \(\pi\) proves the assertions
about the maximal ideal and valuations. The coordinate-norm
equivalence and compact unit sphere argument in
Theorem~\ref{thm:trace-transport-from-primes} make \(O\) compact:
it is a closed subset of a bounded coordinate box, and scalar
multiples of \(\Z_\ell\) are compact by digit extraction.
The quotient modulo \(\pi\) is finite because its open cosets
cover a compact set. The same argument gives finiteness modulo
every \(\pi^h\).

If \(x\equiv a\bmod\pi^h\), each factor in the displayed product
is \(1+(x-a)/(a-c)\), congruent to one modulo \(\pi\).
If \(x\equiv b\bmod\pi^h\) for \(b\in A\setminus\{a\}\),
then \(v(x-c)=v(b-c)\) for \(c\ne b\) and \(v(x-b)\ge h\),
with \(v(0)=+\infty\). The sums
\(\sum_{c\in A,\ c\ne b}v(b-c)\) and
\(\sum_{c\in A,\ c\ne a}v(a-c)\) are equal. Each counts, for
\(1\le t<h\), the number of other residues in the same
class modulo \(\pi^t\), a number independent of its center by
translation. Subtracting the common term \(v(a-b)\) from both
sums yields \(v(H_a(x))\ge h-v(a-b)\ge1\).
If \(v(u)\ge1\), then \(v(u^{\ell^j})\ge\ell^j\).
If \(v(u-1)\ge1\), the binomial expansion of \(u^\ell\)
gives \(v(u^\ell-1)\ge v(u-1)+1\), since the middle binomial
coefficients are divisible by \(\ell\). Iteration proves the
uniform convergence on both sets of residues.
\end{proof}

\begin{proposition}\label{prop:local-compact-polynomial-density}
Let \(C\subseteq L^m\) be compact. For any continuous
\(\phi:C\to K_\ell\) and \(\delta>0\), there is a polynomial
\(P\in K_\ell[X_1,\ldots,X_m]\) with
\(\sup_{x\in C}|\phi(x)-P(x)|_\ell\le\delta\).
\end{proposition}
\begin{proof}
The empty set or a product with zero coordinates is immediate.
After rescaling each coordinate by a power of \(\pi\), compactness
puts \(C\) in \(O^m\) of Lemma~\ref{lem:local-residue-indicators}.
For each point of \(C\), continuity gives a product of residue
balls on whose intersection with \(C\) the oscillation is at most
\(\delta\). By finite subcover, there is a common \(h\ge1\) for
which the boxes modulo \(\pi^h\) refine the chosen local balls
on \(C\). Here two such balls that meet are nested, by the
ultrametric inequality. On each box meeting \(C\), take the
value of \(\phi\) at a representative point of \(C\).
The resulting finite linear combination of box indicators
approximates \(\phi\) within \(\delta\). By the preceding lemma
each box indicator is the uniform limit of the products of the
one-variable polynomials there, which are bounded by one.
Finite linear combinations also converge uniformly by the
ultrametric inequality. Choose one polynomial approximation
within \(\delta\) of the indicator combination and undo the
coordinate rescaling.
\end{proof}

\begin{lemma}\label{lem:generic-local-flag}
For every \(n\ge1\) there is \(Y\in{\rm GL}_n(K_\ell)\) with
entries in a single finite extension \(L'/L\) such that for every
\(g\in{\rm GL}_n(L)\) and \(1\le j\le n\) the upper left
\(j\)-by-\(j\) minor of \(gY\) is nonzero. If \(u\) is the first
column of \(Y\), then \(g\mapsto gu\) is injective on \(M_n(L)\).
\end{lemma}
\begin{proof}
There are elements \(\alpha\) in \(K_\ell\) of arbitrarily
large degree over \(L\). Indeed a generator of
\(\Q_\ell(\mu_{\ell^a-1})\) in
Lemma~\ref{lem:marked-unramified-tower} has degree at least
\(a/[L:\Q_\ell]\) over \(L\), by comparing the degrees of
the compositum over \(\Q_\ell\). Choose \(\alpha\) of degree
greater than \(2^{n^2}\), put \(L'=L(\alpha)\), and set
\[
              Y_{ab}=\alpha^{\,2^{(a-1)n+b-1}}
                      \qquad(1\le a,b\le n).
\]
In any \(j\)-row minor of the first \(j\) columns of \(Y\),
a determinant monomial has exponent whose binary digits give
exactly its chosen positions. Distinct terms have distinct
exponents, and the exponent sets for distinct row subsets are
disjoint. All exponents are smaller than \(2^{n^2}\), so
these minors, for fixed \(j\), are linearly independent over
\(L\) by the degree of \(\alpha\). In particular \(\det Y\ne0\).
The first \(j\) rows of an invertible \(g\) have rank \(j\):
elimination gives some nonzero \(j\)-column minor. Expanding
the leading minor of \(gY\) by multilinearity as in
Proposition~\ref{prop:standard-flag-minors} gives a linear
combination of the independent minors of \(Y\) with these
minors of \(g\) as coefficients. It cannot vanish. At \(j=1\)
the entries of \(u\) are independent over \(L\). Each row of a
matrix over \(L\) is therefore determined by its product with
\(u\), giving the last assertion.
\end{proof}

\section{Density of algebraic coefficients with prescribed weight gaps}

\begin{definition}\label{def:bounded-degree-free-weights}
For \(n\ge1\), call \(u_1,\ldots,u_n\in\Z\) \emph{free up to
degree \(B\)} if for each \(0\le h\le B\) the integers
\(\sum_i\alpha_i u_i\) are distinct as
\(\alpha\in\Z_{\ge0}^n\) ranges over \(\sum_i\alpha_i=h\).
For a dominant integral weight \(\lambda_1\ge\cdots\ge\lambda_n\)
put
\[
                 u_i(\lambda)=-(\lambda_i-\lambda_n+n-i).
\]
If \(m_j=\lambda_j-\lambda_{j+1}\), this is
\(-\sum_{i\le j<n}(m_j+1)\). A matrix coefficient of its
algebraic representation is \(g\mapsto\beta(W_\lambda(g)\xi)\)
for a vector \(\xi\) and a linear functional \(\beta\).
Over \(K_\ell\) we use the exterior tensor realization in
Theorem~\ref{thm:compact-highest-vector}, transported algebraically
through Lemma~\ref{lem:place-fields}. Its tensor basis, minors and
action formulas have integer coefficients, allowing determinant
inverses for negative weights.
\end{definition}

\begin{lemma}\label{lem:one-variable-free-weights}
Fix \(n\ge2\), \(B=(n+1)^2\), and put
\(y_i=(B+1)^{n-i}\) for \(i<n\), \(y_n=0\).
Take \(m_i=y_i-y_{i+1}-1\) for \(1\le i<n\).
For any \(t\ge0\), \(t\in\Z\), increase \(m_1\) by \(t\),
leaving the other \(m_i\) fixed. These nonnegative multiplicities
of diagram columns, with \(\lambda_n=0\), give weights
\(u_i(\lambda)\) free up to degree \(B\).
For \(n=1\) every integral weight has this property.
\end{lemma}
\begin{proof}
The negatives of the shifted weights are \(y_1+t,y_2,\ldots,y_n\)
by summation in Definition~\ref{def:bounded-degree-free-weights}.
Each entry except the last exceeds \(B\) times its successor.
For two distinct exponent lists \(\alpha,\alpha'\) of equal
degree \(h\), let \(j\) be their first different entry, ordered
so that \(\alpha_j>\alpha'_j\). This index is less than \(n\),
since the total degrees agree. The gain at \(j\) on pairing
with the positive weights is at least that entry, and the
loss at later indices is at most \(h\) times its successor
since the total of the negative differences in absolute value
is at most \(h\). Thus the pairings differ. In one variable
there is exactly one exponent list at each given degree.
\end{proof}

\begin{theorem}\label{thm:local-algebraic-weight-density}
Let \(X\) be a compact topological space, \(L/\Q_\ell\) finite
inside \(K_\ell\), \(T\) a finite set, and
\(z:X\to\prod_{\tau\in T}{\rm GL}_n(L)\) continuous.
For every continuous \(\phi:X\to K_\ell\) constant on the
fibres of \(z\), and every \(\delta>0\), there is a finite
linear combination of functions
\[
       x\longmapsto
       \beta\!\left(\Big(\bigotimes_{\tau\in T}
                  W_{\lambda^\tau}(z_\tau(x))\Big)\xi\right)
\]
within \(\delta\) of \(\phi\) uniformly on \(X\). At each
\(\tau\) the weights can all be taken from the one-parameter
family of Lemma~\ref{lem:one-variable-free-weights} when
\(n\ge2\), and from the nonnegative integers when \(n=1\).
In particular the shifted weights at every \(\tau\) are free
up to \((n+1)^2\).
\end{theorem}
\begin{proof}
For empty \(X\) there is no approximation condition. For
empty \(T\) the given function is constant and a scalar multiple
of the coefficient of the empty tensor. In the other cases
take \(Y,L',u\) of Lemma~\ref{lem:generic-local-flag}. Put
\[
       z'(x)=(z_\tau(x)u)_{\tau\in T}\in (L')^{n|T|},
                         \qquad C=z'(X).
\]
This has precisely the same fibres as \(z\) by that lemma.
Every continuous function on \(X\) constant on these fibres
defines a continuous function on \(C\). Indeed inverse images
of closed sets are compact in \(X\), and their images in the
Hausdorff space \(C\) are closed.

For \(n\ge2\) write \(D_j(A)\) for the leading minor of order
\(j\) and form the continuous function
\[
         b(x)=\prod_{\tau\in T}\prod_{1\le j<n}
                     D_j(z_\tau(x)Y)^{m_j}
\]
with the baseline multiplicities of
Lemma~\ref{lem:one-variable-free-weights}. For \(n=1\) put \(b=1\).
It never vanishes by Lemma~\ref{lem:generic-local-flag}.
Thus \(\phi/b\) is continuous and factors through \(C\).
The maximum \(M=\max_X |b|_\ell\) is finite and positive.
Proposition~\ref{prop:local-compact-polynomial-density} gives a
polynomial \(H\) in the entries of \(z'\) such that
\(|\phi-b\,H(z')|_\ell\le\delta\) everywhere, by approximating
\(\phi/b\) within \(\delta/M\).

Consider one monomial of \(H\). Its factors at \(\tau\) are
entries of \(z_\tau u\), hence minors of height one of \(z_\tau Y\)
in its first column. Let \(t_\tau\) be their total degree.
After multiplication by \(b\) the factor at \(\tau\) is
a product of flag minors with the baseline numbers of columns
and \(t_\tau\) extra columns of height one. For \(n\ge2\) this
is the diagram of precisely the weight of
Lemma~\ref{lem:one-variable-free-weights} at \(t_\tau\).
More explicitly, for the exterior tensor representation \(T_\lambda\)
and distinguished tensor \(v_\lambda\) of
Theorem~\ref{thm:compact-highest-vector}, this product evaluates
a tensor coordinate of
\[
           T_\lambda(z_\tau(x))T_\lambda(Y)v_\lambda .
\]
The vector \(T_\lambda(Y)v_\lambda\) lies in \(W_\lambda\)
by that theorem's equality of the group orbit spans. The
coordinate restricted to \(W_\lambda\) is a linear functional,
so this is a coefficient of that representation over \(K_\ell\).
For \(n=1\) the product is a constant times \(z_\tau(x)^{t_\tau}\),
a coefficient of the one-dimensional weight \(t_\tau\).
Products over \(\tau\) are coefficients of the tensor product
by tensoring the vectors and functionals. Summing the finitely
many monomials proves the assertion.
\end{proof}

\section{Local-field coordinates for definite automorphic forms}

\begin{definition}\label{def:definite-padic-form}
Let \(H\) be the standard definite unitary group of
Definition~\ref{def:definite-comparison-forms} in any positive
rank \(n\), and put \(\Gamma_H=H(F)\). Fix a rational prime
\(\ell\), set \(k=K_\ell\), and fix a field isomorphism
\(\iota:\C\to k\) of Lemma~\ref{lem:place-fields}. At the embedding
\(\tau\in\Phi\) use \(\sigma_\tau(a)=\overline{\tau(a)}\) on \(E\),
in agreement with the conjugation in
Definition~\ref{def:definite-comparison-forms}. Put
\(L_\tau=\Q_\ell(\iota\sigma_\tau(E))\subseteq k\).
Write \(z_\tau(h)\) for the matrix of the component of \(h\in H_f\)
at the place of \(E\) induced by \(\iota\sigma_\tau\), with
entries in \(L_\tau\). Put \(z=(z_\tau)_{\tau\in\Phi}\).
For a tuple of dominant weights \(\kappa\) write
\(W_\kappa=\bigotimes_{\tau\in\Phi}W_{\kappa^\tau}\) over \(k\),
with the action in Definition~\ref{def:bounded-degree-free-weights}.

An algebraic form of weight \(\kappa\) over \(k\) is a function
\(f:H_f\to W_\kappa\) fixed on the right by some open subgroup
of the integral stabilizer \(K^0\), with
\(f(\gamma g)=W_\kappa(z(\gamma))f(g)\) for \(\gamma\in\Gamma_H\).
Its local-field realization is
\[
                   F_f(g)= W_\kappa(z(g)^{-1}) f(g).
\]
Given an open subgroup \(J\subseteq K^0\) and a character
\(\theta:J\to k^\times\) with open kernel, forms of level
\((J,\theta)\) satisfy \(f(gj)=\theta(j)f(g)\).
\end{definition}

\begin{lemma}\label{lem:definite-field-coordinates}
Each \(z_\tau\) in Definition~\ref{def:definite-padic-form} is a
well-defined continuous homomorphism. All its entries lie in one
finite extension \(L/\Q_\ell\) for all \(g,\tau\), and
\(z_\tau(\gamma)\) is the entrywise application of
\(\iota\sigma_\tau\) on rational matrices. Transporting coefficients
by \(\iota\) identifies the complex definite finite-level forms
of those weights with the corresponding forms over \(k\).
For every such \(f\), \(F_f\) is continuous and left
\(\Gamma_H\)-invariant. Some finite extension \(L_f/\Q_\ell\)
contains every coordinate of \(f(g)\) and of \(F_f(g)\), for every
\(g\), in algebraic flag bases.
\end{lemma}
\begin{proof}
The field generated by the image of an integral generator of
\(E/\Q\) over \(\Q_\ell\) has finite degree and is \(L_\tau\).
The closure of \(\iota\sigma_\tau(E)\) there is \(L_\tau\):
it contains the \(\Q_\ell\)-span of the powers of the generator
by density of rational coefficients. The finite extension is
complete by Theorem~\ref{thm:trace-transport-from-primes}.
The preimage in \(O_E\) of its maximal ideal is a prime over
\(\ell\), since integral elements have integral images by
their monic equations. Its normalized valuation is a positive
multiple of the number-field valuation there: elements outside
the prime are units in the local field and denominators outside
the prime can be inverted, while a local uniformizer has positive
valuation by Proposition~\ref{prop:general-ideal-factorization}.
Consequently \(L_\tau\) is canonically the completion used for
the indicated component. Evaluation at one finite place is
continuous in the restricted product topology and multiplicative,
giving the asserted homomorphism. The compositum \(L\) of the
finitely many \(L_\tau\) has finite degree by repeated basis
multiplication. Rational matrices evaluate as claimed.

The exterior-power formulas in Theorem~\ref{thm:compact-highest-vector}
and the basis in Proposition~\ref{prop:standard-flag-minors}
are algebraic. Their matrices in rational bases have entries
polynomial in the matrix entries and determinant inverses, with
rational coefficients: the ambient tensor entries have this form,
and choosing independent rational orbit coefficients or the standard
minor basis only requires inversion of a fixed rational matrix.
Thus applying \(\iota\) carries the complex rational-point law
to the displayed left law. The finite-level law is transported
by the same scalar isomorphism, also in the reverse direction.
Right invariance by an open subgroup makes \(f\) locally constant.
The formula and continuity of \(z\) give continuity of \(F_f\);
the factors of \(W_\kappa(z(\gamma))\) cancel on the left.

For any open \(J_0\subseteq K^0\) fixing \(f\), the set
\(\Gamma_H\backslash H_f/J_0\) is finite. Indeed the trace
lattice argument of Lemma~\ref{lem:trace-finite} applies with
the standard rank \(n\) in place of the column rank. It counts
the isometry classes of the fixed-discriminant positive definite
trace lattices with their \(O_E\)-action, hence the classes at
\(K^0\). The subgroup \(J_0\) has finite index by the compactness
of \(K^0\) and its open cosets. Choose one \(c\) in each fine
class. Adjoin all coordinates of \(f(c)\) to \(L\). There are
finitely many and each belongs to \(k=\overline{\Q}_\ell\), so
the resulting \(L_f\) is finite. At \(g=\gamma c j\) the
transformation law applies a matrix over \(L\) to \(f(c)\).
Applying \(W_\kappa(z(g)^{-1})\) also uses entries in \(L\).
This proves the assertion on the whole group.
\end{proof}

\begin{proposition}\label{prop:finite-class-weight-extension}
Fix \(J,\theta,\kappa\) in Definition~\ref{def:definite-padic-form}.
For each representative \(c\) of
\(\Gamma_H\backslash H_f/J\) put \(D_c=J\cap c^{-1}\Gamma_H c\).
Evaluation \(f\mapsto (F_f(c))_c\) identifies the space of these
forms with the direct sum over \(c\) of the spaces
\[
    \{v\in W_\kappa: W_\kappa(z(t))v=\theta(t)v
                                      \text{ for every }t\in D_c\}.
\]
The projection onto the summand condition on a single copy of
\(W_\kappa\) is
\[
           v\longmapsto |D_c|^{-1}
                  \sum_{t\in D_c}\theta(t)W_\kappa(z(t)^{-1})v .
\]
\end{proposition}
\begin{proof}
The double classes are finite in number by
Lemma~\ref{lem:definite-field-coordinates}. The stabilizer \(D_c\)
is finite by the stabilizer part of Lemma~\ref{lem:trace-finite}:
after conjugation it preserves a definite trace lattice, giving
only finitely many choices for the images of its basis.
Writing \(ct=\gamma c\) in the two transformation laws gives the
claimed condition on \(F_f(c)\).
Conversely for one admissible \(v\) define
\[
       f(\gamma c j)=W_\kappa(z(\gamma c))\,\theta(j)v
\]
and set \(f=0\) on the other classes. If \(\gamma c j=\gamma'c j'\),
then \(t=j'j^{-1}=c^{-1}\gamma'^{-1}\gamma c\in D_c\);
inserting its law on \(v\) proves equality of the two proposed
values. This gives all prescribed values and fixes the form
uniquely. Its kernel of right invariance contains \(\ker\theta\),
which is open. Reindexing the finite sum by right multiplication
shows the indicated average satisfies the required law, and it
fixes every already admissible vector. This proves the projection
formula.
\end{proof}

\section{Changing weight and level on the definite group}

\begin{theorem}\label{thm:actual-definite-weight-approximation}
Use \(f,\kappa,J,\theta,z\) of Definition~\ref{def:definite-padic-form},
with \(f\) an algebraic form of weight \(\kappa\), initially of
any finite level. Fix \(T\subseteq\Phi\) and put \(U=\Phi\setminus T\).
Suppose for each \(d\in J\) with \(z_T(d)=1\) that
\[
                    \theta(d)=1,\qquad f(gd)=f(g)\quad(g\in H_f).
\]
For every \(\delta>0\), \(F_f\) is uniformly within \(\delta\) of a
finite sum of functions of the form
\[
        g\longmapsto (A\otimes I)W_{\lambda}(z(g)^{-1})f'(g).
\]
Here \(f'\) is an algebraic form of level \((J,\theta)\);
\(\lambda_U=\kappa_U\), and at every \(\tau\in T\) the weights
\(\lambda^\tau\) come from Theorem~\ref{thm:local-algebraic-weight-density}.
The linear map \(A:W_{\lambda_T}\to W_{\kappa_T}\) is allowed to
depend on the summand. Norms of vectors are the maxima of the
absolute values of their coordinates in the tensor flag bases.
\end{theorem}
\begin{proof}
Choose representatives \(c\) at level \(J\) as in
Proposition~\ref{prop:finite-class-weight-extension}. On the
compact group \(J\) put \(F_c(j)=F_f(cj)\) and
\[
  B_c(j)=\theta(j)^{-1}
                 \bigl(I\otimes W_{\kappa_U}(z_U(j))\bigr)F_c(j).
\]
This is continuous since \(\theta\) has open kernel. Expanding
\(F_c(j)\) leaves
\(\theta(j)^{-1}(W_{\kappa_T}(z_T(j)^{-1})\otimes I)
 W_\kappa(z(c)^{-1}) f(cj)\).
Thus \(B_c\) is constant on fibres of \(j\mapsto z_T(j^{-1})\):
two points in the same fibre differ by multiplication on the
right by an element \(d\) in the hypothesis.

Fix a positive provisional tolerance \(\varepsilon\).
Apply Theorem~\ref{thm:local-algebraic-weight-density} component
by component to \(B_c\) and \(j\mapsto z_T(j^{-1})\). Its
matrix entries are in a common finite local extension by
Lemma~\ref{lem:definite-field-coordinates}. The result is a
finite sum approximating \(B_c\) within \(\varepsilon\), with
terms
\[
          (A W_{\lambda_T}(z_T(j)^{-1})\otimes I)v,
           \qquad v\in W_{\lambda_T}\otimes W_{\kappa_U}.
\]
Indeed a coefficient \(\beta(W_{\lambda_T}(\cdot)\xi)\)
in an output coordinate \(e_a\otimes e_b\) has this form with
\(A=e_a\beta\) and \(v=\xi\otimes e_b\). Multiply back by
\(\theta(j)(I\otimes W_{\kappa_U}(z_U(j)^{-1}))\). Its
operator norm is bounded by a constant \(C\ge1\) on the compact
\(J\), independent of the approximant. We obtain for \(F_c\)
an error at most \(C\varepsilon\) and terms
\[
              G_v(j)=\theta(j)(A\otimes I)
                                  W_\lambda(z(j)^{-1})v .
\]

Use \(D_c\) of Proposition~\ref{prop:finite-class-weight-extension}.
For \(t\in D_c\) one has \(F_c(tj)=F_c(j)\), by the left
\(\Gamma_H\)-invariance of Lemma~\ref{lem:definite-field-coordinates}.
Averaging the approximating sum at \(tj\) therefore leaves a
uniform error at most \(|\,|D_c|\,|_\ell^{-1} C\varepsilon\).
For each term the average changes \(G_v(j)\) to \(G_{v^\circ}(j)\),
where \(v^\circ\) is precisely the finite stabilizer projection
of Proposition~\ref{prop:finite-class-weight-extension}. Indeed
\(z(tj)^{-1}=z(j)^{-1}z(t)^{-1}\) and \(\theta(tj)=\theta(t)\theta(j)\).
The form \(f'\) in that proposition supported on \(\Gamma_H cJ\)
with \(F_{f'}(c)=v^\circ\) satisfies
\(F_{f'}(cj)=\theta(j)W_\lambda(z(j)^{-1})v^\circ\).
This realizes each averaged term by the required automorphic form.

Choose \(\varepsilon\) with
\(C\varepsilon\max_c |\,|D_c|\,|_\ell^{-1}\le\delta\), possible
by finiteness of the classes. Sum the constructions on the
distinct classes, which have disjoint supports, and use left
\(\Gamma_H\)-invariance to extend each estimate from \(cJ\)
to all of \(\Gamma_H cJ\). The map \(A\otimes I\) commutes with
the action on the unchanged \(U\)-factors, as follows directly
from the tensor product. This proves all the assertions.
\end{proof}

\begin{corollary}\label{cor:definite-approximating-eigenforms}
Let \(\mathcal V\) be a set of pairwise distinct finite places
of \(F\) split in \(E\) at which \(J\) contains the full local
integral unitary factor, and suppose \(\theta\) is trivial on
each of those factors. In Theorem~\ref{thm:actual-definite-weight-approximation}
the \(f'\) can all be chosen as simultaneous eigenvectors for
the spherical Hecke algebras at the places in \(\mathcal V\),
transported to \(k\) by \(\iota\).
\end{corollary}
\begin{proof}
At each fixed weight the complex form space is finite dimensional
by evaluation and the finite stabilizers in
Proposition~\ref{prop:finite-class-weight-extension}, transported
by Lemma~\ref{lem:definite-field-coordinates}.
We give the unitary inner product explicitly. Pass to the open
subgroup \(J_1=\ker\theta\). Its index in \(J\) is finite, since
\(J\) is compact, so \(\theta\) has finite image and under
\(\iota^{-1}\) takes values on the complex unit circle. At
level \(J_1\) take the sum of the complex positive definite
coefficient pairings over the double classes, each divided by
the order of its stabilizer, as in
Theorem~\ref{thm:definite-comparison-rank}. These pairings are
unitary on the definite real group by
Theorem~\ref{thm:compact-highest-vector}. Equivalently this is
integration with Haar mass one on \(J_1\), counting a rational
orbit with the reciprocal stabilizer. Partitioning into
\(J_1\)-classes gives precisely the sum by the orbit count in
Theorem~\ref{thm:definite-comparison-rank}. On the full
space before level restriction each finite right translation
preserves that pairing by invariance of the quotient measure,
as in Lemma~\ref{lem:quotient-measure}. Hence convolution has
adjoint given by conjugated inversion. The spherical operators
at each place and their adjoints preserve the desired
\((J,\theta)\)-space: remove that local factor from any \(j\in J\)
by multiplication by its inverse component, reducing the remaining
level action to disjoint factors. At the local factor the
operators are bi-invariant. They commute there and across
distinct places by Lemma~\ref{lem:row-spherical-hecke}, whose
local diagonalization and transpose proof holds in any rank.
Thus they are commuting normal endomorphisms. If one is nonscalar,
diagonalize it and restrict all the others to its mutually
orthogonal proper eigenspaces. The adjoints commute there too.
Induction on dimension gives a simultaneous eigenbasis for the
whole family. For a normal endomorphism this diagonalization
follows by taking an eigenvector, which is also an eigenvector
of the adjoint with conjugate eigenvalue by equality of the two
squared norms for the shifted normal operator, and proceeding
on its orthogonal complement. In the scalar case any basis works.
Apply \(\iota\) to this decomposition of each approximating
form and distribute the finite sum. The spaces and the specified
weights and level are unchanged by taking eigencomponents.
\end{proof}

\section{An auxiliary local character for the definite eigenforms}

\begin{proposition}\label{prop:actual-wild-split-character}
Let \(\ell\) be a prime, \(n\ge1\),
\(J_1=1+\ell M_n(\Z_\ell)\subset {\rm GL}_n(\Q_\ell)\), and
\(J_2=1+\ell^2M_n(\Z_\ell)\). There is a finite-order character
\(\theta:J_1\to K_\ell^\times\), trivial on \(J_2\), with
the following cuspidality property. For every idempotent
\(e\in M_n(\Q_\ell)\) other than zero and one there is a
\(j\in J_1\cap(1+eM_n(\Q_\ell)(1-e))\) with \(\theta(j)\ne1\).
\end{proposition}
\begin{proof}
Choose \(A\in M_n(\mathbf F_\ell)\) to be multiplication by a
generator of \(\mathbf F_{\ell^n}^\times\) in a field basis.
The field and generator exist by Lemma~\ref{lem:torus-primary}.
There is no nonzero proper \(A\)-invariant subspace: stability
under \(A\) gives stability under all polynomials in \(A\),
and these include multiplication by every element of the field.
Choose \(\zeta\in K_\ell\) of order \(\ell\), possible by
Lemma~\ref{lem:place-fields}. Put
\[
              \theta(1+\ell X)=\zeta^{\,\operatorname{tr}(A\bar X)}
                         \qquad (X\in M_n(\Z_\ell)).
\]
The exponent is read modulo \(\ell\). Multiplication adds
\(X\bmod\ell\), since the cross term has a further factor of
\(\ell\). Hence this is a character, trivial on \(J_2\).

Let \(W=\operatorname{im} e\) have dimension \(a\), \(0<a<n\).
The subspace \(eM_n(1-e)\) consists exactly of maps into \(W\)
vanishing on \(W\). In fact \(e\) is the identity on its image,
and any such map \(B\) satisfies \(eB(1-e)=B\).
Choose a basis matrix \(P\in{\rm GL}_n(\Z_\ell)\) whose first
\(a\) columns span \(W\) over \(\Q_\ell\). To construct it, scale
a nonzero vector in \(W\) until it has integral entries and a
unit coordinate. Elementary elimination using that unit extends
it to an integral basis. In this basis \(W\) splits as the first
line and a subspace of the last \(n-1\) coordinates by subtracting
the first coordinate. Induct on \(a\). The lower left block of
\(\bar P^{-1}A\bar P\) is nonzero by the invariant-subspace
property. Choose an elementary matrix \(D\) in the opposite
upper right block for which
\(\operatorname{tr}(\bar P^{-1}A\bar P\bar D)\ne0\), using
the transposed position of one nonzero entry. Then
\(B=PDP^{-1}\) is integral, takes values in \(W\), and
annihilates \(W\). The element \(j=1+\ell B\) is invertible
with inverse \(1-\ell B\) and has nontrivial character by the
displayed formula. For \(n=1\) there is no idempotent to test.
\end{proof}

\begin{lemma}\label{lem:auxiliary-split-embeddings}
For the field \(E\), imaginary quadratic subfield \(K_0\) and
CM type \(\Phi\) of Proposition~\ref{prop:places}, and any finite
set of rational primes, there is an odd prime \(\ell\) outside
the set split completely in \(E\). For such primes outside
finitely many further exceptions, the maps
\(\iota\sigma_\tau:E\to K_\ell\) of
Definition~\ref{def:definite-padic-form} have images in
\(\Q_\ell\). They specify exactly one place from each conjugate
pair above \(\ell\), all above one of the two places of \(K_0\).
\end{lemma}
\begin{proof}
Apply Theorem~\ref{thm:actual-finite-chebotarev} to the identity
of a normal closure of \(E/\Q\). Omit \(2\), the specified
primes, and those dividing the discriminant or power-basis
index of an integral primitive element \(u\) of \(E\) as in
Lemma~\ref{lem:number-field-lattice}. At the resulting primes
all conjugates of \(u\) reduce to distinct elements of
\(\mathbf F_\ell\) by Lemma~\ref{lem:finite-extension-residue-actions}.
Each simple root of its monic minimal polynomial modulo
\(\ell\) lifts to \(\Z_\ell\) by the digit correction proved
in Lemma~\ref{lem:marked-unramified-tower}. Thus all embeddings
of \(E\) into \(K_\ell\) take values in \(\Q_\ell\), with
distinct prime kernels of reduction on \(O_E\), by the
power-basis index condition. Each completion is \(\Q_\ell\)
by Lemma~\ref{lem:definite-field-coordinates}. The chosen complex
embeddings all restrict to the same embedding of \(K_0\)
and select one from each conjugate pair by
Proposition~\ref{prop:places}. Applying conjugation and then
the same \(\iota\) to all of them preserves those two
properties. This proves the claims.
\end{proof}

\begin{theorem}\label{thm:definite-wild-approximants}
Use the arithmetic field of Lemma~\ref{lem:auxiliary-split-embeddings}
with \(|\Phi|\ge2\). Fix a finite family of complex algebraic
forms on the standard definite groups, of any positive finite
ranks and finite levels. Choose any finite set of rational primes
to exclude. There is one auxiliary prime \(\ell\) outside this
set, and two distinct places \(w_0,w_1\) above the same place
of \(K_0\) at \(\ell\), with the following property.
For each given form \(f\), transported to \(K_\ell\), its
realization \(F_f\) is uniformly approximable by finite sums of
the forms and coefficient maps of
Theorem~\ref{thm:actual-definite-weight-approximation}.
Each approximating form is of the cuspidal type of
Proposition~\ref{prop:actual-wild-split-character} at \(w_0\),
is spherical at \(w_1\), and is a simultaneous spherical
Hecke eigenform away from a single finite set of places.
Its shifted weights at the embeddings belonging to the two
places are free up to \((n+1)^2\) in rank \(n\).
\end{theorem}
\begin{proof}
An open subgroup fixing one form contains a product neighborhood
of the identity whose local groups are the full integral group
outside a finite set of rational primes, by the restricted
product topology. At its finitely many smaller factors such
a neighborhood contains compact open subgroups, for example
sufficiently small principal lattice congruence groups. Thus
there is a product level fixing the form. Omit the union of the
exceptional primes for all the given forms when choosing
\(\ell\) by Lemma~\ref{lem:auxiliary-split-embeddings}. Choose
\(w_0,w_1\) from the distinct places in that lemma and let
\(T\) consist of the two corresponding embeddings. For each
form choose \(J\) equal to its product level except at the split
factor for \(w_0\), where it is \(J_1\). The full split integral
group there and at \(w_1\) is \({\rm GL}_n(\Z_\ell)\): with the
standard definite form a matrix on one member of a pair determines
the inverse transpose on the conjugate member, as in
Lemma~\ref{lem:split-level}. Extend the local \(\theta\) of
Proposition~\ref{prop:actual-wild-split-character} to \(J\)
by trivial characters on the other factors. Its value is one
when \(z_T(d)=1\), because the component at \(w_0\) is then the
identity. Also all of \(J\) fixes \(f\) on the right.
Theorem~\ref{thm:actual-definite-weight-approximation} gives
the required approximation, with the prescribed freedom of
weights. The full factor at \(w_1\) fixes each approximant
since its level character is trivial there. Apply
Corollary~\ref{cor:definite-approximating-eigenforms} at every
split place outside the union of the finitely many exceptional
places, and also at \(w_1\). Zero summands can be discarded.
These choices impose the given wild character on every
eigenform and give the claimed common exceptional set.
\end{proof}

\section{Hecke polynomials of the definite approximants}

\begin{definition}\label{def:actual-spherical-polynomial}
Let \(v\) be a finite place of \(E\) distinct from its conjugate,
and put \(q_v=|O_E/v|\). Use the factor of the unitary group
identified with \({\rm GL}_n(E_v)\) by its component at \(v\).
Put \(K_v={\rm GL}_n(O_v)\) for the valuation ring \(O_v\)
of \(E_v\). For \(0\le h\le n\) let \(T_{v,h}\) denote the
sum of right translations by representatives of the cosets
\(uK_v\) in
\[
                    K_v\,{\rm diag}(\pi I_h,I_{n-h})\,K_v ,
\]
where \(\pi\) generates the maximal ideal of \(O_v\), and put
the identity at every other unitary factor. The cosets form a
finite set, since the double coset is compact and \(K_v\) open.
The sum on right-\(K_v\)-invariant forms is independent of
these choices. On a nonzero spherical eigenform \(f\) with
\(T_{v,h}f=t_{v,h}f\) define
\[
 P_v(X)=\sum_{h=0}^n(-1)^h q_v^{h(h-1)/2}\,t_{v,h}X^{n-h}.
\]
Here \(T_{v,0}=1\) on the invariant forms, so \(P_v\) is monic.
\end{definition}

\begin{proposition}\label{prop:away-hecke-trace-contractions}
Let \(f\) over \(k=K_\ell\) be a nonzero eigenform for the
operators of Definition~\ref{def:actual-spherical-polynomial}
at all degree-one places outside a finite set. For its
approximants from Corollary~\ref{cor:definite-approximating-eigenforms}
use any family with the same property outside a common finite
set. In their common output coordinates and at each such \(v\nmid\ell\),
the eigenvalues \(-[X^a]P_v\) and \(-[X^a]P_{i,v}\) arise from
one linear contraction on scalar functions, for each
\(0\le a<n\). All coefficients of the \(P_v\) for \(f\) lie
in a single finite extension of \(\Q_\ell\) inside \(k\).
\end{proposition}
\begin{proof}
For an approximating form \(f'\) the function in question is
\((A\otimes I)F_{f'}\) of
Theorem~\ref{thm:actual-definite-weight-approximation}, and the
target function is \(F_f\). View all as scalar functions on
\(H_f\times\mathcal B_0\), where \(\mathcal B_0\) is the finite
set indexing a basis of \(W_\kappa\), by evaluating one coordinate.
For \(a=n-h\), fix representatives \(u\) of the finite cosets
in Definition~\ref{def:actual-spherical-polynomial} and use
the operator
\[
         B_{v,a}F(g,j)=-(-1)^h q_v^{h(h-1)/2}\sum_{u}F(gu,j).
\]
The ultrametric inequality and \(|q_v|_\ell=1\) show that
it preserves every uniform pointwise upper bound in absolute
value. The components \(z_\tau(u)\) are the identity since
\(u\) is supported over a rational prime different from
\(\ell\). Thus \(z(gu)^{-1}=z(g)^{-1}\).
The common coefficient map and this weight action can be taken
outside the finite sum. Applying the eigenform relation inside
it gives the eigenvalue \(-[X^a]P_{i,v}\) on each approximating
function and \(-[X^a]P_v\) on the target. This also verifies
the signs when \(h=1\).

Choose \(g,j\) with \(f(g)_j\ne0\). The numerator
\((T_{v,h}f)(g)_j\) and denominator \(f(g)_j\) belong to
\(L_f\) from Lemma~\ref{lem:definite-field-coordinates}.
Their quotient is \(t_{v,h}\); all the rational normalizing
factors belong to \(L_f\). Finitely many exceptional
polynomials, if specified over \(k\), can be included by a
finite extension since each has finitely many algebraic
coefficients over \(\Q_\ell\).
\end{proof}

\begin{theorem}\label{thm:definite-seed-attachment}
Fix a nonzero spherical eigensystem of a definite form in
Theorem~\ref{thm:definite-wild-approximants}, with polynomials
\(P_v\) at the degree-one primes of \(E\) outside a finite
set. There is a countable list of definite eigenforms of the
same rank \(n\), with the cuspidal type at \(w_0\), spherical
at \(w_1\), and the free shifted weights specified there,
for which the following implication holds. If each form in
the list admits a continuous
\(r_i:G_E\to{\rm GL}_n(K_\ell)\) with its spherical polynomials
\(P_{i,v}\) in the approximation sense of
Definition~\ref{def:frobenius-approximation}, then the given
form admits a continuous semisimple representation with
polynomials \(P_v\) in that sense. Its trace at every \(s\in G_E\)
obeys the tuple condition of
Theorem~\ref{thm:reconstruct-attached-from-forms} for the same list.
\end{theorem}
\begin{proof}
Outside finitely many rational primes, degree-one places of
\(E\) split over \(F\). More precisely the place below one
such \(v\) has the same prime norm by its residue subfield.
Lemma~\ref{lem:finite-extension-residue-actions} for \(E/F\)
shows that above this place an inert prime would have squared
norm and a split one has a distinct conjugate. Delete the
ramified primes, the rational primes of the common exceptional
levels, and \(\ell\). Use the full spherical algebras at every
remaining split factor for the approximating eigenforms.
Choosing either member of a conjugate pair just chooses local
coordinates on that same split factor. Thus the simultaneous
eigenbasis of Corollary~\ref{cor:definite-approximating-eigenforms}
provides both orientations.

The nonzero scalar function given by the coordinates of
\(F_f\) in Proposition~\ref{prop:away-hecke-trace-contractions}
is uniformly approximable by finite linear combinations of the
corresponding coefficient functions of actual eigenforms by
Theorem~\ref{thm:definite-wild-approximants}. Take the finite
lists for errors tending to zero in
Lemma~\ref{lem:trace-contraction-eigenvalues}, discarding zero
functions. For each listed function fix a form realizing it.
The countable union of these lists is the required one.
For every remaining \(v\) and \(a<n\),
Proposition~\ref{prop:away-hecke-trace-contractions} and that
lemma put \((-[X^a]P_v,(-[X^a]P_{i,v})_i)\) in the resulting
closed monoid \(N\). The target trace coefficient belongs to
the fixed finite local extension in that proposition, which is
complete by Theorem~\ref{thm:trace-transport-from-primes}.
Conditions holding at all places outside a finite set hold
for most Frobenius data by
Definition~\ref{def:frobenius-approximation}, taking the same
open set in its definition. Hence the representation hypotheses
on the listed forms allow
Theorem~\ref{thm:reconstruct-attached-from-forms} to be applied
with \(K=E,d=n\). It gives the asserted representation and the
trace tuple at each \(s\).
\end{proof}

\begin{corollary}\label{cor:auxiliary-prime-label-test}
In Corollary~\ref{cor:marked-label-from-local-lifts} the field
\(K_2\) may everywhere be replaced by \(K_\ell\) for any prime
\(\ell\ne g_0\), with \(T\) transported by the corresponding
\(\iota\). For its uniqueness conclusion it suffices that,
for each occurring label \(c\), the representation supplied
there has, at the specified \(s\), all eigenvalues in
\(Z_c\cup\mu_A(K_\ell)\) and satisfies
\[
              (\det r_c(s))^A=\left(\prod_{z\in Z_c}z\right)^A .
\]
When the representation is obtained by
Theorem~\ref{thm:definite-seed-attachment}, these two conditions
follow if they hold on every approximating representation for
that label.
\end{corollary}
\begin{proof}
Theorem~\ref{thm:common-inertia-characteristic} applies in the
field \(K_\ell\), giving all occurring labels the same multiset
of eigenvalues, provided their target spherical polynomials agree
as in Corollary~\ref{cor:marked-label-from-local-lifts}. The
displayed identity puts the quotient of the two products in
\(\mu_A(K_\ell)\). Lemma~\ref{lem:large-root-label} now gives
uniqueness. For the last assertion the union of the stated
root sets is finite and its prescribed determinant power is
fixed for the particular label. Apply
Corollary~\ref{cor:inertia-root-containment-limits}.
\end{proof}

\section{Powers of a matrix along a diagonal family}

\begin{lemma}\label{lem:homogeneous-monic-divisor}
Let \(k\) be a field, \(R=k[y_1,\ldots,y_n]\), \(K\) its
fraction field, and \(Y\) a square \(M\)-by-\(M\) matrix of
homogeneous linear forms over \(R\). Its monic polynomial \(p\)
of least degree in \(K[X]\) with \(p(Y)=0\) belongs to \(R[X]\).
It divides \(\det(XI-Y)\) in \(R[X]\) and is homogeneous of
total degree \(\deg_X p\) in \(X,y_1,\ldots,y_n\).
\end{lemma}
\begin{proof}
The adjugate identity \((XI-Y)\operatorname{adj}(XI-Y)
=\det(XI-Y)I\) yields \(\det(XI-Y)|_{X=Y}=0\):
write the adjugate in powers of \(X\) and evaluate each
coefficient with the powers of \(Y\) on its left. The two
terms on the left then cancel. Division with remainder by the
least-degree monic annihilator shows that it divides every
annihilator in \(K[X]\).

We give the descent of monic factors. In a unique factorization
domain \(A\), remove from each list of coefficients the common
factor with the minimum exponent at each prime. The resulting
polynomial is primitive. Products of primitive polynomials are
primitive since reduction modulo each prime gives two nonzero
polynomials over a domain. Every polynomial over the fraction
field is a fractional scalar times a primitive polynomial over
\(A\), uniquely up to a unit of \(A\), by those minimum
exponents. The one-variable polynomial ring over a field has
unique factorization: induction on degree gives existence,
and division with remainder gives a Bezout identity with each
irreducible coprime to a given polynomial, so irreducibles
are prime. Apply this factorization over \({\rm Frac}(A)\)
to the primitive part of a polynomial in \(A[X]\), clearing
each factor to its primitive representative. The product is
primitive, so the remaining scalar for the primitive part is a
unit. Uniqueness over the fraction field, primitive normalization
and uniqueness on constant factors give unique factorization
in \(A[X]\). Starting with \(k\) proves it for \(R\).
Apply primitive normalization to both monic factors
of \(\det(XI-Y)\) in \(K[X]\). Their primitive representatives
have product equal to the determinant up to a unit. The product
of their leading coefficients is thus a unit, so each can
be made monic over \(R\). They are then the original monic factors.

The determinant is homogeneous of total degree \(M\).
In a product of nonzero polynomials over a domain the lowest
and highest total degrees both add, by multiplying the respective
nonzero homogeneous parts. A homogeneous product thus has only
homogeneous factors. The term \(X^{\deg_X p}\) determines the
total degree of \(p\).
\end{proof}

\begin{definition}\label{def:power-minor}
For a square matrix \(A\) of size \(M\) and an integer \(m\ge1\),
form the \(M^2\)-by-\(m\) matrix whose columns list the entries
of \(I,A,\ldots,A^{m-1}\) in a fixed common order. Its
\(m\)-row determinants are called the power minors of order
\(m\). If \(m>M^2\) the set of such minors is empty.
Over a field their simultaneous vanishing means that the powers
are linearly dependent: elimination finds \(m\) pivot rows
exactly when the columns are independent.
\end{definition}

\begin{theorem}\label{thm:diagonal-block-polynomial-test}
Fix a field \(k\), integers \(M\ge n\ge1\), and a surjection
\(c:\{1,\ldots,M\}\to\{1,\ldots,n\}\). Write
\(D(y)_{aa}=y_{c(a)}\) with other entries zero. Let \(x\in
{\rm GL}_M(k)\) and \(1\le m\le n+1\). Suppose \(e_1,\ldots,e_n\)
are free in the sense of Definition~\ref{def:bounded-degree-free-weights}
through degree \(m(m-1)/2\). If all power minors of order \(m\)
of \(D(T^{e_1},\ldots,T^{e_n})x\) vanish in \(k[T,T^{-1}]\),
then \(m=n+1\). Moreover, for each \(i\) there is a scalar
\(\lambda_i\in k\) with
\[
       (x^{-1})_{ab}=\lambda_i\delta_{ab}
               \qquad(c(a)=c(b)=i).
\]
\end{theorem}
\begin{proof}
Set \(Y=D(y)x\) with independent indeterminates \(y_i\).
Every entry of \(Y^j\) is homogeneous of degree \(j\), so each
power minor is homogeneous of degree
\(h=0+1+\cdots+(m-1)\). Distinct monomials of that degree
have distinct Laurent exponents under \(y_i\mapsto T^{e_i}\).
Their coefficients cannot cancel. Hence the minors already
vanish over \(k[y_1,\ldots,y_n]\). Definition~\ref{def:power-minor}
implies that the degree \(r\) of the least monic annihilator
over \(k(y_1,\ldots,y_n)\) is less than \(m\).

Write this annihilator as
\(p(X)=X^r+a_1 X^{r-1}+\cdots+a_r\), with \(a_0=1\).
By Lemma~\ref{lem:homogeneous-monic-divisor}, \(a_j\) is
homogeneous of degree \(j\). Its constant term \(a_r\)
divides \((-1)^M\det(x)\prod_a y_{c(a)}\), hence is a
nonzero scalar times a monomial by unique factorization.
If \(y_i=0\), a row with \(c(a)=i\) is zero in every positive
power of \(Y\). The \((a,a)\)-entry of \(p(Y)=0\) at this
specialization shows \(y_i\mid a_r\). Thus
\(n\le r<m\le n+1\), giving \(r=n\) and
\(a_n=u\,y_1\cdots y_n\) with \(u\in k^\times\).

Multiply \(p(Y)=0\) on the right by \(x^{-1}\). For row
\(a\) and column \(b\) of the same block \(i\), each entry
\((Y^j x^{-1})_{ab}\) with \(j\ge2\) is divisible by
\(y_i^2\), since \(Y^jx^{-1}\) starts and ends with \(D(y)\).
The term with \(j=1\) equals \(a_{n-1}D(y)\) after multiplying
by its polynomial coefficient. Extracting the coefficient of
\(y_1\cdots y_n\) gives
\[
  u(x^{-1})_{ab}+
     \delta_{ab}\,[\,\textstyle\prod_{h\ne i} y_h\,]a_{n-1}=0 .
\]
Here brackets denote coefficient extraction, including the
constant coefficient for an empty product. This proves the
assertion for every block.
\end{proof}

\section{The rank determined by spherical annihilators}

\begin{definition}\label{def:local-polynomial-cocharacter}
For a number field \(E\), let \(G_E\) be the group of
Definition~\ref{def:frobenius-data}. Let \(k_0\subset k\) be
Hausdorff topological fields with continuous inclusion and
\(r:G_E\to{\rm GL}_M(k_0)\) continuous. Fix \(Q\in{\rm GL}_M(k)\)
and write \(s(g)=Qr(g)Q^{-1}\). Integers \(\kappa_a\), \(1\le a\le M\),
give a diagonal family in the local polynomial closure in this
basis if the following holds. For every polynomial \(F\) in the
\(M^2\) matrix entries over \(k\), vanishing of \(F(s(g))\)
throughout some neighborhood of the identity implies the Laurent
polynomial identity
\[
        F(\operatorname{diag}(T^{\kappa_a})_{a=1}^M)=0.
\]
A representation has diagonal matrices in its linear span if
in some basis that span contains all diagonal matrices over its
coefficient field.
\end{definition}

\begin{lemma}\label{lem:distinct-constituents-diagonal-span}
For finitely many finite-dimensional continuous representations
of a group \(G\) over an algebraically closed topological field
\(k_0\), form their direct sum \(W\). In a chain of invariant
subspaces from \(0\) to \(W\) with simple successive quotients,
keep one copy of each isomorphism class of quotient. Their
direct sum \(V\) is continuous, and its image has diagonal
matrices in its linear span. For each \(g\in G\), every
polynomial annihilating the action of \(g\) on all the original
representations annihilates its action on \(V\).
\end{lemma}
\begin{proof}
Repeatedly taking a largest-dimensional proper subrepresentation
gives the chain. Bases adapted to the chain give block triangular
matrices with continuous entries on every diagonal block. Their
values there are the quotient actions. Polynomials act on these
blocks by the same polynomials of the quotient actions, proving
annihilation. Apply Lemma~\ref{lem:simultaneous-matrix-image} to
the action of the group algebra on one representative of each
simple type. The resulting image is the full product of the
endomorphism algebras. In concatenated bases of the simple
modules its span thus contains every diagonal matrix, including
each separate coordinate projection.
\end{proof}

\begin{theorem}\label{thm:spherical-annihilator-rank}
Let \(M,n\ge1\), and let \(E,k_0,k,r,Q,\kappa\) be as in
Definition~\ref{def:local-polynomial-cocharacter}. Suppose that
\(r\) has diagonal matrices in its linear span and that the
indicated diagonal family belongs to the local polynomial
closure. Suppose that among the \(\kappa_a\) occur \(n\)
weights \(e_i\) free through degree \(n(n+1)/2\) in
Definition~\ref{def:bounded-degree-free-weights}. At the
degree-one primes \(v\) fix polynomials \(P_v\in k_0[X]\).
Suppose that for most Frobenius data
\((v,h,\gamma)\), in the sense of
Definition~\ref{def:frobenius-approximation}, \(P_v\) is
monic of degree \(n\) and annihilates \(r(\gamma)\).
Then \(M=n\) and
\[
          P_v(X)=\det(XI-r(\gamma))
               \quad\text{for most Frobenius data.}
\]
\end{theorem}
\begin{proof}
By Definition~\ref{def:power-minor}, all power minors of order
\(n+1\) of \(s(\gamma)\) vanish on these data. Their vanishing
is a closed condition on \(G_E\). Every nonempty open set
contains a Frobenius element satisfying it: first take the
smaller open set in Definition~\ref{def:frobenius-approximation},
then apply Lemma~\ref{lem:lift-residue-frobenius} with the
indicated finite exclusions. Thus the minors vanish on \(G_E\).
For any \(g\in G_E\), substitute \(As(g)\) for the variable
matrix in each minor. At \(A=s(h)\) the resulting polynomial
vanishes for every \(h\). Definition~\ref{def:local-polynomial-cocharacter}
puts \(A=\operatorname{diag}(T^{\kappa_a})\) in this identity.

Take \(g=1\). The \(n\) prescribed weights are distinct by
freeness in degree one. If there were a further distinct
weight, use the \(n+1\) diagonal positions with distinct weights
as rows of the power minor. In the field \(k(T)\) the resulting
matrix is invertible: a polynomial of degree at most \(n\)
vanishing at \(n+1\) distinct powers of \(T\) is zero by
division by its linear factors. Consequently
\(\kappa_a=e_{c(a)}\) for a surjection
\(c:\{1,\ldots,M\}\to\{1,\ldots,n\}\).
Theorem~\ref{thm:diagonal-block-polynomial-test} applied at
every \(g\) now says that the compression of \(s(g)^{-1}\) to
each weight block is scalar. Inversion permutes the group
elements. Thus the compression of every matrix in their
\(k\)-linear span is scalar on that block.

In a basis giving the diagonal-span hypothesis, the \(M\)
coordinate projections have rank one and sum to \(I\). After
extension to \(k\) and change to the \(s\)-basis they still
belong to the span and retain those properties. On a block
of dimension at least two, the compression of any rank-one
matrix has rank at most one. Being scalar, each of these
compressions would therefore vanish, whereas their sum is
the block identity. All blocks have dimension one, giving \(M=n\).

Finally the set \(B\) where all order-\(n\) power minors of
\(s(g)\) vanish is closed. If it contained a nonempty open
set, choose \(g\) in that open set. For all \(h\) near \(1\),
\(hg\) belongs to it, and the same polynomial substitution
now puts \(\operatorname{diag}(T^{\kappa_a})s(g)\) in all the
order-\(n\) zero sets. Theorem~\ref{thm:diagonal-block-polynomial-test}
excludes \(m=n\). Every nonempty open set therefore contains a
nonempty open subset disjoint from \(B\). On the latter the
first \(n\) powers are independent. Applying the hypothesis
on most Frobenius data within this open subset gives a monic
annihilator \(P_v\) of degree \(n\). Its difference from the
monic characteristic polynomial, also an annihilator by the
adjugate argument of Lemma~\ref{lem:homogeneous-monic-divisor},
has degree less than \(n\) if nonzero. Independence forces
this difference to vanish. The choices of open subsets give
the asserted quantifier for most data.
\end{proof}

\begin{corollary}\label{cor:rank-test-for-definite-approximants}
In Theorem~\ref{thm:definite-seed-attachment}, it suffices for
each listed eigenform to produce a continuous representation
over \(K_\ell\) of some positive finite dimension, with
diagonal matrices in its linear span, having its spherical
polynomials as annihilators for most Frobenius data and a
diagonal family as in Theorem~\ref{thm:spherical-annihilator-rank}
over an extension of \(K_\ell\) containing among its weights
the \(n\) free shifted weights of that eigenform.
\end{corollary}
\begin{proof}
The bound \((n+1)^2\) in Lemma~\ref{lem:one-variable-free-weights}
dominates \(n(n+1)/2\). The polynomial of
Definition~\ref{def:actual-spherical-polynomial} is monic of
degree \(n\). Theorem~\ref{thm:spherical-annihilator-rank}
therefore gives rank \(n\) and exact characteristic polynomials
on most data. Exact equality gives every coefficient tolerance
in Definition~\ref{def:frobenius-approximation}, as required
in Theorem~\ref{thm:definite-seed-attachment}.
\end{proof}

\section{Fixed elements and cyclotomic periods}

\begin{definition}\label{def:completed-local-action}
Let \(p\) be a prime, \(K_p\) the algebraic closure of
\(\Q_p\) of Lemma~\ref{lem:place-fields}, and \(C_p\) its
metric completion. Put \(G_p={\rm Aut}(K_p/\Q_p)\) with the
subgroups fixing finite normal extensions as a neighborhood
basis at the identity. The absolute value is invariant under
these automorphisms by Lemma~\ref{lem:marked-unramified-tower}.
Each acts on \(C_p\) by applying it to a Cauchy sequence, with
the same norm and with its inverse supplied by the inverse
automorphism. The completion is a field since the norms of
a Cauchy sequence with nonzero limit eventually have fixed
positive value, making the inverse sequence Cauchy. Choose
primitive roots \(\zeta_j\) of order \(p^j\), \(j\ge1\),
with \(\zeta_{j+1}^p=\zeta_j\). Such a choice follows by taking
successive \(p\)-th roots in \(K_p\). The cyclotomic character
\(\chi_p:G_p\to\Z_p^\times\) is determined by
\(g(\zeta_j)=\zeta_j^{\,\chi_p(g)\bmod p^j}\). The compatible
exponents on these cyclic groups multiply under composition.
\end{definition}

\begin{theorem}\label{thm:completed-fixed-field}
For every subgroup \(J\subseteq G_p\) one has
\(C_p^J=\overline{K_p^J}\) inside \(C_p\). More precisely
if \(x\in C_p^J\) and \(z\in K_p\), there is a \(b\in K_p^J\)
with \(|x-b|_p\le16|x-z|_p\).
\end{theorem}
\begin{proof}
The \(J\)-orbit of \(z\) is a finite set, contained in the
roots of its minimal polynomial over \(\Q_p\). Start with this
set \(S\) and the selected \(z\in S\). Suppose at a step
that \(S\subset K_p\) is nonempty and \(J\)-invariant with
\(|S|\le2^k\), \(k\ge1\), and a point \(z\in S\) has been
selected. Write \(\delta=|x-z|_p\) and
\(U=\{t\in S:|x-t|_p=\delta\}\). It is nonempty and
\(J\)-invariant by Definition~\ref{def:completed-local-action}.
Set \(m=2^{k-1}\). The polynomial
\[
                  f(X)=\sum_{t\in U}(X-t)^m
\]
is fixed by \(J\), has degree \(m\), and has leading coefficient
\(|U|\) as an integer. Also \(|f(x)|_p\le\delta^m\).
Factoring \(f\) over \(K_p\), one root \(z'\) therefore satisfies
\[
 |x-z'|_p\le \delta\,\big|\,|U|\,\big|_p^{-1/m}
                 \le \delta\,2^{k/2^{k-1}},
\]
since for any positive integer \(b\) one has
\(|b|_p^{-1}\le b\). This includes \(f(x)=0\). The distinct
roots form a nonempty \(J\)-invariant set of size at most
\(2^{k-1}\). Continue from that set and \(z'\). After \(k\)
steps the set has one element \(b\), necessarily fixed by \(J\).
The factors accumulated have product at most
\(2^{\sum_{i\ge1}i/2^{i-1}}=16\), since summing the geometric
tails \( \sum_{i\ge j}2^{1-i}\) over \(j\ge1\) gives \(4\).
This proves the bound, including the case of an initial
singleton. Density of \(K_p\) gives one inclusion for fixed
fields. Continuity of each \(g\in J\) gives the other.
\end{proof}

\begin{lemma}\label{lem:cyclotomic-trace-projections}
Put \(L_j=\Q_p(\zeta_j)\), \(L_\infty=\bigcup_{j\ge1}L_j\)
and let \(\widehat L_\infty\) be its closure in \(C_p\).
Then \(K_p^{\ker\chi_p}=L_\infty\). There are continuous
\(\Q_p\)-linear maps \(R_j:\widehat L_\infty\to L_j\)
commuting with \(G_p\), with
\[
 |R_j(x)|_p\le p^{1/(p-1)}|x|_p,\qquad
                  \lim_{j\to\infty}R_j(x)=x .
\]
\end{lemma}
\begin{proof}
The kernel fixes precisely the tower among the algebraic
elements. Indeed an element outside \(L_\infty\) has a
separable minimal polynomial of degree greater than one over
\(L_\infty\). Sending it to another root extends to an
automorphism over \(L_\infty\) by the extension argument of
Lemma~\ref{lem:place-fields}, and that automorphism lies in
the kernel.

Let \(d_s=p^{s-1}(p-1)\) and \(\pi_s=\zeta_s-1\). Every
primitive \(p^s\)-th root has the form \(\zeta_s^a\) with \(a\)
prime to \(p\). Factoring \(X^a-1\) and then using an inverse
exponent modulo \(p^s\) shows that all their differences from
one have the same absolute value. The product of these
differences, up to sign, is \(p\), by evaluating
\((X^{p^s}-1)/(X^{p^{s-1}}-1)\) at \(1\) using the finite
geometric sum. Hence \(|\pi_s|_p=p^{-1/d_s}\).
The \(d_s\) terms \(c_i\pi_s^i\), \(0\le i<d_s\), have pairwise
different valuations when nonzero and \(c_i\in\Q_p\). In
particular the powers are independent, and the cyclotomic
equation gives the matching upper degree bound. They are a
basis of \(L_s\). The norm of a sum in this basis is the largest
term norm, which also proves that the valuation group is
\(d_s^{-1}\Z\) in the normalization \(v(p)=1\).
All maps \(\zeta_s\mapsto\zeta_s^a\) for \(a\) prime to \(p\)
are therefore automorphisms: the \(d_s\) distinct primitive
roots are the entire set of conjugates and all lie in \(L_s\).

For \(s\ge j\) put \(q=p^{s-j}\). The maps of \(L_s\) fixing
\(L_j\) act on \(\zeta_s\) by exponents \(1+up^j\),
\(0\le u<q\). On \(L_s\) define \(R_j\) as the average of these
\(q\) maps. This agrees with the average on every intermediate
layer since restriction has equal-sized fibers. The binomial
formula and the geometric sum on the \(q\)-th roots give
\(R_j(\pi_s^i)=(-1)^i\) for \(0\le i<q\). The powers \(\pi_s^i\)
in this range are a basis over \(L_j\), since their valuations
are distinct modulo \(d_j^{-1}\Z\) and \([L_s:L_j]=q\).
Thus for \(z=\sum_{i<q}c_i\pi_s^i\) the norm bound is
\[
      |R_j(z)|_p\le\max_i |c_i|_p
       \le p^{(q-1)/d_s}|z|_p
       \le p^{1/d_j}|z|_p .
\]
All values lie in \(L_j\). The finite field extension \(L_j\)
is complete by the coordinate-norm proof of
Theorem~\ref{thm:trace-transport-from-primes}. This gives
the continuous extensions. They commute with \(G_p\) since
the finite cyclotomic actions commute. Each \(R_j\) fixes
\(L_j\). Approximating any \(x\) by one element of a finite
layer and using the uniform norm bound for all larger layers
proves convergence.
\end{proof}

\begin{theorem}\label{thm:cyclotomic-period-vanishing}
Let \(a\ne0\) be an integer. If \(x\in C_p\) satisfies
\(g(x)=\chi_p(g)^a x\) for all \(g\) in a neighborhood of
\(1\) in \(G_p\), then \(x=0\).
\end{theorem}
\begin{proof}
First assume the equation holds for all \(g\). By
Theorem~\ref{thm:completed-fixed-field} for \(\ker\chi_p\)
and Lemma~\ref{lem:cyclotomic-trace-projections},
\(x\in\widehat L_\infty\). For each \(j\), the assignments
\(\zeta_s\mapsto\zeta_s^{1+p^j}\) are compatible automorphisms
of all the layers by the same lemma and extend to \(K_p\) by
Lemma~\ref{lem:place-fields}. The resulting \(g\) fixes \(L_j\)
and has \(\chi_p(g)=1+p^j\). Applying \(R_j\) to the equation
gives \(R_j(x)=(1+p^j)^aR_j(x)\). The rational multiplier
differs from one for every nonzero \(a\), so each \(R_j(x)=0\).
Their convergence proves \(x=0\).

In the general case, \(g*x=\chi_p(g)^{-a}g(x)\) defines a
group action on the underlying set of \(C_p\), and its stabilizer
at \(x\) contains the subgroup fixing some finite normal
extension by Definition~\ref{def:completed-local-action}.
The latter has finite index since restriction to the extension
takes values in a finite permutation group of its conjugates.
Thus the twisted orbit \(B\) is finite and nonempty. If \(x\ne0\)
its members are all nonzero. The product \(y=\prod_{b\in B}b\)
satisfies \(g(y)=\chi_p(g)^{a|B|}y\) for every \(g\), by the
permutation of \(B\) under the twisted action. The first case
applied to \(a|B|\ne0\) gives \(y=0\), a contradiction.
\end{proof}

\section{A single local frame and the spherical rank}

\begin{lemma}\label{lem:diagonal-semi-linear-subspace}
Let \(w_1,\ldots,w_b\in\Z\) and \(W\subset C_p^b\) be a
linear subspace. Suppose an open subgroup \(H\subseteq G_p\)
preserves \(W\) under the transformations
\[
               S_g(v)_i=\chi_p(g)^{w_i}g(v_i)
                    \qquad (g\in H).
\]
There is a basis of \(W\) whose individual vectors are supported
on positions of one weight and have all coordinates in \(K_p\)
fixed by \(H\). The projection of every \(v\in W\) onto any
one weight still belongs to \(W\).
\end{lemma}
\begin{proof}
Choose positions \(I\subseteq\{1,\ldots,b\}\) such that
projection \(W\to C_p^I\) is an isomorphism, by selecting
independent rows in a matrix for any basis of \(W\). This
includes \(I=\varnothing\) when \(W=0\). Let \(v^{(j)}\)
be the inverse image of the \(j\)-th unit vector for \(j\in I\).
On \(I\), the vector \(S_g(v^{(j)})\) agrees with
\(\chi_p(g)^{w_j}v^{(j)}\), hence the two agree everywhere.
A nonzero entry \(v_i^{(j)}\) therefore has \(w_i=w_j\) by
Theorem~\ref{thm:cyclotomic-period-vanishing}. It is then fixed
by \(H\). Every element of \(C_p\) fixed by \(H\) is algebraic
over \(\Q_p\): choose a finite normal extension \(L/\Q_p\)
whose fixing subgroup lies in \(H\). Its fixed elements in
\(K_p\) are precisely \(L\) by the extension-of-embeddings
argument of Lemma~\ref{lem:place-fields}. The field \(L\) is
complete by Theorem~\ref{thm:trace-transport-from-primes},
so is closed in \(C_p\). Apply
Theorem~\ref{thm:completed-fixed-field} to its fixing subgroup.
The resulting basis of \(W\) has the stated support, which
also proves the assertion about projections.
\end{proof}

\begin{proposition}\label{prop:polynomial-on-hodge-orbit}
Let \(x_i\in C_p\) and \(a_i\in\Z\) for \(1\le i\le b\).
For a polynomial \(F\in C_p[X_1,\ldots,X_b]\), if
\[
       F\big((\chi_p(g)^{-a_i}g(x_i))_i\big)=0
                         \qquad(g\text{ near }1\text{ in }G_p),
\]
then \(F((T^{a_i}x_i)_i)=0\) in \(C_p[T,T^{-1}]\).
\end{proposition}
\begin{proof}
Expand \(F=\sum_\alpha f_\alpha X^\alpha\) over its finite set
of monomials and set \(v_\alpha=\sum_i\alpha_i a_i\).
Choose an open subgroup \(H\) on which the hypothesis holds.
In the coordinate space indexed by the monomials, take the
span of \((\chi_p(g)^{-v_\alpha}g(x^\alpha))_\alpha\) over
all \(g\in H\). It is preserved by the transformations of
Lemma~\ref{lem:diagonal-semi-linear-subspace} with weights
\(-v_\alpha\), since \(h\) carries the vector indexed by
\(g\) to the one indexed by \(hg\). The linear form with
coefficients \(f_\alpha\) vanishes on this span. Project the
vector indexed by \(1\) to the positions with \(v_\alpha=d\).
It remains in the span by that lemma, so
\(\sum_{v_\alpha=d}f_\alpha x^\alpha=0\). These are the
coefficients of the required Laurent identity.
\end{proof}

\begin{definition}\label{def:one-local-hodge-frame}
Fix a continuous homomorphism \(j:G_p\to G_E\). In particular,
for a chosen embedding \(E\subset\Q_p\) at a completely split
prime, view its algebraic closure as the elements of \(K_p\)
algebraic over \(E\), and use restriction. That subfield is
algebraically closed by transitivity of algebraicity and the
algebraic closure in Lemma~\ref{lem:place-fields}. Restriction
is continuous since fixing a finite list in that subfield
follows from fixing its finite normal closure over \(\Q_p\).
For a representation \(r:G_E\to{\rm GL}_M(K_p)\), a local
Hodge--Tate frame of weights \(\kappa_1,\ldots,\kappa_M\)
at \(j\) is an invertible matrix \(P\) over \(C_p\) satisfying
\[
       r(j(g))\,g(P)
         = P\,\operatorname{diag}(\chi_p(g)^{-\kappa_i})
                         \qquad (g\text{ near }1).
\]
The action \(g(P)\) is applied to the entries.
\end{definition}

\begin{theorem}\label{thm:one-frame-cocharacter}
A frame as in Definition~\ref{def:one-local-hodge-frame}
gives the local polynomial closure of
Definition~\ref{def:local-polynomial-cocharacter} in the basis
\(Q=P^{-1}\) over \(k=C_p\). The multiset of its weights
depends only on the representation up to \(K_p\)-isomorphism.
For a \(G_E\)-invariant subspace over \(K_p\), the subspace
and quotient admit such frames, and their multisets together
are the multiset of \(r\).
\end{theorem}
\begin{proof}
The frame identity gives
\(Qr(j(g))=\operatorname{diag}(\chi_p(g)^{-\kappa_i})g(Q)\).
If \(F(Qr(t)Q^{-1})=0\) for \(t\) near \(1\) in \(G_E\),
continuity of \(j\) gives this vanishing at \(t=j(g)\) near
\(1\) in \(G_p\). Apply Proposition~\ref{prop:polynomial-on-hodge-orbit}
to the entries of \(Q\), assigning weight \(\kappa_i\) to
every entry of its \(i\)-th row, with polynomial \(B\mapsto
F(BQ^{-1})\). Its Laurent identity is
\(F(\operatorname{diag}(T^{\kappa_i}))=0\).

For frames \(P,P'\) of the same representation,
\(A=P^{-1}P'\) satisfies
\(g(A_{ij})=\chi_p(g)^{\kappa_i-\kappa'_j}A_{ij}\) near \(1\).
Theorem~\ref{thm:cyclotomic-period-vanishing} kills entries
linking unequal weights. Invertibility of \(A\) makes each
weight block an isomorphism between its column and row
coordinate spaces, giving equality of multiplicities. An
algebraic change of basis and its inverse have finitely many
entries fixed by an open subgroup of \(G_p\), so the same
conclusion applies across a \(K_p\)-isomorphism.

For an invariant subspace \(U\) choose a \(K_p\)-basis adapted
to \(U\), shrinking the subgroup to fix the entries of the
change of basis. Write \(U_C\) for its \(C_p\)-span. The space
\(P^{-1}U_C\) is stable under
\(z\mapsto\operatorname{diag}(\chi_p(g)^{-\kappa_i})g(z)\)
on this subgroup, since under \(P\) this becomes
\(v\mapsto r(j(g))g(v)\). Take the basis supplied by
Lemma~\ref{lem:diagonal-semi-linear-subspace} with its pivot
positions \(I\). Complete it by the coordinate vectors at
the nonpivot positions, which span the kernel of projection
onto \(I\). Each basis vector of \(P^{-1}U_C\) belongs to
the weight of its pivot and has entries fixed on the subgroup.
Multiplying these vectors by \(P\) gives a frame of \(U\).
The images of \(P\) times the nonpivot vectors form a frame
in the quotient. The frame identities follow on applying
\(r(j(g))g\) to these columns, in the adapted coordinates.
Each original position supplies exactly one weight in the
resulting two frames.
\end{proof}

\begin{corollary}\label{cor:one-embedding-geometric-criterion}
For each form in Theorem~\ref{thm:definite-seed-attachment},
suppose there exists a finite-dimensional continuous
\(K_\ell\)-representation of \(G_E\) annihilated by its
monic spherical polynomials on most Frobenius data. Suppose
it has a single local Hodge--Tate frame at the chosen split
embedding for \(w_1\), containing among its weights the \(n\)
specified shifted weights of the form. Then the given definite
eigensystem in that theorem has a continuous semisimple
rank-\(n\) representation with its spherical polynomials in
the approximation sense.
\end{corollary}
\begin{proof}
Take a chain of invariant subspaces for the representation of
each form and keep one copy of each distinct simple quotient
as in Lemma~\ref{lem:distinct-constituents-diagonal-span}.
The annihilators persist. By Theorem~\ref{thm:one-frame-cocharacter}
the successive frames account for the original multiset.
Discarding repeated isomorphism classes retains every weight
that occurs, since isomorphic factors have the same weights.
The frames of the retained factors form a block diagonal frame
of their direct sum, on the intersection of the finitely many
neighborhoods of \(1\). It yields the required diagonal family
by the same theorem. The sum has diagonal matrices in its
span and positive dimension, as at least one prescribed weight
occurs. Apply Corollary~\ref{cor:rank-test-for-definite-approximants}.
\end{proof}

\section{Local weights in a simultaneous Hecke quotient}

\begin{lemma}\label{lem:evaluation-on-weight-algebra}
Let \(D\) be a finite-dimensional vector space over a field \(k\),
let \((T_s)_{s\in S}\) be commuting endomorphisms, and fix
\(\theta_s\in k\). Polynomials in these variables always use
only finitely many of them. Then
\[
 D\Big/\sum_{s\in S}(T_s-\theta_s I)D\ne0
 \quad\Longleftrightarrow\quad
 F(T)|_D=0\ \Longrightarrow\ F(\theta)=0
             \quad\text{for every }F\in k[(X_s)_{s\in S}].
\]
\end{lemma}
\begin{proof}
On the quotient each \(T_s\) acts by \(\theta_s\) since the
subspace being divided out is stable under all \(T_s\) by
commutativity. Thus a nonzero quotient gives the implication.
If \(D=0\), the constant polynomial one violates it.
In positive dimension write \(B=k[(T_s)_s]\subseteq{\rm End}(D)\).
Under the indicated implication evaluation at \(\theta\) descends
to a unital algebra map \(B\to k\). Let \(\mathfrak m\) be its
kernel. If the quotient were zero, express every basis vector
\(d_i\) of \(D\) as a sum of \((T_s-\theta_s I)\) applied to
vectors, expanding those vectors in the same basis. This gives
\(d_i=\sum_j a_{ij}d_j\) for \(a_{ij}\in\mathfrak m\).
The adjugate identity over the commutative algebra \(B\)
gives \(\det(I-(a_{ij}))d_i=0\) for all \(i\). Hence that
determinant is the zero endomorphism in \(B\), whereas its
image in \(k\) is \(1\). This contradiction proves the converse.
\end{proof}

\begin{definition}\label{def:local-weight-spaces}
Given a frame \(P\) of Definition~\ref{def:one-local-hodge-frame}
for \(V\), put \(D_u(V)\subseteq C_p\otimes_{K_p} V\) equal
to the \(C_p\)-span of its columns of weight \(u\), using
the given basis of \(V\) to view these columns as vectors.
For a space over \(\Q_p\) first extend it to \(K_p\).
\end{definition}

\begin{proposition}\label{prop:local-weight-exactness}
The spaces of Definition~\ref{def:local-weight-spaces} are
intrinsic. A \(G_E\)-equivariant \(K_p\)-linear endomorphism
preserves each \(D_u(V)\). For an invariant \(K_p\)-subspace
\(U\subseteq V\) and quotient \(V/U\), the map on the weight
\(u\) spaces is onto \(D_u(V/U)\) with kernel \(D_u(U)\).
\end{proposition}
\begin{proof}
The change-of-frame calculation in
Theorem~\ref{thm:one-frame-cocharacter} makes the matrix linking
any two frames zero between different weights. This gives the
same column spans. An algebraic change of basis transports
them as subspaces of the tensor product by the same proof.
For an equivariant endomorphism \(A\), restrict to an open
subgroup fixing its finitely many entries and satisfying the
frame identity. For \(B=P^{-1}AP\), commuting \(A\) past
\(r(j(g))\) gives \(g(B_{ij})=\chi_p(g)^{\kappa_i-\kappa_j}
B_{ij}\), as in the change-of-frame calculation. Thus the same
vanishing by Theorem~\ref{thm:cyclotomic-period-vanishing}
shows preservation of each \(D_u(V)\).

Use the basis for the subspace and its coordinate complement
constructed in Theorem~\ref{thm:one-frame-cocharacter}. The
subspace has a frame by vectors in the respective \(D_u(V)\),
and the remaining frame vectors project to a frame of the
quotient. In each weight these vectors together are a basis
of \(D_u(V)\), by the basis decomposition with pivot positions
within the same weight. Quotienting gives the stated kernel
and image.
\end{proof}

\begin{theorem}\label{thm:hecke-quotient-weight-detection}
Let \(V=\bigoplus_i V_i\) be a finite sum of finite-dimensional
\(K_p\)-spaces with continuous \(G_E\)-actions. Let \(T_s\),
\(s\in S\), be a commuting family of \(K_p\)-linear operators
preserving each summand and commuting with \(G_E\). For
\(\theta_s\in K_p\) there is a continuous \(G_E\)-quotient
\(V_\theta\) universal among linear quotients on which every
\(T_s\) acts by \(\theta_s\). Assume \(V\) has a frame as in
Definition~\ref{def:one-local-hodge-frame}. The quotient admits
a frame, and \(u\) occurs in its weights exactly when
\[
     F(T)D_u(V)=0\ \Longrightarrow\ F(\theta)=0
                \quad\text{for every }F\in C_p[(X_s)_{s\in S}].
\]
\end{theorem}
\begin{proof}
Use \(V_\theta=V/U\) with
\(U=\sum_s(T_s-\theta_s I)V\). This is an invariant subspace
under the group and all the operators. Every linear map
turning the operators into their scalar values vanishes on
\(U\) and hence factors uniquely through this quotient.
In a basis adapted to \(U\) the quotient action is a diagonal
block of the original continuous matrices, proving continuity.

Over \(C_p\), the span \(U_C\) equals
\(\sum_s(T_s-\theta_s I)(C_p\otimes V)\): both sides are the
span of the images of any original basis for all \(s\).
Each operator preserves \(D_w(V)\) by
Proposition~\ref{prop:local-weight-exactness}. Splitting
a vector into its weight components therefore gives
\[
   U_C\cap D_u(V)=\sum_s(T_s-\theta_s I)D_u(V).
\]
The quotient has a frame by Theorem~\ref{thm:one-frame-cocharacter},
and its weight space is the quotient of \(D_u(V)\) by this
intersection by Proposition~\ref{prop:local-weight-exactness}.
Apply Lemma~\ref{lem:evaluation-on-weight-algebra} over \(C_p\).
\end{proof}

\begin{corollary}\label{cor:congruence-quotient-spherical}
For an approximating form in Theorem~\ref{thm:definite-seed-attachment},
put \(p=\ell\). Let \(V,(T_s)_s,\theta\) satisfy the hypotheses
of Theorem~\ref{thm:hecke-quotient-weight-detection} with a frame
at \(w_1\). Assume its polynomial test there holds for each of
the \(n\) specified shifted weights. At almost all degree-one
primes \(v\) let \(T_{v,h}\), \(1\le h\le n\), be among the
operators with \(\theta_{v,h}\) the eigenvalues defining the
form's \(P_v\) in Definition~\ref{def:actual-spherical-polynomial}.
Write \(T_{v,0}=I\), \(q_v\) for the norm of \(v\), and
\(r_V\) for the action. If for most Frobenius data
\[
       \sum_{h=0}^n(-1)^h q_v^{h(h-1)/2}
                      T_{v,h}\,r_V(\gamma)^{\,n-h}=0
                      \quad\text{on }V,
\]
then \(V_\theta\) supplies the representation for this form
in Corollary~\ref{cor:one-embedding-geometric-criterion}.
\end{corollary}
\begin{proof}
Theorem~\ref{thm:hecke-quotient-weight-detection}
puts all the indicated weights in a frame of \(V_\theta\)
and gives the continuous quotient action. On this quotient
each term in the sum descends with \(T_{v,h}\) replaced
by \(\theta_{v,h}\). The result by
Definition~\ref{def:actual-spherical-polynomial} is exactly
\(P_v(r_{V_\theta}(\gamma))=0\). Removing finitely many
exceptional primes respects the quantifier in
Definition~\ref{def:frobenius-approximation}.
\end{proof}

\section{Integral polarization data for the cyclotomic compositum}

\begin{lemma}\label{lem:actual-signs-and-padding}
Let \(K,E,N_1,f\) be as in Definition~\ref{def:levels} and
Proposition~\ref{prop:places}. Write \(F=E^{\mathfrak c}\)
for the fixed field of complex conjugation, \(\beta=(1+\sqrt{-d})/2\),
and \(b=(d+1)/4\). Fix an embedding \(\sigma:E\to\C\), a finite
set \(B\) of rational primes, and \(P\) the product of the
distinct primes in \(B\cup\{2\}\). There are a set \(\Psi\) of
one embedding from each conjugate pair, an integer \(R\ge3\)
and \(u\in O_F\) with the following properties.
\begin{enumerate}
\item \(\sigma\in\Psi\) and each other \(\tau\in\Psi\) restricts
to the conjugate of \(\sigma|_K\). In particular
\(\tau(\beta)=\sigma(\bar\beta)\) there, while
\(\sigma(\beta)\ne\sigma(\bar\beta)\).
\item \(u\equiv1\bmod PO_F\), \(R\equiv1\bmod P\), and
\(\Im\tau(u/\varpi)>0\) at every \(\tau\in\Psi\), with
\(\varpi\) the different generator of Proposition~\ref{prop:different}.
\item For each \(c\in\{1,\beta\}\) there exist \(R-3\)
elements \(a_t\in O_E\) spanning \(O_E\) over \(\Z\) and satisfying
\[
             c\bar c+3+\sum_{t=1}^{R-3}a_t\bar a_t=R.
\]
\end{enumerate}
\end{lemma}
\begin{proof}
The two embeddings in each pair restrict differently to the
imaginary quadratic field \(K\), giving exactly the required
choice outside the pair of \(\sigma\). The two values of
\(\beta\) differ by \(\sigma(\sqrt{-d})\ne0\). Conjugation has
fixed field \(F\) of degree \(f\): every \(x\in E\) is the sum
of \((x+\bar x)/2\) and \(\sqrt{-d}\) times
\((x-\bar x)/(2\sqrt{-d})\), both coefficients in \(F\).
The restrictions of the \(f\) pairs give its \(f\) real
embeddings, since \(E=F(\sqrt{-d})\). Under these embeddings
\(PO_F\) is a full lattice in \(\R^f\) by
Lemma~\ref{lem:number-field-lattice}. Every point is within a
fixed distance of \(1+PO_F\): subtract \(1\), expand in a
lattice basis, and round its coordinates. The nonzero
\(\tau(1/\varpi)\) are purely imaginary by
Proposition~\ref{prop:different}. Choose \(u\) near a vector
whose coordinates are a common large positive length times
the signs of \(\Im\tau(1/\varpi)\). Taking the length larger
than the error proves the second assertion about \(u\).

Take any \(R\equiv1\bmod P\) with \(R\ge3+b+f(1+b)\).
The basis \(\zeta_{N_1}^j,\beta\zeta_{N_1}^j\), \(0\le j<f\),
of Theorem~\ref{thm:integer-ring} has \(2f\) entries and sum
of relative norms \(f(1+b)\), since \(\beta\bar\beta=b\).
Start with this basis and append \(R-3-c\bar c-f(1+b)\)
entries equal to one. This number is nonnegative. The
remaining number of positions equals \(c\bar c+f(b-1)\ge0\).
Filling them with zeros proves the exact length, span and norm.
\end{proof}

\begin{theorem}\label{thm:actual-pel-torsion-pairing}
Fix \(n\ge1\) and the choices of
Lemma~\ref{lem:actual-signs-and-padding}. Put
\(\delta=R^2u\in O_F\), \(\xi=\delta/\varpi\in E\),
\(h=|N_{F/\Q}(\delta)|\), and
\(L=(\prod_p O_p)^n\subseteq\mathbb A_{E,f}^n\),
where the rings and the character \(\psi_f\) are as in
Proposition~\ref{prop:additive}. For finite adelic column
vectors define
\[
     e(x,y)=\psi_f\!\left(\xi\sum_{a=1}^n\bar y_a x_a\right).
\]
This pairing is biadditive with multiplicative values and is
alternating. The scalar \(\xi\) is purely imaginary and has
positive imaginary part throughout \(\Psi\).
Writing \(L^\perp=\{x:e(x,L)=1\}\), one has
\[
      L^\perp=\delta^{-1}L,\qquad
      T=L^\perp/L,\qquad |T|=h^{2n},\qquad hT=0.
\]
Here \(h\) and \(R\) are prime to \(P\). There is a subgroup
\(\Lambda\) with \(R^{-1}L\subseteq\Lambda\subseteq L^\perp\)
and \(\Lambda^\perp=\Lambda\). Its image in
\(\mathbb A_{E,f}^n/L\) contains the full \(R\)-torsion.
\end{theorem}
\begin{proof}
The character is additive, so \(e\) is multiplicative under
addition in either argument. Locally the trace of
\(\xi\sum_a\bar x_a x_a\) is its own negative by conjugation,
hence zero. This gives \(e(x,x)=1\), and evaluation at \(x+y\)
gives \(e(y,x)=e(x,y)^{-1}\). The sign of \(\xi\) follows from
Lemma~\ref{lem:actual-signs-and-padding} and \(R^2>0\).
Proposition~\ref{prop:additive}, applied in one coordinate at a
time and allowing the inverse character, identifies the dual
of \(L\) as \(\delta^{-1}L\). It contains \(L\) since
\(\delta\) is integral.

Multiplication by \(\delta\) on a \(\Z\)-basis of \(O_F\)
has determinant of absolute value \(h\), by
Lemma~\ref{lem:number-field-lattice}. Its adjugate shows that
\(h/\delta\in O_F\), by applying the inverse map times \(h\)
to \(1\). As \(\delta\equiv1\bmod PO_F\), the multiplication
matrix is congruent to the identity modulo \(P\). Thus
\(h\) is prime to \(P\). Multiplication by \(\delta\) identifies
\(T\) with \(L/\delta L\), and \(h\) kills it. To count it,
the diagonal \(O_E\) surjects onto \((\prod_p O_p)/h(\prod_p O_p)\)
with kernel \(h O_E\) by the Chinese remainder theorem on each
coordinate of the integral basis in Theorem~\ref{thm:integer-ring}.
As \(h/\delta\) is integral, reducing further gives
\((\prod_p O_p)/\delta(\prod_p O_p)\simeq O_E/\delta O_E\).
Indeed modulo \(h\), multiplication by \(\delta\) corresponds
under this isomorphism on both sides. The index in \(O_E\)
is \(|N_{E/\Q}\delta|=h^2\): multiplication scales its lattice
covolume by its real determinant and the embeddings of \(E\)
occur in pairs giving the same value on \(\delta\). Taking
the \(n\) coordinates proves the asserted order.

The pairing on \(L^\perp\) descends to \(T\) by its definition
and skew symmetry. The subgroup \(R^{-1}L/L\) is isotropic:
on \(R^{-1}L\) both denominators cancel the \(R^2\) in \(\xi\),
and \(uL\subseteq L\) in the pairing of Proposition~\ref{prop:additive}.
Take a largest isotropic subgroup of the finite \(T\) containing
it. Any class pairing trivially with that subgroup could be
adjoined with all its integer multiples and preserve isotropy
by alternation. Thus it already belongs to the subgroup. Its
preimage \(\Lambda\) contains \(L\), so any vector pairing
trivially with \(\Lambda\) is in \(L^\perp\). The argument in
\(T\) gives \(\Lambda^\perp=\Lambda\). Vectors representing
\(R\)-torsion modulo \(L\) are precisely \(R^{-1}L\).
\end{proof}

\section{Degree, level and good primes of the polarized data}

\begin{lemma}\label{lem:finite-degree-cut}
Let \(J\) be a finite set, \(R\ge1\) an integer, and
\(\sigma:\Z_{\ge0}\to\Z_{\ge0}\) any function. A monomial
ideal of exponents is a subset \(I\subseteq\Z_{\ge0}^{J}\)
closed under addition by all nonnegative exponent vectors.
There is an integer \(i\ge1\) such that, for every such \(I\),
\[
  \#\{\alpha\notin I:\textstyle\sum_t\alpha_t=R^j\}\le\sigma(j)
   \quad(0\le j\le i)
\]
implies the same inequalities for every \(j\ge0\).
\end{lemma}
\begin{proof}
Every such ideal is the union of finitely many translates of
the nonnegative orthant. Otherwise choose successively an
element outside the union of the translates of all earlier
choices. Any infinite sequence of natural numbers has an
infinite nondecreasing subsequence: some value repeats
infinitely or the sequence eventually exceeds each fixed
bound. Selecting successively in each of the finitely many
coordinates gives two selected vectors comparable in their
original order, a contradiction.

If the asserted cutoff failed, choose for every \(t\ge1\) an
ideal \(I_t\) satisfying the counts through \(t\) and violating
one subsequently. There are at most countably many exponent
vectors. Passing to a diagonal subsequence makes membership
of each vector eventually constant. The vectors eventually
present form a monomial ideal \(I\). At any one degree the
possible exponent vectors form a finite set, so \(I\) satisfies
every count, by eventual agreement and the increasing indices.
All finitely many generators of \(I\) are eventually present
simultaneously in the subsequence. Every later ideal there
contains \(I\) and hence satisfies all counts, a contradiction.
\end{proof}

\begin{proposition}\label{prop:actual-integral-moduli-level}
Use \(E,n,P,R,\delta,h,L,T,\Lambda\) of
Theorem~\ref{thm:actual-pel-torsion-pairing}, and set
\(A=\mathbb A_{E,f}^n/L\). For all \(j\ge0\),
\[
 \#\{y\in A:R^j y\in T\}=h^{2n}R^{2fnj}.
\]
There exist \(i\ge1\) and a positive integer \(N=R^i h\)
prime to \(P\) with both of the following properties. The
cutoff of Lemma~\ref{lem:finite-degree-cut} holds with \(J=T\),
\(\sigma(j)=h^n R^{fnj}\), and this \(i\). For any nonzero
\(\mathfrak n\in NO_E\), every \(y\in A\) with \(R^j y\in T\)
for some \(0\le j\le i\) is killed by \(\mathfrak n\).
\end{proposition}
\begin{proof}
Multiplication by \(R^j\) is onto \(A\), since division by
the rational integer is possible in the finite adeles. Its
kernel \(R^{-j}L/L\) has cardinal \(R^{2fnj}\): in the
\(2f\) integral basis coordinates of Theorem~\ref{thm:integer-ring}
use the Chinese remainder argument of
Theorem~\ref{thm:actual-pel-torsion-pairing} with \(R^j\).
All fibres have this size. The same theorem gives
\(|T|=h^{2n}\), hence the count and its integer square root
\(\sigma(j)\). Choose \(i\) by Lemma~\ref{lem:finite-degree-cut}.
Since \(hT=0\), \(R^j y\in T\) implies \(hR^j y=0\).
For \(j\le i\) the integer \(hR^j\) divides \(N\), proving
the claim for all its integral multiples. The coprimality
follows from Theorem~\ref{thm:actual-pel-torsion-pairing}.
\end{proof}

\begin{lemma}\label{lem:second-quadratic-prime}
Let \(w\) be any finite prime of \(E\), above a prime \(u_0\)
of \(K=\Q(\sqrt{-d})\) over a rational prime \(p\).
There exists a prime \(v\ne w\) of \(E\) above the same
\(u_0\).
\end{lemma}
\begin{proof}
Choose two distinct primes from \(3,5,7\) different from \(p\)
and denote their product by \(m\). Then \(m\mid N_1\) by
Definition~\ref{def:levels}. The group
\((\Z/m\Z)^\times\) has at least four elements of square one
by independent choices of signs at the two primes. No cyclic
group has more than two such elements, since in a cyclic
group of order \(a\) the solutions correspond to
\(2t=0\bmod a\). Choose a unit class outside the subgroup
generated by \(p\bmod m\), and lift it to a unit \(c\bmod N_1\)
by the surjectivity proved in Lemma~\ref{lem:induce}.
Proposition~\ref{prop:places} gives \(\gamma\in{\rm Gal}(E/K)\)
acting on \(\zeta_m\) by exponent \(c\).

In the residue field \(k_w\) the image of \(\zeta_m\) has
order exactly \(m\). Indeed reduction is injective on the
\(m\)-th roots for \(p\nmid m\) by
Proposition~\ref{prop:local-cyclo-different}. Every
automorphism of the finite residue field is a power of the
\(p\)-power map. To see this for a field of \(p^a\) elements,
each nonzero element has \((p^a-1)\)-st power one by multiplying
over its permutation of the multiplicative group. The first
\(a\) powers of Frobenius are distinct by the root bound
for \(X^{p^b}-X\), \(0<b<a\). A degree-\(a\) extension
has at most \(a\) embeddings, by adjoining generators one
by one and counting roots of their minimal polynomials.
This proves the assertion, including \(a=1\). If \(\gamma\)
preserved \(w\), its induced automorphism would send the
primitive residue root to both its \(c\)-th and its
\(p^j\)-th power for some \(j\), contrary to the choice
of \(c\). Thus \(v=\gamma w\ne w\), with contraction \(u_0\).
\end{proof}

\begin{theorem}\label{thm:actual-pel-safe-level}
Fix the choices of Proposition~\ref{prop:actual-integral-moduli-level}
and \(0\ne\mathfrak n\in NO_E\). For a rational prime \(p\)
split into two distinct primes \(u_0,\bar u_0\) in \(K\),
call \(w\) over \(u_0\) safe if \(\mathfrak n\notin w\) and
\(\mathfrak n\notin v\) for every \(v\) over \(\bar u_0\).
For every such \(w\) there is an integer \(b_0\) satisfying
\[
 \begin{gathered}
  p,\ \beta-b_0\in w,\qquad
  (\beta-b_0,\bar\beta-b_0)=O_E,\qquad R,\, |T|\notin w,\\
  (\mathfrak n)+w\,(p,\bar\beta-b_0)=O_E .
 \end{gathered}
\]
A second prime \(v\ne w\) contains \(p,\beta-b_0\).
Here parentheses denote the ideal generated by their entries.
\end{theorem}
\begin{proof}
By Lemma~\ref{lem:quadratic-norm}, \(O_K=\Z[\beta]\),
with polynomial \(X^2-X+(d+1)/4\). Modulo \(p\) its two
distinct prime factors correspond to two distinct roots in
\(\mathbf F_p\), since the quotient is the polynomial quotient
and has degree two. In particular \(p\ne d\), since the
discriminant is \(-d\). Take \(b_0\) congruent to \(\beta\bmod u_0\)
and to \(0\bmod d\), by Lemma~\ref{lem:basic-arith}.
A common prime divisor of \(\beta-b_0\) and \(\bar\beta-b_0\)
would contain their difference \(\sqrt{-d}\), so would be
above \(d\). There \(2\beta=1+\sqrt{-d}\) reduces to \(1\),
whereas \(2b_0\) reduces to zero. This contradiction gives
the unit ideal, since a proper ideal lies in a maximal ideal
as in Proposition~\ref{prop:general-ideal-factorization}.
The other root mod \(p\) belongs to the conjugate prime,
so every prime of \(E\) containing \(p,\bar\beta-b_0\) lies
over \(\bar u_0\). By safety neither these primes nor \(w\)
contains \(\mathfrak n\). A prime containing a product of
two ideals contains one of them, as seen by multiplying two
elements outside it. Thus the indicated ideal sum is \(O_E\).
Since \(\mathfrak n\) is a multiple of \(R^i h\) with \(i\ge1\),
its absence from \(w\) excludes both \(R\) and \(|T|=h^{2n}\).
Finally Lemma~\ref{lem:second-quadratic-prime} gives \(v\)
over \(u_0\), with both required elements in its contraction.
\end{proof}

\section{Weights at conjugate embeddings and scalar twists}

\begin{definition}\label{def:oriented-definite-jumps}
For a dominant weight \(\lambda=(\lambda_1,\ldots,\lambda_n)\)
of a definite form as in Definition~\ref{def:definite-padic-form},
write \(m_h=\lambda_h-\lambda_{h+1}\) for \(h<n\). At the
chosen CM representative \(\tau\) of the pair the oriented
jumps are \(j-1-\lambda_j\), \(1\le j\le n\). At its conjugate
\(\bar\tau\) they are \(\lambda_j+n-j\). Equivalently set
\[
      v_j=\lambda_n+\sum_{h=j}^{n-1}m_h+n-j .
\]
The multiset at \(\tau\) consists of the \(n-1-v_j\);
at \(\bar\tau\) it consists of the \(v_j\). The determinant
exponent of the weight is \(\lambda_n\).
\end{definition}

\begin{lemma}\label{lem:affine-free-jumps}
For any \(B\ge0\), a list of integers is free through degree
\(B\) in Definition~\ref{def:bounded-degree-free-weights}
if and only if its image under \(t\mapsto at+c\) is, for
any fixed integers \(c\) and \(a\ne0\). In particular both
oriented lists of Definition~\ref{def:oriented-definite-jumps}
at either modified place of
Theorem~\ref{thm:definite-wild-approximants}, with any one
integer added to all entries, are free through \((n+1)^2\).
\end{lemma}
\begin{proof}
For an exponent list of total degree \(h\), the pairing with
the transformed weights equals \(a\) times the original pairing
plus \(ch\). Equality at fixed degree is unchanged since
\(a\ne0\). A permutation also preserves freeness by permuting
the exponent lists. For \(u_j=u_j(\lambda)\) in
Definition~\ref{def:bounded-degree-free-weights},
telescoping the \(m_h\) gives
\[
  j-1-\lambda_j=u_j+n-1-\lambda_n,\qquad
  \lambda_j+n-j=-u_j+\lambda_n .
\]
Apply the first assertion to the free \(u_j\) of
Theorem~\ref{thm:local-algebraic-weight-density}. This includes
\(n=1\), where each total degree has one exponent list.
\end{proof}

\begin{theorem}\label{thm:rank-before-scalar-twist}
Let \(P_v\in K_p[X]\) be monic of degree \(n\ge1\) at the
degree-one primes of a number field \(E\) outside a finite set,
and take \(c_v\in K_p^\times\) there. Suppose a continuous
finite-dimensional \(K_p\)-representation \(V\) of \(G_E\)
has a frame at some \(j:G_p\to G_E\) in the sense of
Definition~\ref{def:one-local-hodge-frame}. Suppose some
\(n\) weights in that frame are free through \(n(n+1)/2\).
Assume
\[
               Q_v(X)=c_v^{-n}P_v(c_vX)
\]
annihilates the action on \(V\) for most Frobenius data.
Then the direct sum \(r_0\) of one copy of each distinct
simple constituent of \(V\) has rank \(n\) and characteristic
polynomial \(Q_v\) for most data. If a continuous character
\(a:G_E\to K_p^\times\) satisfies \(a(\gamma)=c_v\) for
most data, the representation \(r(g)=a(g)r_0(g)\) has
characteristic polynomial \(P_v\) for most data. For every
\(s\in G_E\) its eigenvalues are those of \(r_0(s)\)
multiplied by \(a(s)\), and
\(\det r(s)=a(s)^n\det r_0(s)\).
\end{theorem}
\begin{proof}
Lemma~\ref{lem:distinct-constituents-diagonal-span} yields a
continuous direct sum with diagonal matrices in its span and
the prescribed annihilators. The proof of
Corollary~\ref{cor:one-embedding-geometric-criterion} gives a
frame retaining each occurring weight. The proof applies to
any list of weights since it only uses the frames of subspaces
and quotients and their uniqueness from
Theorem~\ref{thm:one-frame-cocharacter}. That theorem supplies
the diagonal family, so
Theorem~\ref{thm:spherical-annihilator-rank} gives the claimed
rank and characteristic polynomial. For the second assertion
intersect the two conditions on most data: refine a nonempty
open set first for one condition, then its nonempty open
subset for the other, and take the union of the two finite
exclusions in Definition~\ref{def:frobenius-approximation}.
For \(a(\gamma)=c_v\) factoring \(c_v\) out of each row of
\(XI-c_vr_0(\gamma)\) gives \(c_v^nQ_v(X/c_v)=P_v(X)\).
This determinant identity at any single group element also
gives both assertions at \(s\). Scalar multiplication gives
a continuous group homomorphism since \(a\) is a continuous
character.
\end{proof}

\begin{corollary}\label{cor:twisted-hecke-quotient-seed}
For each approximating form of
Theorem~\ref{thm:definite-seed-attachment}, use the situation
of Corollary~\ref{cor:congruence-quotient-spherical}, allowing
as the \(n\) tested weights either of its oriented jump lists
at \(w_1\) with an integer \(\mu\) added throughout.
Allow the Hecke eigenvalues of the geometric quotient at \(v\)
to be \(c_v^{-h}t_{v,h}\) for \(0\le h\le n\), where
\(t_{v,h}\) are the coefficients' eigenvalues in
Definition~\ref{def:actual-spherical-polynomial}. If a
continuous character with the Frobenius values \(c_v\) is
given for each form, these conditions suffice for the
attachment conclusion of Theorem~\ref{thm:definite-seed-attachment}.
\end{corollary}
\begin{proof}
For each form, Lemma~\ref{lem:affine-free-jumps} gives
the freeness bound at \(w_1\). The polynomial test of
Theorem~\ref{thm:hecke-quotient-weight-detection} supplies
all these weights in the quotient. Descending the operator
congruence as in Corollary~\ref{cor:congruence-quotient-spherical}
gives coefficients
\((-1)^h q_v^{h(h-1)/2}c_v^{-h}t_{v,h}\) at \(X^{n-h}\).
These form precisely \(c_v^{-n}P_v(c_v X)\). Apply
Theorem~\ref{thm:rank-before-scalar-twist} before using
Theorem~\ref{thm:definite-seed-attachment}. Exact equality
for most data implies each tolerance of
Definition~\ref{def:frobenius-approximation}.
\end{proof}

\section{Principal levels with specified local conditions}

\begin{definition}\label{def:moduli-signed-frame}
Continue with the CM representatives \(\Psi\) and distinguished
\(\sigma\) in Lemma~\ref{lem:actual-signs-and-padding}. Put
\(\widehat O=\prod_p O_p\) and \(\widehat{\Z}=\prod_p\Z_p\).
For a matrix write \(a^*=\bar a^{\mathsf T}\). For
\(0\ne\mathfrak n\in O_E\) define
\[
 K(\mathfrak n)=
 \{a\in {\rm GL}_n(\widehat O): a^*a=tI
       \text{ for some }t\in\widehat{\Z}^{\,\times},\
                a-I\in\mathfrak n M_n(\widehat O)\}.
\]
A signed frame consists of \(x\in{\rm GL}_n(\mathbb A_{E,f})\),
\(J\in M_n(\C)\), and a matrix \(A\in M_n(E)\) with
\(x^*Ax=\varepsilon I\) for some \(\varepsilon\in
\widehat{\Z}^{\,\times}\), such that \(\tau(A)\) is positive
definite Hermitian for all \(\tau\in\Psi\setminus\{\sigma\}\)
and \(\sigma(A)J\) is positive definite Hermitian. Its stabilizer
at this level consists of the \(y\in{\rm GL}_n(E)\) with
\(x^{-1}yx\in K(\mathfrak n)\) and \(\sigma(y)J=J\sigma(y)\).
\end{definition}

\begin{proposition}\label{prop:neat-at-wild-place}
If \(\mathfrak n\) belongs to any prime \(w\) of \(E\) above
an odd rational prime \(\ell\) with \(E_w=\Q_\ell\), every
stabilizer in Definition~\ref{def:moduli-signed-frame} is trivial.
\end{proposition}
\begin{proof}
For \(y\) in the stabilizer set \(a=x^{-1}yx\). Multiplying
the two frame identities gives \(y^*Ay=tA\) at the finite
adeles. Some entry of the invertible \(A\) is nonzero.
The quotient of the corresponding entries in this equation
is an element of \(E\) equal diagonally to \(t\). In a rational
basis of \(E\) starting with \(1\), equality in \(E\otimes\Q_p\)
forces all its other rational coordinates to vanish. Thus
\(t\in\Q^\times\) and is a unit at every prime, giving
\(t=\pm1\) by integer factorization. There is a
\(\tau\ne\sigma\) in \(\Psi\), since \(f=\varphi(N_1)>1\)
by Proposition~\ref{prop:places}. Positivity of \(\tau(A)\)
there implies \(t=1\). Thus \(y\) is unitary for each
\(\tau(A)\) at these embeddings. At \(\sigma\), commuting
with \(J\) gives
\(\sigma(y)^*\sigma(A)J\sigma(y)=\sigma(A)J\), so it is
also unitary for this positive matrix.

For a fixed positive definite form the entries of a unitary
matrix are bounded: its columns have fixed squared lengths
for that form, and the minimum of that form on the Euclidean
unit sphere is positive by compactness. Hence the entries
of \(y\) at all embeddings lie in a bounded box depending
on the frame. Choose a positive integer \(b\) clearing the
entries of both \(x\) and \(x^{-1}\) into \(\widehat O\),
possible by their finite adelic denominators. Then all entries
of \(b^2y=(bx)a(bx^{-1})\) belong diagonally to \(\widehat O\)
and hence to \(O_E\), by rational coordinates in the basis
of Theorem~\ref{thm:integer-ring}. There are only finitely
many such matrices in the box, by the lattice property of
Lemma~\ref{lem:number-field-lattice}. The stabilizer is a
group, so each of its elements has finite order.
At \(w\), its conjugate \(a_w\) belongs to
\(I+\ell M_n(\Z_\ell)\), which is torsion-free for odd
\(\ell\) by the binomial argument in Proposition~\ref{prop:column-dim}.
Thus \(a_w=I\). The embedding \(E\to E_w\) is injective,
giving \(y=I\).
\end{proof}

\begin{lemma}\label{lem:principal-level-support}
Let \(D\) be a finite set of finite primes of \(E\), and
assign an integer \(k_w\ge1\) to each \(w\in D\). For any
positive integer \(N\) there is a nonzero
\(\mathfrak n\in NO_E\) such that
\[
   \mathfrak n\in w^{k_w}\ (w\in D),\qquad
   \{v:\mathfrak n\in v\}=D\cup\{v:N\in v\}.
\]
Let \(K^0\) be the standard definite lattice stabilizer of
Definition~\ref{def:definite-comparison-forms}. If a subgroup
\(H\) contains every element \(g\in K^0\) congruent to \(I\)
modulo \(w^{k_w}\) at the places of \(D\), then for every finite
standard unitary isometry \(g\) with \(g^{-\mathsf T}\in
K(\mathfrak n)\) one has \(g\in H\).
\end{lemma}
\begin{proof}
The product \(\mathfrak a=\prod_{w\in D}w^{k_w}\) has a
positive power principal by Lemma~\ref{lem:bounded-ideal-classes}.
Write \(\mathfrak a^m=zO_E\), \(z\ne0\). Integrality of the ideal
gives \(z\in O_E\). Use \(\mathfrak n=Nz\); its support and
the lower bounds on its valuations follow by
Proposition~\ref{prop:general-ideal-factorization}, also for
the unit ideal when \(D\) is empty. If \(g^{-\mathsf T}\)
is integral and invertible over \(\widehat O\), so is \(g\);
since it is a unitary isometry, \(g\in K^0\). Taking inverse
transpose preserves congruence to \(I\) over every quotient
ring, hence \(g\) satisfies the required conditions in \(D\).
\end{proof}

\begin{corollary}\label{cor:actual-moduli-places}
For one approximating form \(f\) of
Theorem~\ref{thm:definite-wild-approximants}, choose the
coefficient prime \(\ell\) odd, as permitted by
Lemma~\ref{lem:auxiliary-split-embeddings}. Use \(w_0\)
for the place of its cuspidal character and \(w_1\) for the
free-weight spherical place above the same prime \(u_0\)
of \(K\). Choose another completely split prime \(p_*\ne\ell\)
at which \(f\) is spherical at every place, and a place
\(w_*\mid p_*\). There are polarization data of
Theorem~\ref{thm:actual-pel-torsion-pairing} and a principal
level \(\mathfrak n\) as in Proposition~\ref{prop:actual-integral-moduli-level}
for which all signed-frame stabilizers are trivial and \(w_1\)
is safe. The level is divisible by \(w_0^2\) and \(w_*\)
and avoids \(\bar w_0,\bar w_*\). Every finite unitary
\(g\) with \(g^{-\mathsf T}\in K(\mathfrak n)\) fixes \(f\)
by right translation.
\end{corollary}
\begin{proof}
Lemma~\ref{lem:auxiliary-split-embeddings} supplies \(p_*\)
avoiding the finitely many exceptional rational primes and
\(\ell\). Include \(p_*,\ell\) in \(B\) of
Lemma~\ref{lem:actual-signs-and-padding}, so that \(N\)
from Proposition~\ref{prop:actual-integral-moduli-level}
is a unit above both. At \(\ell\), the level of \(f\) in
Theorem~\ref{thm:definite-wild-approximants} contains the full
integral groups except for the factor indexed by the pair of
\(w_0\); its character of Proposition~\ref{prop:actual-wild-split-character}
is trivial on \(1+\ell^2 M_n(\Z_\ell)\). Impose congruence
of depth two at \(w_0\), and of depth one at \(w_*\).
At each of the finitely many other exceptional factors impose
a sufficiently deep principal congruence fixing \(f\), possible
by openness of its stabilizer in the restricted product. Use
one representative of each split pair, since for a unitary
matrix its other component is the inverse transpose after
conjugating the local field, as in Lemma~\ref{lem:split-level}.
At a nonsplit prime impose the congruence in its one factor.
By taking a product neighborhood inside the stabilizer,
these conditions and integrality elsewhere suffice to fix
\(f\). Apply Lemma~\ref{lem:principal-level-support} with
these places and depths. Over \(\ell\) the only selected
prime is \(w_0\). In particular the level meets neither \(w_1\)
nor a prime over \(\bar u_0\); it is also a unit at \(\bar w_*\).
Proposition~\ref{prop:neat-at-wild-place} applies at \(w_0\).
Theorem~\ref{thm:actual-pel-safe-level} applies at \(w_1\)
because \(\ell\) splits in \(K\).
\end{proof}

\section{Extension of a definite vector to the general linear group}

\begin{definition}\label{def:definite-vector-extension}
Let \(U\) denote the finite adelic isometries of the standard
Hermitian form in rank \(n\), with rational subgroup \(U(F)\)
as in Definition~\ref{def:definite-comparison-forms}, and set
\(G={\rm GL}_n(\mathbb A_{E,f})\), \(\Gamma={\rm GL}_n(E)\).
Write \(\alpha(g)=g^{-\mathsf T}\). It is an involutive
automorphism, and on \(U\) it equals entrywise conjugation.
Use \(C=K(\mathfrak n)\) from Definition~\ref{def:moduli-signed-frame}.
For a representation \(\Xi:\Gamma\to{\rm GL}(W)\) on a
finite-dimensional complex space, a definite function at this
level is a function \(\phi:G\to W\) with
\(\phi(y g k)=\Xi(y)\phi(g)\) for \(y\in\Gamma,\ k\in C\),
whose nonzero values occur only at \(g\) admitting a signed
frame at \(J=I\). On either group write
\((R(a)\phi)(g)=\phi(ga)\) for right translation.
\end{definition}

\begin{lemma}\label{lem:rational-weight-extension}
For a complex standard definite form \(f\) of weight
\(\lambda=(\lambda^\tau)_\tau\) as in
Definition~\ref{def:definite-padic-form} before coefficient
transport, there is a representation \(\Xi\) such that
\[
             f(tu)=\Xi(\alpha(t))f(u)
                           \qquad (t\in U(F),\ u\in U).
\]
It may be taken to act by a scalar on all scalar rational
matrices. More generally choose, for each conjugate pair of
embeddings, its original representative \(\tau\). In the
\(\tau\)-factor one may use either \(\tau(y)\) or
\((\bar\tau(y))^{-\mathsf T}\) as the argument of
\(W_{\lambda^\tau}\). One may also multiply the tensor
representation by \(N_{E/\Q}(\det y)^m\) for any \(m\in\Z\).
\end{lemma}
\begin{proof}
Use the algebraic representations on the flag tensors in
Theorem~\ref{thm:compact-highest-vector}, with integral
determinant powers as in Definition~\ref{def:bounded-degree-free-weights}.
Each displayed matrix operation gives a group homomorphism.
For \(y=\alpha(t)=\bar t\) with \(t\in U(F)\), both arguments
are \(\overline{\tau(t)}\), since \(\tau(t)^{-\mathsf T}
=\overline{\tau(t)}\). Also
\(\det(t)\overline{\det(t)}=1\), giving rational norm one
by pairing the embeddings of Lemma~\ref{lem:number-field-lattice}.
Thus the tensor acts by exactly the prescribed covariance of
Definition~\ref{def:definite-comparison-forms}. A scalar \(zI\)
acts in \(W_{\lambda^\tau}\) by \(z^{\sum_j\lambda^\tau_j}\):
each column of height \(h\) in the exterior tensor contributes
\(z^h\), and the last determinant power contributes its
exponent times \(n\) to the total exponent. Scalar rational
matrices therefore act by scalars in every indicated choice.
\end{proof}

\begin{theorem}\label{thm:actual-frame-extension}
Let \(f:U\to W\) be a nonzero form with the rational covariance
of Lemma~\ref{lem:rational-weight-extension}. Let \(\Xi\) be
any representation giving that covariance. Suppose
\(R(u)f=f\) whenever \(u\in U\) and \(\alpha(u)\in C\).
There is a definite function \(\phi\) of
Definition~\ref{def:definite-vector-extension} such that
\[
       \phi(y\alpha(u)k)=\Xi(y)f(u)
          \qquad(y\in\Gamma,\ u\in U,\ k\in C),
\]
with value zero on the complement of \(\Gamma\alpha(U)C\).
In particular it is nonzero.
\end{theorem}
\begin{proof}
Suppose \(y_1\alpha(u_1)k_1=y_2\alpha(u_2)k_2\). For
\(y=y_2^{-1}y_1\) this gives
\(y=\alpha(u_2)k_2 k_1^{-1}\alpha(u_1)^{-1}\).
Since both outside factors are unitary, \(y^*y=tI\) at
the finite adeles for some \(t\in\widehat{\Z}^{\,\times}\).
As in Proposition~\ref{prop:neat-at-wild-place}, a diagonal
entry gives \(t\in\Q^\times\cap\widehat{\Z}^{\,\times}\),
and positivity of \(y^*y\) at any complex embedding gives
\(t=1\). Thus \(y\) is a rational isometry and \(k_2k_1^{-1}\)
is unitary. Applying \(\alpha\) gives
\(\alpha(y)u_1=u_2\alpha(k_2k_1^{-1})\).
The hypothesis makes the last factor fix \(f\). Rational
covariance yields \(f(u_2)=\Xi(y)f(u_1)\), proving equality
of the two proposed values. The support condition and the
formula give left covariance and right invariance since
\(\Gamma\) and \(C\) are groups. On a coset represented by
\(g=y\alpha(u)k\), use the rational matrix
\(A=(y^{-1})^*y^{-1}\). It is positive definite at every
complex embedding, and \(g^*Ag=k^*k\) is a profinite scalar
unit. This gives the signed frame at \(J=I\). Finally
\(\phi(\alpha(u))=f(u)\).
\end{proof}

\begin{proposition}\label{prop:local-translation-frame-extension}
At a split place \(w\ne\bar w\) put
\(u_w(a)\in U\) equal to \(a\in{\rm GL}_n(E_w)\) at \(w\),
the conjugate inverse transpose at its mate and the identity
elsewhere. Put \(b_w(a)=\alpha(u_w(a))\in G\). If
\(\mathfrak n\notin\bar w\), then for each \(g\in G\)
there exists a complex linear map \(L_g\) from functions
\(U\to W\) to \(W\) such that
\[
           \phi(g b_w(a))=L_g(R(u_w(a))f)
                      \qquad(a\in{\rm GL}_n(E_w)).
\]
Consequently every finite complex linear relation among the
local translates of \(f\) is a relation among the corresponding
translates of \(\phi\).
\end{proposition}
\begin{proof}
The finite scalar similitudes normalize \(U=\alpha(U)\), by the
matrix identity \(k^*k=tI\). Thus \(\alpha(U)C\) is a subgroup,
and \(\Gamma\alpha(U)C\) and its complement are stable under
right multiplication by \(b_w(a)\). Outside this set take
\(L_g=0\). Inside it choose \(g=y\alpha(u)k\) as in
Theorem~\ref{thm:actual-frame-extension}. Write
\(d=k_w^{-\mathsf T}\). One has
\(k b_w(a)k^{-1}=b_w(dad^{-1})\). At \(w\) this follows by
applying inverse transpose, and at \(\bar w\) it follows from
the unitary relation for the conjugated element; all other
components are the identity. Also \(b_w(d)\in C\): its
component at \(w\) equals \(k_w\), its component at the mate
is integral and invertible with no congruence condition
there, and it is unitary. Hence \(R(u_w(d)^{-1})f=f\).
The extension formula now gives
\[
  \phi(g b_w(a))
   =\Xi(y) f(u\,u_w(dad^{-1}))
   =\Xi(y)(R(u_w(a))f)(u\,u_w(d)).
\]
Evaluation at the last point followed by \(\Xi(y)\) is linear.
Applying it at each \(g\) to a linear relation proves the
last assertion.
\end{proof}

\section{The central character of the definite test vector}

\begin{proposition}\label{prop:central-character-actual-form}
Write \(\mathbb A_{K,f}^{\times}\) for the finite ideles of
the imaginary quadratic field \(K\). Define \(s(t)=\bar t I\)
in \(G\) via the inclusion of its finite adeles in those of
\(E\). Let \(\phi\) be as in Theorem~\ref{thm:actual-frame-extension},
with \(\Xi(\bar z I)=\eta(z)I\) for \(z\in K^\times\).
The span of the functions \(R(s(t))\phi\) is finite dimensional
and consists of definite functions at level \(\mathfrak n\).
It contains a nonzero \(\psi\) and a character
\(\omega:\mathbb A_{K,f}^\times\to\C^\times\) with open kernel
such that
\[
  R(s(t))\psi=\omega(t)\psi,\qquad
                 \omega(z)=\eta(z)\quad(z\in K^\times).
\]
For the choices of \(\Xi\) in Lemma~\ref{lem:rational-weight-extension},
\(\eta(z)\) is a product of integral powers of the two complex
embeddings of \(z\). All local linear relations of
Proposition~\ref{prop:local-translation-frame-extension}
hold for \(\psi\).
\end{proposition}
\begin{proof}
Choose a positive rational integer \(m\) with \(m\in\mathfrak n O_E\),
for example \(|N_{E/\Q}\mathfrak n|\) by the adjugate argument
of Theorem~\ref{thm:actual-pel-torsion-pairing}. Let \(H_m\) be
the units of \(\prod_u O_{K_u}\) congruent to one modulo \(m\)
in this product. Their scalar images lie in \(C=K(\mathfrak n)\):
the finite multiplier \(t\bar t\) is a rational profinite unit
and \(\bar t\equiv1\bmod m\). Thus \(H_m\) fixes \(\phi\), and
the rational \(z\) act by \(\eta(z)\) by left covariance and
centrality.

The quotient \(\mathbb A_{K,f}^{\times}/(K^\times H_m)\)
is finite. Taking the valuations at all primes maps finite
ideles onto fractional ideals, by choosing local uniformizers
at the finite support, using the factorization of
Proposition~\ref{prop:general-ideal-factorization}. The kernel
is the product of local units. The ordinary class group is
finite by Lemma~\ref{lem:bounded-ideal-classes}. The product
of units modulo \(H_m\) is also finite: only the factors
over \(m\) occur, and each has finitely many residues modulo
\(m\) by Lemma~\ref{lem:local-residue-indicators}. Hence
finitely many central translates span the orbit, since a
change by \(K^\times H_m\) acts by a scalar.

These translates are right \(C\)-invariant and left
\(\Xi\)-equivariant by centrality. If \(\phi(g s(t))\ne0\),
let \(A\) be a positive definite rational matrix giving its
frame at \(J=I\). Then \(g^*Ag=\varepsilon/(t\bar t)\) times
the identity, with \(\varepsilon\) a profinite unit. Choose
\(a=\prod_p p^{v_p(t\bar t)}\in\Q_{>0}\), a finite product
since \(t\) is an idele. The matrix \(aA\) stays positive and
gives a profinite unit multiplier at \(g\). This proves the
support requirement of Definition~\ref{def:definite-vector-extension}.

A commuting family of invertible endomorphisms on a nonzero
finite-dimensional complex space has a common eigenvector.
Indeed, if some endomorphism is nonscalar choose a nonzero
proper eigenspace of it, using the complex root property of
Lemma~\ref{lem:place-fields}; all other operators preserve
that eigenspace. Iterate in smaller dimension. The scalar
case is immediate. Applied to the central orbit this gives
\(\psi\); the eigenvalues form a character since the operators
are invertible and compose multiplicatively. Its kernel
contains \(H_m\), open in the restricted product. The scalar
calculation of Lemma~\ref{lem:rational-weight-extension}
gives the asserted description of \(\eta\), including the
rational norm factor by the product of embeddings. Finally
each central translate commutes with the local translations,
so all the homogeneous linear relations persist under their
finite linear combinations.
\end{proof}

\begin{theorem}\label{thm:actual-central-test-vector}
Take \(f,\mathfrak n,w_0,w_1,w_*\) from
Corollary~\ref{cor:actual-moduli-places}, transporting the
form back to \(\C\) and using a rational weight extension
of Lemma~\ref{lem:rational-weight-extension}. The corresponding
\(\psi\) of Proposition~\ref{prop:central-character-actual-form}
obeys the cuspidal character relation at \(w_0\) through
\(b_{w_0}\) and is fixed by \(b_{w_*}({\rm GL}_n(O_{w_*}))\).
The inverse image \(b_{w_0}^{-1}(C)\) is contained in
\(I+\ell M_n(\Z_\ell)\); the inverse image \(b_{w_*}^{-1}(C)\)
is contained in the subgroup whose reduction is upper triangular
with every diagonal entry one.

Let \(w\) be safe in the sense of
Theorem~\ref{thm:actual-pel-safe-level}, with contraction \(u\)
to \(K\). Put \(z_w\in\mathbb A_{K,f}^\times\) equal to
\(N_{E_w/K_u}(\pi_w)\) at \(u\) and to one elsewhere, for
a uniformizer \(\pi_w\). Set \(c_w=\omega(z_w)\). The
approximants may be chosen so that at every such \(w\) the
following finite sum acts on \(\psi\) through the indicated
eigenvalue:
\[
  \sum_{a} R\big(s(z_w^{-h})\,b_w(a)\big)\psi
              =c_w^{-h}t_{w,h}\psi,\qquad 0\le h\le n .
\]
Here \(a\) runs over the right-factor cosets \(aK_w\) of
Definition~\ref{def:actual-spherical-polynomial},
\(K_w={\rm GL}_n(O_w)\). The scalar \(c_w\) is independent of
the choice of uniformizer.
\end{theorem}
\begin{proof}
Choose simultaneous spherical eigenvectors at every split
factor where the level character is trivial on its full
integral group by Corollary~\ref{cor:definite-approximating-eigenforms}.
Mark in the support choice of
Corollary~\ref{cor:actual-moduli-places} one member of every
other exceptional prime pair. If \(w\) is safe, neither \(w\)
nor its conjugate is marked: the former avoids the level
and the latter lies over \(\bar u\), where the level is a unit.
The pair is split since its contractions to \(K\) differ.
Thus \(f\) has all the indicated spherical eigenvalues at \(w\).
Theorem~\ref{thm:actual-frame-extension} applies to \(f\)
by the level property of Corollary~\ref{cor:actual-moduli-places}.
That corollary excludes \(\bar w_0,\bar w_*\) from the level.
Applying Proposition~\ref{prop:local-translation-frame-extension}
at \(w_0,w_*\), and then
Proposition~\ref{prop:central-character-actual-form}, transfers
both local relations. At \(w_0\) the character satisfies the
idempotent test of Proposition~\ref{prop:actual-wild-split-character}.
At each of \(w_0,w_*\) the level exponent is positive, so
\(a^{-\mathsf T}\) at that place in \(b_w(a)\in C\) reduces
to \(I\). Taking the inverse transpose over the residue field
proves both assertions about inverse images.

At each safe \(w\), the conjugate place also avoids the level.
The same two propositions transfer the \(K_w\)-invariance
and the relation for the sum of \(R(u_w(a))\) on \(f\).
Right multiplying its terms by the same central scalar
multiplies the eigenvalue by \(\omega(z_w)^{-h}\).
This proves the displayed formula and its independence of
the right coset representatives. A unit idele supported
at \(u\) has \(s\)-image integral, with unit multiplier, and
supported above \(\bar u\), where no level congruence is
imposed. It fixes \(\psi\), so \(\omega\) is trivial on these
units. Replacing \(\pi_w\) by a unit multiple changes its
local norm by a unit: in the product of local conjugates
each has absolute value one by the invariance proved in
Lemma~\ref{lem:marked-unramified-tower}. This proves the
independence of \(c_w\).
\end{proof}

\section{The coefficient representation and a marked eigenvalue}

\begin{lemma}\label{lem:alternant-coefficient-degree}
For strictly increasing nonnegative integers
\(l_1<\cdots<l_n\), put
\[
  S_l(z)=\frac{\det[z_j^{l_i}]_{i,j=1}^n}
                    {\det[z_j^{i-1}]_{i,j=1}^n}.
\]
It is a symmetric polynomial with integer coefficients. As a
polynomial in the elementary symmetric functions \(e_1,\ldots,e_n\)
of the roots \(z_1,\ldots,z_n\), it has a representative of
ordinary total degree at most \(l_n-n+1\).
For distinct nonzero complex numbers \(z_j\) put
\(L_i(Y)=\prod_{k\ne i}(Y-z_k)/(z_i-z_k)\). Then
\[
  \sum_{i=1}^n L_i(Xz_i)
     S_l(z_1,\ldots,z_{i-1},Xz_i,z_{i+1},\ldots,z_n)
                  =S_l(z)\sum_{j=1}^n X^{l_j}.
\]
\end{lemma}
\begin{proof}
The denominator is \(\prod_{a<b}(z_b-z_a)\): each collision
forces a zero, its total degree is \(n(n-1)/2\), and the
coefficient of \(\prod_j z_j^{j-1}\) is one. Each linear
factor divides the numerator by monic division and
specialization at the collision. It is prime because quotienting
identifies two variables in an integral polynomial ring, and
the factors for different pairs do not divide one another.
Thus their product divides the numerator over \(\Z\). Both
determinants change by the same sign under permutations.
The quotient has degree at most \(l_n-(n-1)\) in each
individual variable, by additivity of degrees in a product.

For a homogeneous symmetric polynomial with this bound \(D\),
take its lexicographically largest monomial. Its exponents
\(b_1,\ldots,b_n\) must be nonincreasing by symmetry.
With \(b_{n+1}=0\), the product
\(\prod_{h=1}^n e_h^{b_h-b_{h+1}}\) has that leading monomial
with coefficient one, since \(z_1\cdots z_h\) is the leading
monomial of \(e_h\). It uses \(b_1\le D\) factors, and its
other monomials also have degree at most \(D\) in each variable.
Subtracting the leading coefficient times this product and
iterating terminates, since there are finitely many monomials
of the given homogeneous degree. This proves the coefficient
bound for \(S_l\), which is homogeneous by its determinant
formula.

Upon replacing \(z_i\) by \(Xz_i\), the ratio of the new
denominator to the old is \(L_i(Xz_i)\). After cancellation
the \(i\)-th summand has just the \(i\)-th column of the
numerator modified, over the original denominator. In each
permutation term the successive column modifications use
every exponent \(l_j\) once. Summing proves the identity in
\(\C(X)\), hence in \(\C[X]\) by the polynomial assertion.
\end{proof}

\begin{definition}\label{def:geometric-coefficient-system}
Use the original CM representatives \(\tau_v\) of the definite
weight at the real places \(v\) of \(F\). Fix a distinguished
embedding \(\sigma:E\to\C\), set \(s_0=\sigma|_K\), and let
\(\nu_v\) denote the embedding in the \(v\)-pair which restricts
to \(s_0\). Write \(v_0\) for the place of \(\sigma\). Let
\(k_{v,1}<\cdots<k_{v,n}\) be the oriented jumps of
Definition~\ref{def:oriented-definite-jumps} at \(\nu_v\),
listed increasingly. The inequalities are strict by dominance
of \(\lambda^v\). Put
\[
 m_v=\max(0,-k_{v,1}),\qquad
 l_{v,j}=m_v+k_{v,j},\qquad \mu=n\sum_v m_v .
\]
In the \(v\)-factor of Lemma~\ref{lem:rational-weight-extension}
take \(H_v(y)=\nu_v(y)\) if \(\nu_v=\tau_v\) and
\(H_v(y)=\nu_v(y)^{-\mathsf T}\) otherwise. Put
\[
 \Xi(y)=\left(\prod_v
     \big(\nu_v(\det y)\bar\nu_v(\det y)\big)^{-m_v}\right)
             \bigotimes_v W_{\lambda^v}(H_v(y)).
\]
Define \(A_{\nu_v}(z)=S_{l_v}(z)\) and
\(A_{\bar\nu_v}(z)=(\prod_i z_i)^{m_v}\).
\end{definition}

\begin{theorem}\label{thm:actual-jump-character}
The representation \(\Xi\) of Definition~\ref{def:geometric-coefficient-system}
has the rational unitary covariance of
Lemma~\ref{lem:rational-weight-extension}. On the scalars
\(\bar a I\), \(a\in K^\times\), its character of scalar
action is
\[
  \eta(a)=s_0(a)^{-\mu}\bar s_0(a)^{A-\mu},\qquad
  A=\sum_v \epsilon_v\sum_{j=1}^n\lambda^v_j ,
\]
where \(\epsilon_v=1\) at \(\nu_v=\tau_v\) and \(-1\)
otherwise. Let \(y\in{\rm GL}_n(E)\) have characteristic
polynomial with distinct roots. Write \(z_{\tau}\) for the
list of its roots under any embedding \(\tau\), with conjugate
lists numbered by complex conjugation. Then
\[
              {\rm tr}\,\Xi(y^{-1})=\prod_{\tau}A_\tau(z_\tau).
\]
\end{theorem}
\begin{proof}
The scalar factor defines a character on \(\Gamma\). On
unitary rational matrices each pair of determinant values
has product one, so the proof of
Lemma~\ref{lem:rational-weight-extension} applies separately
at each \(v\). At \(\bar a I\) that pair contributes
\((s_0(a)\bar s_0(a))^{-nm_v}\), and the \(v\)-representation
contributes \(\bar s_0(a)^{\epsilon_v\sum_j\lambda^v_j}\)
by the scalar calculation there. Multiplication proves
the central formula.

At \(\nu_v=\tau_v\) the tensor factor on \(y^{-1}\) has
eigenvalue arguments \(z_i^{-1}\) for \(z=z_{\nu_v}\).
Theorem~\ref{thm:compact-alternant} applies at these distinct
nonzero numbers: its determinant identity on diagonal matrices
is a Laurent-polynomial identity by the tableau proof there,
and the matrix is diagonalizable since its roots are distinct.
In both numerator and denominator of that formula multiply
each column \(j\) by \(z_j^{n-1}\). The new exponent list
in the numerator is \(i-1-\lambda^v_i\), and the denominator
list is \(i-1\) in increasing order. At \(\nu_v\ne\tau_v\)
the tensor argument on \(y^{-1}\) has eigenvalues \(z_i\).
Reverse the rows of both alternants. The denominator again has
list \(i-1\) and the numerator list is the increasing ordering
of \(\lambda^v_i+n-i\). These are precisely the two oriented
jump lists. Multiplication by \((\prod_j z_j)^{m_v}\) adds
\(m_v\) to the numerator exponents. The conjugate factor of
the determinant twist is \((\prod_j z_{\bar\nu_v,j})^{m_v}\).
Taking the product over all pairs gives the trace identity,
since the trace of a tensor product is the product of its
traces in product bases.
\end{proof}

\begin{theorem}\label{thm:actual-marked-jump-polynomial}
In Theorem~\ref{thm:actual-jump-character}, for \(1\le i\le n\)
form lists \(z^{(i)}_\tau(X)\) as follows. At
\(\nu_{v_0}=\sigma\) multiply \(z_{\sigma,i}\) by \(X\),
at \(\bar\sigma\) multiply every root except
\(\overline{z_{\sigma,i}}\) by \(X\). For \(v\ne v_0\)
multiply all roots at \(\bar\nu_v\) by \(X\) and keep the
list at \(\nu_v\). Take \(L_i\) of
Lemma~\ref{lem:alternant-coefficient-degree} on \(z_\sigma\).
Then
\[
 \sum_i L_i(X z_{\sigma,i})\prod_\tau A_\tau(z^{(i)}_\tau(X))
          = {\rm tr}\,\Xi(y^{-1})
                     \sum_{j=1}^n X^{\mu+k_{v_0,j}} .
\]
In particular each indicated coefficient is exactly the
trace. The exponents \(\mu+k_{v_0,j}\) are nonnegative and
pairwise distinct.
\end{theorem}
\begin{proof}
Each of the non-distinguished conjugate determinant factors
contributes \(X^{nm_v}\) to the modification. At the conjugate
of the distinguished embedding the factor contributes
\(X^{(n-1)m_{v_0}}\). These are independent of \(i\). All
remaining changes belong to the distinguished alternant.
Lemma~\ref{lem:alternant-coefficient-degree} gives its original
value times \(\sum_j X^{l_{v_0,j}}\). The total exponent outside
this sum is \(\mu-m_{v_0}\). Substitute
\(l_{v_0,j}=m_{v_0}+k_{v_0,j}\) and apply
Theorem~\ref{thm:actual-jump-character}. The exponents are
distinct by Definition~\ref{def:geometric-coefficient-system}
and nonnegative since \(\mu\ge m_{v_0}\).
\end{proof}

\section{A finite family of local terms for the jump coefficients}

\begin{lemma}\label{lem:finite-grid-root-coefficients}
Let \(E\) be the field of Definition~\ref{def:levels}, \(d_E=[E:\Q]\),
and let \(b_1,\ldots,b_{d_E}\) be an integral basis. Let
\(n,\nu\ge1\) and choose, for each \(\tau:E\to\C\), a polynomial
\(Q_\tau\) in \(n\) variables of ordinary total degree at most
\(\nu\). Define
\[
  \mathcal B_n=\left\{\sum_{t=1}^{d_E}h_tb_t:
                    h_t\in\{0,\ldots,n d_E\}\right\}.
\]
There are complex constants \(d_\alpha\), indexed by
\(\alpha\in\mathcal B_n^\nu\), with the following simultaneous
identity for all lists \(x_{\tau,1},\ldots,x_{\tau,n}\)
in any commutative complex algebra:
\[
 \prod_\tau Q_\tau\big(e_1(x_\tau),\ldots,e_n(x_\tau)\big)
   =\sum_{\alpha\in\mathcal B_n^\nu} d_\alpha
       \prod_{a=1}^\nu\prod_{\tau}\prod_{j=1}^n
                   (1-2\tau(\alpha_a)x_{\tau,j}).
\]
The constants depend only on the polynomials, the basis, \(n\)
and \(\nu\).
\end{lemma}
\begin{proof}
First take the \(x_{\tau,j}\) as independent indeterminates.
In the product on the right use independent variables
\(Y_{a,\tau}\) in place of \(2\tau(\alpha_a)\). The coefficient
of \(\prod_{a,\tau}Y_{a,\tau}^{s_{a,\tau}}\), for
\(0\le s_{a,\tau}\le n\), is
\(\prod_{a,\tau}(-1)^{s_{a,\tau}}e_{s_{a,\tau}}(x_\tau)\)
with \(e_0=1\), by finite multiplication. For each monomial
in \(Q_\tau\), assign each of its at most \(\nu\) elementary
symmetric factors to a separate \(a\), filling unused
positions by \(s=0\). The assignments for distinct \(\tau\)
are independent. Expanding the product of the \(Q_\tau\)
shows that the left side is a fixed complex linear combination
of coefficients of this polynomial in the \(Y\)'s.

For each \(a\) substitute
\(Y_{a,\tau}=2\sum_t \tau(b_t)H_{a,t}\).
The matrix \((\tau(b_t))_{\tau,t}\) is nonsingular over \(\C\):
a power basis of a primitive element has a nonsingular
Vandermonde matrix by Lemma~\ref{lem:number-field-lattice},
and the change to the given rational basis is invertible.
The substitution thus has a linear inverse, so each original
coefficient is a complex linear combination of coefficients
in the \(H\)'s. For each \(a\) the total degree in its block
is at most \(n d_E\), as each of its \(d_E\) former variables
had degree at most \(n\). Every coefficient in the \(H\)'s
is recovered linearly from the grid \(0,\ldots,n d_E\) in
each variable. Indeed in one variable the interpolation basis
at \(j\) is \(\prod_{0\le t\le nd_E,\ t\ne j}(X-t)/(j-t)\);
subtracting the interpolant leaves a polynomial of bounded
degree vanishing at more points than its degree. Apply this
successively to each variable. The resulting identity of
complex polynomials in the \(x_{\tau,j}\) holds upon every
indicated substitution into a complex algebra.
\end{proof}

\begin{definition}\label{def:actual-finite-local-terms}
Use the data of Definition~\ref{def:geometric-coefficient-system}
and take
\[
 \nu=\max\left(1,\ \max_v m_v,\ \max_v(l_{v,n}-n+1)\right).
\]
Let \(y\in{\rm GL}_n(E)\) have distinct eigenvalues at each
embedding. For its \(i\)-th marked root at \(\sigma\), form the
lists \(z_{\tau}^{(i)}(X)\) of
Theorem~\ref{thm:actual-marked-jump-polynomial}. For
\(\alpha\in O_E^\nu\) define the polynomial
\[
 B_{\alpha,y,i}(X)=L_i(Xz_{\sigma,i})
    \prod_{a=1}^\nu\prod_\tau\prod_{j=1}^n
          \big(1-2\tau(\alpha_a)z^{(i)}_{\tau,j}(X)\big).
\]
The non-marked lists use the CM representatives \(\Psi\)
of Lemma~\ref{lem:actual-signs-and-padding}.
\end{definition}

\begin{theorem}\label{thm:finite-geometric-local-terms}
There are constants \(d_\alpha\) for the finite grid in
Lemma~\ref{lem:finite-grid-root-coefficients}, the same for
every regular \(y\) and every jump \(k_{v_0,j}\), such that
\[
 {\rm tr}\,\Xi(y^{-1})=
   \sum_{\alpha\in\mathcal B_n^\nu} d_\alpha
       \sum_{i=1}^n [X^{\,\mu+k_{v_0,j}}]B_{\alpha,y,i}(X).
\]
Here regular means that the characteristic polynomial of
\(y\) is separable. In particular, for any finite complex
weighted list \((c_t,y_t)\) of such matrices, a nonzero sum
\(\sum_t c_t\,{\rm tr}\,\Xi(y_t^{-1})\) forces for each
indicated jump a nonzero sum of the corresponding marked
coefficients for at least one \(\alpha\) in the grid.
\end{theorem}
\begin{proof}
At \(\nu_v\), use Lemma~\ref{lem:alternant-coefficient-degree}
to express \(A_{\nu_v}\) in the elementary symmetric functions
by a polynomial of ordinary total degree at most \(l_{v,n}-n+1\).
At \(\bar\nu_v\), take \(Q_{\bar\nu_v}=e_n^{\,m_v}\), meaning
the \(m_v\)-th power of the last elementary symmetric variable.
All degrees are at most \(\nu\) by
Definition~\ref{def:actual-finite-local-terms}. Choose the
constants of Lemma~\ref{lem:finite-grid-root-coefficients}
for these polynomials. Substitute the modified root lists in
\(\C[X]\), multiply by \(L_i(Xz_{\sigma,i})\), and sum over
\(i\). By Theorem~\ref{thm:actual-marked-jump-polynomial}
the coefficient on the left at each stated degree is exactly
\({\rm tr}\,\Xi(y^{-1})\). On the right the operations produce
the stated finite sum. Summing this equality against \(c_t\)
and using that a sum of zeros is zero proves the final assertion.
\end{proof}

\section{An ordinary Iwahori operator with regular support}

\begin{definition}\label{def:ordinary-iwahori-average}
Let \(L\) be a finite extension of \(\Q_p\), \(O\) its
integers, \(\pi\) a uniformizer, \(k_0=O/\pi O\) its residue
field of size \(q\), and \(v(\pi)=1\), as in
Lemma~\ref{lem:local-residue-indicators}. Let \(I\subset
{\rm GL}_n(O)\) be the inverse image of the upper triangular
subgroup over \(k_0\). For a complex smooth representation
\(\rho\) of \({\rm GL}_n(L)\), meaning each vector has an open
stabilizer, denote by \(P_H\) the probability average over a
compact open subgroup \(H\). This means the finite average of
the orbit of a vector, with equal weight on the cosets of its
stabilizer in \(H\). On the space \(V^I\) put
\(U_g=P_I\rho(g)|_{V^I}\).
\end{definition}

\begin{lemma}\label{lem:iwahori-finite-permutations}
If \(IgI\,IhI=IghI\) then \(U_g U_h=U_{gh}\).
For every permutation matrix \(w\), the operator \(U_w\) is
invertible on \(V^I\).
\end{lemma}
\begin{proof}
For a compact subgroup, finitely many cosets occur in an open
quotient by compactness, and passage to a smaller open subgroup
refines every coset equally. The average in
Definition~\ref{def:ordinary-iwahori-average} is therefore
linear and a projection onto invariants. The operator \(U_g\)
averages the right-factor cosets \(aI\) in \(IgI\) uniformly.
On composing two such averages one obtains positive weights
of total one on the right-factor cosets in \(IgIhI\), invariant
under left \(I\). Indeed left multiplication permutes the
cosets in the first average and the resulting changes of
representative act on the second average through left \(I\).
If the product is one double coset, left \(I\) is transitive
on these right-factor cosets. Thus the composite has exactly
the uniform weights of \(U_{gh}\).

Write \(s_i\) for the adjacent swap of positions \(i,i+1\)
and \(P_i\) for the subgroup of \({\rm GL}_n(O)\) preserving
the standard flag modulo \(\pi\) with just its dimension-\(i\)
subspace omitted. The quotient \(P_i/I\) parametrizes lines
in the two-dimensional residue quotient between its
neighboring flag members. Every such line is realized by
an invertible two-by-two integral block lifting a residue
matrix. One is the original flag, and the other \(q\) are
represented by \((1+\tilde a E_{i,i+1})s_i I\), \(a\in k_0\),
where \(\tilde a\) is any integral lift. Thus
\(P_i=I\sqcup Is_iI\) and on \(V^I\) one has
\(P_{P_i}=(1+qU_{s_i})/(q+1)\). Since \(P_{P_i}\) is a
projection, \(U_{s_i}\) acts by \(1\) on its image and by
\(-1/q\) on its kernel. It is invertible.

If \(w(i)<w(i+1)\), conjugation by \(w\) takes every
\(1+\tilde a E_{i,i+1}\) above into \(I\), whence
\(IwI\,Is_iI=Iws_i I\) by the displayed cosets. A permutation
with a descent loses one inversion on right multiplication by
the corresponding adjacent swap. Induct on the number of
inversions, applying the product formula to the shorter
permutation followed by that swap. This proves invertibility
for every \(w\).
\end{proof}

\begin{theorem}\label{thm:iwahori-dominant-unit}
For integers \(a_1\ge\cdots\ge a_n\) put
\(t_a=\operatorname{diag}(\pi^{a_1},\ldots,\pi^{a_n})\).
One has \(It_aI\,It_bI=It_{a+b}I\) whenever both lists are
nonincreasing. The operator \(U_{t_a}\) of
Definition~\ref{def:ordinary-iwahori-average} is invertible
on \(V^I\), and \(U_{t_a}^m=U_{t_{ma}}\) for every \(m\ge1\).
\end{theorem}
\begin{proof}
Every \(k\in I\) factors as \(u_+d u_-\) with \(u_+\) upper
unitriangular integral, \(d\) diagonal with unit entries,
and \(u_-\) lower unitriangular with entries below the diagonal
in \(\pi O\). Eliminate the last column above its diagonal
using its unit last pivot by row operations, then the last
row to the left using column operations with coefficients
in \(\pi O\). The remaining block stays upper triangular with
nonzero diagonal modulo \(\pi\), so induction gives the
factorization. One has \(t_a u_+d t_a^{-1}\in I\) and
\(t_b^{-1}u_-t_b\in I\) by comparing the differences of the
exponents. Hence \(t_a k t_b\in It_{a+b}I\), proving equality
of the two sets. The power formula follows from
Lemma~\ref{lem:iwahori-finite-permutations}.

Define \(\Omega e_1=\pi e_n\), \(\Omega e_i=e_{i-1}\) for
\(i>1\), on the standard column basis. This matrix normalizes
\(I\). Indeed it carries the lattice with its last \(h\)
coordinates in \(\pi O\) and others in \(O\) to the analogous
lattice with \(h+1\) such coordinates for \(0\le h<n\).
For \(h=n\) continue periodically by scaling all lattices
by \(\pi\). The group \(I\) consists exactly of simultaneous
stabilizers of this periodic chain, by reduction of the
intermediate lattices modulo \(\pi\). Thus \(U_{\Omega^j}\)
is the invertible translation by \(\Omega^j\) on \(V^I\).
For \(1\le j<n\) a permutation \(w_j\) undoes the cyclic
permutation in \(\Omega^j\), giving
\(w_j\Omega^j=\operatorname{diag}(\pi I_j,I_{n-j})\).
Since \(\Omega^j\) normalizes \(I\), its product double coset
with \(Iw_jI\) is single. Both factors act invertibly by the
preceding lemma. Finally \(t_a\) is the scalar \(\pi^{a_n}I\)
times these diagonal matrices to powers \(a_j-a_{j+1}\).
The products have nonincreasing exponent lists throughout.
The product formula proves invertibility, also for \(n=1\).
\end{proof}

\begin{proposition}\label{prop:iwahori-simple-local-roots}
If \(a_1>\cdots>a_n\), every element of \(It_aI\) has \(n\)
distinct eigenvalues in \(L\) with valuations
\(a_1,\ldots,a_n\). This holds on \((It_aI)^m\) with these
valuations multiplied by \(m\).
\end{proposition}
\begin{proof}
A product \(k_1t_a k_2\) is conjugate to \(t_a k_2k_1\).
For \(k\in I\), every principal minor of \(k\) reduces to
the product of the corresponding nonzero diagonal entries,
hence is a unit. In \(\det(YI-t_a k)=\sum_{j=0}^n c_jY^{n-j}\),
expanding the determinant gives
\(c_j=(-1)^j\sum_{|S|=j}\pi^{\sum_{i\in S}a_i}\det k_{S,S}\).
There is a unique summand of least valuation, indexed by
the last \(j\) positions. Thus \(v(c_0)=0\) and
\(v(c_j)-v(c_{j-1})=a_{n-j+1}\), strictly increasing in \(j\).

For one of these successive differences \(b\), substitute
\(Y=\pi^b Z\). The valuations of the terms' coefficients
now decrease strictly to two adjacent equal minima and then
increase strictly, since the new successive differences
are \(a_{n-j+1}-b\). Divide by the power of \(\pi\)
of the minimum. The reduction is \(Z^k(A+BZ)\) with both
\(A,B\) nonzero, so its nonzero residue root is simple.
Lift it to a unit root over \(O\): at stage \(l\ge1\), if
a polynomial value lies in \(\pi^l O\) and its derivative is
a unit, adding \(\pi^l z\) corrects the next digit by choosing
\(z\bmod\pi\) from the linear term. Iterating converges in
the complete field by Theorem~\ref{thm:trace-transport-from-primes}.
This gives an eigenvalue of valuation \(b\). All \(n\) values
are distinct and exhaust the roots by degree. The last
assertion follows from Theorem~\ref{thm:iwahori-dominant-unit}.
\end{proof}

\section{Regular tests for the definite eigensystem}

\begin{definition}\label{def:actual-positive-hecke}
For \(G,\Gamma,C\) of Definition~\ref{def:definite-vector-extension}
put
\[
  S=\{g\in G:g^*g=tI,\ t\in\mathbb A_{\Q,f}^\times\},
  \qquad S^+=\{g\in S:g^{-1}\in M_n(\widehat O)\}.
\]
The multiplier is unique since \(n\ge1\); membership in \(S^+\)
is preserved by multiplication. The group \(C\) is compact open
in \(S\): it is the intersection with the integral congruence
subgroup, and on this intersection the multiplier is a unit at
every prime by invertibility. Compactness follows from the closed
matrix equalities for a scalar in \(\widehat{\Z}^{\,\times}\)
inside the compact integral matrix group, using the finite residue
quotients of Lemma~\ref{lem:local-residue-indicators}.
Let \(\mathcal H^+\) be the complex functions on \(S/C\) of finite
support in \(S^+/C\) invariant under left \(C\). On the \(C\)-fixed
vectors of a representation write their action as
\(\sum_{aC} H(aC)R(a)\). The product is composition of these
finite sums, grouping identical right-factor cosets. This is
well defined because changing \(a\) to \(ak\), \(k\in C\),
permutes the second cosets preserving their coefficients.
The product is again left \(C\)-invariant, and multiplication
is associative by finite summation. We extend \(H\) by zero
to \(G\) off the indicated cosets. Call \(H\) regular if
\(H(x^{-1}y^{-1}x)\ne0\), \(x\in G,\ y\in\Gamma\), implies
that \(y\) has separable characteristic polynomial over \(E\).
\end{definition}

\begin{lemma}\label{lem:safe-operators-away-from-markers}
For the level of Corollary~\ref{cor:actual-moduli-places}, the
safe spherical sums of Theorem~\ref{thm:actual-central-test-vector}
give elements of \(\mathcal H^+\). Each has representatives
\(g\in S^+\) with \(d_{w_0}(g)=d_{w_*}(g)=I\), where
\(d_v(g)=g_v^{-\mathsf T}\). Finite products and linear
combinations of the sums have expansions by representatives
with the same two identities.
\end{lemma}
\begin{proof}
Let \(w\) be safe, \(u=w|K\), and \(g=s(z_w^{-h})b_w(a)\)
as in Theorem~\ref{thm:actual-central-test-vector}. Both
factors are similitudes. The inverse of \(b_w(a)\) is integral
at \(w\), since \(a\) is integral, and at its mate its
denominators divide the conjugate uniformizer when \(h\ge1\).
This follows by taking inverses in the integral double coset
of \(\operatorname{diag}(\pi_w I_h,I_{n-h})\). At \(h=0\)
it is already integral. The inverse scalar \(s(z_w^h)\)
is integral everywhere and divisible at \(\bar w\) by its
uniformizer when \(h\ge1\): the local norm of \(\pi_w\) has
positive valuation, by taking the product of its conjugates
with the invariant local absolute value of
Lemma~\ref{lem:marked-unramified-tower}. Thus \(g\in S^+\).
The safety condition excludes \(w,\bar w\) from the level
and excludes every prime above \(\bar u\) from the level.
Both markers divide the level, hence neither occurs in the
support of this local factor or of its central scalar.

For \(k\in C\) one has \(d_w(k)\in K_w\) and
\(b_w(d_w(k))\in C\), since both places of this pair have
unit level. By the conjugation formula of
Proposition~\ref{prop:local-translation-frame-extension},
left multiplication by \(k\) acts on the right cosets of the
sum by left multiplication on \(aK_w\) by \(d_w(k)\).
The scalar is central. This operation permutes the cosets,
proving the required invariance. The projection from local
cosets to global cosets is well defined since \(b_w(K_w)\subset C\).
In a product one may keep the ungrouped representatives,
multiplying the chosen representatives of each factor. The
two dual component maps are homomorphisms, giving the last
assertion.
\end{proof}

\begin{theorem}\label{thm:actual-regular-hecke-test}
Let \(\psi\) be the definite vector of
Theorem~\ref{thm:actual-central-test-vector}. Write \(J=
1+\ell M_n(\Z_\ell)\) for the group of its character \(\chi\)
at \(w_0\). Let \(c\) be a polynomial in the safe spherical
operators, with constant terms allowed, such that \(c\psi=\psi\).
There is \(b\in\mathcal H^+\) with the following properties
for each \(m\ge1\). The element \(H_m=(bc)^m\) is regular,
\(H_m\psi\ne0\), and its coset function satisfies
\[
 H_m(b_{w_0}(j)g)=\chi(j)^{-1}H_m(g)
                 \qquad(j\in J,\ g\in G).
\]
It is a left multiple of \(c\) in \(\mathcal H^+\).
\end{theorem}
\begin{proof}
In the group at \(w_*\), use \(I\) of
Definition~\ref{def:ordinary-iwahori-average}, \(\pi=p_*\)
and \(t=\operatorname{diag}(\pi^{n-1},\pi^{n-2},\ldots,1)\).
In any smooth \(S\)-representation, let
\[
 E_\chi=\int_J\chi(j)^{-1}R(b_{w_0}(j))\,dj,\qquad
 U=P_I R(b_{w_*}(t))P_I,\qquad z=p_*^{-(n-1)}I .
\]
Here \(P_I\) uses the copy \(b_{w_*}(I)\). Probability
averages of these locally constant functions are the finite
averages of Lemma~\ref{lem:iwahori-finite-permutations};
the character has open kernel by
Proposition~\ref{prop:actual-wild-split-character}. Translation
of the first sum by \(b_{w_0}(j)\) multiplies it by \(\chi(j)\),
and \(E_\chi\) fixes vectors with that character. All factors
at distinct marked pairs commute. For \(k\in C\), its action
by conjugation on either local subgroup is conjugation by
\(d_v(k)\) in local coordinates by
Proposition~\ref{prop:local-translation-frame-extension}.
Here \(d_{w_0}(k)\in 1+\ell^2 M_n(\Z_\ell)\subset\ker\chi\),
and \(d_{w_*}(k)\in I\) by
Corollary~\ref{cor:actual-moduli-places}. Thus \(R(k)\)
commutes with \(E_\chi\) by invariance of the character under
conjugation in \(J\), and with \(U\) by its two \(I\)-averages.
The scalar \(z\) commutes with all the operators.

Act with \(R(z)E_\chi U\) on the basis vector \([C]\) in
the permutation representation with basis \(S/C\) and left
translation. This representation is smooth since the stabilizer
of each basis vector is open and a vector uses finitely many
basis vectors. The result is a finite sum of cosets fixed by
left \(C\), and is the desired \(b\). Indeed all its terms
have representatives in
\(z\,b_{w_0}(J)b_{w_*}(ItI)\). The inverse of the last local
factor has denominators at worst \(p_*^{n-1}\) at the two
places, since \(I\) is integral; the inverse of the scalar
clears them. Elsewhere all inverse factors are integral. Thus
\(b\in\mathcal H^+\). For any \(C\)-fixed vector \(v\), the
assignment \([gC]\mapsto R(g)v\) intertwines the translations
and their finite averages, so \(bv=R(z)E_\chi Uv\). Right
translations of a definite function through \(S\) generate a
smooth space since the function is \(C\)-fixed.

By Lemma~\ref{lem:safe-operators-away-from-markers} the indicated
expansion of \(c\) commutes term by term with the local operators:
a similitude with \(d_v(g)=I\) centralizes \(b_v\) by the
conjugation identity. On the \(C\)-fixed vectors the products
therefore commute with \(U,E_\chi\) and \(R(z)\). Now
\(E_\chi\psi=\psi\) and \(P_I\psi=\psi\) by
Theorem~\ref{thm:actual-central-test-vector}. Consequently
\(H_m\psi=R(z)^m U^m\psi\ne0\), since \(U\) is invertible
on \(I\)-invariants by Theorem~\ref{thm:iwahori-dominant-unit}.
Likewise the leftmost factor \(b\) sends \(C\)-invariants into
the \(\chi\)-space, since \(E_\chi\) commutes with \(U,R(z)\).
Acting on \([C]\) gives exactly the stated coset identity.
It also holds outside \(S\) since \(b_{w_0}(J)\subset S\).

Every coset in \(H_m\) has \(d_{w_*}\) in
\((I(p_*^{n-1}t)I)^m=I(p_*^{n-1}t)^m I\) by
Theorem~\ref{thm:iwahori-dominant-unit}, including the final
right factor from \(C\). Its eigenvalues are distinct by
Proposition~\ref{prop:iwahori-simple-local-roots}. At
\(x^{-1}y^{-1}x\) the dual component is conjugate to the
transpose of \(y_{w_*}\). A repeated root of the rational
characteristic polynomial would persist on embedding its
coefficients in this local field, since a common factor with its
derivative stays a common factor. This proves regularity.
Finally \(H_m=((bc)^{m-1}b)c\).
\end{proof}

\section{A nonzero definite orbital sum for each prescribed weight}

\begin{proposition}\label{prop:definite-full-class-set}
Let \(\mathcal D\) be the set of double cosets in
\(\Gamma\backslash G/C\) containing an \(x\) with a signed
frame at \(J=I\), as in Definition~\ref{def:moduli-signed-frame}.
At the level of Corollary~\ref{cor:actual-moduli-places},
\(\mathcal D\) is finite. Evaluation at one representative per
class identifies the definite functions of
Definition~\ref{def:definite-vector-extension} with
\(\bigoplus_{i\in\mathcal D} W\). The algebra
\(\mathcal H^+\) of Definition~\ref{def:actual-positive-hecke}
preserves this space.
\end{proposition}
\begin{proof}
Write \(A_x\) for the rational positive Hermitian matrix with
\(x^*A_x x=\varepsilon I\), \(\varepsilon\in\widehat{\Z}^{\,\times}\).
It is unique: a second matrix is a scalar multiple of the first
at the finite adeles, with their multiplier ratio a profinite
unit. Taking the quotient of corresponding entries at a nonzero
entry shows that the ratio belongs to \(E\) and is a rational
profinite scalar. In a rational basis of \(E\) starting with
one its other coordinates vanish, so it is a rational number.
Positivity at infinity and the unit condition at every finite
prime force it to be one, as in Proposition~\ref{prop:neat-at-wild-place}.

Associate to \(x\) the \(O_E\)-lattice
\(\Lambda_x=E^n\cap x\widehat O^n\) and the rational symmetric
form \(\langle v,w\rangle_x=\operatorname{Tr}_{E/\Q}(v^*A_xw)\).
By the lattice and completion argument in Lemma~\ref{lem:trace-finite}
this is a full lattice with the prescribed completions. The
form is positive because its squared lengths are the sums over
the embeddings of positive Hermitian squared lengths. On each
completion its Gram matrix in the basis transported by \(x\)
is the standard trace Gram matrix multiplied by the unit
\(\varepsilon_p\). Thus the form is integral and its Gram
determinant is the same positive integer for every \(x\), by
comparison at each prime as in that lemma. Its bounded-basis
argument yields finitely many possible isometry classes of
trace lattices. Multiplication by each \(b\) in the basis of
\(O_E\) has operator norm at most \(\max_\tau|\tau(b)|\),
since the summands for squared lengths are scaled by
\(|\tau(b)|^2\). There are therefore only finitely many sets
of multiplication matrices on the representative lattices, as
proved in the same lemma. An isometry respecting these matrices
extends by rational linearity to an \(E\)-linear map, and
recovers the Hermitian form by varying an argument through
all \(E\)-multiples and using nondegeneracy of the field trace.
Consequently the pairs \((\Lambda_x,A_x)\) have finitely many
classes under \(x\mapsto yx\), \(A_x\mapsto y^{-*}A_xy^{-1}\)
with \(y\in\Gamma\), where \(y^{-*}=(y^*)^{-1}\).

Within any one such class arrange that the lattice and the
matrix agree by a rational change of frame. Two resulting
frames differ on the right by \(k\in{\rm GL}_n(\widehat O)\)
because their local lattices agree. Their Gram identities give
\(k^*k\) equal to a profinite scalar unit. The group of these
integral similitudes has finite index over \(C\): reduction
modulo \(\mathfrak n\widehat O\) has kernel \(C\) and finite
image by Lemma~\ref{lem:local-residue-indicators} and
Proposition~\ref{prop:general-ideal-factorization}.
This proves the finiteness of \(\mathcal D\).
Each \(\Gamma\cap x Cx^{-1}\) is trivial by
Proposition~\ref{prop:neat-at-wild-place} at \(J=I\). Arbitrary
prescribed values on representatives therefore extend uniquely
by the required two equivariances and by zero away from these
classes.

For \(a\in S\) with multiplier \(t\), set
\(b=\prod_p p^{v_p(t_p)}>0\), a positive rational number.
If \(x\) has a positive signed frame \(A_x\), the point \(xa\)
has frame \(b^{-1}A_x\), since \(t/b\) is a profinite unit.
Applying this argument to \(a^{-1}\) proves the converse.
Every coset operator in \(\mathcal H^+\) preserves the support
condition and left covariance, and its full sum preserves right
\(C\)-invariance by Definition~\ref{def:actual-positive-hecke}.
\end{proof}

\begin{theorem}\label{thm:definite-marked-orbit-trace}
For \(H\in\mathcal H^+\) let \(\mathcal O(H)\) denote the
\(\Gamma\)-orbits of pairs \((xC,y)\), where \(x\) has a
positive signed frame, \(y\in\Gamma\) and
\(H(x^{-1}y^{-1}x)\ne0\). The action is
\((xC,y)\mapsto (z x C,zyz^{-1})\). The set is finite, and
on the definite space of Proposition~\ref{prop:definite-full-class-set}
one has
\[
 {\rm tr}\,H=\sum_{[(xC,y)]\in\mathcal O(H)}
             H(x^{-1}y^{-1}x)\,{\rm tr}\,\Xi(y^{-1}).
\]
If \(H\) is regular, all characteristic polynomials in this
sum are separable.
\end{theorem}
\begin{proof}
The summand is invariant: changing \(x\) on the right by \(k\in C\)
conjugates its argument by \(k\), which preserves \(H\) by its
two \(C\)-invariances. Simultaneous rational change of variables
leaves this argument unchanged and conjugates the coefficient
endomorphism. On representatives \(x_i\) from the preceding
proposition the diagonal block of \(H\) at \(i\) includes exactly
the cosets \(aC\) of its finite support with
\(x_i a=\gamma x_i k\), \(\gamma\in\Gamma,\ k\in C\).
Each contributes \(H(a)\Xi(\gamma)\) by covariance. For a
fixed such coset \(\gamma\) is unique by triviality of
\(\Gamma\cap x_i Cx_i^{-1}\). The associated pair is
\((x_i C,\gamma^{-1})\). Every orbit in the statement has a
representative of this form, with \(aC=x_i^{-1}y^{-1}x_i C\).
Two such pairs coincide in the orbit set only if they coincide
on the representatives, again by the trivial stabilizers.
This identifies a finite set and summing traces of the diagonal
blocks gives the formula. Regularity gives separability by
Definition~\ref{def:actual-positive-hecke}.
\end{proof}

\begin{corollary}\label{cor:definite-orbit-sum-detects-jump}
For each \(c\) of Theorem~\ref{thm:actual-regular-hecke-test}
there is a single regular left multiple \(H=(bc)^m\), with
\(1\le m\le |\mathcal D|\dim W\), for which the following holds.
For every jump \(k_{v_0,j}\) there is an
\(\alpha\in\mathcal B_n^\nu\) of Theorem~\ref{thm:finite-geometric-local-terms}
with
\[
  \sum_{[(xC,y)]\in\mathcal O(H)} H(x^{-1}y^{-1}x)
      \sum_{i=1}^n[X^{\,\mu+k_{v_0,j}}]B_{\alpha,y,i}(X)\ne0 .
\]
Moreover \(H(b_{w_0}(j')g)=\chi(j')^{-1}H(g)\) for all \(j'\in J\).
\end{corollary}
\begin{proof}
The finite-dimensional operator \(T=bc\) has nonzero positive
powers by Theorem~\ref{thm:actual-regular-hecke-test} and
Proposition~\ref{prop:definite-full-class-set}. Triangularize
it over \(\C\) by choosing an eigenvector using
Lemma~\ref{lem:place-fields} and inducting on the quotient
dimension. If all diagonal entries were zero, the power of
order \(|\mathcal D|\dim W\) of the strictly upper triangular
matrix would be zero. Let \(\lambda_1,\ldots,\lambda_r\) be
the distinct nonzero diagonal entries with multiplicities
\(h_i\ge1\). If the traces of the first \(r\) powers vanished,
\(\sum_i(h_i\lambda_i)\lambda_i^{k}=0\) for \(0\le k<r\).
The Vandermonde evaluation matrix is invertible: a polynomial
of degree below \(r\) vanishing at the \(r\) distinct nodes
is zero by division by linear factors. Thus every
\(h_i\lambda_i=0\), a contradiction. Choose \(m\le r\) with
nonzero trace and put \(H=T^m\). The preceding theorem gives
a finite weighted sum of traces of regular matrices with
this nonzero value. Theorem~\ref{thm:finite-geometric-local-terms}
gives the asserted \(\alpha\) for each jump. All other properties
of \(H\) follow from Theorem~\ref{thm:actual-regular-hecke-test}.
\end{proof}

\section{A positive trace form for each marked polynomial}

\begin{definition}\label{def:actual-signed-orbits}
For a separable characteristic polynomial \(P\) of degree \(n\)
over \(E\), order the roots \(z_i\) of \(\sigma P\) by their real
parts and then their imaginary parts. Use that ordering for the
index of a root in Definition~\ref{def:actual-finite-local-terms}.
For \(y\in\Gamma\) with polynomial \(P\), put
\[
 J_{0,i}(y)=I,\qquad J_{1,i}(y)=2L_i(\sigma y)-I,\qquad
 L_i(T)=\prod_{j\ne i}(T-z_j)/(z_i-z_j).
\]
The polynomial \(L_i(\sigma y)\) acts by one on the \(i\)-th
eigenline and by zero on the others: eigenvectors of distinct
eigenvalues are independent, by successively applying the
operators \(\sigma y-z_j I\), and here these lines span by
degree. Thus \(J_{1,i}\) has sign \(+1\) on that eigenline
and \(-1\) on its complement. Both matrices commute with
\(\sigma y\) and conjugate equivariantly when \(y\) does.

For a regular \(H\in\mathcal H^+\) and
\(\epsilon\in\{0,1\}\), let \(\mathcal O_\epsilon(H)\) consist
of the \(\Gamma\)-orbits of \((x C,y,i)\) for which
\(H(x^{-1}y^{-1}x)\ne0\) and \(x\) has a signed frame at
\(J_{\epsilon,i}(y)\). The action takes \(x C\) to \(z x C\)
and \(y\) to \(zyz^{-1}\), preserving the index. The definition
is independent of representatives: the frame matrix is taken
to \((z^{-1})^* A z^{-1}\) and right multiplication by \(C\)
changes its finite multiplier by a unit.
\end{definition}

\begin{lemma}\label{lem:marked-similitude-root-bounds}
For every term in Definition~\ref{def:actual-signed-orbits},
put \(g=x^{-1}y^{-1}x\) and let \(A\) be its signed-frame
matrix. This matrix is unique, and \(y^* A y=qA\) for a
positive integer \(q\). The multiplier of \(g\) is \(q^{-1}\).
The polynomial \(P\) is monic in \(O_E[T]\), and at every
embedding of \(E\) each of its roots has absolute value
\(\sqrt q\). For a fixed \(H\), the pairs \((q,P)\) occurring
in either orbit set form a finite set.
\end{lemma}
\begin{proof}
The finite frame identity \(x^* A x=\varepsilon I\) implies
\(A=A^*\) by injectivity of \(E\) into its completions. As
in Proposition~\ref{prop:definite-full-class-set}, two frame
matrices for \(x\) differ by a rational scalar which is a unit
at every prime. Positivity at an embedding in
\(\Psi\setminus\{\sigma\}\), nonempty by
Proposition~\ref{prop:places}, selects \(+1\). Since
\(g\in S^+\), substitution of \(yx=xg^{-1}\) gives
\(y^*Ay=t^{-1}A\), where \(g^*g=tI\). Dividing corresponding
entries at a nonzero entry of \(A\) makes \(q=t^{-1}\) a
rational scalar by the same rational basis argument. At an
embedding other than \(\sigma\) the equation is a similitude
of a positive matrix, giving \(q>0\).
Also \(g^{-1}\) is integral, so
\((g^{-1})^*g^{-1}=q I\) is integral and \(q\in\Z\).
The characteristic polynomial of \(g^{-1}\) equals \(P\).
Its coefficients, in \(E\) and integral at every finite
place, belong to \(O_E\) by
Proposition~\ref{prop:general-ideal-factorization}.

At \(\tau\in\Psi\) write \(J=J_{\epsilon,i}(y)\) if
\(\tau=\sigma\), and \(J=I\) otherwise. Because \(J\) commutes
with \(\tau y\), the equation also exhibits \(\tau y\) as
a similitude of the positive Hermitian matrix \(\tau(A)J\).
Evaluating at a nonzero eigenvector of eigenvalue \(r\) gives
\(|r|^2=q\). The conjugate embeddings have the conjugate
roots, so this holds at all embeddings.
On a right coset \(g_0 C\), all prime valuations of the
similitude multiplier agree with those for \(g_0\). They
determine the positive rational \(q\) uniquely, and \(H\)
has finite coset support. For fixed \(q\), the coefficient
of \(T^{n-j}\) in \(P\) has absolute value at most
\(\binom{n}{j}q^{j/2}\) at every embedding by the elementary
symmetric formula. There are finitely many possibilities in
the integral lattice \(O_E\) by Lemma~\ref{lem:number-field-lattice}.
\end{proof}

\begin{proposition}\label{prop:marked-polynomial-positive-form}
For each triple \((q,P,i)\) in Lemma~\ref{lem:marked-similitude-root-bounds}
and each \(\epsilon\), there is \(D\in O_E[T]\), depending only
on these data and \(\epsilon\), such that for every corresponding
frame the matrix \(A'=A D(y)\) is Hermitian positive definite
at all embeddings. Moreover \(D(y)\) is invertible, \(y\) is a
similitude of \(A'\) of multiplier \(q\), and \(x^*A'x\) has
integral entries at every finite place.
\end{proposition}
\begin{proof}
Factor \(P\) over \(E\) into its distinct monic irreducible
factors \(P_b\) by polynomial division, and put
\(L_b=E[T]/(P_b)\) and \(\lambda_b=T\bmod P_b\).
Every embedding of \(L_b\) in \(\C\) takes \(\lambda_b\) to
a number of squared absolute value \(q\), by the preceding lemma.
In one such embedding actual complex conjugation therefore
sends \(\lambda_b\) to \(q/\lambda_b\) and preserves \(E\).
It preserves \(L_b\), acts as an involution \(^{*}\), and has
the same description under every embedding by its values on
the generators. Its fixed field \(K_b\) is totally real,
since every embedding extends to \(L_b\) by
Lemma~\ref{lem:place-fields}. As \(i_0=\sqrt{-1}\in E\) by
Lemma~\ref{lem:skew-covol}, the two parts
\((z+z^*)/2\) and \((z-z^*)/(2i_0)\) give
\(L_b=K_b\oplus i_0 K_b\). Each real embedding of \(K_b\)
extends in exactly two ways, by the two images of \(i_0\).
Precisely one extends an embedding in the chosen type \(\Psi\).
Thus these real embeddings, as \(b\) varies, index exactly the
roots at all the embeddings in \(\Psi\).

At these embeddings prescribe the sign \(+1\), except that
for \(\epsilon=1\) prescribe \(-1\) at the \(n-1\) unmarked
roots over \(\sigma\). Choose \(a_b\in O_{K_b}\) having the
prescribed signs. Indeed the full lattice of
Lemma~\ref{lem:number-field-lattice} has a bounded fundamental
parallelotope, so a point sufficiently far inside that orthant
has a lattice point in its translate by this parallelotope.
The polynomials \(P_b\) are pairwise coprime. For each \(b\)
the product of the other factors has an inverse modulo \(P_b\)
by the Euclidean algorithm. Multiplying by this inverse
and summing gives a polynomial in \(E[T]\) with residues
\(a_b\) in all \(L_b\). A common positive integer denominator
clears its coefficients into \(O_E\) without changing any
sign. Call the resulting polynomial \(D\).

For a complex embedding in \(\Psi\), the eigenlines of \(y\)
are pairwise orthogonal for \(A\): on two distinct roots
\(r,s\) one has \(\bar r s\ne q\) by \(|r|^2=q\), and the
similitude equation forces their pairing to vanish.
Nonsingularity makes each diagonal entry in an eigenbasis
nonzero. Positivity of \(AJ\) means these entries have exactly
the prescribed signs. The values of \(D\) are real, nonzero,
and have these same signs. In this basis \(AD(y)\) is a
positive Hermitian form. The conjugate embeddings inherit
positivity. Hermitian symmetry over \(E\) follows already
at one embedding by injectivity. Each eigenvalue of \(D(y)\)
is nonzero and \(D(y)\) commutes with \(y\). This proves
invertibility and \(y^*A'y=q A'\). Finally
\(x^*A' x=\varepsilon D(g^{-1})\) is integral by the
coefficient choice and \(g^{-1}\in M_n(\widehat O)\).
\end{proof}

\section{Finiteness of the marked regular terms}

\begin{lemma}\label{lem:marked-positive-discriminant}
In Proposition~\ref{prop:marked-polynomial-positive-form} let
\(\Delta_E\) be the positive Gram determinant of
\(\operatorname{Tr}_{E/\Q}(\bar u v)\) on \(O_E\) in any
integral basis, and put \(\Lambda_x=E^n\cap x\widehat O^n\).
The positive trace lattice \((\Lambda_x,B')\), with
\(B'(v,w)=\operatorname{Tr}_{E/\Q}(v^*A' w)\), is integral
of rank \(n[E:\Q]\) and its Gram determinant is
\[
 \Delta_E^n\,
 \left|N_{E/\Q}\big(\det\nolimits_E D(y)\big)\right|.
\]
This number depends only on \(P,D\). The endomorphism \(y\)
preserves \(\Lambda_x\) by inclusion and multiplies its
squared lengths by \(q\).
\end{lemma}
\begin{proof}
The lattice, rank and completion statements follow exactly
as in Lemma~\ref{lem:trace-finite}: the common integer
denominators of \(x,x^{-1}\) bound the local lattices between
two completions of rational lattices. The finite quotient of
the two bounds decomposes into its prime parts, so imposing
the required local submodules commutes with completion.
The trace of an integral local element lies in \(\Z_p\);
in an integral basis it is the trace of its integral
multiplication matrix. Integrality of \(x^*A'x\) in
Proposition~\ref{prop:marked-polynomial-positive-form} now
gives integrality of \(B'\) at every prime and hence over
\(\Z\).

Over \(\Z_p\) use the \(x\)-transport of the product of \(n\)
integral bases of \(O_E\). Write \(Y=x^{-1}y x\).
Its trace Gram matrix is the standard trace Gram matrix
multiplied on the right by the matrix over \(\Q_p\) of
\(D(Y)\), and scaled by \(\varepsilon_p\), because
\(x^*A'x=\varepsilon_p D(Y)\) at \(p\).
The scalar is a unit and the initial Gram determinant is
\(\Delta_E^n\). On tensoring with \(\Q_p\) the linear maps
\(D(Y)\) and \(D(y)\) are conjugate by \(x\), also as
\(\Q_p\)-linear maps. The rational linear determinant of
\(D(y)\) equals \(N_{E/\Q}(\det_E D(y))\): after extending
scalars to \(\C\), the basis of embeddings in
Lemma~\ref{lem:number-field-lattice} separates the map into
one conjugate matrix for each embedding. Their determinants
multiply to the field norm. Also \(\det_E D(y)\) is
determined by \(P,D\), being the product of \(D(r)\) over
the roots of \(P\) in any splitting field. The local computation
therefore gives the claimed prime valuation of the global
Gram determinant at every prime. Both sides are positive
rational numbers, so they are equal. The inclusion
\(y\Lambda_x\subset\Lambda_x\) follows at all completions
from \(Y\in M_n(\widehat O)\), and the squared-length formula
follows by taking the trace of \(y^*A'y=qA'\).
\end{proof}

\begin{theorem}\label{thm:all-marked-regular-finite}
For any regular \(H\in\mathcal H^+\) and either
\(\epsilon\in\{0,1\}\), the entire orbit set
\(\mathcal O_\epsilon(H)\) of Definition~\ref{def:actual-signed-orbits}
is finite. Every orbit has trivial stabilizer in the action of
\(\Gamma\) on triples \((xC,y,i)\).
\end{theorem}
\begin{proof}
There are finitely many \(q,P,i\) by
Lemma~\ref{lem:marked-similitude-root-bounds}; fix one choice.
Use the single polynomial \(D\) of
Proposition~\ref{prop:marked-polynomial-positive-form} for this
choice and \(\epsilon\). Lemma~\ref{lem:marked-positive-discriminant}
gives the same determinant and rank for all the positive
trace lattices under consideration. Lemma~\ref{lem:trace-finite}
gives finitely many isometry classes. On any fixed representative
lattice a ball for its positive form meets the lattice in a
finite set, by boundedness in a real basis. For each member
\(b\) of the fixed integral basis of \(O_E\), multiplication
by \(b\) has norm at most \(\max_\tau|\tau(b)|\): the positive
embedding summands in the trace norm are each scaled by
\(|\tau(b)|^2\). Its values on each basis vector thus have
finitely many possibilities. The same holds independently for
\(y\), since it preserves the lattice by inclusion and has
operator norm \(\sqrt q\) by the preceding lemma. Hence the
trace lattice, its full \(O_E\)-action and its endomorphism
\(y\) form a finite list up to isomorphism.

A lattice isometry respecting this data extends over \(\Q\)
to an \(E\)-linear map intertwining the two endomorphisms.
It preserves the full \(E\)-valued form \(A'\), since one
can multiply its second argument by each element of \(E\)
and use nondegeneracy of field trace as in
Proposition~\ref{prop:definite-full-class-set}. It therefore
also preserves the original form recovered by
\[
          h_A(v,w)=h_{A'}(v,D(y)^{-1}w),
          \qquad h_M(v,w)=v^*Mw .
\]
The mark is transported because it is specified at a fixed
root of \(\sigma P\) and its reflection conjugates with \(y\).
After these rational identifications the same lattice,
endomorphism and original form are fixed. Their finite
frames differ on the right by an integral similitude, by
their equal local lattices and their scalar unit Gram
matrices. Such frames have finitely many right \(C\)-classes
by the finite reduction argument of
Proposition~\ref{prop:definite-full-class-set}. This proves
finiteness for the fixed data and hence for all the orbits.
Finally a stabilizing rational matrix commutes with \(y\),
hence at \(\sigma\) with \(J_{\epsilon,i}(y)\), and belongs
to \(x C x^{-1}\). Proposition~\ref{prop:neat-at-wild-place}
then gives the identity.
\end{proof}

\section{Matching the irreducible marked terms}

\begin{lemma}\label{lem:marked-similitude-trace-field}
Fix \(q,P,i\) occurring on either side in
Definition~\ref{def:actual-signed-orbits}, and suppose \(P\)
is irreducible over \(E\). Put \(L=E[T]/(P)\), let \(\lambda\)
be the residue of \(T\), and let \(K'\) be the fixed field of
the involution of Proposition~\ref{prop:marked-polynomial-positive-form}.
There is a basis identifying each rational \(y\) with multiplication
by \(\lambda\) on \(L\). In these coordinates the nonsingular
Hermitian forms \(A\) satisfying \(y^*Ay=qA\) are exactly
\[
     h_a(v,w)=\operatorname{Tr}_{L/E}(a\,v^* w),
                           \qquad a\in(K')^\times .
\]
The rational changes of coordinates commuting with \(y\)
are precisely the multiplications \(m_z\), \(z\in L^\times\),
and \(m_z^*h_a m_z=h_{az^*z}\). The signs on the eigenlines
are the signs of \(a\) at the corresponding real embeddings of \(K'\).
\end{lemma}
\begin{proof}
The adjugate identity gives \(P(y)=0\) by
Lemma~\ref{lem:torus-primary}. Evaluation defines a unital
map of the field \(L\) into \({\rm End}_E(E^n)\). Its action
on any nonzero vector identifies \(L\) with \(E^n\), by
injectivity and equality of dimensions. An endomorphism
commuting with \(y\) is \(L\)-linear and hence is multiplication
by its value on \(1\). The adjoint of \(y\) for the given
Hermitian form is \(q y^{-1}\). On \(E\) the adjoint is
conjugation, so on \(L\) it is exactly the involution described
in Proposition~\ref{prop:marked-polynomial-positive-form}.
For any one of these forms write \(f(w)=h(1,w)\).
Then \(h(v,w)=f(v^*w)\). The field trace pairing identifies
\(L\) with its \(E\)-linear dual because a nonzero \(b\)
paired with \(b^{-1}\) has nonzero trace \(n\). Thus
\(f(w)=\operatorname{Tr}_{L/E}(aw)\) for a unique \(a\in L\).
Taking a star commutes with trace up to conjugation on \(E\),
as seen by summing over the embeddings above a given embedding
of \(E\), using Lemma~\ref{lem:number-field-lattice}. Hermitian
symmetry now gives \(a^*=a\) by the same nondegenerate pairing.
Nonsingularity gives \(a\ne0\). Conversely such \(a\) gives
a nonsingular form and \(\lambda^*\lambda=q\) gives the
required multiplier. Multiplication by \(z\) acts by the
displayed pullback. At an embedding in \(\Psi\), the power
basis of \(L\) evaluated at the \(n\) distinct roots has
invertible Vandermonde matrix, as in
Theorem~\ref{thm:elliptic-trace-fields}. Trace in those
coordinates is addition and the star is complex conjugation.
The resulting Hermitian form is diagonal with entries the
specified real values of \(a\). This proves the signs.
\end{proof}

\begin{theorem}\label{thm:irreducible-marked-terms-match}
Assume \(n\) is odd. For each regular \(H\in\mathcal H^+\)
there is a bijection between the terms of \(\mathcal O_0(H)\)
and \(\mathcal O_1(H)\) with irreducible \(P\). It preserves
\(P,i,q\) and the value of \(H(x^{-1}y^{-1}x)\) at each term.
\end{theorem}
\begin{proof}
Fix \(P,i,q\) occurring on either side. Take \(L,K',\lambda\)
as in Lemma~\ref{lem:marked-similitude-trace-field}. The real
embeddings of \(K'\) correspond exactly to the eigenlines in
the CM type. The change from the signs of \(J_{0,i}\) to
those of \(J_{1,i}\) reverses \(n-1\) signs at \(\sigma\)
and none elsewhere. Their product is one because \(n\) is
odd. The field \(L\) contains the twelfth root of unity of
Lemma~\ref{lem:skew-covol} and \(L=K'(i_0)\) with \(K'\)
totally real, by Proposition~\ref{prop:marked-polynomial-positive-form}.
Theorem~\ref{thm:even-sign-norms} supplies \(d\in (K')^\times\)
with exactly those signs and a finite idele \(z\) of \(L\)
such that \(z^*z=d\) at every finite place. Fix this choice
for the given data.

Use the coordinates of the lemma, in which a triple is
represented by \(a\) and an adelic isomorphism
\(x:\mathbb A_{E,f}^n\to L\otimes_E\mathbb A_{E,f}\).
The finite Gram equation reads \(x^*h_ax=\varepsilon I\)
with \(\varepsilon\) any rational profinite unit. Changing the
rational coordinates of that fixed \(y\) and changing the
frame representative identify
\[
     (a,x)\ \sim\ (a/(w^*w),\,m_w x k),
                       \qquad w\in L^\times,\ k\in C .
\]
These are all identifications by the centralizer calculation
in the lemma. The matrix \(A\), and hence its parameter \(a\)
in a chosen coordinate system, is uniquely determined for a
signed frame by Lemma~\ref{lem:marked-similitude-root-bounds}.
The signed condition says exactly that \(a\) has the designated
signs, by the lemma and the definition of \(J_{\epsilon,i}\).

Take \((a,x)\) to \((da,m_z^{-1}x)\). Both finite Gram matrices
are identical, by \(z^*z=d\), and the real signs change as
required. The finite Hecke argument is
\(x^{-1}m_{\lambda}^{-1}x\). It stays identical because the
multiplications \(m_z,m_\lambda\) commute. The map descends
to equivalence classes since it commutes with all the displayed
changes of representative, as in Theorem~\ref{thm:elliptic-class-bijection}.
Its inverse uses \(d^{-1},z^{-1}\). Thus it matches the two
sets for the given data and preserves every indicated weight.
Taking the union over the data gives the assertion.
\end{proof}

\section{Cancellation at the place of the cuspidal character}

\begin{lemma}\label{lem:marked-sylvester-translation}
Let \(G_0\in{\rm GL}_n(\Q_\ell)\) have separable characteristic
polynomial \(P=P_1P_2\), with \(P_1,P_2\) coprime and of
positive degree. Write \(V_a=\ker P_a(G_0)\) for \(a=1,2\).
These give a direct sum decomposition of \(\Q_\ell^n\).
Let \(M\) be the space of maps with image in \(V_1\) vanishing
on \(V_1\), and write \(u_B=1+B\) for \(B\in M\).
The map
\[
      \mathcal L:M\longrightarrow M,\qquad
                 \mathcal L(B)=G_0 B-BG_0
\]
is invertible, \(u_B^{-1}G_0u_B=G_0+\mathcal L(B)\), and
for \(D\in M\) there exists a single \(T\in M\) such that
\[
 u_{B+T}^{-1}G_0u_{B+T}=u_D\,u_B^{-1}G_0u_B
                           \qquad(B\in M).
\]
The set of \(B\) with \(u_B^{-1}G_0u_B\) integral is compact.
\end{lemma}
\begin{proof}
An identity \(a_1P_1+a_2P_2=1\) and \(P(G_0)=0\)
decompose every vector into its two kernel components and
show their intersection is zero. Each component is nonzero
by evaluating at a root of its corresponding factor in a
splitting field; extension of scalars preserves a matrix
kernel's dimension by row elimination. The space \(M\)
identifies with \({\rm Hom}(V_2,V_1)\). Its members have
pairwise zero products, and \(G_0\) preserves both subspaces.
Thus both left and right multiplication by \(G_0\) preserve
\(M\), and \(B G_0 B'=0\) for \(B,B'\in M\).
If \(G_0B=BG_0\), then \(P_1(G_0)B=BP_1(G_0)=0\).
The restriction of \(P_1(G_0)\) to \(V_2\) is invertible by
the displayed Bezout identity. Thus \(B=0\), giving
invertibility of \(\mathcal L\) in finite dimension.
Use \(u_B^{-1}=1-B\) and the zero products to obtain the
conjugation formula. Taking \(T=\mathcal L^{-1}(DG_0)\)
gives the required translation since \(D\mathcal L(B)=0\).
Finally \(G_0+M\) is closed, being an affine subspace
defined by linear equations; its intersection with the
compact integral matrices is compact. The inverse affine
map to \(M\) is continuous by its finite matrix formula in
bases. This proves the last assertion.
\end{proof}

\begin{theorem}\label{thm:reducible-marked-terms-cancel}
Let \(H\in\mathcal H^+\) be regular and suppose its coset
function satisfies the character relation of
Theorem~\ref{thm:actual-regular-hecke-test}. For
\(\epsilon\in\{0,1\}\) and any complex weight \(w(P,i)\),
\[
 \sum_{\substack{[(xC,y,i)]\in\mathcal O_\epsilon(H)\\
                  P=\det(TI-y)\ {\rm reducible\ over}\ E}}
       H(x^{-1}y^{-1}x)\,w(P,i)=0 .
\]
\end{theorem}
\begin{proof}
The whole sum is finite by Theorem~\ref{thm:all-marked-regular-finite}.
Fix one occurring reducible \(P\) and a root index \(i\).
Choose a factorization \(P=P_1P_2\) in \(E[T]\) with both
factors nonconstant and coprime, possible by separability.
It remains coprime after base change to \(E_{w_0}=\Q_\ell\),
by the Bezout identity. Work before the rational quotient
with the signed triples \((xC,y,i)\) of this polynomial and
index, allowing the weight to be zero, and imposing just
\(s=x^{-1}y^{-1}x\in S\) on the Hecke argument. All nonzero
terms occur in this set. At one triple let \(G_0=d_{w_0}(s)\).
It is conjugate to \(y_{w_0}^{\mathsf T}\), hence has
polynomial \(P\). Use \(M,u_B\) of
Lemma~\ref{lem:marked-sylvester-translation} for \(G_0\),
and write \(b=b_{w_0}\). Consider the set of triples
\[
                  (x b(u_B) C,\ y,\ i),\qquad B\in M .
\]
The finite Gram matrix stays scalar with the same multiplier
since \(b(u_B)\) is unitary, and the rational signed form
is unchanged. All new Hecke arguments lie in \(S\).

These sets partition the triples under consideration. Indeed
a change \(x\mapsto xk\) with \(k\in C\) conjugates \(G_0\)
and its two kernels by \(d_{w_0}(k)^{-1}\), and so conjugates
the corresponding \(M\) by the same matrix. The identity
\(k b(v)k^{-1}=b(d_{w_0}(k)v d_{w_0}(k)^{-1})\) from
Proposition~\ref{prop:local-translation-frame-extension}
makes the displayed set independent of the representative.
Changing the base point within it by \(b(u_B)\) conjugates its
local group \(1+M\) by \(u_B^{-1}\), leaving that group
unchanged since it is abelian. The map \(B\mapsto u_B\) is
an isomorphism from the additive group of \(M\) to this
group, with \(u_Bu_{B'}=u_{B+B'}\). Thus every member gives
the same set. Rational change of coordinates conjugating \(y\)
and multiplying \(x\) on the left leaves \(s\) unchanged,
hence permutes these sets.

In a fixed set, the displayed triples are indexed without
repetition by
\[
         M/M_C,\qquad M_C=\{B\in M:b(u_B)\in C\}.
\]
The subgroup \(M_C\) is open by continuity and openness of
\(C\) in \(S\). For a nonzero term the dual component of
\(s_B=b(u_B)^{-1}s b(u_B)\) is integral because \(s_B^{-1}\)
is integral. This component is \(u_B^{-1}G_0u_B\).
Lemma~\ref{lem:marked-sylvester-translation} confines the
representatives of nonzero terms to a compact subset of \(M\).
Its image in the discrete quotient by \(M_C\) is finite.

Take the projection \(e\) onto \(V_1\) along \(V_2\).
It is neither zero nor the identity and \(M=eM_n(\Q_\ell)(1-e)\).
Proposition~\ref{prop:actual-wild-split-character} gives
\(u_D\in J\cap(1+M)\) with \(\chi(u_D)\ne1\).
For the \(T\) given by the preceding lemma one has
\(s_{B+T}=b(u_D)s_B\). Indeed, for a similitude \(s\),
\(s b(v) s^{-1}=b(G_0 v G_0^{-1})\) by the same two-place
matrix identity as for \(k\) above. Consequently
\(s_B=b(u_B^{-1}G_0u_B G_0^{-1})s\), giving the claim
from the local equation. The translation by \(T\) permutes
\(M/M_C\) and multiplies every summand by \(\chi(u_D)^{-1}\).
The sum on this set has finite support, hence is zero.

It remains to pass to rational orbits. For a set in this
partition meeting the nonzero support of \(H\), its rational
setwise stabilizer permutes the finite nonempty subset on
which \(H(s_B)\ne0\), since \(H\) is invariant under a
change of representatives. On that subset point stabilizers
are trivial by Theorem~\ref{thm:all-marked-regular-finite}.
The setwise stabilizer is therefore finite and acts freely
there. Rational orbits in the rational saturation of this
set are in bijection with its orbits under that stabilizer:
two of its elements in the same rational orbit differ by a
map taking intersecting parts of the partition to each other,
hence preserving that part. Each nonzero rational-orbit weight
thus occurs exactly the stabilizer order times in the local sum.
Dividing the zero sum by that order gives zero on its saturation.
Only finitely many such saturations contain contributing rational
orbits, again by Theorem~\ref{thm:all-marked-regular-finite}.
Summing over them and over \(P,i\) proves the assertion.
\end{proof}

\section{The definite eigensystem in the marked orbital formula}

\begin{theorem}\label{thm:actual-signed-orbit-transfer}
Let \(n\) be odd and let \(H\in\mathcal H^+\) be regular
with \(H(b_{w_0}(j)g)=\chi(j)^{-1}H(g)\) as in
Theorem~\ref{thm:actual-regular-hecke-test}. For any complex
function \(\omega(P,i)\) of the characteristic polynomial and
its marked root, the two finite sums
\[
 \sum_{[(xC,y,i)]\in\mathcal O_\epsilon(H)}
       H(x^{-1}y^{-1}x)\,\omega(P,i),\qquad \epsilon=0,1,
\]
are equal.
\end{theorem}
\begin{proof}
Theorem~\ref{thm:all-marked-regular-finite} gives finiteness
on both sides. On irreducible polynomials the bijection of
Theorem~\ref{thm:irreducible-marked-terms-match} preserves
each summand. For reducible polynomials each sum vanishes
by Theorem~\ref{thm:reducible-marked-terms-cancel}.
\end{proof}

\begin{definition}\label{def:marked-weight-test}
Let \(\tau_{w,r}\in\mathcal H^+\) be the safe spherical
operator of Lemma~\ref{lem:safe-operators-away-from-markers}
for \(0\le r\le n\), extended by zero for \(r>n\).
A polynomial in these operators may be a finite complex
linear combination of ordered words, including the empty
word. For regular \(H\in\mathcal H^+\), \(\alpha\in O_E^\nu\)
and integer \(h\ge0\), write
\[
 M_{\alpha,h}(H)=
  \sum_{[(xC,y,i)]\in\mathcal O_1(H)}
       H(x^{-1}y^{-1}x)\,[X^h]B_{\alpha,y,i}(X)
\]
with \(B_{\alpha,y,i}\) from Definition~\ref{def:actual-finite-local-terms}.
Here \(B_{\alpha,y,i}\) depends only on \(P,i\): the unordered
root lists at the embeddings and the marked root at \(\sigma\)
and its conjugate determine every factor.
Say that a system \(\theta=(\theta_{w,r})\) occurs in the
marked orbital formula in weight \(h\) if, for every polynomial
\(F\) with \(F(\theta)=1\), there is \(B\in\mathcal H^+\) with
\(BF(\tau)\) regular and \(M_{\alpha,h}(BF(\tau))\ne0\) for
some integral tuple \(\alpha\).
\end{definition}

\begin{theorem}\label{thm:actual-marked-weight-occurrence}
For any one approximating form, choose \(\mathfrak n,\psi,\Xi\)
with the oriented coefficient system of
Definition~\ref{def:geometric-coefficient-system}, as in
Theorems~\ref{thm:actual-central-test-vector} and~\ref{thm:actual-jump-character}.
Set \(\theta_{w,r}=c_w^{-r}t_{w,r}\) at safe places for
\(0\le r\le n\) and zero for \(r>n\).
Then \(\theta\) occurs in the marked orbital formula in
each weight \(h_j=\mu+k_{v_0,j}\), \(1\le j\le n\).
For a fixed \(F\) the same \(B\) can be used in all of
these weights and each testing tuple can be taken in the finite
grid \(\mathcal B_n^\nu\) of Theorem~\ref{thm:finite-geometric-local-terms}.
\end{theorem}
\begin{proof}
The weights \(h_j\) are nonnegative:
\(k_{v_0,j}+m_{v_0}\ge0\) and
\(\mu=n\sum_v m_v\ge m_{v_0}\) by
Definition~\ref{def:geometric-coefficient-system}.
By Theorem~\ref{thm:actual-central-test-vector}, the vector
\(\psi\ne0\) is a common eigenvector with the stated
eigenvalues. Thus \(c=F(\tau)\) fixes \(\psi\) when
\(F(\theta)=1\), by applying each word successively.
Corollary~\ref{cor:definite-orbit-sum-detects-jump} yields a
regular \(H=(bc)^m=Bc\) whose definite trace is nonzero,
with \(B=(bc)^{m-1}b\), and the indicated character relation.
That corollary gives for each \(h_j\) a grid tuple \(\alpha\)
for which the sum over \(\mathcal O(H)\), with internal sum
over marks of \([X^{h_j}]B_{\alpha,y,i}\), is nonzero.
The expanded sum is exactly the sum on \(\mathcal O_0(H)\):
for each rational orbit of pairs, all \(n\) marks are preserved
under a change of representative by their common ordering of
\(\sigma P\). Apply Theorem~\ref{thm:actual-signed-orbit-transfer}
with \(\omega(P,i)=[X^{h_j}]B_{\alpha,y,i}\). Its marked sum
is \(M_{\alpha,h_j}(H)\) and has the same nonzero value. The
choices of \(H,B\) preceded the choice of \(j\), proving the
simultaneous assertion.
\end{proof}

\section{The complex period family of a signed frame}

\begin{definition}\label{def:signed-complex-period-space}
Let \(V_\R=\prod_{\tau\in\Psi}\C^n\) with the underlying real
linear structure; embed \(E^n\) by its embeddings in \(\Psi\).
For \(x\in G\) put \(\Lambda_x=E^n\cap x\widehat O^n\) in
\(V_\R\), as in Lemma~\ref{lem:marked-positive-discriminant}.
A reflection \(J\) is a complex-linear involution on \(\C^n\)
whose \(+1\) eigenspace is a line and whose \(-1\) eigenspace
has dimension \(n-1\). Use triples \((x,J,(P_a)_{a=1}^\nu)\)
with \(P_a\in V_\R/\Lambda_x\) and \(x,J\) admitting a
signed frame in Definition~\ref{def:moduli-signed-frame}.
For the complex structure on \(V_\R\) use multiplication by
\(iJ\) at \(\sigma\) and by \(i\) at all other places of
\(\Psi\). Each of these real maps squares to \(-I\),
giving complex scalar multiplication by
\((u+iv)w=uw+v(iJ)w\) at \(\sigma\) and the usual rule
elsewhere. The finite level \(C\) identifies \(x\) and \(xk\),
\(k\in C\), by the identity on the lattice and on the remaining
data. A rational \(t\in\Gamma\) acts by \(x\mapsto tx\),
\(J\mapsto \sigma(t)J\sigma(t)^{-1}\) and \(P_a\mapsto tP_a\).
The lattice then becomes \(t\Lambda_x\). Denote the classes
by \(\mathcal X(C,\nu)\).
\end{definition}

\begin{proposition}\label{prop:signed-ball-period-coordinates}
Fix an \(x\) admitting some reflection and let \(A\) be its
unique rational signed-frame matrix. At the other embeddings
of \(\Psi\) it is positive definite, and \(H=\sigma(A)\)
has signature \((1,n-1)\). The allowable reflections are in
bijection with the positive complex lines for \(H\). In a
basis with \(H=\operatorname{diag}(1,-I_{n-1})\) they are
parametrized by
\[
       u_b=(1,\bar b),\qquad
       b\in\mathbb B=\{b\in\C^{n-1}:\|b\|^2<1\}.
\]
For \(w=(w_0,w')\) at \(\sigma\) there is a period map
\[
       Q_b(w)=\big(w_0-b^{\mathsf T}w',\
                           \bar w'-b\bar w_0\big).
\]
At the remaining embeddings use the identity in defining
\(Q_b:V_\R\longrightarrow \C^{fn}\).
It is a real-linear isomorphism, carrying the complex
structure of Definition~\ref{def:signed-complex-period-space}
to the standard one. For fixed \(w\) it is holomorphic in
\(b\). In particular the action of \(a\in E\) in these complex
coordinates has characters \(\sigma(a)\) on the first line,
\(\bar\sigma(a)\) with multiplicity \(n-1\), and \(\tau(a)\)
with multiplicity \(n\) at each remaining \(\tau\in\Psi\).
\end{proposition}
\begin{proof}
The uniqueness proof of Lemma~\ref{lem:marked-similitude-root-bounds}
uses only \(x^*Ax=\varepsilon I\) and positivity at another
embedding; it applies to all allowable reflections for \(x\).
The form \(HJ\) is positive Hermitian, hence \(J^*H=HJ\).
The two eigenspaces of \(J\) are \(H\)-orthogonal. On the
positive line of \(J\) the form \(H\) is positive and on
its complement it is negative. Conversely a positive line
determines \(J\) by acting as one there and minus one on
its \(H\)-orthogonal complement. That complement is negative:
in a diagonal basis with one positive coordinate, subtracting
the projection onto a positive unit vector removes the only
positive square, and the restriction is nondegenerate.
A diagonal basis with entries \(1,-1,\ldots,-1\) exists from
any initial reflection by orthogonalizing separately in its
positive and negative eigenspaces and rescaling. In these
coordinates each positive line has a nonzero first coordinate.
Normalizing it to one gives \(u_b\) with \(\|b\|<1\), and
every such vector spans a positive line. More explicitly
its orthogonal complement consists of \(w_0=b^{\mathsf T}w'\)
and is negative there by Cauchy--Schwarz and \(\|b\|<1\).

The first component of \(Q_b\) is complex linear in \(w\)
and kills the negative complement. The second is conjugate
linear and kills the positive line. They therefore carry
\(iJ\) to ordinary multiplication by \(i\): on the negative
summand complex conjugation changes the sign of \(i\).
If both components vanish, \(w'=\bar b w_0\) and
\(w_0=\|b\|^2w_0\), hence \(w=0\).
Equality of the real dimensions gives invertibility.
The formula is affine in \(b\) at each fixed \(w\).
Its linearity in the positive coordinate and conjugate
linearity in the negative coordinates give the displayed
characters. These calculations also apply at \(n=1\), with
\(\mathbb B\) a single point.
\end{proof}

\begin{theorem}\label{thm:analytic-framed-torus-family}
For fixed \(x\), the triples in
Definition~\ref{def:signed-complex-period-space} before
rational identifications form a holomorphic family over
\(\mathbb B\) of \(\nu\)-fold products of compact complex tori.
In the coordinates of Proposition~\ref{prop:signed-ball-period-coordinates}
the quotient manifold of dimension \(n-1+\nu fn\) is
\[
 \big(\mathbb B\times(\C^{fn})^\nu\big)/\Lambda_x^\nu ,
 \qquad (b,z_a)_a\longmapsto (b,z_a+Q_b(\lambda_a))_a .
\]
A rational change \(x\mapsto tx\) induces a biholomorphic
isomorphism of these families, with the induced transformation
of the base. The equivalence classes \(\mathcal X(C,\nu)\)
have complex manifold charts locally induced by those families.
\end{theorem}
\begin{proof}
The lattice is full by Lemmas~\ref{lem:trace-finite}
and~\ref{lem:number-field-lattice}. Its image by the
real isomorphism \(Q_b\) is full, so the quotient in a
fiber is compact by its bounded closed fundamental
parallelotope. All lattice translations in the statement
are holomorphic by the preceding proposition, with inverses
given by \(-\lambda_a\). They act freely. For \(b\) in any
compact subset, the inverse real matrices \(Q_b^{-1}\)
are bounded by the cofactor formula and the nonvanishing
of their determinants. On a compact subset of the total
space, a translation meeting this subset again has bounded
\(\lambda_a\), hence there are only finitely many such
translations. Near a point choose a relatively compact
neighborhood, and then shrink it to avoid all its nonidentity
translates on this finite list, using freeness and continuity.
This gives the quotient complex charts. Two distinct orbits
have disjoint image neighborhoods by the same compactness
argument with a neighborhood of each point, proving the
Hausdorff property.

A rational \(t\) takes the associated form to
\((t^{-1})^*At^{-1}\). The transformation \(b\mapsto b'\)
is holomorphic since \((1,b')\) is obtained from the image
of \((1,b)\) by the conjugate matrix of \(\sigma(t)\) in
the chosen two bases, dividing by its nonzero first coordinate.
It carries \(J_b\) to \(J_{b'}\), so
\(Q'_{b'} t Q_b^{-1}\) is complex linear. This matrix is
holomorphic in \(b\): at any fixed \(b_0\) choose real
vectors \(v_j\) whose \(Q_{b_0}\)-images are a complex
basis. Locally the matrix is the product of the matrix
with columns \(Q'_{b'}t v_j\) and the inverse of the one
with columns \(Q_b v_j\). Both column matrices are holomorphic
and the latter stays invertible. Translation of lattices
is respected since \(t\Lambda_x=\Lambda_{tx}\). Applying
the argument to \(t^{-1}\) gives the inverse.

It remains to pass to rational orbits. Take one \(xC\) in
each rational orbit of such frames. For \(t\) with
\(x^{-1}tx\in C\), uniqueness of the frame matrix from
the finite scalar Gram equation and positivity at a second
embedding gives \(t^*At=A\). Indeed the transformed rational
form has these same properties. Such \(t\) act properly on
the ball in the following explicit sense. If both \(b\) and
\(tb\) lie in a fixed compact subset, then
\(\sigma(t)^* HJ_{tb}\sigma(t)=HJ_b\). The matrices \(HJ_b\)
vary continuously: the projection onto \(u_b\) is
\(w\mapsto u_b(u_b^*Hw)/(u_b^*Hu_b)\), with nonzero
denominator. On this compact subset the maxima and the
positive minima on unit vectors bound \(\sigma(t)\).
At all the other embeddings, preservation of the positive
matrix \(\tau(A)\) also bounds the entries of \(t\).
Its finite entries have common bounded denominators since
\(t\in x C x^{-1}\). Only finitely many rational matrices
satisfy all these bounds by Lemma~\ref{lem:number-field-lattice}.
The stabilizer of each \(b\) is trivial by
Proposition~\ref{prop:neat-at-wild-place}. Thus these rational
actions on the base, and hence on the torus families, are
free and have the same locally disjoint neighborhood property
as the lattice translations. Taking the union of the quotient
charts over the frame representatives gives the claimed charts.
This union is second countable. There are countably many possible
rational matrices \(A\), and for fixed \(A\) every finite frame
differs from any given one on the right by an element of \(S\),
by the two scalar Gram identities. The set \(S/C\) is countable:
the finite adeles have a countable base formed from finitely
many local balls and the integral factors elsewhere. Local balls
have a countable base by uniformizers and finite residue digit
representatives as in Lemma~\ref{lem:local-residue-indicators}.
Finite matrices and their subspaces inherit a countable base,
and each open coset in \(S\) contains a distinct basic set.
Finally each ball and vector space is second countable, and
images of a countable base under an open quotient map form a base.
\end{proof}

\section{Compactness and polarization of the signed family}

\begin{lemma}\label{lem:period-metric-covolume}
In the setting of Theorem~\ref{thm:analytic-framed-torus-family}
write \(A_\tau=\tau(A)\), \(J_\sigma=J\) and \(J_\tau=I\)
for \(\tau\ne\sigma\). On \(V_\R\) put
\[
 B_{A,J}(v,w)=2\Re\sum_{\tau\in\Psi}
                       v_\tau^* A_\tau J_\tau w_\tau .
\]
It is a positive real inner product. In a \(\Z\)-basis of
\(\Lambda_x\) its determinant is \(\Delta_E^n\) of
Lemma~\ref{lem:marked-positive-discriminant}. Every nonzero
lattice vector satisfies \(B_{A,J}(v,v)\ge2\).
\end{lemma}
\begin{proof}
Positivity and symmetry follow from the positive Hermitian
matrices \(A_\tau J_\tau\). Write \(B_A\) for the real form
with each \(J_\tau\) replaced by \(I\). On the rational space
it is the trace of \(v^*Aw\) by Lemma~\ref{lem:number-field-lattice}.
The finite Gram equation \(x^*A x=\varepsilon I\) shows that
\(B_A\) is integral, with determinant of absolute value
\(\Delta_E^n\), by the local basis proof of
Lemma~\ref{lem:marked-positive-discriminant} applied directly
to this Gram equation. In a real basis
\(B_{A,J}(v,w)=B_A(v,\mathcal Jw)\) where \(\mathcal J\) has
components \(J_\tau\). Its real determinant is one, since
the complex-linear involution \(J\) has eigenvalues \(+1,-1\)
on complex subspaces and each real multiplicity is even.
This gives the assertion on the positive determinant.

For \(0\ne v\in\Lambda_x\) the element \(a=v^*Av\) belongs
to \(O_F\): it is conjugation invariant and integral in
every completion by the same finite Gram equation. Its value
at any of the other places is strictly positive, since \(v\)
has a nonzero image under each field embedding and \(A_\tau\)
is positive there. Thus \(a\ne0\) and \(|N_{F/\Q}a|\) is
a positive integer by Lemma~\ref{lem:number-field-lattice}.
Some real embedding has \(|\tau(a)|\ge1\), since the product
of these absolute values is at least one. At each place
\(|\tau(a)|\le v_\tau^* A_\tau J_\tau v_\tau\): in the
orthogonal decomposition from
Proposition~\ref{prop:signed-ball-period-coordinates} the
left side is the absolute difference and the right side the
sum of two nonnegative squared lengths. This proves the bound.
\end{proof}

\begin{theorem}\label{thm:compact-signed-complex-family}
The complex manifold \(\mathcal X(C,\nu)\) of
Definition~\ref{def:signed-complex-period-space} is compact
for every integer \(\nu\ge0\).
\end{theorem}
\begin{proof}
The lattice basis bound proved in Lemma~\ref{lem:trace-finite}
applies to the Euclidean metric \(B_{A,J}\): its covolume
and a lower bound for its shortest length are fixed by
Lemma~\ref{lem:period-metric-covolume}. Choose for each point
an ordered lattice basis with all lengths bounded by one
constant. Put \(D=n[E:\Q]\). In this basis the positive
Gram matrices \(M\) have bounded entries by Cauchy--Schwarz,
and determinant \(\Delta_E^n>0\). Hence their inverses also
have bounded entries by cofactors. Coordinates of a vector
of bounded length are bounded uniformly: its pairings with
the basis vectors are bounded and its coordinate column is
\(M^{-1}\) times the column of those pairings.

The integral trace form \(B_A\) also has bounded entries
in the chosen basis, since \(B_A(v,w)=B_{A,J}(v,\mathcal J w)\)
and \(\mathcal J\) is an isometry for the positive metric
by the orthogonal decomposition. For each integral basis
element \(b\in O_E\), multiplication by \(b\) is integral
on the lattice and has norm at most \(\max_\tau|\tau(b)|\)
by the positive embedding summands. Its integer coordinate
matrix therefore has bounded entries. There are only
finitely many combined possibilities for the matrix of
\(B_A\) and this list of multiplication matrices.
Identify any two ordered bases with identical data. The
resulting rational linear isomorphism is \(E\)-linear and
recovers \(A\) by nondegeneracy of the field trace, as in
Proposition~\ref{prop:definite-full-class-set}. Thus, on
one fixed representative per list, the ordered lattice
basis and rational form can be kept fixed. The possible
\(x\) with this lattice and form have only finitely many
right \(C\)-classes by the same finite reduction argument
in that proposition, which uses only the local scalar unit
Gram equations. We have made only rational and level
identifications allowed in \(\mathcal X(C,\nu)\).

For any one of these finitely many representatives,
the possible positive matrices \(M\) lie in a compact set
of positive matrices: take the closure with the uniform
entry bounds and fixed positive determinant. Limits of
positive matrices are positive semidefinite by evaluation
on real vectors, and a nonsingular such matrix is positive
definite by orthogonal diagonalization, or completing
squares. In the fixed real basis the operator \(\mathcal J\)
is \(B_A^{-1} M\), with \(B_A\) here its fixed nonsingular
Gram matrix. The conditions of being identity off \(\sigma\),
complex linear at \(\sigma\), having \(J^2=I\), and
\({\rm tr}_{\C}J=2-n\), are closed. The last two equations
imply the required multiplicities since the projectors
\((I\pm J)/2\) split the space. Positive definiteness of
\(M\) implies positivity of \(A_\sigma J\): restriction
to this complex factor tests the real part of the form,
and symmetry of \(M\), also applied after multiplication
by \(i\), yields Hermitian symmetry. Hence the allowed
reflections in this compact set form a compact set in
the ball coordinates: \(J\mapsto\operatorname{im}(I+J)\)
is continuous locally by choosing a nonzero column, and
the positive line has nonzero first coordinate in the
coordinates of Proposition~\ref{prop:signed-ball-period-coordinates}.

For each of the \(\nu\) torus points use a representative in
the fixed closed fundamental parallelotope of \(\Lambda_x\)
in \(V_\R\). The maps taking the reflection and vectors in
these compact sets to the family are continuous by the period
charts of Theorem~\ref{thm:analytic-framed-torus-family}.
Their images, for the finitely many representatives, cover
the whole manifold. It is compact.
\end{proof}

\begin{proposition}\label{prop:signed-period-riemann-form}
Let \(\xi=\delta/\varpi\) be as in
Theorem~\ref{thm:actual-pel-torsion-pairing}. For the complex
torus at \((x,J)\) the real extension of
\[
         E_A(v,w)=\operatorname{Tr}_{E/\Q}(\xi\,w^*Av)
                            \qquad(v,w\in E^n)
\]
is an integral Riemann form: it is alternating, takes integer
values on the period lattice, is invariant under simultaneous
application of the complex structure \(I_J=(iJ_\tau)_\tau\),
and \(E_A(v,I_Jv)>0\) for \(v\ne0\).
The quotient of the dual lattice for this form by \(\Lambda_x\)
has order \(h^{2n}\) and exponent dividing \(h\), with
\(h=|N_{F/\Q}\delta|\).
\end{proposition}
\begin{proof}
Alternation follows from \(\bar\xi=-\xi\), \(A^*=A\) and
invariance of field trace under conjugation. At each prime
the form in frame coordinates has matrix pairing
\(\operatorname{Tr}(\xi\varepsilon_p w^*v)\).
The unit \(\varepsilon_p\in\Z_p^\times\) keeps the integral
dual unchanged. Since \(\varpi\) generates the different,
the integral dual of the local standard lattice is precisely
\(\delta^{-1} O_{E,p}^n\), by Proposition~\ref{prop:different}
or the trace-dual computation in
Theorem~\ref{thm:actual-pel-torsion-pairing}.
Thus \(E_A\) is integral on \(\Lambda_x\). The global dual
is determined by these local duals: apply the local Gram
equations in a rational lattice basis, or the completion
argument of Lemma~\ref{lem:trace-finite}.
After transport by \(x\), the finite quotients therefore
have the same orders and exponents as for the standard
lattice in Theorem~\ref{thm:actual-pel-torsion-pairing}.
Their direct product is the global quotient by prime
decomposition of a finite abelian group, giving the
stated size and exponent.

In the real coordinates field trace is \(2\Re\sum_{\tau\in\Psi}\).
Since \(J_\tau^*A_\tau J_\tau=A_\tau\), simultaneous
application of \(iJ_\tau\) in these coordinates preserves
the alternating form. Finally
\[
 E_A(v,I_Jv)=2\sum_{\tau\in\Psi}\Im(\tau(\xi))\,
                            v_\tau^* A_\tau J_\tau v_\tau .
\]
All coefficients are strictly positive by
Lemma~\ref{lem:actual-signs-and-padding}, and each Hermitian
form in the sum is positive definite. This is the required
positivity.
\end{proof}

\section{Hecke legs and their coincidence sets on the complex family}

\begin{proposition}\label{prop:complex-actual-hecke-legs}
Let \(D_0\) be a finite set of right cosets \(gC\) in
\(S^+/C\). For all sufficiently large integers \(m\), put
\(C_m=\{k\in C:k\equiv I\bmod m!\widehat O\}\). It is normal
of finite index in \(C\), satisfies \(g^{-1}C_m g\subseteq C\)
for all these cosets, and forgetting the finer level is a
holomorphic covering
\(\pi_m:\mathcal X(C_m,\nu)\to\mathcal X(C,\nu)\).
The group \(C/C_m\) acts simply transitively on its fibers
by right translation on the finite frame.
For each \(gC\in D_0\) and \(\mathbf a=(a_j)\in O_E^\nu\)
there is a holomorphic map
\[
 u_{g,\mathbf a}:\mathcal X(C_m,\nu)\longrightarrow\mathcal X(C,\nu),
 \qquad [x,J,(P_j)]\longmapsto[xg,J,(a_jP_j)] .
\]
\end{proposition}
\begin{proof}
The reduction map has finite image by the finite local residue
rings of Lemma~\ref{lem:local-residue-indicators}, proving
normality and finite index. Each \(g C g^{-1}\cap C\) is an
open neighborhood of one in \(C\). The congruences \(C_m\)
are a base there, since any such neighborhood contains given
congruences at finitely many primes and the ordinary integral
conditions elsewhere, and \(v_p(m!)\to\infty\) for each fixed
\(p\). This gives the inclusions simultaneously. They depend
only on \(gC\).

The construction of the charts of
Theorem~\ref{thm:analytic-framed-torus-family} applies also at
\(C_m\) by openness, finite index and the inherited trivial
point stabilizers. Changing \(x\) to \(xk\) for \(k\in C\)
induces an action on the finer classes by normality,
with the same lattice and reflection.
Every lift of one coarse class has a representative of
this kind, with its original torus points, by undoing a
coarse rational change. Two such lifts are equal only when
the cosets \(k C_m\) are the same: a rational equivalence
between them belongs to \(xCx^{-1}\) and fixes the reflection,
so it is the identity by Proposition~\ref{prop:neat-at-wild-place}.
The map is locally the identity in the disjoint period charts
of Theorem~\ref{thm:analytic-framed-torus-family}. More explicitly,
choose a neighborhood of the line meeting none of its other
translates under the coarse rational stabilizer of \(xC\).
At each possible \(xk C_m\) this gives a copy of the same
local chart, using the same lattice translates, and the
preceding fiber description exhausts its inverse image. This
proves the covering assertion.

For the second map one has \(\Lambda_x\subseteq\Lambda_{xg}\)
by \(g^{-1}\in M_n(\widehat O)\). Thus \(a_jP_j\) defines
a point mod \(\Lambda_{xg}\) since \(\Lambda_x\) is
\(O_E\)-stable. The form at \(xg\) is a positive rational
multiple of the form at \(x\): divide by the positive rational
with the valuations of the multiplier of \(g\), as in
Proposition~\ref{prop:definite-full-class-set}. All signatures,
the line and the complex structure are preserved. Replacing
\(x\) by \(xk\) for \(k\in C_m\) changes \(xg\) on the right
by \(g^{-1} k g\in C\); replacing \(g\) by \(gk'\) with
\(k'\in C\) also leaves its class fixed. Rational changes
commute with the scalar endomorphisms \(a_j\). This proves
well-definedness. Use the same period coordinates at \(x\)
and \(xg\), since scaling the form by a positive scalar
keeps its line and negative complement unchanged. On those
charts the base map is the identity, and the fiber maps are
the complex linear scalar actions in
Proposition~\ref{prop:signed-ball-period-coordinates}, followed
by reduction modulo the larger lattice. Hence they are holomorphic.
\end{proof}

\begin{lemma}\label{lem:fixed-torus-odd-determinant}
For a term of \(\mathcal O_1(H)\) with \(H\) regular and
any \(\alpha\in O_E\), the endomorphism
\(D_\alpha=I-2\alpha y\) preserves \(\Lambda_x\) by
inclusion. Its kernel on \(V_\R/\Lambda_x\) has order the
positive odd integer
\[
       d_\alpha(y)=\prod_{\tau:E\hookrightarrow\C}
                    \det(I-2\tau(\alpha)\tau(y)).
\]
For \(\boldsymbol\alpha=(\alpha_j)\in O_E^\nu\),
\(B_{\boldsymbol\alpha,y,i}(1)=\prod_j d_{\alpha_j}(y)\).
\end{lemma}
\begin{proof}
The matrix \(x^{-1}yx\) is integral because the inverse Hecke
argument is integral, so \(y\Lambda_x\subseteq\Lambda_x\).
In an integral basis the endomorphism \(D_\alpha\) reduces
to the identity modulo two. Its determinant as an integer
matrix is therefore odd, in particular nonzero. Over \(\R\)
it is invertible, and multiplication by \(D_\alpha\) identifies
its torus kernel \(D_\alpha^{-1}\Lambda_x/\Lambda_x\) with
\(\Lambda_x/D_\alpha\Lambda_x\). This has order the absolute
determinant by the lattice index argument of
Lemma~\ref{lem:class-covol}. The rational determinant is the
displayed product over embeddings by
Lemma~\ref{lem:marked-positive-discriminant}, and is positive
since the determinants occur in nonzero conjugate pairs.
At \(X=1\), every root in
Definition~\ref{def:actual-finite-local-terms} takes its
original value and \(L_i(z_i)=1\), so that product over all
embeddings is precisely the contribution from each \(j\).
\end{proof}

\begin{theorem}\label{thm:complex-hecke-coincidence-count}
Let \(H\in\mathcal H^+\) be regular, take \(D_0=\operatorname{supp}H\)
in Proposition~\ref{prop:complex-actual-hecke-legs}, and fix
\(\boldsymbol\alpha\in O_E^\nu\). The coincidences are the pairs
\[
  (gC,p),\qquad gC\in D_0,\quad p\in\mathcal X(C_m,\nu),
                 \quad u_{g,2\boldsymbol\alpha}(p)=\pi_m(p).
\]
They map onto \(\mathcal O_1(H)\). The fiber above
\([(xC,y,i)]\) has size
\([C:C_m]B_{\boldsymbol\alpha,y,i}(1)\), and
\(H(g)=H(x^{-1}y^{-1}x)\) throughout it. In particular
there are finitely many coincidences.
\end{theorem}
\begin{proof}
Write a finer point using \(x,J,P_j\). Equality of the two
coarse classes is equivalent to the existence of a rational
\(y\) such that
\[
 gC=x^{-1}y^{-1}x C,\qquad \sigma(y)J=J\sigma(y),\qquad
                   (I-2\alpha_j y)P_j=0 \pmod{\Lambda_x}
                        \quad\hbox{for every }j .
\]
Indeed \(y\) takes the frame \(xg\) to \(x\) modulo \(C\),
transports the reflection to itself and identifies the
displayed torus points, exactly the equivalence in
Definition~\ref{def:signed-complex-period-space}. Such \(y\)
is unique: for two choices \(y_1y_2^{-1}\in x Cx^{-1}\),
and it fixes \(J\), so is trivial by
Proposition~\ref{prop:neat-at-wild-place}. Regularity gives
distinct roots of \(\sigma y\). An involution commuting
with this matrix acts by a sign on each eigenline and
a reflection takes sign \(+1\) on exactly one. Thus
\(J=J_{1,i}(y)\) for a unique \(i\).
This defines the claimed orbit. Under rational change of
the finer point \(y\) conjugates accordingly, and under
right \(C_m\) change the same \(y\) works since
\(g^{-1}C_mg\subseteq C\).

Conversely fix \(x,y,i\) for an orbit. For each \(k C_m\)
in \(C/C_m\) use the finer frame \(xk\), reflection
\(J_{1,i}(y)\), and \(gC=k^{-1}x^{-1}y^{-1}xkC\).
This lies in \(D_0\) with its original weight by the two
\(C\)-invariances. Replacing \(k\) by \(kk'\) with
\(k'\in C_m\) keeps that coset: left multiplication on it
by \((k')^{-1}\) fixes it by the choice of \(C_m\) for
all \(D_0\). The fiber coordinates can be any element of
\(\prod_j D_{\alpha_j}^{-1}\Lambda_x/\Lambda_x\), since
\(\Lambda_{xk}=\Lambda_x\). These choices give coincidences,
and every one above the orbit has such a representative
by a rational change. If two choices give the same pair,
the rational equivalence between their points is in
\(x C x^{-1}\) and fixes the common reflection. It is
trivial by the same neatness, after which equality forces
equal \(k C_m\) and equal torus points. The finite count
is therefore as stated by Lemma~\ref{lem:fixed-torus-odd-determinant}
and Theorem~\ref{thm:all-marked-regular-finite}.
\end{proof}

\section{The tangent determinant at a Hecke coincidence}

\begin{theorem}\label{thm:actual-complex-tangent-determinant}
At a coincidence \((gC,p)\) of
Theorem~\ref{thm:complex-hecke-coincidence-count} put
\(u=u_{g,2\boldsymbol\alpha}\), \(z=\pi_m(p)=u(p)\) and
\(A_p=du_p(d\pi_{m,p})^{-1}\) on \(T_z\mathcal X(C,\nu)\).
Let \(x,y,i\) be the corresponding marked data. Then
\[
 \det(I-XA_p)=
       B_{\boldsymbol\alpha,y,i}(X)/B_{\boldsymbol\alpha,y,i}(0),
 \qquad \det(I-A_p)\ne0 .
\]
\end{theorem}
\begin{proof}
By Lemma~\ref{lem:marked-similitude-root-bounds} all the
eigenvalues \(r\) of all conjugates of \(y\) have
\(|r|^2=q\) for the same \(q>0\). At \(\sigma\) they
are distinct, and their eigenlines are pairwise orthogonal
for the rational signed form as in
Proposition~\ref{prop:marked-polynomial-positive-form}. Choose
eigenvectors of squared lengths \(+1,-1,\ldots,-1\) for that
form, the first for the marked root \(z_i\). Such rescaling
is possible since \(AJ_{1,i}\) is positive. The fixed line
then has ball coordinate zero of
Proposition~\ref{prop:signed-ball-period-coordinates}.
Comparing \(u\) to \(\pi_m\) in the target chart entails
the rational coordinate change \(y\) taking \(xg\) to \(x\)
modulo \(C\). On the positive line it sends \((1,\bar b)\)
to \((z_i,(z_j\bar b_j)_{j\ne i})\). Taking the conjugates
of the normalized coordinates gives diagonal action on
the base with eigenvalues \(\bar z_j/\bar z_i=z_i/z_j\)
for \(j\ne i\). The last equality uses the common squared
absolute value.

The derivative preserves the vertical tangent space by the
base map calculation. On its \(a\)-th summand it is induced
by \(2\alpha_a y\) in the period coordinates. At zero
\(Q_0(w)=(w_0,\bar w')\) at \(\sigma\), so on this factor
the eigenvalues are \(2\sigma(\alpha_a)z_i\) on the line and
\(2\bar\sigma(\alpha_a)\bar z_j\) for \(j\ne i\) on the
complement. At each \(\tau\in\Psi\setminus\{\sigma\}\) they
are all the \(2\tau(\alpha_a)r\) for roots \(r\) of \(\tau P\).
This uses the maps of Proposition~\ref{prop:complex-actual-hecke-legs}
and their local charts. Changing a torus lift by a lattice
vector adds the holomorphic section \(Q_b(\lambda)\), whose
vertical derivative is the identity for the resulting
translation. Its derivatives in the base do not affect the
actions on the vertical space or on its quotient. The
determinant thus multiplies the base factor
\(\prod_{j\ne i}(1-X z_i/z_j)\) and the factors
\(1-X(2\tau(\alpha_a)r)\) over exactly the listed fiber roots,
by the determinant of a block triangular matrix.

In Definition~\ref{def:actual-finite-local-terms} these
fiber roots occur exactly in the lists multiplied by \(X\):
outside \(v_0\) the representative of \(\Psi\) is
\(\bar\nu_v\) by Lemma~\ref{lem:actual-signs-and-padding}.
At the base,
\[
          \frac{L_i(Xz_i)}{L_i(0)}
                         =\prod_{j\ne i}(1-X z_i/z_j).
\]
Here \(L_i(0)\ne0\) since \(y\) is invertible and its roots
are distinct. Each constant factor from the unmodified
roots is nonzero because \(B_{\boldsymbol\alpha,y,i}(1)>0\)
by Lemma~\ref{lem:fixed-torus-odd-determinant}. Consequently
\(B_{\boldsymbol\alpha,y,i}(0)\ne0\), and after division by
this constant the factors give the asserted polynomial.
At \(X=1\) the numerator is the same positive integer.
\end{proof}

\begin{corollary}\label{cor:actual-holomorphic-local-sum}
Let \(H\) be regular, and use the finite set of correspondences
and the level in Theorem~\ref{thm:complex-hecke-coincidence-count}.
For \(h\ge0\), let \(A_p^\vee\) denote the ordinary linear
dual of \(A_p\), acting by precomposition on covectors.
Set \(\bigwedge^h A_p^\vee=0\) when \(h\) exceeds the tangent
dimension. Then
\[
 \frac{1}{[C:C_m]}
   \sum_{gC\in\operatorname{supp}H} H(g)
      \sum_{p:\,u_{g,2\boldsymbol\alpha}(p)=\pi_m(p)}
      \frac{(-1)^h\,{\rm tr}(\bigwedge^h A_p^\vee)}
           {\det(I-A_p)}
                   = M_{\boldsymbol\alpha,h}(H).
\]
In particular for every \(F\) normalized as in
Theorem~\ref{thm:actual-marked-weight-occurrence} there is
one regular left multiple \(H\) of \(F(\tau)\) for which the
displayed complex sum is nonzero at each specified \(h_j\)
for some \(\boldsymbol\alpha\in\mathcal B_n^\nu\).
\end{corollary}
\begin{proof}
For a complex endomorphism \(A\), triangularize it by induction
using an eigenvector as in
Corollary~\ref{cor:definite-orbit-sum-detects-jump}.
The diagonal entries of its dual are the same eigenvalues.
On a basis of exterior monomials of degree \(h\) the
diagonal entries are the products over \(h\)-element subsets.
Thus \([X^h]\det(I-XA)=(-1)^h{\rm tr}(\bigwedge^h A^\vee)\),
also at \(h=0\) with the empty product one. Theorem~\ref{thm:actual-complex-tangent-determinant}
therefore makes each fraction
\([X^h]B_{\boldsymbol\alpha,y,i}(X)/B_{\boldsymbol\alpha,y,i}(1)\).
On the fiber of a marked orbit the number of terms is precisely
\([C:C_m]B_{\boldsymbol\alpha,y,i}(1)\) by
Theorem~\ref{thm:complex-hecke-coincidence-count}.
The common Hecke value and finiteness in that theorem now give
the equality after summation over the marked orbits, by
Definition~\ref{def:marked-weight-test}. Its nonvanishing in
the stated applications is Theorem~\ref{thm:actual-marked-weight-occurrence}.
\end{proof}

\section{Theta coordinates for the signed complex fibers}

\begin{lemma}\label{lem:signed-period-elementary-divisors}
Put \(g=fn\), \(L=\Lambda_x\) and \(E_0=E_A\) from
Proposition~\ref{prop:signed-period-riemann-form}.
There is a lattice basis \(e_1,\ldots,e_g,f_1,\ldots,f_g\)
with
\[
  E_0(e_j,e_k)=E_0(f_j,f_k)=0,\qquad
  E_0(e_j,f_k)=d_j{\bf1}_{j=k},\qquad
  0<d_1\mid d_2\mid\cdots\mid d_g .
\]
Each \(d_j\) is a multiple of \(R^2\) from
Theorem~\ref{thm:actual-pel-torsion-pairing};
\(\prod_j d_j=h^n\).
\end{lemma}
\begin{proof}
The form is nonsingular by its positivity on \((v,I_Jv)\).
Among its positive integer values on pairs \(u,v\) in the
lattice choose the smallest \(d=E_0(u,v)\). For every \(w\),
both \(E_0(u,w)\) and \(E_0(w,v)\) are divisible by \(d\):
subtract integer multiples of \(v\) or \(u\) respectively
and apply division with remainder and minimality.
Thus the integral projection
\[
   w\longmapsto d^{-1}\big(E_0(w,v)u+E_0(u,w)v\big)
\]
splits off \(\Z u+\Z v\) orthogonally. The complementary
lattice is free by Lemma~\ref{lem:rd} and its form remains
nonsingular over \(\Q\). On it each pairing \(E_0(w,w')\)
is divisible by \(d\): replacing the second member of
\(E_0(u+w,w')\) by \(w'+kv\) adds \(kd\) to that pairing.
Its remainder cannot lie strictly between zero and \(d\).
Inducting gives all pairs and the stated divisibility chain.
The form \(E_0/R^2\) is integral since \(\xi/R^2=u/\varpi\)
with \(u\in O_F\), by the local trace-dual argument of
Proposition~\ref{prop:signed-period-riemann-form}. This
gives the common divisibility. In the paired basis the
dual modulo the lattice is the product of two copies of
\(\Z/d_j\Z\) per pair. The order \(h^{2n}\) in that
proposition then gives \(\prod_jd_j=h^n\).
\end{proof}

\begin{proposition}\label{prop:signed-normalized-period-matrix}
Put \(D=\operatorname{diag}(d_1,\ldots,d_g)\) and use the
basis of Lemma~\ref{lem:signed-period-elementary-divisors}.
At each point \(b\) of the period ball, in the complex basis
\(Q_b(e_1),\ldots,Q_b(e_g)\) the lattice has the form
\[
                \Z^g+\Omega(b)D\Z^g,
 \qquad \Omega(b)^{\mathsf T}=\Omega(b),\qquad
                   \Im\Omega(b)>0 .
\]
The matrix \(\Omega(b)\) is holomorphic on the ball.
\end{proposition}
\begin{proof}
For the real span \(U\) of the \(e_j\), if \(v,I_Jv\in U\),
then \(E_0(v,I_Jv)=0\) by isotropy, giving \(v=0\) by
Proposition~\ref{prop:signed-period-riemann-form}. Thus
\(U\) and \(I_JU\) span a direct sum of dimension \(2g\).
This proves the complex basis assertion through \(Q_b\) of
Proposition~\ref{prop:signed-ball-period-coordinates}.
Define the matrix by
\(\Omega=(Q_b(e_1),\ldots,Q_b(e_g))^{-1}
          (Q_b(f_1),\ldots,Q_b(f_g))D^{-1}\).
The entries are holomorphic by the period formula and inverse
matrix cofactors with nonvanishing determinant.

The Hermitian form in the complex structure of this fiber is
\(H(v,w)=E_0(v,I_Jw)+iE_0(v,w)\). It is conjugate linear
in the first variable and linear in the second by alternation
and \(I_J\)-invariance, and is positive by the preceding
proposition on the Riemann form. Its matrix \(M\) on the
\(e_j\) is real symmetric positive definite. Write
\(\Omega=X+iY\). The equations \(E_0(e_j,f_k/d_k)={\bf1}_{j=k}\)
give \(MY=I\). Thus \(Y=M^{-1}\) is positive definite and
symmetric. The imaginary part of \(\bar\Omega^{\mathsf T}M\Omega\)
vanishes since all the \(f_k/d_k\) pair to zero.
Consequently \(X^{\mathsf T}MY-Y^{\mathsf T}MX=0\),
giving \(X^{\mathsf T}=X\).
\end{proof}

\begin{theorem}\label{thm:signed-theta-sections}
For \(k\ge1\), define \(\Theta_k(\Omega,D)\) to be the entire
functions \(S\) on \(\C^g\), periodic under \(\Z^g\), with
\[
 S(z+\Omega Dv)=
 \exp\!\big(-\pi i k(Dv)^{\mathsf T}\Omega Dv
                 -2\pi i k(Dv)^{\mathsf T}z\big) S(z)
                       \qquad(v\in\Z^g).
\]
These are the sections of a holomorphic line bundle, the
\(k\)-th tensor power of the one defined by the \(k=1\)
factors. A basis indexed by \(r\in\Z^g/kD\Z^g\) is
\[
 \Theta_{k,r}(z)=
  \sum_{a\in r+kD\Z^g}
       \exp\!\big((\pi i/k)a^{\mathsf T}\Omega a+2\pi i a^{\mathsf T}z\big).
\]
In particular the dimension is \(k^g h^n\). These sections
vary holomorphically in \(b,z\) on the period ball.
\end{theorem}
\begin{proof}
The factors are multiplicative for two successive lattice
translations: the cross term is determined by symmetry of
\(\Omega\), and translations by \(m\in\Z^g\) contribute
\(\exp(-2\pi i k(Dv)^{\mathsf T}m)=1\). They thus give a
lattice action on \(\C^g\times\C\) by scaling the fiber,
forming a line bundle by the disjoint local lattice charts
of Theorem~\ref{thm:analytic-framed-torus-family}. Its
sections lift to exactly the entire functions with these
transformation laws. The factors in degree \(k\) are the
\(k\)-th powers of the degree-one factors.

For a compact set of \(b,z\), the least value on real unit
vectors of \(\Im\Omega(b)\) has a positive uniform lower
bound by Proposition~\ref{prop:signed-normalized-period-matrix}.
The absolute values of the summands are bounded by
\(\exp(-c\|a\|^2+C\|a\|)\), \(c>0\), summable by comparison
with a product of one-dimensional geometric tails. The same
bound with a polynomial factor applies to derivatives on
smaller compact sets. Termwise differentiation, or the
Cauchy formula on coordinate discs of
Lemma~\ref{lem:analytic-tools}, proves holomorphy.
Translation by \(\Omega Dv\) reindexes \(a\) by \(a+kDv\)
and gives the displayed factor. The sections for different
classes \(r\) have disjoint Fourier supports and nonzero
coefficients, so are linearly independent.

Every periodic entire function has a unique normally
convergent Fourier expansion \(\sum_{a\in\Z^g}c_a
\exp(2\pi i a^{\mathsf T} z)\): the contour shift and
product Fourier argument in Theorem~\ref{thm:period-dimension}
uses just periodicity and holomorphy. Substitution in the
period equations gives
\[
 c_{a+kDv}=
   \exp\!\big(2\pi i a^{\mathsf T}\Omega Dv+
                   \pi i k(Dv)^{\mathsf T}\Omega Dv\big)c_a .
\]
Thus the coefficients on each class are the coefficients of
\(\Theta_{k,r}\) times one constant. The functions form a
basis and \( |\Z^g/kD\Z^g|=k^g\prod_jd_j=k^g h^n\) by
Lemma~\ref{lem:signed-period-elementary-divisors}.
\end{proof}

\section{Separating points and tangents by translated theta products}

\begin{proposition}\label{prop:products-translated-theta}
Use any positive diagonal integral \(D\) and symmetric
\(\Omega\) with \(\Im\Omega>0\) in the construction of
Theorem~\ref{thm:signed-theta-sections}. For \(k\ge3\) the
degree-\(k\) theta sections span a system without common
zeros, separating points and nonzero tangent vectors on
\(\C^g/(\Z^g+\Omega D\Z^g)\). More precisely these properties
already hold for the span of
\[
        \prod_{j=1}^k S_j(z+t_j),\qquad
         S_j\in\Theta_1(\Omega,D),\quad t_j\in\C^g,\quad
                                        \sum_jt_j=0 .
\]
\end{proposition}
\begin{proof}
Zero total shift makes the product's automorphy factors exactly
the degree-\(k\) factors by Theorem~\ref{thm:signed-theta-sections}.
Choose one nonzero section \(S\) there. At fixed \(z\) each
\(S(z+t_j)\) is a nonzero entire function of the \(k-1\)
independent shifts, also for the last shift
\(t_k=-\sum_{j<k}t_j\), since each shift ranges over
\(\C^g\). Its nonzero set is open dense. Indeed vanishing
on a ball gives vanishing along each complex line through
the ball by the one-variable identity theorem of
Lemma~\ref{lem:analytic-tools}. Finitely many such open
dense sets have nonempty intersection. This gives a product
nonzero at \(z\).

Suppose \(z,w\in\C^g\) have proportional evaluations on all
these products. Both evaluation vectors are nonzero, so
the constant of proportionality is nonzero. Taking all
\(S_j=S\) gives, wherever the ratios are defined,
\[
       \prod_{j=1}^k f(t_j)=c,\qquad
                    f(t)=S(w+t)/S(z+t).
\]
Differentiate with respect to each coordinate of \(t_1\),
holding \(t_2,\ldots,t_{k-1}\) fixed. Logarithmic differentiation
gives \((\partial_a f/f)(t_1)=(\partial_a f/f)(t_k)\).
The two arguments \(t_1,t_k\) vary independently because
\(k\ge3\). This identity of meromorphic functions follows
everywhere after clearing denominators from the open dense
set on which the original products are nonzero. Consequently
each \(\partial_a f/f\) is constant: fix a generic value
of its second argument and vary the first. It follows that
\(f(t)=C\exp(\ell(t))\) for a complex linear functional
\(\ell\) and \(C\ne0\). Indeed multiplication by
\(\exp(-\ell(t))\) gives zero derivative on a nonempty open
ball of holomorphy; constancy there extends by clearing
denominators and the identity theorem. Periodicity under
\(\Z^g\) gives \(\ell(t)=2\pi i m^{\mathsf T}t\) with
\(m\in\Z^g\).

Replacing only the first section in the proportionality
relation by another \(S'\in\Theta_1\) shows, by cancelling
the other nonzero factors on an open set, that
\(S'(w+t)S(z+t)=S'(z+t)S(w+t)\). This entire identity
holds for all \(t\). Hence with \(\eta=w-z\), there is
a constant \(C_1\ne0\), independent of \(S'\), such that
\[
        S'(v+\eta)=C_1\exp(2\pi i m^{\mathsf T}v)\,S'(v)
                          \qquad(v\in\C^g).
\]
Use as \(S'\) each basis section in degree one. Its Fourier
support is a whole class modulo \(D\Z^g\) by
Theorem~\ref{thm:signed-theta-sections}. The shifted support
on the right must agree with the original one, so
\(m\in D\Z^g\). Comparing the coefficient of frequency \(a\)
for all such classes yields
\[
 \exp(2\pi i a^{\mathsf T}\eta)=C_1
       \exp\!\big(\pi i(m^{\mathsf T}\Omega m-2a^{\mathsf T}\Omega m)\big)
                                  \qquad(a\in\Z^g).
\]
Taking \(a=0\) and then the coordinate vectors gives
\(\eta+\Omega m\in\Z^g\). Thus \(z,w\) are the same torus
point, proving point separation.

If a vector \(v\) at \(z\) has zero projective derivative
there is a constant \(c'\) with \(\partial_v P(z)=c'P(z)\)
for all products \(P\). This follows by differentiating ratios
to one nonvanishing product at \(z\). With \(S=\Theta_{1,0}\)
in every factor, put \(u(t)=(\partial_v S)(z+t)/S(z+t)\).
The product rule gives \(\sum_j u(t_j)=c'\) on the same
open dense set. Differentiating in \(t_1\) now gives
\((\partial_a u)(t_1)=(\partial_a u)(t_k)\).
By the preceding independent-arguments reasoning each
partial derivative is constant, so \(u\) is affine linear.
Its periodicity under \(\Z^g\) makes it constant. In
\(\partial_v S=c'' S\), comparison of Fourier coefficients
gives \(2\pi i a^{\mathsf T}v=c''\) for all \(a\in D\Z^g\).
Taking zero and the columns of \(D\) gives \(v=0\).
This proves tangent separation.
\end{proof}

\begin{theorem}\label{thm:signed-fiber-projective-embedding}
For the actual \(D,\Omega(b)\) of
Proposition~\ref{prop:signed-normalized-period-matrix}, the
degree-one coordinates of Theorem~\ref{thm:signed-theta-sections}
embed every fiber in \(\mathbf P^{h^n-1}(\C)\).
Over a fixed period ball they give a closed holomorphic
embedding of the family of one torus into
\(\mathbb B\times\mathbf P^{h^n-1}(\C)\).
\end{theorem}
\begin{proof}
By Lemma~\ref{lem:signed-period-elementary-divisors},
\(D'=D/R^2\) has positive integral diagonal. Put
\(\Omega'=R^2\Omega\), so \(\Omega'D'=\Omega D\) and
\(\Im\Omega'>0\). The automorphy factors for degree \(R^2\)
of \((\Omega',D')\) are exactly those for degree one of
\((\Omega,D)\): substitute \(D',\Omega'\) into the quadratic
and linear terms of Theorem~\ref{thm:signed-theta-sections}.
Since \(R^2\ge3\), Proposition~\ref{prop:products-translated-theta}
gives all three separation properties for these actual
degree-one sections.

Their coordinates depend holomorphically on the base and
on the complex torus parameter by Theorem~\ref{thm:signed-theta-sections}.
The coordinate change from \(Q_b(w)\) to \(z\) in the
basis \(Q_b(e_j)\) is holomorphic in \(b\) by
Proposition~\ref{prop:signed-normalized-period-matrix}.
The common scalar factor on a lattice shift leaves the
projective point unchanged. With the base recorded, fiber
separation gives injectivity. A tangent vector in the
kernel must have zero base component and is zero by tangent
separation. The map is proper: over any compact set of the
base the torus family is compact, since the base set times
a fixed closed real fundamental parallelotope of \(\Lambda_x\)
maps continuously onto that inverse image, as in
Theorem~\ref{thm:compact-signed-complex-family}. The inverse
image of a compact set in the target is closed in such a
compact set. This turns the continuous injection into a
homeomorphism onto its closed image: for any convergent
sequence of image points, its values and its limit form a
compact set in the target; the preimages have a convergent
subsequence, and injectivity fixes the possible limit
whenever the target limit is in the image.

In complex coordinates select a square submatrix of the
derivative of full rank; the corresponding target coordinates
are holomorphic local coordinates on the source. Indeed
after linear normalization their map is \(x\mapsto x+e(x)\)
with \(e(0)=0\) and derivative zero at zero. On a small
closed convex ball its derivative norm is less than \(1/2\),
so \(x\mapsto t-e(x)\) takes that ball into itself for small
\(t\) and contracts distances by \(1/2\). Iteration from zero
gives a locally uniformly convergent sequence of holomorphic
functions of \(t\) and hence a holomorphic inverse by the
Cauchy formula of Lemma~\ref{lem:analytic-tools}. In these
coordinates the remaining components give a holomorphic
graph. This proves the embedding assertion.
\end{proof}

\section{Finite polynomial projections of coordinate rings}

\begin{lemma}\label{lem:complex-polynomial-normalization}
Let \(R=\C[r_1,\ldots,r_N]\) be a unital domain and let
\(R_{\le m}\) be the linear span of products of at most \(m\)
of the \(r_j\), including one. There are linear combinations
\(t_1,\ldots,t_d\) of the generators that are algebraically
independent, meaning they satisfy no nonzero complex polynomial
relation, such that \(R\) is a finite module over
\(A=\C[t_1,\ldots,t_d]\). For some positive constants \(c,C\)
one has
\[
            c(m+1)^d\le \dim_\C R_{\le m}\le C(m+1)^d
                         \qquad(m\ge0).
\]
Each \(r_j\) satisfies a monic polynomial over \(A\).
\end{lemma}
\begin{proof}
If the current list is dependent, take a nonzero relation of
total degree \(e\ge1\). Choose constants \(c_1,\ldots,c_{N-1}\)
where its degree-\(e\) homogeneous part evaluated at
\((c_1,\ldots,c_{N-1},1)\) is nonzero. Such constants exist:
dehomogenizing a nonzero homogeneous polynomial keeps distinct
monomials distinct, and a nonzero polynomial over the infinite
field \(\C\) cannot vanish on every tuple, by induction using
univariate division. Replace the first \(N-1\) generators by
\(u_j=r_j-c_j r_N\). Substituting in the relation and dividing
by that nonzero scalar gives a degree-\(e\) monic equation
for \(r_N\) over the algebra of the \(u_j\). Hence powers
through \(e-1\) span \(R\) over the smaller algebra.
Repeat on its list until independent or empty. Finite
module generation passes down the tower by multiplying the
two finite lists of generators at each step. This constructs
\(A\). For any one endomorphism of the finite generating
module given by multiplication by \(r_j\), write its matrix
on a generating list using coefficients in \(A\). Multiplying
the resulting equations by the adjugate of \(r_j I\) minus
that matrix shows its monic characteristic polynomial kills
all generators, hence kills \(1\). This is the claimed equation.

The monomials of degree at most \(m\) in \(t\) are independent
and lie in \(R_{\le m}\). For the upper bound choose finite
module generators including one, and express \(r_j\) times
each generator in those generators with polynomial coefficients
in \(t\) of degrees at most a common integer \(b\ge1\). Every
word of length at most \(m\) applied to \(1\) then uses
coefficients of degrees at most \(bm\), by induction. There
are at most \((bm+1)^d\) monomials in that range and at least
\((\lfloor m/d\rfloor+1)^d\) independent monomials in the
lower bound if \(d>0\), by bounding each exponent separately.
When \(d=0\) the entire ring is finite dimensional and
contains one. This proves both inequalities.
\end{proof}

\begin{lemma}\label{lem:polynomial-growth-punctured-space}
Let \(b\) be a nonzero polynomial on \(\C^d\), \(d\ge1\).
A holomorphic function on \(U_b=\{t:b(t)\ne0\}\) bounded
there by \(C(1+\|t\|)^m\), \(m\) a nonnegative integer,
extends uniquely to a polynomial of degree at most \(m\).
Also \(U_b\) is connected.
\end{lemma}
\begin{proof}
Any two distinct points of \(U_b\) lie on an affine complex
line where \(b\) has finitely many zeros, since its values
at the endpoints are nonzero. In that line one can join
the two points by a polygonal path avoiding the finite set,
proving connectedness. At a point outside \(U_b\) choose
one complex line direction where the polynomial on the line
is not identically zero, by joining it to a point of \(U_b\).
Use coordinates \((v,z)\) in that direction, with the given
point at the origin. On some circle \(|z|=\rho>0\) at \(v=0\)
the polynomial has no zero. The same holds uniformly for
all sufficiently small \(v\) by continuity and compactness.
Use the Cauchy integral of the function on this circle
with kernel \(1/(2\pi i(\zeta-z))\) for \(|z|<\rho\).
It is holomorphic in all the variables by the local Cauchy
formula of Lemma~\ref{lem:analytic-tools}, applied successively
in the coordinate variables on smaller closed polydiscs.
For fixed \(v\) the original function is bounded near its
finitely many missing points in the disc. A bounded holomorphic
function at a puncture extends: multiply it by \((z-a)^2\),
extend by zero at \(a\), differentiate there by the difference
quotient to obtain zero, and divide its Taylor series back.
Thus the integral agrees with the original function wherever
defined. These local extensions glue by their agreement on
the dense open complement and continuity. The global polynomial
growth bound also extends by continuity. On polydiscs of
radius \(T\) around zero the repeated Cauchy coefficient
bound is \(C(1+\sqrt d\,T)^m T^{-|\beta|}\) for a coefficient
of multi-index \(\beta\). Letting \(T\) tend to infinity kills
all coefficients of total degree greater than \(m\), proving
the conclusion and uniqueness.
\end{proof}

\begin{proposition}\label{prop:finite-projection-simple-sheets}
Use \(R,A,t\) from Lemma~\ref{lem:complex-polynomial-normalization}
with \(d\ge1\), and let \(Z\subset\C^N\) be the common zero
set of the polynomial relations among the \(r_j\). Evaluation
gives \(Z=\operatorname{Hom}_{\C\text{-alg}}(R,\C)\) in the
coordinates \(r_j\). The map \(t:Z\to\C^d\) has finite fibers
and bounded inverse image over each bounded subset. Given
a particular fiber, there is a complex linear combination
\(l\) of the \(r_j\) separating its distinct points, a monic
\(P(t,T)\in A[T]\) irreducible over \(\C(t_1,\ldots,t_d)\),
and a nonzero \(b\in A\) such that
\[
  Z|_{U_b}\ \longrightarrow\ \{(t,\lambda):t\in U_b,\
                P(t,\lambda)=0\},\qquad z\longmapsto(t(z),l(z))
\]
is bijective with coordinates of its inverse polynomial in
\(\lambda,t,1/b(t)\). The roots on \(U_b\) are all simple and
give local holomorphic graphs. One has \(P(t(z),l(z))=0\)
on all of \(Z\).
\end{proposition}
\begin{proof}
The evaluation description follows from the presentation by
generators and relations. The monic equations of
Lemma~\ref{lem:complex-polynomial-normalization} give finitely
many choices for every coordinate over fixed \(t\).
Roots of \(T^s+\sum_{j<s}a_jT^j\) have magnitude at most
\(1+\max_j|a_j|\), by dividing by \(|T|^s\) and summing a
geometric series if this bound is exceeded. Thus coordinates
remain bounded on bounded \(t\)-sets, and individually satisfy
a polynomial growth bound there and at infinity.

Put \(k=\C(t_1,\ldots,t_d)\). Inverting nonzero elements of
\(A\) in \(R\) gives a finite-dimensional \(k\)-domain, hence
a field \(L\): multiplication by a nonzero member is injective
and therefore surjective. In characteristic zero its finite
extensions are separable by the polynomial derivative.
There are \([L:k]\) embeddings into an algebraic closure,
by extending over each of the finitely many generators and
counting the distinct roots of successive minimal polynomials,
as in Lemma~\ref{lem:number-field-lattice}. Choose a constant
linear combination \(l\) of the coordinate generators taking
different values under each pair of embeddings and at each
pair of distinct points of the given finite fiber. Each
equality cuts out a proper linear condition: the generators
separate the two homomorphisms in question. Finitely many
such conditions do not exhaust the constant tuples, by the
polynomial nonvanishing argument in the normalization lemma,
also for polynomials with coefficients in a larger field.
Then \(L=k(l)\) since the minimal polynomial of \(l\) has
at least \([L:k]\) roots. Denote that monic polynomial by \(P\).

Express every \(r_j\) as a polynomial of degree below
\([L:k]\) in \(l\) over \(k\). Invert one nonzero \(b\in A\)
clearing the denominators of these expressions, of \(P\),
and of a Bezout identity between \(P\) and its derivative.
The resulting presentation is
\(R[1/b]=A[1/b][T]/(P)\): evaluation at \(l\) is onto by the
expressions, and monic division followed by independence
of the lower powers of \(l\) gives its kernel. Over \(U_b\)
this is exactly the asserted bijection with simple roots.
The holomorphic inverse function argument of
Theorem~\ref{thm:signed-fiber-projective-embedding}, applied
to \((t,T)\mapsto(t,P(t,T))\), makes each simple root a local
holomorphic graph. All roots meet the polynomial growth bound
in \(t\) from the coordinate bound, because \(l\) is their
fixed linear combination of coordinates of points of \(Z\).
Thus the coefficients of \(P\), the elementary symmetric
functions of its roots, are polynomially bounded on \(U_b\).
They are regular there by construction. Applying
Lemma~\ref{lem:polynomial-growth-punctured-space} to each gives
polynomials in \(A\), equal as rational functions to their
original values. The equation \(P(t,l)=0\) holds in \(L\)
and hence in its subring \(R\), so also at every point of \(Z\).
\end{proof}

\section{Projective algebraicity from the growth of sections}

\begin{theorem}\label{thm:section-growth-algebraic-image}
Let a connected compact complex manifold \(X\) of dimension
\(d\ge1\) be holomorphically embedded in \(\mathbf P^N(\C)\)
by sections \(S_0,\ldots,S_N\) of a holomorphic line bundle.
Suppose, for every \(m\ge1\), the span of the products of
\(m\) of these sections has dimension at most \(C(m+1)^d\)
for a fixed \(C\). Then the image is exactly the common
zero set of a finite list of homogeneous complex polynomials.
Its homogeneous ideal of relations is prime.
\end{theorem}
\begin{proof}
In an affine chart \(S_j\ne0\) which meets \(X\), use
ratios \(r_i=S_i/S_j\), \(i\ne j\), and set \(R=\C[(r_i)_i]\)
as an algebra of functions there. A section of a holomorphic
line bundle that vanishes on an open set of \(X\) vanishes
everywhere: in trivializing coordinates use the identity
theorem on complex lines as in
Proposition~\ref{prop:products-translated-theta}, continuing
over intersecting coordinate balls by connectedness.
Thus the chart is dense in \(X\), and \(R\) is a domain:
if the product of two ratio polynomials is zero, multiply
by a power of \(S_j\) to obtain an identity of products of
two global sections. Where one is nonzero the other vanishes
on an open set. Multiplication by \(S_j^m\) identifies
\(R_{\le m}\) with the span of homogeneous degree-\(m\)
coordinate sections, since every monomial can be padded by
\(S_j\). On the other hand, at a point in this chart some
\(d\) of the ratios have independent differentials by the
embedding. Their joint values contain a complex open set by
the local inverse proof of
Theorem~\ref{thm:signed-fiber-projective-embedding}.
No nonzero polynomial can vanish there, so their monomials
of degree at most \(m\) give a lower bound of order
\((m+1)^d\). Lemma~\ref{lem:complex-polynomial-normalization}
therefore gives exactly \(d\) independent base coordinates
\(t\).

Let \(Y\subseteq\C^N\) be this affine part of the image, and
\(Z\) the common zero set of its affine polynomial relations.
The set \(Y\) is closed in the affine chart since \(X\) is
compact and projective space is Hausdorff. We prove \(Z=Y\).
Fix a point \(z_0\in Z\). Apply
Proposition~\ref{prop:finite-projection-simple-sheets} with
its fiber \(t=t(z_0)\) distinguished. Use the resulting data
\(l,P,b\). Over \(U_b\), \(Z\) is locally the union of
disjoint graphs over an open ball, one for each simple root
of \(P\), with holomorphic coordinates in both directions.
The intersection with \(Y\) is relatively closed and open
there. Indeed \(Y\) is an embedded manifold of the same
dimension as each graph, so at any common point its inclusion
into the graph has full rank and is open by the same local
inverse argument. There is a point of \(Y\) over \(U_b\)
because \(b(t)\ne0\) in the algebra of functions \(R\).

For each \(t\in U_b\), multiply \(T-\lambda\) over the
sheets at \(t\) lying in \(Y\), with \(\lambda\) their
\(l\)-coordinates. The sheets admitted are locally constant
in every evenly covered ball, by the preceding open and
closed property. Their number is positive and constant
since \(U_b\) is connected by
Lemma~\ref{lem:polynomial-growth-punctured-space}. The
coefficients of this monic polynomial are thus holomorphic
on \(U_b\). They have polynomial growth since all the roots
have the bound in Proposition~\ref{prop:finite-projection-simple-sheets}.
The same two facts hold for the product over the omitted
sheets, with an empty product equal to one. Both products
therefore have polynomial coefficients in \(t\) by
Lemma~\ref{lem:polynomial-growth-punctured-space}. Their
product equals \(P\) on \(U_b\) and hence over \(\C(t)\).
Irreducibility of \(P\) forces the first, nonconstant
factor in \(T\), to have the entire degree. Every point of
\(Z\) over \(U_b\) belongs to \(Y\).

To include \(z_0\), take \(t_j\in U_b\) tending to
\(t(z_0)\), possible since the zero set of a nonzero polynomial
has no interior. One can choose roots \(\lambda_j\) of
\(P(t_j,T)\) tending to \(l(z_0)\). In fact this target is
a root at the limit by the preceding proposition. If along
a subsequence all roots stayed at positive distance from
it, their ordered tuples have a further convergent subsequence
by boundedness of roots of a monic polynomial with bounded
coefficients. Multiplication of the limiting linear factors
would give \(P(t(z_0),T)\) with no such root. Over \(U_b\)
each chosen root gives a point \(z_j\in Y\). Its coordinates
are bounded by Proposition~\ref{prop:finite-projection-simple-sheets}.
Any subsequential limit is in \(Z\) by its polynomial equations
and has \(t=t(z_0)\), \(l=l(z_0)\). The choice of \(l\)
separates points of that fiber, so the limit is \(z_0\).
Closedness of \(Y\) in the affine chart gives \(z_0\in Y\).

Any affine relation on \(Y\) homogenizes using \(S_j\) to a
relation on all of \(X\) by the identity theorem above.
Conversely every homogeneous relation restricts to an affine
one. Charts with \(S_j\) identically zero are excluded by that
coordinate's own homogeneous relation. This proves the asserted
description of the projective zero set. To see primeness, form
the graded ring of the coordinate sections with their degrees.
A product of nonzero homogeneous sections is nonzero by the
identity theorem. The highest nonzero homogeneous terms of
two finite sums therefore have nonzero product as well.
The homogeneous relation ideal, the kernel of the map to this
graded ring, is prime.

Finally order monomials by total degree and then
lexicographically, taking the largest monomial of each
nonzero homogeneous relation. The exponents so obtained
are closed under addition by nonnegative vectors. By the
finite generation argument of Lemma~\ref{lem:finite-degree-cut},
finitely many of them generate this set; choose homogeneous
relations with those leading monomials, of coefficient one.
From any homogeneous relation subtract an appropriate scalar
and monomial multiple of one of the chosen relations to cancel
its leading term. Each step decreases the leading monomial
among finitely many of that degree, so it ends at zero.
Thus the finite list generates the homogeneous ideal.
\end{proof}

\begin{corollary}\label{cor:actual-projective-algebraic-tori}
Each theta-embedded complex fiber of
Theorem~\ref{thm:signed-fiber-projective-embedding} is
the zero set of finitely many homogeneous polynomials with
complex coefficients and is irreducible. In each of its
nonempty affine coordinate charts its coordinate ring is
a finite algebra over a polynomial ring in \(g=fn\) variables.
\end{corollary}
\begin{proof}
The torus is connected as the continuous image of its complex
vector space and is compact by
Theorem~\ref{thm:analytic-framed-torus-family}. A product
of \(m\) coordinate sections belongs to the degree-\(m\)
theta space by multiplication of their automorphy factors.
That space has dimension \(m^g h^n\) by
Theorem~\ref{thm:signed-theta-sections}. Hence
Theorem~\ref{thm:section-growth-algebraic-image} applies and
gives both the finite list and the finite projection in
each chart. For irreducibility, if two proper projective
algebraic subsets covered the image, choose a homogeneous
polynomial for each subset which is nonzero at some point
of the image and vanishes on that subset. Their product
would vanish identically, contradicting primeness.
\end{proof}

\section{Finite jets and the algebraic graphs of the torus operations}

\begin{lemma}\label{lem:compact-section-jet-bound}
Let \(X\) be a compact complex manifold of dimension \(d\ge1\)
and \(L\) a holomorphic line bundle on it. There are finitely
many points \(x_1,\ldots,x_s\), trivializations near those points,
and \(A\ge1\), such that every holomorphic section of
\(L^{\otimes m}\), \(m\ge1\), is determined by its Taylor
coefficients of degrees below \(1+\lfloor Am\rfloor\) at
these points. In particular their space has dimension at most
\(C(m+1)^d\) for some \(C\) independent of \(m\).
\end{lemma}
\begin{proof}
Take polydisc coordinates and nowhere vanishing local sections
of \(L\) whose domains contain the closed polydisc of radius
three about each center, after rescaling. By compactness choose
finitely many for which the open polydiscs of radius one cover.
On each such domain take a continuous real function which is
one on radius one and zero outside radius two, linear in the
maximum coordinate radius between those radii. Extend it
by zero to \(X\). Use these functions to form a continuous
positive squared norm on \(L\): add the squared absolute
values of a fiber vector's local coefficients multiplied by
the corresponding functions wherever they are supported.
At each point at least one term has positive weight.
The norm of each chosen local basis vector on its closed
radius-two polydisc has a positive minimum and a finite
maximum. Choose common bounds \(0<c\le C_0\) for all those
basis norms. The induced norm of a local section \(f e^{\otimes m}\)
is \(|f|\,\|e\|^m\).

If a nonzero section vanished to order at least \(k\) at every
chosen center, scale it to have global maximum norm one. Its
local coefficient has absolute value at most \(c^{-m}\) on
radius two. For each point \(z\) of radius at most one, the
function \(\lambda\mapsto f(\lambda z)\) on the radius-two
disc has a zero of order at least \(k\) at zero. Divide by
\(\lambda^k\) and apply the maximum principle on its boundary
to obtain \(|f(z)|\le2^{-k}c^{-m}\). The maximum principle
here follows from the circle mean in the Cauchy formula of
Lemma~\ref{lem:analytic-tools}: an interior maximum of a
nonconstant holomorphic function is impossible since equality
with its circle average in absolute value forces constancy
on every small circle, and hence on a neighborhood.
Thus on the covering inner discs the global norm is at most
\(2^{-k}(C_0/c)^m\). Choose \(A\ge1\) with
\(A\ge\log(C_0/c)/\log2\), and \(k=1+\lfloor Am\rfloor\).
This would give a global bound strictly below one. Therefore
simultaneous jets of orders below \(k\) separate sections.
There are at most \(s k^d\) coefficients, bounding the number
of monomials by allowing each exponent from zero to \(k-1\).
This proves the dimension bound.
\end{proof}

\begin{corollary}\label{cor:compact-holomorphic-graphs-algebraic}
The image of every holomorphic embedding of a connected compact
complex manifold of positive dimension in projective space
is cut out by finitely many homogeneous complex polynomials.
If \(X\) is such a manifold with a projective embedding and
\(f:X\to\mathbf P^M(\C)\) is holomorphic, its graph in the
product is cut out by finitely many bihomogeneous equations.
This also holds for a product of finitely many projectively
embedded connected compact manifolds as source.
\end{corollary}
\begin{proof}
For a projective embedding, take the pulled-back line bundle
whose local degree-one sections in a chart \(x_j\ne0\) are
the functions \(x_i/x_j\). On overlaps the changes of
trivialization are the nonzero ratios of chart coordinates.
The corresponding global sections give the original embedding.
Their homogeneous products are sections in the indicated
tensor powers. Lemma~\ref{lem:compact-section-jet-bound}
therefore verifies the growth hypothesis of
Theorem~\ref{thm:section-growth-algebraic-image}.

For two projective spaces use the map
\(([x],[y])\mapsto [x_i y_j]_{i,j}\). Its image consists of
the nonzero matrices up to scalar with all two-by-two minors
zero: any nonzero column spans all columns by those equations.
On a patch where \(x_i y_j\ne0\), column and row coordinate
ratios recover the inverse holomorphically. Thus it is a
holomorphic embedding, and iterations embed any finite
product. The graph \(x\mapsto(x,f(x))\) is a holomorphic
embedding because the first projection recovers its inverse
and identifies its differential with a graph. After the
product embedding use the first assertion and pull each
homogeneous equation back via the products \(x_i y_j\).
The resulting equations are bihomogeneous and cut out
exactly the graph in the product.
\end{proof}

\begin{corollary}\label{cor:algebraic-signed-fiber-graphs}
For any fiber of Theorem~\ref{thm:signed-fiber-projective-embedding},
the graphs of addition, inversion and multiplication by every
element of \(O_E\) are complex projective algebraic sets.
The same holds for the graph of every homomorphism between
such fibers induced by a real-linear map intertwining their
complex structures and taking their lattices into one another.
This applies to the fiber legs of
Proposition~\ref{prop:complex-actual-hecke-legs}.
\end{corollary}
\begin{proof}
In the period charts addition and inversion lift to vector
addition and negation and preserve the lattice equivalence.
The integral scalars act complex linearly and preserve the
lattice by Proposition~\ref{prop:signed-ball-period-coordinates}.
An intertwining real-linear map is complex linear in the two
period coordinates, so descends holomorphically precisely
when it takes the source lattice into the target. Apply
Corollary~\ref{cor:compact-holomorphic-graphs-algebraic}
to these maps and to the product torus as source for
addition. For the indicated legs the lattice inclusion and
intertwining were proved in Proposition~\ref{prop:complex-actual-hecke-legs}.
\end{proof}

\section{The canonical quadratic multiplier on a period fiber}

\begin{definition}\label{def:signed-canonical-shift}
Fix a finite frame \(x\), its signed matrix \(A\) and a
point of the ball. Use \(E_0=E_A\) and write
\(H(v,w)=E_0(v,I_Jw)+iE_0(v,w)\), the positive Hermitian
form of Proposition~\ref{prop:signed-normalized-period-matrix}.
Put \(\beta=(1+\sqrt{-d})/2\in O_E\) as in
Lemma~\ref{lem:actual-signs-and-padding}. For \(t\in V_\R\)
regarded as a point of the complex vector space with structure
\(I_J\), define
\[
 \begin{split}
 \eta_\beta(t)&=\exp\big(\pi i E_0(\beta t,t)\big),\\
 (W_tS)(z)&=\exp\big(-\pi H(t,z)-\tfrac\pi2 H(t,t)\big)\,S(z+t),
             \qquad U_t=\eta_\beta(t)^{-1}W_t .
 \end{split}
\]
The canonical section space consists of the entire functions
\(S\) with \(U_\lambda S=S\) for every \(\lambda\in\Lambda_x\).
We use \(z,t\) either as vectors or in the complex basis of
Proposition~\ref{prop:signed-normalized-period-matrix}.
\end{definition}

\begin{lemma}\label{lem:signed-canonical-cocycle}
For \(t,s\in V_\R\) and entire functions,
\[
  W_tW_s=e^{\pi i E_0(t,s)}W_{t+s},\qquad
  U_tU_s=e^{2\pi i E_0(\beta t,s)}U_{t+s}.
\]
For \(t\) in the dual lattice
\(\Lambda_x^\#=\{t:E_0(t,\Lambda_x)\subseteq\Z\}\), \(U_t\)
preserves the canonical section space and there depends only
on \(t\bmod\Lambda_x\). On that space
\(U_tU_s=e^{2\pi i E_0(t,s)}U_sU_t\) for dual \(s,t\).
\end{lemma}
\begin{proof}
The alternating trace expression of
Proposition~\ref{prop:signed-period-riemann-form} gives
\(E_0(\beta t,s)=E_0(t,\bar\beta s)\), and
\(\beta+\bar\beta=1\). These imply
\(E_0(\beta t,s)+E_0(t,\beta s)=E_0(t,s)\).
In the product of the two \(W\)-factors the excess exponent
over the \(W_{t+s}\)-factor is
\(\pi(H(t,s)-H(s,t))/2=\pi i E_0(t,s)\).
For the \(U\)-factors this is increased by
\(\pi i(E_0(\beta t,s)+E_0(\beta s,t))\) on expanding
\(\eta_\beta(t+s)/(\eta_\beta(t)\eta_\beta(s))\). The adjoint identity gives the
stated exponent. Interchanging \(s,t\) in this formula
gives the commutation law. Since \(\beta,\bar\beta\) preserve
\(\Lambda_x\), the numbers \(E_0(\beta\lambda,t)\) and
\(E_0(\beta t,\lambda)\) are integers for every lattice
\(\lambda\) and dual \(t\). Thus \(U_\lambda U_t=U_tU_\lambda
=U_{t+\lambda}\), proving both assertions on sections.
\end{proof}

\begin{theorem}\label{thm:canonical-theta-basis-on-ball}
Write \(D,\Omega\) for the data of
Proposition~\ref{prop:signed-normalized-period-matrix}, and
\(M=(\Im\Omega)^{-1}\). Choose real columns \(a,c\) with
entries in \(\{0,\tfrac12\}\) by
\[
     e^{2\pi i a_j}=\eta_\beta(e_j),\qquad e^{2\pi i c_j}=\eta_\beta(f_j)
\]
where the expressions on the right mean the multiplier in
Definition~\ref{def:signed-canonical-shift}. The columns \(a,c\) are constant throughout the ball. Put
\(d_0=\Omega a-D^{-1}c\). A basis of the canonical space
in the complex coordinates \(z\) is
\[
 S_r(z)=
   \exp\big(\tfrac{\pi}{2}z^{\mathsf T}Mz+2\pi i a^{\mathsf T}z\big)
                \Theta_{1,r}(z+d_0),\qquad r\in\Z^g/D\Z^g .
\]
In particular the canonical space has dimension \(h^n\),
and its projective evaluations embed the fibers holomorphically
over the ball.
\end{theorem}
\begin{proof}
On a lattice vector \(\lambda\), \(E_0(\beta\lambda,\lambda)\)
is an integer since \(\beta\Lambda_x\subseteq\Lambda_x\).
The multiplier is therefore a sign, determined by the fixed
real form and lattice endomorphism independent of \(b\).
By Lemma~\ref{lem:signed-canonical-cocycle}, the \(U_\lambda\)
on lattice vectors multiply by addition, so it suffices to
impose the conditions on \(e_j,f_j\). In the complex \(e\)-basis
one has \(H(t,z)=\bar t^{\mathsf T}Mz\).
For a section write \(S(z)=G(z)h(z)\) with
\(G(z)=\exp(\pi z^{\mathsf T}Mz/2+2\pi i a^{\mathsf T}z)\).
For a general vector \(t=u+\Omega v\) with \(u,v\in\R^g\),
expansion of the quadratic gives
\[
 \frac{(W_tS)(z)}{G(z)}
 =\exp\big(2\pi i v^{\mathsf T}z+\pi i v^{\mathsf T}t+
                       2\pi i a^{\mathsf T}t\big)\,h(z+t).
\]
Indeed \((t-\bar t)^{\mathsf T}M=2i v^{\mathsf T}\) by the
formula for \(M\). At \(e_j\) the equation \(W_{e_j}S=\eta_\beta(e_j)S\)
is just periodicity of \(h\) in coordinate \(j\).
At \(f_j\) one has \(v=D\mathbf e_j\), \(u=0\), with
\(\mathbf e_j\) the \(j\)-th coordinate column.
The multiplier of \(h(z+\Omega v)\) relative to \(h(z)\) becomes
\[
    \exp\big(-\pi i v^{\mathsf T}\Omega v
                      -2\pi i v^{\mathsf T}(z+d_0)\big),
\]
since \(v^{\mathsf T}d_0=a^{\mathsf T}\Omega v-c_j\).
Hence \(h(z)=H_0(z+d_0)\) with \(H_0\) satisfying exactly
the standard degree-one theta equations of
Theorem~\ref{thm:signed-theta-sections}. The converse follows
from the same equations. That theorem gives the basis and
dimension. In projective coordinates the common nonzero
factor \(G(z)\) cancels. The translation by \(d_0\) is
holomorphic in \(b\), and is a biholomorphic translation on
each torus. The embedding result now follows from
Theorem~\ref{thm:signed-fiber-projective-embedding}.
\end{proof}

\section{Descent of the canonical theta coordinates to the level quotient}

\begin{theorem}\label{thm:signed-constant-translation-matrices}
In the symplectic basis of
Lemma~\ref{lem:signed-period-elementary-divisors} the dual
lattice vectors have complex coordinates \(t=u+\Omega v\)
where \(u\in D^{-1}\Z^g,\ v\in\Z^g\).
With the canonical basis of Theorem~\ref{thm:canonical-theta-basis-on-ball}
one has
\[
 U_t S_r=\eta_\beta(t)^{-1}
   \exp\!\left(\pi i v^{\mathsf T}u+
        2\pi i\big(a^{\mathsf T}u+v^{\mathsf T}D^{-1}c+
                                      r^{\mathsf T}u\big)\right) S_{r+v}.
\]
The matrices in this basis of the translations indexed by
\(\Lambda_x^\#/\Lambda_x\) are independent of the period
point and form a basis of the full endomorphism algebra.
\end{theorem}
\begin{proof}
A vector with real coordinates \(\sum e_j u_j+\sum f_j s_j\)
is dual exactly if \(Du,Ds\) are integral by the paired
basis. Its complex coordinate is \(u+\Omega Ds\), proving
the description with \(v=Ds\). In particular \(t\), as a
real vector, and the value of its multiplier
\(\eta_\beta(t)=\exp(\pi i E_0(\beta t,t))\) are independent
of the complex period when \(u,v\) are fixed.
The Fourier summation in Theorem~\ref{thm:signed-theta-sections}
gives
\[
 \Theta_{1,r}(z+d_0+u+\Omega v)=
 e^{2\pi i r^{\mathsf T}u-\pi i v^{\mathsf T}\Omega v
                          -2\pi i v^{\mathsf T}(z+d_0)}
                              \Theta_{1,r+v}(z+d_0).
\]
Indeed for each frequency \(m\in r+D\Z^g\) one has
\(e^{2\pi i m^{\mathsf T}u}=e^{2\pi i r^{\mathsf T}u}\),
and completing the square replaces \(m\) by \(m+v\).
Multiply by the Gaussian translation factor in the proof of
Theorem~\ref{thm:canonical-theta-basis-on-ball} and by
\(\eta_\beta(t)^{-1}\). Since \(d_0=\Omega a-D^{-1}c\),
the \(z,\Omega\) terms in the scalar cancel and give exactly
the stated matrix. The formula respects changing \(r\) by
\(D\Z^g\), and the class of \(t\) suffices by
Lemma~\ref{lem:signed-canonical-cocycle}.

For a nonzero class \(t\) there exists a class \(s\) with
\(E_0(t,s)\notin\Z\): the dual of the dual in the paired
basis is the original lattice, as follows by testing at the
vectors \(e_j/d_j,f_j/d_j\). The commutator formula of
Lemma~\ref{lem:signed-canonical-cocycle} implies that
conjugating \(U_t\) by \(U_s\) multiplies it by a scalar
different from one. All these operators are invertible by
the same formula with the negative shift. Trace invariance
under conjugation gives \({\rm tr}(U_t)=0\) off the zero class.
In any linear relation among the class representatives multiply
by \(U_{-t}\) and take the trace. Only the coefficient at
\(t\) can contribute, with a nonzero scalar times the dimension.
Thus every coefficient vanishes. The number of these operators
is \(h^{2n}\) by Proposition~\ref{prop:signed-period-riemann-form},
equal to the square of the dimension in
Theorem~\ref{thm:canonical-theta-basis-on-ball}. They are a basis.
\end{proof}

\begin{corollary}\label{cor:level-invariant-theta-coordinates}
Fix a representative \(x C\) of one rational orbit of finite
frames of the signed complex family. Use one symplectic basis
on \(\Lambda_x\) throughout its period ball. The canonical
projective evaluation
\[
           (b,P)\longmapsto [S_r(P)]_r
\]
descends to a holomorphic map from this component of
\(\mathcal X(C,1)\) to \(\mathbf P^{h^n-1}(\C)\), giving a
projective embedding on each torus. If \(M\) is any positive
integer with \(\mathfrak n\in M O_E\), then the points of
\(M^{-1}\Lambda_x/\Lambda_x\) give individual holomorphic
sections over that base component. The projective coordinates
of each such section descend there as well.
\end{corollary}
\begin{proof}
Let \(y\in\Gamma\) with \(k=x^{-1}y x\in C\).
Theorem~\ref{thm:analytic-framed-torus-family} gives
\(y^* A y=A\). The rational map commutes with \(\beta\),
preserves \(E_0\) by its trace formula, and carries the complex
structure at \(b\) to the one at \(y b\). It carries the positive
Hermitian form \(H\) to the corresponding form, by the
defining formula in Definition~\ref{def:signed-canonical-shift}.
Thus pullback of entire functions by this complex-linear map
carries the canonical space at \(yb\) isomorphically to the
one at \(b\), and intertwines the translations at \(yt\)
and \(t\) by their explicit Gaussian formulas.

Here \(y\) acts trivially on \(\Lambda_x^\#/\Lambda_x\).
Indeed \(\Lambda_x^\#\subseteq h^{-1}\Lambda_x\) by
Proposition~\ref{prop:signed-period-riemann-form}, and
\(\mathfrak n\in h O_E\) by Proposition~\ref{prop:actual-integral-moduli-level}
and Corollary~\ref{cor:actual-moduli-places}. Since
\(k-I\in hM_n(\widehat O)\), local frame coordinates give
\((y-I)h^{-1}\Lambda_x\subseteq\Lambda_x\).
In the chosen canonical bases, the pullback matrix from
\(yb\) to \(b\) commutes with every translation matrix
by Theorem~\ref{thm:signed-constant-translation-matrices}.
Its commutant is scalar: commuting with the full matrix
algebra entails commuting with each diagonal coordinate
projection and then with each elementary matrix between
two coordinates. Thus the evaluation lists at \((yb,yP)\)
and \((b,P)\) differ by a common nonzero factor. Changing
the torus representative likewise multiplies all sections
by the same factor by Definition~\ref{def:signed-canonical-shift}.
Their projective evaluation is invariant and holomorphic by
Theorem~\ref{thm:canonical-theta-basis-on-ball} and the local
quotient charts in Theorem~\ref{thm:analytic-framed-torus-family}.

For the second assertion replace \(h\) by \(M\) in the local
inclusion. This shows \(y\) fixes every labelled point mod
\(\Lambda_x\). In the ball charts the section corresponding
to a fixed real vector \(v\) is the holomorphic section
\(b\mapsto Q_b(v)\) modulo its periods, by
Proposition~\ref{prop:signed-ball-period-coordinates}.
Hence it descends, as do its projective coordinates.
\end{proof}

\section{Torsion evaluations as coordinates on a compact period set}

\begin{definition}\label{def:fixed-real-torus-evaluation}
Fix \(x\) and the symplectic basis of
Corollary~\ref{cor:level-invariant-theta-coordinates}.
Before taking the rational quotient let \(T=V_\R/\Lambda_x\)
with its fixed real differentiable structure and let
\(F_b:T\to\mathbf P^{h^n-1}(\C)\) be its canonical projective
theta embedding at \(b\in\mathbb B\). Set \(T[m]=m^{-1}\Lambda_x/\Lambda_x\).
Use any ordered integer basis to identify
\(T=\R^{2g}/\Z^{2g}\), with the distance given by the least
Euclidean length of a lift of the difference. On projective
space use the distance \(\|P(z)-P(z')\|\) with the Euclidean
matrix norm, where \(P([u])=u u^*/(u^*u)\). This is independent
of the nonzero vector representative \(u\). The matrix recovers
each affine coordinate \(u_i/u_j\) where \(u_j\ne0\) by dividing
its \(ij\)-entry by its \(jj\)-entry. In particular this gives
a smooth embedding of projective space into a real matrix space.
\end{definition}

\begin{proposition}\label{prop:finite-period-evaluation-tests}
For a compact \(K\subseteq\mathbb B\), there is \(m\ge1\)
such that
\[
             b\longmapsto (F_b(t))_{t\in T[m]}
\]
is injective on \(K\) and its derivative is injective at all
points of \(K\). It is holomorphic on \(\mathbb B\).
Also there are \(0<c_K\le C_K<\infty\) with
\[
 c_K\,{\rm dist}(v,w)\le {\rm dist}(F_b(v),F_b(w))
               \le C_K\,{\rm dist}(v,w)
                  \qquad(b\in K,\ v,w\in T).
\]
\end{proposition}
\begin{proof}
The period coordinates of Proposition~\ref{prop:signed-ball-period-coordinates}
and the theta functions of Theorem~\ref{thm:canonical-theta-basis-on-ball}
show that \((b,v)\mapsto F_b(v)\) is smooth in real coordinates
and holomorphic in \(b\) for fixed \(v\). Indeed fixed real
vectors have holomorphic periods, and in projective coordinates
only the translated Fourier sections remain. In the real torus
\(\bigcup_m T[m]\) is dense, by approximating real basis
coordinates by rational numbers. If \(F_b(t)=F_{b'}(t)\) at
all torsion points, continuity gives equality at all \(v\).
As both maps are holomorphic embeddings for their respective
complex structures by Theorem~\ref{thm:canonical-theta-basis-on-ball},
the identity of the real torus is biholomorphic from one
structure to the other. Differentiating gives \(I_b=I_{b'}\).
Thus the positive lines, and hence their coordinates \(b,b'\),
are identical by Proposition~\ref{prop:signed-ball-period-coordinates}.

Suppose a real tangent vector \(w\) at \(b\) gives zero
derivative of \(F_b(t)\) for every torsion \(t\). Continuity
of the derivative gives \(\partial_w F_b(v)=0\) at all \(v\).
In a local complex target chart the differential in \(v\)
satisfies
\[
               d_v F_b\, I_b=i\,d_v F_b ,
\]
since the embedding of that torus is holomorphic. In these
coordinates \(\partial_w(d_v F_b)=d_v(\partial_w F_b)=0\)
at the chosen \(b\) since the two-variable function is smooth
and its \(w\)-derivative vanishes identically in \(v\).
Differentiating the equation now gives
\(d_vF_b\,\partial_w I_b=0\). Fiber immersion implies
\(\partial_w I_b=0\). Thus \(\partial_w J_b=0\) at the
distinguished embedding. Differentiate \(J_b(1,\bar b)=(1,\bar b)\)
along \(w\). This puts \((0,\bar w)\) in the positive line
of \((1,\bar b)\), forcing \(w=0\).
For \(n=1\) the parameter space has zero-dimensional tangent
space and a single point, giving the same two assertions.

At any \(b\) a finite number of torsion evaluations suffice
for an injective derivative: while a common kernel of the
chosen real differentials has positive dimension, the preceding
assertion supplies a further differential reducing that dimension.
By the holomorphic inverse argument of
Theorem~\ref{thm:signed-fiber-projective-embedding}, these
evaluations are locally injective on a neighborhood where
the derivative keeps full rank. Finitely many such neighborhoods
\(U_a\) cover \(K\). Outside \(\bigcup_a U_a\times U_a\),
the pairs of \(K\times K\) form a compact set of distinct
points. Each pair is distinguished by some torsion evaluation,
and that same evaluation distinguishes pairs on a neighborhood
by continuity. A finite subcover supplies a combined finite
list of tests also separating those pairs. Take \(m\) a
common multiple of the orders of the tests.

For the distance bounds compose \(F_b\) with the real matrix
map \(P\) of Definition~\ref{def:fixed-real-torus-evaluation}.
The derivative in the torus directions is injective, and on
the compact set of \(b\in K\), torus points and unit tangent
vectors its norm has positive minimum and finite maximum.
Integration along a least-length lifted segment gives the
upper bound. The derivative is uniformly continuous in the
torus variable on this compact set. Thus for sufficiently
short segments, integration of its difference from the
initial derivative gives an error at most half of that
positive minimum times the segment length. Subtracting the
error gives the lower bound at short distances. At distances
bounded below, the compact set of triples has a positive
minimum image distance by fiber injectivity; divide this
by the finite diameter of the real torus. Taking the smaller
lower bound proves the assertion.
\end{proof}

\section{A finite level embedding by labelled torsion points}

\begin{lemma}\label{lem:bounded-lifts-torsion-labels}
For a component representative \(x C\) let
\(\Gamma_x=\{y\in\Gamma:x^{-1}yx\in C\}\), acting on
\(\Lambda_x\) and the ball. In an ordered basis of the lattice
write \(D_1=2g\) for its rank. Fix a compact
\(K\subset\mathbb B\). There exist a finite subset
\(S_0\subset\Gamma_x\) and positive integer \(m_0\) such
that whenever \(m\) is a multiple of \(m_0\) and
\[
        F_b(t)=F_{b'}(\gamma t)
                 \quad(t\in T[m]),\qquad
        b,b'\in K,\quad\gamma\in\Gamma_x ,
\]
the action of \(\gamma\) on \(T[m]\) is that of a member
of \(S_0\).
\end{lemma}
\begin{proof}
Write \(C_1,c_1\) for the distance constants in
Proposition~\ref{prop:finite-period-evaluation-tests}. Compare
the equality at zero and at each \(t=\mathbf e_j/m\), where
the \(\mathbf e_j\) are the coordinate lattice basis elements.
Their images at \(b\) have distance at most \(C_1/m\). Thus
\(\gamma\mathbf e_j/m\) has torus distance at most
\(C_1/(c_1m)\) from zero, by the lower bound at \(b'\).
Choose a shortest lift in \(\R^{D_1}\) and multiply it by
\(m\). This gives an integral matrix \(A_0\), congruent to
\(\gamma\) modulo \(m\), whose columns have Euclidean lengths
at most \(C_1/c_1\). Since \(\gamma\) permutes \(T[m]\), the
same argument with \(b,b'\) interchanged gives a matrix
\(B_0\equiv\gamma^{-1}\bmod m\) with the same bound.
There are finitely many such integral pairs.

Each \(\gamma\) commutes with multiplication by the chosen
integral basis of \(O_E\) and preserves the integral Gram
matrix of the rational trace form \(B_A\), as proved in
Theorem~\ref{thm:analytic-framed-torus-family} and
Lemma~\ref{lem:period-metric-covolume}. For every possible
pair \(A_0,B_0\), take all entries of \(A_0B_0-I,B_0A_0-I\),
the commutators of \(A_0\) with these multiplication matrices,
and \(A_0^{\mathsf T} B_A A_0-B_A\) in the chosen basis.
Their absolute values have a common finite bound independent
of \(m\). For \(m\) larger than it all these entries vanish
whenever the hypotheses hold, because reduction modulo \(m\)
gave exactly zero. Thus \(A_0\) is an invertible integral
matrix commuting with \(O_E\) and preserving \(B_A\).
Rational extension then gives an \(E\)-linear isometry of
the Hermitian \(A\): varying the second argument through
all \(E\)-multiples and taking traces recovers its values,
as in Proposition~\ref{prop:definite-full-class-set}.

Choose \(m_0\) also divisible by the positive integer
\(|N_{E/\Q}\mathfrak n|\) and larger than the bound.
Multiplication by this integer maps \(O_E\) into
\(\mathfrak n O_E\) by the adjugate of multiplication by
\(\mathfrak n\), whose determinant is its norm.
Hence \(m\Lambda_x\subseteq\mathfrak n\Lambda_x\).
One has \((\gamma-I)\Lambda_x\subseteq\mathfrak n\Lambda_x\)
by completion in the frame \(x\). Consequently the same
inclusion holds for \(A_0-I\) since \(A_0\equiv\gamma\bmod m\).
In frame coordinates \(x^{-1}A_0x\) is integral invertible
by its action on local lattices, unitary by \(A_0^*AA_0=A\)
and \(x^*Ax=\varepsilon I\), and congruent to \(I\bmod\mathfrak n\).
It lies in \(C\), so \(A_0\in\Gamma_x\).
The finite collection of members of \(\Gamma_x\) with the
specified bound on columns gives \(S_0\).
\end{proof}

\begin{theorem}\label{thm:finite-torsion-level-base-embedding}
Let \(\Gamma_x(m)\) be the kernel of the action of \(\Gamma_x\)
on \(T[m]\). For every multiple of some positive integer
\(m_x\), the map
\[
  \Gamma_x(m)\backslash\mathbb B\longrightarrow
          \big(\mathbf P^{h^n-1}(\C)\big)^{T[m]},
          \qquad [b]\longmapsto (F_b(t))_{t\in T[m]}
\]
is a closed holomorphic embedding. The group
\(\Gamma_x/\Gamma_x(m)\) permutes the projective factors
by its action on the labelled points.
\end{theorem}
\begin{proof}
The map is well defined and equivariant for the indicated
permutations because \(F_{\gamma b}(\gamma t)=F_b(t)\)
by Corollary~\ref{cor:level-invariant-theta-coordinates}.
It is holomorphic by Proposition~\ref{prop:finite-period-evaluation-tests}.
The coarse quotient \(\Gamma_x\backslash\mathbb B\) is the
component at \(xC\) in Theorem~\ref{thm:analytic-framed-torus-family}
with no marked torus coordinates. It is compact by
Theorem~\ref{thm:compact-signed-complex-family}, since these
components form a disjoint open union and are therefore closed.
It has a compact set \(K\subset\mathbb B\) meeting all
orbits: images of relatively compact ball neighborhoods cover
the quotient and a finite subcover suffices. Use
Lemma~\ref{lem:bounded-lifts-torsion-labels} for this \(K\).
Apply Proposition~\ref{prop:finite-period-evaluation-tests}
to the compact set \(K'=K\cup\bigcup_{s\in S_0}s^{-1}K\).
Choose \(m_x\) a common multiple of the resulting testing
integer and \(m_0\) in the preceding lemma. The whole list
at \(m\) then separates all points of \(K'\) and is immersive
there. Equivariance and the coverage by \(K\) give immersion
everywhere.

Suppose \(b_1,b_2\) give the same labelled evaluations. Choose
\(u_1,u_2\in\Gamma_x\) with \(b=u_1 b_1,\ b'=u_2 b_2\) in \(K\).
Equivariance yields \(F_b(t)=F_{b'}(\gamma t)\) for all
\(t\in T[m]\) and \(\gamma=u_2u_1^{-1}\).
The preceding lemma gives \(s\in S_0\) with the same
action as \(\gamma\) on these labels. Therefore
\(F_b(t)=F_{s^{-1}b'}(t)\) for all labels. Both base points
belong to \(K'\), so \(b=s^{-1}b'\).
It follows that \(b_2=u_2^{-1}s u_1 b_1\), with
\(u_2^{-1}s u_1\in\Gamma_x(m)\) by its label action.
This proves injectivity.

The finer quotient is compact because \(\Gamma_x(m)\) has
finite index, being the kernel of a permutation action on a
finite set. The translates of \(K\) by coset representatives
therefore exhaust it after projection. Both groups act freely
with locally disjoint charts by
Theorem~\ref{thm:analytic-framed-torus-family}. A holomorphic
injective immersion of this compact quotient into the
Hausdorff target is a homeomorphism onto a closed image,
and is a local holomorphic graph by the inverse argument
of Theorem~\ref{thm:signed-fiber-projective-embedding}.
This proves the result.
\end{proof}

\section{Projective equations for the complex level correspondences}

\begin{lemma}\label{lem:free-projective-finite-quotient}
Let a finite group \(D_2\) act freely on a compact complex
manifold \(Y\). Suppose \(Y\to\mathbf P(V)\) is a holomorphic
embedding intertwining its action with projective linear
transformations. Then \(Y/D_2\) admits a closed holomorphic
embedding into a projective space.
\end{lemma}
\begin{proof}
Identify \(v\in V\) with the linear polynomial
\(\lambda\mapsto\lambda(v)\) on its dual. Choose linear
representatives for the projective transformations. At the
point represented by \([v]\) send its orbit to the line of
the polynomial \(\prod_{s\in D_2}(sv)\) in the homogeneous
polynomials of degree \(|D_2|\) on \(V^\vee\). Nonzero factors
have nonzero product; rescaling any lifts changes only the
overall scalar. Thus this is a well-defined holomorphic
projective map on \(Y\), constant on orbits. Equality of two
products up to scalar forces equality of their sets of linear
factor lines with multiplicities. Indeed a nonzero linear
polynomial generates a prime ideal, since after an invertible
coordinate change its quotient ring is the polynomial ring
in the remaining variables. Comparing factors and cancelling
proves the assertion by induction. Hence equality identifies
one orbit.

Because the action is free the factor lines at a point are
distinct. Differentiate the product in a local representative
\(v\) of the projective embedding. If the derivative is a scalar
multiple of the product, reduce the differentiated equation
modulo the ideal of the identity factor \(v\).
All terms except its derivative vanish, and none of the other
factors vanishes in this integral domain. Therefore the
derivative of \(v\) is a scalar multiple of \(v\); projective
immersion on \(Y\) makes the original tangent zero.
A free finite action has local quotient charts by shrinking
a neighborhood away from its finitely many nonidentity
translates. The quotient map is locally biholomorphic, and
the descended map is an injective immersion. Compactness and
the local inverse argument of
Theorem~\ref{thm:signed-fiber-projective-embedding} prove it
is a closed embedding.
\end{proof}

\begin{theorem}\label{thm:complex-level-projective-equations}
The complex spaces \(\mathcal X(C,\nu)\) of
Definition~\ref{def:signed-complex-period-space} and
\(\mathcal X(C',\nu)\) for every open finite-index \(C'\subset C\)
admit closed holomorphic embeddings in projective spaces whose
images are cut out by finitely many homogeneous complex
polynomials. For all levels and finite supports in
Proposition~\ref{prop:complex-actual-hecke-legs}, the cover
map and each of its indicated Hecke legs have graphs cut
out by finitely many bihomogeneous complex polynomials.
\end{theorem}
\begin{proof}
First take the base without the marked torus coordinates and
fix a frame \(xC\) representing one of its components. Write
\(\Gamma_{x,C'}=\{z\in\Gamma:x^{-1}zx\in C'\}\) and
\(\Gamma_x=\Gamma_{x,C}\). Over this component every fine
frame can be represented by \(xkC'\), \(k\in C\), by undoing
a rational change; finitely many suffice by the finite index.
For each such representative one may substitute \(xk\) for
\(x\) below. The open subgroup \(C'\) contains
\(\{k\in C:k-I\in m_1 M_n(\widehat O)\}\) for an integer
\(m_1\) by the neighborhood argument of
Proposition~\ref{prop:complex-actual-hecke-legs}.
For \(m\) divisible by \(m_1\), the subgroup \(\Gamma_x(m)\)
of Theorem~\ref{thm:finite-torsion-level-base-embedding}
is contained in \(\Gamma_{x,C'}\): its congruence on \(T[m]\)
means \((z-I)\Lambda_x\subseteq m\Lambda_x\), the displayed
matrix congruence upon using the local lattice frames.
Take \(m\) also divisible by the integer \(m_x\) of that
theorem. The group \(\Gamma_x(m)\) is normal in
\(\Gamma_{x,C'}\) and the quotient group is finite.

Embed the deep quotient by the labelled theta evaluations of
Theorem~\ref{thm:finite-torsion-level-base-embedding} followed
by the product embedding of
Corollary~\ref{cor:compact-holomorphic-graphs-algebraic}.
Permuting these factors acts linearly by permuting coordinates
in the ambient tensor product. The finite quotient group
acts freely on the deep quotient: if an element fixes a
deep class, multiplication by an element of \(\Gamma_x(m)\)
makes it fix a ball point, whose stabilizer in \(\Gamma_x\)
is trivial by Theorem~\ref{thm:analytic-framed-torus-family}.
Lemma~\ref{lem:free-projective-finite-quotient} therefore
gives a projective embedding of the base at \(C'\), component
by component, including \(C'=C\).

For \(\nu\) marked points use the joint map of the base
embedding just constructed and the \(\nu\) corresponding
fiber evaluations in Corollary~\ref{cor:level-invariant-theta-coordinates}.
They descend for the subgroup \(\Gamma_{x,C'}\) and are
holomorphic in the family charts. If two points have the
same image, the base embedding makes the two bases agree
in the quotient; after choosing the same frame and base
representative, fiber injectivity identifies each of their
marked points. Similarly a tangent vector killed by the
joint derivative is vertical by the base embedding, and
then each of its vertical components is zero by fiber immersion.
The domain is compact. For example, above any fixed component
use a compact set of ball representatives for the finite-index
subgroup, by the compact-set argument of
Theorem~\ref{thm:finite-torsion-level-base-embedding}. For
each torus factor add a closed real fundamental parallelotope
as in Theorem~\ref{thm:compact-signed-complex-family}. There
are finitely many components, since the coarse compact base
is a disjoint union of open components and the finer frames
over each have the finite list given above. Therefore the
joint injection is a closed immersion in complex analytic
charts by the local inverse argument already used. Compose
with the product embedding. If desired a single projective
space for all components is obtained by placing their spaces
in disjoint projective coordinate subspaces of their direct
sum.

Corollary~\ref{cor:compact-holomorphic-graphs-algebraic}
applies on each positive-dimensional component and to each
holomorphic map in the statement; the zero-dimensional
components are individual projective points, cut out by
linear equations. Finite unions of the resulting closed
sets are cut out by the finite lists of products formed
by taking one equation from each list, since any point
outside all the sets fails some equation from each list.
The same argument applies in the product of the source and
target projective spaces to the graphs, giving the bihomogeneous
assertion.
\end{proof}

\section{Simple algebraic coordinates at a chosen period point}

\begin{lemma}\label{lem:finite-projection-dense-open}
In Proposition~\ref{prop:finite-projection-simple-sheets}
let \(a\ne0\) in \(R\). For any point \(z_0\in Z\) and
any additional prescribed nonzero polynomial on the base,
there exist points \(z_j\to z_0\) with \(a(z_j)\ne0\),
lying over the complement of that polynomial's zero set
and over the simple-sheet locus for a choice of \(l\) separating
the fiber of \(z_0\). For a finite-dimensional unital complex
domain \(R\), every nonzero element is nonzero at its unique
complex algebra homomorphism to \(\C\).
\end{lemma}
\begin{proof}
Make the indicated choice of \(l,P,b\) in
Proposition~\ref{prop:finite-projection-simple-sheets}.
Its fraction field calculation expresses \(a^{-1}\) in powers
of \(l\) with coefficients rational in \(t\). Enlarge \(b\)
by their nonzero denominators and the extra polynomial.
Then \(a\) is invertible in \(R[1/b]\), hence nonzero on
every sheet above \(U_b\). Choose values \(t_j\in U_b\)
tending to \(t(z_0)\). As in the root-limit argument of
Theorem~\ref{thm:section-growth-algebraic-image}, roots
\(\lambda_j\) of the monic \(P(t_j,T)\) can be chosen to
converge to \(l(z_0)\). The sheets yield \(z_j\) with
these values. All their coordinates are bounded by the
same proposition, and any convergent subsequence has the
same unique limit \(z_0\) by separation of the fiber.
This gives the claimed approximating points. In dimension
zero a finite-dimensional domain over \(\C\) is a field
by injectivity of multiplication, and its elements are
algebraic over \(\C\). Hence it is \(\C\) by the complex
root property in Lemma~\ref{lem:place-fields}.
\end{proof}

\begin{theorem}\label{thm:smooth-algebraic-projection-chart}
Let \(Y\) be a nonempty affine coordinate chart of a connected
compact complex manifold holomorphically embedded in projective
space, of positive complex dimension \(d\).
Use its algebraic coordinate ring \(R\) and the identification
of \(Y\) with its affine zero set in
Theorem~\ref{thm:section-growth-algebraic-image} and
Corollary~\ref{cor:compact-holomorphic-graphs-algebraic}.
For every \(y_0\in Y\) there are a finite linear projection
\(t:Y\to\C^d\), a linear function \(l\) and a monic
\(P(t,T)\in\C[t,T]\) for which \(P(t,l)=0\),
\(s(y_0)\ne0\) with \(s=P_T(t,l)\) the derivative in \(T\),
and
\[
      R[1/s]=\big(\C[t,T]/(P)\big)[1/P_T],\qquad T\mapsto l .
\]
Here brackets denote localization, adjoining the reciprocal
of the indicated element. In particular the points have
algebraically nonsingular charts of dimension \(d\).
\end{theorem}
\begin{proof}
Begin the generator elimination of
Lemma~\ref{lem:complex-polynomial-normalization} with the
affine coordinate ratios of \(Y\), whose differentials have
combined rank \(d\) at \(y_0\) by its analytic embedding.
At each dependent stage, the number \(N'\) of current generators
is greater than \(d\): at rank \(d\) any \(d\) with independent
differentials are algebraically independent by the local
inverse theorem proved in Theorem~\ref{thm:signed-fiber-projective-embedding}.
In the choice of projection along \((c_1,\ldots,c_{N'-1},1)\)
impose also that this vector lie outside the \(d\)-dimensional
image of the tangent space under the current differential.
There is such a vector satisfying the leading polynomial
condition: this tangent image is a proper subspace, so a
nonzero linear functional annihilates it, and the product
of that functional with the nonzero top homogeneous polynomial
is nonzero. The same dehomogenization argument as in the lemma
applies. The linear projection with this kernel preserves
the tangent rank. At termination the number of coordinates
is \(d\) by the growth bounds in
Theorem~\ref{thm:section-growth-algebraic-image}.
Thus \(t\) is locally biholomorphic at \(y_0\).

Apply Proposition~\ref{prop:finite-projection-simple-sheets}
with \(l\) separating the fiber of \(t(y_0)\).
The root \(l(y_0)\) is simple at this base value. Indeed if
it had multiplicity at least two, choose \(t_j\) on the
simple-sheet locus tending to \(t(y_0)\). At least two
distinct roots there tend along a subsequence to \(l(y_0)\):
take a convergent subsequence of all ordered root lists and
multiply their limiting linear factors. Points of both
sheets are bounded, and their only possible limit is \(y_0\)
by separation of the fiber. They would belong eventually
to the same neighborhood on which \(t\) is injective,
contradicting the distinctness of the roots. Hence
\(P_T(t(y_0),l(y_0))\ne0\).

For each affine coordinate \(r\) of \(Y\), on the simple-sheet
locus \(U_b\) of that proposition form the polynomial
\[
  Q_r(t,T)=\sum_{z:\,t(z)=t} r(z)
                       \prod_{\substack{z':\,t(z')=t\\z'\ne z}}
                                         (T-l(z')) .
\]
The coefficients are holomorphic since permutation of the
locally holomorphic sheets does not change the sum.
They have polynomial growth on \(U_b\) since all coordinates
and \(l\) satisfy such bounds by the proposition. They are
therefore polynomials in \(t\) by
Lemma~\ref{lem:polynomial-growth-punctured-space}. Substitution
at \(T=l\) gives \(Q_r(t,l)=r P_T(t,l)\) above \(U_b\),
and on all of \(Y\) by continuity and
Lemma~\ref{lem:finite-projection-dense-open} (with \(a=1\)
and the original \(b\) as additional polynomial).
These equalities put every coordinate in the stated localization
of \(\C[t,l]\). Monic division and independence of the lower
powers of \(l\) over \(\C(t)\) identify \(\C[t,l]\) with
\(\C[t,T]/(P)\), and \(t,l\in R\). This proves the ring
equality. On this chart there is one equation whose derivative
in \(T\) is invertible, giving the asserted nonsingular
hypersurface coordinates. Explicitly the linearized equation
determines the increment of \(T\) uniquely from those of \(t\);
even over a square-zero ideal the same holds for a correction
of a tentative lift by Taylor expansion modulo the square
of that ideal.
\end{proof}

\section{Algebraic regularity of the holomorphic maps}

\begin{lemma}\label{lem:bounded-rational-local-regular}
Let \(r\in\C(t_1,\ldots,t_d)\) be bounded on the complement
of finitely many polynomial zero sets in some open ball about
\(t_0\), with all excluded polynomials nonzero. It is a quotient
of polynomials with denominator nonzero at \(t_0\).
\end{lemma}
\begin{proof}
The polynomial ring is factorial. For the needed argument,
if \(B\) is a unique factorization domain, call a polynomial
in \(B[T]\) primitive when its coefficients have no common
prime divisor. Products of primitive polynomials are primitive,
since reduction modulo any prime is a product of nonzero
polynomials over a domain. Clearing denominators and removing
the coefficient gcd therefore identifies factorizations of
primitive polynomials with those over the fraction field up
to scalar units: the product of the primitive representatives
has coefficient gcd one, and a fraction relating two primitive
polynomials must have numerator and denominator units by prime
divisibility. Factorization with uniqueness over the fraction
field follows by Euclidean division and Bezout. Combining with
factorization of the coefficients gives unique factorization
in \(B[T]\). Induction from \(\C\) proves the assertion, and
unique factorization makes irreducible factors prime.

Write \(r=N/D\) with relatively prime numerator and denominator,
taking \(D=1\) if \(r=0\). If an irreducible factor \(g_0\)
of \(D\) vanished at \(t_0\), the domain
\(\C[t_1,\ldots,t_d]/(g_0)\) would contain a nonzero residue
of \(N\). Its complex points are exactly the zeros of \(g_0\)
by its single generator relation. By
Lemmas~\ref{lem:complex-polynomial-normalization}
and~\ref{lem:finite-projection-dense-open}, arbitrarily near
\(t_0\) on this hypersurface there is a point with \(N\ne0\),
including the zero-dimensional case. At such a point in
the ball the denominator vanishes. Approaching it from the
dense complement of the excluded sets and of \(D=0\) makes
the absolute value of \(N/D\) unbounded, a contradiction.
Thus \(D(t_0)\ne0\).
\end{proof}

\begin{theorem}\label{thm:holomorphic-map-algebraic-locally}
For a connected compact complex manifold \(X\) with projective
embedding, every holomorphic map \(f:X\to\mathbf P^M(\C)\)
is an algebraic morphism for the reduced algebraic structure
of Theorem~\ref{thm:section-growth-algebraic-image}. Here
algebraic regularity means that locally in affine source and
target charts each target coordinate is a quotient of
polynomials in the source coordinates with denominator
nonvanishing at the given point.
\end{theorem}
\begin{proof}
For dimension zero \(X\) is a point. Otherwise fix \(y_0\)
and an affine source chart \(Y\) with domain \(R\) as in
Theorem~\ref{thm:smooth-algebraic-projection-chart}.
Use the \(t,l,P\) chosen there and their fraction fields
\(k=\C(t)\subseteq K=\operatorname{Frac}R=k(l)\).
The fiber over \(t_0=t(y_0)\) is finite by
Proposition~\ref{prop:finite-projection-simple-sheets}.
Choose affine target coordinates \(h_1,\ldots,h_M\) using
a linear denominator nonzero at the images of all points
in that fiber. Such a linear form exists by avoiding finitely
many proper linear subspaces. Regard the \(h_j\) as functions
pulled back by \(f\). They are bounded on the entire inverse
image in \(Y\) of a sufficiently small ball about \(t_0\).
Indeed they are holomorphic and bounded on some neighborhood
of each fiber point after shrinking it. If points outside
these neighborhoods projected to a sequence tending to \(t_0\),
the coordinate bounds of Proposition~\ref{prop:finite-projection-simple-sheets}
and closedness of \(Y\) in its affine chart would give a
subsequential limit in that fiber and outside the neighborhoods.

First all \(h_j\) belong to \(K\). On the affine product chart
of the graph let \(S=R[h_1,\ldots,h_M]\) be its coordinate
ring, viewing the functions on their common domain. Its complex
points are exactly that part of the graph by
Corollary~\ref{cor:compact-holomorphic-graphs-algebraic}.
It is a domain by the identity theorem in
Theorem~\ref{thm:section-growth-algebraic-image}. In the
product embedding of the graph all source and target affine
coordinates occur as ratios of degree-one sections with a
common denominator, by holding the other factor's chart
coordinate fixed. The filtered span of words of length at
most \(m\) in them has dimension \(O((m+1)^d)\) by
Lemma~\ref{lem:compact-section-jet-bound}, with
\(d=\dim_\C X\), after multiplication by the denominator
to power \(m\). Any \(h_j\) transcendental over \(K\)
would give \(d+1\) algebraically independent coordinates
\(t_1,\ldots,t_d,h_j\); their monomials would contradict
this bound by the counting argument of
Lemma~\ref{lem:complex-polynomial-normalization}.
Thus \(\operatorname{Frac}S/K\) is a finite separable
extension. The primitive-element and monic-division proof
of Proposition~\ref{prop:finite-projection-simple-sheets}
applies with \(R\) in place of its base algebra: take a
linear combination of the \(h_j\) separating embeddings,
express each \(h_j\) in powers of it, and invert one nonzero
element of \(R\) clearing the coefficients, its minimal
polynomial and the derivative Bezout identity. Over the
complement in \(Y\) of that element's zero set the fibers
of the affine graph then have as many distinct points as
the extension degree. The complement is nonempty because
\(R\) is the algebra of functions on \(Y\). Each fiber of
a graph has at most one point, forcing degree one.

Now take \(h=h_j\in K\) and form over \(k\) the polynomial
\[
 Q_h(t,T)=\operatorname{Tr}_{K/k}\!
        \left(h\,\frac{P(t,T)-P(t,l)}{T-l}\right),
\]
with trace taken on coefficients. Over the simple sheets,
after excluding finitely many polynomial zero sets in the base
to clear the denominators in \(h\), this is
\(\sum_{z:\,t(z)=t}h(z)\prod_{z'\ne z}(T-l(z'))\).
In detail, evaluation at the distinct roots of \(P\) in
its splitting field diagonalizes multiplication on \(k[l]\)
by the Vandermonde argument of
Lemma~\ref{lem:number-field-lattice}. The rational expression
for \(h\) agrees with the target coordinate wherever its
denominator and the target chart denominator are nonzero,
by the graph identity. Near \(t_0\) the latter holds for
all sheets by the first paragraph, and denominators in \(R\)
can be inverted on all sheets after inverting a nonzero
base polynomial as in Lemma~\ref{lem:finite-projection-dense-open}.
Consequently every coefficient of \(Q_h\), a rational function
of \(t\), is bounded near \(t_0\) off those zero sets. This
uses the uniform bound for \(h\) on all these sheets and
the bound for \(l\) from
Proposition~\ref{prop:finite-projection-simple-sheets}.
By Lemma~\ref{lem:bounded-rational-local-regular}, these
coefficients have polynomial denominators nonzero at \(t_0\).
Evaluating the interpolation identity in the field \(K\)
gives \(Q_h(t,l)=hP_T(t,l)\), since exactly the sheet for
\(l\) contributes there. As \(P_T(t,l)\) is nonzero at \(y_0\)
by Theorem~\ref{thm:smooth-algebraic-projection-chart},
this represents \(h\) by an algebraically regular fraction
near \(y_0\). It agrees there with the given holomorphic
function by continuity from the dense open where the two
expressions agree. Applying this to all target coordinates
proves the statement.
\end{proof}

\begin{corollary}\label{cor:complex-level-algebraic-morphisms}
At all levels in Theorem~\ref{thm:complex-level-projective-equations}
the complex spaces are nonsingular projective algebraic varieties
component by component. The covers and the Hecke legs are
algebraic morphisms over \(\C\). Each projective period fiber
is an abelian variety over \(\C\), meaning here a connected
nonsingular projective variety with commutative algebraic group
operations. Its \(O_E\)-action and the homomorphisms of
Corollary~\ref{cor:algebraic-signed-fiber-graphs} are algebraic.
\end{corollary}
\begin{proof}
Use the projective embeddings of
Theorem~\ref{thm:complex-level-projective-equations}
and the nonsingular charts of
Theorem~\ref{thm:smooth-algebraic-projection-chart}, with
points as charts in dimension zero. Every map in question
is holomorphic by Proposition~\ref{prop:complex-actual-hecke-legs}
or Corollary~\ref{cor:algebraic-signed-fiber-graphs}.
Apply the preceding theorem on each source component,
using the product embedding for addition. The regular maps
take values in their specified algebraic targets since the
target equations vanish on their images. The group identities
hold on the torus by its vector-space quotient construction.
\end{proof}

\section{Relative smoothness and finiteness at complex levels}

\begin{proposition}\label{prop:complex-relative-jacobian}
Let \(f:X\to Y\) be an algebraic morphism of the nonsingular
complex projective varieties in
Corollary~\ref{cor:complex-level-algebraic-morphisms}. If its
complex derivative at a point is onto, it has a smooth
algebraic presentation there of relative dimension
\(\dim X-\dim Y\). This means a presentation over an affine
chart of \(Y\) by \(k\) variables and \(r\) equations after
inverting functions nonzero at the point, with an invertible
\(r\)-rowed square minor of their Jacobian in those variables
and \(k-r=\dim X-\dim Y\). At equal dimensions this is an
étale presentation.
\end{proposition}
\begin{proof}
Use Theorem~\ref{thm:smooth-algebraic-projection-chart} on source
and target. Write the source in coordinates \(x_1,\ldots,x_a,w\)
with \(G(x,w)=0,\ G_w\ne0\), and target in \(t_1,\ldots,t_b,l\)
with \(P(t,l)=0,\ P_l\ne0\). Denote the regular coordinate
expressions of \(f\) by \(F_i(x,w)\), \(L(x,w)\), inverting
their denominators on the source. The graph presentation over
the target has equations
\[
     G=0,\qquad t_i=F_i(x,w)\ (1\le i\le b),\qquad l=L(x,w),
\]
with all the chart denominators inverted. Substituting \(t_i=F_i\)
makes \(P(t,L)=0\) modulo \(G\), since the regular map takes
values in the target. The difference \(P(t,l)-P(t,L)\) is
\((l-L)\) times a polynomial in \(t,l,L\) whose value at
the given point is \(P_l\ne0\). Localize also at this polynomial.
There \(l=L\) follows from the first \(b+1\) equations and
the target relation. The Jacobian rows \(dG,dF_1,\ldots,dF_b\)
in the source coordinates are independent: \(G_w\ne0\) and
the last \(b\) rows restricted to the tangent kernel of \(dG\)
are independent by surjectivity of the derivative and the
local target coordinates \(t\). One minor is consequently
nonzero, hence a unit after a further localization. Clearing
unit denominators of the equations preserves this minor's
invertibility on their common zero set. This gives the
claimed presentation with \(k=a+1,r=b+1\).
If a chart is a point the same proof omits its variables and
equations; a map from a point is submersive only to a point.
\end{proof}

\begin{lemma}\label{lem:complex-projective-finite-fibers}
A projective algebraic morphism over \(\C\) whose complex
fibers are finite is finite. Here projective means that its
source is closed in a product of the base with some projective
space, followed by projection, and finite means that over
an affine neighborhood of every base point its preimage is
affine with coordinate ring a finite module over the base ring.
\end{lemma}
\begin{proof}
We give the homogeneous elimination argument. First, every
proper ideal in a finitely generated complex algebra has a
complex point annihilating it. Extend it to a maximal ideal
by the maximal-element principle of
Definition~\ref{def:field-conventions}. The residue field is
a finite-type complex domain; apply
Lemma~\ref{lem:complex-polynomial-normalization}. A field finite
as a module over a polynomial subring forces that subring
to be a field: for any nonzero \(a\) in the subring, \(a^{-1}\)
is integral by the same finite-module determinant argument.
Multiplying its monic equation by the necessary power of
\(a\) puts \(a^{-1}\) in the subring. The polynomial ring
must therefore have zero variables. The residue field is
then finite over \(\C\) and equals \(\C\) by
Lemma~\ref{lem:place-fields}. Furthermore, if a polynomial
\(p\) vanishes at all common zeros of an ideal \(I\), adjoining
an extra variable \(v\) gives \(1\in(I,1-vp)\) by this fact.
Substituting \(v=1/p\) in a finite expression for \(1\) and
clearing denominators gives \(p^a\in I\) for some \(a\ge1\),
with the zero polynomial immediate.

Work in an affine chart of the base near \(y_0\), and use
fiber coordinates \(z_0,\ldots,z_N\) in the projective space,
choosing \(z_0\) nonzero at each point of the finite fiber
by avoiding finitely many hyperplanes in the dual. The closed
source has homogeneous equations in these variables with
coefficients in the affine base ring; for the graphs considered
here they are the dehomogenized equations of
Theorem~\ref{thm:complex-level-projective-equations}.
At \(y_0\) and \(z_0=0\) these equations have no projective
common zero. The power containment just proved supplies a
positive power of each of \(z_1,\ldots,z_N\) in the specialized
ideal. In sufficiently large degree \(m\) every monomial in
these variables is divisible by one of these powers, by the
sum of their exponents. Thus homogeneous multiples of finitely
many of the specialized equations span the entire space of
degree \(m\) polynomials in these variables. One can take
homogeneous multiples because collecting the degree-\(m\)
part of any ideal expression preserves its degree-\(m\)
target.

The coefficient matrix for those multiples has full row
rank at \(y_0\). Invert a nonzero maximal minor on the base.
Solving by this minor now expresses every such monomial
over the base as a combination of the equations evaluated
at \(z_0=0\), in that same degree. It follows that throughout
this affine base neighborhood there are no projective points
of the source with \(z_0=0\). It also holds scheme-theoretically:
on each chart \(z_j=1,\ j>0\), the degree-\(m\) expression
for \(z_j^m\) makes the ideal with \(z_0=0\) the unit ideal.
The source over the neighborhood therefore lies in the affine
chart \(z_0=1\), closed there and hence affine. Lifting the
same degree-\(m\) expressions to the original homogeneous
equations, the differences all contain \(z_0\).
After dehomogenizing, every degree-\(m\) monomial is therefore
a linear combination of monomials of smaller total degree
in the affine coordinates over the base. Multiplying these
reductions and inducting on degree expresses all monomials
in those of degree below \(m\). Their finite list generates
the coordinate algebra as a module. The case \(N=0\) is a
closed immersion and its algebra is a quotient of the base
ring, generated as a module by one. This proves finiteness.
\end{proof}

\begin{theorem}\label{thm:actual-complex-finite-etale}
The projections \(\mathcal X(C',\nu)\to\mathcal X(C',0)\)
are smooth projective algebraic morphisms over \(\C\) of
relative dimension \(\nu fn\). The cover
\(\pi_m:\mathcal X(C_m,\nu)\to\mathcal X(C,\nu)\) of
Proposition~\ref{prop:complex-actual-hecke-legs} is finite
étale, of degree \([C:C_m]\) on every complex fiber. Each
of its unscaled Hecke legs \(u_{g,\mathbf 1}\) is finite
étale as well.
\end{theorem}
\begin{proof}
In the charts of Theorem~\ref{thm:analytic-framed-torus-family}
the first map is projection from the ball and its torus
coordinates to the ball, with surjective derivative and the
stated relative dimension. The cover is locally the identity
with the fiber count of
Proposition~\ref{prop:complex-actual-hecke-legs}. The unscaled
leg is locally the identity in the base and on its vector-space
fiber coordinates, with a change from the period lattice to
a larger lattice. Thus it too has invertible derivative at
each point. A local isomorphism from a compact manifold has
finite fibers: each fiber is closed and compact, and is
discrete by its local charts, hence finite. All these maps
are algebraic by Theorem~\ref{thm:holomorphic-map-algebraic-locally}.
They are projective since their graphs are closed in products
with a projective source embedding, by
Corollary~\ref{cor:compact-holomorphic-graphs-algebraic}.
Apply Proposition~\ref{prop:complex-relative-jacobian}
for the smooth and étale presentations, and
Lemma~\ref{lem:complex-projective-finite-fibers} for finiteness
of the last two maps.
\end{proof}

\section{Subspace kernels and the pointwise Frobenius polynomial}

\begin{lemma}\label{lem:actual-subspace-kernel-labels}
Fix a safe place \(w\) in
Theorem~\ref{thm:actual-central-test-vector}. Write \(O=O_w\),
\(\pi=\pi_w\), \(\kappa=O/\pi O\), \(q=|\kappa|\).
The summands in \(\tau_{w,r}\), \(0\le r\le n\), can be
labelled \(g_A C\) by the \(r\)-dimensional subspaces
\(A\subseteq\kappa^n\). In finite frame coordinates the
\(w\)-torsion on a complex torus is naturally the module
\(\pi^{-1}O^n/O^n\), identified with \(\kappa^n\) by the
chosen \(\pi\). The unscaled leg from \(x\) to \(xg_A\)
kills precisely \(A\) on this torsion.
\end{lemma}
\begin{proof}
Let \(d_r=\operatorname{diag}(\pi I_r,I_{n-r})\).
Choose \(k_A\in{\rm GL}_n(O)\) with
\(\ker(\bar d_r\bar k_A)=A\), by taking a residue basis
adapted to \(A\) and lifting an invertible matrix entrywise;
its determinant remains a unit. Set \(m_A=d_r k_A\) and
\(a_A=m_A^{\mathsf T}\). The lattice \(a_A O^n\) lies between
\(\pi O^n\) and \(O^n\), with residue subspace
\({\rm im}(\bar m_A^{\mathsf T})=A^\perp\) for the standard
dot product. This identity follows by pairing against the
kernel and comparing dimensions by matrix rank. The lattice
is the full preimage of that subspace since it contains
\(\pi O^n\). Every matrix in the double coset of \(d_r\)
similarly gives the full preimage of a codimension-\(r\)
subspace, and two matrices give the same lattice exactly
when they differ on the right by \({\rm GL}_n(O)\). The
orthogonal complement labels all such subspaces exactly once,
proving the parametrization of the local sum. Put
\(g_A=s(z_w^{-r})b_w(a_A)\).
The central scalar is supported over the conjugate quadratic
place, so at \(w\) the inverse of \(g_A\) is \(a_A^{\mathsf T}=m_A\).
Distinct labels also give distinct global cosets: equality
would give \(a_A^{-1}a_{A'}\in{\rm GL}_n(O)\) on applying
\(d_w\), by Definition~\ref{def:actual-positive-hecke} and
integrality at \(w\).

The torsion of \(V_\R/\Lambda_x\) is \(E^n/\Lambda_x\),
since the lattice spans \(E^n\) over \(\Q\). Its map
\(t\mapsto x^{-1}t\bmod\widehat O^n\) identifies it with
the finite adelic quotient \(\mathbb A_{E,f}^n/\widehat O^n\).
The kernel is zero by definition of the lattice. For surjectivity,
any target tuple has bounded denominators. After multiplication
by \(x\), its class modulo \(x\widehat O^n\) lies in
\(N^{-1}x\widehat O^n/x\widehat O^n\) for an integer \(N\ge1\).
This is the quotient of the completions of
\(N^{-1}\Lambda_x,\Lambda_x\), hence equals their rational
quotient by the finite Chinese remainder argument in
Lemma~\ref{lem:trace-finite}. The quotient splits as the
direct sum of its prime components. Killing it by all
elements of \(w\) leaves just \(\pi^{-1}O^n/O^n\) at \(w\):
locally \(w\) generates \((\pi)\) there and the unit ideal
at every other prime by Proposition~\ref{prop:general-ideal-factorization}.
The map on torsion for the inclusion
\(\Lambda_x\subseteq\Lambda_{xg_A}\) multiplies these finite
coordinates by \(g_A^{-1}\). The class of \(\tilde v/\pi\)
supported at \(w\) dies exactly when \(m_A\tilde v\equiv0
\bmod\pi\), the defining condition \(\bar v\in A\).
\end{proof}

\begin{proposition}\label{prop:line-cancellation-frobenius}
Let \(V\) be an \(n\)-dimensional vector space over a
field of \(q\) elements, \(n\ge1\), and fix a line \(W\subseteq V\).
Let \(Z(A)\) be an element of an abelian group \(B\) for
each subspace \(A\subseteq V\). If \(Z(A)=Z(A+W)\) for
every \(A\), then
\[
           \sum_{A\subseteq V}(-1)^{\dim A}
                      q^{\binom{\dim A}{2}} Z(A)=0 .
\]
The same implication holds when the value depends only on
\(A+W'\) for a specified nonzero subspace \(W'\subseteq V\).
\end{proposition}
\begin{proof}
Group the sum by \(D=A+W\). For fixed \(D\) of dimension
\(r\ge1\) the unique member containing \(W\) is \(A=D\).
The others are the complements of \(W\) in \(D\), of
dimension \(r-1\). There are precisely \(q^{r-1}\): relative
to one complement \(U\) each other one is the graph of a
linear map \(U\to W\), with independent values on a basis.
All the values in this group equal \(Z(D)\). Its coefficient is
\[
       (-1)^r q^{\binom r2}
        +q^{r-1}(-1)^{r-1}q^{\binom{r-1}2}=0
\]
since \(\binom r2=\binom{r-1}2+r-1\).
If \(W'\ne0\), choose any line inside it; adjoining that
line leaves the sum with \(W'\) unchanged.
\end{proof}

\begin{corollary}\label{cor:pointwise-congruence-single-line}
In the labeling of Lemma~\ref{lem:actual-subspace-kernel-labels},
let \(X_0,Z_0\) be sets, \(\phi:Z_0\to Z_0\) a map and
\(u_A:X_0\to Z_0\) maps. If for each \(x\in X_0\) there
is a line \(W_x\subseteq\kappa^n\) such that
\[
   \phi^{\,n-\dim A}u_A(x)=
           \phi^{\,n-\dim(A+W_x)}u_{A+W_x}(x)
                        \qquad(A\subseteq\kappa^n),
\]
then, in the free abelian group on \(Z_0\),
\[
      \sum_{r=0}^n(-1)^r q^{\binom r2}
            \sum_{\dim A=r}\big[\phi^{\,n-r}u_A(x)\big]=0
                         \qquad(x\in X_0).
\]
In particular this is the pointwise Frobenius congruence for
a reduction of the actual legs whenever the displayed hypothesis
is established there with \(\phi\) the \(q\)-power Frobenius
on geometric point coordinates.
\end{corollary}
\begin{proof}
The free abelian group on a set is the group of finite
integer-valued formal sums of its points, with \([z]\)
denoting its basis element at \(z\). Apply
Proposition~\ref{prop:line-cancellation-frobenius} separately
at each \(x\) to \(Z(A)=[\phi^{n-\dim A}u_A(x)]\), and
collect summands by dimension. The inner indexing agrees
with all summands of the safe spherical operator by
Lemma~\ref{lem:actual-subspace-kernel-labels}. The coefficients
and exponents are those of the operator polynomial in
Corollary~\ref{cor:congruence-quotient-spherical}.
\end{proof}

\section{Fields of definition for the projective Hecke diagrams}

\begin{lemma}\label{lem:fixed-projective-equations}
Fix ordered homogeneous coordinates on a projective space over
\(\C\), a closed algebraic set \(Z\), and a subfield
\(k\subseteq\C\). Take the homogeneous vanishing ideal of
\(Z\), that is, the direct sum of the spaces of homogeneous
forms vanishing on it. The following conditions are equivalent.
\begin{enumerate}
\item The ideal is generated over \(\C\) by homogeneous
polynomials with coefficients in \(k\).
\item Coordinatewise application of every
\(\sigma\in{\rm Aut}(\C/k)\) preserves \(Z\) as a set.
\end{enumerate}
In the chosen coordinates there is a smallest subfield of
\(\C\) containing \(k\) over which this ideal is generated.
It is finitely generated as a field over \(k\).
\end{lemma}
\begin{proof}
The fixed field of \({\rm Aut}(\C/k)\) is \(k\). Indeed, for
an element algebraic over \(k\) outside \(k\), its minimal
polynomial in characteristic zero has a distinct conjugate.
Sending the element to that conjugate extends to an
automorphism of the algebraic closure of \(k\) in \(\C\) by
the embedding argument of Lemma~\ref{lem:place-fields}.
Fix a transcendence basis of \(\C\) over this closure and
extend over its algebraic closure by the same argument.
Here a maximal algebraically independent set exists by
Definition~\ref{def:field-conventions}, and maximality
makes the remaining extension algebraic. For an element
transcendental over \(k\), include it in such a basis over
\(k\), replace it by its sum with one, fix the other basis
members and extend the resulting rational function field
automorphism to \(\C\). This moves the element as required.

In each degree \(m\), list the monomials in a fixed order.
The degree-\(m\) part of the ideal is a vector subspace
\(I_m\) of the corresponding finite complex coefficient
space. A subspace has a unique matrix of nonzero reduced
row echelon rows spanning it. Successively using the first
nonzero column on the subspace with all preceding pivot
coordinates zero determines its pivot positions. Projection
to the pivot coordinates is an isomorphism, and the inverse
images of the coordinate basis are the rows. These facts
also show that entrywise field automorphisms carry the
matrix to the corresponding matrix for the conjugate
subspace. Invariance of \(Z\) gives invariance of \(I_m\),
since evaluating a conjugated polynomial at a conjugated
point conjugates its value. The entries of its reduced
matrix then belong to \(k\) by the fixed-field assertion.

There is \(d\) such that the homogeneous ideal is generated
in degrees at most \(d\). For completeness, the exponents
of the leading monomials of its nonzero homogeneous elements
in degree-lexicographic order form an upward closed set.
By Lemma~\ref{lem:finite-degree-cut} it is generated by
finitely many exponents. Choose monic homogeneous elements
attaining them. Subtracting suitable monomial multiples
cancels the greatest monomial of any homogeneous element,
and terminates since there are only finitely many monomials
of its degree. This proves the assertion, with \(d=0\) for
the zero ideal. The echelon rows through degree \(d\) thus
suffice. Conversely, equations over \(k\) are carried to
themselves as equations by every such automorphism, giving
setwise invariance in both directions. Finally adjoin to
\(k\) the finitely many entries of the indicated rows.
They give generators over the resulting field. Any field of
definition containing \(k\) fixes the set under its own
automorphisms, so its fixed-field argument puts every entry
in that field. This proves minimality.
\end{proof}

\begin{proposition}\label{prop:finite-presentation-spread}
Let a finite diagram over \(\C\) consist of projective
algebraic sets with their reduced structures, regular
maps, and finitely many identities between composites of
these maps. Let \(R_0\subseteq\C\) be a subring. There is
a subring \(R\subseteq\C\) generated by finitely many
elements over \(R_0\) on which projective schemes and
morphisms with all those identities are defined, recovering
the given diagram by extending coefficients to \(\C\).
One may preserve any finite list of the relative smooth
or étale presentations of
Proposition~\ref{prop:complex-relative-jacobian} covering
specified maps, as well as finiteness for maps certified
by the projective elimination of
Lemma~\ref{lem:complex-projective-finite-fibers}.
\end{proposition}
\begin{proof}
Use finitely many homogeneous equations for each projective
source and target, as in
Lemma~\ref{lem:fixed-projective-equations}. A standard affine
chart sets a nonzero projective coordinate equal to one and
uses the corresponding dehomogenized equations. Its algebra
is the quotient by those equations. Within such a chart,
an open set \(D(a)\) means the locus where the regular
function \(a\) is invertible. Its algebra is presented by
adjoining \(v\) and \(av-1=0\). Such opens form a basis for
the polynomial topology since an open complement of common
zeros is a union of the opens of the individual equations.
Near each point a regular map into a target chart is given
by regular coordinate functions on some such open of a
source chart. On shrinking, write these functions as
polynomials in the coordinates and \(v\). The domains may
be chosen from a finite list. Indeed, in each affine chart
the functions specifying any covering by principal opens
generate the unit ideal by the maximal ideal argument in
Lemma~\ref{lem:complex-projective-finite-fibers}, and an
expression for one uses finitely many of the functions.
Refine the covers for the finitely many maps, composites
and identities by intersections, using that inverse images
of principal opens in an affine target chart are again
principal on each of these domains. A finite affine
cover of each necessary overlap suffices by the same
argument. Transitions between the standard projective
charts invert one additional coordinate and divide the
other coordinates by it.

All data are now finite polynomial lists. Include the
coefficients of the equations, of the maps and of the
covering functions. For each cover include coefficients of
an expression for one witnessing the cover. For each map
include a witness of ideal membership for every pulled-back
target relation in the localized source presentation.
Include the same kind of witnesses for equality of
coordinate formulas on overlaps, for transition identities
and for the specified composite identities. They exist
since equality in the quotient polynomial ring is exactly
ideal membership. Agreement of two maps to projective
space can also be tested by the cross products of their
homogeneous coordinates on their target chart domains,
since each list there has a unit coordinate. Pointwise
identities on the reduced complex domains are identities
of regular functions by the power-containment argument of
Lemma~\ref{lem:complex-projective-finite-fibers}. Take
\(R\) to contain all these coefficients. The very same
polynomial identities give covers on the affine schemes
over \(R\), maps into the designated target charts, and
agreement on their overlaps. Gluing gives the projective
diagram, with the stated base change and identities.

A specified relative Jacobian chart is an isomorphism of
localizations with a finite polynomial presentation in
which a specified determinant is a unit. Give the formulas
for that isomorphism and its inverse, including the inverse
of the determinant, and give the finite witnesses for both
composites and relations as above. This preserves the
charts and their covering property over \(R\), with exactly
the same relative dimensions. For a specified finite map
use the closed projective graph over the target. Include
its equations and the isomorphism from the original source
to the graph and back. Over the finitely many target opens
in the elimination proof of
Lemma~\ref{lem:complex-projective-finite-fibers}, retain the
coordinate changes, the unit minor witnesses and the
homogeneous reductions of all monomials in its specified
degree. These identities again establish that the graph
over each open lies in one affine chart and that lower
degree monomials generate its coordinate algebra as a
module. Enlarging \(R\) for this finite list gives finiteness
on the same covering. All choices involve only finitely
many actual complex coefficients or their indicated inverses,
so adjoining them produces a subring of the claimed form.
\end{proof}

\begin{corollary}\label{cor:actual-finite-type-hecke-model}
For every finite selection of levels, powers, covers and
Hecke legs of Theorem~\ref{thm:actual-complex-finite-etale},
together with any finite selection of their algebraic
operations and diagram identities, there is a simultaneous
projective model over a finitely generated \(O_E\)-subalgebra
of \(\C\). Its indicated projections are smooth, and the
indicated covers and unscaled legs are finite étale. Here
\(O_E\) is embedded by the chosen inclusion \(E\subseteq\C\).
In fixed projective coordinates, the equations of the
varieties and their graphs have field of definition contained
in \(E\) if and only if each of those embedded point sets
is invariant under \({\rm Aut}(\C/E)\).
\end{corollary}
\begin{proof}
The projective equations are supplied by
Theorem~\ref{thm:complex-level-projective-equations}, the
algebraic maps by
Corollary~\ref{cor:complex-level-algebraic-morphisms}, and
the covering Jacobian and finite presentations by
Theorem~\ref{thm:actual-complex-finite-etale}. Select finite
subcovers using the principal-open argument of
Proposition~\ref{prop:finite-presentation-spread} and apply
it with \(R_0=O_E\). For the last assertion place each
graph in a projective space by the product embedding of
Corollary~\ref{cor:compact-holomorphic-graphs-algebraic}.
That embedding has coordinates \(x_i y_j\), and hence
commutes with field automorphisms. Apply
Lemma~\ref{lem:fixed-projective-equations} with \(k=E\)
to each of the finitely many embedded sets.
\end{proof}


\section{Specialization of the labelled torsion subgroups}

\begin{lemma}\label{lem:finite-fiber-local-algebras}
Let \(k\) be an algebraically closed field, in any
characteristic, and \(D\) a nonzero finite-dimensional
commutative \(k\)-algebra with identity. Its points are the
\(k\)-algebra maps \(x:D\to k\). The algebra is the product
of finitely many local algebras \(D_x\), one at each point,
whose maximal ideals \(\mathfrak m_x\) are nilpotent.
The tangent space at \(x\), consisting of lifts of \(x\) to
\(k[\varepsilon]/(\varepsilon^2)\), identifies with the
\(k\)-dual of \(\mathfrak m_x/\mathfrak m_x^2\). If its
dimension is at most one and \(e_x=\dim_k D_x\), then
\[
                    D_x\simeq k[T]/(T^{e_x}).
\]
When \(e_x>1\) one may use any element of \(\mathfrak m_x\)
nonzero modulo its square as \(T\). Any quotient of this
local algebra of dimension \(j\), \(0\le j\le e_x\), has
kernel \((T^j)\).
\end{lemma}
\begin{proof}
Every maximal ideal gives a finite-dimensional residue
field over \(k\), hence exactly \(k\), since every minimal
polynomial splits. Distinct maximal ideals are comaximal.
For two comaximal ideals with \(a+b=1\) in their sum, the
elements \(a,b\) lift the two coordinate idempotents in
the residues, and their intersection is their product
by multiplying an element of both by \(a+b\). Induction
gives the Chinese remainder map onto the product for
any finite set of pairwise comaximal ideals. In this induction
an ideal comaximal to each of a list is comaximal to their
product by multiplying elements congruent to one modulo
that ideal. The number of distinct maximal ideals is thus
at most \(\dim_k D\).

Write \(J\) for their product. The descending powers of
\(J\) stabilize by dimension. If \(N=J^a=J N\), on a
finite generating list for \(N\) one obtains a matrix with
entries in \(J\) expressing the generators in terms of
themselves. As in Lemma~\ref{lem:place-fields}, the
determinant of the identity minus this matrix annihilates
\(N\). It is one modulo every maximal ideal, hence a unit
since a nonunit lies in a maximal ideal by
Definition~\ref{def:field-conventions}. Thus \(J^a=0\)
for sufficiently large \(a\). The \(a\)-th powers of
distinct maximal ideals are still comaximal by expanding
\((a_1+a_2)^{2a-1}=1\) for \(a_1+a_2=1\) in the two
ideals. Their product is zero. The Chinese remainder map
therefore identifies \(D\) with the product of the
quotients by these powers. Each quotient has nilpotent
maximal ideal and residue \(k\), and all elements outside
that ideal are units by a finite geometric sum.

The ring of dual numbers is local and its only idempotents
are zero and one, as is seen by reducing modulo its
nilpotent ideal. A lift of \(x\) therefore factors through
the \(x\)-factor. Writing its value at \(c+v\) as
\(c+\lambda(v)\varepsilon\) for \(c\in k\),
\(v\in\mathfrak m_x\), multiplicativity is exactly
\(\lambda(\mathfrak m_x^2)=0\) for the linear functional
\(\lambda\). If \(\mathfrak m_x/\mathfrak m_x^2=0\),
nilpotence gives \(\mathfrak m_x=0\). Otherwise choose
the stipulated \(T\). In \(D_x/(T)\) the image of the
maximal ideal equals its square, so nilpotence again
makes it zero. Consequently \(\mathfrak m_x=(T)\).
Successive subtraction of a residue and division in
this principal ideal express every element as a polynomial
in \(T\). If \(T^b=0\) with \(b\) minimal, a nonzero
polynomial of degree less than \(b\) has value \(T^j\)
times a unit for some \(j<b\), by its first nonzero
coefficient. It cannot vanish. Thus \(b=e_x\), and
each nonzero ideal is generated by the lowest power
occurring among its nonzero elements. The zero ideal
is \((T^{e_x})\). This proves the quotient assertion,
including \(e_x=1\) with \(T=0\).
\end{proof}

\begin{proposition}\label{prop:valuation-special-cluster-lengths}
Let \(V\) be a valuation ring with fraction field \(K\),
maximal ideal \(\mathfrak m\), and algebraically closed
residue field \(k\). A valuation ring here is a domain
in which, for nonzero \(a,b\in K\), at least one of
\(a/b,b/a\) belongs to \(V\). Let \(S\) be a finite free
commutative \(V\)-algebra such that
\(K\otimes_V S=K^\Omega\) with a specified finite set
\(\Omega\). Every projection indexed by \(\omega\in\Omega\)
restricts on \(S\) to a \(V\)-valued point. Let
\(\bar\omega\) be its reduction. For \(A\subseteq\Omega\)
put \(S_A={\rm im}(S\to K^A)\), allowing the zero algebra
when \(A\) is empty. Then \(S_A\) is finite free of rank
\(|A|\), and for each point \(x\) of the special algebra
\(D=S/\mathfrak m S\) the local factor at \(x\) of
\(S_A/\mathfrak m S_A\) has dimension
\[
                     |\{\omega\in A:\bar\omega=x\}|.
\]
An absent factor is counted as zero. In particular reduction
from \(\Omega\) onto the points of \(D\) is surjective.
\end{proposition}
\begin{proof}
Multiplication on the finite free module gives a monic
equation for every element of \(S\) by the adjugate
argument of Lemma~\ref{lem:complex-polynomial-normalization}.
An element of \(K\) integral over \(V\) belongs to \(V\).
Otherwise its inverse would be a nonunit of \(V\), and
division of its monic equation by the highest power
would put one in \(\mathfrak m\). Evaluation therefore
takes values in \(V\). The finite \(V\)-module \(S_A\)
is torsion free as a submodule of \(K^A\). Any finite
generating list in a torsion-free module can be reduced
to a basis here. A nontrivial \(K\)-linear relation after
inverting nonzero elements of \(V\) can be divided by
one of its nonzero coefficients so all resulting
coefficients lie in \(V\) and one equals one, by total
divisibility in a valuation ring. The relation already
holds in the module by torsion freeness, so that generator
can be removed. Iteration leaves an independent spanning
list. Inverting nonzero elements of \(V\) in \(S_A\)
gives all of \(K^A\), by the given generic product
and coordinate projection. This proves its rank.

By Lemma~\ref{lem:finite-fiber-local-algebras}, choose
\(\bar t\in D\) with distinct residue values \(c_x\in k\)
at all its points, by taking these constants in the
product factors. The field \(k\) is infinite, since on
a finite field the polynomial one plus the product of
\(X-c\) over all \(c\) has no root. Lift \(\bar t\) to
\(t\in S\). On \(S_A\) the characteristic polynomial
of multiplication by \(t\) becomes
\(\prod_{\omega\in A}(Y-t(\omega))\) over \(K\).
Both sides already have coefficients in \(V\), so they
are equal there. The special algebra of \(S_A\) is a
quotient of \(D\). It is a product of the corresponding
quotients of the local factors of \(D\), by the images
of their idempotents. In the factor at \(x\) multiplication
by \(\bar t-c_x\) is nilpotent. A nilpotent map has a
basis in which its matrix is strictly triangular, by
extending bases along its increasing sequence of kernels.
The characteristic polynomial on this factor is
\((Y-c_x)^{d_x}\), where \(d_x\) is its dimension.
Reducing the characteristic polynomial over \(V\) and
comparing multiplicities of the pairwise distinct roots
gives the formula. Take \(A=\Omega\) for surjectivity,
using the injectivity of \(S\to K\otimes_V S\).
\end{proof}

\begin{theorem}\label{thm:line-in-specializing-torsion}
Assume the hypotheses of
Proposition~\ref{prop:valuation-special-cluster-lengths}.
Suppose \({\rm Spec}(S)\) is a commutative group over \(V\)
with a scalar action of a finite field \(\kappa\), and
identify its generic points as a \(\kappa\)-space
\(\Omega\simeq\kappa^n\). A group over \(V\) here means that
multiplication, inverse and unit are morphisms over \(V\)
satisfying the group identities, with the scalar maps
satisfying the module identities. If the special tangent
space at the identity has dimension one over \(k\), then
the kernel \(W\) of reduction on generic labels is a
nonzero \(\kappa\)-subspace. For every point \(x\) of the
special group and every subspace \(A\subseteq\Omega\)
the ideal of the closure of \(A\) on the special local
factor at \(x\) is
\[
                 (T_x^{j_A(x)}),\qquad
 j_A(x)=
 \begin{cases}
       |A\cap W|,&x\in\bar A,\\
       0,&x\notin\bar A ,
 \end{cases}
\]
where \(\bar A\) is the reduction image, and \(T_x\) is
any generator of the local maximal ideal. Each such
local algebra is \(k[T_x]/(T_x^{|W|})\).
\end{theorem}
\begin{proof}
Composition of the \(V\)-points and reduction commute with
the group operations and scalar maps since these morphisms
are defined over \(V\). Thus reduction is linear on labels,
and each nonempty fiber on \(A\) is a translate of
\(A\cap W\). By
Proposition~\ref{prop:valuation-special-cluster-lengths}
all special points occur, and each full local factor has
dimension \(|W|\). Translation by a special point, with
inverse translation by its negative, identifies its tangent
space with that at the identity. By
Lemma~\ref{lem:finite-fiber-local-algebras} each local
factor is the asserted truncated polynomial algebra. At
the identity its maximal ideal has a nonzero class modulo
its square, so its length is at least two. Hence \(W\)
contains a nonzero vector, and therefore a \(\kappa\)-line.
The closure of the generic labels \(A\) is cut out by the
elements vanishing on all those labels, giving exactly
the quotient \(S_A\). Its local quotient dimension is
\(j_A(x)\) by
Proposition~\ref{prop:valuation-special-cluster-lengths}.
Lemma~\ref{lem:finite-fiber-local-algebras} determines
its ideal from that dimension, including the unit ideal
when the local factor is absent.
\end{proof}


\section{Frobenius powers on a torsion component}

\begin{lemma}\label{lem:frobenius-pulls-kernel-ideal}
Let \(k\) be an algebraically closed field of characteristic
\(p>0\), \(P=p^a\) with \(a\ge0\), \(Y\) a scheme over \(k\),
and \(y\in Y(k)\) a closed point. Write \(Y^{(P)}\) for
change of scalars by \(c\mapsto c^P\) on \(k\). The relative
Frobenius \(F_{Y,P}:Y\to Y^{(P)}\) is defined on affine
coordinate rings by
\[
     R\otimes_{k,\,c\mapsto c^P}k\longrightarrow R,\qquad
                       r\otimes c\longmapsto r^P c.
\]
For a \(k\)-map \(v:H\to Y\) from a finite \(k\)-scheme
and a point \(x\in H(k)\), let \(I_x\) be the ideal in
the local factor \(D_x\) of \(H\) cutting out \(v^{-1}(y)\).
The ideal in \(D_x\) cutting out
\((F_{Y,P}v)^{-1}(F_{Y,P}(y))\) is generated by the
\(P\)-th powers of the elements of \(I_x\).
In particular, if \(I_x=(T^j)\), this ideal is \((T^{Pj})\),
also when \(j=0\).
\end{lemma}
\begin{proof}
The formula respects the tensor relations, addition and
products by the characteristic-\(p\) binomial identity,
iterated \(a\) times. It commutes with localization since
inverting a function is equivalent to inverting its \(P\)-th
power, so defines the indicated morphism. The field map
\(c\mapsto c^P\) is bijective by algebraic closedness and
injectivity of a field homomorphism. The map on the underlying
point spaces is bijective. In an affine chart the scalar
twist is isomorphic as a ring to \(R\), by
\(r\otimes c\mapsto r c^{1/P}\). In that identification
the displayed map raises elements to their \(P\)-th powers,
and \(r^P\) belongs to a prime ideal precisely when \(r\)
does. The argument respects the affine opens.

If \(v(x)\ne y\), both local ideals are the unit ideal.
Otherwise an affine neighborhood of \(y\) contains the
whole image of \({\rm Spec}(D_x)\), which has just one
point by Lemma~\ref{lem:finite-fiber-local-algebras}. If
\(\mathfrak n\) is the ideal of \(y\) in this chart, the
ideal of \(F_{Y,P}(y)\) in its twist is generated by
\(r\otimes1\) for \(r\in\mathfrak n\). The quotient is
\(k\otimes_{k,\,c\mapsto c^P}k\simeq k\) by the displayed
evaluation. The new pullback ideal is generated by
\(v^*(r)^P\) for these \(r\). This also contains the
\(P\)-th power of any element of \(I_x\), by raising its
expression as a linear combination of the \(v^*(r)\)
to that power. Conversely each generator is already of
this form. A principal ideal gives the asserted generator.
\end{proof}

\begin{proposition}\label{prop:flat-subkernel-generic-labels}
Use \(V,K,k\) as in
Proposition~\ref{prop:valuation-special-cluster-lengths}.
Let \(G\) be a commutative group scheme over \(V\) and
let \(H\subseteq G\) be a finite free closed subgroup
whose generic algebra is a product of copies of \(K\).
Let \(\psi:G\to G'\) be a homomorphism. If its kernel
on \(H\) is finite flat over \(V\) and its generic points
in \(H\) form a subset \(A\), then that kernel on \(H\)
is the closure of \(A\). On the special fiber its ideal
at \(x\in H(k)\) is therefore determined by the cluster
length in Proposition~\ref{prop:valuation-special-cluster-lengths}
whenever the tangent dimension at \(x\) is at most one.
\end{proposition}
\begin{proof}
Write \(S\) for the algebra of \(H\). The kernel is assumed
a closed subgroup, with coordinate quotient \(S/J\).
Flatness means that tensoring with this module preserves
injections. Applying this to multiplication by a nonzero
element of \(V\) on \(V\) proves torsion freeness, hence
\(S/J\) injects into its generic algebra. An ideal in
a product of finitely many fields is the product of its
coordinate ideals, by multiplying by the coordinate
idempotents. Its quotient here is exactly \(K^A\).
Thus \(J=\ker(S\to K^A)\), which is the stated closure
ideal. Reduction gives its local quotient lengths by
Proposition~\ref{prop:valuation-special-cluster-lengths}.
The assertion about the ideal follows from
Lemma~\ref{lem:finite-fiber-local-algebras}.
\end{proof}

\begin{theorem}\label{thm:one-place-frobenius-kernel}
Suppose \(k\) has characteristic \(p>0\). Let \(H\) be a
finite free commutative group scheme over \(V\) with split
generic algebra as in
Proposition~\ref{prop:valuation-special-cluster-lengths},
and special tangent dimension at the identity at most one.
Inside its generic point group fix a subgroup
\(L\simeq\kappa^n\) for a finite field \(\kappa\) of
cardinality \(q=p^e\), \(e\ge1\). Assume reduction
\(\nu:L\to H(k)\) respects the \(\kappa\)-scalar action on
its image. For each subspace \(A\subseteq L\) let
\(\psi_A:H\to Y_A\) be a homomorphism to a group scheme
over \(V\) with closed identity, having finite flat kernel
and generic kernel labels \(A\).
Put \(r=\dim_\kappa A\), let \(\bar\psi_A\) denote the
special map, and set
\[
               f_A=F_{(Y_A)_k,\,q^{\,n-r}}\bar\psi_A .
\]
Its kernel on the special fiber of \(H\) has, at \(x\in H(k)\),
local ideal
\[
   \begin{cases}
      (T_x^{\,q^n/|\nu(A)|}),&x\in\nu(A),\\
      (1),&x\notin\nu(A),
   \end{cases}
\]
where \(T_x\) generates the maximal ideal in the local
algebra at \(x\). In particular the special kernels of
\(f_A,f_B\) agree as closed subschemes whenever
\(\nu(A)=\nu(B)\). If \(W=\ker\nu\) contains a line
\(W_0\), this gives kernel equality for the pair
\(A,A+W_0\) occurring in
Corollary~\ref{cor:pointwise-congruence-single-line}.
\end{theorem}
\begin{proof}
Translations by \(k\)-points identify the special tangent
spaces at all points with the one at the identity. Fix a
generator \(T_x\) in the resulting truncated polynomial
algebra of Lemma~\ref{lem:finite-fiber-local-algebras},
taking zero when the algebra is \(k\).
Proposition~\ref{prop:flat-subkernel-generic-labels}
gives the ideal \((T_x^j)\) of \(\bar\psi_A^{-1}(0)\)
with \(j=|\{a\in A:\nu(a)=x\}|\). For a nonempty
fiber, translation by a preimage identifies it with
\(A\cap W\). Partitioning \(A\) into its fibers therefore
gives \(j=|A|/|\nu(A)|\). Since \(|A|=q^r\), application
of Lemma~\ref{lem:frobenius-pulls-kernel-ideal} gives
the exponent \(q^{n-r}j=q^n/|\nu(A)|\). A power exceeding
the nilpotence index simply gives the zero ideal. An empty
fiber gives the unit ideal before and after Frobenius.
Equality of these ideals in all the factors is equality
in their product, the entire finite coordinate algebra by
Lemma~\ref{lem:finite-fiber-local-algebras}. Finally
\(\nu(A+W_0)=\nu(A)\) for any \(W_0\subseteq\ker\nu\).
\end{proof}


\section{The signature and the split factors of the isogenies}

\begin{lemma}\label{lem:degree-one-split-lie-polynomial}
In the notation of Lemma~\ref{lem:actual-signs-and-padding},
identify \(E\) with its image under \(\sigma\), and write
\(g_\tau=\sigma^{-1}\tau\). Choose an integral primitive
element \(\theta\) of \(E/\Q\). Outside a finite set of
rational primes, if \(w\) is a degree-one prime of \(E\)
over \(p\), then \(p\) splits completely in \(E\), that is,
\(pO_E\) is a product of \(2f\) distinct primes with
residue field \(\mathbf F_p\). Write \(\rho:O_E\to
\mathbf F_p\subseteq k\) for reduction at \(w\) into an
algebraically closed field \(k\) of characteristic \(p\).
For \(m\ge1\), the factors of \(O_E/p^mO_E\) are copies
of \(\Z/p^m\Z\), indexed by
\(v_g=\ker(\rho\circ g)\), \(g\in{\rm Gal}(E/\Q)\).
On the tangent space at the identity of every complex
torus of Proposition~\ref{prop:signed-ball-period-coordinates},
the polynomial of \(\theta\) is the image in \(\C[T]\) of
\[
    P_\theta(T)=(T-\theta)(T-\bar\theta)^{n-1}
              \prod_{\tau\in\Psi\setminus\{\sigma\}}
                             (T-g_\tau\theta)^n\in O_E[T].
\]
If \(u=w|_{K_0}\) for the imaginary quadratic subfield,
the roots indexed by \(v_g\) above \(u\) have multiplicity
one at \(w\) and zero at all other such primes. Above
\(\bar u\) the multiplicities are \(n-1\) at \(\bar w\)
and \(n\) at each remaining prime.
\end{lemma}
\begin{proof}
The generators \(\sqrt{-d},\zeta_{N_1}\) have all their
conjugates in \(E\). Every embedding of this finite field
preserves it and the \(2f\) embeddings are the automorphisms,
as follows by adjoining successively the roots of their
separable polynomials as in Lemma~\ref{lem:place-fields}.
They commute since on the first generator each changes at
most its sign and on the second each raises it to a unit
power. An integral primitive element and the finite index
of its power basis are provided by
Lemma~\ref{lem:number-field-lattice}. Exclude \(2\), the
primes dividing that index and the nonzero discriminant
of the minimal polynomial. Its \(2f\) conjugate roots
then have distinct residues under \(\rho\), since the
product of squared differences is the discriminant.
All residues belong to \(\mathbf F_p\) by the degree-one
assumption. For every \(m\) these roots lift uniquely to
roots modulo \(p^m\) by the digit correction in
Lemma~\ref{lem:marked-unramified-tower}. Their pairwise
differences are units, so the minimal polynomial modulo
\(p^m\) is the product of the corresponding linear
factors. Indeed division by each monic factor preserves
the other roots by invertibility of the differences.
Since the power basis index is prime to \(p\), it is a
basis modulo \(p^m\) by its invertible determinant.
The Chinese remainder argument of
Lemma~\ref{lem:finite-fiber-local-algebras} gives the
stated product with its residue maps indexed by \(g\).
In particular \(p\) is a uniformizer in each completion
at these primes by the valuation factorization of
Proposition~\ref{prop:general-ideal-factorization}.

Proposition~\ref{prop:signed-ball-period-coordinates}
computes the polynomial on the complex tangent space.
The normality just proved makes every \(g_\tau\) an
automorphism of \(E\); it preserves integrality by
transporting monic integral equations. For
\(\tau\ne\sigma\) in \(\Psi\), \(g_\tau\) acts by
conjugation on \(K_0\) by Lemma~\ref{lem:actual-signs-and-padding}.
Together with complex conjugation these are all the \(f\)
automorphisms with that restriction, since \(\Psi\) has
one member from each conjugate pair. The maps \(\rho\)
and \(\rho\circ{\mathfrak c}\) are distinct on \(O_{K_0}\)
outside the finite set of primes ramifying there, which
we also exclude. This identifies \(u,\bar u\) and proves
the multiplicity assertions.
\end{proof}

\begin{lemma}\label{lem:valuative-torsion-idempotents}
Let \(V\) be as in Proposition~\ref{prop:valuation-special-cluster-lengths},
with \(K=\C\), equipped with \(O_E\to V\) inducing the
chosen inclusion into \(\C\) and the map \(\rho\) above
in its residue field \(k\). Let \(G\) be a smooth
commutative group scheme over \(V\) with an \(O_E\)-action,
whose generic group with action is one of the complex
tori of Lemma~\ref{lem:degree-one-split-lie-polynomial}.
Suppose the kernel of multiplication by \(p^m\) is
finite flat over \(V\). Then
\[
                   G[p^m]=\prod_{v\mid p} H_{v,m}
\]
with finite free closed factors cut out by the respective
idempotents of \(O_E/p^mO_E\). The special tangent dimension
of \(H_{v,m}\) equals the multiplicity attached to \(v\)
in Lemma~\ref{lem:degree-one-split-lie-polynomial}.
Here smoothness is used in the relative Jacobian
presentation of Proposition~\ref{prop:complex-relative-jacobian}.
\end{lemma}
\begin{proof}
At the identity point of the special fiber take a smooth
chart. The identity section factors through it over \(V\):
the inverse image is an open containing the closed point
of a local spectrum, hence contains a principal neighborhood
of it defined by a unit. Tangent vectors along that section
are solutions of the linearized chart equations over \(V\).
The unit Jacobian minor uniquely solves for the corresponding
coordinates in terms of the others. Thus the tangent module
is finite free, and its formation commutes with passage
to both \(K\) and \(k\). The derivative of \(\theta\) on
it has characteristic polynomial \(P_\theta\) over \(V\),
since both sides agree over \(K\) by
Lemma~\ref{lem:degree-one-split-lie-polynomial} and
\(V\subset K\). In particular the special polynomial is
its reduction.

On the kernel the scalar action factors through
\(O_E/p^mO_E\). For each ring idempotent \(e_v\) of the
product in Lemma~\ref{lem:degree-one-split-lie-polynomial}
let \(H_{v,m}\) be the equalizer of that endomorphism and
the identity on the kernel. The kernel is affine finite,
so the equalizer is closed, defined on its algebra by all
differences \(s-e_v^*s\). This quotient identifies with
the image of the idempotent algebra map \(e_v^*\): the
kernel of that map consists of differences with zero
image, and every difference is killed. Hence its algebra
is a finite torsion-free \(V\)-module, using flatness of
the original kernel as in
Proposition~\ref{prop:flat-subkernel-generic-labels}.
It is free by the generating-list argument of
Proposition~\ref{prop:valuation-special-cluster-lengths}.
The maps \(z\mapsto(e_vz)_v\) and \((z_v)_v\mapsto\sum_v z_v\)
are mutual inverse morphisms by the orthogonal idempotent
identities. This proves the product on the whole group
scheme, after any base change as well.

The derivative of the group operation at the identity
is addition. Its restrictions to both coordinate axes
are the identity, and its derivative is linear by the
linearized equations on dual numbers. Hence multiplication
by \(p^m\) has derivative zero on the special tangent
space. All tangent vectors there belong to the kernel,
by evaluating that map on the dual numbers. The product
idempotents split the tangent into their corresponding
images. They act through \(O_E/pO_E\) on it, since it is
a \(k\)-vector space and the derivative of an integer
scalar is that integer times the identity. On the \(v_g\)
image the derivative of \(\theta\) is scalar
\(\rho(g\theta)\) by the product description.
These scalars are pairwise distinct. Their multiplicities
in the reduced characteristic polynomial are therefore
exactly the dimensions of the summands, proving the claim.
\end{proof}

\begin{theorem}\label{thm:actual-quadratic-half-kernels}
Take the data of Lemma~\ref{lem:valuative-torsion-idempotents}
with \(G\) separated, \(w\) safe, and finite flat multiplication kernels
for \(m=1,n\). Suppose all the unscaled complex fiber isogenies
labelled \(A\subseteq\mathbf F_p^n\) in
Lemma~\ref{lem:actual-subspace-kernel-labels} extend over
\(V\) to \(O_E\)-linear homomorphisms
\(\psi_A:G\to G_A\) with finite flat kernels and separated
target groups. Put \(r=\dim_{\mathbf F_p}A\) and
\(f_A=F_{(G_A)_k,\,p^{\,n-r}}(\psi_A)_k\).
There exists a line \(W_0\subseteq\mathbf F_p^n\) such that
the kernels of \(f_A,f_{A+W_0}\) agree on
\(\prod_{v\mid u}(H_{v,n})_k\) for every \(A\), where
\(u=w|_{K_0}\).
\end{theorem}
\begin{proof}
Use uniformizer \(p\) at \(w\) in
Lemma~\ref{lem:actual-subspace-kernel-labels}.
The generic \(w\)-part killed by the isogeny is exactly
\(A\) even among all \(p^n\)-torsion labels. In the finite
coordinates of that lemma its matrix is \(d_r k_A\), whose
inverse has denominators dividing \(p\). Thus any killed
class is already killed by \(p\), and the computation
modulo \(p\) applies. On the other components above \(u\)
the matrix is the identity, as follows from
Lemma~\ref{lem:safe-operators-away-from-markers}: the
local \(b_w(a_A)\) is supported on \(w,\bar w\) and the
central scalar on the places above \(\bar u\). At
\(\bar w\) the inverse frame matrix is \(p^r\bar a_A^{-1}\).
Its own inverse
has denominators dividing \(p^r\) since \(a_A\) is
integral. At the other places above \(\bar u\) the inverse
frame is scalar \(p^r\), and it is a unit at all remaining
places. These formulas use that the local norm of the
chosen uniformizer is \(p\) in the degree-one completion.
They show that \(p^n\) kills the entire generic kernel
on points, including \(r=0\). This holds on the generic
kernel scheme: its local tangent spaces vanish because
the complex isogeny is locally invertible by
Theorem~\ref{thm:actual-complex-finite-etale}, and its
finite complex algebra is reduced by
Lemma~\ref{lem:finite-fiber-local-algebras}.

Multiplication by \(p^n\) on the integral kernel of
\(\psi_A\) is consequently zero. Separatedness of \(G\)
makes the identity section closed, by pulling back its
closed diagonal along the identity map and the constant
identity-section map. Thus \(G[p^n]\) is closed.
Pulling back the
ideal of \(G[p^n]\) to the finite flat kernel gives
zero by injection of its algebra into its generic fiber.
Decomposing this kernel by the same \(O_E/p^nO_E\)
idempotents as in Lemma~\ref{lem:valuative-torsion-idempotents}
gives finite free factors. Their intersections with
\(H_{v,n}\) are exactly these factors, by the fixed-point
description for the idempotents. Their generic labels
above \(u\) were just computed.

The special tangent dimension of \(H_{w,1}\) and of
\(H_{w,n}\) is one by
Lemma~\ref{lem:valuative-torsion-idempotents}. The generic
points of \(H_{w,1}\) are \(\mathbf F_p^n\) by
Lemma~\ref{lem:actual-subspace-kernel-labels}; their
algebra splits since \([p]\) on the generic group has
invertible derivative and \(\C\) is algebraically closed.
Now Theorem~\ref{thm:line-in-specializing-torsion}
supplies a line \(W_0\) reducing to zero. Use
Theorem~\ref{thm:one-place-frobenius-kernel} on \(H_{w,n}\)
with \(L=H_{w,1}(\C)\), the generic labels just determined
and their finite flat kernels on this component. This
identifies the kernels for \(A,A+W_0\) there. At the other
\(v\mid u\) each generic component kernel is the identity.
Its closure is the identity section, by
Proposition~\ref{prop:flat-subkernel-generic-labels}.
The special tangent dimension there is zero by
Lemma~\ref{lem:valuative-torsion-idempotents}, so all local
algebras are copies of \(k\) by
Lemma~\ref{lem:finite-fiber-local-algebras}. Applying
Lemma~\ref{lem:frobenius-pulls-kernel-ideal} preserves
both the zero and the unit ideals in these fields.
Finally the kernel of a linear map on this product
splits into its component kernels: apply each idempotent
to an element in the kernel over any \(k\)-algebra,
using the \(O_E\)-linearity of the reduced map and the
naturality of relative Frobenius for the scalar maps.
This proves the asserted equality.
\end{proof}


\section{The local algebras of the Frobenius kernels}

\begin{lemma}\label{lem:smooth-group-frobenius-coordinates}
Let \(G\) be a separated smooth commutative group scheme
of finite type and dimension \(g\) over an algebraically
closed field \(k\) of characteristic \(p>0\). For \(a\ge1\)
put \(P=p^a\) and \(J_a=\ker F_{G,P}\) for the relative
Frobenius of Lemma~\ref{lem:frobenius-pulls-kernel-ideal}.
Then \(J_a\) is supported at the identity and has coordinate
algebra
\[
                  k[t_1,\ldots,t_g]/(t_1^P,\ldots,t_g^P).
\]
Multiplication by \(P\) is zero on \(J_a\) as a morphism.
In particular \(J_a\subseteq G[P]\) as a closed subgroup.
\end{lemma}
\begin{proof}
Take a relative Jacobian chart as in
Proposition~\ref{prop:complex-relative-jacobian}, now with
base \(k\). Subtract the coordinates of the identity and
write the remaining variables as \(t,z\), with \(g\) free
variables \(t\) and a square invertible derivative in \(z\)
for the defining equations. In each quotient by the \(b\)-th
power of the maximal ideal, for \(b\ge1\), the equations
have a unique solution \(z=z_b(t)\) with zero constant term
in \(k[t]/(t)^b\). Induct from \(b=1\). A correction of
degree \(b\) multiplies the constant derivative matrix,
so its inverse gives the unique correction killing the
homogeneous degree-\(b\) error. The same inductive identities
in the original chart modulo its maximal ideal powers show
that its variables \(z\) equal these polynomials in \(t\)
to each order. Denominators inverted in the chart have
nonzero constant and so have inverses modulo each such
power by finite geometric sums. The inverse maps on
truncated rings identify their inverse limits. Thus at
the identity the formal functions are \(k[[t_1,\ldots,t_g]]\).
The same construction for products gives independent tuples
of formal parameters for each factor.

The Frobenius has a single point above the identity by
Lemma~\ref{lem:frobenius-pulls-kernel-ideal}. Its ideal in
an affine neighborhood of that point is generated by the
\(P\)-th powers of the elements vanishing there. This
already cuts out a finite local algebra: the maximal ideal
in an affine chart is generated by its finitely many affine
coordinates minus their values, with any presentation
denominators treated as additional affine coordinates.
If there are \(h\) such generators, every product of
\(h(P-1)+1\) generators contains a \(P\)-th power. Thus
one can compute the quotient in the truncated rings
just described. In formal coordinates its ideal is
\((t_1^P,\ldots,t_g^P)\), since \(P\)-powering is additive
and all parameter multiples have their \(P\)-th powers
in this ideal. This proves the algebra formula and shows
that \(J_a\) is finite of length \(P^g\).

In the parameters at the identity, express the sum of
\(p\) points by \(g\) formal series in \(p\) independent
tuples \(t^{(i)}\). The constant terms are zero and the
series are invariant under cyclic permutation of the
tuples, by associativity and commutativity of the group
law. Identify all the tuples with \(t\). Monomials in a
nonconstant orbit of the cyclic permutation occur in
orbits of size \(p\) with equal coefficients and hence
cancel in characteristic \(p\). In a fixed monomial
the full multi-index is the same on every tuple, giving
a monomial in \(t_1^p,\ldots,t_g^p\). This argument is
degree by degree with finite sums. Each coordinate of
\([p]\) therefore lies in \(k[[t_1^p,\ldots,t_g^p]]\)
with zero constant term. Iterating the substitution \(a\)
times puts each coordinate of \([P]\) in the series ring
in \(t_1^P,\ldots,t_g^P\) with zero constant term.
It vanishes in the algebra calculated above. Thus the
morphism \([P]\) restricted to the entire Frobenius kernel
is zero. Since \(G\) is separated its identity is closed,
and both kernels are closed subgroups.
\end{proof}

\begin{theorem}\label{thm:split-frobenius-kernel-ranks}
Let \(G/V\) satisfy Lemma~\ref{lem:valuative-torsion-idempotents}
with \(m=a\ge1\) and \(G\) separated. Put \(P=p^a\).
For \(v\mid p\) let \(d_v\) be the tangent multiplicity
in Lemma~\ref{lem:degree-one-split-lie-polynomial}, and
\(J_{v,a}=J_a\cap (H_{v,a})_k\) in \(G_k\). Then
\[
  J_a=\prod_{v\mid p}J_{v,a},\qquad
  \mathcal O(J_{v,a})\simeq
       k[t_1,\ldots,t_{d_v}]/(t_1^P,\ldots,t_{d_v}^P).
\]
The full special factor \((H_{v,a})_k\) has length \(P^n\).
For \(d_v=n\) it equals \(J_{v,a}\); for \(d_v=0\)
the Frobenius kernel there is the identity section.
\end{theorem}
\begin{proof}
The group \(J_a\) in the statement is the kernel for \(G_k\).
It is killed by \(P\) by
Lemma~\ref{lem:smooth-group-frobenius-coordinates} and is
stable under each scalar endomorphism by the affine formula
for relative Frobenius in
Lemma~\ref{lem:frobenius-pulls-kernel-ideal}. Apply the
idempotents of \(O_E/PO_E\) on this subgroup, using the
equalizer construction and the sum inverse in
Lemma~\ref{lem:valuative-torsion-idempotents}. Its factors
are exactly the indicated intersections. Each coordinate
algebra \(B_v\) is a quotient of the local algebra of \(J_a\)
and is local with residue \(k\). Every element of its
maximal ideal has \(P\)-th power zero by the explicit
algebra of Lemma~\ref{lem:smooth-group-frobenius-coordinates}.
The tangent of \(J_a\) is the full tangent space of \(G_k\),
since \(P\ge2\) and all the defining powers vanish on dual
numbers. Thus the tangent of its \(v\)-factor has dimension
\(d_v\), by the idempotent decomposition and the proof of
Lemma~\ref{lem:valuative-torsion-idempotents}.

Choose \(d_v\) lifts of a basis of the maximal ideal modulo
its square in \(B_v\). They generate the ideal: in the
quotient by them the remaining maximal ideal equals its
square and is nilpotent by Lemma~\ref{lem:finite-fiber-local-algebras}.
They also generate \(B_v\) as a \(k\)-algebra. Subtracting
residues expresses an element successively as a polynomial
of degree less than \(j\) in the lifts plus an element of
the \(j\)-th power of the maximal ideal. Each next power
is generated by products of \(j\) lifts, so subtracting
the residue of each coefficient supplies the next step.
Nilpotence terminates the process. Sending the variables
of the displayed truncated algebra to the lifts is now a
surjection, which gives \(\dim_k B_v\le P^{d_v}\). The
products of all these dimensions equal \(P^{fn}\):
coordinate algebras of the product of finite affine
schemes tensor over \(k\), as follows by specifying the
algebra maps from the factors to any commutative \(k\)-algebra;
dimensions of tensor products multiply by taking bases.
On the other hand \(\sum_v d_v=fn\) by
Lemma~\ref{lem:degree-one-split-lie-polynomial}. Equality
in every factor bound follows, so every stated surjection
is an isomorphism.

The generic \(v\)-factor in \(G[P]\) consists of the
classes \(p^{-a}\Z_p^n/\Z_p^n\) in the finite frame
coordinates of Lemma~\ref{lem:actual-subspace-kernel-labels}.
It is reduced by invertibility of the derivative \(P I\)
and Lemma~\ref{lem:finite-fiber-local-algebras}, hence its
coordinate dimension equals its \(P^n\) points.
Finite freeness in Lemma~\ref{lem:valuative-torsion-idempotents}
gives the same special dimension. For \(d_v=n\) the
surjection from its algebra to \(B_v\) has equal source
and target dimensions and is an isomorphism. For \(d_v=0\)
the displayed algebra is \(k\) at the identity.
\end{proof}

\begin{corollary}\label{cor:total-twisted-hecke-killed}
Under the hypotheses of Theorem~\ref{thm:actual-quadratic-half-kernels},
assume also that the targets \(G_A\) are smooth groups of
finite type. Then every full \(\ker f_A\) is a closed
subgroup of \(G_k[p^n]\), with the same \(p^n\) for every
subspace \(A\).
\end{corollary}
\begin{proof}
The generic matrix computation in
Theorem~\ref{thm:actual-quadratic-half-kernels} gives that
\(p^r\) kills the entire generic kernel of \(\psi_A\) when
\(r=\dim A\ge1\): at \(w\) a single \(p\) suffices and
at the conjugate factors \(p^r\) suffices. For \(r=0\)
the generic kernel is already the identity. Finite
flatness and the closed identity of \(G\) give this same
annihilation on the integral kernel by injection into
its generic coordinate algebra, as in that theorem.
For any special point over any \(k\)-algebra killed by
\(f_A\), its image under \((\psi_A)_k\) belongs to the
Frobenius kernel of exponent \(p^{n-r}\) on the target.
It is killed by \(p^{n-r}\) by
Lemma~\ref{lem:smooth-group-frobenius-coordinates};
for \(r=n\) this kernel is the identity.
Consequently \(p^{n-r}\) times the original point is in
the kernel of \((\psi_A)_k\), and multiplication by
\(p^r\) kills it. This gives the assertion as morphisms
of schemes, by applying the functorial identity on affine
charts of the kernel itself. Closedness follows from
separatedness of the targets and the source.
\end{proof}


\section{Flatness of Frobenius and descent through an isogeny}

\begin{proposition}\label{prop:smooth-frobenius-free}
Let \(k\) be a perfect field of characteristic \(p>0\), so
\(p\)-powering is an automorphism of \(k\). On a smooth
\(k\)-scheme \(X\) of pure dimension \(g\), relative
Frobenius \(F_{X,P}\) for \(P=p^a\), \(a\ge0\), is
finite locally free of rank \(P^g\) and surjective.
Finite locally free of positive rank means that on a
cover by affine target opens the source is affine with
coordinate algebra a finite free module of positive rank.
\end{proposition}
\begin{proof}
Use a smooth chart
\[
 R=(B[z_1,\ldots,z_h]/(f_1,\ldots,f_h))[1/d],
                \qquad B=k[t_1,\ldots,t_g],
\]
with \(J=\det(\partial f_i/\partial z_j)\) a unit in \(R\),
as in Proposition~\ref{prop:complex-relative-jacobian}.
Inverting several denominators is equivalent to inverting
their product. Write \(R^{(P)},B^{(P)}\) for the scalar
twists of Lemma~\ref{lem:frobenius-pulls-kernel-ideal}.
Use variables \(t'_i,z'_j\) in the twists. Their equations
are those of the original presentation with all coefficients
raised to \(P\), since \(c\otimes1=1\otimes c^P\) for a
scalar. Form
\[
    T=R^{(P)}\otimes_{B^{(P)}}B,\qquad t_i'\longmapsto t_i^P .
\]
There is a homomorphism \(T\to R\) sending \(t_i\) to
\(t_i\) and \(z'_j\) to \(z_j^P\). We prove it is an
isomorphism.

Adjoin to \(T\) variables \(v_j\) with \(v_j^P=z'_j\),
obtaining a finite free algebra \(S\) with basis
\(\prod_jv_j^{e_j}\), \(0\le e_j<P\), by division by
the successive monic equations. The basis contains one
(including the case \(P=1\)). The elements \(f_i(t,v)\)
have \(P\)-th power zero by the equations in \(T\). The
ideal \(I\) they generate is nilpotent by the pigeonhole
argument in Lemma~\ref{lem:smooth-group-frobenius-coordinates}.
Every element of \(I\) has \(P\)-th power zero by additivity
of \(P\)-powering. Both \(d(t,v)\) and \(J(t,v)\) are
units since their \(P\)-th powers are the corresponding
units in the twisted presentation.

There exists a tuple \(b\equiv v\bmod I\) with \(f(t,b)=0\).
Indeed, at an approximate solution with error in \(I^j\)
multiply the negative error column by the inverse Jacobian
matrix and add the resulting correction. Polynomial
expansion makes the new error belong to \(I^{2j}\).
The determinant remains a unit under a nilpotent correction.
Iteration terminates, and \(b_j^P=z'_j\). We use also the
following uniqueness of such lifts. For two exact solutions
with invertible Jacobian whose differences generate a
nilpotent ideal \(N\), the linear Taylor term gives their
differences in \(N^2\) after multiplication by the inverse
Jacobian. Thus \(N=N^2\), forcing \(N=0\).

In \(S\otimes_T S\) the two images of \(b_j\) have equal
\(P\)-th powers. Their differences generate a nilpotent
ideal, again by the pigeonhole argument. The uniqueness
just proved gives \(b_j\otimes1=1\otimes b_j\).
The \(T\)-linear functional selecting the coefficient of
one in the free basis of \(S\) takes one to one. Applying
it on the first factor gives \(b_j\in T\). The assignment
\(z_j\mapsto b_j\) defines \(R\to T\): the equations hold
in \(S\), which contains \(T\), and \(d(t,b)^P\) is the
corresponding unit in \(T\). The composite on \(T\) fixes
each \(z_j'\) by \(b_j^P=z_j'\), and fixes \(t\) and the
scalars. The two tuples on the \(R\)-side also solve the
original equations with the same \(P\)-th powers and unit
Jacobians, so uniqueness proves the other composite is
the identity.

Finally \(B\) is free over \(B^{(P)}\) with basis
\(\prod_i t_i^{e_i}\), \(0\le e_i<P\), by partitioning
the monomials by their exponent residues. Thus \(R\simeq T\)
is free of rank \(P^g\) over \(R^{(P)}\). These affine
charts and their twists correspond under Frobenius on
underlying opens, by the prime-ideal argument of
Lemma~\ref{lem:frobenius-pulls-kernel-ideal}, valid over
any perfect field. It gives surjectivity as well.
\end{proof}

\begin{lemma}\label{lem:finite-flat-morphism-descent}
Let \(b:X\to Y\) be a surjective finite locally free
morphism of schemes and \(h:X\to Z\) a morphism. If \(h\)
has the same composite with both projections \(X\times_Y X
\rightrightarrows X\), it factors uniquely through \(b\).
If \(b,h\) are homomorphisms of commutative group schemes
over the same base with closed identity sections, the
same conclusion holds whenever \(\ker b\subseteq\ker h\)
as closed subgroup schemes, and the factor is a homomorphism.
\end{lemma}
\begin{proof}
Finite maps are closed. In fact over an affine target
\({\rm Spec}(A)\) a closed part of the finite source has
coordinate ring \(M\) which is a finite \(A\)-module.
Its image is the common zero set of its annihilator.
To check this at a prime \(\mathfrak p\), the localization
\(M_{\mathfrak p}\) is zero exactly if an element outside
\(\mathfrak p\) kills its finite generating list, by taking
the product of individual denominators. When the
localization is nonzero, its residue algebra is nonzero:
otherwise the determinant argument of
Lemma~\ref{lem:finite-fiber-local-algebras} applied to
the local maximal ideal would kill that module by a unit.
A nonzero commutative ring admits a maximal ideal by
Definition~\ref{def:field-conventions}, giving a point of
the fiber. This proves closedness.

Two points over the same point \(y\) of \(Y\) have a
common lift to their fiber product over \(y\). Indeed
the tensor product of their residue fields over \(k(y)\)
is nonzero by vector space bases, and a prime ideal of it
gives such a lift. The hypothesis therefore gives all
points of \(b^{-1}(y)\) the same image in \(Z\). Choose
an affine neighborhood of that image, remove from \(Y\)
the closed image under \(b\) of its inverse complement
under \(h\), and take an affine neighborhood of \(y\)
in the remaining open. Shrink it so that the finite
algebra \(S\) over its ring \(A\) is free. Because the
fiber at \(y\) is nonempty, the vector for one in a module
basis of \(S\) has a coordinate that is a unit after
shrinking again: otherwise its residue at \(y\) would be
zero in the nonzero fiber algebra. There is consequently
an \(A\)-linear \(\lambda:S\to A\) with \(\lambda(1)=1\).
The ring map \(A\to S\) is injective, and
\(s\otimes1=1\otimes s\) in \(S\otimes_A S\) implies
\(s=\lambda(s)1\) on applying \(\lambda\) to the first
factor. Every coordinate function of \(h\) on the chosen
target chart thus belongs to \(A\), and its relations
hold there by injectivity. This defines the unique local
factor. Such factors agree on overlaps: the common value
on underlying points is fixed by surjectivity, and in
any common affine chart equality of pulled-back functions
implies equality by the same injection. Gluing gives the
first assertion.

For group maps the map \((x,x')\mapsto(x,x'-x)\) identifies
\(X\times_Y X\) with \(X\times\ker b\) over the original
base, with inverse \((x,z)\mapsto(x,x+z)\). The kernel
containment thus gives the required equality. Base change
of a surjective finite locally free map is again so, by
tensoring its free basis locally and by nonzero residue
fiber algebras as above. Its square on products has the
same property by composing two base changes. Uniqueness
of descent for this square verifies addition for the
factor, and the identity section and inverse satisfy the
corresponding identities by surjectivity and uniqueness.
\end{proof}

\begin{corollary}\label{cor:equal-kernels-frobenius-factor}
Let \(U,Y\) be smooth commutative groups of dimension \(g\)
over an algebraically closed field \(k\) of characteristic
\(p>0\), and \(\phi:U\to Y\) a surjective finite locally
free homomorphism. If
\(\ker\phi=\ker F_{U,p}\) as closed subgroups of \(U\),
there is a unique isomorphism \(i:U^{(p)}\to Y\) of
groups over \(k\) satisfying \(\phi=i F_{U,p}\). If
\(\phi\) is linear for an action of a commutative ring,
then \(i\) is linear for the twisted action on \(U^{(p)}\).
Any marked points carried by \(\phi\) are carried by \(i\)
after Frobenius on their source coordinates.
\end{corollary}
\begin{proof}
Proposition~\ref{prop:smooth-frobenius-free} makes the
Frobenius map surjective finite locally free. Apply
Lemma~\ref{lem:finite-flat-morphism-descent} in the two
directions to obtain factorizations of \(\phi\) by
Frobenius and of Frobenius by \(\phi\). Uniqueness
applied to the identity factorizations shows the two
resulting homomorphisms are inverse. The affine formula of
Lemma~\ref{lem:frobenius-pulls-kernel-ideal} commutes with
all scalar morphisms after twisting the target action.
Equality after composition with \(F_{U,p}\) therefore
gives the scalar compatibility by the same uniqueness.
The statement on marked points is evaluation of
\(\phi=i F_{U,p}\).
\end{proof}


\section{Consecutive Hecke neighbours along a specializing line}

\begin{lemma}\label{lem:relative-hecke-line-lattices}
Use a safe degree-one place as in
Lemma~\ref{lem:degree-one-split-lie-polynomial}, with
uniformizer \(p\). For \(A\subset B\subseteq\mathbf F_p^n\)
with \(\dim B=\dim A+1\), let \(K_A,K_B\) denote the generic
complex fiber kernels of the unscaled Hecke maps with
labels \(A,B\). Then
\[
                  K_A\subseteq K_B,\qquad pK_B\subseteq K_A.
\]
The orders on the factors of the quotient \(K_B/K_A\) at
\(w,\bar w\) are \(p,p^{n-1}\), respectively. At the other
primes above \(\bar u\) they are \(p^n\), and at all
remaining primes they are one. Its \(w\)-points are the
image of \(B/A\) in the target of the map with label \(A\).
\end{lemma}
\begin{proof}
Put \(r=\dim A\), \(L=\Z_p^n\), and for a residue subspace
\(D\) let \(L(D)\) be its full inverse image in \(L\).
At \(w\) the two kernels are \(A,B\) in \(p^{-1}L/L\)
by Lemma~\ref{lem:actual-subspace-kernel-labels} and the
inverse denominator calculation in
Theorem~\ref{thm:actual-quadratic-half-kernels}.
At \(\bar w\), identify the coordinates using conjugation.
The kernel for \(A\) is \(p^{-r}L(A^\perp)/L\), since its
inverse frame is \(p^r\bar a_A^{-1}\) and the lattice
\(a_A L\) is \(L(A^\perp)\) by
Lemma~\ref{lem:actual-subspace-kernel-labels}. For \(r=0\)
this formula gives the zero group. For \(B\) replace
\(r,A\) by \(r+1,B\). We have
\(pL(A^\perp)\subseteq pL\subseteq L(B^\perp)\) and
\(L(B^\perp)\subseteq L(A^\perp)\). These prove both
inclusions at this place. The index of \(L(B^\perp)\)
in \(L(A^\perp)\) is \(p\) by their residue dimensions,
so the quotient of the two enlarged lattices has order
\(p^n/p=p^{n-1}\). At a remaining place above \(\bar u\)
the two lattices are \(p^{-r}L,p^{-(r+1)}L\), giving
the inclusions and order \(p^n\). Elsewhere the kernels
are zero by Theorem~\ref{thm:actual-quadratic-half-kernels}.
The final assertion follows from the map with kernel \(A\)
on the \(w\)-labels.
\end{proof}

\begin{proposition}\label{prop:integral-adjacent-hecke-kernel}
In the notation of Theorem~\ref{thm:actual-quadratic-half-kernels}
suppose each \(\psi_A:G\to G_A\) is surjective finite
locally free. Suppose \(G_A\) is smooth and separated,
with the given generic complex group and scalar action,
and \(G_A[p]\) is finite flat over \(V\). For every pair
of Lemma~\ref{lem:relative-hecke-line-lattices} there is a
unique \(O_E\)-linear homomorphism
\(\phi_{A,B}:G_A\to G_B\) with
\(\phi_{A,B}\psi_A=\psi_B\). Its kernel \(K_{A,B}\)
is a finite free closed subgroup of \(G_A[p]\). In each
prime factor it has rank equal to the order given by
Lemma~\ref{lem:relative-hecke-line-lattices}.
\end{proposition}
\begin{proof}
The inclusion of the generic kernels extends to the
integral kernels: the ideal of the closed identity of
\(G_B\) pulled back along the restriction of \(\psi_B\)
to \(\ker\psi_A\) vanishes generically. The coordinate
algebra of that kernel injects into its generic algebra
by flatness, as in
Proposition~\ref{prop:flat-subkernel-generic-labels}.
The same argument applied to \([p]\) on \(\ker\psi_B\)
followed by \(\psi_A\) gives
\([p]\ker\psi_B\subseteq\ker\psi_A\). In both arguments
the generic algebras are products of fields by
Theorem~\ref{thm:actual-quadratic-half-kernels}. Now use
Lemma~\ref{lem:finite-flat-morphism-descent} on \(\psi_A\)
and \(\psi_B\). It gives the homomorphism and its uniqueness.
Scalar compatibility follows by equality after composition
with \(\psi_A\).

Pulling back \(\psi_A\) along the closed immersion
\(K_{A,B}\hookrightarrow G_A\) gives
\(\ker\psi_B\to K_{A,B}\), surjective finite locally free
by base change. Multiplication by \(p\) on \(K_{A,B}\)
vanishes after this cover by the second kernel inclusion.
It vanishes on \(K_{A,B}\) by uniqueness of descent.
Thus \(K_{A,B}\) is a closed subgroup of the finite
affine \(G_A[p]\), hence finite over \(V\). The coordinate
map of the cover is injective by the local functional
argument of Lemma~\ref{lem:finite-flat-morphism-descent}.
Torsion freeness of the coordinate algebra of \(\ker\psi_B\)
over \(V\) therefore implies torsion freeness for the
finite coordinate algebra of \(K_{A,B}\); it is free by
Proposition~\ref{prop:valuation-special-cluster-lengths}.
The idempotent factors are also finite free by the image
argument of Lemma~\ref{lem:valuative-torsion-idempotents}.
Its generic subgroup is reduced since it is closed in
the split reduced group \(G_A[p]_\C\), by the product-of-fields
calculation in Proposition~\ref{prop:flat-subkernel-generic-labels}.
Every generic point lifts through the displayed cover
over \(\C\), since a nonzero finite-dimensional complex
fiber algebra has a complex point by
Lemma~\ref{lem:finite-fiber-local-algebras}. Thus its
point group is \(K_B/K_A\) with the indicated factors.
\end{proof}

\begin{theorem}\label{thm:adjacent-all-except-conjugate}
Take the hypotheses of
Proposition~\ref{prop:integral-adjacent-hecke-kernel}.
There is a line \(W_0\subseteq\mathbf F_p^n\), independent
of \(A\), with the following property. For \(B=A+W_0\)
and \(W_0\not\subseteq A\), put \(U=(G_A)_k\).
On all factors of \(U[p]\) except possibly the factor at
\(\bar w\), the special subgroup \((K_{A,B})_k\) equals
\(\ker F_{U,p}\). On the \(\bar w\)-factor the two subgroups
both have length \(p^{n-1}\). If \(n=1\), equality also
holds on this factor.
\end{theorem}
\begin{proof}
Choose a line of generic \(w\)-labels in \(G[w]\) reducing
to the identity, as in
Theorem~\ref{thm:actual-quadratic-half-kernels}. By
Lemma~\ref{lem:relative-hecke-line-lattices}, the \(p\)
generic points on the \(w\)-factor of \(K_{A,B}\) are
the images of \(W_0\) under \(\psi_A\), since \(W_0\)
maps bijectively to \(B/A\). Integral sections and their
reductions commute with \(\psi_A\), so every one of them
reduces to zero. The special subgroup on that factor is
therefore supported at zero and has length \(p\), by
Proposition~\ref{prop:valuation-special-cluster-lengths}
and Proposition~\ref{prop:integral-adjacent-hecke-kernel}.
The \(w\)-factor of \(U[p]\) has one-dimensional tangent
at zero by Lemma~\ref{lem:valuative-torsion-idempotents}.
Its local algebra there has a unique quotient ideal of
colength \(p\) by Lemma~\ref{lem:finite-fiber-local-algebras}.
The \(w\)-factor of the Frobenius kernel has precisely
that support and length by
Theorem~\ref{thm:split-frobenius-kernel-ranks}, so the
two are equal.

At a remaining place over \(u\) both have length one
and contain the identity, by
Lemma~\ref{lem:relative-hecke-line-lattices} and
Theorem~\ref{thm:split-frobenius-kernel-ranks}. At each
remaining place over \(\bar u\), both are the full
\(p\)-torsion factor, since both are closed subgroups
of its length \(p^n\) and already have that length.
At \(\bar w\) the rank and dimension formulas in the
same results give the common length \(p^{n-1}\). This
is one when \(n=1\), proving equality there in that case.
\end{proof}


\section{Characters of finite groups and the first Frobenius kernels}

\begin{lemma}\label{lem:finite-dual-verschiebung}
Let \(H\) be a finite commutative group scheme over a perfect
field \(k\) of characteristic \(p>0\). There is a finite
commutative group \(H^\vee\) over \(k\) whose points over
a \(k\)-algebra \(S\) are the group characters
\(H_S\to\mathbb G_{m,S}\), with pointwise multiplication.
Here \(\mathbb G_m(S)=S^\times\). The natural evaluation
identifies \(H\) with its double dual. There is a natural
map \(V_H:H^{(p)}\to H\) satisfying
\[
              V_H F_{H,p}=[p]_H,\qquad
 \langle V_H t,\chi\rangle=
        \langle t,F_{H^\vee,p}\chi\rangle_{(p)} .
\]
The brackets on the right evaluate the twisted perfect
duality between \(H^{(p)}\) and \((H^\vee)^{(p)}\).
\end{lemma}
\begin{proof}
Let \(D\) be the finite-dimensional coordinate algebra
of \(H\). Addition gives a coassociative cocommutative
algebra map \(\Delta:D\to D\otimes D\), the identity gives
evaluation \(\epsilon:D\to k\), and inverse gives the
map on functions for negation. These assertions are the
group identities under pullback. On the linear dual
\(D^*={\rm Hom}_k(D,k)\) define multiplication by
\((b c)(x)=(b\otimes c)\Delta(x)\) and unit \(\epsilon\).
Use the transpose of ordinary multiplication of \(D\) as
comultiplication of \(D^*\), evaluation at \(1\in D\)
as counit and the transpose of negation as inverse.
The ring, multiplicativity and group identities on these
maps follow by transposing the corresponding identities
on \(D\), using the natural tensor dual identification
in finite dimensions. Thus they define a finite commutative
group on \({\rm Spec}(D^*)\).

A linear map \(D^*\to S\) is a tensor \(v\in S\otimes D\).
It is a unital algebra map exactly if
\(\Delta(v)=v\otimes_S v\) after extending scalars to
\(S\), and \(\epsilon(v)=1\). This follows by evaluating
on all pairs of dual basis elements and on the dual unit.
Such \(v\) is a unit, its inverse being the pullback by
negation by the group inverse identity. Sending the
coordinate on \(\mathbb G_m\) to \(v\) is therefore
exactly a character. This proves the representation
assertion. Dualizing twice restores \(D\) with all its
operations, and under the described identification the
resulting map is evaluation. A map of finite groups
dualizes by precomposition of characters, equivalently
by transposing its coordinate ring map.

There is a natural identification
\((D^*)^{(p)}=(D^{(p)})^*\) from the perfect pairing
\[
  (b\otimes c,\ x\otimes d)\longmapsto c d\, b(x)^p .
\]
Twisted dual bases pair as dual bases. Define \(V_H\) to
be the dual of \(F_{H^\vee,p}\) in this identification.
The evaluation identity and naturality in \(H\) follow
from precomposition of characters and naturality of
relative Frobenius in Lemma~\ref{lem:frobenius-pulls-kernel-ideal}.

For the remaining identity take a basis \(u_i\) of \(D\)
and dual basis \(b_i\). The transpose defining \(V_H\)
sends \(x\in D\) to
\(\sum_i b_i^p(x)(u_i\otimes1)\in D^{(p)}\). Here
\(b_i^p\) is the \(p\)-th power in the dual multiplication.
After pulling back by Frobenius this is
\(\sum_i b_i^p(x)u_i^p\). On the other hand the pullback
of \([p]\) is the iterated comultiplication into \(p\)
copies of \(D\) followed by multiplying all \(p\) factors.
The tensor coefficients in the basis \(u_i\) are invariant
under cyclic rotation by commutativity of addition on \(H\).
Nonfixed cyclic orbits of indices have length \(p\), so
their terms cancel after multiplication. The tuple
\((i,\ldots,i)\) has coefficient \(b_i^p(x)\), by the
dual multiplication formula. This gives exactly the
previous sum and proves the identity.
\end{proof}

\begin{lemma}\label{lem:flat-lift-split-p-square}
Take the hypotheses of
Lemma~\ref{lem:valuative-torsion-idempotents} for \(m=1,2\)
and fix a prime \(v\mid p\). On special fibers put
\(H_1=(H_{v,1})_k\), \(H_2=(H_{v,2})_k\). Then
\([p]:H_2\to H_1\) is surjective finite locally free of
rank \(p^n\), with kernel \(H_1\).
\end{lemma}
\begin{proof}
The idempotents modulo \(p^2\) reduce to those at the
same prime modulo \(p\) by the residue labels of
Lemma~\ref{lem:degree-one-split-lie-polynomial}.
The functor of points therefore identifies the kernel
of \([p]\) on this factor with \(H_1\). Every \(k\)-point
of \(H_1\) specializes from an integral point labelled
by one of its complex points by
Proposition~\ref{prop:valuation-special-cluster-lengths}.
Each label lifts under \([p]\) to the degree-two group
of labels in \(p^{-2}\Z_p^n/\Z_p^n\), using the finite
coordinates of Lemma~\ref{lem:actual-subspace-kernel-labels}.
Evaluating at this generic point gives an integral point
by the same valuation result. Its reduction lifts the
chosen \(k\)-point. Each fiber is thus isomorphic to
\(H_1\) by translation by such a lift, including the
scheme structure since the group map commutes with
translations.

The ranks over \(k\) of the two finite coordinate algebras
are \(d=p^n\) and \(d^2\), respectively, by
Theorem~\ref{thm:split-frobenius-kernel-ranks}. Decompose
the first algebra as \(\prod_x B_x\) by
Lemma~\ref{lem:finite-fiber-local-algebras}; the second
decomposes as a module over it into factors \(M_x\) by
the idempotents. Each \(M_x/\mathfrak m_x M_x\) has
dimension \(d\) by the fiber isomorphism.
Lifting a basis gives a surjection \(B_x^d\to M_x\),
since its cokernel equals its product by the nilpotent
ideal \(\mathfrak m_x\) and must be zero by iteration.
Summing the dimensions gives equal total dimensions
\(d^2\) for the sources and targets of these surjections.
Thus every one is an isomorphism. This gives the asserted
local freeness and surjectivity.
\end{proof}

\begin{proposition}\label{prop:paired-first-frobenius-kernels}
Let \(H_1,H_2\) be as in
Lemma~\ref{lem:flat-lift-split-p-square}, let \(H'\) be
a finite commutative group over \(k\), and let
\(\gamma:H_1\times H'\to\mathbb G_m\) be bilinear as a
morphism of group functors in its two variables. Then
\(\gamma\) is identically one on
\(\ker F_{H_1,p}\times\ker F_{H',p}\).
\end{proposition}
\begin{proof}
By Lemma~\ref{lem:finite-dual-verschiebung},
\(\gamma\) is equivalently a homomorphism
\(\lambda:H'\to H_1^\vee\); on an affine test scheme
the first variable gives the character for each value
of the second variable. Naturality gives a morphism
by evaluating at the universal point of the finite
affine scheme \(H'\). Thus a point \(y\in\ker F_{H',p}\)
maps to a character in \(\ker F_{H_1^\vee,p}\).
For any test scheme and \(x\in\ker F_{H_1,p}\), pull
back the finite locally free cover in
Lemma~\ref{lem:flat-lift-split-p-square} along \(x\).
On this cover choose \(t\in H_2\) with \(pt=x\).
Naturality of Frobenius gives
\(p F_{H_2,p}(t)=F_{H_1,p}(x)=0\), viewing \(H_1^{(p)}\)
in \(H_2^{(p)}\). Thus \(F_{H_2,p}(t)\) belongs to
\(H_1^{(p)}\), since formation of the kernel in that
lemma commutes with twisting the scalars. Call it \(z\).
Naturality of \(V\) for \(H_1\hookrightarrow H_2\) and
Lemma~\ref{lem:finite-dual-verschiebung} give
\(x=V_{H_1}z\). Its value with \(\lambda(y)\) is one by
the evaluation identity and \(F_{H_1^\vee,p}\lambda(y)=0\).
Equality to one descends across this cover by uniqueness
in Lemma~\ref{lem:finite-flat-morphism-descent}. Applying
the assertion to the universal pair of points proves the
morphism identity.
\end{proof}


\section{The annihilator at the conjugate prime}

\begin{lemma}\label{lem:one-tangent-dual-annihilator}
Let \(H\) be a finite commutative group over an algebraically
closed field \(k\) of characteristic \(p>0\), with tangent
dimension one at its identity. Put \(J=\ker F_{H,p}\) and
write \(d=\dim_k\mathcal O(H)\). Suppose \(J\) has length
\(p\). The subgroup \(J^\perp\subseteq H^\vee\) of
characters trivial on \(J\) is closed of length \(d/p\).
\end{lemma}
\begin{proof}
Write \(D=\mathcal O(H)\), \(C=\mathcal O(J)\), and
\(\pi:D\to C\) for restriction. In the dual algebras of
Lemma~\ref{lem:finite-dual-verschiebung}, the restriction
map on characters \(H^\vee\to J^\vee\) corresponds to
the inclusion \(C^*\to D^*\). Its kernel at the identity
is defined by the ideal \(D^* C^*_+\), where
\(C^*_+=\{\lambda:\lambda(1)=0\}\).
Its quotient dimension equals the dimension of the
orthogonal vector subspace of \(D\) to this ideal.
By dual multiplication, this is
\[
       D^J=\{b\in D:({\rm id}\otimes\pi)\Delta(b)=b\otimes1\}.
\]
Indeed pairing the second tensor component with all
functionals vanishing at \(1\in C\) forces that component
into \(k\,1\), as is seen by extending \(1\) to a basis.
The counit then gives \(b\) for the first component.
These are exactly the functions invariant under translation
by \(J\).

By Lemma~\ref{lem:finite-fiber-local-algebras} and translation
each local factor at \(x\in H(k)\) is \(D_x=k[t]/(t^e)\),
where \(e\) is independent of \(x\). The subgroup \(J\)
is supported at zero by Lemma~\ref{lem:frobenius-pulls-kernel-ideal};
its local algebra is \(k[t]/(t^p)\), since it has length
\(p\) and is a quotient of \(D_0\). Every \(p\)-th power
of an element \(b\in D\) is invariant. In fact the
difference of the coaction on \(b\) and \(b\otimes1\)
lies in \(D\otimes\mathfrak m_C\), by the counit. Its
\(p\)-th power vanishes using the algebra just calculated.

Choose a nonzero tangent vector of \(J\) at zero and
evaluate the coaction on that dual-number point in its
second factor. The resulting map is \(b\mapsto
b+\partial(b)\varepsilon\), with \(\partial\) a \(k\)-derivation
which kills every invariant. It preserves each local
factor, since a derivation kills an idempotent \(a\):
\((1-2a)\partial a=0\) and \((1-2a)^2=1\). On \(D_x\)
the value \(\partial(t)\) is a unit. Indeed the chosen
tangent is nonzero in \(H\), by injectivity of dual-number
points for a closed immersion, and translation by \(x\)
is an isomorphism taking it to the tangent at \(x\).
The parameter \(t\) detects any nonzero vector there by
Lemma~\ref{lem:finite-fiber-local-algebras}. Applying
\(\partial\) to \(t^e=0\) gives
\(e t^{e-1}\partial(t)=0\), so \(p\mid e\).
Written on the basis \(1,t,\ldots,t^{e-1}\), its kernel
is exactly the span of \(t^j\) with \(p\mid j\), by
the power rule and invertibility of \(\partial(t)\).
These are precisely the \(p\)-th powers in \(D_x\)
since \(k\) is perfect. Thus on the product the invariants
are all the \(p\)-th powers and have dimension \(d/p\).
This computes the length of \(J^\perp\).
\end{proof}

\begin{proposition}\label{prop:complex-prime-edge-pairing}
At the safe place of
Lemma~\ref{lem:relative-hecke-line-lattices}, omit also the
rational primes dividing \(h\) of
Proposition~\ref{prop:signed-period-riemann-form}. For a
label \(A\) denote the period lattice of its target by
\(\Lambda_A\) and its integral Riemann form by \(E_A\).
On \(p^{-1}\Lambda_A/\Lambda_A\) the formula
\[
                 b_A(v,t)=\exp(2\pi i\,p E_A(v,t))
\]
defines a perfect alternating pairing with values in
the \(p\)-th roots of unity. It pairs the factor at \(w\)
perfectly with the one at \(\bar w\). For
\(A\subset B\), \(\dim B=\dim A+1\), it is trivial on
pairs in the generic kernel \(K_{A,B}(\C)\) of
Proposition~\ref{prop:integral-adjacent-hecke-kernel}.
\end{proposition}
\begin{proof}
Changing a lift by a lattice element changes \(pE_A(v,t)\)
by an integer. Its product by \(p\) is also integral,
and alternation and bilinearity follow from
Proposition~\ref{prop:signed-period-riemann-form}.
In a lattice basis, the Gram determinant of \(E_A\)
has absolute value \(h^{2n}\) by the dual quotient
in that proposition and the determinant index of
Lemma~\ref{lem:class-covol}. Thus its reduction modulo
\(p\) is nonsingular. Writing lifts as \(v=v_0/p\),
\(t=t_0/p\), gives \(b_A(v,t)=\exp(2\pi i E_A(v_0,t_0)/p)\).
The exponential identifies the additive group of
\(\mathbf F_p\) with the \(p\)-th roots of unity, so
this pairing identifies the full space with its character
dual. The trace formula defining \(E_A\) gives
\(E_A(av,t)=E_A(v,\bar a t)\) for \(a\in O_E\).
Passing to the idempotents modulo \(p\) in
Lemma~\ref{lem:degree-one-split-lie-polynomial} shows the
\(w\)-factor pairs trivially with all except \(\bar w\);
global perfection and their equal dimensions imply
perfection on the pair.

Let the initial signed form be \(E_0\). The normalization
of the signed frame in
Proposition~\ref{prop:complex-actual-hecke-legs} gives
\(E_A=p^{\dim A}E_0\). Indeed the two conjugate inverse
frame matrices in
Theorem~\ref{thm:actual-quadratic-half-kernels} have
scalar Gram product of valuation \(\dim A\) at \(p\),
and unit Gram product away from \(p\). Multiplying the
initial form by the positive rational \(p^{\dim A}\)
therefore gives a unit Gram multiplier for the new frame
at every finite prime, preserving the line and signatures.
The sign and the scalar are unique by the frame argument
of Lemma~\ref{lem:marked-similitude-root-bounds}.
Consequently \(E_B=pE_A\).
The complex intermediate isogeny is the identity on the
common covering vector space modulo the enlarged lattice
\(\Lambda_B\), by the construction in
Proposition~\ref{prop:complex-actual-hecke-legs} and the
generic inclusions of Lemma~\ref{lem:relative-hecke-line-lattices}.
Represent its kernel points by elements of \(\Lambda_B\).
Their pairing exponent \(pE_A(v,t)=E_B(v,t)\) is an
integer by integrality on this lattice.
\end{proof}

\begin{theorem}\label{thm:polarized-edge-first-frobenius}
In Theorem~\ref{thm:adjacent-all-except-conjugate}, suppose
also that \(G_A[p^2]\) is finite flat for every \(A\).
At each \(A\) assume a bilinear morphism over \(V\)
\[
      (G_A[p])_w\times(G_A[p])_{\bar w}\longrightarrow
                                                        \mathbb G_m
\]
whose generic fiber is \(b_A\) on those factors and whose
special fiber induces an isomorphism from the \(\bar w\)-factor
onto the character dual of the \(w\)-factor.
Then for every \(B=A+W_0\) with \(W_0\not\subseteq A\)
there is a unique \(O_E\)-linear special group isomorphism
\[
        i_{A,B}:(G_A)_k^{(p)}\ \longrightarrow\ (G_B)_k,
       \qquad (\psi_B)_k=i_{A,B}F_{(G_A)_k,p}(\psi_A)_k .
\]
In particular if sections of \(G_A,G_B\) are the images
of common sections of \(G\), \(i_{A,B}\) carries the
Frobenius of the former reductions to the latter reductions.
\end{theorem}
\begin{proof}
Write \(H,H'\) for the two special torsion factors at
\(w,\bar w\) of \(U=(G_A)_k\) and \(J=\ker F_{H,p}\).
This kernel is the intersection of \(\ker F_{U,p}\) with
\(H\), since the twisted inclusion is a closed immersion
and relative Frobenius commutes with the inclusion. Its
length is \(p\) by Theorem~\ref{thm:split-frobenius-kernel-ranks}.
It is already the \(w\)-factor of the special \(K_{A,B}\)
by Theorem~\ref{thm:adjacent-all-except-conjugate}.
Under the assumed special duality the annihilator of
\(J\) in \(H'\) has length \(p^{n-1}\) by
Lemma~\ref{lem:one-tangent-dual-annihilator}. The special
pairing is trivial on the two factors of \(K_{A,B}\).
Indeed their product over \(V\) is finite free by
Proposition~\ref{prop:integral-adjacent-hecke-kernel},
so its coordinate algebra injects into its generic
algebra. On that algebra the pairing equals one by
Proposition~\ref{prop:complex-prime-edge-pairing}, the
generic factors being split reduced. The \(\bar w\)-factor
of the special kernel is therefore contained in the
annihilator. So is \(\ker F_{H',p}\), by
Proposition~\ref{prop:paired-first-frobenius-kernels}.
Both have length \(p^{n-1}\) by
Theorem~\ref{thm:adjacent-all-except-conjugate}; the
surjections from the coordinate algebra of the annihilator
are isomorphisms by dimension. Combining the factors
proves \(\ker(\phi_{A,B})_k=\ker F_{U,p}\).

Thus \(F_{U,p}(\psi_A)_k\) and \((\psi_B)_k\) have
the same kernel, by
Proposition~\ref{prop:integral-adjacent-hecke-kernel}.
Both are surjective finite locally free. Indeed
Proposition~\ref{prop:smooth-frobenius-free} gives such
a presentation for the Frobenius. Its map on underlying
open sets is a bijection by
Lemma~\ref{lem:frobenius-pulls-kernel-ideal}. Near each
point of its target we may consequently shrink a free
affine Frobenius chart to a principal open whose inverse
image lies inside a free affine chart for \((\psi_A)_k\).
This inverse image is affine since Frobenius is finite.
Restricting the two free bases to these opens and multiplying
their elements gives a free basis for the composite.

Apply Lemma~\ref{lem:finite-flat-morphism-descent} in
both directions to the two maps with equal kernels.
Their factorizations are inverse by uniqueness, as in
Corollary~\ref{cor:equal-kernels-frobenius-factor}.
Their scalar compatibility and statement on sections
follow by that same uniqueness and by evaluation.
\end{proof}


\section{Constructing the torsion pairing from a theta line bundle}

\begin{definition}\label{def:normalized-line-class}
For a commutative group scheme \(G\) over a base and a
line bundle \(L\) on \(G\), write \({\rm Pic}(T)\) for
the group of isomorphism classes of line bundles on \(T\),
with tensor product denoted by addition. A line bundle
is a locally free sheaf of rank one, with inverse its
dual. For \(T\)-points \(x,y\) of \(G\) put
\[
   l(x)=[x^*L]-[0^*L],\qquad
   d(x,y)=l(x+y)-l(x)-l(y).
\]
Say that \(L\) satisfies the parallelogram law if
\(l(x+y)+l(x-y)=2l(x)+2l(y)\) for every \(T,x,y\).
Equivalently it suffices to test the two universal
points on \(G\times G\) over the base, since every other
pair gives a pullback of that identity. For a \(T\)-point
\(a\), write \(\pi:G_T\to T\) for projection,
\(t_a\) for translation by \(a\), and define the line bundle
\[
   D_a=t_a^*L_T\otimes L_T^{-1}\otimes
                     \pi^*((a^*L)^{-1}\otimes 0^*L).
\]
Its pullback by a section \(z\) has class \(d(a,z)\).
Its fiber at the identity is canonically trivial by
contracting each line with its dual.
\end{definition}

\begin{proposition}\label{prop:actual-theta-parallelogram}
On every complex torus \(G\) of
Proposition~\ref{prop:signed-period-riemann-form}, let
\(L\) be the pullback of the hyperplane line bundle by
the canonical theta embedding of
Theorem~\ref{thm:canonical-theta-basis-on-ball}. It satisfies
the parallelogram law of Definition~\ref{def:normalized-line-class}
as an algebraic line bundle. Its coordinate sections are
a basis of \(H^0(G,L)\), the space of global algebraic
sections.
\end{proposition}
\begin{proof}
Use the period lattice and \(H,E_0,\eta_\beta\) of
Definition~\ref{def:signed-canonical-shift}. The line
on the vector space descends to the torus by the factors
\[
 j_\lambda(z)=\eta_\beta(\lambda)
       \exp\big(\pi H(\lambda,z)+\tfrac{\pi}{2}H(\lambda,\lambda)\big),
 \qquad (z,u)\longmapsto(z+\lambda,j_\lambda(z)u).
\]
They satisfy the translation rule by
Lemma~\ref{lem:signed-canonical-cocycle}. The canonical
theta sections are precisely the holomorphic sections of
this line, by their defining translation law. They have
no common zero by Theorem~\ref{thm:canonical-theta-basis-on-ball}.
On each chart where one is nonzero their ratios give the
embedding coordinates, so the line is holomorphically
the pullback of the hyperplane line. Its algebraic
coordinate sections identify in this way with the theta
basis. Any algebraic section is holomorphic and hence
is a complex linear combination of this basis by that
theorem. Since all its members are algebraic, the basis
assertion follows.

On the universal cover of \(G\times G\) one has
\[
 j_{\lambda+\mu}(z+w)j_{\lambda-\mu}(z-w)
                =j_\lambda(z)^2 j_\mu(w)^2 .
\]
The \(H\)-terms agree by expansion and the two
\(\eta_\beta\)-products agree by expanding the quadratic
expression \(E_0(\beta v,v)\). Thus there is a holomorphic
isomorphism between the product of the pullbacks of \(L\)
under sum and difference and the product of the squares
of its pullbacks under the two projections. The fixed
fiber factors in \(l\) are constant complex lines and
can be trivialized.

This identity of line bundles is algebraic. Indeed both
sides, \(M,N\), are generated by finite lists of algebraic
sections, obtained by pullback and multiplication of the
coordinates, since the group operations are algebraic by
Corollary~\ref{cor:complex-level-algebraic-morphisms}.
Given the holomorphic isomorphism, regard the two lists as
one list of holomorphic sections of \(M\), using the
inverse isomorphism on the list of \(N\). It has no common
zero and defines a holomorphic map from the compact
projectively embedded \(G\times G\) to a projective space.
This map is algebraic by
Theorem~\ref{thm:holomorphic-map-algebraic-locally}.
On an open where the chosen original sections of both
bundles are nonzero, the isomorphism is given, in these
two frames, by an invertible ratio of coordinates of this
map. The ratio and its inverse are regular there.
The opens cover, proving algebraicity and hence the
universal Picard identity.
\end{proof}

\begin{lemma}\label{lem:pullback-translation-line-trivial}
If \(L\) satisfies the parallelogram law, then
\(d(ma,z)=m d(a,z)=d(a,mz)\) in \({\rm Pic}(T)\) for
all \(m\ge0\) and \(T\)-points \(a,z\). In particular,
if \(ma=0\), the pullback of \(D_a\) by \([m]\) on
\(G_T\) is a trivial line bundle.
\end{lemma}
\begin{proof}
Expanding by Definition~\ref{def:normalized-line-class}
and applying the parallelogram identity at \(x+z,y\) and
at \(x,y\) gives
\[
               d(x+y,z)+d(x-y,z)=2d(x,z).
\]
Taking \(x=ja,y=a\) gives a two-step recurrence for
\(d(ja,z)\). Its values at zero and one are \(0,d(a,z)\),
so induction gives the identity in the first variable.
The formula is symmetric and hence gives the other
variable as well. Apply it on \(G_T\) itself with the
universal variable point \(z\) and the pulled-back
section \(a\). The pullback has class \(d(a,mz)=d(ma,z)=0\),
which by definition of \({\rm Pic}\) gives a trivialization.
\end{proof}

\begin{theorem}\label{thm:integral-torsion-pairing-construction}
Let \(V\) be a valuation ring of \(\C\) with algebraically
closed residue field of characteristic \(p\). Let \(G\)
be a flat projective commutative group over \(V\) with
generic fiber one of the complex tori in
Proposition~\ref{prop:actual-theta-parallelogram}. Suppose
\(Q=G[p]\) is finite flat. Suppose a line bundle \(L\)
over \(G\) satisfies the parallelogram law and on the
generic fiber is the canonical theta line of that proposition.
There is a unique morphism \(b:Q\times Q\to\mathbb G_m\)
extending
\[
               b(v,t)=\exp(2\pi i\,p E_0(v,t))
\]
on the generic torus torsion, in the lift convention of
Proposition~\ref{prop:complex-prime-edge-pairing}.
It is alternating and bilinear over \(V\).
\end{theorem}
\begin{proof}
First \(\Gamma(G,\mathcal O_G)=V\). A regular function
restricts to a holomorphic function on the compact
connected complex torus, hence is constant there by
the maximum principle (apply the one-variable principle
of Lemma~\ref{lem:analytic-tools} to complex lines in
a chart of a maximum). Its value on the identity section
belongs to \(V\), and subtracting it gives zero on every
affine chart by flatness over the domain \(V\). Indeed
a flat module over \(V\) injects into its localization
at the fraction field by the torsion-free argument of
Proposition~\ref{prop:flat-subkernel-generic-labels}.
For each flat \(V\)-algebra \(S\) the same equality gives
\(\Gamma(G_S,\mathcal O)=S\). Take a finite affine cover
of the projective scheme \(G\). Intersections of its
affines are affine since they are obtained by pulling
back its closed diagonal to the affine product over
\(V\). The ring of global functions is the kernel of
the difference of restrictions from the finite product
of chart rings to the finite product on intersections.
After tensoring with the flat algebra \(S\) this remains
the kernel, giving the claim. In particular it holds
for finite products of \(Q\) over \(V\), whose coordinate
algebras are finite free by
Proposition~\ref{prop:valuation-special-cluster-lengths}'s
torsion-free module argument.

Let \(a\) be the universal point on \(Q\). On \(G_Q\),
choose a nowhere-zero section \(s\) of \([p]^*D_a\) using
Lemma~\ref{lem:pullback-translation-line-trivial}. Over
\(Q\times Q\) use the two universal torsion points \(a,t\).
Since \([p]t=0\), translation by \(t\) preserves the
map \([p]\), so identifies the translated pullback of
\(D_a\) canonically with the same pullback. In this
identification the quotient \(t_t^*s/s\) is an invertible
regular function on \(G_{Q\times Q}\), hence an invertible
function on \(Q\times Q\) by the preceding paragraph.
It defines \(b\). Any change of the chosen section
multiplies it by a unit from \(Q\), which cancels in the
quotient. This construction commutes with pulling back
the section and its scalar identity to any base scheme.

On the complex fiber let \(v\) represent \(a\) modulo
the period lattice and let \(w\) represent the point
\(t\). In the factors of
Proposition~\ref{prop:actual-theta-parallelogram} the
translation difference \(D_a\) has lattice multiplier
\(\exp(\pi H(\lambda,v))\), by dividing
\(j_\lambda(z+v)\) by \(j_\lambda(z)\). The constant
fiber normalization does not affect this multiplier.
After pullback by \([p]\), a nowhere-zero holomorphic
section is \(s(z)=\exp(\pi H(pv,z))\). In fact
\(H(pv,\lambda)-H(p\lambda,v)=2iE_0(pv,\lambda)\)
and the latter pairing is integral before multiplication
by \(2i\). Any other nowhere-zero section differs by
a constant by compactness. In translating \(z\) by
\(w\) we must identify the fibers of \(D_a\) at \(pz+pw\)
and \(pz\), using the inverse of the multiplier at the
lattice vector \(pw\). The quotient is therefore
\[
    \frac{s(z+w)}{s(z)}\exp(-\pi H(pw,v))
               =\exp(2\pi i\,p E_0(v,w)).
\]
The coordinate algebra of every finite product of \(Q\)
is flat over \(V\), hence injects into its generic
algebra. The generic torsion is split and reduced
since \([p]\) has invertible derivative there, by
Lemma~\ref{lem:finite-fiber-local-algebras}. Its point
formula proves uniqueness, and verifies bilinearity and
alternation on \(Q^3\) and \(Q\) by that same injection.
\end{proof}


\section{Detecting the radical in the complete theta coordinates}

\begin{lemma}\label{lem:projective-special-section-base-change}
Let \(U\) be a connected reduced projective scheme over an
algebraically closed field \(k\), with a \(k\)-point.
Then \(\Gamma(U,\mathcal O_U)=k\). For a line bundle \(M\)
on \(U\) and any \(k\)-algebra \(S\), the natural map
\[
              H^0(U,M)\otimes_k S\longrightarrow
                                  H^0(U_S,M_S)
\]
is an isomorphism. In particular all global units of \(U_S\)
come from units of \(S\).
\end{lemma}
\begin{proof}
We recall the closed-image argument for the constants.
The affine point and power-containment argument of
Lemma~\ref{lem:complex-projective-finite-fibers} holds
over every algebraically closed field. Indeed the only
field properties used for finite polynomial normalization
in Lemma~\ref{lem:complex-polynomial-normalization} are
infinitude, to choose a nonzero leading value, and the
field operations. For the residue at a maximal ideal one
uses algebraic closedness. Such a field is infinite by
the polynomial argument of
Proposition~\ref{prop:valuation-special-cluster-lengths}.
Thus in each affine chart vanishing at all \(k\)-points
means membership in the radical of the ideal.

For a global function, consider its graph in the product
of \(\mathbf P^1\) with the projective space containing
\(U\), viewing its values in the affine chart of
\(\mathbf P^1\). The graph is closed since the projective
line has closed diagonal, defined by the homogeneous
cross product. Its image on \(k\)-points is closed in
\(\mathbf P^1(k)\). At any \(k\)-point outside the image,
the homogeneous fiber equations of the graph generate
an ideal containing powers of every homogeneous variable
of the source by the preceding power-containment argument.
Their homogeneous multiples consequently span every
monomial of a sufficiently large common degree. A
nonzero minor witnessing this spanning at the chosen
point remains invertible on an affine neighborhood
of it, so the graph has empty fibers there, as in
Lemma~\ref{lem:complex-projective-finite-fibers}.
The closed image avoids infinity and is therefore finite,
by the univariate root bound on an affine chart. The
product of the corresponding linear polynomials in the
global function vanishes in every reduced affine chart
by the radical assertion above. The closed inverse
images of its finitely many distinct values are disjoint
and cover \(U\), so each is open as well. Connectedness
gives just one value, and reducedness gives equality
to that constant as a regular function. Evaluation at
the assumed point identifies the constants with \(k\).

Choose finitely many affine trivializing charts for
\(M\), refining the projective affine cover. They can be
chosen from a finite collection since a covering by
principal opens of an affine scheme has a finite subcover
by the unit ideal argument in
Proposition~\ref{prop:finite-presentation-spread}.
Their intersections are affine by separatedness, and
the transition functions describe global sections as
the kernel of the differences of the restrictions there.
Tensoring with \(S\) preserves the kernel since vector
spaces over a field have bases and injections split
on a complementary basis. Sections on each trivializing
chart and its intersections tensor to the sections there
after base change. This proves the displayed isomorphism;
applying it to \(M=\mathcal O_U\) gives the last assertion.
\end{proof}

\begin{lemma}\label{lem:finite-group-injection-point-tangent}
A homomorphism \(j:H\to H'\) between finite commutative
groups over an algebraically closed field \(k\) is a closed
immersion if it is injective on \(k\)-points and on tangent
spaces at the identity. If the two coordinate algebras
also have equal dimension over \(k\), it is an isomorphism.
\end{lemma}
\begin{proof}
At each point of \(H\) the differential is injective by
translation to the identity. In the local factors of
Lemma~\ref{lem:finite-fiber-local-algebras}, the induced
cotangent maps are therefore surjective by duality of
finite-dimensional vector spaces. For a local source
point \(x\), choose images of elements of the target
maximal ideal spanning the source maximal ideal modulo
its square. These elements generate the local source
algebra over \(k\), by the nilpotence argument of
Theorem~\ref{thm:split-frobenius-kernel-ranks}. Thus the
map of local algebras at \(j(x)\) and \(x\) is surjective.
The source points have distinct target points on
underlying sets, and their local factors correspond
to the images of the respective target idempotents.
The coordinate ring map on the products is surjective,
which gives a closed immersion. Equal dimensions
make it an isomorphism.
\end{proof}

\begin{theorem}\label{thm:theta-rigidity-prime-pairing}
Use \(V,G,Q,L,b\) of
Theorem~\ref{thm:integral-torsion-pairing-construction}.
Let \(U=G_k\) be connected and reduced, \(M=L_k\).
Assume \([p]:U\to U\) is surjective finite locally free.
Suppose the complete space \(D=H^0(U,M)\) gives a
projective embedding of \(U\). Suppose there are
translations by \(k\)-points \(c_i\) killed by a positive
integer \(h\) prime to \(p\), with isomorphisms
\(t_{c_i}^*M\simeq M\) whose matrices \(P_i\) on \(D\)
span \({\rm End}_k(D)\).
Then \(a\mapsto b(a,-)\) identifies \(U[p]\) with its
character dual over \(k\). With the action and good
split conditions of
Proposition~\ref{prop:complex-prime-edge-pairing}, the
induced pairing between \(w\) and \(\bar w\) in
Theorem~\ref{thm:polarized-edge-first-frobenius} is perfect
as a pairing of finite group schemes.
\end{theorem}
\begin{proof}
Write \(H=U[p]\) and let \(j:H\to H^\vee\) be the
homomorphism defined by the special pairing using
Lemma~\ref{lem:finite-dual-verschiebung}. Take either
\(S=k\) or \(S=k[\varepsilon]/(\varepsilon^2)\), and
any point \(a\in H(S)\) in its kernel. By the construction
of Theorem~\ref{thm:integral-torsion-pairing-construction},
there is a nowhere-zero section \(s\) of \([p]^*D_a\)
on \(U_S\) for which \(t_y^*s/s=b(a,y)\) after the
natural identification for any kernel point \(y\) over
any \(S\)-algebra. This follows by pulling back the
universal section used there. The right side is one
for all such \(y\). Thus the two pullbacks of \(s\)
to \(U_S\times_{[p],\,U_S,\,[p]}U_S\) coincide, using
the identification with \(U_S\times H_S\) by subtraction
in Lemma~\ref{lem:finite-flat-morphism-descent}.
It descends to a nowhere-zero section of \(D_a\).
In fact, on a trivializing open of \(D_a\) it is a unit
in the covering algebra satisfying the equalizer identity
for that finite locally free map. The proof of
Lemma~\ref{lem:finite-flat-morphism-descent} gives a
descended function there, with inverse obtained by the
same argument. The sections agree on overlaps by injectivity.

The lines from \({\rm Spec}(S)\) in
Definition~\ref{def:normalized-line-class} are trivial,
since any open containing the closed point of the local
spectrum is the whole spectrum. Thus there is an
isomorphism \(t_a^*M_S\simeq M_S\). Pulling back sections
and applying this isomorphism gives an invertible matrix
\(P\) on \(S\otimes_k D\) by
Lemma~\ref{lem:projective-special-section-base-change}.
Its projective action expresses translation on the
embedded \(U_S\), with the usual dual convention for
point coordinates. An automorphism of a line bundle
over the identity is multiplication by a global unit:
in each rank-one frame it is a unit and these units
agree on overlaps. Such units on \(U_S\) belong to \(S\)
by the same lemma. It follows that \(P^p\) is scalar,
each \(P_i^h\) is scalar, and \(P P_i=\lambda_i P_i P\)
for some \(\lambda_i\in S^\times\), since the underlying
translations commute. Iterating this last identity with
the two scalar powers gives \(\lambda_i^p=\lambda_i^h=1\).
Their exponents are coprime, so \(\lambda_i=1\). Hence
\(P\) commutes with the full matrix algebra over \(S\)
by the spanning hypothesis. Commuting with all diagonal
coordinate projections makes it diagonal, and commuting
with the elementary off-diagonal matrices makes all
diagonal entries equal. Thus \(P\) is scalar. The
projective embedding gives \(t_a={\rm id}\), and evaluating
at zero gives \(a=0\).

We have proved injectivity on \(k\)-points and on dual-number
tangents, the latter by the kernel assertion for the
second \(S\). Both coordinate algebras have the same
dimension since they are linear duals by
Lemma~\ref{lem:finite-dual-verschiebung}. Apply
Lemma~\ref{lem:finite-group-injection-point-tangent}.
For the split assertion the adjoint identity
\(b(ca,y)=b(a,\bar c y)\) holds over \(V\) by its complex
formula in Proposition~\ref{prop:complex-prime-edge-pairing}
and injection of the finite flat product algebra into
its generic algebra. Reduce this identity and apply it
to lifts in \(O_E\) of the idempotents modulo \(p\).
Thus the isomorphism \(j\) pairs each prime summand
only with its conjugate. The dual of a product of groups
is the product of their duals, since a character is the
product of its restrictions to the summands on every
test scheme. The corresponding individual maps are
isomorphisms since their product is. Alternation allows
either order of the two factors.
\end{proof}


\section{Theta coordinates indexed by a fixed finite group}

\begin{lemma}\label{lem:odd-pairing-section-basis}
Let \(T\) be a finite abelian group of odd order \(m\),
\(k\) an algebraically closed field whose characteristic
does not divide \(m\), and \(e:T\times T\to k^\times\)
a perfect alternating bicharacter. Perfect means that
the induced map to the group of characters of \(T\) is
bijective. Let \(D_0=D_0^\perp\) be a subgroup, where
the perpendicular uses \(e\). Then \(|D_0|^2=m\).
Put \(I=T/D_0\), \(c=-(m-1)/2\), and choose a section
\(j:I\to T\) with \(j(0)=0\), \(j(-s)=-j(s)\).
Suppose a nonzero vector space \(W\) of dimension
\(\sqrt m\) has operators \(U_t\) with \(U_0=1\) and
\[
                   U_t U_v=e(t,v)^c U_{t+v}.
\]
There is a basis \(s_b\), \(b\in I\), unique up to one
common scalar, with \(s_0\) fixed by \(D_0\) and
\(s_b=U_{j(b)}s_0\). In this basis, for \(b'=b+t+D_0\)
and \(l=t+j(b)-j(b')\),
\[
          U_t s_b=e(t,j(b))^c e(j(b'),l)^{-c}\,s_{b'}.
\]
The matrices \(U_t\) are a basis of \({\rm End}_k(W)\).
Any invertible \(J\) satisfying \(J U_t J^{-1}=U_{-t}\)
acts projectively as the permutation \(b\mapsto-b\).
\end{lemma}
\begin{proof}
Characters on a subgroup extend one element at a time,
choosing a root of the order of the cyclic extension, as
in Lemma~\ref{lem:norm-gauss}. These orders divide \(m\)
and are prime to the characteristic, so the root polynomials
have exactly their degrees in distinct roots. Hence the
dual of each subgroup has its order. Restricting the
perfect pairing to \(D_0\) gives a surjection \(T\to
D_0^\vee\) with kernel \(D_0^\perp=D_0\); here the dual
denotes ordinary characters of the finite abstract group.
This gives the order and identifies \(I\) with that dual.
No nonzero \(b\) equals \(-b\) in the odd group \(I\),
so the required section exists by choosing representatives
in opposite pairs.

The \(U_l\) for \(l\in D_0\) form an ordinary group
representation because \(e\) is trivial on these pairs.
They commute and diagonalize simultaneously: each is
killed by \(X^m-1\) with simple roots, and the eigenspaces
of one are stable under all the others. Since \(2c=1-m\),
the commutator of \(U_t,U_v\) is \(e(t,v)\). The invertible
\(U_t\) with inverse \(U_{-t}\) therefore carry the
\(D_0\)-eigenspaces amongst themselves by multiplication
of their characters by the restrictions of \(e(\cdot,t)\).
These run through every character of \(D_0\).
Since \(W\ne0\), all eigenspaces are nonzero with a common
dimension. Their number equals \(\dim W\), so each is a
line. This proves the asserted uniqueness and the basis
property. Applying the product rule twice and using
\(U_l s_0=s_0\) gives the displayed action. For nonzero
\(t\) some \(U_v\) conjugates \(U_t\) to a nontrivial
scalar multiple of itself, so its trace is zero. Multiply
a proposed relation among the \(U_t\) by \(U_{-v}\) and
take the trace. The trace of the identity is \(\sqrt m\),
nonzero in \(k\), leaving only the coefficient of \(U_v\).
All vanish, and there are \((\dim W)^2\) matrices.
Finally \(J\) preserves the invariant line and carries
\(U_{j(b)}\) to \(U_{-j(b)}=U_{j(-b)}\), proving the
permutation statement.
\end{proof}

\begin{proposition}\label{prop:frame-independent-odd-theta-basis}
Let \(T=L^\perp/L\) and \(D_0=\Lambda/L\) be the finite
group and its self-perpendicular subgroup in
Theorem~\ref{thm:actual-pel-torsion-pairing}. Fix one
section \(j\) as in Lemma~\ref{lem:odd-pairing-section-basis}.
For every framed complex torus, identify \(T\) with
\(\Lambda_x^\#/\Lambda_x\) using the finite frame.
On \(T\) write \(e_x(t,v)=\exp(2\pi i E_x(\tilde t,\tilde v))\),
where \(E_x\) is its integral Riemann form and the tildes
are lifts to its dual lattice. The theta embedding has
a basis indexed by this same \(I=T/D_0\) in every frame,
well defined up to a scalar. The projective coordinates
are holomorphic on each period family and are respected
by rational changes of frame and by the level group.
Translations have the matrices of
Lemma~\ref{lem:odd-pairing-section-basis} for \(e=e_x\);
inversion is its indicated permutation in projective
coordinates.
\end{proposition}
\begin{proof}
The local dual lattices are computed in
Proposition~\ref{prop:signed-period-riemann-form}.
Transport by the finite frame and reduction of rational
lifts as in Lemma~\ref{lem:actual-subspace-kernel-labels}
give exactly \(T\). Changing a lift changes \(E_x\) by
an integer, so \(e_x\) is well defined. It has trivial
radical by the dual-of-the-dual calculation of
Theorem~\ref{thm:signed-constant-translation-matrices},
and hence is perfect by the finite character count of
Lemma~\ref{lem:norm-gauss}. The subgroup \(D_0\) is
self-perpendicular also for \(e_x\). Indeed every subgroup
of the finite abelian \(T\) is the product of its primary
parts: the corresponding primary projections are integer
multiplications by the Chinese remainder identity on
the exponent. Pairings between different primary parts
are one since their orders are coprime. In local frame
coordinates each alternating trace pairing is scaled by a
unit \(\varepsilon_p\in\Z_p^\times\) by
Proposition~\ref{prop:signed-period-riemann-form}. A
local fractional trace is integral exactly when its
product by this unit is integral. Thus its annihilator
on each primary part is precisely that of
Theorem~\ref{thm:actual-pel-torsion-pairing}, giving
\(D_0=D_0^\perp\). Here the agreement of the rational
and local integrality tests follows by taking all
completions of the rational trace, as in
Lemma~\ref{lem:trace-finite}. The order \(m=h^{2n}\)
is odd since \(2\mid P\) in
Lemma~\ref{lem:actual-signs-and-padding} and \(h\) is
prime to \(P\).

Replace \(\beta\) in
Definition~\ref{def:signed-canonical-shift} by
\(\beta'=m\beta-(m-1)/2\). It still satisfies
\(\beta'+\bar\beta'=1\). On a lattice vector the multiplier
\(\eta_{\beta'}\) equals \(\eta_\beta^m=\eta_\beta\):
\(E_x(\lambda,\lambda)=0\) and \(\eta_\beta(\lambda)\)
is a sign. Hence the section space and its embedding
are the same. On the dual lattice take
\(U'_t=\eta_{\beta'}(\tilde t)^{-1}W_{\tilde t}\).
The proof of Lemma~\ref{lem:signed-canonical-cocycle}
applies to \(\beta'\), which preserves the lattice and
its dual. Since \(m\beta\tilde t\in\Lambda_x\), it gives
\(U'_tU'_v=e_x(t,v)^c U'_{t+v}\). The operators depend
only on the indicated classes. Their matrices are constant
over the period ball by
Theorem~\ref{thm:signed-constant-translation-matrices}:
the further scalar adjustment uses only \(E_x,\beta'\)
and fixed real lifts. The invariant line and the other
new basis vectors of Lemma~\ref{lem:odd-pairing-section-basis}
can therefore be chosen by a fixed constant change of the
old basis on that ball. Projective evaluations are then
holomorphic by Theorem~\ref{thm:canonical-theta-basis-on-ball}.
For inversion, \(S(z)\mapsto S(-z)\) acts on the section
space: the multiplier at \(-\lambda\) is the same sign
as at \(\lambda\) and the two linear arguments change
sign together. It conjugates \(U'_t\) to \(U'_{-t}\)
by the defining formula. Apply the last assertion of
Lemma~\ref{lem:odd-pairing-section-basis}.

Rational equivalences carry the trace form, the lattice,
the complex structure and the labelled torsion to those
of the other representative. Pullback therefore
intertwines the \(U'_t\) with the same label by their
defining formula. A level change fixes the \(T\)-labels
by the divisibility argument in
Corollary~\ref{cor:level-invariant-theta-coordinates}.
Such an intertwiner carries the invariant line and the
basis specified from it to their counterparts, up to
a common scalar. This proves the descent assertions.
\end{proof}

\begin{corollary}\label{cor:odd-theta-matrices-specialize}
For an odd prime \(p\nmid h\), and a valuation ring of
\(\C\) of residue characteristic \(p\), the matrices in
Proposition~\ref{prop:frame-independent-odd-theta-basis}
are integral and invertible over the valuation ring.
They specialize to a basis of the full matrix algebra
over the residue field. When the finite pairing \(e_x\)
is replaced by \(e_x^p\), all entries of the new matrices
are the \(p\)-th powers of the previous entries.
\end{corollary}
\begin{proof}
Each nonzero entry is a product of integral powers of
values of \(e_x\), hence a root of unity of order dividing
\(m=h^{2n}\). Roots of unity are units in a valuation ring
by integrality of them and their inverses, as in
Proposition~\ref{prop:valuation-special-cluster-lengths}.
The matrices permute basis vectors up to these units.
If a nonidentity root \(\xi\) of order \(s\mid m\) reduced
to one, reducing \(1+\xi+\cdots+\xi^{s-1}=0\) would give
\(s=0\) in the residue field, impossible. Thus the reduced
pairing remains perfect and the commutator trace argument
of Lemma~\ref{lem:odd-pairing-section-basis} still gives
linear independence and spanning, also directly over the
residue field. The matrix formula of that same lemma
commutes termwise with raising the nonzero entries to \(p\).
\end{proof}


\section{Equality of the marked projective coordinates after Frobenius}

\begin{lemma}\label{lem:odd-isogeny-symmetric-line}
Let \(b:X\to Y\) be a surjective finite locally free
homomorphism of connected reduced projective commutative
groups over an algebraically closed field \(k\). Suppose
\(\ker b\) is killed by an odd positive integer \(d\).
If a symmetric line bundle \(M\) on \(Y\), meaning
\([-1]^*M\simeq M\), has trivial pullback under \(b\),
then \(M\) is trivial.
\end{lemma}
\begin{proof}
Fix a nowhere-zero section \(s\) of \(b^*M\).
Translation by a kernel point \(z\) identifies its pullback
with \(b^*M\), because \(b\,t_z=b\). The ratio \(t_z^*s/s\)
is thus a scalar \(\lambda(z)\) on the test scheme by
Lemma~\ref{lem:projective-special-section-base-change},
applied to the universal point of the affine finite kernel.
Composition of translations gives a character \(\lambda\),
independent of multiplying \(s\) by a constant. The
line \([-1]^*M\) gives the inverse character by pulling
the section back under inversion, since that map conjugates
\(t_z\) to \(t_{-z}\). Isomorphic lines give the same
character by transporting the trivializing section.
Thus \(\lambda^2=1\), while \(\lambda^d=1\) since the
kernel is killed by \(d\). A Bezout identity gives
\(\lambda=1\). As in the proof of
Theorem~\ref{thm:theta-rigidity-prime-pairing}, the section
now has identical pullbacks to \(X\times_Y X\) and descends
to a nowhere-zero section of \(M\) by
Lemma~\ref{lem:finite-flat-morphism-descent}.
\end{proof}

\begin{lemma}\label{lem:projective-theta-centralizer}
Use the prime-to-characteristic data and basis of
Lemma~\ref{lem:odd-pairing-section-basis}.
The only projective linear transformation commuting
projectively with all \(U_t\) and with the inversion
permutation \(J\) is the identity. Here \(J s_b=s_{-b}\).
\end{lemma}
\begin{proof}
Take a linear representative \(P\). The scalars in
\(P U_t P^{-1}=\lambda_t U_t\) form a character of \(T\),
by multiplying the displayed equations and using the
product law of Lemma~\ref{lem:odd-pairing-section-basis}.
By perfection \(\lambda_t=e(v,t)\) for some \(v\in T\).
Thus \(U_{-v}P\) commutes linearly with all \(U_t\) and
is scalar by their spanning the full endomorphism algebra.
The basis formula, \(j(-b)=-j(b)\), and bilinearity give
\(J U_v J^{-1}=U_{-v}\), since simultaneously negating
both arguments of the bicharacter preserves it.
Projective commutation with \(J\) now gives that \(U_{2v}\)
is scalar. Its commutators and perfection imply \(2v=0\).
The odd order of \(T\) gives \(v=0\), as required.
\end{proof}

\begin{theorem}\label{thm:marked-theta-frobenius-tuples}
Assume the integral group and isogeny hypotheses of
Theorem~\ref{thm:polarized-edge-first-frobenius}, including
\(p\nmid 2h\). Its special perfect pairing condition
can be supplied by the following theta data. Let
\(L,L_A\) on \(G,G_A\) be symmetric line bundles with
the parallelogram identity and the stipulated generic
theta lines. Assume on the special fibers that
\begin{enumerate}
\item the groups are connected, reduced and projective,
\([p]\) is surjective finite locally free, and the
labelled homogeneous coordinates of
Proposition~\ref{prop:frame-independent-odd-theta-basis}
give the complete linear embeddings for these lines,
with their indicated matrices for translations and
inversion,
\item \((\psi_A)_k^*(L_A)_k\simeq (L_k)^{\otimes p^{\dim A}}\)
for every \(A\), and the integral groups are projective
and flat over \(V\).
\end{enumerate}
Mark any finite list of sections of \(G/V\) and use their
images under each \(\psi_A\). Include the labelled
\(T\)-translations as common sections with the specified
labels. Write \(z_A\) for the tuple of projective coordinates
of the reduced marked sections on \(G_A\).
Writing \({\rm Fr}\) for coordinatewise \(p\)-powering,
there is a line \(W_0\subseteq\mathbf F_p^n\) such that
\[
      {\rm Fr}^{\,n-\dim A}z_A
           ={\rm Fr}^{\,n-\dim(A+W_0)}z_{A+W_0}
                     \qquad(A\subseteq\mathbf F_p^n).
\]
In particular the sum in
Corollary~\ref{cor:pointwise-congruence-single-line}
is zero for these tuples.
\end{theorem}
\begin{proof}
By Corollary~\ref{cor:odd-theta-matrices-specialize}
the special translation matrices span the endomorphism
algebra. Their points have order dividing \(h\).
Theorems~\ref{thm:integral-torsion-pairing-construction}
and~\ref{thm:theta-rigidity-prime-pairing} provide the
integral pairing and special perfection required in
Theorem~\ref{thm:polarized-edge-first-frobenius}.
Take its line \(W_0\). When \(W_0\subseteq A\) the asserted
identity is immediate. Otherwise put \(B=A+W_0\), \(r=\dim A\),
\(U=(G_A)_k\), and take the isomorphism \(i:U^{(p)}\to(G_B)_k\)
of that theorem. Let \(M_A,M_B\) denote the two special
theta lines. Then
\[
                 i^*M_B\simeq M_A^{(p)} .
\]
Indeed pull back their ratio by the surjective finite
locally free map \(b=F_{U,p}(\psi_A)_k\) in that theorem.
For any line bundle \(N\), \(F_{U,p}^*N^{(p)}=N^{\otimes p}\)
as classes, since its transition functions are raised
to \(p\). Thus both pullbacks are \(L_k^{\otimes p^{r+1}}\)
by the assumed identities. The ratio is symmetric since
both lines are symmetric and \(i\) is a group map.
The kernel of \(b\) equals that of \((\psi_B)_k\), killed
by \(p^{r+1}\) in the proof of
Corollary~\ref{cor:total-twisted-hecke-killed}. Apply
Lemma~\ref{lem:odd-isogeny-symmetric-line}.

For any two labels of \(T\) the finite alternating
pairings on the generic fibers satisfy \(e_B=e_A^p\).
Indeed the two points on the \(B\)-fiber are the images
of those on the \(A\)-fiber, and their lifts in the common
vector space may be kept. Their Riemann form is multiplied
by \(p\) by Proposition~\ref{prop:complex-prime-edge-pairing}.
The two labels on reduction are still matched by \(i\)
after Frobenius by the section assertion in
Theorem~\ref{thm:polarized-edge-first-frobenius}.
In the chosen section bases the translation matrices for
the twisted \(A\)-fiber and the \(B\)-fiber therefore agree
by Corollary~\ref{cor:odd-theta-matrices-specialize}.
Their inversion permutations agree as well. An isomorphism
of the two lines above gives a linear isomorphism on their
complete spaces of sections. It intertwines all these
matrices projectively: the actions lift the same translated
points under \(i\), and two isomorphisms of the same lines
above the same map differ by a global unit, a scalar by
Lemma~\ref{lem:projective-special-section-base-change}.
The same applies to inversion. Consequently the intertwiner
is projectively the identity by
Lemma~\ref{lem:projective-theta-centralizer}. Evaluation
at each marked section and its matched image gives
\(z_B={\rm Fr}(z_A)\). Here the twist of a projective
coordinate list on a Frobenius image is its coordinatewise
\(p\)-th power by Lemma~\ref{lem:frobenius-pulls-kernel-ideal}.
Applying \({\rm Fr}^{\,n-r-1}\) proves the displayed identity.
The formal cancellation now follows from
Corollary~\ref{cor:pointwise-congruence-single-line}.
\end{proof}


\section{One projective coordinate space for all level components}

\begin{lemma}\label{lem:all-torsion-recovers-frame}
Let \(\mathcal T=\mathbb A_{E,f}^n/L\), and for a frame
\(x\) write \(P_x(t)\) for its torus torsion point of label
\(t\in\mathcal T\) via Lemma~\ref{lem:actual-subspace-kernel-labels}.
Write \(F_{x,b}\) for its theta embedding in the common
coordinates \(I\) of Proposition~\ref{prop:frame-independent-odd-theta-basis},
with period parameter \(b\). Suppose \(x,x'\) are two frames
with parameters \(b,b'\) and \(k\in C\), and
\[
       F_{x,b}(P_x(t))=F_{x',b'}(P_{x'}(k t))
                                       \qquad(t\in\mathcal T).
\]
Then there is \(y\in{\rm GL}_n(E)\) with \(yx=x'k\) and
\(y b=b'\), where \(y\) acts on the period structures.
\end{lemma}
\begin{proof}
The torsion points of each real torus are dense by
rational approximation in a lattice basis, as in
Proposition~\ref{prop:finite-period-evaluation-tests}.
The two images of the entire tori are compact and equal,
since they are the closures of the same set of labelled
point images and \(k\) permutes all labels. The composite
\(Y=F_{x',b'}^{-1}F_{x,b}\) is a biholomorphic map of
tori matching those labels. The label maps \(P_x,P_{x'}\)
are \(O_E\)-linear by their finite frame formula, and
\(k\) is \(O_E\)-linear. On the dense torsion subgroup
\(Y\) is therefore additive and respects every element
of \(O_E\). Continuity gives these identities on the
whole torus, using also density of the product of the
two copies of the torsion.

Let \(y\) be its real derivative at zero in the covering
space coordinates. The curve \(s\mapsto Y(sv+\Lambda_x)\)
has constant velocity \(y v\) in translated target
coordinates, by additivity and differentiation at zero.
Subtracting the curve \(s\mapsto s yv+\Lambda_{x'}\)
has derivative zero and initial value zero, so the two
agree for every real \(s\), by local vector-space charts.
The inverse map gives \(y\Lambda_x=\Lambda_{x'}\).
It follows that \(y\) preserves their common rational
span \(E^n\). Its \(O_E\)-linearity there implies
\(E\)-linearity by taking fractions, so \(y\in{\rm GL}_n(E)\).
Biholomorphy at zero gives \(yI_b=I_{b'}y\), hence
\(y b=b'\) by Proposition~\ref{prop:signed-ball-period-coordinates}.

In finite frame coordinates the action on torsion says
that \(x'^{-1}yx-k\) carries \(m^{-1}L\) into \(L\)
for every positive integer \(m\). Thus each matrix entry
lies in \(m\widehat O_E\) for every \(m\). At a prime
over \(p\), taking \(m=p^a\) and letting \(a\) increase
forces it to be zero by the local valuation. This proves
\(yx=x'k\).
\end{proof}

\begin{theorem}\label{thm:common-coarse-torsion-coordinates}
Fix an open subgroup \(C'\) of finite index in \(C\). For
each \(M\ge1\), let \(D_M\) be the permutation group on
\(\mathcal T[M]\) induced by \(C'\). For a framed base
point put
\[
    \mathbf v_M(x,b)=
       \bigotimes_{t\in\mathcal T[M]} F_{x,b}(P_x(t)),
       \qquad
    \mathbf z_M(x,b)=\prod_{\gamma\in D_M}\gamma\mathbf v_M(x,b).
\]
The first expression is the line of a simple tensor of
point representatives in the fixed ordered coordinates.
The second is the line of the product of the indicated
linear forms on the dual of that tensor space. There is
an \(M_0\) such that for every multiple \(M\) of \(M_0\)
the map \(\mathbf z_M\) is a closed holomorphic embedding
of the entire base \(\mathcal X(C',0)\). For \(\nu\ge0\)
the joint map with the \(\nu\) individual \(F_{x,b}(P_a)\)
is a closed holomorphic embedding of \(\mathcal X(C',\nu)\)
in a product of projective spaces. Its image is algebraic
over \(\C\). The tensor, permutation and product operations
above are given by homogeneous polynomials with integer
coefficients.
\end{theorem}
\begin{proof}
Use any fixed order for the finite labels and bases.
The tensor product is the projective product embedding
of Corollary~\ref{cor:compact-holomorphic-graphs-algebraic}.
A rational change preserves the evaluations with their
labels by Proposition~\ref{prop:frame-independent-odd-theta-basis}.
Replacing \(x\) by \(xk\) for \(k\in C'\) changes
\(P_x(t)\) to \(P_x(kt)\) in the same underlying torus,
by its finite-coordinate definition. It preserves the
preferred projective embedding since it fixes the smaller
finite group of theta translations and its labelled basis.
Thus the orbit product descends. It is holomorphic on
the period charts, since for each fixed torsion label
one is evaluating at the periods of a fixed real vector;
see Proposition~\ref{prop:finite-period-evaluation-tests}.
Equality of two orbit products means equality of the
two ordered lists up to an element of \(D_M\), by linear
factor uniqueness in Lemma~\ref{lem:free-projective-finite-quotient}
and injectivity of the projective product embedding.
The factor lines of a single such product are distinct.
Indeed the points \(P_x(t)\) for these labels are distinct,
and \(F_{x,b}\) embeds the torus, so no nonidentity
effective permutation fixes the ordered list.

There are finitely many component representatives \(x\)
and a compact set \(K_x\) in each period ball meeting
every orbit for that component by
Theorem~\ref{thm:complex-level-projective-equations} and
the compact representative argument of
Theorem~\ref{thm:finite-torsion-level-base-embedding}.
On each \(K_x\), finitely many torsion evaluations have
jointly injective derivative and suffice for local
injection near every point of \(K_x\), by
Proposition~\ref{prop:finite-period-evaluation-tests}.
Its proof applies in the present coordinates since the
change of section basis in
Proposition~\ref{prop:frame-independent-odd-theta-basis}
is constant on this ball. For multiples of a common
integer the tensors therefore immerse each base chart.
For the orbit product, reduce its differentiated expression
modulo the linear factor of the identity as in
Lemma~\ref{lem:free-projective-finite-quotient}.
Distinctness of those factors shows that a vanishing
projective derivative of the product forces that of
the ordered tensor to vanish. Thus the orbit map is
immersive at every representative.

To prove separation, suppose injectivity fails for
arbitrarily large factorial levels \(M=j!\). Choose pairs
of distinct base classes with equal product. After taking
subsequences, their representatives have fixed frames
\(x,x'\) and parameters \(b_j\in K_x,\ b'_j\in K_{x'}\)
converging to \(b,b'\). Take \(k_j\in C'\) carrying one
ordered list to the other in the direction
\[
 F_{x,b_j}(P_x(t))=F_{x',b'_j}(P_{x'}(k_j t))
                                      \quad(t\in\mathcal T[j!]).
\]
After a further subsequence \(k_j\to k\in C'\).
For clarity, extract successively subsequences constant
modulo each integer factorial in the integral matrix
coordinates. The compatible limits of the matrices and
their inverses define an invertible integral matrix by
the finite local residue description of
Lemma~\ref{lem:local-residue-indicators}. The limit lies
in \(C'\) since \(C\) is compact closed and an open
finite-index subgroup is closed by its open cosets.
The labels of any fixed order are eventually fixed
under \(k^{-1}k_j\). Passing to the limit hence gives
the hypothesis of Lemma~\ref{lem:all-torsion-recovers-frame}.
Take \(y\) supplied by that lemma.

Pulling the second parameter back by \(y^{-1}\) now
puts it in the first ball, tending to \(b\). In this
frame the displayed equality of the lists replaces
\(k_j\) by \(k^{-1}k_j\), since \(yx=x'k\).
Choose fixed torsion labels giving a local injection
near \(b\) as above. Their orders divide \(j!\) and
\(k^{-1}k_j\) fixes them for all sufficiently large
indices. The two parameters must then be identical
in that neighborhood: \(b_j=y^{-1}b'_j\).
Together with \(yx=x'k\), this says their original
classes coincide at level \(C'\), a contradiction.
Hence some factorial level large enough for immersion
also separates all points. An equality of ordered lists
up to \(C'\) at any multiple restricts to one at the
chosen level, proving the assertion for all its multiples.

The space is compact and a holomorphic injective
immersion there is a closed embedding by the local
inverse argument of
Theorem~\ref{thm:signed-fiber-projective-embedding}.
For the marked total space the individual evaluations
descend by Proposition~\ref{prop:frame-independent-odd-theta-basis}.
After the bases are matched, the fiber embeddings
separate the marked points and detect their vertical
tangents. The same compactness and immersion proof as
in Theorem~\ref{thm:complex-level-projective-equations}
now applies. Algebraicity of the image follows from
Corollary~\ref{cor:compact-holomorphic-graphs-algebraic}
and the projective product embedding. Finally each
coordinate operation takes products of entries,
permutes factors, or takes coefficients of a product
of linear forms. These are the stated integral polynomials.
\end{proof}


\section{Complex periods under conjugation of the coefficients}

\begin{lemma}\label{lem:projective-group-period-uniformization}
Let \(Y\) be a smooth irreducible projective commutative group
over \(\C\), of dimension \(g\). Put \(W=T_0Y(\C)\).
There is a unique holomorphic additive map
\(e:W\to Y(\C)\) with derivative the identity at zero.
It induces an isomorphism \(W/\Lambda\simeq Y(\C)\) of
complex Lie groups, where \(\Lambda\) has a basis over \(\Z\)
which is a real basis of \(W\). Every algebraic homomorphism
between two such groups lifts along these maps to its
complex derivative at the identity.
\end{lemma}
\begin{proof}
The invertible Jacobian minor in a smooth chart gives
complex manifold coordinates by the local inverse argument
of Theorem~\ref{thm:holomorphic-map-algebraic-locally}.
Projective space is compact as the quotient of the unit
sphere; hence its closed algebraic subset \(Y(\C)\) is
compact. Its manifold components are open and closed and
there are finitely many by compactness. Each component
is a closed embedded complex submanifold of projective
space and is algebraic by
Corollary~\ref{cor:compact-holomorphic-graphs-algebraic}.
Irreducibility forces a single component.

Take a holomorphic coordinate \(z\) at zero with identity
derivative. Write \(F(z,w)\) for the local sum, and set
\(M(z)=\partial_2F(z,0)\). It is invertible near zero.
The columns \(X_j(z)=M(z)e_j\) are carried to themselves
by local translations, by differentiating associativity in
the last variable. They commute: the mixed derivatives in
\(s,t\) of \(F(F(z,s e_i),t e_j)=F(F(z,t e_j),s e_i)\)
give \(D X_j(z)X_i(z)=D X_i(z)X_j(z)\).
Thus the rows of \(\omega=M(z)^{-1}dz\) are closed.
Indeed their values on the \(X_j\) are constant, and
differentiating these values and subtracting the two
orders gives \(d\omega(X_i,X_j)=0\). On a ball put
\(h(z)=\int_0^1\omega(tz)z\,dt\). Closedness and
differentiation under this integral give \(dh=\omega\)
by integrating \(d(t\omega(tz))/dt\).
Translation invariance gives \(h(F(z,w))=h(z)+h(w)\)
near \((0,0)\), by differentiating in \(w\) and setting
\(w=0\). The local inverse of \(h\) therefore supplies a
local additive map \(e_0\) of derivative one. Shrink to
a convex ball around zero on which the local sum rule
holds whenever both summands and their sum lie in that ball.

For \(v\in W\) put \(e(v)=m e_0(v/m)\) for a sufficiently
large positive integer \(m\). Subdividing the segment
from \(0\) to \(v/m\) shows that this does not change
on replacing \(m\) by a multiple; a common multiple
proves independence. The local sum rule at
\(v/m,w/m,(v+w)/m\) gives global additivity.
Translations of \(e_0\) give holomorphy everywhere.
Its image is open, hence closed by the openness of its
cosets, and connectedness gives surjectivity. Its kernel
is discrete and the quotient has the local charts of
\(e_0\). Any two additive maps with the same derivative
have a difference with derivative zero at every point
by translation. Their difference is constant, and is
zero at the origin. The same argument with a homomorphism
composed with \(e\) proves functoriality.

The discrete kernel is closed. Its intersection with
a bounded set is finite, since two members of any
convergent subsequence would have arbitrarily small
nonzero difference. Its real span is \(W\), as the compact
quotient maps onto the real vector space modulo that span.
Choose a real basis from the kernel; its integral span
has finite index in the kernel by translating into the
bounded basis parallelepiped. Thus the kernel lies between
this basis lattice and its multiple by \(1/d\), where
\(d\) is the finite index (multiplication by \(d\) kills
the quotient, by partition into translates of a cyclic
subgroup). A subgroup of \(\Z^r\) is free: induct on
\(r\) by projecting to the last coordinate, adjoining
a lift of the positive generator of the image when it
is nonzero. Applied in the coordinates of the enlarged
basis lattice this gives a basis of rank \(2g\), since
the kernel spans \(W\).
\end{proof}

\begin{proposition}\label{prop:intrinsic-pairings-all-orders}
Suppose \(Y\) is as in
Lemma~\ref{lem:projective-group-period-uniformization} and
an algebraic line bundle \(B\) satisfies the parallelogram
law of Definition~\ref{def:normalized-line-class}.
For \(m\ge1\) and \(a,t\in Y[m](\C)\) the scalar
\[
 b_{m,B}(a,t)=t_t^*s/s,\qquad
        0\ne s\text{ a nowhere-zero section of }[m]^*D_a ,
\]
is well defined, using the equality \([m]t_t=[m]\) to
identify fibers. If \(B\) is symmetric, let \(q_B(a)\)
for \(2a=0\) be the scalar on \(B_a\) of the isomorphism
\([-1]^*B\to B\) normalized to be one on \(B_0\).
These scalars commute with conjugation by every field
automorphism of \(\C\), acting on all the algebraic data.
For the actual period tori and lines of
Proposition~\ref{prop:actual-theta-parallelogram}, with
\(a=v+\Lambda,\ t=w+\Lambda\), one has
\[
 b_{m,B}(a,t)=\exp(2\pi i\,m E_0(v,w)),\qquad
 q_B(a)= b_{2,B}(\beta a,a)\quad(2a=0).
\]
\end{proposition}
\begin{proof}
Lemma~\ref{lem:pullback-translation-line-trivial} gives
the section. Regular functions on \(Y\) are constant by
Lemma~\ref{lem:projective-special-section-base-change}, so
its ratio with a translated section is scalar, independent
of the section. Two isomorphisms of the two symmetric
lines differ by a global unit, giving the uniqueness of
the normalization. Double inversion then gives the identity
on the whole line, in particular \(q_B(a)^2=1\).
Transport of local defining equations, algebraic sections,
and their pullbacks by a field automorphism preserves
these ratios and the origin normalization. This proves
the asserted naturality.

For the first evaluation use the multiplier of \(D_a\)
from Theorem~\ref{thm:integral-torsion-pairing-construction}.
On the vector cover a trivializing section after \([m]\)
is \(\exp(\pi H(mv,z))\), since \(mv\in\Lambda\) and
\(E_0(mv,\lambda)\in\Z\) for \(\lambda\in\Lambda\).
Dividing by the multiplier for \(mw\) on translation by
\(w\) gives the displayed exponential, exactly as in that
theorem. A holomorphic nowhere-zero section differs from
the algebraic one by a constant by compactness.

The canonical factors satisfy \(j_{-\lambda}(-z)=j_\lambda(z)\);
the normalized inversion lift in the vector cover is
\((z,u)\mapsto(-z,u)\). It is algebraic by the joint
section argument of
Proposition~\ref{prop:actual-theta-parallelogram}, applied
to \(B\) and \([-1]^*B\). At \(v=\lambda/2\) the
fiber scalar is \(j_\lambda(-\lambda/2)=\eta_\beta(\lambda)\)
(the opposite direction gives the same sign). By
Definition~\ref{def:signed-canonical-shift} this is
\(\exp(\pi i E_0(\beta\lambda,\lambda))\), which equals
\(b_{2,B}(\beta a,a)\) by the first formula.
\end{proof}

\section{A positive curvature form with constant coefficients}

\begin{definition}\label{def:line-metric-conventions}
On a complex vector space with \(z_j=x_j+i y_j\), let
\(\partial f=\sum_j(f_{x_j}-i f_{y_j})\,dz_j/2\) and
\(\bar\partial f=\sum_j(f_{x_j}+i f_{y_j})\,d\bar z_j/2\).
On forms extend by the signed product rule, with the
coordinate differentials constant. Thus
\(d=\partial+\bar\partial\), and commuting real partial
derivatives gives \(\partial^2=\bar\partial^2=0\) and
\(\bar\partial\partial=-\partial\bar\partial\).
A real \((1,1)\)-form \(\omega\) is positive when
\(\omega(v,iv)>0\) for all nonzero real tangent vectors.
The curvature of a smooth Hermitian metric on a holomorphic
line, in a nowhere-zero holomorphic frame of squared norm
\(h\), is \(-i\partial\bar\partial\log h/(2\pi)\).
Changing the frame multiplies \(h\) by \(|u|^2\) for a
holomorphic unit \(u\). The logarithm of this factor is
locally the real part of a holomorphic function (use the
power series for the logarithm at a nonzero value).
Consequently the formula gives a well-defined real closed
\((1,1)\)-form.
\end{definition}

\begin{lemma}\label{lem:torus-fourier-curvature-correction}
Let \(W\) be a finite-dimensional complex vector space,
\(\Lambda\) a full real lattice, and \(\omega\) a smooth
real closed \((1,1)\)-form on \(W/\Lambda\).
Average its coefficients over a lattice parallelepiped
with normalized volume and denote the constant form by
\(E_0\). There is a real smooth periodic function \(\phi\)
such that
\[
       \omega-E_0=\frac{i}{2\pi}\partial\bar\partial\phi .
\]
If \(\omega\) is positive, then \(E_0\) is a positive
\(i\)-invariant alternating real form.
\end{lemma}
\begin{proof}
Write \(\Lambda^\vee=\{k\in W_\R^*:k(\Lambda)\subseteq\Z\}\).
In a real lattice basis the modes are
\(u_k(v)=\exp(2\pi i k(v))\). The Fourier coefficients
\(\omega_k\) in fixed constant differentials decay faster
than each negative power of \(1+|k|\). In fact for a
nonzero integral frequency one coordinate has size at
least a fixed positive multiple of \(|k|\), and repeated
integration by parts in that coordinate moves every
derivative to the smooth periodic coefficient with no
boundary term. The sum of \((1+|k|)^{-s}\) over the
lattice converges for \(s>2\dim_\C W\) by counting
successive dyadic cubes. The Fourier series and all
its derivatives converge absolutely and represent
the stated coefficients: continuous periodic functions
with zero Fourier coefficients vanish by the product
kernel argument of Lemma~\ref{lem:poisson-finite-torus}.

Write \(k=\alpha+\bar\alpha\) for its parts of type
\((1,0),(0,1)\). Closedness at each nonzero frequency
gives separately \(\alpha\wedge\omega_k=0\) and
\(\bar\alpha\wedge\omega_k=0\), by their distinct types.
Extend \(\alpha\) to a complex covector basis; comparison
of exterior basis coefficients in the two identities
makes \(\omega_k\) a scalar multiple of
\(\alpha\wedge\bar\alpha\). The latter has norm proportional
to \(|k|^2\) in a Hermitian orthonormal basis, by taking
the first covector in the direction of \(\alpha\).
Consequently the equation for this individual frequency
is solved uniquely by a coefficient \(\phi_k\), since
\[
   \partial\bar\partial u_k
      =-4\pi^2 u_k\,\alpha\wedge\bar\alpha .
\]
The solution satisfies \(|\phi_k|\le C|\omega_k|/|k|^2\).
Give \(\phi\) mean zero. Its series is smooth by the
same derivative estimates. Complex conjugation of
\(i\partial\bar\partial\phi\) gives
\(i\partial\bar\partial\bar\phi\), so reality follows
by uniqueness in the nonzero modes. The zero mode of
\(\omega\) is its stated average. A real \((1,1)\)-form
is invariant under simultaneous multiplication by \(i\)
in its two arguments, as each elementary
\(dz_j\wedge d\bar z_l\) is. Averaging preserves this
identity and gives a strictly positive value on each
fixed nonzero pair \(v,iv\) if the integrand has that
property.
\end{proof}

\begin{proposition}\label{prop:embedded-torus-invariant-hyperplane-metric}
Let \(F:W/\Lambda\to\mathbb P^{d-1}(\C)\) be a holomorphic
embedding. The line \(B=F^*\mathcal O(1)\) has a smooth
Hermitian metric whose curvature is a constant positive
real alternating form \(E_0\). This constant is independent
of the original metric used to compute its average curvature.
\end{proposition}
\begin{proof}
The line \(\mathcal O(1)\) is the dual of the tautological
line in \(\C^d\). Its pullback has the metric dual to
the Euclidean metric. For a local nonzero holomorphic
lift \(f\) of \(F\) the dual frame has squared norm
\(1/\sum|f_j|^2\). The curvature of
Definition~\ref{def:line-metric-conventions} is positive.
Indeed in real tangent directions \(v,iv\) its value is
\[
 \frac{\|f\|^2\|df(v)\|^2-
                    |\langle f,df(v)\rangle|^2}{\pi\|f\|^4}.
\]
This follows by differentiating the logarithm of the
squared norm in the complex line of direction \(v\).
The numerator is the squared length of the component
of \(df(v)\) perpendicular to \(f\), times \(\|f\|^2\).
It is strictly positive for \(v\ne0\), since the zero
derivative of the chart ratios \(f_j/f_l\) would mean
that \(df(v)\) is a multiple of \(f\), contrary to immersion.
Lemma~\ref{lem:torus-fourier-curvature-correction} now
provides \(\phi\) for this curvature. Multiplying the
squared norm by \(\exp(\phi)\) subtracts
\(i\partial\bar\partial\phi/(2\pi)\), as required.
Two smooth metrics on a line differ by a globally defined
positive smooth function. The second derivatives of its
logarithm have zero average by periodicity in lattice
coordinates, proving independence.
\end{proof}

\section{Period factors and the signs at half-periods}

\begin{lemma}\label{lem:flat-line-frame-on-vector-space}
A holomorphic line bundle on a complex vector space
with a smooth Hermitian metric of zero curvature has a
global holomorphic frame of constant unit norm.
Any two such frames differ by a constant of modulus one.
\end{lemma}
\begin{proof}
In a holomorphic frame with squared norm \(h\), write
\(u=\log h\). By Definition~\ref{def:line-metric-conventions},
\(\bar\partial\partial u=0\), so the closed form
\(\partial u\) has holomorphic coefficients (the
Cauchy--Riemann equations in each variable, or the
one-variable Cauchy formula successively in a polydisc).
On a ball it has a holomorphic primitive \(v\), given
by the radial integration formula in
Lemma~\ref{lem:projective-group-period-uniformization}.
As \(u\) is real, \(d(u-v-\bar v)=0\).
The original frame times \(\exp(-v)\) therefore has
constant norm, which can be scaled to one. Two unit
frames differ by a holomorphic function of modulus one,
constant on each connected overlap by the maximum
principle of Lemma~\ref{lem:analytic-tools}.
Starting with a frame value at zero, continue it along
a polygonal path in these unit frames, matching values
at each subdivision point. Refinements do not change
the result because transition constants agree on each
connected overlap used by the path. Two paths to the
same point are homotopic with fixed endpoints by linear
interpolation. Subdivide their compact parameter square
into sufficiently small rectangles that each rectangle
maps into one frame patch: a finite covering by smaller
open sets with closures in the patches gives a uniform
radius by compactness and uniform continuity. Transport
around each such boundary is the identity in its one
frame; cancellation of paired edges gives equality along
the outside paths. Continuation thus gives a single
global unit frame, holomorphic by the local construction.
The constant ratio argument proves uniqueness up to a unit
complex scalar.
\end{proof}

\begin{theorem}\label{thm:recover-projective-period-factors}
Let \(F,B,W,\Lambda,E_0\) be as in
Proposition~\ref{prop:embedded-torus-invariant-hyperplane-metric}.
Define the Hermitian form, conjugate linear in its first
argument, by
\[
             H(v,w)=E_0(v,iw)+i E_0(v,w).
\]
It is positive definite, \(E_0(\Lambda,\Lambda)\subseteq\Z\),
and \(B\) on \(W/\Lambda\) is the quotient of the trivial
line on \(W\) by
\[
 (z,u)\longmapsto
 (z+\lambda,\ \eta(\lambda)
      \exp(\pi H(\lambda,z)+\tfrac{\pi}{2}H(\lambda,\lambda))u),
\]
where \(|\eta(\lambda)|=1\) and
\(\eta(\lambda+\mu)=(-1)^{E_0(\lambda,\mu)}
\eta(\lambda)\eta(\mu)\).
If \(B\) is symmetric then \(\eta(\lambda)=q_B(\lambda/2+\Lambda)\)
with \(q_B\) as in Proposition~\ref{prop:intrinsic-pairings-all-orders}.
If \(B\) also satisfies the parallelogram law, then
\[
 b_{m,B}(\lambda/m+\Lambda,\mu/m+\Lambda)
         =\exp\big(2\pi i E_0(\lambda,\mu)/m\big)
       \qquad(\lambda,\mu\in\Lambda,\ m\ge1).
\]
\end{theorem}
\begin{proof}
The invariance \(E_0(iv,iw)=E_0(v,w)\) and alternation
give Hermitian symmetry and complex linearity in \(w\).
Its squared norm \(E_0(v,iv)\) is strictly positive
off zero by Lemma~\ref{lem:torus-fourier-curvature-correction}.
The trivial line on \(W\) with squared norm
\(\exp(-\pi H(z,z))\) in its constant frame has curvature
\(E_0\): in a complex orthonormal basis the formula of
Definition~\ref{def:line-metric-conventions} gives
\(i\sum dz_j\wedge d\bar z_j/2=\sum dx_j\wedge dy_j\).
Pull back the metric of
Proposition~\ref{prop:embedded-torus-invariant-hyperplane-metric}
to \(W\) and tensor by the inverse of this model line.
Curvatures subtract and vanish, so
Lemma~\ref{lem:flat-line-frame-on-vector-space} gives
a holomorphic isometric identification of the pullback
of \(B\) with the model.

In this trivialization the deck transformation for
\(\lambda\) multiplies the fiber by a holomorphic unit
\(j_\lambda(z)\). Its absolute value is
\(\exp(\pi\Re H(\lambda,z)+\pi H(\lambda,\lambda)/2)\)
since the deck map preserves the metric. Dividing by
the holomorphic exponential in the theorem gives a
constant \(\eta(\lambda)\) of modulus one. Applying the
transformations for \(\lambda,\mu\) in the two possible
orders gives \(\exp(2\pi i E_0(\lambda,\mu))=1\).
Thus this real number \(E_0(\lambda,\mu)\) is an integer,
and composition in either order now yields the asserted
sign in the formula for \(\eta(\lambda+\mu)\).

For a symmetric line, compare its normalized inversion
lift with the metric pulled back by inversion.
The logarithm of the ratio of the squared norms on
the torus has \(\partial\bar\partial=0\), since inversion
preserves the constant curvature. It is constant by the
nonzero-frequency calculation in
Lemma~\ref{lem:torus-fourier-curvature-correction}; its
value at zero is zero. Thus the normalized lift is
isometric. Upstairs inversion preserves \(H(z,z)\),
so its holomorphic fiber multiplier is a constant of
modulus one and equals one by normalization. At the
half-period \(\lambda/2\), translation from
\(-\lambda/2\) back to \(\lambda/2\) multiplies the
fiber by \(j_\lambda(-\lambda/2)=\eta(\lambda)\).
The square of the normalized involution is one, so
these constants are signs.
Finally, the first calculation in the proof of
Proposition~\ref{prop:intrinsic-pairings-all-orders}
uses only the displayed transformation factors and the
integrality of their imaginary form. In detail the
constant multiplier of \(D_{\lambda/m}\) at \(\gamma\in\Lambda\)
is \(\exp(\pi H(\gamma,\lambda/m))\). The section
\(\exp(\pi H(\lambda,z))\) after \([m]\) has this
multiplier at \(m\gamma\) because \(E_0(\lambda,\gamma)\)
is integral. Translate by \(\mu/m\) and identify the
resulting fiber by the inverse multiplier at \(\mu\).
The resulting ratio is the displayed pairing.
\end{proof}

\section{The rational form and the finite frame after conjugation}

\begin{lemma}\label{lem:conjugated-full-period-lattice}
Fix a framed torus \(Y\) of
Proposition~\ref{prop:frame-independent-odd-theta-basis}
with lattice \(D=\Lambda_x\), hyperplane line \(B\), and
alternating form \(E_A\). Let \(\iota\in{\rm Aut}(\C/E)\);
apply it to projective coordinates, algebraic maps
and equations to form \(Y',B'\). Denote the lattice and
curvature form of \(Y',B'\) by \(D',E'\) using
Lemma~\ref{lem:projective-group-period-uniformization}
and Theorem~\ref{thm:recover-projective-period-factors}.
There is an \(O_E\)-linear isomorphism
\[
  c:\widehat D\longrightarrow\widehat {D'},\qquad
  E'_p(cu,cv)=\epsilon_p E_{A,p}(u,v)
                      \quad(u,v\in D\otimes\Z_p)
\]
at every prime \(p\). Here \(\widehat D=\prod_p(D\otimes\Z_p)\),
the subscript \(p\) denotes scalar extension, and
\(\epsilon=(\epsilon_p)_p\in\widehat\Z^\times\) is uniquely
specified by \(\iota(\exp(2\pi i/m))=\exp(2\pi i\epsilon_m/m)\)
for each \(m\ge1\), with \(\epsilon_m=\epsilon\bmod m\).
On \(m\)-torsion the map on the classes of \(u/m\) is exactly
reduction of \(c\) modulo \(m\) in the two lattice coordinates.
\end{lemma}
\begin{proof}
The original torus is a compact connected complex manifold
algebraic by Corollary~\ref{cor:compact-holomorphic-graphs-algebraic}.
A proper closed algebraic subset has empty analytic
interior: a homogeneous equation nonzero somewhere
cannot vanish on an open set, by the identity theorem
in connected holomorphic coordinate charts (apply the
one-variable identity theorem successively on a polydisc).
Two proper closed subsets therefore cannot cover it.
Its homogeneous vanishing ideal is prime, as the zero
set of a product is the union of the zero sets of its factors.
Conjugation preserves primality, the smooth chart
Jacobian identities, the algebraic group and \(O_E\)-action
of Corollary~\ref{cor:complex-level-algebraic-morphisms}.
It gives a conjugate smooth irreducible projective group,
and transports the symmetric isomorphism and the
parallelogram law by the algebraicity proved in
Propositions~\ref{prop:actual-theta-parallelogram}
and~\ref{prop:intrinsic-pairings-all-orders}.
The cited uniformization and period factors thus apply.
Functoriality of uniformization gives the action on
the new lattice by derivatives.

Via the two period maps, the \(\iota\)-image of each
point of order dividing \(m\) gives an additive
\(O_E\)-linear bijection \(c_m:D/mD\to D'/mD'\), because
division of a lattice vector by \(m\) identifies each
quotient with the whole corresponding torsion group.
For \(m\mid l\), multiplication of points by \(l/m\)
corresponds to reduction from \(l\) to \(m\) and commutes
with conjugation. The integral lattices have finite bases;
in these bases the compatible matrices and their inverses
give mutually inverse matrices over \(\widehat\Z\).
Indeed at \(p\) the entries are compatible residues
at all \(p^j\), hence \(p\)-adic integers, and the Chinese
remainder identity recovers the original entries modulo
each \(m\). This gives \(c\) with the stated action.
The powers \(\epsilon_m\) exist and are units modulo \(m\)
since the automorphism takes a primitive \(m\)-th root
to a primitive \(m\)-th root. Raising roots to \(l/m\)
proves their compatibility and the same residue argument
gives \(\epsilon\).
Proposition~\ref{prop:intrinsic-pairings-all-orders}
and Theorem~\ref{thm:recover-projective-period-factors}
now give, on all lattice classes modulo \(m\),
\[
       E'(c_m u,c_m v)\equiv\epsilon_m E_A(u,v)\pmod m .
\]
Taking \(m=p^j\) for every \(j\) gives the claimed
identity over \(\Z_p\), since the intersection of its
ideals \(p^j\Z_p\) is zero.
\end{proof}

\begin{proposition}\label{prop:conjugated-hermitian-adelic-frame}
With the preceding notation, choose any \(E\)-basis of
\(D'_\Q=D'\otimes_\Z\Q\). It identifies \(D'_\Q\) with
\(E^n\). There are a frame \(x'\in{\rm GL}_n(\mathbb A_{E,f})\)
and a nonsingular Hermitian matrix \(A'\) over \(E\) with
\[
  \widehat {D'}=x'L,\qquad D'=E^n\cap x'L,\qquad
  E'(v,w)=\operatorname {Tr}_{E/\Q}(\xi\, w^*A'v),
  \qquad x'^*A'x'=\epsilon\,x^*Ax .
\]
Here the middle form identity is on \(D'_\Q\), the last
is over the finite adeles, and \(L=\widehat O_E^n\).
The labels \(t\in\mathbb A_{E,f}^n/L\) are carried by
\(\iota\) from the frame \(x\) to the identical labels
for \(x'\).
\end{proposition}
\begin{proof}
The ring action transported from \(Y\) is unital and
additive, so the action on \(D'_\Q\) extends from
\(O_E\) to \(O_E\otimes_\Z\Q=E\). The last equality
follows by multiplying the rational coordinates in an
integral basis by a common denominator. Since the
\(\Z\)-rank of \(D'\) is \(2fn\) and \([E:\Q]=2f\),
the \(E\)-dimension is \(n\). Extend \(c\) of
Lemma~\ref{lem:conjugated-full-period-lattice} by
rational scalars in each completion and put \(x'=cx\).
The chosen rational lattices, including \(O_E^n\),
are commensurable by clearing the finitely many basis
denominators in both directions. Thus \(c\) is expressed
by an \(E\otimes\Q_p\)-linear matrix at \(p\), integral
with integral inverse away from finitely many primes.
It defines an adelic matrix. The completion of \(D\)
is \(xL\) by the lattice approximation in
Lemma~\ref{lem:actual-subspace-kernel-labels}
(equivalently use rational approximation modulo each
fixed multiple of \(xL\)). Hence the completion and
intersection identities follow, the latter since a
rational number integral in all \(\Z_p\) is integral
in a lattice basis. On torsion of order dividing \(m\),
\(t\) corresponds to the class of \(mx t\) modulo \(mD\)
in completed lattice coordinates; replacing \(t\) by
any lift changes the class by \(m\widehat D\). Applying
\(c_m\) replaces \(x\) by \(x'\), proving the label assertion.

The adjoint relation \(E_A(av,w)=E_A(v,\bar a w)\)
holds for \(a\in O_E\) by the trace formula of
Proposition~\ref{prop:signed-period-riemann-form}.
It passes through \(c\) and \(\epsilon_p\) to each
completion of \(D'\), hence holds rationally for \(E'\)
and all \(a\in E\) by linearity. For \(v,w\in D'_\Q\)
define \(H_0(v,w)\in E\) by
\[
          E'(av,w)=\operatorname{Tr}_{E/\Q}(\xi aH_0(v,w))
                    \qquad(a\in E).
\]
The trace map from \(H_0\) to such functionals is
injective: if \(H_0\ne0\), take \(a=1/(\xi H_0)\)
and the trace is \(2f\). Equality of the finite rational
dimensions gives existence. Uniqueness and the adjoint
relation make \(H_0\) linear in \(v\) and conjugate
linear in \(w\). Alternation and \(\bar\xi=-\xi\)
give \(H_0(w,v)=\overline{H_0(v,w)}\) by testing against
all \(a\) and using invariance of trace under conjugation.
It thus equals \(w^*A'v\) for a Hermitian matrix.
Degeneracy would contradict the strict positivity of
the curvature on \(v,iv\) in the real scalar extension.
Finally the equality of alternating forms in
Lemma~\ref{lem:conjugated-full-period-lattice}, extended
to the full rational \(p\)-adic spaces, gives equality
of the two corresponding Hermitian matrices in frame
coordinates. Indeed one may still test with every \(a\);
the trace coefficient matrix was invertible over \(\Q\)
and remains invertible over \(\Q_p\). This proves the last
identity and in particular its scalar is a unit at every
finite place.
\end{proof}

\section{Descent of the embedded marked level varieties}

\begin{proposition}\label{prop:conjugate-is-actual-signed-period}
With \(x',A'\) as in
Proposition~\ref{prop:conjugated-hermitian-adelic-frame},
the conjugate \(Y'\), with its \(O_E\)-action, labelled
torsion and hyperplane line \(B'\), is holomorphically
isomorphic with these structures to the actual framed
torus for \(x'\) and a reflection \(J'\) of
Definition~\ref{def:signed-complex-period-space}.
\end{proposition}
\begin{proof}
In the chosen \(E\)-basis extend the identification of the
rational periods to \(D'_\Q\otimes\R\simeq V_\R\); this
is the real tangent space by
Lemma~\ref{lem:projective-group-period-uniformization}.
The complex structure \(I'\) on it commutes with
\(E\otimes\R=\prod_{\tau\in\Psi}\C\) by functoriality
for the \(O_E\)-action. The product decomposition here
is given by field embeddings: for a primitive element
the distinct roots of its minimal polynomial give the
coordinate decomposition over \(\C\) by polynomial
interpolation; the real vectors are the conjugation-fixed
pairs. Thus \(I'\) is a complex-linear matrix on each
ordinary embedding factor, with eigenvalues \(i,-i\).
On its \(i\)-eigenspace the complex Lie action of \(a\)
has eigenvalue \(\tau(a)\), and on its \(-i\)-eigenspace
it has eigenvalue \(\bar\tau(a)\) using the new complex
structure. The characteristic polynomial for an integral
primitive \(\theta\) is \(P_\theta\) of
Lemma~\ref{lem:degree-one-split-lie-polynomial}: the
differential of an algebraic map in a smooth chart is given
by its linear terms and therefore conjugates with the
coefficients. All coefficients of \(P_\theta\) belong
to \(E\), fixed by \(\iota\). The \(2f\) embeddings take
distinct values on \(\theta\); comparison of multiplicities
therefore makes \(I'=i\) at \(\tau\ne\sigma\) in \(\Psi\),
and \(I'=iJ'\) at \(\sigma\) with positive-one multiplicity
one and negative-one multiplicity \(n-1\) for \(J'\).

Write \(J_\sigma=J'\) and \(J_\tau=1\) at the other places.
The real form of the reconstructed trace formula is
\(2\Re\sum_\tau \tau(\xi)\,w_\tau^*A'_\tau v_\tau\).
Invariance under \(I'\) in both arguments implies
\(J_\tau^*A'_\tau J_\tau=A'_\tau\): test vectors supported
at that place and replace one vector by its ordinary
\(i\)-multiple to recover both real and imaginary parts.
Consequently \(A'_\tau J_\tau\) is Hermitian. Evaluating
at \(v,I'v\) gives
\(2\Im(\tau(\xi))v_\tau^*A'_\tau J_\tau v_\tau\) for
vectors supported there, as in
Proposition~\ref{prop:signed-period-riemann-form}.
Both the curvature value and \(\Im(\tau(\xi))\) are
strictly positive for nonzero \(v\), by
Theorems~\ref{thm:recover-projective-period-factors}
and~\ref{thm:actual-pel-torsion-pairing}.
This verifies the required sign at every place.
The finite equation of
Proposition~\ref{prop:conjugated-hermitian-adelic-frame}
verifies the unit similitude condition there. We have
the signed frame and its full period lattice and torsion.

It remains to identify the line, since the alternating
form alone leaves its semicharacter undetermined.
For \(\lambda\in D'\), let \(v=\lambda/2+D'\).
The algebraic identity \(q_B(a)=b_{2,B}(\beta a,a)\)
and its naturality in
Proposition~\ref{prop:intrinsic-pairings-all-orders}
give \(q_{B'}(v)=b_{2,B'}(\beta v,v)\), as the
conjugated endomorphism acts as the same \(\beta\) in
the new period coordinates. By
Theorem~\ref{thm:recover-projective-period-factors}
this is \(\exp(\pi i E'(\beta\lambda,\lambda))\) and
is exactly the multiplier sign of
Definition~\ref{def:signed-canonical-shift}. The Hermitian
form in that theorem is uniquely specified by \(E',I'\).
All period factors therefore agree with those of the
actual framed theta line. Identifying the vector-cover
trivializations gives the required line isomorphism.
\end{proof}

\begin{theorem}\label{thm:actual-common-level-equations-over-e}
Fix an open finite-index level \(C'\subseteq C\) and
\(\nu\ge0\). For each admissible multiple \(M\) of the
cutoff in Theorem~\ref{thm:common-coarse-torsion-coordinates},
the homogeneous ideal of its embedded \(\mathcal X(C',\nu)\)
is generated over \(E\), in the common projective coordinates
there (using the projective product embedding if necessary).
For each individual framed torus and each
\(\iota\in{\rm Aut}(\C/E)\), the conjugated evaluation of
every point is the canonical theta evaluation on the
reconstructed frame of
Proposition~\ref{prop:conjugate-is-actual-signed-period}.
This holds simultaneously for all torsion labels and
for arbitrary marked points.
\end{theorem}
\begin{proof}
Use \(T=L^\perp/L\), of odd order \(m=h^{2n}\), and
the labels, \(D_0\), and section \(j\) fixed in
Proposition~\ref{prop:frame-independent-odd-theta-basis}.
The conjugation of the finite labels is identified in
Proposition~\ref{prop:conjugated-hermitian-adelic-frame}.
The equality of alternating forms up to a local unit
sends the integral dual lattice to the integral dual
at every prime. Consequently on the new frame precisely
the same \(T\)-labels are the theta translations.

For \(t\in T\), the operator \(U_t\) on sections from
Proposition~\ref{prop:frame-independent-odd-theta-basis}
is induced by a lift of translation on the line.
It satisfies \(U_t^m=1\) by the product rule and
alternation, and inversion conjugates it to \(U_{-t}=U_t^{-1}\).
These properties characterize the lift operator uniquely:
two lifts differ by an automorphism of the line over
the identity, hence by a nonzero scalar on the connected
compact torus. Its \(m\)-th power is one, and conjugation
by inversion gives that it equals its inverse. Since
\(m\) is odd the scalar is one. The lifts used here are
algebraic: their holomorphic line isomorphisms between
a line and its translate are algebraic by the joint
globally generated sections argument of
Proposition~\ref{prop:actual-theta-parallelogram}.
This also applies to the line comparison of the preceding
proposition after identifying the compact projective tori
by their biholomorphic map; that map is algebraic by
Theorem~\ref{thm:holomorphic-map-algebraic-locally}.
Changing the scalar of a line identification leaves its
conjugation of operators unchanged.

Thus \(\iota\) carries the original \(U_t\) to the new
canonical \(U_t\), by the two characterizing identities.
The \(D_0\)-fixed line of sections is unique and the
whole basis is given from it by \(s_b=U_{j(b)}s_0\),
by Lemma~\ref{lem:odd-pairing-section-basis}. Conjugating
the old basis gives the new one up to a common scalar.
For evaluations the scalar and the choice of frame in
each one-dimensional value line disappear on passing to
projective coordinates. This proves the claimed identity
pointwise on each torus, with the same correspondence
for all its labels. The orbit product over \(D_M\) and
the projective product with the individual marked points
are defined by integer polynomials by
Theorem~\ref{thm:common-coarse-torsion-coordinates}, so
they commute with \(\iota\). Each conjugated point
belongs to the image, and applying \(\iota^{-1}\) proves
set equality. The closed algebraic image therefore has
its homogeneous ideal generated over \(E\), by
Lemma~\ref{lem:fixed-projective-equations}.
\end{proof}

\section{The field of definition of the Hecke maps}

\begin{lemma}\label{lem:equivariant-projective-map-descent}
Let \(X,Y\) be reduced complex projective algebraic sets
in specified coordinates with their homogeneous ideals
generated over a subfield \(k\subseteq\C\). If a regular
map \(f:X\to Y\) on complex points commutes with all
\({\rm Aut}(\C/k)\) acting on both sets, it is obtained
by scalar extension from a unique regular map of their
models over \(k\). Its projective graph has homogeneous
ideal generated over \(k\).
\end{lemma}
\begin{proof}
Use the ideal of all homogeneous equations over \(k\)
in each case. Its extension to \(\C\) is the given ideal.
The graph is closed since the diagonal in a product of
projective spaces is closed by the coordinate cross
products; pull this diagonal back along \((f,{\rm id})\).
Equivariance makes the graph invariant, so its ideal
descends by Lemma~\ref{lem:fixed-projective-equations}
using the projective product embedding.

For each target coordinate \(y_i\), the closed set
\(Z_i=f^{-1}(y_i=0)\) in \(X\) has the same invariance.
Its homogeneous ideal therefore has finite generators
over \(k\) by Lemma~\ref{lem:fixed-projective-equations}.
On any standard source chart, dehomogenize these generators.
The principal opens \(U\) of their nonvanishing cover
that chart outside \(Z_i\), with \(U\) defined over \(k\).
Take all \(i\); these opens cover over \(k\) also. Indeed
their defining functions generate the unit ideal in each
complex affine chart by the cover argument of
Proposition~\ref{prop:finite-presentation-spread}; ideal
membership for elements of a \(k\)-algebra descends from
\(\C\) by taking the quotient and observing that any
\(k\)-vector space injects into its scalar extension
(use a basis).
Let \(A\) be the localized affine chart ring over \(k\)
on one such \(U\). Each function \(f^*(y_j/y_i)\) is a
member of \(A\otimes_k\C\) on this open and is fixed
there by the coefficient action. Here equality as
functions implies equality in the reduced affine chart
ring by Lemma~\ref{lem:complex-projective-finite-fibers}'s
power-containment argument. In a \(k\)-basis of \(A\)
the finitely many coefficients of this member are fixed.
They lie in \(k\) by Lemma~\ref{lem:fixed-projective-equations}.
The resulting elements of \(A\) satisfy the target
relations and the projective transition identities on
overlaps, by the same injection into scalar extension.
These formulas glue and are unique by that injection.
\end{proof}

\begin{theorem}\label{thm:hecke-diagram-over-number-field}
In the common coordinates of
Theorem~\ref{thm:actual-common-level-equations-over-e},
forgetting level or marked slots, the group operations
on the fibers, scalar endomorphisms by \(O_E\), and every
map \(u_{g,\mathbf a}\) of
Proposition~\ref{prop:complex-actual-hecke-legs} are
defined over \(E\). This also holds for right translations
by \(k\in C\) between levels for which the translation is well defined.
All identities between their complex composites hold
over \(E\).
\end{theorem}
\begin{proof}
For a representative \(x\), conjugation by
\(\iota\in{\rm Aut}(\C/E)\) on all its marked evaluations
is the single reconstructed group with frame \(\tilde x\)
of Theorem~\ref{thm:actual-common-level-equations-over-e}.
This description already commutes with forgetting labels,
with group operations, and with the specified scalars,
since these were transported as algebraic morphisms on
the individual complex tori in
Lemma~\ref{lem:conjugated-full-period-lattice}. The same
rational period basis may be used at each level on forgetting
level. Changing \(x\) to \(xk\), for an integral level
translation \(k\) preserving the underlying lattice,
changes the associated labelling by \(t\mapsto kt\).
The reconstructed frame then changes to \(\tilde x k\)
using the same period basis, by the finite formula
in Proposition~\ref{prop:conjugated-hermitian-adelic-frame}.

For a Hecke leg first omit the scalars \(\mathbf a\).
Denote by \(\tilde u\) a frame reconstructed for the
target \(xg\) in any rational period basis. The complex
isogeny sends \(P_x(t)\) to \(P_{xg}(g^{-1}t)\) for every
finite torus label \(t\), by the lattice quotient formula
in Lemma~\ref{lem:actual-subspace-kernel-labels}; the
underlying vector is unchanged. The conjugated isogeny
therefore sends \(P_{\tilde x}(t)\) to
\(P_{\tilde u}(g^{-1}t)\). Its complex derivative sends
source periods into target periods by
Lemma~\ref{lem:projective-group-period-uniformization}.
In rational period coordinates this derivative is a
matrix \(y\) over \(E\), since it intertwines the
\(O_E\)-actions. Thus \(\tilde u^{-1}y\tilde x-g^{-1}\)
takes \(l^{-1}L\) into \(L\) at every integer \(l\ge1\),
by the label identity. Taking \(l=p^j\) at each prime
forces every entry to vanish there. Hence
\[
                  y\,\tilde x g=\tilde u.
\]
In particular \(y\) is invertible, since it is invertible
in every finite completion. Its complex-linearity in the
period structures carries the reflection at \(\tilde x\)
to that at \(\tilde u\), by
Proposition~\ref{prop:signed-ball-period-coordinates}.
Thus \(y^{-1}\) identifies the conjugated target with the
actual frame \(\tilde xg\), and in these coordinates the
conjugated isogeny is the vector identity modulo the larger
lattice. It agrees on all marked points by uniformization.
Composing with the scalars \(a_j\) gives equivariance of
the complete map, including the case \(a_j=0\).

All the maps are regular over \(\C\) by
Corollary~\ref{cor:complex-level-algebraic-morphisms} and
Theorem~\ref{thm:holomorphic-map-algebraic-locally},
using the product embedding for their domains where needed.
Apply Lemma~\ref{lem:equivariant-projective-map-descent}
with \(k=E\) to obtain the maps and their graphs. An equality
of composites can be checked on chart coordinate expressions,
which inject into their scalar extensions as in that lemma.
\end{proof}

\section{Integral models of fixed diagrams}

\begin{lemma}\label{lem:number-field-polynomial-certificates}
Let \(E\subset\C\) be a number field. A finite system of
polynomial equalities and nonvanishing conditions over \(E\)
having a complex solution has a solution in a number field
\(K\subset\C\) finite over \(E\). All entries of that
solution and all required inverse entries belong to
\(O_K[1/N]\) for some positive integer \(N\).
\end{lemma}
\begin{proof}
Replace each \(b\ne0\) by \(bv=1\) with a new variable.
Over the algebraic closure of \(E\) in \(\C\) the ideal
of these equations is proper, as evaluation at the given
complex solution takes one to one. This closure is
algebraically closed: roots in \(\C\) of polynomials with
algebraic coefficients are algebraic over a finite
extension of \(E\), hence over \(E\). The affine point
argument of Lemma~\ref{lem:complex-projective-finite-fibers}
in the generality checked in
Lemma~\ref{lem:projective-special-section-base-change}
therefore supplies a solution in that closure. The field
generated by its finitely many entries is finite over
\(E\), by successively using their minimal polynomials.
For any algebraic number \(z\), choose a positive integer
\(d\) clearing the denominators in a monic rational
equation for \(z\); substituting \(dz\) into the equation
multiplied by the corresponding power of \(d\) shows
that \(dz\) is an algebraic integer. Taking a common
multiple for all the entries, including the inverse
variables, proves the assertion.
\end{proof}

\begin{theorem}\label{thm:fixed-integral-level-model}
Fix a finite list of the levels, marked powers and maps
in Theorem~\ref{thm:hecke-diagram-over-number-field}.
One may include identifications of powers as fiber products
over their bases, together with finitely many group or
diagram identities. There is a common positive integer
\(N\) such that the diagram admits projective models and
regular maps over \(O_E[1/N]\) satisfying those identities.
For any valuation ring \(V\) with fraction field \(\C\)
and residue characteristic \(p\nmid N\), the base changes
to \(V\) admit the smooth and finite étale presentations
of Theorem~\ref{thm:actual-complex-finite-etale}, for the
maps specified there. The same holds in their geometric
fibers at every such prime of \(O_E\). In particular the
base and its marked powers are smooth there, and the
one-point family is a smooth projective commutative group
over the base, of relative dimension \(fn\).
\end{theorem}
\begin{proof}
For the last assertion take the diagram enlarged by the
zero, inversion, addition, projections, and the identities
expressing the group laws on at most three factors. The
fiber product of two powers embeds in their product by
equating their maps to the same base, a closed condition
given by projective cross products. Over \(\C\) the
forgetful projections identify the combined power with
this fiber product. Indeed the source and target have
the same local smooth period charts by
Theorem~\ref{thm:analytic-framed-torus-family}; the chosen
level has trivial stabilizer there by the neatness in
Proposition~\ref{prop:complex-actual-hecke-legs}.
The fiber product is reduced: in smooth chart coordinates
for the two projections from
Proposition~\ref{prop:complex-relative-jacobian}, their
equations after base change use disjoint variable lists
over the smooth base, with invertible minors in diagonal
blocks. This gives a smooth chart, hence a reduced chart
by the local calculation in
Theorem~\ref{thm:smooth-algebraic-projection-chart}.
The inverse, which concatenates the two point lists on
these charts, is holomorphic, hence algebraic by
Theorem~\ref{thm:holomorphic-map-algebraic-locally}.
Both maps commute with coefficient automorphisms by
Theorem~\ref{thm:actual-common-level-equations-over-e}.
Lemma~\ref{lem:equivariant-projective-map-descent} descends
the inverse as well. The product presentation itself
extends from \(E\), since its explicit diagonal equations
and their chart rings commute with extension of scalars.
We may retain that presentation and both inverse identities.

All projective equations and map formulas can now be
taken over \(E\), by
Theorem~\ref{thm:hecke-diagram-over-number-field} and the
chart proof of Lemma~\ref{lem:equivariant-projective-map-descent}.
Take their covers, relations, and transition and composite
identities as in Proposition~\ref{prop:finite-presentation-spread}.
One may take all the ideal membership witnesses for these
formulas over \(E\): in a complex polynomial identity with
coefficients to be found, keep the finitely many monomials
of the multiplier polynomials in one witness. The linear
system for their coefficients has entries in \(E\);
elimination over \(E\) gives a solution if one exists over
\(\C\). This applies to cover witnesses for one as well.
Clearing the finitely many number denominators as in
Lemma~\ref{lem:number-field-polynomial-certificates}
gives the models with their identities over \(O_E[1/N]\).

For a projection or leg retain the additional certificates
in the last paragraph of
Proposition~\ref{prop:finite-presentation-spread}. These
may use complex coefficients for auxiliary chart changes
and the nonzero minors. Fix all the finite lists of monomials
in one collection of these certificates, including the
maps between charts, both inverse identities, covering
unit expressions and all ideal membership witnesses.
With the already chosen equations and maps held fixed,
equality of coefficients in those polynomial identities
gives a finite polynomial system over \(E\) for the
additional entries. Unit conditions are expressed by
their displayed inverse identities. This includes the
relative dimensions and the monomial degree used for a
finite morphism, since the certificate fixes those
integer sizes. Apply
Lemma~\ref{lem:number-field-polynomial-certificates}.
Over a single \(O_K[1/N]\), increasing \(N\), the same
identities yield the smooth charts, finite presentations
as modules, and étale charts, exactly as in
Proposition~\ref{prop:finite-presentation-spread}.
The smooth charts for the bases over a point are included
by Theorem~\ref{thm:smooth-algebraic-projection-chart}.
The projective models and maps to which the certificates
apply still have their original \(O_E[1/N]\) equations.

The given \(V\) contains \(O_K\): a root of a monic
integer polynomial cannot have negative valuation, since
then its leading power has strictly smaller valuation
than every lower term. Also \(N\) is a unit of \(V\).
The certificates thus remain valid over \(V\). To verify
the statement in a geometric fiber at any prime over
\(p\nmid N\), extend its residue map to \(O_K[1/N]\)
in an algebraic closure of the residue field. For
completeness, localize \(O_E\) at that prime. The
localized \(O_K\) is a nonzero finite module there:
an integral basis over \(\Z\) generates already over
\(O_E\), and the inclusion is injective upon localization.
Its reduction at the maximal ideal is nonzero, by the
finite generator determinant argument \(M=\mathfrak m M
\Rightarrow M=0\) in
Lemma~\ref{lem:finite-fiber-local-algebras}. The nonzero
finite residue algebra has a maximal ideal, and its
residue field, finite over that of \(O_E\), embeds in the
chosen algebraic closure by adjoining roots successively.
This extends the residue map, and base change of the same
certificates proves the fiber statement. The commutative
group identities, including the indicated fiber products,
were kept over \(O_E[1/N]\) in the first construction.
\end{proof}

\section{Graded presentations on finitely many base strata}

\begin{lemma}\label{lem:graded-generic-free-monomials}
Let \(A\) be a noetherian domain, that is, a domain whose
ideals are finitely generated. Set \(P=A[X_1,\ldots,X_d]\)
with \(\deg X_i=1\). Let \(M\) be a finitely generated
graded \(P\)-module. Then after inverting one nonzero
element of \(A\), \(M\) is free as an \(A\)-module with a
basis given by homogeneous monomials times specified
generators. In each degree the list of basis elements is
finite. This basis remains a basis degree by degree after
every further scalar extension of the localized base.
\end{lemma}
\begin{proof}
The polynomial ring over a noetherian ring is noetherian.
For one variable and an ideal, consider the ideals of
coefficients of \(X^j\) in members of degree at most \(j\).
They increase by multiplication by \(X\) and stabilize:
the union is finitely generated and its generators occur
in finitely many stages. Through a stabilization index
choose polynomials witnessing finite coefficient generators
at each stage. Subtract multiples of these, shifted by a
power of \(X\) in higher degrees, to lower the degree of any
member of the ideal. This generates the ideal. Induct on
the number of variables. A submodule of a finite free
module is finitely generated by projecting to one coordinate
and inducting on the number of remaining coordinates.
Thus a presentation by finite homogeneous generators has
a finitely generated homogeneous relation module \(D\);
take the homogeneous parts of any finite generating list.

Use a free module \(F=\bigoplus_j P(-a_j)\) mapping
onto \(M\), where \(P(-a)_t=P_{t-a}\).
Order its terms \(X^\alpha e_j\) by total degree of the
monomial, then lexicographically by exponents, then by
\(j\); multiplication by a fixed monomial preserves this
order and every strictly decreasing sequence terminates.
Over \(K={\rm Frac}(A)\), the leading monomials of the
relations in \(D_K\) are closed under multiplication by
monomials, separately in each coordinate. By the first
argument of Lemma~\ref{lem:finite-degree-cut}, finitely
many of them generate each such set. Realize the chosen
terms by monic homogeneous relations: the homogeneous
component of a relation containing its leading term has
the same leading term. After inverting one nonzero element
of \(A\), these finitely many monic relations belong to
\(D\), by clearing the denominators of their expressions
in the original generators. Include \(1\) if the list is
empty.

Subtract multiples of the chosen relations to eliminate
divisible terms, successively taking the largest remaining
term and putting it in the remainder when it is indivisible.
Monicity makes all subtractions possible over the localized
ring. Each step strictly decreases the largest unprocessed
term. The residue classes of indivisible terms thus span
\(M\). No nonzero combination of them can belong to \(D\),
since its leading term over \(K\) would then be divisible
by a chosen leading monomial, and \(A\) injects into \(K\).
They give a homogeneous free basis, finite in each degree
since the variable and generator lists are finite.
Tensoring a free basis preserves it.
\end{proof}

\begin{proposition}\label{prop:finite-strata-graded-partial-resolutions}
Let \(S\) be a scheme covered by finitely many affine schemes
of finite type over \(O_E[1/N]\). On each specified affine
chart \({\rm Spec}(A)\) take finitely generated graded
modules over \(P=A[X_1,\ldots,X_d]\), with \(d\ge1\).
For each fixed positive integer \(l\) there is a finite
cover by locally closed reduced affine strata such that
on each stratum every indicated module \(M\), restricted
by scalar extension, has \(l\) exact sequences
\[
  0\longrightarrow K_{i+1}\longrightarrow F_i
               \longrightarrow K_i\longrightarrow0
       \quad(0\le i<l),\qquad K_0=M.
\]
Here the \(F_i\) are finite direct sums of shifts of the
graded polynomial ring, and the \(K_i\), including \(K_l\),
are free over the base with finite homogeneous bases in
each degree. The sequences stay exact after arbitrary
scalar extension of that stratum. Strata from different
original affine charts may overlap.
\end{proposition}
\begin{proof}
The coefficient rings are noetherian. Indeed \(O_E\)
is a free abelian group of finite rank by
Theorem~\ref{thm:integer-ring}; its ideals, as subgroups
of a finite free abelian group, are finitely generated by
the subgroup basis argument of
Lemma~\ref{lem:projective-group-period-uniformization}.
Finite generation persists under forming polynomial
rings by Lemma~\ref{lem:graded-generic-free-monomials},
quotients by taking images of generators, and localization
by contracting an ideal and extending it again.

Work on a reduced irreducible closed subset of one
affine chart with its domain as base ring, restricting
the module presentations by scalar extension. Choose
finite homogeneous generators to obtain
\(F_0\to K_0=M\). The kernel is a finite graded module
by the polynomial and submodule argument of
Lemma~\ref{lem:graded-generic-free-monomials}. Repeating
gives all \(l\) sequences. Apply that lemma simultaneously
to all the \(K_i\), for all indicated modules, multiplying
the finitely many nonzero denominators. On the resulting
nonempty principal open, each surjection \(F_i\to K_i\)
splits as a map of graded modules over the coefficient
ring, by choosing homogeneous preimages of its free
basis. Tensoring by any algebra over the coefficient ring
therefore preserves exactness, including the injections.

This construction covers a dense open in each irreducible
closed subset. In a noetherian affine chart descending
chains of closed subsets stabilize by the ascending chain
condition on their radical ideals. Every closed subset
is a finite union of irreducible ones: otherwise a minimal
counterexample is reducible and the two proper closed
parts have finite decompositions. The coordinate ring
of an irreducible reduced affine closed set is a domain,
since the zero set of a product is the union of the
factor zero sets. Apply induction on proper closed
subsets to the complements of the exhibited opens in
the irreducible pieces. This gives finitely many locally
closed affine strata covering every chart, with the stated
properties. Nilpotents of the original base vanish under
every map to a field and do not change this covering of
its points.
\end{proof}

\section{A finite kernel for projective sections}

\begin{definition}\label{def:graded-chart-cochains}
For \(P=A[X_1,\ldots,X_d]\), a graded \(P\)-module \(M\),
and \(J\subseteq\{1,\ldots,d\}\), let \(M_J\) be the
localization inverting \(X_j\) for \(j\in J\).
Fix a degree \(t\in\Z\). Form the augmented cochain
complex with one summand \((M_J)_t\) in degree \(|J|-1\)
for each \(J\), including \(M_t\) in degree \(-1\).
Its differential on the summand of the result indexed
by \(J=\{j_0<\cdots<j_r\}\) is the sum of the restrictions
from \(J\setminus\{j_a\}\), with signs \((-1)^a\).
Two deletions in either order have opposite signs, so
the composite is zero. Cohomology means the kernel of
the next differential modulo the image of the previous
one. Write \(H^0_{\rm proj}(M,t)\) for the kernel at
degree zero \emph{before} division by the augmented
image. It is precisely the tuples of local expressions
which agree on overlaps of the standard projective
charts. For \(M=P/I\) with \(I\) homogeneous, these are
the global sections of the hyperplane twist \(\mathcal O(t)\)
on the projective scheme cut out by \(I\): divide the
degree-\(t\) expression by \(X_i^t\) to read it in the
usual frame on \(X_i\ne0\). This convention also describes
sections of the sheaf given by a graded module by its
localized chart modules.
\end{definition}

\begin{lemma}\label{lem:high-degree-chart-section-comparison}
In Proposition~\ref{prop:finite-strata-graded-partial-resolutions}
take \(l=d+1\). There is one integer \(t_0\) for its finite
strata such that on each stratum and after any scalar
extension of its ring,
\[
                  M_t \simeq H^0_{\rm proj}(M,t)
                                      \qquad(t\ge t_0)
\]
by restriction to the charts. In fact the entire augmented
complex of Definition~\ref{def:graded-chart-cochains} is
exact in these degrees.
\end{lemma}
\begin{proof}
For a free summand \(P(-a)\) the complex splits into
subcomplexes indexed by Laurent monomials \(X^\alpha\)
with \(\sum_j\alpha_j=t-a\). Each term is a finite sum
of Laurent monomials and the number of charts is finite,
so this is a direct sum. Put \(D=\{j:\alpha_j<0\}\).
A coefficient in that subcomplex is allowed exactly on
subsets \(J\) containing \(D\). When \(D\) is proper,
choose \(v\notin D\). On a coefficient array define
the homotopy at \(J\) by zero for \(v\in J\), and
otherwise by taking the coefficient at \(J\cup\{v\}\)
times \((-1)^{|\{j\in J:j<v\}|}\). It acts between
successive degrees and includes the augmented term
when \(J\) is empty and \(D\) empty. In the sum of
differential then homotopy and homotopy then differential,
the terms deleting an index other than \(v\) cancel in
pairs, while deleting \(v\) returns the original
coefficient once. Terms outside \(J\supseteq D\)
are read as zero; adjoining \(v\notin D\) preserves
this convention. Thus the complex is exact there.
When \(t\ge a\) there can be no monomial with all
exponents negative. The augmented complex for this
free summand is then exact.

Choose \(t_0\) at least every shift of every \(F_i\) in
the finite collection. On this range the complexes for
all \(F_i\) are exact, including after scalar extension.
The short exact module sequences in
Proposition~\ref{prop:finite-strata-graded-partial-resolutions}
give short exact sequences on each cochain term.
Localization preserves exactness: multiply a fraction
with zero image by a common product of the inverted
elements to get a numerator with zero image, then lift
using exactness before localization. Taking homogeneous
parts preserves the sequences. Each short exact sequence
with exact middle complex identifies the cohomology in
degree \(j\) for \(K_i\) with that in degree \(j+1\)
for \(K_{i+1}\). To see this, lift a closed array to
the middle complex and take its differential, which is
an array in the subcomplex. Changing the lift or adding
a boundary changes its class by a boundary. A zero class
allows adjustment of the lift by a subcomplex array to
make it closed, hence a boundary in the middle complex.
Conversely a closed subcomplex array is a boundary in
the middle one, and its primitive projects to a closed
array in the quotient. Iterating \(d+1\) times from any
\(j\ge-1\) ends in degree at least \(d\), where there is no
term since \(J\) has size at most \(d\). This proves
exactness and identifies \(M_t\) with the compatible
chart sections.
\end{proof}

\begin{proposition}\label{prop:fixed-degree-projective-section-matrix}
With these strata and any integer \(u\), choose \(b\ge1\)
with \(u+b\ge t_0\). On each stratum the sections
\(H^0_{\rm proj}(M,u)\), after any scalar extension there,
are naturally the kernel of
\[
  (M_{u+b})^d\longrightarrow
       (M_{u+2b})^{\binom d2},\qquad
  (f_i)_i\longmapsto
           (X_i^b f_j-X_j^b f_i)_{i<j}.
\]
These are matrices between finite free modules over
the ring of the stratum. The fiber dimension of the
section space at field-valued points is constant on a
further finite cover by locally closed strata.
\end{proposition}
\begin{proof}
Multiply a global degree-\(u\) section \(s\) by \(X_i^b\)
to obtain the global higher-degree section \(f_i\), read
as a polynomial module element by
Lemma~\ref{lem:high-degree-chart-section-comparison}.
The indicated differences vanish. Conversely on \(X_i\ne0\)
use the local fraction \(f_i/X_i^b\). The kernel
identities give equality on all overlaps, thus a section
\(s\). On every \(X_j\ne0\),
\(X_i^bs=X_i^b f_j/X_j^b=f_i\), so the recovered higher
sections are precisely the given ones. These constructions
are inverse, including when \(d=1\).
The finite homogeneous bases and their base-change
properties are supplied by
Proposition~\ref{prop:finite-strata-graded-partial-resolutions}.
The rank of a matrix over a field is the largest order
of a nonzero minor, by row elimination. Vanishing of all
larger minors and nonvanishing of a chosen minor therefore
give finitely many locally closed conditions on which the
kernel dimension is fixed, including rank zero.
\end{proof}

\section{Dimensions of sections on the fixed integral level}

\begin{lemma}\label{lem:remove-strata-without-complex-points}
Let finitely many locally closed affine strata of a
finite-type \(O_E[1/N]\)-scheme be given. After increasing
\(N\), each stratum with a geometric point has a complex
point for the specified embedding \(E\hookrightarrow\C\).
Consequently any collection of numerical data constant
on field-valued points of each stratum, and having the
same prescribed value at all complex points, has that
value on every remaining geometric fiber. The data on a
stratum may form an infinite sequence if the whole sequence
is constant there.
\end{lemma}
\begin{proof}
If the affine ring \(A\) of a stratum has
\(A\otimes_\Z\Q\ne0\), this is a nonzero finite-type
\(E\)-algebra since \(O_E\otimes_\Z\Q=E\). Its scalar
extension to \(\C\) is nonzero by the basis injection in
Lemma~\ref{lem:equivariant-projective-map-descent}. It
has a complex point by
Lemma~\ref{lem:complex-projective-finite-fibers}'s affine
point argument. Conversely, if this localization is zero,
then \(1\) is killed in \(A\) by a nonzero integer, by the
definition of localization. Inverting that integer removes
the stratum. There are only finitely many integers to
invert. A single complex point on each retained stratum
compares all the entries that are constant there.
\end{proof}

\begin{theorem}\label{thm:uniform-good-theta-section-dimensions}
Fix \(C'\subseteq C\) as in
Theorem~\ref{thm:actual-common-level-equations-over-e},
put \(g=fn,\ d=h^n\), and use the fixed one-point family
over \(S=\mathcal X(C',0)\). Choose its integral model in
Theorem~\ref{thm:fixed-integral-level-model} inside
\(S\times\mathbb P^{d-1}\) by the common marked embedding.
Write \(B=\mathcal O(1)|_Y\) on a geometric fiber \(Y\)
over an algebraically closed field \(k\). After increasing
the single integer \(N\) of this model, simultaneously
for all such fibers one has
\[
 \dim_k H^0(Y,\mathcal O_Y)=1,\qquad
 \dim_k H^0(Y,B^{\otimes m})=d\,m^g\qquad(m\ge1).
\]
The \(d\) coordinates are a basis for \(H^0(Y,B)\).
Moreover, on each of a fixed finite affine cover of \(S\),
write \(P/I\) for the homogeneous presentation in the
marked coordinates. There is an integer \(m_0\ge1\) such
that on every geometric fiber the restriction
\((P_k/I_k)_m\to H^0(Y,B^{\otimes m})\) is an isomorphism
for all \(m\ge m_0\). Here \(I_k\) is the ideal generated
by the specialized equations.
\end{theorem}
\begin{proof}
The embedding is the product of the common base coordinates
with the individual marked projective coordinates in
Theorem~\ref{thm:common-coarse-torsion-coordinates}; take
its closed image over \(E\) by
Theorem~\ref{thm:actual-common-level-equations-over-e}.
In constructing the model one can retain the graph of the
base projection inside this product by the chart identities
of Theorem~\ref{thm:fixed-integral-level-model}. On
a standard affine of the base, dehomogenizing only the
base variables gives the required homogeneous presentation
in the marked coordinates. Equivalently one can take
equations in the projective product embedding and
substitute the pairwise products of its coordinates;
the usual charts with one nonzero coordinate in each
factor recover the same closed subscheme by division.
Take the finite affine cover by these base charts.

At a complex point the constant functions form \(\C\)
by Lemma~\ref{lem:projective-special-section-base-change}.
For \(m\ge1\) the theta functions for the \(m\)-th power
of the canonical factors are obtained by multiplying
the gauge \(G(z)\) of
Theorem~\ref{thm:canonical-theta-basis-on-ball} to the
power \(m\) and translating the standard degree-\(m\)
theta functions by the same \(d_0\). In fact raising
the multiplier identity in its proof to the \(m\)-th
power gives exactly the degree-\(m\) transformation
equations of Theorem~\ref{thm:signed-theta-sections};
division by \(G(z)^m\) and reverse translation give the
inverse on all sections. That theorem gives dimension
\(d m^g\). All these holomorphic sections are algebraic.
For any one section adjoin it to the list of all degree-\(m\)
products of the \(d\) original coordinates; the latter
generate \(B^{\otimes m}\) at every point. The resulting
holomorphic map to projective space is algebraic by
Theorem~\ref{thm:holomorphic-map-algebraic-locally}. The
ratios against a nonvanishing original product give
regular coefficients in those local frames, proving the
assertion. The degree-one coordinates themselves form a
basis by Proposition~\ref{prop:frame-independent-odd-theta-basis}.

Apply Proposition~\ref{prop:finite-strata-graded-partial-resolutions}
to \(M=P/I\) with \(l=d+1\) on the chosen affine cover.
Let \(m_0=\max(1,t_0)\) using
Lemma~\ref{lem:high-degree-chart-section-comparison}.
On a stratum the numbers of its standard monomials in
degrees \(m\ge m_0\) specify simultaneously the entire
sequence of dimensions in these degrees. They are
unchanged at any field-valued point there, and the chart
section comparison is an isomorphism. For the finitely
many degrees \(0\le m<m_0\), refine using
Proposition~\ref{prop:fixed-degree-projective-section-matrix}.
Include a finite rank test for the list of degree-one
coordinates: under the matrix for \(u=1\) each \(X_j\)
has column \((X_i^b X_j)_i\) in the finite free domain.
The dimension of their span is the rank of this list of
columns, again determined by minors. Each refinement is
given by vanishing of finitely many minors and inverting
a specified minor, so consists of finitely many locally
closed affine strata after taking their reduced structures.
Now apply Lemma~\ref{lem:remove-strata-without-complex-points}
to all the data. The complex fiber calculation gives
exactly the indicated sequence and the rank \(d\) for
the degree-one coordinate list. This proves every assertion
with one finite exclusion.
\end{proof}

\section{The section lattices along a good valuation arc}

\begin{lemma}\label{lem:smooth-connected-fiber-integral}
Let \(Y\) be a nonempty projective scheme of finite type
over an algebraically closed field \(k\), with smooth
Jacobian charts as in Proposition~\ref{prop:complex-relative-jacobian}.
If \(H^0(Y,\mathcal O_Y)=k\), then \(Y\) is reduced and
irreducible.
\end{lemma}
\begin{proof}
At a \(k\)-point the completed local ring is a formal power
series ring over \(k\) in the free tangent parameters.
Indeed translate the point to zero, split the variables
according to the unit Jacobian minor, and solve for the
dependent variables to each order in the remaining ones.
At each successive degree the linear coefficient matrix
in those variables is the same invertible minor evaluated
at zero. The corrections make all equations vanish modulo
the next power and are unique there. Inverses of the
inverted chart functions are given to each order by the
finite geometric sum after removing the nonzero constant.
This yields inverse coordinate homomorphisms modulo each
power, as in the proof of
Lemma~\ref{lem:smooth-group-frobenius-coordinates}, and
hence the assertion on taking compatible sequences.
The power series ring is a domain: the first homogeneous
parts of two nonzero series have a nonzero polynomial
product.

The noetherian local ring \(D\) at the point injects into
its completion. Put \(\mathfrak m\) for its maximal ideal
and \(I=\bigcap_{r\ge0}\mathfrak m^r\). The graded ring
\(\bigoplus_{r\ge0}\mathfrak m^r t^r\) is generated over
\(D\) by \(t\) times finite generators of \(\mathfrak m\);
it is noetherian by
Lemma~\ref{lem:graded-generic-free-monomials}. The homogeneous
ideal \(\bigoplus_r (I\cap\mathfrak m^r)t^r\) has finitely
many homogeneous generators. In any degree above their
degrees every coefficient in \(I\cap\mathfrak m^r=I\)
then belongs to \(\mathfrak m I\). Thus \(I=\mathfrak m I\)
and \(I=0\) by the determinant argument for finite modules
in Lemma~\ref{lem:finite-fiber-local-algebras}. This proves
the injection and shows that the local ring is a domain.

Every nonempty closed subset of an affine chart has a
\(k\)-point by the affine point argument of
Lemma~\ref{lem:projective-special-section-base-change}.
If a nilpotent in such a chart were nonzero, a maximal ideal
containing its proper annihilator would give a localization
where it was still nonzero, contradicting the domain
assertion. There are finitely many irreducible components
by the noetherian argument of
Proposition~\ref{prop:finite-strata-graded-partial-resolutions}.
If two components met, a \(k\)-point in their intersection
would have in an affine neighborhood two distinct minimal
primes contained in its maximal ideal. Indeed the nonempty
open part of each component is dense in it by irreducibility;
the ideals of these two maximal irreducible closed subsets
are the distinct minimal primes, since primality is exactly
irreducibility of the zero set for a radical ideal. They
remain distinct minimal primes on localization, by the
bijection between primes disjoint from the inverted elements
and primes of the localization, obtained by extending and
contracting ideals. A domain cannot have two minimal primes.
The components are thus disjoint, hence open as well as
closed. More than one would give a nonconstant idempotent
section by assigning zero on one and one on the remaining
components. The hypothesis excludes this, proving the lemma.
\end{proof}

\begin{theorem}\label{thm:arc-free-high-sections}
Let \(V\) be a valuation ring with fraction field \(\C\),
and let \({\rm Spec}(V)\to S\) be a point of the fixed model
in Theorem~\ref{thm:uniform-good-theta-section-dimensions}.
Let \(G/V\) be the pulled-back group, with relative
hyperplane line \(B\). Every geometric fiber is integral.
There is a homogeneous coordinate presentation \(R\)
for \(G\) with the following properties for \(t\ge m_0\).
The module \(R_t\) is finite free over \(V\). Its elements
give global sections of \(B^t=B^{\otimes t}\), and in every
geometric fiber \(R_t\) after scalar extension maps
isomorphically onto the complete space of these sections.
The scheme \(G\) is flat over \(V\). The same assertions
hold for each finite positive product of such groups over
\(V\), embedded by the projective product and the tensor
product of their hyperplane lines, with the largest of
their degree thresholds.
\end{theorem}
\begin{proof}
Select a base chart containing the image of the closed
point of \({\rm Spec}(V)\); its inverse image is all of
\({\rm Spec}(V)\) since an open containing the maximal
ideal of a local ring contains every prime. Take its
presentation \(R=V[X_1,\ldots,X_d]/I\) by pullback from
Theorem~\ref{thm:uniform-good-theta-section-dimensions}.
In each degree it is a quotient of a finite free module
by a finite coefficient matrix, using the fixed homogeneous
generators of \(I\). Over \(V\) any such matrix diagonalizes:
choose a nonzero entry of least valuation from the finite
list, move it to the first position, and subtract multiples
of its column and row to clear them since it divides all
entries. Repeat on the smaller rectangle until it is zero.
The rank over the fraction field equals the number of
nonzero diagonal entries, and over the residue field it
equals the number of unit diagonal entries. For \(t\ge m_0\)
these ranks coincide because the two cokernel dimensions
are the same \(dt^g\) by
Theorem~\ref{thm:uniform-good-theta-section-dimensions}
and do not change under extension of a field. Thus there
are only unit or zero diagonal entries. The quotient is
free of the stated dimension. Restriction of homogeneous
elements to local charts gives global sections, and the
same theorem gives the comparison after base change to
every geometric fiber. Smoothness there is supplied by
Theorem~\ref{thm:fixed-integral-level-model}, also in
characteristic zero by the same Jacobian certificates
over a finite extension. With \(H^0(\mathcal O)=k\),
Lemma~\ref{lem:smooth-connected-fiber-integral} gives
integrality.

Over \(X_i\ne0\) the chart ring is the direct limit of
the \(R_t\), \(t\ge m_0\), by successive multiplication by
\(X_i\): the elements are precisely fractions
\(r/X_i^t\), with equality detected after increasing the
common power. A directed limit of flat modules is flat.
Indeed tensor products commute with such limits by their
bilinear presentation, and an element killed in the
target of a directed limit of injections is already
killed at some stage and hence becomes zero in the source.
Free modules are flat by taking direct sums of any given
injection. This proves flatness on the covering charts.

For a product take the graded ring with degree \(t\) part
the tensor product over \(V\) of the constituent degree
\(t\) parts, in nonnegative degrees. It is generated in
degree one: in each factor lift a homogeneous term to a
sum of products of degree-one entries and collect one entry
per factor at each of the \(t\) positions. Its standard
charts at pure coordinate products are exactly the affine
product charts. Indeed fractions in the different factors
can be brought to one common power of their denominators;
the same common-power rule tests equality, or equivalently
the tensor product of the directed limits above is the
limit along simultaneous powers. Thus this ring presents
the projective product, with free high pieces and flat
charts. Over a field, for line bundles on a finite product
the global sections of their external tensor product are
the tensor product of the individual section spaces.
For two factors take finite standard affine covers with
affine intersections. On a product of two affines the
external line has the tensor product of the two section
modules, by their local trivializations and gluing. First
imposing agreement on overlaps in the second variable
gives the tensor product with the kernel of its restriction
differences, since tensoring over a field is exact. Then
impose the same equalities in the first variable. These
two sets of conditions imply agreement on all product
overlaps, which proves the assertion. Induction gives the
finite product. It applies both in degree \(t\ge m_0\)
and to the trivial line, giving \(H^0(\mathcal O)=k\).
The products have smooth charts by using all the equations
with diagonal Jacobian blocks. The preceding lemma now
gives integral geometric fibers for the products as well.
\end{proof}

\section{Extending the generic homomorphism by translations}

\begin{lemma}\label{lem:flat-cover-projective-descent}
Let \(W\to X\) be flat and surjective, with the inverse
image of each member of an affine open cover of \(X\)
a union of finitely many affine opens. Let \(Z\) be closed
in a projective space over a common affine base. A morphism
\(W\to Z\) whose two pullbacks to \(W\times_X W\) coincide
descends to a unique morphism \(X\to Z\).
\end{lemma}
\begin{proof}
Over a member \({\rm Spec}(A)\) of the affine cover
replace its inverse image by the disjoint union of a
finite affine cover, with coordinate ring \(B\) the
finite product of these affine rings. It is flat over
\(A\) and its spectrum surjects onto \({\rm Spec}(A)\).
Here flatness on affine charts can equivalently be
checked after localization at every point: the kernel
of a tensor injection localizes to its local kernels,
and a nonzero element of a module survives at a maximal
ideal containing its annihilator. This also shows that
the finite product is flat. It is faithfully flat:
a nonzero cyclic submodule \(A/I\) of any nonzero module
remains nonzero upon tensoring, since a maximal ideal
above \(I\) has a point above it and \(IB\ne B\); its
inclusion remains injective by flatness.

The equalizer of the two maps \(B\rightrightarrows
B\otimes_A B\) is \(A\). Indeed the sequence
\(0\to A\to B\to B\otimes_A B\) with final difference
map becomes exact after tensoring by \(B\): the new
first map is \(b\mapsto b\otimes1\), split by multiplication,
and if \(\sum b_i\otimes c_i\otimes1
=\sum b_i\otimes1\otimes c_i\), multiplying the first
two factors gives \(\sum b_i c_i\otimes1=\sum b_i\otimes c_i\).
Flatness preserves kernels and images; faithfulness
detects that their quotients vanish. The same proof
holds after localization in \(A\).

On \({\rm Spec}(B)\), take the ideal \(J\) of the inverse
image of the closed projective hyperplane where a target
coordinate vanishes. The equality on the product says
that its two extended ideals in \(B\otimes_A B\) are
equal. If \(I\) is the inverse image of \(J\) in \(A\),
then \(IB=J\). In fact tensor the kernel of \(A\to B/J\)
with \(B\). For \(b\in J\), \(1\otimes b\) lies in the
extended ideal from the second factor, hence in the one
from the first factor, and is zero in \((B/J)\otimes_A B\).
Thus \(b\in IB\); the reverse containment follows from
the definition of \(I\). Consequently the inverse image
of each standard target chart is the pullback of the
open complement of \(V(I)\) in \({\rm Spec}(A)\). These
opens cover by surjectivity. Over each principal open
inside one of them, the affine coordinate functions of
the map belong to the localized equalizer and so descend
to the source. They satisfy the target relations and
agree on overlaps, by the injection into the flat cover.
They define the desired map, uniquely by the same argument
on smaller affines. Its pullback agrees on the original
\(W\) since it agrees on its covering affines. This
uniqueness also glues the maps obtained over the different
affines of \(X\).
\end{proof}

\begin{theorem}\label{thm:arc-extension-by-translation}
Let \(G,G'\) be two group schemes along the same good
valuation ring \(V\) as in
Theorem~\ref{thm:arc-free-high-sections}, allowing distinct
points of the fixed base. Every algebraic group homomorphism
\(f_\C:G_\C\to G'_\C\) extends uniquely to a homomorphism
\(G\to G'\). Equalities of composites of such homomorphisms
and their compatibility with \(O_E\)-actions extend from
the generic fibers.
\end{theorem}
\begin{proof}
Choose a standard affine chart of the integral generic
source meeting its generic point. The target homogeneous
coordinates of \(f_\C\) there give a nonzero projective
tuple of rational functions: divide by one coordinate
nonzero at that generic point. In the fraction field
of the affine source chart, clear a common denominator;
its affine coordinate functions are polynomials in the
source coordinate ratios. Homogenizing to one degree
and multiplying further by powers of the nonzero chart
coordinate gives a tuple in \(R_t\otimes_V\C\), with
\(t\ge m_0\), representing \(f_\C\) on a dense open.
At least one entry is nonzero there. Write all entries
in a free \(V\)-basis of \(R_t\) from
Theorem~\ref{thm:arc-free-high-sections}. Divide all
coefficients by one nonzero coefficient of smallest
valuation. The resulting entries lie in \(R_t\), with
at least one unit coefficient in their collected list.
They define a map \(f_U:U\to\mathbb P^{d'-1}_V\) on the
open \(U\) where these homogeneous sections do not all
vanish; \(d'\) is the target coordinate count.
Every geometric fiber \(U_s\) is nonempty and dense:
a unit coefficient survives in every residue field
and the high-degree comparison of
Theorem~\ref{thm:arc-free-high-sections} makes the
corresponding section nonzero on that integral fiber.
On the generic part of \(U\) this map equals \(f_\C\)
by the closed diagonal and irreducibility. It maps into
\(G'\) over \(V\) as well, since all pulled-back defining
equations vanish generically and chart rings of \(U\)
are flat over \(V\). Indeed multiplication by each
nonzero scalar of \(V\) is injective on a flat module,
so every such ring injects into its fraction-field
base change.

Let
\[
 W=\{(z,u)\in G\times_VG: u\in U,\ z+u\in U\},
 \qquad h(z,u)=f_U(z+u)-f_U(u).
\]
The projection \(W\to G\) is flat, as restriction of
the base change of the flat map \(G\to{\rm Spec}(V)\).
Flatness is preserved under scalar extension by the
associativity of tensor products and on localizing by
exactness of localization. This projection is surjective:
over an algebraically closed extension of the residue
field of \(z\), the two opens \(U_s\) and \(-z+U_s\)
are dense and nonempty in the integral \(G_s\). Their
intersection has a point by the affine point argument
of Lemma~\ref{lem:projective-special-section-base-change}.
Over each standard affine in the first factor, \(W\)
has a finite affine cover. The product with the second
projective factor has its finitely many standard
affine charts. On it the conditions defining \(W\) are
nonvanishing of at least one of each of two finite lists
of sections of lines, pulled back under projection and
sum. On each affine both lines have finite trivializing
principal covers (principal opens are a basis, and an
affine has a finite principal subcover by the unit ideal
argument of Proposition~\ref{prop:finite-presentation-spread}).
On these the conditions are finite unions of principal
opens, giving the assertion.

The two induced maps \(h\) on \(W\times_G W\) agree
generically, since the generic homomorphism gives
\(h(z,u)=f_\C(z)\). This product is open in \(G^3\);
it is flat over \(V\) by
Theorem~\ref{thm:arc-free-high-sections}. The equalizer is
a closed subscheme by the closed diagonal of \(G'\);
its ideal vanishes because it vanishes after fraction-field
extension in flat charts. Apply
Lemma~\ref{lem:flat-cover-projective-descent}. The resulting
map agrees with the original homomorphism generically by
surjectivity and flat descent. It is unique by the same
flat-chart and closed-diagonal argument on \(G\). Compare
the two group-law maps on \(G^2\) by that argument to see
that the extension is a homomorphism. Identities of
composites and identities with the extended scalar
endomorphisms follow in precisely this way on their
common flat source.
\end{proof}

\section{Extension of identities of line bundles}

\begin{lemma}\label{lem:projective-line-high-global-generation}
Let \(X\) be a closed subscheme of projective space of
finite dimension over an affine ringed base \({\rm Spec}(A)\).
Write \(B=\mathcal O_X(1)\).
For any line bundle \(D\) on \(X\), all sufficiently large
\(t\) give finitely many sections generating \(D\otimes B^t\)
at every point. In particular \(D^{-1}\) embeds into
a finite direct sum of copies of \(B^t\) with locally
free quotient.
\end{lemma}
\begin{proof}
Use the finitely many standard affines \(U_i\) where
the coordinate \(X_i\) does not vanish. On an affine
\(U={\rm Spec}(C)\) a line bundle has a finite principal
cover where it is free. Indeed shrink each trivializing
open to principal opens and choose a finite subcover as
in Proposition~\ref{prop:finite-presentation-spread}.
Its module of sections, described by compatible tuples
in free frames on this cover, localizes at \(a\in C\)
to the module on \(D(a)\). This follows by localizing
the kernel of the finite array of differences on the
pairwise intersections; these intersections are again
principal, the transition functions localize, and
localization is exact as in
Lemma~\ref{lem:high-degree-chart-section-comparison}.
In particular, clearing powers of each covering function
in the unit frame on its principal open gives a section
on \(U\) which generates on that open. These finitely
many sections generate on all of \(U\).

Take any such generator \(s\) for the restriction of
\(D\) to a standard chart \(U_i\), meaning one member
of that chart's finite generating list. On \(U_j\) the
overlap with \(U_i\) is the principal open of \(X_i/X_j\).
By the localization just proved, a suitable power
\((X_i/X_j)^a\) times the restriction of \(s\) extends
to a section \(s_j\) on \(U_j\). Choose one exponent
large enough for all charts. These sections, interpreted
as local sections of \(D\otimes B^a\) in the frames
\(X_j^a\) of \(B^a\), agree on overlaps after restriction
to \(U_i\). On each full overlap \(U_j\cap U_k\), which
is affine, their difference consequently vanishes after
inverting \(X_i/X_j\). Some power of this function kills
the difference there. Increase the common twist by the
largest of these finitely many powers, multiplying all
local sections by that power of the global coordinate
\(X_i\). They now glue to a global section agreeing
with \(sX_i^t\) on \(U_i\). Repeating for the finite
generator lists and multiplying any section by further
powers of its own \(X_i\) makes all their twists equal
to any sufficiently large \(t\). They still generate
on their respective charts since \(X_i\) is nonzero there.

The resulting surjection from a finite trivial bundle
onto \(D\otimes B^t\) splits locally, since near each
point some generator is a frame. Dualizing and tensoring
by \(B^t\) gives the asserted embedding with locally
free quotient, directly in these split charts.
\end{proof}

\begin{theorem}\label{thm:arc-picard-injection}
Let \(X/V\) be a group or a finite nonempty product of
groups in Theorem~\ref{thm:arc-free-high-sections}.
The restriction \({\rm Pic}(X)\to{\rm Pic}(X_\C)\)
is injective. More precisely, a nonzero global generic
section of any line on \(X\) has a scalar multiple
extending over \(V\) with nonzero restriction to the
closed fiber.
\end{theorem}
\begin{proof}
For the section assertion use
Lemma~\ref{lem:projective-line-high-global-generation}
with the inverse of the given line \(D\), obtaining
\(D\hookrightarrow(B^t)^{\oplus r}\) with locally free
quotient. Here \(B\) is the projective hyperplane line,
and we choose \(t\) above the threshold in
Theorem~\ref{thm:arc-free-high-sections}. The generic
section is a nonzero tuple of generic sections of \(B^t\).
Write each entry in the generic basis from the free
homogeneous piece \(R_t\) of that theorem. Dividing by
a coefficient of smallest valuation gives a tuple in
\(R_t^{\oplus r}\) with at least one unit coefficient.
It defines global sections of \((B^t)^{\oplus r}\) over
\(V\). Their image in the quotient bundle is zero,
since it is zero generically and locally the free
quotient over a flat \(V\)-chart injects into the
fraction-field base change.
In the locally split embeddings the tuple therefore
defines a section of \(D\), agreeing on all overlaps
by injectivity. Its special restriction is nonzero:
the tuple reduces to a nonzero section of a coordinate
line \(B^t\) in at least one position, by the closed-fiber
comparison in Theorem~\ref{thm:arc-free-high-sections}.
The same proof permits an algebraically closed extension
of the residue field for that test.

If \(D\) is trivial on \(X_\C\), take a generic trivializing
section and its inverse. Apply the section assertion to each
of them. The product of the resulting integral sections is
a regular function on \(X\) which is generically a nonzero
constant \(c\in\C\). Evaluating on the identity section
of the product shows \(c\in V\). Subtracting this constant
gives zero everywhere by flatness. Over an algebraic closure
of the residue field the two reduced sections are nonzero
and their common nonvanishing locus is nonempty, since
this special fiber is integral by
Theorem~\ref{thm:arc-free-high-sections}. Their product
therefore does not vanish identically there and \(c\)
is a unit in \(V\). The two integral sections multiplied
by an inverse unit give mutual inverses. Hence \(D\)
is trivial over \(V\). Taking the quotient of any two
lines with isomorphic generic restrictions proves
injectivity.
\end{proof}

\section{The polarized maps on good valuation arcs}

\begin{proposition}\label{prop:actual-arc-theta-identities}
Fix the integral level family of
Theorem~\ref{thm:uniform-good-theta-section-dimensions}
and a valuation ring \(V\subset\C\) with fraction field
\(\C\) and residue characteristic \(p\nmid 2hN\).
Every complex point of its base, and every chosen point
of its corresponding torus, extends uniquely over \(V\).
Write \(G,B\) for the resulting integral group and relative
hyperplane line. Then \(B\) is symmetric, satisfies the
parallelogram law of Definition~\ref{def:normalized-line-class},
and \([m]^*B\simeq B^{\otimes m^2}\) for every \(m\ge1\).
The action of each \(O_E\)-scalar and the projective
translations by all \(T\)-labels extend; on the reduced
fiber their coordinates use the matrices of
Corollary~\ref{cor:odd-theta-matrices-specialize}. Inversion
has its prescribed permutation matrix. These coordinates
give the complete linear embedding on that fiber.
\end{proposition}
\begin{proof}
For a complex point of a closed projective model, divide
all homogeneous coordinates by one nonzero coordinate of
smallest valuation. The entries then lie in \(V\) and
one is a unit. They define a section over \(V\), since
every defining polynomial vanishes by its generic value
and \(V\subset\C\). Two such sections agree by the
closed diagonal and injection into the fraction field.
Apply this first to the base and then to the projective
fiber over \(V\); this works for any finite list of
points, with the group-law equalities preserved by the
same uniqueness.

The generic line is symmetric by
Proposition~\ref{prop:intrinsic-pairings-all-orders} and
satisfies the algebraic parallelogram law by
Proposition~\ref{prop:actual-theta-parallelogram}.
Apply Theorem~\ref{thm:arc-picard-injection} to \(G\) and
to \(G\times_V G\) to extend the line isomorphisms for
symmetry and for the two universal points in the
parallelogram law. The origin fibers contributing to
the normalized law are pulled back from a line on \(V\),
trivial over the local base since a frame at its closed
point trivializes on an open containing that point.
The class formula \([m]^*B\simeq B^{\otimes m^2}\)
follows by applying the law with \(x=jz,\ y=z\) at
the universal \(z\): \(l((j+1)z)+l((j-1)z)=2l(jz)+2l(z)\).
Induction with \(l(0)=0\) gives \(l(jz)=j^2l(z)\).

Each scalar action is obtained from the generic
endomorphism by Theorem~\ref{thm:arc-extension-by-translation}.
Their ring identities follow there by uniqueness, including
addition by comparison with the group sum. Each labelled
translation is translation by its extended section.
The equalities between its embedded coordinates and
the projective action of the generic translation matrix
hold over \(V\), since these matrices have unit entries
in their nonzero positions and unit determinant by
Corollary~\ref{cor:odd-theta-matrices-specialize}.
Here the matrix on evaluated homogeneous coordinates is
the action dual to the specified action on sections,
with pullback following the convention of
Proposition~\ref{prop:frame-independent-odd-theta-basis}.
Both sides define maps into projective space; equality
on the generic fiber extends by the closed diagonal and
the flatness of Theorem~\ref{thm:arc-free-high-sections}.
The inversion permutation extends in the same way.
Degree-one completeness on the fiber is
Theorem~\ref{thm:uniform-good-theta-section-dimensions};
the reduced translation matrices span the full endomorphism
algebra by Corollary~\ref{cor:odd-theta-matrices-specialize}.
\end{proof}

\begin{proposition}\label{prop:actual-arc-hecke-polarizations}
Let \(w\) be a safe degree-one place over a completely split
prime \(p\) as in Lemma~\ref{lem:degree-one-split-lie-polynomial},
with \(p\nmid 2hN\). Take a valuation ring as in
Proposition~\ref{prop:actual-arc-theta-identities} inducing
this place of \(E\). For every framed complex base point
and residue subspace \(A\subseteq\mathbf F_p^n\), form
the integral groups \(G,G_A\) by extending that point
and its complex neighbour \(xg_A\) on the fixed base.
The unscaled map extends to a unique \(O_E\)-linear
homomorphism \(\psi_A:G\to G_A\), and on their hyperplane
lines
\[
                      \psi_A^*B_A\simeq B^{\otimes p^{\dim A}}.
\]
For adjacent \(A\subset D\) the complex intermediate
isogeny also extends over \(V\), composing with \(\psi_A\)
to \(\psi_D\).
\end{proposition}
\begin{proof}
Use Theorem~\ref{thm:arc-extension-by-translation}
on the algebraic complex homomorphisms of
Corollary~\ref{cor:complex-level-algebraic-morphisms}.
The nested lattices and their intermediate complex map
are described in Lemma~\ref{lem:relative-hecke-line-lattices}.
The scalar identities and the composite identity persist
by uniqueness on the flat domains. By
Proposition~\ref{prop:complex-prime-edge-pairing} the
alternating form of the target in the common covering
space is \(p^{\dim A}\) times the source form. Its complex
structure is the same, so its Hermitian form has the same
multiple, and its canonical \(\eta_\beta\) is the
\(p^{\dim A}\)-th power of the initial one by
Definition~\ref{def:signed-canonical-shift}.
Restricted to the source period lattice its automorphy
factors in Proposition~\ref{prop:actual-theta-parallelogram}
are therefore those of the claimed tensor power.
This identifies the generic lines holomorphically and
algebraically by the joint section argument of that
proposition: both lines are generated by pullbacks
or products of the homogeneous coordinates.
Theorem~\ref{thm:arc-picard-injection} extends their
equality of classes over \(V\).
\end{proof}

\section{Finiteness and the image of a polarized map}

\begin{lemma}\label{lem:arc-projective-polarized-finite}
Suppose \(G,G'\) are as in
Theorem~\ref{thm:arc-extension-by-translation}, and
\(f:G\to G'\) satisfies \(f^*B'\simeq B^{\otimes\mu}\) for
their relative hyperplane lines and some integer \(\mu\ge1\).
Then \(f\) is finite.
\end{lemma}
\begin{proof}
In any geometric fiber of \(f\), the restriction of
\(B^\mu\) is trivial by the displayed identity. Take a
reduced irreducible component of that fiber, projective
over its algebraically closed field. Its global functions
are constant by
Lemma~\ref{lem:projective-special-section-base-change}.
In a frame of the trivial restriction, all degree-\(\mu\)
monomials in the projective coordinates are constants.
These monomials distinguish the original points: on a
chart with \(X_j\ne0\) recover \(X_i/X_j\) as
\(X_iX_j^{\mu-1}/X_j^\mu\). Thus each component has
one point, and there are finitely many components by
Proposition~\ref{prop:finite-strata-graded-partial-resolutions}'s
noetherian argument over the geometric field. Nilpotent
elements do not change the set of points.

The graph makes \(G\) projective over \(G'\). We give the
finite-presentation argument also for this possibly
nonnoetherian valuation base. Over an affine chart of
\(G'\) it is closed in projective space over its ring.
One can use homogeneous equations for this closed subscheme:
on each standard chart take elements of its ideal. Each
such local element extends, after multiplying by a
sufficient power of that chart coordinate, to a
homogeneous polynomial vanishing on the closed subscheme.
Indeed homogenize a polynomial representing the local
element; its restriction to the quotient on each other
chart is zero after inverting the first coordinate.
Some power kills it there, and a common power works on
the finitely many charts. The homogeneous vanishing ideal
therefore gives exactly the original chart ideals.

At a point \(y\) of the target, its residue field \(K_y\)
is infinite: it contains the residue field of the
localization of \(V\) at the corresponding prime. That
localization is a valuation ring of the algebraically
closed field \(\C\), and its residue field is algebraically
closed because a lift of a monic polynomial splits into
linear factors with integral roots, by the valuation
argument of Theorem~\ref{thm:fixed-integral-level-model}.
An algebraically closed field is infinite by the argument
of Proposition~\ref{prop:valuation-special-cluster-lengths}.
A linear coordinate over \(K_y\) can therefore be chosen
nonzero at all points of the finite geometric fiber:
the product of the finitely many nonzero linear
conditions is a nonzero polynomial on the dual space
over an algebraic closure, and cannot vanish on the
whole grid over an infinite subfield, by induction using
the one-variable root bound. Lift its change of coordinates
over an affine neighborhood of \(y\), shrinking to make
the determinant and all denominators units. Write \(X_0\)
for this first coordinate.

Modulo \(X_0\), over the algebraic closure of \(K_y\),
the remaining homogeneous equations span all forms of a
sufficiently high degree by homogeneous multiples, by the
empty-projective-fiber argument of
Lemma~\ref{lem:complex-projective-finite-fibers} and its
algebraically closed field version in
Lemma~\ref{lem:projective-special-section-base-change}.
Only finitely many multiples are needed. The resulting
full-size minor is nonzero already over \(K_y\), and
lifts to a unit on a smaller affine neighborhood.
Consequently \(X_0\ne0\) everywhere on the graph there,
and the same span relations express every monomial of
that degree on \(X_0=1\) by terms of lower degree over
the base. The graph is closed in this affine chart and
its ring is generated as a module by the finitely many
monomials of smaller degree, exactly as in the cited
finite-fiber lemma. The argument also applies to an
empty fiber, or a projective space of dimension zero
as in that lemma. These neighborhoods prove finiteness.
\end{proof}

\begin{lemma}\label{lem:good-fiber-function-field-degree}
Let \(Y\) be a geometric fiber of the one-point family
in Theorem~\ref{thm:uniform-good-theta-section-dimensions},
over \(k\). Its field of rational functions \(k(Y)\),
the fraction field of a nonempty affine open, has
transcendence degree \(g\) over \(k\). The rational
function field of any nonempty proper irreducible closed
subset has transcendence degree at most \(g-1\).
Consequently a finite map between two such geometric
fibers over \(k\) is surjective.
\end{lemma}
\begin{proof}
For a nonempty integral closed projective set \(Z\)
write \(D_Z\) for its homogeneous coordinate domain
in its chosen projective coordinates, the quotient by
all homogeneous vanishing equations. It is finite as
a graded module over a polynomial ring in some \(s\ge1\)
algebraically independent linear forms. To see this,
if the current degree-one generators are dependent,
take a nonzero homogeneous relation, obtained by
collecting homogeneous parts. Make it monic in the last
variable by a linear change: since \(k\) is infinite its
highest homogeneous part takes a nonzero value in some
direction. Taking this direction as the last basis
vector and dividing by that value gives the required
monic equation. Its bounded powers generate over the
subalgebra in the remaining variables. Repeat on that
subalgebra, stopping at independence as in
Lemma~\ref{lem:complex-polynomial-normalization}. The
module generators may all be homogeneous by this
construction. At least one independent generator remains,
since a nonzero degree-one coordinate is transcendental
over \(k\) by the grading in a domain. Counting monomials
in the \(s\) independent forms and through the finitely
many shifted generators gives positive lower and upper
constant multiples of \(m^{s-1}\) for \(\dim_k(D_Z)_m\)
as \(m\) grows. Here the count for degree \(m\) is
\(\binom{m+s-1}{s-1}\), obtained by separating \(m\)
identical units into \(s\) ordered parts. The fraction
field of \(D_Z\) has transcendence degree \(s\) by the
finite module extension. For a nonzero linear coordinate
\(L\), \(D_Z[1/L]\) is the Laurent polynomial ring in
\(L\) over its degree-zero part: divide each homogeneous
element by the power of \(L\) of its degree, and distinct
degrees cannot cancel. Its degree-zero part is the affine
coordinate ring on \(L\ne0\). Hence \(k(Z)\) has
transcendence degree \(s-1\).

For \(Z=Y\) the high homogeneous pieces have precisely
dimension \(dm^g\) by
Theorem~\ref{thm:uniform-good-theta-section-dimensions}.
Indeed the given ideal and the full homogeneous vanishing
ideal have equal pieces there: a polynomial vanishing
pointwise gives zero as a section on the reduced fiber
by Lemma~\ref{lem:smooth-connected-fiber-integral}
and the affine point argument of
Lemma~\ref{lem:projective-special-section-base-change};
the high-degree restriction is injective. Thus \(s-1=g\).
For a proper closed integral subset choose a homogeneous
equation \(P\) of positive degree \(e\) vanishing on it
but nonzero on \(Y\), from its closed defining equations.
Multiplication by \(P\) in the domain \(D_Y\) is injective,
so in all sufficiently large degrees its quotient has
dimension \(d(m^g-(m-e)^g)=O(m^{g-1})\). The coordinate
ring of the subset is a further quotient, giving the
claimed upper bound by the same monomial count.

The image of the finite map on geometric points is
closed: the graph is projective, and the empty-fiber
minor argument of Lemma~\ref{lem:projective-special-section-base-change}
gives an open neighborhood of each point outside its
image. The closed image is irreducible since the source
is. Over a nonempty affine chart of that image the source
has a finite coordinate algebra, a domain, over the
image coordinate domain. Indeed the finite presentation
as a module restricts from the target, and equations
vanishing on the image are zero on the reduced source.
The ring map is injective by density of the image
there. The source fraction field is therefore finite
over the image fraction field: localize at all nonzero
elements of the base domain, and a finite-dimensional
domain over a field is a field by injectivity of
multiplication by any nonzero element. Both fields
have the same transcendence degree, so the image cannot
be a proper closed subset. This gives a preimage for
every closed point, and for arbitrary scheme points by
the same argument after algebraic closure of the residue
field of the point, since both geometric fibers have
the same section dimensions after that extension.
\end{proof}

\section{The degree of an isogeny along the arc}

\begin{theorem}\label{thm:arc-polarized-isogeny-degree}
Let \(G,G'\) be two groups obtained from the same fixed
level model along \(V\), and let \(u:G\to G'\) be a
homomorphism with \(u^*B'\simeq B^{\otimes\mu}\), where
\(\mu\) is a positive integer. Then \(u\) is surjective
and finite locally free of rank \(\mu^g\), with
\(g=n[F:\Q]\) the common relative dimension. Its kernel is finite free of that
rank over \(V\).
\end{theorem}
\begin{proof}
By Lemma~\ref{lem:arc-projective-polarized-finite} the map
is finite, and by Lemma~\ref{lem:good-fiber-function-field-degree}
it is surjective on each geometric fiber and hence on
the whole scheme. On a fixed geometric base fiber over
an algebraically closed field \(k\), write \(u:Y\to Z\)
for the map. Over a nonempty standard affine of the
integral \(Z\), its finite direct image algebra
\(D=u_*\mathcal O_Y\) is free of some positive rank
\(r\) on a nonempty open. In detail choose among finite
module generators a basis of the generic vector space
and clear denominators in expressions for all other
generators. Equality can be checked in the fraction
field because the coordinate algebra of the preimage is
a domain containing the target ring, as in
Lemma~\ref{lem:good-fiber-function-field-degree}. Shrinking
to that denominator open gives the claimed basis.
Choose a \(k\)-point of this open, using the affine point
argument of Lemma~\ref{lem:projective-special-section-base-change}.
Every other \(k\)-point is its translate by \(u(y)\) for
some \(y\in Y(k)\), since the nonempty finite fibers
contain \(k\)-points. Translating in \(Y\) above the
corresponding translation in \(Z\) is an isomorphism of
the finite morphisms since \(u\) is a group map. The same
freeness and rank therefore hold on neighborhoods of
every closed point, which cover \(Z\) since any nonempty
closed complement has a closed point. Thus \(D\) is a
vector bundle of rank \(r\).

There is an integer \(a\) and injections of sheaves on
\(Z\) with \(L=B'|_Z\),
\[
             (L^{-a})^{\oplus r}\lhook\joinrel\longrightarrow D
                \lhook\joinrel\longrightarrow (L^a)^{\oplus r}.
\]
Indeed the proof of
Lemma~\ref{lem:projective-line-high-global-generation}
applies to a finite locally free sheaf of any rank:
use a finite set of basis sections on each trivializing
principal open, clear powers of the covering function
and extend with the same powers of the chart coordinates.
It gives finite global generators of both \(D\otimes L^a\)
and \(D^\vee\otimes L^a\), increasing \(a\) as needed.
In each list choose \(r\) sections independent at the
generic point. The first list gives the left map, the
second gives the right one by evaluating the linear forms.
Both are generically isomorphisms, hence injective in
the integral charts since the sources are locally free.
Taking sections after tensoring with \(L^m\) gives, for
\(m>a\),
\[
 r\,d(m-a)^g
     \ \le\ \dim_k H^0(Y,(u^*L)^m)
     \ \le\ r\,d(m+a)^g .
\]
Here \(H^0(Z,D\otimes L^m)=H^0(Y,(u^*L)^m)\):
on an affine open trivializing \(L\) these sections are
the functions on its preimage in the pulled-back frame,
and the transition functions are the same. The two outside
dimensions follow from
Theorem~\ref{thm:uniform-good-theta-section-dimensions}.
The middle dimension is \(d(\mu m)^g\) by the pullback
identity and that same theorem. Divide by \(d m^g\)
and let \(m\) increase to get \(r=\mu^g\). This argument
works over each geometric residue field and over \(\C\).

Over an affine chart of \(G'\) with ring \(A\), let \(M\)
be the finite algebra of the preimage. At each chosen
prime of \(A\) the dimension of its fiber algebra is
\(\mu^g\), since extension of the residue field to an
algebraic closure preserves that dimension. Lift a basis
to module generators on an affine neighborhood: first
clear the denominators of residue representatives in the
local ring, and use the finite-module determinant argument
of Lemma~\ref{lem:finite-fiber-local-algebras} to get
generation in that local ring. The finitely many original
module generators are then expressed by the lifts after
inverting a common element outside the prime. Denote the
new affine ring again by \(A\). The surjection
\(A^{\mu^g}\to M\) is an isomorphism after extending
\(V\) to \(\C\), since there it is a surjection between
locally free modules of equal rank. In local bases its
matrix modulo each maximal ideal is invertible, so its
determinant is a unit locally. The kernel over \(A\) is
therefore zero, since \(A^{\mu^g}\) injects into this
generic localization by the flatness of \(G'\) over
\(V\) in Theorem~\ref{thm:arc-free-high-sections}.
This proves local freeness over the entire integral
target without any noetherian hypothesis on \(V\).
Pulling back by the identity section gives a finite
locally free kernel over the local ring \(V\); the open
trivializing its closed point covers its spectrum.
\end{proof}

\begin{corollary}\label{cor:actual-arc-finite-free-hecke-kernels}
For the groups of Proposition~\ref{prop:actual-arc-hecke-polarizations}
each \(\psi_A\) is surjective finite locally free of rank
\(p^{g\dim A}\). On every such group, multiplication by
\(m\ge1\) is surjective finite locally free of rank
\(m^{2g}\), also when \(m\) is divisible by \(p\).
The finite flat multiplication and neighbour kernel
hypotheses of Theorem~\ref{thm:polarized-edge-first-frobenius}
are thus satisfied on these valuation arcs.
\end{corollary}
\begin{proof}
Use Theorem~\ref{thm:arc-polarized-isogeny-degree} with
\(\mu=p^{\dim A}\) for \(\psi_A\) by
Proposition~\ref{prop:actual-arc-hecke-polarizations}.
For \([m]\) use \(\mu=m^2\) from
Proposition~\ref{prop:actual-arc-theta-identities}.
In particular this supplies the kernels of \([p]\)
and \([p^2]\) over \(V\) for every neighbour.
The nested kernel and quotient statements follow by
Proposition~\ref{prop:integral-adjacent-hecke-kernel}
using exactly these maps, the extended scalar action
and the smooth separated models of
Theorem~\ref{thm:fixed-integral-level-model}.
\end{proof}

\section{Frobenius on reductions of the actual marked neighbours}

\begin{theorem}\label{thm:actual-arc-point-congruence}
Fix the coarse level \(C\), the number \(\nu\ge0\) of marked
points, and a single admissible \(M\) for the common
embedding of
Theorem~\ref{thm:common-coarse-torsion-coordinates}.
There is a finite set of rational primes such that the
following holds at every safe degree-one \(w\) outside
it, of residue characteristic \(p\). Suppose \(V\subset\C\)
is a valuation ring with fraction field \(\C\) inducing
\(w\) on \(E\), and write \(k\) for its residue field.
For a framed complex point \((x,J,(P_j))\), let
\(\bar z_A\in\mathcal X(C,\nu)(k)\) be the reduction,
in the fixed projective model, of its neighbour
\((xg_A,J,(\psi_A P_j))\), for each
\(A\subseteq\mathbf F_p^n\). With \({\rm Fr}\) acting on
the projective coordinates by \(p\)-th powers, there
exists a line \(W\subseteq\mathbf F_p^n\) such that
\[
    {\rm Fr}^{\,n-\dim A}\bar z_A
         ={\rm Fr}^{\,n-\dim(A+W)}\bar z_{A+W}
                       \qquad(A\subseteq\mathbf F_p^n).
\]
Consequently in the free abelian group on the geometric
points of this special fiber one has
\[
  \sum_{r=0}^n (-1)^r p^{r(r-1)/2}
                  \sum_{\dim A=r}[{\rm Fr}^{\,n-r}\bar z_A]=0 .
\]
The finite exclusion and the model in these statements
are independent of \(A,w,V\) and the initial framed point.
\end{theorem}
\begin{proof}
Use the projective models for the fixed level with its
base, one-point family, marked power and group operations
given in Theorem~\ref{thm:fixed-integral-level-model}.
Exclude their denominators, the primes in
Theorem~\ref{thm:uniform-good-theta-section-dimensions},
the divisors of \(2hM\), and the finite splitting exceptions
of Lemma~\ref{lem:degree-one-split-lie-polynomial}.
Each remaining degree-one prime is completely split
over its rational prime. The valuation residue field
is algebraically closed by the monic lifting argument
of Lemma~\ref{lem:arc-projective-polarized-finite}.
The fixed model and coordinates reduce over \(\mathbf F_p\),
so coordinate \(p\)-powering preserves the reduced
equations and sends their \(k\)-points to \(k\)-points.

Propositions~\ref{prop:actual-arc-theta-identities}
and~\ref{prop:actual-arc-hecke-polarizations} and
Corollary~\ref{cor:actual-arc-finite-free-hecke-kernels}
supply the integral groups, scalar action, signed theta
lines and polarized isogenies on this same \(V\).
Their special fibers are connected reduced and projective
by Theorem~\ref{thm:arc-free-high-sections}. In particular
the complete matrix and section hypotheses of
Theorem~\ref{thm:marked-theta-frobenius-tuples} hold by
Proposition~\ref{prop:actual-arc-theta-identities}.
Its pairing and the opposite-prime kernel comparison
therefore apply to these very neighbours.

Add to the chosen sections \(P_j\) all source torsion
sections of labels in \(\mathcal T[M]\), as well as the
\(T\)-translations. They are defined over \(V\) by
Proposition~\ref{prop:actual-arc-theta-identities}.
At every place away from \(p\) the chosen \(g_A\) is
the identity, since \(b_w(a_A)\) is supported on the
one conjugate pair and \(z_w\) has its only nonunit
entry at \(w|_{K_0}\) and equals one elsewhere, as in
Theorem~\ref{thm:actual-central-test-vector} and
Lemma~\ref{lem:safe-operators-away-from-markers}.
The component of a prime-to-\(p\) torsion label at a
place over \(p\) is zero, since multiplication by its
order is a local unit. Thus \(g_A^{-1}\) preserves each
of these labels individually, as \(p\nmid hM\).
The label action of the generic isogeny in
Lemma~\ref{lem:actual-subspace-kernel-labels} and
uniqueness of integral point extension identify its
images with the indicated target labels.
Theorem~\ref{thm:marked-theta-frobenius-tuples} now
gives a single line \(W\) and the displayed Frobenius
identities for the individual torus evaluations of
all these sections.

To pass to the actual coarse base coordinate take their
ordered tensor over \(\mathcal T[M]\) and the product of
the corresponding linear forms over the finite permutation
group \(D_M\), as in
Theorem~\ref{thm:common-coarse-torsion-coordinates}.
This procedure commutes with reduction and with Frobenius.
Indeed each evaluation has homogeneous representatives
in \(V\) with a unit coordinate. The reduced tensor is
nonzero by selecting one such coordinate in each factor.
Each permuted linear form is nonzero after reduction
and their product is nonzero since a polynomial ring
over a field is a domain. The resulting coefficient
tuple thus also has a unit entry; its reduction is the
projective specialization, with every operation given
by integer polynomials. This argument is valid even
if some factors coincide after reduction. Keeping the
individual marked coordinates and taking the product
embedding gives exactly the \(\bar z_A\) by uniqueness
of projective specialization. The equality for these
points follows. Finally apply
Corollary~\ref{cor:pointwise-congruence-single-line} with
\(q=p\). All exclusions were made for the fixed objects
and coordinates before choosing the place.
\end{proof}

\section{All points of the reduced correspondence graph}

\begin{lemma}\label{lem:horizontal-graph-valuative-centers}
Let \(R=O_E[1/N]\), and let \(H_E\) be a reduced closed
subscheme of a projective space over \(E\), or in a
finite projective product over \(E\). Take its reduced
closure \(H\) over \(R\), that is, the reduced closed
subscheme on the closure of the underlying generic
subset in the corresponding projective model. For every
point \(z\in H\) there exist a valuation ring \(V\) with
fraction field \(\C\) and a map \({\rm Spec}(V)\to H\)
whose closed point maps to \(z\) and induces an embedding
of residue fields. Its generic point gives a complex
point of \(H_E\) for the prescribed embedding of \(E\).
\end{lemma}
\begin{proof}
The generic scheme has finitely many irreducible components
by the noetherian argument of
Proposition~\ref{prop:finite-strata-graded-partial-resolutions}.
The union of their closures is the closure of their union,
so choose a component whose closure contains \(z\).
The reduced structure on that closure is integral:
the closure of an irreducible set is irreducible, and
the radical ideal on each nonempty affine chart is prime
by irreducibility. Its local ring \(A\) at \(z\) is a
domain whose fraction field \(K\) is finitely generated
over \(E\), because its generic component has a chart
of finite type over \(E\). The field \(K\) embeds in \(\C\)
over the given \(E\): choose a maximal algebraically
independent subset of finitely many field generators,
choose the same number of independent complex numbers
over \(E\), then extend by roots of the remaining algebraic
generators. Such independent numbers exist successively
since the algebraic closure in \(\C\) of a field generated
over the countable \(E\) by finitely many elements is
countable (there are only countably many polynomials),
whereas the reals are uncountable by the diagonal argument
on binary sequences embedded using the ternary digits
zero and two.

Inside \(\C\), extend \(A\) to a maximal local subring
\(V\) dominating it. Domination means inclusion with
contraction of the larger maximal ideal equal to the
smaller one. Local subrings so ordered admit as upper
bound of a chain their union with the union of its maximal
ideals, which is local since every element outside that
union has its inverse at some stage. The maximal-element
principle of Definition~\ref{def:field-conventions} applies.
For \(x\in\C^\times\) at least one of \(x,x^{-1}\) belongs
to \(V\). Indeed if \(\mathfrak m V[x]\) is proper, take
a maximal ideal containing it and localize \(V[x]\)
there. Its contraction is \(\mathfrak m\), the maximal
ideal of \(V\), so maximality gives \(x\in V\).
Otherwise an identity \(1=a_0+\cdots+a_jx^j\) with
\(a_i\in\mathfrak m\) makes \(x^{-1}\) integral over
\(V\) by multiplying by \(x^{-j}\) and dividing by the
unit \(1-a_0\). Thus \(V[x^{-1}]\) is a nonzero finite
\(V\)-module and \(\mathfrak m V[x^{-1}]\) is proper
by the finite-module determinant argument of
Lemma~\ref{lem:finite-fiber-local-algebras}. The same
localization now gives \(x^{-1}\in V\). This proves the
valuation property and also that its fraction field
is \(\C\). The map from the local ring and domination
give the stated center and injection
\(\kappa(z)\hookrightarrow V/\mathfrak m\). The generic
map factors through the chosen generic component by
construction of the embedding of \(K\).
\end{proof}

\begin{theorem}\label{thm:actual-special-graph-point-congruence}
Use the fixed level, \(\nu,M\) and finite exceptional
set of Theorem~\ref{thm:actual-arc-point-congruence}.
At each indicated \(w\), choose a fine complex level
\(C_m\) as in Proposition~\ref{prop:complex-actual-hecke-legs}
on which all the maps with labels \(A\subseteq\mathbf F_p^n\)
are defined. Let \(H_w\) be the reduced closure in the
product of fixed integral marked targets of the joint
generic image of these maps. For every geometric point
\(z\) of its fiber at \(w\), with projections \(z_A\)
to those targets, there is a line \(W\subseteq\mathbf F_p^n\)
such that
\[
     {\rm Fr}^{\,n-\dim A}z_A
       ={\rm Fr}^{\,n-\dim(A+W)}z_{A+W}
                       \qquad(A\subseteq\mathbf F_p^n).
\]
In particular the sum in
Theorem~\ref{thm:actual-arc-point-congruence} vanishes on
this entire graph fiber. The result also holds on projecting
a joint graph that records any further finite list of
coarse-target coordinates defined over \(E\) on a common
refined level.
\end{theorem}
\begin{proof}
The complex joint map is algebraic and defined over \(E\)
by Theorem~\ref{thm:hecke-diagram-over-number-field}.
Its complex image is closed: take its projective graph
and project to the target product. Outside the image
the homogeneous equations in the source variables have
empty fiber, and the maximal minor witnessing this in
one degree remains invertible on a target neighborhood,
as in Lemma~\ref{lem:projective-special-section-base-change}.
The image is invariant under \({\rm Aut}(\C/E)\) by
equivariance, hence its ideal descends by
Lemma~\ref{lem:fixed-projective-equations}. Use this
closed reduced image over \(E\) in the closure; its
complex points are exactly the tuples arising from
the complex fine level. The same statements hold for
the additional maps in the last assertion, taking the
joint map on the common fine level.

Let \(z_0\) be the underlying scheme point of \(z\).
Apply Lemma~\ref{lem:horizontal-graph-valuative-centers}
to \(z_0\) and the joint image, using the product
embedding of Corollary~\ref{cor:compact-holomorphic-graphs-algebraic}.
The generic tuple is an actual complex tuple, so choose
one complex fine-level preimage and a framed representative.
Projection of the \(V\)-point to each fixed projective
target is the unique extension of its generic value by
Proposition~\ref{prop:actual-arc-theta-identities}'s
coordinate scaling argument. Its reductions in
\(V/\mathfrak m\) thus satisfy the line identities of
Theorem~\ref{thm:actual-arc-point-congruence}.
These identities already hold in \(\kappa(z_0)\) for the
same finite-indexed line \(W\): in residue coordinates
each equality is tested by homogeneous cross products
of coordinate powers, and
\(\kappa(z_0)\hookrightarrow V/\mathfrak m\) is injective.
They persist under the given field embedding from
\(\kappa(z_0)\) defining the geometric point \(z\).
This gives the claimed equalities, and the sum vanishes
by Corollary~\ref{cor:pointwise-congruence-single-line}.
The argument treats the whole joint graph at once and
therefore compares its projected coordinates in the
same valuation ring also when extra maps are present.
\end{proof}

\section{Projective models of finite diagrams of correspondences}

\begin{theorem}\label{thm:integral-hecke-cover-refinements}
Fix the coarse marked model \(\mathcal Z=\mathcal X(C,\nu)\)
over \(R=O_E[1/N]\) of
Theorem~\ref{thm:actual-arc-point-congruence}. Suppose
finitely many coarse-target Hecke maps \(u_j\) are defined
on the fine generic level \(X_m=\mathcal X(C_m,\nu)_E\)
of Proposition~\ref{prop:complex-actual-hecke-legs}. Include
its forgetful map \(\pi_m\) in the list. There is a reduced
projective \(R\)-scheme \(\mathcal W\) with generic fiber
\(X_m\), an action of the finite group \(D=C/C_m\), and
maps to \(\mathcal Z\) extending all \(u_j\) and their
translates by \(D\). In characteristic zero the map
\(\pi_m\) is a finite étale cover with \(D\) acting simply
transitively on each geometric fiber. Any finitely many
such models admit a common projective refinement,
surjecting onto each of them and respecting all the
specified coarse maps from their respective generic levels.
The ring \(R\) can be fixed independently of their supports.
\end{theorem}
\begin{proof}
Use any of the closed projective embeddings
\(i:X_m\hookrightarrow\mathbb P^b_E\) of
Theorem~\ref{thm:actual-common-level-equations-over-e}.
On the generic level send \(v\) to the tuple
\[
       \big((i(vk))_{k\in D},\ (u_j(vk))_{(j,k)}\big),
\]
where in the second list \(j\) runs over the prescribed
maps and \(k\) over \(D\). All entries are morphisms over
\(E\) by Theorem~\ref{thm:hecke-diagram-over-number-field}.
The factor \(i(v)\), together with the closed graphs of
the other maps, makes this a closed immersion. Explicitly,
on its first factor the point must lie in \(i(X_m)\),
and equating the other factors with their prescribed
images is a closed diagonal condition. Take its reduced
closure in the product of copies of \(\mathbb P^b_R\) and
\(\mathcal Z\). It is projective under the projective
product embedding. It recovers the original reduced
generic image since that image is already closed there.
Right translation \(v\mapsto va\) permutes the factors
by sending the coordinate at \(k\) to the value originally
at \(ak\). This permutation and its inverse preserve the
closed generic image, hence also its reduced closure, and
their compositions give the right group action.
Projection gives all the desired maps. The equalities
between the \(\pi_m(vk)\) extend to the whole model:
on a reduced closure, a regular function vanishing on the
dense generic subset vanishes on every affine chart, by
restriction to each of its integral components as in
Lemma~\ref{lem:horizontal-graph-valuative-centers}.

The generic cover over \(\C\) has the asserted degree,
fibers and transitive free action by
Theorem~\ref{thm:actual-complex-finite-etale} and
Proposition~\ref{prop:complex-actual-hecke-legs}.
These properties hold after passage to any algebraically
closed extension of \(E\). Indeed finite étale charts can
be certified over a finite extension of \(E\) as in
Theorem~\ref{thm:fixed-integral-level-model}, and hence
over each such algebraically closed field. The assertion
on the action is equivalently that \((v,k)\mapsto(v,vk)\)
from the finite disjoint union of the \(X_m\) is an
isomorphism onto \(X_m\times_{\mathcal Z_E} X_m\).
Over \(\C\) this map is a bijection and is locally the
identity on the analytic covering charts. Its inverse is
holomorphic and thus algebraic by
Theorem~\ref{thm:holomorphic-map-algebraic-locally}.
The schemes in that fiber product are reduced by the étale
charts over the smooth base. The inverse commutes with
field automorphisms by uniqueness, so
Lemma~\ref{lem:equivariant-projective-map-descent} descends
it to \(E\). Cover surjectivity persists after scalar
extension: for any point its nonempty finite algebra
fiber stays nonzero over a field extension. Nonemptiness
at each point over \(E\) itself can be tested after
embedding its finitely generated residue field in \(\C\)
by Lemma~\ref{lem:horizontal-graph-valuative-centers}.
Finiteness over \(E\) follows from the homogeneous
finite-fiber argument of
Lemma~\ref{lem:arc-projective-polarized-finite}: each
residue field contains the infinite field \(E\), and the
geometric fibers are finite. On an affine target chart
its finite algebra is locally free of rank \(|D|\).
Lift a basis at a point to generators on a neighborhood
by the argument of Theorem~\ref{thm:arc-polarized-isogeny-degree};
the resulting surjection becomes an isomorphism after
extending to \(\C\), by the rank and freeness there.
Its kernel is already zero by faithful scalar extension.
In the residue algebra at any point choose an element
whose values at all the \(|D|\) geometric fiber points
are distinct. Such an element is a linear combination
of algebra generators: for each pair at least one
generator distinguishes them, and finitely many proper
linear conditions can be avoided over the infinite
residue field. Its powers of degrees \(0,\ldots,|D|-1\)
form a basis by the Vandermonde determinant over the
split reduced geometric fiber. Lift the element on
an affine neighborhood. Its powers give a basis on
inverting their coordinate determinant. The algebra
is the quotient by one monic polynomial of degree
\(|D|\), from expressing the next power in that basis.
Its derivative is a unit after shrinking again:
multiplication by the derivative has nonzero determinant
at the point since the geometric roots are distinct.
This is an étale chart over \(E\).

For a finite list of models choose a finer factorial level
mapping to all their generic levels, using the nested
subgroups in Proposition~\ref{prop:complex-actual-hecke-legs}.
Include their generic forgetful maps, now with targets
the already constructed projective models, and all their
deck translates as extra factors in the same construction.
The generic forgetful maps are the level projections of
the period quotients and are surjective: each point has
a framed representative giving a lift. Projection of the
new reduced closure to each earlier model is projective
(the graph is closed in a projective product) and has
closed image. At an empty geometric fiber the homogeneous
equations in the projected-away variables span all
monomials of a fixed degree after base change to an
algebraic closure, as in
Lemma~\ref{lem:arc-projective-polarized-finite}; one minor
keeps that fiber empty in an open neighborhood. This
proves the closed-image assertion at arbitrary residue
fields. The image contains the dense generic fiber, so
it is the whole earlier model. The specified composite
identities persist by closed diagonals and density on the
reduced source. All new ambient projective spaces and
products were defined over the original \(R\).
\end{proof}

\section{Cancellation on irreducible components}

\begin{proposition}\label{prop:componentwise-special-hecke-maps}
Fix \(w\) as in Theorem~\ref{thm:actual-special-graph-point-congruence}.
Take any cover \(\mathcal W\) of
Theorem~\ref{thm:integral-hecke-cover-refinements} recording
all of its coarse maps \(u_A\) and the source projection \(c\).
Over an algebraically closed extension \(k\) of \(\mathbf F_p\)
let \(Z\) be the coarse special target and let
\(D\) be any reduced irreducible component of
\(\mathcal W_k\). Write \(\Phi:Z\to Z\) for the map over
\(k\) obtained by base change of the coordinatewise
\(p\)-power Frobenius of the \(\mathbf F_p\)-model.
There is a line \(W_D\subseteq\mathbf F_p^n\) such that,
as morphisms on \(D\),
\[
     \Phi^{\,n-\dim A}u_A|_D
          =\Phi^{\,n-\dim(A+W_D)}u_{A+W_D}|_D
                      \qquad (A\subseteq\mathbf F_p^n).
\]
\end{proposition}
\begin{proof}
The projection of \(\mathcal W\) to the product of the
coarse targets factors through \(H_w\) of
Theorem~\ref{thm:actual-special-graph-point-congruence}.
Indeed it does so on the generic fiber by definition of
the joint image, and the defining homogeneous equations
vanish on the reduced closure by density, as in
Theorem~\ref{thm:integral-hecke-cover-refinements}.
Take the point of \(D\) at the algebraic closure of its
function field induced by its generic point. Apply
Theorem~\ref{thm:actual-special-graph-point-congruence}
to its image in \(H_w\). Its resulting line \(W_D\)
works for all \(A\) simultaneously. Equality of the
morphisms already holds over the function field by
injectivity into its algebraic closure. Their equalizer
is closed since the target is projective, and every
rational function in any integral affine chart that
vanishes generically vanishes there. Thus its ideal
pulls back to zero on all charts of \(D\). This proves
equality as morphisms on the entire component. The
coordinate power map used here is over \(k\): it raises
the homogeneous coordinate variables to \(p\)-th powers
and fixes the scalars in \(k\); the defining equations
over \(\mathbf F_p\) ensure that it maps into \(Z\).
\end{proof}

\begin{definition}\label{def:special-image-cycles}
For a projective scheme \(X\) over an algebraically
closed field, a cycle is an element of the free abelian
group on its integral closed subvarieties. For an integral
projective \(D\) and a morphism \(b:D\to X\), let \(C\)
be its reduced closed image. Define the image cycle of
\([D]\) to be zero if \(k(D)/k(C)\) has positive
transcendence degree, and otherwise put
\[
                  b_*[D]=[k(D):k(C)]\,[C].
\]
The image is closed by the projective empty-fiber
argument of Theorem~\ref{thm:integral-hecke-cover-refinements},
and is irreducible by continuity. Pullback of rational
functions injects \(k(C)\) into \(k(D)\) by density.
The fields have finite type over the base field,
so an algebraic extension here has finite degree,
by adjoining finitely many generators successively.
\end{definition}

\begin{corollary}\label{cor:actual-special-correspondence-cycles}
In Proposition~\ref{prop:componentwise-special-hecke-maps}
assign an integer \(m_D\) to each reduced irreducible
component \(D\) of the special cover. There is an identity
of cycles on \(Z\times_k Z\),
\[
 \sum_D m_D\sum_{A\subseteq\mathbf F_p^n}
      (-1)^{\dim A}p^{\binom{\dim A}{2}}\,
      (c,\Phi^{\,n-\dim A}u_A)_*[D]=0 .
\]
Before taking image cycles, the same inner sum of the
corresponding maps is zero in the free abelian group on
morphisms \(D\to Z\times_k Z\), for each single component.
\end{corollary}
\begin{proof}
On \(D\) the two maps paired by adjoining \(W_D\) are
identical by Proposition~\ref{prop:componentwise-special-hecke-maps};
their first component \(c|_D\) is fixed throughout.
Proposition~\ref{prop:line-cancellation-frobenius} applied
to the free group on these morphisms gives the second
assertion. Sending each generator to its image cycle
in Definition~\ref{def:special-image-cycles} defines a
homomorphism of abelian groups and preserves the zero
relation. There are finitely many components by the
noetherian argument of
Proposition~\ref{prop:finite-strata-graded-partial-resolutions};
multiplying their relations by \(m_D\) and summing gives
the formula, including for the lengths at the generic
points if those multiplicities are chosen.
\end{proof}

\section{Elliptic curves with quadratic multiplication and finite level}

\begin{definition}\label{def:cm-elliptic-finite-frames}
Let \(K=\Q(\sqrt{-d})\) with its fixed embedding of positive
imaginary square root, let \(O=O_K\), and put
\(\widehat O=\prod_v O_v,\ \mathcal T_K=\mathbb A_{K,f}/\widehat O\).
The product is over finite places. For \(x\in\mathbb A_{K,f}^\times\)
let \(\mathfrak a_x=K\cap x\widehat O\), a fractional
\(O\)-ideal, and put \(Y_x=\C/\mathfrak a_x\).
Its torsion label \(P_x(t)\) for \(t\in\mathcal T_K\)
is the unique class of \(z\in K/\mathfrak a_x\) with
\(x^{-1}z=t\bmod\widehat O\). For an integer \(m\ge1\)
set \(H_m=\{b\in\widehat O^\times:b-1\in m\widehat O\}\)
and \(\mathcal C_m=\mathbb A_{K,f}^\times/(K^\times H_m)\).
A label of level \(m\) means an element of \(\mathcal T_K[m]\).
The symbol \(H_m\) here denotes the quadratic idele subgroup
used in Proposition~\ref{prop:central-character-actual-form}.
\end{definition}

\begin{lemma}\label{lem:elliptic-cubic-period-algebraic}
For any full lattice \(\mathfrak a\subset\C\), the functions
\((\wp_{\mathfrak a},\wp'_{\mathfrak a})\) identify
\(\C/\mathfrak a\) biholomorphically with a smooth
projective plane cubic, taking zero to its point at
infinity. The torus law and every holomorphic
homomorphism between two such tori are algebraic in
their cubic coordinates.
\end{lemma}
\begin{proof}
Use the series and its termwise computations in
Proposition~\ref{prop:scaled-cm}. Before the choice of a
scale in that proof they apply to every lattice by taking
an oriented basis \(\lambda,\lambda\tau\).
Write \(x=\wp_{\mathfrak a}, y=\wp'_{\mathfrak a}\), and
\(e_1,e_2,e_3\) for the values at the nonzero half-periods.
Then \(y^2=4\prod_{j=1}^3(x-e_j)\) and the roots are
distinct by the double-zero argument of
Proposition~\ref{prop:half-recurrence}. The homogeneous
cubic \(Y^2Z=4\prod(X-e_jZ)\) has no affine singular
point: both partial derivatives can vanish on it only
at a repeated root. Its only point at infinity is
\([0:1:0]\), where the derivative in \(Z\) is nonzero.
The poles of orders two and three at the origin make
\(x/y\) a local parameter there, with linear term
\(-z/2\), and \(1/y\) holomorphic and vanishing.
Thus the map extends to the indicated point with
nonzero derivative. Elsewhere its derivative is
nonzero, since away from half-periods \(x'=y\ne0\)
and at a half-period \(y\) has a simple zero by the
same period computation.

For each finite \(a\in\C\) the function \(x-a\) has two
zeros counted with multiplicity by the pole count of
Lemma~\ref{lem:elliptic-count}. When \(a\) is not a branch
value these are \(z,-z\) with opposite nonzero \(y\);
at a branch value there is the double zero of the indicated
half-period. This accounts bijectively for every affine
point of the curve. The local holomorphic inverse argument
of Theorem~\ref{thm:signed-fiber-projective-embedding}
then gives the claimed isomorphism. Now use
Theorem~\ref{thm:holomorphic-map-algebraic-locally}
on these compact projective manifolds, and their projective
product for the sum map, to obtain the algebraic morphisms.
\end{proof}

\begin{proposition}\label{prop:cm-elliptic-ray-classification}
The definitions of \(P_x\) give additive \(O\)-linear
bijections from \(\mathcal T_K\) onto the torsion of
\(Y_x\). At level \(m\), the isomorphism classes of
the curves \(Y_x\) equipped with their \(O\)-action
and the individual labels \(P_x(t)\) for \(t\in\mathcal T_K[m]\)
are exactly \(\mathcal C_m\), which is finite. More
generally every smooth connected complex projective
elliptic group with unital \(O\)-action by the given
embedding on its tangent line and an \(O\)-linear
bijection from \(\mathcal T_K[m]\) to its \(m\)-torsion
represents one of these classes. The same classification
for all labels simultaneously has equivalence \(x\sim ax\),
\(a\in K^\times\).
\end{proposition}
\begin{proof}
In the valuation factorization of
Proposition~\ref{prop:general-ideal-factorization},
\(\mathfrak a_x\) has at each finite place the local
valuation of \(x\), so its completion is \(x\widehat O\).
It is a full lattice since it lies between two rational
integer multiples or reciprocal multiples of \(O\).
Modulo \(\mathfrak a_x\), the torsion points are exactly
\(K/\mathfrak a_x\). The assignment \(z\mapsto x^{-1}z
\bmod\widehat O\) is injective there. For surjectivity
multiply a proposed lift in \(\mathbb A_{K,f}\) after
applying \(x\) by an integer to clear its denominators
and those of \(x\widehat O\). Approximation of an integral
tuple modulo a fixed integer multiple of \(\widehat O\)
by a member of \(O\) holds by the Chinese remainder
identity in an integral basis, as in
Lemma~\ref{lem:actual-subspace-kernel-labels}. Choose that
multiple inside the cleared lattice and undo the
denominator. This gives all labels and proves additivity
and \(O\)-linearity.

For an elliptic group as in the statement, connectedness
and smoothness imply irreducibility by
Lemma~\ref{lem:smooth-connected-fiber-integral} and constancy
of global functions on a connected projective manifold
by the maximum principle. Use
Lemma~\ref{lem:projective-group-period-uniformization}:
the periods in its tangent line have rational rank two.
Their rational span is a one-dimensional \(K\)-space
because the action extends rationally to \(K\) and
is ordinary multiplication by the given embedding
on that line. Dividing by one nonzero rational period
therefore identifies the lattice with a fractional
\(O\)-ideal in \(K\). Choose an idele \(x\) of its local
generators; existence follows from the ideal factorization
and the local uniformizers in
Proposition~\ref{prop:general-ideal-factorization}.
The finite labels may differ from \(P_x\) by an automorphism
of \(\mathcal T_K[m]\). This module is \(O/mO\) via
multiplication of representatives by \(m\) and reduction
in \(\widehat O/m\widehat O\). An automorphism is
multiplication by a unit of that finite ring. Lift the
unit locally at primes dividing \(m\) and choose one at
the other primes. This gives \(b\in\widehat O^\times\),
and \(P_{xb}(t)=P_x(bt)\), realizing the given labels.
For compatible labellings on all torsion, the corresponding
units are compatible under reduction, by multiplication
of torsion points by \(l/m\) for \(m\mid l\). They
give one \(b\in\widehat O^\times\) by the same residue
limit argument as in
Lemma~\ref{lem:conjugated-full-period-lattice}.

A pointed \(O\)-group isomorphism lifts to its complex
derivative by the period uniformization. It sends one
fractional ideal onto the other, hence is multiplication
by an element \(a\in K^\times\). The induced finite
unit \(x'^{-1}ax\) fixes the level labels exactly when
it lies in \(H_m\), and it fixes all labels exactly
when it is one, by testing modulo \(p^j\) at every
prime. Conversely any such scalar gives the required
isomorphism, algebraic by
Lemma~\ref{lem:elliptic-cubic-period-algebraic}.
Finiteness of \(\mathcal C_m\) is the ideal-class and
finite unit-residue calculation in
Proposition~\ref{prop:central-character-actual-form}.
\end{proof}

\section{Rigid coordinates for the quadratic classes}

\begin{lemma}\label{lem:rigid-two-three-cubic}
Fix \(6\mid m\), two distinct nonzero labels of order two
\(e_1,e_2\in\mathcal T_K[m]\), their sum \(e_3\), and a
label \(t_3\) of order three. On a labelled complex curve
of Proposition~\ref{prop:cm-elliptic-ray-classification}
there are unique rational functions \(X,Y\) with the
following properties. They are regular off the origin,
\(X\) has pole order at most two and takes the values
zero and one at the first two labels, and \(Y\) has pole
order at most three, is odd under inversion, and takes
the value one at \(t_3\). They embed the affine complement
of the origin as \(bY^2=X(X-1)(X-\lambda)\), where, writing
\(u=X(P_x(t_3))\), one has
\[
          \lambda=X(P_x(e_3)),\qquad b=u(u-1)(u-\lambda).
\]
Here \(b\lambda(\lambda-1)\ne0\). Over any algebraically
closed field of characteristic different from two, on
a smooth cubic of this displayed form marked by the
same coordinate normalizations, with origin at infinity
and inversion \((X,Y)\mapsto(X,-Y)\), the same uniqueness
holds.
\end{lemma}
\begin{proof}
In Lemma~\ref{lem:elliptic-cubic-period-algebraic} let
\(a_i=\wp(P_x(e_i))\). Take
\(X=(\wp-a_1)/(a_2-a_1)\) and
\(Y=\wp'/\wp'(P_x(t_3))\). The denominator in \(Y\)
is nonzero since the finite zeros of \(\wp'\) occur
exactly at the nonzero points of order two. The cubic
equation in that lemma, evaluated at \(t_3\) as well,
gives the displayed relation and the nonvanishing.
The linear affine coordinate changes extend invertibly
to the projective plane and put the origin at \([0:1:0]\).

In homogeneous coordinates \([X_h:Y_h:Z_h]\), put
\(t=X_h/Y_h,\ v=Z_h/Y_h\) near that point. The equation
is \(bv=t(t-v)(t-\lambda v)\). Rewriting it gives
\(v(b+(1+\lambda)t^2-\lambda tv)=t^3\), whose parenthesis
is a unit there. The local maximal ideal is generated
by \(t\), and \(v/t^3\) is a unit. The pole orders of
\(X=t/v\) and \(Y=1/v\) are thus two and three.
On the affine complement every regular function is an
element of its coordinate ring and has a unique expression
\(A(X)+B(X)Y\) for polynomials \(A,B\), by monic division
in \(Y\) after dividing by the nonzero \(b\).
The highest poles in the two parts have even and odd
orders, respectively, hence cannot cancel. It follows
that all functions of pole order at most two there are
spanned by \(1,X\), and the odd ones of pole order at
most three by \(Y\), since the characteristic is different
from two. The imposed values pick one element in each
case. This proof applies in all the indicated fields.
\end{proof}

\begin{theorem}\label{thm:quadratic-labelled-cubic-tuples-algebraic}
For the chosen labels of Lemma~\ref{lem:rigid-two-three-cubic},
the tuple consisting of \(b,\lambda\) and the coordinates
\((X(P_x(t)),Y(P_x(t)))\) for all nonzero
\(t\in\mathcal T_K[m]\) distinguishes precisely the elements
of \(\mathcal C_m\). All its entries are algebraic over
\(K\). The permutations of these tuples by coefficient
automorphisms fixing \(K\) factor through a finite Galois
extension of \(K\). The group laws and the \(O\)-action on
each normalized cubic are defined over the algebraic
closure of \(K\) in \(\C\).
\end{theorem}
\begin{proof}
For clarity, every pointed holomorphic map between compact
complex tori is a group map. In their translation coordinates
the matrix of its derivative is globally holomorphic, since
coordinate translations have identity derivative and
holomorphic maps have holomorphic derivatives by the
Cauchy formula of Lemma~\ref{lem:analytic-tools}.
Each entry is constant by compactness and the maximum
principle. On the vector source cover subtract from the
map the target period map applied to this constant linear
derivative. The difference has zero derivative and is
zero at zero, hence vanishes. This proves the assertion
also when the two group structures on the same underlying
pointed compact curve were given separately, by their
uniformizations in
Lemma~\ref{lem:projective-group-period-uniformization}.

An isomorphism matching the group and labels preserves
the normalized functions by Lemma~\ref{lem:rigid-two-three-cubic}.
Conversely equal tuples identify the same cubic and all
the labelled points, including the origin. Under this
pointed identification the group laws agree by the first
paragraph. Their endomorphisms for \(a\in O\) agree since
both have derivative multiplication by the fixed complex
number \(a\), and an endomorphism is determined by its
derivative by the same argument. This proves separation
by Proposition~\ref{prop:cm-elliptic-ray-classification}.

Conjugate a normalized cubic, its algebraic group law,
endomorphisms and marked points by any
\(\iota\in{\rm Aut}(\C/K)\). Algebraicity of the group
and maps before conjugation is
Lemma~\ref{lem:elliptic-cubic-period-algebraic}. The result
is a smooth irreducible projective group with the same
tangent scalars \(a\in O\), since the derivative on a
one-dimensional tangent space is a scalar and its linear
chart terms conjugate. All the finite torsion labels are
transported bijectively, since multiplication and its
point equations are transported. Thus this is a curve
from the same finite class list by
Proposition~\ref{prop:cm-elliptic-ray-classification}.
The normalized coordinates conjugate by their uniqueness
or directly by the conjugated cubic equation and values.
Each entry therefore has finite orbit under these
automorphisms. The monic polynomial with the distinct
orbit elements as roots has coefficients fixed by all of
them and hence in \(K\), by the fixed-field computation
of Lemma~\ref{lem:fixed-projective-equations}. The finite
set of tuples therefore lies in one finite Galois
extension, taking the splitting fields of the finitely
many orbit polynomials. This proves the assertion on
permutations.

Finally an automorphism fixing the algebraic closure
of \(K\) fixes each normalized cubic and its origin as
embedded data. The conjugate group law must be its same
law by the first paragraph. The conjugate endomorphisms
by \(a\) have the original tangent scalars and hence are
the same maps. Apply
Lemma~\ref{lem:equivariant-projective-map-descent} over
that algebraic closure to their regular maps, also to
the sum on the projective product. This proves the
assertion on fields of definition.
\end{proof}

\section{Conjugation and translation on the quadratic levels}

\begin{lemma}\label{lem:compact-quadratic-full-idele-classes}
Give \(\mathcal C=\mathbb A_{K,f}^\times/K^\times\) the
quotient of the restricted product topology. It is compact,
and reduction identifies it as a topological group with
the inverse limit of the finite groups \(\mathcal C_m\).
The multiples of six suffice as levels.
\end{lemma}
\begin{proof}
Projection to \(\mathcal C_1\), the finite ideal class
group, has kernel \(\widehat O^\times/O^\times\), by the
valuation map of Proposition~\ref{prop:central-character-actual-form}.
The subgroup \(K^\times\cap\widehat O^\times=O^\times\)
is finite: an integral quadratic unit has complex
absolute value one by its norm and there are finitely
many points of the lattice \(O\) in a bounded disk.
The compactness of \(\widehat O^\times\) can be seen
from its finite residues modulo \(j!\). Compatible
residues give unique local integers and inverses by
the local completeness and residue description of
Lemma~\ref{lem:local-residue-indicators}. The kernels
of the residue maps are a neighborhood basis in this
product, since finitely many prescribed local powers
divide some \(j!\). If an open cover had no finite
subcover, one could choose recursively a nested sequence
of factorial residue classes with the same failure,
using the finite subdivision at the next stage. The
compatible residues would define a point with a whole
neighborhood in one covering open, a contradiction.
The quotient of the unit group and each of its finitely
many translates are compact, giving compactness of
\(\mathcal C\).

The maps to finite levels are continuous and surjective.
A class trivial at every level has a representative
\(b\in\widehat O^\times\). For each \(j\) it equals
some unit in \(O^\times\) modulo \(H_{j!}\). Finiteness
of the units makes a single one occur for unbounded
\(j\); equality at all powers gives equality in every
completion. Thus the limit map is injective. For a
compatible sequence of finite classes the nonempty
preimages of its factorial terms are nested closed
sets in a compact space, so their intersection supplies
a representative. The finite level residues of that
representative give the whole sequence. The limit is
Hausdorff since distinct sequences differ in some
finite discrete coordinate; the continuous bijection
from the compact group has continuous inverse as closed
subsets are compact and compact subsets of Hausdorff
spaces are closed. Factorials divisible by six are
cofinal, proving the final assertion.
\end{proof}

\begin{proposition}\label{prop:quadratic-cm-galois-translation}
Let \(\bar K\) be the algebraic closure of \(K\) in
\(\C\) and \(G_K={\rm Gal}(\bar K/K)\), with the finite
Galois quotient topology. There exists a continuous
homomorphism
\[
                         r_K:G_K\longrightarrow\mathcal C
\]
such that for each \(6\mid m\) a Galois element takes the
rigid tuple of \([x]\in\mathcal C_m\) to the rigid tuple
of \(r_K(\gamma)[x]\) at that level.
\end{proposition}
\begin{proof}
Lift \(\gamma\) to an automorphism of \(\C\) by choosing a
transcendence basis and extending over the algebraic closure
as in Lemma~\ref{lem:fixed-projective-equations}. Conjugating
all the \(O\)-action and torsion labels of any \(Y_x\) gives
a full labelled curve of the same tangent type, by
Proposition~\ref{prop:cm-elliptic-ray-classification}.
Write \(\widetilde{x}\) for the class of its reconstructed
frame in \(\mathcal C\). At each level \(6\mid m\) its
class is exactly the conjugated tuple of
Theorem~\ref{thm:quadratic-labelled-cubic-tuples-algebraic},
so it depends only on \(\gamma\). By
Lemma~\ref{lem:compact-quadratic-full-idele-classes}
this holds for the full class as well.

Let \(y\) be any finite idele of \(K\). Multiplying its
inverse by a suitable positive integer clears all its
denominators. Dividing \(y\) by that integer leaves
its class in \(\mathcal C\) unchanged, so we may assume
\(y^{-1}\in\widehat O\). Then
\(\mathfrak a_x\subseteq\mathfrak a_{xy}\); the vector
identity gives a nonzero algebraic homomorphism
\(Y_x\to Y_{xy}\) by
Lemma~\ref{lem:elliptic-cubic-period-algebraic}.
Its effect on labels is \(t\mapsto y^{-1}t\) by
Definition~\ref{def:cm-elliptic-finite-frames}.
Conjugating both curves and this map gives a nonzero
group homomorphism between the reconstructed curves
with exactly this same label formula. Its derivative
in the two rational period coordinates is multiplication
by some \(a\in K^\times\): it carries one full fractional
ideal into the other and a nonzero complex homomorphism
cannot have zero derivative by
Lemma~\ref{lem:projective-group-period-uniformization}.
Use arbitrary representatives of the reconstructed
classes. Then \(\widetilde{xy}^{-1}a\widetilde{x}-y^{-1}\)
maps \(l^{-1}\widehat O\) into \(\widehat O\) for every
positive integer \(l\), since its action on all finite
torsion classes vanishes. Taking \(l=p^j\) above each
rational prime proves this adelic difference is zero.
Thus \([\widetilde{xy}]=[\widetilde{x}]\,[y]\) in
\(\mathcal C\), including the original \(y\) since
rational scalings identify the full labelled groups.

The action is therefore translation by its value on
the identity class; call this value \(r_K(\gamma)\).
Composition of conjugations on the finite tuples gives
composition of their translations, hence the homomorphism
identity in \(\mathcal C\). Its projection to each finite
level \(6\mid m\) is continuous because the tuple entries
are in the finite Galois extension of
Theorem~\ref{thm:quadratic-labelled-cubic-tuples-algebraic}.
The inverse limit topology of the preceding lemma gives
continuity.
\end{proof}

\section{Good reduction on one fixed quadratic level}

\begin{lemma}\label{lem:fixed-quadratic-good-cubic-list}
Fix \(6\mid m_0\) and the normalizing labels of
Lemma~\ref{lem:rigid-two-three-cubic}. Choose one normalized
cubic and its full ordered list of \(m_0\)-labels for each
element of \(\mathcal C_{m_0}\). There are a finite Galois
extension \(L/K\) and an integer \(N>0\), divisible by
\(6m_0\), with the following properties over \(O_L[1/N]\).
Each equation
\[
             bY^2Z=X(X-Z)(X-\lambda Z)
\]
defines a smooth projective commutative group of relative
dimension one with the prescribed origin, inversion and
\(O\)-action. Its geometric fibers are integral. The
affine coordinates of its nonzero marked points are
integral in this ring, and the marked tuples for distinct
classes have distinct reductions at every remaining prime.

For every valuation ring \(V\subset\C\) containing this
ring and every such cubic over \(V\), put \(B=\mathcal O(1)\).
In every geometric fiber, for every \(t\ge1\), the complete
section space of \(B^t\) is obtained by reducing a free
degree-\(t\) coordinate module of rank \(3t\). The section
space of the trivial line on a geometric fiber has dimension
one. The corresponding flatness, integrality and high-degree
section assertions of Theorem~\ref{thm:arc-free-high-sections}
hold for any finite product of the listed cubics over \(V\).
\end{lemma}
\begin{proof}
Theorem~\ref{thm:quadratic-labelled-cubic-tuples-algebraic}
gives algebraic coordinates, group laws and endomorphisms.
Choose the integral basis \(1,\beta\) of \(O\) from
Lemma~\ref{lem:quadratic-norm}. It suffices to retain the single
endomorphism \([\beta]\), its integral quadratic relation
and its additivity. All other scalars are integral
combinations of it and the identity. Include the origin,
inversion, addition, their group equations on products
of at most three copies and the finite label identities.
The finite affine presentations of these morphisms and
their identities in Proposition~\ref{prop:finite-presentation-spread}
require only finitely many algebraic coefficients.
Take one finite Galois extension of \(K\) containing
them and all the tuples. The denominator-clearing argument
in Theorem~\ref{thm:fixed-integral-level-model} gives
all the maps and identities over the indicated ring
after increasing \(N\). Invert also the finitely many
denominators of the tuples and the norms of the nonzero
elements \(b,\lambda,1-\lambda\) after those denominators
are cleared. For each pair of distinct tuples invert
the norm of one nonzero coordinate difference in the
same way. These norms are nonzero rational integers
after clearing denominators. This proves both assertions
on the marks. The inversion formula is \((X,Y)\mapsto
(X,-Y)\) since it is so over the complex numbers.

On the affine chart \(Z=1\) the equation has a unit
partial derivative at every geometric point. If its
\(Y\)-derivative vanishes there, then \(Y=0\) and
the three distinct roots give a nonzero \(X\)-derivative.
At \([0:1:0]\) the derivative in \(Z/Y\) is \(b\).
These charts give smoothness of relative dimension one.
Over the function field of the affine \(X\)-line the
cubic polynomial \(X(X-1)(X-\lambda)/b\) cannot
be a square, by its simple zero at \(X=0\).
Monic division in \(Y\) thus gives a domain on the
affine chart over every algebraically closed residue
field. Its homogeneous closure is the entire cubic:
the homogeneous polynomial has no factor \(Z\), and a
homogeneous factor becoming constant at \(Z=1\)
would be a power of \(Z\) times a scalar. This gives
geometric integrality.

For \(P=V[X,Y,Z]\) the homogeneous equation is monic
up to sign in \(X\) of degree three. Division gives a
free basis of its degree-\(t\) quotient consisting of
monomials with \(X\)-exponent at most two. In positive
degree their number is \(3t\), by summing the counts
at exponents zero, one and two that occur. For the
section comparison use the graded sequence
\(0\to P(-3)\to P\to R\to0\), with \(R\) the cubic
ring. It stays exact after every scalar extension by
the same monic division. In the augmented chart complex
of Definition~\ref{def:graded-chart-cochains}, the
monomial contraction in Lemma~\ref{lem:high-degree-chart-section-comparison}
contracts each monomial unless all three exponents are
negative. The latter monomials occur only in cochain
degree two, since all three charts must be used. Thus
both augmented complexes for \(P_t\) and \(P(-3)_t\)
are exact for \(t\ge1\), since \(t-3\ge-2\).
For \(t=0\) the latter complex still has zero cohomology
in degrees zero and one. The lifting argument for the
short exact sequence in the same lemma identifies
\(R_t\) with the compatible sections for every \(t\ge0\).
In degree zero \(R_0=V\). All computations apply over
each geometric residue field. The direct-limit description
of the affine charts and the tensor-product argument
in Theorem~\ref{thm:arc-free-high-sections} now apply
with these free coordinate pieces. They give flatness,
complete product sections and \(H^0(\mathcal O)=k\)
on products over an algebraically closed field \(k\).
The product Jacobian charts and
Lemma~\ref{lem:smooth-connected-fiber-integral} give
their geometric integrality.
\end{proof}

\begin{lemma}\label{lem:good-cubic-prime-to-p-labels}
Let \(V\) in Lemma~\ref{lem:fixed-quadratic-good-cubic-list}
have fraction field \(\C\) and residue characteristic
\(p>0\). On each listed cubic, specialization of complex
torsion points of order prime to \(p\) is injective.
\end{lemma}
\begin{proof}
Normalize a finite list of homogeneous coordinates of
a complex point by dividing by one of smallest valuation.
They give a \(V\)-point of the same cubic, unique since
two such normalized tuples for the same point differ by
a unit. Group operations commute with this specialization
by Lemma~\ref{lem:fixed-quadratic-good-cubic-list}.
It suffices to show that a point \(Q\) reducing to the
origin and killed by \(n\) prime to \(p\) is the origin.

On the chart at the origin let \(t=X/Y,\ v=Z/Y\).
The equality
\(v(b+(1+\lambda)t^2-\lambda tv)=t^3\)
from Lemma~\ref{lem:rigid-two-three-cubic}
makes the ideal of the identity section \((t)\) on an
open neighborhood of the closed identity. Its conormal
line, the identity ideal modulo its square, has basis
\(t\) over \(V\): reducing the displayed relation to
first order sets \(v=0\) and imposes no linear relation
on \(t\). On that line the pullback by multiplication
by \(n\) is \(n\). Indeed at the pair of origins the
two first-order coefficients of the sum map are both
one by its restriction to each axis; compose with the
diagonal and induct on \(n\). Thus, on a neighborhood
of the closed identity on which the source and image
coordinates are defined,
\([n]^*t=nt+t^2h\) with \(h\) regular there. The
section \(Q\) lies in this neighborhood, since its
inverse image is an open containing the closed point
of the local spectrum. If \([n]Q=0\), then
\(t(Q)(n+t(Q)h(Q))=0\) in \(V\); the parenthesis is
a unit since \(t(Q)\) belongs to the maximal ideal.
Consequently \(t(Q)=0\), and the cubic equation gives
\(v(Q)=0\). This is precisely \([0:1:0]\).
\end{proof}

\section{Integral isogenies of the quadratic cubics}

\begin{proposition}\label{prop:good-cubic-extension-degree}
Use any two cubic models \(G,G'\) and a valuation ring \(V\)
with fraction field \(\C\) as in
Lemma~\ref{lem:fixed-quadratic-good-cubic-list}.
Every algebraic group homomorphism of their generic
fibers extends uniquely over \(V\). Suppose the generic
map \(f\) is induced by a nonzero complex linear map of
the period lines and its finite kernel has odd cardinality
\(d\). Then its extension is surjective finite locally free
of rank \(d\), and its kernel is free of length \(d\) over
\(V\), meaning that its coordinate algebra is a free
\(V\)-module of that rank.
\end{proposition}
\begin{proof}
For the extension apply the argument of
Theorem~\ref{thm:arc-extension-by-translation} with the
free high cubic coordinate pieces and integral fibers
of Lemma~\ref{lem:fixed-quadratic-good-cubic-list}.
In detail, clear denominators in a homogeneous tuple
for the generic map, normalize its coefficients in a
high free piece by one of least valuation, and take the
open where the reduced tuple is defined. This open is
dense in each geometric fiber because at least one
coefficient is a unit and the fiber comparison is
injective there. The sum of this open with itself covers
the group fiberwise, since two dense opens on the same
integral geometric fiber meet. The flat surjection from
pairs of points of this open by their sum therefore
descends the sum of their images, by
Lemma~\ref{lem:flat-cover-projective-descent}.
On its overlap the two formulas agree generically and
hence everywhere by the flat chart and product calculation
in Theorem~\ref{thm:arc-extension-by-translation}.
The same calculation gives uniqueness and the homomorphism
identity. In particular it extends equalities of composites.

Write \(O_0\) for the identity point and \(B,B'\) for the
respective hyperplane lines. The line at infinity cuts
out \(3[O_0]\), by the equation for \(Z/Y\) in
Lemma~\ref{lem:rigid-two-three-cubic}. The generic map is
locally a biholomorphism since its linear coefficient is
nonzero. Its inverse image of the target origin is thus
the reduced set \(H=\ker f\), with multiplicity one at
each point. Every nonzero element of \(H\) belongs to a
two-element pair \(\{Q,-Q\}\) since \(d\) is odd. For such
a pair the function \(X-X(Q)\) has a simple zero at both
points and a double pole at \(O_0\), with no further zeros
or poles. Indeed \(Y(Q)\ne0\) so \(X-X(Q)\) is a parameter
there by the cubic equation, and the two points of that
\(X\)-coordinate have opposite \(Y\)-coordinates.
The product over the pairs has divisor
\(\sum_{Q\in H}[Q]-d[O_0]\), also for \(d=1\) by the
empty product. Cubing this identity gives the generic
line isomorphism \(f^* B'\simeq B^d\).
Apply the proof of Theorem~\ref{thm:arc-picard-injection}
to the quotient of the two lines. The free high pieces,
their comparison on the special fiber and flatness are
those of Lemma~\ref{lem:fixed-quadratic-good-cubic-list}.
Normalizing the generic trivialization and its inverse
gives regular sections whose product is a unit, by
integrality of the closed fiber as in that theorem.
This proves the pullback identity over \(V\).

The ample pullback gives finiteness by the affine chart
argument of Lemma~\ref{lem:arc-projective-polarized-finite}.
On each geometric fiber the image is the entire target.
Indeed its function field and that of the source both
have transcendence degree one, directly from the monic
quadratic equation over \(k(X)\) on the affine chart.
For a finite map the image and source have equal
transcendence degrees and a proper closed subset has
smaller transcendence degree, by
Lemma~\ref{lem:good-fiber-function-field-degree}.
On an integral target fiber the direct image of the
structure sheaf is then locally free by translating
a generic basis, as in Theorem~\ref{thm:arc-polarized-isogeny-degree}.
If its rank is \(r\), sandwich it between sums of \(r\)
negative and positive twists by \(B'\) as in the same
proof. By Lemma~\ref{lem:fixed-quadratic-good-cubic-list}
the two outside bounds on the dimension of the sections
after a twist by \(t\) are \(3r(t-c)\) and \(3r(t+c)\)
for some fixed \(c\) on this fiber and all \(t>c\).
The pullback has \(3dt\) sections, so \(r=d\).
The final flat-chart argument of
Theorem~\ref{thm:arc-polarized-isogeny-degree} now gives
local freeness on the full target: a lifted basis at
any point gives a surjection from the free module of
rank \(d\) locally; the map is an isomorphism on the
generic base fiber, and the source module has no
\(V\)-torsion. Pull back at the identity to obtain
the asserted free kernel algebra.
\end{proof}

\begin{corollary}\label{cor:split-quadratic-ideal-isogeny}
For a split rational prime \(p\nmid N\) write
\(pO=\mathfrak u\overline{\mathfrak u}\), and put
\(t_{\mathfrak u}=p\) at \(\mathfrak u\) and \(1\) at
every other finite place, as an idele of \(K\).
For each finite idele \(x\), the identity on period
lines gives a degree-\(p\) isogeny
\[
        \psi_{\mathfrak u}:Y_x\longrightarrow
                 Y_{x t_{\mathfrak u}^{-1}} .
\]
On normalized models from the fixed level \(m_0\) it
extends surjectively and finite locally freely of rank
\(p\) over each good \(V\). It carries the label \(s\)
to the label \(t_{\mathfrak u}s\); in particular it
preserves every label of order prime to \(p\).
\end{corollary}
\begin{proof}
The corresponding period lattices are \(\mathfrak a_x\)
and \(\mathfrak u^{-1}\mathfrak a_x\) by
Definition~\ref{def:cm-elliptic-finite-frames} and
Proposition~\ref{prop:general-ideal-factorization}.
Their inclusion has index \(p\), since at \(\mathfrak u\)
the quotient is its one-dimensional residue field of
size \(p\), and it is zero at every other completion,
by the ideal residue calculation in
Lemma~\ref{lem:local-residue-indicators}.
The map is algebraic by Lemma~\ref{lem:elliptic-cubic-period-algebraic}.
In the target coordinates of
Definition~\ref{def:cm-elliptic-finite-frames} its
effect on the label is multiplication by \(t_{\mathfrak u}\).
A label of order prime to \(p\) has zero component at
\(\mathfrak u\), proving the last assertion. Identifying
each complex group with the model for its \(m_0\)-tuple
by Theorem~\ref{thm:quadratic-labelled-cubic-tuples-algebraic},
the extension and rank now follow from
Proposition~\ref{prop:good-cubic-extension-degree},
since \(p\) is odd.
\end{proof}

\section{Frobenius on the split quadratic isogeny}

\begin{theorem}\label{thm:quadratic-split-frobenius-kernel}
In Corollary~\ref{cor:split-quadratic-ideal-isogeny}
assume that the maximal ideal of \(V\) contracts to
\(\mathfrak u\) in \(O\). Write \(G_k,G'_k\) for the
reductions of the source and target over an algebraic
closure \(k\) of the residue field. There is a unique
isomorphism of pointed groups
\[
            i:G_k^{(p)}\longrightarrow G'_k,\qquad
                 (\psi_{\mathfrak u})_k=i F_{G_k,p}.
\]
Here \(F_{G_k,p}\) is the relative Frobenius. For every
label of order prime to \(p\) this isomorphism takes
the Frobenius of its source reduction to the target
reduction with the same label.
\end{theorem}
\begin{proof}
Put \(J=\ker F_{G_k,p}\). By
Lemma~\ref{lem:smooth-group-frobenius-coordinates},
it is killed by \([p]\), has one geometric point, and has
coordinate algebra \(k[t]/(t^p)\). Every scalar of \(O\)
preserves \(J\), since relative Frobenius commutes with
the corresponding twisted endomorphism by the coordinate
formula in Lemma~\ref{lem:frobenius-pulls-kernel-ideal}.
Choose \(a\in O\) with residues zero modulo \(\mathfrak u\)
and one modulo \(\overline{\mathfrak u}\). Such \(a\)
exists since these distinct maximal ideals have sum
\(O\): write one as a sum of elements of the two ideals
and take the first summand. Their intersection is their
product \(pO\), since any member of the intersection
times this sum belongs to the product. Thus
\(a^2-a\in pO\), so \([a]|_J\) is idempotent under
composition.

The tangent scalar of \([a]\) on \(G_k\) is the image
of \(a\) in \(k\), hence is zero. Indeed the tangent
line at the origin over \(V\) has a basis dual to the
parameter of Lemma~\ref{lem:good-cubic-prime-to-p-labels}.
In that basis the derivative is a scalar of \(V\);
over the fraction field it equals \(a\) by
Proposition~\ref{prop:cm-elliptic-ray-classification},
so its reduction has the asserted value. On the
augmentation ideal \(I=(t)\) of \(J\) the pullback
by \([a]\) therefore maps \(I\) into \(I^2\).
Indeed \(t\) is the image of a tangent parameter of
\(G_k\). The \(j\)-fold iterate maps \(I\) into
\(I^{2^j}\), which vanishes when \(2^j\ge p\).
Idempotence identifies that iterate with the first
map. Thus \([a]|_J\) is the zero group map and
\([1-a]|_J\) is the identity.

Put \(e=1-a\). The linear map on the period line
given by \(e/p\) carries \(\mathfrak a_x\) into
\(\mathfrak u^{-1}\mathfrak a_x\), since
\(e\in\overline{\mathfrak u}=p\mathfrak u^{-1}\).
It gives an algebraic homomorphism \(\phi:Y_x\to
Y_{x t_{\mathfrak u}^{-1}}\) by
Lemma~\ref{lem:elliptic-cubic-period-algebraic}.
On the generic groups
\(\psi_{\mathfrak u}[e]=\phi[p]\) by the equality of
the two linear maps. By Proposition~\ref{prop:good-cubic-extension-degree}
these homomorphisms extend with their equality over
\(V\). Restrict the equality to \(J\) on the special
fiber. Its left side is \((\psi_{\mathfrak u})_k|_J\)
and its right side is zero. Hence \(J\) is a closed
subgroup of the kernel of \((\psi_{\mathfrak u})_k\).
Both have length \(p\), by the first paragraph and
Corollary~\ref{cor:split-quadratic-ideal-isogeny}.
The induced surjection of their coordinate algebras is
therefore an isomorphism. Apply
Corollary~\ref{cor:equal-kernels-frobenius-factor}.
The final assertion follows from the label formula in
Corollary~\ref{cor:split-quadratic-ideal-isogeny} and
specialization of its extended map.
\end{proof}

\begin{corollary}\label{cor:quadratic-frobenius-rigid-coordinates}
Under Theorem~\ref{thm:quadratic-split-frobenius-kernel},
write \((b_x,\lambda_x)\) for the two coefficients of
the normalized model, and \(Q_x(s)\) for the pair of
affine coordinates of its point of nonzero label \(s\).
For every \(s\) of order prime to \(p\) these coordinates
are integral at \(V\). If a tilde denotes reduction and
\(y=xt_{\mathfrak u}^{-1}\), then
\[
  (\widetilde b_y,\widetilde\lambda_y)
       =(\widetilde b_x^{\,p},\widetilde\lambda_x^{\,p}),
      \qquad
        \widetilde Q_y(s)=\widetilde Q_x(s)^{\,p},
\]
with entrywise powering on the right.
\end{corollary}
\begin{proof}
By Lemma~\ref{lem:good-cubic-prime-to-p-labels} the point
of any such label reduces away from the origin. It thus
lies in the affine chart over \(V\), since its reduction
does, proving integrality. On \(G_k^{(p)}\) the cubic has
coefficients \(\widetilde b_x^{\,p},\widetilde\lambda_x^{\,p}\),
and the Frobenius images of the chosen normalizing
points have the same values zero, one and \(Y=1\)
at their respective labels. All of those labels have
order prime to \(p\). The isomorphism of
Theorem~\ref{thm:quadratic-split-frobenius-kernel}
preserves them, the origin and group inversion. By the
uniqueness in Lemma~\ref{lem:rigid-two-three-cubic}
the pullbacks of the target functions \(X,Y\) are
the coordinate functions on that twisted cubic.
Evaluating \(X\) at the third nonzero two-torsion
point gives the coefficient \(\lambda\), and
evaluating the cubic relation at the chosen order-three
point gives \(b\). These evaluations prove the first
equality. The coordinate formula for relative Frobenius
and the assertion on all prime-to-\(p\) labels in the
theorem prove the second one.
\end{proof}

\section{The quadratic Galois action away from the residue prime}

\begin{lemma}\label{lem:residue-map-to-valuation}
Let \(h:O_{\overline K}\to\overline{\mathbf F}_p\) be a
residue map at a place of \(K\) above \(p\), where
\(\overline K\) is its algebraic closure in \(\C\).
There is a valuation ring \(V\) of \(\C\) such that
reduction on the elements of \(O_{\overline K}\) is
\(h\) followed by an embedding of its residue image.
If \(g\in G_K\) satisfies \(h(gz)=h(z)^b\) on
\(O_{\overline K}\), for \(b=1\) or \(p\), this identity
holds on \(\overline K\cap V\) with its extended residue
map, and \(g\) preserves \(\overline K\cap V\).
\end{lemma}
\begin{proof}
Localize \(O_{\overline K}\) at \(\ker h\). The resulting
ring \(A\) is the directed union of the localizations
of \(O_D\) at the contracted primes, as \(D\) runs through
the finite subextensions of \(\overline K/K\). Indeed
an integral numerator and denominator lie together in
one such \(D\), and a denominator outside the kernel
remains outside it on every larger level. Each finite
localization is a valuation ring of \(D\), by the local
ideal factorization in Proposition~\ref{prop:general-ideal-factorization}.
Thus \(A\) is a valuation ring of \(\overline K\),
by testing an element and its inverse in one finite
field. Extend to a local ring \(V\) of \(\C\) dominating
\(A\) by the maximal-local-subring argument of
Lemma~\ref{lem:horizontal-graph-valuative-centers}.
That argument gives the valuation property in \(\C\).
The contraction is \(A\), since if \(x\in\overline K\)
is outside \(A\), then \(x^{-1}\) is in the maximal
ideal of \(A\), hence in that of \(V\). Its residue
embedding agrees with \(h\) on \(O_{\overline K}\)
upon identifying the residue of \(A\) with the fraction
field of the image of \(h\). If \(h(gz)=h(z)^b\),
the prime is preserved in both directions since a
power is zero exactly at zero and \(g\) is invertible.
The formula extends to fractions whose denominator
is outside the prime. This proves the last assertions.
\end{proof}

\begin{theorem}\label{thm:quadratic-idele-geometric-frobenius}
For the imaginary quadratic field \(K\), take \(m_0=6\)
and \(N_K=N\) of Lemma~\ref{lem:fixed-quadratic-good-cubic-list}
at this level. Let \(p\nmid N_K\) split in \(K\), choose
a place \(\mathfrak u\mid p\), and let \(h\) extend its
residue map. Write \(U_p\subset\mathcal C\) for the image
of \(\prod_{v\mid p}O_v^\times\) supported over \(p\).
Let \(I_h\subset G_K\) consist of the automorphisms
\(s\) with \(h(sz)=h(z)\) for every algebraic integer
\(z\). For \(r_K\) of
Proposition~\ref{prop:quadratic-cm-galois-translation}
one has
\[
       r_K(I_h)\subseteq U_p,\qquad
       r_K(\gamma)\in[t_{\mathfrak u}]\,U_p
\]
for every geometric Frobenius datum
\((\mathfrak u,h,\gamma)\) of
Definition~\ref{def:frobenius-data}. These assertions
hold with the same \(N_K\) at all levels.
\end{theorem}
\begin{proof}
Use the valuation of Lemma~\ref{lem:residue-map-to-valuation};
it contains the localized coefficient ring of
Lemma~\ref{lem:fixed-quadratic-good-cubic-list}.
Reduce the normalized model \(Y_1\) and its labels
there. All labels of finite order have algebraic
coordinates by Theorem~\ref{thm:quadratic-labelled-cubic-tuples-algebraic}
applied to a multiple of their order and six, using the
same normalizing points. The affine coordinates of
every nonzero prime-to-\(p\) label are in the valuation
ring by Corollary~\ref{cor:quadratic-frobenius-rigid-coordinates}.
The map \(\gamma^{-1}\) acts on their residues and
on the model coefficients by \(p\)-th powering by
Definition~\ref{def:frobenius-data} and the preceding lemma.

Put \(y=t_{\mathfrak u}^{-1}\). By
Corollary~\ref{cor:quadratic-frobenius-rigid-coordinates}
the conjugate of the level-six tuple by \(\gamma^{-1}\)
and the level-six tuple of \(Y_y\) have identical
reductions. Both are members of the fixed finite list
by Theorem~\ref{thm:quadratic-labelled-cubic-tuples-algebraic}.
Lemma~\ref{lem:fixed-quadratic-good-cubic-list} separates
that list on reduction. Hence they have identical tuples
in characteristic zero and are the very same normalized
cubic with the same group law, as proved in
Theorem~\ref{thm:quadratic-labelled-cubic-tuples-algebraic}.
For every \(n\) prime to \(p\), the conjugate of each
\(n\)-label and the corresponding \(n\)-label of \(Y_y\)
are now points on that one curve killed by \(n\).
Their reductions agree by the same Frobenius coordinate
formula. Lemma~\ref{lem:good-cubic-prime-to-p-labels}
identifies the points themselves. The zero label is
preserved by their group maps. At every \(n\) divisible
by six and prime to \(p\), the classification in
Proposition~\ref{prop:cm-elliptic-ray-classification}
therefore gives
\(r_K(\gamma^{-1})=[t_{\mathfrak u}^{-1}]\) in
\(\mathcal C_n\). Inverting gives the positive uniformizer
for geometric Frobenius. For \(s\in I_h\) compare the
\(s\)-conjugated tuple with \(Y_1\) itself. The residues
are unchanged, and the identical fixed-list and torsion
injection arguments give \(r_K(s)=1\) in every such
\(\mathcal C_n\).

It remains to identify the common kernel of these finite
projections. Its projection to the ideal class quotient
\(\mathcal C_1\) is trivial, so represent any member
by \(u\in\widehat O^\times\), by the calculation in
Proposition~\ref{prop:central-character-actual-form}.
At each indicated \(n\) triviality means \(u\in O^\times H_n\),
since the principal factor relating two everywhere
integral units must be a global unit. Choose increasingly
divisible \(n_j\) cofinal among those integers, for instance
six times the product of all positive integers up to \(j\)
prime to \(p\). The global unit group is finite by
Lemma~\ref{lem:compact-quadratic-full-idele-classes}.
One unit \(c\) therefore occurs for an unbounded subsequence
in \(u\in cH_{n_j}\). At each finite \(v\nmid p\)
the valuations of these moduli tend to infinity, giving
\(u_v=c\). Thus \(c^{-1}u\) is supported above \(p\).
Conversely every such integral unit lies in all these
kernels by Definition~\ref{def:cm-elliptic-finite-frames}.
This proves both assertions in \(\mathcal C\). A unit
supported above \(p\) representing any given class is
unique, since a principal idele equal to one at another
finite place is the identity.
\end{proof}

\section{The scalar character of the approximating form}

\begin{proposition}\label{prop:corrected-quadratic-idele-character}
Let \(\omega\) be a character from
Proposition~\ref{prop:central-character-actual-form}, and
choose integers \(A,B\) with
\(\omega(z)=z^A\bar z^B\) on \(K^\times\).
Fix the auxiliary prime \(\ell\) and the isomorphism
\(\iota:\C\to K_\ell\) of
Lemma~\ref{lem:auxiliary-split-embeddings}. Denote by
\(u_+,u_-\mid\ell\) the places of \(K\) induced by
\(z\mapsto\iota(z),\ z\mapsto\iota(\bar z)\), respectively,
using the corresponding identifications of their
completions with \(\Q_\ell\). The formula
\[
       \omega_\ell([x])=
       \frac{\iota(\omega(x))}{x_{u_+}^{A}x_{u_-}^{B}}
            \qquad(x\in\mathbb A_{K,f}^\times)
\]
defines a continuous character of \(\mathcal C\) taking
values in the unit group of one finite extension of
\(\Q_\ell\). Choose \(m\ge1\) with \(H_m\subseteq\ker\omega\).
For every \(p\nmid m\ell\) it is trivial on \(U_p\).
\end{proposition}
\begin{proof}
The two identifications with \(\Q_\ell\) and their
distinct induced primes follow from
Lemma~\ref{lem:auxiliary-split-embeddings}. The numerator
is a locally constant character on finite ideles, since
the kernel of \(\omega\) is open. By the proof of
Proposition~\ref{prop:central-character-actual-form}
there is such an \(m\); select representatives \(x_i\)
of the finite quotient \(\mathcal C_m\).
Writing \(x_i^{h_i}=z_i v_i\) for its class order
\(h_i\), where \(z_i\in K^\times,\ v_i\in H_m\), gives
\(\omega(x_i)^{h_i}=z_i^A\bar z_i^B\in K^\times\).
Thus all values of \(\omega\) are in a single number
field \(D=K((\omega(x_i))_i)\): each idele is a
representative times a principal idele and a member
of \(H_m\). The numerator now takes values in
\(D_\ell=\Q_\ell(\iota(D))\), which is finite over
\(\Q_\ell\) by the finite algebraic tower argument in
Lemma~\ref{lem:place-fields}. The denominator is a
continuous multiplicative function into \(\Q_\ell^\times\),
and on a principal idele the numerator equals that
denominator by the choice of the two local embeddings.
The displayed quotient therefore descends continuously
for the quotient topology. Its image is a compact
subgroup by Lemma~\ref{lem:compact-quadratic-full-idele-classes}.
The absolute values on a compact set are bounded,
and both positive and negative powers of any element
in this subgroup still belong to it. Each element
therefore has absolute value one. Finally, for units
supported above \(p\nmid m\ell\) the numerator is one
by membership in \(H_m\), and both local components
in the denominator are one.
\end{proof}

\begin{theorem}\label{thm:actual-scalar-hecke-galois-character}
For the form of Theorem~\ref{thm:actual-central-test-vector}
put
\[
           a_\omega(g)=\omega_\ell(r_K(g)),\qquad
                   g\in G_E\subseteq G_K .
\]
Use the common algebraic closure of the indicated fields
in \(\C\). This is a continuous character into
\(K_\ell^\times\). Outside a finite set of rational
primes, at every degree-one place \(w\) of \(E\) and
every geometric Frobenius datum \((w,h,\gamma)\),
\[
                a_\omega(\gamma)=\iota(\omega(z_w)),
       \qquad     a_\omega(I_h\cap G_E)=\{1\},
\]
where \(z_w\) is the uniformizer norm idele of
Theorem~\ref{thm:actual-central-test-vector}.
In particular this supplies the scalar character required
in Corollary~\ref{cor:twisted-hecke-quotient-seed}
for each approximating form.
\end{theorem}
\begin{proof}
Continuity and multiplicativity follow from
Proposition~\ref{prop:quadratic-cm-galois-translation}
and Proposition~\ref{prop:corrected-quadratic-idele-character}.
Restriction to \(G_E\) is continuous: a prescribed
finite normal restriction over \(K\) pulls back to
a finite normal restriction after compositing with \(E\).
Exclude the finitely many primes from the chosen level
and safe conditions, \(N_K m\ell\), and the ramified
rational primes of \(E\) and \(K\). At each remaining
rational prime \(p\) under a degree-one \(w\), its
contraction \(\mathfrak u\) to \(K\) also has residue
field \(\mathbf F_p\). It therefore splits in the
quadratic field in the sense that there are two distinct
places over \(p\), by the degree formula in
Lemma~\ref{lem:finite-extension-residue-actions}.
Both local fields \(E_w\) and \(K_{\mathfrak u}\) are
\(\Q_p\): their ramification and residue degrees are
one by the same formula and the completion calculation
of Lemma~\ref{lem:definite-field-coordinates}. Thus
\(z_w\) and \(t_{\mathfrak u}\) differ only by a
unit at \(\mathfrak u\). This factor belongs to
\(H_m\), since \(p\nmid m\).

A datum over \(E\) is a datum over \(K\) at
\(\mathfrak u\) under the inclusion, since the exponent
in Definition~\ref{def:frobenius-data} is \(p\) for
both and the algebraic integers in the closure are
the same. By Theorem~\ref{thm:quadratic-idele-geometric-frobenius},
\(r_K(\gamma)\) differs from \([t_{\mathfrak u}]\)
by an element of \(U_p\), which
Proposition~\ref{prop:corrected-quadratic-idele-character}
kills. Both \(\ell\)-components of \(t_{\mathfrak u}\)
are one, so the value is \(\iota(\omega(t_{\mathfrak u}))
=\iota(\omega(z_w))\). The same two results give the
triviality on \(I_h\cap G_E\). The finite exclusions
satisfy the quantifier for most data in
Definition~\ref{def:frobenius-approximation}.
Consequently the exact scalar in
Theorem~\ref{thm:rank-before-scalar-twist} is available
for the indicated spherical polynomial.
\end{proof}

\section{The order of choice of the coarse level and the prime}

\begin{proposition}\label{prop:coarse-model-before-coefficient-prime}
Fix \(E\) and a finite list of definite forms as in
Theorem~\ref{thm:definite-wild-approximants}. For each
of their ranks \(n\) and each distinguished embedding
\(\sigma:E\to\C\), the polarization data of
Theorem~\ref{thm:actual-pel-torsion-pairing} and an integer
\(b_\sigma>0\) can be chosen before the coefficient
prime \(\ell\); the dependence of \(b_\sigma\) on \(n\)
is understood. Put \(C_b=K(b_\sigma)\) as in
Definition~\ref{def:moduli-signed-frame}. All stabilizers
at this level are trivial. One may choose \(\ell,w_0,w_1\)
as in Theorem~\ref{thm:definite-wild-approximants} so that
at every \(w\mid\ell\) each coarse complex base
\(\mathcal X(C_b,0)\) and its one-point group have
the smooth projective models of
Theorem~\ref{thm:fixed-integral-level-model}.
Their equations, group maps and \(O_E\)-action are
defined over a localization of \(O_E\) good at every
such \(w\). Over the completion \(O_{E_w}=\Z_\ell\)
these models admit their smooth Jacobian charts.
Every finite marked power of the coarse group is
smooth projective over \(\Z_\ell\) as well.

For each resulting approximant, the element
\(\mathfrak n\) of Corollary~\ref{cor:actual-moduli-places}
may be chosen in \(b_\sigma O_E\), with \(w_0\) the
only prime dividing \(\mathfrak n\) above \(\ell\).
The auxiliary place \(w_*\) there can be chosen over
a rational prime coprime to \(b_\sigma\ell\).
\end{proposition}
\begin{proof}
There are finitely many ranks and \(2[E^+:\Q]\)
possible embeddings at each rank. For each pair
apply Lemma~\ref{lem:actual-signs-and-padding},
Theorem~\ref{thm:actual-pel-torsion-pairing} and
Proposition~\ref{prop:actual-integral-moduli-level}
with a fixed finite set of primes in the padding lemma,
for example the empty set. Call the resulting multiple
for the level \(N_\sigma\). For an odd rational prime
\(q_b\) split completely in \(E\) outside this finite
list of integer divisors, put \(b_\sigma=q_bN_\sigma\).
Such primes exist by
Lemma~\ref{lem:auxiliary-split-embeddings}. The
principal level with this integer is admissible by
Proposition~\ref{prop:actual-integral-moduli-level};
Proposition~\ref{prop:neat-at-wild-place}, applied at
a place above \(q_b\), gives trivial stabilizers.
All these choices concern only the ranks, the
distinguished embeddings and the arithmetic field.

For each coarse level take the base, one-point family,
projections, group operations and the actions of
integral generators of \(O_E\), including their ring
relations. One may use an integral basis, with its
finite table of sums of basis multiples for the
products, by Theorem~\ref{thm:integer-ring}.
Theorem~\ref{thm:fixed-integral-level-model} gives
models over \(O_E\) with finitely many rational
primes inverted. Its proof gives the charts for
relative smoothness over a finite number field
extension with finitely many more primes inverted.
Keep all these fields and primes for the finite list,
including the exclusions in
Theorem~\ref{thm:uniform-good-theta-section-dimensions}
and Lemma~\ref{lem:degree-one-split-lie-polynomial}.
The marked \(\nu\)-point family on the coarse base
is the \(\nu\)-fold fiber product of the one-point
family for all \(\nu\), by the period coordinates
in Theorem~\ref{thm:analytic-framed-torus-family}
and the descent of marked slots in
Theorem~\ref{thm:hecke-diagram-over-number-field}.
The products are projective by the product embedding.
Their relative Jacobian charts use disjoint lists
of fiber variables with block diagonal relative
minors. Thus \(\nu\) adds no excluded primes.

Choose a finite normal extension \(D/\Q\) containing
\(E\) and all fields needed for the charts. By
Theorem~\ref{thm:actual-finite-chebotarev}, applied to
the identity in its Galois group, one can choose
\(\ell\) split completely in \(D\), outside the
finite exceptions above, the primes dividing all
\(b_\sigma\), and those excluded for the original
forms in Theorem~\ref{thm:definite-wild-approximants}.
One can include any further fixed finite list of
rational primes. Complete splitting outside the
power-basis indices gives completions \(\Q_\ell\)
by the root-lifting proof of
Lemma~\ref{lem:auxiliary-split-embeddings}, applied
also to \(D\). Any embedding of \(E\) at \(w\mid\ell\)
extends to one of \(D\) by Lemma~\ref{lem:place-fields}.
The necessary integral chart coefficients and their
inverted denominators map to \(\Z_\ell\) with the
denominators units. The retained presentations hold
there. Composing a relative smooth chart with a
smooth base chart multiplies their unit minors in
a block triangular matrix. This proves the assertion
on smoothness over \(\Z_\ell\).

The proof of Theorem~\ref{thm:definite-wild-approximants}
uses, after the exclusions for the given forms, only
the split completions and the two distinct indicated
places above the same quadratic prime. It applies
to this choice of \(\ell\). For any resulting approximant
choose \(w_*\) by Lemma~\ref{lem:auxiliary-split-embeddings}
away also from the finitely many divisors of
\(b_\sigma\) and the exceptional primes of the form.
Use the support and depths in the proof of
Corollary~\ref{cor:actual-moduli-places}, with
\(b_\sigma\) in place of \(N\) in
Lemma~\ref{lem:principal-level-support}.
Every resulting \(\mathfrak n\) is still divisible
by \(N_\sigma\). The extra congruence at \(q_b\)
only restricts the subgroup fixing the form, and
\(q_b\) differs from both \(\ell\) and the prime
under \(w_*\). The support identity of that lemma
gives exactly the stated support above \(\ell\).
All coprimalities used for the safe places in
Corollary~\ref{cor:actual-moduli-places} hold, since
the chosen primes avoid \(N_\sigma\).
\end{proof}

\section{The torsion fixed by the finer congruence level}

\begin{lemma}\label{lem:fine-level-invertible-lie-action}
Use the choice in Proposition~\ref{prop:coarse-model-before-coefficient-prime},
with the identification of \(E\) by its distinguished
embedding as in Lemma~\ref{lem:degree-one-split-lie-polynomial}.
Write \(S\) for the coarse base and \(G/S\) for its
one-point group after passage to \(O_{E_{w_1}}=\Z_\ell\).
One may impose finitely many further exclusions before
choosing \(\ell\), independent of the approximants.
For each \(\alpha\in O_E\setminus\{0\}\) divisible at
places above \(\ell\) only possibly by \(w_0\), the
relative differential of \([\alpha]\) is invertible on
the geometric fibers of \(G/S\) at every point. This
applies in particular to \(\alpha=\mathfrak n\) for
each of the fine levels in that proposition.
\end{lemma}
\begin{proof}
Fix the integral primitive element \(\theta\) of
Lemma~\ref{lem:degree-one-split-lie-polynomial}.
Include its characteristic polynomial on the Lie
space at the identity section in the finite data
spread at the coarse level. More explicitly, on
each relative Jacobian chart containing part of the
identity section the differentials of the free
parameters form a basis of the relative cotangent
there: differentiating the chart equations expresses
the remaining differentials in this basis by inverting
the chosen minor. Pullback by \([\theta]\) has a matrix
of regular functions in these frames on the base.
The characteristic polynomial of that matrix equals
\(P_\theta(T)\) of
Lemma~\ref{lem:degree-one-split-lie-polynomial} on every
complex point, since the transpose acts on the tangent
space. Its coefficient identities therefore hold on
the reduced generic base, by the affine point argument
of Lemma~\ref{lem:projective-special-section-base-change}.
They extend after a finite exclusion: an element of
the integral chart ring zero after passage to the
number field is killed by a nonzero integer after
clearing number denominators. Use the finitely many
chart identities and covers of
Theorem~\ref{thm:fixed-integral-level-model} for all
the coarse data in
Proposition~\ref{prop:coarse-model-before-coefficient-prime}.
This imposes the indicated exclusions before the
coefficient prime is chosen.

On any geometric fiber in residue characteristic
\(\ell\) the tangent representation of \(O_E\)
factors through \(O_E/\ell O_E\), since the
differential of multiplication by an integer is that
integer times the identity, by the first-order
group-law calculation of
Lemma~\ref{lem:good-cubic-prime-to-p-labels}, in each
free coordinate. The same calculation makes this
action additive, as well as multiplicative by the
chain rule. The split algebra here is a product of
\(2f\) copies of \(\mathbf F_\ell\); its orthogonal
idempotents split the tangent space into the corresponding
subspaces. By Lemma~\ref{lem:degree-one-split-lie-polynomial}
the residues of \(\theta\) at all those factors are
distinct. Their multiplicities in \(P_\theta(T)\)
are thus the dimensions of those subspaces, also on
the special fiber by the preceding coefficient identity.
In particular \(w_0\ne w_1\) lies over the same prime
of \(K\) as \(w_1\) and has multiplicity zero. On every
other factor the scalar \(\alpha\) is a unit at its
place. This proves invertibility at the identity.
In characteristic zero, \(\alpha\) divides the
nonzero integer \(N_{E/\Q}(\alpha)\) in \(O_E\),
by its monic characteristic equation. The corresponding
two endomorphisms multiply on tangent spaces to this
integer times the identity, proving invertibility
there also. Translations in a geometric group identify
the differential at any point with that at the identity,
since \([\alpha]\) intertwines translation by a point
with translation by its image.
\end{proof}

\begin{proposition}\label{prop:fine-level-torsion-etale-charts}
Under Lemma~\ref{lem:fine-level-invertible-lie-action}
let \(H_\alpha=\ker([\alpha]:G\to G)\), the closed
inverse image of the identity section. It is projective
over \(S\) and has relative Jacobian charts of dimension
zero over \(S\). Its geometric fibers are finite and
reduced. Each finite positive product of \(H_\alpha\)
over \(S\) has the same properties and admits smooth
Jacobian charts over \(\Z_\ell\).
\end{proposition}
\begin{proof}
The identity section is closed since the relative
diagonal of the projective group is closed. Thus
\(H_\alpha\) is closed in the projective \(G/S\).
Consider a relative smooth chart at the identity over
a point of \(S\). Denote the free coordinates by \(t\)
and the remaining coordinates by \(z\), so that
\(F(t,z)=0\) with invertible square \(z\)-Jacobian.
Shrink the base so that the identity section lies in
that chart and write its coordinates as \(t_0,z_0\).
Near the identity the ideal of this section is generated
by \(t-t_0\). Indeed in
\(F(t,z)-F(t_0,z_0)\) telescope the differences one
variable at a time. The coefficient matrix of \(z-z_0\)
at the identity is \(F_z(t_0,z_0)\), hence is invertible
on a neighborhood of it. The equations then express
\(z-z_0\) in the ideal of \(t-t_0\), and setting all
coordinates equal to their section values gives
precisely the section on its chart.

At a point of \(H_\alpha\) add the pullbacks of those
\(g=fn\) equations \([\alpha]^*(t-t_0)=0\) to a
relative smooth chart of the source \(G\). They give
exactly the kernel there by the preceding ideal
identity. Their differentials restrict on the
\(g\)-dimensional fiber tangent to independent forms
by Lemma~\ref{lem:fine-level-invertible-lie-action}.
Together with the independent old Jacobian rows
they therefore give an invertible square matrix in
the source variables over the residue field. Invert
its determinant locally, clearing the unit denominators
of the functions as in
Proposition~\ref{prop:complex-relative-jacobian}. This
is the asserted relative dimension-zero presentation.

On each geometric fiber the local maximal ideal at
a closed point is generated by coordinate differences.
Modulo its square the Jacobian relations just obtained
kill all of them. The finite-module determinant argument
of Lemma~\ref{lem:finite-fiber-local-algebras} in the
local ring therefore kills the maximal ideal itself.
The local ring is its residue field. Every prime in
an affine chart belongs to a maximal ideal there; the
prime correspondence on localization and the preceding
field calculation show that the prime was already
maximal. The noetherian fiber has only finitely many
irreducible components by
Proposition~\ref{prop:finite-strata-graded-partial-resolutions};
each is now a point. A nilpotent nonzero on any chart
would remain nonzero on localization at a maximal
ideal containing its annihilator. The field calculation
rules this out, proving reducedness. For a product take
the disjoint variable lists of the factor charts. The
relative Jacobian is block diagonal, and its composition
with the smooth base charts has a block triangular unit
minor as in Proposition~\ref{prop:coarse-model-before-coefficient-prime}.
The projectivity assertion for products follows from
the projective product embedding.
\end{proof}

\section{Components in a smooth chart over the local integers}

\begin{lemma}\label{lem:mixed-smooth-local-domain}
Let \(Y\) be of finite type over \(\Z_p\) with smooth
Jacobian charts in the sense of
Proposition~\ref{prop:complex-relative-jacobian} over
this base ring. All its local rings are domains and
multiplication by \(p\) is injective on them.
\end{lemma}
\begin{proof}
First consider a closed special point \(y\) with finite
residue field \(k\). There is a complete discrete
valuation ring \(D\), finite free over \(\Z_p\), with
uniformizer \(p\) and residue \(k\). For example adjoin
successive lifts of monic irreducibles for a finite
list of field generators of \(k/\mathbf F_p\); apply
Lemma~\ref{lem:integer-extension} at each step using
the irreducible residue polynomial, which is separable
over a finite field. Separability follows since
\(x\mapsto x^p\) is bijective on that field, so an
irreducible with zero derivative would have all
coefficients and monomials \(p\)-th powers. Finite
freeness follows from the monic power bases.
Over \(D\) choose the point \(y'\) with these same
residue coordinates in \(k\). Its local ring is flat
over \(\mathcal O_{Y,y}\), since scalar extension
by the free \(D\) and localization preserve injections
after tensoring. A flat local homomorphism of local rings
\(A\to A'\) is injective: for \(0\ne c\in A\), the
inclusion \(Ac\hookrightarrow A\) stays injective,
and \(Ac\otimes A'=A'/\operatorname{Ann}(c)A'\)
is nonzero since that annihilator maps into the
maximal ideal. Its generator maps to the image of \(c\).
It therefore suffices to work at \(y'\).

In its chart write \(t,z\) for the \(d\) free and \(r\)
dependent variables with \(r\) equations \(F=0\);
the matrix \(F_z\) is a unit matrix in determinant
at the point. Choose lifts \(t_0,z_0\) of the residue
values in \(D\). Put \(P=D[[T_1,\ldots,T_d]]\) and
\(I=(p,T_1,\ldots,T_d)\). This ring is complete and
separated for the powers of \(I\): in degree \(b\)
modulo \(I^b\), each coefficient of a monomial
\(T^\mu\) with \(|\mu|<b\) is specified modulo
\(p^{b-|\mu|}\), and all these specifications have
unique limits in \(D\). Solve
\(F(t_0+T,z_0+Z)=0\) with \(Z\in I^r\).
Modulo \(I\) this holds with \(Z=0\). Given a solution
modulo \(I^b\), the correction in
\((I^b/I^{b+1})^r\) is uniquely obtained by applying
the inverse of the residue matrix \(F_z\) to the
negative error, since all terms of order two in the
correction vanish there. Passing to the limit gives
the solution in \(P\). Any inverted chart functions
have nonzero residues and therefore are units in \(P\),
by the inverse series after factoring out a unit of \(D\).

Write \(A'\) for the local ring at \(y'\), with maximal
ideal \(\mathfrak m'\) generated by \(p,t-t_0,z-z_0\).
Evaluation at the solution identifies
\(A'/(\mathfrak m')^b\) with \(P/I^b\) for each \(b\).
Indeed substitution \(T=t-t_0\) defines an inverse
on the truncated polynomial representatives. The
two choices of \(z\) in the resulting quotient satisfy
the equations and agree modulo the maximal ideal.
Telescoping their polynomial differences multiplies
\(z-z'\) by a square matrix reducing to \(F_z\);
that matrix is invertible in the local quotient.
Thus the substituted \(z\)'s agree there as well.
The formulas agree on denominators by their unit
residues. Taking limits identifies the completed
local ring with \(P\). The ring \(A'\) is noetherian
by the polynomial-ring argument of
Lemma~\ref{lem:graded-generic-free-monomials}, since
a discrete valuation ring has only principal ideals.
It injects into its completion by the intersection-of-powers
argument in Lemma~\ref{lem:smooth-connected-fiber-integral}.
The power series ring over the domain \(D\) is a
domain by its first nonzero homogeneous term, and
\(p\) is nonzero there. This proves the assertion
at \(y\).

For a prime of a special affine chart select a closed
point of its closure after extension of the fiber
to \(\overline{\mathbf F}_p\). Such a point exists
by the affine point argument in
Lemma~\ref{lem:projective-special-section-base-change},
since extension of scalars on the nonzero quotient
by the prime is faithful. The finitely many coordinates
of the closed point generate a finite field over
\(\mathbf F_p\), so it contracts to a closed point
of finite residue field containing the original prime.
Localization of that closed local domain proves the
assertion at the original prime. On the generic fiber
use an algebraic closure of \(\Q_p\) and a geometric
closed point of the closure in an affine chart by
the same argument. The closed point there has a
power series local completion over the algebraically
closed field by Lemma~\ref{lem:smooth-connected-fiber-integral}'s
local Jacobian computation. Descent to the contracted
local ring is injective by the flat local argument
above, and localization gives a domain at the original
prime. On this fiber \(p\) is invertible.
\end{proof}

\begin{proposition}\label{prop:clopen-extension-local-jacobian}
For \(Y\) of Lemma~\ref{lem:mixed-smooth-local-domain},
restriction to \(\Q_p\) gives a bijection between
open and closed subschemes of \(Y\) and of its generic
fiber. The inverse takes reduced closure. A projective
\(Y\) therefore gives a smooth projective \(\Z_p\)-model
of every such generic subscheme.
\end{proposition}
\begin{proof}
Here an open and closed subscheme has the structure
induced on the open subset. The noetherian scheme
\(Y\) has finitely many irreducible components, by
the argument in Proposition~\ref{prop:finite-strata-graded-partial-resolutions}.
Two different components through one point would
give distinct minimal primes in its local ring as in
Lemma~\ref{lem:smooth-connected-fiber-integral},
contradicting Lemma~\ref{lem:mixed-smooth-local-domain}.
Thus the components are disjoint and each is both
open and closed. All their local rings are domains,
so they are reduced. The element \(p\) is nonzero
in every local ring. On every nonempty affine of a
component it is therefore a nonzero element of a
domain, whose generic point belongs to the open
subset \(p\ne0\). The generic fiber of that component
is consequently an irreducible dense open subset.
These opens are precisely the irreducible components
of the generic fiber, since they form a finite
disjoint closed covering of it. An open and closed
subset is a union of whole components in either
scheme, by irreducibility. The bijection and the
description by reduced closure follow. Such a union
inherits its Jacobian charts and its closed projective
embedding. It is flat over \(\Z_p\): its chart rings
have no \(p\)-torsion by the lemma, and a torsion-free
module over this discrete valuation ring is flat by
the directed union of free finitely generated submodules
in Proposition~\ref{prop:valuation-special-cluster-lengths}.
\end{proof}

\section{Recovering the finer level from marked columns}

\begin{theorem}\label{thm:fine-level-open-closed-torsion-orbit}
Fix the coarse data \(C_b=K(b_\sigma)\) in
Proposition~\ref{prop:coarse-model-before-coefficient-prime}.
Over \(E\), write \(S_b=\mathcal X(C_b,0)\) for the
coarse base and \(G_b=\mathcal X(C_b,1)\) for its
one-point group scheme over \(S_b\). Let
\(0\ne\alpha\in b_\sigma O_E\), put \(C_\alpha=K(\alpha)\),
and put \(S_\alpha=\mathcal X(C_\alpha,0)\).
In \(L=\widehat O_E^n\) let \(e_1,\ldots,e_n\) be the
standard columns and write \(t_j=\alpha^{-1}e_j\bmod L\).
The map sending a finer framed torus to its coarse point
and the ordered points \(P_x(t_1),\ldots,P_x(t_n)\)
is an isomorphism from \(S_\alpha\) onto an open and
closed subscheme of the \(n\)-fold product over \(S_b\)
of \(G_b[\alpha]=\ker[\alpha]\). For each \(\nu\ge0\),
\(\mathcal X(C_\alpha,\nu)\) identifies over \(E\)
with the pullback of the \(\nu\)-fold coarse group
power to this subscheme.
\end{theorem}
\begin{proof}
The full \(\alpha\)-torsion label group is
\(A_\alpha=\alpha^{-1}L/L\). Indeed in the coordinate
formula for \(P_x\) of
Lemma~\ref{lem:actual-subspace-kernel-labels},
the action of \(\alpha\) multiplies the finite
adelic lift by \(\alpha\); the adelic quotient there
labels all torus torsion. Alternatively its
\(\alpha\)-torsion is directly
\(\alpha^{-1}\Lambda_x/\Lambda_x\), with the
identification given by the finite lattice residues
as in Lemma~\ref{lem:all-torsion-recovers-frame}.
This is finite, since a nonzero rational integer
multiple of \(1/\alpha\) belongs to \(O_E\).
The group \(C_b\) acts on these labels by its matrices.
Inside \(C_b\) the simultaneous stabilizer of the
listed \(n\) columns is exactly \(C_\alpha\): fixing
them means \((k-I)e_j\in\alpha L\) for every \(j\),
which is the congruence in
Definition~\ref{def:moduli-signed-frame}. It fixes
every member of \(A_\alpha\), and \(C_\alpha\)
is the kernel of the action there.

For every tuple \(\mathbf s=(s_1,\ldots,s_n)\in
A_\alpha^n\) there is a map \(j_{\mathbf s}\) from
\(S_\alpha\) to the indicated product. In a ball
chart use the coarse frame and the points \(P_x(s_i)\)
as its marked slots. These are holomorphic in the
base since the real torsion periods are fixed linear
combinations with rational coefficients of its
holomorphic period columns in
Proposition~\ref{prop:signed-ball-period-coordinates}.
Right multiplication of a frame by \(k\in C_\alpha\)
does not change these marks; rational left changes
identify the points by the same linear lattice
isomorphism as the underlying group. Thus the maps
are well defined on the quotient. They take values
in the closed kernel product scheme, by the torsion
equalities. They are algebraic over \(\C\) by
Theorem~\ref{thm:holomorphic-map-algebraic-locally},
first as morphisms into the product of the coarse
group over its base. The conjugation of all the
marks for a fixed frame in
Theorem~\ref{thm:actual-common-level-equations-over-e}
is the identical list of labels at its reconstructed
frame. Consequently
Lemma~\ref{lem:equivariant-projective-map-descent}
descends \(j_{\mathbf s}\) to \(E\), compatibly with
the forgetful projection.

Fix one complex coarse point represented by a frame
\(x\) and its complex structure. Its lifts to \(S_\alpha\)
are represented by \(xk\), \(k\in C_b/C_\alpha\),
with that complex structure. The only additional
possible identifications would be rational stabilizers
of the coarse point, trivial by
Proposition~\ref{prop:neat-at-wild-place}.
Thus the tuple values of \(j_{\mathbf s}\) on this fiber
are exactly the \(C_b\)-orbit of \(\mathbf s\), in the
label coordinates of \(x\). For
\(\mathbf t=(t_1,\ldots,t_n)\) each value occurs
once by the stabilizer calculation above. The finitely
many orbits on \(A_\alpha^n\) exhaust the geometric
complex fiber, with disjoint sets of tuple values.

Each of these orbit images is closed: its map is from
a projective variety, so project its closed graph and
apply the fiberwise homogeneous minor argument in
Lemma~\ref{lem:projective-special-section-base-change}.
This describes the closed image over \(E\) as well
and commutes with field extension, since emptiness of
a projective fiber is witnessed by spanning all
monomials in one homogeneous degree, with the same
matrix rank after field extension. The union over the
other orbits therefore gives a closed complement to
the \(\mathbf t\)-image. Indeed any nonempty intersection
or uncovered locally closed subset after extension to
\(\C\) would contain a complex point by the affine point
argument of the same lemma. Extension of fields is
faithful on nonzero affine coordinate rings, so the
same partition holds over \(E\). Give the image its
open subscheme structure. The ambient characteristic-zero
kernel product is reduced: the differential of
\([\alpha]\) on the fibers is invertible by the norm
identity of Lemma~\ref{lem:fine-level-invertible-lie-action}.
The kernel charts in
Proposition~\ref{prop:fine-level-torsion-etale-charts}
over the smooth generic base give smooth charts for
the product; their local rings over \(\C\) are domains
by the chart argument in
Lemma~\ref{lem:smooth-connected-fiber-integral}.
Reducedness over \(E\) follows by faithful scalar
extension.

In the local ball charts of
Theorem~\ref{thm:analytic-framed-torus-family} both the
fine cover and these tuple values consist of disjoint
sheets over the base. For the tuple values this follows
also by applying the local inverse argument of
Theorem~\ref{thm:signed-fiber-projective-embedding}
to the zero-dimensional relative charts and continuing
each labelled section; the complete list above accounts
for the fiber. On each fine sheet \(j_{\mathbf t}\)
is the identity in its base coordinate. The bijection
onto the indicated union thus has holomorphic inverse.
The union is a closed projective smooth variety, so the
inverse is algebraic by
Theorem~\ref{thm:holomorphic-map-algebraic-locally}.
Uniqueness makes it commute with all coefficient
automorphisms fixing \(E\), and
Lemma~\ref{lem:equivariant-projective-map-descent}
descends it.

For the marked powers, retain the same \(\nu\) torus
points when forgetting from \(C_\alpha\) to \(C_b\).
In a ball chart the base sheets change and the
identification on each torus power is the identity.
Thus the resulting map to the stated fiber product is
a biholomorphism, including \(\nu=0\). It is algebraic
with algebraic inverse by the same theorem on
holomorphic maps, using the smooth product charts
of Theorem~\ref{thm:actual-complex-finite-etale}.
Forgetting the level and slots commutes with
coefficient automorphisms by
Theorem~\ref{thm:hecke-diagram-over-number-field};
uniqueness of the inverse again gives descent.
\end{proof}

\section{Descent from the completed local integers}

\begin{lemma}\label{lem:jacobian-chart-flat-descent}
Let \(B\to D\) be flat and let
\(A=B[X_1,\ldots,X_s]/I\) with \(I\) finitely generated.
Let \(y\) be a prime of \(A\) admitting a lift \(y'\) in
\(A'=D\otimes_B A\). If the latter scheme at \(y'\)
has a smooth Jacobian chart of relative dimension \(d\)
over \(D\), then near \(y\) the given presentation itself
has \(s-d\) defining equations with an invertible
\((s-d)\)-rowed minor of relative partial derivatives.
These equations may be chosen from any finite set of
generators of \(I\).
\end{lemma}
\begin{proof}
Put \(P'=D[X_1,\ldots,X_s]\) and \(I'=IP'\).
Localize \(P'\) at the inverse image of \(y'\); write
\(M\) for the resulting local ring, \(J=I'M\),
\(A'_0=M/J\). Consider \(C=M/J^2\), a square-zero
extension of \(A'_0\) with kernel \(V=J/J^2\).
There is a \(D\)-algebra section \(A'_0\to C\).
Indeed use the prescribed smooth chart at \(y'\);
locally it presents \(A'_0\) by free and dependent
variables and equations with invertible dependent-variable
Jacobian. Lift all their values from \(A'_0\) to \(C\).
The errors in their equations belong to \(V\), and
changing the dependent variables by elements of \(V\)
changes the errors by exactly that Jacobian, since
\(V^2=0\). Its determinant is a unit in \(C\), as it is
one in \(A'_0\) and the kernel is nilpotent. Solve for
corrections killing the errors. All elements inverted
in the chart and then in its local ring map to units
for the same reason, giving the required section.

Use the section to lift the images of the original
variables \(X_i\), and denote the difference between
\(X_i\) in \(C\) and this lift by \(v_i\in V\).
The gradient gives an \(A'_0\)-linear map
\[
             J/J^2\longrightarrow (A'_0)^s,\qquad
                       f\longmapsto (\partial_i f)_i .
\]
Here derivatives are reduced modulo \(J\); the
product rule defines them on the localized ring and
kills \(J^2\) after reduction. This map has left inverse
sending the \(i\)-th basis vector to \(v_i\). In fact
evaluate \(f\) at the lifted coordinates from the section;
the value is zero since \(f\in J\). Evaluating instead
at the original \(X_i\) adds
\(\sum_i(\partial_i f)v_i\) by the first-order
polynomial identity. It also holds for the inverted
denominators by \((c+v)^{-1}=c^{-1}-c^{-2}v\) when
\(v^2=0\). This gives the claimed left inverse.

After taking the residue field at \(y'\), the dimension
of \(V\) tensored with that field is therefore the
rank of the gradient rows of generators of \(I\).
Indeed those generators still generate \(J\), and
the split injection remains injective. This rank is
\(s-d\). The kernel of the matrix on column vectors
is the relative tangent space: an infinitesimal lift
of the residue point to its field with a square-zero
element is determined by increments of the \(s\)
variables satisfying precisely the linearized equations.
In the smooth chart the \(d\) free increments are
arbitrary and the other increments are uniquely
determined. Both calculations describe the same
lifts with coefficients from \(D\) fixed.

Flatness identifies \(I\otimes_B D\) inside \(P'\)
with \(I'\) and likewise identifies the quotient by
the square with \((I/I^2)\otimes_B D\), by tensoring
the inclusions and quotient sequences. Consequently
the dimension of \((I/I^2)\otimes_A\kappa(y)\) and
the rank of the gradient matrix over \(\kappa(y)\)
are both \(s-d\), since extension to \(\kappa(y')\)
preserves rank. Choose \(s-d\) generators giving
independent rows there and a corresponding unit minor.
They span \((I/I^2)\otimes\kappa(y)\): the gradient map
on that space has rank equal to its dimension. For
the prime \(\mathfrak q\) over \(y\) in \(B[X]\), this
space is \(I_{\mathfrak q}/\mathfrak q I_{\mathfrak q}\)
because \(I\subseteq\mathfrak q\). The finite-module
argument in Lemma~\ref{lem:finite-fiber-local-algebras}
applied to the quotient of \(I_{\mathfrak q}\) by
the selected equations shows they generate there.
Clear denominators in the expressions for all original
generators and invert the chosen minor. On this
neighborhood the asserted chart results.
\end{proof}

\begin{corollary}\label{cor:local-number-ring-clopen-extension}
Let \(v\) be a degree-one prime with
\(E_v=\Q_p\), let \(R=(O_E)_v\), and let \(Y/R\) be of
finite type with a base change to \(\Z_p\) admitting
smooth Jacobian charts. Such charts already exist over
\(R\), with the same relative dimensions. Its local
rings are domains, \(Y\) is flat over \(R\), and every
open and closed subscheme of \(Y_E\) extends uniquely
to an open and closed subscheme of \(Y\), by reduced
closure.
\end{corollary}
\begin{proof}
By Proposition~\ref{prop:general-ideal-factorization}
and the equality of completions, \(R\) is a discrete
valuation ring with uniformizer \(p\), residue
\(\mathbf F_p\) and completion \(\Z_p\). The latter
is flat over \(R\) by the torsion-free argument of
Proposition~\ref{prop:valuation-special-cluster-lengths}.
Every point of \(Y\) lifts after this extension.
Indeed the tensor product of its residue field with
\(\Z_p\) over \(R\) is that same field on the special
fiber and on the generic fiber it is the tensor
product of two extension fields over \(E\), which
is nonzero by taking bases. A maximal ideal of this
nonzero ring gives the lift. On an affine chart,
finite type is finite presentation since \(R\) and
its polynomial rings are noetherian by
Lemma~\ref{lem:graded-generic-free-monomials}.
Apply Lemma~\ref{lem:jacobian-chart-flat-descent}.
At each lift the induced map on local rings is
flat and local, hence injective by the argument of
Lemma~\ref{lem:mixed-smooth-local-domain}. The ring
upstairs is a domain in which \(p\ne0\) by that
lemma; so is the ring downstairs. This proves
flatness over \(R\) by the same torsion-free
criterion. The component argument of
Proposition~\ref{prop:clopen-extension-local-jacobian}
uses just these properties and noetherianity.
Applying it over \(R\) proves the remaining assertions.
\end{proof}

\section{A smooth model at the free-weight place}

\begin{theorem}\label{thm:free-weight-place-good-actual-model}
For the approximants and ordered choices in
Proposition~\ref{prop:coarse-model-before-coefficient-prime}
and Lemma~\ref{lem:fine-level-invertible-lie-action},
let \(v=w_1\), let \(R=(O_E)_v\), and take any of the
finer principal levels \(\alpha=\mathfrak n\) specified
there. There is a smooth projective \(R\)-model of
\(S_\alpha=\mathcal X(K(\alpha),0)_E\). Over it there
is a smooth projective commutative group with its
\(O_E\)-action, whose \(\nu\)-fold power is a smooth
projective \(R\)-model of \(\mathcal X(K(\alpha),\nu)_E\)
for every \(\nu\ge1\). The projection to the base,
the group operations, the actions of \(O_E\) and the
forgetful maps of slots have their specified generic
values.
\end{theorem}
\begin{proof}
Use the coarse models \(G_b\to S_b\to\operatorname{Spec}R\)
from Proposition~\ref{prop:coarse-model-before-coefficient-prime},
base changed from its localization of \(O_E\).
They have their group law and scalar maps over \(R\).
The two morphisms have smooth relative Jacobian
charts after extension to \(\Z_\ell\), by that
proposition, and over \(R\) by
Lemma~\ref{lem:jacobian-chart-flat-descent}.
For the relative assertion apply the latter result
on affine charts of \(S_b\) with the flat extension
obtained by tensoring the chart ring over \(R\)
with \(\Z_\ell\). The existence of a lift of each
point of \(G_b\) is the residue-field argument of
Corollary~\ref{cor:local-number-ring-clopen-extension}.
Charts of finite type are of finite presentation here
by the noetherian polynomial argument of
Lemma~\ref{lem:graded-generic-free-monomials}.

Form the projective \(R\)-scheme
\(Y=(G_b[\alpha])^n_{/S_b}\), where the notation
denotes the \(n\)-fold fiber product. The closed kernel
and its product commute with base change, since the
kernel is an inverse image of a section and the product
is given by the equalities of the base coordinates.
The prime data of the theorem make \(\alpha\) invertible
on the geometric relative tangent spaces over
\(\Z_\ell\) by Lemma~\ref{lem:fine-level-invertible-lie-action}.
Thus \(Y_{\Z_\ell}\) has smooth charts by
Proposition~\ref{prop:fine-level-torsion-etale-charts}.
Apply Corollary~\ref{cor:local-number-ring-clopen-extension}
to \(Y\). By Theorem~\ref{thm:fine-level-open-closed-torsion-orbit}
its generic fiber contains \(S_\alpha\) as the open
and closed part given by the ordered standard columns.
Its reduced closure \(U\) is open and closed in \(Y\).
Consequently \(U/R\) has the required smooth charts
and is flat. It is projective since it is closed in
the projective \(Y\). This gives the fine base at
the actual place \(v\).

Over \(U\) take \(G_b\times_{S_b}U\). Pullback of
the projective embedding and of the relative Jacobian
charts gives a smooth projective group over \(U\);
its operation and each scalar map are base changes
of those on \(G_b\). For a finite power over \(U\)
the relative equations in the point variables have
block diagonal Jacobian minors. Combining with the
smooth chart of \(U/R\) gives a unit minor over \(R\)
with the expected number of free variables, as in
Proposition~\ref{prop:coarse-model-before-coefficient-prime}.
The product has a projective embedding over \(R\):
embed it as a closed subscheme of the product of
the projective models of \(U\) and the coarse group
powers using the base diagonal equations. It is flat
over \(R\) by Corollary~\ref{cor:local-number-ring-clopen-extension}
applied to these charts after completion.
The generic identifications with the stated powers
and their maps are the last assertion of
Theorem~\ref{thm:fine-level-open-closed-torsion-orbit},
which retains the actual marked points in each slot.
\end{proof}

\begin{corollary}\label{cor:prescribed-place-spread-marked-powers}
For each fine base in Theorem~\ref{thm:free-weight-place-good-actual-model}
one may take its models, its one-point group and all
their finite powers over a single localization
\(O_E[1/c]\) with \(0\ne c\in O_E\setminus w_1\).
There the base and the group are projective and
have smooth relative Jacobian charts. The fibers
over \(w_1\) and the generic fiber agree with those
of the theorem. For any one \(\nu\), at all but
finitely many rational primes the marked power
identifies with the earlier model at its same level
used in Theorem~\ref{thm:actual-special-graph-point-congruence}.
The finite exceptions in this last identification
can depend on the approximant and on \(\nu\).
\end{corollary}
\begin{proof}
Take the projective equations of the fine base and
the projective coarse group with its projection,
and the graphs of the finitely many maps defining
the group operations and an integral basis of its
scalars. All are of finite presentation over the
noetherian \(R\). Include the relative smooth
Jacobian charts, their inverted minors, the group
and scalar identities, and the identifications
with the fiber products on each of their charts
from Theorem~\ref{thm:free-weight-place-good-actual-model}.
Finite covers suffice since the schemes and charts
are quasi-compact, by noetherianity. Each coefficient
and identity in these presentations uses finitely
many denominators of \(O_E\) outside \(w_1\).
For the covering assertions use the expressions
of one as a sum of the covering functions with
chart coefficients, and for the algebra identities
use finite expressions in the defining ideal, as
in Proposition~\ref{prop:finite-presentation-spread}.
Multiplying these finitely many denominators
together gives \(c\notin w_1\), after including the
denominators of the coarse data as well. Keeping
these equations and maps over \(O_E[1/c]\) gives
the asserted models. The powers and their slot
operations are then formed by fiber products, so
no further denominator depends on their number.
The generic and \(w_1\)-local identifications follow
by the original chart equalities after localization at
\(w_1\).

For fixed \(\nu\) the two projective models over
their common generic field are isomorphic by
Theorem~\ref{thm:fine-level-open-closed-torsion-orbit}.
Include both directions of this generic isomorphism,
their composites and chart covers in the
denominator-clearing argument of
Theorem~\ref{thm:fixed-integral-level-model}, over
the intersection of the two initial open bases.
This requires finitely many nonzero denominators
and yields an isomorphism on the resulting open.
Only rational primes dividing their norms or those
of the original exclusions can be lost. The
prescribed-place assertion comes directly from
the first construction.
\end{proof}

\section{Cohomology of differential forms}

\begin{definition}\label{def:affine-cover-hodge-cochains}
For a commutative \(k\)-algebra \(A\) over a field \(k\),
write \(\Omega^1_{A/k}\) for the \(A\)-module on symbols
\(da\), \(a\in A\), with additivity, \(dc=0\) for \(c\in k\)
and \(d(ab)=a\,db+b\,da\). Put
\(\Omega^h_{A/k}=\bigwedge_A^h\Omega^1_{A/k}\) for
\(h\ge0\), with \(\Omega^0_{A/k}=A\); the exterior
power is formed by the alternating multilinear relations.
For a smooth separated variety \(X/k\) of finite type,
take a finite ordered affine open cover \(\mathcal U=(U_i)\).
On an affine intersection \(U_J=\bigcap_{j\in J}U_j\)
use the module \(\Omega^h(U_J)=\Omega^h_{A_J/k}\) with
\(A_J=\Gamma(U_J,\mathcal O_X)\), zero on an empty
intersection. In degree \(q\ge0\) form the direct sum
\[
    C^q(\mathcal U,\Omega^h)
        =\bigoplus_{i_0<\cdots<i_q}\Omega^h(U_{\{i_0,\ldots,i_q\}}).
\]
Its differential \(\delta\) takes alternating restrictions
as in Definition~\ref{def:graded-chart-cochains}, now
without the augmented term. Write
\(H^q(X,\Omega^h_{X/k})\) for its cohomology, under the
canonical cover identifications below. The same cochains
will be used for systems \(\mathcal F\) of modules on
affines glued by restriction, for which restriction to
each principal open is localization. Such a system is
called quasi-coherent; finite modules on the affine
charts over a noetherian scheme give coherent systems.
\end{definition}

\begin{lemma}\label{lem:affine-forms-principal-acyclicity}
The modules of forms in
Definition~\ref{def:affine-cover-hodge-cochains} agree
on localization. On a smooth relative dimension-\(d\)
Jacobian chart \(\Omega^1\) over the base has the
\(d\) free coordinate differentials as a basis.
For a quasi-coherent module on an affine scheme
\(U\) and a finite affine open cover \(\mathcal V\)
of \(U\), augmenting its cover cochains by
\(\mathcal F(U)\) in degree \(-1\) gives an exact
complex.
\end{lemma}
\begin{proof}
A homomorphism \(A\to B\) sends \(da\) to \(d\bar a\)
in the notation of the definition, by its relations.
Inverting any element \(f\) forces
\(d(f^{-1})=-f^{-2}df\). This rule extends derivations
uniquely: the formula \(d(a/f^j)=f^{-j}da-ja f^{-j-1}df\)
respects equality of fractions by differentiating a
cleared equality after inverting \(f\). The universal
presentation therefore identifies the localized
differentials with those of \(A_f\). Localization
commutes with the alternating multilinear presentation
of every exterior power. On a polynomial quotient
with generators \(x_i\) and defining equations \(F_j\),
the \(dx_i\) generate with relations
\(\sum_i(\partial_iF_j)dx_i=0\); the polynomial
product rule gives the expression for every differential,
and assigns a derivation for every list of values satisfying
these relations. The same description holds relative to
a base ring. An invertible dependent-variable minor
solves the relations for the corresponding \(dx_i\)
in the free coordinate differentials, without further
relations. This proves the basis assertion.

A finite intersection of affine opens on a separated
scheme is affine: for two opens it is the inverse
image of the closed diagonal inside their affine
product, hence a closed affine subscheme there.
Induct for more opens. In the augmented cover
complex over \(U=\operatorname{Spec}A\), fix a prime
\(\mathfrak p\) of \(A\). There is a principal
neighborhood \(D(f)\) of it contained in one member
\(V_i\), since principal opens form an affine basis.
Localization at \(f\) restricts every cochain term
to the corresponding intersection with \(D(f)\):
the intersection before restriction is affine,
and \(f\) gives a principal open on it.
On this new cover \(V_i\cap D(f)=D(f)\).
Regard cochain entries as alternating in their
indices, zero on repetition. A contraction is obtained
by evaluating an entry with \(i\) inserted first,
including in degree zero the value at \(i\) for the
augmentation. Restrictions agree on its domains
since that member is all of \(D(f)\).
Deleting \(i\) returns the original entry and the
other deletions cancel in pairs in \(\delta H+H\delta\).
The localized complex is therefore exact, also
after further localization at \(\mathfrak p\).
Localization commutes with kernels and images by
clearing denominators as in
Lemma~\ref{lem:high-degree-chart-section-comparison};
every cohomology module thus vanishes at all primes.
If an element of a module is nonzero, its proper
annihilator belongs to a maximal ideal at which the
element stays nonzero. This proves the stated exactness.
In particular on affine principal covers the degree-zero
gluing identifies the common module of forms, giving
the asserted chart interpretation.
\end{proof}

\begin{proposition}\label{prop:affine-cover-hodge-functoriality}
The groups in Definition~\ref{def:affine-cover-hodge-cochains}
are independent of the finite affine cover. Algebraic
maps between these varieties induce functorial pullbacks
on the groups of forms. For a field extension \(k'/k\)
one has canonical isomorphisms
\[
 k'\otimes_k H^q(X,\Omega^h_{X/k})
       \simeq H^q(X_{k'},\Omega^h_{X_{k'}/k'}),
\]
compatible with those pullbacks.
\end{proposition}
\begin{proof}
For two covers \(\mathcal U,\mathcal V\), take in
bidegree \(a,b\ge0\) the cochains on intersections
using \(a+1\) ordered indices from \(\mathcal U\)
and \(b+1\) from \(\mathcal V\). The total differential
is \(\delta_U+(-1)^a\delta_V\), whose square is
zero by commutation of restrictions. Augment
vertically by the \(\mathcal U\) cochains at \(b=-1\).
Each column is exact by
Lemma~\ref{lem:affine-forms-principal-acyclicity},
applied on the affine intersection from \(\mathcal U\).
The augmented total complex is exact. Given a cocycle,
its part of smallest \(a\) is vertically closed,
so subtract a total boundary using a vertical primitive
with the prescribed sign. At the bottom of an exact
column a closed entry is already zero. Continuing
in increasing \(a\) terminates because the cover is
finite. Thus restriction from the \(\mathcal U\)
complex to the unaugmented total complex at \(b=0\)
induces an isomorphism on cohomology. Explicitly a
cocycle in the total complex has an augmented primitive;
its \(b=-1\) part is a cocycle giving that class by
restriction up to the degree sign, and exactness
applied in the reverse direction gives injectivity.
Interchanging the covers gives the same assertion
for \(\mathcal V\) at \(a=0\); the sign convention
is interchanged by multiplying bidegree \((a,b)\)
by \((-1)^{ab}\).

When \(\mathcal V\) refines \(\mathcal U\), choose
\(i(j)\) with \(V_j\subseteq U_{i(j)}\).
From the total complex to the \(\mathcal V\) complex,
evaluate an \((a,b)\)-array on the two lists
\[
   i(j_0),\ldots,i(j_a)
        \quad\hbox{and}\quad j_a,\ldots,j_{a+b},
\]
and restrict to their common \(\mathcal V\)-intersection.
For a total cochain sum over \(a+b\) equal to its
degree. Alternating extensions of the entries give
a cochain map. In its differential, deleting an index
before the shared index is a horizontal deletion and
deleting one after it is a vertical deletion with
the factor \((-1)^a\). The two terms at the junction
in the sum of these differentials have the same
restrictions and opposite signs. This verifies the
identity including degrees zero and one. On cochains
from \(\mathcal V\) it is the identity; on those from
\(\mathcal U\) it is exactly restriction by \(i(j)\).
Hence every such refinement induces the same canonical
isomorphism on cohomology. A third cover refines any
two by all their intersections, and refinement maps
compose by restriction, giving the asserted compatible
identifications.

For a morphism refine the source cover by affine opens
mapping into members of the chosen target cover. A
finite such refinement exists since affine opens form
a basis and the source is noetherian and quasi-compact,
by the polynomial noetherian argument of
Lemma~\ref{lem:graded-generic-free-monomials}. Pullback
of functions and differentials then gives the cochain
map with those index choices. Two choices for the map
of indices yield homotopic restrictions: on a degree
\(q\) array use the sum over \(0\le l<q\) with sign
\((-1)^l\) of the entries at
\[
       i(j_0),\ldots,i(j_l),i'(j_l),\ldots,i'(j_{q-1}).
\]
Pull back to the indicated common source intersection.
In \(\delta H+H\delta\), all deletions off the doubled
index cancel, and those at that index telescope to
the difference of the two pullbacks. Refining on the
source and target therefore gives the same map; for
composites choose compatible refinements and use
composition on differentials. Finally the base-changed
affine cover and all its intersections are their
scalar extensions. The differential presentation and
exterior powers give the base-changed modules of
relative forms. Tensoring over a field preserves
kernels and images by taking vector space bases,
so the isomorphism follows term by term on the cover
complex and on its cohomology.
\end{proof}

\section{Finite Hodge groups and descent along a level cover}

\begin{proposition}\label{prop:finite-projective-coherent-cohomology}
Let \(X\) be a closed subscheme of \(\mathbf P^N_k\)
over a field \(k\), and let \(\mathcal F\) be a coherent
module in Definition~\ref{def:affine-cover-hodge-cochains}.
The groups \(H^q(X,\mathcal F)\) of that definition
are finite dimensional over \(k\) for all \(q\ge0\).
They vanish for \(q>N\).
\end{proposition}
\begin{proof}
Use the \(N+1\) standard affine charts. The proof
of the cover comparison in
Proposition~\ref{prop:affine-cover-hodge-functoriality}
applies to any of the indicated modules by
Lemma~\ref{lem:affine-forms-principal-acyclicity}.
View the module on \(\mathbf P^N\) by taking on any
affine the sections on its closed intersection with
\(X\). On each standard affine this is a finite
module over its polynomial coordinate ring via the
quotient defining \(X\), and principal restrictions
are still localizations. The two standard cover
complexes are identical. Their terms above \(N\)
are zero.

For a coherent module \(\mathcal M\) on \(\mathbf P^N\)
choose finitely many generators in each standard chart.
Each generator \(s\) on the \(i\)-th chart extends,
after twisting, to a global section: on the \(j\)-th
chart multiply its restriction to the overlap by
a power of \(X_i/X_j\) and clear the localization
denominator. Increase to a common power on all charts
and use the frames of \(\mathcal O(t)\) there. On each
overlap the resulting difference becomes zero after
inverting \(X_i\), and increasing the twist by another
common power of \(X_i\) kills all those differences.
This is precisely the localization argument of
Lemma~\ref{lem:projective-line-high-global-generation},
which uses only finite generators and restriction to
principal opens in this step. The resulting finitely
many sections, with their respective twists, give
a surjection \(\mathcal E\to\mathcal M\) from a
finite sum of lines \(\mathcal O(-t)\). The kernel
\(\mathcal K\) is coherent: on every standard chart
it is a submodule of a finite module over a noetherian
ring by Lemma~\ref{lem:graded-generic-free-monomials},
and these kernels localize. Generation on each chart
is immediate from its own selected generators, whose
powers of \(X_i\) are units in that chart frame.
Thus the three cover complexes are termwise short
exact, including on the intersections by localization.

For any integer \(b\), compute the cover complex
of \(\mathcal O(b)\) in Laurent monomials of total
degree \(b\), by Definition~\ref{def:graded-chart-cochains}.
When the set of negative exponents is proper
and nonempty, the contraction of
Lemma~\ref{lem:high-degree-chart-section-comparison}
leaves no cohomology. When that set is empty, the
same contraction on the augmented complex leaves
in the unaugmented complex just one copy of \(k\)
in degree zero for that monomial. When every exponent
is negative only the top intersection occurs, giving
one copy in degree \(N\). There are finitely many
all-nonnegative and all-negative exponent lists of
a fixed sum. This proves finiteness for each line
and hence for \(\mathcal E\).
From the termwise short exact sequence the groups
\[
     H^q(\mathbf P^N,\mathcal E)
       \longrightarrow H^q(\mathbf P^N,\mathcal M)
       \longrightarrow H^{q+1}(\mathbf P^N,\mathcal K)
\]
are exact at the middle: lift a cocycle to an
\(\mathcal E\)-array and take its differential in
\(\mathcal K\), then use the cocycle lifting argument
of Lemma~\ref{lem:high-degree-chart-section-comparison}.
Starting above \(N\) and descending, suppose degree
\(q+1\) is finite for every coherent module. Both
outer spaces are then finite dimensional, proving
the same at degree \(q\) for every \(\mathcal M\).
Apply the conclusion to the module supported on \(X\).
\end{proof}

\begin{proposition}\label{prop:coherent-finite-level-trace}
Let \(k\) be a field of characteristic zero and let
\(\pi:W\to X\) be a finite locally free cover of
smooth projective \(k\)-varieties with relative
étale Jacobian charts. Suppose a finite group \(D\)
acts over \(X\) and that
\((w,d)\mapsto (w,wd)\) identifies \(W\times D\)
with \(W\times_X W\). For all \(q,h\) there is a
trace map
\[
  \operatorname{Tr}_\pi:H^q(W,\Omega^h_{W/k})
                       \longrightarrow H^q(X,\Omega^h_{X/k})
\]
with \(\operatorname{Tr}_\pi\pi^*=|D|I\) and
\(\pi^*\operatorname{Tr}_\pi=\sum_{d\in D}d^*\).
In particular the pullback identifies the space
on \(X\) with the \(D\)-invariant subspace on \(W\).
This applies to every generic cover \(\pi_m\) of
Theorem~\ref{thm:integral-hecke-cover-refinements}
and after extension of its ground field.
\end{proposition}
\begin{proof}
For an affine \(U=\operatorname{Spec}A\) of \(X\),
write \(\pi^{-1}U=\operatorname{Spec}B\). Then
\(\Omega^1_{B/k}=B\otimes_A\Omega^1_{A/k}\).
To check the natural map locally, a \(k\)-derivation
of \(A\) into a module over a relative étale chart
extends uniquely to its coordinate ring: the derivatives
of its extra variables solve the linear equations
obtained by differentiating the defining relations,
whose square Jacobian is invertible. Existing
derivatives of coefficients give the inhomogeneous
terms. Denominators extend by the quotient rule of
Lemma~\ref{lem:affine-forms-principal-acyclicity}.
This verifies the equality by the universal
differential presentations, and on taking exterior
powers it holds for \(\Omega^h\) as well. A map of
modules is an isomorphism when its kernel and cokernel
vanish at every prime, by the annihilator argument
in the same lemma.

Take the trace of multiplication by \(b\) on the
locally free \(A\)-module \(B\). These matrix traces
are independent of changes of basis, by
\(\operatorname{tr}(PQ)=\operatorname{tr}(QP)\)
obtained by interchanging the two finite sums, and
glue to an \(A\)-linear map \(B\to A\).
Tensoring it with the identity on the forms on \(U\)
gives the trace on forms. These maps commute with
restriction and scalar extension because the local
matrices do. On a finite affine cover of \(X\) and
its inverse-image affine cover of \(W\) they therefore
induce the asserted cohomology map, independent of
the cover by Proposition~\ref{prop:affine-cover-hodge-functoriality}.

Every geometric fiber has \(|D|\) points, each of
length one by the relative étale charts and the
local tangent argument of
Proposition~\ref{prop:fine-level-torsion-etale-charts};
the cover is nonempty over each point, and the
given product isomorphism gives a free transitive
action there. The rank of \(B\) is thus \(|D|\);
multiplication by a scalar from \(A\) has trace that
scalar times \(|D|\). For the other identity, change
base to \(B\) on the first copy in the fiber product.
The given isomorphism identifies the algebra
\(B\otimes_A B\) with \(\prod_{d\in D} B\), the second
\(B\) acting in its \(d\)-th component by \(d^*b\).
The trace after this change is the sum of those
components, by its diagonal multiplication matrix.
Applying the same equalities to forms and then
to cover complexes gives the two identities.
Since \(|D|\) is invertible in \(k\), the first
one gives injectivity and the second writes every
invariant class as the pullback of its normalized
trace. The converse inclusion follows from
\(\pi d=\pi\). The covers cited in the statement
have precisely these properties by
Theorem~\ref{thm:integral-hecke-cover-refinements}.
\end{proof}

\section{Hecke operators on the groups of differential forms}

\begin{theorem}\label{thm:actual-hecke-action-on-forms}
Fix a neat admissible level \(C=K(\mathfrak n)\), a
marked power \(X=\mathcal X(C,\nu)_\C\), and integers
\(h,q\ge0\). Put \(V^{h,q}=H^q(X,\Omega^h_{X/\C})\).
For each \(H\in\mathcal H^+\) of
Definition~\ref{def:actual-positive-hecke} there is
a unique operator \(T_H\) on this finite-dimensional
space satisfying, on each sufficiently fine cover,
\[
       \pi_m^*T_H=\sum_{gC}H(gC)\,u_{g,\mathbf1}^* .
\]
Only the support of \(H\) occurs. One has
\(T_{H K}=T_HT_K\), linearity in \(H\), and the
identity coset operator is the identity. For each
\(\mathbf a\in O_E^\nu\), let \(d_{\mathbf a}\) be
scalar multiplication slot by slot on \(X\). Then
\[
   \sum_{gC}H(gC)\,u_{g,\mathbf a}^*
             =\pi_m^*T_H d_{\mathbf a}^* .
\]
The constructions with coefficients in any subfield of
\(\C\) containing \(E\) extend by scalars to \(\C\).
\end{theorem}
\begin{proof}
Finiteness follows from
Proposition~\ref{prop:finite-projective-coherent-cohomology}
and the smooth-chart description of forms in
Lemma~\ref{lem:affine-forms-principal-acyclicity}.
The sum in the first display is invariant under the
deck transformations of the refined cover. In fact,
right translation by \(k\in C\) on its source changes
the target represented by \(xg\) into the target
represented by \(xkg\). The torus marks agree under
this equality by Proposition~\ref{prop:complex-actual-hecke-legs}.
Left translation by \(k\) permutes the right cosets
in the support and preserves their coefficients by
Definition~\ref{def:actual-positive-hecke}. Thus
Proposition~\ref{prop:coherent-finite-level-trace}
gives a unique descended operator, obtained by applying
\(\operatorname{Tr}_{\pi_m}/[C:C_m]\) to the sum.
Pulling back to a common deeper level identifies both
sums since their maps there have the same frame and
slot formulas. The pullback from \(X\) to that common
level is injective by the same trace proposition,
proving independence. The formula gives linearity,
and for the identity coset the map is \(\pi_m\).

For multiplication choose a level \(m\) fine enough
for the support of \(K\). Choose \(l\ge m\) so fine
that \(g^{-1}C_l g\subseteq C_m\) for the finitely
many chosen representatives for \(H\), and fine enough
for all products of these representatives with those
for \(K\). This is possible by the congruence
neighborhood argument of
Proposition~\ref{prop:complex-actual-hecke-legs}.
The same frame formula gives a map \(v_g\) from the
\(l\)-level to the \(m\)-level with all slot scalars
one: the inclusion just imposed gives well-definedness,
and the period charts give holomorphy. The map is
algebraic over \(E\) by the proof of
Theorem~\ref{thm:hecke-diagram-over-number-field},
using this finer congruence in the target. On the
\(l\)-level the unscaled \(g\)-map to \(X\) is
\(\pi_m v_g\), and composing \(v_g\) with the unscaled
\(s\)-map from the \(m\)-level to \(X\) gives the
\(gs\)-map there. These identities follow by multiplying
the frame and retaining the same torus vectors modulo
the successive enlarged lattices. By pullback
functoriality in Proposition~\ref{prop:affine-cover-hodge-functoriality},
\[
 \pi_l^*(T_HT_K)
    =\sum_{gC}H(gC)v_g^*\pi_m^*T_K
    =\sum_{gC,sC} H(gC)K(sC)\,u_{gs,\mathbf1}^* .
\]
Grouping equal final cosets gives precisely the product
of Definition~\ref{def:actual-positive-hecke}. Injection
of \(\pi_l^*\) proves the claim. The identity for
scalar slots follows from \(u_{g,\mathbf a}=
d_{\mathbf a}\,u_{g,\mathbf1}\) in the period formula.
Descent over \(E\) of all maps is in
Theorem~\ref{thm:hecke-diagram-over-number-field},
and the field-extension statements for pullback
and trace follow from
Propositions~\ref{prop:affine-cover-hodge-functoriality}
and~\ref{prop:coherent-finite-level-trace}.
\end{proof}

\begin{corollary}\label{cor:coherent-regular-left-multiple}
For an approximating form and its \(\theta\) in
Theorem~\ref{thm:actual-marked-weight-occurrence}, fix one
tested index \(h=h_j\). Let \(F\) be a finite polynomial
in words in the safe generators with \(F(\theta)=1\),
and write \(F(\tau)\) for its value in \(\mathcal H^+\)
as in Definition~\ref{def:marked-weight-test}. If
\(F((T_\tau)_\tau)\) acts as zero on \(V^{h,q}\) for
\(0\le q\le\dim X\), then for every left multiple
\(H=B F(\tau)\), every \(\mathbf a\) and every such
\(q\) the map
\[
      \sum_{gC}H(gC)\,u_{g,\mathbf a}^*:
         V^{h,q}\longrightarrow
         H^q(\mathcal X(C_m,\nu)_\C,\Omega^h)
\]
is zero at every sufficiently fine common level.
This includes the regular multiple and the scalar
slots of Corollary~\ref{cor:actual-holomorphic-local-sum}.
\end{corollary}
\begin{proof}
Theorem~\ref{thm:actual-hecke-action-on-forms} turns
word evaluation into composition in its given order
and sends the constant identity to the identity.
Thus \(T_H=T_B F((T_\tau)_\tau)=0\) on each stated
space. Apply its last displayed pullback formula,
including when a slot scalar is zero. The cited
corollary takes \(H\) of this form and
\(\mathbf a=2\boldsymbol\alpha\).
\end{proof}

\section{Exterior products and the coefficient of a scaled graph}

\begin{proposition}\label{prop:coherent-forms-product-kunneth}
Let \(X,Y\) be smooth separated varieties of finite
type over a field \(k\). For \(a,r\ge0\) there is a
natural isomorphism
\[
 \bigoplus_{\substack{i+j=a\\ b+c=r}}
 H^b(X,\Omega^i_{X/k})\otimes_k H^c(Y,\Omega^j_{Y/k})
      \ \xrightarrow{\ \simeq\ }\
 H^r(X\times_kY,\Omega^a_{(X\times Y)/k}).
\]
On the summand \(i,j\) the forms on the first factor
precede those on the second. Denote the image of an
elementary tensor by \(v\boxtimes w\).
\end{proposition}
\begin{proof}
On affine factors of rings \(A,B\), a \(k\)-derivation
on \(A\otimes_k B\) is determined by two derivations
on \(A,B\), by the product rule applied to
\((a\otimes1)(1\otimes b)\). Conversely any such
pair with values in an \(A\otimes B\)-module extends
by that rule; its formula is \(k\)-bilinear on \(a,b\).
The corresponding presentation of
\(\Omega^1_{(A\otimes B)/k}\) is the direct sum
of the two extended differential modules. Its
\(a\)-th exterior power is the direct sum of the
exterior products of degrees \(i+j=a\).
On smooth coordinates the ordered wedges of the
two differential bases in
Lemma~\ref{lem:affine-forms-principal-acyclicity}
give a basis for this equality, and the maps glue
since they are pullback followed by wedge. On an affine
rectangle the sections of the summand are
\(\Omega^i(A)\otimes_k\Omega^j(B)\), by the tensor
presentation. Write \(\mathcal M\) for this summand
and fix the pair \(i,j\).

Choose finite affine covers \((U_i')\), \((V_j')\)
of the two factors. Form the double complex of
\(\mathcal M\) on products of their respective
intersections, with differential
\(\delta_U+(-1)^b\delta_V\) in bidegree \((b,c)\).
Its total complex computes the cohomology of
\(\mathcal M\) on the product. To see this, use a
third index of cochains from the ordinary rectangle
cover \(\mathcal W=(U_i'\times V_j')_{i,j}\).
In tridegree \((b,c,e)\) take sections on the
intersection of the double rectangle with each
\((e+1)\)-fold intersection from \(\mathcal W\).
Use total signs \(1,(-1)^b,(-1)^{b+c}\).
Restriction at \(e=0\) embeds the double complex
in this triple calculation at the level of cocycle
classes. In fact it induces an isomorphism: augment
by the double complex at \(e=-1\). For fixed first
and second index lists the remaining augmented column
is exact by Lemma~\ref{lem:affine-forms-principal-acyclicity}
on the affine rectangle. The total cycle-elimination
argument of Proposition~\ref{prop:affine-cover-hodge-functoriality},
in increasing \(b+c\), proves the assertion.

In the other direction, fix one intersection from
\(\mathcal W\). It is contained in the product of the
first and second members of any one rectangle in
its list. Therefore the augmented complex in the
second cover index contracts by inserting that
second member, for every fixed first index list,
even after intersecting with all the first-index
opens. The same holds for the first index after
dropping the second, by inserting the first member
of the selected rectangle. These are the contractions
in Lemma~\ref{lem:affine-forms-principal-acyclicity};
their augmented terms are sections on the intersection
using no opens from the dropped index. Applying the
same finite cycle-elimination argument successively
to the two augmented indices proves that restriction
from the ordinary \(\mathcal W\) complex gives an
isomorphism onto the cohomology of the triple
complex as well. Equivalently one can first eliminate
at each fixed third degree, then in increasing third
degree, since taking its differential only increases
that index. This proves the comparison with the
ordinary cover in Definition~\ref{def:affine-cover-hodge-cochains}.

The double complex is the tensor product over \(k\)
of the two factor complexes for degrees \(i,j\).
For a complex of vector spaces choose in each degree
a complement of the boundaries in the cycles, and
a complement of the cycles in the whole space. Its
differential identifies the latter complement with
the boundaries in the next degree. The complex is
thus a direct sum of spaces with zero differential
representing its cohomology and of pairs joined by
an isomorphism. Tensoring such a pair with any other
complex gives a contractible complex: use the inverse
on the paired component as the homotopy there, and the
opposite signs in consecutive degrees cancel the
terms from the other differential. Applying this
to both factors shows that tensoring representatives
of cocycles induces the displayed isomorphism.
The map on classes is independent of the complements.
All comparisons above are given by restrictions and
tensoring representative arrays. For a pair of algebraic
maps into \(X,Y\), refine the two source covers to
map into the two target covers as in
Proposition~\ref{prop:affine-cover-hodge-functoriality};
their product rectangles refine the corresponding
product cover. Restriction and pullback commute in
all three indices, so both comparison isomorphisms
commute with these maps on cohomology. This proves
naturality and independence of the chosen covers.
\end{proof}

\begin{corollary}\label{cor:scaled-diagonal-weight-annihilation}
Let \(X,W\) be smooth projective over \(k\), put
\(d=\dim X\), and let \(u_i,c:W\to X\) be morphisms
for \(i\) in a finite set, with coefficients
\(\beta_i\in k\). For \(0\le h\le d\), suppose
\(\sum_i\beta_i u_i^*\) is zero on
\(H^q(X,\Omega^h_{X/k})\) for every \(0\le q\le d\).
On degree-\(d\) forms of \(X\times X\) let \(S_T\)
multiply the summand of first form degree \(j\) by
\(T^j\), where \(T\) is an indeterminate. For every
\(\delta\in H^d(X\times X,\Omega^d_{(X\times X)/k})\)
one then has
\[
      [T^h]\sum_i\beta_i (u_i,c)^*(S_T\delta)=0
                       \quad\hbox{in }H^d(W,\Omega^d_{W/k}).
\]
\end{corollary}
\begin{proof}
By Proposition~\ref{prop:coherent-forms-product-kunneth}
the coefficient at \(T^h\) is a sum of external products
\(v\boxtimes w\) before pullback, with
\(v\in H^q(X,\Omega^h)\) and
\(w\in H^{d-q}(X,\Omega^{d-h})\), \(0\le q\le d\).
Write \(\Delta_W:W\to W\times W\) for the diagonal.
Since \((u_i,c)=(u_i\times c)\Delta_W\), the weighted
sum on each product is
\[
    \Delta_W^*\left(\Big(\sum_i\beta_i u_i^*v\Big)
                                \boxtimes c^*w\right)
\]
by the naturality just proved and
Proposition~\ref{prop:affine-cover-hodge-functoriality}.
It vanishes by the hypothesis. Summing gives the result.
\end{proof}

\section{Cohomology with support and the gluing of local classes}

\begin{definition}\label{def:supported-affine-cochains}
Let \(Z\) be a closed subset of a separated finite-type
\(k\)-scheme \(X\), and let \(\mathcal F\) be a coherent
module. Choose finite ordered affine covers \((U_i)\) of
\(X\) and \((V_j)\) of \(X\setminus Z\); an empty open
may have the empty cover. In bidegree \((a,0)\), \(a\ge0\),
take the ordinary \(U\)-cochains of degree \(a\).
In \((a,b)\), \(b\ge1\), take the sum of the modules on
\(U_I\cap V_J\) for \(|I|=a+1,\ |J|=b\), with ordered
nonempty lists in both indices. The horizontal
differential is the \(U\)-cover differential. In the
vertical direction use restriction from \(b=0\) to
\(b=1\), then the \(V\)-cover differential. The total
differential is \(\delta_U+(-1)^a\delta_V\), with the
augmentation included in \(\delta_V\). Its square is
zero by the restriction identities. Denote its cohomology
in degree \(r\) by \(H_Z^r(X,\mathcal F)\). Projection
to \(b=0\) gives a map to ordinary cohomology, called
forgetting the support.
\end{definition}

\begin{proposition}\label{prop:support-cover-glue}
Definition~\ref{def:supported-affine-cochains} is independent
of both covers. Its groups for differential forms admit
pullback by \(f:X'\to X\) with \(f^{-1}Z\subseteq Z'\),
taking values in the groups with support \(Z'\); these
pullbacks commute with forgetting support. In particular
support can be enlarged. The construction is also
functorial for module homomorphisms on a fixed scheme.

Let \(d\ge1\). Assume that on every nonempty finite
intersection \(U_I\) one has
\(H^b_{Z\cap U_I}(U_I,\mathcal F)=0\) for \(0\le b<d\).
Then restriction identifies \(H^d_Z(X,\mathcal F)\)
with the lists of classes in \(H^d_{Z\cap U_i}(U_i,\mathcal F)\)
whose two restrictions on every \(U_i\cap U_{i'}\) agree.
\end{proposition}
\begin{proof}
For fixed \(U_I\), its vertical complex is the complex
for support on this affine using the single chart
and the cover by \(U_I\cap V_j\) of its complement.
Changing the \(V\)-cover by refinement induces an
isomorphism on the cohomology of this column.
Indeed its degrees \(b\ge1\) comprise the ordinary
cover complex shifted by one, compared in
Proposition~\ref{prop:affine-cover-hodge-functoriality};
the degree-zero term and its restriction are left
compatible with that refinement. In degree zero
the kernel consists of the sections vanishing on
the complement. In degree one take the ordinary
degree-zero cohomology modulo the restricted sections,
while in higher degrees the groups are the ordinary
ones with the indicated shift. This proves the
column claim. For a refinement of the \(U\)-cover,
fix instead the vertical degree. At \(b=0\) this
is the ordinary cover comparison on \(X\), and at
\(b>0\) it is the sum of such comparisons on the
affines \(V_J\), covered by their intersections
with the \(U_i\).

A map of bounded double complexes inducing
isomorphisms on cohomology of all columns induces
one on total cohomology. For clarity, attach to each
column map \(L\to M\) its complex in degree \(j\)
with terms \(M^j\oplus L^{j+1}\) and differential
\((m,l)\mapsto (dm+f(l),-dl)\). This complex is
exact precisely when \(f\) induces an isomorphism:
a closed pair first gives a closed \(l\) mapping to
a boundary, and after removing a boundary for \(l\)
the remaining closed \(m\) tests surjectivity.
Conversely applying exactness to \((m,0)\) and to
pairs lifting a kernel class gives the two implications.
With the total signs these column complexes assemble
to the same construction on totals, up to the sign
change \((-1)^a\) on the shifted component in column
\(a\). Exact columns give exact total complex by
eliminating a cocycle in increasing horizontal degree
as in Proposition~\ref{prop:affine-cover-hodge-functoriality}.
The row version follows by interchanging the indices.
This proves invariance under each refinement.
Choices of assignments for a refinement or a pair
pullback give the homotopy of the same proposition
in a changed index; in the vertical index multiply
by \((-1)^a\) and use zero homotopy on the augmented
entry. Restrictions in the unchanged index commute
with it, so the other terms cancel under the total
signs. Intersecting covers and composing the chosen
restrictions gives independence and functoriality.
The hypothesis on \(f^{-1}Z\) permits independent
affine refinements of the two source covers over
the two target covers. Projection to \(b=0\) is
compatible with all these maps. Module homomorphisms
act on all terms by restriction.

For the gluing assertion represent the given classes
by vertical \(d\)-cocycles at \(a=0\). Their horizontal
differential is a vertical boundary in column one
by agreement on double overlaps. Choose vertical
primitives there with sign giving cancellation in
the total differential. After choosing columns
through \(a-1\), the next uncancelled term in column
\(a\ge2\) is vertically closed since the total
differential squares to zero. Its vertical degree
is \(d+1-a<d\), so it too has a primitive when
that degree is positive; at degree zero a closed
entry is already zero under the hypothesis. Continuing
produces a total cocycle. Conversely any total
cocycle supplies such an agreeing list. If its
initial classes vanish, remove vertical boundaries
in column zero by subtracting total boundaries.
At every subsequent nonzero column the remaining
first entry is vertically closed in degree \(d-a<d\)
and can likewise be removed, with the degree-zero
case necessarily zero. The finite number of columns
gives injectivity. A restriction to \(U_i\) computes
the initial class by the projection to that chart,
using the comparison with its single-open cover.
\end{proof}

\section{Local symbols for a transverse system of equations}

\begin{lemma}\label{lem:supported-regular-parameters}
Let \(U=\operatorname{Spec}A\) be smooth over an algebraically
closed field \(k\), let \(Z\) be the common zero locus of
\(f_1,\ldots,f_d\in A\) with \(d\ge1\), and suppose
their differentials at every closed point of \(Z\) are
linearly independent. Let \(M\) be a finite locally
free \(A\)-module, also denoting its corresponding
sheaf on \(U\) by \(M\). Then \(H^b_Z(U,M)=0\) for
\(b<d\). The top group in degree \(d\) is the quotient of
\(M_{f_1\cdots f_d}\) by the sum of the images of the
localizations omitting one factor, taking the unlocalized
\(M\) when \(d=1\). In this quotient write
\[
    \left[\frac{\omega}{f_1,\ldots,f_d}\right]
        \quad\text{for the class of}\quad
               \frac{\omega}{f_1\cdots f_d},\qquad \omega\in M .
\]
The class depends only on \(\omega\bmod (f_1,\ldots,f_d)M\).
These groups commute with restriction to principal opens.
\end{lemma}
\begin{proof}
Use the single chart \(U\) and the ordered principal opens
\(D(f_j)\) on its complement in
Definition~\ref{def:supported-affine-cochains}.
The term with subset \(J\) is \(M_{\prod_{j\in J}f_j}\)
in degree \(|J|\), also for the empty subset in degree
zero. Hence the stated top quotient follows from its last
differential. Multiplying a numerator by \(f_j\) makes
the fraction come from the term omitting \(j\). The
complex localizes termwise, proving the restriction claim.

For \(N\ge1\) use the finite complex with one copy of
\(M\) in every subset position and differential from
\(J\) to \(J\cup\{j\}\) given by multiplication by
\(f_j^N\) with the same incidence sign. Its square is
zero by commutation of these multiplications. In
passing from \(N\) to \(N+1\) multiply the copy at
\(J\) by \(\prod_{j\in J} f_j\). These maps commute with
the differentials, and their direct limit is the above
localized complex by sending an entry at \(J\) to its
fraction with denominator \((\prod_{j\in J} f_j)^N\).
For a direct sequence cohomology commutes with this
limit: a differential which is zero in the limit is
already zero at a later stage, and an equality to a
boundary can likewise be tested at a common later
stage for its two representatives.

We check the finite complexes at maximal ideals.
Away from \(Z\) some \(f_j\) is a unit, and insertion
of that index followed by division by \(f_j^N\)
contracts the complex as in
Lemma~\ref{lem:affine-forms-principal-acyclicity}.
At a closed point of \(Z\), complete the independent
differentials of the \(f_j\) to a basis of cotangent
coordinates using linear combinations of the smooth
chart parameters. The completed local ring is a
power series ring in those parameters by the local
computation of Lemma~\ref{lem:smooth-connected-fiber-integral}.
It is also a power series ring in the \(f_j\) and
the chosen remaining coordinates, centered at the
point. Indeed this substitution has invertible linear
part, so it has an inverse modulo each successive
power of the parameter ideal by solving for the
homogeneous correction with that same linear part.
The inverses are compatible and give the power series
inverse. The completed quotient by any initial list
\(f_1^N,\ldots,f_i^N\) in this power series ring
has unique expressions with the first \(i\) exponents
less than \(N\), by separating monomials. Thus the
next \(f_{i+1}^N\) acts injectively there.

In the local ring the quotient by that initial ideal
injects into the corresponding quotient of the power
series ring. If a local element lies in the ideal
upstairs, approximate its finitely many coefficient
multipliers modulo every power of the maximal ideal
by local elements, using the truncated chart isomorphisms.
Its class in the noetherian local quotient then lies
in every maximal ideal power and is zero by the
intersection argument of
Lemma~\ref{lem:smooth-connected-fiber-integral}.
Thus each next parameter power acts injectively on
the successive local quotients and on the corresponding
quotients of the free local module \(M\).

Starting with degree zero for no parameters, adding
one parameter gives the two copies of the preceding
complex on subsets omitting and containing its new
index, with multiplication by \(f_i^N\) between them
up to the indicated degree sign. The cocycle and
boundary calculation for the two-term cone in
Proposition~\ref{prop:support-cover-glue} shows that
if the preceding cohomology is concentrated in its
top degree with module \(M/(f_1^N,\ldots,f_{i-1}^N)M\),
the new groups in the last two degrees are its kernel
and cokernel under \(f_i^N\); all lower ones vanish.
The kernel is zero by the preceding injectivity,
and the cokernel is the next quotient. Induction
and passage to the direct limit prove the assertion
at the closed point. All maximal ideals here are
closed \(k\)-points by the affine point argument in
Lemma~\ref{lem:projective-special-section-base-change};
vanishing at all of them implies vanishing over
\(A\) by the annihilator argument of
Lemma~\ref{lem:affine-forms-principal-acyclicity}.
\end{proof}

\begin{lemma}\label{lem:elementary-residue-change}
In the setting of Lemma~\ref{lem:supported-regular-parameters},
if \(g=C f\) for two ordered columns of \(d\) functions
with \(C\) invertible over \(A\), then
\[
          \left[\frac{\omega}{f_1,\ldots,f_d}\right]
               =\left[\frac{\det(C)\,\omega}{g_1,\ldots,g_d}\right].
\]
In particular the local class with numerator
\(df_1\wedge\cdots\wedge df_d\) is independent of the
chosen generating parameters under such a change.
\end{lemma}
\begin{proof}
The localizations of the comparison maps commute
with the common cover refinements in
Proposition~\ref{prop:support-cover-glue}. It is enough
to prove the identity near every maximal ideal, by
the restriction and localization calculations in
Lemma~\ref{lem:supported-regular-parameters}.
Over a local ring an invertible matrix is a product
of row swaps, unit row scalings and additions of
a multiple of one row to another. A unit in the
first column exists by invertibility on the residue
field; after moving it to the first row, normalize
it, clear the column, and continue on the remaining
invertible block, then clear the entries above the
pivots. These finitely many operations already hold
on a neighborhood after inverting their denominators.
A swap changes the sign of the ordered top cover
entry; a unit scaling changes its denominator by
that unit. These give the formula in those cases.

For a row addition let \(g_i=f_i+a f_j\) with
\(j\ne i\) and all other functions fixed.
Refine both covers by replacing their \(i\)-th open
by \(D(f_i g_i)\), keeping all the other members.
These opens still cover the complement of \(Z\):
where none of the other functions is nonzero,
\(g_i=f_i\) in the residue field. On the common top
intersection the difference of the two fractions is
\[
           \frac{a\,\omega}
              {f_i g_i\prod_{k\ne i,j} f_k}.
\]
It extends over the intersection omitting \(j\)
of that common cover, hence is a boundary, by using
its appropriate signed entry in the preceding degree.
This proves equality for an addition. Multiplication
of determinants gives the formula for the original
change. Finally \(dg_i=\sum_j C_{ij}df_j\bmod
(f_1,\ldots,f_d)\Omega^1\). Wedging the \(d\) equalities
thus gives numerator \(\det(C)\,df_1\wedge\cdots\wedge df_d\)
modulo that ideal on the differential module, which
suffices by the last assertion about numerators in
Lemma~\ref{lem:supported-regular-parameters}.
\end{proof}

\section{The differential classes of the diagonal and of points}

\begin{proposition}\label{prop:complete-intersection-differential-class}
Let \(Y\) be a smooth separated variety over an
algebraically closed field \(k\), and let \(D\) be a
closed subscheme. Suppose that near each point of \(D\)
its ideal is generated by \(d\ge1\) functions whose
differentials are independent there. There is a unique
class \(\delta_D\in H^d_{|D|}(Y,\Omega^d_{Y/k})\) whose
restriction to each such affine chart, with its ordered
equations \(f_1,\ldots,f_d\), is
\[
       \left[\frac{df_1\wedge\cdots\wedge df_d}
                         {f_1,\ldots,f_d}\right].
\]
On an open disjoint from \(D\) its restriction is zero.
\end{proposition}
\begin{proof}
Take a finite affine cover by the indicated opens and
by opens in the complement of \(D\). After restricting
to any affine intersection meeting \(D\), the chosen
equations from its first member of the cover still
give the ideal and their differentials remain independent
along \(D\). By Lemma~\ref{lem:supported-regular-parameters}
the supported groups below \(d\) vanish there. On an
intersection disjoint from \(D\) the vertical augmented
complex of Definition~\ref{def:supported-affine-cochains}
is an ordinary augmented cover of the entire affine,
exact by Lemma~\ref{lem:affine-forms-principal-acyclicity}.
Thus it also has zero group in degree \(d\).

On a double intersection meeting \(D\), let \(f,g\)
be the two generating lists. At any closed point of
\(D\) write \(g=Cf\) in the local ring, since the
\(f\)'s generate the ideal. On the cotangent space
there \(dg=C\,df\) since all \(f_j\) vanish. Both
lists of covectors are independent, so the square
matrix \(C\) has invertible determinant there.
Localize on a principal neighborhood to obtain the
equations and unit determinant on that neighborhood.
Lemma~\ref{lem:elementary-residue-change} gives identical
classes there. At the other closed points some generator
is a unit locally, giving zero supported cohomology.
The two restrictions agree on the full affine
intersection, by the localization and annihilator
arguments of Lemma~\ref{lem:supported-regular-parameters}.
Proposition~\ref{prop:support-cover-glue} now gives the
class and its uniqueness. Adding any other such chart
to the cover proves the claim there too.
\end{proof}

\begin{corollary}\label{cor:actual-diagonal-and-point-symbols}
Let \(X/k\) be smooth projective of pure dimension
\(d\ge1\) over an algebraically closed field. Its
diagonal in \(X\times X\) has the class of
Proposition~\ref{prop:complete-intersection-differential-class}.
Write \(\delta_X\) for its image in
\(H^d(X\times X,\Omega^d_{(X\times X)/k})\).
Near any diagonal point, if \(t_1,\ldots,t_d\)
are free smooth-chart parameters on \(X\), the defining
functions in its symbol can be taken to be
\[
            f_j=t_j\circ{\rm pr}_1-t_j\circ{\rm pr}_2 .
\]
For a closed point \(z\) of a smooth projective
\(d\)-dimensional \(W/k\) of pure dimension, let
\(\operatorname{cl}_W(z)\in H^d(W,\Omega^d_{W/k})\)
be the class obtained by forgetting support for the
reduced point subscheme. If \(Z\) is a finite set of
such points, the unique class supported on \(Z\)
with local point symbols multiplied by scalars
\(c_z\) has ordinary image
\(\sum_{z\in Z}c_z\operatorname{cl}_W(z)\).
\end{corollary}
\begin{proof}
Take a smooth chart with coordinates \(t\) free and
\(s\) dependent and equations \(P(t,s)=0\) with unit
\(P_s\) minor (the square matrix in the dependent
variables). For a pair of points in this chart the
ideal of the diagonal is generated by \(t-t',s-s'\),
since equating all coordinates and their inverted
functions identifies the two evaluations. Write the
difference \(P(t,s)-P(t',s')\) as matrices times
\(t-t'\) and \(s-s'\) by changing the variables one
at a time. The second matrix reduces on the diagonal
to \(P_s\). On its determinant open the equations
therefore express \(s-s'\) in the ideal of \(t-t'\).
The differentials \(dt_j-dt'_j\) are independent on
the diagonal, by restricting to tangent vectors in
the first factor. Hence the proposition applies,
with the asserted equations.

For a point the same calculation sets the second
chart coordinates to the constants at \(z\), so the
centered free parameters generate its ideal on a
neighborhood. On the complement the singleton class
has empty support. Enlarge that support to \(Z\)
by Proposition~\ref{prop:support-cover-glue}. Its local
restriction at any other member of \(Z\) is zero,
as seen on a neighborhood disjoint from \(z\).
Choose affine neighborhoods of the finitely many
points each excluding all the others, with their
centered parameters, and complete the cover by
affines in the complement. The vanishing below \(d\)
in Lemma~\ref{lem:supported-regular-parameters} and
Proposition~\ref{prop:support-cover-glue} give uniqueness
for the specified list on this cover. The sum of
the enlarged singleton classes with coefficients
\(c_z\) has that list. Forgetting support gives the
stated linear combination, also for the empty set
where the augmented cover complexes are exact.
\end{proof}

\section{Pulling the diagonal class back to a correspondence}

\begin{theorem}\label{thm:scaled-graph-local-symbol-formula}
Let \(X,W\) be smooth projective of pure dimension
\(d\ge1\) over an algebraically closed field \(k\) of
characteristic zero. Let \(u,c:W\to X\) have finite
set \(Z\) of coincident closed points \(z\), meaning
\(u(z)=c(z)\), and suppose \(du_z-dc_z\) is invertible
there. With \(S_T,\delta_X\) from
Corollaries~\ref{cor:scaled-diagonal-weight-annihilation}
and~\ref{cor:actual-diagonal-and-point-symbols}, one has
in \(H^d(W,\Omega^d_{W/k})[T]\)
\[
       (u,c)^* S_T\delta_X
       =\sum_{z\in Z}
             \frac{\det(Tdu_z-dc_z)}{\det(du_z-dc_z)}
                     \operatorname{cl}_W(z).
\]
The ratio is formed using any bases of \(T_zW\) and
\(T_{u(z)}X\); the two changes of basis cancel in it.
If \(dc_z\) is invertible and \(A_z=du_z(dc_z)^{-1}\),
this ratio is \(\det(I-T A_z)/\det(I-A_z)\).
\end{theorem}
\begin{proof}
The first-factor degree decomposition in
Proposition~\ref{prop:coherent-forms-product-kunneth}
is a decomposition of the differential module itself.
The multiplication by \(T^j\) in degree \(j\) is
therefore a polynomial of module homomorphisms
commuting with restrictions and denominators. Apply
it before forgetting the support of the diagonal
class in Corollary~\ref{cor:actual-diagonal-and-point-symbols}.
The pair map \((u,c)\) pulls this supported class
back to support \(Z\) by
Proposition~\ref{prop:support-cover-glue}. Indeed
the inverse image of the diagonal is closed, and
its set outside the finitely many indicated points
is empty since every nonempty locally closed set has
a closed point by Lemma~\ref{lem:projective-special-section-base-change}'s
affine calculation.

At \(z\) use free target parameters \(t_1,\ldots,t_d\)
near its common image. Put
\(g_j=t_j\circ u-t_j\circ c\).
The pullback of the support symbol in those coordinates
is represented near \(z\) by
\[
   \left[
      \frac{\bigwedge_{j=1}^{d}
              (T\,d(t_j\circ u)-d(t_j\circ c))}
           {g_1,\ldots,g_d}
   \right].
\]
In fact, scaling the first-factor form degree multiplies
each occurrence of a differential from that factor
by \(T\) under wedge expansion. The denominators
are functions, hence are only pulled back. The cover
on the complement by the nonvanishing of the original
coordinate differences pulls back to the corresponding
principal opens for \(g_j\), giving this precise
description in Definition~\ref{def:supported-affine-cochains}.

The differentials of the \(g_j\) at \(z\) form a
basis because \(du_z-dc_z\) is invertible. Thus the
\(g_j\) generate the maximal ideal of the local ring:
their images span it modulo its square, and one uses
the finite-module determinant argument of
Lemma~\ref{lem:finite-fiber-local-algebras}. They generate
the reduced point ideal on a neighborhood, by clearing
denominators in expressions for finite local chart
generators as in Lemma~\ref{lem:jacobian-chart-flat-descent}.
One may shrink the affine neighborhood to exclude
other coincidences. Relative to the basis
\(\bigwedge_j dg_j\), the numerator at that point
is the scalar \(\det(Tdu_z-dc_z)/\det(du_z-dc_z)\):
write both wedges in any ordered local cotangent
basis, using the target rows \(dt_j\). By
Lemma~\ref{lem:supported-regular-parameters} the value
of the numerator modulo the point ideal suffices
coefficient by coefficient. By
Lemma~\ref{lem:elementary-residue-change} the resulting
unscaled point symbol is the intrinsic one of
Corollary~\ref{cor:actual-diagonal-and-point-symbols}.
That corollary now identifies the ordinary image
of the global supported class with the stated sum.
When \(dc_z\) is invertible, factor it on the right
in both determinants. The two common factors and the
two signs \((-1)^d\) cancel, giving the final expression.
\end{proof}

\begin{corollary}\label{cor:coherent-relation-among-fixed-points}
For a finite list \(u_i\) with a common \(c:W\to X\)
satisfying Theorem~\ref{thm:scaled-graph-local-symbol-formula},
let \(Z_i\) be their respective coincidence sets.
If \(\sum_i\beta_i u_i^*\) vanishes on
\(H^q(X,\Omega^h_{X/k})\) for \(0\le q\le d\) with
\(0\le h\le d\), then
\[
  \sum_i\beta_i\sum_{z\in Z_i}
     [T^h]\frac{\det(Tdu_{i,z}-dc_z)}
                    {\det(du_{i,z}-dc_z)}
       \operatorname{cl}_W(z)=0
                    \quad\hbox{in } H^d(W,\Omega^d_{W/k}).
\]
In particular this applies to the regular correspondences
and torus slots of
Corollary~\ref{cor:actual-holomorphic-local-sum} under
the vanishing assumption of
Corollary~\ref{cor:coherent-regular-left-multiple}.
\end{corollary}
\begin{proof}
Apply Corollary~\ref{cor:scaled-diagonal-weight-annihilation}
to \(\delta_X\) and use the theorem on each term.
For the specified Hecke covers \(c=\pi_m\) is étale
by Theorem~\ref{thm:actual-complex-finite-etale} and
the two dimensions agree. Their coincidences are finite
by Theorem~\ref{thm:complex-hecke-coincidence-count},
and \(\det(I-A_z)\ne0\) by
Theorem~\ref{thm:actual-complex-tangent-determinant}.
These verify the required transversality for each
\(u_{g,2\boldsymbol\alpha}\). The stated pullback
vanishing is precisely the conclusion of
Corollary~\ref{cor:coherent-regular-left-multiple}.
\end{proof}

\section{Smooth forms representing the classes with support}

\begin{definition}\label{def:smooth-relative-dolbeault}
Let \(M\) be the complex manifold of a smooth projective
complex variety or any open subset of it. For an open
\(U\subseteq M\) let \(\mathcal A^{a,j}(U)\) be the
smooth complex differential forms of holomorphic degree
\(a\) and antiholomorphic degree \(j\). In local
holomorphic coordinates their basis factors are
\(dz_i,d\bar z_i\) in the specified numbers; put
these spaces equal to zero for negative degrees.
Use \(\bar\partial\) of
Definition~\ref{def:line-metric-conventions} with its
signed product rule and \(\bar\partial^2=0\).
For closed \(Z\subseteq M\) define the relative smooth
complex in degree \(r\) by
\[
  \mathcal A^{a,r}(M)\oplus\mathcal A^{a,r-1}(M\setminus Z),
 \qquad (A,B)\longmapsto
      \bigl(\bar\partial A,\ A|_{M\setminus Z}-\bar\partial B\bigr).
\]
Projection to the first component forgets the support.
For a larger closed \(Z'\) restrict the second component
to \(M\setminus Z'\).
\end{definition}

\begin{lemma}\label{lem:supported-forms-smooth-realization}
For a smooth projective complex \(W\), \(M=W(\C)\),
there is a natural linear map from \(H^r(W,\Omega^a)\)
to the cohomology of \((\mathcal A^{a,\bullet}(M),\bar\partial)\).
From \(H_Z^r(W,\Omega^a)\) for an algebraic closed support
there is likewise a map to the relative smooth cohomology
of Definition~\ref{def:smooth-relative-dolbeault}.
These maps commute with restriction, enlargement and
forgetting of support. On analytic open neighborhoods
a restricted Čech representative may be computed on
any finite analytic refinement by the same construction.
\end{lemma}
\begin{proof}
Every finite open cover on any such complex open
manifold has a smooth partition of unity. Here are
details for the covers used below. For an open subset
of the compact projective manifold use the sets at
distance at least \(1/j\) from the complement, in a
metric induced by a projective embedding as in
Definition~\ref{def:fixed-real-torus-evaluation}.
These compact sets exhaust the open and each lies in
the interior of the next. For the whole compact
manifold one compact set suffices. Cover the successive
compact differences \(K_j\setminus\operatorname{int}K_{j-1}\)
by finitely many smaller coordinate balls, each with
its closed larger ball inside a cover member and in
\(\operatorname{int}K_{j+1}\setminus K_{j-2}\), taking
the first missing lower sets empty. Center a smooth
nonnegative bump in each larger ball, positive on the
smaller one. In scaled real coordinates one can use
\(\exp(-1/(1-|x|^2))\) for \(|x|<1\) and zero outside;
all derivatives extend by zero because an exponential
of \(-1/t\) tends to zero faster than every power of
\(t>0\). The collection is locally finite by the two
compact cutoffs. Its sum is positive. Divide by that
sum and collect the functions assigned to each cover
member. Their supports are closed and contained in
the assigned members by local finiteness.

On any such cover the augmented Čech complex for
smooth forms of fixed bidegree is exact. If the
partition is \(\rho_j\), the contraction of an
alternating array has entry
\(\sum_j\rho_j c_{j,I}\) on the intersection indexed
by \(I\), with summands extended by zero there
outside their \(j\)-th opens. The support condition
makes these extensions smooth. In \(\delta H+H\delta\)
the deletions other than the inserted index cancel
and the remaining sum is multiplication by
\(\sum_j\rho_j=1\), including the augmented term.

In the ordinary cover cochains add a second index
\(j\) for smooth forms of type \((a,j)\) and use
\(\bar\partial\) with the total sign. Holomorphic
\(a\)-forms are \(\bar\partial\)-closed, so inserting
an algebraic array at \(j=0\) gives a cochain map,
by the analytic coordinate pullback formulas for
\(df\). The augmentation from the single global
smooth complex into this total complex gives an
isomorphism on cohomology, by the exact augmented
rows and finite filtration argument of
Proposition~\ref{prop:support-cover-glue}. Invert this
isomorphism on cohomology to obtain the ordinary map.

For support use a third index \(j\) in
Definition~\ref{def:supported-affine-cochains}, with
smooth forms of type \((a,j)\) on each intersection
and \(\bar\partial\) carrying the sign \((-1)^{i+b}\)
for horizontal cover degree \(i\) and vertical
augmented degree \(b\). The same insertion of
holomorphic forms is a cochain map. In the \(U\)
cover index the augmentation on \(M\) at \(b=0\)
and on \(V_J\) at \(b>0\) gives isomorphisms on
total cohomology by the partition contraction.
There remains the complex consisting of global
smooth forms at \(b=0\) and smooth cover cochains
of the complement at \(b\ge1\). Augment the latter
by global forms on the complement, placed at \(b=1\).
This again induces an isomorphism: on the subcomplex
with \(b\ge1\) it is the ordinary smooth cover
comparison just proved, and on the quotient at \(b=0\)
it is the identity. The conclusion on the full
complex follows by lifting cocycles in the two
terms, or by the exact-cone argument of
Proposition~\ref{prop:support-cover-glue} filtered
by these two parts. At \(b=0,1\) the resulting
differential is exactly that of
Definition~\ref{def:smooth-relative-dolbeault}.
All comparison arrows being inverted are restrictions
of global forms onto the cover, and commute with
refinements, restriction to an open and restrictions
on the complement. The partitions prove these arrows
are isomorphisms on cohomology and do not enter their
definition. This also proves compatibility with
forgetting support, since projection to \(b=0\)
commutes with the comparisons. The same smooth
comparisons apply to analytic covers and their
refinements, proving the local assertion.
\end{proof}

\begin{lemma}\label{lem:smooth-relative-cutoff-excision}
Let \(Z\) be compact in \(M\), and \(U\) an open
neighborhood of \(Z\). Restriction from \(M\) to \(U\)
induces an isomorphism on the relative smooth
cohomology of Definition~\ref{def:smooth-relative-dolbeault}.
For a relative cocycle on \(U\) represented by
\((A,0)\) with \(A\) supported in \(Z\), its inverse
is represented by the extension of \((A,0)\) by zero.
\end{lemma}
\begin{proof}
Choose a smooth \(\chi\) supported in \(U\), one on a
neighborhood of \(Z\), by the partition construction
in Lemma~\ref{lem:supported-forms-smooth-realization}
for \(U,M\setminus Z\), first replacing \(U\) by
a relatively compact neighborhood inside it if needed.
On arbitrary pairs on \(U\) send \((A,B)\) to the
extension by zero of
\[
                  (\chi A+\bar\partial\chi\wedge B,\ \chi B).
\]
The term involving \(\bar\partial\chi\) vanishes
near \(Z\), so it is defined smoothly over \(U\).
The formula is a cochain map by the signed product
rule and \(\bar\partial^2=0\). On either side the
identity minus the composite with restriction is
the differential times the homotopy plus that
homotopy times the differential, where the homotopy
sends \((A,B)\) to \(((1-\chi)B,0)\). Indeed
\((1-\chi)B\) extends by zero near \(Z\), and the
two first-component derivatives of \(B\) cancel,
leaving \((1-\chi)A-\bar\partial\chi\wedge B\);
the second component is \((1-\chi)B\).
These formulas prove both assertions.
\end{proof}

\section{A degree functional for classes of points}

\begin{lemma}\label{lem:integration-of-top-hodge-forms}
Let \(W\) be smooth and projective over \(\C\), of pure
dimension \(d\ge1\), and put \(M=W(\C)\).
Integration in the complex orientation gives a linear
functional on \(H^d(W,\Omega^d_{W/\C})\) by the map of
Lemma~\ref{lem:supported-forms-smooth-realization}.
If a supported class has smooth representative \((A,0)\)
on an analytic neighborhood of its compact support, its
image under this functional is the integral of \(A\) on
the neighborhood.
\end{lemma}
\begin{proof}
Use the coordinate order
\((\Re z_1,\Im z_1,\ldots,\Re z_d,\Im z_d)\).
A holomorphic coordinate change with complex derivative
\(D\) has real determinant \(\det D\,\overline{\det D}>0\):
over \(\C\) the complexified real derivative in the
coordinates \(dz_j,d\bar z_j\) has the two conjugate
blocks. The set \(M\) is compact, being closed in complex
projective space, which is the continuous image of a
compact unit sphere. Use a finite coordinate cover with
a smooth partition of unity as in
Lemma~\ref{lem:supported-forms-smooth-realization}.
Multiply a top form by the partition functions, integrate
in the resulting oriented coordinates and add. For two
such choices multiply the two partitions. On their
intersections the local integrals agree by the change of
variables proved in Lemma~\ref{lem:regular-fiber-density},
and summing first over either partition gives independence.

If \(B\) has bidegree \((d,d-1)\), then
\(\dd B=\bar\partial B\), since the component of
holomorphic degree \(d+1\) vanishes. Writing
\(B=\sum_j \rho_j B\) as above reduces the integral of
\(\dd B\) to the differentials of compactly supported
forms in coordinate domains. Extend the coefficients by
zero in the real coordinates. Each integral of a first
partial derivative vanishes by one-variable integration
and Lemma~\ref{lem:product-integration}. Thus integration
descends to the smooth cohomology, giving the asserted
functional by composition. For a relative cocycle
\((A,0)\) the restriction of \(A\) to the complement
of the support is zero. Lemma~\ref{lem:smooth-relative-cutoff-excision}
identifies its global image with the extension by zero,
whose integral is the same by coordinate independence.
\end{proof}

\begin{proposition}\label{prop:integral-of-algebraic-point-symbol}
For \(W,d\) of Lemma~\ref{lem:integration-of-top-hodge-forms}
and any \(z\in W(\C)\), the functional there takes the
value
\[
                    c_d=(-1)^{d(d+1)/2}(2\pi i)^d
\]
on the ordinary class \(\operatorname{cl}_W(z)\) of
Corollary~\ref{cor:actual-diagonal-and-point-symbols}.
Consequently there is a linear functional
\(\deg_W:H^d(W,\Omega^d_{W/\C})\to\C\) with
\(\deg_W(\operatorname{cl}_W(z))=1\) on every component.
\end{proposition}
\begin{proof}
Choose the algebraic coordinates for the point symbol
from Corollary~\ref{cor:actual-diagonal-and-point-symbols},
translated to vanish at \(z\). They give local holomorphic
coordinates by the inverse construction in
Lemma~\ref{lem:regular-fiber-density}: the inverse has
complex-linear derivative everywhere. Use a coordinate
polydisc \(P=\prod_{i=1}^d\{|z_i|<R_i\}\) inside this
neighborhood. Write \(\eta=dz_1\wedge\cdots\wedge dz_d\).
On the ordered cover by \(z_i\ne0\) of the punctured polydisc
the local representative is \(\eta/\prod_i z_i\) in
augmented degree \(d\). This calculation on an analytic
restriction represents the smooth image by
Lemma~\ref{lem:supported-forms-smooth-realization}.
Enlarge the support to the compact polydisc
\(K=\prod_i\{|z_i|\le r_i\}\), with \(0<r_i<R_i\).
The new complement is covered in order by the opens
\(D_i=\{|z_i|>r_i\}\) in \(P\). Restriction to their
intersections represents the enlarged class.

On the \(i\)-th disc choose a smooth radial function \(f_i\)
equal to zero near zero and equal to one for
\(|z_i|\ge r_i\). Such a cutoff is obtained by integrating
a nonnegative smooth bump of positive integral supported
inside \((0,r_i^2)\), normalizing its integral to one,
and using \(|z_i|^2\) as argument. The function
\(e_i=f_i/z_i\), extended by zero near the center, is
smooth. Consider smooth \((0,j_i)\)-forms on this disc
in degree \(b_i=0\) and on its annulus in degree \(b_i=1\).
Use the differential given by restriction in \(b_i\) and
\((-1)^{b_i}\bar\partial\) in \(j_i\). The degree-one
cochains given respectively by \(1/z_i\) on the annulus
and by \(-\bar\partial e_i\) on the disc are cocycles.
Their difference is the total differential of \(e_i\)
on the disc, since \(e_i=1/z_i\) on the annulus.

In the tensor product of the complexes over \(\C\),
the differential on factor \(i\) carries sign
\((-1)^{\sum_{k<i}(b_k+j_k)}\).
There is a cochain map to the augmented cover complex of
smooth forms on \(P\) and the \(D_i\), as follows. The set
\(J=\{i:b_i=1\}\) labels the intersection. Pull back the
antiholomorphic factors in their given order to this
intersection, wedge them, then wedge with \(\eta\) on the
right. Multiply by \((-1)^{\sum_{k<l}j_k b_l}\).
For a restriction in factor \(i\) the change in this last
exponent is \(\sum_{k<i}j_k\), so the resulting sign is
the alternating deletion sign \((-1)^{\sum_{k<i} b_k}\)
of the cover. For a derivative in factor \(i\) the change
is \(\sum_{l>i}b_l\). Combining the tensor sign and
the factor sign \((-1)^{b_i}\) with this gives
\((-1)^{|J|+\sum_{k<i}j_k}\), exactly the total
\(\bar\partial\) and wedge signs. This proves the claim.

Tensor products of the cohomologous cocycles give
cohomologous results. Replacing one factor by its
difference adds the differential of the same tensor
with its primitive in that factor, up to the tensor
sign. At both ends the extra sign of the map is one.
Thus the original class on the enlarged support has the
smooth representative \((A,0)\), where
\[
             A=(-1)^d
               (\bar\partial e_1\wedge\cdots\wedge
                    \bar\partial e_d)\wedge\eta .
\]
Indeed \(A\) vanishes on every \(D_i\) and the augmented
cover map from relative pairs in
Lemma~\ref{lem:supported-forms-smooth-realization}
sends this pair to its degree-zero cover component.
Enlarging support preserves the ordinary class by the
same lemma. By Lemma~\ref{lem:integration-of-top-hodge-forms}
it remains to integrate \(A\) on \(P\).

On the \(i\)-th disc the two-form
\(\bar\partial e_i\wedge dz_i=\dd(e_i\,dz_i)\) has
compact support. In polar coordinates write
\(e_i\,dz_i=a\,dt+b\,d\theta\) with radius \(t\).
Its differential has coefficient
\(\partial_t b-\partial_\theta a\) in \(dt\wedge d\theta\).
The second term integrates to zero by periodicity and
the first gives the outer circle integral, since \(e_i\)
vanishes near zero. On any circle of radius between
\(r_i\) and \(R_i\), the one-form is \(dz_i/z_i=i\,d\theta\).
This gives \(2\pi i\) for the disc integral.
Moving each \(dz_i\) next to \(\bar\partial e_i\) costs
the sign \((-1)^{d(d-1)/2}\). Product integration by
Lemma~\ref{lem:product-integration} now gives \(c_d\)
in the complex orientation. Divide the functional by
this nonzero constant.
\end{proof}

\section{Detection of the prescribed form degrees}

\begin{theorem}\label{thm:numerical-half-holomorphic-lefschetz}
Let \(X,W\) be smooth projective complex varieties of the
same pure dimension \(d\ge1\). Let \(c,u_i:W\to X\) be
morphisms for \(1\le i\le m\), let \(\beta_i\in\C\), and
fix \(0\le h\le d\). Suppose that
\(Z_i=\{z\in W(\C):u_i(z)=c(z)\}\) is finite for each
\(i\), and that \(du_{i,z}-dc_z\) is invertible for every
\(z\in Z_i\). If
\[
   \sum_{i=1}^m\beta_i u_i^*
      :H^q(X,\Omega^h_{X/\C})\longrightarrow
       H^q(W,\Omega^h_{W/\C})
\]
is zero for each \(0\le q\le d\), then
\[
  \sum_{i=1}^m\beta_i
     \sum_{z\in Z_i}
          [T^h]\frac{\det(Tdu_{i,z}-dc_z)}
                         {\det(du_{i,z}-dc_z)}
             =0 .
\]
An empty coincidence set contributes zero. If \(dc_z\)
is invertible the fraction is
\(\det(I-T A_{i,z})/\det(I-A_{i,z})\) for
\(A_{i,z}=du_{i,z}(dc_z)^{-1}\).
\end{theorem}
\begin{proof}
Corollary~\ref{cor:coherent-relation-among-fixed-points}
gives a zero relation with the same coefficients on the
classes of the points in \(H^d(W,\Omega^d_{W/\C})\).
Apply the single linear functional \(\deg_W\) of
Proposition~\ref{prop:integral-of-algebraic-point-symbol}
to this finite relation. Every point on every component
has value one. The determinant expression for invertible
\(dc_z\) follows by right composition with \((dc_z)^{-1}\)
in numerator and denominator, as in
Theorem~\ref{thm:scaled-graph-local-symbol-formula}.
\end{proof}

\begin{theorem}\label{thm:marked-weights-in-actual-coherent-cohomology}
Use the approximating form, the central correction
\(\theta_{w,r}=c_w^{-r}t_{w,r}\) for \(0\le r\le n\), and
the integers \(h_j=\mu+k_{v_0,j}\) of
Theorem~\ref{thm:actual-marked-weight-occurrence}.
Take \(X=\mathcal X(C,\nu)_\C\) at its neat level and
put \(d=\dim X\). Let \(T_{w,r}=T_{\tau_{w,r}}\) be the
operators of Theorem~\ref{thm:actual-hecke-action-on-forms},
with \(T_{w,r}=0\) and \(\theta_{w,r}=0\) for \(r>n\).
For each \(1\le j\le n\) one has \(0\le h_j\le d\).
For every finite complex linear combination \(F\) of
ordered words in the safe generators, constants included,
\[
 F(\theta)\ne0\ \Longrightarrow\
 F((T_{w,r})_{w,r})\ne0
   \quad\hbox{on }\quad
       \bigoplus_{q=0}^d H^q(X,\Omega^{h_j}_{X/\C}).
\]
\end{theorem}
\begin{proof}
Nonnegativity was proved in
Theorem~\ref{thm:actual-marked-weight-occurrence}.
For \(F=1\), Corollary~\ref{cor:actual-holomorphic-local-sum}
gives a nonzero fixed-point sum at each \(h_j\). Its
numerators use the \(h_j\)-th exterior power of an
endomorphism on a \(d\)-dimensional cotangent space.
That power is zero for \(h_j>d\), proving the upper
bound. The dimension is positive since the fiber
contains \(\nu\ge1\) copies of the positive-dimensional
torus in Theorem~\ref{thm:analytic-framed-torus-family}.

Divide the given \(F\) by \(F(\theta)\), so that
\(F(\theta)=1\). Suppose the operator is zero on the
displayed direct sum for some \(j\). By
Corollary~\ref{cor:actual-holomorphic-local-sum}, choose
a regular left multiple \(H=B F(\tau)\) and a tuple
\(\boldsymbol\alpha\) for which its sum at that \(h_j\)
is nonzero. Choose a common level sufficiently fine for
Corollary~\ref{cor:coherent-regular-left-multiple}
and the correspondence calculation. Further refinement
multiplies the local sum by the degree of the new cover
and the normalizing index by the same degree, by the
coincidence count in
Theorem~\ref{thm:complex-hecke-coincidence-count}.
The asserted vanishing of \(F(T)\) implies the
vanishing of the linear combination of the pullbacks
\(u_{g,2\boldsymbol\alpha}^*\), with coefficients \(H(gC)\),
from every \(H^q(X,\Omega^{h_j})\), \(0\le q\le d\),
by Corollary~\ref{cor:coherent-regular-left-multiple}.

Use these maps in
Theorem~\ref{thm:numerical-half-holomorphic-lefschetz}
with \(c=\pi_m\). The dimension, finiteness and
transversality hypotheses were checked in
Corollary~\ref{cor:coherent-relation-among-fixed-points}.
The resulting zero sum is precisely the numerator of
the normalized sum for \(H,\boldsymbol\alpha,h_j\):
the coefficient of \(\det(I-TA_z)\) is
\((-1)^{h_j}\operatorname{tr}(\bigwedge^{h_j} A_z^\vee)\)
by the exterior-power calculation in
Corollary~\ref{cor:actual-holomorphic-local-sum}.
The normalizing index \([C:C_m]\) is a positive integer.
This contradicts the choice of the sum.
\end{proof}

\section{Simultaneous quotients at a detected degree}

\begin{lemma}\label{lem:spherical-commutativity-on-marked-tower}
The elements \(\tau_{w,r}\) of
Definition~\ref{def:marked-weight-test} commute pairwise
in \(\mathcal H^+\). Their operators on the coherent
groups of Theorem~\ref{thm:actual-hecke-action-on-forms}
commute over the field \(E\) of the period model.
\end{lemma}
\begin{proof}
By Theorem~\ref{thm:actual-central-test-vector} each sum
is a central translate of the image, under \(b_w\),
of the right cosets of an integral spherical double
coset for \(K_w={\rm GL}_n(O_w)\).
The central factors commute with all representatives.
If \(\{w,\bar w\}\) and \(\{v,\bar v\}\) are disjoint,
the two noncentral factors have disjoint support in
the product of the local groups and commute term by
term. For a single pair, use \(a\mapsto b_w(a)\) from
Proposition~\ref{prop:local-translation-frame-extension}.
At a safe pair its inverse image of \(C\) is \(K_w\):
membership forces integrality and integral invertibility
at \(w\), and \(b_w(K_w)\subseteq C\) by
Lemma~\ref{lem:safe-operators-away-from-markers}.
The description of \(u_w\) and its conjugate inverse
transpose at the mate shows that
\(b_{\bar w}(a)=b_w(\bar a^{-\mathsf T})\).
This group isomorphism carries \(K_{\bar w}\) onto
\(K_w\). Hence both unshifted sums on this pair are
images of local spherical sums for the same \(K_w\).

Multiplication of such sums in the full group of
similitudes agrees with multiplication of those local
sums followed by \(b_w\) on representatives. Indeed,
sums over pairs of right cosets group according to
\(b_w(a)b_w(b)C=b_w(ab)C\), and equality of the final
cosets is precisely equality modulo \(K_w\) at the
chosen member. The proof of
Lemma~\ref{lem:row-spherical-hecke} gives commutativity
of this local convolution by transpose on each double
coset. For the mass-one subgroup its convolution
coefficients are exactly these coset counts, since
integration over a right-factor coset counts it once.
Restoring the central translates gives the asserted
equality of products in \(\mathcal H^+\), including
the zero generators. Their cohomological actions and
their descent to \(E\) follow from
Theorem~\ref{thm:actual-hecke-action-on-forms}, since
the coset functions have integral coefficients.
\end{proof}

\begin{theorem}\label{thm:actual-coherent-hecke-eigenquotients}
Keep the data of
Theorem~\ref{thm:marked-weights-in-actual-coherent-cohomology}.
View \(E\subset\C\) by the embedding defining \(X\), and
denote its projective model over \(E\) by \(X_E\). The
field \(L=E((\theta_{w,r})_{w,r})\) generated by all the
tested eigenvalues is finite over \(E\).
For every \(h=h_j\) set
\[
                  D_h=\bigoplus_{q=0}^d
                            H^q(X_E,\Omega^h_{X_E/E}).
\]
There is a nonzero quotient \(Q_h\) of \(D_h\) with an
\(L\)-vector space structure for which \(T_{w,r}\)
acts by \(\theta_{w,r}\). In particular,
\([L:E]\le\dim_E D_h\) for every such \(h\). For any
field embedding \(\iota:L\to k\), the simultaneous
quotient of \(D_h\otimes_{E,\iota}k\) with eigenvalues
\(\iota(\theta_{w,r})\) is nonzero. This quotient
has a nonzero summand in at least one cohomological
degree \(0\le q\le d\).
\end{theorem}
\begin{proof}
The space \(D_h\) is finite dimensional by
Proposition~\ref{prop:finite-projective-coherent-cohomology}.
Proposition~\ref{prop:affine-cover-hodge-functoriality}
identifies its complex scalar extension with the sum in
Theorem~\ref{thm:marked-weights-in-actual-coherent-cohomology},
and Theorem~\ref{thm:actual-hecke-action-on-forms}
identifies the operators there. By
Lemma~\ref{lem:spherical-commutativity-on-marked-tower}
their image generates a commutative unital \(E\)-algebra
\(A_h\subseteq {\rm End}_E(D_h)\). If a polynomial
over \(E\) in finitely many generators vanishes on
\(D_h\), it still vanishes over \(\C\), so the detected
degree theorem makes its value at \(\theta\) zero.
Evaluation therefore gives a unital \(E\)-algebra map
\(A_h\to\C\). Its image is \(E[(\theta_{w,r})_{w,r}]\),
independent of \(h\) and finite dimensional over \(E\).
Multiplication by any nonzero member of this image is
injective, hence surjective as an endomorphism of a
finite-dimensional space. The image is a field, equal
to \(L\).

Let \(I_h\) be the kernel of this evaluation and put
\(Q_h=D_h/I_hD_h\). Its action by \(A_h/I_h=L\)
has the stated scalar values. Suppose \(I_hD_h=D_h\).
A basis \(v_1,\ldots,v_a\) of \(D_h\) then satisfies
\(v_i=\sum_j b_{ij}v_j\) for \(b_{ij}\in I_h\), by
expressing the elements of \(I_hD_h\) in that basis.
The space is nonzero by the detected degree theorem
applied to \(F=1\). Multiplying the matrix equation by
the adjugate of \((\delta_{ij}-b_{ij})\) shows that its
determinant kills every \(v_i\). It is therefore zero in
the concrete endomorphism algebra \(A_h\), but its
image in \(L\) is one. This proves \(Q_h\ne0\) and
the degree bound follows from
\(\dim_E Q_h=[L:E]\dim_L Q_h\).

After extension to \(k\), the map
\(D_h\otimes_E k\to Q_h\otimes_{L,\iota}k\),
\(v\otimes t\mapsto[v]\otimes t\), is onto since
the classes \([v]\) span \(Q_h\). Its target is
nonzero by extending an \(L\)-basis. Thus the quotient
by the sum of all
\((T_{w,r}-\iota(\theta_{w,r})) (D_h\otimes_E k)\)
is nonzero as well. The operators preserve each
cohomological degree by
Theorem~\ref{thm:actual-hecke-action-on-forms}.
The sum of their images after subtracting the scalars
therefore splits by those degrees, proving the last
assertion.
\end{proof}

\section{Étale cochains and geometric stalks}

\begin{definition}\label{def:small-etale-point-conventions}
In each residue characteristic, including zero, fix an
algebraically closed field \(K\) large enough to contain
an isomorphic copy of every residue field of the schemes
and charts in question. One can take the algebraic closure
of a purely transcendental field on a set larger than the
cardinalities of the base fields and rings in use. An
embedding of a transcendence basis extends to an embedding
of the algebraic extension by adjoining roots, as in
Lemma~\ref{lem:place-fields}. Use all absolute morphisms
\(\operatorname{Spec}K\to X\) in that characteristic.
Write \(P(X)\) for the union of these sets of points of \(X\).

Use the small étale site of \(X\). Its objects are étale
schemes \(U\to X\), its arrows are morphisms over \(X\),
and covering families are jointly surjective étale maps.
The calculations may use separated objects of finite
presentation and affine open refinements. Here étale
charts are locally given over a base ring by the same
number of equations and polynomial variables, with
the square Jacobian determinant inverted, as in
Proposition~\ref{prop:fine-level-torsion-etale-charts}.
Products over the base and compositions have such charts
by scalar extension and by the block triangular Jacobian.
The diagonal of an étale chart is open near each of its
points. Indeed, subtract its equations in two sets of
variables. The matrix of the differences equals the
Jacobian on the diagonal, so where its determinant is
invertible all coordinate differences vanish. This
applies after the principal localizations of a chart.
It shows also that a map between étale objects is étale,
by factoring it as its graph and a projection.

An abelian sheaf \(F\) assigns abelian groups to these
objects, contravariantly, with unique gluing of matching
sections on covers. For \(x\in P(X)\) let \(F_x\) be the
direct limit of \(F(U)\) over the pointed neighborhoods
\((U,u)\), \(u\in P(U)\) over \(x\), with restriction to
smaller neighborhoods. For an abelian group \(\Lambda\)
write \(\underline\Lambda_X\) for the sheaf of locally
constant sections with values in \(\Lambda\), obtained
by gluing constant sections. All sheaves below are
abelian.
\end{definition}

\begin{lemma}\label{lem:etale-germs-and-point-products}
The pointed neighborhoods in
Definition~\ref{def:small-etale-point-conventions} are
filtered in the direction of restriction. A sheaf map is
injective, surjective as a sheaf map, or an isomorphism
precisely when all its stalk maps have that property.
Kernels and cokernels of sheaf maps have stalks equal to
the respective kernels and cokernels on stalks.
For every \(F\) the rule
\[
           G_X(F)(U)=\prod_{u\in P(U)} F_{x(u)},
           \qquad x(u)\in P(X)\ \text{the image of }u ,
\]
with restriction by reindexing, defines a sheaf, and the
map \(F\to G_X(F)\) taking germs is injective. The functor
\(G_X\) takes a short exact sequence to a sequence exact
on sections on every \(U\). The functor
\(Q_X F=\operatorname{coker}(F\to G_X(F))\) on sheaves
is exact.
\end{lemma}
\begin{proof}
A fiber product with its given pair of lifts refines
two neighborhoods. The equalizer of two arrows agreeing
at a lift contains a neighborhood of it by openness of
the diagonal in Definition~\ref{def:small-etale-point-conventions}.
An étale cover of \(U\) is jointly onto \(P(U)\):
over a chosen point of \(U\) some fiber is nonempty,
and after extension to the algebraically closed field
it has a rational point by the affine calculation of
Lemma~\ref{lem:projective-special-section-base-change}.
Conversely each ordinary point admits a lift from the
chosen field in its characteristic, by its size.
A family of neighborhoods of all the chosen lifts
therefore covers \(U\). Equality of germs gives local
equality on such a cover and thus equality of sections.
Lifting a germ of a section gives a local lift after
another restriction, so surjectivity on stalks is
equivalent to local surjectivity.

For clarity, kernels are computed on sections. For a
cokernel take the quotient presheaf and sheafify. For any
abelian presheaf, form a new presheaf from matching
families on covers, identifying two families when their
restrictions agree on a common refinement. This is
separated because local agreements can be refined on
each member and the resulting covers composed. After
repeating the construction it is a sheaf. Indeed
matching families of sections of the separated
presheaf glue by choosing representatives on covers
of each member. Their restrictions agree on overlaps
by separatedness and the local agreements, so their
union is a matching family on the composite cover.
Compatible maps to any sheaf extend uniquely by
gluing at both steps. Neither step changes the stalks:
a matching family has a representative on a member
through each lift, and equality is detected after
refinement near that lift. Taking stalks commutes
with kernels and ordinary quotients on groups since
a relation among finitely many germs holds on one
common pointed neighborhood. This proves the exactness
tests, including the identification of the stalk of
\(\underline\Lambda_X\) with \(\Lambda\).

A matching family of elements of \(G_X(F)\) gives at
each \(u\in P(U)\) a value independent of a chosen
lift to a cover member. Any two lifts pair to a point
of their fiber product, where the prescribed values
agree. Conversely these values give the unique glued
element. The preceding germ test proves injectivity
of \(F\to G_X(F)\). Products of exact sequences of
abelian groups are exact, by choosing lifts in each
factor, proving the assertion about \(G_X\).

Take stalks of a short exact sequence of sheaves and
of its image under \(G_X\). Both sequences are exact
and the vertical germ maps into \(G_X\) are injective.
Their quotients form a short exact sequence. To check
the middle, lift a representative whose image belongs
to the last germ subgroup, subtract a lift in the
middle germ subgroup, then lift the difference in
the first product group. For injectivity the same
diagram identifies the intersection of the first
product group with the middle germ subgroup as the
first germ subgroup. Surjectivity follows from
surjectivity of the last product map. The stalk test
gives exactness of \(Q_X\).
\end{proof}

\begin{definition}\label{def:etale-godement-cochains}
For \(b\ge0\) write
\[
  C^b_{\rm et}(X,F)=G_X Q_X^b F(X),\qquad
  d^b:\ G_X Q_X^b F\longrightarrow Q_X^{b+1} F
                  \longrightarrow G_X Q_X^{b+1} F .
\]
The first arrow is the quotient and the second is the
germ inclusion. The complex of sheaves begins with
\(F\to G_X F\). It is exact there and in every subsequent
degree, by the quotient definition and injectivity in
Lemma~\ref{lem:etale-germs-and-point-products}. Define
\(H^b_{\rm et}(X,F)\) by the cohomology on global
sections of this resolution, called the Godement
resolution. For constant coefficients abbreviate
\(H^b_{\rm et}(X,\underline\Lambda_X)\) to
\(H^b_{\rm et}(X,\Lambda)\).
\end{definition}

\begin{proposition}\label{prop:etale-cochains-strict-pullback}
For a map \(f:Y\to X\) and an abelian sheaf \(B\) on
\(Y\) put \((f_*B)(U)=B(Y\times_X U)\). A sheaf map
\(F\to f_*B\) induces a cochain map
\(C^\bullet_{\rm et}(X,F)\to C^\bullet_{\rm et}(Y,B)\).
These maps preserve identity and composition. In
particular each \(f\) induces functorial pullback on
the cochains with constant coefficients \(\Lambda\).
\end{proposition}
\begin{proof}
Base change takes covers and their overlaps to covers
and overlaps, so \(f_*B\) is a sheaf. A representative
on a pointed neighborhood of \(x=f(y)\) maps to a
section on its base change with the indicated lift of
\(y\). Passing to germs gives \(F_x\to B_y\).
At \(v\in P(Y\times_X U)\) use the factor of
\(G_X F(U)\) at its projection to \(U\) and apply
this stalk map at its projection to \(Y\). This
defines \(G_X F\to f_*G_Y B\), and on sections from
\(F\) it commutes with the germ inclusions. Composing
with the target quotient therefore factors through
the sheaf quotient \(Q_X F\to f_* Q_Y B\).
Iterating the two constructions gives the maps at
every degree and the quotient and germ squares give
compatibility with \(d^b\). For a composite the
chosen projections compose, the stalk maps compose
by restriction of representatives, and the maps on
quotients compose by the uniqueness of factorization.
This gives the stated identities on global cochains.
For constants use the map carrying a locally constant
value to that same value on the inverse image of its
defining neighborhood.
\end{proof}

\section{Restriction to a union of closed subschemes}

\begin{lemma}\label{lem:closed-immersion-etale-stalks}
Let \(j:D\hookrightarrow W\) be a closed immersion and
\(B\) an abelian sheaf on the small étale site of \(D\).
At a point \(x\in P(W)\) the stalk of \(j_*B\) is zero
if \(x\) lies outside \(D\), and is \(B_x\) at the
unique lift to \(D\) otherwise. In particular \(j_*\)
is exact and there are natural identifications
\[
         G_W j_*B=j_*G_D B,\qquad Q_W j_*B=j_*Q_D B .
\]
\end{lemma}
\begin{proof}
Outside \(D\) a pointed neighborhood can be restricted to
its complement. A sheaf of groups has just the zero section
on the empty object by gluing on its empty cover. Inside
\(D\) an arrow from the spectrum of a field whose point
lies in \(D\) kills its defining ideal, giving the unique
lift. Every pointed étale neighborhood on \(D\) can be
refined, near this lift, by the restriction of one on
\(W\). Choose an affine open of \(W\) around the image,
of ring \(A\), with \(D\) cut out by \(I\). Choose in
the given neighborhood a Jacobian chart over \(A/I\)
at the lift. A further principal restriction of the
base may be lifted by lifting its defining element
from \(A/I\) to \(A\). Lift the coefficients of the
chart equations and its inverted functions to \(A\),
and invert the lifted Jacobian determinant. Reduction
modulo \(I\) then gives exactly the chosen restricted
chart, including its lift from the field. The resulting
scheme is étale over \(W\) by
Definition~\ref{def:small-etale-point-conventions}.

Thus every germ on \(D\) is a germ of \(j_*B\).
For injectivity, start with a section on the restriction
to \(D\) of a given pointed \(U\) over \(W\). If its
germ on \(D\) vanishes, it vanishes on some smaller
étale neighborhood over \(U\times_WD\). Apply the same
chart lifting with \(U\) as base and \(U\times_WD\)
as closed subscheme. This gives a pointed restriction
over \(U\) on whose closed part the section is zero,
which proves injectivity. Two germs can be compared in
one neighborhood by the product in
Lemma~\ref{lem:etale-germs-and-point-products}.

The stalk description and that lemma give exactness.
In \(G_W(j_*B)(U)\) only the points of \(P(U)\) factoring
through \(D\) can have nonzero factors. They identify
with \(P(U\times_W D)\), with the stalks just computed.
This identifies the product on every \(U\) with
\((j_*G_D B)(U)\) and respects the inclusion of germs.
Taking sheaf quotients, using exactness of \(j_*\),
gives the formula for \(Q\).
\end{proof}

\begin{proposition}\label{prop:closed-cover-cochain-monomorphism}
Let \(j_a:D_a\hookrightarrow W\) be a finite set of
closed immersions whose underlying subsets cover \(W\).
Let \(\Lambda\) be an abelian group. For each \(b\ge0\)
the combined restriction
\[
       C^b_{\rm et}(W,\underline\Lambda_W)
          \longrightarrow
       \prod_a C^b_{\rm et}(D_a,\underline\Lambda_{D_a})
\]
is injective, with the restrictions of
Proposition~\ref{prop:etale-cochains-strict-pullback}.
In particular suppose \(f_i:W\to Z\) is a finite list
of morphisms, \(m_i\in\Z\), and for each \(a\)
\[
              \sum_i m_i[ f_i j_a]=0
                       \quad\text{in }\Z[\operatorname{Hom}(D_a,Z)].
\]
Here brackets denote the basis of the free abelian group.
Then \(\sum_i m_i f_i^*=0\) as cochain maps on constant
coefficients and on every \(H^b_{\rm et}(Z,\Lambda)\)
with target \(H^b_{\rm et}(W,\Lambda)\).
\end{proposition}
\begin{proof}
The restriction of the constant sheaf on \(W\) to the
closed pieces gives the sheaf map
\[
              \underline\Lambda_W\longrightarrow
                   \prod_a (j_a)_*\underline\Lambda_{D_a}.
\]
At a chosen geometric point its stalk has at least one
component equal to the identity on \(\Lambda\), by
Lemma~\ref{lem:closed-immersion-etale-stalks}.
Finite products commute with stalks by common pointed
refinement. Thus the map is injective. Apply the exact
functor \(Q_W^b\) of
Lemma~\ref{lem:etale-germs-and-point-products}. This
functor preserves finite direct sums since taking germs,
their products, and the sheaf quotients are additive.
By Lemma~\ref{lem:closed-immersion-etale-stalks}, its
target identifies with
\(\prod_a(j_a)_*Q_{D_a}^b\underline\Lambda_{D_a}\).
Applying \(G_W\) and taking sections on \(W\) preserves
injectivity by its point product formula. The same
identifications give exactly the displayed target.
At a closed piece the map is the restriction defined in
Proposition~\ref{prop:etale-cochains-strict-pullback},
because the stalk maps and their iterated quotients
were induced by that restriction.

Compose the cochain sum with each \(j_a^*\). By strict
functoriality the resulting sum is \(\sum_i m_i(f_i j_a)^*\).
Grouping identical morphisms gives the zero map by the
hypothesis in the free group. The injectivity just proved
forces the original cochain sum to be zero in each degree.
Passing to cohomology proves the final assertion.
\end{proof}

\begin{theorem}\label{thm:special-etale-hecke-congruence}
At every degree-one prime and level covered by
Proposition~\ref{prop:componentwise-special-hecke-maps},
use its \(k,Z,\mathcal W_k,\Phi\) and maps \(u_A\).
For every abelian group \(\Lambda\) and every \(b\ge0\),
\[
   \sum_{A\subseteq\mathbf F_p^n}
       (-1)^{\dim A}p^{\binom{\dim A}{2}}\,
                 u_A^*(\Phi^*)^{n-\dim A}=0
     \quad\text{as maps }\quad
       H^b_{\rm et}(Z,\Lambda)\longrightarrow
       H^b_{\rm et}(\mathcal W_k,\Lambda).
\]
The sum is over linear subspaces, including zero and
the whole space.
\end{theorem}
\begin{proof}
Use the reduced irreducible components of \(\mathcal W_k\)
as the closed pieces. They are finite in number and cover
its underlying space by
Corollary~\ref{cor:actual-special-correspondence-cycles}.
On each such \(D\), that corollary gives cancellation
in the free group on morphisms from \(D\) to
\(Z\times_k Z\). Composing its basis morphisms with the
second projection sends this zero sum to the required
one on the maps \(\Phi^{n-\dim A}u_A|_D\).
Proposition~\ref{prop:closed-cover-cochain-monomorphism}
and the order of pullbacks in
Proposition~\ref{prop:etale-cochains-strict-pullback}
give the result on \(\mathcal W_k\).
\end{proof}

\section{Representatives on finite systems of étale neighborhoods}

\begin{lemma}\label{lem:finite-chart-subcover-etale-site}
Let \(W\) be a scheme of finite type over a field. From a
jointly surjective family of étale maps to \(W\) one may
select finitely many affine chart neighborhoods in the
members which still cover \(W\). More generally a cover
of the underlying space by images of finite-type affine
charts over affine opens of \(W\) has a finite subcover.
\end{lemma}
\begin{proof}
Use noetherian induction on closed subsets of \(W\).
Such induction is valid by finite generation of polynomial
ideals and the descending chain condition on closed subsets
as in Proposition~\ref{prop:finite-strata-graded-partial-resolutions}.
It suffices to treat an irreducible reduced closed subset
\(D\), since there are finitely many irreducible components.
Choose a chart whose image contains the generic point of
\(D\), and work over an affine open in \(D\), of ring \(A\),
containing that point. After base change the chart has a
nonempty affine part with a nonempty fiber over
\(K=\operatorname{Frac}(A)\). Its nonzero coordinate
algebra remains nonzero after tensoring with \(\bar K\),
as is seen on a vector space basis. It has a
\(\bar K\)-point by the affine point argument of
Lemma~\ref{lem:projective-special-section-base-change}.
The finitely many coordinates of this point belong
to a finite extension of \(K\).
The images of the finitely many affine generators, including
any inverted coordinates, are algebraic over \(K\).
Clear the denominators in their monic equations by
localizing \(A\) at one nonzero element. They generate
a finite algebra \(B\) over this localized ring inside
the field extension, because bounded powers of integral
generators span their products. The map from the chart
algebra to \(B\) gives a map back to the chart on spectra.

At every prime of the localized \(A\), the localization
of \(B\) is nonzero since it contains the local base.
Its residue algebra is nonzero since it is finite:
if its maximal base ideal generated the module, a matrix
of coefficients in that ideal would express a finite
generating list in terms of itself. Its identity minus
that matrix has unit determinant and would kill the list.
A maximal ideal in the nonzero residue algebra gives a
point above the chosen prime. The selected chart therefore
covers a nonempty open in \(D\). The complement in \(D\)
is a proper closed subset and the induction finishes.
Affine étale charts over affine base opens give the
indicated finite-type family in the first assertion by
Definition~\ref{def:small-etale-point-conventions}.
\end{proof}

\begin{definition}\label{def:etale-matching-system}
An augmented system over a separated finite-type \(X\)
over a field has \(U_{-1}=X\) and schemes \(U_i\) étale
over \(X\) for \(i\ge0\), with face maps \(\partial_j\)
for \(0\le j\le i\) from \(U_i\) to \(U_{i-1}\).
They satisfy \(\partial_j\partial_l=\partial_{l-1}\partial_j\)
for \(j<l\) whenever the composites are defined. Put
\(M_0=X\). For \(i\ge1\), the matching object \(M_i\)
is the subscheme of \(U_{i-1}^{i+1}\) over \(X\) whose
tuple \((a_0,\ldots,a_i)\) satisfies
\(\partial_j a_l=\partial_{l-1}a_j\) for \(j<l\).
We use systems with \(U_i\to M_i\) jointly covering in
each degree, allowing \(U_i\) to be a disjoint union
of finitely many étale charts. A finite prescribed part
can be continued by \(U_i=M_i\) in higher degrees.
If the earlier objects are separated of finite type and
étale over \(X\), the same holds for \(M_i\).
Indeed each equality in its definition is a pullback
of the diagonal of an earlier étale object over \(X\).
That diagonal is closed by separatedness and open by
Definition~\ref{def:small-etale-point-conventions}.
A cover of \(M_i\) by affine étale charts therefore
continues the system with the same properties.
\end{definition}

\begin{proposition}\label{prop:finite-constant-hypercover-representatives}
For \(X\) as in Definition~\ref{def:etale-matching-system},
an abelian group \(\Lambda\), \(n\ge0\), and
\(y\in H^n_{\rm et}(X,\Lambda)\), one may choose such
a system through degree \(n+1\) and a section
\(\beta\in\underline\Lambda_X(U_n)\), constant on each
of finitely many disjoint chart members, with
\(\sum_j(-1)^j\partial_j^*\beta=0\) on \(U_{n+1}\).
Its class maps naturally to \(y\). Pullback of the
charts and section represents pullback of \(y\), and
a map between augmented systems over \(X\) preserves
the class so represented.
\end{proposition}
\begin{proof}
Put \(\mathcal G^b=G_X Q_X^b\underline\Lambda_X\)
from Definition~\ref{def:etale-godement-cochains}.
For a chosen system use the arrays
\(C^{a,b}=\mathcal G^b(U_a)\), \(a,b\ge0\).
The total differential is \(\delta+(-1)^a d\), where
\(\delta=\sum_j(-1)^j\partial_j^*\) increases \(a\);
the face identities give \(\delta^2=0\).
Restriction into \(a=0\) induces an isomorphism from
the cohomology on \(X\) to that of this total complex.
To prove it, augment the array at \(a=-1\) by
\(\mathcal G^b(X)\). Each horizontal row is exact.
Indeed, by the point product formula in
Lemma~\ref{lem:etale-germs-and-point-products}
it splits into products indexed by \(x\in P(X)\).
For each \(x\), the terms are functions on the sets
of lifts of \(x\) to \(U_a\), valued in the single
group \((Q_X^b\underline\Lambda_X)_x\). Every matching
tuple of lifts has a filler by the covering property.
Choose a lift at degree zero. For a simplex \(t\)
of degree \(a\) choose \(H(t)\) of degree \(a+1\)
with zeroth face \(t\) and \(j\)-th face
\(H(\partial_{j-1}t)\) for \(j>0\). Start at degree
\(-1\) with the chosen lift. Induction using the face
identities makes these a matching tuple. On the free
groups on these sets the alternating boundary satisfies
\(\partial H(t)=t-H(\partial t)\). Evaluation gives a
contraction of the augmented row of functions. The
augmented total is consequently exact by the row
filtration argument of Proposition~\ref{prop:support-cover-glue},
which uses only finitely many bidegrees in any fixed
total degree. This proves the comparison. The section
of \(\underline\Lambda_X\) in degree \(n\) is sent to
\(C^{n,0}\) by germs. A face cocycle there defines a
class by inverting this comparison.

To obtain a representative, start with a global closed
\(c\in\mathcal G^n(X)\) representing \(y\). For \(n>0\)
it lifts locally under \(d:\mathcal G^{n-1}\to
\mathcal G^n\) by exactness of the sheaf resolution.
Choose a cover \(U_0\to X\) by finitely many such charts
using Lemma~\ref{lem:finite-chart-subcover-etale-site}.
Subtract from the cocycle the total coboundary of its
chosen primitive there. Only the term in bidegree
\((1,n-1)\) remains, and it is the signed face sum
of the primitives, already defined on \(M_1\).
In general, in bidegree \((a,n-a)\) the remaining
section on \(M_a\) is vertically closed. For \(a=0\)
this is the original condition and thereafter it
follows from commutation with face restriction and
cancellation of the two alternating faces on matching
tuples. When \(a<n\), lift locally under
\((-1)^a d\), again choosing a finite cover of \(M_a\),
and subtract the resulting total coboundary. Its face
sum gives the sole remaining section, now defined on
\(M_{a+1}\). These subtractions take place in the
eventual total complex since they use only already
chosen faces. When \(a=n\), the vertical kernel is
\(\underline\Lambda_X\). Choose \(U_n\to M_n\) to be
a finite chart cover on each member of which that
section has a constant value. This also applies at
\(n=0\) directly to \(c\) in the vertical kernel.
The final section has alternating face sum zero on
\(M_{n+1}\) by the matching identities (at \(n=0\)
it comes from \(X\)). Take \(U_{n+1}=M_{n+1}\).
The subtractions prove that \(\beta\) represents \(y\).

Under changing a system by a map of augmented systems
over \(X\), restriction on each array commutes with
the augmentation. Under \(f:Y\to X\) use the cochain
maps of Proposition~\ref{prop:etale-cochains-strict-pullback}
on the arrays, defined there on every \(U_a\) by
passage to \(Y\times_X U_a\). They respect the
augmentations, faces, and inclusions of locally
constant sections. The isomorphism on cohomology just
proved therefore gives both naturality assertions.
\end{proof}

\section{The action of geometric Frobenius on the special fiber}

\begin{lemma}\label{lem:absolute-frobenius-etale-cartesian}
Let \(X\) be a scheme in characteristic \(p>0\), with
\(p\) prime, and denote the absolute Frobenius by \(F_X\),
the endomorphism given on functions by \(a\mapsto a^p\).
For every étale \(U\to X\) the canonical map
\[
              U\longrightarrow X\times_{F_X,X}U
\]
over \(X\), whose second projection is \(F_U\), is an
isomorphism.
\end{lemma}
\begin{proof}
Use affine charts \(A\to B\) of
Definition~\ref{def:small-etale-point-conventions}.
Write the chart with variables \(t_1,\ldots,t_m\),
equations \(f_1,\ldots,f_m\), and inverted function \(g\)
including the Jacobian determinant in the localization.
For \(B'=A\otimes_{A,F_A}B\) we use the convention
that scalars of \(B\) move to \(A\) by \(F_A\).
It has variables \(t'_i\), equations \(f_j^{(p)}(t')\)
and inverted \(g^{(p)}(t')\), where the superscript
raises only the coefficients to \(p\)-th powers.
The ring map in question is \(B'\to B\), identity on
the left \(A\), \(t_i'\mapsto t_i^p\).

Form the auxiliary algebra
\[
        S=B'[v_1,\ldots,v_m]/(v_i^p-t_i':1\le i\le m).
\]
Successive division by the monic equations gives a free
\(B'\)-basis \(v^a\) with \(0\le a_i<p\). The ideal
\(J=(f_j(v):1\le j\le m)\) satisfies \(z^p=0\) for all
\(z\in J\), since \(f_j(v)^p=f_j^{(p)}(t')=0\).
It is nilpotent, as every product of more than \(m(p-1)\)
of its generators repeats one at least \(p\) times.
Here \(g(v)^p=g^{(p)}(t')\) is a unit. Eliminating \(t'\)
identifies \(S/J\) with \(B\), by \(v_i\mapsto t_i\).

The identity of \(B=S/J\) lifts to an \(A\)-map \(s:B\to S\).
Start with any lifts of the \(t_i\). At each successive
quotient \(J^a/J^{a+1}\), \(a\ge1\), the errors in the
equations can be corrected by multiplying their negatives
by the inverse Jacobian in that square-zero layer. The
Jacobian and the inverted coordinates are units since
they reduce to units modulo the nilpotent ideal. The
iteration terminates when the ideal power is zero.
This is the coordinate correction used in
Lemma~\ref{lem:jacobian-chart-flat-descent}.

In \(S\) each derivation \(\partial/\partial v_i\) fixing
\(B'\) is well defined, since the derivatives of the
power relations vanish. Applied to the equations for
\(s(t_j)\), it gives the invertible Jacobian times their
derivatives equal to zero. The coefficients from \(A\)
are fixed. Hence all \(s(t_j)\) are in the common kernel
of these derivations. On the free basis \(v^a\) that
kernel consists of \(B'\), since \(1,\ldots,p-1\) are
units and distinct terms of positive \(v_i\)-degree have
distinct derivatives in coordinate \(i\). The inverses
of \(s(g)\) are also killed by the derivations, so \(s\)
lands in \(B'\). In \(S\), the difference \(s(t_i)-v_i\)
lies in \(J\), hence \(s(t_i)^p=v_i^p=t_i'\).
The map \(s:B\to B'\) is therefore inverse on both
sides to the proposed map. The affine identifications
agree on overlaps since they invert the canonical map.
Indeed the preimage of each chart under \(F_U\) is the
same open, as \(a^p\) lies in a prime exactly when \(a\)
does. Thus the local inverses cover both sides.
\end{proof}

\begin{proposition}\label{prop:absolute-frobenius-on-etale-cohomology}
Let \(X\) be separated of finite type over a field of
characteristic \(p>0\). The action \(F_X^*\) is the
identity on \(H^b_{\rm et}(X,\Lambda)\) for every \(b\)
and every abelian group \(\Lambda\). If \(X=X_0\times_{\mathbf F_p}
\overline{\mathbf F}_p\), write \(\Phi\) for the base
change of \(F_{X_0}\), and \(s_\varphi\) for the
endomorphism given on the scalar field by
\(\varphi(a)=a^p\) and on \(X_0\) by the identity.
Then \(\Phi^*=(s_{\varphi^{-1}})^*\).
\end{proposition}
\begin{proof}
Represent a class by \(U_\bullet,\beta\) as in
Proposition~\ref{prop:finite-constant-hypercover-representatives}.
Pullback by \(F_X\) uses the base changes of that
system. Lemma~\ref{lem:absolute-frobenius-etale-cartesian}
identifies the original system with this pullback via
a map over \(X\) commuting with all faces. The composite
on each \(U_i\) with the projection back to \(U_i\)
is \(F_{U_i}\), which preserves its underlying points
and each of its disjoint chart members. The constants
of \(\beta\) on those members are unchanged.
The naturality in the representing proposition now
proves that \(F_X^*\) fixes the class.

On each extended affine chart \(A_0\otimes_{\mathbf F_p}
\overline{\mathbf F}_p\), the two commuting ring maps
giving \(\Phi\) and \(s_\varphi\) apply the \(p\)-th
power on the first and second factors respectively.
Their composite is absolute Frobenius. Consequently
\(\Phi^*(s_\varphi)^*=1\) by the first assertion and
Proposition~\ref{prop:etale-cochains-strict-pullback}.
The inverse to \((s_\varphi)^*\) is
\((s_{\varphi^{-1}})^*\), proving the formula.
\end{proof}

\begin{corollary}\label{cor:special-hecke-geometric-galois-operator}
In Theorem~\ref{thm:special-etale-hecke-congruence}
take \(k=\overline{\mathbf F}_p\) and let
\(r_{\rm sp}(\gamma)=(s_{\varphi^{-1}})^*\) act on the
special target at a geometric Frobenius datum
\((w,h,\gamma)\) of Definition~\ref{def:frobenius-data}.
Then for all \(b\) and \(\Lambda\)
\[
    \sum_{A\subseteq\mathbf F_p^n}
      (-1)^{\dim A}p^{\binom{\dim A}{2}}\,
        u_A^*\,r_{\rm sp}(\gamma)^{\,n-\dim A}=0
   \quad\text{on }H^b_{\rm et}(Z,\Lambda),
\]
with values on the special cover \(\mathcal W_k\).
\end{corollary}
\begin{proof}
The residue field action of \(\gamma^{-1}\) is
\(\varphi\) by Definition~\ref{def:frobenius-data}.
The coarse special target is defined over \(\mathbf F_p\)
and its \(\Phi\) is the base change in
Proposition~\ref{prop:componentwise-special-hecke-maps}.
Substitute Proposition~\ref{prop:absolute-frobenius-on-etale-cohomology}
in Theorem~\ref{thm:special-etale-hecke-congruence}.
\end{proof}

\section{Continuity and integral coefficients}

\begin{proposition}\label{prop:generic-etale-galois-continuity}
Let \(E\) be a number field, \(X\) a separated scheme of
finite type over \(E\), and \(\bar X=X_{\bar E}\) for the
algebraic closure in Definition~\ref{def:frobenius-data}.
For \(g\in G_E\) write \(s_g:\bar X\to\bar X\) for the
map on affine charts induced by \(a\otimes t\mapsto
a\otimes g(t)\). For each abelian group \(\Lambda\)
and \(b\ge0\),
\[
        r_b(g)=s_g^*:H^b_{\rm et}(\bar X,\Lambda)
                   \longrightarrow H^b_{\rm et}(\bar X,\Lambda)
\]
defines an action on the discrete group in which each
element has open stabilizer. Morphisms over \(E\) and
homomorphisms of constant coefficient groups induce
equivariant maps.
\end{proposition}
\begin{proof}
For composition \(gh=g\circ h\) one has
\(s_{gh}=s_h\circ s_g\) on schemes. Pullback reverses
this order by Proposition~\ref{prop:etale-cochains-strict-pullback},
giving the representation property and inverses.
Equivariance follows from the commuting squares for a
morphism over \(E\) and from the functorial stalk and
quotient arrows for a coefficient map.

Choose a representative system and its constants for
one class as in
Proposition~\ref{prop:finite-constant-hypercover-representatives}.
Only degrees through \(b+1\) are needed. The charts and
all their face maps and identities are described by
finitely many coefficients in \(\bar E\). To see this
on each member, use finite affine presentations for
the schemes and refine to affine opens on which a face
map lands in a target chart. All overlaps have finite
affine covers since the schemes are separated of finite
type over a field. The transition maps and the mutual
inverse and composite identities can be written with
finitely many polynomials, inverted denominators and
ideal-membership identities, as in the chart proof of
Proposition~\ref{prop:finite-presentation-spread}. This
applies also to the matching objects since they are
formed by finitely many products and equality conditions
in Definition~\ref{def:etale-matching-system}.
Adjoin all these coefficients to \(E\) and take a
finite normal extension containing them, using the
power-basis argument of Lemma~\ref{lem:place-fields}.
Each automorphism fixing it preserves the chart
members, all their maps, and the individual constants
assigned to them. Equivalently its pullback system
is identified with the original by the isomorphisms
on these presentations, with the same cocycle.
Naturality in the representing proposition fixes the
original class. The subgroup so obtained is open by
Definition~\ref{def:frobenius-data}.
\end{proof}

\begin{definition}\label{def:integral-etale-cohomology-limit}
For a scheme \(Y\) and a prime \(\ell\), put
\[
    L^b_\ell(Y)=\varprojlim_{a\ge1}
           H^b_{\rm et}(Y,\Z/\ell^a\Z),\qquad
    V^b_\ell(Y)=\Q_\ell\otimes_{\Z_\ell} L^b_\ell(Y).
\]
The transition maps are coefficient reduction. Use the
inverse limit topology of the discrete groups on the
left. Compatible residues give its \(\Z_\ell\)-module
structure by scalar multiplication in each coordinate.
\end{definition}

\begin{proposition}\label{prop:finiteness-from-one-residue-degree}
If \(m=\dim_{\mathbf F_\ell}H^b_{\rm et}(Y,\mathbf F_\ell)\)
is finite, then \(L^b_\ell(Y)\) is generated by at most
\(m\) elements over \(\Z_\ell\). Its kernel on passage to
\(H^b_{\rm et}(Y,\Z/\ell^a\Z)\) is \(\ell^aL^b_\ell(Y)\)
for every \(a\ge1\). In particular
\(\dim_{\Q_\ell}V^b_\ell(Y)\le m\).
For \(Y=\bar X\) of
Proposition~\ref{prop:generic-etale-galois-continuity}
the action on \(L^b_\ell(Y)\) is continuous with its
inverse limit topology, and under this finiteness
hypothesis the action on \(V^b_\ell(Y)\) is continuous
in matrix entries.
\end{proposition}
\begin{proof}
Write \(H_a=H^b_{\rm et}(Y,\Z/\ell^a\Z)\) and
\(L=L^b_\ell(Y)\). For positive \(a,t\), the short
exact sequence of coefficients
\[
 0\longrightarrow\Z/\ell^t\Z
 \xrightarrow{\ \ell^a\ }\Z/\ell^{a+t}\Z
 \longrightarrow\Z/\ell^a\Z\longrightarrow0
\]
is exact on sheaf stalks, hence on sheaves. Each functor
\(G_YQ_Y^j\) takes it to a sequence exact on sections,
by Lemma~\ref{lem:etale-germs-and-point-products}.
Consequently \(H_t\to H_{a+t}\to H_a\) is exact in
the middle. Indeed, for a cocycle mapping to a
coboundary, lift the primitive to the middle complex
and subtract its differential. The remaining representative
lifts to the first complex and is closed there by
injectivity. In degree zero no subtraction is needed.
Taking \(a=1\) and inducting gives
\(|H_t|\le |H_1|^t<\infty\).

Suppose \(x=(x_j)\in L\) has \(x_a=0\). For each
\(t\) the set of \(y_t\in H_t\) mapping to
\(x_{a+t}\) under the displayed inclusion is nonempty
by exactness and is finite. Reduction sends each such
set into the preceding one. There is a compatible
choice at every stage: among finitely many elements
at the first stage some admits descendants arbitrarily
far up, and among its finitely many immediate descendants
one has the same property. Inducting gives \(y\in L\).
Reducing the equality at \(a+t\) further to \(t\)
gives \(\ell^a y_t=x_t\). Thus \(x=\ell^a y\);
the converse follows on reducing to \(a\).

It follows that \(L/\ell L\) embeds in \(H_1\).
Choose \(r\le m\) lifts of a basis of its image.
Expand any element \(x\) modulo \(\ell\) in these
lifts, using digits between \(0\) and \(\ell-1\).
The residual difference is \(\ell\) times an element
of \(L\). Repeat on that element. The successive
digits for each basis coefficient give a member of
\(\Z_\ell\), by the convergent digit expansions in
Lemma~\ref{lem:place-fields}. Their linear combination
equals \(x\) in each coordinate and hence in \(L\).
This proves finite generation and the rational dimension
bound.

For \(Y=\bar X\), continuity of the action on each
discrete factor follows from
Proposition~\ref{prop:generic-etale-galois-continuity}.
At \((g,x)\), holding the coordinate \(x_a\) fixed
and restricting \(g\) to its coset of the stabilizer
of \(x_a\) fixes the result at that coordinate.
This gives joint continuity on the inverse limit.
In \(V^b_\ell(Y)\) choose a vector space basis from
the images of a finite generating set of \(L\).
In these coordinates the image of \(L\) lies in
\(\ell^{-c}\Z_\ell^r\) for some \(c\ge0\), using the
actual basis size in place of \(r\) if necessary.
If \(g\) fixes the reductions of the chosen lifts
of the basis modulo level \(a\), its matrix differs
from the identity by entries in
\(\ell^{a-c}\Z_\ell\), by the kernel assertion.
The stabilizers form an open finite intersection.
Letting \(a\) increase proves continuity at the
identity and hence everywhere by translation.
\end{proof}

\section{The Hecke action on geometric étale cohomology}

\begin{proposition}\label{prop:etale-cochain-deck-trace}
Let \(\pi:Y\to X\) be finite étale and surjective with
an action over \(X\) of a finite group \(D\) such that
\(\coprod_{g\in D}Y\to Y\times_XY\), \(y\mapsto(y,gy)\)
on the \(g\)-th copy, is an isomorphism. For every
abelian group \(\Lambda\) the constant-coefficient
complexes of Definition~\ref{def:etale-godement-cochains}
have a trace from \(Y\) to \(X\) such that
\[
        {\rm Tr}_\pi\pi^*=|D|I,\qquad
        \pi^*{\rm Tr}_\pi=\sum_{g\in D}g^* .
\]
It commutes with homomorphisms of coefficients.
Consequently the pullback on \(V_\ell^b\) of
Definition~\ref{def:integral-etale-cohomology-limit}
is injective and identifies the space on \(X\) with
the \(D\)-invariant space on \(Y\).
\end{proposition}
\begin{proof}
Over the étale cover \(Y\) of \(X\), \(\pi\) is the
disjoint union of \(|D|\) copies of the identity by
the product isomorphism. For a sheaf \(B\) on \(Y\)
one therefore has
\[
                     (\pi_*B)_x=\prod_{v\mid x} B_v
                      \qquad(x\in P(X)).
\]
Indeed choose a lift of \(x\) to \(Y\). The pointed
neighborhoods lying over that lift are cofinal by
taking fiber products. On such neighborhoods the
direct image is a finite product over the sheets and
the stalk is the product of their stalks by common
refinement. A pointed neighborhood on one sheet is
also étale over \(X\), and two morphisms of pointed
neighborhoods agreeing at the lift agree there after
restriction by the diagonal calculation in
Definition~\ref{def:small-etale-point-conventions}.
This identifies each stalk with the stated \(B_v\).
In particular \(\pi_*\) is exact and
\(G_X\pi_*B=\pi_*G_Y B\) by taking the product
over the lifts on each object. Taking sheaf quotients
gives \(Q_X\pi_*B=\pi_*Q_Y B\) by exactness and
Lemma~\ref{lem:etale-germs-and-point-products}.

On a splitting neighborhood sum the entries of
\(\pi_*\underline\Lambda_Y\) over its sheets. Changes
of splitting permute the sheets, so the sums glue to
a sheaf map \(\pi_*\underline\Lambda_Y\to
\underline\Lambda_X\). Apply the iterated germ and
quotient functors and use the identifications just
proved to get a map from the cochains on \(Y\) to
those on \(X\), commuting with the differentials by
their construction. Restriction of any germ resolution
to an étale object is its resolution there: stalks are
the same by the pointed-neighborhood argument above,
the \(G\) formula on each further étale object is then
identical, and taking \(Q\) commutes with restriction
by the stalk test. Hence on a split neighborhood this
cochain map is ordinary summation and pullback is the
diagonal on the sheets. Summation after the diagonal
multiplies by \(|D|\), and the diagonal after summation
adds all sheet permutations from this regular group
action. These are identities on the complexes of
sheaves, since they hold on a cover, so on global
cochains as well. The same sum commutes with all
changes of coefficients.

Apply cohomology, take compatible inverse limits, and
extend scalars to \(\Q_\ell\). Division by \(|D|\)
is allowed even when \(\ell\) divides \(|D|\). The
first identity gives injectivity and the second
expresses any invariant as the pullback of its
normalized trace. Every pullback is invariant since
\(\pi g=\pi\).
\end{proof}

\begin{theorem}\label{thm:actual-generic-etale-hecke-operators}
For the marked power \(X_E\) and its neat level in
Theorem~\ref{thm:actual-coherent-hecke-eigenquotients}
put \(V^b=V_\ell^b((X_E)_{\bar E})\), \(b\ge0\).
Define the rational positive Hecke algebra by the
coset sums and multiplication of
Definition~\ref{def:actual-positive-hecke} with rational
coefficients. It acts on \(V^b\) by operators
\(T^{\rm et}_H\) commuting with \(G_E\).
On each sufficiently fine common generic level,
\[
      \pi_m^* T^{\rm et}_H
             =\sum_{gC}H(gC)\,u_{g,\mathbf1}^* .
\]
The action is unital and multiplicative and extends by
scalars to \(\Q_\ell\). The safe operators
\(T^{\rm et}_{w,r}=T^{\rm et}_{\tau_{w,r}}\) commute.
\end{theorem}
\begin{proof}
The maps and deck actions are defined over \(E\) by
Theorem~\ref{thm:integral-hecke-cover-refinements};
the product isomorphism for their generic torsors
is the one verified in
Proposition~\ref{prop:coherent-finite-level-trace}.
The finite coset sum of pullbacks is invariant on its
target under \(C/C_m\), because a deck translate
permutes \(gC\) by left multiplication and preserves
\(H\), by the frame identity in
Theorem~\ref{thm:actual-hecke-action-on-forms}.
Proposition~\ref{prop:etale-cochain-deck-trace} thus
gives a unique descended operator. Injectivity on
common refinements makes it independent of the cover.
The map for the unit coset is \(\pi_m\), so the
unit acts by identity. Every pullback is equivariant
for the \(G_E\)-actions of
Proposition~\ref{prop:generic-etale-galois-continuity}
on each coefficient group and thus on the rational
limit. Equivariance of the descended operator follows
by injectivity of \(\pi_m^*\).

For two coset sums choose successively deep levels
on which the first set of representatives maps to a
cover carrying all the second maps. Their composite
maps have frame labels \(gst\) when the starting label
is \(g\) and the successive right representatives are
\(s,t\), by the algebraic map identity proved in
Theorem~\ref{thm:actual-hecke-action-on-forms}. Thus
after pullback the composition of the two operators
is the double sum of the maps for \(stC\), each with
its counted product coefficient. Functoriality is
Proposition~\ref{prop:etale-cochains-strict-pullback}
on finite coefficients, followed by passage to \(V^b\).
Cancel the injective pullback to obtain multiplication.
Linearity follows from the defining sum and its unique
descent. The relations of
Lemma~\ref{lem:spherical-commutativity-on-marked-tower}
are identities of the rational coset counts, hence
hold on \(V^b\).
\end{proof}

\section{The congruence on functions on the components}

\begin{lemma}\label{lem:constant-etale-functions-on-smooth-components}
Let \(X\) be smooth, separated and of finite type over
an algebraically closed field \(k\). Its connected
components are finitely many open and closed irreducible
subschemes. For any abelian group \(\Lambda\),
\(H^0_{\rm et}(X,\Lambda)\) is the group of
\(\Lambda\)-valued functions on these components.
If \(X\) is the geometric generic fiber of a fixed
projective model over a localization of the integers
of a number field \(E\), then the partition into
components extends over a finite normal extension
\(B/E\) and a further localization \(O_B[1/N]\).
It may be chosen with the action of \({\rm Gal}(B/E)\)
permuting the extended subsets in the same way as
the original components.
\end{lemma}
\begin{proof}
The local calculation in
Lemma~\ref{lem:smooth-connected-fiber-integral} gives
a local domain at each closed point and shows that the
finitely many irreducible components are disjoint and
reduced. Its proof up through disjointness uses only
the smooth affine charts. Each component is therefore
open and connected. By
Definition~\ref{def:etale-godement-cochains} the
degree-zero kernel is the group of global sections of
the constant sheaf. On each component such a section
is constant. Indeed take an étale cover of this component
by charts with constant values. Each nonempty chart
maps dominantly on at least one irreducible piece of
its source to the base component. To verify this,
take a \(k\)-point in the chart using
Lemma~\ref{lem:projective-special-section-base-change}.
Its completed local ring is isomorphic to that at its
image: the extra variables in the étale presentation
are solved, uniquely to each order, in the same free
parameters by the Jacobian computation in
Lemma~\ref{lem:smooth-connected-fiber-integral}.
Both local rings inject into their completions by
that lemma's local argument. On the irreducible chart
piece through the point the map from the local domain
on the base is thus injective, so its generic point
maps to the generic point of the base component.
Any two nonempty chart members now have nonempty
fiber product over the generic point, since the tensor
product of two nonzero algebras over its field is
nonzero by a vector space basis. The matching condition
forces their constants to agree. Conversely an
assignment on the open components is a section by
gluing. This proves the description.

For the extension, use a finite affine cover of the
fixed model. On each extended generic affine chart the
component subsets are defined by idempotents, with
product zero for distinct subsets and sum one.
Indeed on disjoint complementary opens the pair of
regular functions one and zero glues to the required
idempotent. Express these finitely many idempotents
and their overlap identities in the coordinate
presentations of the model. Adjoin their coefficients
to \(E\) and take a finite normal extension \(B\)
containing them. The finitely many identities for
their products, sums, overlaps, and conjugates by
\({\rm Gal}(B/E)\) persist after inverting one
positive integer \(N\): clear the denominators of
coordinates and of the certificates of equality in
the localized quotient algebras, as in
Proposition~\ref{prop:finite-presentation-spread}.
Alternatively, each difference already over the number
ring and zero after extension to its fraction field is
killed by one nonzero integer after clearing the norm
of the denominator. The idempotents partition the model
after this single localization. The conjugate identities
give the asserted action.
\end{proof}

\begin{theorem}\label{thm:degree-zero-generic-frobenius-hecke}
Fix a marked coarse target \(X_E\) in
Theorem~\ref{thm:actual-generic-etale-hecke-operators}.
On the finite-dimensional representation
\(V^0=V^0_\ell((X_E)_{\bar E})\), for every geometric
Frobenius datum \((w,h,\gamma)\) at a safe degree-one
prime of norm \(p\) outside a fixed finite set one has
\[
    \sum_{r=0}^n (-1)^r p^{\binom r2}\,
              T^{\rm et}_{w,r}\,r_0(\gamma)^{\,n-r}=0 .
\]
Here \(r_0\) denotes the action of
Proposition~\ref{prop:generic-etale-galois-continuity}
in degree zero. The exclusions may depend on \(X_E\)
and are independent of the fine covers chosen to
represent these operators.
\end{theorem}
\begin{proof}
Take the coarse model and the fixed exclusions of
Theorem~\ref{thm:actual-special-graph-point-congruence}.
Add the primes for the component partition of
Lemma~\ref{lem:constant-etale-functions-on-smooth-components}
of that model. For a remaining \(w\) let
\(A=(O_{\bar E})_{\ker h}\). This is a valuation
ring of \(\bar E\) stabilized by \(\gamma\), by
Lemma~\ref{lem:residue-map-to-valuation}.
Its residue field identifies via \(h\) with
\(\overline{\mathbf F}_p\). Indeed its image is
algebraic over \(\mathbf F_p\); it is algebraically
closed because a monic polynomial over the residue
field lifts to a monic polynomial over \(A\), and
all roots in \(\bar E\) are integral in \(A\).
A root of negative valuation would make its highest
power the unique term of least valuation. Reduction
thus splits the original polynomial. On this residue
field \(\gamma\) acts by inverse \(p\)-power Frobenius.
The coarse model over \(A\) has the open and closed
component labels extended by the lemma, since it is
the base change of the same model from the localized
number ring.

Take any \(a\ge1\), \(\Lambda=\Z/\ell^a\Z\), and
any degree-zero class on the geometric generic fiber.
By the lemma it is a function of those labels. Its
values extend to a section on the model over \(A\)
by assigning the same values on its open and closed
pieces. Choose a projective fine cover \(\mathcal W\)
and its integral diagram for this \(w\) by
Theorem~\ref{thm:integral-hecke-cover-refinements}.
On \(\mathcal W_A\) pull back the extended function
after the appropriate powers of the scalar automorphism
\(s_\gamma\) of the coarse model over the underlying
integer ring. Sum these functions with the coefficients
of Corollary~\ref{cor:special-hecke-geometric-galois-operator}.
The result is given on a finite partition into open
and closed subsets, obtained by intersecting the
partitions for the summands. Its restriction to the
special fiber is zero by that corollary, including
the indicated residue action of \(\gamma\).

Every \(\bar E\)-point of the generic fine cover
extends to an \(A\)-point. In its projective equations
scale homogeneous coordinates by one of least valuation
so that they are integral and one is a unit. All
defining homogeneous relations still vanish in \(A\).
The spectrum of \(A\) is connected since it is a
domain. An \(A\)-point therefore lies in just one
piece of the finite partition. Its generic and
special values agree, so its generic value is zero.
All degree-zero classes on the smooth generic cover
are detected on such points by the lemma.
Consequently the congruence after fine-level pullback
is zero at each finite coefficient level. In the
integral inverse limit and after extension to \(\Q_\ell\)
it is still zero. By Lemma~\ref{lem:actual-subspace-kernel-labels}
summing the unscaled legs in dimension \(r\) gives
the defining pullback of \(T^{\rm et}_{w,r}\).
Cancel the injective \(\pi_m^*\) by
Proposition~\ref{prop:etale-cochain-deck-trace}
to get the formula. Finally \(V^0\) is a finite sum
of copies of \(\Q_\ell\) indexed by the geometric
components, by the lemma at all coefficient levels.
\end{proof}

\section{Restriction of a relation from a valuation model}

\begin{proposition}\label{prop:generic-hecke-two-restrictions}
Fix \(b\ge0\), a safe degree-one prime \(w\) of norm
\(p\ne\ell\) in Theorem~\ref{thm:special-etale-hecke-congruence},
and a geometric datum \((w,h,\gamma)\).
Use \(A=(O_{\bar E})_{\ker h}\) and the coarse model
\(\mathcal Z_A\) and projective fine model \(\mathcal W_A\)
of Theorem~\ref{thm:integral-hecke-cover-refinements}.
For either model \(Y\), let \(j_Y\) and \(i_Y\) be the
inclusions of its geometric generic and special fibers.
Write \(V(Y)=V^b_\ell(Y)\). Suppose the following two
relations hold for this diagram:
\[
 \begin{split}
  j_{\mathcal Z}^*V(\mathcal Z_A)&=V((\mathcal Z_A)_{\bar E}),\\
  \ker\big(i_{\mathcal W}^*:V(\mathcal W_A)
                \to V((\mathcal W_A)_{\overline{\mathbf F}_p})\big)
                          &\subseteq \ker j_{\mathcal W}^*.
 \end{split}
\]
Then on \(V((\mathcal Z_A)_{\bar E})\),
\[
   \sum_{r=0}^n (-1)^r p^{\binom r2}
                  T^{\rm et}_{w,r}\,r_b(\gamma)^{n-r}=0.
\]
The same conclusion after pullback to the fine cover,
with coefficients \(\Lambda=\Z/\ell^a\Z\) for a single
\(a\ge1\), holds if the two relations are satisfied
with \(H^b_{\rm et}(-,\Lambda)\) in place of \(V\).
\end{proposition}
\begin{proof}
The proof of
Theorem~\ref{thm:degree-zero-generic-frobenius-hecke}
identifies the residue field of \(A\) and the action
of \(\gamma\) on it. Since all models and maps descend
to the original localization of \(O_E\), \(\gamma\)
acts on the scalar ring of each model over \(A\);
write \(s_\gamma\) for the resulting scheme map on
\(\mathcal Z_A\). Restriction to either fiber commutes
with its pullback, using the induced action on that
fiber, by Proposition~\ref{prop:etale-cochains-strict-pullback}.
Each integral \(u_A\) of the Hecke diagram likewise
commutes with restriction. To avoid conflict with the
valuation ring, index these maps here by subspaces
\(D\subseteq\mathbf F_p^n\) and write them as \(u_D\).

Take a generic class \(x=j_{\mathcal Z}^*y\), which is
possible by the first relation. On the integral fine
cover form the class
\[
        z=\sum_{D\subseteq\mathbf F_p^n}
                (-1)^{\dim D}p^{\binom{\dim D}{2}}\,
                    u_D^*(s_\gamma^*)^{n-\dim D}y .
\]
Its restriction to the special fiber is zero by
Corollary~\ref{cor:special-hecke-geometric-galois-operator}
on each finite coefficient group, also after inverse
limits and rational extension by
Definition~\ref{def:integral-etale-cohomology-limit}.
The second relation therefore gives \(j_{\mathcal W}^*z=0\).
This is exactly the vanishing on the generic fine
cover. For rational local coefficients, group the
subspaces by dimension using
Lemma~\ref{lem:actual-subspace-kernel-labels}.
Theorem~\ref{thm:actual-generic-etale-hecke-operators}
identifies the result with the pullback of the claimed
polynomial on \(x\), and
Proposition~\ref{prop:etale-cochain-deck-trace}
makes that pullback injective. At a single finite
coefficient level the argument up to
\(j_{\mathcal W}^*z=0\) already proves its assertion.
\end{proof}

\begin{corollary}\label{cor:finite-coefficients-two-restrictions-suffice}
For fixed \(w,b\) in
Proposition~\ref{prop:generic-hecke-two-restrictions},
suppose its two relations hold at an unbounded set of
finite coefficient levels \(a\), using one fine level
carrying the specified correspondence. Then the rational
generic operator polynomial in that proposition vanishes
on the full \(V^b_\ell((\mathcal Z_A)_{\bar E})\).
\end{corollary}
\begin{proof}
Take an element of the integral inverse limit. At every
specified \(a\), the fine-level sum on its \(a\)-th
coordinate vanishes by the finite-coefficient assertion.
At any smaller level the same sum vanishes by coefficient
reduction, since all its maps are pullbacks and its
coefficients are integers. Unboundedness covers every
coordinate. Thus the fine-level sum vanishes on the
integral limit and, by linearity, on its rational scalar
extension. Group terms and cancel \(\pi_m^*\) as in
Proposition~\ref{prop:generic-hecke-two-restrictions}.
\end{proof}

\section{Finite precision in the rank criterion}

\begin{theorem}\label{thm:approximate-annihilator-rank}
Take \(E,M,n,r,Q,\kappa\) with the diagonal span,
local diagonal family, and \(n\) prescribed free
weights of Theorem~\ref{thm:spherical-annihilator-rank},
over \(k_0=K_\ell\subset k=C_\ell\) as in
Definition~\ref{def:completed-local-action}. Write
\(\|B\|=\max_{ij}|B_{ij}|_\ell\) in the given basis.
Assign a monic polynomial \(P_v\in K_\ell[X]\) of
degree \(n\) at each degree-one prime outside a finite
set. Suppose that for every real \(\eta>0\),
\[
                        \|P_v(r(\gamma))\|\le\eta
                \quad\text{for most Frobenius data }(v,h,\gamma)
\]
in the sense of Definition~\ref{def:frobenius-approximation}.
Then \(M=n\) and \(r\) has polynomials \(P_v\) in the
approximation sense of that definition.
\end{theorem}
\begin{proof}
Conjugating by the fixed \(Q\) multiplies the matrix
norm bound by at most \(\|Q\|\|Q^{-1}\|\), by the
ultrametric inequality on a matrix product. Hence the
hypothesis at all precisions applies to \(s=QrQ^{-1}\).
Fix a power minor \(\Delta\) of order \(n+1\) from
Definition~\ref{def:power-minor}. Adding to its last
column the corresponding linear combination of the
preceding columns leaves its determinant unchanged.
Monicity permits replacing that last column of
\(s(\gamma)^n\) by the same entries of \(P_v(s(\gamma))\).

Suppose \(\Delta(s(g))\ne0\) at some \(g\in G_E\).
On a nonempty open neighborhood of \(g\), its absolute
value stays at least \(c>0\), and the entries of the
first \(n\) columns of powers are bounded by \(B\ge1\),
by continuity. Take \(\eta<c/B^n\) in the conjugated
hypothesis and choose its open subset and finite
exclusions there. Lemma~\ref{lem:lift-residue-frobenius}
supplies a geometric Frobenius datum meeting them.
Expanding the altered determinant, every product has
absolute value at most \(B^n\eta<c\). This contradiction
proves that all these minors vanish everywhere on
\(s(G_E)\), with the conclusion vacuous if there are
no such row selections.

The rank argument in
Theorem~\ref{thm:spherical-annihilator-rank}, starting
with that simultaneous vanishing, uses only the local
diagonal family, freeness, and the diagonal span.
Indeed its polynomial substitution, followed by
Theorem~\ref{thm:diagonal-block-polynomial-test},
makes each weight-block compression scalar. The
coordinate projections in the span then force one
dimension per weight. It gives \(M=n\). The same
argument at order \(n\) in the proof of
Theorem~\ref{thm:spherical-annihilator-rank} shows
that every nonempty open set contains a point at
which \(I,s(g),\ldots,s(g)^{n-1}\) are independent.
These two conclusions precede the use of a polynomial
in the final coefficient comparison of that proof.

For arbitrary nonempty open \(V\subset G_E\) choose
such a point in \(V\) and the \(n\) entries witnessing
independence. On some nonempty open \(U\subseteq V\)
the resulting square submatrix of the power columns
is invertible and every entry of its inverse has
absolute value bounded by a constant \(C\ge1\), by
continuity of its adjugate divided by its determinant.
Write \(D_\gamma(X)=\det(XI-s(\gamma))\). It annihilates
\(s(\gamma)\) by the adjugate argument in
Lemma~\ref{lem:homogeneous-monic-divisor}.
The coefficient vector of \(P_v-D_\gamma\), of length
\(n\) in ascending powers, is sent by that submatrix
to the selected entries of \(P_v(s(\gamma))\).
Given a desired error \(\delta>0\), use the conjugated
hypothesis at \(\delta/C\) inside \(U\). Multiplying
by the bounded inverse makes every coefficient error
at most \(\delta\) for these data. The leading
coefficients already agree. Determinants are invariant
under conjugation, giving exactly the prescribed
quantifier for \(r\).
\end{proof}

\begin{corollary}\label{cor:finite-precision-before-central-twist}
For an approximating form in
Corollary~\ref{cor:twisted-hecke-quotient-seed} put
\(c_v=\iota(\omega(z_v))\) and
\(Q_v(X)=c_v^{-n}P_v(c_vX)\). Suppose a continuous
finite-dimensional \(K_\ell\)-representation \(V\)
has a local frame containing its \(n\) prescribed
weights free through \(n(n+1)/2\). It suffices for
attachment of this approximant that for every
\(\eta>0\),
\(\|Q_v(r_V(\gamma))\|\le\eta\) for most
Frobenius data. If the condition holds for all forms
in the list of Theorem~\ref{thm:definite-seed-attachment},
that theorem yields the representation for the initial
eigensystem.
\end{corollary}
\begin{proof}
Pass to the direct sum \(V_0\) of distinct simple
constituents as in
Lemma~\ref{lem:distinct-constituents-diagonal-span}.
In a fixed basis adapted to the invariant filtration,
the polynomial of each diagonal block is the diagonal
block of the polynomial of the full matrix. Its norm
is at most that of the full matrix in this basis.
Changing to this basis costs at most a fixed product
of two matrix norms, which is absorbed by decreasing
the chosen positive tolerance. The local-frame
argument of Theorem~\ref{thm:rank-before-scalar-twist}
retains all the required weights on \(V_0\), and
Theorem~\ref{thm:one-frame-cocharacter} gives its
diagonal family. The cited constituent lemma gives
the diagonal span. Theorem~\ref{thm:approximate-annihilator-rank}
therefore gives rank \(n\) and Frobenius polynomials
\(Q_v\) in the approximation sense.

Multiply the action on \(V_0\) by the continuous
character \(a_\omega\) of
Theorem~\ref{thm:actual-scalar-hecke-galois-character}.
Its Frobenius value is \(c_v\) outside a fixed finite
set. For any chosen coefficient \(j\le n\), the error
for the new characteristic polynomial is the error
for \(Q_v\) multiplied by \(c_v^{\,n-j}\), by scaling
the determinant. Inside any starting open set shrink
to a nonempty open on which \(|a_\omega|\) is bounded,
by continuity. Apply the approximation within it after
dividing the desired tolerance by a bound at least one
for \(|a_\omega|^{n-j}\). This proves the required
coefficient bound for \(P_v\). Applying the attachment
theorem to the list proves the last assertion.
\end{proof}

\section{A norm identity for the quadratic character}

\begin{definition}\label{def:quadratic-idele-cyclotomic-norm}
Let \(K\) be the imaginary quadratic field of
Definition~\ref{def:cm-elliptic-finite-frames} and let
\(p\) split into \(\mathfrak u,\bar{\mathfrak u}\) in
\(K\). Use the isomorphisms of both completions with
\(\Q_p\), induced by the respective extensions of the
identity on \(\Q\). For a fractional ideal
\(\mathfrak a\) put
\[
       N\mathfrak a=\prod_{\mathfrak v}
        |O_K/\mathfrak v|^{v_{\mathfrak v}(\mathfrak a)}
           \in\Q_{>0},
       \qquad
       \nu_p(x)=
        \frac{x_{\mathfrak u}x_{\bar{\mathfrak u}}}
             {N\mathfrak a_x}\in\Q_p^\times
\]
for a finite idele \(x\) and its fractional ideal
\(\mathfrak a_x\) from that definition. Denote by
\(\chi_{p,K}:G_K\to\Z_p^\times\) the action on all
\(p\)-power roots of unity in \(\bar K\), taking
compatible exponents as in
Definition~\ref{def:completed-local-action}.
\end{definition}

\begin{theorem}\label{thm:quadratic-norm-is-cyclotomic}
The character \(\nu_p\) descends continuously to
\(\mathcal C=\mathbb A_{K,f}^\times/K^\times\).
For the homomorphism \(r_K\) of
Proposition~\ref{prop:quadratic-cm-galois-translation}
one has
\[
                    \nu_p(r_K(g))=\chi_{p,K}(g)
                              \qquad(g\in G_K).
\]
\end{theorem}
\begin{proof}
The local exponents of \(\mathfrak a_x\) are those
of \(x\), by the ideal computation in
Proposition~\ref{prop:cm-elliptic-ray-classification}.
Thus the norm factor is multiplicative and is locally
constant on the finite ideles, being unchanged on a
coset of \(\widehat O^\times\).
For \(a\in K^\times\) the principal-ideal formula of
Proposition~\ref{prop:general-ideal-factorization}
gives \(N(aO)=a\bar a\), since the two embeddings
are complex conjugates and \(a\bar a>0\).
The product of the two \(p\)-adic embeddings is this
same rational number. Hence principal ideles have value
one, and continuity of the numerator gives continuity
on the quotient by its quotient topology.

At a degree-one prime \(\mathfrak v\) of rational norm
\(P\) outside the fixed exceptions in
Theorem~\ref{thm:quadratic-idele-geometric-frobenius},
the ramified primes of \(K\), and \(p\), that theorem
gives \(r_K(\gamma)=[t_{\mathfrak v} b]\) with \(b\)
a unit idele supported above \(P\). Such degree-one
primes split in \(K\) by the residue degree calculation
of Lemma~\ref{lem:finite-extension-residue-actions}.
The numerator of \(\nu_p\) is one and the ideal is
\(\mathfrak v\). Its value is therefore \(P^{-1}\).

The value of \(\chi_{p,K}\) on the same geometric
Frobenius is \(P^{-1}\in\Z_p^\times\). Indeed the
reduction of \(p^a\)-th roots of unity in residue
characteristic \(P\) is injective for all \(a\).
For a root reducing to one, the geometric sum
\(1+\zeta+\cdots+\zeta^{p^a-1}\) reduces to the
nonzero \(p^a\); multiplying it by \(\zeta-1\)
gives zero and forces \(\zeta=1\). Taking the
ratio of two roots proves injectivity. By
Definition~\ref{def:frobenius-data} the residue of
\(\gamma^{-1}\zeta\) is that of \(\zeta^P\).
Hence \(\gamma^{-1}\) has exponent \(P\) at every
level and \(\gamma\) has exponent \(P^{-1}\).
The character \(\chi_{p,K}\) is continuous because
each finite root group is contained in a finite
Galois extension. So both characters in the claimed
identity are continuous into a Hausdorff group.
Their equality set is closed and meets every nonempty
open by Lemma~\ref{lem:lift-residue-frobenius} and the
values just computed. It is the whole group.
\end{proof}

\section{Inertia of the central twist}

\begin{lemma}\label{lem:quadratic-conjugate-torsion-unramified}
Let \(p\nmid N_K\) split in \(K\), choose the prime
\(\mathfrak u\) and the residue map \(h\) as in
Theorem~\ref{thm:quadratic-idele-geometric-frobenius}.
For \(s\in I_h\), write \(r_K(s)=[b]\) with \(b\)
a unit idele supported above \(p\), uniquely as in
that theorem. In the level-six normalized coordinates
of \(Y_1\), the action on every torsion point is
\[
                         sP_1(t)=P_1(bt).
\]
For every \(a\ge1\), specialization at \(h\) is
injective on the subgroup labelled by the
\(\bar{\mathfrak u}\)-component of
\(\mathcal T_K[p^a]\). In particular
\(b_{\bar{\mathfrak u}}=1\).
\end{lemma}
\begin{proof}
The unit \(b\) is one at the primes of six and
\(\mathfrak a_b=O\). Its period group is thus the
same \(Y_1\), and \(P_b(t)=P_1(bt)\) by
Proposition~\ref{prop:cm-elliptic-ray-classification}.
In particular its normalizing labels of orders two
and three coincide with the original ones. The
normalized equations are identical by
Lemma~\ref{lem:rigid-two-three-cubic}. For any given
torsion point choose a level divisible by six and its
order. The action on the finite rigid tuple in
Proposition~\ref{prop:quadratic-cm-galois-translation}
identifies the conjugate point in these coordinates
with \(P_b(t)\). The origin is also preserved.
This proves the first formula on the same cubic.

Choose \(\alpha\in O\) congruent to zero modulo
\(\bar{\mathfrak u}^{\,a}\) and to one modulo
\(\mathfrak u\), using the coprime ideal decomposition
in Proposition~\ref{prop:general-ideal-factorization}.
All labels in the stated component are killed by
\(\alpha\), since \(p\) is a uniformizer at both
primes. On the model of the cubic over the valuation
ring in Lemma~\ref{lem:residue-map-to-valuation},
\([\alpha]\) extends by
Lemma~\ref{lem:fixed-quadratic-good-cubic-list}.
Its derivative on the generic tangent line at the
origin is the scalar \(\alpha\), by the complex
uniformization in
Proposition~\ref{prop:cm-elliptic-ray-classification}.
The identity conormal line is free with coordinate
\(t_0=X/Y\) by the proof of
Lemma~\ref{lem:good-cubic-prime-to-p-labels}.
The same scalar equality holds on this line over the
valuation ring since it embeds in its fraction field.
Consequently, near the closed identity,
\([\alpha]^*t_0=\alpha t_0+t_0^2 H\) with \(H\)
regular. If a point in the kernel reduces to the
identity, its parameter lies in the maximal ideal
and \(\alpha+t_0H\) evaluates to a unit.
The identity \([\alpha]^*t_0=0\) at the point then
forces \(t_0=0\). The local cubic equation identifies
the point with the origin, as in
Lemma~\ref{lem:good-cubic-prime-to-p-labels}. Taking
differences gives the required injectivity.

Every torsion point has algebraic homogeneous
coordinates by
Theorem~\ref{thm:quadratic-labelled-cubic-tuples-algebraic}.
Coordinates may be scaled to be integral at \(h\),
one a unit. The inertia action fixes the residue
of each such coordinate by
Lemma~\ref{lem:residue-map-to-valuation}. For labels
in the indicated component, \(P_1(bt)\) is in the
same subgroup and has identical reduction to \(P_1(t)\).
Injectivity and the bijective labelling imply \(bt=t\).
These labels run through
\(p^{-a}O_{K_{\bar{\mathfrak u}}}/O_{K_{\bar{\mathfrak u}}}\)
in that one component. Thus \(b_{\bar{\mathfrak u}}\)
is one modulo \(p^a\) for every \(a\), giving equality.
\end{proof}

\begin{theorem}\label{thm:quadratic-inertia-and-central-scalar}
For \(s\in I_h\) as in
Lemma~\ref{lem:quadratic-conjugate-torsion-unramified}
let \(e_{\mathfrak u}(z)\) denote the idele equal to
\(z\in\Z_p^\times=O_{K_{\mathfrak u}}^\times\) at
\(\mathfrak u\) and one elsewhere. Then
\[
                     r_K(s)=
                     [e_{\mathfrak u}(\chi_{p,K}(s))].
\]
For a character \(a_\omega\) of
Theorem~\ref{thm:actual-scalar-hecke-galois-character}
with coefficient prime \(\ell\ne p\), one has for
every \(s\in I_h\cap G_E\)
\[
            a_\omega(s)=
                \iota\bigl(\omega(e_{\mathfrak u}
                               (\chi_{p,K}(s)))\bigr).
\]
If no prime of \(E\) above \(\bar{\mathfrak u}\)
divides the defining level \(\mathfrak n\) of the
central test vector, this value is one for all such
\(s\).
\end{theorem}
\begin{proof}
The preceding lemma identifies the representative \(b\)
with one outside \(\mathfrak u\). Its ideal is the unit
ideal. Definition~\ref{def:quadratic-idele-cyclotomic-norm}
therefore gives \(\nu_p([b])=b_{\mathfrak u}\), and
Theorem~\ref{thm:quadratic-norm-is-cyclotomic} identifies
this coordinate with \(\chi_{p,K}(s)\). In the formula
for \(\omega_\ell\) in
Proposition~\ref{prop:corrected-quadratic-idele-character}
the two \(\ell\)-coordinates of \(b\) are one.
The composition defining \(a_\omega\) gives the
second formula, without any exclusion for the
conductor of \(\omega\) at \(p\).

For the last assertion, the scalar similitude
\(s(e_{\mathfrak u}(z))=\overline{e_{\mathfrak u}(z)}I\)
of Proposition~\ref{prop:central-character-actual-form}
has unit entries on its diagonal and equals the
identity except at the places over \(\bar{\mathfrak u}\).
Its multiplier is a rational profinite unit, since
the two conjugate factors have product the norm
at \(p\) and are one away from \(p\).
Under the level assumption this matrix is congruent
to the identity modulo \(\mathfrak n\) at every
dividing prime and belongs to \(K(\mathfrak n)\)
by Definition~\ref{def:moduli-signed-frame}.
It fixes the nonzero level vector with central
eigenvalue \(\omega(e_{\mathfrak u}(z))\), so that
eigenvalue is one.
\end{proof}

\section{A tame representative for the marked scalar}

\begin{proposition}\label{prop:marked-proprime-tame-generator}
In Theorem~\ref{thm:marked-tame-one} put \(p=g_0\),
\(m=p^{r!}-1\), and use its \(\xi,T\). Its inertia
element \(s\) may be chosen to have order prime to
\(p\) on every finite Galois extension of \(E\) in
\(\bar E\), while preserving its asserted ratios for
all divisors of \(m\). For this choice put
\[
   \alpha=T^{-1}\bigl(\exp(2\pi i/(p-1))\bigr)
                                     \in\mathbf F_p^\times .
\]
For \(p>2\), the \(p\)-adic cyclotomic value of \(s\)
is the unique \(b\in\Z_p^\times\) with \(b^{p-1}=1\)
and \(\bar b=\alpha\).
\end{proposition}
\begin{proof}
Take an element \(s_0\) and \(\xi\) supplied by
Theorem~\ref{thm:marked-tame-one}. On a finite normal
subextension where \(s_0\) has order \(N=p^a d\),
\(p\nmid d\), use the power \(s_0^k\) with
\(k\equiv0\bmod p^a\) and \(k\equiv1\bmod d\).
Existence and uniqueness modulo \(N\) follow from
Lemma~\ref{lem:basic-arith}. On any smaller normal
subextension the order divides \(N\), so these choices
are compatible. They define one field automorphism of
the union \(\bar E\), since restrictions and their
inverses agree and field operations occur in finite
subextensions. Its order on the selected extension
divides \(d\). Preservation of the valuation and
identity on residues hold element by element, since
they hold for every power of \(s_0\).

The \(m\)-th roots of unity are fixed by \(s_0\):
reduction on them is injective by the geometric-sum
argument of Theorem~\ref{thm:quadratic-norm-is-cyclotomic},
and inertia fixes their reductions. Thus with
\(\zeta=s_0(\xi)/\xi\) one has
\(s_0^k(\xi)/\xi=\zeta^k\). The root \(\zeta\) has
order \(m\) by the prescribed value under \(T\).
On a normal extension containing \(\xi,\zeta\) that
order divides \(N\) by the power formula, hence divides
\(d\). The chosen exponent preserves the ratio. Taking
powers of \(\xi\) gives all the stated ratios also
for \(s\).

Let \(\epsilon\) be a primitive \(p\)-th root of unity
in \(\bar E\) and normalize the valuation by \(v(p)=1\).
It reduces to one since \(X^p-1=(X-1)^p\) in the
residue field. For \(1\le j<p\) the quotient
\((1-\epsilon^j)/(1-\epsilon)\) is a unit reducing to
\(j\). Evaluating \(\prod_{j=1}^{p-1}(X-\epsilon^j)
=(X^p-1)/(X-1)\) at one gives
\(v(\epsilon-1)=1/(p-1)\).
The element \(t=\xi^{m/(p-1)}\) has the same valuation.
Inertia fixes the residue of the unit \((\epsilon-1)/t\).
Consequently
\[
     \overline{s(t)/t}=
         \overline{(s\epsilon-1)/(\epsilon-1)}
                          =\overline{\chi_{p,K}(s)} ,
\]
where \(K\subseteq E\) and we restrict the global
cyclotomic character of
Definition~\ref{def:quadratic-idele-cyclotomic-norm}.
The second equality is the preceding geometric sum
for the exponent of \(s\epsilon\) modulo \(p\).
The first value is \(\alpha\) by the ratio for
\(p-1\); it belongs to \(\mathbf F_p^\times\) since
its order divides \(p-1\) and that group contains
all the roots of \(X^{p-1}-1\).

On the roots of order \(p^j\), the order of the action
of \(s\) is prime to \(p\) by its finite quotient
property. The kernel of reduction of
\((\Z/p^j\Z)^\times\) onto \(\mathbf F_p^\times\)
has order \(p^{j-1}\), by counting lifts of one.
Thus the prime-to-\(p\) order in every such quotient
divides \(p-1\). It follows on taking compatible
residues that \(b=\chi_{p,K}(s)\) satisfies \(b^{p-1}=1\).
If two such values have the same reduction, their
quotient is a \((p-1)\)-st root reducing to one,
which is one by the geometric-sum calculation.
This proves uniqueness.
\end{proof}

\begin{corollary}\label{cor:actual-tame-central-value}
Fix \(K\) and the marked \(p=g_0\) of
Proposition~\ref{prop:marked-proprime-tame-generator}.
At sufficiently large scales \(p\nmid N_K\). For
the chosen \(\mathfrak u=w|_K\), any approximating
central character with coefficient prime \(\ell\ne p\)
and the above \(s\) satisfies
\[
                  a_\omega(s)
                      =\iota\bigl(\omega(e_{\mathfrak u}(b))\bigr),
        \qquad b^{p-1}=1,\quad \bar b=\alpha .
\]
In particular the order of \(a_\omega(s)\) divides \(p-1\).
\end{corollary}
\begin{proof}
The integer \(N_K\) depends only on the fixed field
and its finite level-six list by
Theorem~\ref{thm:quadratic-idele-geometric-frobenius}.
The lower bound for \(g_0\) in Theorem~\ref{thm:data}
tends to infinity. The selected place is split as
in Theorem~\ref{thm:marked-tame-one}. Apply
Theorem~\ref{thm:quadratic-inertia-and-central-scalar}
to the \(s\) of the proposition, and then raise its
value to the \(p-1\)-st power.
\end{proof}

\section{Coset multiplicities in the finite part of a period}
\begin{definition}\label{def:finite-period-cut}
Let \(H\) be a Hausdorff topological group, \(S\) a closed subgroup,
and \(J\subseteq A\) compact open subgroups of \(H\). Write
\(L=S\cap A\) and \(H/J=\{xJ:x\in H\}\) with its left \(H\)-action.
For a finite \(L\)-stable subset \(T\subseteq H/J\), define
\[
 m_T(t)=\#\{rL\in S/L:r^{-1}t\in T\}\qquad(t\in H/J).
\]
The membership test is independent of the representative because
\(T\) is \(L\)-stable. We call \(T\) a cut for \(W\subseteq H\)
if \(W/J:=\{wJ:w\in W\}\) lies in \(ST\).
\end{definition}

\begin{lemma}\label{lem:period-single-orbit-index}
In Definition~\ref{def:finite-period-cut}, \(m_T\) is constant on
each \(S\)-orbit and zero outside \(ST\). Fix \(t=xJ\) and put
\(S_t=S\cap xJx^{-1}\). If \(T\cap St=L t\), then
\[
                  m_T(t)=[S_t:S_t\cap L]<\infty .
\]
In particular \(m_T\) is finite everywhere for any finite \(T\).
\end{lemma}
\begin{proof}
Translation of \(rL\) on the left by \(s\in S\) identifies the
counting sets for \(t\) and \(st\). If \(T\cap St=Lt\),
then \(r^{-1}t=l t\) for some \(l\in L\) exactly when
\(r\in S_t L\). The cosets \(rL\) represented in this product
are parametrized by the left cosets of \(S_t\cap L\) in \(S_t\).
The latter is compact since \(S\) is closed, and the former
subgroup is open in it. Its disjoint open cosets have a finite
subcover. For arbitrary finite \(T\), split \(T\cap St\) into its
finitely many \(L\)-orbits. The same computation, with a
representative of each such orbit in place of \(t\) and then left
translation, gives a finite count for each piece.
\end{proof}

\begin{theorem}\label{thm:finite-period-cut-denominator}
Let \(J,A,S,H,L\) be as in Definition~\ref{def:finite-period-cut}.
Let \(D\subseteq H\) be compact, let \(B\) be a subgroup containing
\(A\) and \(D\), and suppose \(W\subseteq SD\). There is a cut
\(T\subseteq B/J\) for \(W\), independent of the integer \(Q\) below,
with the following property. If \(Q\ge1\) is an integer and for
each \(x\in D\) there is an integer \(a_x\ge0\) such that
\[
                (xux^{-1})^{Q^{a_x}}\in A\qquad(u\in J),
\]
then there is an integer \(c\ge0\) such that \(m_T(t)\mid Q^c\)
on \(ST\). In particular \(m_T(t)=1\) there when \(Q=1\).
\end{theorem}
\begin{proof}
The image of \(D\) in the discrete space \(H/J\) is finite.
From that image choose \(t_i=x_iJ,\ x_i\in D\), one in each
\(S\)-orbit met. Set \(T=\bigcup_i Lt_i\). Each of these orbits
is finite by compactness of \(L\) and discreteness of \(H/J\).
Since \(L\subseteq A\subseteq B\), the cut lies in \(B/J\).
It covers \(W/J\) because \(W\subseteq SD\).
Lemma~\ref{lem:period-single-orbit-index} computes the multiplicity
on \(St_i\) as the index of \(M_i=S_{t_i}\cap L\) in \(S_{t_i}\).

Let \(P_i\) be the image of \(S_{t_i}\) in the permutation group of
the finite set \(S_{t_i}/M_i\). Every \(g\in S_{t_i}\) is of
the form \(x_iux_i^{-1}\), \(u\in J\); its \(Q^{a_{x_i}}\)-th
power belongs to \(M_i\). Applying this to \(h^{-1}gh\) for
every \(h\in S_{t_i}\) shows that the same power of \(g\) acts
trivially on every coset. Every prime \(r\mid |P_i|\) divides
\(Q\). Indeed, cyclic rotation acts on the \(|P_i|^{r-1}\)
tuples of length \(r\) in \(P_i\) having product one. Orbits
have size \(1\) or \(r\), and fixed tuples are the constant
tuples \((v,\ldots,v)\) with \(v^r=1\). The all-identity tuple
therefore has another fixed tuple beside it. Its entry has
order \(r\), and its \(Q^{a_{x_i}}\)-th power is one.
Transitivity gives \([S_{t_i}:M_i]\mid |P_i|\) by counting
the fibers of evaluation at \(M_i\). Factorization in
Lemma~\ref{lem:basic-arith} and finiteness of the list of \(i\)
give a single \(c\). The empty cut permits \(c=0\).
\end{proof}

\section{Integration over a finite cut}
\begin{lemma}\label{lem:period-orbital-cutoff}
Let \(G_1\) be a second countable locally compact Hausdorff group and
\(S_1\) a closed subgroup equipped with left Haar measure \(\sigma_1\).
Here a left Haar measure is a left translation invariant Borel measure
finite on compact sets and positive on every nonempty open set.
There is a positive Borel function \(\alpha\) on \(G_1\) such that
\[
             \int_{S_1}\alpha(u^{-1}x)\dd\sigma_1(u)=1
                         \qquad(x\in G_1).
\]
\end{lemma}
\begin{proof}
The interiors of a sequence of compact sets \(C_i\), \(i\ge1\),
cover \(G_1\). Indeed choose relatively compact open neighborhoods
covering the group and extract a countable subcover by assigning
a cover member to each basis element contained in one. For given \(x\),
if \(u_0^{-1}x\in C_i\) then all \(u\) satisfying \(u^{-1}x\in C_i\)
belong to \(u_0(S_1\cap C_i C_i^{-1})\). Its measure \(b_i\)
is finite and independent of \(x,u_0\). Put
\(h=\sum_{i\ge1}2^{-i}{\bf1}_{C_i}/(1+b_i)\) and
\(c(x)=\int_{S_1}h(u^{-1}x)\dd\sigma_1(u)\).
Then \(c(x)\le1\), and \(c(x)>0\) because \(x\) belongs to the
interior of some \(C_i\). Product integration as in
Lemma~\ref{lem:product-integration} gives measurability of \(c\).
Changing \(u\) by left translation gives \(c(sx)=c(x)\) for
\(s\in S_1\). Hence \(\alpha=h/c\) works.
\end{proof}

\begin{theorem}\label{thm:period-cut-sum}
In addition to the groups in Definition~\ref{def:finite-period-cut}
and Lemma~\ref{lem:period-orbital-cutoff}, assume \(H\) is locally
compact and second countable. Fix a compact open subgroup \(K\) of
\(H\) and a level \(A\). Use left Haar measures on \(G_1,H,S_1,S\),
normalized at the finite places by \(\mu_H(K)=\sigma_S(L)=1\).
Put \(G=G_1\times H,\ U=S_1\times S\), with their product measures.
Let \(\Gamma\subseteq U\) be a countable subgroup acting on
\(G,U\) by left translation with measurable fundamental domains
\(D_G,D_U\), respectively. Fix a measurable \(a:G_1\to\C\).
There exists \(\kappa\in\C\) depending only on \(a\), the fixed
groups, measures and cutoff such that the following holds.

Suppose \(J\subseteq A\cap K\) is a compact open subgroup,
\(T\subseteq H/J\) is finite and \(L\)-stable, and \(f:G\to\C\)
is measurable, left \(\Gamma\)-invariant and integrable over \(D_G\).
Suppose \(f(x,h)\ne0\) implies \(hJ\in ST\). For a function
\(b:T\to\C\), suppose for almost every \((x,h)\) with \(hJ\in T\),
\[
       \int_{D_U} f(u(x,h))\dd u=[K:J]\,a(x)b(hJ).
\]
Then
\[
                 \int_{D_G} f(y)\dd y
                         =\kappa\sum_{t\in T}\frac{b(t)}{m_T(t)}.
\]
\end{theorem}
\begin{proof}
Write \(\psi(t)=1/m_T(t)\) on \(T\), zero elsewhere. The multiplicities
on \(T\) are positive and finite by Lemma~\ref{lem:period-single-orbit-index}.
For \(hJ\in ST\), the \(s\in S\) such that \(s^{-1}hJ\in T\) form
\(m_T(hJ)\) disjoint left cosets of \(L\), each of measure one.
On them \(\psi(s^{-1}hJ)=1/m_T(hJ)\). Thus the integral of this
function over \(S\) equals \({\bf1}_{ST}(hJ)\) for every \(h\).
Put \(A_0(x,h)=\alpha(x)\psi(hJ)\). The preceding computation and
Lemma~\ref{lem:period-orbital-cutoff} give
\(\int_U A_0(u^{-1}y)\dd u={\bf1}_{ST}(hJ)\) at \(y=(x,h)\).

For any nonnegative measurable left \(\Gamma\)-invariant \(F\),
Lemma~\ref{lem:product-integration} gives
\[
 \int_G A_0(g)\int_{D_U} F(ug)\dd u\,\dd g
       =\int_{D_G} F(y)\int_U A_0(u^{-1}y)\dd u\,\dd y.
\]
Indeed change \(g\) to \(u^{-1}g\) in the outer group integral,
split \(G=\coprod_{\gamma\in\Gamma}\gamma D_G\), and use
\(F(\gamma y)=F(y)\). The remaining integral over \(D_U\)
of \(\sum_\gamma A_0(u^{-1}\gamma y)\) is the integral over
\(U=\coprod_\gamma\gamma^{-1}D_U\), by left translation. This also
proves independence of the domains by the decomposition in
Lemma~\ref{lem:quotient-measure}. Taking \(F=|f|\) gives finiteness
on the left and the right. The same identity follows for \(f\) by
decomposing its real and imaginary parts. The right side now equals
\(\int_{D_G}f\) by the support hypothesis.

Every coset \(hJ\) has measure \([K:J]^{-1}\), by partitioning \(K\)
into translates of \(J\). Insert the given period formula on the
left. If \(b(t)\ne0\) for some \(t\), the integral with \(|f|\)
bounds \(\int_G A_0(g)|\int_{D_U} f(ug)\dd u|\,\dd g\).
The contribution from that coset shows \(\int_{G_1}\alpha|a|<\infty\).
We can take \(\kappa=\int_{G_1}\alpha a\) when this integral is
absolutely convergent, and \(\kappa=0\) otherwise. In the latter
case all \(b(t)\) vanish and both sides are zero. This single
choice of \(\kappa\) gives the formula for all \(J,T,f,b\).
\end{proof}

\begin{corollary}\label{cor:period-localized-coefficients}
Under the hypotheses of Theorem~\ref{thm:finite-period-cut-denominator},
use its cut in Theorem~\ref{thm:period-cut-sum} whenever the support
and period hypotheses there hold. All the finite coefficients
\(1/m_T(t)\) belong to \(\Z[1/Q]=\{z/Q^e:z\in\Z,\ e\ge0\}\).
Their sum is at most \(|T|\), and for \(Q=1\) each is exactly one.
\end{corollary}
\begin{proof}
Theorem~\ref{thm:finite-period-cut-denominator} gives
\(1/m_T(t)=(Q^c/m_T(t))/Q^c\) with integral numerator.
Every positive integer \(m_T(t)\) has reciprocal at most one.
Theorem~\ref{thm:period-cut-sum} uses precisely these reciprocals.
\end{proof}

\section{Local depth in the denominator of a finite period}
\begin{lemma}\label{lem:local-period-denominator}
Let \(F\) be a discretely valued field with valuation ring \(O\),
uniformizer \(\pi\), and residue field of order \(q\).
For \(h\ge1\) write \(K_h=\{g\in{\rm GL}_n(O):g\equiv I\bmod\pi^h\}\).
Then \(K_1/K_h\) has order \(q^{n^2(h-1)}\).
If \(x\in{\rm GL}_n(F)\) and \(d,d'\ge0\) satisfy
\(\pi^d x,\pi^{d'}x^{-1}\in M_n(O)\), then
\(xK_hx^{-1}\subseteq{\rm GL}_n(O)\) for \(h\ge\max(1,d+d')\).
\end{lemma}
\begin{proof}
For \(j\ge1\), reduction of \((g-I)/\pi^j\) modulo \(\pi\)
is an additive homomorphism from \(K_j\) with kernel \(K_{j+1}\).
The quadratic term in multiplication vanishes since \(2j\ge j+1\).
It is onto \(M_n(O/\pi)\) because \(I+\pi^j M\) has unit determinant
for every integral lift \(M\). Multiplying the successive indices
proves the order formula. For \(g\in K_h\) write
\[
 xgx^{-1}=I+\pi^{h-d-d'}(\pi^dx)
                         ((g-I)/\pi^h)(\pi^{d'}x^{-1}).
\]
This is integral and its determinant is a unit, so the adjugate
formula gives an integral inverse.
\end{proof}

\begin{proposition}\label{prop:adelic-period-denominator-depth}
Let \(E\) be a CM field, with conjugation \(w\mapsto\bar w\) on
its finite places. In standard coordinates let \(H\) be the
finite-adelic unitary group for the identity Gram matrix, with
\(K=\{x\in H:x_w\in{\rm GL}_n(O_w)\ \forall w\}\). Thus
\(x_{\bar w}\) is the inverse conjugate transpose of \(x_w\).
Let \(P\) be a finite set of split places with \(P\cap\bar P=\varnothing\),
\(B=\{x\in H:x_w\in{\rm GL}_n(O_w)\ (w\notin P\cup\bar P)\}\),
and \(J\subseteq K\) compact open, reducing to the identity at
every \(w\in P\). Let \(S\subseteq H\) be closed, \(D\subseteq B\)
compact and \(W\subseteq SD\). There is an \((S\cap K)\)-stable
cut \(T\subseteq B/J\) for \(W\) whose multiplicities divide
\[
                  \prod_{w\in P}q_w^{\,n^2(h_w-1)}
                          \qquad\text{on }ST ,
\]
where \(q_w=|O_w/\pi_w|\). Here take any integers \(d_w,d'_w\ge0\)
bounding the denominators of all \(x_w,x_w^{-1}\), respectively,
for \(x\in D\), and \(h_w=\max(1,d_w+d'_w)\).
\end{proposition}
\begin{proof}
For each \(w\), the entries of \(x_w\) and of its inverse range
over compact sets. The increasing open fractional ideals
\(\pi_w^{-r}O_w\), \(r\ge0\), cover them. This gives the bounds.
The local matrix rings are inverse limits of their finite
reductions, with the determinant-unit condition closed and open.
The product of the local integral groups is compact by successive
subsequence extraction in its countably many finite reductions.
The unitary equations are closed. Openness follows from the
restricted product topology of Definition~\ref{def:adelic-row}.

Choose representatives \(x_iJ\) and the union \(T\) of their
\((S\cap K)\)-orbits as in
Theorem~\ref{thm:finite-period-cut-denominator}, with \(A=K\).
Put \(S_i=S\cap x_i Jx_i^{-1}\). The map
\[
       s\longmapsto\big((x_i^{-1}sx_i)_w\bmod K_{h_w}\big)_{w\in P}
\]
is a homomorphism from \(S_i\) into \(\prod_{w\in P} K_1/K_{h_w}\)
in the respective local matrix groups. An element of its kernel
is integral at each \(w\in P\) by Lemma~\ref{lem:local-period-denominator},
and then at each conjugate place by the unitary equation. At all
other places \(x_i\) and \(x_i^{-1}s x_i\in J\) are integral
with integral inverses. Thus the kernel lies in \(S_i\cap K\).
Counting cosets in this chain shows that \([S_i:S_i\cap K]\)
divides the order of the finite target. This is exactly the
claimed bound by Lemmas~\ref{lem:period-single-orbit-index}
and~\ref{lem:local-period-denominator}.
\end{proof}

\begin{corollary}\label{cor:period-cut-first-and-empty}
For \(P=\varnothing\) in Proposition~\ref{prop:adelic-period-denominator-depth},
the cut lies in \(K/J\), every multiplicity on it is one, and
\(|T|\le[K:J]\). At any conductor, if \(J\subseteq A\subseteq K\)
are compact open, \(W\subseteq SA\), and the cut is formed using
\(A\) itself for \(D\) and \(S\cap A\) for \(L\), the same conclusions
hold with \(A/J\) and \([A:J]\). In both cases the finite coefficients
of Theorem~\ref{thm:period-cut-sum} are integers.
\end{corollary}
\begin{proof}
For the first assertion the product bound is one and \(B=K\).
In the second case \(x_i\in A\) and \(J\subseteq A\) imply
\(S_i\subseteq A\), so Lemma~\ref{lem:period-single-orbit-index}
gives multiplicity one. In each case sum over a subset of the
indicated finite coset space and apply Theorem~\ref{thm:period-cut-sum}.
\end{proof}

\section{Extension of primitive hermitian frames}
\begin{lemma}\label{lem:finite-unitary-primitive-frame}
Let \(k\) be a finite field of order \(q\). Let \(R\) be its quadratic
field extension with conjugation \(z\mapsto z^q\), or \(k\times k\)
with the two coordinates exchanged. Write \(A^*=\bar A^{\,t}\).
Let \(V=R^n\) with a nonsingular hermitian matrix \(J=J^*\).
If \(P,Q\in M_{n,r}(R)\) have left inverses and \(P^*JP=Q^*JQ\),
some \(g\in{\rm GL}_n(R)\) satisfies \(g^*Jg=J,\ gP=Q\).
\end{lemma}
\begin{proof}
The trace \(z\mapsto z+\bar z\) is onto \(k\). In the field case
its kernel has at most \(q\) elements by its polynomial and the
image lies in \(k\). In the product case use \((z,0)\).
Consequently each hermitian matrix \(D\) can be written \(M+M^*\):
choose lifts of the diagonal entries under the trace, copy the upper
off-diagonal entries and set the lower ones to zero.

All nonsingular hermitian forms on free \(R\)-modules of the same
rank are isometric. In the field case, a nonzero nonsingular space
contains a vector of nonzero norm. Otherwise expansion of the norm
of \(v+\lambda w\) gives \({\rm Tr}(\lambda v^*Jw)=0\) for all
\(\lambda,v,w\), contradicting trace surjectivity and nonsingularity.
The norm on \(R^\times\) maps onto \(k^\times\), since its kernel
has size at most \(q+1\) by the equation \(z^{q+1}=1\). Scale the
vector to have norm one and split off its orthogonal complement.
Induction gives an orthonormal basis. In the product case the form
is a perfect pairing of the two \(k\)-coordinate spaces, together
with its transpose. A basis on the first and its dual on the
second give the identity matrix.

The proposed isometry on the column spans is well defined because
both matrices have left inverses. In the first span split off a
complement \(U\) of the radical \(D_0\) of the restricted pairing.
The form on \(U\) is nonsingular because the quotient by the radical
has that property. In the product case choose a complement on each
coordinate. The radical dimensions there agree by equality of the
ranks of a pairing matrix and its transpose. Make the corresponding
choices in the second span by the prescribed isometry.
In the ambient orthogonal complement of \(U\), let \(v_1,\dots,v_e\)
be a basis of \(D_0\). The linear equations
\(v_i^*Jw_j=\delta_{ij}\) have solutions there since the form there
is nonsingular and the \(v_i\) are independent in each field
coordinate. Write \(v,w\) for these column matrices. The replacement
\(w\mapsto w-vM\), with \(M+M^*=w^*Jw\), makes the Gram matrix
of \(w\) zero while preserving \(v^*Jw=I\). Thus \(v,w\) form
a nonsingular summand of Gram matrix
\(\left(\begin{smallmatrix}0&I\\I&0\end{smallmatrix}\right)\).
Repeat for the image of the radical, using the ordered image basis.
Identify the two summands, use the given isometry on \(U\), and
identify the remaining nonsingular orthogonal complements by the
equal-rank assertion above. The direct sum extends the prescribed
map on every column.
\end{proof}

\begin{theorem}\label{thm:integral-unitary-primitive-frame}
Let \(O\) be a complete discrete valuation ring with finite residue
field and uniformizer \(\pi\). Let \(B\) be the ring of integers
in an unramified quadratic extension with residue conjugation of
order two, or \(O\times O\) with exchange of factors. Suppose
\(J\in{\rm GL}_n(B)\) is hermitian, \(P,Q\in M_{n,r}(B)\)
have left inverses over \(B\), and \(P^*JP=Q^*JQ\).
There is \(g\in{\rm GL}_n(B)\) with \(g^*Jg=J,\ gP=Q\).
\end{theorem}
\begin{proof}
Complete \(P\) to a basis \(C\) by extending its independent
columns after reduction in each residue field and lifting the
remaining columns. The determinant is then a unit.
Lemma~\ref{lem:finite-unitary-primitive-frame} gives a residue
isometry \(\bar g\) taking the reduced \(P\) to the reduced \(Q\).
Take a matrix \(C_1\) with first \(r\) columns exactly \(Q\)
and other columns lifting those of \(\bar g\bar C\). Thus
\(C_1^*JC_1\equiv C^*JC\bmod\pi\).
Suppose \(C_j\) is constructed with this congruence modulo
\(\pi^j\), \(j\ge1\), and with the prescribed first columns.
Reduce \((C^*JC-C_j^*JC_j)/\pi^j\) modulo \(\pi\) to obtain a
hermitian \(D\) with zero leading \(r\)-square block. The
uniformizer belongs to \(O\) and hence is fixed by conjugation.
Choose \(M+M^*=D\) with its first \(r\) columns zero. Take the
upper right block of \(D\) unchanged and apply the trace construction
of Lemma~\ref{lem:finite-unitary-primitive-frame} to the lower
right square block. Solve \((\bar C_j^*\bar J)\bar Z=M\) in
the residue algebra, then lift \(Z\) with its first \(r\) columns
zero. Expansion gives \(C_{j+1}=C_j+\pi^j Z\) with
\(C_{j+1}^*JC_{j+1}\equiv C^*JC\bmod\pi^{j+1}\), since
\(2j\ge j+1\). The limit \(C_\infty\) exists by completeness,
has the prescribed columns, and is invertible by reduction.
Now \(g=C_\infty C^{-1}\) has both required properties.
\end{proof}

\begin{corollary}\label{cor:finite-adelic-frame-extension}
For the arithmetic field \(E\) of Lemma~\ref{lem:real-place-rings},
put \(\widehat O_E=\prod_w O_w\). If two \(n\)-by-\(r\) matrices
over \(\widehat O_E\) have left inverses and the same hermitian Gram
matrix for the identity form, there exists a single element of
\(K\) in Proposition~\ref{prop:adelic-period-denominator-depth}
taking the first ordered frame to the second.
\end{corollary}
\begin{proof}
Group the local factors by the places of the real subfield.
Lemma~\ref{lem:real-place-rings} gives exactly the two integral
rings of Theorem~\ref{thm:integral-unitary-primitive-frame},
including above two. Apply that theorem at every factor with
\(J=I\). Their product belongs to \(K\) because each local
isometry and its inverse are integral.
\end{proof}

\section{Bessel orbits and the quotient lines of a flag}
\begin{definition}\label{def:string-flag-ideals}
Over the field \(E\) of Lemma~\ref{lem:real-place-rings}, let \(V\)
have a nonsingular hermitian form of matrix \(J_0\) and a basis
consisting of a basis of \(X\) and a string \(e_0,\ldots,e_{s-1}\),
\(s\ge0\). Assume \(X\) is perpendicular to the string and its
string entries \(h(e_i,e_j)\) vanish unless \(i+j=s-1\).
Fix an invertible finite adelic frame \(C\) with \(C^*J_0C=I\).
Use the standard finite-adelic \(H,K\) of
Proposition~\ref{prop:adelic-period-denominator-depth}.
Write \(r=\lceil s/2\rceil\), \(F_{-1}=0\), and
\(F_i=\operatorname{span}_E(e_0,\ldots,e_i)\), \(0\le i<r\).
The finite Bessel subgroup \(S\) consists of \(u\in H\) such that
\(CuC^{-1}e_i-e_i\in F_{i-1}\otimes_E\mathbb A_{E,f}\) for
each \(0\le i<r\). Here \(\mathbb A_{E,f}\) is the restricted
product of the finite completions of \(E\) with respect to their
valuation rings. These triangular conditions define a closed
subgroup, since they give closed linear equations, persist under
products and hold for inverses by triangular substitution.

At a finite place \(w\) put \(\Lambda_x(w)=C_w x_w O_w^n\) for
\(x\in H\), \(n=\dim_E V\). Write \(I_i(x;w)\) for the image of
\(\Lambda_x(w)\cap F_i(E_w)\) in \(F_i(E_w)/F_{i-1}(E_w)=E_w e_i\),
identified with a fractional ideal in \(E_w\) by its coefficient
of \(e_i\). Denote its exponent by \(b_i(x;w)\), so that
\(I_i(x;w)=\pi_w^{b_i(x;w)}O_w\).
\end{definition}

\begin{lemma}\label{lem:projected-string-frame}
The ideals in Definition~\ref{def:string-flag-ideals} are nonzero.
For any generators \(\beta_i(w)\) of them one can choose
\(v_i(x;w)\in\Lambda_x(w)\cap F_i(E_w)\) with coefficient
\(\beta_i(w)\) on \(e_i\). These columns extend to a basis of
\(\Lambda_x(w)\). Their hermitian Gram matrix, assembled over
the conjugate places, is zero except possibly for the last diagonal
entry when \(s\) is odd. That entry is
\(\bar\beta_{r-1}\beta_{r-1} h(e_{r-1},e_{r-1})\).
\end{lemma}
\begin{proof}
The lattice contains \(\pi_w^a O_w^n\) and is contained in
\(\pi_w^{-a}O_w^n\) in the displayed rational basis for some \(a\).
Thus the image is a nonzero bounded fractional ideal. Its least
valuation gives a generator in a discrete valuation ring.
Lifts \(v_i\) exist by definition. Successively subtract multiples
of these lifts, starting at the largest index, to see that they
are a basis of \(\Lambda_x(w)\cap F_{r-1}(E_w)\), with the empty
list used for \(r=0\). Their reductions in
\(\Lambda_x(w)/\pi_w\Lambda_x(w)\) are independent. Indeed if a
combination belongs to \(\pi_w\Lambda_x(w)\), dividing by \(\pi_w\)
places it in the same flag intersection. Uniqueness of its basis
coordinates forces all original coefficients to be divisible by
\(\pi_w\). Extend the reduced list to a basis and lift its additional
columns. The resulting determinant in lattice coordinates is a unit.
Finally every index in the early columns is at most \(r-1\). Two
such indices can sum to \(s-1\) only when \(s\) is odd and both
equal \(r-1\). The claimed Gram formula follows by expansion.
\end{proof}

\begin{theorem}\label{thm:finite-bessel-flag-classification}
In Definition~\ref{def:string-flag-ideals}, two elements \(x,x'\)
lie in the same double coset in \(S\backslash H/K\) if and only
if \(I_i(x;w)=I_i(x';w)\) for all \(i<r,w\).
\end{theorem}
\begin{proof}
Right translation by \(K\) preserves the lattice. Left translation
by \(u\in S\) sends it to its translate by \(CuC^{-1}\), which
preserves \(F_i\) and acts as one on \(F_i/F_{i-1}\).
This proves necessity.

For sufficiency choose the same generators \(\beta_i(w)\) for both
lattices, taking one at almost every place since the frames there
are invertible integral matrices in the displayed basis.
Lemma~\ref{lem:projected-string-frame} gives the early columns \(v,v'\)
and equality of their Gram matrices. In the respective standard
coordinates \(P=(Cx)^{-1}v,\ Q=(Cx')^{-1}v'\), the columns have
left inverses over every \(O_w\), because each frame extends to
a basis of its lattice. Their local left inverses assemble over
\(\widehat O_E\). Both frames use the identity form in these coordinates.
Corollary~\ref{cor:finite-adelic-frame-extension} gives \(k\in K\)
taking \(P\) to \(Q\). Hence \(C(x'kx^{-1})C^{-1}\) takes every
\(v_i\) to \(v'_i\). It preserves the flag by their triangular
bases, and on \(F_i/F_{i-1}\) it sends \(\beta_i e_i\) to the
same element. As each \(\beta_i(w)\ne0\), all these quotient actions
are identity. Thus \(x'kx^{-1}\in S\), as required.
\end{proof}

\section{Compact representatives at the twist places}
\begin{theorem}\label{thm:near-flag-adelic-period-cut}
Use the string flag of Definition~\ref{def:string-flag-ideals}.
Let \(P\) be a finite set of split places disjoint from \(\bar P\),
and \(k\ge0\) an integer. Suppose \(W\subseteq H\) satisfies,
for every \(x\in W\) and \(0\le i<r\),
\[
 |b_i(x;w)-b_i(1;w)|\le
     k\,{\bf1}_{P\cup\bar P}(w)\qquad\text{at every finite place }w.
\]
Then there is a compact set \(D\subseteq B\) as in
Proposition~\ref{prop:adelic-period-denominator-depth} such that
\(W\subseteq SD\). It is the union of at most
\((2k+1)^{r|P\cup\bar P|}\) translates of \(K\), each represented
by a group element which is identity outside \(P\cup\bar P\).
In particular, for every compact open \(J\subseteq K\) reducing to
one at \(P\), the same \(D\) yields a cut with coefficients in
\(\Z[1/Q]\), where \(Q=\prod_{w\in P}q_w\).
\end{theorem}
\begin{proof}
The possible tuples of differences at the indicated places belong
to a product of \(r|P\cup\bar P|\) intervals of integers of length
\(2k+1\). Choose a representative \(x\in W\) of each occurring
tuple. Keep its components at \(P\cup\bar P\), replacing all others
by identity, and denote the result by \(x^\circ\). The replacement
respects the unitary equation because it is made on entire conjugation
orbits. The definition of \(b_i\) uses only \(C_w x_w\). Thus the new
representative has the old tuple at the indicated places and the
identity tuple elsewhere. For any \(y\in W\) of that type all its
line ideals agree with those of \(x^\circ\).
Theorem~\ref{thm:finite-bessel-flag-classification} places it in
\(Sx^\circ K\). Take the union of the \(x^\circ K\) for \(D\).
Proposition~\ref{prop:adelic-period-denominator-depth} and
Corollary~\ref{cor:period-localized-coefficients}, or the displayed
integer product in that proposition directly, give the asserted
localization.
\end{proof}

\begin{proposition}\label{prop:marked-stage-cut-size}
Keep \(E\) and the marked split primes \(g_0,g_1\) of
Proposition~\ref{prop:row-group}. Write \(f=[E:\Q]/2\). For
a unitary group in standard finite coordinates of rank \(n\)
let \(A_g\subseteq K\) be the kernel of reduction at the \(f\)
marked places over each \(g_d\), one in every conjugate pair.
Then
\[
       [K:A_g]=\prod_{d=0}^1\prod_{u=0}^{n-1}
                   (g_d^n-g_d^u)^f\ \le (g_0g_1)^{f n^2}.
\]
For \(P=\varnothing\) in Theorem~\ref{thm:near-flag-adelic-period-cut}
and \(J=A_g\), the sum of the absolute finite coefficients of
Theorem~\ref{thm:period-cut-sum} is at most this index. For a
string of length zero the same assertions hold with \(r=0\).
\end{proposition}
\begin{proof}
The residue fields and the number of marked places are those
of Proposition~\ref{prop:row-group}. At a marked split place choose
any invertible residue matrix and lift its entries to the integers.
The lift has unit determinant. Use its inverse conjugate transpose
at the conjugate place and identity outside the chosen pairs to see
that simultaneous reduction is onto the full product. Its kernel
is \(A_g\). In an \(n\)-dimensional space over a field of \(g_d\)
elements, the next column after \(u\) independent columns can
be any of the \(g_d^n-g_d^u\) vectors outside their span. This
gives the order and its upper bound. For empty \(P\) all representatives
in the proof of Theorem~\ref{thm:near-flag-adelic-period-cut}
are identity. Apply Corollary~\ref{cor:period-cut-first-and-empty}.
If \(s=0\) there are no flag conditions, \(S=H\), and the same
proof applies.
\end{proof}

\section{The arithmetic frames for the row lifts}
\begin{lemma}\label{lem:actual-lift-scale}
Use \(g=g_0g_1,\pi_g,\pi_N,\varpi,\mathcal A\) from
Theorem~\ref{thm:data}, Propositions~\ref{prop:level-prime}
and~\ref{prop:different}, and Definition~\ref{def:frame-lattice}.
Write \(\omega=(1+\sqrt{-d})/2\), \(D_c=P'_{N_1}(\zeta_{N_1})\)
and \(\jmath=-1/\sqrt{-d}\). Let \(\alpha_w=D_c^{-1}\) where
\(\omega\) reduces to zero, and let \(\alpha_w=1\) elsewhere.
Let \(\delta_g\) have component \(\bar\pi_g\) at the marked places
dividing \(\bar\pi_g\), and component one elsewhere.
The finite ideles
\[
     a_f=\alpha/\bar\delta_g,\qquad
     u_j=(a_f/\bar a_f)^{\lfloor(m-j+1)/2\rfloor}
                        \quad(0\le j\le m)
\]
satisfy \(u_j\bar u_j=1\). At each \(w\), \(a_f\bar a_f\)
has the same valuation as \(c=-(\jmath g\varpi)^{-1}\).
\end{lemma}
\begin{proof}
By Proposition~\ref{prop:local-cyclo-different}, \(D_c\) is a
unit away from \(N_1\) and has the same valuation on conjugate
branches. At a prime dividing \(N_1\), exactly one branch reduces
\(\omega\) to zero by Proposition~\ref{prop:frame-lattice-trace}.
Hence \(\alpha\bar\alpha\) generates \(D_c^{-1}\widehat O_E\),
in agreement with Lemma~\ref{lem:frame-ideal-product}.
Lemma~\ref{lem:split-level} gives the same fractional ideal for
\(\delta_g\bar\delta_g\) and \(g\). Proposition~\ref{prop:different}
gives \(\jmath\varpi=-\zeta_{N_1}^{h'_0}D_c\), a unit times \(D_c\).
Dividing the local ideals proves the valuation identity.
Conjugating \(a_f/\bar a_f\) inverts it, which proves the norm claim.
\end{proof}

\begin{definition}\label{def:actual-row-lift-spaces}
Let \(r=k+m-1\) be the row rank of Proposition~\ref{prop:row-group},
and let \(t\) be the block size of the arithmetic data.
For \(0\le j<m\) put \(s=m-1-j\), \(d_j=k+j(t+1)\).
Use a definite space \(X_j=E^{d_j}\) with identity Gram matrix
and a string of length \(s\), perpendicular to \(X_j\).
Its antidiagonal entries are one, except that a middle entry
of an odd string is \(-N\). Denote the resulting rational space
by \(Y_j\), with dimension \(d_j+s=r+jt\). Set
\(Y_m=E^{r+mt}\) with the positive identity form.
In \(Y_j\) let \(b\) be the last basis vector of \(X_j\), and
use \(e_0,\ldots,e_{s-1}\) for the string. Define the following
ordered integral finite-adelic lattice \(\Lambda_j\) by its basis,
with scales from Lemma~\ref{lem:actual-lift-scale}.
For even \(s\), take the standard basis of \(X_j\) and string
columns \(u_j a_f^{-1}e_i\) for \(2i<s\), \(u_j\bar a_f e_i\)
for \(2i\ge s\). For odd \(s\), start with the standard basis of
\(X_j\) and the columns \(u_j e_i\), then replace \(b\) by
\[
                 v=b/\bar\pi_N+u_j e_{(s-1)/2}/N .
\]
For \(j=m\) use the standard integral lattice.
\end{definition}

\begin{theorem}\label{thm:actual-row-lift-frames}
For \(j<m\) the number of negative real directions of \(Y_j\)
at each infinite place is \(\lceil s/2\rceil\).
Its lattice has unimodular hermitian Gram matrix \(G_j\). In
the basis matrix \(B_j\) of Definition~\ref{def:actual-row-lift-spaces}
there is a change \(T_j\in{\rm GL}_{r+jt}(\widehat O_E)\)
such that \(C_j=B_jT_j\) is a finite frame for the identity form.
One may prescribe \((T_j)_w=I\) at every marked place
\(w\mid\bar\pi_g\). For \(0\le i<\lceil s/2\rceil\) the
reference ideals \(I_i(1;w)\) of Definition~\ref{def:string-flag-ideals}
are generated by \((u_j)_w\) if \(s\) is odd, and by
\((u_j/a_f)_w\) if \(s\) is even. At \(j=m\) the negative count
is zero and the identity frame works.
\end{theorem}
\begin{proof}
In each distinct paired string plane \((e_i,e_{s-1-i})\), the
sum and difference divided by \(\sqrt2\) have respective norms
one and minus one. For odd \(s\) the middle vector divided by
\(\sqrt N\) has norm minus one. These vectors together with the
standard \(X_j\)-basis are orthogonal, giving the claimed signature.
In finite coordinates the scaled distinct string pairs still
have pairing one by \(u_j\bar u_j=1\). In the odd case the Gram
block on \(v,u_j e_{(s-1)/2}\) equals
\(\left(\begin{smallmatrix}0&-1\\-1&-N\end{smallmatrix}\right)\)
by \(\pi_N\bar\pi_N=N\). The remaining definite basis vectors
are orthonormal. This gives an integral Gram matrix with unit
determinant, including for the empty string.

At a split pair choose the correction on one member equal to
identity, taking the marked member if present. On its conjugate
choose the correction with conjugate transpose \(G_{j,w}^{-1}\).
Both resulting Gram matrices are identity by hermitian symmetry.
At a nonsplit place Lemma~\ref{lem:real-place-rings} gives an
unramified quadratic integer extension. Choose a residue basis
with Gram matrix identity by the classification in
Lemma~\ref{lem:finite-unitary-primitive-frame}. Lift this basis
with Gram congruent to identity modulo the uniformizer. Apply
the correction in Theorem~\ref{thm:integral-unitary-primitive-frame}
with zero prescribed columns, using \(I\) as target Gram matrix:
at step \(h\) write the reduced error as \(M+M^*\), solve for
the correction and pass to the limit. All the resulting corrections
and their inverses are integral, giving the asserted adelic frame.

The corrected basis generates the original lattice. In the odd
case, a vector in an early flag span must have zero coefficient
of \(v\), since its \(b\)-coordinate vanishes. Its independent
early columns are therefore precisely the \(u_j e_i\).
The even case follows from the diagonal columns directly.
Projection to consecutive lines gives the ideals as stated.
\end{proof}

\section{The column level inside the finite Bessel subgroup}
\begin{proposition}\label{prop:column-level-in-stage-bessel}
For \(0\le j<m\) let \(K_{c,j}\subset U(X_j)(\mathbb A_f)\) consist
of integral unitary matrices reducing to identity at the marked
places over \(g\) and satisfying
\(b^*(h-I)\in\pi_N\widehat O_E^{1\times d_j}\), where \(b\) is
the last coordinate vector in Definition~\ref{def:actual-row-lift-spaces}.
Under extension by the identity on the string, \(K_{c,j}\) lies
in the Bessel subgroup, preserves \(\Lambda_j\), and acts as identity
on its reduction at every marked place over \(g\). These conditions
hold in the frames \(C_j\) of Theorem~\ref{thm:actual-row-lift-frames}.
\end{proposition}
\begin{proof}
The conditions define a subgroup. The row condition states that
\(b^* h=b^*\) in the reduction, hence also \(b^*h^{-1}=b^*\)
there. Taking adjoints gives \((h-I)b\in\bar\pi_N\widehat O_E^{d_j}\).
The ideals of \(\pi_N,\bar\pi_N\) have disjoint support by
Lemma~\ref{lem:split-level} and \(\pi_N\bar\pi_N=N\).
The last entry of \((h-I)b\) therefore belongs to
\(N\widehat O_E\), since it belongs to both ideals. Applying the
same facts to \(h^{-1}\) will give the reverse lattice inclusion.

All string vectors are fixed, giving Bessel membership by
Definition~\ref{def:string-flag-ideals}. For even string length
the standard definite lattice and the scaled string lattice
are direct summands, so preservation follows. For odd length
the pure definite vectors in \(\Lambda_j\) have integral
coordinates, with the last coordinate in \(\pi_N\widehat O_E\).
In fact a multiple \(zv\) of the glued column has last definite
coordinate \(z/\bar\pi_N\). Its string component can be cancelled
by an integral multiple of \(u_j e_{(s-1)/2}\) precisely when
\(z\in N\widehat O_E\). For an ordinary definite column \(e\)
other than \(b\), \((h-I)e\) meets these pure definite conditions
by the row congruence. On \(v\) the difference is
\((h-I)b/\bar\pi_N\), integral by the adjoint condition and
with last entry divisible by \(\pi_N\) by the preceding product
divisibility. This proves both inclusions.

At a marked place over \(g\), \(\pi_N,\bar\pi_N\) are units by
Proposition~\ref{prop:column-dim}. In the scaled lattice basis,
a pure definite difference with coordinates \((d',d_b)\) has,
in the odd case, coefficients \(d',\bar\pi_N d_b,-d_b/\pi_N\)
in the ordinary definite, glued, and middle string columns.
Every such difference coming from \(h-I\) is zero on reduction,
also after division by \(\bar\pi_N\) on \(v\). The even case is
immediate on its direct summands. Changes of integral basis
preserve the identity reduction.
\end{proof}

\begin{lemma}\label{lem:initial-frame-central-translation}
Let \(\Lambda^{\rm old}\) be the row lattice of
Definition~\ref{def:row-space}, using the string order in
Definition~\ref{def:actual-row-lift-spaces}. Let \(A^{\rm old}\)
be its marked principal level and \(A^{\rm new}\) the corresponding
level on \(\Lambda_0\), viewed in the same rational coordinates.
Write
\[
      b_0=\operatorname{diag}(u_0^{-1} I_{X_0},I_{\rm string}).
\]
Then \(b_0\in S_f\), where \(S_f\) denotes the Bessel subgroup in
rational coordinates, and \(A^{\rm new}=b_0 A^{\rm old} b_0^{-1}\).
If \(P(y,h)\) is right \(A^{\rm old}\)-invariant and supported over
\(S_f A^{\rm old}\) in its finite variable, \(P'(y,h)=P(y,h b_0)\)
is right \(A^{\rm new}\)-invariant and supported over \(S_f A^{\rm new}\).
Right translation intertwines left multiplication by every rational
element, so it also commutes with absolutely convergent left
periodizations.
\end{lemma}
\begin{proof}
The string length at zero is odd since \(m=2\nu+2\).
The finite isometry acting by \(u_0\) on the string and by one
on \(X_0\) takes the old basis, including its glued column,
to precisely the basis of \(\Lambda_0\). It factors as the
central scalar \(u_0 I\) times \(b_0\), which fixes every string
vector. The level is characterized by lattice preservation and
trivial action on the specified reduced lattices, so transport of
basis gives the conjugation formula. The two assertions on functions
now follow from \(b_0\in S_f\) and \(A^{\rm old} b_0^{-1}
=b_0^{-1} A^{\rm new}\). Left and right translations commute termwise.
\end{proof}

\section{The initial Bessel function of the row columns}
\begin{theorem}\label{thm:actual-base-bessel-periodization}
Use the row group, rational points \(\Gamma\), old principal
level \(A^{\rm old}=K_Y\), and class representatives \(h_c\)
of Definition~\ref{def:adelic-row}. Let \(S(F)\) be the rational
Bessel subgroup in Corollary~\ref{cor:full-bessel-identity}.
Choose for each column class a complex linear functional
\(\lambda_c\) on the value space of \(F_\infty\) in
Definition~\ref{def:full-arithmetic-amplitudes}. Define
\[
 \begin{split}
 b(y,h)&=\sum_c {\bf1}_{h_c A^{\rm old}}(h)
                   \,\beta_c(y)\lambda_c(F_\infty(y)),\\
 P(y,h)&=\sum_{\delta\in S(F)}
                   b(\delta_\infty y,\delta_f h).
 \end{split}
\]
The latter series converges absolutely, is continuous, is left
\(S(F)\)-invariant and right \(A^{\rm old}\)-invariant, and vanishes
unless \(h\in S_f A^{\rm old}\). Its periodization over
\(S(F)\backslash\Gamma\) equals the automorphic column with these
coordinates in Theorem~\ref{thm:actual-full-poincare-columns}.
For all \(d\) and \(y\),
\[
               P(y,h_d)=\lambda_d(F_\infty(y)).
\]
Moreover \(P(n_\infty y,n_f h)=
\chi_\infty(n_\infty)\chi_f(n_f)P(y,h)\) for the initial
unipotent group and the characters of Proposition~\ref{prop:full-phase}.
\end{theorem}
\begin{proof}
The finite class set is finite by Proposition~\ref{prop:column-dim}.
Absolute convergence follows by restricting the absolutely convergent
sum in Theorem~\ref{thm:actual-full-poincare-columns}. Partition its
terms by the left cosets of \(S(F)\) to obtain the asserted
periodization. Reindexing gives rational left invariance, and the
right level follows from the indicators. Each \(h_c\) belongs
to \(S_f\) by its definition as a middle isometry, proving the
support condition. When \(h=h_d\) the proof of
Corollary~\ref{cor:full-bessel-identity} identifies the contributing
rational elements as precisely \(\Gamma_d\), with class \(c=d\).
On this group the character is one by Proposition~\ref{prop:full-phase},
and \(F_\infty\) transforms by this character by
Theorem~\ref{thm:actual-real-bessel-vector}. The cutoffs sum to
one by Proposition~\ref{prop:boxed-arithmetic-cutoff}. This gives
the displayed evaluation.

Every \(h\in S_f A^{\rm old}\) has the form
\(\delta_f h_d a\), with \(\delta\in S(F)\) and \(a\in A^{\rm old}\).
Indeed, first factor the Bessel component as \(u_f t_f\) for a
unipotent and a definite middle isometry. To see this, the preserved
isotropic flag gives an induced middle isometry on
\(P_0^\perp/P_0\), with \(P_0\) the span of the outer isotropic
vectors. It fixes the class of \(z\), so acts as a definite
isometry on its orthogonal complement. After removing this isometry
the action on the successive flag quotients is identity, also on
the dual lines by unitarity. This is the unipotent group of
Definition~\ref{def:unip}. Decompose the definite part as
\(t_f=\eta_f h_d k_c\) by Proposition~\ref{prop:column-dim}.
The embedded \(k_c\) belongs to \(A^{\rm old}\) by
Proposition~\ref{prop:marked-local-level}. Approximate \(u_f\)
by \(v_f\) for \(v\in U(F)\) using Lemma~\ref{lem:unip-approx},
close enough that \(u_f\eta_f h_d\in v_f\eta_f h_d A^{\rm old}\).
This gives the required \(\delta=v\eta\).

On this coset rational left invariance gives
\[
       P(y,h)=\lambda_d(F_\infty(\delta_\infty^{-1}y)).
\]
Thus \(P\) is continuous there, and it vanishes on the remaining
open level cosets. The character \(\chi_\infty\) is unchanged
under conjugation by \(S(F)\). Unipotent conjugation preserves
every group character. Conjugation by a middle isometry preserves
the simple-root entries defining the character because it fixes
each of the string vectors and their coordinate functionals.
The transformation law for \(F_\infty\) therefore proves the
assertion for \(n_\infty\) on each coset. For \(v\in U(F)\),
rational invariance gives
\(P(y,v_f h)=P(v_\infty^{-1}y,h)=\chi_f(v_f)P(y,h)\),
using Proposition~\ref{prop:full-phase}. Density in
Lemma~\ref{lem:unip-approx} and continuity give the finite assertion
everywhere. The real and finite actions commute.
\end{proof}

\begin{corollary}\label{cor:actual-base-stage-cut}
For the stage-zero function \(P'(y,h)=P(y,h b_0)\) in
Lemma~\ref{lem:initial-frame-central-translation}, one may take
\(D=A^{\rm new}\) in the finite period cut. This works simultaneously
for all ramified columns, all twist levels \(J\subseteq A^{\rm new}\),
and every product \(P'z\) with a function \(z\). The cut lies in
\(A^{\rm new}/J\), and \(m_T=1\) on its saturation by the Bessel group.
\end{corollary}
\begin{proof}
Theorem~\ref{thm:actual-base-bessel-periodization} gives the old
support for all coordinate functionals. Transport it with
Lemma~\ref{lem:initial-frame-central-translation}. Multiplying by
\(z\) preserves support containment. Apply
Corollary~\ref{cor:period-cut-first-and-empty} with \(A=A^{\rm new}\).
\end{proof}

\section{Duality on opposite string lines}
\begin{definition}\label{def:full-flag-line-ideals}
For \(Y_j\) of Definition~\ref{def:actual-row-lift-spaces}, \(j<m\),
order the blocks as the first \(\lceil s/2\rceil\) string lines,
the definite block \(X_j\), and the remaining string lines in
increasing order of their string index. Denote the span of all blocks
preceding \(e_i\) by \(V_i^-\), and let \(V_i^+=V_i^-+E e_i\).
For the local lattice \(L_w=(C_j)_w x_w O_w^{r+jt}\) at \(x\in H\),
define \(I_i(x;w)\) by the image of \(L_w\cap V_i^+(E_w)\)
in \(V_i^+/V_i^-\), and \(b_i(x;w)\) by its valuation as in
Definition~\ref{def:string-flag-ideals}. For \(i<\lceil s/2\rceil\)
these are the ideals and exponents of that definition. For every
\(i\) write \(i^*=s-1-i\).
\end{definition}

\begin{proposition}\label{prop:opposite-string-line-duality}
In Definition~\ref{def:full-flag-line-ideals}, for every \(x,w\)
and \(i\ne i^*\) one has
\[
                 b_i(x;w)+b_{i^*}(x;\bar w)=0.
\]
For \(i=i^*\) one has
\(b_i(x;w)+b_i(x;\bar w)\ge-v_w(N)\). A negative sum is
therefore possible only for odd \(s\) and a conjugate pair over
\(N\), in which case it equals \(-1\).
\end{proposition}
\begin{proof}
The local lattices are in perfect hermitian duality between
\(w,\bar w\), using conjugation to identify coefficient fields.
Indeed their coordinate matrices are those of the frame in
Theorem~\ref{thm:actual-row-lift-frames} multiplied by a standard
unitary \(x\), so their pairing in integral coordinates is the
identity matrix. For \(i\ne i^*\), direct inspection of the
antidiagonal form gives
\[
               (V_i^-)^\perp=V_{i^*}^+,\qquad
               (V_i^+)^\perp=V_{i^*}^- .
\]
The central definite block belongs to a prefix exactly when its
paired prefix excludes it. A possible odd middle vector is handled
by the same observation for these noncentral \(i\). On the
two line quotients the induced pairing of the indicated basis
vectors is one. Pairing any two integral lifts therefore gives
\(b_i(x;w)+b_{i^*}(x;\bar w)\ge0\).

Choose a lattice lift \(v\) of a generator of \(I_i(x;w)\).
Choose a basis of \(L_w\cap V_i^-(E_w)\) by successive projection
onto local coordinate lines, as in Lemma~\ref{lem:finite-quotient-gens}.
Together with \(v\) it is a basis of \(L_w\cap V_i^+(E_w)\).
This list extends to a basis of \(L_w\) by the reduction argument
of Lemma~\ref{lem:projected-string-frame}, since the intersection
with a vector subspace is saturated. Assign value one to \(v\)
and zero to the other members of the extended basis. Perfect
duality gives a vector in the opposite integral lattice representing
this functional, with the conjugate-linear convention on the
appropriate side. It annihilates \(V_i^-\), so lies in
\(V_{i^*}^+\). The pairing of its quotient coordinate with the
generator is one. This gives the opposite inequality.

For the middle line choose the early lifts of
Lemma~\ref{lem:projected-string-frame} in both components.
Their pairing has valuation
\(v_w(N)+b_i(x;w)+b_i(x;\bar w)\) and is integral.
By Lemma~\ref{lem:split-level}, \(N\) is of valuation one at
its split places, and of valuation zero elsewhere. The exponents
are integers.
\end{proof}

\begin{corollary}\label{cor:opposite-stage-errors}
For noncentral pairs the deviations
\(\Delta_i(x;w)=b_i(x;w)-b_i(1;w)\) satisfy
\(\Delta_i(x;w)=-\Delta_{i^*}(x;\bar w)\).
For a middle line \(\Delta_i(x;w)+\Delta_i(x;\bar w)\ge0\)
unless the moved lattice has the exceptional value \(-1\)
from Proposition~\ref{prop:opposite-string-line-duality}.
\end{corollary}
\begin{proof}
Subtract the two identities in
Proposition~\ref{prop:opposite-string-line-duality}. For a middle
line the reference generator is \(u_j\) by
Theorem~\ref{thm:actual-row-lift-frames}, of norm one by
Lemma~\ref{lem:actual-lift-scale}. Its pair of valuations sums to zero.
\end{proof}

\section{A local test at the middle line}
\begin{proposition}\label{prop:exceptional-middle-level-stabilizer}
Let \(j<m\) have odd string length \(s\), and let
\(i_0=(s-1)/2\). Suppose
\(b_{i_0}(x;w)+b_{i_0}(x;\bar w)=-1\) at a pair above \(N\).
Let \(K^{(w)}\subseteq H\) be the subgroup supported on that
pair and integral there in the standard finite coordinates of \(Y_j\).
Let \(K_X^{(w)}\) be the analogous subgroup on \(X_j\), embedded
by acting trivially on the string in rational coordinates and
then conjugating into the standard coordinates. There exists a
single \(r_0\in S\) for which
\[
                     (r_0x)^{-1}K_X^{(w)}(r_0x)\subseteq K^{(w)} .
\]
\end{proposition}
\begin{proof}
Change the lattice only on the indicated pair. In rational
coordinates give \(X_j\) its standard lattice there. For every
early \(i\) choose generators \(\beta_i\) of the two component
ideals of \(x\), and use \(\beta_i e_i\) as a basis vector in
the new lattice. For \(i<i_0\) give its partner the basis
vector \(\bar\beta_i^{-1} e_{i^*}\). The distinct pairs have
Gram block \(\left(\begin{smallmatrix}0&1\\1&0\end{smallmatrix}\right)\).
The middle vector has Gram entry \(-N\beta_{i_0}\bar\beta_{i_0}\),
a unit by the assumed sum and Lemma~\ref{lem:split-level}.
This is a unimodular lattice on the split pair. Choosing a basis
in one component and the dual basis in the other identifies its
pairing with the identity, as in Theorem~\ref{thm:actual-row-lift-frames}.
Together with the original coordinates outside the pair this gives
an element \(x'\in H\) whose moved lattice is the new one.

Its early flag ideals are visibly those of \(x\) at the pair and
are unchanged elsewhere. Theorem~\ref{thm:finite-bessel-flag-classification}
gives \(x'=r_0x k_0\) for \(r_0\in S,\ k_0\in K\).
The new lattice is fixed by the entire \(K_X^{(w)}\): its definite
summand is standard and its string vectors are fixed.
Consequently \(x'^{-1}K_X^{(w)}x'\subseteq K^{(w)}\).
Conjugation by \(k_0\) preserves this local direct factor of \(K\),
giving the assertion.
\end{proof}

\begin{lemma}\label{lem:single-central-difference-test}
For the odd string of Proposition~\ref{prop:exceptional-middle-level-stabilizer},
let \(U\) be the matrices whose difference from identity strictly
raises the grades \(s-1-2i\) on the string, with all grades zero
on \(X_j\). Let \(H_X\subset S\) be the embedded finite definite
group, and \(\chi_U:U\to\C^\times\) a character invariant under
\(H_X\). Fix \(z\in K_X^{(w)}\) central in \(H_X\).
Let \(\phi\mapsto f_\phi\) be a linear family of functions on \(H\)
and \(R\) a linear operator on their parameter space. Assume
\[
  f_\phi(uh)=\chi_U(u) f_\phi(h),\quad
  f_\phi(hk)=f_\phi(h),\quad f_{R\phi}(h)=f_\phi(zh)
       \qquad(u\in U,\ k\in K^{(w)},\ h\in H).
\]
At every exceptional \(x\) in that proposition one has
\(f_\phi(x)=0\) for \(\phi\in{\rm im}(1-R)\).
\end{lemma}
\begin{proof}
Here \(S=U H_X\). The positive string grades span the isotropic
flag, and its perpendicular modulo its span has the orthogonal
middle decomposition \(X_j\oplus E e_{i_0}\). An element of \(S\)
acts as one on the line generated by \(e_{i_0}\) in this quotient.
Removing the induced action on \(X_j\) leaves identity on all
graded pieces, including the negative ones by duality. It preserves
the filtration by its flag conditions and unitarity. This proves
the decomposition, also when the positive flag is empty.
The character \(\widetilde\chi(u h_X)=\chi_U(u)\) is a character
of \(S\) because the decomposition is unique and \(H_X\) normalizes
\(U,\chi_U\). For \(r_0\) of
Proposition~\ref{prop:exceptional-middle-level-stabilizer} the
element \(v=r_0^{-1}zr_0 z^{-1}\) belongs to \(U\) by centrality
in the quotient. Since it is a commutator, \(\chi_U(v)=1\).
There is \(k\in K^{(w)}\) with \(v z x=x k\). Therefore
\(f_{R\psi}(x)=f_\psi(zx)=f_\psi(vzx)=f_\psi(xk)=f_\psi(x)\).
Subtract and use linearity.
\end{proof}

\begin{corollary}\label{cor:ramified-middle-test}
For the ramified columns of Theorem~\ref{thm:ram-dim}, at every
marked pair over \(N\) their inclusion in \({\rm im}(1-R)\) in
Lemma~\ref{lem:single-central-difference-test} holds when \(R\)
is translation by the corresponding residue generator. Thus the
two equivariances and the right level in that lemma suffice to
rule out a middle-line defect for such columns.
\end{corollary}
\begin{proof}
The action of each finite residue unit group on the ramified
columns has no trivial character by Theorem~\ref{thm:ram-dim}.
On a character summand its generator acts by a scalar different
from one, and \(1-R\) is invertible there. The finite character
decomposition follows by Lemma~\ref{lem:characters}.
Apply Lemma~\ref{lem:single-central-difference-test}.
\end{proof}

\section{A lattice inequality for a map between consecutive gradings}
\begin{definition}\label{def:row-string-lowering-map}
Consider \(Y_j,Y_{j+1}\) of Definition~\ref{def:actual-row-lift-spaces}
with \(s=m-1-j\ge2\), so that the next string has length \(s'=s-1\).
Use grades \(g_i=s-1-2i\) on the source string and
\(g'_i=s-2-2i\) on the target string, with grade zero on both definite
blocks. Define \(T:Y_j\to Y_{j+1}\) by zero on \(X_j\) and on
\(e_{s-1}\), and by \(T e_i=t_i e'_i\) for \(0\le i<s'\).
Put \(h_i=h(e_{i^*},e_i)\), equal to \(-N\) when \(g_i=0\) and to
one otherwise. The scalars \(t_i=h_i \tau_i\) are given by
\[
 \tau_i=\begin{cases}
 c,& s\text{ even and }g_i>0,\text{ or } s\text{ odd and }g_i\le0,\\
 1,&\text{otherwise},
 \end{cases}
 \qquad c=-(\jmath g\varpi)^{-1}.
\]
Every nonzero entry of \(T\) lowers the grade by one. Write
\(\mathcal N\) for the space of maps whose entries of negative
grade difference vanish. For a finite twist set \(P\) as in
Theorem~\ref{thm:near-flag-adelic-period-cut}, let \(\delta_P\)
be a finite idele with uniformizer components at \(P\) and one
elsewhere. Put \(d_P=\delta_g\delta_P\) and
\(a_P=\bar d_P/\alpha\), using Lemma~\ref{lem:actual-lift-scale}.
Thus \(v_w(a_P)=-v_w(a_f)+{\bf1}_{\bar P}(w)\).
\end{definition}

\begin{proposition}\label{prop:dual-column-one-step}
Take \(x,x'\) at the two consecutive stages, with moved local lattices
\(L_w,L'_w\) in the notation of Definition~\ref{def:full-flag-line-ideals}.
Suppose a finite adelic map \(Z\in T+\mathcal N(\mathbb A_{E,f})\)
satisfies \(a_P Z L_w\subseteq L'_w\) at every place, with the
local component of the scalar and map understood. Then, for every
\(0\le i<s'\),
\[
 b'_i(x';w)\ \le\ v_w(a_P)+v_w(\tau_i)-b_{s-1-i}(x;\bar w).
\]
Here primes on \(b\) indicate the target space.
\end{proposition}
\begin{proof}
At \(\bar w\) take an adapted lattice vector \(v\in V_{i^*}^+\)
whose \(e_{i^*}\)-coefficient \(\beta\) generates
\(I_{i^*}(x;\bar w)\). Use the extension of a lattice basis in
Proposition~\ref{prop:opposite-string-line-duality} to choose
\(y\in L_w\) with \(h(v,y)=1\), annihilating \(V_{i^*}^-\).
From the antidiagonal form, all components of \(y\) have grade
at least \(g_i\). Indeed all string components of lower grade
pair with earlier string vectors than \(e_{i^*}\), and if \(g_i>0\)
then \(X_j\) also occurs before \(e_{i^*}\).
The \(e_i\)-coefficient of \(y\) is \(1/(h_i\bar\beta)\) because
the earlier terms of \(v\) pair to zero with \(y\). When \(g_i=0\)
there may also be components of \(y\) in \(X_j\).

The map \(Z-T\) has image of grade at least \(g_i\) on \(y\).
In grade \(g_i-1\) of \(Ty\) the sole component is
\(\tau_i e'_i/\bar\beta\), since \(T\) kills \(X_j\). Hence
\(Zy\) lies in the target prefix \((V'_i)^+\), with this coefficient
on \(e'_i\). In particular if \(g_i-1=0\) its grade-zero component
has zero projection onto the target definite block, which follows
the middle string line in the block order.
Since \(a_P Zy\in L'_w\), its projected coefficient belongs to
\(I'_i(x';w)\), proving the upper bound by taking valuations.
\end{proof}

\begin{lemma}\label{lem:initial-line-identity}
For the reference lattices the following equality holds for all
the indices of Proposition~\ref{prop:dual-column-one-step}:
\[
 b'_i(1;w)=v_w(\tau_i)-v_w(a_f)-b_{s-1-i}(1;\bar w).
\]
Consequently for the deviations of
Corollary~\ref{cor:opposite-stage-errors} that proposition gives
\[
            \Delta'_i(x';w)
             \le -\Delta_{s-1-i}(x;\bar w)+{\bf1}_{\bar P}(w).
\]
\end{lemma}
\begin{proof}
The defining basis in Definition~\ref{def:actual-row-lift-spaces}
is triangular for the full block order, with invertible diagonal
blocks over the local fields. Thus each string diagonal generates its
projected ideal: working backward, a vector in a prefix must have
zero coefficients on all later basis columns. Denote these diagonal
generators by \(\gamma_i\) and by \(\gamma'_i\) on the next string.
If \(s=2h+1\), then \(u_{j+1}=u_j\) and each \(\gamma_i=u_j\).
For \(i<h\) the next generator is \(u_j/a_f\) and \(\tau_i=1\).
For \(i\ge h\) the next generator is \(u_j\bar a_f\) and
\(\tau_i=c\). If \(s=2h\), then
\(u_{j+1}=u_j\bar a_f/a_f\), and this generates every target
string ideal. For \(i<h\) the reversed source generator is
\(u_j\bar a_f\) and \(\tau_i=c\). For \(i\ge h\) the reversed
generator is \(u_j/a_f\) and \(\tau_i=1\).
In all four cases the ideal generated by \(\gamma'_i\) equals
the one generated by \(\tau_i/(a_f\bar\gamma_{s-1-i})\), using
\(u_j\bar u_j=1\) and equality of the valuations of \(c\) and
\(a_f\bar a_f\) from Lemma~\ref{lem:actual-lift-scale}.
This gives the first identity. Subtract it from
Proposition~\ref{prop:dual-column-one-step} and insert
\(v_w(a_P)\) from Definition~\ref{def:row-string-lowering-map}.
\end{proof}

\section{One-sided movement of the string ideals}
\begin{theorem}\label{thm:one-sided-stage-line-interlacing}
Under Proposition~\ref{prop:dual-column-one-step}, assume the
middle-line exception of Proposition~\ref{prop:opposite-string-line-duality}
does not occur at either of the two points. Then at every finite
place \(w\) and every \(0\le i<s-1\) one has
\[
 \Delta_{i+1}(x;w)-{\bf1}_{P}(w)
     \ \le\ \Delta'_i(x';w)\
     \le\ \Delta_i(x;w)+{\bf1}_{\bar P}(w).
\]
Consider stages \(0,\ldots,j\) with \(0\le j<m\), starting at
\(x_0\) with all deviations zero. If the map and middle-line
conditions hold at every transition with nonempty target string,
then
\[
       -j\,{\bf1}_{P}(w)\le\Delta_i(x_j;w)
                      \le j\,{\bf1}_{\bar P}(w)
                  \qquad(0\le i<m-1-j).
\]
\end{theorem}
\begin{proof}
Write \(i^*=s-1-i\) in the source and \(i^\dagger=s-2-i\) in the
target. By Corollary~\ref{cor:opposite-stage-errors},
\(\Delta_i(x;w)+\Delta_{i^*}(x;\bar w)\ge0\).
Combining this with Lemma~\ref{lem:initial-line-identity} gives
the stated upper bound. Apply that lemma at the target index
\(i^\dagger\) and place \(\bar w\), and negate the inequality.
The lower bound follows from
\(\Delta'_i(x';w)\ge-\Delta'_{i^\dagger}(x';\bar w)\) by the
same corollary and \(s-1-i^\dagger=i+1\).
The second assertion follows by induction on the number of steps.
At empty target string it is vacuous.
\end{proof}

\begin{corollary}\label{cor:one-sided-chain-period-denominators}
For fixed \(j,P\) let \(W\) be any set of stage-\(j\) finite points
reached by chains satisfying Theorem~\ref{thm:one-sided-stage-line-interlacing}.
Its compact representatives in Theorem~\ref{thm:near-flag-adelic-period-cut}
can be chosen in at most
\[
          (j+1)^{\,\lceil(m-1-j)/2\rceil\,|P\cup\bar P|}
\]
translates of the standard integral subgroup. The cut for the
common stage level, the kernel of reduction at the marked places
and \(P\), has every coefficient in \(\Z[1/Q]\) with
\(Q=\prod_{w\in P}q_w\). When \(P=\varnothing\) the sum of its
absolute coefficients is at most \((g_0g_1)^{f(r+jt)^2}\).
\end{corollary}
\begin{proof}
At each place in the indicated union exactly one indicator in
Theorem~\ref{thm:one-sided-stage-line-interlacing} can be one.
There are at most \(j+1\) possible integer values there and a
single value elsewhere. Apply the representative construction
of Theorem~\ref{thm:near-flag-adelic-period-cut} with this count.
Proposition~\ref{prop:adelic-period-denominator-depth}
controls the coefficients. For empty \(P\) use
Proposition~\ref{prop:marked-stage-cut-size} and the stage rank
in Definition~\ref{def:actual-row-lift-spaces}.
\end{proof}

\section{The annihilator of the positive-grade lattice}
\begin{definition}\label{def:hom-oscillator-pairing}
For consecutive \(Y=Y_j,Y'=Y_{j+1}\) set \(V={\rm Hom}_E(Y,Y')\).
For maps \(u,v\) let
\[
 H(u,v)={\rm tr}(J^{-1}u^*J'v),\qquad
 [u,v]_p={\rm Tr}_{E\otimes\Q_p/\Q_p}\big(\jmath H(u,v)\big),
\]
where \(J,J'\) are their rational Gram matrices. This is a
hermitian form and the second pairing is alternating and
\(\Q_p\)-bilinear. Indeed cyclic permutation under matrix trace
gives \(H(v,u)=\overline{H(u,v)}\), and \(\bar\jmath=-\jmath\).
Write \(\mathcal P\) for the space of strictly positive-grade maps,
and \(\mathcal N\) for the nonnegative-grade maps as in
Definition~\ref{def:row-string-lowering-map}. For a finite adelic
additive subgroup \(D\) put
\[
 D^\circ=\{v\in V(\mathbb A_{E,f}):
            \psi_f(\jmath H(u,v))=1\text{ for all }u\in D\}.
\]
The additive character \(\psi_f\) is defined in
Proposition~\ref{prop:additive}.
\end{definition}

\begin{lemma}\label{lem:hom-oscillator-dual-lattices}
Let \(L,L'\) be the source and target lattices moved by \(x,x'\),
and let \(\mathcal M={\rm Hom}_{\widehat O_E}(L,L')\) in the
rational finite coordinates, meaning the adelic linear maps taking
\(L\) into \(L'\). Then \(\alpha\mathcal M\) is self-dual for
Definition~\ref{def:hom-oscillator-pairing}, and
\((d_P\alpha\mathcal M)^\circ=
(\alpha/\bar d_P)\mathcal M\).
Over \(\Q_p\) the exact orthogonal complement of
\(\mathcal P(E\otimes\Q_p)\) for \([\, ,\,]_p\) is
\(\mathcal N(E\otimes\Q_p)\).
\end{lemma}
\begin{proof}
In the moved isometric coordinates of
Theorem~\ref{thm:actual-row-lift-frames}, \(\mathcal M\) consists
of integral matrices and \(H(u,v)\) is the sum of corresponding
products \(\bar u_{ij}v_{ij}\). By Lemma~\ref{lem:actual-lift-scale}
and Proposition~\ref{prop:different},
\(\jmath\alpha\bar\alpha\widehat O_E=\varpi^{-1}\widehat O_E\),
the trace dual of \(\widehat O_E\).
Proposition~\ref{prop:additive} applied entrywise gives exact
self-duality, and moving the factor \(d_P\) across the pairing
gives the second assertion. In particular the local rational
pairing is nonsingular: if a vector paired to zero against all
vectors, all its rational scalar multiples would lie in a fixed
dual lattice, forcing it to vanish.

In the rational graded coordinates \(J,J'\) pair opposite grades,
including perfect pairings on grade zero. The same holds for
their inverses. Expansion of the trace gives zero between map
grades \(d,e\) unless \(d+e=0\). The pairing between each grade
and its opposite is perfect by the nonsingularity just proved and
the direct-sum decomposition. Thus annihilation of every strictly
positive grade is precisely vanishing of all negative-grade
components.
\end{proof}

\begin{proposition}\label{prop:positive-line-oscillator-annihilator}
With \(D=d_P\alpha\mathcal M\) and
\(P_f=\mathcal P(\mathbb A_{E,f})\), one has
\[
       (D\cap P_f)^\circ
                 =(\alpha/\bar d_P)\mathcal M
                         +\mathcal N(\mathbb A_{E,f}).
\]
In particular, if \(\psi_f(\jmath H(u,T))=1\) for every
\(u\in D\cap P_f\), there exists a \(Z\) satisfying the full
map hypothesis of Proposition~\ref{prop:dual-column-one-step}.
\end{proposition}
\begin{proof}
At each rational prime \(p\) let \(D_p\) be the local lattice,
\(P_p=\mathcal P(E\otimes\Q_p)\). A character test confined to
\(p\) says exactly \([u,v]_p\in\Z_p\) for all \(u\in D_p\cap P_p\),
by Proposition~\ref{prop:additive}. The latter intersection is a
full lattice of \(P_p\), since multiplying any finite basis by
a large power of \(p\) puts it in \(D_p\). It is saturated in
\(D_p\): if \(p y\in P_p\) then \(y\in P_p\). Choose its
\(\Z_p\)-basis by successive coordinate projections as in
Lemma~\ref{lem:finite-quotient-gens}. Extend this to a basis of
\(D_p\) by lifting a complementary residue basis as in
Lemma~\ref{lem:projected-string-frame}. Extend the \(\Z_p\)-valued
functional \(u\mapsto[u,v]_p\) on the intersection by giving it
value zero on the extra basis vectors. Its rational extension is
represented by a vector \(z_p\) via the nonsingular local pairing
of Lemma~\ref{lem:hom-oscillator-dual-lattices}.
Then \(z_p\) is in the dual lattice and \(v_p-z_p\) annihilates
\(P_p\) exactly because it vanishes on its spanning lattice.
Here \(v_p\) is the local component of \(v\).
At almost all \(p\), \(v_p\) itself is in the dual lattice, since
the coordinate frames and all idele scalars are integral units
or invertible integral matrices there. We may use \(z_p=v_p\)
there, so the components assemble into finite adeles.
Lemma~\ref{lem:hom-oscillator-dual-lattices} gives the stated
decomposition. Both summands on the right pass the character test
on the left, proving the converse. For \(v=T\) take \(Z=(z_p)_p\).
Then \(Z-T\in\mathcal N(\mathbb A_{E,f})\) and
\((\bar d_P/\alpha)Z\in\mathcal M\), as required.
\end{proof}

\section{A nonzero constant term and its finite frequency}
\begin{definition}\label{def:adelic-oscillator-translations}
Let \(V,H,\jmath\) be as in Definition~\ref{def:hom-oscillator-pairing}.
Write \(\mathbb A_E\) for the full adeles. Use the additive
character
\[
 \Psi(v)=\exp\!\big(-2\pi i\,{\rm Tr}_{E_\infty/\R}(v_\infty)\big)
                \psi_f(v_f),\qquad E_\infty=E\otimes_\Q\R .
\]
It is one on \(E\) by Proposition~\ref{prop:additive}. With
\(\omega=(1+\sqrt{-d})/2\) put
\(B(x,y)=\Psi(\jmath\bar\omega H(y,x))\) on \(V(\mathbb A_E)\).
Use \(B_f\) for its restriction to finite vectors, identified with
full vectors having zero infinite component.
An adelic theta function here means a continuous \(F\) on
\(V(\mathbb A_E)\) obeying
\(F(v+r)=B(r,v)^{-1}F(v)\) for \(r\in V(E)\).
Define \((D_zF)(v)=B(v,z)F(v+z)\). For the positive-grade
subspace \(\mathcal P\) and rational \(T\) of
Definition~\ref{def:row-string-lowering-map}, its constant term
at \(T\) is
\[
 I_T(F)=\int_{\mathcal P(E)\backslash\mathcal P(\mathbb A_E)}
                             (D_TF)(v)\dd v
\]
with quotient probability measure.
\end{definition}

\begin{lemma}\label{lem:theta-positive-translation}
The integral in Definition~\ref{def:adelic-oscillator-translations}
exists. Each \(D_z\) preserves the transformation condition there and
\[
       D_xD_y=B(x,y)D_{x+y},\qquad
       D_xD_y=\Psi(\jmath H(y,x))\,D_yD_x .
\]
For \(u\in\mathcal P(\mathbb A_E)\), \(I_0(D_uF)=I_0(F)\).
\end{lemma}
\begin{proof}
The function \(B\) is multiplicative under addition in either
variable. Factoring \(B(v+r,z)=B(v,z)B(r,z)\) cancels the added
factor in \(B(r,v+z)^{-1}\), so \(D_zF\) obeys the same condition.
Expansion gives the product formula. To compute the commutator use
\(H(x,y)=\overline{H(y,x)}\), invariance of trace under conjugation,
\(\bar\jmath=-\jmath\) and \(\omega+\bar\omega=1\).
Hence \(B(x,y)/B(y,x)=\Psi(\jmath H(y,x))\).
All the formulas use the integral trace-one element \(\omega\)
of Lemma~\ref{lem:actual-lift-scale}.

The hermitian pairing vanishes between two positive-grade maps
by Lemma~\ref{lem:hom-oscillator-dual-lattices}, or by the same
grade expansion over \(E\). Thus on this subspace \(B=1\),
and the restriction of \(D_TF\) is rationally periodic.
In any rational basis of its underlying \(\Q\)-space, one
coordinate of the additive adelic quotient has representatives
\([0,1)\times\widehat\Z\). To find one, clear the finitely many
finite denominators by a rational translation by the Chinese
remainder construction of Lemma~\ref{lem:unip-approx}, then
translate by an integer to choose the real representative.
A difference of two representatives which is rational is an
integer at the finite places, and equals zero by the real interval.
The quotient is compact as the image of
\([0,1]\times\widehat\Z\) in each coordinate. On this domain
take product of length measure and the measures giving every
residue coset modulo \(n\) in \(\widehat\Z\) mass \(1/n\).
These finite measures are compatible under reduction.
Equivalently use independent uniform digits in the \(p\)-adic
expansions at all primes. Addition permutes the residues at each
finite level and hence preserves this measure on its generating
cylinder sets, by the measure extension argument of
Lemma~\ref{lem:product-integration}. The product measure on
the fundamental domain gives an invariant probability on the
quotient: translation moves the domain to another such domain
and its pieces are translated back by rational elements as in
Lemma~\ref{lem:quotient-measure}. The same construction works
in finitely many coordinates, including the zero-dimensional case.
Continuity and compactness give integrability. Translation on this
quotient now proves the last identity.
\end{proof}

\begin{theorem}\label{thm:constant-term-implies-lattice-map}
With the lattices \(L,L'\) and \(\mathcal M\) of
Lemma~\ref{lem:hom-oscillator-dual-lattices}, suppose
\(D_u F=F\) for \(u\in d_P\alpha\mathcal M\) placed at the finite
adeles. If \(I_T(F)\ne0\), the local lattice inequality of
Proposition~\ref{prop:dual-column-one-step} holds at every place.
\end{theorem}
\begin{proof}
For \(u\in d_P\alpha\mathcal M\cap\mathcal P(\mathbb A_{E,f})\),
Lemma~\ref{lem:theta-positive-translation} gives
\[
 I_T(F)=I_0(D_uD_TF)=\psi_f(\jmath H(T,u)) I_0(D_TD_uF)
                         =\psi_f(\jmath H(T,u)) I_T(F).
\]
Nonvanishing forces the character to be one. Conjugating the
hermitian pairing inverts this character, so the same holds with
\(H(u,T)\). Proposition~\ref{prop:positive-line-oscillator-annihilator}
now gives a map \(Z\) with \(Z-T\) of nonnegative grade and
\((\bar d_P/\alpha)Z L\subseteq L'\). Apply
Proposition~\ref{prop:dual-column-one-step}.
\end{proof}

\section{The inclusion datum at a split level}
\begin{definition}\label{def:split-inclusion-finite-oscillator}
For consecutive stage ranks \(n,n'\), use the reference finite
isometric frames of Theorem~\ref{thm:actual-row-lift-frames}.
In the resulting matrix coordinates let
\(L_0=\alpha M_{n',n}(\widehat O_E)\), with its trace pairing
from Definition~\ref{def:hom-oscillator-pairing}. On compactly
supported locally constant functions of the full finite matrix
variable use \(D_z f(v)=B_f(v,z)f(v+z)\), as in
Definition~\ref{def:adelic-oscillator-translations}. Put
\(v_0={\bf1}_{L_0}\), and let \(\mathcal H_f\) be the complex span
of all \(D_z v_0\).

Let \(M\) be the union of the marked places over \(g\) and the
twist set \(P\), with \(P\cap\bar P=\varnothing\) and
\(\prod_{w\in P}q_w\) prime to \(gN\). Thus \(M\cap\bar M=\varnothing\).
Write \(L_+,L_-\) for \(L_0\) supported on \(M,\bar M\),
respectively, and set
\[
       C_+=L_+/d_P L_+,\qquad
       C_-=\bar d_P^{-1}L_-/L_- .
\]
For \(a\in C_+\), \(z\in C_-\), put
\(\chi_z(a)=\psi_f(\jmath H(z,a))\) using any lifts.
At each marked or twist place identify the two stage index sets
with an initial block of \(n\) and its enlargement by \(t\) indices,
as in their ranks in Definition~\ref{def:actual-row-lift-spaces}.
Denote by \(a_{\rm inc}\in C_+\) the residue matrix of the
standard inclusion \(\binom{I_n}{0}\) in the coordinates divided
by \(\alpha\).
\end{definition}

\begin{proposition}\label{prop:finite-oscillator-residue-name}
The groups \(C_+,C_-\) are in perfect duality by \(\chi\), and have
order \(D=((g_0g_1)^f Q)^{nn'}\), \(Q=\prod_{w\in P}q_w\).
The vector
\[
      f_a=D^{-1}\sum_{z\in C_-}\chi_z(-a)\,D_z v_0
                         \qquad(a\in C_+)
\]
is well defined. On the span of the \(D_zv_0\), \(z\in C_-\),
the map \(D_zv_0\mapsto(\chi_z(b))_{b\in C_+}\) is a linear
isomorphism onto functions on \(C_+\). It sends \(f_a\) to
the indicator of \(a\). Every coefficient in the formula is a
root of unity divided by \(D\). Its coefficient sum in absolute
value is one after any automorphism of \(\C\).
\end{proposition}
\begin{proof}
The lattice \(L_0\) is self-dual by
Lemma~\ref{lem:hom-oscillator-dual-lattices}. The trace character
on \(\jmath H(L_0,L_0)\) is one even after multiplying either
argument by an integral element. In particular \(B_f=1\) on
\(L_0\times L_0\), since \(\omega\) is integral. This gives
\(D_l v_0=v_0\) for \(l\in L_0\).
Hermitian pairings between two vectors supported on \(\bar M\)
vanish: conjugating the first argument moves its support into \(M\).
It follows from the product law of
Lemma~\ref{lem:theta-positive-translation} that \(D_zv_0\) is
unchanged by adding \(L_-\) to \(z\). Distinct classes in \(C_-\)
give disjoint support cosets of \(L_0\), so these vectors are
linearly independent.

The pairing defining \(\chi\) is trivial on adding
\(L_-\) to \(z\) or \(d_P L_+\) to \(a\), by the same duality.
If \(\chi_z\) is identically one, the character tested at \(z\)
is trivial on all of \(L_0\), since only its \(M\)-components pair
with \(z\). Hence \(z\in L_-\). Locally \(d_P\) is a uniformizer
times a unit at each member of \(M\), and invertible elsewhere,
by Lemma~\ref{lem:split-level} and
Definition~\ref{def:row-string-lowering-map}. Both quotients are
products of \(nn'\)-dimensional residue vector spaces of the
same orders. Proposition~\ref{prop:marked-stage-cut-size} counts
the marked places, giving the formula for \(D\).
The resulting \(D\) distinct characters are a basis of functions
on \(C_+\): two distinct ones are orthogonal because translation
by an element where their ratio differs from one multiplies their
inner product by that ratio. Thus the character matrix divided
by \(\sqrt D\) is unitary. Its inverse gives the displayed formula
and its image is the indicated point mass. Character values here
are roots of unity since \(C_+\) is finite. Every automorphism
preserves their orders and fixes the integer \(D\), proving the
size assertion.
\end{proof}

\begin{corollary}\label{cor:finite-inclusion-periods}
The actual residue indicator \(f_{a_{\rm inc}}\) is fixed by \(D_u\)
for all \(u\in d_P L_0\). In its expansion in
Proposition~\ref{prop:finite-oscillator-residue-name}, multiplication
by \((g_0g_1)^{f nn'}\) makes every coefficient integral at all
places away from \(Q\). For integral stage isometries \(k,k'\),
the change of variables \(v\mapsto k'^{-1}v k\) on functions sends
\(f_a\) to \(f_{k'a k^{-1}}\). In particular the inclusion indicator
is fixed when the two reductions commute with the inclusion on \(M\).
\end{corollary}
\begin{proof}
Since \(u\in L_0\), it fixes \(v_0\). The commutator with each
\(D_z\) is one by the duality calculation of
Lemma~\ref{lem:hom-oscillator-dual-lattices}, since
\(z\in\bar d_P^{-1}L_0\). This proves the period assertion termwise.
After the indicated multiplication the denominator is \(Q^{nn'}\),
and roots of unity and their inverses satisfy monic equations
over \(\Z\). Integral isometries preserve \(L_0,H\) and \(B_f\).
They therefore preserve its indicator and carry each translated
vector \(D_z v_0\) to \(D_{k'z k^{-1}}v_0\). They likewise preserve
the pairing of the two residue modules. Reindex the finite sum of
Proposition~\ref{prop:finite-oscillator-residue-name} to get the
last assertion.
\end{proof}

\section{Transport between the finite oscillator lattices}
\begin{theorem}\label{thm:finite-oscillator-lattice-transport}
Use \(\mathcal H_f,L_0,D_z,v_0\) of
Definition~\ref{def:split-inclusion-finite-oscillator}, with the
inner product of squared norm \(\int |f|^2\) for additive Haar
measure giving \(L_0\) mass one. For any finite adelic isometry
\(G\) of \(H\) in Definition~\ref{def:hom-oscillator-pairing},
write \(L_1=GL_0\) and \(h=[L_1:L_0\cap L_1]\). Then
\[
            v_1=h^{-1/2}\sum_{r\in L_1/(L_0\cap L_1)}D_r v_0
\]
is a unit vector. There is a unique unitary linear bijection
\(W_G\) of \(\mathcal H_f\) sending \(v_0\) to \(v_1\) and
satisfying \(W_G D_z=D_{Gz}W_G\) for all \(z\). Every unitary
bijection with this intertwining property is a scalar of modulus
one times \(W_G\). If \(K_0\) is an isometry preserving \(L_0\),
\[
 W_{K_0}f(v)=f(K_0^{-1}v),\qquad
 W_{GK_0}=W_GW_{K_0},\qquad W_{K_0G}=W_{K_0}W_G.
\]
\end{theorem}
\begin{proof}
Each lattice is compact open. Its intersection with the other has
finite index by an open coset cover. Both are self-dual and \(B_f\)
is one on pairs in each, by Lemma~\ref{lem:hom-oscillator-dual-lattices},
Proposition~\ref{prop:finite-oscillator-residue-name} and invariance
under \(G\). Hence \(D_rv_0\) in the displayed sum is independent
of the chosen lift of its coset. Translation preserves additive
measure and all factors of \(B_f\) have modulus one, so the \(D_r\)
are unitary. Distinct summands have disjoint supports and norm
one. Left translation by any \(D_l,\ l\in L_1\), permutes them
with factor one. Thus \(v_1\) has norm one and is fixed by \(L_1\).

For \(a=0,1\), the translates \(D_y v_a\) at distinct classes
modulo \(L_a\) are orthogonal. Indeed the commuting unitary operators
\(D_l,\ l\in L_a\), act on the translate by the character
\(l\mapsto\psi_f(\jmath H(y,l))\) by
Lemma~\ref{lem:theta-positive-translation}. Distinct classes give
distinct characters by self-duality. Applying an operator whose
eigenvalues differ forces the inner product to vanish. Within one
class there is only the relation
\(D_{y+l}v_a=B_f(y,l)^{-1}D_yv_a\).

The span of the \(D_yv_1\) contains \(v_0\). Average \(D_l v_1\)
over \(l\in L_0\). In the displayed expansion the average of the
term \(D_rv_0\) vanishes for \(r\notin L_0\), by translation of
a nontrivial character. The identity coset contributes
\(h^{-1/2}v_0\). This is a finite average: the finitely many
characters appearing have open kernels by the local constancy of
\(\psi_f\) in Proposition~\ref{prop:additive}, and the common
kernel has finite index. The span is invariant under all \(D_y\)
by the product rule. Hence the translates of \(v_1\) span
\(\mathcal H_f\), just as those of \(v_0\) do by definition.
The map \(D_yv_0\mapsto D_{Gy}v_1\) now respects all relations
and sends an orthonormal basis to one, giving the required
bijection. An operator commuting with all translations sends
\(v_0\) into the \(L_0\)-fixed line by the character and spanning
calculation for \(a=0\). It must therefore be scalar on the
entire span. This gives uniqueness with the indicated normalization.

An integral isometry \(K_0\) fixes \(v_0\), and change of variables
by \(K_0^{-1}\) intertwines the translations, giving the first
formula. The lattices of \(GK_0\) and \(G\) coincide. For
\(K_0G\) the displayed average is the image under \(W_{K_0}\)
of that for \(G\), by reindexing cosets. Uniqueness gives the
two product formulas.
\end{proof}

\begin{corollary}\label{cor:moved-inclusion-finite-frequency}
Take \(G\) acting on maps by \(v\mapsto x'v x^{-1}\) in the
reference isometric coordinates of two consecutive stages. In
rational coordinates this action sends the reference module of
integral maps to \(\mathcal M={\rm Hom}_{\widehat O_E}(L,L')\).
The vector \(W_G f_{a_{\rm inc}}\) is fixed by every
\(D_u,\ u\in d_P\alpha\mathcal M\). Multiplication by any
complex scalar preserves that conclusion. If a linear map of this
finite space into the adelic theta functions of
Definition~\ref{def:adelic-oscillator-translations} intertwines
the finite translations, every nonzero \(I_T\) of the image of
this vector gives the map of Proposition~\ref{prop:dual-column-one-step}.
\end{corollary}
\begin{proof}
Conjugation on rectangular matrices preserves the hermitian trace
pairing because both group elements are isometries. For
\(u\in d_P\alpha\mathcal M\), \(G^{-1}u\) belongs to the corresponding
reference period lattice. Apply
Corollary~\ref{cor:finite-inclusion-periods} and the intertwining
identity of Theorem~\ref{thm:finite-oscillator-lattice-transport}.
The period test survives the intertwining map, so
Theorem~\ref{thm:constant-term-implies-lattice-map} applies.
\end{proof}

\section{Adelic realization of the finite period test}
\begin{theorem}\label{thm:adelic-periodized-oscillator}
For the rational space \(V\) in Definition~\ref{def:hom-oscillator-pairing},
write \(V_\infty=V(E\otimes_\Q\R)\). For a Schwartz function \(\phi\)
on this real vector space and \(f\in\mathcal H_f\) in the
reference rational coordinates, set
\[
     \Theta_{\phi,f}(y)
      =\sum_{\gamma\in V(E)} B(\gamma,y)
                    \phi(\gamma_\infty+y_\infty) f(\gamma_f+y_f).
\]
A Schwartz function has continuous derivatives of all orders,
each bounded after multiplication by any polynomial in real
coordinates. This series converges absolutely, locally uniformly
in \(y\), and defines an adelic theta function of
Definition~\ref{def:adelic-oscillator-translations}.
For each fixed \(\phi\), the map \(f\mapsto\Theta_{\phi,f}\)
intertwines all finite \(D_u\). If
\(\phi(y)=\exp(-\|y\|^2)\) in any real Euclidean coordinates,
it is injective.
\end{theorem}
\begin{proof}
Every \(f\) in Definition~\ref{def:split-inclusion-finite-oscillator}
is bounded, locally constant and compactly supported. On a compact
set of \(y\), the rational points contributing to the series have
finite components in a fixed compact set. Hence they lie in
\(D^{-1}O_E^{nn'}\) in a rational matrix basis, for some integer
\(D\ge1\). Indeed the finite adele space is the increasing union
of integer-denominator lattices, for example with factorial
denominators, and integral at all places means integral in \(O_E\)
by Proposition~\ref{prop:frame-lattice-trace}. In real coordinates
this is a lattice by Lemma~\ref{lem:number-field-lattice}. A ball
of radius \(R\ge1\) contains at most \(C(1+R)^{\dim_\R V_\infty}\)
of its points, by applying the inverse of a real basis matrix
and counting integer coordinates. On compact sets of \(y_\infty\)
the Schwartz bound gives an arbitrarily large inverse power
of \(1+\|\gamma_\infty\|\). Summing over dyadic shells proves
absolute locally uniform convergence. Each summand is continuous,
so the sum is continuous.

For \(r,\gamma\in V(E)\), \(B(\gamma,r)=1\) by
Proposition~\ref{prop:additive} and rationality of the trace
pairing. Reindex \(\gamma\) by addition of \(r\) to obtain the
factor \(B(r,y)^{-1}\) under \(y\mapsto y+r\). For a finite \(u\),
the identity \(B(y,u)B(\gamma,y+u)=B(\gamma,y)B(\gamma+y,u)\)
gives \(D_u\Theta_{\phi,f}=\Theta_{\phi,D_u f}\) term by term.

The kernel of any such intertwining map is invariant under all
\(D_u\). Any nonzero invariant subspace of \(\mathcal H_f\) is
the whole space. In fact expand a nonzero vector in finitely many
distinct \(D_zv_0\). Averaging over \(L_0\) against the inverse
of one occurring eigencharacter isolates that basis vector by the
character calculation of Theorem~\ref{thm:finite-oscillator-lattice-transport}.
The average is a finite sum there. Translating then gives all the
basis vectors. For the specified Gaussian the value
\(\Theta_{\phi,v_0}(0)\) is positive, since all terms are nonnegative
and the term at zero equals one. Therefore the kernel is zero.
\end{proof}

\begin{corollary}\label{cor:actual-inclusion-line-constant-term}
For \(x,x'\) at successive stages with nonempty target string let
\(G\) be their action on the finite Hom space in
Corollary~\ref{cor:moved-inclusion-finite-frequency}.
For every Schwartz \(\phi\) take
\(F=\Theta_{\phi,W_G f_{a_{\rm inc}}}\), or any scalar multiple.
If \(I_T(F)\ne0\), the deviations satisfy the upper bound of
Lemma~\ref{lem:initial-line-identity}. If the middle-line condition
holds at \(x,x'\), the two-sided bounds of
Theorem~\ref{thm:one-sided-stage-line-interlacing} hold.
\end{corollary}
\begin{proof}
Theorem~\ref{thm:adelic-periodized-oscillator} supplies the finite
intertwiner in Corollary~\ref{cor:moved-inclusion-finite-frequency}.
Use the resulting lattice map in Proposition~\ref{prop:dual-column-one-step}
and Lemma~\ref{lem:initial-line-identity}. The lower bound is then
exactly the deduction of Theorem~\ref{thm:one-sided-stage-line-interlacing}.
\end{proof}

\section{The split local action on the inclusion datum}
\begin{theorem}\label{thm:split-finite-schroedinger}
Fix a finite set \(M\) of finite places of \(E\) disjoint from
\(\bar M\). In the coordinates of
Definition~\ref{def:split-inclusion-finite-oscillator}, let \(X,Y\)
denote the full spaces of matrices supported on \(M,\bar M\),
respectively. Let \(\Lambda=X\cap L_0,\ \Lambda^\vee=Y\cap L_0\).
Give \(X\) its additive measure with \(\Lambda\) of mass one.
The factor of \(\mathcal H_f\) at \(M\cup\bar M\) identifies
isometrically with the locally constant compactly supported
functions on \(X\), with vacuum \({\bf1}_\Lambda\). For \(x\in X\)
and \(y\in Y\) the identification takes \(D_x,D_y\) to
\[
      (L_x b)(t)=b(t+x),\qquad
      (M_y b)(t)=\psi_f(\jmath H(y,t))b(t).
\]
If \(G\) is an \(E\)-linear isometry in the finite coordinates,
acting trivially off these pairs, and \(A=G|_X\), then
\(W_G\) of Theorem~\ref{thm:finite-oscillator-lattice-transport}
acts on this factor by
\[
      (U_A b)(t)=J_A^{-1/2}b(A^{-1}t),\qquad
      J_A=\prod_{w\in M}|\det A_w|_w .
\]
Here \(|u|_w=q_w^{-v_w(u)}\). In particular \(W_GW_{G'}=W_{GG'}\)
for isometries supported on these pairs.
\end{theorem}
\begin{proof}
Write \(\chi_y(t)=\psi_f(\jmath H(y,t))\). The vanishing of \(H\)
on \(X\times X\) and \(Y\times Y\), and the self-duality in
Lemma~\ref{lem:hom-oscillator-dual-lattices}, give
\(\Lambda^\vee=\{y:\chi_y(\Lambda)=1\}\), with the reverse
annihilator equal to \(\Lambda\). The bicharacter \(B_f\)
on separate pairs of places is one. Accordingly a translated
vacuum splits as the product of its translated vacua on
\(X\oplus Y\) and the complementary places. Products of basis
vectors in the two spaces have disjoint coset supports and norm
one, giving the asserted factor with the product inner product.

For \(x+y\) put
\(\widetilde D_{x+y}=B_f(x,y)^{-1}L_xM_y\).
One has \(L_xM_y=\chi_y(x)M_y L_x\), and
\(\chi_y(x)=B_f(x,y)/B_f(y,x)\) by
Lemma~\ref{lem:theta-positive-translation}. To multiply two such
operators reorder the middle \(M_yL_{x'}\), acquiring the factor
\(\chi_y(x')^{-1}\). After factoring out
\(\widetilde D_{x+x'+y+y'}\), the total factor is
\(B_f(x,y')B_f(y,x')=B_f(x+y,x'+y')\).
The vacuum is fixed by each \(\widetilde D_l,\ l\in\Lambda\oplus
\Lambda^\vee\), since \(B_f\) is one on \(L_0\times L_0\).
The translates with distinct \(X/\Lambda\) coordinates
have disjoint support. On one support coset distinct
\(Y/\Lambda^\vee\) coordinates have orthogonal characters, by
integrating their ratio and translating in \(\Lambda\).
Thus sending \(D_zv_0\) on this factor to
\(\widetilde D_z{\bf1}_\Lambda\) preserves the norms and the
relations within each lattice coset, as listed in
Theorem~\ref{thm:finite-oscillator-lattice-transport}.

A locally constant compactly supported function meets finitely
many cosets of \(\Lambda\). Compactness gives, on each of them,
invariance under \(\prod_{w\in M}\pi_w^{e_w}\Lambda_w\) for
some \(e_w\ge0\). Indeed finitely many open constancy balls
cover the coset, and a common smaller subgroup is contained in
the translation subgroup of each ball. The characters from the
opposite inverse powers of the uniformizers form a basis of
functions on this finite quotient. The annihilator argument and
equal-order count of Proposition~\ref{prop:finite-oscillator-residue-name}
apply with \(e_w\) successive residue quotients at \(w\).
This proves surjectivity.

An invertible matrix over \(E_w\) changes additive measure by
\(|\det A_w|_w\). Gaussian elimination uses coordinate
permutations, row additions and scalings. Permutations preserve
the product measure, additions preserve it by translation in the
changed coordinate, and multiplication by a uniformizer gives
the sublattice of index \(q_w\). A unit permutes every finite
residue quotient. These checks on generators give the formula.
Consequently \(U_A\) is unitary. Direct substitution gives
\(U_A L_x U_A^{-1}=L_{Ax}\) and
\(U_A M_y U_A^{-1}=M_{Gy}\), using invariance of \(H\).
The inner product of \(U_A{\bf1}_\Lambda\) with
\({\bf1}_\Lambda\) is strictly positive, since \(A\Lambda\cap
\Lambda\) is open. The normalized vacuum for \(W_G\) also has
positive inner product with the reference vacuum: only its
identity coset contributes to the average in
Theorem~\ref{thm:finite-oscillator-lattice-transport}.
The two intertwiners differ by a scalar of absolute value one
by that theorem. Positivity forces it to be one.
The product law follows from \(J_{AA'}=J_AJ_{A'}\).
\end{proof}

\begin{corollary}\label{cor:split-moved-indicator}
For \(M\) as in Definition~\ref{def:split-inclusion-finite-oscillator},
the vector \(f_a\) on these pairs is \({\bf1}_{a+d_P\Lambda}\)
in the preceding realization. If the two stage elements act at
the \(w\)-component by invertible matrices \(g_w,g'_w\) of
orders \(n,n'\), set \(A_w t=g'_w t g_w^{-1}\). Their moved
indicator on \(X\) is
\[
   J_A^{-1/2}{\bf1}_{A(a+d_P\Lambda)},\qquad
   J_A=\prod_{w\in M}|\det g'_w|_w^{\,n}|\det g_w|_w^{-n'} .
\]
\end{corollary}
\begin{proof}
On \(\Lambda\) the Fourier sum of
Proposition~\ref{prop:finite-oscillator-residue-name} equals
one on the indicated residue coset and zero elsewhere by
character orthogonality. Every term vanishes outside \(\Lambda\).
Apply Theorem~\ref{thm:split-finite-schroedinger}.
Left multiplication acts on the \(n\) columns and right
multiplication on the \(n'\) rows, giving the determinant factors.
\end{proof}

\section{Coefficients of a moved residue coset}
\begin{proposition}\label{prop:moved-residue-coset-fourier}
Use \(M,\Lambda,A,J_A,a\) of Corollary~\ref{cor:split-moved-indicator}.
Choose a representative of \(a\) in \(\Lambda\). Write
\[
    H_1=A(d_P\Lambda),\qquad H_2=H_1\cap\Lambda,\qquad
    u=[H_1:H_2],\qquad v=[\Lambda:H_2].
\]
For any subgroup \(H\subset X\) write
\(H^\vee=\{y\in Y:\chi_y(H)=1\}\), using the character of
Theorem~\ref{thm:split-finite-schroedinger}.
The moved inclusion vector for \(a\) has an expansion with \(uv\)
distinct translated vacuum cosets. Every coefficient is a root
of unity times \(J_A^{-1/2}/v\). Its sum of absolute coefficient
values under each field automorphism of \(\C\) is \(uJ_A^{-1/2}\).
Moreover
\[
                 \frac{u}{v}=\frac{J_A}{[\Lambda:d_P\Lambda]} .
\]
\end{proposition}
\begin{proof}
The intersection of two compact open lattices has finite index in
each by compactness. The cosets of \(\Lambda\) occupied by
\(Aa+H_1\) have representatives \(b\in Aa+H_1\), one for each
of the \(u\) cosets modulo \(H_2\) there. Inside \(b+\Lambda\)
the desired indicator is
\[
  {\bf1}_{b+H_2}(t)=\frac{1}{v}
    \sum_{y\in H_2^\vee/\Lambda^\vee}
        \chi_y(-b)\,(M_y L_{-b}{\bf1}_{\Lambda})(t).
\]
Indeed \(H_2\) contains a product of uniformizer powers
\(\Lambda'\subset\Lambda\). The proof of
Theorem~\ref{thm:split-finite-schroedinger} supplies a character
basis on \(\Lambda/\Lambda'\). Averaging each character under
the translations by \(H_2/\Lambda'\) kills it unless it is
trivial on that subgroup. The surviving characters are therefore
a basis of the functions on \(\Lambda/H_2\) and number \(v\).
Averaging them gives the indicator of its identity element:
the columns of their square character matrix are orthogonal
by the finite Fourier calculation in
Proposition~\ref{prop:finite-oscillator-residue-name}.
Multiplication by \(\chi_y(-b)\) centers this indicator at \(b\).

Each \(M_y L_{-b}\) is \(D_{y-b}\) times a value of \(B_f\) by
Lemma~\ref{lem:theta-positive-translation}. Such values and the
values \(\chi_y(-b)\) are roots of unity by the finite product of
rational fractional-part characters in Proposition~\ref{prop:additive}.
The first coordinate distinguishes the \(u\) supports and the
second distinguishes the \(v\) modulations on each, giving the
claimed distinct cosets. Multiply by the amplitude \(J_A^{-1/2}\)
from Corollary~\ref{cor:split-moved-indicator}. Its square is a
positive rational number, so any automorphism sends it to a root
with the same absolute value. Coefficient magnitudes now give
the stated sum. Finally \(H_2\) has measure \(1/v\), and \(H_1\)
has measure both \(u/v\) and \(J_A/[\Lambda:d_P\Lambda]\),
which proves the identity.
\end{proof}

\begin{corollary}\label{cor:twist-only-moved-denominators}
Let \(Q=\prod_{w\in P}q_w\) and suppose the finite stage elements
are integral away from \(P\cup\bar P\), in their reference
isometric coordinates. For any \(a\), multiplying the coefficients
of the moved \(f_a\) by \((g_0g_1)^{f nn'}\) makes them integral at
every finite place with residue characteristic prime to \(Q\).
The translating vectors can be supported on \(P\cup\bar P\)
and the conjugates of the marked places. At those conjugates they
lie in \(\bar d_P^{-1}L_0\). For \(Q=1\) the sum of the absolute
coefficient values under every complex field automorphism is one.
\end{corollary}
\begin{proof}
On each integral complementary pair or nonsplit place the lattice
is fixed, so the action preserves the vacuum by
Theorem~\ref{thm:finite-oscillator-lattice-transport}. That theorem
also allows the integral part and the part supported on
\(P\cup\bar P\) to be composed in either order when their factors
are disjoint. At a marked \(w\), \(A_w\Lambda_w=\Lambda_w\) and
\(H_{2,w}=d_P\Lambda_w\). In the preceding expansion the translation
of the support there belongs to \(\Lambda_w\) and can be dropped
at the cost of a root of unity, by the lattice-coset relation in
Theorem~\ref{thm:finite-oscillator-lattice-transport}.
Only the modulation from the opposite
\(\bar d_P^{-1}\Lambda^\vee_{\bar w}\) remains.
The marked factors of \(v\) multiply to \((g_0g_1)^{f nn'}\)
by Proposition~\ref{prop:finite-oscillator-residue-name}.
Every other factor of \(v\), and the numerator and denominator
of \(J_A\) in lowest terms, has prime divisors only among those
of \(Q\): the indices over \(w\) are orders of quotients of
local lattices killed by a uniformizer power. Their sizes are
powers of the residue characteristic by successive residue
quotients. Thus \(J_A^{-1/2}\) has valuation zero away from
\(Q\), by squaring. Roots of unity likewise have valuation zero
by raising to their orders. This proves the denominator assertion
using Proposition~\ref{prop:moved-residue-coset-fourier}.
If \(Q=1\), all changes are integral, so \(u=1,\ J_A=1\)
in that proposition.
\end{proof}

\section{Reduction of the finite oscillator to degree zero}
\begin{lemma}\label{lem:zero-grade-lattice-dual}
In the finite adelic space of Definition~\ref{def:hom-oscillator-pairing}
let \(U\) be the strictly positive-grade space, \(W\) the degree-zero
space, and \(N_+=U\oplus W\) the nonnegative-grade space. Let
\(\pi:N_+\to W\) be the projection. If \(L\) is a moved reference
lattice \(\alpha\mathcal M\) as in
Lemma~\ref{lem:hom-oscillator-dual-lattices}, then
\(\Lambda=\pi(L\cap N_+)\) is a compact open self-dual lattice
in \(W\) for the restricted trace character. The bicharacter
\(B_f\) is one on \(\Lambda\times\Lambda\).
\end{lemma}
\begin{proof}
The same grade pairing as in Lemma~\ref{lem:hom-oscillator-dual-lattices}
gives \(U^\perp=N_+\) and \(N_+^\perp=U\). Here the orthogonals
are taken for the local rational trace pairings and assembled
over the finite places. In particular, the hermitian form on
\(N_+\) is the pullback of the nonsingular form on \(W\).
The projection is open in the direct-sum topology, and \(L\cap N_+\)
is compact open in \(N_+\). Its image is therefore compact open.
The calculation of Proposition~\ref{prop:positive-line-oscillator-annihilator}
also gives
\[
             (L\cap U)^\circ=L+N_+,\qquad
             (L\cap N_+)^\circ=L+U.
\]
Indeed at each prime the intersection with either linear subspace
is saturated in \(L_p\). Extend the tested integral functional
by zero on an integral basis complement and use self-duality to
represent the extension in \(L_p\), taking the tested vector
itself at almost all primes as in that proposition.
For \(w\in W\) annihilating \(\Lambda\), the second equality
gives \(w=l+u,\ l\in L,\ u\in U\); then
\(l\in N_+\) and \(w=\pi l\). Conversely \(L\) annihilates
itself and the pairing descends on \(N_+\). This proves
self-duality. The same descent applies to \(B_f\), which is
one on \(L\times L\) by Proposition~\ref{prop:finite-oscillator-residue-name}
and transport of the reference lattice.
\end{proof}

\begin{theorem}\label{thm:finite-grade-oscillator-quotient}
For \(L,U,W,N_+,\pi,\Lambda\) as in
Lemma~\ref{lem:zero-grade-lattice-dual}, form \(\mathcal H(V,L)\)
as the span of the \(D_y{\bf1}_L\) using \(B_f\).
Similarly form \(\mathcal H(W,\Lambda)\) on \(W\) with the
restricted bicharacter and denote its translations by \(D^0_w\).
There is a unique surjection
\(R:\mathcal H(V,L)\to\mathcal H(W,\Lambda)\) with
\(R{\bf1}_L={\bf1}_\Lambda,\ R D_w=D^0_w R\) for \(w\in W\),
and \(R D_u=R\) for \(u\in U\). Its kernel is
\[
    \operatorname{span}_{\C}\{D_u f-f:u\in U,\ f\in\mathcal H(V,L)\}.
\]
If \(y\notin L+N_+\), then \(R D_y{\bf1}_L=0\).
Otherwise, for any \(y=l+n\) with \(l\in L,\ n\in N_+\),
\[
  R D_y{\bf1}_L=B_f(n,l)^{-1}D^0_{\pi n}{\bf1}_\Lambda .
\]
\end{theorem}
\begin{proof}
Write \(e_y=D_y{\bf1}_L\). As in
Theorem~\ref{thm:finite-oscillator-lattice-transport}, the lines
\(\C e_y\) have disjoint supports at distinct classes in \(V_f/L\),
where \(V_f\) denotes the ambient finite adelic space. For
\(l\in L\), \(e_{y+l}=B_f(y,l)^{-1}e_y\).
First take the quotient by the asserted kernel span and denote
classes by brackets. The operators \(D_u\) for \(u\in U\) represent
addition exactly, since \(B_f\) is one on \(U\times U\).
They permute the basis lines along their translation orbits.
For \(u\in L\cap U\) the action on \(e_y\) is multiplication by
\(\psi_f(\jmath H(y,u))\) by Lemma~\ref{lem:theta-positive-translation}.
Thus the class vanishes unless \(y\in L+N_+\), by
Lemma~\ref{lem:zero-grade-lattice-dual}. On a surviving orbit
take \(y=l+n\) and write \(n=w+u\) with \(w=\pi n\).
The lattice-coset relation and \(B_f(u,w)=1\) give
\([e_y]=B_f(n,l)^{-1}[e_w]\).

Every such orbit contributes exactly one dimension. On the
orbit of \(w\in W\), the vectors \(D_u e_w\) in the same
basis line coincide, since their indices differ by an element
of \(U\cap L\), which fixes \(e_w\). Assigning value one to
each of these vectors gives a well-defined invariant functional
on the orbit span, and hence \([e_w]\ne0\). Distinct orbits
give independent quotient spaces because every generating
difference belongs to the sum of the separate orbit relations,
on expanding \(f\) in basis lines.

The orbits through \(w,w'\in W\) coincide exactly when
\(w'-w\in(L+U)\cap W=\Lambda\).
All \(D_w,\ w\in W\), commute with \(D_u\) and thus descend.
If \(\lambda\in\Lambda\), lift it to \(l\in L\cap N_+\).
Writing \(l=\lambda+u\) gives \([e_\lambda]=[e_0]\).
Thus the relations on the surviving vectors are exactly
\[
      [e_{w+\lambda}]=B_f(w,\lambda)^{-1}[e_w]\qquad
                         (w\in W,\ \lambda\in\Lambda).
\]
These are the relations among \(D^0_w{\bf1}_\Lambda\):
the vacuum is fixed by Lemma~\ref{lem:zero-grade-lattice-dual},
and distinct support cosets are disjoint. Sending the bracketed
vectors to these translates identifies both spaces, respecting
the product law on \(W\). This constructs the claimed surjection
with its exact kernel. Every map satisfying the three conditions
factors through the quotient and has the prescribed value on
each of its spanning translates, giving uniqueness.
\end{proof}

\section{An integral lift of the residual space}
\begin{proposition}\label{prop:integral-adjoint-reduction}
Use the notation of Theorem~\ref{thm:finite-grade-oscillator-quotient}.
Let \(W_s\) be another finite adelic hermitian space with the
same scalar \(\jmath\), bicharacter defined as in
Definition~\ref{def:adelic-oscillator-translations}, and compact
open self-dual lattice \(L_s\) on whose square the bicharacter
is one. Suppose \(h:W_s\to W\) is an isometry onto \(W\), and
\(b:W_s\to U\) is linear over the finite adeles of \(E\). Put
\(B=h+b:W_s\to V_f\) and assume \(B(L_s)\subset L\).
Let \(B^\dagger:V_f\to W_s\) denote its hermitian adjoint,
so \(H_s(t,B^\dagger y)=H(Bt,y)\). Then
\[
       h L_s=\Lambda,\qquad B^\dagger L\subset L_s,\qquad
       B^\dagger n=h^{-1}\pi n\quad(n\in N_+).
\]
If \(Z:\mathcal H(V,L)\to\mathcal H(W_s,L_s)\) is linear and
\(ZD_u=Z,\ ZD_{hw}=D^s_w Z\) for \(u\in U,\ w\in W_s\),
there exists a common scalar \(C\) such that for every \(y\in V_f\)
\[
 ZD_y{\bf1}_L=C\zeta_y D^s_{B^\dagger y}{\bf1}_{L_s}.
\]
One can take \(\zeta_y=0\) for \(y\notin L+N_+\) and
\(\zeta_y\) a root of unity for \(y\in L+N_+\).
\end{proposition}
\begin{proof}
The hermitian adjoint exists since the form on \(W_s\) is
nonsingular; in coordinate bases it is a product of the inverse
Gram matrix, the conjugate transpose of \(B\), and the ambient
Gram matrix. Projecting \(B L_s\) gives \(hL_s\subset\Lambda\).
Both lattices are self-dual by Lemma~\ref{lem:zero-grade-lattice-dual}
and invariance of the pairing under \(h\). Taking annihilators
reverses the inclusion and proves equality.
For \(l\in L\), the functional given by \(B^\dagger l\) has
trivial trace character on \(L_s\), since \(B L_s\subset L\).
This proves the second inclusion. For \(n\in N_+\), the pairing
with \(b\) vanishes, and \(H(ht,n)=H_s(t,h^{-1}\pi n)\),
proving the last equality.

For \(t\in L_s\), \(D_{Bt}{\bf1}_L={\bf1}_L\); since
\(Bt=ht+bt\) with mutually orthogonal components, applying \(Z\)
gives \(D^s_t Z{\bf1}_L=Z{\bf1}_L\).
The fixed space is the vacuum line by the character argument in
Theorem~\ref{thm:finite-oscillator-lattice-transport}: the
eigencharacters on the distinct translated basis lines are
distinct by self-duality. Write \(Z{\bf1}_L=C{\bf1}_{L_s}\).
By its \(U\)-invariance \(Z\) kills the kernel in
Theorem~\ref{thm:finite-grade-oscillator-quotient}.
For \(y=l+n\) as there it sends \(D_y{\bf1}_L\) to
\(C B_f(n,l)^{-1}D^s_{h^{-1}\pi n}{\bf1}_{L_s}\).
Adding \(B^\dagger l\in L_s\) to the index of this last translation
changes the vector by a value of the finite bicharacter.
All such factors are roots of unity by Proposition~\ref{prop:additive},
and \(B^\dagger y=B^\dagger l+h^{-1}\pi n\).
This proves the formula, including when \(C=0\).
\end{proof}

\begin{corollary}\label{cor:shifted-integral-adjoint-reduction}
Under Proposition~\ref{prop:integral-adjoint-reduction}, let
\(\tau\in V_f\) satisfy \(B^\dagger\tau=0\).
There are coefficients \(\epsilon_y\), each zero or a root of
unity, such that
\[
       ZD_\tau D_y{\bf1}_L
                = C\epsilon_y D^s_{B^\dagger y}{\bf1}_{L_s}
                    \qquad(y\in V_f).
\]
If \(y\mapsto c_y\) has finite support, this substitution in the sum
over \(c_y D_y{\bf1}_L\), with the common \(C\) factored outside,
preserves every common algebraic denominator for the \(c_y\)
and every upper bound for their total absolute value under all
field automorphisms of \(\C\).
\end{corollary}
\begin{proof}
Multiply the product formula in Lemma~\ref{lem:theta-positive-translation}
by the formula of Proposition~\ref{prop:integral-adjoint-reduction}
at \(\tau+y\). The new character factor is \(B_f(\tau,y)\);
\(B^\dagger(\tau+y)=B^\dagger y\). Roots of unity are integral
units in each finite valuation ring and have absolute value one
after every field automorphism. Adding any coefficients at
coincident images preserves the denominator in each valuation
ring, and the triangle inequality preserves the bound.
\end{proof}

\section{Combining the two finite oscillator spaces}
\begin{theorem}\label{thm:paired-finite-adjoint-model}
For consecutive stage spaces \(Y,Y',Y''\) of
Definition~\ref{def:actual-row-lift-spaces} put
\[
 V_1={\rm Hom}_E(Y,Y'),\quad V_2={\rm Hom}_E(Y',Y''),\quad
 V_*={\rm Hom}_E(Y',Y^-\oplus Y'').
\]
The superscript minus negates the hermitian form. Give these spaces
their trace hermitian forms of Definition~\ref{def:hom-oscillator-pairing}
and use the associated finite bicharacters \(B_1,B_2,B_*\).
Write \(L_Y,L_{Y'},L_{Y''}\) for the three stage lattices in the
isometric frames of Theorem~\ref{thm:actual-row-lift-frames}, and put
\(L_i=\alpha\,{\rm Hom}(L_{\rm source},L_{\rm target})\) in \(V_i\),
\(i=1,2\), with integral finite-adelic Hom. Define
\[
 A(a,b)=(a^\dagger,b),\qquad
 L_*=A(L_1\oplus L_2)
  =\bar\alpha\,{\rm Hom}(L_{Y'},L_Y)
               \ \oplus\ \alpha\,{\rm Hom}(L_{Y'},L_{Y''}).
\]
Here \(a^\dagger\) is the adjoint with respect to the original
forms on \(Y,Y'\). There is an isometric complex linear
identification of the algebraic tensor product
\[
    \mathcal H(V_1,L_1)\otimes_\C\mathcal H(V_2,L_2)
                           \ \longrightarrow\ \mathcal H(V_*,L_*).
\]
It sends the tensor of the vacua to \({\bf1}_{L_*}\) and intertwines
\(D_a^1\otimes D_b^2\) with \(\lambda(a)D^*_{A(a,b)}\), where
\[
       \lambda(a)=\psi_f\big(-\bar\omega H_1(a,a)\big).
\]
The factor \(\lambda(a)\) is a root of unity. The identification
commutes with the normalized lattice intertwiners for every triple
of finite stage isometries acting on the three Hom spaces.
\end{theorem}
\begin{proof}
In isometric coordinates conjugate transpose carries all integral
rectangular matrices bijectively onto the integral transposed
rectangular matrices. Moving the scalar \(\alpha\) through it
gives \(\bar\alpha\). Thus \(A\) is an additive bijection with
the asserted lattice image. For \(z=(a,b),\ z'=(a',b')\),
matrix trace and cyclic permutation give
\[
 H_*(Az',Az)=-\overline{H_1(a',a)}+H_2(b',b).
\]
In fact the trace of \(a'a^\dagger\) equals the conjugate of
the trace of \(a'^\dagger a\). The minus sign is exactly the
one on the first output space. Conjugation preserves the trace
in \(\psi_f\) and negates \(\jmath\). Hence the alternating trace
pairing is preserved. In particular \(L_*\) is self-dual by
Lemma~\ref{lem:hom-oscillator-dual-lattices}.
For the bicharacters, with \(s=H_1(a',a)\), the quotient is
\[
 \frac{B_*(Az,Az')}{B_1(a,a')B_2(b,b')}
     =\psi_f\big(\jmath(\omega-\bar\omega)s\big)
     =\psi_f(-s)
     =\frac{\lambda(a+a')}{\lambda(a)\lambda(a')}.
\]
The middle equality uses \(\jmath(\omega-\bar\omega)=-1\) from
Lemma~\ref{lem:actual-lift-scale}. For the last, the two mixed
terms are \(s,\bar s\); under the trace their coefficients
\(\bar\omega,\omega\) add to one. All finite character values
are roots of unity by Proposition~\ref{prop:additive}.
For \(a,a'\in L_1\) the pairing belongs to
\(\alpha\bar\alpha\widehat O_E\). This is contained in the trace
dual since \(\jmath\alpha\bar\alpha\widehat O_E\) is the trace
dual and \(\jmath^{-1}\) is integral, by
Lemma~\ref{lem:hom-oscillator-dual-lattices} and
Proposition~\ref{prop:different}. It follows that \(\lambda\) is
one on \(L_1\) and the preceding quotient is one on pairs of
lattice vectors. Thus \(B_*\) is one on \(L_*\times L_*\).

The operators \(\lambda(a)D^*_{Az}\) now obey the same product
law as the tensor operators, and they fix the new vacuum for
\(z\in L_1\oplus L_2\). Defining the map on the translated vacuum
tensors by this rule respects the lattice-coset relations of
Theorem~\ref{thm:finite-oscillator-lattice-transport}.
It takes the orthonormal basis indexed by pairs of lattice cosets
to an orthonormal basis since \(A\) is bijective. This proves
the identification. For a triple of isometries the first transformed
adjoint is \(g a^\dagger g'^{-1}\) if the original map becomes
\(g'a g^{-1}\). Also \(\lambda(g'a g^{-1})=\lambda(a)\).
Consequently the two transported operators intertwine the same
translations and differ by a scalar of absolute value one, by
the vacuum-line argument of
Theorem~\ref{thm:finite-oscillator-lattice-transport}.
Each has positive overlap of the moved vacuum with the original
vacuum by the lattice average of that theorem, taking the product
of the two overlaps for the tensor operator. Hence the scalar is one.
\end{proof}

\section{The rational residual space for a pair of lifts}
\begin{proposition}\label{prop:actual-pair-grade-zero}
Let \(j\) be even with \(0\le j\le m-2\) and set
\(Y=Y_j,\ Y'=Y_{j+1},\ Y''=Y_{j+2}\),
\(s=m-1-j\). In \(V_*={\rm Hom}(Y',Y^-\oplus Y'')\)
of Theorem~\ref{thm:paired-finite-adjoint-model}, grade a matrix
entry by the target grade minus the source grade as in
Definition~\ref{def:row-string-lowering-map}. Put \(X'=X_{j+1}\).
Let \(Q_0\) be the grade-zero space of \(Y^-\oplus Y''\).
The degree-zero subspace of \(V_*\) is \({\rm Hom}(X',Q_0)\),
embedded by zero on the string of \(Y'\).

If \(s\ge3\), write \(z,z''\) for the central string vectors of
\(Y,Y''\). There is an orthogonal decomposition
\[
 Q_0=E\eta\oplus E\xi\ \perp\ (X_j^-\oplus X_{j+2}),
 \qquad
 \xi=z+z'',\qquad \eta=\frac{z-z''}{2N},
\]
where \(\eta,\xi\) have zero norms and pairing one. If \(s=1\),
write \(z\) for the unique string vector of \(Y\). Then
\[
              Q_0=X_j^-\ \perp\ (Y_m\oplus E(z/\pi_N)),
\]
and the last displayed vector has norm one. The action of
\(U(Y_m)\) on the enlarged positive space fixes that vector.
The determinant of the enlarged action equals the original
determinant. Here \(U(S)\) is the group preserving the hermitian
form on the indicated space \(S\).
\end{proposition}
\begin{proof}
The odd strings in the outside spaces have even grades, while
the even string of \(Y'\) has only odd grades. The definite
summands have grade zero. Consequently equality of source and
target grades forces both to be zero. This also applies when
a string is empty, by Definition~\ref{def:actual-row-lift-spaces}.
For \(s\ge3\), negating the first outside form gives norms \(N,-N\)
on \(z,z''\); the vectors are perpendicular. The displayed
linear combinations give the indicated pairing and span those
two lines since \(N\ne0\). Their orthogonal complement is the
sum of the definite blocks. For \(s=1\) the other outside
space is the final positive space \(Y_m\); the norm of
\(z/\pi_N\) in \(Y^-\) is \(N/(\pi_N\bar\pi_N)=1\) by
Proposition~\ref{prop:level-prime}. The extension of the last
action is block diagonal with one added identity entry.
\end{proof}

\begin{proposition}\label{prop:paired-rational-frequency}
For the stages of Proposition~\ref{prop:actual-pair-grade-zero}
let \(T:Y\to Y'\) and \(T':Y'\to Y''\) be the maps of
Definition~\ref{def:row-string-lowering-map} when their target
string is nonempty. Set the map to zero if either string is
empty. The combined map \(\tau=(T^\dagger,T')\in V_*(E)\)
annihilates \(X'\). If \(s=2d+1\ge3\), its projection to the
grade-zero target \(Q_0\) vanishes on every basis vector of \(Y'\)
except \(e'_{d-1}\), where its value is \(c\xi\).
If \(s=1\), \(\tau=0\).
\end{proposition}
\begin{proof}
The forward maps kill the corresponding definite summands and
have image in the string. The adjoint \(T^\dagger\) therefore
kills \(X'\). Each nonzero entry of \(T,T'\) lowers grade by one.
The adjoint \(T^\dagger\) also lowers grade by one: pairing
a vector of grade \(h\) with \(Ta\) can be nonzero only if
\(a\) has grade \(1-h\), and this tests precisely the component
of \(T^\dagger\) in grade \(h-1\). For \(s=2d+1\), the only
basis vector of grade one in \(Y'\) is \(e'_{d-1}\).
The coefficient of \(T'\) there is \(c\), so
\(T'e'_{d-1}=cz''\). In \(Y\) the coefficient at \(z=e_d\)
is \(-Nc\), so that
\[
  H_Y(z,T^\dagger e'_{d-1})
      = H_{Y'}(Tz,e'_{d-1})=-N\bar c .
\]
Here \(H_Y,H_{Y'}\) denote the stage forms, conjugate-linear
in the first argument. The adjoint has no component in \(X_j\),
since \(T\) kills \(X_j\). Its value is thus \(\bar c z\).
The factors \(\jmath,\varpi\) in \(c=-(\jmath g\varpi)^{-1}\)
both change sign under conjugation by
Proposition~\ref{prop:different} and Lemma~\ref{lem:actual-lift-scale}.
This gives \(\bar c=c\), proving the formula. In the final pair
both maps have empty source or target strings.
\end{proof}

\begin{corollary}\label{cor:residual-pair-positive-quotient}
For \(s\ge3\) enlarge the strictly positive-degree subspace
\(U_>\) of \(V_*\) by \({\rm Hom}(X',E\xi)\), and for \(s=1\)
use \(U_>\) alone. Call the resulting rational space \(U_*\).
It is totally isotropic. Its hermitian orthogonal is
\(U_*\oplus W_*\), with
\[
 W_*=
 \begin{cases}
 {\rm Hom}(X',X_j^-\oplus X_{j+2}),&s\ge3,\\
 {\rm Hom}(X',X_j^-\oplus(Y_m\oplus E(z/\pi_N))),&s=1.
 \end{cases}
\]
The finite oscillator reduction of
Theorem~\ref{thm:finite-grade-oscillator-quotient} applies with
\(U_*,W_*,L_*\) of the present pair.
\end{corollary}
\begin{proof}
In the matrix trace pairing the grades occur only in opposite
pairs, by the argument of Lemma~\ref{lem:hom-oscillator-dual-lattices}.
On degree zero the hyperbolic pair in
Proposition~\ref{prop:actual-pair-grade-zero} gives a perfect
pairing between the maps into \(E\eta\) and \(E\xi\), by taking
matrix units on the orthonormal basis of \(X'\). The asserted
orthogonal and its direct decomposition follow. The proofs of
Lemma~\ref{lem:zero-grade-lattice-dual} and
Theorem~\ref{thm:finite-grade-oscillator-quotient} use only
\(U^\perp=U\oplus W\), the vanishing of the hermitian pairing
of \(U\) with that sum, self-duality of the lattice, and
triviality of its bicharacter on the lattice square.
Theorem~\ref{thm:paired-finite-adjoint-model} supplies these
lattice properties for \(L_*\). Thus their proofs apply to
the indicated subspaces, including the case of an empty string.
\end{proof}

\section{The integral image in the residual space}
\begin{definition}\label{def:paired-neighbour-lattice}
Retain \(j,X',s,W_*\) of
Corollary~\ref{cor:residual-pair-positive-quotient}, write
\(\pi=\pi_N\) and let \(O=\widehat O_E\). In the residual target
write \(e\) for the last basis vector of \(X_j^-\).
For \(s\ge3\) let \(e'\) be the last basis vector of \(X_{j+2}\).
For \(s=1\) take \(e'=z/\pi\) in the enlarged positive space
\(Y_m\oplus E(z/\pi)\). Normalize the coordinates of a residual
map \(W_*\) by writing its first block as \(\bar\alpha x_-\)
and its second as \(\alpha x_+\), using the last indicated vector
as a basis element in the final pair.
For one column in these coordinates define \(\mathcal B_\pi\)
by integrality of every entry outside \(e,e'\), and by
\[
  \bar\pi x_e,\ \bar\pi x_{e'}\in O,\qquad
                         x_e-x_{e'}\in\pi O .
\]
The matrix lattice with this condition on every column is denoted
by \(\mathcal B_\pi^{X'}\). Put \(\delta=\delta_g/\bar\delta_g\)
with \(\delta_g\) from Lemma~\ref{lem:actual-lift-scale}.
Let \(D_\delta\) multiply the \(e'\)-entry by \(\delta\) in the
last pair and act as the identity in every middle pair.
\end{definition}

\begin{theorem}\label{thm:paired-neighbour-lattice}
For the reference lattices of Theorem~\ref{thm:paired-finite-adjoint-model}
and the projection
\(p:U_*^\perp=U_*\oplus W_*\to W_*\), the image
\(p(L_*\cap U_*^\perp)\) in the normalized coordinates of
Definition~\ref{def:paired-neighbour-lattice} equals
\(D_\delta\mathcal B_\pi^{X'}\).
There is an adelic linear map \(B:W_*\to U_*^\perp\) with
\(pB=1\) taking this image lattice back into \(L_*\).
It is zero on the string of \(Y'\) as a family of maps from
\(Y'\), that is, \(B(w)|_{\mathrm{string}(Y')}=0\) for all \(w\).
Its adjoint for the Hom pairings satisfies \(B^\dagger\tau=0\)
for the frequency in Proposition~\ref{prop:paired-rational-frequency}.
\end{theorem}
\begin{proof}
Each of the three original stage lattices is a direct sum of its
intersections with the individual grades. For even string length
all columns in Definition~\ref{def:actual-row-lift-spaces} are
homogeneous. For odd string length the changed column uses two
vectors of grade zero and every other column is homogeneous.
For \(Y'\) the degree-zero block \(X'\) has the ordinary integral
basis. The Hom lattice therefore decomposes into its separate
source and target grade blocks, by testing a map on each piece
of the source lattice. By the description of \(L_*\) in
Theorem~\ref{thm:paired-finite-adjoint-model}, after the first
positive-degree reduction, its degree-zero image consists exactly
of independent columns on that basis of \(X'\) in
\(\bar\alpha L_Y^0\oplus\alpha L_{Y''}^0\). Here the superscript
zero denotes intersection with the respective grade-zero space.
All other grade components of a lift of these columns can be
set to zero.

For \(s\ge3\), a lift in \(U_*^\perp\) of one normalized column
has outside components
\(
   \bar\alpha x_-+t z,\ \alpha x_++t z''
\)
for a common \(t\), since the additional isotropic direction is
\(\xi=z+z''\) by Proposition~\ref{prop:actual-pair-grade-zero}.
Set \(A=\bar\alpha u_j,\ A'=\alpha u_{j+2}\).
The column of the original lattice with coefficient \(1/\bar\pi\)
on the distinguished definite coordinate has coefficient \(u_j/N\)
on \(z\). Reading its coordinates and the independent column \(u_jz\)
in Definition~\ref{def:actual-row-lift-spaces} gives the equivalent
conditions
\[
 \begin{gathered}
 \bar\pi x_e,\ \bar\pi x_{e'}\in O,\qquad
        t/A-x_e/\pi\in O,\qquad t/A'-x_{e'}/\pi\in O ,
 \end{gathered}
\]
with ordinary integrality on the other normalized entries.
The same basis formula at \(j+2\) gives the second test on \(t\).
Since \(u_j/u_{j+2}=a_f/\bar a_f\),
\[
      A/A'=(\bar\alpha/\alpha)(a_f/\bar a_f)=\delta
\]
by Lemma~\ref{lem:actual-lift-scale}. Above \(N\) this ratio
equals one, as the marked primes are prime to \(N\) by
Lemma~\ref{lem:split-level}. The two tests there are soluble
exactly when \(x_e-x_{e'}\in\pi O_w\), and one may take
\(t=A x_e/\pi\). Away from \(N\) the scalars \(\pi,\bar\pi\)
are units. Both distinguished entries are then integral and
\(t=0\) passes both tests. These local choices are linear in
the column and give the desired adelic lift. They prove the
claimed lattice equality for the middle pairs.

For \(s=1\) there is no second isotropic direction. In the
original outside first space the \(z\)-coordinate coming from
the positive residual coordinate is \(\alpha x_{e'}/\pi\).
The original lattice conditions are thus
\[
   \bar\pi x_e\in O,\qquad
        \delta^{-1}x_{e'}-x_e\in\pi O,
       \qquad \delta=\bar\alpha u_j/\alpha .
\]
Indeed \(j=m-2\) gives \(u_j=a_f/\bar a_f\) in
Lemma~\ref{lem:actual-lift-scale}. The conditions already imply
\(\bar\pi\delta^{-1}x_{e'}\in O\), since \(\pi\bar\pi=N\) is
integral. The other coordinates are ordinary integral in the
chosen scales. This is precisely \(D_\delta\mathcal B_\pi\).
Take the lift into the grade-zero block itself.

In both cases let the lifted columns vanish on every other
basis vector of \(Y'\). Proposition~\ref{prop:paired-rational-frequency}
says that \(\tau\) vanishes on \(X'\). The trace pairing of a map
supported on \(X'\) with \(\tau\) is zero: the source Gram matrix
and its inverse are block diagonal for \(X'\) and its orthogonal
string, so all diagonal contributions vanish. The definition of
the adjoint and nonsingularity on \(W_*\) give \(B^\dagger\tau=0\).
\end{proof}

\section{The residual vacuum as a sum of translates}
\begin{theorem}\label{thm:paired-neighbour-vacuum-sum}
Use normalized coordinates on \(W_*\) as in
Definition~\ref{def:paired-neighbour-lattice}. Put \(L_{\rm st}\)
equal to the integral matrix lattice in these coordinates,
\(L_{\rm res}=D_\delta\mathcal B_\pi^{X'}\),
\(d'=\dim_E X'\), and let \(v_{\rm st}={\bf1}_{L_{\rm st}}\)
in \(\mathcal H(W_*,L_{\rm st})\). In formulas for the oscillator,
use the pullback of the physical hermitian form under coordinate
normalization. Define finite sets of representatives as follows.
\begin{enumerate}
\item For each column and each \(w\mid\bar\pi_N\), choose one
representative in \(\bar\pi_{N,w}^{-1}O_w\) of each class
modulo \(O_w\), with zero representing zero. Insert the chosen
element into entries \(e,e'\) of that column at \(w\) and take
zero elsewhere. The simultaneous choices form a set \(\mathcal R\).
\item For \(s\ge3\) let \(\mathcal D=\{0\}\).
For \(s=1\) and every column at each \(w\mid\pi_g\), choose one
representative of each class of \(\delta_w O_w/O_w\) in
\(\delta_w O_w\), with zero representing zero. Insert it only
in entry \(e'\). Take zero elsewhere, and call the resulting
set \(\mathcal D\).
\end{enumerate}
Then \(h_N=|\mathcal R|=N^{f d'}\) and \(h_g=|\mathcal D|\)
is one for \(s\ge3\) and \((g_0g_1)^{f d'}\) for \(s=1\).
The unique unit vector fixed by \(L_{\rm res}\) and with positive
real overlap with \(v_{\rm st}\) is
\[
 \begin{split}
 v_{\rm res}
   &=(h_Nh_g)^{-1/2}\sum_{d\in\mathcal D}
                        \sum_{r\in\mathcal R}D_{r+d}v_{\rm st}\\
   &=\left(\frac{h_g}{h_N}\right)^{1/2}
       \frac1{h_g}\sum_{d\in\mathcal D}D_d
                      \left(\sum_{r\in\mathcal R}D_r v_{\rm st}\right).
 \end{split}
\]
\end{theorem}
\begin{proof}
In these coordinates the physical form is \(\alpha\bar\alpha\)
times the trace matrix pairing with signs minus on the first
block and plus on the second. Therefore \(L_{\rm st}\) is
self-dual and \(B_f\) is one on its square, by the trace-dual
calculation of Lemma~\ref{lem:hom-oscillator-dual-lattices}.
The corresponding properties for \(L_{\rm res}\) follow from
Corollary~\ref{cor:residual-pair-positive-quotient} and
Theorem~\ref{thm:paired-neighbour-lattice}.

At \(w\mid\pi_N\), the neighboring lattice has integral entries
whose two distinguished residues agree, as \(\bar\pi_{N,w}\)
is a unit. At \(w\mid\bar\pi_N\) it is the sum of the ordinary
lattice and the diagonal \((t,t)\) in the two entries with
\(t\in\bar\pi_{N,w}^{-1}O_w\). Indeed \(\pi_{N,w}\) is a unit,
so the two entries have integral difference. Off \(N\) the two
conditions in Definition~\ref{def:paired-neighbour-lattice}
give the ordinary lattice. For the last pair the additional
factor \(\delta\) on \(e'\) has valuation one at \(w\mid\bar\pi_g\),
minus one at \(w\mid\pi_g\), and zero at all other places, by
Lemma~\ref{lem:actual-lift-scale}. Hence the enlargements of
the standard lattice modulo its intersection with \(L_{\rm res}\)
are exactly the independent diagonal classes of \(\mathcal R\)
and the one-coordinate classes of \(\mathcal D\). Every \(r+d\)
belongs to \(L_{\rm res}\). The split degree-one residue rings
and their number \(f\) over each selected side in
Lemma~\ref{lem:split-level} give the cardinalities.

For any two self-dual lattices \(L_0,L_1\) with trivial
bicharacters on their squares, the normalized sum of translates
of the first vacuum over \(L_1/(L_1\cap L_0)\) is a unit vector
fixed by \(L_1\). The proof in
Theorem~\ref{thm:finite-oscillator-lattice-transport} used just
these lattice properties: disjoint supports give the norm,
translation permutes the summands, averaging back over \(L_0\)
recovers a nonzero multiple of the reference vacuum, and the
translates of the sum are an orthonormal basis up to the
lattice-coset relations. In that basis self-duality leaves
one fixed line. The identity coset yields positive overlap,
so the stated vector is the unique normalized choice.
The supports of \(d\) and \(r\) are on disjoint conjugate pairs,
whence \(D_dD_r=D_{r+d}\) by Lemma~\ref{lem:theta-positive-translation}.
This proves both expressions.
\end{proof}

\begin{corollary}\label{cor:paired-finite-zero-term-index}
There is a surjective reduction from \(\mathcal H(V_*,L_*)\)
onto \(\mathcal H(W_*,L_{\rm st})\) invariant under \(U_*\)
and equivariant under the physical residual translations, which
takes the reference vacuum to \(v_{\rm res}\) of
Theorem~\ref{thm:paired-neighbour-vacuum-sum}. Any linear map
with these invariances into this same residual space differs
from this reduction by a scalar. On the input
\(D_\tau D_y{\bf1}_{L_*}\) its normalized value has the form
\[
             \epsilon_y D_{B^\dagger y}v_{\rm res},
             \qquad \epsilon_y=0\ \text{or a root of unity},
\]
where the adjoint translation is expressed in the normalized
residual coordinates. The integral lift \(B\) is the one in
Theorem~\ref{thm:paired-neighbour-lattice}.
\end{corollary}
\begin{proof}
Identify the oscillator on \(L_{\rm res}\) with that on \(L_{\rm st}\)
by sending \(D_w{\bf1}_{L_{\rm res}}\) to \(D_w v_{\rm res}\).
The basis and lattice-coset relations in the proof of
Theorem~\ref{thm:paired-neighbour-vacuum-sum} give an isometric
bijection intertwining translations. Compose with the exact
quotient in Corollary~\ref{cor:residual-pair-positive-quotient}.
The vacuum of any map with the stated invariances is fixed
by \(L_{\rm res}\), by lifting it inside \(L_*\) as in
Proposition~\ref{prop:integral-adjoint-reduction}. That fixed
line and the spanning translates of the quotient prove scalar
uniqueness. For the value at the shifted frequency apply
Corollary~\ref{cor:shifted-integral-adjoint-reduction} and the
identity \(B^\dagger\tau=0\) in Theorem~\ref{thm:paired-neighbour-lattice}.
\end{proof}

\section{Marked denominators under the adjoint}
\begin{definition}\label{def:finite-pair-allowed-shifts}
For an arrow with its two reference stage lattices call a
finite map \(y\) an allowed shift for \(P\) if it can be written
\(y=\alpha y_0\), with \(y_0\) zero outside \(P\cup\bar P\)
and the places dividing \(\pi_g\), and
\(\bar\delta_g y_0\) an integral Hom at every \(w\mid\pi_g\).
At the twist places impose no size restriction.
For definite arrows use the standard integral lattices in the
coordinates of Proposition~\ref{prop:actual-pair-grade-zero}.
The normalization \(\alpha\) is that of Lemma~\ref{lem:actual-lift-scale}.
\end{definition}

\begin{proposition}\label{prop:paired-output-shifts}
Let \(y_1,y_2\) be allowed shifts on the two original arrows.
Apply Theorem~\ref{thm:paired-finite-adjoint-model}, the normalized
shifted reduction of Corollary~\ref{cor:paired-finite-zero-term-index},
and the inverse of the adjoint identification on the two definite
arrows. Up to the common factor \((h_g/h_N)^{1/2}\), the image
of their translated vacuum tensor is an average over
\(d\in\mathcal D\), with coefficients zero or roots of unity
divided by \(h_g\), of vectors of the form
\[
 (D^1_{t_1}\otimes D^2_{t_2})\,
       \sum_{r\in\mathcal R} D^1_{\alpha r_-^\dagger}v_1
                           \otimes D^2_{\alpha r_+}v_2 .
\]
Here \(v_i\) are the two standard definite oscillator vacua,
\(r_-,r_+\) are the normalized first and second block matrices
of \(r\), and \(r_-^\dagger\) is their conjugate transpose on
the first block with the original positive forms.
Both \(t_i\) are allowed shifts of
Definition~\ref{def:finite-pair-allowed-shifts}.
\end{proposition}
\begin{proof}
Write \(y=A(y_1,y_2)=(y_1^\dagger,y_2)\). By
Theorem~\ref{thm:paired-finite-adjoint-model} the first
identification introduces only a root of unity. The shifted
reduction gives \(\epsilon_yD_{B^\dagger y}v_{\rm res}\)
by Corollary~\ref{cor:paired-finite-zero-term-index}.
All components of \(y\) above \(N\) vanish, since \(gQ\)
is prime to \(N\). The integral lift in
Theorem~\ref{thm:paired-neighbour-lattice} works place by
place, so \(B^\dagger y\) vanishes above \(N\) as well
(the adjoint can interchange conjugate places, which are both
above the same rational prime). Off \(N\) the lift has zero
correction, so its adjoint restricts to \(X'\) on the source
and projects orthogonally onto the residual target. This follows
also by pairing rectangular matrix units in the grade-zero
blocks. It creates no component away from the conjugate pairs
already present.

At the marked pairs all the original lattices split as the
ordinary definite block plus the indicated string columns,
including a separate \(u_jz\) in odd length. Indeed \(\pi_N\)
and its conjugate are units there, so the two glued generators
in Definition~\ref{def:actual-row-lift-spaces} generate both
\(b\) and \(u_jz\) individually. In the normalized coordinates
of Definition~\ref{def:paired-neighbour-lattice}, the residual
first block of \(B^\dagger y\) thus has entries in
\(\delta_g^{-1}O\) supported on \(w\mid\bar\pi_g\) at the marked
pairs: it comes from the conjugate transpose of an integral
map divided by \(\bar\delta_g\). The ordinary second block has
entries in \(\bar\delta_g^{-1}O\) supported on \(w\mid\pi_g\)
there, directly from \(y_2\).

For \(s=1\), the extra second-block coordinate comes from the
\(z\)-coefficient of \(y_1^\dagger\). At \(w\mid\bar\pi_g\)
this coefficient belongs to \(\bar\alpha u_j\delta_g^{-1}O_w\).
Expressing it in the positive basis vector \(z/\pi_N\) and
dividing by \(\alpha\) puts the extra coordinate in
\[
    \pi_N(\bar\alpha u_j/\alpha)\delta_g^{-1}O_w
                    =\pi_N\bar\delta_g^{-1}O_w=O_w .
\]
At \(w\mid\pi_g\) this contribution vanishes. The equality of
scales is the final-pair computation of
Theorem~\ref{thm:paired-neighbour-lattice}. The integral extra
entries may be removed modulo \(L_{\rm st}\) at the cost of
a root of unity, by the lattice-coset relation of
Theorem~\ref{thm:finite-oscillator-lattice-transport}. This
cost is independent of \(r\), since every \(r\) is supported
above \(N\).

Use the sum of Theorem~\ref{thm:paired-neighbour-vacuum-sum}.
Adding \(d\in\mathcal D\) to \(B^\dagger y\) affects only the
extra positive coordinate at \(w\mid\pi_g\), where its entries
are in \(\bar\delta_g^{-1}O_w\) by the definition of \(\delta\).
Combining those two translations introduces a root of unity.
After the definite inverse adjoint identification the first
normalized block is conjugate transposed again and is supported
on \(w\mid\pi_g\) with denominator \(\bar\delta_g\) there.
Both finite shifts are therefore allowed. Applying that
identification to \(D_r v_{\rm st}\) gives exactly the pure
tensor in the sum: the first original translate is supported
on only one side of the conjugate pairs above \(N\), so its
self-pairing and the exponent in \(\lambda\) of
Theorem~\ref{thm:paired-finite-adjoint-model} are zero.
The shifts away from \(N\) introduce only their common
quadratic phase, independent of \(r\), by the intertwining
identity of the same theorem. The factor \(1/h_g\) is the
one in Theorem~\ref{thm:paired-neighbour-vacuum-sum}.
\end{proof}

\section{A common denominator for the finite tensor}
\begin{theorem}\label{thm:paired-finite-chain-denominator}
Put \(g=g_0g_1,\ n_i=r+it\) for \(0\le i\le m\).
Choose arbitrary finite isometries at each stage integral off
\(P\cup\bar P\). On each successive arrow take the corresponding
moved reference inclusion vector using the normalized intertwiners
of Theorem~\ref{thm:finite-oscillator-lattice-transport}.
For each even \(j<m\), combine the two consecutive vectors, apply
the normalized residual reduction at \(\tau\) and return to the
definite tensor product as in Proposition~\ref{prop:paired-output-shifts}.
Let \(v\) be the tensor product of these results in order.
There is a positive constant \(C_*\), independent of \(P\)
and the stage isometries, such that \(v/C_*\) is a finite
linear combination of tensors of the translated sums in
Proposition~\ref{prop:paired-output-shifts}. Each shift is allowed.
The coefficients \(c\) can be chosen so that \(Dc\) is integral
at every finite place away from \(Q=\prod_{w\in P}q_w\), where
\[
           D=g^{f(\,\sum_{i=0}^{m-1}n_i n_{i+1}+n_{m-1}\,)}
                      \ \mid\ \big(g^{f n_m^2}\big)^m .
\]
When \(P=\varnothing\) the sum of their absolute values is at most
one after every field automorphism of \(\C\).
If the sums above \(N\) are also expanded, the new coefficients
satisfy the same denominator bound, and their corresponding
untwisted sum is at most
\[
      N^{\,f((m/2)k+(t+1)(m/2)^2)} .
\]
The resulting individual shifts may also be supported on the
places above \(N\).
\end{theorem}
\begin{proof}
For arrow \(i\), Corollary~\ref{cor:twist-only-moved-denominators}
gives an expansion in allowed shifts with denominator
\(g^{f n_i n_{i+1}}\) away from \(Q\).
Indeed its bound on the opposite marked components is
\(\bar d_P^{-1}L_0=\bar\delta_g^{-1}L_0\) there, and isometric
coordinates identify \(L_0\) with the integral Hom times
\(\alpha\) by Lemma~\ref{lem:hom-oscillator-dual-lattices}.
When \(P\) is empty the sum of the absolute conjugate coefficients
is one. Expand both arrows by bilinearity and use
Proposition~\ref{prop:paired-output-shifts}. For pair \(j\)
extract its factor \((h_{g,j}/h_{N,j})^{1/2}\), using the indices
of Theorem~\ref{thm:paired-neighbour-vacuum-sum} for that pair.
The remaining average has coefficient sum at most one, and
denominator \(h_{g,j}\). This equals one except for the last
pair, for which \(d'=\dim X_{m-1}=n_{m-1}\) as the string is empty.
Consequently \(\prod_j h_{g,j}=g^{f n_{m-1}}\), where the product
is over the even \(j<m\). Take
\(C_*=\prod_j(h_{g,j}/h_{N,j})^{1/2}\).

The tensor expansion uses products of the individual coefficients.
Multiplying the denominators makes each product integral in the
stated valuation rings. For each complex field automorphism, the
sum of absolute product coefficients factors over the independent
choices. Zeros and roots of unity in the residual evaluation can
only decrease the resulting upper bound. Grouping equal translates
preserves integrality and the upper bound by addition and the
triangle inequality. This proves the asserted estimates before
the final expansion. Since \(t\ge1\),
\[
     n_{m-1}n_m+n_{m-1}
       \le(n_m-1)(n_m+1)<n_m^2 .
\]
Every earlier product \(n_i n_{i+1}\) is at most \(n_m^2\).
This proves the integer divisibility used in the statement.

For pair \(j\) the auxiliary-prime sum has exactly
\(N^{f d_{j+1}}\) terms of unit coefficient by
Theorem~\ref{thm:paired-neighbour-vacuum-sum}.
Multiplying a displacement through any one of these terms gives
a root of unity by Lemma~\ref{lem:theta-positive-translation};
the shifts of Proposition~\ref{prop:paired-output-shifts} in fact
vanish at \(N\). The only newly occupied places are those over
\(N\), so the denominator is unchanged. From
Definition~\ref{def:actual-row-lift-spaces},
\(d_{j+1}=k+(j+1)(t+1)\). The \(m/2\) successive odd integers
from \(1\) to \(m-1\) sum to \((m/2)^2\), by subtracting
consecutive squares. Multiplication of the cardinalities gives
the last bound.
\end{proof}

\begin{corollary}\label{cor:paired-finite-chain-scaling}
Replacing each normalized residual map in
Theorem~\ref{thm:paired-finite-chain-denominator} by any fixed
linear residual map with the same invariances changes its output
by one common scalar independent of the twist and stage isometries.
Suppose further that each moved finite vector is multiplied by an
algebraic scalar integral away from \(Q\), with every conjugate
of absolute value at most one when \(P=\varnothing\).
The coefficient and support conclusions, with a common scalar
factored outside the expansion, still hold.
\end{corollary}
\begin{proof}
Corollary~\ref{cor:paired-finite-zero-term-index} expresses each
fixed residual map as a scalar multiple of its normalized version.
Multiply these scalars over the even indices, along with the
common constant from Theorem~\ref{thm:paired-finite-chain-denominator}.
This statement allows a zero map, in which case the output tensor
is zero. The additional arrow scalars enter multiplicatively
by tensor multilinearity. Their product preserves the stated
finite integrality and the untwisted absolute-value bound.
\end{proof}

\section{Recovering the data at the auxiliary prime}
\begin{proposition}\label{prop:auxiliary-prime-projection}
For each standard definite Hom space use its finite oscillator
with reference lattice \(L_i\), and let \(K_{N,i}\) be the additive
subgroup of vectors in \(L_i\) supported at the places above \(N\).
Let \(P_i\) be averaging \(D_l\) over \(l\in K_{N,i}\) with mass
one. It acts on every vector by a finite average and is the projection
onto the subspace fixed by \(K_{N,i}\). On the tensor of the
two definite oscillators of a pair let \(P=P_1\otimes P_2\).
Let \(S\) be the sum of operators
\[
    S=\sum_{r\in\mathcal R}D^1_{\alpha r_-^\dagger}
                             \otimes D^2_{\alpha r_+}
\]
with representatives from Theorem~\ref{thm:paired-neighbour-vacuum-sum}.
Then \(P S b=b\) for all \(b=Pb\).
Consequently \(S P\) is a projection with image \(S(\operatorname{im} P)\).
It fixes all the translated sums in Proposition~\ref{prop:paired-output-shifts}.
\end{proposition}
\begin{proof}
On \(K_{N,i}\) the bicharacter on two lattice vectors is one by
Proposition~\ref{prop:finite-oscillator-residue-name}, so their
translations represent addition. On a translated vacuum
\(D_yv_i\) they act by the character \(l\mapsto\psi_f(\jmath H_i(y,l))\)
by Lemma~\ref{lem:theta-positive-translation}.
Each vector is a finite sum of such eigenvectors and these
characters have open kernels by Proposition~\ref{prop:additive}.
Averaging therefore factors through a finite quotient: the common
kernel is open of finite index by compactness. A nontrivial
character has zero average by translation in that quotient.
It follows that \(P_i\) retains exactly the basis vectors with
\(y|_N\in L_i|_N\), where \(|_N\) means restriction to all
places above \(N\). This last assertion uses local self-duality
from Lemma~\ref{lem:hom-oscillator-dual-lattices}. The description
proves the projection assertion.

For \(b=Pb\) the two factors of every surviving basis tensor
have integral indices above \(N\). In the sum \(Sb\) the zero
representative gives \(b\). Each nonzero \(r\) has at least one
distinguished entry with a nonintegral residue class above \(N\).
Thus at least one translated index fails the integrality test
after adding \(r_+\) or \(r_-^\dagger\) in its coordinates
divided by \(\alpha\). Its term is killed by \(P\).
This proves \(PSb=b\), first on the basis tensors and then on
their span. It follows that \((SP)^2=SP\) and that the image
is as stated. The allowed shifts of
Definition~\ref{def:finite-pair-allowed-shifts} vanish above \(N\),
hence commute with \(S\) and \(P\), by the product rule for
disjoint conjugate pairs. The unshifted vacuum tensor belongs
to \(\operatorname{im}P\), so the translated sums are fixed.
\end{proof}

\begin{corollary}\label{cor:chain-auxiliary-datum-recovery}
For a tensor product over all even \(j<m\) form
\(\mathcal P=\bigotimes_j P_j\) and
\(\mathcal S=\bigotimes_j S_j\) from
Proposition~\ref{prop:auxiliary-prime-projection}.
Every linear combination \(w\) of the reduced translated tensors
of Theorem~\ref{thm:paired-finite-chain-denominator} satisfies
\(\mathcal S\mathcal P w=w\).
If \(w=\sum_t c_t\mathcal S b_t\) with \(b_t\) allowed translated
pure vacuum tensors, then \(\mathcal P w=\sum_t c_t b_t\).
In particular recovery uses the same coefficients, including
their finite denominators and their untwisted size bound.
\end{corollary}
\begin{proof}
The operators at different tensor factors act independently.
The allowed translates \(b_t\) are all fixed by \(\mathcal P\).
Apply \(P_jS_j=1\) on the image of \(P_j\) from
Proposition~\ref{prop:auxiliary-prime-projection} in each factor
and use tensor multilinearity and then linearity of the sum.
\end{proof}

\section{Summing the finite evaluations over the stage cuts}
\begin{theorem}\label{thm:cut-sum-of-finite-evaluations}
Use \(g,Q,n_i,D,C_*\) of
Theorem~\ref{thm:paired-finite-chain-denominator}.
For \(0\le i<m\) let \(T_i\) be a finite cut with positive
multiplicities \(m_i(t)\) whose reciprocals lie in \(\Z[1/Q]\),
and assume that for \(P=\varnothing\) one has
\(\sum_{t\in T_i}m_i(t)^{-1}\le g^{f n_i^2}\).
For example the cuts of Theorem~\ref{thm:near-flag-adelic-period-cut}
with the principal marked level in the untwisted case satisfy
these bounds. For any subset \(\mathcal T\subseteq\prod_i T_i\)
choose stage representatives for each tuple integral off
\(P\cup\bar P\), with the last stage representative equal to one.
Let \(v(t)\) be its finite tensor in
Theorem~\ref{thm:paired-finite-chain-denominator}. Then
\[
      w=\frac{D}{C_*}\sum_{t\in\mathcal T}
                 \left(\prod_{i=0}^{m-1}\frac1{m_i(t_i)}\right)v(t)
\]
is a combination of the allowed translated sums of that theorem,
with coefficients integral away from \(Q\). In the untwisted case
the sum of their absolute values after each field automorphism
is at most
\[
                 D\,g^{\,f\sum_{i=0}^{m-1}n_i^2}
                       \ \le\ \big(g^{f n_m^2}\big)^{2m}.
\]
It satisfies the recovery identity
\(\mathcal S\mathcal P w=w\) of
Corollary~\ref{cor:chain-auxiliary-datum-recovery}.
If a group permutes \(\mathcal T\) preserving the product of the
multiplicities, and acts linearly on the tensor space sending
\(v(t)\) to \(v(\gamma t)\), then \(w\) is fixed by that group.
\end{theorem}
\begin{proof}
The stated cut estimates for the example follow from
Proposition~\ref{prop:adelic-period-denominator-depth} and
Proposition~\ref{prop:marked-stage-cut-size}. At each tuple the
factor \(D\) clears the required denominators of \(v(t)/C_*\)
by Theorem~\ref{thm:paired-finite-chain-denominator}.
The rational weights with denominators supported on \(Q\)
are integral at all remaining finite places. Thus multiplying
by them and summing preserves integrality and allowed support.
For the untwisted sum, after taking absolute values of the
conjugate coefficients, the sum at each tuple is bounded by
\(D\prod_i m_i(t_i)^{-1}\). These bounds are nonnegative,
so their sum over \(\mathcal T\) is at most the sum over the
whole product, which factors into the product of the given
one-stage bounds. The last estimate uses \(n_i\le n_m\)
and the divisibility in Theorem~\ref{thm:paired-finite-chain-denominator}.
The recovery identity follows by linearity from
Corollary~\ref{cor:chain-auxiliary-datum-recovery}.
For a group element the permutation simply reindexes the
displayed sum, proving invariance.
\end{proof}

\begin{corollary}\label{cor:common-period-scalar-normalization}
Let \(\kappa\in\C\) be independent of \(P\) and all tuple choices,
and let
\[
     z=\kappa\sum_{t\in\mathcal T}
                     \left(\prod_i m_i(t_i)^{-1}\right)v(t)
\]
for a family as in Theorem~\ref{thm:cut-sum-of-finite-evaluations}.
There is a common \(c\in\C^\times\) such that \(cz\) satisfies
all the coefficient, support and recovery conclusions about
\(w\) in that theorem for the whole family. The same holds after
any fixed residual maps and the additional arrow scalars of
Corollary~\ref{cor:paired-finite-chain-scaling} are used,
with the corresponding \(v(t)\).
\end{corollary}
\begin{proof}
When \(\kappa C_*\ne0\) take \(c=D/(\kappa C_*)\); this gives
exactly the vector of Theorem~\ref{thm:cut-sum-of-finite-evaluations}.
When the common factor is zero the family of vectors \(z\)
vanishes, and \(c=1\) works with empty coefficient lists.
For the last assertion multiply \(\kappa C_*\) by the fixed
residual scalars and apply the coefficient bounds of
Corollary~\ref{cor:paired-finite-chain-scaling} term by term.
\end{proof}

\section{Algebraic characters of the norm-one torus}
\begin{definition}\label{def:torus-weight-and-units}
For the CM field \(E\), set
\[
 T(E)=\{a\in E^\times:a\bar a=1\},\qquad
 T_f=\{x\in\mathbb A_{E,f}^\times:x\bar x=1\},\qquad
 K_T=T_f\cap\widehat O_E^\times .
\]
Identify \(T_\infty\) with \((S^1)^\Phi\) in the embeddings of
Proposition~\ref{prop:places}, where \(S^1=\{z\in\C:|z|=1\}\).
For \(e=(e_\tau)_{\tau\in\Phi}\in\Z^\Phi\) put
\(p_e(z)=\prod_\tau z_\tau^{e_\tau}\).
Let \(\theta:K_T\to S^1\) be a continuous character of finite
order, meaning some fixed positive power is identically one.
A torus character of these data is a continuous character on
\(T_\infty\times T_f\) with values in \(S^1\), trivial on the
diagonal \(T(E)\), whose value on \(T_\infty\times K_T\)
is \(p_e(z)\theta(k)\).
\end{definition}

\begin{theorem}\label{thm:torus-characters-elementary}
A torus character of Definition~\ref{def:torus-weight-and-units}
exists exactly when
\[
          \prod_{\tau\in\Phi}\tau(\zeta)^{e_\tau}\,\theta(\zeta)=1
                  \qquad\text{for every root of unity }\zeta\in E .
\]
For any such character \(\mu\) and \(x\in T_f\) there are
\(h\ge1\) and \(a\in T(E)\) such that \(x^h/a\in K_T\) and
\[
                 \mu(1,x)^h=\prod_{\tau\in\Phi}\tau(a)^{-e_\tau}.
\]
In particular \(\mu(1,x)\) is algebraic and every complex
conjugate has absolute value one. If \(x\) is a unit at every
place over a rational prime \(p\), its character value is a
unit at every algebraic place over \(p\).
\end{theorem}
\begin{proof}
The ideal of a finite idele \(x\) is
\(I(x)=\prod_w\mathfrak p_w^{v_w(x_w)}\) by
Proposition~\ref{prop:general-ideal-factorization}.
For \(x\in T_f\), the exponents are opposite at conjugate
places and zero at every fixed place. Hence \(I(x)=J/\bar J\)
for a fractional ideal \(J\), by choosing one member of each
of its finitely many conjugate pairs. The group of all such
ideals \(J/\bar J\), modulo ideals of elements of \(T(E)\),
is an image of the ordinary ideal class group: the formula
\(J\mapsto J/\bar J\) sends a principal ideal \((b)\) to
\((b/\bar b)\). The ordinary ideal class group is finite by
Lemma~\ref{lem:bounded-ideal-classes}. The kernel of \(I\)
on \(T_f\) is \(K_T\), so \(T_f/(K_T T(E))\) injects into
this finite image.

An element of \(T(E)\) integral at every finite place belongs
to \(O_E\) by Proposition~\ref{prop:general-ideal-factorization}.
In every embedding it has complex absolute value one because
conjugation in \(E\) induces complex conjugation at each embedding.
All its nonnegative powers lie in a bounded part of the lattice
\(O_E\) of Lemma~\ref{lem:number-field-lattice}. Thus two powers
coincide and it is a root of unity. Conversely roots of unity
are units of norm one. This proves that the intersection of the
diagonal \(T(E)\) with \(T_\infty\times K_T\) is precisely the
group tested in the statement. The test is necessary. If it holds,
the character \(p_e\theta\) and the trivial character on
the diagonal define a character on the subgroup generated by them:
two decompositions differ by this intersection and give equal
values. That subgroup is open and has finite index by the
finite ideal quotient just proved.

For an abelian group and a subgroup of finite index, a character
with values in \(S^1\) extends to the whole group. Indeed adjoin
one generator of the finite quotient at a time. If its first
power in the current subgroup is its \(n\)-th power, assign to
the generator an \(n\)-th root on \(S^1\) of the value there,
which exists by writing a complex unit as \(e^{i\vartheta}\).
On elements \(ba^j\) multiply the old value on \(b\) by that
root to the \(j\)-th power. Equality of two expressions reduces
to the specified relation at the \(n\)-th power, so the rule is
a homomorphism. Continuity follows on the open subgroup from
the construction on \(T_\infty\times K_T\) and on all its cosets
by translation. This proves existence.

Finiteness provides \(h_0\ge1\) with \(x^{h_0}=a_0 k\),
\(a_0\in T(E),\ k\in K_T\). Raise this equation to the order
of \(\theta\) to get \(x^h=a k'\) with \(\theta(k')=1\).
The diagonal law now gives the displayed character power.
Its right side is algebraic, and each of its complex conjugates
has modulus one because every embedding sends \(a\bar a=1\)
to the corresponding norm-one equation. An \(h\)-th root is
again algebraic and has the same modulus property. At primes
over \(p\) in the last assertion \(a=x^h/k'\) is a local unit
at every place of \(E\) above \(p\). Therefore each conjugate
\(\tau(a)\) has valuation zero over \(p\), by the correspondence
of embeddings and places in Lemma~\ref{lem:place-fields}.
Taking valuations in the power formula proves the claim.
\end{proof}

\section{The cyclotomic unit and the marked characters}
\begin{lemma}\label{lem:actual-root-and-unit-witness}
Write \(n=N_1,\ \zeta=\exp(2\pi i/n)\) and
\(K=\Q(\sqrt{-d})\) as in Definition~\ref{def:levels}. The roots
of unity of \(E=K(\zeta)\) are exactly the powers of \(\zeta\).
The element \(u=i(\zeta-1)\), where \(i=\zeta^{n/4}\),
belongs to \(O_E^\times\) and satisfies
\[
         u/\bar u=\zeta,\qquad N_{E/K}u=1 .
\]
Consequently for every root of unity \(\xi\in E\) there is a
unit \(u_\xi\) with \(u_\xi/\bar u_\xi=\xi\) and \(N_{E/K}u_\xi=1\).
\end{lemma}
\begin{proof}
By Proposition~\ref{prop:places}, \(E\subset\Q(\zeta_{dn})\),
\(d\nmid n\) and \([E:\Q]=2\varphi(n)\). If a root of order \(a\)
lies in \(\Q(\zeta_{dn})\), so does a primitive root of order
\(l={\rm lcm}(a,dn)\). Indeed the powers \(\zeta_a,\zeta_{dn}\)
of \(\zeta_l\) have exponents \(l/a,l/(dn)\) with greatest common
divisor one, and all primitive roots of a given order generate
the same cyclic group. Lemma~\ref{lem:cyclo} gives
\(\varphi(l)=\varphi(dn)\). Since \(dn\) is even, any new prime
would be odd and introduce a factor at least two, and raising
an existing prime exponent introduces its prime as a factor.
This follows by counting units in each prime power quotient
and applying the Chinese remainder calculation of
Lemma~\ref{lem:basic-arith}. Thus \(a\mid dn\). If also \(d\mid a\)
and the root belongs to \(E\), then \(E\) contains \(\zeta_d\)
and \(\zeta_n\), hence \(\zeta_{dn}\). This has degree
\((d-1)\varphi(n)>2\varphi(n)\), since \(d\equiv7\bmod8\) by
Proposition~\ref{prop:places}. Therefore \(a\mid n\).

The relative norm from \(E\) to \(K\) is the product over the
embeddings of Proposition~\ref{prop:places}: over \(\C\)
evaluation at the distinct primitive \(n\)-th roots diagonalizes
multiplication in the \(\zeta\)-power basis, so the product is the
multiplication determinant. Pairing inverse roots gives
\(N_{E/K}\zeta=1\) as \(n>2\). For the cyclotomic polynomials
of Lemma~\ref{lem:cyclo}, factoring \((X^n-1)/(X-1)\) and
evaluating at one gives
\[
                 n=\prod_{h\mid n,\ h>1}P_h(1).
\]
All factors are positive integers by conjugate pairing, with
\(P_2(1)=2\). For \(h=p^a,\ a\ge1\), the quotient
\((X^{p^a}-1)/(X^{p^{a-1}}-1)\) gives \(P_{p^a}(1)=p\).
Their product over all prime-power divisors is already \(n\).
As \(6\mid n\) it follows that \(P_n(1)=1\). The even degree
\(\varphi(n)\) therefore gives \(N_{E/K}(\zeta-1)=1\).
The inverse of \(\zeta-1\) is the product of its other integral
relative conjugates. Since \(i\) is a root of unity, \(u\) is
a unit of norm one over \(K\). Direct conjugation gives
\(\bar u=u/\zeta\). Taking powers of \(u\) proves the last assertion.
\end{proof}

\begin{proposition}\label{prop:actual-marked-torus-characters}
Write \(A=M_\omega\) as in Theorem~\ref{thm:data} and
\(w_0(\tau)\) as in Definition~\ref{def:heights}. At each marked
\(w\mid\bar\pi_g\) choose a residue character \(\omega_w\) of order
\(A\) which is an \(A\)-th power. Set \(\omega_w=1\) elsewhere.
Such choices exist. There are torus characters
\(\mu^\flat,\chi\) in the sense of
Definition~\ref{def:torus-weight-and-units}, the first having weight
\(-w_0\) and trivial character on \(K_T\), the second having weight
zero and character \(k\mapsto\prod_w\omega_w(k_w\bmod w)\)
on \(K_T\). There also exists a character of constant weight one
and trivial character on \(K_T\).
\end{proposition}
\begin{proof}
At a marked \(w\mid g_i\) the residue field has \(g_i\) elements
by Lemma~\ref{lem:split-level}. Its multiplicative group is cyclic
by the proof of that lemma. Theorem~\ref{thm:data} gives
\(A^2\mid g_i-1\). Send a generator to \(\exp(2\pi i/A^2)\)
and take the \(A\)-th power. This gives the required \(\omega_w\).
For any root \(\xi\in E\), Lemma~\ref{lem:actual-root-and-unit-witness}
gives \(\xi^n=1\). Since \(n\mid A\) by Theorem~\ref{thm:data},
every \(\omega_w\) is one on the reduction of \(\xi\).
The same lemma and \(N_{E/K}\zeta=1\) give
\(\prod_{\tau\in\Phi}\tau(\xi)=1\).
Finally each \(w_0(\tau)\) is \(h_2\) plus an integral multiple of
\(n\) by Definition~\ref{def:heights}. Thus the product with those
exponents at \(\xi\) is the \(h_2\)-th power of its relative norm
and also equals one. Theorem~\ref{thm:torus-characters-elementary}
now constructs all three characters.
\end{proof}

\section{The splitting character on the full idele group}
\begin{lemma}\label{lem:real-inert-sign-parity}
Let \(F\) be the fixed field of conjugation in \(E\) of
Definition~\ref{def:levels}. Call a prime of \(F\) inert
if it has just one prime in \(E\). For all \(a\in F^\times\),
\[
        \prod_{\tau\in\Phi}\operatorname{sgn}\tau(a)
             =(-1)^{\sum_{\mathfrak p\ {\rm inert}}v_{\mathfrak p}(a)} .
\]
\end{lemma}
\begin{proof}
One has \(E=F(i)\): the real and imaginary parts
\((z+\bar z)/2,(z-\bar z)/(2i)\) of any \(z\in E\) belong to \(F\).
All embeddings of \(F\) are real and are obtained once by restricting
\(\Phi\), since each extends in two ways to \(F(i)\) as in
Theorem~\ref{thm:elliptic-trace-fields}.
Also \(12\mid n=N_1\), so Lemma~\ref{lem:elliptic-finite-norms}
gives an unramified quadratic field or a split algebra at every
finite place. At an odd prime of residue cardinality \(q\) the
inert case occurs exactly for \(q\equiv3\bmod4\).
Every dyadic prime splits here. In fact \(E=F(\sqrt{-d})\) by
degree two and nonreality. In \(\Q_2\) the polynomial
\(X^2-X+(d+1)/4\) has roots reducing to zero and one, since
\(d\equiv7\bmod8\) by Proposition~\ref{prop:places}; its derivative
is a unit at both roots modulo two. Successive linear correction
as in Lemma~\ref{lem:split-level} gives two distinct roots in
\(\Z_2\), hence in every dyadic completion of \(F\).

At each dyadic prime choose a preimage of \(a\) by the norm
from the split algebra, for instance \((a,1)\). By the simultaneous
approximation by an element \(u\in E^\times\) in the proof of
Theorem~\ref{thm:even-sign-norms} one can arrange
\(a'=a/(u\bar u)\equiv1\bmod4\) in each dyadic local ring.
Explicitly, clear the denominators at two, prescribe
sufficiently many residue digits at the finitely many primes
above two by the Chinese remainder theorem, and divide back by
the same power of two. Then the ratios to the local preimages
are one modulo four. The real signs are unchanged and the
valuations at inert primes change by even integers.
The signed norm \(N_{F/\Q}a'\) is one modulo four by the
multiplication determinant argument of
Theorem~\ref{thm:split-classes-kernel}. Its absolute value is
the positive norm of the fractional ideal \((a')\), whose prime
factors all have odd residue cardinality. Modulo four each
inert prime contributes minus one to that cardinality, and
every split one contributes one. This uses the factorization
and multiplicativity of Proposition~\ref{prop:general-ideal-factorization}
also for negative exponents. The signed norm and its negative
have opposite residues in \((\Z/4\Z)^\times\). The sign of the
norm is thus its asserted parity, and the norm is the product
of all real conjugates. This proves the formula.
\end{proof}

\begin{theorem}\label{thm:actual-global-splitting-character}
There is a continuous character \(\upsilon:\mathbb A_E^\times
\to S^1\) trivial on diagonal \(E^\times\) such that
\[
 \upsilon(x)=\prod_{\tau\in\Phi}\frac{x_{\infty,\tau}}{|x_{\infty,\tau}|}
      \quad\text{if }x_f\in\widehat O_E^\times,\qquad
 \upsilon(\bar x)=\upsilon(x)^{-1}.
\]
At a place fixed by conjugation it takes value \(-1\) on an
idele consisting of a uniformizer there and ones elsewhere.
On a finite idele \(x\) with infinite component one its value
is algebraic, all its complex conjugates have absolute value
one, and it is an algebraic unit above any rational prime
at all of whose places \(x\) is a unit.
\end{theorem}
\begin{proof}
Let \(\mathcal I\) be the group of fractional ideals of \(E\).
On its conjugation-fixed subgroup define
\(\eta(J)=(-1)^{\sum_{w=\bar w}v_w(J)}\).
For \(z\in E^\times\) define
\[
            b(z)=\frac{\overline{N_{E/K}z}}{|N_{E/K}z|}.
\]
If \((z)\) is conjugation-fixed, \(z/\bar z\) is an integral
norm-one unit by Proposition~\ref{prop:general-ideal-factorization}.
It is a root of unity by the argument in
Theorem~\ref{thm:torus-characters-elementary}. Choose \(u_\xi\)
of Lemma~\ref{lem:actual-root-and-unit-witness} for this root.
Then \(a=z/u_\xi\) lies in \(F^\times\), generates the same
ideal, and has \(N_{E/K}a=N_{E/K}z\). Its relative norm is the
product of its real conjugates, by the identification of the
embeddings above. At a fixed place the valuation of \(a\) in \(E\)
equals the valuation in \(F\), as the ramification index is one.
Consequently \(b(z)=\eta((z))\) by
Lemma~\ref{lem:real-inert-sign-parity}. For a unit this also proves
\(b(z)=1\). Thus \(b\) descends to the principal ideals, and
agrees with \(\eta\) on the intersection of the two subgroups.

Multiplying these two characters gives a well-defined character
on the product of the principal and conjugation-fixed subgroups.
Its index in \(\mathcal I\) is finite by
Lemma~\ref{lem:bounded-ideal-classes}. Extend it to a circle
character \(\beta\) of \(\mathcal I\) by the finite generator
and root construction of Theorem~\ref{thm:torus-characters-elementary}.
For any \(J\), the ideal \(J\bar J\) is fixed and has even
valuation at every fixed prime, so \(\beta(J\bar J)=1\).
Define
\[
   \upsilon(x)=\left(\prod_{\tau\in\Phi}
                  \frac{x_{\infty,\tau}}{|x_{\infty,\tau}|}\right)
                          \beta(I(x_f)).
\]
The ideal map \(I\) has open kernel \(\widehat O_E^\times\),
so the character is continuous. Its values on the diagonal
are one by the definition of \(b\) and the relative norm formula
in Lemma~\ref{lem:actual-root-and-unit-witness}. The equation
for \(J\bar J\) gives the conjugation law, and \(\eta\) gives
the asserted fixed-prime sign.

If \(J=I(x_f)\), choose \(h\ge1\) such that \(J^h=(z)\) by
ordinary class finiteness. Then
\[
              \beta(J)^{2h}=\frac{\overline{N_{E/K}z}}{N_{E/K}z}.
\]
Both sides show algebraicity of \(\beta(J)\). Every complex
embedding of the imaginary quadratic field \(K\) sends numerator
and denominator to numbers of equal modulus. If \(J\) has zero
valuation above \(p\), then \(z\) and its inverse are integral
there. The same holds for every conjugate factor in the numerator
and denominator of this formula, by Lemma~\ref{lem:place-fields}.
Taking valuations and complex absolute values as appropriate
gives the last assertions.
\end{proof}

\section{The splitting multiplier on the finite representatives}
\begin{definition}\label{def:finite-splitting-multiplier}
In an isometric reference frame divide oscillator coordinates
by \(\alpha\), so the reference lattice is \(O^b\), with
\(O=\widehat O_E\) and \(b\) the Hom dimension. Let \(G\) be
a finite adelic isometry acting in these coordinates. At a place
\(w\) let \(a_w(G)\) be the valuation of the determinant of a
basis matrix for
\[
       S_w(G)=\{x\in O_w^b:G_wx\in O_w^b\}.
\]
Set \(J_G=\prod_w\mathfrak p_w^{a_w(G)}\) and
\(\Delta(G)=\beta(J_G)(N J_G)^{1/2}\), where \(N J_G\) is
the absolute ideal norm and \(\beta\) the finite ideal character
of Theorem~\ref{thm:actual-global-splitting-character}.
The operator with the splitting multiplier sends \(v_0\) to
\[
        \frac{\Delta(G)}{|\Delta(G)|^2}
                   \sum_{r\in GL_0/(GL_0\cap L_0)}D_r v_0
\]
and intertwines \(D_y\) with \(D_{Gy}\).
Here \(GL_0\) denotes the image of the reference lattice
\(L_0\) under the given operator \(G\).
\end{definition}

\begin{theorem}\label{thm:split-cut-weil-multiplier}
The integers \(a_w(G)\) are independent of the bases, nonnegative
and zero at almost all places. The operator of
Definition~\ref{def:finite-splitting-multiplier} exists uniquely
and equals \(\beta(J_G)W_G\). For isometries integral away from
\(P\cup\bar P\) one has
\[
                 \beta(J_G)=\prod_{w\in P}
                              \beta(\mathfrak p_w)^{-v_w(\det G_w)} .
\]
On this group the operators form a representation. Their phase
\(\beta(J_G)\) is an algebraic unit at every finite place away
from \(Q\), has all complex conjugates of modulus one, and is
one for integral \(G\).
\end{theorem}
\begin{proof}
A local invertible matrix is a product of two integral invertible
matrices and a diagonal one with uniformizer powers \(\pi_w^{k_i}\).
Indeed one can clear denominators first, select an entry of least
valuation as pivot, interchange its row and column to the first
position and clear its row and column by integral multiples.
Iterate on the remaining invertible square block and divide
back by the common power. Diagonal units can be absorbed into
the integral matrices. In these coordinates \(S_w(G)\) is
generated by an integral change of basis applied to the diagonal
powers \(\pi_w^{\max(0,-k_i)}\). Any two of its bases differ by
an integral invertible matrix. This proves well-definedness,
and gives
\[
          a_w(G)=\sum_i\max(0,-k_i),\qquad
                [O_w^b:S_w(G)]=q_w^{a_w(G)} .
\]
At almost all places \(G\) and its inverse are integral, giving
\(a_w(G)=0\). Taking the product of the indices and applying \(G\)
shows \(N J_G=[G L_0:G L_0\cap L_0]\), since scalar division
by \(\alpha\) does not change the index. The positive normalized
average in Theorem~\ref{thm:finite-oscillator-lattice-transport}
therefore gives precisely \(\beta(J_G)W_G\) with the prescribed
value and intertwining law.

In the isometric frames the matrix at the conjugate place is
the inverse conjugate transpose. Its diagonal exponents are
the negatives of those for \(G_w\), by the same decomposition.
As \(\beta(\bar{\mathfrak p}_w)=\beta(\mathfrak p_w)^{-1}\)
by Theorem~\ref{thm:actual-global-splitting-character}, the
joint contribution of this pair is the value at \(w\) raised
to the exponent
\[
  \sum_i\big(\max(0,-k_i)-\max(0,k_i)\big)
                 =-\sum_i k_i=-v_w(\det G_w).
\]
If a fixed Gram sign matrix is used, its integral diagonal
signs on either side leave the same exponents. The integral
complementary places contribute one, giving the displayed
character of the determinants.

On \(P\cup\bar P\) multiplicativity of \(W_G\) is
Theorem~\ref{thm:split-finite-schroedinger}.
On the complementary integral subgroup it is the change of
variables of Theorem~\ref{thm:finite-oscillator-lattice-transport}.
Their actions on the separate factors commute by the tensor
factorization in the split theorem. The phase just computed is
multiplicative too. Finally \(J_G\) has support over \(Q\).
Apply the finite character conclusions of
Theorem~\ref{thm:actual-global-splitting-character} to an idele
with ideal \(J_G\), taking components one outside this support.
For integral \(G\) the ideal is the unit ideal.
\end{proof}

\section{The finite action with the prescribed arrow characters}
\begin{definition}\label{def:actual-twisted-arrow-character}
Let \(\upsilon_T\) be the restriction of
Theorem~\ref{thm:actual-global-splitting-character} to the
norm-one torus, and let \(\mu^\flat,\chi,\omega_w\) be as in
Proposition~\ref{prop:actual-marked-torus-characters}.
A twist here is a torus character \(\vartheta\) of weight zero,
with finite unit character given by residue characters
\(\vartheta_w\) supported on \(P\), where \(P\) is the split
set of Definition~\ref{def:split-inclusion-finite-oscillator}.
Extend the residue characters by the trivial character at other
places. For \(1\le l\le m\), with \(h\) of
Definition~\ref{def:heights}, set
\[
           \mu_l=\mu^\flat(\upsilon_T^{-h}\chi\vartheta)^l,
       \qquad \mu_l^{[n]}=\mu_l\,\upsilon_T^{\lfloor n/2\rfloor}.
\]
On the arrow from rank \(n=n_{l-1}\) to \(n'=n_l\),
for finite stage elements \(g,g'\) integral away from \(P\cup\bar P\),
put \(G(v)=g'vg^{-1}\). Define
\[
   \Omega_l(g,g')=
    \frac{\mu_l^{[n]}(1,\det g')}{\mu_l^{[n']}(1,\det g)}
                        \,\beta(J_G) W_G .
\]
All determinants of stage elements are norm-one finite ideles
by the hermitian isometry equation. For \(b\in\Z\) write
\([b]_2=b-2\lfloor b/2\rfloor\).
\end{definition}

\begin{theorem}\label{thm:actual-twisted-finite-arrow}
The character \(\mu_l\) has weight \(-w_l\) and finite unit
characters \((\omega_w\vartheta_w)^l\).
The operators \(\Omega_l\) form a representation of the product
of the indicated finite stage groups. The scalar multiplying
\(W_G\) equals
\[
   \frac{\mu_l(1,\det g')}{\mu_l(1,\det g)}
       \prod_{w\in P}
          \beta(\mathfrak p_w)^{[n']_2 v_w(\det g)-[n]_2 v_w(\det g')}.
\]
It is an algebraic unit at all finite places away from \(Q\)
and has absolute value one under every complex embedding.
For integral \(k,k'\) whose reductions carry the prescribed
inclusion to itself at the marked and twist places, one has
\[
 \Omega_l(k,k') f_{a_{\rm inc}}=
  \left(\prod_{w\mid\bar\pi_g\ {\rm or}\ w\in P}
  \frac{(\omega_w\vartheta_w)^l(\det k'_w\bmod w)}
       {(\omega_w\vartheta_w)^l(\det k_w\bmod w)}\right) f_{a_{\rm inc}} .
\]
\end{theorem}
\begin{proof}
On \(T_\infty\), the splitting character restricts to weight one
since all coordinates have modulus one. On \(K_T\) it is trivial
by Theorem~\ref{thm:actual-global-splitting-character}. Multiplication
of characters adds their weights and multiplies their finite unit
parts. The claimed weights follow from
Proposition~\ref{prop:actual-marked-torus-characters} and
\(w_l=w_0+lh\) in Definition~\ref{def:heights}.
The determinants and \(G\) are homomorphisms on the product,
where the two actions on Hom commute. Theorem~\ref{thm:split-cut-weil-multiplier}
therefore gives the representation law.

For any such stage element \(g\), the contributions from the two
members of the split pair above \(w\in P\) to
\(\upsilon_T(1,\det g)\) multiply to
\(\beta(\mathfrak p_w)^{2v_w(\det g)}\):
the two valuations are opposite, and the two ideal character
values are inverse. All other finite determinant valuations
are zero. Also \(\det G_w=(\det g'_w)^n(\det g_w)^{-n'}\)
by the column and row calculation of
Corollary~\ref{cor:split-moved-indicator}. Combining these two
formulas with Theorem~\ref{thm:split-cut-weil-multiplier}
gives the stated parity exponents. The full ratio of the two
\(\mu_l\)-values satisfies the unit and absolute-value assertions
by Theorem~\ref{thm:torus-characters-elementary}. The local
factors at \(P\) satisfy them by
Theorem~\ref{thm:actual-global-splitting-character}.
For an integral pair the splitting phase and the finite
\(\upsilon_T\)-factors are one. Corollary~\ref{cor:finite-inclusion-periods}
fixes the inclusion vector under \(W_G\) for the residues in the
statement. Evaluating the remaining unit characters gives the
last formula.
\end{proof}

\begin{corollary}\label{cor:actual-character-finite-chain}
The coefficient, denominator and recovery conclusions of
Theorems~\ref{thm:paired-finite-chain-denominator}
and~\ref{thm:cut-sum-of-finite-evaluations} hold with the
finite actions \(\Omega_l\) of
Definition~\ref{def:actual-twisted-arrow-character}, for the
normalized residual maps and for the fixed scalar changes
specified in Corollary~\ref{cor:paired-finite-chain-scaling}.
The finite arrow data depend only on the right level cosets of
the stage elements, for the level reducing to identity at all
marked and twist places in each stage.
\end{corollary}
\begin{proof}
Apply the scalar conclusions of
Theorem~\ref{thm:actual-twisted-finite-arrow} in
Corollary~\ref{cor:paired-finite-chain-scaling} and in the
weighted sum therefrom in Theorem~\ref{thm:cut-sum-of-finite-evaluations}.
For two integral right multipliers in the stated levels the
residues are identities. Hence the last formula of
Theorem~\ref{thm:actual-twisted-finite-arrow} fixes the original
inclusion vector. The representation law then gives the same
moved datum on each coset.
\end{proof}

\section{Column-level equivariance of the reduced tensor}
\begin{proposition}\label{prop:level-equivariant-paired-reduction}
For an even \(j<m\) and its two successive arrows let
\(k_i\in K_{c,i}\), \(i=j,j+1,j+2<m\), with
\(K_{c,i}\) from Proposition~\ref{prop:column-level-in-stage-bessel}.
When \(j+2=m\) put \(k_m=1\). Extend these elements to the
original stage spaces by identity on the strings, writing
\(\widehat k_i\) for the extensions. On the definite outside
target use \(k_{j+2}\) in the middle pairs, and identity on
the enlarged target in the last pair.
The paired reduction of Proposition~\ref{prop:paired-output-shifts}
intertwines the tensor product of the finite actions
\(\Omega_{j+1}(\widehat k_j,\widehat k_{j+1})\) and
\(\Omega_{j+2}(\widehat k_{j+1},\widehat k_{j+2})\)
with the corresponding actions on the two definite arrows,
using their own ranks and the same \(\mu_{j+1},\mu_{j+2}\).
\end{proposition}
\begin{proof}
First omit the scalar characters and consider the lattice
actions \(W\) alone. The embedded column elements preserve the
original lattices by Proposition~\ref{prop:column-level-in-stage-bessel},
and the definite column elements preserve their standard lattices.
Each such action on the corresponding oscillator is consequently
the ordinary change of variables by the inverse isometry, by
Theorem~\ref{thm:finite-oscillator-lattice-transport}.
Theorem~\ref{thm:paired-finite-adjoint-model} intertwines the
two initial changes with the change on \(V_*\). This change fixes
\(\tau\) because the maps \(T,T'\) in
Proposition~\ref{prop:paired-rational-frequency} vanish on
the definite inputs and take values in string vectors, which are
fixed. It preserves \(U_*\) and \(W_*\): it respects the grades
and in the middle case fixes both \(z,z''\). Thus it commutes
with the projection \(p\) onto the residual part on \(U_*^\perp\).
The finite quotient of Theorem~\ref{thm:finite-grade-oscillator-quotient}
is equivariant under these changes. Indeed compose its map with
the two changes on the respective sides. After transporting the
indices back, both maps annihilate positive-translation differences,
take the vacuum to the residual vacuum, and act by the same
residual translations. Uniqueness in that theorem, applied as in
Corollary~\ref{cor:residual-pair-positive-quotient}, gives equality.

In normalized coordinates the residual action is
\(t\mapsto(k_j\oplus k_{j+2})t k_{j+1}^{-1}\), extending the
last action as specified. The separate factors \(\bar\alpha,\alpha\)
commute with these block actions. The residual isometry preserves
both \(L_{\rm res}=p(L_*\cap U_*^\perp)\) and \(L_{\rm st}\).
It sends the unique normalized \(L_{\rm res}\)-fixed vector of
Theorem~\ref{thm:paired-neighbour-vacuum-sum} to itself: it is
unitary, fixes \(v_{\rm st}\), and therefore preserves the positive
overlap with \(v_{\rm st}\). This proves equivariance after the
identification of Corollary~\ref{cor:paired-finite-zero-term-index},
and shifting by \(D_\tau\) preserves it since \(\tau\) is fixed.
The definite inverse adjoint identification intertwines the
remaining changes by Theorem~\ref{thm:paired-finite-adjoint-model}.

For integral elements the splitting multiplier is one by
Theorem~\ref{thm:split-cut-weil-multiplier}, and all powers of
\(\upsilon_T\) are trivial on integral determinants by
Theorem~\ref{thm:actual-global-splitting-character}. Each original
scalar on arrow \(l\) therefore equals the ratio of \(\mu_l\)
on the two determinants as in
Theorem~\ref{thm:actual-twisted-finite-arrow}. Extending an
element by identity preserves its determinant, also in the final
pair by Proposition~\ref{prop:actual-pair-grade-zero}. Thus the
original and definite products of scalar characters agree.
\end{proof}

\begin{corollary}\label{cor:actual-column-invariant-cut-tensor}
In Theorem~\ref{thm:cut-sum-of-finite-evaluations} take the full
product of cuts from Definition~\ref{def:finite-period-cut},
with the embedded column level in \(L_i=S_i\cap A_i\) at each
stage. This holds for \(A_i=K_i\) at positive \(i\), and for the
principal marked level in the stage-zero frame. Use
\(\Omega_l\) on the initial data. The resulting definite tensor
\(w\) is fixed by the product of the \(K_{c,i}\), \(i<m\),
acting simultaneously at the ends of the definite arrows, with
last group element one. Its recovered tensor
\(\mathcal P w\) is fixed under the same action.
\end{corollary}
\begin{proof}
For membership in \(L_i\) use Proposition~\ref{prop:column-level-in-stage-bessel},
including identity reduction at the marked places in the first
stage. Each such group permutes its cut by the defining \(L_i\)-stability.
Its multiplicity stays fixed by Lemma~\ref{lem:period-single-orbit-index},
since the permutation stays in each \(S_i\)-orbit. Acting on the
left on the stage representatives acts on the moved initial vectors
by the representation law of Theorem~\ref{thm:actual-twisted-finite-arrow}.
Changing representatives of the right level cosets preserves these
vectors by Corollary~\ref{cor:actual-character-finite-chain}.
Proposition~\ref{prop:level-equivariant-paired-reduction} therefore
gives exactly the tuple permutation law in
Theorem~\ref{thm:cut-sum-of-finite-evaluations}, proving invariance.
At \(N\) the definite group actions preserve the integral
translation subgroup in Proposition~\ref{prop:auxiliary-prime-projection}.
Conjugating its finite character average by such an action permutes
the translations and preserves normalized measure (equivalently,
it permutes the classes in each corresponding finite average).
The scalar character factors commute with the projections.
Thus \(\mathcal P\) commutes with the full product action and
\(\mathcal P w\) is fixed.
\end{proof}

\section{Fourier expansion along a rational isotropic subspace}
\begin{theorem}\label{thm:rational-frequency-expansion}
Let \(V\) be a rational hermitian space over \(E\) with nonsingular
form \(H\) and bicharacter of
Definition~\ref{def:adelic-oscillator-translations}. Let \(U\subset V\)
be a totally isotropic \(E\)-subspace and let \(U^\perp\) be its
hermitian orthogonal. Write \([U]=U(E)\backslash U(\mathbb A_E)\),
with quotient probability measure.
Let \(F\) be an adelic theta function whose restriction to \(U\)
is invariant under translation by a compact open subgroup
\(D\subset U(\mathbb A_{E,f})\) and is smooth in its real variables.
For \(w\in V(E)\) put
\[
    A_w=\int_{[U]}(D_w F)(u)\dd u,\qquad
    \chi_w(u)=\Psi\big(\jmath H(w,u)\big).
\]
Then \(A_w\) depends only on \(w\bmod U^\perp(E)\), and is zero
unless \(\chi_w|_D=1\). On all of \([U]\) one has the absolutely
and uniformly convergent expansion
\[
                  F(u)=\sum_{w\in V(E)/U^\perp(E)}
                                      A_w\,\chi_w(u)^{-1}.
\]
For a rational \(E\)-subspace \(U_1\subseteq U\) and \(t\in V(E)\),
\[
        \int_{[U_1]}(D_t F)(u)\dd u
             =\sum_{\substack{w\bmod U^\perp(E)\\w-t\in U_1^\perp(E)}} A_w .
\]
\end{theorem}
\begin{proof}
On \(U\) the rational transformation law and the commutator in
Lemma~\ref{lem:theta-positive-translation} give
\[
       (D_wF)(u)=B(u,w)B(w,u)^{-1}F(u)=\chi_w(u)F(u).
\]
The character is one on rational \(u\) by Proposition~\ref{prop:additive}.
Translation by \(D\) kills the integral unless \(\chi_w|_D=1\).
In a rational basis of the underlying \(\Q\)-space of \(U\),
density of rational vectors in the finite adeles holds by the
Chinese remainder proof in Lemma~\ref{lem:unip-approx}.
Consequently every finite vector is rational modulo \(D\), and
\(\Lambda=U(E)\cap D\) is dense in \(D\). It is a full lattice
in \(U_\infty\): in rational coordinates it contains a positive
integer multiple of the standard integer lattice and is contained
in an integer-denominator lattice, since the compact open \(D\)
contains and is contained in the corresponding finite lattices.
One obtains a smaller and a larger lattice by clearing the finitely
many exceptional finite denominators as in Lemma~\ref{lem:class-covol}.
Sending \(u_\infty\) to the class of \((u_\infty,0)\) identifies
\(U_\infty/\Lambda\) with \([U]/D\), by the preceding density and
the description of the rational kernel. Integration of a function
on this compact real torus agrees with normalized volume on a
lattice parallelepiped, by the domain construction in
Lemma~\ref{lem:theta-positive-translation}, or by translation
invariance and subdivision into equal congruent subtiles.

The map \(w\mapsto\ell_w\), where
\(\ell_w(v)=\operatorname{Tr}_{E/\Q}(\jmath H(w,v))\) on \(U(E)\),
identifies \(V(E)/U^\perp(E)\) with the full rational dual
of that \(\Q\)-space. Indeed the trace pairing on \(V\) is
nonsingular: a nonzero hermitian functional takes every value
in \(E\), including one with nonzero scaled trace. Restriction
of rational linear functionals to a subspace is onto by extending
a basis. Its kernel on restricting is the hermitian orthogonal,
because testing all \(E\)-multiples tests every trace multiple.
By Proposition~\ref{prop:additive}, \(\chi_w\) at the finite part
of \(a\in\Lambda\) equals \(\exp(2\pi i\ell_w(a))\).
Density in \(D\) therefore identifies the surviving indices with
the full lattice dual \(\Lambda^\vee\): every real linear
functional integral on a rational lattice is rational on its
rational span. At infinity \(\chi_w(u)=\exp(-2\pi i\ell_w(u))\).
The coefficients \(A_w\) are consequently the ordinary
normalized Fourier coefficients of a smooth periodic function.

For completeness, integration by parts in a largest frequency
coordinate in a lattice basis gives a bound \(C_k(1+|w|)^{-k}\)
in the dual lattice coordinates, for every \(k\). The boundary
terms cancel by periodicity. Counting in coordinate cubes gives a
convergent majorant for \(k>\dim_\R U_\infty\), also covering the
constant term separately. The product kernel argument of
Lemma~\ref{lem:poisson-finite-torus} identifies the uniform sum
with the periodic function, since their Fourier coefficients agree.
This proves the first formula. For the second multiply by
\(\chi_t\) and integrate over \([U_1]\) term by term. Translation
makes the integral of a nontrivial character zero. Triviality on
\([U_1]\) is exactly \(w-t\in U_1^\perp(E)\), as seen by varying
the real coordinate along \(U_1\) and taking all trace multiples.
This gives the required sum, including \(U_1=0\).
\end{proof}

\begin{corollary}\label{cor:periodized-rational-frequency-expansion}
The functions \(\Theta_{\phi,f}\) of
Theorem~\ref{thm:adelic-periodized-oscillator} satisfy the hypotheses
of Theorem~\ref{thm:rational-frequency-expansion} for every
rational totally isotropic subspace \(U\).
\end{corollary}
\begin{proof}
Every finite vector \(f\) is a finite sum of translates of a
lattice vacuum. It is fixed by a common open subgroup of that
lattice by the finite eigencharacter calculation of
Theorem~\ref{thm:finite-oscillator-lattice-transport}.
Intersect with \(U(\mathbb A_{E,f})\). On \(U\) these Heisenberg
translations are ordinary translations because \(H(U,U)=0\).
Finite equivariance of Theorem~\ref{thm:adelic-periodized-oscillator}
gives the required invariance. Differentiation of each summand
in real coordinates brings derivatives of the Schwartz factor
and polynomials in the rational summation coordinates from the
character. The lattice counting and Schwartz bound in that
theorem still give an absolutely locally uniform sum for every
such derivative. This proves smoothness on the real torus.
\end{proof}

\section{A common partial sum for two target gradings}
\begin{definition}\label{def:two-target-gradings}
Choose consecutive spaces \(Y,Y'\) in the tower of
Definition~\ref{def:actual-row-lift-spaces}. Denote their
homomorphism space by \(V=\operatorname{Hom}_E(Y,Y')\). Use \(T\) from
Definition~\ref{def:row-string-lowering-map} for a source string
of length at least two, and \(T=0\) for length zero or one.
Write \(g_a,g'_b\) for the respective grades and \(a^*,b^*\)
for the indices paired by the Gram matrices. Thus the star
reverses each string and fixes each index of the definite block.
For \(d_b=0\) at every \(b\), or \(d_b=g'_b\) at every \(b\),
write \(U_d\) for the span of the matrix units \(E_{ba}\) with
\(d_b>g_a\), and \(N_d\) for the span of those with \(d_b\ge g_a\).
Let
\[
 C=\{(b,a): g_a>\max(0,g'_b)\},\qquad
 U_C=\operatorname{span}_E\{E_{b^*a^*}:(b,a)\in C\}.
\]
Say \(w\) agrees with \(T\) on \(C\) if \(w_{ba}=T_{ba}\)
for all \((b,a)\in C\).
\end{definition}
\begin{proposition}\label{prop:common-partial-frequency-sum}
Suppose \(F\) satisfies the hypotheses of
Theorem~\ref{thm:rational-frequency-expansion} on both \(U_0\)
and \(U_{g'}\). Write \(I^U_w(F)=\int_{[U]}(D_wF)(u)\dd u\).
For either choice of \(d\),
\[
  \sum_{\substack{w\in V(E)/N_d(E)\\w|_C=T|_C}} I^{U_d}_w(F)
                = I^{U_C}_T(F).
\]
Both sums converge absolutely. In particular the identity applies
to the periodized data of
Corollary~\ref{cor:periodized-rational-frequency-expansion}.
\end{proposition}
\begin{proof}
In \(H(A,B)=\operatorname{tr}(J^{-1}A^*J'B)\), the matrix
unit \(E_{ba}\) pairs only with \(E_{b^*a^*}\), with nonzero
coefficient. This follows by multiplying the two antidiagonal
or diagonal Gram matrices of Definition~\ref{def:actual-row-lift-spaces}.
The signs of both grades reverse under star. Thus \(U_d\) is
totally isotropic and \(U_d^\perp=N_d\). The indices common
to the two positive spaces are precisely the reversals of \(C\),
so \(U_C=U_0\cap U_{g'}\). Orthogonality to \(U_C\) means
vanishing at every entry in \(C\). It is unchanged on adding
any element of \(N_d\). The subspace formula in
Theorem~\ref{thm:rational-frequency-expansion}, with \(t=T\),
now gives the assertion.
\end{proof}
\begin{lemma}\label{lem:constant-term-adelic-substitution}
Let \(U\) be a rational totally isotropic subspace of a
nonsingular hermitian \(E\)-space. For any adelic theta function
\(F\), adelic \(w\), and \(v\in U(\mathbb A_E)\), one has
\[
 I^U_{w+v}(F)=B(v,w)^{-1}I^U_w(F).
\]
For a rational hermitian isometry \(G\) carrying \(U\) onto itself,
\(I^U_w(F\circ G^{-1})=I^U_{G^{-1}w}(F)\).
\end{lemma}
\begin{proof}
Heisenberg multiplication as in Lemma~\ref{lem:theta-positive-translation}
gives \(D_vD_w=B(v,w)D_{w+v}\). On \(U\), \(D_v\) acts by
ordinary translation since the pairing is zero. Invariance of
the probability on \([U]\) gives the first formula. For the second,
\(B(Gx,Gy)=B(x,y)\) gives
\(D_w(F\circ G^{-1})=(D_{G^{-1}w}F)\circ G^{-1}\).
Rationality makes \(G\) an automorphism of the compact quotient.
It preserves the invariant probability: for any two invariant
probabilities \(\mu,\nu\) on the quotient and continuous \(h\),
integrating \(h(x+y)\) against their product gives both
\(\int h\dd\mu\) and \(\int h\dd\nu\), by
Lemma~\ref{lem:product-integration}. Apply this to \(\mu\) and
its transport. Integration proves the claim.
\end{proof}

\section{The moment of a surviving frequency}
\begin{lemma}\label{lem:two-string-moments}
Use Definition~\ref{def:two-target-gradings}. For a string of
length \(l\) define \(S_l\) to be zero on the definite part and
on its last string vector, if present. For \(i+1<l\) put
\[
 S_l e_i=\sigma_l(i)e_{i+1},\qquad
 \sigma_l(i)=
 \begin{cases}
  -Nc,& l-1-2i=0,\\
  -Nc^2,&l-1-2i=1,\\
  c,&\text{otherwise}.
 \end{cases}
\]
For source and target lengths \(s,s'=\max(s-1,0)\), respectively,
\(T^\dagger T=S_s\) and \(TT^\dagger=S_{s'}\), where the dagger
denotes the hermitian adjoint. Every displayed shift coefficient
is nonzero.
\end{lemma}
\begin{proof}
For \(s\le1\) all three maps are zero. Otherwise, with the
\(h_i,\tau_i,t_i\) of Definition~\ref{def:row-string-lowering-map},
write \(h'_i\) for the corresponding target Gram entry. Adjointness
at \(e'_i\) for \(0\le i<s-1\) gives
\[
 T^\dagger e'_i=\frac{h'_i}{h_{i+1}}\bar t_{s-2-i}e_{i+1}
                =h'_i\tau_{s-2-i}e_{i+1}.
\]
Indeed \(h_{s-2-i}=h_{i+1}\) and all the scalars belong to the
real subfield. Both adjoints kill the respective definite blocks.
The coefficient of \(T^\dagger T e_i\) is
\(h_i h'_i\tau_i\tau_{s-2-i}\). The two source grades in the
\(\tau\)'s are \(g_i\) and \(2-g_i\). For even \(s\), exactly
one is positive except when both are one. In the latter case
both factors are \(c\) and \(h'_i=-N\). For odd \(s\) exactly
one is nonpositive since the grades are even, and \(h_i=-N\)
exactly at grade zero. This gives \(\sigma_s(i)\).
For \(i+1<s'\) the coefficient of \(TT^\dagger e'_i\) is
\(h'_i h_{i+1}\tau_{s-2-i}\tau_{i+1}\), with the two source
grades \(1-g'_i,g'_i-1\). When they differ, exactly one factor
is \(c\). The factor \(h'_i\) is \(-N\) just at \(g'_i=0\).
When the two grades are equal they are zero, \(s\) is odd,
both factors are \(c\), and \(h_{i+1}=-N\). This gives
\(\sigma_{s'}(i)\). The last target string vector maps under
\(T^\dagger\) to the last source string vector and is therefore
killed by \(TT^\dagger\). Nonvanishing follows from \(cN\ne0\).
\end{proof}
\begin{definition}\label{def:unipotent-frequency-test}
For a source unitary adelic isometry \(u\), say \(u-1\) has
degree at least two if every nonzero entry raises the grade by
at least two. A rational \(w:Y\to Y'\) \emph{survives} for a
target grading \(d\) of Definition~\ref{def:two-target-gradings}
if for every such \(u\) with \(w(u-1)\in U_d(\mathbb A_E)\),
\[
                    B(Tu,T)=B(wu,w).
\]
Here \(B\) is the bicharacter on the homomorphism space in
Definition~\ref{def:adelic-oscillator-translations}.
\end{definition}
\begin{proposition}\label{prop:eichler-frequency-coordinate}
Suppose \(w\) survives for \(d\). Let \(x\in Y(E)\) and let \(k\)
be a source basis index, of grade \(g_k\). Assume that every grade
in \(x\) is at least \(\max(1,2+g_k)\), every target grade in
\(wx\) is strictly greater than \(\max(g_k,-1)\), and every sum
of a target grade in \(w e_{k^*}\) and a source grade in \(x\)
is positive. A zero vector has empty grade support. Then
\[
                       ((T^\dagger T-w^\dagger w)x)_k=0 .
\]
\end{proposition}
\begin{proof}
Use \(h(v,z)=v^*Jz\), linear in \(z\), and put
\(M=T^\dagger T-w^\dagger w\). For arbitrary \(t\in\mathbb A_E\)
take \(y=t e_{k^*}\) and \(\kappa=h(y,y)/2\). The hypotheses
on grades imply \(h(x,x)=h(x,y)=0\). Write \(v^\flat=h(v,\cdot)\)
and form
\[
 R=x y^\flat-y x^\flat-\kappa x x^\flat,\qquad u=1+R .
\]
Indeed \(R+R^\dagger=-2\kappa x x^\flat\) and
\(R^\dagger R=2\kappa x x^\flat\), so \(u^\dagger u=1\).
The two mixed terms of \(R\) raise grade at least two since
\(g_a-g_k\ge2\) on the support of \(x\). The last term raises
it by at least two since all grades of \(x\) are positive.
In particular \(R\) is nilpotent and \(1+R\) is invertible.
The first term in \(wR\) sends grade \(g_k\) into grades above
\(g_k\). For the second term the last support hypothesis gives
strictly positive degree. For the third, the output grades are
nonnegative and the input grades are negative. Thus \(wR\in U_d\).

Write \(A=h(y,Mx)\) and \(D=h(x,Mx)\). The trace of a rank-one
map \(v\ell\) is \(\ell(v)\) by summation of its diagonal.
Hermitian symmetry gives
\(\operatorname{tr}(MR)=A-\bar A-\kappa D\). If \(g_k\ne0\)
then \(\kappa=0\). If \(g_k=0\), \(x\) has grades at least two,
so \(Tx\) has strictly positive target string grade, and \(wx\)
has strictly positive \(d\)-grade. Both have norm zero since
opposite grades pair. Consequently \(D=0\) in this case.
By rationality \(\Psi(\jmath\bar\omega\operatorname{tr}M)=1\).
The survival identity is therefore
\(\Psi(\jmath\bar\omega(A-\bar A))=1\), by the trace expression
for \(H\) in Proposition~\ref{prop:common-partial-frequency-sum}.
Conjugate invariance of \(\Psi\), the relation
\(\bar\jmath=-\jmath\), and \(\omega+\bar\omega=1\)
turn the left side into \(\Psi(\jmath A)\).
Here \(A=\bar t h(e_{k^*},e_k)(Mx)_k\).
The Gram coefficient is nonzero. If \((Mx)_k\ne0\), varying
\(t\) at one infinite place gives every real trace value and
hence a character value different from one, by
Definition~\ref{def:adelic-oscillator-translations}. This
contradiction proves the assertion.
\end{proof}

\section{The orbit of the positive columns}
\begin{lemma}\label{lem:rational-hermitian-witt-extension}
Let \(L\) be a field of characteristic zero with a nontrivial
involution, and let \(W\) be a finite dimensional nonsingular
hermitian space over \(L\). Every isometry from a subspace
\(A\subset W\) onto a subspace \(A'\subset W\) extends to an
isometry of \(W\), including when the restricted forms are
degenerate.
\end{lemma}
\begin{proof}
Write \(h\) for the form, linear in its second entry. First
suppose the restriction to \(A\) is nonsingular. Any nonzero
hermitian space with nonzero form contains a vector of nonzero
norm. Indeed if every norm vanished, expanding the norm of
\(v+t w\) would give \(\operatorname{Tr}_{L/L^+}(t h(v,w))=0\)
for all \(t\), where \(L^+\) is the fixed field. A nonzero
\(h(v,w)\) would permit the trace of one, which is two.
Splitting off each nonsingular line gives an orthogonal basis of
\(A\). Match the basis vectors one at a time in the nonsingular
orthogonal complement of the lines already matched.

Here is the required operation in that complement. If \(v,w\)
have the same nonzero norm, set \(z=v-w\) and \(a=h(z,v)\).
For \(a\ne0\), the transformation
\[
                 L_z(x)=x-z\,h(z,x)/a
\]
takes \(v\) to \(w\). Expanding \(h(L_zx,L_zy)\) leaves the
coefficient \(-1/a-1/\bar a+h(z,z)/(a\bar a)=0\) of
\(h(x,z)h(z,y)\), since \(h(z,z)=a+\bar a\).
This is an isometry, invertible by nonsingularity and equal
finite dimensions. For \(a=0\), \(h(v+w,v)=2h(v,v)\ne0\).
The same operation now takes \(v\) to \(-w\). Negating the
whole complement then takes \(-w\) to \(w\). Extend the
operation by the identity on the already matched lines.
This proves the nonsingular case.

In general, let \(R\) be the radical of \(A\) and choose a vector
space complement \(A_0\) to \(R\) in \(A\). Its restriction is
nonsingular, since any vector orthogonal to \(A_0\) in \(A_0\)
is orthogonal to all of \(A\). In \(A_0^\perp\) take a basis
\(r_1,\ldots,r_b\) of \(R\). The functionals \(h(r_i,\cdot)\)
are independent by nonsingularity, so choose \(t_j\in A_0^\perp\)
with \(h(r_i,t_j)=\delta_{ij}\). Replace \(t_j\) by
\(t_j-\frac12\sum_i r_i h(t_i,t_j)\).
Their mutual Gram entries become zero and the pairings with
the \(r_i\) are unchanged, by expansion using \(h(R,R)=0\).
These \(2b\) vectors span a nonsingular subspace with Gram
matrix \(\left(\begin{smallmatrix}0&I\\ I&0\end{smallmatrix}\right)\).
Perform this construction for the images of \(A_0,r_i\) in
the target. Pairing the two lists extends the prescribed map
to an isometry on a nonsingular subspace. The first part of
the proof extends this isometry to \(W\).
\end{proof}
\begin{theorem}\label{thm:positive-string-rational-orbit}
Let \(d=0\) in Definition~\ref{def:two-target-gradings} and
let \(w\) survive in the sense of Definition~\ref{def:unipotent-frequency-test}.
There is a rational unitary isometry \(\delta\) of \(Y'\)
with \(\delta w e_i=T e_i\) at every positive source index \(i\).
In particular \(\delta w\) agrees with \(T\) on \(C\).
\end{theorem}
\begin{proof}
All positive indices belong to the string. Fix positive
indices \(i,j\). In Proposition~\ref{prop:eichler-frequency-coordinate}
take \(x=e_i\) and \(k=j^*\). Then \(g_k=-g_j\le-1\),
\(g_i\ge\max(1,2-g_j)\), and every target grade is zero.
All three support conditions hold. Pairing the resulting
coordinate identity with \(e_j\) gives
\(h'(w e_j,w e_i)=h'(T e_j,T e_i)\), where \(h'\) is the
form on \(Y'\).

Let \(x\) be a nonzero vector in the positive span with \(wx=0\).
Among its nonzero coefficients take the one of least grade, at
\(b\). The positive string grades occur with multiplicity one.
The next string index \(k=b+1\) exists because \(b\) has positive
grade. Now \(g_k=g_b-2\), \(x\) is supported in grades at least
\(g_b\ge1\), and \(wx\) has empty support. The last support
condition again holds because \(d=0\). The proposition gives
\((T^\dagger T x)_k=0\). Lemma~\ref{lem:two-string-moments}
identifies this coordinate with \(\sigma_s(b)x_b\ne0\).
Thus \(w\) is injective on the positive span. So is \(T\),
whose positive columns are nonzero multiples of distinct
target string vectors by Definition~\ref{def:row-string-lowering-map}.
For source length zero or one there are no positive columns.

The map \(we_i\mapsto Te_i\) consequently defines an isometry
between the two image subspaces, by the Gram equalities.
Apply Lemma~\ref{lem:rational-hermitian-witt-extension} over
\(E\) with conjugation. Every prescribed entry in \(C\) occurs
in one of these columns. This proves the final claim.
\end{proof}

\section{The stabilizer and the negative target rows}
\begin{proposition}\label{prop:triangular-frequency-stabilizer}
Let \(B_Y,B_{Y'}\) denote the groups of unitary isometries
preserving the ordered initial string flags and acting trivially
on their successive lines, as in Definition~\ref{def:string-flag-ideals}.
The same definition applies over \(E\) or over its adeles in rational
coordinates. Suppose \(w\) agrees with \(T\) on \(C\) in
Definition~\ref{def:two-target-gradings}. Then for \(\gamma\in B_Y\)
and any target isometry \(\delta\),
\[
       (\delta w\gamma)|_C=T|_C\quad\Longleftrightarrow\quad
                                                \delta\in B_{Y'} .
\]
Both Bessel groups preserve the filtrations by decreasing grades.
For \(\gamma\in B_Y,\ \delta\in B_{Y'}\) one has
\(\delta T\gamma-T\in N_{g'}\).
\end{proposition}
\begin{proof}
The positive source indices \(i\) run from zero to
\(\lfloor s/2\rfloor-1\). There are exactly this many initial target
flag lines since \(s'=\max(s-1,0)\). On each such column the
condition on \(C\) is precisely
\[
       w e_i=t_i e'_i+\sum_{h<i}b_{hi}e'_h .
\]
In fact \(g'_i=g_i-1\), and the only target indices of grade
at least \(g_i>0\) are the earlier string vectors.
Right multiplication by \(\gamma\) preserves these triangular
columns and their nonzero leading coefficients. The first
column is fixed modulo its empty prefix precisely when \(\delta\)
fixes \(e'_0\). Inductively, once \(\delta\) preserves the earlier
flag with trivial action on the lines, the condition on the
next column says exactly
\(\delta e'_i-e'_i\in\langle e'_h:h<i\rangle\), by division by
\(t_i\ne0\). The reverse implication follows from the same
calculation. With no positive columns the target flag is empty
and the assertion holds for every target isometry.

For a Bessel isometry the filtration in positive grades is the
initial string filtration. The subspace of grades at least \(a\)
for \(a\le0\) is the orthogonal complement of the subspace of
grades strictly greater than \(-a\), by the Gram pairing.
Thus the entire grade filtration is preserved. On each negative
one-dimensional graded piece the induced action is one because
its pairing with the opposite, fixed graded line is perfect.
In grade zero the definite block is preserved on the graded
quotient: when there is a middle string vector, that vector is
fixed on the quotient and spans a nonsingular line, so its
orthogonal complement there is preserved. In particular
\(\gamma-1\) raises every string vector by at least one grade,
and maps the definite block to itself modulo positive grades.
The same assertions apply to \(\delta\). The map \(T\) lowers
grades by one, kills the definite block and takes its values in
the target string. Accordingly \(T(\gamma-1)\) and \((\delta-1)T\)
have nonnegative degree. Multiplying the former by the
filtration-preserving \(\delta\) gives the final assertion.
\end{proof}
\begin{theorem}\label{thm:unique-triangular-surviving-frequency}
Let \(d=g'\), suppose \(w\) agrees with \(T\) on \(C\), and
suppose \(w\) survives for \(d\) in
Definition~\ref{def:unipotent-frequency-test}. Then
\(w-T\in N_{g'}(E)\). Thus \(T+N_{g'}(E)\) is the unique
surviving frequency class among the classes agreeing on \(C\).
\end{theorem}
\begin{proof}
It remains to determine \(w_{ba}\) when \(g'_b<g_a\le0\).
Prove the required equality by increasing induction on the
negative target grades \(g'_b\). For this fixed \(a\), every
entry of \(w e_a\) in grades less than \(g'_b\) has already
been determined. It vanishes, because \(T e_a\) is zero or of
grade \(g_a-1\ge g'_b\).
The index \(b\) is a negative string index. Let \(i\) be the
source string index bearing the same integer as the positive
target index \(b^*\). Then \(g_i=1-g'_b\). The triangular
formula in Proposition~\ref{prop:triangular-frequency-stabilizer}
expresses \(w e_i\) as \(t_i e'_{b^*}\) plus terms of grades
strictly greater than \(-g'_b\).

Apply Proposition~\ref{prop:eichler-frequency-coordinate} to
\(x=e_i,\ k=a^*\). First \(g_i=1-g'_b\ge2-g_a=2+g_k\)
and \(g_i\ge1\). Next every grade in \(w e_i\) is at least
\(-g'_b>-g_a=g_k\ge0\). Finally every grade in \(w e_a\)
is at least \(g'_b\), so adding \(g_i\) gives at least one.
The proposition gives
\(h'(w e_a,w e_i)=h'(T e_a,T e_i)\). All components of
\(w e_i\) of grade greater than \(-g'_b\) pair to zero with
\(w e_a\) by the previously established vanishing. The sole
remaining target vector is \(t_i e'_{b^*}\), which reads the
\(b\)-coordinate of \(w e_a\) up to conjugation and a nonzero
Gram factor. Since \(T e_i=t_i e'_{b^*}\), we obtain
\(w_{ba}=T_{ba}\). This completes the induction, including
the vacuous case of no negative target grades.

To check that survival is a property of the class, take any
rational \(n\in N_{g'}\) and eligible \(u=1+R\) of
Definition~\ref{def:unipotent-frequency-test}. Eligibility
is unchanged for \(w+n\) since \(nR\) has degree at least two.
Both \(wR\) and \(nR\) are positive and \(n\in U_{g'}^\perp\)
by Proposition~\ref{prop:common-partial-frequency-sum}.
Thus \(H(n,wR)=H(n,nR)=H(wR,nR)=0\).
The identity \(H(wu,nu)=H(w,n)\), obtained by cyclicity of
the trace, gives \(H(w,nR)=0\) as well. Expanding
\(H(w+n,(w+n)u)\) now changes \(H(w,wu)\) only by rational
constants, whose character is one. This proves class invariance.
The representative \(T\) satisfies the survival identity with
equality by its definition and belongs to the triangular set.
\end{proof}

\section{A rational unipotent group for each frequency}
\begin{theorem}\label{thm:frequency-unipotent-lattice}
Fix either grading \(d\) in Definition~\ref{def:two-target-gradings}
and a rational matrix \(w\). Let \(J_w\) be the set of unitary
\(u=1+R\) with \(R\) of degree at least two and \(wR\in U_d\).
Write \(\Delta_w\) for its rational points and
\[
 \Gamma_w=\{\gamma\in B_Y(E):w(\gamma-1)\in N_d(E)\}.
\]
Then \(J_w\) is a group, contained in the source Bessel group.
Its adelic points and rational points are obtained by exponentiating
the corresponding points of the rational vector space
\[
 \mathfrak j_w=\{A:A^\dagger=-A,\quad
              \deg A\ge2,\quad wA\in U_d\}.
\]
The notation \(\deg A\ge h\) means that all nonzero entries
raise the source grade by at least \(h\).
The discrete group \(\Delta_w\) is contained in \(\Gamma_w\),
and \(\Gamma_w\) acts on \(J_w\) and \(\Delta_w\) by conjugation.
There is a left and right invariant measure on \(J_w(\mathbb A)\)
giving a compact set of representatives up to boundary and
giving a measurable fundamental domain of \(\Delta_w\) volume
one. Conjugation by \(\Gamma_w\) preserves this measure.
\end{theorem}
\begin{proof}
All operators of strictly positive degree are nilpotent, with
one common bound on their nilpotence as a matrix algebra since
the number of grades is finite. The finite power series
\(\log(1+R)=R-R^2/2+\cdots\) and \(\exp A=1+A+A^2/2!+\cdots\)
are inverse there. The scalar formal identities follow by
differentiating and comparing constant terms in characteristic
zero, and can be reduced modulo the common nilpotence bound.
For a unitary \(1+R\), the logarithm is skew adjoint since
\(\log((1+R)^{-1})=-\log(1+R)\). Conversely the exponential
of a skew adjoint matrix is unitary. If \(wR\) has positive
degree then so does \(wR^i\) for \(i\ge1\), and the same applies
to \(A\). This proves the two coordinate descriptions. The
linear equations defining \(\mathfrak j_w\) have rational
coefficients after writing the entries in a \(\Q\)-basis of \(E\).
Elimination in that basis gives the same parametrization over
\(\Q\), every completion and the full adeles, since the finitely
many coefficients have bounded rational denominators.
For \(R,S\) as above, \(R+S+RS\) has the required degree
and \(w(R+S+RS)\) is positive. The finite inverse series gives
the same conclusion for \((1+R)^{-1}-1\). Thus \(J_w\) is a
group. Raising grade on the initial string flag implies membership
in \(B_Y\) by Proposition~\ref{prop:triangular-frequency-stabilizer}.

Every member of \(B_Y\) preserves grades in the filtered sense
of that proposition. Hence \(N_d\) and \(U_d\) are stable under
right multiplication by it and its inverse. It follows that
\(\Gamma_w\) is a subgroup. If \(\gamma\in\Gamma_w\), write
\(w\gamma=w+n\), \(n\in N_d\). For \(1+R\in J_w\) one has
\[
      w\gamma R\gamma^{-1}=(wR+nR)\gamma^{-1}\in U_d,
      \qquad \deg(\gamma R\gamma^{-1})\ge2 .
\]
Apply this also to \(\gamma^{-1}\) to obtain normalization.
Rational points are preserved and \(\Delta_w\subseteq\Gamma_w\)
since \(U_d\subseteq N_d\).

Choose a \(\Q\)-basis of \(\mathfrak j_w\) in order of increasing
degree, by successively choosing representatives of bases of its
successive quotients by degree. In these coordinates of the
logarithm, multiplication adds the two coordinates of degree \(h\)
plus a polynomial in the coordinates of strictly lower degree.
Indeed every other term of \(\log(\exp A\exp B)\) contains at
least two matrix factors of degree at least two. Replacing one
argument by itself plus a matrix of degree \(h\) changes every
nonlinear term only in degrees greater than \(h\). Rational
linear projection onto the degree quotient gives the assertion.
In particular product additive measure in these coordinates is
left and right invariant by successive translations of each block,
using Lemma~\ref{lem:product-integration}. Inversion is negation
on logarithms and preserves that measure.

In each coordinate use the fundamental set
\([0,1)\times\widehat\Z\) for \(\Q\) in its adeles from
Lemma~\ref{lem:theta-positive-translation}. Successively left
multiply by rational elements supported in the current block of
logarithmic coordinates, bringing that block into its fundamental
set. Such multiplication leaves all lower blocks fixed. Conversely
two points in the product of these sets differing by rational
left multiplication agree in the first block, forcing that block
of the multiplier to vanish. Repeating proves uniqueness. The
product has compact closure and volume one, and discreteness
follows also from discreteness of each rational additive diagonal.
Conjugation by \(\gamma\) acts linearly on logarithms by
\(A\mapsto\gamma A\gamma^{-1}\), an invertible rational map
of \(\mathfrak j_w\). In any rational basis its modulus is one
on the adeles. Row additions and permutations preserve product
additive measure. Multiplication by \(a\in\Q^\times\) changes
the real measure by \(|a|\) and the finite measure by
\(\prod_p p^{-v_p(a)}=|a|^{-1}\), by successive residue indices
and the factorization in Lemma~\ref{lem:basic-arith}. Elimination
into these operations proves the assertion.
\end{proof}
\begin{corollary}\label{cor:frequency-unipotent-character}
On \(J_w(\mathbb A)\) the ratio
\(\eta_w(u)=B(Tu,T)/B(wu,w)\) is a continuous unitary character
trivial on \(\Delta_w\). It is nontrivial precisely when \(w\)
fails the test of Definition~\ref{def:unipotent-frequency-test}.
\end{corollary}
\begin{proof}
By rationality \(B(w,w)=B(T,T)=1\). Given \(u,v\in J_w(\mathbb A)\),
\(w(u-1)\) and \(w(v^{-1}-1)\) both belong to the totally isotropic
\(U_d\). Right unitary invariance of \(H\) gives
\[
 H(w,w(u-1)(v-1))=H(w(v^{-1}-1),w(u-1))=0 .
\]
Consequently \(u\mapsto B(wu,w)\) is multiplicative by expansion
of \(uv-1\). For \(T\), the identical calculation uses the positive
space for the target grading \(g'\), since \(T\) lowers its grade
by one and \(u-1\) and \(v^{-1}-1\) raise source grades by two.
Continuity and modulus one follow from the bicharacter definition.
Its value on rational pairs is one by Proposition~\ref{prop:additive}.
The last assertion is the stated test.
\end{proof}

\section{Orbit integration and cancellation}
\begin{lemma}\label{lem:countable-frequency-block-unfold}
Let \(\Gamma\) be a countable group acting freely on a measured
space \(Z\), preserving measure, with measurable fundamental
domain \(D\). Let it act on a countable set \(I\), and suppose
the measurable functions \(k_i\) on \(Z\) satisfy
\(k_{\gamma i}(\gamma z)=k_i(z)\). Suppose \(C\subset I\) and
\(S\le\Gamma\) satisfy \(\gamma i\in C\) if and only if
\(\gamma\in S\) for every \(i\in C\). Then
\[
       \int_{\Gamma\backslash Z}\sum_{i\in\Gamma C}k_i(z)\dd z
         =\int_{S\backslash Z}\sum_{i\in C}k_i(z)\dd z
\]
whenever either integral of the corresponding sum of absolute
values is finite. Those two nonnegative integrals always agree.
For \(C=\{w\}\) the subgroup is the stabilizer of \(w\).
\end{lemma}
\begin{proof}
The distinct sets \(\gamma C\) are disjoint and indexed by
\(\gamma\in\Gamma/S\): if two translates meet, their quotient
takes one member of \(C\) into \(C\) and hence lies in \(S\).
For coset representatives \(\gamma\), the sets \(\gamma^{-1}D\)
together form a fundamental domain for \(S\). By equivariance,
integration over \(D\) of the sum on \(\gamma C\) equals
integration over \(\gamma^{-1}D\) of the sum on \(C\). The sum
on \(C\) is \(S\)-invariant. Additivity for nonnegative
functions and then absolute convergence give the formula, by
Lemma~\ref{lem:integral-tools}. Independence of the domains
follows by partitioning into their intersections with group
translates as in Lemma~\ref{lem:quotient-measure}.
\end{proof}
\begin{theorem}\label{thm:normalized-unipotent-orbit-average}
Let \(G\) have a left invariant measure and let the countable
subgroup \(\Gamma\) act by left translations with measurable
fundamental domain \(D\). Let \(J\) be a measurable subgroup
normalized by \(\Gamma\) and put \(\Delta=\Gamma\cap J\).
Assume that \(J\) has a left and right invariant measure,
preserved by inversion and by conjugation under \(\Gamma\),
and a measurable left fundamental domain \(D_J\) for
\(\Delta\) of measure one. All groups and maps are Borel in
second countable locally compact groups with measures finite
on compact sets. For every measurable left \(\Gamma\)-invariant
\(k\) integrable on \(D\), let
\[
                         Qk(x)=\int_{D_J}k(jx)\dd j .
\]
The integral defining \(Qk\) exists almost everywhere on \(D\).
It gives a \(\Gamma\)-invariant integrable function up to null
sets, and
\[
                   \int_D Qk(x)\dd x=\int_D k(x)\dd x .
\]
If \(k(jx)=\eta(j)k(x)\), where \(\eta\) is a nontrivial
unitary character of \(J\) trivial on \(\Delta\), then
\(\int_D k=0\).
\end{theorem}
\begin{proof}
Put \(a={\bf1}_D\), so \(\sum_{\gamma\in\Gamma}a(\gamma y)=1\).
The subgroup \(\Delta\) is normal in \(\Gamma\). Changing
variables by conjugation shows that \(Qk(\gamma x)=Qk(x)\)
where the integral is defined: \(\gamma^{-1}D_J\gamma\) is
another left fundamental domain and \(k(jx)\) is left
\(\Delta\)-invariant in \(j\).
First take \(k\ge0\), allowing infinite integrals.
Left invariance on \(G\) and product integration give
\[
   \int_D Qk(x)\dd x=\int_G k(y)a'(y)\dd y,\qquad
   a'(y)=\int_{D_J}a(j^{-1}y)\dd j .
\]
Here and below nonnegative sums and integrals interchange by
Lemma~\ref{lem:product-integration}. Choose representatives
\(s\) for \(\Delta\backslash\Gamma\). Summing \(a'(\gamma y)\)
first over \(\gamma=\delta s,\ \delta\in\Delta\), tiles \(J\)
by the sets \(\delta^{-1}D_J\), giving
\(\sum_s\int_J a(j^{-1}s y)\dd j\). Conjugation by \(s\)
changes this to \(\sum_s\int_J a(s v^{-1}y)\dd v\).
Use the right fundamental domain \(D_J^{-1}\), which has
measure one by inversion. Write \(v=v_0\delta\).
The elements \(s\delta^{-1}\) run over all of \(\Gamma\) once,
because \(\Delta\) is normal. The last sum equals
\[
    \int_{D_J^{-1}}\sum_{\gamma\in\Gamma}
                        a(\gamma v_0^{-1}y)\dd v_0=1 .
\]
Partitioning the \(G\)-integral by the translates of \(D\)
now gives \(\int_D k\). Applying this equality first to
\(|k|\) and then to positive and negative parts of the real
and imaginary components proves the integrability and equality
claimed for complex \(k\).
For the character assertion the integral of \(\eta\) over
\(D_J\) is unchanged on replacing that domain by \(D_J j_0\).
Right invariance multiplies its value by \(\eta(j_0)\), so
choosing \(\eta(j_0)\ne1\) gives zero. Hence \(Qk=0\).
\end{proof}
\begin{corollary}\label{cor:frequency-orbit-cancellation}
Let \(B=B_Y(\mathbb A)\) and \(\Gamma=B_Y(E)\) in
Proposition~\ref{prop:triangular-frequency-stabilizer}. On
\(I=V(E)/N_d(E)\) act by \(\gamma[w]=[w\gamma^{-1}]\).
Assume \(B\) is given left Haar measure with a measurable
fundamental domain of \(\Gamma\). Suppose \(k_i\) is an
equivariant family as in Lemma~\ref{lem:countable-frequency-block-unfold}
with \(\sum_i|k_i|\) integrable on that domain. On an orbit
suppose one rational representative \(w\) satisfies
\(k_{[w]}(ux)=\eta_w(u)k_{[w]}(x)\) for \(u\in J_w(\mathbb A)\).
If \(w\) fails the survival condition, the integral of the sum
over its entire orbit is zero.
\end{corollary}
\begin{proof}
The stabilizer of \([w]\) is \(\Gamma_w\) of
Theorem~\ref{thm:frequency-unipotent-lattice}, since right
multiplication by \(\gamma\) preserves \(N_d\).
Unfold the single orbit by Lemma~\ref{lem:countable-frequency-block-unfold}.
Its integral is that of \(k_{[w]}\) on \(\Gamma_w\backslash B\)
and converges absolutely. Here
\(\Gamma_w\cap J_w(\mathbb A)=\Delta_w\): a rational matrix
satisfying the defining equations adelically satisfies them
over its field of definition. Theorem~\ref{thm:frequency-unipotent-lattice}
provides the required normalization and measure, and
Corollary~\ref{cor:frequency-unipotent-character} makes
\(\eta_w\) nontrivial. Apply
Theorem~\ref{thm:normalized-unipotent-orbit-average}.
\end{proof}

\section{A rational polarization adapted to the constant term}
\begin{definition}\label{def:rational-polarized-oscillator}
Regard the space \(V\) of Definition~\ref{def:two-target-gradings}
as a rational vector space by restriction of scalars. Denote
its positive, zero and negative degree subspaces for a fixed
\(d\) by \(U,W,U_-\). Define rational bilinear forms
\[
 b(v,z)=\operatorname{Tr}_{E/\Q}(\jmath\bar\omega H(z,v)),
       \qquad a(v,z)=b(v,z)-b(z,v).
\]
Let \({\bf e}\) be the additive character of \(\mathbb A_\Q\) with
real component \(t\mapsto\exp(-2\pi i t)\) and with finite
component \(\exp(2\pi i r)\) when the finite input is congruent
to \(r\in\Q\) modulo \(\widehat\Z\). Such \(r\) exists and is
unique modulo \(\Z\) by Lemma~\ref{lem:theta-positive-translation}.
The trace definition in Proposition~\ref{prop:additive} gives
\(B(v,z)={\bf e}(b(v,z))\) adelically.
\end{definition}
\begin{lemma}\label{lem:rational-graded-polarization}
There exist rational subspaces \(P_W,L_W\subset W\) such that
\(P=U_-\oplus P_W,\ L=U\oplus L_W\) are complementary maximal
isotropic subspaces for \(a\). The form \(a\) pairs \(P,L\)
perfectly. In these coordinates \(a(t,U)=0\) for \(t\in P\)
if and only if \(t\in P_W\).
\end{lemma}
\begin{proof}
As in Lemma~\ref{lem:theta-positive-translation},
\(a(v,z)=\operatorname{Tr}_{E/\Q}(\jmath H(z,v))\).
It is alternating and nonsingular: for nonzero \(v\) the
hermitian functional in \(z\) takes all values in \(E\), including
one with nonzero scaled trace. Proposition~\ref{prop:common-partial-frequency-sum}
pairs opposite basis indices with nonzero Gram coefficients.
Thus opposite degree pieces are perfectly paired also for \(a\),
and \(W\) is a nonsingular alternating space. In \(W\) take a
nonzero \(p\), choose \(l\) with \(a(p,l)=1\), and remove the
plane they span by orthogonal projection. Indeed the two
coefficients in that projection are obtained by pairing with
\(l,p\). Repeat on its nonsingular perpendicular until exhausted.
Taking the spans of the resulting \(p\)'s and \(l\)'s gives
\(P_W,L_W\), also when \(W=0\). Perfect pairing between \(U_-\)
and \(U\) gives the last assertion.
\end{proof}
\begin{theorem}\label{thm:rational-schroedinger-zero-section}
Use \(P,L\) of Lemma~\ref{lem:rational-graded-polarization}.
A test function on \(P(\mathbb A_\Q)\) is a finite sum of products
of a real Schwartz function and a locally constant compactly
supported finite function. Write \(v=x+l\) according to \(P,L\)
and \(z=x'+l'\). Put \(b_0(v,z)=a(x,l')\),
\(q(v)=-(b-b_0)(v,v)/2\), and
\[
  (\rho(v)f)(t)={\bf e}\big(q(v)+a(t,l)\big) f(t+x),\qquad
  \theta_P(f)=\sum_{t\in P(\Q)} f(t).
\]
Then \(\rho(v)\rho(z)=B(v,z)\rho(v+z)\).
The function \(F_f(v)=\theta_P(\rho(v)f)\) is an adelic theta
function in Definition~\ref{def:adelic-oscillator-translations}
and \(D_z F_f=F_{\rho(z)f}\). Its constant term, for every
adelic \(w\), is
\[
             I^U_w(F_f)=\sum_{t\in P_W(\Q)}(\rho(w)f)(t).
\]
In particular the right side uses only the restriction of
\(\rho(w)f\) to \(P_W(\mathbb A_\Q)\).
\end{theorem}
\begin{proof}
The alternating parts of \(b,b_0\) coincide since \(P,L\)
are complementary isotropic spaces. Hence \(b-b_0\) is
symmetric and \(q(v)+q(z)-q(v+z)=(b-b_0)(v,z)\).
In a product of two displayed translations the ordinary
translation and modulation produce the additional exponent
\(a(x,l')=b_0(v,z)\). This proves the multiplication law.

Finite support places all rational summation indices in a single
lattice in \(P(\R)\), by clearing denominators. Uniformly
over compact adelic translating sets this same conclusion
holds by adding the compact set of finite shifts to the finite
support. Polynomial decay of arbitrarily high order on the
real lattice gives locally uniform absolute convergence, as
in Theorem~\ref{thm:adelic-periodized-oscillator}. Finite
characters are locally constant. Real derivatives introduce
only polynomial factors or Schwartz derivatives, giving the
same convergence and continuity. For \(r\in V(\Q)\) its two
components in \(P,L\) are rational. Every character factor
of \(\rho(r)\) on an argument \(t\in P(\Q)\) is therefore one.
Translation of \(P(\Q)\) by \(x_r\) preserves the summation
set, so \(\theta_P(\rho(r)f)=\theta_P(f)\). The multiplication
law now gives \(F_f(v+r)=B(r,v)^{-1}F_f(v)\) and
\(D_zF_f=F_{\rho(z)f}\).

On \(U\) the quadratic term \(q\) is zero since \(H(U,U)=0\)
and \(U\subset L\). Restriction of \(F_{\rho(w)f}\) to \(U\)
is thus a sum of \((\rho(w)f)(t){\bf e}(a(t,u))\).
The compact set of representatives for \([U]\) from
Lemma~\ref{lem:theta-positive-translation} permits termwise
integration by the uniform bound. A character has integral
zero unless it is trivial, by translating the probability
integral. For rational \(t\) triviality means \(a(t,U)=0\),
as is seen by varying the real arguments. These are exactly
\(t\in P_W(\Q)\) by Lemma~\ref{lem:rational-graded-polarization}.
Every remaining character has integral one, proving the formula.
\end{proof}

\section{The parabolic operator in rational coordinates}
\begin{theorem}\label{thm:graded-parabolic-schroedinger}
Let \(G\) be an adelic isometry of \(b\) in
Definition~\ref{def:rational-polarized-oscillator} preserving
\(L\) in Lemma~\ref{lem:rational-graded-polarization}. Write
\[
 G(x+l)=Ax+Cx+Dl,\qquad Ax\in P,\quad Cx,Dl\in L .
\]
Put \(K_G(t,t')=a(t,CA^{-1}t')\) on \(P\). For an adelic
invertible \(A\) write \(|\det A|_{\mathbb A}\) for the product
of its real and normalized finite absolute determinants.
For the test functions of Theorem~\ref{thm:rational-schroedinger-zero-section}
set
\[
 (\mathcal W_G f)(t)=|\det A|_{\mathbb A}^{-1/2}
       {\bf e}\big(-K_G(t,t)/2\big)\, f(A^{-1}t).
\]
These operators form a unitary representation on the test
functions with their additive \(L^2\) norm, and
\(\mathcal W_G\rho(z)=\rho(Gz)\mathcal W_G\). If \(G-1\) raises
the degrees on \(V\) by at least one, it preserves \(L\) and
\[
              I^U_w(F_{\mathcal W_G f})
                    =I^U_{G^{-1}w}(F_f)
                 \qquad(w\in V(\mathbb A)).
\]
For rational \(G\) preserving \(L\) one also has
\(F_{\mathcal W_G f}(v)=F_f(G^{-1}v)\).
\end{theorem}
\begin{proof}
Preservation of \(a\) gives \(a(Ax,Dl)=a(x,l)\) and
\(a(Ax,Cx')=a(Ax',Cx)\). Thus \(K_G\) is symmetric. Matrix
multiplication for two such maps \(G_1,G_2\) gives
\(A_{12}=A_1A_2,\ C_{12}=C_1A_2+D_1C_2\). Consequently
\[
 K_{12}(t,t)=K_1(t,t)+K_2(A_1^{-1}t,A_1^{-1}t),
\]
proving the representation law. A real linear substitution
changes volume by the absolute determinant, and at finite
places the measure change follows by row elimination and the
residue-index argument in Theorem~\ref{thm:split-finite-schroedinger}.
This gives unitarity. The character factor preserves the test
space since at infinity it has polynomially bounded derivatives
and at the finite places it is locally constant. On the compact
finite support only finitely many constancy cosets are required.

With \(q\) from Theorem~\ref{thm:rational-schroedinger-zero-section}
one has \(q(Gz)-q(z)=a(Ax,Cx)/2\) for \(z=x+l\). This follows
from invariance of \(b\) and
\(b_0(Gz,Gz)=a(Ax,Cx)+a(x,l)\).
The ratio of the quadratic phases at \(t\) and \(t+Ax\)
is \({\bf e}(a(t,Cx)+a(Ax,Cx)/2)\). Combining the last two
identities with \(a(A^{-1}t,l)=a(t,Dl)\) gives the asserted
intertwining by evaluating both sides at \(t\).

When \(G-1\) raises degrees, \(G\) and its inverse preserve
\(U\), take each vector in \(W\) to itself modulo \(U\),
and hence preserve \(L=U+L_W\). The induced operator \(A\)
on \(V/L\) is unipotent, so has determinant one: in a basis
ordered by degree it is triangular with diagonal one.
For \(t\in P_W\) one has \(At=t,\ Ct\in U\). In particular
\(A^{-1}t=t\) and \(a(t,Ct)=0\). It follows that
\((\mathcal W_G f)|_{P_W}=f|_{P_W}\).
Move \(\rho(w)\) through \(\mathcal W_G\) by the intertwining
identity and use the constant term formula of
Theorem~\ref{thm:rational-schroedinger-zero-section}. Finally,
for rational \(G\), the adelic modulus is one by the rational
product calculation of Theorem~\ref{thm:frequency-unipotent-lattice},
and the quadratic phase is one at rational \(t\).
Reindexing \(P(\Q)\) by \(A\) gives \(\theta_P\mathcal W_G=\theta_P\).
Intertwining once more gives the rational assertion.
\end{proof}
\begin{corollary}\label{cor:graded-source-target-covariance}
For \(v\in J_w(\mathbb A)\) let \(G\) act on \(V\) by
\(z\mapsto z v^{-1}\). For rational \(w\),
\[
 B(Tv,T)\,I^{U_d}_w(F_{\mathcal W_G f})
                         =\eta_w(v) I^{U_d}_w(F_f).
\]
If \(d=g'\) and \(v'\) is a target unitary with \(v'-1\) of
target degree at least two, let \(G'z=v'z\). Then
\[
 I^{U_{g'}}_T(F_{\mathcal W_{G'} f})
     =\Psi\!\left(\frac{\jmath}{2}\operatorname{tr}(S_{s'}v')\right)
                 I^{U_{g'}}_T(F_f).
\]
For source \(v\), the factor \(B(Tv,T)\) in the first identity
equals \(\Psi(\frac{\jmath}{2}\operatorname{tr}(S_s v))\).
\end{corollary}
\begin{proof}
Both maps preserve \(b\) by cyclicity of the hermitian trace.
Right multiplication by \(v^{-1}-1\) raises Hom degree by
at least two, since it raises the grade of a source argument
before the Hom matrix is applied. The same degree increase
holds for left multiplication by \(v'-1\) with \(d=g'\).
Theorem~\ref{thm:graded-parabolic-schroedinger} therefore changes
the respective frequencies to \(wv\) and \(v'^{-1}T\).
Their differences from the original frequencies belong to the
positive space, in the first case by definition of \(J_w\)
and in the second since \(T\) lowers grade by one.
Lemma~\ref{lem:constant-term-adelic-substitution} now contributes
the inverses of \(B(wv,w)\) and \(B(v'^{-1}T,T)\), respectively,
using \(B(w,w)=B(T,T)=1\) by rationality.

For any shift \(S\) of Lemma~\ref{lem:two-string-moments}
and any unitary \(u=1+R\) of the corresponding space with
\(\deg R\ge2\), one has \(\operatorname{tr}(S)=0\) and
\(\operatorname{tr}(SR^l)=0\) for \(l\ge2\), by strict increase
of grade. The finite inverse series gives
\(\operatorname{tr}(Su^{-1})=-\operatorname{tr}(Su)\).
Since \(S^\dagger=S\), adjointness and cyclicity also identify
the left side with \(\overline{\operatorname{tr}(Su)}\).
For an anti-invariant value \(z\), the traces to \(\Q\) of
\(\jmath\bar\omega z\) and \(\jmath z/2\) agree, since
\(\jmath z\) is real and \(\omega+\bar\omega=1\).
Apply this with \(S=T^\dagger T=S_s\) on the source and
\(S=TT^\dagger=S_{s'}\) on the target. This proves both formulas,
with the inverse in the target factor canceling the sign.
\end{proof}

\section{Changing the rational polarization}
\begin{lemma}\label{lem:rational-partial-poisson}
With the character \({\bf e}\) of Definition~\ref{def:rational-polarized-oscillator},
give \(\mathbb A_\Q\) product additive measure with Lebesgue measure
at infinity and \(\widehat\Z\) of finite measure one.
For a test function of one adelic variable put
\(\widehat f(t)=\int_{\mathbb A_\Q}{\bf e}(t y)f(y)\dd y\).
This is a test function and
\(\sum_{r\in\Q}\widehat f(r)=\sum_{r\in\Q}f(r)\), absolutely.
Fourier transformation in one coordinate of a several-variable
test function preserves the test space and its sum over all
rational coordinates.
\end{lemma}
\begin{proof}
A locally constant compactly supported finite function is a
finite linear combination of indicators of \(u+d\widehat\Z\)
with \(u\in\Q,\ d\in\Q_{>0}\): cover its support by finitely
many constancy cosets, refine to a common open lattice, and take
rational representatives as in Lemma~\ref{lem:theta-positive-translation}.
For \(f=\phi\otimes{\bf1}_{u+d\widehat\Z}\), the finite Fourier
transform is \(d^{-1}{\bf e}_f(tu){\bf1}_{d^{-1}\widehat\Z}(t)\).
Indeed the finite volume is \(d^{-1}\) by the real and finite
modulus calculation in Theorem~\ref{thm:frequency-unipotent-lattice};
translation kills any character nontrivial on \(d\widehat\Z\),
whose dual is \(d^{-1}\widehat\Z\) by the one-coordinate
annihilator calculation of Proposition~\ref{prop:additive}.
The real transform \(\widehat\phi(t)\) is Schwartz: derivatives
insert powers of the real integration variable and repeated
integration by parts gives every required decay bound.
Thus the sum of transforms is
\[
  d^{-1}\sum_{m\in\Z} e^{2\pi i m u/d}\widehat\phi(m/d).
\]
The smooth periodic sum \(\sum_n\phi(x+dn)\) has Fourier
coefficients \(d^{-1}\widehat\phi(m/d)\) by integrating over
one real period. Its Fourier series converges to the function
by the coefficient estimate and the product kernel proof in
Theorem~\ref{thm:rational-frequency-expansion}. Evaluation at
\(u\) gives \(\sum_n\phi(u+dn)=\sum_{r\in\Q}f(r)\). For several
coordinates choose product finite constancy lattices and apply
the one-coordinate argument with the remaining variables fixed.
Partial derivatives in the remaining real variables obey the
same decay estimates, uniformly after multiplication by any
polynomial in those variables. Thus all sums over rational
coordinates converge absolutely, permitting their interchange.
\end{proof}
\begin{theorem}\label{thm:rational-theta-symplectic-action}
Let \(P,L\) be complementary rational maximal isotropic spaces
for \(a\) in Lemma~\ref{lem:rational-graded-polarization}, with
dual rational bases. For \(v=x+l\) set
\(\rho_a(v)={\bf e}(b(v,v)/2)\rho(v)\), and
\(F^a_f(v)=\theta_P(\rho_a(v)f)\). Every rational linear
automorphism \(G\) preserving \(a\) induces a unique operator
\(\mathcal T_G\) on test functions such that
\[
                F^a_{\mathcal T_G f}(v)=F^a_f(G^{-1}v).
\]
The operators are a unitary representation and intertwine
\(\rho_a(z)\) with \(\rho_a(Gz)\).
For \(G\) also preserving \(b\) the same statements hold
for \(F_f,\rho\). In particular if \(G\) preserves \(U_d\),
\(I_w^{U_d}(F_{\mathcal T_G f})=I_{G^{-1}w}^{U_d}(F_f)\).
The space of functions \(F_f\) obtained in this way is independent
of the rational polarization.
\end{theorem}
\begin{proof}
In the dual bases write \(a(x,l)=x\cdot l\). The phase in
\(\rho_a(x+l)f(t)\) is \({\bf e}(x\cdot l/2+t\cdot l)\).
For any fixed adelic \(x\), integration of
\({\bf e}(-x\cdot l/2) F^a_f(x+l)\) over \([L]\) returns \(f(x)\).
Indeed each rational summation index \(r\) contributes its
character \({\bf e}(r\cdot l)\), leaving only \(r=0\).
The absolute convergence is that of
Theorem~\ref{thm:rational-schroedinger-zero-section}.
This gives injectivity. Squaring absolute values and using
the same character orthogonality gives
\(\int_{[L]}|F^a_f(x+l)|^2\dd l=\sum_{r\in P(\Q)}|f(x+r)|^2\).
Integrate over \([P]\) using the product fundamental domains of
Lemma~\ref{lem:theta-positive-translation}. Hence the squared
norm of \(f\) equals the integral of \(|F^a_f|^2\) over
\(V(\Q)\backslash V(\mathbb A_\Q)\). Its integrand is periodic,
by the rational translation law and unit modulus.

For rational \(G\) preserving \(L\), the formula of
Theorem~\ref{thm:graded-parabolic-schroedinger} applies with
the same \(A,C,D,K_G\) and rational determinant of adelic norm
one. For \(\rho_a\) the calculation works for every such
symplectic \(G\), since the change in its constant quadratic
phase is \(a(Ax,Cx)/2\) by \(a(Ax,Dl)=a(x,l)\).
The phase in \(\mathcal W_G\) is one at rational arguments and
\(\theta_P\mathcal W_G=\theta_P\).
For the coordinate exchange \(s_i:(x_i,l_i)\mapsto(l_i,-x_i)\)
use partial Fourier transformation in coordinate \(i\).
Translation by \(x_i\) then gives multiplication by
\({\bf e}(-t_i x_i)\), while multiplication by
\({\bf e}(t_i l_i)\) gives translation by \(l_i\).
Their product together with the half phase above gives precisely
\(\rho_a(s_i v)\). Lemma~\ref{lem:rational-partial-poisson}
makes \(\theta_P\) invariant under this operator as well.
Intertwining and invariance of \(\theta_P\) give the displayed
identity for both types of operation.

These operations generate the rational symplectic group.
Indeed put \(M=G L\), \(\dim(M\cap L)=n-r\), \(n=\dim L\).
A rational change of basis on \(L\) with inverse transpose on
\(P\) can carry \(M\cap L\) to the span of the last \(n-r\)
basis vectors. Orthogonality to that intersection and the
dimension of \(M\) imply that its projection to \(P\) is the
span of the first \(r\) basis vectors. Representatives in
\(M\) of these vectors may be adjusted by \(M\cap L\) to
have \(L\)-components only in the first \(r\) positions.
Their coefficients there form a symmetric matrix by isotropy.
The shear with that matrix, extended by zero on the remaining
positions, preserves \(L\) and sends
\(s_1\cdots s_r L\) to \(M\) in the adjusted basis. Returning
to the original basis, the remaining factor preserves \(L\).
This proves generation. Compose the indicated operators.
Injectivity of \(f\mapsto F^a_f\) proves independence of the
word and the representation law, including inverses. Unitarity
follows from the squared norm identity because \(G\) preserves
quotient probability by the argument of
Lemma~\ref{lem:constant-term-adelic-substitution}.

For \(G\) preserving \(b\), the common factor
\({\bf e}(-b(v,v)/2)\) gives the assertions for \(F_f,\rho\).
Use Lemma~\ref{lem:constant-term-adelic-substitution} for the
constant term. Finally choose a rational symplectic map \(h\)
carrying dual bases of one polarization to those of a second.
Transferring \(k\) by \(k'(t')=k(h^{-1}t')\) gives
\(F^{a,\mathrm{second}}_{k'}(v)=F^{a,\mathrm{first}}_k(h^{-1}v)\)
by reindexing the rational sum. Taking \(k=\mathcal T_{h^{-1}}f\)
identifies the two functions; the reverse inclusion follows using
\(h^{-1}\). Multiplying by the same quadratic factor proves
the last claim.
\end{proof}

\section{Compact representatives and absolute Fourier bounds}
\begin{proposition}\label{prop:all-bessel-quotients-compact}
For every stage of Definition~\ref{def:actual-row-lift-spaces}
the rational Bessel subgroup is discrete in the adelic Bessel
group \(B\) of Proposition~\ref{prop:triangular-frequency-stabilizer}.
It has a measurable left fundamental domain with compact closure
and finite positive Haar measure.
\end{proposition}
\begin{proof}
By the filtration proof in Proposition~\ref{prop:triangular-frequency-stabilizer},
the induced action on grade zero restricts to a unitary operator
on \(X\) and is the identity on a middle string line if present.
On every nonzero graded line the induced action is the identity.
The map to \({\rm U}(X)\) so obtained has a section acting trivially
on the string. Its kernel \(J\) consists precisely of the unitary
matrices with \(u-1\) of degree at least one. These statements
hold rationally and adelically. Thus each element is uniquely
\(nh\) with \(n\in J,\ h\in {\rm U}(X)\).
For \(J\), use the logarithmic proof of
Theorem~\ref{thm:frequency-unipotent-lattice} with the filtration
starting at degree one and the entire space of skew adjoint
strictly grade-raising matrices. The exponential and logarithm
are finite, and every nonlinear term containing a degree \(a\)
argument now has degree at least \(a+1\). The identical blockwise
additive reduction therefore gives a rational fundamental set
with compact closure in \(J(\mathbb A)\).

The real unitary group of the definite \(X\) is compact, being
a closed bounded set of matrices in the identity Gram coordinates.
Theorem~\ref{thm:class-number-bound} gives finitely many rational
classes modulo the finite compact integral unitary group. This
latter group is compact since it is closed in the product of
integral matrix spaces, each the inverse limit of its finite
residue quotients. Compactness of the countable product also
follows by successive diagonal extraction of convergent
subsequences. Thus the real group times a finite union of
compact finite cosets meets every rational orbit on
\({\rm U}(X)(\mathbb A)\). Rank zero gives one point.
First left multiply \(nh\) by a rational definite unitary
bringing \(h\) into this compact set, conjugating \(n\)
in the process. Then left multiply by a rational member of
\(J\) to bring the new \(n\) into its compact set. The product
of the two compact sets meets every rational orbit of \(B\).

For clarity, a left Haar measure on \(B\) in coordinates \(hn\)
is the product of those on \({\rm U}(X)\) and \(J\):
left multiplication changes \(h\) by a left translation and
changes \(n\) by a left translation depending on \(h\).
The measure on \(J\) was constructed in logarithmic coordinates.
At the real unitary factors use the chart density of
Lemma~\ref{lem:quotient-measure}. At the finite group give its
integral subgroup the probability assigning each coset of a
normal principal congruence subgroup the reciprocal of its
index. The nested finite residue quotients give that probability
by uniform successive fiber choices as in
Lemma~\ref{lem:theta-positive-translation}. Sum its left
translates over cosets of the integral subgroup to get left
invariance. Continuity of the two coordinate factorizations
follows from matrix multiplication and the graded block projection.
Discreteness and finiteness of the rational matrices meeting a
given compact set follow by clearing the finite denominators
and bounding the archimedean entries as in
Lemma~\ref{lem:quotient-measure}. Enlarge the compact covering
set slightly to a relatively compact open neighborhood. In that
neighborhood take a countable cover by open sets disjoint from
their distinct rational translates, and keep successive pieces
outside the translates of previous ones as in that lemma. The
resulting measurable left fundamental domain has compact closure.
It has positive measure since countably many translates of it
cover a nonempty open set, and finite measure by compact closure.
\end{proof}
\begin{proposition}\label{prop:rational-frequency-lattice-bound}
Choose the polarization of Lemma~\ref{lem:rational-graded-polarization}
for \(d\). For a test function \(f\) in this realization,
\[
  \sum_{w\in V(E)/N_d(E)} |I^{U_d}_w(F_f)|
                   \le \sum_{t\in P(\Q)} |f(t)|.
\]
For a rational basis of \(P\), set \(n=\dim_\Q P\) and use
the maximum norm on its real coordinates. Suppose the finite
support of \(f\) lies in
\(b_1^{-1}\widehat\Z^n,\ b_1\) a positive integer. For any
integer \(k>n\) the right side is at most
\[
 C_{n,k}\, b_1^n
    \sup_{t_\infty,t_f}(1+\|t_\infty\|)^k|f(t_\infty,t_f)|,
\]
where \(C_{n,k}\) depends only on the displayed indices.
This also holds uniformly over any family with one choice
of \(b_1\) and bounded displayed supremum.
\end{proposition}
\begin{proof}
Use the rational representatives \(w\in U_-(\Q)\) for the
quotient, since \(N_d=W+U\). In the notation of
Theorem~\ref{thm:rational-schroedinger-zero-section},
\(w\in P(\Q)\) acts by pure translation: its \(L\)-component
is zero and the constant phase is one by rationality.
Consequently \(I^{U_d}_w(F_f)=\sum_{\xi\in P_W(\Q)} f(w+\xi)\).
The rational points of \(P=U_-\oplus P_W\) are partitioned
in this way, proving the first bound.
The rational indices in finite support belong to
\(b_1^{-1}\Z^n\) by rational prime factorization. In a real
box of radius \(R\ge1\) there are at most
\((2b_1 R+1)^n\le(3b_1R)^n\) such points. Sum first over the
unit box and then over shells \(2^j\le\|t_\infty\|<2^{j+1}\),
\(j\ge0\), assigning boundary points once. The geometric series
with ratio \(2^{n-k}<1\) proves the claim. For \(n=0\) the
point at zero gives the same assertion directly.
\end{proof}

\section{Unfolding one period}
\begin{theorem}\label{thm:one-rational-bessel-unfold}
Use consecutive spaces \(Y,Y'\) and put
\(B=B_Y(\mathbb A)\), \(\Gamma=B_Y(E)\),
\(G'={\rm U}(Y')(\mathbb A)\), \(\Gamma'={\rm U}(Y')(E)\),
and \(\Gamma'_B=B_{Y'}(E)\), with rational points understood
as unitary isometries over the real subfield. Use a left
fundamental domain \(D_B\) with compact closure as in
Proposition~\ref{prop:all-bessel-quotients-compact}, a left
Haar measure on \(G'\), and a measurable fundamental domain
\(D'\) for \(\Gamma'\) there. Let \(p:B\to\C\) be continuous,
left \(\Gamma\)-invariant, with
\[
          p(ur)=\Psi\!\left(\tfrac{\jmath}{2}
                  \operatorname{tr}(S_s u)\right)p(r)
                         \quad(\deg(u-1)\ge2).
\]
Let \(y:G'\to\C\) be measurable and left \(\Gamma'\)-invariant.

Let \(F_{r,h}\), for \(r\in B,h\in G'\), be theta functions
in the common space of Theorem~\ref{thm:rational-theta-symplectic-action},
jointly continuous in all variables. Suppose that rational translates
satisfy
\[
       F_{\gamma r,\delta h}(v)=F_{r,h}(\delta^{-1}v\gamma)
                           \quad(\gamma\in\Gamma,\ \delta\in\Gamma').
\]
For either grading \(d\), write \(f^d_{r,h}\) for its test
function in the corresponding polarization with translation
space \(P_d\). Under left translation of \(r\) by a unitary
\(u\) with \(\deg(u-1)\ge2\), assume this test function
transforms by \(\mathcal W_{z\mapsto zu^{-1}}\) of
Theorem~\ref{thm:graded-parabolic-schroedinger}. For \(d=g'\)
assume the same for left translation of \(h\) by a target
unitary \(v\) with \(\deg(v-1)\ge2\), using \(\mathcal W_{z\mapsto vz}\).
Suppose \(M_d(r,h)=\sum_{\xi\in P_d(\Q)}|f^d_{r,h}(\xi)|\)
is measurable and bounded on compact sets for both gradings.
Assume the single absolute bound
\[
         \int_{D'} |y(h)|\int_{D_B}|p(r)|M_0(r,h)\dd r\dd h
                                                             <\infty .
\]
Define \(p'(h)=\int_{D_B}p(r)I^{U_{g'}}_T(F_{r,h})\dd r\).
It is continuous and left \(\Gamma'_B\)-invariant and has
the target generic character
\[
      p'(vh)=\Psi\!\left(\tfrac{\jmath}{2}
                  \operatorname{tr}(S_{s'}v)\right)p'(h)
                         \quad(\deg(v-1)\ge2).
\]
The following integrals converge absolutely and agree:
\[
  \int_{D_B} p(r)\int_{D'} F_{r,h}(0)y(h)\dd h\dd r
               =\int_{\Gamma'_B\backslash G'}p'(h)y(h)\dd h .
\]
The inner integral on the left is taken for almost every \(r\)
with respect to \(|p(r)|\dd r\).
\end{theorem}
\begin{proof}
Put \(I_d=V(E)/N_d(E)\), \(C_d=\{[w]\in I_d:w|_C=T|_C\}\).
For rational source and target translations preserving \(U_d\),
Lemma~\ref{lem:constant-term-adelic-substitution} makes
\[
 k^d_{[w]}(r,h)=p(r)y(h)I^{U_d}_w(F_{r,h})
\]
equivariant under \([w]\mapsto[\delta w\gamma^{-1}]\).
This holds for all \(\delta\in\Gamma'\) when \(d=0\), and for
\(\delta\in\Gamma'_B\) when \(d=g'\), by
Proposition~\ref{prop:triangular-frequency-stabilizer}.
The constant terms are continuous in \(r,h\) before multiplication
by \(y\), by integration on the compact additive quotient and
joint continuity. Proposition~\ref{prop:rational-frequency-lattice-bound}
bounds the sum of absolute values of each family by
\(|p(r)y(h)|M_d(r,h)\). On \(J_w(\mathbb A)\) in
Theorem~\ref{thm:frequency-unipotent-lattice} the source
transformation is multiplication of \(k^d_{[w]}\) by \(\eta_w\),
by Corollary~\ref{cor:graded-source-target-covariance}.

At fixed \(h\) start with \(d=0\). The sum over \(I_0\) is
\(p(r)y(h)F_{r,h}(0)\) by Theorem~\ref{thm:rational-frequency-expansion}.
Let \(D_0=\Gamma' C_0\), the union of the target translates.
The set \(C_0\) is stable under \(\Gamma\) by
Proposition~\ref{prop:triangular-frequency-stabilizer}. The source
and target actions commute, so \(I_0\setminus D_0\) is also
stable under \(\Gamma\). Every representative there fails the test,
for a survivor admits a rational target isometry bringing it into
\(C_0\), by Theorem~\ref{thm:positive-string-rational-orbit}.
On each source orbit in the complement, apply
Corollary~\ref{cor:frequency-orbit-cancellation}. Its hypotheses
on absolute integrability hold on \(D_B\) by compact closure
and the bound just proved. Add over the countably many orbits.
Thus the integral over \(D_B\) of the sum over \(I_0\) equals
the integral of the sum over \(D_0\).

Still for \(d=0\), integrate each \(k^0_i\) over the chosen
\(D_B\) to obtain a function of \(h\). These functions are
equivariant under \(\Gamma'\) by the rational covariance with
\(\gamma=1\). For \(i\in C_0\), \(\delta i\in C_0\) holds
precisely for \(\delta\in\Gamma'_B\) by
Proposition~\ref{prop:triangular-frequency-stabilizer}. Apply
Lemma~\ref{lem:countable-frequency-block-unfold} to \(C_0\)
on \(G'\). The assumed integral of \(|yp|M_0\) bounds the
sum of absolute values on the original domain, also if the
absolute value is taken before integrating over \(D_B\).
Thus on \(\Gamma'_B\backslash G'\) the integral of the sum
over \(C_0\) is absolutely convergent and equals the original
integral. The same bound permits the initial interchange of
\(D_B,D'\) by Lemma~\ref{lem:product-integration}.

By Proposition~\ref{prop:common-partial-frequency-sum} the
sum over \(C_0\) equals the sum over \(C_{g'}\) at each \(r,h\).
This use of the two polarizations is justified by
Theorem~\ref{thm:rational-theta-symplectic-action}. In \(C_{g'}\),
the subset obtained by removing \([T]\) is stable under the
rational source group: it preserves \(C_{g'}\) and fixes
\([T]\) by Proposition~\ref{prop:triangular-frequency-stabilizer}.
Every class in that subset fails the test by
Theorem~\ref{thm:unique-triangular-surviving-frequency}.
The cancellation corollary eliminates these source orbits at
each fixed \(h\), with absolute convergence on \(D_B\) from
the bound for \(M_{g'}\). The result of the \(D_B\)-integration
is precisely \(p'(h)y(h)\), proving both the identity and the
asserted integrability on the larger target quotient.

Finally the integrand defining \(p'\) at \(T\) is rationally
invariant in the source, because its frequency changes only
by \(N_{g'}(E)\). The same calculation gives rational
\(\Gamma'_B\)-invariance on the target. Joint continuity
and compact closure of \(D_B\) give continuity of \(p'\)
by bounded convergence on compact parameter neighborhoods.
For the target unipotents of degree two or more, pull the
constant phase of Corollary~\ref{cor:graded-source-target-covariance}
through the compact integral. This gives the claimed character.
\end{proof}

\section{The cocycle on three finite lattices}
\begin{lemma}\label{lem:three-lattice-overlap}
Work at one finite place of the real subfield, taking both
branches together in the split case. Let \(L\) be the reference
lattice of Theorem~\ref{thm:finite-oscillator-lattice-transport}.
For isometries \(g,h\), put \(L_1=gL,\ L_2=ghL\) and
\(m_a=[aL:aL\cap L]\). The unitary scalar \(\xi(g,h)\)
in \(W_g W_h=\xi(g,h)W_{gh}\) satisfies
\[
  \xi(g,h)=\left(\frac{m_{gh}}{m_gm_h}\right)^{1/2}
                    \sum_{x\in D}\Phi(x),\qquad
          D=\frac{L_2\cap(L+L_1)}{L_2\cap L_1}.
\]
For \(x\) represented in its numerator choose \(r\in L_1\)
with \(x+r\in L\). Then \(\Phi(x)=B_f(x,r)\) is well defined
and \(\Phi(0)=1\).
\end{lemma}
\begin{proof}
The scalar exists and has modulus one by the uniqueness and
unitarity in Theorem~\ref{thm:finite-oscillator-lattice-transport},
applied locally by its same lattice proof. For the unit vacuum
\(v_0\), the coefficient of \(v_0\) in \(W_{gh}v_0\) is
\(m_{gh}^{-1/2}\), by its lattice sum. In \(W_gW_h v_0\)
that coefficient is
\[
 (m_gm_h)^{-1/2}
 \sum_{x\in L_2/(L_2\cap L_1)}
 \ \sum_{\substack{r\in L_1/(L_1\cap L)\\x+r\in L}} B_f(x,r).
\]
Indeed the first translations in \(W_hv_0\) are taken over to
the \(x\)'s by \(g\), and \(D_xD_r=B_f(x,r)D_{x+r}\).
The vacuum coefficient of \(D_{x+r}v_0\) is one for \(x+r\in L\)
and zero otherwise by the self-dual lattice basis in that theorem.
There is precisely one admissible \(r\)-class for each \(x\in D\).
For \(j\in L_1\cap L\), \(B_f(x,j)=1\), since
\(x+r\in L,\ r\in L_1\) and \(B_f\) is one on the square of
each of the three lattices. Thus the phase is unchanged by
adding \(j\) to \(r\). For \(k\in L_1\cap L_2\), replace
\((x,r)\) by \((x+k,r-k)\). The phase ratio is
\(B_f(x,k)^{-1}B_f(k,r-k)=1\). This proves the formula and
independence of all representatives. For \(x=0\) choose \(r=0\).
\end{proof}
\begin{theorem}\label{thm:all-finite-weil-splitting}
On the full finite adelic unitary group of the reference hermitian
space in Definition~\ref{def:finite-splitting-multiplier}, the
operators \(\mathcal V_g=\beta(J_g)W_g\) form a unitary
representation. At a nonsplit place with residue cardinality
\(q\) in the real field and \(a_g=v_w(J_g)\) at the unique place
\(w\) of \(E\), the sum of Lemma~\ref{lem:three-lattice-overlap} is
\[
                         (-q)^{\,a_g+a_h-a_{gh}} .
\]
\end{theorem}
\begin{proof}
At a nonsplit place Lemma~\ref{lem:real-inert-sign-parity}
gives an unramified quadratic extension, \(q\equiv3\bmod4\),
and odd residue characteristic. The diagonal reduction of
Theorem~\ref{thm:split-cut-weil-multiplier} gives
\(m_g=q^{2a_g}\), since the index is computed over the
quadratic residue field. Put \(t=a_g+a_h-a_{gh}\).
It is nonnegative: passing from \(L_2\) to \(L_2\cap L_1\)
costs \(m_h\), and intersecting next with \(L\) costs at
most \(m_g\) by inclusion into \(L_1/(L_1\cap L)\).
The index \([L_2:L_2\cap L]\) is thus at most \(m_gm_h\).
The overlap lemma implies \(\left|\sum_D\Phi\right|=q^t\).

The local integer rings act on \(D\). For an integral unit
\(z\), multiplication by \(z\) replaces \(r\) by \(zr\) in
the phase and scales \(H(r,x)\) by \(z\bar z\). There is an
integral unit \(z\) of norm \(-1\). The residue norm onto
the real residue group is surjective by the cyclicity proof in
Lemma~\ref{lem:finite-unitary-primitive-frame}. Choose a lift
of a suitable residue element. If its norm already equals
\(-1\) modulo \(\pi^a\), multiply by \(1+b\pi^a\) with \(b\)
in the real integer ring and \(\pi\) a real uniformizer. Modulo
\(\pi^{a+1}\) the norm correction has linear coefficient \(2b\).
Since two is a unit one can correct each successive residue digit.
Completeness gives the desired norm. This permutation sends
\(\Phi\) to its complex conjugate, so the sum is real.

Multiplication by \(i\in E=E^+(i)\) has norm one, preserves
\(\Phi\), and squares to negation. Every nonzero element of
\(D\) has orbit of length four: two is invertible on this
finite module, so only zero equals its negative. Each value of
\(\Phi\) is a root of unity, since the finite additive character
on any fixed fractional local ideal factors through a finite
quotient by Proposition~\ref{prop:additive}. It follows that
\((\sum_D\Phi-1)/4\) is an algebraic integer. Indeed it is a
sum of one root from each nonzero orbit. All these roots are
powers of a single root \(\zeta\) of finite order \(b\), and
every element of \(\Z[\zeta]\) is integral, by its multiplication
action on the quotient \(\Z[X]/(X^b-1)\) and the monic
characteristic polynomial. The rational values
\(\sum_D\Phi=\pm q^t\) therefore satisfy
\(\sum_D\Phi\equiv1\bmod4\): a rational number satisfying a
monic integral polynomial is an integer by clearing a fraction
in lowest terms. Since \(-q\equiv1\bmod4\), there is only one
such sign. Thus \(\xi(g,h)=(-1)^t\). The ideal phase at this
place is \((-1)^{a_g}\) by Theorem~\ref{thm:actual-global-splitting-character}.
It cancels the scalar in the composition exactly.

At a split pair the proof of
Theorem~\ref{thm:split-cut-weil-multiplier} gives multiplicativity
for arbitrary local isometries there by the
Schr\"odinger change of variables. Finally at almost every place
\(g,h\) are integral, act on the local vacuum by the identity,
and have unit ideal phase. Every translated adelic vacuum is
an elementary tensor with vacuum factors at almost every place.
The lattice sum and the diagonal ideal calculation factor into
the indicated local operations on these vectors. Multiplicativity
at each place consequently proves it on their span, the full
finite space of Theorem~\ref{thm:finite-oscillator-lattice-transport}.
\end{proof}

\section{The finite integral vacuum in a polarization}
\begin{proposition}\label{prop:finite-polarized-vacuum}
At the finite adeles choose a rational symplectic polarization
\(P,L\) as in Theorem~\ref{thm:rational-theta-symplectic-action}.
Let \(\Lambda\) be the reference lattice in the symplectic space,
self-dual for \({\bf e}_f(a(\cdot,\cdot))\), with \(B_f=1\)
on \(\Lambda\times\Lambda\). Denote by \(\rho\) the finite
factor of Theorem~\ref{thm:rational-schroedinger-zero-section}.
The common fixed vectors of \(\rho(\Lambda)\) among the locally
constant compactly supported functions on \(P_f\) form a line.
Each unit vector in this line gives a unique unitary Heisenberg
intertwiner from the finite lattice space of
Theorem~\ref{thm:finite-oscillator-lattice-transport} onto these
test functions, sending the reference vacuum to that vector.
If \(P,L\) are hermitian isotropic and
\(\Lambda=(\Lambda\cap P_f)\oplus(\Lambda\cap L_f)\), the
normalized indicator of \(\Lambda\cap P_f\) is such a unit vector.
\end{proposition}
\begin{proof}
Use the abbreviations \(A=\operatorname{pr}_{P_f}\Lambda\),
\(D=\Lambda\cap L_f\). These are compact open additive lattices.
The annihilator of \(A\) in \(L_f\) is \(D\): a member pairs
integrally with \(A\) if and only if it annihilates \(\Lambda\)
and hence lies in \(\Lambda\). The reciprocal annihilator is
\(A\). For detail, a compact open additive lattice in rational
coordinates lies between \(b\widehat\Z^n\) and
\(b^{-1}\widehat\Z^n\) for an integer \(b>0\). On their finite
quotient the character pairing on the dual coordinates is
perfect by the one-coordinate character in Proposition~\ref{prop:additive}
and a product over the coordinates. Every character of a
subgroup of a finite abelian group extends to that group:
adjoin one generator at a time, choosing a suitable root of
the prescribed character value on its first multiple in the
old subgroup. This proves the reciprocal annihilator assertion
by separating any nonzero coset by a cyclic character and
extending it. It also shows that every continuous character of
an additive compact open lattice with open kernel is the
restriction of \({\bf e}_f(a(v,\cdot))\) for some finite \(v\):
use a small standard lattice in its kernel and a large standard
lattice containing it, with the same finite dual calculation.

For \(l\in D\), the action is multiplication by
\({\bf e}_f(q(l)+a(t,l))\). The factor \({\bf e}_f(q(l))\)
on \(D\) is a character, since \(B_f=1\) on \(\Lambda^2\).
It is locally constant by continuity of the finite character.
Choose \(t_0\) representing its inverse by the preceding
duality. A fixed function must be supported on \(t_0+A\).
For every \(x+l\in\Lambda\), its value at \(t_0+x\) is forced
to be its value at \(t_0\) times
\({\bf e}_f(-q(x+l)-a(t_0,l))\). This prescription is
well defined: two lifts differ by \(D\), whose multiplier at
\(t_0\) is one, and \(\rho\) is an ordinary representation
on \(\Lambda\). The same representation law gives invariance
under every member of \(\Lambda\) on the resulting coset.
The resulting function is locally constant, since the phase
is continuous with open level sets on compact \(\Lambda\) and
its projection to \(A\) is open: the projection of any open
ball in the open subgroup \(\Lambda\) contains an open ball.
With value \({\rm vol}(A)^{-1/2}\) at \(t_0\) the function
has squared norm one. This proves existence and uniqueness
up to a scalar.

Translates indexed by distinct classes modulo \(\Lambda\) are
orthogonal eigenvectors for \(\Lambda\), with distinct characters
by self-duality and the commutator in
Lemma~\ref{lem:theta-positive-translation}. They span the test
space. Indeed any test function is fixed by a small open subgroup
of \(\Lambda\), by compact finite support, local constancy and
the formula for \(\rho\). More explicitly, choose a common
translation lattice for its finitely many support cosets, make
\(q\) integral on the small subgroup and make all modulations
trivial on that compact support. The action on its orbit span
then factors through a finite abelian group. Averaging against
each of its characters decomposes the function into eigenvectors:
summing the averages returns the identity by finite character
orthogonality, with a full list of characters obtained by
extension from cyclic factors as above. Every occurring character
is realized by a translate, and translating back reduces its
eigenspace to the one-dimensional fixed space. The translate
basis has the same phases under \(\rho\) as the lattice basis
under \(D_z\), by the same multiplication law. This gives the
unitary intertwiner. In the split lattice case \(q(x+l)=-b(l,x)\)
by hermitian isotropy of \(P,L\). This has trivial finite
character for integral \(x,l\) in the given summands, so
one may take \(t_0=0\) and the normalized indicator.
\end{proof}
\begin{proposition}\label{prop:hyperbolic-finite-levi-character}
Assume the final split lattice condition of
Proposition~\ref{prop:finite-polarized-vacuum} and that \(P,L\)
are \(E\)-linear. If \(G\) preserves each summand, write \(A\)
for its restriction to \(P\). In this Schr\"odinger realization
\[
        \mathcal V_G
       =\upsilon_f(\det_E A)^{-1}\mathcal W_G ,
\]
where \(\mathcal W_G\) is the finite normalized change of variables
of Theorem~\ref{thm:graded-parabolic-schroedinger} and
\(\upsilon_f\) is the finite factor in
Theorem~\ref{thm:actual-global-splitting-character}.
\end{proposition}
\begin{proof}
On \(\Lambda_P=\Lambda\cap P_f\) and
\(\Lambda_L=\Lambda\cap L_f\) the commutator gives perfect
additive lattice duality by the preceding proposition.
Locally the hermitian pairing identifies the latter lattice
with the module dual of the former, scaled by a fixed fractional
ideal. Indeed integrality of its scaled trace against every
integral multiple of every basis vector says that the corresponding
\(E\)-valued pairing lies in the scalar trace-dual ideal, by the
argument of Proposition~\ref{prop:additive}. The action on \(L\)
is consequently inverse conjugate transpose to \(A\) in dual
lattice bases. Write \(k_{w,i}\) for the diagonal valuations
of \(A_w\) on \(\Lambda_P\) obtained by integral elimination
as in Theorem~\ref{thm:split-cut-weil-multiplier}. At \(w\)
the other summand has valuations \(-k_{\bar w,i}\), up to order.
Writing \(J_-=\prod_{w,i}\mathfrak p_w^{\max(0,-k_{w,i})}\)
and \(J_+=\prod_{w,i}\mathfrak p_w^{\max(0,k_{w,i})}\), we get
\(J_G=J_-\bar J_+\) and the determinant ideal of \(A\) is
\(J_+/J_-\). Since \(\beta(\bar J)=\beta(J)^{-1}\), the
phase \(\beta(J_G)\) is \(\upsilon_f(\det_E A)^{-1}\).
The saturation index ideals are intrinsic to the lattice:
integral changes of basis and their inverses preserve their
determinant valuations by the same elimination proof.
Finally \(\mathcal W_G\) sends the integral indicator to a
positive multiple of the indicator of \(A\Lambda_P\). Its
overlap with the original vacuum is positive. Uniqueness
of the positive-overlap intertwiner in
Theorem~\ref{thm:finite-oscillator-lattice-transport} therefore
identifies it with \(W_G\). Use
Theorem~\ref{thm:all-finite-weil-splitting} to obtain the formula.
\end{proof}

\section{Real symplectic transformations of a Gaussian}
\begin{lemma}\label{lem:schroedinger-real-gaussian}
In real symplectic coordinates \((x,l)\in\R^n\oplus\R^n\)
with \(a((x,l),(y,k))=x^Tk-l^Ty\), use the real factor
\[
   \rho_a(x,l)f(t)=e^{-2\pi i(x^Tl/2+t^Tl)}f(t+x).
\]
Fourier transformation with kernel \(e^{-2\pi i t^Ty}\)
is a unitary bijection of the Schwartz functions with the
\(L^2\) norm. Two unitary bijections implementing the same
change of variables on \(\rho_a\) which agree on
\(\phi_0(t)=2^{n/4} e^{-\pi t^Tt}\) agree on all Schwartz
functions. A unitary bijection commuting with \(\rho_a\)
and preserving the Schwartz space is scalar.
\end{lemma}
\begin{proof}
The transformed function is Schwartz by differentiation and
integration by parts as in Lemma~\ref{lem:rational-partial-poisson}.
The Gaussian integral calculation in Lemma~\ref{lem:theta} gives
\(\widehat{e^{-\pi\epsilon t^Tt}}(v)
=\epsilon^{-n/2}e^{-\pi v^Tv/\epsilon}\) by taking products.
For Schwartz \(f,h\) insert \(e^{-\pi\epsilon t^Tt}\) in
\(\int\widehat f(t)\overline{\widehat h(t)}\dd t\).
Product integration changes the resulting integral to
\(\int f(x)\overline{h(y)}\epsilon^{-n/2}
e^{-\pi\|x-y\|_2^2/\epsilon}\dd x\dd y\). It tends to
\(\int f\bar h\): write \(y=x+\sqrt\epsilon u\) and use
the integrable majorant \(|f(x)|\sup|h|e^{-\pi\|u\|_2^2}\).
The original integral tends to the unweighted one by Schwartz
decay. The same insertion in the inverse transform of
\(\widehat f\) gives convolution with the same Gaussian and tends
to \(f(x)\) for every \(x\). This proves isometry and inversion,
also for a single variable with the remaining ones fixed.

For two implementers as stated their quotient \(A\) commutes
with every \(\rho_a(v)\) and fixes \(\phi_0\). For a Schwartz
function \(h\), unitarity implies that \((A-1)h\) is orthogonal
to every modulation \(e^{-2\pi i t^Tl}\phi_0(t)\). The Fourier
transform of \(((A-1)h)\phi_0\) is zero, hence this product
vanishes by inversion. Since \(\phi_0>0\), equality follows.
For the last assertion, differentiating the commutation identities
along real translations and modulations shows that the commuting
operator commutes on Schwartz functions with \(\partial_j\) and
\(t_j\): the difference quotients converge in \(L^2\) by the
Schwartz estimates and unitarity preserves their limits.
The equations \((\partial_j+2\pi t_j)A\phi_0=0\) imply
\(A\phi_0=c\phi_0\) by multiplying by \(e^{\pi t^Tt}\).
Its modulus is one since \(\|\phi_0\|_2=1\). The preceding
uniqueness now gives \(A=cI\).
\end{proof}
\begin{proposition}\label{prop:real-symplectic-gaussian-action}
Let \(\mathfrak H_n\) consist of complex symmetric \(n\)-by-\(n\)
matrices \(Z\) with positive definite imaginary part. Put
\(\phi_Z(t)=2^{n/4}e^{\pi i t^T Zt}\), so \(\phi_{iI}=\phi_0\)
of Lemma~\ref{lem:schroedinger-real-gaussian}. For each real
symplectic \(G=\left(\begin{smallmatrix}A&B\\C&D\end{smallmatrix}\right)\)
there is a unitary bijection \(M_G\) of the Schwartz functions
with \(M_G\rho_a(v)=\rho_a(Gv)M_G\). It satisfies
\[
 M_G\phi_Z=c_G(Z)\phi_{GZ},\qquad
 GZ=(C+DZ)(A+BZ)^{-1},\qquad
 c_G(Z)^2\det(A+BZ)=\epsilon_G ,
\]
where \(c_G\) is continuous and \(\epsilon_G\) is a constant
of modulus one. Every other such implementer differs by a
unit scalar.
\end{proposition}
\begin{proof}
For vectors \((t,Zt)\) in the complex graph, the alternating
form vanishes on that graph, and
\(a(\bar v,v)/(2i)=\bar t^T\operatorname{Im}(Z)t>0\)
for nonzero \(t\). A real symplectic transformation preserves
both properties. If \((A+BZ)t=0\), its image has zero first
coordinate, contradicting the strict inequality. Hence the new
graph is that of \(GZ\), with symmetry and positive imaginary
part by the same tests. Projecting the graph after two
transformations gives
\[
   \det(A_{12}+B_{12}Z)
     =\det(A_1+B_1(G_2Z))\det(A_2+B_2 Z)
\]
for \(G_{12}=G_1G_2\).

Use the generators in the rational symplectic proof of
Theorem~\ref{thm:rational-theta-symplectic-action}, now over
\(\R\). For \(B=0\) the quadratic change of variables in
Theorem~\ref{thm:graded-parabolic-schroedinger}, with real
measure and the symmetric phase, is unitary, intertwines the
translations and gives \(c_G=|\det A|^{-1/2}\).
For the exchange of \(x_j,l_j\), use the partial Fourier
transform of Lemma~\ref{lem:schroedinger-real-gaussian}.
Completing its scalar Gaussian integral gives
\(c_G(Z)=(-iZ_{jj})^{-1/2}\), with square root positive at
\(Z_{jj}=i\). To verify the integral, the Fourier transform in
Lemma~\ref{lem:theta} gives
\[
 \int_\R e^{-\pi u y^2+2\pi i z y}\dd y
                   =u^{-1/2}e^{-\pi z^2/u}
                \qquad(\operatorname{Re}u>0,\ z\in\C).
\]
It holds first for positive real \(u\) and real \(z\), then for
complex \(z\) and then for \(\operatorname{Re}u>0\) by the identity
principle of Lemma~\ref{lem:analytic-tools}. Gaussian domination
on compact parameter sets justifies holomorphy by differentiation
under the integral. Here \(u=-iZ_{jj}\) and \(z\) is the sum
of the cross terms minus the Fourier variable. The resulting
quadratic exponent is the graph transform: with \(v_j\) the old
coordinate, the new one is \((Zv)_j\), so substituting
\(v_j=((Zv)_j-\sum_{k\ne j}Z_{jk}v_k)/Z_{jj}\) gives its new
dual coordinate \(-v_j\) and the unchanged dual coordinates
\((Zv)_k,\ k\ne j\). The integrated exponent has precisely
these linear derivatives after division by \(2\pi i\).
For this exchange \(\det(A+BZ)=Z_{jj}\), so the constant
\(\epsilon_G\) equals \(i\). For \(B=0\) it is
\(\operatorname{sgn}(\det A)\). Composition using the graph
determinant identity proves the assertion for any word. Scalar
uniqueness is Lemma~\ref{lem:schroedinger-real-gaussian}.
\end{proof}

\section{The archimedean Weil representation of a unitary group}
\begin{theorem}\label{thm:negative-block-real-splitting}
At a complex place write the hermitian form in coordinates
\(z=(z_+,z_-)\) with signs \(+1,-1\) of multiplicities \(p,q\),
and \(n=p+q\). Put \(j_0=\operatorname{Im}(\tau\jmath)>0\) for
the chosen CM embedding, and use real Darboux coordinates
\[
             x+i l=\sqrt{2j_0}\,(z_+,\bar z_-).
\]
Then the trace alternating form in
Definition~\ref{def:rational-polarized-oscillator} is
\(x^T l'-l^T x'\). For the unitary \(g=
\left(\begin{smallmatrix}\alpha&\beta\\\gamma&\delta\end{smallmatrix}\right)\)
in these signed complex coordinates, there is a unique unitary
Schwartz operator \(\mathcal V_g^\infty\) implementing \(g\) on
the real Heisenberg translations and satisfying
\[
            \langle\phi_0,\mathcal V_g^\infty\phi_0\rangle
                                             =(\det_\C\delta)^{-1}.
\]
The inner product is conjugate linear in the first entry, and an
empty determinant is one. The operators form a representation
continuous on Schwartz functions in their Schwartz seminorms.
In the symmetric gauge of Lemma~\ref{lem:schroedinger-real-gaussian}
the coefficient of a translation on the same vacuum is
\(\exp(-\pi j_0\|z\|_2^2)\).
\end{theorem}
\begin{proof}
The trace of \(\jmath H(z',z)\) at this place is
\(2j_0\operatorname{Im}H(z,z')\). Splitting into the two sets
of signs gives the stated real alternating form. The element
\(\jmath=-1/\sqrt{-d}\) in Lemma~\ref{lem:actual-lift-scale}
has positive imaginary part at the embeddings with
\(\sqrt{-d}=i\sqrt d\). In the symmetric gauge, direct Gaussian
integration with \(t\) shifted by \(x/2\), as in Lemma~\ref{lem:theta},
gives \(e^{-\pi(\|x\|_2^2+\|l\|_2^2)/2}\) for the translation
coefficient. Returning to \(\rho\) multiplies each translation
by \({\bf e}_\infty(-b(v,v)/2)\); \(g\) preserves \(b\), so both
gauges have the same implementing operators.

Write \(A,B,C,D\) for the real symplectic blocks of \(g\).
For \(Z\in\mathfrak H_n\) of Proposition~\ref{prop:real-symplectic-gaussian-action}
set \(H_0(Z)=\langle\phi_0,\phi_Z\rangle\). Then
\[
 H_0(Z)^2=\det\big((I-iZ)/2\big)^{-1}.
\]
In fact this is the integral of a Gaussian with complex symmetric
matrix \(I-iZ\) of positive real part. Make that real part the
identity by its positive real square root, and orthogonally
diagonalize the remaining real symmetric imaginary part, by the
orthogonal eigenbasis argument in Lemma~\ref{lem:common-hecke-interpolation}
(real eigenvectors are obtained by taking real and imaginary parts).
The one-dimensional integrals in
Proposition~\ref{prop:real-symplectic-gaussian-action} give the
squared determinant formula with the indicated power of two.
This calculation also gives nonvanishing and continuity of \(H_0\).

The action on \(\zeta=x+i l\) is \(P\zeta+Q\bar\zeta\), where
\(P=\operatorname{diag}(\alpha,\bar\delta)\) and
\(Q=\left(\begin{smallmatrix}0&\beta\\\bar\gamma&0\end{smallmatrix}\right)\).
Complexify and evaluate at \(l=i x\). Adding and subtracting the
two resulting complex coordinates gives
\[
 A+iB=\bar P+Q,\qquad C+iD=i(\bar P-Q),\qquad
             (I-i(g\,iI))/2=\bar P(A+iB)^{-1}.
\]
Both \(\alpha,\delta\) are invertible: the unitary equations give
\(\alpha^*\alpha=I+\gamma^*\gamma\),
\(\delta^*\delta=I+\beta^*\beta\).
The negative block of \(g^{-1}\) is \(\delta^*\), hence block
elimination gives
\(\det_\C g=\det\alpha/\overline{\det\delta}\), of modulus one.
Consequently
\[
 j_g(iI)=\big(\det\delta\,H_0(g\,iI)\big)^{-1}
   \quad\text{satisfies}\quad
 j_g(iI)^2=\frac{(\det_\C g)^{-1}}{\det(A+iB)}.
\]
Extend this value continuously to \(Z\in\mathfrak H_n\) by
multiplying it by
\[
 \exp\!\left[-\frac12\int_0^1
   \operatorname{tr}\{(A+BZ_t)^{-1}B(Z-iI)\}\dd t\right],
             \qquad Z_t=iI+t(Z-iI).
\]
The graph proof in Proposition~\ref{prop:real-symplectic-gaussian-action}
gives invertibility on this segment. Differentiating the determinant
by multilinearity of its columns shows that its logarithmic
derivative is the displayed trace. Thus \(j_g(Z)^2\det(A+BZ)\)
is constantly \((\det_\C g)^{-1}\). The extension is jointly
continuous in \(g,Z\). Its composition quotient
\(j_{g_1}(g_2 Z)j_{g_2}(Z)/j_{g_1g_2}(Z)\) has square one by
the determinant multiplication in that proposition.

The group \({\rm U}(p,q)\) is connected. Indeed, put
\(J=\operatorname{diag}(I_p,-I_q)\), \(H=g^*g\), so
\(JHJ=H^{-1}\). By diagonalizing the positive hermitian matrix
\(H\), its real powers satisfy \(JH^tJ=H^{-t}\). Thus
\(gH^{-t/2}\), \(0\le t\le1\), is a path in the group to an
ordinary unitary member. The latter preserves \(J\) and commutes
with it, so belongs to \({\rm U}(p)\times{\rm U}(q)\).
In each factor an orthonormal eigenbasis from
Lemma~\ref{lem:common-hecke-interpolation} reduces a path to
the identity to paths of individual phases on the circle.
Since \(\mathfrak H_n\) is convex, the continuous composition
quotient equals its value one at the identity pair and \(iI\).

In Proposition~\ref{prop:real-symplectic-gaussian-action}
the ratio \(j_g(Z)/c_g(Z)\) has constant square of modulus one.
Continuity on the connected half space makes it a single scalar.
Multiply the implementer there by this scalar. It acts on all
\(\phi_Z\) by \(j_g(Z)\phi_{gZ}\) and has the specified vacuum
coefficient by the definition at \(iI\). The composition formula
for \(j\) and Lemma~\ref{lem:schroedinger-real-gaussian} now
give the representation law. Scalar uniqueness of implementers
and the nonzero coefficient give uniqueness.

Near the identity the real block \(A\) is invertible. The
symplectic operator factors as a lower shear with coefficient
\(CA^{-1}\), the diagonal transformation \(A,A^{-T}\), and
an upper shear with coefficient \(A^{-1}B\), in this order.
Multiplication verifies the factorization from \(G^TaG=a\),
and both shear coefficients are symmetric. The upper shear
is a Fourier conjugate of a lower shear. The corresponding
Schwartz operators depend continuously on the matrix parameters:
each derivative of a quadratic phase is a polynomial times
that phase, and locally the linear substitutions and their
inverses have bounded norms. Split into a compact real box
and its complement and use an extra Schwartz power there for
convergence in each seminorm. The Fourier transform is continuous
in these seminorms by the integration by parts in
Lemma~\ref{lem:rational-partial-poisson}. The correcting scalar
is continuous by its nonzero Gaussian coefficient and the formula
for \(j_g(iI)\). The representation law gives continuity
everywhere. For \(n=0\) the group is trivial and the assertions
use the constant vacuum one.
\end{proof}

\section{The action on a maximal isotropic parabolic}
\begin{lemma}\label{lem:parabolic-splitting-comparison}
In a real or finite Schr\"odinger realization, let two continuous
unitary representations of a group of isometries implement
the same Heisenberg action. At finite places take continuity
on each locally constant test vector. Their ratio is a unitary
character. If the group preserves a polarization variable \(L\),
one of the two representations may be the operators
\(\mathcal W\) of Theorem~\ref{thm:graded-parabolic-schroedinger}.
The character is one on any subgroup contracted to the identity
by conjugation by members of this group.
\end{lemma}
\begin{proof}
The quotient of two implementers commutes with all translations.
It is scalar by Lemma~\ref{lem:schroedinger-real-gaussian} or by
Proposition~\ref{prop:finite-polarized-vacuum} and the vacuum-line
argument in Theorem~\ref{thm:finite-oscillator-lattice-transport}.
Modulus one follows by unitarity. The two group laws give
multiplicativity, and evaluation of one nonzero matrix
coefficient of the scalar operator gives continuity. A character
has constant value under conjugation, so contraction forces its
value to equal its value at the identity.
For the finite representation of Theorem~\ref{thm:all-finite-weil-splitting},
continuity holds because the integral subgroup fixes its vacuum.
On any particular translated vacuum a sufficiently deep integral
congruence subgroup changes the translating vector by a lattice
vector with trivial commutator and multiplication phases, by
local constancy in Lemma~\ref{lem:theta-positive-translation}.
This works simultaneously for a finite linear combination.
The same continuity for \(\mathcal W\) follows from its explicit
formula on a locally constant function of compact support.
\end{proof}
\begin{theorem}\label{thm:hyperbolic-adelic-parabolic-splitting}
Let a rational hermitian space be the sum of complementary
maximal hermitian isotropic subspaces \(P,L\), each of
\(E\)-dimension \(b\), with a reference lattice splitting
as in Proposition~\ref{prop:finite-polarized-vacuum}.
In its Schr\"odinger model take the product of the real
operators in Theorem~\ref{thm:negative-block-real-splitting}
and the finite operators in Theorem~\ref{thm:all-finite-weil-splitting}.
For every adelic unitary \(G\) preserving \(L\), let \(A\)
be its action on the quotient by \(L\). The product acts by
\[
                       \upsilon(\det_E A)^{-1}\mathcal W_G .
\]
In particular it equals \(\mathcal W_G\) for every rational
unitary preserving \(L\).
\end{theorem}
\begin{proof}
Real Schr\"odinger coordinates can be changed from the signed
coordinates to \(P,L\). More explicitly, a real symplectic
map carrying one dual basis to another transfers ordinary
coordinate functions by its restriction to the position
variables, with the positive constant making the measure change
unitary. Compose with an implementer of the inverse symplectic
map from Proposition~\ref{prop:real-symplectic-gaussian-action}
in the old coordinates, as in the last step of
Theorem~\ref{thm:rational-theta-symplectic-action}. This identifies
the two Heisenberg actions on the same physical vectors. Two such
identifications differ by a scalar by
Lemma~\ref{lem:schroedinger-real-gaussian}, so transporting the
group operators is unambiguous. At finite places use
Proposition~\ref{prop:finite-polarized-vacuum}.

On the parabolic the ratio to \(\mathcal W\) is a character
locally by Lemma~\ref{lem:parabolic-splitting-comparison}.
It is one on all the shears \(x+l\mapsto x+Cx+l\). Indeed
conjugation by the unitary acting as \(t\) on \(P\), \(t^{-1}\)
on \(L\), for \(t\) in the real base field at this place,
multiplies \(C\) by \(t^{-2}\). Take \(|t|\to\infty\).
Every parabolic element is a shear times a map preserving
both summands. At finite places Proposition~\ref{prop:hyperbolic-finite-levi-character}
therefore gives the required factor.

We record explicitly why the real comparison can be computed
using a different complementary isotropic space. A second
complement is the graph \(x+E_0 x\) in \(P+L\) with
\(c(x,y)=a(x,E_0y)\) symmetric. In the symmetric gauge
both position coordinates represent \(t\in V/L\). The
coordinate identification
\(f(t)\mapsto {\bf e}_\infty(c(t,t)/2) f(t)\) intertwines
the same translations: substituting \(l-E_0x\) for \(l\)
in the new phase subtracts \(c(t,x)+c(x,x)/2\), exactly
the change of the quadratic factor under \(t\mapsto t+x\).
It transports the canonical parabolic operators to one another.
In fact the transported and canonical implementers differ by
a scalar, and both have value \(|\det_{\R}A|^{-1/2}f(0)\)
at zero. The common scalar gauge for \(b\) has no effect on
this argument.

In signed complex coordinates the positive and negative
summands have equal dimension. Projection identifies the
isotropic \(L\) with the graph of a unitary map from the
positive to the negative half: projection to the first is
injective by negative definiteness on its kernel, and the
equality of norms gives unitarity. A unitary choice of bases
makes \(L=\{(z,z)\}\). Take the complementary
\(P_1=\{(z,-z)\}\). In the real Euclidean coordinates of
Theorem~\ref{thm:negative-block-real-splitting}, these are
orthogonal real Lagrangians, and the compatible complex
structure interchanges them. Thus an orthonormal basis of
\(P_1\) and its images under that complex structure form a
Euclidean symplectic change of Darboux basis. It preserves the
vacuum Gaussian up to phase under the preceding identification.
Indeed for a Euclidean symplectic map the graph of \(iI\) is
preserved since it commutes with the compatible complex structure.
Proposition~\ref{prop:real-symplectic-gaussian-action} then
takes \(\phi_0\) to a scalar multiple of itself of norm one.
The signed unitary basis changes have the same property.

If the action \(A\) on \(P_1\) is an ordinary unitary matrix,
the signed complex matrix is \(\operatorname{diag}(A,A)\).
The change of variables \(\mathcal W\) fixes the radial Gaussian,
whereas the vacuum coefficient of the prescribed operator is
\(1/\det A\) by Theorem~\ref{thm:negative-block-real-splitting}.
Their ratio is \(1/\det A\). If \(A\) is diagonal positive
real, its action on \(L\) is \(A^{-1}\), and the negative
block in signed coordinates is \((A+A^{-1})/2\), positive
definite. The actual vacuum coefficient is positive. The
canonical \(\mathcal W\) also has positive vacuum coefficient,
since it acts by a positive rescaling of a real Gaussian.
The ratio is therefore one in this case. Every invertible
complex matrix is a product of a unitary, a positive diagonal
matrix and a unitary: diagonalize \(A^*A\) by the proof in
Lemma~\ref{lem:common-hecke-interpolation}, then divide the
orthogonal columns by their positive lengths. Multiplicativity
gives the full real ratio \((\det A/|\det A|)^{-1}\).
Together with the finite factors this is precisely
\(\upsilon(\det_E A)^{-1}\) by
Theorem~\ref{thm:actual-global-splitting-character}. For an
adelic element, all but finitely many shear or block factors
preserve their local integral vacua, so the product formulas
on elementary test tensors extend by linearity. The character
is one on rational determinants, as asserted.
\end{proof}

\section{Rational invariance by doubling}
\begin{lemma}\label{lem:global-heisenberg-scalar}
A unitary bijection of the adelic test space in
Theorem~\ref{thm:rational-schroedinger-zero-section} commuting
with all its Heisenberg translations is scalar. In particular
unitary identifications of two rational models intertwining
the same Heisenberg translations differ by a scalar.
\end{lemma}
\begin{proof}
The \(\Lambda\)-fixed subspace for a self-dual finite reference
lattice as in Proposition~\ref{prop:finite-polarized-vacuum}
consists of \(h_\infty\otimes v_\Lambda\), where \(v_\Lambda\)
spans the finite fixed line. Indeed express each test as a
finite sum with linearly independent real Schwartz coefficients.
Invariance forces each finite coefficient to be fixed.
The operator and its inverse preserve this subspace, and on its
real factor commute with every real translation and modulation
up to the common scalar gauge. Lemma~\ref{lem:schroedinger-real-gaussian}
makes their action there scalar. The finite translates span by
Proposition~\ref{prop:finite-polarized-vacuum}, giving the result
on all tests. Compare two identifications by their quotient.
\end{proof}
\begin{theorem}\label{thm:actual-adelic-rational-theta-invariance}
Let \(V\) be one of the rational hermitian oscillator spaces,
with the signed infinite frames and self-dual reference finite
lattice used in Definition~\ref{def:finite-splitting-multiplier}.
In any rational symplectic polarization let \(\Omega_V(g)\)
be the product of the real splittings of
Theorem~\ref{thm:negative-block-real-splitting} and the finite
splitting of Theorem~\ref{thm:all-finite-weil-splitting},
transported to these coordinates. This is a unitary representation
on adelic tests, and for every rational unitary isometry \(\gamma\)
one has
\[
 \Omega_V(\gamma)=\mathcal T_\gamma,\qquad
 \theta_P(\Omega_V(\gamma) f)=\theta_P(f),\qquad
 F_{\Omega_V(\gamma)f}(v)=F_f(\gamma^{-1}v).
\]
Here \(\mathcal T_\gamma\) denotes the rationally normalized
operator of Theorem~\ref{thm:rational-theta-symplectic-action}.
\end{theorem}
\begin{proof}
Transport the real factors as in
Theorem~\ref{thm:hyperbolic-adelic-parabolic-splitting}, and the
finite factor by Proposition~\ref{prop:finite-polarized-vacuum}.
A real change mixing the archimedean components can also be
made by a real symplectic implementer on their direct sum,
using Proposition~\ref{prop:real-symplectic-gaussian-action}.
Products at the real places act on Schwartz functions of
all real variables: the changes of variables, quadratic phases
and partial Fourier transforms defining each factor act on its
coordinates with the other coordinates as parameters, as in
Lemma~\ref{lem:rational-partial-poisson}. Thus the product has
the asserted representation and intertwining laws on the full
test space. The unitary identification of rational models
preserving \(F_f\) in Theorem~\ref{thm:rational-theta-symplectic-action}
intertwines the physical translations by the formula \(D_zF_f
=F_{\rho(z)f}\) and injectivity there. It is unitary by the
squared norm identity in the same theorem. The local identifications
through the fixed real and finite models also intertwine those
translations. Hence they transport the product operators in
the same way by Lemma~\ref{lem:global-heisenberg-scalar}.

First consider a hermitian space which is hyperbolic over \(E\)
and has a splitting of its reference lattice into complementary
rational maximal isotropic spaces as in
Theorem~\ref{thm:hyperbolic-adelic-parabolic-splitting}.
On its rational unitary group \(\Omega(\gamma)\mathcal T_\gamma^{-1}\)
is a scalar by Lemma~\ref{lem:global-heisenberg-scalar}, and
these scalars form a character since both actions are genuine.
On the stabilizer of the isotropic summand \(L\) this character
is one by Theorems~\ref{thm:hyperbolic-adelic-parabolic-splitting}
and~\ref{thm:graded-parabolic-schroedinger}.

Take rational dual hermitian bases of the complementary isotropic
spaces, so their cross Gram matrix is the identity.
For any rational maximal isotropic \(M=\gamma L\), a rational
block diagonal unitary can take \(M\cap L\) to the span of the
last \(b-r\) coordinates of \(L\), where \(b=\dim_E L\).
Arbitrary changes of basis on \(L\) are available by their
inverse conjugate transposes on the complement. As in the
graph argument of Theorem~\ref{thm:rational-theta-symplectic-action},
the remaining basis vectors of \(M\) can be taken to be the
first \(r\) complementary basis vectors plus columns in the
first \(r\) coordinates of \(L\). Isotropy now says that this
coefficient matrix is skew hermitian. The corresponding lower
shear, extended by zero in the remaining coordinates, stabilizes
\(L\). Thus the parabolic and exchanges of individual dual
lines generate the rational unitary group. For each such pair
and nonzero \(t\in E\) with \(\bar t=-t\), the product
\[
 \begin{pmatrix}1&t\\0&1\end{pmatrix}
 \begin{pmatrix}1&0\\-t^{-1}&1\end{pmatrix}
 \begin{pmatrix}1&t\\0&1\end{pmatrix}
       =\begin{pmatrix}0&t\\-t^{-1}&0\end{pmatrix}
\]
exchanges the two lines. The lower shear is in the parabolic.
The upper shears are conjugate to lower shears by interchanging
the dual basis vectors, a rational isometry. A character trivial
on the parabolic is trivial on all its conjugates. Consequently
the scalar character is one in this hyperbolic case.

For general \(V\) form \(V\oplus V^-\), where \(V^-\) is the
same space with negative form. At finite places the latter can
use the same underlying lattice \(\Lambda\). A frame correcting
its sign is available by a local integral scalar of norm \(-1\):
use the unit construction of Theorem~\ref{thm:all-finite-weil-splitting}
at a nonsplit place and \((1,-1)\) at a split pair.
Choose the primitive sixth root \(\xi=e^{\pi i/3}\in E\),
which exists since \(6\mid N_1\) as in Lemma~\ref{lem:real-inert-sign-parity}.
Then \(\xi-1=\xi^2\) is a unit. The two rational graphs
\[
              L=\{(x,x)\},\qquad P=\{(\xi x,x)\}
\]
are hermitian isotropic and complementary. Moreover their
intersections with \(\Lambda\oplus\Lambda\) split this
lattice at every finite place. Given any \((x,y)\) in it,
the \(P\)-parameter is \((x-y)/(\xi-1)\in\Lambda\) and
the \(L\)-parameter is \(y-(x-y)/(\xi-1)\in\Lambda\).
Thus the hyperbolic case applies.

Return in the doubled space to a product of rational polarizations
of \(V,V^-\). On pure test tensors,
\(\Omega_{V\oplus V^-}(g,1)=\Omega_V(g)\otimes I\).
At finite places the two lattice overlaps multiply and the
saturation determinant ideals multiply, by the direct sum
in Lemma~\ref{lem:three-lattice-overlap} and the definition
of \(J_g\). At each real place the product Gaussian and the
block diagonal negative determinant give the same normalized
operator by Theorem~\ref{thm:negative-block-real-splitting}.
The theta sum in the product polarization factors by absolute
convergence. Choose a test on \(V^-\) with nonzero sum, for
example a positive real Gaussian times the indicator of the
standard integral coordinates. The rational theta identity for
\((\gamma,1)\) on the double then gives
\(\theta_P\Omega_V(\gamma)=\theta_P\), after cancellation of
this nonzero sum. Intertwining yields
\(F_{\Omega_V(\gamma)f}(v)=F_f(\gamma^{-1}v)\).
The uniqueness in Theorem~\ref{thm:rational-theta-symplectic-action}
completes the proof.
\end{proof}

\section{Operator laws and size for the actual arrow}
\begin{theorem}\label{thm:global-actual-arrow-covariance}
For successive \(Y,Y'\) put \(G(g,h)z=hzg^{-1}\) in rational
coordinates, and write \(n,n'\) for their \(E\)-dimensions.
With the characters of Theorem~\ref{thm:actual-twisted-finite-arrow}
use the full adelic operator
\[
       \Omega_l(g,h)=
        \frac{\mu_l^{[n]}(\det h)}{\mu_l^{[n']}(\det g)}
                             \Omega_V(G(g,h)).
\]
For a fixed test \(f\), \(F_{g,h}=F_{\Omega_l(g,h)f}\)
satisfies the rational covariance of
Theorem~\ref{thm:one-rational-bessel-unfold}. On both polarizations
it also satisfies all the unipotent transformation laws in
that theorem. Its test functions in either polarization are
Schwartz-continuous at infinity with locally constant finite
vectors, apart from the displayed character factor.
\end{theorem}
\begin{proof}
The determinant ratio is one on rational unitary points by
Theorem~\ref{thm:actual-twisted-finite-arrow} and on unipotents
since their determinants are one. The action on the Hom space
is a hermitian isometry by cyclicity of the trace.
Theorem~\ref{thm:actual-adelic-rational-theta-invariance}
therefore gives rational covariance.

For either \(d\), each source or target unipotent used in
Theorem~\ref{thm:one-rational-bessel-unfold} acts on the Hom
space by a map \(G\) for which \(G-1\) raises degree by at
least one, by the degree calculation of
Corollary~\ref{cor:graded-source-target-covariance}. It preserves
the real-symplectic polarization variable \(L=U+L_W\).
For a scalar \(t\) in the multiplicative group of the local real
base field, act on each Hom degree \(j\) by \(t^j\).
This preserves the hermitian form, since only opposite degrees
pair. It also preserves \(L\) and acts as the identity on the
degree zero space. Conjugation scales a strictly grade-raising
matrix entry by a positive power of \(t\), so \(|t|\to0\)
contracts \(G\) to the identity. At any rational finite prime
one may take \(t\) to be powers of that rational prime in all
the local factors, and at the combined real places take one
positive real \(t\) tending to zero. This argument also works
for a rational polarization of the degree zero part mixing
those factors, since the dilation is the identity there.
Apply Lemma~\ref{lem:parabolic-splitting-comparison} to the
unitary subgroup preserving \(L\), comparing the actual and
canonical operators. It gives equality for the strictly
grade-raising isometries and hence the desired transformations.
At finite adeles use finite-support truncations of \(G\), equal
to it at a growing finite set of rational primes and identity
elsewhere. They converge in the restricted product, as \(G\)
is integral at almost every place. The continuity in that
lemma gives the full identity. The real and finite continuity
assertions follow from Theorems~\ref{thm:negative-block-real-splitting}
and~\ref{thm:actual-adelic-rational-theta-invariance}.
\end{proof}
\begin{proposition}\label{prop:adelic-schwartz-frequency-height}
Fix dual rational coordinates on the underlying symplectic
space and put \(n_0=\dim_\Q P\). For an adelic unitary \(G\)
in these coordinates let \(b(G)\) be the least positive integer
clearing all finite denominators in the matrices \(G_f,G_f^{-1}\)
over \(\widehat\Z\). Write
\[
        H(G)=b(G)(1+\|G_\infty\|+\|G_\infty^{-1}\|),
\]
using Euclidean operator norms. For a fixed test \(f\), there
are \(C,a>0\), independent of \(G\), such that
\[
           \sum_{\xi\in P(\Q)}|(\Omega_V(G)f)(\xi)|
                                             \le C H(G)^a.
\]
The same bound holds with a unit character multiplying the
operator. Rational changes of basis or multiplication of
\(G\) on either side by a member of a fixed compact set change
the size by at most a constant factor.
\end{proposition}
\begin{proof}
It suffices to treat an elementary real-times-finite test and
sum the resulting bounds. For its finite factor \(v\) choose
a positive integer \(b_0\) such that
\(\rho(b_0\widehat\Z^{2n_0})v=v\), meaning invariance under
every element of that lattice, by Proposition~\ref{prop:finite-polarized-vacuum}.
Enlarge \(b_0\) to clear the rational denominators in \(q\), so
that its finite character is one there. Put \(b=b_0 b(G)\).
Since \(G_f^{-1}b\widehat\Z^{2n_0}\subseteq b_0\widehat\Z^{2n_0}\),
the moved vector \(v'\) is invariant under this smaller lattice
of translations and modulations. Its finite support lies in
\(b^{-1}\widehat\Z^{n_0}\) by testing pure modulations, and it
is constant on the cosets of \(b\widehat\Z^{n_0}\) by the pure
translations. Each such coset has volume \(b^{-n_0}\) in the
standard measure. Hence
\(\sup|v'|\le b^{n_0/2}\|v\|_2\), using unitarity.

At infinity put \(R=1+\|G_\infty^{-1}\|\). Infinitesimal
translations in the symmetric gauge act by \(\partial_i\)
and \(-2\pi i t_i\). Intertwining and unitarity give, for a
word of length \(h\) in these operators applied to the moved
real test, an \(L^2\) bound \(C_{h,f}R^h\). Indeed commute each
generator across the group operator by replacing its direction
by \(G_\infty^{-1}\) times that direction. Expansion gives at
most \((2n_0)^h\) words in the original generators with
coefficients bounded by \(R^h\). Differentiation is valid in
\(L^2\) by the difference-quotient argument in
Lemma~\ref{lem:schroedinger-real-gaussian}. The real model may
be taken in the chosen dual basis by the coordinate identification
in Theorem~\ref{thm:actual-adelic-rational-theta-invariance}.

For a real-variable Schwartz function \(v_1\), integrate
\(\partial_1\cdots\partial_{n_0}|v_1|^2\) from each coordinate
of a point to \(+\infty\). The absolute value bounds
\(|v_1(t)|^2\) by the sum over subsets \(S\) of
\(\|\partial_S v_1\|_2\|\partial_{S^c} v_1\|_2\), by the product
rule and Cauchy--Schwarz. Apply this to the moved real function
and to each of its products with \(t_i^k\), \(k=n_0+1\).
The word bounds give a uniform weighted supremum at most
\(C_f R^{k+n_0}\). Multiplying by the finite supremum and
applying Proposition~\ref{prop:rational-frequency-lattice-bound}
in its estimate for the rational lattice sum gives the result,
also for the zero-dimensional space directly.
Finite denominators in products divide the products of the
individual clearing integers, while operator norms are
submultiplicative. Both quantities are bounded on compact
sets, by finite denominators on compact subsets of the finite
adeles. Applying this in both directions proves the final claim.
\end{proof}

\section{A shortest-vector bound for rational representatives}
\begin{theorem}\label{thm:unitary-linear-shortest-representative}
Let \(Y\) be one of the positive-rank rational hermitian stages,
\(\Lambda\subset Y_f\) its fixed integral adelic lattice, and
\(\iota\) a fixed real isometry from \(Y_\infty\) to the product
of signed complex standard spaces. Write \(G={\rm U}(Y)(\mathbb A)\),
\(\Gamma={\rm U}(Y)(E)\). For \(h\in G\) let
\[
 L(h)=Y(E)\cap h_f\Lambda,\quad
 \ell_h(v)=\max_\tau\|(\iota h_\infty^{-1}v)_\tau\|_2,\quad
 \lambda(h)=\min_{0\ne v\in L(h)}\ell_h(v).
\]
In a fixed rational basis of the underlying \(\Q\)-space take
\(H(h)=b(h)(1+\|h_\infty\|+\|h_\infty^{-1}\|)\) as in
Proposition~\ref{prop:adelic-schwartz-frequency-height} and
\(H_{\rm red}(h)=\inf_{\gamma\in\Gamma}H(\gamma h)\).
For constants \(c,C>0\) depending only on the fixed space,
frame, lattice and coordinates,
\[
 c\,\max(1,\lambda(h)^{-1})\le H_{\rm red}(h)
                          \le C\,\max(1,\lambda(h)^{-1}).
\]
In the upper estimate some \(\gamma\in\Gamma\) attains the
displayed bound with \(H(\gamma h)\). In rank zero the group
is a point and both heights may be taken as one.
\end{theorem}
\begin{proof}
The image of \(L(h)\) in the archimedean coordinates defining
\(\ell_h\) is a full real lattice of fixed covolume \(v_0>0\).
Indeed Lemma~\ref{lem:class-covol} gives the global lattice
and its covolume from the finite determinant moduli. For a
unitary \(h_f\) those moduli are one, since the norm of the
hermitian determinant to the real subfield is one, also on
a split pair. The same determinant calculation applies to
\(h_\infty\), and the reference frame is fixed. The minimum
is positive and attained by discreteness. It is left
\(\Gamma\)-invariant by \(L(\gamma h)=\gamma L(h)\).
In the chosen rational coordinates \(\Lambda\) is contained
in \(b_0^{-1}\widehat\Z^{D}\) for some integer \(b_0\),
\(D=\dim_\Q Y\). Every nonzero \(v\in L(h)\) consequently
has a real coordinate of size at least \(1/(b_0b(h))\),
by rational prime factorization. Equivalence of the fixed
real norms and \(\|h_\infty^{-1}v\|\ge\|v\|/\|h_\infty\|\)
give the lower estimate after taking the infimum.

First find a compact set meeting every class with
\(\lambda(h)\ge1/2\). In ordinary Euclidean norm the moved
lattice then has mutual distances at least \(1/2\).
It has \(D\) independent vectors in a ball of fixed radius.
For detail, the average number of its points in a translated
ball of radius \(R\), as its center runs over a lattice
parallelepiped with probability measure, is
\({\rm vol}(\mathrm{ball}(R))/v_0\), by tiling. Points in a
ball contained in an affine subspace of dimension \(e<D\)
number at most \((4R+1)^e\), by taking disjoint \(e\)-dimensional
balls of radius \(1/4\) centered at them and comparing volumes
in that subspace. A ball of radius \(R\) in the ambient space
cuts the affine subspace in a ball of radius at most \(R\).
For a fixed sufficiently large \(R\) the average, of order
\(R^D\), exceeds \((4R+1)^{D-1}\). Some translate then contains
points affinely spanning the whole space. Their differences
supply independent lattice vectors of norms at most \(2R\).

The corresponding rational vectors span \(Y\) over \(E\):
a proper \(E\)-subspace would have smaller real dimension
after taking all embeddings. Select \(n=\dim_EY\) giving an
\(E\)-basis \(v_1,\ldots,v_n\). Its Gram entries lie in
\(O_E\) because \(h_f\Lambda\) is integral, and all their
complex conjugates are bounded by \(4R^2\) by Cauchy--Schwarz
in the signed coordinates. There are only finitely many such
Gram matrices by the lattice argument of Theorem~\ref{thm:class-number-bound}.
For each occurring Gram matrix choose once a rational basis
\(t_1,\ldots,t_n\) having that Gram matrix. The map sending
\(v_i\) to \(t_i\) is a rational unitary \(\gamma\).
At infinity \((\gamma h)_\infty^{-1}t_i\) have bounded
coordinates. At finite places \((\gamma h)_f^{-1}t_i\in\Lambda\).
The finite list of fixed bases and their inverses thus gives
one uniform real bound and one integer denominator bound
for \((\gamma h)^{-1}\). The same is true for \(\gamma h\)
by the inverse conjugate transpose formula with the fixed
rational Gram matrix. These bounds describe a compact set in
the unitary group: at infinity use closed bounded matrix sets
and at finite places products of fractional integer lattices,
with the unitary equations closed. This proves the thick-part
assertion.

For arbitrary \(h\), all vectors in \(L(h)\) of \(\ell_h\)-length
strictly less than one are mutually hermitian-orthogonal.
Their pairing is an algebraic integer with every conjugate
of modulus strictly less than one, hence zero by its integral
absolute norm. At each infinite place the images of these
vectors span a totally isotropic complex subspace \(I\).
With \(J\) the diagonal matrix of signs, \(I\) and \(JI\)
are Euclidean orthogonal. Let \(D_\tau\) multiply \(I\) by
two, \(JI\) by \(1/2\) and their joint Euclidean orthogonal
complement by one. It preserves the hermitian form since
\(J\) exchanges the first two spaces and preserves the last.
Both its norm and its inverse norm are at most two.
Right multiply \(h_\infty\) by the inverse of this operation,
in the coordinates \(\iota\), leaving \(h_f\) fixed. In the
new coordinates every old vector of length less than one
has doubled in length, and every other vector has length at
least \(1/2\). Thus
\(\lambda_{\rm new}\ge\min(2\lambda(h),1/2)\).
Iterate until the thick part is reached. For \(\lambda(h)<1/2\)
at most \(\lceil\log_2(1/(2\lambda(h)))\rceil\) steps are
needed; their product and inverse in signed coordinates have
norm at most \(\lambda(h)^{-1}\). Choose \(\gamma\) for the
resulting thick element as above. Undoing the right multipliers
retains the same finite denominators and increases the
archimedean bound by at most a fixed constant times
\(\max(1,\lambda(h)^{-1})\). This gives the requested representative.
\end{proof}

\section{Absolute convergence for the prescribed Weil arrows}
\begin{proposition}\label{prop:unitary-reduced-height-integrability}
For the group \(G,\Gamma\) in
Theorem~\ref{thm:unitary-linear-shortest-representative} there
is a measurable left fundamental domain \(D'\) with
\(H(h)\le2H_{\rm red}(h)\) on \(D'\). Its subset
\(H_{\rm red}\le R\) has Haar measure at most
\(C(1+R)^b\) for \(R\ge1\) and fixed \(C,b\).
Call a measurable left \(\Gamma\)-invariant function \(y\)
rapidly decreasing if for every integer \(k\ge0\) there is
\(c_k\) with
\[
          |y(h)|\max(1,\lambda(h)^{-1})^k\le c_k
                                      \qquad(h\in G).
\]
For such a function \(\int_{D'}|y(h)|H(h)^a\dd h<\infty\)
for every \(a\ge0\).
\end{proposition}
\begin{proof}
The rational matrix subgroup is countable and intersects every
compact set in finitely many points by the denominator and
archimedean lattice calculation of Lemma~\ref{lem:quotient-measure}.
Its left action is free. The finite adelic group is second
countable, using integral lattices with congruence neighborhoods,
a countable union of bounds on denominators, and finitely many
residues at any given level. As in that lemma choose a measurable
fundamental domain \(D\). The function \(H\) is Borel: finite
denominator divisibility conditions on matrices and inverses
are compact-open conditions in each bounded region, and the real
norm is continuous. Its countable infimum \(H_{\rm red}\) is
Borel and positive. On \(D\) select the first member \(\gamma\)
of an enumeration with \(H(\gamma h)\le2H_{\rm red}(h)\).
These Borel pieces translated by their respective \(\gamma\)
give a measurable fundamental domain \(D'\).

We bound the Haar measure of \(H(h)\le2R\) in the entire
group. At finite places let \(K\) be the stabilizer of the
standard underlying \(\widehat\Z\)-lattice of rank \(D_0=\dim_\Q Y\).
It is compact open. When \(b(h)=m\le2R\), the moved lattice lies
between \(m\widehat\Z^{D_0}\) and \(m^{-1}\widehat\Z^{D_0}\).
The quotient of these bounds is \((\Z/m^2\Z)^{D_0}\) by the
Chinese remainder formula of Lemma~\ref{lem:class-covol}.
Each subgroup of this quotient has at most \(D_0\) generators:
project onto the first coordinate, lift one generator of its
cyclic image, subtract multiples and induct on the kernel
in the remaining coordinates. Thus the subgroup count is at
most \(m^{2D_0^2}\), by counting padded generator lists.
Two elements with the same lattice differ on the right by \(K\).
Summing the count for \(m\le 2R\) bounds the finite measure
by a polynomial times the mass of \(K\). The finite Haar
measure exists by the compact-subgroup construction in
Proposition~\ref{prop:all-bessel-quotients-compact}, equally
for the finite unitary group considered here.

In real matrix coordinates the elements in question and their
inverses have norms at most \(2R\). Cover the ambient box of
such matrices by cubes of diameter at most \(\epsilon/(2R)\)
in operator norm for a fixed small \(\epsilon>0\), using a
coordinate grid. At most \(C_1(1+R)^{2D_0^2}\) cubes suffice.
If \(x,y\) occur in the same cube then
\(\|x^{-1}y-I\|\le\epsilon\), so the part of the cube in the
group is contained in a translate of a fixed compact
neighborhood of the identity. For example the closure of the
relative ball of sufficiently small radius has bounded inverse
and is compact in matrices. Each such translate has bounded
real Haar measure by the chart construction of
Lemma~\ref{lem:quotient-measure}. The product estimate gives
the sublevel bound on \(D'\).
By Theorem~\ref{thm:unitary-linear-shortest-representative},
\(|y(h)|\le C_k H_{\rm red}(h)^{-k}\). Sum over the dyadic
shells \(2^j\le H_{\rm red}<2^{j+1}\) using \(k>a+b+1\)
and \(H\le2H_{\rm red}\) on the chosen domain. This gives
the stated integral.
\end{proof}
\begin{theorem}\label{thm:rapid-actual-one-step-unfold}
Consider successive spaces \(Y,Y'\), the genuine twisted arrow
\(\Omega_l\) of Theorem~\ref{thm:global-actual-arrow-covariance},
and any fixed test \(f\) in its oscillator space. Fix a source
unitary adelic element \(x\). Let \(p(r)\) satisfy the continuity
and generic character conditions of Theorem~\ref{thm:one-rational-bessel-unfold},
and let \(y\) be a rapidly decreasing automorphic function on
the target in the sense of Proposition~\ref{prop:unitary-reduced-height-integrability}.
For \(F_{r,h}=F_{\Omega_l(rx,h)f}\) the period identity and
the target generic transformation law of
Theorem~\ref{thm:one-rational-bessel-unfold} hold with absolute
convergence. In particular this applies to every continuous
automorphic \(y\) on the definite terminal target.
\end{theorem}
\begin{proof}
Proposition~\ref{prop:all-bessel-quotients-compact} gives compact
closure of the source domain. Theorem~\ref{thm:global-actual-arrow-covariance}
gives rational covariance, unipotent covariance and continuity
of the oscillator vectors. Their theta functions are jointly
continuous by the locally uniform lattice estimates in
Theorem~\ref{thm:rational-schroedinger-zero-section}. Their
rational lattice sums of absolute values in both polarizations
are bounded on compact sets by
Proposition~\ref{prop:adelic-schwartz-frequency-height}.
They are measurable: each evaluation of the continuous
real and locally constant finite test family is continuous,
and there are countably many rational indices.

In the Hom coordinates the matrix for \(z\mapsto h z(rx)^{-1}\)
and its inverse have size \(H\) at most a fixed multiple of
\(H(h)\) on \(D_B\), for each of the two chosen polarizations.
Indeed right multiplication by \((rx)^{-1}\) runs through a
compact set there. Left multiplication by \(h\) consists of
its action on the individual columns, in a fixed multiplicative
field basis. Passing to the chosen rational matrix basis changes
the clearing integer and the two norms by bounded factors, as
in Proposition~\ref{prop:adelic-schwartz-frequency-height}.
The same facts hold for \(h^{-1}\). The scalar character
does not affect absolute values. Hence \(|p|M_0\le C H(h)^a\)
on \(D_B\), because \(p\) is bounded on its compact closure.
Proposition~\ref{prop:unitary-reduced-height-integrability}
gives the single absolute bound required in
Theorem~\ref{thm:one-rational-bessel-unfold}, which proves the claim.

For the terminal target the Bessel flag is empty and the
group is definite. Proposition~\ref{prop:all-bessel-quotients-compact}
gives compact representatives, so \(H_{\rm red}\) and every
continuous automorphic \(y\) are bounded there. The lower bound
in Theorem~\ref{thm:unitary-linear-shortest-representative}
therefore bounds \(\max(1,\lambda^{-1})\) on the whole quotient.
This verifies rapid decrease for every such \(y\).
\end{proof}

\section{The initial character and the successive Bessel periods}
\begin{proposition}\label{prop:actual-poincare-seed-generic}
At stage zero use the Poincar\'e function of
Theorem~\ref{thm:actual-base-bessel-periodization} with the finite
right translation of Lemma~\ref{lem:initial-frame-central-translation}.
Call the resulting function \(P_0\). It is continuous, invariant
under the left rational Bessel subgroup, and, for every source
unitary \(u\) whose difference from one raises grades by at least
two,
\[
       P_0(ux)=\Psi\!\left(\frac{\jmath}{2}
                   \operatorname{tr}(S_su)\right)P_0(x).
\]
\end{proposition}
\begin{proof}
Continuity and rational invariance are in the stated base theorem
and persist under fixed right translation. For the character
calculation \(s=2\nu+1\), and the initial simple string
entries in Proposition~\ref{prop:full-phase} are \(u_{i,i+1}\)
for \(0\le i<\nu\), with \(e_\nu=z\). Write \(R=u-1\).
Between source grades differing by exactly two the inverse
series for \(u\) gives \(R^\dagger=-R\) in those entries,
since \(R^2\) raises degree by at least four. With
\(h_i=h(e_{i^*},e_i)\), the adjoint formula gives, for
\(k=s-2-i\),
\[
 R_{k,k+1}=-(h_i/h_{i+1})\overline{R_{i,i+1}}
                                 \qquad(0\le i<\nu).
\]
The shift \(S_s\) of Lemma~\ref{lem:two-string-moments}
has coefficient \(c\) on the edges \(i<\nu\), and also on
their mirrors except at \(k=\nu\), where it is \(-Nc\).
For \(i=\nu-1\) the Gram ratio is \(-1/N\), exactly
compensating this middle coefficient. Summing the trace along
the two halves therefore gives
\(\operatorname{tr}(S_su)=c(\alpha(u)-\overline{\alpha(u)})\),
\(\alpha(u)=\sum_{i<\nu}u_{i,i+1}\). When \(\nu=0\) both
sides vanish.
Taking the trace of \(\jmath c(\alpha-\bar\alpha)/2\)
is the same as taking the trace of \(\jmath c\alpha\)
because \(c\) is real and \(\bar\jmath=-\jmath\).
Consequently the displayed phase equals
\(\Psi(-\alpha(u)/(g\varpi))\) by Lemma~\ref{lem:actual-lift-scale}.
Its finite part is exactly the character in Proposition~\ref{prop:full-phase}.
At each infinite embedding its equality to the archimedean
character there follows from the same change of frame: the
ratio of consecutive string scales gives
\(\alpha_\tau=-2\pi i\,\tau(-\alpha/(g\varpi))\) on these
entries. This identity of linear coordinate functions holds
for all archimedean entries. The base theorem supplies precisely
this product of phases.
\end{proof}
\begin{theorem}\label{thm:actual-unfolded-stage-functions}
Fix stage elements \(x_i\in{\rm U}(Y_i)(\mathbb A)\), consecutive
arrow tests \(f_i\), and arrow characters \(\Omega_{l_i}\) from
Theorem~\ref{thm:global-actual-arrow-covariance}. With \(P_0\)
as in Proposition~\ref{prop:actual-poincare-seed-generic}, define
for \(0\le i<m\)
\[
 P_{i+1}(h)=\int_{[B_i]} P_i(r x_i)\,
 I^{U_{g'}}_{T_i}\!\left(F_{\Omega_{l_i}(r x_i,h)f_i}\right)\dd r ,
 \qquad [B_i]=B_{Y_i}(E)\backslash B_{Y_i}(\mathbb A).
\]
Here \(g'\) is the grading on the indicated target stage. Each
integral converges absolutely. Each \(P_i\) is continuous,
left rationally Bessel-invariant and has the trace character
of its own shift on every unipotent of degree at least two.
It depends only on the preceding \(x_0,\ldots,x_{i-1}\).
Against any rapidly decreasing target automorphic function
the one-step identity of Theorem~\ref{thm:rapid-actual-one-step-unfold}
holds with \(p(r)=P_i(r x_i)\).

Suppose the finite factor of \(f_i\) is the prescribed inclusion
vector of Definition~\ref{def:split-inclusion-finite-oscillator},
with any real Schwartz factor. At any \(h\) with \(P_{i+1}(h)\ne0\)
there is \(r\in B_{Y_i}(\mathbb A)\) with \(P_i(r x_i)\ne0\)
at which the constant term is nonzero. For a nonempty target
string the local lattice inequality of
Proposition~\ref{prop:dual-column-one-step} holds at the finite
points of \(r x_i,h\), hence gives the same string-ideal
inequalities at \(x_i,h\).
\end{theorem}
\begin{proof}
Induct from the seed proposition. At a fixed \(x_i\), right
translation preserves the rational invariance and left character,
so \(P_i(r x_i)\) satisfies the hypotheses on \(p\).
Theorem~\ref{thm:rapid-actual-one-step-unfold} gives the target
character and invariance, by the compact integral defining its
new period, and gives the identity with every rapidly decreasing
test. Existence and continuity of the new period already follow
from the compact-domain proof of Theorem~\ref{thm:one-rational-bessel-unfold},
independently of that target test. The data in its definition
involve just the preceding stage elements, giving the stated
dependence by induction.

Nonzero integral implies the integrand is nonzero somewhere.
By Theorem~\ref{thm:all-finite-weil-splitting} the finite part
is a scalar multiple of the transported inclusion vector, and
the finite and real Heisenberg actions on the theta function
intertwine by Theorem~\ref{thm:rational-schroedinger-zero-section}.
The inclusion invariance, transported as in
Corollary~\ref{cor:moved-inclusion-finite-frequency},
therefore supplies the period lattice in
Theorem~\ref{thm:constant-term-implies-lattice-map}.
Apply that theorem and Proposition~\ref{prop:dual-column-one-step}.
Finally a Bessel element preserves every string prefix in
Definition~\ref{def:full-flag-line-ideals} and is the identity
on the string line of each quotient. For the negative prefixes
this follows by unitary duality exactly as in
Proposition~\ref{prop:triangular-frequency-stabilizer}, and the
middle line is part of the defining flag. It leaves each
projected string ideal unchanged, proving the last clause.
\end{proof}

\section{The adjoint identification of two adelic arrows}
\begin{theorem}\label{thm:global-two-arrow-adjoint}
Use \(V_1,V_2,V_*,L_1,L_2,L_*,A(a,b)=(a^\dagger,b)\) of
Theorem~\ref{thm:paired-finite-adjoint-model}. There is a fixed
unitary coordinate identification \(\mathcal J\) on the adelic
Schr\"odinger test spaces, with that theorem's identification
on the finite factors, such that on pure tensors
\[
 \begin{split}
 \theta_*\mathcal J(f_1\otimes f_2)&=\theta_1(f_1)\theta_2(f_2),\\
 \mathcal J(\rho_1(a)\otimes\rho_2(b))
       &=\lambda(a)\rho_*(A(a,b))\mathcal J,\qquad
           \lambda(a)=\Psi(-\bar\omega H_1(a,a)).
 \end{split}
\]
For \(g,h,k\) on the three consecutive unitary spaces of ranks
\(n,n',n''\), put \(G_1 a=ha g^{-1}\), \(G_2 b=kb h^{-1}\),
and \(G_*z=(g\oplus k)z h^{-1}\) on the combined space. For
the untwisted unitary splittings of
Theorem~\ref{thm:actual-adelic-rational-theta-invariance},
\[
 \mathcal J(\Omega_1(G_1)\otimes\Omega_2(G_2))
       =\upsilon_T(\det_E G_1)^{-1}\Omega_*(G_*)\mathcal J .
\]
These identities hold on all pure tests, including after their
independent translations by the stage elements.
\end{theorem}
\begin{proof}
The two-arrow computation of Theorem~\ref{thm:paired-finite-adjoint-model}
holds for the rational trace forms before completion. On the
first block adjunction replaces its hermitian value \(z\) by
\(-\bar z\). Thus \(A\) preserves the alternating form and
is a rational symplectic isometry. In the symmetric gauge
carry a product polarization to its image under \(A\) by
changing variables. The rational theta sum is unchanged by
this reindexing, and additive measures of covolume one change
by no global factor by the rational determinant calculation in
Theorem~\ref{thm:frequency-unipotent-lattice}. If another rational
polarization is used on \(V_*\), apply
Theorem~\ref{thm:rational-theta-symplectic-action} to change to it.
The product rational sum factors on pure tensors by absolute
convergence. In the original gauge the change in
\(b(v,v)/2\) gives the indicated \(\lambda(a)\), because
\(\jmath(\omega-\bar\omega)=-1\) and \(H_1(a,a)\) is real.
This is also the direct bicharacter calculation of
Theorem~\ref{thm:paired-finite-adjoint-model}, with the full
character \(\Psi\).

The coordinate identification just constructed factors into
its real and finite operators: each rational change of variables
and each rational partial Fourier transform in
Theorem~\ref{thm:rational-theta-symplectic-action} has this
factorization. It agrees with the finite identification fixed
in the statement up to unit scalar by
Proposition~\ref{prop:finite-polarized-vacuum}, since both
implement the same translations with the same quadratic gauge.
Transfer that scalar to the real factor. This fixes the product
identification with its displayed theta normalization and the
prescribed finite factor.

For the operator formula, the second block is unchanged.
Orthogonal direct sums give products of the genuine actions
by the vacuum and determinant argument in
Theorem~\ref{thm:actual-adelic-rational-theta-invariance}.
In the first block the finite positive-overlap lattice operators
intertwine under adjunction by
Theorem~\ref{thm:paired-finite-adjoint-model}. The saturation
ideal on the image lattice is \(\bar J_{G_1}\), since the
map on that block is conjugate linear and carries the lattice
onto its image. In the original lattice, unitary duality
interchanges the elementary valuations \(k_{w,i}\) with the
negatives at \(\bar w\), exactly as in
Theorem~\ref{thm:split-cut-weil-multiplier}. At a nonsplit
place this also follows from the inverse conjugate transpose
in a unimodular reference frame. Consequently the determinant
ideal of \(G_1\) is \(\bar J_{G_1}/J_{G_1}\), since at \(w\)
the valuation difference is
\(\sum_i(\max(0,k_{w,i})-\max(0,-k_{w,i}))\).
The ratio of the two phases is therefore
\(\beta(J_{G_1})/\beta(\bar J_{G_1})
 =\upsilon_f(\det_E G_1)^{-1}\).

At an infinite place adjunction in signed orthonormal frames
is complex conjugation, transposition and row and column signs.
It exchanges the positive and negative complex axes for the
negated target form, and is a Euclidean symplectic isometry in
the Darboux coordinates of Theorem~\ref{thm:negative-block-real-splitting}.
It therefore preserves the Gaussian line, as in
Theorem~\ref{thm:hyperbolic-adelic-parabolic-splitting}.
If the positive and negative diagonal blocks of \(G_1\) are
\(\alpha,\delta\), the negative block of the transformed
isometry has determinant \(\overline{\det\alpha}\). Indeed
adjunction maps the old positive axes antiunitarily to the
new negative axes, with the fixed signs canceling in conjugation.
Both operators implement the same physical isometry after
the coordinate change. Their scalar ratio is the ratio of their
Gaussian coefficients, namely
\(\overline{\det\alpha}/\det\delta=(\det_\C G_1)^{-1}\)
by Theorem~\ref{thm:negative-block-real-splitting}. This is
the inverse infinite factor of \(\upsilon_T(\det_E G_1)\).
The finite and infinite conclusions give the claimed formula.
\end{proof}
\begin{corollary}\label{cor:paired-global-frequency-product}
For even \(j<m\) use the consecutive pair starting at \(Y_j\).
Let \(U_1,U_2\) be the positive Hom grades for the string
gradings on its two arrows and let \(U_>\) be the positive
combined Hom grades in Corollary~\ref{cor:residual-pair-positive-quotient}.
For all tests \(f_1,f_2\),
\[
 I_{T}^{U_1}(F_{f_1}) I_{T'}^{U_2}(F_{f_2})
     = I_\tau^{U_>}(F_{\mathcal J(f_1\otimes f_2)}),
                     \qquad \tau=(T^\dagger,T').
\]
\end{corollary}
\begin{proof}
Reversing a pair of indices under the adjoint reverses their
individual grades and interchanges source and target. Hence
the grade of the Hom entry is unchanged, so \(A\) sends
\(U_1\oplus U_2\) onto \(U_>\). The quadratic phase in
Theorem~\ref{thm:global-two-arrow-adjoint} is one on a
positive Hom vector since its norm is zero. It is also one
on the rational \(T\). Apply that theorem first to the
translations by \(T,T'\), then to the positive translations
in the two quotient integrals. The rational additive change
of variables preserves quotient probabilities as in
Lemma~\ref{lem:constant-term-adelic-substitution}. Product
integration gives the identity.
\end{proof}

\section{The mean on the Bessel quotient}
\begin{lemma}\label{lem:bessel-product-probability}
For a graded stage of Proposition~\ref{prop:all-bessel-quotients-compact}
let \(H={\rm U}(X)\) act trivially on its string and let \(J\)
consist of the unitaries \(n\) with \(n-1\) strictly grade-raising.
Write \([S]=S(E)\backslash S(\mathbb A)\) for any of the rational
groups \(S=H,J,B\), with \(B=J\rtimes H\). Give \(J\) the
logarithmic measure for which \([J]\) has volume one, and use
the chosen definite quotient measure on \(H\). In coordinates
\(n h\) the product is a left Haar measure on \(B\), and for
this normalization and every continuous left rationally invariant
function \(K\) on \(B(\mathbb A)\),
\[
           \int_{[B]}K(r)\dd r
              =\int_{[H]}\int_{[J]}K(nh)\dd n\dd h .
\]
\end{lemma}
\begin{proof}
The logarithmic description and the rational fundamental sets
were established in Proposition~\ref{prop:all-bessel-quotients-compact}
and Theorem~\ref{thm:frequency-unipotent-lattice}. Conjugation
by \(h\) is linear on the logarithm. In the skew adjoint
strictly positive matrices the entries from \(X\) into each
positive string line are free and their adjoint entries are
then determined. The other free entries have both indices in
the string. Indeed \(X\) has grade zero, so no entry inside \(X\)
occurs, and the adjoint exchanges the block from a negative
string line into \(X\) with the indicated positive one. Thus
conjugation acts on each free row out of \(X\) by right
multiplication by \(h^{-1}\) and fixes the string-only coordinates.
The determinant over \(E\) on each such row is \((\det_E h)^{-1}\).
Its norm to the real subfield is one by taking determinants in
the unitary equation. In particular the absolute determinant
on the underlying rational coordinates at each rational place
is one. This last determinant is the product of the conjugates
of the \(E\)-determinant: elimination into row additions, swaps
and diagonal multiplications reduces it to the definition of
the field norm as the determinant of scalar multiplication.
This argument applies to the product algebras at a split prime
as well. Conjugation by \(h\) therefore preserves the logarithmic
measure by the change of coordinates in
Theorem~\ref{thm:graded-parabolic-schroedinger}.

The group law in the given coordinates is
\((n,h)(n',h')=(n h n'h^{-1},h h')\). The invariance just
proved gives left invariance of product measure. Choose left
fundamental sets \(D_H,D_J\) as in
Proposition~\ref{prop:all-bessel-quotients-compact}. Left rational
translation first brings the \(H\)-coordinate into \(D_H\),
then brings the resulting \(J\)-coordinate into \(D_J\).
If two elements of \(D_JD_H\) are rational translates, comparing
their \(H\)-coordinates first gives the same \(h\) and trivial
rational \(H\)-multiplier. Comparing the \(J\)-coordinates then
gives the same \(n\), up to the null boundaries of the chosen sets.
The integrand is bounded on their compact closures. Product
integration by Lemma~\ref{lem:product-integration} proves the
formula. Invariance of the inner integral under rational changes
of the representative \(h\) follows by the same conjugation,
which also preserves \(J(E)\).
\end{proof}

\begin{proposition}\label{prop:middle-unipotent-additive-mean}
In a pair of consecutive arrows starting at even \(j\), write
\(s=2a+1\) for the initial string length. Use \(J'\) for the
strictly grade-raising unitary group on the middle space \(Y'\).
For \(a>0\) put \(D={\rm Hom}_E(X',E\xi)\), with \(\xi\) from
Proposition~\ref{prop:actual-pair-grade-zero}; for \(a=0\)
put \(D=0\). Define \(b:J'(\mathbb A)\to D(\mathbb A_E)\) by
\[
 b(n)x=c\,\big((n-1)x\big)_{e'_{a-1}}\,\xi\quad(x\in X',\ a>0),
                \qquad b(n)=0\quad(a=0).
\]
Then \(b\) is a continuous surjective homomorphism both adelically
and on rational points. For continuous \(\phi\) on the additive
quotient \([D]=D(\mathbb A_E)/D(E)\),
\[
                    \int_{[J']}\phi(b(n))\dd n
                                =\int_{[D]}\phi(t)\dd t
\]
with the quotient probabilities on both sides.
\end{proposition}
\begin{proof}
A product of two strictly grade-raising differences raises
degree by at least two. It has zero component from \(X'\) into
the grade-one line \(Ee'_{a-1}\). The degree-one row thus
adds under multiplication, which proves the homomorphism
identity. For any row \(\ell:X'\to E\), put \(Ax=c^{-1}
\ell(x)e'_{a-1}\) on \(X'\) and put \(A=0\) on the string.
Then \(M=A-A^\dagger\) raises grade by one, and its adjoint
term is zero on \(X'\) since \(X'\) is perpendicular to this
string line. The unitary \(n=\exp M\) belongs to \(J'\) by
the finite exponential calculation of
Theorem~\ref{thm:frequency-unipotent-lattice}. Its first-degree
row is that of \(M\), whence \(b(n)x=\ell(x)\xi\). This works
over \(E\) and over its adeles. All formulas are polynomial
in the coordinates, hence continuous. The case \(a=0\) is the
unique map onto the zero group.

For adelic \(n_0\), right translation by \(n_0\) preserves the
quotient probability on \([J']\). Indeed the logarithmic measure
is right invariant by the triangular coordinate proof in
Theorem~\ref{thm:frequency-unipotent-lattice}, and the right
translate of a left fundamental set is another such set.
Decomposing one set by left rational translates of the other
gives equal integrals for left rationally invariant functions.
The pushforward under \(b\) is therefore a probability on
\([D]\) invariant under \(b(n_0)\), hence under all translations
by surjectivity. Uniqueness by the double-integration argument
of Lemma~\ref{lem:constant-term-adelic-substitution} proves the
identity.
\end{proof}

\section{The definite quotients in two consecutive periods}
\begin{theorem}\label{thm:actual-double-period-mean}
Use the spaces \(Y,Y',Y''\), tests \(f_1,f_2\), and maps
\(T,T',\tau\) of Corollary~\ref{cor:paired-global-frequency-product}.
Write \(H={\rm U}(X_j),\ H'={\rm U}(X')\). Let \(k\) be a
unitary on the definite block of \(Y''\), extended by the identity
on its string, and in the terminal pair let \(k\) be any unitary
on \(Y_m\). Fix adelic positions \(x,x',x''\) on the three spaces.
Use \(p=P_j\) in Theorem~\ref{thm:actual-unfolded-stage-functions},
or any continuous \(p\) with its rational Bessel invariance and
its character there. Denote the two succeeding period operations
of that theorem at \(x,x'\) by \(p\mapsto p_1\mapsto p_2\).
Use the product measures of Lemma~\ref{lem:bessel-product-probability}.

Write \(\Omega_1,\Omega_2\) for the two twisted arrow operators.
Set
\[
 u=\mathcal J\big(\Omega_1(x,x') f_1
                        \otimes\Omega_2(x',x'') f_2\big),
 \qquad
 \mathcal G(g,h,k)=\mathcal J
       \big(\Omega_1(g,h)\otimes\Omega_2(h,k)\big)\mathcal J^{-1},
\]
using the identification of Theorem~\ref{thm:global-two-arrow-adjoint}.
An element of a definite block in this formula acts trivially on
the attached string. Choose a rational polarization adapted to
\(U_*,W_*\) from Corollary~\ref{cor:residual-pair-positive-quotient}
and write \(Z\) for restriction of its position variable to the
residual position space. With its rational theta functional
\(\theta_{W_*}\), one has
\[
 p_2(kx'')=\int_{[H]}\int_{[H']}p(gx)\,
       \theta_{W_*}\big(Z\rho_*(\tau)\mathcal G(g,h,k)u\big)
                                               \dd h\dd g .
\]
All integrals are absolutely convergent.
\end{theorem}
\begin{proof}
Expanding the two period operations gives an integral over
\(r\in[B],r'\in[B']\) with integrand \(p(r x)\) times the two
lowering-map coefficients at \((r x,r' x')\) and \((r' x',k x'')\).
The group law absorbs the three fixed positions into the two
tests as in the definition of \(u\). Each rational Bessel
translation in either integration variable preserves the
integrand. Indeed rational covariance in
Theorem~\ref{thm:global-actual-arrow-covariance} acts on the
frequency by that translation, and the rational change at \(T\)
lies in the nonnegative space, as in the last paragraph of
the proof of Theorem~\ref{thm:one-rational-bessel-unfold}.
The same holds at \(T'\). Absolute convergence and interchanges
of integration follow from continuity on the compact representative
closures of Proposition~\ref{prop:all-bessel-quotients-compact}
and Lemma~\ref{lem:product-integration}.

Write \(r=n g\) as in Lemma~\ref{lem:bessel-product-probability}.
Every grade in \(Y\) is even. Hence \(n-1\) raises degree by at
least two. Since \(T\) lowers by one, \(T(n-1)\) is positive,
so \(n\in J_T\) of Theorem~\ref{thm:frequency-unipotent-lattice}.
The factor \(p(ngx)\) acquires \(B(Tn,T)\), and the first
coefficient acquires its inverse by
Corollary~\ref{cor:graded-source-target-covariance}, where
\(\eta_T=1\). The second coefficient is independent of \(n\).
The \(n\)-integral thus has the value at \(n=1\).

Write \(r'=n'h\) with \(h\in H'(\mathbb A)\). At \(n'=1\)
the product of coefficients is
\(I_\tau^{U_>}(F_v)\) for \(v=\mathcal G(g,h,k)u\), by
Corollary~\ref{cor:paired-global-frequency-product}. Inserting
\(n'\) acts on \(v\) by the combined untwisted unitary operator
for \(Gz=z(n')^{-1}\): all determinant factors from
Theorems~\ref{thm:global-two-arrow-adjoint}
and~\ref{thm:global-actual-arrow-covariance} are one for this
unipotent. The difference \(G-1\) strictly raises Hom degree.
For the grading with positive part \(U_>\) it follows that
\(I_\tau^{U_>}(F_{\Omega_*(G)v})=I_{\tau n'}^{U_>}(F_v)\).
Indeed the dilation by \(t^i\) on degree \(i\) contracts \(G\)
in the parabolic preserving the polarization variable.
Lemma~\ref{lem:parabolic-splitting-comparison} and the
restricted-product argument of Theorem~\ref{thm:global-actual-arrow-covariance}
identify its operator with \(\mathcal W_G\). The stated
coefficient formula is Theorem~\ref{thm:graded-parabolic-schroedinger}.

Put \(b=b(n')\) as in Proposition~\ref{prop:middle-unipotent-additive-mean}.
The difference \(\tau(n'-1)=b+v_+\) has \(v_+\in U_>\).
Indeed it has degree at least zero, and by
Propositions~\ref{prop:actual-pair-grade-zero}
and~\ref{prop:paired-rational-frequency} its degree-zero part
is precisely the displayed row into \(E\xi\). For the final
pair the difference vanishes. The adjoint in the negated first
block gives
\[
         \tau^\dagger_*\tau=-TT^\dagger+(T')^\dagger T'=0
\]
by Lemma~\ref{lem:two-string-moments} applied to the common
middle space. In particular
\(H_*(\tau,\tau(n'-1))=0\), so \(B(b+v_+,\tau)=1\).
Moreover \(B(v_+,b)=1\) by orthogonality of positive and zero
degrees. The Heisenberg law and invariance of the positive
quotient mean now give
\[
    I_{\tau n'}^{U_>}(F_v)
       =I_0^{U_>}\big(F_{\rho_*(b)\rho_*(\tau)v}\big).
\]
The last expression is continuous and periodic in \(b\) modulo
\(D(E)\). For a rational \(d\in D\), its translation factor
on \(U_>+D\) is one by total isotropy of that sum and rational
theta invariance in Theorem~\ref{thm:rational-schroedinger-zero-section}.
Proposition~\ref{prop:middle-unipotent-additive-mean} replaces
the \(n'\)-integral by the additive mean over \([D]\).
Because \(U_*=U_>\oplus D\) is totally isotropic, averaging
over \([U_>]\times[D]\) gives \(I_0^{U_*}(F_{\rho_*(\tau)v})\).
This holds also when \(D=0\).

Finally a rational hermitian isotropic dual to \(U_*\) is given
by the negative Hom grades and, when \(D\ne0\), the extra
\({\rm Hom}(X',E\eta)\) in Proposition~\ref{prop:actual-pair-grade-zero}.
Its pairing with \(U_*\) is perfect and the remaining space is
\(W_*\). Choose dual rational position and modulation variables
in \(W_*\) by the plane removal of Lemma~\ref{lem:rational-graded-polarization}.
The proof of Theorem~\ref{thm:rational-schroedinger-zero-section}
then identifies \(I_0^{U_*}\) with \(\theta_{W_*}Z\): a rational
position outside the residual position space has a nontrivial
character of \([U_*]\). This gives the asserted integral.
\end{proof}

\section{The group action after restriction to the residual variables}
\begin{lemma}\label{lem:weil-unimodular-frame-independence}
In the Schr\"odinger model of Theorem~\ref{thm:actual-adelic-rational-theta-invariance},
the physical finite unitary action on a hermitian space over
the arithmetic field \(E\) is independent of the choice of
reference lattice \(\alpha M\) among finite adelic
\(\widehat O_E\)-lattices \(M\) unimodular for the hermitian form.
Here unimodular means \(M=\{v:H(M,v)\subset\widehat O_E\}\),
and \(\alpha\) is the fixed idele in Lemma~\ref{lem:actual-lift-scale}.
The real action is independent of the choice of a signed
orthonormal complex basis for the same hermitian space.
\end{lemma}
\begin{proof}
Each reference lattice \(\alpha M\) is self-dual for the
alternating trace character, and its square has trivial
bicharacter, by
\(\jmath\alpha\bar\alpha\widehat O_E=\widehat O_E^\vee\)
in Lemma~\ref{lem:hom-oscillator-dual-lattices}.
Here the right hand side is the trace dual of the integer ring.
Two such \(M,M'\) are carried to one another by a finite adelic
unitary \(k\). To see this, both have integral bases locally,
using elimination over each valuation ring, and their Gram
matrices \(J,J'\) are hermitian and invertible over the ring
by unimodularity. At a split pair choose a basis on one branch
and its dual on the other. At a nonsplit place the extension
is unramified by Lemma~\ref{lem:real-place-rings}. Modulo its
real uniformizer \(\pi\), the diagonalization in
Lemma~\ref{lem:finite-unitary-primitive-frame} makes \(J,J'\)
isometric. Given an invertible \(A\) with \(A^*J'A=J\bmod\pi^t\)
for \(t\ge1\), put \(S=A^*J'A\). Modulo \(\pi\) solve
\[
                C^*\bar S+\bar S C
                       =\big[(J-S)/\pi^t\big]\bmod\pi ,
\]
where \(\bar S\) in this display denotes residue reduction.
The right side is hermitian and can be written \(D+D^*\)
by the trace surjectivity in that same lemma. Take \(C=\bar S^{-1}D\).
Replacing \(A\) by \(A(1+\pi^t\widetilde C)\) for any lift
\(\widetilde C\) increases the precision by one. Completeness
and invertibility of the first residue give an exact isometry.
In rational coordinates the given lattices agree almost everywhere:
each is contained in a common integer-denominator multiple of
the standard lattice and contains a common positive integer
multiple of it, by compactness and openness. Take the local
isometry to be the identity there. This proves the claim about \(k\).

Let \(\phi\) be a unit vector on the fixed line for \(\alpha M\).
The fixed line for \(\alpha M'\) is spanned by
\(\mathcal V_k\phi\) by Proposition~\ref{prop:finite-polarized-vacuum}.
For each unitary \(G\) the old operator \(\mathcal V_G\)
implements the prescribed translations. Its coefficient against
the new vacuum equals the old coefficient for \(k^{-1}Gk\)
by the group law of Theorem~\ref{thm:all-finite-weil-splitting}.
The ideal and lattice index of \(k^{-1}Gk\) for \(\alpha M\)
are those of \(G\) for \(\alpha M'\), since \(k\) identifies
the respective intersections and ambient bases. By
Theorem~\ref{thm:split-cut-weil-multiplier}, the coefficient
is therefore exactly \(\beta(J_G)[G\alpha M':G\alpha M'\cap
\alpha M']^{-1/2}\) with the new reference. This nonzero
coefficient and the intertwining law determine the operator
by Lemma~\ref{lem:parabolic-splitting-comparison}. This proves
independence at the finite places.

At a real place a new signed basis is the image of the old
one by a unitary \(k\) for the indefinite hermitian form. The
Gaussian in the new coordinates, transported into the old
physical model, is proportional to \(\Omega(k)\phi_0\).
Indeed coordinate transport implements the change of the two
symplectic bases and two unitary implementers differ by a
scalar by Lemma~\ref{lem:schroedinger-real-gaussian}. The scalar
has modulus one. The coefficient of \(\Omega(G)\) on this
Gaussian is the old coefficient of \(\Omega(k^{-1}Gk)\). The
negative block of \(G\) in the new basis is the old negative
block of \(k^{-1}Gk\). The uniqueness statement of
Theorem~\ref{thm:negative-block-real-splitting} now identifies
the two operators.
\end{proof}

\begin{theorem}\label{thm:residual-levi-zero-section}
Let \(V=U_-\oplus W\oplus U\) be a rational hermitian space over
\(E\), with \(U_-,U\) isotropic and in perfect duality, both
orthogonal to the nonsingular \(W\). Assume \(W\) has a unimodular
finite lattice after division of its reference by \(\alpha\).
Give \(V,W\) the splittings of
Theorem~\ref{thm:actual-adelic-rational-theta-invariance} with
reference lattices of Lemma~\ref{lem:weil-unimodular-frame-independence}.
For an adelic isometry \(G\) preserving each of the three
summands, put \(A=G|_{U_-}\) and \(G_W=G|_W\). Take a rational
symplectic polarization of \(W\) and use \(U_-\) as the additional
position variable of \(V\). If \(Zf(t)=f(0,t)\) denotes evaluation
in that variable, then
\[
       Z\Omega_V(G)f=
         \upsilon(\det_E A)^{-1}|\det_{\Q} A|_{\mathbb A}^{-1/2}
                                \Omega_W(G_W)Zf .
\]
The map \(Z\) is the product of real and finite restriction on
each real-times-finite test. In particular it applies to
\(U_*,W_*\) of Theorem~\ref{thm:actual-double-period-mean}.
\end{theorem}
\begin{proof}
In a rational \(E\)-basis of \(U_-\) choose the integral
lattice \(M_-\), and take its hermitian dual \(M_+\) in \(U\).
The union of the two corresponding dual bases has hyperbolic
Gram matrix, so \(M_-\oplus M_+\) is unimodular. Use its
orthogonal sum with the indicated unimodular lattice on \(W\).
Lemma~\ref{lem:weil-unimodular-frame-independence} permits this
choice. Also choose real signed orthonormal bases on these two
orthogonal blocks separately. In a product polarization the
operator on the orthogonal sum is then the product of the block
operators. At finite places the reference vacuum and the
index ideals factor, as does the normalized lattice sum in
Definition~\ref{def:finite-splitting-multiplier}. At infinity the
normalized Gaussian factors and its coefficient is the reciprocal
of the product of the two negative block determinants, by
Theorem~\ref{thm:negative-block-real-splitting}. The product
implements the translations on the direct sum, so the scalar
uniqueness arguments of Theorem~\ref{thm:actual-adelic-rational-theta-invariance}
give the assertion in the chosen rational polarization.
On the real Schwartz space the coordinate substitutions,
quadratic phases and partial Fourier transforms on \(W\) act
with the \(U_-\) coordinates as parameters, so they commute with
evaluation there. For the Fourier transforms this follows
directly from their absolutely convergent integral on Schwartz
functions in Lemma~\ref{lem:rational-partial-poisson}.

The hyperbolic factor acts by
\(\upsilon(\det_E A)^{-1}\mathcal W_G\) on its own variables
by Theorem~\ref{thm:hyperbolic-adelic-parabolic-splitting}.
At zero its phase is one and the normalized substitution there
multiplies by \(|\det_\Q A|_{\mathbb A}^{-1/2}\), by
Theorem~\ref{thm:graded-parabolic-schroedinger}. This proves
the formula, including \(U=0\). The evaluation is literally
a restriction of the two variables on a real-times-finite
test. It remains a Schwartz function at infinity and is locally
constant with compact support at finite places. For the
paired application the dual was constructed at the end of
Theorem~\ref{thm:actual-double-period-mean}. The standard
residual lattice in Theorem~\ref{thm:paired-neighbour-vacuum-sum}
is of the required form: after dividing its physical coordinates
by \(\alpha\) the two scalar multiples in
Definition~\ref{def:paired-neighbour-lattice} differ by
\(\bar\alpha/\alpha\) of norm one. Thus the resulting signed
matrix lattice is unimodular.
\end{proof}

\section{Determinant cancellation on the two definite arrows}
\begin{proposition}\label{prop:paired-restriction-determinant}
In Theorem~\ref{thm:actual-double-period-mean} write \(X=X_j\),
and write \(X''=X_{j+2}\) for \(s\ge3\), or
\(X''=Y_m\oplus E(z/\pi_N)\) for \(s=1\).
Let \(k_0\) be \(k\) on the usual definite target, extended
by one on the additional line if present. For \(g,h,k\) in
that theorem write \(A=\det_E g,\ B=\det_E h,\ C=\det_E k\).
The action \(Gz=(g\oplus k)zh^{-1}\) on the combined Hom
space preserves \(U_*,W_*\) and their chosen negative dual.
On either of the two isotropic summands its determinant is
\((AC)^a B^{-2a}\), where \(s=2a+1\). It fixes \(\tau\).
Let \(\mathcal G_0(g,h,k_0)\) on \(W_*={\rm Hom}(X',X^-\oplus
X'')\) be the adjoint combination of the two twisted definite
arrows, with their own ranks and the same two characters
\(\mu_l,\mu_{l+1}\) as on the indefinite arrows. Then
\[
     Z\rho_*(\tau)\mathcal G(g,h,k)
           =\mathcal G_0(g,h,k_0)Z\rho_*(\tau).
\]
\end{proposition}
\begin{proof}
The three definite unitaries fix every string vector and preserve
the grades. They also fix \(\eta,\xi\) of
Proposition~\ref{prop:actual-pair-grade-zero} when present.
The descriptions in Corollary~\ref{cor:residual-pair-positive-quotient}
and Theorem~\ref{thm:actual-double-period-mean} prove preservation
of the three summands. By Proposition~\ref{prop:paired-rational-frequency}
the two components of \(\tau\) vanish on \(X'\) and have image
in the outside strings, so \(G\tau=\tau\). Intertwining of the
Heisenberg operators therefore commutes \(\rho_*(\tau)\)
through \(\mathcal G\).

For \(a>0\), the middle string has \(a\) negative lines.
Positive Hom entries from these lines into the outside definite
blocks contribute determinant \((AC)^a\). The positive outside
string lines number \(a+(a-1)\). Positive Hom entries into
them from \(X'\) contribute \(B^{-(2a-1)}\). The additional
maps from \(X'\) into \(E\xi\) contribute \(B^{-1}\). Every
other positive entry has trivial action since both indices
are in strings. Thus the claimed determinant holds on \(U_*\).
For \(a=0\) this space is zero and the determinant is one.
By hermitian duality the determinant on the opposite summand
is the inverse conjugate of that on \(U_*\), equal to it
since \(A,B,C\) have norm one. The determinant modulus over
the rational adeles is consequently one by the norm calculation
of Lemma~\ref{lem:bessel-product-probability}.

Write \(n,n',n''\) for the original three ranks, and \(d,d',d''\)
for the definite ones, including the additional final line.
By Theorems~\ref{thm:global-two-arrow-adjoint}
and~\ref{thm:global-actual-arrow-covariance}, the scalar on
the untwisted combined action is
\[
 \frac{\mu_l(B)\mu_{l+1}(C)}{\mu_l(A)\mu_{l+1}(B)}\,
 \upsilon_T(A)^{\lceil n'/2\rceil}
 \upsilon_T(C)^{\lfloor n'/2\rfloor}
 \upsilon_T(B)^{-\lceil n/2\rceil-\lfloor n''/2\rfloor}.
\]
Indeed the first Hom action \(v\mapsto hvg^{-1}\) has determinant
\(B^n A^{-n'}\), by its action on columns and rows, and adjunction
supplies the inverse splitting character of this determinant.
The remaining exponents are those in
Definition~\ref{def:actual-twisted-arrow-character}. The same
calculation applies to the three definite spaces: the proof of
Theorem~\ref{thm:global-two-arrow-adjoint} uses just the Hom
pairing and the signs of the orthonormal frames. Moreover
\(\det k_0=C\) even on the enlarged target.

By Definition~\ref{def:actual-row-lift-spaces},
\[
 n=d+2a+1,\qquad n'=d'+2a,\qquad n''=d''+2a-1 .
\]
The last equation for \(a=0\) uses \(d''=n''+1\). The outside
ranks \(n,n''\) differ by \(2t\), and hence \(d,d''\) have the
same parity. The first two exponents involving \(n'\) increase
by \(a\) upon passing from definite to original ranks. The sum
\(\lceil n/2\rceil+\lfloor n''/2\rfloor\) increases by \(2a\):
if \(d,d''\) are even the two increases are \(a+1,a-1\);
if they are odd the two increases are \(a,a\). The ratio of
the original scalar to the definite one is therefore
\(\upsilon_T((AC)^aB^{-2a})\). Theorem~\ref{thm:residual-levi-zero-section}
gives precisely its inverse on applying \(Z\), proving the identity.
\end{proof}

\begin{corollary}\label{cor:actual-reduced-definite-pair}
In the notation of Proposition~\ref{prop:paired-restriction-determinant},
let \(\mathcal J_0\) denote the adjoint identification for
the two definite Hom spaces and put
\[
                          v_0=\mathcal J_0^{-1}Z\rho_*(\tau)u .
\]
Let \(\theta_{12}\) be their rational theta functional, the
product of the two individual theta functionals on pure tensors.
Then the integrand on the two definite quotients in
Theorem~\ref{thm:actual-double-period-mean} equals
\[
 p(gx)\,\theta_{12}\!\left(
       (\Omega^{0}_1(g,h)\otimes\Omega^{0}_2(h,k_0))v_0\right),
\]
where the superscript \(0\) denotes the operators on the
definite arrows. For elementary real-times-finite input tests,
\(v_0\) has one real Schwartz factor and one finite factor.
For the moved finite inclusion tests the latter is the shifted
finite evaluation of Proposition~\ref{prop:paired-output-shifts},
with the actual arrow multipliers and a nonzero scalar depending
only on the fixed models. In particular, for fixed real tests
and fixed infinite positions the real factor can be chosen
independently of all weight-zero twists and finite positions.
\end{corollary}
\begin{proof}
Apply the proposition and the theta normalization of
Theorem~\ref{thm:global-two-arrow-adjoint} on the definite pair.
These identifications extend to its full real Schwartz space:
the rational linear coordinate transports and partial Fourier
transforms in its proof act on joint Schwartz functions,
and the rational sum on the two position spaces is absolutely
convergent by Proposition~\ref{prop:rational-frequency-lattice-bound}.
This defines \(\theta_{12}\) on the joint space. The finite
part is the algebraic tensor product by the translated basis
of Theorem~\ref{thm:paired-finite-adjoint-model}.

All the coordinate identifications, the shift by \(\tau\) and
the restriction factor into real and finite operators by
Theorems~\ref{thm:global-two-arrow-adjoint}
and~\ref{thm:residual-levi-zero-section}. In the adapted finite
Schr\"odinger model \(Z_f\rho_f(v)=Z_f\) for \(v\in U_*\),
since its modulation is one at zero and it is perpendicular
to \(W_*\). For \(w\in W_*\), \(Z_f\rho_f(w)=\rho_{W_*,f}(w)Z_f\):
in the formula of Theorem~\ref{thm:rational-schroedinger-zero-section}
both the translation and its phase use only the residual variables
at that section. This map is onto. Multiply any finite residual
test by the indicator of a compact open lattice containing zero
in the opposite variables to extend it. Under the unitary
finite identifications of Proposition~\ref{prop:finite-polarized-vacuum},
Theorem~\ref{thm:finite-grade-oscillator-quotient} now identifies
\(Z_f\) with a scalar multiple of its normalized quotient map
in the residual lattice model. Indeed the quotient is spanned
by residual Heisenberg translates of the vacuum, and a commuting
endomorphism of that space is scalar by uniqueness of the
fixed line. Surjectivity makes this scalar nonzero. To use the
standard residual lattice, send that vacuum to the unit fixed
vector of Theorem~\ref{thm:paired-neighbour-vacuum-sum}, choosing
its positive overlap with the standard vacuum. The choice of
phase can be transferred to the real factor while preserving
the global identification.

Consequently inserting \(\rho_f(\tau)\) before this map gives
the shifted evaluation of Corollary~\ref{cor:paired-finite-zero-term-index}.
The final inverse adjoint identification is exactly the one
in Proposition~\ref{prop:paired-output-shifts}. The fixed
scalar is independent of which vector undergoes those operations.
Finally \(\vartheta\) in Definition~\ref{def:actual-twisted-arrow-character}
has weight zero and hence trivial infinite factor. The two
moved real tests thus depend only on the fixed infinite tests
and positions, as claimed.
\end{proof}

\section{A central translation through the successive definite pairs}
\begin{lemma}\label{lem:initial-bessel-column-coefficient}
Let \(\pi\) be the product of the symmetric tensor representations
in Theorem~\ref{thm:actual-real-bessel-vector} at the infinite
places and let \(W=\bar\pi\) be its conjugate representation space
from Definition~\ref{def:columns}. Write
\(\mathcal B(\bar v,w)=\langle v,w\rangle\) for the complex bilinear
pairing of \(W\) with \(\pi\), using the tensor Hermitian product
conjugate linear in the first variable. Denote the definite
component of \(b_0\) in Lemma~\ref{lem:initial-frame-central-translation}
by \(b\), a finite central scalar on \(X_0\). In
Theorem~\ref{thm:actual-base-bessel-periodization} use the
functionals
\[
              \lambda_c(v)=\mathcal B(\phi(h_c b^{-1}),v)
\]
for a column \(\phi\), viewing \(h_c\) also as the corresponding
definite representative. Use the right-translated stage-zero
function, denoted \(P_{0,\phi}\). All coordinate families of
that theorem occur in this way. If \(x_{0,f}\in A^{\rm new}\),
then for every \(a=(a_\infty,a_f)\in{\rm U}(X_0)(\mathbb A)\)
and its block embedding \(\widehat a\),
\[
   P_{0,\phi}(\widehat a x_0)
       =\mathcal B\big(\phi(a_f),\pi(a_\infty)F_\infty(x_{0,\infty})\big).
\]
\end{lemma}
\begin{proof}
The pairing is invariant under the simultaneous conjugate and
ordinary actions because \(\pi\) is unitary. It identifies \(W\)
linearly with the full dual of \(\pi\), as is seen in an orthonormal
basis. The classes of \(h_c b^{-1}\) form a complete system since
\(b\) is central. The freedom of evaluation in
Proposition~\ref{prop:column-dim} proves the first assertion.

Right invariance allows \(x_{0,f}\) to be replaced by one.
In the old realization the resulting finite position is the
block embedding of \(a_f b\). Write \(a_f b=\gamma_f h_c k\)
with \(\gamma\) a rational definite unitary and \(k\in K_c\).
Its embedding lies in the corresponding old right level by
Proposition~\ref{prop:marked-local-level}. Rational left
invariance and evaluation at \(h_c\) in
Theorem~\ref{thm:actual-base-bessel-periodization} thus give
\(\lambda_c(F_\infty(\widehat\gamma_\infty^{-1}
\widehat a_\infty x_{0,\infty}))\).
On the definite block the string character is one by its
simple-root formula in Proposition~\ref{prop:full-phase}.
The transformation of Theorem~\ref{thm:actual-real-bessel-vector}
therefore gives \(\pi(\gamma_\infty^{-1}a_\infty)F_\infty(x_{0,\infty})\)
inside \(\lambda_c\). Invariance of the bilinear pairing moves
\(\gamma_\infty\) to the first variable, acting there by the
conjugate representation. Since \(b\) is central,
\(a_f=\gamma_f h_c b^{-1}k\), and the transformation of
\(\phi\) in Definition~\ref{def:columns} gives the formula.
\end{proof}

\begin{theorem}\label{thm:even-period-central-translation}
Take the stage functions of Theorem~\ref{thm:actual-unfolded-stage-functions}
with the preceding seed and the indicated finite inclusion
vectors. Assume \(x_{0,f}\in A^{\rm new}\). Fix a split pair
over \(N\) and the integral norm-one finite idele \(z\)
supported on that pair which lifts the generator in
Theorem~\ref{thm:ram-dim}. Write \(z_i\) for scalar \(z\) on
\(X_i\), embedded by identity on the string in \(Y_i\), and
\((R_z\phi)(v)=\phi(vz)\) on the initial definite group.
For every even \(i\) with \(2\le i<m\), every preceding choice
of positions and every \(v\in{\rm U}(Y_i)(\mathbb A)\),
\[
               P_{i,R_z\phi}(v)=P_{i,\phi}(z_i v).
\]
At stage zero one has the same identity on \(v=\widehat a x_0\)
for \(a\in{\rm U}(X_0)(\mathbb A)\). For every ramified column,
every \(0\le i<m\), and every \(v\) with \(P_{i,\phi}(v)\ne0\),
the exceptional middle-line case of
Proposition~\ref{prop:opposite-string-line-duality} is excluded
at all places.
\end{theorem}
\begin{proof}
The stage-zero identity follows from
Lemma~\ref{lem:initial-bessel-column-coefficient}, since
\(z\) has trivial real component and commutes with \(a_f\).
Prove the claim at \(i=j+2<m\) by induction in pairs. This is
a pair with \(s\ge3\), so all three residual definite spaces
are \(X_j,X_{j+1},X_{j+2}\) without an added line. On either
definite arrow the simultaneous scalar at its two ends acts
on the Hom space by the identity. Its two determinant values
are units supported over \(N\). Every arrow character in
Definition~\ref{def:actual-twisted-arrow-character} is trivial
on these units by Theorem~\ref{thm:actual-twisted-finite-arrow}
because \(N\) is disjoint from the marked and twist places.
The full twisted action of the simultaneous scalars on each
definite arrow is therefore the identity, on every test.
In particular the theta kernel \(K(g,h,k)\) of
Corollary~\ref{cor:actual-reduced-definite-pair} for the given
positions satisfies
\[
 K(z_j g,z_{j+1}h,z_{j+2} k)=K(g,h,k).
\]
Here the scalars in arguments of the definite groups act
on the whole definite spaces. Only their embeddings into
the original stages fix the strings.

Choose \(x,x',x''\) in the paired formula with \(x=x_j\),
\(x'=x_{j+1}\) and arbitrary \(x''=v\). Evaluate it first
with \(k=z_{j+2}\). Change the two integration variables to
\(z_jg,z_{j+1}h\). On each definite quotient this preserves
the measure, since central left translation commutes with
the rational left action and preserves left Haar measure.
The kernel becomes \(K(g,h,1)\). The factor
\(P_{j,\phi}(z_j g x_j)\) becomes \(P_{j,R_z\phi}(g x_j)\)
by the preceding stage, using the restricted stage-zero identity
when \(j=0\). The integral now gives \(P_{j+2,R_z\phi}(v)\)
by the same paired formula. The tests in that formula depend
on the three fixed positions and are independent of \(\phi\).
This proves the claimed central equivariance for arbitrary \(v\).

At any positive stage \(i\), the function \(P_{i,\phi}\) in
its last argument is right invariant under the whole local
integral unitary group \(K^{(w)}\) of
Proposition~\ref{prop:exceptional-middle-level-stabilizer}.
Indeed on the last input arrow \(\Omega(1,k)\) fixes its inclusion
test for every \(k\in K^{(w)}\), by the integral case of
Theorem~\ref{thm:actual-twisted-finite-arrow}. Right multiplication
in the defining integral therefore changes no factor. If \(i\)
is even then the source grades at \(i\) are even, and all the
strictly grade-raising unipotents there have degree at least two.
The character from Theorem~\ref{thm:actual-unfolded-stage-functions}
is invariant under conjugation by the embedded definite group,
since its shift \(S_s\) commutes with that group by
Lemma~\ref{lem:two-string-moments}. At each fixed infinite
argument all three hypotheses of Lemma~\ref{lem:single-central-difference-test}
are therefore satisfied for \(2\le i<m\) even. Apply
Corollary~\ref{cor:ramified-middle-test} to obtain vanishing at
the exceptional points. For odd \(i\) there is no middle line.
At \(i=0\) the support is inside \(S_f A^{\rm new}\) by
Corollary~\ref{cor:actual-base-stage-cut}. All initial string
ideals on this set are those of the reference lattice by
Theorem~\ref{thm:finite-bessel-flag-classification} on the early
lines. In particular its middle exponent sum is zero, since
the middle generator \(u_0\) in Lemma~\ref{lem:initial-line-identity}
has norm one. Proposition~\ref{prop:opposite-string-line-duality}
has now excluded all remaining possibilities.
\end{proof}

\section{Finite cuts for the periods of the ramified columns}
\begin{theorem}\label{thm:actual-ramified-period-cuts}
For fixed arithmetic data and twist set \(P\), use the stage
functions of Theorem~\ref{thm:even-period-central-translation}
for all ramified columns, all real tests, and all choices of
preceding positions with \(x_{0,f}\in A^{\rm new}\) whenever
this position occurs. Let \(W_i\) be the union of their
nonzero supports projected to the finite group of \(Y_i\).
For every \(0\le i<m,\ y\in W_i\), every finite place \(w\)
and every string index \(b<m-1-i\), \(b\ge0\), the deviation
from the reference exponent satisfies
\[
                 -i\,{\bf1}_P(w)\le\Delta_b(y;w)
                               \le i\,{\bf1}_{\bar P}(w).
\]

Let \(K_i\) denote the full integral group of the stage frame
and \(J_i\) the subgroup of \(K_i\) with identity reduction at
the selected places over \(g\) and at \(P\). Put \(A_0=A^{\rm new}\),
\(A_i=K_i\) for \(0<i<m\), \(S_i=B_{Y_i}(\mathbb A_f)\), and
\(Q=\prod_{w\in P}q_w\). There are finite cuts \(T_i\) in
the sense of Definition~\ref{def:finite-period-cut} with
\(W_i/J_i\subseteq S_iT_i\). They are stable under
\(S_i\cap A_i\), contain representatives integral off
\(P\cup\bar P\), and have \(m_i(t)^{-1}\in\Z[1/Q]\) on their
saturations, where \(m_i=m_{T_i}\). Every representative of
\(T_0\) lies in \(A^{\rm new}\). At \(i>0\) the support can be
covered up to left \(S_i\)-translation by at most
\((i+1)^{\lceil(m-1-i)/2\rceil |P\cup\bar P|}\) right \(K_i\)-cosets.
All choices can be made independently of the column and positions.
\end{theorem}
\begin{proof}
Stage zero has its claimed support by
Corollary~\ref{cor:actual-base-stage-cut}. Right level elements
preserve the reference lattice and a Bessel element acts
trivially on the successive string line quotients by
Theorem~\ref{thm:actual-unfolded-stage-functions}, giving zero
deviation for the full string. Suppose the inequalities have
been proved at stage \(i\). For a nonzero period at stage
\(i+1\) that same theorem finds a single source point of
nonzero period at \(r x_i\) which also satisfies the one-step
lattice comparison with the given target. Both points have
the middle-line condition by Theorem~\ref{thm:even-period-central-translation}.
The source finite position is one to which the inductive bound
applies, since all positions in its own prior stages are the
same and the current argument of the function is unrestricted.
If the target string is nonempty, use the first inequality
pair of Theorem~\ref{thm:one-sided-stage-line-interlacing}
at every target index. Its upper endpoint adds at most
\({\bf1}_{\bar P}(w)\) to a source upper bound, and its lower
endpoint subtracts at most \({\bf1}_{P}(w)\) from a source
lower bound. For empty string there are no inequalities.
This proves the displayed conclusion simultaneously at all
places for the union \(W_i\).

At stage zero apply Corollary~\ref{cor:actual-base-stage-cut}
with the given \(J_0\). For each positive stage the proof of
Theorem~\ref{thm:near-flag-adelic-period-cut} uses just the list
of possible tuples of the early deviations. The displayed
one-sided bound permits at most \(i+1\) values at each of the
indicated places for each early index and one value elsewhere.
Thus the same construction gives the indicated number of cosets.
The subgroup of elements integral off \(P\cup\bar P\) contains
\(K_i\) and all the chosen representatives. Applying
Theorem~\ref{thm:finite-period-cut-denominator} with the local
power bounds of Proposition~\ref{prop:adelic-period-denominator-depth}
gives the cuts and their stated multiplicities, also with no
string indices. Multiplying any representative by a member of
\(J_i\) keeps it integral at the complementary places. The cut
uses the whole support union, proving the asserted uniformity.
\end{proof}

\begin{corollary}\label{cor:actual-weighted-cut-tensor-bound}
Use the full product of the cuts of
Theorem~\ref{thm:actual-ramified-period-cuts}. At each tuple
take the tensor of the finite evaluations in
Corollary~\ref{cor:actual-reduced-definite-pair}, with last
finite position one. Divide by the fixed nonzero model scalars
of that corollary, use the standard normalized residual maps,
and call the resulting finite tensor \(v(t)\). With \(D,C_*\)
from Theorem~\ref{thm:paired-finite-chain-denominator} the tensor
\[
      w=\frac{D}{C_*}\sum_{t\in\prod_{i<m}T_i}
                         \frac{v(t)}{\prod_{i<m}m_i(t_i)}
\]
is fixed by the product of the column levels, satisfies
\(\mathcal S\mathcal P w=w\), and has an expansion in allowed
translated sums with coefficients integral away from \(Q\).
When \(P=\varnothing\) the sum of the absolute values of its
coefficients after every field automorphism is at most
\[
 D\prod_{i=1}^{m-1}\prod_{p\in\{g_0,g_1\}}
                  \prod_{u=0}^{n_i-1}(p^{n_i}-p^u)^f
      \ \le\ Dg^{f\sum_{i=1}^{m-1}n_i^2}
      \ \le\ (g^{f n_m^2})^{2m-1}.
\]
Here \(n_i=r+it\) are the original stage ranks, with \(t\)
the block size in this rank formula.
\end{corollary}
\begin{proof}
The finite comparison in Corollary~\ref{cor:actual-reduced-definite-pair}
identifies each normalized tensor with the finite vector of
Corollary~\ref{cor:actual-character-finite-chain}. Its right
level invariance there makes the choice of representatives
immaterial. Apply Theorem~\ref{thm:cut-sum-of-finite-evaluations}.
The embedded column levels lie in the respective \(S_i\cap A_i\)
by Proposition~\ref{prop:column-level-in-stage-bessel}, so
Corollary~\ref{cor:actual-column-invariant-cut-tensor} gives
column invariance and the recovery identity.

In the untwisted case \(J_0=A^{\rm new}=A_0\). The first cut
can be the single identity coset, of multiplicity one by
Corollary~\ref{cor:actual-base-stage-cut}. At each later stage
every multiplicity is one and the cut lies in \(K_i/J_i\) by
Proposition~\ref{prop:marked-stage-cut-size}. Its size is at
most the exact marked index in that proposition. Multiplying
these bounds in the sum-of-absolute-values proof of
Theorem~\ref{thm:cut-sum-of-finite-evaluations} gives the
first expression. Bound each marked index by \(g^{f n_i^2}\)
and use \(D\mid(g^{f n_m^2})^m\) and \(n_i\le n_m\) for the
remaining two bounds.
\end{proof}

\section{Integrating the normalized chain over the finite cuts}
\begin{lemma}\label{lem:actual-period-cut-disintegration}
Use the groups and measures of Theorem~\ref{thm:period-cut-sum}.
Suppose \(f\) satisfies its support, rational invariance and
integrability assumptions and is right \(J\)-invariant in its
finite argument. Choose \(\alpha\) of
Lemma~\ref{lem:period-orbital-cutoff}. For \(t\in T\) fix a
representative \(h_t\). Then
\[
 \int_{D_G} f=
   \frac1{[K:J]}\sum_{t\in T}\frac1{m_T(t)}
     \int_{G_1}\alpha(x)\int_{D_U} f\big(u(x,h_t)\big)\dd u\dd x .
\]
Each integral on the right is absolutely convergent with
\(|f|\) in place of \(f\).
\end{lemma}
\begin{proof}
The cutoff computation in Theorem~\ref{thm:period-cut-sum}
applied first to \(|f|\) bounds the integral of its period
against \(\alpha(x){\bf1}_T(hJ)/m_T(hJ)\) on the full group.
It equals the prescribed finite integral of \(|f|\) on \(D_G\).
The same identity holds for \(f\) by positive and negative
parts. Its inner period is constant in \(h\) on each \(h_tJ\)
by the right invariance. Each such coset has mass \(1/[K:J]\).
The displayed identity and the separate finiteness follow
since \(T\) is finite and its multiplicities there are positive.
\end{proof}

\begin{theorem}\label{thm:normalized-chain-finite-integration}
Use the ramified stage functions and the cuts \(T_i,J_i,K_i\)
of Theorem~\ref{thm:actual-ramified-period-cuts}, taking elementary
tests with the indicated finite inclusion vectors. Write
\(G_i={\rm U}(Y_i)\), \(B_i=B_{Y_i}\) and \([G_i]=G_i(E)\backslash
G_i(\mathbb A)\). Use product Haar measures on \(G_i(\mathbb A)\)
with \({\rm vol}(K_i)=1\) at the finite places and on the Bessel
groups as in Lemma~\ref{lem:bessel-product-probability}.
Consider continuous automorphic functions \(y_i\), rapidly
decreasing in the sense of Proposition~\ref{prop:unitary-reduced-height-integrability},
with
\[
      y_i(v)=[K_i:J_i]\int_{[G_{i+1}]}
         \theta\big(\Omega_i(v,h)f_i\big)y_{i+1}(h)\dd h
                                          \qquad(0\le i<m).
\]
Here \(\theta(f)=F_f(0)\) for the indicated Hom space and
\(\Omega_i\) is its arrow operator. At the definite terminal
space every continuous automorphic function meets the decay
condition. Fix any complex vector space \(M\), a vector \(V(t)\)
in \(M\) for each tuple \(t\in\prod_{i<m}T_i\), and a linear
functional \(L:M\to\C\). Write \(\xi=(\xi_0,\ldots,\xi_{m-1})\)
for infinite positions. Suppose a measurable scalar function
\(a_m(\xi)\) satisfies, for all such tuples and positions,
\[
       \int_{[G_m]} P_{m,\phi}(h)y_m(h)\dd h
                             =a_m(\xi)L(V(t)),
\]
where the preceding stage elements have infinite components
\(\xi_i\) and arbitrary fixed finite representatives of \(t_i\).
Let \(\mathcal P_\phi\) be the full rational periodization of
\(P_{0,\phi}\). There is a scalar \(\kappa\), depending only on
\(a_m\), the stages, measures and fixed real cutoffs, such that
\[
        \int_{[G_0]}\mathcal P_\phi(h)y_0(h)\dd h
          =\kappa L\left(\sum_{t\in\prod_{i<m}T_i}
                            \frac{V(t)}{\prod_{i<m}m_i(t_i)}\right).
\]
Thus \(\kappa\) can be kept the same across twists, columns
and terminal functionals whenever \(a_m\) is the same.
\end{theorem}
\begin{proof}
For a tuple of preceding positions define
\[
 I_i=\int_{B_i(E)\backslash G_i(\mathbb A)}
                              P_{i,\phi}(h)y_i(h)\dd h .
\]
At \(i=m\) the Bessel group is \(G_m\). For \(i>0\), absolute
convergence follows by applying Theorem~\ref{thm:rapid-actual-one-step-unfold}
to the preceding stage with fixed positions, ramified source
period and rapidly decreasing target \(y_i\). For \(i=0\) the
sum of absolute seed terms in the full periodization is bounded
uniformly by Proposition~\ref{prop:full-frame-majorant} as in
Theorem~\ref{thm:actual-full-poincare-columns}, after summing the
finitely many column coordinates. The fixed finite translation of
Lemma~\ref{lem:initial-frame-central-translation} preserves that
bound. Grouping the seed terms by rational Bessel cosets bounds
the sum of the absolute values of the corresponding translates
of \(P_{0,\phi}\) by the same constant. Since \(y_0\) is integrable
by Proposition~\ref{prop:unitary-reduced-height-integrability},
partitioning the Bessel quotient by translates of a full unitary
fundamental domain gives absolute convergence and the equality
of \(I_0\) with the left side in the theorem.

For every \(i<m\), the defining arrows show right \(J_i\)-invariance
of \(y_i\), because \(\Omega_i(j,1)\) fixes \(f_i\) for \(j\in J_i\)
by Theorem~\ref{thm:actual-twisted-finite-arrow}. The same theorem
on the preceding arrow gives right \(J_i\)-invariance of \(P_{i,\phi}\)
for \(i>0\), and at zero use the base right level. The indicated
integrals are consequently unchanged by the representative choices
in the preceding positions, where the two adjacent arrow tests both
have these level invariances. At zero the seed also has this
invariance. In factoring the Bessel product measure one can take
the finite measure of \(S_i\cap A_i\) to be one: divide the finite
factor by its positive finite volume and multiply the real factor
by the same number. The original product measure is preserved.
These normalizations depend on \(A_i\), independent of the twist.

Apply Lemma~\ref{lem:actual-period-cut-disintegration} to
\(f=P_{i,\phi}y_i\) using Theorem~\ref{thm:actual-ramified-period-cuts}
for support. Write \(\alpha_i\) for its real cutoff on \(G_i(\R)\),
with the product over the real embeddings understood. For each
new position \(x_i=(\xi_i,h_{t_i})\) in that formula, the inner
Bessel integral is
\[
      \int_{[B_i]} P_{i,\phi}(r x_i)y_i(r x_i)\dd r
                                =[K_i:J_i]\,I_{i+1}.
\]
Substitute the given formula for \(y_i\) and apply
Theorem~\ref{thm:rapid-actual-one-step-unfold}. The point
\(x_{0,f}\) selected from \(T_0\) lies in \(A^{\rm new}\) as
required at all the subsequent stages. Thus the factor
\([K_i:J_i]\) cancels exactly in the cut integral.

Proceed backward in \(i\). For each fixed preceding \(\xi_0,\ldots,
\xi_{i-1}\) put \(a_i=\int\alpha_i(\xi_i)a_{i+1}\dd\xi_i\)
if absolutely convergent, and put \(a_i=0\) otherwise. These are
measurable functions by Lemma~\ref{lem:product-integration},
using truncation of the integrands. Inductively the value \(I_i\)
is \(a_i\) times the sum of \(L(V(t))\) over all the remaining
indices with their remaining reciprocal multiplicities.
The terminal case is the hypothesis. For the induction step
insert the formula for \(I_{i+1}\) into the preceding cut integral.
For each \(t_i\), the coefficient of \(a_{i+1}\) is a finite
sum independent of \(\xi_i\). If any such coefficient is nonzero,
absolute convergence for that term, supplied by the lemma at
stage \(i\), proves \(\int\alpha_i|a_{i+1}|<\infty\).
If they all vanish, the desired equality holds with either
definition of \(a_i\). This proves the induction without any
change of order in the infinite real integrations.
At \(i=0\), \(a_0\) is a scalar depending only on the indicated
real functions and measures. Take \(\kappa=a_0\) and use
linearity of \(L\).
\end{proof}

\section{Polynomial tests for a positive hermitian oscillator}
\begin{definition}\label{def:positive-hermite-tests}
At one infinite place of a positive hermitian space of complex
dimension \(d\), take the Darboux coordinates \(x+i l\) of
Theorem~\ref{thm:negative-block-real-splitting}. On the real
position space \(\R^d\) put \(\phi_0(t)=2^{d/4}e^{-\pi |t|^2}\).
For \(1\le j\le d\) and any polynomial \(P\) in \(d\) variables put
\[
   C_j=\frac{2\pi t_j-\partial_{t_j}}{\sqrt{4\pi}},\qquad
   A_j=\frac{2\pi t_j+\partial_{t_j}}{\sqrt{4\pi}},\qquad
   H_P=P(C_1,\ldots,C_d)\phi_0 .
\]
Write \(H_\alpha=H_{Z^\alpha}\) for a nonnegative integer
multi-index \(\alpha\), \(|\alpha|=\sum_j\alpha_j\), and
\(\alpha!=\prod_j\alpha_j!\). Products of these constructions
are used at distinct infinite places or on two independent
oscillator spaces.
\end{definition}

\begin{theorem}\label{thm:positive-hermite-circle-projection}
The \(H_\alpha\) have inner products
\(\langle H_\alpha,H_\beta\rangle=\alpha!{\bf1}_{\alpha=\beta}\).
A Schwartz function orthogonal to all of them vanishes.
For \(u\in{\rm U}(d)\) the untwisted positive oscillator satisfies
\[
            \Omega_\infty(u)H_P=H_{P(u^{\mathsf T}Z)} .
\]
For every Schwartz function \(f(t,t')\), with any additional
real variables \(t'\), and every integer \(b\), its projection
under scalar rotations in the first factor is
\[
 \frac{1}{2\pi}\int_0^{2\pi} e^{-ib\theta}
             \big(\Omega_\infty(e^{i\theta}I)\otimes1\big)f
                                            \,\dd\theta
   =\sum_{|\alpha|=b} H_\alpha(t)\,
      \frac{\langle H_\alpha,f(\,\cdot\,,t')\rangle}{\alpha!}.
\]
The sum is taken to be zero for \(b<0\), and every coefficient
on the right is Schwartz in \(t'\).
\end{theorem}
\begin{proof}
Integration by parts on the Schwartz space gives \(C_j^*=A_j\).
The product rule gives \(A_iC_j-C_jA_i=\delta_{ij}\) and
the \(C_j\) commute with each other. Also \(A_j\phi_0=0\).
Moving the adjoints successively through the two products
in an inner product proves the displayed factorial formula:
a derivative beyond the available \(C_j\)'s annihilates the
vacuum, and the same applies with the arguments reversed.
The norm of \(\phi_0\) is one by the product Gaussian integral
of Lemma~\ref{lem:theta}.

For a polynomial \(Q(t)\), \(C_j(Q\phi_0)\) is \(\phi_0\)
times \(2\sqrt\pi\,t_jQ-\partial_jQ/\sqrt{4\pi}\).
Inducting on the total degree shows that the span of the
\(H_\alpha\) is the set of all polynomials times \(\phi_0\),
since the leading term corresponding to \(\alpha\) is a
nonzero multiple of \(t^\alpha\). If a Schwartz function
\(f_0\) is orthogonal to all these, all moments of
\(\phi_0 f_0\) vanish. At any real Fourier frequency expand
the exponential in the Fourier transform in its power series.
The absolute series inside the integral is bounded by
\(\phi_0(t)|f_0(t)|e^{c|t|}\) for a constant depending on the
frequency, which is integrable by completing the square.
Thus the Fourier transform vanishes. Inversion in
Lemma~\ref{lem:schroedinger-real-gaussian} and strict positivity
of \(\phi_0\) give \(f_0=0\).

In the symmetric gauge the two infinitesimal translations are
\(D_{x_j}=\partial_j,\ D_{l_j}=-2\pi i t_j\) by
Lemma~\ref{lem:schroedinger-real-gaussian}. Differentiating
the intertwining law is valid by its difference quotient proof.
Under \(u=U+iV\) the real blocks on \(x,l\) are \(U,-V,V,U\).
It follows that conjugation carries
\(D_{x_j}-iD_{l_j}\) to \(\sum_k u_{kj}(D_{x_k}-iD_{l_k})\).
Multiplying by \(-1/\sqrt{4\pi}\) gives the same transformation
on \(C_j\). The vacuum is fixed: its normalized coefficient is
one by Theorem~\ref{thm:negative-block-real-splitting} with
empty negative block, and the operator is unitary. This proves
the polynomial action.

For the circle integral, real Schwartz continuity in that theorem
gives convergence of Riemann sums in each Schwartz seminorm:
derivatives and polynomial multiples of the integrand are
uniformly continuous as functions of the angle in their
uniform norms. Their limits are the derivatives and polynomial
multiples of the pointwise integral, by differentiation under
the integral. The operators on the first block act with \(t'\)
as a parameter by the partial Fourier construction in
Theorem~\ref{thm:actual-adelic-rational-theta-invariance}.
At each fixed \(t'\), take the Hermite inner products of the
integral. Unitarity and the polynomial action multiply the
coefficient at \(\alpha\) by the mean of
\(e^{i(|\alpha|-b)\theta}\), equal to \({\bf1}_{|\alpha|=b}\)
by direct integration. The claimed finite sum has just these
coefficients and the orthogonality assertion gives equality.
Derivatives and polynomial multiples in \(t'\) of its coefficients
are integrals against a polynomial times a Gaussian of the
corresponding Schwartz derivatives, giving the Schwartz bounds.
\end{proof}

\begin{proposition}\label{prop:positive-hermite-theta-series}
For \(z\in\C^d\) define the Schwartz test
\[
            B_z(t)=\exp(2\sqrt\pi\, z^{\mathsf T}t
                               -z^{\mathsf T}z/2)\phi_0(t).
\]
It has Taylor coefficient \(H_\alpha/\alpha!\) at \(z^\alpha\).
Writing \(v=(x,l)\) and \(\zeta=\sqrt\pi(x+i l)\), the symmetric
real translation satisfies
\[
 \rho_a(v)\phi_0=e^{-|\zeta|^2/2} B_{-\zeta},
 \qquad
 \langle\phi_0,\rho_a(v)H_P\rangle
                     =e^{-|\zeta|^2/2}P(\bar\zeta).
\]
For a definite adelic Hom space take the product of the \(B_z\)
over its infinite places and any fixed finite test \(\Phi\).
For \(c\in\C\) the function of all its complex variables
\(\mathcal T_{c,\Phi}(z)=\theta(B_{cz}\otimes\Phi)\) is entire,
with coefficient \(c^{|\alpha|}\theta(H_\alpha\otimes\Phi)/\alpha!\).
In particular the value of \(P(\partial)\mathcal T_{c,\Phi}\)
at zero is \(c^b\theta(H_P\otimes\Phi)\) for homogeneous \(P\)
of total degree \(b\).
\end{proposition}
\begin{proof}
The defining exponential gives
\(\partial B_z/\partial z_j=(2\sqrt\pi\, t_j-z_j)B_z=C_jB_z\).
At \(z=0\) this yields the coefficients by iteration. For a
bounded complex \(z\)-set, all \(t\)- and \(z\)-derivatives are
bounded by a polynomial in \(|t|\) times \(e^{-\pi|t|^2+c_1|t|+c_2}\)
for fixed constants. These estimates justify complex
differentiation in every Schwartz seminorm by applying the
one-variable integral formula to a difference and its derivative.
Substitution of the symmetric translation of
Lemma~\ref{lem:schroedinger-real-gaussian} gives the first
translation formula after expanding \((t+x)^{\mathsf T}(t+x)\).
The inverse of \(\rho_a(v)\) is \(\rho_a(-v)\). Move it to the
first inner product variable and apply the Taylor coefficients
and orthogonality above. More explicitly,
\(\langle B_z,H_\alpha\rangle=\bar z^\alpha\) by differentiating
the exponential expansion under its Gaussian integral, whose
Taylor series converges on every bounded set by the same Gaussian
majorant. This gives the second formula for all \(P\).

The theta functional is continuous in the real Schwartz
seminorms for fixed finite test. In rational coordinates this
is the lattice bound of Proposition~\ref{prop:rational-frequency-lattice-bound}.
The changes from the real coordinates here to those coordinates
are continuous in Schwartz seminorms by their linear substitutions,
quadratic phases and partial Fourier transforms in
Theorem~\ref{thm:actual-adelic-rational-theta-invariance}.
Thus complex differentiation commutes with the functional.
The preceding coefficient calculation in all variables proves
the last two claims.
\end{proof}

\section{Rectangular substitutions and the definite weights}
\begin{definition}\label{def:rectangular-flag-matrix}
Fix \(1\le d<d'\) and nonnegative counts \(\kappa(h)\) of
columns of heights \(1\le h<d'\), zero for \(d<h<d'\).
For \(a=d,d'\), use the diagram with just its columns of height
\(<a\). Write \(b_a(Y)\) for the column of all its semistandard
minor products of Definition~\ref{def:compact-flag-space},
ordered in any fixed order, and \(E_a\) for the complex coordinate
space with those indices. An empty product is one. Write
\(c_0=\kappa(d)\), \(b_a^*=\sum_{h<a}h\kappa(h)\), and
\(\xi_a^0=b_a(I_a)\).
Use \(b_{d'}(Z)\) also for a \(d'\)-by-\(d\) matrix \(Z\):
all the selected minors use only the first \(d\) columns.
For square \(a\)-by-\(a\) matrices \(U\), define \(S_a(U)\)
by \(b_a(UY)=S_a(U)b_a(Y)\). Define the rectangular matrix
\(S(Z):E_d\to E_{d'}\) by
\[
            b_{d'}(ZY)=(\det Y)^{c_0} S(Z)b_d(Y)
                       \qquad (Y\text{ of size }d\text{ by }d).
\]
\end{definition}
\begin{proposition}\label{prop:rectangular-flag-law}
The matrices in Definition~\ref{def:rectangular-flag-matrix}
exist uniquely, have integer polynomial entries, and satisfy
\[
 S(A Z B)=S_{d'}(A) S(Z)(\det B)^{c_0} S_d(B),
                   \qquad b_{d'}(Z)=S(Z)\xi_d^0 .
\]
Here \(A,B\) are arbitrary square matrices of sizes \(d',d\).
One has \(S_a(AB)=S_a(A)S_a(B)\) for matrices of size \(a\).
Each entry of \(S(Z)\) is homogeneous of degree \(b_{d'}^*\).
For an arbitrary formal power series \(F\) in the entries of \(Z\),
the \((I',I)\) entry of \(S(\partial) F(0)\) is the coefficient
of \((\det Y)^{c_0}b_d(Y)_I\) in
\[
                  b_{d'}(\partial Y)_{I'}F(0).
\]
Here the entries of \(\partial\) are the commuting formal
derivatives with respect to the corresponding entries of \(Z\),
and \(Y\) consists of independent scalar variables.
\end{proposition}
\begin{proof}
Expand each minor of \(ZY\) by multilinearity in its intermediate
indices as in Proposition~\ref{prop:standard-flag-minors}. A
height-\(d\) minor contributes \(\det Y\) times the corresponding
maximal minor of \(Z\). For each smaller height the expansion
gives a sum of minors of \(Z\) times flag minors of \(Y\).
Multiply the expansions and straighten the remaining products of
\(Y\)-minors by Lemma~\ref{lem:flag-straighten}, with the matrix
transpose matching its index convention. This yields integer
polynomial coefficients and their claimed common degree.
Independence of the standard products, also after multiplying
them all by the same nonzero polynomial, gives uniqueness.
For square \(a\) this argument with its diagram below \(a\)
gives \(S_a\); successive substitutions and independence give
its multiplication law and \(S_a(I)=I\).

For the rectangular law apply the square substitution on the
left by \(A\). It holds for rectangular arrays as well since
the products here involve no columns beyond the \(d\)-th,
so one can append zero columns for this calculation.
Substitution of \(BY\) for \(Y\) then multiplies its determinant
by \(\det B\) and replaces \(b_d(Y)\) by \(S_d(B)b_d(Y)\).
Compare coefficients to get the first identity. Putting \(Y=I\)
in the definition gives the second. Finally the same identities
of polynomials hold after replacing every variable \(Z_{bi}\)
by \(\partial_{bi}\), since these operators commute, as can be
checked on formal monomials. Applying to \(F\) at zero gives the
displayed coefficient identity. The determinants in this expression
are exactly the column determinants of the differential construction
in Definition~\ref{def:linear-jet}, up to transposing each small
determinant, which preserves its value.
\end{proof}

\begin{proposition}\label{prop:actual-definite-arrow-weights}
For the definite ranks write \(d_j=k+j(t+1)\), \(0\le j\le m\),
where the last rank includes the extra line in
Proposition~\ref{prop:actual-pair-grade-zero}. At each place \(\tau\)
take the column counts of Definition~\ref{def:stage-diagrams}
on \(d_j\to d_{j+1}\). Set
\[
 e_{0,\tau}=0,\qquad e_{j,\tau}=\lfloor d_{j-1}/2\rfloor-w_j(\tau)
                                                    \quad(1\le j\le m).
\]
Take their tensor products \(S(Z)=\bigotimes_\tau S_\tau(Z_\tau)\)
and \(E_j=\bigotimes_\tau E_{d_j,\tau}\), and define the action
on \(E_j\) for \(g=(g_\tau)_\tau\) in the definite real group by
\[
       R_j(g)=\bigotimes_\tau
                   (\det\bar g_\tau)^{e_{j,\tau}}S_{d_j,\tau}(\bar g_\tau).
\]
For every \(A\in E_{j+1}^\vee,\ v\in E_j\), the polynomial
\(P_{A,v}(Z)=A S(Z)v\) gives the real Hermite test of
Definition~\ref{def:positive-hermite-tests}. For the actual twisted
definite arrow the transformation is
\[
 \Omega_{j+1,\infty}^0(g,h) H_{P_{A,v}}
                = H_{P_{\,A R_{j+1}(h)^{-1},\,R_j(g)v}} .
\]
At stage zero \(E_0,R_0\) identify with the column representation
of Definition~\ref{def:weyl-coord}.
\end{proposition}
\begin{proof}
At every sufficiently large arithmetic scale one has \(k>1\)
and \(h_{\rm gap}>t+1\), where \(h_{\rm gap}=2r^8\) denotes
the integer called \(h\) in Definition~\ref{def:heights}, by
Lemma~\ref{lem:diagram-size}. At \(j=0\), the new full-height
count \(c_\tau=\kappa_\tau(d_j)\) is
\(w_1(\tau)-\lfloor d_1/2\rfloor\).
At \(j\ge1\) it is \(h_{\rm gap}-(t+1)\). Since
\(d_{j+1}-d_{j-1}=2(t+1)\) and \(w_{j+1}-w_j=h_{\rm gap}\),
in both cases
\[
 e_{j,\tau}=\lfloor d_{j+1}/2\rfloor-w_{j+1}(\tau)+c_\tau,\qquad
 e_{j+1,\tau}=\lfloor d_j/2\rfloor-w_{j+1}(\tau).
\]
There are no columns with an intermediate height, so the
preceding proposition applies. At \(j=0\) its evaluation
coordinates are the monomials of degree \(a{\bf1}_A(\tau)\)
in the first column, exactly those in Definition~\ref{def:weyl-coord}.

On the positive Hom space an orthonormal matrix unit
\(E_{bi}\) transforms to \(h E_{bi}g^{-1}\).
The substitution \(u^{\mathsf T}Z\) in
Theorem~\ref{thm:positive-hermite-circle-projection}, written
with the two matrix indices, is thus \(h^{\mathsf T} Z\bar g\)
at each infinite place, since the definite \(g\) is unitary.
The additional character of Theorem~\ref{thm:global-actual-arrow-covariance}
multiplies by
\[
           \prod_\tau
          (\det h_\tau)^{\lfloor d_j/2\rfloor-w_{j+1}(\tau)}
          (\det g_\tau)^{-\lfloor d_{j+1}/2\rfloor+w_{j+1}(\tau)}
\]
by Definition~\ref{def:actual-twisted-arrow-character} and
Proposition~\ref{prop:actual-marked-torus-characters}.
Now use Proposition~\ref{prop:rectangular-flag-law} on
\(S(h^{\mathsf T}Z\bar g)\). Its full-height factor contributes
\((\det\bar g_\tau)^{c_\tau}\), and \(h_\tau^{\mathsf T}=\bar h_\tau^{-1}\).
Combining with the two exponent identities gives exactly the
stated polynomial action.
\end{proof}

\section{The prescribed source representation in the definite oscillator}
\begin{lemma}\label{lem:rectangular-borel-polynomials}
With \(d<d'\) and \(\kappa\) of Definition~\ref{def:rectangular-flag-matrix},
suppose a polynomial \(p\) on \(d'\)-by-\(d\) matrices satisfies
\[
 p(Zb)=\prod_{h\le d}(b_{11}\cdots b_{hh})^{\kappa(h)}p(Z)
\]
for every invertible upper triangular \(d\)-by-\(d\) matrix \(b\).
There exists a unique row \(A\in E_{d'}^\vee\) with
\(p(Z)=A b_{d'}(Z)\). For a product of matrix spaces at several
places the analogous independent covariance in each argument
gives the tensor product of these row spaces.
\end{lemma}
\begin{proof}
Extend \(p\) to square \(d'\)-by-\(d'\) matrices by using just
the first \(d\) columns. On the right upper triangular group
of size \(d'\) it transforms according to the same weight,
since the first \(d\) columns of a product with an upper
triangular matrix depend only on the first \(d\) columns.
Let \(\mathcal D\) be the space of all square-matrix polynomials
with this upper triangular transformation. Right multiplication
by a scalar implies that each is homogeneous of degree
\(\sum_{h\le d}h\kappa(h)\), by comparing powers of the scalar.
Thus \(\mathcal D\) is finite-dimensional. It is stable under
the action \(f(X)\mapsto f(u^{\mathsf T}X)\).
The span of the coordinates of \(b_{d'}\) belongs to \(\mathcal D\)
and is invariant by Proposition~\ref{prop:standard-flag-minors}.
The entries are independent by that same proposition.

Give the homogeneous polynomials the Hermitian inner product
obtained by integrating \(\bar f g\) against the positive
Gaussian density \(e^{-{\rm tr}(X^*X)}\) on all complex matrices.
All polynomial moments are finite by elementary real Gaussian
integration as in Lemma~\ref{lem:theta}, and a nonzero polynomial
is nonzero on an open set, so this form is positive definite.
It is invariant under the indicated left substitutions for
\(u\in{\rm U}(d')\). If the stated span were proper, its orthogonal
complement inside \(\mathcal D\) would therefore be nonzero
and compact-group invariant. It is invariant under the full
complex group by the differentiation and elementary shear
argument in Theorem~\ref{thm:compact-highest-vector}, which
applies to this finite polynomial substitution representation.
As in that theorem, take a lexicographically largest diagonal
weight occurring in the complement. Diagonal interpolation
selects a vector in that weight space, and every strictly upper
elementary root raises the weight. Thus the upper unitriangular
group fixes this vector under the left transpose action.

Any vector in \(\mathcal D\) fixed under that action is determined
by its value at \(I\). On the dense set \(X=l b\), where \(l\)
is lower unitriangular and \(b\) upper triangular invertible,
its value is the given diagonal character of \(b\) times its
value at \(I\). Existence and density of this factorization
are proved by elimination of leading principal minors in
Theorem~\ref{thm:compact-highest-vector}. The product of the
leading minors of heights \(h\) with multiplicities \(\kappa(h)\)
belongs to the stated span, is fixed on the left by lower
unitriangular substitutions, and is one at \(I\). A second
fixed vector in the orthogonal complement is impossible.
It follows that \(\mathcal D\) equals the span. Restricting
back to the first \(d\) columns, which suffice for every minor
in it, proves existence and uniqueness. For multiple places
expand in the monomials of all but the first matrix. Their
coefficients satisfy the first-place law and hence expand in
its independent minor products. The remaining coefficient
polynomials satisfy the other laws by that independence.
Iteration proves the product assertion.
\end{proof}

\begin{theorem}\label{thm:definite-source-oscillator-multiplicity}
Fix an actual definite arrow \(X_j\to X_{j+1}\) and its
prescribed \(R_j\) of Proposition~\ref{prop:actual-definite-arrow-weights}.
Let \(L\) be a complex linear map from \(E_j\) to the real
Schwartz space of the arrow's position variables, with optional
additional real variables \(t'\) acted on trivially. Suppose
\[
                  L(R_j(g)v)=\Omega_{j+1,\infty}^0(g,1)L(v)
                 \qquad(g\in{\rm U}(X_j)(\R),\ v\in E_j).
\]
Here and below the real unitary group includes all archimedean
factors. There is a unique row \(A(t')\in E_{j+1}^\vee\)
of Schwartz functions such that
\[
                        L(v)(\,\cdot\,,t')=H_{P_{A(t'),v}}
                              \qquad(v\in E_j).
\]
With no additional variables \(A\) is an ordinary constant row.
Conversely all these maps satisfy the displayed intertwining law.
\end{theorem}
\begin{proof}
Fix a real place, write \(d=d_j,d'=d_{j+1}\), and let
\(e_b=\lfloor d'/2\rfloor-w_{j+1}(\tau)\). For a scalar
\(g_\tau=\eta I_d\) of modulus one, the action on \(E_j\)
is \(\eta^{-b_d^*-d e_{j,\tau}}\), by homogeneity of the
basis products. Its action on Hermite degree \(a\) is
\(\eta^{-a-d e_b}\) by Theorem~\ref{thm:positive-hermite-circle-projection}
and the arrow character in Proposition~\ref{prop:actual-definite-arrow-weights}.
The assumed equation and the circle projection of that theorem
therefore force the degree in this factor to be
\[
           a=b_d^*+d(e_{j,\tau}-e_b)=b_d^*+d\kappa_\tau(d)=b_{d'}^* .
\]
Apply this successively at every place, regarding all other
variables as additional coordinates. It follows that
\(L(v)(\,\cdot\,,t')=H_{R(Z;t')v}\) for a uniquely determined
row \(R\) of polynomials in the matrices \(Z_\tau\), with
Schwartz coefficients in \(t'\). Indeed there are only finitely
many allowed multidegrees, and application to a basis of \(E_j\)
determines each entry.

Fix \(t'\), suppress it, and act in one place at a time.
The Hermite transformation law and its injectivity from
Theorem~\ref{thm:positive-hermite-circle-projection} yield
\[
        R(Zu)=(\det u)^{\kappa_\tau(d)}R(Z) S_d(u),
                                \qquad u\in{\rm U}(d).
\]
Here \(u=\bar g_\tau\), and the square action in the right hand
side is on its own tensor factor. The identity extends to all
\(u\in{\rm GL}_d(\C)\). For skew Hermitian \(D\), substitute
\(\exp(tD)\). Both sides are entire in \(t\) and agree for real
\(t\), hence for complex \(t\) by Lemma~\ref{lem:analytic-tools}.
This gives it on positive Hermitian invertible matrices by
diagonalizing them and taking the logarithms of their positive
eigenvalues. Every invertible matrix is a unitary times such
a matrix by the diagonalization of \(u^*u\) used in
Theorem~\ref{thm:hyperbolic-adelic-parabolic-splitting}.
The multiplication law of Proposition~\ref{prop:rectangular-flag-law}
gives the extension.

Set \(\xi^0=\bigotimes_\tau b_{d_j}(I)\).
For an invertible upper triangular \(b\) we have
\[
 S_d(b)b_d(I)=b_d(b)
      =\prod_{h<d}(b_{11}\cdots b_{hh})^{\kappa_\tau(h)}b_d(I).
\]
Indeed the only nonzero minor in the first \(h\) columns of an
upper triangular matrix with \(h\) chosen rows uses exactly
the first \(h\) rows. Hence \(p(Z)=R(Z)\xi^0\) satisfies all
the hypotheses of Lemma~\ref{lem:rectangular-borel-polynomials}.
It equals \(A b_{d'}(Z)\) in the tensor product notation, for a
unique row \(A\). By Proposition~\ref{prop:rectangular-flag-law}
this is \(A S(Z)\xi^0\). The vectors
\(\bigotimes_\tau S_d(u_\tau)b_d(I)\) for invertible \(u_\tau\)
span \(E_j\): a linear relation against all these evaluation
vectors would be a linear relation among independent products
of the basis minors, by Proposition~\ref{prop:standard-flag-minors}
and its polynomial continuation to singular matrices.
Covariance under all these substitutions now gives \(R(Z)=A S(Z)\)
as rows on \(E_j\). Finally the coefficients of \(A\) are fixed
linear combinations of the finitely many coefficients of \(p\),
by solving in the independent minor basis. Restoring \(t'\)
shows that they are Schwartz. Uniqueness and the converse follow
from the same evaluation and Proposition~\ref{prop:actual-definite-arrow-weights}.
\end{proof}

\section{The finite class sum of the definite differential kernel}
\begin{definition}\label{def:definite-jet-class-operator}
For \(0\le j\le m\) let \(H_j\) be the unitary group of the
positive hermitian space \(X_j\) of
Proposition~\ref{prop:actual-definite-arrow-weights}. Write
\(H_{j,f}\) for its finite adelic group, \(H_{j,\infty}\) for its
product of real groups and \(\Gamma_j\) for its rational points
over \(E^+\). A complex algebraic modular form of weight \(R_j\)
is a function \(\phi:H_{j,f}\to E_j\) fixed by right translation
under some compact open subgroup and satisfying
\(\phi(\gamma_f g)=R_j(\gamma_\infty)\phi(g)\) for
\(\gamma\in\Gamma_j\). Write \(r_a\phi(g)=\phi(ga)\).
At \(j<m\) normalize finite Haar measure by
\({\rm vol}(K_{c,j})=1\), using the column level in
Proposition~\ref{prop:column-level-in-stage-bessel}.

Fix \(c\in\C^\times\), independent of the finite twist. Put
\(N_j=\sum_{\tau}\sum_{h<d_{j+1}}h\kappa_\tau(h)\), the total
degree of the rectangular matrix \(S(Z)\) on this arrow. For a
finite oscillator test \(\Phi\) define \(J_j^c(\Phi):E_j\to E_{j+1}\)
by
\[
       A J_j^c(\Phi)v
         =\big[P_{A,v}(\partial)\mathcal T_{c,\Phi}\big](0)
       \qquad(A\in E_{j+1}^{\vee},\ v\in E_j).
\]
Here \(P_{A,v}=A S(Z)v\) and \(\mathcal T_{c,\Phi}\) is the entire
function in Proposition~\ref{prop:positive-hermite-theta-series}.
Bilinearity in \(A,v\) defines the map by taking coordinate rows.
Denote the finite factor of the twisted definite arrow
by \(\Omega_f(g,h)\). Its finite class operator is
\[
 (\mathcal L_{j,\Phi}^c\phi)(h)=
 \int_{\Gamma_j\backslash H_{j,f}}
 J_j^c\big(\Omega_f(g,h)\Phi\big)\phi(g)\dd g,\qquad h\in H_{j+1,f}.
\]
\end{definition}

\begin{lemma}\label{lem:definite-class-quotient-measures}
For any of these definite groups \(H\), its rational points
\(\Gamma\) embed discretely in \(H_f\), and
\(\Gamma\backslash H_f/J\) is finite for every compact open
\(J\subset H_f\). Given a compact open \(K\) of volume one and
\(J\subset K\), a right \(J\)-invariant, left \(\Gamma\)-invariant
function \(D\) with values in a finite-dimensional complex vector
space satisfies
\[
 \int_{\Gamma\backslash H_f}D(g)\dd g
  =\frac1{[K:J]}\sum_{[g]\in\Gamma\backslash H_f/J}
             \frac{D(g)}{|\Gamma\cap gJg^{-1}|}.
\]
The quotient measure is right invariant. For the full adelic
Haar measure used on \([H]=\Gamma\backslash (H_\infty\times H_f)\)
there is a number \(b_H>0\), depending only on that measure and
on \(K\), with
\[
 \int_{[H]}F
    = b_H \int_{\Gamma\backslash H_f}
              \int_{H_\infty}F(u,g)\dd u\,\dd g .
\]
The inner measure is probability Haar measure. This last identity
holds for every continuous left rational invariant function \(F\).
\end{lemma}
\begin{proof}
For a compact subset of \(H_f\), clear a common finite denominator
of its matrix entries. All infinite coordinates of a rational
unitary lie in the compact group \(H_\infty\). Bounded sets in
the resulting fractional integer lattice contain only finitely
many elements by Lemma~\ref{lem:trace-finite}. The rational
group therefore meets compact subsets in finite sets, and its
projection to the finite matrices is injective by injectivity
of the completion of \(E\). This gives discreteness. The same
lemma gives finitely many double cosets modulo the full integral
compact \(K_0\). The intersection \(K_0\cap J\) is open in
\(K_0\), and the cosets covering the compact \(K_0\) admit a finite
subcover. Passing first to this intersection and then to \(J\)
gives the asserted finiteness for \(J\).

Construct the left Haar measure as in
Proposition~\ref{prop:all-bessel-quotients-compact}, and rescale
it to give \(K\) volume one. Over the class of \(gJ\), left
translation identifies the quotient with
\((g^{-1}\Gamma g\cap J)\backslash J\). This is the free left
action of a finite group of size \(|\Gamma\cap gJg^{-1}|\).
Choose a measurable representative for each orbit by taking its
first element in a fixed lexicographic order on the countable
separating residue coordinates of \(J\). Two distinct elements
are separated by a finite congruence coordinate, so comparisons
are Borel sets. All translates of the representative set have
equal measure and partition \(J\). Its measure is
\({\rm vol}(J)/|\Gamma\cap gJg^{-1}|\), with
\({\rm vol}(J)=1/[K:J]\) by the left coset partition. These
representative pieces together give a measurable fundamental
domain \(D_f\) with compact closure and finite positive measure.
Two fundamental domains give the same integral of a left
invariant function, by partition into their countably many
intersecting rational translates as in Lemma~\ref{lem:quotient-measure}.

For right translation by \(a\in H_f\), the measure \(B\mapsto
{\rm vol}(Ba)\) is another left invariant measure. It is a
constant multiple of the original one. Indeed primitives of
the form \(hJ'\), for open subgroups \(J'\subset K\) of finite
index, have measures determined by the respective masses of \(K\).
Principal congruence subgroups in a countable nested basis
suffice. Open sets are countable unions of such cosets and
decompose into disjoint such cosets by successive refinement.
Agreement up to scale extends to Borel measures by the
construction from these basic sets, or by applying the
increasing and dominated limits of Lemma~\ref{lem:integral-tools}
to the generating algebra on each compact coset.
Now \(D_fa\) is another fundamental domain. Its mass equals
that of \(D_f\), forcing the constant to be one. This also
proves right invariance on the quotient.

The archimedean group is compact, so its positive smooth density in
Lemma~\ref{lem:quotient-measure} normalizes to a probability.
The full definite Haar measure is a positive constant times
its product with the finite measure, by the product construction
in Proposition~\ref{prop:all-bessel-quotients-compact}.
A fundamental domain for the full rational action is
\(H_\infty\times D_f\), since the finite representative is
unique and its stabilizer under left multiplication at a point
is trivial. Integration on this domain gives the asserted
formula. Invariance of the real probability makes the inner
integral left rational invariant as a function of the finite
variable. All continuous integrands here are bounded on the
closure of the domain.
\end{proof}

\begin{theorem}\label{thm:definite-jet-class-operator}
The operator in Definition~\ref{def:definite-jet-class-operator}
is bilinear in \(\Phi,\phi\) and produces a modular form of
weight \(R_{j+1}\). For every pair of finite elements \(a,b\)
on the source and target,
\[
 \mathcal L^c_{j,\Omega_f(a,b)\Phi}(r_a\phi)
                         =r_b(\mathcal L^c_{j,\Phi}\phi).
\]
For \(J\subset K_{c,j}\) compact open fixing \(\phi\) on the
right and fixing \(\Phi\) in the source action, its value at
\(h\) is the finite sum in Lemma~\ref{lem:definite-class-quotient-measures}
with \(D(g)=J_j^c(\Omega_f(g,h)\Phi)\phi(g)\).
\end{theorem}
\begin{proof}
By Proposition~\ref{prop:rectangular-flag-law} and homogeneity,
the defining formula for the jet is
\(A J_j^c(\Phi)v=c^{N_j}\theta(H_{P_{A,v}}\otimes\Phi)\),
by Proposition~\ref{prop:positive-hermite-theta-series}. For
rational points \(\gamma,\delta\) on source and target, apply
Theorem~\ref{thm:global-actual-arrow-covariance} to this theta
value. The real transformation in
Proposition~\ref{prop:actual-definite-arrow-weights} and cancellation
of the arbitrary row yield
\[
 J_j^c\big(\Omega_f(\gamma_f,\delta_f)\Phi\big)R_j(\gamma_\infty)
                         =R_{j+1}(\delta_\infty)J_j^c(\Phi).
\]
In particular the integrand is left invariant in \(g\) and
the integral transforms on the left as prescribed in \(h\).
The smoothness of the finite operators in
Theorem~\ref{thm:actual-adelic-rational-theta-invariance} gives
compact open subgroups on each side fixing the finite test.
Equivalently, use the finite basis of translated characteristic
functions: the sufficiently small congruence subgroups fix the
vacuum, each of the finitely many translations modulo its open
period subgroup, and the finitely many accompanying phases.
The twisting characters are locally constant on the finite
group by Proposition~\ref{prop:actual-marked-torus-characters}.
Intersect on the source with the right stabilizer of \(\phi\)
and the column level. This gives the asserted \(J\) and a
right invariance subgroup on the target. The lemma now gives
the finite sum and hence the existence and bilinearity of the
operator. Finally the integrand with translated pair and form
is \(J_j^c(\Omega_f(ga,hb)\Phi)\phi(ga)\). Right invariance of
quotient measure proves the translation identity.
\end{proof}

\section{Matrix coefficients of a definite theta lift}
\begin{definition}\label{def:definite-form-coefficient}
For the groups and weights of Definition~\ref{def:definite-jet-class-operator},
a row \(\lambda\in E_j^\vee\) and a modular form \(\phi\) define
\[
   F_{\phi,\lambda}(u,g)=\lambda\big(R_j(u)^{-1}\phi(g)\big)
       \qquad(u\in H_{j,\infty},\ g\in H_{j,f}).
\]
This is a continuous function on \([H_j]\). Indeed a simultaneous
left translation by \(\gamma\in\Gamma_j\) contributes
\(R_j(\gamma_\infty)^{-1}R_j(\gamma_\infty)\) between the two
factors. Smoothness at the finite group and continuity of the
matrix representation at infinity give continuity.
\end{definition}

\begin{lemma}\label{lem:definite-real-coefficient-mean}
Let \(f(t,t')\) be a real-place Schwartz test on the \(j\)-th
definite arrow with possible additional real variables \(t'\).
Use probability Haar measure on \(H_{j,\infty}\). Given a row
\(\lambda\in E_j^\vee\), there is a unique row
\(A(t')\in E_{j+1}^\vee\) of Schwartz functions such that
\[
 \int_{H_{j,\infty}}\!
       \lambda\big(R_j(u)^{-1}v\big)
                 \Omega_\infty(u,1)f(\,\cdot\,,t')\dd u
                         = H_{P_{A(t'),v}}\qquad(v\in E_j).
\]
Without the additional variables the row has constant entries.
The row is bilinear in \(f,\lambda\), continuous for the
Schwartz seminorms, and independent of every finite twist
with the prescribed infinite weight.
\end{lemma}
\begin{proof}
For each fixed \(v\), the integral exists as a Schwartz function.
For the first variables the action is a product of the real
oscillator actions of Theorem~\ref{thm:negative-block-real-splitting}
and of the determinant characters in
Proposition~\ref{prop:actual-definite-arrow-weights}. It acts
only on the first block in a joint Schwartz function as in
Theorem~\ref{thm:positive-hermite-circle-projection}. In every
Schwartz seminorm the integral is a limit of Riemann sums.
Indeed the compact group admits finite measurable partitions of
arbitrarily small diameter by restricting finite ball covers,
and continuity gives uniform approximation of the integrand in
that seminorm. Equivalently, differentiated integrals have the
bounds obtained by integrating the seminorms of the integrands.
Joint continuity of the action gives a finite family of input
seminorms bounding each output seminorm uniformly on the compact
group. To see this, at each \((u,0)\) take a product neighborhood
mapped into the unit ball of that output seminorm, take a finite
subcover in \(u\), intersect the finitely many neighborhoods of
zero and use linearity. The same argument on joint Schwartz
functions follows from the partial linear substitutions, phases
and Fourier transforms in Theorem~\ref{thm:actual-adelic-rational-theta-invariance};
polynomial factors and derivatives in \(t'\) commute through them.
This proves continuity of the average.

Denote the resulting linear map in \(v\) by \(L\). For fixed
\(b\in H_{j,\infty}\), the substitution \(u'=bu\) gives
\(\Omega_\infty(b,1)L(v)=L(R_j(b)v)\). Theorem~\ref{thm:definite-source-oscillator-multiplicity}
therefore gives the stated row and its uniqueness. Its proof
recovers the coefficients by finitely many integrals against
fixed Hermite functions followed by a fixed finite system of
linear equations. These operations are continuous in every
Schwartz seminorm of the additional variables. The averaging
formula is bilinear, and the operators at infinity use only
the common infinite weight, proving the remaining assertions.
\end{proof}

\begin{theorem}\label{thm:definite-adelic-class-lift}
Let \(\Omega\) be any twisted definite arrow at stage \(j<m\) with
the prescribed infinite weight, and let \(f\) be a real-place
Schwartz test on this arrow. For every \(\lambda\in E_j^\vee\)
there is a row \(\lambda'\in E_{j+1}^\vee\) such that
\[
 \int_{[H_j]} F_{\phi,\lambda}(g)
             \theta\big(\Omega(g,h)(f\otimes\Phi)\big)\dd g
      = F_{\mathcal L^c_{j,\Phi}\phi,\,\lambda'}(h)
                 \qquad(h\in H_{j+1}(\mathbb A)).
\]
The same row works for all modular forms \(\phi\), all finite
data \(\Phi\), and all finite twists of this infinite weight.
More precisely \(\lambda'=b_{H_j}c^{-N_j}A\), with \(A\) given
by Lemma~\ref{lem:definite-real-coefficient-mean}.
\end{theorem}
\begin{proof}
The integrand is rational invariant by the definition of the
coefficient and Theorem~\ref{thm:global-actual-arrow-covariance}.
It is continuous by real Schwartz continuity, finite local
constancy and the lattice bound in
Proposition~\ref{prop:rational-frequency-lattice-bound}. Its
absolute value is integrable on the compact representatives of
Lemma~\ref{lem:definite-class-quotient-measures}. Write
\(g=(u,a)\), \(h=(w,b)\). At fixed finite \(a,b\), integrate
over \(u\) using the product measure formula of that lemma.
In the real action \(\Omega_\infty(1,w)\) commutes with
\(\Omega_\infty(u,1)\). Both this operator and theta at the
fixed finite test \(\Omega_f(a,b)\Phi\) are continuous in
the real Schwartz topology, by the same estimates. They
therefore commute with the Riemann sums and the integral
in Lemma~\ref{lem:definite-real-coefficient-mean}, with
\(v=\phi(a)\). By Proposition~\ref{prop:actual-definite-arrow-weights}
the result of the real integration is
\[
 \theta\big(H_{P_{A R_{j+1}(w)^{-1},\,\phi(a)}}
                           \otimes\Omega_f(a,b)\Phi\big)
       =c^{-N_j} A R_{j+1}(w)^{-1}
                   J_j^c\big(\Omega_f(a,b)\Phi\big)\phi(a).
\]
The equality is the homogeneous derivative identity in
Theorem~\ref{thm:definite-jet-class-operator}. Integrate
in \(a\) and multiply by \(b_{H_j}\). The row and the real
representation in the right side are constant in \(a\).
The definition of \(\mathcal L^c_{j,\Phi}\) gives the claimed
coefficient at \(h\). The formula for \(\lambda'\) and the
last assertion of the lemma give the indicated uniformity.
\end{proof}

\section{Composition on two consecutive definite arrows}
\begin{definition}\label{def:paired-definite-class-operator}
At stages \(j,j+1\), put
\[
 \mathcal L^c_{j:j+2,\,\Phi_1\otimes\Phi_2}\phi
       =\mathcal L^c_{j+1,\Phi_2}\big(\mathcal L^c_{j,\Phi_1}\phi\big).
\]
Extend linearly in the algebraic tensor product of the two finite
test spaces. This prescription is well defined by bilinearity in
Theorem~\ref{thm:definite-jet-class-operator}. On a joint real
Schwartz function and a finite tensor write \(\theta_{12}\) for
the functional of Corollary~\ref{cor:actual-reduced-definite-pair}.
\end{definition}

\begin{theorem}\label{thm:double-definite-class-lift}
Fix any \(j\) with \(j+2\le m\), a real Schwartz test \(f\) on the
product of the two positive oscillator spaces, a scale \(c\ne0\)
and \(\lambda\in E_j^\vee\). There is a row \(\lambda''\in
E_{j+2}^\vee\), bilinear in \(f,\lambda\) and continuous in \(f\),
such that for every finite tensor \(y\), every modular form
\(\phi\) of weight \(R_j\), and every \(k\in H_{j+2}(\mathbb A)\),
\[
\begin{aligned}
 F_{\mathcal L^c_{j:j+2,y}\phi,\,\lambda''}(k)
  ={}&\int_{[H_j]}\int_{[H_{j+1}]}F_{\phi,\lambda}(g)\\
 &{}\times\theta_{12}\!\left(
    \big(\Omega_1(g,h)\otimes\Omega_2(h,k)\big)(f\otimes y)\right)
                                                 \dd h\,\dd g .
\end{aligned}
\]
The row depends on the fixed infinite weights and measures
and works simultaneously for all their finite twists.
\end{theorem}
\begin{proof}
Linearity reduces the proof to \(y=\Phi_1\otimes\Phi_2\).
Both groups have compact fundamental representatives by
Lemma~\ref{lem:definite-class-quotient-measures}. Simultaneous
rational invariance in the middle group follows by applying
Theorem~\ref{thm:global-actual-arrow-covariance} on each arrow.
For joint real tests this follows as well from the rational
sum in the two position variables of
Corollary~\ref{cor:actual-reduced-definite-pair}. The integrand
is continuous on the closure of the representatives by the
lattice bound of Proposition~\ref{prop:rational-frequency-lattice-bound}
and local constancy of the finite tests. It is absolutely
integrable, so Lemma~\ref{lem:product-integration} allows us
to integrate first in \(g\).

Write \(h=(w,b)\). Average in the infinite component of \(g\),
keeping the second test variables as the additional variables
in Lemma~\ref{lem:definite-real-coefficient-mean}. Choose any
basis \(\ell_1,\ldots,\ell_D\) of \(E_{j+1}^\vee\). Its
unique row of functions has the form
\(A(t')=\sum_{\nu=1}^D f_\nu(t')\ell_\nu\), with each \(f_\nu\)
a real Schwartz test on the second arrow. Thus the averaged
real vector for the finite argument \(v=\phi(g_f)\) is
\[
           \sum_{\nu=1}^D H_{P_{\ell_\nu,v}}\otimes f_\nu .
\]
All finite operators, the real target operator on the first
arrow, and the second arrow operator commute through this
first real integration. Their actions and the functional are
continuous on joint Schwartz tests, by the same argument as
in Theorem~\ref{thm:definite-adelic-class-lift}. The expression
now contains only a finite sum of product tests, on which
\(\theta_{12}\) is the product of the two theta functionals.

On its first factor the real target \(w\) replaces \(\ell_\nu\)
by \(\ell_\nu R_{j+1}(w)^{-1}\), by
Proposition~\ref{prop:actual-definite-arrow-weights}. The
calculation of the finite integral in
Theorem~\ref{thm:definite-adelic-class-lift} now gives
\[
 b_{H_j}c^{-N_j}\sum_{\nu=1}^D
    F_{\mathcal L^c_{j,\Phi_1}\phi,\,\ell_\nu}(h)
                       \,\theta(\Omega_2(h,k)(f_\nu\otimes\Phi_2))
\]
after integrating over \([H_j]\). Integrate each summand over
\([H_{j+1}]\) using the same theorem with \(f_\nu,\ell_\nu\).
If its row is \(\lambda'_\nu\), take
\(\lambda''=b_{H_j}c^{-N_j}\sum_\nu\lambda'_\nu\).
Construction by these two averages shows bilinearity and
continuity by Lemma~\ref{lem:definite-real-coefficient-mean}.
The row construction uses neither the finite tensor, the form,
the twist nor the target point. It therefore proves the result
on their stated ranges, including arbitrary finite tensor sums.
\end{proof}

\begin{corollary}\label{cor:even-unfolded-algebraic-coefficients}
Equip all real tests in the successive periods of
Theorem~\ref{thm:actual-unfolded-stage-functions} with fixed
infinite data. Write \(x_j=(\xi_j,a_j)\) for the stage positions,
with \(a_0\in A^{\rm new}\). Let \(0\le 2i\le m\).
For each of the \(i\) pairs, take its evaluated finite tensor
\(y_j\) in Corollary~\ref{cor:actual-reduced-definite-pair},
where \(j=0,2,\ldots,2i-2\). For a column \(\phi\) put
\[
    \psi_0=\phi,\qquad
    \psi_{j+2}=\mathcal L^c_{j:j+2,\,y_j}\psi_j .
\]
Then there is a row \(t_{2i}\in E_{2i}^{\vee}\), depending only
on the real positions through \(\xi_{2i}\), the chosen real
tests, \(c\), and the fixed models and measures, such that
\[
        P_{2i,\phi}(\widehat k x_{2i})
                         =F_{\psi_{2i},t_{2i}}(k_0).
\]
For \(2i<m\), \(k\) runs over the adelic definite group on
\(X_{2i}\), \(\widehat k\) extends it by the identity on the
string and \(k_0=k\). At \(2i=m\), \(k\) runs over the adelic
unitary group on \(Y_m\), \(\widehat k=k\), and
\(k_0=k\oplus1\) on \(X_m\). Each row can be chosen continuously
in the real positions, independently of the finite positions,
twist and column.
\end{corollary}
\begin{proof}
At zero take the row \(v\mapsto\mathcal B(v,F_\infty(\xi_0))\)
of Lemma~\ref{lem:initial-bessel-column-coefficient}. Its
invariant pairing moves \(\pi(k_\infty)\) in that lemma to
\(R_0(k_\infty)^{-1}\) on the first argument. This gives
Definition~\ref{def:definite-form-coefficient} and proves the
initial case. The base's finite position hypothesis is exactly
\(a_0\in A^{\rm new}\).

Suppose the formula holds at an even \(j<m\). In the definite
pair formula of Corollary~\ref{cor:actual-reduced-definite-pair}
substitute \(F_{\psi_j,t_j}(g)\) for \(p(gx_j)\). The test is
a joint real Schwartz function times the evaluated finite
tensor \(y_j\). A common nonzero scalar for the finite model
identification, if its normalization requires one, can be
included in the real test. Theorem~\ref{thm:double-definite-class-lift}
evaluates this period with \(\phi=\psi_j,\lambda=t_j\) as
the asserted coefficient at \(k_0\). The target embedding
including the last line is precisely that in
Proposition~\ref{prop:paired-restriction-determinant}.
Real translation, restriction to a coordinate subspace,
linear substitutions and partial Fourier transform in the
pair construction are continuous on Schwartz functions by
Theorems~\ref{thm:global-two-arrow-adjoint}
and~\ref{thm:actual-adelic-rational-theta-invariance}.
The continuous dependence and twist independence of the row
therefore follow at the next stage from the theorem. This
proves the claim by induction.
\end{proof}

\section{An integral operator with image the flag representation}
\begin{definition}\label{def:flag-tensor-coordinates}
Let \(V=\C^n\) with its ordered orthonormal basis. A diagram
\(\alpha\) is a list \(h_1\ge\cdots\ge h_a\ge1\) of column
heights at most \(n\), possibly empty. Put \(L=\sum_i h_i\),
\(T_\alpha(V)=\bigotimes_i\bigwedge^{h_i}V\), with increasing
basis wedges orthonormal. For an ordered array \(z_1,\ldots,z_n\)
of vectors in \(V\), regarded as columns of a matrix \(Z\), put
\[
  v_\alpha(Z)=\bigotimes_{i=1}^a(z_1\wedge\cdots\wedge z_{h_i}),
     \qquad W_\alpha(V)=\operatorname{span}_\C\{v_\alpha(Z):Z\}.
\]
The empty tensor is the scalar one. Use semistandard evaluation
coordinates \(b_\alpha(Z)\) as in
Definition~\ref{def:rectangular-flag-matrix}, allowing columns of
full height to contribute powers of \(\det Z\).
\end{definition}

\begin{lemma}\label{lem:flag-tensor-integral-realization}
There is a unique integer matrix \(I_\alpha\) satisfying
\(I_\alpha b_\alpha(Z)=v_\alpha(Z)\). It identifies the
evaluation coordinate space with \(W_\alpha(V)\) and admits
an integer left inverse given by coordinate selection.
The squared norm on evaluation coordinates is \(x^*G_\alpha x\)
with \(G_\alpha=I_\alpha^{\mathsf T}I_\alpha\). It is uniquely
determined by
\[
 b_\alpha(Y)^*G_\alpha b_\alpha(Z)
   =\prod_{i=1}^a\det\!\big((Y_{\le h_i})^* Z_{\le h_i}\big),
\]
where \(Z_{\le h}\) consists of the first \(h\) columns.
The action on \(W_\alpha(V)\) by the exterior powers of \(g\)
identifies in these coordinates with \(S_\alpha(g)\), the
matrix determined by \(b_\alpha(gZ)=S_\alpha(g)b_\alpha(Z)\).
It is unitary on \(U(n)\) and irreducible there.
\end{lemma}
\begin{proof}
Expanding a wedge in increasing basis wedges gives exactly its
row minors. Thus every tensor coordinate is a product of flag
minors. Lemma~\ref{lem:flag-straighten} expresses each product
as an integral linear combination of the semistandard ones.
Each column of full height is already a determinant and may
be retained as a common factor throughout straightening.
Semistandard products are independent by
Proposition~\ref{prop:standard-flag-minors}, including this
common factor. Their evaluation vectors therefore span the
whole coordinate space. This proves existence and uniqueness,
and selecting the semistandard coordinates among all products
gives the left inverse. The induced norm consequently has the
displayed integer Gram matrix. In a single wedge its polarized
inner product is the determinant of the scalar inner products,
by multilinear expansion in an orthonormal basis as in
Proposition~\ref{prop:standard-flag-minors}. Multiplication
over the factors gives the displayed kernel, whose two sets
of evaluation vectors span and thus determine the matrix.

Applying \(g\) to each column of \(Z\) proves the action formula.
Exterior powers of a unitary matrix are unitary by the determinant
inner product. Invertible \(Z\) suffice to span \(W_\alpha\),
since each \(Z+tI\) is invertible except at finitely many \(t\)
and the tensor entries are polynomials in \(t\). Thus it is
the span of the general linear orbit of the standard tensor.
This coincides with the span of its unitary orbit by the
complexification argument of Theorem~\ref{thm:compact-highest-vector},
which also proves irreducibility. Full-height columns contribute
one-dimensional determinant factors to that theorem.
\end{proof}

\begin{theorem}\label{thm:integral-row-symmetrizer}
There is an integer endomorphism \(D_\alpha\) of the exterior
tensor space \(T_\alpha(V)\), in its increasing wedge bases,
with the following properties. It commutes with the general
linear action and is a positive scalar times orthogonal
projection onto \(W_\alpha(V)\). Its operator norm satisfies
\[
                 \|D_\alpha\|\le L!.
\]
If \(z_i\in\bigwedge^{h_i}V\) are any nonzero decomposable
wedges, \(D_\alpha(\bigotimes_i z_i)\ne0\).
For the empty diagram \(D_\alpha=1\).
\end{theorem}
\begin{proof}
Label the \(L\) tensor slots by the boxes of the diagram.
Let \(I:T_\alpha(V)\to V^{\otimes L}\) send each increasing
wedge to the signed sum of its ordered tensor permutations,
and tensor these maps over the columns. Its Hermitian adjoint
\(I^*\) sends a pure tensor to its columnwise wedges. This
follows by testing on the orthonormal basis words, including
words with repeated letters in a column. Both maps have
integer entries and commute with simultaneous linear changes
of variables. The nonzero words in distinct increasing wedge
basis images are disjoint and each occurs with coefficient
of modulus one. Thus \(\|I\|=\|I^*\|=(\prod_i h_i!)^{1/2}\).

Let \(\mathfrak R\) be the subgroup of the permutations of
the \(L\) slots which separately permutes each row. Set
\(D_\alpha=I^*(\sum_{\pi\in\mathfrak R}\pi)I\).
The sum divided by \(|\mathfrak R|\) is a self-adjoint
idempotent, by inversion and multiplication in the finite
group. Thus \(D_\alpha\) is positive semidefinite, integral,
and commutes with the diagonal general linear action. Its
norm is at most \(|\mathfrak R|\prod_i h_i!\). The column
permutation group \(\mathfrak C\) has order \(\prod_i h_i!\),
and \(\mathfrak R\cap\mathfrak C=\{1\}\) because each row and
column meet in at most one slot. The products of a row and a
column permutation are therefore all distinct, giving the
bound \(L!\).

A symmetrized tensor in \(s\) slots is spanned by powers
\(z^{\otimes s}\). Indeed the sum over subsets \(J\subset
\{1,\ldots,s\}\) of
\((-1)^{s-|J|}(\sum_{j\in J}z_j)^{\otimes s}\) selects
exactly the words using each \(z_j\) once, with labels counted
even when the vectors coincide. Applying this argument to
each row shows that the row symmetrizer has image spanned
by tensors whose row \(r\) is filled by repetitions of a
single vector \(z_r\). The map \(I^*\) sends every such tensor
to \(v_\alpha(Z)\), so \(D_\alpha\) has image in \(W_\alpha\).
It is nonzero there. For the standard highest tensor
\(v_\alpha(I_n)\), the word with all boxes in row \(r\)
filled by the \(r\)-th basis vector occurs in \(I v_\alpha(I_n)\)
with coefficient one. This word is fixed by the row group,
so the same coefficient in its row average is one. Hence the
quadratic form of \(D_\alpha\) on the highest tensor is positive.

By Lemma~\ref{lem:flag-tensor-integral-realization}, an eigenspace
of this commuting self-adjoint operator on \(W_\alpha\) is a
unitary invariant subspace of an irreducible space. Consequently
\(D_\alpha\) is a positive scalar there. Self-adjointness and
the inclusion of its image in \(W_\alpha\) make it zero on the
orthogonal complement. Finally, the inner product of
\(\bigotimes_i z_i\) against \(v_\alpha(Z)\), taken linear in
\(v_\alpha(Z)\), is a product of polynomials in \(Z\).
Each polynomial is nonzero since an appropriate list of
basis columns detects a nonzero wedge coordinate of \(z_i\).
Their product is nonzero since polynomial rings over \(\C\)
have no zero divisors, as follows by multiplying their leading
monomials. It takes a nonzero value by the polynomial
identity argument of Proposition~\ref{prop:standard-flag-minors}.
This shows that the tensor has nonzero projection as asserted.
\end{proof}

\section{Branching with the cost confined to the tall columns}
\begin{lemma}\label{lem:flag-through-subspace}
Let \(U\subset V\) have dimension \(u\). Suppose every height
of a diagram \(\beta\) is at least \(u\), and let \(a_\beta\)
be its number of columns. Write \(d(t)=t\,{\rm id}_U\oplus
{\rm id}_{U^\perp}\) for positive \(t\). The subspace of
\(W_\beta(V)\) on which \(d(t)\) acts by \(t^{ua_\beta}\)
for all \(t\) is spanned by flag tensors whose column spaces
all contain \(U\). If a tensor in \(T_\beta\) has that weight,
its image under \(D_\beta\) lies in this span.
\end{lemma}
\begin{proof}
In an orthonormal basis adapted to \(U\), the exponent on
a basis wedge counts the indices from \(U\). Thus \(u\) is
the largest possible exponent in each column, and
\(t^{-ua_\beta}d(t)\) on the tensor space tends as
\(t\to\infty\) to the projection onto the stated exponent.
It preserves \(W_\beta\) by Lemma~\ref{lem:flag-tensor-integral-realization}.
Since the flag tensors span \(W_\beta\), their limits span
its part of maximal exponent.

Consider an invertible flag array with subspaces \(F_h\)
spanned by its first \(h\) columns. If the orthogonal
projection of \(F_h\) onto \(U\) fails to have rank \(u\)
at a used height, the maximal component of that wedge is
zero, by expansion of the \(u\)-row minors. Otherwise choose
a basis of \(F_h\) projecting in its first \(u\) entries
to a basis of \(U\), and take the remaining vectors in
\(F_h\cap U^\perp\). Its maximal component is a nonzero
wedge spanning \(U\oplus(F_h\cap U^\perp)\). These spaces are
nested as \(h\) increases. Choosing a single ordered basis
adapted to the resulting chain identifies the product of the
limiting wedges with a scalar multiple of a flag tensor.
Invertible arrays suffice by the preceding lemma on tensor
coordinates. Conversely a flag through \(U\) already has
maximal weight in each factor. The last assertion follows
from equivariance and the image statement in
Theorem~\ref{thm:integral-row-symmetrizer}. The argument includes
the empty diagram and \(u=0\).
\end{proof}

\begin{theorem}\label{thm:factorial-integral-branching}
Write \(V^+=V\oplus\C e_*\), with \(e_*\) a last orthonormal
basis vector and \(\dim V=n\ge1\). At each height \(1\le h\le n\),
let a source diagram \(\alpha\) have \(c_h+r_h\) columns.
Retain \(c_h\) of them and shorten \(r_h\) by one box,
dropping the columns of resulting height zero. The target
diagram \(\beta\) has \(c_h+r_{h+1}\) columns of height \(h\),
where \(r_{n+1}=0\). Fix \(1\le q\le n\) such that \(c_h=0\)
for \(h<q\). Define
\[
 L=\sum_{h=q}^n h(c_h+r_{h+1}),\qquad R=\sum_{h=1}^n r_h.
\]
There is a nonzero integer matrix \(B\) between the
semistandard evaluation coordinates of \(\alpha,\beta\) with
\(\|B\|\le L!\) for the norms of
Lemma~\ref{lem:flag-tensor-integral-realization}. For every
field of characteristic zero, every \(g\in{\rm GL}_n\) over
that field and every nonzero scalar \(s\) there, it satisfies
\[
       B S_\alpha(\operatorname{diag}(g,s))
                      =s^R S_\beta(g)B.
\]
The tensor product of these maps over finitely many places
has norm at most \((\sum_\tau L_\tau)!\).
\end{theorem}
\begin{proof}
Let \(p:V^+\to V\) forget the last coordinate, and let
\(\lambda_*\) take that coordinate. In each retained tensor
factor apply \(\bigwedge^h p\). In each shortened one contract
the wedge by \(\lambda_*\) and project its remaining factors
by \(p\). On an increasing basis wedge this second map gives
zero when the last basis vector is absent, and otherwise gives
the signed wedge with that vector deleted. Both operations
therefore have integer matrices and norm at most one, since
distinct surviving basis wedges in each factor have distinct
images up to sign. The contractions at height one yield scalar
factors. Permute the surviving tensor factors into a fixed
order by height. Each contraction contributes \(s\) under
\(\operatorname{diag}(g,s)\); each other operation commutes
with \(g\). This establishes the block transformation law
on the full tensor spaces before the next operation.

Call a resulting factor tall when its height is at least
\(q\), and apply \(D_{\beta_{\rm tall}}\) only to their
product. This integer operator commutes with \(g\), and its
norm is at most \(L!\) by Theorem~\ref{thm:integral-row-symmetrizer}.
It acts with the same bound in the presence of the remaining
factors by expanding in an orthonormal basis there and summing
squared norms. We prove that the resulting map takes
\(W_\alpha(V^+)\) into \(W_\beta(V)\).

Take flags \(F_1\subset\cdots\subset F_n\) in \(V^+\) for
which \(p\) has full rank on \(F_n\) and \(\lambda_*\) is
nonzero on \(F_1\). Such flags suffice to span: these conditions
are nonvanishing of polynomial coordinates of an array, and
they hold, for example, on the columns \(e_1+e_*,e_2,\ldots,e_n\).
They hold on a nonempty Euclidean open set of arrays, so a
polynomial vanishing there is identically zero by successively
varying the entries in small disks. Apply this to any linear
functional on the span of flag tensors. For these flags every
retained wedge spans \(pF_h\), and every shortened wedge of
positive height spans \(F_h\cap V\). Indeed take the first
basis vector of \(F_h\) with \(\lambda_*\)-value one and the
others in its kernel, and contract that basis wedge. All these
wedges and the height-zero scalars are nonzero.

Put \(U=F_q\cap V\), of dimension \(q-1\). Every short
positive-height factor comes from shortening a column with
\(h\le q\), and therefore lies on the nested flag \(F_h\cap V\)
inside \(U\). Every tall factor contains \(U\). For retained
ones this uses \(pF_h\supset F_q\cap V\) when \(h\ge q\).
For shortened ones it follows from \(h>q\) and nesting.
The tall tensor consequently has maximal \(U\)-weight
as in Lemma~\ref{lem:flag-through-subspace}. Its image under
\(D_{\beta_{\rm tall}}\) is a linear combination of flag tensors
through \(U\). Each such flag extends the short flag inside
\(U\), so inserting the fixed short wedges produces a scalar
multiple of a full flag tensor. In detail, choose a basis
successively on the short subspaces, complete it in \(U\),
then in the successive tall subspaces, and finally in all
of \(V\). Each prescribed wedge differs by a scalar from its
basis wedge. This proves the inclusion, also when either group
of factors is empty. The image of a generic flag is nonzero:
the tall wedges are nonzero and have nonzero image together
by the last assertion of Theorem~\ref{thm:integral-row-symmetrizer},
and the short factors and scalar factors are all nonzero.

Precompose by \(I_\alpha\) and postcompose by the integral
coordinate selection for \(I_\beta\) in
Lemma~\ref{lem:flag-tensor-integral-realization}. These changes
are isometries on the respective flag spaces in the specified
norms. The resulting nonzero matrix \(B\) is integral and
satisfies the claimed norm bound. On complex invertible
matrices the asserted equivariance follows from the construction.
The two sides have integer polynomial entries in \(g,s\);
the identities therefore hold as polynomial identities, by
the open-set argument above. This proves the assertion over
each stated field. Finally the tensor product norm is at most
the product of the bounds by successive applications to
orthonormal factor expansions. One has \(\prod_\tau L_\tau!
\le(\sum_\tau L_\tau)!\) by the integral multinomial coefficient.
\end{proof}

\section{Multiplicity and the weight on the terminal subgroup}
\begin{lemma}\label{lem:compact-intertwiner-character}
Let \(K\) be a compact unitary group, or a finite product of
such groups, with bi-invariant probability measure, and let
\(S,T\) be finite-dimensional continuous unitary representations.
Write \(\chi_S,\chi_T\) for their traces. Then
\[
  \dim_\C\operatorname{Hom}_K(S,T)
         =\int_K\overline{\chi_S(k)}\,\chi_T(k)\dd k .
\]
For the flag representations of
Lemma~\ref{lem:flag-tensor-integral-realization} at one place,
twisted by integral powers of the determinant, the right
side is one for equal highest weights and zero for distinct
highest weights. The same assertion holds for independent
tensor products of the indicated types.
\end{lemma}
\begin{proof}
On the space of linear maps average
\(L\mapsto T(k)L S(k)^{-1}\). Left invariance of Haar measure
shows that its image consists of intertwiners, and every
intertwiner is fixed, so the average is an idempotent with
trace its image dimension. The trace of the transformation on
maps is \(\chi_T(k)\operatorname{tr}(S(k)^{-1})\), by expansion
in matrix units. Unitarity identifies the second trace with
\(\overline{\chi_S(k)}\), giving the formula.

At one place the representations are irreducible by
Theorem~\ref{thm:compact-highest-vector}. A nonzero intertwiner
between two irreducibles is an isomorphism by its kernel and
image. It extends its equivariance to invertible complex
matrices by differentiating along the skew Hermitian matrices,
complexifying and exponentiating, as proved in that theorem.
It consequently preserves the upper unitriangular fixed line
and the diagonal weight on that line. The weights must
coincide. Conversely equal weights have the same exterior
model, including the common determinant power. Endomorphisms
are scalar since any eigenspace is invariant. For the conjugate
actions \(g\mapsto S(\bar g)\), replace \(g\) by \(\bar g\)
throughout, an automorphism preserving the probability by
Proposition~\ref{prop:compact-formal-degree}. For product groups
the integrals factor by Lemma~\ref{lem:product-integration};
the multiplicities are the products of the local ones.
\end{proof}

\begin{proposition}\label{prop:rank-one-branch-multiplicity}
Use the diagrams of Theorem~\ref{thm:factorial-integral-branching}.
Let \(\lambda_i=\sum_{h\ge i}(c_h+r_h)\) for \(1\le i\le n+1\)
and \(\mu_i=\sum_{h\ge i}(c_h+r_{h+1})\) for \(1\le i\le n\),
where all undefined counts are zero. Under \(U(n)\) embedded
by \(g\mapsto\operatorname{diag}(g,1)\), the space of
intertwiners from \(W_\alpha(V^+)\) to \(W_\beta(V)\)
has dimension one and is spanned by the integral matrix \(B\)
in that theorem. This holds after twisting both representations
by the same determinant exponent, conjugating the entries of
both unitary arguments, and taking products of these constructions
at the infinite places.
\end{proposition}
\begin{proof}
One has \(\mu_i=\lambda_i-r_i\), so
\(\lambda_i\ge\mu_i\ge\lambda_{i+1}\) for \(1\le i\le n\).
Remove the boxes with letter \(n+1\) from any semistandard
tableau of shape \(\lambda\). Each is the bottom box of its
column and all removed boxes are rightmost in their rows.
The remaining row lengths \(\nu_i\) satisfy
\(\lambda_i\ge\nu_i\ge\lambda_{i+1}\) because there were
no two removed boxes in the same column. Conversely each
such diagram and each semistandard filling of it with
\(1,\ldots,n\) uniquely extends by filling the removed boxes
with \(n+1\). This is the removal bijection of
Lemma~\ref{lem:tableaux-binomial}. By the tableau character in
Theorem~\ref{thm:compact-alternant}, evaluating at
\(\operatorname{diag}(g,1)\) gives
\[
  \chi_\alpha(\operatorname{diag}(g,1))
       =\sum_{\lambda_i\ge\nu_i\ge\lambda_{i+1}}\chi_\nu(g).
\]
Indeed it holds on every diagonal unitary by the bijection,
and on every unitary by orthogonal diagonalization. Full
height \(n\) on the right is allowed by a determinant factor.
Use the preceding lemma on this identity with \(\chi_\mu\).
Exactly one summand contributes. Interchanging source and
target in the character integral conjugates it and hence
gives the same integer multiplicity. Theorem~\ref{thm:factorial-integral-branching}
supplies a nonzero map in the required direction, so it spans.
A common determinant power agrees on the embedded subgroup,
and conjugating the arguments preserves all the identities.
The product assertion follows from the last part of the lemma
and nonvanishing of the tensor product of the maps.
\end{proof}

\begin{definition}\label{def:terminal-branch-normalization}
For the actual last space put \(n=d_m-1=\dim_EY_m\), identifying
\(X_m=Y_m\oplus E\) in the terminal definite frame of
Proposition~\ref{prop:actual-pair-grade-zero}. Let
\(\iota(g)=\operatorname{diag}(g,1)\) in this frame, with
the ordered coordinates of \(Y_m\) used as the first \(n\)
coordinates. For \(1\le h\le n\) retain
\[
 c_h=
 \begin{cases}
 h_2-\lfloor(t+1)/2\rfloor,& h=k+i(t+1),\quad 0\le i<m,\\
 0,&\text{otherwise},
 \end{cases}
\]
where \(t=5r\), \(h_2=r^8\).
Write \(r_{\tau,h}=\kappa_\tau(h)-c_h\) and
\(e_\tau=e_{m,\tau}\) of
Proposition~\ref{prop:actual-definite-arrow-weights}. On the
definite group of \(Y_m\) let \(R_Y\) be the product of
\(S_{\beta_\tau}(\bar g_\tau)(\det\bar g_\tau)^{e_\tau}\),
where \(\beta_\tau\) has columns of counts
\(c_h+r_{\tau,h+1}\) with \(r_{\tau,n+1}=0\).
Denote its space by \(E_Y\). Put
\[
 b_\tau=e_\tau+\sum_{h=1}^n r_{\tau,h},\qquad
 a_N=\prod_{\tau\in\Phi}\tau(\bar\pi_N)^{-b_\tau}.
\]
The symbol \(b_\tau\) here denotes a branching exponent; the
box counts in Definition~\ref{def:stage-diagrams} are used
only in the estimates for the individual arrows. For a
modular form \(\psi\) of weight \(R_m\) on \(X_m\) define
\[
             (\mathcal R\psi)(h)=a_N B\psi(\iota(h))
                     \qquad(h\in U(Y_m)(\mathbb A_f)),
\]
where \(B\) is the product of the integral branching maps
constructed with cutoff \(q=k\).
\end{definition}

\begin{theorem}\label{thm:actual-terminal-branch-map}
For all sufficiently large scales the counts in
Definition~\ref{def:terminal-branch-normalization} give
branching maps as stated there. The space
\(\operatorname{Hom}_{U(Y_m)_\infty}(E_m,E_Y)\) is spanned
by \(B\), and \(\mathcal R\psi\) is a modular form of weight
\(R_Y\). One has \(b_\tau\ge\kappa_\tau(1)-2b\) at every
infinite place, where \(b\) is the scale of
Definition~\ref{def:heights}.
\end{theorem}
\begin{proof}
For large scales, \(k>1\) and \(h_2>t+1\), by
Lemma~\ref{lem:diagram-size} and \(t=5r,\ h_2=r^8\).
At \(k\), Definition~\ref{def:stage-diagrams} gives
\(\kappa_\tau(k)=w_1(\tau)-\lfloor d_1/2\rfloor\).
Here \(w_1(\tau)\ge3h_2\) by Definition~\ref{def:heights}
and \(d_1=k+t+1=O(r)\), so this exceeds \(c_k\).
At \(d_i\) for \(1\le i<m\) the count is \(2h_2-(t+1)\),
at least \(c_{d_i}\). At every other height the retained
count is zero. None of the used heights equals \(n\):
the largest possible source height apart from height one is
\(d_{m-1}=n-t<n\). In particular the target has no full-height
columns and has determinant shift \(e_\tau\) as written.
Theorem~\ref{thm:factorial-integral-branching} applies with
\(q=k\) and Proposition~\ref{prop:rank-one-branch-multiplicity}
gives the claimed span.

The block law at \(s=1\), including the common determinant
factor, gives \(B R_m(\iota(\gamma_\infty))=R_Y(\gamma_\infty)B\)
for rational \(\gamma\). Therefore \(\mathcal R\psi\) is left
covariant. On the right it is invariant under the inverse image
of the open stabilizer of \(\psi\), which is open by continuity
of the block inclusion and contains a compact open subgroup.
This proves it is a modular form. Finally every column of
height one is shortened since \(k>1\). Also
\(e_\tau=\lfloor d_{m-1}/2\rfloor-w_m(\tau)\ge-2b\)
by Proposition~\ref{prop:height-cost}. Summing the nonnegative
shortened counts proves the lower bound.
\end{proof}

\section{The height and the finite gauge of the terminal restriction}
\begin{theorem}\label{thm:terminal-restriction-height}
Denote the operator \(\mathcal R\) of
Definition~\ref{def:terminal-branch-normalization} by
\(\operatorname{Res}_N\). Uniformly on the windows of
Theorem~\ref{thm:mean-covolume}, for all sufficiently large \(x\),
every field embedding \(\sigma:\C\to\C\), every modular form
\(\psi\) on \(X_m\) of weight \(R_m\), and every finite argument
\(h\) on \(Y_m\), one has
\[
 \big\|\sigma\big((\operatorname{Res}_N\psi)(h)\big)\big\|^2
 \le \exp\!\left(x^{59}-\sum_{\tau} b_\tau\log N\right)
                       \big\|\sigma\big(\psi(\iota(h))\big)\big\|^2 .
\]
In this formula embeddings act entrywise in the evaluation
coordinates, and both norms are the Gram norms of
Lemma~\ref{lem:flag-tensor-integral-realization}. More exactly,
\(x^{59}\) in the exponent can be replaced by \(2\log L!\),
where
\[
       L=\sum_\tau\sum_{s=k}^{n}s(c_s+r_{\tau,s+1}).
\]
\end{theorem}
\begin{proof}
All entries of \(B\) are integers by
Theorem~\ref{thm:factorial-integral-branching}. The entries
therefore commute with every field embedding. The element
\(\pi_N\) constructed in Proposition~\ref{prop:level-prime}
is a Gaussian integer of norm \(N\). A complex embedding
of its coefficient field sends \(i\) to \(i\) or \(-i\),
so every \(\sigma(\tau(\bar\pi_N))\) has modulus \(\sqrt N\).
The absolute square of the scalar \(a_N\) after embedding is
thus \(N^{-\sum_\tau b_\tau}\), including when some exponents
are negative. The factorial norm estimate and the defining
restriction give the more exact assertion.

To bound \(L\), recall \(d_i=k+i(t+1)\), \(n=d_m-1\).
In the sum at one place the retained term can occur only
at the \(m\) heights \(d_0,\ldots,d_{m-1}\), with count
at most \(2h_2\) there. In a shortened term \(s\ge k\)
implies \(s+1>k\). The source multiplicity at such a
height is zero except at \(d_1,\ldots,d_{m-1}\), where it is
at most \(2h_2\), by Definition~\ref{def:stage-diagrams}.
In particular the potentially large multiplicity at \(k\)
does not enter this latter sum. Consequently
\(L\le4f m n h_2\). The estimates of
Lemma~\ref{lem:diagram-size} and Proposition~\ref{prop:height-cost}
give \(f\le x^{12},\ m\le2\sqrt{x},\ n=O(x^{11/2})\) and
\(h_2=O(x^{40})\). This proves \(L=O(x^{58})\) uniformly.
For \(L\ge1\), \(2\log(L!)\le2L\log L=O(x^{58}\log x)\le
x^{59}\) eventually, using \(\log x=o(x)\) by integration or
the exponential series. For \(L=0\) the same bound is immediate.
\end{proof}

\begin{proposition}\label{prop:terminal-restriction-gauge}
Let \(F\) be a characteristic-zero field with nonarchimedean
absolute value and \(\jmath_0:\C\to F\) a field embedding.
On evaluation coordinate vectors use the maximum absolute
value \(\|\cdot\|_F\). For a tuple \(g=(g_\tau)\) of invertible
\(n\)-by-\(n\) matrices over \(F\) put
\[
 D_Y(g)=\bigotimes_\tau(\det g_\tau)^{e_\tau} S_{\beta_\tau}(g_\tau).
\]
Define \(D_X\) on tuples of size \(n+1\) in the same way,
using the terminal source diagram and the same exponent.
For scalars \(s_\tau\in F^\times\) satisfying
\(|s_\tau|=|\jmath_0(\tau(\bar\pi_N))|\), put
\(\widetilde g_\tau=\operatorname{diag}(g_\tau,s_\tau)\).
For every finite argument \(h\) on \(Y_m\) one has
\[
 \big\|D_Y(g)^{-1}\jmath_0\big((\operatorname{Res}_N\psi)(h)\big)\big\|_F
      \le \big\|D_X(\widetilde g)^{-1}
                     \jmath_0\big(\psi(\iota(h))\big)\big\|_F .
\]
For a \(p\)-adic absolute value with \(p\ne N\) one may take
every \(s_\tau=1\). At \(p=N\) with \(|N|=N^{-1}\), one may
take \(s_\tau=N\) at the embeddings where
\(\jmath_0(\tau(\bar\pi_N))\) is a nonunit, and one otherwise.
\end{proposition}
\begin{proof}
The block identity of Theorem~\ref{thm:factorial-integral-branching}
and multiplicativity of the determinant give
\[
     B D_X(\widetilde g)=\left(\prod_\tau s_\tau^{b_\tau}\right)
                                  D_Y(g)B .
\]
Since \(\jmath_0(B)=B\), the vector on the left of the proposed
inequality is \(\jmath_0(a_N)\prod_\tau s_\tau^{b_\tau}\)
times \(B\) applied to the vector on the right. The absolute
value of the scalar is one by hypothesis. Every integer is
a sum of units of absolute value one or its negative, hence
has absolute value at most one in \(F\). The ultrametric
inequality applied to each row of \(B\) proves the estimate.
When \(p\ne N\), the Gaussian integers \(\pi_N,\bar\pi_N\)
are integral at \(p\) and their product \(N\) is a unit, so
both are units. At \(N\), Lemma~\ref{lem:split-level}
shows the two factors generate coprime ideals. At any place
above \(N\) one is therefore a unit and the other has absolute
value \(|N|\). This gives the stated choices.
\end{proof}

\section{Separation of the terminal period}
\begin{definition}\label{def:terminal-coefficient-pairing}
Write \(H=U(Y_m)\), \(K=H_\infty\), and use the embedding
\(\iota:H\to H_m=U(X_m)\) of
Definition~\ref{def:terminal-branch-normalization}. Choose
a unit vector \(u_0\in E_Y\) in the Gram norm of
Lemma~\ref{lem:flag-tensor-integral-realization}. For a modular
form \(\eta:H_f\to E_Y\) of weight \(R_Y\) put
\[
    C_\eta(u,g)=\langle R_Y(u)u_0,\eta(g)\rangle,\qquad
    (C_\eta,C_\zeta)_{[H]}
             =\int_{[H]}\overline{C_\eta(h)}C_\zeta(h)\dd h .
\]
The Hermitian product is conjugate linear in its first argument.
In particular the pairing is linear in \(\zeta\). These functions
are continuous and rational invariant by the argument of
Definition~\ref{def:definite-form-coefficient}.
\end{definition}

\begin{theorem}\label{thm:terminal-branching-pairing}
For each row \(t\in E_m^\vee\) there is a scalar \(s(t)\),
linear in \(t\), such that for every modular form \(\psi\) on
\(X_m\) of weight \(R_m\) and \(\eta\) on \(Y_m\) of weight
\(R_Y\) one has
\[
    \int_{[H]} F_{\psi,t}(\iota(h))\,\overline{C_\eta(h)}\dd h
                       =s(t)\,(C_\eta,C_{\operatorname{Res}_N\psi})_{[H]} .
\]
More precisely, let \(D_Y=\dim E_Y\) and let \(z(t)\) be
determined by the compact average
\[
   \int_K R_Y(u)u_0\,t R_m(\iota(u))^{-1}\dd u=z(t) B
\]
with probability Haar measure. Then \(s(t)=D_Y z(t)/a_N\).
\end{theorem}
\begin{proof}
The integrand of the compact average is a linear map from
\(E_m\) to \(E_Y\). Left translation of \(u\) shows that its
average intertwines the two group actions. Existence and
uniqueness of \(z(t)\) follow from the one-dimensional
intertwining space and \(B\ne0\) in
Theorem~\ref{thm:actual-terminal-branch-map}. In particular it
is linear in \(t\), hence continuous in the finite-dimensional
row space. At fixed \(g\in H_f\) integrate first over \(K\)
in the left side of the assertion. Conjugating the coefficient
of \(\eta\) gives \(\langle \eta(g),R_Y(u)u_0\rangle\).
Linearity in the second argument and the displayed average
show that the result equals
\[
     z(t)\,\langle\eta(g),B\psi(\iota(g))\rangle
       =\frac{z(t)}{a_N}
                   \langle\eta(g),(\operatorname{Res}_N\psi)(g)\rangle .
\]
To integrate over the finite group apply
Lemma~\ref{lem:definite-class-quotient-measures} to \(H\),
with its integral compact subgroup of volume one. Both functions
in the product have compact quotient representatives and all
the indicated integrals are absolute.

The average on \(K\) of projections onto \(R_Y(u)u_0\)
commutes with \(K\). Its commutant is one dimensional by
Lemma~\ref{lem:compact-intertwiner-character} applied to the
same product weight on both sides. Its trace is one, so
the average is \(I/D_Y\). Therefore for any two forms of
weight \(R_Y\),
\[
        (C_\eta,C_\zeta)_{[H]}
        = \frac{b_H}{D_Y}\int_{\Gamma\backslash H_f}
                               \langle\eta(g),\zeta(g)\rangle\dd g ,
\]
where \(\Gamma=H(E^+)\). Comparing with the preceding integral
gives precisely \(s(t)\).
\end{proof}

\begin{corollary}\label{cor:separated-terminal-period}
Take \(x_m=1\) in Corollary~\ref{cor:even-unfolded-algebraic-coefficients}
and write \(\mathbf t\) for a tuple of preceding finite cut
positions. Let \(v(\mathbf t)\) be the tensor product in order
of all the evaluated finite pairs in that corollary, in the
same fixed models. Write \(\mathcal L^c_{v(\mathbf t)}\) for
the iterated class operator of
Definition~\ref{def:paired-definite-class-operator}, extended
linearly to the full tensor product. There is a continuous
function \(a_m(\xi)\) of the preceding infinite positions such
that for every column \(\phi\), every terminal modular form
\(\eta\) as above, and all cut tuples for which the initial
finite argument lies in \(A^{\rm new}\),
\[
 \int_{[H]}P_{m,\phi}(h)\overline{C_\eta(h)}\dd h
   = a_m(\xi)\,
       (C_\eta,C_{\operatorname{Res}_N(\mathcal L^c_{v(\mathbf t)}\phi)})_{[H]} .
\]
One function \(a_m\) works simultaneously for the finite twists
and cut tuples, columns and terminal tests.
\end{corollary}
\begin{proof}
The even-period formula identifies \(P_{m,\phi}(h)\) with
\(F_{\psi,t_m(\xi)}(\iota(h))\) for
\(\psi=\mathcal L^c_{v(\mathbf t)}\phi\). Apply
Theorem~\ref{thm:terminal-branching-pairing} and put
\(a_m(\xi)=s(t_m(\xi))\). The continuity and independence
assertions follow from the linearity of \(s\) and those for
\(t_m\) in Corollary~\ref{cor:even-unfolded-algebraic-coefficients}.
\end{proof}

\begin{corollary}\label{cor:conditional-ramified-transfer-pairing}
Let \(\phi\) be any ramified column. Choose the finite inclusion
tests and the stages of Theorem~\ref{thm:normalized-chain-finite-integration},
and fix the real tests. Suppose \(y_m=\overline{C_\eta}\) and
its successive normalized descents \(y_i\) satisfy the
continuity, automorphy and rapidity assumptions in that theorem.
Then there is a scalar \(\kappa\), independent of \(\eta,\phi\)
and of the finite twist, such that
\[
 \int_{[G_0]}\mathcal P_\phi(g)y_0(g)\dd g
     =\kappa\,(C_\eta,C_{\operatorname{Res}_N(\mathcal L^c_{v_*}\phi)})_{[H]},
 \qquad
 v_*=\sum_{\mathbf t}\frac{v(\mathbf t)}
                             {\prod_{i<m}m_i(t_i)} .
\]
The sum ranges over the actual finite cuts. For their normalized
evaluations, its coefficient and support estimates are those
of Corollary~\ref{cor:actual-weighted-cut-tensor-bound}, with the
common normalization specified there.
\end{corollary}
\begin{proof}
In Theorem~\ref{thm:normalized-chain-finite-integration} use
the full finite tensor space for \(M\), \(V(\mathbf t)=v(\mathbf t)\),
and the linear functional
\(v\mapsto(C_\eta,C_{\operatorname{Res}_N(\mathcal L^c_v\phi)})_{[H]}\).
Linearity follows from Theorem~\ref{thm:definite-jet-class-operator}
and Definition~\ref{def:terminal-coefficient-pairing}.
The required terminal separation was just proved. The descent
theorem gives the identity and the dependence of \(\kappa\).
The finite estimates apply to the same product of cuts by
Theorem~\ref{thm:actual-ramified-period-cuts} and the normalization
of Corollary~\ref{cor:actual-reduced-definite-pair}.
\end{proof}

\section{The ordered parameters of the row groups}
\begin{definition}\label{def:ordered-stage-parameters}
Fix an infinite place \(\tau\), suppress it from \(w_l(\tau)\),
and let \(0\le j\le m\). Put \(\epsilon_j=(j\bmod 2)/2\).
Use \(k,t,h_2\) of Definition~\ref{def:terminal-branch-normalization}.
The following list, with \(\epsilon_j\) subtracted from
\emph{every} entry, specifies the half-parameter on \(Y_j\).
Take the \(k-1\) numbers \(k/2-1-i,\ 0\le i<k-1\), then the
\(m\) numbers \(u_l=h_2-w_l,\ 1\le l\le m\). Immediately
after each \(u_l\) with \(1\le l\le j\) insert
\[
            -w_l+(t-1)/2-v,\qquad v=0,\ldots,t-1 .
\]
Give every entry positive sign except the \(u_l\) with
\(l\ge j+2\) and \(l-j\) even, which have negative sign.
The sign of a parameter means the signature of the coordinate
carrying it. Doubled parameters are twice the displayed,
shifted numbers.
\end{definition}

\begin{proposition}\label{prop:terminal-parameter-from-gaps}
The list in Definition~\ref{def:ordered-stage-parameters} is
strictly decreasing, has \(r+jt\) entries and has exactly
\(\lfloor(m-j)/2\rfloor\) negative signs. Any two entries
of opposite signs differ by at least \(2h_2\) in absolute
value. At \(j=m\) write \(p_0>\cdots>p_{n-1}\) for its
entries. With \(\beta_\tau(h)=c_h+r_{\tau,h+1}\) and
\(e_\tau\) as in Definition~\ref{def:terminal-branch-normalization},
one has
\[
        2p_i=2\left(e_\tau+
                \sum_{i<h<n}\beta_\tau(h)\right)+n-1-2i
                       \qquad(0\le i<n).
\]
In particular this gives the exact doubled highest-weight
parameter of the terminal definite representation, including
its determinant exponent.
\end{proposition}
\begin{proof}
Before insertion the list is the row list of
Definition~\ref{def:row-vandermonde}. Its consecutive \(u_l\)
differ by \(2h_2\) since the \(w_l\) differ by this amount.
A block inserted after \(u_l\) starts
\(h_2-(t-1)/2\) below that number, runs down by unit steps,
and ends \(h_2-(t-1)/2\) above \(u_{l+1}\) if the latter
exists. These gaps are positive by the inequalities in
Theorem~\ref{thm:actual-terminal-branch-map}. The last head
entry is \(1-k/2\); it lies above \(u_1\) by
Proposition~\ref{prop:exact-row-degree}. Thus the list is
strictly decreasing. Its length is \(k-1+m+jt=r+jt\).
Its negative signs occur at \(l=j+2,j+4,\ldots\le m\).
Their neighboring entries, when present, are unshifted-list
anchors \(u_{l-1},u_{l+1}\) before applying the common shift,
since all insertions stop with the block following \(u_j\).
The distance to either neighbor is \(2h_2\), giving the
asserted lower bound. The count agrees with the actual signature
in Theorem~\ref{thm:actual-row-lift-frames}. In the original
basis numbering a base entry \(i\) has
\(a=\max(0,i+1-k)\). Its sign is \((-1)^{a+\min(a,j)}\),
which is positive for \(a\le j\) and otherwise negative
exactly when \(a+j\) is odd. Substitution \(a=l-1\) gives
the listed signs on the anchors; the other coordinates have
positive signs.

At \(j=m\), \(m\) is even, so the common shift is zero.
The number \(k\) is even and \(t=5r\) is odd by
Definition~\ref{def:row-vandermonde}. Write
\(c_*=h_2-(t+1)/2\). Each gap from an anchor to its block
or from that block to the next anchor is \(1+c_*\).
The excess over one before \(u_1\) is
\[
       (1-k/2)-(h_2-w_1)-1
           = w_1-\lfloor(k+t+1)/2\rfloor-c_*
           = r_{\tau,k}.
\]
The anchor \(u_l\) has position \(d_{l-1}-1\), counting from
zero. The excesses at heights \(i+1=d_{l-1}\) are thus \(c_*\),
and at heights \(d_{l-1}-1\) with \(l\ge2\) they are also
\(c_*\). At the latter heights this is the shortened count
\(r_{\tau,d_{l-1}}\), since
\(\kappa_\tau(d_{l-1})=2h_2-(t+1)=2c_*\).
At height \(k-1\) the shortened count is the first-gap excess
just computed. Every remaining gap is one. This proves
\(p_i-p_{i+1}=1+\beta_\tau(i+1)\) for \(i<n-1\).
The last entry is \(-w_m-(t-1)/2=e_\tau-(n-1)/2\) because
\(n=d_{m-1}+t\) and \(d_{m-1}\) is even. Summing the gaps
back from this last entry proves the formula. The weights
are \(e_\tau+\sum_{i<h<n}\beta_\tau(h)\) by the diagonal
action on the leading wedges in Theorem~\ref{thm:compact-highest-vector}.
Adding the half-sum of the standard positive roots amounts
to adding \((n-1)/2-i\): among the differences of two coordinate
weights oriented from the smaller index to the larger, index
\(i\) occurs \(n-1-i\) times positively and \(i\) times negatively.
This proves the last assertion with the stated convention.
\end{proof}

\begin{corollary}\label{cor:actual-last-descent-rapid}
For the terminal tests of Definition~\ref{def:terminal-coefficient-pairing}
choose for \(u_0\) the unit flag tensor in the basis ordered
by decreasing parameter at each infinite place. Every scalar
function of this terminal coefficient type and fixed finite
level has the form \(C_\eta\). More precisely, for the kernel
\[
 K(u,v)=\dim(E_Y)\langle R_Y(u)u_0,R_Y(v)u_0\rangle
\]
the range of \(F(u,a)\mapsto\int_K K(u,v)F(v,a)\dd v\)
on continuous automorphic functions of that level is exactly
the family \(C_\eta\) of that level. Moreover for any such
\(\eta\) and any fixed elementary real-times-finite last
arrow test \(f_{m-1}\), the function
\[
 y_{m-1}(g)=[K_{m-1}:J_{m-1}]\int_{[U(Y_m)]}
 \theta\big(\Omega_{m-1}(g,h)f_{m-1}\big)\overline{C_\eta(h)}\dd h
\]
is continuous, automorphic, of open finite level and rapidly
decreasing on \(U(Y_{m-1})\).
\end{corollary}
\begin{proof}
An ordered orthonormal basis gives a pure tensor of unit wedges,
hence a unit vector by Lemma~\ref{lem:flag-tensor-integral-realization}.
For fixed \(a\) the projected function is the coefficient of
\[
       \eta(a)=\dim(E_Y)\int_K R_Y(v)u_0 F(v,a)\dd v .
\]
Rational covariance follows by a left translation in \(K\),
and the right finite level is inherited from \(F\). Conversely
the rank-one averaging identity in
Theorem~\ref{thm:terminal-branching-pairing} gives this vector
integral equal to \(\eta(a)\) on \(C_\eta\). This proves the
range statement. Here \(K\) as an integration domain denotes
the compact infinite group of \(Y_m\).

Both \(Y_m\) and \(Y_{m-1}=X_{m-1}\) are positive definite
by Definition~\ref{def:actual-row-lift-spaces}. Their quotients
have compact fundamental representatives by
Lemma~\ref{lem:definite-class-quotient-measures}. For \(g\) in
a compact neighborhood the theta kernel is continuous and
uniformly bounded on its product with the compact terminal
representatives, by Theorem~\ref{thm:actual-adelic-rational-theta-invariance}
and Proposition~\ref{prop:rational-frequency-lattice-bound}.
Thus the integral exists and is continuous by dominated
convergence. Rational invariance is
Theorem~\ref{thm:global-actual-arrow-covariance}; a right
finite level fixing the source action on the finite test
fixes the result, as in Theorem~\ref{thm:definite-jet-class-operator}.
Finally the reduced height in
Proposition~\ref{prop:unitary-reduced-height-integrability}
is bounded on definite quotient representatives, as proved
at the end of Theorem~\ref{thm:rapid-actual-one-step-unfold}.
A bounded continuous function on these representatives therefore
satisfies every rapid decay estimate in that proposition.
\end{proof}

\section{A rational translation seen by the finite inclusion datum}
\begin{lemma}\label{lem:primitive-short-and-target-residue}
Fix a stage \(Y\), its lattice \(\Lambda\), and a prime ideal
\(w\) of \(E\). For each \(g\in U(Y)(\mathbb A)\) there exists
\(v\in L(g)\) with nonzero reduction in
\((g_f\Lambda)_w/w(g_f\Lambda)_w\) and
\(\ell_g(v)\le C_E\lambda(g)\). Here the lattice and lengths
are those of Theorem~\ref{thm:unitary-linear-shortest-representative}.
For any target stage \(T\), \(h\in U(T)(\mathbb A)\), and
prescribed residue in \((h_f\Lambda_T)_w/w(h_f\Lambda_T)_w\),
there is \(v'\in L(h)\) with that residue and
\(\ell_h(v')\le C_{T,w}H(h)\). The constants are independent
of the group elements and residues.
\end{lemma}
\begin{proof}
Choose a shortest nonzero \(v_1\in L(g)\). Let \(e\) be the
largest integer with \(v_1\in w^e(g_f\Lambda)_w\). It is finite
since the local vector is nonzero. For \(e=0\) take \(v=v_1\).
For \(e\ge1\), choose \(\beta\in w^{-e}\setminus w^{1-e}\)
in the global fractional ideals; these ideals are strictly
nested by Proposition~\ref{prop:general-ideal-factorization}.
Subtract an element of \(O_E\) to put its infinite vector
in a fixed bounded lattice parallelepiped of \(O_E\). Since
\(O_E\subset w^{1-e}\), this keeps its valuation at \(w\)
equal to \(-e\). At every other finite place it is integral.
Thus \(v=\beta v_1\) belongs to \(L(g)\), is primitive at
\(w\), and has length at most the fixed bound for
\(\max_\tau|\tau(\beta)|\) times \(\lambda(g)\).

For the second assertion fix rational \(E\)-coordinates on \(T\).
For some fixed positive integer \(b_0\), using \(b(h)\) of
Proposition~\ref{prop:adelic-schwartz-frequency-height}, the
finite lattice \(h_f\Lambda_T\) lies between the completions of
\(b_0 b(h)O_E^{\dim T}\) and \((b_0 b(h))^{-1}O_E^{\dim T}\).
This follows by applying the clearing integer to both \(h_f\)
and its inverse in the fixed rational coordinates and by
comparability of the fixed reference lattices. The finite
completion argument of Lemma~\ref{lem:trace-finite} applies
between these two lattices. Concretely the intervening quotient
is finite and is the product of its primary parts, so any local
residue can be lifted globally, also after multiplying the
lower bound by the rational prime below \(w\). Choose one such
lift \(u\). If \(p\) is that prime, subtract a member of
\(p b_0 b(h)O_E^{\dim T}\) to place \(u\) inside its bounded
archimedean parallelepiped. This subtraction preserves membership
and the required residue. The resulting \(v'\) has ordinary
archimedean norm at most a fixed multiple of \(b(h)\).
After applying the archimedean frame \(h_\infty^{-1}\) its
norm is at most \(C_{T,w}H(h)\).
\end{proof}

\begin{theorem}\label{thm:rank-one-finite-detection}
Consider the arrow \(Y_j\to Y_{j+1}\) for any \(j<m\). Use the
initial-block inclusion \(a_{\rm inc}\) at all the selected finite
places and write \(\Phi=f_{a_{\rm inc}}\) for the finite test of
Corollary~\ref{cor:finite-inclusion-periods}. Choose one marked
place \(w\) over \(p=g_0\). For all adelic source and target
points \(g,h\) there is a rational linear map \(z:Y_j\to Y_{j+1}\)
of rank one with the following properties. Put
\(x=h^{-1}zg\) in the finite and infinite reference frames.
There is a nontrivial \(p\)-th root of unity \(\zeta\) such that
\[
           D_{x_f}\Phi=\zeta\Phi,\qquad
          \|x_\infty\|\le C\,\lambda(g)\,H(h).
\]
Here the norm is any of the fixed real Euclidean matrix norms,
and \(C\) may depend on the fixed arithmetic data, stage and
twist. The lower bound
\(|1-\zeta|\ge2\sin(\pi/p)>0\) is uniform in \(g,h\).
\end{theorem}
\begin{proof}
In the finite source and target isometric frames
\(\Lambda_j,\Lambda_{j+1}\) are standard integral lattices.
Let \(L_0=\alpha\operatorname{Hom}(\Lambda_j,\Lambda_{j+1})\)
and \(d_P\) be as in
Definition~\ref{def:split-inclusion-finite-oscillator}. Choose
once a scalar \(\beta\in E^\times\) integral in the fractional
ideal \(d_P\alpha\widehat O_E\) everywhere except \(\bar w\),
where its valuation is \(v_{\bar w}(\alpha)-1\). For existence,
take the fractional ideal with these local lower bounds.
Its product by \(\bar w\) is a proper subideal by
Proposition~\ref{prop:general-ideal-factorization}; any element
of the difference has the required properties. All but finitely
many exponents in these specifications are zero.

Choose a source vector \(v\) primitive at \(w\) by
Lemma~\ref{lem:primitive-short-and-target-residue}. Its framed
coordinate \(c=g^{-1}v\) at \(w\) has nonzero reduction
\(\bar c_{\rm res}\). The inclusion sends this reduction to a
nonzero vector in the target residue space. Choose a target
residue at \(\bar w\) such that its conjugate row on the
\(w\)-branch pairs nontrivially with this vector. Such a row
exists by choosing an index with nonzero coordinate; the
target hermitian frames identify conjugate branch lattices
in perfect duality. Lift the residue in \(h_f\Lambda_{j+1}\)
to \(v'\in L(h)\) by the same lemma, using the target frame
\(h_f^{-1}\) to prescribe that residue. Put
\[
           z(y)=\beta v'\,\langle v,y\rangle_{Y_j}.
\]
This is rational of rank one because \(v'\ne0\) and the
hermitian form is nondegenerate. In the infinite signed
frames it is represented after moving by the analogous outer
product of \(h_\infty^{-1}v'\) and the signed adjoint of
\(g_\infty^{-1}v\), times the embeddings of \(\beta\).
The two vector bounds give the stated real estimate.

At the finite places the two framed vectors are integral.
Hence \(x_f=z_-+u\), where \(z_-\) consists solely of its
component at \(\bar w\), in
\((\bar d_P^{-1}L_0)_{\bar w}\), and \(u\in d_P L_0\).
In the product rule for \(D_{z_-}D_u\), the phase is one.
Indeed only the component \(u_w\) can pair with \(z_-\);
its factor \(d_P\) cancels the single denominator in
\(\bar d_P^{-1}\) across conjugation. The remaining trace
is integral even with the trace-one factor \(\bar\omega\),
by Lemma~\ref{lem:hom-oscillator-dual-lattices}.
Corollary~\ref{cor:finite-inclusion-periods} says that \(D_u\)
fixes \(\Phi\). Translation by \(z_-\) on the Fourier sum of
Proposition~\ref{prop:finite-oscillator-residue-name} reindexes
its summation variable and gives the eigenvalue
\(\chi_{z_-}(a_{\rm inc})\).

In residue matrix coordinates this character is nontrivial.
In detail, after division by the reference lattice scales,
the reduction of the denominator-one-step component \(z_-\)
is a unit times the outer product of the chosen target residue
and the conjugate row of \(\bar c_{\rm res}\). The exponent in
the pairing with \(a_{\rm inc}\) is the matrix trace of its
conjugate transpose times the inclusion. This is a unit times
the chosen nonzero row pairing. The scalar pairing of the
single residue field with itself is perfect by
Proposition~\ref{prop:finite-oscillator-residue-name}, or by
the equality with the trace dual in
Lemma~\ref{lem:hom-oscillator-dual-lattices}. The marked prime
splits completely by Theorem~\ref{thm:data}, so its residue
field is \(\mathbb F_p\). A nontrivial additive character of
this group has kernel zero, since any nonzero element generates
the group. Thus the eigenvalue is a nontrivial \(p\)-th root
of unity. Its distance to one is at least the chord length
\(2\sin(\pi/p)\).
\end{proof}

\section{Rapid source decay of the inclusion theta kernel}
\begin{lemma}\label{lem:real-small-heisenberg-differences}
In fixed real symplectic coordinates write \(D_x^\infty\) for
the real Heisenberg translation in the gauge of
Definition~\ref{def:adelic-oscillator-translations}. Put
\[
   p_b(f)=\max_{|\nu|\le b}\sup_t(1+\|t\|)^b
                                      |\partial^\nu f(t)|
                         \qquad(b=0,1,2,\ldots).
\]
For every \(b,B\ge0\) integral there exist \(C,b'\) such that
\[
 p_b\big((D_x^\infty-I)^B f\big)\le
         C\|x\|^B p_{b'}(f) \qquad(\|x\|\le1)
\]
for every Schwartz test \(f\). For a fixed finite inclusion test
\(\Phi\), there also exist an integer \(a\ge0\), a seminorm
\(p_b\) and \(C_0>0\), independent of \(f,g,h\), for which
\[
 |\theta(\Omega_j(g,h)(f\otimes\Phi))|
              \le C_0 H(g)^a H(h)^a p_b(f).
\]
\end{lemma}
\begin{proof}
In the symmetric gauge the translation has the formula of
Lemma~\ref{lem:schroedinger-real-gaussian}. Passing to the
given gauge multiplies the operator by a phase with real
quadratic exponent in \(x\), as in
Theorem~\ref{thm:actual-adelic-rational-theta-invariance}.
Hence, writing \(x=(a_1,l)\), the expression is a phase
\(\exp(iQ(x)+iL(l,t)) f(t+a_1)\), with \(Q\) quadratic and
\(L\) real bilinear and fixed. For \(\|x\|\le1\) subtract
\(f(t)\) and use
\(f(t+a_1)-f(t)=\int_0^1(a_1\cdot\nabla)f(t+s a_1)\dd s\).
For the remaining phase difference use
\(|e^{iv}-1|\le |v|\) from integrating \(i e^{iv}\).
Every term of a derivative striking the phase at least once
contains a factor from \(l\); its other such factors are
bounded for \(\|x\|\le1\). The undifferentiated difference
has a factor at most \(C\|x\|(1+\|t\|)\). The weights at
\(t+s a_1\) and at \(t\) differ by at most a fixed factor.
These facts give \(p_b((D_x^\infty-I)f)\le C_b\|x\|
p_{b+2}(f)\). Iteration proves the first assertion.

For the second estimate inspect the seminorm bounds in
Proposition~\ref{prop:adelic-schwartz-frequency-height}.
The finite supremum there depends only on \(\Phi\) and
the finite clearing integer of the Hom action. In the real
estimate the \(L^2\) norms of finitely many words in the
position and differentiation operators are bounded by \(C p_b(f)\)
for one fixed \(b\): choose \(b\) larger than their derivative
orders and their polynomial degrees plus the real dimension
plus one, and integrate the resulting inverse power of
\(1+\|t\|\) over successive unit shells. Thus the degree in
\(H\) and the seminorm index are independent of the test.
Fixed changes of polarization by linear substitution, polynomial
phase and partial Fourier transform preserve continuity in
these seminorms by Theorem~\ref{thm:actual-adelic-rational-theta-invariance}.
Finally the matrix \(v\mapsto hvg^{-1}\) and its inverse
have clearing integers dividing a fixed integer times
\(b(h)b(g)\), by matrix multiplication in a rational field
basis. Their real norms are bounded by a fixed multiple
of the product of the two group norm bounds. This gives the
asserted estimate, increasing \(a\) to an integer if necessary.
The character in Theorem~\ref{thm:global-actual-arrow-covariance}
has modulus one.
\end{proof}

\begin{theorem}\label{thm:inclusion-kernel-source-decay}
For an actual arrow \(Y_j\to Y_{j+1}\) let
\(U(g)=\max(1,\lambda_{Y_j}(g)^{-1})\) and
\(V(h)=\max(1,\lambda_{Y_{j+1}}(h)^{-1})\). At the finite
places take the inclusion test \(\Phi\) of
Theorem~\ref{thm:rank-one-finite-detection}. There is an
integer \(a\ge0\) such that for every integer \(N_0\ge0\)
there are \(C_{N_0}>0\) and \(b_{N_0}\ge0\) with
\[
 |\theta(\Omega_j(g,h)(f\otimes\Phi))|
   \le C_{N_0}p_{b_{N_0}}(f)\,
                U(g)^{-N_0}V(h)^{N_0+2a}
\]
for every adelic \(g,h\) and every real Schwartz \(f\).
The constants may depend on the fixed stage and arithmetic
data, including its finite test and twist.
\end{theorem}
\begin{proof}
Rational invariance in each group follows from
Theorem~\ref{thm:global-actual-arrow-covariance}. Choose the
rational representatives of
Theorem~\ref{thm:unitary-linear-shortest-representative}
so that \(H(g)\le C U(g)\) and \(H(h)\le C V(h)\), with
the minima unchanged. For brevity let \(T(f)\) denote the
theta value at these points, with the finite test fixed.
Lemma~\ref{lem:real-small-heisenberg-differences} gives
\(|T(f)|\le C p_b(f)U(g)^a V(h)^a\).

If \(U(g)>1\), use the rational \(z\) of
Theorem~\ref{thm:rank-one-finite-detection}, with
\(x=h^{-1}zg\). Then
\(\|x_\infty\|\le C_1 V(h)/U(g)\). Rational translation
by \(z\) preserves the theta value of every test: at zero the
translated value of the periodized function is \(F(z)=F(0)\)
by Definition~\ref{def:adelic-oscillator-translations}.
Intertwining carries \(D_z\) back through \(\Omega_j(g,h)\)
to \(D_x\), by Theorems~\ref{thm:actual-adelic-rational-theta-invariance}
and~\ref{thm:global-actual-arrow-covariance}; the character
commutes with translations. In the product model \(D_x\)
acts by \(D^\infty_{x_\infty}\otimes D_{x_f}\). The resulting
scalar identity, valid for \emph{every} real test, is
\[
             T(f)=\zeta\,T(D^\infty_{x_\infty}f).
\]
Consequently
\(T((D^\infty_{x_\infty}-I)^B f)=(\zeta^{-1}-1)^B T(f)\)
for every integer \(B\ge0\), by induction. There is no
dependence on \(f\) in the choice of \(z\).

If \(U(g)\le\max(1,C_1)V(h)\), the coarse estimate already
implies the conclusion after multiplying its right side by
the bound for \(U(g)^{N_0+a}/V(h)^{N_0+a}\).
Otherwise \(U(g)>1\) and \(\|x_\infty\|\le1\). Apply the
coarse estimate to the \(B\)-fold difference, with
\(B=N_0+a\). The lower bound on \(|1-\zeta|\) in
Theorem~\ref{thm:rank-one-finite-detection} and the real
difference estimate give
\[
   |T(f)|\le C'_B p_{b+2B}(f)
        U(g)^a V(h)^a \big(V(h)/U(g)\big)^B .
\]
This is the claimed inequality. Both arguments include
\(N_0=0\), and rational invariance gives the result at the
original representatives.
\end{proof}

\section{Transposition for the whole chain of inclusion kernels}
\begin{definition}\label{def:moderate-inclusion-chain}
For the row group \(G_i=U(Y_i)\) put \(U_i(g)=
\max(1,\lambda_{Y_i}(g)^{-1})\) as in
Theorem~\ref{thm:inclusion-kernel-source-decay}. Call a
continuous automorphic scalar function \(z\) of moderate growth
if \(|z(g)|\le C U_i(g)^s\) for some \(C,s\ge0\). The finite
level of an automorphic function here is any compact open
subgroup under which it is invariant on the right.

Choose real Schwartz tests \(f_{i,\infty}\) and the fixed finite
inclusion tests \(\Phi_i\) on the consecutive arrows. Write
\[
   \mathcal K_i(g,h)=
     [K_i:J_i]\theta\big(\Omega_i(g,h)(f_{i,\infty}\otimes\Phi_i)\big).
\]
Use the measures of Theorem~\ref{thm:normalized-chain-finite-integration}.
For a source function \(z\) and a target function \(y\) put
\[
  (\mathcal F_i z)(h)=\int_{[G_i]}z(g)\mathcal K_i(g,h)\dd g,
 \qquad
  (\mathcal D_i y)(g)=\int_{[G_{i+1}]}\mathcal K_i(g,h)y(h)\dd h
\]
when absolutely convergent. Here \(\mathcal D_i\) is a
transpose for the bilinear integral pairing.
\end{definition}

\begin{theorem}\label{thm:inclusion-chain-transposition}
The operator \(\mathcal F_i\) takes moderate continuous
automorphic functions to moderate continuous automorphic
functions, and \(\mathcal D_i\) takes rapidly decreasing
continuous automorphic functions to such functions. For every
pair \(z,y\) in the respective domains both integrals below
converge absolutely and
\[
 \int_{[G_{i+1}]}(\mathcal F_i z)(h)y(h)\dd h
            =\int_{[G_i]}z(g)(\mathcal D_i y)(g)\dd g .
\]
The outputs have open finite level.
\end{theorem}
\begin{proof}
For each of the two groups, \(U_i\) is comparable to the
reduced height by
Theorem~\ref{thm:unitary-linear-shortest-representative}. The
polynomial sublevel volume bound in
Proposition~\ref{prop:unitary-reduced-height-integrability}
implies \(\int_{[G_i]}U_i^{-D}<\infty\) for sufficiently
large \(D\). Indeed on \(2^j\le U_i<2^{j+1}\) its integral
is at most \(C 2^{j(b-D)}\) for a fixed \(b\), and \(D>b\)
suffices. Similarly a rapidly decreasing function \(y\) on
the target satisfies
\(\int_{[G_{i+1}]}|y|U_{i+1}^b<\infty\) for any fixed \(b\).

Use the kernel bound in Theorem~\ref{thm:inclusion-kernel-source-decay},
including the finite constant \([K_i:J_i]\). Given a source
bound \(C U_i^s\), choose \(N_0>s+D\). Integration over the
source gives a bound \(C' U_{i+1}(h)^{N_0+2a}\) for the
absolute integral defining \(\mathcal F_i z\). Integrating
the same bound against \(|y|\) proves absolute integrability
of the full product on the product of the quotients. In
the reverse direction the kernel bound gives, for every \(N_0\),
\[
    |(\mathcal D_i y)(g)|\le C_{N_0} U_i(g)^{-N_0}
                   \int_{[G_{i+1}]}U_{i+1}(h)^{N_0+2a}|y(h)|\dd h ,
\]
which proves rapid decrease. Fubini's formula in
Lemma~\ref{lem:product-integration} yields transposition.

The kernels are jointly continuous and rationally invariant by
Theorem~\ref{thm:global-actual-arrow-covariance} and the rational
lattice bounds in Proposition~\ref{prop:rational-frequency-lattice-bound}.
On compact subsets in a variable, the corresponding \(U_i\)
is bounded, since \(U_i\ll H_{\rm red}\le H\) and the clearing
integers and real norms are bounded there. Thus the preceding
majorants can be chosen uniformly there for continuity in that
variable, by Lemma~\ref{lem:integral-tools}. Rational invariance
passes to the integrals. The right levels fixing the input
oscillator vector on the appropriate side contain compact open
subgroups by Corollary~\ref{cor:finite-inclusion-periods} and
Theorem~\ref{thm:actual-twisted-finite-arrow}; their invariance
passes to the outputs.
\end{proof}

\begin{theorem}\label{thm:projected-ramified-chain-transfer}
For a ramified column \(\phi\) let \(z_0=\mathcal P_\phi\) be
its actual Poincar\'e series of
Theorem~\ref{thm:normalized-chain-finite-integration}.
Recursively put \(z_{i+1}=\mathcal F_i z_i\) for \(0\le i<m\).
Let \(\Pi_Y\) be the compact infinite projection of
Corollary~\ref{cor:actual-last-descent-rapid}, with the ordered
unit flag \(u_0\). With the finite sum \(v_*\) of
Corollary~\ref{cor:conditional-ramified-transfer-pairing},
one has the equality of functions on \([U(Y_m)]\)
\[
       \Pi_Y z_m=\kappa\,C_{\operatorname{Res}_N
                                      (\mathcal L^c_{v_*}\phi)} .
\]
The same constant \(\kappa\) works for every ramified column
and every finite twist with the fixed real tests and weights.
All forward and backward integrations used in this identity
are absolutely convergent.
\end{theorem}
\begin{proof}
The initial Poincar\'e series is bounded by
Proposition~\ref{prop:full-frame-majorant} and
Theorem~\ref{thm:actual-full-poincare-columns}, using its
finite set of column coordinates. The finite translation of
Lemma~\ref{lem:initial-frame-central-translation} preserves
boundedness. Thus the forward functions exist and have moderate
growth by induction in Theorem~\ref{thm:inclusion-chain-transposition}.
For every terminal test choose \(y_m=\overline{C_\eta}\) and
let \(y_i=\mathcal D_i y_{i+1}\). At the definite terminal
group the initial \(y_m\) is rapidly decreasing, as in
Corollary~\ref{cor:actual-last-descent-rapid}. The transposition
theorem now gives all the hypotheses on the \(y_i\) in
Corollary~\ref{cor:conditional-ramified-transfer-pairing} and,
by its successive pairing identities,
\[
       (C_\eta,z_m)_{[H]}
          =\int_{[G_0]}z_0 y_0
          =\kappa (C_\eta,C_{\operatorname{Res}_N
                                  (\mathcal L^c_{v_*}\phi)})_{[H]}.
\]
Here \(H=U(Y_m)\) and all finite representatives and tests
are as in the finite descent formula.

In the compact variable the kernel of \(\Pi_Y\) is conjugate
symmetric on interchanging its arguments. Applying Fubini
with the compact representatives of
Lemma~\ref{lem:definite-class-quotient-measures} shows it is
self-adjoint for \((\cdot,\cdot)_{[H]}\). It fixes every
\(C_\eta\) by Corollary~\ref{cor:actual-last-descent-rapid}.
Thus \(z_m\) in the last display can be replaced by
\(\Pi_Y z_m\). The difference of the two functions in the
theorem belongs to the range of \(\Pi_Y\) and has open finite
level, taking the intersection of the two levels when necessary.
It is itself \(C_\eta\) for some \(\eta\) by that corollary.
Testing against this difference gives its squared norm zero.
A continuous automorphic function with zero squared integral
vanishes everywhere: any nonzero value stays bounded away
from zero in some open set, which has positive quotient measure
by discreteness and the Haar construction in
Lemma~\ref{lem:definite-class-quotient-measures}. This proves
the formula. The independence of \(\kappa\) is the independence
already established in the pairing identity.
\end{proof}

\section{Translated vacua for the projected transfer}
\begin{definition}\label{def:chain-translated-vacuum-coefficients}
Fix once for each definite arrow a nonzero scalar \(\lambda_i\)
and replace its unit finite vacuum \(v_i\) by \(v_i^\circ=\lambda_i v_i\).
The scalars are independent of the finite twist. Use physical
finite displacements \(s_i\), so that the corresponding integral
coordinates are measured after division by \(\alpha\) as in
Definition~\ref{def:finite-pair-allowed-shifts}. For a tuple
\(s=(s_i)_{0\le i<m}\) write
\[
              e(s)=\bigotimes_{i=0}^{m-1}D_{s_i}v_i^\circ .
\]
For a finitely supported function \(\Psi\) of such tuples put
\(w_\Psi=\sum_s\Psi(s)e(s)\) and \(B_\Psi\phi=
\operatorname{Res}_N(\mathcal L^c_{w_\Psi}\phi)\), with the
scale and the chain of Definition~\ref{def:definite-jet-class-operator}
and Definition~\ref{def:paired-definite-class-operator}. The
finite operators \(\mathcal P,\mathcal S\) on the paired factors
are those of Corollary~\ref{cor:chain-auxiliary-datum-recovery}.
\end{definition}

\begin{theorem}\label{thm:projected-chain-good-coefficients}
Fix the real data of Theorem~\ref{thm:projected-ramified-chain-transfer}
and the vacuum normalizations of
Definition~\ref{def:chain-translated-vacuum-coefficients}.
There is one \(b\in\C^\times\) such that for each finite twist
with set \(P\) and \(Q=\prod_{w\in P}q_w\) there is a finitely
supported function \(\Psi\) with the following properties.
Its values are algebraic and integral at every finite place of
residue characteristic prime to \(Q\). Each coordinate of a
tuple in its support is an allowed shift of
Definition~\ref{def:finite-pair-allowed-shifts} plus a shift
supported above \(N\) from
Proposition~\ref{prop:paired-output-shifts}. The tensor \(w_\Psi\)
is invariant under the simultaneous actions of \(K_{c,i}\)
for \(0\le i<m\), with the final group element equal to one,
and it satisfies
\[
       w_\Psi=\sum_s\Psi(s)\,\mathcal S\mathcal P e(s).
\]
For every ramified column one has
\[
                b\,\Pi_Y z_m(\phi)=C_{B_\Psi\phi}.
\]
When \(P\) is empty, for every field automorphism \(\sigma\) of
\(\C\) one has, uniformly throughout the windows and for all
sufficiently large \(x\),
\[
 \sum_s|\sigma\Psi(s)|
   \le (g^{f n_m^2})^{2m-1}
            N^{f((m/2)k+(t+1)(m/2)^2)}
   \le \exp(x^{25}).
\]
\end{theorem}
\begin{proof}
Take the normalized tensor \(w\) in
Corollary~\ref{cor:actual-weighted-cut-tensor-bound}. It is
a sum of translated auxiliary-prime sums with algebraic
coefficients integral away from \(Q\). Algebraicity here also
follows directly from the finite Fourier coefficients and arrow
scalars in Theorem~\ref{thm:paired-finite-chain-denominator} and
Corollary~\ref{cor:actual-character-finite-chain}. Expand every
auxiliary-prime sum. On pair \(j,j+1\) this gives
\(N^{f d_{j+1}}\) terms of unit modulus, allowing roots of
unity when displacements are combined by their multiplication
law. No division occurs. The support and the product of these
numbers over even \(j<m\) are computed in
Theorem~\ref{thm:paired-finite-chain-denominator}. Group terms
giving the same displacement tuple, summing their coefficients;
use zero off this finite set. Call the resulting function
\(\Psi_0\). With the original unit vacua its tensor is \(w\),
and with the chosen vacua it is
\((\prod_i\lambda_i)w\) by multilinearity. The coefficients
have not changed under this replacement. The nonarchimedean
triangle inequality preserves integrality under grouping.
The ordinary triangle inequality and the last estimate in
Corollary~\ref{cor:actual-weighted-cut-tensor-bound} give
the first bound for their total size when \(P\) is empty.
The same corollary proves invariance and
\(\mathcal S\mathcal P w=w\), so applying the two linear
operators term by term yields the displayed recovery formula
also for the scaled vacua.

To get the bound uniform in the arithmetic parameters, the
exponent of \(N\) is at most \(f m d_m\), since it is the sum
of \(f d_{j+1}\) over even \(j<m\). Here \(d_m=n_m+1\) by
Definition~\ref{def:actual-row-lift-spaces}. The bounds
\(f\le x^{12}\), \(m=O(\sqrt x)\), \(n_m=O(x^{11/2})\)
are in the proof of Lemma~\ref{lem:diagram-size}, and
\(1+\log(gN)\le x\) in Corollary~\ref{cor:class-exponents}.
The logarithm of the first bound is therefore
\(O_{d,M,\epsilon}(x^{49/2})\), which is at most \(x^{25}\)
eventually.

The common finite model factors in
Corollary~\ref{cor:conditional-ramified-transfer-pairing}
and the factor \(C_*/D\) in
Corollary~\ref{cor:actual-weighted-cut-tensor-bound} convert
its \(v_*\) to a scalar multiple of \(w\), independent of the
twist. Thus Theorem~\ref{thm:projected-ramified-chain-transfer}
gives \(\Pi_Y z_m(\phi)=a C_{B_{\Psi_0}\phi}\) for
one \(a\in\C\), also incorporating the inverse of
\(\prod_i\lambda_i\). The coefficient is linear in the form
by Definition~\ref{def:terminal-coefficient-pairing}. If
\(a\ne0\), take \(b=a^{-1}\) and \(\Psi=\Psi_0\). If \(a=0\),
the projected chain vanishes for every column and twist.
Take \(b=1\) and the zero function for every \(\Psi\).
All the assertions then follow in this case as well.
\end{proof}


\section{Evaluation on tuples of definite classes}
\begin{theorem}\label{thm:joint-column-class-evaluation}
Use the definite groups and measures of
Definition~\ref{def:definite-jet-class-operator}. For \(0\le i<m\)
let \(A_i\subset H_{i,f}\) represent the finite double quotient
\(\Gamma_i\backslash H_{i,f}/K_{c,i}\).
Suppose \(w\) is a finite oscillator tensor fixed by the joint
column action, with the group element at \(m\) set equal to one.
Let \(\phi\) have right level \(K_{c,0}\). Fix \(h\in H_{m,f}\)
and put \(a_m=h\). For pure tensors \(e=\otimes_i e_i\) define
\[
 D_{\mathbf a}(e)=
      J_{m-1}^c\!\left(\Omega_{m-1,f}(a_{m-1},a_m)e_{m-1}\right)
      \cdots
      J_0^c\!\left(\Omega_{0,f}(a_0,a_1)e_0\right)\phi(a_0),
\]
extending linearly to \(e\). Here \(\Omega_{i,f}\) denotes the
finite action on the arrow from \(i\) to \(i+1\). Then
\[
          (\mathcal L^c_w\phi)(h)=
                           \sum_{\mathbf a\in\prod_{i<m}A_i}
                                         D_{\mathbf a}(w).
\]
Each term on the right is independent of the choice of the
representatives in that tuple. The number of terms is at most
\(\exp(m x^{25})\le\exp(x^{26})\) for all sufficiently large \(x\)
in the windows of Corollary~\ref{cor:class-exponents}.
\end{theorem}
\begin{proof}
The product of linear maps is multilinear in the \(e_i\), which
gives the stated linear extension by the defining property of a
tensor product. The rational covariance identity in the proof of
Theorem~\ref{thm:definite-jet-class-operator} gives invariance
of \(D_{\mathbf a}(e)\) under the left rational groups at the
\(m\) integrated variables. Indeed at an interior variable
the factor \(R_i(\gamma_\infty)\) on the output of one matrix
cancels the inverse factor on the input of the next. At \(i=0\)
it cancels against the transformation of \(\phi\). By the same
theorem \(D_{\mathbf a}(e)\) is continuous and has open finite
levels in each integrated variable. Its integral over their
product quotients is obtained by composing the individual
operators. This follows by successive integration for pure
tensors and by linearity for arbitrary tensors. Fubini's lemma
(Lemma~\ref{lem:product-integration}) applies since the finite
product of the compact fundamental closures of
Lemma~\ref{lem:definite-class-quotient-measures} bounds all
these continuous integrands.

For \(\mathbf k=(k_i)_{i<m}\) in the column levels put \(k_m=1\).
The finite representation law changes \(D_{\mathbf a\mathbf k}(w)\)
into \(D_{\mathbf a}((\otimes_i\Omega_{i,f}(k_i,k_{i+1}))w)\),
since \(\phi(a_0 k_0)=\phi(a_0)\). This equals
\(D_{\mathbf a}(w)\). Each quotient integral may now be taken
at \(K_{c,i}\) itself. Its stabilizers are trivial. In fact a
stabilizer is finite by Lemma~\ref{lem:definite-class-quotient-measures}.
At the chosen component of \(g_0\) the level reduces to identity
in the split general linear group by
Proposition~\ref{prop:column-level-in-stage-bessel}. The proof
of Proposition~\ref{prop:column-dim} that \(I+g_0M_d(\Z_{g_0})\)
is torsion free applies in every rank, and the rational group
injects into that component. Thus the quotient measure assigns
mass one to each class by Lemma~\ref{lem:definite-class-quotient-measures}.
This proves the formula, including independence of representatives.

Each \(K_{c,i}\) contains the principal congruence subgroup of
level \(gN\), since that subgroup reduces to identity at all
specified places and satisfies the last-row congruence in
Proposition~\ref{prop:column-level-in-stage-bessel}. Its rank is
\(d_i=k+i(t+1)\), in the range of
Corollary~\ref{cor:class-exponents}. Applying that corollary to
each of the \(m\) factors proves the first counting bound.
The second follows from \(m=O(\sqrt x)\) in
Lemma~\ref{lem:diagram-size}.
\end{proof}

\begin{corollary}\label{cor:joint-column-class-size}
Take \(w=w_\Psi\) as in
Theorem~\ref{thm:projected-chain-good-coefficients}. Embed the
field generated by the coordinates being evaluated, the
coefficients and the arrow matrix entries in a valued field \(F\).
At each representative and at \(h\) choose norms on the
corresponding coordinate spaces over \(F\), ultrametric when
the absolute value is nonarchimedean. Suppose the embedded
\(\phi(a_0)\) have norm at most \(B_0\), and that at every tuple
\(\mathbf a\) and for each \(s\) with \(\Psi(s)\ne0\), the
\(i\)-th displayed jet map applied to \(e_i=D_{s_i}v_i^\circ\)
has operator norm at most \(b_i\) between the chosen spaces.
If \(F\) is nonarchimedean of residue characteristic prime to
\(Q\), with absolute value extending its usual prime-adic
absolute value, then the embedded value of
\((\mathcal L^c_{w}\phi)(h)\) has norm at most \(B_0\prod_i b_i\).
For the complex absolute value after an automorphism of \(\C\),
when \(P=\varnothing\) its norm is at most
\[
                \exp\big((m+1)x^{25}\big)\,B_0\prod_{i<m}b_i .
\]
\end{corollary}
\begin{proof}
Expand each \(D_{\mathbf a}(w)\) of
Theorem~\ref{thm:joint-column-class-evaluation} in the tuples \(s\).
Successive applications of the operator inequalities bound each
\(D_{\mathbf a}(e(s))\) by \(B_0\prod_i b_i\). The sum is finite,
so applying the embedding commutes with all its additions and
multiplications. In the nonarchimedean case each embedded
\(\Psi(s)\) has absolute value at most one by
Theorem~\ref{thm:projected-chain-good-coefficients}. Taking the
maximum over all terms proves the first estimate. For the
complex norm the triangle inequality costs the number of
class tuples times the sum of the absolute coefficient values.
Apply the respective bounds from
Theorems~\ref{thm:joint-column-class-evaluation}
and~\ref{thm:projected-chain-good-coefficients}.
\end{proof}


\section{The holomorphic sections of a definite arrow}
\begin{lemma}\label{lem:coherent-translation-scale}
Identify a standard positive Hom space, by its matrix entries,
with \(E^n\). Use \(r(\gamma)\) of
Proposition~\ref{prop:frame-polarization}. In the real model of
Lemma~\ref{lem:schroedinger-real-gaussian} the symmetric translation
of a coherent vector satisfies
\[
 \rho_a(\gamma_\infty)B_z
     =\exp\!\left(H(r\gamma,z)-\|r\gamma\|^2/2\right)B_{z-r\gamma}.
\]
Here \(H\) in the exponent is the ordinary positive Hermitian
product of Proposition~\ref{prop:frame-polarization}.
For finite physical \(X\) write
\(\widetilde D_X={\bf e}_f(b(X,X)/2)D_X\), with \(b\) from
Definition~\ref{def:rational-polarized-oscillator}. For a finite
test \(\Phi\) put \(S_\Phi(z)=\theta(B_{-z}\otimes\Phi)\).
Then \(S_\Phi\) is entire and for rational \(\gamma\)
\[
        S_{\widetilde D_{\gamma_f}\Phi}
                          = W_{-r\gamma} S_\Phi .
\]
\end{lemma}
\begin{proof}
At each real place Theorem~\ref{thm:negative-block-real-splitting}
gives \(x+i l=\sqrt{2/\sqrt d}\,\tau(\gamma)\) in each positive
coordinate. Hence \(r\gamma=\sqrt\pi(x+i l)\). Substitute
the exponential of Proposition~\ref{prop:positive-hermite-theta-series}
in the formula for \(\rho_a(x,l)\). The coefficient of \(t\) in
the exponential multiplying \(\phi_0(t)\) is
\(2\sqrt\pi(z-r\gamma)^{\mathsf T}t\). After extracting
\(B_{z-r\gamma}\) the constant exponent is
\(\overline{r\gamma}^{\mathsf T}z-\|r\gamma\|^2/2\).
Multiplying over all places gives the first identity.
Entireness follows from Proposition~\ref{prop:positive-hermite-theta-series}.

The simultaneous real and finite symmetric translations at
\(\gamma\) leave the theta functional invariant. Indeed the
unsymmetrized translations do so by
Theorem~\ref{thm:rational-schroedinger-zero-section} and the
additional scalar is \({\bf e}(b(\gamma,\gamma)/2)=1\) on the
diagonal rational number. The intertwiners of
Proposition~\ref{prop:finite-polarized-vacuum} and
Theorem~\ref{thm:actual-adelic-rational-theta-invariance} preserve
this calculation in the indicated models. The inverse symmetric
real translation is the translation by \(-\gamma_\infty\).
Consequently
\[
 \theta(B_{-z}\otimes\widetilde D_{\gamma_f}\Phi)
       =\theta(\rho_a(-\gamma_\infty)B_{-z}\otimes\Phi)
       =e^{H(r\gamma,z)-\|r\gamma\|^2/2}\,S_\Phi(z-r\gamma).
\]
This is precisely the shift of Theorem~\ref{thm:frame-sections}.
\end{proof}

\begin{theorem}\label{thm:definite-holomorphic-test-identification}
For every standard definite arrow the linear map \(\Phi\mapsto S_\Phi\)
of Lemma~\ref{lem:coherent-translation-scale} is an isomorphism of
the finite oscillator onto \(\mathcal F\) of
Definition~\ref{def:finite-frame-space}. It sends the unit vacuum
to a nonzero element of \(\mathcal T(\Gamma_0)\) and
intertwines \(\widetilde D_X\) with \(T_X\) of
Theorem~\ref{thm:finite-frame-action}. In particular for \(u\ne0\),
any row \(A\), and any column vector \(v\) of the prescribed weight,
\[
  A J_j^{-u}(\Phi)v
        =\left[P_{A,v}(\partial)\, S_\Phi(uZ)\right]_{Z=0}.
\]
For the physical translated test \(D_X\Phi\), replace \(S_\Phi\)
on the right by \({\bf e}_f(-b(X,X)/2)T_X S_\Phi\). This scalar
is a root of unity.
\end{theorem}
\begin{proof}
The finite physical reference lattice in one coordinate is
\(\alpha\widehat O_E\), the completion of \(\mathcal A\) by
Definition~\ref{def:frame-lattice},
Proposition~\ref{prop:frame-lattice-trace}, and
Lemma~\ref{lem:actual-lift-scale}. For the unit vacuum \(v_0\),
\(D_l v_0=v_0\) for every element \(l\) of that lattice in all
\(n\) entries, by Proposition~\ref{prop:finite-oscillator-residue-name}.
For \(\gamma\in\mathcal A^n\) the value of \(b(\gamma,\gamma)\)
is \({\rm Tr}_{E/K}\sum_i\bar\gamma_i\gamma_i\). Indeed
\(\jmath\bar\omega=1/2-1/(2\sqrt{-d})\), and summation over
conjugate embeddings cancels the purely imaginary summand for
this real norm. The trace is an integer by
Proposition~\ref{prop:frame-lattice-trace}. Thus
\({\bf e}_f(b(\gamma,\gamma)/2)=\chi_{\Gamma_0}(r\gamma)\)
by Theorem~\ref{thm:frame-sections}. The lemma gives its period
equations for \(S_{v_0}\), since the signs are unchanged by
\(\gamma\mapsto-\gamma\).

Each \(\Phi\) is a finite sum of translated vacua by
Definition~\ref{def:split-inclusion-finite-oscillator}. For some
positive integer \(h\) all physical translations by elements
of \(h\alpha\widehat O_E^n\) fix every summand. To see this, the
commutator of a translation with a fixed displacement is a
continuous additive character with open kernel by
Lemma~\ref{lem:theta-positive-translation} and
Proposition~\ref{prop:additive}. Intersect these finitely many
kernels with the reference lattice and take a sufficiently
large integer multiple of that lattice inside the intersection.
Applying the preceding sign calculation to \(h\mathcal A^n\)
now gives \(S_\Phi\in\mathcal T(h\Gamma_0)\).

Suppose \(S_\Phi=0\) and put \(f=\phi_0\otimes\Phi\).
The coherent translation formula with \(z=0\) gives
\(F^a_f(v_\infty,0)=0\) for every \(v_\infty\).
The rational translation law in
Theorem~\ref{thm:rational-schroedinger-zero-section}, multiplied
by the symmetric phase as in
Theorem~\ref{thm:rational-theta-symplectic-action}, propagates
this vanishing to all \((v_\infty,\gamma_f)\) for rational
\(\gamma\). The finite components of rational points are dense
by the coordinate argument in Theorem~\ref{thm:finite-frame-action}.
Continuity gives \(F^a_f=0\) everywhere. Its squared integral
is the squared norm of \(f\) by
Theorem~\ref{thm:rational-theta-symplectic-action}. The real
and finite model changes are unitary by
Proposition~\ref{prop:finite-polarized-vacuum} and
Theorem~\ref{thm:actual-adelic-rational-theta-invariance}.
Therefore \(\Phi=0\), which proves injectivity and in particular
\(S_{v_0}\ne0\).

For all rational displacements the intertwining statement is
the lemma and Theorem~\ref{thm:finite-frame-action}.
For fixed \(\Phi\) the function \(X\mapsto\widetilde D_X\Phi\)
is locally constant by the lattice period just proved, the
multiplication formula of Lemma~\ref{lem:theta-positive-translation}
and local constancy of the finite phase. The action on the
section is locally constant by Theorem~\ref{thm:finite-frame-action}.
Density extends the identity to all finite \(X\).
The translates of any nonzero member of \(\mathcal T(\Gamma_0)\)
span \(\mathcal F\) by that same theorem. Hence the map is onto.
Finally \(S_\Phi(uZ)=\mathcal T_{-u,\Phi}(Z)\) by the definitions.
Apply Definition~\ref{def:definite-jet-class-operator}.
The physical correction has finite order since the finite
adelic character takes values in the image of \(\Q/\Z\) by
Definition~\ref{def:rational-polarized-oscillator}.
\end{proof}


\section{All complex conjugates of the reduced torsion values}
\begin{lemma}\label{lem:uniform-complex-duplication}
Let \(\lambda,\Lambda,A,a_1\) be the normalized curve data of
Propositions~\ref{prop:scaled-cm} and~\ref{prop:cm-theta}.
Write \(s_y=(U_y\sigma_\Lambda)(0)\) and \(t_y=(U_y\theta_\Lambda)(0)\).
There exist \(0<c_d\le1\le C_d<\infty\), independent of the
torsion order and of the field automorphism \(\nu\) of \(\C\),
such that for every \(y\in\lambda K\), with
\(H_\nu(y)=\max(|\nu s_y|,|\nu t_y|)\), one has
\[
  c_d H_\nu(y)^4\le H_\nu(2y)\le C_d H_\nu(y)^4,\qquad
  C_d^{-1/3}\le H_\nu(y)\le c_d^{-1/3}.
\]
\end{lemma}
\begin{proof}
Theorem~\ref{thm:cm-unit} gives \(|\iota a_1|_p\le1\) for
every finite-place isomorphism. The transcendence-basis and
minimal-polynomial argument in Lemma~\ref{lem:cm-section-sign}
therefore makes \(a_1\) an algebraic integer. In particular
\(a_1\) has finitely many complex conjugates. Each conjugate
\(A'\) satisfies \((A')^4\ne64\), since \(a_1^4-64\) has
absolute value one at every finite place by the same theorem.

After reducing the Gaussian factors, the duplication equations of
Proposition~\ref{prop:cm-duplication} read
\[
 t_{2y}=t_y^4-s_y^4,\qquad
 s_{2y}^2=s_y^2 t_y^2
      \left(4(t_y^4+s_y^4)+a_1^2s_y^2t_y^2\right).
\]
Indeed the Gaussian for \(2y\) is the fourth power of that for
\(y\) by Proposition~\ref{prop:cm-theta}. Apply \(\nu\)
to these two algebraic equalities. For fixed \(A'=\nu(a_1)\)
the maximum of the modulus of the first output and the square
root of the modulus of the second output, on replacing
\((s_y,t_y)\) by \((a,b)\), is a continuous function homogeneous
of degree four in the absolute scale. On
\(\max(|a|,|b|)=1\) it is strictly positive. Indeed if its first
output vanishes, then \(a^4=b^4\ne0\). Vanishing also of the
second would imply \(8+(A')^2(a/b)^2=0\), and squaring
would give \((A')^4=64\). Thus its minimum on this compact
set is positive and its maximum finite. Take the least of
these minima and one, and the greatest of these maxima and
one, over the finitely many \(A'\). This proves the first
pair of inequalities. The normalized curve was fixed with \(d\)
in Proposition~\ref{prop:scaled-cm}. Changing the fourth root
used there gives only finitely many possible coefficients,
which can equally be included in these extrema.

For \(ny\in\Lambda\) and \(l\in\Lambda\), the two reduced
values at \(y+l\) differ from their respective values at
\(y\) by roots of unity. The sigma and theta multipliers in
Propositions~\ref{prop:reduced-sigma} and~\ref{prop:cm-theta}
give the sign of \(l\), an additional sign for theta, and
the phase \(\exp(i\pi\,\operatorname{Im}(\bar l y)/A)\);
the ratio \(\operatorname{Im}(\bar l y)/A\) belongs to
\(n^{-1}\Z\). These multipliers have unit modulus after
every \(\nu\). The two values are never both zero by
Theorem~\ref{thm:local-cm-height}. Hence \(H_\nu\) has positive
finite values on the finite group \(n^{-1}\Lambda/\Lambda\).
Writing \(h_j=H_\nu(2^j y)\), iteration gives
\[
   c_d^{(4^j-1)/3}h_0^{4^j}\ \le\ h_j\
                    \le C_d^{(4^j-1)/3}h_0^{4^j}.
\]
If \(h_0>c_d^{-1/3}\) the lower bound tends to infinity, and
if \(h_0<C_d^{-1/3}\) the upper bound tends to zero. Both
possibilities contradict the finite positive range. This
proves the remaining inequalities.
\end{proof}

\begin{corollary}\label{cor:conjugate-monomial-torsion-bound}
There is \(M_d\ge1\) such that every monomial \(P\) of degree
\(b\) in Lemma~\ref{lem:cm-section-sign}, every \(y\in\lambda K\),
and every field automorphism \(\nu\) of \(\C\) satisfy
\[
            \left|\nu\big((U_y^{[b]}P)(0)\big)\right|
                                      \le M_d^{\,b}.
\]
For the sections \(\mathcal S\) of Theorem~\ref{thm:frame-sections}
formed with multiplicities \(b_i\) in \(n\) coordinates, all
finite displacements \(X\) satisfy
\[
             \left|\nu\big((T_X\mathcal S)(0)\big)\right|
                                      \le M_d^{\,n\sum_i b_i}.
\]
\end{corollary}
\begin{proof}
The generators of the monomials are the sections of degrees
\(1,2,3\) in Lemma~\ref{lem:cm-section-sign}. For the first
two the bounds follow directly from
Lemma~\ref{lem:uniform-complex-duplication} on \(t_y\) and
\(s_y^2\). For the third put \(g_y=(U_y^{[3]}G_\Lambda)(0)\).
The last equation of Proposition~\ref{prop:cm-duplication},
with the homogeneous Gaussian factors inserted, reads
\[
       g_y^2+a_1 s_y^2 t_y g_y=s_y^6+s_y^2t_y^4 .
\]
Let \(K_0\) be the largest modulus of the finitely many
conjugates of \(a_1\). The modulus \(r=|\nu g_y|/H_\nu(y)^3\)
satisfies \(r^2\le K_0 r+2\). For example \(r\le K_0+2\),
since the opposite inequality gives \(r(r-K_0)>2\).
Increasing \(M_d\) makes all three generator bounds at most
\(M_d\) to the power of their degrees. The Gaussian multipliers
multiply with the degrees, so the bound passes to each monomial.
For any finite \(X\), Theorem~\ref{thm:finite-frame-action}
realizes \(T_X\mathcal S\) by a rational displacement on this
section. Apply the product identity of
Theorem~\ref{thm:frame-sections}. Its shifts all belong to
\(\lambda K\), and the degrees in the full product sum to
\(n\sum_i b_i\), which proves the second inequality.
\end{proof}


\section{An integral normalization of the vacuum in every rank}
\begin{proposition}\label{prop:integral-vacuum-character-sum}
For the given \(E=K(\zeta_{N_1})\) there is a nonzero
\(V^\circ\in\mathcal T(r(\mathcal A))\) in rank one such that,
for every finite displacement \(X\), \((T_X V^\circ)(0)\)
is an algebraic integer and
\[
          \sup_{\nu\in{\rm Aut}(\C)}
                  |\nu((T_X V^\circ)(0))|\le \exp(C_d N_1^2).
\]
Here the constant is independent of the exponent \(s\) determining
\(E\) in Definition~\ref{def:levels}.
\end{proposition}
\begin{proof}
Take \(D=N_1\), \(\delta=D\), and the \(D\) elements
\(w_j=\zeta_D^j,\ 0\le j<D\), each with multiplicity \(D\).
The geometric sum gives
\(\sum_j D\,\tau(w_j)\overline{\sigma(w_j)}
               =D^2{\bf1}_{\tau=\sigma}\)
since the characters of the cyclic group here have distinct
exponents by Proposition~\ref{prop:places}. These elements
belong to \(\mathcal A\) by Proposition~\ref{prop:frame-lattice-trace}.
Thus the length is \(L=D^2=\delta\bar\delta\). Use \(1\in O_K\)
for the element of norm \(\delta\bar\delta/L\) in
Theorem~\ref{thm:frame-sections}, and take every one-variable
monomial to be \(\theta_\Lambda^D\). The resulting section \(F\)
belongs to \(\mathcal T(\Gamma)\), where
\(\Gamma=D\Gamma_0\) and \(\Gamma_0=r(\mathcal A)\), and
\(F(0)=1\). Each \((T_XF)(0)\) is an algebraic integer by
Corollary~\ref{cor:integral-adelic-sections}. All its conjugates
have absolute value at most \(M_d^{D^2}\) by
Corollary~\ref{cor:conjugate-monomial-torsion-bound}.

For \(v\in\Gamma_0\) let \(U_v=\chi_{\Gamma_0}(v)^{-1}W_v\)
on \(\mathcal T(\Gamma)\). These shifts preserve the space
since \(\Gamma_0\subset\Gamma^\#\) by
Proposition~\ref{prop:frame-polarization}. Their two multipliers
cancel by the two multiplication equations in
Theorem~\ref{thm:period-dimension}, making \(U_v\) an ordinary
representation of \(\Gamma_0/\Gamma\). On \(\Gamma\) the two
sign prescriptions agree by Theorem~\ref{thm:frame-sections}.
The invariants are exactly \(\mathcal T(\Gamma_0)\).
The unnormalized sum \(P=\sum_{v\in\Gamma_0/\Gamma}U_v\) maps
into the invariants by reindexing. It acts there by the
nonzero integer \(|\Gamma_0/\Gamma|=D^{2f}\), since the
lattice has rank \(2f\) by Proposition~\ref{prop:frame-polarization}.
The invariant space has dimension one by
Theorem~\ref{thm:period-dimension}. Since the translates
\(W_a F\), \(a\in\Gamma^\#/\Gamma\) with any representatives,
span \(\mathcal T(\Gamma)\) by Proposition~\ref{prop:theta-shift-irrep},
some \(P W_aF\) is nonzero. Take this section for \(V^\circ\).

The shifts \(v,a\) correspond to rational displacements by
Proposition~\ref{prop:frame-polarization}. Theorem~\ref{thm:finite-frame-action}
therefore expresses \(T_X V^\circ\) as a sum of \(D^{2f}\)
single translates of \(F\), each with a root of unity as its
coefficient. Indeed each multiplication gives one value of
the finite character and each \(\chi_{\Gamma_0}\) is a sign.
No reciprocal of \(D\) occurs. The values at zero are integral
by the preceding construction and closure under sums and
products, as in Theorem~\ref{thm:frame-sections}. Every absolute
conjugate is bounded by \(D^{2f}M_d^{D^2}\). Since \(f\le D\)
by Proposition~\ref{prop:places}, its logarithm is at most
\((2+\log M_d)D^2\), using \(\log D\le D\). This gives the claim.
\end{proof}

\begin{corollary}\label{cor:integral-vacuum-chain-normalization}
In the standard positive space \(E^n\) the finite unit vacuum
can be scaled by a nonzero scalar so that its section is
\[
                  V_n^\circ(z)=\prod_{j=1}^n V^\circ(z_{\cdot,j}) .
\]
In this normalization all \(\theta(\phi_0\otimes D_Xv_n^\circ)\)
are algebraic integers, and all their complex conjugates have
modulus at most \(\exp(C_d nN_1^2)\), uniformly for finite \(X\).
Some such translated value is nonzero. These scalings can be
used on every definite arrow in
Theorem~\ref{thm:projected-chain-good-coefficients}.
\end{corollary}
\begin{proof}
In the product the multiplier on \(r(\mathcal A^n)\) is the
product of the one-coordinate multipliers by the explicit \(W\)
and sign in Theorem~\ref{thm:frame-sections}. Hence the
nonzero product belongs to the one-dimensional vacuum space
of Theorem~\ref{thm:period-dimension}. Its nonzero vectors
correspond to precisely the nonzero multiples of the finite
unit vacuum by Theorem~\ref{thm:definite-holomorphic-test-identification}.
For rational displacements the translated product factors as
the product of the translated one-coordinate sections by the
same \(W\) formula. For a general finite displacement approximate
all coordinates simultaneously by rational displacements, using
local constancy in Theorem~\ref{thm:finite-frame-action}. The
product identity persists. Apply
Proposition~\ref{prop:integral-vacuum-character-sum} coordinate
by coordinate. The physical correction in
Theorem~\ref{thm:definite-holomorphic-test-identification} is
a root of unity and leaves both conclusions valid. If every
translated value vanished, injectivity of evaluation on translates
in Theorem~\ref{thm:finite-frame-action} would make the product
zero. Finally the scalings depend only on the field, the
model identifications and the ranks, all independent of the
finite twist, as required in
Definition~\ref{def:chain-translated-vacuum-coefficients}.
\end{proof}


\section{Local expansions of the scaled definite kernels}
\begin{definition}\label{def:standard-local-test-module}
Use \(\lambda\) of Proposition~\ref{prop:scaled-cm},
\(\sigma_0=(2\pi/\sqrt d)^{1/2}\), and \(s_0\) of
Theorem~\ref{thm:global-frame-scale}. Fix \(s_1\in\C^\times\)
such that \(|\iota(s_1/s_0)|_p=1\) for every finite-place
isomorphism \(\iota\). This includes the choice \(s_1=s_0\).
Set \(u=\sigma_0s_1/\lambda\). For a prime \(p\), a given
\(\iota:\C\to K_p\), \(0\ne\delta\in O_K\), and \(n\ge1\),
write \(\mathcal D_p(\delta)\) for the module over
\(\iota^{-1}(O)\) generated by the finite vectors \(D_Xv_n^\circ\)
such that \(\delta X_p\in\mathcal A^n\otimes\Z_p\). Here
\(O=\{z\in K_p:|z|_p\le1\}\), \(X_p\) denotes all physical
components at \(p\), and the vacuum is normalized by
Corollary~\ref{cor:integral-vacuum-chain-normalization}.
For a section use its Taylor series in all complex variables
at the origin, with ordinary factorial denominators, and apply
\(\iota\) coefficient by coefficient when taking it over \(K_p\).
\end{definition}

\begin{theorem}\label{thm:standard-local-taylor-expansion}
Fix the data of Definition~\ref{def:standard-local-test-module}.
For every \(a\in\C\) and \(\Phi\in\mathcal D_p(\delta)\), the
transported Taylor series of \(S_\Phi(auZ)\) is a finite
\(O\)-linear combination of products of the form
\[
       \prod_{i,k}
       F_{y_{ik}}\!\left(\sum_{\tau\in G}
                    \iota(a s_1\,\overline{\tau(w_i)}/M)
                                  Z_{\tau k}\right).
\]
Here \(1\le k\le n\), \(F_y\) is from
Proposition~\ref{prop:cm-theta}, \(M\ge1\) is prime to \(p\),
and the \(w_i\in\mathcal A\), allowing repetitions, are a
tight frame of length \(M^2\). For each product one integer
\(b\ge1\), prime to \(p\), satisfies \(b\delta y_{ik}\in\Lambda\)
for all \(i,k\). The integer \(M\) and the frame can be taken
the same for every \(\delta,\Phi,a\) at the given \(p\).
\end{theorem}
\begin{proof}
If \(p\nmid N_1\), take the frame of
Proposition~\ref{prop:integral-vacuum-character-sum} with
\(M=N_1\), repeating each vector according to its multiplicity.
For \(p\mid N_1\), choose two consecutive positive integers
at least \(\sqrt{24dN_1^4(1+\log_2 N_1)}+1\). At least one is
prime to \(p\). Use it for \(M\) and apply
Theorem~\ref{thm:every-frame-length} at length \(M^2\) and \(R=0\).
Repeating vectors turns their multiplicities into ones without
changing the frame identity. In both cases construct a section
\(F\) in rank \(n\) by Theorem~\ref{thm:frame-sections}, using
\(M\) for its parameter \(\delta\), \(1\) for the element of
norm \(M^2/L\), and \(\theta_\Lambda\) at every repeated entry.
Then \(F(0)=1\).

Apply Theorem~\ref{thm:vacuum-sum} with \(\ell(S)=S(0)\),
\(V_0=V_n^\circ\), and the newly constructed \(F\).
With \(N'=M^2,\ g'=fn\), it gives
\[
 V_n^\circ=(N')^{-g'}\sum_{v\in\mathcal R}
                         (W_v V_n^\circ)(0)\,W_{-v}F .
\]
The representatives \(v\) have the form \(r(\xi_v)\) with
\(M\xi_v\in\mathcal A^n\). The coefficient in each summand is
\(p\)-integral by Corollary~\ref{cor:integral-vacuum-chain-normalization}
and Theorem~\ref{thm:finite-frame-action}, since \(p\nmid M\).
It suffices by linearity to treat one generator \(D_Xv_n^\circ\).
Pass from the physical to the symmetric action using
Theorem~\ref{thm:definite-holomorphic-test-identification}, and
apply \(T_X\) to the sum. The multiplication rule of
Theorem~\ref{thm:finite-frame-action} replaces \(T_XW_{-v}F\)
by a root of unity times \(T_{X+\xi_v}F\), with rational
\(\xi_v\) viewed at the finite places. All extra phases are
\(p\)-integral units.

There is a rational \(\gamma\in E^n\) with
\(T_\gamma F=T_{X+\xi_v}F\) and
\(\gamma_p-X_p-(\xi_v)_p\in\mathcal A^n\otimes\Z_p\).
Both conditions specify a neighborhood of \(X+\xi_v\) by
local constancy in Theorem~\ref{thm:finite-frame-action}, and
rational points are dense there by its proof. Since
\(\delta X_p\in\mathcal A^n\otimes\Z_p,\ p\nmid M\) and
\(M\xi_v\in\mathcal A^n\), this implies
\(\delta\gamma_p\in\mathcal A^n\otimes\Z_p\). In a
\(\Z\)-basis of \(\mathcal A^n\) the coordinates of
\(\delta\gamma\) are rational numbers integral at \(p\).
A common positive denominator \(b_0\) prime to \(p\) gives
\(b_0\delta\gamma\in\mathcal A^n\).

Use the product formula of Theorem~\ref{thm:frame-sections} on
this rational translate. The torsion shifts are
\[
 y_{ik}=-(\lambda/M)\,{\rm Tr}_{E/K}(\bar w_i\gamma_k).
\]
Thus \(b_0 M\delta y_{ik}\in\Lambda\) by
Proposition~\ref{prop:frame-lattice-trace}, simultaneously in
the indices. Take \(b=b_0M\). The coefficient of \(Z_{\tau k}\)
in the one-variable argument after scaling by \(au\) is
\((\lambda/(M\sigma_0))au\,\overline{\tau(w_i)}
 =a s_1\,\overline{\tau(w_i)}/M\). The Taylor product statement
of Theorem~\ref{thm:frame-sections} gives the displayed formula.
All sums are finite, and in each degree the coefficient of a
product or of a homogeneous linear substitution is a finite
polynomial expression. Coefficientwise application of \(\iota\)
therefore preserves all the identities.
\end{proof}


\section{Finite valuations of the rectangular jet matrices}
\begin{theorem}\label{thm:standard-test-all-prime-jet}
Use the scale \(u\) of Definition~\ref{def:standard-local-test-module}.
Let \(\Phi\in\mathcal D_p(\delta)\) in any positive rank. Choose
\(\delta_d\in\C^\times\) with \(\delta_d^4=d\).
Every mixed derivative at zero of
\[
                         S_\Phi(\bar\delta\,\delta_d\,u Z)
\]
is integral under \(\iota\). In particular on a definite arrow,
for \(a\in\C^\times\) every entry of \(J_j^{-au}(\Phi)\) satisfies
\[
 v_p\!\left(\iota(J_j^{-au}(\Phi)_{T',T})\right)
    \ \ge\ N_j\left(v_p(\iota a)-v_p(\iota\bar\delta)
                                     -\tfrac14v_p(d)\right)
\]
with \(v_p(0)=+\infty\).
\end{theorem}
\begin{proof}
Apply Theorem~\ref{thm:standard-local-taylor-expansion} with
\(a=\bar\delta\,\delta_d\). For any product in its expansion,
write \(\delta'=b\delta\), with \(p\nmid b\), so that
\(\delta' y_{ik}\in\Lambda\) for every shift. Put
\(t=\iota(\delta_d)\), which has valuation \(v_p(d)/4\).
The coefficient of \(Z_{\tau k}\) in the linear argument,
divided by \(\iota(\overline{\delta'})t\), is
\[
                   \frac{\iota(s_1\overline{\tau(w_i)})}{bM}.
\]
This is integral. Indeed membership in \(\mathcal A\) gives
\(v_p(\iota\overline{\tau(w_i)})\ge-d_\iota\) by
Definition~\ref{def:frame-lattice}. The valuation of \(\iota s_1\)
is \(d_\iota\) by Theorem~\ref{thm:global-frame-scale} and the
choice of \(s_1\). Both integers in the denominator are \(p\)-units.
Thus Theorem~\ref{thm:all-prime-jets} gives integrality of all
mixed derivatives of the product. Evaluation of a fixed derivative
commutes with coefficient transport, since it multiplies one
coefficient by a product of integer factorials. It also commutes
with the finite \(O\)-linear combination, proving the first assertion.

The entries of the rectangular Schur matrix are polynomials
with integer coefficients homogeneous of total degree \(N_j\)
by Proposition~\ref{prop:rectangular-flag-law}, taking the product
over the places as in
Definition~\ref{def:definite-jet-class-operator}. At scale
\(\bar\delta\,\delta_d\,u\) the corresponding jet is an integral
linear combination of the mixed derivatives just treated.
Theorem~\ref{thm:definite-holomorphic-test-identification}
identifies this jet with its class matrix before changing the
scale. Replacing \(\bar\delta\delta_d\) by \(a\) multiplies
every entry by \((a/(\bar\delta\delta_d))^{N_j}\) by homogeneity.
Taking valuations proves the bound.
\end{proof}

\begin{theorem}\label{thm:standard-test-nonsplit-minors}
Fix a prime \(p<x_1\) with \(\eta_d(p)\ne1\). On the \(j\)-th
definite arrow take \(\Phi\in\mathcal D_p(1)\) and
\(a\in\C\) with \(|\iota a|_p\le1\). For every \(V\subset G\)
and every charge system of Theorem~\ref{thm:schur-jet}, one has
\[
       v_p\!\left(\iota(J_j^{-au}(\Phi)_{T',T})\right)
                              \ge J_p(j,V)-m_{T'}
\]
whenever \(a\ne0\), for all semistandard output and input
indices \(T',T\). For \(a=0\) the same inequality holds for
the derivatives defining the entries at scale zero. All the
quantities on the right and the charge conditions are exactly
those of Definition~\ref{def:stage-diagrams} and
Theorem~\ref{thm:schur-jet}, including the exponent zero at
\(p=d\).
\end{theorem}
\begin{proof}
Such a prime is prime to \(N_1\) by the splitting argument of
Proposition~\ref{prop:frame-frobenius} and
Lemma~\ref{lem:inert-prime-fibres}. The frame in
Theorem~\ref{thm:standard-local-taylor-expansion} can therefore
be taken to be the repeated roots of unity, with \(M=N_1\).
Its length \(M^2\) is prime to \(p\). Also \(v_p(\iota s_1)=0\)
by Theorem~\ref{thm:global-frame-scale} and
Definition~\ref{def:frame-lattice}. Thus the scalar
\(\beta=\iota(a s_1/M)\) in the linear forms of that expansion
belongs to \(O\), and with \(\delta=1\) all torsion denominators
in each product can be cleared by one integer prime to \(p\).
Theorem~\ref{thm:actual-theta-normal} and
Theorem~\ref{thm:product-null} give integrality through degree
\(N_j\) after each of the linear substitutions in
Definition~\ref{def:linear-jet}. This integral coefficient
condition is preserved under the finite \(O\)-linear sum of the
expansion. It therefore holds for the transported series of
\(S_\Phi(auZ)\) itself.

Write the entries of the arrow variable as \(Z_{\tau,a',b'}\)
with \(a'\) the input index and \(b'\) the output index, so
this is the transpose display of the rectangular variable in
Definition~\ref{def:rectangular-flag-matrix}. The determinant
for a column in \(F_{T'}(Y)\) in Definition~\ref{def:linear-jet}
has entries \(\sum_{a'}Y_{\tau,i,a'}\partial_{\tau,a',b'}\).
Transposing its determinant puts the output labels in its rows
and the ordered positions \(i\) in its columns. Its product
over the diagram is precisely the polynomial
\(b_{d_{j+1}}(\partial Y^{\mathsf T})_{T'}\) across the places.
Each full input-height column in a semistandard input filling
has the entries \(1,\ldots,d_j\), in order. Removing these
columns is a bijection with fillings of the shorter diagram:
when they are restored on the left their entries in row \(i\)
are \(i\), and every entry in row \(i\) of a strictly increasing
column is at least \(i\). Their minor is \(\det Y\) at that
place. Consequently the standard coefficient indexed by the
remaining \(T\) is exactly the coefficient extraction in
Proposition~\ref{prop:rectangular-flag-law}. It gives the
\((T',T)\) jet entry by
Theorem~\ref{thm:definite-holomorphic-test-identification} for
nonzero scale and by the same differentiation formula at zero.
Theorem~\ref{thm:schur-jet} now proves the asserted bound.
\end{proof}


\section{Integral frames for a finite-place class estimate}
\begin{lemma}\label{lem:definite-one-prime-unitary-approximation}
Let \(H=U(E^h)\) for the positive standard Hermitian form, with
rational points \(\Gamma\) over \(E^+\). At any rational prime
\(p\), \(\Gamma\) is dense in the product \(H_p\) of the groups
over all completions of \(E^+\) above \(p\). Every finite position
\(g\in H_f\) therefore admits \(\gamma\in\Gamma\) such that
\((\gamma_f g)_p\) preserves \((O_E\otimes\Z_p)^h\).
\end{lemma}
\begin{proof}
In rational matrix coordinates put \(\mathfrak u=\{X:X^*=-X\}\).
Over \(\Q_p\) its scalar extension consists of all skew-adjoint
matrices with entries in \(E\otimes\Q_p\). Indeed the map
\(X\mapsto(X-X^*)/2\) is a rational projection whose image is
exactly this subspace over either field, including over \(\Q_2\).
Taking a rational basis of its image, approximate each coefficient
in \(\Q_p\) by a rational number, by truncating its \(p\)-adic
expansion. Thus the rational skew-adjoint matrices are dense
in the local ones.

For \(h_p\in H_p\) choose a rational scalar \(\eta\in K^\times\)
of norm one such that \(1+\eta h_p\) is invertible in every
field component. The scalars
\((1+j\sqrt{-d})/(1-j\sqrt{-d})\), \(j=1,2,\ldots\), are
pairwise distinct in every component by cross multiplication.
At a fixed component \(\det(1+T h_p)\) is a polynomial
with constant term one, so it rules out only finitely many
scalars. There are finitely many components by
Theorem~\ref{thm:integer-ring}. The matrix
\(X=(\eta h_p-I)(\eta h_p+I)^{-1}\) is skew-adjoint, by using
\((\eta h_p)^*=(\eta h_p)^{-1}\). The inverse formula is
\(\eta h_p=(I+X)(I-X)^{-1}\), with \(I-X\) invertible.
Approximate \(X\) by rational skew-adjoint matrices as above.
For sufficiently good approximants the inverse in this formula
exists in every component, by the determinant and adjugate
formulas. It then exists over \(E\) as well. Each resulting
matrix is rational unitary by taking adjoints, and after multiplication
by \(\eta^{-1}\) they approximate \(h_p\). Finally the integral
invertible matrices form an open subgroup at \(p\): requiring
a matrix and its inverse to be integral is an open condition
at every component. Approximate \(g_p^{-1}\) to obtain the
last assertion.
\end{proof}

\begin{proposition}\label{prop:local-test-unitary-preservation}
On a definite arrow let \(g,h\) be finite source and target
elements preserving the standard lattices at \(p\).
Then \(\Omega_{i,f}(g,h)\mathcal D_p(\delta)=\mathcal D_p(\delta)\)
for every \(0\ne\delta\in O_K\).
\end{proposition}
\begin{proof}
Put \(G X=hXg^{-1}\). For the physical lattice
\(L_0=\alpha M_{d_{i+1},d_i}(\widehat O_E)\), one has
\(G L_{0,p}=L_{0,p}\). The index
\(H=[GL_0:GL_0\cap L_0]\) is prime to \(p\): the local factor
at \(p\) is one and the quotient at each other rational prime
\(l\) has order a power of \(l\), by the diagonal lattice
calculation in Theorem~\ref{thm:split-cut-weil-multiplier}.
Take representatives \(r\) for this quotient zero above \(p\).
By Theorem~\ref{thm:finite-oscillator-lattice-transport} the
normalized intertwiner \(W_G\) sends the scaled vacuum
\(v_n^\circ\) to \(H^{-1/2}\sum_r D_r v_n^\circ\). It sends
\(D_Xv_n^\circ\) to the same sum translated by \(D_{GX}\).
Combining two displacements introduces only finite character
values by Lemma~\ref{lem:theta-positive-translation}. Each index
\(GX+r\) satisfies \(\delta(GX+r)_p\in L_{0,p}\) whenever
\(\delta X_p\in L_{0,p}\), and \(H^{-1/2}\) is a \(p\)-unit.

It remains to account for the splitting and determinant factors.
Theorem~\ref{thm:all-finite-weil-splitting} multiplies \(W_G\)
by \(\beta(J_G)\). The ideal \(J_G\) has valuation zero at all
places above \(p\) by Definition~\ref{def:finite-splitting-multiplier},
so its character value is a \(p\)-unit by
Theorem~\ref{thm:actual-global-splitting-character}. Both
determinants in Theorem~\ref{thm:global-actual-arrow-covariance}
are units above \(p\). Theorem~\ref{thm:torus-characters-elementary}
gives the same conclusion for their characters, including all
rank twists. Thus the entire operator preserves the module.
Apply the inclusion also to the inverse pair and use the
representation law to get equality.
\end{proof}

\begin{theorem}\label{thm:local-gauge-class-matrix}
For stage \(i\), let \(D_i(g_p)\) be the matrix over \(K_p\)
obtained by evaluating
\(\bigotimes_\tau(\det A_\tau)^{e_{i,\tau}}S_{d_i}(A_\tau)\)
at the matrices of \(g_p\) under the local embeddings induced
by \(z\mapsto\iota(\overline{\tau(z)})\) on \(E\).
The embedding extends to its completion by limits of Cauchy
sequences, with the valuation defining that place.
For \(v\in E_i\), the complex weight space, define
\[
             |v|_{i,g}=\|D_i(g_p)^{-1}\iota(v)\|_{\max}.
\]
Suppose that every \(\Phi\in\mathcal D_p(\delta)\) satisfies
\(\|\iota J_i^c(\Phi)\|_{\max}\le b_i\) for a fixed scale
\(c\ne0\), where the matrix norm is its largest entry modulus.
Then at all finite positions \(g,h\),
\[
              |J_i^c(\Omega_{i,f}(g,h)\Phi)v|_{i+1,h}
                              \le b_i |v|_{i,g}.
\]
Consequently, in the nonarchimedean conclusion of
Corollary~\ref{cor:joint-column-class-size}, these norms satisfy
the required operator inequalities whenever the \(i\)-th
translated vacua in the support belong to the respective
\(\mathcal D_p(\delta_i)\) and the bounds \(b_i\) hold there.
\end{theorem}
\begin{proof}
The evaluation matrices are multiplicative by
Proposition~\ref{prop:rectangular-flag-law} and multiplicativity
of the determinant. On a rational \(\gamma\) one has
\(D_i(\gamma_p)=\iota(R_i(\gamma_\infty))\) by
Proposition~\ref{prop:actual-definite-arrow-weights}. Hence
\(\left|R_i(\gamma_\infty)v\right|_{i,\gamma_f g}=|v|_{i,g}\).
For an integral unitary frame at \(p\), both \(D_i\) and its
inverse have integral entries by the integer polynomial formulas
and the unit determinant. Such a matrix preserves the max norm
by the ultrametric inequality applied in both directions.

Choose rational \(\gamma,\epsilon\) on the two sides by
Lemma~\ref{lem:definite-one-prime-unitary-approximation} making
\(\gamma_f g,\epsilon_f h\) integral at \(p\). The rational
covariance in Theorem~\ref{thm:definite-jet-class-operator}
identifies the prescribed inequality with that at the new two
positions, using \(R_i(\gamma_\infty)v\) and the preceding
isometries. There the two norms are ordinary max norms.
Proposition~\ref{prop:local-test-unitary-preservation} puts
\(\Omega_{i,f}(\gamma_f g,\epsilon_f h)\Phi\) in the same
module as \(\Phi\). By the hypothesis all matrix entries have
absolute value at most \(b_i\), so each output coordinate has
absolute value at most \(b_i\) times the input max norm.
This proves the inequality. Its product over the finite set of
stages gives the claimed application to the class sum.
\end{proof}


\section{The local estimate away from marked collisions}
\begin{theorem}\label{thm:actual-unmarked-chain-divisibility}
Take the normalized vacua of
Corollary~\ref{cor:integral-vacuum-chain-normalization} on all
arrows. Fix \(s_1,u\) of Definition~\ref{def:standard-local-test-module}
and take \(c=-\bar\pi_g u\), with \(\pi_g\in O_K\) as in
Theorem~\ref{thm:data}. Use the function \(\Psi\) of
Theorem~\ref{thm:projected-chain-good-coefficients} at this scale.
For \(p\nmid NQ\) fix \(\iota\), and suppose
\(|\phi(g)|_{0,g}\le1\) at all finite arguments of a column of
level \(K_{c,0}\), with the norms of
Theorem~\ref{thm:local-gauge-class-matrix}.
On \(U(Y_m)\) define \(|\cdot|_{Y,h}\) by the same maximum norm
in the frame \(D_Y\) of Proposition~\ref{prop:terminal-restriction-gauge},
using the embeddings of that theorem on the finite matrix \(h_p\)
as in Theorem~\ref{thm:local-gauge-class-matrix}. Then for every \(h\)
and every sufficiently large \(x\) in the specified windows the
following bounds hold.
\begin{enumerate}
\item If \(p\ne d\), then \(|(B_\Psi\phi)(h)|_{Y,h}\le1\).
\item If \(p<x_1\) and \(\eta_d(p)\ne1\), then
\[
 \begin{split}
 |(B_\Psi\phi)(h)|_{Y,h}
  &\le p^{-\sum_{j<m} J_p(j,\varnothing)}\\
  &\le p^{-G_p+m C_p(\varnothing)+30 f\ell a/p}.
 \end{split}
\]
In particular if this prime is unserved in
Definition~\ref{def:collisions}, one has
\(|(B_\Psi\phi)(h)|_{Y,h}\le p^{-G_p+2e_p/3}\).
\end{enumerate}
The estimates are uniform in \(\iota,\Psi,\phi,h\) for the indicated
data and include the ramified columns on which the projected
transfer identity holds.
\end{theorem}
\begin{proof}
The support statement of
Theorem~\ref{thm:projected-chain-good-coefficients} places each
translated vacuum in \(\mathcal D_p(\pi_g)\). Indeed the
displacements are zero at places above \(p\) outside the marked
level unless \(p\mid NQ\). At places dividing \(\pi_g\) their
entries divided by \(\alpha\) become integral after multiplication
by \(\bar\delta_g\) in Definition~\ref{def:finite-pair-allowed-shifts},
which equals \(\pi_g\) at those places by
Lemma~\ref{lem:actual-lift-scale}. Thus the defining bound holds
simultaneously above \(p\). The first estimate now follows for
the chain on \(X_m\) by taking \(\delta=\pi_g,\ a=\bar\pi_g\)
in Theorem~\ref{thm:standard-test-all-prime-jet}, where \(a\)
in that theorem denotes its Taylor multiplier. As \(p\ne d\)
its lower valuation is zero at every stage. Use
Theorem~\ref{thm:local-gauge-class-matrix} and then
Corollary~\ref{cor:joint-column-class-size} with every \(b_i=1\).
The condition \(p\nmid Q\) ensures integrality of every
coefficient in this class sum.

For the second bound, \(p\) differs from \(g_0,g_1\), since
\(p<x_1\) while \(\log g_i\ge10\ell>\log x_1\) by
Theorem~\ref{thm:data}. All the displacements in the expansion
are therefore integral at \(p\), and
\(|\iota(\bar\pi_g)|_p=1\) because
\(\pi_g\bar\pi_g=g_0g_1\). Apply
Theorem~\ref{thm:standard-test-nonsplit-minors} with
\(V=\varnothing\) and \(\gamma_{\tau,h,b'}=0\) for all indices.
The charge condition of Theorem~\ref{thm:schur-jet} is then
empty, and \(m_{T'}=0\) for every output. Thus the full matrix
at stage \(j\) has max norm at most
\(p^{-J_p(j,\varnothing)}\) on \(\mathcal D_p(1)\).
The same class and gauge arguments give the product of these
bounds on \(X_m\).

In either case evaluate the chain at \(\iota_H(h)=
\operatorname{diag}(h,1)\), the terminal inclusion in
Definition~\ref{def:terminal-branch-normalization}. The last
coordinate of its evaluated frame is \(1\), as allowed in
Proposition~\ref{prop:terminal-restriction-gauge} because \(p\ne N\).
That proposition passes both bounds to \(B_\Psi\phi\).
Summing Theorem~\ref{thm:diagram-gain} over \(j<m\) with
\(V=\varnothing\) gives the second inequality in the display.
For the unserved case
Proposition~\ref{prop:collision-exact} gives
\(mC_p(\varnothing)\le e_p/3\). Also
\(30 f\ell a/p\le e_p/3\) once \(\ell\ge30\), by
Proposition~\ref{prop:loss}. Theorem~\ref{thm:data} gives this
uniformly eventually. Adding the two losses proves the last bound.
\end{proof}


\section{A one-column Fourier expansion at the auxiliary prime}

\begin{definition}\label{def:level-residue-schroedinger}
Fix an embedding \(\iota:\C\longrightarrow K_N\). Let \(S\)
consist of the places defined by \(z\mapsto\iota(\overline{\tau z})\)
for the \(f\) embeddings \(\tau\) of the chosen CM type, and put
\(R=\{w:w\mid\pi_N\}\).
Write \(\mathfrak o=\prod_{w\in S}\Z_N\),
\(k_N=\prod_{w\in S}\mathbb F_N\), and \(e_j=(0,\ldots,0,1)^t\)
in rank \(d_j\). In each factor use the split matrix variable \(T_w\)
on the \(w\)-side, in units of \(\alpha_w\). The conjugate side
acts by characters of this matrix variable. Normalize the measure of
the integral box to one and write \(P_N\) for projection to its
unit indicator on the factors at \(N\). If \(P_N v={\bf1}_{\rm box}
\otimes v'\), define \([F]v=F\otimes v'\) for each locally constant,
compactly supported function of these matrices. This definition
includes \(v'=0\). For \(y\in k_N^{d_{j+1}}\) set
\[
 b_y(T)={\bf1}_{T\in M_{d_{j+1},d_j}(\mathfrak o)}
                 {\bf1}_{(T e_j\bmod N)=y}.
\]
\end{definition}

\begin{lemma}\label{lem:level-cut-vacuum-fourier}
The set \(S\) consists of one place in each conjugate pair above
\(N\), all of residue degree one. There exists
\(\delta\in O_K\setminus\{0\}\), depending on \(\iota\), such
that \(N^{f d_{j+1}}[b_y]D_Xv^\circ\in\mathcal D_N(\delta)\)
for every finite displacement \(X\), every \(y\), and every \(j<m\).
Moreover \(v_N(\iota\bar\delta)=0\).
\end{lemma}
\begin{proof}
The prime \(N\) splits completely by Lemma~\ref{lem:split-level}.
The \(f\) conjugated embeddings give the \(f\) distinct extensions
of the same place of \(K\), since the original embeddings fix
\(K\) and permute its \(f\) extensions freely at a completely
split prime. Conjugation gives the extensions of the other place.
Put \(\delta=\omega-h\), where \(h\in\Z\) represents the
residue of \(\iota(\omega)\) in \(\mathbb F_N\) and
\(\omega=(1+\sqrt{-d})/2\). The original and conjugate roots
are distinct since \(N\nmid d\). Thus \(\delta\) is divisible
by \(N\) at all the places in \(\bar S\), while
\(\iota\bar\delta\) is a unit.

In the split factor the two matrix summands, supported respectively
on \(S,\bar S\), are isotropic and dual. In fact their
commutator pairs each matrix entry with its conjugate entry by
the trace pairing of Lemma~\ref{lem:hom-oscillator-dual-lattices}.
Its conductor on the normalized coordinates is the integral
lattice in both arguments. The proof of
Proposition~\ref{prop:finite-polarized-vacuum}, applied just
to these local summands, therefore sends the local vacuum to
\({\bf1}_{\rm box}\). A displacement on \(S\) translates the
argument and a displacement on \(\bar S\) multiplies it by the
corresponding character. All finite products of these actions
differ from displacement by their sum by a root of unity, by
the multiplication law of Theorem~\ref{thm:finite-oscillator-lattice-transport}.
The tensor description in Definition~\ref{def:level-residue-schroedinger}
also follows directly by taking a basis of translated vacua
at \(N\) and its complementary factors. The local vacuum
coefficient of \(D_Xv^\circ\) is zero if its normalized
\(N\)-component is nonintegral: a nonintegral translation moves
the support off the box, and a nontrivial character on the box has
integral zero by translation invariance. Otherwise projection
replaces this component by the vacuum with a root of unity.
The scalar used to scale the global vacuum stays on the
complementary factor in both cases.

Choose in each factor a nontrivial character of the additive
group \(\mathbb F_N\), and let \(\psi\) be their product. On
integral matrices the finite geometric sum gives
\[
 N^{f d_{j+1}} b_y(T)
  ={\bf1}_{\rm box}(T)\sum_{t\in k_N^{d_{j+1}}}
                 \psi\big(t^t((Te_j\bmod N)-y)\big).
\]
Indeed in each differing coordinate the sum of all powers of
a nontrivial \(N\)-th root of unity vanishes, and each matching
coordinate contributes \(N\). Outside the box both sides
vanish. Perfect duality realizes each multiplier as a displacement
supported on \(\bar S\) in the dual of the last column, with entries
in \(N^{-1}\mathfrak o\) in normalized coordinates. To see the
denominator explicitly, pairing the quotient
\(N^{-1}\Z_N/\Z_N\) on the dual side with \(\Z_N/N\Z_N\)
is injective in each argument by the conductor calculation.
Both have \(N\) elements, so it realizes all residue characters.
Every summand acting on the projected vacuum is consequently
a translate with \(\delta X'_N\in\widehat{\mathcal A}^{\,d_jd_{j+1}}_N\)
and coefficient zero or a root of unity. This is precisely
membership in Definition~\ref{def:standard-local-test-module}.
\end{proof}

\begin{proposition}\label{prop:one-column-level-cut-bound}
At scale \(c=-\bar\pi_g u\), all finite frames \(g,h\),
all displacements \(X\), vectors \(v\) of the source weight,
and residues \(y\) satisfy
\[
 \big|J_j^c(\Omega_{j,f}(g,h)[b_y]D_Xv^\circ)v\big|_{j+1,h}
       \le N^{f d_{j+1}} |v|_{j,g}.
\]
In particular \(N^{fb}\) suffices for each arrow at all
sufficiently large arithmetic scales.
\end{proposition}
\begin{proof}
In Theorem~\ref{thm:standard-test-all-prime-jet} use the
\(\delta\) of Lemma~\ref{lem:level-cut-vacuum-fourier} and
\(a=\bar\pi_g\). The three valuations in its bound vanish
because \(N\nmid gd\) and \(\iota\bar\delta\) is a unit.
Theorem~\ref{thm:local-gauge-class-matrix} transfers the
resulting integral estimate at the identity to all frames.
Apply it to \(N^{f d_{j+1}}[b_y]D_Xv^\circ\) and divide,
using \(|N|_N=N^{-1}\). Finally \(d_{j+1}\le n_m+1
=O(x^{11/2})\) by Definition~\ref{def:actual-row-lift-spaces}
and Lemma~\ref{lem:diagram-size}, whereas \(b=\exp(\ell)\)
dominates every fixed power of \(x\) by Theorem~\ref{thm:data}.
\end{proof}

\begin{proposition}\label{prop:level-paired-residue-coordinate}
For \(U\subseteq S\) let \(\Delta_j(U)\) be the rank \(d_j\)
diagonal matrix with entries \(1,\ldots,1,N\) on \(U\) and
identity elsewhere. Use \(\psi_w(\bar x)\) equal to the
modulation at \(T=1\) of a displacement with dual coordinate
\(-\bar x/N\), lifting the residue to \(\Z_N\). Put
\(\psi=\prod_w\psi_w\) and
\[
 c_y(Z)={\bf1}_{\rm box}(Z)\,
            \psi\big((e_{j+2}^t Z\bmod N)y\big)
 \qquad (y\in k_N^{d_{j+1}}),
\]
where the second factor is evaluated only in the box. Put
\(A=\Delta_j(S\cap R)\), \(B=\Delta_{j+2}(S\setminus R)\).
On each even pair \(j,j+1\) the projection followed by the
sum of auxiliary translates in Corollary~\ref{cor:chain-auxiliary-datum-recovery}
takes \(v\otimes v'\) to
\[
     \sum_{y\in k_N^{d_{j+1}}}
           [\,b_y(\,\cdot\, A)\,]v \ \otimes\
           [\,c_y(B\,\cdot\,)\,]v'.
\]
\end{proposition}
\begin{proof}
Projection removes all the \(N\)-components except their vacuum
coefficients. Representatives in Theorem~\ref{thm:paired-neighbour-vacuum-sum}
may be changed by integral elements on their supported side:
that side is isotropic and integral displacements fix the vacuum.
Choose representatives with negative numerator divided by \(N\).
By Proposition~\ref{prop:paired-output-shifts} the two
displacements occupy the last column and last row in consecutive
matrices, their entries conjugate to each other. Under the two
maps to \(\Q_N\) related by field conjugation their numerators
therefore agree. The original two definite Hom spaces have
identity Gram matrices. In the same normalized entry the
modulation character on either space is thus the same, by
Lemma~\ref{lem:hom-oscillator-dual-lattices}. It is nontrivial
on \(N^{-1}\Z_N/\Z_N\) by the conductor computation just proved.

Consider one \(w\in S\cap R\), and denote the two matrix arguments
by \(T,Z\). At \(w\) the first displacement translates \(Te_j\)
by \(-z/N\) for a residue vector \(z\), and the second multiplies
an integral \(Z\) by \(\psi_w((e_{j+2}^t Z\bmod N)z)\).
All the other columns of \(T\) must be integral.
Precisely one \(z\) contributes when \(NTe_j\) is integral,
namely its reduction, and none otherwise. This is the claimed
sum with \(y_w=NTe_j\bmod N\). At \(w\in S\setminus R\)
the two roles are reversed. Now \(T\) must be integral, the
rows of \(Z\) other than the last must be integral, and
\(Ne_{j+2}^t Z\) must be integral. The unique translating
numerator is \(Ne_{j+2}^t Z\bmod N\). Its modulation on the
first factor gives precisely the indicated phase with
\(y_w=Te_j\bmod N\). Each computation used a pure translation
and a pure modulation on different factors, so no product
phase occurs. Multiplying over \(w\) proves equality on matrix
arguments. Equality in the algebraic tensor product follows:
in the finite-dimensional spans of the finitely many factors,
point evaluations span the dual since they separate functions.
The complementary components are unchanged by the projection
and these displacements.
\end{proof}


\section{Moving a residue Fourier operator across the arrows}

\begin{definition}\label{def:level-fourier-slide}
Use \(\mathfrak o,k_N,e_j,\psi\) of
Definition~\ref{def:level-residue-schroedinger} and
Proposition~\ref{prop:level-paired-residue-coordinate}. In
rank \(d_j\) put \(G_j^S=\prod_{w\in S}{\rm GL}_{d_j}(\Q_N)\),
\(K_j^S={\rm GL}_{d_j}(\mathfrak o)\), and
\(M_j=\{a\in K_j^S:\bar a e_j=e_j\}\), where the bar on an
integral matrix is reduction modulo \(N\).
Let \(\pi_j\) be a complex representation of \(G_j^S\) on
\(V_j\), and \(u_j=\Delta_j(S)^{-1}\) in this section.
A vector is ramified here if for every \(w\in S\) it belongs
to the image of \(1-\pi_j(x I)\) for a unit \(x\in\mathfrak o^\times\)
equal to \(1\) modulo \(N\) outside \(w\). The unit may depend on
the place and the vector.

For every \(y\in k_N^{d_j}\) with unit last entry \(y_*\)
choose \(\kappa_y\in K_j^S\) with last column \(y\) and last
row \(y^t\) modulo \(N\). Such a choice exists: if \(C\)
is the identity with last column replaced by \(y\), the
matrix \(y_*^{-1}CC^t\) has both prescribed boundaries and
unit determinant. Lifting its entries arbitrarily gives a
matrix with unit determinant over \(\mathfrak o\). Define
\[
     F_j=\pi_j(u_j)
        \sum_{y:\ y_*\in k_N^\times}\psi(y_*)\pi_j(\kappa_y).
\]
For a function \(f\) on residue matrices identify \(f\) with
its inflation to integral matrices extended by zero. Write
\(b_j=b_{e_{j+1}}\), \(c_j=c_{e_j}\), and let
\(r_j\) indicate integral matrices whose last row modulo \(N\)
is \(e_j^t\). Let \(L_j[f]:V_j\to V_{j+1}\) be linear in
both arguments. The two identities required below are
\begin{align*}
 \pi_{j+1}(b)L_j[f(\bar b\,\cdot\,\bar a)]\pi_j(a)&=L_j[f]
                              &&(a\in K_j^S,\ b\in K_{j+1}^S),\\
 \pi_{j+1}(u_{j+1})L_j[r_j]&=\gamma_j L_j[b_j]\pi_j(u_j)
                              &&(\gamma_j\in\C).
\end{align*}
\end{definition}

\begin{lemma}\label{lem:level-fourier-slide-one}
Under the identities of Definition~\ref{def:level-fourier-slide},
for every ramified \(\phi\in V_j\),
\[
 \pi_{j+1}(u_{j+1})L_j[c_j]\phi
    = F_{j+1}L_j[b_j]\phi
    = \gamma_j L_j[b_j]F_j\phi .
\]
Both \(F_j\) and each \(L_j[f]\) preserve ramification in
their respective target representations. The operator \(F_j\)
also preserves the subspace of \(M_j\)-fixed vectors.
\end{lemma}
\begin{proof}
Write \(T_*=e_{j+1}^t T e_j\) on integral matrices. For
\(t\in k_N\) with \(t_w=0\) choose a central difference
\(\phi=(1-\pi_j(xI))\xi\) at \(w\). Multiplication of \(T\)
by \(\bar x\) preserves the indicator \(1_{T_*=t}\):
at \(w\) it preserves the zero residue, elsewhere it is
identity. Take \(a=xI,b=1\) in the first identity of the
definition. Consequently \(L_j[1_{T_*=t}]\phi=0\).
Expanding \(c_j(T)=\psi(T_*)\) in the indicators of \(T_*\)
therefore allows restriction to unit \(T_*\).

Now partition this restricted function by the last column
\(y\). Its weight is \(\psi(y_*)\). The first identity with
\(a=1,b=\kappa_y\) gives
\(\pi_{j+1}(\kappa_y)L_j[b_j]=L_j[b_y]\), since
\(\bar\kappa_y e_{j+1}=y\). Summing and multiplying by
\(\pi_{j+1}(u_{j+1})\) gives the first equality.
Partition instead by the last row \(y^t\); its weight is again
\(\psi(y_*)\). With \(a=\kappa_y,b=1\) the indicator of
last row \(y^t\) gives \(L_j[r_j]\pi_j(\kappa_y)\), as
\(e_j^t\bar\kappa_y=y^t\). The second identity of the
definition gives the last equality.

Every central scalar commutes with each summand of \(F_j\).
Taking \(a=x^{-1}I,b=xI\) in integral covariance gives
\(\pi_{j+1}(xI)L_j[f]=L_j[f]\pi_j(xI)\), so the image
conditions persist. Finally take \(g\in M_j\) and
\(k=u_j^{-1}g u_j\). In the splitting off of the last
coordinate, \(g\) has integral entries, upper-right block
divisible by \(N\), and bottom-right entry \(1\bmod N\).
Thus \(k\) is integral, with integral inverse and last row
\(e_j^t\bmod N\). The permutation \(y\mapsto z=\bar k y\)
fixes \(y_*\). The element \(\kappa_z^{-1} k\kappa_y\)
belongs to \(M_j\) because its last column reduces to \(e_j\).
For any \(M_j\)-fixed vector \(\xi\) we get
\(\pi_j(g u_j\kappa_y)\xi=\pi_j(u_j\kappa_z)\xi\).
Summation gives the last assertion.
\end{proof}

\begin{theorem}\label{thm:level-fourier-slide-chain}
Suppose the identities of Definition~\ref{def:level-fourier-slide}
hold along \(2p\) consecutive arrows starting at stage zero.
Put \(B_j=L_j[b_j]\) and, for \(j=2i<2p\), let
\[
 T_i=\pi_{j+2}(u_{j+2})\sum_{y\in k_N^{d_{j+1}}}
                  L_{j+1}[c_y] L_j[b_y],\qquad
 o_i=(N^{d_{j+1}}-1)^f .
\]
For every ramified \(\phi\in V_0\),
\[
 T_{p-1}\cdots T_0\phi
  =\left(\prod_{i=0}^{p-1}o_i
                   \prod_{l=0}^{2i+1}\gamma_l\right)
       B_{2p-1}\cdots B_0 F_0^p\phi .
\]
Empty products of operators denote the identity.
\end{theorem}
\begin{proof}
If \(y_w=0\) as an entire vector at any \(w\), multiplication
by the residue of the central scalar at \(w\) preserves the
condition \(T e_j=y\). The central difference argument of
Lemma~\ref{lem:level-fourier-slide-one} gives \(L_j[b_y]\xi=0\)
on every ramified \(\xi\).
Otherwise choose \(k\in K_{j+1}^S\) with \(\bar k e_{j+1}=y\):
over each field factor complete the nonzero vector to a basis
by adjoining standard vectors which increase its span, then
lift the resulting invertible matrix. Integral covariance gives
\(L_j[b_y]=\pi_{j+1}(k)B_j\). On the next arrow
\(c_y(T)=c_{j+1}(T\bar k)\), so
\(L_{j+1}[c_y]\pi_{j+1}(k)=L_{j+1}[c_{j+1}]\).
There are \(o_i\) such vectors. The lemma first at \(j+1\)
and then at \(j\) now gives
\[
       T_i\xi=o_i\,\gamma_{j+1}\gamma_j B_{j+1}B_j F_j\xi .
\]
For completeness, moving a single \(F_s\) leftward in the
composition order toward the initial vector gives
\[
 F_s B_{s-1}\cdots B_0\eta
  =\left(\prod_{l<s}\gamma_l\right)
                    B_{s-1}\cdots B_0 F_0\eta
\]
on every ramified \(\eta\), by applying the second equality
of the lemma successively to \(B_l\cdots B_0\eta\).
All intermediate vectors remain ramified by that lemma.
Induct on \(p\), applying the pair identity to the output
of the first \(p-1\) pairs and the last formula with
\(s=2p-2\) to \(\eta=F_0^{p-1}\phi\). This gives exactly
the factors and the initial power in the assertion.
\end{proof}


\section{The auxiliary-prime estimate for the projected transfer}

\begin{definition}\label{def:level-shifted-column-action}
Keep \(S,R\) of Definition~\ref{def:level-residue-schroedinger}.
Embed \(G_j^S\) in the finite unitary group by the indicated
matrices on \(S\), their conjugate inverse transposes on
\(\bar S\), and identity elsewhere. Denote both a matrix and its
embedded element by the same letter when used in right translation
\(r_j\). Put \(U=S\cap R\), \(U'=S\setminus R\),
\(A_j=\Delta_j(U)\), and \(\pi_j(a)=r_j(A_j a A_j^{-1})\)
on smooth forms of the weight of stage \(j\).
For one finite vector \(v_j\) on this arrow put
\[
 L_j[F]=r_{j+1}(A_{j+1})\mathcal L^c_{j,[F]v_j}r_j(A_j)^{-1}
       =\mathcal L^c_{j,\Omega_{j,f}(A_j,A_{j+1})[F]v_j}.
\]
The equality is Theorem~\ref{thm:definite-jet-class-operator}.
Use the last-row level of Proposition~\ref{prop:column-level-in-stage-bessel}
also on \(X_m\) when discussing its local subgroup at \(N\).
\end{definition}

\begin{lemma}\label{lem:actual-level-character-factor}
The operators \(L_j\) satisfy both identities of
Definition~\ref{def:level-fourier-slide} with
\(\gamma_j=N^{f(d_{j+1}-d_j)/2}\). Define the complex character
of the local general linear group, for any even rank parameter \(l\), by
\[
 D_{j,l}(a)=\mu_{j+1}(1,\det a^{\rm emb})
                 \prod_{w\in S}|\det a_w|_w^{-l/2},
\]
where \(a^{\rm emb}\) denotes the unitary embedding and
\(|N|_w=N^{-1}\). For every \(V\subseteq S\),
\[
 \log_N|\iota D_{j,l}(\Delta_i(V))|_N
        = \sum_{\tau:\,w_\tau\in V} w_{j+1}(\tau)-\tfrac12 l|V|,
 \qquad w_\tau=\operatorname{place}(\iota\circ\bar{\phantom z}\circ\tau).
\]
\end{lemma}
\begin{proof}
Write \(T\) for the split coordinate. On named functions the
full finite action supported on the pairs at \(N\) satisfies
\[
 \Omega_{j,f}(a,b)[F]v_j
   = \frac{D_{j,d_j}(b)}{D_{j,d_{j+1}}(a)}
                        [F(b^{-1}\,\cdot\,a)]v_j .
\]
Indeed \(T\mapsto bTa^{-1}\) has modulus
\(\prod_w|\det b_w|_w^{d_j}|\det a_w|_w^{-d_{j+1}}\):
apply row operations on each column and column operations on
each row. Change of variables by its inverse, with the
inverse square root of this modulus, is the unitary intertwiner
\(W_G\): it has positive vacuum overlap, and this property
uniquely specifies the intertwiner by
Theorem~\ref{thm:finite-oscillator-lattice-transport}. For the
remaining phase, the pair of determinant ideals at \(w,\bar w\)
contributes \(\beta(\mathfrak p_w)^{-v_w(\det G_w)}\) by the
diagonal calculation of Theorem~\ref{thm:split-cut-weil-multiplier}.
The rank twists contribute its inverse, exactly as in the
proof of Theorem~\ref{thm:actual-twisted-finite-arrow}: the
possible leftover exponents are the two rank parities, and
\(d_i=k+i(t+1)\) is even since \(k\) is even and \(t\) odd
by Definition~\ref{def:row-vandermonde}. The remaining character
is the ratio of the two \(\mu_{j+1}\)-values. All these operators
act on the designated split factor and leave the complementary
test in the formula fixed.

For integral \(a,b\) their character factors are one by the
finite unit character in Theorem~\ref{thm:actual-twisted-finite-arrow}
and \(N\nmid Qg\). Applying Theorem~\ref{thm:definite-jet-class-operator}
and conjugating the two outer translations by \(A_j,A_{j+1}\)
gives integral covariance in precisely the stated order.
For \(a=u_j,b=u_{j+1}\), observe
\(r_j^{\rm row}(u_{j+1}^{-1} T u_j)=b_j(T)\), where
\(r_j^{\rm row}\) is the last-row indicator denoted \(r_j\)
in Definition~\ref{def:level-fourier-slide}. In block form
the argument scales its upper-right entries by \(N^{-1}\),
its lower-left entries by \(N\), leaving its lower-right entry
fixed. Integrality with the last-row condition after this
change is exactly integrality with the last-column condition
before it. Both determinants have the same norm-one embedding
at the finite places, so the \(\mu_{j+1}\) factors cancel.
The positive modulus factor is \(\gamma_j\), giving diagonal
covariance.

For the valuation formula put \(z=\det\Delta_i(V)^{\rm emb}\).
Its valuations are \(1\) on \(V\), \(-1\) on \(\bar V\)
and zero elsewhere. Theorem~\ref{thm:torus-characters-elementary}
gives \(z^h/a\) integral unit everywhere with \(a\bar a=1\)
and \(\mu_{j+1}(1,z)^h=\prod_\tau\tau(a)^{w_{j+1}(\tau)}\)
for some \(h\ge1\), since the weight is \(-w_{j+1}\).
The place induced by \(\iota\tau\) is \(\bar w_\tau\).
Hence \(v_N(\iota\tau(a))=-h\) when \(w_\tau\in V\) and
zero otherwise. Since \(v_N(N)=1\), including the positive
factor \(N^{l|V|/2}\) gives exactly the displayed equality.
\end{proof}

\begin{proposition}\label{prop:actual-F-chain-level-identity}
Let \(v_j\) be arbitrary finite vectors. On ramified columns
the definite chain of the recovered pairs
\(\mathcal S\mathcal P(\otimes_{j<m}v_j)\) is
\[
      t_N L_{m-1}[b_{m-1}]\cdots L_0[b_0] F_0^{m/2},
        \qquad |\,\iota(t_N)\,|_N\le N^{mfb}.
\]
The scalar \(t_N\) is independent of the \(v_j\) and the column.
The operator \(F_0\) preserves the column level.
\end{proposition}
\begin{proof}
The local level at \(N\) in the embedded coordinates equals
\(A_j M_j A_j^{-1}\). On \(U\) the original last row fixes
\(e_j^t\) modulo \(N\), so its lower-left block is divisible
by \(N\) and its bottom-right entry is \(1\bmod N\).
Conjugation by \(A_j^{-1}\) on the left and \(A_j\) on the
right gives exactly the last-column subgroup. On \(U'\) the
condition is already on the column, by inverse transpose at
the conjugate place. This proves the equality, including the
reverse inclusion, since the block calculations reverse.
By Theorem~\ref{thm:ram-dim}, at each split pair a ramified
column lies in the image of \(1-r_0(z)\) for a supported
central unit. It is also a scalar in the present coordinates
and commutes with \(A_0\), so it witnesses ramification for
Definition~\ref{def:level-fourier-slide}. Lemma~\ref{lem:level-fourier-slide-one}
shows \(F_0\) preserves the level at \(N\), and its summands
commute with translations at all other primes.

At an even \(j=2i\) put \(B=\Delta_{j+2}(U')\). In the pair
formula of Proposition~\ref{prop:level-paired-residue-coordinate}
move \(A_j\) from the first argument and \(B\) from the second
by the change of variables in Lemma~\ref{lem:actual-level-character-factor}
and translation covariance. Before summing in \(y\) the pair is
\[
 d_i'\, r_{j+2}(B^{-1})\,
 \mathcal L^c_{j+1,[c_y]v_{j+1}}\,
 \mathcal L^c_{j,[b_y]v_j}\,r_j(A_j)^{-1},
 \quad
 d_i'=D_{j,d_{j+1}}(A_j)D_{j+1,d_{j+1}}(B).
\]
Indeed the first change uses \((a,b)=(A_j,1)\) and the second
\((1,B^{-1})\). Since \(B^{-1}=A_{j+2}u_{j+2}\), their sum is
\(d_i'T_i\) of Theorem~\ref{thm:level-fourier-slide-chain}
with the present \(L_j,\pi_j\). That theorem yields \(t_N\)
as the product of \(d_i'\), of the positive integers \(o_i\),
and of the factors \(\gamma_l\) occurring there.
Each \(o_i\) is prime to \(N\); every \(\gamma_l\) has
\(N\)-adic absolute value at most one since \(d_{l+1}\ge d_l\).
Lemma~\ref{lem:actual-level-character-factor} bounds
\(\log_N|\iota d_i'|_N\) by \(2b(|U|+|U'|)=2fb\), using
Proposition~\ref{prop:height-cost}. Multiplication over \(m/2\)
pairs proves the assertion.
\end{proof}

\begin{theorem}\label{thm:actual-level-prime-column-bound}
Use the scale, vacua and \(\Psi\) of
Theorem~\ref{thm:actual-unmarked-chain-divisibility}. For a
ramified column \(\phi\) suppose the shifted column
\(\xi=F_0^{m/2}\phi\) satisfies
\(|\xi(g)|_{0,g A_0}\le1\) at every \(g\).
Then for every terminal argument \(h\)
\[
  |(B_\Psi\phi)(h)|_{Y,h}
    \le N^{mfb+f\sum_{j<m}d_{j+1}}\le N^{2mfb}
\]
at all sufficiently large scales, uniformly in the twist.
\end{theorem}
\begin{proof}
Apply Proposition~\ref{prop:actual-F-chain-level-identity}
to the expansion of \(\mathcal S\mathcal P w_\Psi=w_\Psi\)
in Theorem~\ref{thm:projected-chain-good-coefficients}. Each
summand of the resulting chain has modified arrow
\(\Phi\mapsto T_j\Phi=\Omega_{j,f}(A_j,A_{j+1})[b_j]\Phi\).
The tensor \((\otimes T_j)w_\Psi\) is jointly invariant for
the same column level. To verify this assertion at \(N\),
each original level element is integral, acts on the box
by an integral change of variables with modulus one and
trivial character, and hence fixes the box and its coefficient
projection. After conjugation by \(A_j,A_{j+1}\) the two
matrices are in \(M_j,M_{j+1}\); their left and right action
fixes \(b_j\) since each fixes the last basis vector modulo \(N\).
This proves that \(T_j\) intertwines the two level actions
at \(N\). At the complementary places the operations commute
because they concern separate oscillator factors. Tensoring
shows that the original invariance implies the assertion.

By Proposition~\ref{prop:actual-F-chain-level-identity}, \(\xi\)
has the required column level. Theorem~\ref{thm:joint-column-class-evaluation}
therefore expresses this modified chain at \(h'\in H_{m,f}\)
as the sum over tuples of classes. On each arrow at \(g,g'\)
the test is \(\Omega_{j,f}(g A_j,g' A_{j+1})[b_j]D_Xv^\circ\).
Proposition~\ref{prop:one-column-level-cut-bound} bounds its
operator norm in exactly these shifted frames by \(N^{f d_{j+1}}\).
The coefficients of \(\Psi\) are integral at \(N\) since \(N\nmid Q\).
The ultrametric inequality gives the product of these bounds
times the bound for \(t_N\), uniformly in \(h'\).
Finally take \(h'=\operatorname{diag}(h,1)\). In the frame
\(h'A_m\) the last entries at \(w_\tau\) are \(N\) exactly
when \(w_\tau\mid\pi_N\), equivalently when
\(\iota\tau(\bar\pi_N)\) is a nonunit.
Proposition~\ref{prop:terminal-restriction-gauge} now
gives the indicated norm on \(Y_m\). Proposition~\ref{prop:one-column-level-cut-bound}
also proves \(d_{j+1}\le b\) eventually, which gives the
last inequality.
\end{proof}


\section{The index of the initial Fourier condition}

\begin{definition}\label{def:initial-shifted-integral-columns}
Fix an embedding \(\iota\) at \(N\) and use
\(O_\iota=\{z\in\C:|\iota z|_N\le1\}\). For a finite
frame \(C\) in the initial unitary group put
\[
       \mathcal E(C)=
         \{\phi\in{\rm Col}:|\phi(g)|_{0,gC}\le1
                                    \text{ at every finite }g\}.
\]
Write \(c_0=\dim_\C{\rm Col}\) and let \(A_{\rm top}\) be
the set \(A\) of the \(U\) nontrivial symmetric-power factors in
Definition~\ref{def:columns}. For \(V\subseteq S\) set
\[
       h(V)=\#\{\tau\in A_{\rm top}:w_\tau\in V\},\qquad
       t(V)=\frac{a}{k}\,h(V)c_0.
\]
Use \(A_0=\Delta_0(S\cap R)\) from
Definition~\ref{def:level-shifted-column-action}.
\end{definition}

\begin{proposition}\label{prop:initial-shift-frame-exact-index}
For every \(V\subseteq S\), set \(D=\Delta_0(V)\).
The module \(\mathcal E(A_0)\) is contained in
\(\mathcal E(A_0D^{-1})\). In suitable common class coordinates
these two modules are respectively
\(\bigoplus_\lambda N^{s_\lambda(V)}O_\iota\)
and \(O_\iota^{c_0}\). The index exponents are
\[
  s_\lambda(V)=\sum_{\substack{\tau\in A_{\rm top}\\w_\tau\in V}}
                       m_{\tau,k},
       \qquad \sum_\lambda s_\lambda(V)=t(V)\in\Z_{\ge0},
\]
where \(\lambda\) runs over a class and a tensor monomial
\(\prod_{\tau,i} e_{\tau,i}^{m_{\tau,i}}\) of degree \(a\) in
each top factor. Consequently the least integer \(s\ge0\)
for which
\(\mathcal I(\mathcal E(A_0D^{-1}),\mathcal E(A_0);N^s)\)
holds is \(t(V)\).
\end{proposition}
\begin{proof}
In the \(S\)-coordinates the level conjugated by \(A_0^{-1}\)
is \(M_0\) by Proposition~\ref{prop:actual-F-chain-level-identity}.
The level conjugated by \((A_0D^{-1})^{-1}\) is \(DM_0D^{-1}\).
It too is integral: on \(V\) a matrix in \(M_0\) has an
upper-right block divisible by \(N\); conjugation divides
that block by \(N\) and multiplies the lower-left block by
\(N\). Apply the same argument to its inverse.
Hence for either frame \(C\) the gauge at \(gC\) is unchanged
by replacing \(g\) by \(gk'\), \(k'\in K_c\), on a column:
the frame then has an integral unit right factor. Left rational
invariance of the gauge is the calculation of
Theorem~\ref{thm:local-gauge-class-matrix}. Testing one representative
in each class therefore suffices.

For each class choose a representative \(g\) so that
\((g A_0D^{-1})_N\) belongs to the intersection of the maximal
integral unitary subgroup \(K_N^0\) with \(D K_N^0D^{-1}\).
Indeed this intersection is an open neighborhood of the
identity, and Lemma~\ref{lem:definite-one-prime-unitary-approximation}
allows left multiplication of any representative by a rational
element to enter it. Then \((D^{-1}g A_0)_N\) also is integral
unitary. Thus in these representatives the \(A_0D^{-1}\)
condition is the ordinary integral condition. The \(A_0\)
condition is membership in the image under the weight action
of \(D\) of the ordinary integral module, since \(g A_0\)
is \(D\) times an integral frame at \(N\). For arbitrary
complex coordinate vectors this says precisely that their
entries at \(\lambda\) belong to \(N^{s_\lambda(V)}O_\iota\):
the weight action at zero has determinant exponent zero and
\(D\) multiplies just the selected last letters by \(N\).
These coordinate vectors run independently through every
complex tuple by Proposition~\ref{prop:column-dim}.

For each fixed \(\tau\) permutations of the \(k\) letters
give bijections on the monomial indices. The sum of
\(m_{\tau,i}\) over all classes and monomials is therefore
the same for each \(i\); its sum over \(i\) is \(a c_0\).
This evaluates the exponent sum as claimed, including when
\(V\) is empty. The ring \(O_\iota\) is a valuation ring
because one of \(z,z^{-1}\) has absolute value at most one
for any \(0\ne z\in\C\). Lemma~\ref{lem:valuation-index}
now gives the assertion about the least exponent.
\end{proof}

\begin{theorem}\label{thm:actual-F-index-and-transfer}
For \(p_0=m/2\) set
\[
 E_N=t(S\cap R)+p_0t(S),\qquad
 \mathcal E_*=\{\phi\in{\rm Col}:F_0^{p_0}\phi\in\mathcal E(A_0)\}.
\]
Then \(\mathcal I(\mathcal E(1),\mathcal E_*;N^{E_N})\)
holds and, uniformly at the prescribed scales,
\[
 E_N\ \le\ (p_0+1)\frac{aU}{k}c_0
              \ \ll\ x^{-4}mfb\,c_0 .
\]
For every finite list of ramified columns in \(\mathcal E(1)\)
one can therefore recombine the list by a matrix over \(O_\iota\)
of determinant dividing \(N^{E_N}\) so that every resulting
column satisfies the bound in Theorem~\ref{thm:actual-level-prime-column-bound}.
The same matrix works for every permitted finite twist at this
embedding.
\end{theorem}
\begin{proof}
Write \(F=F_0\) and \(u=\Delta_0(S)^{-1}\). For
\(\phi\in\mathcal E(A_0)\) evaluate
\[
      (F\phi)(g)=
          \sum_{y:\,y_*\in k_N^\times}\psi(y_*)
                           \phi(g A_0 u\kappa_y A_0^{-1}).
\]
The input frame for each displayed value is \(g A_0u\kappa_y\).
Right multiplication by \(\kappa_y\) preserves its gauge
because \(\kappa_y\) is integral with integral inverse.
The coefficients are roots of unity. Thus \(F\phi\) is
integral in the frame \(gA_0u\), and it has column level
by Proposition~\ref{prop:actual-F-chain-level-identity}.
We have proved \(F\mathcal E(A_0)\subseteq\mathcal E(A_0u)\).

For any finite list in \(\mathcal E(1)\), first apply
Proposition~\ref{prop:initial-shift-frame-exact-index} with
\(V=S\cap R\) to obtain a list in \(\mathcal E(A_0)\) at
cost \(N^{t(S\cap R)}\). Given a current list whose \(i\)-th
Fourier powers lie in \(\mathcal E(A_0)\), their next powers
lie in \(\mathcal E(A_0u)\). Apply that same proposition with
\(V=S\) to these next powers; its recombination matrix also
acts on the original list by linearity of \(F\). The newly
recombined next powers belong to \(\mathcal E(A_0)\) and the
cost divides \(N^{t(S)}\). Repeating \(p_0\) times and
multiplying determinants gives the property of
Definition~\ref{def:det-index}. We used no inverse to \(F\).
By Definition~\ref{def:initial-shifted-integral-columns},
\(t(S\cap R)\le t(S)=aUc_0/k\).
Theorem~\ref{thm:data} and Definition~\ref{def:heights} give
\(a\le 2xb,\ k\gg x^5,\ U\le f\) for sufficiently large \(x\),
which gives the second bound. Recombination by complex scalars
preserves the ramified subspace. The preceding level-prime
theorem therefore applies to every output. The index calculation
uses only the fixed initial representation and the prime \(N\),
so the same recombination gives all the twists as asserted.
\end{proof}


\section{Split Hecke polynomials on a definite theta chain}

\begin{lemma}\label{lem:finite-rank-subspace-polynomial}
Let \(k_q\) be a field of \(q\) elements. Write \(b(n,i)\)
for the number of codimension-\(i\) subspaces of \(k_q^n\),
with value zero outside \(0\le i\le n\). Then
\[
       \sum_i b(n,i)q^{i(i-1)/2}Z^i
                   =\prod_{\ell=0}^{n-1}(1+q^\ell Z).
\]
The number of codimension-\(i\) subspaces containing a given
\(h\)-dimensional subspace is \(b(n-h,i)\).
\end{lemma}
\begin{proof}
Subspaces containing a fixed subspace correspond bijectively
under quotient to subspaces of the quotient, with codimension
preserved. Split \(k_q^n=H\oplus k_q e\), where \(H\) is
a hyperplane. A codimension-\(i\) subspace in \(H\) has
codimension \(i-1\) there. One outside \(H\) is specified
by its codimension-\(i\) intersection \(W\) in \(H\) and
the class modulo \(W\) of an \(h\) such that \(e+h\) lies
in the subspace. This gives \(b(n,i)=b(n-1,i-1)+q^i b(n-1,i)\).
For the weighted polynomial \(P_n\) it says
\(P_n(Z)=(1+Z)P_{n-1}(qZ)\), including the end coefficients.
Induction from \(P_0=1\) proves the product.
\end{proof}

\begin{proposition}\label{prop:actual-served-hecke-polynomial}
Let \(p\) be served, choose \(\lambda\in\Lambda\) as in
Theorem~\ref{thm:prime-marking} for the residue embeddings
induced by \(\iota\tau\), and assume \(p\nmid Q\).
Its residue cardinality is \(q=p^L\) and \(p\) generates
the local maximal ideal. Write \(T_{j,i}\) for the operators
of Definition~\ref{def:actual-spherical-polynomial} on \(X_j\)
at \(\lambda\), and
\[
       P_j(X)=\sum_{i=0}^{d_j}(-X)^i q^{i(i-1)/2}T_{j,i}.
\]
Write \(z_j\) for the value of \(\mu_{j+1}\) on the norm-one
finite idele with component \(p\) at \(\lambda\), \(p^{-1}\)
at \(\bar\lambda\), identity elsewhere. On the arrow
\(j\to j+1\) put \(\mathcal L_j=\mathcal L^c_{j,\Phi_j}\)
for any translate \(\Phi_j=D_{s_j}v_j^\circ\) occurring in
Theorem~\ref{thm:projected-chain-good-coefficients}.
Then \(|\iota z_j|_p=1\), and on all inputs fixed by the local
integral group at \(\lambda\) one has the polynomial identity
\[
  P_{j+1}(X)\mathcal L_j
    =\prod_{\ell=0}^{d_{j+1}-d_j-1}
          (1-z_j q^{d_j/2+\ell}X)\,
          \mathcal L_j P_j(q^{(d_{j+1}-d_j)/2}X).
\]
The output of \(\mathcal L_j\) is fixed by the local integral group.
\end{proposition}
\begin{proof}
In Lemma~\ref{lem:inert-prime-fibres} the residue field is
generated by a root of degree two and a primitive \(N_1\)-th
root. The first is fixed by the even powers of Frobenius,
the second exactly by powers \(s\) with \(p^s=1\bmod N_1\).
The least positive even such \(s\) equals \(L\), since
\(L\) is even by Theorem~\ref{thm:prime-marking} and
\(p^s=(-p)^s\) for even \(s\). These conditions compute the
degree: in a finite field of \(p^a\) elements the \(a\)-th
Frobenius power is identity by the multiplicative group order,
and no smaller positive power can fix all its elements by
the root bound for \(X^{p^s}-X\). Also \(p\nmid N_1d\),
and the two defining polynomials for \(O_E/pO_E\) have
simple roots over an algebraic closure. They are the quadratic
polynomial for \(\omega\) and the cyclotomic polynomial in
Proposition~\ref{prop:local-cyclo-different}. Factoring the
first over \(\mathbb F_p\) and then the second over each
factor expresses \(O_E/pO_E\) as a product of fields.
The ideal factorization of Proposition~\ref{prop:general-ideal-factorization}
therefore gives valuation one for \(p\) at each factor.

The two sets of embeddings \(\iota\tau\) at \(\lambda,\bar\lambda\)
are the two fibres of equal size in the \(h\)-orbit of
Theorem~\ref{thm:prime-marking}. This orbit avoids the distinguished
set \(D\), so its nominal weight is constant \(a/x_1\); the
rounded \(w_{j+1}\) is constant on it by Definition~\ref{def:heights}.
The two contributions of opposite valuations cancel in the
character power formula of Theorem~\ref{thm:torus-characters-elementary}.
This proves the unit assertion. By \(p\nmid QgN\) and the
displacement description of Theorem~\ref{thm:projected-chain-good-coefficients},
the local component of \(\Phi_j\) is a scalar multiple of
the vacuum. In the split matrix realization of
Proposition~\ref{prop:finite-polarized-vacuum} its local
function is the integral-box indicator. The same entrywise
duality proof as in Lemma~\ref{lem:level-cut-vacuum-fourier}
works over the valuation ring \(O_\lambda\), since the
normalized lattice is self-dual. The character on integral
stage determinants is trivial as \(p\nmid Qg\). Thus the
box, and hence the output, is target-invariant by
Theorem~\ref{thm:definite-jet-class-operator}.

For a rank \(a\), the right cosets of
\(K{\rm diag}(pI_i,I_{a-i})K\), \(K={\rm GL}_a(O_\lambda)\),
are classified by their column lattices. They are exactly the
inverse images of codimension-\(i\) subspaces \(U\) in the
residue space. In fact lift a basis of \(U\) completed to a
residue basis to \(k\in K\), whose first \(a-i\) columns
give the chosen basis of \(U\). The symmetric representative
\(g_U=k\,{\rm diag}(I_{a-i},pI_i)\,k^t\) has this column
lattice and also the same row lattice after transpose.
Two matrices with the same column lattice differ on the
right by an element of \(K\), by applying both inverse basis
changes. This proves the claim and gives representatives with
both required properties.

In the target Hecke sum, translation by \(g_U\) on the output
changes the box to the indicator of \(g_U^{-1}T\) integral,
multiplied by \(z_j^i q^{i d_j/2}\). In the source sum
translation by \(g_U\) on the input changes the box to
the indicator of \(Tg_U^{-1}\) integral, multiplied by
\(z_j^i q^{i d_{j+1}/2}\). These equalities follow by
Theorem~\ref{thm:definite-jet-class-operator} with target
\(g_U\), respectively source \(g_U^{-1}\), and the two
determinant/Jacobian factors proved in
Lemma~\ref{lem:actual-level-character-factor}. The latter
calculation uses only split coordinates, even ranks and the
unramified unit character at the chosen pair, so it gives
these factors with residue cardinality \(q\) and local
uniformizer \(p\) here.

Both conditions imply \(T\) integral. For an integral \(T\)
of residue rank \(h\), the target condition requires \(U\)
to contain its column space, and the source condition requires
\(U\) to contain its transposed row space. The two counts in
degree \(i\) are \(b(d_{j+1}-h,i)\) and \(b(d_j-h,i)\).
By Lemma~\ref{lem:finite-rank-subspace-polynomial} the target
polynomial on the box is
\(\prod_{0\le\ell<d_{j+1}-h}(1-z_jq^{d_j/2+\ell}X)\).
Removing its first \(d_{j+1}-d_j\) factors leaves exactly
the source polynomial at \(q^{(d_{j+1}-d_j)/2}X\).
Outside the box all terms vanish. Equality of local tests
gives the asserted identity by linearity of the lift.
\end{proof}

\begin{corollary}\label{cor:actual-prefix-hecke-coefficients}
Let \(1\le J\le m\), \(\mathcal C_J=\mathcal L_{J-1}\cdots\mathcal L_0\),
and \(i_0=\lceil\log_p m\rceil\). At every such \(\lambda\)
there are \(a_i\in O_\iota=\{z:|\iota z|_p\le1\}\),
independent of all displacements in the support of \(\Psi\), with
\[
  p^{\Delta_J}\mathcal C_J T_{0,i_0}
      =\sum_{i=0}^{i_0}a_i T_{J,i}\mathcal C_J,\qquad
  \Delta_J=L\left(\frac{i_0(d_J-d_0)}2+
                              \binom{i_0}{2}\right)
                 \le \frac{\delta_p}{2}
\]
on \(\lambda\)-spherical inputs, where
\(\delta_p=i_0 f(n_m+i_0)\) as in Proposition~\ref{prop:marked-cost}.
\end{corollary}
\begin{proof}
Iterating Proposition~\ref{prop:actual-served-hecke-polynomial}
gives \(P_J(X)\mathcal C_J=H(X)\mathcal C_J P_0(sX)\)
with \(s=q^{(d_J-d_0)/2}\). Here \(H\) is a product of the
displayed linear factors at arguments \(q^{(d_J-d_{l+1})/2}X\)
for stages \(l<J\). Its coefficients are in \(O_\iota\)
by the unit assertion, and \(H(0)=1\). Its inverse exists in
\(O_\iota[[X]]\): the constant is one, and each successive
coefficient is minus a finite sum of products of earlier ones
with positive-degree coefficients of \(H\). Denote these inverse
coefficients by \(h_n\). Compare coefficients of \(X^{i_0}\)
after multiplying by \(H^{-1}\). This gives the assertion with
\(a_i=(-1)^{i+i_0}h_{i_0-i}q^{i(i-1)/2}\), since
\(s^{i_0}q^{i_0(i_0-1)/2}=p^{\Delta_J}\). Here \(i_0\le k\)
eventually by \(m=O(\sqrt x)\) and \(k\gg x^5\).
Finally \(L\le f=|G|\), \(d_J-d_0\le n_m\) since \(d_m=n_m+1\)
and \(d_0=k\ge1\). These inequalities give the stated bound.
\end{proof}


\section{The jet matrix at marked neighbours}

\begin{definition}\label{def:served-scaled-neighbours}
Let \(p,\iota,\Lambda,V\) have the meanings of
Proposition~\ref{prop:actual-served-hecke-polynomial} and
Theorem~\ref{thm:prime-marking}. Set
\[
       A_\lambda=a\,\#\{\tau\in A:\mathfrak p_\tau=\lambda\},
       \qquad Y_\lambda=p^{A_\lambda}T_{0,i_0}
                       \quad\text{at }\lambda,\quad
       Y=\prod_{\lambda\in\Lambda}Y_\lambda.
\]
Hecke operators in different factors commute on the common
spherical space: each is a finite sum, and each pair of terms
has disjoint local support. Thus the product on columns may
be taken in any order. For \(0\le j<m\) write \(r_\tau(j)\)
for the number \(r_\tau\) of columns of height greater than
one in Definition~\ref{def:stage-diagrams}.
\end{definition}

\begin{lemma}\label{lem:served-scaled-integral-column}
Each \(Y_\lambda\) carries columns integral in the ordinary
\(p\)-adic gauge to columns integral in that gauge.
\end{lemma}
\begin{proof}
At this pair \(K_c\) contains the full integral group by
Definition~\ref{def:columns} since \(p\nmid gN\). A left
integral translation permutes the right cosets of a double
coset, and the resulting right unit factors fix the vector.
Thus the output retains this local level, and all other
level actions commute with the sum. For an individual
representative \(s=k_1 d k_2\), the local diagonal \(d\)
has \(p\) on \(i_0\) entries at \(\lambda\) and the inverse
on the conjugate entries, all other coordinates being one.
For \(\phi\) integral,
\[
 D_0(g_p)^{-1}\iota\phi(gs)
       =D_0(s_p)\,\big(D_0((gs)_p)^{-1}\iota\phi(gs)\big).
\]
Both unit factors of \(D_0(s_p)\) act isometrically. The
negative powers in \(D_0(d)\) come only from the tops for
which \(\mathfrak p_\tau=\lambda\), as the weight evaluates
at the conjugated embedding in Theorem~\ref{thm:local-gauge-class-matrix}.
Each such factor is a symmetric power of degree \(a\) with
determinant exponent zero, so no entry of its diagonal has
valuation below \(-a\). The entire loss is at most
\(A_\lambda\). Multiplying by \(p^{A_\lambda}\) cancels it,
and summing uses the ultrametric inequality.
\end{proof}

\begin{theorem}\label{thm:served-neighbour-jet}
For each \(\lambda\in\Lambda\), let \(s_\lambda\) be any
representative in a basic double coset of degree
\(0\le i_\lambda\le i_0\) in rank \(d_{j+1}\).
Embed \(s=\prod_\lambda s_\lambda\) in the target group
with its inverse transpose at the conjugate factors. Let
\(\Phi\in\mathcal D_p(1)\), \(c=-\bar\pi_g u\).
For all finite frames \(g,h\) and every coefficient vector \(v\),
\[
 \big|J_j^c(\Omega_{j,f}(g,hs)\Phi)v\big|_{j+1,h}
   \le p^{\,\sum_\lambda A_\lambda
               +i_0\sum_{\tau\in V}r_\tau(j)-J_p(j,V)}
          |v|_{j,g}.
\]
\end{theorem}
\begin{proof}
Write the \(p\)-component \(s=k'd k''\) with integral
unitary \(k',k''\), where \(d\) has \(p\) on coordinates
\(I_\lambda\) of size \(i_\lambda\) at \(\lambda\) and
\(p^{-1}\) on those coordinates at \(\bar\lambda\).
Use the identity at the other factors. These decompositions
are the defining basic double cosets.
Choose a rational source element \(\gamma\) with \(\gamma g\)
integral at \(p\). Choose a rational target element \(\epsilon\)
with \(v_0=(\epsilon h)_p k'd\) in
\(K_p^0\cap d^{-1}K_p^0d\) at the target rank. Both choices
are possible by Lemma~\ref{lem:definite-one-prime-unitary-approximation}.
Then \((\epsilon h s)_p=v_0 k''\) is integral, and
\((\epsilon h)_p^{-1}=k'dv_0^{-1}\) lies in \(K_p^0 d\).
Rational covariance as in Theorem~\ref{thm:local-gauge-class-matrix}
therefore reduces the desired bound to the max norm of
\(D_{j+1}(d)\iota J_j^c(\Phi')\), where
\(\Phi'=\Omega_{j,f}(\gamma_f g,\epsilon_f hs)\Phi\) belongs
to \(\mathcal D_p(1)\) by Proposition~\ref{prop:local-test-unitary-preservation}.
The transformed input has ordinary max norm \(|v|_{j,g}\).

For a target tableau \(T'\) write \(n_\tau(I)\) for the
number of its entries at \(\tau\) in \(I\).
The determinant part of \(D_{j+1}(d)\) has valuation zero.
Indeed the two fibres above each pair have equal size and
the same value of \(d_j/2-w_{j+1}(\tau)\), as proved for the
weights in Proposition~\ref{prop:actual-served-hecke-polynomial}.
The tableau part is diagonal on \(T'\): a flag minor acquires
one factor from each selected row, by the definition in
Proposition~\ref{prop:rectangular-flag-law}. Its valuation is
\[
    \sum_{\lambda\in\Lambda}
       \left(\sum_{\bar{\mathfrak p}_\tau=\lambda}
                      n_\tau(I_\lambda)
            -\sum_{\mathfrak p_\tau=\lambda}n_\tau(I_\lambda)\right),
\]
with the signs as in Lemma~\ref{lem:served-scaled-integral-column}.

Choose the charge \(\gamma_{\tau,h',b'}\) of
Theorem~\ref{thm:schur-jet} as follows. When
\(\mathfrak p_\tau=\lambda\in\Lambda\), assign one at
height \(h'=1\) to each \(b'\notin I_\lambda\) and zero
elsewhere. When \(\bar{\mathfrak p}_\tau=\lambda\in\Lambda\),
assign one to \(b'\in I_\lambda\) at every height and zero
elsewhere. Give zero at all other places. For any relevant
odd \(e\), \(\sigma=(-p^e)\tau\) induces the conjugate prime
of \(\tau\) by Lemma~\ref{lem:inert-prime-fibres}.
If \(\tau\in V\), the two charges at height one on its side
and at any height on the other side cover every \(b'\).
If \(\sigma\in V\) and its height is one, the same two sets
in reverse order cover every \(b'\). These are all the
charge conditions since \(p\) is inert.

Apply Theorem~\ref{thm:standard-test-nonsplit-minors} with
\(a=\bar\pi_g\), a unit at \(p\). In the valuation of the
scaled row \(T'\) subtract its charge \(m_{T'}\) and add the
displayed diagonal valuation. All terms from the conjugates of
the selected fibres cancel. At \(\tau\in V\) each height-one
entry now costs exactly one whether or not it is in \(I_\lambda\),
giving \(a{\bf1}_A(\tau)\) in total. A taller column costs
one for each of its entries in \(I_\lambda\). The entries
increase strictly down the column, so this costs at most
\(|I_\lambda|\le i_0\). Summing over columns and places gives
the stated valuation for every matrix entry and therefore
the operator inequality.
\end{proof}


\section{Divisibility on the marked Hecke image}

\begin{theorem}\label{thm:actual-served-hecke-image-gain}
Let \(p<x_1\) be served and fix \(\iota,\Lambda,V\) as
in Definition~\ref{def:served-scaled-neighbours}. Write
\(\mathcal E_p\) for the columns integral in the ordinary
gauge at \(p\), and put
\[
    s_p=\lceil 2mfa/p\rceil,\qquad
    \mathcal E_p' =Y^m\mathcal E_p+p^{s_p}\mathcal E_p .
\]
Then \(\mathcal E_p'\subseteq\mathcal E_p\). For every twist
with \(p\nmid Q\), every coefficient function \(\Psi\) of
Theorem~\ref{thm:projected-chain-good-coefficients} with the
scale and normalization of Theorem~\ref{thm:actual-unmarked-chain-divisibility},
and every \(\xi\in\mathcal E_p'\), one has
\[
          |(B_\Psi\xi)(h)|_{Y,h}
               \le p^{-G_p+(30\ell+3)fa/p}
\]
at every terminal argument \(h\), for all sufficiently large \(x\).
\end{theorem}
\begin{proof}
The module inclusion follows by Lemma~\ref{lem:served-scaled-integral-column}.
Consider a tuple \(s\) in the support of \(\Psi\) and put
\(\Phi_j=D_{s_j}v_j^\circ\). Each such datum is in
\(\mathcal D_p(1)\) at this prime by
Theorem~\ref{thm:projected-chain-good-coefficients} since
\(p\nmid gNQ\). For a marked \(\lambda\) and output rank
\(d_{j+1}\) use the local oscillator representation
\(b\mapsto\Omega_{j,f}(1,b)\). Denote by \(T^{\rm osc}_{j+1,i,\lambda}\)
the sum in Definition~\ref{def:actual-spherical-polynomial} using
this representation. Take any common choices of right coset
representatives. On all the local fixed vectors the operator
is independent of these choices. At each stage put
\[
 \mathcal T_j=\prod_{\lambda\in\Lambda}
   \left(p^{A_\lambda-\Delta_{j+1,\lambda}}
       \sum_{i=0}^{i_0} a_{i,j,\lambda}
                     T^{\rm osc}_{j+1,i,\lambda}\right),
\]
where the coefficients and \(\Delta_{j+1,\lambda}\) come from
Corollary~\ref{cor:actual-prefix-hecke-coefficients}, with
\(J=j+1\). They are independent of the tuple \(s\).
A prefix of the original chain will be denoted by \(\mathcal C_j\)
with \(\mathcal C_0=1\).
Translation covariance of Theorem~\ref{thm:definite-jet-class-operator}
gives
\[
 \mathcal L^c_{j,\mathcal T_j\Phi_j}\mathcal C_j\eta
                         =\mathcal C_{j+1}Y\eta
\]
for every initial column \(\eta\). In detail, moving a single
target Hecke sum outside the last arrow gives the same sum on
the output of \(\mathcal C_{j+1}\). The corollary then moves
it to \(p^{A_\lambda}T_{0,i_0,\lambda}\) on the initial
column after the displayed scalar is included. Operators at
the remaining primes commute with this move, and the resulting
initial column is spherical there. Iteration proves the identity.
Induction on prefix length consequently identifies the chain
with every datum replaced by \(\mathcal T_j\Phi_j\), applied
to \(\eta\), with the original chain applied to \(Y^m\eta\).

We justify evaluation of this replacement at the original class
level. Every \(\Phi_j\) in the support is spherical under both
integral actions at each marked pair, since its local component
is a multiple of the box by
Proposition~\ref{prop:actual-served-hecke-polynomial}.
Source translations commute with the target sum. A target
integral translation permutes its right cosets with integral
right factors; these fix each spherical vector. Thus
\(\mathcal T_j\) commutes with the level actions on the
spherical subspace. Translations of the rest of the level
at other pairs commute term by term. It follows from the
joint invariance in Theorem~\ref{thm:projected-chain-good-coefficients}
that \((\otimes_j\mathcal T_j)w_\Psi\) is jointly invariant,
so Theorem~\ref{thm:joint-column-class-evaluation} applies to it.

For a tuple of representatives in this sum, expand
\(\mathcal T_j\) on each individual \(\Phi_j\). Its Hecke
translations range precisely over the cases in
Theorem~\ref{thm:served-neighbour-jet}. The coefficients
\(a_{i,j,\lambda}\) are \(p\)-integral and multiplication
by \(p^{\sum_\lambda(A_\lambda-\Delta_{j+1,\lambda})}\)
absorbs the \(\sum_\lambda A_\lambda\) in that theorem.
The resulting operator bound on this stage is
\[
  p^{\,\sum_\lambda\Delta_{j+1,\lambda}
                  +i_0\sum_{\tau\in V}r_\tau(j)-J_p(j,V)} .
\]
All sums in the integral class formula are finite and each
\(\Psi(s)\) is integral at \(p\). The product of these bounds
therefore bounds the output on \(\eta\in\mathcal E_p\) by the
ultrametric inequality. Since
\(\Delta_{j+1,\lambda}\le\delta_p\), Proposition~\ref{prop:marked-cost}
gives the exponent \(-G_p+(30\ell+3)fa/p\). Proposition~\ref{prop:terminal-restriction-gauge}
passes this estimate to the restricted chain since \(p\ne N\).

For the second summand of \(\mathcal E_p'\), use
Theorem~\ref{thm:actual-unmarked-chain-divisibility}.
By Definition~\ref{def:collisions} and geometric summation,
\(mC_p(\varnothing)\le mUa/(p-1)\le 2mfa/p\) for \(p\ge2\).
Multiplication by \(p^{s_p}\) cancels this loss, leaving
exponent at most \(-G_p+30f\ell a/p\). Both bounds pass to
the sum of these modules by linearity and the ultrametric
inequality.
\end{proof}

\begin{corollary}\label{cor:local-index-last-served-condition}
Let \(\beta_p=\lceil G_p-e_p\rceil\) with the conventions of
Proposition~\ref{prop:exact-normalization}. At every prime \(p\)
there is an explicit \(O_\iota\)-module of columns on which
\(|(B_\Psi\phi)(h)|_{Y,h}\le p^{-\beta_p}\) for all ramified
\(\phi\) in the module and all twists with \(p\nmid Q\). At
a served prime take \(\mathcal E_p'\), at \(N\) take
\(\mathcal E(1)\cap\mathcal E_*\) of
Theorem~\ref{thm:actual-F-index-and-transfer}, and elsewhere
take \(\mathcal E_p\). The determinant-index assertion
\[
      \mathcal I(\mathcal E_p,\mathcal E'_p;p^{l_p})
            \quad\text{for some integer }0\le l_p\le e_p c_0
\]
at each served \(p\) suffices for the index condition at every
prime in Proposition~\ref{prop:exact-normalization} and its
application to the local recombinations.
\end{corollary}
\begin{proof}
At a served prime \(e_p=3fa\ell^2/p\);
\((30\ell+3)fa/p\le e_p-1\) eventually because
\(fa/p\ge fa/x_1\) tends to infinity by Theorem~\ref{thm:data}
and the scale bounds in Lemma~\ref{lem:diagram-size}.
Theorem~\ref{thm:actual-served-hecke-image-gain} now gives
the stated rounding. At an unserved small nonsplit prime
Theorem~\ref{thm:actual-unmarked-chain-divisibility} costs at
most \(2e_p/3\le e_p-1\) eventually by the same bounds.
At \(N\) one has \(G_N=0,\ e_N=4mfb\) by
Proposition~\ref{prop:loss}. Theorem~\ref{thm:actual-level-prime-column-bound}
costs at most \(2mfb\le e_N-1\); the determinant exponent of
Theorem~\ref{thm:actual-F-index-and-transfer} is at most \(e_N c_0\).
Its recombinations of integral columns also stay in
\(\mathcal E(1)\), giving the asserted intersection.
At all other primes \(G_p=e_p=0\); integrality is
Theorem~\ref{thm:actual-unmarked-chain-divisibility} since \(d\)
is in the small nonsplit range eventually. Module equality gives
index exponent zero at those primes and at the unserved ones.
Definition~\ref{def:det-index} then gives the asserted sufficiency.
\end{proof}


\section{The class matrix of one marked Hecke operator}

\begin{definition}\label{def:served-column-cycle}
Fix \(p,\lambda\) as in Definition~\ref{def:served-scaled-neighbours}
and an embedding \(\iota\) at \(p\). To avoid confusion with the
global auxiliary level, write \(D=L/o,\ P=p^o\), where \(o,L\)
are the orders of \(-p\) in Definition~\ref{def:collisions}.
Choose a top \(\tau_0\) with \(\mathfrak p_{\tau_0}=\lambda\).
For \(s\in\Z/D\Z\), put \(\tau_s=(-p)^{os}\tau_0\bmod N_1\).
Let \(W=(O_E)_\lambda\), the fractions of elements of \(O_E\)
with denominator outside \(\lambda\). Set \(O=\{z\in K_p:
|z|_p\le1\}\). Define maps \(W\to O\) by
\[
 \sigma_s(x)=
 \begin{cases}
   \iota\tau_s(x)&s\text{ even},\\
   \iota\tau_s(\bar x)&s\text{ odd}.
 \end{cases}
\]
On \(k\)-by-\(k\) invertible matrices over \(W\) define the
representation in the integral evaluation coordinates of
Definition~\ref{def:weyl-coord} by
\[
 R_\lambda(g)=
  \bigotimes_{s\ {\rm even}}S(\sigma_s(g^{-t}))
       \ \otimes\
  \bigotimes_{s\ {\rm odd}}S(\sigma_s(g)),
 \qquad g^{-t}=(g^{-1})^t,
\]
with an arbitrary fixed ordering of the factors. Here the matrix entries
and their images in \(K_p\) come from algebraic numbers, so matrices
over \(O\) in this formula are also matrices over \(\iota(O_\iota)\).
\end{definition}

\begin{lemma}\label{lem:served-actual-frobenius-cycle}
The \(D\) embeddings listed are exactly the tops inducing
\(\lambda\) or \(\bar\lambda\). The first prime is induced for
even \(s\) and the second for odd \(s\), and \(D\) is even.
The ring \(W\) has residue field of cardinality \(P^D\) and
uniformizer \(p\). The \(\sigma_s\) are ring homomorphisms and
\(\sigma_{s+1}(x)-\sigma_s(x)^P\in pO\) for every \(x\in W\),
including the index crossing \(s=D-1\).
\end{lemma}
\begin{proof}
By Theorem~\ref{thm:prime-marking} the orbit returns to the
tops exactly at the multiples of \(o\); its full period is
\(L\), with \(m\mid L/o\). Since \(o\) is odd each return
exchanges the two prime fibres by Lemma~\ref{lem:inert-prime-fibres}.
These facts prove the enumeration and the parity assertion.
Each \(\sigma_s\) uses \(\lambda\) on \(E\) by this description,
so every denominator in \(W\) maps to a unit. The assertion
about \(W\) follows from the calculation of the residue degree
and unramifiedness in Proposition~\ref{prop:actual-served-hecke-polynomial}.

To prove the congruence, compare the two induced maps to \(O/pO\).
On \(\zeta_{N_1}\) the two exponents agree even before
reduction: the additional minus in passing to \(\tau_{s+1}\)
is canceled by the change of conjugation. On \(\omega\) they
agree since \(X^2-X+(1+d)/4\) is irreducible in the inert
residue field of \(K\), where odd powers of Frobenius exchange
the two roots. These root identities hold modulo \(pO\):
the maps on \(O_E\) factor there through the field
\(O_E/\lambda\), by the uniformizer calculation, with nonzero
elements mapping to units. Since \(O_E\) is generated over
\(\Z\) by \(\omega,\zeta_{N_1}\) by
Proposition~\ref{prop:local-cyclo-different}, this proves
equality on \(O_E\). It extends to \(W\) by inversion of
the indicated units.
\end{proof}

\begin{theorem}\label{thm:served-actual-hecke-class-matrix}
Let \(r_0=i_0=\lceil\log_p m\rceil\),
\(t={\rm diag}(p I_{r_0},I_{k-r_0})\), and \(K={\rm GL}_k(O_\lambda)\).
Put
\[
 \mathcal P=K\cap tKt^{-1},\qquad
 \mathcal P'=K\cap t^{-1}Kt .
\]
Let \(C=\Gamma\backslash H_{0,f}/K_c\) and
\(H_{\rm rest}\) be the tensor product of the evaluation
coordinate spaces of the tops outside this pair. Give it the
integral lattice in these coordinates.
There are class representatives giving integral coordinates for
\(\mathcal E_p\), and an edge set \(\mathcal B\) with two
bijections
\[
       e\longmapsto(c(e),\,b_e\mathcal P),\qquad
       e\longmapsto(c'(e),\,k_e\mathcal P')
\]
onto \(C\times K/\mathcal P\) and \(C\times K/\mathcal P'\)
respectively, with \(b_e,k_e\in{\rm GL}_k(W)\).
In these coordinates the class matrix is
\[
 \iota(Y_\lambda\phi)(c)=
  \sum_{e:\,c(e)=c}
    \left(R_\lambda(b_e)\,{\rm diag}_{\beta}(p^{d(\beta)})\,
          R_\lambda(k_e^{-1})\otimes u_e\right)
                     \iota\phi(c'(e)),
\]
where \(u_e\) is an invertible integral matrix on \(H_{\rm rest}\).
Here \(\beta=(\beta_{b,s})\) ranges over nonnegative arrays with
each \(\sum_{1\le b\le k}\beta_{b,s}=a\), and
\[
   d(\beta)=\sum_{s\ {\rm even}}\sum_{b>r_0}\beta_{b,s}
                +\sum_{s\ {\rm odd}}\sum_{b\le r_0}\beta_{b,s}.
\]
\end{theorem}
\begin{proof}
Choose representatives \(g_c\) integral at all the places above
\(p\) by Lemma~\ref{lem:definite-one-prime-unitary-approximation}.
Since \(K_c\) is maximal integral above this \(p\), multiply on
the right by the inverse of the \(p\)-component supported above
\(p\). We can thus assume \((g_c)_p=1\) in every class. The
transformation law and Proposition~\ref{prop:column-dim} now
identify \(\mathcal E_p\) with all \(O_\iota\)-valued class
coordinate tuples: both right level and rational left changes
preserve the gauge by Theorem~\ref{thm:local-gauge-class-matrix}.

Embed \(t\) in the unitary group at the split pair.
Let \(K'\) agree with \(K_c\) off the pair and with \(\mathcal P\)
there, and set \(\mathcal B=\Gamma\backslash H_{0,f}/K'\).
For a fixed \(c\) its fibre consists of \(g_c b_e K'\) with
\(b_e\mathcal P\) running once through \(K/\mathcal P\),
since \(\Gamma\cap g_cK_c g_c^{-1}=1\) by
Proposition~\ref{prop:column-dim}. We can choose \(b_e\)
in \({\rm GL}_k(W)\): the condition defining \(\mathcal P\)
is divisibility of the upper-right block by \(p\), so the
coset is determined by reduction. Lift an invertible residue
matrix from \(O_E/\lambda\); its determinant lies outside
\(\lambda\). Right multiplication by \(t\) gives a bijection
of the edge set onto the quotient for \(t^{-1}K't\), with
inverse right multiplication by \(t^{-1}\). The new level
is \(\mathcal P'\) at the pair. The same fibre calculation on
this side gives the second bijection.

In particular \(g_c b_e t=\gamma_e g_{c'} k'_e\) for
\(\gamma_e\in\Gamma,\ k'_e\in K_c\). At \(\lambda\) write
\(k_e=(k'_e)_\lambda\). Using the identity \(p\)-components
of the representatives gives \((\gamma_e)_\lambda=b_e t k_e^{-1}\),
and therefore \(k_e\in{\rm GL}_k(W)\): all its entries and
those of its inverse are rational over \(E\), integral at
\(\lambda\). At other pairs over \(p\), \(\gamma_e\) and its
inverse are integral. Thus on all the other tops its weight
action is an integral automorphism \(u_e\).

The terms of \(T_{0,r_0,\lambda}\) in a given class are
\(\phi(g_c b_e t)=R_0(\gamma_{e,\infty})\phi(g_{c'})\);
\(b_e tK\) runs once over the desired cosets by the definition
of \(\mathcal P\). On the even tops, where \(\iota\tau_s\)
induces \(\lambda\), conjugating the rational unitary matrix
\(\gamma_e\) gives \(\gamma_e^{-t}\). On the odd tops its
ordinary local matrix at \(\lambda\) is already evaluated by
the conjugated embedding. This proves the action
\(R_\lambda(b_e)R_\lambda(t)R_\lambda(k_e^{-1})\) on the
linked factors, interpreting \(R_\lambda\) on \(t\) over
the fraction field. In an even factor \(t\) contributes minus
the number of letters in the first block; in an odd factor
it contributes plus this number. Multiplication by
\(p^{A_\lambda}\), with \(A_\lambda=aD/2\) by
Lemma~\ref{lem:served-actual-frobenius-cycle}, replaces the
negative count at each even factor by the positive count in
the second block. This gives exactly \(p^{d(\beta)}\) on each
monomial and proves the formula.
\end{proof}


\section{Evaluation coordinates and depth in all integral frames}

\begin{lemma}\label{lem:one-row-lattice-transpose}
For a matrix \(g\) of size \(k\), let \(M_a(g)\) denote the
coefficient matrix of the substitution
\(x_j\mapsto\sum_i g_{ij}x_i\) on monomials of degree \(a\);
the input monomial indexes its column. In the integral
evaluation coordinates of Definition~\ref{def:weyl-coord},
\[
                     S(g)=M_a(g^t)^t .
\]
In particular, with the ordinary polynomial representation
\[
  P_\lambda(g)=\bigotimes_{s\ {\rm even}}M_a(\sigma_s(g))
             \ \otimes\
         \bigotimes_{s\ {\rm odd}}M_a(\sigma_s(g^{-t})),
\]
one has \(R_\lambda(g)=P_\lambda(g^{-1})^t\).
The linked factor of each edge in
Theorem~\ref{thm:served-actual-hecke-class-matrix} is exactly
\[
          \left(P_\lambda(k_e)\,
                {\rm diag}_{\beta}(p^{d(\beta)})
                P_\lambda(b_e^{-1})\right)^t
\]
in the gauge used for the local index.
\end{lemma}
\begin{proof}
Expand the defining identity \(b(gx)=S(g)b(x)\). In row
\(\beta\) it reads
\[
       \prod_j\left(\sum_i g_{ji}x_i\right)^{\beta_j}
                       =\sum_\alpha S(g)_{\beta,\alpha}x^\alpha .
\]
The coefficient of \(x^\alpha\) on the left is
\(M_a(g^t)_{\alpha,\beta}\). This proof uses integral polynomial
identities, including when the coefficient characteristic
divides a multinomial. Apply the identity at each even and odd
factor of Definition~\ref{def:served-column-cycle}. Transposing
the resulting matrix product reverses the factors; the middle
diagonal is fixed by transpose. This proves both remaining
claims.
\end{proof}

\begin{definition}\label{def:depth-in-every-frame}
Use the cycle in Lemma~\ref{lem:served-actual-frobenius-cycle}.
Choose \(\varrho\in K_p\) with \(\varrho^2=p\) and set
\(\mathfrak m=\{z\in O:|z|_p<1\}\),
\(\mathbf k=O/\mathfrak m\). Write \(\mathcal V\) for the
free \(O\)-module of coefficients in the linked polynomial
representation \(P_\lambda\), and \(\bar{\mathcal V}\) for
its reduction to \(\mathbf k\). The depth of \(\beta\) is
\(d(\beta)\) in Theorem~\ref{thm:served-actual-hecke-class-matrix}.
Set \(D_\varrho={\rm diag}_\beta(\varrho^{d(\beta)})\) and
\(B=\max_\beta d(\beta)\). For an integer \(t\ge0\) define
the subspace \(V_t\subset\bar{\mathcal V}\) of polynomials
whose transform by every \(\overline{P_\lambda(g)}\),
\(g\in{\rm GL}_k(W)\), has zero coefficients at every
monomial of depth less than \(t\). The same constructions
work over \(\iota(O_\iota)\) and its residue field by using
\(\varrho\) in the image of \(\iota\).
\end{definition}

\begin{proposition}\label{prop:two-half-depth-saturation}
Let \(C'\) be a finite set, let \(s:E'\to C'\) have nonempty
finite fibres, and take \(h_e\in{\rm GL}_k(W)\) so that
in each fibre the left cosets \(\mathcal P h_e\) cover
\({\rm GL}_k(W)\) viewed inside \(K\).
Consider the matrix
\[
       J:\mathcal V^{C'}\longrightarrow\mathcal V^{E'},
           \qquad (JF)_e=D_\varrho P_\lambda(h_e)F_{s(e)} .
\]
It factors as \(J=Q A\) over \(O\), where \(A\) is square
of the source rank and \(Q\) has an integral left inverse.
Its determinant satisfies
\[
                2v_p(\det A)\le |C'|\sum_{t=1}^B
                                      \dim_{\mathbf k} V_t .
\]
The assertion applies both to \(s(e)=c(e),h_e=b_e^{-1}\)
and to \(s(e)=c'(e),h_e=k_e^t\) for the actual edges of
Theorem~\ref{thm:served-actual-hecke-class-matrix}.
\end{proposition}
\begin{proof}
For \(u\in\mathcal P\cap{\rm GL}_k(W)\), its upper-right block
is divisible by \(p\). On an even letter the deep variables
are in the second block. Rescale each deep variable by
\(\varrho\). An entry mapping a deep variable into a shallow
one is then divided by \(\varrho\); this entry lies in
\(pO=\varrho^2O\). On an odd letter the deep variables are in
the first block. The lower-left block of \(u^{-t}\) is now
divisible by \(p\), giving the same integrality statement.
By monomial substitution and tensor product, this proves
\(D_\varrho P_\lambda(u) D_\varrho^{-1}\) is integral.
The same calculation for \(u^{-1}\) proves it is invertible
over \(O\). Consequently if \(JF\) is divisible over \(O\)
by a scalar \(z\), then
\(D_\varrho P_\lambda(g)F_c\) is divisible by \(z\) for
every \(c,g\): choose \(e\) above \(c\) with \(g h_e^{-1}\in\mathcal P\).
The element \(g h_e^{-1}\) also belongs to \({\rm GL}_k(W)\).

The matrix \(J\) has full column rank since each fibre contains
an invertible block over the fraction field. The pivot argument
of Lemma~\ref{lem:valuation-index}, selecting the entry of least
valuation in the remaining rectangle, gives \(J=U{\binom{D'}{0}}V\)
with integral invertible \(U,V\) and nonzero diagonal \(D'\).
Take \(Q\) to be the first columns of \(U\), \(A=D'V\); the
corresponding rows of \(U^{-1}\) are a left inverse. A column
of \(V^{-1}\) is primitive and has a primitive component in
at least one fibre of the source. Applying any \(P_\lambda(h_e)\)
there preserves primitivity. Multiplication by \(D_\varrho\)
then leaves an entry of valuation at most \(B/2\).
Thus every diagonal entry \(a_l\) of \(D'\) has
\(0\le 2v_p(a_l)\le B\).

For an integer \(0\le j<B\) take the columns of \(V^{-1}\)
for which \(2v_p(a_l)>j\). Their reductions are independent,
since \(V\) is invertible over \(O\). For each of these columns
\(F\), \(JF\) is divisible by \(a_l\); applying the first
paragraph shows that the coordinates of
\(P_\lambda(g)F_c\) at all depths at most \(j\) have positive
valuation. Hence their reductions lie in \(V_{j+1}^{C'}\).
There can be at most \(|C'|\dim V_{j+1}\) such columns.
For any real \(0\le x\le B\), \(x\le\sum_{j=0}^{B-1}
{\bf1}_{x>j}\). Apply this to \(2v_p(a_l)\) and sum.
For the two applications the first coset coverage follows by
inverting \(b_e\mathcal P\); transposing \(k_e\mathcal P'\)
gives \(\mathcal P k_e^t\), since
\((\mathcal P')^t=\mathcal P\) by the two block descriptions.
The bijections of the class-matrix theorem give the two
coverings. All operations on this matrix can also be taken
over \(\iota(O_\iota)\), a valuation ring, giving the last
variant of the definition.
\end{proof}


\section{Weighted polynomial interpolation with multiplicity}

\begin{definition}\label{def:weighted-point-power-equations}
Let \(K_0\) be a field, \(n\ge1\), and \(k_i\) positive integers
for \(1\le i\le n\). The weight of a monomial
\(X^\alpha\) in \(K_0[X_1,\ldots,X_n]\) is
\(\sum_i\alpha_i/k_i\). At a point \(z\in K_0^n\), a
polynomial vanishes to order at least \(d\in\Z_{\ge0}\)
if its expansion in the translated variables \(X_i-z_i\)
has zero coefficients in total degrees below \(d\).
\end{definition}

\begin{theorem}\label{thm:weighted-points-multiplicity}
Suppose a finite set \(Z\subseteq K_0^n\) has at least
\(\prod_i k_i\) points. For each \(i\) suppose \(g_i\)
is a polynomial all of whose monomials have weight less
than one, with \(g_i(z)=z_i^{k_i}\) on \(Z\). Then \(Z\)
has exactly \(\prod_i k_i\) points. For each \(d\ge1\) the
ideal of polynomials vanishing to order at least \(d\)
on \(Z\) equals the \(d\)-th power of
\((X_i^{k_i}-g_i:1\le i\le n)\). In particular each nonzero
such polynomial has a monomial of weight at least \(d\).
\end{theorem}
\begin{proof}
Set \(h_i=X_i^{k_i}-g_i\). The polynomials
\[
          X^\alpha h^\gamma,\qquad
       0\le\alpha_i<k_i,\quad \gamma\in\Z_{\ge0}^n,
       \qquad h^\gamma=\prod_i h_i^{\gamma_i},
\]
form a basis of the polynomial ring as a vector space.
Indeed each has a unique highest-weight monomial
\(X^{\alpha+(k_i\gamma_i)_i}\) with coefficient one. These
leading monomials are distinct and cover all monomials by
Euclidean division in each exponent. Comparing the terms
of greatest weight proves independence in any finite linear
combination. For spanning, subtract the corresponding basis
element from a monomial and induct on the weight multiplied
by \(\prod_i k_i\); this is a nonnegative integer and every
remaining monomial has strictly smaller weight.

Write \(I=(h_i)_i\). The basis elements with \(|\gamma|\ge d\)
span \(I^d\). One inclusion is by divisibility; for the reverse
one, expand an arbitrary polynomial multiplying any product
\(h^\eta\), \(|\eta|=d\), in the displayed basis. The product
just adds \(\eta\) to its index \(\gamma\). Hence the codimension
of \(I^d\) is \((\prod_i k_i)\#\{\gamma:|\gamma|<d\}\).
For a point \(z\), its ideal
\(\mathfrak m_z=(X_i-z_i)_i\) satisfies
\(f\in\mathfrak m_z^d\) exactly under the vanishing condition
of Definition~\ref{def:weighted-point-power-equations}, by expansion
in the translated monomials. Its codimension is
\(\#\{\gamma:|\gamma|<d\}\).
For \(z\ne z'\) a differing coordinate gives
\(a\in\mathfrak m_z\), \(1-a\in\mathfrak m_{z'}\) by affine
rescaling. Expanding \(1=(a+(1-a))^{2d-1}\) gives
\(1\in\mathfrak m_z^d+\mathfrak m_{z'}^d\).
For each fixed \(z\), multiply, over the other points,
polynomials equal to \(1\bmod\mathfrak m_z^d\) and to
\(0\bmod\mathfrak m_{z'}^d\) at that other point. These
products interpolate arbitrary classes in all the quotients
by multiplying individual representative polynomials and
adding. Thus the codimension of \(\bigcap_z\mathfrak m_z^d\)
is \(|Z|\#\{\gamma:|\gamma|<d\}\).
Since \(I^d\) is contained in this intersection, the two
codimensions and the assumed point count force equality of
both the cardinalities and the ideals.

For nonzero \(f\in I^d\) use its finite basis expansion.
Every index there has \(|\gamma|\ge d\). At the largest
weight of these terms a leading monomial survives by
independence of the distinct leading monomials. Its weight is
\(\sum_i\alpha_i/k_i+|\gamma|\ge d\), as required.
\end{proof}


\section{Interpolation in the linked symmetric powers}

\begin{theorem}\label{thm:served-deep-mean-cost}
Use \(D,P,a,k\) and \(\mathbf k\) of
Definition~\ref{def:served-column-cycle} and
Definition~\ref{def:depth-in-every-frame}. Write \(r=i_0\)
and number the \(k\) coordinates \(0,\ldots,k-1\) in this
section. Assume \(1\le2r\le k,\ 2r<D\), as holds eventually
by \(D\ge m,\ r\le1+\log_2m,\ k\gg m\).
A coordinate \((b,s)\) is deep if \(b\ge r\) for even \(s\)
or \(b<r\) for odd \(s\). For a monomial \(\beta\) define
\[
 C(\beta)=\frac1P\sum_{\substack{(b,s)\ {\rm deep}\\ b<2r}}\beta_{b,s}
       +\frac1{P^{2r}}\sum_{\substack{(b,s)\ {\rm deep}\\b\ge 2r}}
                                                   \beta_{b,s}.
\]
Every nonzero \(F\in V_d\), \(d\ge1\), has a monomial of
nonzero coefficient and \(C(\beta)\ge d\). In particular,
for \(M_0=\binom{a+k-1}{k-1}^{D}\) one has
\[
 \sum_{d\ge1}\dim_{\mathbf k} V_d
    \le M_0\,\frac{Da}{k}
               \left(\frac rP+\frac{k-2r}{2P^{2r}}\right).
\]
\end{theorem}
\begin{proof}
Let \(k_q=W/pW\) of size \(q=P^D\). Each \(\sigma_s\) reduces
to an embedding \(k_q\to\mathbf k\), denoted \(\bar\sigma_s\);
by Lemma~\ref{lem:served-actual-frobenius-cycle} its successor
is its \(P\)-th power. Write \(y^{(s)}=\bar\sigma_s(y)\).
Replace all variables at shallow coordinates by independent
indeterminates \(x_{b,s}\) and let \(K_0=\mathbf k(x_{b,s}
:(b,s)\ {\rm shallow})\). Keep the deep coordinates as
polynomial variables \(Z_{b,s}\) over \(K_0\). This substitution
is injective, by collecting coefficients of the deep monomials
in the independent shallow variables.

Take a \(k\)-by-\(k\) matrix \(A_y\) over \(k_q\) supported
on the following positions with independent entries \(y_{bc}\):
\[
      (b,r+b)\ (0\le b<r),\qquad
      (b,c)\ (0\le b<r,\ 2r\le c<k).
\]
It is square-zero, and may be lifted with the same support
to \(W\), where its square is still zero. In the action of
\(1+A_y\), the shallow variables are fixed. In each even
letter, a deep variable in column \(c\) gains
\(\sum_b x_{b,s}y_{bc}^{(s)}\); in each odd letter, a deep
variable indexed by \(b<r\) gains
\(-\sum_c x_{c,s} y_{bc}^{(s)}\). Entries outside the chosen
support are read as zero. The second sign comes from
\((1+A_y)^{-t}=1-A_y^t\). Call the vector of these gains
\(z_y\). Since the coefficients of the transform of \(F\) at
every depth below \(d\) are zero, the polynomial obtained
from \(F\) after the substitution vanishes at \(z_y\) to
order at least \(d\): its translate \(F(Z+z_y)\) has no
terms of ordinary deep degree less than \(d\).

Assign costs \(1/P\) to the deep indices \(b<2r\) and
\(P^{-2r}\) to all other deep indices. We show for each \(i\)
that all \(y_{bc}^{(i)}\), as functions on the points, are
represented by polynomials in \(Z\) of weight at most \(1/P\).
For \(c\ge2r\) choose the \(r\) integers \(j\) in
\(0\le j<2r\) for which \(i-j\) is even. Their letters
\(i-j\) are distinct modulo \(D\) since \(D>2r\).
The equations at those even letters give
\[
        (z_y)_{c,i-j}^{\,P^j}
                 =\sum_{b<r}x_{b,i-j}^{P^j}y_{bc}^{(i)} .
\]
Here iterating Frobenius carries the residue value to the
letter \(i\), and taking powers distributes over sums in
characteristic \(p\). The \(r\)-by-\(r\) coefficient matrix
is invertible over \(K_0\): the terms of its determinant
are distinct monomials in independent variables indexed by
the distinct row letters and column indices. Its inverse
expresses the unknowns as linear combinations of
\(Z_{c,i-j}^{P^j}\), each of weight at most \(1/P\).
For an even \(i\), the coordinate \(r+b\) gives
\(y_{b,r+b}^{(i)}=(z_y)_{r+b,i}/x_{b,i}\).
For an odd \(i\), use the coordinate \(b\). Its value is
\(-x_{r+b,i}y_{b,r+b}^{(i)}-\sum_{c\ge2r}x_{c,i}y_{bc}^{(i)}\);
the latter entries have already been recovered. In both
cases division is by a nonzero indeterminate of \(K_0\)
and the claimed weight bound follows. These formulas also
apply when there are no columns \(c\ge2r\), by the empty sum.

Recovery at one letter proves the \(z_y\) distinct. There
are \(q^{r(k-2r+1)}\) of them. For each deep coordinate \(v\)
write its cost as \(1/k_v\). At every letter there are
exactly \(r\) deep indices below \(2r\), and in the \(D/2\)
even letters there are \(k-2r\) more. Hence
\(\prod_v k_v=P^{Dr+Dr(k-2r)}=q^{r(k-2r+1)}\).
Furthermore the function \((z_y)_v^{k_v}\) has a polynomial
of weight less than one on these points. Indeed \(k_v\)
is \(P^j\) with \(j=1\) or \(2r\); raise the defining sum
for \((z_y)_v\) to this power. It becomes a \(K_0\)-linear
combination of entries \(y_{bc}^{(s+j)}\). The polynomials
just obtained have weight at most \(1/P<1\).
Theorem~\ref{thm:weighted-points-multiplicity} therefore gives
a monomial of the substituted \(F\) with weight at least \(d\).
Its nonzero coefficient comes from at least one original
monomial with the same deep exponents, hence the same cost.

Restriction to the coefficients of monomials of cost at
least \(d\) is thus injective on \(V_d\). Sum these dimension
bounds over positive integers \(d\). Each monomial contributes
at most \(\lfloor C(\beta)\rfloor\le C(\beta)\).
There are \(M_0\) monomials. In each letter their average
exponent in any specified coordinate is \(a/k\), by permuting
coordinates and summing all \(k\) exponents. The counts of deep
indices just computed now give the claimed sum.
\end{proof}

\begin{corollary}\label{cor:served-half-cost-actual}
In either half of Proposition~\ref{prop:two-half-depth-saturation},
with the actual classes, set \(z_\lambda= (Da/k)
(r/P+(k-2r)/(2P^{2r}))\). Then
\[
               2v_p(\det A)\le |C|M_0\,z_\lambda,\qquad
               z_\lambda\le\frac{128 A_\lambda\ell}{mp}
\]
for all sufficiently large arithmetic scales at every served
prime. Tensoring the half matrix with the identity of
\(H_{\rm rest}\) multiplies its determinant cost bound by
\(\dim H_{\rm rest}\).
\end{corollary}
\begin{proof}
Use Theorem~\ref{thm:served-deep-mean-cost} in
Proposition~\ref{prop:two-half-depth-saturation}. For the
numerical bound note \(P=p^o\ge p,\ p^r\ge m\), and \(r\ge1\),
so \(P^{2r}\ge mp\). Also
\(r\le1+\log_2m,\ mr/k\le1\) eventually from the scales in
Lemma~\ref{lem:diagram-size} and Definition~\ref{def:heights}.
It follows that \(z_\lambda\le
Da(1/(mp)+1/(2mp))=3A_\lambda/(mp)\), since \(Da=2A_\lambda\).
This even implies the asserted bound as \(\ell\to\infty\).
The tensor statement follows by reordering the basis into
one copy of every matrix for each basis element of the fixed
factor. Determinants multiply over block diagonal matrices.
\end{proof}


\section{The volume of cyclic interlacing inequalities}

\begin{definition}\label{def:cyclic-interlacing-volume}
Fix integers \(k,D\ge1\), a real number \(P>1\), and a sign
\(\epsilon_s\in\{1,-1\}\) for each \(s\in\Z/D\Z\). Put \(q=P^D\).
Number the rows \(0,\ldots,k-1\). For permutations
\(\sigma_s\) of the rows define
\[
 (L_\sigma\nu)_{b,s}
     =\epsilon_s(P\nu_{b,s+1}-\nu_{\sigma_s(b),s}).
\]
Write \(L_1\) for the identity permutations. Define \(\Omega\)
to be the image under \(L_1\) of the following inequalities for all
\(b<k-1,s\):
\[
  \begin{cases}
     P\nu_{b+1,s+1}\le\nu_{b,s}, &\epsilon_s=1,\\
     \nu_{b+1,s}\le P\nu_{b,s+1}, &\epsilon_s=-1 .
  \end{cases}
\]
For \(A\in\R^D\), let \(\Sigma(A)\) consist of the arrays
\(\beta\ge0\) with \(A_s-1\le\sum_b\beta_{b,s}\le A_s\)
for all \(s\). Volume is Euclidean volume in the \(kD\)
coordinates.
\end{definition}

\begin{lemma}\label{lem:cyclic-indicator-determinant}
For strictly decreasing real \(k\)-tuples \(x,y\), assume
\(x_b\ne y_c\) for all \(b,c\). The matrix
\(({\bf1}_{x_b>y_c})_{b,c}\) has determinant one when
\(x_0>y_0>x_1>y_1>\cdots>x_{k-1}>y_{k-1}\). It has determinant
zero otherwise.
\end{lemma}
\begin{proof}
Each row consists of some zeros followed by ones. The numbers of
ones decrease weakly down the rows. A zero row or two equal
counts give a zero determinant. In the remaining case the counts
are the \(k\) distinct positive integers \(k,k-1,\ldots,1\).
This gives exactly the displayed inequalities, and the matrix
is upper triangular with ones on the diagonal.
\end{proof}

\begin{theorem}\label{thm:cyclic-flag-share}
With the notation of Definition~\ref{def:cyclic-interlacing-volume},
\[
    {\rm vol}\bigl(\Omega\cap\Sigma(A)\bigr)
     =\left(\prod_{j=1}^k\frac{q^j-1}{q-1}\right)^{-1}
        {\rm vol}\bigl(\Sigma(A)\bigr).
\]
When \(q\) is the cardinality of a finite field the denominator
counts its complete flags in dimension \(k\).
\end{theorem}
\begin{proof}
First compute the Jacobians. Write \(S_s\) for the permutation
matrix taking \((u_b)_b\) to \((u_{\sigma_s(b)})_b\). Up to
signs and rearrangement of coordinates, the new variables
are the residuals \(P\nu_{s+1}-S_s\nu_s\). Successively replacing
\(\nu_1,\ldots,\nu_{D-1}\) by the first \(D-1\) residuals,
retaining \(\nu_0\), has absolute determinant \(P^{k(D-1)}\)
by triangularity. Holding these residuals fixed, the last residual
\(P\nu_0-S_{D-1}\nu_{D-1}\) has linear part in \(\nu_0\) equal
to \(P^{1-D}(qI-S_{D-1}\cdots S_0)\). This also describes the
one-letter case. Therefore
\[
       |\det L_\sigma|=|\det(qI-S_{D-1}\cdots S_0)|.
\]
It is positive since a permutation preserves the maximum norm
and \(q>1\). In particular \(|\det L_1|=(q-1)^k\).

In integrating indicators below one can discard the finitely
many equality hyperplanes between two entries being compared,
or between two different entries within one letter. Each is
a proper hyperplane, including when \(D=1\) since \(P>1\).
For clarity, the Euclidean volume rules needed here follow
by slicing bounded polyhedra in coordinates. An affine
hyperplane has zero volume by slicing in any coordinate with
nonzero coefficient. The same conclusion follows by covering
its bounded portion with grid boxes of side \(\delta\): over each
choice of the other coordinates it meets a bounded number of
layers of boxes, so their volume tends to zero. Thus inner and
outer grid sums of a bounded polyhedron agree in the limit.
They also give the iterated integrals of its interval slices
by squeezing. Permuting coordinates preserves volume, translation
or shear in one coordinate translates the corresponding slices,
and scaling a coordinate multiplies slice length by its absolute
scale. Elimination writes each invertible linear change as a product
of these elementary changes, which proves the absolute determinant
rule. All sets integrated below are bounded or finite unions of
bounded polyhedra, by invertibility and boundedness of \(\Sigma(A)\).

For a fixed \(\nu\), impose the column-sum intervals on \(L_1\nu\).
Their indicator \(I(\nu)\) is unchanged by independently permuting
the rows at each letter, and the column sums of \(L_\sigma\nu\)
are the same. Form the matrices
\[
  (H_s)_{bc}={\bf1}_{\,\epsilon_s(P\nu_{c,s+1}-\nu_{b,s})>0}.
\]
Expanding and multiplying their determinants gives
\[
 I(\nu)\prod_s\det H_s
   =\sum_\sigma \left(\prod_s{\rm sgn}\,\sigma_s\right)
                    {\bf1}_{L_\sigma\nu\in\Sigma(A)}
\]
off the equality walls. Independently sort the entries in each
letter of \(\nu\) in decreasing order. This permutes the rows of
one matrix and the columns of its predecessor by the same
permutation, preserving the product of determinants. For the
sorted array, apply Lemma~\ref{lem:cyclic-indicator-determinant}
to \(H_s^t\) at a positive sign and to \(H_s\) at a negative
sign. The product is one exactly when the sorted array has
positive \(L_1\)-coordinates and satisfies the defining
interlacings. Conversely these strict interlacings and the
positive coordinates force decreasing order in each letter:
in either orientation, two successive entries of that letter
enclose an entry of its neighbor. Thus \(I(\nu)\prod_s\det H_s\)
is the sum over the \((k!)^D\) choices of row permutations
\(\tau=(\tau_s)_s\) of the indicators
\({\bf1}_{L_1(\tau\nu)\in\Omega\cap\Sigma(A)}\).

Integrate these two evaluations. Every product permutation
\(S=S_{D-1}\cdots S_0\) occurs \((k!)^{D-1}\) times and its
sign is the product of the individual signs. The result is
\[
 {\rm vol}(\Omega\cap\Sigma(A))=K_k(q)\,{\rm vol}(\Sigma(A)),
 \qquad
 K_k(q)=\frac{(q-1)^k}{k!}
      \sum_{S\in{\mathfrak S}_k}
                    \frac{{\rm sgn}(S)}{|\det(qI-S)|}.
\]
Only \(k,q\) occur on the right in the expression for \(K_k(q)\).

Evaluate \(K_k(q)\) with one positive letter, \(P=q\) and \(A=1\).
Here \(\beta=(q-1)\nu\), and the inequalities are
\(\beta_b\ge q\beta_{b+1}\). Put \(\beta_k=0\) and
\(d_j=(q^j-1)/(q-1)\). The change
\[
    z_b=d_{b+1}(\beta_b-q\beta_{b+1}),\qquad 0\le b<k,
\]
maps this part of the slab bijectively onto the whole slab.
Back substitution from \(b=k-1\) shows that \(z\ge0\) forces
all entries of \(\beta\) nonnegative, and \(d_{b+1}-qd_b=1\)
gives \(\sum_b z_b=\sum_b\beta_b\). The determinant is
\(\prod_{j=1}^k d_j\). This slab has positive volume, since it
contains \([0,1/k]^k\), so its volume cancels to give the
asserted factor. Finally a line in an \(h\)-dimensional space
over the field of order \(q\) can be chosen in \((q^h-1)/(q-1)\)
ways by counting its nonzero vectors. After choosing a line,
the rest of a flag is a flag in the quotient. Induction from
dimension zero gives the stated product.
\end{proof}


\section{Thickened inequalities at the boundary between two blocks}

\begin{definition}\label{def:cyclic-root-excess}
Use the notation of Definition~\ref{def:cyclic-interlacing-volume}
with \(D\) even and \(\epsilon_s=1\) at even \(s\).
For \(i,s\in\Z/D\Z\), let \([s-i]_D\) be the representative
in \(\{0,\ldots,D-1\}\) and put \(M_s=P^{[s-i]_D}\).
For \(0\le b<k-1\), denote the two variable indices in rows
\(b,b+1\) at letter \(s\) by \(y_s,x_s\), where \(y_s\) is
the index in row \(b+1\) at even \(s\), and in row \(b\) at
odd \(s\). Define the linear form
\[
  E_{i,b}(\beta)=
      (1-q^{-1})\beta_{y_i}
           +q^{-1}\sum_s M_s(\beta_{y_s}-\beta_{x_s}).
\]
For \(1\le r<k\) and real \(H\) let \(\Omega_H^{(r)}\)
be the set of \(\beta\) with
\(E_{i,b}(\beta)\le H{\bf1}_{b+1=r}\) for every \(i,b\).
\end{definition}

\begin{lemma}\label{lem:cyclic-excess-shift}
One has \(\Omega_0^{(r)}=\Omega\). For \(H\ge0\) there is a
single nonnegative translation taking
\[
 \Omega_H^{(r)}\cap\Sigma(A)\ \text{into}\
      \Omega\cap\Sigma(A'),\qquad
 A'_s=A_s+H{\bf1}_{s\ {\rm even}}\frac{q(P^{2r}-1)}{q-1}.
\]
\end{lemma}
\begin{proof}
For \(\beta=L_1\nu\) put \(\Delta_s=\nu_{b,s}-\nu_{b+1,s}\).
In both parities \(\beta_{y_s}-\beta_{x_s}
=\Delta_s-P\Delta_{s+1}\).
In \(\sum_s M_s(\Delta_s-P\Delta_{s+1})\) every term cancels
except at \(i\), since \(P M_{s-1}=M_s\) there except when
\(s=i\), where the two values are \(q,1\). Hence
\[
       E_{i,b}(L_1\nu)=(1-q^{-1})(\beta_{y_i}-\Delta_i).
\]
The nonpositivity conditions are precisely the interlacings
in Definition~\ref{def:cyclic-interlacing-volume}.
Invertibility of \(L_1\) was proved in
Theorem~\ref{thm:cyclic-flag-share}.

Set \(v_b=(1-q^{-1})^{-1}P^{2(r-1-b)}\) for \(b<r\) and
\(v_b=0\) otherwise. Take \(\nu_{b,s}=v_b\) at even letters
and \(\nu_{b,s}=Pv_b\) at odd ones. Then \(\xi=L_1\nu\) is
\((P^2-1)v_b\) at even letters and zero at odd ones.
At even \(i\) the excess on \(\xi\) is
\((1-q^{-1})(P^2v_{b+1}-v_b)\). This equals
\(-{\bf1}_{b+1=r}\). At odd \(i\) the excess is
\(-(1-q^{-1})P(v_b-v_{b+1})\), which is nonpositive and
equals \(-P\) at the cut. Thus addition of \(H\xi\) takes the
indicated cone into \(\Omega\). Every entry of \(\xi\) is
nonnegative, and summing its entries at each even letter
gives \((1-q^{-1})^{-1}(P^{2r}-1)\) by the finite geometric
sum. Addition translates both endpoints of each column-sum
interval by that same column total multiplied by \(H\).
\end{proof}

\begin{proposition}\label{prop:thick-cyclic-flag-volume}
Take \(a\ge1\) integral, \(H\ge0\), and \(A_s=a+k+1\). Write
\[
    F_k(q)=\prod_{j=1}^k\frac{q^j-1}{q-1},\qquad
    C_r=\frac{P^{2r}-1}{1-q^{-1}},\qquad
    u=\frac{D(k-1)}{a}\left(k+1+\frac{HC_r}{2}\right).
\]
Then
\[
 F_k(q)\,{\rm vol}(\Omega_H^{(r)}\cap\Sigma(A))
        \le \exp(u)\binom{a+k-1}{k-1}^{D}.
\]
In Definition~\ref{def:served-column-cycle} take
\(r=i_0,\ h=x^{40},\ H=h+D(k+1)\). Uniformly at all served
primes, for sufficiently large \(x\) the factor \(\exp(u)\)
is at most \(1+x^{-14}\).
\end{proposition}
\begin{proof}
Translate by Lemma~\ref{lem:cyclic-excess-shift} and apply
Theorem~\ref{thm:cyclic-flag-share} at the new endpoints.
In \(k\) nonnegative coordinates the simplex of sum at most
\(z\ge0\) has volume \(z^k/k!\): successively integrate the
last coordinate from 0 to \(z\), starting with length in
dimension one. For \(z\ge1\), subtracting the volume at
\(z-1\) gives at most \(z^{k-1}/(k-1)!\) by integrating the
derivative over that unit interval. The slabs are independent
in the letters, so their volumes multiply by iterated integration.
Divide each resulting bound by
\(\binom{a+k-1}{k-1}=\prod_{j=1}^{k-1}(a+j)/j
\ge a^{k-1}/(k-1)!\). The binomial expression counts the
placements of \(k-1\) separators among \(a+k-1\) slots,
equivalently monomials of degree \(a\) in \(k\) variables.
The inequality \(1+t\le\exp(t)\) for \(t\ge0\) follows from
the exponential series. Use it at
\(t=(k+1+H C_r{\bf1}_{s\ {\rm even}})/a\) and sum the
exponents, noting that exactly \(D/2\) letters are even.

For the arithmetic assertion, \(D=L/o\le f\le x^{12}\):
the orbit length \(L\) is at most the order \(f\) of the
group modulo \(N_1\), as in
Lemma~\ref{lem:served-actual-frobenius-cycle}. Also \(k=O(x^5)\)
and \(a\ge x\exp(\ell)\), with \(\ell=x^\epsilon\), by
Lemma~\ref{lem:diagram-size} and its proof. At \(p\ge m\),
\(o=r=1\) by Proposition~\ref{prop:marked-cost} and its proof,
and \(\log p\le (10^6\log x)^2\). At \(p<m\), both \(o\le r\)
and \(r\le1+\log_2 m\) by Definition~\ref{def:collisions},
while \(m\le2\sqrt x\). In either case
\(\log C_r=O((1+\log x)^3)\), using \(q\ge2\).
The other factors in the numerator of \(u\) grow at most as a
fixed power of \(x\) times \(C_r\). Hence \(u\le x^{-14}/2\)
eventually: every fixed power of \(\log x\) is \(o(x^\epsilon)\),
as seen from the exponential series with argument
\(\epsilon\log x\). For \(0\le u\le1/2\) the series also gives
\(\exp(u)-1\le\sum_{n\ge1}u^n\le2u\). This proves the
uniform estimate.
\end{proof}


\section{A polynomial detecting an extremal exponent}

\begin{definition}\label{def:root-coefficient-order}
Use \(W,O,p,P,q,D,\sigma_s\) of
Definition~\ref{def:served-column-cycle}, and the row pair and
letter \(i\) of Definition~\ref{def:cyclic-root-excess}.
Use coefficients \(B=O\), or \(B=\mathbf k\) with the reduced
maps \(\sigma_s\) of Definition~\ref{def:depth-in-every-frame}.
For \(z\ne0\) put \(v(z)=\lfloor2v_p(z)\rfloor\) over \(O\),
where \(v_p(p)=1\), and \(v(z)=0\) over \(\mathbf k\).
Set \(v(0)=+\infty\). For \(g\in{\rm GL}_k(W)\) let \(gF\)
denote the polynomial action of
Lemma~\ref{lem:one-row-lattice-transpose}, also on polynomials
of arbitrary multidegrees by the same linear substitutions.
For real coordinate weights \(c\), the cost of a nonzero term
\(zX^\alpha\) is \(v(z)+c\cdot\alpha\).
\end{definition}

\begin{theorem}\label{thm:root-extremum-by-one-polynomial}
Suppose \(P\ge5\) and
\(2/5\le c(y_s)-c(x_s)\le8/5\) for every \(s\).
Let \(G\subseteq{\rm GL}_k(W)\) be a subgroup containing the
invertible matrices on the chosen two rows acting as the
identity on the complementary rows. Let \(F\) be a polynomial
over \(B\) with nonzero coefficient at \(\alpha\) of order \(l\).
Suppose every term of \(gF\), for every \(g\in G\), has cost
at least \(l+c\cdot\alpha\). Then \(E_{i,b}(\alpha)\le0\).
\end{theorem}
\begin{proof}
Write \(\delta_s=c(y_s)-c(x_s)\), \(\omega=\delta_i/q\), and
\(M_s=P^{[s-i]_D}\). Use lower-case \(x_s,y_s\) for polynomial
variables as well as their indices. Make new variables \(X_s,Y_s\)
for \(s\ne i\) and only \(X_i\) at \(i\). Variables outside
the pair are retained. Write \(\pi\alpha\) for the exponents
after merging \(x_i,y_i\) into \(X_i\) and relabelling the rest.
Define integers
\[
   m(\alpha)=q\sum_s M_s\alpha_{x_s},\qquad
   e(\alpha)=q\alpha_{y_i}+\sum_{s\ne i}M_s\alpha_{y_s}.
\]
Let \(\eta_s=-1\) for \(s\) of the same parity as \(i\), and
\(\eta_s=1\) otherwise. With one additional variable \(T\)
use the substitution \(\Phi\) given by
\[
 \begin{array}{c|cc}
     &x_s&y_s\\ \hline
 s=i &X_iT^q&X_i\\
 s\ne i & X_sT^{qM_s}-\eta_sY_s T^{M_s}&Y_s-\eta_sX_s .
 \end{array}
\]
It fixes every variable outside the pair. Put
\(I_l=\{z:v(z)>l\}\). This is an ideal: the order of a sum
is at least the minimum and the order of a product is at
least the sum, by the valuation inequality and
\(\lfloor u+v\rfloor\ge\lfloor u\rfloor+\lfloor v\rfloor\).
Let \(A(T)\) be the coefficient at \(\pi\alpha\) in \(\Phi F\)
over \(B/I_l\), treating \(T\) as scalar. We prove that
its degree is \(m(\alpha)\) and that it vanishes to order
\(e(\alpha)\) at all \(q\) residues.

At infinity give \(T\) weight \(-\omega\), \(X_i\) weight
\(c(y_i)\), and, for \(s\ne i\), give \(X_s,Y_s\) weights
\(c(x_s)+qM_s\omega,\ c(y_s)\), respectively. The first
summand of each substitution has the original weight, including
both entries at \(i\). At \(s\ne i\) each second summand
has strictly greater weight. The two increments are
\(\delta_s-M_s\omega>0\) and \(M_s\delta_i-\delta_s>0\),
since \(5\le M_s\le q/5\). Fixed variables keep their weights.
Every product obtained by expansion therefore has at least its
original cost and gains strictly if it uses a second summand.
The new variables with exponent \(\pi\alpha\) alone have weight
\(c\cdot\alpha+\omega m(\alpha)\). Consequently coefficients
of \(T^j\) there with \(j>m(\alpha)\) have order strictly
greater than \(l\). At \(j=m(\alpha)\) all terms using a
second summand also have order greater than \(l\). Using just
first summands gives exactly \(F_\alpha\): from \(\pi\alpha\)
one knows every exponent except the split of
\(\alpha_{x_i}+\alpha_{y_i}\), and \(m(\alpha)\) determines
that split. Thus \(A\ne0\) and its degree is \(m(\alpha)\).

Fix \(t\in W\), let \(u=\sigma_i(t)\) and \(u_s=\sigma_s(t)\).
Take the matrix on rows \(b,b+1\) which is their swap times the
upper transvection of parameter \(-t\) for even \(i\), \(t\) for
odd \(i\). In the order \(x_s,y_s\) it acts on variables by
\[
               g x_s=y_s,\qquad g y_s=x_s+\eta_s u_s y_s.
\]
Indeed the even action has columns of the swap times an upper
transvection. The odd action takes the inverse transpose and
interchanges the row names, reversing the sign of the parameter.
Use \(F=g(g^{-1}F)\) before substituting, and replace \(T\)
by \(T+u\). For \(h=M_s\) or \(qM_s\) write
\[
         D_{s,h}(T)=(T+u)^h-u_s-T^h .
\]
All its coefficients have order at least two. For positive
intermediate exponents the binomial coefficients are divisible
by \(p\), since successive \(p\)-th powers distribute over
sums modulo \(p\). The constant coefficient has this same
divisibility by iterating the congruence of
Lemma~\ref{lem:served-actual-frobenius-cycle} through
\([s-i]_D\) steps or through \(D+[s-i]_D\) steps. In the
residue field these error polynomials vanish. In particular
this holds at \(s=i,h=q\).

At this finite point give \(T\) weight \(\omega\), \(X_i\)
weight \(c(x_i)\), and for \(s\ne i\) give \(X_s,Y_s\)
weights \(c(x_s),\,c(y_s)-M_s\omega\), respectively. At \(i\)
the substituted images of \(gx_i,gy_i\) are \(X_i\) and
\(X_i(T^q+D_{i,q})\). Their terms have at least the weights of
\(x_i,y_i\), since \(q\omega=\delta_i<2\). At \(s\ne i\) the
corresponding images are
\[
   Y_s-\eta_sX_s,\qquad
     X_s(T^{qM_s}+D_{s,qM_s})
               -\eta_sY_s(T^{M_s}+D_{s,M_s}).
\]
The first has weight at least \(c(x_s)\) because
\(M_s\omega<\delta_s\). For the second, its two pure powers
have weights \(c(x_s)+M_s\delta_i\ge c(y_s)\) and \(c(y_s)\).
Its errors have cost at least \(c(x_s)+2\) and
\(c(y_s)-M_s\omega+2\), both at least \(c(y_s)\).
Thus every substituted term of \(g^{-1}F\) respects its cost
bound. The finite-point weight of the variables at \(\pi\alpha\)
alone is \(c\cdot\alpha-\omega e(\alpha)\).
For \(j<e(\alpha)\), the coefficient at that monomial and \(T^j\)
therefore has order at least \(l+\omega(e(\alpha)-j)>l\).
It vanishes in \(B/I_l\), proving
\(T^{e(\alpha)}\mid A(T+u)\).

Take \(q\) lifts \(t\) of the distinct residues of \(W\).
Their images have pairwise unit differences by
Lemma~\ref{lem:served-actual-frobenius-cycle}.
The factors \(T-u\) for different lifts thus generate the
unit ideal in pairs, by subtracting them. Their positive
\(e(\alpha)\)-th powers are also pairwise comaximal: expand
\(1=(v+w)^{2e(\alpha)-1}\) when \(v+w=1\) comes from the two
ideals. For comaximal ideals \(I,J\), writing \(1=v+w\) in
the same way gives \(I\cap J=IJ\). Products of ideals
comaximal with a fixed ideal remain comaximal with it, by
multiplying representatives of one in the quotient. Hence
the product of the indicated monic powers divides \(A\).
Multiplying a nonzero polynomial by a monic polynomial adds
their degrees even over \(B/I_l\). It follows that
\(q e(\alpha)\le m(\alpha)\), also true for \(e(\alpha)=0\).
As \(E_{i,b}(\alpha)=(q e(\alpha)-m(\alpha))/q^2\), this
is the asserted inequality.
\end{proof}


\section{Divisibility in every frame preserving the two blocks}

\begin{definition}\label{def:served-excess-forms}
Put \(r=i_0\) in the cycle of
Definition~\ref{def:served-column-cycle}, with \(1\le r<k\).
Let \(L_r(W)\) be the block diagonal subgroup with block sizes
\(r,k-r\). For a real \(h\) let \(U(h)\subset\bar{\mathcal V}\)
be the subspace of residue polynomials \(F\) of degree \(a\) in
each letter for which every monomial \(\beta\) in every
\(gF,\ g\in L_r(W)\), satisfies
\(E_{i,r-1}(\beta)\le h\) for all \(i\), with
\(E_{i,r-1}\) as in Definition~\ref{def:cyclic-root-excess}.
Here and below a monomial in a polynomial means one with nonzero
coefficient. The subspace condition follows since vanishing of
each specified coefficient after a fixed linear action is linear.
\end{definition}

\begin{theorem}\label{thm:served-saturation-excess}
In either factorization \(J=QA\) of
Proposition~\ref{prop:two-half-depth-saturation}, put \(F_e=(Qj)_e\)
for any integral column \(j\). For every \(g\in L_r(W)\),
every exponent \(\beta\) and every letter \(i\),
\[
       v\bigl((gF_e)_\beta\bigr)\ \ge\ \tfrac12 E_{i,r-1}(\beta)
\]
with the coefficient order of
Definition~\ref{def:root-coefficient-order}. In particular these
coefficients belong to \(\varrho^c O\) when \(E_{i,r-1}(\beta)>2c\)
for any nonnegative integer \(c\).
\end{theorem}
\begin{proof}
The prime is inert in the quadratic field. The primes \(2,3\)
divide \(M\) and split there by Proposition~\ref{prop:places}.
Thus \(p\ge5\) and \(P=p^o\ge5\).
Take an integer \(b_0\ge0\) such that \(\varrho^{b_0}A^{-1}\)
is integral. This is possible because \(A\) is finite and
invertible over the fraction field. Put \(F^*=\varrho^{b_0}A^{-1}j\)
and fix one source component \(f\) of \(F^*\).
Every entry of \(JF^*=\varrho^{b_0}Qj\) is divisible by
\(\varrho^{b_0}\). By the coset coverage and integrality
argument of Proposition~\ref{prop:two-half-depth-saturation},
in every full frame \(g'\in{\rm GL}_k(W)\) one therefore has
\[
            v((g'f)_\beta)+d(\beta)\ge b_0 .
\]
Here \(v(\varrho^n z)=n+v(z)\) for integral \(z\ne0\) and
integer \(n\) whenever the product is integral, directly
from the order definition.

Fix \(i\) and abbreviate \(E=E_{i,r-1}\).
Write \(d_z\) for the zero or one depth of the individual
variable \(z\) and \(E_z\) for its coefficient in \(E\).
Use coordinate weights \(c_z=d_z-E_z/2\).
For the boundary pair, \(y_s\) is deep and \(x_s\) shallow at
both parities. Hence
\[
 c_{y_s}-c_{x_s}=
 \begin{cases}
       (1-q^{-1})/2,&s=i,\\
       1-M_s/q,&s\ne i .
 \end{cases}
\]
These numbers belong to \([2/5,8/5]\), since \(q\ge P\ge5\)
and \(M_s/q\le 1/P\) for \(s\ne i\).

If \(f=0\) the component contributes zero. Otherwise minimize
order plus \(c\)-weight among all terms of all \(g'f\). A minimum
is attained: there are finitely many exponent vectors of the
fixed multidegree, and for each occurring exponent the
nonnegative integral orders have a least member. Choose an
attaining frame and exponent \(\alpha\). Every translate of the
polynomial in that frame is still a translate of \(f\). Thus
Theorem~\ref{thm:root-extremum-by-one-polynomial} gives
\(E(\alpha)\le0\). The minimum is at least \(b_0\) by the
displayed depth bound, since subtracting \(E(\alpha)/2\)
increases it. This proves, for every \(g',\beta\),
\[
 v((g'f)_\beta)+d(\beta)-b_0\ \ge\ E(\beta)/2 .
\]
At an edge from this source component,
\(\varrho^{b_0}F_e=D_\varrho P_\lambda(h_e)f\).
The action of a block diagonal \(g\) commutes with \(D_\varrho\):
it sends each variable to a sum of variables of the same depth.
Apply \(g\) and use the preceding estimate at \(g'=gh_e\).
This gives the coefficient inequality, and order at least
\(c\) implies divisibility by \(\varrho^c\).
\end{proof}


\section{Counting residue polynomials by flags}

\begin{theorem}\label{thm:full-served-residue-dimension}
Let \(p,\lambda,r,D,P,q,k,a\) and \(U(h)\) be as in
Definition~\ref{def:served-excess-forms}, with \(P\ge5\) and
\(h\ge0\). Let \(d_{\mathcal P}=[K:\mathcal P]\) be the
index of the two-block parahoric of
Theorem~\ref{thm:served-actual-hecke-class-matrix}. For
\(H=h+D(k+1)\) and \(A_s=a+k+1\),
\[
 d_{\mathcal P}\dim_{\mathbf k} U(h)\ \le\
     F_k(q)\,{\rm vol}(\Omega_H^{(r)}\cap\Sigma(A)).
\]
In particular, with \(h=x^{40}\), uniformly at every served
prime for all sufficiently large \(x\),
\[
     d_{\mathcal P}\dim_{\mathbf k}U(x^{40})
                \le (1+x^{-14})\binom{a+k-1}{k-1}^{D}.
\]
\end{theorem}
\begin{proof}
Write \(k_q=W/pW\) and \(L={\rm GL}_r(k_q)\times
{\rm GL}_{k-r}(k_q)\). Any element of \(L\) lifts to
\(L_r(W)\) by lifting entries, since a determinant with
nonzero residue is a unit. The action on residue polynomials
factors through \(L\), by the successive residue embeddings
of Lemma~\ref{lem:served-actual-frobenius-cycle}. Let
\(B_L\subset L\) be the two lower triangular subgroups and let
\(\mathcal T\) be representatives of \(B_L\backslash L\)
including the identity.

For each integer \(j\ge1\) put \(m'=q^j\). Take one polynomial
\(F_u\in U(h)\) for each array of residues
\(0\le u_{b,s}<m'\). Increase \(u_{r-1,s}\) by
\(m'(k-\lfloor\sum_bu_{b,s}/m'\rfloor)\) at each letter \(s\),
and write \(\gamma(u)\) for the resulting array. All its entries
are nonnegative and in each letter its sum lies in
\([km',(k+1)m')\). Form
\[
   B(X)=\sum_u X^{\gamma(u)}F_u(X^{m'}),\qquad
   X^{m'}=(X_{b,s}^{m'})_{b,s}.
\]
This is linear over \(\mathbf k\), including when its scalars
are outside \(k_q\). Substitution changes only the variables.
It is injective on the tuple \((F_u)_u\): distinct summands
have distinct exponent arrays modulo \(m'\), and substitution
is injective on each summand. Every residue matrix entry
under the polynomial action is fixed by raising to \(m'\).
Indeed the entries of each letter lie in an image of \(k_q\)
and \(z^q=z\) in every such image. Taking \(m'\)-th powers
distributes over sums of variables in characteristic \(p\).
It follows by substitution on the variables that for \(g\in L\)
\[
        gB=\sum_u (gX^{\gamma(u)})(gF_u)(X^{m'}).
\]

Give \(X_{b,s}\) weight \(b\) at even \(s\), and \(-b\) at odd
\(s\). Take a nonzero tuple and choose a term \(\alpha\) of
least weight among the polynomials \(g_0B,\ g_0\in\mathcal T\).
This attains the least weight in all \(gB,\ g\in L\).
Indeed left multiplication by an element of \(B_L\) replaces
a variable in an even letter by a sum from its own and lower
rows. In an odd letter its inverse transpose uses its own and
upper rows. Both operations preserve or increase the weight
of every term. They give every \(g\) by coset coverage. For all
adjacent pairs within one block, the weights satisfy
\(c(y_s)-c(x_s)=1\). For each such \(b\), the two-row subgroup
lies in \(L_r(W)\). Apply
Theorem~\ref{thm:root-extremum-by-one-polynomial} to \(g_0B\)
in the residue field. The simultaneous conclusions are
\(E_{i,b}(\alpha)\le0\) for all \(i\) and \(b+1\ne r\).

At the boundary, the nonzero coefficient at \(\alpha\) is
a sum with at least one nonzero contributing term
\(\alpha=m'\delta+\gamma'\), where \(\delta\) occurs in \(g_0F_u\)
and \(\gamma'\) occurs in \(g_0X^{\gamma(u)}\).
Each letter of \(\gamma'\) has the same total degree as that
of \(\gamma(u)\) by the linear substitutions within letters.
The coefficients of \(E_{i,r-1}\) are at most one by
Definition~\ref{def:cyclic-root-excess}.
Therefore \(E_{i,r-1}(\gamma')\le D(k+1)m'\).
By Definition~\ref{def:served-excess-forms} the other
contribution satisfies \(E_{i,r-1}(\delta)\le h\).
Consequently \(\alpha/m'\in\Omega_H^{(r)}\).
The multidegrees also give \(\alpha/m'\in\Sigma(A)\), because
\(\delta\) has degree \(a\) at each letter.

Put \(S=\Omega_H^{(r)}\cap\Sigma(A)\). Restricting the
coefficients of each \(g_0B\), \(g_0\in\mathcal T\), to
\(\alpha\in m'S\cap\Z^{kD}\) thus detects every nonzero tuple.
The dimension inequality for this injective linear map is
\[
      (m')^{kD}\dim U(h)\ \le\
                 |\mathcal T|\,\#(m'S\cap\Z^{kD}).
\]
The ratio of the lattice count to \((m')^{kD}\) tends to
\({\rm vol}(S)\) as \(j\to\infty\). Indeed the cubes of side
\(1/m'\) with a grid vertex in \(S\), using that vertex as
the lower corner, agree with \(S\) up to cubes meeting its
finite union of boundary hyperplanes in a fixed bounded box.
Their total volume tends to zero by the grid argument of
Theorem~\ref{thm:cyclic-flag-share}.

Reduction of \(K\) modulo \(p\) is onto \({\rm GL}_k(k_q)\) by
the same unit-determinant lifting argument. The lower block
parabolic stabilizes the subspace spanned by the last \(k-r\)
coordinate vectors, and \(\mathcal P\) is its inverse image.
Thus \(d_{\mathcal P}\) counts the \((k-r)\)-dimensional
subspaces: every basis of one can be extended to a basis of
the whole space. For each such subspace, a full flag through
it is determined by a full flag in it and one in its
\(r\)-dimensional quotient. The lower triangular groups are
the stabilizers of the corresponding standard flags, and
transitivity follows by taking a basis adapted to a flag.
Hence the number of flags extending each subspace is
\(|\mathcal T|\). By Theorem~\ref{thm:cyclic-flag-share},
\(d_{\mathcal P}|\mathcal T|=F_k(q)\). Multiplication of the
limit inequality by \(d_{\mathcal P}\) proves the first claim.
Proposition~\ref{prop:thick-cyclic-flag-volume} proves the
second one.
\end{proof}


\section{The index of a pairing of two matrices}

\begin{lemma}\label{lem:transverse-pairing-index}
Let \(O\) be a valuation ring with maximal ideal \(\mathfrak m\)
and residue field \(\mathbf k\). Let \(V=O^n,\ W'=O^N\),
let \(Q,M:V\to W'\), \(\beta:W'\to V\) be integral linear
maps, and suppose \(RQ=I_n\) for an integral \(R:W'\to V\).
Let \(a,z\in O\setminus\{0\}\) and \(w\in\mathfrak m\setminus\{0\}\).
Suppose
\[
   \beta M\equiv a I_n\bmod aw,\qquad
   Qu-Mv=awF\quad\Longrightarrow\quad\bar F\in S
\]
for all \(u,v\in V,\ F\in W'\), with one fixed subspace
\(S\subset W'/\mathfrak mW'\). There is an integer
\(0\le b\le\dim_{\mathbf k} S-n\) such that
\[
       \mathcal I(V,\beta QV+zV;\ a^{n-b}z^b).
\]
Here \(\mathcal I\) is the divisibility of determinant indices
in Definition~\ref{def:det-index}.
\end{lemma}
\begin{proof}
Write \(\beta M=a(I+wB)\). The second factor has determinant one
in the residue field and is invertible over \(O\) by the
adjugate formula. Right multiplication of \(M\) by its inverse
preserves the difference hypothesis and gives \(\beta M=aI\).
Write \(W'=Q V\oplus W_0\), where \(W_0=\ker R\).
This is a finite free summand: it is generated by the columns
of \(I-QR\), and the rectangular pivot argument of
Lemma~\ref{lem:valuation-index} gives a basis of their span
by clearing the row and column of each successive nonzero pivot.
Write \(M=QM_1+M_0\) accordingly. By the same argument there
is a basis \(v_1,\ldots,v_n\) of \(V\) with
\(M_0v_j=t_j u_j\) at its nonzero pivots, where those \(u_j\)
form part of a basis of \(W_0\). Use \(t_j=0\) on the kernel.

Call an index exceptional when \(t_j\notin awO\), and let
\(b\) count them. Then \(t_j\ne0\) and \(aw/t_j\in O\)
by total divisibility. With \(v=(aw/t_j)v_j,\ u=M_1v\)
the difference is \(-aw u_j\), so \(\bar u_j\in S\).
Also \(\bar QV\subset S\) by taking \(u\in aw V,\ v=0\).
The \(b\) reductions are independent and independent of the
\(n\) basis images of \(Q\), by the direct-sum bases of \(W'\).
This proves \(n+b\le\dim S\).

For the other indices \(M_0v_j=aw u'_j\) with integral
\(u'_j\), including zero. Therefore
\[
       \beta Q M_1v_j=a(v_j-w\beta u'_j).
\]
Replace these \(v_j\) by \(v'_j=v_j-w\beta u'_j\) and leave
the exceptional basis elements fixed. Their residues still
form a basis, so their determinant is a unit and the resulting
vectors are a basis of \(V\). In this basis \(\beta QV+zV\)
contains \(a v'_j\) in each of the \(n-b\) indicated directions
and \(zv_j\) in each exceptional direction. Applying
Lemma~\ref{lem:valuation-index} to this diagonal submodule
gives the claimed determinant divisibility.
\end{proof}


\section{Weighted traces on the tensor powers of a Levi subgroup}

\begin{lemma}\label{lem:integral-monomial-pairing}
For \(p\) a prime and \(f_0\ge1\), let \(\gamma\) be a vector
of \(u\ge1\) nonnegative exponents of sum \(d<p^{f_0}\).
Set \(B_\gamma=p^{u f_0}\prod_j\gamma_j!/d!\). Then
\(B_\gamma\) and \(p^{u f_0}/B_\gamma\) are integral at \(p\).
With \(M_d\) as in Lemma~\ref{lem:one-row-lattice-transpose},
\[
     B_\alpha M_d(g)_{\alpha\beta}
          =B_\beta M_d(g^t)_{\beta\alpha}
\]
for any \(u\)-by-\(u\) matrix \(g\) in characteristic zero.
For the representation \(P_\lambda\) take \(f_0=oD\) and
\(a<q=p^{f_0}\). The integral diagonal matrix \(B\) whose
entry at \(\beta\) is \(\prod_s B_{\beta_s}\), using \(u=k,d=a\),
satisfies
\[
     B P_\lambda(g^t)=P_\lambda(g)^t B,\qquad
     p^{k f_0D}B^{-1}\text{ is integral}.
\]
\end{lemma}
\begin{proof}
The reciprocal quotient \(d!/\prod_j\gamma_j!\) is the number
of words with the indicated multiplicities. The number of
factors of \(p\) in \(d!\) is \(\sum_{e\ge1}\lfloor d/p^e\rfloor\)
by counting one contribution from each multiple of \(p^e\).
Subtracting the corresponding sums for the \(\gamma_j!\)
gives, for each \(e\), at most \(u-1\) by adding the remainders
modulo \(p^e\). The contributions for \(e\ge f_0\) are zero.
This proves integrality. An array \((n_{ij})\) of nonnegative
integers with row sums \(\alpha\), column sums \(\beta\)
contributes to \(M_d(g)_{\alpha\beta}\) the product of
\(g_{ij}^{n_{ij}}\) times \(\prod_j\beta_j!/\prod_{i,j}n_{ij}!\)
by expansion of the input monomial. Multiplication by
\(B_\alpha\) makes the factor symmetric in the two sets of
marginals. Transpose the array and sum. The tensor identity
follows letter by letter since transpose commutes with taking
the inverse transpose, and the diagonal denominators multiply.
\end{proof}

\begin{theorem}\label{thm:served-invariant-class-mean}
Take \(W,O,p,q=p^{f_0}\) and one embedding \(\sigma:W\to O\)
of Definition~\ref{def:served-column-cycle}. Let \(L=L_r(W)\),
let \(X=(X_{bc})\) be a generic \(k\)-by-\(k\) matrix, and put
\(F^g(X)=F(\sigma(g)^{-1}X\sigma(g))\). Give the rows and
columns real weights \(x_b\). For a monomial \(X^\alpha\) put
\[
   v_R(\alpha)=\sum_{bc}x_b\alpha_{bc},\qquad
   v_C(\alpha)=\sum_{bc}x_c\alpha_{bc}.
\]
Let \(J\subseteq\R\), and let \(d_*\ge0\) be an integer
satisfying \(d_*+k<q\). There is an integral linear functional
\(\kappa_J\) on \(O[X]\) with the following properties.
\begin{enumerate}
\item \(\kappa_J(F^g)=\kappa_J(F)\) for all \(g\in L\).
It vanishes on degrees above \(d_*\) and on monomials
whose row weight differs from their column weight.
\item If \(v_R(\alpha)\in J\) and \(|\alpha|\le d_*\), then
\(\kappa_J(X^\alpha)=p^{k^2 f_0}(X^\alpha)(I_k)\).
\item It is a finite integral linear combination of functionals
\(F\mapsto l(F^g)\), with \(g\in L\) and \(l(X^\alpha)=0\)
whenever \(v_R(\alpha)\notin J\). One may take each \(l\)
supported in a single degree at most \(d_*\).
\end{enumerate}
\end{theorem}
\begin{proof}
Work first in fixed degree \(d\le d_*\).
Let \(T=(O^k)^{\otimes d}\) with basis \(e_t\) indexed by
words of length \(d\) in the row indices, including the empty
word at \(d=0\). Write \(A(g)=\sigma(g)^{\otimes d}\).
Associate to \(X^\alpha\) the endomorphism \(K(X^\alpha)\)
whose entry at output word \(s\) and input word \(t\) is
\(B_\alpha\) from Lemma~\ref{lem:integral-monomial-pairing}
with \(u=k^2\), if both of the following hold: the ordered
pairs \((t_j,s_j)\) have multiplicities \(\alpha\), and the
two entries of each pair belong to the same block. All other
entries are zero. Extend \(K\) linearly. It is integral.
It satisfies
\[
                K(F^g)=A(g)K(F)A(g)^{-1}.
\]
Indeed deleting monomials involving off-block variables amounts
to first restricting \(X\) to a block diagonal matrix, which
commutes with the conjugation. In the fraction field, write a
monomial as \(\prod_{j=1}^d X_{b_j c_j}\). Its matrix \(K\)
on the block variables is \(p^{k^2f_0}/d!\) times the sum
of all permutations of the tensor factors of
\(\bigotimes_j E_{c_j b_j}\), with \(E_{cb}\) the elementary
matrix. Each allowed entry occurs \(\prod_{bc}\alpha_{bc}!\)
times. In degree one the substitution corresponds to
\(E_{cb}\mapsto \sigma(g)E_{cb}\sigma(g)^{-1}\), by matrix
multiplication. Tensoring, permuting and summing gives the
identity in every degree, hence also over \(O\).

We need an integral span assertion. For a block diagonal
matrix \(Y\), the entries of \(Y^{\otimes d}\) are either zero
or the degree-\(d\) monomials in its block entries. The
evaluations of the complete list of these monomials on
invertible residue block matrices span the coordinate space
over the residue field of \(O\). In fact a linear dependence
annihilating all these evaluations gives a polynomial \(P_0\).
Multiply it by the product of the two block determinants.
The product vanishes on all tuples of block entries from
\(\bar\sigma(k_q)\), since a noninvertible tuple has zero
determinant. Its degree in each variable is at most
\(d+k<q\). A univariate polynomial of degree less than \(q\)
vanishing at \(q\) distinct points is zero by division by its
linear factors. Applying this in one variable at a time
proves the same assertion for this grid, even with
coefficients in the larger residue field. The determinant
polynomial is nonzero, so \(P_0=0\). This proves spanning by
ordinary row reduction. Select a square minor with nonzero
residue determinant, lift its evaluation matrices to \(L\)
as in Theorem~\ref{thm:full-served-residue-dimension}, and
invert the minor over \(O\) by its unit determinant. The
coefficient matrices of \(Y^{\otimes d}\) consequently lie
in the integral span of a finite list of \(A(g)\).
These coefficient matrices have ones at a specified set of
word pairs and zero elsewhere. Both all \(A(g)\) and the
\(K(X^\alpha)\) are integral linear combinations of them,
taking account of the transposed indexing of \(K\).

Let \(T_J\) be the span of \(A(g)e_t\) for all \(g\in L\)
and words \(t\) with \(\sum_j x_{t_j}\in J\). It has a
finite generating set by the preceding finite span assertion.
The pivot proof of Lemma~\ref{lem:valuation-index} therefore
gives a free basis in its span. It is stable under all \(A(g)\)
by the group law and under all \(K(F)\) by that same span
assertion. Define \(\kappa_J(F)\) in this degree as the trace
of \(K(F)\) on \(T_J\), the sum of the diagonal entries in a
basis. It is integral. Trace is invariant under change of
basis because \({\rm tr}(UV)={\rm tr}(VU)\), by expanding
the double sum of entries. The covariance identity proves
conjugation invariance. A matrix \(K(X^\alpha)\) acts only
on words of weight \(v_R(\alpha)\) and its outputs have weight
\(v_C(\alpha)\). At distinct weights its square is zero,
so its trace on \(T_J\) vanishes: over the fraction field
take a basis of its image, extend to its kernel and then
to the entire space. Its diagonal in this basis is zero.

If \(v_R(\alpha)\in J\), every basis word on which \(K(X^\alpha)\)
can act is already in \(T_J\). Its image on all of \(T\) is
therefore in \(T_J\). Extending a basis of the latter over the
fraction field gives zero on the quotient, so its trace on
\(T_J\) equals its trace on \(T\). No diagonal entry is
possible for a monomial using an off-diagonal variable. For a
diagonal monomial the number of words giving a diagonal entry
is \(d!/\prod_b\alpha_{bb}!\), and each entry is \(B_\alpha\).
This gives \(p^{k^2f_0}\) exactly as in the stated evaluation.

Finally choose a basis \(v_j\) of \(T_J\) with coordinate
maps \(v_j^*\), and write each \(v_j\) as an integral sum of
\(A(g)e_t\) at the chosen weights. In the trace formula
\(\sum_j v_j^*(K(F)v_j)\) each summand becomes a combination of
\[
          v_j^*\bigl(A(g)K(F^{g^{-1}})e_t\bigr).
\]
The map \(l:H\mapsto v_j^*(A(g)K(H)e_t)\) is integral and
kills monomials of row weight outside \(J\) by the word
description. This gives the required conjugate expression
in degree \(d\), also when \(T_J=0\). Sum the finitely many
functionals over \(0\le d\le d_*\), extending each by zero on
other degrees. Conjugation preserves degree, which proves all
three properties on \(O[X]\).
\end{proof}


\section{Frobenius products and the weight of a matrix coefficient}

\begin{definition}\label{def:served-generic-operator}
Use Definition~\ref{def:served-excess-forms}. Set \(t=D/2,\ B_0=P^t\)
and \(x_b=({\bf1}_{b=r}-{\bf1}_{b=r-1})/B_0\).
Let \(X_s\) be independent generic \(k\)-by-\(k\) matrices.
For \(H\in L_r(W)\) write \(H_s=\sigma_s(H)\). With entrywise
powers denoted by brackets, put, for \(0\le d\le t\),
\[
 \Psi^H_{s,d}=H_s
       (H_{s-d}^{-1}X_{s-d}H_{s-d})^{[P^d]}H_s^{-1},\qquad
 Z^H_{s,j}=\Psi^H_{s,j-1}\cdots\Psi^H_{s,0},\quad Z^H_{s,0}=I.
\]
Set \(j_b=t+d_{b,s}\) with \(d_{b,s}\in\{0,1\}\) the depth
of the variable in row \(b\) of letter \(s\). Define \(A_s^H\)
by taking its row \(b\) from \(Z^H_{s,j_b}\) at even \(s\)
and from \({\rm adj}(Z^H_{s,j_b})^t\) at odd \(s\). Put
\[
       \mathcal F_H(X)=\bigotimes_s M_a(A_s^H),\qquad
       \mathcal F=\mathcal F_1 .
\]
The adjugate is the transpose of the matrix of signed minors.
For any monomial in all entries let \(\nu_s\) denote its
exponent array in \(X_s\). Use the row and column weights
\(v_R,v_C\) of Theorem~\ref{thm:served-invariant-class-mean}
with the above \(x_b\).
\end{definition}

\begin{lemma}\label{lem:generic-operator-covariance}
The matrix \(\mathcal F_H\) equals \(P_\lambda(H)\mathcal F
((H_s^{-1}X_sH_s)_s)P_\lambda(H)^{-1}\), and
\(\mathcal F((I)_s)=I\). Each matrix entry has degree, in every
\(X_s\), at most \(d_*=Da k(t+2)B_0\). At all sufficiently
large served scales one has \(Da k(t+2)+a+k<B_0\), and hence
\(d_*+k<q\).
\end{lemma}
\begin{proof}
In a fixed product the conjugations by \(H_s\) telescope.
Taking the transpose of an adjugate changes conjugation by
\(H_s\) to conjugation by \(H_s^{-t}\). For invertible matrices
this follows from \({\rm adj}(Z)=\det(Z)Z^{-1}\) by the
cofactor expansion. For arbitrary matrices apply the identity
to \(Z+uI\) over the field of rational functions in \(u\)
and specialize the polynomial identity at \(u=0\).
Rows at each length form a block preserved by both \(H_s\)
and \(H_s^{-t}\), so selecting rows commutes with the respective
left change. Right changes are common to the whole letter.
The coefficient action \(M_a\) is multiplicative by composition
of the substitutions. This proves covariance. Setting every
matrix equal to \(I\) gives the stated specialization.
Each \(\Psi^H_{s,d}\) is homogeneous of degree \(P^d\le B_0\).
A product of at most \(t+1\) of these matrices has degree
at most \((t+1)B_0\); its minors of size \(k-1\) have degree
at most \((k-1)(t+1)B_0\). The tensor touches \(Da\) entries
counting repetitions. This proves the asserted degree bound.
At a served prime \(D\ge m\) by
Lemma~\ref{lem:served-actual-frobenius-cycle}, so
\(B_0\ge2^{m/2}\). By Proposition~\ref{prop:thick},
\(m\ge\sqrt{x}\) up to a fixed positive factor and
\(0<\epsilon<1/100\). The bounds \(D\le x^{12}, k=O(x^5)\)
of Proposition~\ref{prop:thick-cyclic-flag-volume} and
\(a\le2x\exp(x^\epsilon)\) of Proposition~\ref{prop:level-prime}
therefore give \(Da k(t+2)+a+k<B_0\) eventually. Multiply
by \(B_0\ge1\) and use \(B_0^2=q\) for the last assertion.
\end{proof}

\begin{theorem}\label{thm:generic-frobenius-weight}
Assume \(Da k(t+2)\le B_0\). Fix a letter \(i\) and a
monomial \(\nu\) with \(v_R(\nu_j)=v_C(\nu_j)\) for all
\(j\ne i\). For any \(H\in L_r(W)\) and exponent vectors
\(\beta,\gamma\) of the coefficient representation,
let \(z=[X^\nu](\mathcal F_H)_{\beta\gamma}\).
Then, with the order of Definition~\ref{def:root-coefficient-order},
\[
     2v(z)+1\ \ge\
       \left|E_{i+t,r-1}(\beta)-v_R(\nu_i)\right| .
\]
\end{theorem}
\begin{proof}
Fix a sign \(\epsilon=\pm1\) and put
\[
 f_n(b)=\frac{\epsilon x_b P^n}{2q},\qquad u=\frac1{2B_0}.
\]
Give the variable \((X_j)_{bc}\) the weight \(f_D(b)\) if
\(j=i\), and \(f_{[j-i]_D}(b)-f_{[j-i]_D}(c)\) otherwise.
Its absolute weight is at most \(1/B_0\). For a fixed output
letter \(s\) put
\(\lambda_d(b)=f_{d+[s-d-i]_D}(b)\).
Each \(\lambda_d\) has sum zero over \(b\). At \(d=0\) its
absolute value is at most \(u\), and for \(0\le d\le t+1\)
it is at most \(1/2\), since \(d+[s-d-i]_D\le D+t\).

Every term of entry \(bc\) of \(\Psi^H_{s,d}\) has cost at
least \(\lambda_{d+1}(b)-\lambda_d(c)-u\), where cost is
coefficient order plus the weight of its variables. To see
this, reduce its formula modulo \(p\). By
Lemma~\ref{lem:served-actual-frobenius-cycle}, conjugation
and the entrywise power now commute with the change of
embedding from \(s-d\) to \(s\). The matrix reduces to
\(X_{s-d}^{[P^d]}\), because raising to \(P^d\) is additive
and multiplicative in characteristic \(p\). Thus the difference
has coefficients of order at least two. Its terms have degree
\(P^d\), with variable weight at least \(-P^d/B_0\ge-1\);
their costs are at least one, which bounds the claimed
difference of the \(\lambda\)'s. For the pure entry put \(j=s-d\).
If \(j\ne i\), then \([j-1-i]_D=[j-i]_D-1\); multiplying
the variable weight by \(P^d\) gives exactly
\(\lambda_{d+1}(b)-\lambda_d(c)\). If \(j=i\), the weight
after multiplication is \(P^d f_D(b)=\lambda_{d+1}(b)\)
and \(|\lambda_d(c)|=|f_d(c)|\le u\). This proves the
entry estimate, also if terms combine by addition.

Multiplication adds the lower bounds for cost. Summing along
a matrix-product path telescopes its intermediate weights:
entry \(bc\) at length \(j\) has cost at least
\(\lambda_j(b)-\lambda_0(c)-ju\). At an odd letter the
signed minor deletes row \(b\) and column \(c\), so summing
the same bounds over its permutations gives at least
\(-\lambda_j(b)+\lambda_0(c)-(k-1)ju\), using the two
zero sums. With \(j=t+d_{b,s}\) both parity cases give the
bound
\[
     (-1)^s\lambda_{t+d_{b,s}}(b)-k(t+2)u
\]
on entry \(bc\) of \(A_s^H\). The first term here equals
\(\epsilon E_{i+t,r-1}({\bf e}_{b,s})/2\), where
\({\bf e}_{b,s}\) is the unit exponent vector. Indeed the
numerator is zero outside the boundary rows. At the shallow
one its signed weight is \(-P^{[s-i-t]_D}/q\), and at the
deep one it is \(P^{1+[s-1-i-t]_D}/q\), before multiplication
by \(\epsilon/2\). These are precisely the coefficients in
Definition~\ref{def:cyclic-root-excess}.

In an entry of \(M_a(A_s^H)\) with output \(\beta_s\), each
term is a product of \(a\) entries using row \(b\) exactly
\(\beta_{b,s}\) times, by polynomial expansion. The full
tensor thus has cost at least
\(\epsilon E_{i+t,r-1}(\beta)/2-Da k(t+2)u\).
By hypothesis the subtracted term is at most \(1/2\).
For the specified monomial the variable weight outside \(i\)
vanishes, since each letter contributes a scalar times its
row weight minus its column weight. The remaining weight
is \(\epsilon v_R(\nu_i)/2\).
Subtract it from the cost bound and use both signs to obtain
the claimed absolute value inequality.
\end{proof}


\section{The integral operator selecting the permitted weights}

\begin{definition}\label{def:integral-boundary-forms}
For \(c\ge1\) integral and \(\theta\in\R\), let
\(\mathcal U(\theta,c)\subset\mathcal V\) be the integral
polynomials for which every coefficient at \(\beta\) of
\(gF,\ g\in L_r(W)\), lies in \(\varrho^c O\) whenever
\(E_{i,r-1}(\beta)>\theta\) for some \(i\).
This is an \(O\)-submodule since the coefficient tests and
the action are integral linear. The group law makes it
stable under \(L_r(W)\).
\end{definition}

\begin{theorem}\label{thm:served-integral-approximate-projector}
Use the arithmetic range of Lemma~\ref{lem:generic-operator-covariance},
put \(f_0=oD,\ K_0=(k^2+k)f_0D\), and let \(B\) be the
diagonal matrix of Lemma~\ref{lem:integral-monomial-pairing}.
For every \(c\ge1\) integral and real \(\theta\) there is an
integral matrix \(Z\) such that
\[
  Z\mathcal V\subset\mathcal U(\theta,c),\qquad
  y^t BZ\equiv p^{K_0}y^t\pmod{\varrho^c}
       \quad\bigl(y\in\mathcal U(\theta-4c-2,c)\bigr).
\]
The congruence is entrywise in the row.
\end{theorem}
\begin{proof}
Use the degree bound \(d_*\) and variable weights of
Definition~\ref{def:served-generic-operator} and
Lemma~\ref{lem:generic-operator-covariance}.
At each letter \(s\), take the two functionals from
Theorem~\ref{thm:served-invariant-class-mean} for the map
\(\sigma_s\), with weight sets \(\R\) and
\((\theta-2c-1,\infty)\). Call them \(a_s\) and \(b_s\),
and set \(l_s=a_s-b_s\). All three are conjugation invariant
and annihilate unbalanced monomials, meaning those with
\(v_R\ne v_C\). Also \(l_s\) vanishes whenever
\(v_R>\theta-2c-1\), by their common scaled evaluation in
the allowed degrees and common zero above the cutoff.

Given a functional \(u_s\) at each letter, denote by
\([u_s]_s F\) the result of applying their tensor product
entrywise to a polynomial matrix \(F\). Explicitly a monomial
\(X^\nu\) is replaced by \(\prod_s u_s(X_s^{\nu_s})\) and
the result is extended by linearity. This also shows
multilinearity in the functionals and compatibility with
constant matrices on each side. Put
\[
          Z_0=[l_s]_s\mathcal F,\qquad B'=p^{k f_0D}B^{-1},
          \qquad Z=Z_0 B'.
\]
All these entries are integral by
Lemma~\ref{lem:integral-monomial-pairing}.

For \(H\in L_r(W)\), invariance in all letters and covariance
in Lemma~\ref{lem:generic-operator-covariance} give
\(P_\lambda(H)Z_0=([l_s]_s\mathcal F_H)P_\lambda(H)\).
Fix a row \(\beta\) for which \(E_{i+t,r-1}(\beta)>\theta\).
Every monomial surviving the functionals is balanced outside
\(i\) and has \(v_R(\nu_i)\le\theta-2c-1\).
Theorem~\ref{thm:generic-frobenius-weight} gives order at
least \(c\) for its coefficient. Consequently the indicated
row of \(P_\lambda(H)Z_0 B'\) is divisible by \(\varrho^c\).
All letters occur as \(i+t\), proving the image assertion.

Evaluation with \([a_s]_s\) gives \(p^{k^2 f_0D}I\) on
\(\mathcal F\), since every letter has degree at most \(d_*\)
and evaluation at \((I)_s\) gives \(I\) by
Lemma~\ref{lem:generic-operator-covariance}. Expand
\([a_s]_s-[l_s]_s\) by replacing one factor at a time in any
fixed order. Each difference has \(b_i=a_i-l_i\) at the
replaced letter, with all other factors invariant and
annihilating unbalanced monomials. In the finite expansion
of \(b_i\) supplied by Theorem~\ref{thm:served-invariant-class-mean},
each term applies an integral functional killing
\(v_R(\nu_i)\le\theta-2c-1\), after a conjugation by some
\(H\in L_r(W)\). Invariance at the other letters permits the
same conjugation everywhere. Covariance therefore writes the
corresponding matrix as
\[
          P_\lambda(H)^{-1} A P_\lambda(H),
\]
where \(A\) is obtained by applying the untransported functionals
to \(\mathcal F_H\). In its entries every surviving monomial
is balanced away from \(i\) and has
\(v_R(\nu_i)>\theta-2c-1\).

For \(y\) in the stated submodule let \(y'=P_\lambda(H^{-t})y\).
The matrix \(H^{-t}\) is block diagonal, so \(y'\) lies in
the same submodule. By Lemma~\ref{lem:integral-monomial-pairing},
\(y^t B P_\lambda(H)^{-1}=(y')^tB\).
For any row \(\beta\) of \(A\) with
\(E_{i+t,r-1}(\beta)>\theta-4c-2\), the factor \(y'_\beta\)
is divisible by \(\varrho^c\). In every other row the surviving
monomials have \(v_R(\nu_i)-E_{i+t,r-1}(\beta)>2c+1\).
Their coefficients are divisible by \(\varrho^c\) by
Theorem~\ref{thm:generic-frobenius-weight}. Hence the product
of \(y^t B\) with this entire conjugate matrix is divisible
by \(\varrho^c\). Sum over the telescoping expansion and
multiply by the integral \(B'\). Since \(BB'=p^{k f_0D}I\),
the contribution from \([a_s]_s\mathcal F\) is exactly
\(p^{K_0}y^t\). This proves the row congruence.
\end{proof}


\section{The elementary divisors of a single marked operator}

\begin{definition}\label{def:served-truncated-cost}
For an integral endomorphism \(T\) of a free module \(V\) of
finite rank over a valuation ring with valuation \(v_p(p)=1\),
and \(s\in\Z_{\ge0}\), let
\[
                 d_s(T)=\sum_j\min\{s,v_p(t_j)\}.
\]
Here \(t_j\) are the diagonal entries in any Smith reduction
of \(T\), and \(v_p(0)=+\infty\). By
Lemma~\ref{lem:valuation-index}, \(\mathcal I(V,TV+p^sV;z)\)
holds precisely when \(v_p(z)\ge d_s(T)\) for \(z\ne0\) in the
ring. Indeed each nonzero generator of \(t_jO+p^sO\) has the
valuation in the sum, and the product of such generators is
in the ring. This also shows independence of the reduction.
\end{definition}

\begin{theorem}\label{thm:single-served-actual-index}
For each sufficiently large arithmetic scale, each served
\(p\), each \(\lambda\in\Lambda\), and every \(s\in\Z_{\ge0}\),
the endomorphism \(Y_\lambda\) of the integral columns
\(\mathcal E_p\) in Definition~\ref{def:served-scaled-neighbours}
satisfies
\[
    d_s(Y_\lambda)\ \le\
      c_0\left(\frac{3A_\lambda}{mp}+x^{40}+\frac{s}{x^{14}}\right),
       \qquad c_0=\dim{\rm Col}.
\]
\end{theorem}
\begin{proof}
Use the two edge matrices \(J_b,J_k\) of
Proposition~\ref{prop:two-half-depth-saturation}, with frames
\(b_e^{-1},k_e^t\) and sources \(c(e),c'(e)\), respectively.
Tensor them with the identity on \(H_{\rm rest}\) and reorder
the factors in the same way on each side. Write \(V\) for the
class-coordinate module of rank \(c_0\) over \(O\), and \(W'\)
for the edge-coordinate module. Each edge has \(M_0\) linked
coordinates with \(M_0=\binom{a+k-1}{k-1}^D\), and
\(h'=\dim H_{\rm rest}\) remaining coordinates, so
\(c_0=|C|h'M_0\) by
Theorem~\ref{thm:served-actual-hecke-class-matrix}. Write
\(J_j=Q_j A_j,\ R_j Q_j=I\) for \(j=b,k\) by tensoring
the displayed factorizations of the proposition in these
same coordinates. Theorem~\ref{thm:served-saturation-excess}
applies on every edge and each fixed remaining coordinate.

Extend the diagonal \(B\) of Lemma~\ref{lem:integral-monomial-pairing}
to each class or edge independently, writing \(B_V,B_{W'}\)
for the two extensions. Let \(U\) act on each edge by its
\(u_e\) on the remaining factor. This is an integral automorphism
and commutes with \(B_{W'}\). In the chosen gauges,
\[
 Y_\lambda B_V
      = J_b^t\, U B_{W'}\, J_k
      = A_b^t\,(Q_b^t U B_{W'} Q_k)\, A_k .
\]
Indeed on edge \(e\) the linked factor on the right is
\[
 P_\lambda(b_e^{-1})^t D_\varrho B D_\varrho P_\lambda(k_e^t)
   =P_\lambda(b_e^{-1})^t D_\varrho^2 P_\lambda(k_e)^t B.
\]
Here \(B\) commutes with \(D_\varrho\). By
Lemma~\ref{lem:one-row-lattice-transpose} this is the
edge term of \(Y_\lambda\) times \(B\). The transposed
restriction sums over precisely the edges entering the
indicated output coordinate, as in
Theorem~\ref{thm:served-actual-hecke-class-matrix}.

Set \(K_0=(k^2+k)(oD)D,\ c=2K_0+2,\ \theta=x^{40}\).
One has \(6c+2\le\theta\) eventually, since \(oD=L\le f\le
x^{12},\ D\le x^{12},\ k=O(x^5)\) by
Proposition~\ref{prop:thick-cyclic-flag-volume} and its proof.
Take \(Z\) of Theorem~\ref{thm:served-integral-approximate-projector}
and extend it to \(Z_{W'}\) on the linked coordinate at every
edge. Apply Lemma~\ref{lem:transverse-pairing-index} with
\[
   Q=Q_k,\qquad \beta=Q_b^t U B_{W'},\qquad
   M=U^{-1} Z_{W'} R_b^t,\qquad
   a_0=p^{K_0},\quad w=\varrho,\quad z=p^s ,
\]
where \(a_0\) is the scalar denoted \(a\) in that lemma.
Each edge slice of each column of \(Q_b\) belongs to
\(\mathcal U(2c,c)\) by Theorem~\ref{thm:served-saturation-excess},
and \(2c\le\theta-4c-2\). Hence the row congruence of
Theorem~\ref{thm:served-integral-approximate-projector} gives
\[
 \beta M=Q_b^t B_{W'}Z_{W'}R_b^t
        \equiv a_0 Q_b^t R_b^t=a_0I\bmod\varrho^c.
\]
Since \(a_0w=\varrho^{c-1}\), this implies the required congruence.

For a difference \(Q_k u-Mv=a_0wF\), each fixed edge slice
on the left lies in \(\mathcal U(\theta,c)\). For \(Q_k\)
use Theorem~\ref{thm:served-saturation-excess} again. For \(M\)
use the image assertion of the projector theorem and the fact
that \(U^{-1}\) makes only integral linear combinations among
the remaining coordinates of a single edge. Thus division by
\(\varrho^{c-1}\) leaves each tested coefficient divisible
by \(\varrho\). Its reduction vanishes, proving that
\(\bar F\) lies in the direct sum of one \(U(\theta)\) per
edge and remaining coordinate. This subspace has dimension
\[
        |C|\,d_{\mathcal P}\,h'\dim U(\theta)
                 \le (1+x^{-14})c_0
\]
by Theorem~\ref{thm:full-served-residue-dimension}. The split
inclusion lemma now bounds \(d_s(\beta Q_k)\) by
\(c_0(K_0+s x^{-14})\), since its exceptional count is at
most \(x^{-14}c_0\).

For the outside square factors the determinants have combined
valuation at most \(3c_0 A_\lambda/(mp)\) by the stronger
inequality proved in Corollary~\ref{cor:served-half-cost-actual}.
Each full determinant bounds its own truncated cost by
Definition~\ref{def:served-truncated-cost}. Apply
Lemma~\ref{lem:product-type} at \(p^s\) to the three factors,
using a generator of the diagonal index for each, to add
their costs. The integral multiplier \(B_V\) gives
\(Y_\lambda B_VV+p^sV\subseteq Y_\lambda V+p^sV\), so the
same bound holds for \(Y_\lambda\). The size estimates above
give \(K_0\le x^{40}\) uniformly.

Finally the entries of \(Y_\lambda\) in the indicated class
basis lie in \(\iota(O_\iota)\) by
Theorem~\ref{thm:served-actual-hecke-class-matrix}.
Perform Smith reduction there. The same invertible changes,
diagonal entries and their divisibilities give a Smith reduction
over \(O\), since the smaller ring is the intersection of
its fraction field with \(O\). Consequently the computed
truncated valuation is exactly the one on the original integral
columns. This proves the assertion in that lattice.
\end{proof}


\section{The local index and divisibility of the transferred columns}

\begin{theorem}\label{thm:all-served-actual-index}
For all sufficiently large arithmetic scales, every served prime
\(p\) and every embedding \(\iota\) above it, put
\[
        s_p=\lceil 2mfa/p\rceil,\qquad
        \mathcal E'_p=Y^m\mathcal E_p+p^{s_p}\mathcal E_p
\]
as in Theorem~\ref{thm:actual-served-hecke-image-gain}. There
is an integer \(l_p\) such that
\[
       \mathcal I(\mathcal E_p,\mathcal E'_p;p^{l_p}),
       \qquad 0\le l_p\le 3c_0fa/p .
\]
\end{theorem}
\begin{proof}
All the \(Y_\lambda\) preserve \(\mathcal E_p\) by
Lemma~\ref{lem:served-scaled-integral-column}.
Apply Lemma~\ref{lem:product-type} at the single modulus
\(p^{s_p}\) to the \(m\) copies of the list of operators
in \(Y\). By Definition~\ref{def:served-truncated-cost},
the sum of their truncated costs bounds that of the whole
product. Theorem~\ref{thm:single-served-actual-index} gives
\[
 \frac{d_{s_p}(Y^m)}{c_0}
 \le \frac3p\sum_{\lambda\in\Lambda} A_\lambda+
              m|\Lambda|\left(x^{40}+\frac{s_p}{x^{14}}\right).
\]
Here \(c_0>0\) since the identity class contributes the
nonzero coefficient representation in
Proposition~\ref{prop:column-dim}.
Theorem~\ref{thm:prime-marking} gives
\(\sum_\lambda A_\lambda=Ua/2,\ m|\Lambda|\le U\le f\).
The right side is thus bounded by
\[
         \frac{3fa}{2p}+Ux^{40}
              +\frac{U}{x^{14}}\left(\frac{2mfa}{p}+1\right).
\]
The last two contributions together are \(o(fa/p)\), uniformly.
Indeed \(a/p\) dominates every fixed power of \(x\) by
Lemma~\ref{lem:diagram-size} and its proof. Also
\(2mU/x^{14}\le4x^{-3/2}\) since \(U\le f\le x^{12}\)
and \(m\le2\sqrt x\).
Eventually \(d_{s_p}(Y^m)\le2c_0fa/p\).
Take \(l_p=\lceil d_{s_p}(Y^m)\rceil\). The cost characterization
gives the claimed index and the rounding costs at most one,
absorbed by \(c_0fa/p\to\infty\).
\end{proof}

\begin{corollary}\label{cor:local-transfer-proved-all-primes}
In Proposition~\ref{prop:exact-normalization} the local
divisibility and index assertions for \(B_\Psi\) hold at
every prime with the modules specified in
Corollary~\ref{cor:local-index-last-served-condition}. They
also hold with the smaller losses
\[
    \widetilde e_p=
      \begin{cases}
       (30\ell+4)fa/p,&p\text{ served},\\
       e_p,&\text{otherwise}.
      \end{cases}
\]
In particular with \(\widetilde\beta_p=\lceil G_p-\widetilde e_p\rceil\),
on these modules \(|(B_\Psi\phi)(h)|_{Y,h}\le
p^{-\widetilde\beta_p}\) for the ramified columns and twists
with \(p\nmid Q\) of that corollary. The exponents of the
local index witnesses are at most \(c_0\widetilde e_p\), and
\(\sum_p\widetilde e_p\log p=o(fa\ell^3)\).
\end{corollary}
\begin{proof}
At a served prime Theorem~\ref{thm:all-served-actual-index}
gives the index with \(3fa/p\) per coordinate. The valuation
cost in Theorem~\ref{thm:actual-served-hecke-image-gain} is
\((30\ell+3)fa/p\). Its difference from \(\widetilde e_p\) is
\(fa/p\ge1\) eventually, so the same estimate implies the
one with the integer ceiling. Also \(3fa/p\le\widetilde e_p\).
Corollary~\ref{cor:local-index-last-served-condition} proves
the two assertions at all other primes. Finally
\(0\le\widetilde e_p\le e_p\) at every prime eventually by
\(\ell\to\infty\) and \(e_p=3fa\ell^2/p\) for served primes.
Proposition~\ref{prop:exact-normalization} gives
\(\sum_p e_p\log p=o(fa\ell^3)\). Decreasing the losses
preserves this upper bound and replacing them by \(e_p\)
again weakens both the local valuation and the index
requirements. This also proves the original assertions.
\end{proof}


\section{Quadratic ideal classes and normalized period lattices}

\begin{lemma}\label{lem:quadratic-inverse-paired-class-reps}
Let \(K=\Q(\sqrt{-d})\) with \(d\) a prime congruent to
\(7\bmod8\), and \(O_K=\Z[\omega]\) as in Lemma~\ref{lem:quadratic-norm}.
Its ideal class group is finite and squaring is a bijection on it.
It has fractional ideal representatives \(I_c\), one per class,
with \(I_1=O_K,\ I_{c^{-1}}=I_c^{-1}\) and \(\prod_c I_c=O_K\).
For a fractional ideal \(I\), define \(NI={\rm area}(I)/
{\rm area}(O_K)\) in the fixed complex embedding. Then
\(I\bar I=(NI)O_K\).
\end{lemma}
\begin{proof}
Use ideal factorization and multiplicativity of the norm from
Proposition~\ref{prop:general-ideal-factorization}, and finiteness
from Lemma~\ref{lem:bounded-ideal-classes}. The area formula for
integral ideals is Lemma~\ref{lem:number-field-lattice}; clearing
a rational integer denominator extends it multiplicatively to
fractional ideals. For an odd prime \(p\), the quotient of \(O_K\)
by \(p\) is the algebra of \(X^2-X+(d+1)/4\) over \(\mathbf F_p\).
It is a quadratic field or a product of two copies of \(\mathbf F_p\)
when \(p\ne d\), by the discriminant \(-d\); at \(p=d\) the
ramified prime is \((\sqrt{-d})\), whose square is \((d)\) and
whose norm is \(d\). At two the polynomial has two distinct roots.
Thus on a split pair conjugation exchanges the primes of norm \(p\),
at an inert prime \((p)\) the norm is \(p^2\), and at the ramified
prime the square and norm are as just computed. Multiplying the
prime factorizations gives \(I\bar I=(NI)O_K\).

If \([I]^2=1\), then \(\bar I=a I\) for some \(a\in K^\times\).
Comparing areas gives \(a\bar a=1\). Choose \(b=1+a\) if
\(a\ne-1\), and \(b=\sqrt{-d}\) otherwise. In both cases
\(b/\bar b=a\), so \(bI\) equals its conjugate. Its two exponents
at a split pair are equal and give a power of the rational
principal ideal \((p)\). All inert primes and the unique ramified
prime are principal. Thus \(bI\) and \(I\) are principal.
The squaring map on the finite abelian class group is now injective
and hence surjective. Its only self-inverse class is the identity.
Take arbitrary representatives from one class in each inverse pair
and the inverse ideals in the other. Their product is as claimed.
\end{proof}

\begin{definition}\label{def:normalized-quadratic-lattice}
For a lattice \(\Gamma\) homothetic to a fractional \(O_K\)-ideal,
let \(t\in(\tfrac12\Gamma)/\Gamma\) be the mixed nonzero class
of Lemma~\ref{lem:half-lattice-cycle}. Call \((\Gamma,t,a_1)\)
normalized if \(\wp_\Gamma''(t)=2,\ a_1^2=12\wp_\Gamma(t)\).
Use \(\sigma_\Gamma,\theta_\Gamma,\xi_\Gamma,\phi_\Gamma,G_\Gamma\)
of Propositions~\ref{prop:reduced-sigma}, \ref{prop:cm-theta},
\ref{prop:scaled-cm} and~\ref{prop:cm-duplication} with this
lattice and half-period. In particular \(\theta_\Gamma(0)=1\),
\(\phi_\Gamma(0)=0,\ \phi_\Gamma'(0)=1\).
As in Proposition~\ref{prop:cm-theta}, changing the lift of \(t\)
multiplies numerator and denominator of the normalized translate
by the same phase after canceling the \(z\)-linear factors.
\end{definition}

\begin{proposition}\label{prop:cm-isogeny-differential-integral}
Every such ideal lattice admits a nonzero homothety making it
normalized. For a normalized lattice, \(a_1\) is an algebraic
integer and \(\Delta_\Gamma=g_2^3-27g_3^2=a_1^4-64\) is an
algebraic unit. If \(a\ne0\) and \(a\Gamma\subseteq\Gamma'\)
for two normalized lattices, \(a\) is an algebraic integer.
\end{proposition}
\begin{proof}
The simple zero of \(\wp'\) at the chosen half-period in the
proof of Proposition~\ref{prop:scaled-cm} gives
\(\wp''(t)\ne0\). Under multiplication of the lattice and \(t\)
by \(z\ne0\), this value changes by \(z^{-4}\), by the defining
series for \(\wp\). Take a fourth root of \(\wp''(t)/2\) and
choose either permitted \(a_1\). With \(e=\wp_\Gamma(t)\) the
differential equation in that proposition gives
\(g_2=12e^2-4,\ g_3=4e^3-g_2e\). Substitution gives
\(\Delta_\Gamma=144e^2-64=a_1^4-64\).
The index-two operation of Lemma~\ref{lem:half-lattice-cycle}
multiplies the lattice by \(O_K+(\omega/2)O_K\).
Some power of this ideal is principal by
Lemma~\ref{lem:quadratic-inverse-paired-class-reps}, so every
initial ideal lattice returns up to homothety. The mixed half-period
is intrinsic and is also carried back by that homothety.
Proposition~\ref{prop:half-recurrence} and Lemma~\ref{lem:cycle-unit}
therefore apply to this cycle. The calculation in
Theorem~\ref{thm:cm-unit} identifies its initial \(D\) with
\(a_1^4-64\), using
\((e_t-e_T)(e_t-e_S)=\wp''(t)/2=1\).
It gives \(|\iota\Delta_\Gamma|_p=1,\ |\iota a_1|_p\le1\)
for every finite place isomorphism. The argument at the end of
Lemma~\ref{lem:cm-section-sign}, applied also to
\(\Delta_\Gamma^{-1}\), gives the algebraic assertions.

Fix a finite place isomorphism and write
\(H_\Gamma(z)=\max(|\iota\sigma_\Gamma(z)|_p,
                      |\iota\theta_\Gamma(z)|_p)\).
The addition and duplication identities of
Proposition~\ref{prop:cm-duplication} use only the simple
half-period, its normalization and the common sigma multiplier,
so hold on \(\Gamma,\Gamma'\). With the unit just proved,
the proof of Theorem~\ref{thm:local-cm-height} gives
\(H_\Gamma(2z)=H_\Gamma(z)^4>0\) and
\(|\iota G_\Gamma(z)|_p\le H_\Gamma(z)^3\) for both curves.
Write \(n=[\Gamma':a\Gamma]=|a|^2{\rm area}(\Gamma)/
{\rm area}(\Gamma')\). For \(f=\sigma_{\Gamma'},\theta_{\Gamma'}\),
the quotients \(f(az)^2/\theta_\Gamma(z)^{2n}\) are periodic.
Indeed the exponential multipliers agree by the area formula and
all signs disappear on squaring. In the proof of
Lemma~\ref{lem:section-polynomials} the pole orders two and three
of \(\xi_\Gamma,\phi_\Gamma\) give the same polynomial expression
and the same bound on \(\Gamma\).
Applying it to both numerators shows
\(H_{\Gamma'}(az)/H_\Gamma(z)^n\) is bounded by a constant
independent of \(z\). Doubling raises this ratio to its fourth
power, so it is at most one everywhere.

Apply Proposition~\ref{prop:bounded-sections-integral} on \(\Gamma\)
to \(\theta_{\Gamma'}(az)^2\) in degree \(2n\) and
\(G_{\Gamma'}(az)^2\) in degree \(6n\). Its proof applies because
the monic quadratic for \(\phi_\Gamma,\xi_\Gamma\), the local
parameter \(\phi_\Gamma=z+O(z^2)\), and the unit discriminant
are precisely those used in that proposition. The periodicity
in the second case follows from \(G_{\Gamma'}=\phi_{\Gamma'}
\theta_{\Gamma'}^3\), and the two height inequalities follow
from the preceding paragraph. After division by the respective
powers of \(\theta_\Gamma\), both series belong to
\(O[[\Phi_\Gamma]]\), where \(\Phi_\Gamma\) is the transported
Taylor series of \(\phi_\Gamma\). The first series has constant
term one, so divide the second by its cube. This proves that
the transported \(\phi_{\Gamma'}(az)^2\) lies in
\(O[[\Phi_\Gamma]]\). Its coefficient at \(\Phi_\Gamma^2\)
is \(\iota(a)^2\), since both local parameters have derivative
one and constant term zero. Hence \(a\) is integral at every
such place, and Lemma~\ref{lem:cm-section-sign} again gives
algebraic integrality.
\end{proof}


\section{Conjugating a normalized elliptic differential}

\begin{proposition}\label{prop:normalized-cm-differential-conjugation}
Let \(C=(\Gamma,t,a_1)\) be as in
Definition~\ref{def:normalized-quadratic-lattice}. On the cubic
with coordinates \(X=\wp_\Gamma,\ Y=\wp'_\Gamma\) mark the origin
at infinity, the point \((a_1^2/12,0)\), and the differential
\(\eta=dX/Y\). For \(\nu\in{\rm Aut}(\C/K)\) there is a normalized
lattice \(C^\nu=(\Gamma^\nu,t^\nu,\nu a_1)\) such that the
conjugate cubic, points, group law and differential are given by
these same formulas on \(C^\nu\). A rational function regular at
the origin has the Taylor series of its conjugated function on
\(C^\nu\) obtained by applying \(\nu\) to the coefficients in
the period variable \(z\). For normalized \(C,D\) one has
\[
       a\Gamma_C\subseteq\Gamma_D
         \quad\Longleftrightarrow\quad
       (\nu a)\Gamma_{C^\nu}\subseteq\Gamma_{D^\nu}.
\]
The induced maps of torsion are conjugated by the same marked
group identifications.
\end{proposition}
\begin{proof}
By Lemma~\ref{lem:elliptic-cubic-period-algebraic} the smooth
cubic, group operations and complex multiplications are algebraic.
In these coordinates the equation is
\(Y^2=4X^3-g_2X-g_3\) and \(\eta\) pulls back to \(dz\), so it
is regular, nowhere vanishing and invariant under translations.
These statements are algebraic: regularity and nonvanishing
are tested in regular local coordinates, and invariance follows
from the identity for the pullback under addition, namely the
sum of the two pullbacks to the product. They hold after applying
\(\nu\) to the coefficients. The new cubic is smooth and
irreducible since \(\nu\Delta_\Gamma\ne0\) and the cubic
polynomial in \(X\) has three distinct roots. The transported
group remains projective and commutative. Use its invariant
differential to identify its tangent line with \(\C\), and
use Lemma~\ref{lem:projective-group-period-uniformization} to
obtain a lattice \(\Gamma^\nu\) with pullback of the conjugate
differential equal to \(dz\). The transported scalar action
of \(O_K\) has the same tangent scalars, because \(\nu\) fixes
\(K\). Hence the lattice is stable under \(O_K\) and is
homothetic to a fractional ideal by the rational-span proof
of Proposition~\ref{prop:cm-elliptic-ray-classification}.

In this variable \(X'=Y\). At the identity \(X\) has a double
pole and is even, since inversion on the cubic is \((X,-Y)\).
Its leading term is \(z^{-2}\) by comparing the leading
coefficients in \(X'^2=4X^3-\nu(g_2)X-\nu(g_3)\).
The coefficient of \(z^{-4}\) in that equation makes its
constant term zero. There are no other poles, so
\(X-\wp_{\Gamma^\nu}\) is entire elliptic with constant term
zero and vanishes by Lemma~\ref{lem:elliptic-count}. This also
identifies \(Y\). The marked point is the mixed half-period
because its order and its two nonzero images under
\(\omega,\bar\omega\) are preserved by conjugation. At that
point \(X=(\nu a_1)^2/12\). Since \(\nu g_2=12(\nu(a_1^2/12))^2-4\),
the second derivative of \(X\) there is two.
This gives exactly the asserted normalized lattice.

For a regular rational function \(f\) near a point \(P\),
define the rational derivative \(\partial f=df/\eta\).
It is regular wherever \(f\) is regular, since \(\eta\) is
nowhere vanishing. In the period variable it is \(d/dz\).
Evaluation at \(P\) and algebraic differentiation commute
with conjugation, which thus carries each
\((\partial^j f)(P)/j!\) to the corresponding coefficient
at the conjugate point. Finally multiplication by \(a\) under
the left inclusion is a holomorphic group homomorphism and
therefore algebraic by Lemma~\ref{lem:elliptic-cubic-period-algebraic}.
Its pullback of the target differential is \(a\) times the
source differential. Conjugation makes the derivative \(\nu a\);
the lifted map in Lemma~\ref{lem:projective-group-period-uniformization}
is multiplication by that derivative. Conversely, undo the
construction by \(\nu^{-1}\). Applying it twice restores the
original group and differential, and the period uniformization
with a fixed differential is unique by the same lemma. This
proves the equivalence and its compatibility with torsion.
\end{proof}

\begin{theorem}\label{thm:normalized-cm-theta-full-germ}
Write \(\Theta_C(Z)\) for the ordinary Taylor series of
\(\theta_\Gamma\) at zero in the period variable and let \(\nu\)
act on series coefficientwise. Then
\[
                          \nu\Theta_C=\Theta_{C^\nu}
                  \qquad(\nu\in{\rm Aut}(\C/K)).
\]
\end{theorem}
\begin{proof}
The logarithmic second derivative
\(-(\theta'_\Gamma/\theta_\Gamma)'\) equals
\(\xi_\Gamma+c_C\) for a constant \(c_C\): its multiplier
gives periodicity after two logarithmic derivatives, and its
double pole at \(t\) has principal part \((z-t)^{-2}\), the
same as \(\xi_\Gamma\). Subtract and use
Lemma~\ref{lem:elliptic-count}, as in
Proposition~\ref{prop:theta-log}. Now
\(\xi_\Gamma=\wp_\Gamma(z+t)-a_1^2/12\) is a rational
function on the cubic by its algebraic translation law.
It transports to \(\xi_{\Gamma^\nu}\) by
Proposition~\ref{prop:normalized-cm-differential-conjugation}.
Both theta series have constant term one and are even by
Proposition~\ref{prop:cm-theta}, whose translation argument
uses only \(2t\in\Gamma\). The preceding identity, applied
to formal logarithms at zero, now shows that
\(\Theta_{C^\nu}/(\nu\Theta_C)=\exp(cZ^2)\) for some \(c\in\C\).
Indeed all coefficients of degree greater than two in its
logarithm vanish, as do its constant and linear terms.

Take \(\alpha=\sqrt{-d}\), with \(\alpha\bar\alpha=d\).
The function \(H_\alpha(z)=\theta_\Gamma(\alpha z)/
\theta_\Gamma(z)^d\) is elliptic. In fact \(\alpha=2\omega-1\)
acts identically on \(\Gamma/2\Gamma\). Both theta signs depend
only on that quotient by the multiplier in
Proposition~\ref{prop:cm-theta}, and \(d\) is odd. The exponential
factors cancel since \(|\alpha|^2=d\). The quotient is rational
on the cubic: it is a polynomial in \(\xi_\Gamma,\phi_\Gamma\)
by subtracting its possible poles at \(t\) as in
Lemma~\ref{lem:section-polynomials}. Its divisor consists of
the inverse images of \(t\) under \([\alpha]\), with
multiplicity one since the derivative is nonzero, less
\(d\) copies of \(t\). Its value at the origin is one.
The conjugate function on \(C^\nu\) has these same descriptions
using \(t^\nu\), by
Proposition~\ref{prop:normalized-cm-differential-conjugation}.
Thus it equals the corresponding \(H_\alpha\) there: the
quotient of two functions with this divisor is a holomorphic
nowhere-zero function on the compact torus, hence is constant,
and the value at zero fixes it. Comparing their Taylor series
now gives \(\exp(c(\alpha^2-d)Z^2)=1\). The coefficient of
\(Z^2\), with \(\alpha^2=-d\ne d\), forces \(c=0\).
\end{proof}


\section{The phases of conjugate torsion translates}

\begin{lemma}\label{lem:torsion-theta-exact-phase}
Let \(C,D=C^\nu\) be as in
Proposition~\ref{prop:normalized-cm-differential-conjugation}.
Write \([z]_C\) for the point represented by \(z\) modulo the
period lattice, and \(U_y\) for the translate in
Proposition~\ref{prop:cm-theta} with the area of its own lattice.
Let \(n\ge1,\ ny\in\Gamma_C,\ ny'\in\Gamma_D\). If
\(\nu[-y/(2n)]_C=[-y'/(2n)]_D\), then the Taylor series at
zero satisfy
\[
       \nu\,{\rm Tay}_0(U_y\theta_C)
                      ={\rm Tay}_0(U_{y'}\theta_D).
\]
\end{lemma}
\begin{proof}
Put \(\epsilon_C(l)=s(l)\exp(2\pi i\Im(\bar l t)/A_C)\)
for \(l\in\Gamma_C\), with \(s(l)\) as in
Proposition~\ref{prop:reduced-sigma}. Since \(2t\in\Gamma_C\),
this sign depends only on \(l\bmod 2\Gamma_C\) and
\(\epsilon_C(nl)=\epsilon_C(l)^{n^2}\). Define
\[
 H_n(z)=\frac{\theta_C(nz)}{\theta_C(z)^{n^2}},\qquad
 P_{y,n}(z)=\frac{(U_y\theta_C)(nz)}{\theta_C(z)^{n^2}} .
\]
Their numerators and denominators have the same multiplier
under \(z\mapsto z+l\). For \(P_{y,n}\) the extra factor is
\(\exp(2\pi i\Im(\bar l ny)/A_C)=1\), since the oriented
area of two lattice vectors is an integral multiple of \(A_C\).
Both quotients are rational on \(C\) by the pole subtraction
in Lemma~\ref{lem:section-polynomials}. Denote the analogous
functions on \(D\) by primes. The divisor of \(H_n\) is
\([n]^*[t]-n^2[t]\), where \([n]^*\) sums the preimages of
the indicated point. All preimages are simple since \([n]\)
has nonzero derivative. This divisor transports to that of \(H_n'\).
The quotient of the two functions with that divisor is constant
by Lemma~\ref{lem:elliptic-count}, and their values at zero
are one. Thus \(\nu H_n=H_n'\).

Likewise the divisor of \(P_{y,n}\) is
\([n]^*[t-y]-n^2[t]\). By the hypothesis, multiplication by
\(-2n\) on the two curves identifies \(\nu[y]_C=[y']_D\).
Thus \(\nu P_{y,n}\) and \(P_{y',n}'\) have identical divisors.
If \(y\notin\Gamma_C\), put \(z_0=-y/(2n)\). The denominator
is nonzero there, since \([z_0]_C=t\) would give \([y]_C=0\).
The value of \(\theta_C(nz_0)=\theta_C(-y/2)\) is also nonzero
for the same reason. As \(nz_0=-y/2\), the exponent defining
\(U_y\) vanishes there and evenness gives
\(P_{y,n}(z_0)=H_n(z_0)\). These facts hold at \(z_0'=-y'/(2n)\)
on \(D\) too. The conjugate point and both values of \(H_n\)
correspond, proving \(\nu P_{y,n}=P_{y',n}'\).

If \(y\in\Gamma_C\), the translate is
\(\epsilon_C(y)\theta_C\). The sign is negative exactly when
\([y/2]_C=t\). Indeed \(s(y)=1\) on the zero class modulo
\(2\Gamma_C\) and is negative on each of its three nonzero
classes. The alternating area pairing modulo two is nonzero
precisely on distinct nonzero vectors, as is seen in an oriented
lattice basis. This proves the sign rule by its defining
formula. Conjugation carries \([y/2]_C\) to \([y'/2]_D\),
and \(t\) to \(t^\nu\). Hence the two signs agree, and again
\(\nu P_{y,n}=P_{y',n}'\). Take Taylor coefficients at zero
by Proposition~\ref{prop:normalized-cm-differential-conjugation}
and multiply by \({\rm Tay}_0(\theta_D)^{n^2}\) using
Theorem~\ref{thm:normalized-cm-theta-full-germ}. Substituting
\(Z/n\) for \(Z\) proves the claim.
\end{proof}

\begin{theorem}\label{thm:finite-multiplier-for-conjugate-theta}
Suppose \(\Gamma_C=\lambda O_K\), and \(\nu\in{\rm Aut}(\C/K)\).
There exist \(\lambda_\nu\in\C^\times\), a fractional ideal
\(I\) of norm one, and an invertible finite adele \(b_\nu\) of
\(K\) such that
\[
 \Gamma_{C^\nu}=\lambda_\nu I,\qquad
 |\lambda_\nu|^2=2\,{\rm area}(\Gamma_{C^\nu})/\sqrt d,\qquad
 b_\nu\widehat O_K=\widehat I.
\]
Here \(\widehat I=\varprojlim_m I/mI\), with \(m\) running
through positive integers under divisibility. In particular
\(e\lambda_\nu\in\Gamma_{C^\nu}\) for some positive integer \(e\).
For every finite \(S\subset K\), every integer \(n\ge1\) with
\(nS\subset O_K\), and every \(a\in I\) congruent to \(b_\nu\)
modulo \(2n^2\widehat I\), one has
\[
  \nu\,{\rm Tay}_0(U_{\lambda x}\theta_C)
      ={\rm Tay}_0(U_{\lambda_\nu a x}\theta_{C^\nu})
                    \qquad (x\in S).
\]
Such \(a\) can at the same time be chosen in any prescribed
finite-adelic neighbourhood of \(b_\nu\).
\end{theorem}
\begin{proof}
By Proposition~\ref{prop:normalized-cm-differential-conjugation}
the new lattice is homothetic to a fractional ideal. By
Lemma~\ref{lem:quadratic-inverse-paired-class-reps} its class
is a square. Write \(\Gamma_{C^\nu}=vJ^2\) with \(J\) a
fractional ideal and \(v\ne0\). Take
\(\lambda_\nu=vNJ,\ I=J^2/NJ=J/\bar J\). Then \(NI=1\).
The area identity follows from \({\rm area}(O_K)=\sqrt d/2\)
in the basis \(1,\omega\). Clearing denominators in a basis
of \(I\) gives an integer \(e\in I\), for \(I\) spans \(K\)
over \(\Q\).

For each integer \(m\ge1\), identify \(I/mI\) with the
\(m\)-torsion of \(C^\nu\) by \(v_0\mapsto
[\lambda_\nu v_0/m]_{C^\nu}\), and similarly identify \(O_K/mO_K\)
with that of \(C\). Algebraicity of multiplication and the
tangent scalar computation of
Proposition~\ref{prop:normalized-cm-differential-conjugation}
make conjugation an \(O_K\)-linear isomorphism between these
two finite modules. It sends \(1\) to a residue \(b_m\).
For \(m\mid M\), multiplying the points by \(M/m\) gives the
reduction of \(b_M\) to \(b_m\). Hence \(b=(b_m)_m\in\widehat I\)
and multiplication by \(b\) gives an isomorphism
\(\widehat O_K\to\widehat I\). Indeed the finite-level inverse
maps commute with the same reductions. The Chinese remainder
calculation of Proposition~\ref{prop:general-ideal-factorization}
identifies these completions with the products of the local
integer rings and the completed ideals, respectively: factor
each \(mO_K\) into prime powers, split its quotient by these
powers, and pass to the inverse limit. Each local component
of \(b\) generates the corresponding nonzero local ideal, and
outside the finitely many primes dividing \(I\) it is a unit.
Thus \(b_\nu=b\) is a finite idele.

Choose \(a\) in the residue \(b_{2n^2}\). The transported
point of \(\lambda/(2n^2)\) is then represented by
\(\lambda_\nu a/(2n^2)\). Multiplication by \(-nx\in O_K\)
gives the hypothesis of Lemma~\ref{lem:torsion-theta-exact-phase}
for \(y=\lambda x,\ y'=\lambda_\nu ax\). This proves the
formula simultaneously for all \(x\). A neighbourhood of
\(b_\nu\) contains \(b_\nu+m\widehat I\) for some \(m\):
only finitely many local congruences need tightening. Choose
a representative in \(I\) of \(b_{2mn^2}\) to meet both conditions.
\end{proof}


\section{The absolute discriminant of a period lattice}

\begin{lemma}\label{lem:reduced-sigma-infinite-product}
Let \(\Gamma=\lambda(\Z+\tau\Z)\) with \(y=\Im\tau>0\),
\(A=|\lambda|^2y\) and \(Q=\exp(2\pi i\tau)\). Then
\[
 \sigma_\Gamma(\lambda z)=\frac{\lambda}{\pi}
   e^{\pi z^2/(2y)}\sin(\pi z)
   \prod_{j\ge1}
       \frac{(1-Q^j e^{2\pi iz})(1-Q^j e^{-2\pi iz})}{(1-Q^j)^2}.
\]
The function \(S(w)=|\sigma_\Gamma(w)|\exp(-\pi |w|^2/(2A))\)
is \(\Gamma\)-periodic. In the modulus of the preceding product
the net exponential for \(S(\lambda z)\) is
\(\exp(-\pi(\Im z)^2/y)\).
\end{lemma}
\begin{proof}
On each compact set \(e^{\pm2\pi iz}\) are bounded, so the
factors of the product differ from one by \(O(|Q|^j)\) uniformly.
For all sufficiently large \(j\) the logarithms obtained by
the power series of \(\log(1+w)\) converge uniformly and
absolutely when summed over \(j\). Exponentiation gives a
locally uniform product and a nonvanishing tail. The disk
formula of Lemma~\ref{lem:analytic-tools} makes the result
holomorphic. The product is even, and has value one at zero.
The displayed expression is odd with derivative one in \(w=\lambda z\).

Put \(W=e^{2\pi iz}\), and write \(Q^{1/2}=e^{\pi i\tau}\).
For translation by \(1\), the sine changes sign and the product
stays fixed. Under \(z\mapsto z+\tau\) the ratio of the sines is
\(Q^{-1/2}(QW-1)/(W-1)\). The ratio of the products is
\((1-W^{-1})/(1-QW)\), by cancellation in the finite products
followed by a limit. Multiplication gives \(-Q^{-1/2}W^{-1}\).
The identities extend through zeros by continuity.
With the Gaussian these are exactly the two generating
multipliers in Proposition~\ref{prop:reduced-sigma}, as
\(\tau-2iy=\bar\tau\). Its uniqueness proves the formula.
Taking absolute values in that proposition cancels the Gaussian
in \(S\) under every period. Finally
\(\Re(z^2)-|z|^2=-2(\Im z)^2\), giving the stated exponential.
\end{proof}

\begin{proposition}\label{prop:period-discriminant-product}
In Lemma~\ref{lem:reduced-sigma-infinite-product} set
\(\Delta_\Gamma=g_2^3-27g_3^2\), with \(g_2,g_3\) from
the differential equation of Proposition~\ref{prop:scaled-cm}.
Then
\[
       |\Delta_\Gamma|=(2\pi/|\lambda|)^{12}
                          |Q|\prod_{j\ge1}|1-Q^j|^{24}.
\]
\end{proposition}
\begin{proof}
Let \(u,v,w\) represent the three distinct nonzero half periods
and put \(e_h=\wp_\Gamma(h)\). They are the three distinct
roots of \(4X^3-g_2X-g_3\): the zeros of \(\wp'\) at the
three half periods are simple by the proof of
Proposition~\ref{prop:scaled-cm}, and an equal value at two
of them would give four zeros counted with multiplicities
for an elliptic function \(\wp-e_h\) with just one double pole.
The addition identity in the proof of
Proposition~\ref{prop:cm-duplication} used only the sigma
multiplier, its simple zero and oddness. Taking moduli in it,
the Gaussians cancel by the parallelogram identity. Since
\(u+v\) and \(u-v\) both represent \(w\), this gives
\[
       |e_u-e_v|=\frac{S(w)^2}{S(u)^2 S(v)^2}.
\]
The discriminant identity
\(\Delta_\Gamma=16\prod_{h<k}(e_h-e_k)^2\), with any ordering
of the roots, follows by expansion: for the monic polynomial
\(X^3+pX+r\) the squared product of differences is
\(-4p^3-27r^2\), on substituting \(e_w=-e_u-e_v\).
Thus \(|\Delta_\Gamma|=16/(S(u)S(v)S(w))^4\).

Use \(u=\lambda/2,\ v=\lambda\tau/2,\
w=\lambda(1+\tau)/2\). In the product of the three formulas
of Lemma~\ref{lem:reduced-sigma-infinite-product}, the three
sines and Gaussians together have modulus
\(|Q|^{-1/4}|1-Q|/4\). Indeed the first sine is one. For
the other two write the sine as
\(e^{-\pi iz}(e^{2\pi iz}-1)/(2i)\). Their two arguments
\(e^{2\pi iz}\) are \(Q^{1/2},-Q^{1/2}\), and their Gaussian
factor together is \(|Q|^{1/4}\).
The numerators of the infinite products give
\[
 \prod_{j\ge1}
       |1+Q^j|^2\,|1-Q^{2j+1}|\,|1-Q^{2j-1}|.
\]
Multiplying this by \(|1-Q|\) gives one, since
\(\prod_{j\ge1}(1+Q^j)
  =1/\prod_{j\ge1}(1-Q^{2j-1})\) by factoring \(1-Q^{2j}\).
All rearrangements are permitted by the absolute summability
of logarithms of the tails. We obtain
\[
 S(u)S(v)S(w)
   =\frac{(|\lambda|/\pi)^3 |Q|^{-1/4}}
            {4\prod_{j\ge1}|1-Q^j|^6}.
\]
Substitution in the discriminant gives the assertion.
\end{proof}


\section{A first limit formula for lattice sums}

\begin{lemma}\label{lem:lattice-row-finite-part}
Let \(\tau=u+iy\), \(y>0\), and \(Q=e^{2\pi i\tau}\).
For \(n\ge1,\ 1\le s\le3/2\) put
\[
 J(s)=\int_{\R}(1+x^2)^{-s}\dd x,\qquad
 D_n(s)=\sum_{m\in\Z}((m+nu)^2+(ny)^2)^{-s}
                         -(ny)^{1-2s}J(s).
\]
Then \(\sum_{n\ge1}D_n(s)\) converges uniformly on this interval
and is continuous there. Its value at \(s=1\) is
\[
                      -\frac{2\pi}{y}\sum_{j\ge1}\log|1-Q^j|.
\]
Also \(J(1)=\pi,\ J'(1)=-2\pi\log 2\).
\end{lemma}
\begin{proof}
Write \(b=ny\) and \(f(x)=(x^2+b^2)^{-s}\). On
\([t,t+1]\) the absolute difference between \(f(t)\) and the
integral of \(f\) is at most \(\int_t^{t+1}|f'(x)|\dd x\)
by the fundamental theorem of calculus. Sum over the consecutive
intervals starting at \(t=m+nu\). The total variation is
\(2b^{-2s}\), because \(f\) increases to \(f(0)\) and decreases
afterwards. The integral is \(b^{1-2s}J(s)\) by substitution.
For \(ny\ge1\) the variation is at most \(2/(ny)^2\). For each
fixed \(n\) the series over \(m\) has a uniform summable tail
bounded by a constant times \(\sum_{m\ne0}m^{-2}\). These
estimates and Lemma~\ref{lem:integral-tools} prove continuity
and uniform convergence.

For \(b>0\), the periodization of \(1/(x^2+b^2)\) has Fourier
coefficients \(\pi e^{-2\pi b|j|}/b\) at \(j\in\Z\).
Indeed absolute uniform convergence on a period allows unfolding.
For \(t\ge0\) the integral of \(e^{itx}/(x^2+b^2)\) over
the real axis is \(\pi e^{-bt}/b\). On closing the contour
in the upper half-plane, the pole at \(ib\) has residue
\(e^{-bt}/(2ib)\) and the semicircle of radius \(R>b\) costs
at most \(\pi R/(R^2-b^2)\). Use the contour formula of
Lemma~\ref{lem:analytic-tools} and let \(R\to\infty\).
Complex conjugation gives the integral for negative \(t\).
The Fej\'er kernel argument in Lemma~\ref{lem:theta} now
identifies the continuous periodization with its absolutely
summable Fourier series. Evaluation at \(nu\) and subtraction
of the constant term give
\[
            D_n(1)=\frac{2\pi}{ny}
                      \sum_{j\ge1}\Re(Q^{nj}).
\]
The double series with weights \(|Q|^{nj}/n\) is bounded by
\(\sum_n |Q|^n/(n(1-|Q|))\). Summing first over \(n\) gives
the asserted logarithms by the power series of \(-\log(1-z)\).

Differentiation of \(J\) near one is justified by the integrable
majorant \((1+x^2)^{-3/4}(1+\log(1+x^2))\).
Set \(x=\tan v\). This gives \(J(1)=\pi\) and
\(J'(1)=4\int_0^{\pi/2}\log(\cos v)\dd v\). The last integral
equals \(I=\int_0^{\pi/2}\log(\sin v)\dd v\) by reflection.
Both converge since \(\sin v/v\to1\) at zero.
Adding the two integrals and using
\(\sin(2v)=2\sin v\cos v\) gives
\(2I=I-(\pi/2)\log2\) by changing \(2v\) to the integration
variable and reflecting about \(\pi/2\). This proves the derivative.
\end{proof}

\begin{theorem}\label{thm:real-kronecker-lattice-limit}
Let \(\Gamma\) be a rank-two complex lattice of area \(A\),
and \(k>0\). For real \(s>1\) put
\[
       Z_\Gamma(s;k)=\sum_{0\ne l\in\Gamma}
                                  (|l|^2/(kA))^{-s}.
\]
The sum converges. With \(\gamma\) as in
Definition~\ref{def:heights}, as \(s\downarrow1\) it has the expansion
\[
 Z_\Gamma(s;k)=\frac{\pi k}{s-1}
  +\pi k\left(2\gamma-\log\frac{A}{k\pi^2}
                         -\frac16\log|\Delta_\Gamma|\right)+o(1).
\]
\end{theorem}
\begin{proof}
Take a positively oriented lattice basis \(\lambda,\lambda\tau\)
and the notation of Lemma~\ref{lem:lattice-row-finite-part}.
The comparison with the maximum norm in the proof of
Proposition~\ref{prop:scaled-cm} bounds the tail, for fixed
\(s>1\), by a constant times \(\sum_{j\ge1}j^{1-2s}\).
Separating \(n=0,\ n>0,\ n<0\) in \(l=\lambda(m+n\tau)\) gives
for \(1<s\le3/2\)
\[
  Z_\Gamma(s;k)=2(ky)^s
       \left(\zeta(2s)+\sum_{n\ge1}D_n(s)
                   +y^{1-2s}J(s)\zeta(2s-1)\right).
\]
For use at \(2s=2\), the Fourier cosine coefficients of the
continuous \(2\pi\)-periodic extension of \(x^2\) on
\([-\pi,\pi]\) are \(4(-1)^j/j^2\) for \(j\ge1\), by two
integrations by parts. The constant term is \(\pi^2/3\).
The same Fej\'er kernel argument as in Lemma~\ref{lem:theta}
recovers the function from its absolutely summable series.
Setting \(x=\pi\) gives \(\zeta(2)=\pi^2/6\).
By Proposition~\ref{prop:cont} and the computation of the
regularized derivative in Proposition~\ref{prop:height-cost},
\(\zeta(1+\delta)=1/\delta+\gamma+o(1)\) for \(\delta\downarrow0\).

Put \(h=s-1\) and expand the last summand including its prefactor.
Since \(k^s y^{1-s}=k(1+h\log(k/y)+O(h^2))\),
Lemma~\ref{lem:lattice-row-finite-part} gives its contribution
divided by \(\pi k\) as
\(h^{-1}+2\gamma+\log(k/(4y))+o(1)\).
For the first two summands the contributions divided by
\(\pi k\) tend respectively to \(\pi y/3=-(\log|Q|)/6\)
and \(-4\sum_{j\ge1}\log|1-Q^j|\). Finally
\(A=|\lambda|^2 y\), and
Proposition~\ref{prop:period-discriminant-product} gives
\[
 \log\frac{k}{4y}-\frac16\log|Q|
          -4\sum_{j\ge1}\log|1-Q^j|
       =-\log\frac{A}{k\pi^2}-\frac16\log|\Delta_\Gamma|.
\]
This is the desired constant term.
\end{proof}


\section{A common period correction over the quadratic field}

\begin{proposition}\label{prop:ideal-class-average-lattice-limit}
Let \(K=\Q(\sqrt{-d})\) as in
Lemma~\ref{lem:quadratic-inverse-paired-class-reps},
let \(h=|{\rm Cl}(O_K)|,\ \eta_d(m)=(m/d)\), and
\(k=2/\sqrt d\). For each ideal class choose one complex
lattice \(\Gamma_c\) homothetic to a fractional ideal of
that class, and let \(A_c\) be its area. Then
\[
 \left(\prod_c A_c|\Delta_{\Gamma_c}|^{1/6}\right)^{1/h}
      =\frac{2\pi^2}{\sqrt d}
            \exp\!\left(\gamma-\frac{L'(1,\eta_d)}{L(1,\eta_d)}\right).
\]
\end{proposition}
\begin{proof}
Write \(\Gamma_c=v_cJ_c\). For \(0\ne b\in J_c\),
\[
       \frac{|v_cb|^2}{kA_c}=\frac{b\bar b}{NJ_c}
                               =N\big((b)J_c^{-1}\big),
\]
because \({\rm area}(O_K)=1/k\).
The ideal on the right is integral and has class \(c^{-1}\).
Conversely, for each integral ideal \(B\ne0\) in that class,
\(BJ_c\) is principal and contained in \(J_c\). Its nonzero
generators differ by precisely the units of \(O_K\).
There are finitely many units, since their norms equal one
and the lattice \(O_K\) has only finitely many points in the
closed unit disk. Put \(w=|O_K^\times|\).
Thus for real \(s>1\),
\[
       \sum_c Z_{\Gamma_c}(s;k)
              =w\sum_{0\ne B\subset O_K}(NB)^{-s}
              =w\zeta(s)L(s,\eta_d),
\]
where the middle sum is over integral ideals. Indeed by
Proposition~\ref{prop:general-ideal-factorization} ideals
are unique products of prime ideals. At rational \(p\) the
local factors are \((1-p^{-s})^{-2}\) for a split prime,
\((1-p^{-2s})^{-1}\) for an inert prime, and
\((1-p^{-s})^{-1}\) for the ramified prime \(d\).
The splitting calculation of
Lemma~\ref{lem:quadratic-inverse-paired-class-reps} and
Lemma~\ref{lem:reciprocity} identify the three cases with
\(\eta_d(p)=1,-1,0\), including \(p=2\) in the first case.
All products converge absolutely since the positive
ideal factors are at most those of \(\zeta(s)^2\).
This also justifies the multiplication of the two Dirichlet
Euler products by finite factorization and passage to the limit.

The character \(\eta_d\) is nonprincipal since exactly half
the nonzero residues are squares, as shown in
Lemma~\ref{lem:reciprocity}. Proposition~\ref{prop:cont}
therefore gives its Taylor expansion at one. Compare the
coefficient of \(1/(s-1)\) in the preceding identity by
Theorem~\ref{thm:real-kronecker-lattice-limit}. This gives
\(h\pi k=wL(1,\eta_d)>0\). Comparing the constants, using
the expansion of \(\zeta\) from the same theorem's proof,
and dividing by \(h\pi k\), gives
\[
  2\gamma+\log(k\pi^2)
       -\frac1h\sum_c\log(A_c|\Delta_{\Gamma_c}|^{1/6})
                 =\gamma+\frac{L'(1,\eta_d)}{L(1,\eta_d)}.
\]
Exponentiation proves the formula.
\end{proof}

\begin{theorem}\label{thm:common-unit-conjugate-cm-area}
Let \(C\) be a normalized lattice with \(\Gamma=\lambda O_K\).
There exists an algebraic unit \(U\ne0\), depending on \(C\),
such that for every \(\nu\in{\rm Aut}(\C/K)\),
\[
 {\rm area}(\Gamma_{C^\nu})|\nu U|^2
      =\frac{2\pi^2}{\sqrt d}
              \exp\!\left(\gamma-\frac{L'(1,\eta_d)}{L(1,\eta_d)}\right).
\]
\end{theorem}
\begin{proof}
Take the inverse-paired fractional ideals \(I_c\) of
Lemma~\ref{lem:quadratic-inverse-paired-class-reps}.
Normalize the lattice \(I_c\Gamma\) by a complex factor
\(u_c\ne0\) using Proposition~\ref{prop:cm-isogeny-differential-integral}.
Call this curve \(C_c\), so \(\Gamma_c=u_c I_c\Gamma\),
where multiplication by a fractional ideal means the sums
of products with lattice vectors. For all \(z\in I_c\)
and \(b\in I_c^{-1}\) one has
\((u_c z)\Gamma\subseteq\Gamma_c\) and
\((u_c^{-1}b)\Gamma_c\subseteq\Gamma\).
Thus \(u_c z,u_c^{-1}b\) are algebraic integers by that
proposition, with zero included. Taking a nonzero rational
integer \(z\in I_c\) proves \(u_c\) algebraic.
Set \(u_*=\prod_c u_c\). Since \(\prod_c I_c=O_K\), there is
an expression \(1=\sum_j\prod_c z_{cj}\), with
\(z_{cj}\in I_c\), by the definition of ideal multiplication.
Multiply by \(u_*\) to express it as a sum of products of
algebraic integers. The same argument for \(\prod_c I_c^{-1}=O_K\)
makes \(u_*^{-1}\) integral. The ring property follows by
the finite-module argument in the proof of
Theorem~\ref{thm:frame-sections}, which applies to sums as
well as products. Thus \(u_*\) is a unit.

Choose any \(12h\)-th root
\[
                 U^{12h}=u_*^{12}\prod_c \Delta_{\Gamma_c}.
\]
All factors on the right are algebraic units by
Proposition~\ref{prop:cm-isogeny-differential-integral}.
If a power of a complex number is an algebraic integer,
substitution of that power into a monic integer polynomial
gives a monic integer polynomial for the number itself.
Applying this also to \(U^{-1}\) shows \(U\) is an algebraic unit.

Fix \(\nu\) as in the statement and write \(\Gamma^\nu=\Gamma_{C^\nu}\).
By Proposition~\ref{prop:normalized-cm-differential-conjugation},
the two inclusions for \(u_c z,u_c^{-1}b\) become the corresponding
inclusions for \((\nu u_c)z,(\nu u_c)^{-1}b\) on the conjugate
curves. Summing over \(z\) gives
\((\nu u_c)I_c\Gamma^\nu\subseteq\Gamma_{C_c^\nu}\).
The reverse follows from the second family by multiplying
by \((\nu u_c)I_c\) and using \(I_c I_c^{-1}=O_K\).
Consequently the lattices \(\Gamma_{C_c^\nu}\) run once
through all ideal classes up to homothety: multiplication
by the fixed class of \(\Gamma^\nu\) permutes the classes.
Their areas are \(|\nu u_c|^2 NI_c\,{\rm area}(\Gamma^\nu)\).
Moreover \(\Delta_{\Gamma_{C_c^\nu}}=\nu\Delta_{\Gamma_c}\)
by the conjugation of the smooth cubic. Multiply the areas
times the sixth roots of the absolute discriminants and take
their geometric mean. Since \(\prod_c NI_c=1\), the result is
\({\rm area}(\Gamma^\nu)|\nu U|^2\) by the defining equation
for \(U\). Apply Proposition~\ref{prop:ideal-class-average-lattice-limit}
to these conjugate lattices.
\end{proof}

\begin{corollary}\label{cor:absolute-conjugate-jet-scale}
In Definition~\ref{def:standard-local-test-module} one may take
\(s_1=s_0/U\) simultaneously at every place. For the period
\(\lambda_\nu\) of Theorem~\ref{thm:finite-multiplier-for-conjugate-theta}
put \(t_\nu=\nu(s_1)\sigma_0/\lambda_\nu\). Then
\[
                         |t_\nu|^2=\Omega^{-1}.
\]
\end{corollary}
\begin{proof}
Every finite place gives \(U\) and \(U^{-1}\) absolute values
at most one by their monic equations. Hence
\(|\iota(s_1/s_0)|_p=1\), as required in
Definition~\ref{def:standard-local-test-module}.
Using \(\sigma_0^2=2\pi/\sqrt d\), the period-area identity
of Theorem~\ref{thm:finite-multiplier-for-conjugate-theta}
gives \(|t_\nu|^2=\pi|\nu s_0|^2/
({\rm area}(\Gamma_{C^\nu})|\nu U|^2)\).
Theorem~\ref{thm:global-frame-scale} gives
\(|\nu s_0|^2={\rm rd}_E/\sqrt d\).
Now substitute Theorem~\ref{thm:common-unit-conjugate-cm-area}
and Definition~\ref{def:heights}.
\end{proof}


\section{Conjugate jets of the cyclotomic frame}

\begin{definition}\label{def:renamed-conjugate-taylor}
For the embeddings \(G={\rm Hom}_K(E,\C)\) of
Proposition~\ref{prop:places} and \(\nu\in{\rm Aut}(\C/K)\),
write \(\mathcal C_\nu\) for coefficientwise application of
\(\nu\) to a series followed by the variable permutation
\(Z_{\tau k}\mapsto Z_{\nu\circ\tau,k}\). This is an invertible
additive and multiplicative map on formal series, semilinear
over \(\nu\). For \(a\in\C^\times\) and \(S\in\mathcal F\) write
\(J_a S={\rm Tay}_0(S(aZ))\) in these variables, as in
Definition~\ref{def:standard-local-test-module}.
\end{definition}

\begin{theorem}\label{thm:cyclic-conjugate-jet-generator}
Take \(D=N_1,\ w_j=\zeta_D^j\) for \(0\le j<D\), and
\(u=\sigma_0s_1/\lambda\) with \(s_1=s_0/U\) from
Corollary~\ref{cor:absolute-conjugate-jet-scale}. For every
standard positive rank \(n\ge1\) there is an \(F\in\mathcal F\)
with \(F(0)=1\) whose finite translates span \(\mathcal F\).
For every \(\nu\in{\rm Aut}(\C/K)\) there is \(F_\nu\in\mathcal F\)
with \(F_\nu(0)=1\), also cyclic, such that for all \(X\in V_f\),
\[
     \mathcal C_\nu(J_u T_XF)=J_{t_\nu} T_{b_\nu X}F_\nu,
                \qquad |t_\nu|^2=\Omega^{-1}.
\]
The scalar \(b_\nu\) in Theorem~\ref{thm:finite-multiplier-for-conjugate-theta}
acts through the diagonal inclusion of the finite adeles of \(K\)
in those of \(E\), on all coordinates. It is independent of \(n\).
\end{theorem}
\begin{proof}
Use the frame of Proposition~\ref{prop:integral-vacuum-character-sum}
in each of the \(n\) coordinates. For either
\((C_*,\lambda_*)=(C,\lambda)\) or
\((C_*,\lambda_*)=(C^\nu,\lambda_\nu)\) put
\[
 B^*_{jk}(Z)=\frac{\lambda_*}{D\sigma_0}
    \sum_{\tau\in G}\overline{\tau(w_j)}Z_{\tau k},\qquad
 F_*(Z)=\prod_{k=1}^n\prod_{j=0}^{D-1}
                           \theta_{C_*}(B^*_{jk}(Z))^D.
\]
In the first case this is the section \(F\) of
Theorem~\ref{thm:frame-sections}, with \(\delta=D,\ \alpha=1\).
In the second case use
\({\rm area}(\Gamma_{C^\nu})=|\lambda_\nu|^2\sqrt d/2\).
The frame identity
\(\sum_j D\,\tau(w_j)\overline{\tau'(w_j)}
       =D^2{\bf1}_{\tau=\tau'}\) gives in both cases
\[
 \frac{\pi}{{\rm area}(\Gamma_{C_*})}
       \sum_{jk}D\,\overline{B^*_{jk}(v)}B^*_{jk}(z)=H(v,z).
\]
Multiplying the elementary exponential factors in
Proposition~\ref{prop:cm-theta} thus expresses each \(W_vF_*\)
as a product with every theta factor translated by \(B^*_{jk}(v)\).
Both functions have value one at zero.

Choose \(e\ge1\) with \(e\lambda_\nu\in\Gamma_{C^\nu}\) as in
Theorem~\ref{thm:finite-multiplier-for-conjugate-theta}. For
\(h\in 2De\,\mathcal A^n\), the displacement
\(B^\nu_{jk}(rh)=(\lambda_\nu/D){\rm Tr}_{E/K}(h_k\bar w_j)\)
belongs to \(2\Gamma_{C^\nu}\) by
Proposition~\ref{prop:frame-lattice-trace} and \(O_K\)-stability.
Here \(B^\nu\) denotes the second choice of the displayed forms.
All theta signs are one on twice a period by
Lemma~\ref{lem:torsion-theta-exact-phase}. The prescribed sign
\(\chi_{\Gamma_0}(rh)\) is also one, since \(h\) is twice an
integral frame vector and its norm trace is a multiple of four.
Thus \(F_\nu=F_*\) in the second case belongs to
\(\mathcal T(2De\Gamma_0)\).
For either \(F_*\), choose any level divisible by its own.
It is a nonzero vector at that level, and its translates
by the dual lattice span the whole level by
Proposition~\ref{prop:theta-shift-irrep}. These are finite
translations by Theorem~\ref{thm:finite-frame-action} since
the dual lattice is rational in the coordinates of
Proposition~\ref{prop:frame-polarization}. The eligible levels
are cofinal in Definition~\ref{def:finite-frame-space}.
This proves both cyclicity assertions.

First take \(X=\gamma\in E^n\), so \(T_\gamma=W_{-r\gamma}\).
Put \(x_{jk}=-{\rm Tr}_{E/K}(\gamma_k\bar w_j)/D\in K\).
The source translation has one-variable shifts \(\lambda x_{jk}\).
By local constancy of \(T_YF_\nu\) in
Theorem~\ref{thm:finite-frame-action}, all finite scalars
sufficiently close to \(b_\nu\) act on \(\gamma\) with the
same resulting translate of \(F_\nu\). Multiplication by
\(\gamma\) is continuous coordinatewise in the finite adeles.
Use Theorem~\ref{thm:finite-multiplier-for-conjugate-theta}
for the finite set of \(x_{jk}\), choosing \(a\in I\subset K\)
in that neighbourhood. The target action equals \(T_{a\gamma}\)
on \(F_\nu\) and has shifts \(\lambda_\nu a x_{jk}\).
Thus the translated factors have exactly conjugate
one-variable Taylor series. Under \(\mathcal C_\nu\) the
linear form \(B_{jk}(uZ)\) becomes \(B^\nu_{jk}(t_\nu Z)\):
the coefficient before renaming is
\(\nu(s_1)\overline{(\nu\circ\tau)(w_j)}/D\).
Here \(\overline{w_j}=w_j^{-1}\), and
\(t_\nu=\sigma_0\nu(s_1)/\lambda_\nu\).
Finite products and substitution of linear forms with zero
constant term commute with coefficientwise conjugation in
each fixed degree. This proves the identity for rational
\(\gamma\). Both sides are locally constant functions of
\(X\) into the series ring taken as a set, since the
sections themselves are locally constant. Rational points
are dense by Theorem~\ref{thm:finite-frame-action}, which
proves the identity everywhere. The modulus of \(t_\nu\)
is Corollary~\ref{cor:absolute-conjugate-jet-scale}.
\end{proof}


\section{A semilinear bijection on the finite oscillator}

\begin{theorem}\label{thm:semilinear-conjugation-finite-heisenberg}
With the notation of Theorem~\ref{thm:cyclic-conjugate-jet-generator},
there is a unique \(\nu\)-semilinear bijection
\(\mathcal L_\nu:\mathcal F\to\mathcal F\) satisfying
\[
       J_{t_\nu}\mathcal L_\nu S=\mathcal C_\nu(J_uS)
                                          \qquad(S\in\mathcal F).
\]
For all \(X,Y\in V_f\),
\[
   \nu c(X,Y)=c(b_\nu X,b_\nu Y),\qquad
                  \mathcal L_\nu T_X=T_{b_\nu X}\mathcal L_\nu .
\]
If \(\kappa_\nu\in\widehat\Z^\times\) is the exponent by which
\(\nu\) acts on all roots of unity, then
\(b_\nu\bar b_\nu=\kappa_\nu\) in the rational finite ideles.
\end{theorem}
\begin{proof}
The map \(J_t\) on entire functions is injective for \(t\ne0\).
Indeed vanishing of all Taylor coefficients implies vanishing
on a polydisk by the successive disk formula of
Lemma~\ref{lem:analytic-tools}, then on every complex line
through zero by the one-variable identity principle.
The scaled Taylor map in the assertion has \(t_\nu\ne0\) by
Corollary~\ref{cor:absolute-conjugate-jet-scale}.
On the spanning family \(T_XF\), its prescribed preimages
are \(T_{b_\nu X}F_\nu\) by
Theorem~\ref{thm:cyclic-conjugate-jet-generator}. Every linear
relation between the first family gives the relation with
conjugated coefficients between the second by this identity
and injectivity of \(J_{t_\nu}\). Extending by semilinearity
defines the unique map. Its kernel is zero by injectivity of
\(J_u\) and \(\mathcal C_\nu\). Its image contains all finite
translates of \(F_\nu\), since \(b_\nu\) is an idele, so the
same theorem gives surjectivity.

By local constancy in Theorem~\ref{thm:finite-frame-action}
there is an open subgroup \(H\) of \(V_f\) such that
\(T_xF=F\) and \(T_{b_\nu x}F_\nu=F_\nu\) for \(x\in H\).
For any \(y\in V_f,\ x\in H\), write the product \(T_yT_xF\)
as \(c(y,x)T_{y+x}F\). Apply \(\mathcal L_\nu\) to both
expressions. Since \(T_{b_\nu y}F_\nu\ne0\), the target
addition law gives
\(\nu c(y,x)=c(b_\nu y,b_\nu x)\).
For arbitrary \(x\) choose an integer \(m\ge1\) with \(mx\in H\),
using the lattice neighbourhoods of
Definition~\ref{def:finite-frame-space}. Replace \((y,x)\) by
\((y/m,mx)\). The bilinear trace in the defining character
gives the same values on both sides, proving the identity
for all \(y,x\). Apply it to \(T_yT_xF\) again, now with
arbitrary \(x\). Cyclicity proves the operator identity.

For the last assertion write \({\bf e}_f(h)=\prod_p{\bf e}_p(h_p)\)
on the rational finite adeles as in Proposition~\ref{prop:additive}.
The trace definition in Proposition~\ref{prop:frame-polarization}
gives \(\mathcal E_p(b_\nu X,b_\nu Y)=(b_\nu\bar b_\nu)_p
\mathcal E_p(X,Y)\), with the product in the fixed scalar
field \(\Q_p\). The pairing \(\mathcal E(X,Y)/2\) takes every
finite rational adelic value. In fact choose two integral
symplectic basis vectors with pairing one from
Lemma~\ref{lem:period-coordinates} at level one, and multiply
the second by twice the desired finite adele. Thus
\[
       \nu{\bf e}_f(h)={\bf e}_f((b_\nu\bar b_\nu)h)
                 \qquad (h\in\mathbb A_{\Q,f}).
\]
On the \(p^j\)-th roots of unity \(\nu\) acts by an invertible
exponent modulo \(p^j\), by applying it to a primitive root.
The exponents are compatible under powering, giving
\(\kappa_{\nu,p}\in\Z_p^\times\). As \({\bf e}_p\) depends
only on the finite polar part modulo \(\Z_p\), its values are
\(p\)-power roots and
\(\nu{\bf e}_p(v)={\bf e}_p(\kappa_{\nu,p}v)\).
A scalar \(a\in\Q_p\) with \({\bf e}_p(av)=1\) for all \(v\)
must be zero: otherwise set \(v=1/(pa)\). This character at
\(1/p\) is nontrivial by Proposition~\ref{prop:additive}.
Test \(h\) supported at \(p\) to conclude
\((b_\nu\bar b_\nu)_p=\kappa_{\nu,p}\) for every \(p\).
\end{proof}

\begin{corollary}\label{cor:semilinear-conjugation-physical-translation}
Identify the finite oscillator \(\mathcal H_f\) of a standard
positive space with the sections \(\mathcal F\) by
Theorem~\ref{thm:definite-holomorphic-test-identification},
and denote the transported \(\mathcal L_\nu\) by the same symbol.
For the physical translations of
Definition~\ref{def:adelic-oscillator-translations} one has
\[
               \mathcal L_\nu D_X=D_{b_\nu X}\mathcal L_\nu
                                 \qquad(X\in V_f).
\]
\end{corollary}
\begin{proof}
The scalar correcting \(T_X\) to \(D_X\) in
Theorem~\ref{thm:definite-holomorphic-test-identification}
is \({\bf e}_f(-b(X,X)/2)\), with \(b(X,X)\) the rational
finite trace of \(\jmath\bar\omega H(X,X)\) in
Definition~\ref{def:adelic-oscillator-translations}.
Multiplying \(X\) by \(b_\nu\) scales this trace by
\(b_\nu\bar b_\nu\). The additive character identity of
Theorem~\ref{thm:semilinear-conjugation-finite-heisenberg}
therefore conjugates the correction to the target correction.
Combine this with that theorem's intertwining of \(T_X\).
\end{proof}


\section{Conjugate vacuum lines and the finite Weil operator}

\begin{theorem}\label{thm:finite-weil-jet-conjugation}
Use the physical realization of
Corollary~\ref{cor:semilinear-conjugation-physical-translation}.
Let \(R\) be any finite set of places of \(E\) with
\(R\cap\bar R=\varnothing\), let \(v_0={\bf1}_{L_0}\) be
the unit reference vacuum, and let \(\mathcal V_G\) be the finite
Weil representation of Theorem~\ref{thm:all-finite-weil-splitting}.
There exist a scalar \(z\) in the finite ideles of \(E\) with
\(z\bar z=1\), an integral unit idele \(a\), and \(c\in\C^\times\)
such that
\[
     b_\nu=za,\qquad a_w=1\ (w\in R),\qquad
     \mathcal L_\nu v_0=c\mathcal V_z v_0,\qquad
     \mathcal L_\nu\mathcal V_G
          =\frac{\nu\Delta(G)}{\Delta(G)}\,\mathcal V_G\mathcal L_\nu
\]
for every finite \(E\)-linear isometry \(G\), with \(\Delta(G)\)
as in Definition~\ref{def:finite-splitting-multiplier}.
The last quotient has complex modulus one.
Let \([f]_R\) denote a vector given by a locally constant
compactly supported function \(f\) in the coordinates at \(R\)
of Theorem~\ref{thm:split-finite-schroedinger}, and the reference
vacuum off \(R\cup\bar R\). Then
\[
                   \mathcal L_\nu[f]_R
                              =c\,\mathcal V_z[\nu\circ f]_R .
\]
\end{theorem}
\begin{proof}
By Theorem~\ref{thm:semilinear-conjugation-finite-heisenberg},
\(b_\nu\bar b_\nu=\kappa_\nu\in\widehat\Z^\times\), also
after extending to \(E\). At a place fixed by conjugation
this forces valuation zero. At each pair \(w,\bar w\) meeting
\(R\), choose \(w\in R\) and put
\(z_w=(b_\nu)_w,\ z_{\bar w}=\overline{(b_\nu)_w}^{-1}\),
where the bar identifies the two completions. Do the same,
with either orientation, at the finitely many other pairs
where \(b_\nu\) has nonzero valuation. Give \(z\) value one
elsewhere. Then \(z\bar z=1\), and \(a=b_\nu/z\) is a unit
at every place with \(a_w=1\) on \(R\).
In particular \(aL_0=L_0,\ a\bar a=\kappa_\nu\).

Set \(A=\mathcal V_z^{-1}\mathcal L_\nu\), viewing \(z\)
as the scalar isometry. By
Corollary~\ref{cor:semilinear-conjugation-physical-translation}
and Theorem~\ref{thm:finite-oscillator-lattice-transport},
\(A D_X=D_{aX}A\). Since \(v_0\) is fixed by \(D_l\) for all
\(l\in L_0\), so is \(A v_0\). The joint fixed space is the
line spanned by \(v_0\): the translates indexed by \(V_f/L_0\)
span and have distinct characters under \(L_0\), by that same
lattice transport theorem. Bijectivity gives \(A v_0=cv_0\)
for \(c\ne0\). Thus \(K_\nu=c^{-1}A\) is \(\nu\)-semilinear,
takes \(D_Xv_0\) to \(D_{aX}v_0\), and fixes \(v_0\).

Put \(h=[G L_0:G L_0\cap L_0]=|\Delta(G)|^2\). By
Definition~\ref{def:finite-splitting-multiplier} the value
of \(\mathcal V_G\) on \(v_0\) is the sum of the \(D_r v_0\)
over this quotient multiplied by \(\Delta(G)/h\). The scalar
\(a\) commutes with \(G\), so it permutes this quotient.
Changing a representative inside its coset leaves \(D_rv_0\)
unchanged by Theorem~\ref{thm:finite-oscillator-lattice-transport}.
Hence
\[
 K_\nu\mathcal V_G v_0
           =(\nu\Delta(G)/\Delta(G))\mathcal V_G v_0 .
\]
The identical identity with \(v_0\) replaced by \(D_Xv_0\)
and \(K_\nu\) applied on the right follows by moving the
translation through the operators, since \(aGX=GaX\).
These vectors span. Undo \(K_\nu=c^{-1}\mathcal V_z^{-1}
\mathcal L_\nu\) and use that the representation sends
commuting isometries \(z,G\) to commuting operators.
The factor \(\Delta(G)=\beta(J_G)\sqrt h\) has conjugates of
modulus \(\sqrt h\): Theorem~\ref{thm:actual-global-splitting-character}
gives modulus one for every conjugate of \(\beta(J_G)\), and
\(\nu\sqrt h=\pm\sqrt h\).

In the split coordinates, \(D_x\), supported at \(R\), translates
a function by \(x\), which is fixed by \(a\). A displacement
\(D_y\) supported at \(\bar R\) multiplies by the character
\(t\mapsto\psi_f(\jmath H(y,t))\). Conjugating this character
replaces \(y\) by \(ay\), since \(\bar a_{\bar w}\), transported
to \(w\in R\), equals \(\kappa_{\nu,p}\) for \(w\mid p\) by
\(a\bar a=\kappa_\nu,\ a_w=1\). This uses the additive character
identity in Theorem~\ref{thm:semilinear-conjugation-finite-heisenberg}
after taking the local rational trace. Pointwise conjugation
fixes the box indicator. It therefore agrees with \(K_\nu\)
on translates and character multiples of that indicator, since
both actions have the same semilinearity and intertwine these
two sets of generators. They span the stated functions by
the finite Fourier construction of Theorem~\ref{thm:split-finite-schroedinger}.
Undoing \(K_\nu\) proves the last formula.
\end{proof}


\section{A mean-value estimate on the moved period torus}

\begin{lemma}\label{lem:moved-positive-period-geometry}
Let \(n\ge1,\ g=fn,\ \Gamma_0=r(\mathcal A^n)\) as in
Proposition~\ref{prop:frame-polarization}, and let
\(M_0=\widehat{\mathcal A}^{\,n}\) in physical finite coordinates.
For a finite \(E\)-linear isometry \(G\) of the standard positive
form let \(M=G M_0\) and \(\Gamma_G=r(E^n\cap M)\).
In the ordinary Euclidean metric on \(\C^g\) one has
\[
       {\rm vol}(\C^g/\Gamma_G)=\pi^g,\qquad
       \|l\|^2\in \sigma_0^2\Z_{>0}
                                  \quad(0\ne l\in\Gamma_G).
\]
\end{lemma}
\begin{proof}
Each of \(M,M_0\) is compact open, and their intersection has
finite index in both. The rational-point quotients by this
intersection have the same indices as the adelic lattices:
every coset is open and contains a rational vector by the
density in Theorem~\ref{thm:finite-frame-action}. The rational
points of \(M_0\) are \(\mathcal A^n\), by that theorem's
coordinate argument. Thus \(\Gamma_G\) is commensurable
with \(\Gamma_0\). The two finite indices are equal,
since the isometry preserves additive volume locally.
In fact \(G\) preserves the nondegenerate rational alternating
matrix of Proposition~\ref{prop:frame-polarization}, so at
each rational prime its rational determinant squared is one.
The local change of additive measure is the absolute determinant
by the elimination argument of
Theorem~\ref{thm:split-finite-schroedinger} applied over \(\Q_p\).
Consequently the two real covolumes are equal.
In ordinary real coordinates on \(\C^g\), the matrix of \(E_H\)
of Proposition~\ref{prop:frame-polarization} has determinant
\(\pi^{-2g}\), from its \(g\) alternating two by two blocks.
In a lattice basis of \(\Gamma_0\) it has determinant one by
Lemma~\ref{lem:period-coordinates} at level one. Change of
basis gives the square of the real covolume as \(\pi^{2g}\).

For \(l=r(v)\) put \(q(v)={\rm Tr}_{E/K}\sum_i\bar v_i v_i\).
It is rational, since it belongs to \(K\) and is fixed by
its nontrivial conjugation. On the reference lattice
\(\mathcal A^n\) it is integral by
Proposition~\ref{prop:frame-lattice-trace}. It remains in
\(\Z_p\) on the completion at \(p\): approximate all entries
by lattice vectors and use continuity of this quadratic
polynomial and closedness of \(\Z_p\). The finite isometry
preserves the hermitian norm, hence preserves its local trace.
Thus for rational \(v\in M\) the number \(q(v)\) lies in
every \(\Z_p\) and is an integer by rational prime factorization.
In the positive real coordinates \(\|r(v)\|^2=\sigma_0^2 q(v)\)
by Proposition~\ref{prop:frame-polarization}. It is positive
when \(v\ne0\), proving the second assertion.
\end{proof}

\begin{theorem}\label{thm:moved-vacuum-ball-bound}
Let \(n,g,G\) be as in Lemma~\ref{lem:moved-positive-period-geometry},
let \(X\) be any finite displacement and \(\Phi\) a scalar
multiple of \(D_X\mathcal V_G v_0\). Write \(\|\Phi\|_2\)
for the finite oscillator norm of
Theorem~\ref{thm:finite-oscillator-lattice-transport}.
For its holomorphic section \(S_\Phi\) put
\(A_\Phi(z)=|S_\Phi(z)|e^{-\|z\|^2/2}\). Then \(A_\Phi\)
is periodic under \(\Gamma_G\) and
\[
 \frac1{\pi^g}\int_{\C^g/\Gamma_G} A_\Phi(z)^2\dd z=\|\Phi\|_2^2,
 \qquad
 \sup_{z\in\C^g} A_\Phi(z)
      \le e^{r_d^2/2}\,r_d^{-g}\sqrt{g!}\,\|\Phi\|_2,
       \qquad r_d=\min(1,\sigma_0/3).
\]
Here the integral uses ordinary Euclidean volume on the period torus.
\end{theorem}
\begin{proof}
For \(l\in M=GM_0\), \(D_l\mathcal V_G v_0=\mathcal V_G v_0\)
by intertwining and the reference vacuum condition. The
displacement \(D_X\) adds only its unit-modulus commutator
from Lemma~\ref{lem:theta-positive-translation}. Thus every
\(D_l\), \(l\in M\), acts on \(\Phi\) by a phase.

Consider the modulus of the adelic function \(F^a_{\phi_0\otimes\Phi}\)
in Theorem~\ref{thm:rational-theta-symplectic-action}.
Use the real and finite unitary model changes of
Theorem~\ref{thm:actual-adelic-rational-theta-invariance}
and Proposition~\ref{prop:finite-polarized-vacuum}, as in
Theorem~\ref{thm:definite-holomorphic-test-identification}.
The modulus is rationally periodic, and it is unchanged
under addition of \((0,l)\), \(l\in M\), by the phase property
and the Heisenberg product law.
On the slice with finite coordinate zero this modulus is
\(A_\Phi(r x_\infty)\) by Lemma~\ref{lem:coherent-translation-scale}
at coherent parameter zero: the translated real Gaussian is
\(e^{-\|r x_\infty\|^2/2}B_{-r x_\infty}\). The real formula
there extends from rational coordinates to all real coordinates
by continuity. The slice maps onto the quotient by rational
vectors and \((0,M)\), since any finite vector can be placed in
\(M\) by subtracting a rational one, by
Theorem~\ref{thm:finite-frame-action}. The kernel on the slice
is precisely \(E^n\cap M\) in these coordinates. This proves
the asserted periodicity and identifies that quotient with
\(\C^g/\Gamma_G\). Indeed the induced map from this compact
torus is a continuous bijection, and the target is Hausdorff
since the image of the compact group \(M\) in the compact
adelic quotient is closed.

The probability Haar measure of the rational adelic quotient
pushes to normalized volume on this torus by translation
invariance. Uniqueness follows by averaging any continuous
test function against both probability measures in either
order, using Lemma~\ref{lem:integral-tools}.
The squared integral of \(F^a_{\phi_0\otimes\Phi}\) on the
rational quotient equals \(\|\phi_0\otimes\Phi\|_2^2\)
by the character orthogonality calculation in
Theorem~\ref{thm:rational-theta-symplectic-action}.
The Gaussian \(\phi_0\) of Lemma~\ref{lem:schroedinger-real-gaussian}
has norm one by the Gaussian integral in Lemma~\ref{lem:theta}.
Now use the covolume in Lemma~\ref{lem:moved-positive-period-geometry}
to obtain the integral identity.

For fixed \(z\) the entire function \(W_z S_\Phi(w)\) has
absolute value \(A_\Phi(z)\) at \(w=0\), and modulus
\(e^{\|w\|^2/2}A_\Phi(z+w)\) at every \(w\), by the defining
formula in Theorem~\ref{thm:frame-sections}. Its squared
modulus at zero is at most its average squared modulus on
the ball \(\|w\|<r_d\). To see this, first average the function
on \(e^{i\theta}w\) over \(0\le\theta\le2\pi\): this gives its
value at zero by the circle formula of
Lemma~\ref{lem:analytic-tools}. Cauchy--Schwarz and integration
over the rotation-invariant ball give the squared inequality.
The ball about \(z\) embeds in the torus by the period lower
bound in Lemma~\ref{lem:moved-positive-period-geometry}, since
\(2r_d<\sigma_0\). On it the squared Gaussian factor costs at
most \(e^{r_d^2}\). Its Euclidean volume is
\(\pi^g r_d^{2g}/g!\): for the unit ball integrate the last
complex coordinate in polar coordinates, giving volume
\(2\pi\int_0^1 t(1-t^2)^{g-1}\dd t=\pi/g\) times the preceding
unit volume, starting from dimension zero. Scaling gives the
general radius. The integral over the embedded ball is at
most the full integral of \(A_\Phi^2\). Substitution and a
square root complete the proof.
\end{proof}


\section{An absolute bound for conjugates of differential kernels}

\begin{theorem}\label{thm:uniform-weighted-conjugate-section}
In the notation of Theorems~\ref{thm:finite-weil-jet-conjugation}
and~\ref{thm:moved-vacuum-ball-bound}, use \(v_n^\circ\) of
Corollary~\ref{cor:integral-vacuum-chain-normalization}. Choose
\(C_d\ge0\) as in that corollary and put
\[
    H_n=\exp(C_d nN_1^2)\,e^{r_d^2/2}r_d^{-fn}\sqrt{(fn)!}.
\]
For a finite sum \(\Phi=\sum_j c_jD_{X_j}\mathcal V_{G_j}v_n^\circ\)
and every \(\nu\in{\rm Aut}(\C/K)\) one has
\[
  \sup_{z\in\C^{fn}}e^{-\|z\|^2/2}|S_{\mathcal L_\nu\Phi}(z)|
                       \le H_n\sum_j|\nu c_j|.
\]
There is \(C'_d\), independent of the exponent determining the
field \(E\), for which \(\log H_n\le C'_d(N_1 n+1)^2\).
\end{theorem}
\begin{proof}
Write \(S=S_{\mathcal L_\nu v_n^\circ}\). For every finite \(X\),
the evaluation \((T_{b_\nu X}S)(0)\) is
\(\nu((T_X S_{v_n^\circ})(0))\) by
Theorem~\ref{thm:semilinear-conjugation-finite-heisenberg}
and its series identity at zero. The correcting phase for
the physical \(D_X\) is a root of unity by
Theorem~\ref{thm:definite-holomorphic-test-identification}.
Corollary~\ref{cor:integral-vacuum-chain-normalization} bounds
the modulus of this value by \(B_n=\exp(C_d nN_1^2)\).
The finite scalar \(b_\nu\) is invertible, so this holds at
every finite displacement. At a rational displacement
\(\gamma\) the modulus in question is
\(e^{-\|r\gamma\|^2/2}|S(-r\gamma)|\) by
Theorem~\ref{thm:finite-frame-action}. The vectors \(r(E^n)\)
are dense in the real space \(\C^{fn}\), since \(r\) extends
to a real isomorphism in Proposition~\ref{prop:frame-polarization}
and rational coordinates are dense in real coordinates.
Continuity gives \(e^{-\|z\|^2/2}|S(z)|\le B_n\) for every \(z\).

By Theorem~\ref{thm:finite-weil-jet-conjugation} the finite
vector \(\mathcal L_\nu v_n^\circ\) is a multiple of
\(\mathcal V_z v_0\) for a norm-one finite scalar \(z\).
The mean identity of Theorem~\ref{thm:moved-vacuum-ball-bound}
therefore gives \(\|\mathcal L_\nu v_n^\circ\|_2\le B_n\).
By that same conjugation theorem and
Corollary~\ref{cor:semilinear-conjugation-physical-translation},
\(\mathcal L_\nu D_X\mathcal V_G v_n^\circ\) is obtained from
\(\mathcal L_\nu v_n^\circ\) by a finite isometry, a translation,
and a scalar of modulus one. The finite representation and the
translations preserve Hilbert norm by
Theorems~\ref{thm:all-finite-weil-splitting}
and~\ref{thm:finite-oscillator-lattice-transport}.
It is still a scalar multiple of a single displaced moved
vacuum, so Theorem~\ref{thm:moved-vacuum-ball-bound} bounds
its weighted section by \(H_n\). Semilinearity and the triangle
inequality give the bound for the sum.
Finally \(f\le N_1\) by Proposition~\ref{prop:places},
\(\log(g!)\le g\log g\le g^2\) for \(g=fn\ge1\), and \(r_d\)
depends only on \(d\). These inequalities give the stated constant.
\end{proof}

\begin{lemma}\label{lem:gaussian-pairing-of-jets}
Let \(S\) be entire on \(\C^g\), \(g\ge1\), and
\(|S(z)|\le B e^{\|z\|^2/2}\) everywhere for some \(B\ge0\).
For every polynomial \(P\) in \(g\) variables and \(a\in\C\),
\[
  \left[P(\partial)S(aZ)\right]_{Z=0}
     =\pi^{-g}\int_{\C^g}S(z)P(a\bar z)e^{-\|z\|^2}\dd z .
\]
If \(L\ge0,\ A\ge0,\ w\in\R\) and
\(|P(a\bar z)|\le A\exp(w+L\|z\|^2/(2L+2))\) for all \(z\),
the absolute value of this integral with its prefactor is at
most \(BAe^w(2L+2)^g\).
\end{lemma}
\begin{proof}
By linearity take \(P(Z)=Z^\alpha\) for the identity. Use polar
coordinates in each complex variable. On each product of circles
of fixed radii the Taylor series of \(S\) converges absolutely
and uniformly by the successive disk formulas of
Lemma~\ref{lem:analytic-tools}. Integration over angles retains
exactly its coefficient of \(z^\alpha\). In one complex dimension
\(\pi^{-1}\int_{\C}|z|^{2j}e^{-|z|^2}\dd z=j!\), by the change
\(t=|z|^2\) and \(j\) integrations by parts. Multiplying these
moments yields the factorial converting that coefficient to
the derivative. Absolute integrability for changing the order
of integrations follows from the Gaussian bound times the
polynomial and \(e^{-\|z\|^2}\), by
Lemma~\ref{lem:integral-tools}. For the estimate, insert both
absolute bounds. The remaining exponent is
\(-\|z\|^2/(2L+2)\). The product of the one-variable Gaussian
integrals, divided by \(\pi^g\), is \((2L+2)^g\).
\end{proof}

\begin{corollary}\label{cor:conjugate-polynomial-jet-bound}
Use \(\Phi,\nu,H_n\) of Theorem~\ref{thm:uniform-weighted-conjugate-section}.
For \(a\in\C,\ P\) a polynomial, put \(P_\nu=\mathcal C_\nu P\)
and \(\alpha=(\nu a)t_\nu\). Suppose \(L\ge0,\ A\ge0,\ w\in\R\)
are such that for all real \(R\ge0\) and all \(y\in\C^{fn}\),
\[
       \|y\|^2\le\frac{|\nu a|^2}{\Omega}(2L+2)R
             \quad\Longrightarrow\quad |P_\nu(y)|\le Ae^{w+LR}.
\]
Then
\[
 \left|\nu\!\left(\left[P(\partial)S_\Phi(auZ)\right]_{Z=0}\right)\right|
             \le H_n\left(\sum_j|\nu c_j|\right)Ae^w(2L+2)^{fn}.
\]
To bound the modulus under an arbitrary automorphism of \(\C\),
one may use either that automorphism or its composition on
the left with complex conjugation, choosing the one fixing \(K\).
\end{corollary}
\begin{proof}
The derivative at zero is a finite sum of products of polynomial
coefficients, factorials and Taylor coefficients. Hence by
Theorem~\ref{thm:semilinear-conjugation-finite-heisenberg}
its conjugate is
\([P_\nu(\partial) S_{\mathcal L_\nu\Phi}(\alpha Z)]_{Z=0}\).
The same permutation of variable indices in both series and
polynomial leaves this contraction unchanged. Set
\(R=\|z\|^2/(2L+2)\), \(y=\alpha\bar z\), and use
\(|t_\nu|^2=1/\Omega\) from Corollary~\ref{cor:absolute-conjugate-jet-scale}.
The bound follows from Lemma~\ref{lem:gaussian-pairing-of-jets}
and Theorem~\ref{thm:uniform-weighted-conjugate-section}.
Finally an automorphism preserves the set of the two roots of
\(X^2+d\), so its action on \(K\) is the identity or complex
conjugation. Left composition with complex conjugation preserves
the absolute value of every number.
\end{proof}


\section{Rectangular differential matrices in the tensor norm}

\begin{lemma}\label{lem:rectangular-exterior-operator-bound}
Use the rectangular matrix \(S(Z)\) of
Definition~\ref{def:rectangular-flag-matrix} in a single
block of size \(n'\)-by-\(n\), \(1\le n<n'\), with
\(\kappa(h)=0\) for \(n<h<n'\). Define
\(\lambda_i=\sum_{i\le h<n'}\kappa(h)\) for \(1\le i\le n\).
Use the tensor norms from Lemma~\ref{lem:flag-tensor-integral-realization}.
If \(s_1\ge\cdots\ge s_n\ge0\) are the singular values of \(Z\), then
\[
            \|S(Z)\|_{\rm op}\le\prod_{i=1}^n s_i^{\lambda_i},
      \qquad \sum_i s_i^2=\sum_{bi}|Z_{bi}|^2
\]
with \(s_i^0=1\), including \(s_i=0\).
\end{lemma}
\begin{proof}
For every height \(h\) of the target diagram take one
\(\bigwedge^h Z\) per column. On the source tensor space
adjoin, for each full-height column \(h=n\), an extra factor
\(e_1\wedge\cdots\wedge e_n\) of norm one. On evaluating an
invertible array \(Y\), the source flag vector with these
extra factors equals the full flag tensor at \(Y\) divided
by \((\det Y)^{\kappa(n)}\). The product of the exterior maps
therefore agrees with \(S(Z)\) on these spanning vectors, by
Proposition~\ref{prop:rectangular-flag-law} and
Lemma~\ref{lem:flag-tensor-integral-realization}. They span
since invertible arrays are dense and their evaluations span
the same finite-dimensional space as those of all arrays.
In particular the identification is isometric at both ends.

The Hermitian matrix \(Z^*Z\) has an orthonormal eigenbasis
with nonnegative eigenvalues. For completeness maximize
\(v^*Z^*Zv\) on the compact unit sphere. Varying the
maximizer along each perpendicular unit vector, and along
its multiple by \(i\), shows that the perpendicular component
of \(Z^*Zv\) vanishes. Its perpendicular complement is invariant
by the Hermitian identity. Induct there. In the resulting
basis \(Ze_i\) are orthogonal with lengths \(s_i\), the
nonnegative square roots of the eigenvalues, ordered decreasingly.
Their squared sum is the trace of \(Z^*Z\), which gives the
stated entry sum by direct multiplication.
The image of an increasing orthonormal basis wedge of size \(h\)
has length the product of its \(s_i\), and distinct nonzero
images are orthogonal. On the tensor product of these bases
the same orthogonality persists. Hence the norm of the full
tensor map is bounded by
\(\prod_h(s_1\cdots s_h)^{\kappa(h)}\).
Restriction to the source flag subspace proves the claim.
\end{proof}

\begin{theorem}\label{thm:optimized-absolute-rectangular-jet}
Take one such rectangular block at each \(\tau\in G\), with
the same \(n,n'\) and possibly different column counts.
Use their tensor product \(S(Z)\), put \(b=nn'\), \(g=fb\),
\(N=\sum_{\tau,i}\lambda_{\tau i}\), and set
\(E(0)=0,\ E(k)=k\log k-k\) for \(k\ge1\).
Let \(\Phi=\sum_j c_jD_{X_j}\mathcal V_{G_j}v_b^\circ\),
and \(a\ne0\). For any \(\nu\in{\rm Aut}(\C/K)\) the
entrywise conjugate differential matrix, between its tensor
norm spaces, satisfies
\[
 \begin{split}
 \left\|\nu\!\left([S(\partial)S_\Phi(auZ)]_{Z=0}\right)
                       \right\|_{\rm op}
 \ \le\ &H_b\left(\sum_j|\nu c_j|\right)
            \bigl(e(2+N/g)\bigr)^g\\
 &{}\cdot\exp\!\left(\frac12\sum_{\tau,i}
     \left\{ E(\lambda_{\tau i})
            +\lambda_{\tau i}\log\frac{|\nu a|^2}{\Omega}\right\}\right).
 \end{split}
\]
When \(N=0\) the factor \((e(2+N/g))^g\) can be replaced by \(2^g\).
In every case \(E(k)\le\log(k!)\).
\end{theorem}
\begin{proof}
The coefficients of \(S\) are integers by
Proposition~\ref{prop:rectangular-flag-law}. Renaming the
blocks by \(\nu\circ\tau\) only permutes the input variables
of the polynomial. Let \(A,v\) be arbitrary target row and
source column with \(\|\nu A\|_{\rm dual}\le1,\ \|\nu v\|\le1\);
all complex rows and columns arise after applying \(\nu\).
The polynomial \(\mathcal C_\nu(A S(Z)v)\) at \(y\) therefore
has absolute value at most
\(\prod_{\tau,i}s_{\tau i}^{\lambda_{\tau i}}\), with singular
values taken in the correspondingly renamed blocks, by
Lemma~\ref{lem:rectangular-exterior-operator-bound}.
Tensor operator norms multiply, as follows either by using the
orthogonal tensor bases in that lemma or by diagonalizing the
positive Gram matrix of each operator in an orthonormal basis.

Put \(\theta=|\nu a|^2/\Omega\). For \(L>0\) and
\(\sum_{\tau,i}s_{\tau i}^2=\|y\|^2\le\theta(2L+2)R\), set
\(v_L=\theta(1+1/L)>0\). For all \(x\ge0,\ k\ge0\) integral,
\(x^k e^{-x}\le e^{E(k)}\). For \(k=0\) this follows from
\(e^{-x}\le1\), and for \(k>0\) differentiation gives the
maximum at \(x=k\). Apply the inequality with \(x=s_{\tau i}^2/v_L\)
and take square roots. Multiplication bounds the polynomial by
\[
       e^{LR}(1+1/L)^{N/2}
          \exp\!\left(\frac12
               \sum_{\tau,i}\{E(\lambda_{\tau i})
                                   +\lambda_{\tau i}\log\theta\}\right),
\]
because \(\sum x/2\le LR\). Use this bound in
Corollary~\ref{cor:conjugate-polynomial-jet-bound}.
For \(N>0\) choose \(L=N/(2g)\). This minimizes
\((2L+2)^g(1+1/L)^{N/2}\), since the logarithmic derivative is
\((g-N/(2L))/(L+1)\). Its value is at most
\((2+N/g)^g e^g\) by \(\log(1+x)\le x\).
If \(N=0\) the polynomial is constant of modulus at most one
by the same exterior estimate, and \(L=0,\ w=0\) in the corollary
give the sharper assertion.
Taking the supremum over the two unit balls now gives the operator
norm: a row of dual norm one attains the length of any fixed
nonzero image by the Hermitian inner product. Finally
\(\int_1^k\log t\dd t\le\sum_{j=1}^k\log j\) for \(k\ge1\)
by monotonicity, and \(k=0\) gives equality separately.
\end{proof}


\section{Cancellation of the terminal weight and the first-row overlap}

\begin{proposition}\label{prop:exact-terminal-shortened-exponent}
For all sufficiently large scales, put
\(X_\tau=a{\bf1}_{A}(\tau)\), \(t_1=5r\), and
\(d_i=k+i(t_1+1)\) as in Definition~\ref{def:stage-diagrams}.
The terminal exponent of Definition~\ref{def:terminal-branch-normalization}
is
\[
                         b_\tau=X_\tau-mh_2 .
\]
Consequently the logarithmic contribution of the squared scalar
\(|\sigma a_N|^2\) at every complex automorphism is
\(-Ua\log N+fmh_2\log N\).
\end{proposition}
\begin{proof}
There is one group of \(X_\tau\) height-one columns, all shortened.
The number of tall columns in the last diagram is
\(w_1(\tau)-\lfloor d_1/2\rfloor+(m-1)(2h_2-(t_1+1))\).
Of these, exactly \(m(h_2-\lfloor(t_1+1)/2\rfloor)\) are
retained by Definition~\ref{def:terminal-branch-normalization}.
Add the determinant exponent \(\lfloor d_{m-1}/2\rfloor-w_m(\tau)\)
of Proposition~\ref{prop:actual-definite-arrow-weights}.
By Definition~\ref{def:heights}, \(w_m=w_1+(m-1)2h_2\).
Also \(t_1+1\) is even as \(r\) is a power of three. Thus
\(\lfloor d_{m-1}/2\rfloor-\lfloor d_1/2\rfloor
=(m-2)(t_1+1)/2\). The total coefficient of \(t_1+1\) is
\((m-2)/2-(m-1)+m/2=0\), leaving \(X_\tau-mh_2\).
Each conjugate of \(\bar\pi_N\) has modulus \(\sqrt N\)
by Theorem~\ref{thm:terminal-restriction-height}.
Summing over \(\tau\), with \(|A|=U,\ |G|=f\), proves the claim.
\end{proof}

\begin{lemma}\label{lem:sum-nominal-rectangular-jet-cost}
Use \(J\) of Definition~\ref{def:heights} and let
\(\lambda_{j,\tau i}\) be the length of the \(i\)-th row,
counting from \(1\), on the \(j\)-th arrow of
Definition~\ref{def:stage-diagrams}. Throughout the windows of
Theorem~\ref{thm:mean-covolume} one has, for all sufficiently
large \(x\),
\[
 \sum_{j<m}\sum_{\tau}\sum_{i=1}^{d_j}J(\lambda_{j,\tau i})
 \le mk\sum_\tau J(w_0(\tau))+mU J(a)
          +x^{75}+\frac{2mUa}{x_1}\bigl(1+\log(1+x_1)\bigr).
\]
\end{lemma}
\begin{proof}
Let \(r_{j,\tau}\) count the tall columns on this arrow,
and \(h=2h_2\). Then its first row has length \(X_\tau+r_{j,\tau}\),
its next \(k-1\) rows all have length \(r_{j,\tau}\), and each
row below the \(k\)-th has length at most \(jh\).
Indeed the columns at height \(k\) stop there, while the higher
columns have total multiplicity \(j(h-t_1-1)\).
The factorial identity gives exactly
\[
 J(X_\tau+r_{j,\tau})=J(X_\tau)+J(r_{j,\tau})
                       +\log \binom{X_\tau+r_{j,\tau}}{r_{j,\tau}} .
\]
By Lemma~\ref{lem:diagram-size}, \(r_{j,\tau}\le w(\tau)+x^{49}
\le w_0(\tau)+x^{49}\) and \(r_{j,\tau}\ge w(\tau)\).
Put \(\theta=g_0g_1/\Omega\).
Theorem~\ref{thm:data}, Lemma~\ref{lem:rd} and
Definition~\ref{def:heights} give \(1\le\theta\) eventually
and \(\log\theta=O_{d,M,\epsilon}(\ell)\):
the \(L\)-value and logarithmic derivative are fixed,
\(b^{20}\le g_0g_1\le b^{22}\) and
\(1\le{\rm rd}_E\ll_d f\le O_d(x^{12})\).
Thus \(J\) is increasing on nonnegative integers and each
increment up to \(2a+1\) costs at most \(\log(2a+1)+\log\theta
=O_{d,M,\epsilon}(\ell)\). All row lengths and \(w_0\) in
this comparison are in this range by
Proposition~\ref{prop:height-cost} and Lemma~\ref{lem:diagram-size},
as \(a\ge xb,\ w\le a/x,\ b=e^\ell\).
For the first \(k\) rows the change from \(J(r_{j,\tau})\)
to \(J(w_0(\tau))\) therefore costs at most
\(O_{d,M,\epsilon}(k x^{49}\ell)\) per \((j,\tau)\).
The rows below \(k\) together cost \(O_{d,M,\epsilon}(d_j jh\ell)\)
there. The sum of these two errors is at most \(x^{75}\)
eventually, since \(m\le2\sqrt x,\ d_m=O(x^{11/2}),\
h=O(x^{40}),\ f\le x^{12},\ \ell=x^\epsilon,\ \epsilon<1/100\)
by Proposition~\ref{prop:thick} and Lemma~\ref{lem:diagram-size}.

The binomial term vanishes off \(A\). At \(\tau\in A\) one
has \(w(\tau)=a/x_1\): \(A\) reduces to \(1\bmod M\)
and \(D\) to \(-1\bmod M\), distinct since \(16\mid M\),
by Definition~\ref{def:collisions}. Throughout the windows
\(v=\log x_1\le\ell/10^6\) by Theorem~\ref{thm:data}, so
\(a/x_1\ge x\exp((1-10^{-6})\ell)\ge x^{49}\) eventually.
It follows that \(a/x_1\le r_{j,\tau}\le2a/x_1\).
For a positive integer \(r'\) one has
\(\binom{a+r'}{r'}\le(e(a+r')/r')^{r'}\), using the
numerator bound \((a+r')^{r'}\) and
\(\log(r'!)\ge r'\log r'-r'\) from
Theorem~\ref{thm:optimized-absolute-rectangular-jet}.
Thus the binomial term is at most
\(2a(1+\log(1+x_1))/x_1\) at each of the \(mU\) occurrences.
Summing proves the inequality.
\end{proof}


\section{The archimedean estimate for the untwisted transfer}

\begin{theorem}\label{thm:actual-untwisted-transfer-size}
Take the trivial twist in Theorem~\ref{thm:projected-chain-good-coefficients},
the normalized vacua of Corollary~\ref{cor:integral-vacuum-chain-normalization},
and \(c=-\bar\pi_g u\) as in Theorem~\ref{thm:standard-test-all-prime-jet}
and Corollary~\ref{cor:absolute-conjugate-jet-scale}, with
\(\pi_g\bar\pi_g=g_0g_1\). Write \(B_\Psi\) for the projected
definite class operator with terminal restriction in
Definition~\ref{def:chain-translated-vacuum-coefficients}.
Fix \(\sigma\in{\rm Aut}(\C)\) and a column \(\phi\) of
level \(K_{c,0}\). Suppose \(\|\sigma\phi(y)\|^2\le e^V\)
for every finite argument \(y\), for some real \(V\).
Then throughout the windows of Theorem~\ref{thm:mean-covolume},
for all sufficiently large \(x\), at every finite argument \(h\),
\[
 \begin{split}
 \|\sigma((B_\Psi\phi)(h))\|^2
     &\le \exp\!\left(V+h_0+2\sum_p G_p\log p+E_x\right),\\
 E_x&=x^{76}+\frac{2mUa}{x_1}\bigl(1+\log(1+x_1)\bigr)
                                      =o_{d,M,\epsilon}(fa\ell^3).
 \end{split}
\]
The norms are the tensor norms of
Lemma~\ref{lem:flag-tensor-integral-realization}.
In particular for the spanning columns of Corollary~\ref{cor:short-ram}
this gives the bound \(\exp(h_0+2fa\ell^3+2\sum_p G_p\log p)\)
required in Proposition~\ref{prop:raw-to-short}.
\end{theorem}
\begin{proof}
First replace \(\sigma\), if necessary, by
\(\nu=\mathfrak c\circ\sigma\) fixing \(K\), as in
Corollary~\ref{cor:conjugate-polynomial-jet-bound}.
Entrywise complex conjugation preserves these norms, since
their Gram matrices are real integer matrices by
Lemma~\ref{lem:flag-tensor-integral-realization}.
For a fixed tuple of definite representatives on arrow \(j\)
write \(n=d_j,\ n'=d_{j+1},\ b_j=nn',\ g_j=fb_j\).
The finite operator on the physical test
\(D_{s_j}v_{b_j}^\circ\) is, up to a scalar all of whose
complex conjugates have modulus one, a single translation
of \(\mathcal V_G v_{b_j}^\circ\).
Indeed \(G\) acts on Hom matrices by multiplication on the
left and right by the stage isometries. The finite arrow
in Definition~\ref{def:actual-twisted-arrow-character}
multiplies its Weil operator by a quotient of determinant
characters. Their values at finite norm-one ideles have
all conjugates of absolute value one by
Theorems~\ref{thm:torus-characters-elementary}
and~\ref{thm:actual-global-splitting-character}. The Weil
action on the full finite group is
Theorem~\ref{thm:all-finite-weil-splitting}. Its translation
intertwining moves \(s_j\) to \(G s_j\).

By Theorem~\ref{thm:definite-holomorphic-test-identification},
the jet at \(c=-\bar\pi_g u\) in this arrow is the rectangular
differential at scale \(\bar\pi_g u\). Apply
Theorem~\ref{thm:optimized-absolute-rectangular-jet} with
\(a\) there equal to \(\bar\pi_g\). Its squared modulus
under \(\nu\) is \(g_0g_1\), so its logarithmic weight is
\(\log((g_0g_1)/\Omega)\). Since \(E(k)\le\log(k!)\)
for every row length, the logarithm of the squared operator
norm is at most
\[
       2\log H_{b_j}+2g_j\log\bigl(e(2+N_j/g_j)\bigr)
                       +\sum_{\tau,i}J(\lambda_{j,\tau i}),
\]
with the same inequality for a zero operator interpreted as
a direct exponential upper bound.
Use Corollary~\ref{cor:joint-column-class-size} and then
Theorem~\ref{thm:terminal-restriction-height}.
The logarithm of the squared bound for the whole transfer
is thus at most \(V\) plus the sum over these arrows, plus
\(2(m+1)x^{25}+x^{59}-\sum_\tau b_\tau\log N\).

Here \(\sum_{j<m}(2\log H_{b_j}
       +2g_j\log(e(2+N_j/g_j)))\le x^{59}\) for sufficiently
large \(x\). In fact Theorem~\ref{thm:uniform-weighted-conjugate-section}
gives \(\log H_{b_j}\le C'_d(N_1b_j+1)^2\);
\(N_1\le x^{12}\), \(f\le x^{12}\), \(d_m\ll x^{11/2}\),
\(m\le2\sqrt x\), and \(\log N_j\le3\ell\) by
Corollary~\ref{cor:class-exponents} and Lemma~\ref{lem:diagram-size}.
Proposition~\ref{prop:exact-terminal-shortened-exponent} gives
the terminal term \(-Ua\log N+fmh_2\log N\). The added positive
quantity is at most \(x^{59}\) eventually, since
\(h_2=r^8\ll x^{40},\ \log N\le30m\ell\) by
Propositions~\ref{prop:thick} and~\ref{prop:level-prime}.
Combining these polynomial errors with
Lemma~\ref{lem:sum-nominal-rectangular-jet-cost} gives the
stated \(E_x\), because \(x^{75}+3x^{59}+2(m+1)x^{25}\le x^{76}\)
eventually.

It remains to identify the three main terms from the row
weights and the level. Proposition~\ref{prop:local-ask} gives
\[
 2\sum_p G_p\log p
   =m\sum_{\tau\in A}P(a,\tau)+mk\sum_\tau P(w(\tau),\tau).
\]
Substitution in Definition~\ref{def:heights} yields exactly
\[
  mk\sum_\tau J(w_0(\tau))+mU J(a)-Ua\log N
                                  =h_0+2\sum_p G_p\log p .
\]
For the error term, \(fa\ell^3\ge x e^\ell\) and
\(\ell=x^\epsilon\). Also \(U\le f,\ m\le2\sqrt x\) and
\(v=\log x_1\ge10^6\log x\). For \(v\ge1\),
\((1+\log(1+e^v))e^{-v}\le(2+v)e^{-v}\), a decreasing
upper bound. These inequalities give \(E_x/(fa\ell^3)\to0\)
uniformly. Finally take \(V=fa\ell^3\) for
Corollary~\ref{cor:short-ram} and \(E_x\le fa\ell^3\)
eventually to obtain the last assertion.
\end{proof}



\section{The finite right action on the projected row chain}

\begin{definition}\label{def:row-chain-terminal-action}
Fix a twist \(\vartheta\) with set \(P\) and put \(M_g=\{w:w\mid\bar\pi_g\}\)
in the notation of Definition~\ref{def:split-inclusion-finite-oscillator}.
At stage \(i\) write \(K_i\) for the stabilizer of the integral lattice
of Definition~\ref{def:actual-row-lift-spaces}. Use its isometric frame
from Theorem~\ref{thm:actual-row-lift-frames} to reduce at the places
of \(M_g\cup P\). The first block has size \(r\); the next \(i\)
blocks have size \(t\). Let \(K_0(g)\subset K_0\) be the subgroup
reducing to the identity at \(M_g\). On right \(K_0(g)\)-invariant
functions the element \(s\in\mathcal G\) acts by \(R_0(s)z(h)=z(hk_0)\),
where \(k_0\in K_0\) has marked reduction \(s\).
At each later stage and for each element \(k\) of its finite adelic group
use the right-translation convention \(R(k)F(h)=F(hk)\). For a bounded continuous
automorphic \(z\) at level \(K_0(g)\) set
\[
 C_{U_\vartheta z}=\Pi_Y\mathcal F_{m-1}\cdots\mathcal F_0 z ,
\]
using Definition~\ref{def:moderate-inclusion-chain} with the real
tests of Theorem~\ref{thm:projected-ramified-chain-transfer}.
Theorem~\ref{thm:inclusion-chain-transposition} gives the
successive absolutely convergent integrals and an open right level.
Corollary~\ref{cor:actual-last-descent-rapid} gives a unique
vector-valued coefficient \(U_\vartheta z\), recovered by its
explicit compact integral. In particular \(U_\vartheta\) is linear.
\end{definition}

\begin{theorem}\label{thm:actual-terminal-parabolic-action}
Suppose \(k_m\in K_m\) is block upper triangular after reduction at
every \(w\in M_g\cup P\). Let \(A_{\nu,w}\) denote diagonal block
\(\nu\) of its reduction, numbered from \(0\) to \(m\), and set
\(s=(A_{0,w})_{w\in M_g}\in\mathcal G\).
For every \(z\) in Definition~\ref{def:row-chain-terminal-action}
and finite terminal argument \(h\),
\[
 (U_\vartheta z)(hk_m)=
   \left(\prod_{i=1}^{m}\prod_{w\in M_g\cup P}
           (\omega_w\vartheta_w)^i(\det A_{i,w})\right)
            (U_\vartheta R_0(s)z)(h).
\]
In particular \(U_1z\) is invariant under the terminal integral
subgroup reducing to the identity at \(M_g\).
\end{theorem}
\begin{proof}
Let \(G_{i,w}\) be the upper left block submatrix up to and
including block \(i\) of the reduction of \(k_m\).
The determinant of an upper block triangular matrix is the product
of the block determinants: expanding in permutations, the last
block of rows must use the last block of columns, and induction
removes one block at a time. Thus all \(G_{i,w}\) are invertible.
For \(i<m\) lift these matrices to one \(k_i\in K_i\). At a chosen
split place lift each entry to the local integers. The determinant
is a unit by its nonzero residue, so the inverse given by the
adjugate is integral. In the conjugate component use the inverse
adjoint, and take the identity at the remaining pairs and nonsplit
places. This is an integral isometry in the chosen frames since
the sets of chosen places and their conjugates are disjoint by
Definition~\ref{def:split-inclusion-finite-oscillator}.
At \(i=0\) this also proves that the marked reduction is onto.
Its kernel \(K_0(g)\) is normal, so its action through \(R_0(s)\)
is independent of the lift and is a representation.

For \(1\le i\le m\), block triangularity gives
\(G_{i,w}\binom I0=\binom I0G_{i-1,w}\).
Theorem~\ref{thm:actual-twisted-finite-arrow} multiplies the
inclusion test on arrow \(i-1\) by
\(d_i=\prod_w(\omega_w\vartheta_w)^i(\det A_{i,w})\), since
\(\det G_{i,w}/\det G_{i-1,w}=\det A_{i,w}\).
The real test is unchanged by these finite elements. The representation
law in Theorem~\ref{thm:global-actual-arrow-covariance} therefore gives
\(\mathcal K_{i-1}(ak_{i-1},hk_i)=d_i\mathcal K_{i-1}(a,h)\).
In the forward integral at \(hk_i\) replace its source variable by
\(ak_{i-1}\). Right invariance of the quotient measure gives
\[
 R(k_i)\mathcal F_{i-1}y
          =d_i\,\mathcal F_{i-1}R(k_{i-1})y .
\]
All these changes of variables on the successive forward images
are legitimate by Theorem~\ref{thm:inclusion-chain-transposition};
equivalently the covariance identity bounds the translated integrals
by the original absolute integrals. Iterating gives the displayed
product. The infinite compact projection and the recovery integral
of Corollary~\ref{cor:actual-last-descent-rapid} commute with
translation in the independent finite variable. For the trivial
twist \(P=\varnothing\), and identity reductions make every \(d_i=1\)
and \(s=1\).
\end{proof}

\begin{corollary}\label{cor:actual-translated-projection-formula}
For \(s\in\mathcal G\) choose \(k_\vartheta(s)\in K_m\) having
marked reduction \({\rm diag}(s,I_{mt})\) and identity reduction
at \(P\). Let \(e=\sum_w a_w[w]\) be any sum of distinct central
idempotents of Lemma~\ref{lem:central-proj}. For a ramified column
\(\phi\) put
\[
     Z_s\phi=|\mathcal G|R_0(e)R_0(s)(b\mathcal P_\phi)
\]
with the common \(b\ne0\) of
Theorem~\ref{thm:projected-chain-good-coefficients}.
Then for every twist and every finite \(h\),
\[
   (U_\vartheta Z_s\phi)(h)
       =|\mathcal G|\sum_{w\in\mathcal G} a_w
                      (B_{\Psi_\vartheta}\phi)(h k_\vartheta(ws)).
\]
\end{corollary}
\begin{proof}
The local lifts exist by the matrix construction in
Theorem~\ref{thm:actual-terminal-parabolic-action}. On these lifts
every block after the first has determinant one, so their action is
\(U_\vartheta R_0(s)=R(k_\vartheta(s))U_\vartheta\).
The boundedness and initial principal level of \(\mathcal P_\phi\)
follow from Theorem~\ref{thm:actual-full-poincare-columns} and
Lemma~\ref{lem:initial-frame-central-translation}. By recovery
of coefficients the identity in
Theorem~\ref{thm:projected-chain-good-coefficients} reads
\(U_\vartheta(b\mathcal P_\phi)=B_{\Psi_\vartheta}\phi\).
Finally \(R_0(w)R_0(s)=R_0(ws)\) by the right-translation convention.
Expanding \(R_0(e)\) gives the asserted equality.
\end{proof}



\section{Short rows and their simultaneous local index}

\begin{lemma}\label{lem:terminal-class-projection-isometry}
In the notation of Corollary~\ref{cor:actual-translated-projection-formula}
let \(T_{\vartheta,s}\eta=|\mathcal G|\sum_w a_w
R(k_\vartheta(ws))\eta\). It acts on coefficient forms invariant
under the subgroup \(J^0_\vartheta\subset K_m\) reducing to the
identity at \(M_g\cup P\). For every prime \(p\) and embedding
\(\iota\) of Definition~\ref{def:place-height} it preserves the
module given by \(|\eta(h)|_{Y,h}\le1\) for all \(h\), with the
gauges of Theorem~\ref{thm:actual-unmarked-chain-divisibility}.

For \(P=\varnothing\) take \(J=J^0_1\), write
\(\Gamma=U(Y_m)(E^+)\), where \(E^+\) is the real subfield of \(E\),
and use the terminal class norm
\[
 \|\eta\|_{\sigma,J}^2=
  \sum_{h\in\Gamma\backslash U(Y_m)_f/J}
        \frac{\|\sigma(\eta(h))\|^2}
             {|\Gamma\cap hJh^{-1}|}.
\]
Then \(\|T_{1,s}\eta\|_{\sigma,J}^2\le
|\mathcal G|^2\|\eta\|_{\sigma,J}^2\) for every
\(\sigma\in{\rm Aut}(\C)\).
\end{lemma}
\begin{proof}
The kernel \(J^0_\vartheta\) is normal in \(K_m\), and the prescribed
residues \(({\rm diag}(s,I),1)\) multiply as \(s\) does. Thus the
right translations give a representation of \(\mathcal G\) on this
space. All \(U_\vartheta z\) belong to it by
Theorem~\ref{thm:actual-terminal-parabolic-action}.
At \(p\) the weight matrices \(D_Y(k_p)^{\pm1}\) of any \(k\in K_m\)
have integral entries. Indeed the terminal frame is the standard
one by Definition~\ref{def:actual-row-lift-spaces}, and these
matrices are products of determinant units and integral evaluation
polynomials, as in Theorem~\ref{thm:local-gauge-class-matrix}
and Proposition~\ref{prop:terminal-restriction-gauge}.
Multiplicativity gives \(|v|_{Y,hk}=|v|_{Y,h}\) by applying the
ultrametric inequality to both matrices. Each coefficient
\(|\mathcal G|a_w\) is integral by Proposition~\ref{prop:integral-proj},
so the same inequality applied to the sum proves preservation.

For the complex assertion every \(\gamma\in\Gamma\) acts by a
unitary matrix at every embedding of \(E\), including after
application of \(\sigma\): the embeddings on \(E\) commute with
complex conjugation as shown in Proposition~\ref{prop:short-col}.
Their actions on exterior powers, determinant powers and subspaces
of tensor products preserve the induced inner products by
Lemma~\ref{lem:flag-tensor-integral-realization}.
Thus the norm is independent of the rational and level choices
of representatives. The stabilizers are finite by
Lemma~\ref{lem:definite-class-quotient-measures}. Right
multiplication by \(k_1(s)\) permutes the classes and preserves
their stabilizers up to rational conjugacy, since it normalizes
\(J\). It acts unitarily also on the space of entrywise
\(\sigma\)-transported forms with this norm. Now
\(\sigma(e)\) is a self-adjoint idempotent by
Proposition~\ref{prop:integral-proj}. In this unitary action
its norm is at most one by orthogonal decomposition. The
remaining translation is an isometry, proving the asserted
squared factor.
\end{proof}

\begin{theorem}\label{thm:actual-short-rows-simultaneous-index}
Let \(B=fa\ell^3,\ \ell=\log b_{\rm sc}\), where \(b_{\rm sc}\)
denotes the scale \(b\) of Theorem~\ref{thm:data}. Use the rational
\(c>0\) of Proposition~\ref{prop:exact-normalization} and the maps
\(Z_s\) of Corollary~\ref{cor:actual-translated-projection-formula}.
For all sufficiently large \(x\) on the windows of
Theorem~\ref{thm:mean-covolume} there is a spanning family of
everywhere integral ramified columns such that, for every member
\(\phi\), every \(s\in\mathcal G\) and \(\sigma\in{\rm Aut}(\C)\),
\[
       \|U_1(cZ_s\phi)\|_{\sigma,J}^2
                   \le \exp\bigl(H-m\mathcal R+o_{d,M,\epsilon}(B)\bigr),
\]
with uniform error. In addition, take \emph{any} finite list
\(y_i=cZ_{s_i}\phi_i\) of everywhere integral ramified columns
so mapped, with at most \(u\) distinct \(s_i\).
For each \(p,\iota\) there is one square matrix \(A_\iota\) of the
size of the list over \(O_\iota=\{z:|\iota z|_p\le1\}\) satisfying
\[
 |\iota(\det A_\iota)|_p\ge e^{-u c_0e_p\log p},\qquad
 \left|\left(U_\vartheta\sum_j (A_\iota)_{ij}y_j\right)(h)\right|_{Y,h}
 \le 1
\]
for every \(i,h\) and every allowed twist with \(p\nmid Q\).
Here \(c_0\) is the dimension of the column space and the \(e_p\)
are the losses of Proposition~\ref{prop:loss}.
\end{theorem}
\begin{proof}
Take the family constructed in Corollary~\ref{cor:short-ram}.
Its proof gives an upper bound \(V=f\log4+2aU\log(kaC_0)\)
for the logarithm of its squared pointwise norm at every conjugate.
Since \(U\le f\), \(\log a\le 2\ell\) and
\(\log(kC_0)=O_{d,M}(1+\log x)\) there,
\(V=O_{d,M,\epsilon}(fa\ell)\) uniformly.
For the terminal group the rank \(r+mt=k+m(t+1)-1\)
is less than \(2n_m\) in Corollary~\ref{cor:class-exponents}.
Moreover \(J\) contains the principal level modulo \(gN\):
at the marked places reduction modulo \(gN\) implies identity
modulo the residue prime. Hence there are at most \(\exp(x^{25})\)
classes and each of their weights is at most one.
Theorem~\ref{thm:actual-untwisted-transfer-size} bounds the
pointwise squared norm of the raw definite image at every
representative. Apply Lemma~\ref{lem:terminal-class-projection-isometry}
and the identity of Corollary~\ref{cor:actual-translated-projection-formula}.
After also multiplying by \(c\), the logarithmic upper bound is
\[
 V+h_0+E_x+2L_{\rm loc}+x^{25}+2\log|\mathcal G| ,
\]
because \(2\log c+2\sum_pG_p\log p\le2L_{\rm loc}\) by
Proposition~\ref{prop:exact-normalization}.
The proof of Proposition~\ref{prop:height-cost} gives
\(h_0-H+m\mathcal R\le U\{m(1+\log a)+a(G_0+15\ell/2)\}
=O(fa\ell)\), since \(m\le a,\ G_0\le23\ell\).
By Proposition~\ref{prop:row-group} and Theorem~\ref{thm:data},
\(\log|\mathcal G|\le22fr^2\ell\) and \(r=O(x^5)\).
The polynomial costs are \(o(B)\) by \(a\ge x\exp(x^\epsilon)\).
The same holds for \(fa\ell\), for \(L_{\rm loc}\) by
Proposition~\ref{prop:exact-normalization}, and for \(E_x\)
by Theorem~\ref{thm:actual-untwisted-transfer-size}.
These bounds give the first assertion.

For the second assertion take the column modules and their
index witnesses of Corollary~\ref{cor:local-transfer-proved-all-primes}.
For the embedding \(\iota\) these witnesses depend on the column
module and on \(p\); its guaranteed gain holds simultaneously for
all the indicated \(B_{\Psi_\vartheta}\). On the ramified submodule
the identity of Corollary~\ref{cor:actual-translated-projection-formula}
reads \(U_\vartheta Z_s=T_{\vartheta,s}B_{\Psi_\vartheta}\),
with \(T_{\vartheta,s}\) preserving the integral module by
Lemma~\ref{lem:terminal-class-projection-isometry}. For the common
target of the \(Z_s\) take the linear span of their images as
bounded continuous functions of principal level. Its boundedness
follows from Theorem~\ref{thm:actual-full-poincare-columns}.
Proposition~\ref{prop:exact-normalization} now supplies the single
matrix for the entire list and all allowed twists. For the empty
list the empty matrix has determinant one.
\end{proof}



\section{Independent short rows}

\begin{proposition}\label{prop:identify-actual-short-row-spans}
Let \(W,D,R,Q_s\) be as in Definition~\ref{def:actual-marked-span}.
Identify old and new marked row groups by transport of the initial
lattice in Lemma~\ref{lem:initial-frame-central-translation}.
Use the transported good central idempotent for \(e\) in
Corollary~\ref{cor:actual-translated-projection-formula}.
Write \(S_{\rm row}z(h)=z(hb_0)\) for the same initial change of
frame. For every subset \(T\) of the marked group, identifying its
elements by this transport, one has
\[
  \sum_{s\in T}cZ_s({\rm Col}_{\rm ram})
       = S_{\rm row}\!\left(\sum_{s\in T} RQ_s(V_{\rm ram})\right)
       \ \subseteq\ S_{\rm row}W .
\]
Here \(c\) is the positive rational of
Theorem~\ref{thm:actual-short-rows-simultaneous-index}.
\end{proposition}
\begin{proof}
Let \(b_X\) be the central finite definite component of \(b_0\).
The substitution \(\phi\mapsto[h\mapsto\phi(hb_X^{-1})]\)
is an automorphism of the column space: it commutes with the left
rational covariance and, by centrality, with the right level. Its
inverse uses \(b_X\). It commutes with each central difference
operator at \(N\) in Theorem~\ref{thm:ram-dim}, hence carries the
ramified space onto itself. Use an orthonormal basis of the complex weight space of the real group
and its conjugate basis for the coefficient
representation. The bilinear pairing of
Lemma~\ref{lem:initial-bessel-column-coefficient} sends the latter
basis to the dual of the former. In particular
\(\lambda_d(v)=\mathcal B(\phi(h_d b_X^{-1}),v)\) has as row of
coordinates exactly the coordinates of the translated column at
\(h_d\). This describes the entire space \(V_{\rm ram}\) in
Definition~\ref{def:actual-marked-span} when \(\phi\) varies
over the ramified columns.

By Theorem~\ref{thm:actual-base-bessel-periodization} the rational
periodization with these rows of coordinates is their linear
combination of the original full columns. By
Definition~\ref{def:normalized-actual-poincare} that sum is
\(\sqrt{\mathfrak v}\,Qv\) in the old frame for the resulting
coordinate vector \(v\); the normalizing volume is positive.
Lemma~\ref{lem:initial-frame-central-translation} identifies its
right translate by \(b_0\) with \(\mathcal P_\phi\).
If \(k\) is an old integral representative, then
\[
 R(b_0 k b_0^{-1})S_{\rm row}=S_{\rm row}R(k).
\]
The lattice transport sends the full stabilizer to the full
stabilizer since the basis transport there differs from \(b_0\)
only by a central scalar. It also identifies their principal
kernels, so this equality descends to both marked group actions.
The idempotent inducing \(R\) is the product of the marked
projections of Theorem~\ref{thm:actual-marked-first-moments}.
Their product is a sum of distinct central projections: the
regular unitary action decomposes into irreducible spaces by
Lemma~\ref{lem:central-proj}; on each irreducible both commuting
central projections act by \(0\) or \(1\), as does their product.
Transport therefore gives an allowed choice of \(e\).
Multiplication by \(c|\mathcal G|b\sqrt{\mathfrak v}\) is by a
nonzero scalar, by Theorem~\ref{thm:projected-chain-good-coefficients}
and Proposition~\ref{prop:exact-normalization}. Taking spans proves
the equality and Theorem~\ref{thm:actual-ramified-rank-lower}
gives the inclusion.
\end{proof}

\begin{corollary}\label{cor:independent-actual-short-rows}
Write \(b_{\rm sc}\) for the arithmetic scale. There is a set \(T\)
with \(|T|c_0\le3D\log(1+b_{\rm sc}^2)\) and an independent family
\(y_i=c Z_{s_i}\phi_i\) of length
\(n\ge(1-4/b_{\rm sc}^2)D\), with \(s_i\in T\) and short
everywhere integral \(\phi_i\), for which
\[
 \|U_1 y_i\|_{\sigma,J}^2
       \le \exp(H-m\mathcal R+o_{d,M,\epsilon}(B)).
\]
For every \(p,\iota\) a simultaneous integral recombination for
all twists prime to \(p\) has determinant with absolute value
at least \(\exp(-|T|c_0e_p\log p)\). The total logarithmic cost is
at most \(3D\log(1+b_{\rm sc}^2)L_{\rm loc}\).
If \(\dim W\le D\), the same assertions hold with \(D\) replaced
throughout this corollary by \(\dim W\).
\end{corollary}
\begin{proof}
Apply Proposition~\ref{prop:actual-numerical-selection} to the old
column spaces and transport its set by
Proposition~\ref{prop:identify-actual-short-row-spans}.
The resulting sum has the same dimension, since the inverse
translation uses \(b_0^{-1}\). By Corollary~\ref{cor:short-ram}
the short columns span the entire ramified source. Their images
under the selected maps therefore span this sum. Extract a basis
from those images by adjoining a vector outside the current span
until its full dimension has been reached. This uses only spanning
by short integral columns in the new domain.
Theorem~\ref{thm:actual-short-rows-simultaneous-index} gives both
bounds for the chosen rows. Sum the logarithmic costs using
\(L_{\rm loc}=\sum_p e_p\log p\). Under the dimension hypothesis
use Corollary~\ref{cor:actual-selection-upper-only} to choose the
old translates with its improved count and dimension. The same
transport and extraction apply.
\end{proof}



\section{The terminal level and its inner product}

\begin{theorem}\label{thm:explicit-terminal-level-index}
At the marked places \(M_g\) require identity reduction on \(Y_m\);
at the twist places \(P\) require upper block triangular reduction
of type \((r,t,\ldots,t)\), with \(m\) blocks of size \(t\).
Let \(J_\vartheta\) be this subgroup of \(K_m\).
Write \(q_w=|\mathbf F_w|\), \(Q=\prod_{w\in P}q_w\),
\(g_h(q)=\prod_{a=0}^{h-1}(q^h-q^a)\), and set
\[
 A_h=\prod_{w\in P}g_h(q_w),\qquad
 W_i(P)=\prod_{w\in M_g\cup P}g_{r+it}(q_w),\qquad
 k_m=mrt+\binom m2t^2.
\]
Empty products are one. If \(I_m(P)=[K_m:J_\vartheta]\), then
\[
 I_m(P)=\frac{W_m(P)}{A_r A_t^m Q^{k_m}} .
\]
At rank \(r+t\) let \(I_{\rm aux}(P)\) be the integral level index
where at each marked place the first \(r\) coordinate vectors are
fixed individually, and at each twist place their complementary
\(t\)-plane is preserved. Put
\[
  \delta(P)=\frac{I_{\rm aux}(P) A_r^2 Q^{rt}}{W_0(P)},\qquad
  C=W_0(\varnothing)
      \left(\frac{W_1(\varnothing)}{\delta(\varnothing)}\right)^m .
\]
Then \(\delta(P)>0\) and the normalized index is
\[
            I_m(P)=C Q^{-k_m}
             \frac{\delta(P)^m W_m(P)}{W_0(P) W_1(P)^m}.
\]
\end{theorem}
\begin{proof}
The frame and split lifts in
Theorem~\ref{thm:actual-terminal-parabolic-action} show that
reduction is onto the full product in every rank in question.
The number of invertible \(h\)-by-\(h\) matrices over \(\mathbf F_q\)
is \(g_h(q)\) by choosing successive columns outside their
predecessors' span, as in Proposition~\ref{prop:row-group}.
An upper block triangular matrix is invertible precisely when
all its diagonal blocks are invertible. All entries strictly
above the diagonal blocks are free. For the \(r\)-block and \(m\)
successive \(t\)-blocks these provide \(mrt+\binom m2t^2\)
entries. Multiplying the orders at \(P\), with trivial subgroup
at \(M_g\), gives \(A_r A_t^m Q^{k_m}\). Division gives the
first formula.

Put \(V_g=\prod_{w\in M_g}g_t(q_w)q_w^{rt}\).
At a marked place of the auxiliary level the matrices have
diagonal blocks \(I_r,D\), lower left block zero and arbitrary
upper right block. At a twist place the upper right block is zero
and both diagonal blocks are arbitrary invertibles, with all
\(rt\) lower left entries free. The orders therefore give
\[
 I_{\rm aux}(P)=\frac{W_1(P)}{V_g A_r A_t Q^{rt}},\qquad
 \frac{\delta(P)}{W_1(P)}
       =\frac1{V_g W_0(\varnothing) A_t},
\]
using \(W_0(P)=W_0(\varnothing)A_r\). All factors are positive.
At \(P=\varnothing\) the second ratio equals
\(1/(V_g W_0(\varnothing))\). Dividing by that value and substituting
in the first formula for \(I_m(P)\) proves the identity.
\end{proof}

\begin{proposition}\label{prop:projected-class-inner-normalization}
Use the terminal coefficient space \(E_Y\) of
Definition~\ref{def:terminal-coefficient-pairing}, put
\(d_Y=\dim E_Y\), and use \(b_H>0\) of
Lemma~\ref{lem:definite-class-quotient-measures} with finite
\(K_m\)-volume one. Let \(z,z'\) satisfy
Definition~\ref{def:row-chain-terminal-action}. Set
\(\eta=U_\vartheta z,\ \eta'=U_\vartheta z'\), and write
\(F,F'\) for their scalar forward chains before terminal projection.
Then
\[
 \sum_{h\in\Gamma\backslash U(Y_m)_f/J_\vartheta}
      \frac{\langle\eta(h),\eta'(h)\rangle}
           {|\Gamma\cap hJ_\vartheta h^{-1}|}
       =\frac{I_m(P)d_Y}{b_H}\,(F,\Pi_Y F')_{[H]} .
\]
\end{proposition}
\begin{proof}
At \(k\in J_\vartheta\) the two coefficients transform by the
same character by Theorem~\ref{thm:actual-terminal-parabolic-action}.
It has modulus one since each tame factor has finite image.
Rational left multipliers are unitary at infinity by
Lemma~\ref{lem:terminal-class-projection-isometry}. Hence their
pointwise Hermitian product is a function on the indicated double
quotient. Compact averaging in
Theorem~\ref{thm:terminal-branching-pairing} gives
\(\int_{H_\infty}\overline{C_\eta(u,h)}C_{\eta'}(u,h)\dd u
=\langle\eta(h),\eta'(h)\rangle/d_Y\), with probability Haar
measure. Indeed the averaged rank-one orthogonal projection
commutes with the irreducible compact action and has trace one,
which determines its scalar. Apply both integral formulas of
Lemma~\ref{lem:definite-class-quotient-measures} and the index
of Theorem~\ref{thm:explicit-terminal-level-index}. Finally
\(C_\eta=\Pi_Y F,\ C_{\eta'}=\Pi_Y F'\) and \(\Pi_Y\) is a
self-adjoint idempotent by Corollary~\ref{cor:actual-last-descent-rapid}
and the compact-integral argument of
Theorem~\ref{thm:projected-ramified-chain-transfer}. This gives
the displayed pairing.
\end{proof}



\section{A real test for the projected terminal coefficient}

\begin{proposition}\label{prop:last-compact-projection-on-test}
Let \(K=U(Y_m)_\infty\), \(d_Y=\dim E_Y\), and take the unit vector
\(u_0\) of Corollary~\ref{cor:actual-last-descent-rapid}.
For the last arrow denote its action at infinity by
\(\Omega_\infty\). If \(f\) is its real Schwartz test, set
\[
 f^\natural=d_Y\int_K
         \langle u_0,R_Y(k)u_0\rangle\,\Omega_\infty(1,k)f\,\dd k ,
\]
with probability Haar measure. This is a Schwartz test independent
of the finite twist. For any moderate automorphic function \(y\)
at stage \(m-1\), the last forward integral with real factor
\(f^\natural\) equals \(\Pi_Y\mathcal F_{m-1}y\) with original
real factor \(f\). Consequently replacing the last test by
\(f^\natural\) realizes \(C_{U_\vartheta z}\) directly as the
forward chain, for every \(z\) of
Definition~\ref{def:row-chain-terminal-action}.
\end{proposition}
\begin{proof}
The real operator on the consecutive spaces is the real Weil
representation multiplied by determinant characters, as in
Theorem~\ref{thm:global-actual-arrow-covariance}.
Theorem~\ref{thm:negative-block-real-splitting} gives continuity
of the Weil action in every Schwartz seminorm. The determinant
characters on real isometries are determined by their infinity
weights, which are independent of the finite twist by
Theorem~\ref{thm:actual-twisted-finite-arrow}. Their scalar
functions are continuous and have modulus one. Thus on the compact
group \(K\) the seminorms \(p_b\) of
Lemma~\ref{lem:real-small-heisenberg-differences} of the integrand
are uniformly bounded for each \(b\). Define \(f^\natural\)
by the pointwise integral. On each compact real coordinate box,
derivatives pass through the integral: a coordinate difference
quotient is an average of the next derivative along a segment,
which has a uniform bound there, so apply
Lemma~\ref{lem:integral-tools} and iterate. Multiplying each
resulting derivative by its defining power of \(1+\|x\|\)
gives its required uniform seminorm estimate. This proves the
Schwartz and independence claims.

Writing \(v=uk\) in the projection kernel of
Corollary~\ref{cor:actual-last-descent-rapid} gives, for every
terminal function \(F\),
\[
 (\Pi_Y F)(u,h)=d_Y\int_K
     \langle u_0,R_Y(k)u_0\rangle F(uk,h)\,\dd k .
\]
By the representation law of
Theorem~\ref{thm:global-actual-arrow-covariance} the substitution
\(u\mapsto uk\) in the last forward kernel simply acts on its
Schwartz test by \(\Omega_\infty(1,k)\). For justification of
the interchange, \(Y_{m-1}=X_{m-1}\) and \(Y_m\) are both definite
by Definition~\ref{def:actual-row-lift-spaces}. The source
quotient has compact fundamental representatives by
Lemma~\ref{lem:definite-class-quotient-measures}, and \(y\) is
bounded there by continuity. With \(u,h\) fixed the kernel on
these representatives times \(K\) is uniformly bounded by the
seminorm estimate in Lemma~\ref{lem:real-small-heisenberg-differences}.
Integration in \(k\) commutes with evaluation of the theta
functional on the test: the seminorm estimate gives continuity
of this linear functional; finite partitions of compact \(K\)
with representative in each part approximate the continuous
integrand uniformly in that seminorm. Such partitions follow from
a finite open cover on which its variation is arbitrarily small,
made disjoint by successive differences. The corresponding sums
converge to the pointwise test integral in the seminorm by the
derivative estimates above. Finally exchange the source and
compact integrals by Lemma~\ref{lem:product-integration}.
This gives precisely the kernel with real test \(f^\natural\).
Applying the assertion to the preceding moderate forward image,
given by Theorem~\ref{thm:inclusion-chain-transposition}, proves
the conclusion.
\end{proof}

\begin{corollary}\label{cor:actual-terminal-weighted-gram}
Let \(z_1,\ldots,z_d\) be bounded continuous automorphic functions
of initial principal level. Write \(F_i^\natural\) for the scalar
forward image with last test as in
Proposition~\ref{prop:last-compact-projection-on-test}.
The determinant of the \(d\)-by-\(d\) class Gram matrix of
\(U_\vartheta z_i\) at level \(J_\vartheta\) equals
\[
 \left(\frac{d_Y C}{b_H}
        Q^{-k_m}\frac{\delta(P)^m W_m(P)}
                         {W_0(P)W_1(P)^m}\right)^{\!d}
       \det\big((F_i^\natural,F_j^\natural)_{[H]}\big)_{i,j}.
\]
For \(d\ge1\) it is positive exactly when the functions
\(F_i^\natural\) are linearly independent. In particular this
criterion applies to the independent rows of
Corollary~\ref{cor:independent-actual-short-rows}.
\end{corollary}
\begin{proof}
Proposition~\ref{prop:last-compact-projection-on-test} gives
\(F_i^\natural=C_{U_\vartheta z_i}\), fixed by the self-adjoint
terminal projection. Apply
Proposition~\ref{prop:projected-class-inner-normalization}
entrywise and use Theorem~\ref{thm:explicit-terminal-level-index}.
Factoring out the positive scalar from each row gives its \(d\)-th
power. The integral of the squared modulus of any nonzero
continuous function is positive, as proved at the end of
Theorem~\ref{thm:projected-ramified-chain-transfer}. Thus on
the finite span of these functions the Petersson form is positive
definite. Expressing the family in an orthonormal basis, the
kernel of its Gram matrix is the kernel of its column map,
because the associated quadratic form is the squared norm of
the linear combination. All Gram eigenvalues are nonnegative,
so positivity of its determinant is equivalent to independence.
\end{proof}



\section{Injectivity from descent along the triangular array}

\begin{definition}\label{def:row-triangle-operators}
Fix \(m\ge1\), a complex vector space \(V\) and vector spaces
\(F_{i,n}\) for \(i,n\ge0,\ i+n\le m\). For \(i+n<m\)
let
\[
 a_{i,n}:F_{i,n}\longrightarrow F_{i,n+1},\qquad
 d_{i,n}:F_{i,n+1}\longrightarrow F_{i+1,n}
\]
be linear. Write \(A_{i,0}=1_{F_{i,0}}\) and
\(A_{i,n+1}=a_{i,n}A_{i,n}\). At level zero take injections
\(j_i:V\longrightarrow F_{i,0}\) for \(0\le i\le m\).
Write \(W\) for any subspace of \(V\). For \(i<m\) suppose
there is an injective \(H_i:W\longrightarrow W\) with
\[
             d_{i,0}a_{i,0}j_i x=j_{i+1}H_i x\qquad(x\in W).
\]
\end{definition}

\begin{theorem}\label{thm:one-sided-row-triangle-injectivity}
In Definition~\ref{def:row-triangle-operators}, suppose for each
\(i+n+2\le m\) there is a scalar \(t_{i,n}\) with
\[
 a_{i+1,n}d_{i,n}y=t_{i,n}\,d_{i,n+1}a_{i,n+1}y
            \qquad \bigl(y\in A_{i,n+1}j_i(W)\bigr).
\]
Then for all \(i+n+1\le m\) and \(x\in W\),
\[
 A_{i+1,n}j_{i+1}H_i x
    =\left(\prod_{u=0}^{n-1}t_{i,u}\right)
                    d_{i,n} A_{i,n+1}j_i x .
\]
For every \(i+n\le m\), the restriction of \(A_{i,n}j_i\) to \(W\)
is injective. The assertions also hold with the arrows in any
row multiplied by individually prescribed nonzero scalar weights,
in the sense that their successive composites have the same
kernels on \(W\).
\end{theorem}
\begin{proof}
For \(n=0\) the identity is the relation defining \(H_i\).
If it holds at \(n\), apply \(a_{i+1,n}\) on the left of both
sides. On the right the vector to which \(a_{i+1,n}d_{i,n}\)
is applied is exactly \(A_{i,n+1}j_i x\). The square relation
in the statement may therefore be used at that vector. Its
new source after applying \(a_{i,n+1}\) is
\(A_{i,n+2}j_i x\), proving the identity at \(n+1\).

Prove injectivity simultaneously in all rows by induction on
the length. For length zero it follows from injectivity of
\(j_i\). If \(A_{i,n+1}j_i x=0\), the displayed identity gives
\(A_{i+1,n}j_{i+1}H_i x=0\), irrespective of the value of the
scalar product in that identity. By the assertion for length
\(n\) and \(H_i x\in W\) we have \(H_i x=0\); then \(x=0\).
Multiplying individual arrows by nonzero scalars multiplies
their composite by their product, which preserves its kernel.
\end{proof}

\begin{corollary}\label{cor:terminal-gram-first-arrow-test}
For one twist take any finite-dimensional space \(W\) of bounded
initial principal-level forms as in
Definition~\ref{def:row-chain-terminal-action}. Suppose there
is an array as in Theorem~\ref{thm:one-sided-row-triangle-injectivity}
for this \(W\) such that its weighted first-row composite on \(W\)
is the scalar forward chain with the last test
\(f^\natural\) of Proposition~\ref{prop:last-compact-projection-on-test}.
Then \(U_\vartheta\) is injective on \(W\) and its class Gram
determinant in every basis is strictly positive.
\end{corollary}
\begin{proof}
Any vanishing linear combination of the forward images is the
image of the same linear combination in \(W\). Its injectivity
follows from Theorem~\ref{thm:one-sided-row-triangle-injectivity}.
The equivalence of independence with positivity of the determinant
is Corollary~\ref{cor:actual-terminal-weighted-gram}.
Coefficient recovery in Definition~\ref{def:row-chain-terminal-action}
then gives the assertion for \(U_\vartheta\).
\end{proof}



\section{Gram forms along an interval of the triangular array}

\begin{theorem}\label{thm:exact-triangle-gram-product}
Endow the spaces of Definition~\ref{def:row-triangle-operators}
with complex inner products conjugate linear in the first variable.
Let \(W\ne0\) be finite dimensional and suppose the \(j_i\) are
isometries on \(W\). Choose positive weights \(w_0,\ldots,w_m\).
Suppose the square multipliers of
Theorem~\ref{thm:one-sided-row-triangle-injectivity} are
\(t_{i,n}=w_{n+2}/(w_n\overline{\beta_{i,n}})\) with
\(\beta_{i,n}\ne0\). Impose the following two pairing identities.
For every \(x\in W,\ z\in F_{i,1}\),
\[
    \langle z,a_{i,0}j_i x\rangle
                 =\langle d_{i,0}z,j_{i+1}x\rangle .
\]
For each \(i+n+2\le m\) and \(z\in F_{i,n+2}\), suppose there is
\(z'\in F_{i,n+1}\) such that for all \(v\in F_{i,n+1}\) and
\(y\in F_{i+1,n}\),
\[
 \langle z,a_{i,n+1}v\rangle=\langle z',v\rangle,\qquad
 \langle d_{i,n}z',y\rangle
       =\kappa_{i,n}\langle d_{i,n+1}z,a_{i+1,n}y\rangle ,
\]
where \(\kappa_{i,n}\) is independent of \(z,v,y\).
Write \(S_i=w_0 H_i\) on \(W\) and
\(X_{i,n}=(\prod_{u<n}w_u)A_{i,n}j_i|_W\), where products
over \(u<n\) start at zero. Then \(S_i\) is positive definite
and self-adjoint. With the operators \(S_i,\ldots,S_{i+n-1}\)
in this order, for \(i+n\le m\) one has
\[
 \begin{split}
  \langle X_{i,n}x,X_{i,n}y\rangle
   &= L_{i,n}\frac{w_0w_1^n}{w_n}
       \langle x,S_i\cdots S_{i+n-1}y\rangle,\\
  L_{i,n}&=\prod_{h=0}^{n-2}\;
              \prod_{u=0}^{n-h-2}
                    \kappa_{i+h,u}\beta_{i+h,u}.
 \end{split}
\]
The scalar product \(L_{i,n}\) has exactly \(\binom n2\) factors
of the form \(\kappa\beta\). All empty products are identities.
\end{theorem}
\begin{proof}
Putting \(z=a_{i,0}j_i x\) in the first pairing identity gives
\(\langle a_{i,0}j_i x,a_{i,0}j_i y\rangle=\langle H_i x,y\rangle\).
Its conjugate symmetry gives self-adjointness of \(H_i\).
If its quadratic form is zero then \(a_{i,0}j_i x=0\), hence
\(H_i x=0\) by Definition~\ref{def:row-triangle-operators};
thus \(x=0\). Multiplication by \(w_0>0\) proves the assertion
about \(S_i\).

For \(i+n+1\le m\), \(x\in W\) and \(z\in F_{i,n+1}\) one has
\[
 \langle z,A_{i,n+1}j_i x\rangle
       =\left(\prod_{u<n}\kappa_{i,u}\right)
          \langle d_{i,n}z,A_{i+1,n}j_{i+1}x\rangle .
\]
The assertion at \(n=0\) is the first pairing identity.
For the induction apply the first equality involving \(z'\)
to \(v=A_{i,n+1}j_i x\), then apply the formula for \(n\)
with \(z'\), and use the second equality involving \(z'\)
on \(A_{i+1,n}j_{i+1}x\). This gives the formula with factor
\(\kappa_{i,n}\) appended.

Put \(z=A_{i,n+1}j_i x\) and use \(y\) for the argument in
the second slot. By
Theorem~\ref{thm:one-sided-row-triangle-injectivity},
\(d_{i,n}z=(\prod_{u<n}t_{i,u})^{-1}
 A_{i+1,n}j_{i+1}H_i x\).
The conjugate of the inverse product, from the first slot, is
\((w_0w_1)/(w_nw_{n+1})\prod_{u<n}\beta_{i,u}\);
this follows by cancellation of consecutive positive weights,
including \(n=0\). Multiplying by the squared weights of
\(X_{i,n+1}\), and accounting for those already in \(X_{i+1,n}\),
gives
\[
 \langle X_{i,n+1}x,X_{i,n+1}y\rangle
   =\frac{w_1w_n}{w_{n+1}}
       \left(\prod_{u<n}\kappa_{i,u}\beta_{i,u}\right)
             \langle X_{i+1,n}S_i x,X_{i+1,n}y\rangle .
\]
Here the factor \(1/w_0\) from \(H_i=S_i/w_0\) is included.
For length zero the norm identity is that \(j_i\) is isometric.
For positive length iterate this recurrence, moving \(S_i\)
from the first slot to the second by its self-adjointness after
applying the formula for the shorter interval. The positive
weights telescope to \(w_0 w_1^n/w_n\), and the numbers of
new square factors are \(n-1,n-2,\ldots,1\). Their sum is
\(\binom n2\).
\end{proof}

\begin{corollary}\label{cor:first-arrow-grams-commute}
Under the hypotheses of Theorem~\ref{thm:exact-triangle-gram-product}
all \(S_i\) commute pairwise and \(L_{i,n}>0\). In an orthonormal
basis of \(W\), with \(d=\dim W\),
\[
       \det(X_{i,n}^*X_{i,n})
       =\left(L_{i,n}\frac{w_0w_1^n}{w_n}\right)^{d}
                                  \prod_{h=0}^{n-1}\det S_{i+h}.
\]
\end{corollary}
\begin{proof}
The Gram form on the left of the theorem is positive definite by
Theorem~\ref{thm:one-sided-row-triangle-injectivity}. We prove
commutation of \(S_i,S_j\) by induction on \(j-i>0\).
The intermediate product \(M=S_{i+1}\cdots S_{j-1}\) (the identity
for adjacent indices) consists of commuting positive definite
matrices by the shorter distances. Each factor commutes with both
endpoints by the same induction. The product \(S_i M\) is positive
definite: simultaneously diagonalize its commuting factors by
decomposing into the eigenspaces of one Hermitian matrix and
repeating on each preserved eigenspace, as proved in
Lemma~\ref{lem:common-hecke-interpolation}.
In an orthonormal eigenbasis of \(S_i M\) its positive eigenvalues
multiply the strictly positive diagonal entries of \(S_j\) in
the trace of \(S_i M S_j\). This trace is real and positive.
The interval Gram matrix is a scalar multiple of this matrix;
its trace is also real and positive. Its scalar, and thus \(L\)
on this interval, is positive real. The same matrix is consequently
self-adjoint, giving \(S_i M S_j=S_j M S_i\). Cancel the invertible
commuting \(M\) to obtain \(S_iS_j=S_jS_i\). This proves the
induction; at lengths zero and one the empty scalar product is one.
Taking determinants in Theorem~\ref{thm:exact-triangle-gram-product}
now gives the formula.
\end{proof}



\section{Cancellation in the comparison of chain determinants}

\begin{proposition}\label{prop:actual-level-cancelled-gram}
Fix one twist and a nonzero finite-dimensional Hilbert space \(W\) of
bounded initial forms. Suppose its scalar forward map with terminal
test \(f^\natural\) in Proposition~\ref{prop:last-compact-projection-on-test}
is \(X_{0,m}\) of Theorem~\ref{thm:exact-triangle-gram-product}
as a map with the indicated terminal Petersson inner product and
\(w_n=W_n(P)\) of Theorem~\ref{thm:explicit-terminal-level-index}.
With \(d=\dim W\), let \(Z_\vartheta\) be the determinant of the
terminal class Gram matrix of \(U_\vartheta\) in an orthonormal
basis of \(W\). Then
\[
 Z_\vartheta=
    \left(\frac{d_Y C}{b_H}Q^{-k_m}\delta(P)^m L_{0,m}\right)^d
      \prod_{i=0}^{m-1}\det S_i>0.
\]
Suppose also that \(W_i^+\) has dimension \(d\), that
\(v_{i,\vartheta}:W_i^+\to H_i^+\) is a linear map to a Hilbert
space, and that with nonzero complex \(c_i,a_i\) and endomorphisms
\(T_i,T_i^+\) of \(W,W_i^+\), respectively,
\[
   S_i=c_i T_i,\qquad
   v_{i,\vartheta}^*v_{i,\vartheta}=a_i\delta(P) T_i^+,\qquad
   \det T_i=\det T_i^+ .
\]
Writing \(\lambda_i(\vartheta)=\det(v_{i,\vartheta}^*v_{i,\vartheta})\),
one has \(\lambda_i(\vartheta)>0\) and
\[
 Z_\vartheta=|c(\vartheta)|^d\prod_{i=0}^{m-1}\lambda_i(\vartheta),
 \qquad
 |c(\vartheta)|=
   \frac{d_Y C}{b_H}Q^{-k_m} L_{0,m}
                      \prod_{i=0}^{m-1}|c_i/a_i|.
\]
\end{proposition}
\begin{proof}
By Corollary~\ref{cor:actual-terminal-weighted-gram} the class
Gram determinant equals \((d_Y I_m(P)/b_H)^d\) times the
Petersson Gram determinant. Use
Corollary~\ref{cor:first-arrow-grams-commute} to compute the
second determinant. By Theorem~\ref{thm:explicit-terminal-level-index},
\[
    I_m(P)\frac{W_0(P)W_1(P)^m}{W_m(P)}
                  = C Q^{-k_m}\delta(P)^m .
\]
Each \(S_i\) is positive definite, and \(L_{0,m}>0\),
proving the first equality and positivity. In the second
assertion \(\det T_i\ne0\) because \(S_i\) is invertible. Hence
\(v_{i,\vartheta}^*v_{i,\vartheta}\) is a positive semidefinite
matrix with nonzero determinant, so it is positive definite.
Taking absolute determinants in the hypotheses gives
\[
       \delta(P)^d\det S_i
                 = |c_i/a_i|^d\lambda_i(\vartheta).
\]
Substitute in the first equality. The determinant equality
hypothesis follows in particular from equality of characteristic
polynomials by taking constant coefficients, even when those
operators act on different spaces.
\end{proof}

\begin{corollary}\label{cor:exact-square-count-height}
Suppose the data of Proposition~\ref{prop:actual-level-cancelled-gram}
are given for the trivial twist and for twists
\(\vartheta_\iota\) indexed by the places
\(\iota=(p,j)\) of Definition~\ref{def:place-height}.
Assume \(p\nmid Q(\vartheta_\iota)\), and that \(d_Y,C,b_H\)
and the product of the absolute ratios \(\prod_i |c_i/a_i|\)
are common to these twists. For each square of
Theorem~\ref{thm:exact-triangle-gram-product} set
\(D_{h,u}(\vartheta)=(\kappa_{h,u}\beta_{h,u})(\vartheta)\).
Suppose
\[
 {\rm ht}\!\left(
     |D_{h,u}(\vartheta_\iota)/D_{h,u}(1)|^{-2}\right)
                         \le E_{h,u}\qquad(h+u+2\le m).
\]
Then the scalar family in the volume factorization satisfies
\[
 {\rm ht}\!\left(|c(\vartheta_\iota)/c(1)|^{-2}\right)
          \le\sum_{h+u+2\le m} E_{h,u},
\]
where \(h,u\ge0\). In particular a uniform nonnegative per-square
bound \(E_*\) gives \(\binom m2 E_*\).
\end{corollary}
\begin{proof}
All square factors in \(L_{0,m}\) are nonzero by
Corollary~\ref{cor:first-arrow-grams-commute}. The trivial twist
has \(Q=1\). Substitute the formula of
Proposition~\ref{prop:actual-level-cancelled-gram}; after division
the common factors cancel. The inverse squared modulus is the
product of the indicated families and \(Q(\vartheta_\iota)^{2k_m}\).
This rational power is a \(p\)-unit at \(\iota=(p,j)\) by the
coprimality hypothesis. The constant witnesses \(h=1,A=0,e_p=0\)
give that family height at most zero in
Definition~\ref{def:place-height}. Apply
Proposition~\ref{prop:height-product} to all factors and count
the square indices as in Theorem~\ref{thm:exact-triangle-gram-product}.
\end{proof}



\section{Finite matrix coefficients for the first arrows}

\begin{theorem}\label{thm:first-arrow-finite-coefficient}
Use a hermitian source and target of ranks \(r,h\) with finite
isometric frames for the standard lattice, \(h\ge r\), and the
finite arrow of Theorem~\ref{thm:global-actual-arrow-covariance}.
Write \(\chi_H\) for the character in its denominator on the source
determinant. Suppose its unit character is a product of residue
characters at \(M=M_g\cup P\). Set \(D_M=\prod_{w\in M}q_w\).
In the split coordinates divided by \(\alpha\) at \(M\) choose a
complex function \(F\) on \(\prod_{w\in M}M_{h,r}(\mathbf F_w)\).
Take its extension by \(F(\bar X)\) on integral matrices and by
zero elsewhere. Take the unit vacuum at all other pairs and
nonsplit places. Denote the resulting finite vector by \(v_F\).
Assume that for each \(w\in M\) a fixed \(r\)-by-\(h\) residue
matrix \(L_w\) satisfies \(L_w \bar X_w=I_r\) on the support
of \(F\).

For \(g\) in the finite source group let \(g_*\) be \(g\) away
from \(M\cup\bar M\) and the identity on these pairs. In the
source frame let \(\mathfrak a(g_*)\) be the integral ideal whose
valuation at \(w\) is the determinant valuation of the row lattice
\[
      \{v\in O_w^{1\times r}:v(g_*)_w^{-1}\in O_w^{1\times r}\}.
\]
Define
\[
 \rho_h(g_*)=
 \frac{\beta(\mathfrak a(g_*))^h}
      {\chi_H(1,\det g_*)\, (N\mathfrak a(g_*))^{h/2}} .
\]
The matrix coefficient \(\langle v_F,\Omega_f(g,1)v_F\rangle\)
vanishes unless \(g_w\in I_r+\mathfrak p_w M_r(O_w)\) at
every \(w\in M\). When this condition holds it equals
\[
      \rho_h(g_*)\,\|v_F\|^2,\qquad
      \|v_F\|^2=D_M^{-hr}\sum_{\bar X}|F(\bar X)|^2 .
\]
\end{theorem}
\begin{proof}
The isometry of Theorem~\ref{thm:split-finite-schroedinger}
gives the coordinate description at the split pairs, with integral
lattice of mass one. Reduction there has \(D_M^{hr}\) equally
measured fibers and the other vacuum factors have norm one,
proving the stated norm. The row and column action in
Corollary~\ref{cor:split-moved-indicator} changes \(F\) on this
set to a scalar times the function evaluated at \(Xg_w\) in
component \(w\). The coordinates here include all the chosen
places simultaneously; no factorization of \(F\) among them is
needed. If the two supports intersect, \(X_w\) and \(X_wg_w\)
are integral with their indicated reductions. Lift \(L_w\)
entrywise to \(\widetilde L_w\) over \(O_w\). Both
\(\widetilde L_wX_w\) and \(\widetilde L_w X_wg_w\) are integral
and congruent to \(I_r\). The adjugate formula gives an integral
inverse to the first and hence \(g_w\equiv I_r\bmod\mathfrak p_w\).
This proves disjointness of the supports if the condition fails.

When the condition holds, the action at the chosen pairs fixes
the entire vector: the matrix substitution and its inverse
preserve the support and its reduction, the Jacobian is one,
the ideal multiplier has valuation zero there, and the
determinant unit characters are trivial on principal units.
These statements follow from Theorems~\ref{thm:split-finite-schroedinger}
and~\ref{thm:all-finite-weil-splitting} and
Definition~\ref{def:finite-splitting-multiplier}.
For \(g_*\) the lattice map is right multiplication by \(g_*^{-1}\)
on \(h\) rows. Its saturation lattice \(S_w(G)\) in that
definition is the direct sum of \(h\) copies of the displayed
row lattice. Hence \(J_G=\mathfrak a(g_*)^h\). Also
\([GL_0:GL_0\cap L_0]=[L_0:L_0\cap G^{-1}L_0]=N J_G\):
locally diagonalize the integral sublattice over the discrete
valuation ring as in Lemma~\ref{lem:row-spherical-hecke}, then
multiply the residue cardinalities raised to the diagonal
valuations. In the vacuum average of
Theorem~\ref{thm:finite-oscillator-lattice-transport} only the
identity coset pairs with the vacuum, by disjoint support.
Its coefficient is \((N J_G)^{-1/2}\). The full finite action
adds \(\beta(J_G)\) and the inverse determinant character by
Theorem~\ref{thm:global-actual-arrow-covariance}. The chosen
pairs and their complement split as tensor factors because the
lattice, pairing and actions split there, as in
Theorem~\ref{thm:split-finite-schroedinger}. Multiplication gives
the desired formula.
\end{proof}

\begin{corollary}\label{cor:first-arrow-level-masses}
For \(h=r+t\) let \(F_{\rm row}\) be the inclusion indicator
at all \(w\in M\). Let \(F_{\rm aux}\) equal \(A_r\) times the
indicator of matrices reducing to the inclusion at \(M_g\)
and having first \(r\) rows \(I_r\) at \(P\).
The two norm squares are, respectively,
\[
       D_M^{-hr},\qquad D_M^{-hr} A_r^2 Q^{rt}.
\]
Both finite matrix coefficients on the subgroup integral off
\(M\cup\bar M\) are supported on the principal source level,
and on it they equal their respective norm squares.
The ratio of the auxiliary to row finite masses, each obtained
by dividing the norm square by the principal source level
index \(W_0(P)\), is \(A_r^2 Q^{rt}\).
Multiplying the auxiliary norm square divided by \(W_0(P)\)
by its target index \(I_{\rm aux}(P)\) gives \(D_M^{-hr}\delta(P)\).
\end{corollary}
\begin{proof}
Use the projection onto the top \(r\) rows for the common
left inverse. At \(M_g\) there is a single permitted matrix.
At each \(w\in P\) the auxiliary lower \(t\)-by-\(r\) block
is arbitrary, giving \(Q^{rt}\) permitted tuples. The norm
and coefficient assertions follow from
Theorem~\ref{thm:first-arrow-finite-coefficient}; for an element
integral off the chosen pairs the complementary ideal is the
unit ideal and the complementary determinant character is
trivial by its prescribed unit parts. Reduction onto the full
product on the source was proved in
Theorem~\ref{thm:actual-terminal-parabolic-action}, so the
index is \(W_0(P)\). The last claim is the definition of
\(\delta(P)\) in Theorem~\ref{thm:explicit-terminal-level-index}.
\end{proof}



\section{The nonidentity finite cosets of a first arrow}

\begin{lemma}\label{lem:two-index-local-coset-count}
Let \(O\) be the integer ring of a nonarchimedean local field \(L\)
with uniformizer \(\pi\) and residue size \(q\), and
\(M=O^{1\times r}\). For another lattice \(M'\) of rank \(r\)
write \(q^a=[M':M\cap M']\), \(q^b=[M:M\cap M']\).
For fixed \(a,b\ge0\) the number of possibilities for \(M'\)
is at most \(q^{r(a+b)}\). The lattices \(Mg^{-1}\) parametrize
right cosets \(g\,{\rm GL}_r(O)\) in \({\rm GL}_r(L)\).
Replacing \(g\) by its inverse transpose interchanges its two
lengths \(a,b\); applying an isomorphism of valued fields to
all the entries preserves them.
\end{lemma}
\begin{proof}
Diagonal reduction over \(O\), after clearing a denominator as
in Lemma~\ref{lem:row-spherical-hecke}, gives elementary powers
\(\pi^{u_1},\ldots,\pi^{u_r}\) for \(g^{-1}\), up to integral
invertible factors. Thus \(a=\sum\max(0,-u_i)\) and
\(b=\sum\max(0,u_i)\). Negating all the \(u_i\) interchanges the
lengths. Transposition preserves the powers by transposing the
diagonalization. A valued-field isomorphism preserves uniformizers
up to units, proving the last claims. Every lattice has a basis
over \(O\) by the same diagonal reduction, so is \(Mg^{-1}\).
Two matrices give the same row lattice precisely when they differ
on the left by an integral invertible row operation, which is
right coset equivalence on their inverses.

Every sublattice of length \(b\) is reached through \(b\)
successive index-\(q\) sublattices: in diagonal coordinates
increase the nonnegative exponents of the desired sublattice
one unit at a time. At each step the choices are the hyperplanes
modulo \(\pi\), because a quotient of size \(q\) of \(O\)-length
one is its residue field. Equivalently each step in this construction
is the kernel of a nonzero residue linear form. There are
\((q^r-1)/(q-1)\le q^r\) such hyperplanes, counting the linear
forms up to nonzero scalar. Every superlattice of length \(a\)
is likewise reached through \(a\) steps. The superlattices of
length one of \(M_0\) correspond to lines in
\(\pi^{-1}M_0/M_0\), with the same bound. Choose first the
intersection with \(M\), then a superlattice of this intersection.
The product of the two bounds is \(q^{r(a+b)}\).
\end{proof}

\begin{theorem}\label{thm:finite-vacuum-tail-first-arrow}
Use the standard finite unitary source frames of
Theorem~\ref{thm:first-arrow-finite-coefficient}. Let \(G_*\) be
the restricted product of its local unitary groups off the chosen
pairs \(M\cup\bar M\) and \(K_*\) their product of integral
stabilizers. Use the function \(\rho_h\) of that theorem on \(G_*\),
with \(\chi_H\) trivial on its determinant units, and put
\(a_0=h/2-2r\). If \(a_0\ge2\), then
\[
 \sum_{gK_*\in G_*/K_*,\ gK_*\ne K_*}|\rho_h(g)|
                \le \exp(24 f\,2^{-a_0})-1 .
\]
In particular, for the first arrows with ranks \(r,h=r+t=6r\)
on the arithmetic windows this is \(O(f\,2^{-r})=o(1)\),
uniformly in the twist.
\end{theorem}
\begin{proof}
By Theorems~\ref{thm:first-arrow-finite-coefficient},
\ref{thm:actual-global-splitting-character} and
\ref{thm:torus-characters-elementary}, the absolute weight is
\((N\mathfrak a(g))^{-h/2}\). The ideal and this weight are
constant on integral double cosets by the row lattice definition
and integral changes of basis. Each local coset determines the
lattice of Lemma~\ref{lem:two-index-local-coset-count}. At a
split pair the group is the full general linear group on one
branch. The other matrix is its conjugate inverse transpose
in the isometric frame, so its length \(a\) is the first
branch's \(b\). The contribution of this pair to
\(N\mathfrak a(g)\) is \(q^{a+b}\), where \(q\) is the common
residue size. The sum of absolute local weights, including the
identity coset, is at most
\((1-q^{-(h/2-r)})^{-2}\), by summing two geometric series.

At a nonsplit place let \(q_E\) denote the residue size in \(E\).
The map from the unitary cosets to row lattices is injective,
since an integral general linear matrix belonging to the unitary
group lies in its integral stabilizer. The determinant of a
unitary has valuation zero: its product with its conjugate is
one and the place is fixed by conjugation. Hence \(a=b\) by
the elementary-power calculation in
Lemma~\ref{lem:two-index-local-coset-count}. The local contribution
to \(N\mathfrak a(g)\) is now \(q_E^a\). Summing the bound
\(q_E^{2ra} q_E^{-ha/2}\) gives at most
\((1-q_E^{-a_0})^{-1}\).

Above a rational prime \(p\) there are at most \(2f=[E:\Q]\)
places of \(E\), by the prime decomposition and degree formula
of Proposition~\ref{prop:general-ideal-factorization}, or by
counting the factors of \(O_E/pO_E\), an algebra of
dimension \(2f\) over \(\mathbf F_p\). Each residue size is
at least \(p\). For \(0\le y\le1/4\),
\(-\log(1-y)=\int_0^y(1-u)^{-1}\dd u\le2y\).
Both local logarithmic bounds are therefore at most
\(4p^{-a_0}\). To sum over all rational primes use
\[
 \sum_p p^{-a_0}\le 2^{-a_0}+\int_2^\infty u^{-a_0}\dd u
                  \le 3\,2^{-a_0}.
\]
Multiplying the local bounds over finite sets of places gives
at most \(\exp(24f\,2^{-a_0})\). Every coset has nonidentity
components at finitely many places, and the summands are
nonnegative. Increasing the finite sets exhausts their sum,
whose identity term is one. This proves the asserted tail bound.
For \(h=6r\) apply Proposition~\ref{prop:thick}:
\(r\) grows on the order of \(x^5\), whereas
\(f=O_M(x^{12})\). Thus \(f2^{-r}\to0\), and the exponential
series gives the final estimate uniformly.
\end{proof}



\section{Uniform positivity of the finite first-arrow operator}

\begin{theorem}\label{thm:positive-finite-first-arrow-convolution}
Use \(G_*,K_*,h,r\) of Theorem~\ref{thm:finite-vacuum-tail-first-arrow}.
Let \(\pi\) be a strongly continuous unitary representation of \(G_*\)
on a complex Hilbert space \(\mathcal H\). For \(K_*\)-fixed \(y\)
put
\[
         T_hy=\sum_{gK_*\in G_*/K_*}\rho_h(g)\pi(g)y,\qquad
         \eta_h=\exp(24 f\,2^{-h/2+2r})-1 .
\]
If \(h/2-2r\ge2\), the series converges absolutely in Hilbert
norm. It gives a self-adjoint operator on \(\mathcal H^{K_*}\)
with \(\|T_h-I\|\le\eta_h\). For \(\eta_h<1\) one has
\[
 (1-\eta_h)\|y\|^2\le\langle y,T_h y\rangle
                      \le(1+\eta_h)\|y\|^2,\qquad
              \|T_h^{-1}\|\le(1-\eta_h)^{-1}.
\]
If \(A\) is a \(d\)-dimensional subspace invariant under \(T_h\),
then \((1-\eta_h)^d\le\det(T_h|_A)\le(1+\eta_h)^d\) in this range.
\end{theorem}
\begin{proof}
With the vacuum of norm one at the omitted complements in
Theorem~\ref{thm:first-arrow-finite-coefficient}, \(\rho_h\) on
\(G_*\) is precisely the matrix coefficient of the vacuum away
from the chosen pairs, evaluated in the full unitary finite action.
The vacuum is fixed by \(K_*\), since the matrices preserve its
lattice and the determinant unit characters are trivial there.
It follows that \(\rho_h\) is bi-\(K_*\)-invariant,
\(\rho_h(1)=1\), and \(\rho_h(g^{-1})=\overline{\rho_h(g)}\)
by unitarity. The vector summands are well defined on the stated
right cosets. Reindexing by \(g\mapsto kg\) for \(k\in K_*\)
fixes their sum. The sum of the norms off the identity coset
is at most \(\eta_h\|y\|\) by
Theorem~\ref{thm:finite-vacuum-tail-first-arrow}. This gives
absolute convergence, invariance, and the asserted norm estimate.

Give Haar measure on \(G_*\) mass one on \(K_*\).
It is invariant under inversion: on each local unitary factor,
conjugation has Jacobian of modulus one on the unitary Lie
algebra by the trace-pairing and local-coordinate argument
of Proposition~\ref{prop:marked-poincare-kernel}.
The assertion passes to finite products and then the restricted
product by integration of compactly supported functions on finite
sets of factors times the remaining integral groups. On fixed
vectors the sum equals integration against \(\rho_h(g)\),
by its absolute convergence and the coset volumes. Inverting
\(g\), using unitarity of \(\pi\) and the preceding identity
for \(\rho_h\), interchanges the two arguments of the associated
Hermitian pairing with conjugation. This is self-adjointness.

Cauchy--Schwarz bounds
\(|\langle y,(T_h-I)y\rangle|\le\eta_h\|y\|^2\).
For \(\eta_h<1\), the geometric sum
\(\sum_{n\ge0}(I-T_h)^n\) converges in operator norm by summation
of \(\eta_h^n\); multiplication on either side by \(T_h\)
and telescoping gives the identity. The inverse estimate follows.
On a finite-dimensional invariant subspace diagonalize the
self-adjoint restriction by
Lemma~\ref{lem:common-hecke-interpolation} and use the bounds on
each unit eigenvector. Their product is its determinant.
\end{proof}

\begin{corollary}\label{cor:first-arrow-injectivity-from-norm-identity}
At \(h=6r\) and all sufficiently large arithmetic scales let
\(W\subset\mathcal H^{K_*}\) be a nonzero subspace and \(v:W\to H'\)
a linear map to a complex Hilbert space. If
\[
       \langle vx,vy\rangle=A\langle x,T_h y\rangle
               \qquad(x,y\in W),\qquad A\in\C^\times ,
\]
then \(A>0\) and \(v\) is bounded below by
\(\sqrt{A(1-\eta_h)}\) on \(W\).
If two finite-dimensional spaces of this kind have equal
characteristic polynomials for every pair of corresponding
finite double-coset convolutions and their finite linear
combinations, their \(T_h\) have the same characteristic
polynomial, provided the same \(\rho_h\) is used and both spaces
are invariant under those convolutions.
\end{corollary}
\begin{proof}
Theorem~\ref{thm:finite-vacuum-tail-first-arrow} gives
\(\eta_h<1\) eventually. At any nonzero \(y\), the quadratic
form of \(T_h\) is then strictly positive by
Theorem~\ref{thm:positive-finite-first-arrow-convolution}.
The identity at \(x=y\) and \(A\ne0\) imply that \(\|vy\|^2>0\)
and \(A>0\). The lower bound follows from the same identity.

Each integral double coset is a finite union of right cosets:
\(K_*\cap gK_*g^{-1}\) is open in the compact group \(K_*\),
so it has finite index by an open coset cover. Take finite
unions of double cosets exhausting \(G_*\); one may enumerate
the cosets because local lattices of bounded indices are finite
by Lemma~\ref{lem:two-index-local-coset-count}, the places are
countable, and only finitely many local components are nonintegral.
The truncated sums converge in operator norm on both fixed spaces
by the same absolute tail estimate. In a fixed basis each
coefficient of a characteristic polynomial is a polynomial in
the entries, by determinant expansion, hence continuous. Passing
to the limit in the polynomial equality gives the claim.
\end{proof}



\section{Graph terms of the doubled theta kernel}

\begin{definition}\label{def:doubled-graph-zero-term}
Let \(S,T\) be nonsingular hermitian \(E\)-spaces of positive
ranks \(r,h\), and put \(V={\rm Hom}_E(S,T)\) with the trace
hermitian pairing of Definition~\ref{def:hom-oscillator-pairing}.
Set \(D=S\oplus S^{-}\), using the negative form on the second
copy. For an \(r\)-dimensional totally isotropic rational
\(E\)-subspace \(A\subset D\) put
\[
   \mathcal N_A=\{X\in{\rm Hom}_E(D,T): X|_A=0\}\subset V\oplus V^- .
\]
For Schwartz tests \(f_1,f_2\) in the rational Schr\"odinger
realization of Theorem~\ref{thm:rational-schroedinger-zero-section}
write \(F_{f_i}\) for the corresponding theta sections. The
zero term of their doubled tensor along \(A\) is
\[
       Z_A(f_2,f_1)=
       \int_{\mathcal N_A(E)\backslash\mathcal N_A(\mathbb A_E)}
             F_{f_1}(v_1)\overline{F_{f_2}(v_2)}\,\dd(v_1,v_2),
             \qquad X=(v_1,v_2).
\]
The quotient has its additive probability measure.
\end{definition}

\begin{lemma}\label{lem:graph-zero-term-oscillator-pairing}
The space \(\mathcal N_A\) is maximal totally isotropic in
\({\rm Hom}(D,T)\) and the defining integral exists. If
\(A=\{(x,-\gamma x):x\in S\}\) for a rational isometry
\(\gamma\) of \(S\), then
\[
           Z_A(f_2,f_1)=\langle f_2,\Omega(\gamma,1)f_1\rangle
\]
with the oscillator inner product of the rational
Schr\"odinger model.
\end{lemma}
\begin{proof}
The hermitian identification of \(D\) with the conjugate dual
carries \(A=A^\perp\) onto the linear covectors killing \(A\);
indeed \(A\subset A^\perp\) and both have dimension \(r\)
by nonsingularity. Thus \(\mathcal N_A\) is spanned by maps
\(z\mapsto b\langle a,z\rangle_D\), \(a\in A,b\in T\).
Two such maps have trace pairing a product with the factor
\(\langle a,a'\rangle_D\), hence zero. Its dimension is
\(hr\), half that of \({\rm Hom}(D,T)\), so it is maximal.
On it the two bicharacters in the product section cancel.
The rational translation law of
Definition~\ref{def:adelic-oscillator-translations} then gives
an ordinary continuous function on the compact additive quotient
of Lemma~\ref{lem:theta-positive-translation}. This proves
existence.

For the displayed graph the condition on \(X\) reads
\(v_1=v_2\gamma\).
The additive rational parametrization \(v\mapsto(v\gamma,v)\)
identifies quotient probabilities. The rational covariance of
Theorem~\ref{thm:actual-adelic-rational-theta-invariance} gives
\(F_{f_1}(v\gamma)=F_{\Omega(\gamma,1)f_1}(v)\):
the Hom action of \((\gamma,1)\) is right multiplication by
\(\gamma^{-1}\), and rational determinant characters are one
by Theorem~\ref{thm:global-actual-arrow-covariance}.
The norm identity in the proof of
Theorem~\ref{thm:rational-theta-symplectic-action} identifies
the integral of \(|F_f|^2\) on the additive quotient with
\(\|f\|^2\). Polarization gives the pairing identity, since
the map \(f\mapsto F_f\) is complex linear.
\end{proof}

\begin{theorem}\label{thm:marked-test-kills-nongraph-zero-terms}
If \(S\) is definite at some infinite place, every rational
\(A\) in Definition~\ref{def:doubled-graph-zero-term} is the
graph of a unique isometry in the convention of
Lemma~\ref{lem:graph-zero-term-oscillator-pairing}. Without the
definiteness assumption, suppose the first finite test at a
split place \(w\) in the realization of
Theorem~\ref{thm:split-finite-schroedinger} is supported on
injective \(h\)-by-\(r\) matrices over \(E_w\).
Then \(Z_A(f_2,f_1)=0\) for every nongraph \(A\).
The assertion applies to either first-arrow residue test of
Corollary~\ref{cor:first-arrow-level-masses}, also after arbitrary
source and target adelic isometries act on the tests.
\end{theorem}
\begin{proof}
If \(A\cap(S\oplus0)=0\), the second projection is a
bijection onto \(S\) by dimension. Thus
\(A=\{(c y,y):y\in S\}\). Isotropy gives
\(\langle c y,c y'\rangle_S=\langle y,y'\rangle_S\).
Nonsingularity makes \(c\) injective and hence bijective,
showing that \(A\) is the indicated graph with
\(\gamma=-c^{-1}\). The graph determines \(\gamma\) uniquely.
If \(S\) is definite at one place, every rational \((x,0)\in A\)
has \(\langle x,x\rangle_S=0\) there and hence \(x=0\).
This proves the first assertion.

For a nongraph \(A\) choose \(x\ne0\) with \((x,0)\in A\).
At the chosen split place choose a matrix \(C\) of a basis of
\(A_w\), with first column \((x,0)\). In the local matrix model
write \(X=(X_1,X_2)\). The compact support of a local component
of the doubled test has \(XC\ne0\), since its first column is
\(X_1x\ne0\). We explain why its zero term annihilates such
a test. Under \(\mathcal N_A(\mathbb A_E)\)-translations the zero
term is invariant, by translation on its additive quotient and
triviality of the bicharacter on this subspace, exactly as in
Lemma~\ref{lem:theta-positive-translation}. At \(\bar w\)
these translations act by all the characters of \(XC\).
Indeed the pure opposite-branch action is modulation by the
perfect trace pairing on the matrix space at \(w\), by
Theorem~\ref{thm:split-finite-schroedinger}, with its sign
reversed on the conjugated second tensor. The local subspace
\((\mathcal N_A)_{\bar w}\) pairs trivially with
\(\{X:XC=0\}\) by total isotropy. Its dimension \(hr\)
equals the codimension of this kernel, since restriction to
a basis of \(A_w\) is surjective. It is therefore the whole
annihilator. The trace pairing identifies its parameters with
all linear residue characters in the entries of \(XC\).
The factor \(\alpha\) in the split coordinates may be absorbed
in the dual parameters by its invertibility.

The image \(XC\) of the compact support misses a sufficiently
small open lattice \(D'\) around zero: otherwise points of the
compact image would approach zero, so that zero belonged to it.
All values lie in a larger compact open lattice \(D''\).
Average the modulations associated with the characters of the
finite group \(D''/D'\), extending them by the local trace
characters as in Proposition~\ref{prop:moved-residue-coset-fourier}.
The average multiplies the test by the indicator of \(XC\in D'\)
on its support, by finite character orthogonality. It is zero,
whereas every translation in the average gives the same zero
term. This proves vanishing. A test with finitely many local
tensor components follows by linearity, or by taking the union
of their compact supports. The top \(r\) rows of both named
first-arrow tests reduce to \(I_r\), so those matrices are
injective. Isometries on the source and target act locally by
invertible right and left substitutions with scalar factors,
as in Corollary~\ref{cor:split-moved-indicator}, preserving
injectivity of all matrices in the moved support.
\end{proof}



\section{Absolute integrability at the real places}

\begin{lemma}\label{lem:real-unitary-matrix-ball-measure}
Let \(p,q\ge0,\ r=p+q>0\), \(G=U(p,q)\),
\(K=U(p)\times U(q)\), for the diagonal form \((I_p,-I_q)\).
Write \(\mathbb B=\{Z\in M_{p,q}(\C):I_q-Z^*Z>0\}\).
Then \(G/K=\mathbb B\), and, up to a fixed positive scalar,
Haar integration on \(G\) is integration of \(g_Z k\)
against
\[
       \det(I_q-Z^*Z)^{-r}\dd Z\,\dd k ,
\]
with Lebesgue measure in real and imaginary entries and
probability Haar measure on \(K\). If \(s_j\) are the singular
values of \(Z\), \(1\le j\le\min(p,q)\), and
\(\tanh u_j=s_j\), \(u_j\ge0\), then \(g_Z k\) is a product
of two compact factors surrounding the simultaneous real
two-by-two boosts with entries \(\cosh u_j,\sinh u_j\) on paired
positive and negative lines. The other lines are fixed by the
middle factor.
\end{lemma}
\begin{proof}
Every negative \(q\)-plane projects isomorphically to the negative
coordinate summand by definiteness of the other summand and the
dimension. Its graph has columns \(\binom ZI\), and is negative
precisely under the displayed matrix inequality. Put
\(D_Z=(I-Z^*Z)^{-1/2}\), \(A_Z=(I-ZZ^*)^{-1/2}\).
The matrices \(\binom{I}{Z^*}A_Z,\ \binom ZI D_Z\) give
orthonormal positive and negative columns orthogonal to each
other. Their concatenation \(g_Z\in G\) sends \(0\) to \(Z\).
The positive square roots exist by diagonalization as in
Lemma~\ref{lem:common-hecke-interpolation}; their dependence on
\(Z\) is continuous, for example by the power series for
\((1-u)^{-1/2}\) converging uniformly on compact subintervals
of \([0,1)\). The stabilizer of the negative coordinate plane
preserves its orthogonal complement, and is \(K\).
Singular-value coordinates follow by diagonalizing \(Z^*Z\),
taking \(Zv/s_j\) for the corresponding positive-side unit
vectors when \(s_j>0\), and completing both orthonormal bases.
The section on each pair has diagonal \(1/\sqrt{1-s_j^2}\)
and off-diagonal \(s_j/\sqrt{1-s_j^2}\), as required.

For \(g=\left(\begin{smallmatrix}A&B\\C&D\end{smallmatrix}\right)\),
the transformed graph has coordinate
\(Y=(AZ+B)(CZ+D)^{-1}\); the denominator is invertible by the
negative-plane projection argument. Taking the forms on the graph
columns gives
\(I-Y^*Y=(CZ+D)^{-*}(I-Z^*Z)(CZ+D)^{-1}\), where
\(L^{-*}=(L^*)^{-1}\). Differentiation gives
\(\dd Y=(A-YC)\dd Z(CZ+D)^{-1}\) as a linear differential
on the matrix variable. Block elimination in
\(g\left(\begin{smallmatrix}I&Z\\0&I\end{smallmatrix}\right)\)
gives \(\det(A-YC)\det(CZ+D)=\det g\) and \(|\det g|=1\).
Taking the determinants of left multiplication on each column and
right multiplication on each row gives the real Jacobian
\(|\det(CZ+D)|^{-2r}\). Hence the proposed density on the ball
is invariant by change of variables. In the coordinates \(g_Z k\),
left translation changes \(k\) by a left multiplier depending
on \(Z\). Its Haar probability is preserved. This proves left
invariance and the Haar formula up to the chosen normalization,
as in Lemma~\ref{lem:quotient-measure}. When either rank is zero
the ball has one point with mass one and all the formulas reduce
to compact Haar integration.
\end{proof}

\begin{theorem}\label{thm:archimedean-first-arrow-integrability}
At an infinite place let the source have signature \((p,q)\)
and the target have dimension \(h>0\) and any hermitian
signature. For any two Schwartz oscillator tests \(f_1,f_2\),
the coefficient of the actual unitary pair satisfies
\[
   \left|\langle f_2,\Omega_\infty(g,1)f_1\rangle\right|
          \le C(f_1,f_2)\det(I-Z^*Z)^{h/2},\qquad g=g_Z k.
\]
For the two unit Gaussian vacua of
Theorem~\ref{thm:negative-block-real-splitting} the two sides
are equal with \(C=1\). In particular the coefficient is
integrable on the source when \(h\ge2r\). The same holds on
the product of the infinite groups, including Schwartz tests
on their direct sum.
\end{theorem}
\begin{proof}
Absorb the two compact factors of
Lemma~\ref{lem:real-unitary-matrix-ball-measure} into the two
tests. Their Schwartz seminorms are uniformly bounded on those
compact groups by Theorem~\ref{thm:negative-block-real-splitting},
also after a fixed real change of polarization by
Proposition~\ref{prop:real-symplectic-gaussian-action}. The
additional determinant characters have modulus one by
Theorem~\ref{thm:global-actual-arrow-covariance}.

On a paired source plane the sum and difference of its two
basis vectors diagonalize the boost with eigenvalues \(e^{u_j}\)
and \(e^{-u_j}\). Each eigenline is hermitian isotropic and
they pair perfectly with each other. On the corresponding
Hom block the action is right multiplication by the inverse
source matrix. Take its \(h\) complex coordinates expanding
by \(e^{u_j}\) as real configuration coordinates, with the
opposite \(h\) complex coordinates as their momentum duals.
Their alternating pairing is the real trace of the hermitian
pairing as in Lemma~\ref{lem:rational-graded-polarization}.
This choice depends only on the fixed coordinate pairs and is
also valid at \(u_j=0\). On the remaining fixed symplectic
block take any complementary real Lagrangians, obtained by
successively splitting off paired basis vectors as in that lemma.
The boost in these coordinates is \((A,A^{-{\mathsf T}})\)
on position and momentum, with
\(|\det_\R A|=\exp(2h\sum_j u_j)\).
In the symmetric gauge its unitary implementer on the positions
is \(f(x)\mapsto|\det A|^{-1/2}f(A^{-1}x)\) up to a scalar of
absolute value one, by direct Heisenberg substitution and uniqueness
in Proposition~\ref{prop:real-symplectic-gaussian-action}.
The absolute pairing is therefore bounded by
\(\|f_2\|_1\|f_1\|_\infty\exp(-h\sum_j u_j)\), with the compactly
moved tests in these norms. Their bounds are uniform since a
sufficiently high Schwartz seminorm bounds the \(L^1\)-norm
by integration of \((1+\|x\|)^{-b}\) for \(b\) greater than
the real dimension. As \(e^{-u_j}\le(\cosh u_j)^{-1}\), the
upper bound follows.

For the Gaussian at each compact factor the negative-to-negative
block is unitary. Its vacuum coefficient has modulus one by
Theorem~\ref{thm:negative-block-real-splitting}; a unit vector
whose inner product with another unit vector has modulus one
differs by a unit scalar. On the boost its negative block in
the signed Hom coordinates has one factor \(\cosh u_j\) per
paired source plane and per target row of either sign; the
unpaired coordinates give ones. The coefficient in that theorem
therefore has absolute value
\(\prod_j(\cosh u_j)^{-h}=\det(I-Z^*Z)^{h/2}\).

Multiplication by the density of
Lemma~\ref{lem:real-unitary-matrix-ball-measure} leaves
\(\det(I-Z^*Z)^{h/2-r}\le1\) when \(h\ge2r\).
The ball lies in a bounded Euclidean box by
\(\|Z\|_{\rm op}<1\); hence its Lebesgue measure is finite.
For the product of infinite places choose the direct sum of
the preceding real polarizations. The determinant of its dilation
is the product of their determinants. The same \(L^1\)-supremum
bound for the two Schwartz functions on this sum yields the
product of the determinant bounds, even if the tests are not
pure tensors. Product integration finishes the proof.
\end{proof}



\section{Unfolding the graph contribution to the doubled kernel}

\begin{theorem}\label{thm:absolute-doubled-graph-unfolding}
Let \(G=U(S)(\mathbb A_{E^+})\), \(\Gamma=U(S)(E^+)\),
\([G]=\Gamma\backslash G\), and use compatible Haar and quotient
measures, with the finite integral stabilizer of mass one.
Take ranks \(r,h\) with \(h/2-2r\ge2\) and \(M_g\ne\varnothing\).
Use the same nonzero residue test \(v_F\) of
Theorem~\ref{thm:first-arrow-finite-coefficient} on both sides,
and arbitrary real Schwartz factors \(f_\infty,f'_\infty\).
Write \(f=f_\infty\otimes v_F,\ f'=f'_\infty\otimes v_F\) and
\[
    \Phi(g)=\langle f',\Omega(g,1)f\rangle,\qquad
    \Phi_\infty(a)=\langle f'_\infty,\Omega_\infty(a,1)f_\infty\rangle.
\]
Then \(\Phi\in L^1(G)\). For \(y,y'\in L^2([G])\) the integral
of the sum
\[
 \int_{[G]^2}\overline{y'(g')}y(g)
        \sum_{\gamma\in\Gamma}\Phi(g'^{-1}\gamma g)\,\dd g\,\dd g'
       =\int_G\Phi(u)\langle y',R(u)y\rangle_{[G]}\,\dd u
\]
is absolutely convergent, including termwise absolute values
on the left. The sum inside the pairing is the sum of the
rational zero terms of Definition~\ref{def:doubled-graph-zero-term}
on \(\Omega(g',1)f',\Omega(g,1)f\), with equality wherever
the absolute graph series converges.

If \(y\) is fixed at finite places by the principal source
level at \(M\) and the integral level elsewhere, the right side
also equals
\[
   \frac{\|v_F\|^2}{W_0(P)}
       \int_{G_\infty}\!
       \Phi_\infty(a)\,
          \langle y', R(a) T_h y\rangle_{[G]}\,\dd a .
\]
Here \(T_h\) is the operator of
Theorem~\ref{thm:positive-finite-first-arrow-convolution}
on the finite factors off \(M\cup\bar M\).
\end{theorem}
\begin{proof}
The real and finite actions on tensor tests factor by
Theorem~\ref{thm:global-actual-arrow-covariance}.
The real coefficient is integrable by
Theorem~\ref{thm:archimedean-first-arrow-integrability}.
At the finite places the coefficient is supported on the
principal subgroup of the chosen pairs, with value
\(\|v_F\|^2\rho_h(g_*)\) there, by
Theorem~\ref{thm:first-arrow-finite-coefficient}. Its absolute
integral is finite by summation of integral cosets in
Theorem~\ref{thm:finite-vacuum-tail-first-arrow}.
This proves absolute integrability of their product.

Take a measurable fundamental domain for \(\Gamma\) on \(G\).
The construction of Lemma~\ref{lem:quotient-measure} applies to
the adelic group: a compact set has bounded denominators at
finite places and bounded real matrix entries, so meets only
finitely many rational matrices by
Lemma~\ref{lem:number-field-lattice}; remove earlier rational
translates from a countable cover by neighborhoods with disjoint
nontrivial translates. For the nonnegative majorant with
\(|\Phi|,|y|,|y'|\), sum the translates of the inner domain by
\(\gamma\). Rational invariance of \(y\) gives integration on
\(G\), and \(v=g'u\) converts the result to
\[
 \int_G |\Phi(u)|\int_{[G]}
               |y'(g')|\,|y(g'u)|\,\dd g'\,\dd u
        \le \|\Phi\|_1\|y'\|_2\|y\|_2 .
\]
Here nonnegative integration is justified by
Lemma~\ref{lem:product-integration}. Cauchy--Schwarz and right
invariance of quotient measure give the inequality, using
Proposition~\ref{prop:marked-poincare-kernel} for unimodularity.
Thus the identical unfolding with the complex factors is valid.
Theorem~\ref{thm:marked-test-kills-nongraph-zero-terms} removes
every nongraph rational subspace on the moved tests. For a
graph the zero term of
Lemma~\ref{lem:graph-zero-term-oscillator-pairing} is
\(\langle\Omega(g',1)f',\Omega(\gamma g,1)f\rangle
=\Phi(g'^{-1}\gamma g)\) by the unitary representation law.
This proves the identification.

Finally integrate on the chosen principal subgroup, of mass
\(1/W_0(P)\) by Corollary~\ref{cor:first-arrow-level-masses}.
Off these pairs each right integral coset has mass one and
\(y\) is invariant under its integral subgroup. The remaining
sum acts on \(y\) by \(T_h\), commuting with the real operators
by their separate group variables. The sum and integral
interchange by the established absolute bound, or by norm
convergence of \(T_h\) together with the integrability of
\(\Phi_\infty\).
\end{proof}

\begin{corollary}\label{cor:remaining-doubled-mean-for-first-arrow}
In Theorem~\ref{thm:absolute-doubled-graph-unfolding} use \(f'=f\)
and \(h=6r\). Put \(H=U(T)\),
\(\Theta(g,h')=\theta(\Omega(g,h')f)\) and
\(v y(h')=\int_{[G]}y(g)\Theta(g,h')\dd g\).
Let \(W\ne0\) be a space of initial forms with the stipulated
finite level. Suppose the following assertions hold on these forms.
The target pairing of the integral lifts expands by product
integration, and for almost all \(g,g'\) on the two source
quotients the target mean satisfies
\[
 \frac1{{\rm vol}[H]}
       \int_{[H]}\overline{\Theta(g',h')}\Theta(g,h')\dd h'
              =\sum_{\gamma\in\Gamma}\Phi(g'^{-1}\gamma g).
\]
Suppose the weak real convolution
\(L_\infty=\int_{G_\infty}\Phi_\infty(a)R(a)\dd a\)
acts by a nonzero scalar \(\alpha\) on \(W\).
Then on \(W\) the Gram form of \(v\) is
\[
 \langle vx,vy\rangle_{[H]}
       ={\rm vol}[H]\,\alpha\,\frac{\|v_F\|^2}{W_0(P)}
                       \langle x,T_h y\rangle_{[G]} .
\]
Thus \(v\) is injective for all sufficiently large scales.
\end{corollary}
\begin{proof}
The asserted expansion and target mean reduce the Gram pairing
to the graph integral in
Theorem~\ref{thm:absolute-doubled-graph-unfolding}.
Absolute integrability of the real coefficient makes its
weak integral bounded by \(\|\Phi_\infty\|_1\) on unit vectors.
The finite operator commutes with this convolution: every
partial sum of its right translations commutes with every
real operator, and the tail is bounded in norm on the fixed
space by Theorem~\ref{thm:positive-finite-first-arrow-convolution}.
Consequently the integrand paired after real integration gives
\(\langle x,T_h L_\infty y\rangle=\alpha\langle x,T_hy\rangle\).
The target quotient volume is positive and finite in the stated
mean, and \(\|v_F\|>0\); apply
Corollary~\ref{cor:first-arrow-injectivity-from-norm-identity}
to obtain injectivity.
\end{proof}


\section{A height series over rational subspaces}

\begin{definition}\label{def:doubled-pluecker-height}
Let \(E\) be a CM field, \(O\) its ring of integers,
\(\Sigma\) one embedding per conjugate pair, and \(b=2|\Sigma|\).
Give \(E_\infty^n=\prod_{\tau\in\Sigma}\C^n\) the sum of the ordinary
squared Euclidean norms. For \(1\le d\le n\), let \(A\subset E^n\)
be a rational \(d\)-plane and \(Y\in M_{n,d}(E)\) a basis matrix
of \(A\). At every finite place \(w\) put
\(|u|_w=q_w^{-v_w(u)}\) for \(u\ne0\), and \(|0|_w=0\), where
\(q_w\) is the residue cardinality. Define
\[
 H(Y)=
 \prod_{\tau\in\Sigma}\det(Y_\tau^*Y_\tau)
 \prod_{w\ {\rm finite}}\ \max_{|I|=d}|\det Y_I|_w .
\]
Here \(Y_I\) consists of the rows indexed by \(I\). Let \(L_A=A\cap O^n\)
and take its covolume in \(\prod_\tau A_\tau\) with the induced
Euclidean measure. Write \(V_E\) for the covolume of \(O\) in \(E_\infty\).
\end{definition}

\begin{proposition}\label{prop:rational-subspace-covolume-height}
For the data of Definition~\ref{def:doubled-pluecker-height},
\(L_A\) is a full lattice in \(\prod_\tau A_\tau\), and
\[
                    \operatorname{covol}(L_A)=V_E^{d} H(Y).
\]
In particular \(H(A):=H(Y)\) is independent of the basis.
\end{proposition}
\begin{proof}
A change \(Y\mapsto YB\) multiplies every minor by \(\det B\)
and every infinite Gram determinant by \(|\tau(\det B)|^2\).
The product of these factors is one. Indeed the principal ideal
norm of a nonzero integral element is the absolute determinant
of multiplication on \(O\), by Lemma~\ref{lem:number-field-lattice},
and is also the product of its prime-power norms by
Proposition~\ref{prop:general-ideal-factorization}. Division treats
a fractional element. Only finitely many local factors differ from
one for a fixed rational frame, by clearing denominators and choosing
one nonzero minor. We may therefore multiply \(Y\) by a nonzero
integer and assume it integral.

Let \(u\in O\setminus\{0\}\) be a nonzero maximal minor. The module
\(Y^{-1}L_A=\{c\in E^d:Yc\in O^n\}\) lies between \(O^d\) and
\(u^{-1}O^d\), by the adjugate formula on the selected rows.
Hence \(YO^d\) has finite index in \(L_A\). In these same coordinates
the quotient is the kernel of multiplication by \(Y\) from
\((O/uO)^d\) to \((O/uO)^n\), by sending a class \(c\) to \(uc\bmod u\).
Factor \(uO\) and take the product of the residue rings of its
prime powers. This follows by the Chinese remainder identity:
for coprime \(I,J\) write \(1=i+j\), \(i\in I,j\in J\);
multiplication by \(j,i\) lifts the two prescribed residues and
\(I\cap J=IJ\). Iteration gives the assertion. The factors are the
same residue rings as those for the completed valuation rings, since
completion preserves residues modulo each power by its definition.

Over a valuation ring, move a nonzero entry of smallest valuation
to the top left. It divides every entry, so elementary integral
row and column operations clear its column and row. Repeat on
the remaining rectangle. This gives a rectangular diagonal
matrix with entries units times \(\pi_w^{a_1},\ldots,\pi_w^{a_d}\).
Integral invertible operations preserve the ideal of maximal
minors by multilinearity in both directions. Its valuation is
therefore \(\sum_j a_j\le v_w(u)\). The local kernel just described
has size \(q_w^{\sum_j a_j}\), by counting each of the diagonal
congruences modulo \(\pi_w^{v_w(u)}\). Thus the reciprocal
index is exactly the finite product in \(H(Y)\).

At \(\tau\), the real volume expansion of the complex injection
\(Y_\tau\) is \(\det(Y_\tau^*Y_\tau)\): orthogonalize its columns,
then each orthogonal length expands a real two-plane in both
coordinates. Lemma~\ref{lem:number-field-lattice} gives \(V_E^d\)
for the domain lattice. Multiplying these expansions gives the
covolume of \(YO^d\), and dividing by the index gives the formula.
\end{proof}

\begin{theorem}\label{thm:rational-plane-height-series}
For fixed \(E,n,d\) as in Definition~\ref{def:doubled-pluecker-height}
there is \(C>0\) such that for \(X\ge1\)
\[
        \#\{A\subset E^n:\dim_E A=d,\ H(A)\le X\}
              \le C X^n(1+\log X)^d .
\]
Consequently \(\sum_{\dim_E A=d}H(A)^{-s}<\infty\) for every
real \(s>n\). The sum over rational maximal isotropic subspaces
in a hyperbolic space of rank \(2r\) converges at \(s=h/2\)
whenever \(h>4r\), in particular at \(h=6r\).
\end{theorem}
\begin{proof}
Write \(N=bd\). For any rank \(N\) Euclidean lattice \(L\) the
proof of Lemma~\ref{lem:reduced-real-basis} gives
\(\lambda_j(L)\le C_N r_j\), where \(r_j\) are the successive
orthogonalized lengths of an integral basis. Hence
\(\prod_{j=1}^N\lambda_j(L)\le C_N^N\operatorname{covol}(L)\).
In \(L_A\) choose successively a shortest vector \(v_j\) outside
the \(E\)-span of the preceding \(j-1\) choices. This is possible
through \(j=d\) by Proposition~\ref{prop:rational-subspace-covolume-height}
and discreteness. The real spans at all embeddings of that excluded
space have total dimension \(b(j-1)\). Among \(b(j-1)+1\) real-linearly
independent short vectors at least one lies outside it. Thus
\[
        \|v_j\|\le \lambda_{b(j-1)+1}(L_A),\qquad
        \prod_{j=1}^d\|v_j\|^b
            \le\prod_{i=1}^N\lambda_i(L_A)\le C' H(A),
\]
where monotonicity of the successive minima gives the middle
inequality.

Let \(\delta>0\) be the shortest length of a nonzero vector of
\(O^n\). Balls of radius \(\delta/2\) centered at its lattice
points have disjoint interiors. Volume comparison in real
dimension \(bn\) bounds the number of these points of length
at most \(R\), for \(R\ge\delta\), by \((1+2R/\delta)^{bn}\).
Choose \(k_j\ge0\) with
\(\delta 2^{k_j}\le\|v_j\|<\delta 2^{k_j+1}\).
For \(H(A)\le X\) we have \(2^{b\sum_j k_j}\le C'' X\),
so at most \(C'''(1+\log X)^d\) tuples can occur. For a fixed
tuple the number of possible ordered vectors is at most a constant
times \(2^{bn\sum_j k_j}\le (C'' X)^n\). Each independent list
determines its span. This proves the count. In particular only
finitely many heights are below one, and they are positive.
Grouping the others between \(2^k\) and \(2^{k+1}\) gives a
majorant proportional to \(\sum_{k\ge0}(1+k)^d2^{-k(s-n)}\),
which converges by the ratio test. Take \(n=2r,d=r\) and retain
the isotropic subspaces to obtain the final assertion.
\end{proof}

\section{The exterior height of a theta zero term}

\begin{definition}\label{def:hyperbolic-height-action}
Use the arithmetic CM field of Lemma~\ref{lem:real-place-rings}.
Let \(D\) be a split hermitian space of dimension \(2r\), \(r\ge1\),
and \(T\) a hermitian target of dimension \(h\) with the finite
frames of Theorem~\ref{thm:global-actual-arrow-covariance}. Use a
rational hyperbolic basis of \(D\), so its Gram matrix is
\(J=\left(\begin{smallmatrix}0&I_r\\ I_r&0\end{smallmatrix}\right)\).
Write \(A_0\) for the first \(r\)-plane and \(Y_0=\binom{I_r}{0}\).
For a basis matrix \(Y\) of any rational maximal isotropic plane
\(A\) and a source adelic isometry \(g\), define
\[
          H_g(A)=H(g^{-1}Y)
\]
by the local products in Definition~\ref{def:doubled-pluecker-height}.
The basis independence proof in
Proposition~\ref{prop:rational-subspace-covolume-height}
applies also with \(g^{-1}\) present. On a Schwartz test \(u\) of
the oscillator on \({\rm Hom}_E(D,T)\), let \(Z_A(u)\) denote
the mean of its theta section over \([\mathcal N_A]\), where
\(\mathcal N_A=\{X:X|_A=0\}\) as in
Definition~\ref{def:doubled-graph-zero-term}. Denote an oscillator
action of the source by \(\Omega(g)\). In absolute values its
unitary scalar normalization is immaterial.
\end{definition}

\begin{lemma}\label{lem:hyperbolic-frame-height-factor}
Let \(K\) be the product of the integral unitary stabilizers of
\(O_w^{2r}\) at finite places (taken in conjugate pairs when split)
and the isometries preserving the ordinary Euclidean metric at all
infinite places. Every rational maximal isotropic plane has a
basis \(Y=\gamma Y_0\) for a rational unitary \(\gamma\).
For every adelic unitary \(g\) there is a factorization
\(g^{-1}\gamma=kp\), \(k\in K\), \(pA_0=A_0\). If
\(pY_0=Y_0 a\), then
\[
                  H_g(A)=|\det a|_{\mathbb A_E}.
\]
The idele norm uses squared complex absolute values at infinity.
\end{lemma}
\begin{proof}
For any rational basis \(Y\) solve \(Y^*JW=I_r\) by full
row rank. Replacing \(W\) by \(W-Y(W^*JW)/2\) preserves this
identity and makes the Gram matrix of \(W\) zero, by expansion.
The combined columns have Gram matrix \(J\) and thus define a
rational isometry \(\gamma\) with \(\gamma Y_0=Y\). For
\(D=S\oplus S^{-}\), the diagonal copy of \(S\) is isotropic
of half dimension. The same construction produces the
hyperbolic basis in Definition~\ref{def:hyperbolic-height-action}.

Consider locally the span of \(g^{-1}\gamma Y_0\). Its intersection
with the integral lattice has a basis with an integral left inverse:
clear denominators and use the valuation-ring row and column
reduction in Proposition~\ref{prop:rational-subspace-covolume-height};
saturation replaces each resulting diagonal ideal by the whole
coordinate ring. In the split case perform this in each field factor.
The basis so obtained is isotropic and its Gram matrix is zero.
By Theorem~\ref{thm:integral-unitary-primitive-frame} it is the
image of \(Y_0\) by an integral unitary matrix. The local algebras
satisfy that theorem's unramified or split hypothesis by
Lemma~\ref{lem:real-place-rings}. At every archimedean place,
pass by the Euclidean unitary change \((v,w)\mapsto
((v+w)/\sqrt2,(v-w)/\sqrt2)\) to the form
\({\rm diag}(I,-I)\). A maximal isotropic plane projects
bijectively to the positive summand and is the graph of a unitary
map onto the negative summand, by equality of the two norms.
Acting independently by unitary matrices on the two summands
therefore carries any one such graph to any other. This supplies
\(k\) at every place. Their product belongs to \(K\), and
\(p=k^{-1}g^{-1}\gamma\) preserves \(A_0\).

Integral invertible matrices preserve the maximum finite absolute
value of the maximal minors by multilinearity and the inverse
inequality. Euclidean unitary matrices preserve the infinite Gram
determinants by direct multiplication. Evaluating the local factors
on \(Y_0a\) leaves just the absolute determinant at each place.
This proves the formula.
\end{proof}

\begin{theorem}\label{thm:adelic-zero-term-exterior-height}
For each fixed Schwartz test \(u\) of
Definition~\ref{def:hyperbolic-height-action} there is a finite
constant \(C_u\) such that for every \(g,A\)
\[
                   |Z_A(\Omega(g)u)|\le C_u H_g(A)^{-h/2}.
\]
One constant works for a compact family of real Schwartz tests
with finite parts in a fixed finite list. For \(h>4r\) the
series of absolute zero terms converges locally uniformly in \(g\).
In particular the graph series of
Theorem~\ref{thm:absolute-doubled-graph-unfolding} at \(h=6r\)
converges absolutely at every pair of source points, uniformly
on compact sets of the pair.
\end{theorem}
\begin{proof}
Use rational \(\gamma\) and \(g^{-1}\gamma=kp\) of
Lemma~\ref{lem:hyperbolic-frame-height-factor}. Transport of the
additive probability by the rational isometry identifies the
mean along \(\mathcal N_A\) with that along \(\mathcal N_{A_0}\)
on the test acted on by \(\gamma^{-1}g=p^{-1}k^{-1}\), in modulus.
Indeed the rational theta action of
Theorem~\ref{thm:rational-theta-symplectic-action} is pullback
on sections, and preserves the quotient probability; any
unitary normalization differs by a unit scalar, by
Theorem~\ref{thm:finite-oscillator-lattice-transport} and
Proposition~\ref{prop:real-symplectic-gaussian-action}.

In the rational polarization with modulation space
\(\mathcal N_{A_0}\), the configuration coordinates are the
restrictions \(X|_{A_0}\in{\rm Hom}_E(A_0,T)\), using the second
hyperbolic summand to extend a restriction by zero. The zero
term is evaluation at the origin: integration over the compact
modulation quotient kills every nonzero rational frequency, as in
Theorem~\ref{thm:rational-schroedinger-zero-section}. The action
on \({\rm Hom}(D,T)\) of \(p^{-1}\) is \(X\mapsto Xp\).
On configuration coordinates this is right multiplication by \(a\),
of module \(|\det a|_{\mathbb A_E}^{h}\) from its \(h\) rows.
The unitary parabolic formula of
Theorem~\ref{thm:graded-parabolic-schroedinger} multiplies
evaluation at zero by the inverse square root of this module
in modulus; its quadratic phase equals one there. By the
height formula it remains to bound evaluation at zero on
\(\Omega(k^{-1})u\) uniformly for \(k\in K\).

At infinity this follows from continuity in Schwartz seminorms
in Theorem~\ref{thm:negative-block-real-splitting}, after a fixed
polarization change by Proposition~\ref{prop:real-symplectic-gaussian-action}.
At the finite places use the lattice model of
Theorem~\ref{thm:finite-oscillator-lattice-transport}. The subgroup
of \(K_f\) preserving its reference lattice is open: two full
adelic lattices are commensurable, and an integral congruence
subgroup preserves the finitely many differing local lattices.
On this subgroup the action up to phase is pullback on the
lattice model. The finite test is a finite linear combination
of reference-vacuum translates, by the basis argument in that
theorem. An open subgroup fixes the finitely many translation
classes and their phases: choose the congruence depth so that the
finitely many differences lie in the reference lattice and their
pairing with every listed representative lies in the kernel of
the finite additive character. Thus the compact group gives only
finitely many test vectors up to phases. Their evaluations in
the fixed polarization are bounded. Products, then finite sums,
give the required \(C_u\). The same proof with the given compact
real test family uses boundedness of the continuous seminorms on
its product with the compact group.

For \(g\) in a compact set, \(H_g(A)\ge c H(A)\) for a common \(c>0\).
Indeed the norms of \(\bigwedge^r g\) on each local exterior space
are bounded and are at most one at almost every finite place
uniformly: a compact subset of a restricted product is contained
in a product with integral factors outside a fixed finite set,
by a finite subcover. Apply \(g\) to the exterior coordinates
of \(g^{-1}Y\), using the squared Euclidean norm at each complex
place. Theorem~\ref{thm:rational-plane-height-series} now gives a
summable majorant. For the last assertion take the doubled
tensor test, with conjugation on the negative Hom summand,
and act by the block diagonal pair. Tensoring the translation
formulas gives its product theta section; this also follows by
absolute convergence of the two rational sums in
Theorem~\ref{thm:rational-schroedinger-zero-section}.
Lemma~\ref{lem:graph-zero-term-oscillator-pairing} identifies
each graph term and Theorem~\ref{thm:marked-test-kills-nongraph-zero-terms}
annihilates the others.
\end{proof}

\section{Removing one hyperbolic plane at a boundary}

\begin{definition}\label{def:positive-hyperbolic-dilations}
Write the hyperbolic basis of Definition~\ref{def:hyperbolic-height-action}
as \(e_1,\dots,e_r,e'_1,\dots,e'_r\), with \(A_0=\sum_i Ee_i\).
In the rational polarization with configurations \(X\in T^r\)
given by the \(r\) columns \(X(e_i)\), write \(s=|\Sigma|\).
For \(x=(x_1,\dots,x_r)\in\Z_{\ge0}^r\), let \(d_x\) be identity
at finite places and \({\rm diag}(2^{x_i})\) at each infinite
place. The source element \(t_x\) has blocks \(d_x,d_x^{-1}\).
On tests set \(D_x u(X)=u(Xd_x)\). For \(1\le j\le r\),
let \(D'\) be the span of the other hyperbolic pairs, \(x'\)
the tuple omitting \(x_j\), and \(u^{(j)}\) the restriction
\(X(e_j)=0\) at all places. For \(r=1\) the smaller test is
a scalar and its only zero term and theta value are that scalar.
\end{definition}

\begin{theorem}\label{thm:uniform-dilated-zero-term-boundary}
If \(h>4r\), then for every \(x,A\) as in
Definitions~\ref{def:hyperbolic-height-action} and
\ref{def:positive-hyperbolic-dilations},
\[
           |Z_A(D_xu)|\le C_u H(A)^{-h/2}.
\]
For \(e_j\notin A\) there is also a bound
\(|Z_A(D_xu)|\le C_{u,A}2^{-sh x_j}\). For \(e_j\in A\)
there is a unique maximal isotropic plane \(B\subset D'\)
with \(A=Ee_j\oplus B\), and in that case
\[
                      Z_A(D_xu)=Z_B(D_{x'}u^{(j)}).
\]
Consequently the difference between the two zero-term series
for \(D_xu\) and \(D_{x'}u^{(j)}\) tends to zero as
\(x_j\to\infty\), uniformly in the other \(x_i\). The bounds
and the limit are uniform over the compact test families in
Theorem~\ref{thm:adelic-zero-term-exterior-height}.
\end{theorem}
\begin{proof}
The action of \(t_x\) on configurations is \(X\mapsto Xd_x^{-1}\)
and has module \(2^{-2sh\sum_i x_i}\). Hence \(D_x\) is
\(2^{-sh\sum_i x_i}\Omega(t_x)\) up to a unit scalar, by
Theorem~\ref{thm:graded-parabolic-schroedinger}; there is no
quadratic term in this block diagonal change. Fix a rational
frame \(Y\) of \(A\). At each infinite place all coordinates of
\(\bigwedge^r t_x^{-1}\) in the exterior basis are scaled by
positive real numbers at least \(2^{-\sum_i x_i}\): in a
chosen minor each negative exponent can appear only once.
The squared norm of an exterior vector is the sum of the squared
absolute values of its maximal minors (expand the determinant
of the Gram matrix by multilinearity). At finite places \(t_x=1\).
Thus \(H_{t_x}(A)\ge 2^{-2s\sum_i x_i}H(A)\). The normalization
cancels this factor to power \(-h/2\) in
Theorem~\ref{thm:adelic-zero-term-exterior-height}.

If \(e_j\notin A\), deletion of the \(j\)-th coordinate is
injective on \(A\), so one nonzero rational minor lacks that row.
Its scale is at least \(2^{-\sum_i x_i+x_j}\) at each place.
Its absolute value is positive at every embedding. Using just this
minor in the infinite factors and retaining the fixed finite
factors gives the second bound. If \(e_j\in A\), orthogonality
to \(e_j\) forces the \(e'_j\) coordinate of every member of
\(A\) to vanish. Subtract its \(e_j\) component to see that
\(B=A\cap D'\) has dimension \(r-1\); this gives both uniqueness
and the bijection with the smaller isotropic planes.

For such a direct sum, \(\mathcal N_A\) is the direct sum
of the modulation space on the \(j\)-th hyperbolic pair and
\(\mathcal N_B\) on the other pairs. The trace pairing and its
translation phases have no cross terms between these two Hom
spaces. Integrate the theta section first over that modulation
space. Character orthogonality kills exactly the frequencies
with nonzero \(j\)-th column, leaving the smaller theta section
of the restriction \(X(e_j)=0\), as in
Theorem~\ref{thm:rational-schroedinger-zero-section}. Uniform
absolute convergence on compact representatives in that theorem
permits the integration for a test coupling different columns.
The product of the two additive quotient probabilities is
the quotient probability on their direct sum. Its remaining
integral is \(Z_B(D_{x'}u^{(j)})\), proving the exact identity
with the explicit normalization of \(D_x\). The difference of
the two series is therefore the sum over planes missing \(e_j\).
Theorem~\ref{thm:rational-plane-height-series} and the first bound
control a tail of this sum uniformly. The second bound on the
finitely many remaining terms tends to zero uniformly in the
other exponents and over the indicated tests.
\end{proof}

\begin{proposition}\label{prop:pointwise-dilated-theta-boundary}
For \(b\in U(T)(\mathbb A_{E^+})\) and a test \(u\) in the
polarization of Definition~\ref{def:positive-hyperbolic-dilations}
put
\[
   \Theta_r(u;b)=\sum_{X\in T(E)^r}u(b^{-1}X).
\]
For \(b\) in a fixed compact subset and \(u\) in a compact real
Schwartz family with fixed finite parts as in
Theorem~\ref{thm:adelic-zero-term-exterior-height}, for each \(m>0\)
there is \(C_m\) such that
\[
  |\Theta_r(D_xu;b)-\Theta_{r-1}(D_{x'}u^{(j)};b)|
          \le C_m 2^{-m x_j}
\]
for every \(x\). In particular this convergence holds pointwise
throughout the target quotient.
\end{proposition}
\begin{proof}
By clearing denominators on the compact set of finite parts of
\(b\) times the finite supports, all contributing columns belong
to one fractional integer lattice \(L\) in \(T_\infty\).
Within the compact set of infinite \(b\), there is \(c>0\)
with \(\|b^{-1}_\infty v\|\ge c\|v\|\). Since all the diagonal
factors at infinity are at least one,
\(\|b_\infty^{-1}Xd_x\|\ge c\|X\|\); if \(X(e_j)\ne0\),
the same length is at least \(c\,2^{x_j}\lambda_1(L)\).
Uniform Schwartz decay of order \(m+[E:\Q]rh+2\) bounds the sum on the nonzero \(j\)-th columns by
a constant times \(2^{-m x_j}\sum_{X\in L^r}(1+\|X\|)^{-[E:\Q]rh-2}\).
The last sum converges by lattice ball counting in the proof
of Theorem~\ref{thm:rational-plane-height-series}. On the zero
column the summands are precisely those of the smaller theta
value. Reindexing rational target translates leaves each sum
invariant, so the pointwise limit descends to the quotient.
\end{proof}

\section{Reduction along a shortest isotropic line}

\begin{definition}\label{def:unitary-reduction-tower}
Let \(F=E^+\) for the arithmetic CM field \(E\), and let \(V\) be
a hermitian \(E\)-space of rank \(n>0\) admitting a lattice
self-dual for its hermitian form at the finite places.
Split off a rational hyperbolic pair whenever a nonzero rational
isotropic vector exists. For such a vector \(e\), choose \(f\)
with \(\langle e,f\rangle=1\) and subtract
\(e\langle f,f\rangle/2\) from \(f\). This gives an isotropic
partner. Continue on the nonsingular perpendicular. Denote the
resulting pairs by \(e_i,f_i\), \(1\le i\le w\), in the order
removed, and the final space by \(V_{\rm an}\). Thus \(2w\le n\);
\(V_{\rm an}\) has no nonzero rational isotropic vector.

Choose a self-dual finite lattice in \(V_{\rm an}\) and adjoin the
standard integral lines on all pairs to obtain \(\Lambda\) on \(V\).
At infinity choose the Euclidean metric orthogonal on these
summands, of unit length on each displayed basis vector, and
coming from signed orthonormal coordinates on \(V_{\rm an}\).
The same choices restrict to every stage. Let \(K\) be the
product of the unitary stabilizers of the finite lattice and
the infinite Euclidean metrics. For \(g\in U(V)(\mathbb A_F)\)
set
\[
 \lambda_V(g)=\min\{\max_{\tau\in\Sigma}\|g_\tau^{-1}v_\tau\|:
             0\ne v\in V(E),\ g_f^{-1}v\in\Lambda\}.
\]
At a positive-rank stage write \(n\) for its current rank.
For \(t>0\), \(d(t)\) is the idele of \(E\) equal to \(t\)
at every infinite place and to one at finite places, and
\(b=[E:\Q]\).
\end{definition}

\begin{lemma}\label{lem:unitary-tower-compact-parts}
The lattices and minima in Definition~\ref{def:unitary-reduction-tower}
exist. For each fixed \(\epsilon>0\) the classes with
\(\lambda_V(g)\ge\epsilon\) have representatives in a compact
subset of the unitary group. A compact set of norm-one ideles meets
every rational multiple class in the norm-one ideles of \(E\).
For any open neighborhood \(U\) of zero in \(\mathbb A_E\) there
is an integer \(M\ge1\) such that, for every \(\alpha\in\mathbb A_E\),
there exist \(1\le m\le M,\ c\in E\) with \(m\alpha+c\in U\).
\end{lemma}
\begin{proof}
At a finite place of \(F\), in the original self-dual lattice choose
a primitive generator \(e_*\) of its intersection with the local
isotropic line of \(e_1\). In a split pair choose a primitive
generator on each branch. Such generators exist by valuation-ring
elimination as in Proposition~\ref{prop:rational-subspace-covolume-height}.
Self-duality supplies a lattice vector \(y\) paired to one with
\(e_*\): in a lattice basis primitiveness and invertibility of
the Gram matrix make its pairing row unimodular. There is an
integral \(c_0\) with \(c_0+\bar c_0=\langle y,y\rangle\), since
the arithmetic field has an integral element of relative trace
one by Lemma~\ref{lem:heisenberg-layer}. Subtract \(e_*c_0\).
The two vectors now span an integral hyperbolic plane with
unimodular Gram matrix. Orthogonal projection splits it off
integrally. Its self-dual perpendicular realizes a lattice in
\(e_1^\perp/Ee_1\), identified rationally with the prescribed next
perpendicular. These operations can use the original rational
pair at almost every place, since finitely many basis coordinates
and pairings suffice to test integrality. Iteration supplies the
announced adelic lattice on \(V_{\rm an}\).

For the minimum, the moved real lattice is full and has fixed
covolume by the determinant argument of
Theorem~\ref{thm:unitary-linear-shortest-representative}: a unitary
determinant has absolute norm one at every real place and on
each group of finite places over \(F\). Lemma~\ref{lem:class-covol}
computes the finite contribution and the real change of variables
computes the infinite one. A lattice minimum is attained, and
is invariant under rational left translation. For completeness,
on the part with minimum at least \(\epsilon\),
Lemma~\ref{lem:reduced-real-basis} gives an integral basis of the
underlying real lattice of bounded lengths, since its covolume is
fixed. Choose an \(E\)-basis from the corresponding rational
vectors; they span over \(E\) since they span over \(\Q\).
The hermitian Gram entries are integral and bounded at all
embeddings by the Euclidean bounds. Hence only finitely many
Gram matrices occur by Lemma~\ref{lem:number-field-lattice}.
Match to one fixed rational basis for each occurring matrix by
a rational isometry. The inverse of the resulting representative
takes that basis into a compact adelic box. Thus its inverse
matrix lies in a compact box; the unitary inverse formula bounds
the forward matrix in another compact box. The defining matrix
equations are closed in their product. This proves the compact
representative assertion at every stage, including \(V_{\rm an}\).

For norm-one ideles use the lattice
\(\{v a_\infty:v\in E,\ v a_f\in\widehat O_E\}\).
Its covolume is \(V_E|a|_{\mathbb A_E}=V_E\), by the same local
index calculation as in Lemma~\ref{lem:real-volume-compact},
with squared complex moduli at the infinite places. An averaging
of a fixed sufficiently large ball over a lattice parallelepiped
gives two points congruent modulo the lattice, hence a nonzero
lattice point of bounded length. Take the corresponding \(v\).
The product formula of Proposition~\ref{prop:rational-subspace-covolume-height}
bounds the integral norm of the finite ideal of \(va\), and
bounds each infinite modulus away from zero as all are bounded
above. Finitely many ideals have bounded norm, by finite
factorization in Proposition~\ref{prop:general-ideal-factorization}
and finiteness of the residue rings over any bounded rational
prime. At finite places the permissible ideles therefore lie
in finitely many translates of the compact integral units.
This gives the asserted compact set. Finally \(\mathbb A_E/E\)
is compact by the additive lattice argument of
Lemma~\ref{lem:theta-positive-translation}. Choose an open set
\(U'\) about zero with \(U'-U'\subset U\). Finitely many translates
of its image cover the quotient. Among one more multiple
\(0,\alpha,\ldots,M\alpha\) than there are translates two
are in the same translate; their difference gives \(m,c\).
\end{proof}

\begin{theorem}\label{thm:unitary-short-line-reduction}
At each stage in Definition~\ref{def:unitary-reduction-tower},
if \(\lambda_V(g)<1/2\), the stage has a rational hyperbolic pair.
Write it as \(V=Ee\oplus W\oplus Ef\) using its fixed first pair.
Write \(N\) for the unipotent isometries upper block triangular with
identity on all three diagonal blocks, and put
\(m(a,z)={\rm diag}(a,z,\bar a^{-1})\). There exist rational
\(\gamma\in U(V)(F)\), \(k\in K\), \(n_0\in N(\mathbb A_F)\), and
\[
       \gamma g=n_0 m(a,z)k,\qquad
       a=d(t)u,\qquad z\in U(W)(\mathbb A_F),
\]
where \(u\) lies in a fixed compact set of norm-one ideles,
\(t\ge \lambda_V(g)^{-1}\), and \(t\le C\lambda_V(g)^{-1}\).
If \(W\ne0\), then \(\lambda_W(z)\ge c\lambda_V(g)\).
The positive constants depend only on the stage. The unipotent
factor can be taken in a fixed compact set after rational
unipotent left translation.
\end{theorem}
\begin{proof}
Take a rational vector \(v\) attaining the minimum.
Its norm pairing is an algebraic integer, since \(g_f^{-1}v\)
is integral. At every embedding its absolute value is at
most \(\lambda_V(g)^2<1\): the matrix of the form for the
chosen metric is a self-adjoint involution. The norm of a
nonzero algebraic integer has absolute value at least one
by Lemma~\ref{lem:number-field-lattice}. Thus \(v\) is isotropic.

A rational isometry can carry \(v\) to the specified \(e\).
On the same line, scale it and its dual inversely with conjugation.
For distinct nonorthogonal lines the two vectors span a hyperbolic
plane: interchange them there with a scalar on the second image
chosen to preserve their nonzero mutual pairing, and take identity
on the perpendicular. For distinct orthogonal lines solve the two
independent linear equations prescribing pairing one against
both vectors. Subtract a multiple of \(v\) from the solution to
make it isotropic, using half its norm as in
Definition~\ref{def:unitary-reduction-tower}. The two pairings stay
one. Two applications of the nonorthogonal case finish.
After this rational move, locally \(K\) acts transitively on
isotropic lines. At finite places apply
Theorem~\ref{thm:integral-unitary-primitive-frame} to primitive
generators in the self-dual lattice and \(e\); at infinity the
positive and negative projections of a unit isotropic vector
both have length \(1/\sqrt2\) and their two unitary groups are
transitive on the corresponding spheres. Thus choose \(k\)
so that the remaining left factor preserves \(Ee\). Its action
on \(e^\perp/Ee=W\) is unitary, and pairing with \(e\) forces
the last diagonal entry to be \(\bar a^{-1}\). Removing these
diagonal blocks gives \(n_0\). In particular the inverse image
of \(e\) has finite part \(k_f^{-1}a_f^{-1}e\) integral at this
point and infinite component of length \(|a_\tau|^{-1}\le
\lambda_V(g)\) at each \(\tau\). Hence
\(t=|a|_{\mathbb A_E}^{1/b}\ge1/\lambda_V(g)\).
Left multiply by a rational \(m(c,1)\) to put \(u=a/d(t)\)
in the compact set of Lemma~\ref{lem:unitary-tower-compact-parts}.
This leaves \(t\) unchanged by the product formula. For these
new factors some fixed nonzero integer \(B\) has \(Ba_f^{-1}\)
integral for all the allowed \(u\). The rational vector \(Be\)
in the new coordinates is then admissible with real lengths
at most \(C/t\). Rational invariance gives the reverse bound.

Suppose \(W\ne0\), and choose \(v'\) realizing \(\lambda_W(z)\).
There is an adele \(\alpha\) with \(n_0^{-1}v'=v'+\alpha e\).
By compactness of \(u\), \(\bigcap_u u_f\widehat O_E\)
contains \(B'\widehat O_E\) for some positive integer \(B'\),
and \(|a_\tau|^{-1}\le C_1/t\). Apply the last assertion of
Lemma~\ref{lem:unitary-tower-compact-parts} to the neighborhood
with finite part \(B'\widehat O_E\) and infinite moduli less
than \((2C_1)^{-1}\). The rational vector \(mv'+ce\) is nonzero
and after the inverse group action is integral at finite places;
its infinite lengths are bounded by
\(M\lambda_W(z)+1/(2t)\), by orthogonality or the triangle
inequality. Since \(1/t\le\lambda_V(g)\), the desired bound
follows with \(c=1/(2M)\).

Finally the unipotent coordinates consist of a row on \(W\) and a
central entry with specified relative trace, as in
Lemma~\ref{lem:heisenberg-layer}. Use rational rows to reduce
the first coordinate modulo a compact additive set of representatives.
A rational row has a rational lift, by taking minus half its
hermitian norm for the required trace. Then reduce the remaining
skew central coordinate by the rational skew line. Both additive
quotients have compact representatives by
Lemma~\ref{lem:theta-positive-translation}. The lifts of the
compact row set using the same quadratic formula are compact.
This gives the claimed compact unipotent representatives.
\end{proof}

\section{Ordered scales on the unitary quotient}

\begin{theorem}\label{thm:ordered-unitary-cusp-cover}
Use the tower of Definition~\ref{def:unitary-reduction-tower} and
write \(W_j\) for the perpendicular after \(j\) pairs,
\(0\le j\le w\). Write \(N_j\) for the upper triangular unipotent
block group for the ordered blocks
\(e_1,\ldots,e_j,W_j,f_j,\ldots,f_1\), with diagonal blocks
the identity, and write \(m_j(a_1,\ldots,a_j,z)\) for diagonal
blocks \(a_1,\ldots,a_j,z,\bar a_j^{-1},\ldots,\bar a_1^{-1}\).
There are compact sets in \(N_j(\mathbb A_F)\), in the norm-one
ideles of \(E\), and in \(U(W_j)(\mathbb A_F)\), and a fixed
constant \(C\ge1\), such that every \(g\in U(V)(\mathbb A_F)\)
has a representative modulo \(U(V)(F)\) of the form
\[
      n_0 m_j(d(t_1)u_1,\ldots,d(t_j)u_j,z)k,\qquad k\in K,
\]
with \(n_0,u_i,z\) in the respective compact sets,
\(1\le t_i\) and \(t_{i+1}\le Ct_i\). For \(j=0\) there are
no scales. When \(V\) is the split source of
Definition~\ref{def:positive-hyperbolic-dilations}, one compact
set \(\mathcal C\) suffices in the covering
\[
         U(V)(\mathbb A_F)=
            \bigcup_{\gamma\in U(V)(F),\ x\in\Z_{\ge0}^r}
                          \gamma t_x\mathcal C .
\]
\end{theorem}
\begin{proof}
If the minimum is at least \(1/2\), use the compact representatives
of Lemma~\ref{lem:unitary-tower-compact-parts}. Otherwise
Theorem~\ref{thm:unitary-short-line-reduction} gives \(t_1\)
and the smaller adelic group element. Apply the same procedure
on successive perpendiculars until a thick stage or rank zero.
It must stop at the anisotropic stage at the latest, since
another short vector there would be rational isotropic by that
theorem. If consecutive stages are both reduced, their minima
satisfy \(\lambda_{\rm next}\ge c\lambda_{\rm prev}\), and
their scales lie between the inverse minima and a constant
times these inverses. Hence \(t_{i+1}\le Ct_i\) with one \(C\),
by finiteness of the tower.

Embed each rational move on a perpendicular by identity on
the earlier pairs. It conjugates the earlier unipotent layers
into themselves. Once the later coordinates have been chosen,
reduce the earlier conjugated unipotents back into their compact
sets by rational translations, from the innermost layer out.
These last left translations affect only the outer layer being
reduced. The scale on each earlier pair commutes with all
elements of its perpendicular. The compact factors on the
successive right sides combine in \(K\) by the compatible
metrics and finite lattices of the tower. Thus the unipotents
in their fixed compact sets give one compact product in \(N_j\);
the remaining perpendicular coordinates meet their compact set
by the thick assertion. This proves the first covering.

Move \(m_j(d(t_i),1)\) to the left of \(n_0\). In its conjugation
of \(n_0\) the entry in row block \(a\) and column block \(b\)
is multiplied at infinity by the scale of \(b\) divided by
that of \(a\), for \(a<b\). Along the displayed ordered block
list every such ratio is bounded in terms of \(C,w\), using
\(1\le t_i,\ t_{i+1}\le Ct_i\). The inverse is bounded by the
same argument on the inverse unipotents; at finite places the
scale is identity. Thus the new right factor belongs to a fixed
compact set. Replace \(t_i\) by \(2^{\lfloor\log_2 t_i\rfloor}\),
absorbing ratios in \([1,2]\) into that compact set. For a split
source set the other exponents equal to zero and take the union
of finitely many compact sets. This proves the second covering.
\end{proof}

\begin{lemma}\label{lem:parabolic-cusp-integration}
Let \(G=U(V)(\mathbb A_F)\), \(\Gamma=U(V)(F)\); use compatible
Haar measures on the groups and on the left quotient. In the
coordinates of Theorem~\ref{thm:ordered-unitary-cusp-cover},
for every nonnegative measurable left \(\Gamma\)-invariant
function \(\Phi\), its quotient integral is bounded by a
constant times the sum, \(0\le j\le w\), of
\[
 \int \Phi\big(n_0 m_j(d(t_i)u_i,z)k\big)
         \prod_{i=1}^j t_i^{-b(n-2i+1)}\,\frac{\dd t_i}{t_i}\,
                   \dd n_0\,\dd u_1\cdots\dd u_j\,\dd z\,\dd k .
\]
Here \(n_0,u_i,z,k\) range over fixed compact sets (enlarging
those in the theorem is allowed), \(\dd u_i\) is Haar measure
on the norm-one idele group, and the scale domain is
\(t_i\ge1,\ t_{i+1}\le Ct_i\).
\end{lemma}
\begin{proof}
Write \(P_j=N_jM_j\) for the parabolic preserving the successive
coordinate isotropic subspaces. Its Levi group \(M_j\) consists
of the displayed \(m_j\), by preservation of the hermitian form.
The local transitivity argument in
Theorem~\ref{thm:unitary-short-line-reduction}, repeated on each
perpendicular, gives \(G=P_jK\), also when applied without any
condition on the minimum. The group \(P_j\cap K\) is compact.
For nonnegative functions on \(G\), Haar integration equals a
fixed positive constant times integration on \(pk\) for left
Haar measure on \(P_j\) and probability Haar on \(K\).
Indeed average a compactly supported continuous function on
the right by \(K\). Haar measure on \(G\) is right invariant
under \(K\), since the modular character has compact image there
in the positive reals. The resulting integral on \(G/K\) is
left \(P_j\)-invariant. Under
\(G/K=P_j/(P_j\cap K)\) there is just one such nonzero measure
up to scale: lifting it by averaging over right translates of
the compact stabilizer gives left Haar measure on \(P_j\).
The homogeneous-space identification respects the topology,
since inverse images of compact sets in \(G/K\) are compact in
\(G\) (cover by projections of finitely many relatively compact
open translates), and intersecting with the closed \(P_j\)
preserves compactness. The formula now follows for Borel
integrals by equality of measures.

For one layer of current rank \(n'=n-2(i-1)\), the row variable
of Lemma~\ref{lem:heisenberg-layer} has \(n'-2\) entries over
\(E\); the central skew coordinate is one-dimensional over \(F\).
Product additive measure is Haar by the triangular translation
law in Proposition~\ref{prop:unipotent-covol}. Conjugating by
\(m(a_i,z')\) multiplies the row by \(a_i\) and the inverse
unitary transformation of the perpendicular, and multiplies
the center by \(a_i\bar a_i\). The full module on the row is
\(|a_i|_{\mathbb A_E}^{n'-2}\), since the perpendicular
determinant has norm one on every group of conjugate places.
On the center the module is \(|a_i|_{\mathbb A_E}\) by the
local norm identity, including the two factors at a split place.
Thus the module on \(N_j\) is
\(\delta(m)=\prod_i |a_i|_{\mathbb A_E}^{n-2i+1}\), by
multiplying the successive modules. Left Haar on \(P_j\)
is \(\delta(m)^{-1}\dd n_0\,\dd m\): substitution by left
translation cancels exactly the conjugation module.
The idele norm splits by \(a_i=d(t_i)u_i\) and
\(\dd^\times a_i\) is a constant times \(\dd t_i/t_i\,\dd u_i\),
by the direct product law. This gives the displayed powers.

Let \(\mathcal S_j\) be the image of the compact factors and
scale domain in the first covering. It is a countable union of
compact sets by bounding the scales. Rational translates of
these images cover \(G\); unfolding their indicators in a
fundamental domain as in Lemma~\ref{lem:quotient-measure}
bounds the quotient integral by \(\sum_j\int_{\mathcal S_j}\Phi\).
For every decomposition \(pk\) in one such image the possible
\(p\)'s differ by right multiplication by \(P_j\cap K\).
At infinity that compact subgroup is block diagonal: it
preserves each member of the flag and its Euclidean perpendicular,
and the hermitian involution interchanges the two paired flags.
Its scale moduli on the lines are one. At finite places it has
integral unit diagonals on the lines, and its unipotent and
perpendicular projections are compact, since block projection
and diagonal removal are continuous. Conjugating these finite
unipotents by the allowed finite Levi coordinates stays in a
compact set, since the finite \(d(t_i)\) are identity. Thus
all such \(p\)'s satisfy the same scale inequalities and
enlarged compact restrictions. Apply the Haar formula just
proved to the images. This gives the stated bound.
\end{proof}

\section{An integrable bound for rational tuples}

\begin{theorem}\label{thm:moving-tuple-mean-unitary}
Let \(V,n,w,b,G,\Gamma\) be as in
Definition~\ref{def:unitary-reduction-tower},
and fix a compact subset \(\mathcal B_f\subset V(\mathbb A_{E,f})^d\).
For \(d\ge0,\ R\ge1\) let \(N_R(g)\) count \(d\)-tuples
\(X\in V(E)^d\) satisfying \(g_f^{-1}X\in\mathcal B_f\) and
\(\|(g_\infty^{-1}X)_\infty\|\le R\) in the sum Euclidean metric.
If \(d<n-w\), the quotient has finite positive Haar volume and
\[
                 \int_{\Gamma\backslash G}N_R(g)\,\dd g
                          \le C(1+R)^{bnd}.
\]
The constant depends on the field, space, compact set, measures
and \(d\). In particular the condition holds when \(2d<n\).
\end{theorem}
\begin{proof}
The count is measurable by the countable sum of the indicators
of the stated Borel conditions. Use the coordinates in
Lemma~\ref{lem:parabolic-cusp-integration}. Absorb the right compact
factor into the allowed compact finite set and real ball, since
its action and inverse are uniformly bounded there. In
\(g^{-1}X\) one now first applies an upper block unipotent
\(n_0^{-1}\) to \(X\), then the inverse block diagonal element.
Count the rows of \(X\) from the last block to the first.
In a fixed row block the unipotent adds a shift determined by
the later, already chosen rows. After multiplying the box by the
diagonal block, the possible rational rows lie in a translated
finite compact box and in a translated real ball. If any rational
row is present, subtract one such row. The differences have finite
part in one fixed fractional integral lattice: the finite
parts of \(u_i,z\) lie in compact sets and \(d(t_i)_f=1\).
The real differences have lengths at most \(C R t_i\) on the
\(e_i\) block, \(C R/t_i\) on \(f_i\), and \(CR\) on \(W_j\).
The use of differences makes these estimates independent of all
the shifts. They also hold if a block has zero possible rows.

For a block of dimension \(k\) over \(E\), rational differences
with these finite constraints form a subset of a fixed real lattice
of dimension \(bkd\). A ball of radius \(L\) contains at most
\(C'(1+L)^{bkd}\) such points by packing as in
Theorem~\ref{thm:rational-plane-height-series}. Multiplication
of the successive bounds, using \(t_i\ge1\), gives
\[
                 N_R(g)\le C''(1+R)^{bnd}
                                       \prod_{i=1}^j t_i^{bd}.
\]
All compact integrations have finite measure. It remains to
integrate \(\prod_i t_i^{-b(n-2i+1-d)}\dd t_i/t_i\) in the scale
domain of Lemma~\ref{lem:parabolic-cusp-integration}. Replace
\(t_i\) by \(s_i=C^{j-i}t_i\), paying a fixed constant, and
enlarge to \(s_1\ge\cdots\ge s_j\ge1\). With
\(v_i=\log s_i-\log s_{i+1}\) and \(s_{j+1}=1\), the Jacobian
in the logarithms is one. The coefficient of \(-b v_k\)
in the exponential is
\[
                   \sum_{i=1}^k(n-2i+1-d)=k(n-k-d)>0
                                      \qquad(1\le k\le j\le w).
\]
The integrals of the resulting decaying exponentials on
\([0,\infty)\) are finite. This proves the count. Using the
same estimate at \(d=0\) on the singleton box gives finite
quotient volume for every \(n>0\); positivity follows from
an identity neighborhood disjoint from its other rational left
translates as in Lemma~\ref{lem:quotient-measure}.
\end{proof}

\begin{corollary}\label{cor:integrable-dilated-theta-majorant}
Let \(V=T\) have rank \(h>2r\), let it have a self-dual finite
lattice as above, and consider \(0\le d\le r\) hyperbolic pairs
in the source. There is a single integrable function of
\(b'\in[U(T)]\) bounding
\[
       \sum_{X\in T(E)^d}|(D_xu)(b'^{-1}X)|
\]
for all \(x\in\Z_{\ge0}^d\) and all \(u\) in a compact family
of real Schwartz tests with fixed finite parts as in
Theorem~\ref{thm:adelic-zero-term-exterior-height}. Define
\(I_d(u)\) to be the integral of \(\Theta_d(u;b')\) divided
by the quotient volume. For \(1\le j\le d\), as \(x_j\to\infty\),
\[
                  I_d(D_xu)-I_{d-1}(D_{x'}u^{(j)})\longrightarrow0
\]
uniformly over that family and the other coordinates \(x'\).
\end{corollary}
\begin{proof}
For a common compact finite support take the union of the finite
supports of the family and sum the bounds for its elementary
factors. The dilation expands all real columns and fixes the
finite part. For \(M>bhd+2\), uniform Schwartz decay bounds
the summands by \(C(1+\|b_\infty'^{-1}X\|)^{-M}\) times
the finite indicator. Separating the unit ball and the shells
\(2^{k-1}<\|b_\infty'^{-1}X\|\le2^k\), \(k\ge1\), gives
the quotient majorant
\[
              C N_1(b')+2^M C\sum_{k\ge1}2^{-kM}N_{2^k}(b').
\]
Its integral is finite by Theorem~\ref{thm:moving-tuple-mean-unitary}
and nonnegative summation. The assertion for \(d=0\) is just
absolute value of the scalar test on a finite-volume quotient.
The smaller restricted tests again form a compact family:
restriction at zero is continuous in each real Schwartz seminorm
and preserves local constancy and compact support on the finite
arguments. Proposition~\ref{prop:pointwise-dilated-theta-boundary}
gives pointwise convergence along every sequence with
\(x_j\to\infty\), even if the other coordinates and the test
vary along the sequence. The two majorants are integrable and
independent of that sequence, so dominated convergence gives
the integral limit. If uniform convergence failed, selecting
a sequence on which the difference stayed bounded below would
contradict this conclusion.
\end{proof}

\section{The target action and a split matrix model}

\begin{proposition}\label{prop:commuting-rational-source-target}
In the split source of Definition~\ref{def:hyperbolic-height-action}
write \(d\) for the number of hyperbolic pairs. On the configuration
space \(T(\mathbb A_E)^d\) let \(V(b)u(X)=u(b^{-1}X)\) for
\(b\in U(T)(\mathbb A_F)\). For rational source isometries \(\gamma\)
let \(\mathcal T_\gamma\) be the operator normalized on rational theta
sections by Theorem~\ref{thm:rational-theta-symplectic-action}.
These actions commute. If \(h>4d>0\), put
\[
    S_d(u)=\sum_A Z_A(u),\qquad
    \mathcal L_d(u)=I_d(u)-S_d(u),
\]
where \(I_d\) is the normalized integral of
Corollary~\ref{cor:integrable-dilated-theta-majorant}.
Set \(I_0(u)=S_0(u)=u\) on scalar tests. The functionals \(I_d,S_d\)
are invariant under both \(V(b)\) and \(\mathcal T_\gamma\).
\end{proposition}
\begin{proof}
The target determinant has module one on its full adelic vector
space, since its product with its conjugate is one. Thus \(V(b)\)
is unitary. It implements the target action on translations and
modulations by change of variable and preservation of the trace
pairing in Definition~\ref{def:hom-oscillator-pairing}.

We give generators over the rational field. Let \(A=\gamma A_0\)
and transform \(A\cap A_0\) by a change of basis of \(A_0\),
extended by the conjugate inverse transpose on the opposite
summand, into the first \(k=\dim(A\cap A_0)\) coordinate lines.
Pass to their orthogonal space modulo these lines. There \(A\)
projects isomorphically to the primed summand of the remaining
pairs, by disjointness from the unprimed summand and dimension.
It is the graph of a skew-hermitian matrix \(Z\), since its
columns have Gram matrix \(Z+Z^*=0\). The upper unipotent shear
by \(-Z\) removes the graph. In the unreduced space the first
\(k\) coordinates of the lifted columns can be removed using
the \(k\) lines already in \(A\). Hence elements preserving
\(A_0\) and interchanges \(e_i\leftrightarrow e'_i\) generate
the rational source group: the operations just described and
the indicated interchanges return \(A\) to \(A_0\).

For the stabilizer of \(A_0\) the operator in
Theorem~\ref{thm:graded-parabolic-schroedinger} consists of
a column change and a quadratic phase. The phase is obtained
by pairing the columns with each other in \(T\), since its
bilinear expression pairs \(X\) with the source shear applied
to \(X\) in the trace Hom form. Both operations commute with
\(V(b)\), by target isometry. For the interchange of one pair
the operator is a Fourier transform on that column for the
nondegenerate adelic trace pairing on \(T\), with the sign and
the identification with its dual fixed by the symplectic form.
Indeed this transform interchanges translations and the prescribed
modulations by direct change of variable; the normalization on
rational theta sums follows from Lemma~\ref{lem:rational-partial-poisson}.
A unit scalar in this description has no effect on commutation.
The target preserves both the pairing and additive measure, so
it commutes with the partial Fourier transform. This proves
commutation for all rational source isometries.

Evaluation at the origin of the theta section is the rational
configuration sum. The rational action preserves it by
Theorem~\ref{thm:rational-theta-symplectic-action}; applying
this after \(V(b)\) proves rational invariance of the integrand.
Right invariance of quotient Haar measure gives target invariance
of its integral, using the absolute integrability in
Corollary~\ref{cor:integrable-dilated-theta-majorant}. Here the
unitary group is unimodular: conjugation on its Lie algebra has
absolute determinant one at real and finite places, since after
splitting the quadratic algebra it is conjugation on all square
matrices with determinant \((\det b)^h(\det b^{-1})^h=1\).
Equivalently use the matrix Haar calculation of
Proposition~\ref{prop:marked-poincare-kernel}. On \([\mathcal N_A]\)
rational changes carry quotient probabilities to quotient
probabilities, permuting the zero terms. The zero term of \(A_0\)
is evaluation of the test at zero and hence is target invariant.
Lemma~\ref{lem:hyperbolic-frame-height-factor} and the commutation
already proved give the same invariance for every \(A\). Sum
absolutely by Theorem~\ref{thm:adelic-zero-term-exterior-height}.
\end{proof}

\begin{lemma}\label{lem:good-split-place-doubled-matrix}
For each fixed adelic test \(u\) and \(d\ge1\) there is a split
finite place of \(F\), with completion \(L\), integer ring
\(O_L\), uniformizer \(\pi\) and residue cardinality \(Q\),
where the local test in the following realization is
\(\phi_0={\bf1}_{M_{h,2d}(O_L)}\). The source and target
groups there are respectively \({\rm GL}_{2d}(L)\) and
\({\rm GL}_h(L)\). On matrix tests the local action \(V(b)\)
of Proposition~\ref{prop:commuting-rational-source-target} and
a unitary choice \(B_c\) of source oscillator operators are
\[
    (V(b)\phi)(Z)=|\det b|^{-d}\phi(b^{-1}Z),\qquad
    (B_c\phi)(Z)=|\det c|^{h/2}\phi(Zc).
\]
Here the target action has the same normalization as the global
one. Every other unitary source implementer differs from \(B_c\)
by a unit scalar.
\end{lemma}
\begin{proof}
There are infinitely many split places available. The field
contains \(i\) and \(E=F(i)\), since \(F\) is real and the
extension has degree two. For any finite list of rational primes,
take a positive even multiple \(m\) of them. Every prime \(p\)
dividing \(m^2+1>1\) is odd and outside the list. The class of
\(m\) is a square root of \(-1\) modulo \(p\) with nonzero
derivative \(2m\); lifting one digit at a time, solving its
linear congruence at each power of \(p\), gives a root in
\(\Q_p\). Thus \(X^2+1\) splits into distinct linear factors
in every \(F_v\) above \(p\). We can omit finitely many places,
including denominators of the rational frames, trace character
and their duals. At all remaining places a finite test factors
with one common integral indicator: its compact support and
local constancy make it a finite sum of compact open boxes,
each unrestricted integral at almost all places, by a finite
subcover and common refinement.

In the two field branches of a split place use the target form
to identify the second branch with the dual of the first.
The two configuration matrices are \(M,N\), each with \(h\)
rows and \(d\) columns, and \(V(b)\) acts on functions by
substitution \((b^{-1}M,b^{\mathsf T}N)\). Take the Fourier
transform in \(N\) only, with dual variable \(Y\), using
self-dual additive measure and the pairing on the columns.
Substitution \(N'=b^{\mathsf T}N\) contributes
\(|\det b|^{-d}\) and substitutes \(b^{-1}Y\) in the transform.
The resulting \(Z\) consists of \(M,Y\); in the dual coordinates
fixed by the alternating Hom pairing these are the two blocks
of columns of a map on the first source branch. This follows
also by Heisenberg translation: a displacement on that branch
translates its unprimed columns and, after the transform, its
primed columns, while the dual branch modulates them.
The source action on these displacements is \(Z\mapsto Zc^{-1}\)
by the definition of the Hom action. Change of configuration
variables in the induced oscillator therefore gives the stated
\(B_c\), whose squared determinant factor makes it unitary.
It implements the corresponding dual action on modulations by
the same substitution. Uniqueness up to unit scalars is the finite
translation-basis argument of
Theorems~\ref{thm:split-finite-schroedinger} and
\ref{thm:finite-oscillator-lattice-transport}. At the chosen
good place integral coordinates are integral duals, the
Fourier transform of their integral indicator is the integral
dual indicator by additive character orthogonality (normalize
the integral box to mass one). Thus the local matrix test is
\(\phi_0\) as claimed, with any fixed constant absorbed into the
factors at the other places.
\end{proof}

\section{One split Hecke operator on both sides}

\begin{definition}\label{def:doubled-hyperplane-neighbours}
Use the split place of Lemma~\ref{lem:good-split-place-doubled-matrix}
and write \(N=2d\). For every residue hyperplane \(D\) of row
vectors of length \(N\), choose \(c_D\in M_N(O_L)\) with row
lattice the inverse image of \(D\) and determinant \(\pi\).
Such a matrix is obtained by a residue basis lift and scaling a
complement by \(\pi\), as in Lemma~\ref{lem:split-hecke-eigen}.
Let \(s_D\) be the adelic source element given by \(c_D^{-1}\)
on the first split branch, the paired isometry on its conjugate,
and identity elsewhere. Put \([k]_Q=(Q^k-1)/(Q-1)\), and for
\(h>4d\) put
\[
       P_d(g)=\sum_{A}H_g(A)^{-h/2},\qquad
       \Lambda=[d]_Q\big(Q^{h/2}+Q^{d-h/2}\big),
\]
where \(A\) runs over rational maximal isotropic planes.
\end{definition}

\begin{lemma}\label{lem:positive-height-hyperplane-eigenvalue}
The function \(P_d\) is positive, rationally left invariant,
and bounded on compact sets of the source adelic group. One has
\[
       \sum_D P_d(gs_D)=\Lambda P_d(g),\qquad
                                      \Lambda>[2d]_Q .
\]
\end{lemma}
\begin{proof}
There is at least \(A_0\) in the sum. The height transforms by
\(H_{\gamma g}(A)=H_g(\gamma^{-1}A)\) for rational \(\gamma\),
by the product formula in
Proposition~\ref{prop:rational-subspace-covolume-height}.
The local norm bound on compact sets proved in
Theorem~\ref{thm:adelic-zero-term-exterior-height} and
Theorem~\ref{thm:rational-plane-height-series} give compact
boundedness and convergence. Fix \(A,g\). At the chosen split
place let \(B\subset L^{2d}\) be the first branch of \(g^{-1}A\)
and let \(\bar B_0\) be the reduction of its saturated lattice
\(B\cap O_L^{2d}\). It has residue dimension \(d\) by the
diagonal reduction in Proposition~\ref{prop:rational-subspace-covolume-height}.
The reduction of \(c_D\) has kernel the line annihilated by the
row hyperplane \(D\); these are all lines, each once.
The finite exterior norm on \(B\) is multiplied by
\(\alpha=Q^{-1}\) if that line is contained in \(\bar B_0\),
and by \(\alpha=1\) otherwise. Indeed on a primitive basis
the reduction loses exactly one or zero in rank, respectively.
The elementary divisors of the image in its saturation all
divide \(\pi\), since \(\pi c_D^{-1}\) is integral.
Their product gives precisely the norm multiplier.

The other branch is the annihilator of \(B\) under the perfect
split pairing, and its multiplier is \(\alpha/|\det c_D|\).
To verify this, complete a primitive basis of \(B\) by columns
to an integral invertible matrix. The final \(d\) columns of
its inverse transpose span the annihilator in dual coordinates.
An exterior coordinate of these columns is, up to sign, the
complementary minor of the first \(d\) columns divided by the
determinant. This follows by expanding the top exterior pairing
between the two complementary lists, or by applying the
determinant to the matrix with complementary standard columns
and multiplying by the inverse. After left multiplication of
the original matrix by \(c_D\), the dual matrix is left multiplied
by \(c_D^{-{\mathsf T}}\). Taking maxima of the minors gives
the asserted multiplier. The identification by \(J\) uses an
integral invertible matrix. Ratios of norms are unchanged if
the starting local bases of either branch are multiplied on
the right by any invertible \(d\)-by-\(d\) matrix.

The total height multiplier is thus \(Q\alpha^2\). There are
\([d]_Q\) lines in \(\bar B_0\), and
\([2d]_Q-[d]_Q=Q^d[d]_Q\) outside it. Summing
\((Q\alpha^2)^{-h/2}\) over them gives exactly \(\Lambda\).
All terms are positive, so summing over \(A\) proves the identity.
Finally the difference of \(\Lambda\) and \([2d]_Q\) is
\([d]_Q(Q^{h/2}-1)(1-Q^{d-h/2})>0\).
\end{proof}

\begin{proposition}\label{prop:target-invariant-hyperplane-inequality}
Let \(\ell\) be any complex linear functional on the adelic
tests invariant under the target action \(V(b)\) of
Proposition~\ref{prop:commuting-rational-source-target}.
For fixed \(u\), use its good split place of
Lemma~\ref{lem:good-split-place-doubled-matrix}, and put
\(F_u(g)=|\ell(\Omega(g)u)|\). Then
\[
                        \Lambda F_u(g)\le\sum_D F_u(gs_D).
\]
\end{proposition}
\begin{proof}
At the split place write \(\phi_0\) for its matrix indicator.
For each target residue hyperplane \(H\subset (O_L/\pi)^h\)
of column vectors choose \(b_H\) with column lattice its
preimage and determinant \(\pi\). At an integral \(h\)-by-\(N\)
matrix of residue rank \(j\), exactly \([h-j]_Q\) such lattices
contain every column, and exactly \([N-j]_Q\) of the source
row lattices contain every row. This follows by counting
nonzero linear forms on the quotient by the corresponding
span, modulo scalars. Since
\([h-j]_Q=[h-N]_Q+Q^{h-N}[N-j]_Q\), the two actions in
Lemma~\ref{lem:good-split-place-doubled-matrix} give the
identity of functions
\[
       Q^{-d}\sum_H V(b_H)\phi_0
        =[h-N]_Q\phi_0+
                       Q^{h-N-h/2}\sum_D B_{c_D^{-1}}\phi_0 .
\]
All indicators on both sides vanish off the integral matrices.
Tensor with the other factors of the fixed test, apply
\(\Omega(g)\), then \(\ell\). On the chosen branch the source
action of \(g\) is \(B_{g_v}\) up to phase, hence commutes
with the displayed target action, and its other local factors
commute by separate variables. The sum on the left is
\(Q^{-d}[h]_Q\,\ell(\Omega(g)u)\) by target invariance.
Solving for the source sum, the scalar on the original value
is
\[
       Q^{N-h/2}\big(Q^{-d}[h]_Q-[h-N]_Q\big)
                    =[d]_Q\big(Q^{h/2}+Q^{d-h/2}\big)
\]
by substitution of \(N=2d\). Each translated value on the
right is the value for \(gs_D\) up to a unit scalar, by
uniqueness of the implementers in
Lemma~\ref{lem:good-split-place-doubled-matrix}. The triangle
inequality finishes the proof.
\end{proof}

\section{The mean of the split-source theta series}

\begin{theorem}\label{thm:high-rank-unitary-siegel-weil}
Let \(E/F\) be the arithmetic CM extension of
Lemma~\ref{lem:real-place-rings}, let \(D\) be rationally hyperbolic
of dimension \(2d\), and let \(T\) be a nonsingular hermitian
space of rank \(h>4d\) admitting a self-dual finite hermitian
lattice. For every Schwartz test \(u\) in the rational split
polarization, the normalized theta integral equals the sum of
its rational zero terms:
\[
   \frac{1}{{\rm vol}[U(T)]}\int_{[U(T)]}
          \sum_{X\in T(E)^d} u(b^{-1}X)\,\dd b
             =\sum_{\substack{A\subset D\\ A=A^\perp}} Z_A(u).
\]
Here \([U(T)]=U(T)(F)\backslash U(T)(\mathbb A_F)\) and the
additive mean defining \(Z_A\) has total mass one. Both sides
are absolutely convergent, including the integral of the sum
of the absolute values on the left.
\end{theorem}
\begin{proof}
Corollary~\ref{cor:integrable-dilated-theta-majorant} at \(x=0\)
and Theorem~\ref{thm:adelic-zero-term-exterior-height} give the
absolute convergence, together with finite positive quotient
volume by Theorem~\ref{thm:moving-tuple-mean-unitary}.
Proceed by induction in \(d\), with fixed target. At \(d=0\)
both sides equal the scalar test. Suppose \(d\ge1\) and write
\(\mathcal L_d=I_d-S_d\) as in
Proposition~\ref{prop:commuting-rational-source-target}.
The induction hypothesis applies to every restricted test on
\(d-1\) pairs since \(h>4(d-1)\). For every compact family of
tests of the kind in Theorem~\ref{thm:adelic-zero-term-exterior-height},
\[
           \sup_u|\mathcal L_d(D_xu)|\longrightarrow0
                         \quad\text{as }\max_i x_i\to\infty .
\]
Indeed for each fixed \(j\), subtract the smaller expressions
on \(D_{x'}u^{(j)}\) from \(I_d\) and \(S_d\), respectively.
They cancel by induction; the two differences tend to zero
uniformly in the other coordinates and tests, by
Corollary~\ref{cor:integrable-dilated-theta-majorant} and
Theorem~\ref{thm:uniform-dilated-zero-term-boundary}. There
are only \(d\) possible indices attaining the maximum.

Fix \(u\), and set \(F(g)=|\mathcal L_d(\Omega(g)u)|\) on
the source adelic group. Absolute values allow any unitary
normalization of the local operators, since the implementations
of the same Heisenberg change differ by a unit scalar as in
Theorem~\ref{thm:adelic-zero-term-exterior-height}. By rational
invariance in Proposition~\ref{prop:commuting-rational-source-target},
\(F\) is rationally left invariant. Choose \(\mathcal C\) from
Theorem~\ref{thm:ordered-unitary-cusp-cover}. For
\(g=\gamma t_x k,\ k\in\mathcal C\), one has
\[
  F(g)=2^{sh\sum_i x_i}\big|\mathcal L_d(D_x\Omega(k)u)\big|,
  \qquad
  P_d(g)\ge 2^{sh\sum_i x_i}\,H_k(A_0)^{-h/2},
                  \qquad s=|\Sigma|.
\]
The first equality is the dilation formula in the proof of
Theorem~\ref{thm:uniform-dilated-zero-term-boundary}. For the
second remove \(\gamma\) by rational left invariance and retain
\(A_0\) in the positive sum. The columns of \(t_x^{-1}Y_0\) are
multiplied by \(d_x^{-1}\) on the right, so the exterior height
of \(k^{-1}t_x^{-1}Y_0\) equals
\(2^{-2s\sum_i x_i}H_k(A_0)\). The latter undilated height is
bounded above on \(\mathcal C\), by bounded exterior norms at
the finitely varying places and at infinity. The tests
\(\Omega(k)u\) range, up to phases, in a compact real Schwartz
family with finitely many finite factors. For the real part use
continuity in Schwartz seminorms in
Theorem~\ref{thm:negative-block-real-splitting}. At finite
places the test line has open stabilizer by the lattice argument
in Theorem~\ref{thm:adelic-zero-term-exterior-height}; a compact
set is covered by finitely many corresponding cosets.
For finitely many small \(x\), both functionals on these tests
are bounded by their uniform summable and integrable majorants.
The boundary limit just proved handles all other \(x\).
Consequently for each \(\epsilon>0\) the quantity
\[
                  M_\epsilon=\sup_g\big(F(g)-\epsilon P_d(g)\big)
\]
is a finite real number.

At a good split place for \(u\), apply
Proposition~\ref{prop:target-invariant-hyperplane-inequality}
to \(\mathcal L_d\), which is target invariant by
Proposition~\ref{prop:commuting-rational-source-target}.
Subtract \(\epsilon\) times the positive identity of
Lemma~\ref{lem:positive-height-hyperplane-eigenvalue}. For every
\(g\) this gives
\[
     \Lambda\big(F(g)-\epsilon P_d(g)\big)
       \le \sum_D\big(F(gs_D)-\epsilon P_d(gs_D)\big)
       \le [2d]_Q M_\epsilon .
\]
Taking the supremum and using \(\Lambda>[2d]_Q\) gives
\(M_\epsilon\le0\). In particular \(F(1)\le\epsilon P_d(1)\);
\(P_d(1)\) is finite, so \(\epsilon\downarrow0\) proves
\(\mathcal L_d(u)=0\). This closes the induction.
\end{proof}

\section{The global norm pairing at the first arrow}

\begin{theorem}\label{thm:actual-doubled-global-mean}
Use the spaces and tests \(S,T,f,f'\) of
Theorem~\ref{thm:absolute-doubled-graph-unfolding}, and suppose
\(h>4r\), with the given self-dual finite frames. Put
\(\Theta_f(g,b)=\theta(\Omega(g,b)f)\), \(H=U(T)\).
For all source adelic points \(g,g'\),
\[
 \frac{1}{{\rm vol}[H]}
     \int_{[H]}\overline{\Theta_{f'}(g',b)}\,\Theta_f(g,b)\,\dd b
          =\sum_{\gamma\in U(S)(F)}
                    \langle f',\Omega(g'^{-1}\gamma g,1)f\rangle .
\]
The integral and the series converge absolutely. In particular
this identity applies at \(h=6r\) for all sufficiently large
arithmetic scales.
\end{theorem}
\begin{proof}
In \(D=S\oplus S^-\) take the rational hyperbolic summand \(A_0\)
to be the graph \(\{(x,-x):x\in S\}\). Extend it to a rational
hyperbolic basis by Lemma~\ref{lem:hyperbolic-frame-height-factor}.
For two tests \(\xi,\xi'\) on \({\rm Hom}(S,T)\), the section
\(F_\xi(v_1)\overline{F_{\xi'}(v_2)}\) comes from a test \(u\)
on the doubled Hom space. In the product polarization use
\(\xi\otimes\overline{\xi'}\); the second symplectic form and
bicharacter are reversed, and the absolutely convergent
configuration sum factors. Pass to the polarization with
modulation space \(\mathcal N_{A_0}\) by
Theorem~\ref{thm:rational-theta-symplectic-action}, which preserves
the section as a function of the doubled phase variable.
In particular its value at the origin is
\(\theta(\xi)\overline{\theta(\xi')}\).

The simultaneous actual target actions on \(\xi,\xi'\), with
the second action conjugated, implement on the doubled space
the same Heisenberg change as \(V(b)\) in
Proposition~\ref{prop:commuting-rational-source-target}.
Both operators are unitary. They can differ only by one
scalar of modulus one: use the real uniqueness in
Proposition~\ref{prop:real-symplectic-gaussian-action} and finite
uniqueness in Theorem~\ref{thm:finite-oscillator-lattice-transport},
or their tensor product on the dense elementary tests.
The scalar equals one. Indeed on the tensor of a fixed nonzero
\(\xi\) with its conjugate, Lemma~\ref{lem:graph-zero-term-oscillator-pairing}
gives \(Z_{A_0}(u)=\|\xi\|^2\). The tensor of its simultaneous
target translates gives the same value by unitarity of the
actual pair. The action \(V(b)\) also preserves the value,
by Proposition~\ref{prop:commuting-rational-source-target}.
A nonzero value fixes the scalar for every test.

Now use \(\xi=\Omega(g,1)f,\ \xi'=\Omega(g',1)f'\).
The preceding identification and the pair representation law of
Theorem~\ref{thm:global-actual-arrow-covariance} identify the
theta value on \(V(b)u\) with the product in the displayed integral.
Theorem~\ref{thm:high-rank-unitary-siegel-weil} computes its
normalized mean as \(\sum_A Z_A(u)\), with absolute convergence.
By Theorem~\ref{thm:marked-test-kills-nongraph-zero-terms}
only graphs contribute and, by
Lemma~\ref{lem:graph-zero-term-oscillator-pairing}, their values
are the displayed source coefficients. The rank conditions and the
marked tests at the actual scales are those of
Theorem~\ref{thm:finite-vacuum-tail-first-arrow}.
\end{proof}

\begin{theorem}\label{thm:actual-first-arrow-global-norm}
Under the preceding hypotheses take \(f'=f\), let
\(G=U(S)(\mathbb A_F)\), and put
\(\Phi(u)=\langle f,\Omega(u,1)f\rangle\).
On compactly supported \(L^2([G])\) functions define
\((v_fy)(b)=\int_{[G]}y(g)\Theta_f(g,b)\,\dd g\).
This map has a unique bounded extension
\(L^2([G])\to L^2([H])\), and that extension satisfies
\[
   \langle v_f x,v_f y\rangle_{[H]}
      ={\rm vol}[H]\int_G\Phi(u)\langle x,R(u)y\rangle_{[G]}\,\dd u .
\]
Whenever the original lift integral converges absolutely for
almost every \(b\), its value agrees almost everywhere with
this extension. This holds in particular for the rapidly
decreasing functions of Proposition~\ref{prop:unitary-reduced-height-integrability}.
For \(y\) at the finite level of
Theorem~\ref{thm:absolute-doubled-graph-unfolding}, the right side
equals
\[
 \frac{{\rm vol}[H]\,\|v_F\|^2}{W_0(P)}
      \,\langle x,T_h L_\infty y\rangle_{[G]},\qquad
 L_\infty y=\int_{G_\infty}
       \langle f_\infty,\Omega_\infty(a,1)f_\infty\rangle R(a)y\,\dd a .
\]
\end{theorem}
\begin{proof}
For \(g\) on a compact subset of the source quotient, choose
representatives in a compact subset of the group. This is
possible by a finite cover by projections of relatively compact
open neighborhoods. The diagonal case of
Theorem~\ref{thm:actual-doubled-global-mean} bounds
\(\int_{[H]}|\Theta_f(g,b)|^2\dd b\) uniformly there:
local uniform absolute convergence of the graph series is
Theorem~\ref{thm:adelic-zero-term-exterior-height}. Cauchy--Schwarz
therefore bounds the target integral of
\(|\Theta_f(g,b)\Theta_f(g',b)|\) uniformly for both points
in prescribed compact sets. Compactly supported \(L^2\) functions
have finite \(L^1\) norm. Thus the triple product with
\(|x(g')y(g)|\) is integrable. Expanding the two lift integrals
and using Theorems~\ref{thm:actual-doubled-global-mean} and
\ref{thm:absolute-doubled-graph-unfolding} gives the stated
identity on these functions. The same unfolding bounds the
squared operator norm by \({\rm vol}[H]\|\Phi\|_1\).

The quotient has an exhaustion by compact sets since the matrix
group is a countable union of compact sets, by bounding real
coordinates and finite denominators for matrices and inverses.
Truncating to these sets gives density by \(L^2\) convergence.
The norm bound extends the map, and the integral on the right
is continuous in both arguments since
\(|\langle x,R(u)y\rangle|\le\|x\|_2\|y\|_2\) and \(\Phi\in L^1\)
by Theorem~\ref{thm:absolute-doubled-graph-unfolding}. For an
absolutely convergent pointwise integral the truncations converge
to that integral. They also have a subsequence converging almost
everywhere to the \(L^2\) extension: choose a subsequence with
summable squared errors times \(2^{2k}\) and bound the measures
of the sets with error exceeding \(2^{-k}\) by their squared
integrals times \(2^{2k}\); the sums of the measures of the
tails tend to zero. This proves the identification. For a
rapidly decreasing input and fixed \(b\), the theta value is
bounded by a fixed power of the source height \(H(g)\) on the
reduced domain by Proposition~\ref{prop:adelic-schwartz-frequency-height},
applied to the Hom action with fixed target as in the proof
of Theorem~\ref{thm:global-actual-arrow-covariance}; the Hom
matrix acts by right multiplication by \(g^{-1}\).
Proposition~\ref{prop:unitary-reduced-height-integrability} now
gives absolute convergence there. Finally the finite-level identity
is the second formula of Theorem~\ref{thm:absolute-doubled-graph-unfolding},
using commutation of finite and real translates. The real integral
is bounded on \(L^2\) by Theorem~\ref{thm:archimedean-first-arrow-integrability}.
\end{proof}

\section{The real factor after projection to the source type}

\begin{lemma}\label{lem:real-coefficient-compression}
Let \(G_\infty\) be the product of the real unitary groups with
the measures in Theorem~\ref{thm:row-real-orthogonality}.
At each place suppose the moment hypothesis of that theorem
and take its \(e_\tau\). Put \(e=\prod_\tau e_\tau\) and
\(d_e=e(1)>0\). For a strongly continuous unitary representation
\(R\) of \(G_\infty\) on a Hilbert space put
\[
                  P_e^R y=\int e(a)R(a)y\,\dd a .
\]
Then \(P_e^R\) and \(P_{\bar e}^R\) are orthogonal projections and
\[
       P_e^R R(a)P_e^R=\frac{\overline{e(a)}}{d_e}P_e^R .
\]
In particular for \(x,y\) fixed by \(P_e^R\), the matrix
coefficient \(\langle x,R(a)y\rangle\) is
\(\overline{e(a)}\,\langle x,y\rangle/d_e\).
\end{lemma}
\begin{proof}
The kernels are integrable by Theorem~\ref{thm:polar-row-decay}
and product integration, so these integrals of vectors define
bounded operators of norm at most \(\|e\|_1\). Symmetry under
inversion \(e(z^{-1})=\overline{e(z)}\) and unitarity give
self-adjointness. Apply the second identity of
Theorem~\ref{thm:row-real-orthogonality} separately at each place,
and multiply. In \(P_e^R R(a)P_e^R\) the variable of the right
translate is \(v=zaw\), so its integrated coefficient is
\[
              \int e(z)e((za)^{-1}v)\,\dd z
                      =\frac{e(a^{-1})e(v)}{d_e}.
\]
The double integral converges in vector norm by \(\|e\|_1^2\).
This proves the formula and gives idempotence on setting \(a=1\).
Conjugate the kernel identities for the same assertions about
\(\bar e\). A self-adjoint idempotent projects orthogonally
onto its range. Taking the inner product with two fixed vectors
gives the last identity.
\end{proof}

\begin{theorem}\label{thm:first-arrow-real-projected-norm}
At the finite level and ranks of
Theorem~\ref{thm:actual-first-arrow-global-norm} let \(W_e\)
be the subspace of \(L^2([G])\) fixed by that finite level and
by \(P_e^R\) for the right real action of \(G_\infty\).
On the real oscillator let \(Q_e=P_{\bar e}^{\Omega_\infty}\)
for the actual source action \(a\mapsto\Omega_\infty(a,1)\).
For \(x,y\in W_e\),
\[
 \langle v_f x,v_f y\rangle_{[H]}
     =\frac{{\rm vol}[H]\ \|v_F\|^2\,\|Q_e f_\infty\|^2}
            {W_0(P)\,e(1)}\,\langle x,T_h y\rangle_{[G]} .
\]
If \(h=6r\), the arithmetic scale is sufficiently large, and
\(W_e\ne0\), then the lift is injective on \(W_e\) if and only if
\(Q_e f_\infty\ne0\). If the integral defining \(v_f\) is
multiplied by a scalar \(a_0\), its Gram factor is multiplied
by \(|a_0|^2\).
\end{theorem}
\begin{proof}
The finite translates in \(T_h\) commute with every real
operator, hence with \(P_e^R\), by the norm convergence in
Theorem~\ref{thm:positive-finite-first-arrow-convolution}.
Thus \(P_e^R T_hy=T_hy\). The integral in
Theorem~\ref{thm:actual-first-arrow-global-norm} consequently
uses \(\langle x,R(a)T_hy\rangle=\overline{e(a)}
\langle x,T_hy\rangle/e(1)\) at every \(a\), by
Lemma~\ref{lem:real-coefficient-compression}. Its real scalar is
\[
 \frac1{e(1)}\int \overline{e(a)}
        \langle f_\infty,\Omega_\infty(a,1)f_\infty\rangle\,\dd a
     =\frac{\langle f_\infty,Q_e f_\infty\rangle}{e(1)}
     =\frac{\|Q_e f_\infty\|^2}{e(1)} .
\]
The last equality is orthogonal projection. All finite
coefficients outside this scalar are positive on the diagonal
for nonzero \(y\) at sufficiently large scales by
Theorem~\ref{thm:positive-finite-first-arrow-convolution};
the target volume and the finite test norm are positive.
If the projection vanishes the squared norm is zero for every
\(y\); if the projection is nonzero every nonzero \(y\) has
positive squared norm. Scalar multiplication of the lift gives
the stated factor by sesquilinearity.
\end{proof}

\section{A polar test for nonvanishing of the real projection}

\begin{proposition}\label{prop:polar-double-compact-positivity}
At one real place use \(G,K,\Lambda,\mu,m=m_\Lambda,\psi\)
of Definition~\ref{def:polar-row-coefficient}, and put
\(c=\int_K|m(k)|^2\dd k\). Suppose the moment condition of
Theorem~\ref{thm:row-real-orthogonality}. For a strongly
continuous unitary representation \(R\) and a vector \(v\), put
\[
       v_K=c^{-1}\int_K m(k)R(k)v\,\dd k,\qquad
                            \Phi_v(g)=\langle v,R(g)v\rangle .
\]
For \(g=k_0P\) its ordinary polar decomposition, the mean of
\((\psi\Phi_v)(k_1gk_2)\) on \(K^2\) equals
\[
        c\int_K
           \frac{\langle v_K,R(a^{-1}Pa)v_K\rangle}
                       {m_{-\mu}(a^*Pa)}\,\dd a .
\]
If \(v_K\ne0\) and \(\langle v_K,R(P)v_K\rangle\) is real
and nonnegative for every positive definite hermitian \(P\in G\),
then \(P_{\bar e}^{R}v\ne0\), where
\(e=(\int_G|\psi|^2)^{-1}\bar\psi\).
\end{proposition}
\begin{proof}
Absorb \(k_0\) in \(k_1\). The polar factors of \(k_1Pk_2\)
are \(k_1k_2\) and \(k_2^{-1}Pk_2\), since the latter matrix
is positive hermitian. Insert the integral defining \(\psi\)
and denote its compact integration variable by \(u\).
Set \(a=k_2u,\ b=k_1a\); the three product Haar probabilities
in \(k_1,k_2,u\) become those in \(b,a,u\), by successive
translations and inversion in the compact group. The numerator
and denominator of the integrand, apart from the common \(1/c\),
are then
\[
 m(b)\,\overline{m(u)}
       \langle v,R(b a^{-1}P a u^{-1})v\rangle ,
                     \qquad m_{-\mu}(a^*Pa).
\]
The principal minors of a positive hermitian matrix are positive,
so the denominator is a positive continuous function on the
compact set in question. The compact coefficient satisfies
\(m(u^{-1})=\overline{m(u)}\) by
Definition~\ref{def:cyclic-row-module}. The \(u\)-integral
replaces \(R(u^{-1})v\) with its weighted integral \(c v_K\).
The operator \(\int m(b)R(b)\dd b\) is self-adjoint by inversion
and the same symmetry of \(m\), so the \(b\)-integral on the
other side gives another \(c v_K\). This proves the formula.

The double average is nonnegative under the displayed hypothesis
and equals \(c\|v_K\|^2>0\) at the identity. Continuity makes
it positive on a neighborhood of positive Haar measure.
The product \(\psi\Phi_v\) is integrable by
Theorem~\ref{thm:polar-row-decay} and
\(|\Phi_v|\le\|v\|^2\). Its integral is unchanged by the two
compact averages by bi-invariance of Haar and absolute
integration. Hence \(\int_G\psi(g)\Phi_v(g)\dd g>0\).
By Lemma~\ref{lem:real-coefficient-compression},
\(\|P_{\bar e}^R v\|^2=e(1)\int_G\psi\Phi_v\), proving the
assertion.
\end{proof}

\begin{corollary}\label{cor:polar-real-first-arrow-test}
For a tensor of real oscillator tests at the infinite places in
Theorem~\ref{thm:first-arrow-real-projected-norm}, it suffices
that the two hypotheses on \(v_K\) of
Proposition~\ref{prop:polar-double-compact-positivity} hold for
the local test and the actual local source action at every
place. Under those hypotheses the first-arrow norm on \(W_e\)
is a positive real constant times its finite Gram form.
\end{corollary}
\begin{proof}
For a tensor test the inner product of its translates factors
into the local coefficients by the real place product of
Theorem~\ref{thm:global-actual-arrow-covariance}. The integrals
against \(\bar e\) factor by absolute integrability of each
local \(e_\tau\) and the unitary norm bound. The proof of
Proposition~\ref{prop:polar-double-compact-positivity} gives
a strictly positive real result at each place. Their product
is \(\|Q_e f_\infty\|^2>0\); now apply
Theorem~\ref{thm:first-arrow-real-projected-norm}.
\end{proof}

\section{Positive oscillator coefficients for polynomial tests}

\begin{lemma}\label{lem:signed-oscillator-coherent-coefficient}
Use the Darboux coordinates of
Theorem~\ref{thm:negative-block-real-splitting} and the creation
and annihilation operators \(C_i,A_i\) of
Definition~\ref{def:positive-hermite-tests} on their real
position space. Write \(B_z=\exp(\sum_i z_i C_i)\phi_0\) for
the coherent vector of Proposition~\ref{prop:positive-hermite-theta-series}.
If a signed complex isometry acts in the Darboux complex
coordinate \(\xi=x+i l\) by \(\xi\mapsto U\xi+V\bar\xi\),
put \(B=(U^*)^{-1}\). Here \(U\) is invertible. For its
unitary oscillator implementer \(\mathcal W\), put
\(c_0=\langle\phi_0,\mathcal W\phi_0\rangle\). Then
\[
 \langle B_{\bar z},\mathcal W B_y\rangle =
 c_0\exp\!\left(\tfrac12 z^{\mathsf T}BV^{\mathsf T}z
             +z^{\mathsf T}B y
             -\tfrac12 y^{\mathsf T}\bar V^{\mathsf T}B y\right).
\]
\end{lemma}
\begin{proof}
In the positive and negative signed summands write the original
complex isometry in blocks \(\alpha,\beta,\gamma,\delta\).
In the Darboux coordinate the negative variables are conjugated,
so \(U={\rm diag}(\alpha,\bar\delta)\) and \(V\) consists of the
off-diagonal blocks \(\beta,\bar\gamma\). Both diagonal blocks
are invertible: the isometry equations give
\(\alpha^*\alpha=I+\gamma^*\gamma\) and
\(\delta^*\delta=I+\beta^*\beta\). The vacuum coefficient \(c_0\)
is nonzero by Theorem~\ref{thm:negative-block-real-splitting},
also after a unit scalar twist.

Write \(A(t)=\sum t_i A_i,\ C(t)=\sum t_i C_i\) with bilinear
coefficient convention. Differentiating the symmetric Heisenberg
translation in Lemma~\ref{lem:schroedinger-real-gaussian} at
zero gives \(\sqrt\pi(A(\bar\xi)-C(\xi))\).
Intertwine this derivative by \(\mathcal W\); using first
\(\xi\) and then \(i\xi\) separates the two parts, giving
\[
 \mathcal W C(t)\mathcal W^{-1}=C(Ut)-A(\bar V t),
 \qquad
 \mathcal W A(t)\mathcal W^{-1}=A(\bar U t)-C(Vt).
\]
The difference quotients are valid on Schwartz vectors and
the implementer preserves them, by
Theorem~\ref{thm:negative-block-real-splitting}. Also
\(A(t)B_y=t^{\mathsf T}y B_y\), since
\([A_i,C_j]=\delta_{ij}\) and \(A_i\phi_0=0\), and
\(\partial_{y,t} B_y=C(t)B_y\), where \(\partial_{y,t}\)
denotes directional differentiation. These identities also
follow by differentiating the explicit Gaussian formula of
Proposition~\ref{prop:positive-hermite-theta-series}. That formula
gives entire dependence in Schwartz seminorms: on compact
parameter sets every parameter or position derivative is bounded
by a polynomial times a decaying real Gaussian after completing
the square.

Let \(K(z,y)\) denote the left side. Moving a creation operator
from the second position to the first gives multiplication by
its coefficient paired bilinearly with \(z\), since its adjoint
is annihilation with conjugate coefficients. Moving annihilation
to the first position differentiates in \(z\). Applying the
second intertwining identity to \(B_y\) thus gives
\[
           \nabla_z K=B(y+V^{\mathsf T}z)K .
\]
The first intertwining identity at \(z=0\) gives
\(\nabla_y K(0,y)=-\bar V^{\mathsf T}B y\,K(0,y)\).
Integrate each of these first-order equations along the straight
line from zero to its argument, first in \(y\) at \(z=0\),
then in \(z\) at fixed \(y\). The scalar equation on each
line has solution the exponential of the integral of its
coefficient, by differentiating its product with the inverse
exponential. The initial value \(K(0,0)=c_0\) yields the formula.
\end{proof}

\begin{theorem}\label{thm:harmonic-positive-source-coefficient}
At one real place index the coordinates of \({\rm Hom}(S,T)\)
by a target row \(a\) and a source column \(i\), with the
hermitian signs in Theorem~\ref{thm:negative-block-real-splitting}.
Let \(P\ne0\) be a complex polynomial in these coordinates and
suppose
\[
       \partial_{a i}\partial_{a j}P=0
       \quad\text{whenever the source signs of \(i,j\) differ}.
\]
For the Hermite test \(H_P=P(C)\phi_0\) one then has
\[
       \langle H_P,\Omega_\infty(g,1)H_P\rangle>0
       \quad\text{for every positive hermitian source isometry }g .
\]
The inequality means a strictly positive real number. The space
of polynomials satisfying the derivative condition is invariant
under the source maximal compact group. Hence the nonnegativity
hypothesis of Proposition~\ref{prop:polar-double-compact-positivity}
holds for the compact average of such a test; it is strictly
positive if the compact average is nonzero.
\end{theorem}
\begin{proof}
The Hom action is right multiplication by \(g^{-1}\), which
on each target row is represented on its column coordinates by
\(g^{-{\mathsf T}}\). For \(g\) positive hermitian this matrix
is also positive hermitian. The full Hom isometry is its block
repetition and has positive hermitian diagonal signed blocks
\(\alpha,\delta\). The \(U\) of
Lemma~\ref{lem:signed-oscillator-coherent-coefficient} is thus
positive hermitian, as is \(B=(U^*)^{-1}\). It preserves each
target row and each sign block, while \(V\) only connects
opposite signs in the same target row. Both pure quadratic
forms in that lemma therefore consist entirely of products
of variables of opposite signs in one row. The vacuum factor
is positive: its untwisted value is \(1/\det\delta>0\) by
Theorem~\ref{thm:negative-block-real-splitting}, and the
determinant of a positive hermitian unitary-group element is
one, since it is positive and has modulus one. The determinant
twist of Theorem~\ref{thm:global-actual-arrow-covariance} is
therefore one.

Apply \(\overline P(\partial_z)P(\partial_y)\) at \(z=y=0\)
to the coherent coefficient. Multiplication of any power series
by \(z_i\) under evaluation by \(\overline P(\partial_z)\)
is equivalent to replacing \(\overline P\) by
\(\partial_i\overline P\), as follows on monomials including
their factorials. The same holds on the \(y\) side. Thus any
term containing a factor from either pure quadratic form
vanishes by the stated second derivative condition. Only finitely
many Taylor coefficients enter. What remains is the same
differentiation of \(c_0\exp(z^{\mathsf T}B y)\). Let
\(A=B^{1/2}\) be the positive hermitian square root from
Lemma~\ref{lem:common-hecke-interpolation}. The last exponential
is \(\langle B_{A\bar z},B_{Ay}\rangle\): the coherent inner
product is the exponential of the hermitian pairing by
Theorem~\ref{thm:positive-hermite-circle-projection} and
the factorial expansion. Its differentiated value is the squared
norm of \(P((\sum_j A_{ji}C_j)_i)\phi_0\). The substituted
polynomial is nonzero since \(A\) is invertible; the Hermite
norm formula in the same theorem makes the squared norm positive.
This proves the assertion on positive elements.

On the source compact group the block \(V\) vanishes. The
creation-operator formula of
Lemma~\ref{lem:signed-oscillator-coherent-coefficient} then acts
on a polynomial by linear changes preserving each target row
and the two source-sign blocks, times the scalar on the vacuum.
The chain rule preserves the derivative condition.
A compact average stays in the same finite-dimensional polynomial
space of bounded degree and obeys the same linear conditions.
If it is zero the coefficient is zero, and otherwise the
strict positivity just proved applies.
\end{proof}

\section{A determinant polynomial with the required compact coefficient}

\begin{theorem}\label{thm:one-sided-harmonic-minor-test}
At one infinite place let \(S,T\) have signed orthonormal bases,
with source signs \(s_i\in\{\pm1\}\) and target signs \(t_a\).
Write \(h_\pm=\#\{a:t_a=\pm1\}\).
Use the local source action of Theorem~\ref{thm:global-actual-arrow-covariance}
with determinant twist \((\det g)^\alpha\) on the source, where
\(\alpha\in\Z\). Let \(\Lambda_i\) be the compact row weight
of Definition~\ref{def:polar-row-coefficient}, in the given
parameter order, weakly decreasing separately on the two sign
sets. Define
\[
                    b_i=\Lambda_i+\alpha+h_{-s_i}.
\]
Suppose the indices with \(b_i\ne0\) admit an injection
\(i\mapsto a(i)\) into target rows with
\(t_{a(i)}=s_i\,{\rm sign}(b_i)\).
Then there is an explicit nonzero Hermite test \(H_P\) with
\[
     c^{-1}\int_K m_\Lambda(k)\Omega_\infty(k,1)H_P\,\dd k=H_P,
     \qquad
     \partial_{ai}\partial_{aj}P=0\quad(s_i\ne s_j).
\]
Under the moment condition of Theorem~\ref{thm:row-real-orthogonality}
this test has \(P_{\bar e}^{\Omega_\infty}H_P\ne0\).
\end{theorem}
\begin{proof}
For each sign \(s\) and each integer \(j\ge1\) form the initial
and final segments
\[
       I_{s,j}^+=\{i:s_i=s,\ b_i\ge j\},\qquad
       I_{s,j}^-=\{i:s_i=s,\ b_i\le-j\}.
\]
They have the stated segment property by the compact weight
ordering. For every nonempty segment \(I\) take the determinant
\(\det(Z_{a(i),l})_{i,l\in I}\) in the source order.
Let \(P\) be the product of all these finitely many determinants,
including repetitions. Put \(Z_{a(i),i}=1\) and all other entries
zero to see \(P\ne0\), also for the empty product. A target
row used in the product occurs only against source columns of
one sign by the injection and the definitions of the segments.
This proves the derivative assertion.

We compute the compact projection before multiplying the minors.
For \(k=(k_+,k_-)\in K\), its Hom action on the column vector
of a fixed target row of entries is \(k_s^{-{\mathsf T}}=\bar k_s\)
in a source-sign block \(s\). In signed Darboux coordinates
the creation labels of that row therefore transform by
\(\bar k_s\) when \(s=t_a\), and by \(k_s\) when \(s=-t_a\),
by Lemma~\ref{lem:signed-oscillator-coherent-coefficient}.
The vacuum transforms by \(\prod_s(\det k_s)^{\alpha+h_{-s}}\).
Indeed the negative Hom block contains \(h_{-s}\) copies
of \(\bar k_s\) and its inverse determinant gives the untwisted
vacuum factor by Theorem~\ref{thm:negative-block-real-splitting}.
Since \(k\) is compact the negative block is unitary; the
vacuum coefficient of absolute value one gives the actual
vacuum character.

For a segment of length \(l\), use the \(l\)-th exterior power
of the indicated label representation, and take its unit
coordinate wedge on the segment \(I\). The alternating map
from \(l\) label vectors to the determinant of their entries
in the assigned target rows is equivariant in the source:
apply the same \(k_s\) or \(\bar k_s\) to each label and
expand multilinearly. Equivalently its value on a coordinate
wedge \(i_1\wedge\cdots\wedge i_l\) is the minor with those
columns in the chosen rows. Tensor all the wedges with the
one-dimensional vacuum character and multiply the resulting
polynomials acting on \(\phi_0\) by creation operators.
This gives a linear compact intertwiner sending their tensor
of unit wedges \(v\) to \(H_P\).

The coefficient of \(v\) is the vacuum character times the
product of the conjugate leading minors on the positive
segments and the ordinary trailing minors on the negative
segments. For a unitary matrix a trailing principal minor
equals its determinant times the conjugate of the complementary
leading minor. This follows from the complementary minor and
inverse-transpose identity proved in
Lemma~\ref{lem:positive-height-hyperplane-eigenvalue}.
Thus on each block the coefficient is a determinant character
times conjugate leading minors. On a diagonal element its
exponent at \(i\) is \(\alpha+h_{-s_i}-b_i=-\Lambda_i\).
Such exponents determine the product on the whole block:
the adjacent differences of exponents give the multiplicities
of the proper leading minors and the last exponent fixes the
determinant character. Hence the coefficient is precisely
\(\overline{m_\Lambda(k)}\), by the compact formula of
Theorem~\ref{thm:compact-highest-vector} on both sign blocks.

We verify that this coefficient suffices for the projection.
For its irreducible compact representation with unit vector \(v_0\)
and coefficient \(m_\Lambda\), average the rank-one orthogonal
projection onto \(v_0\) over conjugation by the compact group.
It commutes with the representation, so is scalar: diagonalize
the averaged self-adjoint matrix, whose eigenspaces are invariant,
and use irreducibility. Its \(v_0\) entry is
\(c=\int_K|m_\Lambda|^2\). Taking the coefficient at an arbitrary
group element gives
\(\int_K m_\Lambda(k)m_\Lambda(k^{-1}u)\dd k=c\,m_\Lambda(u)\).
The same computation on the product compact group follows by
multiplying the two block identities. Thus
\(c^{-1}\int_K m_\Lambda(k)R(k)\dd k\) is a self-adjoint
idempotent in any unitary action \(R\), by convolution and
inversion as in Lemma~\ref{lem:real-coefficient-compression}.
In the tensor of wedges its inner product with the unit
vector \(v\) on both sides is one by the coefficient just
computed; the projection therefore fixes \(v\). Passing through
the intertwiner proves it fixes \(H_P\). Theorem~\ref{thm:harmonic-positive-source-coefficient}
and Proposition~\ref{prop:polar-double-compact-positivity} now
give the claimed nonzero real projection.
\end{proof}

\section{Insertion of a positive block of parameters}

\begin{theorem}\label{thm:joint-insertion-harmonic}
Let the source have rank \(n\), signs \(s_i\), and decreasing
parameters \(\lambda_i\), \(0\le i<n\). Let the target have
rank \(H=n+D\), where \(D\ge1\). Fix \(0\le M\le n\).
List the target indices in order as the first \(M\) old indices,
\(D\) new indices, and the remaining old indices. The signs are
\(s_i\) in the prefix, positive on the new block, and \(-s_i\)
in the suffix. Write \(\alpha,\beta\in\Z\) for the exponents
of the extra source and target determinant characters in the
Hom oscillator. Set \(\Delta=\alpha+\beta+D/2\). On the
target take the parameters
\[
   \lambda'_{\rm old(i)}=\lambda_i+\Delta,\qquad
   \lambda'_{\rm new(j)}=\beta-(H+1)/2+D-j,\quad 0\le j<D.
\]
Suppose they are strictly decreasing in the indicated order,
that the compact weights \(\Lambda,\Lambda'\) of
Definition~\ref{def:polar-row-coefficient} are integral and
dominant within their sign blocks, and that
\(b_i=\Lambda_i+\alpha+h_{-s_i}\) is positive for \(i<M\)
and negative for \(i\ge M\), with \(h_t\) the target sign counts.
Use the injection in Theorem~\ref{thm:one-sided-harmonic-minor-test}
sending an index to its old copy in the target. The resulting
nonzero test \(H_P\) is fixed by both compact coefficient averages,
using \(m_\Lambda\) for the source and \(\overline{m_{\Lambda'}}\)
for the target, each divided by its own squared coefficient integral.
Its polynomial satisfies
\[
  \partial_{ai}\partial_{aj}P=0\ \ (s_i\ne s_j),\qquad
  \partial_{ai}\partial_{bi}P=0\ \ (t_a\ne t_b).
\]
\end{theorem}
\begin{proof}
Write \(n_s\) for the number of source indices of sign \(s\),
\(P_s\) for their count in the prefix, and \(U_s=n_s-P_s\)
for their count in the suffix. Then
\(h_t=P_t+U_{-t}+D{\bf1}_{t=+}\), where \({\bf1}\) denotes
the stated indicator. For an index \(i\) let \(B_i\) count the
preceding indices of the same source sign. The weight formula
of Definition~\ref{def:polar-row-coefficient}, now indexed from
zero, reads
\[
           \Lambda_i=\lambda_i+(n+1)/2-i-n_{s_i}+2B_i .
\]
Use the identical formula on the target. Its compact vacuum
exponent for a sign \(t\) is \(a'_t=\beta-n_{-t}\), since
the target acts by left multiplication on the Hom entries and
the negative block for \(k_t\) contains \(n_{-t}\) copies
of \(k_t\). Its inverse determinant gives the untwisted
vacuum character in Theorem~\ref{thm:negative-block-real-splitting}.

For a prefix position \(i\) the position, sign and count of
preceding equal signs stay the same. Hence
\(\Lambda'_{\rm old(i)}-\Lambda_i
=\Delta+D/2-h_{s_i}+n_{s_i}
=\alpha+\beta+h_{-s_i}-n_{-s_i}\).
For a suffix position of old sign \(s\), new sign \(t=-s\),
the new position is \(i+D\) and the preceding equal signs
number \(P_t+D{\bf1}_{t=+}+B_i-P_s\).
Using \(h_t=n_s+P_t-P_s+D{\bf1}_{t=+}\) gives
\[
    \Lambda'_{\rm old(i)}-\Lambda_i
       =\Delta-D/2+h_t-n_s=\alpha+\beta+h_{-s}-n_{-t}.
\]
Thus on every old position \(\Lambda'_{\rm old(i)}=a'_{t}+b_i\).
The \(j\)-th new index has position \(M+j\) and preceding
positive count \(P_++j\). Since \(M+h_+-2P_+=n_-+D\),
substitution of its parameter gives \(\Lambda'_{\rm new(j)}=a'_+\).

The injection satisfies the sign condition in
Theorem~\ref{thm:one-sided-harmonic-minor-test}, which supplies
the source assertion and the first derivative condition. For
the second, a fixed source column can occur only in segments
of its own source sign and its own sign of \(b_i\). All target
rows of those segments therefore have the same target sign.
This proves the other vanishing.

For the target compact action, the creation label transforms
by \(k_t\) on rows of the same sign as the source column and
by \(\bar k_t\) on rows of opposite sign, by the signed
coordinate calculation in
Lemma~\ref{lem:signed-oscillator-coherent-coefficient}.
The rows of the segments with \(b_i\ge j>0\) are initial
in their target-sign set: every row of that sign in the prefix
has that same source sign, and the \(b_i\) decrease within it.
The rows with \(b_i\le-j\) are final in their target-sign set
by the identical order in the suffix, whose source sign is the
opposite one. In the exterior tensor construction of
Theorem~\ref{thm:one-sided-harmonic-minor-test}, now alternating
the row labels for fixed columns, the compact coefficient is
the vacuum character times ordinary leading minors from the
prefix and conjugate trailing minors from the suffix. Its
torus exponents are \(a'_t+b_i\) on old and \(a'_+\) on new
indices. By the weight identities proved above and the same
complementary-minor calculation, the coefficient equals
\(m_{\Lambda'}\). The weighted compact projection argument
of that theorem, with complex conjugation in the averaging
kernel, therefore fixes the polynomial test.
\end{proof}

\section{Real first arrows at the arithmetic parameters}

\begin{definition}\label{def:first-arrow-insertion-parameters}
Fix a sufficiently large arithmetic scale,
\(k,r,m,h_2=r^8,\ t=5r\) and \(w_j(\tau)\) as in
Definition~\ref{def:heights}. For an offset \(0\le o<m\)
and \(z\in\{0,1\}\) put \(n=r+z,\ D=t-z,\ H=n+D=6r\),
\(M=k+o\). At an infinite place write \(W=W_\tau\) for the
rounded weight of Definition~\ref{def:row-vandermonde} and
\(w=w_{o+1}(\tau)\). For \(0\le i<n\) set
\[
 a_i=\max(0,i+1-k),\qquad
 \lambda_i=\frac z2+
 \begin{cases}
 k/2-1-i,& i<k-1,\\
 -W-(i+2-k)2h_2,& i\ge k-1 .
 \end{cases}
\]
For the offset grid choose \(s_i=(-1)^{o+a_i}\).
For the auxiliary group at \(z=0\) choose \(s_i=1\) for
\(i<M\) and \(s_i=-1\) for \(i\ge M\).
Use the target signs of Theorem~\ref{thm:joint-insertion-harmonic}.
The target old parameters are \(\lambda_i-(1-z)/2\).
At its inserted block they are \((t-2)/2-w-j\), \(0\le j<D\).
The two determinant exponents in the actual Hom action are
\[
                   \alpha=w-H/2,\qquad
                   \beta=(n+z-1)/2-w .
\]
In particular the auxiliary target is definite. In the grid
the target sign of an old index is
\((-1)^{o+a_i+\min(\max(a_i-o,0),1)}\); every inserted
direction is positive.
\end{definition}

\begin{theorem}\label{thm:arithmetic-first-arrow-minor-nonzero}
Uniformly in the parameters of
Definition~\ref{def:first-arrow-insertion-parameters}, at
all sufficiently large scales the construction in
Theorem~\ref{thm:joint-insertion-harmonic} applies at every
infinite place. For the resulting tensor of Hermite tests
\(f_\infty\) and the source kernel \(e\) of those parameters,
\[
                             \|Q_e f_\infty\|^2>0 .
\]
The test can be chosen the same for all finite twists at the
fixed signatures, weights and offset. For the initial row
arrow \(Y_0\to Y_1\) of Definition~\ref{def:actual-row-lift-spaces}
these are the data at \(o=z=0\), in the ordered signed frames.
\end{theorem}
\begin{proof}
The exponents are integers since \(r\) is odd. For the ordinary
row and auxiliary arrows the character of weight \(-w\) is
rank-twisted by \(\lfloor H/2\rfloor\) on the source determinant
in the denominator and by \(\lfloor n/2\rfloor\) on the target
in the numerator, as in
Definition~\ref{def:actual-twisted-arrow-character} and
Theorem~\ref{thm:actual-twisted-finite-arrow}.
For the second grid the upper rounding of these half ranks
has the same value as the lower rounding, since \(n=r+1\)
and \(H=6r\) are even. This gives exactly the exponents above.
The source signs of the offset-zero row at \(z=0\) are positive
on the initial block, then alternate starting with a negative
sign at \(i=k\), as in Proposition~\ref{prop:actual-bessel-weights}.
Its target removes that negative sign and alternates thereafter,
with the \(t\) new directions positive; the signature counts
agree with Theorem~\ref{thm:actual-row-lift-frames}.
More generally \(a_i\le o\) holds exactly on \(i<M\);
hence the stated target formula preserves the prefix signs
and reverses the suffix signs at every offset.

Writing \(\Delta=\alpha+\beta+D/2\) gives
\(\Delta=-(1-z)/2\), and
\(\beta-(H+1)/2+D=(t-2)/2-w\). Thus the target parameters
have exactly the insertion formula. For \(i\ge k-1\),
\[
                  \lambda_i+w=z/2+2h_2(o-a_i)+h_2 .
\]
Its values are at least \(h_2\) in the prefix and at most
\(-h_2+1/2\) in the suffix. For \(i<k-1\) one has
\(w=W+2(o+1)h_2+h_2\), \(W\ge0\), and
\(\lambda_i\ge-r\); hence the prefix bound also holds there.
In particular the old prefix lies strictly above all \(D\)
inserted target parameters and the suffix strictly below them:
relative to \(-w\) the inserted parameters are at most \(t\)
in absolute value, and the old shift from \(\lambda_i\) is
at most \(1/2\) in absolute value. Within the old and inserted
lists the parameters decrease by at least one by their formulas.
Both formal parameter lists have the required half-integral
parity for their ranks by the parity of \(r,t\).

For a list of rank \(n_0\) its sign correction in
Definition~\ref{def:polar-row-coefficient} satisfies
\(|\Lambda_i-\lambda_i|\le 3n_0/2\), by bounding the full
half sum and the same-sign difference of counts separately.
For successive indices \(i<j\) within one sign block their
weight difference equals \(\lambda_i-\lambda_j+j-i-2\)
by the formula in Theorem~\ref{thm:joint-insertion-harmonic}.
It is nonnegative since \(\lambda_i-\lambda_j\ge j-i\).
Thus the compact weights are integral and dominant on both
groups. For the source,
\[
    |b_i-(\lambda_i+w)|
       \le 3n/2+H/2 ,
\]
since \(-H/2\le h_{-s_i}-H/2\le H/2\). This is strictly
smaller than \(h_2-1/2\) at large scales. Hence precisely the
prefix has positive \(b_i\). Mapping each index to its old
copy in the target gives the required injection.

Finally, adjacent opposite source signs occur only outside
the consecutive first block or at its end; their parameter
gaps are at least \(2h_2-O(r)\). In the target the only
additional possible changes of sign occur at the two ends
of the inserted block, with gaps at least \(h_2-O(r)\)
by the preceding separation. The constants here are absolute:
the first consecutive block has magnitude at most \(r\),
the inserted block has magnitude at most \(t\) after adding
\(w\), and the old shifts have magnitude at most one.
Thus Lemma~\ref{lem:sign-run-determinant} and
Theorem~\ref{thm:polar-row-decay} apply with \(B=0\),
since \(h_2-O(r)>1+4H^2\) at large scales. If all signs
at a place agree, integrability follows directly by compactness.
Apply Theorems~\ref{thm:joint-insertion-harmonic} and
\ref{thm:harmonic-positive-source-coefficient}, then
Corollary~\ref{cor:polar-real-first-arrow-test} or its
local proof at each place. The determinant data used here
involve only real parameters and signs, so the same choice
works for every allowed finite twist.
\end{proof}

\section{The first transferred rows}

\begin{proposition}\label{prop:deep-first-arrow-norm}
At a nonempty finite set \(M\) of split pairs, in the
first-branch coordinates of Theorem~\ref{thm:first-arrow-finite-coefficient}
use a nonzero matrix test depending on residues modulo
\(\mathfrak p_w^{l_w}\), \(l_w\ge1\), supported on integral
\(H\)-by-\(n\) matrices. Assume that on its support
\(L_w X_w\equiv I_n\bmod\mathfrak p_w^{l_w}\) for a fixed
integral \(n\)-by-\(H\) matrix \(L_w\). Use the vacuum
elsewhere, a unitary determinant character trivial on the
principal subgroups at these depths, and unramified on
determinant units elsewhere. Put
\[
 I_l=\prod_{w\in M}|{\rm GL}_n(O_w/\mathfrak p_w^{l_w})|,
       \qquad D_l=\prod_{w\in M}q_w^{l_w}.
\]
The squared norm of the finite test \(F\) is
\[
                 N_F=D_l^{-Hn}\sum_X|F(X)|^2 .
\]
For \(H/2-2n\ge2\), on the indicated principal source level
and source real type \(e\) one has
\[
       \langle v_f x,v_f y\rangle_{[U(T)]}
          =\frac{{\rm vol}[U(T)]\,N_F\,\|Q_e f_\infty\|^2}
                       {I_l e(1)}\,\langle x,T_H y\rangle .
\]
Here \(T_H\) and its tail estimate use rank \(n\) on the source.
\end{proposition}
\begin{proof}
Each integral residue class in the split matrix model has
volume \(D_l^{-Hn}\) by its number of additive cosets, proving
the norm statement. If the support at \(X\) meets that at \(Xg\),
then \(L_w X_w\) and \(L_w X_w g_w\) are both integral and
congruent to \(I_n\) at the indicated depth. The former has
integral inverse and the quotient gives
\(g_w\equiv I_n\bmod\mathfrak p_w^{l_w}\). Conversely this
congruence preserves every residue of integral matrices and
their integrality in both directions. The determinant character
and other scalar factors are one there by the unitary integral
frame calculation of Theorem~\ref{thm:first-arrow-finite-coefficient}.
Thus the finite coefficient is supported there and equals
\(N_F\rho_H(g_*)\). Reduction of the integral group onto the
invertible residue matrices is surjective: lift the entries
and use the unit determinant to invert the lift by its adjugate.
Consequently its principal subgroup has mass \(I_l^{-1}\).

On this support \(X_w\) is injective over the field since
\(L_wX_w\) is invertible. Hence the nongraph vanishing in
Theorem~\ref{thm:marked-test-kills-nongraph-zero-terms} applies.
Theorem~\ref{thm:high-rank-unitary-siegel-weil} and the proof
of Theorem~\ref{thm:actual-doubled-global-mean} give the same
graph formula, since \(H/2-2n\ge2\) implies \(H>4n\).
The real coefficient is integrable by
Theorem~\ref{thm:archimedean-first-arrow-integrability}; the
absolute tail of the unchanged finite complement is bounded
by Theorem~\ref{thm:finite-vacuum-tail-first-arrow}. The absolute
unfolding of Theorem~\ref{thm:absolute-doubled-graph-unfolding}
therefore applies with the just-computed marked mass \(N_F/I_l\).
Use the compact truncation argument of
Theorem~\ref{thm:actual-first-arrow-global-norm} and compress
on the real type by Lemma~\ref{lem:real-coefficient-compression}
to obtain the asserted formula.
\end{proof}

\begin{theorem}\label{thm:independent-actual-first-images}
For an offset and signature in
Definition~\ref{def:first-arrow-insertion-parameters}
use the real test of Theorem~\ref{thm:arithmetic-first-arrow-minor-nonzero}.
At all sufficiently large scales the lift of
Proposition~\ref{prop:deep-first-arrow-norm} is bounded below
on its prescribed real and finite source subspace. More
precisely its squared norm at \(y\) is at least
\[
 \frac{{\rm vol}[U(T)]\,N_F\,\|Q_e f_\infty\|^2}
      {I_l e(1)}(1-\eta_H)\|y\|^2,\qquad
         \eta_H=\exp(24|\Sigma|\,2^{-H/2+2n})-1 .
\]
In particular choose this real test in the first arrow
\(\mathcal F_0\) of Definition~\ref{def:moderate-inclusion-chain}
for \(Y_0\to Y_1\). For every allowed twist the images
\(\mathcal F_0 y_i\) of the independent source rows
\(y_i=cZ_{s_i}\phi_i\) in Corollary~\ref{cor:independent-actual-short-rows}
are linearly independent. Their number is at least
\((1-4/b_{\rm sc}^2)D\) for the numerical rank \(D\) there.
\end{theorem}
\begin{proof}
Here \(H/2-2n=r-2z\), which is at least two eventually.
The field degree factor \(|\Sigma|\) has polynomial growth in
the arithmetic scale and \(r\) grows at least as a positive
power by Proposition~\ref{prop:thick}. Hence \(\eta_H=o(1)\),
uniformly in all finite twists and levels. By
Theorem~\ref{thm:positive-finite-first-arrow-convolution},
\(\langle y,T_Hy\rangle\ge(1-\eta_H)\|y\|^2\).
All the displayed prefactors are strictly positive by
Theorem~\ref{thm:arithmetic-first-arrow-minor-nonzero} and
Proposition~\ref{prop:deep-first-arrow-norm}. This proves
the norm bound.

At \(o=z=0\) every selected row is in the actual real source
type. Indeed Theorem~\ref{thm:actual-full-poincare-columns}
reproduces the Poincar\'e functions by the product of the
actual \(e_J\), which are transports of \(e\) by
Theorem~\ref{thm:actual-real-bessel-vector}. The coordinate
combinations, central frame change and marked-group operators
forming the rows in Proposition~\ref{prop:identify-actual-short-row-spans}
are linear combinations and translations at finite places, so
commute with that real projection. These rows have principal
level at the original marked places and integral invariance
at the twist places, thus also the required principal invariance
on their union. The first finite test is the inclusion test of
Corollary~\ref{cor:first-arrow-level-masses}. The scalar
\([K_0:J_0]\) multiplying the theta kernel in the definition
of \(\mathcal F_0\) is finite and positive. The bounded
source rows have absolutely convergent lift integrals by
Theorem~\ref{thm:inclusion-chain-transposition}. They agree
with the \(L^2\) extension by
Theorem~\ref{thm:actual-first-arrow-global-norm}. Applying the
lower bound to every linear combination of them and using their
source independence proves the conclusion and the stated count.
\end{proof}

\section{Complementary frames over a finite ring}

\begin{definition}\label{def:complementary-finite-frames}
Let \(R\) be a finite commutative ring with unit, of order \(Q\).
Write \(M_{i,j}\) for the \(i\)-by-\(j\) matrices over \(R\)
and \(G_i={\rm GL}_i(R)\). Fix a coordinate decomposition
\(J=R^a\oplus R^d\) of a free module, with inclusion
\(y_0:R^a\to J\) and projection \(x_0:J\to R^d\).
The two summands may occupy interleaved positions in an ordered
basis of \(J\). For each \(k\in{\rm GL}(J)\) put
\(Y_k=ky_0,\ X_k=x_0k^{-1}\).
Fix a character \(\chi:R^\times\to\C^\times\).
A pair of functions \(\Phi\) on \({\rm Hom}(J,R^d)\) and
\(\Psi\) on \({\rm Hom}(R^a,J)\) is matched if they are
supported on the \(X_k\) and \(Y_k\), respectively, and
\[
                 \Phi(X_k)=\chi(\det k)\Psi(Y_k)
                    \qquad(k\in{\rm GL}(J)).
\]
Given \(c\in\C\) and functions \(L,L'\) on the same respective
sets as \(\Phi,\Psi\), the exchange identity on this finite
ring means
\[
              \sum_Y\Psi(Y)L'(Y)=c\sum_X\Phi(X)L(X)
                 \qquad\text{for every matched pair}.
\]
\end{definition}

\begin{proposition}\label{prop:complementary-frame-orbit-exchange}
In Definition~\ref{def:complementary-finite-frames} assume
\(L(hX)=\chi(\det h)L(X)\) and \(L'(Yg)=\chi(\det g)L'(Y)\)
for all \(h\in G_d,\ g\in G_a\).
The exchange identity is equivalent to the pointwise equations
\[
       |G_a| L'(Y_k)=c\,|G_d|\,\chi(\det k)L(X_k)
                        \qquad(k\in{\rm GL}(J)).
\]
For the exchange identity it suffices to test one pair per
submodule \(S=k y_0(R^a)\).
\end{proposition}
\begin{proof}
Fix \(k_*\) and put \(Y_*=Y_{k_*},X_*=X_{k_*}\).
An invertible \(k\) gives the same submodule precisely when
\(k_*^{-1}k\) stabilizes the first summand. In the decomposition
of Definition~\ref{def:complementary-finite-frames} it then has
blocks \(\left(\begin{smallmatrix}g&v\\0&u\end{smallmatrix}\right)\)
with \(g\in G_a,\ u\in G_d\). Indeed both diagonal blocks are
invertible since their determinants have invertible product,
and the converse follows by block inversion. Multiplication gives
\(Y_k=Y_*g,\ X_k=u^{-1}X_*\). The two orbits determine the same
submodule since \({\rm im}\,Y_k=\ker X_k\). Their respective
parameters \(g,h=u^{-1}\) are unique since \(Y_*\) is split
injective and \(X_*\) split surjective.

Take \(\Psi(Y_*g)=\chi(\det g)^{-1}\) and
\(\Phi(hX_*)=\chi(\det k_*)\chi(\det h)^{-1}\),
and give both functions value zero off these orbits.
The block determinant verifies matching, including its support
condition. Every matched pair on these orbits is a multiple
of this one. Indeed one can choose \(g,u\) independently in
the block formula and then set first \(u=1\) and next \(g=1\).
Running through the finitely many \(S\) therefore gives a
basis of matched pairs. Equivariance cancels each character
in the two sums. The orbit for \(Y_*\) contributes \(|G_a|L'(Y_*)\)
and that for \(X_*\) contributes \(|G_d|\chi(\det k_*)L(X_*)\).
This gives both implications, including \(c=0\).
\end{proof}

\begin{theorem}\label{thm:finite-frame-descent-step}
Let \(a,n,t\) be nonnegative integers, \(d=n+t\), and write
\(J=R^n\oplus R^a\oplus R^t\) in the displayed order. In
Definition~\ref{def:complementary-finite-frames} take the
middle summand first and its ordered complement \(R^n\oplus R^t\)
second. Let \(L,L',c\) satisfy the equivariances and exchange
identity of Proposition~\ref{prop:complementary-frame-orbit-exchange}.
For \(y\in M_{n,a}\), \(Z\in M_{d,a+t}\), put
\[
     Y(y)=\begin{pmatrix}y\\ I_a\\0\end{pmatrix},\qquad
     X(Z)=\left(\begin{pmatrix}I_n\\0\end{pmatrix}\ \ Z\right).
\]
For \(p:M_{n,a}\to\C\), define \(p^+\) on \(M_{d,a+t}\)
by block rows \(n,t\) and columns \(a,t\):
\[
 p^+\begin{pmatrix}-y&u\\0&I_t\end{pmatrix}
       =Q^{-nt}p(y) \quad(u\in M_{n,t}),
\]
with value zero on all other matrices. Then
\[
       |G_a|\sum_y p(y)L'(Y(y))
          = c\,|G_d|\sum_Z p^+(Z)L(X(Z)).
\]
Only the pointwise equations of
Proposition~\ref{prop:complementary-frame-orbit-exchange}
for submodules \(Y(y)R^a\) are needed.
\end{theorem}
\begin{proof}
The substitution \(k(v,b,w)=(v+yb,b,w)\) has determinant one,
first columns \(Y(y)\) in the specified decomposition, and quotient
map \(X_k=\left(\begin{smallmatrix}I_n&-y&0\\0&0&I_t\end{smallmatrix}\right)\).
The pointwise equation gives \(|G_a|L'(Y(y))=c|G_d|L(X_k)\).
Multiplication of \(X_k\) on the left by
\(\left(\begin{smallmatrix}I_n&u\\0&I_t\end{smallmatrix}\right)\)
inserts \(u\) into its last columns. Its determinant is one,
so this leaves the value of \(L\) unchanged. Average over the
\(Q^{nt}\) choices of \(u\), multiply by \(p(y)\), and sum over
\(y\). The sign in the definition of \(p^+\) is the one in
the inverse of \(k\).
\end{proof}

\section{The Fourier scalar for a complementary pair}

\begin{definition}\label{def:residue-matrix-fourier}
Use the finite ring of Definition~\ref{def:complementary-finite-frames}.
An additive character \(\psi:R\to\C^\times\) is primitive if for
each \(v\ne0\) there is \(z\in R\) with \(\psi(vz)\ne1\).
A character \(\chi:R^\times\to\C^\times\) is primitive here if
for each \(v\ne0\) some \(z\) gives \(1+zv\in R^\times\) and
\(\chi(1+zv)\ne1\). Extend \(\chi\) and \(\chi^{-1}\) by
zero on nonunits. All nonzero character values have absolute
value one since the underlying groups are finite. If
\(f:M_{a,b}\to\C\), put
\[
        \widehat f(Y)=\sum_{V\in M_{a,b}} f(V)\psi({\rm tr}(YV))
                      \qquad (Y\in M_{b,a}).
\]
All finite Fourier transforms in this section use this normalization.
\end{definition}

\begin{lemma}\label{lem:finite-determinant-fourier-stabilizer}
Let \(\psi,\chi\) be primitive in
Definition~\ref{def:residue-matrix-fourier}. For \(t\ge0\) set
\(D(V)=\chi(\det V)\) on \(M_{t,t}\) and \(G=\widehat D(I_t)\).
If \(S(gV)=\chi(\det g)S(V)\) for all \(g\in G_t\), then
\[
 S(V)=S(I_t)D(V),\qquad
 \widehat D(Y)=\chi^{-1}(\det Y)G,\qquad |G|^2=Q^{t^2}.
\]
\end{lemma}
\begin{proof}
When \(V\) is invertible the first equation is the given transformation
applied to the identity. Otherwise right multiplication by \(V\)
on the finite set of row vectors cannot be injective. If it were
injective it would be onto, giving a left inverse to \(V\) by
taking preimages of the standard rows. Taking determinants and
using the adjugate would make \(V\) invertible. Choose a nonzero
row \(u\) with \(uV=0\), and choose \(j\) with \(u_j\ne0\).
For the \(j\)-th coordinate column \(e_j\), the matrix
\(g=I+ze_j u\) has determinant \(1+zu_j\) by expansion in row
\(j\). Primitivity gives a choice of \(z\) making \(g\) invertible
and \(\chi(\det g)\ne1\). Since \(gV=V\), the transformation
forces \(S(V)=0\).

Changing variables in the Fourier sum gives
\(\widehat D(Yg)=\chi(\det g)^{-1}\widehat D(Y)\).
Apply the first assertion to \(S(V)=\widehat D(V^{\mathsf T})\)
with character \(\chi^{-1}\), which has the same primitivity.
This gives the middle formula. For any rectangle, summing
\(\psi({\rm tr}(YZ))\) over all \(Y\) gives the size of
that matrix set when \(Z=0\), and zero otherwise. In the
second case, translate \(Y\) in one entry paired with a nonzero
entry of \(Z\) to multiply the sum by a scalar different from
one. Expanding the square of the transform and using this
orthogonality gives
\(\sum_Y|\widehat f(Y)|^2=Q^{ab}\sum_V|f(V)|^2\).
For \(D\) the latter sum before the factor is \(|G_t|\);
the middle formula makes the former \(|G|^2|G_t|\).
This proves the norm, also when \(t=0\), in which case there
is one matrix and its determinant is one.
\end{proof}

\begin{theorem}\label{thm:finite-fourier-descent-step}
Let \(n,t\ge0,\ a=n+t\), and use \(X(Z),Y(y),L,L',c\) of
Theorem~\ref{thm:finite-frame-descent-step} with \(d=a\).
Let \(\chi,\psi\) satisfy Lemma~\ref{lem:finite-determinant-fourier-stabilizer}.
In the notation of that lemma put \(\gamma=(Q^{nt}G)^{-1}\).
For every \(p:M_{a,n}\to\C\), define its extension \(p^+\) on
\(M_{a+t,n+t}\) by the rule of
Theorem~\ref{thm:finite-frame-descent-step} with \(a,n\)
interchanged. Then
\[
 \begin{split}
     \sum_{y\in M_{n,a}}\widehat{\bar p}(y)L'(Y(y))
         &=\gamma c\sum_{Z\in M_{a,a+t}}
                         \widehat{\overline{p^+}}(Z)L(X(Z)),\\
              |\gamma|^2 Q^{\,t(t+2n)}&=1.
 \end{split}
\]
Here too it suffices to assume the pointwise incidence equations
on the graph submodules listed in
Theorem~\ref{thm:finite-frame-descent-step}.
\end{theorem}
\begin{proof}
Write \(q(y)=\widehat{\bar p}(y)\). At
\(Z=\left(\begin{smallmatrix}A&B\\C&D_0\end{smallmatrix}\right)\)
use block rows \(n,t\) and columns \(a,t\). The trace after
multiplication by
\(\left(\begin{smallmatrix}-v&w\\0&I_t\end{smallmatrix}\right)\)
is \(-{\rm tr}(Av)+{\rm tr}(Cw)+{\rm tr}(D_0)\).
Averaging over \(w\in M_{a,t}\) cancels the factor \(Q^{-at}\)
on the support of \(p^+\) when \(C=0\) and gives zero otherwise,
by the entrywise orthogonality in
Lemma~\ref{lem:finite-determinant-fourier-stabilizer}.
Thus the transform is \({\bf1}_{C=0}q(-A)\psi({\rm tr}D_0)\).

For fixed \(A,B\), the function
\[
 S(D_0)= L\begin{pmatrix}I_n&A&B\\0&0&D_0\end{pmatrix}
\]
satisfies \(S(gD_0)=\chi(\det g)S(D_0)\). This follows by
acting on the bottom rows with \(g\).
Lemma~\ref{lem:finite-determinant-fourier-stabilizer}
evaluates its sum against \(\psi({\rm tr}D_0)\) as \(GS(I_t)\).
At \(D_0=I_t\) the block \(B\) can be removed by a left shear
of determinant one. The sum over \(B\) therefore gives \(Q^{nt}\)
copies of the same value. Changing \(A\) to \(-y\), the full
sum over \(Z\) on the right before multiplication by \(\gamma c\)
equals
\[
       GQ^{nt}\sum_y q(y)
                  L\begin{pmatrix}I_n&-y&0\\0&0&I_t\end{pmatrix}.
\]
The pointwise equation in the proof of
Theorem~\ref{thm:finite-frame-descent-step} now gives the
assertion, since \(|G_a|=|G_d|\). The displayed norm of
\(\gamma\) follows from \(|G|^2=Q^{t^2}\).
\end{proof}

\section{The marked Fourier normalization and its conductor}

\begin{definition}\label{def:marked-residue-transform}
Let \(S\) be a nonempty finite set of chosen branches of split
places, with valuation rings \(O_w\), uniformizers \(\pi_w\),
residue field orders \(q_w\) and integers \(l_w\ge1\).
Write \(O_S=\prod_{w\in S}O_w,\ \delta=(\pi_w^{l_w})_w\),
\(R=O_S/\delta O_S\), and \(Q_R=|R|=\prod_w q_w^{l_w}\).
The cardinality follows by the successive quotients of the
powers of a uniformizer, each identified with the residue field
by division by that power. Let \(\psi_w^0\) be the normalized
additive character \(\psi_f^0\) of Proposition~\ref{prop:additive}
restricted to the component at \(w\). Fix arbitrary
\(\epsilon_w\in O_w^\times\), using \(\epsilon_w=1\)
unless specified, and put
\[
        \psi_R(v)=\prod_{w\in S}
                 \psi_w^0(\epsilon_w\pi_w^{-l_w}\widetilde v_w),
                     \qquad v\in R ,
\]
where \(\widetilde v_w\) is any integral lift. For
\(p:M_{a,b}(R)\to\C\) let \(\widetilde p(X)\) take the
value \(p(X\bmod\delta)\) on \(M_{a,b}(O_S)\) and zero
outside it. Define
\[
       \mathcal U_{a,b}p=Q_R^{-ab/2}\widehat{\bar p}
                   \quad\text{on } M_{b,a}(R)
\]
using \(\psi_R\) in Definition~\ref{def:residue-matrix-fourier}.
\end{definition}

\begin{proposition}\label{prop:marked-lattice-unitary-fourier}
The character \(\psi_R\) in Definition~\ref{def:marked-residue-transform}
is well defined and primitive. With the integral matrix box of
mass one, the additive Fourier integral of \(\overline{\widetilde p}\)
at a dual matrix \(Z\) is zero outside \(\delta^{-1}M_{b,a}(O_S)\).
At \(Z=\delta^{-1}Y\) for integral \(Y\) it equals
\(Q_R^{-ab}\widehat{\bar p}(Y\bmod\delta)\).
After the norm-preserving change from \(Z\) to \(\delta Z\)
the transformed function is the integral lift of
\(\mathcal U_{a,b}p\). In the finite counting norm and in the
finite probability norm this last transform is an antilinear
isometry onto the reversed matrix space.
\end{proposition}
\begin{proof}
The annihilator of \(O_w\) under multiplication followed by
\(\psi_w^0\) is \(O_w\). To see this, vary only the \(w\)-entry
of the integral argument in Proposition~\ref{prop:additive};
conjugating the other variable permutes the branches and
their integral lattices. Hence changes of integral lift do
not alter \(\psi_R\). For \(v_w\notin\pi_w^{l_w}O_w\), the
annihilator assertion gives \(z_w\in O_w\) for which
\(\psi_w^0(\pi_w^{-l_w}\widetilde v_w z_w)\ne1\).
Use zero at all other entries to obtain primitivity.

The Fourier integral uses
\(\prod_w\psi_w^0(\epsilon_w{\rm tr}(Z_wX_w))\). Multiplication
by a unit preserves each of the annihilator lattices just computed. Translating \(X\)
by \(\delta B,\ B\in M_{a,b}(O_S)\), fixes the test. If
\(\delta Z\) has a nonintegral entry the annihilator assertion
in that entry makes the integral zero. For integral
\(Y=\delta Z\) split \(X\) into its \(Q_R^{ab}\) equal-volume
residue classes. The phase on each class is constant, giving
the stated sum. Replacing \(Z\) by \(\delta^{-1}Y\) multiplies
the measure by \(Q_R^{ab}\), by the index of the new box.
Consequently the norm-preserving substitution multiplies
the function value at \(\delta^{-1}Y\) by \(Q_R^{ab/2}\).
This proves the normalization. In the finite spaces the norm
identity follows directly from the orthogonality computation
in Lemma~\ref{lem:finite-determinant-fourier-stabilizer},
which uses only additive primitivity. Injectivity gives
surjectivity by equal finite dimensions. Dividing both
counting measures by \(Q_R^{ab}\) preserves the isometry.
\end{proof}

\begin{proposition}\label{prop:primitive-square-level-characters}
In Definition~\ref{def:actual-twisted-arrow-character}, let
\(M_g\) be the marked places and let the twist at each \(w\in P\)
satisfy \(\vartheta_w^j\ne1\) for \(1\le j\le m\). This is
the \(m\)-fold twist condition. Use the squarefree ring on
\(S=M_g\cup P\) in Definition~\ref{def:marked-residue-transform}.
For \(1\le j\le m\) and any integer \(u\) with
\(0<|u|\le8m\), the character with components
\(\omega_w^u\) on \(M_g\) and \(\vartheta_w^{\,j}\) on \(P\)
is primitive for all sufficiently large scales. Its inverse
is primitive as well.
\end{proposition}
\begin{proof}
Proposition~\ref{prop:actual-marked-torus-characters} gives
\(\omega_w\) order \(M_\omega=(Md)^{s+6}\).
Proposition~\ref{prop:thick} gives \(2^s=m^{22}\);
since \(Md\ge2\), eventually \(M_\omega>8m\).
Thus every component is nontrivial, including those in \(P\)
by hypothesis. The sets are disjoint by
Definition~\ref{def:split-inclusion-finite-oscillator}.
For \(v\ne0\) in the product of residue fields choose \(w\)
with \(v_w\ne0\) and choose a unit \(b\) there on which its
character is different from one. Put \(z_w=(b-1)/v_w\)
and put \(z_{w'}=0\) elsewhere. Then \(1+zv\) is a unit
and its character value equals that at \(b\).
This is exactly the required primitivity, also after inversion.
\end{proof}

\begin{theorem}\label{thm:rank-independent-unitary-square-scalar}
With \(a=n+t\) in Theorem~\ref{thm:finite-fourier-descent-step},
use the normalized transform of
Definition~\ref{def:marked-residue-transform} on both data.
Then
\[
 \sum_y (\mathcal U_{a,n}p)(y)L'(Y(y))
   =c\,\Gamma\sum_Z(\mathcal U_{a+t,n+t}p^+)(Z)L(X(Z)),
 \qquad |\Gamma|^2=Q_R^{t^2}.
\]
In particular the squared modulus of the product of these
Fourier multipliers over the pairs \(i,j\ge0,\ i+j+2\le m\)
is \(Q_R^{\,t^2\binom m2}\), independently of the intermediate
ranks. On the squarefree arithmetic level the ratio of this
squared modulus to its value at the trivial twist is
\(Q_P^{\,t^2\binom m2}\), where \(Q_P=\prod_{w\in P}q_w\).
\end{theorem}
\begin{proof}
Multiply the Fourier identity of
Theorem~\ref{thm:finite-fourier-descent-step} by \(Q_R^{-an/2}\).
Proposition~\ref{prop:marked-lattice-unitary-fourier} gives the
new multiplier
\[
 \Gamma=\gamma Q_R^{((a+t)(n+t)-an)/2}
              =\gamma Q_R^{t(n+t)} .
\]
Since \(|\gamma|^2=Q_R^{-t(t+2n)}\), the exponent in
\(|\Gamma|^2\) is \(2t(n+t)-t(t+2n)=t^2\).
There are \(\sum_{i=0}^{m-2}(m-i-1)=\binom m2\) pairs,
including zero pairs when \(m=1\). Proposition~\ref{prop:primitive-square-level-characters}
checks multiplicative primitivity for the indicated squarefree
characters, and Proposition~\ref{prop:marked-lattice-unitary-fourier}
checks the additive character. Removing \(P\) divides \(Q_R\)
by \(Q_P\), giving the claimed ratio. The proportionality
constant \(c\) in the exchange identity remains in the formula.
\end{proof}

\section{Adjunction on rational and real theta data}

\begin{definition}\label{def:negative-hermitian-adjoint}
Let \(S,T\) be nonsingular hermitian spaces over the CM field
\(E\), with matrices \(J_S,J_T\). Use
\(V={\rm Hom}_E(S,T),\ V'={\rm Hom}_E(T,S)\),
with trace hermitian forms \(H,H'\) as in
Definition~\ref{def:hom-oscillator-pairing}. Put
\[
 X^\dagger=J_S^{-1}X^*J_T,\qquad \sigma X=-X^\dagger .
\]
On \(V,V'\) denote by \(b,b'\) the rational bilinear forms of
Definition~\ref{def:rational-polarized-oscillator} and by
\(a,a'\) their alternating parts. Write
\(G_{g,h}(X)=hXg^{-1}\) for \(g\in U(S),\ h\in U(T)\);
on \(V'\) use \(G'_{h,g}(Z)=gZh^{-1}\).
\end{definition}

\begin{proposition}\label{prop:rational-adjoint-conjugates-theta}
The map \(\sigma\) in Definition~\ref{def:negative-hermitian-adjoint}
is a rational linear isomorphism with
\[
 H'(\sigma X,\sigma Y)=\overline{H(X,Y)},\qquad
 a'(\sigma X,\sigma Y)=-a(X,Y),\qquad
 \sigma G_{g,h}=G'_{h,g}\sigma .
\]
It induces an antiunitary bijection \(C_\sigma\) on rational
Schr\"odinger test spaces, characterized in the symmetric
coordinates of Theorem~\ref{thm:rational-theta-symplectic-action}
by
\[
             F^{a'}_{C_\sigma f}(\sigma v)
                         =\overline{F^a_f(v)} .
\]
In particular the theta value at zero is conjugated. The map
intertwines the symmetric translations at \(v\) and \(\sigma v\).
\end{proposition}
\begin{proof}
Taking the adjoint twice gives the original map. Moreover
\({\rm tr}(XY^\dagger)={\rm tr}(Y^\dagger X)\), by expanding
both traces entry by entry. This proves the first identity,
and \(\bar\jmath=-\jmath\) gives the second by the trace formula
in Lemma~\ref{lem:rational-graded-polarization}. The adjoints
of \(g,h\) are their inverses, proving the third identity.

Choose complementary rational isotropic spaces \(P,L\) for
\(a\), with dual bases as in Lemma~\ref{lem:rational-graded-polarization}.
Use \(\sigma P,\sigma L\) on \(V'\) with the transported basis
on the first and the negative of the transported basis on the
second. These are dual by the alternating-form identity. On
positions define \((C_\sigma f)(\sigma x)=\overline{f(x)}\).
In these coordinates it preserves the \(L^2\) norm and takes
tests bijectively onto tests. The position coordinate of
\(\sigma(x+l)\) is \(x\) and its momentum coordinate is \(-l\).
Consequently the symmetric translation formula in
Theorem~\ref{thm:rational-theta-symplectic-action} intertwines
under conjugation. The bijection of rational positions gives
\(\theta_{\sigma P}(C_\sigma f)=\overline{\theta_P(f)}\) by
absolute convergence from
Theorem~\ref{thm:rational-schroedinger-zero-section}. Apply this
identity after translation to get the section formula. The
injectivity and change of polarization in
Theorem~\ref{thm:rational-theta-symplectic-action} prove
characterization and independence of the rational coordinates.
\end{proof}

\begin{theorem}\label{thm:real-reversal-of-polynomial-test}
At one real place choose signed orthonormal coordinates on
\(S,T\), with signs \(s_j,t_i\), respectively. In the reversed
coordinates define the real orthogonal signed permutation
\(B\) by \((Bz)_{ji}=-s_jt_i z_{ij}\).
If \(H_P\) is the Hermite test of
Definition~\ref{def:positive-hermite-tests}, conjugate the
coefficients of \(P\) to form \(\bar P\) and put
\[
              P^\#(Z')=\bar P(B^{\mathsf T}Z').
\]
There is an antiunitary real-place intertwiner \(C_\infty\)
implementing \(\sigma\) on the symmetric translations,
sending \(H_P\) to \(H_{P^\#}\). It intertwines the actual
untwisted real splittings for \(G_{g,h},G'_{h,g}\).
For unitary determinant characters \(a_0,b_0\), let
\(\lambda=a_0(\det h)/b_0(\det g)\). The same operator
intertwines the first splitting multiplied by \(\lambda\)
with the second multiplied by \(\bar\lambda\).
\end{theorem}
\begin{proof}
The adjoint formula of Definition~\ref{def:negative-hermitian-adjoint}
in these bases is the signed conjugation \(z\mapsto B\bar z\).
A position \(ij\) and its reversed position \(ji\) have identical
Hom sign \(s_jt_i\). Thus even after the negative coordinates
are conjugated in the Darboux convention of
Theorem~\ref{thm:negative-block-real-splitting}, the action is
\(\xi=x+i l\mapsto B\bar\xi\). On real position functions set
\((C_\infty f)(x')=\overline{f(B^{\mathsf T}x')}\).
It is antiunitary, fixes the standard Gaussian and interchanges
the translated coordinates \((x,l)\) with \((Bx,-Bl)\), by
Lemma~\ref{lem:schroedinger-real-gaussian}. Each creation
operator is a real linear combination of its position and
the corresponding derivative. Changing signs and permuting
variables in that formula gives the substitution \(B^{\mathsf T}Z'\)
and conjugates each coefficient. This proves the Hermite formula.

By Proposition~\ref{prop:rational-adjoint-conjugates-theta},
conjugation by this operator transports an implementer of \(G_{g,h}\)
to one of \(G'_{h,g}\). In the original signed complex coordinates
the negative-to-negative block of the latter is \(B_-\bar\delta
B_-^{-1}\), where \(\delta\) is that block for \(G_{g,h}\)
and \(B_-\) is the restriction of \(B\) to the negative positions.
Its inverse determinant is the conjugate of the vacuum coefficient
for \(G_{g,h}\). Since \(C_\infty\) fixes the vacuum, uniqueness
with this coefficient in Theorem~\ref{thm:negative-block-real-splitting}
gives equality of the transported and reversed real splittings.
The determinant ratios are complex conjugates because the
characters are unitary, which proves the last assertion.
\end{proof}

\section{The finite adjoint and the marked shift}

\begin{proposition}\label{prop:finite-adjoint-conjugates-splitting}
Use the self-dual finite stage frames of
Theorem~\ref{thm:actual-row-lift-frames} on \(S,T\), the scale
\(\alpha\) of Lemma~\ref{lem:actual-lift-scale}, and \(\sigma\)
of Definition~\ref{def:negative-hermitian-adjoint}. Set
\(\eta=(\alpha/\bar\alpha)\sigma\) at the finite places.
On the finite oscillator spaces there is an antiunitary map
\(C_\eta\) taking the reference vacuum to the reversed
reference vacuum and intertwining symmetric translations through
\(\eta\). For every finite pair \(g,h\) one has
\[
        C_\eta\,\mathcal V_{G_{g,h}}
              =\mathcal V'_{G'_{h,g}}\,C_\eta ,
\]
where the splittings are those of
Theorem~\ref{thm:all-finite-weil-splitting}.
\end{proposition}
\begin{proof}
In the finite isometric frames the adjoint is ordinary conjugate
transpose. Hence \(\eta(\alpha X)=-\alpha X^*\), and \(\eta\)
takes the reference lattice \(L\) bijectively onto \(L'\).
It still conjugates the hermitian forms and negates the alternating
form, since \(\alpha/\bar\alpha\) has norm one. In particular
\(b'(\eta v,\eta v)=b(v,v)\). The symmetric bilinear form
\(b(v,z)+b'(\eta v,\eta z)\) has diagonal \(2b(v,v)\);
its symmetry follows because the alternating parts cancel.

Write \({\bf e}_f\) for the finite part of the additive character of
Definition~\ref{def:rational-polarized-oscillator}. Use the full
finite lattice realization of
Theorem~\ref{thm:finite-oscillator-lattice-transport} and set
\[
       (C_\eta F)(\eta v)
                  ={\bf e}_f(-b(v,v))\,\overline{F(v)} .
\]
With \(D_z^a={\bf e}_f(b(z,z)/2)D_z\), expand
\(b(v+z,v+z)-b(v,v)-b(z,z)\) using the symmetric sum above.
Its value is \(b(v,z)+b'(\eta v,\eta z)\). Substitution into
the translation formula of Lemma~\ref{lem:theta-positive-translation}
gives \(C_\eta D_z^a=D_{\eta z}^{a'} C_\eta\).
The phase is one on \(L\), since \(B_f\) is one on \(L\times L\)
by Theorem~\ref{thm:finite-oscillator-lattice-transport}.
Thus the initial vacuum goes to the new vacuum. Additive
measure in each space is normalized on its vacuum lattice;
\(\eta\) preserves this normalization. The translates span both
test spaces by Definition~\ref{def:split-inclusion-finite-oscillator},
so the displayed map is an antiunitary bijection.

The transported \(W_{G_{g,h}}\) and the reversed \(W'_{G'_{h,g}}\)
implement the same Heisenberg change by
Proposition~\ref{prop:rational-adjoint-conjugates-theta}; the
extra scalar in \(\eta\) commutes with the Hom action. Both have
positive real vacuum coefficient, since only the identity coset
in the vacuum formula of Theorem~\ref{thm:finite-oscillator-lattice-transport}
has support meeting the initial box. Its uniqueness up to a unit
scalar identifies the operators. To include the splitting phase,
the definition of \(J_G\) in
Definition~\ref{def:finite-splitting-multiplier} counts the
valuation of the domain lattice on which \(G\) is integral.
On coordinates divided by \(\alpha\), the map
\(X\mapsto-X^*\) exchanges the two conjugate places and permutes
integral entries. It therefore sends the domain lattice for \(G\)
to that for \(G'\), exchanging \(w\) and \(\bar w\) in its index.
Consequently \(J_{G'}=\overline{J_G}\). The ideal phase of
Theorem~\ref{thm:actual-global-splitting-character} takes conjugate
values on these two ideals. This proves the asserted identity.
\end{proof}

\begin{theorem}\label{thm:finite-adjoint-residue-test}
In Proposition~\ref{prop:finite-adjoint-conjugates-splitting} let
\(a=\dim_E T,\ b=\dim_E S\). At a marked set \(M\) of split branches
disjoint from \(\bar M\) use the depths of
Definition~\ref{def:marked-residue-transform}. Extend its \(\delta\)
to a finite idele by one off \(M\), and let \(s_\delta=\bar\delta/\delta\).
Write \(v_D\) for the finite test with first-branch matrix function
\(\widetilde D\) in coordinates divided by \(\alpha\), and vacuum
elsewhere. Use analogous notation \(v'_D\) in the reversed pair.
For \(D:M_{a,b}(R)\to\C\) one has
\[
                  C_\eta v_D=W'_{s_\delta}v'_{\mathcal U_{a,b}D}.
\]
Here \(W'_{s_\delta}\) implements multiplication by \(s_\delta\)
on the reversed Hom space with positive vacuum coefficient,
and \(\mathcal U_{a,b}\) uses
\(\epsilon_w=(-\varpi\jmath\alpha\bar\alpha)_w\) in
Definition~\ref{def:marked-residue-transform}.
\end{theorem}
\begin{proof}
The specified \(\epsilon_w\) is a unit since
\(\jmath\alpha\bar\alpha\,\widehat O_E=\varpi^{-1}\widehat O_E\)
by Lemma~\ref{lem:hom-oscillator-dual-lattices}. Thus the local
character in the trace pairing in these coordinates is
\(\zeta_w(t)=\psi_w^0(\epsilon_w t)\). Indeed
\(\psi_w^0(u)=\psi_{f,w}(-u/\varpi)\) by
Proposition~\ref{prop:additive}.
In the split model of Theorem~\ref{thm:split-finite-schroedinger}
a first-branch displacement \(x\) translates the argument
by \(x\). A displacement \(y\) on the conjugate branch
multiplies the value at \(X\) by
\(\prod_w\zeta_w({\rm tr}(y^*X))\). On either pure branch the
quadratic correction from symmetric translations is one.
The normalized adjoint map on these displacements is
\(x\mapsto-x^*\), \(y\mapsto-y^*\).

Consider the integral transform
\[
 \mathfrak F f(Z)=\int \overline{f(X)}
                      \prod_{w\in M}\zeta_w({\rm tr}(Z_wX_w))\,\dd X .
\]
Translating \(X\) by \(x\) in the input multiplies the result
by the character at \(-{\rm tr}(Zx)\); modulation by the
character at \({\rm tr}(y^*X)\) translates the output by \(-y^*\).
The transform takes the integral indicator to the integral indicator,
by the entrywise annihilator in
Proposition~\ref{prop:marked-lattice-unitary-fourier}.
The two intertwining properties and the value on the vacuum
determine its values on the span of the translates. This span
is the whole split test space by Theorem~\ref{thm:split-finite-schroedinger}.
It agrees there with \(C_\eta\); at all other places the vacuum is
taken to the vacuum by Proposition~\ref{prop:finite-adjoint-conjugates-splitting}.
These assertions on products follow by using the spanning vacuum
translates separately on each factor.

On the first branches the scalar \(s_\delta\) is \(\delta^{-1}\).
The normalized substitution for \(W'_{s_\delta}\) in
Theorem~\ref{thm:split-finite-schroedinger} sends \(f(Z)\) to
\(Q_R^{-ab/2} f(\delta Z)\); it acts as the identity away
from the pairs. Apply Proposition~\ref{prop:marked-lattice-unitary-fourier}
to \(\mathfrak F(\widetilde D)\). It equals that substitution
on the integral lift of \(\mathcal U_{a,b}D\), proving the claim.
\end{proof}

\section{Reversal of the global theta kernel}

\begin{definition}\label{def:two-order-theta-kernel}
Let \(S,T\) and their frames satisfy
Proposition~\ref{prop:finite-adjoint-conjugates-splitting}.
Write \(G=U(S)(\mathbb A_F),\ H=U(T)(\mathbb A_F)\) with
\(F=E^+\), and write \([G],[H]\) for their rational left
quotients. Choose unitary characters \(a_0,b_0\) of the
adelic norm-one torus, trivial on its rational points.
On \(V={\rm Hom}_E(S,T)\) set
\[
      \Omega(g,h)=\frac{a_0(\det h)}{b_0(\det g)}
                         \Omega_V(G_{g,h}).
\]
On \(V'\) use \(\Omega'(h,g)=b_0(\det g)
\Omega_{V'}(G'_{h,g})/a_0(\det h)\).
Use the splittings of
Theorem~\ref{thm:actual-adelic-rational-theta-invariance}.
Let \(P\) denote a polynomial in the real Hermite variables at
all infinite places and \(D\) the residue matrix function of
Theorem~\ref{thm:finite-adjoint-residue-test}. Put
\[
        K_{P,D}(g,h)=\theta\big(\Omega(g,h)(H_P\otimes v_D)\big).
\]
Here the Hermite polynomial may couple variables at different
places, by expansion into products of monomials. Write
\(K'_{P',D'}(h,g)\) for the corresponding kernel in the
reversed order. The rational theta invariance and the
triviality of \(a_0,b_0\) on rational points make both
kernels functions on the indicated quotients.
\end{definition}

\begin{theorem}\label{thm:actual-adjoint-reversed-kernel}
With the finite ideles \(\delta,\alpha\) of
Theorem~\ref{thm:finite-adjoint-residue-test}, put
\[
                  u_\delta=\frac{\overline{\alpha\delta}}{\alpha\delta},
\]
and give it real component one. Regard it as a central scalar
in \(G\). For fixed spaces, realizations, characters and modulus
there is \(c_\delta\) of modulus one such that for every \(P,D,g,h\),
\[
   \overline{K_{P,D}(g,h)}
      =c_\delta\,K'_{P^\#,\mathcal U_{a,b}D}(h,g u_\delta),
                         \qquad a=\dim_E T,\quad b=\dim_E S .
\]
The substitution \(P^\#\) of Theorem~\ref{thm:real-reversal-of-polynomial-test}
is applied in all places, conjugating the coefficients once.
For the unnormalized transform \(\widehat{\bar D}\) the
corresponding multiplier has squared modulus \(Q_R^{-ab}\).
\end{theorem}
\begin{proof}
Transport both rational models to the local models used for their
splittings as in Theorem~\ref{thm:actual-adelic-rational-theta-invariance}.
On elementary test tensors combine \(C_\infty\) at the infinite
places with \(C_\eta\) at the finite places to give an antiunitary
\(C_0\). This is well defined on the complex tensor product,
since moving a scalar between tensor factors conjugates it on
either side. It is bijective on tests, and multiplication of
the local inner products proves antiunitarity. In the real
variables the product substitution also acts on arbitrary Schwartz
functions, by conjugation and one signed permutation in all real
variables. Theorem~\ref{thm:real-reversal-of-polynomial-test}
and Proposition~\ref{prop:finite-adjoint-conjugates-splitting}
give \(C_0\Omega(g,h)=\Omega'(h,g)C_0\) because the determinant
multipliers are conjugate.

Let \(s_\alpha=\bar\alpha/\alpha\), extended by one at infinity.
For symmetric translations the operator \(\Omega'(1,s_\alpha)C_0\)
implements the rational map \(\sigma\) at every place. Indeed at
finite places \(s_\alpha\eta=\sigma\), and the archimedean
implementation is the one in
Theorem~\ref{thm:real-reversal-of-polynomial-test}. Its quotient
with the antiunitary \(C_\sigma\) of
Proposition~\ref{prop:rational-adjoint-conjugates-theta} is a linear
unitary commuting with all symmetric translations. These differ
from ordinary translations only by scalars, so
Lemma~\ref{lem:global-heisenberg-scalar} gives
\(C_\sigma=c_0\Omega'(1,s_\alpha)C_0\) for a fixed \(|c_0|=1\).

By Theorem~\ref{thm:finite-adjoint-residue-test}, the finite
part of \(C_0\) on \(v_D\) equals \(W'_{s_\delta}\) on
\(v'_{\mathcal U_{a,b}D}\), where \(s_\delta=\bar\delta/\delta\).
The operator \(\Omega'(1,s_\delta)\) differs from this finite
operator by the unitary ideal phase of
Theorem~\ref{thm:all-finite-weil-splitting} and a unitary
determinant value. Both are independent of \(P,D,g,h\).
Apply the equality for \(C_\sigma\) to
\(\Omega(g,h)(H_P\otimes v_D)\). Commute the central scalars,
use the representation law, and evaluate the resulting rational
theta value at zero. Proposition~\ref{prop:rational-adjoint-conjugates-theta}
conjugates that evaluation. The two central elements multiply to
\(u_\delta\), and the two fixed phases give \(c_\delta\).
Finally the normalization of
Definition~\ref{def:marked-residue-transform} replaces the
unitary transform by \(Q_R^{-ab/2}\widehat{\bar D}\).
\end{proof}

\begin{corollary}\label{cor:actual-first-reversed-hilbert-adjoint}
For the bounded first lift \(A\) of
Theorem~\ref{thm:actual-first-arrow-global-norm} take
\(K=K_{P,D}\) of Definition~\ref{def:two-order-theta-kernel}.
For compactly supported \(z\in L^2([H])\) put
\[
            Bz(g)=\int_{[H]} K'_{P^\#,\mathcal U_{a,b}D}(h,g)
                                          z(h)\,\dd h .
\]
With \(R(u)v(g)=v(gu)\), its Hilbert adjoint satisfies
\[
                    A^*z=c_\delta R(u_\delta)Bz .
\]
In particular the right side is square integrable. It extends
from such \(z\) to the adjoint on all of \(L^2([H])\).
\end{corollary}
\begin{proof}
The rational theta sum converges locally uniformly in \(g,h\).
Indeed the group action is continuous on real Schwartz tests and
locally constant on the finite test, as in
Theorem~\ref{thm:global-actual-arrow-covariance}, and the
rational lattice tail on a bounded family of Schwartz tests
is summable by the proof of
Theorem~\ref{thm:rational-schroedinger-zero-section}.
Quotient compact sets have lifts in compact sets by the finite
open cover argument in Theorem~\ref{thm:actual-first-arrow-global-norm}.
Thus the reverse integral is locally bounded. For compactly
supported \(y\in L^2([G])\), both \(y,z\) are integrable by
Cauchy--Schwarz on their finite-volume supports. The kernel is
bounded on their product, so integration in either order is
allowed. Theorem~\ref{thm:actual-adjoint-reversed-kernel} gives
\(\langle z,Ay\rangle=\langle v,y\rangle\) with
\(v=c_\delta R(u_\delta)Bz\).

For a compact exhaustion \(K_j\) of \([G]\), test with
\(y={\bf1}_{K_j}v\). Boundedness of \(A\) gives
\(\|{\bf1}_{K_j}v\|_2^2\le\|z\|_2\|A\|\|{\bf1}_{K_j}v\|_2\).
Increasing the compact sets and applying
Lemma~\ref{lem:integral-tools} gives \(\|v\|_2\le\|A\|\|z\|_2\).
The same exhaustion argument as in
Theorem~\ref{thm:actual-first-arrow-global-norm} gives density
of compactly supported functions on both sides. Thus the
assignment on \(z\) extends boundedly and its pairing identity
extends to all \(y,z\), which characterizes the Hilbert adjoint.
\end{proof}

\section{Descent to the opposite source}

\begin{lemma}\label{lem:opposite-source-frame-and-density}
Let \(S\) have an actual stage frame as in
Theorem~\ref{thm:actual-row-lift-frames}. There is a finite
integral idele \(\nu\in\widehat O_E^\times\) with
\(\nu\bar\nu=-1\). On \(S^-\), obtained by negating \(J_S\),
take the same ordered real frame with all signs reversed and
multiply the old finite frame by \(\nu\). This finite frame
again has identity Gram matrix. The two rational groups,
their adelic quotients, and the Haar measures with the stage
normalizations coincide under the identity on rational matrix
coordinates. For a fixed real parameter sequence the polar
coefficient of Definition~\ref{def:polar-row-coefficient}
is also identical.
\end{lemma}
\begin{proof}
At a split pair choose local components \(1,-1\).
At every nonsplit finite place use the residue norm lift in the
proof of Theorem~\ref{thm:all-finite-weil-splitting}: there the
extension is unramified of degree two, the characteristic of
the residue field is odd, and the successive unit corrections
lift a residue norm of \(-1\) to an exact norm of \(-1\).
All these components are units, so their product is a finite
idele. The norm of the scalar cancels the sign of the new Gram
matrix. The defining equation \(g^*J_Sg=J_S\) is unchanged by
negation. The scalar change of finite frame commutes with all
matrices and preserves the integral levels. At infinity
\(g^*(-J_S)\dd g=-g^*J_S\dd g\) gives the same density in
Definition~\ref{def:unitary-density}. Thus the identical rational
and adelic groups have the same quotient measure. For the polar
coefficient, equality of two signs at specified indices is
unaffected by a simultaneous sign change. Consequently \(\Lambda\),
\(\mu\), the compact subgroup, its probability measure, and the
ordinary polar decomposition used in
Definition~\ref{def:polar-row-coefficient} are the same.
\end{proof}

\begin{theorem}\label{thm:opposite-source-conjugate-kernel}
Use \(a=\dim_E T,\ b=\dim_E S\) and the kernel of
Definition~\ref{def:two-order-theta-kernel}. Write \(\upsilon_T\)
for the character of Theorem~\ref{thm:actual-global-splitting-character}
restricted to the norm-one torus. On \(S^-\to T\) use the
characters
\[
        a_-=a_0^{-1}\upsilon_T^{\,b},
               \qquad b_-=b_0^{-1}\upsilon_T^{\,a}
\]
in the same numerator and denominator convention.
Define \(D^-(X)=\overline{D(\nu^{-1}X)}\) on the residue
matrices, reducing \(\nu\) in the first branches. For a unit
scalar \(c_-\) independent of \(P,D,g,h\),
\[
              \overline{K_{P,D}(g,h)}
                       =c_- K^-_{\bar P,D^-}(g,h).
\]
The variables \(g,h\) on the right use the identification of
Lemma~\ref{lem:opposite-source-frame-and-density}. In particular
a point mass at the standard inclusion gives a point mass at
\(\nu\) times that inclusion. On compactly supported target
functions the first Hilbert adjoint is integration against this
opposite-source kernel with the scalar \(c_-\).
\end{theorem}
\begin{proof}
In rational coordinates the identity from
\({\rm Hom}(S,T)\) to \({\rm Hom}(S^-,T)\) sends \(H\) to
\(-H\). It therefore negates both the bilinear form and its
alternating part in Definition~\ref{def:rational-polarized-oscillator}.
Use the same position and modulation subspaces on both sides,
negating the modulation basis. The proof of
Proposition~\ref{prop:rational-adjoint-conjugates-theta}
then gives an antiunitary map conjugating theta at zero and
intertwining symmetric translations through this rational identity.

At infinity each positive Hom coordinate becomes negative and
each negative coordinate becomes positive. The Darboux coordinates
of Theorem~\ref{thm:negative-block-real-splitting} are therefore
conjugated, with a relabeling if the signs were grouped.
Complex conjugation on real position functions, with that
relabeling, sends \(H_P\) to \(H_{\bar P}\), by
Definition~\ref{def:positive-hermite-tests}. It intertwines the
translations. If \(A_+,A_-\) are the original positive and
negative diagonal blocks of the Hom isometry \(G_0=G_{g,h}\),
the conjugated old vacuum coefficient divided by the new
one is \(\det A_+/\overline{\det A_-}=\det_\C G_0\), by
the determinant calculation in
Theorem~\ref{thm:negative-block-real-splitting}. Uniqueness with
a specified vacuum coefficient shows that this ratio is exactly
the multiplier on the new real operator.

At finite places the rational identity sends normalized old
matrix coordinates \(X\) to \(\nu X\) in the new source frame.
The conjugate pullback on the full lattice model takes the old
vacuum to the new vacuum and intertwines the Heisenberg
translations. This follows directly from negation of \(b\) in
rational coordinates, so in this case no quadratic gauge factor
is needed. By the positive vacuum coefficient argument of
Proposition~\ref{prop:finite-adjoint-conjugates-splitting}, the
bare operators \(W_{G_0}\) agree under this transport.
The finite ideal \(J_{G_0}\) is unchanged since multiplication
by \(\nu\) preserves the lattice. The conjugate ideal phase divided
by the new one is consequently \(\beta(\bar J_{G_0}/J_{G_0})\).
The valuation calculation in the proof of
Theorem~\ref{thm:global-two-arrow-adjoint} gives
\(\bar J_{G_0}/J_{G_0}=(\det_E G_0)_f\) as fractional
ideals, where the right side is the ideal of the finite idele.
Together with the real determinant just found this gives the
factor \(\upsilon_T(\det_E G_0)\). In the split matrix realization
the same conjugate pullback is \(f(X)\mapsto\overline{f(\nu^{-1}X)}\).
Indeed this formula intertwines pure translations under
\(x\mapsto\nu x\), and the modulation character is conjugated
since \(\bar\nu\nu=-1\). It fixes the integral indicator, so
Theorem~\ref{thm:split-finite-schroedinger} determines it on the
test. It gives precisely \(D^-\).

The determinant of \(X\mapsto hXg^{-1}\) is
\((\det h)^b(\det g)^{-a}\), by multiplying the column and row
determinants. Thus multiplication of the conjugated old
character ratio by \(\upsilon_T(\det_E G_0)\) gives the stated
new ratio. The product of the local antiunitaries and the
rational antiunitary implement the same physical change. Their
quotient is a unit scalar by Lemma~\ref{lem:global-heisenberg-scalar}.
Evaluating at the origin proves the formula. Substitution gives
the assertion about the point mass. For the adjoint repeat the
compact-support integration and exhaustion argument of
Corollary~\ref{cor:actual-first-reversed-hilbert-adjoint}
with this expression for \(\bar K\).
\end{proof}

\section{The return operator on the prescribed source type}

\begin{lemma}\label{lem:compact-lowest-line-in-cyclic-copy}
Use the real product group \(G_\infty\), kernel \(e\), positive
constant \(d_e=e(1)\), and projection \(P_e^R\) of
Lemma~\ref{lem:real-coefficient-compression} in a strongly
continuous unitary representation \(R\). Write \(K=\prod_\tau K_\tau\),
\(m(k)=\prod_\tau m_{\Lambda_\tau}(k_\tau)\) and
\(s=\prod_\tau\dim T_{\Lambda_\tau}\), with the conventions of
Definition~\ref{def:cyclic-row-module}. Let \(\Omega\) be any
strongly continuous unitary representation of \(G_\infty\).
Suppose \(f\) satisfies
\[
 f=s\int_K m(k)\Omega(k)f\,\dd k,\qquad
          \phi(g)=\langle f,\Omega(g)f\rangle\in L^1(G_\infty).
\]
On every \(y=P_e^R y\),
\[
        \int_{G_\infty}\phi(g)R(g)y\,\dd g
          =\frac{\|P_{\bar e}^{\Omega}f\|^2}{d_e}\,y .
\]
\end{lemma}
\begin{proof}
It suffices to consider \(\|y\|=1\). Its matrix coefficient
is \(\bar e(g)/d_e=\prod_\tau\psi_\tau(g_\tau)\) by
Lemma~\ref{lem:real-coefficient-compression}. In the product
of the cyclic row representations of
Lemma~\ref{lem:cyclic-row-isometry}, the tensor \(v\) of their
unit generators has exactly this coefficient. The map
\(R(g)y\mapsto\pi(g)v\) therefore preserves all Gram entries
by taking the coefficient at \(g^{-1}g'\). Extending first
linearly and then by norm limits identifies the two closed cyclic
spans isometrically, intertwining translations. The latter
span is the completed tensor product of the single-place cyclic
spans, since products of independent translates give all their
elementary spanning tensors.

On either cyclic span put \(P_K=s\int_K\overline{m(k)}R(k)\dd k\),
using \(\pi\) in the tensor realization. It has rank one by
Theorem~\ref{thm:row-compact-finite-rank}: on elementary tensors
it is the tensor product of the one-place projections, and the
formula extends by boundedness. Its image contains \(v\).
Indeed the compact coefficient of each generator is \(m_{\Lambda_\tau}\)
by Definition~\ref{def:polar-row-coefficient} at positive polar
part \(I\) and the compact convolution computation in
Theorem~\ref{thm:row-compact-finite-rank}. Its squared integral
on the compact group is the reciprocal compact dimension by
that computation. Hence the expectation of the orthogonal
projection at the unit vector \(v\) is one. Transporting back,
the image of \(P_K\) on the original cyclic span is \(\C y\).

Write \(v_f=\int\phi(g)R(g)y\dd g\); absolute norm convergence
places it in the cyclic span. After applying \(P_K\), the coefficient
of \(R(g)y\) in the integral becomes
\[
       s\int_K\overline{m(k)}
              \langle\Omega(k)f,\Omega(g)f\rangle\,\dd k
                       =\phi(g).
\]
This follows by left translation in \(g\) and the assumed compact
identity for \(f\). Absolute integrability on the product follows
from \(\phi\in L^1\), compactness and translation invariance.
Thus \(v_f\in\C y\). Its coefficient on \(y\) is
\[
 \langle y,v_f\rangle
     =d_e^{-1}\int_{G_\infty}\bar e(g)\langle f,\Omega(g)f\rangle\,\dd g
     =d_e^{-1}\|P_{\bar e}^{\Omega}f\|^2
\]
by the orthogonal projection in
Lemma~\ref{lem:real-coefficient-compression}. Scaling proves
the formula for arbitrary \(y\).
\end{proof}

\begin{theorem}\label{thm:actual-first-descent-return}
At the arithmetic first arrows of
Theorem~\ref{thm:independent-actual-first-images}, use the
Hermite test of Theorem~\ref{thm:arithmetic-first-arrow-minor-nonzero}
and the marked finite test of Proposition~\ref{prop:deep-first-arrow-norm}.
Write \(A\) for the first lift and \(d_0\) for the bounded opposite-source
descent of Theorem~\ref{thm:opposite-source-conjugate-kernel}, including
\(c_-\) in its test. Let \(j\) be identification of opposite-source
functions by Lemma~\ref{lem:opposite-source-frame-and-density}, from
original to opposite. On every vector \(y\) of the finite principal
source level and real type \(e\),
\[
       d_0 A y=j\,C T_H y,\qquad
 C=\frac{{\rm vol}[U(T)]\,N_F\,\|Q_e f_\infty\|^2}{I_l e(1)}>0 .
\]
This is an identity of vectors in the full opposite-source
\(L^2\) space. The return preserves the initial finite and real
conditions under \(j\). It preserves any additional integral
invariance or finite group-algebra projection supported at
the marked pairs that is imposed on this source subspace.
It is injective there for all sufficiently large scales.
\end{theorem}
\begin{proof}
The integral for \(d_0\) on compactly supported target vectors
is \(j A^*\) there by Theorem~\ref{thm:opposite-source-conjugate-kernel}.
Its unique bounded extension is therefore \(j A^*\).
For arbitrary \(x\in L^2([G])\) and any \(y\) at the indicated
finite level, Theorem~\ref{thm:actual-first-arrow-global-norm}
gives
\[
       \langle Ax,Ay\rangle
          ={\rm vol}[U(T)]\,(N_F/I_l)
                         \langle x,T_H L_\infty y\rangle .
\]
At the marked prime-power depths the same formula follows before
real compression in Proposition~\ref{prop:deep-first-arrow-norm}:
the finite coefficient is zero off the principal level at
the marked places, constant there with mass \(N_F/I_l\), and
the remaining factor is \(\rho_H\). Only \(y\) must be fixed
under the compact groups to replace each right-coset integral
by its translation on \(y\). At the real places
\(\phi(g)=\langle f_\infty,\Omega(g,1)f_\infty\rangle\) is
integrable by Theorem~\ref{thm:archimedean-first-arrow-integrability}.
The compact fixed-vector identity for \(f_\infty\) holds factor
by factor by Theorems~\ref{thm:one-sided-harmonic-minor-test}
and~\ref{thm:arithmetic-first-arrow-minor-nonzero}.
Lemma~\ref{lem:compact-lowest-line-in-cyclic-copy} evaluates
\(L_\infty y\) as an actual multiple of \(y\), with the displayed
positive numerator by the latter arithmetic theorem. Testing
against every \(x\) identifies \(A^*Ay=CT_H y\).

The finite sum defining \(T_H\) in
Theorem~\ref{thm:positive-finite-first-arrow-convolution}
converges in operator norm on the fixed level by its summable
coset bounds. It preserves the complementary unmarked level
by left invariance of its coefficient. Each translate in the
sum commutes with the marked group actions and with all real
translations because it is supported away from those factors.
Thus the finite and real conditions are preserved, including
the additional marked projections. The same theorem and
\(H/2-2(r+z)=r-2z\) give a strictly positive lower bound on
the real quadratic form of \(T_H\) at sufficiently large
scales, as in Theorem~\ref{thm:independent-actual-first-images}.
This proves injectivity.
\end{proof}

\section{Convergence of the two complementary periods}
\begin{definition}\label{def:complementary-theta-periods}
Use the field, isometric frames and finite oscillator of
Definitions~\ref{def:split-inclusion-finite-oscillator}
and~\ref{def:marked-residue-transform}. Let \(U_1,U_2,A,A'\)
be positive-rank hermitian spaces with those frames, including
opposites as in Lemma~\ref{lem:opposite-source-frame-and-density}.
Put \(W=U_1\perp U_2\), with the sum frames, and assume
\(\dim W=\dim A+\dim A'\). Retain a marked branch \(w\ne\bar w\) above the rational prime
\(p\). Write \(k_w\) for its residue field and denote reduction
there modulo its maximal ideal by a subscript \(0\).
For any of these spaces \(M\), put
\(h_M(s)=\max(1,\lambda_M(s)^{-1})\), where the length and minimum
are those of Theorem~\ref{thm:unitary-linear-shortest-representative},
and use \([M]\) for the rational quotient of its adelic unitary group.
For a residue matrix \(B\) write \(v_B\) for the oscillator vector
named by the point mass at \(B\), with the normalizing index of the
finite Fourier sum in Proposition~\ref{prop:finite-oscillator-residue-name}.
For Schwartz real data \(f,f'\), use the kernels of
Definition~\ref{def:two-order-theta-kernel} on the pairs \((W,A)\)
and \((A',W)\), respectively, permitting independent unitary
determinant characters in the two kernels. Write
\[
 \begin{split}
 I_B(g)&=\int_{[A]}K^{W,A}_{f,B}(g,h)\dd h,\\
 I'_C(g)&=\int_{[A']}K^{A',W}_{f',C}(h,g)\dd h .
 \end{split}
\]
Here \(B:W_R\to A_R\) and \(C:A'_R\to W_R\) are matrices in the
finite reference frames over the residue ring \(R\); a framed space
of dimension \(j\) has coordinate module \(R^j\). The kernels
use \(v_B,v_C\) and the indicated Schwartz vectors in place of
the Hermite vectors there.
\end{definition}

\begin{lemma}\label{lem:two-sided-rank-detection}
In Definition~\ref{def:complementary-theta-periods} consider one
pair \((S,T)\), group arguments \((s,t)\), and a named point \(D\).
When \(D_0\) is injective there is a rational rank-one
\(z:S\to T\) such that \(x=t^{-1}zs\) satisfies
\[
   D_{x_f}v_D=\zeta v_D,\qquad
   \|x_\infty\|\le C\lambda_S(s)H(t).
\]
When \(D_0\) is surjective the same assertion holds with bound
\(C H(s)\lambda_T(t)\). In each case \(\zeta\) is a nontrivial
\(p\)-th root of unity. The constant is uniform in \(s,t\).
\end{lemma}
\begin{proof}
Choose the finite scalar convention and the lattice
\(L_0=\alpha\operatorname{Hom}(\Lambda_S,\Lambda_T)\)
of Definition~\ref{def:split-inclusion-finite-oscillator}.
Choose \(\beta\in E^\times\) in \(\delta\alpha\widehat O_E\)
at every place except \(\bar w\), with valuation
\(v_{\bar w}(\alpha)-1\) there.
The difference of the two fractional ideals obtained by increasing
this single valuation bound by one is nonempty by
Proposition~\ref{prop:general-ideal-factorization}.
Here \(\delta\) has components \(\pi_v^{l_v}\) at the marked
branches and one elsewhere.

For vectors \(u\in L(s)\), \(v\in L(t)\) take
\(z(y)=\beta v\langle u,y\rangle_S\) with the hermitian form
linear in \(y\). In finite coordinates let \(a\) be the reduction
of \(s_w^{-1}u\) and \(b^*\) the row given by the conjugate of
the reduction of \(t_{\bar w}^{-1}v\). The frames are isometries
with identity finite Gram matrices, so both residues can be
specified independently. Apart from its \(\bar w\)-component,
\(x_f\) lies in \(\delta L_0\). Translation by that part fixes
\(v_D\) by the Fourier sum in
Proposition~\ref{prop:finite-oscillator-residue-name}: it fixes
the vacuum and commutes with the translations indexed by
\(\bar\delta^{-1}L_-/L_-\), since their pairings are integral.
This argument works at depth \(l_v\) by the same trace duality
as in Proposition~\ref{prop:marked-lattice-unitary-fourier}.
The product-rule phase between the two parts of \(x_f\) is one,
as the \(w\)-part has valuation at least \(l_w\ge1\) relative
to \(L_0\). The \(\bar w\)-part acts on the Fourier sum by
reindexing. Its character at \(D\) is the additive character
of the residue field evaluated on a fixed nonzero scalar times
\(b^*D_0a\). Indeed the outer product at \(\bar w\) pairs by
matrix trace with \(D\) at \(w\), and has at most one denominator
relative to \(L_0\). The character on this first residue layer
is nontrivial by Lemma~\ref{lem:hom-oscillator-dual-lattices}.
The additive group of \(k_w\) has exponent \(p\).
Choose \(t_0\in k_w\) at which this additive character, with
the fixed nonzero scalar in its argument, is nontrivial.
Its value there is a nontrivial \(p\)-th root.

In the injective case take \(u\) short and primitive at \(w\)
by Lemma~\ref{lem:primitive-short-and-target-residue}. Then
\(D_0a\ne0\). Choose \(b^*\) taking value \(t_0\) on that
vector. This is possible by assigning its value in one nonzero
coordinate. Lift this residue to \(v\) of length at most
\(C H(t)\) by the other part of the same lemma.
In the surjective case take \(v\) short and primitive at
\(\bar w\). Then \(b^*D_0\ne0\). This row maps onto \(k_w\),
so choose \(a\) with \(b^*D_0a=t_0\) and lift it to \(u\) of
length at most \(C H(s)\).
The proof of the short-vector lemma applies to each space here:
it uses integral lattices and the fixed real frames, both preserved
under orthogonal sum and change of sign. The two vectors are
nonzero, giving rank one, and the norm of their framed real outer
product is bounded by a constant times the product of their
lengths. This proves both estimates.
\end{proof}

\begin{theorem}\label{thm:complementary-period-convergence}
If \(B_0\) is onto and \(C_0\) is one-to-one, both periods in
Definition~\ref{def:complementary-theta-periods} converge absolutely
and are continuous in \(g\). For each integer \(N\ge0\) there are
\(C_1>0\) and integers \(b,k\ge0\) such that
\[
 \begin{split}
 |K^{W,A}_{f,B}(g,h)|&\le C_1p_k(f)\,h_A(h)^{-N}h_W(g)^b,\\
 |K^{A',W}_{f',C}(h',g)|&\le C_1p_k(f')\,h_{A'}(h')^{-N}h_W(g)^b .
 \end{split}
\]
Here \(p_k\) is the real Schwartz seminorm of
Lemma~\ref{lem:real-small-heisenberg-differences}, and constants
can be chosen for all named points satisfying the stated ranks
at fixed spaces and level. If \(y_i\) are rapidly decreasing on
\([U_i]\) as in Proposition~\ref{prop:unitary-reduced-height-integrability},
then both kernels at \(g=g_1\oplus g_2\), multiplied by
\(y_1(g_1)y_2(g_2)\), are absolutely integrable on the respective
triple quotients. Finite matched sums and all three integrations
can be interchanged.
\end{theorem}
\begin{proof}
Use rational representatives with \(H\ll h_M\) by
Theorem~\ref{thm:unitary-linear-shortest-representative}.
Its proof for an integral hermitian lattice applies also to
the orthogonal sums here. Write \(u=h_M(h)\) for the height
of the variable of integration, where \(M=A\) or \(A'\), and
\(v=h_W(g)\). Lemma~\ref{lem:real-small-heisenberg-differences}
gives a coarse bound \(C p_j(f)u^a v^a\) for either kernel:
its proof uses only the finite datum, the height of the Hom action
and the real Schwartz seminorms. Use \(f'\) on the second pair.
For \(u>1\), Lemma~\ref{lem:two-sided-rank-detection} gives a
translation of real norm at most \(C_2v/u\); use its surjective
case on the first pair and its injective case on the second.

Write \(T(f)\) for the value of the fixed kernel. The rational
translation preserves theta evaluation, and its finite part
acts by \(\zeta\). Thus
\(T((D^\infty_{x_\infty}-1)^j f)=(\zeta^{-1}-1)^j T(f)\)
for every integer \(j\ge0\), by the intertwining in
Theorem~\ref{thm:actual-adelic-rational-theta-invariance}.
If \(u>\max(1,C_2)v\), apply the coarse bound to the difference
of order \(N+a\). The real difference estimate of
Lemma~\ref{lem:real-small-heisenberg-differences} and
\(|1-\zeta|\ge2\sin(\pi/p)\) give a bound by
\(C_Np_k(f)u^{-N}v^{N+2a}\). For the other \(u\), the coarse
estimate itself gives this bound with an increased constant.
There are only finitely many named points at the fixed level,
so maxima of the constants and indices suffice.

Proposition~\ref{prop:unitary-reduced-height-integrability}
bounds the measure of \(h_M\le R\) by \(C(1+R)^d\), using
comparability with the reduced height. With \(N>d+1\) the integral
of \(h_M^{-N}\) converges by summing over
\(2^j\le h_M<2^{j+1}\). This proves absolute convergence.
The endpoint height is bounded on compact sets since
\(h_W\ll H\); the kernels are continuous by
Theorem~\ref{thm:actual-adelic-rational-theta-invariance}
and its Schwartz convergence estimate. The same majorant
therefore gives continuity by Lemma~\ref{lem:integral-tools}.
For the sum frames the finite lattice is a direct sum and the
minimum on the sum is the smaller of its two minima. In fact
the norm of a vector dominates the norm of each component and
vectors with one component zero realize the two minima.
Thus \(h_W(g_1\oplus g_2)=\max_i h_{U_i}(g_i)\).
Bound its \(b\)-th power by \(\prod_i h_{U_i}(g_i)^b\).
The absolute integral is now bounded by a constant times the
integral just estimated and the product of the two finite integrals
\(\int_{[U_i]}|y_i|h_{U_i}^b\) from
Proposition~\ref{prop:unitary-reduced-height-integrability}.
Each matched sum is finite and every matrix in it has the
required rank by Definition~\ref{def:complementary-finite-frames}.
Lemmas~\ref{lem:integral-tools} and~\ref{lem:product-integration}
now permit every interchange.
\end{proof}

\section{A fixed pair of complementary frames}
\begin{definition}\label{def:complementary-period-characters}
In Definition~\ref{def:complementary-theta-periods} write
\(a=\dim A'\), \(d=\dim A\) and \(N=a+d\).
Choose a character \(\nu\) of the adelic norm-one torus, trivial
on rational points, whose restriction to its finite integral
units is given by a residue character \(\chi:R^\times\to\C^\times\).
Explicitly, on a norm-one finite unit its value is \(\chi\) of
the tuple of the reductions at the chosen branches.
Put \(\gamma=\nu^{-1}\upsilon_T^{\lfloor N/2\rfloor}\), using
Definition~\ref{def:actual-twisted-arrow-character}.
In \(I_B\) take the two characters \((\gamma,1)\), and in \(I'_C\)
take \((\gamma,\gamma)\), in the order of
Definition~\ref{def:two-order-theta-kernel}. Identify \(W_R\)
with \(J=R^a\oplus R^d\), with the chosen ordering of the
coordinates, and use \(x_0,y_0,X_k,Y_k,G_i\) from
Definition~\ref{def:complementary-finite-frames}.
For residue functions supported on the indicated frames set
\(I_\Phi=\sum_B\Phi(B)I_B\) and \(I'_\Psi=\sum_C\Psi(C)I'_C\).
\end{definition}

\begin{lemma}\label{lem:finite-covariance-of-periods}
For \(k\in G_N\), let \(\widetilde k\) be any integral finite
unitary lift, supported over the marked split pairs with real
component one. Such lifts exist. For all \(g\) and for all
quotient and subspace frames \(X,Y\) one has
\[
 \begin{aligned}
 I_{Xk^{-1}}(g)&=I_X(g\widetilde k),&
 I'_{kY}(g)&=\chi(\det k) I'_Y(g\widetilde k),\\
 I_{hX}(g)&=\chi(\det h) I_X(g),&
 I'_{Yj}(g)&=\chi(\det j) I'_Y(g)
 \end{aligned}
\]
where \(h\in G_d,\ j\in G_a\).
\end{lemma}
\begin{proof}
At each chosen branch lift the entries of \(k\) to the local
integers. Its determinant is a unit since its image modulo the
maximal ideal is a unit. At the conjugate branch use the inverse
conjugate transpose in the finite isometric frame. This preserves
the hermitian form, and using the identity elsewhere gives
\(\widetilde k\). The same construction applies on \(A,A'\).
The splitting ideal phase in
Theorem~\ref{thm:all-finite-weil-splitting} is one for these
integral isometries, since they preserve the reference lattice.
Likewise \(\upsilon_T\) is trivial on integral units by
Theorem~\ref{thm:actual-global-splitting-character}, so
\(\gamma(\det\widetilde k)=\chi(\det k)^{-1}\).

Acting on a point mass in the finite oscillator by the change of
variables for source \(s\) and target \(t\) sends its name \(D\)
to \(tDs^{-1}\), before multiplication by the determinant
character. This follows by reindexing the finite Fourier sum as in
Corollary~\ref{cor:finite-inclusion-periods}; at a higher depth
the same reindexing is on the two dual boxes of
Proposition~\ref{prop:marked-lattice-unitary-fourier}.
Consequently right multiplication of the ambient argument by
\(\widetilde k\) produces \(Xk^{-1}\) in the large kernel, and
\(\chi(\det k)^{-1}\) times the point mass \(kY\) in the small
kernel. This proves the first two equations using the integral
in Theorem~\ref{thm:complementary-period-convergence}.

In the large integral right translate its integration variable
by \(\widetilde h\). In the small integral right translate that
variable by \(\widetilde j^{-1}\). These compact translations
preserve the rational quotient measure of
Lemma~\ref{lem:quotient-measure}: their modular function is one,
since a compact subgroup of the positive real multiplicative
group is trivial. The resulting factors are
\(\chi(\det h)^{-1}\) on the name \(hX\) and
\(\chi(\det j)^{-1}\) on the name \(Yj\), respectively.
This gives the second line.
\end{proof}

\begin{theorem}\label{thm:single-frame-global-period-criterion}
For any \(c\in\C\) in
Definition~\ref{def:complementary-period-characters}, the identity
\(I_\Phi(g)=c I'_\Psi(g)\) for every matched pair and every \(g\)
is equivalent to the single identity of functions
\[
                    |G_d| I_{x_0}=c\,|G_a|I'_{y_0}.
\]
For the graph sums of Theorem~\ref{thm:finite-frame-descent-step}
on any subset \(B\) of ambient arguments it suffices that this
last identity hold on the right translates
\(B\widetilde k_y\), where \(k_y(v,b,w)=(v+yb,b,w)\)
in the coordinates of that theorem. For \(c\ne0\) the multiplier
in those finite frame sums in the orientation from \(I'\) to
\(I\) is \(c^{-1}\).
\end{theorem}
\begin{proof}
The basis of matched pairs in the proof of
Proposition~\ref{prop:complementary-frame-orbit-exchange}, applied
to \(k=1\), consists on its supporting orbits of the functions
\(hx_0\mapsto\chi(\det h)^{-1}\) and
\(y_0j\mapsto\chi(\det j)^{-1}\). The second line of
Lemma~\ref{lem:finite-covariance-of-periods} evaluates their
periods as \(|G_d|I_{x_0}\) and \(|G_a|I'_{y_0}\), which
proves necessity. For a general orbit representative \(k\)
the same basis calculation gives contributions
\[
       |G_d|\chi(\det k) I_{X_k}(g),\qquad |G_a| I'_{Y_k}(g).
\]
By the first line of that lemma these are
\(\chi(\det k)\) times the respective expressions for \(k=1\)
evaluated at \(g\widetilde k\). Thus the single identity implies
the one for every orbit, and addition proves sufficiency without
division by \(c\). In the graph calculation only the displayed
shears occur, each with determinant one. Using these translates
and the proof of Theorem~\ref{thm:finite-frame-descent-step}
gives the restricted assertion and its reciprocal when \(c\ne0\).
\end{proof}

\section{Inversion of the standard determinant tests}
\begin{definition}\label{def:integral-level-determinant-sections}
Use \(J=R^a\oplus R^d\), \(N=a+d\), \(Q=|R|\) and the ordered
coordinate identification in
Definition~\ref{def:complementary-period-characters}. Write
\(\mathcal G=\{k y_0(R^a):k\in G_N\}\). For an \(a\)-row matrix
\(S:J\to R^a\), define the pair
\[
 D_S(Y)=\chi^{-1}(\det(SY)),\qquad
 E_S(X)=\chi^{-1}\left(\det\binom{S}{X}\right).
\]
Determinants are relative to the ordered decomposition, and the
character has value zero on nonunits. Denote projection onto the
first summand by \(s_0\).
\end{definition}

\begin{lemma}\label{lem:standard-determinant-frame-pair}
Every \((E_S,D_S)\) in
Definition~\ref{def:integral-level-determinant-sections} is
matched. The pair vanishes unless \(S=s_0k^{-1}\) for some \(k\in G_N\).
For such \(k\),
\[
 \begin{aligned}
 E_{s_0k^{-1}}(X)&=\chi(\det k) E_{s_0}(Xk),&
 D_{s_0k^{-1}}(Y)&=D_{s_0}(k^{-1}Y).
 \end{aligned}
\]
\end{lemma}
\begin{proof}
If \(SY\) is invertible then the reduction of \(Y\) at every
residue field has independent columns. Extend these to a basis
at each field by successively adjoining coordinate vectors
outside the previous span. Lift the new columns to the factors
\(O_w/\pi_w^{l_w}\) of \(R\) and combine the entries into
columns over \(R\). The completed determinant is a unit in every
factor, so \(Y=Y_k\). Applying this argument to the transpose
of \(S\) shows \(S=s_0k^{-1}\) whenever \(SY\) is invertible.
If \(\binom{S}{X}\) is invertible, its inverse itself completes
\(X\) as a quotient frame and its first rows give \(S\) as a
quotient frame. For \(Y_k,X_k\),
\(\binom{S}{X_k}k\) has diagonal blocks \(SY_k,I_d\) and zero
lower-left block. Its determinant is \(\det(SY_k)\). This proves
matching even when the common determinant factor is a nonunit.
Finally \(\binom{s_0k^{-1}}{X}=\binom{s_0}{Xk}k^{-1}\);
this proves both transformation formulas.
\end{proof}

\begin{theorem}\label{thm:determinant-section-reconstruction}
Suppose the characters \(\chi,\psi\) on the multiplicative units
and additive group of \(R\), respectively, are primitive in
Definition~\ref{def:residue-matrix-fourier}. Choose one frame
\(Y_V\) for each \(V\in\mathcal G\). For any matched pair
\((\Phi,\Psi)\), put
\[
 b_S=\sum_{V\in\mathcal G}
           \Psi(Y_V)\overline{D_S(Y_V)} .
\]
Then \(b_S\) is independent of the choice of these frames, and
\[
       (\Phi,\Psi)=\frac{1}{Q^{ad}|G_a|}
              \sum_{S:J\to R^a} b_S(E_S,D_S).
\]
\end{theorem}
\begin{proof}
For \(Y=ky_0\) consider \(F_Y(S)=D_S(Y)\) as a function of \(S\).
For its Fourier transform use the trace pairing with
\(Z:R^a\to J\). Write \(Sk=(S_1,S_2)\). In the exponent
\({\rm tr}(SZ)={\rm tr}((Sk)(k^{-1}Z))\) the second component
of \(k^{-1}Z\) pairs with the free matrix \(S_2\), while
\(F_Y(S)=\chi^{-1}(\det S_1)\). The character sum in \(S_2\)
vanishes unless this second component is zero, in which case
it equals \(Q^{ad}\), by the entrywise orthogonality in
Lemma~\ref{lem:finite-determinant-fourier-stabilizer}. When
\(Z=YB\) the remaining factor is
\(G_-\,\chi(\det B)\), where \(G_-\) is the transform of
\(\chi^{-1}\circ\det\) at \(I_a\) in that lemma. Thus
\(\widehat F_Y\) is supported precisely on \(YG_a\), since
\(|G_-|^2=Q^{a^2}\). Transforms attached to different submodules
have disjoint supports. The same character orthogonality,
expanding the product of one transform and the conjugate of
the other, gives orthogonality of the original functions.
At a single frame exactly \(|G_a|Q^{ad}\) matrices \(S\)
satisfy \(SY\in G_a\), by the coordinates \(S_1,S_2\).
Consequently
\[
  \sum_S\overline{D_S(Y_V)}D_S(Y_U)
        =Q^{ad}|G_a|\,{\bf1}_{U=V}.
\]

Every matched \(\Psi\) satisfies
\(\Psi(Y_Vj)=\chi(\det j)^{-1}\Psi(Y_V)\) for \(j\in G_a\),
as proved by the block calculation in
Proposition~\ref{prop:complementary-frame-orbit-exchange}.
The corresponding law for \(D_S\) follows from its determinant.
Their product in \(b_S\) is independent of \(j\), as claimed.
The orthogonality formula recovers \(\Psi\) on each representative,
and equivariance recovers it on all frames. Both sides vanish
off the frames by Lemma~\ref{lem:standard-determinant-frame-pair}
and the support condition. The first component of a matched
pair is determined by its second via
Definition~\ref{def:complementary-finite-frames}. This proves
the formula for both components.
\end{proof}

\begin{corollary}\label{cor:standard-single-determinant-period}
Under the hypotheses of Theorem~\ref{thm:determinant-section-reconstruction}
the identity of functions \(I_{E_{s_0}}=c I'_{D_{s_0}}\) is
equivalent to \(I_\Phi=c I'_\Psi\) for every matched pair at
this level. This equivalence applies to the squarefree levels and
characters of Proposition~\ref{prop:primitive-square-level-characters}.
\end{corollary}
\begin{proof}
In Lemma~\ref{lem:finite-covariance-of-periods} reindex the finite
sum in \(X\) by \(Xk^{-1}\) and the one in \(Y\) by \(kY\).
The formulas of Lemma~\ref{lem:standard-determinant-frame-pair}
then give
\[
 \begin{aligned}
 I_{E_{s_0k^{-1}}}(g)&=\chi(\det k) I_{E_{s_0}}(g\widetilde k),\\
 I'_{D_{s_0k^{-1}}}(g)&=\chi(\det k) I'_{D_{s_0}}(g\widetilde k).
 \end{aligned}
\]
All other standard pairs vanish. The one equality therefore gives
the identity on every standard pair; the identical coefficients
in the reconstruction formula give it on all matched pairs.
The converse follows by choosing the matched pair at \(s_0\).
On the stated levels the multiplicative primitivity is precisely
Proposition~\ref{prop:primitive-square-level-characters} and the
additive primitivity is Proposition~\ref{prop:marked-lattice-unitary-fourier}.
\end{proof}

\begin{proposition}\label{prop:inflated-determinant-section-obstruction}
If \(R=\mathbb Z/p^2\mathbb Z\), \(p\) prime, \(a=d=1\), and
\(\chi\) is inflated from \(\mathbb F_p^\times\), the standard
pairs need not span all matched pairs.
\end{proposition}
\begin{proof}
The columns \(Y=(1,0)^{\mathsf T}\) and \(Y'=(1,p)^{\mathsf T}\)
give distinct summands since their first entries are one. For
every \(S=(u,v)\) the numbers \(u\) and \(u+pv\) have the same
reduction modulo \(p\). They are either both nonunits, or units
with the same character value. Thus \(D_S(Y)=D_S(Y')\) for
every \(S\). The single-orbit matched pair constructed in
Proposition~\ref{prop:complementary-frame-orbit-exchange} for
the summand of \(Y\) takes respective values \(1,0\) on these
two frames, which is impossible for a combination of standard pairs.
\end{proof}

\section{Fourier coefficients indexed by hermitian forms}
\begin{definition}\label{def:gram-indexed-mixed-period}
Let the ambient space have a rational splitting
\(W=L\oplus W_0\oplus L'\), with \(L,L'\) of dimension \(b>0\),
both totally isotropic, mutually in duality, and orthogonal to
the nonsingular \(W_0\). Choose dual ordered bases of \(L,L'\).
Write \(J_0,J_T\) for the hermitian matrices of \(W_0\) and a
nonsingular space \(T\), and assume both have unimodular finite
lattices as in Lemma~\ref{lem:weil-unimodular-frame-independence}.
For \(X\in{\rm Hom}_E(W,T)\) use coordinates \((x,m,u)\) on those
three summands, so \(x,u\) are matrices with \(b\) columns.
In \({\rm Hom}(W_0,T)\) choose a rational symplectic polarization
with position space \(P_0\) by the construction in
Lemma~\ref{lem:rational-graded-polarization}. In the full oscillator
take \(x\) and \(P_0\) as positions, and \(u\) and the complementary
space of the residual polarization as momenta. Use the trace
pairing and the additive character \({\bf e}\) of
Definition~\ref{def:rational-polarized-oscillator}.
For a test \(f\) in the product of these rational coordinates put
\(f_{g,h}=\Omega(g,h)f\), with any unitary pair characters as in
Definition~\ref{def:two-order-theta-kernel}. For rational \(x\) set
\[
      t_x(g,h;f)=\sum_{\xi\in P_0(\Q)} f_{g,h}(x,\xi).
\]
Write \(\mathcal Z\) for the additive rational space of
skew-hermitian \(b\)-by-\(b\) matrices over \(E\) and
\([\mathcal Z]=\mathcal Z(\mathbb A_\Q)/\mathcal Z(\Q)\).
For \(z\in\mathcal Z(\mathbb A_\Q)\) set \(n(z)(l,w,l')
=(l+zl',w,l')\) in \(W\). For each hermitian rational
\(b\)-by-\(b\) matrix \(B\), set
\[
       \xi_B(z)={\bf e}\left(-\tfrac12
        {\rm Tr}_{E/\Q}\big(\jmath\,{\rm tr}(zB)\big)\right).
\]
For a continuous automorphic function \(I\) on \([W]\) define its
\(B\)-coefficient at \(g\) by
\[
        \int_{[\mathcal Z]}\overline{\xi_B(z)}I(n(z)g)\dd z
\]
with additive probability measure. The space \([\mathcal Z]\)
is compact since any rational basis identifies it with a finite
product of \(\mathbb A_\Q/\Q\), as in
Lemma~\ref{lem:theta-positive-translation}. Both factors in
the integrand are invariant under rational translations of \(z\).
\end{definition}

\begin{lemma}\label{lem:gram-shear-character-on-slices}
In Definition~\ref{def:gram-indexed-mixed-period},
\[
 K_f(n(z)g,h)=
      \sum_{x\in{\rm Hom}_E(L,T)}\xi_{x^*J_Tx}(z)t_x(g,h;f).
\]
The sum converges absolutely, uniformly in \(z\) for fixed \(g,h\).
Moreover \(t_{\delta x}(g,\delta h;f)=t_x(g,h;f)\) for
\(\delta\in U(T)(E)\).
\end{lemma}
\begin{proof}
In the dual bases the off-diagonal Gram block of \(W\) is \(I_b\).
Thus \(n(z)\) preserves the form since \(z+z^*=0\), and on the
Hom space its action is \((x,m,u)\mapsto(x,m,u-xz)\).
The factors on \({\rm Hom}(L\oplus L',T)\) and
\({\rm Hom}(W_0,T)\) split in the oscillator: choose orthogonal
real frames and a direct sum of the hyperbolic dual lattices and
the residual lattice, as in the proof of
Theorem~\ref{thm:residual-levi-zero-section}. There the vacuum
coefficient, the splitting character and the translation action
give precisely the product operators, and
Lemma~\ref{lem:weil-unimodular-frame-independence} allows the
original reference frame. The shear is the identity on the
residual factor. On the hyperbolic part the position block is
the identity and the lower block \(C\) sends \(x\) to \(-xz\).
By Theorems~\ref{thm:hyperbolic-adelic-parabolic-splitting}
and~\ref{thm:graded-parabolic-schroedinger} its action on a
position value is multiplication by
\({\bf e}(-a(x,Cx)/2)\). All determinant factors are one.
Indeed \(a(x,Cx)={\rm Tr}_{E/\Q}(\jmath H(Cx,x))\);
the trace in \(H\) is
\({\rm tr}((-xz)^*J_T x)={\rm tr}(z x^*J_Tx)\)
by Definition~\ref{def:hom-oscillator-pairing}. This gives the
specified sign and character. The rational theta sum in
Theorem~\ref{thm:rational-schroedinger-zero-section} now gives
the asserted expansion. Its sum of absolute values before
grouping by \(x\) is finite by
Proposition~\ref{prop:rational-frequency-lattice-bound}. The
phase has modulus one, giving uniformity in \(z\).

A rational target isometry \(\delta\) acts on the hyperbolic
part by \(x\mapsto\delta x\) and the dual action on \(u\).
The same parabolic formula gives the substitution
\(x\mapsto\delta^{-1}x\) on functions, with factor one: the
adelic determinant modulus of a rational invertible matrix is
one and \(\upsilon\) is trivial on rational determinants by
Theorem~\ref{thm:actual-global-splitting-character}. The residual
operator preserves its rational theta evaluation by
Theorem~\ref{thm:actual-adelic-rational-theta-invariance}.
The pair character is trivial on the rational point as well.
Evaluating at \(\delta x\) proves the covariance.
\end{proof}

\begin{theorem}\label{thm:truncated-gram-period-formula}
In Definition~\ref{def:gram-indexed-mixed-period} take \(f\) to be
a finite sum of real Schwartz tests times point data mapping onto
\(T\) at the chosen branch of
Definition~\ref{def:complementary-theta-periods}. Set
\(I(g)=\int_{[T]}K_f(g,h)\dd h\). For any rational hermitian
matrix \(B\) its coefficient is
\[
    \int_{[T]} \ \sum_{\substack{x\in{\rm Hom}_E(L,T)\\x^*J_Tx=B}}
                     t_x(g,h;f)\,\dd h .
\]
The grouped integrand is absolutely integrable.
Truncating the outer integral to any increasing compact exhaustion
of the quotient gives the same value in the limit. In particular
the coefficient vanishes if \(B\) has no representation over \(E\)
by \(T\). For nonsingular \(B\) with a representation \(x_*\), all
the indices of the sum lie in the single orbit \(U(T)(E)x_*\);
its rational stabilizer acts identically on \(x_*(L)\) and is the
full unitary group of its orthogonal complement there.
\end{theorem}
\begin{proof}
The proof of Theorem~\ref{thm:complementary-period-convergence}
for a surjective matrix uses only \(W,T\), its ranks at the chosen
prime and its two heights. Thus it applies to this pair. A compact
set of \(z\) represents \([\mathcal Z]\); its elements \(n(z)g\)
have bounded height for fixed \(g\) by
Theorem~\ref{thm:unitary-linear-shortest-representative}. The bound
for the full theta kernel by \(C_Nh_T(h)^{-N}\) on this set
is integrable for large \(N\). Lemma~\ref{lem:product-integration}
therefore interchanges the two integrations. For fixed \(h\),
Lemma~\ref{lem:gram-shear-character-on-slices} permits termwise
integration on the compact group. Characters with different
hermitian indices are distinct. Indeed for a nonzero difference
\(D\), at a complex place take \(z=r\jmath^{-1}D\) there and zero
at the other places, with \(r\in\R\) arbitrary. This component
is skew-hermitian, and \({\rm tr}(D^2)=\sum_{ij}|D_{ij}|^2>0\)
at that place. By varying \(r\) the value of the character differs
from one. Translation of additive probability by such a \(z\)
shows its integral is zero; the trivial character has integral one.
This proves the formula. Pointwise the absolute value of the
retained sum is bounded by the integral of the absolute value
of the full theta kernel on \([\mathcal Z]\), which has the same
integrable majorant. Covariance in
Lemma~\ref{lem:gram-shear-character-on-slices} makes the
grouped sum a function on \([T]\).
Lemma~\ref{lem:integral-tools} gives the asserted exhaustion and
the empty index set gives the vanishing.

For nonsingular \(B\), every representing matrix is injective
since its kernel is in the kernel of \(B\). The rule \(x_*v\mapsto xv\)
defines an isometry between nonsingular rational subspaces.
Lemma~\ref{lem:rational-hermitian-witt-extension} extends this
map to a rational isometry of \(T\). A unitary fixing all columns
of \(x_*\) preserves their perpendicular; conversely the direct
sum of the identity on their span and any unitary on that
perpendicular fixes the matrix. This proves the orbit and
stabilizer assertions.
\end{proof}

\section{Absolute unfolding of a nonsingular Gram coefficient}
\begin{definition}\label{def:residual-frame-rank-condition}
Use Definition~\ref{def:gram-indexed-mixed-period} with the marked
branch \(w\) of Definition~\ref{def:complementary-theta-periods}.
In the marked finite frame of \(W\) put
\[
 M_{g,w}=g_w^{-1}(L^\perp\otimes_E E_w)\cap\Lambda_{W,w}.
\]
Here \(L^\perp=L\oplus W_0\). Write \(\overline M_{g,w}\) for the
image of \(M_{g,w}\) in the residue vector space of the lattice.
For a named point \(D:W_R\to T_R\) the condition at \(g\) is
that \(D_0\) maps \(\overline M_{g,w}\) onto the full target
residue space. For a finite linear combination of point data
impose the condition separately on each point with nonzero
coefficient.
\end{definition}

\begin{theorem}\label{thm:absolute-mixed-frame-terms}
Fix \(g\) and a named point \(D\) satisfying
Definition~\ref{def:residual-frame-rank-condition}. Write
\(f=v\otimes v_D\), with \(v\) a real Schwartz function. For
each integer \(M\ge0\) there are \(C_{g,M}>0\) and an integer
\(j\ge0\), independent of \(h,v\), such that
\[
   \sum_{x\in{\rm Hom}_E(L,T)}|t_x(g,h;f)|
              \le C_{g,M}\,p_j(v)\,h_T(h)^{-M}.
\]
In particular the sum of the absolute values is integrable on
\([T]\). These conclusions apply by addition to finite sums of
the allowed elementary data.
\end{theorem}
\begin{proof}
Fix \(g\). Each residue \(a\in\overline M_{g,w}\) lifts to a
rational \(u\in L^\perp\) with \(g_f^{-1}u\) integral at all
places and residue \(a\) at \(w\). Indeed pull the compact open
lattice \(g_f\Lambda_W\cap L^\perp(\mathbb A_{E,f})\) back to a
rational basis of \(L^\perp\). It contains an integer multiple
of the standard integral lattice and lies in an inverse-integer
multiple. Increasing the first integer by a factor divisible
by \(p\) if needed gives a contained lattice mapping to zero in
the required residue. The quotient of these two bounding lattices
is finite and every class lifts to the rational space, by the
integer Chinese remainder construction in Lemma~\ref{lem:trace-finite}.
Lift a local representative in the specified intersection, setting
the other components to zero. Fix one rational lift for every
residue. Their \(g\)-lengths are bounded by a constant depending
on \(g\).

Choose \(v'\in L(h)\) short and primitive at \(\bar w\) by
Lemma~\ref{lem:primitive-short-and-target-residue}; here \(L(h)\)
is the target lattice of
Theorem~\ref{thm:unitary-linear-shortest-representative}. Its
conjugate residue row \(b^*\) is nonzero. By the surjectivity
hypothesis the row \(b^*D_0\) on \(\overline M_{g,w}\) maps
onto \(k_w\). Take an \(a\) on which its value is the
\(t_0\) selected in Lemma~\ref{lem:two-sided-rank-detection}.
Use its fixed lift \(u\) to form
\(z=\beta v'\langle u,\cdot\rangle_W\) with the scalar of
Lemma~\ref{lem:two-sided-rank-detection}. That lemma's finite
calculation gives \(D_{e_f}v_D=\zeta v_D\) for
\(e=h^{-1}zg\), with \(|1-\zeta|\ge2\sin(\pi/p)\);
the outer product bound gives
\(\|e_\infty\|\le C_g\lambda_T(h)\).
Crucially \(z|_L=0\) since \(u\in L^\perp\).

A rational translation by this \(z\) preserves \emph{each}
functional \(F\mapsto\sum_{\eta\in P_0(\Q)}F(x,\eta)\).
In the formula of Theorem~\ref{thm:rational-schroedinger-zero-section}
its position increment in \(x\) is zero, its residual position
increment is rational, and all its phase arguments at the
rational positions \((x,\eta)\) are rational. The additive
character there is one. Moving the translation through
\(\Omega(g,h)\) therefore gives, simultaneously for all \(x\),
\[
 t_x\bigl(g,h;(D^\infty_{e_\infty}-1)^j v\otimes v_D\bigr)
        =(\zeta^{-1}-1)^j\,t_x(g,h;v\otimes v_D)
                                    \qquad(j\ge0).
\]
This follows for \(j=1\) from the intertwining in
Theorem~\ref{thm:actual-adelic-rational-theta-invariance}, and
iteration is valid for every Schwartz vector.

The sum on the left of the proposed estimate is invariant under
rational left translation of \(h\), by
Lemma~\ref{lem:gram-shear-character-on-slices}. Take a representative
with \(H(h)\ll h_T(h)\). Before using the translation, the sum is
at most \(\sum_{x,\eta}|f_{g,h}(x,\eta)|\); this in turn is at
most \(C_g p_i(v)h_T(h)^k\) for fixed integers \(i,k\).
This is Proposition~\ref{prop:adelic-schwartz-frequency-height}
in the present rational polarization, with the seminorm dependence
computed in Lemma~\ref{lem:real-small-heisenberg-differences}.
When \(h_T(h)>\max(1,C_g)\), increase \(C_g\) if needed so
\(\|e_\infty\|\le1\). Apply this coarse sum bound to the
difference of order \(j=M+k\), and the preceding identity
and the small difference estimate of the same lemma to get
\(C_{g,M}p_{i+2j}(v)h_T(h)^{k-j}\). At smaller \(h_T(h)\) the
coarse estimate suffices with a larger constant.
Finally Proposition~\ref{prop:unitary-reduced-height-integrability}
and its dyadic sublevel estimate give integrability on the quotient.
\end{proof}

\begin{corollary}\label{cor:absolute-stabilizer-period}
Fix a nonsingular represented rational hermitian \(B=x_*^*J_Tx_*\),
and let \(\Gamma_*=\{\delta\in U(T)(E):\delta x_*=x_*\}\).
For the data of Theorem~\ref{thm:absolute-mixed-frame-terms},
the \(B\)-coefficient of \(I\) at \(g\) is the absolutely
convergent integral
\[
        \int_{\Gamma_*\backslash U(T)(\mathbb A_F)}
                    t_{x_*}(g,h;f)\,\dd h ,\qquad F=E^+ .
\]
The measure is induced by the same Haar measure used for \([T]\).
When \(T\) is the orthogonal sum of the framed space \(x_*L\)
and \(T_1\), the group \(\Gamma_*\) is \(U(T_1)(E)\) embedded
as the identity on \(x_*L\); when \(\dim T=b\) it is trivial.
\end{corollary}
\begin{proof}
The group and orbit are precisely those of
Theorem~\ref{thm:truncated-gram-period-formula}. Use the covariance
of Lemma~\ref{lem:gram-shear-character-on-slices} and apply
Lemma~\ref{lem:countable-frequency-block-unfold} to the singleton
\(x_*\). Its integrability requirement follows from
Theorem~\ref{thm:absolute-mixed-frame-terms} since the sum on the
orbit is bounded in absolute sum by the sum over all rational
matrices. The lemma integrates over the disjoint translates of a
domain for the rational group, so the Haar normalizations agree.
The description of the rational stabilizer gives both special cases.
\end{proof}

\begin{proposition}\label{prop:standard-test-satisfies-slice-rank}
In Definition~\ref{def:integral-level-determinant-sections} take
\(a=b,\ d=\dim W-b=\dim T\). There is a quotient frame \(S\)
for which the standard test \(E_S\) satisfies
Definition~\ref{def:residual-frame-rank-condition} at every
adelic \(g\) preserving \(L\). The criterion of
Corollary~\ref{cor:standard-single-determinant-period} may use
this \(S\) in place of \(s_0\).
\end{proposition}
\begin{proof}
At \(w\) the intersection \(L^\perp(E_w)\cap\Lambda_{W,w}\)
is a direct summand of rank \(d\) of the valuation-ring lattice.
To see this, clear denominators of a basis of \(L^\perp(E_w)\)
and diagonalize its column matrix by integral row and column
elimination, successively using an entry of least valuation
as pivot. In the resulting ambient basis the field subspace
is a coordinate subspace, so its full intersection is the
corresponding integral coordinate summand. Choose \(S_w\) to
project away from this summand, then reduce modulo \(\pi_w^{l_w}\).
At other marked branches use any coordinate projection of rank
\(b\). Their product is a quotient frame by the completion
argument of Lemma~\ref{lem:standard-determinant-frame-pair}.
An isometry preserving \(L\) preserves \(L^\perp\), so
\(\overline M_{g,w}=\ker(S_w\bmod\pi_w)\).
Whenever \(E_S(D)\ne0\), the matrix \(\binom{S}{D}\) is
invertible over the residue ring. Its reduction shows that
\(D_0\) maps the indicated kernel isomorphically to the
\(d\)-dimensional target residue space. This proves the rank
condition. Finally \(S=s_0k^{-1}\) for some \(k\); the last
two period transformation formulas in the proof of
Corollary~\ref{cor:standard-single-determinant-period} are
invertible and have the same character factor.
\end{proof}

\section{The small period in the same polarization}
\begin{lemma}\label{lem:adjoint-on-conjugate-names}
In Proposition~\ref{prop:finite-adjoint-conjugates-splitting}
use the chosen set \(M\) for naming the original datum and
\(\bar M\) for naming the reversed datum. At conjugate places
take conjugate uniformizer powers, so entrywise field conjugation
identifies the two residue rings. Write \(Z^*\) for conjugation
between these rings followed by transposition. The finite
antiunitary \(C_\eta\) takes \(v_D\) to the reversed test
named at \(\bar M\) by \(Z\mapsto\overline{D(-Z^*)}\).
\end{lemma}
\begin{proof}
In the split model at \(M\), a displacement supported on \(M\)
is a translation in the position variable \(X\), normalized by
\(\alpha\), and a displacement on \(\bar M\) is a modulation,
by Theorem~\ref{thm:split-finite-schroedinger}. In the reversed
model take the position variable \(Z\) on \(\bar M\). Under
\(\eta\) a position displacement \(x\) goes to \(-x^*\), now
again a position displacement, and the same formula sends the
modulation variable to a modulation variable. The pairing
character of the two transformed displacements is the conjugate
of the original one, since \(\eta\) negates the alternating
pairing by Proposition~\ref{prop:finite-adjoint-conjugates-splitting}.
On each pure branch the quadratic change to symmetric
translations is one. It follows that
\(\phi(X)\mapsto\overline{\phi(-Z^*)}\) implements both
intertwining relations. It fixes the integral indicator since
conjugation and transposition take one integral box onto the other.
The vacuum translates span the split test space by
Theorem~\ref{thm:split-finite-schroedinger}, so the formula
agrees with \(C_\eta\). Off the marked pairs the vacuum maps to
the vacuum. Reduction at the prescribed conjugate depths gives
the stated residue function.
\end{proof}

\begin{theorem}\label{thm:opposite-branch-small-period}
Use Definition~\ref{def:complementary-period-characters},
with \(a=\dim A'\), \(N=\dim W\).
On \(T=(A')^-\) use the finite frame \(\epsilon C_{A'}\) of
Lemma~\ref{lem:opposite-source-frame-and-density}, where
\(\epsilon\bar\epsilon=-1\). Given small-side residue data \(D\)
read at \(M\), read the following data at \(\bar M\) on the pair
\((W,T)\):
\[
    D^\flat(Z)=D(-\epsilon^{-1} Z^*),\qquad
    a^\flat=\gamma^{-1}\upsilon_T^{\,N},\qquad
    b^\flat=\gamma^{-1}\upsilon_T^{\,a}.
\]
Here the symbol \(\upsilon_T\) denotes the splitting character
on the norm-one torus from Definition~\ref{def:actual-twisted-arrow-character}.
For a Schwartz real test \(v\) on the small pair there is a
Schwartz test \(v^\flat\), obtained by a fixed linear coordinate
identification of real oscillator spaces. For some
\(|c_0|=1\) independent of the tests and group arguments,
\[
 K^{A',W}_{v,D}(h,g)
   =c_0 K^{W,T}_{v^\flat,D^\flat}(g,h s_\alpha),\qquad
           s_\alpha=\bar\alpha/\alpha
              \quad\hbox{at finite places},\quad
           s_{\alpha,\infty}=1 .
\]
The new kernel uses the characters \(a^\flat,b^\flat\).
If \(D\) is supported on subspace frames, its period is \(c_0\)
times the period of the new kernel with arguments \((g,h)\).
\end{theorem}
\begin{proof}
Conjugation to the opposite source in
Theorem~\ref{thm:opposite-source-conjugate-kernel} first changes
the name to \(D^-(X)=\overline{D(\epsilon^{-1}X)}\) and changes
the two characters to
\(\gamma^{-1}\upsilon_T^{a},\gamma^{-1}\upsilon_T^N\).
Now apply the rational adjoint conjugation of
Proposition~\ref{prop:rational-adjoint-conjugates-theta} on this
opposite pair. As in the proof of
Theorem~\ref{thm:actual-adjoint-reversed-kernel}, its real
implementation is \(C_\infty\), and at finite places the
normalized adjoint \(\eta\) differs from the rational adjoint
by \(\alpha/\bar\alpha\). Multiplication by \(s_\alpha\) in
the reversed target group corrects this difference. The quotient
of these implementations is a unit scalar by
Lemma~\ref{lem:global-heisenberg-scalar}. The adjoint interchanges
the two characters. Lemma~\ref{lem:adjoint-on-conjugate-names}
gives directly \(\overline{D^-(-Z^*)}=D(-\epsilon^{-1}Z^*)\).
Both uses of conjugation are antiunitary, and their product
is linear. At infinity they act on real position functions by
conjugation and fixed signed coordinate changes, as shown in
Theorems~\ref{thm:real-reversal-of-polynomial-test}
and~\ref{thm:opposite-source-conjugate-kernel}. These formulas
preserve Schwartz tests and give \(v^\flat\). The two scalar
comparison factors give \(c_0\).

The groups of \(A'\) and \(T\) and their quotient measures agree
by Lemma~\ref{lem:opposite-source-frame-and-density}.
The old period converges absolutely by
Theorem~\ref{thm:complementary-period-convergence}. On the new
side translation by \(s_\alpha\) preserves quotient Haar measure:
this element is central, so its right translation is also left
translation upstairs. Consequently the shift disappears in the
period integral.
\end{proof}

\begin{corollary}\label{cor:small-period-trivial-stabilizer}
In Proposition~\ref{prop:standard-test-satisfies-slice-rank}
choose the standard test at a frame \(S\) whose integral lift
at \(w\) has kernel \(L^\perp(E_w)\cap\Lambda_{W,w}\). For every
\(g\) preserving \(L\), every named point in the conjugate-branch
test for \(D_S\) satisfies
Definition~\ref{def:residual-frame-rank-condition} at \(\bar w\).
If a nonsingular \(B\) is represented by \(x_*\) in \(T=(A')^-\),
then the small-side \(B\)-coefficient equals
\[
           c_0\int_{U(T)(\mathbb A_F)}
                  t_{x_*}(g,h;v^\flat\otimes v_{D_S^\flat})\,\dd h .
\]
Here a finite combination of point vectors has the linear meaning
of Definition~\ref{def:complementary-period-characters}.
For a nonsingular \(B\) unrepresented by \(T\), this coefficient
is zero.
\end{corollary}
\begin{proof}
At the split pair the finite frames identify the lattice on
\(\bar w\) with the conjugate dual of the one on \(w\).
The columns of the conjugate transpose of the chosen integral
lift of \(S_w\) form a basis of the annihilator of its kernel:
a split surjection onto the free module identifies that annihilator
with the dual of the quotient, by extending a basis of the
kernel to a basis of the lattice. The field span of the
annihilator is \(L(E_{\bar w})\), by the perfect hermitian
pairing and the identity \((L^\perp)^\perp=L\).
It is the full intersection with the lattice, by the same
dual description. In particular its reduced columns lie in
\(\overline M_{g,\bar w}\) for every \(g\) preserving \(L\),
since \(L\subset L^\perp\).

For a point \(Z\) of the support the determinant of
\(S(-\epsilon^{-1}Z^*)\) is a unit by
Definition~\ref{def:integral-level-determinant-sections}.
Conjugation, transposition and multiplication by the unit scalar
show that \(Z S^*\) is invertible at the conjugate branch.
Thus \(Z_0\) maps the reduced columns just described onto the
target residue space. Corollary~\ref{cor:absolute-stabilizer-period}
applies to every term of this finite combination. The target
has dimension \(b\), so the stabilizer of a frame of nonsingular
Gram matrix is trivial. Using
Theorem~\ref{thm:opposite-branch-small-period} gives the displayed
formula, and Theorem~\ref{thm:truncated-gram-period-formula} gives
the vanishing in the absence of a frame.
\end{proof}

\section{Support of the individual frame terms}
\begin{lemma}\label{lem:frame-slice-support-transport}
Use the mixed polarization of
Definition~\ref{def:gram-indexed-mixed-period}, with \(gL=L\)
and \(g_L=g|_L\). At each place the map induced on the frame
coordinate by the Hom action is \(R(x)=hxg_L^{-1}\).
At a real place suppose the original Schwartz test vanishes
for frame coordinates outside a closed set \(C\). The moved
test vanishes for frame coordinates outside \(R(C)\).

At a chosen branch \(v\) of a split pair write
\(\mathfrak l=L(E_v)\cap\Lambda_{W,v}\). For a finite point
test \(v_D\) choose any integral lift \(\widetilde D\) of the
residue matrix there in the full finite frames. If
\(t_x(g,h;f)\ne0\), then in these frames
\[
  (h^{-1}xg_L)_v\big|_{\mathfrak l}\ \in\
  \alpha_v\widetilde D\big|_{\mathfrak l}
      +\alpha_v\pi_v^{l_v}
                    {\rm Hom}_{O_v}(\mathfrak l,\Lambda_{T,v}).
\]
A linear combination of point tests uses the union of their
cosets.
\end{lemma}
\begin{proof}
Restriction of Hom matrices to \(L\) intertwines the two linear
maps since \((hXg^{-1})|_L=h(X|_L)g_L^{-1}\).
For its dual variables use displacements in the momentum
summand \({\rm Hom}(L',T)\), extended by zero on \(L^\perp\).
The alternating form pairs them perfectly with the frame
position and pairs them trivially with the residual variables,
as seen from the off-diagonal Gram blocks in
Lemma~\ref{lem:gram-shear-character-on-slices}. Such displacements
remain in that summand under the Hom action since \(gL^\perp=L^\perp\).
On functions they multiply by the additive character of a
linear functional on the frame coordinate, by
Theorem~\ref{thm:rational-schroedinger-zero-section}. Their
quadratic gauge factor is one on the isotropic momentum summand.
The intertwining law in
Theorem~\ref{thm:actual-adelic-rational-theta-invariance}
therefore transports the multiplier of a functional \(\ell\)
to the multiplier of \(\ell\circ R^{-1}\).

At infinity let \(b\) be a smooth function of compact support
on the frame coordinates. Fourier inversion expresses its
multiplication operator as the integral of the character
multipliers against \(\widehat b\) in real dual coordinates.
Indeed integration by parts in each coordinate gives an
integrable Fourier transform, and the inversion formula follows
by the Gaussian insertion in Lemma~\ref{lem:schroedinger-real-gaussian},
one coordinate at a time. The identity can be integrated in
\(L^2\) because the multipliers are unitary and
\(\int|\widehat b|<\infty\). Thus conjugation carries
multiplication by \(b\) to multiplication by \(b\circ R^{-1}\).
Take \(b\) supported in the complement of \(C\) and nonzero
at the preimage of a proposed position outside \(R(C)\).
Such a bump on a small ball follows by scaling
\(\exp(-1/(1-|y|^2))\) for \(|y|<1\), extended by zero:
each derivative approaches zero at the boundary since an
exponential dominates any fixed power. The old product is
zero, so the moved product is zero in \(L^2\) and pointwise by
continuity. This proves the real assertion.

At the split place the branch opposite \(v\) contains the dual
of the whole position matrix space on \(v\), by
Theorem~\ref{thm:split-finite-schroedinger}. Restrict its
modulations to the dual of the frame coordinate just described.
The original support on the chosen branch is
\(\alpha_v(\widetilde D+\pi_v^{l_v}{\rm Hom}_{O_v}
(\Lambda_{W,v},\Lambda_{T,v}))\). Its restriction to
\(\mathfrak l\) is exactly the coset in the assertion:
\(\mathfrak l\) is an integral direct summand by the elimination
argument in Proposition~\ref{prop:standard-test-satisfies-slice-rank}
applied to \(L\), so every integral map on it extends. Every
character trivial on the difference lattice of that coset acts
on the named vector by its value on the coset. After transport
it acts on the moved test by its value on \(R^{-1}x\), independently
of the residual position. If \(R^{-1}x\) lies outside the coset
some such character distinguishes it. Indeed choose bases for
the two lattices, divide by \(\alpha_v\pi_v^{l_v}\), and take a
nonintegral entry of the difference. By the additive annihilator
calculation of Proposition~\ref{prop:marked-lattice-unitary-fourier}
an integral multiple of this entry has nontrivial normalized
character. The corresponding dual functional is trivial on all
integral entries. Hence the moved slice at that \(x\) is zero
pointwise. Summation over residual positions proves the claim
for \(t_x\), and linearity proves the last assertion.
\end{proof}

\begin{theorem}\label{thm:singular-small-gram-coefficients-vanish}
For the standard test \(S\) of
Corollary~\ref{cor:small-period-trivial-stabilizer} all singular
Gram coefficients of the small period vanish at \(g\) preserving
\(L\). Every nonzero coefficient has index rationally isometric
to \((A')^-\). For the corresponding large-side test \(E_S\),
every contributing rational matrix has rank \(b\). If the large
real test, in its mixed position variables, is supported at each
infinite place in a closed set of frames whose Gram matrices
are nonsingular of the signature of \((A')^-\) there, every
large-side coefficient of a different signature or of singular
index vanishes at these \(g\).
\end{theorem}
\begin{proof}
On the small side use the conjugate marked branch of
Theorem~\ref{thm:opposite-branch-small-period}.
The proof of Corollary~\ref{cor:small-period-trivial-stabilizer}
shows that each support matrix restricts invertibly to the
reduction of \(\mathfrak l\): the columns of the conjugate
transpose of the lifted \(S\) are an integral basis there.
Lemma~\ref{lem:frame-slice-support-transport} now shows that
\(\alpha_v^{-1}(h^{-1}xg_L)_v\) is an integral isomorphism
of lattices whenever \(t_x\ne0\). Its determinant is a unit
by reduction. Since \(h,g_L\) and \(\alpha_v\) are invertible
in the field, \(x\) is invertible at this place and hence
over \(E\). Thus \(x^*J_Tx\) is a nonsingular rational form
isometric to \(T=(A')^-\). The coefficient formula of
Theorem~\ref{thm:truncated-gram-period-formula} gives the two
assertions for the small period.

On the large side, every named point in \(E_S\) is invertible
on the reduced lattice of \(L^\perp\) by
Proposition~\ref{prop:standard-test-satisfies-slice-rank}.
The lattice of \(L\) has independent reductions inside that
of \(L^\perp\), since it is a direct summand there by the
same basis elimination. Thus the restriction to
\(\mathfrak l\) has rank \(b\) modulo the maximal ideal.
The coset condition of Lemma~\ref{lem:frame-slice-support-transport}
implies full column rank at \(w\), hence rational rank \(b\).
At infinity that lemma puts \(h^{-1}xg_L\) in the prescribed
closed set whenever the slice is nonzero. Its Gram matrix is
\(g_L^*(x^*J_Tx)g_L\), since \(h\) is unitary. Congruence by
an invertible matrix preserves nonsingularity and signature:
it bijects the subspaces on which the form is positive definite
and also those on which it is negative definite. The maximal
dimensions of these two kinds of subspaces are the respective
sign counts by diagonalization as in
Lemma~\ref{lem:rational-hermitian-witt-extension}. Thus the
retained rational indices have the claimed signatures, and the
same coefficient formula finishes the proof.
\end{proof}

\section{Real data on a neighborhood of a nondegenerate plane}
\begin{lemma}\label{lem:explicit-smooth-hermitian-plane-section}
Let \(J\) be an invertible hermitian complex \(m\)-by-\(m\)
matrix, \(X_0\) an \(m\)-by-\(b\) frame with invertible
\(J_b=X_0^*JX_0\), and \(P_0=X_0J_b^{-1}X_0^*J\).
For \(X\) with \(X^*JX\) invertible set
\(P_X=X(X^*JX)^{-1}X^*J\).
On the neighborhood \(\|P_X-P_0\|<1\) in these frames there
are smooth functions \(k(P_X)\in U(J)\) and
\(c(X)\in{\rm GL}_b(\C)\) with
\[
 k(P_0)=I,\qquad X=k(P_X)X_0c(X),\qquad
        c(X)^*J_b c(X)=X^*JX.
\]
The norm is the Euclidean operator norm. For every invertible
\(D\) one has \(P_{XD}=P_X\) and \(c(XD)=c(X)D\).
\end{lemma}
\begin{proof}
Write \(M^\dagger=J^{-1}M^*J\). The formula for \(P=P_X\)
gives \(P^2=P=P^\dagger\) and \({\rm im}\,P={\rm im}\,X\);
the corresponding identities hold for \(P_0\). Put
\[
   A=PP_0+(1-P)(1-P_0),\quad Z=(P-P_0)^2,\qquad
   k(P)=A\sum_{n\ge0}4^{-n}\binom{2n}{n}Z^n .
\]
Then \(PA=AP_0\), \(Z^\dagger=Z\), \(ZP_0=P_0Z\) and
\(A^\dagger A=1-Z\), by multiplication and idempotence.
Since \(\|Z\|<1\), the series is convergent and smooth there.
Indeed its coefficient is at most one, and on \(\|Z\|\le r<1\)
any derivative of order \(s\) has summable bounds by
\(n^s r^{n-s}\) times the norms of the \(s\) directions.
For \(u(z)=\sum 4^{-n}\binom{2n}{n}z^n\), the ratio of
successive coefficients gives \(2(1-z)u'=u\), \(u(0)=1\);
hence \((1-z)u^2=1\). Multiplication of the absolutely
convergent matrix series thus gives \(k(P)^\dagger k(P)=I\).
This identity implies invertibility by determinants. Since
\(Pk(P)=k(P)P_0\), \(k(P)\) carries the original plane onto
\({\rm im}\,X\), and at \(P=P_0\) it is the identity. Take
\(c(X)=J_b^{-1}X_0^*J k(P)^{-1}X\). This is the unique
coordinate matrix of the frame in the transported basis and is
invertible. Taking Gram matrices gives the asserted equality.
The projection formula is invariant under \(X\mapsto XD\);
the coordinate formula then gives the final transformation.
\end{proof}

\begin{lemma}\label{lem:real-oscillator-smooth-near-identity}
For the real oscillator on a hermitian Hom space used in
Definition~\ref{def:two-order-theta-kernel}, with determinant
characters of integer weights, the function \(k\mapsto\pi(k)v\)
near the identity in either real unitary factor is smooth in
every Schwartz seminorm when \(v\) is Schwartz.
\end{lemma}
\begin{proof}
In real symplectic position and dual coordinates write the
matrix of the Hom action as \(\left(\begin{smallmatrix}
A&B\\ C&D\end{smallmatrix}\right)\). Near the identity it factors
as the lower shear \(CA^{-1}\), the diagonal with blocks
\(A,A^{-\mathsf T}\), and the upper shear \(A^{-1}B\), in that
order. Indeed preservation of the alternating pairing makes
both shear matrices symmetric and gives
\(D=CA^{-1}B+A^{-\mathsf T}\). By
Theorem~\ref{thm:graded-parabolic-schroedinger} and
Lemma~\ref{lem:schroedinger-real-gaussian}, unitary implementers
of these factors are respectively multiplication by a quadratic
phase, normalized linear substitution, and a quadratic phase
conjugated by a fixed Fourier transform. The corresponding
product \(V(k)\) acts smoothly in Schwartz seminorms. To see
this, each parameter derivative of the phase inserts a polynomial
in the position variables. Differentiation of a linear substitution
inserts a first derivative multiplied by a linear function.
On a compact neighborhood the substitution and its inverse
have bounded norms. These observations bound each differentiated
expression by finitely many original seminorms; the integral
mean-value formula in the parameter verifies convergence of
the difference quotients in each seminorm. Fourier transform is
continuous in these seminorms by differentiation and integration
by parts in Lemma~\ref{lem:schroedinger-real-gaussian}.

The true implementer differs by a scalar. Its vacuum coefficient
is the inverse determinant of the negative diagonal block by
Theorem~\ref{thm:negative-block-real-splitting}, before multiplication
by the pair character. The vacuum coefficient of \(V(k)\) is
smooth by the preceding seminorm bounds and equals one at
identity. In a smaller neighborhood the ratio of the prescribed
coefficient to the computed one is therefore smooth. The
determinant characters there are integer powers of smooth unit
complex determinants by Definition~\ref{def:torus-weight-and-units}.
Multiplying by these scalars proves the claim.
\end{proof}

\begin{theorem}\label{thm:matched-real-gram-slices}
At the infinite places take \(J=J_b\oplus J_n\), with signed
diagonal forms of ranks \(b,n>0\), and \(X_0=\binom{I_b}{0}\)
in Lemma~\ref{lem:explicit-smooth-hermitian-plane-section}.
Let \(n=2j\) and \(\mu_\tau\) be nonnegative integer weights.
Use the mixed positions of Definition~\ref{def:gram-indexed-mixed-period}
for targets of rank \(b+n\) and \(b\). In the residual oscillator
for the first target identify the two orthogonal factors as in
Theorem~\ref{thm:residual-levi-zero-section}. Let \(v,v_1\)
be arbitrary real Schwartz vectors on the two respective
factors. Denote by \(\pi\) the residual representation of the
real target group with any fixed unitary determinant characters
of integer weights. There exist real oscillator data \(F,F'\)
of Schwartz class and a nonnegative smooth function \(\rho\)
on \(\prod_\tau U(J)\), positive near identity, with these
properties. The opposite slice of \(F'\) at zero equals \(v\)
in self-dual Fourier coordinates. The frame support of \(F\)
is compact and has nonsingular Gram matrices of signature \(J_b\)
at each place. For every \(h\in\prod_\tau U(J)\) some
\(p\in\prod_\tau(U(J_b)\times U(J_n))\) satisfies, for every
\(B\in\prod_\tau{\rm GL}_b(\C)\),
\[
 \pi(p h) F((p h)^{-1}X_0B)
  =\rho(h)\prod_\tau\left[
      \left(\frac{\det B_\tau}{\overline{\det B_\tau}}\right)^{-\mu_\tau}
      |\det B_\tau|^{-2j}\right]\big(F'(B)\otimes v_1\big).
\]
Here evaluation in a frame variable denotes the remaining
Schwartz vector. The last product, including both determinant
factors at each place, is the inverse of the archimedean
base-change character of weight \(\mu\) times
\(\prod_\tau|\det B_\tau|^{-n}\).
\end{theorem}
\begin{proof}
Choose smooth nonnegative bumps \(\beta(P)\) and \(\eta(C)\)
on the finite products of matrix spaces, positive near \(P_0\)
and \(I\), respectively. For \(\beta\) take compact support so
small that \(\|P-P_0\|<1\) placewise and the \(k(P)\)'s on
the support are in the neighborhoods of
Lemma~\ref{lem:real-oscillator-smooth-near-identity}. For \(\eta\)
take compact support inside the invertible matrices. Scaled
translates of the ball bump in
Lemma~\ref{lem:frame-slice-support-transport} give these choices.
Use self-dual real position measure \(\dd C\) for the mixed
coordinates on the small space and put
\[
  m(C)=\eta(C)\prod_\tau|\det C_\tau|^{2(j+\mu_\tau)},
       \qquad A=\left(\int m(C)\dd C\right)^{-1}.
\]
The integral is finite and positive by compact support and
positivity near the identity. Set \(F'(C)=A m(C)v\). Exchanging
the two isotropic spaces is partial Fourier transformation in
\(C\); evaluation at zero is integration in \(C\), by
Lemma~\ref{lem:schroedinger-real-gaussian} in dual coordinates.
This gives the stated opposite slice.

For \(P=P_X\) and \(c=c(X)\) in the plane chart put
\[
    F(X)=A\,\beta(P)\eta(c)\!
          \prod_\tau\overline{\det c_\tau}^{\,2\mu_\tau}
                   \,\pi(k(P))(v\otimes v_1),
\]
and extend by zero outside the chart. Its frame support lies
in the image under \((P,c)\mapsto k(P)X_0c\) of a compact
set: use the intersection of the support of \(\beta\) with
the closed set of self-adjoint idempotents and the support of
\(\eta\). On the component used here the rank equals \(b\),
since \(k(P)\) carries \(P_0\) onto \(P\).
These image frames have invertible Gram matrices of the desired
signature and lie inside the chart. The extended function is
smooth by Lemmas~\ref{lem:explicit-smooth-hermitian-plane-section}
and~\ref{lem:real-oscillator-smooth-near-identity}; all derivatives
in the frame variables take bounded values in every residual
Schwartz seminorm on this compact support. Thus it is Schwartz
in the joint position variables and defines real oscillator data.

Put \(P=P_{h^{-1}X_0}\) and \(\rho(h)=\beta(P)\). This plane is
always nondegenerate. If \(\rho(h)=0\), take \(p=1\);
right multiplication of the frame by \(B\) leaves its plane
unchanged, so both sides vanish. Otherwise
\(h^{-1}X_0=k(P)X_0c_0\), with \(c_0\in U(J_b)\) by its
Gram matrix. Hence \(h k(P)\) preserves the original plane
and its perpendicular, and acts there as \(c_0^{-1}\) and
some \(q\in U(J_n)\). Take \(p=\operatorname{diag}(c_0,q^{-1})\)
at each place. Then \(p h=k(P)^{-1}\). The argument of \(F\)
has plane \(P\) and coordinate \(c=B\). The two group operators
cancel by the representation law. On substituting \(F'(B)\)
the remaining quotient is the product of
\(\bar z^{2\mu}/|z|^{2(j+\mu)}
=(z/\bar z)^{-\mu}|z|^{-2j}\), where \(z=\det B_\tau\).
A torus character of weight \(\mu_\tau\) evaluates on
\(z/\bar z\) as its \(\mu_\tau\)-th power by
Definition~\ref{def:torus-weight-and-units}. The support and
normalization already obtained prove all assertions.
\end{proof}

\section{The Hecke values at the positive integers}
\begin{definition}\label{def:complementary-hecke-values}
Let \(E\) be the CM field, \(D=[E:\Q]=2|\Phi|\), and let
\(\nu\) be a character of the norm-one torus with weight
\(\mu=(\mu_\tau)_{\tau\in\Phi}\) as in
Definition~\ref{def:torus-weight-and-units}. On the ideles of
\(E\) put \(\lambda(x)=\nu(x/\bar x)\); write \(\lambda_f\)
for its finite part. Let \(S\) be a finite set of places,
disjoint from its conjugate, such that \(\lambda_f\) is trivial
on the local unit groups outside \(S\cup\bar S\).
For each such unmarked place \(v\) take any uniformizer
\(\pi_v\), inserted as an idele with other components one.
Define in \(\Re s>1\)
\[
  L^S(s,\nu)=\prod_{v\notin S\cup\bar S}
             (1-\lambda_f(\pi_v)q_v^{-s})^{-1},
                \qquad Q_S=\prod_{w\in S}q_w .
\]
The values of \(\lambda_f(\pi_v)\) are independent of the
uniformizers by triviality on units. The character on principal
ideles is one by rational triviality of \(\nu\), and its infinite
value on \(x\) is \(\prod_\tau (x_\tau/|x_\tau|)^{2\mu_\tau}\).
\end{definition}

\begin{lemma}\label{lem:angular-log-sector-count}
For a nonzero tuple \(k\in\Z^\Phi\) and \(x_\tau\ne0\) put
\(u_k(x)=\prod_\tau(x_\tau/|x_\tau|)^{k_\tau}\).
Use the sector \(\mathcal R\) of Lemma~\ref{lem:log-sector-count}
in the totally imaginary field \(E\). For each full real lattice
\(\Lambda\subset E_\infty\), each translate \(z+\Lambda\), and
\(Y\ge1\), one has
\[
       \left|\sum_{x\in(z+\Lambda)\cap Y^{1/D}\mathcal R}
                         u_k(x)\right|\ll_{k,\Lambda,\mathcal R}
                              1+Y^{1-1/D}.
\]
\end{lemma}
\begin{proof}
The integral of \(u_k\) over the dilated sector is zero. A rotation
in a complex coordinate with \(k_\tau\ne0\) preserves the sector
and measure and may be chosen to multiply the integrand by a
scalar different from one. Put \(T=Y^{1/D}\).
Tile by the fixed lattice cells based at the points of
\(z+\Lambda\) as in Lemma~\ref{lem:log-sector-count}. The boundary
meets \(O(1+T^{D-1})\) cells by its proof, each contributing a
bounded error when the lattice sum is compared with the integral
divided by the cell volume. The finitely many cells in a fixed
ball also give a bounded error, uniformly in the translate.

On all other interior cells, for \(\|x\|\) between \(R\) and \(2R\)
at a corner, every \(|x_\tau|\gg R\) on that cell. Indeed the
logarithmic coordinates of the corner lie in the bounded
parallelepiped of the sector and each cell has bounded diameter.
The real gradient of \(u_k\) on the cell has size \(O_k(R^{-1})\),
by differentiation of each unit complex phase. Its variation
is therefore \(O_k(R^{-1})\) by integration along segments.
There are \(O_\Lambda(R^D)\) such cells for \(R\ge1\), by comparing
the volumes of their disjoint interiors with a ball of radius
\(2R+O_\Lambda(1)\). The sector lies in a ball of radius \(O(T)\).
Summing the resulting \(O(R^{D-1})\) over successive dyadic
radii gives \(O(1+T^{D-1})\) since \(D\ge2\). This proves
the estimate.
\end{proof}

\begin{theorem}\label{thm:torus-hecke-nonvanishing-at-one}
If \(\mu\ne0\) and the finite unit character of \(\nu\) factors
through residues at the indicated split places, the product in
Definition~\ref{def:complementary-hecke-values} converges
absolutely, with nonzero value, in \(\Re s>1\). Its continuation
is holomorphic on \(\Re s>1-1/D\) and satisfies \(L^S(1,\nu)\ne0\).
In particular all its positive integer values are well defined
and nonzero.
\end{theorem}
\begin{proof}
Temporarily omit all primes over a positive rational integer
\(m\) divisible by the rational primes below \(S\cup\bar S\);
take \(m\) divisible by their required conductor powers.
At every divisor of \(m\), the character of finite units
congruent to one modulo \(m\) is then one by the residue
hypothesis. For an ideal \(\mathfrak a\) prime to \(m\)
write \(u(\mathfrak a)\) for the multiplicative product of
the specified uniformizer values, with powers the ideal
valuations. By Proposition~\ref{prop:general-ideal-factorization},
expanding the Euler factors gives the ideal Dirichlet series
with coefficients \(u(\mathfrak a)\) in \(\Re s>1\). In fact
the counting estimate of Theorem~\ref{thm:partial-ideal-counts}
gives \(\sum_{\mathfrak a}(N\mathfrak a)^{-\sigma}<\infty\)
for \(\sigma>1\), by integration. The series of prime-power
logarithms is absolutely convergent there, by comparison with
\(\sum_v q_v^{-\sigma}/(1-2^{-\sigma})\), so its exponential
is nonzero.

For a nonzero integer \(r\), apply the ray class enumeration
in Proposition~\ref{prop:general-ray-norm-count}. In a fixed
class choose an integral ideal \(\mathfrak b\) in the inverse
class prime to \(m\). The integral ideals \(\mathfrak a\) in
the class correspond to \(\mathfrak a\mathfrak b=(x)\),
with \(0\ne x\in\mathfrak b,\ x\equiv1\bmod mO_E\), modulo
the units \(U\) congruent to one there. Since \(E\) is totally
imaginary there is no sign condition. On such generators
\[
   u(\mathfrak a)^r
       =u(\mathfrak b)^{-r}\!
          \prod_{\tau\in\Phi}
             \left(\frac{\tau(x)}{|\tau(x)|}\right)^{-2r\mu_\tau}.
\]
This follows by evaluating \(\lambda\) on the principal idele
of \(x\): its components above \(m\) contribute one, and every
other component gives the uniformizer value to its valuation.
The right side is constant on \(U\)-orbits by the same argument
on units. The congruence and membership in \(\mathfrak b\)
specify a translate of \(m\mathfrak b\), by the Chinese
remainder argument in Proposition~\ref{prop:general-ray-norm-count}.
In the sector of a logarithmic basis subgroup \(U_0\subset U\)
every ideal occurs exactly \([U:U_0]\) times by
Lemma~\ref{lem:unit-log-domain} and the orbit argument of
Theorem~\ref{thm:partial-ideal-counts}. Its norm bound \(Y\)
becomes \(|N_{E/\Q}x|\le YN\mathfrak b\).
Apply Lemma~\ref{lem:angular-log-sector-count} with \(k=-2r\mu\),
which is nonzero. Summing the bounds over finitely many classes
gives
\[
       \sum_{\substack{N\mathfrak a\le Y\\(\mathfrak a,m)=1}}
             u(\mathfrak a)^r =O_{E,m,\nu,r}(1+Y^{1-1/D}).
\]
Partial summation, first for \(\Re s>1\), writes the corresponding
series \(L_m(s,u^r)\) as \(s\) times the integral of this counting
function against \(Y^{-s-1}\dd Y\) from one to infinity.
Lemma~\ref{lem:analytic-tools} gives holomorphy on
\(\Re s>1-1/D\).

Write \(\zeta_m(s)\) for the series of the trivial ideal character
with the same deletions. It has a simple pole at one by
Proposition~\ref{prop:general-ray-norm-count}. For each real
\(\sigma>1\) the prime-power logarithms give
\[
       \zeta_m(\sigma)^3 |L_m(\sigma,u)|^4
                          |L_m(\sigma,u^2)|\ge1.
\]
Indeed the coefficient in the logarithm at the \(l\)-th power
of a prime is
\((3+4\Re t+\Re(t^2))/(l q_v^{l\sigma})\), where
\(t=u(\mathfrak p_v)^l\) has modulus one. Its numerator is
\(2(1+\Re t)^2\). If \(L_m(1,u)=0\), holomorphy bounds the fourth
power by \(O((\sigma-1)^4)\), and holomorphy for \(r=2\) bounds
the last factor near one. The left side would tend to zero.
This proves nonvanishing at one. Restoring the finitely many
unramified factors over \(m\) uses factors holomorphic and
nonzero on \(\Re s>0\), proving every assertion for \(L^S\).
\end{proof}

\begin{corollary}\label{cor:positive-exchange-lvalue-product}
Suppose the ambient rank \(N\), the number \(q\) of its negative
directions, and the integer weights satisfy
\(2q+2\le N\le2\mu_\tau+2q\) for every \(\tau\).
For \(e\in\Z_{\ge0}\) put
\[
       \Pi_S(e,\nu)=Q_S^{e^2}\prod_{j=1}^e|L^S(j,\nu)|^2 .
\]
Then \(\Pi_S(0,\nu)=1\) and \(\Pi_S(e,\nu)>0\) for every
\(e\), under the finite hypotheses of
Theorem~\ref{thm:torus-hecke-nonvanishing-at-one}. If a comparison
satisfies \(|c_\nu|^2=C\Pi_S(e,\nu)\) with \(C>0\), division
by \(c_\nu\) in Theorem~\ref{thm:single-frame-global-period-criterion}
is valid.
\end{corollary}
\begin{proof}
The two rank inequalities give \(\mu_\tau\ge1\). Each factor is
positive by Theorem~\ref{thm:torus-hecke-nonvanishing-at-one},
and the empty product and \(Q_S^0\) equal one. The indicated
equality then gives a positive squared modulus.
\end{proof}

\section{A finite Mellin transform with one displaced unit}
\begin{lemma}\label{lem:triangular-line-gauss-integral}
Let \(O\) be a discrete valuation ring in a nonarchimedean local
field \(K\), \(\pi\) a uniformizer, and \(k=O/\pi O\) of order \(q\).
Give \(O\) additive mass one. Let \(\psi\) be an additive character
whose annihilator of \(O\) under multiplication is \(O\).
Let \(\chi\) be a nontrivial character of \(k^\times\), and put
\(d(x)=\chi(x\bmod\pi)^{-1}\) on \(O^\times\), zero elsewhere.
For \(t\in K\) define \(g_\chi=\int_O d(z)\psi(\pi^{-1}z)\dd z\).
Then \(|g_\chi|^2=q^{-1}\) and
\[
  \int_O d(z)\psi(tz)\dd z =
       \begin{cases}
       g_\chi\chi(\pi t\bmod\pi),&\pi t\in O^\times,\\
       0,&\text{otherwise}.
       \end{cases}
\]
For a fixed square matrix \(Y\), \(a\in K\), take the test
on a block upper triangular matrix
\[
  H=\begin{pmatrix}a&u&v\\0&Y&w\\0&0&z\end{pmatrix},
       \qquad {\bf1}_{H\ {\rm integral}}\,d(\det H).
\]
Its integral in the three off-diagonal blocks and in \(z\),
against \(\psi(tz)\) and additive measures giving each integral
box mass one, is \({\bf1}_{Y\ {\rm integral}}d(\det Y)d(a)\)
times the displayed scalar. The same conclusion holds after
interchanging the two scalar diagonal entries.
\end{lemma}
\begin{proof}
On a residue coset \(r+\pi O\) the multiplier \(d(z)\) is constant.
If \(t\pi\notin O\), translation by an element of \(\pi O\)
multiplies the character integral there by a number different
from one. If \(t\in O\), the character is trivial and the sum
of the nontrivial multiplicative character over \(k^\times\)
vanishes by translation in that group. In the remaining case
\(t=\pi^{-1}u_0\) with \(u_0\in O^\times\), use the substitution
\(y=u_0z\); it preserves additive measure and gives the factor
\(\chi(u_0\bmod\pi)\). The character
\(r\mapsto\psi(\pi^{-1}\widetilde r)\) on \(k\) is nontrivial
by the annihilator assumption. Thus \(qg_\chi\) is the
one-by-one determinant Gauss sum of
Lemma~\ref{lem:finite-determinant-fourier-stabilizer} with
multiplicative character \(\chi^{-1}\). Its primitivity on a
field follows by taking \(1+zr\) to be any prescribed unit for
\(r\ne0\). This gives the claimed norm. The full matrix test
factors on the integral box because
\(\det H=a(\det Y)z\), and an integral product is a unit exactly
when each factor is a unit. The off-diagonal integrals have mass
one. The determinant equality holds in the other order as well.
\end{proof}

\begin{definition}\label{def:ramified-line-mellin}
In Definition~\ref{def:complementary-hecke-values} assume the finite
unit part of \(\nu\) is \(\prod_{w\in S}\chi_w\), with every
\(\chi_w\) a nontrivial residue character at a split branch and
with trivial remaining unit characters. Use conjugate uniformizers
on conjugate branches, and local additive characters as in
Proposition~\ref{prop:marked-lattice-unitary-fourier}. Put
\(G_S=\prod_{w\in S}g_{\chi_w}\) and take a finite idele \(c\)
with \(c_{\bar w}=\bar\pi_w^{-1}\) for \(w\in S\) and one
elsewhere. For a finite idele \(x\) let \(D_S(x)\) be zero
unless \(x_v\in O_v\) off \(S\cup\bar S\), and
\(x_w,\pi_w\overline{x_{\bar w}}\in O_w^\times\) for all
\(w\in S\). On this support put
\[
 D_S(x)=\prod_{w\in S}\chi_w(x_w\bmod\pi_w)^{-1}
           \chi_w(\pi_w\overline{x_{\bar w}}\bmod\pi_w).
\]
Give the finite idele group the invariant multiplicative
measure for which \(\prod_v O_v^\times\) has mass one. This
means summation over the valuation vectors of finite support,
with the translated unit probability measure on each coset.
Write \(|x|_f=\prod_v q_v^{-v(x_v)}\). The finite Mellin integral
for real \(\sigma>1\) is
\[
             M_S(\sigma)=G_S\int D_S(x)\lambda_f(x)|x|_f^\sigma
                                                    \dd^\times x .
\]
The \(D_S\) in this definition is an idele function; the
determinant tests in Definition~\ref{def:integral-level-determinant-sections}
are functions on matrices.
\end{definition}

\begin{theorem}\label{thm:finite-mellin-one-plane-ratio}
In Definition~\ref{def:ramified-line-mellin}, for \(\sigma>1\)
the integral converges absolutely and
\[
      M_S(\sigma)=G_S\,\lambda_f(c)\,Q_S^\sigma L^S(\sigma,\nu).
\]
If \(\mu\ne0\), define \(M_S(1)=\lim_{\sigma\downarrow1}M_S(\sigma)\).
This limit exists and is nonzero. At \(\sigma=1\) the integral
of the absolute value of the original finite integrand is infinite.
For any integer \(e\), put
\(\Pi^{\rm int}(e)=\Pi_S(e,\nu)\) when \(e\ge0\) and
\(\Pi^{\rm int}(e)=\Pi_S(-e,\nu)^{-1}\) otherwise. With
\(s_e=\max(e+1,-e)\) and \(\mu\ne0\), one has
\[
          |M_S(s_e)|^2=\frac{\Pi^{\rm int}(e+1)}{\Pi^{\rm int}(e)} .
\]
\end{theorem}
\begin{proof}
Lemma~\ref{lem:triangular-line-gauss-integral} at each chosen branch
with \(a=x_w\) and \(t=\overline{x_{\bar w}}\) gives the factor
\(G_SD_S(x)\) there with \(|G_S|^2=Q_S^{-1}\).
Write \(x=c y\). Now \(y\) has unit components on all the
marked pairs and integral components elsewhere. Write \(y=u p\)
with a unit idele \(u\) and an idele \(p\) given by nonnegative
uniformizer powers off the marked pairs and by one there.
On units, \(\lambda_f(u)=\nu(u/\bar u)\) takes the value
\(\prod_w\chi_w(u_w)\chi_w(\overline{u_{\bar w}})^{-1}\),
with reductions understood. Thus \(D_S(x)\lambda_f(u)=1\).
Since \(|c|_f=Q_S\), every allowed valuation coset has
constant integrand after the unit phases cancel. Summing over
these cosets gives the ideal series of
Theorem~\ref{thm:torus-hecke-nonvanishing-at-one}, with the
factor \(G_S\lambda_f(c)Q_S^\sigma\). For \(\sigma>1\)
the absolute sum is bounded by the unrestricted positive ideal
series in Corollary~\ref{cor:ray-series-at-one}, so this calculation
is valid. For \(\mu\ne0\), continuation at one and nonvanishing
are given by Theorem~\ref{thm:torus-hecke-nonvanishing-at-one}.

For absolute values at one the same coset calculation gives
\(|G_S|Q_S\sum_{\mathfrak a}(N\mathfrak a)^{-1}\) with the
finitely many marked divisors excluded. As \(\sigma\downarrow1\)
the unrestricted positive series has a simple pole by
Corollary~\ref{cor:ray-series-at-one}; deleting finitely many
Euler factors multiplies by a factor with positive limit.
Monotonicity of the terms proves divergence at one.
Finally if \(e\ge0\), the ratio in the assertion is
\(Q_S^{2e+1}|L^S(e+1,\nu)|^2\), by
Corollary~\ref{cor:positive-exchange-lvalue-product} and its
definition. For \(e=-k-1<0\) it is instead
\(\Pi_S(k+1,\nu)/\Pi_S(k,\nu)\), giving
\(Q_S^{2k+1}|L^S(k+1,\nu)|^2\). In both cases \(s_e\ge1\)
and the evaluated Mellin transform has precisely that squared
modulus because \(|\lambda_f(c)|=1\) and \(|G_S|^2=Q_S^{-1}\).
\end{proof}

\section{Finite slices on an integral Witt decomposition}
\begin{definition}\label{def:adapted-finite-frame-slices}
Use the finite places of the arithmetic field \(E\). Suppose that the
splitting \(W=L\oplus W_0\oplus L'\) of
Definition~\ref{def:gram-indexed-mixed-period} splits the finite
lattice as the sum of lattices on the three spaces. The first and
last lattices are in perfect hermitian duality, and the middle
lattice is unimodular. Let \(b=\dim L>0\), \(n=\dim W_0\),
\(T=T_b\perp W_0\), where \(T_b\) has dimension \(b\).
Use unimodular finite lattices on both summands and their direct
sum on \(T\). At each marked split pair choose ordered integral
frames in the order \(L,W_0,L'\) and conjugate dual frames with
the pairing between \(L,L'\) equal to the identity. Do the same
in the order \(T_b,W_0\) on the target. At the branch \(w\in S\)
take \(S_w\) in Definition~\ref{def:integral-level-determinant-sections}
to be projection on the last \(b\) coordinates. Its conjugate
adjoint inserts the first \(b\) columns at \(\bar w\).
Here and below coordinates on Hom lattices are divided by
the finite scale \(\alpha\) of
Lemma~\ref{lem:hom-oscillator-dual-lattices}.

In the mixed model use both branches of the restriction to \(L\)
as frame position \(Y\). At the pair \(\{w,\bar w\}\) the
remaining position is the restriction \(Z\) to \(W_0\) at
the branch used to name the test. Normalize the finite models
to take each reference vacuum to its lattice indicator, with
position lattice of mass one. Write \(\mathbf 1_{\rm int}\)
for the indicator that all specified entries are integral.
Extend \(\chi_w^{-1}\) by zero from residue units, and denote this
extension by \(d_w\), also applied to integral lifts. The depth
is the one in the determinant datum. Write \(F^+\) for the finite
large datum \(E_S\), and \(F^-\) for the transformed small datum
\(D_S^\flat\) of Theorem~\ref{thm:opposite-branch-small-period}.
\end{definition}

\begin{lemma}\label{lem:finite-witt-frame-partial-fourier}
In Definition~\ref{def:adapted-finite-frame-slices}, an integral-box
test at a split pair named at \(w\), whose value on the box has
the form \(H(Y_w,Z_w)\), becomes
\[
  {\bf1}_{\rm int}(Y_w,Y_{\bar w},Z_w)\,H(Y_w,Z_w)
\]
in the mixed model. The same formula with the branches interchanged
holds for a test named at \(\bar w\). Off the marked pairs the
vacuum slice is the integral-frame indicator times the residual
vacuum.
\end{lemma}
\begin{proof}
At \(w\) the three source blocks of a position matrix are
\((Y_w,Z_w,U_w)\) by Theorem~\ref{thm:split-finite-schroedinger}.
In the conjugate branch the block paired with \(U_w\) is
\(Y_{\bar w}\), because \(L'\) pairs with \(L\) and is orthogonal
to \(W_0\) and to itself. The change of polarization therefore
takes the Fourier transform in \(U_w\) alone. It implements the
exchange of translations and multipliers on these coordinates
by substitution in the Fourier integral, as in
Lemma~\ref{lem:rational-partial-poisson}. The kernel is the trace
character from Lemma~\ref{lem:hom-oscillator-dual-lattices}.
The annihilator of the integral \(U_w\)-box is exactly the
integral \(Y_{\bar w}\)-box by the trace-dual identity there
and perfect duality of the two frames. Fourier transformation
of the first box is consequently the indicator of the second
box. Indeed it integrates to one for a trivial character and
translation by an element on which the character is nontrivial
forces the integral to vanish otherwise. This also fixes the
normalization of the intertwiner, since it takes the vacuum to
the vacuum. The other variables retain their factors. Interchanging
the places gives the second assertion. At an unmarked place the
hyperbolic position and momentum lattices split and are dual,
so their vacuum is the normalized frame indicator by
Proposition~\ref{prop:finite-polarized-vacuum}. Translations
in the residual Hom space act independently, giving the tensor
product with its vacuum by the direct sum argument of
Theorem~\ref{thm:residual-levi-zero-section}.
\end{proof}

\begin{theorem}\label{thm:block-frame-finite-exchange}
In Definition~\ref{def:adapted-finite-frame-slices} the large
and small finite slices at a marked pair are respectively
\[
\begin{aligned}
  F^+(Y;Z)&={\bf1}_{\rm int}(Y_w,Y_{\bar w},Z_w)\,
      d_w\!\left((-1)^{b(b+n)}\det[\,Y_w\; Z_w\,]\right),\\
  F^-(X;Z')&={\bf1}_{\rm int}(X_w,X_{\bar w},Z'_{\bar w})\,
      d_w\!\left((-\epsilon_w^{-1})^b
                      \overline{\det X_{\bar w}}\right).
\end{aligned}
\]
Here \(Y\) has target \(T\), \(X\) has target \(T_b\), and
\(\epsilon\) is the integral frame unit of
Lemma~\ref{lem:opposite-source-frame-and-density}. The bar takes
scalars from \(E_{\bar w}\) to \(E_w\).

Let \(X=(X_v)_v\) be a finite frame in \(T_b\) which is an
integral isomorphism of the frame lattice onto the lattice of
\(T_b\) at every finite place. Set \(Y=\binom X0\) and
\(\lambda_f(x)=\nu_f(x/\bar x)\), with unit character as in
Definition~\ref{def:ramified-line-mellin}. In the compatible
frames, the finite residual vector satisfies
\[
   F^+\!\left(\tbinom X0\right)
         =u_S\,\lambda_f(\det X)^{-1}\,F^-(X)\otimes F_{\rm eq},
\]
where \(F_{\rm eq}\) is the integral determinant test
\(d_w(\det Z_1)\) on \({\rm Hom}(W_0,W_0)\) at \(w\in S\)
and the vacuum elsewhere. The fixed scalar \(u_S\) has modulus
one and is the product, over \(w\), of
\[
        d_w((-1)^{b(b+n)})\,/\,d_w((-\epsilon_w^{-1})^b).
\]
The determinant of \(X\) here is the finite unit idele in the
chosen ordered frames.
\end{theorem}
\begin{proof}
On the large box the determinant is that of the matrix with
first block row \((0,0,I_b)\) and remaining block rows
\((Y_w,Z_w,U_w)\). Moving its last \(b\) columns ahead of its
first \(b+n\) contributes the displayed sign, after which the
matrix is block triangular. On the small box the determinant
of \(S(-\epsilon^{-1}Z^*)\) in
Theorem~\ref{thm:opposite-branch-small-period} is the determinant
of the conjugate transpose of the restriction to \(L\), times
\((-\epsilon_w^{-1})^b\). In each case the box value is
independent of the \(L'\)-block, so
Lemma~\ref{lem:finite-witt-frame-partial-fourier} proves both
slice formulas.

At \(Y=\binom X0\), write \(Z_w=\binom{Z_b}{Z_1}\). Then
\(\det[Y_w\;Z_w]=\det X_w\det Z_1\). The free integral
\(Z_b\)-block is the reference vacuum on
\({\rm Hom}(W_0,T_b)\). The small test has its vacuum on the
same space in the conjugate-branch model. These give the same
abstract vector since both identifications take the reference
vacuum to the appropriate box, by
Theorem~\ref{thm:split-finite-schroedinger}. On comparing the
two scalar unit factors at the marked pair one obtains, besides
the displayed constant,
\(\chi_w(\overline{\det X_{\bar w}})/\chi_w(\det X_w)\).
This is the inverse of the \(w\)-factor of
\(\nu_f(\det X/\overline{\det X})\) on units. All frame
indicators are one and at the unmarked places
Lemma~\ref{lem:finite-witt-frame-partial-fourier} leaves only
the two residual vacua. Their tensor product gives the original
residual vacuum by the same splitting argument. Multiplication
over places gives the asserted equality. All the factors of
\(u_S\) are unit character values.
\end{proof}

\section{The finite orbit in the nonsingular coefficient}
\begin{lemma}\label{lem:finite-unitary-frame-coset}
Let \(T=T_b\perp T_n\) have the orthogonal direct sum of two
unimodular finite hermitian lattices over the arithmetic CM field
\(E\). Let \(F=E^+\), \(G_f=U(T)(\mathbb A_{F,f})\),
\(U_i=U(T_i)(\mathbb A_{F,f})\) for \(i=b,n\), and
\(K_T\subset G_f\) be the stabilizer of the sum lattice.
Write \(\iota:T_b\to T\) for the first inclusion. In integral
frames let \(Y=u A\), where \(u\in U_b\) and
\(A\in{\rm GL}_b(\widehat O_E)\). Then
\[
 \{h\in G_f:h^{-1}\iota Y\ \text{is integral}\}
                   ={\rm diag}(u,U_n)K_T .
\]
For \(y\in U_b\), if \(y^{-1}Y\) is integral then
\(y^{-1}Y\in{\rm GL}_b(\widehat O_E)\).
\end{lemma}
\begin{proof}
Multiplication on the right by \(A\) and its integral inverse
shows that the first integrality condition is equivalent to
integrality of \(B=h^{-1}\iota u\). Its Gram matrix in the
chosen source frame is \(J_b\), the unimodular Gram matrix of
the source lattice. If \(B\) is integral then
\(J_b^{-1}B^*J_T\) is an integral left inverse, where \(J_T\)
is the Gram matrix of the sum lattice. The inclusion \(\iota\)
also has this Gram matrix and an integral left inverse. Apply
Theorem~\ref{thm:integral-unitary-primitive-frame} at each
finite place of \(F\). Its two possibilities for the local
integer algebra exhaust all places, including those above two,
by Lemma~\ref{lem:real-place-rings}. The product of the
resulting isometries is \(k\in K_T\) with \(kB=\iota\), as
in Corollary~\ref{cor:finite-adelic-frame-extension}.
Consequently \(kh^{-1}{\rm diag}(u,1)\) fixes \(\iota\)
pointwise. An isometry with this property restricts to an
isometry on the orthogonal complement, so it has the form
\({\rm diag}(1,l)\). Hence \(h={\rm diag}(u,l^{-1})k\).
Conversely each element of the indicated set takes \(\iota Y\)
by its inverse to \(k^{-1}\iota A\), which is integral.
For the last assertion, \(y^{-1}u\) is integral if \(y^{-1}Y\)
is integral. Its inverse is \(J_b^{-1}(y^{-1}u)^*J_b\),
also integral by unimodularity. Multiplication by \(A\)
gives the assertion.
\end{proof}

\begin{lemma}\label{lem:standard-determinant-fixed-target}
With characters \((\gamma,1)\) as in
Definition~\ref{def:complementary-period-characters}, the entire
standard large datum \(F^+\) of Definition~\ref{def:adapted-finite-frame-slices}
is fixed by \((1,k)\) for every \(k\in K_T\), acting at finite
places and as the identity at infinity.
\end{lemma}
\begin{proof}
The finite lattice action of \(k\) takes a named point \(D\)
to \(kD\) on the chosen branches, as shown in
Lemma~\ref{lem:finite-covariance-of-periods}. This holds for
the full integral group because the complementary places carry
the vacuum, which an integral isometry fixes by
Theorem~\ref{thm:all-finite-weil-splitting}. This theorem also
gives phase one on an integral lattice stabilizer.
Thus the action on a function of the residue names replaces
\(D\) by \(k^{-1}D\) in the function. In \(E_S(k^{-1}D)\)
the lower block rows are multiplied by the reduction of
\(k^{-1}\), so this replacement contributes
\(\chi(\det k)\). It preserves integrality and the support
on unit determinants. Finally the target twist in
Definition~\ref{def:two-order-theta-kernel} is
\(\gamma(\det k)=\chi(\det k)^{-1}\) by
Theorem~\ref{thm:actual-global-splitting-character} and
Definition~\ref{def:complementary-period-characters}.
The factors multiply to one.
\end{proof}

\begin{proposition}\label{prop:actual-coset-finite-coefficient}
Use the adapted integral coordinates of
Definition~\ref{def:adapted-finite-frame-slices}. Consider a
term \(t_{x_*}(g,h;f)\) of the large standard test in
Corollary~\ref{cor:absolute-stabilizer-period}, with \(gL=L\).
Suppose the normalized finite coordinates of \(x_*g|_L\) are
\(\iota Y\), where \(Y=uA\) as in
Lemma~\ref{lem:finite-unitary-frame-coset}. The finite coordinate
\(h_f\) of every nonzero term lies in \({\rm diag}(u,U_n)K_T\).
Changing \(h_f\) on the right by any \(k\in K_T\) leaves the
term unchanged. In particular at a block element
\(h_f={\rm diag}(y,l)\) the only contributing small frames
\(y^{-1}Y\) belong to \({\rm GL}_b(\widehat O_E)\), so that
Theorem~\ref{thm:block-frame-finite-exchange} applies there.
\end{proposition}
\begin{proof}
The initial large test is supported on integral normalized
frame coordinates at every finite place by
Lemma~\ref{lem:finite-witt-frame-partial-fourier}. We check
that this support condition is transported on the whole finite
frame and remains valid for the residual theta sum. All
modulations in the annihilator of the frame lattice fix the
initial test, since they multiply it by a character equal to
one on this lattice. The inverse image of this annihilator law
on the transformed test is obtained by composing its frame
variable with \(x\mapsto h^{-1}xg|_L\), by the intertwining
and pure momentum calculation in
Lemma~\ref{lem:frame-slice-support-transport}. A finite frame
outside the lattice is detected by one such modulation supported
at a single place. Indeed the dual lattice has that original
lattice as annihilator by the character argument in
Proposition~\ref{prop:finite-polarized-vacuum}. The resulting
nontrivial scalar multiplies every residual value of the slice,
so they all vanish. Thus \(h_f^{-1}\iota Y\) is integral
whenever the term is nonzero. The coset description follows
from Lemma~\ref{lem:finite-unitary-frame-coset}.
Right invariance holds before taking the frame slice by the
representation law and Lemma~\ref{lem:standard-determinant-fixed-target}.
At a block element the same integrality condition reads
\(y^{-1}Y\) integral. The last assertion of
Lemma~\ref{lem:finite-unitary-frame-coset} now proves the
hypothesis on \(X=y^{-1}Y\) in
Theorem~\ref{thm:block-frame-finite-exchange}.
\end{proof}

\section{Integrating the coefficient in two stages}
\begin{definition}\label{def:transverse-gram-chart}
Continue with \(T=T_b\perp T_n\) and the finite groups of
Lemma~\ref{lem:finite-unitary-frame-coset}. Use the rational
block decomposition, set \(G=U(T)(\mathbb A_F)\), and let
\(H=U(T_b)(\mathbb A_F)\times U(T_n)(\mathbb A_F)\) act
block diagonally. Write \(H_K=H_f\cap K_T\) and
\(\Gamma_n={\rm diag}(1,U(T_n)(F))\), with rational points here
meaning isometries over \(E\) defined over \(F\). Give the
groups product Haar measures, using the densities of
Definition~\ref{def:unitary-density} at infinity and any fixed
positive Haar normalizations at the finite places. On a discrete
subgroup use counting measure.

In the real coordinates of Lemma~\ref{lem:explicit-smooth-hermitian-plane-section}
let \(P_0\) be projection on the first block. For a tuple \(Z\)
of complex \(n\)-by-\(b\) matrices let \(P_Z\) be orthogonal
projection for \(J_T\) onto the graph of \(Z\), when that
graph is nondegenerate. Choose an open neighborhood
\(\mathcal U\) of zero on which these graphs are nondegenerate
and \(\|P_Z-P_0\|<1\) placewise. Take the support of the real
plane bump \(\beta\) of Theorem~\ref{thm:matched-real-gram-slices}
inside this chart, with its graph coordinates compactly
contained in \(\mathcal U\). Write \(k_Z=k(P_Z)\) and put
\[
  \mathcal V_\infty=\{s_\infty:
       P_{s_\infty^{-1}\iota}=P_Z \text{ for } Z\in\mathcal U\},
     \qquad \mathcal V=\mathcal V_\infty\times H_f K_T .
\]
Use ordinary coordinate Lebesgue measure \(\dd Z\) on the
product of matrix spaces.
\end{definition}

\begin{lemma}\label{lem:transverse-block-product-measure}
There is a positive smooth density factor \(J(Z)\), determined
by the real Haar measures in Definition~\ref{def:transverse-gram-chart},
such that for each nonnegative measurable \(A\) supported on
\(\mathcal V\),
\[
 \int_G A(s)\dd s
 =c_f\int_{\mathcal U}J(Z)
    \int_{K_T}\int_H A\big(h(k_Z^{-1},k)\big)
                          \dd h\,\dd k\,\dd Z,
 \qquad
        c_f=\frac{{\rm vol}_{G_f}(K_T)}{{\rm vol}_{H_f}(H_K)}.
\]
Here \(\dd k\) is probability on \(K_T\) and \((k_Z^{-1},k)\)
has the indicated infinite and finite parts. If \(A\) is left
\(\Gamma_n\)-invariant, the same formula holds with the left
group \(G\) and the inner group \(H\) both replaced by their
left quotients by \(\Gamma_n\). Each version holds for
absolutely integrable complex functions.
\end{lemma}
\begin{proof}
The map \((h_\infty,Z)\mapsto h_\infty k_Z^{-1}\) is a smooth
bijection from \(H_\infty\times\mathcal U\) onto
\(\mathcal V_\infty\) with smooth inverse. Indeed if
\(s_\infty^{-1}\iota=\binom{V_1}{V_2}\), its plane is a graph
exactly when \(V_1\) is invertible. Then \(Z=V_2V_1^{-1}\),
and \(s_\infty k_Z\) preserves the first plane by
Lemma~\ref{lem:explicit-smooth-hermitian-plane-section}.
It preserves the perpendicular as well and belongs to
\(H_\infty\). This gives the inverse and proves smoothness
in both directions. Pull back the real group density in this
chart and divide by the product of the density on \(H_\infty\)
and \(\dd Z\). The result is smooth and positive since the
coordinate differentials are inverse isomorphisms. Left
invariance under \(H_\infty\) on both sides, by
Lemma~\ref{lem:unitary-cayley-density}, makes the ratio
independent of the \(H_\infty\) coordinate. This is \(J(Z)\).

The finite-adelic region \(H_fK_T\) is the union of
the disjoint translates \(tK_T\), where \(t\) ranges over
representatives for \(H_f/H_K\). This union can be countably
infinite. The translates are open and the local fields have
countable bases, as does their restricted product since only
finitely many integrality conditions may be changed at a
time. Each translate has measure \({\rm vol}_{G_f}(K_T)\);
in the integral on \(H_f\) the translate \(tH_K\) has measure
\({\rm vol}_{H_f}(H_K)\). Inserting the normalized average on
\(K_T\) after the \(H_f\) coordinate gives exactly the
claimed ratio, since translating that average on the left
by any element of \(H_K\) preserves it. Applying
Lemma~\ref{lem:product-integration} gives the first formula.

For the quotient, take Borel fundamental domains as in
Lemma~\ref{lem:quotient-measure}. Its construction applies
on both adelic groups: a compact set of translating rational
matrices has common finite denominators and bounded real entries,
hence is finite by Lemma~\ref{lem:trace-finite}. Cover by a
countable family of neighborhoods disjoint from their nonidentity
rational translates and keep in each only the orbits which did
not meet an earlier neighborhood. The resulting domains meet
each orbit once. Multiply \(A\) upstairs by the indicator
of the domain on \(G\). In the formula already proved,
decompose the \(H\)-integral over translates of the domain
on \(H\). At fixed \(Z,k,h\), the sum of the domain indicators
at \(\delta h(k_Z^{-1},k)\), \(\delta\in\Gamma_n\), is one.
The value of \(A\) is invariant under these translations.
Nonnegative summation gives the quotient formula, and applying
it first to the modulus gives the complex case.
\end{proof}

\begin{theorem}\label{thm:two-stage-nonsingular-gram-coefficient}
In Proposition~\ref{prop:actual-coset-finite-coefficient},
let \(T_n=W_0\), \(x_*=\iota\) in rational coordinates, and
take the real test \(F\) of
Theorem~\ref{thm:matched-real-gram-slices} with the plane bump
of Definition~\ref{def:transverse-gram-chart}. Take \(g\)
at which the specified finite frame condition holds and \(gL=L\).
Let \(A_g(s)=t_{x_*}(g,s;F\otimes F^+)\).
Write \([U(T_n)]=U(T_n)(F)\backslash U(T_n)(\mathbb A_F)\),
and let \(B_*=x_*^*J_Tx_*\). The \(B_*\)-coefficient of the
large period at \(g\) equals
\[
 c_f\int_{\mathcal U} J(Z)\!
    \int_{U(T_b)(\mathbb A_F)}\int_{[U(T_n)]}
      A_g\big({\rm diag}(y,l)(k_Z^{-1},1)\big)
                       \dd l\,\dd y\,\dd Z .
\]
The integral of the modulus in this expression is finite.
The transverse constant
\[
              R=c_f\int_{\mathcal U}J(Z)\beta(P_Z)\dd Z
\]
is strictly positive, finite, and independent of \(S,\chi,\nu\)
when the real bump, spaces, lattices and measures are fixed.
\end{theorem}
\begin{proof}
By Corollary~\ref{cor:absolute-stabilizer-period} the coefficient
is the integral of \(A_g\) over \(\Gamma_n\backslash G\) and
is absolutely integrable. Its finite support lies in \(H_fK_T\)
by Proposition~\ref{prop:actual-coset-finite-coefficient}.
At infinity the transported frame coordinate for a nonzero
term is \(s_\infty^{-1}\iota g|_{L,\infty}\), in the real
frames and scales. Its column plane is the plane of
\(s_\infty^{-1}\iota\), since the last factor is invertible
on \(L\) and the real scale is scalar. The mixed-test support
and its transport in Lemma~\ref{lem:frame-slice-support-transport}
therefore put \(s_\infty\) in \(\mathcal V_\infty\).
Apply the quotient formula in
Lemma~\ref{lem:transverse-block-product-measure}. Right
\(K_T\)-invariance in Proposition~\ref{prop:actual-coset-finite-coefficient}
removes the last finite compact average. The quotient of \(H\)
is the direct product displayed since \(\Gamma_n\) acts only
in the second factor. This proves the formula and the absolute
integrability assertion. The continuity and positivity of
\(J\), the compact support of \(Z\mapsto\beta(P_Z)\), its
nonnegativity and its positivity in a neighborhood of zero
give \(0<R<\infty\). Its defining factors involve only the
data listed in the statement.
\end{proof}

\section{Levi action and the characters of the two periods}
\begin{definition}\label{def:normalized-complementary-levi}
In Definition~\ref{def:adapted-finite-frame-slices} let
\(n=2j\), \(N=2b+n\) and \(d=b+n\), and identify \(T_b\)
with the opposite target of Theorem~\ref{thm:opposite-branch-small-period}.
Write \(\upsilon_T=\upsilon|_{U(1)}\) with \(\upsilon\)
from Theorem~\ref{thm:actual-global-splitting-character}, and
put
\[
  \gamma_0=\nu^{-1}\upsilon_T^{j},\qquad
  \eta_0=\nu\upsilon_T^{j},\qquad \gamma=\gamma_0\upsilon_T^b,
       \qquad \lambda(t)=\nu(t/\bar t).
\]
Write \(\delta(s)=\det_E s\) for an isometry, and
\(|t|_{\mathbb A}\) for the idele modulus of \(E\), using
\(\prod_\tau|t_\tau|^2\) at infinity. For an adelic
\(q\in{\rm GL}(L)\), \(g_0\in U(W_0)(\mathbb A_F)\),
let \(m(q,g_0)\) act on \(L,W_0,L'\) by
\(q,g_0,(q^*)^{-1}\), respectively. For the large and the
flipped small period multiply both functions at \(g\)
by \(\gamma(\delta(g))^{-1}\). Denote the resulting kernels
by the corresponding normalized kernels, using for the small
one its directly transformed finite datum \(F^-\).
\end{definition}

\begin{lemma}\label{lem:levi-operator-every-frame}
In the mixed position model for \({\rm Hom}(W,T')\) write
\(f[x]\) for evaluation in the full adelic frame variable
at \(x\), taking values in the residual oscillator on
\({\rm Hom}(W_0,T')\). Put \(t=\dim T'\), \(D=\det q\) and
let \(h\in U(T')(\mathbb A_F)\). For unitary characters \(a,B\)
of the norm-one torus use \(a\upsilon_T^b,B\) on the full pair.
Then
\[
\begin{split}
   \big(\Omega_{a\upsilon_T^b,B}(m(q,g_0),h)f\big)[x]
       &=p_{B,t}(D)\,\Omega_{a,B}(g_0,h)\big(f[h^{-1}xq]\big),\\
 p_{B,t}(D)&=\upsilon(D)^t\,B(D/\bar D)^{-1}|D|_{\mathbb A}^{t/2}.
\end{split}
\]
The symbols \(x,h^{-1}xq\) in this identity are physical frame
coordinates. Expressing them in the same normalized frames
on the two sides gives the identical formula.
\end{lemma}
\begin{proof}
The map on the hyperbolic position is \(X\mapsto hXq^{-1}\).
Its \(E\)-determinant is \(\delta(h)^bD^{-t}\), since left
multiplication acts separately on the columns and right
multiplication separately on the rows. The hyperbolic and
residual actions are independent factors by the direct sum
calculation of Theorem~\ref{thm:residual-levi-zero-section}.
The whole hyperbolic operator, at any position, is
substitution by the inverse map times the inverse splitting
character of its determinant and the inverse square root of
its modulus. This follows from
Theorems~\ref{thm:hyperbolic-adelic-parabolic-splitting}
and~\ref{thm:graded-parabolic-schroedinger}, since an operator
preserving both isotropic summands has zero shear.
As \(\delta(h)\overline{\delta(h)}=1\) and conjugation
preserves the idele modulus, \(|\delta(h)|_{\mathbb A}=1\).
Thus the two factors are
\(\upsilon_T(\delta(h))^{-b}\upsilon(D)^t\) and
\(|D|_{\mathbb A}^{t/2}\). The target determinant twist cancels
the first power of \(\upsilon_T(\delta(h))\), leaving \(a(\delta(h))\).
The source determinant is \((D/\bar D)\delta(g_0)\), giving the
other displayed scalar and \(B(\delta(g_0))^{-1}\).
Together the remaining action and characters give
\(\Omega_{a,B}(g_0,h)\) by Definition~\ref{def:two-order-theta-kernel}.
This calculation applies in a dense family of pure tests and also
pointwise for arbitrary Schwartz tests: in a product polarization
the block substitutions act on separate variables and the residual
Fourier integrals commute with evaluation in the frame variables,
as shown in Theorem~\ref{thm:residual-levi-zero-section}.
\end{proof}

\begin{proposition}\label{prop:complementary-levi-factor-quotient}
After the common normalization of Definition~\ref{def:normalized-complementary-levi},
the full characters on the large pair are \((\gamma,\gamma)\),
and on the small pair they are
\((\eta_0\upsilon_T^b,\upsilon_T^b)\). The ratio of their
Levi scalars in Lemma~\ref{lem:levi-operator-every-frame} is
\[
       p_{\gamma,d}(D)/p_{\upsilon_T^b,b}(D)
                    =\lambda(D)|D|_{\mathbb A}^{n/2}.
\]
On the residual pair at a block element \({\rm diag}(y,l)\)
of the target, the normalized large action on a tensor \(u\otimes v\)
is
\[
 \nu(\delta(y))^{-2}\,
    \Omega_{\eta_0,\upsilon_T^b}(g_0,y)u
                   \ \otimes\ \Omega_{\gamma_0,\gamma_0}(g_0,l)v .
\]
\end{proposition}
\begin{proof}
The original large characters are \((\gamma,1)\) by
Definition~\ref{def:complementary-period-characters}, and the
flipped small ones are
\((\gamma^{-1}\upsilon_T^N,\gamma^{-1}\upsilon_T^b)\) by
Theorem~\ref{thm:opposite-branch-small-period}. Multiplication
by the common scalar on the ambient group multiplies the
denominator character by \(\gamma\).
Since \(N=2(b+j)\), \(\gamma^{-1}\upsilon_T^N=\eta_0\upsilon_T^b\).
For the ratio of Levi scalars the unit character before
simplification is
\(\upsilon(D)^n\gamma_0(D/\bar D)^{-1}\).
The identity
\(\upsilon_T(D/\bar D)=\upsilon(D)^2\) follows from the
conjugation rule of Theorem~\ref{thm:actual-global-splitting-character}.
Use \(n=2j\) and the definition of \(\gamma_0\) to get \(\lambda(D)\).
The ratio of moduli is \(|D|_{\mathbb A}^{(d-b)/2}\).

The residual target character is \(\gamma_0\) on both
blocks. The original residual source character is
\(\gamma=\upsilon_T^b\gamma_0\) and the untwisted action on the
orthogonal Hom sum is the tensor product, by the product
calculation in Theorem~\ref{thm:residual-levi-zero-section}.
Putting \(\upsilon_T^b\) on the first source denominator
and \(\gamma_0\) on the second accounts for this character
at \(g_0\). Changing the first target numerator to \(\eta_0\)
multiplies it by \(\eta_0(\delta(y))/\gamma_0(\delta(y))
=\nu(\delta(y))^2\), giving the displayed correcting factor.
\end{proof}

\section{The two period coefficients on the Levi subgroup}
\begin{definition}\label{def:actual-equal-rank-integral}
With the hypotheses of Theorem~\ref{thm:two-stage-nonsingular-gram-coefficient}
and Definition~\ref{def:normalized-complementary-levi},
assume \(S\ne\varnothing\) and that its marked branches are
distinct from their conjugates, as in
Definition~\ref{def:ramified-line-mellin}.
Let \(v_1\) be the real residual vector on \({\rm Hom}(W_0,W_0)\)
occurring in Theorem~\ref{thm:matched-real-gram-slices}. Define
\[
  J_\nu(g_0)=\int_{[U(W_0)]}
       \theta\big(\Omega_{\gamma_0,\gamma_0}(g_0,l)
                              (v_1\otimes F_{\rm eq})\big)\dd l ,
\]
where \(F_{\rm eq}\) is the finite determinant vector of
Theorem~\ref{thm:block-frame-finite-exchange}. This integral is
absolutely convergent. Indeed the vector is a finite linear
combination of named residue points, all surjective at the
chosen branch. The large-side estimate and the height integral
in Theorem~\ref{thm:complementary-period-convergence} apply
with source and target \(W_0\). The unitary twists leave its
estimate unchanged.

Write \(A_{\nu}^{+}(m)\) and \(A_{\nu}^{-}(m)\) for the
\(B_*\)-coefficients of the two normalized periods in
Definition~\ref{def:normalized-complementary-levi}, using
\(F\otimes F^+\) and \(F'\otimes F^-\) respectively, with
the real data of Theorem~\ref{thm:matched-real-gram-slices}.
The small period uses the directly transformed model, so the
phase \(c_0\) of Theorem~\ref{thm:opposite-branch-small-period}
does not enter these symbols. Identify \(L\) with \(T_b\)
by the rational frame \(x_*\). Write \(Y(q)\) for \(x_*q\)
in the chosen normalized frame coordinates and \(a=(\det Y(1))^{-1}\).
Then \(\det q=a\det Y(q)\): changing from rational to normalized
frames multiplies the frame matrix on the left and right by
fixed invertible adelic matrices.
\end{definition}

\begin{lemma}\label{lem:fixed-finite-frame-vacuum-factor}
Use rational product polarizations on the hyperbolic and residual
Hom summands in Lemma~\ref{lem:levi-operator-every-frame}, and
on the two residual summands for \(T=T_b\perp W_0\). In any fixed
physical position measures the normalized finite slice formulas
of Theorem~\ref{thm:block-frame-finite-exchange} acquire one
nonzero constant \(c_t\) for a target of rank \(t\). It is fixed
by the space, frames, measures and choice of vacuum phase,
independently of the set of branches and of the residue function.
With positive frame vacuum and unit residual vacuum one has
\(c_t={\rm vol}(\Lambda_t)^{-1/2}\), where \(\Lambda_t\) is the
physical hyperbolic frame lattice.
\end{lemma}
\begin{proof}
The full finite vacuum in the orthogonal sum splits as the
hyperbolic vacuum times the residual vacuum, since the factors
are independently fixed by their lattice translations. Both
fixed spaces are lines by Proposition~\ref{prop:finite-polarized-vacuum}.
In the hyperbolic frame model this line is spanned by the
lattice indicator. Its squared norm is the volume of the lattice.
These assertions hold before specifying any marked branches.
Normalize the tensor factors to be unit vectors and take the
positive frame indicator. This gives the stated \(c_t\). With
any other fixed phase multiply it by that phase. A named residue
function acts by a finite linear combination of the physical
translations in Proposition~\ref{prop:finite-oscillator-residue-name};
each acts in both normalizations by the same substitutions and
characters. Thus their common scalar persists for the full
linear combination. In particular the transformation of the
unused \(L'\)-block in
Lemma~\ref{lem:finite-witt-frame-partial-fourier} introduces
no further scalar varying with the marked set, because its
intertwiner takes the unit vacuum to the unit vacuum. The
residual vacuum on the direct sum is the product of the two
residual vacua with the compatible product convention.
\end{proof}

\begin{theorem}\label{thm:actual-levi-complementary-coefficients}
Suppose \(Y(q)_f\in U_b\,{\rm GL}_b(\widehat O_E)\) as in
Lemma~\ref{lem:finite-unitary-frame-coset}. For every \(g_0\),
with \(m=m(q,g_0)\) of Definition~\ref{def:normalized-complementary-levi},
\[
       A_\nu^+(m)=K_\nu\,J_\nu(g_0)\,A_\nu^-(m),\qquad
       K_\nu=R\,(c_d/c_b)\,u_S\,\lambda(a)|a|_{\mathbb A}^{n/2}.
\]
Here \(R>0\) is as in Theorem~\ref{thm:two-stage-nonsingular-gram-coefficient}.
In particular \(|K_\nu|\) is positive and independent of the
chosen twist and its marked set for fixed spaces, real bump,
lattices and measure choices.
\end{theorem}
\begin{proof}
Use the integral over \(Z,y,l\) of
Theorem~\ref{thm:two-stage-nonsingular-gram-coefficient} and
put \(B=y^{-1}Y(q)\) in target coordinates. At the finite
places a nonintegral \(B_f\) annihilates both frame slices by
Lemma~\ref{lem:finite-witt-frame-partial-fourier}. Otherwise
\(B_f\) is unimodular by Lemma~\ref{lem:finite-unitary-frame-coset}.
At an integral \(B_f\), the finite slice ratio of the large to
the tensor of the small and equal-rank vectors means the scalar
in the equality of vectors in Theorem~\ref{thm:block-frame-finite-exchange}.
By Lemma~\ref{lem:fixed-finite-frame-vacuum-factor} this scalar is
\((c_d/c_b)u_S\lambda_f(\det B_f)^{-1}\) in physical coordinates.
Also \(|\det B_f|_f=1\).

At infinity \((k_Z^{-1})^{-1}\iota B_\infty=k_Z\iota B_\infty\).
In the construction of Theorem~\ref{thm:matched-real-gram-slices}
the plane of this frame is \(P_Z\) and its coordinate \(c\)
is \(B_\infty\). Apply the real target action of \(k_Z^{-1}\)
to the explicit \(F\) in its proof, taking \(\pi\) to be the
residual target action of \(\gamma_0\). The factor \(\pi(k_Z)\)
in \(F\) cancels. This gives the scalar
\(\beta(P_Z)\lambda_\infty(\det B_\infty)^{-1}
|\det B_\infty|_\infty^{-n/2}\) times \(F'(B_\infty)\otimes v_1\).
Here \(|\ |_\infty\) is the product of the squared complex
absolute values. The infinite character of \(\gamma_0\) depends
only on the fixed weight by Definition~\ref{def:torus-weight-and-units},
so this choice of the real data works for all these twists.

Apply Lemma~\ref{lem:levi-operator-every-frame} and
Proposition~\ref{prop:complementary-levi-factor-quotient} to
the remaining block \({\rm diag}(y,l)\). In addition to
\((c_d/c_b)u_S\beta(P_Z)\) the total scalar on comparing the
large slice to the small slice and the equal-rank vector is
\[
  \lambda(\det q)|\det q|_{\mathbb A}^{n/2}\,
  \lambda(\det B)^{-1}|\det B|_{\mathbb A}^{-n/2}
                                  \nu(\delta(y))^{-2}.
\]
Since \(\delta(y)\overline{\delta(y)}=1\),
\(\lambda(\delta(y))=\nu(\delta(y))^2\) and
\(|\delta(y)|_{\mathbb A}=1\). As
\(\det B=\delta(y)^{-1}\det Y(q)\), the expression is
\(\lambda(a)|a|_{\mathbb A}^{n/2}\), even at zeros of the
remaining theta functions. These computations establish a
vector equality before taking theta values.

In the product rational position polarization the theta sum of
a tensor is the product of the theta sums. Indeed its index set
is the direct product of the two rational position spaces and
the sums converge absolutely by
Theorem~\ref{thm:rational-schroedinger-zero-section}.
The two sums now depend on \(y\) and \(l\) respectively.
The integral of the first over \(U(T_b)(\mathbb A_F)\) is the
small coefficient with trivial frame stabilizer by
Corollary~\ref{cor:small-period-trivial-stabilizer}, in its
directly transformed model. The second integrates to
\(J_\nu(g_0)\) by Definition~\ref{def:actual-equal-rank-integral}.
Both are absolutely integrable by those results. The common
normalization commutes with taking the coefficient because
the central shears in Definition~\ref{def:gram-indexed-mixed-period}
have determinant one. The last transverse integration gives
\(R\), proving the formula. The constants \(c_b,c_d\) and the
idele \(a\) are fixed, \(R>0\), and both \(\lambda(a)\) and
\(u_S\) have modulus one.
\end{proof}

\section{The unipotent radical in the coefficient comparison}
\begin{definition}\label{def:full-frame-radical}
In the ordered dual frames \(L,W_0,L'\) use the Gram matrix
with identity off-diagonal blocks between \(L,L'\), and block
\(J_0\) on \(W_0\). For an adelic \(b\)-by-\(n\) matrix \(C\)
and an adelic \(b\)-by-\(b\) matrix \(S_0\) with \(S_0^*=-S_0\),
put
\[
 u(C,S_0)=\begin{pmatrix}
 I&C&S_0-\tfrac12 CJ_0^{-1}C^*\\
 0&I&-J_0^{-1}C^*\\ 0&0&I
 \end{pmatrix},\qquad U=\{u(C,S_0)\}.
\]
These are precisely the isometries preserving \(L\), acting as
identity on \(L\), \(L^\perp/L\) and \(W/L^\perp\). Indeed such a
matrix is upper triangular with identity diagonal blocks. In the
Gram equation the middle-right block must equal \(-J_0^{-1}C^*\),
and the upper-right block plus its adjoint must equal
\(-CJ_0^{-1}C^*\). Conversely substitution gives the Gram
equation. Every adelic isometry preserving \(L\) is a unique
product \(u\,m(q,g_0)\) with \(u\in U\), by taking its diagonal
blocks and using the same equation.
\end{definition}

\begin{lemma}\label{lem:radical-residual-frame-action}
For a target \(T'\) let the pair action of
\(u\in U\) on \({\rm Hom}(W,T')\) be the one induced by
precomposition with \(u^{-1}\). In the frame polarization of
Lemma~\ref{lem:levi-operator-every-frame} it induces, at each
frame \(x\), an operator \(D_{T'}(u,x)\) on residual Schwartz
vectors satisfying
\[
        (\Omega(u,1)f)[x]=D_{T'}(u,x)f[x],\qquad
                              D_{T'}(u,0)=1.
\]
For compatible product polarizations on \(T'=T_b\perp T_n\),
\[
 D_{T'}(u,(x_b,x_n))=
          D_{T_b}(u,x_b)\otimes D_{T_n}(u,x_n).
\]
These identities hold for all the pair characters in
Definition~\ref{def:normalized-complementary-levi}.
\end{lemma}
\begin{proof}
Write a Hom matrix in three column blocks \((x,M,V)\) according
to \(L,W_0,L'\). Under precomposition by \(u^{-1}\), its first
block remains \(x\), its middle block changes by \(-x C\), and
the third changes by a linear function of \(x,M\). Thus if
\(x=0\) and \(M\) lies in the chosen residual momentum, all
position coordinates remain zero. The whole momentum space
is preserved. This also holds for \(d=m(2I,1)\). On the group
generated by \(U,d\) use the operators \(\mathcal W\) of
Theorem~\ref{thm:graded-parabolic-schroedinger}. Both these
operators and the physical \(\Omega\) represent this group and
intertwine the same Heisenberg translations by
Theorem~\ref{thm:actual-adelic-rational-theta-invariance}.
Their ratio is scalar by Lemma~\ref{lem:global-heisenberg-scalar}
and is therefore a character.

Conjugation by \(d\) carries \((C,S_0)\) to \((2C,4S_0)\).
On the central subgroup \(C=0\), invariance of the character
under this conjugation makes it trivial on every \(3S_0\),
hence on all \(S_0\). Modulo that central subgroup the radical
is additive in \(C\), by block multiplication in
Definition~\ref{def:full-frame-radical}. The descended character
has equal values at \(2C\) and \(C\), so is one. This proves
\(\Omega(u,1)=\mathcal W_u\). Pair characters do not change
the conclusion since \(\delta(u)=1\).

In the explicit formula of Theorem~\ref{thm:graded-parabolic-schroedinger}
the quotient substitution fixes the frame coordinate, so the
answer at a fixed \(x\) depends only on \(f[x]\). At \(x=0\)
the residual position is unchanged. Its change in momentum
lies entirely in the third column block and therefore pairs
trivially with the residual position, by the Gram blocks.
The quadratic phase is consequently one there. The quotient
substitution has determinant one since it fixes \(x\) and adds
a linear function of \(x\) to the residual position. This proves
the zero-frame assertion. For an orthogonal sum of targets, the
linear maps are direct sums and the two summands pair trivially.
Thus the quotient substitution acts on the two factors separately
and the exponent of the quadratic phase is the sum of the two
exponents. The displayed formula for \(\mathcal W_u\) gives the
tensor identity.
\end{proof}

\begin{theorem}\label{thm:parabolic-unimodular-coefficient-product}
In Theorem~\ref{thm:actual-levi-complementary-coefficients},
for every \(u\in U\) and \(m=m(q,g_0)\) in its range,
\[
           A_\nu^+(um)=K_\nu J_\nu(g_0) A_\nu^-(um).
\]
Here each side uses the coefficient of the same index \(B_*\)
with the same left central integration as in
Definition~\ref{def:gram-indexed-mixed-period}.
\end{theorem}
\begin{proof}
In the two-stage formula of
Theorem~\ref{thm:two-stage-nonsingular-gram-coefficient} the
finite support condition is unchanged because \((um)|_L=q\).
The same statement holds for the real plane chart.
In the proof of Theorem~\ref{thm:actual-levi-complementary-coefficients}
the equality at \(m\) holds on residual vectors at the frames
\((x_*,0)\) and \(x_*\) in the two targets, writing \(x_*\)
here for the identification with the first summand. It holds
at each of the indicated \(Z,y,l\) before taking any residual
theta sum. Apply \(\Omega(u,1)\) on the left of each transformed
kernel and evaluate at the same frames. By
Lemma~\ref{lem:radical-residual-frame-action} the large residual
action is \(D_{T_b}(u,x_*)\otimes 1\), and the small residual
action is \(D_{T_b}(u,x_*)\). Hence the first vector is still
the same scalar times the tensor of the new small vector and
the previous equal-rank vector. Taking rational theta sums as
in Theorem~\ref{thm:actual-levi-complementary-coefficients}
and integrating gives the asserted identity. Its absolute
convergence at \(um\) follows from
Corollary~\ref{cor:absolute-stabilizer-period} on the large side,
Corollary~\ref{cor:small-period-trivial-stabilizer} on the small
side, and Definition~\ref{def:actual-equal-rank-integral} on
the residual factor.
\end{proof}

\section{The equal-rank period at a marked split place}
\begin{lemma}\label{lem:relative-haar-invertible-matrices}
Let \(K\) be a nonarchimedean local field with valuation ring
\(O\) and normalized absolute value. Put \(G={\rm GL}_n(K)\)
for \(n>0\). On locally constant compactly supported functions
on \(G\) write \(L_hf(X)=f(h^{-1}X)\), \(R_sf(X)=f(Xs)\).
Give \(G\) the measure \(\dd^\times X\) of mass one on
\({\rm GL}_n(O)\) obtained by normalizing
\(|\det X|^{-n}\dd X\), where \(\dd X\) is additive measure
on the matrix entries. For a locally constant character
\(\omega:G\to\C^\times\), any linear functional \(\ell\)
with \(\ell(\omega(h)L_h f)=\ell(f)\) satisfies
\[
      \ell(f)=c\int_G f(X)\omega(X)^{-1}\dd^\times X,\qquad
      \ell\big(\omega(s)^{-1}R_sf\big)=\ell(f).
\]
Here \(c=\ell(\omega\,{\bf1}_{{\rm GL}_n(O)})\). Continuity of
\(\ell\) is unnecessary.
\end{lemma}
\begin{proof}
Left and right multiplication of a matrix by \(h\) each change
additive measure by \(|\det h|^n\), using the columns and the
rows, respectively, and the determinant change of variables in
Theorem~\ref{thm:split-finite-schroedinger}. Thus \(\dd^\times X\)
is invariant on both sides. The functional
\(T(f)=\ell(\omega f)\) is left invariant because
\(\omega L_hf=\omega(h)L_h(\omega f)\).
Let \(K_m=1+\pi^m M_n(O)\) for \(m\ge1\) and a uniformizer
\(\pi\). It is a subgroup: the inverse series for \(1+\pi^m A\)
converges entrywise by completeness. Its finite index in
\(K_0={\rm GL}_n(O)\) follows by reduction modulo \(\pi^m\).
Partitioning \(K_0\) into its \(K_m\) cosets gives
\(T({\bf1}_{K_m})=T({\bf1}_{K_0})/[K_0:K_m]\).
Every test is constant on the cosets of some \(K_m\): cover
its compact open support by finitely many constancy cosets
and refine to the largest depth. Outside their union it vanishes.
Translation invariance proves \(T(f)=c\int f\,\dd^\times X\)
on all such tests. Substituting \(\omega^{-1}f\) gives the
formula for \(\ell\), and right invariance of the measure gives
the last assertion upon substituting \(Y=Xs\).
\end{proof}

\begin{theorem}\label{thm:marked-equal-rank-period-invariance}
Use Definition~\ref{def:actual-equal-rank-integral} and fix
a marked branch \(w\) to which its rank estimate applies.
In split integral frames identify \(G_w={\rm GL}_n(E_w)\)
with the entire unitary factor over the place below \(w\)
by taking the conjugate dual inverse on \(\bar w\).
Write \(\widetilde h\) for the corresponding group element
supported on this pair. Replacing the determinant test at \(w\)
by \(f\in C_c^\infty(G_w)\), extended by zero on singular
matrices, defines an absolutely convergent equal-rank period
\(\ell_{g_0}(f)\). If \(e_w=\beta(\mathfrak p_w)\) is the ideal
character of Theorem~\ref{thm:split-cut-weil-multiplier}, set
\[
 \omega(h)=\gamma_0(\delta(\widetilde h))\,
        e_w^{-n v_w(\det h)}|\det h|_w^{-n/2}.
\]
With the measure of Lemma~\ref{lem:relative-haar-invertible-matrices},
\[
       \ell_{g_0}(f)=J_\nu(g_0)
                \int_{G_w} f(X)\omega(X)^{-1}\dd^\times X,
          \qquad J_\nu(g_0\widetilde s)=J_\nu(g_0)
                                  \quad(s\in G_w).
\]
At every unmarked finite place \(J_\nu\) is right invariant
under the integral unitary group there.
\end{theorem}
\begin{proof}
In the split Hom coordinates \(X\) the target and source
act by \(X\mapsto h Xs^{-1}\). On the chosen branch the
determinant of the Hom action is
\((\det h)^n(\det s)^{-n}\). The normalized substitution of
Theorem~\ref{thm:split-finite-schroedinger}, the splitting
phase of Theorem~\ref{thm:split-cut-weil-multiplier}, and the
two identical determinant twists of
Definition~\ref{def:actual-equal-rank-integral} therefore give
on functions the action
\(\omega(h)\omega(s)^{-1}L_h R_s\).
The character \(\omega\) is locally constant because valuation
is locally constant and on integral units the torus character
factors through the finite residue group by
Definition~\ref{def:complementary-period-characters}.

Any such \(f\) is a finite sum of indicators of translates
of \(K_m\) in the proof of Lemma~\ref{lem:relative-haar-invertible-matrices}.
Extended by zero on the full matrix space these indicators are
locally constant and compactly supported, since each coset is both
compact and open there. The test \({\bf1}_{K_m}\) is a residue
test at depth \(m\) with invertible reduction. Hence the target
decay proof of Theorem~\ref{thm:complementary-period-convergence}
applies: the detecting displacement in
Lemma~\ref{lem:two-sided-rank-detection} uses the first residue
layer and the test is a finite sum of named points at this
depth. For a left translate \({\bf1}_{hK_m}\), the target action
already computed amounts to right translation of the integration
variable by \(\widetilde h\) and multiplication by a fixed
nonzero scalar. This translation preserves the quotient measure,
since it preserves Haar measure upstairs by
Lemma~\ref{lem:relative-haar-invertible-matrices} and transports
a rational fundamental domain to another one. We thus obtain
absolute convergence for every \(f\), even when this is the only
marked pair. The same change of variables gives
\(\ell_{g_0}(\omega(h)L_hf)=\ell_{g_0}(f)\).

On \(K_0\) the scalar \(\omega(h)\) equals
\(\chi_w(\det h)^{-1}\) by integral triviality of \(\upsilon_T\)
in Theorem~\ref{thm:actual-global-splitting-character} and the
unit restriction of \(\nu\). The defining determinant test of
\(J_\nu\) at \(w\) is therefore \(\omega{\bf1}_{K_0}\).
Lemma~\ref{lem:relative-haar-invertible-matrices} gives the
first equality with exactly the stated normalization. Translating
\(g_0\) on the right by \(\widetilde s\) acts on this test by
\(\omega(s)^{-1}R_s\). The second equality follows from the
same lemma. At an unmarked place the integral group fixes the
vacuum by Theorem~\ref{thm:all-finite-weil-splitting}, and
all the finite determinant twists there are trivial on integral
units. This gives the last invariance.
\end{proof}

\section{Removing one place in additive approximation}
\begin{lemma}\label{lem:additive-approximation-away-from-place}
Let \(F\) be a number field and \(v\) a finite place. The diagonal
copy of \(F\) is dense in its additive adeles with the component
at \(v\) omitted. The same holds for any finite-dimensional
\(F\)-vector space in rational coordinates.
\end{lemma}
\begin{proof}
Write \(\mathfrak p\) for the prime at \(v\). Some positive power
\(\mathfrak p^h\) is principal by Lemma~\ref{lem:bounded-ideal-classes};
choose its generator \(a\). Its absolute norm is
\((N\mathfrak p)^h>1\) by Proposition~\ref{prop:general-ideal-factorization}.
There is an element \(\varpi\) generating a positive power of
\(\mathfrak p\) with \(|\sigma\varpi|>1\) at every infinite
place. Indeed put \(d=[F:\Q]\),
\(A=d^{-1}\log|N_{F/\Q}a|>0\). The vector
\((A-\log|\sigma a|)_\sigma\) has weighted sum zero. By
Lemma~\ref{lem:unit-log-domain} the multiples of this vector
can be approximated by unit logarithms with bounded error
in every coordinate, by subtracting a vector in the bounded
fundamental parallelepiped. For a sufficiently large integer
\(m\), choose the corresponding unit \(u\) and set
\(\varpi=a^m u\). Then \(\log|\sigma\varpi|=mA+O(1)>0\)
and \(\varpi\) is a local unit at every finite place other
than \(v\).

A basic approximation problem prescribes balls at finitely
many places outside \(v\), arbitrary open neighborhoods at
infinity and integrality at the remaining finite places other
than \(v\). Choose \(x_0\in F\) satisfying the finite conditions.
For detail, first multiply the finite targets by a positive
integer \(D\) so that all prescribed centers and ball ideals
become integral. Include every other prime dividing \(D\)
among the conditions, requiring integrality after division
by \(D\), and impose any integral condition at \(v\) if
needed. Using residues of the local centers, prescribe the
resulting congruences for \(D x_0\) in \(O_F\). Local residues
modulo powers of a prime ideal lift to \(O_F\): denominators
in its localization are invertible modulo that power, and
completion has the same finite quotients. Powers of distinct
prime ideals are coprime by Proposition~\ref{prop:general-ideal-factorization},
so their congruences can be solved together by taking a
combination of elements equal to one in a single quotient and
zero in the others, using the unit ideal generated by coprime
ideals. This produces \(x_0\).

There is a nonzero integral ideal \(I\subset O_F\) whose addition
to \(x_0\) preserves all the finite conditions. Take the product
of sufficiently high prime powers at the specified places.
Since \(\varpi\) is a unit away from \(v\), every lattice
\(x_0+\varpi^{-k}I\), \(k\ge0\), still satisfies them.
The ideal \(I\) is a full real lattice under the archimedean
embeddings by Lemma~\ref{lem:number-field-lattice}. A bounded
fundamental parallelepiped for it shrinks in diameter to zero
after multiplication by \(\varpi^{-k}\), since each
\(|\sigma\varpi|>1\). Thus for large \(k\) a point of that
translate meets the prescribed real and complex conditions.
Applying the scalar assertion to the coordinates of a basis
gives the vector-space case.
\end{proof}

\begin{theorem}\label{thm:equal-rank-invariant-under-radicals}
Let \(W_0\) of Definition~\ref{def:actual-equal-rank-integral}
admit a rational Witt splitting and let \(U_0,U_0^-\) be its two
opposite radicals as in Definition~\ref{def:full-frame-radical}.
For all \(g\in U(W_0)(\mathbb A_F)\) and
\(u\in U_0(\mathbb A_F)\) or \(U_0^-(\mathbb A_F)\),
\[
                        J_\nu(ug)=J_\nu(g).
\]
In particular the integral of \(J_\nu(ug)\) over a rational
quotient of either radical with probability measure is \(J_\nu(g)\).
\end{theorem}
\begin{proof}
The period is continuous by Theorem~\ref{thm:complementary-period-convergence}
and is left rationally invariant by
Theorem~\ref{thm:actual-adelic-rational-theta-invariance} and
rational triviality of both twists. Choose the finite place \(v\)
of \(F\) below the branch of
Theorem~\ref{thm:marked-equal-rank-period-invariance}. Invariance
on the right under the entire local factor at \(v\) also permits
changing that coordinate independently at any point, since two
points agreeing at all other places differ on the right by an
element of that local factor.

The coordinates \(C,S_0\) in Definition~\ref{def:full-frame-radical}
range over finite-dimensional \(F\)-vector spaces. Skew-hermitian
matrices form such a rational space, with the diagonal entries
in the anti-invariant line and one free \(E\)-entry in each
off-diagonal pair. The formulas from those coordinates to the
group and back are polynomial and preserve integrality outside
finitely many places. Lemma~\ref{lem:additive-approximation-away-from-place}
therefore gives a sequence \(u_i\in U_0(F)\) tending to \(u\)
in every adelic coordinate away from \(v\), in the restricted
product topology. The topology is first countable by the countable
basis argument of Lemma~\ref{lem:transverse-block-product-measure}.
In \(u_i g\) replace the \(v\)-component by \((ug)_v\).
The modified sequence tends adelically to \(ug\) and each
modified point has value \(J_\nu(u_i g)=J_\nu(g)\). Continuity
gives the identity. The identical argument with \(L,L'\)
interchanged applies to \(U_0^-\). The last assertion is the
integral of a constant against a measure of mass one.
\end{proof}

\section{Generation by the two radicals}
\begin{definition}\label{def:local-unitary-line-generators}
Let \(k\) be a field of characteristic zero and \(K/k\) a quadratic
field extension with conjugation, or \(K=k\times k\) with its exchange
involution. Let \(W=K e\oplus V\oplus K f\) have the hyperbolic
pair \(e,f\) with \(\langle e,f\rangle=1\), perpendicular to \(V\)
of invertible hermitian matrix \(J\). For \(v\in V,\ z\in K\)
with \(z+\bar z=-v^*Jv\), write
\[
  N(v,z)=\begin{pmatrix}1&-v^*J&z\\0&I&v\\0&0&1\end{pmatrix},
 \qquad N^-(v,z)=\begin{pmatrix}1&0&0\\v&I&0\\z&-v^*J&1\end{pmatrix}.
\]
Let \(D\) be the subgroup generated by all these matrices and
\(M(t,r)={\rm diag}(t,r,\bar t^{-1})\) for \(t\in K^\times\),
\(r\in U(V)(k)\). The two displayed families are the radicals of
the opposite isotropic lines by Definition~\ref{def:full-frame-radical}.
\end{definition}

\begin{lemma}\label{lem:quasireflection-in-two-radicals}
In Definition~\ref{def:local-unitary-line-generators}, if
\(\theta\in K^\times\) has \(\bar\theta=-\theta\), then
\[
 w_\theta=\begin{pmatrix}0&0&\theta\\0&I&0\\-\theta^{-1}&0&0\end{pmatrix}
                                                        \in D.
\]
If \(v\in V\), \(\sigma\in K^\times\) and
\(\sigma+\bar\sigma=v^*Jv\), then
\(r=I-v\sigma^{-1}v^*J\) is an isometry of \(V\) and
\(M(-\theta\sigma,r)\in D\).
\end{lemma}
\begin{proof}
Multiply \(N(0,\theta)N^-(0,-\theta^{-1})N(0,\theta)\) to get
\(w_\theta\). Put \(z=-\sigma\) and \(w=-v/z\), so that
\(1/\bar z+1/z=-w^*Jw\). In the product
\[
                  g=w_\theta N(w,1/\bar z)N^-(v,z)
\]
the rightmost factor takes \(e\) to \((1,v,z)\). The middle
factor sends it to \((0,0,z)\): its middle coordinate is
\(v+wz=0\), and its first is
\(1+(v^*Jv)/\bar z+z/\bar z=0\). Thus \(ge=\theta z e\).
On a vector \((0,x,0)\), the middle coordinate after the last
two factors is \(x-wv^*Jx=rx\); \(w_\theta\) fixes this
coordinate. This product is an isometry preserving \(Ke\), so
its induced map on \(e^\perp/Ke\) is an isometry and equals
\(r\). Factoring the parabolic matrix as in
Definition~\ref{def:full-frame-radical} gives \(g=u M(\theta z,r)\)
with \(u\) in the upper radical. Hence \(M(-\theta\sigma,r)\in D\).
\end{proof}

\begin{theorem}\label{thm:local-unitary-two-radical-generation}
In Definition~\ref{def:local-unitary-line-generators} one has
\(D=\{g\in U(W)(k):\det_Kg=1\}\). Every \(g\in U(W)(k)\)
has the form \(dM(t,1)\), \(d\in D,\ t\in K^\times\).
\end{theorem}
\begin{proof}
First suppose \(K\) is a field. Every nonsingular hermitian
space over \(K\) is spanned by an orthogonal basis of nonzero
norms. Indeed if all norms vanish, expanding the norm of
\(x+cy\) for \(c=1\) and for a nonzero anti-invariant \(c\)
shows that all pairings vanish. Split off a nonzero norm vector
by orthogonal projection and repeat. The same reasoning supplies
an anti-invariant \(\theta\ne0\) by subtracting the conjugate of
an element outside \(k\).

Every isometry of \(V\) is a product of the maps \(r\) of
Lemma~\ref{lem:quasireflection-in-two-radicals}. To see this,
let \(x,y\) have the same nonzero norm and set \(v=x-y\),
\(\sigma=v^*Jx\). Then \(\sigma+\bar\sigma=v^*Jv\), and when
\(\sigma\ne0\) the indicated map sends \(x\) to \(y\).
If \(\sigma=0\), replace \(y\) by \(-y\), for which the
new scalar is \(2x^*Jx\ne0\). A second such map with vector
\(y\) and scalar \(y^*Jy/2\) sends \(-y\) to \(y\).
Use this procedure to align the image of the first member of
an orthogonal basis with that member, then repeat in its
perpendicular. Reflections chosen in the perpendicular fix
the aligned member. Induction on the dimension proves generation,
also in dimension zero. By the lemma and multiplication, for
each such isometry \(r\) some \(t_r\) satisfies \(M(t_r,r)\in D\).

Now take \(g e=(a,v,c)\). If \(c\ne0\), multiply on the left by
\(N(-v/c,\bar a/\bar c)\). This is a permissible parameter
since \(v^*Jv=-\bar ac-\bar ca\), and sends the vector to
\((0,0,c)\) by direct substitution. Applying \(w_\theta\)
moves it onto \(Ke\). If \(c=0\) and \(a\ne0\), a first
\(w_\theta\) makes the last entry nonzero. If \(c=a=0\),
then \(v\ne0\). Choose \(w\) with \(w^*Jv\ne0\) by
nonsingularity. Applying \(N(w,-w^*Jw/2)\) makes the first
entry nonzero. We have thus left multiplied by elements of
\(D\) to reach the parabolic of \(Ke\). Its upper radical lies
in \(D\), and its Levi part \(M(t,r)\) lies in \(D M(K^\times,1)\)
by the preceding paragraph. This proves the asserted representation.

Both radicals have determinant one. If \(dM(t,1)\) has determinant
one, then \(t=\bar t\), so \(t\in k^\times\). In this case
\(M(t,1)=w_{t\theta}w_\theta^{-1}\in D\), proving equality.
For the split algebra use either branch to identify the unitary
group with \({\rm GL}_{\dim V+2}(k)\), the other branch being
forced by the pairing. In the upper and lower radicals the
rectangular blocks adjacent to the end indices can be varied
independently on this branch; their conjugate entries belong to
the other branch. Thus \(I+xE_{ij}\) with \(i\ne j\) belongs
to \(D\) if at least one index is an end index. For two interior
indices use the commutator of \(I+xE_{i1}\) with \(I+E_{1j}\),
which equals \(I+xE_{ij}\). These row additions reduce an
invertible matrix to diagonal form, exchanging two pivots
up to sign by the two-by-two product giving \(w_\theta\)
with entries \(t,-t^{-1}\) in those two positions. Every
two-by-two diagonal block \({\rm diag}(t,t^{-1})\) is a product
of two such pivot exchanges. They reduce any determinant-one
diagonal to identity by successively clearing diagonal entries.
Consequently \(D\) contains the special linear group, and is
contained in it by determinants. The determinant of \(M((t,1),1)\)
on the chosen branch is \(t\), giving the representation for
every determinant in this case as well.
\end{proof}

\section{Equal-rank invariance at unmarked split places}
\begin{lemma}\label{lem:singular-frame-stratum-dilation}
Let \(K,O,\pi\) be as in Lemma~\ref{lem:relative-haar-invertible-matrices},
and let \(Q=|O/\pi O|\). Use the topology induced by the
matrix entries. Let \(\omega:{\rm GL}_n(K)\to\C^\times\)
be a locally constant character of the determinant, and put
\(a=\omega({\rm diag}(\pi,1,\ldots,1))\). Suppose
\[
                      a\ne Q^r\qquad(0\le r<n).
\]
Any linear functional \(\Delta\) on \(C_c^\infty(M_n(K))\)
which satisfies \(\Delta(\omega(h)L_h f)=\Delta(f)\) for
all \(h\) and vanishes on tests compactly supported on
\({\rm GL}_n(K)\) is zero.
\end{lemma}
\begin{proof}
For a closed subset \(Y\) of a locally closed subset \(X\) of
matrix space, restriction gives an exact sequence of test spaces
\[
 0\longrightarrow C_c^\infty(X\setminus Y)
 \longrightarrow C_c^\infty(X)\longrightarrow C_c^\infty(Y)
 \longrightarrow0 ,
\]
where the first arrow is extension by zero. All the spaces have a
basis of compact open neighborhoods: intersect small matrix balls
with a closed set in an open neighborhood where the locally closed
set is closed. To extend a compact open indicator on \(Y\), cover
its compact support by finitely many such neighborhoods in \(X\)
whose intersections with \(Y\) lie inside the indicator set.
Their union works. A test is a finite sum of these indicators.
A test in the kernel has compact open support disjoint from \(Y\),
giving exactness. For subsets preserved by left multiplication
the maps commute with that action.

Let \(X_r\) be the closed locus of rank at most \(r\), defined
by vanishing of its larger minors. Descend from \(X_n\). The
hypothesized vanishing on the invertible matrices makes
\(\Delta\) factor through \(X_{n-1}\). At each step it suffices
by exactness to prove that the induced functional is zero on
compactly supported tests of the rank \(r\) stratum. This
stratum is covered by the finitely many left invariant open
charts where a prescribed set of \(r\) columns is independent.
After placing those columns first, the matrix on the chart is
\([V\;VC]\), with \(V\) an \(n\)-by-\(r\) frame and \(C\)
arbitrary of the indicated size. These are homeomorphic
coordinates since on a nonzero minor the entries of \(C\) are
obtained by inverting that minor. Left multiplication acts only
on \(V\). Compactly supported tests can be partitioned among
the charts by a finite compact open refinement. On each chart
they are sums of indicators of products of compact open sets
in \(V,C\), by the product topology and compactness. Fixing
a test in \(C\), it thus suffices to treat a linear functional
on frames with the same transformation law.

Let \(V_0=\binom{I_r}{0}\), and set
\(B_m=V_0+\pi^m M_{n,r}(O)\) for \(m\ge1\). Each element
is \(kV_0\) with \(k\in1+\pi^mM_n(O)\), by adjoining the
unchanged final columns of the identity. Take \(h\) which
scales just row \(r+1\) by \(\pi\). Then
\(B_m\) is the union of \(Q^r\) disjoint translates of \(hB_m\)
by unipotent matrices. Indeed the row \(r+1\) in \(hB_m\)
lies in \(\pi^{m+1}O^r\). Add to it \(\pi^m t\) times the
first \(r\) rows, for \(t\) running over row vectors modulo
\(\pi\). The leading square block is \(I_r\bmod\pi\);
these additions give exactly the \(Q^r\) possible residues of
the row divided by \(\pi^m\). Subtracting the same row multiple
proves surjectivity onto each piece. Applying the functional to
the partition gives
\(\Delta({\bf1}_{B_m})=Q^r a^{-1}\Delta({\bf1}_{B_m})\),
since the row additions have determinant one. Thus its value
is zero. For \(r=0\) the frame space is a point, the empty
row partition has one piece, and the same conclusion holds.

Every full frame is \(gV_0\) by extension to a basis. The
sets \(gB_m\) give compact open neighborhoods at every frame.
Inside a fixed \(gB_m\), smaller boxes of the form \(gk B_{m'}\),
\(k\in1+\pi^m M_n(O)\), for \(m'\ge m\), partition it into
the ordinary residue boxes after applying \(g^{-1}\).
Their indicators have zero functional by covariance. Cover a
compact open test support with finitely many such chart boxes,
subtract earlier boxes to get disjoint compact open portions,
and subdivide each portion in its box to a common sufficient
depth by compactness. The functional is zero on every portion.
This proves the required vanishing on rank \(r\) and completes
the descent to the empty locus.
\end{proof}

\begin{theorem}\label{thm:unramified-split-equal-rank-invariance}
Let \(J_\nu\) satisfy Definition~\ref{def:actual-equal-rank-integral}.
Let \(w\notin S\cup\bar S\) lie above a split finite place of \(F\).
If \(u_w=\lambda_f(\pi_w)\ne1\) for a uniformizer idele on
the chosen branch, then
\[
                     J_\nu(g_0\widetilde s)=J_\nu(g_0)
                     \qquad(s\in{\rm GL}_n(E_w)).
\]
At every finite place, including the marked pairs, \(J_\nu\)
is right invariant under its full integral unitary group.
\end{theorem}
\begin{proof}
First the integral assertion holds off the marked set by
Theorem~\ref{thm:marked-equal-rank-period-invariance}. On a
marked pair an integral change of variables on the source
multiplies \(d_w(\det X)\) by \(\chi_w(\det s)^{-1}\).
The inverse source twist \(\gamma_0(\delta(\widetilde s))^{-1}\)
is \(\chi_w(\det s)\) and the splitting phase is one by
Theorem~\ref{thm:split-cut-weil-multiplier}. Hence the entire
vector is fixed. Products over all integral factors act in the
same way since off the finite set the vector is a tensor of
vacua.

At the unmarked \(w\) vary the finite test over all
\(C_c^\infty(M_n(E_w))\), and denote the period by \(\ell(f)\).
Absolute convergence for every such test follows using the
fixed marked place of Definition~\ref{def:actual-equal-rank-integral}.
In detail the test at \(w\) is invariant under finite Heisenberg
translations in a sufficiently small lattice on the two branches:
for positions use compact support and local constancy, and for
momenta use the annihilator of its compact support in
Theorem~\ref{thm:split-finite-schroedinger}. In the rational
displacement proof of Lemma~\ref{lem:two-sided-rank-detection},
multiply the displacement coefficient by an integer of \(E\)
that is one modulo the marked prime and its conjugate and
divisible by a sufficiently high power of both primes here.
Such an integer exists by the coprime ideal argument in
Lemma~\ref{lem:additive-approximation-away-from-place}. It
preserves the detecting character at the marked place and
increases the finite depths here to those fixing \(f\).
Its fixed infinite norms change the constant in the difference
estimate only. The polynomial theta estimate in the proof of
Theorem~\ref{thm:complementary-period-convergence} holds for a
fixed compactly supported locally constant finite test by clearing
its finite support denominators as in
Proposition~\ref{prop:rational-frequency-lattice-bound}.
The same iterated real difference and height integration now
give absolute convergence.

The source and target action remains
\(\omega(h)\omega(s)^{-1}L_h R_s\) by the split operator
calculation of Theorem~\ref{thm:marked-equal-rank-period-invariance}.
Here
\[
       \omega(h)=u_w^{-v_w(\det h)}|\det h|_w^{-n/2}.
\]
Indeed \(\delta(\widetilde h)=t/\bar t\) for the idele \(t\)
supported on \(w\) with value \(\det h\). In \(\gamma_0\) its
\(\nu^{-1}\)-value is \(\lambda_f(t)^{-1}\), while its
\(\upsilon_T^j\)-value is \(\upsilon(t)^{2j}\) by
Theorem~\ref{thm:actual-global-splitting-character}.
This last factor cancels the split ideal phase since \(n=2j\)
and \(\upsilon\) is unramified here. The extension of the
period satisfies \(\ell(\omega(h)L_h f)=\ell(f)\) by the
quotient change of variables of
Theorem~\ref{thm:marked-equal-rank-period-invariance}.

For fixed \(s\) subtract the two functionals by setting
\(\Delta(f)=\ell(\omega(s)^{-1}R_s f)-\ell(f)\).
It has the same left transformation since the two translations
commute. On compactly supported tests of the invertible locus
it is zero by Lemma~\ref{lem:relative-haar-invertible-matrices}.
The number \(a\) of Lemma~\ref{lem:singular-frame-stratum-dilation}
is \(u_w^{-1}Q^{n/2}\). If \(a=Q^r\), absolute values force
\(r=n/2\) and then \(u_w=1\), contrary to the hypothesis.
The lemma makes \(\Delta\) vanish on the whole matrix space.
Apply it to the unmarked full integral matrix indicator.
This is the finite vacuum in the split model of
Theorem~\ref{thm:split-finite-schroedinger}. Right translating
the source point acts by \(\omega(s)^{-1}R_s\), which proves
the desired equality.
\end{proof}

\section{Split primes generating the ideal class group}
\begin{lemma}\label{lem:angular-unit-value-prime-density}
Let \(E,F,\nu,\lambda,S\) be as in
Definition~\ref{def:complementary-hecke-values}, with nonzero
weight \(\mu\) and finite angular conductor as in
Theorem~\ref{thm:torus-hecke-nonvanishing-at-one}. For an
unmarked prime ideal \(\mathfrak p\) write
\(u_{\mathfrak p}=\lambda_f(\pi_{\mathfrak p})\) and
\(Q_{\mathfrak p}=N\mathfrak p\). As real \(\sigma\downarrow1\),
\[
 \sum_{\mathfrak p\ {\rm unmarked}}Q_{\mathfrak p}^{-\sigma}
       =\log\frac1{\sigma-1}+O(1),\qquad
 \sum_{\substack{\mathfrak p\ {\rm unmarked}\\u_{\mathfrak p}=1}}
      Q_{\mathfrak p}^{-\sigma}
       =o\!\left(\log\frac1{\sigma-1}\right).
\]
Exclude any additional fixed finite set in either sum if desired.
\end{lemma}
\begin{proof}
In the logarithm of an absolutely convergent Euler product
of unitary factors the terms with prime power at least two
have bounded total absolute value uniformly for \(\sigma>1\).
Indeed their sum is bounded by \(2\sum_{\mathfrak p}Q_{\mathfrak p}^{-2}\),
finite by Corollary~\ref{cor:ray-series-at-one}. The simple pole
of its trivial ideal series then gives the first assertion, since
a finite set contributes a bounded amount. For every fixed
positive integer \(l\), apply
Theorem~\ref{thm:torus-hecke-nonvanishing-at-one} to \(\nu^l\).
Its weight \(l\mu\) is nonzero, and one may omit the same finite
set even if the actual conductor decreases. Its continued value
at one is finite and nonzero, so taking the real logarithm of
its absolute value in the Euler region gives
\[
        \sum_{\mathfrak p\ {\rm unmarked}}
             \Re(u_{\mathfrak p}^{\,l})\,Q_{\mathfrak p}^{-\sigma}
                                                          =O_l(1).
\]
For \(m\ge1\) and \(|t|=1\) one has
\[
 {\bf1}_{t=1}\ \le\ m^{-2}|1+t+\cdots+t^{m-1}|^2
  =\frac1m+\frac{2}{m^2}\sum_{l=1}^{m-1}(m-l)\Re(t^l).
\]
The inequality holds since the right side equals one at \(t=1\)
and is everywhere nonnegative. Sum with prime weights,
divide by \(\log(1/(\sigma-1))\), and let \(\sigma\downarrow1\)
with \(m\) fixed. The upper limit is at most \(1/m\). Letting
\(m\) grow proves the second assertion.
\end{proof}

\begin{theorem}\label{thm:nontrivial-split-primes-generate-class-group}
Under the hypotheses of Lemma~\ref{lem:angular-unit-value-prime-density}
the ordinary ideal class group of \(E\) is generated by the
unmarked primes \(\mathfrak p\) lying over split places of \(F\)
with \(u_{\mathfrak p}\ne1\). The same holds after omitting
any fixed finite set of primes.
\end{theorem}
\begin{proof}
The sum of \(Q_{\mathfrak p}^{-\sigma}\) over primes lying
over inert finite places of \(F\) is bounded for \(\sigma>1\):
their absolute residue degrees are at least two. Above a rational
prime \(p\) there are at most \([E:\Q]\) primes: their norms
are at least \(p\) and their product divides \(N(pO_E)=p^{[E:\Q]}\)
by Proposition~\ref{prop:general-ideal-factorization}. The bound
\([E:\Q]\sum_{p}p^{-2}\) suffices. The arithmetic CM field has
only split and unramified quadratic finite local algebras by
Lemma~\ref{lem:real-place-rings}. Consequently the primes excluded
from the generating set have total weight
\(o(\log(1/(\sigma-1)))\) by
Lemma~\ref{lem:angular-unit-value-prime-density}.

If the generated subgroup is proper, the finite abelian class
group in Lemma~\ref{lem:bounded-ideal-classes} has a nontrivial
character \(\xi\) trivial on this subgroup. Indeed its nontrivial
quotient has a maximal proper subgroup. The resulting simple
abelian quotient is cyclic of prime order, by taking any nonzero
element and its cyclic subgroup, and admits a nontrivial
root-of-unity character. Write \(L(s,\xi)\) for the ideal
series of this character. It is holomorphic near one by the
ordinary ideal count in Theorem~\ref{thm:partial-ideal-counts}:
all classes have the same linear main term and their character
sum vanishes. Summing its error \(O(1+Y^{1-1/[E:\Q]})\) by
partial summation gives the asserted continuation as in
Corollary~\ref{cor:ray-series-at-one}. In the Euler region the
convergent prime-power logarithm therefore gives
\[
       \sum_{\mathfrak p}
             \Re\xi([\mathfrak p])Q_{\mathfrak p}^{-\sigma}
                  =\log|L(\sigma,\xi)|+O(1)\le O(1).
\]
On the generating primes the real part in this sum is one.
Everywhere else its difference from one is at most two, so
the sum equals
\(\log(1/(\sigma-1))+o(\log(1/(\sigma-1)))\) by the preceding
weight estimate. This contradiction proves generation.
\end{proof}

\section{Constancy of the equal-rank period}
\begin{lemma}\label{lem:equal-rank-central-invariance}
The period \(J_\nu\) of Definition~\ref{def:actual-equal-rank-integral}
is invariant on the right under central isometries \(zI\), where
\(z\bar z=1\) adelically.
\end{lemma}
\begin{proof}
The action of \((zI,zI)\) on \({\rm Hom}(W_0,W_0)\) is the
identity, and the determinant twists of the two equal-rank
groups agree by Definition~\ref{def:actual-equal-rank-integral}.
Consequently \(\Omega(g zI,l)=\Omega(g,l z^{-1}I)\).
Right translation by the central point \(z^{-1}I\) is also
left translation on the full target group and preserves Haar
measure there. It takes a rational fundamental domain to another
one, so also preserves its quotient integral by the argument
of Lemma~\ref{lem:quotient-measure}. The defining integral is
absolutely convergent, giving the identity.
\end{proof}

\begin{theorem}\label{thm:actual-constant-equal-rank-period}
In Definition~\ref{def:actual-equal-rank-integral} let the
weight of \(\nu\) be nonzero. Suppose \(W_0\) contains a
rational hyperbolic plane whose splitting over its perpendicular
splits the reference lattice at every finite place, with dual
line lattices. For every Schwartz \(v_1\) used in that definition,
\[
                     J_\nu(g)=J_\nu(1)
             \qquad\big(g\in U(W_0)(\mathbb A_F)\big).
\]
\end{theorem}
\begin{proof}
Use \(M(t,1)\) of Definition~\ref{def:local-unitary-line-generators}
on this particular plane, defined over \(F\) and on all its
adeles. At each place of \(F\) regard local points as supported
there. By Theorem~\ref{thm:equal-rank-invariant-under-radicals},
\(J_\nu\) is invariant on the left under the two radicals
supported at that place. It is therefore invariant on the left
under the entire local special unitary group, by
Theorem~\ref{thm:local-unitary-two-radical-generation}.
At an infinite place the determinant of a local isometry lies on
the complex unit circle. It is the \(n\)-th power of some point
\(z\) on this circle, since \(n=\dim W_0>0\), by division of
its real argument by \(n\). Removing \(zI\) leaves determinant
one. The preceding left invariance and the central invariance
of Lemma~\ref{lem:equal-rank-central-invariance} thus account
for the full archimedean unitary factors.

Define a function \(\varphi\) on the ideles of \(E\) by
\(\varphi(t)=J_\nu(M(t,1))\). It is invariant under multiplying
\(t\) by \(E^\times\), since a rational \(M\) acts on the left
and \(J_\nu\) is rationally invariant by
Theorem~\ref{thm:equal-rank-invariant-under-radicals} and its
proof. It is invariant under every infinite component by the
preceding paragraph. It is invariant under \(\widehat O_E^\times\):
the corresponding \(M\) preserves the split finite lattice
because both its line factors are units, and all integral unitary
points fix the function on the right by
Theorem~\ref{thm:unramified-split-equal-rank-invariance}.
Consequently \(\varphi\) factors through the ordinary ideal
class group of \(E\). Explicitly \(t\mapsto\prod_w\mathfrak p_w^{v_w(t_w)}\)
is onto the fractional ideals by choosing uniformizer powers
at finitely many places. Its kernel consists of finite units
and arbitrary infinite components. On rational ideles it gives
the principal ideals by Proposition~\ref{prop:general-ideal-factorization}.

At every prime \(\mathfrak p\) of
Theorem~\ref{thm:nontrivial-split-primes-generate-class-group}
the matrix \(M(\pi_{\mathfrak p},1)\) belongs to the full local
factor fixing \(J_\nu\) on the right by
Theorem~\ref{thm:unramified-split-equal-rank-invariance}.
Thus the function on classes is invariant under multiplication
by \([\mathfrak p]\), hence also by its inverse. These classes
generate the class group, so \(\varphi(t)=J_\nu(1)\) for
every idele \(t\).

Finally consider \(g\) supported at infinity and finitely many
finite places. Apply Theorem~\ref{thm:local-unitary-two-radical-generation}
at each of these places to write \(g=dM(t,1)\), where \(d\)
is a finite product of supported special unitary elements.
Left invariance gives \(J_\nu(g)=\varphi(t)\). Such elements
\(g\) are dense in the full restricted product: for any point
and any basic neighborhood one may keep its components where
conditions are prescribed and replace the others by identity
in the prescribed integral groups. Continuity of the period
in Definition~\ref{def:actual-equal-rank-integral}, proved by
Theorem~\ref{thm:complementary-period-convergence}, extends the
equality to every \(g\).
\end{proof}

\begin{corollary}\label{cor:parabolic-scalar-independent-of-argument}
Under the splitting and weight assumptions of
Theorem~\ref{thm:actual-constant-equal-rank-period}, put
\(j_\nu=J_\nu(1)\). Throughout the unimodular parabolic range
of Theorem~\ref{thm:parabolic-unimodular-coefficient-product},
\[
                     A_\nu^+(um)=K_\nu j_\nu\,A_\nu^-(um).
\]
The same scalar applies at every such \(u,q,g_0\) and depends
linearly on \(v_1\) for fixed remaining data and measures.
\end{corollary}
\begin{proof}
Substitute Theorem~\ref{thm:actual-constant-equal-rank-period}
in Theorem~\ref{thm:parabolic-unimodular-coefficient-product}.
By Definition~\ref{def:actual-equal-rank-integral} the dependence
of \(j_\nu\) on \(v_1\) is linear, as absolute convergence permits
addition under the integral. The factors defining \(K_\nu\)
in Theorem~\ref{thm:actual-levi-complementary-coefficients}
are independent of \(v_1\).
\end{proof}

\section{Fiber integrals of complementary frame tests}
\begin{definition}\label{def:marked-fiber-integrals}
Let \(K=\prod_{w\in S}E_w\), \(O=\prod_{w\in S}O_w\),
\(G_i={\rm GL}_i(K)\), \(K_i={\rm GL}_i(O)\), and \(N=b+d\)
be even with \(b,d>0\). Use the original large and small
orientations of Definition~\ref{def:complementary-period-characters}.
For \(t\in K^\times\) let \(\lambda_S(t)\) be \(\lambda_f\)
on \(t\) supported at the chosen branches, equal to one as an
idele outside them. Put
\[
      \omega(t)=\lambda_S(t)^{-1}\prod_{w\in S}|t_w|_w^{-N/2},
        \qquad \omega(k)=\omega(\det k)\quad(k\in G_i).
\]
Give \(G_i\) the product measures of
Lemma~\ref{lem:relative-haar-invertible-matrices}. A row test
\(\Phi\) is locally constant with compact support on the
surjections \(K^N\to K^d\); a column test \(\Psi\) has the
same properties on the injections \(K^b\to K^N\), in each
case of full rank at every factor. Both extend by zero.
Write \(x_0=[0\;I_d]\) and \(y_0=\binom{I_b}{0}\).
For \(k\in G_N\) put
\[
  B^+_\Phi(k)=\int_{G_d}\Phi(r x_0 k^{-1})\omega(r)^{-1}\dd^\times r,
 \qquad B^-_\Psi(k)=\int_{G_b}\Psi(k y_0 t)\omega(t)^{-1}\dd^\times t .
\]
The preimages of the compact supports on each fiber are compact.
Indeed the row and column coordinates describe the frames of a
fixed quotient or subspace, with continuous inverses obtained
by any full minor. This fiber is closed in the respective
full-rank locus, being defined by vanishing on a fixed plane
or taking values in it.
\end{definition}

\begin{lemma}\label{lem:periods-factor-through-frame-integrals}
With fixed real tests and all unmarked factors as in
Definition~\ref{def:complementary-period-characters}, the
original periods \(I_\Phi(g),I'_\Psi(g)\) converge absolutely
for all the tests of Definition~\ref{def:marked-fiber-integrals}.
For fixed \(g\), \(B^+_\Phi=0\) implies \(I_\Phi(g)=0\), and
\(B^-_\Psi=0\) implies \(I'_\Psi(g)=0\).
\end{lemma}
\begin{proof}
For a surjection \(X\) the image \(XO^N\) is a full lattice
in \(K^d\). Multiply \(X\) on the left by an invertible matrix
taking that lattice onto \(O^d\). This makes \(X\) integral
and its reduction surjective at each branch. Such matrices form
an open set, since integrality and a unit minor at each branch
are open conditions. For a column test apply the same procedure
to its transpose. By compactness, partition the support into
finitely many compact open pieces on each of which a single
such row or column operation works. Each transformed test
then depends on a finite residue depth inside the integral box.
The surjective or injective estimate of
Theorem~\ref{thm:complementary-period-convergence} applies on
these pieces. Undoing the transformation translates the integration
variable at the split places and multiplies by a constant. It
preserves finiteness because their Haar measures are right
invariant by Lemma~\ref{lem:relative-haar-invertible-matrices}
and the change of quotient domains in
Theorem~\ref{thm:marked-equal-rank-period-invariance}.

More precisely an integrated large target element \(r\in G_d\)
acts on \(\Phi\) by \(\omega(r)\Phi(r^{-1}\,\cdot)\), and an
integrated small source element \(t\in G_b\) acts on \(\Psi\)
by \(\omega(t)^{-1}\Psi(\cdot\,t)\). Indeed each acts on
\(N\) columns or rows, and has target character \(\gamma\)
or inverse source character \(\gamma^{-1}\). The split
substitution and ideal phase of
Theorems~\ref{thm:split-finite-schroedinger}
and~\ref{thm:split-cut-weil-multiplier} give the modulus
powers \(-N/2,N/2\) and the ideal phases
\(e_w^{-N v_w(\det r)},e_w^{N v_w(\det t)}\), respectively,
where \(e_w=\beta(\mathfrak p_w)\). The \(\upsilon_T^{N/2}\)
part of \(\gamma\) cancels these phases by
Theorem~\ref{thm:actual-global-splitting-character}. Thus both
claimed actions follow from the definition of \(\omega\).
The corresponding changes of integration variable leave the
periods invariant.

For surjections cover the Grassmannian of their kernels by the
charts on which a prescribed complementary set of \(d\) columns
is invertible. After ordering the columns the chart on frames
has coordinates \(X=[r C\;r]\), \(r\in G_d\), with arbitrary
\(C\). Order columns separately at different branches if needed.
For fixed \(C\) this is exactly a fiber \(r x_0 k(C)^{-1}\),
by taking an invertible shear and the indicated permutations.
If the row fiber integral vanishes, multiply the test by a
locally constant cutoff in \(C\); its fiber integral still
vanishes. A finite compact open partition of the compact
projected support allows these cutoffs inside individual charts.
Write a test there as a sum of \(\phi_i(r){\bf1}_{U_i}(C)\)
with the \(U_i\) disjoint compact open, by finite rectangle
refinement as in Lemma~\ref{lem:singular-frame-stratum-dilation}.
For each nonempty \(U_i\) the integral of \(\phi_i\omega^{-1}\)
vanishes. Holding the cutoff fixed, the large-period functional
on \(\phi_i\) satisfies the left transformation law of
Lemma~\ref{lem:relative-haar-invertible-matrices}. That proof
also works on the finite product \(K\), by taking products of
its compact open cosets. The functional therefore vanishes
on each piece. For columns use coordinates \(Y=\binom{I_b}{C}t\),
with the invertible chosen rows
first. Translation on \(t\) is now on the right. The same
coset proof with right translates, or the substitution
\(t\mapsto t^{-1}\), gives the functional as a scalar times
\(\int\phi(t)\omega(t)^{-1}\dd^\times t\). Thus the column
assertion follows as well.
\end{proof}

\begin{proposition}\label{prop:semilocal-matched-period-difference}
Let \(\mathcal M\) be the column tests for which a row test exists
satisfying \(B^-_\Psi(k)=\omega(k)B^+_\Phi(k)\) for all
\(k\in G_N\). For a fixed scalar \(c\), the automorphic
function \(D\Psi=I_\Phi-cI'_\Psi\) depends only on \(\Psi\)
and is linear on \(\mathcal M\). Every original matched pair
at the squarefree residue level belongs to this relation after
integral lifting. The kernel of \(D\) is stable under
\(\Psi\mapsto\Psi(q\,\cdot)\) for every \(q\in G_N\).
\end{proposition}
\begin{proof}
Two choices for the row test have difference with zero fiber
integral, so Lemma~\ref{lem:periods-factor-through-frame-integrals}
proves the first assertion. On \(K_N\), for an original residue
pair the row integrand can be nonzero only for \(r\in K_d\),
since the induced matrix of a primitive quotient frame generates
the entire target lattice. The column integrand can be nonzero
only for \(t\in K_b\) by the transposed assertion. The determinant
equivariance on ordered frames in
Definition~\ref{def:complementary-finite-frames} cancels the
weight \(\omega^{-1}\) there, and the groups have mass one.
Thus on \(K_N\) the relation is exactly
\(\Psi(ky_0)=\chi(\det k)^{-1}\Phi(x_0k^{-1})\) of that
definition. To extend it, for
\(m=\left(\begin{smallmatrix}a&u\\0&d_0\end{smallmatrix}\right)\)
the functions \(B^+_\Phi(km)\), \(B^-_\Psi(km)\) acquire
the factors \(\omega(d_0)^{-1},\omega(a)\), respectively,
by changing their integration variables. These give the same
law after inclusion of \(\omega(k)\). Moreover \(K_N\)
is transitive on \(b\)-planes over \(K\): intersect the plane
with \(O^N\) and extend a basis of this primitive summand as in
Proposition~\ref{prop:standard-test-satisfies-slice-rank}.
Hence every \(k\) is a product of an integral element with
such a block \(m\).

The pair transformed by \(q\) is
\((\omega(q)\Phi(\cdot\,q^{-1}),\Psi(q\,\cdot))\).
Its two fiber functions use the original functions at \(qk\),
with the extra \(\omega(q)\) on the row side. This preserves the
relation. On both periods this transformation is right translation
by the ambient element \(\widetilde q^{-1}\) times the same nonzero
scalar. Indeed, before division by the small-side multiplier,
a right translation by \(\widetilde k\) multiplies
\(\Phi(\cdot\,k)\) by
\(\prod_w e_w^{d v_w(\det k)}|\det k_w|_w^{d/2}\), and
\(\Psi(k^{-1}\,\cdot)\) by
\(\gamma(\delta(\widetilde k))\prod_w
e_w^{-b v_w(\det k)}|\det k_w|_w^{-b/2}\).
These follow from the same split substitutions, now acting on
the ambient variable in the two original orientations. Their
ratio at \(k=q^{-1}\) is \(\omega(q)\) since \(N=b+d\) and
the splitting characters cancel as above. This proves kernel
stability for arbitrary \(q\).
\end{proof}

\section{The determinant sum of shears}
\begin{definition}\label{def:marked-central-shear-projector}
In Definition~\ref{def:marked-fiber-integrals} let \(N=2b+n\)
and write the column frame \(Y\) in three successive blocks
\(P,H,Q\) of respective row sizes \(b,n,b\).
Let \(R=\prod_{w\in S} O_w/\pi_wO_w\), \(\pi=(\pi_w)\), and
\(d(M)=\chi^{-1}(\det M)\) on invertible residue matrices and
zero elsewhere. Use the same \(d\) on integral matrices by
reduction. Assume multiplicative primitivity as in
Lemma~\ref{lem:finite-determinant-fourier-stabilizer}.
Write \(D_Q(Y)={\bf1}_{\rm int}(Y)d(Q)\) and
\(D_P(Y)={\bf1}_{\rm int}(Y)d(P)\), putting either expression
zero off the integral box. The upper-right shear \(v(t)\)
adds \(tQ\) to \(P\), leaving \(H,Q\) fixed. Set
\[
       S_t(Y)=D_Q(v(t)Y),\qquad
       G(Y)=\sum_{t\in M_b(R)}d(t)\,S_{\pi^{-1}\widetilde t}(Y),
\]
where \(\widetilde t\) is any integral lift of \(t\). The
functions belong to \(\mathcal M\) of
Proposition~\ref{prop:semilocal-matched-period-difference}.
Indeed \(D_Q,D_P\) are standard partners of
Lemma~\ref{lem:standard-determinant-frame-pair}, and the action
of \(G_N\) there preserves \(\mathcal M\).
\end{definition}

\begin{theorem}\label{thm:finite-shears-recover-standard-test}
In Definition~\ref{def:marked-central-shear-projector}, the
function \(G\) is zero unless \(H\) is integral and
\(Q,\pi P\in{\rm GL}_b(O)\); on that set it equals \(d(-\pi P)\).
For \(a={\rm diag}(-\pi^{-1}I_b,I_n,I_b)\), let \(w(t)\)
add \(tP\) to \(Q\), preserving \(P,H\). Then
\[
         \sum_{t\in M_b(R)}G(a w(\widetilde t)Y)
                         =|{\rm GL}_b(R)|\,D_P(Y).
\]
In particular, if the global difference \(DG\) in
Proposition~\ref{prop:semilocal-matched-period-difference} is
identically zero, the original comparison holds for every
matched pair at this residue level.
\end{theorem}
\begin{proof}
If a summand is nonzero, \(Q\) is an integral unit matrix and
\(H\) is integral. The remaining integrality condition is
\(P+\pi^{-1}\widetilde t Q\) integral, equivalent to
\(\pi P\) integral and \(tQ\equiv-\pi P\bmod\pi\).
This gives exactly one allowable residue matrix \(t\); changes
of its lift do not affect the condition. The weight is zero
unless \(t\) is invertible. Otherwise its product with the
initial datum is \(d(t)d(Q)=d(-\pi P)\). This gives the first
formula. Upon substituting \(aw(\widetilde t)Y\) it says that
the summand is \(d(P)\) exactly when \(P\in{\rm GL}_b(O)\),
\(H\) is integral, and \(Q+\widetilde tP\in{\rm GL}_b(O)\).
In particular \(Q\) must be integral. Conversely for integral
\(Y\) with \(P\) invertible, \(t\mapsto Q+tP\bmod\pi\) is a
bijection of \(M_b(R)\). Exactly \(|{\rm GL}_b(R)|\) terms
therefore contribute, proving the formula, including all zero
cases outside that support.

The kernel in Proposition~\ref{prop:semilocal-matched-period-difference}
is a linear space stable under all the changes just used. Thus
\(D(G)=0\) implies \(D(D_P)=0\). The first-block projection
is an integral
quotient frame. By Corollary~\ref{cor:standard-single-determinant-period}
the comparison for this one standard pair implies it for every
matched pair at the level.
\end{proof}

\begin{lemma}\label{lem:primitive-marked-shear-frequency}
Suppose \(L:\mathcal M\to\C\) is linear and there are matrices
\(T_w\in M_b(E_w)\) and additive characters \(\psi_w\) with
annihilator \(O_w\) such that, for every \(z\in M_b(K)\),
\[
        L(S_z)=\prod_{w\in S}\psi_w({\rm tr}(T_w z_w))\,L(D_Q).
\]
If \(L(G)\ne0\), then \(L(D_Q)\ne0\) and every \(T_w\) belongs
to \({\rm GL}_b(O_w)\).
\end{lemma}
\begin{proof}
If \(L(D_Q)=0\), the stated law annihilates every summand of
\(G\). Otherwise \(S_z=D_Q\) for integral \(z\), since on
the support where \(Q,H\) are integral, \(P+zQ\) is integral
exactly when \(P\) is. The product character is then one on
all integral matrices. Varying a single entry at a single
place and using the annihilator shows that \(T_w\) has integral
entries there. On integral \(z\) divided by \(\pi\) the product
character factors through the finite matrix pairing with
\(\overline T=(T_w\bmod\pi_w)_w\). The residue additive character
is primitive because at each field the original conductor is
the integer ring. The ratio \(L(G)/L(D_Q)\) is the finite
matrix Fourier sum of \(d(t)\) at \(\overline T\), with its
sign adjusted if the opposite Fourier convention is used.
Lemma~\ref{lem:finite-determinant-fourier-stabilizer} applied
to \(\chi^{-1}\) makes this sum zero on noninvertible matrices.
This proves the necessity of integral unit determinant at
every marked branch.
\end{proof}

\section{A global criterion at the marked branches}
\begin{lemma}\label{lem:adelic-central-fourier-uniqueness}
For the compact additive quotient \([\mathcal Z]\) in
Definition~\ref{def:gram-indexed-mixed-period}, a continuous
function all of whose coefficients against the \(\xi_B\) vanish
is zero. Here \(B\) ranges over the rational hermitian
\(b\)-by-\(b\) matrices.
\end{lemma}
\begin{proof}
The rational bilinear pairing
\[
       (B,z)\longmapsto
         -\tfrac12{\rm Tr}_{E/\Q}(\jmath\,{\rm tr}(zB))
\]
identifies the hermitian matrix space with the rational dual of
the skew-hermitian one. Their dimensions agree by multiplication
by \(\jmath\), and a nonzero \(B\) pairs nontrivially with
\(\jmath^{-1}B\): the trace of \(B^2\) in each complex embedding
is a sum of squared moduli of its entries. Thus in a rational
basis of dimension \(h\) the \(\xi_B\) include the characters
for every rational frequency in \(\Q^h\).

For a positive integer \(m\) average the function over the
image of \(m\widehat\Z^h\) at the finite adeles, with probability
measure. The average is continuous and all of its same
coefficients vanish, by translation and compact integration as
in Lemma~\ref{lem:product-integration}. The quotient by this
compact subgroup is \(\R^h/m\Z^h\): subtract a rational vector
to put the finite part into \(m\widehat\Z^h\) by the finite
congruence argument of Lemma~\ref{lem:additive-approximation-away-from-place},
and the rational choices then differ by \(m\Z^h\). The induced
map from the real torus is a continuous bijection from a compact
space to a Hausdorff group, hence a homeomorphism. Pushing
forward the invariant probability measure gives uniform measure
there: translated equal boxes of side \(m/k\) partition the
torus up to their faces, parallel faces having measure zero
since infinitely many disjoint translates exist. Approximation
by such boxes gives the asserted measure on continuous tests.
All ordinary Fourier coefficients of the average vanish since
their frequencies belong to \(m^{-1}\Z^h\). The product
Fej\'er-kernel argument of Lemma~\ref{lem:gaussian-cover} shows
that the average is zero. Along \(m=j!\), the averaging groups
shrink into every neighborhood of zero at the finite adeles.
Uniform continuity on the compact quotient makes the averages
converge uniformly to the function, proving the claim.
\end{proof}

\begin{definition}\label{def:marked-index-at-standard-shear}
At the chosen branches use the integral frame of
Definition~\ref{def:adapted-finite-frame-slices}. For
\(t\in M_b(K)\) let \(z(t)\) be the adelic skew-hermitian
matrix giving the shear \(v(t)\) there, its forced conjugate
on the opposite branches and zero elsewhere. Thus
\(\widetilde{v(t)}=n(z(t))\) in
Definition~\ref{def:gram-indexed-mixed-period}.
For a rational hermitian \(B\) define its marked index
\(T(B)=(T_w(B))\) by
\[
    \xi_B(-z(t))=\prod_{w\in S}\psi_w({\rm tr}(T_w(B)t_w)).
\]
Use the conductor-\(O_w\) characters from the trace-dual
identity in Lemma~\ref{lem:hom-oscillator-dual-lattices}.
This specifies a unique matrix at each branch. Indeed the two
terms of the trace in \(\xi_B\) at a split pair are conjugate
field transports with the same local trace: the two minus signs
from \(\jmath\) and the skew matrix cancel. In the adapted
coordinates the resulting character is the local additive
character of a trace linear form in the entries of \(t_w\).
Divide its coefficients by the nonzero scalar defining
\(\psi_w\) in the trace character to obtain \(T_w(B)\), with
indices transposed for the matrix trace. Uniqueness follows
by testing a single matrix entry at arbitrary multiples, since
\(\psi_w\) has annihilator \(O_w\).
\end{definition}

\begin{theorem}\label{thm:global-primitive-marked-criterion}
Use \(D\) of Proposition~\ref{prop:semilocal-matched-period-difference}
and the data of Definition~\ref{def:marked-central-shear-projector}.
The exchange equality at this twist and scalar \(c\), for all
matched pairs at the prescribed residue level and every ambient
point, is equivalent to the following condition. For every \(g\)
whose components at the marked pairs are identity and each rational
hermitian \(B\) with \(T_w(B)\in{\rm GL}_b(O_w)\) for every
\(w\in S\), the coefficient at \(B,g\) of \(D(D_Q)\) is zero.
\end{theorem}
\begin{proof}
At fixed \(B,g\) integration against \(\overline{\xi_B}\)
on the compact center defines a linear functional \(L\) on
\(\mathcal M\) by taking this coefficient of \(D(\cdot)\).
The defining integrals exist since the periods are continuous:
in the compact partitions in
Lemma~\ref{lem:periods-factor-through-frame-integrals}, each
piece is a fixed translate of primitive tests, for which
Theorem~\ref{thm:complementary-period-convergence} gives
continuity. By the transformation in
Proposition~\ref{prop:semilocal-matched-period-difference}
and \(\det v(t)=1\), the function \(D(S_t)\) equals the
right translate of \(D(D_Q)\) by \(n(-z(t))\). If \(g\)
is identity at the marked pairs, it commutes with this
supported element. Moving the translate through \(g\) and
changing variables in the central integral gives
\(L(S_t)=\xi_B(-z(t))L(D_Q)\). Thus if \(L(G)\ne0\),
Lemma~\ref{lem:primitive-marked-shear-frequency} gives
\(T_w(B)\in{\rm GL}_b(O_w)\) at all marked branches and
\(L(D_Q)\ne0\). Under the stated condition every coefficient
of \(D(G)\) at this \(g\) therefore vanishes.
Lemma~\ref{lem:adelic-central-fourier-uniqueness} applied to
\(z\mapsto D(G)(n(z)g)\) gives \(D(G)(g)=0\).

We verify that such points and their left rational translates
are dense. At the finitely many marked places the images of
rational isometries are dense in the product of the local
unitary groups. To see this, the Cayley map
\(Z\mapsto(1+Z)(1-Z)^{-1}\) and its inverse
\(g\mapsto(g-1)(g+1)^{-1}\) give inverse rational maps between
the open skew-adjoint and unitary charts, by the Gram equation
as in Lemma~\ref{lem:unitary-cayley-density}. At a split pair
the chart where \(g+1\) is invertible is dense: on one branch
multiply by a scalar arbitrarily close to one, avoiding the
finitely many roots of \(\det(tg+1)\); the inverse conjugate
branch then also avoids eigenvalue \(-1\). The skew-adjoint
matrices form an \(F\)-vector space, whose rational points
approximate any prescribed tuple in these charts by
Lemma~\ref{lem:additive-approximation-away-from-place},
omitting a place outside this finite set. Such a place exists
since each rational prime has a prime of \(F\) above it by
Proposition~\ref{prop:general-ideal-factorization}. Inverting
the rational Cayley transform is legitimate for sufficiently
close approximants because their denominators are invertible
at these places, hence nonzero over the field. This proves density.

For arbitrary \(g\) choose rational \(\delta_i\) tending to its
marked components, and take \(g_i\) equal to identity there
and to \(\delta_i^{-1}g\) elsewhere. Then \(\delta_i g_i\to g\).
The function \(D(G)\) is continuous and left rationally invariant
by Theorem~\ref{thm:actual-adelic-rational-theta-invariance}
and rational triviality of the twists, so it is identically zero.
Theorem~\ref{thm:finite-shears-recover-standard-test} gives
exchange on all the required pairs. Conversely that exchange
makes \(D(D_Q)\) identically zero by using its standard partner,
and consequently gives every coefficient in the condition.
\end{proof}



\section{Transport of unimodular Fourier indices}
\begin{definition}\label{def:finite-transported-gram}
Use the data of Corollary~\ref{cor:parabolic-scalar-independent-of-argument}.
Identify \(L\) rationally with \(T_b\) by the reference frame
of Definition~\ref{def:actual-equal-rank-integral}. Its Gram matrix
on \(L\) is \(B_*\). On the finite lattice of \(L\) choose an
ordered adelic basis with columns \(C_L\) in rational coordinates.
Let \(C_T\) be such a basis for the finite lattice of \(T_b\),
and \(J_f=C_T^*B_*C_T\), in the identified rational coordinates.
For \(g=u\,m(q,g_0)\) as in Definition~\ref{def:full-frame-radical}
and any rational hermitian matrix \(B\) on \(L\), put
\[
   H_f(B,q)=(\alpha\bar\alpha)^{-1}C_L^*q_f^* Bq_f C_L .
\]
Here \(\alpha\) is the finite scale in
Definition~\ref{def:adapted-finite-frame-slices}. Call \(H_f(B,q)\)
unimodular if it and its inverse have integral entries at every
finite place. The definition is unchanged by replacing \(C_L\)
by \(C_L k\) with \(k\in{\rm GL}_b(\widehat O_E)\), since this
conjugates the form by \(k^*(\cdot)k\).
For either normalized period \(P^\pm\) of
Definition~\ref{def:normalized-complementary-levi} write
\(P^\pm_B(g)\) for its coefficient in
Definition~\ref{def:gram-indexed-mixed-period}.
\end{definition}

\begin{lemma}\label{lem:rational-levi-transports-gram-index}
For any continuous left rationally invariant function \(P\)
on the ambient group and any \(c\in{\rm GL}_b(E)\), one has
\[
       P_{B_*}\big(m(c,1)g\big)=P_{c^*B_*c}(g).
\]
This applies to both \(P^\pm\).
\end{lemma}
\begin{proof}
Conjugation by \(m(c,1)\) sends \(n(z)\) to \(n(czc^*)\) by
the three blocks of Definition~\ref{def:normalized-complementary-levi}.
The linear map \(z\mapsto czc^*\) on the skew-hermitian space
is invertible and rational. It induces an automorphism of its
compact additive adelic quotient preserving probability measure:
the pushed-forward measure is again translation invariant with
mass one. Uniqueness here follows directly by integrating
\(f(x+y)\) against two invariant probabilities in the two orders,
as in Lemma~\ref{lem:product-integration}. Now change the
integration variable to \(czc^*\) in the coefficient at \(m(c,1)g\).
The rational left factor \(m(c,1)\) drops out. The trace formula
of Definition~\ref{def:gram-indexed-mixed-period} gives
\(\xi_{B_*}(czc^*)=\xi_{c^*B_*c}(z)\) by cyclicity of
matrix trace, which follows by summing products of entries.
The two periods are rationally invariant by
Theorem~\ref{thm:actual-adelic-rational-theta-invariance},
since all determinant characters in the normalization are
trivial on rational points.
\end{proof}

\begin{theorem}\label{thm:all-unimodular-parabolic-indices}
Assume the hypotheses of
Definition~\ref{def:finite-transported-gram}, including the real
tests of Theorem~\ref{thm:matched-real-gram-slices}. For every
\(g=u\,m(q,g_0)\) preserving \(L\) and every rational hermitian
\(B\) for which \(H_f(B,q)\) is unimodular,
\[
                  P^+_B(g)=K_\nu j_\nu P^-_B(g).
\]
The same scalar works for all these \(B,g\).
\end{theorem}
\begin{proof}
Unimodularity forces \(B\) to be nonsingular. If at any real
place its signature differs from that of \(B_*\), both coefficients
vanish by Theorem~\ref{thm:singular-small-gram-coefficients-vanish}.
Indeed the real test on the large side has compact frame
support inside the specified nonsingular signature locus by
Theorem~\ref{thm:matched-real-gram-slices}.
The character used to normalize either period is constant on
\(n(z)g\), since \(\det n(z)=1\).

Suppose the signatures agree at every real place. At each
finite place the two integral unimodular matrices \(H_f(B,q)\)
and \(J_f\) are integrally isometric. The argument in
Lemma~\ref{lem:weil-unimodular-frame-independence} applies to
these two matrices directly: at a split prime use a basis
and its dual on the two branches. At an inert prime reduce
the two matrices modulo the uniformizer, apply the finite-field
diagonalization in Lemma~\ref{lem:finite-unitary-primitive-frame},
and lift the isometry one digit at a time by the trace correction
in Lemma~\ref{lem:weil-unimodular-frame-independence}. This gives
\(A\in{\rm GL}_b(\widehat O_E)\) with \(A^*J_f A=H_f(B,q)\):
each local choice and its inverse are integral, hence they form
adelic matrices. It follows by undoing both changes of basis,
including the scalar \(\alpha^{-1}\), that \(B,B_*\) are isometric
over every finite completion. Our field \(E=F(i)\) has
\(\sqrt3\in F\) by the twelfth root of unity in
Lemma~\ref{lem:skew-covol}. Theorem~\ref{thm:hermitian-local-global}
therefore gives \(c\in{\rm GL}_b(E)\) with \(c^*B_*c=B\).

In the identified coordinates the normalized finite matrix of
the moved reference frame is
\[
    Y(cq)_f=\alpha^{-1}C_T^{-1}c q_f C_L,\qquad
         Y(cq)_f^*J_fY(cq)_f=H_f(B,q)=A^*J_fA .
\]
Thus \(Y(cq)_f A^{-1}\) is unitary. The element
\(m(c,1)g\) still preserves \(L\), with Levi frame map \(cq\),
since conjugating \(u\) by \(m(c,1)\) preserves the radical
in Definition~\ref{def:full-frame-radical}.
Corollary~\ref{cor:parabolic-scalar-independent-of-argument}
now compares the two \(B_*\)-coefficients there.
Lemma~\ref{lem:rational-levi-transports-gram-index} gives the
stated equality at \(B,g\). The factors \(K_\nu,j_\nu\) in
that corollary are fixed before the rational change of frame.
\end{proof}

\begin{corollary}\label{cor:nonintegral-parabolic-gram-vanishing}
Under the same hypotheses, if \(H_f(B,q)\) has a nonintegral
entry then \(P_B^+(g)=P_B^-(g)=0\). This also holds for
every singular \(B\) and every \(B\) with an incorrect real signature.
\end{corollary}
\begin{proof}
For the latter two conclusions use
Theorem~\ref{thm:singular-small-gram-coefficients-vanish}
as in the preceding theorem. For the first, any contributing
rational frame \(x\) has \(x^*J_{T'}x=B\) for its target \(T'\),
by Theorem~\ref{thm:truncated-gram-period-formula}.
The transported finite frame \(h_f^{-1}xq_f C_L/\alpha\)
has columns in the target lattice: at every marked place both
slices in Theorem~\ref{thm:block-frame-finite-exchange} carry
the frame integrality indicator, and elsewhere the vacuum does
so by Lemma~\ref{lem:finite-witt-frame-partial-fourier}. Transport
of this support for a parabolic element follows from the dual
modulation argument in Proposition~\ref{prop:actual-coset-finite-coefficient}:
a frame outside the lattice is detected by a dual character
trivial on the lattice; that argument precedes any use of the
Gram determinant. Since \(h_f\) is unitary, the Gram matrix of
these columns is precisely \(H_f(B,q)\). Entries of the form
on the finite target lattice are integral by
Definition~\ref{def:adapted-finite-frame-slices}. Thus no frame
contributes at a nonintegral entry. The truncated formula and
its limiting assertion give the vanishing for both periods.
\end{proof}

\section{A rank bound from one split place}
\begin{definition}\label{def:column-partial-transform-at-place}
Fix a branch \(w\) in the nonempty marked set of
Definition~\ref{def:marked-fiber-integrals}. Write \(K=E_w\),
\(O=O_w\), take a uniformizer \(\pi\), and choose the additive
character \(\psi\) of conductor \(O\) whose
annihilator on \(O\) is \(O\), as in
Proposition~\ref{prop:marked-lattice-unitary-fourier}.
Fix a rational isotropic splitting as in
Definition~\ref{def:gram-indexed-mixed-period}, now with
isotropic dimension \(k\), \(1\le k\le N/2\), allowing \(k\)
to differ from the column count \(b\). Choose dual rational
bases on the two isotropic spaces. At \(w\) write the columns
in their three ambient blocks as \((P,H,Q)\), with \(P,Q\)
of size \(k\)-by-\(b\). The component at \(\bar w\) of a
central shear is determined by the component at \(w\).
On these columns the latter acts as
\(v(z)(P,H,Q)=(P+zQ,H,Q)\) for \(z\in M_k(K)\).
These coordinates may be obtained from the original coordinates
on the \(w\)-component by an invertible change of basis.
For a locally constant compactly supported column test \(f\),
extended by zero off its full-rank support, put
\[
  \widehat f(C,H,Q)=\int_{M_{k,b}(K)}
       f(P,H,Q)\psi(-{\rm tr}(CP))\dd P,\qquad C\in M_{b,k}(K).
\]
Other marked variables are kept as parameters. Use additive
measure giving the integral matrices mass one in each matrix space.
\end{definition}

\begin{lemma}\label{lem:partial-column-frequency-support}
Let \(f\) belong to a linear space of such tests stable under
\(f\mapsto f(v(z)\,\cdot)\) for every \(z\in M_k(K)\).
Suppose a linear functional \(\Lambda\) on this space satisfies
\[
 \Lambda(f(v(z)\,\cdot))=\psi({\rm tr}(Tz))\Lambda(f)
\]
for every test and every \(z\), with \(T\in M_k(K)\).
If \(\widehat f(C,H,Q)=0\) whenever \(QC=T\), for all \(H\)
and the other parameters, then \(\Lambda(f)=0\).
In particular \(\Lambda=0\) when \({\rm rank}_K T>b\).
\end{lemma}
\begin{proof}
A compactly supported locally constant function on a product of
matrix spaces is a finite linear combination of indicators of
products of compact open boxes. Indeed cover its compact support
by boxes on which it is constant, and subdivide at a common
precision in each variable. Its zero extension from the open
full-rank locus is locally constant also on the complement
because the support is compact and closed in the ambient space.
The transform of a \(P\)-box \(P_0+\pi^aM_{k,b}(O)\) is its
volume times \(\psi(-{\rm tr}(CP_0))\) on
\(\pi^{-a}M_{b,k}(O)\) and zero elsewhere. Integration of a
nontrivial character over a compact group vanishes by translation,
and the trace pairs the entries in opposite positions. The
annihilator assertion in
Definition~\ref{def:column-partial-transform-at-place} gives
exactly this dual box. Transforming back with the positive
character recovers the indicator: the two box volumes are
reciprocals by counting the cosets at successive precisions.
Thus the partial transform is injective and its result is again
locally constant and compactly supported in all the variables.

For \(m\ge0\) form the probability average of
\(\psi(-{\rm tr}(Tz))f(v(z)\,\cdot)\) over
\(z\in\pi^{-m}M_k(O)\). This is a finite linear combination in
the original space. A sufficiently small open subgroup fixes the
test by substitution, since \(Q\) is bounded on its support,
translations of \(P\) by a sufficiently small box preserve it,
and the inverse shear has the same property. Shrink the subgroup
to kill the scalar character there too. The functional of the
average equals \(\Lambda(f)\). Substituting \(P+zQ\) in the
partial transform gives its original value multiplied by
\(\psi({\rm tr}(QCz))\). Hence the transform of the average is
\[
            {\bf1}_{QC-T\in\pi^m M_k(O)}\,\widehat f(C,H,Q).
\]
The support of \(\widehat f\) is compact and equals its nonzero
set, since it is locally constant. The continuous image under
\((C,H,Q)\mapsto QC\) is compact and disjoint from \(T\)
by hypothesis. A sufficiently small box about \(T\) misses this
image, so the transform and hence the average are zero for all
large \(m\). This proves the first assertion. The product
\(QC\) factors through \(K^b\), proving the rank assertion.
\end{proof}

\begin{theorem}\label{thm:matched-difference-has-bounded-fourier-rank}
Let \(D\) and \(\mathcal M\) be as in
Proposition~\ref{prop:semilocal-matched-period-difference}, for
any scalar of comparison. For every rational isotropic splitting
of dimension \(k\) as in
Definition~\ref{def:column-partial-transform-at-place}, every
\(\Psi\in\mathcal M\), every adelic argument \(g\) and every
rational hermitian \(k\)-by-\(k\) index \(R\) with
\({\rm rank}_E R>b\), the \(R\)-coefficient of \(D\Psi\)
at \(g\) is zero.
\end{theorem}
\begin{proof}
First let \(g\) have identity component at the split pair of \(w\).
Taking the \(R\)-coefficient at \(g\) gives a linear functional
\(\Lambda\) on \(\mathcal M\). A shear supported at this split
pair has determinant one. By the transformation formula in the
proof of Proposition~\ref{prop:semilocal-matched-period-difference},
\(D(\Psi(v(z)\,\cdot))\) is the right translate of \(D\Psi\)
by the inverse shear. This supported shear commutes with \(g\).
In the compact central integral of
Definition~\ref{def:gram-indexed-mixed-period} its effect is
therefore multiplication by \(\xi_R(-z_{\mathbb A})\), where
\(z_{\mathbb A}\) denotes \(z\) at \(w\), its negative adjoint
at \(\bar w\), and zero elsewhere. In dual rational bases,
\(\xi_R(-z_{\mathbb A})=\psi({\rm tr}(Tz))\) with \(T\) a
nonzero local scalar times \(R_w\). Indeed the two branch traces
are equal: both \(\jmath\) and the skew-hermitian \(z_{\mathbb A}\)
change sign under conjugation, whereas \(R^*=R\). The two copies
cancel the factor \(1/2\) in the defining trace. Expressing the
resulting field trace character in the fixed normalization
\(\psi\) multiplies its argument by a nonzero scalar, as in
Proposition~\ref{prop:marked-lattice-unitary-fourier}. This also shows
directly that \({\rm rank}_K T={\rm rank}_E R\), since a rational
minor is nonzero exactly when its embedding in \(K\) is nonzero.
The linear space of tests is stable under every such shear by
Proposition~\ref{prop:semilocal-matched-period-difference}.
Lemma~\ref{lem:partial-column-frequency-support} applies.

For arbitrary \(g\), write \(g=g^{(w)}h_w\), with \(g^{(w)}\)
identity at the split pair and \(h_w\) supported there.
Right translation of \(D\Psi\) by \(h_w\) is a nonzero scalar
times \(D\) applied to a transformed column test, by the same
proposition, which applies to the entire split general linear
group. This translation passes through the defining left central
integral. Applying the first case at \(g^{(w)}\) proves the result.
\end{proof}

\section{Invariance at a bordered index of maximal rank}
\begin{definition}\label{def:bordered-null-line-radical}
Use a rational isotropic space \(L_+=Ee_0\oplus L\) and its
dual isotropic space with bases \(e'_0\) and a basis dual to
\(L\), where \(\dim_E L=b\). Denote the orthogonal complement of
the two dual spaces by \(W_2\), and put
\(M=L\oplus W_2\oplus L'\), where \(L'\) is the chosen dual of
\(L\). Thus \(W=Ee_0\oplus M\oplus Ee'_0\) is also a rational
isotropic splitting. Write \(U_e\) for the unipotent radical of
the line \(Ee_0\), given explicitly by
Definition~\ref{def:full-frame-radical} with frame dimension one.
It acts as the identity on the three successive quotients for
\(Ee_0\subset e_0^\perp\subset W\).
For any nonsingular rational hermitian \(B\) on \(L\) use the
index
\[
                 R_0=\begin{pmatrix}0&0\\0&B\end{pmatrix}
\]
on \(L_+\). For \(\Psi\in\mathcal M\) and fixed comparison
scalar let \(F_\Psi(g)=(D\Psi)_{R_0}(g)\) denote the coefficient
along \(L_+\) as in Theorem~\ref{thm:matched-difference-has-bounded-fourier-rank}.
\end{definition}

\begin{lemma}\label{lem:marked-maximal-rank-radical-invariance}
For any \(u\) supported at the split pair of one marked branch
and belonging there to \(U_e\), one has
\(F_\Psi(ug)=F_\Psi(g)\) at every \(g\).
\end{lemma}
\begin{proof}
First take \(g\) with identity component at this pair. As in
Theorem~\ref{thm:matched-difference-has-bounded-fourier-rank},
the coefficient there is a functional \(\Lambda\) with local
frequency \(T={\rm diag}(0,T_1)\), \(T_1\) invertible, in the
variables of Definition~\ref{def:column-partial-transform-at-place}.
Its value at \(ug\) equals its value at \(g\) on the test
\(\Psi(u^{-1}\cdot)\), because \(g\) commutes with the supported
element and \(\det u=1\). Work at the selected branch, leaving
all other marked variables fixed. In the enlarged splitting
write \(P=(P_0,P_1)^{\mathsf t}\), \(Q=(Q_0,Q_1)^{\mathsf t}\)
in row blocks and \(C=(C_0\;C_1)\) in column blocks, with index
zero referring to the first row or column. On \(QC=T\), both
\(Q:K^b\to K^{b+1}\) and \(C:K^{b+1}\to K^b\) have rank \(b\).
The zero first row of \(T\) gives \(Q_0 C=0\), hence \(Q_0=0\)
by surjectivity. The zero first column gives \(Q C_0=0\), hence
\(C_0=0\) by injectivity.

The row ordering for the line splitting is
\((P_0;(P_1,H,Q_1);Q_0)\), with \(H\) the positions in \(W_2\).
In these coordinates \(u^{-1}\) is upper triangular in three
blocks with identity diagonal by
Definition~\ref{def:bordered-null-line-radical}. At \(Q_0=0\) it
therefore fixes \(P_1,H,Q_1\) and adds to \(P_0\) a row depending
only on \(P_1,H,Q_1\). In the partial transform integral,
\({\rm tr}(CP)\) is independent of \(P_0\) when \(C_0=0\).
For fixed \(P_1,H,Q_1\), the indicated translation of \(P_0\)
preserves its measure. The integrals for the moved and original
test consequently agree on the whole fiber \(QC=T\), also when
the original test couples the various variables. The difference
of the two tests belongs to \(\mathcal M\) by
Proposition~\ref{prop:semilocal-matched-period-difference}.
Lemma~\ref{lem:partial-column-frequency-support} gives the
asserted equality at \(g\). For general \(g=g^{(w)}h_w\)
absorb the rightmost \(h_w\) into the datum in both coefficients
as in Theorem~\ref{thm:matched-difference-has-bounded-fourier-rank}.
\end{proof}

\begin{theorem}\label{thm:adelic-maximal-rank-radical-invariance}
For every \(u\in U_e(\mathbb A_F)\), every adelic \(g\), and
every test in Definition~\ref{def:bordered-null-line-radical},
\[
                              F_\Psi(ug)=F_\Psi(g).
\]
The assertion applies to any rational rank-\(b\) hermitian
index on an isotropic \((b+1)\)-plane, with \(Ee_0\) its kernel.
\end{theorem}
\begin{proof}
A rational \(v\in U_e(F)\) preserves \(L_+\). Its action on
that plane has matrix \(a=\left(\begin{smallmatrix}1&\ell\\0&I_b
\end{smallmatrix}\right)\) since it fixes \(e_0\) and acts trivially
on \(e_0^\perp/Ee_0\). Therefore \(a^*R_0a=R_0\). Conjugation
by \(v\) on the center of \(L_+\) sends \(n(z)\) to \(n(aza^*)\).
Indeed \(n(z)-1\) maps \(W/L_+^\perp\) to \(L_+\), and on the
former quotient \(v\) acts by \((a^*)^{-1}\) by the perfect
pairing. The rational change of variables and probability measure
argument of Lemma~\ref{lem:rational-levi-transports-gram-index}
therefore shows \(F_\Psi(vg)=F_\Psi(g)\).

Let \(v_0\) be the finite place of \(F\) below the chosen marked
branch. By the rational coordinates of
Definition~\ref{def:full-frame-radical} and
Lemma~\ref{lem:additive-approximation-away-from-place}, there
are \(v_i\in U_e(F)\) tending to \(u\) on the adeles away from
\(v_0\). Replace the \(v_0\)-component of \(v_i\) by that of \(u\),
obtaining \(u_i\). This changes \(v_i g\) by left multiplication
by an element of the supported local radical, so
\(F_\Psi(u_i g)=F_\Psi(v_i g)=F_\Psi(g)\) by
Lemma~\ref{lem:marked-maximal-rank-radical-invariance}.
Since \(u_i g\to ug\), take limits. The coefficient is continuous:
the period is continuous as in
Theorem~\ref{thm:global-primitive-marked-criterion}; translation
is continuous uniformly on the compact center quotient over a
compact neighborhood of \(g\). Integrating on that quotient
preserves this convergence.

For the final assertion choose a rational basis of the kernel
line and extend it to a basis of the given plane. The restriction
of the form to the chosen complement is nonsingular: a vector
of the complement orthogonal to it pairs also to zero with
the kernel, hence belongs to the kernel. The index is \(R_0\)
in this basis. Change the opposite isotropic basis by the
inverse adjoint, giving a rational splitting of the required
kind, and use the result just proved.
\end{proof}

\section{Two means on a hyperbolic plane}
\begin{definition}\label{def:two-line-means-on-gram-coefficient}
In Definition~\ref{def:bordered-null-line-radical} use the labels
\(e=e_0,\ e'=e'_0\), with \(\langle e,e'\rangle=1\).
In the rational frame of \(L\) identify \(V=E^b\) with \(L\)
and its dual coordinates with \(L'\). For \(x,y\in V_{\mathbb A}
=V\otimes_{\mathbb Q}\mathbb A_{\mathbb Q}\) let \(X(x)\) add
\(x\) to \(e\) and send \(l'\in L'\) to \(l'-e'x^*l'\);
let \(Y(y)\) add \(y\) to \(e'\) and send \(l'\) to
\(l'-e\,y^*l'\). Each fixes all other source summands. Thus, for
example, \(X(x)e=e+x\) and \(X(x)e'=e'\). Let
\(Z(z)=n(z)\) be the center for the isotropic space \(L\).
For fixed \(\Psi\in\mathcal M\) let \(F(g)=(D\Psi)_B(g)\).
With probability on \([V]=V_{\mathbb A}/V\), put
\[
 \begin{aligned}
 M_+(g)&=\int_{[V]}F(Y(y)g)\dd y,&
 \phi_+(x;g)&=M_+(X(x)g),\\
 M_-(g)&=\int_{[V]}F(X(x)g)\dd x,&
 \phi_-(y;g)&=M_-(Y(y)g).
 \end{aligned}
\]
\end{definition}

\begin{lemma}\label{lem:line-shift-character-and-averages}
The preceding formulas define continuous functions. Write
\(\eta_B(x,y)={\bf e}(-{\rm Tr}_{E/\mathbb Q}(\jmath x^*By))\)
with the adelic trace and character of
Definition~\ref{def:gram-indexed-mixed-period}. Then
\[
 \phi_+(x;Y(y)g)=\eta_B(x,y)\phi_+(x;g),\qquad
 \phi_-(y;X(x)g)=\eta_B(x,y)^{-1}\phi_-(y;g).
\]
Moreover \(M_+\) is the bordered coefficient at
\({\rm diag}(0,B)\) for \(Ee\oplus L\), and \(M_-\) is the
corresponding one for \(Ee'\oplus L\). In particular \(M_\pm\)
are left invariant under the respective adelic line radicals.
\end{lemma}
\begin{proof}
The maps \(X(x),Y(y)\) are isometries: the added pairing
of \(x\) with a vector of \(L'\) cancels the pairing of \(e\)
with \(-e'x^*l'\), and likewise on interchanging \(e,e'\).
All pairings of two added vectors in either map vanish.
Both maps are additive in their parameters under multiplication.
They fix \(L\) pointwise and commute with \(Z(z)\).
The latter acts on the left on \(F\) by \(\xi_B(z)\), by a
translation in the defining central integral. A rational \(X\)
or \(Y\) fixes \(F\) on the left by rational invariance of
\(D\Psi\) and commutation with \(Z\). This proves periodicity
of each integrand. Continuity follows from the compact integration
argument of Theorem~\ref{thm:adelic-maximal-rank-radical-invariance}.

The definitions on \(L'\) give
\[
    X(x)Y(y)=Z(yx^*-xy^*)Y(y)X(x).
\]
Indeed the mixed products of the two added maps on \(L'\) are
\(-xy^*\) and \(-yx^*\), respectively, and all longer products
vanish. The central character on their difference is
\(\eta_B(x,y)\), by moving the row to the front in each
trace and using \(\bar\jmath=-\jmath\), \(B^*=B\).
Substitute this commutation into each mean and translate
its integration variable on \([V]\). This gives both formulas.

For the first bordering, the skew matrices on \(Ee\oplus L\)
have independent coordinates consisting of the skew scalar on
\(Ee\), the off-diagonal column \(y\), and the skew block on
\(L\). Their central elements are the products of the shear
\(Z_e(s)\) sending \(e'\) to \(e'+se\), the element \(Y(y)\),
and \(Z(z)\), by the displayed actions. The compact quotient
and its probability measure are the products in these rational
coordinates, by uniqueness of invariant probability as in
Lemma~\ref{lem:rational-levi-transports-gram-index}.
Therefore the coefficient at \({\rm diag}(t,B)\), \(t\in F_{\rm re}\),
is the \(t\)-th Fourier coefficient of \(s\mapsto M_+(Z_e(s)g)\);
here \(F_{\rm re}=E^{\mathfrak c}\) denotes the base field
to distinguish it from the function \(F\).
For every \(t\ne0\) this index has rank \(b+1\), so the coefficient
vanishes by Theorem~\ref{thm:matched-difference-has-bounded-fourier-rank}.
Subtract the constant coefficient and use
Lemma~\ref{lem:adelic-central-fourier-uniqueness} for the one
dimensional hermitian index space. The function of \(s\) is
constant. Its value at zero is \(M_+(g)\), so \(M_+\) equals
the coefficient for \(t=0\). The computation with \(e,e'\)
exchanged gives the assertion for \(M_-\).
Theorem~\ref{thm:adelic-maximal-rank-radical-invariance}
gives both radical invariances.
\end{proof}

\begin{theorem}\label{thm:opposed-means-determine-each-other}
At any two points \(g,h\) and for any complex number \(c\), each
of the conditions \(\phi_+(\cdot;g)=c\phi_+(\cdot;h)\) and
\(\phi_-(\cdot;g)=c\phi_-(\cdot;h)\) implies the other. Either
one implies \(F(g)=cF(h)\).
\end{theorem}
\begin{proof}
Suppose the plus functions agree with this factor. For rational
\(r\in V\), move \(X(r)\) to the far left using the commutator
in Lemma~\ref{lem:line-shift-character-and-averages}. Left
rational invariance of \(F\) then gives
\[
   \phi_+(x+r;g)=
      \int_{[V]}\eta_B(r,y)^{-1} F(Y(y)X(x)g)\dd y .
\]
The characters \(y\mapsto\eta_B(r,y)\), \(r\in V\), supply all
rational frequencies on the compact vector quotient. Indeed
\(r\mapsto(y\mapsto-{\rm Tr}_{E/\mathbb Q}(\jmath r^*By))\)
is an injective rational linear map to its rational dual.
For \(r\ne0\) the \(E\)-linear row \(\jmath r^* B\) is onto
\(E\), so it takes the value one whose trace is \([E:\mathbb Q]\).
The two rational dimensions coincide, giving bijectivity.
The Fourier uniqueness proof of
Lemma~\ref{lem:adelic-central-fourier-uniqueness} in rational
coordinates therefore applies on \([V]\). Subtract \(c\) times
the same formula at \(h\) and use uniqueness for every \(x\).
It follows that \(F(Y(y)X(x)g)=cF(Y(y)X(x)h)\) for every
\(y,x\). Taking both zero gives the assertion for \(F\).
Commuting \(X(x)\) and \(Y(y)\) by the same common unit scalar
and integrating in \(x\) gives the equality of the minus functions.
Starting with the minus equality gives the reverse implication:
the commutator pairing is the inverse under reversing the group
order, and \(x\mapsto{\rm Tr}_{E/\mathbb Q}(\jmath x^*By)\)
is a nonzero rational linear functional for every \(y\ne0\)
by the same row argument and conjugation. Thus Fourier uniqueness
applies with the two parameters exchanged.
\end{proof}

\section{The projected lattice at the unmarked places}
\begin{definition}\label{def:unmarked-projected-lattices}
Use Definition~\ref{def:two-line-means-on-gram-coefficient}.
Suppose the finite unimodular lattice \(\Lambda\) of \(W\)
is the direct sum in the ordered splitting
\(Ee,L,W_2,L',Ee'\). Its paired isotropic summands are perfect
integral duals, and the lattice on \(W_2\) is unimodular.
Use the single standard pair \(E_S,D_Q\) in
Proposition~\ref{prop:semilocal-matched-period-difference},
and write \(f=D(D_Q)\). Let \(\mathcal U\) consist of all finite
places of \(E^{\mathfrak c}\) away from the marked pairs.
Call \(g\) admissible here if it fixes \(L\) pointwise at every
place outside \(\mathcal U\), including infinity. For a sign
\(\epsilon=+\) or \(-\) write \(e_+=e,\ e_-=e'\) and
\(L_\epsilon=Ee_\epsilon\oplus L\). Define, at
\(v\in\mathcal U\),
\[
       \Gamma_\epsilon(g)_v
       ={\rm pr}_{L}\big(g_v\Lambda_v\cap L_{\epsilon,v}\big),
\]
where projection has kernel the first line. Call an admissible
\(p\) a point preserving the decomposition at the unmarked
places if, at every \(v\in\mathcal U\), it preserves each of
\(Ee,Ee',M\) and also \(L\subset M\). Here \(M\) is as in
Definition~\ref{def:bordered-null-line-radical}.
At such a point the two projected lattices coincide and equal
\(\Gamma(p)_v=p_v\Lambda_{L,v}\), where
\(\Lambda_{L,v}=\Lambda_v\cap L_v\). Write \(\tau_v O_{E,v}e\)
for \(p_v\Lambda_{e,v}\). Its dual \(p_v\Lambda_{e',v}\)
is \(\bar\tau_v^{-1}O_{E,v}e'\), by the perfect pairing.
Here \(O_{E,v}\) is the full local integer algebra, including
both factors at split places. These assertions follow by applying
\(p_v\) to the indicated original intersections, since \(p_v\)
preserves both larger planes and their common subspace \(L_v\).
\end{definition}

\begin{lemma}\label{lem:straighten-unmarked-isotropic-flag}
For every admissible \(g\) and either sign there is an admis\-sible
\(p\) preserving the decomposition at the unmarked places, with
\[
       M_\epsilon(p)=M_\epsilon(g),\qquad
                 \Gamma(p)_v=\Gamma_\epsilon(g)_v
                           \quad(v\in\mathcal U).
\]
It can be chosen equal to \(g\) at every place outside \(\mathcal U\).
\end{lemma}
\begin{proof}
At every unmarked place the original periods are fixed on the
right by the stabilizer \(K_v=U(\Lambda_v)\) in their ambient
variable. Indeed the local test there is the vacuum of
Lemma~\ref{lem:finite-witt-frame-partial-fourier}; the integral
pair action fixes its unit vector by the lattice coefficient in
Theorem~\ref{thm:all-finite-weil-splitting}, for which the
intersection lattice is the original one and the ideal phase
is one. The determinant twists are trivial there on integral
units by Definition~\ref{def:complementary-period-characters}
and Theorem~\ref{thm:actual-global-splitting-character}.
This proves invariance under the product of the \(K_v\): it
also follows from invariance under finite subproducts and
continuity on the restricted product. The coefficient and both
its means retain right invariance by their definitions.

For a place \(v\in\mathcal U\), the integral isometries act
transitively on isotropic flags of a line in a \((b+1)\)-plane.
Indeed the intersection of a local lattice with a field subspace
is a saturated submodule: if a nonzero scalar multiple of a
lattice vector lies in the subspace, the vector does also.
Over a discrete valuation ring a saturated submodule is spanned
by part of a lattice basis. For detail, move an entry of least
valuation in any finite generating matrix to the first position
and clear its row and column by integral elementary operations.
Iterating puts the submodule in diagonal form, and saturation
forces each nonzero diagonal entry to be a unit. Such a matrix
of generators exists since multiplying a subspace basis by
suitable powers of a uniformizer gives a full sublattice, and
its intersection with a bounded box is contained in a fixed
denominator multiple of that lattice. The intervening quotient
is finite. Apply this first inside the plane for the line,
then inside the ambient lattice for the plane. At split places
perform this argument on each field factor. The resulting frame
has an integral left inverse and zero Gram matrix.
Theorem~\ref{thm:integral-unitary-primitive-frame} carries any
such ordered frame to another, including its distinguished
first vector. This proves transitivity.

Choose \(k_v\in K_v\) taking the reference flag to the inverse
image of the flag by \(g_v\). Then \(g_vk_v\) preserves the
flag. In the line parabolic it uniquely decomposes as \(u_v m_v\)
by Definition~\ref{def:full-frame-radical}, where \(m_v\)
preserves both hyperbolic lines and their common orthogonal
complement \(M_v\). The radical acts trivially on \(e_\epsilon^\perp/Ee_\epsilon\).
Flag preservation therefore implies that \(m_v\) preserves \(L_v\).
If \(g_v\) is integral one may choose \(k_v=g_v^{-1}\).
Thus we may replace \(g_v\) by \(m_v=u_v^{-1}g_vk_v\) at the
unmarked places to form an adelic \(p\). Right compact invariance
and left radical invariance from
Lemma~\ref{lem:line-shift-character-and-averages} preserve the
mean. The compact factor preserves \(\Lambda_v\), and \(u_v\)
preserves \(L_{\epsilon,v}\) and acts trivially on its quotient
by the line. These facts preserve the projected lattice.
\end{proof}

\begin{lemma}\label{lem:discriminant-order-projected-lattice}
On \(L_v\) let \(h_v(u,w)=(\alpha\bar\alpha)_v^{-1}u^*Bw\).
For a product of full lattices \(\Gamma=(\Gamma_v)_{v\in\mathcal U}\)
equal to \(\Lambda_{L,v}\) at almost all places, write
\(\Gamma^\#_v=\{u:h_v(u,\Gamma_v)\subset O_{E,v}\}\).
If each \(\Gamma_v\) is integral for \(h_v\), set
\[
                    \Delta(\Gamma)=\prod_{v\in\mathcal U}
                                  [\Gamma_v^\#:\Gamma_v].
\]
This is a positive integer. If integral \(\Gamma'\) contains
\(\Gamma\) at every place, then
\[
        \Delta(\Gamma)=\Delta(\Gamma')
                          \prod_{v\in\mathcal U}[\Gamma'_v:\Gamma_v]^2 .
\]
In particular a strict enlargement at one place strictly lowers
this integer. For a point \(p\) of
Definition~\ref{def:unmarked-projected-lattices} preserving the
decomposition, \(M_\epsilon(p)\ne0\) implies that \(\Gamma(p)\)
is integral.
\end{lemma}
\begin{proof}
In an integral basis with nonsingular Gram matrix \(H\) the
dual lattice has coordinates \(H^{-1}O_{E,v}^b\), since \(H^*=H\).
The index is the cardinality of \(O_{E,v}^b/H O_{E,v}^b\).
By the diagonal elimination in
Lemma~\ref{lem:straighten-unmarked-isotropic-flag}, on a field
factor of residue size \(q_w\) this is
\(q_w^{v_w(\det H)}\). At almost all places the rational \(B\)
and its inverse are integral in a fixed rational basis,
\(\alpha\) is a unit, and the lattice has that basis, so the
index is one. For a change from \(\Gamma'\) to its sublattice
\(\Gamma\) by an integral basis matrix \(A\), the new Gram
matrix is \(A^* H' A\). Its determinant gains the factor
\(\overline{\det A}\det A\). In the product over all branches
at \(v\), conjugation permutes branches of equal residue size.
The additional index is consequently exactly
\([\Gamma'_v:\Gamma_v]^2\), by the same elimination for \(A\).
Every strict local inclusion has index at least two.

For the last assertion nonvanishing of the mean supplies a
shift \(Y(y)\), or \(X(x)\) for the opposite sign, at which
the original \(B\)-coefficient \(F\) is nonzero. The shift fixes
\(L\) pointwise. The shifted point is therefore in the parabolic
at every place and induces on \(L\) the same frame action as \(p\).
Corollary~\ref{cor:nonintegral-parabolic-gram-vanishing} excludes
every nonintegral entry of the transported Gram matrix for
either period. It applies also to their difference before
normalization: the determinants are constant along the central
integration and the small-side change of model is the constant
multiple of Theorem~\ref{thm:opposite-branch-small-period}.
At the unmarked places the transported matrix is the matrix
of \(h_v\) on \(\Gamma(p)_v\). This proves integrality.
\end{proof}

\section{A least discriminant argument}
\begin{lemma}\label{lem:two-supports-at-least-discriminant}
Keep the data of Definition~\ref{def:unmarked-projected-lattices}
and fix a nonsingular rational hermitian \(B\). If \(F(g)\ne0\)
at some admissible \(g\), there exist a choice of order of
\(e,e'\) and an admissible \(p\) preserving the decomposition
at the unmarked places, such that \(M_+(p)\ne0\) and
\[
 \begin{aligned}
 \phi_+(x;p)\ne0&\ \Longrightarrow\
             \tau_v x_v\in\Gamma(p)_v &&(v\in\mathcal U),\\
 \phi_-(y;p)\ne0&\ \Longrightarrow\
             \bar\tau_v^{-1} y_v\in\Gamma(p)_v &&(v\in\mathcal U).
 \end{aligned}
\]
\end{lemma}
\begin{proof}
By Theorem~\ref{thm:opposed-means-determine-each-other} with
\(c=0\), some \(\phi_+(x;g)\) is nonzero. The point \(X(x)g\)
is admissible because the shift fixes \(L\) pointwise.
Lemma~\ref{lem:straighten-unmarked-isotropic-flag} gives a
point preserving the decomposition with a nonzero mean at the
origin. Among all such points for either sign choose one \(p\)
minimizing the positive integer \(\Delta(\Gamma(p))\), defined
by Lemma~\ref{lem:discriminant-order-projected-lattice}.
Exchange the names of the two lines if necessary, so that
\(M_+(p)\ne0\).

If \(\phi_+(x;p)\ne0\), straighten \(X(x)p\) for the plus
line. Its projected lattice at \(v\in\mathcal U\) is
\[
             \Gamma_+(X(x)p)_v
                 =\Gamma(p)_v+O_{E,v}\tau_v x_v.
\]
Indeed \(p_v\Lambda_v\cap L_{+,v}\) is the direct sum of
\(\tau_v O_{E,v}e\) and \(\Gamma(p)_v\).
The shift preserves \(L_{+,v}\), takes \(e\) to \(e+x_v\)
and fixes \(L_v\). Intersecting its moved lattice with this
plane and projecting gives the equality. The new point has
nonzero mean and the same projected lattice by
Lemma~\ref{lem:straighten-unmarked-isotropic-flag}; hence this
lattice is integral by Lemma~\ref{lem:discriminant-order-projected-lattice}.
A strict enlargement anywhere would lower \(\Delta\), contrary
to its choice. This gives the first implication. For the second
use \(Y(y)p\) and the minus line. Its projected lattice is
\(\Gamma(p)_v+O_{E,v}\bar\tau_v^{-1}y_v\) by the same
projection, using the actual dual line lattice in
Definition~\ref{def:unmarked-projected-lattices}. Minimality
was taken over both signs, so again there is equality.
\end{proof}

\begin{theorem}\label{thm:least-discriminant-self-dual}
At the point \(p\) of Lemma~\ref{lem:two-supports-at-least-discriminant},
one has \(\Gamma(p)_v^\#=\Gamma(p)_v\) for every \(v\in\mathcal U\).
\end{theorem}
\begin{proof}
Fix \(v\in\mathcal U\), write \(\Gamma=\Gamma(p)_v\), and let
\(u\in\Gamma^\#\). Let \(a\in V_{\mathbb A}\) be supported
at \(v\), with \(a_v=\tau_v^{-1}u\), using both field factors
when \(v\) is split. If \(\phi_-(y;p)\ne0\), then
\(w=\bar\tau_v^{-1}y_v\in\Gamma\) by
Lemma~\ref{lem:two-supports-at-least-discriminant}. Thus
\[
   a_v^* B y_v=u^*Bw\in(\alpha\bar\alpha)_v O_{E,v}.
\]
The trace dual identity
\(\jmath\alpha\bar\alpha\widehat O_E=\widehat O_E^\vee\)
in Lemma~\ref{lem:hom-oscillator-dual-lattices} gives
\(\eta_B(a,y)=1\), since the pairing in the exponent is
supported at \(v\). Where the minus function is zero the
same multiplier leaves its value unchanged. The commutator
formula of Lemma~\ref{lem:line-shift-character-and-averages}
therefore gives \(\phi_-(\cdot;X(a)p)=\phi_-(\cdot;p)\).
By Theorem~\ref{thm:opposed-means-determine-each-other} the
plus functions at these two points agree. At zero their equality
reads \(\phi_+(a;p)=\phi_+(0;p)=M_+(p)\ne0\).
The plus support implication of
Lemma~\ref{lem:two-supports-at-least-discriminant} now gives
\(u=\tau_v a_v\in\Gamma\). The reverse containment follows
from integrality in Lemma~\ref{lem:discriminant-order-projected-lattice}.
\end{proof}

\begin{corollary}\label{cor:primitive-indices-every-admissible-point}
Take the real tests and the scalar \(K_\nu j_\nu\) of
Theorem~\ref{thm:all-unimodular-parabolic-indices}, with
the nested integral splitting of
Definition~\ref{def:unmarked-projected-lattices}. In the original
orientations put \(c=K_\nu j_\nu/c_0\), where \(c_0\) is the
unit scalar of Theorem~\ref{thm:opposite-branch-small-period}.
If \(H_f(B,1)\) is unimodular at all marked pairs, then the
\(B\)-coefficient of \(f=D(D_Q)\) vanishes at every admissible point.
\end{corollary}
\begin{proof}
There is a marked pair by Definition~\ref{def:actual-equal-rank-integral},
so unimodularity there implies \(B\) is nonsingular over \(E\).
If the coefficient is nonzero, choose \(p\) by
Lemma~\ref{lem:two-supports-at-least-discriminant}.
Since \(M_+(p)\ne0\), some \(F(Y(y)p)\ne0\), where \(Y\)
denotes the opposite-line shift for the chosen order. This
point preserves \(L\) everywhere and has the same action \(q\)
as \(p\) on \(L\). At a marked pair this action is identity,
by admissibility. At every unmarked place, its transported
Gram matrix is the matrix of \(h_v\) on \(\Gamma(p)_v\) in
Definition~\ref{def:unmarked-projected-lattices} and
Lemma~\ref{lem:discriminant-order-projected-lattice}. This is
an invertible integral matrix with integral inverse by
Theorem~\ref{thm:least-discriminant-self-dual} and the basis
formula for the dual lattice. Hence \(H_f(B,q)\) is unimodular
at all finite places.
Theorem~\ref{thm:all-unimodular-parabolic-indices} compares the
two normalized coefficients at this point with factor \(K_\nu j_\nu\).
Restoring the common determinant normalization multiplies both
sides by the same character, since the central shear has determinant
one. By Theorem~\ref{thm:opposite-branch-small-period}, the
original small period is \(c_0\) times its directly transformed
one. Thus \(c=K_\nu j_\nu/c_0\) makes \(F(Y(y)p)=0\), a
contradiction.
\end{proof}

\section{The complementary identity on the enlarged spaces}
\begin{lemma}\label{lem:marked-index-is-transported-gram}
For the standard integral shear of
Definition~\ref{def:marked-index-at-standard-shear} and the
finite Gram matrix of Definition~\ref{def:finite-transported-gram},
a rational \(B\) satisfies
\[
 T_w(B)\in{\rm GL}_b(O_w)\quad(w\in S)
 \quad\Longleftrightarrow\quad
 H_f(B,1)\text{ is unimodular at every marked pair}.
\]
\end{lemma}
\begin{proof}
In rational coordinates on the skew center, the shear specified
by \(t\) in integral dual bases is \(C_L z C_L^*\), where
the coordinate \(z_w=t_w\), \(z_{\bar w}=-t_w^*\) at a split
pair. This follows by applying it to the dual basis, whose
basis matrix is the inverse adjoint of \(C_L\).
Cyclicity in the trace of
Definition~\ref{def:gram-indexed-mixed-period} replaces \(B\)
by \(C_L^*B C_L\). As in
Theorem~\ref{thm:matched-difference-has-bounded-fourier-rank},
the two branch traces are equal, canceling the \(1/2\);
the remaining character at \(w\) is the ordinary local trace
character applied to \(\jmath_w{\rm tr}(t_w (C_L^*B C_L)_w)\)
with the sign prescribed by \(\xi_B(-z)\).
The trace dual identity in
Lemma~\ref{lem:hom-oscillator-dual-lattices} says that the
character with coefficient \(\jmath_w(\alpha\bar\alpha)_w\)
has annihilator \(O_w\) on \(O_w\).
The trace character chosen in
Definition~\ref{def:marked-index-at-standard-shear} has the same
annihilator. Its coefficient therefore differs from this one
by a unit, since the annihilator for a scalar multiple is
the inverse scalar multiple of the original annihilator.
It follows by the entry tests of that definition that
\(T_w(B)\) is \(H_f(B,1)_w\) times a unit. A sign in the
choice of additive character also contributes a unit.
Since \(H_f(B,1)\) is hermitian its component at the opposite
branch is the conjugate transpose, preserving integrality
of the matrix and its inverse. This gives the equivalence.
\end{proof}

\begin{lemma}\label{lem:primitive-coefficients-at-finite-supported-points}
In Theorem~\ref{thm:global-primitive-marked-criterion} it
suffices to check the stated primitive indices only at points
whose infinite components and marked components are all identity.
\end{lemma}
\begin{proof}
At every such point the shear calculation in
Theorem~\ref{thm:global-primitive-marked-criterion} and
Lemma~\ref{lem:adelic-central-fourier-uniqueness} give
\(D(G)(g)=0\), where \(G\) is the finite shear sum of
Definition~\ref{def:marked-central-shear-projector}.
We prove density of the left rational translates of these points.
It suffices to approximate a prescribed element at all infinite
and marked places by one rational isometry. On each local group
the locus \(\det(h+1)\ne0\) is dense. For a split pair multiply
the matrix on one branch by a scalar close to one and avoid
the finitely many roots of \(\det(th+1)\); the other branch
is the inverse adjoint. At a real place multiply by a unit
complex scalar tending to one, avoiding those same finitely
many roots. This preserves the hermitian form for every signature.
On the locus in question the skew-adjoint Cayley coordinates
of Lemma~\ref{lem:unitary-cayley-density} and their inverse
are rational and continuous. The skew-adjoint matrices form
a rational vector space over the totally real field. Apply
Lemma~\ref{lem:arch-finite-approx} to the coordinates in a basis
of this space at all the chosen places together. Sufficiently
close rational coordinates have nonzero denominator determinant
at those places, hence over the field, and yield a rational
isometry. This proves the required approximation.

For arbitrary \(g\), choose rational points \(r_i\) approaching
its components at the specified places and let \(g_i\) be identity
there and \(r_i^{-1}g\) at the other finite places.
Then \(r_i g_i\to g\) since the other finite components are
already exactly those of \(g\). Continuity and rational invariance
of \(D(G)\), established in
Theorem~\ref{thm:global-primitive-marked-criterion}, give its
vanishing everywhere. Theorem~\ref{thm:finite-shears-recover-standard-test}
then gives the original exchange identity.
\end{proof}

\begin{theorem}\label{thm:global-nested-isotropic-comparison}
Assume the nonzero weight and isotropic hypotheses of
Corollary~\ref{cor:parabolic-scalar-independent-of-argument}
and the nested integral hyperbolic splitting in
Definition~\ref{def:unmarked-projected-lattices}.
Take the real tests of Theorem~\ref{thm:matched-real-gram-slices},
using the inverse of the real identification of
Theorem~\ref{thm:opposite-branch-small-period} for the original
small orientation. For every allowed twist \(\nu\) of the fixed
weight with squarefree marked conductor as in
Corollary~\ref{cor:standard-single-determinant-period},
for all original matched residue tests \((\Phi,\Psi)\) at
that level and at every adelic ambient argument \(g\), one has
\[
            I_\Phi(g)=c_\nu I'_\Psi(g),\qquad
                 c_\nu=K_\nu j_\nu/c_0 .
\]
Here \(j_\nu=J_\nu(1)\) for the chosen residual real vector,
\(K_\nu\) is the constant of
Theorem~\ref{thm:actual-levi-complementary-coefficients}, and
\(|c_0|=1\). In particular \(|c_\nu|=C_*|j_\nu|\) with
\(C_*>0\) independent of the twist and its marked set.
\end{theorem}
\begin{proof}
An argument \(g\) supported at the unmarked finite places is
admissible in Definition~\ref{def:unmarked-projected-lattices}.
For every marked primitive index,
Lemma~\ref{lem:marked-index-is-transported-gram} gives
unimodularity of the Gram matrix at the marked pairs. Therefore
Corollary~\ref{cor:primitive-indices-every-admissible-point}
gives the required vanishing of the \(B\)-coefficient for the
single fixed test \(D_Q\) and the scalar displayed here.
Lemma~\ref{lem:primitive-coefficients-at-finite-supported-points}
proves the identity on all tests and at every \(g\).
The real data are chosen using only the fixed weight, real
frames and the two residual real vectors by
Theorem~\ref{thm:matched-real-gram-slices} and the real action
in Theorem~\ref{thm:actual-levi-complementary-coefficients}.
Taking absolute values in the formula for \(c_\nu\) gives
\(C_*=|K_\nu|>0\), independent of the marked set and twist
by that last theorem.
\end{proof}

\section{Averaging a theta period along one line}
\begin{definition}\label{def:one-line-radical-mean}
Write \(W'=Ee\oplus W\oplus Ee'\) as a rational hermitian direct
sum of a hyperbolic plane and its perpendicular, with
\(\langle e,e'\rangle=1\); the order displayed is used for
matrix blocks. Use a finite lattice adapted to this splitting
as in Definition~\ref{def:adapted-finite-frame-slices}.
For any rational nonsingular target \(T\) with unimodular
finite lattice, identify a frame \(x:Ee\to T\) with its value
\(x(e)\in T(E)\). Write \(J_T\) for its hermitian matrix.
Let \(N_e\) consist of \(u(C,S)\) in
Definition~\ref{def:full-frame-radical}, with \(C\) a row on
\(W\) and \(S+\bar S=0\) scalar. For a function \(P\) on the
ambient quotient put
\[
        {\mathcal C}_e P(g)=
            \int_{N_e(F_0)\backslash N_e(\mathbb A_{F_0})}
                                     P(ng)\dd n,\qquad F_0=E^{\mathfrak c},
\]
using right invariant probability. For a theta test \(f\)
let \(t_x(g,h;f)\) have the meaning of
Definition~\ref{def:gram-indexed-mixed-period}, and put
\[
             d_x(g,h;f)=\int_{[N_e]}t_x(ng,h;f)\dd n ,
 \quad [N_e]=N_e(F_0)\backslash N_e(\mathbb A_{F_0}).
\]
\end{definition}

\begin{lemma}\label{lem:one-line-radical-measure}
The quotient \([N_e]\) has compact representatives. Integration
against its right invariant probability is obtained by taking
the mean first in \(S\) modulo rational skew scalars, then in
\(C\) modulo rational rows, using \(u(C,S)\).
For a fixed \(x\in T(E)\), the integrand defining \(d_x\) is a
left rationally invariant function of the radical variable.
\end{lemma}
\begin{proof}
By multiplying the three blocks in
Definition~\ref{def:full-frame-radical}, the product law adds
the two rows \(C,D\) and adds their two skew scalars with the
additional skew scalar
\((D J_W^{-1}C^*-C J_W^{-1}D^*)/2\). Either left or right
translation is thus a translation on the row variable and
a translation depending on that row on the scalar variable.
Both rational additive quotients have compact representatives:
a rational basis reduces this to \(\mathbb A_\mathbb Q/\mathbb Q\)
as in Definition~\ref{def:gram-indexed-mixed-period}. Concretely
the finite part can be made integral by a rational translation,
and then an integer translation puts the real coordinate in
\([0,1]\). First put \(C\) in its representatives by a rational
radical element and then put \(S\) there by a rational central
element. The two means are independent of the choice of lifts
of the row because the additional translation in \(S\) preserves
its mean. The same product law proves invariance on the right.
This is the normalized quotient measure: in these coordinates
ordinary additive Haar product is both left and right invariant
by successive translation, and the indicated representatives
give a fundamental domain up to null boundaries, using
half-open intervals in the rational coordinates.

A rational radical element preserves the momentum polarization
used in Lemma~\ref{lem:radical-residual-frame-action}. Its
Schr\"odinger action is the rational substitution and quadratic
phase in Theorem~\ref{thm:graded-parabolic-schroedinger}, with
trivial scalar by that lemma. It fixes the frame position \(x\).
On rational positions its phase is one because the global
additive character is trivial on rational arguments, and its
invertible rational substitution permutes the residual rational
positions at this \(x\). Summing as in
Definition~\ref{def:gram-indexed-mixed-period} proves the
claimed rational invariance.
\end{proof}

\begin{proposition}\label{prop:one-line-constant-term-frame-expansion}
For a period \(I_f\) satisfying the estimates of
Theorem~\ref{thm:complementary-period-convergence} in the
orientation of Definition~\ref{def:one-line-radical-mean},
\[
 {\mathcal C}_e I_f(g)
   =\int_{[U(T)]}\ \sum_{\substack{x\in T(E)\\x^*J_Tx=0}}
                                      d_x(g,h;f)\dd h .
\]
For fixed \(h,g\) the sum is absolutely convergent. Its grouped
value is absolutely integrable in \(h\). Moreover
\(d_0(g,h;f)=t_0(g,h;f)\).
\end{proposition}
\begin{proof}
For fixed \(g,h\), the full rational position sum after moving
the test by \(n\) is absolutely convergent uniformly on the
compact representatives in Lemma~\ref{lem:one-line-radical-measure}.
Indeed in a rational polarization the finite position supports
are contained in one bounded box after clearing denominators.
The moved real tests range in a set with uniformly bounded
Schwartz seminorms: \(n\) acts by the polynomial phase and
triangular substitution of Lemma~\ref{lem:radical-residual-frame-action},
whose coefficients and inverse coefficients range over a compact
set. The finite substitutions and phases give a common finite
supremum bound on this compact set as well.
In a rational basis the resulting sum of absolute values is
bounded by a constant times the sum of
\((1+\|a\|)^{-m}\) on one integer-denominator lattice, with
\(m\) greater than its real dimension plus one; this converges
by counting integer points in concentric boxes as in
Proposition~\ref{prop:adelic-schwartz-frequency-height}.
Consequently the frame sum may be integrated term by term for
fixed \(h\). The inner central translation acts at \(x\) by
\(\xi_{x^*J_Tx}\) by Lemma~\ref{lem:gram-shear-character-on-slices}.
Its mean vanishes unless \(x^*J_Tx=0\), by character orthogonality
as in Theorem~\ref{thm:truncated-gram-period-formula}. At \(x=0\)
the entire residual radical operator is identity by
Lemma~\ref{lem:radical-residual-frame-action}, giving \(d_0=t_0\).

For the integral in \(h\), \(ng\) ranges in a compact set
of representatives at fixed \(g\). The full kernel is bounded
there by an integrable function of \(h\), by
Theorem~\ref{thm:complementary-period-convergence} and
Proposition~\ref{prop:unitary-reduced-height-integrability}.
Thus the radical and target integrations may be interchanged.
At each \(h\) the grouped inner sum is the mean of that full
kernel. Bounding its modulus by the mean of the full kernel
modulus gives the asserted integrability and proves the formula.
\end{proof}

\begin{lemma}\label{lem:one-line-large-zero-frame-vanishes}
Let \(\dim W=N\), let \(b\ge0\), let \(V'\) have rank
\(N+2-b\), and at a marked branch choose an integral matrix
\(s:W\to E_w^b\) in the reference frames. On the enlarged
space use the standard large determinant test with upper rows
\(\widetilde s=s\,{\rm pr}_W\), as in
Lemma~\ref{lem:standard-determinant-frame-pair}. When read at
several marked branches impose the same form at each.
For this test \(f\) and every \(g\) preserving \(Ee\), one has
\(t_0(g,h;f)=0\) for every \(h\).
\end{lemma}
\begin{proof}
At the chosen branch let \(e_v\) generate the integral line
\(\Lambda_{e,v}\). A matrix \(D\) contributing to this test has
\(\left(\begin{smallmatrix}\widetilde s\\D\end{smallmatrix}\right)\)
invertible on reduction, and \(\widetilde s(e_v)=0\).
Since \(e_v\) has nonzero reduction by the adapted direct sum,
\(D(e_v)\) has nonzero reduction. This holds for every contributing
residue matrix. By Lemma~\ref{lem:frame-slice-support-transport}
the transported value of the zero frame would have to reduce to
\(D(e_v)\) in scaled integral coordinates. It is zero for every
\(g,h\) in the stated range. Thus the frame slice is zero already
at this branch, and its residual theta term is zero.
\end{proof}

\section{Descent of the small period}
\begin{definition}\label{def:lifted-small-line-test}
Write \(W'=Ee\oplus W\oplus Ee'\) as in
Definition~\ref{def:one-line-radical-mean}, \(N=\dim W\) even,
and let \(T_b=(A')^-\) be the directly transformed small
target of rank \(b\). Use compatible finite frames for the
summands at both ranks and the same opposite-target unit
\(\epsilon\) of Theorem~\ref{thm:opposite-branch-small-period}.
For integral \(s\) of size \(b\)-by-\(N\) at the marked branches,
take the small standard determinant of \(s\,{\rm pr}_W\) at
rank \(N+2\), and that of \(s\) at rank \(N\).
Let \(P_{W'}^-(v',s\,{\rm pr}_W)\) and \(P_W^-(v,s)\) be their
directly transformed normalized periods as in
Definition~\ref{def:normalized-complementary-levi},
with real vectors in this transformed orientation. Their
pair characters are respectively
\[
  (\nu\upsilon_T^{(N+2)/2},\upsilon_T^b),
       \qquad(\nu\upsilon_T^{N/2},\upsilon_T^b).
\]
Indeed this is Proposition~\ref{prop:complementary-levi-factor-quotient}
at each ambient rank, or follows from the two characters in
Theorem~\ref{thm:opposite-branch-small-period} by the common
ambient normalization. For \(v'\) written in the line frame
polarization, write \(v'[0]\) for its residual real vector at
frame zero. Use the same rational residual polarization at rank \(N\).
\end{definition}

\begin{lemma}\label{lem:nonzero-small-frame-line-average-vanishes}
For the small test of Definition~\ref{def:lifted-small-line-test}
and every \(x\ne0\) rational in \(T_b\),
\[
                    d_x(m(1,g),h;f)=0
\]
for all \(g\in U(W)(\mathbb A_{F_0})\), \(h\in U(T_b)(\mathbb A_{F_0})\).
\end{lemma}
\begin{proof}
Select one marked split pair and work in the position model on
the conjugate branch, where the directly transformed test is
named. In that model write the matrix as \((L_1,Z,L_2)\) in source
column blocks \(e,W,e'\), where \(Z\) has \(b\) rows.
The finite test before making the frame polarization is the
integral-box indicator times the determinant character from
the residual block \(Z\). This is because the original
\(s\,{\rm pr}_W\) kills both line summands and the split adjoint
in Theorem~\ref{thm:opposite-branch-small-period} identifies the
residual block with its own adjoint. On its support \(Z\) has
row rank \(b\) over the local field: the \(b\)-by-\(b\) determinant
of its adjoint composed with \(s\), up to the fixed units,
is nonzero on reduction. Applying \(m(1,g)\) and \(h\) here
only multiplies \(Z\) on its two sides by invertible matrices,
by Theorem~\ref{thm:split-finite-schroedinger}. It also acts
separately on the two outside columns. The transformed test
has compact support, and its residual block still has row
rank \(b\) throughout that support.

For each local column \(d\) of length \(N\) choose the source
line radical element whose matrix at this branch adds the
\(W\)-column \(d\) to \(e'\) and fixes \(e,W\). At the other
branch use the inverse adjoint, elsewhere the identity. The
line radical in Definition~\ref{def:full-frame-radical}
allows this arbitrary block in a split pair. These elements
form an additive subgroup; on positions before Fourier change
they act by \(L_2\mapsto L_2+Zd\), up to the common sign from
inversion, which can be absorbed by replacing \(d\) by \(-d\).
The split character and modulus for this matrix are one since
its determinant is one. In the line frame polarization one
Fourier transforms in \(L_2\), by the duality argument of
Lemma~\ref{lem:finite-witt-frame-partial-fourier}. With rational
frame \(x\) fixed, its Fourier variable is the opposite-branch
component of \(x\) in the dual target coordinates, including
the invertible trace scale. It gives a nonzero local row \(r_x\):
the rational nonzero vector has nonzero image at every embedding,
and the target pairing with the opposite branch is perfect.
Thus the supported radical acts on the slice by the character
\(d\mapsto\psi(r_x Z d)\) with nontrivial additive \(\psi\).

The set of \(Z\)'s in the projection of the original transformed
support is compact and consists of full-row-rank matrices.
The rows \(r_x Z\) on it form a compact set avoiding zero.
Consequently the average of these characters over a sufficiently
large compact additive box in \(d\) vanishes for every such
\(Z\), by the shrinking-annihilator computation of
Lemma~\ref{lem:partial-column-frequency-support}. Fourier
transformation in \(L_2\) preserves the projection bound for
\(Z\). The averaged test therefore has identically zero residual
slice at \(x\). To insert this local average in the defining
integral for \(d_x\), translate the integrating radical on the
right by each of the supported elements and use
Lemma~\ref{lem:one-line-radical-measure}. The operations commute
with the integral since the average on a compact box of a
locally constant moved finite test is a finite sum. For the
remaining left radical element the residual action at \(x\)
depends only on the slice there, by
Lemma~\ref{lem:radical-residual-frame-action}. Thus the entire
mean vanishes. If \(b=0\) there is no nonzero \(x\).
\end{proof}

\begin{theorem}\label{thm:small-standard-period-one-plane-descent}
With compatible physical measures in
Definition~\ref{def:lifted-small-line-test} there is a fixed
nonzero constant \(c_{\rm sm}\) such that for every real
Schwartz vector \(v'\), every allowed twist of the fixed weight,
every integral standard \(s\) at the squarefree level, and every \(g\),
\[
   {\mathcal C}_e P_{W'}^-(v',s\,{\rm pr}_W)(m(1,g))
                    =c_{\rm sm}\,P_W^-(v'[0],s)(g).
\]
With positive unit finite frame vacuum, \(c_{\rm sm}\) is the
inverse square root of the volume of the finite line frame
lattice for this target. It is independent of the twist and
marked set.
\end{theorem}
\begin{proof}
When the standard test is zero both sides vanish. Otherwise
the original small data at each rank are supported on full
column frames since their determinant is a unit. Convergence
and the compact uniform estimate in
Proposition~\ref{prop:one-line-constant-term-frame-expansion}
follow by Theorem~\ref{thm:complementary-period-convergence}
and the fixed central change of variables in
Theorem~\ref{thm:opposite-branch-small-period}.
That proposition and Lemma~\ref{lem:nonzero-small-frame-line-average-vanishes}
leave only \(t_0\) under the target integral.

At the finite places the slice at zero is
\(c_{\rm sm}\) times the smaller finite test. At a marked pair
the original function of \((L_1,Z,L_2)\) is the integral
indicator in the two outside columns times the smaller function
of \(Z\). Fourier transforming the last column takes its box
to its dual box, leaving the central block unchanged, by
Lemma~\ref{lem:finite-witt-frame-partial-fourier}. Both frame
boxes are one at zero. Away from the marked pairs the frame
vacuum and residual vacuum factor as in the same lemma.
In physical coordinates the positive unit frame vacuum is the
box indicator divided by the square root of its volume, by
Lemma~\ref{lem:fixed-finite-frame-vacuum-factor}. The same
fixed physical normalization therefore works at every marked
set. At infinity the zero slice is \(v'[0]\). Evaluation at
zero is continuous on the real Schwartz space by its defining
seminorms, so the residual vector is Schwartz.
Apply Lemma~\ref{lem:levi-operator-every-frame} with one
source line and \(q=1\). Its frame scalar is one, and the
target determinant character loses one power of \(\upsilon_T\);
these are precisely the two pairs of characters in
Definition~\ref{def:lifted-small-line-test}. The residual theta sum
and its integral over the same target quotient are consequently
the claimed smaller period.
\end{proof}

\section{The two scalar sections on a target line}
\begin{definition}\label{def:large-two-scalar-sections}
Let \(W'=Ee\oplus W\oplus Ee'\) and \(V'=Ea\oplus V\oplus Ea'\)
be rational hyperbolic splittings with the indicated dual
line pairings one. Write \(N=\dim W,\ r=\dim V\), and denote
the Gram matrices on these summands by \(J_W,J_V\). In these
bases an element of \({\rm Hom}_E(W',V')\) has block form
\[
                  \begin{pmatrix}
                       A&B_1&C_1\\ D&Z&H_1\\ I&J&K
                  \end{pmatrix}.
\]
Use as positions the whole first column \((A,D,I)\), the row
\(J\) on \(W\), and rational positions \(\zeta\) for \(Z\) as
in Lemma~\ref{lem:rational-graded-polarization}. The paired
momenta belong to the last column, the row \(B_1\), and the
remaining coordinates of \(Z\). This is a polarization: in the
hermitian trace form of Definition~\ref{def:hom-oscillator-pairing}
the first and last columns pair, \(B_1\) pairs with \(J\)
through \(J_W^{-1}\), and the \(Z\) block is orthogonal to all
these blocks. For a real-times-finite Schwartz test \(f\)
written in this model and \(t\in\mathbb A_E\), set
\[
 \begin{aligned}
  \mathcal L_+(t)f(\zeta)&= f(t,0,0,0,\zeta),\\
  \mathcal L_-(t)f(\zeta)&=\int_{\mathbb A_E}
     {\bf e}\big({\rm Tr}_{E/\mathbb Q}(\jmath\bar t u)\big)
                           \mathcal L_+(u)f(\zeta)\dd u .
 \end{aligned}
\]
Use the measure on the scalar adeles given in a rational basis
by Lemma~\ref{lem:rational-partial-poisson}; in particular the
rational scalar lattice has covolume one. Write \(\theta_{W,V}\)
for the residual rational position sum on \(\zeta\).
\end{definition}

\begin{proposition}\label{prop:isotropic-vector-gives-scalar-section}
For every \(\xi\in E^\times\), the mean of the frame term at
\(x=\xi a\) of Definition~\ref{def:one-line-radical-mean} is
\[
 d_{\xi a}(g,h;f)=\theta_{W,V}\big(\mathcal L_+(\xi)
                                  \Omega(g,h)f\big).
\]
For any \(t\in\mathbb A_E\), taking this plus section of a
test is unchanged by first applying the radical of the
target line \(Ea\).
\end{proposition}
\begin{proof}
Write \(f_1=\Omega(g,h)f\) and consider the mean for \(f_1\) at
the identity. For \(x=\xi a\) the central source phase is one
by Lemma~\ref{lem:gram-shear-character-on-slices}. We compute
the action of the remaining source radical \(u(C,0)\), where
\(C\) is a row on \(W\). On columns \((x,M,U)\) its underlying
linear map by precomposition with the inverse sends
\[
  (x,M,U)\longmapsto
  \big(x,\ M-xC,\ U+M J_W^{-1}C^*-\tfrac12 xC J_W^{-1}C^*\big).
\]
At \(x=\xi a\) the correction to \(M\) is purely in row \(B_1\),
which is momentum. Hence the map fixes every position with this
\(x\). On such a position put the original momenta equal to zero.
Let \(\Delta\) be the change under the displayed map and let
\(p\) denote the original position. Put \(t=x^*J_{V'}M J_W^{-1}C^*
=\bar\xi J J_W^{-1}C^*\). In the hermitian trace \(H(\Delta,p)\)
the last-column correction contributes \(\bar t\) by the pairing
with \(x\), its quadratic term contributing zero since
\(x^*J_{V'}x=0\). The middle-column correction contributes
\(-t\). Thus the total is \(\bar t-t\).
The phase of Theorem~\ref{thm:graded-parabolic-schroedinger}
is \({\bf e}(-{\rm Tr}_{E/\mathbb Q}(\jmath(\bar t-t))/2)
={\bf e}({\rm Tr}_{E/\mathbb Q}(\jmath t))\); the radical
scalar is trivial by Lemma~\ref{lem:radical-residual-frame-action}.
For a rational residual row \(J\ne0\), as \(C\) varies this is
a nontrivial character of the compact row quotient: \(J_W\)
is invertible, \(\xi\ne0\), and at one infinite place a suitable
component of \(C\) realizes every real trace value. For rational
\(C\) the trace is rational and its global character is one.
The mean kills the terms with \(J\ne0\) and retains every term
with \(J=0\). The interchange with the rational position sum is
justified by the fixed-\(g,h\) compact majorant of
Proposition~\ref{prop:one-line-constant-term-frame-expansion}.
The remaining variables are precisely \(A=\xi,\ D=I=J=0,\zeta\),
proving the formula.

For the target radical, its linear action preserves separately
the first and last source columns. On those hyperbolic variables
it acts by ordinary substitution with determinant one, by
Theorem~\ref{thm:hyperbolic-adelic-parabolic-splitting}, and
fixes \(ta\). In the residual middle column its upper triangular
target blocks leave \(J\) fixed and change \(Z\) only by a
multiple of \(J\). They preserve the momentum space. At \(J=0\)
the residual position is fixed and its change in momentum
is entirely in \(B_1\), pairing only with \(J=0\).
The quadratic phase is therefore one. To see that the
Schr\"odinger parabolic formula carries no further scalar on
this radical, conjugate by the target element scaling \(a\) by
two and \(a'\) by one half; this scales its central parameters
by four and its quotient row by two. The character-ratio
argument in Lemma~\ref{lem:radical-residual-frame-action}
applies verbatim, since here also the group preserves the
chosen residual momentum and the determinant twist is trivial
on the radical. This proves invariance of the residual section.
\end{proof}

\begin{lemma}\label{lem:scalar-sections-poisson-and-levi}
The two sections have the same rational theta sum:
\[
   \sum_{\xi\in E}\theta_{W,V}(\mathcal L_+(\xi)f)
           =\sum_{\xi\in E}\theta_{W,V}(\mathcal L_-(\xi)f),
\]
both absolutely convergent for fixed \(f\). Let \(\ell(c,h)\)
act on \(Ea,V,Ea'\) by \(c^{-1},h,\bar c\), where \(c\) is an
idele of \(E\), \(h\in U(V)(\mathbb A_{F_0})\), and take
\(g\in U(W)(\mathbb A_{F_0})\). For even \(N\) put
\(\gamma=\nu^{-1}\upsilon_T^{N/2}\), \(\lambda(c)=\nu(c/\bar c)\).
Use the characters \((\gamma\upsilon_T,\gamma\upsilon_T)\)
on the large pair at the upper rank, \((\gamma,\gamma)\)
on the residual pair. Writing \(|c|=|c|_{\mathbb A_E}\), one has
\[
\begin{aligned}
 \mathcal L_+(t)\Omega(m(1,g),\ell(c,h)) f
       &=\lambda(c)|c|^{-N/2}
                \Omega(g,h)\mathcal L_+(ct)f,\\
 \mathcal L_-(t)\Omega(m(1,g),\ell(c,h)) f
       &=\lambda(c)|c|^{-N/2-1}
                \Omega(g,h)\mathcal L_-(t/\bar c)f .
\end{aligned}
\]
\end{lemma}
\begin{proof}
The rational linear pairing \((t,u)\mapsto{\rm Tr}(\jmath\bar t u)\)
is perfect over \(\mathbb Q\): a nonzero coefficient may be
multiplied by its inverse to obtain the nonzero trace of one.
Thus in rational scalar coordinates its matrix gives a bijection
of the rational dual lattice of frequencies with \(E\).
Lemma~\ref{lem:rational-partial-poisson} gives the theta identity
at each residual rational position. On taking the sum there,
absolute convergence on both sides follows from the bounded
finite support and all real Schwartz seminorms, preserved by
the partial transform as proved in that lemma.

For the first transformation formula use
Lemma~\ref{lem:levi-operator-every-frame} on the first source
column with line multiplier one. The target inverse takes
\(ta\) to \(cta\), and the residual pair acting on the middle
column has characters \((\gamma,\gamma\upsilon_T)\).
Within that column the position row \(J\) is taken by the
underlying linear map to \(\bar c J g^{-1}\).
Write \(\delta(g)=\det_E g\). The determinant on this row
space is \(\bar c^{\,N}/\delta(g)\). By
Theorem~\ref{thm:residual-levi-zero-section} at \(J=0\), the
untwisted scalar is
\(\upsilon(\bar c)^{-N}\upsilon_T(\delta(g))|c|^{-N/2}\);
here \(|\delta(g)|=1\). Its factor at \(g\) cancels the extra
power in the source determinant character. The line contribution
to the target character is \(\gamma(c/\bar c)^{-1}\).
Since \(\upsilon(\bar c)=\upsilon(c)^{-1}\) and
\(\upsilon_T(c/\bar c)=\upsilon(c)^2\) by
Theorem~\ref{thm:actual-global-splitting-character}, the remaining
unit scalar is \(\lambda(c)\), with residual characters
\((\gamma,\gamma)\). This proves the first formula. Fourier
transformation in \(t\) multiplies the expression by \(|c|^{-1}\)
by the additive change of variables \(u\mapsto cu\), and
transports the dual argument to \(t/\bar c\) by its trace
pairing. The residual pair operator acts in independent variables
and commutes with this integral as in
Theorem~\ref{thm:residual-levi-zero-section}. This gives the
second formula.
\end{proof}

\section{Finite and infinite factors of the line sections}
\begin{definition}\label{def:compatible-standard-line-lattices}
Use the two splittings of
Definition~\ref{def:large-two-scalar-sections}. Endow each added
hyperbolic plane with the direct sum of the standard finite
lines \(\widehat O_E e,\widehat O_E e'\), respectively
\(\widehat O_E a,\widehat O_E a'\), with dual pairings one.
Keep the given unimodular lattices on the two perpendiculars.
Thus all four scalar Hom entries between the added planes have
the same physical scale \(\alpha\) of
Lemma~\ref{lem:hom-oscillator-dual-lattices}. At split places
use bases putting the line pairings equal to one and use the
same complementary bases at the upper and lower ranks.
Let \(b=N-r\ge0\). For the integral \(b\)-by-\(N\) matrix \(s\)
of Lemma~\ref{lem:one-line-large-zero-frame-vanishes} let
\(F^{\rm up}_{s}\) be the large test whose top rows are
\(s\,{\rm pr}_W\), and let \(F^{\rm low}_{s}\) be the smaller
large test with top rows \(s\). At a marked branch both use the
same ordering of the fixed rows and operator rows in the
determinant. Use positive normalized finite vacua in the two
Schr\"odinger models. Write \(\mathcal L_{\beta,f}\) and
\(\mathcal L_{\beta,\infty}\) for the finite and infinite factors
of Definition~\ref{def:large-two-scalar-sections}, with compatible
product measures. Thus on elementary tensors the full section
is the tensor of its two factors, by successive integration.
\end{definition}

\begin{proposition}\label{prop:finite-large-standard-scalar-sections}
For each \(\beta=+,-\) there is a fixed nonzero \(b_\beta\),
depending only on the framed spaces, the measures and phases.
At every allowable squarefree marked set \(S\) there is a
number \(G_{\nu,\beta}\), \(|G_{\nu,\beta}|^2=Q_S^{-1}\),
such that for finite ideles \(t\)
\[
        \mathcal L_{\beta,f}(t)F^{\rm up}_{s}
            =b_\beta G_{\nu,\beta}\,
                        D_S(t/\alpha)F^{\rm low}_{s}.
\]
Here \(D_S,Q_S\) are as in Definition~\ref{def:ramified-line-mellin}.
Both sections vanish at the scalar adele zero when \(S\ne\varnothing\).
\end{proposition}
\begin{proof}
First use at each finite place the lattice models in which the
vacua are their integral-box indicators. At an unmarked place the
orthogonal paired blocks in Definition~\ref{def:large-two-scalar-sections}
separate the vacuum into the frame box and the residual vacuum.
A partial transform of a box is its dual indicator in these
normalized models, by Proposition~\ref{prop:finite-polarized-vacuum}
and Lemma~\ref{lem:hom-oscillator-dual-lattices}. Evaluating at
the indicated zeros leaves the smaller vacuum times the condition
\(t/\alpha\) integral, for either scalar coordinate.

At a marked pair use the full split matrix position at \(w\).
Making the plus polarization of
Definition~\ref{def:large-two-scalar-sections} retains \(A,D,I,J,Z\)
at this branch and transforms \(K,H_1,C_1,B_1\) into the entries
at the opposite branch paired to them. This follows by pairing
the blocks in the hermitian trace as in
Lemma~\ref{lem:finite-witt-frame-partial-fourier}; the residual
\(Z\) is kept in its original local model. In units of \(\alpha\)
set \(A=u_w,\ D=I=J=0\), where \(u=t/\alpha\).
In the transformed direction integrate the compact boxes of
\(H_1,C_1,B_1\) at zero frequency. The top rows of the stacked
matrix are \((0,s,0)\), and its remaining rows are those of
the Hom matrix with the displayed entries. Expansion in its
first column and last row gives determinant
\[
                 (-1)^b u_w K\det\binom{s}{Z},
\]
where \(K\) is also in normalized coordinates. The ordinary
integral in \(K\) has the additive multiplier at
\(\overline{u_{\bar w}}\), up to a fixed unit and sign: the
paired scalar coordinates have line pairing one, and their
trace coefficient is \(\jmath\alpha\bar\alpha\), a generator
of the trace dual. Indeed a scalar supported in the \(K_w\) block pairs only with
the \(A_{\bar w}\) block. The hermitian product is supported
on a single branch. Moving its trace to \(w\) conjugates
\(\jmath\) and introduces only a sign, as in the normalized
split transform of Proposition~\ref{prop:marked-lattice-unitary-fourier}.
Lemma~\ref{lem:triangular-line-gauss-integral} gives \(g_{\chi_w}\)
times the local factor of \(D_S(u)\) and the residual determinant
test. The sign, and any unit in the chosen conductor normalization,
change the Gauss factor by a character of a unit, hence preserve
its absolute value. The three other box integrals have mass one.

For the minus section one additional Fourier transform is made in
the full scalar adele \(A\). At this split pair it transforms the
old \(A_w\) into the variable paired to it at \(\bar w\), and
inverts the preceding transform of \(K_w\), up to a sign.
Indeed the Fourier transform of the Fourier transform with
opposite pairing recovers a function by the box computation in
Lemma~\ref{lem:partial-column-frequency-support}; using the
same sign just negates its argument in the symmetric convention.
Thus one sets \(K_w\) equal to \(u_w\) up to a unit and
integrates \(A_w\) with the same unit-normalized frequency
\(\overline{u_{\bar w}}\). The other three frequencies are
again zero. The product of the two scalar diagonal entries in
the stacked determinant has the same character law, and the
interchanged-entry calculation in
Lemma~\ref{lem:triangular-line-gauss-integral} applies.
For both sections multiplying the local Gauss moduli gives
\(\prod_{w\in S}q_w^{-1}=Q_S^{-1}\) after squaring. The
calculation also gives zero at \(t=0\).

Passing back to fixed physical measures and unit vacua multiplies
each full finite section by the fixed factor \(b_\beta\).
For detail, evaluate the same polarizations on the unmarked
vacuum at frame zero. Each normalized box is there replaced
by its nonzero unit-vector value; each transform has a fixed
measure conversion. The chosen local normalized transforms map
the box to its dual box, so these factors are exactly the same
for every test processed in that polarization. They are one at
almost all places, where the scale is a unit and the physical
self-dual lattice measure is one. This is the product normalization
argument of Lemma~\ref{lem:fixed-finite-frame-vacuum-factor},
now also applied to the scalar Fourier variable. The two
choices of polarization may have different fixed factors.
\end{proof}

\begin{theorem}\label{thm:real-line-cone-schwartz}
Let \(K_\infty\) be the product of the compact
\(U(p)\times U(q)\) in the real target group for \(V'\),
using signed orthonormal bases. Give it invariant probability
and give \((\mathbb C^\times)^\Phi\) product multiplicative Haar
measure \(d^\times t=\prod_\tau (dr_\tau/r_\tau)d\theta_\tau/(2\pi)\),
\(t_\tau=r_\tau e^{i\theta_\tau}\). Write
\(|t|_\infty=\prod_\tau r_\tau^2\) and let
\(\lambda_\infty(t)=\nu_\infty(t/\bar t)\) for the prescribed
weight. For either \(\beta\), a real Schwartz vector \(v'\) on the
upper pair and real \(\sigma>0\), define
\[
 V_\beta(\sigma;v')=\int_{K_\infty}\int_{(\mathbb C^\times)^\Phi}
    \lambda_\infty(t)|t|_\infty^\sigma\,\mathcal L_{\beta,\infty}(t)
             \Omega_\infty(1,k)v'\,\dd^\times t\,\dd k .
\]
Use the upper character pair of
Lemma~\ref{lem:scalar-sections-poisson-and-levi}.
This integral converges in every real Schwartz seminorm, gives
a Schwartz vector on the smaller pair, and depends continuously
on \(\sigma\) in those seminorms. The construction depends only
on the archimedean character.
\end{theorem}
\begin{proof}
For either section, \(v'\) is written after a fixed linear
change of polarization and, for the minus section, a partial
Fourier transform. These are continuous Schwartz operators by
Lemma~\ref{lem:rational-partial-poisson} and the substitution
and phase formulas of Lemma~\ref{lem:real-oscillator-smooth-near-identity}.
Restricting a Schwartz function to the coordinates \((t,0,0,0,\zeta)\)
bounds any residual Schwartz seminorm by
\(C_m\prod_\tau(1+r_\tau)^{-m}\), for every \(m\).
Indeed multiply its residual derivatives by the required
residual coordinate powers and also by sufficiently many powers
of the scalar coordinates before taking the supremum. The
constants are uniform for \(\Omega_\infty(1,k)v'\), \(k\in K_\infty\):
near each fixed \(k\) use the representation law and the
continuity in every seminorm of
Lemma~\ref{lem:real-oscillator-smooth-near-identity}, and cover
the compact group by finitely many such neighborhoods.
Fix \(0<\sigma_0\le\sigma\le\sigma_1\). For \(r_\tau\le1\)
the radial majorant is \(r_\tau^{2\sigma_0-1}\dd r_\tau\);
for \(r_\tau\ge1\) it is \(r_\tau^{2\sigma_1-m-1}\dd r_\tau\).
Both are integrable for \(m>2\sigma_1\). Their product
dominates each prescribed seminorm. Differentiating in
\(\zeta\) under the integral is justified by the same majorants
for the next derivatives. Tail truncation gives convergence
in each seminorm, and dominated convergence gives continuity
in \(\sigma\) by applying the bounds to the difference of the
scalar factors. The unitary phase and the compact action at
infinity use only the fixed archimedean weight.
\end{proof}

\section{Integration over rational isotropic lines}
\begin{lemma}\label{lem:target-line-parabolic-quotient-measure}
Use the splittings of Definition~\ref{def:large-two-scalar-sections}
and put \(G=U(V')(\mathbb A_{F_0})\), \(\Gamma=U(V')(F_0)\).
Let \(P\) be the adelic stabilizer of \(Ea\), \(N_a\) its unipotent
radical, \(\Gamma_P=\Gamma\cap P\), and \(K=K_\infty K_f\),
where \(K_f\) preserves the given unimodular lattice and \(K_\infty\)
is as in Theorem~\ref{thm:real-line-cone-schwartz}. Use probability
on \(K\) and on \([N_a]=N_a(F_0)\backslash N_a(\mathbb A_{F_0})\).
Let \(C_E=E^\times\backslash\mathbb A_E^\times\), and normalize
its quotient measure by counting on \(E^\times\) and the product of
the multiplicative measures of Definition~\ref{def:ramified-line-mellin}
and Theorem~\ref{thm:real-line-cone-schwartz}. Write
\([U(V)]=U(V)(F_0)\backslash U(V)(\mathbb A_{F_0})\).
There is a positive \(C_P\), depending only on the spaces, compact
groups and measures, such that for any nonnegative measurable
\(A\) on \(\Gamma_P\backslash G\),
\[
 \int_{\Gamma_P\backslash G}A(y)\dd y
  =C_P\int_{C_E\times[U(V)]\times K}
     |c|^{r+1}\int_{[N_a]} A(n\ell(c,h)k)\dd n\,
                           \dd^\times c\,\dd h\,\dd k .
\]
The inner mean is independent of the choices of quotient
representatives for \(c,h\). The same equality holds for
absolutely integrable complex functions.
\end{lemma}
\begin{proof}
At a finite place the intersection of an isotropic field line
with the integral lattice has a primitive generator. This follows
by clearing denominators and then dividing a generator by the
largest common uniformizer power in its coordinates, separately
in both factors at a split place. The local algebra is split or
unramified by Lemma~\ref{lem:real-place-rings}.
Theorem~\ref{thm:integral-unitary-primitive-frame} sends such a
generator to the reference generator. At a real place an isotropic
vector in the orthogonal sum of the positive and negative spaces
has components of equal nonzero length. After rescaling, the
unitary groups of these two definite spaces can carry the two
components to those of a fixed generator of \(Ea\). To see this,
extend each unit vector to an orthonormal basis by successively
subtracting projections from a spanning list and normalizing its
nonzero remainders. Thus at all places \(K\) is transitive on
the respective isotropic lines. For \(y\in G\), choose compact
elements carrying the reference line to its inverse image by
\(y\). At almost all places take \(y_v^{-1}\). This proves
\(G=PK\). Every member of \(P\) is uniquely \(n\ell(c,h)\):
its actions on the line and on \(a^\perp/Ea\) determine \(c,h\),
and the remaining triangular matrix belongs to the radical by
Definition~\ref{def:full-frame-radical}. The same assertion
with rational entries describes \(\Gamma_P\).

Conjugation by \(\ell(c,h)\) on the radical row is
\(C\mapsto c^{-1}Ch^{-1}\). Its additive modulus is
\(|c|^{-r}\), since the determinant of \(h\) has conjugate
product one. On the skew scalar the action is multiplication
by \((c\bar c)^{-1}\), of additive modulus \(|c|^{-1}\):
at split places the two ordinary local absolute values multiply,
at inert places the residue field degree is two, and at a real
place \(c\bar c\) is the square of the complex modulus.
The triangular product law in Lemma~\ref{lem:one-line-radical-measure}
therefore gives the factor \(|c|^{-(r+1)}\) on radical Haar measure.
The measure
\[
                dp=|c|^{r+1}\dd n\,\dd^\times c\,\dd h
\]
is left invariant on \(P\), by substitution in these coordinates.
We take radical Haar with the compact quotient of mass one.
Right translation by a radical element preserves \(dp\).
For right translation by \(\ell(c_0,h_0)\) the integral of a
translated test acquires the factor \(|c_0|^{-(r+1)}\), since
the two Levi measures are right invariant. Right invariance
for the unitary group follows at infinity from
Lemma~\ref{lem:quotient-measure}, at finite places from
Lemma~\ref{lem:finite-unitary-level-haar}, and on finite
restricted products by normalization on the integral compact.

We justify proportionality of measures in \(G=PK\).
The natural bijection \(P/(P\cap K)\to G/K\) is a homeomorphism:
a closed saturated set in \(P\) is closed in \(G\) since \(P\)
is closed, and its product with the compact \(K\) is closed.
Indeed a convergent sequence of products has a subsequence
whose compact factors converge, and the remaining factors
converge in the closed set. These second countable groups
are metrizable in their local coordinates. Push the measure
on \(G\) to \(G/K\). The result is locally finite because the
quotient by a compact subgroup has compact preimages of compact
sets: translates of the image of a relatively compact open
identity neighborhood give a finite cover of any such set,
and their closed lifts multiplied by \(K\) are compact.
Lift this measure from \(P/(P\cap K)\) to \(P\) by integrating
over right \(P\cap K\) against its invariant probability.
This defines a nonzero left invariant locally finite measure
\(\Lambda\). For completeness, for nonnegative continuous
compactly supported \(f,b\) integrate \(f(x)b(x^{-1}y)\)
against \(\dd\Lambda(x)\,dp(y)\). Integrating in \(y\) first
gives \(\Lambda(f)p(b)\). In the opposite order write \(x=yz\)
by left invariance of \(\Lambda\). The result is
\(p(f)\int b(z^{-1})R(z)\dd\Lambda(z)\), where \(R(z)>0\)
is the right translation factor of \(dp\) computed above.
Choose \(b\) positive near one with \(p(b)>0\).
Thus \(\Lambda\) is a fixed positive multiple of \(dp\).
Equality on positive compact tests gives equality on open sets
by exhaustion with compactly supported continuous cutoffs,
then on Borel sets by countable monotone approximation of
the measures. Integrating over right \(K\), which preserves
\(G\)-measure, now gives the formula on \(G\) with this
constant \(C_P\).

For its quotient form first take
\(A(y)=\sum_{\delta\in\Gamma_P} A_0(\delta y)\), \(A_0\ge0\).
Use the rational semidirect product just obtained to sum first
over the rational radical and then over the rational Levi.
Rational Levi conjugation preserves the rational radical and
its Haar measure: the adelic modulus of a rational scalar
is one, by the rational determinant calculation in the proof of
Theorem~\ref{thm:frequency-unipotent-lattice}, applied to
multiplication on \(E\) in a rational basis. The rational \(h\) has determinant of norm
one as above. The inner mean is consequently unchanged on
replacing \(c,h\) by rational translates. Unfolding the
three sums against their fundamental domains gives the
formula on \(G\) just proved. Conversely every nonnegative
quotient function is a periodization: restrict its pullback
to a measurable fundamental domain and put it zero outside.
For these countable rational groups such a domain is obtained
by the disjoint chart construction of Lemma~\ref{lem:quotient-measure}.
Indeed a compact set in any of the adelic matrix groups contains
only finitely many rational matrices, by real bounds and fixed
finite denominators, so a sufficiently small neighborhood has
disjoint translates. Taking a countable chart cover and assigning
each orbit to its first chart gives a measurable selector.
For complex functions apply the nonnegative identity to the
positive and negative parts of their real and imaginary parts.
\end{proof}

\begin{proposition}\label{prop:large-orbit-unfold-on-one-line}
Use the standard upper large test \(f=v'\otimes F_s^{\rm up}\)
of Proposition~\ref{prop:finite-large-standard-scalar-sections},
its upper normalized character pair of
Lemma~\ref{lem:scalar-sections-poisson-and-levi}, and a nonempty
marked set. For \(g\in U(W)(\mathbb A_{F_0})\) write \(g'=m(1,g)\).
On \(C_E\times[U(V)]\times K_\infty\) define the function
\[
 B(c,h,k)=|c|^{r+1}\sum_{\xi\in E^\times}
       \theta_{W,V}\big(\mathcal L_+(\xi)
                 \Omega(g',\ell(c,h)k) f\big).
\]
If \(\int |B(c,h,k)|\dd^\times c\,\dd h\,\dd k<\infty\), then
\[
 {\mathcal C}_e P^+_{W',V'}(v',s\,{\rm pr}_W)(g')
       = C_P\int B(c,h,k)\dd^\times c\,\dd h\,\dd k .
\]
Here \(P^+\) denotes the normalized large period, obtained from
Definition~\ref{def:complementary-theta-periods} by the common
ambient determinant character in
Definition~\ref{def:normalized-complementary-levi}.
\end{proposition}
\begin{proof}
Each rational isotropic line in \(V'\) is the image of \(Ea\)
by a member of \(\Gamma\). Indeed any map sending the reference
generator to a nonzero vector on that line is an isometry on
the line, so Lemma~\ref{lem:rational-hermitian-witt-extension}
extends it to the whole space. For \(y\in G\) set
\(A(y)=\sum_{\xi\in E^\times}d_{\xi a}(g',y;f)\), using
Definition~\ref{def:one-line-radical-mean}.
These are absolutely convergent sums for fixed \(y\), by
Proposition~\ref{prop:one-line-constant-term-frame-expansion}.
For rational \(\delta\in\Gamma\), rational substitution on the
frame gives \(d_{\delta x}(g',\delta y;f)=d_x(g',y;f)\) by
Lemma~\ref{lem:gram-shear-character-on-slices} followed by
the source radical mean. Thus \(A\) is left \(\Gamma_P\)-invariant
and the sum of \(A(\delta y)\) over \(\Gamma_P\backslash\Gamma\)
is the sum over all nonzero rational isotropic frames. The zero
frame contributes zero by Lemma~\ref{lem:one-line-large-zero-frame-vanishes}.

By Proposition~\ref{prop:isotropic-vector-gives-scalar-section},
\(A\) is left invariant by the adelic target radical and its
value at \(\ell(c,h)k\) is the indicated plus sum. The finite
part of \(k\) fixes \(f\): in the proof of
Lemma~\ref{lem:standard-determinant-fixed-target} left
multiplication of the operator rows contributes the residue
determinant character, whose inverse is exactly the finite
target twist. This proof uses no condition on the fixed upper
rows. The additional power of the unramified character in the
enlarged rank is trivial on integral units. The unmarked vacuum
is fixed by Theorem~\ref{thm:all-finite-weil-splitting}.
Consequently \(|c|^{r+1}A(n\ell(c,h)k)=B(c,h,k_\infty)\).
Lemma~\ref{lem:target-line-parabolic-quotient-measure} proves that
\(A\) is absolutely integrable under the stated condition.
The sum over lines therefore unfolds in the target quotient,
by tiling a \(\Gamma_P\)-domain with translates of a
\(\Gamma\)-domain as in Lemma~\ref{lem:countable-frequency-block-unfold}.
Absolute integrability of \(A\) justifies this last exchange of
sum and integral. Proposition~\ref{prop:one-line-constant-term-frame-expansion}
identifies the original grouped integral with the constant term.
The quotient formula yields the asserted equality.
\end{proof}


\section{Removing a plane from the large period}
\begin{lemma}\label{lem:two-real-mellin-ranges}
Use Proposition~\ref{prop:large-orbit-unfold-on-one-line}
with even \(N\), \(0\le r\le N\), and put \(d=r-N/2\).
For real \(t\) write \(B_t(c,h,k)=|c|^t B(c,h,k)\).
For \(\sigma>1\) set \(t_+(\sigma)=\sigma-d-1\),
\(t_-(\sigma)=-\sigma-d\). For either sign \(\beta\),
\(B_{t_\beta(\sigma)}\) is absolutely integrable on the product
quotient in that proposition. Its integral equals
\[
 \begin{split}
 &b_\beta\,\frac{G_{\nu,\beta}}{G_S}\,
      \lambda_f(\alpha)|\alpha|_f^\sigma M_S(\sigma)\\[-2pt]
 &\hspace{28mm}\cdot
       P^+_{W,V}\big(V_\beta(\sigma;v'),s\big)(g).
 \end{split}
\]
Here the real vector and all finite factors use
Theorem~\ref{thm:real-line-cone-schwartz} and
Proposition~\ref{prop:finite-large-standard-scalar-sections},
and the normalized lower period uses \((\gamma,\gamma)\) of
Lemma~\ref{lem:scalar-sections-poisson-and-levi}. For a
zero-dimensional target its group has just one element and its
Hom space is zero. Its theta evaluation is evaluation on the
unique position, with the indicated group characters.
\end{lemma}
\begin{proof}
By Lemma~\ref{lem:scalar-sections-poisson-and-levi}, the sum in
\(B\) may be computed with either section. Indeed first add
the summand at \(\xi=0\), apply the rational theta identity
there, then remove the summand at zero on the other side.
For either sign the zero summand vanishes by transporting the
section past \(g',\ell(c,h)\) using that lemma and evaluating
the finite section of Proposition~\ref{prop:finite-large-standard-scalar-sections}
at zero. Here \(K_f\) fixes the finite datum by
Proposition~\ref{prop:large-orbit-unfold-on-one-line}.
Write \(j_+(c)=c,\ j_-(c)=1/\bar c\), and put
\[
 T_\beta(u,h,k)=
 \theta_{W,V}\big(\Omega(g,h)
     \mathcal L_\beta(u)\Omega(1,k)f\big).
\]
The same Levi formulas give the two identities
\[
 B_{t_\beta(\sigma)}
     =\sum_{\xi\in E^\times}
       \lambda(c)|j_\beta(c)|^\sigma
                       T_\beta(\xi j_\beta(c),h,k).
\]
For plus the total exponent of \(|c|\) is \(t+d+1\), and for
minus it is \(t+d\). Taking the indicated values of \(t\)
gives respectively \(\sigma,-\sigma\).

Both \(j_\beta\) preserve the multiplicative measure on
\(\mathbb A_E^\times\) and carry \(E^\times\) onto itself.
Conjugation permutes equally normalized finite factors and
preserves the real radial and angular measures. Inversion
negates the finite valuation vectors and preserves unit
probabilities, and negates the real log radii and angles.
It follows from the defining quotient integral that \(j_\beta\)
preserves the quotient measure. Moreover
\(\lambda(j_\beta(c))=\lambda(c)\), since
\(j_-(c)/\overline{j_-(c)}=c/\bar c\).
The character \(\lambda(u)=\nu(u/\bar u)\) and the adelic
modulus are trivial on rational scalars, by
Definition~\ref{def:complementary-hecke-values} and the rational
determinant computation of Theorem~\ref{thm:frequency-unipotent-lattice}.
Consequently the fundamental domain for \(C_E\) unfolds the
sum of the absolute values to the integral over full ideles
of \(|u|^\sigma|T_\beta(u,h,k)|\) for fixed \(h,k\).
The domain calculation holds for nonnegative functions without
any finiteness assumption, by Lemma~\ref{lem:integral-tools}.

To bound this last integral, put
\(w(u_\infty,k)=\mathcal L_{\beta,\infty}(u_\infty)
                         \Omega_\infty(1,k)v'\).
The finite section formula separates \(T_\beta\) as
\(b_\beta G_{\nu,\beta}D_S(u_f/\alpha)\) times the residual
theta kernel with real test \(w\) and finite test \(F_s^{\rm low}\).
Each residue matrix contributing to \(F_s^{\rm low}\) has
independent operator rows, since the stacked determinant is
a unit. For fixed \(g\) its theta kernel is therefore bounded
by \(C_gp_m(w) h_V(h)^{-L}\), with \(L\) chosen so that the
last factor is integrable, by
Theorem~\ref{thm:complementary-period-convergence} and
Proposition~\ref{prop:unitary-reduced-height-integrability}.
More explicitly one applies the single-point estimate there
to each of the finitely many named matrices in this standard
test and adds their bounds. The estimate in the integration
variable requires only that the named map onto \(V\) be
surjective. Its proof uses the height of \(g\) only as a
multiplicative constant for fixed \(g\), so it also applies
when the entire source is presented as one space. For \(r=0\)
take \(h_V=1\) and use the bound on evaluation on the unique
position, with the one-point quotient measure.
The proof of Theorem~\ref{thm:real-line-cone-schwartz} gives
\(\int_{K_\infty}\int |u_\infty|^\sigma
          p_m(w(u_\infty,k))\dd^\times u_\infty\dd k<\infty\).
The finite factor has absolute integral
\[
 \int |u_f|_f^\sigma |D_S(u_f/\alpha)|\dd^\times u_f
   =|\alpha|_f^\sigma Q_S^\sigma
       \sum_{\substack{\mathfrak a\ {\rm integral}\\
                       (\mathfrak a,S\cup\bar S)=1}}
                               (N\mathfrak a)^{-\sigma}<\infty
\]
by the valuation-coset computation of
Theorem~\ref{thm:finite-mellin-one-plane-ratio}. The sum is
over nonzero ideals coprime to the primes in the indicated
set. These three bounds prove absolute integrability of the
unfolded terms, hence also of their grouped sum on the quotient.

We may now unfold the complex integrand by the same translations.
Lemma~\ref{lem:product-integration} separates the real and
finite variables. Substitution of \(u_f=\alpha x\) gives
\[
 b_\beta G_{\nu,\beta}\lambda_f(\alpha)|\alpha|_f^\sigma
                          M_S(\sigma)/G_S
\]
for the finite integral.
The residual theta functional at fixed \(g,h\) is continuous
in the real test by the seminorm estimate just proved. Its
integral in \(h\) is likewise continuous, by integrating that
estimate. Thus the real and compact integrals may be taken
in the real Schwartz vector first: on compact truncations
approximate the continuous Schwartz-valued integrand by finite
sums in the controlling seminorm, then use the preceding
integrable tail bound. This gives precisely \(V_\beta(\sigma;v')\)
and proves the equality.
\end{proof}

\begin{theorem}\label{thm:large-standard-period-one-plane-descent}
In Lemma~\ref{lem:two-real-mellin-ranges} assume that the prescribed
weight of \(\nu\) is nonzero. Put \(s_d=\max(d+1,-d)\), and take
\(\beta=+\) if \(d\ge0\), \(\beta=-\) if \(d<0\). There is a
fixed \(C_{\rm lg}>0\) such that each allowed twist has a scalar
\(\Gamma_\nu\ne0\) for which, for every real Schwartz \(v'\),
integral standard \(s\) and adelic \(g\),
\[
 \begin{aligned}
 {\mathcal C}_e P^+_{W',V'}(v',s\,{\rm pr}_W)(m(1,g))
     &=\Gamma_\nu P^+_{W,V}(V_\beta(s_d;v'),s)(g),\\
 |\Gamma_\nu|^2
     &= C_{\rm lg}\,
                 \frac{\Pi^{\rm int}(d+1)}{\Pi^{\rm int}(d)} .
 \end{aligned}
\]
The sign, exponent, and operation on real tests depend only on
the fixed spaces, polarizations and archimedean weight.
\end{theorem}
\begin{proof}
Choose two real numbers \(a<0<b\) of the form
\(a=t_-(\sigma_-)\), \(b=t_+(\sigma_+)\) with
\(\sigma_-,\sigma_+>1\), possible by taking both large.
For \(a\le t\le b\),
\[
          |B_t|\le |B_a|+|B_b|.
\]
Indeed use \(|c|^t\le|c|^a\) for \(|c|\le1\) and
\(|c|^t\le|c|^b\) for \(|c|\ge1\), multiplying by the modulus
of the same summed function \(B\). Both bounding functions
are integrable by Lemma~\ref{lem:two-real-mellin-ranges}.
At \(t=0\) the absolute integrability hypothesis of
Proposition~\ref{prop:large-orbit-unfold-on-one-line} follows.
The integral of \(B_t\) is continuous for \(a<t<b\) by
dominated convergence, Lemma~\ref{lem:integral-tools}.

For the stated sign, \(t_\beta(s_d)=0\). If \(s_d>1\) use
Lemma~\ref{lem:two-real-mellin-ranges} directly. If \(s_d=1\),
let \(\sigma\downarrow1\) in that lemma. The left integrals
converge by the domination above and the real datum on the
right converges in each seminorm by
Theorem~\ref{thm:real-line-cone-schwartz}. The residual period
is continuous in the controlling seminorm as in the preceding
proof. Finally \(M_S(\sigma)\to M_S(1)\ne0\) by
Theorem~\ref{thm:finite-mellin-one-plane-ratio}. This uses a
one-sided limit of the finite integrals. In both cases the scalar is
\[
 \Gamma_\nu=C_P b_\beta (G_{\nu,\beta}/G_S)\,
           \lambda_f(\alpha)|\alpha|_f^{s_d}M_S(s_d).
\]
The two Gauss factors have the same nonzero absolute value by
Proposition~\ref{prop:finite-large-standard-scalar-sections}
and Theorem~\ref{thm:finite-mellin-one-plane-ratio}. Since
\(\lambda_f\) is unitary, this last theorem gives the required
squared modulus with
\(C_{\rm lg}=C_P^2|b_\beta|^2|\alpha|_f^{2s_d}>0\).
The nonzero weight gives \(M_S(s_d)\ne0\), so division by
\(\Gamma_\nu\) at any allowed twist is legitimate.
\end{proof}


\section{Successive removal of hyperbolic planes}
\begin{lemma}\label{lem:joint-zero-frame-slice}
Take \(p\ge0\) mutually perpendicular hyperbolic planes in a
source and a fixed nonsingular small target \(T\).
In each plane choose one of the two isotropic generators as
the frame position, and use the dual generator for the momentum.
In compatible real polarizations the \(p\) successive zero
slices on \({\rm Hom}(\cdot,T)\) equal the joint zero slice
in the chosen \(p\) frame variables, in any removal order.
Every real Schwartz vector on the perpendicular occurs as such
a joint slice. These assertions include \(T=0\).
\end{lemma}
\begin{proof}
Use a polarization on the perpendicular from
Lemma~\ref{lem:rational-graded-polarization}. In the trace
symplectic form the Hom spaces on distinct orthogonal source
summands are orthogonal, by
Definition~\ref{def:hom-oscillator-pairing}. For one plane
its columns on the two isotropic generators pair perfectly
with one another and have zero internal pairings. Their direct
sum with the residual polarization consequently gives positions
\((x_1,\ldots,x_p,z)\), with \(z\) the residual positions.
By Definition~\ref{def:lifted-small-line-test}, taking a zero
slice sets just the corresponding \(x_i\) to zero. In this
model the assertion on their composition is ordinary evaluation
in independent variables. Reversing the choice on any plane
uses a partial Fourier transform on that factor, with the dual
coordinate and measure there, by Lemma~\ref{lem:rational-partial-poisson}.
Such transforms leave all other factors alone and commute
with evaluation there. For a given residual vector \(v(z)\)
take \(\prod_i\exp(-\|x_i\|^2)v(z)\) in the chosen coordinates.
It is Schwartz and evaluation gives \(v\). Transport to any
of the indicated polarizations preserves the Schwartz property
by that lemma and Lemma~\ref{lem:real-oscillator-smooth-near-identity}.
For the zero target every variable has dimension zero and all
the evaluation maps are identity.
\end{proof}

\begin{theorem}\label{thm:iterate-standard-period-plane-descent}
Let \(N\) be even, \(0\le r\le N\), \(b=N-r\) and \(d=r-N/2\).
Write \(W_0=W,\ V_0=V\) for framed spaces of ranks \(N,r\).
For \(1\le i\le p\) let \(W_i,V_i\) be obtained from
\(W_{i-1},V_{i-1}\), respectively, by adjoining one physical
hyperbolic plane as in
Proposition~\ref{prop:finite-large-standard-scalar-sections}.
Keep the fixed small target \(T_b\) of rank \(b\), the admissible
squarefree levels, and the fixed nonzero archimedean weight.
At each rank use the large normalized period and the directly
transformed normalized small period, as in
Theorems~\ref{thm:large-standard-period-one-plane-descent}
and~\ref{thm:small-standard-period-one-plane-descent}.
Let \(x_p,y_p\) be Schwartz real data at the upper rank.
Suppose for each allowed \(\nu\) there is \(a_{\nu,p}\in\C\)
with
\[
 P^+_{W_p,V_p}(x_p,s_p)=a_{\nu,p} P^-_{W_p}(y_p,s_p)
\]
as functions at every integral standard \(s_p\).
There are fixed continuous linear operations \(x_p\mapsto x_0\)
and \(y_p\mapsto y_0\) on real Schwartz data, a fixed \(D>0\)
and scalars \(a_{\nu,0}\) such that
\[
 \begin{aligned}
 P^+_{W,V}(x_0,s)&=a_{\nu,0}P^-_W(y_0,s)
                         &&\text{for every integral standard }s,\\
 |a_{\nu,0}|^2&=
 D |a_{\nu,p}|^2\,\frac{\Pi^{\rm int}(d)}{\Pi^{\rm int}(d+p)} .
 \end{aligned}
\]
The vector \(y_0\) is the joint zero slice for the chosen
source lines. The scalar formula holds also when \(a_{\nu,p}=0\).
\end{theorem}
\begin{proof}
For the step \(i\to i-1\), put \(d_i=d+i\). In fact the
large rank grows by two and half the ambient rank grows by
one at each step. Set \(\beta_i=+\) if \(d_{i-1}\ge0\) and
\(\beta_i=-\) otherwise, and \(z_i=\max(d_{i-1}+1,-d_{i-1})\).
Take \(x_{i-1}=V_{\beta_i}(z_i;x_i)\) on the indicated spaces,
\(y_{i-1}=y_i[0]\).
Both operations are linear. For the integral operation, the
seminorm majorants in the proof of Theorem~\ref{thm:real-line-cone-schwartz}
are bounded by finitely many input seminorms for fixed exponent:
partial Fourier transforms are continuous, and the explicit
near-identity substitution, phase and Fourier formulas in
Lemma~\ref{lem:real-oscillator-smooth-near-identity} bound these
seminorms uniformly on each compact coordinate patch.
A finite cover handles the compact group. For evaluation
the same assertion follows directly by restricting the
supremum defining each seminorm. The resulting constructions
therefore have the stated continuity and depend only on
the fixed weight and spaces.

Given any integral standard \(s_{i-1}\), set
\(s_i=s_{i-1}{\rm pr}_{W_{i-1}}\) on \(W_i\). It is integral
in the split lattice on \(W_i\). In the upper identity average
both sides along the chosen source line at \(m(1,g)\), for
arbitrary lower \(g\). The two theorems in the statement give
\[
 \Gamma_{\nu,i}P^+_{W_{i-1},V_{i-1}}(x_{i-1},s_{i-1})(g)
  =a_{\nu,i}c_{{\rm sm},i}
                 P^-_{W_{i-1}}(y_{i-1},s_{i-1})(g).
\]
Here \(\Gamma_{\nu,i}\ne0\), \(c_{{\rm sm},i}\ne0\), and
\[
 |\Gamma_{\nu,i}|^2=C_{{\rm lg},i}
                         \Pi^{\rm int}(d_i)/\Pi^{\rm int}(d_{i-1}).
\]
Thus one may take
\(a_{\nu,i-1}=a_{\nu,i}c_{{\rm sm},i}/\Gamma_{\nu,i}\).
Only \(\Gamma_{\nu,i}\) occurs in the denominator. Starting
from \(p\) proves the lower equality at every rank, including
on all tests needed for the next step. The product of the
squared moduli telescopes, with
\(D=\prod_{i=1}^p|c_{{\rm sm},i}|^2/C_{{\rm lg},i}>0\).
Lemma~\ref{lem:joint-zero-frame-slice} identifies the small
datum. Empty products and identity operations give the
assertion for \(p=0\).
\end{proof}

\begin{corollary}\label{cor:lower-descent-all-squarefree-partners}
In Theorem~\ref{thm:iterate-standard-period-plane-descent} take
\(T_b=(A')^-\) and let \(y\) be the real datum on \((A',W)\)
whose opposite transform is \(y_0\). Then for all matched
residue data at the indicated squarefree level and all \(g\),
\[
                     I_\Phi(g)=c_{\nu,0}I'_\Psi(g),\qquad
                          |c_{\nu,0}|=|a_{\nu,0}|.
\]
\end{corollary}
\begin{proof}
Undo the common nonzero ambient determinant character in
Definition~\ref{def:normalized-complementary-levi}.
Theorem~\ref{thm:opposite-branch-small-period} identifies the
original small period with \(c_0\) times the directly
transformed one, where \(|c_0|=1\), uniformly on the standard
tests. Its real coordinate identification is invertible, since
each step is an isometric coordinate change or the conjugation
identification twice in the proof there. The original
comparison on standard tests therefore holds with
\(c_{\nu,0}=a_{\nu,0}/c_0\).
For \(b,r>0\), Theorem~\ref{thm:determinant-section-reconstruction}
expresses every matched pair as a common linear combination
of those tests, proving the assertion.
For \(b=0\) the small matrix is unique and the large matrix
runs through the invertible \(N\)-square matrices on its
support. The relation of Definition~\ref{def:complementary-finite-frames}
with \(X_k=k^{-1}\) determines both functions by the value on
the unique small matrix. The standard pair whose upper row
matrix is empty has that value one. For \(r=0\) the large
matrix is unique, and the same relation at \(Y_k=k\)
determines the pair by its value on the large matrix.
The standard pair at \(s=I_N\) takes value one there.
This also gives the zero-dimensional case by the empty
determinant convention.
\end{proof}


\section{Changing the framed target of a theta period}
\begin{lemma}\label{lem:change-framed-target-large-period}
Use the arithmetic field and physical oscillator of
Theorem~\ref{thm:actual-adelic-rational-theta-invariance}.
Fix the framed source \(W'\). Let \(T_1,T_2\) be two nonsingular
hermitian targets of the same rank with unimodular finite lattices
and the same signatures at every infinite place. Use standard
large data at the same nonempty squarefree marked set, with
the same ordered coordinate convention on the operator rows.
Let \(P_i\) be the normalized large period over \([U(T_i)]\).
There is a fixed linear isomorphism \(A_\infty\) between the
real Schwartz spaces, from the second pair to the first, a
fixed \(m_T>0\), and for each twist a scalar
\(|u_\nu|=1\), such that for all real \(v\), standard \(s\)
and group arguments \(g\),
\[
                P_2(v,s)(g)=m_T u_\nu\,P_1(A_\infty v,s)(g).
\]
The identity holds as well for any common function of the
marked matrix entries supported on surjective residue maps,
using the corresponding point names in each framing.
\end{lemma}
\begin{proof}
At a split pair two unimodular hermitian lattices of the same
rank are integrally isometric: choose a basis on one branch,
then take its pairing-dual basis on the other. At a nonsplit
finite place the residue extension is quadratic and the
lifting argument in Lemma~\ref{lem:weil-unimodular-frame-independence}
gives the integral isometry, including at residue characteristic
two. Thus the rational hermitian forms are locally isometric
everywhere. The arithmetic field satisfies the hypotheses of
Theorem~\ref{thm:hermitian-local-global}: its cyclotomic field
contains \(\zeta_{12}\), so its real subfield contains
\(\zeta_{12}+\zeta_{12}^{-1}=\sqrt3\) and the field contains
\(i=\zeta_{12}^3\). The CM extension of its real subfield has
degree two. That theorem gives a rational isometry
\(j:T_1\to T_2\). Use it to pull back the rational coordinates,
forms and target group of the second pair to the first.

We explain its effect on the oscillator. Choose a rational
polarization for the first Hom space and its image for the
second. In transported rational bases the trace pairing and
the rational position sum agree term by term. Passage from
these coordinates to the stipulated physical models intertwines
the group action by Theorem~\ref{thm:actual-adelic-rational-theta-invariance}
and Lemma~\ref{lem:weil-unimodular-frame-independence}.
At infinity this changes the real test by a fixed invertible
Schwartz operator \(A_\infty\): it is a change of position
coordinates, partial Fourier transforms and quadratic phases
as in Theorem~\ref{thm:rational-theta-symplectic-action}.
The determinant twists are preserved under conjugation by
\(j\). These choices are independent of the finite test and
of the finite characters.

Write \(L_1,L_2\) for the two finite target lattices in these
common rational coordinates. Choose \(k_f\in U(T_1)(\mathbb A_f)\)
with \(k_fL_1=L_2\), by
Lemma~\ref{lem:weil-unimodular-frame-independence}. At each
marked pair choose it to send the first ordered marked branch
frame to the pulled-back second one. A map on this split
branch extends by its inverse adjoint on the opposite branch,
which carries the corresponding dual lattice to its prescribed
lattice. Elsewhere any such lattice isometry suffices and at
almost all places both lattices coincide. In particular \(k_f\)
is independent of \(s\) at the given level.

Left multiplication by \(k_f\) on Hom is an isometry carrying
the two scaled Hom lattices and their branch subquotients to one
another. Its untwisted finite oscillator operator \(U_f\) carries
the first unit vacuum to a unit multiple of the second: it
intertwines the translations and the fixed line is unique by
Proposition~\ref{prop:finite-polarized-vacuum}. For a residue
point use the translation sum of
Proposition~\ref{prop:finite-oscillator-residue-name}.
The isometry relabels each translation by its image, preserves
the pairing in the coefficient, and induces a bijection of the
finite summation groups. Since the marked bases were transported,
a point with a given matrix name is carried to the same matrix
name in the second frame, with precisely the vacuum phase.
Adding these formulas with arbitrary coefficients proves
\(F_{2,s}=\eta\,U_f F_{1,s}\), \(|\eta|=1\), for the standard
determinant and for the general finite datum in the statement.
Replacing \(U_f\) by the normalized pair action at \((1,k_f)\)
multiplies the scalar by another number of modulus one, since
its target character is unitary.

In the target integral this group operator right translates
the variable by \(k_f\). Right translation preserves the
quotient measure by the bi-invariance in
Lemmas~\ref{lem:finite-unitary-level-haar} and~\ref{lem:quotient-measure}.
Pullback under the rational isometry identifies the discrete
subgroups and their counting measures. The two group Haar
measures in common coordinates have a fixed positive ratio
\(m_T\): at infinity compare the invariant densities at the
identity and at finite places use the compact normalization
in Lemma~\ref{lem:finite-unitary-level-haar}. The induced
quotient measures have the same ratio upon integration over
a common fundamental domain. The possible choice of \(k_f\)
has no effect on that ratio. The integrals converge for the
stated supports by Theorem~\ref{thm:complementary-period-convergence};
for rank zero use the one-point interpretation in
Lemma~\ref{lem:two-real-mellin-ranges}. This proves the formula.
\end{proof}

\begin{proposition}\label{prop:matching-top-and-descending-targets}
Let \(W,V,U\) be framed nonsingular rational hermitian spaces
over the same field, of ranks \(N,r,b\) with \(r+b=N\).
Write \(q_M(\tau)\) for the negative index of a space \(M\)
at a real embedding \(\tau\), and assume
\(q_W(\tau)=q_V(\tau)+q_U(\tau)\) for every \(\tau\).
Give \(T_b=U^-\) the opposite frame of
Lemma~\ref{lem:opposite-source-frame-and-density}.
Let \(\mathbb H\) denote the rational plane on \(e,e'\)
with \(\langle e,e\rangle=\langle e',e'\rangle=0\) and
\(\langle e,e'\rangle=1\), with the two free finite line
lattices generated by \(e,e'\). For \(p\ge1\) put
\[
 X=\mathbb H^{p}\perp W,\qquad Z=\mathbb H^{b}\perp X,\qquad
 T=T_b\perp X,\qquad \widetilde T=\mathbb H^{p+b}\perp V .
\]
All powers here are orthogonal sums. On the common source \(Z\)
the identity of Lemma~\ref{lem:change-framed-target-large-period}
holds between \(T\) and \(\widetilde T\).
Also \(X\) has an adapted rational hyperbolic plane, and the
\(b\) opposite pairs in \(Z\) and that plane together split
the finite lattice as required in
Definition~\ref{def:unmarked-projected-lattices} when \(b>0\).
\end{proposition}
\begin{proof}
At every real place a hyperbolic plane has one sign of each
kind, as seen in the basis
\((e+e')/\sqrt2,(e-e')/\sqrt2\). Thus the negative index
of \(T\) is \(b-q_U+p+q_W=b+p+q_V\), that of \(\widetilde T\).
Their dimensions are respectively \(b+2p+N\) and \(2p+2b+r\),
which agree. The plane lattice is unimodular since its pairing
matrix interchanges its two basis vectors. Orthogonal sums of
unimodular lattices are unimodular: the condition on a vector
pairing integrally with the sum tests each component separately.
The hypotheses on the two target lattices in the preceding
lemma now hold. Since \(p\ge1\), reserve one of the explicit
planes in \(X\) for the nested splitting. Its two lines and
all the lines in \(\mathbb H^b\) lie in the prescribed lattice
as direct summands with integral duality, by its construction.
\end{proof}


\section{Exchange at the prescribed complementary ranks}
\begin{definition}\label{def:normalized-equal-rank-size}
Let \(X\) be a framed even-rank space with a rational hyperbolic
plane splitting its finite lattice and let \(z\) be any real
Schwartz vector on \({\rm Hom}(X,X)\). For a nonzero fixed
weight and a tame twist of the squarefree kind in
Definition~\ref{def:ramified-line-mellin} write
\[
 R_\nu(X,z)=\frac{|J_\nu(X,z;1)|^2}
                     {\Pi^{\rm int}(\dim X/2)} .
\]
Here \(J_\nu(X,z;1)\) is the value of the normalized
equal-rank determinant period of
Definition~\ref{def:actual-equal-rank-integral} with real
datum \(z\), at the identity. This definition uses only the
space \(X\), its own frames, measures and the twist.
The finite datum is the determinant character in the split
matrix coordinates and the vacuum off the marked pairs,
as in Theorem~\ref{thm:block-frame-finite-exchange}.
Its integral converges by Theorem~\ref{thm:complementary-period-convergence},
since the finite maps supporting that determinant are invertible.
If an auxiliary positive-rank \(T_b\) is used to write
Definition~\ref{def:actual-equal-rank-integral}, it does not
occur in the defining integral over \(U(X)\).
The denominator is positive by
Theorem~\ref{thm:finite-mellin-one-plane-ratio} and
Corollary~\ref{cor:positive-exchange-lvalue-product}.
\end{definition}

\begin{theorem}\label{thm:complementary-exchange-up-to-equal-rank}
Use \(W,V,U,X\), \(p\ge1\) and the hypotheses on ranks and
signatures of Proposition~\ref{prop:matching-top-and-descending-targets}.
Assume \(N\) even and let the integer weights \(\mu_\tau\)
be positive at all infinite places. Write \(d=r-N/2\).
For any prescribed Schwartz real test \(y\) on \((U,W)\)
and any Schwartz \(z\) on \((X,X)\) there are a Schwartz
real test \(x\) on \((W,V)\) and a constant \(D_*>0\) such
that for every allowed tame twist \(\nu\) of weight \(\mu\)
and its squarefree marked set there is \(c_\nu\in\C\) with
\[
 \begin{aligned}
 I_\Phi(x;g)&=c_\nu I'_\Psi(y;g)
       &&\text{for all matched residue tests and all adelic }g,\\
 |c_\nu|^2&=D_*\,\Pi^{\rm int}(d)R_\nu(X,z).
 \end{aligned}
\]
The periods have their original orientations and determinant
characters of Definition~\ref{def:complementary-period-characters}.
In particular \(x\) and \(D_*\) may be fixed before the finite
twist and its marked set are chosen.
\end{theorem}
\begin{proof}
Transform \(y\) to the opposite model on \((W,T_b)\) by
Theorem~\ref{thm:opposite-branch-small-period}, and denote
its real vector by \(y^\flat\). Use Lemma~\ref{lem:joint-zero-frame-slice}
on \(X=\mathbb H^p\perp W\) to choose a real \(y_X\) whose
joint zero slice in those \(p\) planes is \(y^\flat\).
Suppose first that \(b>0\). On \(Z=\mathbb H^b\perp X\)
take the real vectors \(F,F'\) from
Theorem~\ref{thm:matched-real-gram-slices} with residual
real vectors \(y_X,z\). On the \(b\) outside planes choose
for subsequent descent the isotropic lines opposite
those used for the matched frame slices in that theorem.
It gives exactly \(y_X\) as the simultaneous zero slice
of \(F'\) on these opposite lines.

Use the spaces \(T=T_b\perp X\) and \(T_b\) on this top
source. Proposition~\ref{prop:matching-top-and-descending-targets}
supplies the adapted nested splitting.
Every nonempty allowed \(S\) contains a branch where the
finite determinant has full rank on its support, as needed
in Definition~\ref{def:actual-equal-rank-integral}.
Theorem~\ref{thm:global-nested-isotropic-comparison} gives
the comparison for \(F,F'\). In the directly transformed
normalized periods the top scalar is \(K_\nu J_\nu(X,z;1)\):
the phase \(c_0\) in that theorem came from converting to
the original small period, and the common ambient character
multiplies both sides equally. It follows for every standard
upper row test that the squared modulus of the top scalar
is \(C_*^2|J_\nu(X,z;1)|^2\), with \(C_*>0\) independent
of \(\nu,S\). The opposite lines used for descent merely
specify which upper tests to specialize next. At present
the identity holds for every one of them.

If \(b=0\), then \(Z=X,\ T=X\) and the small Hom space is
zero. In the normalized direct orientation its source
character is \(\upsilon_T^0=1\) by
Definition~\ref{def:lifted-small-line-test}. The small
period with the unique standard datum is consequently
\(m_0y_X\), where \(m_0>0\) is the one-point measure
times its positive vacuum value. Take \(F=y_X z,\ F'=y_X\).
The large period is \(y_X J_\nu(X,z;1)\) at every argument
by Theorem~\ref{thm:actual-constant-equal-rank-period}.
The required top scalar is \(J_\nu(X,z;1)/m_0\), which
has the same asserted form of squared modulus. This argument
also covers \(y_X=0\).

Apply Lemma~\ref{lem:change-framed-target-large-period} to
replace \(T\) by \(\widetilde T=\mathbb H^{p+b}\perp V\)
on the large side. It changes the upper real vector by its
fixed map, and the squared modulus of the resulting top
scalar is multiplied by the fixed factor \(m_T^{-2}\).
Now \(Z\) and \(\widetilde T\) are both obtained from \(W,V\)
by adjoining \(p+b\) explicit planes. Apply
Theorem~\ref{thm:iterate-standard-period-plane-descent} to
these spaces, removing first the \(b\) outside source planes
in the chosen opposite orientations and then the \(p\) planes
of \(X\). Its \(y_0\) equals \(y^\flat\) by
Lemma~\ref{lem:joint-zero-frame-slice}. The theorem gives
a large lower test \(x\) and multiplies the squared modulus
of the scalar by
\(D\,\Pi^{\rm int}(d)/\Pi^{\rm int}(d+p+b)\).
But
\[
            d+p+b=r-N/2+p+b=N/2+p=\dim X/2.
\]
Together with the top value this is the claimed size, with
\(D_*=D m_T^{-2}C_*^2\), where \(C_*=1/m_0\) for \(b=0\).
Finally Corollary~\ref{cor:lower-descent-all-squarefree-partners}
gives the equality for the original \(y\) and every matched
residue pair, preserving the modulus. The lifts, real slices
and real cone operations in this proof involve only the
spaces, real data and the prescribed infinite weight. Thus
the same \(x,D_*\) apply to all the stated twists.
\end{proof}


\section{Independence of the normalized size from the twist}
\begin{lemma}\label{lem:definite-group-finite-mass-for-exchange}
Let \(J\) be a positive definite hermitian space at every real
place, of positive rank, with unimodular integral finite lattice.
The reduced height \(h_J\) of Definition~\ref{def:complementary-theta-periods}
is bounded, and \(0<{\rm vol}([U(J)])<\infty\).
For every squarefree determinant datum and real Schwartz vector
on \((J,J)\), its kernel with either unitary determinant twist
is absolutely integrable on \([U(J)]^2\).
\end{lemma}
\begin{proof}
For \(h\in U(J)(\mathbb A)\) and \(0\ne v\in L(h)\) in
Theorem~\ref{thm:unitary-linear-shortest-representative},
the rational scalar \(\langle v,v\rangle\) is integral at
every finite place, because \(h_f^{-1}v\) is in the integral
lattice. It is positive at every real place. Its product over
real places is a positive integer: multiplication by this
nonzero algebraic integer on an integral basis of \(O_{F_0}\)
has integer determinant. Conjugating this multiplication
matrix by the invertible embedding matrix in
Lemma~\ref{lem:number-field-lattice} diagonalizes it with the
conjugates as diagonal entries. At least one conjugate is at
least one. In the signed real frames all signs are positive,
so this conjugate is the squared Euclidean length of the
corresponding component of \(h_\infty^{-1}v\). The real norm
in the shortest length dominates a fixed positive multiple of
each of these component norms by equivalence of coordinate
norms in the fixed finite dimension. Thus \(\lambda_J(h)\)
has a fixed positive lower bound. This bounds \(h_J\).
Proposition~\ref{prop:unitary-reduced-height-integrability}
gives finite measure to bounded reduced height sublevels.
The identity has an open neighborhood disjoint from its other
rational translates by Lemma~\ref{lem:quotient-measure}, so
the quotient has positive measure. Both height variables in
Theorem~\ref{thm:complementary-period-convergence} are bounded.
All points on the support of the finite determinant are
invertible. Adding their kernel estimates gives a bounded
function on this finite-measure product, proving the last claim.
\end{proof}

\begin{proposition}\label{prop:equal-rank-size-independent-definite-anchor}
Let \(J\) be the space of rank \(2k>0\) with identity hermitian
form and standard unimodular finite lattice. Fix \(p\ge1\),
\(X=\mathbb H^p\perp J\), and positive integer weights \(\mu_\tau\).
For every real Schwartz \(z\) on \((X,X)\),
\(R_\nu(X,z)\) of Definition~\ref{def:normalized-equal-rank-size}
is independent of the allowed tame twist of this weight,
including its marked set.
\end{proposition}
\begin{proof}
Take the zero-dimensional space with quotient measure \(m_0>0\)
and unit finite vacuum on the unique position. For ambient
\(W=J\) apply Theorem~\ref{thm:complementary-exchange-up-to-equal-rank}
first to \(V=J,\ U=0\), with small real datum one and residual
datum \(z\). This gives a real datum \(y\) on \((J,J)\) and
scalars \(b_\nu\) with
\[
                  |b_\nu|^2=D_1\Pi^{\rm int}(k)R_\nu(X,z),
                       \qquad D_1>0 .
\]
Use its standard matched pair on the residue module of rank
\(2k\): the large value is \(\chi^{-1}(\det)\) on the invertible
matrices, extended by zero, and the small value is one on the
unique matrix. Write \(K_\nu(g,h)\) for the \((J,J)\) kernel
on this determinant datum and \(y\), with both characters
\(\gamma=\nu^{-1}\upsilon_T^k\).
Definition~\ref{def:two-order-theta-kernel} shows that the
unnormalized large kernel with source denominator one is
\(\gamma(\det g)K_\nu(g,h)\). On the small side the action
on the zero Hom space is identity and its target character
gives the same \(\gamma(\det g)\). Cancelling this nonzero
factor gives
\[
                        \int_{[U(J)]}K_\nu(g,h)\dd h=m_0 b_\nu .
\]

Apply the same theorem to \(V=0,\ U=J\), with this very
\(y\) as small datum, the same \(X,z\), and the same ordered
bases on the two copies of \(J\). Its large real datum is a
scalar \(x\) on the zero Hom space independent of \(\nu,S\).
This time use the standard pair with \(s=I_{2k}\): the
large value is one on the unique matrix and the small value
is exactly the determinant function used in \(K_\nu\).
The large kernel has trivial source character, while on
the small kernel both characters equal \(\gamma\). Thus
\[
 m_0x=c_\nu\int_{[U(J)]}K_\nu(h,g)\dd h,\qquad
 |c_\nu|^2=D_2\Pi^{\rm int}(-k)R_\nu(X,z),\quad D_2>0 .
\]
The same zero-dimensional Haar measure occurs in the two
formulas. Integrate each equality in \(g\), using the finite
positive volume in Lemma~\ref{lem:definite-group-finite-mass-for-exchange}.
That lemma permits the two integrations of \(K_\nu\) in
either order. Substitution gives \(x=c_\nu b_\nu\), after
cancelling \(m_0{\rm vol}([U(J)])>0\). Since
\(\Pi^{\rm int}(k)\Pi^{\rm int}(-k)=1\), this proves
\[
                R_\nu(X,z)^2=\frac{|x|^2}{D_1D_2}.
\]
Both sides are nonnegative. The positive square root is
uniquely determined when \(x\ne0\), and when \(x=0\) every
such \(R_\nu\) is zero. This proves the assertion for all \(z\).
\end{proof}

\begin{lemma}\label{lem:move-source-and-target-equal-rank}
Let \(X_1,X_2\) be even positive-rank spaces having adapted
hyperbolic planes as in Theorem~\ref{thm:actual-constant-equal-rank-period}.
Suppose they have the same signatures at all real places,
equal rank and unimodular finite lattices. There are an
invertible linear change \(A\) of their real equal-rank
Schwartz data and \(m>0\), independent of \(\nu,S\), such that
\[
                    R_\nu(X_2,z)=m^2 R_\nu(X_1,Az).
\]
\end{lemma}
\begin{proof}
Choose a rational isometry as in the proof of
Lemma~\ref{lem:change-framed-target-large-period}, and use it
simultaneously on the source and target. Conjugation identifies
their rational endomorphism spaces, preserves the trace pairing
and carries their rational position sums bijectively to one
another in transported polarizations. Pullback at infinity
and passage to these coordinates give the invertible real
Schwartz change \(A\) by the same proof. On the finite part
choose one \(k_f\) aligning the pulled-back lattice and the
marked bases with those on \(X_1\), using it on both sides.
The resulting map on Hom is conjugation by \(k_f\), taking
both paired lattices for finite translations to the new ones.
The finite point names transform with one vacuum phase by
Proposition~\ref{prop:finite-oscillator-residue-name}, as
in the target change lemma. In particular the determinant
datum is preserved in the corresponding bases. The two
equal-rank determinant factors at \((k_f,k_f)\) cancel.
In the target variable \(k_f\) can be removed by right
translation in the integration. In the source variable
the identity is carried to \(k_f\) or its inverse, according
to the direction used for aligning the lattices. Both give
the value at one by
Theorem~\ref{thm:actual-constant-equal-rank-period}.
Hence the two values have absolute values in the fixed ratio
\(m\) of target quotient measures, exactly as in
Lemma~\ref{lem:change-framed-target-large-period}.
Their denominators in Definition~\ref{def:normalized-equal-rank-size}
agree because their ranks agree. Squaring gives the result.
\end{proof}

\begin{corollary}\label{cor:complementary-size-independent-all-twists}
In Theorem~\ref{thm:complementary-exchange-up-to-equal-rank}
suppose the negative index \(q_W=q\) is constant over the real
places and \(k=N/2-q\ge1\). For every \(z\) the ratio
\(|c_\nu|^2/\Pi^{\rm int}(d)\) for the constructed scalars is
independent of the allowed tame twist. It is positive if
\(J_{\nu_0}(X,z;1)\ne0\) for any one such twist \(\nu_0\).
\end{corollary}
\begin{proof}
The space \(X=\mathbb H^p\perp W\) has the same rank and
signatures as \(\mathbb H^{p+q}\perp J\) with
\(J\) the positive definite space of rank \(2k\).
Indeed their negative indices are \(p+q\) and their total
ranks are \(2p+N\). Since \(p\ge1\), both have an adapted
plane. Lemma~\ref{lem:move-source-and-target-equal-rank} and
Proposition~\ref{prop:equal-rank-size-independent-definite-anchor}
therefore prove independence of \(R_\nu(X,z)\). Multiplication
by \(D_*>0\) in
Theorem~\ref{thm:complementary-exchange-up-to-equal-rank}
gives the claim, including positivity under the stated condition.
\end{proof}


\section{Compact stabilizers of frames with a positive direction}
\begin{lemma}\label{lem:coherent-generators-dense-in-schwartz}
In the real Darboux coordinates of
Theorem~\ref{thm:negative-block-real-splitting}, let
\(\phi(u)=2^{m/4}e^{-\pi\|u\|^2}\), \(u\in\R^m\).
The complex linear span of its symmetric Heisenberg translates
is dense in the real-place Schwartz space with all its usual
seminorms. With \(a,l\in\R^m\), \(\xi=a+i l\), these translates
have the form
\[
 \rho_a(a,l)\phi(u)=e^{-\pi\|\xi\|^2/2}\,
       \phi(u)\exp\bigl(-2\pi u^{\mathsf T}\xi
                              -\pi\xi^{\mathsf T}\xi/2\bigr).
\]
The second factor after the scalar Gaussian is entire in
\(\xi\) in every Schwartz seminorm. The same statements apply
on the product of all infinite places, using their direct
sum of real coordinates.
\end{lemma}
\begin{proof}
Insert \(\phi\) in the translation formula of
Lemma~\ref{lem:schroedinger-real-gaussian} and expand
\((u+a)^{\mathsf T}(u+a)\). This gives the formula, with
\(\|\xi\|^2=\bar\xi^{\mathsf T}\xi\).
On a bounded complex parameter set, every fixed derivative
in \(u,\xi\) of the second factor is bounded by a polynomial
in \(\|u\|\) times \(C e^{-\pi\|u\|^2/2}\), by completing
the square. Taylor expansion of the exponentials is therefore
valid in each Schwartz seminorm. More explicitly, bounds
on a larger parameter disk and the scalar circle coefficient
formula give a geometrically summable bound on each Taylor
tail after any prescribed \(u\)-derivatives and coordinate
powers, as in Proposition~\ref{prop:positive-hermite-theta-series}.

First approximate any Schwartz \(f\) by compactly supported
smooth functions. Multiply by \(\eta(u/M)\), where \(\eta\)
is a smooth cutoff of compact support equal to one near zero.
Such a cutoff follows by integrating the ball bump of
Lemma~\ref{lem:frame-slice-support-transport} against the
indicator of a larger ball and normalizing: convolution with
a probability bump equals one sufficiently far inside the ball.
In the difference from \(f\) every fixed seminorm tends to zero,
by the product rule and rapid decay where \(\|u\|\gg M\).
For compactly supported smooth \(f\), the function \(f/\phi\)
is compactly supported and smooth. Fourier inversion of
Lemma~\ref{lem:schroedinger-real-gaussian} writes \(f\) as
the integral of modulations of \(\phi\), with coefficient the
Fourier transform of \(f/\phi\) and the sign chosen for
inversion. This coefficient is rapidly decreasing by repeated
integration by parts. Each seminorm of a modulation grows at
most polynomially in its frequency by the product rule.
Tail truncation and finite Riemann sums on compact frequency
boxes now give approximation in any finite list of seminorms.
Modulations are the translates with \(a=0\), proving density.
\end{proof}

\begin{proposition}\label{prop:negative-line-compact-mean-zero}
At one infinite place let the source \(X\) have signature
\((P,Q)\) and let the target \(Y\) have negative subspace
\(Y_-\) in a signed orthonormal decomposition. Choose
\(D\subset Y_-\) of complex dimension at least \(Q+1\).
Embed \(U(D)\) in the target by the identity on its orthogonal
complement. On the oscillator of \({\rm Hom}(X,Y)\) let this
group act with its untwisted Weil action times \((\det h)^a\)
for an integer \(a<P\). Its probability mean annihilates
every Schwartz vector. This conclusion is valid on the whole
real product, with identity action at the other places, and
for a conjugate of this subgroup in the same target group.
\end{proposition}
\begin{proof}
For any physical Heisenberg parameter \(T\in{\rm Hom}(X,Y)\),
project the \(Q\) negative source columns of \(T\) onto \(D\).
Their span has dimension at most \(Q\). Choose a unit
vector \(l\in D\) perpendicular to all those projections.
For \(|\eta|=1\) let \(h_\eta\) act by \(\eta^{-1}\) on
\(\C l\) and identity on its perpendicular. The varying pieces
of \(h_\eta T\) involve only the positive source columns and
the negative target line \(\C l\). These are negative Hom
coordinates. Thus in the Darboux coordinates of
Theorem~\ref{thm:negative-block-real-splitting} the image of the
parameter is \(\xi_0+\eta\xi_1\), where \(\xi_1\) corresponds
to those pieces and the two parts are orthogonal. This follows
from conjugation of the negative physical coordinates in the
formula \(x+i l_{\rm mom}=\sqrt{2j_0}(T_+,\bar T_-)\).
In particular its Hermitian squared norm is constant on the
circle. The other real places contribute only to the fixed part.

On the negative Hom subspace the determinant of the physical
action of \(h_\eta\) is \(\eta^{-P}\), from the \(P\) copies
of the rotated line against positive source columns.
The vacuum coefficient is its inverse by
Theorem~\ref{thm:negative-block-real-splitting}. Since the
operator is unitary and this coefficient has modulus one,
the untwisted action on the vacuum is multiplication by
\(\eta^P\). Including the target twist gives \(\eta^{P-a}\).
Intertwining symmetric translations therefore yields for the
circle orbit of the coherent vector at \(T\) the expression
in Lemma~\ref{lem:coherent-generators-dense-in-schwartz} at
\(\xi_0+\eta\xi_1\), multiplied by \(\eta^{P-a}\).
Its initial scalar Gaussian is constant on the circle.
The remaining entire function of \(\eta\) has only
nonnegative powers there and \(P-a\ge1\).
Integrating its uniformly convergent Taylor series as proved
in that lemma gives zero, since the mean of every positive
power of \(\eta\) is zero. This equality holds in every
Schwartz seminorm.

Right invariance of Haar probability on \(U(D)\) allows
insertion of this circle average before taking the group
mean. Indeed for fixed \(\eta\), change \(h\) to \(h h_\eta\)
in the compact integral, then integrate in \(\eta\).
Thus the group mean kills every coherent vector although
the chosen circle may depend on \(T\). The compact mean
is continuous in Schwartz seminorms by the finite coordinate
cover argument of Theorem~\ref{thm:real-line-cone-schwartz}.
Lemma~\ref{lem:coherent-generators-dense-in-schwartz} proves
the assertion on all tests. Conjugation of the group conjugates
the averaging operator by its invertible real pair operator,
giving the last assertion.
\end{proof}

\begin{corollary}\label{cor:nonnegative-frame-no-compact-invariant}
At a real place suppose \(Y\) has \(b+Q\) negative directions
and \(x:\C^b\to Y\) has nonsingular Gram matrix with a
positive eigenvalue. For \(X\) of signature \((P,Q)\) and
an integral target weight \(a<P\), the stabilizer of \(x\)
in \(U(Y)\) contains a compact subgroup whose mean on
\({\rm Hom}(X,Y)\) with this twist annihilates every real
Schwartz test.
\end{corollary}
\begin{proof}
Choose a vector in the span of \(x\) of positive norm and
normalize it to norm one. A unitary transform takes it to
a coordinate vector in the positive space: split off the
positive line in each case and identify the complements
in signed orthonormal bases as in the archimedean part
of Theorem~\ref{thm:hermitian-local-global}. In the new
coordinates the negative projection of the column span has
dimension at most \(b-1\), since that projection kills the
chosen line. Its ordinary orthogonal complement inside
\(Y_-\) has dimension at least \(Q+1\). The full unitary
group on this complement fixes each transformed column.
Apply Proposition~\ref{prop:negative-line-compact-mean-zero}
and conjugate the compact group back.
\end{proof}


\section{The stabilizer integral at an incorrect signature}
\begin{lemma}\label{lem:frame-stabilizer-fiber-integral}
Let \(G\) be an adelic unitary group here, \(H\) the pointwise
stabilizer of a nonsingular rational frame, and
\(\Delta=H(F_0)\). Equip \(G,H\) with their bi-invariant
Haar measures and use counting measure on \(\Delta\).
There is a nonnegative Borel \(u\) on \(G\) such that
\(\int_H u(hy)\dd h=1\) for every \(y\).
For \(A\in L^1(\Delta\backslash G)\) one has
\[
 \int_{\Delta\backslash G}A(g)\dd g
      =\int_G u(y)\int_{\Delta\backslash H} A(hy)\dd h\,\dd y .
\]
The inner integral is absolutely convergent for
\(u(y)\dd y\)-almost every \(y\). If a compact subgroup
\(C\subset H\) has \(\int_C A(hcy)\dd c=0\) for all
\(h,y\), then the integral on the left is zero.
\end{lemma}
\begin{proof}
The nonsingular frame splits off its span. Thus \(H\) is the
unitary group of the perpendicular by
Theorem~\ref{thm:truncated-gram-period-formula} and it is
closed in \(G\), as fixing the frame is a closed matrix
condition. Bi-invariance of the unitary measures follows as
in Lemma~\ref{lem:target-line-parabolic-quotient-measure}.
Take increasing compact \(E_j\subset G\), \(j\ge1\), whose
interiors cover \(G\). Such sets are obtained by bounding
the real coordinates and their inverses and bounding finite
denominators of matrices and their inverses by \(j!\).
The local lattice conditions give compact products at
finite places with the integral compact prescribed almost
everywhere. Increasing \(j\) places each point in an interior.
Put \(v_j={\rm vol}_H(H\cap E_jE_j^{-1})<\infty\).
For any \(y\) the set of \(h\in H\) with \(hy\in E_j\)
has measure at most \(v_j\): if it contains \(h_0\), right
translation by \(h_0^{-1}\) puts it inside \(H\cap E_jE_j^{-1}\).
The function
\[
 b(y)=\sum_{j\ge1}2^{-j}\frac{{\bf1}_{E_j}(y)}{1+v_j}
\]
therefore satisfies \(0<b_H(y):=\int_H b(hy)\dd h\le1\).
For positivity choose \(j\) with \(y\) interior to \(E_j\);
the corresponding set of \(h\)'s contains an open neighborhood
of one. By right invariance \(b_H\) is constant on \(Hy\).
Then \(u=b/b_H\) has the required integral.

Take fundamental domains \(D_H,D_G\) for \(\Delta\) acting
on \(H,G\) as in Lemma~\ref{lem:quotient-measure}.
For a nonnegative \(\Delta\)-invariant \(A\), substitute
\(g=hy\) on the right. The result is the integral on \(G\)
of \(A(g)\int_{D_H}u(h^{-1}g)\dd h\).
Partition \(G\) into \(\delta D_G\).
The sum of the inner weights at \(\delta g\) is one: the
sets \(D_H^{-1}\delta\) partition \(H\), and inversion
and right translation preserve its measure by bi-invariance.
This gives the formula for nonnegative \(A\) by
Lemma~\ref{lem:product-integration}. Apply it to \(|A|\)
and then to the real and imaginary parts to get the assertion
for complex functions. For almost every fiber with finite
absolute integral, right translation by \(c\in C\) preserves
the quotient integral on \(\Delta\backslash H\). Integrating
over \(C\) and using the same absolute integral bound and
compact probability gives zero on that fiber if the indicated
mean vanishes. The formula now proves the last assertion.
\end{proof}

\begin{proposition}\label{prop:positive-frame-gram-period-vanishes}
Let \(b>0\), let the ambient source be \(Z=\mathbb H^b\perp X\)
with its adapted \(b\)-dimensional frame space \(L\), and let
\(Y\) be a target of rank \(b+\dim X\) in the large normalized
period. Use the standard upper rows projecting onto \(L'\),
as in Definition~\ref{def:adapted-finite-frame-slices}.
At a chosen real place let \(X\) have signature \((P,Q)\),
let \(Y\) have \(Q+b\) negative directions, and suppose the
residual target determinant character at this place has
weight \(a<P\). Here the residual character is the upper
target character divided by \(\upsilon_T^b\).
At every \(g\) preserving \(L\), the coefficient of this period
at a nonsingular rational hermitian \(B\) with a positive
direction at this real place is zero for every real Schwartz test.
\end{proposition}
\begin{proof}
On the support of the standard test the operator rows restricted
to \(L^\perp\) form an invertible residue matrix. This follows
by expanding the stacked determinant in the fixed projection
rows. At a parabolic \(g\) the reduced subspace in
Definition~\ref{def:residual-frame-rank-condition} equals the
reduction of \(L^\perp\) in its adapted lattice. Hence
Theorem~\ref{thm:absolute-mixed-frame-terms} and
Corollary~\ref{cor:absolute-stabilizer-period} apply. If no
rational frame has Gram matrix \(B\), the coefficient is zero
by Theorem~\ref{thm:truncated-gram-period-formula}.
Otherwise fix one \(x\) with that Gram matrix and put
\(G=U(Y)(\mathbb A_{F_0})\). The coefficient is the integral
of \(A(y)=t_x(g,y;f)\) on \(\Delta\backslash G\), with
\(H\) the stabilizer of \(x\) and \(\Delta=H(F_0)\).
It is absolutely convergent.

For a fixed \(y\), move the test first by \(\Omega(g,y)\)
and then restrict it to \(x\) in the frame positions.
On this residual vector, left application of \((1,h)\) for
\(h\in H\) acts by the residual pair operator, by
Lemma~\ref{lem:levi-operator-every-frame} with \(q=1\).
Indeed \(h^{-1}x=x\) and the frame scalar for \(q=1\) is
one. This gives \(A(hy)\) by the rational residual theta
sum. At the chosen real place the stabilizer of \(x\) contains
the compact subgroup \(C\) of
Corollary~\ref{cor:nonnegative-frame-no-compact-invariant};
extend it by identity elsewhere. For every \(h,y\) its
average in \(A(hcy)\) applies the zero operator to the
real residual vector before applying \((1,h)\).
The rational residual theta sum can be taken after that
compact mean. At fixed \(h,y\) it is continuous in the
Schwartz vector, since only a bounded-denominator lattice
occurs in finite support and all sufficiently high real
decay norms control its absolute sum, as in
Proposition~\ref{prop:adelic-schwartz-frequency-height}.
A fixed finite test remains fixed during this real averaging.
Lemma~\ref{lem:frame-stabilizer-fiber-integral} now gives zero.
\end{proof}

\begin{corollary}\label{cor:wrong-sign-anchor-period-zero}
Write \(J\) for a positive definite space of rank \(b=2k>0\),
\(X=\mathbb H^p\perp J\) with \(p\ge1\),
\(Z=\mathbb H^b\perp X\), and \(Y=J^-\perp X\).
Use the upper normalized pair on \((Z,Y)\) of
Definition~\ref{def:normalized-complementary-levi} and
positive weights \(\mu_\tau\).
For each represented nonsingular \(B\) which fails to be
negative definite at some real place, the large coefficient
in Proposition~\ref{prop:positive-frame-gram-period-vanishes}
is zero.
\end{corollary}
\begin{proof}
At the exceptional place a nonsingular hermitian matrix
which is not negative definite has a positive direction, by
diagonalization in Theorem~\ref{thm:hermitian-local-global}.
Here \(P=p+2k,\ Q=p\), and \(Y\) has \(p+b\) negative signs.
The residual target character is
\(\nu^{-1}\upsilon_T^{\dim X/2}\) by
Lemma~\ref{lem:levi-operator-every-frame}. Its real weight
is \(a=p+k-\mu_\tau\) since \(\upsilon_T\) has weight one
by Theorem~\ref{thm:actual-global-splitting-character}.
Thus \(P-a=k+\mu_\tau>0\), as required in the proposition.
\end{proof}


\section{Consequences of a zero equal-rank functional}
\begin{lemma}\label{lem:zero-equal-rank-kills-correct-gram}
Use the spaces \(J,X,Z,Y\) and characters of
Corollary~\ref{cor:wrong-sign-anchor-period-zero}. Fix an
allowed tame twist \(\nu\) and suppose that
\(J_\nu(X,z;1)=0\) for every real Schwartz vector \(z\).
Use the finite standard test with upper rows the projection
onto \(L'\) in Definition~\ref{def:adapted-finite-frame-slices}.
For every real Schwartz test on \((Z,Y)\), every parabolic
argument \(g=u\,m(q,g_0)\) preserving \(L\), and every
rational \(B\) for which \(H_f(B,q)\) of
Definition~\ref{def:finite-transported-gram} is unimodular
at all finite places, the large \(B\)-coefficient is zero.
\end{lemma}
\begin{proof}
A unimodular transported matrix is nonsingular. If \(B\) has
a positive direction at any real place apply
Corollary~\ref{cor:wrong-sign-anchor-period-zero}, or use
Theorem~\ref{thm:truncated-gram-period-formula} if there is
no rational frame. It remains to consider \(B\) negative
definite at every real place. In a rational basis of \(L\)
identify it with the underlying vector space of \(J^-\).
Write \(B_*\) for the pullback of the form of \(J^-\) by
that identification. The finite lattice isometries and
rational local-global argument in the proof of
Theorem~\ref{thm:all-unimodular-parabolic-indices} depend
only on the two forms: they give \(c\in{\rm GL}(L)(E)\) with
\(B=c^*B_*c\) and the normalized finite frame for \(cq\)
in \(U(J^-)(\mathbb A_f){\rm GL}_b(\widehat O_E)\).
By Lemma~\ref{lem:rational-levi-transports-gram-index}, the
coefficient in question is the \(B_*\)-coefficient at \(m(c,1)g\).
This is still parabolic, and we use \(q,g_0,u\) for its
Levi coordinates and radical in what follows.

The frame \(x_*=\binom{I_b}{0}\) is now rational in the
identified coordinates. Its stabilizer is \(H=U(X)(\mathbb A_{F_0})\)
acting in the second block of \(Y=J^-\perp X\).
By the full residue rank on \(L^\perp\) established in
Proposition~\ref{prop:positive-frame-gram-period-vanishes},
Corollary~\ref{cor:absolute-stabilizer-period} unfolds
the coefficient to \(t_{x_*}(g,y;f)\) on the rational
stabilizer quotient of \(G=U(Y)(\mathbb A_{F_0})\).
Apply Lemma~\ref{lem:frame-stabilizer-fiber-integral} to
integrate first on \([U(X)]\) at fixed \(y\).

We describe the integrand on that fiber. If the finite part
of \(y\) is outside the block coset in
Proposition~\ref{prop:actual-coset-finite-coefficient}, no point
on the fiber contributes, since left multiplication by \(H_f\)
preserves that coset. In the coset, remove the right integral
factor, which fixes the standard finite test by
Lemma~\ref{lem:standard-determinant-fixed-target}.
Write what remains as \({\rm diag}(y_b,l_f)\) at finite places.
At \(x_*\) the finite residual test on
\({\rm Hom}(X,J^-)\perp{\rm Hom}(X,X)\) then has, up to a scalar,
the form
\[
           f_b\otimes
             \Omega_{\gamma_X,\gamma_X}(g_{0,f},l_f)F_{\rm eq},
              \qquad \gamma_X=\nu^{-1}\upsilon_T^{\dim X/2}.
\]
Indeed by Lemma~\ref{lem:levi-operator-every-frame} the
frame is evaluated at \(\binom{y_b^{-1}Y(q)_f}{0}\), whose
upper block is an integral isomorphism by
Proposition~\ref{prop:actual-coset-finite-coefficient}.
Theorem~\ref{thm:block-frame-finite-exchange} factors this
finite slice with equal-rank factor \(F_{\rm eq}\).
Proposition~\ref{prop:complementary-levi-factor-quotient}
puts the equal-rank characters on this factor. The remaining
radical acts only on the first factor at the frame \(x_*\),
by Lemma~\ref{lem:radical-residual-frame-action}.
All other finite scalars and the first finite vector \(f_b\)
are independent of the integration variable \(h\in[U(X)]\).

At infinity move the original arbitrary real test by its given
arguments and restrict to the indicated rational frame. This
gives a possibly coupled Schwartz vector on the two residual
summands: evaluation at fixed frame coordinates and the real
Weil action preserve Schwartz class by
Lemma~\ref{lem:real-oscillator-smooth-near-identity}. For
\(h\in H\) the same Levi formula with frame multiplier one
and \(hx_*=x_*\) now acts only on the equal-rank summand, with
target character \(\gamma_X\). Evaluate the rational theta
sum on the first summand, putting its finite component \(f_b\)
there. This gives a real Schwartz vector \(z_y\) on the
second summand. In rational product positions the first sum
is over a fixed denominator lattice with bounded finite
coefficients. For each residual seminorm its absolute sum
is bounded by multiplying the original Schwartz seminorm
by the sum of \((1+\|n\|)^{-M}\) on this lattice, finite
when \(M\) exceeds its dimension plus one. Applying this
also to derivatives proves the Schwartz assertion and the
continuity of the partial theta operator. It commutes with
the second group action: fixed-coordinate substitutions and
phases in the second variables commute term by term, and
partial Fourier integrals do so by the same Schwartz bounds,
as in Lemma~\ref{lem:rational-partial-poisson}.

Consequently, after a scalar independent of \(h\), this fiber
integrates the equal-rank kernel for \(z_y,F_{\rm eq}\), with
source argument \(g_{0,f}\) extended by one at infinity and
target variable \(h l_f\). Translating the latter variable
on the right preserves its quotient measure. The resulting
integral is \(J_\nu(X,z_y;g_{0,f})\), absolutely convergent by
Theorem~\ref{thm:complementary-period-convergence}. It vanishes
by Theorem~\ref{thm:actual-constant-equal-rank-period} and the
hypothesis on every \(z\). This applies on all contributing
fibers with arbitrary real \(y_\infty\), so
Lemma~\ref{lem:frame-stabilizer-fiber-integral} gives zero.
\end{proof}

\begin{proposition}\label{prop:zero-equal-rank-kills-all-top-periods}
Under the hypotheses of Lemma~\ref{lem:zero-equal-rank-kills-correct-gram},
every real Schwartz test on \((Z,Y)\) gives a zero large period
on every integral standard datum and at every group argument,
for the chosen twist. The same holds with \(Y\) replaced by
\(\mathbb H^{p+b}\), the \(p+b\) planes above the zero target.
\end{proposition}
\begin{proof}
Form the difference of Proposition~\ref{prop:semilocal-matched-period-difference}
with scalar zero, any fixed real small vector, and the given
real large vector. Thus its value on a column test is just
the original large period on a matching row test. The statement
and proof of that proposition apply to arbitrary real data.
We verify the primitive-coefficient criterion for this difference.
The coefficient of the large standard test at a parabolic
point vanishes whenever its transported \(H_f(B,q)\) has
a nonintegral entry: the finite frame integrality part
of Corollary~\ref{cor:nonintegral-parabolic-gram-vanishing}
uses only the finite test. Indeed a contributing frame is
integral in the transported lattice by
Lemma~\ref{lem:finite-witt-frame-partial-fourier} and
Proposition~\ref{prop:actual-coset-finite-coefficient}'s
annihilator test for frame support. Its Gram matrix would
then be integral. This contradiction in the rational frame
sum of Theorem~\ref{thm:truncated-gram-period-formula} gives
the vanishing without a condition on the real test.

Fix a marked primitive \(B\) and an admissible argument
in Definition~\ref{def:unmarked-projected-lattices}; in the
final criterion it suffices to take arguments supported at
unmarked finite places by
Lemma~\ref{lem:primitive-coefficients-at-finite-supported-points}.
The index is nonsingular and its transported Gram matrix is
unimodular at the marked pairs for a point trivial there, by
Lemma~\ref{lem:marked-index-is-transported-gram}.
If the coefficient were nonzero, the proof of
Lemma~\ref{lem:two-supports-at-least-discriminant} would
produce a point \(p_*\) with nonzero plus mean after possibly
interchanging the two line labels. Its hypotheses on the
period are the marked rank bound and radical invariance of
its bordered coefficient, valid for the period difference
by Theorems~\ref{thm:matched-difference-has-bounded-fourier-rank},
\ref{thm:adelic-maximal-rank-radical-invariance}
and~\ref{thm:opposed-means-determine-each-other}, and the
integrality of the projected
lattice at a contributing parabolic point. The latter
follows by the finite-support argument just given, exactly
as in Lemma~\ref{lem:discriminant-order-projected-lattice}.
Theorem~\ref{thm:least-discriminant-self-dual} consequently
makes the projected lattice self-dual at every unmarked place.
As in the proof of Corollary~\ref{cor:primitive-indices-every-admissible-point},
the nonzero plus mean gives a nonzero coefficient at a shift
of \(p_*\) preserving \(L\). This shift has the same Levi
action on \(L\), hence a transported Gram matrix unimodular
at all unmarked places and unchanged at the marked pairs.
At this parabolic point Lemma~\ref{lem:zero-equal-rank-kills-correct-gram}
gives zero, a contradiction. The common determinant
normalization has no effect on vanishing since the center
has determinant one. The primitive criterion just cited
therefore makes the difference zero at every argument and
on every matched residue pair, in particular every standard pair.

The proposed replacement target has the rank and real signature
of \(J^-\perp X\), by
Proposition~\ref{prop:matching-top-and-descending-targets}
with original \(W=U=J,\ V=0\). Both have unimodular lattices.
The real change of tests in
Lemma~\ref{lem:change-framed-target-large-period} is onto
and its finite scalar is nonzero. This proves the same
vanishing for every real test on the replacement.
\end{proof}


\section{Hermite polynomials on the two scalar sections}
\begin{lemma}\label{lem:physical-complex-position-coherent-evaluation}
Work at a single real place, with the real Darboux complex
coordinate \(\xi=x+i l\) of Theorem~\ref{thm:negative-block-real-splitting}.
Denote its multiplication by \(i\) by \(I\), its Hermitian
inner product by \(\langle\ ,\ \rangle_D\), conjugate linear
in the first variable. On a signed hermitian Hom space,
let \(J_{\rm ph}\) mean multiplication of every physical
matrix entry by \(i\). In Darboux coordinates \(J_{\rm ph}\)
acts by \(i\) on positive Hom entries and by \(-i\) on
negative Hom entries.
Suppose the space is an orthogonal sum of external and
residual matrix blocks for both the physical hermitian form
and the positive Darboux metric. On the external space take
a real position subspace \(P\) which is physically complex,
with complementary momentum \(IP=P^\perp\) for the real
Darboux metric. Use a real orthonormal basis of \(P\) for
its position coordinates, identifying it with a standard
real coordinate space.
Let \(z=z_E+z_R\) be a coherent parameter in Darboux coordinates
and assume \(z_E\) uses only positive Hom entries. Write
\(B_z\) for the coherent vector of
Proposition~\ref{prop:positive-hermite-theta-series}. Evaluation
in the external position \(v\in P\), retaining the full residual
vector, gives
\[
              B_z[v]=c_P e^{-\pi\|v\|_D^2}
                      e^{\,2\sqrt\pi\,\langle v,z_E\rangle_D}
                          B_{z_R},
               \qquad c_P\ne0.
\]
Here \(v\) is measured in Darboux units and the fixed scalar
includes normalization of the position model. With a fixed
linear position rescaling, use its Darboux image for \(v\).
\end{lemma}
\begin{proof}
The coordinates satisfy \(J_{\rm ph}=i\) or \(-i\) on the
indicated entries because negative physical coordinates are
conjugated when forming \(\xi\).
Choose the unitary change of external Darboux coordinates
\(U\) carrying the chosen real orthonormal basis of \(P\)
to the standard real basis and its \(I\)-multiples to the
standard imaginary basis. These vectors have the required
complex orthonormality: the imaginary pairings on \(P\)
are zero since \(IP\) is orthogonal to \(P\) for the real
metric. Extend \(U\) by identity on the residual. The
change of real polarization carries \(B_z\) to a fixed
unit multiple of \(B_{Uz_E}\otimes B_{z_R}\) in these
coordinates for normalized Lebesgue measures. Indeed the
transformed vacuum is a scalar multiple of the standard
Gaussian: its translation coefficient is the same radial
Gaussian by Theorem~\ref{thm:negative-block-real-splitting}.
Equivalently use the Gaussian action of
Proposition~\ref{prop:real-symplectic-gaussian-action};
a complex unitary Darboux matrix fixes \(iI\) in that
formula. Intertwining the symmetric translations and
using Lemma~\ref{lem:coherent-generators-dense-in-schwartz}
then transports the coherent parameter by the same complex
linear \(U\). The relation between a translate and \(B_z\)
is \(z=-\sqrt\pi\,\xi\), times its scalar
\(e^{-\pi\|\xi\|^2/2}\), by the explicit formulas in that
lemma and Proposition~\ref{prop:positive-hermite-theta-series}.
A change of position volume by a constant only rescales
the unitary coordinate identification by a fixed nonzero factor.

Let \(C_P\) be conjugation fixing the real subspace \(P\):
it is identity on \(P\) and minus identity on \(IP\).
Physical multiplication \(J_{\rm ph}\) preserves \(P\)
and commutes with \(I\), so preserves \(IP\) and commutes
with \(C_P\) as a real operator. Since \(C_P\) is antilinear
in Darboux coordinates, \(J_{\rm ph}C_P z_E=-iC_P z_E\).
Thus \(C_Pz_E\) is orthogonal to \(z_E\), whose coordinates
are positive Hom entries. In the \(U\)-coordinates this gives
\[
     (Uz_E)^{\mathsf T}Uz_E=\langle C_Pz_E,z_E\rangle_D=0 .
\]
Now evaluate the formula
\(B_w(u)=\phi_0(u)\exp(2\sqrt\pi\,u^{\mathsf T}w-w^{\mathsf T}w/2)\)
from Proposition~\ref{prop:positive-hermite-theta-series}.
The linear term at \(Uv\) is \(2\sqrt\pi\langle v,z_E\rangle_D\),
since \(Uv\) is real. The external quadratic term vanishes
by the preceding identity. This proves the formula.
\end{proof}

\begin{proposition}\label{prop:two-sections-same-sign-polynomial}
At a real place take \(W'=Ee\oplus W\oplus Ee'\),
\(V'=Ea\oplus V\oplus Ea'\) as in
Definition~\ref{def:large-two-scalar-sections}, now viewed
over \(\C\). Use the signed orthonormal vectors
\((e+e')/\sqrt2,(e-e')/\sqrt2\) and
\((a+a')/\sqrt2,(a-a')/\sqrt2\) for the new positive and
negative directions. Keep corresponding signed bases on \(W,V\).
Call an entry of \({\rm Hom}(W',V')\) a same-sign entry
when its source and target signs agree. In the Hermite test
\(H_P=P(C)\phi_0\) of Definition~\ref{def:positive-hermite-tests}
let \(P\) involve only same-sign creation labels. For either
\(\beta=+,-\) there are \(c_\beta,a_\beta\ne0\) and \(A_\beta>0\),
independent of \(P,t\), such that
\[
  \mathcal L_{\beta,\infty}(t)H_P
        = c_\beta e^{-A_\beta|t|^2}
                   H_{P_\beta(t)} \qquad(t\in\C).
\]
In forming the residual polynomial \(P_\beta(t)\), retain
every residual label. Substitute \(a_\beta\bar t\) for
each of the two entries from new positive source to new
positive target and from new negative source to new negative
target. Substitute zero for every other same-sign entry
involving a new row or column.
\end{proposition}
\begin{proof}
Use the external blocks \(A,B_1,C_1,D,H_1,I,J,K\) in
Definition~\ref{def:large-two-scalar-sections}; the middle
matrix is residual. In the Witt bases displayed there the
real positive metric is the Frobenius metric up to its
common Darboux scale \(2j_0\), since the real orthogonal
change between the two new signed vectors and the two Witt
vectors has entries \(1/\sqrt2,\,-1/\sqrt2\). In physical
matrices the Darboux complex structure is
\[
                           I M=iJ_{V'}M J_{W'},
\]
where the \(J\)'s in these bases interchange each of the
two Witt vectors and are diagonal sign matrices on the
old spaces. This follows entrywise in the signed bases
from the two rules on positive and negative Hom entries
in Theorem~\ref{thm:negative-block-real-splitting}.
Consequently the plus external positions \(A,D,I,J\) span
a physically complex subspace, and their \(I\)-images fill
the complementary blocks \(K,H_1,C_1,B_1\), respectively.
Here \(I\) in the list of blocks is the matrix entry called
\(I\) in Definition~\ref{def:large-two-scalar-sections}.
Orthogonality for the real metric follows from disjoint
matrix entries. The minus section exchanges the \(A\) and
\(K\) scalar blocks, with the same properties. Indeed the
trace symplectic pairing between \(K=t'\) and \(A=u\)
is the real trace of \(\jmath\overline{t'}u\), up to sign
from the order of the blocks. This is the kernel defining
the scalar Fourier transform at \(t\) in that definition.
Thus the physical coordinate on the new position block
is \(t'=\pm t\). More generally an equivalent fixed scaling
of the measures and the dual coordinate changes these by
a nonzero constant and the corresponding linear rescaling.
The trace pairing is nondegenerate, so this change is
complex linear in \(t\): both kernels pair \(u\) with a
constant times \(\bar t\).

In either section the selected external position has just
one nonzero physical entry, \(A=t\) or \(K=t'\), in the
two-by-two Witt block. If this entry is \(\kappa_\beta t\)
with \(\kappa_\beta\ne0\), its two same-sign entries in
signed bases both equal \(\kappa_\beta t/2\): for the \(K\)
block the two minus signs on the new negative vectors
multiply to one. All other external same-sign entries
are zero. Its squared Darboux norm is a positive constant
times \(|t|^2\). For coherent parameters using only same-sign
entries the substitution in the linear exponent of
Lemma~\ref{lem:physical-complex-position-coherent-evaluation}
therefore has the form
\(a_\beta\bar t(z_{++}+z_{--})\), with \(a_\beta\ne0\);
the common positive scale \(\sqrt{2j_0}\) multiplies both
entries. Apply \(P(\partial_z)\) at \(z=0\) to the coherent
evaluation formula. On residual coordinates the derivatives
give the corresponding creation operators by the power
series for \(B_z\) in
Proposition~\ref{prop:positive-hermite-theta-series}.
On external coordinates they replace the label by the
indicated scalar. All differentiations commute with the
section by the Schwartz estimates of
Lemma~\ref{lem:coherent-generators-dense-in-schwartz} and
the continuity of partial Fourier transformation. This
proves the asserted formula for both sections.
\end{proof}


\section{A nonzero equal-rank functional}
\begin{definition}\label{def:determinant-chain-on-definite-source}
Let \(J\) have rank \(2k>0\), with positive identity form,
and let the integral weights satisfy \(\mu_\tau\ge k\) at all
real places. Write \(m_\tau=\mu_\tau-k\).
Set \(W_i=\mathbb H^i\perp J,\ V_i=\mathbb H^i\), for
integers \(i\ge0\), in the reference coordinates of
Proposition~\ref{prop:two-sections-same-sign-polynomial}.
Order the added planes by the order in which they are
adjoined. Among the positive source columns, use just the
\(i\) positive columns in these planes. Use all \(i\) negative
source columns. Let \(D_{i,+}\) denote the determinant of
the \(i\)-by-\(i\) matrix of creation labels on the positive
target rows and the selected positive source columns.
Similarly let \(D_{i,-}\) use the negative target rows
and negative source columns. Within each sign use the same
ordering of added planes in source and target. Define the
Schwartz vector \(v_i\) as the product over infinite places of
\[
                   H_{D_{i,+}^{\,m_\tau}
                              D_{i,-}^{\,m_\tau+2k}} .
\]
The notation is that of Definition~\ref{def:positive-hermite-tests}
in signed Darboux coordinates. For \(i=0\) both determinants
are one and \(v_0=1\) on the zero Hom space. In positive
ranks the indicated constant-coefficient polynomial in
creation operators sends a Gaussian to a polynomial times
that Gaussian, which is Schwartz.
\end{definition}

\begin{lemma}\label{lem:determinant-chain-compact-and-cone}
For \(i\ge1\) the real test \(v_i\) of
Definition~\ref{def:determinant-chain-on-definite-source}
is fixed by the real compact target group in the large
normalized pair on \(W_i,V_i\).
For both signs and all real \(\sigma>0\) the cone
operation from \(i\) to \(i-1\) in
Theorem~\ref{thm:real-line-cone-schwartz} satisfies
\[
                 V_\beta(\sigma;v_i)
                            =d_{i,\beta,\sigma} v_{i-1},
                       \qquad d_{i,\beta,\sigma}\ne0 .
\]
\end{lemma}
\begin{proof}
Work first at one real place. For the target compact block
\((k_+,k_-)\in U(i)\times U(i)\), the negative physical
Hom entries comprise positive target rows against the \(i\)
negative source columns and negative target rows against
the \(2k+i\) positive source columns. The untwisted action
on the vacuum is therefore
\((\det k_+)^{-i}(\det k_-)^{-(2k+i)}\), by
Theorem~\ref{thm:negative-block-real-splitting}; its vacuum
coefficient has modulus one. The normalized target character
\(\nu^{-1}\upsilon_T^{\dim W_i/2}\) adds the power
\(k+i-\mu_\tau=i-m_\tau\) to each exponent by
Theorem~\ref{thm:actual-global-splitting-character}.
A same-sign creation label in one source column transforms
by \(k_+\) on positive rows, or \(k_-\) on negative rows.
Indeed its Darboux coordinate undergoes just this complex
linear action, so the compact implementer sends \(B_z\)
to its vacuum character times \(B_{kz}\), by the translation
identity in Lemma~\ref{lem:physical-complex-position-coherent-evaluation}.
Differentiating at zero gives the asserted law on the labels.
Taking the determinant of labels in \(i\) ordered columns
multiplies it by the determinant of the target block:
expand each new row label linearly, and the alternating sum
retains only distinct indices. The polynomial therefore
adds \(m_\tau\) on positive rows and \(m_\tau+2k\) on
negative rows. The two total exponents are
\[
     -i+(i-m_\tau)+m_\tau=0,\qquad
     -(2k+i)+(i-m_\tau)+(m_\tau+2k)=0 .
\]
This proves compact invariance at every place.

In either scalar section the last row and last column of
each determinant have zeros outside their intersection and
the common value \(a_\beta\bar t\) at that intersection,
by Proposition~\ref{prop:two-sections-same-sign-polynomial}.
Thus each determinant goes to \(a_\beta\bar t\) times
its preceding determinant. Their powers sum to
\(2m_\tau+2k=2\mu_\tau\). After including the Gaussian
factor, the local section of \(v_i\) equals
\[
 C_{i,\tau,\beta}\,
             e^{-A_{i,\tau,\beta}|t|^2}
                  \bar t^{\,2\mu_\tau}\,(v_{i-1})_\tau ,
          \qquad C_{i,\tau,\beta}\ne0,\quad A_{i,\tau,\beta}>0 .
\]
The compact variable drops out by the first assertion.
The real character
\((t/|t|)^{2\mu_\tau}\) in the cone integral cancels the
angle of this power. For \(t=re^{i\theta}\), angular
integration gives one with its probability measure, and the
radial integral is
\[
           \int_0^\infty e^{-A_{i,\tau,\beta}r^2}
                             r^{\,2\sigma+2\mu_\tau-1}\dd r .
\]
It is strictly positive, finite near zero since
\(\sigma+\mu_\tau>0\), and finite at infinity by Gaussian
decay. The product of these radial integrals and the nonzero
fixed constants is the asserted \(d_{i,\beta,\sigma}\).
The product over the places is justified by
Theorem~\ref{thm:real-line-cone-schwartz}.
\end{proof}

\begin{theorem}\label{thm:one-twist-nonzero-equal-rank-functional}
Let \(X=\mathbb H^p\perp J\) with \(p\ge1,\ \dim J=2k>0\)
and the positive standard form on \(J\). For positive integer
weights \(\mu_\tau\ge k\) and each allowed tame twist \(\nu\),
there is a real Schwartz \(z\) on \((X,X)\) with
\(J_\nu(X,z;1)\ne0\).
\end{theorem}
\begin{proof}
Fix the twist and suppose the functional vanishes on every
real Schwartz test. Put \(b=2k\) and \(h=p+b\). By
Proposition~\ref{prop:zero-equal-rank-kills-all-top-periods},
on \((W_h,V_h)=(\mathbb H^{p+b}\perp J,\mathbb H^{p+b})\)
every large period is identically zero on every integral
standard test, for every upper real vector. Use the test
\(v_h\) of Definition~\ref{def:determinant-chain-on-definite-source}.
We show this forces the large period for \(v_0\) to vanish
on every standard test. Suppose this assertion of vanishing
holds at floor \(i\ge1\) with \(v_i\).
For any \(s\) at floor \(i-1\), apply it at
\(s\,{\rm pr}_{W_{i-1}}\) upstairs and take the source-line
constant term. Its value at every \(m(1,g)\) is zero.
Theorem~\ref{thm:large-standard-period-one-plane-descent}
gives the lower period there multiplied by \(\Gamma_{\nu,i}\ne0\),
with real vector \(V_{\beta_i}(s_{d_{i-1}};v_i)\). Here
\(d_{i-1}=\dim V_{i-1}-\dim W_{i-1}/2=i-1-k\), and the
exponent is at least one even when \(d_{i-1}=0,-1\).
Lemma~\ref{lem:determinant-chain-compact-and-cone} makes
this vector a nonzero scalar multiple of \(v_{i-1}\).
Hence the period for \(v_{i-1}\) vanishes at every \(g,s\).
Induction reaches \(i=0\).

At the bottom the large target is zero and the source is \(J\).
Evaluate at \(g=1\) with standard \(s=I_b\) on every marked
branch. The stacked determinant at the unique zero-dimensional
map is one, so the finite test is the unit vacuum, and \(v_0=1\).
The theta sum on the zero Hom space has a single term one,
and the determinant characters at the identities are one.
Its target integral is the positive one-point mass. This
contradicts the forced vanishing. The functional at the
fixed twist is therefore nonzero.
\end{proof}


\section{The nonzero complementary-period comparison}
\begin{theorem}\label{thm:actual-complementary-period-exchange}
Let \(W,V,U\) be framed nonsingular rational hermitian spaces
over the arithmetic CM field, with unimodular finite lattices,
and put \(N=\dim W,\ r=\dim V,\ b=\dim U\).
Suppose \(N=r+b\), \(r=b+2e\) with \(e\) a nonnegative integer.
Write \(q\) for the number of negative signs of \(W\), assumed
the same at each infinite place. Assume
\[
 q=q_V(\tau)+q_U(\tau),\qquad
       2q+2\le N\le2q+2\mu_\tau\quad\text{for every }\tau,
\]
where \(\mu_\tau\) are integer weights. Fix a real Schwartz
test \(y\) on \({\rm Hom}(U,W)\). There are a real Schwartz
test \(x\) on \({\rm Hom}(W,V)\) and \(C>0\) with the
following property. At every nonempty allowed tame squarefree
marked set \(S\) and character \(\nu\) of this weight as in
Definitions~\ref{def:ramified-line-mellin}
and~\ref{def:complementary-period-characters}, put
\(\gamma=\nu^{-1}\upsilon_T^{N/2}\). There is \(c_\nu\ne0\)
such that for all matched pairs \((\Phi,\Psi)\) over the
residue ring and every \(g\in U(W)(\mathbb A_{F_0})\),
\[
 \begin{aligned}
 \int_{[U(V)]} K^{W,V;\gamma,1}_{x,\Phi}(g,h)\dd h
    &=c_\nu \int_{[U(U)]}K^{U,W;\gamma,\gamma}_{y,\Psi}(h,g)\dd h,\\
 |c_\nu|^2
    &=C\left(\prod_{w\in S}q_w\right)^{e^2}
                          \prod_{j=1}^e |L^S(j,\nu)|^2 .
 \end{aligned}
\]
Here \(K^{S_1,T_1;a_0,b_0}\) takes the target determinant
character \(a_0\) and inverse source determinant character
\(b_0^{-1}\) as in Definition~\ref{def:two-order-theta-kernel}.
The matching law is Definition~\ref{def:complementary-finite-frames}.
The Euler product in Definition~\ref{def:complementary-hecke-values}
omits both \(w\) and \(\bar w\), and each chosen \(w\) occurs
once in the conductor product. When \(e=0\) the entire explicit
factor multiplying \(C\) equals one.
\end{theorem}
\begin{proof}
Since \(N=2b+2e\), \(k=N/2-q\) is an integer. The two
inequalities in the hypotheses say \(1\le k\le\mu_\tau\)
for every \(\tau\). If there is no twist in the stated family
the universal assertion holds with \(x=0,C=1\). Otherwise
choose one member \(\nu_0\) and take
\(X=\mathbb H\perp W\). It has the signatures of
\(X_1=\mathbb H^{q+1}\perp J\) with \(J\) positive definite
of rank \(2k\), and both have an adapted plane.
Theorem~\ref{thm:one-twist-nonzero-equal-rank-functional}
supplies real \(z_1\) on \((X_1,X_1)\) with \(R_{\nu_0}(X_1,z_1)>0\).
The real change in Lemma~\ref{lem:move-source-and-target-equal-rank}
is invertible, so take \(z\) on \((X,X)\) corresponding to
this \(z_1\). The positive quotient-measure ratio gives
\(R_{\nu_0}(X,z)>0\).
Proposition~\ref{prop:equal-rank-size-independent-definite-anchor}
and that same change show that \(R_\nu(X,z)=R_{\nu_0}(X,z)\)
for every member of the family.

Apply Theorem~\ref{thm:complementary-exchange-up-to-equal-rank}
with \(p=1\), the given \(y\), and this single \(z\).
Its real test \(x\) has the required equality on all matched
pairs and all arguments. Its squared scalar is
\(C\Pi^{\rm int}(r-N/2)\) with
\(C=D_*R_{\nu_0}(X,z)>0\). Since
\(r-N/2=(r-b)/2=e\ge0\), the definition in
Theorem~\ref{thm:finite-mellin-one-plane-ratio} gives exactly
the displayed product. All its factors are strictly positive
by Corollary~\ref{cor:positive-exchange-lvalue-product}.
Consequently \(c_\nu\ne0\) for every twist. All choices
involving \(\nu_0\) have been made once for the entire family.
\end{proof}

\begin{corollary}\label{cor:integrated-complementary-exchange}
In Theorem~\ref{thm:actual-complementary-period-exchange} suppose
\(W=W_1\perp W_2\) with the product framing, all four spaces
\(W_1,W_2,V,U\) of positive rank. Let \(f_i\) be rapidly
decreasing functions on \([U(W_i)]\) as in
Proposition~\ref{prop:unitary-reduced-height-integrability},
and \(\xi_i\) be unitary rationally trivial adelic determinant
characters there. In the period identity put \(g=g_1\oplus g_2\),
multiply by
\(\prod_{i=1}^2 f_i(g_i)\xi_i(\det g_i)\), and integrate
in \(g_1,g_2\) over their rational quotients. The resulting
two triple integrals converge absolutely and satisfy the
same equality with exactly \(c_\nu\) and the real tests just
constructed. In either integral the finite point expansion
and the three integrations can be taken in any order.
\end{corollary}
\begin{proof}
On the support of \(\Phi\) the named residue point is a
surjection, and on that of \(\Psi\) it is an injection, by
Definition~\ref{def:complementary-finite-frames}. For each
point Theorem~\ref{thm:complementary-period-convergence}
gives the absolute triple integral bound. The added characters
have modulus one. Both finite sums of these bounds are finite,
and Lemma~\ref{lem:product-integration} permits all stated
exchanges. Now integrate the pointwise equality of
Theorem~\ref{thm:actual-complementary-period-exchange}.
\end{proof}


\section{Factoring the complementary kernels on an orthogonal sum}
\begin{lemma}\label{lem:orthogonal-block-point-names}
Use the pair kernels \(K^{S,T;a,b}_{z,D}\) of
Definitions~\ref{def:two-order-theta-kernel}
and~\ref{def:complementary-theta-periods}. In a sum of framed
spaces use product vacuum vectors, product position measures,
and the concatenation of the two rational position bases.
Write \(\delta(g)=\det_E g\). In a source sum \(W=W_1\perp W_2\)
write \(X=(X_1\ X_2)\) for a residue point mapping to \(T\).
In a target sum write \(Y=(Y_1;Y_2)\) for a point mapping
from \(S\) into \(W\). In either formula let \(z_i\) be
Schwartz vectors on the respective Hom factors, and put
\(g=g_1\oplus g_2\). With rationally trivial unitary torus
characters \(a,a_1,a_2,b,b_1,b_2\) one has
\[
\begin{aligned}
 K^{W,T;a_1a_2,b}_{z_1\otimes z_2,X}(g,h)
    &=\prod_{i=1}^2 K^{W_i,T;a_i,b}_{z_i,X_i}(g_i,h),\\
 K^{S,W;a,b_1b_2}_{z_1\otimes z_2,Y}(h,g)
    &=\prod_{i=1}^2 K^{S,W_i;a,b_i}_{z_i,Y_i}(h,g_i).
\end{aligned}
\]
The subscript of a kernel at a residue point denotes its point mass.
\end{lemma}
\begin{proof}
In either sum the physical Hom space, its trace symplectic
pairing and its finite reference lattice are orthogonal direct
sums of the two separate ones. This follows from the formula
for the Hom pairing in Definition~\ref{def:hom-oscillator-pairing}
by taking the trace in block coordinates. Both groups in the
finite sum defining a name in
Proposition~\ref{prop:finite-oscillator-residue-name} are direct
products. Their orders multiply, their dual characters multiply,
and translation by \((z,z')\) on a tensor acts by the two
translations separately because the cross pairings are zero.
Thus the name of a point is the tensor of its two point names.

On the underlying block actions the untwisted oscillator is the
tensor of the two separate actions. Indeed, in the real splitting
of Theorem~\ref{thm:negative-block-real-splitting} the positive
and negative diagonal blocks split after reordering coordinates.
The vacuum coefficients therefore multiply. In the finite
splitting of Theorem~\ref{thm:all-finite-weil-splitting} the
lattice intersections are direct sums, so their indices and
determinant ideals multiply. These equalities determine the
operator since tensoring the two operators implements the same
translation action, with the same vacuum coefficient, as in
the direct-sum calculation of Theorem~\ref{thm:residual-levi-zero-section}.
In rational position coordinates the sum at zero of a tensor
is the product of its two sums. Indeed the lattice of rational
positions is the Cartesian product and the individual series
converge absolutely by
Theorem~\ref{thm:actual-adelic-rational-theta-invariance}.
Finally \(\delta(g)=\delta(g_1)\delta(g_2)\). The two target
characters in the first product multiply to \(a_1a_2\) at \(h\);
the two source denominators in the second multiply to \(b_1b_2\)
there. This proves both identities.
\end{proof}

\begin{definition}\label{def:four-arrows-block-characters}
Use the positive-rank spaces \(W_i,V,U\) and the hypotheses of
Corollary~\ref{cor:integrated-complementary-exchange}. Denote their
ranks by \(n_i,r,b\), so \(n_1+n_2=r+b=N\). For the chosen twist
fix torus characters \(a_j,b_j\) at the ranks used below, satisfying
\(a_p b_q=\gamma\) whenever \(p+q=N\) among those ranks. All are
unitary and trivial on rational points. To a test on \((W_1,V)\)
attach numerator and denominator \((a_{n_1},a_r)\), and on
\((W_2,V)\) attach \((b_{n_2},b_r)\). On \((U,W_1)\) attach
\((b_b,b_{n_1})\), and on \((U,W_2)\) attach \((a_b,a_{n_2})\).

Given rapidly decreasing \(f_i\) on \([W_i]\), define
\(\mathcal B_X(z)\) as the integral on
\([W_1]\times[W_2]\times[V]\) of
\[
 \frac{f_1(g_1)f_2(g_2)}
      {a_r(\delta(g_1))b_r(\delta(g_2))}
                  K^{W,V;\gamma,1}_{z,X}(g_1\oplus g_2,h).
\]
Define \(\mathcal D_Y(y)\) by the same formula on
\([W_1]\times[W_2]\times[U]\) with kernel
\(K^{U,W;\gamma,\gamma}_{y,Y}(h,g_1\oplus g_2)\).
We use \(\mathcal B_X\) for split surjections \(X\) and
\(\mathcal D_Y\) for split injections \(Y\). The quotient notation
is that of Definition~\ref{def:complementary-theta-periods}.
For any single twisted kernel \(K\) on a pair \((S,T)\) put
\[
      (L_K f)(h)=\int_{[S]} f(g)K(g,h)\dd g,\qquad
      (D_K F)(g)=\int_{[T]} F(h)K(g,h)\dd h
\]
whenever these integrals converge. These pairings are bilinear.
\end{definition}

\begin{proposition}\label{prop:bilinear-two-routes-through-middle}
In Definition~\ref{def:four-arrows-block-characters}, let
\(K_1,K_2\) be the two kernels through \(V\) with the indicated
characters, point names \(X_1,X_2\), and real data \(z_1,z_2\).
Let \(K'_1,K'_2\) be the kernels through \(U\) with point names
\(Y_1,Y_2\) and real data \(y_1,y_2\). Then
\[
\begin{aligned}
 \mathcal B_X(z_1\otimes z_2)
     &=\int_{[V]}(L_{K_1}f_1)(h)(L_{K_2}f_2)(h)\dd h,\\
 \mathcal D_Y(y_1\otimes y_2)
     &=\int_{[U]}(D_{K'_1}f_1)(h)(D_{K'_2}f_2)(h)\dd h .
\end{aligned}
\]
The two integrals converge absolutely. For any fixed
\(y=y_1\otimes y_2\), the \(x\) and \(c_\nu\) of
Theorem~\ref{thm:actual-complementary-period-exchange} satisfy
\[
             \sum_X\Phi(X)\mathcal B_X(x)
                =c_\nu\sum_Y\Psi(Y)\mathcal D_Y(y)
\]
for every matched pair, with sums over the indicated full-rank points.
\end{proposition}
\begin{proof}
For the large kernels the two numerators multiply to \(\gamma\)
at \(\delta(h)\), since \(n_1+n_2=N\); the separate denominators
supply exactly the outside character in \(\mathcal B_X\).
For the small kernels the two denominators at \(\delta(h)\)
multiply to \(\gamma\), since \(a_{n_2}b_{n_1}=\gamma\).
At \(g_1,g_2\), respectively, their numerators are
\(b_b=\gamma/a_r\) and \(a_b=\gamma/b_r\), since \(r+b=N\).
Lemma~\ref{lem:orthogonal-block-point-names} therefore identifies
the two triple integrands with products of the two individual
kernels and \(f_1f_2\).

For each fixed \(h\), either kernel in its endpoint variable
has at most polynomial growth in the reduced height by
Lemma~\ref{lem:real-small-heisenberg-differences} and
Theorem~\ref{thm:unitary-linear-shortest-representative}.
In those estimates the other group argument is fixed.
Proposition~\ref{prop:unitary-reduced-height-integrability}
then gives existence of each endpoint integral. The joint
rank assumption gives absolute integrability of each triple
product by Theorem~\ref{thm:complementary-period-convergence}.
Integrating first in the independent endpoint variables and
applying Lemma~\ref{lem:product-integration} proves the two
bilinear formulas and their absolute convergence. The last
identity is Corollary~\ref{cor:integrated-complementary-exchange}
with outer determinant characters \(a_r^{-1},b_r^{-1}\),
using linearity in the finitely many point names. Its real
test on the large group may depend on the whole tensor \(y\).
\end{proof}


\section{A simultaneous scalar on the prescribed real types}
\begin{lemma}\label{lem:two-index-translation-pairing}
Let \(R(u)F(h)=F(hu)\) be right real translation on
\(L^2([V])\) with its invariant measure. Take \(e,d_e,P_e^R\)
as in Lemma~\ref{lem:real-coefficient-compression}. For two
indices \(j,l\), let \(F_i,G_i,A_i,B_i\in L^2([V])\), \(i=j,l\),
with \(\bar F_i,G_i\) fixed by \(P_e^R\). Write
\(B(F,G)=\int_{[V]} FG\). Suppose that for each \(\epsilon>0\)
there are finite sums \(S=\sum_u s_uR(u)\), \(T=\sum_v t_vR(v)\)
satisfying
\[
    \|A_i-SF_i\|_2+\|B_i-TG_i\|_2<\epsilon
                 \quad(i=j,l).
\]
Then
\[
       B(A_j,B_j)B(F_l,G_l)=B(A_l,B_l)B(F_j,G_j).
\]
If \(B(F_i,G_i)=0\) at either index, then \(B(A_i,B_i)=0\)
at that index.
\end{lemma}
\begin{proof}
Right invariance of measure and
Lemma~\ref{lem:real-coefficient-compression} give
\[
   B(R(u)F_i,R(v)G_i)
     =\langle\bar F_i,R(u^{-1}v)G_i\rangle
     =\overline{e(u^{-1}v)}B(F_i,G_i)/d_e.
\]
For each index this gives
\[
   B(SF_i,TG_i)=m B(F_i,G_i),\qquad
       m=d_e^{-1}\sum_{u,v}s_ut_v\overline{e(u^{-1}v)} .
\]
Given the stated approximation, Cauchy--Schwarz bounds
\(\left|B(A_i,B_i)-B(SF_i,TG_i)\right|\) by
\(\epsilon\|B_i\|_2+(\|A_i\|_2+\epsilon)\epsilon\).
Multiply the equality at each index by the canonical pairing
at the other index, subtract, and let \(\epsilon\) tend to zero.
When a canonical pairing is zero, each approximating pairing
there is itself zero. Passing to the limit proves the last assertion.
\end{proof}

\begin{theorem}\label{thm:two-index-criterion-large-real-tests}
Fix the spaces in Definition~\ref{def:four-arrows-block-characters}.
For each allowed twist choose the endpoint characters there. At
each twist take any specified families of rapidly decreasing
endpoint functions \(f_1,f_2\). An index \(j\) consists of the
twist, a split surjection \(X\), and a choice of these functions.
For fixed Hermite polynomials \(P_1^0,P_2^0\) on the two large
Hom factors let \(F_j,G_j\) be the two lifts through \(V\) in
Proposition~\ref{prop:bilinear-two-routes-through-middle}
with the respective real vectors \(H_{P_1^0},H_{P_2^0}\).
For general polynomials \(P,Q\) on the separate factors write
these lifts as \(F_j(P),G_j(Q)\), respectively.

Assume all these lifts lie in \(L^2([V])\), and that a single
kernel \(e\) as in Lemma~\ref{lem:two-index-translation-pairing}
fixes \(\bar F_j,G_j\) by its projection, for every \(j\).
Assume that for each \(P,Q\) and every two indices the finite
translation approximations of that lemma hold with
\(A_i=F_i(P), B_i=G_i(Q)\). Then for every Schwartz \(x\)
on \({\rm Hom}(W,V)\) there exists \(\kappa(x)\in\C\)
independent of all indices such that
\[
                \mathcal B_{X}(x)=\kappa(x) B(F_j,G_j).
\]
Here \(\mathcal B_X\) uses the functions and characters at \(j\).
If one \(B(F_j,G_j)\ne0\), the scalar is unique and is a
continuous linear functional of \(x\) in the Schwartz topology.
If all are zero one may take \(\kappa(x)=0\) for every \(x\).
\end{theorem}
\begin{proof}
Abbreviate the left side by \(\mathcal B_j(x)\), the canonical
pairing by \(b_j\). At \(x=H_P\otimes H_Q\),
Proposition~\ref{prop:bilinear-two-routes-through-middle}
identifies \(\mathcal B_j(x)\) with \(B(F_j(P),G_j(Q))\).
Lemma~\ref{lem:two-index-translation-pairing} gives
\(\mathcal B_j(x)b_l=\mathcal B_l(x)b_j\) for any two indices,
and gives \(\mathcal B_j(x)=0\) whenever \(b_j=0\)
(take the two indices equal if necessary). Both identities
extend by linearity to the span of such tensors.

This span is dense in the Schwartz space. Indeed a coherent
vector \(B_z\) on the full space factors over the two Hom
summands by the explicit exponential in
Proposition~\ref{prop:positive-hermite-theta-series}. Its
Taylor polynomials along the parameter line \(t\mapsto tz\)
are finite sums of the indicated Hermite tensors, by multiplication
of the two exponential series and the definition of the creation
operators in Definition~\ref{def:positive-hermite-tests}.
They converge at \(t=1\) in every Schwartz seminorm by the
disk estimate in Lemma~\ref{lem:coherent-generators-dense-in-schwartz}.
That lemma also gives density of the span of coherent vectors,
since they are nonzero scalar multiples of translated vacua.

For each individual \(j\), the absolute integral estimate
of Theorem~\ref{thm:complementary-period-convergence} bounds
\(|\mathcal B_j(x)|\) by a constant times a fixed Schwartz
seminorm of \(x\). The constant and seminorm may depend on
\(j\), which is sufficient for continuity of each of the
preceding two-index identities. Hence they hold on the entire
Schwartz space. If \(b_l\ne0\), take
\(\kappa(x)=\mathcal B_l(x)/b_l\). Its continuity, linearity
and uniqueness follow from this formula. Otherwise every
individual functional vanishes by the second identity.
\end{proof}

\begin{corollary}\label{cor:two-index-real-transfer-matched-sums}
Under the hypotheses of
Theorem~\ref{thm:two-index-criterion-large-real-tests}, fix the
two small polynomials \(P'_1,P'_2\) and let
\(y=H_{P'_1}\otimes H_{P'_2}\). For all twists, endpoints and
matched point functions simultaneously,
\[
    \sum_Y\Psi(Y)\int_{[U]}(D_{K'_1}f_1)(D_{K'_2}f_2)
      =\frac{\kappa}{c_\nu}
          \sum_X\Phi(X)\int_{[V]}(L_{K_1}f_1)(L_{K_2}f_2).
\]
The kernels here use the prescribed four polynomials
\(P'_1,P'_2,P_1^0,P_2^0\) and the corresponding blocks of the
summation variable. The integrals use the measures of
Definition~\ref{def:four-arrows-block-characters}.
The complex number \(\kappa\) is independent of the twist,
the finite point functions and the two endpoint forms. The
nonzero \(c_\nu\) has the squared modulus in
Theorem~\ref{thm:actual-complementary-period-exchange}.
\end{corollary}
\begin{proof}
Use the common \(x\) supplied by
Theorem~\ref{thm:actual-complementary-period-exchange}
for this fixed \(y\). In the last identity of
Proposition~\ref{prop:bilinear-two-routes-through-middle},
replace each \(\mathcal B_X(x)\) by
\(\kappa(x)B(F_j,G_j)\) using
Theorem~\ref{thm:two-index-criterion-large-real-tests}.
Replace each \(\mathcal D_Y(y)\) by the small bilinear pairing
in that proposition. Division by \(c_\nu\ne0\) gives the formula.
\end{proof}


\section{The real tests throughout the triangular array}
\begin{definition}\label{def:triangular-real-parameter-lists}
Use \(r,k,m,t=5r,h_2=r^8\), the rounded weight \(W=W_\tau\)
and \(w_l=W+2lh_2+h_2\) at an infinite place \(\tau\) from
Definitions~\ref{def:heights}, \ref{def:row-vandermonde}
and~\ref{def:first-arrow-insertion-parameters}. In particular
\(r=k+m-1\) is odd, \(k\) is even, \(W\ge0\).
Fix \(z=0\) or \(1\), put \(D=t-z\), and for \(0\le j\le m\) set
\[
         n_j=r+z+jD,\qquad d_j=\{z-(1-z)(j\bmod 2)\}/2 .
\]
For offset \(o\ge0\) and \(j\ge0\) with \(o+j\le m\) form a
list \(T_{o,j}\) of \(n_j\) parameters with signs. For the
\(r+z\) base indices \(0\le i<r+z\) use
\(c_i=\max(0,i+1-k)\) and the following parameter and sign:
\[
 \lambda_i=d_j+
 \begin{cases}
 k/2-1-i,&i<k-1,\\
 -W-2(c_i+1)h_2,&i\ge k-1,
 \end{cases}
 \qquad
 s_i=(-1)^{o+c_i+\min(\max(c_i-o,0),j)} .
\]
For \(1\le l\le j,\ 0\le v<D\) include the parameter
\[
                  (t-1)/2-w_{o+l}-z-v+d_j
\]
with positive sign. Identify these entries by the labels
\((o+l,v)\), disjoint from the base labels.
Order each list by decreasing parameter. For \(-T_{o,j}\)
negate every parameter, retain each sign, and order again.

For \(o+j<m\) consider the following two Hom pairs. Put
\(n=n_j,\ H=n_{j+1}\). In the first row \(w=w_{o+j+1}\);
in the second \(w=w_{o+1}\). The last columns are the source and
target determinant exponents in the real action.
\[
\begin{array}{c|c|c|c}
 \text{source}&\text{target}&\alpha&\beta\\ \hline
 T_{o,j}&T_{o,j+1}&w-\lfloor H/2\rfloor&\lfloor n/2\rfloor-w\\
 -T_{o+1,j}&-T_{o,j+1}&-\lceil H/2\rceil-w&\lceil n/2\rceil+w
\end{array}
\]
For the first row the torus character at rank \(p\) has weight
\(\lfloor p/2\rfloor-w\); for \(z=1\) both ranks are even.
In the second row take the quotient of the \(p\)-th power of
\(\upsilon_T\) by the first character of weight
\(\lfloor p/2\rfloor-w_{o+1}\). Its weight is
\(\lceil p/2\rceil+w_{o+1}\) by
Theorem~\ref{thm:actual-global-splitting-character}. In each
case the numerator uses rank \(n\) and the denominator rank
\(H\), as in Definition~\ref{def:two-order-theta-kernel}.
\end{definition}

\begin{proposition}\label{prop:triangular-two-insertion-orders}
For each row in Definition~\ref{def:triangular-real-parameter-lists},
at sufficiently large scale, put \(u_i=\lambda_i+\alpha+H/2\)
for a source coordinate and \(v_a=\lambda'_a-\beta+n/2\) for a
target coordinate. Indices \(i,a\) in these formulas denote
labels in their respective ordered lists. There is exactly
one target entry with the old label of each source entry. They
satisfy
\[
            v_{\rm old(i)}=u_i,\qquad
                   s_i t_{\rm old(i)}u_i>D/2 ,
\]
where \(t_a\) denotes the target sign. The remaining \(D\)
entries have positive sign and, in decreasing order, centered
values \((D-1)/2-v,\ 0\le v<D\). In particular all lie strictly
between \(-D/2\) and \(D/2\). In both lists distinct parameters
are at least one apart; parameters with opposite signs are
separated by more than \(500(n+H)^3\), uniformly in the row,
offset, level and place.
\end{proposition}
\begin{proof}
In the first row the common shift on the old labels is
\(d_{j+1}-d_j=\alpha+\beta+D/2\); in the second it is its
negative, again \(\alpha+\beta+D/2\) for the exponents there.
Indeed \(H-n=D\), and \(n_j\bmod2\) is \(1-(j\bmod2)\)
if \(z=0\), zero if \(z=1\). Substituting these parities gives
both identities. The new block in the first row is labeled
\(o+j+1\). In the second the common added labels run from
\(o+2\) through \(o+j+1\), and the new block has label \(o+1\).
Its internal order is reversed by negation. Thus the new
parameters, at internal position \(v\) starting at zero, are
\[
 \begin{cases}
  (t-1)/2-w-z+d_{j+1}-v,&\text{first row},\\
  (t-1)/2+w-d_{j+1}-v,&\text{second row}.
 \end{cases}
\]
In both cases the constant term equals
\(\beta-(H+1)/2+D\) by the same parities. This proves the
two assertions on centered values.

In the first row increasing \(j\) changes the parity of the
base sign exactly when \(c_i>o+j\). For \(i\ge k-1\) one has
\(\lambda_i+w=2h_2(o+j-c_i)+h_2+d_j\). For \(i<k-1\) this
sum is at least \(h_2-r-1\), since \(W\ge0\) and the central
base parameters before the shift have magnitude at most \(r\).
For every old added block it is at least \(2h_2-t-2\), by
the spacing \(w_{s+1}-w_s=2h_2\). Since
\(u_i=\lambda_i+w+(H\bmod2)/2\), the sign of \(u_i\) is
positive on the unchanged signs and negative on the flipped signs.

In the second row the two sign exponents on a base label are
\(o+1+c_i+\min(\max(c_i-o-1,0),j)\) and
\(o+c_i+\min(\max(c_i-o,0),j+1)\). They have the same parity
if \(c_i\ge o+1\), and opposite parity if \(c_i\le o\).
Here \(u_i=\lambda_i-w-(H\bmod2)/2\), with \(\lambda_i\)
the negated source parameter. For \(i\ge k-1\), before
the last rank correction this equals
\(2h_2(c_i-o)-h_2-d_j\). On the central base it is negative
with magnitude at least \(h_2-r-1\). Every old added block
has label at least \(o+2\), so its negated parameter exceeds
\(w_{o+1}\) by at least \(2h_2-t-2\). Thus the same description
by unchanged and flipped signs holds. In either row we have
\(|u_i|\ge h_2-2t-r-4>D/2\) at large scale.

Within each consecutive block all signs agree and successive
parameters differ by one. The base entries outside the central
block before shifting lie at \(-W-2(c_i+1)h_2\). The added
blocks lie within \(t+1\) of \(-W-(2(o+l)+1)h_2\).
The central block lies within \(r\) of zero and its sign is
constant. Hence different blocks, treating each outer base
entry as its own block, are separated by at least
\(h_2-2t-r-4\). Negation leaves distances unchanged.
The bounds \(n,H\le r+1+5r^2\) follow from \(j+1\le m\le r\).
As \(h_2=r^8\), the stated separation exceeds \(500(n+H)^3\)
for all sufficiently large \(r\). Within a list all entries
have the same fractional part, since \(k\) is even and
\(t\) is odd. The separation and the internal differences
now prove distinctness in both orders.
\end{proof}

\begin{theorem}\label{thm:joint-determinants-all-triangular-arrows}
For every row in Definition~\ref{def:triangular-real-parameter-lists}
at sufficiently large scale there is an explicit nonzero Hermite
polynomial test \(H_P\) on its Hom space which is fixed by both
compact coefficient averages of Theorem~\ref{thm:joint-insertion-harmonic}.
Its polynomial satisfies both second-derivative vanishing
conditions in that theorem, and its projection by the conjugate
source kernel is nonzero. At all infinite places the tensor
of these choices works for every finite twist with the prescribed
weights.
\end{theorem}
\begin{proof}
In the source order, the condition \(u_i>0\) picks out a prefix,
since \(u_i\) differs from \(\lambda_i\) by a fixed scalar.
Proposition~\ref{prop:triangular-two-insertion-orders} places
the new block just after this prefix in the target order.
Old signs are preserved in the prefix and reversed in the
suffix. The common shift on old parameters and the formula on
the new block proved there are exactly those of
Theorem~\ref{thm:joint-insertion-harmonic}.

Before negation the fractional part \(d_j\bmod1\) equals
\((n_j-1)/2\bmod1\) by the rank parity, and negation preserves
this congruence. Thus the compact weights in
Definition~\ref{def:polar-row-coefficient} are integral. On a
list of size \(s\), the formula there gives
\(|\Lambda_i-\lambda_i|\le 3s/2\) by bounding the half sum of
all differences of basis weights and the same-sign count
separately. For two consecutive positions \(i<i'\) of the same
sign within its sign set the weight difference is
\(\lambda_i-\lambda_{i'}+i'-i-2\), by the count of preceding
equal signs in Theorem~\ref{thm:joint-insertion-harmonic}.
Successive parameters in the whole list differ by at least one,
so this difference is nonnegative. This verifies dominance
on both lists.

The source integer \(b_i=\Lambda_i+\alpha+h_{-s_i}\) of
Theorem~\ref{thm:one-sided-harmonic-minor-test} differs from
\(u_i\) by at most \(3n/2+H/2\), since \(0\le h_{-s_i}\le H\).
The lower bound \(h_2-2t-r-4\) in
Proposition~\ref{prop:triangular-two-insertion-orders} exceeds
that error. Hence \(b_i\) is positive in the prefix and negative
in the suffix. Apply Theorem~\ref{thm:joint-insertion-harmonic}
using each old label and its target copy. Its construction is
the product of the determinants on the initial and final segments
defined by the positive and negative integer levels of \(b_i\).
It supplies both compact assertions and both differential identities.

For either list the \(\mu\) in
Definition~\ref{def:polar-row-coefficient} differs from its
parameter by the decreasing full half sum. At an adjacent change
of sign its gap is therefore at least the parameter gap.
The bound in Proposition~\ref{prop:triangular-two-insertion-orders}
supplies the moment condition of Theorem~\ref{thm:polar-row-decay}
with \(B=0\). A definite group is compact and has the same
integrability directly. The last assertion of
Theorem~\ref{thm:one-sided-harmonic-minor-test} now proves
nonvanishing of the source projection at each place. The
product projection acts on the tensor by the tensor of the
projected vectors, by absolute integration of the kernel as
in Lemma~\ref{lem:real-coefficient-compression}. Its norm is
the product of their positive norms. All lists and polynomials
depend only on the infinite data.
\end{proof}


\appendix
\section{Elementary tools for the analytic reductions}
\begin{lemma}\label{lem:basic-arith}
For positive integers \(a,b\), their greatest common divisor is an
integer linear combination of \(a,b\). In particular congruences
modulo pairwise coprime integers can be solved simultaneously, and
prime factorizations of positive integers are unique up to order.
\end{lemma}
\begin{proof}
The least positive member \(h\) of \(a\Z+b\Z\) divides each member,
since division by \(h\) leaves a remainder in the same set. Hence
\(h=(a,b)\). For coprime \(a,b\), an integer inverse of \(a\) modulo
\(b\) solves \(u+ak\equiv v\bmod b\), giving the two-modulus assertion.
Iteration proves the finite case. If a prime \(p\) divides \(ab\)
and does not divide \(a\), multiply \(1=pa'+ab'\) by \(b\) to see
that \(p\mid b\). Factor a composite by induction on its size. To see
uniqueness, match a prime on one side against a factor on the other by
the divisibility property just proved, then cancel and repeat.
\end{proof}

\begin{lemma}\label{lem:integral-tools}
Increasing nonnegative measurable functions converge in integral to
the integral of their pointwise limit. A pointwise convergent sequence
dominated in absolute value by one integrable function converges in
integral. A series can be integrated term by term if the integrals of
its absolute values are summable.
\end{lemma}
\begin{proof}
Define the nonnegative integral as the supremum over simple functions
below the integrand. If \(f_n\uparrow f\), \(0\le u\le f\) is simple
and \(0<c<1\), the sets \(\{f_n\ge cu\}\) increase to the whole
space up to points where \(u=0\), which do not affect the integral of
\(u\). Countable additivity on disjoint differences gives convergence
of their measures on every level set of \(u\). Hence
\(\lim\int f_n\ge c\int u\). Take the supremum and let \(c\uparrow1\).
The opposite inequality is monotonicity. For \(|g_n|,|g|\le h\)
with \(h\) integrable and \(g_n\to g\), apply this result to
\(2h-\sup_{m\ge n}|g_m-g|\uparrow2h\). The integral of the supremum
tends to zero. For the series, apply the first assertion to the partial
sums of absolute values, then the second to the partial sums of the
original series.
\end{proof}

\begin{lemma}\label{lem:analytic-tools}
A holomorphic function on a neighbourhood of \(|z-z_0|\le r\) has a
Taylor series about \(z_0\) throughout the open disk. Its \(k\)-th
coefficient has modulus at most \(r^{-k}\max_{|z-z_0|=r}|f(z)|\).
Holomorphic functions agreeing on a set with an interior accumulation
point in a connected domain agree throughout it. The same uniqueness
holds for meromorphic functions. At a zero of finite order \(a\) of
\(f\), the logarithmic derivative \(f'/f\) has principal part
\(a/(z-z_0)\). A locally integrably dominated integral of holomorphic
functions is holomorphic, provided the parameter dependence is
measurable and the domination holds on every compact subdisk.
\end{lemma}
\begin{proof}
For the disk formula, first integrate a holomorphic function around
a triangle. Dividing into four similar triangles cancels the interior
edges. If the original integral has modulus \(A\), a nested sequence
of triangles has integrals of modulus at least \(4^{-n}A\) at stage \(n\).
At their common point \(w\), subtract the affine function
\(f(w)+f'(w)(z-w)\), whose contour integral is zero. The remaining
integral is \(o(2^{-n})\) times a perimeter of order \(2^{-n}\).
Thus \(A=0\). Triangulating polygonal annuli and approximating their
circular boundaries by inscribed polygons applies this result to
\(f(w)/(w-z)\) away from \(z\). The integral along the inner circle
shrunk about \(z\) tends to \(2\pi i f(z)\), by continuity and explicit
integration of \(1/(w-z)\). On any fixed circle the polygonal integrals
converge by uniform continuity on a surrounding compact annulus
and ordinary parametrisation of the circle. This proves the disk
integral formula. Expanding \(1/(w-z)\) in its uniformly convergent
geometric series on that boundary gives the expansion and coefficient
bound. If the series is nonzero, its first nonzero coefficient exhibits
a factorisation \(f(z)=(z-z_0)^a h(z)\) with \(h(z_0)\ne0\). This
gives isolated zeros and the asserted principal part. For a function
vanishing on an accumulating set, the points at which its germ vanishes
form a nonempty open and relatively closed set. Connectedness gives
uniqueness. Clearing local denominators proves the meromorphic version.

For the integral assertion, on any closed disk of domination the coefficient
bound just proved dominates the \(k\)-th coefficient by the integrable
majorant times \(r^{-k}\). On a smaller disk these bounds sum to an
integrable geometric series. Lemma~\ref{lem:integral-tools} therefore
allows integration of the coefficients. The integral is given locally
by the resulting power series.
\end{proof}

\begin{thebibliography}{9}
\bibitem{compact} L. Sernau, collaborators at Anthropic, and Claude,
\emph{GRH for Dirichlet \(L\)-functions, in Lean 4}, compact edition.
The repository credits the underlying argument to L. Alp\"oge and Claude.
\href{https://github.com/lsernau/grh-lean/tree/72fead102f663d5c7a407964a13dc0cf9d977207/grh-compact}{Source snapshot \(72\mathrm{fead}1\)}.
\bibitem{dlmf} NIST Digital Library of Mathematical Functions,
\emph{Dirichlet \(L\)-functions}, \href{https://dlmf.nist.gov/25.15}{\S25.15},
equations (25.15.4) and (25.15.5).
\end{thebibliography}
\end{document}
